% !TEX program = xelatex
\documentclass[UTF8,zihao=-4,oneside,openany]{ctexbook}
\usepackage[a4paper,top=20mm,bottom=22mm,left=19mm,right=19mm,headheight=25pt]{geometry}
\usepackage[table]{xcolor}
\usepackage{amsmath,amssymb,mathtools,bm}
\usepackage{siunitx}
\usepackage{booktabs,tabularx,array,multicol,longtable}
\usepackage{enumitem}
\usepackage{graphicx}
\usepackage{float}
\usepackage{caption}
\usepackage{tikz}
\usepackage{pgfplots}
\pgfplotsset{compat=1.18}
\usetikzlibrary{arrows.meta,positioning,calc,angles,quotes,patterns,patterns.meta,
  decorations.pathmorphing,decorations.markings,intersections,shapes.geometric,
  backgrounds,3d,fit,spy}
\usepackage[most]{tcolorbox}
\usepackage{fancyhdr}
\usepackage{hyperref}
\usepackage{bookmark}
\usepackage{xparse}

\definecolor{Ink}{HTML}{202326}
\definecolor{DeepBlue}{HTML}{294E69}
\definecolor{Blue}{HTML}{47687D}
\definecolor{PaleBlue}{HTML}{F3F5F6}
\definecolor{Gold}{HTML}{6E777D}
\definecolor{PaleGold}{HTML}{F6F6F5}
\definecolor{RuleGray}{HTML}{D3D7DA}
\definecolor{TextGray}{HTML}{61686D}
\definecolor{SoftGray}{HTML}{F6F6F6}
\color{Ink}

\sisetup{per-mode=symbol,inter-unit-product=\ensuremath{{}\cdot{}},exponent-product=\times}
\hypersetup{
  colorlinks=true,linkcolor=DeepBlue,urlcolor=Blue,
  bookmarksdepth=2,
  pdftitle={高中物理竞赛一轮讲义},
  pdfsubject={数学工具、力学、热学、电磁学、光学与近代物理的一轮竞赛讲义},
  pdfauthor={物理竞赛一轮资料编写组}
}

\setlength{\parindent}{2em}
\setlength{\parskip}{0.18em}
\setlist{nosep,leftmargin=2em}
\setcounter{secnumdepth}{2}
\setcounter{tocdepth}{1}
\allowdisplaybreaks[2]
\emergencystretch=2em

\ctexset{
  chapter={format=\raggedright\LARGE\bfseries\color{Ink},beforeskip=0pt,afterskip=18pt},
  section={format=\Large\bfseries\color{DeepBlue},beforeskip=1.8ex,afterskip=1.0ex},
  subsection={format=\large\bfseries\color{Ink},beforeskip=1.5ex,afterskip=0.8ex},
  subsubsection={format=\normalsize\bfseries\color{Ink}}
}

\pagestyle{fancy}
\fancyhf{}
\fancyhead[L]{\small\color{TextGray}高中物理竞赛一轮讲义}
\fancyhead[R]{\small\color{TextGray}\nouppercase{\leftmark}}
\fancyfoot[C]{\small\color{TextGray}\thepage}
\renewcommand{\headrulewidth}{0.4pt}
\renewcommand{\headrule}{\hbox to\headwidth{\color{RuleGray}\leaders\hrule height \headrulewidth\hfill}}

\newcommand{\dd}{\mathop{}\!\mathrm d}
\newcommand{\e}{\mathrm e}
\newcommand{\ii}{\mathrm i}
\newcommand{\vect}[1]{\bm{#1}}
\newcommand{\uvec}[1]{\hat{\vect{#1}}}
\newcommand{\dimens}[1]{\operatorname{dim}\!\left(#1\right)}
\newcommand{\order}[1]{\mathcal O\!\left(#1\right)}
\DeclareMathOperator{\proj}{proj}
\DeclareMathOperator{\sinc}{sinc}

\renewcommand{\theequation}{\thechapter.\arabic{equation}}
\renewcommand{\thefigure}{\thechapter.\arabic{figure}}
\renewcommand{\thetable}{\thechapter.\arabic{table}}
\makeatletter
\@addtoreset{equation}{chapter}
\@addtoreset{figure}{chapter}
\@addtoreset{table}{chapter}
\makeatother

\newcounter{bookexample}[chapter]
\renewcommand{\thebookexample}{\thechapter.\arabic{bookexample}}
\NewDocumentEnvironment{origexamplebox}{m}{%
  \refstepcounter{bookexample}%
  \begin{tcolorbox}[
    enhanced,breakable,colback=white,colframe=RuleGray,boxrule=0.5pt,arc=0mm,
    title={例~\thebookexample\quad #1},fonttitle=\bfseries,
    coltitle=Ink,colbacktitle=PaleBlue,borderline west={1.4pt}{0pt}{DeepBlue}]
}{\end{tcolorbox}}


% FinalReviewSolutionHeading20261024
\makeatletter
\newcommand{\LectureSolutionHeading}{%
  \par\noindent{\bfseries\color{DeepBlue}解答}\par\nobreak
  \@afterindentfalse\@afterheading}
\makeatother

\NewDocumentEnvironment{xesolution}{+b}{%
  \par\smallskip\begingroup\small\noindent
  \LectureSolutionHeading #1\par\endgroup\smallskip
}{}

\NewDocumentEnvironment{thinkanswer}{+b}{%
  \par\smallskip\begingroup\small\noindent
  {\bfseries\color{DeepBlue}思考辨析}\quad #1\par\endgroup\smallskip
}{}

% Used when a long worked solution is split by a page or a figure.  The
% continuation deliberately carries no second heading, so that it remains one
% answer both typographically and semantically.
\NewDocumentEnvironment{solutioncontinuation}{+b}{%
  \par\begingroup\small #1\par\endgroup\smallskip
}{}

% Numbering and restrained boxes for the independent advanced-reading volume.
% Advanced units deliberately do not consume the 32-chapter counter: their
% public identifiers are A1.1, A2.1, ... and remain stable whether a unit is
% compiled by itself or appended to the complete book.
\newcounter{advunit}
\newcounter{advexample}[advunit]
\newcounter{advdefinition}[advunit]
\newcounter{advtheorem}[advunit]
\renewcommand{\theadvexample}{A\arabic{advunit}.\arabic{advexample}}
\renewcommand{\theadvdefinition}{A\arabic{advunit}.\arabic{advdefinition}}
\renewcommand{\theadvtheorem}{A\arabic{advunit}.\arabic{advtheorem}}

\NewDocumentCommand{\AdvancedUnit}{m m}{%
  \clearpage
  \setcounter{advunit}{#1}%
  \setcounter{advexample}{0}%
  \setcounter{advdefinition}{0}%
  \setcounter{advtheorem}{0}%
  \setcounter{equation}{0}%
  \setcounter{figure}{0}%
  \setcounter{table}{0}%
  \setcounter{section}{0}%
  \setcounter{subsection}{0}%
  \renewcommand{\thesection}{A\arabic{advunit}.\arabic{section}}%
  \renewcommand{\thesubsection}{\thesection.\arabic{subsection}}%
  \renewcommand{\theequation}{A\arabic{advunit}.\arabic{equation}}%
  \renewcommand{\thefigure}{A\arabic{advunit}.\arabic{figure}}%
  \renewcommand{\thetable}{A\arabic{advunit}.\arabic{table}}%
  % Hyperref destinations must remain unique after the appendix and after
  % each advanced unit resets its visible section/figure counters.
  \renewcommand{\theHsection}{advA#1.section.\arabic{section}}%
  \renewcommand{\theHsubsection}{advA#1.section.\arabic{section}.\arabic{subsection}}%
  \renewcommand{\theHequation}{advA#1.equation.\arabic{equation}}%
  \renewcommand{\theHfigure}{advA#1.figure.\arabic{figure}}%
  \renewcommand{\theHtable}{advA#1.table.\arabic{table}}%
  \chapter*{\texorpdfstring{A#1\quad #2}{A#1 #2}}%
  \markboth{A#1\quad #2}{A#1\quad #2}%
  \addcontentsline{toc}{chapter}{\texorpdfstring{A#1\quad #2}{A#1 #2}}%
}

\NewDocumentEnvironment{advexamplebox}{m}{%
  \refstepcounter{advexample}%
  \begin{tcolorbox}[
    enhanced,breakable,colback=white,colframe=RuleGray,boxrule=0.5pt,arc=0mm,
    title={例~\theadvexample\quad #1},fonttitle=\bfseries,
    coltitle=Ink,colbacktitle=PaleBlue,borderline west={1.4pt}{0pt}{DeepBlue}]
}{\end{tcolorbox}}

\NewDocumentEnvironment{advsolution}{+b}{%
  \par\smallskip\begingroup\small\noindent
  \LectureSolutionHeading #1\par\endgroup\smallskip
}{}

\NewDocumentEnvironment{advdefinition}{m +b}{%
  \refstepcounter{advdefinition}%
  \begin{tcolorbox}[
    enhanced,breakable,colback=white,colframe=RuleGray,boxrule=0.5pt,arc=0mm,
    title={定义~\theadvdefinition\quad #1},fonttitle=\bfseries,
    coltitle=Ink,colbacktitle=PaleGold,borderline west={1.2pt}{0pt}{Gold}]
  #2
  \end{tcolorbox}
}{}

\NewDocumentEnvironment{advtheorem}{m +b}{%
  \refstepcounter{advtheorem}%
  \begin{tcolorbox}[
    enhanced,breakable,colback=white,colframe=RuleGray,boxrule=0.5pt,arc=0mm,
    title={定理~\theadvtheorem\quad #1},fonttitle=\bfseries,
    coltitle=Ink,colbacktitle=PaleBlue,borderline west={1.2pt}{0pt}{DeepBlue}]
  #2
  \end{tcolorbox}
}{}

\newtcolorbox{starread}[1]{
  enhanced,breakable,colback=PaleGold,colframe=RuleGray,boxrule=0.5pt,arc=0mm,
  title={选读：#1},fonttitle=\bfseries,coltitle=Ink,colbacktitle=PaleGold,
  borderline west={1.2pt}{0pt}{Gold}}

\newtcolorbox{methodbox}[1]{
  enhanced,breakable,colback=PaleBlue,colframe=RuleGray,boxrule=0.5pt,arc=0mm,
  title={方法：#1},fonttitle=\bfseries,coltitle=Ink,colbacktitle=PaleBlue,
  borderline west={1.2pt}{0pt}{DeepBlue}}
\newtcolorbox{warningbox}[1]{
  enhanced,breakable,colback=SoftGray,colframe=RuleGray,boxrule=0.5pt,arc=0mm,
  title={易错提醒：#1},fonttitle=\bfseries,coltitle=Ink,colbacktitle=SoftGray,
  borderline west={1.2pt}{0pt}{TextGray}}
\newtcolorbox{extensionbox}[1]{
  enhanced,breakable,colback=PaleGold,colframe=RuleGray,boxrule=0.5pt,arc=0mm,
  title={拓展：#1},fonttitle=\bfseries,coltitle=Ink,colbacktitle=PaleGold,
  borderline west={1.2pt}{0pt}{Gold}}

\newcommand{\Needspace}[1]{%
  \par\ifdim\dimexpr\pagegoal-\pagetotal\relax<#1\relax\newpage\fi}

\newcommand{\bookcaption}[1]{%
  \refstepcounter{figure}%
  \par\nopagebreak[4]\smallskip
  {\centering\small\color{TextGray}图~\thefigure：#1\par}\smallskip}
\newcommand{\booktablecaption}[1]{%
  \refstepcounter{table}%
  \par\nopagebreak[4]\smallskip
  {\centering\small\color{TextGray}表~\thetable：#1\par}\smallskip}
\newcommand{\diagramcaption}[1]{\bookcaption{#1}}
\newcommand{\figcap}[1]{\bookcaption{#1}}
\newcommand{\origcaption}[1]{\bookcaption{#1}}
% Keep newly authored diagrams and their complete captions on the same page.
\NewDocumentEnvironment{qualityfigure}{}{%
  \par\begin{center}\begin{minipage}{\linewidth}\centering
}{\end{minipage}\end{center}}
\newcommand{\historyfigure}[3][0.9\linewidth]{%
  \begin{center}\begin{minipage}{#1}\centering
  \includegraphics[width=\linewidth]{#2}\par\smallskip
  \bookcaption{#3}
  \end{minipage}\end{center}}
\newcommand{\sourceimage}[3][0.9\linewidth]{\historyfigure[#1]{#2}{#3}}
\newcommand{\redrawimage}[2][0.72\linewidth]{%
  \begin{center}\includegraphics[width=#1]{#2}\end{center}}
\newcommand{\unitheading}[2]{}
\newcommand{\examplesource}[1]{}

\tikzset{
  every node/.append style={fill=none,inner sep=1.2pt},
  clearlabel/.style={fill=none,inner sep=1pt,text=Ink},
  point/.style={circle,fill=Ink,draw=none,inner sep=1.7pt},
  vector/.style={-{Latex[length=2.5mm]},very thick,draw=DeepBlue},
  component/.style={-{Latex[length=2.2mm]},thick,draw=TextGray},
  axis/.style={-{Latex[length=2mm]},thin,draw=TextGray},
  guide/.style={densely dashed,draw=TextGray},
  curve/.style={very thick,draw=DeepBlue},
  ground/.style={very thick,draw=Ink},
  spring/.style={decorate,decoration={zigzag,segment length=4pt,amplitude=2pt},thick,draw=DeepBlue},
  tangent/.style={thick,draw=TextGray},
  secant/.style={thick,draw=DeepBlue},
  kaxis/.style={-{Latex[length=2mm]},thin,draw=TextGray},
  kcurve/.style={very thick,draw=DeepBlue},
  kguide/.style={densely dashed,draw=TextGray}
}

\hypersetup{
  pdftitle={高中物理竞赛一轮讲义：2026年10月24日终审版},
  pdfsubject={32章、附录与高阶选读，含完整解答及可编辑TikZ配图}
}

\begin{document}
\frontmatter
\begin{titlepage}
  \centering
  \vspace*{0.22\textheight}
  {\Huge\bfseries 高中物理竞赛一轮讲义\par}
  \vspace{1.2em}
  {\Large 2026年10月24日终审版\par}
  \vfill
  {\color{TextGray}数学工具·力学·热学·电磁学·光学·近代物理\par}
\end{titlepage}
\chapter*{阅读与使用}
\addcontentsline{toc}{chapter}{阅读与使用}

本书以理论建模、数学推导和图像分析为主线。阅读一章的目标，是能说明一个模型适用的条件，能从图中建立方程，并能检查结论，而不只是记住最后一个公式。题图给出对象、约束和符号；解答图进一步标出分量、微元、过程或极限。阅读时应把文字、公式和图来回对照，特别留意正方向、参考系、积分区域以及状态与过程的区别。

\section*{先建立基础，再进入选读}
数学工具章可按当前专题的需要反复查阅。学习运动学与力学时，先熟悉位置、速度、受力图和守恒条件，再使用微分、积分表达连续变化。接着学习振动、引力、波和流体，把前面的方法用于新的约束与边界。热学中先分清状态量、过程量和功热符号，再学习概率、分布与涨落；电磁学中先建立场、势、电路与感应的基本关系，再进入传播和材料响应。光学与近代物理既延续这些方法，也会引入新的模型假设。

标有“选读”的部分用于扩展理论视野。第一遍可以先掌握相应章节的基本模型和常规例题，再回头阅读多自由度振动、信息论或场论等较深内容。高阶选读需要更熟练的微积分、线性代数和相关物理基础。高级定义出现得多，并不意味着基础已经掌握；能否清楚处理一个新的简单情形，往往比能否复述复杂公式更能检验理解。

\section*{用例题训练独立推导}
先遮住解答，把题设改写成自己的图：哪些量已知，哪些待求，哪些边界保持固定，是否允许接触、脱离、滑动或穿越。再选研究系统与参考系，写出控制关系，并说明采用的近似。做到这一步后再阅读解答，比较的是模型选择与推导次序，不能只对照末行数值。

完成计算后，至少检查量纲、符号与一个能直接理解的极限。求到代数根时还要检查定义域与原方程；求到平衡点时还要区别存在性与稳定性；使用无穷过程或小量展开时，应说明极限的取法与误差的阶。发现题图与算式不一致时，先回到定义核对，而不是为了迎合一个熟悉答案而临时改变正方向。

% CompletenessTail20261024 guide-and-notation
\par\Needspace{.72\textheight}
\section*{把理论理解接到实际训练}
例题的图和连续曲线表示模型中的理想关系，实际测量还包含仪器分辨率、随机起伏、系统偏差和模型失配。理论学习应配合独立作答、限时表达与实验数据处理。涉及实验的内容，需要实际进行器材读数、方案搭建、重复测量和异常分析；只阅读原理图不能代替操作训练。

教师可根据学生基础安排学习次序与选读深度。每次学习应留下可检查的产物：一张标清条件的图、一段有依据的推导，以及对结果有效范围的说明。这样才能把读懂例题逐步变成在陌生问题中建立模型的能力。

% Completeness20261024: put the notation beside the reading guide.
\section*{记号与单位约定}
\phantomsection
\addcontentsline{toc}{section}{记号与单位约定}
除题目另有说明外，本书采用国际单位制。数值计算所用物理常数和近似值以题目或所在节的说明为准。

矢量通常以粗体表示，部分图示与推导使用箭头记号；两者均表示具有大小和方向的量。
微分记号写作 $\dd$，虚数单位写作 $\ii$。
同一字母在不同专题中可能表示不同物理量，应结合本节定义及单位辨认；
例如 $I$ 可表示电流、转动惯量或光强，$L$ 可表示长度、角动量大小或电感。

除专门说明角度单位的题目外，微积分与小角展开中的角度以弧度计。
对数的底、正方向及参考系均以局部定义为准。各章例题、图和公式均按“章号.顺序号”编号。


\begingroup
\setlength{\parskip}{0pt}
\makeatletter
\renewcommand{\@pnumwidth}{2.6em}
\renewcommand{\@tocrmarg}{3.6em}
\makeatother
\tableofcontents
\clearpage
\endgroup
\mainmatter

\begingroup
\tcbset{breakable/.default=false}
\chapter{数学工具：量纲、矢量、坐标与微积分}

\section{量纲、尺度估计与无量纲化}

物理等式两端必须具有相同量纲。以质量、长度、时间为基本量纲，记作
$M,L,T$。若某物理量 $Q$ 的量纲为
\begin{equation}
 \dimens Q=M^\alpha L^\beta T^\gamma,
\end{equation}
则 $\alpha,\beta,\gamma$ 只描述它随单位选择怎样缩放，并不包含数值系数。
量纲分析能排除错误公式、确定幂律形式，却通常不能给出无量纲常数，也不能分辨具有相同量纲而物理意义不同的量。

若问题含 $n$ 个有量纲变量，而它们只涉及 $r$ 个相互独立的基本量纲，则可组成 $n-r$ 个独立无量纲组合。
把控制方程改写成这些组合的关系称为无量纲化。它不仅减少参数数目，还会显露决定不同物理区间的尺度比。

\begin{origexamplebox}{单摆周期的量纲结构}
长度为 $\ell$ 的单摆在重力场中作小振动。只用 $\ell$、重力加速度 $g$ 和摆球质量 $m$，由量纲分析确定周期的幂律形式，并说明量纲分析不能确定什么。
\end{origexamplebox}
\Needspace{6\baselineskip}\begin{xesolution}
设 $T_0=C\ell^a g^b m^c$，其中 $C$ 无量纲。比较质量量纲先得 $c=0$；再由
\[
 [T_0]=T,\qquad [\ell^ag^b]=L^{a+b}T^{-2b}
\]
得到 $a+b=0$、$-2b=1$，故 $a=1/2$、$b=-1/2$：
\[
 \boxed{T_0=C\sqrt{\frac{\ell}{g}}}.
\]
质量不进入结果。量纲分析不能确定 $C=2\pi$，也不能判断大振幅时还会出现怎样的无量纲振幅修正。
单位检查：$\ell/g$ 的量纲为 $T^2$，开方后确为时间。
\end{xesolution}

\begin{origexamplebox}{黏性流体中小球的终端速度}
半径为 $R$、密度为 $\rho_s$ 的小球，在密度为 $\rho_f$、动力黏度为 $\eta$ 的静止流体中由静止释放，且 $\rho_s>\rho_f$。为判断黏性是否主导流动，定义以球直径 $2R$ 为特征长度的雷诺数
\[
 \mathrm{Re}=\frac{\rho_fv(2R)}{\eta}.
\]
当 $\mathrm{Re}\ll1$ 时，黏滞阻力方向与速度相反，大小可写成 $F_D=C\eta Rv$，其中 $C$ 是与球形几何有关的无量纲常数。

(1) 先用尺度分析估计终端速度 $v_t$；

(2) 若进一步采用斯托克斯阻力 $F_D=6\pi\eta Rv$，求 $v_t$ 的准确表达式；

(3) 把求得的 $v_t$ 代入 $\mathrm{Re}$，写出斯托克斯阻力适用的参数条件。
\end{origexamplebox}
\Needspace{6\baselineskip}\begin{xesolution}
先把物理过程分成两个阶段：刚释放时小球加速下落；速度增大后，向上的黏滞阻力也增大。当加速度降为零时，小球保持恒定速度，这个恒定速度才叫终端速度。

图~\ref{ch01:fig:terminal-speed} 中，小球向下运动，重力向下，浮力与黏滞阻力向上。终端状态满足
\begin{equation}
 mg=F_b+F_D.
\end{equation}
小球体积为 $V=4\pi R^3/3$，所以重力减浮力为
\[
 mg-F_b=(\rho_s-\rho_f)Vg
 =\frac{4\pi}{3}(\rho_s-\rho_f)R^3g.
\]

只做尺度分析时，忽略 $4\pi/3$ 以及阻力定律中的无量纲系数。用
$F_D\sim\eta Rv_t$ 与重力减浮力平衡，得到
\begin{equation}
 \boxed{v_t\sim\frac{(\rho_s-\rho_f)gR^2}{\eta}}.
\end{equation}
符号“$\sim$”表示两边具有相同的尺度和参数依赖关系，并不表示数值严格相等。

若采用准确的斯托克斯阻力，终端平衡式成为
\[
 \frac{4\pi}{3}(\rho_s-\rho_f)R^3g=6\pi\eta Rv_t,
\]
因而
\begin{equation}
 \boxed{v_t=\frac{2(\rho_s-\rho_f)gR^2}{9\eta}}.
\end{equation}
这说明量纲分析能够确定 $R^2$、$g$、密度差和黏度怎样进入结果，却不能得到数值因子 $2/9$。

雷诺数不是另一条凭空给出的经验公式。惯性项的尺度为
$\rho_fv^2/L$，黏性项的尺度为 $\eta v/L^2$，二者之比为
\[
 \frac{\rho_fv^2/L}{\eta v/L^2}=\frac{\rho_fvL}{\eta}.
\]
取小球直径 $L=2R$，就得到题设中的 $\mathrm{Re}$。它确为无量纲量；若改用半径 $R$ 作特征长度，只会相差无关紧要的常数因子 $2$。

把准确的终端速度代入，终端状态的雷诺数为
\begin{equation}
 \mathrm{Re}_t=\frac{\rho_f v_t(2R)}{\eta}
 =\frac{4\rho_f(\rho_s-\rho_f)gR^3}{9\eta^2}.
\end{equation}
所以斯托克斯解答自洽的参数条件是
\begin{equation}
 \boxed{\frac{4\rho_f(\rho_s-\rho_f)gR^3}{9\eta^2}\ll1}.
\end{equation}
题目没有给出具体数值，因此不能预先断言该条件一定成立，只能在得到参数或实验数据后检验。量纲检查还给出
$[(\rho_s-\rho_f)gR^2/\eta]=\mathrm{m/s}$。若 $\mathrm{Re}_t$ 不再远小于 $1$，阻力不再与速度成正比，本题的模型必须更换。
小球由静止单调加速到终端速度的过程中始终有 $0\le v\le v_t$，而雷诺数与 $v$ 成正比，
因此终端状态的 $\mathrm{Re}_t$ 是整个下落过程中的最大值；只要 $\mathrm{Re}_t\ll1$，此前各时刻也都满足低雷诺数条件。

\begin{center}
\begin{tikzpicture}[scale=.9,>=Latex]
  \fill[DeepBlue!8] (-3.2,-2.0) rectangle (3.2,2.15);
  \draw[DeepBlue!35] (-3.2,1.55)--(3.2,1.55);
  \shade[ball color=DeepBlue!24] (0,0) circle (.72);
  \draw[vector] (0,-.05)--(0,-1.65) node[midway,left=5pt] {$mg$};
  \draw[vector] (0,.05)--(0,1.65) node[midway,left=5pt] {$F_b+F_D$};
  \draw[vector,orange!85!black] (1.35,.55)--(1.35,-1.15)
    node[midway,right=5pt] {$v_t$};
  \node[anchor=west] at (1.75,.55) {向下匀速};
  \node[anchor=west,TextGray] at (-3.0,1.87) {静止流体：$\rho_f,\eta$};
  \node[anchor=east] at (-.9,.05) {小球：$R,\rho_s$};
\end{tikzpicture}
\figcap{终端状态的受力图：速度向下，浮力与黏滞阻力的合力向上。}
\label{ch01:fig:terminal-speed}
\end{center}
\end{xesolution}

\section{点乘、投影、叉积与方向判断}

\subsection{点乘与投影}

两个矢量的点乘（scalar product）定义为
\begin{equation}
 \boxed{\vect a\cdot\vect b=ab\cos\theta
 =a_xb_x+a_yb_y+a_zb_z},
\end{equation}
其中 $0\le\theta\le\pi$ 是两矢量的夹角。点乘的结果是标量：同向时为正，垂直时为零，夹角为钝角时为负。物理中它用来取“沿某方向的分量”，也用来表示功
$W=\vect F\cdot\vect s$。点乘满足交换律，且 $\vect a\cdot\vect a=a^2$。

令 $\uvec b=\vect b/b$ 是 $\vect b$ 方向的单位矢量，则 $\vect a$ 在该方向上的标量投影为
\begin{equation}
 a_b=\vect a\cdot\uvec b=a\cos\theta.
\end{equation}
$a_b$ 可以为负，它只给出带符号的分量大小。把这个大小再乘方向单位矢量，才得到矢量投影：
\begin{equation}
 \boxed{\proj_{\vect b}\vect a
 =(\vect a\cdot\uvec b)\uvec b
 =\frac{\vect a\cdot\vect b}{b^2}\vect b}.
\end{equation}
这里要求 $\vect b\ne\vect 0$，否则“沿 $\vect b$ 的方向”没有定义。

\begin{center}
\begin{tikzpicture}[scale=.95,>=Latex]
  \coordinate (O) at (0,0);
  \draw[axis] (-.3,0)--(6,0) node[right]{$\vect b$方向};
  \draw[vector] (O)--(4.6,2.4) node[above right]{$\vect a$};
  \draw[component] (O)--(4.6,0) node[midway,below]{$\proj_{\vect b}\vect a$};
  \draw[guide] (4.6,2.4)--(4.6,0);
  \draw[->,DeepBlue] (1.2,0) arc[start angle=0,end angle=27,radius=1.2];
  \node at (1.45,.32) {$\theta$};
  \node[point,label={[clearlabel]below left:$O$}] at (O) {};
\end{tikzpicture}
\figcap{矢量投影的几何意义；虚线表示从 $\vect a$ 端点向 $\vect b$ 所在直线所作的垂线。}\label{ch01:fig:projection}
\end{center}
图~\ref{ch01:fig:projection} 中垂足把 $\vect a$ 分成平行、垂直于 $\vect b$ 的两部分。若 $\theta<90^\circ$，投影与 $\vect b$ 同向；若 $\theta>90^\circ$，投影反向。后续的受力分解、功和通量计算都会使用这一结构。

\subsection{叉积与右手定则}

叉积（vector product）的结果仍是矢量，其模和方向分别为
\begin{equation}
 |\vect a\times\vect b|=ab\sin\theta,
 \qquad \vect a\times\vect b\perp\vect a,
 \quad \vect a\times\vect b\perp\vect b.
\end{equation}
方向必须使用右手定则判断：右手四指先指向 $\vect a$，再沿两矢量之间不超过 $180^\circ$ 的较小转角弯向 $\vect b$，此时伸直的拇指方向就是 $\vect a\times\vect b$。
选定右手直角坐标系后，用 $\uvec i,\uvec j,\uvec k$ 分别表示沿 $x,y,z$ 轴正方向的单位矢量，
其长度均为 $1$。叉积的顺序不能交换：
\begin{equation}
 \vect a\times\vect b=-\vect b\times\vect a,
 \qquad
 \uvec i\times\uvec j=\uvec k,
 \quad \uvec j\times\uvec k=\uvec i,
 \quad \uvec k\times\uvec i=\uvec j.
\end{equation}
同向或反向矢量的叉积为零。叉积是有向面积、力矩、角动量和洛伦兹力的共同语言。

\begin{center}
\begin{tikzpicture}[scale=.9,>=Latex]
  \coordinate (O) at (0,0);
  \draw[vector] (O)--(4.2,0) node[right] {$\vect a$};
  \draw[vector] (O)--(2.45,2.45) node[above right] {$\vect b$};
  \draw[->,DeepBlue,thick] (1.1,0) arc[start angle=0,end angle=45,radius=1.1];
  \node at (1.35,.48) {$\theta$};
  \node[font=\Large,DeepBlue] at (-.48,.45) {$\boldsymbol\odot$};
  \node[anchor=east] at (-.78,.45) {$\vect a\times\vect b$};
  \node[anchor=west,align=left,TextGray] at (4.85,1.35)
    {四指：$\vect a\rightarrow\vect b$\\拇指：指向纸外};
\end{tikzpicture}
\figcap{从纸面正面观察，$\vect a$ 逆时针转向 $\vect b$，故 $\vect a\times\vect b$ 指向纸外；圆点符号 $\odot$ 表示箭头朝向读者。}
\label{ch01:fig:right-hand-rule}
\end{center}

\begin{origexamplebox}{力矩的坐标计算与方向}
在 $xy$ 平面内，质点位置为 $\vect r=(2.0\,\uvec i+1.0\,\uvec j)\,\mathrm m$，受力
$\vect F=(3.0\,\uvec i-4.0\,\uvec j)\,\mathrm N$。求力对原点的力矩，并解释正负号。
\end{origexamplebox}
\Needspace{6\baselineskip}\begin{xesolution}
先把力矩写成基矢的叉积，使用分配律逐项计算。平行基矢的叉积为零，
$\uvec i\times\uvec j=\uvec k$，而交换顺序后
$\uvec j\times\uvec i=-\uvec k$，所以
\begin{align*}
 \vect r\times\vect F
 &=(2\uvec i+\uvec j)\times(3\uvec i-4\uvec j)\,\mathrm{N\,m}\\
 &=\bigl(6\uvec i\times\uvec i-8\uvec i\times\uvec j
       +3\uvec j\times\uvec i-4\uvec j\times\uvec j\bigr)\,\mathrm{N\,m}\\
 &=(-8-3)\uvec k\,\mathrm{N\,m}.
\end{align*}
一般地，把三维分量分别相乘并合并同一基矢，得到
\[
 \vect a\times\vect b
 =(a_yb_z-a_zb_y)\uvec i+(a_zb_x-a_xb_z)\uvec j
 +(a_xb_y-a_yb_x)\uvec k.
\]
常把这三项缩写成下面的三阶行列式。沿第一行展开时，各项符号依次为正、负、正；
划去该项所在的行与列，剩下的二阶式按“左上乘右下，减去右上乘左下”计算。
对本题补上两个为零的 $z$ 分量，便有
\[
 \vect\tau=\vect r\times\vect F
 =\begin{vmatrix}
 \uvec i&\uvec j&\uvec k\\2&1&0\\3&-4&0
 \end{vmatrix}
 =(-8-3)\uvec k\,\mathrm{N\,m}.
\]
故
\[
 \boxed{\vect\tau=-11\,\uvec k\,\mathrm{N\,m}}.
\]
$-\uvec k$ 指向纸内；从 $+z$ 轴正向朝原点观察，力使质点呈顺时针转动趋势。量纲为力乘力臂，即 $\mathrm{N\,m}$；它与能量虽同量纲，物理对象不同。
\end{xesolution}

\section{常用坐标系与微元}

这里的“微元”是把一条很短的线、一小块曲面或一个很小的体积写成可积分的量。$\dd s$ 表示长度微元，$\dd A$ 表示面积微元，$\dd V$ 表示体积微元；它们分别具有长度、面积、体积的量纲。

先约定坐标的含义：直角坐标为 $(x,y,z)$；平面极坐标 $(r,\varphi)$ 中，
$r$ 是到平面原点的距离，$\varphi$ 是从 $+x$ 轴朝 $+y$ 轴量起的方位角。
柱坐标 $(\rho,\varphi,z)$ 中，$\rho$ 是到 $z$ 轴的垂直距离，
不是前面量纲分析中使用的密度；$\varphi$ 为水平投影的方位角。
球坐标 $(r,\theta,\varphi)$ 的角度约定见下面的球面图。
这四种坐标的线元分别为
\begin{align}
 \dd s^2&=\dd x^2+\dd y^2+\dd z^2,\\
 \dd s^2&=\dd r^2+r^2\dd\varphi^2,\\
 \dd s^2&=\dd \rho^2+\rho^2\dd\varphi^2+\dd z^2,\\
 \dd s^2&=\dd r^2+r^2\dd\theta^2+r^2\sin^2\theta\,\dd\varphi^2.
\end{align}
相应面积或体积微元来自各坐标方向的尺度因子。采用图~\ref{ch01:fig:spherical} 的球坐标约定：$r$ 是点到原点的距离，$\theta$ 从 $+z$ 轴向下量，$\varphi$ 是绕 $z$ 轴、从 $+x$ 轴朝 $+y$ 轴量的方位角。

在半径固定的球面上，把 $\theta$ 增加 $\dd\theta$，沿经线得到的小弧长是 $r\dd\theta$；把 $\varphi$ 增加 $\dd\varphi$，所在纬圈的半径为 $r\sin\theta$，所以另一条小弧长是 $r\sin\theta\dd\varphi$。两者相乘就是图中黄色小面片的面积：
\begin{equation}
 \boxed{\dd A_{r=\text{常量}}
 =(r\dd\theta)(r\sin\theta\dd\varphi)
 =r^2\sin\theta\,\dd\theta\,\dd\varphi}.
\end{equation}
$\dd A$ 是面积大小，是标量；若计算电通量或流量，需要带方向的面积矢量
\begin{equation}
 \dd\vect A=\uvec e_r\,\dd A,
 \qquad
 \dd V=\dd r\,\dd A
 =r^2\sin\theta\,\dd r\,\dd\theta\,\dd\varphi.
\end{equation}
这里 $\uvec e_r$ 是球面外法线方向。$\sin\theta$ 不是额外的物理假设，而是纬圈弧长随极角缩短产生的几何因子。

\begin{center}
\begin{tikzpicture}[scale=.86,>=Latex]
 \begin{scope}[xshift=-3.7cm]
  \coordinate (O) at (0,0);
  \coordinate (X) at (-2.75,-1.55);
  \coordinate (Y) at (3.3,-1.3);
  \coordinate (Z) at (0,3.45);
  \coordinate (Q) at (1.05,-1.18);
  \coordinate (P) at (1.05,2.05);
  \draw[axis] (O)--(X) node[below left] {$x$};
  \draw[axis] (O)--(Y) node[right] {$y$};
  \draw[axis] (O)--(Z) node[above] {$z$};
  \draw[guide] (O)--(Q);
  \draw[guide] (Q)--(P);
  \draw[vector] (O)--(P) node[above right] {$P(r,\theta,\varphi)$};
  \node[below right] at (Q) {$P_{xy}$};
  \draw[->,DeepBlue,thick] (0,1.08) arc[start angle=90,end angle=62,radius=1.08];
  \node[DeepBlue] at (.35,1.50) {$\theta$};
  \draw[->,TextGray,thick] (-.85,-.48) .. controls (-.15,-1.22) .. (.68,-.77);
  \node[TextGray] at (-1.10,-1.30) {$\varphi$};
  \draw[TextGray,thin] (-.90,-1.21)--(-.22,-.86);
  \node[point,label={[clearlabel]above left:$O$}] at (O) {};
  \node[align=center,TextGray] at (0,-2.85)
    {$\varphi$ 的正向为 $+x\rightarrow +y$\\$\uvec i\times\uvec j=\uvec k$（右手系）};
 \end{scope}
 \begin{scope}[xshift=4.65cm,yshift=.05cm]
  \coordinate (C) at (0,0);
  \filldraw[draw=DeepBlue,fill=yellow!55]
    (-1.75,-.52)--(1.18,-.28)--(1.55,.7)--(-1.38,.46)--cycle;
  \draw[vector] (C)--(0,2.12) node[above] {$\dd\vect A=\uvec e_r\,\dd A$};
  \draw[<->,TextGray] (-1.62,-.78)--(1.23,-.55)
    node[midway,below=5pt] {$r\sin\theta\,\dd\varphi$};
  \draw[<->,TextGray] (1.42,-.3)--(1.78,.62)
    node[midway,right=6pt] {$r\,\dd\theta$};
  \node[align=center] at (0,-1.72)
    {图中黄色小面片的面积就是 $\dd A$};
  \node[align=center,TextGray] at (0,-2.3)
    {$\dd A=(r\,\dd\theta)(r\sin\theta\,\dd\varphi)$};
 \end{scope}
\end{tikzpicture}
\figcap{左：球坐标建立在右手直角坐标系上；右：球面面积微元的切平面放大图。}\label{ch01:fig:spherical}
\end{center}
图~\ref{ch01:fig:spherical} 的角度约定决定公式中的 $\sin\theta$；若把 $\theta$ 改为从赤道面量起，尺度因子的写法也必须同步改变。特别要区分 $\dd A$ 与 $\dd\vect A$：前者只有面积大小，后者还带球面外法线方向。

\begin{origexamplebox}{球对称密度的总质量}
半径为 $R$ 的球体密度为 $\rho(r)=\rho_0(1-r/R)$。求总质量，并检查结果的符号和量纲。
\end{origexamplebox}
\Needspace{6\baselineskip}\begin{xesolution}
\textbf{先把“累加微元”写清楚。}把球体分成许多很薄的球壳，第 $j$ 层的质量近似为
$\rho(r_j)\Delta V_j$，总质量近似为这些质量之和。
当每层厚度趋于零时，这个和的极限称为定积分，符号 $\int$ 表示连续累加；
上下限表示累加范围，$\dd r$ 对应逐渐变小的层厚，而不是一个额外的未知量。
对于半径 $r$、厚度 $\Delta r$ 的球壳，
\[
 \Delta V=\frac{4\pi}{3}\bigl[(r+\Delta r)^3-r^3\bigr]
 =4\pi r^2\Delta r+4\pi r(\Delta r)^2+\frac{4\pi}{3}(\Delta r)^3.
\]
与第一项相比，后两项在分层无限变细时不再贡献有限总量，故 $\dd V=4\pi r^2\dd r$。
本题所需的两个累加公式也可先由有限求和得到。令 $\Delta r=R/n$、$r_j=jR/n$，利用
\[
 \sum_{j=1}^n j^2=\frac{n(n+1)(2n+1)}6,\qquad
 \sum_{j=1}^n j^3=\frac{n^2(n+1)^2}4,
\]
分别乘以 $(R/n)^3$、$(R/n)^4$ 并令 $n\to\infty$，便得到
\[
 \int_0^R r^2\dd r=\frac{R^3}{3},\qquad
 \int_0^R r^3\dd r=\frac{R^4}{4}.
\]
因此球壳累加已经足够完成本题，不必先掌握多重积分的一般运算。

也可使用前文的球坐标体积元
$\dd V=r^2\sin\theta\,\dd r\,\dd\theta\,\dd\varphi$。因为密度只依赖 $r$，
两个角变量可以先积分：
\[
 \int_0^{2\pi}\!\dd\varphi
 \int_0^\pi\!\sin\theta\,\dd\theta=4\pi.
\]
这里第一项是在长度为 $2\pi$ 的角区间上累加 $1$，结果为 $2\pi$；
第二项等于 $[-\cos\theta]_0^\pi=-\cos\pi+\cos0=2$。
方括号表示把上限与下限分别代入后相减；由于 $-\cos\theta$ 的变化率是
$\sin\theta$，它的总变化就等于相应的累加量。这种“由变化率求总变化”的计算法将在积分部分系统说明。
因此半径在 $r$ 与 $r+\dd r$ 之间的薄球壳体积是 $4\pi r^2\dd r$，总质量为
\begin{align*}
 M&=\int_V\rho\,\dd V
 =4\pi\rho_0\int_0^R\left(1-\frac rR\right)r^2\dd r\\
 &=4\pi\rho_0\left(\frac{R^3}{3}-\frac{R^3}{4}\right)
 =\boxed{\frac{\pi}{3}\rho_0R^3}.
\end{align*}
被积函数在 $0\le r\le R$ 非负，所以总质量为正；$\rho_0R^3$ 的量纲为质量。
这里的 $r^2\sin\theta$ 正是把三个坐标微元换成实际体积的几何因子；它给出薄球壳体积
$4\pi r^2\dd r$，从而把本题的三重积分化成只对半径 $r$ 的一重积分。
\end{xesolution}


{
\graphicspath{{/Users/dongxingxi/Documents/物理竞赛一轮资料/sample/xes-original/fourteenth-section/}}
\section{导数及其应用}
% TikZ figures redrawn from the source.  Blue annotations in solution figures
% are added for the present edition.

\newcommand{\extremaDiagram}{%
\begin{center}
\begin{tikzpicture}[x=0.78cm,y=0.62cm,>=Latex]
  \foreach \s in {0,4,8,12}{
    \draw[axis] (\s,0)--++(3,0) node[right] {$x$};
    \draw[axis] (\s+0.25,-0.35)--++(0,2.75) node[above] {$y$};
  }
  \draw[curve,domain=-1.05:1.05,smooth,variable=\t]
    plot ({1.65+\t},{0.25+1.35*\t*\t});
  \draw[tangent] (0.9,0.25)--(2.4,0.25);
  \node[below] at (1.6,-0.35) {(1)};
  \draw[curve,domain=-1.05:1.05,smooth,variable=\t]
    plot ({5.65+\t},{1.9-1.2*\t*\t});
  \draw[tangent] (4.9,1.9)--(6.4,1.9);
  \node[below] at (5.6,-0.35) {(2)};
  \draw[curve,domain=-1.15:1.15,smooth,variable=\t]
    plot ({9.65+\t},{1.05+0.42*\t*\t*\t});
  \draw[tangent] (8.9,1.05)--(10.4,1.05);
  \node[below] at (9.6,-0.35) {(3)};
  \draw[curve,smooth] plot coordinates
    {(12.65,1.1) (12.9,1.85) (13.25,1.2) (13.55,1.2)
     (13.85,0.45) (14.2,1.9) (14.6,1.1)};
  \draw[tangent] (12.75,1.85)--(13.35,1.85);
  \draw[tangent] (13.2,1.2)--(13.75,1.2);
  \draw[tangent] (13.55,0.45)--(14.1,0.45);
  \draw[tangent] (13.9,1.9)--(14.45,1.9);
  \node[below] at (13.6,-0.35) {(4)};
\end{tikzpicture}
\figcap{函数的导数为零的各种情况}\label{ch01:fig:s14-1-1}
\end{center}}

\newcommand{\probdiagOneFour}{%
\begin{center}
\begin{tikzpicture}[scale=0.9,>=Latex]
  \draw[thick] (-2,0) arc[start angle=180,end angle=0,radius=2];
  \draw[thick] (-2,0)--(2,0);
  \coordinate (C) at (0,0);
  \coordinate (P) at ({2*sin(38)},{2*cos(38)});
  \draw[guide] (C)--(0,2);
  \draw[guide] (C)--(P);
  \draw[thick] (P)--++(0,2.0);
  % 两只固定导向块分列杆的两侧，二者之间留缝让杆竖直滑动。
  \draw[thick,fill=DeepBlue!6]
    ($(P)+(-0.78,1.04)$) rectangle ($(P)+(-0.10,1.30)$);
  \draw[thick,fill=DeepBlue!6]
    ($(P)+(0.10,1.04)$) rectangle ($(P)+(0.78,1.30)$);
  \draw[->,very thick] ($(P)+(0,2.0)$)--++(0,0.65);
  \draw[vector] (2.15,0.25)--++(0.9,0) node[right] {$v_0$};
  \coordinate (Top) at (0,2);
  \draw pic[draw,->,"$\theta$",angle radius=8mm,angle eccentricity=1.35]
    {angle=P--C--Top};
  \node at (-0.95,0.25) {$R$};
\end{tikzpicture}
\end{center}}

\newcommand{\soldiagOneFour}{%
\begin{center}
\begin{tikzpicture}[scale=0.88,>=Latex]
  \draw[thick] (-2,0) arc[start angle=180,end angle=0,radius=2];
  \draw[thick] (-2,0)--(2,0);
  \coordinate (C) at (0,0);
  \coordinate (P) at ({2*sin(38)},{2*cos(38)});
  \draw[guide] (C)--(P);
  \node[below left=2pt] at ($(C)!0.43!(P)$) {$R$};
  \draw[guide] (C)--(0,2);
  \draw[thick] (P)--++(0,2);
  \draw[thick,fill=DeepBlue!6]
    ($(P)+(-0.78,1.04)$) rectangle ($(P)+(-0.10,1.30)$);
  \draw[thick,fill=DeepBlue!6]
    ($(P)+(0.10,1.04)$) rectangle ($(P)+(0.78,1.30)$);
  \draw[vector] (-2.4,0.25)--++(1.1,0) node[above] {$\dot X=v_0$};
  \draw[component] (P)--++(0,1.05);
  \node[right=4pt] at ($(P)+(0,0.78)$) {$\dot y$};
  \draw[component] (P)--++(-1.18,0);
  \node[right=7pt] at (P) {$R\cos\theta\,|\dot\theta|$};
  \coordinate (Top) at (0,2);
  \draw pic[draw,->,"$\theta$",angle radius=8mm,angle eccentricity=1.35]
    {angle=P--C--Top};
  \node[align=center,TextGray] at (0,-0.75)
    {$x_{\rm rod}=X+R\sin\theta$\\$y=R\cos\theta$};
\end{tikzpicture}
\end{center}}

\newcommand{\probdiagOneSix}{%
\begin{center}
\begin{tikzpicture}[scale=0.88,>=Latex]
  \draw[thick] (-3,0)--(3,0);
  \node[right] at (3,0.28) {$n_1$};
  \node[right] at (3,-0.28) {$n_2$};
  \coordinate (A) at (-2.2,2.0);
  \coordinate (P) at (0.35,0);
  \coordinate (B) at (2.25,-1.55);
  \draw[thick] (A)--(P)--(B);
  \draw[guide] (P)--++(0,1.35);
  \draw[guide] (P)--++(0,-1.25);
  \node[left] at (A) {$A$}; \node[above right] at (P) {$P$}; \node[right] at (B) {$B$};
  \coordinate (PN) at (0.35,1);
  \draw pic[draw,->,"$\alpha$",angle radius=7mm,angle eccentricity=1.35]
    {angle=PN--P--A};
  \coordinate (PS) at (0.35,-1);
  \draw pic[draw,<- ,"$\beta$",angle radius=7mm,angle eccentricity=1.35]
    {angle=PS--P--B};
  \draw[<->] (-2.2,-0.35)--(0.35,-0.35);
  \node[above=1pt] at (-0.925,-0.35) {$x$};
  \draw[<->] (-2.2,-0.72)--(2.25,-0.72);
  \node[below=1pt] at (0.025,-0.72) {$L$};
  \draw[<->] (-2.55,0)--node[left] {$h_1$}(-2.55,2);
  \draw[<->] (2.55,0)--node[right] {$h_2$}(2.55,-1.55);
\end{tikzpicture}
\end{center}}

\newcommand{\probdiagOneSeven}{%
\begin{center}
\begin{tikzpicture}[x=1.05cm,y=0.8cm,>=Latex]
  \begin{scope}
    \draw[axis] (0,0)--(3.3,0) node[right] {$V$};
    \draw[axis] (0,0)--(0,3.0) node[above] {$p$};
    \draw[curve] (1,2)--(2,1);
    \draw[guide] (1,0)--(1,2)--(0,2);
    \draw[guide] (2,0)--(2,1)--(0,1);
    \node[above right] at (1,2) {$A$}; \node[above right] at (2,1) {$B$};
    \node[below] at (1,0) {$V_0$}; \node[below] at (2,0) {$2V_0$};
    \node[left] at (0,1) {$p_0$}; \node[left] at (0,2) {$2p_0$};
  \end{scope}
  \begin{scope}[xshift=6.0cm]
    \draw[axis] (0,0)--(3.3,0) node[right] {$V$};
    \draw[axis] (0,0)--(0,3.0) node[above] {$p$};
    \draw[curve] (1.55,1.55) circle[radius=0.78];
    \draw[guide] (1.55,0)--(1.55,1.55)--(0,1.55);
    \node[above right] at (1.55,1.55) {$(p_0,V_0)$};
    \node[below] at (1.55,0) {$V_0$}; \node[left] at (0,1.55) {$p_0$};
  \end{scope}
\end{tikzpicture}
\end{center}}

\newcommand{\probdiagOneEight}{%
\begin{center}
\begin{tikzpicture}[scale=0.86,>=Latex]
  \coordinate (U) at (0,1.45); \coordinate (D) at (0,-1.45); \coordinate (P) at (3.2,0);
  \fill[white] (U) circle (0.21); \draw[thick] (U) circle (0.21);
  \fill[white] (D) circle (0.21); \draw[thick] (D) circle (0.21);
  \fill[white] (P) circle (0.13); \draw[thick] (P) circle (0.13);
  \draw[guide] (U)--(D); \draw[guide] (U)--(P); \draw[guide] (D)--(P);
  \draw[thick] (0,0)--(P);
  \node[above left=4pt] at (U) {$M$};
  \node[below left=4pt] at (D) {$M$};
  \node[left=5pt] at (0,0.72) {$L$};
  \node[left=5pt] at (0,-0.72) {$L$};
  \node[above=4pt] at (1.55,0) {$x$};
  \draw[vector] (P)--++(-1.1,0);
  \node[below=5pt] at (2.70,0) {$\vec F$};
\end{tikzpicture}
\end{center}}

\newcommand{\soldiagOneNine}{%
\begin{center}
\begin{tikzpicture}[scale=0.82,>=Latex]
  \pgfmathsetmacro{\throwD}{(sqrt(3)+3*sqrt(11))/8}
  \draw[thick] (1.9,0)--(6.5,0);
  \coordinate (S) at ({4.6-1.25*\throwD},0); \coordinate (C) at (4.6,1.25);
  \coordinate (K) at ({4.6-1.25*sqrt(3)/2},1.875); \coordinate (T) at (4.6,2.5);
  \draw[thick] (C) circle[radius=1.25];
  \draw[curve,domain=-\throwD:0,samples=90,variable=\x]
    plot ({4.6+1.25*\x},{1.25*(2-\x/sqrt(3)-4*\x*\x/3)});
  \draw[tangent] ($(K)+(-0.38,-0.658)$)--($(K)+(0.38,0.658)$);
  \fill[DeepBlue] (K) circle (1.7pt) node[above left] {相切点};
  \draw[guide] (4.6,0)--(T);
  \draw[<->] ({4.6-1.25*\throwD},-.4)--(4.6,-.4);
  \node[below=1pt] at (3.6875,-0.4) {$d_0$};
  \draw[<->] (6.15,0)--node[right] {$2R$}(6.15,2.5);
  \draw[guide] (T)--(6.3,2.5);
  \draw[vector] (S)--++(.42,{.42*sqrt(11)}) node[left] {$\vec v_0$};
  \coordinate (SR) at ($(S)+(1,0)$); \coordinate (SV) at ($(S)+(.42,{.42*sqrt(11)})$);
  \draw pic[draw,->,"$\theta_0$",angle radius=8mm,angle eccentricity=1.35]
    {angle=SR--S--SV};
  \node[anchor=west] at (5.65,3.10) {球顶目标};
  \draw[TextGray,thin] (5.60,3.02)--(T);
  \node[TextGray,align=center] at (4.1,-1.15)
    {临界轨迹先与球面相切，再到达球顶};
\end{tikzpicture}
\end{center}}

\newcommand{\potentialDiagram}{%
\begin{center}
\begin{tikzpicture}[x=0.9cm,y=0.65cm,>=Latex]
  \begin{scope}
    \coordinate (O) at (0,2.2); \coordinate (M) at (-1.25,0.2);
    \fill (O) circle (1.5pt) node[above right] {$O$};
    \draw[thick] (O)--(M) node[midway,left] {$l$};
    \fill (M) circle (2.5pt) node[left] {$m$};
    \draw[guide] (O)--++(0,-1.45);
    \coordinate (OD) at (0,0.75);
    \draw pic[draw,->,"$\theta$",angle radius=7mm,angle eccentricity=1.35]
      {angle=M--O--OD};
  \end{scope}
  \begin{scope}[xshift=4.4cm]
    \draw[axis] (-0.2,0)--(5.0,0) node[right] {$x$};
    \draw[axis] (0,-0.3)--(0,3.1) node[above] {$V(x)$};
    \draw[curve,smooth] plot coordinates
      {(0.2,2.5) (0.75,1.0) (1.2,0.42) (1.75,1.15) (2.35,2.25)
       (2.95,1.25) (3.6,0.62) (4.15,1.15) (4.7,2.35)};
    \draw[guide] (1.2,0)--(1.2,0.42);
    \draw[guide] (2.35,0)--(2.35,2.25);
    \draw[guide] (3.6,0)--(3.6,0.62);
    \fill[DeepBlue] (1.2,0.42) circle (1.8pt);
    \fill[DeepBlue] (2.35,2.25) circle (1.8pt);
    \fill[DeepBlue] (3.6,0.62) circle (1.8pt);
    \node[below=2pt] at (1.2,0) {$S_1$};
    \node[below=2pt] at (2.35,0) {$U$};
    \node[below=2pt] at (3.6,0) {$S_2$};
    \node[DeepBlue] at (1.0,3.25) {稳定};
    \draw[DeepBlue,thin] (1.0,2.98)--(1.2,.57);
    \node[above=5pt,DeepBlue] at (2.35,2.25) {不稳定};
    \node[DeepBlue] at (3.9,3.25) {稳定};
    \draw[DeepBlue,thin] (3.9,2.98)--(3.6,.77);
  \end{scope}
\end{tikzpicture}
\figcap{左：单摆的角度约定；右：势能曲线的局部极小点与局部极大点。}
\end{center}}

\newcommand{\probdiagOneTen}{%
\begin{center}
\begin{tikzpicture}[scale=0.9,>=Latex]
  \coordinate (B) at (0,0); \coordinate (A) at (3.0,0); \coordinate (M) at (0,-2.0);
  \draw[thick] (0,0.45)--(0,-2.45);
  \draw[thick] (B)--(A);
  \draw[decorate,decoration={coil,aspect=0.35,segment length=4pt,amplitude=3pt},thick]
    (M)--(A);
  \fill (M) circle (2.8pt) node[left] {$m$};
  \fill (A) circle (1.7pt) node[right] {$A$}; \node[left] at (B) {$B$};
  \draw[<->] (0,0.35)--(3,0.35);
  \node[above=1pt] at (1.5,0.35) {$h$};
  \draw[<->] (-0.45,0)--(-0.45,-2);
  \node[left=1pt] at (-0.45,-1) {$x$};
  \node[below right] at (1.55,-1.05) {$\ell=\sqrt{h^2+x^2}$};
\end{tikzpicture}
\end{center}}

\newcommand{\probdiagOneEleven}{%
\begin{center}
\begin{tikzpicture}[scale=0.86,>=Latex]
  \coordinate (O) at (0,3.1);
  \draw[thick] (O)--(-2.05,0); \draw[thick] (O)--(2.05,0);
  \draw[thick] (-2.05,0) arc[start angle=180,end angle=360,x radius=2.05,y radius=0.32];
  \draw[guide] (-2.05,0) arc[start angle=180,end angle=0,x radius=2.05,y radius=0.32];
  \draw[very thick,DeepBlue] (-1.15,1.35) arc[start angle=180,end angle=360,x radius=1.15,y radius=0.22];
  \draw[guide] (-1.15,1.35) arc[start angle=180,end angle=0,x radius=1.15,y radius=0.22];
  \draw[axis] (O)--(0,0.55);
  \draw[guide] (O)--(2.55,3.1);
  \draw[guide] (0,1.35)--(2.55,1.35);
  \draw[<->] (2.4,3.1)--node[right=4pt] {$x$}(2.4,1.35);
  \coordinate (OD) at (0,1.8); \coordinate (RB) at (2.05,0);
  \draw pic[draw,->,"$\alpha$",angle radius=8mm,angle eccentricity=1.62]
    {angle=OD--O--RB};
  \draw[TextGray] (-1.05,1.35)--(-2.25,1.0);
  \node[anchor=east,DeepBlue] at (-2.3,1.0) {$\ell=2\pi x\tan\alpha$};
  \node[above] at (O) {$O$};
\end{tikzpicture}
\end{center}}

\newcommand{\probdiagOneTwelve}{%
\begin{center}
\begin{tikzpicture}[scale=0.9,>=Latex]
  \coordinate (L) at (-1.75,0); \coordinate (R) at (1.75,0); \coordinate (T) at (0,2.35);
  \draw[thick] (-2.5,0)--(2.5,0);
  \draw[thick] (L)--(T)--(R);
  \draw[decorate,decoration={coil,aspect=0.35,segment length=4pt,amplitude=3pt},thick]
    (L)--(R);
  \fill (T) circle (2.8pt) node[above right] {$m$};
  \node[left] at (-1.12,1.15) {$l$}; \node[right] at (1.12,1.15) {$l$};
  \node[above=5pt] at (0,0) {$k$};
  \coordinate (RL) at (-2.5,0);
  \draw pic[draw,->,"$\theta$",angle radius=8mm,angle eccentricity=1.35]
    {angle=T--R--RL};
  \draw[<->] (-1.75,-0.72)--(1.75,-0.72);
  \node[below=3pt] at (0,-0.72) {$2l\cos\theta$};
\end{tikzpicture}
\end{center}}

\newcommand{\incrementDiagram}{%
\begin{center}
\begin{tikzpicture}[x=1.05cm,y=0.85cm,>=Latex]
  \draw[axis] (0,0)--(6,0) node[right] {$x$};
  \draw[axis] (0,0)--(0,4.0) node[above] {$y$};
  \draw[curve,domain=0.7:5.2,smooth,variable=\x]
    plot ({\x},{0.55+0.13*exp(0.62*\x)});
  \coordinate (P) at (4.1,{0.55+0.13*exp(0.62*4.1)});
  \coordinate (Q) at (4.8,{0.55+0.13*exp(0.62*4.8)});
  \draw[guide] (P)--(4.1,0);
  \draw[guide] (Q)--(4.8,0);
  \node[below left=3pt] at (4.1,0) {$x_0$};
  \node[below right=3pt] at (4.8,0) {$x_0+\Delta x$};
  \draw[guide] (P)--(0,{0.55+0.13*exp(0.62*4.1)}) node[left] {$y_0$};
  \draw[guide] (Q)--(0,{0.55+0.13*exp(0.62*4.8)});
  \draw[<->] (4.1,1.0)--(4.8,1.0);
  \node[below=1pt] at (4.45,1.0) {$\Delta x$};
  \draw[<->] (5.2,{0.55+0.13*exp(0.62*4.1)})--node[right] {$\Delta y$}
    (5.2,{0.55+0.13*exp(0.62*4.8)});
  \node[above left] at (3.0,1.7) {$y(x)$};
\end{tikzpicture}
\end{center}}

\newcommand{\probdiagOneSixteen}{%
\begin{center}
\begin{tikzpicture}[scale=0.9,>=Latex]
  \coordinate (A) at (0,-2.2); \coordinate (O) at (0,0); \coordinate (X) at (2.6,0); \coordinate (DX) at (3.35,0);
  \draw[thick,->] (-0.5,0)--(4.2,0);
  \draw[guide] (A)--(O);
  \draw[thick] (A)--(X); \draw[thick] (A)--(DX);
  \fill (A) circle (2.5pt) node[below left] {$A$};
  \node[above] at (O) {$O$};
  \node[above left=1pt] at (X) {$x$};
  \node[above right=1pt] at (DX) {$x+\Delta x$};
  \draw[<->] (-0.35,0)--node[left] {$h$}(-0.35,-2.2);
  \draw pic[draw,->,"$\theta$",angle radius=8mm,angle eccentricity=1.3]
    {angle=X--A--O};
  \draw pic[draw,->,angle radius=17mm]
    {angle=DX--A--X};
  \coordinate (DeltaMid) at (1.36,-1.18);
  \node (DeltaLabel) at (3.8,-1.50) {$\Delta\theta$};
  \draw[guide] (DeltaLabel.west)--(DeltaMid);
\end{tikzpicture}
\end{center}}

\newcommand{\soldiagOneSixteen}{%
\begin{center}
\begin{tikzpicture}[scale=0.9,>=Latex]
  \coordinate (A) at (0,-2.2); \coordinate (O) at (0,0);
  \coordinate (X) at (2.6,0); \coordinate (DX) at (3.35,0);
  \draw[thick,->] (-0.5,0)--(4.2,0);
  \draw[guide] (A)--(O);
  \draw[thick] (A)--(X); \draw[thick] (A)--(DX);
  \fill (A) circle (2.5pt) node[below left] {$A$};
  \node[above] at (O) {$O$};
  \node[above left=1pt] at (X) {$x$};
  \node[above right=1pt] at (DX) {$x+\Delta x$};
  \draw[<->] (-0.35,0)--node[left] {$h$}(-0.35,-2.2);
  \draw pic[draw,->,"$\theta$",angle radius=8mm,angle eccentricity=1.3]
    {angle=X--A--O};
  \draw pic[draw,->,DeepBlue,angle radius=17mm]
    {angle=DX--A--X};
  \node[DeepBlue] (DeltaLabel) at (3.25,-1.75) {$\Delta\theta$};
  \draw[guide] (DeltaLabel.west)--(1.36,-1.18);
  \draw[component] (2.6,.85)--(3.35,.85) node[midway,above] {$\Delta x$};
  \draw[guide] (X)--(2.6,.95) (DX)--(3.35,.95);
  \node[align=left,TextGray] at (6.05,-.75)
    {$\displaystyle \tan\theta=\frac{x}{h}$\\[3pt]
     $\displaystyle \Delta\theta\simeq
       \frac{h\,\Delta x}{h^2+x^2}$};
\end{tikzpicture}
\end{center}}


在过去的学习中，我们已经学习了许多物理定律及其应用。对于一大类问题，通过已知定律
并结合系统分析，已经可以得到结果，其中也包括一些相当复杂的问题。但在学习过程中会
发现，处理某些问题时会遇到数学上的困难。例如，位置随时间的变化率被定义为速度，可
我们还没有掌握由任意给定运动求速度的普遍方法；又如，速度—时间图像与横轴所围面积
表示位移，却还不能对一般的速度函数求任意一段时间内的位移。

历史上，人们在研究引力及其性质时也遇到了类似困难。牛顿（I. Newton）和莱布尼茨
（G. W. Leibniz）分别独立地发明了新的数学工具——微积分（calculus），使这类问题的
处理成为可能。微积分及随之发展起来的数学理论极大地推动了物理学的发展，不仅使我们
能够处理许多与曲线、曲面性质有关的问题，其中也包含着重要的物理思想。

具有光滑图像的单自变量函数在某点处切线的斜率称为函数在该点的\textbf{导数}
（derivative）。若函数记为 $y=y(x)$，曲线在 $x$ 点处的切线可以看作割线的极限，其
斜率为
\begin{equation}
 y'(x)=\lim_{\Delta x\to0}\frac{y(x+\Delta x)-y(x)}{\Delta x}.
\end{equation}
导数有多种表示方法，例如 $y'(x)$、$\dfrac{\dd y}{\dd x}(x)$。物理中遇到的绝大部分
函数都可以求导；除了会计算特定函数的导数，还需要理解导数的物理意义。

直线运动质点的速度是坐标对时间的导数：
\begin{equation}
 v(t)=\frac{\dd x}{\dd t}(t),
\end{equation}
加速度则是速度对时间的导数：
\begin{equation}
 a(t)=\frac{\dd v}{\dd t}(t)=\frac{\dd^2x}{\dd t^2}(t).
\end{equation}
因此加速度也是坐标对时间的二阶导数（second order derivative），常记作 $x''$ 或
$\dfrac{\dd^2x}{\dd t^2}$；类似地还可以定义更高阶导数。

许多物理量本来就是借助导数定义的。例如实验给出一般弹性体的外力与形变量的函数
$F(x)$，它在形变 $x$ 处的切线刚度（局部劲度）可写成
\begin{equation}
 k(x)=\frac{\dd F}{\dd x}(x).
\end{equation}
胡克定律只是切线刚度为常数的特殊情况。又如实验测得某物体吸收的热量与温度之间的关系
$Q(T)$，则该物体在温度 $T$ 处的热容量为
\begin{equation}
 C(T)=\frac{\dd Q}{\dd T}(T).
\end{equation}
对于真实物体，热容量一般会随温度改变。

利用导数的定义和基本运算法则，可以求出物理中绝大部分常见函数的导数。下面通过例题
熟悉基本求导规律。

\subsection{局部变化率、链式法则与偏导数}

前述割线斜率的极限说明，导数不是“两个无穷小量的普通商”，而是一个局部变化率。位置对时间的一阶导数是速度，二阶导数是
加速度；电量对时间的导数是电流；弹性体的切线刚度可写成 $\dd F/\dd x$；若物体吸热量为
$Q(T)$，则温度为 $T$ 时的热容 $C(T)$ 为 $\dd Q/\dd T$。这些定义都在问同一件事：系统处于某个状态时，
再作极小改变会得到怎样的即时响应。

复合关系要用链式法则。若 $y=f(u)$、$u=g(x)$，则
\[
 \frac{\dd y}{\dd x}=\frac{\dd y}{\dd u}\frac{\dd u}{\dd x}.
\]
例如势能 $U(x)$ 随位置变化，而位置又随时间变化，就有
$\dd U/\dd t=(\dd U/\dd x)v$。若函数由多个物理量共同决定，则应说明求偏导时保持哪些量不变；
热学中 $(\partial p/\partial V)_T$ 与 $(\partial p/\partial V)_S$ 就代表不同的物理过程。

\subsection{极值判别、物理约束与近似展开}

若函数在定义域内部的 $x_0$ 处可导并取得局部极值，则必有 $f'(x_0)=0$；但驻点不一定是极值，
还要比较左右导数或二阶导数。若 $f'(x_0)=0$ 且 $f''(x_0)>0$，则 $x_0$ 是严格局部极小点；
若 $f'(x_0)=0$ 且 $f''(x_0)<0$，则 $x_0$ 是严格局部极大点。端点和不可导点也可能
给出最大、最小值，故完整求极值必须连同定义域一起检查。物理答案还要满足长度为正、接触不脱离、
速度可达等约束，纯代数驻点未必可实现。

在 $x_0$ 附近，足够光滑的函数可逐阶展开：
\[
 f(x_0+\xi)=f(x_0)+f'(x_0)\xi+
 \frac12f''(x_0)\xi^2+\frac1{3!}f'''(x_0)\xi^3+\cdots .
\]
近似保留到哪一阶取决于所需精度。若 $x_0$ 是势能的平衡位置，一阶项消失，最低非零项通常是二次项，
系统因而近似简谐振动；若二次项也为零，就必须保留更高阶，回复力和振动周期会表现出不同规律。
小角近似 $\sin\theta\simeq\theta$、$\cos\theta\simeq1-\theta^2/2$ 也来自同一展开，使用前要说明
$|\theta|\ll1$ 且角度以弧度计。

对时间求导时还常使用点记号：$\dot X=\dd X/\dd t$，
$\ddot X=\dd^2X/\dd t^2$；对一般自变量仍写撇号或 $\dd/\dd x$。
导数还提供变化率的量级判断。设某量的特征值为 $X$、变化时间为 $T$，则
$\dot X$ 的量级约为 $X/T$；两个机制谁占主导，可比较相应导数项的无量纲比值。竞赛推导中，
这类先估量级、再保留主导项的做法，常比直接处理完整方程更清楚。


\begin{origexamplebox}{}\label{ch01:ex:s14-1-1}
根据定义求下列函数的导数：
\[
 y(x)=C,\qquad y(x)=kx,\qquad y(x)=ax^n,\qquad
 y(x)=A\sin(kx),\qquad y(x)=\e^x,
\]
其中除自变量 $x$ 外，其余字母均为常数；幂函数一式先取非负整数 $n$。
\end{origexamplebox}

\Needspace{6\baselineskip}\begin{xesolution}
直接使用导数定义。例如幂函数有
\[
 \frac{\dd}{\dd x}(ax^n)
 =a\lim_{\Delta x\to0}
 \frac{(x+\Delta x)^n-x^n}{\Delta x}.
\]
由有限二项式展开
\[
 (x+\Delta x)^n-x^n
 =nx^{n-1}\Delta x+\binom n2x^{n-2}(\Delta x)^2+\cdots+(\Delta x)^n,
\]
除以 $\Delta x$ 后，第一项为 $nx^{n-1}$，其余各项都至少含一个 $\Delta x$，在
$\Delta x\to0$ 时趋于零，故得到 $anx^{n-1}$。对实数指数，在 $x>0$ 的区间可由广义二项式展开得到同一公式。同理可得
\[
 \boxed{C'=0},\qquad
 \boxed{(kx)'=k},\qquad
 \boxed{(ax^n)'=anx^{n-1}},
\]
\[
 \boxed{[A\sin(kx)]'=Ak\cos(kx)},\qquad
 \boxed{(\e^x)'=\e^x}.
\]

下面把三角函数和指数函数两式的极限过程完整写出。对 $y=A\sin(kx)$，由导数定义
\begin{align*}
 \frac{\dd y}{\dd x}
 &=A\lim_{\Delta x\to0}
 \frac{\sin[k(x+\Delta x)]-\sin(kx)}{\Delta x}\\
 &=A\lim_{\Delta x\to0}
 \frac{2\cos\!\left(kx+\dfrac{k\Delta x}{2}\right)
       \sin\!\left(\dfrac{k\Delta x}{2}\right)}{\Delta x}.
\end{align*}
第二步使用了恒等式
$\sin\alpha-\sin\beta
=2\cos[(\alpha+\beta)/2]\sin[(\alpha-\beta)/2]$。
把分式改写成
\[
 Ak\cos\!\left(kx+\frac{k\Delta x}{2}\right)
 \frac{\sin(k\Delta x/2)}{k\Delta x/2}.
\]
当 $\Delta x\to0$ 时，余弦因子趋于 $\cos(kx)$，而令
$u=k\Delta x/2$ 后，需要使用基本极限 $\sin u/u\to1$。这个极限可由单位圆直接证明：
当 $0<u<\pi/2$ 时，扇形面积与内、外两个三角形面积的比较给出
\[
 \sin u<u<\tan u.
\]
各项除以 $u$ 并整理，得到
\[
 \cos u<\frac{\sin u}{u}<1.
\]
当 $u\to0^+$ 时两侧都趋于 $1$，由夹逼定理可得右极限为 $1$；又因为
$\sin(-u)/(-u)=\sin u/u$，左极限相同。因此
\[
 \frac{\dd}{\dd x}[A\sin(kx)]=Ak\cos(kx).
\]

对 $y=\e^x$，先把公因子 $\e^x$ 提出：
\begin{align*}
 \frac{\dd}{\dd x}\e^x
 &=\lim_{\Delta x\to0}
 \frac{\e^{x+\Delta x}-\e^x}{\Delta x}\\
 &=\e^x\lim_{\Delta x\to0}
 \frac{\e^{\Delta x}-1}{\Delta x}.
\end{align*}
把自然指数函数写成它的收敛级数定义
\[
 \e^u=1+u+\frac{u^2}{2!}+\frac{u^3}{3!}+\cdots,
\]
便有
\[
 \frac{\e^u-1}{u}
 =1+\frac{u}{2!}+\frac{u^2}{3!}+\cdots.
\]
当 $|u|<1/2$ 时，余项的绝对值满足
\[
 \left|\frac{u}{2!}+\frac{u^2}{3!}+\cdots\right|
 \le |u|+|u|^2+\cdots=\frac{|u|}{1-|u|}\longrightarrow0.
\]
因此 $\lim_{u\to0}(\e^u-1)/u=1$。取 $u=\Delta x$，上式中的极限等于 $1$，最终得到
$\dd(\e^x)/\dd x=\e^x$。
\end{xesolution}

\begin{origexamplebox}{}\label{ch01:ex:s14-1-2}
函数 $f(x),g(x)$ 对 $x$ 的导数已知。根据导数定义证明以下求导运算的基本规律：
\[
 \frac{\dd}{\dd x}[a f(x)]=a\frac{\dd f}{\dd x},
\]
\[
 \frac{\dd}{\dd x}[f(x)g(x)]
 =f\frac{\dd g}{\dd x}+g\frac{\dd f}{\dd x},
\]
\[
 \frac{\dd}{\dd x}f(g(x))=\frac{\dd f}{\dd g}\frac{\dd g}{\dd x}.
\]
\end{origexamplebox}

\Needspace{6\baselineskip}\begin{xesolution}
对第一式，常数 $a$ 可直接移到极限号外：
\[
 \lim_{\Delta x\to0}\frac{af(x+\Delta x)-af(x)}{\Delta x}
 =a f'(x).
\]
对乘积，在分子中加上再减去 $f(x+\Delta x)g(x)$：
\begin{align*}
 &\frac{f(x+\Delta x)g(x+\Delta x)-f(x)g(x)}{\Delta x}\\
 &\quad=f(x+\Delta x)\frac{g(x+\Delta x)-g(x)}{\Delta x}
 +g(x)\frac{f(x+\Delta x)-f(x)}{\Delta x}.
\end{align*}
令 $\Delta x\to0$，由连续性 $f(x+\Delta x)\to f(x)$，便得乘积法则。

对复合函数令 $\Delta g=g(x+\Delta x)-g(x)$。当 $\Delta g\ne0$ 时
\[
 \frac{f(g+\Delta g)-f(g)}{\Delta x}
 =\frac{f(g+\Delta g)-f(g)}{\Delta g}\,
  \frac{\Delta g}{\Delta x}.
\]
不过 $\Delta g$ 可能在某些 $\Delta x\ne0$ 时恰为零，不能把上式中的除法当成总能进行。
为同时包含这个分支，利用 $f$ 在 $g(x)$ 处可导，把它的增量写为
\[
 f(g+\Delta g)-f(g)=f'(g)\Delta g+\varepsilon(\Delta g)\Delta g,
 \qquad \varepsilon(h)\to0\quad(h\to0).
\]
余项表示“实际增量减去切线估计”的差比 $\Delta g$ 更小；当 $\Delta g=0$ 时，
取 $\varepsilon(0)=0$，这条等式仍成立。只除以 $\Delta x\ne0$，得到
\[
 \frac{f(g+\Delta g)-f(g)}{\Delta x}
 =[f'(g)+\varepsilon(\Delta g)]\frac{\Delta g}{\Delta x}.
\]
$g$ 可导因而连续，所以 $\Delta x\to0$ 时 $\Delta g\to0$，
右侧两因子分别趋于 $f'(g(x))$ 和 $g'(x)$。这就严格得到链式法则
\[
 \boxed{\frac{\dd f(g(x))}{\dd x}=f'(g(x))g'(x)}.
\]
\end{xesolution}

\subsection{物理量的变化率}

导数在物理中最直接的应用，是求一个物理量对另一个物理量的变化率。这种计算广泛存在
于物理学的各个分支：速度是坐标随时间的变化率，加速度是速度随时间的变化率，保守力
也与势能随坐标的变化率直接相关。若能写出两个物理量之间的函数关系，直接求导便可
得到一个物理量改变时另一个物理量的变化率。

\begin{origexamplebox}{}\label{ch01:ex:s14-1-3}
一个一维运动质点的坐标随时间的关系分别为下列函数，求瞬时速度、加速度与时间的函数：
\begin{enumerate}[label=\arabic*.,leftmargin=2.2em]
 \item $x(t)=v_0t+\dfrac12at^2$；
 \item $x(t)=A\e^{-\beta t}\cos(\omega t)$，$\beta,\omega$ 为已知数；
 \item $x(t)=\dfrac{mv_0}{\beta}\left(1-\e^{-\frac{\beta}{m}t}\right)$，
       $\beta,m$ 为已知数。
\end{enumerate}
\end{origexamplebox}

\Needspace{6\baselineskip}\begin{xesolution}
逐式对 $t$ 求一阶、二阶导数。

\textbf{(1)}
\[
 \boxed{v(t)=v_0+at},\qquad \boxed{a(t)=a}.
\]

\textbf{(2)} 先使用乘积法则，再使用链式法则：
\[
 \boxed{v(t)=-A\e^{-\beta t}
   [\beta\cos(\omega t)+\omega\sin(\omega t)]},
\]
\[
 \boxed{a(t)=A\e^{-\beta t}
   [(\beta^2-\omega^2)\cos(\omega t)
    +2\beta\omega\sin(\omega t)]}.
\]

\textbf{(3)} 指数项中的 $-\beta/m$ 还要再乘一次：
\[
 \boxed{v(t)=v_0\e^{-\frac{\beta}{m}t}},\qquad
 \boxed{a(t)=-\frac{\beta v_0}{m}\e^{-\frac{\beta}{m}t}}.
\]
由前两式可得
\[
 a=\frac{\dd v}{\dd t}=-\frac{\beta}{m}v.
\]
在本题给出的运动中 $v>0$ 而 $a<0$，所以加速度方向与速度方向相反，速率随时间减小。
这种“与速度反向、使速率减小的加速度分量”常被称为减速度；它不是一种新的物理量，
仍然是加速度。若把这一运动解释为物体只受线性阻力，或其他恒力已经彼此平衡，则由
牛顿第二定律
\[
 F_{\rm net}=ma=-\beta v
\]
可知阻力为 $F_{\rm D}=-\beta v$：负号表示阻力与速度反向，大小 $\beta v$ 随速率增大。
若没有这一受力模型，仅由给定的 $x(t)$ 只能判断物体在减速，不能凭空断定具体是哪一种力造成的。
\end{xesolution}

\Needspace{0.38\textheight}
\begin{origexamplebox}{}\label{ch01:ex:s14-1-4}
一个半径为 $R$ 的半圆柱体在垂直于柱轴的水平方向上以匀速 $v_0$ 运动。如图，在圆柱面
上搁一竖直杆，此杆只能沿竖直方向上下运动。求杆上任意一点的速度、加速度与图示角
$\theta$ 的关系。
\probdiagOneFour
\end{origexamplebox}

\Needspace{6\baselineskip}\begin{xesolution}
设半圆柱圆心的横坐标为 $X$，杆的横坐标为 $x_{\rm r}$，接触点相对圆心的高度为 $y$。
由解图可见
\[
 x_{\rm r}=X+R\sin\theta,\qquad y=R\cos\theta.
\]
杆只能竖直运动，所以 $x_{\rm r}$ 为常数，而 $\dot X=v_0$。对第一式求导：
\[
 0=v_0+R\cos\theta\,\dot\theta,\qquad
 \dot\theta=-\frac{v_0}{R\cos\theta}.
\]
再对 $y$ 求导，得到杆的竖直速度
\[
 v_y=\dot y=-R\sin\theta\,\dot\theta
 =\boxed{v_0\tan\theta}.
\]
\soldiagOneFour
继续求导：
\[
 a_y=v_0\sec^2\theta\,\dot\theta
 =\boxed{-\frac{v_0^2}{R\cos^3\theta}}.
\]
若取竖直向上为正，速度向上，加速度向下。刚性杆作平动，所以杆上任意一点的速度和
加速度都相同。当 $\theta\to\pi/2$ 时理想模型给出发散结果，说明接触状态不可能一直
保持到圆柱侧缘。
\end{xesolution}

\begin{origexamplebox}{}\label{ch01:ex:s14-1-5}
一个发光点在焦距为 $f$ 的凸透镜主光轴上，以恒定不变的速度 $v$ 向着透镜运动。求其
像点移动速度与物距 $u$ 的关系。
\end{origexamplebox}

\Needspace{6\baselineskip}\begin{xesolution}
设物体在透镜左侧，$u>0$ 是它到透镜的距离；取像坐标 $q$ 向右为正，
所以右侧实像 $q>0$、左侧虚像 $q<0$。在近轴近似下，对焦距 $f>0$ 的薄凸透镜，
\[
 \frac1f=\frac1u+\frac1q,\qquad q=\frac{fu}{u-f}.
\]
光点向透镜运动，故 $\dot u=-v$。对 $q(u)$ 求导：
\[
 \frac{\dd q}{\dd u}=-\frac{f^2}{(u-f)^2},
\]
从而
\[
 \boxed{\dot q=\frac{f^2}{(u-f)^2}v}.
\]
当 $u>f$ 时，像位于透镜另一侧，$\dot q>0$ 表示实像远离透镜；物体越接近焦点，像点
移动越快。当 $u=f$ 时，出射光平行，不存在有限位置的像点，因此此时不能给像点一个有限速度。
当 $0<u<f$ 时，$q<0$，形成同侧虚像；虽然仍有 $\dot q>0$，
其含义却是负的像坐标向零增加，虚像向透镜靠近，而不是远离。
例如 $u=f/2$ 时 $q=-f$、$\dot q=4v$，这与有向坐标的判断一致。
不能把 $u=f$ 两侧的像当作以有限速度连续越过无穷远的同一个物点。
\end{xesolution}

\subsection{求极值}

对于光滑函数，函数曲线上各点的切线斜率就是该点的导数。当某一点处切线水平时，该点
可能是局部极大值、局部极小值，也可能既非极大又非极小。导数为零的点统称为\textbf{驻点}
（stationary point）。若在驻点处二阶导数为正，则是局部极小值；二阶导数为负，则是
局部极大值；二阶导数为零时，还需进一步判断。

\extremaDiagram

当感兴趣的物理量随其他因素变化时，先写出它们之间的函数关系，再令一阶导数为零，便
可寻找变化过程中的极值；必要时再借助二阶导数或端点条件判断极值类型。

\Needspace{0.38\textheight}
\begin{origexamplebox}{}\label{ch01:ex:s14-1-6}
利用费马原理证明光的折射定律：两个均匀介质的折射率分别为 $n_1,n_2$，分界面为一
平面。证明由平面外一点 $A$ 到达另一点 $B$ 的光线中，满足折射定律的光线的光程为
极值。
\probdiagOneSix
\end{origexamplebox}

\Needspace{6\baselineskip}\begin{xesolution}
设入射点 $P$ 到 $A$ 在界面上的垂足的距离为 $x$，两垂足间的水平距离为 $L$，而 $A,B$
到界面的垂直距离分别为 $h_1,h_2$。光程为
\[
 \mathcal L(x)=n_1\sqrt{h_1^2+x^2}
 +n_2\sqrt{h_2^2+(L-x)^2}.
\]
费马原理要求实际光路使光程取稳定值，因此
\[
 \frac{\dd\mathcal L}{\dd x}
 =n_1\frac{x}{\sqrt{h_1^2+x^2}}
 -n_2\frac{L-x}{\sqrt{h_2^2+(L-x)^2}}=0.
\]
由图中几何关系
\[
 \sin\alpha=\frac{x}{\sqrt{h_1^2+x^2}},\qquad
 \sin\beta=\frac{L-x}{\sqrt{h_2^2+(L-x)^2}},
\]
所以
\[
 \boxed{n_1\sin\alpha=n_2\sin\beta},
\]
这就是斯涅耳折射定律。此外
\[
 \frac{\dd^2\mathcal L}{\dd x^2}
 =\frac{n_1h_1^2}{(h_1^2+x^2)^{3/2}}
 +\frac{n_2h_2^2}{[h_2^2+(L-x)^2]^{3/2}}>0,
\]
故该稳定光路给出的确是光程的极小值。
\end{xesolution}

\Needspace{0.40\textheight}
\begin{origexamplebox}{}\label{ch01:ex:s14-1-7}
求理想气体在下面两个热力学过程中的温度最大值和最小值：
\begin{enumerate}[label=(\alph*),leftmargin=2.2em]
 \item 气体从状态 $(2p_0,V_0)$ 沿连接两点的直线段变化到状态 $(p_0,2V_0)$；
 \item $p$--$V$ 图中的一个圆（椭圆），过程满足
 \[
  \left(\frac p{p_0}-1\right)^2+
  \left(\frac V{V_0}-1\right)^2=\left(\frac12\right)^2.
 \]
\end{enumerate}
\probdiagOneSeven
\end{origexamplebox}

\Needspace{6\baselineskip}\begin{xesolution}
对一定量理想气体，$T=pV/(nR_g)$，因此只需在给定曲线上求乘积 $pV$ 的极值，其中
$R_g$ 为气体常量。

\textbf{(a)} 连接两状态的直线为
\[
 p=3p_0-\frac{p_0}{V_0}V.
\]
于是
\[
 pV=3p_0V-\frac{p_0}{V_0}V^2,\qquad
 \frac{\dd(pV)}{\dd V}=3p_0-2\frac{p_0}{V_0}V.
\]
在题设线段上 $V_0\le V\le2V_0$。驻点位于 $V=3V_0/2$，此时 $p=3p_0/2$，为抛物线的最大值；
最小值只能出现在两个端点，而两端点的温度相同，故
\[
 \boxed{T_{\max}=\frac{9p_0V_0}{4nR_g}},\qquad
 \boxed{T_{\min}=\frac{2p_0V_0}{nR_g}}.
\]

\textbf{(b)} 令 $P=p/p_0, X=V/V_0$，并参数化圆：
\[
 P=1+\frac12\cos\phi,\qquad X=1+\frac12\sin\phi.
\]
于是 $T=(p_0V_0/nR_g)PX$。求导可得
\[
 \frac{\dd(PX)}{\dd\phi}
 =\frac12(\cos\phi-\sin\phi)+\frac14\cos2\phi
 =(\cos\phi-\sin\phi)
 \left(\frac12+\frac14(\cos\phi+\sin\phi)\right).
\]
第二因子恒正，故驻点满足 $\sin\phi=\cos\phi$。比较两组对径点：
\[
 P=X=1\pm\frac{1}{2\sqrt2}.
\]
因此
\[
 \boxed{T_{\max}=\frac{p_0V_0}{nR_g}
  \left(1+\frac1{2\sqrt2}\right)^2},
\]
\[
 \boxed{T_{\min}=\frac{p_0V_0}{nR_g}
  \left(1-\frac1{2\sqrt2}\right)^2}.
\]
\end{xesolution}

\Needspace{0.30\textheight}
\begin{origexamplebox}{}\label{ch01:ex:s14-1-8}
宇宙中有两个天体，质量均为 $M$，相距 $2L$。求在它们连线的垂直平分线上何处的小
天体受到两个大天体的万有引力合力能够达到最大值。
\probdiagOneEight
\end{origexamplebox}

\Needspace{6\baselineskip}\begin{xesolution}
设小天体质量为 $m$，到连线中点的距离为 $x$。到每个大天体的距离为
$r=\sqrt{x^2+L^2}$。两个引力的垂直分量相消，沿垂直平分线方向的分量相加：
\[
 F(x)=2\frac{GMm}{r^2}\frac{x}{r}
 =\frac{2GMm\,x}{(x^2+L^2)^{3/2}}.
\]
求导：
\[
 F'(x)=2GMm\frac{L^2-2x^2}{(x^2+L^2)^{5/2}}.
\]
令 $F'(x)=0$，得 $x=L/\sqrt2$。由于 $F(0)=0$ 且 $x\to\infty$ 时 $F\to0$，这必为
最大值。因此在垂直平分线两侧各有一个对称位置：
\[
 \boxed{|x|=\frac{L}{\sqrt2}},
 \qquad
 \boxed{F_{\max}=\frac{4GMm}{3\sqrt3\,L^2}}.
\]
\end{xesolution}

\Needspace{0.43\textheight}
\begin{origexamplebox}{}\label{ch01:ex:s14-1-9}
\begin{enumerate}[label=(\alph*),leftmargin=2.2em]
 \item 从坐标原点以给定初速度 $v_0$ 沿不同角度抛出物体。证明能击中的目标与因初速度
 不足而无法击中的目标之间，分界线满足
 \[
  y=\frac{v_0^2}{2g}-\frac{g}{2v_0^2}x^2.
 \]
 \item 平面上有一个半径为 $R$ 的球体，希望从地面斜向上抛出物体并击中球体顶部。
 在命中球顶以前不接触球体的条件下，求抛出速度的下界，并判断最小值是否能取到；
 再讨论允许提前相切但不穿入球体时的最小抛出速度、抛出角及抛出点到球底部的距离。
\end{enumerate}
\end{origexamplebox}

\Needspace{6\baselineskip}\begin{xesolution}
\textbf{(a)} 设 $q=\tan\theta$，抛物线轨迹为
\[
 y=xq-\frac{gx^2}{2v_0^2}(1+q^2).
\]
对于给定横坐标 $x$，令右端对 $q$ 取最大值：
\[
 \frac{\partial y}{\partial q}=x-\frac{gx^2}{v_0^2}q=0,\qquad
 q=\frac{v_0^2}{gx}.
\]
代回即得所有同速抛物线的包络线
\[
 \boxed{y_{\max}(x)=\frac{v_0^2}{2g}-\frac{g}{2v_0^2}x^2}.
\]
包络线下方的目标至少对应一个发射角，线上只有一个临界角，线上方不能击中。

\textbf{(b)} 取球底部为原点，球面截线满足
\[
 x^2+(y-R)^2=R^2,
\]
目标是球顶 $(0,2R)$。令抛出点在 $(-d,0)$，以
\[
 X=\frac{x}{R},\quad Y=\frac{y}{R},\quad D=\frac dR,
 \quad Q=\tan\theta,
 \quad A=\frac{gR}{2v_0^2}(1+Q^2)
\]
作无量纲化，则轨迹可写为
\[
 Y=Q(X+D)-A(X+D)^2.
\]
它通过球顶 $(0,2)$，故
\begin{equation}\label{ch01:eq:sphere-top}
 QD-AD^2=2.
\end{equation}
先研究允许提前相切的边界轨迹，再检查它是否符合严格不接触的条件。设接触点
横坐标为 $s$，并记 $r=\sqrt{1-s^2}$，则上半圆在该点的纵坐标和斜率分别为
$1+r$ 与 $-s/r$。相切条件为
\begin{equation}\label{ch01:eq:sphere-tangent-point}
 Q(s+D)-A(s+D)^2=1+r,
\end{equation}
\begin{equation}\label{ch01:eq:sphere-tangent-slope}
 Q-2A(s+D)=-\frac{s}{r}.
\end{equation}
把临界抛物线改写成以相切点为展开中心的形式。因为它在
$X=s$ 处的函数值为 $1+r$、斜率为 $-s/r$，且二次项系数为 $-A$，所以
\[
 Y(X)=1+r-\frac{s}{r}(X-s)-A(X-s)^2.
\]
令 $X=0$，利用轨迹通过球顶 $Y=2$，得到
\[
 2=1+r+\frac{s^2}{r}-As^2.
\]
因此
\[
 A=\frac{r-1+s^2/r}{s^2}
 =\frac{1-r}{s^2r}
 =\frac{1}{r(1+r)},
\]
最后一步使用了 $s^2=1-r^2$。

接下来求发射方向。记
\[
 u=s+D>0,
\]
它是抛出点 $X=-D$ 到相切点 $X=s$ 的水平距离。把 $X=-D$、$Y=0$ 代入以相切点为中心的抛物线，得到
\[
 0=1+r+\frac{s}{r}u-Au^2.
\]
代入 $A=1/[r(1+r)]$ 并整理：
\[
 u^2-s(1+r)u-r(1+r)^2=0.
\]
取 $u>0$ 的根，
\[
 u=\frac{1+r}{2}\left(s+\sqrt{s^2+4r}\right).
\]
在发射点，轨迹斜率就是 $Q=\tan\theta$。由上面的相切点形式求导并令 $X=-D$，
\[
 Q=-\frac{s}{r}+2Au
 =\frac{\sqrt{s^2+4r}}{r}
 =\frac{\sqrt{1+4r-r^2}}{r}.
\]
这同时说明所取的是向右上方发射的 $Q>0$ 分支。于是
\[
 1+Q^2=\frac{1+4r}{r^2}.
\]
再代入 $A=gR(1+Q^2)/(2v_0^2)$，临界速度便成为只含 $r$ 的函数：
\[
 \frac{v_0^2}{gR}=\frac{1+Q^2}{2A}
 =\frac{(r+1)(4r+1)}{2r},\qquad 0<r\le1.
\]
对 $r$ 求导：
\[
 \frac{\dd}{\dd r}\frac{(r+1)(4r+1)}{2r}
 =\frac{(2r-1)(2r+1)}{2r^2}.
\]
最小值在 $r=1/2$ 处取得。此时
\[
 s=-\frac{\sqrt3}{2},\qquad A=\frac43,\qquad
 Q=\sqrt{11},\qquad
 D=\frac{\sqrt3(1+\sqrt{33})}{8}.
\]
因此
\[
 \boxed{v_*=3\sqrt{\frac{gR}{2}}},\qquad
 \boxed{\theta_0=\arctan\sqrt{11}},
\]
\[
 \boxed{d_0=\frac{\sqrt3(1+\sqrt{33})}{8}R}.
\]
\soldiagOneNine
在这组参数下，抛物线在球体左上方
$(x/R,y/R)=(-\sqrt3/2,3/2)$ 处与球面相切，随后到达球顶；因此解图中的“相切”发生在
球顶之前，而不是球顶本身。

还需证明没有另一类轨迹给出更小速度。将任意从左侧到达球顶的轨迹写成
$Y=2-qX-AX^2$。为避免在球顶左邻域穿入球面，必须有 $q\ge0$；
在球顶的水平、竖直速度分别为 $w_x>0$、$-qw_x$，而 $A=gR/(2w_x^2)$。
由竖直运动的速度关系，初速度平方为
\[
 \frac{v_0^2}{gR}=4+\frac{1+q^2}{2A}.
\]
令 $z=-X\in(0,1]$。要求轨迹不低于球的上半圆，等价于
\[
 q\ge h_A(z):=Az-\frac{z}{1+\sqrt{1-z^2}}.
\]
若 $A\le1/2$，可取 $q=0$，但速度平方不小于 $5gR$。
若 $A>1/2$，令 $r=\sqrt{1-z^2}$，直接求导得
\[
 h_A'(z)=A-\frac1{r(1+r)}.
\]
右端随 $z$ 严格减小，由正变负，因此 $h_A$ 恰有一个最大值，
最小可行 $q$ 正是在该处取等号，即前面已经求出的相切轨迹。
这就补足了对所有可行轨迹的比较，而不是只在相切轨迹内部求极值。

\textbf{严格不接触时只有下确界。}等号轨迹提前碰到球面，不能作为严格条件下的解。
固定 $A=4/3$，取 $q=1/\sqrt3+\varepsilon$，其中 $\varepsilon>0$。
在每个 $X<0$ 处，轨迹都比临界轨迹高出 $-\varepsilon X$，仍通过球顶；
令 $\varepsilon\to0^+$，初速度从上方趋于 $v_*$。所以
\[
 \boxed{\inf v_0=v_*,\quad\text{严格不提前接触时最小值不能取到。}}
\]
若允许提前相切而不穿入，才有 $v_{\min}=v_*$，相应角度和距离为上式的 $\theta_0,d_0$。
$\sqrt{gR}$ 具有速度量纲，$d_0$ 与 $R$ 成正比，符合几何相似性。
\end{xesolution}

\subsection{平衡位置及其稳定性}

对一维保守系统，力与势能满足 $F=-\dd U/\dd x$，所以平衡位置首先要求 $U'(x)=0$。
局部极小点受到小扰动后会出现指向平衡点的回复力，因而是稳定平衡；局部极大点则是不稳定平衡。
必须注意：没有阻尼时，物体从局部极小点附近释放后通常会在其两侧振动，并不会自行停在最低点；
只有存在摩擦等耗散时，机械能逐渐减少，系统才可能最终趋近较低势能的稳定平衡。过去常由受力分析
判断平衡；利用势能曲线也能得到相同答案，而且系统越复杂，能量方法往往越简便。

\potentialDiagram

右侧势能曲线中，$S_1,S_2$ 是局部极小点，对应稳定平衡；$U$ 是局部极大点，对应不稳定平衡。
图线其余斜率不为零的位置受力不为零，不能作为静止平衡位置。

以单摆为例，长度为 $l$ 的轻杆一端与 $O$ 点光滑连接，另一端连接质量为 $m$ 的质点。
若以摆竖直向下为 $\theta=0$，系统的重力势能为
\begin{equation}
 E_p=-mgl\cos\theta.
\end{equation}
势能取稳定值的条件为
\begin{equation}
 \frac{\dd E_p}{\dd\theta}=mgl\sin\theta=0,
\end{equation}
故 $\theta=0,\pi$ 都是平衡位置。再求二阶导数
\begin{equation}
 \frac{\dd^2E_p}{\dd\theta^2}=mgl\cos\theta.
\end{equation}
在 $\theta=0$ 处二阶导数为正，是稳定平衡；在 $\theta=\pi$ 处二阶导数为负，是不稳定
平衡。更一般地，势能曲线的局部极小点是稳定平衡，局部极大点是不稳定平衡；导数不为
零的点不能平衡，二阶导数为零的平衡位置须作进一步判断。

\Needspace{0.37\textheight}
\begin{origexamplebox}{}\label{ch01:ex:s14-1-10}
如图，在固定点 $A$ 上系有一根质量可忽略、原长为 $h$、劲度系数为 $k$ 的弹簧，另一
端连接质量为 $m$ 的重物。重物被限制在竖直光滑直线上，$A$ 点到约束直线的距离恰为
$h$。求平衡时重物到 $B$ 点的距离 $x$。
\probdiagOneTen
\end{origexamplebox}

\Needspace{6\baselineskip}\begin{xesolution}
当重物下降 $x$ 时，弹簧长度
\[
 \ell=\sqrt{h^2+x^2},
\]
伸长量为 $\ell-h$。取 $B$ 点处的重力势能为零，系统势能为
\[
 U(x)=\frac12k(\sqrt{h^2+x^2}-h)^2-mgx.
\]
求导：
\[
 U'(x)=k(\ell-h)\frac{x}{\ell}-mg.
\]
平衡条件为
\[
 \boxed{kx\left(1-\frac{h}{\sqrt{h^2+x^2}}\right)=mg}.
\]
这已经给出 $x$ 的确定方程。若希望写成无根式的多项式，令
$\lambda=mg/(kh)$、$z=x/h$，则
\[
 z\left(1-\frac1{\sqrt{1+z^2}}\right)=\lambda,
\]
等价于
\[
 \boxed{z^4-2\lambda z^3+\lambda^2z^2-2\lambda z+\lambda^2=0},
\]
并须取满足 $z>\lambda$ 的正根，以排除平方带来的伪根。也可以看出
\[
 U''(x)=k\left(1-\frac{h^3}{(h^2+x^2)^{3/2}}\right)>0\quad(x>0),
\]
所以该平衡位置唯一且稳定。一般参数下没有更简洁的初等显式根式，保留隐式方程
更能体现其物理结构。
\end{xesolution}

\Needspace{0.36\textheight}
\begin{origexamplebox}{}\label{ch01:ex:s14-1-11}
如图，锥角为 $\alpha$ 的光滑圆锥体上放着一段质量为 $m$、等效劲度系数为 $k$、原长为
$L_0$ 的弹簧，弹簧始终保持水平。求平衡时弹簧到锥体顶端的距离 $x$。
\probdiagOneEleven
\end{origexamplebox}

\Needspace{6\baselineskip}\begin{xesolution}
题图中的 $\alpha$ 是圆锥母线与轴线的夹角。距顶端竖直距离为 $x$ 的水平截面半径为
$r=x\tan\alpha$，所以环形弹簧长度
\[
 \ell=2\pi r=2\pi x\tan\alpha.
\]
取顶点高度处为重力势能零点，系统势能为
\[
 U(x)=-mgx+\frac12k(2\pi x\tan\alpha-L_0)^2.
\]
令 $U'(x)=0$：
\[
 -mg+k(2\pi x\tan\alpha-L_0)(2\pi\tan\alpha)=0.
\]
因此
\[
 \boxed{x=\frac{L_0}{2\pi\tan\alpha}
 +\frac{mg}{4\pi^2k\tan^2\alpha}}.
\]
又因
$U''(x)=k(2\pi\tan\alpha)^2>0$，该平衡位置稳定。这里的 $k$ 是整段环形弹簧在周长方向的等效劲度系数。
\end{xesolution}

\Needspace{0.37\textheight}
\begin{origexamplebox}{}\label{ch01:ex:s14-1-12}
两根质量可忽略、长度为 $l$ 的杆在顶部无摩擦连接，质量为 $m$ 的重物 $A$ 放在连接
处。两杆底部放在光滑水平地面上，底端之间由原长为零、劲度系数为 $k$ 的弹簧连接。
求平衡时两杆与地面的夹角 $\theta$。
\probdiagOneTwelve
\end{origexamplebox}

\Needspace{6\baselineskip}\begin{xesolution}
由对称性，两杆与地面的夹角相同。重物高度为 $l\sin\theta$，两杆底端距离即弹簧长度为
$2l\cos\theta$。系统势能
\[
 U(\theta)=mgl\sin\theta+\frac12k(2l\cos\theta)^2.
\]
求导：
\[
 U'(\theta)=mgl\cos\theta-4kl^2\sin\theta\cos\theta
 =l\cos\theta(mg-4kl\sin\theta).
\]
非退化平衡位置满足
\[
 \boxed{\sin\theta=\frac{mg}{4kl}}.
\]
作为 $0<\theta<\pi/2$ 的内点解，存在条件是 $0<mg<4kl$。在该点
\[
 U''(\theta)=-4kl^2\cos^2\theta<0,
\]
故这一内点平衡是\textbf{不稳定平衡}：角度稍增或稍减，系统都会继续偏离。两个几何
边界 $\theta=0$ 与 $\theta=90^\circ$ 中哪一个是较低势能状态，取决于 $2kl^2$ 与
$mgl$ 的大小；实际能否到达边界，还须结合杆、铰链和重物的接触约束判断。
特别地，$mg=4kl$ 时该解到达 $\theta=\pi/2$，上面的二阶导数为零，不能继续用
$U''<0$ 判断。令 $\delta=\pi/2-\theta$，展开得
$U=4kl^2-\tfrac12kl^2\delta^4+O(\delta^6)$。
若允许向 $\delta>0$ 偏离，这仍是退化的单侧极大值，而不是稳定内点。
\end{xesolution}

\subsection{做展开}

在光滑函数图像上，曲线在某点附近与该点的切线很接近。设函数 $f(x)$ 在 $x_0$ 处的
取值为 $f(x_0)$，切线斜率为 $f'(x_0)$，切线方程为
\begin{equation}
 y=f(x_0)+f'(x_0)(x-x_0).
\end{equation}
当 $x$ 与 $x_0$ 足够接近时，可以近似写为
\begin{equation}
 f(x)\simeq f(x_0)+f'(x_0)(x-x_0).
\end{equation}
若令 $\Delta f=f(x)-f(x_0)$、$\Delta x=x-x_0$，则
\begin{equation}
 \Delta f\simeq f'(x_0)\Delta x.
\end{equation}
这说明连续函数在某点附近的函数值变化量，近似正比于自变量的变化量，比例系数就是
函数在该点的导数。

下面令 $\varepsilon$ 表示无量纲小量，且 $|\varepsilon|\ll1$。常见的小量展开有
\[
 \sin\varepsilon\simeq\tan\varepsilon\simeq\varepsilon,
\]
\[
 (1+\varepsilon)^n\simeq1+n\varepsilon,\qquad
 \e^{\varepsilon}\simeq1+\varepsilon,\qquad
 \ln(1+\varepsilon)\simeq\varepsilon.
\]
它们在 $\varepsilon\to0$ 时与原函数越来越接近。若实际变化量 $\Delta x$ 具有长度量纲，
应先用问题中的特征长度 $L$ 构造 $\varepsilon=\Delta x/L$，不能直接把有量纲的 $\Delta x$ 与数字 $1$ 比较。
小量展开在研究物理系统某点附近的运动行为时非常有效。

\begin{origexamplebox}{}\label{ch01:ex:s14-1-13}
相对论中运动物体的能量随速度的关系不是经典力学中的动能 $\frac12mv^2$，而是
\[
 E(v)=\frac{mc^2}{\sqrt{1-v^2/c^2}}.
\]
证明当物体速度远小于光速时，相对论性能量与经典力学动能给出的结果一致。
\end{origexamplebox}

\Needspace{6\baselineskip}\begin{xesolution}
令 $\varepsilon=v^2/c^2\ll1$，使用二项式展开
\[
 (1-\varepsilon)^{-1/2}
 =1+\frac12\varepsilon+\frac38\varepsilon^2+\cdots.
\]
因此
\[
 E=mc^2+\frac12mv^2+\frac38m\frac{v^4}{c^2}+\cdots.
\]
第一项 $mc^2$ 是与速度无关的静能，不影响机械运动中的能量差；扣除静能后，相对论动能
为
\[
 E_k=E-mc^2=\frac12mv^2+\mathcal O\!\left(\frac{v^4}{c^2}\right).
\]
所以当 $v\ll c$ 时
\[
 \boxed{E_k\simeq\frac12mv^2},
\]
与经典力学一致。严格地说，题设所给 $E(v)$ 是总能量，不是单独的动能；这正是为何须
扣除 $mc^2$。
\end{xesolution}

\begin{origexamplebox}{}\label{ch01:ex:s14-1-14}
相对论中，质量为 $m$ 的质点在恒力 $F$ 作用下由静止出发做直线运动，位移和时间的关系为
\[
 x(t)=\frac{mc^2}{F}
 \left(\sqrt{1+\left(\frac{Ft}{mc}\right)^2}-1\right).
\]
证明在 $\dfrac{Ft}{mc}\ll1$ 时，相对论公式与经典力学一致。
\end{origexamplebox}

\Needspace{6\baselineskip}\begin{xesolution}
令 $\varepsilon=Ft/(mc)\ll1$。由
$\sqrt{1+\varepsilon^2}=1+\varepsilon^2/2+\order{\varepsilon^4}$，有
\begin{align*}
 x(t)
 &\simeq\frac{mc^2}{F}\cdot\frac12
 \left(\frac{Ft}{mc}\right)^2\\
 &=\frac12\frac{F}{m}t^2.
\end{align*}
经典力学中恒力产生的加速度为 $a=F/m$，故
\[
 \boxed{x(t)\simeq\frac12at^2}.
\]
下一阶修正为 $-F^3t^4/(8m^3c^2)$，说明相对论位移略小于经典结果。
\end{xesolution}

当一个物理量随另一个物理量改变时，在自变量产生微小变化的情况下，函数值的微小变化
也由给定点处的导数决定。对函数 $y(x)$，当 $x$ 从 $x_0$ 变到 $x_0+\Delta x$ 时，
\begin{equation}
 \Delta y\simeq\frac{\dd y}{\dd x}(x_0)\Delta x.
\end{equation}

\incrementDiagram

上式右边含有两个量：自变量的取值 $x_0$ 与其改变量 $\Delta x$。它们共同给出 $x_0$
附近的函数值变化量，且 $\Delta y$ 正比于该点的导数和 $\Delta x$。

\begin{origexamplebox}{}\label{ch01:ex:s14-1-15}
已知抛体运动的射程
\[
 S=\frac{v^2}{g}\sin2\theta,
\]
其中 $v$ 为初速度，$\theta$ 为抛出角。求速度或角度有微小变化时，射程如何变化。
\end{origexamplebox}

\Needspace{6\baselineskip}\begin{xesolution}
把 $S$ 看成 $v,\theta$ 的二元函数，其一阶微小变化为全微分
\[
 \Delta S\simeq
 \frac{\partial S}{\partial v}\Delta v
 +\frac{\partial S}{\partial\theta}\Delta\theta.
\]
分别求偏导：
\[
 \frac{\partial S}{\partial v}=\frac{2v}{g}\sin2\theta,\qquad
 \frac{\partial S}{\partial\theta}=\frac{2v^2}{g}\cos2\theta.
\]
所以当两者都可能微变时
\[
 \boxed{\Delta S\simeq
 \frac{2v}{g}\sin2\theta\,\Delta v
 +\frac{2v^2}{g}\cos2\theta\,\Delta\theta}.
\]
若只改变速度，取 $\Delta\theta=0$；若只改变角度，取 $\Delta v=0$，其中角度增量须以
弧度计。特别地，
$\theta=45^\circ$ 时角度的一阶影响为零，这正对应最大射程；此时角度误差的主要影响是
二阶量，而速度误差仍产生一阶影响。
\end{xesolution}

\Needspace{0.40\textheight}
\begin{origexamplebox}{}\label{ch01:ex:s14-1-16}
如图，在 $A$ 点有一个粒子发射源，它在上方 $180^\circ$ 的范围内独立、均匀随机地发射出 $N$ 个
粒子。粒子被与 $A$ 相距 $h$ 的壁接收。建立图示坐标系，求在 $x$ 到 $x+\Delta x$
范围内接收到的粒子数的期望值 $\mathbb E[\Delta n]$，其中 $\Delta x>0$。
假设该小区间足够窄，使方向角在其中可作一阶近似；
例如在 $x=0$ 附近要求 $|\Delta x|\ll h$。
\probdiagOneSixteen
\end{origexamplebox}

\Needspace{6\baselineskip}\begin{xesolution}
“随机”本身并不意味着均匀；本题另行给出的均匀条件表示等角宽的区间具有相同概率。
长度为 $\pi$ 的角区间中，角宽 $\Delta\theta$ 对应概率 $\Delta\theta/\pi$，
故接收粒子数的期望为
\[
 \mathbb E[\Delta n]=\frac{N}{\pi}\Delta\theta.
\]
由解图的直角三角形
\[
 \tan\theta=\frac{x}{h},\qquad
 \theta=\arctan\frac{x}{h}.
\]
两条边界射线对应的精确角宽，是两个端点方向角之差，因此
\[
 \mathbb E[\Delta n]=\frac N\pi
 \left[\arctan\frac{x+\Delta x}{h}-\arctan\frac{x}{h}\right].
\]
因此
\[
 \frac{\dd\theta}{\dd x}
 =\frac{h}{h^2+x^2},\qquad
 \Delta\theta\simeq\frac{h}{h^2+x^2}\Delta x.
\]
这里的小区间条件不是 $\Delta x\ll x$，而是角度变化率在区间内近似不变。一个简明的充分条件是
$|\Delta x|\ll\sqrt{h^2+x^2}$；在 $x=0$ 处它退化为 $|\Delta x|\ll h$，因此公式仍然完全适用。
于是
\[
 \boxed{\mathbb E[\Delta n]\simeq
 \frac{Nh}{\pi(h^2+x^2)}\Delta x}.
\]
\soldiagOneSixteen
该结果显示墙上粒子分布不均匀：离 $O$ 点越远，单位长度接收的粒子越少。若把墙的两侧
全部计入，积分平均粒子数密度 $\rho(x)=Nh/[\pi(h^2+x^2)]$ 从 $-\infty$ 到 $+\infty$ 恰好得到
总数 $N$，也为系数提供了检验。
一次实验的实际粒子数为整数，会围绕期望值波动；期望值本身不必是整数。
\end{xesolution}
}

\section{泰勒展开与近似的阶}

函数在 $x_0$ 附近足够光滑时，
\begin{equation}
 f(x_0+h)=f(x_0)+f'(x_0)h+\frac{f''(x_0)}{2!}h^2+\cdots.
\end{equation}
写“$h$ 很小”还不够，必须说明被舍项相对保留项小到什么阶。常用展开包括
\begin{align}
 \sin x&=x-\frac{x^3}{6}+\order{x^5},&
 \cos x&=1-\frac{x^2}{2}+\order{x^4},\\
 (1+x)^\alpha&=1+\alpha x+\frac{\alpha(\alpha-1)}2x^2+\order{x^3},&
 \ln(1+x)&=x-\frac{x^2}{2}+\order{x^3}.
\end{align}
展开参数必须无量纲；不能直接声称“长度 $x$ 很小”，而应比较 $x/L$。

\begin{origexamplebox}{单摆非线性修正的量级}
单摆精确方程为 $\ddot\theta+(g/\ell)\sin\theta=0$。若最大摆角为 $\theta_0\ll1$，估计线性化所忽略项相对于恢复项的量级。
\end{origexamplebox}
\Needspace{6\baselineskip}\begin{xesolution}
展开 $\sin\theta=\theta-\theta^3/6+\cdots$。线性恢复项是 $\theta$，首个被忽略项为 $-\theta^3/6$，二者比值的最大量级为
\[
 \left|\frac{-\theta^3/6}{\theta}\right|\lesssim\frac{\theta_0^2}{6}.
\]
因此小角近似的相对方程误差为 $\order{\theta_0^2}$。角度必须用弧度，弧度是无量纲量；若用度数直接代入展开，数值和阶数判断都会失去意义。
\end{xesolution}


{
\graphicspath{{/Users/dongxingxi/Documents/物理竞赛一轮资料/sample/xes-original/fifteenth-section/}}
\section{积分及其应用}
% TikZ figures redrawn from the source; annotations used only in solutions are
% drawn in the restrained blue accent of the common style.

\newcommand{\partitionDiagram}{%
\begin{center}
\begin{tikzpicture}[x=0.66cm,y=0.58cm,>=Latex]
  \begin{scope}
    \draw[axis] (0,0)--(7.0,0) node[right] {$t$};
    \draw[axis] (0,0)--(0,4.2) node[above] {$x$};
    \draw[curve,domain=0.45:6.2,smooth,variable=\t]
      plot ({\t},{0.55+0.075*\t*\t});
    \foreach \i in {1,...,6}{
      \pgfmathsetmacro{\xa}{0.65+0.82*(\i-1)}
      \pgfmathsetmacro{\xb}{0.65+0.82*\i}
      \pgfmathsetmacro{\ya}{0.55+0.075*\xa*\xa}
      \pgfmathsetmacro{\slope}{0.15*\xa}
      \draw[TextGray] (\xa,\ya)--(\xb,{\ya+\slope*(\xb-\xa)});
      \draw[guide] (\xa,0)--(\xa,\ya);
    }
    \node[below] at (3.1,-0.12) {$\Delta t\quad\Delta t\quad\cdots$};
  \end{scope}
  \begin{scope}[xshift=5.9cm]
    \draw[axis] (0,0)--(7.0,0) node[right] {$t$};
    \draw[axis] (0,0)--(0,4.2) node[above] {$v$};
    \draw[curve,domain=0.45:6.2,smooth,variable=\t]
      plot ({\t},{0.15*\t});
    \foreach \i in {1,...,6}{
      \pgfmathsetmacro{\xa}{0.65+0.82*(\i-1)}
      \pgfmathsetmacro{\xb}{0.65+0.82*\i}
      \pgfmathsetmacro{\ya}{0.15*\xa}
      \draw[fill=DeepBlue!16,draw=TextGray] (\xa,0) rectangle (\xb,\ya);
    }
    \node[below] at (3.1,-0.12) {$\Delta t\quad\Delta t\quad\cdots$};
  \end{scope}
\end{tikzpicture}

\figcap{同一运动的 $x(t)$ 与 $v(t)=x'(t)$：左图的切线增量等于右图对应矩形的面积，均为 $v(t_i)\Delta t$。}\label{ch01:fig:s15-1-5}
\end{center}}

\newcommand{\integralDiagram}{%
\begin{center}
\begin{tikzpicture}[x=0.9cm,y=0.7cm,>=Latex]
  \draw[axis] (0,0)--(7.2,0) node[right] {$x$};
  \draw[axis] (0,0)--(0,4.1) node[above] {$y$};
  \pgfmathdeclarefunction{areaheight}{1}{\pgfmathparse{3.1-.23*(##1-3.2)^2}}
  \fill[DeepBlue!14] (1.1,0)
    -- plot[domain=1.1:5.45,samples=90] (\x,{areaheight(\x)})
    --(5.45,0)--cycle;
  \draw[curve,domain=.35:6.2,samples=100] plot (\x,{areaheight(\x)});
  \draw[thick] (1.1,0)--(1.1,{areaheight(1.1)})
    (5.45,0)--(5.45,{areaheight(5.45)});
  \node[below] at (1.1,0) {$a$}; \node[below] at (5.45,0) {$b$};
  \node[above] at (3.2,3.1) {$f(x)$};
  \node[align=center,anchor=north] at (3.45,-.65)
    {$\displaystyle \frac{\dd F}{\dd x}=f(x)$\\[2pt]
     $\displaystyle \int_a^b f(x)\dd x=F(b)-F(a)$};
\end{tikzpicture}

\figcap{积分的代数和几何意义}\label{ch01:fig:s15-1-6}
\end{center}}

\newcommand{\areaLineDiagram}{%
\begin{center}
\begin{tikzpicture}[x=1cm,y=0.72cm,>=Latex]
  \draw[axis] (0,0)--(4.2,0) node[right] {$x$};
  \draw[axis] (0,0)--(0,3.2) node[above] {$y$};
  \fill[DeepBlue!16] (0,0)--(3.1,0)--(3.1,2.35)--cycle;
  \draw[curve] (0,0)--(3.6,2.73) node[above left] {$y=kx$};
  \draw[thick] (3.1,0)--(3.1,2.35);
  \node[below] at (3.1,0) {$a$};
\end{tikzpicture}
\end{center}}

\newcommand{\areaParabolaDiagram}{%
\begin{center}
\begin{tikzpicture}[x=1cm,y=0.65cm,>=Latex]
  \draw[axis] (0,0)--(4.2,0) node[right] {$x$};
  \draw[axis] (0,0)--(0,3.2) node[above] {$y$};
  \fill[DeepBlue!16] (0,0) plot[domain=0:3.1,smooth,variable=\x]
    ({\x},{0.27*\x*\x}) --(3.1,0)--cycle;
  \draw[curve,domain=0:3.45,smooth,variable=\x]
    plot ({\x},{0.27*\x*\x}) node[above left] {$y=kx^2$};
  \draw[thick] (3.1,0)--(3.1,{0.27*3.1*3.1});
  \node[below] at (3.1,0) {$a$};
\end{tikzpicture}
\end{center}}

\newcommand{\circlePartitionsDiagram}{%
\begin{center}
\begin{tikzpicture}[scale=0.78,>=Latex]
  \begin{scope}
    \draw[thick] (0,0) circle (1.55);
    \foreach \a in {0,10,20,30}{\draw[thick] (0,0)--(\a:1.55);}
    \coordinate (Oa) at (0,0);
    \coordinate (A20) at (20:1);
    \coordinate (A30) at (30:1);
    \draw pic[draw,->,angle radius=10mm] {angle=A20--Oa--A30};
    \node (anglelabel) at (1.15,2.12) {$\dd\theta$};
    \draw[TextGray,thin] (anglelabel.south)--(.91,.43);
    \fill (0,0) circle (1.4pt) node[below left] {$O$};
  \end{scope}
  \begin{scope}[xshift=4.5cm]
    \draw[thick] (0,0) circle (1.55);
    \draw[DeepBlue,thick] (0,0) circle (0.95);
    \draw[DeepBlue,thick] (0,0) circle (1.05);
    \draw[thick] (0,0)--(25:1.0) node[midway,above] {$r$};
    \node[DeepBlue,anchor=west] (drlabel) at (1.85,1.20) {$\dd r$};
    \draw[DeepBlue,thin] (drlabel.west)--(.94,.34);
    \fill (0,0) circle (1.4pt) node[below left] {$O$};
  \end{scope}
  \begin{scope}[xshift=9.0cm]
    \draw[thick] (0,0) circle (1.55);
    \draw[DeepBlue,thick] (0.8,-1.33)--(0.8,1.33);
    \draw[DeepBlue,thick] (0.92,-1.25)--(0.92,1.25);
    \draw[thick] (0,0)--(0.8,0) node[midway,below] {$x$};
    \node[DeepBlue,anchor=west] (dxlabel) at (2.0,.65) {$\dd x$};
    \draw[DeepBlue,thin] (dxlabel.west)--(.86,.65);
    \fill (0,0) circle (1.4pt) node[below left] {$O$};
  \end{scope}
\end{tikzpicture}
\end{center}}

\newcommand{\massElementsDiagram}{%
\begin{center}
\begin{tikzpicture}[scale=0.88,>=Latex]
  \begin{scope}
    \draw[axis] (0,0)--(0,3.5) node[above] {$x$};
    \draw[curve] plot[smooth] coordinates {(0,0) (1.15,0.5) (1.55,1.5) (1.0,2.55) (0,3.15)};
    \draw[curve] plot[smooth] coordinates {(0,0) (-1.15,0.5) (-1.55,1.5) (-1.0,2.55) (0,3.15)};
    \fill[DeepBlue!24] (-1.5,1.4) rectangle (1.5,1.55);
    \draw[DeepBlue,thick] (-1.5,1.4)--(1.5,1.4) (-1.5,1.55)--(1.5,1.55);
    \draw[<->] (0,1.18)--(1.5,1.18);
    \node[below=1pt] at (0.75,1.18) {$r(x)$};
    \node[DeepBlue,anchor=west] (dxlabel) at (2.0,1.95) {$\dd x$};
    \draw[DeepBlue,thin] (dxlabel.west)--(1.5,1.475);
    \node[below] at (0,-0.45) {$\dd m=\rho(x)\pi r^2(x)\dd x$};
  \end{scope}
  \begin{scope}[xshift=6.0cm]
    \draw[thick] (0,0) circle (1.55);
    \draw[DeepBlue,thick] (0,0) circle (1.0);
    \draw[DeepBlue,thick] (0,0) circle (1.1);
    \draw[<->] (0,0)--(0.72,0.72);
    \node[below right=1pt] at (0.36,0.36) {$r$};
    \node[DeepBlue,anchor=west] (drlabel) at (1.90,.70) {$\dd r$};
    \draw[DeepBlue,thin] (drlabel.west)--(.96,.43);
    \node[below] at (0,-2.0) {$\dd m=4\pi r^2\rho(r)\dd r$};
  \end{scope}
\end{tikzpicture}
\end{center}}

\newcommand{\pvWorkDiagram}{%
\begin{center}
\begin{tikzpicture}[x=0.9cm,y=0.68cm,>=Latex]
  \begin{scope}
    \draw[axis] (0,0)--(4.4,0) node[right] {$V$};
    \draw[axis] (0,0)--(0,3.2) node[above] {$p$};
    \fill[DeepBlue!16] (0.8,0)
      -- plot[domain=0.8:3.6,smooth,variable=\x] ({\x},{2.16/\x})--(3.6,0)--cycle;
    \draw[curve,domain=0.65:4.0,smooth,variable=\x]
      plot ({\x},{2.16/\x});
    \node at (2.35,3.65) {$p=nR_gT/V$};
    \draw[thick] (0.8,0)--(0.8,2.7) (3.6,0)--(3.6,0.6);
    \node[below] at (0.8,0) {$V_1$}; \node[below] at (3.6,0) {$V_2$};
  \end{scope}
  \begin{scope}[xshift=6.2cm]
    \draw[axis] (0,0)--(4.4,0) node[right] {$V$};
    \draw[axis] (0,0)--(0,3.2) node[above] {$p$};
    \fill[DeepBlue!16] (0.8,0)
      -- plot[domain=0.8:3.6,smooth,variable=\x] ({\x},{1.95/(\x^1.35)})--(3.6,0)--cycle;
    \draw[curve,domain=0.65:4.0,smooth,variable=\x]
      plot ({\x},{1.95/(\x^1.35)});
    \node at (2.35,3.65) {$pV^\gamma=C$};
    \draw[thick] (0.8,0)--(0.8,{1.95/(.8^1.35)})
      (3.6,0)--(3.6,{1.95/(3.6^1.35)});
    \node[below] at (0.8,0) {$V_1$}; \node[below] at (3.6,0) {$V_2$};
  \end{scope}
\end{tikzpicture}
\end{center}}

\newcommand{\radialFallDiagram}{%
\begin{center}
\begin{tikzpicture}[scale=0.9,>=Latex]
  \fill (0,0) circle (4pt) node[above left=2pt] {$A(M)$};
  \draw[axis] (0,0)--(5.35,0) node[right] {径向坐标};
  \draw[DeepBlue,dashed] (4.75,0) circle (3.0pt);
  \node[above=8pt,DeepBlue] at (4.75,0) {$m(0)$};
  \fill[DeepBlue] (3.25,0) circle (2.5pt);
  \node[below left=3pt,DeepBlue] at (3.25,0) {$m(t)$};
  \draw[vector] (3.1,0.68)--(2.05,0.68) node[midway,above] {$\vec v$};
  \draw[<->] (0,-0.72)--(3.25,-0.72);
  \node[below=1pt] at (1.63,-0.72) {$r$};
  \draw[<->,TextGray] (0,-1.28)--(4.75,-1.28);
  \node[below=1pt,TextGray] at (2.38,-1.28) {$r_0$};
  \draw[guide] (3.25,-0.9)--(3.25,0);
  \draw[guide] (4.75,-1.44)--(4.75,0);
\end{tikzpicture}
\end{center}}

\newcommand{\arcLengthDiagram}{%
\begin{center}
\begin{tikzpicture}[x=1cm,y=0.72cm,>=Latex]
  \draw[axis] (0,0)--(5.8,0) node[right] {$x$};
  \draw[axis] (0,0)--(0,3.5) node[above] {$y$};
  \draw[curve,domain=.55:4.6,samples=80] plot (\x,{1.4+.27*(\x-2.3)^2});
  \draw[thick] (1.2,0)--(1.2,{1.4+.27*(1.2-2.3)^2})
    (4.35,0)--(4.35,{1.4+.27*(4.35-2.3)^2});
  \node[below] at (1.2,0) {$x_1$}; \node[below] at (4.35,0) {$x_2$};
  \coordinate (P) at (7.0,1.2);
  \coordinate (H) at (9.0,1.2);
  \coordinate (Q) at (9.0,{1.2+2*.54*(3.7-2.3)});
  \draw[guide] (3.7,{1.4+.27*(3.7-2.3)^2})--(6.7,2.9);
  \node[TextGray,above] at (8,3.2) {局部切线三角形};
  \draw[DeepBlue,very thick] (P)--(Q) node[midway,above left] {$\dd s$};
  \draw[component] (P)--(H) node[midway,below] {$\dd x$};
  \draw[component] (H)--(Q) node[midway,right] {$\dd y$};
  \draw[TextGray,line width=.45pt] ($(H)+(-.13,0)$)--($(H)+(-.13,.13)$)--($(H)+(0,.13)$);
  \node[above left] at (4.55,2.78) {$y(x)$};
\end{tikzpicture}
\end{center}}

\newcommand{\trackTimeDiagram}{%
\begin{center}
\begin{tikzpicture}[x=0.95cm,y=0.72cm,>=Latex]
  \draw[axis] (0,0)--(6.4,0) node[right] {$x$};
  \draw[axis] (0,0)--(0,4.2) node[above] {$y$};
  \draw[curve] plot[smooth] coordinates {(0,3.6) (0.55,2.7) (1.5,2.0) (3.0,1.15) (4.8,0.35) (5.8,0)};
  \fill[DeepBlue] (1.7,1.88) circle (2.3pt);
  \draw[vector] (1.7,1.88)--++(0.75,-0.45) node[above right] {$\vec v$};
  \draw[guide] (1.7,1.88)--(1.7,3.6);
  \draw[<->] (2.05,1.88)--node[right] {$H-y$}(2.05,3.6);
  \node[left] at (0,3.6) {$A(0,H)$};
  \node[below] at (5.8,0) {$B(L,0)$};
  \node[above right] at (3.05,1.15) {$y(x)$};
\end{tikzpicture}
\end{center}}

\newcommand{\trackTimeSolutionDiagram}{%
\begin{center}
\begin{tikzpicture}[x=0.95cm,y=0.72cm,>=Latex]
  \draw[axis] (0,0)--(6.4,0) node[right] {$x$};
  \draw[axis] (0,0)--(0,4.2) node[above] {$y$};
  \draw[curve] plot[smooth] coordinates
    {(0,3.6) (0.55,2.7) (1.5,2.0) (3.0,1.15) (4.8,0.35) (5.8,0)};
  \coordinate (P) at (1.70,1.88);
  \fill[DeepBlue] (P) circle (2.3pt);
  \draw[guide] (P)--(1.70,3.6);
  \draw[<->] (2.02,1.88)--(2.02,3.6);
  \node[right] at (2.02,2.78) {$H-y$};
  \draw[vector] (P)--++(0.74,-0.43) node[below right] {$\vec v$};
  \node[left] at (0,3.6) {$A(0,H)$};
  \node[below] at (5.8,0) {$B(L,0)$};
  \node[align=left,TextGray] at (4.65,3.8)
    {$\displaystyle v=\sqrt{2g(H-y)}$};
  \node[anchor=north,TextGray] at (3.0,-.45)
    {$\displaystyle \dd t=\frac{\dd s}{v}$};
  \begin{scope}[shift={(3.75,2.25)}]
    \coordinate (U) at (0,0);
    \coordinate (V) at (1.35,0);
    \coordinate (W) at (1.35,-0.78);
    \draw[component] (U)--(V) node[midway,above] {$\dd x$};
    \draw[component] (V)--(W) node[midway,right] {$|\dd y|$};
    \draw[DeepBlue,very thick] (U)--(W) node[midway,below left] {$\dd s$};
    \draw[TextGray,line width=.45pt]
      ($(V)+(-.14,0)$)--($(V)+(-.14,-.14)$)--($(V)+(0,-.14)$);
  \end{scope}
\end{tikzpicture}
\end{center}}


回到运动学中，假设一艘船沿给定航线行驶。规定一个正方向后，为了确定船的位置，需要知道
它相对港口的带符号位移。船员能够记录船在不同时刻沿航线的速度，并据此推断当前位置。
如果船始终作匀速直线运动，只需用速度乘时间即可得到位移；但实际速度会不断变化，因而
需要由速度随时间的函数计算变速运动中的位移。若另求航程，才应把各段路程作为正量相加。

\partitionDiagram

把时间分成许多小段。设物体在时刻 $t$ 的速度为 $v(t)$，若时间段 $\Delta t$ 足够小，
可近似认为速度在这段时间内不变，于是从 $t$ 到 $t+\Delta t$ 的位移近似为
\begin{equation}
 \Delta x\simeq v(t)\Delta t.
\end{equation}
在 $x$--$t$ 图像上，这段位移是以 $v(t)$ 为斜率的短直线段在竖直方向的投影；在
$v$--$t$ 图像上，它是宽为 $\Delta t$、高为 $v(t)$ 的窄矩形面积。把所有小段相加，
并令 $\Delta t$ 趋于零，折线趋近真实的 $x(t)$ 曲线，矩形面积之和趋近 $v$--$t$
曲线下的面积。

已知位置函数时，速度由
\begin{equation}
 v(t)=\frac{\dd x}{\dd t}(t)
\end{equation}
给出。现在要解决的是它的逆问题：已知一条曲线各点的导数，怎样求曲线本身。解决这类
问题的方法就是积分。

若函数 $F(x)$ 的导数为 $f(x)$，则称 $F(x)$ 是 $f(x)$ 的一个原函数，并写成
\begin{equation}
 \int f(x)\dd x=F(x)+C,
\end{equation}
这称为 $f(x)$ 的不定积分（indefinite integration），其中 $f(x)$ 为被积函数，$x$ 为
积分变量，$C$ 为积分常数。若原函数已知，则 $f(x)$ 在 $x_1$ 到 $x_2$ 之间的定积分为
\begin{equation}
 \int_{x_1}^{x_2}f(x)\dd x
 =\left.F(x)\right|_{x_1}^{x_2}=F(x_2)-F(x_1).
\end{equation}

\integralDiagram

积分可以处理曲线面积、复杂形体体积以及许多连续求和问题。基本思想是把复杂对象分割
成简单的小对象，先算出每个小对象的相应量，再在小对象趋于无穷小时把它们积分起来。
积分的线性规律为
\begin{equation}
 \int af(x)\dd x=a\int f(x)\dd x,
\end{equation}
\begin{equation}
 \int[f(x)+g(x)]\dd x
 =\int f(x)\dd x+\int g(x)\dd x.
\end{equation}

\Needspace{0.32\textheight}
\subsection{定积分的方向、运算与检验}

前述分割求和给出了定积分。它在 $v$--$t$ 图上表示有向面积：横轴下方的部分给负位移，而路程应对速率 $|v|$ 积分。类似地，
密度不均匀细棒的质量为 $m=\int\lambda(x)\dd x$，变力做功为
$W=\int\vec F\cdot\dd\vec r$，积分的共同思想都是“局部近似简单、整体取极限求和”。

若 $F'(x)=f(x)$，则
\[
 \int_a^b f(x)\,\dd x=F(b)-F(a).
\]
积分上限、下限和积分变量缺一不可；换元时微分和边界必须同时变换。分部积分
\[
 \int_a^b u\,\dd v=[uv]_a^b-\int_a^b v\,\dd u
\]
来自乘积求导法则，适于含对数、反三角函数或两个不同类型函数乘积的积分。结果求出后可通过
求导核对原函数，并用量纲和边界极限检查物理答案。

\subsection{积分建立物理量的空间分布}

连续分布的质心坐标满足
\[
 x_C=\frac1M\int x\,\dd m.
\]
关键是先把质量元写成合适坐标的函数，如细线 $\dd m=\lambda\dd l$、薄片
$\dd m=\sigma\dd S$、立体 $\dd m=\rho\dd V$。转动惯量
$I=\int r_\perp^2\dd m$ 同样依赖转轴到质量元的垂直距离。选择对称坐标后，许多奇函数部分积分为零；
在结果中保留这种对称性，比从一开始展开全部坐标更简洁。

微分方程把“局部变化规律”还原成运动。若 $m\dot v=F(v)$，可写成
$\dd v/F(v)=\dd t/m$ 再积分；更一般的 $F(v,t)$ 只有在它能够分解为速度函数与时间函数的乘积时才能分离变量。
由 $v\dd v/\dd x=a$ 可在不显含时间的问题中把二阶运动方程降阶。
积分常数由初始条件确定，不能因只求到一个原函数而省略。遇到分段力或分段约束时，应在每一段
分别积分，并在转换时刻用位置、速度的连续性或冲量条件衔接。

积分也能给出曲线长度、面积和体积。平面极坐标的一般面积元为
$\dd S=r\dd r\dd\theta$，柱坐标体积元为 $\dd V=r\dd r\dd\theta\dd z$；这些微元中的几何因子
不是记号装饰，而是坐标网格随位置伸缩的结果。若区域从原点延伸到边界 $r=R(\theta)$，先对半径积分才得到
$\int_0^{R(\theta)}r\dd r\dd\theta=\tfrac12R^2(\theta)\dd\theta$；后者是微小扇形面积，不是一般面积元。
先画出被积区域、再确定积分次序，能够避免漏算或重复。

\subsection{近似求积与物理积分的检查}

有些积分没有初等原函数，但定积分仍然有清楚的物理意义。可以把区间划分为小段，用矩形、梯形或
中点值近似每段面积；分割越细，数值和越接近真实积分。实验数据本来就是离散的，此时
\[
 \int_{t_1}^{t_2}F(t)\,\dd t
 \simeq\sum_i\frac{F_i+F_{i+1}}2(t_{i+1}-t_i)
\]
可直接由测得的力—时间数据估计冲量。数值积分给出的是近似值，应同时说明时间步长和误差怎样随
分割改变，而不能把计算器输出的许多小数位都当成有效数字。

广义积分还要求检查端点或无穷远是否收敛。例如总质量
$M=\int_0^\infty\lambda(x)\dd x$ 只有在密度衰减足够快时才有限；势能或电势把无穷远取作零点时，
也隐含了相应积分收敛。若积分发散，往往表示理想化模型不能延伸到无限远，而不是简单地说物理量
“等于无穷”。无限长导线、无限大平面等模型通常只能可靠给出场强或两个有限点之间的势差。

含参数的积分能够表示一族物理状态。若积分区间固定并满足适当光滑条件，可在积分号下求导：
\[
 \frac{\dd}{\dd a}\int_{x_1}^{x_2}f(x,a)\,\dd x
 =\int_{x_1}^{x_2}\frac{\partial f}{\partial a}\,\dd x.
\]
这个方法可把难积分转化为较容易的积分，再对参数积分回来；积分常数则由某个易求的参数值确定。
若上下限也依赖参数，还要加上端点移动产生的边界项。

任何物理积分至少应通过三项检查：被积式乘微元的量纲是否正确，改变积分方向时符号是否合理，
以及在均匀分布、恒力或恒速极限下能否退化为熟悉的初等公式。先写清微元代表什么，再写积分限，
往往比先套积分表更能避免错误。


\begin{origexamplebox}{}\label{ch01:ex:s15-1-17}
设 $a>0,k>0$。计算通过坐标原点的直线 $f(x)=kx$ 在 $x=0$ 到 $x=a$ 之间与横轴所围成的几何面积。
\areaLineDiagram
\end{origexamplebox}

\Needspace{6\baselineskip}\begin{xesolution}
把区间 $[0,a]$ 分成许多宽为 $\dd x$ 的小段，每个窄条的面积为
\[
 \dd S=f(x)\dd x=kx\dd x.
\]
将它们相加：
\[
 S=\int_0^a kx\dd x
 =\left.\frac12kx^2\right|_0^a
 =\boxed{\frac12ka^2}.
\]
这正是底为 $a$、高为 $ka$ 的三角形面积 $\frac12a\cdot ka$。
\end{xesolution}

\Needspace{0.32\textheight}
\begin{origexamplebox}{}\label{ch01:ex:s15-1-18}
设 $a>0,k>0$。计算抛物线 $f(x)=kx^2$ 在 $x=0$ 到 $x=a$ 之间与横轴所围成的几何面积。
\areaParabolaDiagram
\end{origexamplebox}

\Needspace{6\baselineskip}\begin{xesolution}
取宽为 $\dd x$ 的竖直窄条，其面积元为
\[
 \dd S=kx^2\dd x.
\]
因此
\[
 S=\int_0^a kx^2\dd x
 =\left.\frac{kx^3}{3}\right|_0^a
 =\boxed{\frac13ka^3}.
\]
若以外接矩形面积 $a\cdot ka^2=ka^3$ 作比较，抛物线下的面积恰为它的三分之一。
\end{xesolution}

\Needspace{0.42\textheight}
\begin{origexamplebox}{}\label{ch01:ex:s15-1-19}
把下列计算圆面积的分割思路用积分表示出来并计算圆的面积，与熟悉的公式
$S=\pi R^2$ 比较。
\begin{enumerate}[label=(\alph*),leftmargin=2.2em]
 \item 从圆心出发把圆分成许多扇形，每个扇形的顶角为 $\dd\theta$；
 \item 把圆分成许多厚度为 $\dd r$ 的圆环，每个圆环的内半径为 $r$；
 \item 把圆切成许多矩形，每个矩形到圆心的距离为 $x$，宽为 $\dd x$。
\end{enumerate}
\circlePartitionsDiagram
\end{origexamplebox}

\Needspace{6\baselineskip}\begin{xesolution}
三种分割方式都在计算同一个圆面积，只是所选面积元不同。

\textbf{(a) 扇形分割。}\quad 半径为 $R$、圆心角为 $\dd\theta$ 的小扇形面积为
\[
 \dd S=\frac12R^2\dd\theta.
\]
绕圆一周：
\[
 S=\int_0^{2\pi}\frac12R^2\dd\theta
 =\boxed{\pi R^2}.
\]

\textbf{(b) 圆环分割。}\quad 半径为 $r$、厚度为 $\dd r$ 的薄圆环可近似展开成长为
$2\pi r$、宽为 $\dd r$ 的矩形：
\[
 \dd S=2\pi r\dd r,
 \qquad
 S=\int_0^R2\pi r\dd r=\boxed{\pi R^2}.
\]

\textbf{(c) 竖条分割。}\quad 距圆心横坐标为 $x$ 的竖条，其高度由
$x^2+y^2=R^2$ 给出：
\[
 h(x)=2\sqrt{R^2-x^2}.
\]
所以
\[
 S=\int_{-R}^{R}2\sqrt{R^2-x^2}\dd x
 =4\int_0^R\sqrt{R^2-x^2}\dd x.
\]
令 $x=R\sin\phi$，$\dd x=R\cos\phi\dd\phi$：
\[
 S=4R^2\int_0^{\pi/2}\cos^2\phi\dd\phi
 =4R^2\cdot\frac\pi4
 =\boxed{\pi R^2}.
\]
三种积分都得到熟悉的圆面积公式。
\end{xesolution}

\begin{origexamplebox}{}\label{ch01:ex:s15-1-20}
利用积分证明圆锥体体积是同底等高圆柱体积的 $1/3$。
\end{origexamplebox}

\Needspace{6\baselineskip}\begin{xesolution}
设圆锥高为 $H$、底面半径为 $R$。从锥顶起沿轴线建立坐标 $x$，在 $x$ 处取厚度为
$\dd x$ 的圆盘。由相似三角形
\[
 r(x)=\frac{R}{H}x.
\]
圆盘体积元为
\[
 \dd V=\pi r^2(x)\dd x
 =\pi\frac{R^2}{H^2}x^2\dd x.
\]
从锥顶累加到底面：
\[
 V=\int_0^H\pi\frac{R^2}{H^2}x^2\dd x
 =\pi\frac{R^2}{H^2}\cdot\frac{H^3}{3}
 =\boxed{\frac13\pi R^2H}.
\]
同底等高圆柱体积为 $\pi R^2H$，故圆锥体积恰为其三分之一。
\end{xesolution}

在物理中经常会遇到复杂的求和过程。若一个物理量发生变化，与它相关的量也随之改变，
直接求和往往并不现实；但把变化过程分成无穷多足够小的部分时，各小段中的物理量变化
通常可以近似表达。

例如物体在某时刻速度为 $v$，在无限小时间 $\dd t$ 内可近似看作匀速运动：
\begin{equation}
 \Delta x\simeq v(t)\Delta t,
\end{equation}
在 $\Delta t\to0$ 的极限下写成
\begin{equation}
 \dd x=v\dd t.
\end{equation}
已知 $v(t)$ 时，从 $t_1$ 到 $t_2$ 的位移为
\begin{equation}
 x(t_2)-x(t_1)=\int_{t_1}^{t_2}v(t)\dd t.
\end{equation}
这里得到的是带符号的位移；若要求路程，应在速度变号处分段积分 $|v(t)|$。

又如某物体的热容量随温度变化为 $C(T)$，在小温区 $\Delta T$ 内吸收的热量近似为
\begin{equation}
 \Delta Q\simeq C(T)\Delta T.
\end{equation}
取极限得 $\dd Q=C(T)\dd T$，所以物体从 $T_1$ 升温至 $T_2$ 所需热量为
\begin{equation}
 Q=\int_{T_1}^{T_2}C(T)\dd T.
\end{equation}

\begin{origexamplebox}{}\label{ch01:ex:s15-1-21}
利用积分写出下列两个物体质量的表达式：
\begin{enumerate}[label=(\alph*),leftmargin=2.2em]
 \item 轴对称物体的底部为 $x=0$、顶部为 $x=H$，半径随 $x$ 的关系为 $r(x)$，
 密度随高度的关系为 $\rho(x)$；
 \item 半径为 $R$ 的球对称物体，密度与距球心距离 $r$ 的关系由函数 $\rho(r)$ 给出。
\end{enumerate}
\end{origexamplebox}

\Needspace{6\baselineskip}\begin{xesolution}
质量元总是 $\dd m=\rho\dd V$，关键在于按对称性选择体积元。

\textbf{(a)} 在高度 $x$ 处取厚度 $\dd x$ 的薄圆盘：
\[
 \dd V=\pi r^2(x)\dd x,
 \qquad
 \dd m=\rho(x)\pi r^2(x)\dd x.
\]
所以
\[
 \boxed{M=\int_0^H\rho(x)\pi r^2(x)\dd x}.
\]

\textbf{(b)} 在距球心 $r$ 处取厚度 $\dd r$ 的薄球壳，其体积元为
$4\pi r^2\dd r$，故
\[
 \boxed{M=4\pi\int_0^R\rho(r)r^2\dd r}.
\]
\massElementsDiagram
两个式子的上下限都从对称中心或底部开始，直到物体边界；量纲均为
$[\rho][\text{体积}]=[\text{质量}]$。
\end{xesolution}

\Needspace{0.38\textheight}
\begin{origexamplebox}{}\label{ch01:ex:s15-1-22}
一维运动中，外力与物体位置的关系为 $F(x)$。物体从 $x_1$ 运动到 $x_2$ 时，外力的功为
\[
 W_{12}=\int_{x_1}^{x_2}F(x)\dd x.
\]
计算下列情况下给定外力所做的功：
\begin{enumerate}[label=(\alph*),leftmargin=2.2em]
 \item 弹性力 $F=-kx$ 从 $x_1$ 到 $x_2$ 的功；
 \item 万有引力 $F=-\dfrac{Gm_1m_2}{r^2}$ 从 $r_1$ 到 $r_2$ 的功；
 \item 分子力 $F=\dfrac{A}{r^6}-\dfrac{B}{r^{12}}$ 从 $r_1$ 到 $r_2$ 的功；
 \item 理想气体分别经过准静态等温过程和准静态可逆绝热过程，从体积 $V_1$ 膨胀到 $V_2$ 对外所做的功。
\end{enumerate}
\pvWorkDiagram
\end{origexamplebox}

\Needspace{6\baselineskip}\begin{xesolution}
\textbf{(a)}
\[
 W=\int_{x_1}^{x_2}(-kx)\dd x
 =\boxed{\frac12k(x_1^2-x_2^2)}.
\]
它等于弹性势能的减少量。

\textbf{(b)} 取径向向外为正：
\[
 W=\int_{r_1}^{r_2}-\frac{Gm_1m_2}{r^2}\dd r
 =\boxed{Gm_1m_2\left(\frac1{r_2}-\frac1{r_1}\right)}.
\]
若物体向内运动，$r_2<r_1$，引力做正功。

\textbf{(c)}
\begin{align*}
 W&=\int_{r_1}^{r_2}\left(\frac A{r^6}-\frac B{r^{12}}\right)\dd r\\
 &=\boxed{\left[-\frac{A}{5r^5}+\frac{B}{11r^{11}}\right]_{r_1}^{r_2}}.
\end{align*}

\textbf{(d)} 气体准静态膨胀时，气体对外做功 $W=\int p\dd V$。
这里 $n$ 是气体的物质的量，$R_g$ 是摩尔气体常量，$T$ 是热力学温度。
以下取一定量理想气体，并把相应过程方程作为本题积分练习的已知规律；
过程方程的物理推导在热学中讨论。

等温过程中 $pV=nR_gT$，温度不变：
\[
 W_{\rm iso}=\int_{V_1}^{V_2}\frac{nR_gT}{V}\dd V
 =nR_gT[\ln V]_{V_1}^{V_2}
 =\boxed{nR_gT\ln\frac{V_2}{V_1}}.
\]
对于可逆绝热过程，还要求所讨论温区内的定压、定容摩尔热容 $C_p,C_V$
可视为常量。记热容比 $\gamma=C_p/C_V>1$，过程方程为 $pV^\gamma=C$，
其中 $C$ 是沿这条绝热线保持不变的常数，不是热容。于是
\begin{align*}
 W_{\rm ad}
 &=\int_{V_1}^{V_2}CV^{-\gamma}\dd V\\
 &=\frac{C}{1-\gamma}
   \left(V_2^{1-\gamma}-V_1^{1-\gamma}\right)\\
 &=\boxed{\frac{p_1V_1-p_2V_2}{\gamma-1}}.
\end{align*}
图中曲线下的阴影面积就是相应过程的功；同一初态膨胀到同一体积时，绝热线下降更快，
其曲线下面积小于等温过程。
\end{xesolution}

积分还能处理物理量之间更复杂的依赖关系。通常先用小量分析写出过程中的物理关系，再
取极限得到微分方程，最后通过积分得到所需结果。

\begin{origexamplebox}{}\label{ch01:ex:s15-1-23}
质量为 $m$ 的质点受到正比于速度的阻力 $f=-\alpha v$，在其作用下作减速运动。设
$t=0$ 时速度为 $v_0$，且位于坐标原点 $x=0$。求速度随时间的规律 $v(t)$、前进距离
随时间的关系 $x(t)$，以及速度与位移的关系 $v(x)$。
\end{origexamplebox}

\Needspace{6\baselineskip}\begin{xesolution}
取初速度方向为正。牛顿第二定律给出
\[
 m\frac{\dd v}{\dd t}=-\alpha v.
\]
分离变量并代入 $v(0)=v_0$：
\[
 \int_{v_0}^{v}\frac{\dd v'}{v'}
 =-\frac\alpha m\int_0^t\dd t',
\]
故
\[
 \boxed{v(t)=v_0\e^{-\alpha t/m}}.
\]
再由 $\dd x=v\dd t$：
\[
 x(t)=\int_0^tv_0\e^{-\alpha t'/m}\dd t'
 =\boxed{\frac{mv_0}{\alpha}
 \left(1-\e^{-\alpha t/m}\right)}.
\]
若直接求 $v(x)$，利用 $\dd v/\dd t=v\dd v/\dd x$：
\[
 mv\frac{\dd v}{\dd x}=-\alpha v.
\]
在 $v>0$ 的运动阶段约去 $v$ 并积分，得到
\[
 \boxed{v(x)=v_0-\frac\alpha m x}.
\]
当 $t\to\infty$ 时速度渐近于零，极限位移为
$x_\infty=mv_0/\alpha$，与 $v(x)=0$ 给出的结果一致。
\end{xesolution}

\begin{origexamplebox}{}\label{ch01:ex:s15-1-24}
物体在重力作用下下落，同时受到正比于速度的阻力 $f=-\alpha v$。已知 $t=0$ 时它在
竖直向下坐标系的原点处静止，求此后速度随时间的关系 $v(t)$。
\end{origexamplebox}

\Needspace{6\baselineskip}\begin{xesolution}
取竖直向下为正，运动方程为
\[
 m\frac{\dd v}{\dd t}=mg-\alpha v.
\]
引入终端速度 $v_\infty=mg/\alpha$，则
\[
 \frac{\dd v}{\dd t}=\frac\alpha m(v_\infty-v).
\]
分离变量并用 $v(0)=0$：
\[
 \int_0^v\frac{\dd v'}{v_\infty-v'}
 =\frac\alpha m\int_0^t\dd t'.
\]
因此
\[
 \boxed{v(t)=\frac{mg}{\alpha}
 \left(1-\e^{-\alpha t/m}\right)}.
\]
初始时 $v=0$、加速度为 $g$；时间足够长时 $v\to mg/\alpha$、加速度趋于零，符合受力
平衡的终端速度图景。
\end{xesolution}

\begin{origexamplebox}{}\label{ch01:ex:s15-1-25}
质点在 $t=0$ 时位于坐标原点，初速度为 $v_0$。它受到正比于速度平方的阻力
$f=-\beta v^2$。求速度随时间衰减的规律 $v(t)$。
\end{origexamplebox}

\Needspace{6\baselineskip}\begin{xesolution}
在速度尚为正的阶段，运动方程为
\[
 m\frac{\dd v}{\dd t}=-\beta v^2.
\]
分离变量：
\[
 \int_{v_0}^{v}\frac{\dd v'}{v'^2}
 =-\frac\beta m\int_0^t\dd t'.
\]
于是
\[
 -\frac1v+\frac1{v_0}=-\frac\beta m t,
\]
即
\[
 \boxed{v(t)=\frac{v_0}{1+\frac{\beta v_0}{m}t}}.
\]
与线性阻力的指数衰减不同，平方阻力下速度按 $1/t$ 的量级衰减。阻力写作
$-\beta v^2$ 默认 $v\ge0$；一般一维形式应写成 $-\beta v|v|$，才能始终与速度反向。
\end{xesolution}

\Needspace{0.30\textheight}
\begin{origexamplebox}{}\label{ch01:ex:s15-1-26}
坐标原点处有质量为 $M$ 的物体 $A$，它始终保持静止。在距它 $r_0$ 处有另一个质量远小
于它的质点，在二者万有引力作用下由静止出发落向 $A$。求它们相撞的时间。
\end{origexamplebox}

\Needspace{6\baselineskip}\begin{xesolution}
把小质点视为试探质点，并把中心物体视为质点；“相撞”即 $r\to0$。机械能守恒给出
\[
 \frac12mv^2-\frac{GMm}{r}=-\frac{GMm}{r_0},
\]
因此向内运动的速率为
\[
 v(r)=\sqrt{2GM\left(\frac1r-\frac1{r_0}\right)}.
\]
\radialFallDiagram
图中虚圈表示 $t=0$ 时位于 $r_0$ 的初始位置，实点表示当前位于 $r$ 的位置；速度箭头指向
半径减小的方向，所以 $\dd r/\dd t<0$。
由于 $\dd r/\dd t=-v$，下落时间为
\[
 t=\int_0^{r_0}\frac{\dd r}
 {\sqrt{2GM(1/r-1/r_0)}}.
\]
令 $r=r_0\sin^2\phi$，则 $\dd r=2r_0\sin\phi\cos\phi\dd\phi$，积分化为
\[
 t=2r_0\sqrt{\frac{r_0}{2GM}}
 \int_0^{\pi/2}\sin^2\phi\dd\phi.
\]
利用 $\int_0^{\pi/2}\sin^2\phi\dd\phi=\pi/4$，得到
\[
 \boxed{t=\frac{\pi}{2\sqrt2}
 \sqrt{\frac{r_0^3}{GM}}}.
\]
小质点的质量被约去，体现了引力场中自由落体时间与试探质量无关。若中心物体有有限
半径 $R_A$，真正碰到表面的时间应把积分下限改为 $R_A$，会小于上述结果。
\end{xesolution}

\begin{origexamplebox}{}\label{ch01:ex:s15-1-27}
放射性衰变是随机过程。对含有大量同类原子的样品，用 $N(t)$ 表示尚未衰变原子数的期望值
（实际宏观样品中也可视为平均数）。在无穷小时间间隔中，期望衰变数正比于当前期望总数，即
\[
 \frac{\dd N}{\dd t}=-\lambda N,
\]
其中 $\lambda$ 为衰变常数。若 $t=0$ 时有 $N_0$ 个放射性原子，求任意时刻 $t$ 尚未
衰变的原子数期望值；并给出衰变常数与半衰期 $\tau_{1/2}$ 的关系。
\end{origexamplebox}

\Needspace{6\baselineskip}\begin{xesolution}
分离变量并积分：
\[
 \int_{N_0}^{N(t)}\frac{\dd N'}{N'}
 =-\lambda\int_0^t\dd t'.
\]
所以
\[
 \ln\frac{N(t)}{N_0}=-\lambda t,
 \qquad
 \boxed{N(t)=N_0\e^{-\lambda t}}.
\]
半衰期定义为 $N(\tau_{1/2})=N_0/2$，代入上式：
\[
 \frac12=\e^{-\lambda\tau_{1/2}},
\]
因此
\[
 \boxed{\tau_{1/2}=\frac{\ln2}{\lambda}},
 \qquad
 \boxed{\lambda=\frac{\ln2}{\tau_{1/2}}}.
\]
衰变常数的量纲是时间的倒数；样品经过 $n$ 个半衰期后剩余比例为 $2^{-n}$。
\end{xesolution}

\Needspace{0.36\textheight}
\begin{origexamplebox}{}\label{ch01:ex:s15-1-28}
用定积分给出任意光滑曲线 $y(x)$ 在 $x_1$ 到 $x_2$ 之间的部分曲线长度 $L_{12}$。
\arcLengthDiagram
\end{origexamplebox}

\Needspace{6\baselineskip}\begin{xesolution}
在曲线上取相邻两点。图中右侧把局部三角形放大：横、纵两条直角边分别表示
$\dd x,\dd y$，斜边表示弧长微元 $\dd s$。局部曲线段在极限
下可看作切线上的小直线段，由勾股定理
\[
 \dd s=\sqrt{(\dd x)^2+(\dd y)^2}.
\]
把 $\dd x$ 提出：
\[
 \dd s=\sqrt{1+\left(\frac{\dd y}{\dd x}\right)^2}\,|\dd x|.
\]
若沿 $x$ 增大的方向从 $x_1$ 走到 $x_2$，且 $x_2>x_1$，则
\[
 \boxed{L_{12}=\int_{x_1}^{x_2}
 \sqrt{1+[y'(x)]^2}\dd x}.
\]
被积函数始终不小于 $1$，所以弧长不小于两端点的水平距离 $x_2-x_1$；这是公式的一个
直接检验。
\end{xesolution}

\Needspace{0.39\textheight}
\begin{origexamplebox}{}\label{ch01:ex:s15-1-29}
坐标系中有 $A,B$ 两点，坐标分别为 $(0,H)$ 和 $(L,0)$。一条光滑轨道 $y(x)$ 连接
$A,B$。求物体由 $A$ 点从静止出发、沿该轨道下落到 $B$ 点所需的时间。
\trackTimeDiagram
\end{origexamplebox}

\Needspace{6\baselineskip}\begin{xesolution}
假设轨道光滑且没有能量损失。物体到达横坐标 $x$ 处时下降高度为 $H-y(x)$。由机械能
守恒
\[
 \frac12mv^2=mg[H-y(x)],
\]
故速率为
\[
 v(x)=\sqrt{2g[H-y(x)]}.
\]
沿轨道的弧长元由例~\ref{ch01:ex:s15-1-28} 给出：
\[
 \dd s=\sqrt{1+[y'(x)]^2}\dd x.
\]
而 $\dd t=\dd s/v$，所以总时间为
\[
 \boxed{t_{AB}=\int_0^L
 \frac{\sqrt{1+[y'(x)]^2}}
 {\sqrt{2g[H-y(x)]}}\dd x}.
\]
\trackTimeSolutionDiagram
积分在起点处的分母趋于零，因此这是需要检查收敛性的反常积分；其是否收敛取决于轨道离开 $A$ 点的
方式。若起点切线具有非零向下斜率，则 $H-y(x)$ 与 $x$ 同阶，积分像
$\int\dd x/\sqrt{x}$ 一样收敛；若起点是严格水平且曲线足够光滑，物体可能不能自行
离开静止点，相应积分会发散。该式也是最速降线问题的出发点，但这里只要求给出任意
已知轨道的下落时间。
\end{xesolution}
}

\section{复数相量}

\subsection{复数如何记录振幅与相位}
虚数单位 $\ii$ 满足 $\ii^2=-1$。复数 $z=a+\ii b$ 的实部记作 $\Re z=a$，虚部记作 $\Im z=b$，
可以用平面上的点 $(a,b)$ 表示。它的模为 $|z|=\sqrt{a^2+b^2}$；
从实轴正向量到向量 $(a,b)$ 的角称为辐角 $\phi$，应根据象限确定，而不能只取
$\arctan(b/a)$ 的主值。因此
\[
 z=|z|(\cos\phi+\ii\sin\phi)=|z|e^{\ii\phi}.
\]
最后一步使用欧拉公式 $e^{\ii\phi}=\cos\phi+\ii\sin\phi$。
将指数级数中的偶次项与奇次项分别合并，因 $\ii^{2n}=(-1)^n$、
$\ii^{2n+1}=\ii(-1)^n$，恰好得到余弦和正弦的级数，因而这不是新的物理假设。

简谐量可写成复指数的实部：
\begin{equation}
 x(t)=A\cos(\omega t+\phi)=\Re\!\left(\widetilde x\,e^{\ii\omega t}\right),
 \qquad \widetilde x=Ae^{\ii\phi}.
\end{equation}
其中 $\Re$ 表示取实部。复平面向量随时间以角速度 $\omega$ 转动，
其在实轴上的投影才是实际位移 $x(t)$；常量 $\widetilde x$ 则称为相量，记录振幅 $A\ge0$ 和初相 $\phi$。
相量本身不是额外的空间位移。

\subsection{相量的运算}
同频信号共有时间因子 $e^{\ii\omega t}$，因此先把相量相加，最后取实部即可：
\[
 x_1+x_2=\Re\!\left[(\widetilde x_1+\widetilde x_2)e^{\ii\omega t}\right].
\]
不同频率没有共同时间因子，不能直接当作静态相量相加。对固定相量求导，
\[
 \frac{\dd}{\dd t}\bigl(\widetilde x e^{\ii\omega t}\bigr)
 =\ii\omega\widetilde x e^{\ii\omega t},\qquad
 \int\widetilde x e^{\ii\omega t}\dd t
 =\frac{\widetilde x}{\ii\omega}e^{\ii\omega t}+C\quad(\omega\ne0).
\]
所以微分对应乘 $\ii\omega$：振幅乘 $|\omega|$，当 $\omega>0$ 时相位超前 $\pi/2$；
积分对应除 $\ii\omega$，但还必须保留积分常数 $C$。例如位移的导数是速度，
再次求导时乘 $(\ii\omega)^2=-\omega^2$，故简谐运动满足 $\ddot x=-\omega^2x$。

\begin{origexamplebox}{同频简谐量的合成}
$x_1=A\cos\omega t$，$x_2=A\cos(\omega t+2\pi/3)$。求 $x=x_1+x_2$ 的振幅和初相。
\end{origexamplebox}
\Needspace{6\baselineskip}\begin{xesolution}
相量为 $\widetilde x_1=A$、$\widetilde x_2=Ae^{\ii2\pi/3}$，所以
\[
 \widetilde x=A\left(1-\frac12+\ii\frac{\sqrt3}{2}\right)
 =Ae^{\ii\pi/3}.
\]
合相量的实部、虚部分别为 $A/2$、$\sqrt3A/2$，故其模为
$\sqrt{A^2/4+3A^2/4}=A$；它在第一象限，且正切为 $\sqrt3$，所以初相为 $\pi/3$。
图~\ref{ch01:fig:phasor-sum} 中，两相量首尾相接得到合相量，虚线表示平移后的第二相量。
\begin{center}
\begin{tikzpicture}[>=Latex,scale=1]
 \draw[axis] (-1.9,0)--(3.3,0) node[right]{$\Re$};
 \draw[axis] (0,-.4)--(0,2.65) node[above]{$\Im$};
 \draw[vector] (0,0)--(2.4,0) node[midway,below=5pt]{$\widetilde x_1=A$};
 \draw[vector] (0,0)--(-1.2,2.07846) node[above left]{$\widetilde x_2$};
 \draw[guide,->] (2.4,0)--(1.2,2.07846);
 \draw[vector,orange!85!black] (0,0)--(1.2,2.07846) node[above right]{$\widetilde x$};
 \draw[->,TextGray] (.75,0) arc (0:60:.75);
 \node at (.87,.60) {$\pi/3$};
\end{tikzpicture}
\figcap{同频相量相加：合相量长度为 $A$，辐角为 $\pi/3$。}\label{ch01:fig:phasor-sum}
\end{center}
因此
\[
 \boxed{x=A\cos(\omega t+\pi/3)}.
\]
合振幅仍为 $A$，不是 $2A$。极限检查：两相量夹角为 $120^\circ$，等腰三角形的合矢量长度确为一边长。
\end{xesolution}

\section{常微分方程的物理读法}

\subsection{一阶方程与初始条件}
微分方程把未知函数与它的变化率联系起来；上方一点表示对时间求一次导数，
两点表示二次导数，即 $\dot x=\dd x/\dd t$、$\ddot x=\dd^2x/\dd t^2$。
“常微分”指未知量只依赖一个自变量。以常数 $a$ 的一阶线性方程为例：
\begin{equation}
 \dot x+ax=f(t)
\end{equation}
两边乘积分因子 $e^{at}$，左端由乘积求导法则变成一个完整导数：
\[
 e^{at}\dot x+ae^{at}x=\frac{\dd}{\dd t}(e^{at}x)=e^{at}f(t).
\]
从 $0$ 积分到 $t$，并使用 $x(0)=x_0$，得到
\[
 e^{at}x(t)-x_0=\int_0^t e^{as}f(s)\dd s,\qquad
 \boxed{x(t)=e^{-at}\left[x_0+\int_0^t e^{as}f(s)\dd s\right]}.
\]
积分变量 $s$ 与上限 $t$ 必须区分。$f=0,a>0$ 时，$x=x_0e^{-at}$ 衰减；
$a=0$ 时退化为 $x=x_0+\int_0^t f(s)\dd s$。

\subsection{二阶方程、特征根与阻尼}
弹簧振子的线性模型为 $m\ddot x+b\dot x+kx=0$。令
$\gamma=b/(2m)$、$\omega_0=\sqrt{k/m}>0$，得到二阶常系数齐次方程
\begin{equation}
 \ddot x+2\gamma\dot x+\omega_0^2x=0
\end{equation}
试取 $x=e^{rt}$，代入并约去不为零的 $e^{rt}$，得到特征方程及其根
\[
 r^2+2\gamma r+\omega_0^2=0,\qquad r_\pm=-\gamma\pm\sqrt{\gamma^2-\omega_0^2}.
\]
下述三种阻尼分类以 $\gamma>0$ 为前提。给定 $x(0)=x_0$、$\dot x(0)=v_0$：
\begin{itemize}
 \item $0<\gamma<\omega_0$：令 $\omega_d=\sqrt{\omega_0^2-\gamma^2}$。
 共轭复根给出
 \[
 x=e^{-\gamma t}\left[x_0\cos\omega_dt+
 \frac{v_0+\gamma x_0}{\omega_d}\sin\omega_dt\right].
 \]
 括号内振荡，外面的正因子衰减，这是欠阻尼。
 \item $\gamma=\omega_0$：重根下两个独立解是 $e^{-\gamma t}$ 和
 $te^{-\gamma t}$，代回方程即可逐一验证。代入两条初始条件得
 \[
 x=[x_0+(v_0+\gamma x_0)t]e^{-\gamma t},
 \]
 称为临界阻尼。
 \item $\gamma>\omega_0$：两个实根均为负，
 \[
 x=C_+e^{r_+t}+C_-e^{r_-t},\quad
 C_+=\frac{v_0-r_-x_0}{r_+-r_-},\quad
 C_-=\frac{r_+x_0-v_0}{r_+-r_-}.
 \]
 这是过阻尼。它没有周期振荡，但特定初始速度仍可能使物体越过平衡点一次，不能一概称为单调返回。
\end{itemize}
当 $\gamma=0$ 时，根为 $\pm\ii\omega_0$，运动无衰减；当 $\gamma<0$ 时，
模型中的速度项向系统供能，非零解通常增长，不属于上述被动阻尼衰减。
根的实部决定振幅随时间增长还是衰减，虚部决定是否具有振荡因子。
初始条件的个数应与方程阶数相同；只写通解而不代入初始条件，物理问题尚未求完。

\begin{origexamplebox}{一阶弛豫过程}
温度为 $T(t)$ 的小物体置于恒温环境 $T_a$ 中，满足牛顿冷却模型
$\dot T=-k(T-T_a)$，其中 $k>0$，且 $T(0)=T_0$。求 $T(t)$，并检验长时间极限。
\end{origexamplebox}
\Needspace{6\baselineskip}\begin{xesolution}
令温差 $\Theta=T-T_a$，因环境温度不变，方程化为 $\dot\Theta=-k\Theta$。
乘 $e^{kt}$ 后，$\dd(e^{kt}\Theta)/\dd t=0$，说明 $e^{kt}\Theta$ 为常量；
代入 $t=0$ 得该常量为 $T_0-T_a$。因此
\[
 \Theta(t)=(T_0-T_a)e^{-kt},\qquad
 \boxed{T(t)=T_a+(T_0-T_a)e^{-kt}}.
\]
当 $T_0\ne T_a$ 时也可分离变量，积分为
$\ln|\Theta(t)|-\ln|\Theta(0)|=-kt$；温差始终同号，才可化为
$\ln[\Theta(t)/\Theta(0)]=-kt$。当 $T_0=T_a$ 时不能除以温差，
但上面的积分因子解仍直接给出 $T(t)=T_a$。
$k$ 的量纲为 $\mathsf T^{-1}$（例如 $\mathrm{s^{-1}}$），所以指数 $kt$ 无量纲。$t\to\infty$ 时 $T\to T_a$；若 $T_0>T_a$，温度单调下降且不会在有限时间越过环境温度，均与模型的物理含义一致。
\end{xesolution}

\endgroup

\begingroup
\tcbset{breakable/.default=false}
\chapter{质点运动的描述}

{
\graphicspath{{/Users/dongxingxi/Documents/物理竞赛一轮资料/sample/xes-original/kinematics-supplement/}{/Users/dongxingxi/Documents/物理竞赛一轮资料/sample/xes-original/kinematics-supplement/assets/}}
\section{理想物理模型：质点}
\newcommand{\EarthMoonScaleDiagram}{%
\begin{center}\begin{tikzpicture}[x=1cm,y=1cm,>=Latex]
  % 主图按同一线性比例绘制：中心距 10.5 cm，对应天体半径分别
  % 为 1.74 mm 和 0.48 mm。这样保留题图“真实尺度”的物理作用。
  \coordinate (E) at (-5.25,0);
  \coordinate (M) at (5.25,0);
  \shade[ball color=Blue!70] (E) circle (.174);
  \shade[ball color=TextGray!45] (M) circle (.048);
  \draw[<->,TextGray] (-5.25,-.45)--(5.25,-.45)
    node[midway,clearlabel,below=2pt]
      {平均中心距 $3.84\times10^5\,\mathrm{km}$};
  \draw[TextGray] (E)--++(-.42,.58);
  \node[clearlabel,above right,align=left] at (-5.72,.55)
    {地球 $R_E\approx6.37\times10^3\,\mathrm{km}$\\$M_E\approx5.97\times10^{24}\,\mathrm{kg}$};
  \draw[TextGray] (M)--++(.35,.58);
  \node[clearlabel,above left,align=right] at (5.67,.55)
    {月球 $R_M\approx1.74\times10^3\,\mathrm{km}$\\$M_M\approx7.35\times10^{22}\,\mathrm{kg}$};
  \node[clearlabel,below] at (0,-1.05)
    {半径与中心距采用同一比例；$M_M/M_E\approx0.0123$};
\end{tikzpicture}\figcap{月球、地球及月地距离的真实比例。研究月球绕地球公转时可把月球近似为质点；研究着陆时则不能。}\label{ch02:fig:k01-2-1}\end{center}}


自然界中的物体都处在运动过程中；所谓静止，只能相对于某个指定参考系而言。在经典力学的运动学描述中，不把“绝对静止”作为可直接判定的状态。学习力学的第一步，是建立一套能够准确描写运动的方法：先说明研究对象，再规定参考系和坐标，最后用位置随时间的变化表达运动。只有在运动学描述清楚之后，才有可能进一步研究运动状态变化的原因。

物理研究通常从最简单的对象开始。把一个具有确定质量、但形状和内部结构可以忽略的运动物体抽象成几何点，称为\textbf{质点}（mass point）。质点不是客观存在的“无大小物体”，而是针对具体问题建立的理想模型。

能否把物体看成质点，不能只比较物体自身大小，还要比较它的大小与研究尺度，并判断姿态、转动和形变是否影响所问物理量。例如：

\begin{itemize}
  \item 研究跳水运动员入水姿态时，身体各部分的位置关系直接影响结果，不能把运动员看成质点；研究他从北京到上海的行程时，身体尺寸相对于路程可以忽略，可用质点表示。
  \item 研究月球绕地球公转时，月球半径远小于地月距离，通常可把月球当成质点；研究探测器在月面着陆时，月球表面的曲率、地形和自转均可能重要，不能再用同一个简化。
  \item 研究乒乓球质心的抛体轨迹时可以暂时忽略球的大小；研究旋转引起的弧线或球与球拍接触时，必须保留半径和自转。
\end{itemize}

\EarthMoonScaleDiagram

建模时应明确写出“忽略了什么”。质点模型通常忽略物体的形状、姿态、转动和内部自由度，只保留质量与位置。若后续推导需要角速度、转动惯量、接触几何或形变能，就说明质点模型已经不足，必须升级为刚体、连续介质或可变形体模型。

判断模型是否合理的可靠方法，是做尺度和极限检查。设物体特征尺寸为 $\ell$、运动轨迹的特征尺度为 $L$。若 $\ell/L\ll1$，并且观测精度不要求分辨 $\ell$ 量级的差异，质点近似通常可靠；若结果对 $\ell$ 的一阶修正已经不可忽略，就不能仅凭“看起来很小”继续使用质点模型。
}

{
\graphicspath{{/Users/dongxingxi/Documents/物理竞赛一轮资料/sample/xes-original/kinematics-supplement/}{/Users/dongxingxi/Documents/物理竞赛一轮资料/sample/xes-original/kinematics-supplement/assets/}}
\section{时间、位移、速度与加速度}
\newcommand{\SprintGraph}{%
\begin{center}\begin{tikzpicture}
\begin{axis}[
  width=.78\linewidth,height=6.0cm,
  axis lines=left,xlabel={$t/\mathrm{s}$},ylabel={$x/\mathrm{m}$},
  xmin=0,xmax=12.5,ymin=0,ymax=105,
  xtick={0,...,12},ytick={0,20,...,100},
  grid=major,grid style={draw=RuleGray},
  every axis label/.style={fill=none},tick label style={font=\small}
]
\addplot[kcurve,mark=*,mark size=1.3pt] coordinates {
  (0,0)(1,1)(2,4)(3,9)(4,16)(5,25)(6,35)(7,45)(8,55)(9,65)(10,75)(11,85)(12,100)
};
\end{axis}
\end{tikzpicture}\figcap{短跑选手位置随时间变化。曲线越陡，位置变化越快。}\label{ch02:fig:k02-2-2}\end{center}}

\newcommand{\PositionSlopeDiagram}{%
\begin{center}\begin{tikzpicture}[scale=.94,>=Latex]
  \foreach \dx/\dy/\mode in {-3/0/1,3/0/2,0/-5.2/3}{
    \begin{scope}[shift={(\dx,\dy)}]
      \draw[kaxis] (-1.8,0)--(1.9,0) node[right,clearlabel]{$t$};
      \draw[kaxis] (-1.55,-1.35)--(-1.55,1.55) node[above,clearlabel]{$x$};
      \ifnum\mode=1
        \draw[kcurve] (-1.35,-.9)--(1.55,1.05);
        \node[clearlabel,below,align=center] at (0,-1.55){匀速运动\\直线斜率恒定};
      \else\ifnum\mode=2
        \draw[kcurve] plot[smooth] coordinates{(-1.35,.12)(-.7,.35)(0,.64)(.7,1.0)(1.55,1.27)};
        \coordinate (A) at (-.7,.35);\coordinate (B) at (.7,1.0);
        \fill (A) circle (.04);
        \fill (B) circle (.04);
        \draw[secant] (A)--(B);
        \draw[kguide] (-.7,0)--(A) (.7,0)--(B);
        \draw[kguide] (-1.55,.35)--(A) (-1.55,1.0)--(B);
        \draw[<->,TextGray] (-.7,-.45)--(.7,-.45)
          node[midway,clearlabel,below=1pt]{$\Delta t$};
        \draw[<->,TextGray] (1.18,.35)--(1.18,1.0)
          node[midway,clearlabel,right=1pt]{$\Delta x$};
        \draw[TextGray] (A)--(.38,.35);
        \draw[TextGray] (-.30,.35) arc[start angle=0,end angle=24.90,radius=.40];
        \node[clearlabel] at (.20,.58){$\alpha$};
        \node[clearlabel,above left=2pt] at (A) {$A$};
        \node[clearlabel,above right=2pt] at (B) {$B$};
        \node[clearlabel,below=1pt] at (-.7,0) {$t_1$};
        \node[clearlabel,below=1pt] at (.7,0) {$t_2$};
        \node[clearlabel,left=1pt] at (-1.55,.35) {$x$};
        \node[clearlabel,left=1pt] at (-1.55,1.0) {$x+\Delta x$};
        \node[clearlabel,below,align=center] at (0,-1.55){割线斜率：平均速度\\$\Delta t=t_2-t_1$};
      \else
        \draw[kcurve,domain=-1.35:1.5,samples=80] plot (\x,{.12+.65*\x-.1*\x*\x});
        \coordinate (A) at (0,.12);\fill (A) circle (.04);
        \draw[tangent] (-1.0,-.53)--(1.1,.835);
        \node[clearlabel,below right=3pt] at (A) {$A$};
        \node[clearlabel,below,align=center] at (0,-1.55){切线斜率\\瞬时速度};
      \fi\fi
    \end{scope}
  }
\end{tikzpicture}\figcap{$x$-$t$ 图像中平均速度与瞬时速度的几何意义。}\label{ch02:fig:k02-2-3}\end{center}}

\newcommand{\VelocityAreaDiagram}{%
\begin{center}\begin{tikzpicture}[scale=.82,>=Latex]
  \foreach \dx/\mode in {-4.7/1,0/2,4.7/3}{
    \begin{scope}[xshift=\dx cm]
      \draw[kaxis] (-1.8,0)--(1.9,0) node[right,clearlabel]{$t$};
      \draw[kaxis] (-1.55,-.15)--(-1.55,1.65) node[above,clearlabel]{$v$};
      \ifnum\mode=1
        \fill[Blue!12] (-1.1,0) rectangle (1.3,.9);
        \draw[kcurve] (-1.25,.9)--(1.5,.9);
        \draw[kguide] (-1.1,0)--(-1.1,.9) (1.3,0)--(1.3,.9);
        \node[clearlabel,below=2pt] at (-1.1,0){$A$};
        \node[clearlabel,below=2pt] at (1.3,0){$B$};
        \node[clearlabel,above left=1pt] at (-1.1,.9){$C$};
        \node[clearlabel,above right=1pt] at (1.3,.9){$D$};
        \node[clearlabel,below] at (0,-.55){匀速：矩形面积};
      \else
        \path[fill=Blue!12] (-1.1,0)--
          plot[domain=-1.1:1.3,samples=60]
            ({\x},{0.55+0.75*(\x+1.1)-0.22*(\x+1.1)^2})--(1.3,0)--cycle;
        \draw[kcurve]
          plot[domain=-1.1:1.3,samples=60]
            ({\x},{0.55+0.75*(\x+1.1)-0.22*(\x+1.1)^2});
        \draw[kguide] (-1.1,0)--(-1.1,.55) (1.3,0)--(1.3,1.08);
        \node[clearlabel,below=2pt] at (-1.1,0){$A$};
        \node[clearlabel,below=2pt] at (1.3,0){$B$};
        \node[clearlabel,above left=1pt] at (-1.1,.55){$C$};
        \node[clearlabel,above right=1pt] at (1.3,1.08){$D$};
        \ifnum\mode=2
          \foreach \x in {-1.05,-.85,...,1.25}{
            \draw[Blue!65] (\x,0)--
              (\x,{0.55+0.75*(\x+1.1)-0.22*(\x+1.1)^2});}
          \node[clearlabel,below] at (0,-.55){窄矩形求和};
        \else
          \node[clearlabel,below] at (0,-.55){极限：曲线下面积};
        \fi
      \fi
    \end{scope}
  }
\end{tikzpicture}\figcap{$v$-$t$ 图像下的有向面积等于位移；若求路程，应对 $|v|$ 积分。}\label{ch02:fig:k02-2-4}\end{center}}

\newcommand{\ExampleTwoSixGraphs}{%
\begin{center}\begin{tikzpicture}[scale=.78,>=Latex]
  \foreach \dx/\mode in {-4.6/1,0/2,4.6/3}{
    \begin{scope}[xshift=\dx cm]
      \draw[kaxis] (-1.65,0)--(1.75,0) node[right,clearlabel]{$t$};
      \draw[kaxis] (-1.45,-.55)--(-1.45,2.05) node[above,clearlabel]{$x$};
      \ifnum\mode=1
        \draw[kcurve] (-1.35,0)--(0,1.35)--(1.35,0);
        \draw[vector] (-.92,.43)--(-.42,.93);
        \draw[vector] (.42,.93)--(.92,.43);
        \node[clearlabel,below] at (0,-.62){先正向、后反向};
      \else\ifnum\mode=2
        \draw[kcurve,domain=-1.35:1.30,samples=80] plot (\x,{.25*(\x+1.35)^2});
        \draw[vector] (-.65,.12)--(-.15,.34);
        \draw[vector] (.35,.73)--(.82,1.25);
        \node[clearlabel,below] at (0,-.62){斜率为正且增大};
      \else
        \draw[kcurve,domain=-1.35:1.35,samples=100] plot (\x,{1.45*sqrt(max(0,1-(\x/1.35)^2))});
        \draw[vector] (-.85,.99)--(-.35,1.34);
        \draw[vector] (.35,1.34)--(.85,.99);
        \node[clearlabel,below] at (0,-.62){斜率由正变为负};
      \fi\fi
    \end{scope}}
\end{tikzpicture}\figcap{三幅 $x$-$t$ 图的切线斜率分别给出速度的方向与变化趋势。}\label{ch02:fig:k02-ex-2-6}\end{center}}

\newcommand{\ExampleTwoSevenGraphs}{%
\begin{center}\begin{tikzpicture}[scale=.9,>=Latex]
  \begin{scope}[shift={(-4.9,-1.1)}]
    \draw[kaxis] (0,0)--(2.75,0) node[right,clearlabel]{$t$};
    \draw[kaxis] (0,-1.55)--(0,1.85) node[above,clearlabel]{$v$};
    \draw[kcurve] (0,1.3)--(2.4,-1.3);
    \node[clearlabel,left] at (0,1.3){$v_0$};
    \node[clearlabel,left] at (0,-1.3){$-v_0$};
    \draw[kguide] (0,-1.3)--(2.4,-1.3);
    \node[clearlabel,below left] at (0,0){$O$};
    \node[clearlabel,above right=2pt] at (1.2,0){$T/2$};
    \node[clearlabel,above] at (2.4,0){$T$};
    \node[clearlabel,below] at (1.2,-1.7){已知速度图};
  \end{scope}
  \foreach \dx/\dy/\mode/\tag in {-.7/0/1/A,3.3/0/2/B,-.7/-3.3/3/C,3.3/-3.3/4/D}{
    \begin{scope}[shift={(\dx,\dy)}]
      \draw[kaxis] (0,0)--(2.7,0) node[right,clearlabel]{$t$};
      \draw[kaxis] (0,-.25)--(0,1.65) node[above,clearlabel]{$x$};
      \ifnum\mode=1
        \draw[kcurve] (0,0)--(1.2,1.15)--(2.4,0);
      \else\ifnum\mode=2
        \draw[kcurve,domain=0:1,samples=70]
          plot ({2.4*\x},{4.6*\x*(1-\x)});
      \else\ifnum\mode=3
        \draw[kcurve,domain=0:1,samples=70] plot ({2.4*\x},{1.15*(2*\x-\x*\x)});
      \else
        \draw[kcurve] (0,1.15)--(1.2,0)--(2.4,.85);
      \fi\fi\fi
      \node[clearlabel,below left] at (0,0){$O$};
      \node[clearlabel,below=3pt] at (1.2,0){$T/2$};
      \node[clearlabel,below=3pt] at (2.4,0){$T$};
      \node[clearlabel,below] at (1.2,-.65){$\tag$};
    \end{scope}}
\end{tikzpicture}\figcap{左图为给定的 $v$-$t$ 图，右侧四图为待判断的位移图；正确图必须连续且斜率等于给定速度。}\label{ch02:fig:k02-ex-2-7}\end{center}}

\newcommand{\ExampleTwoEightGraphs}{%
\begin{center}\begin{tikzpicture}[scale=.78,>=Latex]
  \foreach \dx/\mode in {-4.4/1,0/2,4.4/3}{
    \begin{scope}[xshift=\dx cm]
      \draw[kaxis] (-1.55,0)--(1.75,0) node[right,clearlabel]{时间};
      \draw[kaxis] (-1.35,-.25)--(-1.35,1.85) node[above,clearlabel]{$v$};
      \ifnum\mode=1
        \fill[Blue!12] (-1.25,0)--(1.25,0)--(1.25,1.35)--cycle;
        \draw[kcurve] (-1.25,0)--(1.25,1.35);
        \draw[kguide] (1.25,0)--(1.25,1.35);
        \draw[kguide] (-1.35,1.35)--(1.25,1.35);
        \node[clearlabel,left] at (-1.35,1.35){$v$};
      \else\ifnum\mode=2
        \fill[Blue!12] (-1.25,0)--(-1.25,.55)--(1.25,1.35)--(1.25,0)--cycle;
        \draw[kcurve] (-1.25,.55)--(1.25,1.35);
        \draw[kguide] (1.25,0)--(1.25,1.35);
        \draw[kguide] (-1.35,.55)--(-1.25,.55) (-1.35,1.35)--(1.25,1.35);
        \node[clearlabel,left] at (-1.35,.55){$v_0$};
        \node[clearlabel,left] at (-1.35,1.35){$v_1$};
      \else
        \fill[Blue!12] (-1.25,0)--(-.42,1.35)--(1.25,0)--cycle;
        \draw[kcurve] (-1.25,0)--(-.42,1.35)--(1.25,0);
        \draw[kguide] (-.42,0)--(-.42,1.35);
        \draw[kguide] (-1.35,1.35)--(-.42,1.35);
        \node[clearlabel,left] at (-1.35,1.35){$v_0$};
        \node[clearlabel,below] at (-.42,0){$t/3$};
      \fi\fi
      \node[clearlabel,below] at (1.25,0){$t$};
    \end{scope}}
\end{tikzpicture}\figcap{总路程等于三幅 $v$-$t$ 图中阴影面积；梯形或三角形面积直接给出结果。}\label{ch02:fig:k02-ex-2-8}\end{center}}

\newcommand{\FreeFallMeetingDiagram}{%
\begin{center}\begin{tikzpicture}[scale=.9,>=Latex]
  \draw[very thick] (-1.4,0)--(1.4,0);
  \node[point,label={[clearlabel]below right=3pt:$B$}] (B) at (0,0) {};
  \node[point,label={[clearlabel]left:$A$}] (A) at (0,3.2) {};
  \draw[<->,TextGray] (.8,0)--(.8,3.2) node[midway,clearlabel,right]{$h$};
  \draw[DeepBlue,very thick,->] (B)--(0,1.15) node[midway,clearlabel,left]{$v_0$};
  \draw[Blue,very thick,->] (A)--(0,2.45) node[midway,clearlabel,left]{$g$};
\end{tikzpicture}\figcap{解图：两球具有相同的重力加速度，因此相对加速度为零，初始间距以相对速度 $v_0$ 关闭。}\label{ch02:fig:k02-ex-2-17}\end{center}}

\newcommand{\BounceDiagram}{%
\begin{center}\begin{tikzpicture}[scale=.85,>=Latex]
  \draw[very thick] (-1.7,0)--(1.7,0);
  \fill (0,2.7) circle (.12);
  \fill (0,.12) circle (.12);
  \draw[kguide] (0,.27)--(0,2.55);
  \draw[<->,TextGray] (.88,0)--(.88,2.7)
    node[midway,clearlabel,right=2pt]{$H$};
  \draw[DeepBlue,very thick,->] (-.35,1.0)--(-.35,.25) node[midway,clearlabel,left]{$u$};
  \draw[Blue,very thick,->] (.35,.25)--(.35,1.0) node[midway,clearlabel,right]{$eu$};
\end{tikzpicture}\figcap{解图：每次反弹速度按 $e$ 成等比衰减，飞行时间按 $e$ 衰减，路程按 $e^2$ 衰减。}\label{ch02:fig:k02-ex-2-19}\end{center}}

\newcommand{\KQuestionAB}{%
\begin{center}\begin{minipage}{.97\linewidth}\centering
\begin{tikzpicture}[>=Latex,scale=.9]

\draw[kaxis] (-.25,0)--(6.6,0) node[right]{$t/\mathrm s$};
\draw[kaxis] (0,-2.7)--(0,1.2) node[above]{$x/\mathrm m$};
\foreach \t in {1,...,6}{\draw (\t,0)--(\t,.08);\node[above=2pt] at (\t,0){\t};}
\foreach \y/\v in {-2.4/-60,-1.6/-40,-.8/-20,.8/20}{\draw (0,\y)--(.08,\y);\node[left] at (0,\y){$\v$};}
\node[below left] at (0,0){$O$};

\end{tikzpicture}
\figcap{题图：按表中数据在空坐标系中作图。}\label{ch02:ref:question-figure-2-1}
\end{minipage}\end{center}
}

\newcommand{\KQuestionAC}{%
\begin{center}\begin{minipage}{.97\linewidth}\centering
\begin{tikzpicture}[>=Latex,scale=.85]

\foreach \dx/\i in {0/1,4.3/2,8.6/3}{\begin{scope}[xshift=\dx cm]
\draw[kaxis] (-.2,0)--(3.1,0) node[right]{$t$};
\draw[kaxis] (0,-.6)--(0,1.8) node[above]{$x$};
\node[below left] at (0,0){$O$};\node at (1.4,-.85){(\i)};
\end{scope}}

\end{tikzpicture}
\figcap{题图：分别画出三个运动函数的图像。}\label{ch02:ref:question-figure-2-2}
\end{minipage}\end{center}
}

\newcommand{\KQuestionAD}{%
\begin{center}\begin{minipage}{.97\linewidth}\centering
\begin{tikzpicture}[>=Latex,scale=.9]

\foreach \dx/\mode in {0/1,6.2/2}{\begin{scope}[xshift=\dx cm]
\draw[kaxis] (-.2,0)--(4.6,0) node[right]{$t$};
\draw[kaxis] (0,-1.25)--(0,1.3) node[above]{$x$};
\ifnum\mode=1
\draw[kcurve,domain=0:4.25,samples=200] plot(\x,{.8*sin(280*\x)});
\else
\draw[kcurve,domain=0:4.25,samples=240] plot(\x,{1.05*exp(-.85*\x)*cos(450*\x)});
\fi
\node[below left] at (0,0){$O$};\node at (2,-1.6){(\mode)};
\end{scope}}

\end{tikzpicture}
\figcap{题图：两个质点的一维位置随时间变化，坐标轴未给定数值比例。}\label{ch02:ref:question-figure-2-3}
\end{minipage}\end{center}
}

\newcommand{\KQuestionAG}{%
%
\begin{center}\begin{tikzpicture}[scale=.78,>=Latex]
  \foreach \dx/\mode in {-4.6/1,0/2,4.6/3}{
    \begin{scope}[xshift=\dx cm]
      \draw[kaxis] (-1.65,0)--(1.75,0) node[right,clearlabel]{$t$};
      \draw[kaxis] (-1.45,-.55)--(-1.45,2.05) node[above,clearlabel]{$x$};
      \ifnum\mode=1
        \draw[kcurve] (-1.35,0)--(0,1.35)--(1.35,0);
      \else\ifnum\mode=2
        \draw[kcurve,domain=-1.35:1.30,samples=80] plot (\x,{.25*(\x+1.35)^2});
      \else
        \draw[kcurve,domain=-1.35:1.35,samples=100] plot (\x,{1.45*sqrt(max(0,1-(\x/1.35)^2))});
      \fi\fi
    \end{scope}}
\end{tikzpicture}\figcap{题图：从左到右为三幅给定的位置—时间图像。}\label{ch02:ref:question-figure-2-6}\end{center}
}

\newcommand{\KQuestionAH}{%
%
\begin{center}\begin{tikzpicture}[scale=.9,>=Latex]
  \begin{scope}[shift={(-4.9,-1.1)}]
    \draw[kaxis] (0,0)--(2.75,0) node[right,clearlabel]{$t$};
    \draw[kaxis] (0,-1.55)--(0,1.85) node[above,clearlabel]{$v$};
    \draw[kcurve] (0,1.3)--(2.4,-1.3);
    \node[clearlabel,left] at (0,1.3){$v_0$};
    \node[clearlabel,left] at (0,-1.3){$-v_0$};
    \draw[kguide] (0,-1.3)--(2.4,-1.3);
    \node[clearlabel,below left] at (0,0){$O$};
    \node[clearlabel,above right=2pt] at (1.2,0){$T/2$};
    \node[clearlabel,above] at (2.4,0){$T$};
    \node[clearlabel,below] at (1.2,-1.7){已知速度图};
  \end{scope}
  \foreach \dx/\dy/\mode/\tag in {-.7/0/1/A,3.3/0/2/B,-.7/-3.3/3/C,3.3/-3.3/4/D}{
    \begin{scope}[shift={(\dx,\dy)}]
      \draw[kaxis] (0,0)--(2.7,0) node[right,clearlabel]{$t$};
      \draw[kaxis] (0,-.25)--(0,1.65) node[above,clearlabel]{$x$};
      \ifnum\mode=1
        \draw[kcurve] (0,0)--(1.2,1.15)--(2.4,0);
      \else\ifnum\mode=2
        \draw[kcurve,domain=0:1,samples=70]
          plot ({2.4*\x},{4.6*\x*(1-\x)});
      \else\ifnum\mode=3
        \draw[kcurve,domain=0:1,samples=70] plot ({2.4*\x},{1.15*(2*\x-\x*\x)});
      \else
        \draw[kcurve] (0,1.15)--(1.2,0)--(2.4,.85);
      \fi\fi\fi
      \node[clearlabel,below left] at (0,0){$O$};
      \node[clearlabel,below=3pt] at (1.2,0){$T/2$};
      \node[clearlabel,below=3pt] at (2.4,0){$T$};
      \node[clearlabel,below] at (1.2,-.65){$\tag$};
    \end{scope}}
\end{tikzpicture}\figcap{题图：左图给出速度随时间的变化，右侧 A—D 为待选的位置图像。}\label{ch02:ref:question-figure-2-7}\end{center}
}

\newcommand{\KQuestionAI}{%
%
\begin{center}\begin{tikzpicture}[scale=.78,>=Latex]
  \foreach \dx/\mode in {-4.4/1,0/2,4.4/3}{
    \begin{scope}[xshift=\dx cm]
      \draw[kaxis] (-1.55,0)--(1.75,0) node[right,clearlabel]{时间};
      \draw[kaxis] (-1.35,-.25)--(-1.35,1.85) node[above,clearlabel]{$v$};
      \ifnum\mode=1
        \draw[kcurve] (-1.25,0)--(1.25,1.35);
        \draw[kguide] (1.25,0)--(1.25,1.35);
        \draw[kguide] (-1.35,1.35)--(1.25,1.35);
        \node[clearlabel,left] at (-1.35,1.35){$v$};
      \else\ifnum\mode=2
        \draw[kcurve] (-1.25,.55)--(1.25,1.35);
        \draw[kguide] (1.25,0)--(1.25,1.35);
        \draw[kguide] (-1.35,.55)--(-1.25,.55) (-1.35,1.35)--(1.25,1.35);
        \node[clearlabel,left] at (-1.35,.55){$v_0$};
        \node[clearlabel,left] at (-1.35,1.35){$v_1$};
      \else
        \draw[kcurve] (-1.25,0)--(-.42,1.35)--(1.25,0);
        \draw[kguide] (-.42,0)--(-.42,1.35);
        \draw[kguide] (-1.35,1.35)--(-.42,1.35);
        \node[clearlabel,left] at (-1.35,1.35){$v_0$};
        \node[clearlabel,below] at (-.42,0){$t/3$};
      \fi\fi
      \node[clearlabel,below] at (1.25,0){$t$};
    \end{scope}}
\end{tikzpicture}\figcap{题图：三个运动过程的速度—时间图像，求相应时间内的路程。}\label{ch02:ref:question-figure-2-8}\end{center}
}

\newcommand{\KQuestionAR}{%
%
\begin{center}\begin{tikzpicture}[scale=.9,>=Latex]
  \draw[very thick] (-1.4,0)--(1.4,0);
  \node[point,label={[clearlabel]below right=3pt:$B$}] (B) at (0,0) {};
  \node[point,label={[clearlabel]left:$A$}] (A) at (0,3.2) {};
  \draw[<->,TextGray] (.8,0)--(.8,3.2) node[midway,clearlabel,right]{$h$};
  \draw[DeepBlue,very thick,->] (B)--(0,1.15) node[midway,clearlabel,left]{$v_0$};
  \draw[Blue,very thick,->] (A)--(0,2.45) node[midway,clearlabel,left]{$g$};
\end{tikzpicture}\figcap{题图：A 球从上方释放，B 球从下方以初速度 $v_0$ 竖直抛出，初始间距为 $h$。}\label{ch02:ref:question-figure-2-17}\end{center}
}

\newcommand{\KQuestionAT}{%
%
\begin{center}\begin{tikzpicture}[scale=.85,>=Latex]
  \draw[very thick] (-1.7,0)--(1.7,0);
  \fill (0,2.7) circle (.12);
  \fill (0,.12) circle (.12);
  \draw[kguide] (0,.27)--(0,2.55);
  \draw[<->,TextGray] (.88,0)--(.88,2.7)
    node[midway,clearlabel,right=2pt]{$H$};
  \draw[DeepBlue,very thick,->] (-.35,1.0)--(-.35,.25) node[midway,clearlabel,left]{$u$};
  \draw[Blue,very thick,->] (.35,.25)--(.35,1.0) node[midway,clearlabel,right]{$eu$};
\end{tikzpicture}\figcap{题图：初始高度为 $H$；一次碰撞前、后的速率分别用 $u$、$eu$ 表示。}\label{ch02:ref:question-figure-2-19}\end{center}
}


把运动物体抽象为质点以后，质点的位置代表整个物体的位置。我们从一维直线运动开始：在运动直线上取原点、正方向和长度单位，建立数轴 $Ox$。质点在时刻 $t$ 的位置由坐标 $x(t)$ 唯一确定。坐标的符号表示点位于原点的哪一侧，并不表示“路程为负”。

这种简化只是描写运动的第一步。即使研究对象已经缩成一个点，要完整说明它怎样运动，仍须知道它在每一时刻处在什么位置。先从直线运动建立方法，再推广到曲线运动和真实物体的复杂运动，是本章的基本顺序；等运动本身被描述清楚之后，才进一步追问支配运动变化的原因。

建立一维坐标系时，原点、正方向和长度单位都可按问题需要选择，通常以米或千米为长度单位。质点在某时刻的坐标又称它在该坐标系中的位置。改变原点或正方向会改变坐标的数值或符号，却不会改变实际运动。若已知 $t=0,1,2,\ldots$ 各时刻的坐标，就已掌握一部分运动信息；记录时刻越密，所知道的运动细节越多，但只有连续的函数才能原则上给出每一时刻的位置。

描写同一运动可以采用叙述、列表、图像和函数四种方式。叙述便于日常交流；列表适合记录实验数据；图像最直观；函数则能把连续运动在任意时刻的位置完整表达出来。以下短跑数据同时说明四种表达之间的联系。

若用叙述表达表~\ref{ch02:tab:k02-2-1}中的比赛过程，就要依次说出“第 $1\,\mathrm{s}$ 跑到 $1\,\mathrm{m}$、第 $2\,\mathrm{s}$ 跑到 $4\,\mathrm{m}$、……、第 $12\,\mathrm{s}$ 到达 $100\,\mathrm{m}$”。这种说法能够传递信息，却繁琐且不容易立即比较各阶段的快慢。列表节省空间，适合实验记录，但仍不能一眼看出变化趋势。把时间作为横轴、离开起点的距离作为纵轴，数据点连成图~\ref{ch02:fig:k02-2-2}，比赛开始阶段曲线越来越陡，说明选手在加速；中段斜率近似不变；终点前曲线又略变陡，可看出冲刺过程。

\begin{center}\small
\begin{tabular}{c*{12}{r}}
\toprule
$t/\mathrm{s}$&1&2&3&4&5&6&7&8&9&10&11&12\\
$x/\mathrm{m}$&1&4&9&16&25&35&45&55&65&75&85&100\\
\bottomrule
\end{tabular}
\booktablecaption{不同时刻短跑选手的位置。}\label{ch02:tab:k02-2-1}
\end{center}

\SprintGraph

若函数 $x=x(t)$ 已知，就能在定义域内求任意时刻的位置。图像并不是质点在空间中的轨迹，而是“位置”这个物理量随“时间”变化的关系图；横轴和纵轴的物理意义不能混淆。

从实验得到函数通常要经历“测量—作图—猜测规律—检验”的过程。有限个数据点本身并不能保证点间一定按某条光滑曲线运动；把点连起来，实际上加入了质点连续运动以及测量误差较小等物理判断。若在某段出现突然跳点，应先检查是否发生碰撞、换向或记录错误，而不能为了图线好看强行作光滑连接。反过来，给定解析函数后可以任意加密取点，所画曲线只是函数关系的可视化。

还要区分“时刻”和“时间间隔”。$t$ 是相对所选计时零点的读数，改变计时起点会改变它的数值；$\Delta t=t_2-t_1$ 是两个事件间隔，通常不随计时零点改变。位置 $x$ 依赖坐标原点，位置变化 $\Delta x=x_2-x_1$ 则与平移坐标原点无关。许多运动公式实际联系的是间隔和变化量，列式前把初始时刻平移到 $t=0$ 只是方便，并非物理规律只在零时刻开始成立。

叙述、表格和有限个实验点都不能完全表现运动的连续性。若在两个已记录时刻之间质点仍连续运动，就还存在无穷多个没有列出的时刻。函数表达式
\begin{equation}
  x=x(t)
\end{equation}
把任意允许时刻映射为一个位置，因此最适合精确分析。例如 $x=3t$ 表示相等时间内坐标增加量相等；$x=5t^2$ 的图线随时间变陡，表示相同时间内的坐标增加量逐渐增大。函数和图像各有所长：函数便于定量运算，图像便于观察单调性、转折、周期和极值。

\begin{origexamplebox}{由测量数据作位置图像}\label{ch02:ex:k02-2-1}
建立竖直向上的坐标系表示从原点竖直上抛物体的位置。实验在 $t=0,0.5,1,\ldots,6\,\mathrm{s}$ 测得位置依次为
\[
0,\ 8.75,\ 15,\ 18.75,\ 20,\ 18.75,\ 15,\ 8.75,\ 0,\ -11.25,\ -25,\ -41.25,\ -60\ \mathrm{m}.
\]
在坐标系中标出数据并连接成光滑曲线。
\par 题图见图~\ref{ch02:ref:question-figure-2-1}。
\KQuestionAB
\end{origexamplebox}
\Needspace{6\baselineskip}\begin{xesolution}
\textbf{第一步：确定坐标轴和比例。}
横轴表示时刻 $t$，范围取 $0$ 到 $6\,\mathrm{s}$；纵轴表示相对所选原点的竖直坐标 $x$，
范围至少要覆盖 $-60\,\mathrm{m}$ 到 $20\,\mathrm{m}$。把题中每一对 $(t,x)$ 作为一个点标在
$x$--$t$ 坐标系中。这里画的是“位置随时间变化”的图像，不是物体在竖直空间中走出的轨迹。

\textbf{第二步：由数据判断最简单的函数形式。}
相邻测量时刻的间隔均为
\[
 h=0.5\,\mathrm{s}.
\]
相邻位置的一阶差分依次为
\[
\begin{aligned}
\Delta x_i={}&8.75,\ 6.25,\ 3.75,\ 1.25,\ -1.25,\ -3.75,\\
              &-6.25,\ -8.75,\ -11.25,\ -13.75,\ -16.25,\ -18.75
              \quad(\mathrm m).
\end{aligned}
\]
这些数每次都减少 $2.50\,\mathrm m$，所以二阶差分为常量
\[
 \Delta^2x_i=x_{i+2}-2x_{i+1}+x_i=-2.50\,\mathrm m.
\]
等时间间隔数据具有恒定二阶差分，最简单的光滑模型是二次函数
\[
 x(t)=At^2+Bt+C.
\]
对二次函数，步长为 $h$ 时有
\[
 x(t+2h)-2x(t+h)+x(t)=2Ah^2.
\]
代入 $h=0.5\,\mathrm s$ 和二阶差分 $-2.50\,\mathrm m$，得到
\[
 -2.50\,\mathrm m=2A(0.5\,\mathrm s)^2,
 \qquad A=-5.0\,\mathrm{m/s^2}.
\]
又由 $x(0)=0$ 得 $C=0$；把 $x(0.5\,\mathrm s)=8.75\,\mathrm m$ 代入，
\[
 8.75\,\mathrm m
 =(-5.0\,\mathrm{m/s^2})(0.5\,\mathrm s)^2
 +B(0.5\,\mathrm s),
\]
解得 $B=20\,\mathrm{m/s}$。因此所有测量点都满足
\[
 \boxed{x(t)=(20\,\mathrm{m/s})t-(5.0\,\mathrm{m/s^2})t^2}.
\]

\textbf{第三步：确定图像的关键点和运动含义。}
把二次函数配方：
\[
 x(t)=20\,\mathrm m
 -(5.0\,\mathrm{m/s^2})(t-2.0\,\mathrm s)^2.
\]
所以图像是开口向下的抛物线，对称轴为 $t=2.0\,\mathrm s$，顶点为
$(2.0\,\mathrm s,20\,\mathrm m)$。令 $x=0$，得到 $t=0$ 和 $t=4.0\,\mathrm s$：物体
在 $4.0\,\mathrm s$ 时回到坐标原点。$t>4.0\,\mathrm s$ 后 $x<0$，表示物体位于所选原点下方，
并不是“路程为负”。

图线斜率给出速度。由
\[
 v(t)=\frac{\dd x}{\dd t}
 =20\,\mathrm{m/s}-(10\,\mathrm{m/s^2})t
\]
可见：$0<t<2\,\mathrm s$ 时斜率为正，物体上升；$t=2\,\mathrm s$ 时斜率为零，到达最高点；
$t>2\,\mathrm s$ 时斜率为负，物体下降。加速度恒为
$a=-10\,\mathrm{m/s^2}$，与竖直向上为正的坐标约定一致。量纲检查中，两项
$(20\,\mathrm{m/s})t$ 与 $(5.0\,\mathrm{m/s^2})t^2$ 都是长度。

最后需要说明：有限个测量点本身不能唯一确定点与点之间的真实运动。这里所有点具有严格恒定的二阶差分，
因此采用匀加速度对应的二次函数作最简单的光滑连接；若没有这一物理建模假设，还可以画出无穷多条经过同一组数据点的光滑曲线。
\begin{center}\begin{tikzpicture}
\begin{axis}[width=.72\linewidth,height=5.8cm,axis lines=middle,
  xlabel={$t/\mathrm{s}$},ylabel={$x/\mathrm{m}$},
  xlabel style={at={(axis description cs:1.01,.49)},anchor=west},
  ylabel style={rotate=0,at={(axis description cs:.035,1.02)},anchor=south west},
  xmin=0,xmax=6.2,ymin=-65,ymax=25,grid=major,grid style={draw=RuleGray},
  every axis label/.style={fill=none}]
\addplot[kcurve,domain=0:6,samples=100]{20*x-5*x^2};
\addplot[only marks,mark=*,mark size=1.5pt,DeepBlue] coordinates{(0,0)(.5,8.75)(1,15)(1.5,18.75)(2,20)(2.5,18.75)(3,15)(3.5,8.75)(4,0)(4.5,-11.25)(5,-25)(5.5,-41.25)(6,-60)};
\addplot[guide] coordinates{(2,-65)(2,20)};
\node[anchor=south west,DeepBlue] at (axis cs:2,20) {最高点};
\end{axis}\end{tikzpicture}\end{center}
\end{xesolution}

\begin{origexamplebox}{三种运动函数的图像}\label{ch02:ex:k02-2-2}
在坐标系中分别画出
\[
  (1)\ x=5t^2;\qquad (2)\ x=\frac{3}{t+1};\qquad
  (3)\ x=5\sin(2\pi t).
\]
\par 题图见图~\ref{ch02:ref:question-figure-2-2}。
\KQuestionAC
\end{origexamplebox}
\Needspace{6\baselineskip}\begin{xesolution}
作函数图像不能只凭“形状印象”，应先确定定义域、关键点和变化趋势。若 $t$ 表示从计时开始后的时间，物理上通常只取 $t\ge0$；下面同时说明函数在完整数学定义域上的特征。

\textbf{（1）$x=5t^2$。}
它是开口向上的抛物线，顶点为 $(0,0)$。在 $t\ge0$ 上，
\[
 \frac{\dd x}{\dd t}=10t\ge0,
\]
所以位置随时间单调增加；且斜率 $10t$ 越来越大，图线越来越陡。可先取
$(0,0),(0.5,1.25),(1,5)$ 等点，再用光滑曲线连接。

\textbf{（2）$x=3/(t+1)$。}
完整定义域不含 $t=-1$。当 $t\to-1^+$ 时 $x\to+\infty$，当
$t\to-1^-$ 时 $x\to-\infty$，故 $t=-1$ 是竖直渐近线；当
$t\to\pm\infty$ 时 $x\to0$，故 $x=0$ 是水平渐近线。若只画物理区间
$t\ge0$，则从 $(0,3)$ 出发，并因
\[
 \frac{\dd x}{\dd t}=-\frac{3}{(t+1)^2}<0
\]
而单调下降，逐渐逼近但不会到达 $x=0$。

\textbf{（3）$x=5\sin(2\pi t)$。}
振幅为 $5$，角频率为 $2\pi\,\mathrm{rad/s}$，所以周期
\[
 T=\frac{2\pi}{2\pi\,\mathrm{rad/s}}=1\,\mathrm s.
\]
一个周期内的关键点为
\[
 (0,0),\quad (0.25\,\mathrm s,5),\quad (0.50\,\mathrm s,0),
 \quad (0.75\,\mathrm s,-5),\quad (1.00\,\mathrm s,0).
\]
按这些点画出正弦曲线，再以 $1\,\mathrm s$ 为周期重复。三幅图的纵轴都是位置 $x$，曲线形状不是质点在空间中的轨迹。
\begin{center}\begin{tikzpicture}
\foreach \dx/\mode in {-4.6/1,0/2,4.6/3}{
\begin{scope}[xshift=\dx cm]
\draw[kaxis] (-1.7,0)--(1.75,0) node[right,clearlabel]{$t$};
\draw[kaxis] (-1.45,-1.35)--(-1.45,1.5) node[above,clearlabel]{$x$};
\ifnum\mode=1 \draw[kcurve] plot[domain=0:1.55,samples=60] ({\x-1.45},{.45*\x*\x});\node[clearlabel,below] at (0,-1.6){$5t^2$};\fi
\ifnum\mode=2 \draw[kcurve] plot[domain=0:3.1,samples=80] ({\x-1.45},{1.15/(\x+1)});\node[clearlabel,below] at (0,-1.6){$3/(t+1)\ (t\ge0)$};\fi
\ifnum\mode=3 \draw[kcurve] plot[domain=0:3.1,samples=100] ({\x-1.45},{.85*sin(360*\x)});\node[clearlabel,below] at (0,-1.6){$5\sin2\pi t$};\fi
\end{scope}}
\end{tikzpicture}\end{center}
\end{xesolution}

\begin{origexamplebox}{由位置图判断运动特征}\label{ch02:ex:k02-2-3}
下图分别给出两个质点的一维位置随时间的变化。指出各自运动最明显的特征。
\par 题图见图~\ref{ch02:ref:question-figure-2-3}。
\KQuestionAD
\end{origexamplebox}
\Needspace{6\baselineskip}\begin{xesolution}
先明确读图规则：纵坐标给位置 $x$，曲线斜率给速度 $v=\dd x/\dd t$，极大值和极小值之间的竖直距离不能直接当成路程。

\textbf{第一图。}相邻同向过零点之间的时间间隔不变，说明周期 $T$ 不变；正、负极值到平衡位置 $x=0$ 的距离相同且不随时间改变，说明振幅不变。每到极值点，切线水平，故 $v=0$，质点在这一瞬间换向；经过 $x=0$ 附近时图线最陡，故速度大小最大。一个周期内，质点从一侧极值出发，经过平衡位置，到达另一侧极值，再返回原位置。

\textbf{第二图。}质点仍反复经过 $x=0$ 并在正、负两侧换向，但各次极值的绝对值逐渐减小，所以振幅随时间衰减。若相邻同向过零点的间隔近似相同，只能说周期近似不变；仅凭示意图不能读出精确周期。极值越来越接近零，表示运动逐渐趋向平衡位置，而不是“波形沿水平方向传播”。

两幅图都是同一质点的 $x$--$t$ 图：横向代表时间流逝，不能把图线误认为质点在空间中走出的波浪形轨迹。若坐标轴没有给比例，还只能判断上述定性特征，不能计算具体速度数值。
题图~\ref{ch02:ref:question-figure-2-3}中，左图振幅保持不变；右图各次极值逐渐靠近平衡位置。
\end{xesolution}

上述方法也能描写沿弯曲公路行驶的汽车。只需在公路上选定里程原点和前进的正方向，用沿公路量得的里程标记位置；我国公路上的里程标志就是这种一维坐标思想的实际应用。此时坐标沿道路累计，并不等于汽车到原点的直线距离。

\subsection{速率与速度}

路程 $s$ 是质点沿轨迹实际走过的长度。匀速运动中，速率可写成 $v=s/t$。一般运动在时间区间 $[t,t+\Delta t]$ 内的平均速率为
\begin{equation}
  \bar v_{\mathrm{speed}}=\frac{\Delta s}{\Delta t},
\end{equation}
瞬时速率是时间间隔趋于零时的极限：
\begin{equation}
  v_{\mathrm{speed}}(t)=\lim_{\Delta t\to0}\frac{\Delta s}{\Delta t}.
\end{equation}
在极短时间内，路程可近似为 $\Delta s\simeq v_{\mathrm{speed}}\Delta t$；时间间隔越小，忽略的高阶项越小。

速率是运动快慢的量度，可以把汽车速率表的读数看成直观例子。对匀速运动，用总路程除以总时间即可得到任何时段相同的速率；对变速运动，同一个过程取不同时间段会得到不同的平均速率，所以必须进一步引入瞬时速率。以 $x=5t^2$ 且质点在所考察时段始终沿正方向运动为例，从 $t=1\,\mathrm{s}$ 起选取不同的 $\Delta t$，有
\[
\begin{array}{c|ccccc}
\Delta t/\mathrm s&1&0.1&0.01&0.001&10^{-n}\\\toprule
\Delta s/\Delta t\,/\mathrm{(m/s)}&15&10.5&10.05&10.005&10+5\times10^{-n}
\end{array}
\]
当 $\Delta t\to0$ 时，这些平均值趋于 $10\,\mathrm{m/s}$，即 $t=1\,\mathrm{s}$ 的瞬时速率。这个极限并不是说实际测量要取“等于零的时间”，而是说明时间窗越小，区间平均量越逼近该时刻的运动状态。

同一例子也说明 $\Delta s\simeq v\Delta t$ 的近似。取 $t=1\,\mathrm{s}$ 时的瞬时速率 $10\,\mathrm{m/s}$，当 $\Delta t=1,0.1,0.01\,\mathrm{s}$ 时，线性估计的路程依次为 $10,1,0.1\,\mathrm{m}$，真实路程依次为 $15,1.05,0.1005\,\mathrm{m}$；时间窗缩小后，两者差别迅速减小。数学上这正是切线对曲线的局部近似。

一维坐标轴已经规定了正方向。为了同时表达快慢和方向，引入有正负号的速度
\begin{equation}
  \bar v=\frac{\Delta x}{\Delta t},\qquad
  v(t)=\lim_{\Delta t\to0}\frac{\Delta x}{\Delta t}=\frac{\dd x}{\dd t}.
\end{equation}
瞬时速率等于 $|v|$。若速度在区间内恒定，则
\begin{equation}
  x=x_0+v\Delta t.
\end{equation}

例如 $x=5t$ 的平均速度与所取区间无关，始终为 $5\,\mathrm{m/s}$；而 $x=5t^2$ 在 $t=1\,\mathrm{s}$ 附近的平均速度为 $10+5\Delta t$，当 $\Delta t\to0$ 时趋于 $10\,\mathrm{m/s}$。这正是从平均量过渡到瞬时量的极限思想。

必须区分路程、位移、速率和速度。路程记录轨迹累计长度，始终非负；位移是末位置减初位置，在一维中可正可负。速率只表示快慢，速度还包含方向。若质点先沿正方向走一段再返回，路程继续增加，而位移可能减小甚至回到零；相应地，平均速率一般不等于平均速度的绝对值。式 $x=x_0+v\Delta t$ 对正、负速度都成立，符号已经把方向包含在内。

例如质点从 $x=0$ 走到 $x=10\,\mathrm m$，再返回 $x=4\,\mathrm m$，全过程路程为 $16\,\mathrm m$，位移却只有 $+4\,\mathrm m$。若用时 $4\,\mathrm s$，平均速率是 $4\,\mathrm{m/s}$，平均速度是 $+1\,\mathrm{m/s}$。只有运动过程中不换向时，路程才等于位移大小，平均速率才等于平均速度的绝对值。判断是否换向应看速度符号，而不能只看位置始终为正还是为负。

瞬时速度为零也不等于质点将在此后保持静止。位置图像在极值点切线水平，表示该时刻速度为零；若曲线随后斜率改变符号，质点正是在这一瞬间换向。相反，若位置图在一段时间内保持水平，才表示质点在这一整段都静止。这一区别在竖直上抛最高点和往复振动中尤其重要。

\begin{origexamplebox}{相邻测量间的平均速率}\label{ch02:ex:k02-2-4}
表中 $t=0,0.5,\ldots,6\,\mathrm{s}$ 时的位置分别为
\[
0,8.75,15,18.75,20,18.75,15,8.75,0,-11.25,-25,-41.25,-60\ \mathrm{m}.
\]
求相邻两次测量之间的平均速率。
\end{origexamplebox}
\Needspace{6\baselineskip}\begin{xesolution}
相邻两次测量的时间间隔均为
\[
 \Delta t=0.5\,\mathrm s.
\]
先计算每一小段的位移 $\Delta x=x_{i+1}-x_i$，再除以 $\Delta t$。例如前两段为
\[
 \bar v_{0\sim0.5}=\frac{8.75-0}{0.5}=17.5\,\mathrm{m/s},
 \qquad
 \bar v_{0.5\sim1}=\frac{15-8.75}{0.5}=12.5\,\mathrm{m/s}.
\]
依次计算得到有符号平均速度
\[
17.5,12.5,7.5,2.5,-2.5,-7.5,-12.5,-17.5,-22.5,-27.5,-32.5,-37.5\ \mathrm{m/s}.
\]
由例~\ref{ch02:ex:k02-2-1}的抛物线可知，物体在每个相邻测量区间内都不发生额外换向，唯一的换向点恰在区间端点 $t=2\,\mathrm s$。因此各段路程等于 $|\Delta x|$，平均速率就是上述数值的绝对值：
\[
\boxed{17.5,12.5,7.5,2.5,2.5,7.5,12.5,17.5,22.5,27.5,32.5,37.5\ \mathrm{m/s}}.
\]
$t=2\,\mathrm{s}$ 前平均速度为正，之后为负；负号表示沿坐标轴负方向运动，不表示运动“比静止更慢”。还要注意：若只知道相邻时刻的位置，却不知道区间内是否换向，一般只能求平均速度，不能仅凭端点确定平均速率；本题能这样做，是因为前一例已经确定了区间内的运动模型。
\end{xesolution}

\begin{origexamplebox}{相向运动}\label{ch02:ex:k02-2-5}
质点 $A,B$ 在同一直线上匀速相向运动。初始位置 $x_A=-5\,\mathrm m$、$x_B=10\,\mathrm m$；$A$ 沿正方向以 $1.5\,\mathrm{m/s}$ 运动，$B$ 沿负方向的速率为 $0.5\,\mathrm{m/s}$。求相遇时间和位置。
\end{origexamplebox}
\Needspace{6\baselineskip}\begin{xesolution}
取题目给定的正方向，令 $t=0$ 为两质点处在初始位置的时刻。匀速直线运动满足“位置等于初位置加上速度乘时间”，因此
\[
x_A(t)=-5\,\mathrm m+(1.5\,\mathrm{m/s})t,
\qquad
x_B(t)=10\,\mathrm m-(0.5\,\mathrm{m/s})t.
\]
相遇的物理条件是同一时刻坐标相同：
\[
-5+1.5t=10-0.5t.
\]
整理得 $(2.0\,\mathrm{m/s})t=15\,\mathrm m$，所以
\[
\boxed{t=7.5\,\mathrm s},\qquad
\boxed{x=6.25\,\mathrm m}.
\]
其中相遇位置可分别代回两式检查：
\[
x_A=-5+1.5\times7.5=6.25\,\mathrm m,
\qquad
x_B=10-0.5\times7.5=6.25\,\mathrm m.
\]
再用相对运动复核：初始间距为 $15\,\mathrm m$，接近速率为
$1.5+0.5=2.0\,\mathrm{m/s}$，故闭合时间为 $15/2.0=7.5\,\mathrm s$，与坐标法一致。时间量纲也由“长度除以速度”正确得到。
\end{xesolution}

\subsection{图像的斜率与面积}

\PositionSlopeDiagram

$x$-$t$ 图像中两点割线斜率是平均速度，切线斜率是瞬时速度。图线向上并不一定表示质点在空间中“向上”，它只表示坐标随时间增大；具体方向由坐标轴的正方向决定。

图~\ref{ch02:fig:k02-2-3}把这一点作了几何化说明。匀速运动的位置图像是一条直线，任意两点割线都与原直线重合，斜率恒定。一般运动中，时刻 $t$ 与 $t+\Delta t$ 对应曲线上的 $A,B$ 两点，割线 $AB$ 的斜率是 $\Delta x/\Delta t$。让 $\Delta t$ 逐渐缩小，$B$ 沿曲线趋近 $A$，割线趋近于 $A$ 点切线，其斜率的极限就是瞬时速度。因而位置图像的几何信息可以直接翻译成速度：斜率为正表示沿正方向，斜率为负表示沿负方向，斜率为零表示该时刻瞬时静止。

\begin{origexamplebox}{从 $x$-$t$ 图判断速度趋势}\label{ch02:ex:k02-2-6}
根据下列三幅位置—时间图像，分别判断质点的运动方向和速度大小怎样变化。
\par 题图见图~\ref{ch02:ref:question-figure-2-6}。
\KQuestionAG
\end{origexamplebox}
\Needspace{6\baselineskip}\begin{xesolution}
沿横轴箭头方向读取时间，并在每一段比较切线斜率 $v=\dd x/\dd t$。


\ExampleTwoSixGraphs
\textbf{第一图：先线性上升后线性下降。}上升直线的斜率为正常量，表示质点先以恒定正速度运动；下降直线的斜率为负常量，表示随后以恒定负速度运动。折点处左右斜率不同，理想模型中速度在同一瞬间跳变，普通加速度不存在；真实换向若由有限力造成，折角附近应是一小段光滑曲线。

\textbf{第二图：向上凹且越来越陡。}任一点切线斜率均为正，所以质点始终沿正方向运动；斜率随时间增大，所以正速度增大，因而加速度为正。这里“图线向上弯”是对 $x$--$t$ 曲线凹凸性的描述，不代表质点在空间中沿弯曲轨迹运动。

\textbf{第三图：上半圆形。}左半段斜率为正，但逐渐减小到零，表示质点沿正方向运动并逐渐减速；最高点切线水平，故该时刻 $v=0$；右半段斜率为负且绝对值逐渐增大，表示质点换向后沿负方向越来越快。若把理想半圆写成
\[
 \left(\frac{t-t_c}{\tau}\right)^2+
 \left(\frac{x-x_c}{X}\right)^2=1,\qquad x\ge x_c,
\]
则
\[
 v=\frac{\dd x}{\dd t}
   =-\frac X\tau\frac{(t-t_c)/\tau}{\sqrt{1-[(t-t_c)/\tau]^2}}.
\]
这里 $\tau>0$ 是时间尺度，$X>0$ 是长度尺度，不能把时间平方与位置平方直接相加。
图中“半圆”指归一化坐标 $((t-t_c)/\tau,(x-x_c)/X)$ 下的几何形状；
实际位置为 $x=x_c+X\sqrt{1-[(t-t_c)/\tau]^2}$，求导后含有速度尺度 $X/\tau$。
这验证了顶点 $t=t_c$ 时 $v=0$，而两端 $t=t_c\pm\tau$ 的竖直切线处 $|v|\to\infty$。
无限速度没有现实意义，说明把一整段真实位置图精确画成半圆通常只是一种几何理想化。
\end{xesolution}

速度也可以画成 $v$-$t$ 图。时间区间内的位移为
\begin{equation}
  \Delta x=\int_{t_1}^{t_2}v(t)\,\dd t,
\end{equation}
即曲线与时间轴围成的\textbf{有向面积}；路程则为 $\int|v|\dd t$。

图~\ref{ch02:fig:k02-2-4}还说明了面积结论的来由。图中 $A,B$ 表示时间区间的两个端点，$C,D$ 是相应时刻速度曲线上的点；因此封闭区域 $ABDC$ 正是这一时间段的累计位移。匀速时，$v$-$t$ 图是一条平行于时间轴的直线，矩形面积 $v\Delta t$ 就是位移。变速时，把时间区间分成许多很窄的小段，每段内用左端或某个时刻的瞬时速度近似不变，小矩形面积近似该小段位移。分割越细，所有小矩形面积之和越接近曲线与时间轴围成的面积，极限就是积分。曲线位于时间轴下方时，速度为负，对位移贡献负面积；计算总路程则要把这部分面积取正。

\VelocityAreaDiagram

\begin{origexamplebox}{速度图与位移图}\label{ch02:ex:k02-2-7}
速度从正值随时间线性减小，在 $T/2$ 时变为零，$T$ 时变为等大的负值。判断从 $x(0)=0$ 开始的位置图像。
\par 题图见图~\ref{ch02:ref:question-figure-2-7}。
\KQuestionAH
\end{origexamplebox}
\Needspace{6\baselineskip}\begin{xesolution}
设初始速度为 $v_0>0$，末速度为 $-v_0$。速度随时间线性变化，所以连接 $(0,v_0)$ 与 $(T,-v_0)$ 的直线方程是
\[
 v(t)=v_0-\frac{2v_0}{T}t.
\]
\ExampleTwoSevenGraphs

其斜率即加速度，
\[
 a=-\frac{2v_0}{T}<0,
\]
为常量。因此位置由积分得到
\[
 x(t)-x(0)=\int_0^t v(t')\,\dd t'
 =v_0t-\frac{v_0}{T}t^2.
\]
又因 $x(0)=0$，
\[
 \boxed{x(t)=v_0t-\frac{v_0}{T}t^2},
\]
这是开口向下的抛物线。它在 $t=T/2$ 处斜率为零并达到最大位置
\[
 x_{\max}=x(T/2)=\frac{v_0T}{4},
\]
在 $t=T$ 时回到 $x(T)=0$。

同一结论也可由 $v$--$t$ 图的有向面积看出：$0$ 到 $T/2$ 的正三角形面积与 $T/2$ 到 $T$ 的负三角形面积大小相等，故总位移为零。右侧四图中应选光滑的抛物线拱形 B 图；折线 A、D 对应速度突变，C 图的末位置不回到零。式中 $v_0t$ 与 $(v_0/T)t^2$ 都是长度，量纲一致。
\end{xesolution}

\begin{origexamplebox}{由 $v$-$t$ 图求路程}\label{ch02:ex:k02-2-8}
求三种运动在 $0$ 到 $t$ 内的路程：（1）速度从 $0$ 线性增至 $v$；（2）速度从 $v_0$ 线性增至 $v_1$；（3）速度从 $0$ 线性增至 $v_0$，在 $t/3$ 达到峰值，再线性降到 $0$。
\par 题图见图~\ref{ch02:ref:question-figure-2-8}。
\KQuestionAI
\end{origexamplebox}
\Needspace{6\baselineskip}\begin{xesolution}
三幅 $v$--$t$ 图都位于时间轴上方，故速度始终非负，路程与位移相等，均可直接用图线下的面积计算。


\ExampleTwoEightGraphs
\textbf{（1）}图形是底为 $t$、高为 $v$ 的三角形，
\[
 s_1=\frac12tv.
\]

\textbf{（2）}图形是两条平行边长分别为 $v_0,v_1$、高为 $t$ 的梯形，
\[
 s_2=\frac{v_0+v_1}{2}t.
\]

\textbf{（3）}峰值左、右两部分分别是底为 $t/3$ 与 $2t/3$、共同高为 $v_0$ 的三角形，
\[
 s_3=\frac12v_0\frac t3+\frac12v_0\frac{2t}{3}
 =\frac12v_0t.
\]
综上，
\[
  \boxed{s_1=\frac12vt},\qquad
  \boxed{s_2=\frac{v_0+v_1}{2}t},\qquad
  \boxed{s_3=\frac12v_0t}.
\]
第三问说明：只要速度从零线性升到同一峰值再线性降到零，峰值时刻只是在左右两个三角形之间重新分配底边，总底边仍为 $t$，所以总面积不变。各答案都是“速度乘时间”，量纲为长度。若某段图线落在时间轴下方，求位移要保留负面积，求路程则必须把该部分面积取正；本题没有这种情况。
\end{xesolution}

\subsection{加速度与匀变速直线运动}

速度变化量为 $\Delta v=v(t+\Delta t)-v(t)$。平均加速度和瞬时加速度分别为
\begin{equation}
  \bar a=\frac{\Delta v}{\Delta t},\qquad
  a(t)=\lim_{\Delta t\to0}\frac{\Delta v}{\Delta t}=\frac{\dd v}{\dd t}.
\end{equation}
加速度描述速度变化的快慢和方向。$a$ 与 $v$ 同号时速率增加，异号时速率减小；不能把“加速度为负”直接说成“减速”。

与速度的定义完全类似，若 $\Delta t$ 取一个有限区间，$\Delta v/\Delta t$ 是该区间的平均加速度；若让时间间隔趋于零，就得到某一时刻的瞬时加速度。相同时间内速度变化越大，加速度的绝对值越大。加速度的正负由坐标正方向决定，它表示速度在数轴上的变化方向，而不是单独表示“加快”或“减慢”。例如 $v<0,a<0$ 时，速度的绝对值增大，质点实际在加速向负方向运动。

速度图像为正只说明运动方向沿坐标正向，为负只说明沿负向；加速度图像同理表示速度的变化方向。判断速率增减，可以比较 $v$ 与 $a$ 的符号：同号时 $|v|$ 增大，异号时 $|v|$ 减小。若 $v=0$ 而 $a\ne0$，质点可能正处在换向的临界时刻。把这些符号关系与 $x$-$t$、$v$-$t$、$a$-$t$ 三类图像同时核对，能排除许多仅凭“图线向上或向下”产生的误解。

三个运动学量之间存在成对的局部与累计关系：位置图的切线斜率给速度，速度图的切线斜率给加速度；速度图的有向面积给位移，加速度图的有向面积给速度变化量。所谓有向面积必须保留横轴上下的正负。若只求路程，则须把速度负值部分翻到横轴上方后再累计，而不能在最后对总位移简单取绝对值。

实际测量的“瞬时值”总是在很短但有限的时间窗内取得的近似。速度表、光电门和高速摄影并不会真的用零秒作分母，而是让采样时间足够短，使这段时间内速度变化小于所要求的精度。数学极限给出理想定义，实验误差则决定时间窗能缩小到什么程度；时间窗过大不能代表瞬时状态，过小又可能放大计时和位置测量误差。

从函数观点看，$x(t)$ 可导时，$v(t)$ 就是其导函数；$v(t)$ 可导时，$a(t)$ 是速度的导函数。因此位置连续不保证速度连续，速度连续也不保证加速度连续。理想碰撞或分段控制常使速度图出现折点甚至跳变，这时某一瞬间的普通加速度可能不存在，需要按碰撞前后分段讨论。讲义中的匀变速公式默认所考察时间段内加速度有限且恒定。

消去时间得到的关系
\begin{equation}
 v^2-v_0^2=2a(x-x_0)
\end{equation}
特别适合题目不给时间的情形。它含有速度平方，求最终速度时仍要结合运动方向选择正、负根；只写 $v=\sqrt{v_0^2+2a\Delta x}$ 会默认最终沿正方向。若物体先减速到零再反向，同一段过程还应以换向时刻为界分段，避免把一条带符号的位移误当成全程路程。

匀变速运动中，任意时间段的平均速度等于首末速度的算术平均，但这个结论依赖加速度恒定。一般变速运动不能只用 $(v_1+v_2)/2$ 代替区间平均速度，应回到 $v$-$t$ 图面积或积分。相应地，列车中点通过定点的速度不是首尾速度的算术平均，正是因为题目给的是相等位移而非相等时间。

列运动方程时还应先统一单位。若距离用米、时间用秒，速度和加速度应分别化为 $\mathrm{m/s}$ 与 $\mathrm{m/s^2}$；千米每小时不能直接与米、秒混算。答案完成后，用量纲检查每一项：位置方程各项必须都是长度，速度方程各项必须都是速度。量纲正确不能保证数值一定正确，却能迅速发现漏乘时间、误用平方或单位换算错误。

若 $a$ 恒定，初始位置和速度为 $x_0,v_0$，则
\begin{equation}
  v=v_0+at,\qquad
  x=x_0+v_0t+\frac12at^2.
\end{equation}
$v$-$t$ 图的斜率是 $a$，$a$-$t$ 图的面积是速度变化量。

匀加速运动的 $v$-$t$ 图是直线，斜率为常量 $a$；直线下的梯形面积给出位移，得到位置的二次函数。一般运动的 $v$-$t$ 曲线在某点的切线斜率给出该时刻加速度；反过来，给出 $a$-$t$ 图后，两时刻间曲线的有向面积等于速度改变量 $v_2-v_1$。这与“位置图的斜率给速度、速度图的面积给位移”组成一组相互对应的微分与积分关系。

\begin{origexamplebox}{用速度图推出位置公式}\label{ch02:ex:k02-2-9}
根据匀加速运动的速度-时间图像，由 $v=v_0+at$ 推出位置随时间的关系。
\end{origexamplebox}
\Needspace{6\baselineskip}\begin{xesolution}
匀加速运动满足 $v(t)=v_0+at$，所以 $v$--$t$ 图是从 $(0,v_0)$ 到
$(t,v)$ 的直线。位置变化等于速度图下的有向面积：
\[
 x-x_0=\int_0^t v(t')\,\dd t'.
\]
几何上，这一面积是一个梯形；也可拆成高为 $v_0$ 的矩形和高为 $at$ 的三角形：
\[
 x-x_0=v_0t+\frac12(at)t
 =v_0t+\frac12at^2.
\]
若直接用梯形面积，同样有
\[
  x-x_0=\frac{v_0+v}{2}t
  =\frac{v_0+(v_0+at)}{2}t
  =v_0t+\frac12at^2.
\]
因此
\[
 \boxed{x=x_0+v_0t+\frac12at^2}.
\]
检查：对右式求时间导数，得到 $\dd x/\dd t=v_0+at=v$，与原速度式一致；令 $a=0$，又退化为匀速公式 $x=x_0+v_0t$。量纲方面，$v_0t$ 与 $at^2$ 都是长度。图线若部分位于时间轴下方，仍须按有向面积计算，公式中的速度和加速度都要保留符号。
\end{xesolution}

\begin{origexamplebox}{消去时间}\label{ch02:ex:k02-2-10}
匀加速质点由 $(x_1,v_1)$ 运动到 $(x_2,v_2)$，证明
\[
  v_2^2-v_1^2=2a(x_2-x_1).
\]
\end{origexamplebox}
\Needspace{6\baselineskip}\begin{xesolution}
记从状态1到状态2所经历的时间为 $\Delta t$。由于加速度恒定，
\[
 v_2-v_1=a\Delta t,
\]
而这段时间内的速度图是一条直线，因此平均速度是首、末速度的算术平均：
\[
  x_2-x_1=\frac{v_1+v_2}{2}\Delta t.
\]
将第二式乘以 $2a$，再用 $a\Delta t=v_2-v_1$ 消去时间：
\[
\begin{aligned}
2a(x_2-x_1)
 &=a(v_1+v_2)\Delta t\\
 &=(v_2-v_1)(v_2+v_1)\\
 &=v_2^2-v_1^2.
\end{aligned}
\]
故
\[
 \boxed{v_2^2-v_1^2=2a(x_2-x_1)}.
\]
两边量纲都是 $\mathrm{m^2/s^2}$。该关系不含时间，适用于所考察整段加速度恒定的情形；若由它反求 $v_2$，必须结合运动方向选择平方根的正负号，不能总写成正根。
\end{xesolution}

\begin{origexamplebox}{列车中部通过定点的速度}\label{ch02:ex:k02-2-11}
一列火车以恒定加速度沿直线前进。车头通过点 $O$ 时速度为 $v_1$，车尾通过 $O$ 时速度为 $v_2$。求列车中部通过 $O$ 时的速度 $v$。
\end{origexamplebox}
\Needspace{6\baselineskip}\begin{xesolution}
设车长为 $L$，取列车前进方向为正。把整列车看成刚体，用车头、车中部或车尾上任一点的位置描述列车平移，三点在同一时刻具有相同速度。

车头刚通过 $O$ 时，把该时刻记为状态1。等车尾通过 $O$ 时，整列车恰向前平移了一个车长 $L$，于是由无时间公式
\[
  v_2^2-v_1^2=2aL.
\]
当列车中部通过 $O$ 时，列车相对状态1只平移了 $L/2$，故
\[
 v^2-v_1^2=2a\frac L2=aL=\frac{v_2^2-v_1^2}{2}.
\]
整理得
\[
 v^2=\frac{v_1^2+v_2^2}{2}.
\]
列车始终沿正方向前进，所以取正根：
\[
  \boxed{v=\sqrt{\frac{v_1^2+v_2^2}{2}}}.
\]
这不是 $(v_1+v_2)/2$，因为中部对应的是总位移的一半，不是总时间的一半。检查：若 $a=0$，则 $v_1=v_2$，上式正确退化为 $v=v_1$；若 $v_2>v_1$，所得 $v$ 确在 $v_1$ 与 $v_2$ 之间。
\end{xesolution}

\begin{origexamplebox}{通过路口或制动停车}\label{ch02:ex:k02-2-12}
汽车以 $v_0=20\,\mathrm{m/s}$ 行驶，在距路口 $60\,\mathrm m$ 处发现倒计时还剩 $2\,\mathrm s$。忽略司机反应时间，并假定采取方案后加速度保持恒定。若加速通过，所需加速度至少多大？若停车不闯红灯，制动加速度大小至少多大？
\end{origexamplebox}
\Needspace{6\baselineskip}\begin{xesolution}
两种选择要分别建立临界条件。“至少”的边界都对应恰好满足要求；比临界值更小就会失败。

\textbf{加速通过。}设加速度沿行驶方向为 $a_{\rm pass}>0$。要在倒计时结束时刚好到达路口，$2\,\mathrm s$ 内的位移应恰为 $60\,\mathrm m$：
\[
60\,\mathrm m=(20\,\mathrm{m/s})(2\,\mathrm s)
+\frac12a_{\rm pass}(2\,\mathrm s)^2.
\]
所以
\[
\boxed{a_{\rm pass}=10\,\mathrm{m/s^2}}.
\]
此时末速度为 $20+10\times2=40\,\mathrm{m/s}>0$，运动方向没有改变。

\textbf{制动停车。}设制动加速度的大小为 $a_{\rm stop}>0$，实际加速度为 $-a_{\rm stop}$。最小制动能力对应汽车恰在路口前速度降为零、制动距离恰为 $60\,\mathrm m$。使用无时间公式：
\[
0^2-(20\,\mathrm{m/s})^2
=2(-a_{\rm stop})(60\,\mathrm m),
\]
故
\[
\boxed{a_{\rm stop}=\frac{10}{3}\,\mathrm{m/s^2}}.
\]
复核制动时间为 $t=v_0/a_{\rm stop}=6\,\mathrm s$；停车方案不要求在剩余 $2\,\mathrm s$ 内停下，只要求不越过路口。两种计算都忽略了司机反应时间：若反应阶段仍以 $20\,\mathrm{m/s}$ 前进，剩余制动距离会缩短，实际所需制动加速度将更大。
\end{xesolution}

\begin{origexamplebox}{追尾避免条件}\label{ch02:ex:k02-2-13}
$A$ 车以 $20\,\mathrm{m/s}$ 行驶，发现前方同轨道上相距 $100\,\mathrm m$ 的另一列车以 $10\,\mathrm{m/s}$ 行驶，立即以大小为 $a$ 的恒定加速度制动。求不相撞所需的最小 $a$。
\end{origexamplebox}
\Needspace{6\baselineskip}\begin{xesolution}
取行驶方向为正，设后车为 $A$、前车为 $B$。从开始制动计时，
\[
x_A=20t-\frac12at^2,
\qquad
x_B=100+10t.
\]
两车间距为“前车位置减后车位置”：
\[
d(t)=x_B-x_A=100-10t+\frac12at^2.
\]
只要 $d(t)\ge0$，后车就没有越过前车。对时间求导，
\[
\frac{\dd d}{\dd t}=-10+at.
\]
所以间距在
\[
t_*=\frac{10}{a}
\]
时最小；这也正是两车速度相等、相对速度降为零的时刻。代回得
\[
d_{\min}=100-10\frac{10}{a}
+\frac12a\left(\frac{10}{a}\right)^2
=100-\frac{50}{a}.
\]
不相撞要求最小间距非负：
\[
100-\frac{50}{a}\ge0,\qquad \boxed{a_{\min}=0.50\,\mathrm{m/s^2}}.
\]
也可用相对运动的“制动距离”快速复核：初始接近速度为
$20-10=10\,\mathrm{m/s}$，相对减速度大小为 $a$，消除相对速度所需的闭合距离为
\[
\frac{(10\,\mathrm{m/s})^2}{2a}\le100\,\mathrm m,
\]
得到同一条件。临界情况下 $a=0.50\,\mathrm{m/s^2}$、$t_*=20\,\mathrm s$，两车以共同速度 $10\,\mathrm{m/s}$ 恰好相接；若要求保留正的安全间距，则应取 $a>0.50\,\mathrm{m/s^2}$。
\end{xesolution}

\begin{origexamplebox}{直道追及条件}\label{ch02:ex:k02-2-14}
两辆赛车进入长度为 $L$ 的直道，并在所考察时间内始终沿直道正方向运动。以两车车头所在点代表各自位置：初始时后车位于直道起点，前车在其前方 $d$ 处；后车初速为 $v_1$，前车初速为 $v_2>v_1$；后车加速度为 $a_1>0$，前车加速度为 $a_2$，且 $a_1>a_2$。求后车车头在到达直道末端以前追上前车车头的条件。若要求车身完全超越，还需另给两车车长，本题不讨论该情形。
\end{origexamplebox}
\Needspace{6\baselineskip}\begin{xesolution}
取后车进入直道的位置为 $x=0$，令前车初始位置为 $x=d$。两车的位置分别为
\[
x_1(t)=v_1t+\frac12a_1t^2,
\qquad
x_2(t)=d+v_2t+\frac12a_2t^2.
\]
设
\[
\Delta v=v_2-v_1>0,
\qquad
\Delta a=a_1-a_2>0.
\]
前车与后车的间距是
\[
 r(t)=x_2-x_1=d+\Delta v\,t-\frac12\Delta a\,t^2.
\]
开始时 $r(0)=d>0$，且前车初速度较大，间距会先增大；但后车加速度更大，最终有可能追上。相遇条件为 $r(t)=0$：
\[
\frac12\Delta a\,t^2-\Delta v\,t-d=0.
\]
两个根中只有下式为正，故第一次相遇时刻为
\[
 t_*=\frac{\Delta v+\sqrt{\Delta v^2+2\Delta a d}}{\Delta a}.
\]
相遇点坐标等于后车走过的距离
\[
x_*=x_1(t_*)=v_1t_*+\frac12a_1t_*^2.
\]
因此，后车能在直道末端以前追上前车的充要条件是
\[
  \boxed{v_1t_*+\frac12a_1t_*^2<L}.
\]
另一种检验方式是不先求相遇点。令后车到达直道末端的时间为 $t_L$，即
\[
L=v_1t_L+\frac12a_1t_L^2.
\]
因题设 $a_1>0$，取正根得
\[
t_L=\frac{-v_1+\sqrt{v_1^2+2a_1L}}{a_1},
\]
则条件等价为
\[
\boxed{r(t_L)=d+\Delta v\,t_L-\frac12\Delta a\,t_L^2<0}.
\]
因为 $r(t)$ 的正根只有一个，$r(t_L)<0$ 恰好表示相遇已经在 $t_L$ 之前发生；若 $r(t_L)=0$，只能说明两车车头恰在直道末端重合，不满足“末端以前”。量纲检查中，$\Delta vt$、$\Delta at^2$、$d$ 和 $L$ 都是长度。若计算得到 $x_*\ge L$，即使无限长道路上最终能追上，也不能算在本段直道内完成追及。
\end{xesolution}

\subsection{落体运动}

取竖直向上为正方向，忽略空气阻力时，地表附近所有自由落体具有相同加速度
\begin{equation}
  a=-g,\qquad g\approx9.8\,\mathrm{m/s^2}.
\end{equation}
若 $t=0$ 时位置和速度为 $x_0,v_0$，则
\begin{equation}
  v=v_0-gt,\qquad x=x_0+v_0t-\frac12gt^2.
\end{equation}
公式中的正负号由坐标正方向决定；若改取竖直向下为正，加速度分量改为 $a=+g$，
相应重力项也取正号。$g$ 始终表示重力加速度的大小，不随坐标反向变成负数。

把物体由静止释放后，在地表附近小范围内可近似认为重力加速度大小不随高度改变，这种运动是匀加速运动的一种。实验表明，在忽略空气阻力时，不同质量和材料的物体具有相同的自由落体加速度。若由原点静止释放，竖直向上为正，则 $v=-gt$、$x=-gt^2/2$；若初始位置或初速度不为零，直接使用上面的通式。所谓“上抛运动”和“下抛运动”并不是新的公式，只是 $v_0$ 符号不同。

空气中羽毛与石块落得不同，主要来自阻力相对重力的影响不同，并不否定真空中自由落体加速度相同。使用恒定 $g$ 的模型还隐含高度变化远小于地球半径；若运动范围扩展到高空，$g$ 的大小和方向都需要随位置修正。在本节范围内，把竖直运动统一写成带符号的匀变速方程，比为“自由落下、上抛、下抛”分别背三套公式更可靠。

\begin{origexamplebox}{三种竖直运动的位置}\label{ch02:ex:k02-2-15}
以抛出点为原点、竖直向上为正，取 $g=10\,\mathrm{m/s^2}$。分别求：（1）静止释放后 $1,5,10\,\mathrm s$ 的位置；（2）以 $10\,\mathrm{m/s}$ 竖直上抛后 $1,2,3\,\mathrm s$ 的位置；（3）以 $10\,\mathrm{m/s}$ 竖直下抛后相同时刻的位置。
\end{origexamplebox}
\Needspace{6\baselineskip}\begin{xesolution}
统一取抛出点为 $x=0$、竖直向上为正。因为 $a=-g=-10\,\mathrm{m/s^2}$，三种运动都使用
\[
x(t)=v_0t-\frac12gt^2=v_0t-5t^2,
\]
区别只在初速度 $v_0$。

\textbf{（1）静止释放：$v_0=0$。}
\[
x=-5t^2.
\]
分别代入 $t=1,5,10\,\mathrm s$，得到
\[
\begin{array}{c|ccc}
t/\mathrm{s}&1&5&10\\\toprule
 x/\mathrm m&-5&-125&-500\\\bottomrule
\end{array}
\]

\textbf{（2）竖直上抛：$v_0=+10\,\mathrm{m/s}$。}
\[
x=10t-5t^2.
\]
在 $t=1,2,3\,\mathrm s$ 时，
\[
\begin{array}{c|ccc}
t/\mathrm{s}&1&2&3\\\toprule
x/\mathrm m&5&0&-15\\\bottomrule
\end{array}
\]
$t=1\,\mathrm s$ 时速度 $v=10-10t=0$，物体处在最高点；$t=2\,\mathrm s$ 时回到原点；随后位置变为负。

\textbf{（3）竖直下抛：$v_0=-10\,\mathrm{m/s}$。}
\[
x=-10t-5t^2,
\]
故在 $t=1,2,3\,\mathrm s$ 时，
\[
\begin{array}{c|ccc}
t/\mathrm{s}&1&2&3\\\toprule
x/\mathrm m&-15&-40&-75\\\bottomrule
\end{array}
\]
负号只说明物体位于原点下方。三种位置公式的每一项都是长度；上抛、下抛无需另背公式，只要把初速度按坐标方向正确取正负即可。
\end{xesolution}

\begin{origexamplebox}{到达指定高度的时间}\label{ch02:ex:k02-2-16}
质点从原点以 $v_0=20\,\mathrm{m/s}$ 竖直上抛，取 $g=10\,\mathrm{m/s^2}$。求到达 $x=15,20,25,0,-25\,\mathrm m$ 所需时间。
\end{origexamplebox}
\Needspace{6\baselineskip}\begin{xesolution}
位置方程为
\[
x=20t-5t^2.
\]
对给定位置 $x$，整理为
\[
5t^2-20t+x=0,
\]
其根为
\[
t=2\pm\sqrt{4-\frac{x}{5}},
\]
并且只保留 $t\ge0$ 的根。

逐项计算：
\[
\begin{aligned}
x=15\,\mathrm m:&\quad t=2\pm1=1,3\,\mathrm s;\\
x=20\,\mathrm m:&\quad t=2\,\mathrm s\quad\text{（重根）};\\
x=25\,\mathrm m:&\quad 4-x/5=-1<0\quad\text{（无实根）};\\
x=0:&\quad t=0,4\,\mathrm s;\\
x=-25\,\mathrm m:&\quad t=2\pm3=-1,5\,\mathrm s,
\end{aligned}
\]
最后一项舍去负时间 $-1\,\mathrm s$。汇总为
\[
\begin{array}{c|ccccc}
x/\mathrm m&15&20&25&0&-25\\\toprule
t/\mathrm s&1,3&2&\text{无实根}&0,4&5\\\bottomrule
\end{array}
\]
为什么 $x=15\,\mathrm m$ 有两个时刻？速度
\[
v=20-10t
\]
在 $t=1\,\mathrm s$ 为正、在 $t=3\,\mathrm s$ 为负，分别对应上升和下降经过同一高度。最高点由 $v=0$ 得 $t=2\,\mathrm s$，最大高度为
\[
x_{\max}=\frac{v_0^2}{2g}=20\,\mathrm m,
\]
所以不可能到达 $25\,\mathrm m$。$x=0$ 的两个时刻分别是抛出瞬间和回到原点时。
\end{xesolution}

\begin{origexamplebox}{一落一上抛}\label{ch02:ex:k02-2-17}
高度 $h$ 处的质点 $A$ 自由下落，同时从地面以初速度 $v_0$ 竖直上抛质点 $B$。允许临界情况下两者恰在 $A$ 落地时相遇，求所需的最小 $v_0$。
\par 题图见图~\ref{ch02:ref:question-figure-2-17}。
\KQuestionAR
\end{origexamplebox}
\Needspace{6\baselineskip}\begin{xesolution}
取地面为 $y=0$、竖直向上为正。两球的位置分别为
\[
y_A=h-\frac12gt^2,
\qquad
y_B=v_0t-\frac12gt^2.
\]
两式中的重力项完全相同，这是因为两球都具有加速度 $-g$。令 $y_A=y_B$，得到
\[
h=v_0t,
\]
所以相遇时刻为
\FreeFallMeetingDiagram
\[
t_m=\frac{h}{v_0}.
\]
$A$ 从高度 $h$ 落到地面的时间由 $h-gt_f^2/2=0$ 得
\[
t_f=\sqrt{\frac{2h}{g}}.
\]
要在 $A$ 落地以前或恰好落地时相遇，必须有 $t_m\le t_f$，即
\[
\frac{h}{v_0}\le\sqrt{\frac{2h}{g}}.
\]
因 $h,v_0>0$，整理得临界值
\[
  \boxed{v_{0,\min}=\sqrt{\frac{gh}{2}}}.
\]
也可从相对运动理解：两球相对加速度为零，初始间距 $h$ 以恒定相对速度 $v_0$ 闭合。临界情况下
$t_m=t_f$，两球恰在地面相遇，所以题设所允许的最小值能够取到。若另加“必须在离地面高度大于零处相遇”的严格条件，则必须取
$v_0>\sqrt{gh/2}$，此时不存在能取到的最小值，上式只是允许初速度的下确界。答案的量纲为
$\sqrt{(\mathrm{m/s^2})\mathrm m}=\mathrm{m/s}$。
\end{xesolution}

\begin{origexamplebox}{超人追赶自由落体}\label{ch02:ex:k02-2-18}
一名女子从高为 $H$ 的楼顶落下。经过时间 $T$ 后，超人从楼顶以向下初速度 $v_0$ 出发，并同样以重力加速度向下运动。允许临界情况下恰在女子落地时追上她，求所需的最小 $v_0$。
\end{origexamplebox}
\Needspace{6\baselineskip}\begin{xesolution}
取楼顶为原点、竖直向下为正，并从女子开始下落时计时。女子的位置为
\[
y_w(t)=\frac12gt^2,
\]
故她的落地时刻为
\[
t_f=\sqrt{\frac{2H}{g}}.
\]
超人只有在 $T<t_f$ 时还有追赶机会。令他出发后的时间为
$\tau=t-T$，则
\[
y_s(\tau)=v_0\tau+\frac12g\tau^2.
\]
要使所需初速度最小，应把允许追赶的时间用到最长，即在临界情况下恰于女子落地时追上。此时
\[
\Delta=t_f-T
\]
是超人的飞行时间，并满足
\[
H=v_0\Delta+\frac12g\Delta^2.
\]
解得
\[
\boxed{v_{0,\min}=\frac{H-\frac12g(t_f-T)^2}{t_f-T}
=\frac{gT(2t_f-T)}{2(t_f-T)}}.
\]
第二个形式可由 $H=gt_f^2/2$ 化简得到。

用相对运动复核更能看清条件。超人出发时，女子已经领先
\[
s_0=\frac12gT^2
\]
并具有向下速度 $gT$。两人以后加速度相同，所以相对加速度为零；超人相对女子的接近速度恒为 $v_0-gT$。临界条件是
\[
(v_0-gT)(t_f-T)=\frac12gT^2,
\]
解得同一结果。题设允许临界状态，所以等号给出可取得的最小值；若另要求严格在落地前追上，则应取 $v_0>v_{0,\min}$，此时上式只是下确界。当 $T\to0$ 时该下界趋于零，当 $T\to t_f^-$ 时所需初速度趋于无穷，均符合直觉。
\end{xesolution}

\begin{origexamplebox}{恢复系数与无限次反弹}\label{ch02:ex:k02-2-19}
质点从高度 $H$ 自由下落，每次碰地后的反弹速率为碰前速率的 $e$ 倍，$0<e<1$。求从第一次反弹后算起到最终停止的总时间和总路程。
\par 题图见图~\ref{ch02:ref:question-figure-2-19}。
\KQuestionAT
\end{origexamplebox}
\Needspace{6\baselineskip}\begin{xesolution}
从高度 $H$ 自由下落到第一次碰地前，由
\[
u^2=2gH
\]
得碰前速率
\[
u=\sqrt{2gH}.
\]
恢复系数为 $e$ 表示每次碰地后的反弹速率等于碰前速率的 $e$ 倍。由于每次飞行都回到同一地面高度，下一次碰地前的速率又等于该次反弹初速，所以历次反弹初速度为
\[
eu,\ e^2u,\ e^3u,\ldots.
\]

若某次反弹初速度为 $w$，上升到最高点用时 $w/g$，再下落回地面也用时 $w/g$，故一轮飞行时间为
\[
\frac{2w}{g}.
\]
最高高度为 $w^2/(2g)$，上、下路程之和为
\[
2\frac{w^2}{2g}=\frac{w^2}{g}.
\]
题目从第一次反弹后开始计时和计路程，因此从 $w=eu$ 对应的项开始求和：
\BounceDiagram
\[
\begin{aligned}
T&=\frac{2u}{g}\sum_{n=1}^{\infty}e^n
=\frac{2u}{g}\frac{e}{1-e}\\
&=\boxed{\frac{2e\sqrt{2gH}}{g(1-e)}},
\end{aligned}
\]
\[
\begin{aligned}
S&=\frac{u^2}{g}\sum_{n=1}^{\infty}e^{2n}
=\frac{u^2}{g}\frac{e^2}{1-e^2}\\
&=\boxed{\frac{2He^2}{1-e^2}}.
\end{aligned}
\]
时间与速度成正比，所以按 $e$ 求和；路程与速度平方成正比，所以按 $e^2$ 求和，二者不能混用。检查极限：$e\to0$ 时第一次碰撞后不再反弹，$T,S\to0$；$e\to1^-$ 时能量几乎不损失，总时间和总路程都发散。若题目要求从最初释放算起，还应分别再加第一次下落时间 $\sqrt{2H/g}$ 和路程 $H$。
\end{xesolution}

\begin{origexamplebox}{逐次增强的回向加速度}\label{ch02:ex:k02-2-20}
质点以速度大小 $v_0$ 撞向一固定板。第一次及以后每次碰板时，速度均立即反向且大小不变。第一次反弹后，它以大小为 $a$、方向始终指向板的加速度运动；以后每完成一次碰板，回向加速度的大小变为前一次的 $k>1$ 倍。求第一次反弹后的总时间和总路程。
\end{origexamplebox}
\Needspace{6\baselineskip}\begin{xesolution}
把板的位置取为 $x=0$，离开板的方向取为正。第 $n$ 个“离板—回板”循环中，加速度大小为
\[
a_n=ak^{n-1}\qquad(n=1,2,\ldots),
\]
方向始终指向板，所以在 $x>0$ 区域实际加速度为 $-a_n$。每次碰板后，质点都以速度 $+v_0$ 离板。

在第 $n$ 个循环内，离板后经过时间 $\tau$ 的位置和速度为
\[
x=v_0\tau-\frac12a_n\tau^2,
\qquad
v=v_0-a_n\tau.
\]
令 $v=0$，得转向时刻 $\tau_1=v_0/a_n$；此时最远距离为
\[
x_{\max,n}=\frac{v_0^2}{2a_n}.
\]
令 $x=0$ 并舍去 $\tau=0$，得回到板的时刻
\[
\tau_n=\frac{2v_0}{a_n}.
\]
往返路程是最远距离的两倍：
\[
s_n=2x_{\max,n}=\frac{v_0^2}{a_n}.
\]
所以总时间和总路程分别为几何级数
\[
\begin{aligned}
T&=\frac{2v_0}{a}\sum_{n=0}^{\infty}k^{-n}
=\boxed{\frac{2v_0}{a}\frac{k}{k-1}},
\end{aligned}
\]
\[
\begin{aligned}
S&=\frac{v_0^2}{a}\sum_{n=0}^{\infty}k^{-n}
=\boxed{\frac{v_0^2}{a}\frac{k}{k-1}}.
\end{aligned}
\]
级数收敛依赖题设 $k>1$。单位检查：$v_0/a$ 是时间，$v_0^2/a$ 是长度。再检查极限：$k\to\infty$ 时，后续循环的加速度极大、耗时和路程趋零，结果只剩第一次循环的 $2v_0/a$ 与 $v_0^2/a$；$k\to1^+$ 时各循环几乎一样长，总和发散。
\end{xesolution}

\begin{origexamplebox}{周期运动的频闪取样}\label{ch02:ex:k02-2-21}
竖直真空玻璃管内，小球从距管底高度 $H$ 处由静止下落，与管底完全弹性碰撞后不停上下运动。照相机每隔固定时间 $T$ 拍摄一次，曝光时间可忽略；所有照片中小球均处于同一位置。求可能的 $H$，以及相应照片中小球距管底的可能高度。
\end{origexamplebox}
\Needspace{6\baselineskip}\begin{xesolution}
设从最高点落到底部的时间
\[
\tau=\sqrt{\frac{2H}{g}},
\]
完整运动周期为 $P=2\tau$。同一高度除转折点外在一个周期内出现两次，且两次相位关于最高点对称。

\begin{center}
\begin{tikzpicture}[x=1.25cm,y=2.05cm,>=Latex]
  \draw[axis,->] (-.15,0)--(4.35,0) node[right] {$t$};
  \draw[axis,->] (0,-.06)--(0,1.18) node[above] {$y/H$};
  \draw[curve,domain=0:2,samples=80] plot (\x,{1-(\x/2)^2});
  \draw[curve,domain=2:4,samples=80] plot (\x,{1-((4-\x)/2)^2});
  \draw[guide] (0,.75)--(4,.75);
  \fill (1,.75) circle (1.7pt);
  \fill (3,.75) circle (1.7pt);
  \node[below left=3pt] at (1,.75) {$t_1$};
  \node[below right=3pt] at (3,.75) {$t_2$};
  \draw[guide] (1,.75)--(1,.92) (3,.75)--(3,.92);
  \draw[<->,DeepBlue] (1,.92)--(3,.92) node[midway,above=1pt] {$\tau$};
  \foreach \x/\lab in {2/$\tau$,4/$P$}{\draw (\x,.025)--(\x,-.025) node[below] {\lab};}
  \node[left] at (0,1) {$1$};
\end{tikzpicture}
\figcap{一个周期内的高度—时间图。除最高点和碰撞点外，同一高度对应下降、上升两个相位；二者只有在图示特殊位置才相差半周期。}
\end{center}

为什么不会遗漏其他取样周期？把拍摄时刻按周期 $P$ 取模。若 $T/P$ 为无理数，取样相位会在整个周期内稠密分布，
高度显然不可能恒定；若 $T/P=p/q$ 为最简分数，一个循环会访问 $q$ 个不同相位。
而图中任一高度在一个周期内至多对应两个相位，所以只能有 $q=1$ 或 $q=2$。
它们分别对应相位完全不变和每次跨越半个周期，恰好就是下面两类。

若相机开始拍摄的相位未指定，能使无限序列始终停留在同一高度的取样只有两类：
\begin{enumerate}
  \item $T=nP=2n\tau$，$n=1,2,\ldots$。此时每次相位完全相同，
  \[
    \boxed{H=\frac{gT^2}{8n^2}},
  \]
  拍摄高度可为 $0\le y\le H$，由第一张照片的相位决定。
  \item $T=(2n+1)\tau$，$n=0,1,2,\ldots$。相位每次相差半个周期；要使两相位高度相同，它们必须位于上、下行关于中点对称的位置，即距最高点的等效时间均为 $\tau/2$。于是
  \[
    \boxed{H=\frac{gT^2}{2(2n+1)^2}},\qquad
    \boxed{y=H-\frac12g\left(\frac{\tau}{2}\right)^2=\frac34H}.
  \]
\end{enumerate}

第一类中，$T=nP$ 表示相邻两张照片之间恰好经过整数个完整周期，所以不仅高度相同，速度的大小和方向也完全相同。由
$T=2n\tau$ 得 $\tau=T/(2n)$，代入 $H=g\tau^2/2$，立即得到
$H=gT^2/(8n^2)$。由于首次曝光可以发生在周期内任意相位，照片中的高度也可以取闭区间 $[0,H]$ 内的任意值；一旦第一张照片确定，后续照片就重复同一状态。

第二类中，相邻照片的相位相差半个周期，速度方向相反，因而只能利用上升段和下降段在同一高度处的对称性。把两次经过该高度的时刻记为 $t_1,t_2$，图中有
$t_2-t_1=\tau=P/2$。两时刻分别位于下降、上升过程的中点，因此从最高点计的飞行时间都是 $\tau/2$，高度为
\[
 y=H-\frac12g\left(\frac{\tau}{2}\right)^2
  =H-\frac14\left(\frac12g\tau^2\right)=\frac34H.
\]
这里使用了 $H=g\tau^2/2$。两类结果的量纲均正确：$gT^2$ 是长度；当 $n$ 增大时，在固定拍摄间隔 $T$ 下允许的 $H$ 按 $1/n^2$ 缩小，与单程时间 $\tau\propto\sqrt H$ 相符。以上讨论还默认碰撞瞬间完成且完全弹性；若恢复系数小于 $1$，运动不再严格周期，结论不能直接沿用。

若题意另规定第一张照片就在释放瞬间拍摄，则只有第一类，且所有照片都拍到最高点 $y=H$。把“首次拍摄相位”说清楚，是这道题条件完整性的关键。
\end{xesolution}
}

\endgroup

% MechanicsNumberedSolutionFigures20261023-v1
% MechanicsQualityRevision20261023-v1
% QualityRevision20261023
\begingroup
\tcbset{breakable/.default=false}
\chapter{合运动、抛体、圆周与一般曲线运动}

{
\graphicspath{{../../../sample/xes-original/kinematics-supplement/}{../../../sample/xes-original/kinematics-supplement/assets/}}
\section{运动的合成与分解}
\newcommand{\CoordinateSystemsDiagram}{%
\begin{center}\begin{tikzpicture}[scale=.9,>=Latex]
  \begin{scope}[xshift=-3.2cm]
    \draw[kaxis] (-2,0)--(2.2,0) node[right,clearlabel]{$x$};
    \draw[kaxis] (0,-1.7)--(0,1.9) node[above,clearlabel]{$y$};
    \coordinate (P) at (1.35,1.15);\node[point,label={[clearlabel]above:$P(x,y)$}] at (P) {};
    \draw[kguide] (P)--(1.35,0) (P)--(0,1.15);
    \node[below left] at (0,0){$O$};
  \end{scope}
  \begin{scope}[xshift=3.2cm]
    \draw[kaxis] (0,0)--(2.3,0) node[right,clearlabel]{$N$};
    \draw[kaxis] (0,0)--(0,1.9);
    \coordinate (P) at (1.45,1.25);\draw[vector] (0,0)--(P) node[midway,clearlabel,above left]{$r$};
    \node[point,label={[clearlabel]above:$P(r,\theta)$}] at (P) {};
    \draw[->,TextGray] (.65,0) arc[start angle=0,end angle=41,radius=.65];
    \node[clearlabel] at (.82,.28){$\theta$};
    \node[below left] at (0,0){$O$};
  \end{scope}
\end{tikzpicture}\figcap{平面直角坐标系与平面极坐标系。}\label{ch03:fig:k03-2-5}\end{center}}

\newcommand{\PathChoicesDiagram}{%
\begin{center}\begin{minipage}{.96\linewidth}\centering
\begin{tikzpicture}[scale=.9,>=Latex]
 \coordinate (A) at (0,0);\coordinate (B) at (8,1.5);
 \draw[kcurve] (A)--(B) node[pos=.72,above=3pt]{$P_1$};
 \draw[kcurve] (A)..controls (1,2.9) and (5.9,3.0)..(B) node[pos=.46,above=3pt]{$P_2$};
 \draw[kcurve] (A)..controls (1.3,1.4) and (2.1,1.65)..(3.2,.25)
 ..controls (5,-1.0) and (7.7,0)..(B);
 \draw[kcurve] (A)..controls (-2.5,4) and (3.6,4.3)..(3.0,-.1)
 ..controls (2.8,-1.8) and (8.9,-1.1)..(B);
 \draw[kcurve] (A)..controls (3.6,-2.4) and (8.8,-2.0)..(B);
 \node at (5,.3){$P_3$};\node at (.8,3.35){$P_4$};\node at (5.5,-1.85){$P_5$};
 \node[point,label=below left:$A$] at (A){};\node[point,label=above right:$B$] at (B){};
\end{tikzpicture}
\figcap{五条路径 $P_1,\ldots,P_5$ 均从 $A$ 到 $B$；路径不同，位移相同。}\label{ch03:fig:k03-2-6}
\end{minipage}\end{center}}

\newcommand{\DisplacementFanDiagram}{%
\begin{center}\begin{minipage}{.98\linewidth}\centering
\begin{tikzpicture}[scale=.95,>=Latex]
 \coordinate (A) at (0,0);
 \foreach \i in {1,...,10}{
  \pgfmathsetmacro{\xx}{.95*\i}\pgfmathsetmacro{\yy}{.8*\xx-.06*\xx*\xx}
  \draw[DeepBlue!60,->,thin] (A)--(\xx,\yy);
 }
 \draw[kcurve,domain=0:9.75,samples=100] plot (\x,{.8*\x-.06*\x*\x});
 \foreach \i in {1,...,10}{
  \pgfmathsetmacro{\xx}{.95*\i}\pgfmathsetmacro{\yy}{.8*\xx-.06*\xx*\xx}
  \node[point] at (\xx,\yy){};\node[above=5pt,font=\small] at (\xx,\yy){$B_{\i}$};
 }
 \node[point,label=below left:$A$] at (A){};
\end{tikzpicture}
\figcap{粗曲线是实际路径；从 $A$ 指向各 $B_i$ 的直箭头是对应位移。}\label{ch03:fig:k03-2-7}
\end{minipage}\end{center}}

\newcommand{\VectorOperationsDiagram}{%
\begin{center}\begin{tikzpicture}[scale=.82,>=Latex]
  \begin{scope}[xshift=-4.4cm]
    \draw[vector] (0,0)--(1.5,.35) node[midway,clearlabel,below]{$\vect A$};
    \draw[vector] (1.5,.35)--(2.15,1.55) node[midway,clearlabel,right]{$\vect B$};
    \draw[vector] (0,0)--(2.15,1.55) node[midway,clearlabel,above left]{$\vect A+\vect B$};
  \end{scope}
  \begin{scope}
    \draw[vector] (0,0)--(1.55,.4) node[midway,clearlabel,below]{$\vect A$};
    \draw[vector] (0,0)--(.55,1.55) node[midway,clearlabel,left]{$\vect B$};
    \draw[vector] (.55,1.55)--(1.55,.4) node[midway,clearlabel,above right]{$\vect A-\vect B$};
  \end{scope}
  \begin{scope}[xshift=4.4cm]
    \draw[vector] (0,0)--(1.7,.35) node[midway,clearlabel,below]{$\vect A$};
    \draw[vector] (0,0)--(.75,1.45) node[midway,clearlabel,left]{$\vect B$};
    \draw[->,TextGray] (.65,.13) arc[start angle=12,end angle=63,radius=.67];
    \node[clearlabel] at (.9,.75){$\theta$};
    \node[clearlabel,below] at (.85,-.55){$\vect A\cdot\vect B=AB\cos\theta$};
  \end{scope}
\end{tikzpicture}\figcap{矢量加法、减法与内积。}\label{ch03:fig:k03-2-8}\end{center}}

\newcommand{\VectorComponentsDiagram}{%
\begin{center}\begin{tikzpicture}[scale=.9,>=Latex]
  \draw[kaxis] (-.4,0)--(4.4,0) node[right,clearlabel]{$x$};
  \draw[kaxis] (0,-.35)--(0,3.0) node[above,clearlabel]{$y$};
  \coordinate (P) at (3.4,2.2);\draw[vector] (0,0)--(P) node[midway,clearlabel,above left]{$\vect A$};
  \draw[kguide] (P)--(3.4,0) node[clearlabel,below]{$A_x$};
  \draw[kguide] (P)--(0,2.2) node[clearlabel,left]{$A_y$};
\end{tikzpicture}\figcap{矢量平移到原点后，终点坐标就是矢量分量。}\label{ch03:fig:k03-2-9}\end{center}}

\newcommand{\CurveVelocityDiagram}{%
\begin{center}\begin{minipage}{.95\linewidth}\centering
\begin{tikzpicture}[>=Latex]
 \draw[kcurve,domain=-3:3,samples=100] plot (\x,{.09*\x^3-.55*\x});
 \foreach \xx in {-2.3,-.9,.8,2.2}{
  \pgfmathsetmacro{\yy}{.09*\xx^3-.55*\xx}
  \pgfmathsetmacro{\slope}{.27*(\xx)*(\xx)-.55}\pgfmathsetmacro{\dx}{1/sqrt(1+(\slope)*(\slope))}
  \node[point] at (\xx,\yy){};\draw[vector] (\xx,\yy)--++(\dx,{\slope*\dx});
 }
\end{tikzpicture}
\figcap{每个速度箭头都从所在位置出发，并沿该点的轨迹切线。}\label{ch03:fig:k03-2-10}
\end{minipage}\end{center}}

\newcommand{\AccelerationVectorDiagram}{%
\begin{center}\begin{tikzpicture}[scale=.9,>=Latex]
  \coordinate (O) at (0,0);\coordinate (A) at (2.3,1.6);\coordinate (B) at (3.2,.75);
  \draw[vector] (O)--(A) node[midway,clearlabel,above left]{$\vect v(t)$};
  \draw[vector] (O)--(B) node[midway,clearlabel,below=7pt]{$\vect v(t+\Delta t)$};
  \draw[vector] (A)--(B) node[midway,clearlabel,above right]{$\Delta\vect v$};
  \node[clearlabel,right] at (3.65,1.2){$\displaystyle \vect a=\lim_{\Delta t\to0}\frac{\Delta\vect v}{\Delta t}$};
\end{tikzpicture}\figcap{速度矢量首端重合，终端之差给出 $\Delta\vect v$。}\label{ch03:fig:k03-2-11}\end{center}}

\newcommand{\RescueRouteDiagram}{%
\begin{center}\begin{minipage}{.96\linewidth}\centering
\begin{tikzpicture}[>=Latex]
 \coordinate (A) at (-2.8,2.3);\coordinate (C) at (.2,0);\coordinate (B) at (1.7,-1.9);
 \coordinate (U) at (.2,1.4);\coordinate (D) at (.2,-1.4);
 \draw[thick] (-3.5,0)--(3.6,0) node[right]{岸线};
 \draw[kguide] (A)--(-2.8,0) (B)--(1.7,0) (U)--(D);
 \draw[vector] (A)--(C) node[pos=.35,above right=2pt]{$v_1$};
 \draw[vector] (C)--(B) node[pos=.45,above right=2pt]{$v_2$};
 \draw pic[draw=TextGray,angle radius=7mm] {angle=U--C--A};
 \draw pic[draw=TextGray,angle radius=7mm] {angle=D--C--B};
 \node at (-.24,1.03){$\theta_1$};\node at (.55,-1.1){$\theta_2$};
 \draw[<->,TextGray] (-3.3,0)--(-3.3,2.3) node[midway,left]{$d_1$};
 \draw[<->,TextGray] (2.15,0)--(2.15,-1.9) node[midway,right]{$d_2$};
 \draw[<->,TextGray] (-2.8,-2.5)--(.2,-2.5) node[midway,below]{$x$};
 \draw[<->,TextGray] (.2,-2.5)--(1.7,-2.5) node[midway,below]{$L-x$};
 \node[point,label=above left:$A$] at (A){};\node[point,label=below right:$B$] at (B){};
 \node[point,label=above right:$C$] at (C){};
\end{tikzpicture}
\figcap{入水点为 $C$；$\theta_1,\theta_2$ 从岸线法线量起，两垂足间距为 $L$。}\label{ch03:fig:k03-ex-2-31}
\end{minipage}\end{center}}

\newcommand{\FootballInterceptionDiagram}{%
\begin{center}\begin{minipage}{.96\linewidth}\centering
\begin{tikzpicture}[scale=1.1,>=Latex]
 \coordinate (A) at (0,0);\coordinate (B) at (4,2);\coordinate (C) at (5,0);\coordinate (X) at (6,0);
 \draw[kaxis] (A)--(X) node[right]{$X$};
 \draw[thick] (A)--(B) node[midway,above left]{$d$};\draw[kguide] (B)--(C);
 \draw[vector] (B)--($(B)!.43!(C)$) node[midway,right=5pt]{$\vect v_p$};
 \draw[vector] (A)--(1.25,0) node[midway,below=4pt]{$\vect v$};
 \draw pic[draw=TextGray,angle radius=8mm] {angle=X--A--B};\node at (1.03,.27){$\theta$};
 \draw pic[draw=TextGray,angle radius=4mm] {right angle=A--B--C};
 \foreach \p/\pos in {A/below left,B/above,C/below}\node[point,label=\pos:$\p$] at (\p){};
\end{tikzpicture}
\figcap{最低奔跑速率对应 $BC\perp AB$；队员速度沿 $BC$，球的速度沿 $AC$。}\label{ch03:fig:k03-ex-2-29}
\end{minipage}\end{center}}

\newcommand{\TwoParticlesMinimumDistanceDiagram}{%
\begin{center}\begin{tikzpicture}[scale=.82,>=Latex]
  \coordinate (P) at (-3.7,-.75);\coordinate (Q) at (2.25,-.75);\coordinate (R) at (-.3,2.25);
  \draw[very thick] (P)--(Q);
  \draw[kguide] (Q)--(R);
  \draw[vector] (-3.25,-.75)--(-1.8,-.75) node[midway,clearlabel,above]{$\vect v_1$};
  \draw[vector] (1.93,-.38)--(.95,.78) node[midway,clearlabel,right]{$\vect v_2$};
  \draw pic[draw=TextGray,->,angle radius=8mm,"$\alpha$"{clearlabel}] {angle=R--Q--P};
  \draw[<->,TextGray] (P)--(Q) node[midway,clearlabel,below=3pt]{$l$};
  \node[point,label={[clearlabel]left:$A,P$}] at (P) {};
  \node[point,label={[clearlabel]right:$B,Q$}] at (Q) {};
  \node[clearlabel,above] at (R){$R$};
\end{tikzpicture}\figcap{初始连线为 $PQ$；$A$ 沿 $PQ$，$B$ 沿 $QR$，最近距离由相对速度直线确定。}\label{ch03:fig:k03-ex-2-30}\end{center}}

\newcommand{\KQuestionAW}{%
\begin{center}\begin{minipage}{.97\linewidth}\centering
\begin{tikzpicture}[>=Latex,scale=.70]

\foreach \x in {-4,-3,-2,-1,1,2,3,4,5}{\draw[RuleGray,thin] (\x,-5.5)--(\x,-.8) (\x,.1)--(\x,5);}
\foreach \y in {-5,-4,-3,-2,-1,1,2,3,4,5}{\draw[RuleGray,thin] (-4,\y)--(-1.05,\y) (.12,\y)--(5.5,\y);}
\draw[kaxis] (-4.2,0)--(5.8,0) node[right]{$x$};
\draw[kaxis] (0,-5.7)--(0,5.3) node[above]{$y$};
\foreach \x in {-3,-2,-1,1,2,3,4,5}\node[below=2pt] at (\x,0){$\x$};
\foreach \y in {-5,-4,-3,-2,-1,1,2,3,4}\node[left=2pt] at (0,\y){$\y$};
\node[above left=3pt] at (0,0){$O$};

\end{tikzpicture}
\figcap{题图：在直角坐标网格中标出题设各点。}\label{ch03:ref:question-figure-2-22}
\end{minipage}\end{center}
}

\newcommand{\KQuestionAX}{%
\begin{center}\begin{minipage}{.97\linewidth}\centering
\begin{tikzpicture}[>=Latex,scale=.9]

\draw[kaxis] (0,0)--(6,0) node[right]{$r$};
\foreach \r in {1,...,5}{\draw (\r,0)--(\r,.1);\node[below] at (\r,0){$\r$};}
\node[left] at (0,0){$O$};
\path[use as bounding box] (-.5,-1.4) rectangle (6.4,1.4);

\end{tikzpicture}
\figcap{题图：以向右射线为极轴，逆时针为正角方向。}\label{ch03:ref:question-figure-2-23}
\end{minipage}\end{center}
}

\newcommand{\KQuestionBC}{%
\begin{center}\begin{minipage}{.97\linewidth}\centering
\begin{tikzpicture}[>=Latex,scale=.9]

\coordinate (A) at (3,0);\coordinate (M) at (0,3);\coordinate (N) at (6,3);
\coordinate (B) at (1.125,3);\coordinate (H) at (3,3);
\draw[kguide] (M)--(N) (A)--(H) (A)--(B);
\node[point,label=below:$A$] at (A){};
\node[left] at (M){$M$};\node[right] at (N){$N$};
\draw[vector] (B)--++(1.1,0) node[midway,above]{$v$};
\draw[vector] (A)--++(0,1.7) node[midway,right]{$u$};
\node[below left] at (B){$B$};\node[above right] at (H){$H$};
\draw[TextGray] (A)++(90:.8) arc(90:122:.8);\node at (2.64,1.13){$\theta$};
\draw[<->] (3.7,0)--(3.7,3) node[midway,right]{$L$};

\end{tikzpicture}
\figcap{题图：潜艇 A 与航线 MN 的垂直距离为 $L$，$\theta$ 从航线法线量起。}\label{ch03:ref:question-figure-2-28}
\end{minipage}\end{center}
}

\newcommand{\KQuestionBD}{%
\begin{center}\begin{minipage}{.97\linewidth}\centering
\begin{tikzpicture}[>=Latex,scale=.9]

\coordinate (A) at (0,0);\coordinate (B) at (4,2);\coordinate (X) at (6,0);
\draw[kaxis] (A)--(X) node[right]{$X$};\draw[thick] (A)--(B) node[midway,above left]{$d$};
\draw[vector] (A)--(1.3,0) node[midway,below]{$v$};
\draw[vector] (B)--++(.5,-1.1) node[right]{$v_p$};
\draw (0.8,0) arc(0:26.565:.8);\node at (1.12,.3){$\theta$};
\node[point,label=below left:$A$] at (A){};\node[point,label=above:$B$] at (B){};

\end{tikzpicture}
\figcap{题图：球沿 AX 运动，队员从 B 开始奔跑；奔跑方向待选。}\label{ch03:ref:question-figure-2-29}
\end{minipage}\end{center}
}

\newcommand{\KQuestionBE}{%
%
\begin{center}\begin{tikzpicture}[scale=.82,>=Latex]
  \coordinate (P) at (-3.7,-.75);\coordinate (Q) at (2.25,-.75);\coordinate (R) at (-.3,2.25);
  \draw[very thick] (P)--(Q);
  \draw[kguide] (Q)--(R);
  \draw[vector] (-3.25,-.75)--(-1.8,-.75) node[midway,clearlabel,above]{$\vect v_1$};
  \draw[vector] (1.93,-.38)--(.95,.78) node[midway,clearlabel,right]{$\vect v_2$};
  \draw pic[draw=TextGray,->,angle radius=8mm,"$\alpha$"{clearlabel}] {angle=R--Q--P};
  \draw[<->,TextGray] (P)--(Q) node[midway,clearlabel,below=3pt]{$l$};
  \node[point,label={[clearlabel]left:$A,P$}] at (P) {};
  \node[point,label={[clearlabel]right:$B,Q$}] at (Q) {};
  \node[clearlabel,above] at (R){$R$};
\end{tikzpicture}\figcap{题图：A、B 初始位于 P、Q，分别沿 PQ、QR 匀速运动。}\label{ch03:ref:question-figure-2-30}\end{center}
}

\newcommand{\KQuestionBF}{%
%
\begin{center}\begin{minipage}{.96\linewidth}\centering
\begin{tikzpicture}[>=Latex]
 \coordinate (A) at (-2.8,2.3);\coordinate (C) at (.2,0);\coordinate (B) at (1.7,-1.9);
 \coordinate (U) at (.2,1.4);\coordinate (D) at (.2,-1.4);
 \draw[thick] (-3.5,0)--(3.6,0) node[right]{岸线};
 \draw[kguide] (A)--(-2.8,0) (B)--(1.7,0) (U)--(D);
 \draw[vector] (A)--(C) node[pos=.35,above right=2pt]{$v_1$};
 \draw[vector] (C)--(B) node[pos=.45,above right=2pt]{$v_2$};
 \draw pic[draw=TextGray,angle radius=7mm] {angle=U--C--A};
 \draw pic[draw=TextGray,angle radius=7mm] {angle=D--C--B};
 \node at (-.24,1.03){$\theta_1$};\node at (.55,-1.1){$\theta_2$};
 \draw[<->,TextGray] (-3.3,0)--(-3.3,2.3) node[midway,left]{$d_1$};
 \draw[<->,TextGray] (2.15,0)--(2.15,-1.9) node[midway,right]{$d_2$};
 \draw[<->,TextGray] (-2.8,-2.5)--(.2,-2.5) node[midway,below]{$x$};
 \draw[<->,TextGray] (.2,-2.5)--(1.7,-2.5) node[midway,below]{$L-x$};
 \node[point,label=above left:$A$] at (A){};\node[point,label=below right:$B$] at (B){};
 \node[point,label=above right:$C$] at (C){};
\node[below left=3pt] at (-3.5,0){$M$};\node[above left=6pt] at (3.6,0){$N$};
\end{tikzpicture}
\figcap{题图：A 在岸上、B 在水中，C 为待选择的入水点；两垂足间距为 $L$。}\label{ch03:ref:question-figure-2-31}
\end{minipage}\end{center}
}


一维运动只需一个坐标。质点在平面或空间中运动时，位置、位移、速度和加速度都需要同时包含多个方向的信息。多数常见运动先可限制在一个平面内讨论；掌握二维情况后，只需增加第三个分量即可推广到三维。

这里的困难并不只是“多写一个坐标”。在直线上，正、负两个方向已经由数轴规定；到了平面内，同样大小的位移可以指向无穷多个方向，速度的大小不变时方向也可能不断改变。因此，描述平面运动必须把“大小”和“方向”同时纳入运算。抛体、车辆转弯和行星绕中心运动都属于这一类问题，它们虽然轨迹不同，却可以用同一套坐标、矢量和分量方法处理。

\subsection{位置的确定：坐标系}

平面直角坐标系由原点 $O$ 和两条互相垂直、带正方向的坐标轴组成，点 $P$ 由 $(x,y)$ 唯一确定。平面极坐标系由极点 $O$、参考方向 $ON$ 和坐标 $(r,\theta)$ 组成，其中 $r\ge0$；角度相差 $2\pi$ 的整数倍表示同一点，且 $r=0$ 时角度没有实际意义。

日常语言里也常用这两种定位思路。例如“从某地向西若干千米、再向北若干千米”对应直角坐标分量；“位于某地北偏西某角度、距离若干千米”对应极坐标。直角坐标中的 $x<0$ 只说明点在 $y$ 轴左侧，$y<0$ 只说明点在 $x$ 轴下方；极坐标中的 $r$ 是距离，通常取非负值，而 $\theta$ 的零方向和正向必须事先约定。

\CoordinateSystemsDiagram

坐标系必须在列方程前说明。直角坐标适合方向固定或边界平直的问题；极坐标适合具有明显中心的运动，如行星公转和圆周运动。坐标系不同不会改变实际运动，但会改变分量表达式的复杂程度。

完整的坐标系约定至少包括原点位置、各轴正方向和单位；极坐标还必须说明参考射线及角度正向。只有这些信息给全，$(x,y)$ 或 $(r,\theta)$ 才能被还原为唯一的空间位置。选择坐标系具有任意性，却不是随意性：同一问题可以有许多合法选法，其中与边界、运动方向或对称中心一致的选法往往使后续方程最简单。

两种坐标表示可以互相转换。在约定 $x$ 轴正向为极角零方向、逆时针为正时，
\begin{equation}
 x=r\cos\theta,\qquad y=r\sin\theta,\qquad
 r=\sqrt{x^2+y^2}.
\end{equation}
由 $x,y$ 求 $\theta$ 时要根据象限确定，而不能只使用 $\arctan(y/x)$ 的主值。极点处 $r=0$，所有角度都给出同一点，所以此处的 $\theta$ 不可由位置唯一确定。这种坐标奇点是表示方法的特点，不是实际空间出现了异常。

直角坐标适合分量方向固定的问题，极坐标适合“离中心多远、绕中心转了多少”的问题。例如同一圆周在直角坐标中需写 $x^2+y^2=R^2$，在极坐标中只需写 $r=R$；但匀速直线运动若不通过极点，用极坐标表示反而可能更复杂。选系的标准是让约束和运动规律尽量简单，而不是固定偏爱某一种。

\begin{origexamplebox}{在直角坐标系中标点}\label{ch03:ex:k03-2-22}
在平面直角坐标系中标出
\[
A(3,4),\ B(5,0),\ C(-3,2),\ D(2,-5),\ E(-3,-2),\ F(0,3),\ G(-1,0),\ H(0,-3).
\]
\par 题图见图~\ref{ch03:ref:question-figure-2-22}。
\KQuestionAW
\end{origexamplebox}
\Needspace{6\baselineskip}\begin{xesolution}
以下推导所用的解图为图~\ref{mech:legacy:ch03:k03-2-22}。

横坐标决定左右，纵坐标决定上下；坐标为零表示点在相应坐标轴上。图中 $A,C,E,D$ 分别位于第一、第二、第三、第四象限，$B,G$ 在 $x$ 轴上，$F,H$ 在 $y$ 轴上。
\begin{center}\begin{qualityfigure}
\begin{tikzpicture}[scale=.52]
\draw[kaxis] (-5.8,0)--(5.8,0) node[right,clearlabel]{$x$};\draw[kaxis] (0,-5.8)--(0,5.8) node[above,clearlabel]{$y$};
\draw[step=1,RuleGray!65,thin] (-5,-5) grid (5,5);
\foreach \name/\x/\y/\pos in {A/3/4/above,B/5/0/above,C/-3/2/above,D/2/-5/below,E/-3/-2/below,F/0/3/right,G/-1/0/below,H/0/-3/right}{\node[point,label={[clearlabel]\pos:$\name$}] at (\x,\y) {};}
\end{tikzpicture}
\figcap{本题解答中的位置、分量与辅助量。}\label{mech:legacy:ch03:k03-2-22}
\end{qualityfigure}\end{center}
\end{xesolution}


\begin{origexamplebox}{在极坐标系中标点}\label{ch03:ex:k03-2-23}
在平面极坐标系中标出
\[
A(5,\pi/3),\ B(2,\pi/3),\ C(5,\pi/2),\ D(4,\pi),\ E(4,3\pi/2),\ F(3,0).
\]
\par 题图见图~\ref{ch03:ref:question-figure-2-23}。
\KQuestionAX
\end{origexamplebox}
\Needspace{6\baselineskip}\begin{xesolution}
以下推导所用的解图为图~\ref{mech:legacy:ch03:k03-2-23}。

先约定从正 $x$ 轴逆时针量角，再按角度画射线，沿射线量取半径。
$A,B$ 都在 $60^\circ$ 射线上，距原点分别为 $5,2$。
$C$ 的极角为 $\pi/2=90^\circ$，在正 $y$ 轴；$E$ 的极角为 $3\pi/2=270^\circ$，在负 $y$ 轴。
$D$ 在负 $x$ 轴，$F$ 在正 $x$ 轴。也可用 $y_C=5\sin(\pi/2)=5$ 检查 $C$ 的上下位置。
\begin{center}\begin{qualityfigure}
\begin{tikzpicture}[scale=.55]
\draw[kaxis] (-5.8,0)--(5.8,0) node[right,clearlabel]{$x$};\draw[kaxis] (0,-5.8)--(0,5.8) node[above,clearlabel]{$y$};
\foreach \r in {1,...,5}{\draw[RuleGray] (0,0) circle (\r);}
\foreach \name/\r/\a/\pos in {A/5/60/above,B/2/60/above,C/5/90/right,D/4/180/above,E/4/270/left,F/3/0/above}{\node[point,label={[clearlabel]\pos:$\name$}] at (\a:\r) {};}
\end{tikzpicture}
\figcap{本题解答中的位置、分量与辅助量。}\label{mech:legacy:ch03:k03-2-23}
\end{qualityfigure}\end{center}
\end{xesolution}

\subsection{位移、矢量及其分量}

质点从起点到终点的\textbf{位移}是由起点指向终点的有向线段。位移只依赖始末位置，路程则依赖实际轨迹。沿曲线前进一小段时，位移大小与路程之差是高阶小量；经过较长路程后，两者可能相差很大。

因此，给定初位置和一段过程的位移就能唯一确定末位置，却不能由此还原质点实际经过的道路。反过来，即使两次运动的轨迹完全不同，只要起点和终点相同，它们的位移就相同。把初、末坐标写成 $(x_1,y_1)$ 与 $(x_2,y_2)$，坐标增量 $\Delta x=x_2-x_1$、$\Delta y=y_2-y_1$ 正好给出位移在两轴上的分量。

\PathChoicesDiagram

图~\ref{ch03:fig:k03-2-6}中的五条路径具有相同的起点 $A$ 和终点 $B$，却给出不同路程；图~\ref{ch03:fig:k03-2-7}中的每根箭头都从同一初始点 $A$ 指向一个时刻的位置 $B_i$，直接表示该时刻相对初态的位移。后一图中连接各 $B_i$ 的粗曲线是实际路径，直箭头并不表示沿这条路径走过的路程。

\DisplacementFanDiagram

既有大小又有方向的量称为矢量，只需大小的量称为标量。两个矢量相等，必须大小相等且方向相同。标量 $c$ 与矢量 $\vect A$ 相乘得到 $c\vect A$；$c<0$ 时方向反转。矢量加法遵循首尾相接法和平行四边形法：
\begin{equation}
\vect C=\vect A+\vect B,
\qquad
\vect A-\vect B=\vect A+(-\vect B).
\end{equation}
内积是标量：
\begin{equation}
  \vect A\cdot\vect B=AB\cos\theta.
\end{equation}

\VectorOperationsDiagram

图~\ref{ch03:fig:k03-2-8}中，加法的首尾相接次序可以交换，最后的合矢量不变；减法则等于加上反矢量。内积只留下两个矢量彼此平行方向上的信息：夹角为锐角时为正，直角时为零，钝角时为负。矢量运算不能只对大小套用普通代数，例如 $|\vect A+\vect B|$ 一般既不等于 $A+B$，也不等于 $|A-B|$，还必须知道夹角。

在直角坐标系中，$\vect A=(A_x,A_y)$。标量乘法、加减法和内积可逐分量计算：
\begin{align}
(c\vect A)_x&=cA_x,&(c\vect A)_y&=cA_y,\\
(\vect A\pm\vect B)_x&=A_x\pm B_x,&
(\vect A\pm\vect B)_y&=A_y\pm B_y,\\
\vect A\cdot\vect B&=A_xB_x+A_yB_y.
\end{align}

\VectorComponentsDiagram

分量表示与箭头表示的是同一个矢量，而不是两套不同的物理量。把矢量平移到原点不会改变它，因为自由矢量只关心大小和方向；终点坐标便给出 $(A_x,A_y)$。反过来，已知两个分量，可以沿坐标轴画出直角三角形，合矢量大小和方向为
\begin{equation}
 A=\sqrt{A_x^2+A_y^2},\qquad
 \tan\phi=\frac{A_y}{A_x},
\end{equation}
其中象限不能只由反正切的数值判断，还必须结合 $A_x,A_y$ 的符号。分量可以为负，而矢量大小始终非负。

逐分量运算的优势在于方向信息已经编码在代数符号中。两矢量相加时，分别把所有 $x$ 分量相加、所有 $y$ 分量相加；内积则把对应分量乘积相加。若更换坐标轴，分量数值会改变，但合矢量的大小、两矢量夹角和内积等几何量不变。计算后可用这种不变性作交叉检查。

若初、末位置分别为 $\vect r_1=(x_1,y_1)$、$\vect r_2=(x_2,y_2)$，则
\begin{equation}
\Delta\vect r=\vect r_2-\vect r_1=(x_2-x_1,\,y_2-y_1).
\end{equation}

位置矢量从坐标原点指向质点，位移矢量从初位置指向末位置。前者会随坐标原点改变，后者在只平移原点时不变。书写 $\Delta\vect r=\vect r_2-\vect r_1$ 时，减法次序决定箭头方向；交换始末点会得到大小相同、方向相反的位移。分段运动中每一段位移按时间顺序首尾相接，矢量和等于全过程位移，这个结论与各段轨迹是否直线无关。

\begin{origexamplebox}{三个时刻的位移关系}\label{ch03:ex:k03-2-24}
质点在 $t_1,t_2,t_3$ 时刻依次位于 $A,B,C$。画出 $t_1\to t_2$、$t_2\to t_3$ 和 $t_1\to t_3$ 的位移，并说明结论。

\begin{qualityfigure}
\begin{tikzpicture}[>=Latex,font=\small]
\coordinate (A) at (0,0);\coordinate (B) at (2.5,2);\coordinate (C) at (5,.5);
\foreach \p/\t in {A/A(t_1),B/B(t_2),C/C(t_3)}{\fill (\p) circle(2pt) node[above]{$\t$};}
\draw[dashed] (A)..controls(.5,2.5)and(1.5,3)..(B)..controls(4,3)and(3,-.5)..(C);
\end{tikzpicture}
\figcap{三个时刻的位置与真实路径}\label{mech:03:24:q}
\end{qualityfigure}
\end{origexamplebox}
\Needspace{6\baselineskip}\begin{xesolution}
以下推导所用的解图为图~\ref{mech:legacy:ch03:k03-2-24}。

三个位移分别为 $\overrightarrow{AB}$、$\overrightarrow{BC}$、$\overrightarrow{AC}$。按矢量首尾相接法，
\[
  \boxed{\overrightarrow{AB}+\overrightarrow{BC}=\overrightarrow{AC}}.
\]
这说明分段位移的矢量和等于全过程位移，与中间经过何种轨迹无关。
\begin{center}\begin{qualityfigure}
\begin{tikzpicture}[scale=.8]
\coordinate (A) at (-2.4,-.7);\coordinate (B) at (.2,1.2);\coordinate (C) at (2.8,.2);
\node[point,label={[clearlabel]left:$A$}] at (A) {};\node[point,label={[clearlabel]above:$B$}] at (B) {};\node[point,label={[clearlabel]right:$C$}] at (C) {};
\draw[vector] (A)--(B);\draw[vector] (B)--(C);\draw[Blue,very thick,->] (A)--(C);
\end{tikzpicture}
\figcap{本题解答中的位置、分量与辅助量。}\label{mech:legacy:ch03:k03-2-24}
\end{qualityfigure}\end{center}
\end{xesolution}

\begin{origexamplebox}{由两个累计位移求分段位移}\label{ch03:ex:k03-2-25}
已知质点在 $t_1\to t_2$ 的位移为 $\vect D_1$，在 $t_1\to t_3$ 的位移为 $\vect D_2$，求 $t_2\to t_3$ 的位移。

\begin{qualityfigure}
\begin{tikzpicture}[>=Latex,font=\small]
\coordinate (A) at (0,0);\coordinate (B) at (2.5,2);\coordinate (C) at (5,.5);
\foreach \p/\t in {A/A(t_1),B/B(t_2),C/C(t_3)}{\fill (\p) circle(2pt) node[above]{$\t$};}
\draw[dashed] (A)..controls(.5,2.5)and(1.5,3)..(B)..controls(4,3)and(3,-.5)..(C);\draw[->,DeepBlue,thick] (A)--(B) node[midway,above left]{$\vect D_1$};\draw[->,DeepBlue,thick] (A)--(C) node[midway,below]{$\vect D_2$};
\end{tikzpicture}
\figcap{累计位移共用起点，分段位移连接终点}\label{mech:03:25:q}
\end{qualityfigure}
\end{origexamplebox}
\Needspace{6\baselineskip}\begin{xesolution}
\begin{qualityfigure}
\begin{tikzpicture}[>=Latex,font=\small]
\coordinate (A) at (0,0);\coordinate (B) at (2.5,2);\coordinate (C) at (5,.5);
\foreach \p/\t in {A/A(t_1),B/B(t_2),C/C(t_3)}{\fill (\p) circle(2pt) node[above]{$\t$};}
\draw[dashed] (A)..controls(.5,2.5)and(1.5,3)..(B)..controls(4,3)and(3,-.5)..(C);\draw[->,DeepBlue,thick] (A)--(B) node[midway,above left]{$\vect D_1$};\draw[->,DeepBlue,thick] (A)--(C) node[midway,below]{$\vect D_2$};\draw[->,Gold,thick] (B)--(C) node[midway,above right]{$\vect D_{23}$};
\end{tikzpicture}
\figcap{累计位移共用起点，分段位移连接终点（解答用辅助图）}\label{mech:03:25:a}
\end{qualityfigure}
图~\ref{mech:03:25:a}中，$A$ 是两个累计位移的共同起点；从 $B$ 指向 $C$ 才是待求分段位移，因此应作矢量差 $\vect D_2-\vect D_1$。

由全过程等于分段之和，$\vect D_2=\vect D_1+\vect D_{23}$，故
\[
  \boxed{\vect D_{23}=\vect D_2-\vect D_1}.
\]
矢量相减必须同时处理方向，不能只把两个大小相减。
\end{xesolution}

\begin{origexamplebox}{坐标与位移相加}\label{ch03:ex:k03-2-26}
质点初始坐标为 $(2,3)\,\mathrm m$，随后位移为 $(4,5)\,\mathrm m$，求最终位置。

\begin{qualityfigure}
\begin{tikzpicture}[>=Latex,font=\small]
\draw[->] (0,0)--(4.3,0) node[right]{$x$};\draw[->] (0,0)--(0,5) node[above]{$y$};
\coordinate (A) at (1.2,1.8);\coordinate (B) at (3.6,4.8);
\fill (A) circle(2pt) node[left]{$A(2,3)$};
\draw[->,DeepBlue,thick] (A)--(B) node[midway,right]{$\Delta\vect r=(4,5)\,\mathrm m$};
\node[below left] at (0,0){$O$};\fill (B) circle(2pt) node[above]{$B$};
\end{tikzpicture}
\figcap{初始位置、位移与最终位置}\label{mech:03:26:q}
\end{qualityfigure}
\end{origexamplebox}
\Needspace{6\baselineskip}\begin{xesolution}
\begin{qualityfigure}
\begin{tikzpicture}[>=Latex,font=\small]
\draw[->] (0,0)--(4.3,0) node[right]{$x$};\draw[->] (0,0)--(0,5) node[above]{$y$};
\coordinate (A) at (1.2,1.8);\coordinate (B) at (3.6,4.8);
\fill (A) circle(2pt) node[left]{$A(2,3)$};
\draw[->,DeepBlue,thick] (A)--(B) node[midway,right]{$\Delta\vect r=(4,5)\,\mathrm m$};
\node[below left] at (0,0){$O$};\fill (B) circle(2pt) node[above]{$B(6,8)$};\draw[->,Gold] (0,0)--(A) node[midway,left]{$\vect r_1$};\draw[->,Gold] (0,0)--(B) node[pos=.70,above left=10pt]{$\vect r_2$};\draw[dashed] (B)--(3.6,0) node[below]{$6$};\draw[dashed] (B)--(0,4.8) node[left]{$8$};
\end{tikzpicture}
\figcap{初始位置、位移与最终位置（解答用辅助图）}\label{mech:03:26:a}
\end{qualityfigure}
图~\ref{mech:03:26:a}中，$O\!A$ 是初始位置矢量，$A\!B$ 是随后位移，二者首尾相接得到 $O\!B$；横、纵坐标须分别相加。

\[
\vect r_2=\vect r_1+\Delta\vect r=(2,3)+(4,5)=(6,8)\,\mathrm m.
\]
故最终位置为 $\boxed{(6,8)\,\mathrm m}$。
\end{xesolution}

\subsection{速度}

平面运动的平均速度和瞬时速度为
\begin{equation}
\bar{\vect v}=\frac{\Delta\vect r}{\Delta t},\qquad
\vect v=\lim_{\Delta t\to0}\frac{\Delta\vect r}{\Delta t}=\frac{\dd\vect r}{\dd t}.
\end{equation}
瞬时速度方向沿轨迹切线，大小等于瞬时速率。已知某时刻位置和速度，在极短时间内有 $\vect r(t+\Delta t)\simeq\vect r(t)+\vect v(t)\Delta t$。若速度的大小和方向始终不变，运动是匀速直线运动，且
\begin{equation}
  \vect r(t)=\vect r_0+\vect vt.
\end{equation}

平均速度只由总位移除以总时间得到。一个质点沿闭合曲线回到起点时，平均速度为零，但平均速率显然不为零；这与一维运动中往返的区别完全一致。瞬时速度的两个分量分别为 $v_x=\dd x/\dd t$、$v_y=\dd y/\dd t$，速率为 $\sqrt{v_x^2+v_y^2}$。只要任一分量随时间改变，速度矢量就可能改变；即使速率恒定，分量仍可因方向转动而不断变化。

匀速直线运动要求速度的大小和方向同时保持不变。若只知道速率恒定，轨迹仍可能弯曲；圆周运动就是典型反例。矢量方程 $\vect r=\vect r_0+\vect vt$ 表明位置矢量的终点沿固定方向均匀推进，所以轨迹必为直线，并且相等时间内的位移矢量完全相同。

\CurveVelocityDiagram

平均速度的方向由这段时间的总位移决定，不一定与轨迹中任一点的切线重合。让时间间隔逐渐缩小，末点沿轨迹逼近初点，位移弦方向才趋向切线方向。式 $\vect r(t+\Delta t)\simeq\vect r(t)+\vect v(t)\Delta t$ 是局部直线近似：它只在足够短的时间内成立，时间拉长后，速度方向或大小的变化会累积成不可忽略的偏差。

\begin{origexamplebox}{飞机紧急改航}\label{ch03:ex:k03-2-27}
飞机以 $800\,\mathrm{km/h}$ 沿北偏东 $53^\circ$ 飞行 $1\,\mathrm h$ 后，需在 $0.5\,\mathrm h$ 内到达出发机场正东 $1000\,\mathrm{km}$ 的机场 $B$。求改航后的速度大小和方向。取 $\sin53^\circ=4/5$。

\begin{qualityfigure}
\begin{tikzpicture}[>=Latex,font=\small]
\draw[->] (0,0)--(6,0) node[right]{东 $x$};\draw[->] (0,0)--(0,3.8) node[above]{北 $y$};
\coordinate (P) at (3.2,2.4);\coordinate (B) at (5,0);
\fill (0,0) circle(2pt) node[below left]{$O$};\fill (P) circle(2pt) node[above]{$P$};\fill (B) circle(2pt) node[below]{$B$};
\draw[->,DeepBlue,thick] (0,0)--(P) node[midway,above left]{$800\,\mathrm{km}$};
\draw (0,.8) arc(90:36.87:.8) node[midway,above]{$53^\circ$};
\draw[<->] (0,-.7)--(5,-.7) node[midway,below]{$1000\,\mathrm{km}$};\draw[dashed] (P)--(B) node[midway,right]{改航};
\end{tikzpicture}
\figcap{改航前后的端点与东西、南北分量}\label{mech:03:27:q}
\end{qualityfigure}
\end{origexamplebox}
\Needspace{6\baselineskip}\begin{xesolution}
\begin{qualityfigure}
\begin{tikzpicture}[>=Latex,font=\small]
\draw[->] (0,0)--(6,0) node[right]{东 $x$};\draw[->] (0,0)--(0,3.8) node[above]{北 $y$};
\coordinate (P) at (3.2,2.4);\coordinate (B) at (5,0);
\fill (0,0) circle(2pt) node[below left]{$O$};\fill (P) circle(2pt) node[above]{$P$};\fill (B) circle(2pt) node[below]{$B$};
\draw[->,DeepBlue,thick] (0,0)--(P) node[midway,above left]{$800\,\mathrm{km}$};
\draw (0,.8) arc(90:36.87:.8) node[midway,above]{$53^\circ$};
\draw[<->] (0,-.7)--(5,-.7) node[midway,below]{$1000\,\mathrm{km}$};\draw[->,Gold,thick] (P)--(B) node[midway,right]{$\Delta\vect r$};\draw[dashed] (P)--(3.2,0) node[below]{$640$};\node[left] at (0,2.4){$480$};\draw[dashed] (0,2.4)--(P);
\end{tikzpicture}
\figcap{改航前后的端点与东西、南北分量（解答用辅助图）}\label{mech:03:27:a}
\end{qualityfigure}
图~\ref{mech:03:27:a}中，$P$ 是改航点，$B$ 是目标；$\overrightarrow{PB}$ 的东向分量是 $1000-640$，北向分量是 $0-480$。不能把从 $O$ 到 $B$ 的总位移除以最后半小时。

取东为 $x$ 正向、北为 $y$ 正向。“北偏东 $53^\circ$”的分量为
\[
\vect v_1=(800\sin53^\circ,800\cos53^\circ)=(640,480)\,\mathrm{km/h}.
\]
飞行 $1\,\mathrm h$ 后位置为 $(640,480)\,\mathrm{km}$，目标为 $(1000,0)\,\mathrm{km}$，余下位移 $(360,-480)\,\mathrm{km}$。除以 $0.5\,\mathrm h$：
\[
\vect v_2=(720,-960)\,\mathrm{km/h},\qquad
\boxed{|\vect v_2|=1200\,\mathrm{km/h}}.
\]
方向为东偏南 $\arctan(960/720)\approx53.1^\circ$。
\end{xesolution}

\begin{origexamplebox}{潜艇发射鱼雷}\label{ch03:ex:k03-2-28}
潜艇在距敌船航线 $MN$ 的垂直距离 $L$ 处。敌船沿 $MN$ 以速率 $v$ 匀速航行，鱼雷以速率 $u$ 沿直线航行。若鱼雷要垂直击中船体，应在视线与航线垂线夹角为 $\theta$ 时发射。求 $\theta$。
\par 题图见图~\ref{ch03:ref:question-figure-2-28}。
\KQuestionBC
\end{origexamplebox}
\Needspace{6\baselineskip}\begin{xesolution}
以下推导所用的解图为图~\ref{mech:legacy:ch03:k03-2-28}。

如图，$A$ 为潜艇，$B$ 为发射时敌船的位置，$H$ 为 $A$ 到航线的垂足。
鱼雷沿 $AH$ 走距离 $L$，飞行时间 $t=L/u$。这段时间敌船沿 $BH$ 前进
$vt=vL/u$，所以发射时必须有 $BH=vL/u$。由直角三角形 $ABH$，
\[
  \tan\theta=\frac{vL/u}{L}=\frac vu,\qquad
  \boxed{\theta=\arctan\frac vu}.
\]
图中 $u$ 与 $v$ 必须分别对应鱼雷和船的地面速度；交换二者会得到互为倒数的错误结果。
\begin{center}\begin{qualityfigure}
\begin{tikzpicture}[scale=.92,>=Latex]
\coordinate (A) at (0,-1.5);\coordinate (B) at (-2,1.8);\coordinate (H) at (0,1.8);
\draw[very thick] (-3,1.8)--(2.3,1.8) node[right,clearlabel]{$MN$};
\draw[thick] (A)--(H);
\draw[kguide] (A)--(B);
\node[point,label={[clearlabel]below:$A$}] at (A) {};
\node[point,label={[clearlabel]below left:$B$}] at (B) {};
\node[point,label={[clearlabel]above right:$H$}] at (H) {};
\draw[TextGray] (-.19,1.8)--(-.19,1.61)--(0,1.61);
\draw[vector] (-2,2.35)--(-.8,2.35) node[midway,above]{$v$};
\draw[vector] (.42,-1.0)--(.42,.2) node[midway,right]{$u$};
\draw[<->,TextGray] (1.55,-1.5)--(1.55,1.8) node[midway,right]{$L=AH$};
\draw[<->,TextGray] (-2,3.0)--(0,3.0) node[midway,above]{$BH=vt$};
\draw[->,TextGray] (0,-.63) arc[start angle=90,end angle=121.22,radius=.87];
\node[clearlabel] (th) at (-1.45,-.95){$\theta$};
\draw[TextGray,thin] (th.east)--(-.23,-.66);
\end{tikzpicture}
\figcap{本题解答中的位置、分量与辅助量。}\label{mech:legacy:ch03:k03-2-28}
\end{qualityfigure}\end{center}
\end{xesolution}

\begin{origexamplebox}{足球抢断条件}\label{ch03:ex:k03-2-29}
球从 $A$ 沿射线 $AX$ 以速率 $v$ 匀速滚动。守方队员初始在 $B$，$AB=d$，$\angle BAX=\theta<90^\circ$，队员以速率 $v_p$ 沿直线奔跑。忽略边界，求队员能抢到球的必要条件。
\par 题图见图~\ref{ch03:ref:question-figure-2-29}。
\KQuestionBD
\end{origexamplebox}
\Needspace{6\baselineskip}\begin{xesolution}
以下推导所用的解图为图~\ref{ch03:fig:k03-ex-2-29}。

若在距 $A$ 为 $r\ge0$ 的点相遇，球用时 $r/v$；队员路程由余弦定理为
\[
\ell_p=\sqrt{d^2+r^2-2dr\cos\theta}.
\]
\FootballInterceptionDiagram

必须存在 $r>0$ 使 $\ell_p/v_p\le r/v$。令 $y=d/r>0$，平方后所需速度比满足
\[
\frac{v_p^2}{v^2}\ge y^2-2y\cos\theta+1
=(y-\cos\theta)^2+\sin^2\theta.
\]
因 $\theta<90^\circ$，可取 $y=\cos\theta$，即在图示射线上 $r=d/\cos\theta$ 的点 $C$ 相遇，从而右端最小值为 $\sin^2\theta$。因此必要且充分条件为
\[
  \boxed{v_p\ge v\sin\theta}.
\]
等号成立时，$BC\perp AB$，并非垂直于球的路线 $AX$。这是“用时相同条件下所需速率最小”，不等于“奔跑距离最短”。
具体地，取 $A=(0,0)$、$B=(d\cos\theta,d\sin\theta)$，最优相遇点为 $C=(d/\cos\theta,0)$，则
\[
 \overrightarrow{AB}\cdot\overrightarrow{BC}
 =d\cos\theta\left(\frac d{\cos\theta}-d\cos\theta\right)
   -d^2\sin^2\theta=0.
\]
所以图中的直角应画在 $B$ 点。此时 $AC=d/\cos\theta$、
$BC=d\tan\theta$，两段距离之比为 $\sin\theta$，再次得到 $v_p/v=\sin\theta$。
速度比无量纲，且当 $\theta\to90^\circ$ 时所需速率趋于球速；当 $\theta\to0^+$ 时，队员接近球的路线，所需速率趋于零。
这里取 $d>0$、$0<\theta<90^\circ$；$\theta=0$ 时队员已在球前方路线上，等待即可截球。
若 $v_p<v\sin\theta$，任何候选相遇点都来不及到达。
\end{xesolution}

\begin{origexamplebox}{两匀速质点的最短距离}\label{ch03:ex:k03-2-30}
质点 $A,B$ 同时从相距 $l$ 的 $P,Q$ 出发。$A$ 以速度 $v_1$ 沿 $PQ$ 方向运动，$B$ 以速度 $v_2$ 沿 $QR$ 运动，$QR$ 与 $QP$ 的夹角为 $\alpha$。求以后两质点的最短距离。
\par 题图见图~\ref{ch03:ref:question-figure-2-30}。
\KQuestionBE
\end{origexamplebox}
\Needspace{6\baselineskip}\begin{xesolution}
以下推导所用的解图为图~\ref{ch03:fig:k03-ex-2-30}。

取 $P\to Q$ 为 $x$ 正向。初始相对位置 $\vect r_0=(l,0)$，相对速度
\[
\vect u=\vect v_B-\vect v_A=(-v_2\cos\alpha-v_1,\ v_2\sin\alpha).
\]
\TwoParticlesMinimumDistanceDiagram

记两质点间距离为 $d(t)$。经过同一时间 $t$ 后，相对位置为 $\vect r_0+\vect ut$，因而
\[
 d^2(t)=[l-(v_1+v_2\cos\alpha)t]^2+(v_2\sin\alpha\,t)^2
 =l^2-2l(v_1+v_2\cos\alpha)t+Dt^2,
\]
其中 $D=v_1^2+v_2^2+2v_1v_2\cos\alpha=|\vect u|^2$。
若 $D=0$，两质点速度完全相同，相对位置不变，最短距离就是 $l$。
以下取 $D>0$。距离与距离平方在非负区间同时取最小值，所以可以对二次式配方：
\[
 d^2(t)=D(t-t_*)^2+l^2-\frac{l^2(v_1+v_2\cos\alpha)^2}{D}.
\]
若 $v_1+v_2\cos\alpha>0$，顶点时刻在允许范围 $t\ge0$ 内，最近时刻为
\[
t_*=-\frac{\vect r_0\cdot\vect u}{|\vect u|^2}
=\frac{l(v_1+v_2\cos\alpha)}{D},
\]
代回并使用
$D-(v_1+v_2\cos\alpha)^2=v_2^2\sin^2\alpha$，得到
\[
\boxed{d_{\min}=\frac{lv_2|\sin\alpha|}{\sqrt D}}.
\]
若 $v_1+v_2\cos\alpha\le0$，则 $t_*\le0$，二次式在 $t\ge0$ 上不减，最短距离为初始值 $\boxed l$。
几何上，内部极小值满足 $(\vect r_0+\vect ut_*)\cdot\vect u=0$：
最近时刻的相对连线垂直于相对速度。时间式量纲为秒，距离式为米。
当两者沿初始连线相向而行时 $\alpha=0$，结果退化为碰撞距离零。
\end{xesolution}

\begin{origexamplebox}{奔跑与游泳的最短时间路径}\label{ch03:ex:k03-2-31}
岸上点 $A$ 到直岸的垂直距离为 $d_1$，水中落水点 $B$ 到岸的垂直距离为 $d_2$，两垂足沿岸相距 $L$。救援者奔跑速率 $v_1$、游泳速率 $v_2$。求最短时间路径。
\par 题图见图~\ref{ch03:ref:question-figure-2-31}。
\KQuestionBF
\end{origexamplebox}
\Needspace{6\baselineskip}\begin{xesolution}
以下推导所用的解图为图~\ref{ch03:fig:k03-ex-2-31}。

设入水点 $C$ 到 $A$ 的垂足沿岸距离为 $x$。总时间
\RescueRouteDiagram
\[
T(x)=\frac{\sqrt{d_1^2+x^2}}{v_1}
+\frac{\sqrt{d_2^2+(L-x)^2}}{v_2}.
\]
令 $\dd T/\dd x=0$：
\[
\frac{x}{v_1\sqrt{d_1^2+x^2}}
=\frac{L-x}{v_2\sqrt{d_2^2+(L-x)^2}}.
\]
图中的 $\theta_1,\theta_2$ 分别从 $C$ 处的岸线法线量起。
跑步段长为 $\sqrt{d_1^2+x^2}$，其中沿岸边为 $x$，所以
$\sin\theta_1=x/\sqrt{d_1^2+x^2}$；同理，
$\sin\theta_2=(L-x)/\sqrt{d_2^2+(L-x)^2}$。代入导数方程即得
\[
  \boxed{\frac{\sin\theta_1}{v_1}=\frac{\sin\theta_2}{v_2}}.
\]
这就把求极值转化为确定入水点 $C$。还需证明该驻点确实给出最短时间。
取 $d_1,d_2,L,v_1,v_2>0$。若 $x<0$ 或 $x>L$，把入水点向两垂足之间移近可同时缩短两段路程，因此只需考虑 $0\le x\le L$。
进一步，
\[
 T''(x)=\frac{d_1^2}{v_1(d_1^2+x^2)^{3/2}}
       +\frac{d_2^2}{v_2[d_2^2+(L-x)^2]^{3/2}}>0.
\]
而 $T'(0)<0$、$T'(L)>0$，故区间内恰有一个驻点，并且它是全局极小点。
通常 $v_1>v_2$，于是 $\sin\theta_1>\sin\theta_2$，在两角均为锐角时有 $\theta_1>\theta_2$：
应让较慢的水中路段更接近法线。若 $v_1=v_2$，则两角相等，最优路线恢复为连接 $A,B$ 的直线。
两边 $\sin\theta_i/v_i$ 的量纲均为时间除以长度。
\end{xesolution}

\subsection{速度的变化：加速度}

二维加速度同样是矢量：
\begin{equation}
\bar{\vect a}=\frac{\Delta\vect v}{\Delta t},\qquad
\vect a=\lim_{\Delta t\to0}\frac{\Delta\vect v}{\Delta t}=\frac{\dd\vect v}{\dd t}.
\end{equation}
速度的大小改变或方向改变都会产生非零加速度。短时间内
\begin{equation}
\vect v(t+\Delta t)\simeq\vect v(t)+\vect a(t)\Delta t.
\end{equation}
加速度方向是速度矢量终点移动的方向，不一定与质点位置变化方向相同。

在固定直角坐标系中，$\vect a=(\ddot x,\ddot y)=(\dot v_x,\dot v_y)$。若加速度为常矢量，各分量都是一维匀变速运动；若加速度方向随时间变化，则即使大小不变，也不能直接套用常加速度的位移公式。当 $\vect v\ne\vect0$ 时，判断质点是“加速”还是“减速”要看 $\vect a$ 在 $\vect v$ 方向上的投影：由 $\dd v/\dd t=(\vect a\cdot\vect v)/v$，内积大于零时速率增加，小于零时速率减小，等于零时瞬时速率变化率为零，非零加速度只有法向分量。在瞬时静止处速度方向未定，不能根据内积为零断言“只改变方向”；例如由静止做匀加速直线运动，初始内积虽为零，随后速率仍增加。

速度空间的图形有助于避免把加速度画错。把不同时刻的速度箭头都从同一点画出，箭头终点会形成一条“速度轨迹”；加速度方向就是这个终点运动的切线方向。它与质点在真实空间中的轨迹是两种不同的图，不能因为真实轨迹向右弯，就未经计算断定加速度一定向右。

平面匀加速运动可以直接积分。若 $\vect a$ 为常矢量，
\begin{equation}
 \vect v=\vect v_0+\vect at,
 \qquad
 \vect r=\vect r_0+\vect v_0t+\frac12\vect at^2.
\end{equation}
它们与一维公式形式相同，但每个量都是矢量。沿任何固定坐标轴投影后，都得到相应分量的一维匀变速方程。重力场中的抛体正是取 $\vect a=\vect g$ 的特例，下一节将用这一统一形式处理。

多阶段路线问题中，应先为每一段写出位移矢量，再按时间顺序相加；若求最短时间，还需把转折点作为变量建立总时间函数。
例~\ref{ch03:ex:k03-2-31}的结果形式与光的折射定律相似：两者都来自不同区域传播速率不同时的时间驻值条件。
一般有边界的优化问题仍需检查端点；本题则已由 $T''>0$ 和端点导数异号证明最优点位于两垂足之间。

所有矢量公式都应带着适用的参考系理解。同一位移、速度在旋转或平动的另一组坐标轴中分量可能不同，但由同一坐标系算出的分量必须成套使用。不能把一个矢量的 $x$ 分量取自地面系、$y$ 分量取自运动系再拼成所谓“合速度”。在后续相对运动中，先统一下标和参考系，再作矢量加减，是保证方向正确的第一步。

比较两个时刻的速度时，应先把两矢量的起点平移到同一点，再由旧速度的终点指向新速度的终点画出 $\Delta\vect v$。加速度沿 $\Delta\vect v$ 在极限中的方向，而不是简单指向轨迹的内侧或运动方向。速度大小改变、方向不变会产生切向加速度；速度大小不变、方向改变同样会产生加速度；二者也可以同时发生。

\AccelerationVectorDiagram
}

{
\graphicspath{{../../../sample/xes-original/kinematics-supplement/}{../../../sample/xes-original/kinematics-supplement/assets/}}
\section{抛体运动}
\newcommand{\ProjectileVectorsDiagram}{%
\begin{center}\begin{minipage}{.98\linewidth}\centering
\begin{tikzpicture}[scale=.9,>=Latex]
 \begin{scope}[xshift=-3.6cm]
  \coordinate (O) at (0,0);\coordinate (A) at (2,3);\coordinate (B) at (2,-1);
  \draw[vector] (O)--(A) node[midway,above left]{$\vect v_0$};
  \draw[vector] (A)--(B) node[midway,right=3pt]{$\vect g t$};
  \draw[vector] (O)--(B) node[midway,below left]{$\vect v$};
  \node[below] at (1,-1.45){(a) 速度三角形};
 \end{scope}
 \begin{scope}[xshift=3cm]
  \coordinate (O) at (0,0);\coordinate (A) at (2,3);\coordinate (B) at (2,1);\coordinate (C) at (2,-1);
  \draw[vector] (O)--(A) node[midway,above left]{$\vect v_0t$};
  \draw[vector] (A)--(B) node[midway,right=4pt]{$\tfrac12\vect g t^2$};
  \draw[vector] (O)--(B) node[pos=.6,below right=3pt]{$\vect r$};
  \draw[TextGray,dashed,->] (B)--(C);
  \draw[TextGray,dashed,->] (O)--(C) node[midway,below left]{$t\vect v$};
  \node[right=4pt,align=left,font=\small] at (2,-.2){再接一段\\$\tfrac12\vect g t^2$};
  \node[below] at (1,-1.45){(b) 位移及辅助三角形};
 \end{scope}
\end{tikzpicture}
\figcap{此刻速度向下、位移仍向上；右图加入虚线段后，外侧三角形才与左图相似。}\label{ch03:fig:k04-2-12}
\end{minipage}\end{center}}

\newcommand{\TennisServeDiagram}{%
\begin{center}\begin{minipage}{.95\linewidth}\centering
\begin{tikzpicture}[>=Latex]
 \draw[thick] (-.15,0)--(6.7,0);\draw[thick] (4,0)--(4,1.2);
 \draw[<->,TextGray] (4.25,0)--(4.25,1.2) node[midway,right=2pt]{$h$};
 \draw[kguide] (0,0)--(0,3) (6,0)--(6,.3);
 \draw[kcurve,->,domain=0:5.6,samples=90] plot(\x,{3-3*\x*\x/31.36});
 \draw[vector] (0,3)--(1.3,3) node[midway,above]{$v_0$};
 \draw[<->,TextGray] (-.5,0)--(-.5,3) node[midway,left]{$H$};
 \draw[<->,TextGray] (0,-.75)--(4,-.75) node[midway,below]{$L$};
 \draw[<->,TextGray] (4,-.75)--(6,-.75) node[midway,below]{$l$};
 \node[above left] at (0,3){$A$};\node[below=3pt] at (4,0){球网};\node[below=3pt] at (6,0){发球线};
\end{tikzpicture}
\figcap{合法轨迹从 $A$ 水平出发，从球网上方通过，并落在发球线之前。}\label{ch03:fig:k04-ex-2-32}
\end{minipage}\end{center}}

\newcommand{\HorizontalThrowImpactDiagram}{%
\begin{center}\begin{minipage}{.95\linewidth}\centering
\begin{tikzpicture}[>=Latex]
 \coordinate (A) at (0,2);\coordinate (B) at (4,0);\coordinate (V) at (5.1,-1.1);\coordinate (X) at (5.3,0);
 \draw[thick] (-.2,0)--(4.35,0);\draw[kguide] (A)--(0,0) (B)--(X);
 \draw[kcurve,dashed,domain=0:4,samples=80] plot(\x,{2-\x*\x/8});
 \draw[vector] (A)--(1.3,2) node[midway,above]{$\vect v_0$};
 \draw[vector] (B)--(V) node[right=4pt]{$\vect v$};
 \draw pic[draw=TextGray,angle radius=6mm] {angle=V--B--X};
 \node[right] at (4.7,-.21){$45^\circ$};
 \draw[<->,TextGray] (-.5,0)--(-.5,2) node[midway,left]{$h$};
 \node[point] at (A){};\node[point] at (B){};
\end{tikzpicture}
\figcap{轨迹在抛出点水平，在落地点的切线与水平线成 $45^\circ$。}\label{ch03:fig:k04-ex-2-33}
\end{minipage}\end{center}}

\newcommand{\ProjectileRangeDiagram}{%
\begin{center}\begin{minipage}{.95\linewidth}\centering
\begin{tikzpicture}[scale=.95,>=Latex]
 \coordinate (A) at (0,0);\coordinate (B) at (7.2,0);\coordinate (V) at (52:1.5);
 \draw[thick] (-.4,0)--(7.6,0);
 \draw[kcurve,->,domain=0:7.2,samples=100] plot(\x,{tan(52)*\x*(1-\x/7.2)});
 \draw[vector] (A)--(V) node[above left]{$v_0$};
 \draw pic[draw=TextGray,angle radius=7mm] {angle=B--A--V};\node at (.97,.39){$\theta$};
 \draw[<->,TextGray] (0,-.55)--(7.2,-.55) node[midway,below]{$S$};
\end{tikzpicture}
\figcap{斜抛轨迹先升后降，回到出发高度时的水平距离为射程 $S$。}\label{ch03:fig:k04-ex-2-34}
\end{minipage}\end{center}}

\newcommand{\ComplementaryProjectileDiagram}{%
\begin{center}\begin{minipage}{.96\linewidth}\centering
\begin{tikzpicture}[scale=.9,>=Latex]
 \coordinate (A) at (0,0);\coordinate (B) at (6.4,0);
 \draw[thick] (-.3,0)--(6.8,0);
 \foreach \ang in {22,68}{\draw[kcurve,dashed,domain=0:6.4,samples=100] plot(\x,{tan(\ang)*\x*(1-\x/6.4)});
 \draw[vector] (A)--(\ang:1.3);}
 \draw[TextGray] (.7,0) arc[start angle=0,end angle=22,radius=.7];
 \draw[TextGray] (1.1,0) arc[start angle=0,end angle=68,radius=1.1];
 \node at (1.39,.95){$\theta_1$};\node at (1.55,.22){$\theta_2$};
 \node[above] at (3.2,4){高弧线};\node[above] at (3.2,.7){低弧线};
 \node[below] at (3.2,-.4){$\theta_1+\theta_2=\pi/2$，两次初速率相同};
\end{tikzpicture}
\figcap{互余抛射角对应同一落点；高弧线飞行时间更长。}\label{ch03:fig:k04-ex-2-35}
\end{minipage}\end{center}}

\newcommand{\ObstacleClearanceDiagram}{%
\begin{center}\begin{minipage}{.96\linewidth}\centering
\begin{tikzpicture}[>=Latex]
 \pgfmathsetmacro{\uu}{(sqrt(25+2.25)+1.5)/5}
 \pgfmathsetmacro{\kk}{(5*\uu-1.5)/25}
 \draw[thick] (-.3,0)--(5.9,0) (5,0)--(5,1.5);
 \draw[kcurve,dashed,->,domain=0:5,samples=90] plot(\x,{\uu*\x-\kk*\x*\x});
 \draw[vector] (0,0)--({1.25/sqrt(1+\uu^2)},{1.25*\uu/sqrt(1+\uu^2)}) node[above left]{$v_0$};
 \draw[TextGray] (.7,0) arc[start angle=0,end angle={atan(\uu)},radius=.7];
 \node at (.95,.42){$\theta$};\node[below left] at (0,0){$A$};
 \draw[<->,TextGray] (0,-.55)--(5,-.55) node[midway,below]{$L$};
 \draw[<->,TextGray] (5.5,0)--(5.5,1.5) node[midway,right]{$H$};
\end{tikzpicture}
\figcap{临界轨迹经过障碍顶端；严格越过时须留出正的高度余量。}\label{ch03:fig:k04-ex-2-36}
\end{minipage}\end{center}}

\newcommand{\ProjectileReachabilityDiagram}{%
\begin{center}\begin{minipage}{.96\linewidth}\centering
\begin{tikzpicture}
\begin{axis}[width=.90\linewidth,height=7cm,axis lines=middle,xlabel={$x$},ylabel={$y$},
 xmin=0,xmax=7.2,ymin=-1,ymax=3.7,ticks=none,clip=true]
 \path[fill=Blue!7] (axis cs:0,-1)--plot[domain=0:6.6,samples=100]
 (axis cs:\x,{3-\x*\x/12})--(axis cs:6.6,-1)--cycle;
 \foreach \ang in {20,35,50,65,78}{
  \addplot[Blue!75,thick,domain=0:6.6,samples=140]{x*tan(\ang)-x*x/(12*cos(\ang)^2)};
 }
 \addplot[DeepBlue,very thick,domain=0:6.6,samples=100]{3-x*x/12};
 \node at (axis cs:5.7,2.8){不可达区};
 \node at (axis cs:3.9,-.5){可达区};
 \node at (axis cs:2.35,3.35){分界线};
 \draw[TextGray] (axis cs:2.35,3.08)--(axis cs:2.35,2.54);
\end{axis}
\end{tikzpicture}
\figcap{同一初速率的各条轨迹均在包络线下方，且分别与它相切。}\label{ch03:fig:k04-ex-2-37}
\end{minipage}\end{center}}

\newcommand{\HillsideTargetDiagram}{%
\begin{center}\begin{minipage}{.96\linewidth}\centering
\begin{tikzpicture}[>=Latex]
 \coordinate (C) at (0,0);\coordinate (T) at (5,1.5);\coordinate (X) at (5.8,0);
 \draw[thick] (C)--(X) (C)--(5.6,1.68);
 \draw[kcurve,dashed,->,domain=0:5,samples=100] plot(\x,{tan(55)*\x-(5*tan(55)-1.5)*\x*\x/25});
 \draw[vector] (C)--(55:1.3) node[above left]{$v_0$};
 \draw[TextGray] (.75,0) arc[start angle=0,end angle={atan(.3)},radius=.75];
 \draw[TextGray] (1.3,0) arc[start angle=0,end angle=55,radius=1.3];
 \node at (.99,.10){$\alpha$};\node at (1.48,.75){$\beta$};
 \draw[kguide] (T)--(5,0);
 \draw[<->,TextGray] (0,-.55)--(5,-.55) node[midway,below]{$L$};
 \node[point,label=below left:$C$] at (C){};\node[point,label=above right:$T$] at (T){};
\end{tikzpicture}
\figcap{目标坐标为 $(L,L\tan\alpha)$；炮弹初速度沿轨迹起点的切线。}\label{ch03:fig:k04-ex-2-39}
\end{minipage}\end{center}}

\newcommand{\HorizontalAtObstacleDiagram}{%
\begin{center}\begin{minipage}{.96\linewidth}\centering
\begin{tikzpicture}[>=Latex]
 \draw[thick] (-.3,0)--(6.5,0) (5,0)--(5,2);
 \draw[kcurve,dashed,domain=0:5,samples=100] plot(\x,{2-2*(\x-5)^2/25});
 \draw[vector] (0,0)--({1.3/sqrt(1.64)},{1.04/sqrt(1.64)}) node[above left]{$v_0$};
 \draw[vector] (5,2)--(6.3,2) node[midway,above]{$v_T$};
 \draw[TextGray] (.7,0) arc[start angle=0,end angle={atan(.8)},radius=.7];
 \node at (1.02,.31){$\theta$};
 \draw[<->,TextGray] (0,-.55)--(5,-.55) node[midway,below]{$L$};
 \draw[<->,TextGray] (5.5,0)--(5.5,2) node[midway,right]{$H$};
 \node[below left] at (0,0){$A$};\node[above left] at (5,2){$T$};
\end{tikzpicture}
\figcap{轨迹在障碍顶端 $T$ 达到最高点，此处速度为水平方向。}\label{ch03:fig:k04-ex-2-40}
\end{minipage}\end{center}}

\newcommand{\InclinedPlaneNormalImpactDiagram}{%
\begin{center}\begin{minipage}{.96\linewidth}\centering
\begin{tikzpicture}[scale=.95,>=Latex]
 \pgfmathsetmacro{\thetaLaunch}{20+atan(1/(2*tan(20)))}
 \pgfmathsetmacro{\uu}{tan(\thetaLaunch)}
 \pgfmathsetmacro{\kk}{(\uu-tan(20))/4.5}
 \coordinate (A) at (0,0);\coordinate (P) at (4.5,{4.5*tan(20)});
 \coordinate (B) at (5.6,{5.6*tan(20)});\coordinate (V) at ($(P)+(-70:1.2)$);
 \draw[thick] (A)--(5.9,0) (A)--(B)--(5.6,0);
 \draw[kcurve,dashed,domain=0:4.5,samples=100] plot(\x,{\uu*\x-\kk*\x*\x});
 \draw[vector] (A)--(\thetaLaunch:1.3) node[above left]{$v_0$};
 \draw[vector] (P)--(V) node[right=3pt]{$v_P$};
 \draw pic[draw=TextGray,angle radius=4mm] {right angle=A--P--V};
 \draw[TextGray] (.7,0) arc[start angle=0,end angle=20,radius=.7];
 \draw[TextGray] (1.3,0) arc[start angle=0,end angle=\thetaLaunch,radius=1.3];
 \node at (.96,.14){$\alpha$};\node at (1.42,.9){$\theta$};
 \node[point,label=above right:$P$] at (P){};
\end{tikzpicture}
\figcap{碰撞点 $P$ 在斜面上；轨迹末端切线与斜面垂直。}\label{ch03:fig:k04-ex-2-41}
\end{minipage}\end{center}}

\newcommand{\StairBounceDiagram}{%
\begin{center}\begin{minipage}{.98\linewidth}\centering
\begin{tikzpicture}[scale=1,>=Latex]
 \pgfmathsetmacro{\hh}{1/(1-.45^2)}\pgfmathsetmacro{\xa}{1.4-1.8/1.45}
 \draw[very thick] (-.1,3)--(2.3,3)--(2.3,2)--(4.1,2)--(4.1,1)--(5.9,1)--(5.9,0)--(7.8,0);
 \draw[kcurve,domain=0:1,samples=70] plot({\xa+(1.4-\xa)*\x},{3+\hh-\hh*\x*\x});
 \foreach \n in {0,1,2}{
  \draw[kcurve,domain=0:1,samples=90] plot({1.4+1.8*\n+1.8*\x},{3-\n+2*.45/(1-.45)*\x-(1+.45)/(1-.45)*\x*\x});
 }
 \foreach \n in {0,1,2,3}{\node[point,label=below:{$C_{\n}$}] at ({1.4+1.8*\n},{3-\n}){};}
 \draw[vector] (\xa,{3+\hh})--++(1.1,0) node[midway,above]{$v_0$};
 \node[above left] at (\xa,{3+\hh}){$A$};
 \draw[<->,TextGray] (-.45,3)--(-.45,{3+\hh}) node[midway,left]{$H$};
 \draw[<->,TextGray] (3.2,3.55)--(5,3.55) node[midway,above]{$d$};
 \draw[<->,TextGray] (7.4,0)--(7.4,1) node[midway,right]{$l$};
\end{tikzpicture}
\figcap{$A$ 在首个碰撞台面上方；各 $C_i$ 均在水平踏面上，反弹段可平移重合。}\label{ch03:fig:k04-ex-2-43}
\end{minipage}\end{center}}

\newcommand{\SphereClearanceDiagram}{%
\begin{center}\begin{minipage}{.98\linewidth}\centering
\begin{tikzpicture}[scale=1.2,>=Latex]
 \foreach \q/\shift/\lab in {.577350269/-2.8/{临界：先与球面相切},.70/2.8/{可行：到球顶前不接触}}{
  \begin{scope}[xshift=\shift cm]
   \pgfmathsetmacro{\ddist}{(\q+sqrt(\q*\q+32/3))/(8/3)}
   \draw[thick] (-2.05,0)--(1.3,0);
   \draw[thick] (0,1) circle (1);
   \draw[kguide] (0,0)--(0,2);
   \draw[kcurve,domain=-\ddist:0,samples=120] plot(\x,{2-\q*\x-4*\x*\x/3});
   \pgfmathsetmacro{\sl}{-\q+8*\ddist/3}
   \draw[vector] (-\ddist,0)--++({.65/sqrt(1+\sl^2)},{.65*\sl/sqrt(1+\sl^2)});
   \draw[vector] (0,2)--++({.7/sqrt(1+\q^2)},{-.7*\q/sqrt(1+\q^2)});
   \node[above right] at (.58,2.02){$\vect v_T$};
   \node[point,label=above left:$T$] at (0,2){};
   \node[point,label=below:$O$] at (0,1){};
   \draw[<->,TextGray] (-\ddist,-.4)--(0,-.4) node[midway,below]{$d$};
   \node[font=\small,align=center] at (-.4,-1.05){\lab};
   \ifdim\q pt<.6pt
    \node[point,label=left:$P$] at (-.8660254,1.5){};
   \fi
  \end{scope}
 }
\end{tikzpicture}
\figcap{两条轨迹均在下降时到达球顶 $T$；左图的提前切点 $P$ 决定速度下界。}\label{ch03:fig:k04-ex-2-45}
\end{minipage}\end{center}}

\newcommand{\KQuestionBG}{%
%
\begin{center}\begin{minipage}{.95\linewidth}\centering
\begin{tikzpicture}[>=Latex]
 \draw[thick] (-.15,0)--(8.2,0);\draw[thick] (4,0)--(4,1.2);
 \draw[<->,TextGray] (4.25,0)--(4.25,1.2) node[midway,right=2pt]{$h$};
 \draw[kguide] (0,0)--(0,3) (6,0)--(6,.3);
 \draw[kcurve,->,domain=0:5.6,samples=90] plot(\x,{3-3*\x*\x/31.36});
 \draw[vector] (0,3)--(1.3,3) node[midway,above]{$v_0$};
 \draw[<->,TextGray] (-.5,0)--(-.5,3) node[midway,left]{$H$};
 \draw[<->,TextGray] (0,-.75)--(4,-.75) node[midway,below]{$L$};
 \draw[<->,TextGray] (4,-.75)--(6,-.75) node[midway,below]{$l$};
 \node[above left] at (0,3){$A$};\node[below=3pt] at (4,0){球网};\node[below=3pt] at (6,0){发球线};
\draw[thick] (8,-.12)--(8,.25);
\draw[<->,TextGray] (0,-1.4)--(8,-1.4) node[midway,below]{$2L$};
\node[above] at (8,.25){底线};
\end{tikzpicture}
\figcap{题图：A 为水平击球点；全场长 $2L$，球网到对方发球线的距离为 $l$。}\label{ch03:ref:question-figure-2-32}
\end{minipage}\end{center}
}

\newcommand{\KQuestionBH}{%
%
\begin{center}\begin{minipage}{.95\linewidth}\centering
\begin{tikzpicture}[>=Latex]
 \coordinate (A) at (0,2);\coordinate (B) at (4,0);\coordinate (V) at (5.1,-1.1);\coordinate (X) at (5.3,0);
 \draw[thick] (-.2,0)--(4.35,0);\draw[kguide] (A)--(0,0) (B)--(X);
 \draw[kcurve,dashed,domain=0:4,samples=80] plot(\x,{2-\x*\x/8});
 \draw[vector] (A)--(1.3,2) node[midway,above]{$\vect v_0$};
 \draw[vector] (B)--(V) node[right=4pt]{$\vect v$};
 \draw pic[draw=TextGray,angle radius=6mm] {angle=V--B--X};
 \node[right] at (4.7,-.21){$45^\circ$};
 \draw[<->,TextGray] (-.5,0)--(-.5,2) node[midway,left]{$h$};
 \node[point] at (A){};\node[point] at (B){};
\end{tikzpicture}
\figcap{题图：从高度 $h$ 水平抛出，落地速度与水平线成 $45^\circ$。}\label{ch03:ref:question-figure-2-33}
\end{minipage}\end{center}
}

\newcommand{\KQuestionBI}{%
%
\begin{center}\begin{minipage}{.95\linewidth}\centering
\begin{tikzpicture}[scale=.95,>=Latex]
 \coordinate (A) at (0,0);\coordinate (B) at (7.2,0);\coordinate (V) at (52:1.5);
 \draw[thick] (-.4,0)--(7.6,0);
 \draw[kcurve,->,domain=0:7.2,samples=100] plot(\x,{tan(52)*\x*(1-\x/7.2)});
 \draw[vector] (A)--(V) node[above left]{$v_0$};
 \draw pic[draw=TextGray,angle radius=7mm] {angle=B--A--V};\node at (.97,.39){$\theta$};
 \draw[<->,TextGray] (0,-.55)--(7.2,-.55) node[midway,below]{$S$};
\end{tikzpicture}
\figcap{题图：初速度为 $v_0$、仰角为 $\theta$，抛出点与落点等高。}\label{ch03:ref:question-figure-2-34}
\end{minipage}\end{center}
}

\newcommand{\KQuestionBJ}{%
%
\begin{center}\begin{minipage}{.96\linewidth}\centering
\begin{tikzpicture}[scale=.9,>=Latex]
 \coordinate (A) at (0,0);\coordinate (B) at (6.4,0);
 \draw[thick] (-.3,0)--(6.8,0);
 \foreach \ang in {22,68}{\draw[kcurve,dashed,domain=0:6.4,samples=100] plot(\x,{tan(\ang)*\x*(1-\x/6.4)});
 \draw[vector] (A)--(\ang:1.3);}
 \draw[TextGray] (.7,0) arc[start angle=0,end angle=22,radius=.7];
 \draw[TextGray] (1.1,0) arc[start angle=0,end angle=68,radius=1.1];
 \node at (1.39,.95){$\theta_1$};\node at (1.55,.22){$\theta_2$};
 \node[above] at (3.2,4){高弧线};\node[above] at (3.2,.7){低弧线};
 \node[below] at (3.2,-.4){$\theta_1+\theta_2=\pi/2$，两次初速率相同};
\end{tikzpicture}
\figcap{题图：两次初速率相同，两个抛射角互余。}\label{ch03:ref:question-figure-2-35}
\end{minipage}\end{center}
}

\newcommand{\KQuestionBK}{%
\begin{center}\begin{minipage}{.97\linewidth}\centering
\begin{tikzpicture}[>=Latex,scale=.9]

\draw[thick] (-.3,0)--(6,0) (5,0)--(5,1.7);
\draw[vector] (0,0)--(52:1.5) node[above left]{$v_0$};
\draw (.7,0) arc(0:52:.7);\node at (1.03,.4){$\theta$};
\node[below left] at (0,0){$A$};
\draw[<->] (0,-.55)--(5,-.55) node[midway,below]{$L$};
\draw[<->] (5.5,0)--(5.5,1.7) node[midway,right]{$H$};

\end{tikzpicture}
\figcap{题图：A 到障碍物的水平距离为 $L$，障碍物高 $H$。}\label{ch03:ref:question-figure-2-36}
\end{minipage}\end{center}
}

\newcommand{\KQuestionBL}{%
%
\begin{center}\begin{minipage}{.96\linewidth}\centering
\begin{tikzpicture}
\begin{axis}[width=.90\linewidth,height=7cm,axis lines=middle,xlabel={$x$},ylabel={$y$},
 xmin=0,xmax=7.2,ymin=-1,ymax=3.7,ticks=none,clip=true]
 
 \foreach \ang in {20,35,50,65,78}{
  \addplot[Blue!75,thick,domain=0:6.6,samples=140]{x*tan(\ang)-x*x/(12*cos(\ang)^2)};
 }
\end{axis}
\end{tikzpicture}
\figcap{题图：从同一点以同一初速率、不同仰角抛出的轨迹族。}\label{ch03:ref:question-figure-2-37}
\end{minipage}\end{center}
}

\newcommand{\KQuestionBN}{%
%
\begin{center}\begin{minipage}{.96\linewidth}\centering
\begin{tikzpicture}[>=Latex]
 \coordinate (C) at (0,0);\coordinate (T) at (5,1.5);\coordinate (X) at (5.8,0);
 \draw[thick] (C)--(X) (C)--(5.6,1.68);
 \draw[kcurve,dashed,->,domain=0:5,samples=100] plot(\x,{tan(55)*\x-(5*tan(55)-1.5)*\x*\x/25});
 \draw[vector] (C)--(55:1.3) node[above left]{$v_0$};
 \draw[TextGray] (.75,0) arc[start angle=0,end angle={atan(.3)},radius=.75];
 \draw[TextGray] (1.3,0) arc[start angle=0,end angle=55,radius=1.3];
 \node at (.99,.10){$\alpha$};\node at (1.48,.75){$\beta$};
 \draw[kguide] (T)--(5,0);
 \draw[<->,TextGray] (0,-.55)--(5,-.55) node[midway,below]{$L$};
 \node[point,label=below left:$C$] at (C){};\node[point,label=above right:$T$] at (T){};
\end{tikzpicture}
\figcap{题图：坡脚 C、坡上目标 T；$\alpha$ 为坡角，$\beta$ 为发射仰角。}\label{ch03:ref:question-figure-2-39}
\end{minipage}\end{center}
}

\newcommand{\KQuestionBO}{%
%
\begin{center}\begin{minipage}{.96\linewidth}\centering
\begin{tikzpicture}[>=Latex]
 \draw[thick] (-.3,0)--(6.5,0) (5,0)--(5,2);
 \draw[kcurve,dashed,domain=0:5,samples=100] plot(\x,{2-2*(\x-5)^2/25});
 \draw[vector] (0,0)--({1.3/sqrt(1.64)},{1.04/sqrt(1.64)}) node[above left]{$v_0$};
 \draw[vector] (5,2)--(6.3,2) node[midway,above]{$v_T$};
 \draw[TextGray] (.7,0) arc[start angle=0,end angle={atan(.8)},radius=.7];
 \node at (1.02,.31){$\theta$};
 \draw[<->,TextGray] (0,-.55)--(5,-.55) node[midway,below]{$L$};
 \draw[<->,TextGray] (5.5,0)--(5.5,2) node[midway,right]{$H$};
 \node[below left] at (0,0){$A$};\node[above left] at (5,2){$T$};
\end{tikzpicture}
\figcap{题图：要求抛体在障碍顶端 T 处的速度平行于水平地面。}\label{ch03:ref:question-figure-2-40}
\end{minipage}\end{center}
}

\newcommand{\KQuestionBP}{%
%
\begin{center}\begin{minipage}{.96\linewidth}\centering
\begin{tikzpicture}[scale=.95,>=Latex]
 \pgfmathsetmacro{\thetaLaunch}{20+atan(1/(2*tan(20)))}
 \pgfmathsetmacro{\uu}{tan(\thetaLaunch)}
 \pgfmathsetmacro{\kk}{(\uu-tan(20))/4.5}
 \coordinate (A) at (0,0);\coordinate (P) at (4.5,{4.5*tan(20)});
 \coordinate (B) at (5.6,{5.6*tan(20)});\coordinate (V) at ($(P)+(-70:1.2)$);
 \draw[thick] (A)--(5.9,0) (A)--(B)--(5.6,0);
 \draw[kcurve,dashed,domain=0:4.5,samples=100] plot(\x,{\uu*\x-\kk*\x*\x});
 \draw[vector] (A)--(\thetaLaunch:1.3) node[above left]{$v_0$};
 \draw[vector] (P)--(V) node[right=3pt]{$v_P$};
 \draw pic[draw=TextGray,angle radius=4mm] {right angle=A--P--V};
 \draw[TextGray] (.7,0) arc[start angle=0,end angle=20,radius=.7];
 \draw[TextGray] (1.3,0) arc[start angle=0,end angle=\thetaLaunch,radius=1.3];
 \node at (.96,.14){$\alpha$};\node at (1.42,.9){$\theta$};
 \node[point,label=above right:$P$] at (P){};
\end{tikzpicture}
\figcap{题图：从斜面底端抛出，要求再次碰到斜面时速度垂直于斜面。}\label{ch03:ref:question-figure-2-41}
\end{minipage}\end{center}
}

\newcommand{\KQuestionBR}{%
%
\begin{center}\begin{minipage}{.98\linewidth}\centering
\begin{tikzpicture}[scale=1,>=Latex]
 \pgfmathsetmacro{\hh}{1/(1-.45^2)}\pgfmathsetmacro{\xa}{1.4-1.8/1.45}
 \draw[very thick] (-.1,3)--(2.3,3)--(2.3,2)--(4.1,2)--(4.1,1)--(5.9,1)--(5.9,0)--(7.8,0);
 \draw[kcurve,domain=0:1,samples=70] plot({\xa+(1.4-\xa)*\x},{3+\hh-\hh*\x*\x});
 \foreach \n in {0,1,2}{
  \draw[kcurve,domain=0:1,samples=90] plot({1.4+1.8*\n+1.8*\x},{3-\n+2*.45/(1-.45)*\x-(1+.45)/(1-.45)*\x*\x});
 }
 \foreach \n in {0,1,2,3}{\node[point,label=below:{$C_{\n}$}] at ({1.4+1.8*\n},{3-\n}){};}
 \draw[vector] (\xa,{3+\hh})--++(1.1,0) node[midway,above]{$v_0$};
 \node[above left] at (\xa,{3+\hh}){$A$};
 \draw[<->,TextGray] (-.45,3)--(-.45,{3+\hh}) node[midway,left]{$H$};
 \draw[<->,TextGray] (3.2,3.55)--(5,3.55) node[midway,above]{$d$};
 \draw[<->,TextGray] (7.4,0)--(7.4,1) node[midway,right]{$l$};
\end{tikzpicture}
\figcap{题图：A 在首个踏面上方 $H$；台阶宽 $d$、高 $l$，小球逐次反弹。}\label{ch03:ref:question-figure-2-43}
\end{minipage}\end{center}
}

\newcommand{\KQuestionBT}{%
\begin{center}\begin{minipage}{.97\linewidth}\centering
\begin{tikzpicture}[>=Latex,scale=.9]

\draw[thick] (-1.8,0)--(4.1,0);\draw[thick] (0,1.3) circle(1.3);
\coordinate (A) at (3.2,0);\coordinate (T) at (0,2.6);\coordinate (O) at (0,1.3);
\draw[kguide] (0,0)--(T);
\draw[kcurve,dashed,domain=0:3.2,samples=100] plot(\x,{2.6+.75*\x-(5/10.24)*\x*\x});
\pgfmathsetmacro{\launchangle}{180-atan(2.375)}
\draw[vector] (A)--++(\launchangle:1.3) node[above right]{$v_0$};
\draw[vector] (T)--++(-.8,-.6) node[above left]{$v_T$};
\draw[<->] (0,-.55)--(3.2,-.55) node[midway,below]{$d$};
\draw (O)--(-1.3,1.3) node[midway,above]{$R$};
\draw (A)++(\launchangle:.75) arc(\launchangle:180:.75);\node at (2.2,.38){$\theta$};
\node[point,label=below:$O$] at (O){};\node[point,label=above left:$T$] at (T){};
\node[point,label=below right:$A$] at (A){};

\end{tikzpicture}
\figcap{题图：A 为抛出点，T 为球顶；虚线仅示意一条可行轨迹，并非唯一轨迹。}\label{ch03:ref:question-figure-2-45}
\end{minipage}\end{center}
}


加速度的大小和方向均不变的运动称为匀加速运动。若 $t=0$ 时质点位于 $\vect r_0$、速度为 $\vect v_0$，加速度为常矢量 $\vect a$，则
\begin{equation}
 \vect v(t)=\vect v_0+\vect a t,\qquad
 \vect r(t)=\vect r_0+\vect v_0t+\frac12\vect a t^2.
\end{equation}

这两个矢量方程同时包含各坐标方向的运动规律。任选固定直角坐标系，把 $\vect r_0,\vect v_0,\vect a$ 分解到各轴，便可得到相互独立的分量方程；再把同一时刻各分量合起来，就还原出空间中的位置和速度。这种“先分解、分别研究、再合成”的方法并非抛体运动专用，而是处理一般匀加速运动的基本工具。

在地球表面附近抛出的物体，若空气阻力可忽略，运动过程中加速度恒为重力加速度 $\vect g$，方向竖直向下。这样的运动称为抛体运动。若初速度仅有竖直分量，轨迹是直线；若初速度有水平分量，轨迹是抛物线。抛体运动容易观察，它所用的分解、合成方法也可推广到一般匀加速运动。

理想模型中的“只受重力”包括两个条件：运动范围内可把 $g$ 视为常量，并且空气阻力相对重力足够小。此时水平分速度不受重力改变，竖直分速度每秒减少 $g$。物体是否正在上升、位于最高点或已经下降，只改变 $v_y$ 的符号，不改变加速度始终向下这一事实。

取抛出点为原点，$x$ 轴水平向前，$y$ 轴竖直向上。初速率为 $v_0$，与水平面夹角为 $\theta$，则
\begin{align}
v_x(t)&=v_0\cos\theta,&x(t)&=v_0\cos\theta\,t,\\
v_y(t)&=v_0\sin\theta-gt,&y(t)&=v_0\sin\theta\,t-\frac12gt^2.
\end{align}
这就是把平面曲线运动分解为水平方向的匀速运动和竖直方向的匀变速运动。由 $t=x/(v_0\cos\theta)$ 消去时间，得到轨迹方程
\begin{equation}
 \boxed{y=x\tan\theta-\frac{gx^2}{2v_0^2\cos^2\theta}}.
\end{equation}
它是开口向下的抛物线。$v_0=0$ 时为自由落体；$\theta=\pi/2$ 或 $-\pi/2$ 时为竖直上抛或竖直下抛；$\theta=0$ 时为平抛。

分量运动必须使用同一个时间变量：不能先独立算出两个方向各自经历的时间，再把它们任意拼接。轨迹方程由消去这个共同的 $t$ 得到，它回答“到达水平坐标 $x$ 时高度是多少”；分量方程则能进一步回答到达该点的时刻和速度。轨迹的对称轴、最高点和宽窄都由 $v_0$ 与 $\theta$ 决定，同一抛物线的上升段和下降段在等高处具有大小相同、竖直分量方向相反的速度。

当落点与抛出点等高且 $0<\theta<\pi/2$ 时，除去 $t=0$ 的根，由 $y(t)=0$ 得飞行时间
\begin{equation}
 T=\frac{2v_0\sin\theta}{g}.
\end{equation}
最高点满足 $v_y=0$，所以到达最高点的时间为 $T/2$，最大高度为 $v_0^2\sin^2\theta/(2g)$。水平射程为 $v_0\cos\theta\,T$。这些结果都依赖“始末等高”；从高台落地或落在斜面上时，不能未经修改直接使用。

在任意高度相同的上升点和下降点，机械条件相同而时刻关于最高点对称。水平分速度相同，竖直分速度大小相同、方向相反，因而总速率相同。若题目给出落地速度方向、飞行时间或最高点条件，优先利用这些对称与分量性质往往比从完整二次方程重新求解更直接。

固定初速率、改变抛射角时，同一目标点可能对应两个角度，也可能只有一个临界角或完全不可达。把 $u=\tan\theta$ 代入轨迹方程后，目标条件成为关于 $u$ 的二次方程；判别式大于零对应两条轨迹，等于零对应包络线上的唯一轨迹，小于零则不可达。图~\ref{ch03:fig:k04-ex-2-37}的可达区域正是把这一判别条件对所有目标点统一表达出来。

障碍问题不仅要让轨迹到达目标，还要保证到达前不穿过障碍。对有体积的球面、墙面或台阶，应比较整段轨迹与边界，而不是只让终点坐标相等。
极限轨迹可能在目标之前与障碍相切，切点不一定是目标点。例~\ref{ch03:ex:k04-2-45}中，下降时命中球顶也满足题意；
不能把“球顶”自动当作“抛物线顶点”。还需区分允许擦边的临界轨迹与要求完全不提前接触的实际轨迹。

最终还应检查时间为正、根号内非负、发射角处在题设范围，并把所得轨迹代回边界条件。任何一个代数根若对应反向发射、碰撞发生在抛出前或提前穿过障碍，都不是物理解。

\begin{origexamplebox}{网球发球的最低击球高度}\label{ch03:ex:k04-2-32}
网球场总长为 $2L$，球网高 $h$，发球区远端距球网为 $l$。把发球简化为从高度 $H$ 处以速率 $v_0$ 水平抛出。证明：要使球越过球网并落入规定区域，$H$ 必须大于某个值 $H_0$，并求 $H_0$。
\par 题图见图~\ref{ch03:ref:question-figure-2-32}。
\KQuestionBG
\end{origexamplebox}
\Needspace{6\baselineskip}\begin{xesolution}
以下推导所用的解图为图~\ref{ch03:fig:k04-ex-2-32}。

令 $k=g/(2v_0^2)$，则球的高度为 $y=H-kx^2$。越过 $x=L$ 处球网要求
\[
 k<\frac{H-h}{L^2};
\]
\TennisServeDiagram

在发球区远端 $x=L+l$ 之前落地要求
\[
 \sqrt{\frac Hk}\le L+l,\qquad k\ge\frac{H}{(L+l)^2}.
\]
两条件能由同一个 $k$ 满足，必须有
\[
 \frac{H}{(L+l)^2}<\frac{H-h}{L^2}.
\]
整理得
\[
 \boxed{H>H_0=\frac{h(L+l)^2}{l(2L+l)}}.
\]
等号对应擦网的临界情形，故合法发球需取严格大于。落点在球网之后的条件可与上述不等式同时满足。
\end{xesolution}


\begin{origexamplebox}{落地速度方向}\label{ch03:ex:k04-2-33}
石块从 $20\,\mathrm m$ 高处水平抛出，落地速度与水平面成 $45^\circ$。求抛出速度。
\par 题图见图~\ref{ch03:ref:question-figure-2-33}。
\KQuestionBH
\end{origexamplebox}
\Needspace{6\baselineskip}\begin{xesolution}
以下推导所用的解图为图~\ref{ch03:fig:k04-ex-2-33}。

落地时竖直分速度大小 $v_y=\sqrt{2gh}$，水平分速度始终为 $v_0$。速度与水平面成 $45^\circ$，故 $v_0=v_y$：
\[
 \boxed{v_0=\sqrt{2gh}}.
\]
\HorizontalThrowImpactDiagram

取 $g=10\,\mathrm{m/s^2}$、$h=20\,\mathrm m$，得 $\boxed{v_0=20\,\mathrm{m/s}}$。
\end{xesolution}

\begin{origexamplebox}{斜抛射程}\label{ch03:ex:k04-2-34}
证明：以初速率 $v_0$、仰角 $\theta$ 抛出的物体，回到抛出点同一高度时的射程为
\[
 S=\frac{v_0^2}{g}\sin2\theta,
\]
并证明同一初速率下 $45^\circ$ 斜抛射程最大。
\par 题图见图~\ref{ch03:ref:question-figure-2-34}。
\KQuestionBI
\end{origexamplebox}
\Needspace{6\baselineskip}\begin{xesolution}
以下推导所用的解图为图~\ref{ch03:fig:k04-ex-2-34}。

除 $t=0$ 外，由 $y(t)=0$ 得飞行时间 $T=2v_0\sin\theta/g$，于是
\[
 S=v_0\cos\theta\,T
 =\frac{2v_0^2\sin\theta\cos\theta}{g}
 =\boxed{\frac{v_0^2}{g}\sin2\theta}.
\]
\ProjectileRangeDiagram

因 $\sin2\theta\le1$，最大值在 $2\theta=90^\circ$ 即 $\boxed{\theta=45^\circ}$ 处取得，$S_{\max}=v_0^2/g$。
\end{xesolution}


\begin{origexamplebox}{互余抛射角}\label{ch03:ex:k04-2-35}
同一初速率 $v_0$ 下，仰角 $\theta_1$ 的射程为 $S$。证明仰角 $\theta_2=\pi/2-\theta_1$ 的射程也为 $S$，并求两次飞行时间乘积 $t_1t_2$。
\par 题图见图~\ref{ch03:ref:question-figure-2-35}。
\KQuestionBJ
\end{origexamplebox}
\Needspace{6\baselineskip}\begin{xesolution}
以下推导所用的解图为图~\ref{ch03:fig:k04-ex-2-35}。

$\sin2\theta_2=\sin(\pi-2\theta_1)=\sin2\theta_1$，故射程相同。又
\[
t_1=\frac{2v_0\sin\theta_1}{g},\qquad
t_2=\frac{2v_0\cos\theta_1}{g},
\]
\ComplementaryProjectileDiagram

所以
\[
 t_1t_2=\frac{2v_0^2\sin2\theta_1}{g^2}
 =\boxed{\frac{2S}{g}}.
\]
\end{xesolution}

\begin{origexamplebox}{越过障碍物的最小速度}\label{ch03:ex:k04-2-36}
从 $A$ 点抛出的物体要越过水平距离为 $L$、高度为 $H$ 的障碍物，求所需的最小初速率。
\par 题图见图~\ref{ch03:ref:question-figure-2-36}。
\KQuestionBK
\end{origexamplebox}
\Needspace{6\baselineskip}\begin{xesolution}
以下推导所用的解图为图~\ref{ch03:fig:k04-ex-2-36}。

令 $u=\tan\theta$。轨迹通过临界点 $(L,H)$ 时
\[
H=Lu-\frac{gL^2}{2v_0^2}(1+u^2),
\quad
v_0^2=\frac{gL^2(1+u^2)}{2(Lu-H)}.
\]
\ObstacleClearanceDiagram

由于速度平方必须为正，先限定 $Lu-H>0$，即 $u>H/L$。
对 $u$ 求导，保留决定符号的部分：
\[
 \frac{\dd v_0^2}{\dd u}
 =\frac{gL^2}{2}\,
 \frac{2u(Lu-H)-L(1+u^2)}{(Lu-H)^2}
 =\frac{gL^2}{2}\frac{Lu^2-2Hu-L}{(Lu-H)^2}.
\]
分子为零有两根。负号根小于 $H/L$，不满足可达条件；取正号根得
\[
u=\frac{H+\sqrt{H^2+L^2}}{L}.
\]
代回可得
\[
 \boxed{v_{\min}=\sqrt{g\bigl(\sqrt{L^2+H^2}+H\bigr)}}.
\]
在 $u\to(H/L)^+$ 与 $u\to+\infty$ 两端，$v_0^2$ 均趋于无穷，且允许范围内只有这一个驻点，因此它确为全局极小值。
相应最优角满足 $\tan\theta=(H+\sqrt{H^2+L^2})/L$。
式中 $g$ 乘长度后开方，量纲为速度；$H=0$ 时恢复 $\theta=45^\circ$ 和 $v_0=\sqrt{gL}$。
图中所画的是经过顶端的边界轨迹。若“越过”要求不碰顶端，则速度必须略大于该界限，严格说该值为下确界。
\end{xesolution}

\begin{origexamplebox}{固定速率的可达区域}\label{ch03:ex:k04-2-37}
质点从给定点 $A$ 以固定速率 $v_0$、不同方向抛出。求所有可能到达区域与不可到达区域的分界线。
\par 题图见图~\ref{ch03:ref:question-figure-2-37}。
\KQuestionBL
\end{origexamplebox}
\Needspace{6\baselineskip}\begin{xesolution}
以下推导所用的解图为图~\ref{ch03:fig:k04-ex-2-37}。

把轨迹写成 $u=\tan\theta$ 的二次式：
\[
y=xu-\frac{gx^2}{2v_0^2}(1+u^2).
\]
\ProjectileReachabilityDiagram

先取 $x>0$。记 $K=gx^2/(2v_0^2)>0$，按 $u$ 配方：
\[
 y=-K\left(u-\frac{v_0^2}{gx}\right)^2
    +\frac{v_0^2}{2g}-\frac{gx^2}{2v_0^2}.
\]
平方项不能为负，所以在给定横坐标处，任何抛射角能到达的高度都不超过最后两项。
当 $u=v_0^2/(gx)$ 时平方项为零，这个最高位置确实能到达。因此包络线为
\[
 \boxed{y=\frac{v_0^2}{2g}-\frac{gx^2}{2v_0^2}}.
\]
在图示 $x>0$ 半平面内，包络线以下有两个关于 $u$ 的实根，边界上两根重合，边界上方无实根。
各实际轨迹只能与包络线相切，不能穿到它的上方。
若题意允许向左右各方向抛射，则 $x<0$ 的情形由左右对称得到，分界线方程不变。
$x=0$ 必须单独处理：竖直上抛的最高高度是 $v_0^2/(2g)$，正好是包络线顶点，不能在此处使用含 $1/x$ 的最优角公式。
这里没有引入地面等额外障碍；若存在地面，还须把可达集合与允许运动空间相交。
\end{xesolution}


\begin{origexamplebox}{铅球的最佳出手角}\label{ch03:ex:k04-2-38}
运动员能使铅球以最大速率 $v_0$ 从离地 $H$ 处出手。求最大水平射程对应的出手角；当 $v_0=12.37\,\mathrm{m/s}$、$H=1.5\,\mathrm m$ 时，求角度与成绩。

\begin{qualityfigure}
\begin{tikzpicture}[>=Latex,font=\small]
\draw[thick] (-.5,-1)--(6,-1) node[right]{地面};
\coordinate (O) at (0,1);\coordinate (P) at (5,-1);
\fill (O) circle(2pt) node[left]{$O$};\fill (P) circle(2pt) node[below]{$P$};
\draw[<->] (-.4,-1)--(-.4,1) node[midway,left]{$H$};
\draw[dashed] (O)--(5,1);\draw[->,DeepBlue,thick] (O)--(1.3,2) node[right]{$v_0$};\draw (1,1) arc(0:37.57:1) node[midway,right]{$\theta$};\draw[DeepBlue,domain=0:5,samples=60] plot(\x,{1+.8*\x-.24*\x*\x});
\end{tikzpicture}
\figcap{从高于地面的出手点射向落点}\label{mech:03:38:q}
\end{qualityfigure}
\end{origexamplebox}
\Needspace{6\baselineskip}\begin{xesolution}
\begin{qualityfigure}
\begin{tikzpicture}[>=Latex,font=\small]
\draw[thick] (-.5,-1)--(6,-1) node[right]{地面};
\coordinate (O) at (0,1);\coordinate (P) at (5,-1);
\fill (O) circle(2pt) node[left]{$O$};\fill (P) circle(2pt) node[below]{$P$};
\draw[<->] (-.4,-1)--(-.4,1) node[midway,left]{$H$};
\draw[dashed] (O)--(5,1);\draw[->,DeepBlue,thick] (O)--(1.3,2) node[right]{$v_0$};\draw (1,1) arc(0:37.57:1) node[midway,right]{$\theta$};\draw[DeepBlue,domain=0:5,samples=60] plot(\x,{1+.8*\x-.24*\x*\x});\draw[<->,Gold] (0,-1.6)--(5,-1.6) node[midway,below]{$R$};\draw[->] (0,1)--(0,2.7) node[above]{$y$};\node[above right] at (5,-1){$(R,-H)$};
\end{tikzpicture}
\figcap{从高于地面的出手点射向落点（解答用辅助图）}\label{mech:03:38:a}
\end{qualityfigure}
图~\ref{mech:03:38:a}中，以出手点 $O$ 为坐标原点，地面是 $y=-H$，落点 $P$ 的坐标为 $(R,-H)$。把这一边界代入轨迹方程，才可对出手角优化。

落地时间为
\[
T=\frac{v_0\sin\theta+\sqrt{v_0^2\sin^2\theta+2gH}}{g},
\]
这里取了落地时间的正根，负根没有飞行意义。射程为 $R=v_0\cos\theta\,T$。
直接对含根式的表达式求导较繁琐，可以把“给定 $R$ 是否能到达地面”写成同一个问题。
以出手点为原点，落点为 $(R,-H)$，令 $u=\tan\theta$，则
\[
 -H=Ru-\frac{gR^2}{2v_0^2}(1+u^2)
 =-\frac{gR^2}{2v_0^2}
   \left(u-\frac{v_0^2}{gR}\right)^2
   +\frac{v_0^2}{2g}-\frac{gR^2}{2v_0^2}.
\]
存在实数抛射角的条件为
\[
 -H\le\frac{v_0^2}{2g}-\frac{gR^2}{2v_0^2}
 \quad\Longleftrightarrow\quad
 g^2R^2\le v_0^4+2gHv_0^2.
\]
最大射程对应等号；此时平方项必须为零，于是
$u_{\rm opt}=v_0^2/(gR_{\max})$。分别求出 $R_{\max}$ 和 $u_{\rm opt}$，得到
\[
 \boxed{\tan\theta_{\rm opt}=\frac{v_0}{\sqrt{v_0^2+2gH}}},
 \qquad
 \boxed{R_{\max}=\frac{v_0}{g}\sqrt{v_0^2+2gH}}.
\]
取 $g=9.8\,\mathrm{m/s^2}$，得 $\theta_{\rm opt}\approx42.5^\circ$，$R_{\max}\approx17.1\,\mathrm m$。出手点高于地面，所以最佳角小于 $45^\circ$。
\end{xesolution}

\begin{origexamplebox}{山坡靶的炮口初速}\label{ch03:ex:k04-2-39}
炮位看靶子的仰角为 $\alpha$，炮与靶的水平距离为 $L$，炮身仰角为 $\beta$。求击中目标所需的初速率。
\par 题图见图~\ref{ch03:ref:question-figure-2-39}。
\KQuestionBN
\end{origexamplebox}
\Needspace{6\baselineskip}\begin{xesolution}
以下推导所用的解图为图~\ref{ch03:fig:k04-ex-2-39}。

目标坐标为 $(L,L\tan\alpha)$。代入轨迹方程：
\[
L\tan\alpha=L\tan\beta-
\frac{gL^2}{2v_0^2\cos^2\beta}.
\]
\HillsideTargetDiagram

所以
\[
 \boxed{v_0=\sqrt{\frac{gL}{2\cos^2\beta(\tan\beta-\tan\alpha)}}},
\]
并要求 $\beta>\alpha$，否则根号内不为正，理想抛体轨迹不能击中该点。
\end{xesolution}

\ProjectileVectorsDiagram

除逐分量列方程外，抛体运动还可作矢量合成：
\begin{equation}
 \vect v(t)=\vect v_0+\vect g t,
 \qquad
 \vect r(t)=\vect v_0t+\frac12\vect g t^2.
\end{equation}
图~\ref{ch03:fig:k04-2-12}中，$\vect v_0$ 与 $\vect v_0t$ 平行，但两个实线三角形一般并不相似：重力项的相对长度差一个因子 $2$。
图示时刻物体正在下降，因此 $\vect v$ 向右下；但它仍高于抛出点，因此 $\vect r$ 向右上。
从右图 $\vect r$ 的终点再接一段 $\frac12\vect g t^2$，得到的新合矢量为
\[
 \vect r+\frac12\vect g t^2
 =\vect v_0t+\vect g t^2=t(\vect v_0+\vect gt)=t\vect v.
\]
这个补上虚线的外侧三角形，其三条边才分别是速度三角形三边的 $t$ 倍。
几何方向相同不代表速度与位移量纲相同；乘上同一时间后才可逐边比较。

几何法的关键是保留矢量方向。$\vect g t$ 必须从 $\vect v_0$ 的终点接出，才能得到时刻 $t$ 的速度；$\frac12\vect g t^2$ 同理从 $\vect v_0t$ 的终点接出，得到相对抛出点的位移。若只把三个矢量的大小写成普通加法，便会丢失上升与下降、锐角与钝角等方向信息。遇到障碍、斜面或固定瞄准线时，矢量图常能直接揭示相似三角形和投影关系，减少繁琐的坐标消元。

\begin{origexamplebox}{在障碍物顶端水平飞过}\label{ch03:ex:k04-2-40}
水平地面上距 $A$ 点 $L$ 处有高为 $H$ 的障碍物。从 $A$ 抛出的物体恰在通过障碍物上端时速度水平，求抛射仰角 $\theta$ 和初速率 $v_0$。
\par 题图见图~\ref{ch03:ref:question-figure-2-40}。
\KQuestionBO
\end{origexamplebox}
\Needspace{6\baselineskip}\begin{xesolution}
以下推导所用的解图为图~\ref{ch03:fig:k04-ex-2-40}。

障碍物顶端即轨迹最高点。最高点条件给出
\[
H=\frac{v_0^2\sin^2\theta}{2g},\qquad
L=\frac{v_0^2\sin\theta\cos\theta}{g}.
\]
\HorizontalAtObstacleDiagram

相除得
\[
 \boxed{\tan\theta=\frac{2H}{L}},
\]
再代回任一式，得
\[
 \boxed{v_0=\sqrt{\frac{g(L^2+4H^2)}{2H}}}.
\]
\end{xesolution}

\begin{origexamplebox}{垂直撞上斜面}\label{ch03:ex:k04-2-41}
从倾角为 $\alpha$ 的斜面底部斜向上抛出物体。仰角 $\theta$ 为何值时，物体击中斜面时的速度恰与斜面垂直？
\par 题图见图~\ref{ch03:ref:question-figure-2-41}。
\KQuestionBP
\end{origexamplebox}
\Needspace{6\baselineskip}\begin{xesolution}
以下推导所用的解图为图~\ref{ch03:fig:k04-ex-2-41}。

以抛出点为原点，水平向右为 $x$ 轴、竖直向上为 $y$ 轴。
图示构型取 $0<\alpha<\pi/2$、$\alpha<\theta<\pi/2$，即物体先进入斜面上方的空间。
抛体的位置为
\[
 x=v_0\cos\theta\,t,\qquad y=v_0\sin\theta\,t-\frac12gt^2.
\]
\InclinedPlaneNormalImpactDiagram

再次击中斜面时满足 $y=x\tan\alpha$，代入后整理为
\[
 t\left[v_0(\sin\theta-\cos\theta\tan\alpha)-\frac12gt\right]=0.
\]
其中 $t=0$ 只是抛出瞬间，不是所求的再次碰撞。利用
$\sin\theta-\cos\theta\tan\alpha=\sin(\theta-\alpha)/\cos\alpha$，
得到正的飞行时间
\[
t=\frac{2v_0\sin(\theta-\alpha)}{g\cos\alpha}.
\]
斜面切向单位矢量为 $(\cos\alpha,\sin\alpha)$。速度垂直斜面要求切向分量为零：
\[
v_0\cos(\theta-\alpha)-gt\sin\alpha=0.
\]
代入 $t$ 后
\[
\cot(\theta-\alpha)=2\tan\alpha,
\]
因此
\[
 \boxed{\theta=\alpha+\arctan\!\left(\frac{1}{2\tan\alpha}\right)}.
\]
这里反正切取 $(0,\pi/2)$ 内的值。由于
$1/(2\tan\alpha)<\cot\alpha$，所得 $\theta-\alpha<\pi/2-\alpha$，
确实满足 $\alpha<\theta<\pi/2$。此时法向速度指向斜面，表示撞入而不是离开斜面；
飞行时间中的 $v_0/g$ 具有时间量纲，角度关系与初速率无关。
\end{xesolution}

因为抛体运动非常普遍，上述标准公式并不能自动替代建模。处理新题时应先确定坐标系、初始位置、落点或碰撞点的几何条件，再把同一时刻的水平和竖直方程联立。若题目含台阶、反弹或球面障碍，还要逐段检查碰撞后的速度规则以及整段轨迹是否提前与障碍相交，不能只验证终点坐标。


\begin{origexamplebox}{轰炸机瞄准镜}\label{ch03:ex:k04-2-42}
轰炸机在离地高度 $H$ 处以速率 $v$ 水平匀速飞行。设计瞄准镜，使飞行员在目标位于十字线正中时投弹即可命中。求瞄准镜轴线与竖直线的夹角。

\begin{qualityfigure}
\begin{tikzpicture}[>=Latex,font=\small]
\draw[thick] (-.5,-1)--(6,-1) node[right]{地面};
\coordinate (O) at (0,1);\coordinate (P) at (5,-1);
\fill (O) circle(2pt) node[left]{$O$};\fill (P) circle(2pt) node[below]{$P$};
\draw[<->] (-.4,-1)--(-.4,1) node[midway,left]{$H$};
\draw[dashed] (O)--(5,1);\draw[->,DeepBlue,thick] (O)--(1.7,1) node[above]{$v$};\draw[dashed] (O)--(P);\draw[dashed] (O)--(0,-1);\draw (0,.3) arc(-90:-21.8:.7) node[midway,below right]{$\phi$};
\end{tikzpicture}
\figcap{投弹点、竖直线与目标视线}\label{mech:03:42:q}
\end{qualityfigure}
\end{origexamplebox}
\Needspace{6\baselineskip}\begin{xesolution}
\begin{qualityfigure}
\begin{tikzpicture}[>=Latex,font=\small]
\draw[thick] (-.5,-1)--(6,-1) node[right]{地面};
\coordinate (O) at (0,1);\coordinate (P) at (5,-1);
\fill (O) circle(2pt) node[left]{$O$};\fill (P) circle(2pt) node[below]{$P$};
\draw[<->] (-.4,-1)--(-.4,1) node[midway,left]{$H$};
\draw[dashed] (O)--(5,1);\draw[->,DeepBlue,thick] (O)--(1.7,1) node[above]{$v$};\draw[dashed] (O)--(P);\draw[dashed] (O)--(0,-1);\draw (0,.3) arc(-90:-21.8:.7) node[midway,below right]{$\phi$};\draw[Gold,thick,domain=0:5,samples=50] plot(\x,{1-2*\x*\x/25});\draw[<->] (0,-1.6)--(5,-1.6) node[midway,below]{$D=vt$};
\end{tikzpicture}
\figcap{投弹点、竖直线与目标视线（解答用辅助图）}\label{mech:03:42:a}
\end{qualityfigure}
图~\ref{mech:03:42:a}中，$OP$ 是投弹瞬间的瞄准线，$H$ 是竖直落差，$D$ 是水平超前距离；夹角 $\phi$ 从竖直线量起，因此用 $\tan\phi=D/H$。

炸弹下落时间 $t=\sqrt{2H/g}$，在此期间水平前进
\[
D=vt=v\sqrt{\frac{2H}{g}}.
\]
投弹瞬间目标应位于飞机前下方，视线与竖直方向夹角 $\phi$ 满足
\[
 \boxed{\tan\phi=\frac DH=v\sqrt{\frac{2}{gH}}}.
\]
因此瞄准镜应沿此前下方视线设置；忽略空气阻力时炸弹保留飞机的水平速度。
\end{xesolution}

\begin{origexamplebox}{台阶上的全同反弹}\label{ch03:ex:k04-2-43}
每级台阶长 $d$、高 $l$。球从 $A$ 点以水平速度 $v_0$ 抛出；碰撞后水平分速度不变，竖直分速度反向并变为碰撞前的 $e$ 倍，$e<1$。若每次碰撞前与台阶的相对速度相同，轨迹形成全同的往复过程，求 $v_0$ 和每次反弹后的上升高度。
\par 题图见图~\ref{ch03:ref:question-figure-2-43}。
\KQuestionBR
\end{origexamplebox}
\Needspace{6\baselineskip}\begin{xesolution}
以下推导所用的解图为图~\ref{ch03:fig:k04-ex-2-43}。

设每次碰撞前竖直分速度大小为 $V$，则碰撞后向上为 $eV$。下一碰撞点比前一点低 $l$，由竖直运动的速度—位移关系
\[
V^2=e^2V^2+2gl,
\quad
V=\sqrt{\frac{2gl}{1-e^2}}.
\]
\StairBounceDiagram

相邻碰撞间的时间等于竖直速度从 $+eV$ 变到 $-V$ 所需时间：
\[
T=\frac{(1+e)V}{g}.
\]
水平位移为 $d=v_0T$，故
\[
 \boxed{v_0=d\sqrt{\frac{g(1-e)}{2l(1+e)}}}.
\]
每次反弹后的上升高度为
\[
 \boxed{h_b=\frac{e^2V^2}{2g}=\frac{e^2l}{1-e^2}}.
\]
若题图另以 $H$ 表示首次下落高度，全同条件还要求 $H=l/(1-e^2)$，于是 $h_b=e^2H$。
\end{xesolution}

\begin{origexamplebox}{反弹轨迹的面积和}\label{ch03:ex:k04-2-44}
小球从距水平地面 $h$ 高处，以水平速度 $\sqrt{2gh}$ 抛出。每次落地反弹时水平速度不变，竖直速度反向且按固定比率 $e$ 减小，其中 $0<e<1$。若从第一次反弹开始，各段轨迹与地面所围面积之和为 $\frac{8}{21}h^2$，求 $e$。提示：一次从地面出发又落回地面的斜抛轨迹所围面积，等于其最大高度与射程乘积的 $2/3$。

\begin{qualityfigure}
\begin{tikzpicture}[>=Latex,font=\small]
\begin{scope}[y=2cm]\draw[thick](-.5,0)--(6.4,0);\fill(0,1)circle(2pt);\draw[<->](-.4,0)--(-.4,1)node[midway,left]{$h$};
\draw[->,DeepBlue](0,1)--(1.2,1)node[right]{$\sqrt{2gh}$};
\draw[DeepBlue,domain=0:2,samples=40]plot(\x,{1-\x*\x/4});
\draw[DeepBlue,domain=2:4.4,samples=40]plot(\x,{.25*(\x-2)*(4.4-\x)});
\draw[DeepBlue,domain=4.4:5.84,samples=40]plot(\x,{.25*(\x-4.4)*(5.84-\x)});\end{scope}
\end{tikzpicture}
\figcap{反弹后的抛物线弓形面积}\label{mech:03:44:q}
\end{qualityfigure}
\end{origexamplebox}
\Needspace{6\baselineskip}\begin{xesolution}
\begin{qualityfigure}
\begin{tikzpicture}[>=Latex,font=\small]
\begin{scope}[y=2cm]\fill[PaleBlue,domain=2:4.4,samples=40](2,0)--plot(\x,{.25*(\x-2)*(4.4-\x)})--cycle;\draw[thick](-.5,0)--(6.4,0);\fill(0,1)circle(2pt);\draw[<->](-.4,0)--(-.4,1)node[midway,left]{$h$};
\draw[->,DeepBlue](0,1)--(1.2,1)node[right]{$\sqrt{2gh}$};
\draw[DeepBlue,domain=0:2,samples=40]plot(\x,{1-\x*\x/4});
\draw[DeepBlue,domain=2:4.4,samples=40]plot(\x,{.25*(\x-2)*(4.4-\x)});
\draw[DeepBlue,domain=4.4:5.84,samples=40]plot(\x,{.25*(\x-4.4)*(5.84-\x)});\draw[<->,Gold](2,-.3)--(4.4,-.3)node[midway,below]{$S_n$};\draw[<->,Gold](3.2,0)--(3.2,.36)node[midway,right]{$H_n$};\node at(3.2,.85){$A_n=\frac23H_nS_n$};\end{scope}
\end{tikzpicture}
\figcap{反弹后的抛物线弓形面积（解答用辅助图）}\label{mech:03:44:a}
\end{qualityfigure}
图~\ref{mech:03:44:a}中，阴影从某次反弹点起、到下一次落地点止，对应高度 $H_n$ 与射程 $S_n$；题目从第一次反弹开始计面积，故级数下标从 $n=1$ 起。

第一次落地前竖直速度大小为 $\sqrt{2gh}$。第 $n$ 次反弹后的竖直初速度为 $e^n\sqrt{2gh}$，故最大高度 $H_n=e^{2n}h$；水平速度为 $\sqrt{2gh}$，飞行时间 $2e^n\sqrt{2h/g}$，射程 $S_n=4e^nh$。于是第 $n$ 个面积
\[
A_n=\frac23H_nS_n=\frac83h^2e^{3n}.
\]
从 $n=1$ 起求和：
\[
\frac83h^2\frac{e^3}{1-e^3}=\frac8{21}h^2.
\]
解得 $7e^3=1-e^3$，所以
\[
 \boxed{e=\frac12}.
\]
\end{xesolution}

\begin{origexamplebox}{击中球顶的速度界限}\label{ch03:ex:k04-2-45}
地面上放有半径为 $R$ 的球体。从地面斜向上抛出一个质点，要求击中球顶，且在此之前不碰球体。
求所需初速率的下界及相应临界抛射角、抛出点到球底的距离，并判断这个下界能否取到。
同时给出允许“提前与球面相切但不穿入”时的临界结果。
\par 题图见图~\ref{ch03:ref:question-figure-2-45}。
\KQuestionBT
\end{origexamplebox}
\Needspace{6\baselineskip}\begin{xesolution}
以下推导所用的解图为图~\ref{ch03:fig:k04-ex-2-45}。

\textbf{先判断命中时的速度方向。}
物体不可能上升着从球体内部穿到球顶；但既可以在球顶水平掠过，也可以在下降时命中。
因此不能预先令球顶的竖直速度为零。下面把这两类运动统一处理。
\SphereClearanceDiagram
图中 $T$ 是球顶；左图临界轨迹在到达 $T$ 之前还与球面在 $P$ 点相切，右图则从球面上方经过后在下降阶段到达 $T$。
两个点的作用不同：$T$ 决定命中，$P$ 决定是否提前碰球。

\textbf{把全部候选轨迹写成两个参数。}
以球底为水平坐标原点，以地面为 $y=0$，抛出点在 $x=-d$。
设水平分速度为 $u>0$，到达 $T=(0,2R)$ 时的竖直分速度为 $w\le0$。
记无量纲量
\[
 s=-\frac{x}{R},\qquad q=-\frac wu\ge0,\qquad
 k=\frac{gR}{2u^2}>0.
\]
以球顶为已知点写轨迹，得到
\[
 y=2R+\frac wu x-\frac{g}{2u^2}x^2,
 \qquad \boxed{\frac{y}{R}=2+qs-ks^2}.
\]
其中 $s>0$ 表示位于球顶左侧。$q=0$ 才是球顶速度水平的特殊情况。
从抛出点到球顶上升的净高度为 $2R$。由竖直运动公式
$v_{y0}^2-w^2=4gR$，可得
\[
 \frac{v_0^2}{gR}
 =4+\frac{u^2+w^2}{gR}
 =\boxed{4+\frac{1+q^2}{2k}}.
\]
这一式同时检查了机械能：初动能减去到球顶的动能恰为 $2mgR$。

\textbf{检查到顶之前的整段球面。}
在左半球的水平投影 $0<s\le1$ 内，球体上边界为
\[
 \frac{y_s}{R}=1+\sqrt{1-s^2}.
\]
所以不穿入球体要求
\[
 2+qs-ks^2\ge1+\sqrt{1-s^2},
 \quad\text{即}\quad
 q\ge h_k(s):=ks-\frac{1-\sqrt{1-s^2}}{s}.
\]
严格“不提前接触”时，不等式在 $0<s\le1$ 必须严格成立。
在球体投影以外没有球面障碍；轨迹与地面的左交点就是抛出点。
给定 $k$ 后，速度平方随 $q\ge0$ 增加，因此应先找使整段不等式成立的最小 $q$。

若 $k\le1/2$，由于
\[
 1-\sqrt{1-s^2}=\frac{s^2}{1+\sqrt{1-s^2}}\ge\frac{s^2}{2}\ge ks^2,
\]
取 $q=0$ 已足够使轨迹位于球面上方，且
\[
 \frac{v_0^2}{gR}=4+\frac1{2k}\ge5.
\]
这给出“在球顶速度水平”一类轨迹的下界 $\sqrt{5gR}$，但它不一定是所有轨迹的下界。

若 $k>1/2$，函数 $h_k$ 在球顶附近先增加，随后下降，最大值出现在一个内部点。
为求这个点，令 $c=\sqrt{1-s^2}$。对 $s$ 求导，得到
\[
 h_k'(s)=k-\frac1{c(1+c)}.
\]
当 $s$ 从 $0$ 增加到 $1$，$c$ 从 $1$ 减少到 $0$，故第二项从 $1/2$ 单调增至无穷。
因此恰有一个最大值点，满足
\[
 k=\frac1{c(1+c)},\qquad s=\sqrt{1-c^2},\qquad 0<c<1.
\]
在该点取 $q=h_k(s)$，化简为
\[
 q=\frac{s(1-c)}{c(1+c)}.
\]
这两个等式正是提前切点 $P$ 的位置相同、斜率相同条件；并非只在球顶作局部比较。

\textbf{求所有临界轨迹中的最小速度。}
将上述 $k,q$ 代入速度式，注意 $s^2=1-c^2$：
\begin{align*}
 \frac{1+q^2}{2k}
 &=\frac{c^2(1+c)+(1-c)^3}{2c}\\
 &=2c-\frac32+\frac1{2c}.
\end{align*}
于是
\[
 \frac{v_0^2}{gR}
 =\frac52+2c+\frac1{2c}\ge\frac52+2=\frac92,
\]
其中用了 $2c>0$ 与 $(2c)^{-1}>0$ 的和不小于 $2$。
等号在 $c=1/2$ 取得，故
\[
 k_*=\frac43,\quad q_*=\frac1{\sqrt3},\quad
 u_*^2=\frac38gR,\quad w_*^2=\frac18gR,\quad w_*<0.
\]
初始竖直分速度满足 $v_{y0,*}^2=w_*^2+4gR=33gR/8$。
由 $y(-d)=0$，取正的距离根：
\[
 \frac{d}{R}=\frac{q+\sqrt{q^2+8k}}{2k}.
\]
所以临界结果为
\[
 \boxed{v_*=\frac3{\sqrt2}\sqrt{gR}},\qquad
 \boxed{\tan\theta_*=\sqrt{11}},\qquad
 \boxed{d_*=\frac{\sqrt3+3\sqrt{11}}8R}.
\]
切点位于 $x_P=-\sqrt3R/2$、$y_P=3R/2$，早于球顶 $T$。

\textbf{区分最小值和下确界。}
上述等号轨迹先碰到 $P$，因此不满足严格的“不提前接触”。
固定 $k=4/3$，取 $q=1/\sqrt3+\varepsilon$，$\varepsilon>0$，轨迹在每个 $s>0$ 处都比临界轨迹高出 $R\varepsilon s$；
它避开提前接触，仍在 $T$ 命中。令 $\varepsilon\to0^+$，其初速率可无限趋近 $v_*$。
因此：
\[
 \boxed{\text{严格不提前接触：}\quad v_0>v_*,
 \quad \inf v_0=v_*\ \text{但不存在能取到的最小值。}}
\]
这里的不等式表示任何可行轨迹的必要速度界限，且可行速度可以从上方任意逼近此界；
不能据此认为任意抛射角都能命中。若允许提前相切但不穿入，则 $v_*$ 可以取到，临界角和距离即上式。
$\sqrt{gR}$ 具有速度量纲，$d_*\propto R$、$\theta_*$ 与尺度无关；半径缩小时速度界限按 $\sqrt R$ 减小，符合几何相似缩放。
\end{xesolution}
}

{
\graphicspath{{../../../sample/xes-original/kinematics-supplement/}{../../../sample/xes-original/kinematics-supplement/assets/}}
\section{圆周运动}
\newcommand{\CirclePositionDiagram}{%
\begin{center}\begin{tikzpicture}[scale=.9,>=Latex]
 \draw[very thick] (0,0) circle (2);
 \coordinate (O) at (0,0);\coordinate (N) at (2,0);\coordinate (P) at (52:2);
 \draw[very thick] (O)--(N) node[midway,clearlabel,below]{$R$};
 \draw[very thick] (O)--(P);
 \draw[->,TextGray] (.75,0) arc[start angle=0,end angle=52,radius=.75];
 \node[clearlabel] at (1,.43){$\theta$};
 \node[point,label={[clearlabel]below left:$O$}] at (O) {};
 \node[point,label={[clearlabel]right:$N$}] at (N) {};
 \node[point,label={[clearlabel]above right:$P$}] at (P) {};
\end{tikzpicture}\figcap{角位置 $\theta$ 由固定射线 $ON$ 转到半径 $OP$ 的有向角确定；逆时针为正。}\label{ch03:fig:k05-2-13}\end{center}}

\newcommand{\AngularSignDiagram}{%
\begin{center}\begin{tikzpicture}[scale=.75,>=Latex]
 \begin{scope}[xshift=-2.4cm]
  \draw[very thick] (0,0) circle (1.25);\draw[DeepBlue,very thick,->] (1.45,-.2) arc[start angle=-8,end angle=95,radius=1.45];
  \node[clearlabel,below] at (0,-1.65){$\omega>0$（逆时针）};
 \end{scope}
 \begin{scope}[xshift=2.4cm]
  \draw[very thick] (0,0) circle (1.25);\draw[Blue,very thick,->] (-1.45,-.2) arc[start angle=188,end angle=85,radius=1.45];
  \node[clearlabel,below] at (0,-1.65){$\omega<0$（顺时针）};
 \end{scope}
\end{tikzpicture}\figcap{角速度的符号只表示转动方向，不表示快慢。}\label{ch03:fig:k05-2-14}\end{center}}

\newcommand{\CentripetalDerivationDiagram}{%
\begin{center}\begin{minipage}{.98\linewidth}\centering
\begin{tikzpicture}[scale=1,>=Latex]
 \begin{scope}[xshift=-3.5cm]
  \draw[thick] (0,0) circle (1.7);
  \coordinate (A) at (110:1.7);\coordinate (B) at (70:1.7);
  \draw[TextGray] (0,0)--(A) (0,0)--(B);
  \draw[vector] (A)--++(20:1.2) node[pos=.85,above=3pt]{$\vect v_A$};
  \draw[vector] (B)--++(-20:1.2) node[right=3pt]{$\vect v_B$};
  \draw[TextGray] (70:.65) arc[start angle=70,end angle=110,radius=.65];
  \node at (0,1.22){$|\Delta\theta|$};
  \node[point,label=above left:$A$] at (A){};
  \node[point] at (B){};\node[above right] at ($(B)+(.12,.17)$){$B$};
  \node[below] at (0,-2){(a) 同一圆周上的两速度};
 \end{scope}
 \begin{scope}[xshift=2.4cm]
  \coordinate (S) at (0,0);\coordinate (P) at (20:2.5);\coordinate (Q) at (-20:2.5);
  \draw[vector] (S)--(P) node[midway,above left]{$\vect v_A$};
  \draw[vector] (S)--(Q) node[midway,below left]{$\vect v_B$};
  \draw[vector] (P)--(Q) node[midway,right=4pt]{$\Delta\vect v$};
  \draw[TextGray] (-20:.8) arc[start angle=-20,end angle=20,radius=.8];
  \node[right] at (.85,0){$|\Delta\theta|$};
  \node[left] at (S){$S$};\node[above] at (P){$P$};\node[below] at (Q){$Q$};
  \node[below] at (1.1,-2){(b) 平移后作矢量差};
 \end{scope}
\end{tikzpicture}
\figcap{右图保持左图两速度的方向；$\Delta\vect v$ 由旧速度终点指向新速度终点。}\label{ch03:fig:k05-2-15}
\end{minipage}\end{center}}

\newcommand{\RadialCollisionBase}[1]{%
\begin{tikzpicture}[scale=.75,>=Latex]
 \draw[very thick] (0,0) circle (2);\coordinate (P) at (0,2);\coordinate (A) at (25:2);\coordinate (B) at (0,.72);
 \node[point,label={[clearlabel]below left:$O$}] at (0,0) {};
 \node[point,label={[clearlabel]above:$P$}] at (P) {};
 \node[point,label={[clearlabel]below left:$A$}] at (A) {};
 \node[point,label={[clearlabel]right:$B$}] at (B) {};
 \draw[kguide] (0,0)--(P);
 \node[clearlabel,left,xshift=-2pt] at (0,.35){$R$};
 \draw[vector] (B)--(0,1.60) node[midway,clearlabel,right,xshift=3pt]{$#1$};
 \draw[DeepBlue,very thick,->] (4:2.25) arc[start angle=4,end angle=55,radius=2.25];
 \node[right] at (2.42,1.28){$\omega$};
\end{tikzpicture}}

\newcommand{\RadialUniformCollisionDiagram}{%
\begin{center}\RadialCollisionBase{\vect v_B}
\figcap{$B$ 以恒定速度沿射线 $OP$ 运动，故与圆周质点 $A$ 的相撞点只能是 $P$。}\label{ch03:fig:k05-ex-2-47}\end{center}}

\newcommand{\RadialAcceleratedCollisionDiagram}{%
\begin{center}\RadialCollisionBase{\vect a_B}
\figcap{$B$ 从 $O$ 由静止沿 $OP$ 匀加速；几何相撞点仍只能是 $P$，但到达时间条件改变。}\label{ch03:fig:k05-ex-2-48}\end{center}}

\newcommand{\SameCirclePursuitDiagram}{%
\begin{center}\begin{tikzpicture}[scale=.8,>=Latex]
  \draw[very thick] (0,0) circle (2.25);
  \coordinate (O) at (0,0);\coordinate (P) at (90:2.25);\coordinate (Q) at (135:2.25);
  \draw[very thick] (O)--(P);\draw[kguide] (O)--(Q);
  \draw[DeepBlue,very thick,->] (125:2.5) arc[start angle=125,end angle=175,radius=2.5]
    node[midway,clearlabel,left]{$\omega_1$};
  \draw[Blue,very thick,->] (-2.5,-.55) arc[start angle=192,end angle=230,radius=2.5]
    node[midway,clearlabel,left]{$\omega_2$};
  \draw pic[draw=TextGray,->,angle radius=8mm,"$\theta$"{clearlabel}] {angle=P--O--Q};
  \node[point] at (Q) {};\node[left] at ($(Q)+(-.7,.32)$){$Q$};
  \node[point,label={[clearlabel]above:$P$}] at (P) {};
  \node[point,label={[clearlabel]below right:$O$}] at (O) {};
\end{tikzpicture}\figcap{两质点从同一点出发沿同一圆周逆时针运动，第二次相遇位置相对 $OP$ 转过角 $\theta$。}\label{ch03:fig:k05-ex-2-49}\end{center}}

\newcommand{\ConcentricCollinearityDiagram}{%
\begin{center}\begin{tikzpicture}[scale=.8,>=Latex]
  \draw[very thick] (0,0) circle (1.35);\draw[very thick] (0,0) circle (2.35);
  \coordinate (O) at (0,0);\coordinate (A) at (28:1.35);\coordinate (B) at (28:2.35);
  \draw[kguide] (O)--(B);
  \draw[vector] (A)--++(118:1.05) node[midway,clearlabel,right]{$\vect v_A$};
  \draw[vector] (B)--++(118:.82);
  \node at (2.12,2.17){$\vect v_B$};
  \draw[TextGray] (O)--(145:1.35) node[midway,clearlabel,above=3pt]{$R_1$};
  \draw[TextGray] (O)--(225:2.35);
  \node at (-1.75,-.65){$R_2$};
  \node[point] at (O) {};\node at (.25,-.30){$O$};
  \node[point,label={[clearlabel]right:$A$}] at (A) {};
  \node[point,label={[clearlabel]below right:$B$}] at (B) {};
\end{tikzpicture}\figcap{$A,B$ 分别在半径 $R_1,R_2$ 的同心圆上；题目要求两点与圆心再次共线。}\label{ch03:fig:k05-ex-2-50}\end{center}}

\newcommand{\ProjectileCircleCollisionDiagram}{%
\begin{center}\begin{minipage}{.96\linewidth}\centering
\begin{tikzpicture}[scale=1,>=Latex]
 \draw[TextGray,dashed] (3,0) arc[start angle=0,end angle=180,x radius=3,y radius=.75];
 \draw[thick] (-3,0) arc[start angle=180,end angle=360,x radius=3,y radius=.75];
 \coordinate (O) at (0,0);\coordinate (A) at (3,0);\coordinate (S) at (0,3);
 \draw[kguide] (O)--(A) (O)--(S);
 \draw[vector] (S)--(1.25,3) node[midway,above]{$v_0$};
 \draw[kcurve,dashed,->,domain=0:3,samples=90] plot(\x,{3-\x*\x/3});
 \draw[Blue,->,thick] (3.3,0) arc[start angle=0,end angle=52,x radius=3.3,y radius=.83];
 \node[above right] at (2.6,.7){$\omega$};
 \draw[<->,TextGray] (-.45,0)--(-.45,3) node[midway,left]{$H$};
 \node[below] at (1.5,0){$R$};
 \node[point,label=below left:$O$] at (O){};
 \node[point,label=right:$A$] at (A){};
 \node[point,label=above left:{$S\;(B\text{的初始位置})$}] at (S){};
\end{tikzpicture}
\figcap{$S$ 在圆心 $O$ 正上方；抛体从 $S$ 水平出发，在圆周上与整周返回的质点相遇。}\label{ch03:fig:k05-ex-2-51}
\end{minipage}\end{center}}

\newcommand{\AccelerationVelocityAngleDiagram}{%
\begin{center}\begin{minipage}{.96\linewidth}\centering
\begin{tikzpicture}[scale=.95,>=Latex]
 \draw[thick] (0,0) circle (2.0);\coordinate (P) at (35:2);
 \coordinate (V) at ($(P)+(125:1.4)$);\coordinate (A) at ($(P)+(170:1.4)$);
 \draw[kguide] (0,0)--(P) node[pos=.4,below right]{$R$};
 \draw[vector] (P)--(V) node[above right]{$\vect v$};
 \draw[vector] (P)--(A) node[above left]{$\vect a$};
 \draw pic[draw=TextGray,angle radius=7mm] {angle=V--P--A};
 % The circumference crosses this narrow wedge. Show its label separately,
 % retaining the physical vectors and the main angle arc without a mask.
 \begin{scope}[shift={(4.7,1.3)}]
 \draw[vector] (0,0)--(125:1.6) node[above right]{$\vect v$};
 \draw[vector] (0,0)--(170:1.6) node[above left]{$\vect a$};
 \draw[TextGray] (125:.6) arc(125:170:.6);
 \node at (147.5:.95){$\theta$};
 \node[below] at (-.5,-.4){夹角放大};
 \end{scope}
 \node[point,label=below left:$O$] at (0,0){};\node[point] at (P){};
\end{tikzpicture}
\figcap{右侧单独放大速度与加速度的夹角 $\theta$；该角取值在 $0$ 与 $\pi$ 之间。}\label{ch03:fig:k05-ex-2-59}
\end{minipage}\end{center}}

\newcommand{\ConcentricSightDiagram}{%
\begin{center}\begin{minipage}{.96\linewidth}\centering
\begin{tikzpicture}[scale=.9,>=Latex]
 \pgfmathsetmacro{\ax}{1.5^2/2.8}\pgfmathsetmacro{\ay}{sqrt(1.5^2-\ax^2)}
 \coordinate (O) at (0,0);\coordinate (B) at (2.8,0);\coordinate (A) at (\ax,\ay);
 \draw[thick] (O) circle (1.5) (O) circle (2.8);
 \draw[thick] (O)--(B) (B)--(A);\draw[kguide] (O)--(A);
 \draw pic[draw=TextGray,angle radius=4mm] {right angle=O--A--B};
 \draw pic[draw=TextGray,angle radius=8mm] {angle=A--B--O};
 \node at (1.75,.27){$\theta$};
 \draw[TextGray] (O)--(145:1.5) node[midway,below left]{$R_1$};
 \draw[TextGray] (O)--(230:2.8) node[pos=.78,below right=4pt]{$R_2$};
 \node[point,label=below:$O$] at (O){};\node[point,label=right:$B$] at (B){};
 \node[point,label=above:$A$] at (A){};
\end{tikzpicture}
\figcap{$A$ 为从 $B$ 引向内圆的切点，因而 $OA\perp AB$。}\label{ch03:fig:k05-ex-2-52}
\end{minipage}\end{center}}

\newcommand{\RacetrackDiagram}{%
\begin{center}\begin{tikzpicture}[scale=.72,>=Latex]
 \draw[very thick] (-2.4,-2.4)--(-2.4,0) arc[start angle=180,end angle=0,radius=2.4]--(2.4,-2.4);
 \draw[very thick] (-1.4,-2.4)--(-1.4,0) arc[start angle=180,end angle=0,radius=1.4]--(1.4,-2.4);
 \draw[->,DeepBlue,very thick] (-1.4,-1.9)--(-1.4,-.85);
 \draw[->,Blue,very thick] (-2.4,-1.9)--(-2.4,-.85);
 \draw[<->,TextGray] (0,0)--(0,1.4) node[midway,clearlabel,right]{$R_1$};
 \draw[<->,TextGray] (0,0)--(135:2.4) node[pos=.85,clearlabel,above right=3pt]{$R_2$};
 \node[clearlabel,below] at (0,-2.75){内、外路线均为半圆弯道};
\end{tikzpicture}\figcap{解图：外道允许更高速度，但半圆弧也更长。}\label{ch03:fig:k05-ex-2-57}\end{center}}

\newcommand{\KQuestionBV}{%
%
\begin{center}\RadialCollisionBase{\vect v_B}
\figcap{题图：A 沿半径为 $R$ 的圆周运动，B 沿 OP 匀速运动。}\label{ch03:ref:question-figure-2-47}\end{center}
}

\newcommand{\KQuestionBW}{%
%
\begin{center}\RadialCollisionBase{\vect a_B}
\figcap{题图：A 沿圆周匀速运动，B 从 O 由静止沿 OP 匀加速运动。}\label{ch03:ref:question-figure-2-48}\end{center}
}

\newcommand{\KQuestionBX}{%
%
\begin{center}\begin{tikzpicture}[scale=.8,>=Latex]
  \draw[very thick] (0,0) circle (2.25);
  \coordinate (O) at (0,0);\coordinate (P) at (90:2.25);\coordinate (Q) at (135:2.25);
  \draw[very thick] (O)--(P);\draw[kguide] (O)--(Q);
  \draw[DeepBlue,very thick,->] (125:2.5) arc[start angle=125,end angle=175,radius=2.5]
    node[midway,clearlabel,left]{$\omega_1$};
  \draw[Blue,very thick,->] (-2.5,-.55) arc[start angle=192,end angle=230,radius=2.5]
    node[midway,clearlabel,left]{$\omega_2$};
  \draw pic[draw=TextGray,->,angle radius=8mm,"$\theta$"{clearlabel}] {angle=P--O--Q};
  \node[point] at (Q) {};\node[left] at ($(Q)+(-.7,.32)$){$Q$};
  \node[point,label={[clearlabel]above:$P$}] at (P) {};
  \node[point,label={[clearlabel]below right:$O$}] at (O) {};
\end{tikzpicture}\figcap{题图：两质点从正上方 P 同时出发，逆时针运动；$\theta$ 标记再次相遇的位置。}\label{ch03:ref:question-figure-2-49}\end{center}
}

\newcommand{\KQuestionBY}{%
%
\begin{center}\begin{tikzpicture}[scale=.8,>=Latex]
  \draw[very thick] (0,0) circle (1.35);\draw[very thick] (0,0) circle (2.35);
  \coordinate (O) at (0,0);\coordinate (A) at (28:1.35);\coordinate (B) at (28:2.35);
  \draw[kguide] (O)--(B);
  \draw[vector] (A)--++(118:1.05) node[midway,clearlabel,right]{$\vect v_A$};
  \draw[vector] (B)--++(118:.82);
  \node at (2.12,2.17){$\vect v_B$};
  \draw[TextGray] (O)--(145:1.35) node[midway,clearlabel,above=3pt]{$R_1$};
  \draw[TextGray] (O)--(225:2.35);
  \node at (-1.75,-.65){$R_2$};
  \node[point] at (O) {};\node at (.25,-.30){$O$};
  \node[point,label={[clearlabel]right:$A$}] at (A) {};
  \node[point,label={[clearlabel]below right:$B$}] at (B) {};
\end{tikzpicture}\figcap{题图：A、B 分别在半径为 $R_1$、$R_2$ 的同心圆上同向运动。}\label{ch03:ref:question-figure-2-50}\end{center}
}

\newcommand{\KQuestionBZ}{%
%
\begin{center}\begin{minipage}{.96\linewidth}\centering
\begin{tikzpicture}[scale=1,>=Latex]
 \draw[TextGray,dashed] (3,0) arc[start angle=0,end angle=180,x radius=3,y radius=.75];
 \draw[thick] (-3,0) arc[start angle=180,end angle=360,x radius=3,y radius=.75];
 \coordinate (O) at (0,0);\coordinate (A) at (3,0);\coordinate (S) at (0,3);
 \draw[kguide] (O)--(A) (O)--(S);
 \draw[vector] (S)--(1.25,3) node[midway,above]{$v_0$};
 \draw[kcurve,dashed,->,domain=0:3,samples=90] plot(\x,{3-\x*\x/3});
 \draw[Blue,->,thick] (3.3,0) arc[start angle=0,end angle=52,x radius=3.3,y radius=.83];
 \node[above right] at (2.6,.7){$\omega$};
 \draw[<->,TextGray] (-.45,0)--(-.45,3) node[midway,left]{$H$};
 \node[below] at (1.5,0){$R$};
 \node[point,label=below left:$O$] at (O){};
 \node[point,label=right:$A$] at (A){};
 \node[point,label=above left:{$S\;(B\text{的初始位置})$}] at (S){};
\end{tikzpicture}
\figcap{题图：水平圆心 O 上方为 S；B 从 S 水平抛出，A 沿圆周运动。}\label{ch03:ref:question-figure-2-51}
\end{minipage}\end{center}
}

\newcommand{\KQuestionCA}{%
\begin{center}\begin{minipage}{.97\linewidth}\centering
\begin{tikzpicture}[>=Latex,scale=.9]

\draw[thick] (0,0) circle(1.5) (0,0) circle(2.6);
\coordinate (O) at (0,0);\coordinate (A) at (75:1.5);\coordinate (B) at (27:2.6);
\draw[kguide] (O)--(27:2.6) (O)--(A);\draw[thick] (A)--(B);
\pgfmathsetmacro{\lineangle}{atan2(1.5*sin(75)-2.6*sin(27),1.5*cos(75)-2.6*cos(27))}
\draw (B)++(\lineangle:.6) arc(\lineangle:207:.6);
\node at ($(B)+(190:.91)$){$\theta$};
\draw[->,DeepBlue] (80:1.78) arc(80:122:1.78) node[above left]{$\omega_1$};
\draw[->,DeepBlue] (5:2.87) arc(5:39:2.87) node[above right]{$\omega_2$};
\draw (O)--(150:1.5);\node at (-1.72,1.05){$R_1$};
\draw (O)--(230:2.6) node[pos=.7,below right]{$R_2$};
\node[point,label=below:$O$] at (O){};\node[point,label=above:$A$] at (A){};\node[point,label=right:$B$] at (B){};

\end{tikzpicture}
\figcap{题图：B 观察 A，视线 BA 与指向圆心的 BO 夹角为 $\theta$；图中未预设最大夹角位置。}\label{ch03:ref:question-figure-2-52}
\end{minipage}\end{center}
}

\newcommand{\KQuestionCF}{%
%
\begin{center}\begin{tikzpicture}[scale=.72,>=Latex]
 \draw[very thick] (-2.4,-2.4)--(-2.4,0) arc[start angle=180,end angle=0,radius=2.4]--(2.4,-2.4);
 \draw[very thick] (-1.4,-2.4)--(-1.4,0) arc[start angle=180,end angle=0,radius=1.4]--(1.4,-2.4);
 \draw[->,DeepBlue,very thick] (-1.4,-1.9)--(-1.4,-.85);
 \draw[->,Blue,very thick] (-2.4,-1.9)--(-2.4,-.85);
 \draw[<->,TextGray] (0,0)--(0,1.4) node[midway,clearlabel,right]{$R_1$};
 \draw[<->,TextGray] (0,0)--(135:2.4) node[pos=.85,clearlabel,above right=3pt]{$R_2$};
 \node[clearlabel,below] at (0,-2.75){内、外路线均为半圆弯道};
\end{tikzpicture}\figcap{题图：内、外半圆弯道半径分别为 $R_1,R_2$。}\label{ch03:ref:question-figure-2-57}\end{center}
}

\newcommand{\KQuestionCH}{%
%
\begin{center}\begin{minipage}{.96\linewidth}\centering
\begin{tikzpicture}[scale=.95,>=Latex]
 \draw[thick] (0,0) circle (2.0);\coordinate (P) at (35:2);
 \coordinate (V) at ($(P)+(125:1.4)$);\coordinate (A) at ($(P)+(170:1.4)$);
 \draw[kguide] (0,0)--(P) node[pos=.4,below right]{$R$};
 \draw[vector] (P)--(V) node[above right]{$\vect v$};
 \draw[vector] (P)--(A) node[above left]{$\vect a$};
 \draw pic[draw=TextGray,angle radius=7mm] {angle=V--P--A};
 % The circumference crosses this narrow wedge. Show its label separately,
 % retaining the physical vectors and the main angle arc without a mask.
 \begin{scope}[shift={(4.7,1.3)}]
 \draw[vector] (0,0)--(125:1.6) node[above right]{$\vect v$};
 \draw[vector] (0,0)--(170:1.6) node[above left]{$\vect a$};
 \draw[TextGray] (125:.6) arc(125:170:.6);
 \node at (147.5:.95){$\theta$};
 \node[below] at (-.5,-.4){夹角放大};
 \end{scope}
 \node[point,label=below left:$O$] at (0,0){};\node[point] at (P){};
\end{tikzpicture}
\figcap{题图：速度 $\vect v$ 与加速度 $\vect a$ 的夹角为 $\theta$。}\label{ch03:ref:question-figure-2-59}
\end{minipage}\end{center}
}


质点沿圆的运动称为圆周运动。最简单的是速率不变的匀速圆周运动；它的速度大小虽然不变，方向却时刻变化，因而加速度不为零。圆周运动普遍存在于机械、游乐设施和天体运动的近似模型中。

圆周运动的特殊性在于轨迹已经由圆心和半径确定，质点的位置只剩沿圆周的一个自由度。因而，与其反复给出平面直角坐标，不如用圆心到质点的半径相对某条固定射线的转角表示位置。历史上人们也常用圆周运动近似描述天体；即使真实轨道未必是圆，角位置、周期和切向—法向分解仍然是重要工具。

\subsection{角位置、角速度与周期}

\CirclePositionDiagram

以圆心 $O$、固定方向 $ON$ 为基准，半径 $OP$ 与 $ON$ 的有向夹角 $\theta$ 称为角位置。规定逆时针为正。角位置相差 $2\pi$ 的整数倍对应同一空间位置，这一点不同于直线坐标。

角位置的数值可以连续累积而不必每转一周就归零，例如 $\theta=0,2\pi,4\pi$ 虽对应同一几何点，却记录了不同的转动圈数。若只关心空间位置，可把角度按 $2\pi$ 取模；若关心总转角和经过时间，则必须保留完整的累积值。顺时针运动使 $\theta$ 减小，角速度因此为负，负号表示转向而不是“转得较慢”。

在时间 $\Delta t$ 内转过 $\Delta\theta$，平均角速度为
\begin{equation}
 \bar\omega=\frac{\Delta\theta}{\Delta t}.
\end{equation}
匀速圆周运动中 $\omega$ 为常量，故
\begin{equation}
 \theta(t)=\theta_0+\omega t.
\end{equation}

\AngularSignDiagram

匀速圆周运动中的“匀速”只指线速率和角速度大小不变，并不表示速度矢量恒定。经过相同时间转过相同角度，角位置随时间呈线性关系；但由于半径方向持续转动，切线方向也持续转动，速度矢量每一瞬间都在变化。

转一周所需时间为周期 $T$，单位时间内的转数为频率 $f$：
\begin{equation}
 |\omega|T=2\pi,\qquad T=\frac{2\pi}{|\omega|},\qquad
 f=\frac1T=\frac{|\omega|}{2\pi}.
\end{equation}
半径为 $R$ 的圆周上，有向弧坐标的变化为 $R\Delta\theta$；在这一小段不反向的运动中，实际路程是 $R|\Delta\theta|$。
若区间内转向，应逐段累计路程，不能只取总转角的绝对值。所以线速度大小为
\begin{equation}
 v=R|\omega|,
\end{equation}
方向沿轨迹切线，并与半径垂直。

频率的单位是赫兹，表示每秒完成的圈数，不要求取整数。周期和频率只描述转动的快慢；线速度还与半径有关。因此，在同一刚体上角速度相同的两点，离转轴越远，线速度越大。公式 $v=R|\omega|$ 中使用弧度制，若角速度以每秒多少度给出，必须先换算为弧度每秒。

角位置公式也能直接处理追及和相遇。两个质点在同一圆周上相遇，要求它们的角位置之差为 $2\pi$ 的整数倍；若只要求两点与圆心共线，则角位置之差为 $\pi$ 的整数倍。写方程时应给顺、逆时针角速度带上符号，而不是一律先取正数再口头判断。相对角速度为零时，相对角位置保持不变，可能始终重合、始终共线或永不满足条件，必须单独讨论。

对不同半径的同心圆，角速度和线速度不能混为一谈。已知线速度时要先除以各自半径得到角速度，再比较角位置；已知周期或频率时则可以直接转成角速度。涉及“第几次相遇”时，正整数既记录追上了多少圈，也排除了初始时刻可能已经相遇的情形。

从圆周外一点观察内圆上的质点时，视线方向并不等于质点的极角。视角达到极值的临界位置可由几何切线判断：从观察点引向内圆的视线与半径垂直。图~\ref{ch03:fig:k05-ex-2-52}把 $B$ 到 $A$ 的线画成内圆切线，因而最大视角只由 $R_1/R_2$ 决定；但左右两个极值在相对圆周上的弧间隔通常不同，所以相继出现的时间不一定相等。若只报告一个“周期”，必须说明是同一侧极值再次出现，还是左右极值交替出现。

圆周与抛体相遇的题目要同时满足空间和时间条件。例~\ref{ch03:ex:k05-2-51}中抛体落到圆所在水平面只决定一个时刻，水平位移等于 $R$ 才保证它落到圆周；圆周质点还必须在同一时刻回到该碰撞点。这三项条件缺一不可。图中初速度方向、圆心正上方高度以及圆周转向若有任何改变，整数周条件也应相应修改。

\begin{origexamplebox}{转速、角速度与线速度}\label{ch03:ex:k05-2-46}
质点沿半径 $r=10\,\mathrm{cm}$ 的圆周匀速运动，每分钟转 $30$ 圈。求角速度、线速度和周期。

\begin{qualityfigure}
\begin{tikzpicture}[>=Latex,font=\small]
\draw[DeepBlue,thick] (0,0) circle(1.5);\fill (0,0) circle(2pt) node[below left]{$O$};
\fill (1.5,0) circle(2pt) node[below right]{$P$};\draw[thin] (0,0)--(1.5,0) node[midway,below]{$r$};
\draw[->,DeepBlue,thick] (1.5,.1)--(1.5,1.3) node[right]{$v$};\node[align=center] at (0,-2.2){$r=10\,\mathrm{cm}$\\每分钟 $30$ 圈};
\end{tikzpicture}
\figcap{圆轨道半径、转动周期和切向速度}\label{mech:03:46:q}
\end{qualityfigure}
\end{origexamplebox}
\Needspace{6\baselineskip}\begin{xesolution}
\begin{qualityfigure}
\begin{tikzpicture}[>=Latex,font=\small]
\draw[DeepBlue,thick] (0,0) circle(1.5);\fill (0,0) circle(2pt) node[below left]{$O$};
\fill (1.5,0) circle(2pt) node[below right]{$P$};\draw[thin] (0,0)--(1.5,0) node[midway,below]{$r$};
\draw[->,DeepBlue,thick] (1.5,.1)--(1.5,1.3) node[right]{$v$};\node[align=center] at (0,-2.2){$r=10\,\mathrm{cm}$\\每分钟 $30$ 圈};\node at(3,-1.5){$v=\omega r$};
\end{tikzpicture}
\figcap{圆轨道半径、转动周期和切向速度（解答用辅助图）}\label{mech:03:46:a}
\end{qualityfigure}
图~\ref{mech:03:46:a}中，半径 $OP=r$，速度沿 $P$ 点切线；一周路程为 $2\pi r$，每秒转数是 $f$，故 $v=2\pi rf$、$|\omega|=2\pi f$。

频率 $f=30/60=0.5\,\mathrm{Hz}$，所以
\[
 \boxed{|\omega|=2\pi f=\pi\,\mathrm{rad/s}},\qquad
 \boxed{v=|\omega|r=0.1\pi\,\mathrm{m/s}},\qquad
 \boxed{T=1/f=2\,\mathrm s}.
\]
题目没有指定转动方向，所以只能确定角速度大小；若按逆时针为正，顺时针时应取负号。
弧度本身无量纲，但写出 $\mathrm{rad/s}$ 能明确这是角速度。
\end{xesolution}

\begin{origexamplebox}{径向匀速质点与圆周质点相撞}\label{ch03:ex:k05-2-47}
质点 $A$ 在半径 $R$ 的圆周上以角速度 $\omega$ 匀速运动。当 $A$ 到达给定点 $P$ 时，质点 $B$ 从圆心 $O$ 出发，沿 $OP$ 方向以速率 $v$ 匀速运动。若两质点能相撞，求 $v$ 的所有可能值。
\par 题图见图~\ref{ch03:ref:question-figure-2-47}。
\KQuestionBV
\end{origexamplebox}
\Needspace{6\baselineskip}\begin{xesolution}
以下推导所用的解图为图~\ref{ch03:fig:k05-ex-2-47}。

$B$ 只经过圆周上的 $P$ 点，到达时间为 $R/v$。$A$ 从 $P$ 出发后第 $n$ 次回到 $P$ 的时间为 $2\pi n/|\omega|$，其中 $n=1,2,\ldots$。令两者相等：
\RadialUniformCollisionDiagram
\[
 \boxed{v_n=\frac{|\omega|R}{2\pi n}},\qquad n=1,2,\ldots
\]
不能取 $n=0$，因为初始时 $B$ 位于圆心而非 $P$。
\end{xesolution}

\begin{origexamplebox}{径向匀加速质点与圆周质点相撞}\label{ch03:ex:k05-2-48}
条件同例~\ref{ch03:ex:k05-2-47}，但 $B$ 从圆心由静止沿 $OP$ 做加速度为 $a$ 的匀加速运动。若能相撞，求 $a$ 的所有可能值。
\par 题图见图~\ref{ch03:ref:question-figure-2-48}。
\KQuestionBW
\end{origexamplebox}
\Needspace{6\baselineskip}\begin{xesolution}
以下推导所用的解图为图~\ref{ch03:fig:k05-ex-2-48}。

$B$ 到达 $P$ 时 $R=at^2/2$。相撞时间仍须为 $t_n=2\pi n/|\omega|$，所以
\[
 \boxed{a_n=\frac{2R}{t_n^2}=\frac{R\omega^2}{2\pi^2n^2}},
 \qquad n=1,2,\ldots
\]
\RadialAcceleratedCollisionDiagram

单位为 $\mathrm{m/s^2}$，且式中 $\omega^2$ 消除了转向符号。
\end{xesolution}

\begin{origexamplebox}{同一圆周上的追及}\label{ch03:ex:k05-2-49}
两个质点在半径 $R$ 的同一圆周上同向匀速运动，角速度分别为 $\omega_1,\omega_2$，且 $\omega_2>\omega_1$。$t=0$ 时两者在圆心正上方的 $P$ 点相遇，均逆时针运动。求再次相遇的时间以及相遇位置与 $OP$ 的夹角 $\theta$。
\par 题图见图~\ref{ch03:ref:question-figure-2-49}。
\KQuestionBX
\end{origexamplebox}
\Needspace{6\baselineskip}\begin{xesolution}
以下推导所用的解图为图~\ref{ch03:fig:k05-ex-2-49}。

快者相对慢者多转一周时再次相遇：
\[
 (\omega_2-\omega_1)t=2\pi,
 \qquad
 \boxed{t=\frac{2\pi}{\omega_2-\omega_1}}.
\]
\SameCirclePursuitDiagram

相遇位置从 $OP$ 逆时针转过
\[
 \boxed{\theta\equiv\omega_1t
 =\frac{2\pi\omega_1}{\omega_2-\omega_1}\pmod{2\pi}}.
\]
若题目只要 $0\le\theta<2\pi$，应把右式减去适当的 $2\pi$ 整数倍。
\end{xesolution}

\begin{origexamplebox}{同心圆上两点共线}\label{ch03:ex:k05-2-50}
质点 $A,B$ 分别在半径 $R_1,R_2$ 的同心圆上匀速运动，线速度为 $v_A,v_B$。$t=0$ 时 $A,B$ 位于圆心同侧且其连线通过圆心。求以后两点连线通过圆心的所有时刻。
\par 题图见图~\ref{ch03:ref:question-figure-2-50}。
\KQuestionBY
\end{origexamplebox}
\Needspace{6\baselineskip}\begin{xesolution}
以下推导所用的解图为图~\ref{ch03:fig:k05-ex-2-50}。

两点和圆心共线等价于角位置之差为 $k\pi$。令
\[
\omega_A=\frac{v_A}{R_1},\qquad \omega_B=\frac{v_B}{R_2}
\]
\ConcentricCollinearityDiagram

（若方向相反，应给其中一个角速度带负号），则
\[
 \boxed{t_k=\frac{k\pi}{|\omega_A-\omega_B|}},\qquad k=0,1,2,\ldots
\]
当 $k$ 为偶数时两点在圆心同侧，为奇数时在异侧。若 $\omega_A=\omega_B$，两点在所有时刻都共线。
\end{xesolution}

\begin{origexamplebox}{平抛质点与水平圆周质点相撞}\label{ch03:ex:k05-2-51}
半径为 $R$ 的圆水平放置，质点 $A$ 在圆周上以角速度 $\omega$ 匀速运动。在圆心正上方高度 $H$ 的 $S$ 点水平抛出质点 $B$，初速方向正对 $A$ 的初始位置，初速率为 $v_0$。若 $B$ 与 $A$ 相撞，$v_0$ 和 $H$ 应满足什么条件？
\par 题图见图~\ref{ch03:ref:question-figure-2-51}。
\KQuestionBZ
\end{origexamplebox}
\Needspace{6\baselineskip}\begin{xesolution}
以下推导所用的解图为图~\ref{ch03:fig:k05-ex-2-51}。

$B$ 下落到圆所在水平面的时间为 $t=\sqrt{2H/g}$。它要到达圆周，须有
\[
 v_0t=R.
\]
\ProjectileCircleCollisionDiagram

同时，碰撞点是 $A$ 的初始位置，故 $A$ 必须整周返回：$|\omega|t=2\pi n$，$n=1,2,\ldots$。因此条件可写成
\[
 \boxed{H=\frac{2\pi^2n^2g}{\omega^2}},\qquad
 \boxed{v_0=\frac{|\omega|R}{2\pi n}},\qquad n=1,2,\ldots
\]
等价地，$v_0=R\sqrt{g/(2H)}$ 且 $|\omega|\sqrt{2H/g}$ 必须是 $2\pi$ 的正整数倍。
\end{xesolution}

\begin{origexamplebox}{外圈观察者看到的角摆动}\label{ch03:ex:k05-2-52}
$A,B$ 分别在半径 $R_1,R_2$ 的同心圆轨道上同向运动，$R_1<R_2$，角速度 $\omega_1>\omega_2$。从外圈的 $B$ 看，$A$ 在圆心 $O$ 左右摆动。求 $\angle ABO$ 的最大值，以及 $A$ 相继出现于最大夹角位置的时间间隔。
\par 题图见图~\ref{ch03:ref:question-figure-2-52}。
\KQuestionCA
\end{origexamplebox}
\Needspace{6\baselineskip}\begin{xesolution}
以下推导所用的解图为图~\ref{ch03:fig:k05-ex-2-52}。

夹角最大时，从 $B$ 引向内圆的视线 $BA$ 与内圆相切，所以 $OA\perp BA$。在直角三角形 $OAB$ 中
\ConcentricSightDiagram
\[
 \boxed{\theta_{\max}=\arcsin\frac{R_1}{R_2}}.
\]
设 $\delta=\angle AOB$，切点满足 $\cos\delta=R_1/R_2$。相对角速度为 $\Omega=\omega_1-\omega_2$。若把左右两个极值位置都算作“最大夹角”，相邻出现的时间间隔交替为
\[
 \boxed{\Delta t_1=\frac{2\delta}{\Omega},\qquad
 \Delta t_2=\frac{2\pi-2\delta}{\Omega}},
 \quad \delta=\arccos\frac{R_1}{R_2}.
\]
同一侧最大位置再次出现的周期为 $\boxed{2\pi/\Omega}$。这样写可避免把左右极值误当成等时间间隔。
\end{xesolution}

\begin{origexamplebox}{经典自转电子模型的矛盾}\label{ch03:ex:k05-2-53}
20世纪早期，实验表明电子具有一种类似“绕自身轴转动”的内禀性质。当时有人尝试用经典小球的自转来理解这一现象：把电子看成半径 $R=r_e\simeq3\times10^{-15}\,\mathrm m$ 的小球，并要求自转满足 $\omega R^2\simeq10^{-5}\,\mathrm{m^2/s}$。指出这个经典模型的致命缺陷。

\begin{qualityfigure}
\begin{tikzpicture}[>=Latex,font=\small]
\draw[DeepBlue,thick] (0,0) circle(1.5);\fill (0,0) circle(2pt) node[below left]{$O$};
\fill (1.5,0) circle(2pt) node[below right]{$P$};\draw[thin] (0,0)--(1.5,0) node[midway,below]{$R$};
\draw[->,DeepBlue,thick] (1.5,.1)--(1.5,1.3) node[right]{$v$};\draw[->,Gold] (-.5,1.65) arc(105:40:1.72);\node at (3,.3){$R=r_e$};\node at (3,-.5){$\omega R^2\simeq10^{-5}$};
\end{tikzpicture}
\figcap{经典自转模型中表面点的速度}\label{mech:03:53:q}
\end{qualityfigure}
\end{origexamplebox}
\Needspace{6\baselineskip}\begin{xesolution}
\begin{qualityfigure}
\begin{tikzpicture}[>=Latex,font=\small]
\draw[DeepBlue,thick] (0,0) circle(1.5);\fill (0,0) circle(2pt) node[below left]{$O$};
\fill (1.5,0) circle(2pt) node[below right]{$P$};\draw[thin] (0,0)--(1.5,0) node[midway,below]{$R$};
\draw[->,DeepBlue,thick] (1.5,.1)--(1.5,1.3) node[right]{$v$};\draw[->,Gold] (-.5,1.65) arc(105:40:1.72);\node at (3,.3){$R=r_e$};\node at (3,-.5){$\omega R^2\simeq10^{-5}$};\node at(3,-1.5){$v=\omega R$};
\end{tikzpicture}
\figcap{经典自转模型中表面点的速度（解答用辅助图）}\label{mech:03:53:a}
\end{qualityfigure}
图~\ref{mech:03:53:a}中，表面点 $P$ 到轴心的距离为 $R=r_e$，故切向速率是 $\omega R$；把题给 $\omega R^2$ 除以同一个 $R$ 后再与光速比较。

模型中电子表面线速度为
\[
 v=\omega R=\frac{\omega R^2}{R}
 \simeq\frac{10^{-5}}{3\times10^{-15}}
 \simeq3.3\times10^9\,\mathrm{m/s}.
\]
这约为光速的 $11$ 倍，违反相对论的速度上限。因此电子不能被当作具有该经典半径、作刚性自转的小球；“自旋”不是经典物体绕自身轴转动。
\end{xesolution}

\subsection{向心加速度}

\CentripetalDerivationDiagram

相隔 $\Delta t$ 的两速度大小均为 $v=|\omega|R$，方向转过 $\Delta\theta=|\omega|\Delta t$。速度三角形在 $\Delta t\to0$ 时给出
\[
 |\Delta\vect v|\simeq v\Delta\theta=v|\omega|\Delta t.
\]
因此匀速圆周运动的加速度指向圆心，大小
\begin{equation}
 \boxed{a_r=\omega^2R=\frac{v^2}{R}}.
\end{equation}

图~\ref{ch03:fig:k05-2-15}的推导还说明了方向。把 $A,B$ 两点的速度平移到同一首端，二者等长，夹角等于半径转过的角 $\Delta\theta$；连接两个终端得到 $\Delta\vect v$。当 $\Delta t$ 趋于零，速度三角形的顶角趋于零，底边方向趋于半径内侧，所以加速度与瞬时速度垂直并指向圆心。底边长度的一阶近似为 $v\Delta\theta$，除以 $\Delta t$ 得 $v|\omega|=\omega^2R$。这里的“向心”只说明加速度方向，并不意味着质点正在向圆心移动。

公式可以用两个极限作检查。在相同线速度下，半径越大，轨迹越平缓，$a_r=v^2/R$ 越小；在相同角速度下，半径越大，线速度也按 $R$ 增大，因而 $a_r=\omega^2R$ 反而越大。两个看似相反的结论来自所保持不变的量不同。量纲方面，$\omega^2R$ 为 $\mathrm{s^{-2}}\cdot\mathrm m=\mathrm{m/s^2}$，与加速度一致。

若周期已知，$a_r=4\pi^2R/T^2$；若频率已知，$a_r=4\pi^2f^2R$。这些写法在比较天体、转盘或同轴转动物体时很方便，但仍只适用于瞬时轨迹半径为 $R$ 的圆周运动，不能把普通曲线的坐标到某个任意点的距离代作 $R$。

向心加速度不是一种额外的“运动类型”，而是圆周运动总加速度沿半径方向的分量。研究受力时，不能把“向心力”再作为某个独立物体施加的新力；应先列出真实相互作用力，再把它们指向圆心的合力与 $ma_r$ 对应。本节只讨论运动学，但明确这一点可以防止后续动力学中重复计算力。

在用经典电子自转模型作估算时，表面速度 $v=\omega R$ 超过光速，说明模型本身失效，而不是说明电子真的能超光速。数量级与极限检查是圆周运动中很重要的审查手段：给定极小半径却要求固定的 $\omega R^2$，会使 $\omega R$ 反比于半径急剧增大，立刻暴露经典刚性小球图像的矛盾。

圆周问题中的几个几何关系可以互相联系：速度沿切线，向心加速度沿半径指向圆心；
径向运动质点只在其射线与圆周的交点处可能发生碰撞；圆外观察者的最大视角对应视线与圆相切。
先从这些几何约束确定可能的位置，再用时间方程检验是否同时到达，能减少无意义的代数根。

\begin{origexamplebox}{向心加速度的两个反算}\label{ch03:ex:k05-2-54}
（1）质点以 $v=10\,\mathrm{m/s}$ 在半径 $R=10\,\mathrm m$ 的圆轨道上匀速运动，求加速度大小。\par
（2）半径 $R=10\,\mathrm m$，向心加速度 $a_r=10\,\mathrm{m/s^2}$，求频率。

\begin{qualityfigure}
\begin{tikzpicture}[>=Latex,font=\small]
\draw[DeepBlue,thick] (0,0) circle(1.5);\fill (0,0) circle(2pt) node[below left]{$O$};
\fill (1.5,0) circle(2pt) node[below right]{$P$};\draw[thin] (0,0)--(1.5,0) node[midway,below]{$R$};
\draw[->,DeepBlue,thick] (1.5,.1)--(1.5,1.3) node[right]{$v$};
\end{tikzpicture}
\figcap{切向速度与指向圆心的加速度}\label{mech:03:54:q}
\end{qualityfigure}
\end{origexamplebox}
\Needspace{6\baselineskip}\begin{xesolution}
\begin{qualityfigure}
\begin{tikzpicture}[>=Latex,font=\small]
\draw[DeepBlue,thick] (0,0) circle(1.5);\fill (0,0) circle(2pt) node[below left]{$O$};
\fill (1.5,0) circle(2pt) node[below right]{$P$};\draw[thin] (0,0)--(1.5,0) node[midway,below]{$R$};
\draw[->,DeepBlue,thick] (1.5,.1)--(1.5,1.3) node[right]{$v$};\draw[->,Gold,very thick] (1.4,0)--(.4,0) node[midway,above]{$a_r$};\node at (3,-1.5){$v=\omega R$};
\end{tikzpicture}
\figcap{切向速度与指向圆心的加速度（解答用辅助图）}\label{mech:03:54:a}
\end{qualityfigure}
图~\ref{mech:03:54:a}中，$OP=R$，$\vect v$ 与半径垂直，向心加速度沿 $PO$ 指向圆心。匀速只说明速率不变，速度方向仍在改变，所以 $a_r=v^2/R=\omega^2R$。

（1）
\[
\boxed{a_r=\frac{v^2}{R}=10\,\mathrm{m/s^2}}.
\]
（2）由 $a_r=\omega^2R$ 得 $\omega=1\,\mathrm{rad/s}$，故
\[
 \boxed{f=\frac{\omega}{2\pi}=\frac1{2\pi}\,\mathrm{Hz}\approx0.159\,\mathrm{Hz}}.
\]
\end{xesolution}

\begin{origexamplebox}{开普勒型周期律对应的加速度}\label{ch03:ex:k05-2-55}
若多个质点绕共同圆心匀速运动，并满足 $T^2/R^3=C$，其中 $C$ 为常量，求向心加速度与半径的关系。

\begin{qualityfigure}
\begin{tikzpicture}[>=Latex,font=\small]
\draw[DeepBlue,thick] (0,0) circle(1.5);\fill (0,0) circle(2pt) node[below left]{$O$};
\fill (1.5,0) circle(2pt) node[below right]{$P_1$};\draw[thin] (0,0)--(1.5,0) node[midway,below]{$R_1$};
\draw[->,DeepBlue,thick] (1.5,.1)--(1.5,1.3) node[right]{$v$};\draw[dashed] (0,0) circle(2.2);\draw (0,0)--(-1.8,1.265) node[midway,above]{$R_2$};\fill (-1.8,1.265) circle(2pt) node[above left]{$P_2$};\node[align=left] at (4,.2){$T_1^2/R_1^3=C$\\$T_2^2/R_2^3=C$};
\end{tikzpicture}
\figcap{同心轨道满足同一个周期律常量}\label{mech:03:55:q}
\end{qualityfigure}
\end{origexamplebox}
\Needspace{6\baselineskip}\begin{xesolution}
\begin{qualityfigure}
\begin{tikzpicture}[>=Latex,font=\small]
\draw[DeepBlue,thick] (0,0) circle(1.5);\fill (0,0) circle(2pt) node[below left]{$O$};
\fill (1.5,0) circle(2pt) node[below right]{$P_1$};\draw[thin] (0,0)--(1.5,0) node[midway,below]{$R_1$};
\draw[->,DeepBlue,thick] (1.5,.1)--(1.5,1.3) node[right]{$v$};\draw[dashed] (0,0) circle(2.2);\draw (0,0)--(-1.8,1.265) node[midway,above]{$R_2$};\fill (-1.8,1.265) circle(2pt) node[above left]{$P_2$};\node[align=left] at (4,.2){$T_1^2/R_1^3=C$\\$T_2^2/R_2^3=C$};\draw[->,Gold,very thick] (1.4,0)--(.4,0) node[midway,above]{$a_r$};\node at (3,-1.5){$v=\omega R$};
\end{tikzpicture}
\figcap{同心轨道满足同一个周期律常量（解答用辅助图）}\label{mech:03:55:a}
\end{qualityfigure}
图~\ref{mech:03:55:a}中，两个轨道均以 $O$ 为圆心，但各自的周期随半径变化。对任一轨道先写 $a_r=(2\pi/T)^2R$，再把该轨道的 $T^2=CR^3$ 代入。

\[
a_r=\frac{4\pi^2R}{T^2}
=\boxed{\frac{4\pi^2}{CR^2}}.
\]
因此 $a_r\propto R^{-2}$，即随半径平方反比衰减。
\end{xesolution}

\begin{origexamplebox}{周期相同的同心转动}\label{ch03:ex:k05-2-56}
若多个质点绕共同圆心匀速运动，周期 $T$ 与各自轨道半径 $R$ 无关，求向心加速度与半径的关系。

\begin{qualityfigure}
\begin{tikzpicture}[>=Latex,font=\small]
\draw[DeepBlue,thick] (0,0) circle(1.5);\fill (0,0) circle(2pt) node[below left]{$O$};
\fill (1.5,0) circle(2pt) node[below right]{$P_1$};\draw[thin] (0,0)--(1.5,0) node[midway,below]{$R_1$};
\draw[->,DeepBlue,thick] (1.5,.1)--(1.5,1.3) node[right]{$v$};\draw[dashed] (0,0) circle(2.2);\draw (0,0)--(-1.8,1.265) node[midway,above]{$R_2$};\fill (-1.8,1.265) circle(2pt) node[above left]{$P_2$};\node[align=left] at (4,.2){$T_1=T_2=T$\\半径不同，周期相同};
\end{tikzpicture}
\figcap{半径不同而周期相同的同心轨道}\label{mech:03:56:q}
\end{qualityfigure}
\end{origexamplebox}
\Needspace{6\baselineskip}\begin{xesolution}
\begin{qualityfigure}
\begin{tikzpicture}[>=Latex,font=\small]
\draw[DeepBlue,thick] (0,0) circle(1.5);\fill (0,0) circle(2pt) node[below left]{$O$};
\fill (1.5,0) circle(2pt) node[below right]{$P_1$};\draw[thin] (0,0)--(1.5,0) node[midway,below]{$R_1$};
\draw[->,DeepBlue,thick] (1.5,.1)--(1.5,1.3) node[right]{$v$};\draw[dashed] (0,0) circle(2.2);\draw (0,0)--(-1.8,1.265) node[midway,above]{$R_2$};\fill (-1.8,1.265) circle(2pt) node[above left]{$P_2$};\node[align=left] at (4,.2){$T_1=T_2=T$\\半径不同，周期相同};\draw[->,Gold,very thick] (1.4,0)--(.4,0) node[midway,above]{$a_r$};\node at (3,-1.5){$v=\omega R$};
\end{tikzpicture}
\figcap{半径不同而周期相同的同心轨道（解答用辅助图）}\label{mech:03:56:a}
\end{qualityfigure}
图~\ref{mech:03:56:a}中，两个轨道在同一周期内各转一周，故角速度大小均为 $2\pi/T$；向心加速度箭头指向 $O$，其大小由 $\omega^2R$ 决定。

周期相同意味着角速度的大小相同，即 $|\omega|=2\pi/T$；转动方向可以不同。故
\[
 \boxed{a_r=\frac{4\pi^2}{T^2}R\propto R}.
\]
这与刚体绕固定轴转动时各点具有相同角速度的情况一致。
\end{xesolution}

\begin{origexamplebox}{赛车走弯道内侧还是外侧}\label{ch03:ex:k05-2-57}
赛车在圆弧赛道上允许的最大向心加速度为 $a$，直线加、减速的最大加速度也为 $a$。
弯道为半圆，内外半径分别为 $R_1<R_2$，直道最高速率为 $v_0>\sqrt{aR_1}$。
按“直道完成加减速、半圆内匀速”的模型，比较两条路线的过弯时间；
再把入弯制动和出弯加速纳入比较。两端取相同横断面，直道总长为 $S$ 且足够完成加减速，起终速率均为 $v_0$。
\par 题图见图~\ref{ch03:ref:question-figure-2-57}。
\KQuestionCF
\end{origexamplebox}
\Needspace{6\baselineskip}\begin{xesolution}
以下推导所用的解图为图~\ref{ch03:fig:k05-ex-2-57}。

先看半圆弧。由 $v^2/R\le a$，加上不能超过最高速率的限制，弯道所用速率为
\[
 u(R)=\min\{v_0,\sqrt{aR}\}.
\]
内道一定需要制动，但外道是否需要制动，取决于 $v_0$ 是否大于 $\sqrt{aR_2}$。
\RacetrackDiagram
图中的内、外圆弧长分别为 $\pi R_1,\pi R_2$。对一般半径 $R$，
\[
 t_{\rm arc}(R)=\frac{\pi R}{u(R)}
 =\begin{cases}\pi\sqrt{R/a},&R<v_0^2/a,\\
 \pi R/v_0,&R\ge v_0^2/a.\end{cases}
\]
两段都随 $R$ 增大，分界处连续，所以只计圆弧时内道更快。

再计加减速。若弯道速率为 $u<v_0$，以最大加速度制动和加速，每段分别需要
\[
 t_b=\frac{v_0-u}{a},\qquad s_b=\frac{v_0^2-u^2}{2a}.
\]
直道余下的 $S-2s_b$ 以 $v_0$ 行驶。先扣除两段距离，避免重复计算直道时间：
\begin{align*}
 T(R)&=\frac{S-2s_b}{v_0}+2t_b+\frac{\pi R}{u}\\
 &=\frac S{v_0}+\frac{\pi R}{u}+\frac{(v_0-u)^2}{av_0}.
\end{align*}
无需制动时取 $u=v_0$，最后一项自然为零。代入两种可能的弯道速率：
\[
 T(R)-\frac S{v_0}
 =\begin{cases}
 \displaystyle\frac{v_0}{a}+(\pi-2)\sqrt{\frac Ra}+\frac R{v_0},
 &R<v_0^2/a,\\[2mm]
 \displaystyle\frac{\pi R}{v_0},&R\ge v_0^2/a.
 \end{cases}
\]
第一段中含 $R$ 的两项都递增，第二段也递增，分界处同为 $\pi v_0/a$。因此
\[
 \boxed{T(R_1)<T(R_2)}.
\]
内道较快，但理由不能只说“路程短”：内道还承担更大的制动、加速时间损失，必须一并比较。
各项量纲都是时间；当 $u\to v_0$ 时附加损失趋于零。
本模型不含弯内变速、横向变道或切弯；允许这些操作时需另建加速度和路线约束。
\end{xesolution}

\subsection{变速圆周运动}

一般圆周运动的瞬时角速度和角加速度定义为
\begin{equation}
 \omega(t)=\lim_{\Delta t\to0}\frac{\Delta\theta}{\Delta t},
 \qquad
 \beta(t)=\lim_{\Delta t\to0}\frac{\Delta\omega}{\Delta t}.
\end{equation}
任一时刻仍有
\begin{equation}
 v=R|\omega|,\qquad a_r=\omega^2R=\frac{v^2}{R},
 \qquad a_\theta=R\beta.
\end{equation}
$a_\theta$ 是沿角度增加方向的有符号切向分量，$a_r$ 是指向圆心的法向分量大小，二者互相垂直。
若改用沿实际速度方向的自然切向，则在 $\omega\ne0$ 时
\[
 a_t=\frac{\dd v}{\dd t}=R\,\operatorname{sgn}(\omega)\beta.
\]
逆时针转动时 $a_t=a_\theta$，顺时针转动时 $a_t=-a_\theta$；减速或加速应看 $a_t$ 的符号。
若 $\beta$ 为常量，
\begin{equation}
 \omega(t)=\omega_0+\beta t,\qquad
 \theta(t)=\theta_0+\omega_0t+\frac12\beta t^2.
\end{equation}

变速圆周运动中，每一瞬间仍可把运动看成“沿切线改变速率”和“向圆心改变方向”的合成。即便 $\omega$ 随时间改变，法向分量仍为当时的 $v^2/R$，不能把它漏掉；切向分量则由角速度变化产生。总加速度大小为 $\sqrt{a_t^2+a_r^2}$，方向由两分量共同决定。只有 $\beta=0$ 时切向分量消失，运动才退化为匀速圆周运动。

角加速度的定义与直线运动的加速度完全类比：在有限时间间隔中先定义平均角加速度 $\Delta\omega/\Delta t$，再令间隔趋于零得到瞬时值。若 $\beta$ 恒定，角速度图是直线，图下有向面积给角位移；若 $\beta$ 随时间变化，则 $\omega(t)=\omega_0+\int\beta\,\dd t$，$\theta(t)=\theta_0+\int\omega\,\dd t$。这些积分关系与直线运动中的 $a$-$t$、$v$-$t$ 图完全平行。

注意 $a_\theta=R\beta$ 带有方向符号，而 $a_r=\omega^2R$ 的大小总为非负。
角速度过零时，向心加速度瞬间为零，但 $a_\theta$ 仍可非零，质点随后反向转动。
在瞬时静止处速度方向未定，不能直接定义“沿速度的单位切向”；仍可使用预先规定的角度正向处理。

\begin{origexamplebox}{圆周边缘达到声速}\label{ch03:ex:k05-2-58}
质点在半径 $R=11\,\mathrm m$ 的圆周上运动，$t=0$ 时 $\omega_0=0$，角加速度 $\beta=1\,\mathrm{rad/s^2}$。求边缘速度达到声速 $340\,\mathrm{m/s}$ 所需时间。

\begin{qualityfigure}
\begin{tikzpicture}[>=Latex,font=\small]
\draw[DeepBlue,thick] (0,0) circle(1.5);\fill (0,0) circle(2pt) node[below left]{$O$};
\fill (1.5,0) circle(2pt) node[below right]{$P$};\draw[thin] (0,0)--(1.5,0) node[midway,below]{$R$};
\draw[->,DeepBlue,thick] (1.5,.1)--(1.5,1.3) node[right]{$v$};\draw[->,Gold] (-.5,1.7) arc(105:40:1.8);\node at (3,-.5){$\omega_0=0,\ \beta>0$};
\end{tikzpicture}
\figcap{角加速度引起边缘切向速率增长}\label{mech:03:58:q}
\end{qualityfigure}
\end{origexamplebox}
\Needspace{6\baselineskip}\begin{xesolution}
\begin{qualityfigure}
\begin{tikzpicture}[>=Latex,font=\small]
\draw[DeepBlue,thick] (0,0) circle(1.5);\fill (0,0) circle(2pt) node[below left]{$O$};
\fill (1.5,0) circle(2pt) node[below right]{$P$};\draw[thin] (0,0)--(1.5,0) node[midway,below]{$R$};
\draw[->,DeepBlue,thick] (1.5,.1)--(1.5,1.3) node[right]{$v$};\draw[->,Gold] (-.5,1.7) arc(105:40:1.8);\node at (3,-.5){$\omega_0=0,\ \beta>0$};\draw[->,Gold,very thick](2,0)--(2,1)node[right]{$a_t$};\node at(3,-1.5){$a_t=R\beta$};
\end{tikzpicture}
\figcap{角加速度引起边缘切向速率增长（解答用辅助图）}\label{mech:03:58:a}
\end{qualityfigure}
图~\ref{mech:03:58:a}中，点 $P$ 的速率沿切线，$v=R\omega$；题给的 $\beta$ 控制角速度变化，先求 $\omega=\beta t$，再由线速度达到声速确定时间。

\[
v=R\omega=R\beta t,
\]
所以
\[
 \boxed{t=\frac{340}{11\times1}\,\mathrm s\approx30.9\,\mathrm s}.
\]
这里不能把 $\beta$ 直接当作线加速度；切向加速度为 $R\beta=11\,\mathrm{m/s^2}$。
\end{xesolution}


\begin{origexamplebox}{加速度与速度夹角恒定}\label{ch03:ex:k05-2-59}
质点沿半径为 $R$ 的圆轨道运动，初速率为 $v_0$，其加速度方向与速度方向的夹角 $\theta$ 恒定。求速度大小随时间的关系。
\par 题图见图~\ref{ch03:ref:question-figure-2-59}。
\KQuestionCH
\end{origexamplebox}
\Needspace{6\baselineskip}\begin{xesolution}
以下推导所用的解图为图~\ref{ch03:fig:k05-ex-2-59}。

把加速度沿切向和指向圆心的法向分解。若图示加速度的法向分量指向圆心，则
\[
a_t=a\cos\theta=\frac{\dd v}{\dd t},
\qquad
a_r=a\sin\theta=\frac{v^2}{R}.
\]
\AccelerationVelocityAngleDiagram

消去 $a$：
\[
\frac{\dd v}{\dd t}=\frac{v^2}{R}\cot\theta.
\]
由 $v(0)=v_0$ 积分得
\[
 -\frac1v+\frac1{v_0}=\frac{\cot\theta}{R}t,
\]
即
\[
 \boxed{v(t)=\frac{v_0}{1-(v_0t/R)\cot\theta}}.
\]
$\theta=90^\circ$ 时 $v=v_0$；$\theta<90^\circ$ 时加速，$\theta>90^\circ$ 时减速。分母为零前模型才有效。
\end{xesolution}
}

{
\graphicspath{{../../../sample/xes-original/kinematics-supplement/}{../../../sample/xes-original/kinematics-supplement/assets/}}
\section{平面内的一般曲线运动}
\newcommand{\FixedCoordinatesDiagram}{%
\begin{center}\begin{tikzpicture}[scale=.78,>=Latex]
 \draw[kaxis] (-.5,0)--(6.4,0) node[right,clearlabel]{$x$};
 \draw[kaxis] (0,-.4)--(0,4.2) node[above,clearlabel]{$y$};
 \draw[kcurve] plot[smooth] coordinates{(.5,.5)(1.4,2.8)(2.8,3.4)(4.1,2.0)(5.7,3.2)};
 \coordinate (A) at (1.4,2.8);\coordinate (B) at (4.1,2.0);
 \node[point,label={[clearlabel]above left:$A$}] at (A) {};
 \node[point,label={[clearlabel]below right:$B$}] at (B) {};
 \draw[kguide] (A)--(1.4,0) node[clearlabel,below]{$x(t_1)$};\draw[kguide] (A)--(0,2.8) node[clearlabel,left]{$y(t_1)$};
 \draw[kguide] (B)--(4.1,0) node[clearlabel,below]{$x(t_2)$};\draw[kguide] (B)--(0,2) node[clearlabel,left]{$y(t_2)$};
\end{tikzpicture}\figcap{分量运动方程 $x(t),y(t)$ 在每个时刻唯一确定曲线上的位置。}\label{ch03:fig:k06-2-16}\end{center}}

\newcommand{\ParametricEllipseDiagram}{%
\begin{center}\begin{minipage}{.96\linewidth}\centering
\begin{tikzpicture}[scale=.85,>=Latex]
 \draw[kaxis] (-3.8,0)--(3.9,0) node[right]{$x$};\draw[kaxis] (0,-2.3)--(0,2.5) node[above]{$y$};
 \draw[kcurve] (0,0) ellipse (3 and 1.65);
 \draw[vector,domain=10:40,samples=25] plot({3*cos(\x)},{1.65*sin(\x)});
 \node[point,label=below right:{$t=0$}] at (3,0){};
 \draw[<->,TextGray] (0,-.3)--(3,-.3) node[midway,below]{$a$};
 \draw[<->,TextGray] (.3,0)--(.3,1.65) node[midway,right]{$b$};
\end{tikzpicture}
\figcap{闭合椭圆的半轴为 $a,b$；短箭头表示 $\omega>0$ 时的运动方向。}\label{ch03:fig:k06-ex-2-61}
\end{minipage}\end{center}}

\newcommand{\CycloidDiagram}{%
\begin{center}\begin{tikzpicture}[scale=.72,>=Latex]
 \draw[very thick] (-.4,0)--(10,0);
 \draw[kaxis] (0,0)--(10.2,0) node[right,clearlabel]{$x$};\draw[kaxis] (0,0)--(0,3.4) node[above,clearlabel]{$y$};
 \draw[kcurve,domain=0:6.283,samples=120] plot ({1.35*(\x-sin(\x r))},{1.35*(1-cos(\x r))});
 \draw[very thick] (0,1.35) circle (1.35);\node[point,label={[clearlabel]below left:$P$}] at (0,0) {};
 \draw[vector] (0,1.35)--(1.1,1.35) node[pos=.22,clearlabel,above=3pt]{$v$};
\end{tikzpicture}\figcap{无滑动滚动圆周点的摆线。尖点处圆周点与地面瞬时相对静止。}\label{ch03:fig:k06-ex-2-62}\end{center}}

\newcommand{\HelixDiagram}{%
\begin{center}\begin{tikzpicture}[scale=.76,>=Latex,
  x={(1cm,0cm)},y={(-.50cm,-.28cm)},z={(0cm,.80cm)}]
 \draw[kaxis] (0,0,0)--(3.1,0,0) node[right,clearlabel]{$x$};
 \draw[kaxis] (0,0,0)--(0,2.7,0) node[below left,clearlabel]{$y$};
 \draw[kaxis] (0,0,0)--(0,0,5.2) node[above,clearlabel]{$z$};
 \draw[TextGray,dashed,domain=0:360,samples=90] plot ({1.4*cos(\x)},{1.4*sin(\x)},0);
 \draw[TextGray,dashed,domain=0:360,samples=90] plot ({1.4*cos(\x)},{1.4*sin(\x)},4.4);
 \foreach \a in {0,90,180,270}{\draw[TextGray!65] ({1.4*cos(\a)},{1.4*sin(\a)},0)--({1.4*cos(\a)},{1.4*sin(\a)},4.4);}
 \draw[kcurve,->,domain=0:720,samples=220]
   plot ({1.4*cos(\x)},{1.4*sin(\x)},{.0061*\x});
 \draw[<->,TextGray] (0,0,-.75)--(1.4,0,-.75) node[midway,clearlabel,below]{$R$};
 \draw[<->,TextGray] (1.72,0,0)--(1.72,0,2.2) node[midway,clearlabel,right]{$2\pi v/\omega$};
 \node[clearlabel,right,align=left] at (2.15,0,3.8){水平投影为圆\\竖直方向匀速上升};
\end{tikzpicture}\figcap{解图：$xy$ 投影是圆，$z$ 坐标匀速增加，合成轨迹为圆柱螺旋线。}\label{ch03:fig:k06-ex-2-63}\end{center}}

\newcommand{\CircleComponentsDiagram}{%
\begin{center}\begin{tikzpicture}[scale=.78,>=Latex]
  \draw[kaxis] (-3.1,0)--(3.25,0) node[right,clearlabel]{$x$};
  \draw[kaxis] (0,-3.0)--(0,3.15) node[above,clearlabel]{$y$};
  \draw[very thick] (0,0) circle (2.25);
  \coordinate (O) at (0,0);\coordinate (P) at (38:2.25);
  \draw[very thick] (O)--(P) node[midway,clearlabel,below right]{$R$};
  \draw[kguide] (P)--(P|-O) node[clearlabel,below]{$x$};
  \draw[kguide] (P)--(P-|O) node[clearlabel,left]{$y$};
  \draw[vector] (P)--++(128:1.25) node[midway,clearlabel,right]{$\vect v$};
  \draw[vector] (P)--++(218:1.15) node[midway,clearlabel,above left]{$\vect a$};
  \node[point,label={[clearlabel]below left:$O$}] at (O) {};
  \node[point,label={[clearlabel]above right:$P$}] at (P) {};
\end{tikzpicture}\figcap{圆周运动的位置投影为 $x,y$；速度沿切线，加速度沿半径指向圆心。}\label{ch03:fig:k06-ex-2-64}\end{center}}

\newcommand{\RollingContactDiagram}{%
\begin{center}\begin{tikzpicture}[scale=.8,>=Latex]
  \draw[very thick] (-3.7,0)--(3.7,0);
  \draw[very thick] (0,1.55) circle (1.55);
  \coordinate (O) at (0,1.55);\coordinate (P) at (0,0);
  \draw[vector] (O)--++(1.45,0) node[midway,clearlabel,above]{$\vect v_O$};
  \draw[vector] (P)--++(-1.45,0) node[midway,clearlabel,below=2pt]{$\vect v_{P/O}$};
  \draw[vector] (P)--++(1.45,0);
  \node at (1.7,.3){$\vect v_O$};
  \draw[DeepBlue,very thick,->] ({1.85*cos(30)},{1.55+1.85*sin(30)}) arc[start angle=30,end angle=-20,radius=1.85];
  \node[clearlabel,right] at (2.05,1.15){$\omega=v/R$};
  \draw[very thick] (O)--(P) node[midway,clearlabel,right]{$R$};
  \node[point] at (P) {};\node[below=4pt] at (P){$P$};
  \node[point,label={[clearlabel]above left:$O$}] at (O) {};
  \node[clearlabel] at (0,-1.1){$\vect v_P=\vect0$};
\end{tikzpicture}\figcap{接地点 $P$ 的平动速度与绕圆心转动速度大小相等、方向相反。}\label{ch03:fig:k06-ex-2-65}\end{center}}

\newcommand{\ArchimedeanSpiralDiagram}{%
\begin{center}\begin{tikzpicture}[scale=.64,>=Latex]
  \begin{scope}[xshift=-5.7cm]
    \coordinate (O) at (0,0);\coordinate (P) at (38:2.5);
    \node[point,label={[clearlabel]below left:$O$}] at (O) {};
    \draw[kaxis] (O)--(3,0) node[right]{$N$};
    \draw[vector] (O)--(P) node[midway,above left=3pt]{$r=vt$};
    \node[point,label=above:$P$] at (P){};
    \draw[TextGray] (.9,0) arc[start angle=0,end angle=38,radius=.9];
    \node at (1.35,.32){$\theta$};
    \node[clearlabel,below] at (1,-1.0){$\theta=\omega t$};
  \end{scope}
  \begin{scope}[xshift=2.7cm]
    \draw[kaxis] (-4.3,0)--(4.5,0) node[right,clearlabel]{$x$};
    \draw[kaxis] (0,-4.0)--(0,4.2) node[above,clearlabel]{$y$};
    \draw[kcurve,->,domain=0:18.3,samples=180] plot ({.21*\x*cos(\x r)},{.21*\x*sin(\x r)});
    \node[point] at (0,0){};\node at (.3,-.35){$O$};
    \node[clearlabel,below] at (0,-4.55){$r=(v/\omega)\theta$};
  \end{scope}
\end{tikzpicture}\figcap{左图在一个时刻确定 $P$ 的极坐标；右图以 $O$ 为极点，连续描出 $r\propto\theta$ 的螺线。}\label{ch03:fig:k06-ex-2-66}\end{center}}

\newcommand{\NaturalCoordinatesDiagram}{%
\begin{center}\begin{minipage}{.98\linewidth}\centering
\begin{tikzpicture}[scale=1.25,>=Latex]
 \draw[kcurve,domain=-2.7:2.7,samples=120] plot(\x,{.15*\x^3-.6*\x});
 \coordinate (P) at (-1.4,.4284);\coordinate (Q) at (-.2,.1188);
 \pgfmathsetmacro{\ang}{atan(.282)}
 \draw[vector] (P)--++(\ang:1.5) node[above=5pt]{$\vect v(t),\ \vect e_t$};
 \draw[vector] (Q)--++({atan(-.582)}:1.25) node[below right=2pt]{$\vect v(t+\Delta t)$};
 \draw[vector] (P)--++({\ang-90}:1.15) node[below right]{$a_n,\ \vect e_n$};
 \draw[vector] (P)--++({\ang+180}:.8) node[above left=3pt]{$a_t<0$};
 \node[point,label=above left:$P$] at (P){};\node[point] at (Q){};\node[above=5pt] at (Q){$Q$};
\end{tikzpicture}
\figcap{切向沿速度，法向指向 $P$ 处曲线的凹侧；图中切向加速度取减速情形。}
\label{ch03:fig:k06-natural-frame}\end{minipage}\end{center}}

\newcommand{\VelocityIncrementNaturalDiagram}{%
\begin{center}\begin{tikzpicture}[scale=.82,>=Latex]
  \coordinate (A) at (-4.0,0);\coordinate (B) at (0,0);\coordinate (C) at (2.25,2.7);
  \coordinate (D) at (2.25,0);\coordinate (E) at (4.1,0);
  \draw[vector] (A)--(B) node[midway,clearlabel,below]{$\vect v(t)$};
  \draw[vector] (B)--(C);\node at (.2,1.2){$\vect a\,\Delta t$};
  \draw[vector] (A)--(C) node[midway,clearlabel,above left]{$\vect v(t+\Delta t)$};
  \draw[kguide] (C)--(D);
  \draw[kguide] (B)--(D);
  \draw[<->,TextGray] (B)--(D) node[midway,clearlabel,below=3pt]{$a\cos\alpha\,\Delta t$};
  \draw[<->,TextGray] (D)--(C) node[midway,clearlabel,right=3pt]{$a\sin\alpha\,\Delta t$};
  \draw pic[draw=TextGray,->,angle radius=7mm] {angle=E--B--C};
  \node[clearlabel] at ($(B)+(1.1,.45)$){$\alpha$};
  \draw pic[draw=TextGray,->,angle radius=10mm] {angle=B--A--C};
  \node at (-2.45,.34){$\Delta\theta$};
  \node[clearlabel,below left] at (A){$A$};\node[clearlabel,below left] at (B){$B$};
  \node[clearlabel,above] at (C){$C$};\node[clearlabel,below right] at (D){$D$};
\end{tikzpicture}\figcap{速度三角形把 $\vect a\Delta t$ 分解为平行、垂直于原速度的两部分，分别决定速率变化和方向转角。}\label{ch03:fig:k06-2-17}\end{center}}

\newcommand{\OsculatingCircleDiagram}{%
\begin{center}\begin{minipage}{.96\linewidth}\centering
\begin{tikzpicture}[scale=1,>=Latex]
 \draw[TextGray,thick] (0,{5/3}) circle ({5/3});
 \draw[kcurve,domain=-2.4:2.4,samples=100] plot(\x,{.3*\x*\x});
 \draw[thick] (-2.5,0)--(2.5,0) node[right]{$S$};
 \draw[TextGray,->] (0,0)--(0,{5/3}) node[midway,right=4pt]{$\rho$};
 \draw[TextGray] (.18,0)--(.18,.18)--(0,.18);
 \node[point,label=below:$P$] at (0,0){};
 \node[point,label=above left:$O$] at (0,{5/3}){};
\end{tikzpicture}
\figcap{圆与曲线在 $P$ 处共切线 $S$、共曲率；半径 $PO$ 垂直于 $S$。}\label{ch03:fig:k06-2-18}
\end{minipage}\end{center}}

\newcommand{\FoldedCurvePath}{(-2.80000000,1.40000000)
 ..controls (-1.73540121,0.68166038) and (-0.67080242,-0.03667924)..(0.00000000,0.00000000)
 ..controls (0.67080242,0.03667924) and (0.94780848,0.82837734)..(1.45000000,1.55000000)
 ..controls (1.80153407,2.05513586) and (2.26340901,2.52593477)..(2.65000000,2.65000000)
 ..controls (2.81568185,2.70317081) and (2.96753604,2.69265585)..(3.00000000,2.65000000)
 ..controls (3.04328528,2.59312553) and (2.87432126,2.47911170)..(2.75000000,2.30000000)
 ..controls (2.56351810,2.03133245) and (2.47748239,1.61619468)..(2.52000000,1.30000000)
 ..controls (2.59086268,0.77300887) and (3.01881791,0.52085968)..(3.45000000,0.40000000)
 ..controls (3.88118209,0.27914032) and (4.31559105,0.28957016)..(4.75000000,0.30000000)}

\newcommand{\FoldedCurvePoints}{%
\node[point,label={[label distance=3pt]below left:{$P_1$}}] at (0.00000000,0.00000000) {};
\node[point,label={[label distance=3pt]above left:{$P_2$}}] at (1.45000000,1.55000000) {};
\node[point,label={[label distance=3pt]above right:{$P_3$}}] at (3.00000000,2.65000000) {};
\node[point,label={[label distance=3pt]left:{$P_4$}}] at (2.52000000,1.30000000) {};
}

\newcommand{\FoldedCurveCircles}{%
\begin{scope}
\clip (-3.0,-1.1) rectangle (5,3.9);
\draw[TextGray,thin] (-0.04766861,0.87177976) circle (0.87308204);
\draw[TextGray,thin] (5.68213522,-1.39522683) circle (5.15609635);
\draw[TextGray,thin] (2.97003756,2.62719657) circle (0.03765295);
\draw[TextGray,thin] (3.59629028,1.44472504) circle (1.08597703);
\end{scope}
}

\newcommand{\FoldedCurveTipZoom}{%
\begin{scope}[scale=10]
\clip (2.88,2.54) rectangle (3.08,2.78);
\draw[kcurve] \FoldedCurvePath;
\draw[TextGray,thin] (2.97003756,2.62719657) circle (0.03765295);
\end{scope}
}

\newcommand{\FourCurvaturesDiagram}{%
\begin{center}\begin{minipage}{.98\linewidth}\centering
\begin{tikzpicture}[scale=1.15,>=Latex]
 \FoldedCurveCircles
 \draw[kcurve] \FoldedCurvePath;
 \FoldedCurvePoints
\end{tikzpicture}
\figcap{题图的回折曲线与对应密切圆；大圆仅显示局部圆弧。}\label{ch03:fig:k06-ex-2-69}
\par\medskip
\begin{tikzpicture}[>=Latex]
\FoldedCurveTipZoom
\end{tikzpicture}
\par\smallskip {\small $P_3$ 邻域放大十倍：灰色小圆为此处的密切圆。}
\end{minipage}\end{center}}

\newcommand{\ProjectileCurvatureDiagram}{%
\begin{center}\begin{tikzpicture}[scale=.78,>=Latex]
  \begin{scope}[xshift=-3.4cm]
    \draw[kcurve,->,domain=-2:2,samples=70] plot (\x,{1.75-.48*\x*\x});
    \coordinate (U) at (-.9,{1.75-.48*.9*.9});
    \draw[vector] (U)--++({atan(.864)}:1.15) node[midway,clearlabel,above left]{$\vect v$};
    \draw[vector] (U)--++(0,-1.2) node[midway,clearlabel,right]{$\vect g$};
    \node[clearlabel,below] at (0,-1.35){上升：$\vect g\cdot\vect v<0$};
  \end{scope}
  \begin{scope}[xshift=3.4cm]
    \draw[kcurve,->,domain=-2:2,samples=70] plot (\x,{1.75-.48*\x*\x});
    \coordinate (D) at (.9,{1.75-.48*.9*.9});
    \draw[vector] (D)--++({-atan(.864)}:1.15) node[midway,clearlabel,above right]{$\vect v$};
    \draw[vector] (D)--++(0,-1.2) node[midway,clearlabel,left]{$\vect g$};
    \node[clearlabel,below] at (0,-1.35){下降：$\vect g\cdot\vect v>0$};
  \end{scope}
  \begin{scope}[yshift=-3.8cm]
    \draw[kaxis] (-2.4,0)--(2.4,0) node[right,clearlabel]{$x$};
    \draw[kaxis] (0,-.3)--(0,1.7) node[above,clearlabel]{$y$};
    \draw[kcurve,domain=-1.7:1.7,samples=70] plot (\x,{.45*\x*\x});
    \draw[TextGray] (0,{1/.9}) circle ({1/.9});
    \node[point,label={[clearlabel]below right:$x=0$}] at (0,0) {};
    \draw[<->,TextGray] (0,0)--(0,{1/.9}) node[midway,clearlabel,right]{$\rho(0)$};
  \end{scope}
\end{tikzpicture}\figcap{上升、下降阶段的速度—重力夹角决定速率增减；下图在抛物线顶点画出密切圆。}\label{ch03:fig:k06-ex-2-68}\end{center}}

\newcommand{\EllipseCurvatureDiagram}{%
\begin{center}\begin{tikzpicture}[scale=.74,>=Latex]
  \draw[kaxis] (-4.0,0)--(4.0,0) node[right,clearlabel]{$x$};
  \draw[kaxis] (0,-2.6)--(0,2.6) node[above,clearlabel]{$y$};
  \draw[very thick] (0,0) ellipse (3.25 and 1.65);
  \coordinate (A) at (3.25,0);\coordinate (B) at (0,1.65);
  \draw[TextGray] ({3.25-1.65*1.65/3.25},0) circle ({1.65*1.65/3.25});
  \begin{scope}
    \clip (-4.0,-2.55) rectangle (4.0,2.55);
    \draw[TextGray] (0,{1.65-3.25*3.25/1.65}) circle ({3.25*3.25/1.65});
  \end{scope}
  \draw[<->,TextGray] (0,-2.05)--(3.25,-2.05) node[midway,clearlabel,below]{$a$};
  \draw[<->,TextGray] (.25,0)--(.25,1.65) node[midway,clearlabel,right]{$b$};
  \node[point] at (A) {};\node[below right=4pt] at (A){$A$};
  \node[point,label={[clearlabel]above:$B$}] at (B) {};
  \draw[TextGray] (2.62,.78)--(3.28,1.45);
  \node[clearlabel,above] at (3.15,1.52){$\rho_A=b^2/a$};
  \draw[TextGray] (-1.45,1.86)--(-2.25,2.18);
  \node[clearlabel,left] at (-2.18,2.20){$\rho_B=a^2/b$};
\end{tikzpicture}\figcap{椭圆长轴端 $A$ 的密切圆较小，短轴端 $B$ 的密切圆较大。}\label{ch03:fig:k06-ex-2-70}\end{center}}

\newcommand{\KQuestionCI}{%
\begin{center}\begin{minipage}{.97\linewidth}\centering
\begin{tikzpicture}[>=Latex,scale=.9]

\draw[kaxis] (-3,0)--(3.2,0) node[right]{$x$};
\draw[kaxis] (0,-1.8)--(0,2) node[above]{$y$};\node[below left] at (0,0){$O$};

\end{tikzpicture}
\figcap{题图：根据给定参数方程在空坐标系中画出轨迹。}\label{ch03:ref:question-figure-2-60}
\end{minipage}\end{center}
}

\newcommand{\KQuestionCJ}{%
\begin{center}\begin{minipage}{.97\linewidth}\centering
\begin{tikzpicture}[>=Latex,scale=.9]

\draw[kaxis] (-3,0)--(3.2,0) node[right]{$x$};
\draw[kaxis] (0,-1.8)--(0,2) node[above]{$y$};\node[below left] at (0,0){$O$};

\end{tikzpicture}
\figcap{题图：根据给定参数方程在空坐标系中画出轨迹。}\label{ch03:ref:question-figure-2-61}
\end{minipage}\end{center}
}

\newcommand{\KQuestionCK}{%
%
\begin{center}\begin{tikzpicture}[scale=.72,>=Latex]
 \draw[very thick] (-.4,0)--(10,0);
 \draw[kaxis] (0,0)--(10.2,0) node[right,clearlabel]{$x$};\draw[kaxis] (0,0)--(0,3.4) node[above,clearlabel]{$y$};
 \draw[kcurve,domain=0:6.283,samples=120] plot ({1.35*(\x-sin(\x r))},{1.35*(1-cos(\x r))});
 \draw[very thick] (0,1.35) circle (1.35);\node[point,label={[clearlabel]below left:$P$}] at (0,0) {};
 \draw[vector] (0,1.35)--(1.1,1.35) node[pos=.22,clearlabel,above=3pt]{$v$};
\end{tikzpicture}\figcap{题图：圆沿水平地面无滑动滚动，P 为圆周上的标记点，初始位于接触处。}\label{ch03:ref:question-figure-2-62}\end{center}
}

\newcommand{\KQuestionCL}{%
\begin{center}\begin{minipage}{.97\linewidth}\centering
\begin{tikzpicture}[>=Latex,scale=.85,x={(1cm,0cm)},y={(-.50cm,-.28cm)},z={(0cm,.80cm)}]

\draw[kaxis] (0,0,0)--(3,0,0) node[right]{$x$};
\draw[kaxis] (0,0,0)--(0,2.8,0) node[below left]{$y$};
\draw[kaxis] (0,0,0)--(0,0,3.5) node[above]{$z$};
\node[below right] at (0,0,0){$O$};

\end{tikzpicture}
\figcap{题图：在右手直角坐标系中，根据三个坐标随时间的变化画出轨迹。}\label{ch03:ref:question-figure-2-63}
\end{minipage}\end{center}
}

\newcommand{\KQuestionCM}{%
%
\begin{center}\begin{tikzpicture}[scale=.78,>=Latex]
  \draw[kaxis] (-3.1,0)--(3.25,0) node[right,clearlabel]{$x$};
  \draw[kaxis] (0,-3.0)--(0,3.15) node[above,clearlabel]{$y$};
  \draw[very thick] (0,0) circle (2.25);
  \coordinate (O) at (0,0);\coordinate (P) at (38:2.25);
  \draw[very thick] (O)--(P) node[midway,clearlabel,below right]{$R$};
  \draw[kguide] (P)--(P|-O) node[clearlabel,below]{$x$};
  \draw[kguide] (P)--(P-|O) node[clearlabel,left]{$y$};
  \node[point,label={[clearlabel]below left:$O$}] at (O) {};
  \node[point,label={[clearlabel]above right:$P$}] at (P) {};
\end{tikzpicture}\figcap{题图：圆周上 P 点在两个坐标轴上的投影分别为 $x,y$。}\label{ch03:ref:question-figure-2-64}\end{center}
}

\newcommand{\KQuestionCN}{%
%
\begin{center}\begin{tikzpicture}[scale=.8,>=Latex]
  \draw[very thick] (-3.7,0)--(3.7,0);
  \draw[very thick] (0,1.55) circle (1.55);
  \coordinate (O) at (0,1.55);\coordinate (P) at (0,0);
  \draw[vector] (O)--++(1.45,0) node[midway,clearlabel,above]{$\vect v_O$};
  \draw[DeepBlue,very thick,->] ({1.85*cos(30)},{1.55+1.85*sin(30)}) arc[start angle=30,end angle=-20,radius=1.85];
  \node[clearlabel,right] at (2.05,1.15){$\omega$};
  \draw[very thick] (O)--(P) node[midway,clearlabel,right]{$R$};
  \node[point] at (P) {};\node[below=4pt] at (P){$P$};
  \node[point,label={[clearlabel]above left:$O$}] at (O) {};
\end{tikzpicture}\figcap{题图：圆无滑动向右滚动，P 为此刻与地面接触的圆周点。}\label{ch03:ref:question-figure-2-65}\end{center}
}

\newcommand{\KQuestionCO}{%
\begin{center}\begin{minipage}{.97\linewidth}\centering
\begin{tikzpicture}[>=Latex,scale=.9]

\draw[kaxis] (0,0)--(6,0) node[right]{$r$};\node[below left] at (0,0){$O$};
\path[use as bounding box] (-.5,-1.35) rectangle (6.4,1.35);

\end{tikzpicture}
\figcap{题图：以 O 为极点、向右为极轴，描出给定极坐标运动。}\label{ch03:ref:question-figure-2-66}
\end{minipage}\end{center}
}

\newcommand{\KQuestionCQ}{%
\begin{center}\begin{minipage}{.97\linewidth}\centering
\begin{tikzpicture}[>=Latex,scale=.9]

\draw[kaxis] (-2.2,0)--(2.3,0) node[right]{$x$};
\draw[kaxis] (0,-.7)--(0,2.7) node[above]{$y$};
\draw[kcurve,domain=-1.8:1.8,samples=100] plot(\x,{.65*\x*\x});
\node[below left] at (0,0){$O$};\node[right] at (1.8,2.1){$y=ax^2$};

\end{tikzpicture}
\figcap{题图：确定抛物线在原点处的曲率半径；这里示意 $a>0$。}\label{ch03:ref:question-figure-2-68}
\end{minipage}\end{center}
}

\newcommand{\KQuestionCR}{%
\begin{center}\begin{minipage}{.97\linewidth}\centering
\begin{tikzpicture}[>=Latex,scale=1.15]
\draw[kcurve] \FoldedCurvePath;\FoldedCurvePoints
\end{tikzpicture}
\figcap{题图：比较 $P_1,P_2,P_3,P_4$ 的局部弯曲，注意 $P_3$ 后向左下回折。}\label{ch03:ref:question-figure-2-69}
\end{minipage}\end{center}
}

\newcommand{\KQuestionCS}{%
\begin{center}\begin{minipage}{.97\linewidth}\centering
\begin{tikzpicture}[>=Latex,scale=.9]

\draw[thick] (0,0) ellipse(3 and 1.6);\draw (0,0)--(3,0) node[midway,below]{$a$};
\draw (0,0)--(0,1.6) node[midway,left]{$b$};\node[below left] at (0,0){$O$};

\end{tikzpicture}
\figcap{题图：椭圆的半长轴为 $a$、半短轴为 $b$。}\label{ch03:ref:question-figure-2-70}
\end{minipage}\end{center}
}


抛体和圆周运动提供了分析一般曲线运动的基本工具。对平面运动常用两种方法：一是在固定的直角或极坐标系中研究各坐标分量；二是跟随质点，在瞬时速度方向和与之垂直的方向上分解运动学量。前者便于写运动方程，后者能直接揭示速率与轨迹弯曲的变化。

前面两类曲线运动都有特殊结构：抛体具有恒定加速度，圆周运动具有固定曲率。一般二维轨迹未必满足这些条件，但描述它所需的语言已经齐备。固定坐标系把复杂运动拆成若干一维分运动，适合从已知函数计算；自然坐标系随着质点一起转动，把加速度直接分成改变速率和改变方向的两部分，适合解释轨迹为什么在某处弯得更急。

\subsection{固定坐标系}

\FixedCoordinatesDiagram

在平面直角坐标系 $Oxy$ 中，任意运动可由
\begin{equation}
 x=x(t),\qquad y=y(t)
\end{equation}
描述。两个分量分别看似一维运动，合在一起唯一确定质点的位置。联立二式并消去 $t$，有时可以得到轨迹方程；但并非所有参数曲线都能方便地写成显式函数 $y(x)$。

坐标轴的取向没有统一规定，应服从问题的几何结构。研究抛体时常取水平、竖直方向；研究地面导航时可取东西、南北方向。给定 $t$，先由两个运动方程分别求得 $x(t),y(t)$，图~\ref{ch03:fig:k06-2-16}中的交点就是质点此刻的位置。随着 $t$ 连续改变，这些位置连成轨迹。消去参数 $t$ 得到的轨迹方程只保留空间几何形状，通常会丢失运动方向、快慢以及同一点被经过的时刻，因此不能代替完整的运动方程。


\begin{origexamplebox}{圆的参数方程}\label{ch03:ex:k06-2-60}
已知 $x(t)=A\cos\omega t$，$y(t)=A\sin\omega t$。求轨迹方程并画出示意图。
\par 题图见图~\ref{ch03:ref:question-figure-2-60}。
\KQuestionCI
\end{origexamplebox}
\Needspace{6\baselineskip}\begin{xesolution}
以下推导所用的解图为图~\ref{mech:legacy:ch03:k06-2-60}。

平方相加消去 $t$：
\[
 \boxed{x^2+y^2=A^2}.
\]
轨迹是以原点为圆心、半径 $A$ 的圆。$t=0$ 时在 $(A,0)$；若 $\omega>0$，随后 $y$ 增大，故逆时针运动。
\begin{center}\begin{qualityfigure}
\begin{tikzpicture}[scale=.65]
\draw[kaxis] (-2.5,0)--(2.5,0) node[right,clearlabel]{$x$};\draw[kaxis] (0,-2.5)--(0,2.5) node[above,clearlabel]{$y$};
\draw[kcurve] (0,0) circle (2);
\draw[vector] (10:2) arc[start angle=10,end angle=40,radius=2];
\node[point,label={[clearlabel]below right:$t=0$}] at (2,0) {};
\end{tikzpicture}
\figcap{本题解答中的位置、分量与辅助量。}\label{mech:legacy:ch03:k06-2-60}
\end{qualityfigure}\end{center}
\end{xesolution}


\begin{origexamplebox}{椭圆的参数方程}\label{ch03:ex:k06-2-61}
已知 $x(t)=a\cos\omega t$，$y(t)=b\sin\omega t$。求轨迹方程并画出示意图。
\par 题图见图~\ref{ch03:ref:question-figure-2-61}。
\KQuestionCJ
\end{origexamplebox}
\Needspace{6\baselineskip}\begin{xesolution}
以下推导所用的解图为图~\ref{ch03:fig:k06-ex-2-61}。

\[
 \left(\frac xa\right)^2+\left(\frac yb\right)^2
 =\cos^2\omega t+\sin^2\omega t=1,
\]
\ParametricEllipseDiagram

故
\[
 \boxed{\frac{x^2}{a^2}+\frac{y^2}{b^2}=1}.
\]
若 $a>b$，长轴沿 $x$ 轴；若 $b>a$，长轴沿 $y$ 轴。运动周期为 $2\pi/|\omega|$，但沿椭圆的速率一般不恒定。
\end{xesolution}

\begin{origexamplebox}{滚动圆周点的坐标}\label{ch03:ex:k06-2-62}
半径为 $R$ 的圆在平面上无滑动地向右匀速滚动，圆心速度为 $v$。圆周固定点 $P$ 在 $t=0$ 时位于接地点。求 $P$ 在图示坐标系中的 $x,y$ 坐标。
\par 题图见图~\ref{ch03:ref:question-figure-2-62}。
\KQuestionCK
\end{origexamplebox}
\Needspace{6\baselineskip}\begin{xesolution}
以下推导所用的解图为图~\ref{ch03:fig:k06-ex-2-62}。

无滑动条件给出圆转过的角度 $\phi=vt/R$。圆心坐标为 $(vt,R)$；从圆心指向 $P$ 的半径由竖直向下顺时针转过 $\phi$，分量为 $(-R\sin\phi,-R\cos\phi)$。相加得
\CycloidDiagram
\[
 \boxed{x(t)=vt-R\sin\frac{vt}{R}},\qquad
 \boxed{y(t)=R\left(1-\cos\frac{vt}{R}\right)}.
\]
这条轨迹称为摆线。$t=0$ 时两式均为零，与题图原点一致。
\end{xesolution}


\begin{origexamplebox}{空间螺旋线}\label{ch03:ex:k06-2-63}
质点在直角坐标系 $Oxyz$ 中满足
\[
x(t)=R\cos\omega t,\qquad y(t)=R\sin\omega t,
\qquad z(t)=vt,
\]
其中各参数均为正。画出轨迹示意图。
\par 题图见图~\ref{ch03:ref:question-figure-2-63}。
\KQuestionCL
\end{origexamplebox}
\Needspace{6\baselineskip}\begin{xesolution}
以下推导所用的解图为图~\ref{ch03:fig:k06-ex-2-63}。

$x^2+y^2=R^2$，说明质点始终在半径 $R$ 的圆柱面上；同时 $z$ 随时间匀速增加。每当 $\omega t$ 增加 $2\pi$，质点转一周而高度增加
\HelixDiagram
\[
 \boxed{p=v\frac{2\pi}{\omega}}.
\]
因此轨迹是绕 $z$ 轴上升的圆柱螺旋线，$p$ 为螺距。
\end{xesolution}

曲线运动的速度和加速度均可逐分量求导：
\begin{equation}
 \vect v=(\dot x,\dot y),\qquad
 \vect a=(\ddot x,\ddot y).
\end{equation}
反过来，给定矢量也可投影到各坐标轴。运动的合成与分解对速度、加速度同样成立。

上述等式来自固定基矢不随时间改变。若 $\vect r=x\vect e_x+y\vect e_y$，则 $\vect v=\dot x\vect e_x+\dot y\vect e_y$，再求导得到加速度。各分量可以分别使用一维运动的方法，但最后必须在同一时刻按矢量合成。三维情形只需再加入 $z(t)$ 及其一、二阶导数，方法没有本质变化；例~\ref{ch03:ex:k06-2-63}正是这种推广。

极坐标中的分量要更谨慎，因为径向和横向单位矢量本身会随 $\theta$ 转动。短时间内，半径长度改变 $\dot r\Delta t$ 产生径向速度，半径方向转过 $\dot\theta\Delta t$ 产生长度约为 $r\dot\theta\Delta t$ 的横向位移。因此
\begin{equation}
 \vect v=\dot r\,\vect e_r+r\dot\theta\,\vect e_\theta.
\end{equation}
例~\ref{ch03:ex:k06-2-66}中 $\dot r=v$、$r\dot\theta=vt\omega$，随时间增长横向分量越来越大；例~\ref{ch03:ex:k06-2-67}中两分量相等，才得到速度与半径恒成 $45^\circ$。不能把 $\dot\theta$ 本身当作线速度，因为它的量纲是角速度，还必须乘当地半径。

固定坐标法的计算流程通常是：先写位置分量，再求导得到速度分量与速率，继续求导得到加速度分量，最后根据需要消去时间或投影到特定方向。若只给轨迹方程而没有参数随时间的规律，通常只能知道几何形状，不能唯一确定速度和加速度；同一条轨迹可以被不同速率、不同方向甚至往返地经过。


\begin{origexamplebox}{圆周运动的分量速度和加速度}\label{ch03:ex:k06-2-64}
质点 $P$ 在半径 $R$ 的圆上以角速度 $\omega$ 逆时针匀速运动，$t=0$ 时位于 $x$ 轴正向。已知 $x=R\cos\omega t$、$y=R\sin\omega t$，求两坐标投影的速度和加速度。
\par 题图见图~\ref{ch03:ref:question-figure-2-64}。
\KQuestionCM
\end{origexamplebox}
\Needspace{6\baselineskip}\begin{xesolution}
以下推导所用的解图为图~\ref{ch03:fig:k06-ex-2-64}。

逐次求导：
\begin{align*}
v_x&=-R\omega\sin\omega t,&v_y&=R\omega\cos\omega t,\\
a_x&=-R\omega^2\cos\omega t,&a_y&=-R\omega^2\sin\omega t.
\end{align*}
因此
\[
\vect a=-\omega^2(x,y),
\]
\CircleComponentsDiagram

始终指向圆心，大小 $R\omega^2$；速度大小为 $R\omega$，且与半径垂直。
\end{xesolution}

\begin{origexamplebox}{滚动圆的接地点瞬时速度}\label{ch03:ex:k06-2-65}
证明：半径为 $R$ 的圆以圆心速度 $v$ 无滑动滚动时，圆周固定点 $P$ 恰与地面接触的瞬时速度为零。
\par 题图见图~\ref{ch03:ref:question-figure-2-65}。
\KQuestionCN
\end{origexamplebox}
\Needspace{6\baselineskip}\begin{xesolution}
以下推导所用的解图为图~\ref{ch03:fig:k06-ex-2-65}。

质点速度等于圆心平动速度与绕圆心转动速度的矢量和。圆向右滚动时角速度大小 $\omega=v/R$、方向为顺时针。接地点的转动速度水平向左，大小
\[
\omega R=v,
\]
\RollingContactDiagram

恰与圆心向右的平动速度抵消。因此
\[
 \boxed{\vect v_P=\vect v_O+\vect v_{P/O}=\vect0}.
\]
“瞬时速度为零”不表示该点有一段时间静止；接触前后它仍沿摆线运动。
\end{xesolution}

\begin{origexamplebox}{阿基米德螺线}\label{ch03:ex:k06-2-66}
平面极坐标中质点满足 $r=vt$、$\theta=\omega t$，所有参数为正。定性画出轨迹。
\par 题图见图~\ref{ch03:ref:question-figure-2-66}。
\KQuestionCO
\end{origexamplebox}
\Needspace{6\baselineskip}\begin{xesolution}
以下推导所用的解图为图~\ref{ch03:fig:k06-ex-2-66}。

消去 $t$：
\[
 \boxed{r=\frac v\omega\theta}.
\]
半径随转角线性增加，质点从极点出发逆时针向外盘旋，相邻两圈的径向间距为
\[
\Delta r=\frac v\omega(2\pi)=\frac{2\pi v}{\omega}.
\]
这是一条阿基米德螺线。
\ArchimedeanSpiralDiagram
\end{xesolution}

\begin{origexamplebox}{对数螺线的等角性质}\label{ch03:ex:k06-2-67}
平面极坐标中质点满足 $r=r_0e^\theta$、$\theta=\omega t$，其中 $r_0>0,\omega>0$。
证明速度与从极点指向质点的径向射线夹角恒定，并求该角。可用 $e^{x+\Delta x}\simeq e^x(1+\Delta x)$。

\begin{qualityfigure}
\begin{tikzpicture}[>=Latex,font=\small]
\draw[->] (-.2,0)--(4,0) node[right]{$x$};\draw[->] (0,-.2)--(0,3.3) node[above]{$y$};
\draw[DeepBlue,thick,domain=-1:1.1,samples=80] plot({exp(\x)*cos(\x*180/pi)},{exp(\x)*sin(\x*180/pi)});
\coordinate (P) at (1.55,1.36);\fill (P) circle(2pt) node[left]{$P$};
\draw[dashed] (0,0)--(P) node[midway,above left]{$r$};\node[below left] at (0,0){$O$};\draw[->,DeepBlue] (P)--(1.64,2.73) node[right]{$\vect v$};
\end{tikzpicture}
\figcap{对数螺线上径向与横向速度分解}\label{mech:03:67:q}
\end{qualityfigure}
\end{origexamplebox}
\Needspace{6\baselineskip}\begin{xesolution}
\begin{qualityfigure}
\begin{tikzpicture}[>=Latex,font=\small]
\draw[->] (-.2,0)--(4,0) node[right]{$x$};\draw[->] (0,-.2)--(0,3.3) node[above]{$y$};
\draw[DeepBlue,thick,domain=-1:1.1,samples=80] plot({exp(\x)*cos(\x*180/pi)},{exp(\x)*sin(\x*180/pi)});
\coordinate (P) at (1.55,1.36);\fill (P) circle(2pt) node[left]{$P$};
\draw[dashed] (0,0)--(P) node[midway,above left]{$r$};\node[below left] at (0,0){$O$};\draw[->,Gold] (P)--(2.68,2.35) node[right]{$v_r$};\draw[->,Gold] (P)--(.56,2.49) node[left]{$v_\theta$};\draw[->,DeepBlue,very thick] (P)--(1.69,3.48) node[above]{$\vect v$};\node at (3.3,3){$v_r=v_\theta$};
\end{tikzpicture}
\figcap{对数螺线上径向与横向速度分解（解答用辅助图）}\label{mech:03:67:a}
\end{qualityfigure}
图~\ref{mech:03:67:a}中，$\vect v_r$ 沿向外径向，$\vect v_\theta$ 垂直于径向且沿逆时针方向；两分量相等，合速度位于两者之间，由直角速度三角形确定夹角。

极坐标速度的径向、横向分量分别是
\[
v_r=\dot r=\frac{\dd r}{\dd\theta}\dot\theta=r\omega,
\qquad
v_\theta=r\dot\theta=r\omega.
\]
两分量大小相等且互相垂直，所以速度与径向夹角 $\alpha$ 满足
\[
 \tan\alpha=\frac{v_\theta}{v_r}=1,
 \qquad \boxed{\alpha=45^\circ}.
\]
由于 $\omega>0$，径向和横向分量都为正，速度位于两正向基矢之间，故取锐角 $45^\circ$，不能仅凭 $\tan\alpha=1$ 忽略象限。
如果反向运动，轨迹仍相同，但速度方向反转，与向外径向射线的夹角变为 $135^\circ$；
若只说与无向直线的锐夹角，则仍为 $45^\circ$。
这就是对数螺线的等角性质。
\end{xesolution}

\subsection{自然坐标系与密切圆}

\NaturalCoordinatesDiagram

把瞬时速度方向称为切向，把与其垂直且指向轨迹弯曲一侧的方向称为法向。速度只有切向分量；加速度一般可写成
\[
\vect a=a_t\vect e_t+a_n\vect e_n.
\]
自然坐标系的方向随质点运动而改变，因此其物理意义直观，但计算时要注意基矢本身也在转动。

切向单位矢量 $\vect e_t$ 沿轨迹前进方向，法向单位矢量 $\vect e_n$ 指向曲线当前的凹侧。它们不像固定坐标轴那样处处平行：质点从一点运动到另一点，$\vect e_t$ 和 $\vect e_n$ 都会随轨迹转动。这正是为什么即使速率恒定，$\vect v=v\vect e_t$ 的导数仍可能非零。

在短时间 $\Delta t$ 内，速度由 $\vect v(t)$ 变为 $\vect v(t)+\vect a\Delta t$。设加速度与速度夹角为 $\alpha$。由速度三角形并保留 $\Delta t$ 的一阶项，速率变化为
\begin{equation}
 \Delta v=a\cos\alpha\,\Delta t,
 \qquad \boxed{a_t=\frac{\dd v}{\dd t}}.
\end{equation}
所以切向分量同向时速率增加，反向时减小，垂直时速率不变。速度方向转角 $\Delta\theta$ 则满足
\begin{equation}
 \Delta\theta\simeq\frac{a\sin\alpha}{v}\Delta t.
\end{equation}
若以半径 $\rho$ 的圆周运动模拟这一瞬时弯曲，$\Delta\theta/\Delta t=v/\rho$，故
\begin{equation}
 \boxed{a_n=\frac{v^2}{\rho}}.
\end{equation}

\VelocityIncrementNaturalDiagram

这一结论可由图~\ref{ch03:fig:k06-2-17}逐步看出。速度三角形的一条边为 $\vect v(t)$，第二条边为 $\vect a\Delta t$，合成后得到 $\vect v(t+\Delta t)$。把 $\vect a\Delta t$ 投影到旧速度方向，得到速率的一阶改变量 $a\cos\alpha\,\Delta t$；投影到垂直方向，得到速度方向的小转角 $\Delta\theta\simeq a\sin\alpha\,\Delta t/v$。这里舍去的是 $\Delta t^2$ 及更高阶项，所以不能把有限时间的整段曲线直接当成同一个圆。

\OsculatingCircleDiagram

$\rho$ 称为曲率半径，对应圆称为密切圆。直线由两个无限接近的点决定切线；类似地，曲线上一个给定点及其两侧无限接近的点决定密切圆。$\rho$ 越小，轨迹弯曲越强。

密切圆只在所考察点附近与原曲线具有相同的切线和二阶弯曲趋势，它不要求与整条曲线重合。圆心必在该点法线上并位于凹侧。曲线接近拐点时曲率趋于零，曲率半径趋于无穷，局部运动接近直线；尖锐转弯处曲率大，密切圆半径小。用 $a_n=v^2/\rho$ 判断时，必须同时使用当地速率和当地曲率，不能只凭图上“看起来弯”来决定法向加速度大小。

下面推导曲率公式，而不把它当成记忆规则。设切线与 $x$ 轴的夹角为 $\varphi$，
取沿曲线的弧长为 $s$。在足够短的一段上，密切圆的弧长关系是
$|\dd\varphi|=\dd s/\rho$，故曲率 $1/\rho=|\dd\varphi/\dd s|$。
对显式曲线 $y=y(x)$，有 $\tan\varphi=y'$。两边对 $x$ 求导，得到
\[
 \frac{\dd\varphi}{\dd x}=\frac{y''}{1+(y')^2},
 \qquad \dd s=\sqrt{1+(y')^2}\,|\dd x|.
\]
把转角变化除以弧长变化，再取绝对值，就有
\[
 \frac1\rho=\frac{|y''|}{[1+(y')^2]^{3/2}}.
\]
所以曲率半径为
\begin{equation}
 \rho=\frac{[1+(y')^2]^{3/2}}{|y''|};
\end{equation}
对参数曲线 $x=x(t),y=y(t)$，在 $\dot x\ne0$ 的局部，
链式法则给出
\[
 y'=\frac{\dot y}{\dot x},\qquad
 y''=\frac1{\dot x}\frac{\dd}{\dd t}\frac{\dot y}{\dot x}
 =\frac{\dot x\ddot y-\dot y\ddot x}{\dot x^3}.
\]
代回上式，分子、分母中的 $|\dot x|^3$ 抵消，得到
\begin{equation}
 \rho=\frac{(\dot x^2+\dot y^2)^{3/2}}
 {|\dot x\ddot y-\dot y\ddot x|}.
\end{equation}
这个结果关于 $x,y$ 对称；在 $\dot x=0,\dot y\ne0$ 的位置，可以交换两坐标重复推导，公式仍成立。
它要求曲线二次可微，并且 $\dot x^2+\dot y^2>0$，即此处参数化是正则的。
在这一前提下，分母为零对应曲率为零，可将曲率半径理解为趋于无穷；它不一定是拐点，还需看凹向是否改变。
若速度也为零，可能出现 $0/0$，不能据此断定曲率为零。例如摆线尖点处就需要从相邻点取极限另行分析。
绝对值保证半径为正，曲线凹向由法线方向判断。
分母越大不单独保证曲率越大，因为分子也随参数化改变；整个比值表达的几何曲率与正则参数选择无关。

自然坐标表达把一般平面加速度统一写为
\begin{equation}
 \boxed{\vect a=\frac{\dd v}{\dd t}\vect e_t+\frac{v^2}{\rho}\vect e_n}.
\end{equation}
第一项决定速度大小，第二项决定速度方向。直线运动是 $\rho\to\infty$ 的极限，法向项消失；匀速圆周运动是 $\dd v/\dd t=0$ 且 $\rho=R$ 的特例。这样，前面两类运动都成为一般曲线运动公式的自然极限。

图~\ref{ch03:fig:k06-2-18}中，切线 $S$ 经过接触点 $P$，半径 $PO$ 与它垂直。圆与曲线在 $P$ 处有相同的二阶弯曲，
但远离 $P$ 后可以分开。因此密切描述的是局部关系，不是用一个圆代替整条轨迹。

四点曲率比较时，先确定切线与凹向，再比较局部切线方向改变的快慢。
接近拐点时，单位弧长引起的转角很小，密切圆半径很大；转弯急处则相反。
仅仅过点、相切还不够，密切圆还必须具有相同的曲率，才能用于 $a_n=v^2/\rho$。

阻力型直线运动的两个例题也可放进自然坐标框架：轨迹为直线，$a_n=0$，加速度全在切向。$a=-kv$ 给出指数衰减，速度在任何有限时刻都不严格为零但总路程有限；$a=-\beta v^2$ 给出反比例衰减，速度同样渐近于零，却因 $v\sim1/t$ 使无限时间总路程呈对数发散。不同阻力规律会产生不同的长时间行为，不能仅凭“最终越来越慢”认为结果相同。

\begin{origexamplebox}{抛体的速率变化与抛物线曲率}\label{ch03:ex:k06-2-68}
填空：抛体上升时加速度和速度夹角为（锐、钝）角，速率（增加、减小）；下降时夹角为（锐、钝）角，速率（增加、减小）。并求抛物线 $y=ax^2$ 在 $x=0$ 处的曲率半径。
\par 题图见图~\ref{ch03:ref:question-figure-2-68}。
\KQuestionCQ
\end{origexamplebox}
\Needspace{6\baselineskip}\begin{xesolution}
以下推导所用的解图为图~\ref{ch03:fig:k06-ex-2-68}。

重力加速度竖直向下。上升时速度有向上分量，故 $\vect a\cdot\vect v<0$，夹角为钝角，速率减小；下降时 $\vect a\cdot\vect v>0$，夹角为锐角，速率增加。


\ProjectileCurvatureDiagram
平面曲线 $y(x)$ 的曲率半径为
\[
\rho=\frac{[1+(y')^2]^{3/2}}{|y''|}.
\]
对 $y=ax^2$，在 $x=0$ 有 $y'=0$、$y''=2a$，所以
\[
 \boxed{\rho(0)=\frac1{2|a|}}.
\]
\end{xesolution}


\begin{origexamplebox}{比较曲线上四点的弯曲程度}\label{ch03:ex:k06-2-69}
评估题图曲线上 $P_1,P_2,P_3,P_4$ 四点的弯曲情况，并尽可能画出这些点的密切圆。
\par 题图见图~\ref{ch03:ref:question-figure-2-69}。
\KQuestionCR
\end{origexamplebox}
\Needspace{6\baselineskip}\begin{xesolution}
以下推导所用的解图为图~\ref{ch03:fig:k06-ex-2-69}。

先沿题图从左向右上读到 $P_3$，再沿曲线向左下回折至 $P_4$：
$P_4$ 位于 $P_3$ 左下方，不能把这段回折曲线当成单值函数 $y(x)$ 的下降支。
作密切圆时先在各点画切线，圆心必须在该点法线上，并处在曲线局部凹向的一侧。
$P_1$ 位于宽缓谷底，圆心在上方；$P_2$ 附近弯曲较弱，密切圆半径较大；
$P_3$ 位于急剧转弯处，半径很小；$P_4$ 附近切线近乎竖直，曲线向右弯，圆心在右侧。
这些判断给出各圆的位置和相对大小的定性图示，不能从没有精确函数或比例尺的示意图中读出精确半径，
也不应据此强行断言所有四点的严格曲率排序。


\FourCurvaturesDiagram
若已知参数曲线 $(x(u),y(u))$，记 $\vect r'=(x',y')$、$\vect r''=(x'',y'')$，
则在 $x'y''-y'x''\ne0$ 的点，半径和圆心分别为
\[
 \rho=\frac{(x'^2+y'^2)^{3/2}}{|x'y''-y'x''|},\qquad
 \vect C=\vect r+\frac{x'^2+y'^2}{x'y''-y'x''}(-y',x').
\]
第二式中的 $(-y',x')$ 与切线 $(x',y')$ 点积为零，保证半径垂直切线；
分母的符号决定圆心在哪一侧。解图中的圆按同一条参数曲线的导数确定，
各圆经过对应点、与其共切线并有相同曲率。密切圆是局部近似，不要求它与远处曲线不相交；
若某点恰为曲率为零的拐点，则有限密切圆退化为半径趋于无穷的极限。
\end{xesolution}

\begin{origexamplebox}{椭圆轴端点的曲率半径}\label{ch03:ex:k06-2-70}
求半长轴、半短轴分别为 $a,b$（$a>b$）的椭圆在长轴端点和短轴端点处的曲率半径。
\par 题图见图~\ref{ch03:ref:question-figure-2-70}。
\KQuestionCS
\end{origexamplebox}
\Needspace{6\baselineskip}\begin{xesolution}
以下推导所用的解图为图~\ref{ch03:fig:k06-ex-2-70}。

取参数方程 $x=a\cos t$、$y=b\sin t$。参数曲线的曲率半径为
\[
\rho=\frac{(\dot x^2+\dot y^2)^{3/2}}{|\dot x\ddot y-\dot y\ddot x|}
=\frac{(a^2\sin^2t+b^2\cos^2t)^{3/2}}{ab}.
\]
\EllipseCurvatureDiagram

长轴端点对应 $t=0,\pi$，短轴端点对应 $t=\pi/2,3\pi/2$，所以
\[
 \boxed{\rho_{\text{长轴端}}=\frac{b^2}{a}},\qquad
 \boxed{\rho_{\text{短轴端}}=\frac{a^2}{b}}.
\]
这里 $a,b$ 指半轴长度；若给的是全轴长，须先各除以二。
\end{xesolution}

\begin{origexamplebox}{加速度与速度成正比的减速}\label{ch03:ex:k06-2-71}
质点沿直线运动，取初速度方向为正向，$t=0$ 时速度为 $v_0>0$，此后加速度满足 $a=-kv$，$k>0$。求完全停止所需时间和总路程。

\begin{qualityfigure}
\begin{tikzpicture}[>=Latex,font=\small]
\draw[->] (0,0)--(6.2,0) node[right]{$x$};\fill (1,0) circle(2pt) node[below]{$P$};
\draw[->,DeepBlue,thick] (1,.5)--(2.6,.5) node[right]{$v>0$};
\draw[->,Gold,thick] (2.6,1.3)--(1,1.3) node[midway,above]{$a<0$};
\end{tikzpicture}
\figcap{一次阻力下的速度方向和衰减曲线}\label{mech:03:71:q}
\end{qualityfigure}
\end{origexamplebox}
\Needspace{6\baselineskip}\begin{xesolution}
\begin{qualityfigure}
\begin{tikzpicture}[>=Latex,font=\small]
\draw[->] (0,0)--(6.2,0) node[right]{$x$};\fill (1,0) circle(2pt) node[below]{$P$};
\draw[->,DeepBlue,thick] (1,.5)--(2.6,.5) node[right]{$v>0$};
\draw[->,Gold,thick] (2.6,1.3)--(1,1.3) node[midway,above]{$a<0$};\begin{scope}[xshift=3.8cm,yshift=.8cm,x=.55cm,y=.9cm]\draw[->] (0,0)--(4,0) node[right]{$t$};\draw[->] (0,0)--(0,1.5) node[above]{$v/v_0$};\draw[DeepBlue,thick,domain=0:3.5,samples=50] plot(\x,{exp(-\x)});\end{scope}
\end{tikzpicture}
\figcap{一次阻力下的速度方向和衰减曲线（解答用辅助图）}\label{mech:03:71:a}
\end{qualityfigure}
图~\ref{mech:03:71:a}中，正向速度与负向加速度反向，右侧衰减曲线在任意有限时间都位于时间轴上方；总路程要由曲线下的无穷区间面积判断。

负号表示加速度与正向运动的速度相反，速率不断减小，而不是说质点已反向运动。
由加速度定义，
\[
 \frac{\dd v}{\dd t}=-kv,\qquad \frac{\dd v}{v}=-k\,\dd t.
\]
在速度仍为正的阶段，从初态 $(0,v_0)$ 积分到 $(t,v)$：
\[
 \int_{v_0}^{v}\frac{\dd u}{u}=-k\int_0^t\dd\tau,
 \qquad \ln\frac{v}{v_0}=-kt.
\]
两边取指数，得到
\[
 \boxed{v(t)=v_0e^{-kt}}.
\]
任意有限 $t$ 时 $v>0$，所以数学模型中完全停止所需时间为
\[
 \boxed{t_{\rm stop}=\infty}.
\]
该解在任意有限时刻均为正，说明推导中使用正速度自洽。
取初始位置为零，先求有限时间内的路程：
\[
 s(t)=\int_0^t v_0e^{-k\tau}\dd\tau
 =\left[-\frac{v_0}{k}e^{-k\tau}\right]_0^t
 =\frac{v_0}{k}(1-e^{-kt}).
\]
再令 $t\to\infty$，得到总路程收敛：
\[
 \boxed{s_\infty=\int_0^\infty v_0e^{-kt}\,\dd t=\frac{v_0}{k}}.
\]
这正是“速度永不严格为零，但总路程有限”的极限现象。
$[k]=\mathrm{s^{-1}}$，所以 $kt$ 无量纲，$v_0/k$ 是长度；当 $kt\ll1$，
$e^{-kt}\simeq1-kt$，有 $s(t)\simeq v_0t$，恢复短时间内近似匀速的结果。
若用一个小阈值 $v_{\rm tol}$ 表示实际可视为停止，则所需时间为
$t_{\rm tol}=k^{-1}\ln(v_0/v_{\rm tol})$，是有限的。
\end{xesolution}

\begin{origexamplebox}{加速度与速度平方成正比的减速}\label{ch03:ex:k06-2-72}
质点沿直线运动，取初速度方向为正向，$t=0$ 时速度为 $v_0>0$，此后加速度满足 $a=-\beta v^2$，$\beta>0$。求速度随时间的关系。

\begin{qualityfigure}
\begin{tikzpicture}[>=Latex,font=\small]
\draw[->] (0,0)--(6.2,0) node[right]{$x$};\fill (1,0) circle(2pt) node[below]{$P$};
\draw[->,DeepBlue,thick] (1,.5)--(2.6,.5) node[right]{$v>0$};
\draw[->,Gold,thick] (2.6,1.3)--(1,1.3) node[midway,above]{$a<0$};
\end{tikzpicture}
\figcap{平方阻力下的正速度与长时尾部}\label{mech:03:72:q}
\end{qualityfigure}
\end{origexamplebox}
\Needspace{6\baselineskip}\begin{xesolution}
\begin{qualityfigure}
\begin{tikzpicture}[>=Latex,font=\small]
\draw[->] (0,0)--(6.2,0) node[right]{$x$};\fill (1,0) circle(2pt) node[below]{$P$};
\draw[->,DeepBlue,thick] (1,.5)--(2.6,.5) node[right]{$v>0$};
\draw[->,Gold,thick] (2.6,1.3)--(1,1.3) node[midway,above]{$a<0$};\begin{scope}[xshift=3.8cm,yshift=.8cm,x=.55cm,y=.9cm]\draw[->] (0,0)--(4,0) node[right]{$t$};\draw[->] (0,0)--(0,1.5) node[above]{$v/v_0$};\draw[DeepBlue,thick,domain=0:3.5,samples=50] plot(\x,{1/(1+\x)});\end{scope}
\end{tikzpicture}
\figcap{平方阻力下的正速度与长时尾部（解答用辅助图）}\label{mech:03:72:a}
\end{qualityfigure}
图~\ref{mech:03:72:a}中，速度始终沿正向，加速度反向；右侧曲线的 $1/t$ 尾部与指数尾部不同，不能仅凭速度趋零便断言总路程有限。

分离变量：
\[
\frac{\dd v}{v^2}=-\beta\,\dd t.
\]
按初始条件取积分上下限，
\[
 \int_{v_0}^{v}\frac{\dd u}{u^2}=-\beta\int_0^t\dd\tau.
\]
由于 $\int u^{-2}\dd u=-u^{-1}$，得到
\[
-\frac1v+\frac1{v_0}=-\beta t,
\]
故
\[
 \boxed{v(t)=\frac{v_0}{1+\beta v_0t}}.
\]
代回可检验 $v(0)=v_0$ 及 $\dot v=-\beta v^2$，且始终 $v>0$，没有发生反向。
$[\beta]=\mathrm{m^{-1}}$，故 $\beta v_0t$ 无量纲。
速度同样只在 $t\to\infty$ 时趋于零；与指数衰减不同，大时间行为为 $v\sim1/(\beta t)$。
积分给出 $s(t)=\beta^{-1}\ln(1+\beta v_0t)$，其极限发散，说明“速度趋于零”并不保证总路程有限。
\end{xesolution}
}

\endgroup

% MechanicsNumberedSolutionFigures20261023-v1
% MechanicsQualityRevision20261023-v1
% QualityRevision20261023
\begingroup
\tcbset{breakable/.default=false}
\chapter{相对运动、约束和刚体运动学}

{
\graphicspath{{../../../sample/xes-original/kinematics-supplement/}{../../../sample/xes-original/kinematics-supplement/assets/}}
\section{相对运动与参考系}
\newcommand{\CrossRiverRoundTripDiagram}{%
\begin{center}\begin{tikzpicture}[scale=.78,>=Latex]
 \fill[Blue!7] (-3,-1.45) rectangle (3,1.45);
 \draw[very thick] (-3,-1.45)--(3,-1.45) (-3,1.45)--(3,1.45);
 \coordinate (A) at (0,-1.45);\coordinate (B) at (0,1.45);
 \node[point,label={[clearlabel]below:$A$}] at (A) {};
 \node[point,label={[clearlabel]above:$B$}] at (B) {};
 \draw[DeepBlue,very thick,->] (.10,-1.25)--(.10,1.20);
 \draw[DeepBlue,very thick,->] (-.10,1.20)--(-.10,-1.25);
 \draw[vector] (-2.35,.15)--(-1.05,.15) node[midway,clearlabel,above]{$\vect u$};
 \draw[vector] (A)--++(145:1.25) node[near end,clearlabel,below left]{$\vect v$};
 \draw[vector] (B)--++(215:1.25) node[near end,clearlabel,above left]{$\vect v$};
 \draw[<->,TextGray] (2.45,-1.45)--(2.45,1.45) node[midway,clearlabel,right]{$L$};
\end{tikzpicture}\figcap{解图：$A,B$ 位于两岸正对位置；往返时都用船对水速度的上游分量抵消水流。}\label{ch04:fig:k07-ex-2-73}\end{center}}

\newcommand{\EqualSpeedFrameDiagram}{%
\begin{center}\begin{tikzpicture}[scale=.82,>=Latex]
 \draw[very thick] (-3,0)--(3,0);
 \coordinate (A) at (-2.1,0);\coordinate (B) at (2.1,0);
 \node[point,label={[clearlabel]below:$A$}] at (A) {};
 \node[point,label={[clearlabel]below:$B$}] at (B) {};
 \draw[vector] (A)--++(1.55,0) node[midway,clearlabel,above]{$v$};
 \draw[vector] (-.75,1.0)--(.75,1.0) node[midway,clearlabel,above]{$S'\text{ 的速度 }u$};
 \draw[->,TextGray] (-1.55,-.58)--(-.2,-.58);
 \draw[->,TextGray] (1.55,-.58)--(.2,-.58);
 \node[below] at (0,-.65){在 $S'$ 中迎面接近};
\end{tikzpicture}\figcap{解图：在地面系中 $A$ 向右、$B$ 静止，参考系 $S'$ 也向右平动。}\label{ch04:fig:k07-ex-2-74}\end{center}}

\newcommand{\VelocityTriangleDiagram}{%
\begin{center}\begin{tikzpicture}[scale=.8,>=Latex]
 \coordinate (O) at (0,0);\coordinate (A) at (2.7,0);\coordinate (B) at (2.7,1.8);
 \draw[vector] (O)--(A) node[midway,clearlabel,below]{$\vect v_{P/A}$};
 \draw[vector] (A)--(B) node[midway,clearlabel,right]{$\vect v_{A/O}$};
 \draw[vector] (O)--(B) node[midway,clearlabel,above left]{$\vect v_{P/O}$};
\end{tikzpicture}\figcap{相对速度三角形：先走“物体相对载体”，再走“载体相对地面”，得到物体相对地面。}\end{center}}

\newcommand{\KinematicKeepSpace}[1]{%
\par\begingroup
\dimen0=#1\relax
\dimen2=\lastskip
\advance\dimen0 by \dimen2
\vskip-\dimen2
\vskip0pt plus\dimen0
\penalty-100
\vskip0pt plus-\dimen0
\vskip\dimen0
\penalty9999
\vskip-\dimen0
\vskip\dimen2
\endgroup}

\newcommand{\ReferenceFramesDiagram}{%
\begin{center}\begin{minipage}{\linewidth}\centering\begin{tikzpicture}[>=Latex,scale=.88]
\draw[kaxis] (-.3,0)--(7,0) node[right]{$S:x$};
\draw[kaxis] (2,1.6)--(7,1.6) node[right]{$S':x'$};
\coordinate (P) at (5.7,2.4);
\draw[kguide] (2,-.6)--(2,1.6) (P)--(5.7,-.6);
\fill (0,0) circle(1.5pt) (2,1.6) circle(1.5pt) (P) circle(1.5pt);
\node[below left] at (0,0){$O$};\node[above left] at (2,1.6){$O'$};
\node[above] at (P){$P$};
\draw[<->] (0,-.65)--(2,-.65) node[midway,below]{$x_{O'}$};
\draw[<->,DeepBlue] (0,-1.25)--(5.7,-1.25) node[midway,below]{$x$};
\draw[<->,DeepBlue] (2,.8)--(5.7,.8) node[midway,above]{$x'$};
\draw[vector] (2.6,2.35)--(3.9,2.35) node[midway,above]{$u$};
\draw (5.7,0)--(5.7,.12) (5.7,1.6)--(5.7,1.72);
\end{tikzpicture}\figcap{同一点的两个横向投影坐标满足 $x=x_{O'}+x'$；轴的上下错开不表示额外横向距离。}\label{ch04:fig:k07-2-19}\end{minipage}\end{center}}

\newcommand{\AircraftWindDiagram}{%
\begin{center}\begin{minipage}{\linewidth}\centering\begin{tikzpicture}[>=Latex,scale=.88]
\begin{scope}
\draw[kaxis] (0,0)--(4,0) node[right]{东};\draw[kaxis] (0,0)--(0,2.4) node[above]{北};
\draw[vector] (0,0)--(3,0) node[midway,below]{$\vect V$};
\draw[vector] (3,0)--(3,1.5) node[midway,right]{$\vect u$};
\draw[vector] (0,0)--(3,1.5) node[midway,above left]{$\vect v_{\text{地}}$};
\node at (1.7,-.85){(a) 机头正东};
\end{scope}
\begin{scope}[xshift=7cm]
\draw[kaxis] (0,0)--(3.8,0) node[right]{东};\draw[kaxis] (0,0)--(0,2.4) node[above]{北};
\draw[vector] (0,0)--(2.598,-1.5) node[midway,below left]{$\vect V$};
\draw[vector] (2.598,-1.5)--(2.598,0) node[midway,right]{$\vect u$};
\draw[vector] (0,0)--(2.598,0) node[midway,above]{$\vect v_{\text{地}}$};
\draw (.8,0) arc(0:-30:.8);\node at (1.15,-.3){$\theta$};
\node at (1.5,-2.1){(b) 航迹正东};
\end{scope}
\end{tikzpicture}\figcap{解图：风速向北；两图均按 $\vect v_{\text{地}}=\vect V+\vect u$ 首尾相接。}\label{ch04:fig:k07-ex-2-81}\end{minipage}\end{center}}

\newcommand{\RiverCrossingDiagram}{%
\begin{center}\begin{minipage}{\linewidth}\centering\begin{tikzpicture}[>=Latex,scale=.9]
\begin{scope}
\fill[Blue!5] (-.5,0) rectangle (3.9,3);
\draw[thick] (-.5,0)--(3.9,0) (-.5,3)--(3.9,3);
\coordinate (A) at (0,0);\draw[kguide] (A)--(2.25,3);
\fill (A) circle(1.5pt);\node[below] at (A){$A$};
\draw[vector] (A)--(0,1.6) node[midway,left]{$\vect v_1$};
\draw[vector] (0,1.6)--(1.2,1.6) node[midway,above]{$\vect v_2$};
\draw[vector] (A)--(1.2,1.6);
\node[right] at (1.55,.65){$\vect v_{\text{岸}}$};\draw[thin] (1.55,.78)--(.82,1.02);
\draw[<->] (3.5,0)--(3.5,3) node[midway,left]{$d$};
\node at (1.7,-.8){(a) 最短时间的实际航迹};
\end{scope}
\begin{scope}[xshift=7cm,yshift=.35cm]
\coordinate (W) at (2.5,0);\coordinate (T) at (1.6,1.2);
\draw[kaxis] (-.3,0)--(4.4,0) node[right]{$v_x$};
\draw[kaxis] (0,-1.75)--(0,2.3) node[above]{$v_y$};
\draw (W) circle(1.5);\draw[kguide] (W)--(T);
\draw[vector] (0,0)--(T);\fill (W) circle(1.5pt) (T) circle(1.5pt);
\node[below left] at (0,0){$O$};\node[below] at (W){$W$};\node[above] at (T){$T$};
\node[above right] at (2.2,.65){$v_1$};
\draw[<->] (0,-1.9)--(2.5,-1.9) node[midway,below]{$v_2$};
\draw (.4,0) arc(0:36.87:.4);\node at (.67,.2){$\psi$};
\node at (1.8,-2.65){(b) $v_1<v_2$ 的速度圆};
\end{scope}
\end{tikzpicture}\figcap{解图：左图出发点在岸上，航迹平行于合速度；右图 $WT\perp OT$，切线给出最小位移方向。}\label{ch04:fig:k07-ex-2-82}\end{minipage}\end{center}}

\newcommand{\AircraftEqualSpeedDiagram}{%
\begin{center}\begin{minipage}{\linewidth}\centering\begin{tikzpicture}[>=Latex,scale=1]
\foreach \s/\shift/\txt in {1/0/东,-1/6/西}{
\begin{scope}[xshift=\shift cm]
\coordinate (O) at (0,0);\coordinate (G) at (0,3.6);
\coordinate (A) at ({\s*1.74284},3.15);\coordinate (F) at ($(A)!.5!(G)$);
\draw[vector] (O)--(A);\draw[vector] (A)--(G);\draw[vector] (O)--(G);
\draw[kguide] (O)--(F);
\node[below] at (O){$O$};\node[above] at (G){$G$};
\ifnum\s>0 \node[right] at (A){$A$};\node[left] at (0,1.45){$60$};
\else \node[left] at (A){$A$};\node[right] at (0,1.45){$60$};\fi
\node[above] at (F){$F$};\node at ({\s*1.35},1.6){$60$};
\draw (90:1.4) arc[start angle=90,end angle={90-\s*28.955},radius=1.4];
\node at ({\s*.61},1.77){$\theta$};
\node at ({\s*.9},3.9){$30$};
\node at (0,-.65){机头北偏\txt};
\end{scope}}
\node at (3,2.1){正北 $\uparrow$};
\end{tikzpicture}\figcap{解图：$OA$ 为空速，$AG$ 为风速，$OG$ 为地速；$F$ 为风速边中点，速度单位均为 $\mathrm{m/s}$。}\label{ch04:fig:k07-ex-2-83}\end{minipage}\end{center}}

\newcommand{\SmokeWindDiagram}{%
\begin{center}\begin{minipage}{\linewidth}\centering\begin{tikzpicture}[>=Latex,scale=.95]
\foreach \vv/\shift/\txt in {1/0/第一次,3/6/第二次}{
\begin{scope}[xshift=\shift cm]
\draw[kaxis] (-.3,0)--(3.3,0) node[right]{东};\draw[kaxis] (0,-1.6)--(0,2.6) node[above]{北};
\draw[vector] (0,0)--(2,-1) node[midway,below left]{$\vect w$};
\draw[vector] (2,-1)--(2,{\vv-1});
\node[right] at (2.1,{.5*\vv-1}){$-\vect v_{\text{船}}$};
\draw[vector] (0,0)--(2,{\vv-1});
\node[above left] at (1,{.5*(\vv-1)+.2}){$\vect v_{\text{烟/船}}$};
\node at (1.4,-2.1){\txt};
\end{scope}}
\end{tikzpicture}\figcap{解图：同一风速加上两种反向船速，分别得到向东、向东北的烟对船速度。}\label{ch04:fig:k07-ex-2-84}\end{minipage}\end{center}}

\newcommand{\InterceptTriangleDiagram}{%
\begin{center}\begin{minipage}{\linewidth}\centering\begin{tikzpicture}[>=Latex,x=.045cm,y=.045cm]
\coordinate (A) at (0,0);\coordinate (B) at (-93.9693,34.2020);
\coordinate (C1) at (-17.3065,34.2020);\coordinate (C2) at (79.9527,34.2020);
\draw[thick] (A)--(B);\draw[vector] (B)--(C2);
\draw[vector] (A)--(C1);\draw[vector] (A)--(C2);
\foreach \p/\lab/\pos in {A/A/below,B/B/above left,C1/C_1/above,C2/C_2/above right}{\fill (\p) circle(1.6pt);\node[\pos] at (\p){$\lab$};}
\node[left] at (-50,12){$L$};\node[above] at (30,38){敌船沿 $+x$};
\node[left] at (-8,20){$v_2t_1$};\node[below right] at (42,17){$v_2t_2$};
\draw (B)++(0:22) arc(0:-20:22);\node at (-64,29){$\alpha$};
\draw (23.16018:14) arc(23.16018:160:14);\node[below] at (0,-12){$\beta_1=\angle BAC_1,\quad\beta_2=\angle BAC_2$};
\end{tikzpicture}\figcap{解图：按文中数值例得到的两个未来交点；发射方向由各自的 $AC_i$ 决定，而不只由正弦值决定。}\label{ch04:fig:k07-ex-2-85}\end{minipage}\end{center}}

\newcommand{\BombInterceptDiagram}{%
\begin{center}\begin{minipage}{\linewidth}\centering\begin{tikzpicture}[>=Latex,scale=.92]
\draw[kaxis] (-5,0)--(3.8,0) node[right]{$x$};
\coordinate (B) at (-4,3);\coordinate (Q) at (-2,2.25);
\draw[kcurve,domain=0:1,samples=60,variable=\t] plot({-4+4*\t},{3*(1-\t*\t)});
\draw[kguide] (B)--(-4,0);\draw[vector] (B)--(-2.8,3) node[midway,above]{$v_1$};
\draw[thick,dashed] (-2,0)--(Q) (2,0)--(Q);
\foreach \x/\lab/\pos in {0/0/below,-2/-D/below,2/+D/below}{\fill (\x,0) circle(1.6pt);\node[\pos] at (\x,0){$\lab$};}
\node[above left] at (B){$(-L,H)$};\fill (Q) circle(1.6pt);\node[above right] at (Q){$Q$};
\draw[<->] (-4.65,0)--(-4.65,3) node[midway,left]{$H$};
\draw[<->] (-4,-.8)--(0,-.8) node[midway,below]{$L$};
\node at (1.6,3){两炮位分别讨论};
\end{tikzpicture}\figcap{解图：蓝线是炸弹轨迹；两条虚线仅连接炮位与相遇点，不是炮弹的抛物线轨迹。}\label{ch04:fig:k07-ex-2-86}\end{minipage}\end{center}}

\newcommand{\ThreeBoatsDiagram}{%
\begin{center}\begin{minipage}{\linewidth}\centering\begin{tikzpicture}[>=Latex,scale=.95]
\begin{scope}
\draw[kguide] (0,0)--(4.8,0);
\foreach \x/\lab in {0/A,1.8/B,3.6/C}{\fill (\x,0) circle(1.5pt);\node[below] at (\x,0){$\lab$};}
\draw[vector] (0,0)--(.6,0);\node[above] at (.3,0){$u$};
\draw[vector] (1.8,0)--(1.8,1.2) node[above]{$2u$};
\draw[vector] (3.6,0)--(4.8,1.2) node[above]{$2\sqrt2u$};
\draw[<->] (0,-.65)--(1.8,-.65) node[midway,below]{$a$};
\draw[<->] (1.8,-.65)--(3.6,-.65) node[midway,below]{$a$};
\node at (2,2.1){(a) $t=0$，速度同比例};
\end{scope}
\begin{scope}[xshift=7cm]
\coordinate (M) at (2.7,2.2);\coordinate (A) at (0,0);
\draw[kguide] (A)--(2.7,0)--(M);\draw[thick,->] (M)--(A);
\fill (M) circle(1.5pt) (A) circle(1.5pt);
\node[above] at (M){$M(t)$};\node[left] at (A){$A(t)$};
\node[below] at (1.35,0){$3a/2$};\node[right] at (2.7,1){$2ut$};
\draw[vector] (A)--(0,-1.2) node[below]{$\vect v_{A/M}=(0,-2u)$};
\node[above left] at (1.5,1.3){$\vect r_{A/M}$};
\node at (1.5,3.05){(b) 随 $M$ 平移的坐标};
\end{scope}
\end{tikzpicture}\figcap{解图：$M$ 与 $A$ 的横向速度相同，$MA$ 的水平分量固定、竖直分量随时间增加。}\label{ch04:fig:k07-ex-2-87}\end{minipage}\end{center}}

\newcommand{\PursuitDiagram}{%
\begin{center}\begin{minipage}{\linewidth}\centering\begin{tikzpicture}[>=Latex,scale=1.25]
\draw[thick] (-.3,0)--(3.5,0);
\draw[kguide] (0,3)--(3.1,3);
\draw[kcurve,domain=1:0,samples=100,variable=\u] plot({3*(2/3-\u+\u^3/3)},{3*(1-\u^2)});
\coordinate (D) at (.625,2.25);\coordinate (P) at (1.1875,3);
\draw[thick] (D)--(P);\draw[kguide] (D)--++(1.3,0);
\draw[vector] (D)--++(.27,.36);\draw[vector] (P)--++(.5,0) node[midway,above]{$v_0$};
\fill (D) circle(1.5pt) (P) circle(1.5pt);
\node[left] at (D){$D$};\node[above left] at (P){$P$};
\node[left] at (.86,2.7){$r$};
\draw (D)++(.38,0) arc(0:53.1301:.38);\node at (1.18,2.52){$\theta$};
\draw[vector] (0,0)--(0,.7) node[left]{$v_1$};
\draw[vector] (0,0)--(.65,0) node[below]{$a_0$};
\draw[<->] (-.65,0)--(-.65,3) node[midway,left]{$H$};
\node[below left] at (0,0){$A$};\node[above right] at (2,3){命中};
\end{tikzpicture}\figcap{解图：取 $v_1=2v_0$ 展示追击曲线；中途导弹 $D$ 的切线沿同一时刻飞机 $P$ 的视线，初始竖直、末端水平。}\label{ch04:fig:k07-ex-2-88}\end{minipage}\end{center}}

\newcommand{\KQuestionCV}{%
%
\begin{center}\begin{tikzpicture}[scale=.78,>=Latex]
 \fill[Blue!7] (-3,-1.45) rectangle (3,1.45);
 \draw[very thick] (-3,-1.45)--(3,-1.45) (-3,1.45)--(3,1.45);
 \coordinate (A) at (0,-1.45);\coordinate (B) at (0,1.45);
 \node[point,label={[clearlabel]below:$A$}] at (A) {};
 \node[point,label={[clearlabel]above:$B$}] at (B) {};
 \draw[vector] (-2.35,.15)--(-1.05,.15) node[midway,clearlabel,above]{$\vect u$};
 \draw[vector] (A)--++(145:1.25) node[near end,clearlabel,below left]{$\vect v$};
 \draw[<->,TextGray] (2.45,-1.45)--(2.45,1.45) node[midway,clearlabel,right]{$L$};
\end{tikzpicture}\figcap{题图：两岸正对的 A、B 距离为 $L$，水流速度向右，船速指船相对水的速度。}\label{ch04:ref:question-figure-2-73}\end{center}
}

\newcommand{\KQuestionCW}{%
%
\begin{center}\begin{tikzpicture}[scale=.82,>=Latex]
 \draw[very thick] (-3,0)--(3,0);
 \coordinate (A) at (-2.1,0);\coordinate (B) at (2.1,0);
 \node[point,label={[clearlabel]below:$A$}] at (A) {};
 \node[point,label={[clearlabel]below:$B$}] at (B) {};
 \draw[vector] (A)--++(1.55,0) node[midway,clearlabel,above]{$v$};
 \draw[vector] (-.75,1.0)--(.75,1.0) node[midway,clearlabel,above]{$S'\text{ 的速度 }u$};
\end{tikzpicture}\figcap{题图：地面系中 A 向右运动、B 静止；参考系 $S\prime$ 以速度 $u$ 向右平动。}\label{ch04:ref:question-figure-2-74}\end{center}
}

\newcommand{\KQuestionDE}{%
\begin{center}\begin{minipage}{.97\linewidth}\centering
\begin{tikzpicture}[>=Latex,scale=.9]

\draw[thick] (-2,0)--(3,0) (-2,2.7)--(3,2.7);
\node[point,label=below:$A$] at (0,0){};
\draw[vector] (-1.7,1.5)--(-.3,1.5) node[midway,above]{$v_2$};
\draw[vector] (0,0)--(110:1.4) node[pos=.5,right=7pt]{$v_1$};
\draw[<->] (2.4,0)--(2.4,2.7) node[midway,left]{$d$};

\end{tikzpicture}
\figcap{题图：水流平行于河岸，小船对水速度的大小为 $v_1$，船头方向可以选择。}\label{ch04:ref:question-figure-2-82}
\end{minipage}\end{center}
}

\newcommand{\KQuestionDH}{%
\begin{center}\begin{minipage}{.97\linewidth}\centering
\begin{tikzpicture}[>=Latex,scale=.9]

\coordinate (A) at (4,0);\coordinate (B) at (-2,2.4);\coordinate (C) at (2.4,2.4);
\draw[thick] (A)--(B) node[midway,below left]{$L$};
\draw[kguide] (B)--(5.1,2.4) (A)--(4,3.2) (A)--(C);
\draw[vector] (B)--(-.7,2.4) node[midway,above]{$v_1$};
\draw[vector] (A)--($(A)!.43!(C)$) node[left]{$v_2$};
\draw (B)++(-21.801:.85) arc(-21.801:0:.85);\node at (-.7,2.1){$\alpha$};
\draw (A)++(123.690:.65) arc(123.690:158.199:.65);\node at (3.23,.69){$\beta$};
\foreach \p/\pos in {A/below right,B/above left,C/above}\node[point,label=\pos:$\p$] at (\p){};

\end{tikzpicture}
\figcap{题图：A 为潜艇、B 为敌船初始位置，C 表示可能的相遇点；发射方向待定。}\label{ch04:ref:question-figure-2-85}
\end{minipage}\end{center}
}

\newcommand{\KQuestionDJ}{%
\begin{center}\begin{minipage}{.97\linewidth}\centering
\begin{tikzpicture}[>=Latex,scale=.9]

\draw[kguide] (0,0)--(4.8,0);
\foreach \x/\lab in {0/A,1.8/B,3.6/C}{\fill (\x,0) circle(1.5pt);\node[below] at (\x,0){$\lab$};}
\draw[vector] (0,0)--(.6,0);\node[above] at (.3,0){$u$};
\draw[vector] (1.8,0)--(1.8,1.2) node[above]{$2u$};
\draw[vector] (3.6,0)--(4.8,1.2) node[above]{$2\sqrt2u$};
\draw[<->] (0,-.65)--(1.8,-.65) node[midway,below]{$a$};
\draw[<->] (1.8,-.65)--(3.6,-.65) node[midway,below]{$a$};
\end{tikzpicture}
\figcap{题图：三船初始位置及对地速度，图中 A、B、C 等间距。}\label{ch04:ref:question-figure-2-87}
\end{minipage}\end{center}
}


地球自转并绕太阳公转，太阳也在银河系中运动。不存在可供一切运动绝对参照的静止物体，因此观察者可以选择合适的物体作为参考系。在参考系中取原点和一组方向建立坐标系；同一运动在不同参考系中的坐标、速度乃至加速度表达可以不同，但它们之间有确定的变换关系。

参考系与坐标系相关却并不等同。选择地面、车厢或船舱作为“静止的物体集合”，是在选择参考系；随后在其中指定原点、坐标轴和单位，才建立了坐标系。不同观察者都可以把自己所在的参考物体看作静止，并据此描述其他物体。运动学不要求寻找一个无法实现的“绝对静止”背景，而是要求把参考对象说清，并能在不同描述之间进行转换。

这种相对性与日常经验完全一致：站在站台上看列车前进，坐在列车里看站台后退；两种叙述的速度数值不同，却对应同一组相遇事件。只要两个参考系之间的相对位置随时间已知，它们对任意第三个物体的坐标、速度和加速度描述就能互相换算，变换关系本身与第三个物体怎样运动无关。

\subsection*{平动参考系的坐标变换}

\ReferenceFramesDiagram

若两个一维坐标系轴向相同、原点相距 $x_0$，则
\begin{equation}
x'=x-x_0.
\end{equation}
若 $S'$ 的原点相对 $S$ 沿正方向以常速 $u$ 运动，且 $t=0$ 两原点重合，则经典力学中取共同时间 $t$，有
\begin{equation}
x'=x-ut,\qquad v'=v-u,\qquad a'=a.
\end{equation}
共同时间只适用于远低于光速的经典极限；接近光速时应使用相对论的时空变换。

图~\ref{ch04:fig:k07-2-19}中两条横轴上下错开是为了分别标出读数；比较的是同一点 $P$ 在平行轴上的投影坐标，不是斜线 $OP$、$O'P$ 的长度。图中两个投影落点在同一竖直线上，横向有向距离满足 $x=x_{O'}+x'$，即 $x'=x-x_{O'}$。当两个原点相对静止时，$x_{O'}=x_0$ 为常数，对时间求导后速度、加速度相同；当 $x_{O'}=ut$ 时，速度相差 $u$，而因 $u$ 恒定，加速度仍相同。

经典力学把不同参考系的时钟视为同步，认为同一事件在各系中的时间坐标相同。这在普通车辆、飞机和抛体的速度范围内精度足够高，并构成伽利略变换的基础。接近光速时，这一共同时间假设不再成立；此处指出适用范围，是为了避免把经典公式错误外推到相对论情形。

\begin{origexamplebox}{横渡流动河水往返}\label{ch04:ex:k07-2-73}
河宽为 $L$，两岸正对位置分别为 $A,B$。船在静水中的速率为 $v$，河水沿河岸以速率 $u$ 匀速流动，且 $v>u$。船从 $A$ 出发驶向 $B$，再沿同一直线由 $B$ 返回 $A$。求往返时间。若要使河中往返时间等于船以速率 $v$ 在静水中按同一路线往返的时间，河中船对水速率应提高到多少？
\par 题图见图~\ref{ch04:ref:question-figure-2-73}。
\KQuestionCV
\end{origexamplebox}
\Needspace{6\baselineskip}\begin{xesolution}
以下推导所用的解图为图~\ref{ch04:fig:k07-ex-2-73}。

要使船对岸速度沿 $AB$，船头在两个航程中都必须偏向上游，使船对水速度的沿河分量恰好抵消水流：
\CrossRiverRoundTripDiagram
\[
v_{\parallel}=u.
\]
船对水速度大小为 $v$，所以横河分量，也就是船对岸沿 $AB$ 的速率，为
\[
v_{\perp}=\sqrt{v^2-u^2}.
\]
往、返两程的横河速率相同，故
\[
\boxed{T=\frac{2L}{\sqrt{v^2-u^2}}}.
\]
设提高后的对水速率为 $V$。要求
\[
\frac{2L}{\sqrt{V^2-u^2}}=\frac{2L}{v},
\]
因此
\[
\boxed{V=\sqrt{v^2+u^2}}.
\]
当 $u\to0$ 时，$T\to2L/v$、$V\to v$，退化为静水结果；同时 $v>u$ 保证横渡可以完成。
\end{xesolution}


\begin{origexamplebox}{选择迎面等速的参考系}\label{ch04:ex:k07-2-74}
在地面系 $S$ 中，质点 $A$ 以速率 $v$ 撞向静止质点 $B$。参考系 $S'$ 沿 $A$ 的运动方向以速率 $u$ 匀速运动。$u$ 取何值时，$S'$ 看见两质点以相同速率迎面碰撞？
\par 题图见图~\ref{ch04:ref:question-figure-2-74}。
\KQuestionCW
\end{origexamplebox}
\Needspace{6\baselineskip}\begin{xesolution}
以下推导所用的解图为图~\ref{ch04:fig:k07-ex-2-74}。

地面系中的位置和速度方向如图。变换到向右运动的 $S'$ 后，$A$ 的速度为 $v-u$，$B$ 的速度为 $-u$：
\EqualSpeedFrameDiagram
迎面且等速要求 $v-u=u$，所以
\[
\boxed{u=\frac v2}.
\]
这正是两质点速度的算术平均值。
\end{xesolution}

\begin{origexamplebox}{匀速平动系中的两质点速度}\label{ch04:ex:k07-2-75}
在地面系 $S$ 中，$A$ 以速率 $v$ 沿正方向撞向静止的 $B$。求相对 $S$ 沿正方向以速率 $u$ 匀速运动的参考系中两质点的速度。

\begin{qualityfigure}
\begin{tikzpicture}[>=Latex,font=\small]
\draw[->] (-.5,0)--(6.5,0) node[right]{$x$};\node[below left] at (0,0){$O$};
\fill (1,0) circle(2pt) node[below]{$A$};\fill (4,0) circle(2pt) node[below]{$B$};
\draw[->,DeepBlue,thick] (1,.4)--(2.3,.4) node[above]{$v$};
\node at (4,.7){$v_B=0$};\draw[->,Gold] (.4,-1)--(2,-1) node[right]{$S':u$};
\end{tikzpicture}
\figcap{地面系与向正方向平移的参考系}\label{mech:04:75:q}
\end{qualityfigure}
\end{origexamplebox}
\Needspace{6\baselineskip}\begin{xesolution}
\begin{qualityfigure}
\begin{tikzpicture}[>=Latex,font=\small]
\draw[->] (-.5,0)--(6.5,0) node[right]{$x$};\node[below left] at (0,0){$O$};
\fill (1,0) circle(2pt) node[below]{$A$};\fill (4,0) circle(2pt) node[below]{$B$};
\draw[->,DeepBlue,thick] (1,.4)--(2.3,.4) node[above]{$v$};
\node at (4,.7){$v_B=0$};\draw[->,Gold] (.4,-1)--(2,-1) node[right]{$S':u$};\node[align=left] at (4,-1.4){每个速度均减去 $u$\\相对速度仍为 $v$};
\end{tikzpicture}
\figcap{地面系与向正方向平移的参考系（解答用辅助图）}\label{mech:04:75:a}
\end{qualityfigure}
图~\ref{mech:04:75:a}中，参考系 $S'$ 相对地面沿正向运动，必须从两个地面速度中扣除同一个 $u$。两点速度之差在这种平移中保持不变。

直接用伽利略速度变换：
\[
\boxed{v_A'=v-u},\qquad \boxed{v_B'=-u}.
\]
若 $u>v$，两者在 $S'$ 中都沿负方向运动，但 $B$ 的负向速度更大，二者仍相互接近。
\end{xesolution}

一般平动参考系的原点坐标为 $x_0(t)$ 时，
\begin{equation}
x'=x-x_0(t),\qquad v'=v-v_0(t),\qquad a'=a-a_0(t).
\end{equation}
这些式子要求两坐标系的轴始终平行；若参考系转动，还会出现额外项。选自由落体系时，所有同受重力加速度的抛体相对它都作匀速直线运动。

一般变换可从位置关系连续求导得到：原点速度 $v_0=\dot x_0$ 从物体在 $S$ 中的速度中扣除，原点加速度 $a_0=\ddot x_0$ 同样从物体加速度中扣除。若坐标轴还在转动，仅减去原点平动已不够，因为同一个空间方向在两系中的分量也随时间改变。本节公式因此只用于轴始终平行的平动参考系。

选择参考系的价值不只是换一组数值，而是可能把复杂运动变简单。两个物体具有相同加速度时，选其中一个为平动参考系，另一物体的相对加速度为零，原本的两段匀加速运动就转化为匀速相对运动。自由落体系正是这一思想：所有忽略空气阻力的抛体都含共同的 $\vect g$，相减后重力项消失。相遇、追及和碰撞题常可用这种方法直接得到相对距离何时关闭。

不过，参考系变换不会改变事件是否相遇。若两个物体在某时刻同处一点，任何合法坐标系都应给出相同的相遇时刻和相遇事件；若不同解法给出互相矛盾的结论，通常是速度下标、正方向或初始位置转换出了问题。位置、速度的数值可以变，物理事件之间的对应关系不能变。


\begin{origexamplebox}{加速火车上的人}\label{ch04:ex:k07-2-76}
火车启动时相对地面以加速度 $a_1$ 匀加速；车上的人沿列车前进方向，以相对列车初速度 $v_2$、加速度 $a_2$ 匀加速跑动。证明人相对地面也作匀加速运动，并求加速度。

\begin{qualityfigure}
\begin{tikzpicture}[>=Latex,font=\small]
\draw[thick,fill=PaleBlue] (0,0) rectangle (6,1.8);\draw[thick] (-.4,-.4)--(6.6,-.4);
\fill (2,.8) circle(2pt) node[below]{人};\draw[->,DeepBlue] (2,1.1)--(3.6,1.1) node[above]{$v_2,a_2$（对车）};
\draw[->,Gold] (4,-.9)--(5.6,-.9) node[right]{$a_1$（对地）};\node[left] at (0,0){$O'$};\node[left] at (0,-.4){$O$};
\end{tikzpicture}
\figcap{人、车和地面三个位置的相加}\label{mech:04:76:q}
\end{qualityfigure}
\end{origexamplebox}
\Needspace{6\baselineskip}\begin{xesolution}
\begin{qualityfigure}
\begin{tikzpicture}[>=Latex,font=\small]
\draw[thick,fill=PaleBlue] (0,0) rectangle (6,1.8);\draw[thick] (-.4,-.4)--(6.6,-.4);
\fill (2,.8) circle(2pt) node[below]{人};\draw[->,DeepBlue] (2,1.1)--(3.6,1.1) node[above]{$v_2,a_2$（对车）};
\draw[->,Gold] (4,-.9)--(5.6,-.9) node[right]{$a_1$（对地）};\node[left] at (0,0){$O'$};\node[left] at (0,-.4){$O$};\draw[<->] (0,-1.6)--(2,-1.6) node[midway,below]{$x_{\text{人/车}}$};\node at (3,2.5){$x=x_{\text{车}}+x_{\text{人/车}}$};
\end{tikzpicture}
\figcap{人、车和地面三个位置的相加（解答用辅助图）}\label{mech:04:76:a}
\end{qualityfigure}
图~\ref{mech:04:76:a}中，人的坐标先相对车内原点 $O'$ 给出，再加 $O'$ 相对地面原点 $O$ 的位置。两次求导须针对同一个地面时间。

人的地面位置为 $x=x_{\text{车}}+x_{\text{人/车}}$。两次求导：
\[
v=v_{\text{车}}+v_{\text{人/车}},\qquad
a=a_1+a_2.
\]
因此
\[
\boxed{a_{\text{人/地}}=a_1+a_2},
\]
为常量，故确为匀加速运动。若人反向跑，$a_2$ 应按选定正方向带符号。
\end{xesolution}

\begin{origexamplebox}{各向抛出小球形成的球面}\label{ch04:ex:k07-2-77}
$t=0$ 时从同一点以相同速率 $v_0$ 向各个方向抛出小球。证明任意时刻这些小球位于一个球面上，球心以自由落体加速度下落，并求球半径。

\begin{qualityfigure}
\begin{tikzpicture}[>=Latex,font=\small]
\fill (0,2) circle(2pt) node[above left]{$O$};\draw[->,Gold] (0,1.9)--(0,.8) node[midway,left]{$\vect g$};
\draw[DeepBlue] (0,0) circle(1.5);\draw[dashed] (-1.5,0) arc(180:360:1.5 and .35);\draw (-1.5,0) arc(180:0:1.5 and .35);
\fill (0,0) circle(2pt) node[below left]{$C$};\fill (1.06,1.06) circle(2pt) node[right]{$P$};\draw[dashed] (0,0)--(1.06,1.06) node[midway,below right]{$R(t)$};
\end{tikzpicture}
\figcap{自由下落的球心与各向扩展的球面}\label{mech:04:77:q}
\end{qualityfigure}
\end{origexamplebox}
\Needspace{6\baselineskip}\begin{xesolution}
\begin{qualityfigure}
\begin{tikzpicture}[>=Latex,font=\small]
\fill (0,2) circle(2pt) node[above left]{$O$};\draw[->,Gold] (0,1.9)--(0,.8) node[midway,left]{$\vect g$};
\draw[DeepBlue] (0,0) circle(1.5);\draw[dashed] (-1.5,0) arc(180:360:1.5 and .35);\draw (-1.5,0) arc(180:0:1.5 and .35);
\fill (0,0) circle(2pt) node[below left]{$C$};\fill (1.06,1.06) circle(2pt) node[right]{$P$};\draw[->,DeepBlue,thick] (0,0)--(1.06,1.06) node[midway,below right=4pt]{$v_0t$};\draw[->] (0,2)--(0,0) node[pos=.65,left=14pt]{$\frac12\vect g t^2$};\draw[->,Gold] (0,2)--(1.06,1.06) node[midway,above right=6pt]{$\vect r$};
\end{tikzpicture}
\figcap{自由下落的球心与各向扩展的球面（解答用辅助图）}\label{mech:04:77:a}
\end{qualityfigure}
图~\ref{mech:04:77:a}中，$O$ 是共同抛出点，$C$ 是自由下落的球心，任意球 $P$ 相对 $C$ 的矢量正是 $\vect v_0t$，其长度只由初速率决定。

任一小球的位移为
\[
\vect r=\vect v_0t+\frac12\vect g t^2,
\qquad |\vect v_0|=v_0.
\]
取球心位置 $\vect r_C=\frac12\vect g t^2$，则
\[
|\vect r-\vect r_C|=|\vect v_0|t=v_0t.
\]
所以球心由静止以加速度 $\vect g$ 下落，球面半径为
\[
\boxed{R(t)=v_0t}.
\]
在随球心自由下落的参考系中，各小球只是沿初始方向匀速远离原点。
这里使用匀强重力、忽略空气阻力的自由飞行模型；球面结论只适用于各球尚未碰到地面或其他障碍物的阶段。
\end{xesolution}

\subsection*{质心及其运动}

由质量 $m_i$、位置 $\vect r_i$ 的质点组成的系统，总质量 $M=\sum_i m_i$，质心为
\begin{equation}
\vect r_C=\frac1M\sum_i m_i\vect r_i.
\end{equation}
把 $\vect r_i=(x_i,y_i,z_i)$ 按三个坐标方向展开，就是
\begin{equation}
x_C=\frac{\sum_i m_ix_i}{M},\qquad
y_C=\frac{\sum_i m_iy_i}{M},\qquad
z_C=\frac{\sum_i m_iz_i}{M}.
\end{equation}
例如两质点沿 $x$ 轴分布时，$x_C=(m_1x_1+m_2x_2)/(m_1+m_2)$；它不是两位置的简单平均，除非两质量相等。
连续物体可分割为许多微元并取极限。质心位置随各部分运动而改变，是连接运动学与后续动力学的重要量。

质心不是必须由真实物质占据的点，而是质量分布的加权平均位置。质量越大的部分对质心位置影响越大；两个质点的质心位于连线上，并更靠近质量较大的一个。对连续物体，可把它划分为许多足够小的质量元，对各质量元使用同一加权平均，最后取分割无限细的极限。由于质量在刚体内部的位置不变，刚体质心相对刚体也保持固定。

对质心定义分别求一次和二次时间导数，就得到质心速度和加速度分别是各质点速度、加速度的质量加权平均。这些纯运动学关系不要求各质点速度相同，也不要求系统是刚体；只要各部分质量在考察过程中不发生交换，就可以使用。
具体地，因 $m_i$ 和 $M$ 不随时间改变，求导可移入求和号：
\begin{align}
\vect v_C&=\frac{\dd}{\dd t}\left(\frac{\sum_i m_i\vect r_i}{M}\right)
=\frac1M\sum_i m_i\frac{\dd\vect r_i}{\dd t}
=\frac1M\sum_i m_i\vect v_i,\\
\vect a_C&=\frac{\dd\vect v_C}{\dd t}
=\frac1M\sum_i m_i\frac{\dd\vect v_i}{\dd t}
=\frac1M\sum_i m_i\vect a_i.
\end{align}
若质量随时间交换，就不能直接把质量当常数移出求导；那属于变质量系统问题。

以质心为原点的参考系称为质心系。在该系中，各质点质量加权的位置和速度之和都为零：
\begin{equation}
 \sum_i m_i\vect r_i'=0,\qquad
 \sum_i m_i\vect v_i'=0.
\end{equation}
对两质点系统，这意味着它们在质心系中的速度方向相反，大小与对方质量成反比。质心系并不保证每个质点静止，只保证整个系统的质量加权平均运动被扣除；它在以后处理碰撞和系统动力学时尤其重要。

\begin{origexamplebox}{匀速质点系的质心速度}\label{ch04:ex:k07-2-78}
证明：若质点系中每个质点都作匀速运动，质心也作匀速运动。若各质点都沿 $x$ 轴，速度为 $v_i$，写出质心速度。

\begin{qualityfigure}
\begin{tikzpicture}[>=Latex,font=\small]
\draw[->] (-.5,0)--(6.5,0) node[right]{$x$};\node[below left] at (0,0){$O$};
\fill (1,0) circle(2pt) node[below]{$m_1$};\fill (4,0) circle(2pt) node[below]{$m_2$};
\draw[->,DeepBlue,thick] (1,.4)--(2.3,.4) node[above]{$v_1$};
\draw[->,DeepBlue,thick] (4,.4)--(5,.4) node[above]{$v_2$};\fill[Gold] (2.5,0) circle(2pt) node[below]{$C$};
\end{tikzpicture}
\figcap{质量加权位置及质心的平移}\label{mech:04:78:q}
\end{qualityfigure}
\end{origexamplebox}
\Needspace{6\baselineskip}\begin{xesolution}
\begin{qualityfigure}
\begin{tikzpicture}[>=Latex,font=\small]
\draw[->] (-.5,0)--(6.5,0) node[right]{$x$};\node[below left] at (0,0){$O$};
\fill (1,0) circle(2pt) node[below]{$m_1$};\fill (4,0) circle(2pt) node[below]{$m_2$};
\draw[->,DeepBlue,thick] (1,.4)--(2.3,.4) node[above]{$v_1$};
\draw[->,DeepBlue,thick] (4,.4)--(5,.4) node[above]{$v_2$};\fill[Gold] (2.5,0) circle(2pt) node[below]{$C$};\node at (3,-1.2){对 $Mx_C=m_1x_1+m_2x_2$ 求导};
\end{tikzpicture}
\figcap{质量加权位置及质心的平移（解答用辅助图）}\label{mech:04:78:a}
\end{qualityfigure}
图~\ref{mech:04:78:a}中，两点示意质量加权坐标的构造，多质点时对所有点求和。各质量不变时先求导位置加权和，才能得到速度加权和；常速度的加权和仍是常量。

对质心坐标求导：
\[
\boxed{v_C=\frac{\sum_i m_iv_i}{\sum_i m_i}}.
\]
各 $v_i$ 和质量均为常量，所以 $v_C$ 也是常量，质心作匀速直线运动。矢量形式为 $\vect v_C=M^{-1}\sum_i m_i\vect v_i$。
\end{xesolution}

\begin{origexamplebox}{两质点在质心系中的速度}\label{ch04:ex:k07-2-79}
质量为 $m_1,m_2$ 的两质点在一条直线上相对地面以速度 $v_1,v_2$ 匀速运动。求在质心参考系中的速度。

\begin{qualityfigure}
\begin{tikzpicture}[>=Latex,font=\small]
\draw[->] (-.5,0)--(6.5,0) node[right]{$x$};\node[below left] at (0,0){$O$};
\fill (1,0) circle(2pt) node[below]{$m_1$};\fill (4,0) circle(2pt) node[below]{$m_2$};
\draw[->,DeepBlue,thick] (1,.4)--(2.3,.4) node[above]{$v_1$};
\draw[->,DeepBlue,thick] (4,.4)--(5,.4) node[above]{$v_2$};\fill[Gold] (2.5,0) circle(2pt) node[below]{$C$};
\end{tikzpicture}
\figcap{地面速度与随质心平移的速度}\label{mech:04:79:q}
\end{qualityfigure}
\end{origexamplebox}
\Needspace{6\baselineskip}\begin{xesolution}
\begin{qualityfigure}
\begin{tikzpicture}[>=Latex,font=\small]
\draw[->] (-.5,0)--(6.5,0) node[right]{$x$};\node[below left] at (0,0){$O$};
\fill (1,0) circle(2pt) node[below]{$m_1$};\fill (4,0) circle(2pt) node[below]{$m_2$};
\draw[->,DeepBlue,thick] (1,.4)--(2.3,.4) node[above]{$v_1$};
\draw[->,DeepBlue,thick] (4,.4)--(5,.4) node[above]{$v_2$};\fill[Gold] (2.5,0) circle(2pt) node[below]{$C$};\draw[->,Gold] (2.5,1.2)--(3.7,1.2) node[right]{$V_C$};\node at (3,-1.2){$m_1v_1'+m_2v_2'=0$};
\end{tikzpicture}
\figcap{地面速度与随质心平移的速度（解答用辅助图）}\label{mech:04:79:a}
\end{qualityfigure}
图~\ref{mech:04:79:a}中，$C$ 随系统质心平移，质心系中每个速度均减去 $V_C$。解出的相对速度方向相反，且质量加权和应为零。

质心速度 $V_C=(m_1v_1+m_2v_2)/(m_1+m_2)$。在随质心平移的坐标系中，要从每个地面速度扣除同一个 $V_C$，而不是分别扣除两速率的绝对值。例如
\[
v_1'=v_1-V_C
=\frac{(m_1+m_2)v_1-m_1v_1-m_2v_2}{m_1+m_2}.
\]
同理计算 $v_2'=v_2-V_C$，得到
\[
\boxed{v_1'=\frac{m_2}{m_1+m_2}(v_1-v_2)},\qquad
\boxed{v_2'=-\frac{m_1}{m_1+m_2}(v_1-v_2)}.
\]
并有 $m_1v_1'+m_2v_2'=0$，这是质心系定义的直接检验。
\end{xesolution}

\begin{origexamplebox}{质心加速度}\label{ch04:ex:k07-2-80}
质量为 $m_1,m_2$ 的两质点沿直线以加速度 $a_1,a_2$ 匀加速运动，求质心加速度。若质心作匀速直线运动，已知 $a_1$，求 $a_2$。

\begin{qualityfigure}
\begin{tikzpicture}[>=Latex,font=\small]
\draw[->] (-.5,0)--(6.5,0) node[right]{$x$};\node[below left] at (0,0){$O$};
\fill (1,0) circle(2pt) node[below]{$m_1$};\fill (4,0) circle(2pt) node[below]{$m_2$};
\draw[->,DeepBlue,thick] (1,.4)--(2.3,.4) node[above]{$a_1$};
\draw[->,DeepBlue,thick] (4,.4)--(5,.4) node[above]{$a_2$};\fill[Gold] (2.5,0) circle(2pt) node[below]{$C$};
\end{tikzpicture}
\figcap{两质点加速度的质量加权}\label{mech:04:80:q}
\end{qualityfigure}
\end{origexamplebox}
\Needspace{6\baselineskip}\begin{xesolution}
\begin{qualityfigure}
\begin{tikzpicture}[>=Latex,font=\small]
\draw[->] (-.5,0)--(6.5,0) node[right]{$x$};\node[below left] at (0,0){$O$};
\fill (1,0) circle(2pt) node[below]{$m_1$};\fill (4,0) circle(2pt) node[below]{$m_2$};
\draw[->,DeepBlue,thick] (1,.4)--(2.3,.4) node[above]{$a_1$};
\draw[->,DeepBlue,thick] (4,.4)--(5,.4) node[above]{$a_2$};\fill[Gold] (2.5,0) circle(2pt) node[below]{$C$};\node at (3,-1.2){对 $Mx_C=m_1x_1+m_2x_2$ 求导};
\end{tikzpicture}
\figcap{两质点加速度的质量加权（解答用辅助图）}\label{mech:04:80:a}
\end{qualityfigure}
图~\ref{mech:04:80:a}中，$C$ 的位置由两端质量加权；对同一时间连续求导两次，便把权重原样带入加速度关系。质心匀速时，加权加速度总和为零。

从 $x_C=(m_1x_1+m_2x_2)/(m_1+m_2)$ 连续求导两次，先得到 $v_C=(m_1v_1+m_2v_2)/(m_1+m_2)$，再得到
\[
\boxed{a_C=\frac{m_1a_1+m_2a_2}{m_1+m_2}}.
\]
质心匀速意味着 $a_C=0$，故
\[
\boxed{a_2=-\frac{m_1}{m_2}a_1}.
\]
负号表示两质点加速度方向相反。
\end{xesolution}

\subsection*{相对运动的矢量运算}

若 $A$ 相对 $O$ 的位置、速度、加速度分别为 $\vect r_{A/O},\vect v_{A/O},\vect a_{A/O}$，而 $B$ 相对 $A$ 的相应量已知，则
\begin{align}
\vect r_{B/O}&=\vect r_{B/A}+\vect r_{A/O},\\
\vect v_{B/O}&=\vect v_{B/A}+\vect v_{A/O},\\
\vect a_{B/O}&=\vect a_{B/A}+\vect a_{A/O}.
\end{align}

\VelocityTriangleDiagram

下标应按“前者相对后者”阅读。例如 $\vect v_{B/O}$ 表示 $B$ 相对 $O$ 的速度。式子对应一条首尾相接的矢量链：先从 $O$ 到 $A$，再从 $A$ 到 $B$，总效果就是从 $O$ 到 $B$。位置关系对时间求导，便得到同样形式的速度和加速度关系。船—水—岸、飞机—空气—地面等问题都应先写清这条链，再画速度三角形，避免仅凭“顺流相加、逆流相减”的口号误判二维方向。

相对速度还有两个常用几何解释。其一，$\vect v_{B/A}=\vect v_B-\vect v_A$ 是在速度平面中由 $\vect v_A$ 终点指向 $\vect v_B$ 终点的矢量；其二，在短时间 $\Delta t$ 内，它给出 $B$ 相对 $A$ 的位置变化 $\Delta\vect r_{B/A}\simeq\vect v_{B/A}\Delta t$。因此两匀速质点的最短距离可以转化为初始相对位置到相对运动直线的垂距，拦截问题则转化为相对位置能否沿相对速度方向缩短到零。

二维风、河流问题中，速度三角形必须保持同一矢量链。例如“船对岸”$=$“船对水”$+$“水对岸”；若改求“船对水”，应移项得到“船对岸”减“水对岸”。下标顺序一旦颠倒，矢量方向也会反向。画图时让相加矢量首尾相接，并把题目给定的速率标在对应边上，是最稳妥的方向检查。

三维平动参考系的关系与一维完全相同，只需把标量坐标改成矢量：
\begin{equation}
 \vect r'=\vect r-\vect r_0(t),\qquad
 \vect v'=\vect v-\vect v_0(t),\qquad
 \vect a'=\vect a-\vect a_0(t).
\end{equation}
各分量都使用同一个时间，坐标轴保持平行。所谓“迎面”“横风”“垂直击中”等文字条件，最终都应翻译为相对速度在特定方向上的分量条件。

追踪问题与普通直线拦截不同：若追踪者速度方向始终指向目标，它的方向会随视线转动，轨迹通常是曲线。速率恒定只意味着切向加速度为零，视线角速度仍会产生法向加速度。例~\ref{ch04:ex:k07-2-88}发射瞬间，导弹速度竖直而飞机横向运动，视线立刻转动，所以导弹从一开始就有水平法向加速度，不能把它当成先直飞一段再转弯。

相对运动答案应做极限检查：河水或风速趋于零时，应退化为静止介质中的结果；追踪者速率趋近目标速率时，捕获时间应显著增大；两个参考系相对速度为零时，坐标差只剩常量平移。这些极限能快速暴露加减号或速度下标写反的问题。

\begin{origexamplebox}{飞机与风的速度合成}\label{ch04:ex:k07-2-81}
机载罗盘指出飞机机头向正东，空气相对地面向正北运动。画出飞机相对地面的速度。飞行员应怎样驾驶，才能使地面航迹为正东？本题的风方向指空气速度的去向，不指气象术语中的来向。

\begin{qualityfigure}
\begin{tikzpicture}[>=Latex,font=\small]
\draw[->] (0,0)--(4,0) node[right]{东};\draw[->] (0,0)--(0,3.4) node[above]{北};\draw[->,DeepBlue,thick] (0,0)--(2.8,0) node[midway,below]{$\vect v_{\text{机/气}}$};\draw[->,Gold,thick] (2.8,0)--(2.8,2) node[midway,right]{$\vect v_{\text{气/地}}$};
\end{tikzpicture}
\figcap{机头正东、空气向北的给定速度}\label{mech:04:81:q}
\end{qualityfigure}
\end{origexamplebox}
\Needspace{6\baselineskip}\begin{xesolution}
以下推导所用的解图为图~\ref{ch04:fig:k07-ex-2-81}。

飞机对地速度为飞机对空气速度与风对地速度的矢量和。机头正东、风向正北时，合速度指向东北。若要地面航迹正东，机头必须向东偏南，使其向南的空气速度分量恰好抵消风的北向分量。设空速为 $V$、风速为 $u$，偏南角为 $\theta$，则
\AircraftWindDiagram
左图中蓝色合速度为 $(V,u)$，与正东的夹角为 $\arctan(u/V)$。右图把机头改为东偏南 $\theta$，于是合速度的北向分量为 $-V\sin\theta+u$；令它等于零，得到
\[
\boxed{V\sin\theta=u},\qquad u<V.
\]
地面向东速率为 $V\cos\theta=\sqrt{V^2-u^2}$。必须严格要求 $u<V$：$u=V$ 时机头只能正南，地速为零，并没有向东飞行；$u>V$ 时更不可能抵消北向漂移。
\end{xesolution}


\begin{origexamplebox}{渡河的最短时间和最小位移}\label{ch04:ex:k07-2-82}
河宽 $d$，水流对岸速率为 $v_2$。小船相对水保持速率 $v_1$，船头方向可选。求最短渡河时间和最小位移。求最优航向时先取航向恒定；随后说明改变航向能否进一步改善结果。
\par 题图见图~\ref{ch04:ref:question-figure-2-82}。
\KQuestionDE
\end{origexamplebox}
\Needspace{6\baselineskip}\begin{xesolution}
以下推导所用的解图为图~\ref{ch04:fig:k07-ex-2-82}。

渡河时间由对岸速度的横河分量决定。船头垂直河岸时该分量最大为 $v_1$，所以
\RiverCrossingDiagram
\[
\boxed{t_{\min}=d/v_1}.
\]
取沿河向下游为 $x$、横河为 $y$。船对水速度与 $+x$ 夹角为 $\eta$ 时，$v_x=v_2+v_1\cos\eta$、$v_y=v_1\sin\eta>0$，到岸时间 $t=d/v_y$，横向漂移 $\Delta x=dv_x/v_y$，位移大小为 $s=\sqrt{d^2+\Delta x^2}$。

（1）若 $v_1>v_2$，可取 $v_1\cos\eta=-v_2$，把船头偏向上游，使沿河分量抵消水流，位移就是河宽：
\[
\boxed{s_{\min}=d},\qquad
t=\frac{d}{\sqrt{v_1^2-v_2^2}}.
\]
（2）若 $v_1=v_2$，令 $v_x=0$ 就会同时令 $v_y=0$，不能完成横渡。任意有限时间的横渡都有 $s>d$；取 $\eta=\pi-\varepsilon$，则
\[
\frac{\Delta x}{d}=\frac{1-\cos\varepsilon}{\sin\varepsilon}
=\tan\frac\varepsilon2\longrightarrow0,
\qquad t=\frac{d}{v_1\sin\varepsilon}\longrightarrow\infty.
\]
因此 $d$ 是位移的下确界，而不是能取得的最小值。

（3）若 $v_1<v_2$，速度圆圆心 $W=(v_2,0)$、半径 $v_1$，圆外原点 $O$ 发出的切线 $OT$ 给出合速度最大的偏转角 $\psi$。在直角三角形 $OWT$ 中，$WT=v_1$、$OW=v_2$，所以
\[
\sin\psi_{\max}=\frac{v_1}{v_2},\qquad
\boxed{s_{\min}=\frac{d}{\sin\psi_{\max}}=\frac{dv_2}{v_1}}.
\]
此时 $v_y=v_1\sqrt{v_2^2-v_1^2}/v_2$，有限渡河时间为 $dv_2/[v_1\sqrt{v_2^2-v_1^2}]$。图中的 $OT$ 是速度，不是河中的位置线；二者仅方向相同。

若中途改变航向，平均对水速度仍位于同一速度圆的内部或边界。这个圆盘给出的合速度方向范围不比圆周更大，横河平均分量也不超过 $v_1$，所以以上最优结果不会被改进。特别在 $v_1=v_2$ 时，每一瞬间 $v_x\ge0$；总漂移为零要求始终逆流，仍无法横渡。
\end{xesolution}


\begin{origexamplebox}{空速、风速与地速}\label{ch04:ex:k07-2-83}
飞机在静止空气中的速率为 $60\,\mathrm{m/s}$，风速为 $30\,\mathrm{m/s}$，飞机仍能以 $60\,\mathrm{m/s}$ 的地速正北飞行。求风向和机头指向。

\begin{qualityfigure}
\begin{tikzpicture}[>=Latex,font=\small]
\draw[->] (0,0)--(4,0) node[right]{东};\draw[->] (0,0)--(0,3.4) node[above]{北};\draw[->,DeepBlue,thick] (0,0)--(0,2.5) node[midway,left]{$\vect v_{\text{机/地}}$};\node[align=left] at (3,2){地速 $60$\\空速 $60$\\风速 $30$};
\end{tikzpicture}
\figcap{正北地速与给定空速、风速}\label{mech:04:83:q}
\end{qualityfigure}
\end{origexamplebox}
\Needspace{6\baselineskip}\begin{xesolution}
以下推导所用的解图为图~\ref{ch04:fig:k07-ex-2-83}。

设空速矢量与正北夹角为 $\theta$。空速与地速大小同为 $60$，两矢量之差是大小 $30$ 的风速，因此等腰三角形给出
\[
30=2(60)\sin\frac\theta2,
\quad
\boxed{\theta=2\arcsin\frac14\approx28.96^\circ}.
\]
\AircraftEqualSpeedDiagram
图中 $OG=OA=60$、$AG=30$；作风速边 $AG$ 的中点 $F$，等腰三角形的 $OF$ 平分顶角且垂直 $AG$，所以 $\sin(\theta/2)=AF/OA=15/60$。这说明半角公式来自哪一个直角三角形。

取东为 $x$ 正向、北为 $y$ 正向。机头北偏东 $28.96^\circ$ 时，空速分量约为 $(29.05,52.5)\,\mathrm{m/s}$，风速必须为 $(0,60)-(29.05,52.5)=(-29.05,7.5)\,\mathrm{m/s}$，即空气向西偏北 $14.48^\circ$ 运动。镜像方案为机头北偏西 $28.96^\circ$、空气向东偏北 $14.48^\circ$ 运动。两种风的速度都朝向北侧；若换用气象学的风来向，要再反向 $180^\circ$。
\end{xesolution}

\begin{origexamplebox}{由烟迹判断风速}\label{ch04:ex:k07-2-84}
轮船先以速率 $v_1$ 向南航行，船上人看到烟囱冒出的烟向东；船速增至 $v_2$ 后，看到烟向东北。假设烟一离开烟囱就立即随空气以风速运动，忽略上升与喷出过程，且两次风速不变。求风速。

\begin{qualityfigure}
\begin{tikzpicture}[>=Latex,font=\small]
\draw[->] (0,0)--(4,0) node[right]{东};\draw[->] (0,0)--(0,3.4) node[above]{北};\draw[->,Gold] (1.5,0)--(1.5,-1.3) node[below]{$\vect v_{\text{船}}$};\draw[->,DeepBlue] (0,1)--(2.4,1) node[right]{第一次烟对船};\draw[->,DeepBlue] (0,1.8)--(1.6,3.4) node[right]{第二次烟对船};
\end{tikzpicture}
\figcap{船向南航行时的两次相对烟迹}\label{mech:04:84:q}
\end{qualityfigure}
\end{origexamplebox}
\Needspace{6\baselineskip}\begin{xesolution}
以下推导所用的解图为图~\ref{ch04:fig:k07-ex-2-84}。

取东为 $x$ 正向、北为 $y$ 正向，风对地速度为 $(w_x,w_y)$。烟立即随风运动，故烟对船速度等于风速减去船速。图中两次船速分别为 $(0,-v_1)$、$(0,-v_2)$：
\SmokeWindDiagram
第一次烟对船速度 $(w_x,w_y+v_1)$ 正东，纵分量必须为零、横分量为正，故 $w_y=-v_1$、$w_x>0$。第二次烟对船速度 $(w_x,w_y+v_2)$ 指东北，两个正分量相等，故 $w_x=w_y+v_2=v_2-v_1$。于是
\[
\boxed{\vect w=(v_2-v_1,-v_1)},\qquad
\boxed{|\vect w|=\sqrt{(v_2-v_1)^2+v_1^2}}.
\]
风向为东偏南 $\arctan[v_1/(v_2-v_1)]$，并要求 $v_2>v_1$。
\end{xesolution}

\begin{origexamplebox}{鱼雷拦截匀速敌船}\label{ch04:ex:k07-2-85}
潜艇静止于 $A$，敌船在 $B$，$AB=L$。敌船以速率 $v_1$ 匀速直线航行，其方向与 $BA$ 的夹角为 $\alpha$。鱼雷速率为 $v_2$ 且发射后沿直线运动。求能击中敌船的发射方向。
\par 题图见图~\ref{ch04:ref:question-figure-2-85}。
\KQuestionDH
\end{origexamplebox}
\Needspace{6\baselineskip}\begin{xesolution}
以下推导所用的解图为图~\ref{ch04:fig:k07-ex-2-85}。

取潜艇 $A$ 为原点，敌船速度沿 $+x$，并把初始船位放在上半平面；另一侧是镜像。$0\le\alpha\le\pi$，初始 $B=(-L\cos\alpha,L\sin\alpha)$。若在时刻 $t>0$ 相撞，则交点
\[
C=(-L\cos\alpha+v_1t,L\sin\alpha).
\]
鱼雷恰好走过 $AC=v_2t$，因此
\[
(-L\cos\alpha+v_1t)^2+L^2\sin^2\alpha=v_2^2t^2,
\]
展开整理为
\[
\boxed{(v_2^2-v_1^2)t^2+2Lv_1\cos\alpha\,t-L^2=0}.
\]
求出这个方程的\emph{全部正根}，才能列全发射方向。图中 $C_1,C_2$ 展示较慢的鱼雷也可能有两次截击机会：
\InterceptTriangleDiagram
若 $v_2\ne v_1$，两个候选时间为
\[
t_{\pm}=\frac{L[-v_1\cos\alpha\pm\sqrt{v_2^2-v_1^2\sin^2\alpha}]}{v_2^2-v_1^2}.
\]
逐一保留实数且严格为正的根。每个有效根 $t_i$ 给出唯一的单位发射方向
\[
\boxed{\widehat{\vect n}_i=
\left(\frac{-L\cos\alpha+v_1t_i}{v_2t_i},
\frac{L\sin\alpha}{v_2t_i}\right)}.
\]
两分量的平方和为 $1$，是对速率条件的检验。发射线与 $AB$ 的夹角 $\beta_i\in[0,\pi]$ 可由点积确定：
\[
\cos\beta_i=\frac{L-v_1t_i\cos\alpha}{v_2t_i}.
\]
对非退化三角形，正弦定理也给出 $\sin\beta_i=(v_1/v_2)\sin\alpha$，但它不能区分锐角与钝角，不能只取反正弦主值。

可以把存在性分成三种情况（两速率均为正）：
\begin{itemize}
\item $v_2>v_1$：二次项系数为正，常数项为负，两根异号，恰有一次正时间截击。
\item $v_2=v_1$：方程降为一次式。仅当 $\cos\alpha>0$ 时有 $t=L/(2v_1\cos\alpha)$；否则不能截击。
\item $v_2<v_1$：必须同时满足 $\cos\alpha>0$ 与 $v_2\ge v_1\sin\alpha$。严格大于时有两个正根，等于时两根重合，几何上是相切；其余情况无解。
\end{itemize}
例如 $L=100\,\mathrm m$、$v_1=20\,\mathrm{m/s}$、$v_2=10\,\mathrm{m/s}$、$\alpha=20^\circ$，得到 $t_1=3.8331\,\mathrm s$、$t_2=8.6961\,\mathrm s$，对应 $\beta_1=43.160^\circ$、$\beta_2=136.840^\circ$。两次均在敌船未来的位置相遇，不应舍弃后一个钝角解。若 $\alpha=0$，三点共线，直接使用坐标方程仍有效，而不必使用退化三角形的正弦定理。
\end{xesolution}

\begin{origexamplebox}{炸弹与炮弹的相遇}\label{ch04:ex:k07-2-86}
轰炸机以速率 $v_1$ 在高度 $H$ 水平飞行。（1）投弹时距地面目标多远？（2）距目标水平距离 $D$ 处的大炮同时发射炮弹，使其击中飞行中的炸弹，炮弹初速率 $v_2$ 至少多大？（3）取最小值时发射角多大？所有物体只受重力。

\begin{qualityfigure}
\begin{tikzpicture}[>=Latex,font=\small]
\draw[thick] (-3,0)--(4,0);\fill (0,0) circle(2pt) node[below]{目标};
\fill (-2,2.5) circle(2pt) node[left]{投弹点};\draw[->,DeepBlue] (-2,2.5)--(-.5,2.5) node[right]{$v_1$};
\fill (2.5,0) circle(2pt) node[below]{炮位 $s=D$};\fill (-2.5,0) circle(2pt) node[below]{炮位 $s=-D$};
\draw[<->] (3.2,0)--(3.2,2.5) node[midway,right]{$H$};
\draw[<->] (0,-.8)--(2.5,-.8) node[midway,below]{$D$};
\end{tikzpicture}
\figcap{轰炸机、地面目标与两种炮位}\label{mech:04:86:q}
\end{qualityfigure}
\end{origexamplebox}
\Needspace{6\baselineskip}\begin{xesolution}
以下推导所用的解图为图~\ref{ch04:fig:k07-ex-2-86}。

炸弹落地时间
\[
T=\sqrt{\frac{2H}{g}},
\]
故投弹点在目标前方水平距离
\[
\boxed{L=v_1T=v_1\sqrt{\frac{2H}{g}}}.
\]

题面未指定大炮在目标哪一侧，故分别考虑两侧。取目标为 $x=0$，轰炸机沿 $+x$ 飞行，投弹点为 $(-L,H)$；大炮坐标记为 $(s,0)$，其中 $s=+D$ 或 $-D$。图中的 $Q$ 是任一可能的空中相遇点。
\BombInterceptDiagram
炮弹和炸弹具有相同重力加速度，因此相减时重力项消失。具体写出相遇条件：
\[
s+u_xt=-L+v_1t,\qquad
u_yt-\tfrac12gt^2=H-\tfrac12gt^2.
\]
严格要求击中“飞行中的炸弹”时 $0<t<T$，炮弹初速度必须为
\[
u_x(t)=v_1-\frac{L+s}{t},\qquad u_y(t)=\frac Ht,
\]
所以
\[
v_2^2(t)=\left(v_1-\frac{A}{t}\right)^2+\left(\frac Ht\right)^2,
\qquad A=L+s.
\]
若 $A>0$，无约束极小点为
\[
t_0=\frac{A^2+H^2}{v_1A}.
\]
这是由
\[
\frac{\dd v_2^2}{\dd t}
=\frac{2[v_1At-(A^2+H^2)]}{t^3}
\]
求零点得到的：$t<t_0$ 时函数递减，$t>t_0$ 时递增。

当 $A>0$ 且 $t_0<T$ 时，最小值在空中取得，取 $t_*=t_0$，有
\[
\boxed{v_{2,\min}=\sqrt{\left(v_1-\frac{A}{t_*}\right)^2+\left(\frac H{t_*}\right)^2}},
\]
其值还可化为 $v_{2,\min}=v_1H/\sqrt{A^2+H^2}$。若 $A\le0$，导数始终为负；若 $A>0$ 但 $t_0\ge T$，函数在空中相遇的允许区间仍递减。因此这两类情形\emph{没有可取得的最小速率}，只有
\[
\boxed{\inf v_2=\sqrt{\left(v_1-\frac A T\right)^2+\left(\frac H T\right)^2}}.
\]
它只能在 $t\to T^-$ 时无限接近；若另允许落地瞬间相撞，才可令 $t_*=T$ 并称该值为最小值。这两种题意要分开。

发射角 $\alpha$ 从 $+x$ 轴逆时针量起，$0<\alpha<\pi$。对任一允许的相遇时刻，令 $u_x=v_1-A/t$、$u_y=H/t>0$，则
\[
\cos\alpha=\frac{u_x}{\sqrt{u_x^2+u_y^2}},\qquad
\sin\alpha=\frac{u_y}{\sqrt{u_x^2+u_y^2}}.
\]
$u_x>0$ 时向右上方发射，$u_x=0$ 时竖直发射，$u_x<0$ 时向左上方发射；按这三个象限条件取角不会误用反正切主值。若最小值只为下确界，发射角同样只能谈 $t\to T^-$ 的极限角。代入 $s=\pm D$ 即可逐侧判断。各速度分量均为长度除以时间，量纲一致。
\end{xesolution}

\begin{origexamplebox}{三船中点连线的角速度}\label{ch04:ex:k07-2-87}
船 $A$ 以速率 $u$ 正东航行；船 $B$ 以 $2u$ 正北航行；船 $C$ 以 $2\sqrt2u$ 沿东偏北 $45^\circ$ 航行。$t=0$ 时 $B,C$ 同在 $A$ 航线前方，距 $A$ 分别为 $a,2a$。设 $M$ 为 $BC$ 中点，求连线 $MA$ 绕 $M$ 的角速度。
\par 题图见图~\ref{ch04:ref:question-figure-2-87}。
\KQuestionDJ
\end{origexamplebox}
\Needspace{6\baselineskip}\begin{xesolution}
以下推导所用的解图为图~\ref{ch04:fig:k07-ex-2-87}。

图中同时给出了三船的初始相对位置和速度方向。取 $t=0$ 时 $A=(0,0)$，$B=(a,0)$，$C=(2a,0)$：
\ThreeBoatsDiagram
速度分别为
\[
\vect v_A=(u,0),\quad\vect v_B=(0,2u),\quad\vect v_C=(2u,2u).
\]
故 $M(0)=(3a/2,0)$，$\vect v_M=(u,2u)$。从 $M$ 指向 $A$ 的相对位置和相对速度为
\[
\vect r_{A/M}=\left(-\frac{3a}{2},-2ut\right),\qquad
\vect v_{A/M}=(0,-2u).
\]
图中右侧的 $M(t)$ 与 $A(t)$ 的横坐标差始终为 $-3a/2$，纵坐标差为 $-2ut$。若 $\vect r=(x,y)=r(\cos\phi,\sin\phi)$，分别求导后相减得
\[
x\dot y-y\dot x
=r^2(\cos^2\phi+\sin^2\phi)\dot\phi=r^2\dot\phi.
\]
因此方向角变化率为 $(x\dot y-y\dot x)/(x^2+y^2)$。代入 $x=-3a/2$、$y=-2ut$、$\dot x=0$、$\dot y=-2u$：
\[
\Omega=\frac{(-3a/2)(-2u)-(-2ut)0}{9a^2/4+4u^2t^2},
\]
于是
\[
\boxed{\Omega(t)=\frac{12au}{9a^2+16u^2t^2}}.
\]
正号表示按所取坐标方向逆时针转动。
\end{xesolution}

\begin{origexamplebox}{导弹追踪匀速飞机}\label{ch04:ex:k07-2-88}
敌机以水平速率 $v_0$ 在高度 $H$ 飞行。当它位于地面 $A$ 点正上方时，从 $A$ 发射速率恒为 $v_1>v_0$ 的导弹；导弹速度方向始终指向飞机。求发射瞬间的加速度和击中所需总时间。

\begin{qualityfigure}
\begin{tikzpicture}[>=Latex,font=\small]
\draw[thick] (-.5,0)--(6,0);\fill (0,0) circle(2pt) node[below]{$A$};
\fill (0,2.8) circle(2pt) node[left]{飞机};\draw[->,DeepBlue] (0,2.8)--(2,2.8) node[above]{$v_0$};
\draw[dashed] (0,0)--(0,2.8);\draw[<->] (-.6,0)--(-.6,2.8) node[midway,left]{$H$};
\draw[->,Gold] (0,0)--(0,1.3) node[right]{$v_1$（初始）};
\draw[DeepBlue,->] (0,.2)..controls(.5,2)and(2,2.4)..(4,2.8);
\node[right] at (1.7,1){速度始终沿视线};
\end{tikzpicture}
\figcap{恒速导弹的初始视线与飞机高度}\label{mech:04:88:q}
\end{qualityfigure}
\end{origexamplebox}
\Needspace{6\baselineskip}\begin{xesolution}
以下推导所用的解图为图~\ref{ch04:fig:k07-ex-2-88}。

初始视线竖直，飞机的横向速度使视线角速度大小为 $v_0/H$。导弹速率恒定，故初始加速度纯属法向，水平指向飞机运动方向：
\PursuitDiagram
\[
\boxed{a_0=v_1\frac{v_0}{H}}.
\]

设 $r$ 为导弹到飞机的距离，$\theta$ 为视线与飞机速度方向的夹角。相对运动满足
\[
\dot r=v_0\cos\theta-v_1,
\qquad r\dot\theta=-v_0\sin\theta,
\qquad r(\pi/2)=H.
\]
第一式是飞机与导弹在视线上的相对速度；第二式是飞机垂直视线的分量，负号表示 $\theta$ 随追踪减小。令 $q=v_1/v_0>1$，两式相除消去时间：
\[
\frac{\dd r}{\dd\theta}
=r\frac{q-\cos\theta}{\sin\theta},\qquad
\frac{\dd r}{r}=q\frac{\dd\theta}{\sin\theta}
-\frac{\cos\theta}{\sin\theta}\dd\theta.
\]
利用 $\int\dd\theta/\sin\theta=\ln\tan(\theta/2)$ 与 $\int\cot\theta\,\dd\theta=\ln\sin\theta$，得
\[
\ln r=q\ln\tan(\theta/2)-\ln\sin\theta+C.
\]
初始 $r=H$、$\theta=\pi/2$，两三角函数都为 $1$，故 $C=\ln H$，于是
\[
r=H\frac{\tan^q(\theta/2)}{\sin\theta}.
\]
命中时 $r\to0$，对应 $\theta\to0$。由 $\dd t=-r\,\dd\theta/(v_0\sin\theta)$，把反向上下限交换：
\[
T=\frac H{v_0}\int_0^{\pi/2}
\frac{\tan^q(\theta/2)}{\sin^2\theta}\,\dd\theta.
\]
令 $u=\tan(\theta/2)$，则 $\sin\theta=2u/(1+u^2)$、$\dd\theta=2\dd u/(1+u^2)$，上下限变为 $0,1$。因此
\[
T=\frac H{2v_0}\int_0^1(u^{q-2}+u^q)\dd u
=\frac H{2v_0}\left(\frac1{q-1}+\frac1{q+1}\right).
\]
第一项在 $u=0$ 收敛要求 $q>1$，恰好是题设“导弹比飞机快”。合并分式，得到
\[
T=\frac{H}{v_0}\frac{q}{q^2-1}
=\boxed{\frac{Hv_1}{v_1^2-v_0^2}}.
\]
$v_1\downarrow v_0$ 时捕获时间趋于无穷，符合极限判断。
\end{xesolution}
}

{
\graphicspath{{../../../sample/xes-original/kinematics-supplement/}{../../../sample/xes-original/kinematics-supplement/assets/}}
\section{受约束的运动}
\newcommand{\KinematicKeepSpace}[1]{%
\par\begingroup
\dimen0=#1\relax
\dimen2=\lastskip
\advance\dimen0 by \dimen2
\vskip-\dimen2
\vskip0pt plus\dimen0
\penalty-100
\vskip0pt plus-\dimen0
\vskip\dimen0
\penalty9999
\vskip-\dimen0
\vskip\dimen2
\endgroup}

\newcommand{\SlidingRodDiagram}{%
\begin{center}\begin{minipage}{\linewidth}\centering\begin{tikzpicture}[>=Latex,scale=.88]
\begin{scope}
\coordinate (O) at (0,0);\coordinate (A) at (0,3.6);\coordinate (B) at (4,0);\coordinate (P) at (1.52,2.232);
\draw[thick] (0,4.1)--(O)--(4.9,0);\draw[kcurve] (A)--(B);
\foreach \p/\pos in {A/left,B/below,P/above right}{\fill (\p) circle(1.5pt);\node[\pos] at (\p){$\p$};}
\draw[vector] (A)--++(0,-.8) node[midway,left]{$v_A$};\draw[vector] (B)--++(.8,0) node[above]{$v_B$};
\draw (0,2.9) arc(-90:-41.987:.7);\node at (.42,2.7){$\theta$};
\node[above right] at (.75,3.25){$aL$};\node[above right] at (2.7,1.4){$(1-a)L$};
\end{scope}
\begin{scope}[xshift=7cm]
\draw[kaxis] (0,0)--(3.3,0) node[right]{$x$};\draw[kaxis] (0,0)--(0,4.1) node[above]{$y$};
\draw[kcurve,domain=0:90,samples=60] plot({.38*sqrt(28.96)*sin(\x)},{.62*sqrt(28.96)*cos(\x)});
\coordinate (P) at (1.52,2.232);\fill (P) circle(1.5pt);
\draw[vector] (P)--++(.342,-.62);\node[above right] at (P){$P$};
\node[right] at (2,1.55){$\vect v_P$};
\node[above left] at (0,3.337){$(1-a)L$};\node[below] at (2.045,0){$aL$};
\end{scope}
\end{tikzpicture}\figcap{解图：左图杆长固定、端点接触墙与地面；右图使用同一 $P$，速度沿四分之一椭圆的切线。}\label{ch04:fig:k08-ex-2-89}\end{minipage}\end{center}}

\newcommand{\BoatTowDiagram}{%
\begin{center}\begin{minipage}{\linewidth}\centering\begin{tikzpicture}[>=Latex,scale=.93]
\coordinate (O) at (0,2.5);\coordinate (B) at (4,0);
\draw[thick] (-3.8,2.3)--(-.15,2.3)--(-.15,-.4);
\fill[Blue!5] (0,-.7) rectangle (5.3,0);\draw[thin,Blue] (0,0)--(5.3,0);
\draw[thick] (-3.2,2.5) rectangle (-2.2,3);\draw (-3,2.4) circle(.1) (-2.4,2.4) circle(.1);
\draw[kcurve] (-2.2,2.5)--(O)--(B);
\fill (O) circle(1.8pt);\node[above right] at (O){$O$};
\draw[thick] (B)--(5.1,0)--(4.9,-.4)--(4.2,-.4)--cycle;
\fill (B) circle(1.5pt);\node[above right] at (B){$B$};
\draw[vector] (-2.3,3.45)--(-3.4,3.45) node[midway,above]{$v$};
\draw[vector] (4.9,-1)--(3.9,-1) node[midway,below]{$u$};
\node[above] at (-1.2,2.5){$l_c$};\node[above right] at (1.9,1.32){$l$};
\draw[<->] (-4.25,0)--(-4.25,2.5) node[midway,left]{$h$};
\draw[kguide] (-4.25,0)--(0,0) (-4.25,2.5)--(-3.5,2.5);
\draw (B)++(148:.9) arc(148:180:.9);\node at (2.82,.31){$\theta$};
\draw[<->] (0,-.85)--(4,-.85) node[midway,below]{$x$};
\end{tikzpicture}\figcap{解图：$O$ 表示半径忽略的定滑轮转向点；汽车侧增长 $\dot l_c=v$，船侧缩短 $\dot l=-v$。}\label{ch04:fig:k08-ex-2-90}\end{minipage}\end{center}}

\newcommand{\MovingConstraintDiagram}{%
\begin{center}\begin{minipage}{\linewidth}\centering\begin{tikzpicture}[>=Latex,scale=.9]
\begin{scope}
\draw[thick] (0,0)--(4,0)--(0,2)--cycle;
\coordinate (P) at (2.2,.9);\draw[very thick] (P)--(2.2,3.1);
\foreach \x in {2.06,2.34}{\draw[thick] (\x,2)--(\x,2.6);}
\draw (1.6,2)--(2.06,2) (2.34,2)--(2.8,2);
\fill (P) circle(1.5pt);\node[above right=3pt] at (P){$P$};
\draw[vector] (.6,-.5)--(1.7,-.5) node[midway,below]{$v_0$};
\draw[vector] (2.7,2.65)--(2.7,3.5) node[right]{$v$};
\draw (3.2,0) arc(180:153.435:.8);\node at (2.95,.24){$\theta$};
\node at (1.8,3.9){(a) 斜面与固定导向};
\end{scope}
\begin{scope}[xshift=7cm]
\coordinate (C) at (0,0);\coordinate (P) at (1.2,1.6);
\draw[thick] (-2,0)--(2,0) arc(0:180:2);
\draw[kguide] (C)--(0,2.2) (C)--(P) (P)--(1.2,0);
\draw[very thick] (P)--(1.2,3.2);
\foreach \x in {1.06,1.34}{\draw[thick] (\x,2.4)--(\x,3);}
\draw (.6,2.4)--(1.06,2.4) (1.34,2.4)--(1.8,2.4);
\fill (P) circle(1.5pt) (C) circle(1.5pt);
\node[below left] at (C){$C=(X,0)$};\node[above right=3pt] at (P){$P$};
\node at (.7,.46){$R$};\node[below] at (1.2,0){$x_0$};
\draw (0,.75) arc(90:53.1301:.75);\node at (.29,1.13){$\theta$};
\draw[vector] (1.8,2.7)--(1.8,3.5) node[right]{$v_y$};
\draw[vector] (.5,-.8)--(1.6,-.8) node[midway,below]{$v_0$};
\node at (0,3.9){(b) 半圆柱与固定导向};
\end{scope}
\end{tikzpicture}\figcap{解图：两侧固定导向只允许杆竖直平移；接触点由同一斜面方程或圆方程确定。}\label{ch04:fig:k08-ex-2-91}\end{minipage}\end{center}}

\newcommand{\RodPulleyDiagram}{%
\begin{center}\begin{minipage}{\linewidth}\centering\begin{tikzpicture}[>=Latex,scale=1]
\coordinate (O) at (0,0);\coordinate (B) at (0,3.4);\coordinate (A) at (2.3,1.2);\coordinate (C) at (-2.2,3.4);
\draw[very thick] (O)--(A) node[midway,below=5pt]{$R$};
\draw[kcurve] (A)--(B)--(C)--(-2.2,.65);
\draw[thick] (-2.55,0) rectangle (-1.85,.65);\node at (-2.2,.32){$M$};
\draw[kguide] (O)--(B);
\foreach \p/\pos in {O/below,B/above,C/above,A/right}{\fill (\p) circle(1.7pt);\node[\pos] at (\p){$\p$};}
\draw[<->] (-.65,0)--(-.65,3.4) node[midway,left]{$H$};
\draw (B)++(-90:.8) arc(-90:-43.727:.8);\node at (.42,2.4){$\alpha$};
\draw[vector] (A)--++(-.48,.92) node[above]{$\vect v_A$};
\draw[->,DeepBlue] (35:.75) arc(35:120:.75);\node[above] at (.3,1.2){$\omega$};
\draw[vector] (-2.85,.6)--(-2.85,-.3) node[below]{$v_M$};
\node[left] at (.83,2.0){$l$};
\end{tikzpicture}\figcap{解图：零半径转向点 $B,C$ 与绳折点重合；正角速度下 $AB$ 缩短而物块下降。}\label{ch04:fig:k08-ex-2-93}\end{minipage}\end{center}}

\newcommand{\KQuestionDL}{%
\begin{center}\begin{minipage}{.97\linewidth}\centering
\begin{tikzpicture}[>=Latex,scale=.9]

\coordinate (O) at (0,0);\coordinate (A) at (0,3.6);\coordinate (B) at (4,0);\coordinate (P) at (1.52,2.232);
\draw[thick] (0,4.1)--(O)--(4.9,0);\draw[kcurve] (A)--(B);
\foreach \p/\pos in {A/left,B/below,P/above right}{\fill (\p) circle(1.5pt);\node[\pos] at (\p){$\p$};}
\draw[vector] (A)--++(0,-.8) node[midway,left]{$v_A$};\draw[vector] (B)--++(.8,0) node[above]{$v_B$};
\draw (0,2.9) arc(-90:-41.987:.7);\node at (.42,2.7){$\theta$};
\node[above right] at (.75,3.25){$aL$};\node[above right] at (2.7,1.4){$(1-a)L$};
\end{tikzpicture}
\figcap{题图：杆的 A、B 端接触墙、地；$\theta$ 从竖直墙量起，$AP=aL$。}\label{ch04:ref:question-figure-2-89}
\end{minipage}\end{center}
}

\newcommand{\KQuestionDM}{%
%
\begin{center}\begin{minipage}{\linewidth}\centering\begin{tikzpicture}[>=Latex,scale=.93]
\coordinate (O) at (0,2.5);\coordinate (B) at (4,0);
\draw[thick] (-3.8,2.3)--(-.15,2.3)--(-.15,-.4);
\fill[Blue!5] (0,-.7) rectangle (5.3,0);\draw[thin,Blue] (0,0)--(5.3,0);
\draw[thick] (-3.2,2.5) rectangle (-2.2,3);\draw (-3,2.4) circle(.1) (-2.4,2.4) circle(.1);
\draw[kcurve] (-2.2,2.5)--(O)--(B);
\fill (O) circle(1.8pt);\node[above right] at (O){$O$};
\draw[thick] (B)--(5.1,0)--(4.9,-.4)--(4.2,-.4)--cycle;
\fill (B) circle(1.5pt);\node[above right] at (B){$B$};
\draw[vector] (-2.3,3.45)--(-3.4,3.45) node[midway,above]{$v$};
\draw[vector] (4.9,-1)--(3.9,-1) node[midway,below]{$u$};
\node[above] at (-1.2,2.5){$l_c$};\node[above right] at (1.9,1.32){$l$};
\draw[<->] (-4.25,0)--(-4.25,2.5) node[midway,left]{$h$};
\draw[kguide] (-4.25,0)--(0,0) (-4.25,2.5)--(-3.5,2.5);
\draw (B)++(148:.9) arc(148:180:.9);\node at (2.82,.31){$\theta$};
\draw[<->] (0,-.85)--(4,-.85) node[midway,below]{$x$};
\end{tikzpicture}\figcap{题图：汽车经岸上定滑轮拉船；绳不可伸长，$h$ 为滑轮相对水面的高度。}\label{ch04:ref:question-figure-2-90}\end{minipage}\end{center}
}

\newcommand{\KQuestionDN}{%
\begin{center}\begin{minipage}{.97\linewidth}\centering
\begin{tikzpicture}[>=Latex,scale=.9]

\draw[thick] (0,0)--(4,0)--(0,2)--cycle;
\coordinate (P) at (2.2,.9);\draw[very thick] (P)--(2.2,3.1);
\foreach \x in {2.06,2.34}{\draw[thick] (\x,2)--(\x,2.6);}
\draw (1.6,2)--(2.06,2) (2.34,2)--(2.8,2);
\fill (P) circle(1.5pt);\node[above right=3pt] at (P){$P$};
\draw[vector] (.6,-.5)--(1.7,-.5) node[midway,below]{$v_0$};
\draw[vector] (2.7,2.65)--(2.7,3.5) node[right]{$v$};
\draw (3.2,0) arc(180:153.435:.8);\node at (2.95,.24){$\theta$};
\end{tikzpicture}
\figcap{题图：斜面向右运动，竖杆位于两侧固定导向之间，杆底 P 与斜面保持接触。}\label{ch04:ref:question-figure-2-91}
\end{minipage}\end{center}
}

\newcommand{\KQuestionDO}{%
\begin{center}\begin{minipage}{.97\linewidth}\centering
\begin{tikzpicture}[>=Latex,scale=.9]

\coordinate (C) at (0,0);\coordinate (P) at (1.2,1.6);
\draw[thick] (-2,0)--(2,0) arc(0:180:2);
\draw[kguide] (C)--(0,2.2) (C)--(P) (P)--(1.2,0);
\draw[very thick] (P)--(1.2,3.2);
\foreach \x in {1.06,1.34}{\draw[thick] (\x,2.4)--(\x,3);}
\draw (.6,2.4)--(1.06,2.4) (1.34,2.4)--(1.8,2.4);
\fill (P) circle(1.5pt) (C) circle(1.5pt);
\node[below left] at (C){$C=(X,0)$};\node[above right=3pt] at (P){$P$};
\node at (.7,.46){$R$};\node[below] at (1.2,0){$x_0$};
\draw (0,.75) arc(90:53.1301:.75);\node at (.29,1.13){$\theta$};
\draw[vector] (.5,-.8)--(1.6,-.8) node[midway,below]{$v_0$};
\end{tikzpicture}
\figcap{题图：半圆柱向右运动；竖杆横坐标固定，P 为接触点，$\theta$ 从竖直方向量起。}\label{ch04:ref:question-figure-2-92}
\end{minipage}\end{center}
}

\newcommand{\KQuestionDP}{%
%
\begin{center}\begin{minipage}{\linewidth}\centering\begin{tikzpicture}[>=Latex,scale=1]
\coordinate (O) at (0,0);\coordinate (B) at (0,3.4);\coordinate (A) at (2.3,1.2);\coordinate (C) at (-2.2,3.4);
\draw[very thick] (O)--(A) node[midway,below=5pt]{$R$};
\draw[kcurve] (A)--(B)--(C)--(-2.2,.65);
\draw[thick] (-2.55,0) rectangle (-1.85,.65);\node at (-2.2,.32){$M$};
\draw[kguide] (O)--(B);
\foreach \p/\pos in {O/below,B/above,C/above,A/right}{\fill (\p) circle(1.7pt);\node[\pos] at (\p){$\p$};}
\draw[<->] (-.65,0)--(-.65,3.4) node[midway,left]{$H$};
\draw (B)++(-90:.8) arc(-90:-43.727:.8);\node at (.42,2.4){$\alpha$};
\draw[->,DeepBlue] (35:.75) arc(35:120:.75);\node[above] at (.3,1.2){$\omega$};
\node[left] at (.83,2.0){$l$};
\end{tikzpicture}\figcap{题图：杆 OA 绕 O 转动，绳从 A 经固定转向点 B、C 连到物块 M。}\label{ch04:ref:question-figure-2-93}\end{minipage}\end{center}
}


前面的多数质点可以在空间中自由运动。真实物体却常受几何条件限制，称为受约束运动。刚杆保持两端距离恒定；不可伸长的绳在绷紧时长度恒定，但松弛时两端距离可以小于绳长；固体表面阻止物体进入内部，若保持接触，则表面形状直接限定运动。约束形式很多，解决问题的核心是先把不变量写成几何方程，再对时间求导。

常见约束的物理含义必须分别判断：
\begin{itemize}
\item 刚杆约束要求杆上任意两点的距离保持不变。若一端固定，另一端只能在以固定端为圆心、杆长为半径的圆或球面上运动；
\item 绳约束只有在绳已经绷紧时才给出“总绳长不变”。柔绳可以弯折并经滑轮改变方向，也可能松弛，所以不能无条件把两端距离写成绳长；
\item 固体表面首先给出“不允许进入内部”的不等式。只有题目明确物体始终贴着表面时，接触点才满足表面方程的等式约束。
\end{itemize}

约束并不直接说明各点速度相等，而是说明某个长度、坐标关系或接触方程在全过程中成立。正确做法是先用坐标把这个不变量写出来，再对时间求导得到速度关系；若要求加速度，再求导一次。求导后出现的三角函数、投影和符号都要回到题图解释，特别要检查接触点是否确实落在两个物体的边界上。

长度约束求导时，最容易出现的错误是把“绳端速度”直接等同于“物体速度”。绳长变化率只等于物体速度在绳方向上的投影；当绳方向本身倾斜或转动时，物体还可能具有垂直绳方向的速度分量。例~\ref{ch04:ex:k08-2-90}中船速大于收绳速率，正是因为只有 $u\cos\theta$ 用于缩短斜绳。固定滑轮只改变绳的方向，若忽略滑轮半径，绕轮部分长度不变。

移动曲面约束可写成 $y=f(x-X(t))$，其中 $X(t)$ 是曲面整体的水平位移，竖杆横坐标 $x$ 固定。对时间求导得
\begin{equation}
 \dot y=-f'(x-X)\dot X,
\end{equation}
它说明竖直速度由接触处斜率与曲面平移速度共同决定；再求导还会出现曲率 $f''$。所以图中杆与曲面的接触点、局部切线方向必须准确，否则即使代数形式看似正确，速度和加速度方向也会随之画反。

刚杆问题也可用距离平方避免根号。若两端坐标为 $(x_A,y_A)$、$(x_B,y_B)$，约束为
\begin{equation}
 (x_A-x_B)^2+(y_A-y_B)^2=L^2.
\end{equation}
求导后得到两端相对速度在杆方向上的投影为零。这不表示两端速度相同，而只表示它们不会沿杆方向彼此靠近或远离。把这一投影关系与端点各自允许的运动方向结合，才能唯一确定速度比例。

约束方程求出的速度关系只在约束有效的阶段成立。绳一旦松弛，原来的等长方程要改成不等式；物体一旦脱离表面，接触高度方程也不再适用。求得结果后应检查绳张紧趋势、接触方向和几何可达范围。若模型给出接近某位置时速度趋于无穷，通常说明理想驱动、刚性或无滑动假设在该极限前已经失效，而不是现实物体真能达到无限速度。

接触约束规定两物体在接触处具有同一个空间位置，且不能沿公法线互相穿入；绳端系于物体时，绳端与系点的位置相同。导向结构则限制杆的横坐标。把这些位置条件与长度不变条件共同使用，才能建立完整的运动方程。

\begin{origexamplebox}{墙角中滑动的刚杆}\label{ch04:ex:k08-2-89}
长为 $L$ 的刚杆两端 $A,B$ 始终分别接触竖直墙和水平地面。（1）杆与墙夹角为 $\theta$，$A$ 以速率 $v_A$ 向下运动，求 $B$ 的瞬时速度。（2）杆从竖直到水平，求距 $A$ 为 $aL$（$0<a<1$）的点 $P$ 的轨迹。（3）求该时刻 $P$ 的速度大小和方向。
\par 题图见图~\ref{ch04:ref:question-figure-2-89}。
\KQuestionDL
\end{origexamplebox}
\Needspace{6\baselineskip}\begin{xesolution}
以下推导所用的解图为图~\ref{ch04:fig:k08-ex-2-89}。

取墙角为原点，$B$ 的横坐标和 $A$ 的纵坐标为
\SlidingRodDiagram
\[
x_B=L\sin\theta,\qquad y_A=L\cos\theta.
\]
因 $A$ 向下，$v_A=-\dot y_A=L\sin\theta\,\dot\theta$。于是
\[
\boxed{v_B=\dot x_B=L\cos\theta\,\dot\theta=v_A\cot\theta}.
\]

$P$ 从 $A$ 向 $B$ 占杆长比例 $a$，故
\[
x_P=aL\sin\theta,\qquad y_P=(1-a)L\cos\theta.
\]
消去 $\theta$：
\[
\boxed{\frac{x_P^2}{(aL)^2}+\frac{y_P^2}{[(1-a)L]^2}=1},
\]
轨迹是第一象限内的四分之一椭圆。
左图的 $AP=aL$、$PB=(1-a)L$ 决定右图两条半轴。杆竖直时 $P=(0,(1-a)L)$，杆水平时 $P=(aL,0)$；从前者向后者运动，$P$ 的速度沿椭圆切线指向右下方。

对坐标求导：
\[
v_{Px}=a v_A\cot\theta,\qquad v_{Py}=-(1-a)v_A.
\]
所以
\[
\boxed{v_P=v_A\sqrt{a^2\cot^2\theta+(1-a)^2}},
\]
方向向右下方，与水平正向的俯角 $\phi$ 满足
\[
\boxed{\tan\phi=\frac{(1-a)\tan\theta}{a}}.
\]
\end{xesolution}

\begin{origexamplebox}{汽车通过定滑轮牵船}\label{ch04:ex:k08-2-90}
岸边汽车以恒定速率 $v$ 沿水平绳段远离定滑轮，拉动不可伸长的绳。滑轮的转向位置高出船上系绳点所在水平面 $h$；忽略滑轮半径，小船系绳点始终处于该水平面。斜绳与水平面夹角为 $\theta$ 时，求船速 $u$。
\par 题图见图~\ref{ch04:ref:question-figure-2-90}。
\KQuestionDM
\end{origexamplebox}
\Needspace{6\baselineskip}\begin{xesolution}
以下推导所用的解图为图~\ref{ch04:fig:k08-ex-2-90}。

设滑轮到船的斜绳长为 $l$，汽车侧水平绳长为 $l_c$。汽车远离转向点，$\dot l_c=v$；总长 $l+l_c$ 不变，故 $\dot l=-v$。两侧长度都在改变，只是变化率之和为零。船向岸速度为 $u$，其在绳方向的分量为 $u\cos\theta$，故
\BoatTowDiagram
\[
\dot l=-u\cos\theta.
\]
两式相等：
\[
\boxed{u=\frac{v}{\cos\theta}}.
\]
也可由几何直接验证：设船与转向点的水平距离为 $x$，则 $l^2=x^2+h^2$。求导得 $l\dot l=x\dot x$，代入 $\dot x=-u$、$x/l=\cos\theta$，仍得 $\dot l=-u\cos\theta$。
船越接近岸边，$\theta$ 越大，理想模型给出的船速越大；现实中绳、动力和阻力会限制这一增长。
\end{xesolution}


\begin{origexamplebox}{平移斜面顶起竖杆}\label{ch04:ex:k08-2-91}
倾角为 $\theta$ 的斜面在水平地面上以速率 $v_0$ 匀速平移，一根杆被限制为只能竖直运动并与斜面接触。求杆的竖直速度。
\par 题图见图~\ref{ch04:ref:question-figure-2-91}。
\KQuestionDN
\end{origexamplebox}
\Needspace{6\baselineskip}\begin{xesolution}
以下推导所用的解图为图~\ref{ch04:fig:k08-ex-2-91}。

在杆的固定横坐标处，斜面水平移动 $\Delta x$ 会使接触高度改变 $\Delta y=\Delta x\tan\theta$。按题图斜面向右移动时，杆被向上顶起，因此
\MovingConstraintDiagram
\[
\boxed{v=v_0\tan\theta}
\]
且方向竖直向上。若斜面反向运动，速度符号相反。
\end{xesolution}

\begin{origexamplebox}{平移半圆柱顶起竖杆}\label{ch04:ex:k08-2-92}
半径为 $R$ 的半圆柱体以水平速率 $v_0$ 平移，圆柱面上搁一根只能竖直运动的杆。接触点半径与竖直方向夹角为 $\theta$。求杆上任一点的速度和加速度。
\par 题图见图~\ref{ch04:ref:question-figure-2-92}。
\KQuestionDO
\end{origexamplebox}
\Needspace{6\baselineskip}\begin{xesolution}
\begin{qualityfigure}
\begin{tikzpicture}[>=Latex,font=\small]
\draw[thick] (-2,0) arc(180:0:2);\draw[thick] (-2,0)--(2,0);\fill (0,0) circle(2pt) node[below]{$C(X,0)$};
\coordinate (P) at (1,1.732);\draw[very thick] (P)--(1,3.5);\fill (P) circle(2pt) node[right]{$P$};
\draw[dashed] (0,0)--(0,2.3) (0,0)--(P) node[midway,right]{$R$};
\draw (.0,.8) arc(90:60:.8) node[midway,above]{$\theta$};
\draw[->,DeepBlue] (-1,-.7)--(.5,-.7) node[right]{$v_0$};\draw[->,Gold] (1.4,2.5)--(1.4,3.3) node[right]{$v_y$};
\node[align=left] at (4,1.7){杆的横坐标固定为 $x_0$\\$x_0-X=R\sin\theta$\\$y=R\cos\theta$};
\end{tikzpicture}
\figcap{圆面接触点的几何约束和杆的竖直速度（解答用辅助图）}\label{mech:04:92:a}
\end{qualityfigure}
图~\ref{mech:04:92:a}中，杆沿固定竖直线运动，而圆心 $C$ 的横坐标为 $X$；$CP=R$ 将接触点 $P$ 的纵坐标与圆心位移联系起来。对约束求导时不能把 $\theta$ 误当作常量。

本题题图~\ref{ch04:ref:question-figure-2-92}中的两条导向边限制杆的横坐标为 $x_0$。以半圆柱底边为 $y=0$，圆心为 $C=(X,0)$。接触点 $P$ 在右半圆上，$CP=R$，因此
\[
x_0-X=R\sin\theta,\qquad y=R\cos\theta.
\]
$\dot X=v_0$，故 $-v_0=R\cos\theta\,\dot\theta$。于是
\[
\boxed{v_y=\dot y=-R\sin\theta\,\dot\theta=v_0\tan\theta}.
\]
再求导：
\[
\boxed{a_y=v_0\sec^2\theta\,\dot\theta
=-\frac{v_0^2}{R\cos^3\theta}}.
\]
速度向上而加速度向下；杆作平动，所以杆上各点具有相同的 $v_y,a_y$。这一区分能避免把“当前向上”误判成“加速度也向上”。
\end{xesolution}


\begin{origexamplebox}{转杆端点牵引物块}\label{ch04:ex:k08-2-93}
杆 $OA$ 长 $R$，可绕 $O$ 在竖直平面内转动。端点 $A$ 系着跨过固定滑轮 $B,C$ 的不可伸长轻绳，另一端连接物块 $M$；滑轮半径忽略，绳始终绷紧。$B$ 在 $O$ 正上方，$OB=H$。当 $BA$ 与 $BO$ 的夹角为 $\alpha$、杆角速度为 $\omega$ 时，求物块速率 $v_M$。取逆时针转动为 $\omega$ 的正方向，图示 $A$ 位于 $OB$ 右侧。
\par 题图见图~\ref{ch04:ref:question-figure-2-93}。
\KQuestionDP
\end{origexamplebox}
\Needspace{6\baselineskip}\begin{xesolution}
以下推导所用的解图为图~\ref{ch04:fig:k08-ex-2-93}。

图中用固定转向点 $B,C$ 表示半径忽略的滑轮。$BC$ 长度不变，故物块竖直绳段的长度变化率与 $AB$ 段相反，速率大小等于 $|\dot{AB}|$。$AB$ 长度的变化率是 $A$ 点速度在 $BA$ 方向的投影。
\RodPulleyDiagram

取 $O=(0,0)$、$B=(0,H)$，从 $B$ 指向 $A$ 的单位矢量为 $(\sin\alpha,-\cos\alpha)$。设 $AB=l$，则
\[
A=(l\sin\alpha,H-l\cos\alpha).
\]
$A$ 绕 $O$ 的速度为 $\vect v_A=\omega(-y_A,x_A)$。投影得
\[
\dot l=\vect v_A\cdot(\sin\alpha,-\cos\alpha)
=-\omega(H-l\cos\alpha)\sin\alpha
-\omega l\sin\alpha\cos\alpha
=-\omega H\sin\alpha.
\]
因此物块速率为
\[
\boxed{v_M=H|\omega|\sin\alpha}.
\]
$R$ 决定哪些 $\alpha$ 几何上可达，但在给定瞬时 $\alpha$ 后投影结果中恰好消去。
在图示 $\omega>0$ 的情况下 $AB$ 缩短、悬垂段增长，所以物块向下；$\omega<0$ 时方向相反。$H\omega$ 的单位是速度，$\sin\alpha$ 无量纲。
\end{xesolution}
}

{
\graphicspath{{../../../sample/xes-original/kinematics-supplement/}{../../../sample/xes-original/kinematics-supplement/assets/}}
\section{刚体运动学}
\newcommand{\KinematicKeepSpace}[1]{%
\par\begingroup
\dimen0=#1\relax
\dimen2=\lastskip
\advance\dimen0 by \dimen2
\vskip-\dimen2
\vskip0pt plus\dimen0
\penalty-100
\vskip0pt plus-\dimen0
\vskip\dimen0
\penalty9999
\vskip-\dimen0
\vskip\dimen2
\endgroup}

\newcommand{\EulerStepsDiagram}{%
\begin{center}\begin{minipage}{\linewidth}\centering\begin{tikzpicture}[>=Latex,scale=.87]
\begin{scope}
\draw[kaxis] (-.2,0)--(2,0) node[right]{$x$};\draw[kaxis] (0,-.2)--(0,2) node[above]{$y$};
\draw[vector] (0,0)--(35:2) node[above right]{$N$};
\draw[->] (.85,0) arc(0:35:.85);\node at (17.5:1.16){$\alpha$};
\node[below left] at (0,0){$O$};\node at (.7,-.85){(1) 绕 $z$};\node at (.7,-1.35){从 $+z$ 端看};
\end{scope}
\begin{scope}[xshift=5cm]
\draw[TextGray,thick] (-1.6,0)--(1.7,0) node[right]{$xy$};
\draw[DeepBlue,thick] (215:1.7)--(35:1.7) node[right]{$XY$};
\draw[kaxis] (0,0)--(90:2.1) node[above]{$z$};
\draw[vector] (0,0)--(125:2.1) node[above left]{$Z$};
\draw[->] (90:1.05) arc(90:125:1.05);\node at (107.5:1.45){$\beta$};
\draw (0,0) circle(.085);\fill (0,0) circle(.02);
\node[right] at (.12,-.4){$N$ 向外};
\node at (0,-1.4){(2) 绕 $N$};\node at (0,-1.9){从 $+N$ 端看};
\end{scope}
\begin{scope}[xshift=9.8cm]
\draw[kguide] (-.2,0)--(2,0) node[right]{$N$};
\draw[vector] (0,0)--(35:2) node[above right]{$X$};
\draw[vector] (0,0)--(125:1.9) node[above left]{$Y$};
\draw[->] (.85,0) arc(0:35:.85);\node at (17.5:1.16){$\gamma$};
\node[below left] at (0,0){$O$};\node at (.5,-.85){(3) 绕 $Z$};\node at (.5,-1.35){从 $+Z$ 端看};
\end{scope}
\end{tikzpicture}\figcap{右手正转的三步欧拉角；第二图中 $XY$ 截线转过 $\beta$，其法线 $Z$ 同转过 $\beta$，二者始终垂直。}\label{ch04:fig:k09-2-20}\end{minipage}\end{center}}

\newcommand{\RigidMotionTypesDiagram}{%
\begin{center}\begin{minipage}{\linewidth}\centering\begin{tikzpicture}[>=Latex,scale=.9]
\begin{scope}
\foreach \dy in {0,1}{\draw[kguide,domain=0:1,samples=45,variable=\t] plot({3*\t},{\dy+.9*sin(180*\t)});}
\draw[thick] (-.3,-.25) rectangle (.3,1.25);
\draw[thick] (2.7,-.25) rectangle (3.3,1.25);
\foreach \x in {0,3}{\foreach \y in {0,1}{\fill (\x,\y) circle(1.5pt);}}
\draw[vector] (0,0)--(.6,.565);\draw[vector] (0,1)--(.6,1.565);
\node[left] at (-.45,0){$P_1$};\node[left] at (-.45,1){$Q_1$};
\node[right] at (3.45,0){$P_2$};\node[right] at (3.45,1){$Q_2$};
\node at (1.5,2.4){(a) 曲线平动};
\node at (1.5,-1){两姿态的对应边平行};
\end{scope}
\begin{scope}[xshift=8cm]
\draw[kguide] (0,-1.2)--(0,1.2);
\draw[thick] (0,-2)--(0,-1.2) node[pos=0,below]{$B$};
\draw[kaxis] (0,1.2)--(0,2.4) node[above]{$A$};
\draw[thick,dashed] (2,0) arc(0:180:2 and 1.2);
\draw[thick] (-2,0) arc(180:360:2 and 1.2);
\coordinate (P) at ({2*cos(315)},{1.2*sin(315)});
\draw[thick] (0,0)--(P);\draw[kguide] (0,0)--(2,0);
\draw[vector] (P)--++(.78,.468);
\draw (.7,0) arc(0:-45:.7 and .42);
\node at (1,-.19){$\theta$};\node at (.46,-.62){$R$};
\fill (P) circle(1.5pt);\node[below] at (P){$P$};\node[left] at (0,0){$O$};
\node[above right] at (2.2,-.38){$\vect v_P$};
\node at (0,-2.7){(b) 定轴转动的圆平面投影};
\end{scope}
\end{tikzpicture}\figcap{平动可以沿曲线；定轴转动中 $OP=R$，点 $P$ 的切向速度为 $|\omega|R$。}\label{ch04:fig:k09-2-21}\end{minipage}\end{center}}

\newcommand{\RigidBarCasesDiagram}{%
\begin{center}\begin{minipage}{\linewidth}\centering\begin{tikzpicture}[>=Latex,scale=.95]
\begin{scope}[xscale=1.5]
\draw[thick] (0,0)--(2.4,0);\draw[kguide] (2.4,0)--(4.8,0);
\draw[vector] (0,0)--(0,1.2) node[above]{$v_0$};\draw[vector] (2.4,0)--(2.4,.6) node[above]{$v_0/2$};
\foreach \x/\lab in {0/A,2.4/B,4.8/I}{\fill (\x,0) circle(1.5pt);\node[below] at (\x,0){$\lab$};}
\fill (1.2,0) circle(1.5pt);\node[above] at (1.2,0){$G$};
\draw[<->] (0,-.7)--(2.4,-.7) node[midway,below]{$l$};\draw[<->] (2.4,-.7)--(4.8,-.7) node[midway,below]{$l$};
\node at (2.4,2){(1) $x_G=l/2,\ x_I=2l$};
\end{scope}
\begin{scope}[yshift=-4.1cm,xscale=1.5]
\draw[thick] (0,0)--(2.4,0);\draw[vector] (0,0)--(0,1.2) node[above]{$v_0$};\draw[vector] (2.4,0)--(2.4,-1.2) node[below]{$-v_0$};
\node[left] at (0,0){$A$};\node[right] at (2.4,0){$B$};\fill (1.2,0) circle(1.6pt);\node[above] at (1.2,0){$G=I$};
\node at (1.2,2){(2) 等质量};
\end{scope}
\begin{scope}[xshift=7cm,yshift=-4.1cm]
\draw[thick] (0,0)--(3.6,0);\draw[vector] (0,0)--(0,1.2) node[above]{$v_0$};\draw[vector] (3.6,0)--(3.6,-1.2) node[below]{$-v_0$};
\node[left] at (0,0){$A$};\node[right] at (3.6,0){$B$};
\fill (1.2,0) circle(1.6pt) (1.8,0) circle(1.6pt);\node[above] at (1.2,0){$G$};\node[below] at (1.8,0){$I$};
\node at (1.8,2){(3) $m_A=2m_B$};
\draw[<->] (0,-.65)--(1.2,-.65) node[midway,below]{$l/3$};
\end{scope}
\end{tikzpicture}\figcap{解图：三问分别标出实际速度、质心 $G$ 与瞬心 $I$；质量改变不改变给定速度场的瞬心。}\label{ch04:fig:k09-ex-2-94}\end{minipage}\end{center}}

\newcommand{\RollingBallDiagram}{%
\begin{center}\begin{minipage}{\linewidth}\centering\begin{tikzpicture}[>=Latex,scale=1]
\draw[thick] (-2,0)--(4,0);\draw[thick] (0,1) circle(1);
\fill (0,1) circle(1.5pt) (0,0) circle(1.5pt);\node[below left] at (0,1){$C$};\node[below left] at (0,0){$P$};
\draw (0,1)--(-.707,1.707);\node[left] at (-.58,1.28){$R$};
\draw[vector] (0,1)--(2.6,1);\node[above] at (2,1){$v_0$};
\draw[->,DeepBlue,thick] (40:1.5)++(0,1) arc(40:145:1.5);
\node[above] at (0,2.65){$\omega_0>0$（逆时针）};
\draw[vector] (0,-.65)--(2.5,-.65) node[midway,below]{$v_P=v_0+R\omega_0$};
\end{tikzpicture}\figcap{解图：接触点的平动速度和转动速度初始都向右；纯滚动要求它们相消。}\label{ch04:fig:k09-ex-2-95}\end{minipage}\end{center}}

\newcommand{\KnifeOutline}{%
\draw[thick,fill=TextGray!16] (-.15,-.012) rectangle (-.045,.012);
\draw[thick,fill=Blue!8] (-.045,-.012)--(.10,0)--(-.045,.012)--cycle;}

\newcommand{\KnifeFlightDiagram}{%
\begin{center}\begin{minipage}{\linewidth}\centering\begin{tikzpicture}[>=Latex]
\begin{scope}[x=2.1cm,y=2.1cm]
\draw[thick] (5,-.1)--(5,1.1);\node[right] at (5,.7){木靶};
\draw[kguide,domain=0:1,samples=100,variable=\t] plot({4.9*\t},{4.9*tan(30)*\t*(1-\t)});
\foreach \t/\n in {0/1,.25/2,.5/3,.75/4,1/5}{
\begin{scope}[shift={({4.9*\t},{4.9*tan(30)*\t*(1-\t)})},rotate={60+300*\t}]
\KnifeOutline\fill (0,0) circle(.7pt);
\end{scope}
\node[above] at ({4.9*\t},{4.9*tan(30)*\t*(1-\t)+.17}){\n};}
\draw[<->] (0,-.35)--(5,-.35) node[midway,below]{$s=5.0\,\mathrm m$};
\node at (2.5,-.8){(a) 质心抛物线与五个等时姿态};
\end{scope}
\begin{scope}[xshift=2.1cm,yshift=-3.8cm]
\begin{scope}[scale=15]\KnifeOutline\end{scope}
\fill (0,0) circle(1.5pt);\node[below] at (0,-.25){$C$};\node[right] at (1.5,0){刀尖};
\draw[<->] (-2.25,-.65)--(1.5,-.65) node[midway,below]{$L=25\,\mathrm{cm}$};
\draw[<->] (0,.65)--(1.5,.65) node[midway,above]{$l=10\,\mathrm{cm}$};
\node at (0,-1.45){(b) 长度定义（局部放大）};
\end{scope}
\begin{scope}[xshift=8cm,yshift=-3.8cm]
\draw[kaxis] (-.3,0)--(2.5,0) node[right]{$x$};
\begin{scope}[rotate=60,scale=10]\KnifeOutline\end{scope}
\draw[vector] (0,0)--(30:2.1) node[above right]{$\vect v_C$};
\draw (.65,0) arc(0:30:.65);\node at (15:.95){$30^\circ$};
\draw (60:1.85) arc(60:360:1.85);\draw[->] (335:1.85) arc(335:360:1.85);
\node[left] at (-1.85,.45){$300^\circ$};
\node[right] at (.25,-.45){$C$};\draw[thin] (.25,-.45)--(0,0);
\node[above] at (60:2.05){$\theta=60^\circ$};
\node at (.3,-2.45){(c) 初姿态与最小正转角};
\end{scope}
\end{tikzpicture}\figcap{解图：尖端在刀轴正向；第5姿态刀尖恰达木靶，质心仍距靶 $l$，不是整段刀身穿靶。}\label{ch04:fig:k09-ex-2-96}\end{minipage}\end{center}}

\newcommand{\KnifeImpactDiagram}{%
\begin{center}\begin{minipage}{\linewidth}\centering\begin{tikzpicture}[>=Latex]
\begin{scope}[x=2cm,y=2cm]
\draw[kguide] (0,0)--(5,0);\draw[thick] (5,-.15)--(5,.5);
\begin{scope}[shift={(4.9,0)}]\KnifeOutline\end{scope}
\fill (0,0) circle(1.5pt);\node[above] at (0,0){$C_0$};\node[above] at (4.9,.15){$C_1$};
\draw[<->] (0,-.35)--(4.9,-.35) node[midway,below]{$d_1=s-l$};
\node at (2.5,.65){(a) 第一把刀击中靶点 $Q$};
\end{scope}
\begin{scope}[xshift=5cm,yshift=-3.9cm,x=15cm,y=15cm]
\begin{scope}[shift={(.10,0)}]\KnifeOutline\end{scope}
\begin{scope}[shift={(-.15,0)}]\KnifeOutline\end{scope}
\draw[thick] (.20,-.08)--(.20,.09);
\fill (-.05,0) circle(1.4pt) (-.15,0) circle(1.4pt) (.10,0) circle(1.4pt);
\node[above] at (-.05,.025){$P$};\node[below] at (-.15,-.015){$C_2$};\node[below] at (.10,-.015){$C_1$};\node[right] at (.20,0){$Q$};
\draw[<->] (-.05,.065)--(.20,.065) node[midway,above]{$L$};
\draw[<->] (-.15,-.085)--(-.05,-.085) node[midway,below]{$l$};
\node at (-.06,-.15){(b) 接触局部放大：$C_2P=l,\ PQ=L$};
\end{scope}
\end{tikzpicture}\figcap{碰撞几何：两刀仅在 $P$ 接触；第一把木柄远端就是第二把刀尖，故 $C_0C_2=s-L-l$。}\label{ch04:fig:k09-ex-2-96-dup2}\end{minipage}\end{center}}

\newcommand{\KQuestionDQ}{%
\begin{center}\begin{minipage}{.97\linewidth}\centering
\begin{tikzpicture}[>=Latex,scale=.9]

\draw[thick] (0,0)--(4,0);\node[point,label=left:$A$] at (0,0){};\node[point,label=right:$B$] at (4,0){};
\draw[vector] (0,0)--(0,1.6) node[above]{$v_A$};\draw[vector] (4,0)--(4,1.6) node[above]{$v_B$};
\draw[<->] (0,-.55)--(4,-.55) node[midway,below]{$l$};

\end{tikzpicture}
\figcap{题图：A、B 用轻杆连接；图示规定竖直向上为端点速度的正方向，各问按给定符号取向。}\label{ch04:ref:question-figure-2-94}
\end{minipage}\end{center}
}

\newcommand{\KQuestionDR}{%
%
\begin{center}\begin{minipage}{\linewidth}\centering\begin{tikzpicture}[>=Latex,scale=1]
\draw[thick] (-2,0)--(4,0);\draw[thick] (0,1) circle(1);
\fill (0,1) circle(1.5pt) (0,0) circle(1.5pt);\node[below left] at (0,1){$C$};\node[below left] at (0,0){$P$};
\draw (0,1)--(-.707,1.707);\node[left] at (-.58,1.28){$R$};
\draw[vector] (0,1)--(2.6,1);\node[above] at (2,1){$v_0$};
\draw[->,DeepBlue,thick] (40:1.5)++(0,1) arc(40:145:1.5);
\node[above] at (0,2.65){$\omega_0>0$（逆时针）};
\end{tikzpicture}\figcap{题图：球初始向右平动并逆时针转动，地面为水平面。}\label{ch04:ref:question-figure-2-95}\end{minipage}\end{center}
}

\newcommand{\KQuestionDS}{%
%
\begin{center}\begin{minipage}{\linewidth}\centering\begin{tikzpicture}[>=Latex]
\begin{scope}[x=2.1cm,y=2.1cm]
\draw[thick] (5,-.1)--(5,1.1);\node[right] at (5,.7){木靶};
\draw[kguide,domain=0:1,samples=100,variable=\t] plot({4.9*\t},{4.9*tan(30)*\t*(1-\t)});
\foreach \t/\n in {0/1,.25/2,.5/3,.75/4,1/5}{
\begin{scope}[shift={({4.9*\t},{4.9*tan(30)*\t*(1-\t)})},rotate={60+300*\t}]
\KnifeOutline\fill (0,0) circle(.7pt);
\end{scope}
\node[above] at ({4.9*\t},{4.9*tan(30)*\t*(1-\t)+.17}){\n};}
\draw[<->] (0,-.35)--(5,-.35) node[midway,below]{$s=5.0\,\mathrm m$};
\node at (2.5,-.8){(a) 质心抛物线与五个等时姿态};
\end{scope}
\begin{scope}[xshift=2.1cm,yshift=-3.8cm]
\begin{scope}[scale=15]\KnifeOutline\end{scope}
\fill (0,0) circle(1.5pt);\node[below] at (0,-.25){$C$};\node[right] at (1.5,0){刀尖};
\draw[<->] (-2.25,-.65)--(1.5,-.65) node[midway,below]{$L=25\,\mathrm{cm}$};
\draw[<->] (0,.65)--(1.5,.65) node[midway,above]{$l=10\,\mathrm{cm}$};
\node at (0,-1.45){(b) 长度定义（局部放大）};
\end{scope}
\begin{scope}[xshift=8cm,yshift=-3.8cm]
\draw[kaxis] (-.3,0)--(2.5,0) node[right]{$x$};
\begin{scope}[rotate=60,scale=10]\KnifeOutline\end{scope}
\draw[vector] (0,0)--(30:2.1) node[above right]{$\vect v_C$};
\draw (.65,0) arc(0:30:.65);\node at (15:.95){$30^\circ$};
\node[right] at (.25,-.45){$C$};\draw[thin] (.25,-.45)--(0,0);
\node[above] at (60:2.05){$\theta=60^\circ$};
\node at (.3,-2.45){(c) 初姿态与速度方向};
\end{scope}
\end{tikzpicture}\figcap{题图：飞刀的飞行姿态、长度及初始方向；木靶与初始质心等高。}\label{ch04:ref:question-figure-2-96}\end{minipage}\end{center}
}


刚体是内部任意两点距离始终不变的质点系统，是对形变可以忽略的真实物体的理想化。仅给出刚体质心位置还不足以确定其状态，还必须给出姿态。三维姿态可由三个参数描述，欧拉角是常用选择之一。

任何真实物体受力后都会产生程度不同的形变；只有当形变量相对物体尺寸足够小、对所研究问题没有显著影响时，才把它理想化为刚体。由于内部质量分布相对刚体不变，质心在刚体上的位置也固定不动。描述质点只需给出一个位置，描述刚体则既要说明代表其整体平移的位置，通常取质心，又要说明物体朝向，因此自由度更多。

\EulerStepsDiagram

欧拉角不是三根固定轴上的三个独立转角，而是按顺序进行的三次转动：先绕固定轴 $z$ 定出两平面的交线（节线）$N$，再绕 $N$ 改变固连轴 $Z$ 的倾角，最后绕 $Z$ 决定刚体自身方位。高中阶段不要求用它作复杂计算，但应知道三维刚体姿态需要三个自由度。

图~\ref{ch04:fig:k09-2-20}中每一步都从正转轴的一端朝原点看，正转向在纸面上表现为逆时针。第二图中 $N$ 垂直纸面向外，$xy$、$XY$ 两平面各表现为一条截线；$z$、$Z$ 分别垂直于它们，所以平面倾角与法线夹角同为 $\beta$。第三步才绕已经倾斜的 $Z$ 改变 $X,Y$ 的方位。三次转轴依次为 $z,N,Z$，不是始终固定的三根坐标轴；交换转动次序通常会改变最终姿态。

常见刚体运动包括：

\begin{itemize}
\item 平动：刚体姿态不变，所有点在同一时刻速度和加速度相同，质心运动即可代表整体运动；
\item 定轴转动：刚体绕空间中固定轴线转动，可用角位置、角速度和角加速度描述；
\item 平面平行运动：平动与转动同时存在、转轴方向保持平行，轮的滚动是典型例子。
\end{itemize}

平面平行运动可以分解为选定参考点的平动与刚体绕该点的转动，通常选质心最方便；也可以在某一瞬间改选速度为零的瞬心，把当时的速度场看成绕瞬心的纯转动。两种分解描述的是同一瞬时运动，参考点的选择不同不会改变刚体各点的实际速度。

定轴转动中，轴上各点速度为零，离轴垂直距离为 $r$ 的点速率为 $|\omega|r$，方向沿当地圆周切线。平面运动在某一瞬间的瞬心具有同样作用：任一点速度垂直于该点到瞬心的连线，大小与距离成正比。因此，已知两个不平行的点速度方向时，可分别作速度的垂线，其交点就是瞬心。若各点速度平行且相等，垂线不相交，瞬心可视为在无穷远处，对应纯平动。

瞬心只由当时的速度场决定，可能随时间移动，也不必与质心重合。质心由质量分布决定，瞬心由运动状态决定；图~\ref{ch04:fig:k09-ex-2-94}的不同质量和端点速度组合正是为了把两者区分开。求瞬心时不应预先把质心标在杆的几何中点，更不能因为某点在图上居中就假定它的速度为零。

\RigidMotionTypesDiagram
图中左侧两个姿态的对应边始终平行，各点沿相同形状、彼此平移的曲线运动，说明平动不一定是直线运动。右侧 $P$ 到固定轴 $AB$ 的垂距为 $R$；从参考半径转过角位置 $\theta$ 后，$P$ 的速度仍垂直 $OP$，速率为 $|\omega|R$。后半圆用虚线表示，半径 $OP$ 与点 $P$ 处切线属于同一圆平面。

平面刚体在任一瞬间都可视为绕一个速度为零的点转动，这个点称为瞬时转动中心，简称瞬心。它可以在刚体内部、边界或外部；纯平动的瞬心可视为在无穷远处。


\begin{origexamplebox}{两质点轻杆刚体的质心与瞬心}\label{ch04:ex:k09-2-94}
质点 $A,B$ 由长 $l$ 的刚性轻杆连接，此刻杆水平，约定有符号速度的正方向均为竖直向上（图中箭头另按各问实际方向画出）。分别求下列情形中质心位置、速度和瞬心位置：\par
（1）$m_A=m_B=m$，$v_A=2v_B=v_0$；\par
（2）$m_A=m_B=m$，$v_A=-v_B=v_0$；\par
（3）$m_A=2m_B=m$，$v_A=-v_B=v_0$。
\par 题图见图~\ref{ch04:ref:question-figure-2-94}。
\KQuestionDQ
\end{origexamplebox}
\Needspace{6\baselineskip}\begin{xesolution}
以下推导所用的解图为图~\ref{ch04:fig:k09-ex-2-94}。

图中三种情形使用同一长度标度；$G$ 表示质心，$I$ 表示瞬心。取 $A$ 为 $x=0$、$B$ 为 $x=l$：
\RigidBarCasesDiagram
由质心的质量加权定义，坐标和速度分别为
\[
x_G=\frac{m_Bl}{m_A+m_B},\qquad
v_G=\frac{m_Av_A+m_Bv_B}{m_A+m_B}.
\]
以 $A$ 为参考点，把运动分解成 $A$ 的平动和绕 $A$ 的转动。转动部分的竖直速度为 $\omega x$，而在 $B$ 点有 $\omega l=v_B-v_A$，所以沿杆各点的竖直速度随 $x$ 线性变化：
\[
v(x)=v_A+\frac{v_B-v_A}{l}x.
\]
瞬心 $I$ 满足 $v(x_I)=0$。

（1）$x_G=l/2$，$v_G=3v_0/4$ 向上；
\[
x_I=-\frac{lv_A}{v_B-v_A}=2l,
\]
故瞬心在 $B$ 外侧、距 $B$ 为 $l$。

（2）$x_G=l/2$，$v_G=0$，且 $x_I=l/2$。刚体此刻绕质心纯转动，角速度大小 $2v_0/l$，方向为顺时针。

（3）由 $m_A=m$、$m_B=m/2$，得
\[
\boxed{x_G=l/3},\qquad \boxed{v_G=v_0/3\text{，向上}}.
\]
端点速度仍等大反向，所以瞬心仍在 $\boxed{x_I=l/2}$，不与质量分布一起移动。三种情况的结果分别说明瞬心由速度场决定，而质心由质量分布决定。
\end{xesolution}


\begin{origexamplebox}{滑动乒乓球转为纯滚动}\label{ch04:ex:k09-2-95}
半径为 $R$ 的球初始质心速度为 $v_0$、向右，同时以角速度 $\omega_0$ 逆时针转动。转为纯滚动前，质心加速度大小为 $a$、向左，角加速度大小为 $\beta$、顺时针。接触点瞬时速度为零后，质心加速度和角加速度均变为零，球保持匀速纯滚动。求最终质心速度的大小和方向。
\par 题图见图~\ref{ch04:ref:question-figure-2-95}。
\KQuestionDR
\end{origexamplebox}
\Needspace{6\baselineskip}\begin{xesolution}
以下推导所用的解图为图~\ref{ch04:fig:k09-ex-2-95}。

题图把平动方向、转动方向和接触点放在同一幅图中。取向右和逆时针为正：
\RollingBallDiagram
滑动阶段
\[
v(t)=v_0-at,\qquad \omega(t)=\omega_0-\beta t.
\]
球底接触点速度为 $v+R\omega$；纯滚动临界条件是
\[
v_0-at_*+R(\omega_0-\beta t_*)=0.
\]
故
\[
t_*=\frac{v_0+R\omega_0}{a+R\beta},
\]
最终有符号速度
\[
\boxed{v_f=v_0-at_*
=\frac{R(\beta v_0-a\omega_0)}{a+R\beta}}.
\]
因此速率为 $|v_f|$；$\beta v_0>a\omega_0$ 时向右，反之向左，等号时质心恰停下。此时必满足 $v_f=-R\omega_f$，可作方向校验。
这里归零的是 $a_C$ 和角加速度，而不是所有物质点的加速度。匀速纯滚动时，取通过质心、平行于角速度的轴线，距此轴垂直距离为 $r$ 的球上点仍有大小 $\omega_f^2r$、指向此轴的向心加速度。这里的轴不是经过接触点的瞬时转动轴；底部瞬时接触点速度虽为零，加速度也不必为零。只有 $v_f=\omega_f=0$ 的静止情形才全部为零。
\end{xesolution}


\begin{origexamplebox}{飞刀的质心抛体与绕质心转动}\label{ch04:ex:k09-2-96}
在一次飞刀表演中，表演者先后甩出两把飞刀：第一把钉在竖直木靶上，到达时刀身水平；第二把的刀尖恰好扎在第一把的木柄上。观察飞刀姿态序列可把运动分解为两部分：质心 $C$ 作抛体运动，飞刀整体同时绕质心 $C$ 匀速转动。要准确命中，必须同时控制质心出手速度的大小和方向，以及绕质心的角速度。\par
飞刀质心 $C$ 与靶点初始等高，水平距离 $s=5.0\,\mathrm m$；刀轴与水平夹角 $\theta=60^\circ$，刀长 $L=25\,\mathrm{cm}$，质心到刀尖距离 $l=10\,\mathrm{cm}$，$g=9.8\,\mathrm{m/s^2}$。\par
（1）质心按 $30^\circ$ 仰角抛出，使刀尖击中靶点且刀身最终水平、刀尖朝向木靶。刀从初始刀尖朝右上方 $60^\circ$ 的姿态逆时针转动，取达到上述姿态的最小正转角。求质心初速率 $v_C$ 和匀速转动角速度 $\omega$。\par
（2）第二把飞刀从同一位置、同一初始姿态出手，刀尖击中第一把刀木柄的远端；碰撞瞬间第二把刀仍为水平姿态，并取从初始姿态转到该姿态的最小正转角。若两次运动的质心最大高度相同，求 $v_C'$、$\omega'$ 分别比 $v_C$、$\omega$ 小百分之几。
\par 题图见图~\ref{ch04:ref:question-figure-2-96}。
\KQuestionDS
\end{origexamplebox}
\Needspace{6\baselineskip}\begin{xesolution}
以下推导所用的解图为图~\ref{ch04:fig:k09-ex-2-96}。

第一把刀最终水平时，质心在刀尖后方 $l$，所以质心的水平射程是
\KnifeFlightDiagram
\[
d_1=s-l.
\]
同高斜抛射程 $d_1=v_C^2\sin60^\circ/g$，故
\[
\boxed{v_C=\sqrt{\frac{g(s-l)}{\sin60^\circ}}
\approx7.45\,\mathrm{m/s}}.
\]
飞行时间
\[
T=\frac{2v_C\sin30^\circ}{g}=\frac{v_C}{g}\approx0.760\,\mathrm s.
\]
刀从 $60^\circ$ 逆时针转到下一次“刀尖朝向木靶”的水平姿态时，方向角为 $360^\circ$，角位移为 $2\pi-\theta=5\pi/3$。途中 $180^\circ$ 虽然刀轴也水平，但刀尖朝后，不满足命中姿态。因此
\[
\boxed{\omega=\frac{2\pi-\theta}{T}\approx6.89\,\mathrm{rad/s}}.
\]
若允许额外整周，角速度还可加 $2\pi n/T$；题设取最小正值。图中的等时姿态由 $x_C=d_1t/T$、$y_C=v_C\sin30^\circ\,t-gt^2/2$、$\varphi=60^\circ+300^\circ t/T$ 同时决定。

第二把刀的刀尖击中第一把刀木柄远端，且碰撞时刀身水平。下图的局部放大部分表明：从木靶 $Q$ 到撞击点 $P$ 要先退回第一把刀的全长 $L$，第二把刀的质心 $C_2$ 还要再退回 $l$：
\KnifeImpactDiagram
所以第二把刀质心的水平射程唯一确定为
\[
d_2=s-L-l=4.65\,\mathrm m.
\]
最大高度相同意味着初始竖直分速度相同；始末质心仍等高，所以飞行时间仍为 $T$。因此
\[
v_C'=\sqrt{\left(\frac{d_2}{T}\right)^2+\left(v_C\sin30^\circ\right)^2},
\]
并可写成比值
\[
\frac{v_C'}{v_C}=\frac{\cos30^\circ}{d_1}
\sqrt{d_2^2+d_1^2\tan^2 30^\circ}\approx0.9620.
\]
故 $v_C'$ 约小 $\boxed{3.8\%}$。题设已经规定相同的最终水平姿态、最小正转角和相同飞行时间，故
\[
\boxed{\omega'=\omega},
\]
即减小 $\boxed{0\%}$。木柄上的撞击点与碰撞姿态均已在题设中固定，因此结果没有多解。
\end{xesolution}
}

{
\graphicspath{{../../../sample/xes-original/kinematics-supplement/}{../../../sample/xes-original/kinematics-supplement/assets/}}
\section{运动学综合}
\newcommand{\KinematicKeepSpace}[1]{%
\par\begingroup
\dimen0=#1\relax
\dimen2=\lastskip
\advance\dimen0 by \dimen2
\vskip-\dimen2
\vskip0pt plus\dimen0
\penalty-100
\vskip0pt plus-\dimen0
\vskip\dimen0
\penalty9999
\vskip-\dimen0
\vskip\dimen2
\endgroup}

\newcommand{\MarsCoplanarPath}{\draw[kcurve] plot coordinates {(-3.3333333,-0.5510407) (-3.3055556,-0.5222965) (-3.2777778,-0.4937709) (-3.2500000,-0.4654674) (-3.2222222,-0.4373898) (-3.1944444,-0.4095421) (-3.1666667,-0.3819279) (-3.1388889,-0.3545514) (-3.1111111,-0.3274165) (-3.0833333,-0.3005273) (-3.0555556,-0.2738881) (-3.0277778,-0.2475030) (-3.0000000,-0.2213764) (-2.9722222,-0.1955128) (-2.9444444,-0.1699166) (-2.9166667,-0.1445923) (-2.8888889,-0.1195447) (-2.8611111,-0.0947784) (-2.8333333,-0.0702983) (-2.8055556,-0.0461093) (-2.7777778,-0.0222163) (-2.7500000,0.0013755) (-2.7222222,0.0246611) (-2.6944444,0.0476352) (-2.6666667,0.0702925) (-2.6388889,0.0926277) (-2.6111111,0.1146351) (-2.5833333,0.1363092) (-2.5555556,0.1576444) (-2.5277778,0.1786350) (-2.5000000,0.1992751) (-2.4722222,0.2195588) (-2.4444444,0.2394801) (-2.4166667,0.2590329) (-2.3888889,0.2782111) (-2.3611111,0.2970085) (-2.3333333,0.3154188) (-2.3055556,0.3334355) (-2.2777778,0.3510522) (-2.2500000,0.3682624) (-2.2222222,0.3850595) (-2.1944444,0.4014369) (-2.1666667,0.4173877) (-2.1388889,0.4329052) (-2.1111111,0.4479827) (-2.0833333,0.4626132) (-2.0555556,0.4767897) (-2.0277778,0.4905054) (-2.0000000,0.5037533) (-1.9722222,0.5165262) (-1.9444444,0.5288173) (-1.9166667,0.5406194) (-1.8888889,0.5519254) (-1.8611111,0.5627284) (-1.8333333,0.5730214) (-1.8055556,0.5827972) (-1.7777778,0.5920490) (-1.7500000,0.6007698) (-1.7222222,0.6089529) (-1.6944444,0.6165913) (-1.6666667,0.6236784) (-1.6388889,0.6302076) (-1.6111111,0.6361725) (-1.5833333,0.6415667) (-1.5555556,0.6463840) (-1.5277778,0.6506184) (-1.5000000,0.6542641) (-1.4722222,0.6573154) (-1.4444444,0.6597670) (-1.4166667,0.6616138) (-1.3888889,0.6628508) (-1.3611111,0.6634734) (-1.3333333,0.6634776) (-1.3055556,0.6628593) (-1.2777778,0.6616150) (-1.2500000,0.6597415) (-1.2222222,0.6572363) (-1.1944444,0.6540970) (-1.1666667,0.6503219) (-1.1388889,0.6459097) (-1.1111111,0.6408597) (-1.0833333,0.6351718) (-1.0555556,0.6288464) (-1.0277778,0.6218845) (-1.0000000,0.6142879) (-0.9722222,0.6060591) (-0.9444444,0.5972011) (-0.9166667,0.5877179) (-0.8888889,0.5776141) (-0.8611111,0.5668950) (-0.8333333,0.5555671) (-0.8055556,0.5436373) (-0.7777778,0.5311136) (-0.7500000,0.5180049) (-0.7222222,0.5043209) (-0.6944444,0.4900722) (-0.6666667,0.4752704) (-0.6388889,0.4599281) (-0.6111111,0.4440586) (-0.5833333,0.4276764) (-0.5555556,0.4107969) (-0.5277778,0.3934361) (-0.5000000,0.3756115) (-0.4722222,0.3573410) (-0.4444444,0.3386438) (-0.4166667,0.3195397) (-0.3888889,0.3000495) (-0.3611111,0.2801948) (-0.3333333,0.2599979) (-0.3055556,0.2394821) (-0.2777778,0.2186712) (-0.2500000,0.1975897) (-0.2222222,0.1762629) (-0.1944444,0.1547165) (-0.1666667,0.1329767) (-0.1388889,0.1110704) (-0.1111111,0.0890246) (-0.0833333,0.0668668) (-0.0555556,0.0446247) (-0.0277778,0.0223265) (0.0000000,0.0000000) (0.0277778,-0.0223265) (0.0555556,-0.0446247) (0.0833333,-0.0668668) (0.1111111,-0.0890246) (0.1388889,-0.1110704) (0.1666667,-0.1329767) (0.1944444,-0.1547165) (0.2222222,-0.1762629) (0.2500000,-0.1975897) (0.2777778,-0.2186712) (0.3055556,-0.2394821) (0.3333333,-0.2599979) (0.3611111,-0.2801948) (0.3888889,-0.3000495) (0.4166667,-0.3195397) (0.4444444,-0.3386438) (0.4722222,-0.3573410) (0.5000000,-0.3756115) (0.5277778,-0.3934361) (0.5555556,-0.4107969) (0.5833333,-0.4276764) (0.6111111,-0.4440586) (0.6388889,-0.4599281) (0.6666667,-0.4752704) (0.6944444,-0.4900722) (0.7222222,-0.5043209) (0.7500000,-0.5180049) (0.7777778,-0.5311136) (0.8055556,-0.5436373) (0.8333333,-0.5555671) (0.8611111,-0.5668950) (0.8888889,-0.5776141) (0.9166667,-0.5877179) (0.9444444,-0.5972011) (0.9722222,-0.6060591) (1.0000000,-0.6142879) (1.0277778,-0.6218845) (1.0555556,-0.6288464) (1.0833333,-0.6351718) (1.1111111,-0.6408597) (1.1388889,-0.6459097) (1.1666667,-0.6503219) (1.1944444,-0.6540970) (1.2222222,-0.6572363) (1.2500000,-0.6597415) (1.2777778,-0.6616150) (1.3055556,-0.6628593) (1.3333333,-0.6634776) (1.3611111,-0.6634734) (1.3888889,-0.6628508) (1.4166667,-0.6616138) (1.4444444,-0.6597670) (1.4722222,-0.6573154) (1.5000000,-0.6542641) (1.5277778,-0.6506184) (1.5555556,-0.6463840) (1.5833333,-0.6415667) (1.6111111,-0.6361725) (1.6388889,-0.6302076) (1.6666667,-0.6236784) (1.6944444,-0.6165913) (1.7222222,-0.6089529) (1.7500000,-0.6007698) (1.7777778,-0.5920490) (1.8055556,-0.5827972) (1.8333333,-0.5730214) (1.8611111,-0.5627284) (1.8888889,-0.5519254) (1.9166667,-0.5406194) (1.9444444,-0.5288173) (1.9722222,-0.5165262) (2.0000000,-0.5037533) (2.0277778,-0.4905054) (2.0555556,-0.4767897) (2.0833333,-0.4626132) (2.1111111,-0.4479827) (2.1388889,-0.4329052) (2.1666667,-0.4173877) (2.1944444,-0.4014369) (2.2222222,-0.3850595) (2.2500000,-0.3682624) (2.2777778,-0.3510522) (2.3055556,-0.3334355) (2.3333333,-0.3154188) (2.3611111,-0.2970085) (2.3888889,-0.2782111) (2.4166667,-0.2590329) (2.4444444,-0.2394801) (2.4722222,-0.2195588) (2.5000000,-0.1992751) (2.5277778,-0.1786350) (2.5555556,-0.1576444) (2.5833333,-0.1363092) (2.6111111,-0.1146351) (2.6388889,-0.0926277) (2.6666667,-0.0702925) (2.6944444,-0.0476352) (2.7222222,-0.0246611) (2.7500000,-0.0013755) (2.7777778,0.0222163) (2.8055556,0.0461093) (2.8333333,0.0702983) (2.8611111,0.0947784) (2.8888889,0.1195447) (2.9166667,0.1445923) (2.9444444,0.1699166) (2.9722222,0.1955128) (3.0000000,0.2213764) (3.0277778,0.2475030) (3.0555556,0.2738881) (3.0833333,0.3005273) (3.1111111,0.3274165) (3.1388889,0.3545514) (3.1666667,0.3819279) (3.1944444,0.4095421) (3.2222222,0.4373898) (3.2500000,0.4654674) (3.2777778,0.4937709) (3.3055556,0.5222965) (3.3333333,0.5510407)};}

\newcommand{\MarsInclinedPath}{\draw[TextGray,thick] plot coordinates {(-3.3295465,-2.0955284) (-3.2802682,-2.0811303) (-3.2311171,-2.0666671) (-3.1820946,-2.0521385) (-3.1332019,-2.0375442) (-3.0844402,-2.0228836) (-3.0358108,-2.0081564) (-2.9873151,-1.9933622) (-2.9389542,-1.9785006) (-2.8907295,-1.9635712) (-2.8426423,-1.9485736) (-2.7946940,-1.9335072) (-2.7468859,-1.9183719) (-2.6992194,-1.9031670) (-2.6516959,-1.8878922) (-2.6043168,-1.8725471) (-2.5570835,-1.8571312) (-2.5099975,-1.8416442) (-2.4630601,-1.8260855) (-2.4162729,-1.8104548) (-2.3696374,-1.7947516) (-2.3231551,-1.7789754) (-2.2768275,-1.7631259) (-2.2306561,-1.7472026) (-2.1846425,-1.7312051) (-2.1387883,-1.7151328) (-2.0930950,-1.6989855) (-2.0475643,-1.6827625) (-2.0021979,-1.6664635) (-1.9569973,-1.6500881) (-1.9119642,-1.6336357) (-1.8671004,-1.6171059) (-1.8224075,-1.6004983) (-1.7778873,-1.5838123) (-1.7335415,-1.5670477) (-1.6893719,-1.5502038) (-1.6453803,-1.5332802) (-1.6015685,-1.5162765) (-1.5579383,-1.4991923) (-1.5144916,-1.4820270) (-1.4712303,-1.4647802) (-1.4281562,-1.4474514) (-1.3852712,-1.4300402) (-1.3425774,-1.4125461) (-1.3000767,-1.3949687) (-1.2577710,-1.3773075) (-1.2156624,-1.3595620) (-1.1737529,-1.3417318) (-1.1320446,-1.3238164) (-1.0905395,-1.3058154) (-1.0492397,-1.2877284) (-1.0081475,-1.2695548) (-0.9672648,-1.2512942) (-0.9265939,-1.2329463) (-0.8861370,-1.2145104) (-0.8458963,-1.1959863) (-0.8058741,-1.1773734) (-0.7660726,-1.1586713) (-0.7264942,-1.1398796) (-0.6871411,-1.1209979) (-0.6480158,-1.1020257) (-0.6091205,-1.0829626) (-0.5704577,-1.0638082) (-0.5320299,-1.0445622) (-0.4938394,-1.0252240) (-0.4558888,-1.0057933) (-0.4181806,-0.9862697) (-0.3807174,-0.9666528) (-0.3435015,-0.9469423) (-0.3065358,-0.9271377) (-0.2698228,-0.9072387) (-0.2333651,-0.8872450) (-0.1971654,-0.8671562) (-0.1612265,-0.8469720) (-0.1255510,-0.8266920) (-0.0901418,-0.8063159) (-0.0550015,-0.7858435) (-0.0201331,-0.7652744) (0.0144606,-0.7446083) (0.0487768,-0.7238450) (0.0828125,-0.7029842) (0.1165648,-0.6820258) (0.1500307,-0.6609693) (0.1832073,-0.6398147) (0.2160915,-0.6185618) (0.2486803,-0.5972103) (0.2809706,-0.5757601) (0.3129592,-0.5542111) (0.3446432,-0.5325631) (0.3760191,-0.5108160) (0.4070840,-0.4889698) (0.4378346,-0.4670243) (0.4682675,-0.4449796) (0.4983796,-0.4228356) (0.5281674,-0.4005923) (0.5576278,-0.3782497) (0.5867573,-0.3558080) (0.6155525,-0.3332672) (0.6440100,-0.3106273) (0.6721263,-0.2878886) (0.6998980,-0.2650512) (0.7273216,-0.2421154) (0.7543935,-0.2190813) (0.7811103,-0.1959492) (0.8074682,-0.1727194) (0.8334638,-0.1493923) (0.8590933,-0.1259683) (0.8843532,-0.1024477) (0.9092398,-0.0788311) (0.9337493,-0.0551189) (0.9578781,-0.0313116) (0.9816224,-0.0074100) (1.0049784,0.0165855) (1.0279425,0.0406742) (1.0505107,0.0648552) (1.0726793,0.0891278) (1.0944445,0.1134912) (1.1158024,0.1379444) (1.1367491,0.1624866) (1.1572809,0.1871168) (1.1773937,0.2118338) (1.1970837,0.2366365) (1.2163470,0.2615238) (1.2351798,0.2864944) (1.2535780,0.3115470) (1.2715378,0.3366802) (1.2890553,0.3618925) (1.3061265,0.3871823) (1.3227475,0.4125482) (1.3389144,0.4379884) (1.3546232,0.4635011) (1.3698701,0.4890844) (1.3846513,0.5147366) (1.3989627,0.5404554) (1.4128005,0.5662389) (1.4261608,0.5920848) (1.4390399,0.6179909) (1.4514339,0.6439547) (1.4633391,0.6699738) (1.4747515,0.6960456) (1.4856676,0.7221674) (1.4960836,0.7483365) (1.5059959,0.7745500) (1.5154007,0.8008047) (1.5242946,0.8270978) (1.5326740,0.8534258) (1.5405354,0.8797856) (1.5478754,0.9061736) (1.5546905,0.9325863) (1.5609776,0.9590199) (1.5667332,0.9854707) (1.5719543,1.0119348) (1.5766377,1.0384079) (1.5807804,1.0648861) (1.5843795,1.0913649) (1.5874320,1.1178398) (1.5899353,1.1443064) (1.5918866,1.1707597) (1.5932834,1.1971951) (1.5941232,1.2236074) (1.5944036,1.2499916)};
\draw[DeepBlue,thick] plot coordinates {(1.5944036,1.2499916) (1.5941294,1.2760148) (1.5933054,1.3020001) (1.5919299,1.3279421) (1.5900008,1.3538355) (1.5875167,1.3796745) (1.5844759,1.4054533) (1.5808772,1.4311661) (1.5767192,1.4568067) (1.5720010,1.4823691) (1.5667215,1.5078468) (1.5608801,1.5332334) (1.5544761,1.5585223) (1.5475092,1.5837068) (1.5399791,1.6087800) (1.5318859,1.6337349) (1.5232296,1.6585643) (1.5140107,1.6832610) (1.5042297,1.7078176) (1.4938874,1.7322266) (1.4829847,1.7564803) (1.4715230,1.7805711) (1.4595036,1.8044909) (1.4469283,1.8282319) (1.4337989,1.8517860) (1.4201176,1.8751449) (1.4058869,1.8983004) (1.3911094,1.9212441) (1.3757880,1.9439675) (1.3599260,1.9664621) (1.3435267,1.9887192) (1.3265940,2.0107302) (1.3091318,2.0324862) (1.2911445,2.0539784) (1.2726366,2.0751981) (1.2536130,2.0961362) (1.2340789,2.1167838) (1.2140397,2.1371319) (1.1935012,2.1571716) (1.1724694,2.1768938) (1.1509508,2.1962896) (1.1289519,2.2153500) (1.1064797,2.2340659) (1.0835415,2.2524286) (1.0601450,2.2704290) (1.0362979,2.2880583) (1.0120084,2.3053078) (0.9872852,2.3221687) (0.9621369,2.3386324) (0.9365727,2.3546905) (0.9106020,2.3703343) (0.8842345,2.3855558) (0.8574802,2.4003466) (0.8303493,2.4146987) (0.8028525,2.4286044) (0.7750006,2.4420557) (0.7468048,2.4550453) (0.7182764,2.4675657) (0.6894271,2.4796097) (0.6602689,2.4911705) (0.6308139,2.5022413) (0.6010746,2.5128157) (0.5710636,2.5228872) (0.5407939,2.5324501) (0.5102784,2.5414984) (0.4795306,2.5500267) (0.4485639,2.5580300) (0.4173921,2.5655031) (0.3860290,2.5724416) (0.3544888,2.5788412) (0.3227855,2.5846979) (0.2909335,2.5900080) (0.2589474,2.5947682) (0.2268418,2.5989754) (0.1946313,2.6026272) (0.1623307,2.6057210) (0.1299549,2.6082551) (0.0975190,2.6102277) (0.0650378,2.6116375) (0.0325264,2.6124838) (0.0000000,2.6127660) (-0.0325264,2.6124838) (-0.0650378,2.6116375) (-0.0975190,2.6102277) (-0.1299549,2.6082551) (-0.1623307,2.6057210) (-0.1946313,2.6026272) (-0.2268418,2.5989754) (-0.2589474,2.5947682) (-0.2909335,2.5900080) (-0.3227855,2.5846979) (-0.3544888,2.5788412) (-0.3860290,2.5724416) (-0.4173921,2.5655031) (-0.4485639,2.5580300) (-0.4795306,2.5500267) (-0.5102784,2.5414984) (-0.5407939,2.5324501) (-0.5710636,2.5228872) (-0.6010746,2.5128157) (-0.6308139,2.5022413) (-0.6602689,2.4911705) (-0.6894271,2.4796097) (-0.7182764,2.4675657) (-0.7468048,2.4550453) (-0.7750006,2.4420557) (-0.8028525,2.4286044) (-0.8303493,2.4146987) (-0.8574802,2.4003466) (-0.8842345,2.3855558) (-0.9106020,2.3703343) (-0.9365727,2.3546905) (-0.9621369,2.3386324) (-0.9872852,2.3221687) (-1.0120084,2.3053078) (-1.0362979,2.2880583) (-1.0601450,2.2704290) (-1.0835415,2.2524286) (-1.1064797,2.2340659) (-1.1289519,2.2153500) (-1.1509508,2.1962896) (-1.1724694,2.1768938) (-1.1935012,2.1571716) (-1.2140397,2.1371319) (-1.2340789,2.1167838) (-1.2536130,2.0961362) (-1.2726366,2.0751981) (-1.2911445,2.0539784) (-1.3091318,2.0324862) (-1.3265940,2.0107302) (-1.3435267,1.9887192) (-1.3599260,1.9664621) (-1.3757880,1.9439675) (-1.3911094,1.9212441) (-1.4058869,1.8983004) (-1.4201176,1.8751449) (-1.4337989,1.8517860) (-1.4469283,1.8282319) (-1.4595036,1.8044909) (-1.4715230,1.7805711) (-1.4829847,1.7564803) (-1.4938874,1.7322266) (-1.5042297,1.7078176) (-1.5140107,1.6832610) (-1.5232296,1.6585643) (-1.5318859,1.6337349) (-1.5399791,1.6087800) (-1.5475092,1.5837068) (-1.5544761,1.5585223) (-1.5608801,1.5332334) (-1.5667215,1.5078468) (-1.5720010,1.4823691) (-1.5767192,1.4568067) (-1.5808772,1.4311661) (-1.5844759,1.4054533) (-1.5875167,1.3796745) (-1.5900008,1.3538355) (-1.5919299,1.3279421) (-1.5933054,1.3020001) (-1.5941294,1.2760148) (-1.5944036,1.2499916)};
\draw[TextGray,thick] plot coordinates {(-1.5944036,1.2499916) (-1.5941232,1.2236074) (-1.5932834,1.1971951) (-1.5918866,1.1707597) (-1.5899353,1.1443064) (-1.5874320,1.1178398) (-1.5843795,1.0913649) (-1.5807804,1.0648861) (-1.5766377,1.0384079) (-1.5719543,1.0119348) (-1.5667332,0.9854707) (-1.5609776,0.9590199) (-1.5546905,0.9325863) (-1.5478754,0.9061736) (-1.5405354,0.8797856) (-1.5326740,0.8534258) (-1.5242946,0.8270978) (-1.5154007,0.8008047) (-1.5059959,0.7745500) (-1.4960836,0.7483365) (-1.4856676,0.7221674) (-1.4747515,0.6960456) (-1.4633391,0.6699738) (-1.4514339,0.6439547) (-1.4390399,0.6179909) (-1.4261608,0.5920848) (-1.4128005,0.5662389) (-1.3989627,0.5404554) (-1.3846513,0.5147366) (-1.3698701,0.4890844) (-1.3546232,0.4635011) (-1.3389144,0.4379884) (-1.3227475,0.4125482) (-1.3061265,0.3871823) (-1.2890553,0.3618925) (-1.2715378,0.3366802) (-1.2535780,0.3115470) (-1.2351798,0.2864944) (-1.2163470,0.2615238) (-1.1970837,0.2366365) (-1.1773937,0.2118338) (-1.1572809,0.1871168) (-1.1367491,0.1624866) (-1.1158024,0.1379444) (-1.0944445,0.1134912) (-1.0726793,0.0891278) (-1.0505107,0.0648552) (-1.0279425,0.0406742) (-1.0049784,0.0165855) (-0.9816224,-0.0074100) (-0.9578781,-0.0313116) (-0.9337493,-0.0551189) (-0.9092398,-0.0788311) (-0.8843532,-0.1024477) (-0.8590933,-0.1259683) (-0.8334638,-0.1493923) (-0.8074682,-0.1727194) (-0.7811103,-0.1959492) (-0.7543935,-0.2190813) (-0.7273216,-0.2421154) (-0.6998980,-0.2650512) (-0.6721263,-0.2878886) (-0.6440100,-0.3106273) (-0.6155525,-0.3332672) (-0.5867573,-0.3558080) (-0.5576278,-0.3782497) (-0.5281674,-0.4005923) (-0.4983796,-0.4228356) (-0.4682675,-0.4449796) (-0.4378346,-0.4670243) (-0.4070840,-0.4889698) (-0.3760191,-0.5108160) (-0.3446432,-0.5325631) (-0.3129592,-0.5542111) (-0.2809706,-0.5757601) (-0.2486803,-0.5972103) (-0.2160915,-0.6185618) (-0.1832073,-0.6398147) (-0.1500307,-0.6609693) (-0.1165648,-0.6820258) (-0.0828125,-0.7029842) (-0.0487768,-0.7238450) (-0.0144606,-0.7446083) (0.0201331,-0.7652744) (0.0550015,-0.7858435) (0.0901418,-0.8063159) (0.1255510,-0.8266920) (0.1612265,-0.8469720) (0.1971654,-0.8671562) (0.2333651,-0.8872450) (0.2698228,-0.9072387) (0.3065358,-0.9271377) (0.3435015,-0.9469423) (0.3807174,-0.9666528) (0.4181806,-0.9862697) (0.4558888,-1.0057933) (0.4938394,-1.0252240) (0.5320299,-1.0445622) (0.5704577,-1.0638082) (0.6091205,-1.0829626) (0.6480158,-1.1020257) (0.6871411,-1.1209979) (0.7264942,-1.1398796) (0.7660726,-1.1586713) (0.8058741,-1.1773734) (0.8458963,-1.1959863) (0.8861370,-1.2145104) (0.9265939,-1.2329463) (0.9672648,-1.2512942) (1.0081475,-1.2695548) (1.0492397,-1.2877284) (1.0905395,-1.3058154) (1.1320446,-1.3238164) (1.1737529,-1.3417318) (1.2156624,-1.3595620) (1.2577710,-1.3773075) (1.3000767,-1.3949687) (1.3425774,-1.4125461) (1.3852712,-1.4300402) (1.4281562,-1.4474514) (1.4712303,-1.4647802) (1.5144916,-1.4820270) (1.5579383,-1.4991923) (1.6015685,-1.5162765) (1.6453803,-1.5332802) (1.6893719,-1.5502038) (1.7335415,-1.5670477) (1.7778873,-1.5838123) (1.8224075,-1.6004983) (1.8671004,-1.6171059) (1.9119642,-1.6336357) (1.9569973,-1.6500881) (2.0021979,-1.6664635) (2.0475643,-1.6827625) (2.0930950,-1.6989855) (2.1387883,-1.7151328) (2.1846425,-1.7312051) (2.2306561,-1.7472026) (2.2768275,-1.7631259) (2.3231551,-1.7789754) (2.3696374,-1.7947516) (2.4162729,-1.8104548) (2.4630601,-1.8260855) (2.5099975,-1.8416442) (2.5570835,-1.8571312) (2.6043168,-1.8725471) (2.6516959,-1.8878922) (2.6992194,-1.9031670) (2.7468859,-1.9183719) (2.7946940,-1.9335072) (2.8426423,-1.9485736) (2.8907295,-1.9635712) (2.9389542,-1.9785006) (2.9873151,-1.9933622) (3.0358108,-2.0081564) (3.0844402,-2.0228836) (3.1332019,-2.0375442) (3.1820946,-2.0521385) (3.2311171,-2.0666671) (3.2802682,-2.0811303) (3.3295465,-2.0955284)};
\draw[->,DeepBlue,thick] plot coordinates {(-0.8732643,-1.2086051) (-0.8296331,-1.1884459) (-0.7862609,-1.1681816) (-0.7431505,-1.1478119) (-0.7003050,-1.1273361) (-0.6577273,-1.1067537) (-0.6154204,-1.0860642) (-0.5733875,-1.0652670) (-0.5316316,-1.0443616)};
\draw[->,DeepBlue,thick] plot coordinates {(0.6894671,2.4795935) (0.6573750,2.4922850) (0.6249252,2.5043828) (0.5921344,2.5158783) (0.5590194,2.5267631) (0.5255976,2.5370295) (0.4918864,2.5466700) (0.4579039,2.5556774) (0.4236680,2.5640451)};
\draw[->,DeepBlue,thick] plot coordinates {(0.8732643,-1.2086051) (0.9171514,-1.2286600) (0.9612918,-1.2486111) (1.0056825,-1.2684589) (1.0503207,-1.2882040) (1.0952037,-1.3078471) (1.1403288,-1.3273886) (1.1856933,-1.3468291) (1.2312945,-1.3661693)};}

\newcommand{\LightTimingDiagram}{%
\begin{center}\begin{minipage}{\linewidth}\centering\begin{tikzpicture}[>=Latex,scale=.95]
\coordinate (A) at (-3,0);\coordinate (B) at (3,0);\coordinate (O) at (-.7,2.5);
\draw[thick] (A)--(B);\draw[kguide] (A)--(O)--(B);
\foreach \p/\pos in {A/below left,B/below right,O/above}{\fill (\p) circle(1.5pt);\node[\pos] at (\p){$\p$};}
\node[above left] at (-1.95,1.2){$r_A$};\node[above right] at (1.4,1.2){$r_B$};
\draw[vector] (-.8,-.7)--(.8,-.7) node[midway,below]{运动员};
\node[above] at (0,0){$L=100\,\mathrm m$};
\end{tikzpicture}\figcap{解图：观察者 $O$ 接收起点 $A$、终点 $B$ 的光信号；传播路程一般不相等。}\label{ch04:fig:k11-ex-2-102}\end{minipage}\end{center}}

\newcommand{\ApparentTransverseDiagram}{%
\begin{center}\begin{minipage}{\linewidth}\centering\begin{tikzpicture}[>=Latex,scale=1]
\coordinate (E) at (0,0);\coordinate (P1) at (0,4);\coordinate (P2) at (2,2.6);\coordinate (F) at (0,2.6);
\draw[thick] (E)--(P1);\draw[thick] (E)--(P2);
\draw[vector] (P1)--(P2);\draw[kguide] (F)--(P2);
\fill (E) circle(1.6pt) (P1) circle(1.6pt) (P2) circle(1.6pt);
\node[below] at (E){地球};\node[above] at (P1){$P_1$};\node[right] at (P2){$P_2$};
\draw[<->] (-.9,0)--(-.9,4) node[midway,left]{$L$};
\node[above right] at (.85,3.55){$v\Delta t$};\node[below right] at (1.15,1.4){$D_2$};
\draw[kguide] (0,4.35)--(0,4.85) (P2)--(2,4.85);
\draw[<->] (0,4.85)--(2,4.85) node[midway,above]{$v\sin\theta\,\Delta t$};
\draw[<->] (-.4,2.6)--(-.4,4);\node[left,align=right] at (-1.3,3.3){接近量\\$v\cos\theta\,\Delta t$};\draw[thin] (-1.25,3.3)--(-.42,3.3);
\draw (0,3.35) arc(-90:-34.992:.65);\node at (.3,3.1){$\theta$};
\end{tikzpicture}\figcap{解图：横向与视线方向的位移分开标注；三角形给出 $D_2$，再对 $v\Delta t/L$ 作一阶展开。}\label{ch04:fig:k11-ex-2-103}\end{minipage}\end{center}}

\newcommand{\MarsRetrogradeDiagram}{%
\begin{center}\begin{minipage}{\linewidth}\centering\begin{tikzpicture}[>=Latex,scale=.9]
\foreach \t/\n/\dx/\dy in {-90/1/0/0,-60/2/3.5/0,-36.37/3/7/0,0/4/10.5/0,36.37/5/1.75/-4,60/6/5.25/-4,90/7/8.75/-4}{
\begin{scope}[shift={(\dx,\dy)}]
\draw[TextGray] (0,0) circle(.8);\draw[thick] (0,0) circle({.8*1.524});
\fill (0,0) circle(1.3pt);\node[above] at (0,0){$S$};
\pgfmathsetmacro{\ea}{.985609*\t-90}
\pgfmathsetmacro{\ma}{.985609*\t*pow(1.524,-1.5)-90}
\pgfmathsetmacro{\ex}{.8*cos(\ea)}\pgfmathsetmacro{\ey}{.8*sin(\ea)}
\pgfmathsetmacro{\mx}{.8*1.524*cos(\ma)}\pgfmathsetmacro{\my}{.8*1.524*sin(\ma)}
\draw[DeepBlue,dashed,thin,->] (\ex,\ey)--({\ex+1.7*(\mx-\ex)},{\ey+1.7*(\my-\ey)});
\fill (\ex,\ey) circle(1.5pt) (\mx,\my) circle(1.5pt);
\node[font=\scriptsize] at (\ea:.5){$E_{\n}$};
\ifdim\t pt<0pt \node[font=\scriptsize,left] at ({\mx-.13},{\my-.16}){$M_{\n}$};
\else \node[font=\scriptsize,right] at ({\mx+.13},{\my-.16}){$M_{\n}$};\fi
\node[font=\footnotesize] at (0,1.7){\n：$t=\t\,\mathrm d$};
\end{scope}}
\node at (5.25,-6){虚线由地球经过火星，指向背景星空；各小图比例相同};
\begin{scope}[xshift=5.25cm,yshift=-8.6cm]
\draw[kaxis] (-3.7,0)--(3.8,0) node[right]{$t/\mathrm d$};
\draw[kaxis] (0,-1.9)--(0,1.45) node[above]{$\lambda$};
\MarsCoplanarPath
\foreach \t/\n/\pos in {-90/1/above left,-60/2/above left,-36.37/3/above,0/4/below left,36.37/5/below,60/6/below right,90/7/below right}{
\pgfmathsetmacro{\ll}{atan2(1.524*sin(.985609*\t*pow(1.524,-1.5))-sin(.985609*\t),1.524*cos(.985609*\t*pow(1.524,-1.5))-cos(.985609*\t))/12}
\fill (\t/27,\ll) circle(1.5pt);\node[\pos] at (\t/27,\ll){\n};}
\draw[<->] (-1.347,-1.75)--(1.347,-1.75) node[midway,below]{逆行};
\node at (0,2.15){同平面模型：七个时刻的视经度};
\node at (0,-2.65){第3、5点：$\dot\lambda=0$};
\end{scope}
\end{tikzpicture}\figcap{解图：七组位置与七个视方向一一对应；冲日取 $t=0$，第3至第5点之间视经度随时间减小。}\label{ch04:fig:k11-ex-2-105}\end{minipage}\end{center}}

\newcommand{\MarsInclinedDiagram}{%
\begin{center}\begin{minipage}{\linewidth}\centering\begin{tikzpicture}[>=Latex]
\draw[kaxis] (-4.4,-2.1)--(4.6,-2.1) node[right]{视经度增加};
\draw[kaxis] (-4.4,-2.1)--(-4.4,3.05) node[above]{视纬度增加};
\MarsInclinedPath
\node[above] at (0,3.25){三维圆轨道示例：轨道倾角 $1.85^\circ$};
\node at (0,-2.8){箭头表示时间先后；蓝色区间才是逆行};
\end{tikzpicture}\figcap{回环与计算模型的区别：加入非零倾角后视纬度也会变化；蓝段按视经度减小的条件确定，不把整个回环都称为逆行。}\end{minipage}\end{center}}

\newcommand{\RomerDiagram}{%
\begin{center}\begin{minipage}{\linewidth}\centering\begin{tikzpicture}[>=Latex,scale=1]
\coordinate (S) at (0,0);\coordinate (J) at (5.2,0);
\pgfmathsetmacro{\ang}{acos(1/5.2)}
\draw[thick] (S) circle(1);\draw[TextGray] ({5.2*cos(-10)},{5.2*sin(-10)}) arc(-10:10:5.2);
\foreach \a/\lab/\pos in {\ang/E_+/above left,-\ang/E_-/below left}{
\coordinate (E) at (\a:1);\draw[kguide] (S)--(E);\draw[thick] (E)--(J);\fill (E) circle(1.6pt);\node[\pos] at (E){$\lab$};
\draw[vector] (E)--++({-.65*sin(\a)},{.65*cos(\a)});}
\fill (S) circle(1.6pt) (J) circle(1.6pt);\node[left] at (S){$S$};\node[right] at (J){$J$};
\draw[kguide] (S)--(J);
\draw[<->] (0,-2.25)--(5.2,-2.25) node[midway,below]{$R_J=5.2R_E$};
\node at (2,1.65){$E_+$：距离增加，周期最长};\node at (2,-1.65){$E_-$：距离减小，周期最短};
\begin{scope}[xshift=9cm]
\draw (0,0) circle(.65);\fill (0,0) circle(2pt);\node[below] at (0,0){$J$};
\fill (45:.65) circle(1.6pt);\node[above right] at (45:.65){木卫I};
\draw[->] (50:.9) arc(50:170:.9);\node at (0,-1.2){卫星轨道放大};
\end{scope}
\end{tikzpicture}\figcap{解图：$SE_\pm\perp E_\pm J$，两切点为周期极值位置；木卫I单独放大，不与日心轨道混用比例。}\label{ch04:fig:k11-ex-2-106}\end{minipage}\end{center}}

\newcommand{\AimLocusDiagram}{%
\begin{center}\begin{minipage}{\linewidth}\centering\begin{tikzpicture}[>=Latex,scale=.9]
\begin{scope}
\draw[thick] (-.4,0)--(4.1,0);\draw[thick] (3.6,0)--(3.6,4.4);
\coordinate (O) at (0,.6);\coordinate (Q) at (3.6,3.7);\coordinate (P) at (3.6,1.6);
\draw[kguide] (O)--(Q);\draw[kcurve,domain=0:1,samples=70,variable=\t] plot({3.6*\t},{.6+3.1*\t-2.1*\t^2});
\foreach \p in {O,Q,P}{\fill (\p) circle(1.5pt);}
\node[above left=3pt] at (O){$O$};\node[right] at (Q){$Q$};\node[below left] at (3.5,1.45){$P\ (z=h)$};
\draw[<->] (0,-.5)--(3.6,-.5) node[midway,below]{$r$};
\draw[<->] (4.65,.6)--(4.65,3.7) node[midway,right]{$Z$};
\draw[kguide] (0,.6)--(4.65,.6);
\node[left] at (0,.25){$z_0$};
\node at (1.8,5){(a) 包含 $OQ$ 的竖直截面};
\end{scope}
\begin{scope}[xshift=9cm]
\draw[kaxis] (-2.4,0)--(2.6,0) node[right]{$x$};\draw[kaxis] (0,0)--(0,4.6) node[above]{$z$};
\draw[DeepBlue,thick] (0,2.7) circle(1.5);\draw[thick] (-2.4,.7)--(2.4,.7);
\coordinate (Q) at (1.2,3.6);\coordinate (P) at (1.2,.7);
\draw[kguide] (Q)--(P);\fill (Q) circle(1.5pt) (P) circle(1.5pt) (0,2.7) circle(1.5pt);
\node[above right] at (Q){$Q$};\node[below right] at (P){$P$};
\node[right] at (2.45,.7){$h$};
\node[left] at (-.1,2.7){$K$};
\node at (0,5.25){(b) 正视墙面};
\end{scope}
\end{tikzpicture}\figcap{解图：瞄准点 $Q$ 位于以 $K$ 为圆心的墙面圆上，实际命中点 $P$ 位于水平线；两图分别说明投掷平面与墙面轨迹。}\label{ch04:fig:k11-ex-2-109}\end{minipage}\end{center}}

\newcommand{\KQuestionDZ}{%
\begin{center}\begin{minipage}{.97\linewidth}\centering
\begin{tikzpicture}[>=Latex,scale=.9]

\coordinate (E) at (0,0);\coordinate (P) at (0,3.3);
\draw[thick] (E)--(P) node[midway,right]{$L$};
\draw (E) circle(.25);\node[below] at (0,-.3){地球 E};
\draw[vector] (P)--++(-135:1.5) node[below left]{$v$};
\draw (P)++(-135:.7) arc(-135:-90:.7);\node at (-.4,2.34){$\theta$};
\node[point,label=above:$P$] at (P){};

\end{tikzpicture}
\figcap{题图：天体初始距地球 $L$，运动方向与指向地球的视线成角 $\theta$。}\label{ch04:ref:question-figure-2-103}
\end{minipage}\end{center}
}

\newcommand{\KQuestionEB}{%
%
\begin{center}\begin{minipage}{\linewidth}\centering\begin{tikzpicture}[>=Latex,scale=.9]
\foreach \t/\n/\dx/\dy in {-90/1/0/0,-60/2/3.5/0,-36.37/3/7/0,0/4/10.5/0,36.37/5/1.75/-4,60/6/5.25/-4,90/7/8.75/-4}{
\begin{scope}[shift={(\dx,\dy)}]
\draw[TextGray] (0,0) circle(.8);\draw[thick] (0,0) circle({.8*1.524});
\fill (0,0) circle(1.3pt);\node[above] at (0,0){$S$};
\pgfmathsetmacro{\ea}{.985609*\t-90}
\pgfmathsetmacro{\ma}{.985609*\t*pow(1.524,-1.5)-90}
\pgfmathsetmacro{\ex}{.8*cos(\ea)}\pgfmathsetmacro{\ey}{.8*sin(\ea)}
\pgfmathsetmacro{\mx}{.8*1.524*cos(\ma)}\pgfmathsetmacro{\my}{.8*1.524*sin(\ma)}
\draw[DeepBlue,dashed,thin,->] (\ex,\ey)--({\ex+1.7*(\mx-\ex)},{\ey+1.7*(\my-\ey)});
\fill (\ex,\ey) circle(1.5pt) (\mx,\my) circle(1.5pt);
\node[font=\scriptsize] at (\ea:.5){$E_{\n}$};
\ifdim\t pt<0pt \node[font=\scriptsize,left] at ({\mx-.13},{\my-.16}){$M_{\n}$};
\else \node[font=\scriptsize,right] at ({\mx+.13},{\my-.16}){$M_{\n}$};\fi
\node[font=\footnotesize] at (0,1.7){\n：$t=\t\,\mathrm d$};
\end{scope}}
\node at (5.25,-6){虚线由地球经过火星，指向背景星空；各小图比例相同};
\end{tikzpicture}\figcap{题图：七个时刻的日、地、火位置与观测视线；E 为地球、M 为火星、S 为太阳。}\label{ch04:ref:question-figure-2-105}\end{minipage}\end{center}
%
\begin{center}\begin{minipage}{\linewidth}\centering\begin{tikzpicture}[>=Latex]
\draw[kaxis] (-4.4,-2.1)--(4.6,-2.1) node[right]{视经度增加};
\draw[kaxis] (-4.4,-2.1)--(-4.4,3.05) node[above]{视纬度增加};
\MarsInclinedPath
\node[above] at (0,3.25){三维圆轨道示例：轨道倾角 $1.85^\circ$};
\node at (0,-2.8){箭头表示观测的时间先后};
\end{tikzpicture}\figcap{同一现象在背景星空上的视位置回环示意。}\end{minipage}\end{center}
}


前面各节广泛讨论了运动学的内容，建立了坐标、矢量、分运动、相对运动和刚体运动等描述方法。现实世界中物体的运动具有更大的复杂性，往往需要把多个观念结合在一起。下面针对几种对未来学习具有启发性、同时又很有趣的问题展开讨论：光传播延迟、天体视运动、观测周期、射程修正和瞄准轨迹。它们既是运动学方法的综合，也为后续相对论、天体力学和波动理论留下接口。


\begin{origexamplebox}{不同位置看到的百米成绩}\label{ch04:ex:k11-2-102}
赛场不同位置的人看到短跑选手从出发到冲过终点所需的时间是否相同？光速为 $c$。

\begin{qualityfigure}
\begin{tikzpicture}[>=Latex,font=\small]
\fill (0,0) circle(2pt) node[below]{$A$（起点）};\fill (5,0) circle(2pt) node[below]{$B$（终点）};\draw[->,thick] (0,0)--(5,0);
\fill (1.5,2.7) circle(2pt) node[above]{观察者 $O$};\draw[dashed] (0,0)--(1.5,2.7) node[midway,left]{$r_A$};\draw[dashed] (5,0)--(1.5,2.7) node[midway,right]{$r_B$};
\end{tikzpicture}
\figcap{观察者到起点和终点的不同光程}\label{mech:04:102:q}
\end{qualityfigure}
\end{origexamplebox}
\Needspace{6\baselineskip}\begin{xesolution}
以下推导所用的解图为图~\ref{ch04:fig:k11-ex-2-102}。

设运动员真实比赛时间为 $T$，观察者到起点、终点的距离分别为 $r_A,r_B$。起跑光信号在 $r_A/c$ 后到达，冲线光信号在 $T+r_B/c$ 后到达，所以观察者由视觉得到的时间间隔为
\LightTimingDiagram
\[
\boxed{T_{\rm obs}=T+\frac{r_B-r_A}{c}}.
\]
因此一般不相同；只有两观察者的 $r_B-r_A$ 相同，或先把光传播延迟校正掉，才会得到同一成绩。设跑道长为 $L=100\,\mathrm m$，三角形的两边之差不大于第三边，所以 $|r_B-r_A|\le L$。这给出单个观察结果相对真实时长的偏差界：
\[
|T_{\rm obs}-T|\le\frac Lc\approx0.334\,\mu\mathrm s.
\]
而起点观察者得到 $T+L/c$，终点观察者得到 $T-L/c$；两位观察者之间的最大差异是
\[
\boxed{\Delta T_{\rm observers,max}=\frac{2L}{c}\approx0.667\,\mu\mathrm s}.
\]
这两个界限比较的对象不同。它们在日常计时中很小，但概念上不能混为一谈。
\end{xesolution}


\begin{origexamplebox}{发光天体的视横向速度}\label{ch04:ex:k11-2-103}
遥远发光天体距地球约为 $L$，以速率 $v$ 运动，速度方向与天体指向地球的连线夹角为 $\theta$。考虑光速 $c$ 有限，求地球观察者看到的横向运动速度。可用小量展开 $(1+ax+bx^2)^{1/2}\simeq1+ax/2$。
这个展开式的意思是：当 $x\ll1$ 时，只保留常数项和 $x$ 的一次项，$x^2$ 及更高次项在数值上远小于前两项，可在一阶近似中忽略。
\par 题图见图~\ref{ch04:ref:question-figure-2-103}。
\KQuestionDZ
\end{origexamplebox}
\Needspace{6\baselineskip}\begin{xesolution}
以下推导所用的解图为图~\ref{ch04:fig:k11-ex-2-103}。

取地球所在惯性系，两次发光的坐标时刻间隔为 $\Delta t$，它不是天体随身钟测得的固有时间。初次发光位置为 $P_1$、到地球距离 $L$；第二次发光位置为 $P_2$。图中沿初始视线方向的接近量为 $v\cos\theta\,\Delta t$，横向位移为 $v\sin\theta\,\Delta t$：
\ApparentTransverseDiagram
由直角三角形，第二段传播路程的精确值为
\begin{align*}
D_2&=\sqrt{(L-v\cos\theta\,\Delta t)^2+(v\sin\theta\,\Delta t)^2}\\
&=L\sqrt{1-2\varepsilon\cos\theta+\varepsilon^2},
\qquad \varepsilon=\frac{v\Delta t}{L}.
\end{align*}
当 $\varepsilon\ll1$，使用题给展开得 $D_2=L-v\cos\theta\,\Delta t+O(L\varepsilon^2)$。两信号的接收时刻分别为 $t_1+L/c$、$t_1+\Delta t+D_2/c$，所以
\[
\Delta t_{\rm obs}=\Delta t+\frac{D_2-L}{c}
\simeq\Delta t\left(1-\frac vc\cos\theta\right).
\]
横向位移除以接收间隔，再取 $\Delta t\to0$ 的瞬时极限，故视横向速度
\[
\boxed{v_{\perp,\rm app}=\frac{v\sin\theta}{1-(v/c)\cos\theta}}.
\]
若 $\theta>90^\circ$，$\cos\theta<0$，天体有远离地球的分量，分母自然大于 $1$，不必另改角度定义。$v<c$ 保证分母始终为正；$v/c\to0$ 时退化为 $v\sin\theta$。当天体接近光速、又以小而非零的角度朝观察者运动时，该“视速度”可超过 $c$；这来自信号到达间隔缩短，并不表示物质或信息真的超光速。
\end{xesolution}


\begin{origexamplebox}{新轨道上的太阳日与一年}\label{ch04:ex:k11-2-104}
可以近似认为地球在一条圆轨道上绕日公转，取日心为参考系。地球公转一周的原恒星年为 $365.2564$ 个太阳日；地球上的人连续两次看见太阳在天空中同一位置的时间间隔称为一个太阳日，即 $24\,\mathrm h$。某种作用改变地球公转轨道半径，但不改变地球自转周期。若新的一年恰为 $360$ 个新太阳日，求（1）一个新太阳日是多少原小时；（2）新的一年是多少原小时。

\begin{qualityfigure}
\begin{tikzpicture}[>=Latex,font=\small]
\fill (0,0) circle(3pt) node[below]{太阳};\draw[DeepBlue] (0,0) circle(2);
\fill (2,0) circle(2pt) node[right]{$E_1$};\fill (1.532,1.286) circle(2pt) node[above right]{$E_2$};
\draw[dashed] (0,0)--(2,0) (0,0)--(1.532,1.286);\draw[->,Gold] (1.5,0) arc(0:40:1.5) node[midway,right]{$\Omega D$};
\draw[->] (2,-.5) arc(-90:170:.5);\node at (4,-.2){自转与公转同向};
\end{tikzpicture}
\figcap{地球自转与日地连线的相对转角}\label{mech:04:104:q}
\end{qualityfigure}
\end{origexamplebox}
\Needspace{6\baselineskip}\begin{xesolution}
\begin{qualityfigure}
\begin{tikzpicture}[>=Latex,font=\small]
\fill (0,0) circle(3pt) node[below]{太阳};\draw[DeepBlue] (0,0) circle(2);
\fill (2,0) circle(2pt) node[right]{$E_1$};\fill (1.532,1.286) circle(2pt) node[above right]{$E_2$};
\draw[dashed] (0,0)--(2,0) (0,0)--(1.532,1.286);\draw[->,Gold] (1.5,0) arc(0:40:1.5) node[midway,right]{$\Omega D$};
\draw[->] (2,-.5) arc(-90:170:.5);\node at (4,-.2){自转与公转同向};\node[align=center] at (0,-3){一天内自转角 $2\pi+\Omega D$\\相对日地连线转角为 $2\pi$};
\end{tikzpicture}
\figcap{地球自转与日地连线的相对转角（解答用辅助图）}\label{mech:04:104:a}
\end{qualityfigure}
图~\ref{mech:04:104:a}中，$E_1,E_2$ 表示相邻太阳日同一时刻的地球位置。日地连线也在公转中转过 $\Omega D$，因此地球自转必须额外多转这段角度。图中角度为辨认而放大。

设不变的恒星自转周期为 $T_r$，原太阳日 $D=24\,\mathrm h$，原恒星年 $Y=365.2564D$。地球公转、自转同向，一天内地面子午线必须比日地连线多转一整圈，太阳才回到同一方位。所以 $(2\pi/T_r-2\pi/Y)D=2\pi$，约去 $2\pi D$ 得
\[
\frac1D=\frac1{T_r}-\frac1Y.
\]
由此
\[
T_r=\left(\frac1{24\,\mathrm h}+\frac1{365.2564\times24\,\mathrm h}\right)^{-1}
\approx23.934472\,\mathrm h.
\]
新年 $Y'=360D'$，且
\[
\frac1{D'}=\frac1{T_r}-\frac1{Y'}
=\frac1{T_r}-\frac1{360D'}.
\]
整理得 $D'=(361/360)T_r$。因此
\[
\boxed{D'\approx24.000957\,\mathrm h},
\qquad
\boxed{Y'=360D'\approx8640.344\,\mathrm h}.
\]
结果中的“361”来自一年内地球自转圈数比太阳日数多一圈。
\end{xesolution}


% FinalReviewQuestionSplit20261024 火星的顺行与逆行时间
\tcbset{breakable/.default=true}
\begin{origexamplebox}{火星的顺行与逆行时间}\label{ch04:ex:k11-2-105}
太阳系内行星的公转方向基本相同，地球观测者看到的行星位置是该行星在遥远背景星空中的投影。由于各行星公转速度和轨道位置不同，地球上看到的视运动方向并不总与地球公转方向相同。视运动方向与地球公转方向相同称为“顺行”，相反称为“逆行”；从顺行转为逆行或从逆行转为顺行的转折位置，行星看起来短暂停在星空中，称为“留”。地外行星、地球和太阳位于同一直线上，且太阳与地外行星分居地球两侧时，称为“冲日”。

把地球和火星近似为在同一平面内绕太阳作同向圆周运动，火星轨道半径为地球的 $1.524$ 倍，不计地球自转。求相邻两次火星冲日之间，火星视运动逆行和顺行的时间。
\par 题图见图~\ref{ch04:ref:question-figure-2-105}。
\KQuestionEB
\end{origexamplebox}
\tcbset{breakable/.default=false}
\Needspace{6\baselineskip}\begin{xesolution}
以下推导所用的解图为图~\ref{ch04:fig:k11-ex-2-105}。

记 $k=R_M/R_E=1.524$。下面使用绕同一中心天体的圆轨道关系 $T^2\propto R^3$，即开普勒第三定律；其动力学依据见后续万有引力部分。先看七个时刻的轨道与视方向：
\MarsRetrogradeDiagram
角速度与周期成反比，所以
\[
q=\frac{\omega_M}{\omega_E}=k^{-3/2}\approx0.53152.
\]
相邻两次冲日的会合周期
\[
T_{\rm syn}=\frac{2\pi}{\omega_E-\omega_M}
=\frac{T_E}{1-q}\approx779.67\,\mathrm d.
\]

设火星与地球的日心方位差为 $\delta$。火星相对地球的位置为 $\vect r=\vect r_M-\vect r_E$。
视线方向角 $\lambda$ 的导数为 $\dot\lambda=(\vect r\times\dot{\vect r})_z/r^2$，它的正负决定顺行或逆行。为算叉积，可在考察的那一瞬间把地球日心方位取为零，则
\begin{align*}
\vect r/R_E&=(k\cos\delta-1,k\sin\delta),\\
\dot{\vect r}/R_E&=(-k\omega_M\sin\delta,k\omega_M\cos\delta-\omega_E).
\end{align*}
两分量交叉相乘并相减：
\begin{align*}
\frac{(\vect r\times\dot{\vect r})_z}{R_E^2}
&=(k\cos\delta-1)(k\omega_M\cos\delta-\omega_E)
+k^2\omega_M\sin^2\delta\\
&=k^2\omega_M+\omega_E-k(\omega_E+\omega_M)\cos\delta.
\end{align*}
也就是
\[
(\vect r\times\dot{\vect r})_z
=R_E^2\bigl[k^2\omega_M+\omega_E-k(\omega_E+\omega_M)\cos\delta\bigr]
\]
其中 $r^2$ 为正，所以“留”满足上述分子为零：
\[
\cos\delta_0=\frac{k^2q+1}{k(1+q)}\approx0.95735,
\quad \delta_0\approx16.793^\circ.
\]
冲日前后 $|\delta|<\delta_0$ 为逆行。因此
\[
\boxed{T_{\text{逆}}=T_{\rm syn}\frac{\delta_0}{\pi}\approx72.74\,\mathrm d},
\qquad
\boxed{T_{\text{顺}}=T_{\rm syn}-T_{\text{逆}}\approx706.93\,\mathrm d}.
\]
这里 $\delta_0$ 代入时间公式时用弧度。两次“留”之间日心相对角变化 $2\delta_0$，除以相对角速度 $\omega_E-\omega_M$ 就得到逆行时长。

图中的七个时刻在日心轨道、背景视方向与 $\lambda(t)$ 曲线上逐一对应；第3、5点是两次“留”，第4点是冲日。比较背景位置时要比较视线方向，不能把有限距离处的直线当作无限远星空。严格同平面模型的视方向在黄道同一大圆上往返，逆行段不是二维闭合圈。带有轨道倾角的三维模型还会产生视纬度变化，因而可能形成通常所说的回环；下方单独展示这种情况，不把它当作同平面计算的结果。
\MarsInclinedDiagram
\end{xesolution}


\begin{origexamplebox}{由木卫I观测周期估算光速}\label{ch04:ex:k11-2-106}
历史上曾有人比较地球位于公转轨道不同位置时测得的木星卫星周期，由此首次求得光的传播速度。已知木卫I真实周期 $T_0=42.5\,\mathrm h$；在地球轨道各处测得的所有周期中，最大值比 $T_0$ 多 $15\,\mathrm s$，最小值比 $T_0$ 少 $15\,\mathrm s$。近似认为地球、木星及木卫I都在同一平面内的圆轨道上运动：$R_E=1.5\times10^8\,\mathrm{km}$，$R_J=7.8\times10^8\,\mathrm{km}$，木卫I轨道半径 $R_I=4.2\times10^5\,\mathrm{km}$。（1）分析测得周期最大、最小时地球在公转轨道上的大致位置，以及地球、木星和木卫I的运动对测量周期的影响；（2）作合理近似，利用给定数据估算光速。

\begin{qualityfigure}
\begin{tikzpicture}[>=Latex,font=\small]
\draw[DeepBlue] (0,0) circle(1.7);\fill (0,0) circle(3pt) node[below]{太阳};
\fill (.578,1.5987) circle(2pt) node[above]{$E$};\fill (5,0) circle(3pt) node[right]{木星};
\draw[dashed] (.578,1.5987)--(5,0) node[midway,above]{视线};\draw[->,Gold] (.578,1.5987)--(1.7,1.193) node[below]{$\vect v_E$};
\draw[thin] (5,0) circle(.45);\node at (5,-.8){木卫I轨道};\node at (1,-2.5){距离变化使相邻信号的传播时间不同};
\end{tikzpicture}
\figcap{地球轨道切向速度和木星视线}\label{mech:04:106:q}
\end{qualityfigure}
\end{origexamplebox}
\Needspace{6\baselineskip}\begin{xesolution}
以下推导所用的解图为图~\ref{ch04:fig:k11-ex-2-106}。

连续两个同相位事件之间，地球—木星距离若增加 $\Delta D$，后一信号多传播 $\Delta D/c$，观测周期就变长；距离减小时观测周期变短。极值出现在地球相对木星的径向速度最大附近，即地球—木星视线近似与地球轨道相切的位置，而不是合日或冲日。木星也在公转，但因 $R_J\gg R_E$ 且其公转较慢，视线方向上的修正较小；木卫I在一个真实周期后回到同一相位，其轨道位移一阶抵消。$R_I$ 给出的卫星速率约 $17.2\,\mathrm{km/s}$，主要影响周期内的瞬时位置，不主导周期极值。
\RomerDiagram

地球轨道速率
\[
v_E=\frac{2\pi R_E}{365.2564\times86400\,\mathrm s}
\approx29.86\,\mathrm{km/s}.
\]
在上述极值位置取 $|\Delta D|\simeq v_ET_0$，而 $|\Delta T_{\rm obs}|=15\,\mathrm s$，于是
\[
c\simeq\frac{v_ET_0}{15\,\mathrm s}
=\frac{29.86\times(42.5\times3600)}{15}\,\mathrm{km/s}
\approx\boxed{3.05\times10^5\,\mathrm{km/s}}.
\]
这一层是忽略木星公转的数量级估算。进一步保留地球和木星的公转，令 $\delta$ 为地球相对木星的日心方位角，则
\[
D^2=R_J^2+R_E^2-2R_ER_J\cos\delta,\qquad
\dot D=\frac{R_ER_J(\omega_E-\omega_J)\sin\delta}{D}.
\]
为求 $|\dot D|$ 最大值，可令 $c_\delta=\cos\delta$，对 $(1-c_\delta^2)/(R_J^2+R_E^2-2R_ER_Jc_\delta)$ 求导。约去正因子后，极值条件为
\[
R_ER_Jc_\delta^2-(R_J^2+R_E^2)c_\delta+R_ER_J=0.
\]
两根为 $R_E/R_J$ 和 $R_J/R_E$；后者大于 $1$，舍去。有效位置满足
\[
\cos\delta=\frac{R_E}{R_J},\qquad
\delta\approx78.91^\circ\ \text{或}\ 281.09^\circ,
\qquad |\dot D|_{\max}=R_E(\omega_E-\omega_J).
\]
图中上侧 $E_+$ 处 $\dot D>0$，观测周期最长；下侧 $E_-$ 处 $\dot D<0$，观测周期最短。这两处 $SE$ 与 $EJ$ 垂直，确实是切点。

用开普勒第三定律估计 $\omega_J/\omega_E=(R_E/R_J)^{3/2}\approx0.0843$，同样由周期偏差 $15\,\mathrm s$ 得
\[
c\simeq\frac{R_E(\omega_E-\omega_J)T_0}{15\,\mathrm s}
\approx2.79\times10^5\,\mathrm{km/s}.
\]
保留木星公转使前一估计降低约 $8.4\%$，不能只称为很小的尾数修正。题给周期差本身是近似数据，不能为了更贴近标准光速而选择近似层级。木卫I半径相对地木距离很小，且同一轨道相位事件间其相对木星的位置重复；因此估计中不应把卫星速率 $2\pi R_I/T_0$ 直接加到长期地木径向速度上。
\end{xesolution}

\begin{origexamplebox}{用两次试射修正炮口角}\label{ch04:ex:k11-2-107}
第一次以 $30^\circ$ 发射，落点比目标远 $20\,\mathrm m$；第二次以 $29^\circ$ 发射，落点比目标近 $30\,\mathrm m$。确定再次发射的角度，使炮弹尽可能命中目标。

\begin{qualityfigure}
\begin{tikzpicture}[>=Latex,font=\small]
\draw[->] (0,0)--(5,0) node[right]{$\theta$};\draw[->] (0,-1.7)--(0,1.7) node[above]{落点误差};
\fill (1,-1.2) circle(2pt) node[left]{$-30\,\mathrm m$};\fill (4,0.8) circle(2pt) node[right]{$+20\,\mathrm m$};
\draw[dashed] (1,0) node[above]{$29^\circ$}--(1,-1.2);\draw[dashed] (4,0) node[below]{$30^\circ$}--(4,.8);
\end{tikzpicture}
\figcap{两次试射误差与零误差角度}\label{mech:04:107:q}
\end{qualityfigure}
\end{origexamplebox}
\Needspace{6\baselineskip}\begin{xesolution}
\begin{qualityfigure}
\begin{tikzpicture}[>=Latex,font=\small]
\draw[->] (0,0)--(5,0) node[right]{$\theta$};\draw[->] (0,-1.7)--(0,1.7) node[above]{落点误差};
\fill (1,-1.2) circle(2pt) node[left]{$-30\,\mathrm m$};\fill (4,0.8) circle(2pt) node[right]{$+20\,\mathrm m$};
\draw[dashed] (1,0) node[above]{$29^\circ$}--(1,-1.2);\draw[dashed] (4,0) node[below]{$30^\circ$}--(4,.8);\draw[DeepBlue,thick] (1,-1.2)--(4,.8);\fill[Gold] (2.8,0) circle(2pt) node[above left]{$\theta_*$};
\end{tikzpicture}
\figcap{两次试射误差与零误差角度（解答用辅助图）}\label{mech:04:107:a}
\end{qualityfigure}
图~\ref{mech:04:107:a}中，两个实测点分别位于零误差轴的两侧；连接两点后的横轴交点给出线性插值角度，横向距离按误差的 $30:20$ 分配。

在 $1^\circ$ 小区间内把落点误差对角度作线性插值：从 $29^\circ$ 的 $-30\,\mathrm m$ 增加到 $30^\circ$ 的 $+20\,\mathrm m$，零点位于
\[
29^\circ+\frac{30}{30+20}\times1^\circ
=\boxed{29.6^\circ}.
\]
若假设炮口速度不变、发射与落地点等高，射程 $S=K\sin2\theta$。由
\[
K\sin60^\circ=D+20,
\qquad K\sin58^\circ=D-30
\]
先求 $K,D$，再令 $K\sin2\theta=D$，得到 $\theta\approx29.593^\circ$，与线性插值一致。因此实际设定 $29.6^\circ$ 已是合理最优值。
\end{xesolution}

\begin{origexamplebox}{步枪准线与枪管的夹角}\label{ch04:ex:k11-2-108}
步枪的瞄准遵循“三点一线”原理：当缺口、准星和目标三点连成一线时击发。子弹速度虽然很大，通常约为 $800$--$1000\,\mathrm{m/s}$，但子弹在飞行中仍作抛体运动，因此必须在射击前校准步枪。校准的结果是使缺口与准星的连线并不与枪管平行，而是保留一个很小的夹角。忽略空气阻力，子弹初速率为 $v_0$，求为命中水平距离 $L$、与枪口等高的目标，瞄准线与枪管应有多大夹角 $\alpha$。

\begin{qualityfigure}
\begin{tikzpicture}[>=Latex,font=\small]
\fill (0,0) circle(2pt) node[below]{$O$（枪口）};\fill (6,0) circle(2pt) node[below]{$P$（目标）};
\draw[dashed] (0,0)--(6,0) node[midway,below]{水平瞄准线};
\draw[->,very thick,DeepBlue] (0,0)--(2,1) node[above]{$v_0$（枪管方向）};\draw (.8,0) arc(0:26.565:.8) node[midway,right]{$\alpha$};
\draw[<->] (0,-.7)--(6,-.7) node[midway,below]{$L$};
\end{tikzpicture}
\figcap{枪管方向、水平准线和同高落点}\label{mech:04:108:q}
\end{qualityfigure}
\end{origexamplebox}
\Needspace{6\baselineskip}\begin{xesolution}
\begin{qualityfigure}
\begin{tikzpicture}[>=Latex,font=\small]
\fill (0,0) circle(2pt) node[below]{$O$（枪口）};\fill (6,0) circle(2pt) node[below]{$P$（目标）};
\draw[dashed] (0,0)--(6,0) node[midway,below]{水平瞄准线};
\draw[->,very thick,DeepBlue] (0,0)--(2,1) node[above]{$v_0$（枪管方向）};\draw (.8,0) arc(0:26.565:.8) node[midway,right]{$\alpha$};
\draw[<->] (0,-.7)--(6,-.7) node[midway,below]{$L$};\draw[Gold,thick,domain=0:6,samples=60] plot(\x,{.5*\x-\x*\x/12});
\end{tikzpicture}
\figcap{枪管方向、水平准线和同高落点（解答用辅助图）}\label{mech:04:108:a}
\end{qualityfigure}
图~\ref{mech:04:108:a}中，$OP$ 为水平瞄准线，初速度沿枪管与它夹角为 $\alpha$。命中点 $P$ 与枪口等高，应在抛体轨迹中使用同高射程条件，并取实际步枪所对应的小角分支。

瞄准线指向目标，近似水平；枪管需略向上抬高 $\alpha$。同高斜抛射程条件
\[
L=\frac{v_0^2}{g}\sin2\alpha
\]
给出
\[
\boxed{\alpha=\frac12\arcsin\frac{gL}{v_0^2}}.
\]
实数命中角要求 $gL/v_0^2\le1$。射程方程通常还有高抛分支，本题步枪校准取接近水平的小角分支。步枪中 $gL/v_0^2\ll1$，故
\[
\boxed{\alpha\simeq\frac{gL}{2v_0^2}}.
\]
角度很小但不能取零；否则子弹到达目标时已因重力下落。
\end{xesolution}


\begin{origexamplebox}{击中墙上水平线的瞄准点轨迹}\label{ch04:ex:k11-2-109}
水平地面上有一面竖直墙，墙上离地 $h$ 处有一条水平直线。投掷点离地高度为 $z_0$、距墙 $L$，人以固定初速率 $v_0$ 投掷物体，要击中这条直线。求相应瞄准点在墙面上构成的轨迹。

\begin{qualityfigure}
\begin{tikzpicture}[>=Latex,font=\small]
\draw[fill=PaleBlue] (1,-.5)--(6,-.5)--(6,3.7)--(1,3.7)--cycle;
\draw[thick,DeepBlue] (1,1.3)--(6,1.3) node[right]{目标线 $z=h$};
\fill (-1,0) circle(2pt) node[left]{$O(z_0)$};\fill (4,1.3) circle(2pt) node[below right]{$P$};\fill (4,2.8) circle(2pt) node[above right]{$Q$};
\draw[dashed] (-1,0)--(4,2.8);\draw[<->] (-1,-1)--(1,-1) node[midway,below]{$L$};\node at (3.5,3.35){墙面};
\end{tikzpicture}
\figcap{墙面上的实际命中点与准线延长交点}\label{mech:04:109:q}
\end{qualityfigure}
\end{origexamplebox}
\Needspace{6\baselineskip}\begin{xesolution}
以下推导所用的解图为图~\ref{ch04:fig:k11-ex-2-109}。

设投掷点 $O$ 离地高度为 $z_0$，先保留这一参数，最后再给出 $z_0=0$ 的特例。在墙面内取水平方向坐标 $x$，$x=0$ 对应 $O$ 在墙上的法向投影；瞄准点 $Q$ 的离地高度为 $z$，实际命中点 $P$ 的离地高度为 $h$。定义相对高度 $Z=z-z_0$、$H=h-z_0$。
\AimLocusDiagram
墙面俯视的水平直角三角形给出 $OP$ 的水平投影长度
\[
r=\sqrt{L^2+x^2}.
\]
初速度方向沿 $OQ$，故 $Z=r\tan\theta$。水平位移 $r=v_0\cos\theta\,t$，所以 $t=r/(v_0\cos\theta)$。把它代入竖直位移 $H=v_0\sin\theta\,t-gt^2/2$：
\[
H=r\tan\theta-\frac{gr^2}{2v_0^2\cos^2\theta}
=Z-\frac{g}{2v_0^2}(r^2+Z^2).
\]
其中使用 $1/\cos^2\theta=1+\tan^2\theta=1+Z^2/r^2$。代入 $r^2=L^2+x^2$ 得 $x^2+Z^2-2v_0^2Z/g+L^2+2v_0^2H/g=0$，配方：
\[
\boxed{x^2+\left(z-z_0-\frac{v_0^2}{g}\right)^2
=\frac{v_0^4}{g^2}-L^2-\frac{2v_0^2(h-z_0)}{g}}.
\]
右端记作 $\mathcal D$，各项量纲均为长度平方。$\mathcal D>0$ 时，瞄准点轨迹是墙面上的圆，图中圆心 $K$ 的坐标为 $(0,z_0+v_0^2/g)$，半径为
\[
\boxed{R_{\rm aim}=\sqrt{\frac{v_0^4}{g^2}-L^2-\frac{2v_0^2(h-z_0)}{g}}}.
\]
$\mathcal D=0$ 时只有圆心这一个瞄准点；$\mathcal D<0$ 时给定速率不能到达墙上的命中线。若采用出手点在地面的理想模型，令 $z_0=0$，便得到特例
\[
x^2+(z-v_0^2/g)^2=v_0^4/g^2-L^2-2v_0^2h/g.
\]
图中实际命中点 $P$ 总在水平线 $z=h$ 上，瞄准点 $Q$ 则在圆上变化，二者不能混同。在不限制投掷方向、墙无限延伸的理想模型中，完整圆均对应可能方向；若只许向上抛，则需进一步取 $z>z_0$ 的圆弧，墙高、开口等限制也要另取交集。$h,z_0\ge0$ 时抛物线在两端点连线的上方，因而在到墙之前不会先穿入地面。
\end{xesolution}
}

\endgroup

\begingroup
\tcbset{breakable/.default=false}
\chapter{力、受力分析与静力平衡}

{
\graphicspath{{/Users/dongxingxi/Documents/物理竞赛一轮资料/sample/xes-original/first-section/}}
\section{力的基本性质}
同学们可能有所耳闻，自然界中有四种基本的相互作用力，它们分别是万有引力、电磁力、
强和弱相互作用力。力（force）是物体之间相互作用的一种表现形式，当有力作用在任何一个
物体上时，它将会对力的作用产生响应，响应的形式有很多，其中最明显的就是运动状态或
形状的改变。我们将首先学习力的基本的一面，即如何对力进行测量，如何描写物质系统的
受力，在已知力作用下处于静止状态力学系统的性质。在随后的课程中将逐步深入，研究在
已知的力作用下如何预测物体未来的运动的方法，从力的宏观性质出发进入到物质世界的深处
进一步探索力的起源。


对宏观物体受力的研究是学习描写各种力最初的一步，将来可以看到，这里所学的技术对于
那些无法直观体验的对象之间相互作用的问题提供了丰富的术语和方法。为了完整描写宏观
物体所受的一个力，需要指出该力的大小、方向和作用在物体上的位置，也就是作用点，它们
合起来称为力的三要素。从更基本的观点来看，力可以用空间中的一个矢量来描写，力的三要素
一起给出了作为一个矢量的大小、方向和起点的位置。方向和作用点由传统的几何方法给出，
而力的大小则需要引入一个新的单位，为了纪念力学的奠基人，今天将力的单位取为牛顿，
用 N 来表示，将在动力学中给出其准确含义。

虽然从最基本的角度来看自然界只有四种相互作用力，而且强、弱相互作用只存在于亚原子
物理当中，也就是说宏观物体的受力全部起源于万有引力和电磁相互作用力，原则上可以根据
它们的性质以及物质的结构从微观理论出发预测出各种情况下宏观物体的受力。但是实践上以
这种方法得到宏观物体的受力实际上是不现实的，因为物质世界极大的复杂性。所以在实际的
研究过程中首先根据力的表现形式和效果对其进行分类，研究力和那些已知的、容易观察的物理量
之间的关系，最后逐步地深入探讨力的本质和起源。基于经验和观察人们发现了几种常见的力，
下面我们逐一了解一下它们。

\clearpage
\subsection{重力}

本质上讲重力是万有引力在地球表面附近的一种表现形式，它来自于地球对所有物体的吸引。
观测表明，处于地球表面附近的物体都会受到重力作用，在地球的任何地方，其方向均竖直
向下指向地面，大小正比于物体的质量，作用于物体的重心处。关于重力有几点需要说明的地方：

\begin{itemize}
  \item 将地表任何一点处物体所受重力大小 $G$ 与其质量 $m$ 的比例系数称为地表的重力
  加速度，用 $g$ 表示
  \begin{equation}
    G=mg,\qquad g\approx\SI{9.8}{\metre\per\second\squared}.
  \end{equation}

  \item 真实测量表明，地表各处的 $g$ 并不完全相同。纬度、海拔、地球形状和自转
  都会带来修正；在大多数近地竞赛题中，这些差异可以忽略。

  \item 近地问题中把重力方向定义为当地铅垂向下。它一般不严格指向地心，但在小尺度问题中
  可视为方向和大小都不变的均匀场。

  \item 在均匀重力场中，各部分所受重力的合力作用线通过质心，因而重心与质心重合。
  如果物体可以看成
  是由多个质点构成，将这些质点编号，第 $i$ 个质点的质量为 $m_i$，在固定的直角坐标系中
  的位置为 $(x_i,y_i,z_i)$ 时，质心的坐标 $(x_c,y_c,z_c)$ 为
  \begin{equation}
    x_c=\frac{\sum_i m_ix_i}{\sum_i m_i},\qquad
    y_c=\frac{\sum_i m_iy_i}{\sum_i m_i},\qquad
    z_c=\frac{\sum_i m_iz_i}{\sum_i m_i}.
  \end{equation}
\end{itemize}

\begin{origexamplebox}{}\label{ch05:ex:s01-1-1}
求一块三边长分别为 $a,b,c$ 的均匀三角板的重心位置。

\begin{center}
\begin{tikzpicture}[scale=0.9]
  \coordinate (A) at (0.6,2.7);
  \coordinate (B) at (4.5,0.2);
  \coordinate (C) at (-0.4,0);
  \draw[thick] (A)--node[above right] {$c$}(B)--node[below] {$a$}(C)--node[above left] {$b$}cycle;
  \node[above] at (A) {$A$};
  \node[right] at (B) {$B$};
  \node[left] at (C) {$C$};
\end{tikzpicture}
\end{center}
\end{origexamplebox}

\Needspace{6\baselineskip}\begin{xesolution}
取三角形三个顶点的位置矢量分别为 $\vec r_A,\vec r_B,\vec r_C$。先作一条中线，例如从
$A$ 连到 $BC$ 的中点 $D$。把三角板分成平行于 $BC$ 的许多极细长条，每个长条的质心都在
$AD$ 上，因而整块板的质心也在 $AD$ 上。同理，它还在其他中线上，所以质心就是三条中线的交点 $G$。

\begin{center}
\begin{tikzpicture}[scale=0.72]
  \coordinate (A) at (0.6,2.7); \coordinate (B) at (4.5,0.2); \coordinate (C) at (-0.4,0);
  \coordinate (D) at ($(B)!0.5!(C)$); \coordinate (E) at ($(C)!0.5!(A)$);
  \coordinate (F) at ($(A)!0.5!(B)$); \coordinate (G) at (barycentric cs:A=1,B=1,C=1);
  \draw[thick] (A)--(B)--(C)--cycle;
  \draw[guide] (A)--(D) (B)--(E) (C)--(F);
  \fill[Ink] (D) circle (1.3pt) node[below] {$D$};
  \fill[DeepBlue] (G) circle (2pt)
    node[above=6pt,xshift=4pt,inner sep=0pt] {$G$};
  \node[above] at (A) {$A$}; \node[right] at (B) {$B$}; \node[left] at (C) {$C$};
\end{tikzpicture}
\end{center}

中线交点将每条中线按从顶点起 $2:1$ 分割，等价地
\[
  \boxed{\vec r_G=\frac{\vec r_A+\vec r_B+\vec r_C}{3}}.
\]
边长 $a,b,c$ 只决定三角形的形状，不改变“重心为中线交点”这一结论。
\end{xesolution}

\begin{origexamplebox}{}\label{ch05:ex:s01-1-2}
求由三根均匀杆构成的三角形的重心位置，其中三杆长度分别为 $a,b,c$，满足
$a^2+b^2=c^2$。分别用 $\lambda_a,\lambda_b,\lambda_c$ 表示各杆的线密度，先求一般结果，
再讨论三杆线密度相同的情形。

\begin{center}
\begin{tikzpicture}[scale=0.9]
  \coordinate (C) at (0,0);
  \coordinate (B) at (0,2.6);
  \coordinate (A) at (4.2,0);
  \draw[very thick] (B)--node[left] {$a$}(C)--node[below] {$b$}(A)--node[above right] {$c$}(B);
  \draw[thin] (0,0.34)--(0.34,0.34)--(0.34,0);
  \node[above] at (B) {$B$};
  \node[below left] at (C) {$C$};
  \node[right] at (A) {$A$};
\end{tikzpicture}
\end{center}
\end{origexamplebox}

\Needspace{6\baselineskip}\begin{xesolution}
“均匀”表示每根杆内部的线密度不随位置变化，并不表示三根杆的线密度一定相等。
取直角顶点 $C$ 为原点，$CA$ 沿 $x$ 轴、$CB$ 沿 $y$ 轴。
三根杆的质量及各自中点为
\[
\begin{array}{c|c|c}
\text{杆} & \text{质量} & \text{质心坐标}\\ \hline
BC & \lambda_a a &(0,a/2)\\
CA & \lambda_b b &(b/2,0)\\
AB & \lambda_c c &(b/2,a/2)
\end{array}
\]
记总质量为 $M=\lambda_a a+\lambda_b b+\lambda_c c$。各杆的质量乘其质心坐标后相加，得到
\[
 Mx_G=\lambda_a a\cdot0+\lambda_b b\cdot\frac b2+\lambda_c c\cdot\frac b2,
\]
\[
 My_G=\lambda_a a\cdot\frac a2+\lambda_b b\cdot0+\lambda_c c\cdot\frac a2.
\]
所以一般结果为
\[
 \boxed{x_G=\frac{b(\lambda_b b+\lambda_c c)}{2M},\qquad
 y_G=\frac{a(\lambda_a a+\lambda_c c)}{2M}}.
\]
只有进一步给定 $\lambda_a=\lambda_b=\lambda_c=\lambda$，才可约去共同线密度，得到
\[
  \boxed{x_G=\frac{b(b+c)}{2(a+b+c)},\qquad
  y_G=\frac{a(a+c)}{2(a+b+c)}} ,
\]
其中 $c=\sqrt{a^2+b^2}$。这里不能套用“三角板的中线交点”：面密度均匀与线密度均匀是两个不同的质量分布。
式中分子量纲为质量乘长度、分母为质量，坐标的量纲正确。若某一杆的质量远大于另外两杆，
一般式便趋于该杆中点的坐标，也符合质心的含义。
\end{xesolution}

\begin{origexamplebox}{}\label{ch05:ex:s01-1-3}
求下面阴影的重心。

\begin{center}
\begin{tikzpicture}[scale=0.92]
  \coordinate (O) at (0.70,0);
  % Leave a visible annular gap: the sketch has R=2.20, r=0.65, l=1.20.
  \coordinate (Op) at (-0.50,0);
  \fill[TextGray!35,even odd rule] (O) circle (2.20) (Op) circle (0.65);
  \draw[thick] (O) circle (2.20);
  \draw[thick] (Op) circle (0.65);
  \fill[Ink] (O) circle (1.8pt) node[below right=2pt] {$O$};
  \fill[Ink] (Op) circle (1.8pt)
    node[below=3pt,inner sep=0pt] {$O'$};
  \draw[->,draw=Ink] (O)--++(45:2.20) node[midway,above left] {$R$};
  \draw[->,draw=Ink] (Op)--++(135:.65) node[midway,above right] {$r$};
  \draw[draw=Ink] (Op)--(O) node[pos=0.68,below=3pt] {$l$};
\end{tikzpicture}
\end{center}
\end{origexamplebox}

\Needspace{6\baselineskip}\begin{xesolution}
设薄片面密度为 $\sigma$，以大圆圆心 $O$ 为原点。把挖去的小圆看成位于 $O'$ 的“负质量”，则
\[
  M=\sigma\pi R^2,\qquad m=\sigma\pi r^2.
\]
剩余部分的质量为 $M-m$，对 $O$ 取质量矩，
\[
  (M-m)\vec r_G=M\vec 0-m\,\overrightarrow{OO'}.
\]
所以
\[
  \boxed{\vec r_G=-\frac{r^2}{R^2-r^2}\,\overrightarrow{OO'}}.
\]
\begin{center}
\begin{tikzpicture}[scale=0.62]
  \coordinate (O) at (0,0);
  \coordinate (Op) at (-0.92,0);
  \coordinate (G) at (0.42,0);
  \fill[TextGray!30,even odd rule] (O) circle (1.75) (Op) circle (0.52);
  \draw[thick] (O) circle (1.75);
  \draw[thick] (Op) circle (0.52);
  \draw[guide] (Op)--(G);
  \fill[Ink] (Op) circle (1.6pt);
  \draw[thin,TextGray] (Op)--(-1.92,0);
  \node[left=2pt,inner sep=0pt] at (-1.92,0) {$O'$};
  \fill[Ink] (O) circle (1.6pt) node[below=3pt] {$O$};
  \fill[DeepBlue] (G) circle (2pt) node[above right=2pt] {$G$};
  \draw[-{Latex[length=2.2mm]},thick,draw=DeepBlue]
    ($(O)+(0,0.34)$)--($(G)+(0,0.34)$)
    node[midway,above] {$\vec r_G$};
  \node[font=\scriptsize,text=TextGray,anchor=west] at (1.95,0.35)
    {质心在小孔的反方向};
  \draw[guide] (1.85,0.20)--(0.55,0.06);
\end{tikzpicture}\\[-0.4ex]
{\scriptsize\color{TextGray}方向示意；为便于辨认，图中 $OG$ 的长度未按比例绘制。}
\end{center}
质心位于 $OO'$ 直线上、在小孔的反方向，与 $O$ 的距离为
$lr^2/(R^2-r^2)$。这一结果要求小圆孔完全位于大圆内，即 $l+r\le R$。
\end{xesolution}

\begin{origexamplebox}{}\label{ch05:ex:s01-1-4}
已知一个质量为 $M$ 的没有盖子的圆柱形的杯子。里面开始是空的，然后往里面慢慢地注水。
进入了杯子的水和杯子的总重心会发生变化。请问重心如何变化？是否有最值？是多少？
（已知杯子底面积为 $S$，水的密度为 $\rho$，原来空杯子重心高度为 $H$，杯内允许的
最大水深为 $h_{\max}$。）
\end{origexamplebox}

\Needspace{6\baselineskip}\begin{xesolution}
设水面高度为 $h$，从杯底向上量取坐标。水的质量为 $\rho Sh$，其质心高度为 $h/2$，
因而系统质心高度为
 \begin{equation}
  z(h)=\frac{MH+\frac12\rho Sh^2}{M+\rho Sh}.
  \label{ch05:eq:cup-centre-mass}
\end{equation}
\begin{center}
\begin{tikzpicture}[scale=0.68]
  \fill[PaleBlue] (0.06,0.06) rectangle (2.64,1.45);
  \draw[very thick] (0,3.10)--(0,0)--(2.70,0)--(2.70,3.10);
  \draw[thick,draw=DeepBlue] (0.06,1.45)--(2.64,1.45);
  \draw[guide] (1.35,0)--(1.35,3.10);
  \draw[<->,draw=Ink] (-0.36,0)--(-0.36,1.45)
    node[midway,left] {$h$};
  \fill[Ink] (1.35,0.725) circle (1.7pt);
  \draw[guide] (1.45,0.725)--(3.08,0.725);
  \node[anchor=west,inner sep=0pt] at (3.12,0.725)
    {水的质心：$h/2$};
  \draw[Ink,thick] (1.27,2.48)--(1.43,2.64) (1.27,2.64)--(1.43,2.48);
  \draw[guide] (1.45,2.56)--(3.08,2.56);
  \node[anchor=west,inner sep=0pt] at (3.12,2.56)
    {空杯重心：$H$};
  \fill[DeepBlue] (1.35,1.86) circle (2pt);
  \draw[guide] (1.45,1.86)--(3.08,1.86);
  \node[anchor=west,inner sep=0pt,text=DeepBlue] at (3.12,1.86)
    {总质心：$z(h)$};
\end{tikzpicture}\\[-0.4ex]
{\scriptsize\color{TextGray}某一注水状态示意；$z(h)$ 总在 $H$ 与 $h/2$ 之间，图中具体高度不代表一般大小关系。}
\end{center}
注水时新加入的一层水位于高度 $h$，由式 \eqref{ch05:eq:cup-centre-mass} 也可直接得到
\[
  \frac{\dd z}{\dd h}=\frac{\rho S}{M+\rho Sh}\,[h-z(h)].
\]
因此，当新加水层低于当前质心时，质心下降；高于当前质心时，质心上升。极小值满足 $z(h_*)=h_*$，即
\[
  \frac12\rho S h_*^2+Mh_*-MH=0.
\]
取正根得
\[
  \boxed{h_*=z_{\min}
  =\frac{\sqrt{M^2+2M\rho SH}-M}{\rho S}}.
\]
结合题设给出的最大水深，可达的最低重心为
\[
  \boxed{z_{\min,\text{可达}}=
  \begin{cases}
    h_*, & h_{\max}\ge h_*,\\[2pt]
    \displaystyle\frac{MH+\frac12\rho S h_{\max}^2}
    {M+\rho S h_{\max}}, & h_{\max}<h_*.
  \end{cases}}
\]
第二种情形中，$0\le h\le h_{\max}$ 内有 $h<z(h)$，所以 $z(h)$ 始终随注水下降，
最小值确在装到 $h_{\max}$ 时取得。
\end{xesolution}

\begin{origexamplebox}{【思考题】}\label{ch05:ex:s01-1-5}
已知一个均匀质量的圆球壳，半径为 $R$，现在将球壳上端高度为 $h$ 的部分切下，然后
粘贴在圆球的最下面，形成一个“杯子”。求这个杯子的重心。
\end{origexamplebox}

\Needspace{6\baselineskip}\begin{xesolution}
以原球心为原点，向上为 $z$ 轴。这里的“粘贴在圆球的最下面”是指：
球冠保持原来的朝向，不作翻转，整体向下平移，使球冠顶点与原球壳最低点粘接；
球冠的圆形边缘位于最下方，构成杯子的底座。

\begin{center}
\begin{tikzpicture}[scale=0.82]
  \begin{scope}[shift={(-3.30,0)}]
    % 取示意半径 1.5、球冠高度 0.5；切口在 y=1.0。
    \draw[thick]
      (-1.12,1.00) arc[start angle=138.2,end angle=401.8,radius=1.50];
    \draw[very thick,DeepBlue]
      (-1.12,1.00) arc[start angle=138.2,end angle=41.8,radius=1.50];
    \draw[guide] (-1.12,1.00)--(1.12,1.00);
    \draw[guide] (-1.62,0)--(1.62,0);
    \draw[guide] (0,-1.65)--(0,1.75) node[above=1pt]{$z$};
    \draw[<->] (-1.72,1.50)--(-1.72,1.00) node[midway,left]{$h$};
    \fill[Ink] (0,0) circle (1.5pt) node[below right]{$O$};
    \node[DeepBlue,font=\small] at (0,-1.90) {(1) 切下上球冠};
  \end{scope}
  \draw[-{Latex[length=2.3mm]},thick,draw=TextGray] (-1.25,0.05)--(1.25,0.05);
  \node[font=\scriptsize,align=center,inner sep=0pt] at (0,0.72)
    {不翻转\\整体下移 $2R$};
  \begin{scope}[shift={(3.30,0.45)}]
    % 原球壳上方开口；灰虚线仅标出球冠原来的位置。
    \draw[thick]
      (-1.12,1.00) arc[start angle=138.2,end angle=401.8,radius=1.50];
    \draw[TextGray,dashed]
      (-1.12,1.00) arc[start angle=138.2,end angle=41.8,radius=1.50];
    \draw[guide] (-1.12,1.00)--(1.12,1.00);
    \draw[guide] (-1.62,0)--(1.62,0);
    \draw[guide] (0,-2.30)--(0,1.70) node[above=1pt]{$z$};
    % 同一球冠向下平移 2R：顶点恰好接触原球壳最低点。
    \draw[very thick,DeepBlue]
      (-1.12,-2.00) arc[start angle=138.2,end angle=41.8,radius=1.50];
    % 圆形切口的透视示意：前半实线、后半虚线。
    \draw[DeepBlue,thick]
      (-1.12,-2.00) arc[start angle=180,end angle=360,x radius=1.12,y radius=0.12];
    \draw[DeepBlue,thick,densely dashed]
      (-1.12,-2.00) arc[start angle=180,end angle=0,x radius=1.12,y radius=0.12];
    \fill[Ink] (0,0) circle (1.5pt) node[below right]{$O$};
    \fill[DeepBlue] (0,-1.50) circle (1.7pt);
    \draw[TextGray,thin] (0,-1.50)--(1.68,-1.12)--(1.92,-1.12);
    \node[font=\scriptsize,anchor=west,inner sep=0pt]
      at (1.96,-1.12) {粘接点};
    \draw[DeepBlue,thin] (1.12,-2.00)--(1.48,-2.25);
    \node[DeepBlue,font=\scriptsize,anchor=west,inner sep=0pt]
      at (1.50,-2.25) {圆形支承边缘};
    \node[DeepBlue,font=\small] at (0,-2.82) {(2) 作为底座粘贴};
  \end{scope}
\end{tikzpicture}
\\[-0.4ex]
{\scriptsize\color{TextGray}蓝色曲线表示同一球冠；右图中球冠在原球壳外部，顶点与球壳最低点相接。}
\end{center}

设球壳面密度为 $\sigma$。高为 $h$ 的薄球冠面积为 $2\pi Rh$，而对球壳水平分层有
$\dd A=2\pi R\,\dd z$，所以球冠质心高度为
\[
  z_c=\frac{\int_{R-h}^{R}z\,\dd z}{\int_{R-h}^{R}\dd z}
  =R-\frac h2.
\]
球冠原来的顶点坐标为 $R$。平移后该顶点与原球壳最低点 $-R$ 重合，
所以整个球冠向下平移了 $2R$，其质心的新坐标为
\[
  z_c'=z_c-2R=-R-\frac h2.
\]
球冠质量 $m_c=2\pi Rh\sigma$，原球壳总质量 $M_0=4\pi R^2\sigma$。
取下原位置的球冠，再把它粘到新位置，系统质量不变；相对于原球壳，质量矩的改变量为
\[
  M_0z_G=-m_cz_c+m_cz_c'
  =m_c(z_c'-z_c)=-2Rm_c.
\]
因此
\[
  z_G=\frac{-2R(2\pi Rh\sigma)}{4\pi R^2\sigma}
  =\boxed{-h}.
\]
负号表示杯子的重心在原球心下方，距原球心为 $h$。这个结果也可直接理解为：
占总质量 $h/(2R)$ 的球冠整体下降了 $2R$，所以系统质心下降 $h$。
\end{xesolution}

\subsection{弹力}

当相互接触的物体发生形变时所产生的恢复形变的力称为弹性力，胡克定律表明，当物体
形变不太大时，弹性力与形变成正比，弹簧的弹性力 $F$ 与弹簧相对于原长的形变（拉伸
或压缩）$x$ 成正比。取沿弹簧轴的有符号形变量 $x$，弹力分量与形变反向，即
\begin{equation}
  F=-kx,
\end{equation}
式中比例系数 $k$ 称为弹簧的倔强系数，也叫劲度系数，负号表示弹性力与形变反方向。
这里形变为零的位置对应弹簧原长，不一定是物体受所有外力后的平衡位置。例如竖直弹簧
悬挂物体时，平衡伸长量为 $mg/k$，该处弹力仍向上且大小为 $mg$，并不为零。
关于弹力也有几点需要注意的地方：

\begin{itemize}
  \item 当把多根弹簧并排放置共同作用时，称两弹簧并联，并联弹簧在任意位置处形变量
  相同。这里取共同连接到刚性横杆上的理想轻弹簧，讨论沿弹簧轴的小形变：共同位移为 $x$，
  第 $i$ 根弹簧提供的力为 $F_i=k_i x$，故总力为
  \[
    F=\sum_iF_i=\left(\sum_i k_i\right)x.
  \]
  与 $F=kx$ 比较，多根弹簧并联的等效劲度系数为
  \begin{equation}
    k=k_1+k_2+\cdots,
  \end{equation}
  为每根弹簧劲度系数之和。和电路类比，除了并联以外多根弹簧之间也可以串联，处于串联
  状态的弹簧在受力以后形变量并不相同；连接点质量可忽略且无其他外力时，各弹簧的张力相同。
  把各段伸长量相加，可得串联弹簧的等效劲度
  系数为
  \begin{equation}
    \frac1k=\frac1{k_1}+\frac1{k_2}+\cdots.
  \end{equation}

  \item 绳子的张力是一种弹性力，绳子和与之连接的物体之间有相互作用时，不仅绳子与
  物体之间有弹性力，而且在绳子内部也因发生相对形变而出现弹性力。这时，绳子上任一
  横截面两边互施作用力，这对作用力和反作用力称为绳子的张力，一般情况下，与绳子相应
  的比例系数 $k$ 很大，因而形变很小，可以忽略。所以绳子的张力不是由绳子的形变规律
  确定，而是由求解力学问题时确定。因而我们一般抽象出柔软不可伸长的轻绳。

  \item 物体与支撑面接触时，微观形变会产生接触力。接触力的法向分量称为支持力 $N$；
  只在光滑面模型中，合接触力才必定沿曲面法线。若表面粗糙，还要加上沿切平面的摩擦力，
  此时合接触力通常不与表面垂直。支持力和摩擦力的大小不是预先给定的，而由运动状态与约束共同决定。

  \item 所有弹力的起源均为电磁相互作用力，这一点只有从微观角度才能够发现。
\end{itemize}

\begin{origexamplebox}{}\label{ch05:ex:s01-1-6}
证明：串联弹簧的等效弹性系数小于其中任何一根弹簧的弹性系数，并解释形成这一现象的原因。
\end{origexamplebox}

\Needspace{6\baselineskip}\begin{xesolution}
两根弹簧串联时张力相同，设为 $F$。各弹簧伸长量为 $x_i=F/k_i$，总伸长量为
\[
  x=\sum_i x_i=F\sum_i\frac1{k_i}.
\]
令 $F=k_{\rm eq}x$，得
\[
  \frac1{k_{\rm eq}}=\sum_i\frac1{k_i}.
\]
因为对任意一根弹簧 $j$ 都有
\[
  \frac1{k_{\rm eq}}=\frac1{k_j}+\sum_{i\ne j}\frac1{k_i}>\frac1{k_j},
\]
所以
\[
  \boxed{k_{\rm eq}<k_j\qquad(\text{对所有 }j)}.
\]
物理上，同一个力使每根弹簧都伸长，而串联组的总伸长是各伸长量之和；与任意一根单独弹簧相比，
同力下形变更大，等效劲度自然更小。
\end{xesolution}

\subsection{滑动摩擦力、静摩擦力}

摩擦力也是一种接触力，当相互接触的物体作相对运动或有相对运动趋势时，接触面间会
产生一种阻碍相对运动或相对运动趋势的力，这种力称为摩擦力，前者称为滑动摩擦力，
后者称为静摩擦力。测量结果表明：

\begin{enumerate}[label=\arabic*.]
  \item 在库仑摩擦的近似模型中，滑动摩擦力大小与法向支持力成正比，方向与接触面上的相对滑动速度相反：
  \begin{equation}
    \vec f_k=-\mu_kN\frac{\vec v_{\rm rel}}{|\vec v_{\rm rel}|},
  \end{equation}
  其中 $\mu_k$ 称为动摩擦因数。“与表观接触面积无关”和“$\mu_k$ 为常数”都是一定条件下的理想化，
  不是对所有材料和速度都精确成立的定律。

  \item 在高中竞赛的常规库仑模型中，滑动摩擦力大小取为 $\mu_kN$，不再显式依赖速度大小。

  \item 静摩擦力是一个由平衡或共同运动条件决定的约束力，其大小满足
  \begin{equation}
    0\le |f_s|\le \mu_sN.
  \end{equation}
  只有在临界滑动时才取等号。$\mu_s$ 是最大静摩擦因数，与 $\mu_k$ 应当区分；若题目只给一个
  $\mu$，需根据所讨论的阶段明确它表示哪一个摩擦因数。
\end{enumerate}

\begin{origexamplebox}{}\label{ch05:ex:s01-1-7}
如图所示，木板 $A$ 质量为 $M$，以相对地面的速度 $v$ 在水平面上向北运动，木板上放
一质量为 $m$ 的板 $B$，各接触面间滑动摩擦因数均为 $\mu$，当木块 $B$ 也有相对地面
向东的速度 $v$ 时，试分析木块 $B$ 的受摩擦力的情况。

\begin{center}
\begin{tikzpicture}[x=1.0cm,y=1.0cm]
  \draw[fill=PaleBlue,draw=Ink,thick] (-2.5,-1.55) rectangle (2.5,1.55);
  \draw[fill=white,draw=Ink,thick] (-0.65,-0.4) rectangle (0.65,0.4);
  \node at (0,0) {$B$};
  \node[anchor=north west] at (-2.30,1.34) {$A$};
  \draw[vector] (2.9,-0.45)--(2.9,1.15) node[above] {$v$};
  \draw[vector] (-0.15,-0.72)--(1.45,-0.72) node[right] {$v$};
  \node[right] at (3.35,0.7) {北};
  \node[below] at (1.25,-1.02) {东};
\end{tikzpicture}
\end{center}
\end{origexamplebox}

\Needspace{6\baselineskip}\begin{xesolution}
取向东为 $x$ 正方向、向北为 $y$ 正方向。板 $A$ 和块 $B$ 对地速度分别为
\[
  \vec v_A=v\,\hat{\boldsymbol j},\qquad
  \vec v_B=v\,\hat{\boldsymbol i}.
\]
因而 $B$ 相对 $A$ 的滑动速度为
\[
  \vec v_{B/A}=\vec v_B-\vec v_A
  =v(\hat{\boldsymbol i}-\hat{\boldsymbol j}),
\]
方向为东南。$A$ 对 $B$ 的滑动摩擦力必须阻碍这个相对滑动，所以指向西北。

\noindent\begin{minipage}[c]{0.42\linewidth}
\centering
\begin{tikzpicture}[scale=0.64]
  \draw[axis] (-2.2,0)--(2.3,0) node[right]{东};
  \draw[axis] (0,-1.8)--(0,2.1) node[above]{北};
  \draw[vector] (0,0)--(1.4,0) node[below right]{$\vec v_B$};
  \draw[vector] (0,0)--(0,1.4) node[above left]{$\vec v_A$};
  \draw[component] (0,0)--(1.4,-1.4) node[below right]{$\vec v_{B/A}$};
  \draw[vector] (0,0)--(-1.15,1.15) node[above left]{$\vec f_{A\to B}$};
\end{tikzpicture}
\end{minipage}\hfill
\begin{minipage}[c]{0.55\linewidth}
若 $B$ 在竖直方向无加速度，则 $N=mg$，故
\[
  \boxed{\vec f_{A\to B}
  =\frac{\mu_kmg}{\sqrt2}
  (-\hat{\boldsymbol i}+\hat{\boldsymbol j})},
  \qquad |\vec f_{A\to B}|=\mu_kmg.
\]
摩擦力的方向要根据\textbf{相对速度}判断，不能只看 $B$ 的对地速度。
\end{minipage}
\end{xesolution}
}

{
\graphicspath{{/Users/dongxingxi/Documents/物理竞赛一轮资料/sample/xes-original/second-section/}}
\section{受力分析}
当一个物体受多个力作用时，需要给定所有力的大小、方向和作用点才能够完全了解它的
受力情况。对质点，多个力可合成为一个力；对刚体，若各力的作用线通过同一点，称为共点力，
也可用通过该点的合力代替。其合力为几个力的矢量和：
\begin{equation}
  \vec{F}=\sum_i\vec{F}_i.
\end{equation}
一般刚体上的非共点力还可能产生力偶；此时仅知道矢量和不足以说明转动效果，必须同时
考虑合力矩。例如两个等大反向、作用线不重合的力，合力为零，却仍能使物体转动。

当在坐标系中分析受力时，将一个给定的力向着坐标轴的方向分解对处理问题也是有帮助
的。这里采用沿正交坐标轴的投影分解；沿某个正轴的分量是有符号的数，不能把它与非负的力的大小混为一谈。

若所选单位轴为 $\hat{\boldsymbol e}$，则力沿该轴的\textbf{有符号分量}可统一写成
\[
  F_e=\vec F\cdot\hat{\boldsymbol e}=F\cos\theta.
\]
当投影方向与正轴相反时，$F_e<0$。这个负号已包含方向信息，不应再人为补一个负号。
相应的分矢量是 $F_e\hat{\boldsymbol e}$，它的大小为 $|F_e|$。若使用非正交方向作斜分解，
则需联立矢量方程，不能把每个分力都直接写成余弦投影。

\begin{origexamplebox}{}\label{ch05:ex:s02-1-8}
如图所示，有五个力 $\vec{F}_1,\vec{F}_2,\vec{F}_3,\vec{F}_4,\vec{F}_5$ 作用于一点 $O$，
构成一个正六边形的两邻边和三条对角线。设 $F_3=\SI{10}{\newton}$，试求这五个力的合力。

\begin{center}
\begin{tikzpicture}[scale=0.64]
  \coordinate (O)  at (-3,0);
  \coordinate (F1) at (-1.5,2.598);
  \coordinate (F2) at (1.5,2.598);
  \coordinate (F3) at (3,0);
  \coordinate (F4) at (1.5,-2.598);
  \coordinate (F5) at (-1.5,-2.598);

  \draw[guide,thick] (F1)--(F2)--(F3)--(F4)--(F5);
  \foreach \P in {F1,F2,F3,F4,F5}{
    \draw[-{Latex[length=2.4mm]},very thick,draw=Ink] (O)--(\P);
  }

  \node[left] at (O) {$O$};
  \node[above left] at (F1) {$F_1$};
  \node[above right] at (F2) {$F_2$};
  \node[right] at (F3) {$F_3$};
  \node[below right] at (F4) {$F_4$};
  \node[below left] at (F5) {$F_5$};
\end{tikzpicture}
\end{center}
\end{origexamplebox}

\Needspace{6\baselineskip}\begin{xesolution}
图中各力的大小就是相应矢量的几何长度。沿 $F_3$ 方向取 $x$ 轴，正六边形的边长设为 $s$，则
\[
  F_3=2s.
\]
关于 $x$ 轴对称的两对力的竖直分量相互抵消；水平分量分别为
\[
  \vec F_1+\vec F_5=s\,\hat{\boldsymbol i}
  =\frac12\vec F_3,
\]
\[
  \vec F_2+\vec F_4=3s\,\hat{\boldsymbol i}
  =\frac32\vec F_3.
\]
因而
\[
  \vec F_{\rm sum}
  =\left(\frac12+\frac32+1\right)\vec F_3
  =3\vec F_3.
\]
已知 $F_3=\SI{10}{\newton}$，所以
\[
  \boxed{F_{\rm sum}=\SI{30}{\newton}},
\]
方向与 $\vec F_3$ 相同。
\end{xesolution}

找出力学系统中每个物体受力的大小、方向和作用点的过程称为受力分析，一般来说在
静力学中主要考虑上面提到过的几种典型的受力。除了特别的假设，地面附近所有物体都受
重力，发生形变的弹性物体有弹力作用，相互接触的表面之间有可能有压力，当有压力并且
有相对运动或相对运动趋势时还可能有摩擦力。除此之外，在有的问题中还可能出现大气压
力、浮力等其它种类的力，做受力分析的时候要把所有重要的因素全部考虑方可。

虽然力的效果在物理上是一致的，但不同的力也有细微的区别。在做受力分析时，可以
将力分为两大类，一类称做主动力，例如重力、弹簧弹力等。这些力的共同特点就是它们可
由系统的状态直接确定，只要知道了物体的质量、弹簧的弹性系数和形变量以后就可以确定
其大小和方向。另一类称约束力或被动力，例如支持力、两个物体连接处的作用力等，它们
无法直接确定，是由主动力和限制物体运动约束的性质共同决定。

从建立未知量的角度看，主动力通常可由相互作用律直接写出，例如近地重力 $mg$ 和线性
弹簧力 $-kx$；约束力则通常需要与平衡方程或运动方程联立求出。这一划分是建模工具，
不是对力的微观来源进行分类。

\begin{origexamplebox}{}\label{ch05:ex:s02-1-9}
\begin{minipage}[c]{0.62\linewidth}
均匀长棒一端搁在地面上，另一端用细线系在天花板上，如图所示，若细线竖直，
试分析棒的受力情况。
\end{minipage}\hfill
\begin{minipage}[c]{0.34\linewidth}
\centering

\begin{tikzpicture}[scale=0.48]
  \draw[thick] (-3.1,-1.7)--(3.0,-1.7);
  % Model the long rod by its centre line, so its lower endpoint has
  % single-point contact with the floor rather than face contact.
  \draw[draw=Ink,line width=2.2pt,line cap=round]
    (-1.76,-1.70)--(1.67,1.69);
  \draw[thick] (1.67,1.69)--(1.67,3.05);
  \draw[thick] (0.95,3.05)--(2.39,3.05);
  \foreach \x in {1.02,1.22,...,2.32}{
    \draw[thin] (\x,3.05)--++(0.22,0.22);
  }
\end{tikzpicture}
\end{minipage}
\end{origexamplebox}

\Needspace{6\baselineskip}\begin{xesolution}
隔离木棒。对图示的非竖直木棒，需要考虑中点处的重力 $mg$、上端细线的竖直张力 $T$、
下端地面的法向力 $N$ 和水平摩擦力未知量 $f$。图中灰色虚线只表示 $f$ 的预设正方向，
并不表示摩擦力一定非零。

\begin{center}
\begin{tikzpicture}[scale=0.68]
  \coordinate (A) at (-1.8,-1.2); \coordinate (B) at (1.6,1.8); \coordinate (C) at ($(A)!0.5!(B)$);
  \draw[very thick] (A)--(B);
  \draw[vector] (B)--++(0,1.25) node[above]{$T$};
  \draw[vector] (C)--++(0,-1.35) node[below]{$mg$};
  \draw[vector] (A)--++(0,1.2) node[left]{$N$};
  \draw[-{Latex[length=2.1mm]},thick,draw=TextGray,densely dashed]
    (A)--++(1.05,0) node[pos=0.72,below=3pt,inner sep=0pt]{$f$};
  \node[below left] at (A) {$A$}; \node[right] at (B) {$B$};
\end{tikzpicture}
\end{center}

水平方向只有 $f$，由平衡条件立即得
\[
  \boxed{f=0}.
\]
力矩衡量力使物体绕参考点转动的作用，其大小为力乘参考点到作用线的垂直距离（力臂）。
取逆时针为正，平衡时各力矩的代数和为零。设木棒与水平面夹角为 $\theta$\ ($\cos\theta\ne0$)，
对下端 $A$ 取矩：$N$ 和 $f$ 的作用线都通过 $A$，力臂为零；$T$ 和 $mg$ 均竖直，
它们的力臂分别为水平距离 $L\cos\theta$ 和 $(L/2)\cos\theta$，转动趋势相反。因此
\[
  T(L\cos\theta)-mg\left(\frac L2\cos\theta\right)=0,
\]
所以
\[
  \boxed{T=\frac{mg}{2}},\qquad
  \boxed{N=mg-T=\frac{mg}{2}}.
\]
灰色虚线仅标明 $f$ 的预设正方向；求得 $f=0$ 说明接触并不意味着摩擦力非零。
若木棒恰好竖直，三个竖直力共线，力矩方程不能分别确定 $T$ 与 $N$，只能得到 $T+N=mg$。
\end{xesolution}
}

{
\graphicspath{{/Users/dongxingxi/Documents/物理竞赛一轮资料/sample/xes-original/third-section/}}
\section{质点与刚体的静力平衡}
% TikZ diagrams written for the added solutions.  They expose the forces
% and geometric quantities used in each derivation and are visually checked
% against the independently redrawn problem figures.

\tikzset{
  solbody/.style={draw=Ink,line width=0.8pt},
  solforce/.style={-{Latex[length=2.0mm]},draw=DeepBlue,line width=0.9pt},
  solgeom/.style={draw=TextGray,densely dashed,line width=0.55pt},
  soldim/.style={{Latex[length=1.4mm]}-{Latex[length=1.4mm]},draw=TextGray,line width=0.55pt},
  solpoint/.style={circle,fill=Ink,inner sep=1.2pt},
  % Labels must not carry a rectangular white mask: place them in open
  % space and use anchors/leader lines to avoid the drawing instead.
  sollabel/.style={font=\footnotesize,text=Ink,fill=none,inner sep=1pt,outer sep=1pt},
  solpanel/.style={font=\footnotesize\bfseries,text=TextGray}
}

\newcommand{\soldiagcaption}[1]{%
  \par\vspace{-0.3em}{\centering\scriptsize\color{TextGray}解图：#1\par}\vspace{0.2em}%
}

\newcommand{\soldiagOneTen}{%
\begin{center}
\begin{tikzpicture}[scale=0.88]
  \node[solpanel] at (-2.3,1.65) {(a) 铰链 $C$};
  \node[solpoint,label={[sollabel]below left:$C$}] (C) at (-2.3,0) {};
  \draw[solbody] (C) -- ++(120:1.35);
  \draw[solbody] (C) -- ++(0:1.35);
  \draw[solforce] (C) -- ++(120:1.05) node[left=5pt,sollabel] {$S$};
  \draw[solforce] (C) -- ++(0:1.10) node[above,sollabel] {$H$};
  \draw[solforce] (C) -- ++(-90:1.20) node[below,sollabel] {$mg$};
  \draw[solgeom] (C) -- ++(180:0.82);
  \draw[solgeom] ($(C)+(120:0.50)$) arc[start angle=120,end angle=180,radius=0.50];
  \node[sollabel] at (-3.03,0.30) {$60^\circ$};

  \node[solpanel] at (2.15,1.65) {(b) 铰链 $D$};
  \node[solpoint,label={[sollabel]below left=4pt:$D$}] (D) at (2.15,0) {};
  \draw[solbody] (D) -- ++(60:1.45)
    node[pos=.68,right=4pt,sollabel] {$DB$};
  \draw[solbody] (D) -- ++(180:1.25);
  \draw[solforce] (D) -- ++(180:1.10) node[above,sollabel] {$H$};
  \draw[solforce] (D) -- ++(-30:1.15) node[right,sollabel] {$F_{\min}$};
  \draw[solgeom] (D) -- ++(150:1.05);
  \draw[solgeom] ($(D)+(0.35,0.61)$) -- ++(-30:0.22) -- ++(-120:0.22);
\end{tikzpicture}
\soldiagcaption{先由 $C$ 点平衡求出杆 $CD$ 的轴力，再在 $D$ 点取垂直于 $DB$ 的分量。}
\end{center}}

\newcommand{\soldiagOneEleven}{%
\begin{center}
\begin{tikzpicture}[scale=1]
  \node[solpanel] at (-2.3,2.2) {(a) 小球受力};
  \coordinate (O) at (-2.3,0);
  \draw[solbody] (O) circle (.42);
  \draw[solforce] (O)--++(102:1.65) node[above left,sollabel] {$N_1$};
  \draw[solforce] (O)--++(-18:1.65) node[right,sollabel] {$N$};
  \draw[solforce] (O)--++(-90:1.60) node[below,sollabel] {$mg$};
  \node[solpanel] at (2.35,2.2) {(b) 沿下斜面投影};
  \draw[solgeom,->] (.55,.8)--(4.15,.8) node[right,sollabel] {向上};
  \draw[solforce] (1.1,.20)--(3.8,.20) node[midway,above=4pt,sollabel] {$N\sin\theta$};
  \draw[solforce] (3.8,-.65)--(1.1,-.65) node[midway,below=4pt,sollabel] {$mg\sin\alpha$};
  \node[sollabel] at (2.4,-1.55) {$N_1$ 的投影为零};
\end{tikzpicture}
\soldiagcaption{$N_1$ 垂直于下斜面，故沿下斜面只需平衡 $N$ 和 $mg$ 的分量。}
\end{center}}

\newcommand{\soldiagOneTwelve}{%
\begin{center}
\begin{tikzpicture}[scale=1]
  \coordinate (O) at (0,0);
  \draw[solbody] (O) circle (.35);
  \draw[solgeom] (0,-2.25)--(0,2.25);
  \draw[solgeom,->] (-1.55,0)--(2.45,0) node[right,sollabel] {$x$};
  \draw[solforce] (O)--++(103:2.10) node[above left,sollabel] {$N_1$};
  \draw[solforce] (O)--++(-38:2.10) node[below right,sollabel] {$N_2$};
  \draw[solforce] (O)--++(-90:2.00) node[below left,sollabel] {$mg$};
  \draw[solgeom] (90:1.25) arc[start angle=90,end angle=103,radius=1.25];
  \draw[TextGray,thin] (96.5:1.25)--(-1.05,1.35);
  \node[left,sollabel] at (-1.05,1.35) {$\theta_1$};
  \draw[solgeom] (-90:.95) arc[start angle=-90,end angle=-38,radius=.95];
  \draw[TextGray,thin] (-64:.95)--(1.55,-.58);
  \node[right,sollabel] at (1.55,-.58) {$\theta_2$};
\end{tikzpicture}
\soldiagcaption{$N_1$ 与竖直向上的夹角为 $\theta_1$；$N_2$ 与竖直向下的夹角为 $\theta_2$，因为支持力分别垂直于对应斜面。}
\end{center}}

\newcommand{\soldiagOneThirteen}{%
\begin{center}
\begin{tikzpicture}[scale=0.82]
  \node[solpanel] at (-2.45,2.15) {(a) 接触几何};
  \coordinate (O) at (-2.45,1.55);
  \coordinate (A) at (-3.25,-0.85);
  \coordinate (C) at (-2.45,0.05);
  \fill[Ink] (O) circle (1.4pt) node[above,sollabel] {$O$};
  \fill[Ink] (A) circle (1.4pt) node[below left,sollabel] {$A$};
  \fill[Ink] (C) circle (1.4pt) node[right,sollabel] {$C$};
  \draw[solbody] (O)--(A) node[midway,left,sollabel] {$d_0$};
  \draw[solbody] (A)--(C);
  \draw[TextGray,thin] ($(A)!.5!(C)$)--(-1.85,-.12);
  \node[right,sollabel] at (-1.85,-.12) {$R+r$};
  \draw[solgeom] (A)--(-1.55,-0.85);
  \draw[solgeom] (O)--(-2.45,-1.10);
  \draw[solgeom] ($(O)+(0,-0.55)$) arc[start angle=-90,end angle=-108.5,radius=0.55];
  \draw[TextGray,thin] ($(O)+(-99.25:.55)$)--(-3.15,1.04);
  \node[left,sollabel] at (-3.15,1.04) {$\theta$};
  \draw[solgeom] ($(A)+(0.55,0)$) arc[start angle=0,end angle=48,radius=0.55];
  \node[sollabel] at (-2.57,-0.56) {$\beta$};

  \node[solpanel] at (2.15,2.15) {(b) 大球 $A$ 受力};
  \node[solpoint,label={[sollabel]right=4pt:$A$}] (P) at (2.15,0) {};
  \draw[solforce] (P)--++(71.5:1.45) node[above,sollabel] {$T$};
  \draw[solforce] (P)--++(-132:1.18) node[below left,sollabel] {$N$};
  \draw[solforce] (P)--++(180:1.05) node[above left,sollabel] {$Q$};
  \draw[solforce] (P)--++(-90:1.30) node[below,sollabel] {$Mg$};
  \draw[solgeom] (P)--++(180:1.25);
  \draw[solgeom] ($(P)+(180:0.54)$)
    arc[start angle=180,end angle=228,radius=0.54];
  \node[sollabel] at ($(P)+(204:0.78)$) {$\beta$};
\end{tikzpicture}
\soldiagcaption{$\beta$ 由三球接触几何确定；两大球刚分离时令接触力 $Q=0$。}
\end{center}}

\newcommand{\soldiagOneFourteen}{%
\begin{center}
\begin{tikzpicture}[scale=0.92]
  \node[solpanel] at (-2.4,2.05) {(a) 切线几何};
  \coordinate (O) at (-2.40,1.55);
  \coordinate (C) at (-3.05,-0.55);
  % P is the tangent point for R=.96; hence CP is perpendicular to OP.
  \coordinate (P) at (-2.101,-0.405);
  \draw[solbody] (C) circle (0.96);
  \fill[Ink] (O) circle (1.4pt) node[above,sollabel] {$O$};
  \fill[Ink] (C) circle (1.4pt) node[below left,sollabel] {$C$};
  \fill[Ink] (P) circle (1.2pt) node[right,sollabel] {$P$};
  \draw[solbody] (O)--(C);
  \node[sollabel] at (-3.30,.08) {$l$};
  \draw[solbody] (O)--(P);
  \draw[solbody] (C)--(P) node[pos=.42,below=5pt,sollabel] {$R$};
  \draw[solgeom] (O)--++(0,-2.45);
  \draw[solgeom] ($(P)+(171.3:.16)$)--++(81.3:.16)--++(-8.7:.16);
  \draw[draw=TextGray] ($(O)+(-107.2:.60)$) arc[start angle=-107.2,end angle=-90,radius=.60];
  \draw[draw=TextGray] ($(O)+(-90:.80)$) arc[start angle=-90,end angle=-81.3,radius=.80];
  \draw[draw=TextGray] ($(O)+(-107.2:1.10)$) arc[start angle=-107.2,end angle=-81.3,radius=1.10];
  \node[sollabel] at (-2.90,.82) {$\theta$};
  \node[sollabel] at (-1.92,.79) {$\psi$};
  \draw[TextGray,thin] ($(O)+(-94.25:1.10)$)--(-3.85,.40);
  \node[left,sollabel] at (-3.85,.40) {$\delta$};

  \node[solpanel] at (2.05,2.05) {(b) 球的合力图};
  \node[solpoint] (G) at (2.05,0) {};
  \draw[solforce] ($(G)+(-.04,0)$)--++(72.8:1.35) node[above right,sollabel] {$S$};
  % At P the upper-rope tension points from P to O, i.e. along the
  % tangent, almost vertically upward in panel (a).
  \draw[solforce] ($(G)+(.04,0)$)--++(98.7:1.13) node[above left,sollabel] {$T$};
  \draw[solforce] (G)--++(-90:1.22) node[below,sollabel] {$m_1g$};
  \draw[solforce] ($(G)+(.18,-.04)$)--++(-90:.82) node[right=3pt,sollabel] {$T$};
  \draw[solgeom] (G)--++(0,1.25);
\end{tikzpicture}
\soldiagcaption{$\delta=\theta+\psi$ 且 $\sin\delta=R/l$；绕绳对球的作用等于两端张力的矢量和。}
\end{center}}

\newcommand{\soldiagOneFifteen}{%
\begin{center}
\begin{tikzpicture}[scale=0.78]
  \draw[solbody] (-2,-1.55)--(-2,2.00);
  \draw[solbody] (2,2.00)--(2,-1.55);
  \draw[solbody] (-2.35,-1.55)--(2.35,-1.55);
  \coordinate (A) at (-0.90,-0.45);
  \coordinate (B) at (0.90,0.815);
  \draw[solbody] (A) circle (1.10);
  \draw[solbody] (B) circle (1.10);
  \draw[solgeom] (A)--(B);
  \draw[solgeom] (A)--++(2.00,0);
  \draw[solgeom] ($(A)+(0:.55)$) arc[start angle=0,end angle=35.1,radius=.55];
  \node[sollabel] at ($(A)+(17.6:.78)$) {$\beta$};
  \draw[solforce] (-2,-0.45)--++(-0.95,0) node[left,sollabel] {$N_2$};
  \draw[solforce] (2,0.815)--++(0.95,0) node[right,sollabel] {$N_2$};
  \draw[soldim] (2.48,-0.45)--(2.48,0.815);
  \node[sollabel,right] at (2.48,0) {$2r\sin\beta$};
  \draw[solforce] (0,-0.25)--(0,-2.02) node[below,sollabel] {$Mg$};
  \draw[TextGray,thin] (2,-1.55)--(2.72,-1.72);
  \node[right,sollabel] at (2.72,-1.72) {倾覆边缘};
\end{tikzpicture}
\soldiagcaption{两侧壁受到一对等大反向水平力，形成力偶 $N_2(2r\sin\beta)$。}
\end{center}}

\newcommand{\soldiagOneSixteen}{%
\begin{center}
\begin{tikzpicture}[scale=0.72]
  \coordinate (A) at (-3,-1.35);
  \pgfmathsetmacro{\sone}{.52*(1+sin(25))/cos(25)}
  \draw[solbody] (-3,-1.65)--(-3,2.05);
  \draw[solbody] (A)--++(25:6.90);
  \foreach \i/\lab in {0/1,1/2,2/3}{
    \coordinate (P\lab) at ($(A)+(25:{\sone+1.04*\i})$);
    \coordinate (C\lab) at ($(P\lab)+(115:0.52)$);
    \draw[solbody] (C\lab) circle (0.52);
    \draw[solforce] (P\lab)--++(115:0.92);
    \ifnum\lab=1
      \draw[TextGray,thin] ($(P\lab)+(115:.92)$)--(-3.30,.40);
      \node[left,sollabel] at (-3.30,.40) {$N_1$};
    \else
      \node[above left,sollabel] at ($(P\lab)+(115:.92)$) {$N_{\lab}$};
    \fi
  }
  \fill[Ink] (A) circle (1.2pt) node[below left,sollabel] {$A$};
  \draw[soldim] ($(A)+(25:0.05)+(115:-0.32)$)--($(A)+(25:\sone)+(115:-0.32)$);
  \node[sollabel] at (-2.54,-1.82) {$s_1$};
  \draw[solgeom] ($(A)+(25:.30)$) arc[start angle=25,end angle=90,radius=.30];
  \draw[TextGray,thin] ($(A)+(57.5:.30)$)--(-3.58,-1.06);
  \node[left,sollabel] at (-3.58,-1.06) {$\theta$};
  \node[sollabel] at (0.55,-0.65) {$s_2=s_1+2r$};
  \node[sollabel] at (2.05,-0.05) {$s_3=s_1+4r$};
\end{tikzpicture}
\soldiagcaption{接触点沿板的位置为 $s_1,s_2,s_3$，板对各圆柱的支持力均垂直于板。}
\end{center}}

\newcommand{\soldiagOneSeventeen}{%
\begin{center}
\begin{tikzpicture}[scale=0.82]
  \coordinate (A) at (-0.40,-1.25);
  \coordinate (B) at ($(A)+(60:3.15)$);
  \draw[solbody,line width=1.2pt] (A)--(B);
  \draw[solbody] (-2.75,-1.25)--(2.55,-1.25);
  \draw[solbody] (B).. controls ($(B)+(210:1.10)$) and (-1.15,0.63)..(-2.55,0.55);
  \draw[solgeom] (B)--++(180:1.75);
  \draw[solgeom] ($(A)+(0.72,0)$) arc[start angle=0,end angle=60,radius=0.72];
  \node[sollabel] at (0.50,-0.88) {$\beta$};
  \draw[solgeom] ($(B)+(-0.68,0)$) arc[start angle=180,end angle=210,radius=0.68];
  \draw[TextGray,thin] ($(B)+(195:.68)$)--(-.40,1.82);
  \node[above left,sollabel] at (-.40,1.82) {$\gamma$};
  \draw[solgeom] ($(B)+(210:0.82)$) arc[start angle=210,end angle=240,radius=0.82];
  \draw[TextGray,thin] ($(B)+(225:.82)$)--(1.35,.95);
  \node[right,sollabel] at (1.35,.95) {$30^\circ$};
  \draw[solforce] (B)--++(210:1.45) node[below left=3pt,sollabel] {$T$};
  \coordinate (G) at ($(A)!0.5!(B)$);
  \draw[solforce] (G)--++(-90:0.76)
    node[pos=0.62,right=4pt,sollabel] {$4mg$};
  \fill[Ink] (A) circle (1.3pt) node[below,sollabel] {铰链};
  \node[sollabel] at (-1.82,0.90) {钢索切线};
\end{tikzpicture}
\soldiagcaption{一侧斜壁重为 $4mg$；斜壁的力矩平衡与屋面整体的竖直平衡必须同时满足。}
\end{center}}

\newcommand{\soldiagOneEighteen}{%
\begin{center}
\begin{tikzpicture}[scale=0.88]
  \draw[solbody] (-2.65,-0.85)--(2.75,-0.85);
  \draw[solbody,fill=white] (-1.00,-0.85) rectangle (1.00,0.55);
  \coordinate (G) at (0,-0.15);
  \draw[solforce] (1.00,0.12)--++(30:1.55) node[above right,sollabel] {$F$};
  \draw[solforce] (G)--++(90:1.35) node[above,sollabel] {$N$};
  \draw[solforce] (G)--++(-90:1.35) node[below,sollabel] {$Mg$};
  \draw[solforce] (-0.80,-0.85)--++(180:1.25)
    node[above left=2pt,sollabel] {$f=\mu N$};
  \draw[solgeom] (1.00,0.12)--++(0:1.25);
  \draw[solgeom] ($(1.00,0.12)+(0:0.68)$) arc[start angle=0,end angle=30,radius=0.68];
  \node[sollabel] at (1.82,0.33) {$\alpha$};
  \draw[solgeom,-{Latex[length=1.5mm]}] (1.55,-1.15)--(2.55,-1.15)
    node[midway,below,sollabel] {匀速};
\end{tikzpicture}
\soldiagcaption{沿运动方向有 $F\cos\alpha=f$；拉力的竖直分量使支持力减为 $N=Mg-F\sin\alpha$。}
\end{center}}

\newcommand{\soldiagOneNineteen}{%
\begin{center}
\begin{tikzpicture}[scale=0.82]
  \node[solpanel] at (-2.20,1.70) {(a) 两接触点};
  \coordinate (O) at (-3.55,0);
  \draw[solbody] (O)--(-0.55,1.15);
  \draw[solbody] (O)--(-0.55,-1.15);
  \draw[solbody] (-1.55,0) circle (0.72);
  \coordinate (U) at (-1.808,0.672); \coordinate (L) at (-1.808,-0.672);
  \draw[solforce] (U)--++(-69:0.92);
  \draw[solforce] (U)--++(201:0.88);
  \draw[solforce] (L)--++(69:0.92);
  \draw[solforce] (L)--++(159:0.88);
  \node[sollabel] at (-1.25,.22) {$N$};
  \node[sollabel] at (-1.25,-.22) {$N$};
  \node[sollabel] at (-2.67,.63) {$\mu N$};
  \node[sollabel] at (-2.67,-.63) {$\mu N$};
  \draw[solgeom] (O)--(-0.65,0);
  \draw[solgeom] ($(O)+(0.78,0)$) arc[start angle=0,end angle=21,radius=0.78];
  \draw[TextGray,thin] ($(O)+(10.5:.78)$)--(-3.53,-.72);
  \node[left,sollabel] at (-3.53,-.72) {$\theta/2$};

  \node[solpanel] at (2.55,1.70) {(b) 沿角平分线的合力};
  \draw[solgeom] (0.70,0)--(4.45,0);
  \fill[Ink] (2.55,0) circle (1.2pt);
  \draw[solforce] (2.55,0.30)--++(0:1.55)
    node[right,sollabel] {$2N\sin(\theta/2)$};
  \draw[solforce] (2.55,-0.30)--++(180:1.55);
  \node[sollabel,below=3pt] at (1.78,-.30) {$2\mu N\cos(\theta/2)$};
  \node[sollabel] at (3.45,0.70) {向外};
  \node[sollabel] at (1.65,-1.20) {向内};
\end{tikzpicture}
\soldiagcaption{正压力的合力向外，极限摩擦力的合力沿角平分线指向 $O$。}
\end{center}}

\newcommand{\soldiagOneTwenty}{%
\begin{center}
\begin{tikzpicture}[scale=0.82]
  \coordinate (A) at (-2.7,0);
  \coordinate (B) at (2.7,0);
  \draw[solbody] (-2.7,-1.1)--(-2.7,2.0);
  \draw[solbody,line width=1.15pt] (A)--(B);
  \draw[solbody] (B)--(-2.7,1.65);
  \fill[Ink] (A) circle (1.2pt) node[below left,sollabel] {$A$};
  \fill[Ink] (B) circle (1.2pt) node[below right,sollabel] {$B$};
  \coordinate (P) at (-0.90,0);
  \fill[Ink] (P) circle (1.2pt) node[above=4pt,sollabel] {$P$};
  \draw[solforce] (A)--++(0:1.05) node[below,sollabel] {$N$};
  \draw[solforce] (A)--++(90:1.10) node[left,sollabel] {$f$};
  \draw[solforce] (B)--++(163:1.32) node[above,sollabel] {$T$};
  \draw[solforce] (0,0)--++(-90:1.05) node[below,sollabel] {$Mg$};
  \draw[solforce] (P)--++(-90:1.05) node[below,sollabel] {$mg$};
  \draw[solgeom] ($(B)+(-0.72,0)$) arc[start angle=180,end angle=163,radius=0.72];
  \draw[TextGray,thin] ($(B)+(171.5:.72)$)--(2.05,-.40);
  \node[right,sollabel] at (2.05,-.40) {$\theta$};
  \draw[soldim] ($(A)+(0,-0.58)$)--($(P)+(0,-0.58)$);
  \node[sollabel] at (-1.46,-0.78) {$y$};
\end{tikzpicture}
\soldiagcaption{附加重物放在距 $A$ 为 $y$ 处；对 $A$ 取矩可直接消去 $N$、$f$。}
\end{center}}

\newcommand{\soldiagOneTwentyOne}{%
\begin{center}
\begin{tikzpicture}[scale=0.86]
  \coordinate (A) at (-2.2,-0.85);
  \coordinate (B) at (2.05,0.95);
  \draw[solbody] (-2.75,-0.85)--(2.75,-0.85);
  \draw[solbody,line width=1.2pt] (A)--(B);
  \draw[solforce] (B)--++(113:1.25) node[above,sollabel] {$F\perp$ 杆};
  \coordinate (G) at ($(A)!0.5!(B)$);
  \draw[solforce] (G)--++(-90:1.18) node[below,sollabel] {$mg$};
  \draw[solforce] (A)--++(90:1.18) node[left,sollabel] {$N$};
  \draw[solforce] (A)--++(0:1.08) node[below,sollabel] {$f$};
  \draw[solgeom] ($(A)+(0.72,0)$) arc[start angle=0,end angle=23,radius=0.72];
  \draw[TextGray,thin] ($(A)+(11.5:.72)$)--(-.70,-.42);
  \node[right,sollabel] at (-.70,-.42) {$\theta$};
\end{tikzpicture}
\soldiagcaption{外力始终与杆垂直；地板处的 $f/N$ 随抬起角度变化。}
\end{center}}

\newcommand{\soldiagOneTwentyTwo}{%
\begin{center}
\begin{tikzpicture}[scale=0.82]
  \node[solpanel] at (-2.15,1.95) {(a) 中间结点};
  \node[solpoint] (C) at (-2.15,0) {};
  \draw[solforce] (C)--++(145:1.25) node[above left,sollabel] {$T$};
  \draw[solforce] (C)--++(35:1.25) node[above right,sollabel] {$T$};
  \draw[solforce] (C)--++(-90:1.30) node[below,sollabel] {$Mg$};
  \draw[solgeom] (C)--++(1.45,0);
  \node[sollabel] at (-1.45,0.25) {$\beta$};

  \node[solpanel] at (2.15,1.95) {(b) 圆环 $A$};
  \draw[solbody] (0.70,0.72)--(3.65,0.72);
  \draw[solbody] (2.15,0.72) circle (0.20);
  \draw[solforce] (2.15,0.62)--++(-35:1.25) node[below right,sollabel] {$T$};
  \draw[solforce] (2.15,0.92)--++(90:0.82) node[midway,right,sollabel] {$N$};
  \draw[solforce] (2.15,0.52)--++(-90:1.02) node[below,sollabel] {$mg$};
  \draw[solforce] (2.15,0.76)--++(180:1.00)
    node[above left=2pt,sollabel] {$f$};
\end{tikzpicture}
\soldiagcaption{结点给出 $2T\sin\beta=Mg$；圆环的支持力还要承担细线的竖直分量。}
\end{center}}

\newcommand{\soldiagOneTwentyThree}{%
\begin{center}
\begin{tikzpicture}[scale=0.82]
  \node[solpanel,align=center] at (-3.00,2.35) {(a) 绳圈在圆锥\\上的位置};
  \coordinate (A) at (-3.00,1.60);
  \pgfmathsetmacro{\solConeHalfAngle}{atan(1.55/2.95)}
  \pgfmathsetmacro{\solConeRingRadius}{1.78*1.55/2.95}
  \draw[solbody,line width=1.05pt] (-4.55,-1.35)--(A)--(-1.45,-1.35);
  \draw[solgeom] (A)--(-3.00,-1.55) node[below,sollabel] {轴};
  \draw[solbody,line width=1.05pt] ({-3-\solConeRingRadius},-.18)
    arc[start angle=180,end angle=360,x radius=\solConeRingRadius,y radius=.20];
  \draw[solgeom] ({-3-\solConeRingRadius},-.18)
    arc[start angle=180,end angle=0,x radius=\solConeRingRadius,y radius=.20];
  \draw[solgeom] (-3.00,-0.18)--({-3+\solConeRingRadius},-.18)
    node[midway,above,sollabel] {$\rho$};
  \draw[<->,draw=TextGray] (-4.28,-.18)--(-4.28,1.60)
    node[midway,left,sollabel] {$z=\rho\cot\delta$};
  \draw[draw=TextGray] ($(A)+(-90:.66)$)
    arc[start angle=-90,end angle={-90+\solConeHalfAngle},radius=.66];
  \draw[TextGray,thin] ($(A)+({-90+\solConeHalfAngle/2}:.66)$)--(-1.65,1.43);
  \node[right,sollabel] at (-1.65,1.43) {$\delta=\alpha/2$};

  \node[solpanel,align=center] at (2.10,2.35) {(b) 纬圈微元\\径向--竖直截面};
  \coordinate (P) at (2.10,-.20);
  \draw[solgeom] (P)--($(P)+(2.05,0)$) node[right,sollabel] {水平径向};
  \draw[solgeom] ($(P)+(0,-1.65)$)--($(P)+(0,1.65)$);
  \node[sollabel,right=4pt] at ($(P)+(0,1.55)$) {竖直};
  % Numerical parameters never redefine a mathematical symbol such as delta.
  \draw[solbody,line width=1.05pt]
    ($(P)+({90+\solConeHalfAngle}:1.55)$)--($(P)+({-90+\solConeHalfAngle}:1.55)$);
  \fill[Ink] (P) circle (1.3pt);
  \draw[solforce] (P)--++(180:1.28) node[left,sollabel] {$T\,\mathrm d\varphi$};
  \draw[solforce] (P)--++(\solConeHalfAngle:1.55) node[above right,sollabel] {$\mathrm dN$};
  \draw[solforce] (P)--++(-90:1.20) node[below,sollabel] {$\mathrm dm\,g$};
  \draw[solgeom] ($(P)+(0:.78)$) arc[start angle=0,end angle=\solConeHalfAngle,radius=.78];
  \node[sollabel] at ($(P)+({\solConeHalfAngle/2}:1.03)$) {$\delta$};
\end{tikzpicture}
\soldiagcaption{左图先标出水平纬圈；右图把该圈微元放入径向--竖直截面，张力合力水平向内，支持力向外且向上。}
\end{center}}

\newcommand{\soldiagOneTwentyFour}{%
\begin{center}
\begin{tikzpicture}[scale=0.86]
  \coordinate (O) at (0,0);
  \draw[solbody] (O) ++(0:1.65) arc[start angle=0,end angle=90,radius=1.65];
  \draw[solbody] (-0.15,-0.08)--(1.82,-0.08);
  \fill[Ink] (O) circle (1.2pt) node[below left,sollabel] {$O$};
  \coordinate (P) at ($(O)+(48:1.65)$);
  \fill[Ink] (P) circle (1.2pt);
  \draw[solgeom] (O)--(P) node[midway,above left=5pt,sollabel] {$R$};
  \draw[solgeom] (O)--++(0:1.15);
  \draw[solgeom] ($(O)+(0:0.62)$) arc[start angle=0,end angle=48,radius=0.62];
  \node[sollabel] at (0.85,0.20) {$\varphi$};
  \draw[solforce] (P)--++(138:1.20) node[above left,sollabel] {$T+\mathrm dT$};
  \draw[solforce] (P)--++(-42:1.08) node[below right,sollabel] {$T$};
  \draw[solforce] (P)--++(48:0.95) node[above right,sollabel] {$\mathrm dN$};
  \draw[solforce] (P)--++(-90:1.15);
  % 重力标记单列在图下方，避免同时压住竖直力箭头和斜向张力箭头。
  \node[sollabel] at (1.08,-0.42) {$\lambda gR\,\mathrm d\varphi$};
  \node[sollabel] at (-1.60,0.95) {$\mathrm ds=R\,\mathrm d\varphi$};
\end{tikzpicture}
\soldiagcaption{微元两端张力之差平衡重力的切向分量；沿 $\varphi$ 增大方向张力增加。}
\end{center}}

\newcommand{\soldiagOneTwentyFive}{%
\begin{center}
\begin{tikzpicture}[scale=0.80]
  \draw[solbody] (0,0.45) circle (1.25);
  \draw[solbody,line width=1.15pt] (-1.25,0.45)--(-1.25,-2.05);
  \draw[solbody,line width=1.15pt] (1.25,0.45)--(1.25,-2.05);
  \draw[solbody,line width=1.15pt] (-1.25,0.45) arc[start angle=180,end angle=0,radius=1.25];
  \draw[soldim] (1.68,-2.05)--(1.68,0.45);
  \node[sollabel,right] at (1.68,-0.80) {$h=(l-\pi R)/2$};
  \draw[solforce] (1.25,0.38)--++(90:0.95) node[right,sollabel] {$T_0=\lambda gh$};
  \draw[solforce] (0,1.70)--++(0.95,0) node[right,sollabel] {$T_{\max}$};
  \draw[solgeom] (0,0.45)--(0,1.70) node[midway,left,sollabel] {$R$};
  \node[sollabel] at (-2.65,1.45) {$T_{\max}=T_0+\lambda gR$};
\end{tikzpicture}
\soldiagcaption{从自由端向最高点累加下方绳段的切向重力，最高点张力最大。}
\end{center}}

\newcommand{\soldiagOneTwentySix}{%
\begin{center}
\begin{tikzpicture}[scale=0.82]
  \node[solpanel,align=center] at (-3.05,2.65) {(a) 球面上的\\水平纬圈};
  \coordinate (O) at (-3.05,-0.25);
  \draw[solbody] (O) circle (1.65);
  \draw[solgeom] (-3.05,-2.05)--(-3.05,1.72);
  \draw[solbody,line width=1.05pt] (-4.22,.917)
    arc[start angle=180,end angle=360,x radius=1.17,y radius=.25];
  \draw[solgeom] (-4.22,.917)
    arc[start angle=180,end angle=0,x radius=1.17,y radius=.25];
  \fill[Ink] (O) circle (1.2pt) node[below left,sollabel] {$O$};
  \draw[solgeom] (O)--(-3.05,.917) node[midway,left,sollabel] {$z$};
  \draw[solgeom] (-3.05,.917)--(-1.88,.917) node[midway,above,sollabel] {$b$};
  \draw[solgeom] (O)--(-1.88,.917) node[midway,below right,sollabel] {$R$};

  \node[solpanel,align=center] at (1.38,2.65) {(b) 纬圈微元\\的受力};
  \coordinate (C) at (-.05,-.55);
  \coordinate (P) at (1.15,.65);
  \draw[solgeom] (C)--(P) node[midway,above left,sollabel] {$R$};
  \draw[solgeom] (C)--(1.15,-.55);
  \node[sollabel] at (.50,-.82) {$b$};
  \draw[solgeom] (1.15,-.55)--(P) node[midway,right=4pt,sollabel] {$z$};
  \fill[Ink] (C) circle (1.2pt) node[below left,sollabel] {$O$};
  \fill[Ink] (P) circle (1.2pt) node[above left=5pt,sollabel] {$P$};
  \draw[solforce] (P)--++(45:1.18) node[above right,sollabel] {$\mathrm dN$};
  \draw[solforce] (P)--++(180:1.18) node[left,sollabel] {$T\,\mathrm d\varphi$};
  \draw[solforce] (P)--++(-90:1.60) node[below=4pt,sollabel] {$\dfrac{Mg}{2\pi}\,\mathrm d\varphi$};
\end{tikzpicture}
\soldiagcaption{左图明确绳圈是球面纬圈；右图显示支持力沿球面半径，水平分量平衡张力合力，竖直分量平衡重力。}
\end{center}}

\newcommand{\soldiagOneTwentySeven}{%
\begin{center}
\begin{tikzpicture}[scale=1]
  \def\cc{2.7}\def\xp{2.2}
  \pgfmathsetmacro{\uu}{\xp/\cc}
  \pgfmathsetmacro{\yp}{\cc*(cosh(\uu)-1)}
  \pgfmathsetmacro{\slope}{sinh(\uu)}
  \pgfmathsetmacro{\angleT}{atan(\slope)}
  \pgfmathsetmacro{\xg}{\xp-\yp/\slope}
  \pgfmathsetmacro{\yg}{\xp/(2*sinh(\uu))+\cc*cosh(\uu)/2-\cc}
  \draw[axis] (-1.65,0)--(3.85,0) node[right,sollabel] {$x$};
  \draw[axis] (0,-.65)--(0,2.35) node[above,sollabel] {$y$};
  \draw[TextGray,thick,domain=-1.65:0,samples=41] plot (\x,{\cc*(cosh(\x/\cc)-1)});
  \draw[curve,domain=0:2.2,samples=61] plot (\x,{\cc*(cosh(\x/\cc)-1)});
  \fill[Ink] (0,0) circle (1.3pt) node[below left=3pt,sollabel] {$O$};
  \coordinate (P) at (\xp,\yp);
  \fill[Ink] (P) circle (1.3pt) node[above left=3pt,sollabel] {$P$};
  \draw[solforce] (P)--++(\angleT:1.40) node[above right,sollabel] {$T$};
  \coordinate (G) at (\xg,\yg);
  \draw[solforce] (G)--(\xg,-1.10) node[below=4pt,sollabel] {$\lambda gs$};
  \draw[solforce] (0,0)--++(180:1.40) node[below left=3pt,sollabel] {$H$};
  \draw[solgeom] (P)--++(0:1.40);
  \draw[solgeom] ($(P)+(0:.76)$) arc[start angle=0,end angle=\angleT,radius=.76];
  \node[sollabel] at ($(P)+({\angleT/2}:1.05)$) {$\theta$};
  \draw[solgeom] (P)--(\xp,-1.85) (0,-1.55)--(0,-1.85);
  \draw[soldim] (0,-1.75)--(\xp,-1.75) node[midway,below=3pt,sollabel] {$x$};
  \node[sollabel] at (.75,1.62) {弧长 $s$};
  \draw[TextGray,thin] (.75,1.40)--(.65,{\cc*(cosh(.65/\cc)-1)});
\end{tikzpicture}
\soldiagcaption{对最低点到 $P$ 的绳段整体受力：左端受 $H$，右端受沿切线的 $T$，绳重为 $\lambda gs$。}
\end{center}}

\newcommand{\soldiagOneTwentyEight}{%
\begin{center}
\begin{tikzpicture}[scale=0.82]
  \coordinate (O) at (0,-1.25);
  \coordinate (A0) at (0,1.45);
  \def\rodlen{2.70}\def\tilt{17}
  \coordinate (A) at ($(O)+({90-\tilt}:\rodlen)$);
  \coordinate (L) at (-2.85,1.45);
  \coordinate (R) at (2.85,1.45);
  \draw[solbody] (-3.15,-1.25)--(3.15,-1.25);
  \fill[Ink] (O) circle (1.4pt) node[below left=2pt,sollabel] {$O$};
  \draw[solgeom] (O)--(A0);
  \draw[solbody,line width=1.2pt] (O)--(A);
  \fill[Ink] (A) circle (1.3pt) node[below right=5pt,sollabel] {$A$};
  \draw[solbody] (L)--++(0,-0.42);
  \draw[solbody] (R)--++(0,-0.42);
  \draw[solbody,decorate,decoration={zigzag,segment length=4pt,amplitude=2pt}]
    (L)--(A);
  \draw[solbody,decorate,decoration={zigzag,segment length=4pt,amplitude=2pt}]
    (A)--(R);
  \draw[solgeom] (L)--(R);
  \coordinate (DimBase) at (0,2.24);
  \draw[solgeom] (A0)--(A0|-DimBase) (A)--(A|-DimBase);
  \draw[soldim] (A0|-DimBase)--(A|-DimBase)
    node[midway,above=3pt,sollabel] {$x=L\sin\varphi$};
  \draw[solgeom] ($(O)+(90:0.72)$) arc[start angle=90,end angle={90-\tilt},radius=0.72];
  \draw[TextGray,thin] ($(O)+({90-\tilt/2}:.72)$)--(-.52,-.28);
  \node[left,sollabel] at (-.52,-.28) {$\varphi$};
  \coordinate (G) at ($(O)!0.5!(A)$);
  \draw[solforce] (G)--++(-90:0.94) node[right=3pt,sollabel] {$mg$};
  \node[sollabel] at (-1.72,1.92) {$\Delta\ell_{\rm L}\simeq+x$};
  \node[sollabel] at (1.95,1.92) {$\Delta\ell_{\rm R}\simeq-x$};
\end{tikzpicture}
\soldiagcaption{杆顶水平位移为 $x$；小角度下两弹簧一伸一缩，形变量的一阶项为 $\pm x$。}
\end{center}}

\newcommand{\soldiagOneTwentyNine}{%
\begin{center}
\begin{tikzpicture}[scale=0.82]
  \node[solpanel] at (-2.15,1.80) {(a) 顶端质点 $M$};
  \coordinate (M) at (-2.15,.25);
  \draw[solgeom] (M)--++(235:1.15) (M)--++(305:1.15);
  \draw[solgeom] (M)--++(90:1.18);
  \fill[Ink] (M) circle (2.0pt) node[right=6pt,sollabel] {$M$};
  \draw[solforce] (M)--++(125:1.20) node[above left,sollabel] {$P$};
  \draw[solforce] (M)--++(55:1.20) node[above right,sollabel] {$P$};
  \draw[solforce] (M)--++(-90:1.15) node[below,sollabel] {$Mg$};
  \draw[solgeom] ($(M)+(90:.56)$) arc[start angle=90,end angle=55,radius=.56];
  \node[sollabel] at ($(M)+(72.5:.78)$) {$\theta$};

  \node[solpanel] at (2.05,1.80) {(b) 下端结点 $A$};
  \coordinate (A) at (2.05,-.25);
  \draw[solgeom] (A)--++(55:1.38);
  \fill[Ink] (A) circle (2.0pt) node[below right=6pt,sollabel] {$A$};
  \draw[solforce] (A)--++(235:1.10) node[below left,sollabel] {$P$};
  \draw[solforce] (A)--++(0:1.12) node[right,sollabel] {$T$};
  \draw[solforce] (A)--++(90:1.12) node[above left,sollabel] {$N_A$};
\end{tikzpicture}
\soldiagcaption{两个面板分别只画一个研究对象：顶端质点满足 $2P\cos\theta=Mg$；下端结点的杆力、弹簧力和支持力作用于同一点。}
\end{center}}

\newcommand{\soldiagOneThirty}{%
\begin{center}
\begin{tikzpicture}[scale=.92]
  \def\rodlen{4.6}\def\xp{3.8}
  \pgfmathsetmacro{\yp}{(\rodlen-sqrt(\rodlen^2-\xp^2))/2}
  \coordinate (A) at (0,{\rodlen-\yp});
  \coordinate (P) at (\xp,\yp);
  \coordinate (G) at ($(A)!0.5!(P)$);
  \coordinate (C) at (G|-A);
  \draw[solbody] (0,0)--(0,4.20);
  \draw[curve,domain=0:4.58,samples=101] plot (\x,{(\rodlen-sqrt(\rodlen^2-\x^2))/2});
  \draw[solbody,line width=1.2pt] (A)--(P);
  \fill[Ink] (A) circle (1.3pt) node[left,sollabel] {$A(0,Y)$};
  \fill[Ink] (P) circle (1.3pt) node[below right,sollabel] {$P(x,y)$};
  \fill[Ink] (G) circle (1.3pt) node[below left,sollabel] {$G$};
  \fill[Ink] (C) circle (1.3pt) node[above right=4pt,sollabel] {$C$};
  \draw[solgeom] (A)--($(C)+(1.25,0)$);
  \draw[solgeom] (G)--($(C)+(0,.40)$);
  \draw[solgeom] (P)--(C);
  \draw[solforce] (A)--++(0:1.15) node[below,sollabel] {$N_A$};
  \draw[solforce] (G)--++(-90:1.25) node[below,sollabel] {$mg$};
  \draw[solforce] (P)--($(P)!0.60!(C)$) node[midway,right,sollabel] {$N_P$};
  \node[sollabel,above] at (2.2,4.30) {三力作用线交于 $C$};
\end{tikzpicture}
\soldiagcaption{墙面支持力、重力和曲面法向力三条作用线必须共点。}
\end{center}}

\newcommand{\soldiagOneThirtyOne}{%
\begin{center}
\begin{tikzpicture}[scale=0.95]
  \def\Rbow{3.45}\def\rball{0.68}\def\phideg{12}
  \pgfmathsetmacro{\aradius}{\Rbow-\rball}
  \pgfmathsetmacro{\psideg}{-\aradius/\rball*\phideg}
  \pgfmathsetmacro{\centerangle}{-90+\phideg}
  \pgfmathsetmacro{\gangle}{90+\psideg}
  \coordinate (C) at (0,1.15);
  \coordinate (O) at ($(C)+(\centerangle:\aradius)$);
  \coordinate (K) at ($(C)+(\centerangle:\Rbow)$);
  \draw[solbody] ($(C)+(180:\Rbow)$) arc[start angle=180,end angle=360,radius=\Rbow];
  \fill[Ink] (C) circle (1.3pt) node[above left=2pt,sollabel] {碗心 $C$};
  \draw[solgeom] (C)--(0,-1.05);
  \draw[solgeom] (C)--(O) node[pos=0.48,left=6pt,sollabel] {$a=R-r$};
  \draw[solbody] (O) circle (\rball);
  \fill[Ink] (O) circle (1.3pt) node[below left,sollabel] {$O$};
  \fill[Ink] (K) circle (1.0pt);
  \draw[solgeom] (K)--++(.48,-.14) node[right,sollabel] {接触点};
  \draw[solgeom] ($(C)+(0,-0.66)$) arc[start angle=-90,end angle=\centerangle,radius=0.66];
  \draw[TextGray,thin] ($(C)+(-84:.66)$)--(-.64,.38);
  \node[left,sollabel] at (-.64,.38) {$\varphi$};
  \coordinate (G) at ($(O)+(\gangle:0.46)$);
  \fill[Ink] (G) circle (1.3pt) node[above right=4pt,sollabel] {$G$};
  \draw[solbody] (O)--(G);
  \coordinate (Hm) at ($(O)!.5!(G)$);
  \draw[TextGray,thin] (Hm)--($(O)+(2.30,-.12)$);
  \node[right,sollabel] at ($(O)+(2.30,-.12)$) {$OG=h$};
  \draw[solgeom] (O)--++(90:1.18);
  \draw[solgeom,-{Latex[length=1.7mm]}] ($(O)+(90:1.02)$)
    arc[start angle=90,end angle=\gangle,radius=1.02];
  \node[sollabel] at ($(O)+(66:1.32)$) {$\psi$};
\end{tikzpicture}
\soldiagcaption{球心与接触点沿同一半径；示意转角由无滑动关系 $\psi=-(a/r)\varphi$ 同步生成。}
\end{center}}

\newcommand{\soldiagConvexRolling}{%
\begin{center}
\begin{tikzpicture}[scale=1]
  \def\Sradius{2.10}\def\rball{.65}\def\phideg{15}
  \pgfmathsetmacro{\aradius}{\Sradius+\rball}
  \pgfmathsetmacro{\psideg}{-\aradius/\rball*\phideg}
  \coordinate (C) at (0,-1.8);
  \coordinate (O) at ($(C)+({90-\phideg}:\aradius)$);
  \coordinate (K) at ($(C)+({90-\phideg}:\Sradius)$);
  \coordinate (G) at ($(O)+({270+\psideg}:.30)$);
  \draw[solbody] ($(C)+(20:\Sradius)$) arc[start angle=20,end angle=160,radius=\Sradius];
  \draw[solbody] (O) circle (\rball);
  \draw[solgeom] (C)--($(C)+(90:3.25)$) (C)--(O);
  \fill[Ink] (C) circle (1.4pt) node[below,sollabel] {$C$};
  \fill[Ink] (K) circle (1.4pt);
  \draw[TextGray,thin] (K)--(1.90,-.12) node[right,sollabel] {接触点};
  \fill[Ink] (O) circle (1.4pt) node[right=4pt,sollabel] {$O$};
  \fill[Ink] (G) circle (1.4pt);
  \draw[TextGray,thin] (G)--(1.85,.60);
  \node[right,sollabel] at (1.85,.60) {$G$};
  \draw[solbody] (O)--(G);
  \draw[solgeom] ($(C)+(90:1.04)$) arc[start angle=90,end angle={90-\phideg},radius=1.04];
  \node[sollabel] at ($(C)+(82.5:1.32)$) {$\varphi$};
  \draw[solgeom] (O)--++(90:1.24);
  \draw[solgeom,->] ($(O)+(90:.95)$) arc[start angle=90,end angle={90+\psideg},radius=.95];
  \node[sollabel] at ($(O)+({90+\psideg/2}:1.25)$) {$\psi$};
  \node[sollabel,align=left] at (-1.80,.82) {$CO=S+r$\\$OG=|h|$，$h<0$};
  \node[sollabel] at (0,-2.28) {凸面：重心须在球心下方};
\end{tikzpicture}
\soldiagcaption{$C$、接触点和 $O$ 共线；球心向右偏转时球顺时针转动，图示 $G$ 位于球心下侧。}
\end{center}}

% Independently redrawn replacements for the 26 scanned source diagrams in Section 1.3.
% Each registration returns only a tikzpicture; the original requested width is retained.
\newcommand{\registerfaithful}[2]{\expandafter\def\csname faithful@#1\endcsname{#2}}
\newcommand{\sourcefigure}[2][0.52\linewidth]{%
  \begin{center}\resizebox{#1}{!}{\csname faithful@#2\endcsname}\end{center}}
\newcommand{\sourcefiguretrim}[3][0.52\linewidth]{\sourcefigure[#1]{#2}}
\newcommand{\sourcefigurefixB}{\sourcefigure[0.42\linewidth]{ex-1-10.png}}

\registerfaithful{fig-1-1.png}{%
\begin{tikzpicture}[x=1cm,y=1cm,>=Latex]
  % B is the orthogonal projection of O onto the force line through A.
  \coordinate (O) at (0,0); \coordinate (A) at (5.2,0); \coordinate (B) at (2.148,2.560);
  \draw[Ink,very thick] (O)--(A) node[midway,below] {$r$};
  \draw[TextGray,dashed] (O)--(B) node[midway,above left] {$r_\perp$};
  \draw[Ink,very thick,->] (A)--(B) node[midway,below left] {$F$};
  \draw[TextGray,dashed] (B)--++(-0.92,0.77) (B)--(5.2,2.560) (A)--(5.2,2.560);
  \draw[Ink,thick,->] (A)--(5.2,2.560) node[right] {$F_\perp$};
  \draw[Ink,thick,->] (5.2,2.560)--(B);
  \draw[Ink,thin] ($(B)+(-40:.25)$)--($(B)+(-40:.25)+(50:.25)$)--($(B)+(50:.25)$);
  \fill[Ink] (O) circle (2.2pt); \node[below left] at (O) {$O$}; \node[below right] at (A) {$A$}; \node[above right=8pt] at (B) {$B$};
  \draw[Ink] (4.25,0) arc[start angle=180,end angle=140,radius=0.95]; \node at (4.00,0.42) {$\theta$};
\end{tikzpicture}}

\registerfaithful{ex-1-10.png}{%
\begin{tikzpicture}[x=1cm,y=1cm,>=Latex]
  \draw[Ink,thick] (-3.1,2.295)--(3.1,2.295); \foreach \x in {-2.9,-2.5,...,2.9}{\draw[Ink] (\x,2.295)--++(0.28,0.30);}
  % AC=CD=DB=l=2.65 and AB=2l.
  \coordinate (A) at (-2.65,2.295); \coordinate (B) at (2.65,2.295); \coordinate (C) at (-1.325,0); \coordinate (D) at (1.325,0);
  \draw[Ink,very thick] (A)--(C)--(D)--(B); \draw[Ink,thick] (C)--++(0,-1.00);
  \draw[fill=SoftGray,draw=Ink,thick] (-1.725,-1.70) rectangle (-0.925,-1.00);
  \foreach \p/\lab/\pos in {A/A/above left,B/B/above right,C/C/below left,D/D/below right}{\fill[Ink] (\p) circle (1.7pt);\node[\pos] at (\p) {$\lab$};}
\end{tikzpicture}}

\registerfaithful{ex-1-11.png}{%
\begin{tikzpicture}[x=1cm,y=1cm,>=Latex]
  % The two boards meet at V; the circle is tangent to the 12 and 72 degree rays.
  \coordinate (V) at (-2.60,-1.20);
  \coordinate (C) at ($(V)+(42:{.85/sin(30)})$);
  \coordinate (E) at ($(V)+(12:5.75)$);
  \draw[Ink,thick] (-3.0,-1.20)--(3.25,-1.20);
  \draw[Ink,very thick] (V)--++(72:4.55);
  \draw[Ink,very thick] (V)--(E);
  \draw[Ink,thick] (E)--(E|-V);
  \draw[Ink,thick] (C) circle (0.85);
  \draw[Ink] ($(V)+(0:0.72)$) arc[start angle=0,end angle=12,radius=0.72];
  \draw[TextGray,thin] ($(V)+(6:.74)$)--(-1.30,-1.50);
  \node[below,inner sep=1pt] at (-1.30,-1.50) {$\alpha$};
  % Put the included-angle arc next to the vertex, clear of the ball.
  \draw[Ink] ($(V)+(12:0.52)$) arc[start angle=12,end angle=72,radius=0.52];
  \node[font=\small,inner sep=0pt] at ($(V)+(42:0.65)$) {$\theta$};
\end{tikzpicture}}

\registerfaithful{ex-1-12.png}{%
\begin{tikzpicture}[x=1cm,y=1cm,>=Latex]
  \coordinate (V) at (-3.0,-1.2); \coordinate (C) at ($(V)+(32.5:{.92/sin(19.5)})$);
  % The fixed faces are 13 and 52 degrees above horizontal; the ball is tangent to both.
  \draw[Ink,thick] (V)--(3.4,-1.2); \draw[Ink,very thick] (V)--++(13:6.1); \draw[Ink,very thick] (V)--++(52:5.0);
  \draw[Ink,thick] (C) circle (0.92);
  \draw[Ink] ($(V)+(0:.85)$) arc[start angle=0,end angle=13,radius=.85];
  \draw[TextGray,thin] ($(V)+(6.5:.86)$)--(-1.75,-1.55);
  \node[below,font=\small] at (-1.75,-1.55) {$\theta_1$};
  % Both inclinations start from the same horizontal ray; theta2 is not the wedge angle.
  \draw[Ink] ($(V)+(0:1.15)$) arc[start angle=0,end angle=52,radius=1.15];
  \node[font=\small,inner sep=0pt] at ($(V)+(34:1.42)$) {$\theta_2$};
\end{tikzpicture}}

\registerfaithful{ex-1-13.png}{%
\begin{tikzpicture}[x=1cm,y=1cm]
  % R=1.15, r=R/2.  The two large balls and the small ball are drawn
  % exactly tangent.  O is 3R from each large-ball centre, so the visible
  % string segment from O to the surface has length 2R and its extension
  % passes through the corresponding centre.
  \draw[Ink,thick] (-2.4,2.703)--(2.4,2.703); \foreach \x in {-2.2,-1.8,...,2.2}{\draw[Ink] (\x,2.703)--++(0.25,0.28);}
  \coordinate (O) at (0,2.703);
  \coordinate (A) at (-1.15,-0.55); \coordinate (B) at (1.15,-0.55); \coordinate (C) at (0,0.73574);
  \coordinate (As) at ($(A)!0.333333!(O)$); \coordinate (Bs) at ($(B)!0.333333!(O)$);
  \draw[Ink,thick] (O)--(As) (O)--(Bs);
  \draw[Ink,thick] (A) circle (1.15); \draw[Ink,thick] (B) circle (1.15); \draw[Ink,thick] (C) circle (0.575);
  \fill[Ink] (A) circle (1.8pt) node[below left] {$A$}; \fill[Ink] (B) circle (1.8pt) node[below right] {$B$}; \node at (C) {$C$};
\end{tikzpicture}}

\registerfaithful{ex-1-14.png}{%
\begin{tikzpicture}[x=1cm,y=1cm,>=Latex]
  \draw[Ink,thick] (-1.35,3.35)--(1.35,3.35); \foreach \x in {-1.2,-0.8,...,1.2}{\draw[Ink] (\x,3.35)--++(0.25,0.28);}
  \coordinate (O) at (0,3.35); \coordinate (C) at (-0.75,0.10);
  \coordinate (P) at (0.664,0.421); \coordinate (Q) at (0.70,0.10);
  \draw[TextGray,dashed] (O)--(C); \draw[Ink,thick] (O)--(C);
  \draw[Ink,thick] (C) circle (1.45) node[below left=11pt] {$m_1$};
  % OP is tangent; the rope then follows the surface to Q and leaves vertically.
  \draw[Ink,very thick] (O)--(P);
  \draw[Ink,very thick] (P) arc[start angle=12.8,end angle=0,radius=1.45];
  \draw[Ink,very thick] (Q)--(0.70,-1.45);
  \draw[fill=SoftGray,draw=Ink] (0.35,-1.96) rectangle (1.05,-1.45);
  \node[right=4pt] at (1.05,-1.70) {$m_2$};
  \draw[TextGray,dashed] (O)--++(0,-1.15);
  % theta is the angle between OC and the downward vertical.
  \draw[Ink] ($(O)+(-103:0.70)$) arc[start angle=-103,end angle=-90,radius=0.70];
  \draw[TextGray,thin] ($(O)+(-96.5:.70)$)--(-.72,2.68)--(-1.03,2.68);
  \node[left,font=\small] at (-1.03,2.68) {$\theta$};
  \fill[Ink] (O) circle (1.5pt) node[above=11pt] {$O$};
\end{tikzpicture}}

\registerfaithful{ex-1-15.png}{%
\begin{tikzpicture}[x=1cm,y=1cm,>=Latex]
  % Choose a=2 and r=1.1, satisfying a<2r<2a.  Both balls are tangent
  % to their adjacent wall and exactly tangent to each other.
  \draw[Ink,very thick] (-2.00,-1.55)--(-2.00,3.45) (2.00,-1.55)--(2.00,3.45); \draw[Ink,thick] (-2.50,-1.55)--(2.50,-1.55);
  \foreach \x in {-2.4,-2.0,...,2.4}{\draw[Ink] (\x,-1.55)--++(-0.22,-0.24);}
  \coordinate (O) at (-0.90,-0.45); \coordinate (Op) at (0.90,0.815);
  \draw[Ink,thick] (O) circle (1.10); \draw[Ink,thick] (Op) circle (1.10);
  \fill[Ink] (O) circle (1.6pt) node[below] {$O$}; \fill[Ink] (Op) circle (1.6pt) node[above right] {$O'$};
  \draw[<->,Ink] (-2.00,3.02)--node[above] {$2a$}(2.00,3.02); \node[below=9pt] at (-2.00,-1.55) {$A$}; \node[below=9pt] at (2.00,-1.55) {$B$};
  \draw[Ink,->] (Op)--++(130:1.10);
  \node at ($(Op)+(130:.55)+(220:.25)$) {$r$};
\end{tikzpicture}}

\registerfaithful{ex-1-16.png}{%
\begin{tikzpicture}[x=1cm,y=1cm,>=Latex]
  \draw[Ink,thick] (-3.4,-2.1)--(-3.4,2.8); \foreach \y in {-1.9,-1.5,...,2.7}{\draw[Ink] (-3.4,\y)--++(-0.28,0.22);}
  \coordinate (A) at (-3.4,-1.55); \coordinate (E) at ($(A)+(25:6.60)$); \draw[Ink,very thick] (A)--(E);
  \draw[Ink,very thick] (-3.4,{6.60*sin(25)-1.55})--(E); \draw[->,Ink,thick] (E)--++(180:3.1) node[midway,above] {$F$};
  % r=.55; s1=r(1+sin25)/cos25, so cylinder 1 is tangent to both wall and board.
  \pgfmathsetmacro{\sone}{.55*(1+sin(25))/cos(25)}
  \foreach \i in {0,1,2}{\coordinate (C\i) at ($(A)+(25:{\sone+1.10*\i})+(115:.55)$); \draw[Ink,thick] (C\i) circle (.55);}
  \node[below left=7pt] at (A) {$A$};
  % Attach r to an actual radius, so it cannot be read as an F subscript.
  \draw[Ink,->] (C2)--++(70:.55) node[midway,left] {$r$};
  \node[below] at (1.75,-0.02) {$l$};
  \draw[Ink] ($(A)+(90:.34)$) arc[start angle=90,end angle=25,radius=.34];
  \draw[TextGray,thin] ($(A)+(57.5:.34)$)--(-4.00,-1.20);
  \node[left] at (-4.00,-1.20) {$\theta$};
\end{tikzpicture}}

\registerfaithful{ex-1-17.png}{%
\begin{tikzpicture}[x=1cm,y=1cm]
  \draw[Ink,thick] (-4.2,-1.35)--(4.2,-1.35);
  % Each wall is a 60-degree line; the end tangent differs from it by exactly 30 degrees.
  \draw[Ink,very thick] (-2.20,-1.35)--(-3.63,1.12) (2.20,-1.35)--(3.63,1.12);
  \draw[Ink,thick,domain=-3.63:3.63,samples=90] plot (\x,{.07948*\x*\x+.073});
  \draw[Ink,thick,domain=-3.63:3.63,samples=90] plot (\x,{1.65-.0400*\x*\x});
  \foreach \x in {-3.0,-2.0,-1.0,0,1.0,2.0,3.0}{\draw[Ink] (\x,{1.65-.0400*\x*\x})--(\x,{.07948*\x*\x+.073});}
  \coordinate (Ewall) at (3.63,1.12);
  \draw[Ink] ($(Ewall)+(210:.58)$) arc[start angle=210,end angle=240,radius=.58];
  \node at ($(Ewall)+(225:.82)$) {$30^\circ$};
  \node[above] at (0,1.70) {屋面}; \node[below] at (0,0.073) {钢索}; \node[right] at (3.34,0.45) {斜壁};
\end{tikzpicture}}

\registerfaithful{fig-1-2.png}{%
\begin{tikzpicture}[x=1cm,y=1cm,>=Latex]
  \draw[Ink,thick] (-2.8,-1.15)--(2.8,-1.15); \foreach \x in {-2.6,-2.2,...,2.6}{\draw[Ink] (\x,-1.15)--++(-0.22,-0.22);}
  \draw[Ink,thick] (-0.50,-1.15) rectangle (0.72,-0.05);
  \coordinate (P) at (0.10,-0.05); \draw[->,Ink,thick] (P)--(0.10,1.72) node[right] {$N$}; \draw[->,Ink,thick] (P)--(-1.38,-0.05) node[below] {$f$};
  \draw[->,DeepBlue,very thick] (P)--(-1.38,1.72) node[above left] {$F$}; \draw[TextGray,dashed] (-1.38,1.72)--(0.10,1.72) (P)--(-1.38,-0.05);
  \pgfmathsetmacro{\forceangle}{atan2(1.77,-1.48)}
  \draw[Ink] ($(P)+(90:.92)$) arc[start angle=90,end angle=\forceangle,radius=.92];
  \node at ($(P)+({(90+\forceangle)/2}:1.20)$) {$\alpha$};
\end{tikzpicture}}

\registerfaithful{ex-1-18.png}{%
\begin{tikzpicture}[x=1cm,y=1cm,>=Latex]
  \draw[Ink,thick] (-3,-1.0)--(3,-1.0); \foreach \x in {-2.8,-2.4,...,2.8}{\draw[Ink] (\x,-1)--++(-0.22,-0.22);}
  \draw[Ink,thick,fill=SoftGray] (-1.30,-1) rectangle (0.25,0.25); \node at (-0.55,-0.35) {$M$};
  \draw[->,DeepBlue,very thick] (0.25,0.15)--(2.28,1.45) node[right] {$F$}; \draw[TextGray,dashed] (0.25,0.15)--(2.28,0.15);
  \draw[Ink] (1.05,0.15) arc[start angle=0,end angle=33,radius=0.80]; \node at (1.15,0.48) {$\alpha$};
\end{tikzpicture}}

\registerfaithful{ex-1-19.png}{%
\begin{tikzpicture}[x=1cm,y=1cm]
  \coordinate (O) at (-3.4,0); \coordinate (C) at (0.851,0);
  \draw[Ink,very thick] (O)--++(15:7.1) (O)--++(-15:7.1);
  \draw[Ink,thick] (C) circle (1.10); \fill[Ink] (O) circle (1.8pt) node[above left] {$O$};
  \draw[Ink] ($(O)+(-15:1.45)$) arc[start angle=-15,end angle=15,radius=1.45];
  \node at ($(O)+(0:1.78)$) {$\theta$};
  % The radius to the upper contact is normal to the 15-degree face.
  \draw[Ink,->] (C)--++(105:1.10) node[midway,right] {$r$};
\end{tikzpicture}}

\registerfaithful{ex-1-20.png}{%
\begin{tikzpicture}[x=1cm,y=1cm]
  \draw[Ink,thick] (-3.1,-1.55)--(-3.1,2.35); \foreach \y in {-1.4,-1.0,...,2.2}{\draw[Ink] (-3.1,\y)--++(-0.25,0.22);}
  \coordinate (A) at (-3.1,-0.72); \coordinate (B) at (2.65,-0.72); \coordinate (C) at (-3.1,1.80);
  \draw[Ink,very thick] (A)--(B); \draw[Ink,thick] (C)--(B); \fill[Ink] (0.0,-0.72) circle (1.8pt) node[below] {$P$};
  \node[left=9pt] at (A) {$A$}; \node[right] at (B) {$B$}; \node[left=9pt] at (C) {$C$};
  \draw[<->,Ink] (-3.0,-0.35)--node[above] {$x$}(0.0,-0.35); \draw[Ink] (1.72,-0.72) arc[start angle=180,end angle=156,radius=0.93]; \node at (1.55,-0.40) {$\theta$};
\end{tikzpicture}}

\registerfaithful{ex-1-21.png}{%
\begin{tikzpicture}[x=1cm,y=1cm,>=Latex]
  \draw[Ink,thick] (-3.1,-1.55)--(3.1,-1.55); \foreach \x in {-2.9,-2.5,...,2.9}{\draw[Ink] (\x,-1.55)--++(-0.22,-0.22);}
  \draw[Ink,very thick] (-2.45,-1.55)--(2.15,0.70); \draw[->,DeepBlue,very thick] (2.15,0.70)--(1.40,2.23) node[above left] {$F$};
  \draw[Ink] (-1.45,-1.55) arc[start angle=0,end angle=26,radius=1.0]; \node at (-1.16,-1.30) {$\theta$};
\end{tikzpicture}}

\registerfaithful{ex-1-22.png}{%
\begin{tikzpicture}[x=1cm,y=1cm]
  \draw[Ink,thick] (-3.4,1.10)--(3.4,1.10); \fill[Ink] (-2.60,1.10) circle (2.4pt) node[above] {$A$} node[below left] {$m$}; \fill[Ink] (2.60,1.10) circle (2.4pt) node[above] {$B$} node[below right] {$m$};
  \coordinate (C) at (0,-0.55); \draw[Ink,thick] (-2.60,1.10)--(C)--(2.60,1.10); \draw[Ink,thick] (C)--(0,-1.35);
  \draw[fill=SoftGray,draw=Ink] (-0.43,-1.90) rectangle (0.43,-1.35) node[midway] {$M$};
  \draw[<->,Ink] (-2.60,1.52)--node[above] {$x$}(2.60,1.52);
\end{tikzpicture}}

\registerfaithful{fig-1-3.png}{%
\begin{tikzpicture}[x=1cm,y=1cm,>=Latex]
  \begin{scope}[shift={(-4.1,0)}]
    \node[font=\small] at (0,3.00) {(a) 竖直绳};
    \draw[Ink,thick] (-0.55,2.0)--(0.55,2.0); \draw[Ink,very thick] (0,2.0)--(0,-2.0); \node[above] at (0,2.05) {$O$};
    % Separate the position coordinate, total length and rope mass so none
    % of the three labels shares the same line or label box.
    \draw[<->,Ink] (-0.62,2.0)--node[midway,left] {$x$}(-0.62,0.35);
    \draw[TextGray,dashed] (-0.62,0.35)--(0,0.35);
    \fill[Ink] (0,0.35) circle (1.5pt);
    \draw[<->,Ink] (-1.30,2.0)--node[midway,left] {$L$}(-1.30,-2.0);
    \node[right] at (0.12,0) {$m$};
  \end{scope}
  \begin{scope}[shift={(2.2,0)}]
    \node[font=\small] at (0,3.00) {(b) 弯绳微元};
    \draw[TextGray,thick,domain=-2.2:1.65,samples=80] plot (\x,{.18*(\x+2.2)^2-1.35});
    \draw[Ink,very thick,domain=-1.55:.70,samples=60] plot (\x,{.18*(\x+2.2)^2-1.35});
    \coordinate (Aseg) at (-1.55,-1.274); \coordinate (Bseg) at (.70,.164);
    \draw[->,Ink,thick] (Aseg)--++(193.2:1.15) node[below left] {$T_A$};
    \draw[->,Ink,thick] (Bseg)--++(46.2:1.25) node[above right] {$T_B$};
    % Place the differential element's weight on the highlighted rope,
    % rather than on the straight chord joining its endpoints.
    \coordinate (Gseg) at (-.425,-.783);
    \draw[->,DeepBlue,thick] (Gseg)--++(0,-1.15) node[right] {$\Delta mg$};
    \fill[Ink] (Aseg) circle (1.6pt) node[above left] {$A$};
    \fill[Ink] (Bseg) circle (1.6pt) node[below right] {$B$};
  \end{scope}
\end{tikzpicture}}

\registerfaithful{ex-1-23.png}{%
\begin{tikzpicture}[x=1cm,y=1cm]
  \draw[TextGray,dashed] (0,3.2)--(0,-2.35); \draw[Ink,thick] (0,3.2)--(-1.75,-1.70) (0,3.2)--(1.75,-1.70);
  \draw[TextGray,dashed] (1.75,-1.70) arc[start angle=0,end angle=180,x radius=1.75,y radius=0.42];
  \draw[Ink,thick] (-1.75,-1.70) arc[start angle=180,end angle=360,x radius=1.75,y radius=0.42];
  \draw[TextGray,dashed,thick] (1.054,0.25) arc[start angle=0,end angle=180,x radius=1.054,y radius=0.27];
  \draw[DeepBlue,very thick] (-1.054,0.25) arc[start angle=180,end angle=360,x radius=1.054,y radius=0.27];
  \pgfmathsetmacro{\coneangle}{atan(1.75/4.90)}
  \draw[Ink] (0,2.15) arc[start angle=-90,end angle={-90+\coneangle},radius=1.05];
  \draw[TextGray,thin] ({1.05*sin(\coneangle/2)},{3.2-1.05*cos(\coneangle/2)})--(1.40,2.40)--(1.65,2.40);
  \node[right,font=\small] at (1.65,2.40) {$\delta=\alpha/2$};
\end{tikzpicture}}

\registerfaithful{ex-1-24.png}{%
\begin{tikzpicture}[x=1cm,y=1cm,>=Latex]
  \draw[Ink,thick] (-3.4,-1.3)--(3.4,-1.3); \foreach \x in {-3.2,-2.8,...,3.2}{\draw[Ink] (\x,-1.3)--++(-0.22,-0.22);}
  \draw[Ink,thick] (-2.25,-1.3) arc[start angle=180,end angle=0,radius=2.25];
  \draw[DeepBlue,line width=1.5pt] (0,0.95) arc[start angle=90,end angle=0,radius=2.25];
  \draw[->,DeepBlue,very thick] (0,0.95)--(-1.85,0.95) node[left] {$F$};
  \coordinate (O) at (0,-1.3); \coordinate (Q) at ($(O)+(135:2.25)$);
  \draw[Ink,->] (O)--(Q);
  \node at ($(O)!.5!(Q)+(225:.30)$) {$R$};
  \fill[Ink] (O) circle (1.8pt) node[below=9pt] {$O$};
\end{tikzpicture}}

\registerfaithful{ex-1-25.png}{%
\begin{tikzpicture}[x=1cm,y=1cm,>=Latex]
  \draw[Ink,thick] (0,0.75) circle (1.65); \fill[Ink] (0,0.75) circle (2pt);
  \draw[Ink,very thick] (-1.65,0.75)--(-1.65,-2.55) (1.65,0.75)--(1.65,-2.55);
  \coordinate (O) at (0,.75); \coordinate (Q) at ($(O)+(135:1.65)$);
  \draw[Ink,->] (O)--(Q); \node at ($(O)!.5!(Q)+(225:.26)$) {$R$};
\end{tikzpicture}}

\registerfaithful{ex-1-26.png}{%
\begin{tikzpicture}[x=1cm,y=1cm]
  \draw[Ink,thick] (-3.0,-2.40)--(3.0,-2.40);
  \foreach \x in {-2.8,-2.4,...,2.8}{\draw[Ink] (\x,-2.40)--++(-.22,-.22);}
  \draw[Ink,thick] (0,0) circle (2.40);
  % The centre line is a construction axis, not a physical rod.
  \draw[TextGray,dashed] (0,2.65)--(0,-2.40);
  % Back and front halves of the horizontal ring must both be visible.
  \draw[TextGray,dashed,thick] (1.697,1.697) arc[start angle=0,end angle=180,x radius=1.697,y radius=0.38];
  \draw[DeepBlue,very thick] (-1.697,1.697) arc[start angle=180,end angle=360,x radius=1.697,y radius=0.38];
  \coordinate (O) at (0,0); \coordinate (Q) at (155:2.40);
  \draw[Ink,->] (O)--(Q); \node at ($(O)!.5!(Q)+(245:.28)$) {$R$};
  \fill[Ink] (O) circle (2pt) node[below right] {$O$};
  \draw[<->,Ink] (0,1.697)--node[above] {$b$}(1.697,1.697);
\end{tikzpicture}}

\registerfaithful{ex-1-27.png}{%
\begin{tikzpicture}[x=1cm,y=1cm]
  \draw[Ink,thick] (-3.4,1.25)--(3.4,1.25); \node[above] at (-3.0,1.25) {$A$}; \node[above] at (3.0,1.25) {$B$};
  \draw[Ink,very thick,domain=-3:3,samples=101] plot (\x,{1.25+3.4*(cosh(\x/3.4)-cosh(3/3.4))});
  \node[above] at (0,1.25) {$d$};
\end{tikzpicture}}

\registerfaithful{fig-1-4.png}{%
\begin{tikzpicture}[x=1cm,y=1cm]
  \begin{scope}[shift={(-4.5,0)}] \draw[Ink,thick] (-0.7,2.2)--(0.7,2.2); \foreach \x in {-0.6,-0.2,0.2,0.6}{\draw[Ink] (\x,2.2)--++(0.22,0.25);} \fill[SoftGray,draw=Ink] (0,1.95) circle (0.15); \draw[Ink,very thick] (0,1.80)--(0,-1.75); \end{scope}
  \begin{scope}[shift={(-1.5,0)}] \draw[Ink,thick] (-0.7,-1.75)--(0.7,-1.75); \foreach \x in {-0.6,-0.2,0.2,0.6}{\draw[Ink] (\x,-1.75)--++(-0.22,-0.25);} \fill[SoftGray,draw=Ink] (0,-1.55) circle (0.15); \draw[Ink,very thick] (0,-1.40)--(0,2.15); \end{scope}
  \begin{scope}[shift={(1.5,0)}] \draw[Ink,thick] (-0.7,2.2)--(0.7,2.2); \foreach \x in {-0.6,-0.2,0.2,0.6}{\draw[Ink] (\x,2.2)--++(0.22,0.25);} \fill[SoftGray,draw=Ink] (0,1.95) circle (0.15); \draw[TextGray,dashed] (0,1.95)--(0,-0.7); \draw[Ink,very thick,rotate around={16:(0,1.95)}] (0,1.80)--(0,-1.75); \end{scope}
  \begin{scope}[shift={(4.5,0)}] \draw[Ink,thick] (-0.7,-1.75)--(0.7,-1.75); \foreach \x in {-0.6,-0.2,0.2,0.6}{\draw[Ink] (\x,-1.75)--++(-0.22,-0.25);} \fill[SoftGray,draw=Ink] (0,-1.55) circle (0.15); \draw[TextGray,dashed] (0,-1.55)--(0,0.85); \draw[Ink,very thick,rotate around={-16:(0,-1.55)}] (0,-1.40)--(0,2.15); \end{scope}
\end{tikzpicture}}

\registerfaithful{ex-1-28.png}{%
\begin{tikzpicture}[x=1cm,y=1cm,decoration={coil,aspect=0.45,segment length=4pt,amplitude=3pt}]
  \draw[Ink,very thick] (0,-1.75)--(0,1.50); \fill[Ink] (0,-1.75) circle (2.2pt); \node[right] at (0,0) {$m,L$};
  \draw[Ink,thick] (-3.9,2.10)--(-3.9,0.95) (3.9,2.10)--(3.9,0.95); \foreach \y in {1.05,1.35,...,2.05}{\draw[Ink] (-3.9,\y)--++(-0.25,0.20);\draw[Ink] (3.9,\y)--++(0.25,0.20);}
  \draw[Ink,thick,decorate] (-3.9,1.50)--(0,1.50); \draw[Ink,thick,decorate] (0,1.50)--(3.9,1.50); \node[above] at (-2.0,1.72) {$k$}; \node[above] at (2.0,1.72) {$k$};
  \draw[Ink,thick] (-0.85,-1.75)--(0.85,-1.75); \foreach \x in {-0.75,-0.35,0.05,0.45}{\draw[Ink] (\x,-1.75)--++(-0.22,-0.22);}
\end{tikzpicture}}

\registerfaithful{ex-1-29.png}{%
\begin{tikzpicture}[x=1cm,y=1cm,decoration={coil,aspect=0.45,segment length=4pt,amplitude=3pt}]
  \coordinate (T) at (0,2.55); \coordinate (A) at (-2.25,-1.45); \coordinate (B) at (2.25,-1.45);
  \draw[Ink,very thick] (A)--(T)--(B);
  \filldraw[fill=SoftGray,draw=Ink] (T) circle (0.18);
  \node[above=4pt,text=Ink,inner sep=0pt] at (T) {$M$};
  \draw[Ink,thick,decorate] (A)--(B); \node[above=5pt] at (0,-1.45) {$k$}; \node[left=5pt] at (-1.15,0.45) {$L$}; \node[right=5pt] at (1.15,0.45) {$L$};
  % The lower ends are sliders on a smooth horizontal support.  Drawing
  % the support here keeps the problem figure and the solution model equal.
  \foreach \X in {-2.37,-2.13,2.13,2.37}{\draw[Ink] (\X,-1.64) circle (0.08);}
  \draw[Ink,thick] (-2.70,-1.76)--(2.70,-1.76);
  \foreach \x in {-2.55,-2.15,...,2.55}{\draw[Ink] (\x,-1.76)--++(-0.18,-0.18);}
  \draw[TextGray,dashed] (T)--(0,-0.10); \draw[Ink] (0,1.55) arc[start angle=-90,end angle=-60,radius=1.0]; \node at (0.36,1.38) {$\theta$}; \node[left] at (A) {$A$}; \node[right] at (B) {$B$};
\end{tikzpicture}}

\registerfaithful{ex-1-30.png}{%
\begin{tikzpicture}[x=1cm,y=1cm,>=Latex]
  \draw[axis] (-0.25,0)--(4.65,0) node[right] {$x$}; \draw[axis] (0,-0.25)--(0,4.60) node[above] {$y$};
  \def\rodlength{4.4}\def\xp{3.20}
  \pgfmathsetmacro{\yp}{(\rodlength-sqrt(\rodlength^2-\xp^2))/2}
  \draw[Ink,very thick,domain=0:4.38,samples=101] plot (\x,{(\rodlength-sqrt(\rodlength^2-\x^2))/2}) node[right] {$y(x)$};
  \coordinate (P) at (\xp,\yp); \coordinate (A) at (0,{\rodlength-\yp});
  \draw[Ink,very thick] (A)--(P) node[midway,above right=4pt] {$L$};
\end{tikzpicture}}

\newcommand{\faithfulExOneThirtyOne}{%
\begin{tikzpicture}[x=0.72cm,y=0.72cm,>=Latex]
  % The problem figure shows the reference equilibrium position.  The
  % displaced rolling state is reserved for the solution diagram.
  \def\Rbow{3.05}\def\rball{0.86}
  \pgfmathsetmacro{\aradius}{\Rbow-\rball}
  \coordinate (C) at (0,1.15);
  \coordinate (O) at ($(C)+(-90:\aradius)$);
  \coordinate (K) at ($(C)+(-90:\Rbow)$);
  % At the reference equilibrium position G lies vertically above O.
  \coordinate (G) at ($(O)+(90:0.54)$);
  \draw[Ink,very thick] ($(C)+(180:\Rbow)$) arc[start angle=180,end angle=360,radius=\Rbow];
  \draw[Ink,thick] (O) circle (\rball);
  \fill[Ink] (O) circle (1.8pt) node[below left=4pt] {$O$};
  \fill[Ink] (G) circle (1.8pt);
  \draw[TextGray,thin] (G)--($(O)+(-1.25,1.10)$);
  \node[left,inner sep=1pt] at ($(O)+(-1.25,1.10)$) {$G$};
  \draw[TextGray,dashed] (O)--($(O)+(1.18,0)$);
  \draw[TextGray,dashed] (G)--($(G)+(1.18,0)$);
  \draw[<->,Ink] ($(O)+(1.18,0)$)--($(G)+(1.18,0)$)
    node[midway,right=2pt,inner sep=0pt] {$h$};
  \coordinate (Qr) at ($(O)+(135:\rball)$);
  \draw[Ink,->] (O)--(Qr);
  \node[font=\small] at ($(O)!.53!(Qr)+(225:.23)$) {$r$};
  \fill[Ink] (K) circle (1.3pt);
  \draw[Ink,->] (C)--++(180:\Rbow) node[midway,above] {$R$};
  \fill[Ink] (C) circle (1.8pt) node[above right=2pt] {$C$};
\end{tikzpicture}}


在多个力作用下处于平衡状态的系统，所受到各个力的大小和方向必然有一定的联系。因
为力的效果总是使物体的运动状态发生变化，而处于静力平衡的物体的运动状态没有发生
变化，所以它受到的合力必然为零。对这一事实的描写有两种方式。一种是直接的矢量计算，
画出各个力的大小和方向，对于处于平衡状态系统中的每一个独立的部分，处于静力平衡时
所受外力的矢量和必然为零，利用此方法更多地需要几何知识和技巧。第二种常用的方法则
是根据问题的特点将各个物体所受的外力向两到三个彼此正交的方向分解，平衡时每个方
向上的合力分量必然为零，即各外力在该方向上的分量之和为零，而不是每一个分力均为零。
当以这种方式处理问题时，解方程的能力则是必不可少的。

\sourcefigure[0.54\linewidth]{fig-1-1.png}
\origcaption{力矩的定义}\label{ch05:fig:s03-1-1}

力的效果除了让物体有沿力的方向运动的趋势以外，对于刚体来说还会引起转动的趋势，
力学中用力矩定量地描写力引起的转动效果。对于如图所示的刚性杆，其一端固定于 $O$ 点，
在杆上 $A$ 点处有外力 $\vec F$ 的作用。令 $\vec r=\overrightarrow{OA}$，力对 $O$ 点的力矩先按
矢量定义为
\begin{equation}
  \vec M_O=\vec r\times\vec F.
\end{equation}
若只讨论图示平面内的转动，令 $\theta$ 为杆的作用线与力的作用线的较小夹角，则力矩大小为
\begin{equation}
  M_O=Fr_{\perp}=F_{\perp}r=Fr\sin\theta.
\end{equation}

各个量已在图中标出，其中 $O$ 称为定义力矩的参考点，$A$ 为力的作用点；从 $O$ 向力
$\vec F$ 的作用线所作垂线段的长度 $OB=r_\perp$ 称为力臂。除了具有一定的大
小，力矩其实也是有方向的，对于定轴转动的物体来说，力矩的方向可以用它使物体有顺时
针或逆时针转动的趋势来描写。当一个物体受到多个力作用时，合力矩则是各个分力力矩之
和。因为力矩使物体有转动的趋势，所以处于静力平衡的系统中每个物体所受的合力矩也必须
为零。需要注意的是和力不同，力矩的定义除了要指定力的大小和方向以外还有一个转动的
参考点 $O$，$O$ 点的选取并不一定必须在转动轴，甚至也不一定必须在受力物体上，无论将
力矩的参考点选在何处，静力平衡均要求合力矩为零。合理地选择参考点可以极大地简化受
力分析的计算量。

总的来说，处于平衡状态的力学系统中，每个可独立研究的部分都需要满足合外力为零，相对任意
参考点的合力矩为零。

\begin{origexamplebox}{}\label{ch05:ex:s03-1-10}
如图三根长度均为 $l$ 的轻杆用铰链连接并固定在水平天花板上的 $A$、$B$ 两点，$AB$
两点相距 $2l$，若在铰链 $C$ 上悬挂一个质量为 $m$ 的重物，要使 $CD$ 杆保持水平，则在 $D$ 点
上应施加的最小力为多少？
\sourcefigurefixB
\end{origexamplebox}

\Needspace{6\baselineskip}\begin{xesolution}
由 $AB=2l$、$CD=l$ 以及三杆等长可知，左右两根斜杆与水平方向的夹角均为
$60^\circ$。先研究铰链 $C$。设杆 $AC$ 中的轴力为 $S$，杆 $CD$ 中的轴力为 $H$。解图把
$C$、$D$ 两个铰链的受力分开画出：
\soldiagOneTen
由解图 (a) 可得
\[
  S\sin60^\circ=mg,
  \qquad
  H=S\cos60^\circ=\frac{mg}{\sqrt3}.
\]
在铰链 $D$ 处，杆 $DB$ 只能提供沿 $DB$ 方向的力。因此，外力中垂直于 $DB$ 的分量
必须抵消 $H$ 垂直于 $DB$ 的分量。外力与 $DB$ 垂直时，其大小最小，故
\[
  F_{\min}=H\sin60^\circ=\boxed{\frac{mg}{2}}.
\]
\end{xesolution}

\begin{origexamplebox}{}\label{ch05:ex:s03-1-11}
质量为 $m$ 的小球静止在固定斜面与可调整方向的挡板之间，两处接触均光滑。下斜面与水平面
夹角 $\alpha$ 固定，且 $0<\alpha<90^\circ$；挡板与下斜面之间的夹角为 $\theta$。
在能够保持两处接触的范围内，求挡板在何角度时对小球的压力最小。
\sourcefiguretrim[0.40\linewidth]{ex-1-11.png}{0 0 0 65}
\end{origexamplebox}

\Needspace{62mm}
\Needspace{6\baselineskip}\begin{xesolution}
设下方斜面和挡板对球的压力分别为 $N_1$、$N$。三个力的方向及投影方向如解图所示：
\soldiagOneEleven
$N_1$ 垂直于下方斜面，沿该斜面方向的分量为零。因此沿斜面向上取正方向，
$N$ 的分量 $N\sin\theta$ 与重力分量 $mg\sin\alpha$ 平衡：
\[
  N\sin\theta=mg\sin\alpha,
  \qquad
  N=\frac{mg\sin\alpha}{\sin\theta}.
\]
当 $\sin\theta$ 取最大值 $1$ 时压力最小，即
\[
  \boxed{\theta=90^\circ},
  \qquad
  \boxed{N_{\min}=mg\sin\alpha}.
\]
还需检查这个角度能够实现：$\theta=90^\circ$ 时挡板压力恰沿下斜面向上，下斜面压力
为 $N_1=mg\cos\alpha>0$，两处接触均不需要“拉住”小球，所以最小值确实可取。
\end{xesolution}

\begin{origexamplebox}{}\label{ch05:ex:s03-1-12}
一质量为 $m$ 的匀质球静止于两固定斜面之间，如图所示。下斜面与水平面的夹角为
$\theta_1$，上斜面与同一水平面的夹角为 $\theta_2$，其中 $0<\theta_1<\theta_2<\pi/2$。
所有接触面都光滑，求斜面作用于球上的力。
\sourcefigure[0.43\linewidth]{ex-1-12.png}
\end{origexamplebox}

\Needspace{6\baselineskip}\begin{xesolution}
按题图，两个倾角都从水平面量起，因而两斜面之间的夹角是 $\theta_2-\theta_1$。
设下、上斜面对球的压力分别为 $N_1$、$N_2$，三个力
的方向如解图所示：
\soldiagOneTwelve
取水平向右、竖直向上为正方向，有
\[
  -N_1\sin\theta_1+N_2\sin\theta_2=0,
\]
\[
  N_1\cos\theta_1-N_2\cos\theta_2-mg=0.
\]
由第一式得 $N_2=N_1\sin\theta_1/\sin\theta_2$，代入第二式。利用
\[
 \cos\theta_1\sin\theta_2-\sin\theta_1\cos\theta_2=\sin(\theta_2-\theta_1),
\]
得
\[
  \boxed{N_1=\frac{mg\sin\theta_2}{\sin(\theta_2-\theta_1)}},
  \qquad
  \boxed{N_2=\frac{mg\sin\theta_1}{\sin(\theta_2-\theta_1)}}.
\]
两力均为正，符合光滑接触只能推不能拉的条件；量纲都是力。
当下斜面趋于水平，即 $\theta_1\to0$ 且 $\theta_2$ 保持不变时，
$N_1\to mg$、$N_2\to0$，与水平面独自承重相符。
\end{xesolution}

\begin{origexamplebox}{}\label{ch05:ex:s03-1-13}
两个质量为 $M$，半径为 $R$ 的相同圆球 $A$ 和 $B$，用两根长为 $l$（$l=2R$）的绳悬
挂于 $O$ 点。这里绳长 $l$ 指从 $O$ 到球面连接点的长度，连接点处的半径与绳共线，
所以 $O$ 到大球球心的距离为 $l+R$。在两球上另有一质量为 $m$（$m=nM$），半径为
$r$（$r=R/2$）的圆球 $C$，如图，
已知三球的表面光滑，试讨论此系统处于平衡时，绳与竖直线的夹角 $\theta$ 与 $n$ 的关系。
\sourcefigure[0.43\linewidth]{ex-1-13.png}
\end{origexamplebox}

\Needspace{6\baselineskip}\begin{xesolution}
把图中 $O$ 点到大球球心的距离记为 $d_0$。由题设的绳长定义和 $l=2R$，有
$d_0=l+R=3R$。

接触几何和一个大球的受力分别画在解图 (a)、(b) 中：
\soldiagOneThirteen

两大球接触时，其球心间距为 $2R$，所以
\[
  2d_0\sin\theta_0=2R,
  \qquad
  \theta_0=\arcsin\frac{R}{d_0}.
\]
按图示取 $d_0=3R$，有 $\theta_0=\arcsin(1/3)$。设小球与大球球心连线同水平方向
的夹角为 $\beta$。当两大球接触时
\[
  \cos\beta_0=\frac{R}{R+r}=\frac23,
  \qquad
  \sin\beta_0=\frac{\sqrt5}{3}.
\]
设小球与每个大球间的压力为 $N$，则小球平衡给出
\[
  2N\sin\beta=nMg.
\]
再研究一个大球。若两大球之间的压力为 $Q$，绳中张力为 $T$，则
\[
  T\cos\theta=Mg+N\sin\beta=Mg\left(1+\frac n2\right),
\]
\[
  T\sin\theta=N\cos\beta+Q.
\]
在两大球刚好失去接触时 $Q=0$，于是临界质量比为
\[
  n_{\rm c}
  =\frac{2\tan\theta_0}{\cot\beta_0-\tan\theta_0}.
\]
按 $d_0=3R$ 计算，
\[
  \boxed{n_{\rm c}=\frac{2\sqrt5}{4\sqrt2-\sqrt5}\approx1.31}.
\]
因此 $0\le n\le n_{\rm c}$ 时，两大球可以保持接触，且
\[
  \boxed{\theta=\arcsin\frac13}.
\]
若两大球分离，则 $Q=0$，几何关系与平衡关系分别为
\[
  \cos\beta=\frac{d_0\sin\theta}{R+r},
  \qquad
  \boxed{n=\frac{2\tan\theta}{\cot\beta-\tan\theta}}.
\]
按 $d_0=3R$、$r=R/2$，这可写成
\[
  \boxed{n=\frac{2\sqrt{1-4\sin^2\theta}}
  {2\cos\theta-\sqrt{1-4\sin^2\theta}}},
  \qquad
  \arcsin\frac13<\theta<30^\circ.
\]
下面确定分离支的存在范围。记 $w=\sqrt{1-4\sin^2\theta}>0$，则
\[
 n=\frac{2w}{2\cos\theta-w},\qquad
 w'=-\frac{4\sin\theta\cos\theta}{w}.
\]
分母为正；对 $n$ 求导并合并分子，得到
\[
 \frac{\mathrm dn}{\mathrm d\theta}
 =\frac{4(w'\cos\theta+w\sin\theta)}{(2\cos\theta-w)^2}
 =-\frac{12\sin\theta}{w(2\cos\theta-w)^2}<0.
\]
当 $\theta$ 从 $\theta_0$ 增至 $30^\circ$ 时，$n$ 从 $n_{\rm c}$ 单调减至零。
所以每个 $0<n<n_{\rm c}$ 都有一个接触构型和一个分离构型，而非只有一个解。
$n=n_{\rm c}$ 时分离支到达接触边界；$n>n_{\rm c}$ 时两支都不能托住小球。

还要判断分离构型能否抵抗小扰动。取两大球对称地分开、但小球仍同时接触两大球的扰动，
以悬点 $O$ 为高度原点。解图 (a) 中 $OA=d_0$，因此两个大球球心的高度均为
$z_A=-3R\cos\theta$。小球球心比大球球心高
\[
 (R+r)\sin\beta=\frac{3R}{2}\sqrt{1-4\sin^2\theta}=\frac{3R}{2}w,
\]
所以 $z_C=-3R\cos\theta+3Rw/2$。把三球重力势能相加：
\[
 U=2Mgz_A+nMgz_C
 =MgR\left[-3(2+n)\cos\theta+\frac{3n}{2}w\right].
\]
其中绳长、半径和质量比在扰动中保持不变。定义
$A(\theta)=2\cos\theta/w-1$，则
\[
 \frac{U'}{MgR}=3\sin\theta[2-nA(\theta)],\qquad
 A'(\theta)=\frac{6\sin\theta}{w^3}.
\]
在分离平衡点 $2-nA=0$，故
\[
 \frac{U''}{MgR}=-3n\sin\theta A'
 =-\frac{18n\sin^2\theta}{w^3}<0.
\]
势能在该方向上为局部极大值，稍有偏离便可继续降低势能，因此分离支是不稳定平衡。
在 $n=n_{\rm c}$ 的接触边界，从允许的分离方向扰动也会降低势能，临界构型同样不能抵抗这种扰动。
$n>n_{\rm c}$ 时，小球会把两大球推开并从中间落下。
\end{xesolution}

\begin{origexamplebox}{}\label{ch05:ex:s03-1-14}
如图所示一个均匀的质量为 $m_1$、半径为 $R$ 的球挂在天花板上 $O$ 点处，从 $O$ 点
到球心的距离已知为 $l$。另一根轻且不可伸长的绳从同一点引出，与球面光滑接触后
竖直悬挂质量为 $m_2$ 的重物。求图中 $\theta$ 角的大小。
\sourcefigure[0.36\linewidth]{ex-1-14.png}
\end{origexamplebox}

\Needspace{6\baselineskip}\begin{xesolution}
设悬挂 $m_2$ 的绳中张力为 $T=m_2g$，球心到 $O$ 的连线与绳的上方切线之间夹角为
$\delta$。由切线性质
\[
  \sin\delta=\frac{R}{l}.
\]
再设上方切线与竖直方向的夹角为 $\psi$，由图中几何关系 $\delta=\theta+\psi$。
球与这段绕绳之间的合力，等于两端张力矢量之和。设悬球绳的张力为 $S$，对球列水平、
竖直方向的平衡方程。相应的几何关系与合力图如下：
\soldiagOneFourteen
于是
\[
  S\sin\theta=T\sin\psi,
\]
\[
  S\cos\theta=m_1g+T(1-\cos\psi).
\]
消去 $S$ 并代入 $T=m_2g$，得到
\[
  (m_1+m_2)\sin\theta
  =m_2\sin(\theta+\psi)
  =m_2\sin\delta
  =m_2\frac{R}{l}.
\]
故
\[
  \boxed{\theta=\arcsin\!\left(\frac{m_2R}{(m_1+m_2)l}\right)}.
\]
\end{xesolution}

\begin{origexamplebox}{}\label{ch05:ex:s03-1-15}
有一个半径为 $a$，高为 $4a$，质量为 $M$ 的两端开口的薄壁圆筒，现将筒竖放在光
滑的水平面上，之后将半径为 $r$，质量为 $m$ 的两个完全相同的光滑圆球放入筒内而呈叠放
状态，当 $a<2r<2a$ 时，试求使圆筒不翻倒的条件。
\sourcefigure[0.46\linewidth]{ex-1-15.png}
\end{origexamplebox}

\Needspace{6\baselineskip}\begin{xesolution}
设两球球心连线与水平方向的夹角为 $\beta$。由几何关系
\[
  2r\cos\beta=2(a-r),
  \qquad
  \cos\beta=\frac{a-r}{r}.
\]
筒壁所受的两水平力及其高度差如解图所示：
\soldiagOneFifteen
研究上方小球。设两球间压力为 $N$、右壁压力为 $N_2$，则
\[
  N\sin\beta=mg,
  \qquad
  N_2=N\cos\beta=mg\cot\beta.
\]
下方小球对左壁的水平压力也等于 $N_2$。这两个大小相等、方向相反的水平力作用高度
相差 $2r\sin\beta$，对圆筒形成的力偶矩为
\[
  \mathcal M=N_2(2r\sin\beta)
  =2mgr\cos\beta=2mg(a-r).
\]
圆筒将要翻倒时，地面对筒的支持力作用线移到筒底边缘，筒重能够提供的最大抗倾覆
力矩为 $Mga$。所以不翻倒的条件为
\[
  2mg(a-r)\le Mga,
  \qquad
  \boxed{M\ge \frac{2m(a-r)}{a}}.
\]
\end{xesolution}

\begin{origexamplebox}{}\label{ch05:ex:s03-1-16}
有一轻木板，其自重可忽略，长为 $l$，$A$ 端用铁链固定在竖直墙面上，另一端用水
平绳拉住，板上依次放着三个圆柱体，其半径均为 $r$，质量均为 $m$，木板与墙面的夹角为 $\theta$，
一切摩擦均忽略不计，求水平绳对板的拉力 $F$ 多大？
\sourcefigure[0.44\linewidth]{ex-1-16.png}
\end{origexamplebox}

\Needspace{6\baselineskip}\begin{xesolution}
沿木板向上的单位方向记为 $\boldsymbol t$，板的法线方向记为 $\boldsymbol n$。木板与竖直
墙面的夹角为 $\theta$。设三个圆柱与板的接触点到 $A$ 的距离依次为 $s_1,s_2,s_3$。
由第一个圆柱同时与墙、板相切可得
\[
  s_1\sin\theta-r\cos\theta=r,
  \qquad
  s_1=r\frac{1+\cos\theta}{\sin\theta}=r\cot\frac\theta2.
\]
又因相邻圆柱相切，故
\[
  s_2=s_1+2r,
  \qquad
  s_3=s_1+4r.
\]
这些接触位置及法向力如解图所示：
\soldiagOneSixteen
设板对三个圆柱的支持力依次为 $N_1,N_2,N_3$。后两个圆柱所受的圆柱间作用力沿板
方向，故沿板法线列平衡方程有
\[
  N_2=N_3=mg\sin\theta.
\]
再把三个圆柱看成整体。墙的支持力无竖直分量，所以
\[
  (N_1+N_2+N_3)\sin\theta=3mg,
  \qquad
  N_1+N_2+N_3=\frac{3mg}{\sin\theta}.
\]
对木板取 $A$ 点力矩。水平绳的力臂为 $l\cos\theta$，从而
\[
  Fl\cos\theta=N_1s_1+N_2s_2+N_3s_3
  =\frac{3mg}{\sin\theta}s_1+6mgr\sin\theta.
\]
代入 $s_1$ 后得到
\[
  \boxed{F=\frac{3mgr}{l\cos\theta}
  \left(\frac{1+\cos\theta}{\sin^2\theta}+2\sin\theta\right)}.
\]
上述构型除了要求 $s_3\le l$，还要保证最上方圆柱不与水平绳相交。第三个圆柱的圆心
高度为 $s_3\cos\theta+r\sin\theta$，而水平绳的高度为 $l\cos\theta$，故还需
\[
  l\cos\theta-\bigl(s_3\cos\theta+r\sin\theta\bigr)\ge r.
\]
也就是
\[
  \boxed{l-s_3\ge r\frac{1+\sin\theta}{\cos\theta}}.
\]
这些条件只限制图示构型能否实现，不改变上面的拉力表达式。
\end{xesolution}

\begin{origexamplebox}{}\label{ch05:ex:s03-1-17}
某机场候机楼结构简化图如图所示：候机楼侧壁是倾斜的，用钢索将两边斜壁系
住，在钢索上有许多竖直短钢棒将屋面支撑在钢索上。假设每边斜壁的质量为 $4m$，质量分布
均匀；钢索与屋面（包括短钢棒）的总质量为 $2m$，在地面处用铰链与水平地面连接，钢索固
定于斜壁上端以支撑整个屋面，钢索上端与斜壁的夹角为 $30^\circ$；整个系统左右对称。求斜壁
对钢索的拉力的大小和斜壁与地面的夹角。
\sourcefigure[0.54\linewidth]{ex-1-17.png}
\end{origexamplebox}

\Needspace{6\baselineskip}\begin{xesolution}
设斜壁与地面的夹角为 $\beta$，钢索端部切线与地面的夹角为 $\gamma$。由图示几何关系
\[
  \beta-\gamma=30^\circ.
\]
取一侧斜壁，关键力和角度如解图所示：
\soldiagOneSeventeen
把钢索、屋面和短钢棒作为一个整体。其总重为 $2mg$，由左右两端对称支承，因此每一端
拉力 $T$ 的竖直分量应为 $mg$：
\begin{equation}\label{ch05:ref:roof-vertical}
  T\sin\gamma=mg.
\end{equation}
再研究一侧斜壁。斜壁重力为 $4mg$，对底部铰链取矩；其力臂为斜壁长度的
$\cos\beta/2$，钢索拉力与斜壁的夹角为 $30^\circ$，故
\[
  T\sin30^\circ=\frac{4mg}{2}\cos\beta.
\]
因 $\sin30^\circ=1/2$，得
\begin{equation}\label{ch05:ref:wall-moment}
 T=4mg\cos\beta.
\end{equation}
联立式~\eqref{ch05:ref:roof-vertical}、\eqref{ch05:ref:wall-moment}，并代入 $\gamma=\beta-30^\circ$，得
\[
  4\cos\beta\sin(\beta-30^\circ)=1.
\]
利用积化和差公式，左边为
\[
  4\cos\beta\sin(\beta-30^\circ)
  =2\sin(2\beta-30^\circ)-1.
\]
在图示实际构型 $30^\circ<\beta<90^\circ$ 内，方程化为
$\sin(2\beta-30^\circ)=1$，故
\[
  \boxed{\beta=60^\circ},\qquad \gamma=30^\circ.
\]
最后由式~\eqref{ch05:ref:roof-vertical} 得
\[
  \boxed{T=\frac{mg}{\sin30^\circ}=2mg}.
\]
代回式~\eqref{ch05:ref:wall-moment} 也得到 $T=4mg\cos60^\circ=2mg$，两部分平衡条件一致。
\end{xesolution}

\section{摩擦约束与平衡临界条件}

对于存在静摩擦力的平衡系统，求解平衡问题时有一个需要注意的地方就是静摩擦力
并不是可以无限增加的，而是有一最大值。通常来说，静摩擦力的最大值与相互接触且有相对
滑动趋势的物体之间接触面的性质以及相互之间的压力有关，可写成
\[
  f_{\max}=\mu_{\rm s}N,
\]
其中 $\mu_{\rm s}$ 为静摩擦系数。

在处理此类问题时，一般来说首先假设两者间存在静摩擦力，在静摩擦力和其它所有力
共同作用下保持平衡。但最后还要判断平衡所需的摩擦力是否超过上限：若
$|f|\le\mu_{\rm s}N$，则可能保持静止，其中等号对应临界静止；若
$|f|>\mu_{\rm s}N$，假设的静力平衡才不可能实现，物体将发生相对滑动。

\sourcefigure[0.43\linewidth]{fig-1-2.png}
\origcaption{全反力示意图}\label{ch05:fig:s03-1-2}
图中将接触力矢量平移到共同起点以说明合成关系，不表示这些力实际作用在物块顶面。

除了用摩擦力和压力比值的方法判断是否滑动以外，还可以考虑可能滑动物体与接触
表面之间压力和静摩擦力的合力，有时将其称为全反力。如图~\ref{ch05:fig:s03-1-2} 所示的 $F$ 就是全反力，它
与接触面垂线的夹角 $\alpha$ 是标志全反力的一个重要的物理量。从图中可以看出，当 $\alpha$ 为已知
时，压力 $N=F\cos\alpha$，摩擦力 $f=F\sin\alpha$，所以不发生滑动的条件此时就变成了
\begin{equation}
  \sin\alpha\le\mu_{\rm s}\cos\alpha,\qquad \text{或}\qquad \tan\alpha\le\mu_{\rm s},
\end{equation}
也就是说，全反力和法线夹角的正切不大于静摩擦系数时，接触面可能保持相对静止；等号
对应摩擦角达到极限，超过这一角度才不可能相对静止。

\begin{origexamplebox}{}\label{ch05:ex:s03-1-18}
木箱质量为 $M$，与地面间的动摩擦因数为 $\mu$，用斜向上的力 $F$ 拉木箱使之沿水平
地面匀速前进，如图所示，问角 $\alpha$ 为何值时拉力 $F$ 最小？这个最小值为多大？
\sourcefigure[0.42\linewidth]{ex-1-18.png}
\end{origexamplebox}

\Needspace{6\baselineskip}\begin{xesolution}
木箱受拉力 $F$、重力 $Mg$、地面支持力 $N$ 和滑动摩擦力 $f$，方向如解图所示：
\soldiagOneEighteen
由解图可见，木箱水平向右匀速运动，摩擦力向左，且 $f=\mu N$。水平方向平衡给出
\[
  F\cos\alpha=\mu N.
\]
竖直方向上，拉力的向上分量减小了地面支持力，故
\[
  N+F\sin\alpha=Mg.
\]
消去 $N$，得到
\[
  F=\frac{\mu Mg}{\cos\alpha+\mu\sin\alpha}.
\]
要使 $F$ 最小，应使分母最大。写成
\[
  \cos\alpha+\mu\sin\alpha
  =\sqrt{1+\mu^2}\cos(\alpha-\arctan\mu),
\]
可知
\[
  \boxed{\alpha=\arctan\mu},
  \qquad
  \boxed{F_{\min}=\frac{\mu Mg}{\sqrt{1+\mu^2}}}.
\]
\end{xesolution}

\begin{origexamplebox}{}\label{ch05:ex:s03-1-19}
筷子夹鸡蛋（只考虑在水平光滑桌面上）的摩擦紧锁问题：设筷子与鸡蛋的摩擦
系数为 $\mu$，求筷子多大角度时在水平桌面上始终不会把鸡蛋滑出。
\sourcefigure[0.44\linewidth]{ex-1-19.png}
\end{origexamplebox}

\Needspace{6\baselineskip}\begin{xesolution}
取关于角平分线对称的两个接触点。设每根筷子对鸡蛋的正压力为 $N$。两正压力沿角平分
线向外，而极限摩擦力沿筷子指向 $O$：
\soldiagOneNineteen
由解图可知，两正压力沿角平分线向外的合力为
\[
  2N\sin\frac\theta2.
\]
鸡蛋将要向外滑出时，两处摩擦力均沿筷子指向 $O$，其沿角平分线向内的最大合力为
\[
  2\mu N\cos\frac\theta2.
\]
发生摩擦紧锁要求后者不小于前者，即
\[
  2\mu N\cos\frac\theta2\ge 2N\sin\frac\theta2.
\]
所以
\[
  \boxed{\tan\frac\theta2\le\mu},
  \qquad\text{即}\qquad
  \boxed{\theta\le2\arctan\mu}.
\]
\end{xesolution}

\begin{origexamplebox}{}\label{ch05:ex:s03-1-20}
有一长 $l$ 质量为 $M$ 的均匀杆 $AB$，$A$ 端顶在竖直的粗糙墙壁上，杆端与墙间的摩擦
系数为 $\mu$，$B$ 端用一强度足够而不可伸长的绳悬挂，绳的另一端固定在墙壁 $C$ 点，木杆呈
水平状态，绳与杆的夹角为 $\theta$，求杆能保持平衡时 $\mu$ 与 $\theta$ 应满足的条件。杆保持平衡时，杆
上有一点 $P$ 存在，若 $A$ 与 $P$ 点间挂一质量足够大的重物，则可以破坏平衡，而在 $PB$ 间
任一点悬挂任意重物均不能破坏平衡。求 $PA$ 距离 $x$。
\sourcefigure[0.46\linewidth]{ex-1-20.png}
\end{origexamplebox}

\Needspace{6\baselineskip}\begin{xesolution}
杆、墙面和细绳的受力关系如解图所示；附加重物的位置暂记为 $y$：
\soldiagOneTwenty
先考虑没有附加重物的木杆。设绳张力为 $T$，墙面对杆的支持力和摩擦力分别为 $N$、
$f$。对 $A$ 点取矩：
\[
  T\sin\theta\,l=Mg\frac l2,
  \qquad
  T\sin\theta=\frac{Mg}{2}.
\]
再由水平、竖直方向平衡得到
\[
  N=T\cos\theta=\frac{Mg}{2}\cot\theta,
  \qquad
  f=Mg-T\sin\theta=\frac{Mg}{2}.
\]
静摩擦条件 $f\le\mu N$ 因而给出
\[
  \boxed{\mu\ge\tan\theta}.
\]

现在在距 $A$ 点为 $y$ 的位置悬挂质量为 $m$ 的重物。重复上述计算：
\[
  T\sin\theta=\frac{Mg}{2}+mg\frac yl,
\]
\[
  f=\frac{Mg}{2}+mg\left(1-\frac yl\right),
  \qquad
  N=\left(\frac{Mg}{2}+mg\frac yl\right)\cot\theta.
\]
若无论 $m$ 多大都不能破坏平衡，则除原先的 $\mu\ge\tan\theta$ 外，还必须使 $m$ 的
系数满足
\[
  1-\frac yl\le\mu\frac yl\cot\theta.
\]
临界点 $P$ 满足等号，于是
\[
  \boxed{x=AP=\frac{l\tan\theta}{\mu+\tan\theta}}.
\]
当悬挂点位于 $AP$ 内时，可用足够重的物体破坏平衡；位于 $PB$ 内时，任意重物都不能
使接触处超过静摩擦极限。
\end{xesolution}

\begin{origexamplebox}{}\label{ch05:ex:s03-1-21}
如图，一人对一均匀细杆的一端施力，力的方向总与杆垂直，要将杆从地板上无
滑动地慢慢抬到竖直位置，试求杆与地板间的静摩擦因数至少应为多大？
\sourcefigure[0.40\linewidth]{ex-1-21.png}
\end{origexamplebox}

\Needspace{6\baselineskip}\begin{xesolution}
设杆长为 $L$、质量为 $m$，杆与地板夹角为 $\theta$。外力 $F$ 始终垂直于杆。对杆的下端
取矩最方便，因为地板处的两个力都不产生力矩：
\soldiagOneTwentyOne
由解图有
\[
  FL=mg\frac L2\cos\theta,
  \qquad
  F=\frac{mg}{2}\cos\theta.
\]
由水平、竖直方向平衡可得地板处的摩擦力和支持力：
\[
  f=F\sin\theta=\frac{mg}{2}\sin\theta\cos\theta,
\]
\[
  N=mg-F\cos\theta
  =\frac{mg}{2}\left(1+\sin^2\theta\right).
\]
整个抬起过程中均不滑动，要求
\[
  \mu\ge\max_{0\le\theta\le\pi/2}\frac fN
  =\max_{0\le\theta\le\pi/2}
  \frac{\sin\theta\cos\theta}{1+\sin^2\theta}.
\]
在 $0\le\theta<\pi/2$ 内令 $t=\tan\theta\ge0$，上式化为 $t/(1+2t^2)$。
由非负平方 $(\sqrt2t-1)^2\ge0$ 得
\[
 1+2t^2\ge2\sqrt2t.
\]
当 $t>0$ 时，两边除以正数 $1+2t^2$，得
$t/(1+2t^2)\le1/(2\sqrt2)$，等号恰在 $t=1/\sqrt2$ 时成立。
$t=0$ 时此比值为零；当 $\theta\to\pi/2$、$t\to\infty$ 时也趋于零，
所以两端都不提供更大的摩擦要求。因此
\[
  \boxed{\mu_{\min}=\frac{1}{2\sqrt2}}.
\]
\end{xesolution}

\begin{origexamplebox}{}\label{ch05:ex:s03-1-22}
质量均为 $m$ 的两环 $A$、$B$ 用长为 $a$ 的细线相连在水平杆上，在细线的中点拴有
一质量为 $M$ 的物块 $C$，$A$、$B$ 环与杆间的静摩擦系数为 $\mu$，求平衡情况下的两环的最大距离
$x$。
\sourcefigure[0.50\linewidth]{ex-1-22.png}
\end{origexamplebox}

\Needspace{6\baselineskip}\begin{xesolution}
系统关于中点对称。每一侧细线长度为 $a/2$。设细线与水平方向夹角为 $\beta$，则
\[
  \cos\beta=\frac xa,
  \qquad
  \sin\beta=\frac{\sqrt{a^2-x^2}}{a}.
\]
中间结点和任一圆环的受力应分开研究：
\soldiagOneTwentyTwo
在中间结点处列竖直方向平衡方程：
\[
  2T\sin\beta=Mg,
  \qquad
  T=\frac{Mg}{2\sin\beta}.
\]
对任一圆环，细线的竖直分量为 $T\sin\beta=Mg/2$，所以杆对环的支持力为
\[
  N=\left(m+\frac M2\right)g.
\]
环将要滑动时
\[
  T\cos\beta=\mu N.
\]
代入 $T$ 和几何关系得
\[
  Mx=\mu(2m+M)\sqrt{a^2-x^2}.
\]
因此两环能够保持平衡的最大距离为
\[
  \boxed{x_{\max}=
  \frac{a\mu(2m+M)}{\sqrt{M^2+\mu^2(2m+M)^2}}}.
\]
\end{xesolution}

\section{有质量绳的平衡}

之前问题中的绳一般取无质量的极限。对于没有沿程外力、只在两端受力的自由绳段，平衡且
张力不为零时绳段为直线，两端张力大小相等；多数情况下还忽略伸长。有滑轮或曲面支承时，
轻绳可以弯曲，在光滑支承下张力大小仍可相等，但方向沿绳改变。
实际上很多问题中绳的质量和形变不能忽略，例如桥梁上的斜
拉索，它一方面负责承担桥梁其它部分的重量，另一方面必须采用足够强度的材料以保证其
可靠性，这时其自重就不可忽略不计。

当考虑可变形物体的质量时，其内部各部分之间的作用力将不再相同，除了承担外部的
受力以外，其自重也会改变不同位置处的受力。最简单的例子如图~\ref{ch05:fig:s03-1-3}(a) 所示，一根质量为
$m$，长度为 $L$ 的均匀绳静止悬挂于 $O$ 点，建立一个原点在 $O$ 点、竖直向下的坐标系 $Ox$，
这时绳中张力 $T$ 就不再是常数，而是坐标 $x$ 的函数。在某一 $x$ 点处的张力用来平衡处于其
下部绳的重力，简单的分析可知绳中张力满足
\[
  T(x)=\frac{mg}{L}(L-x).
\]

在绳的底端，也就是 $x=L$ 处很明显张力为零。在更复杂的情况下不但绳中张力不同，绳
还会发生弯曲。这时任何一点处的张力的方向是绳在该处的切线方向，此时如果绳的质量不
可忽略，需要取绳中极小一段作为研究对象，两端的张力在水平方向上平衡，而在竖直方向
上与该段绳的重力平衡，如图~\ref{ch05:fig:s03-1-3}(b) 中那样。

\sourcefigure[0.76\linewidth]{fig-1-3.png}
\origcaption{(a) 有质量的绳中各部分张力不同，是位置的函数；(b) 有质量且弯曲绳的受力分析}\label{ch05:fig:s03-1-3}

\begin{origexamplebox}{}\label{ch05:ex:s03-1-23}
如图所示，一个顶角为 $\alpha$ 的光滑圆锥体，一个质量为 $m$、原长为 $L_0$、劲度系数为
$k$ 的弹性绳处于平衡状态，所有部分均处于同一水平面上。求平衡绳圈到锥顶的母线距离
$s$，并求绳圈相对锥顶的竖直高度 $z$。
\sourcefigure[0.27\linewidth]{ex-1-23.png}
\end{origexamplebox}

\Needspace{6\baselineskip}\begin{xesolution}
令圆锥的半顶角为 $\delta=\alpha/2$，平衡绳圈的水平半径为 $\rho$，绳中张力为 $T$。
取绳圈上圆心角为 $\mathrm d\varphi$ 的一小段，微元受力如解图所示：
\soldiagOneTwentyThree
其两端张力的合力沿水平方向指向圆锥轴，
大小为 $T\,\mathrm d\varphi$。圆锥支持力的竖直分量与水平分量之比为 $\tan\delta$，所以
这段绳所受支持力的竖直分量为
\[
  T\tan\delta\,\mathrm d\varphi.
\]
这应当等于该段绳的重力 $(mg/2\pi)\,\mathrm d\varphi$，故
\[
  T=\frac{mg}{2\pi}\cot\delta.
\]
弹性绳伸长后的总长为 $2\pi\rho$，由胡克定律
\[
  T=k(2\pi\rho-L_0).
\]
因此
\[
  \rho=\frac{L_0}{2\pi}+\frac{mg}{4\pi^2k}\cot\delta.
\]
由 $\rho=s\sin\delta$，母线距离为
\[
  \boxed{s=\frac{1}{\sin(\alpha/2)}
  \left[\frac{L_0}{2\pi}+\frac{mg}{4\pi^2k}\cot\frac\alpha2\right]}.
\]
相应的竖直高度为
\[
  \boxed{z=\rho\cot\frac\alpha2}.
\]
\end{xesolution}

\begin{origexamplebox}{}\label{ch05:ex:s03-1-24}
有一个固定在水平面上半径为 $R$ 的半圆柱形光滑物体和一根质量为 $m$、长为 $\pi R/2$
的均匀绳。它们处于如图所示位置的平衡状态，求维持平衡作用于绳顶端水平向左力 $F$ 的
大小。
\sourcefigure[0.46\linewidth]{ex-1-24.png}
\end{origexamplebox}

\Needspace{6\baselineskip}\begin{xesolution}
绳的线密度为
\[
  \lambda=\frac{m}{\pi R/2}=\frac{2m}{\pi R}.
\]
以圆柱右端为 $\varphi=0$、最高点为 $\varphi=\pi/2$。沿绳切线方向研究圆心角为
$\mathrm d\varphi$ 的小段，切向受力如解图所示：
\soldiagOneTwentyFour
因此
\[
  \mathrm dT=\lambda gR\cos\varphi\,\mathrm d\varphi.
\]
右端为自由端，故 $T(0)=0$。积分到最高点：
\[
  F=T\!\left(\frac\pi2\right)
  =\lambda gR\int_0^{\pi/2}\cos\varphi\,\mathrm d\varphi
  =\lambda gR.
\]
所以
\[
  \boxed{F=\frac{2mg}{\pi}}.
\]
\end{xesolution}

\begin{origexamplebox}{}\label{ch05:ex:s03-1-25}
质量为 $m$、长为 $l$ 的均匀光滑细绳，穿过半径为 $R$ 的光滑滑轮并搭在轮上。设
$l\ge\pi R$，绳关于滑轮的竖直中线对称，两端自由段均竖直且等长。求绳上的最大张力。
\sourcefigure[0.34\linewidth]{ex-1-25.png}
\end{origexamplebox}

\Needspace{6\baselineskip}\begin{xesolution}
设 $l\ge\pi R$，绳关于滑轮的竖直中线对称。搭在滑轮上的部分长为 $\pi R$，两段竖直
部分各长
\[
  h=\frac{l-\pi R}{2}.
\]
张力从自由端向上逐渐增加，其关键位置如解图所示：
\soldiagOneTwentyFive
绳的线密度为 $\lambda=m/l$。从任一自由端向上，竖直段顶端的张力为 $\lambda gh$。
再从侧面切点沿圆弧走到滑轮最高点，张力增加
\[
  \lambda gR\int_0^{\pi/2}\cos\varphi\,\mathrm d\varphi=\lambda gR.
\]
故最大张力出现在滑轮最高点：
\[
  T_{\max}=\lambda g(h+R)
  =\frac{mg}{l}\left(\frac{l-\pi R}{2}+R\right).
\]
即
\[
  \boxed{T_{\max}=\frac{mg}{2l}\bigl[l-(\pi-2)R\bigr]}.
\]
\end{xesolution}

\begin{origexamplebox}{}\label{ch05:ex:s03-1-26}
半径为 $R$ 的刚性球固定在水平桌面上，有一个质量为 $M$ 的圆环状均匀弹性绳圈，
原长 $2\pi a$，$a=R/2$。绳圈的弹性系数为 $k$（绳圈伸长 $s$ 时，绳中弹性张力为 $ks$）。将绳圈从
球的正上方轻放到球上，并用手扶着绳圈使之保持水平并最后停留在某个静力平衡位置。
扶持仅维持水平姿态，不承担绳圈的重量，也不施加径向拉力。
设此时绳圈的长度为 $2\pi b$，$b=\sqrt{2}a$，考虑重力，忽略摩擦，求绳圈的弹性系数 $k$？（用 $M$、$R$、
$g$ 表示，$g$ 为重力加速度）
\sourcefigure[0.39\linewidth]{ex-1-26.png}
\end{origexamplebox}

\Needspace{6\baselineskip}\begin{xesolution}
设绳圈所在水平面的高度相对球心为 $z$。由球面几何关系
\[
  b^2+z^2=R^2.
\]
已知 $a=R/2$、$b=\sqrt2a=R/\sqrt2$，所以 $z=R/\sqrt2=b$。

取绳圈上圆心角为 $\mathrm d\varphi$ 的一小段，设球面对它的支持力为 $\mathrm dN$。
微元所用的几何量和三个力如解图所示：
\soldiagOneTwentySix
竖直方向平衡给出
\[
  \mathrm dN\frac zR=\frac{Mg}{2\pi}\,\mathrm d\varphi.
\]
支持力的水平分量与两端张力的合力平衡：
\[
  T\,\mathrm d\varphi=\mathrm dN\frac bR.
\]
两式相除得
\[
  T=\frac{Mgb}{2\pi z}=\frac{Mg}{2\pi}.
\]
绳圈的总伸长量为
\[
  s=2\pi(b-a)=\pi R(\sqrt2-1).
\]
由 $T=ks$，
\[
  \boxed{k=\frac{Mg}{2\pi^2R(\sqrt2-1)}
  =\frac{(\sqrt2+1)Mg}{2\pi^2R}}.
\]
\end{xesolution}

\begin{origexamplebox}{【思考题】}\label{ch05:ex:s03-1-27}
一根长度为 $L$、质量为 $m$ 的均匀绳两端固定在处于同一水平面上相距
$d<L$ 的 $A$、$B$ 两点，你能够用何种方法决定平衡以后绳的形状。
\sourcefigure[0.54\linewidth]{ex-1-27.png}
\end{origexamplebox}

\Needspace{6\baselineskip}\begin{xesolution}
平衡后的绳是悬链线。取最低点为坐标原点，$x$ 轴水平、$y$ 轴竖直向上。设绳的线密度
$\lambda=m/L$，最低点处的水平张力为 $H$，并定义
\[
  c=\frac{H}{\lambda g}.
\]
对最低点到横坐标 $x$ 的一段绳列平衡方程。记这段绳的弧长为 $s$，研究对象如解图所示：
\soldiagOneTwentySeven
设 $P$ 点绳的切线与水平方向的夹角为 $\theta$。解图中三个力分别沿水平和竖直方向平衡：
\[
  T\cos\theta=H,
  \qquad
  T\sin\theta=\lambda gs.
\]
两式相除，并利用曲线切线斜率 $y'=\tan\theta$，得
\[
  y'=\tan\theta=\frac{\lambda gs}{H}=\frac{s}{c}.
\]
另一方面，弧长元满足
\[
  \frac{\mathrm ds}{\mathrm dx}=\sqrt{1+y'^2}.
\]
对 $y'=s/c$ 关于 $x$ 求导，再代入上式，才得到由解图决定的微分方程
\[
  y''=\frac{1}{c}\sqrt{1+y'^2},
  \qquad y(0)=0,\quad y'(0)=0.
\]
先把这个二阶方程降为一阶。令 $p=y'$，并只看 $x\ge0$ 的右半段，则
\[
 \frac{\mathrm dp}{\sqrt{1+p^2}}=\frac{\mathrm dx}{c}.
\]
由于
\[
 \frac{\mathrm d}{\mathrm dp}\ln\!\left(p+\sqrt{1+p^2}\right)
 =\frac{1+p/\sqrt{1+p^2}}{p+\sqrt{1+p^2}}
 =\frac1{\sqrt{1+p^2}},
\]
从 $x=0,p=0$ 积分到 $x,p$，得到
\[
 \ln\!\left(p+\sqrt{1+p^2}\right)=\frac xc.
\]
令 $E=\mathrm e^{x/c}$，则 $p+\sqrt{1+p^2}=E$；两因子
$(\sqrt{1+p^2}+p)(\sqrt{1+p^2}-p)=1$，故
\[
 p=\frac{E-E^{-1}}2=\frac{\mathrm e^{x/c}-\mathrm e^{-x/c}}2.
\]
为简写这类指数式，定义双曲正弦和双曲余弦：
\[
 \sinh u=\frac{\mathrm e^u-\mathrm e^{-u}}2,\qquad
 \cosh u=\frac{\mathrm e^u+\mathrm e^{-u}}2.
\]
由指数函数求导可得 $(\cosh u)'=\sinh u$、$(\sinh u)'=\cosh u$，直接展开还得到
$\cosh^2u-\sinh^2u=1$。现在 $y'=\sinh(x/c)$，再次积分并使用 $y(0)=0$，得
\[
 y(x)=\int_0^x\sinh\frac{\xi}{c}\,\mathrm d\xi
 =c\left[\cosh\frac{\xi}{c}\right]_0^x.
\]
因此平衡曲线为悬链线
\[
  \boxed{y(x)=c\cosh\frac{x}{c}-c}.
\]
它是偶函数，左半段由对称性给出，两固定点位于 $x=\pm d/2$。
利用 $\sqrt{1+\sinh^2u}=\cosh u$，弧长条件为
\[
  L=\int_{-d/2}^{d/2}\sqrt{1+y'^2}\,\mathrm dx
  =2c\sinh\frac{d}{2c}.
\]
因此先由
\[
  \boxed{2c\sinh\frac{d}{2c}=L}
\]
求出 $c>0$，再代入悬链线方程。为什么一定有且只有一个解？令 $q=d/(2c)>0$，则
\[
 f(q)=\frac{\sinh q}{q}=\frac Ld>1,\qquad
 f'(q)=\frac{q\cosh q-\sinh q}{q^2}.
\]
分子在 $q=0$ 时为零，其导数为 $q\sinh q>0$，所以 $f'(q)>0$。
同时 $f(q)\to1$（$q\to0^+$），且 $f(q)\to\infty$（$q\to\infty$），因此每个 $L/d>1$
都对应唯一的 $q$，也就对应唯一的 $c=d/(2q)$。可先用二分法求 $q$，再计算曲线。

由 $T\cos\theta=H$ 及 $\tan\theta=\sinh(x/c)$，有
$1/\cos\theta=\sqrt{1+\tan^2\theta}=\cosh(x/c)$，所以绳中张力为
\[
  T(x)=H\cosh\frac{x}{c}=\lambda gc\cosh\frac{x}{c}.
\]
$c$ 的量纲是长度，指数中的 $x/c$ 和弧长方程中的 $d/(2c)$ 均无量纲。当 $d\to L^-$ 时，
$q\to0$、$c\to\infty$，曲线近似为 $y\simeq x^2/(2c)$，下垂趋于零，
而水平张力 $H=\lambda gc$ 发散。因而有自重的柔绳不能用有限水平张力拉成完全水平的直线。
\end{xesolution}

\section{平衡的稳定性}

受力平衡方程只能确定可能的平衡位置；要判断物体受到微小扰动后能否回到原位置，还要
研究平衡的稳定性。如下图所示，同样的两根竖直杆，第一根由顶部悬挂，第二根以下端支承
在水平桌面上。理想情况下，两者都能满足静力平衡条件，但前者是稳定平衡，后者是
不稳定平衡。

\sourcefigure[0.69\linewidth]{fig-1-4.png}
\origcaption{稳定和不稳定平衡}\label{ch05:fig:s03-1-4}

如果系统稍微偏离平衡位置后出现使偏离量减小的回复力或回复力矩，该平衡就是稳定的。
例如悬挂杆偏离竖直方向后，重力矩使偏角减小。反之，如果扰动后的力或力矩使偏离量
继续增大，该平衡就是不稳定的；竖直支在桌面上的杆正是这种情形。对保守系统，也可以
用势能判断：稳定平衡对应势能的严格局部极小值，不稳定平衡通常对应局部极大值。

不稳定平衡仍然是严格意义上的平衡，并非“不可能出现”；只是它不能抵抗微小扰动。
这类平衡在倒立摆、结构屈曲和临界状态分析中都有重要意义。工程应用中必须把“满足平衡
方程”和“具有稳定性”分开检验；稳定性既可以定性判断，也可以通过力矩展开或势能展开
作定量分析。

\begin{origexamplebox}{}\label{ch05:ex:s03-1-28}
一根质量为 $m$ 的均匀杆，长为 $L$，下端可绕固定水平轴转动。有两根水平弹簧，
劲度系数相同，拴在杆的上端，使其处于竖直位置。问：弹簧的劲度系数 $k$ 为何值，才能使
杆处于稳定平衡？两弹簧按图对称布置，外端固定在竖直杆顶点两侧等高、等距的位置；
每根弹簧的原长记为 $\ell_0>0$，初始长度记为 $d_0>0$。按理想轻弹簧模型讨论小幅平面偏转。
\sourcefigure[0.49\linewidth]{ex-1-28.png}
\end{origexamplebox}

\Needspace{6\baselineskip}\begin{xesolution}
设杆从竖直方向偏转一个小角 $\varphi$。杆上端的水平位移为
\[
  x=L\sin\varphi=L\varphi+O(\varphi^3).
\]
\Needspace{66mm}
偏转后的杆、重力与两弹簧形变如解图所示：
\soldiagOneTwentyEight
在竖直位置两弹簧关于杆对称。由解图可见，小角度下左、右弹簧的长度改变量一阶项分别为
$+x$ 和 $-x$。因而两弹簧的总弹性势能相对竖直位置的变化为
\[
  \Delta U_{\rm s}=kL^2\varphi^2+O(\varphi^4).
\]
杆的重力势能变化为
\[
  \Delta U_{\rm g}=mg\frac L2(\cos\varphi-1)
  =-\frac{mgL}{4}\varphi^2+O(\varphi^4).
\]
因此总势能在 $\varphi=0$ 附近满足
\[
  U=U(0)+\left(kL^2-\frac{mgL}{4}\right)\varphi^2+O(\varphi^4).
\]
二次项系数为正时，势能在竖直位置为严格局部极小值，有线性回复力矩。因此
\[
  \boxed{k>\frac{mg}{4L}}.
\]
若 $k<mg/(4L)$，二次项为负，竖直位置不稳定。等号处二次项消失，必须继续检查，不能只据二次项下结论。
杆顶下降量为 $v=L(1-\cos\varphi)$，根据两固定点到杆顶的距离，弹簧长度为
\[
 \ell_\pm=\sqrt{(d_0\pm x)^2+v^2}.
\]
将两根弹簧的势能分别写成 $k(\ell_\pm-\ell_0)^2/2$，减去初态势能，有
\[
 \Delta U_{\rm s}
 =k\left[x^2+v^2-\ell_0(\ell_++\ell_--2d_0)\right].
\]
这里用了 $\ell_+^2+\ell_-^2=2d_0^2+2x^2+2v^2$。再利用
$x^2+v^2=2L^2(1-\cos\varphi)$，加上重力势能，得精确关系
\[
 \Delta U=\left(2kL^2-\frac{mgL}{2}\right)(1-\cos\varphi)
 -k\ell_0(\ell_++\ell_--2d_0).
\]
当 $k=mg/(4L)$ 时，第一项恰好消失。只要 $\varphi\ne0$，杆顶就低于两固定点的连线；
三角形两边长度之和大于第三边，故 $\ell_++\ell_->2d_0$。由于 $\ell_0>0$，此时
$\Delta U<0$，所以临界点仍不稳定，稳定范围确为上面的严格不等式。
若另取原长恰为零的理想极限，则临界时 $\Delta U\equiv0$，为随遇平衡，而不是有回复力的稳定平衡。
\end{xesolution}

\begin{origexamplebox}{}\label{ch05:ex:s03-1-29}
两根长度均为 $L$、质量忽略不计的刚性轻杆固定在竖直平面内，顶端由铰链连接
并固定有一个质量为 $M$ 的质点。两杆下端置于同一光滑水平支承面上，可以水平滑动；
下端之间连接一根劲度系数为 $k$、原长为 0 的轻弹簧。求平衡时杆与垂线的夹角 $\theta$。
\sourcefigure[0.39\linewidth]{ex-1-29.png}
\end{origexamplebox}

\Needspace{6\baselineskip}\begin{xesolution}
题图中的光滑水平支承面只提供竖直支持力，$A$、$B$ 可以沿水平方向滑动。
解图分别画顶点 $M$ 和下端 $A$ 的受力，其中 $N_A$ 是支承面对 $A$ 的竖直支持力；
$B$ 端按对称性处理，$N_B=N_A$：
\soldiagOneTwentyNine

弹簧原长为零，平衡时长度为 $2L\sin\theta$，所以弹簧张力
\[
  T=2kL\sin\theta.
\]
设每根杆中的压力为 $P$。顶端质点的竖直平衡给出
\[
  2P\cos\theta=Mg.
\]
在下端 $A$ 或 $B$ 处列水平方向平衡方程：
\[
  P\sin\theta=T=2kL\sin\theta.
\]
因此完整的平衡条件可以写成
\[
  \sin\theta\bigl(4kL\cos\theta-Mg\bigr)=0.
\]
它包含两个分支。第一支是 $\theta=0$，此时 $A$、$B$ 重合、两杆重叠，是图示构型的退化
极限。第二支为非退化构型 $\sin\theta\ne0$，此时 $P=2kL$，并有
\[
  \boxed{\cos\theta=\frac{Mg}{4kL}},
  \qquad
  \boxed{\theta=\arccos\frac{Mg}{4kL}}.
\]
非退化分支存在的条件为
\[
  \boxed{k>\frac{Mg}{4L}}.
\]

由于本题位于“平衡的稳定性”一节，还应继续检查稳定性。在上述补充模型下，系统势能为
\[
  U(\theta)=MgL\cos\theta+2kL^2\sin^2\theta.
\]
在非退化平衡角处
\[
  U''(\theta)=-4kL^2\sin^2\theta<0,
\]
所以方框中的非零夹角分支是\textbf{不稳定平衡}。退化的 $\theta=0$ 分支在
$k>Mg/(4L)$ 时稳定；等号处二次项消失而四次项为负，仍不稳定。
\end{xesolution}

\begin{origexamplebox}{}\label{ch05:ex:s03-1-30}
一根长为 $L$、质量为 $m$ 的均匀杆一端靠于竖直墙面，另一端靠在曲面上，墙面与曲面均光滑。
取水平向右为 $x$ 轴、竖直向上为 $y$ 轴，曲线过原点。要杆在图示构型的每个可实现位置都平衡，
求曲面的截面满足的条件，并说明适用的位置范围。
\sourcefigure[0.42\linewidth]{ex-1-30.png}
\end{origexamplebox}

\Needspace{6\baselineskip}\begin{xesolution}
设杆与曲面的接触点为 $P(x,y)$，杆在墙上的端点为 $A(0,Y)$。由杆长为 $L$，
\[
  Y-y=\sqrt{L^2-x^2}.
\]
三个力的作用线及其交点 $C$ 如解图所示：
\soldiagOneThirty
假设墙面和曲面均光滑。杆只受墙的水平支持力、曲面的法向支持力和重力三个力。三力
平衡时，其作用线必须交于一点。墙面支持力的作用线是过 $A$ 的水平线，重力作用线是过
杆中点的竖直线，两者交点为
\[
  C\left(\frac x2,Y\right).
\]
所以曲面在 $P$ 点的法线必须通过 $C$。直线 $PC$ 的斜率为
\[
  \frac{Y-y}{x/2-x}=-\frac{2\sqrt{L^2-x^2}}{x}.
\]
另一方面，曲线 $y(x)$ 的法线斜率为 $-1/y'$，于是
\[
  y'=\frac{x}{2\sqrt{L^2-x^2}}.
\]
利用 $\mathrm d\sqrt{L^2-x^2}/\mathrm dx=-x/\sqrt{L^2-x^2}$，有
\[
 y(x)=-\frac12\sqrt{L^2-x^2}+C_0.
\]
曲线过原点，代入 $x=0,y=0$ 得 $C_0=L/2$，于是
\[
  \boxed{y(x)=\frac{L-\sqrt{L^2-x^2}}{2}}.
\]
等价地，曲线满足
\[
  \boxed{x^2+4\left(y-\frac L2\right)^2=L^2},
\]
即它是半长轴为 $L$、半短轴为 $L/2$ 的椭圆下半支中位于墙右侧的一段。
杆的质心高度为
\[
 y_G=\frac{Y+y}{2}=y+\frac12\sqrt{L^2-x^2}=\frac L2,
\]
不随位置改变，因此在可保持接触的位置上为随遇平衡。

还需检验接触力和端点。曲面法线的水平、竖直分量之比为 $-y':1$；由竖直平衡和水平平衡，
\[
 N_{P,y}=mg,\qquad
 N_A=mg\,y'=\frac{mgx}{2\sqrt{L^2-x^2}},\qquad
 N_P=mg\sqrt{1+y'^2}.
\]
在 $0<x<L$ 内两处支持力均不要求拉力。$x=0$ 是竖直杆极限，此时 $N_A=0$、$N_P=mg$，
可单独由受力平衡验证，不能在前面含 $1/x$ 的斜率式中直接代入零。
当 $x\to L^-$ 时，$N_A$ 和 $N_P$ 都发散；在 $x=L$ 处杆水平，两处接触法线均水平，
不能平衡重力。因此范围是
\[
 \boxed{0\le x<L},
\]
不包括最右端点。实际材料强度有限时，可用范围还会比这一理想范围更小。
\end{xesolution}

\begin{origexamplebox}{【思考题】}\label{ch05:ex:s03-1-31}
\begin{minipage}[t]{0.62\linewidth}
如图，一个固定圆形“碗”的半径为 $R$，一个小球半径为 $r$；接触面足够粗糙，小球
作纯滚动而不滑动，初始时小球重心位于球心正上方。小球不是匀质
的，设 $R>r$，且重心在球内、$0<h=OG<r$。请问 $h$ 为多大时形成稳定平衡。如果半径 $R$ 不断增大，会怎样？
增大到无穷大会怎样？再进一步增大会怎样？
\end{minipage}\hfill
\begin{minipage}[t]{0.34\linewidth}
\vspace{0pt}\centering
\resizebox{\linewidth}{!}{\faithfulExOneThirtyOne}
\end{minipage}
\end{origexamplebox}

\Needspace{6\baselineskip}\begin{xesolution}
设 $OG=h$，并取 $h>0$ 表示重心位于球心上方。假设小球在碗内滚动而不滑动，记
$a=R-r$。小球偏离最低位置后的两个转角和重心位置如解图所示：
\soldiagOneThirtyOne
球心从最低点沿碗转过小角 $\varphi$ 时，平动弧长为 $a\varphi$。无滑动条件给出小球
相对地面的转角
\[
  \psi=-\frac{a}{r}\varphi.
\]

略去与 $\varphi$ 无关的常数后，重心的高度可以写成
\[
  y_G=-a\cos\varphi+h\cos\frac{a\varphi}{r}.
\]
在 $\varphi=0$ 附近展开：
\[
  y_G=y_G(0)+\frac12\left(a-\frac{ha^2}{r^2}\right)\varphi^2
  +O(\varphi^4).
\]
二次项系数为正时，最低位置是势能的严格局部极小点，具有线性回复力。因此
\[
  \boxed{h<\frac{r^2}{R-r}}.
\]
当 $h>r^2/(R-r)$ 时二次项为负，最低位置不稳定。等号不能只凭二次项判断。
展开两个余弦至四阶：
\[
 y_G-y_G(0)=\frac12\left(a-\frac{ha^2}{r^2}\right)\varphi^2
 +\frac1{24}\left(-a+\frac{ha^4}{r^4}\right)\varphi^4+O(\varphi^6).
\]
在 $h=r^2/a$ 处，四次项系数化为
\[
 \frac1{24}\left(-a+\frac{a^3}{r^2}\right)
 =\frac{a(a^2-r^2)}{24r^2}.
\]
$R>2r$ 时 $a>r$，该系数为正，临界点仍是严格势能极小点，只是回复力矩从三次项开始，
称为四阶弱稳定。$R=2r$ 时 $a=r,h=r$，两项余弦完全抵消，为随遇平衡；
$r<R<2r$ 时形式上的临界四次项为负，但所需 $h=r^2/a>r$ 已不可能位于球内。

结合题设 $0<h<r$，物理上可实现的稳定范围为
\[
 \boxed{\begin{cases}
  0<h<r, & r<R\le2r,\\
  0<h\le r^2/(R-r), & R>2r.
 \end{cases}}
\]
第二行的等号包含弱稳定；若专指有非零线性回复力的稳定，则把等号去掉。
例如 $R=3r,h=r/2$ 恰是可实现的弱稳定临界状态。

当 $R$ 增大时，$h_{\rm c}=r^2/(R-r)$ 单调减小；$R\to\infty$ 时碗面趋于水平面，
$h_{\rm c}\to0$，任何固定的 $h>0$ 最终都会变得不稳定。平面上球心高度不变，
重心向上时滚动会使其降低，所以这个极限也可直接理解。

“再进一步增大”不是把无穷大当作普通数继续相加，而是让表面的曲率从凹面经零曲率变成凸面。
下面另取凸面半径 $S>0$，仍令 $h>0$ 表示初态重心在球心上方，$h<0$ 表示在下方。
\Needspace{65mm}
凸面圆心、接触点和小球球心的关系如解图所示：
\soldiagConvexRolling
记球心轨道半径为 $a_+=S+r$，取球心向右偏转的角度为 $\varphi$。
球心速度沿凸面切线向右下方；接触点无滑动要求小球顺时针转动，故
\[
 \psi=-\frac{a_+}{r}\varphi.
\]
以凸面圆心为高度原点，球心高度为 $a_+\cos\varphi$，重心相对球心的竖直分量为
$h\cos(a_+\varphi/r)$，因而
\[
 y_G=a_+\cos\varphi+h\cos\frac{a_+\varphi}{r},
\]
\[
 y_G-y_G(0)=-\frac12\left(a_++\frac{ha_+^2}{r^2}\right)\varphi^2+O(\varphi^4).
\]
二次项为正给出
\[
 \boxed{h<-\frac{r^2}{S+r}}.
\]
即重心必须位于球心下方，且下移足够多。在等号处，四次项系数为
$a_+(1-a_+^2/r^2)/24<0$，所以临界点不稳定，等号不取。
该条件还需与重心在球内的限制 $-r<h<r$ 同时满足。
\end{xesolution}
}

\endgroup

% MechanicsNumberedSolutionFigures20261023-v1
% MechanicsQualityRevision20261023-v1
% QualityRevision20261023
\begingroup
\tcbset{breakable/.default=false}
\chapter{牛顿定律、非惯性系及其应用}

{
\graphicspath{{../../../sample/xes-original/fourth-section/}}
\section{牛顿运动定律}
% TikZ diagrams for the added solutions.  Each diagram contains only the
% geometric or force quantities used immediately in the accompanying text.

\tikzset{
  solbody/.style={draw=Ink,line width=0.8pt},
  solforce/.style={-{Latex[length=2.0mm]},draw=DeepBlue,line width=0.9pt},
  solgeom/.style={draw=TextGray,densely dashed,line width=0.55pt},
  solpoint/.style={circle,fill=Ink,inner sep=1.2pt},
  sollabel/.style={font=\footnotesize,fill=none,inner sep=0.7pt,outer sep=1.2pt,text=Ink},
  solpanel/.style={font=\scriptsize\bfseries,text=TextGray}
}

\newcommand{\soldiagcaption}[1]{\figcap{#1}}

\newcommand{\diagcaption}[2]{\figcap{#1：#2}}

\newcommand{\soldiagTwoOne}{%
\begin{center}\begin{minipage}{\linewidth}\centering
\begin{tikzpicture}[scale=0.90]
  \draw[axis] (-3.0,0)--(3.2,0) node[right,sollabel] {$x$};
  \fill[Ink] (-2.1,0) circle (1.3pt) node[above left,sollabel] {$x<0$};
  \fill[Ink] (2.1,0) circle (1.3pt) node[above right,sollabel] {$x>0$};
  \draw[solforce] (-2.1,0)--++(0.95,0) node[above,sollabel] {$F>0$};
  \draw[solforce] (2.1,0)--++(-0.95,0) node[above,sollabel] {$F<0$};
  \fill[Ink] (0,0) circle (1.3pt) node[below,sollabel] {$x=0$};
\end{tikzpicture}
\soldiagcaption{$F=-m\omega^2x$ 始终指向 $x=0$，是与位移成正比的回复力。}
\end{minipage}\end{center}}

\newcommand{\soldiagTwoTwo}{%
\begin{center}\begin{minipage}{\linewidth}\centering\begin{tikzpicture}[scale=.95]
 \coordinate (O) at (0,0); \coordinate (P) at (35:1.8);
 \draw[axis] (-2.4,0)--(2.7,0) node[right,sollabel] {$x$};
 \draw[axis] (0,-2.05)--(0,2.55) node[above,sollabel] {$y$};
 \draw[solbody] (O) circle (1.8);
 \fill (O) circle (1.3pt) node[below left,sollabel] {$O$};
 \fill (P) circle (1.5pt) node[right=5pt,sollabel] {$P$};
 \draw[solgeom] (O)--(P) node[pos=.45,below right=4pt,sollabel] {$A$};
 \draw[solforce] (P)--($(P)!.65!(O)$) node[pos=.65,above left=4pt,sollabel] {$\vec F$};
 \draw[solgeom,-{Latex[length=2mm]}] (P)--++(125:1.05)
   node[above left=3pt,sollabel] {$\vec v$};
\end{tikzpicture}
\soldiagcaption{$\omega>0$ 时逆时针运动，速度是位置函数的导数，合力沿半径向内。}
\end{minipage}\end{center}}

\newcommand{\soldiagTwoThree}{%
\begin{center}\begin{minipage}{\linewidth}\centering
\begin{tikzpicture}[scale=.90,x={(1cm,0cm)},y={(-.48cm,-.28cm)},z={(0cm,1cm)}]
 \draw[axis] (0,0,0)--(2.35,0,0) node[right,sollabel] {$x$};
 \draw[axis] (0,0,0)--(0,2.1,0) node[below left,sollabel] {$y$};
 \draw[axis] (0,0,0)--(0,0,3.95) node[above,sollabel] {$z$};
 \fill (0,0,0) circle(1.3pt);
 \node[sollabel] at (.45,0,.55) {$O$};
 \draw[solgeom,domain=0:360,samples=100,variable=\t] plot ({1.4*cos(\t)},{1.4*sin(\t)},0);
 \draw[curve,domain=0:630,samples=230,variable=\t] plot ({1.4*cos(\t)},{1.4*sin(\t)},{.0055*\t});
 \coordinate (P) at ({1.4*cos(480)},{1.4*sin(480)},{.0055*480});
 \coordinate (H) at (0,0,{.0055*480});
 \coordinate (Pz) at ({1.4*cos(480)},{1.4*sin(480)},0);
 \coordinate (V) at ({1.4*cos(480)-.55*1.4*sin(480)},{1.4*sin(480)+.55*1.4*cos(480)},{.0055*480+.55*.0055*180/pi});
 \draw[solgeom] (P)--(Pz);
 \fill (Pz) circle (1pt) node[left=7pt,sollabel] {$P_0$};
 \fill (P) circle (1.4pt) node[below left=6pt,sollabel] {$P$};
 \fill (H) circle (1pt) node[right=5pt,sollabel] {$H$};
 \draw[solforce] (P)--(H) node[pos=.55,below=9pt,sollabel] {$\vec F$};
 \draw[solgeom,-{Latex[length=2mm]}] (P)--(V) node[above left=3pt,sollabel] {$\vec v$};
\end{tikzpicture}
\soldiagcaption{图取 $q(t)=ut,\ u>0$；底部虚圆为水平投影，$PH$ 是同一高度上的径向线段。}
\end{minipage}\end{center}}

\newcommand{\soldiagTwoFour}{%
\begin{center}\begin{minipage}{\linewidth}\centering\begin{tikzpicture}[scale=.86]
 \begin{scope}[xshift=-3.5cm]
  \draw[axis] (0,0)--(3.5,0) node[right,sollabel] {$x$};
  \draw[axis] (0,0)--(0,2.25) node[above,sollabel] {$v$};
  \draw[curve] (0,1.75)--(3,0);
  \node[left=3pt,sollabel] at (0,1.75) {$v_0$};
  \node[below=4pt,sollabel] at (3,0) {$v_0/\alpha$};
  \node[sollabel] at (1.9,1.6) {$v=v_0-\alpha x$};
 \end{scope}
 \begin{scope}[xshift=2.3cm]
  \draw[axis] (-1.3,0)--(2,0) node[right,sollabel] {$x$};
  \fill (0,0) circle (2pt);
  \draw[solforce] (0,0)--(-1.1,0) node[pos=.6,above=4pt,sollabel] {$\vec F$};
  \draw[solgeom,-{Latex[length=2mm]}] (0,.65)--(1.2,.65) node[right,sollabel] {$\vec v$};
  \node[sollabel] at (.25,-.65) {质点受力};
 \end{scope}
\end{tikzpicture}
\soldiagcaption{左图描述状态变化；右图在空间中表示速度向右、阻力向左。}
\end{minipage}\end{center}}

\newcommand{\soldiagTwoFive}{%
\begin{center}\begin{minipage}{\linewidth}\centering
\begin{tikzpicture}[scale=0.88]
  \draw[axis] (-0.25,0)--(4.45,0) node[right,sollabel] {$x$};
  \draw[solforce] (0.40,0.30)--++(1.10,0) node[above,sollabel] {$\vec F$};
  \draw[solgeom,-{Latex[length=1.7mm]}] (0.40,-0.30)--++(0.90,0)
    node[midway,below,sollabel] {$v_0$};
  \draw[solgeom,-{Latex[length=1.7mm]}] (3.00,-0.30)--++(1.20,0)
    node[midway,below,sollabel] {$v_1$};
  \fill[Ink] (0.40,0) circle (1.3pt);
  \fill[Ink] (3.00,0) circle (1.3pt);
  \draw[<->,draw=TextGray] (0.40,0.65)--(3.00,0.65) node[midway,above,sollabel] {$s$};
\end{tikzpicture}
\soldiagcaption{取恒力方向为正，则 $a=F/m$，位移和速度均可使用匀加速公式。}
\end{minipage}\end{center}}

\newcommand{\soldiagTwoSix}{%
\begin{center}\begin{minipage}{\linewidth}\centering\begin{tikzpicture}[scale=.88]
 \draw[axis] (-.3,0)--(4.7,0) node[right,sollabel] {$x$};
 \draw[axis] (0,.5)--(0,-3.35) node[below,sollabel] {$y$};
 \node[above left,sollabel] at (0,0) {$O$};
 \draw[curve,domain=0:4.3,samples=90] plot (\x,{-.38*\x-.075*\x*\x});
 \coordinate (P) at (2.5,{-.38*2.5-.075*2.5*2.5});
 \draw[solgeom,-{Latex[length=2mm]}] (0,0)--(1.1,{-.38*1.1})
   node[pos=.65,above=14pt,sollabel] {$\vec v_0$};
 \fill (P) circle (1.4pt);
 \draw[solforce] (P)--++(0,-1.1) node[below,sollabel] {$\vec F$};
 \node[sollabel] at (3.3,-.65) {$y=Ax+Bx^2$};
\end{tikzpicture}
\soldiagcaption{轨迹从 $O$ 出发，$y$ 轴沿恒力方向；初速度箭头沿起点切线。}
\end{minipage}\end{center}}

\newcommand{\soldiagTwoSeven}{%
\begin{center}\begin{minipage}{\linewidth}\centering
\begin{tikzpicture}[scale=0.86]
  \coordinate (O) at (0,0);
  \coordinate (P) at (35:1.80);
  \draw[solbody] (O) circle (1.80);
  \fill[Ink] (O) circle (1.3pt) node[below left,sollabel] {$O$};
  \fill[Ink] (P) circle (1.4pt);
  \draw[solgeom,-{Latex[length=1.7mm]}] (P)--++(-55:1.00) node[right,sollabel] {$\vec v$};
  \draw[solforce] (P)--++(215:1.10) node[pos=.7,above left=5pt,sollabel] {$\vec F$};
  \draw[solgeom] (O)--(P) node[midway,below right,sollabel] {$R$};
  \draw[solgeom,-{Latex[length=1.5mm]}] ($(O)+(90:2.22)$)
    arc[start angle=90,end angle=35,radius=2.22];
  \node[sollabel] at (2.4,2.35) {向右转};
\end{tikzpicture}
\soldiagcaption{恒定大小且始终垂直于速度的力只改变速度方向，不改变速率。}
\end{minipage}\end{center}}

\newcommand{\soldiagTwoEight}{%
\begin{center}\begin{minipage}{\linewidth}\centering
\begin{tikzpicture}[scale=0.82]
  \node[solpanel] at (-2.25,1.50) {(整体)};
  \draw[solbody] (-3.80,-0.55)--(-0.60,-0.55);
  \draw[solbody,fill=white] (-3.20,-0.55) rectangle (-2.45,0.20);
  \draw[solbody,fill=white] (-1.75,-0.55) rectangle (-1.00,0.20);
  \draw[solbody] (-2.45,-0.17)--(-1.75,-0.17);
  \draw[solforce] (-1.00,-0.17)--++(0.90,0) node[right,sollabel] {$F$};
  \node[sollabel] at (-2.82,-0.16) {$m_2$};
  \node[sollabel] at (-1.38,-0.16) {$m_1$};

  \node[solpanel] at (2.10,1.50) {(分离 $m_2$)};
  \draw[solbody] (0.65,-0.55)--(3.80,-0.55);
  \draw[solbody,fill=white] (1.50,-0.55) rectangle (2.35,0.30);
  \draw[solforce] (2.35,-0.10)--++(0.95,0) node[right,sollabel] {$T$};
  \node[sollabel] at (1.92,-0.10) {$m_2$};
\end{tikzpicture}
\soldiagcaption{先对整体求共同加速度，再分离 $m_2$ 求绳张力，可避免同时解两个方程。}
\label{mech:legacy:ch06:s04-2-8}\end{minipage}\end{center}}

\newcommand{\soldiagTwoNine}{%
\begin{center}\begin{minipage}{\linewidth}\centering
\begin{tikzpicture}[scale=.9]
\foreach \x/\title/\fr in {0/共同运动/f,5.4/超过极限/\mu mg}{
 \begin{scope}[xshift=\x cm]
 \node[solpanel] at (0,2.1) {\title};
 \fill (0,1) circle(2pt) node[below=5pt,sollabel] {$m$};
 \draw[solforce] (0,1)--(1.35,1) node[midway,above=4pt,sollabel] {$\fr$};
 \fill (0,-.6) circle(2pt) node[below=5pt,sollabel] {$M$};
 \draw[solforce] (0,-.6)--(-1.35,-.6) node[midway,above=4pt,sollabel] {$\fr$};
 \draw[solforce] (0,-.6)--(1.5,-.6) node[midway,above=4pt,sollabel] {$F$};
 \end{scope}
}
\end{tikzpicture}
\soldiagcaption{分离两物块，只列水平力；上块摩擦向右，下块受到向左的反作用力。超过极限后用 $\mu mg$ 替代 $f$。}
\label{mech:legacy:ch06:s04-2-9}\end{minipage}\end{center}}

\newcommand{\soldiagTwoTen}{%
\begin{center}\begin{minipage}{\linewidth}\centering
\begin{tikzpicture}[scale=0.82]
  \draw[solbody] (-3.20,1.00)--(0.65,1.00);
  \draw[solbody,fill=white] (-2.45,1.00) rectangle (-1.55,1.90);
  \node[sollabel] at (-2.00,1.45) {$m_1$};
  \draw[solbody] (-1.55,1.90)--(1.15,1.90);
  \draw[solbody] (1.15,1.90) arc[start angle=90,end angle=0,radius=0.55];
  \draw[solbody] (1.70,1.35)--(1.70,-0.35);
  \draw[solbody] (1.15,1.35) circle (0.55);
  \draw[solbody,fill=white] (1.25,-1.20) rectangle (2.15,-0.35);
  \node[sollabel] at (1.70,-0.77) {$m_2$};
  \draw[solforce] (-2.00,1.90)--++(0.95,0) node[above,sollabel] {$T$};
  \draw[solforce] (1.70,-0.35)--++(0,0.90) node[right,sollabel] {$T$};
  \draw[solforce] (1.70,-1.20)--++(0,-.85) node[below,sollabel] {$m_2g$};
  \draw[solgeom,-{Latex[length=1.6mm]}] (-2.00,0.68)--++(0.85,0) node[below,sollabel] {$a$};
  \draw[solgeom,-{Latex[length=1.6mm]}] (2.45,-0.35)--++(0,-0.85) node[right,sollabel] {$a$};
\end{tikzpicture}
\soldiagcaption{轻绳不可伸长，故 $m_1$ 向右和 $m_2$ 向下的加速度大小相等。}
\label{mech:legacy:ch06:s04-2-10}\end{minipage}\end{center}}


质点是对运动物体的简化模型，它的运动学性质集中体现为它作为“点”的一面，可以用质
点的位置随时间的变化关系给出，在运动学的学习中已经广泛讨论了质点和由多个质点构
成的质点系统的运动。与此同时，质点还具有质量。质量是物体的一个基本属性；对于均匀的物质
来说，它的质量由物质的密度和体积的乘积共同决定；如果一个物体由多个部分组成，则其质量
为组成它的各部分质量之和。

具有相同速度但质量不同的质点，表现出的运动性质并不相同。为了体现这种不同，定义一个
新的物理量——动量。它是一个矢量，定义为质点的质量和速度的乘积，方向与质点的速度方向相同：
\begin{equation}
  \vec p=m\vec v.
\end{equation}
在相同外力作用下，质量越大的物体运动状态的变化越不明显；反之，力作用于质量较小的物体
时，则会引起较大的速度变化。

在大量经验和实验的基础上，英国物理学家牛顿给出了支配物体运动的三大定律：
\begin{description}[style=nextline,leftmargin=2em,labelindent=0pt]
  \item[第一定律]不受外力作用的物体必将保持静止或做匀速直线运动。
  \item[第二定律]物体的动量随时间的变化率等于所受的合外力；质量不变时，加速度由合外力决定。
  \item[第三定律]任意两个物体之间的相互作用力大小相等、方向相反，并且在一条直线上。
\end{description}
牛顿的三个运动定律是动力学的基础，相当精确地给出了宏观物体的运动规律。作为最基本的运动
定律，有必要深入讨论它们适用的条件，以及它们对自然现象的解释。

\subsection{第一定律：惯性}

牛顿第一定律又被称为惯性定律，指出了物体在不受外力作用时“应有”的运动状态，即运动学
中最简单的运动——匀速直线运动。惯性是任何物体所固有的属性，它使物体有保持原有运动状态，
也就是保持速度的趋势。惯性本身也可以度量，其大小的量度正是物体本身的质量。质量越大，物体
的惯性越大；在相同力的作用下，其运动状态的改变反而越小。

在运动学中可以在任意参考系中建立合适的坐标系来描写物体的运动，但在动力学中，牛顿定律并不
能适用于所有参考系，而只能适用于一类特殊的参考系，称为惯性参考系或惯性系。只有在惯性系中，
不受力物体的运动才是匀速直线运动。

惯性系的选取并不唯一。当参考系 $S$ 是惯性系时，所有相对于 $S$ 做匀速直线运动的参考系都是
惯性系。同一个不受力质点在不同惯性系中均做匀速直线运动，只不过运动速度不同。而相对于
$S$ 做加速运动的参考系则不是惯性系，牛顿定律若要在其中使用，必须加以改写。

关于是否存在真正的惯性参考系、在哪里能够找到惯性系，长期以来一直是极具争论的话题。绝大多数
力学现象都以地球表面为参考系来描写；实践表明，在地表参考系中应用牛顿定律，已经能够以相当高的
精度描述日常物体的运动。但是，对运动范围较广、速度较快或观测精度较高的过程，地表就不能再当作
严格的惯性系。地球不但自转，还绕太阳公转，这些转动都意味着地表参考系具有加速度。

以太阳为参考系时，牛顿定律能够在更高的精度上成立；然而天文观测又表明，太阳带着太阳系中的天体
以很高的速度绕银河系中心运动，虽然它的加速度很小，却也不是严格意义上的惯性参考系。关于惯性系
的争论始终存在，爱因斯坦的相对论给出了部分答案，同时又面临宇宙学问题的进一步检验。在实践中，
如果没有特别说明，描写地球上速度不高、范围不大、持续时间不长的物理过程时，通常近似取地表为
惯性系；相对于地面做加速运动的参考系则是非惯性系。

\noindent\textbf{【思考】}你能够发挥想象力，或通过已知的物理知识给出一个惯性参考系吗？

\begin{thinkanswer}
严格惯性系是一种理想化概念，实际中只能根据问题精度选择足够好的近似。在普通地面力学问题中，
地面参考系就是实用的近似惯性系；研究行星运动时，可选以太阳系质心为原点、坐标轴指向遥远天体的
参考系。判断关键不在于某个参考系的名字，而在于其非惯性效应是否小到可在当前精度下忽略。
\end{thinkanswer}

\subsection{第二定律：加速}

第一定律指出，不受力的物体在惯性系中运动速度不会发生变化；第二定律进一步告诉我们，
当物体受到力的作用时，其运动状态如何发生变化。利用动量，第二定律可写成
\begin{equation}
  \vec F=\frac{\dd\vec p}{\dd t}.
\end{equation}
在牛顿力学的常见问题中，物体的质量在运动过程中保持不变，根据动量的定义，上式化为更常用的形式
\begin{equation}
  \vec F=m\vec a.
\end{equation}
即质量与加速度的乘积等于物体所受的合外力，加速度的方向与合外力方向相同。使用牛顿第二定律
时需要注意：
\begin{itemize}
  \item 和第一定律一样，牛顿第二定律也只在惯性系中按上述形式成立。在非惯性系中，必须引入惯性力后才能使用类似形式。
  \item 经典形式适用于宏观、低速物体。接近光速的运动和微观粒子的运动需要其它物理理论。
  \item 当一个物体同时受到多个力时，决定加速度的 $\vec F$ 是所有外力的矢量和。
\end{itemize}

\noindent\textbf{【思考】}牛顿第一定律是否仅仅是第二定律在 $\vec F=0$ 时的特例？

\begin{thinkanswer}
不是简单的重复。如果已经预先确定了惯性系，由 $\vec F=m\vec a$ 确实可在 $\vec F=0$ 时得到 $\vec a=0$。
但第一定律同时给出了惯性系的判别标准，也说明物体保持匀速直线运动并不需要力来维持。第二定律要有明确物理意义，
已经隐含了“在惯性系中”这一前提；这个前提正是第一定律所确立的。
\end{thinkanswer}

第二定律有两类典型应用：一类是通过观察到的运动确定物体所受外力的性质；另一类是在已知所有外力
时预测物体未来的运动。例如：

研究物体运动的早期工作，常常从观察具体运动入手，再利用第二定律推测外力的大小、方向与物体性质、
外界条件之间的关系。运动学中出现的各种运动形式，从动力学角度看，都是力改变运动状态的结果；质量
为 $m$ 的物体若做某种已知运动，它在每一时刻所受合外力必为 $m$ 与该时刻加速度的乘积。通过分析
运动过程，便可反过来认识多种相互作用力。
\begin{itemize}
  \item 地表附近抛出的物体均做抛体运动，加速度大小为 $g$、方向竖直向下，因此质量为 $m$ 的物体在飞行过程中始终受到重力
  \begin{equation}
    F=mg.
  \end{equation}
  重力指向地面，其作用点可等效取在物体的重心。
  \item 做匀速圆周运动的质点具有指向圆心的加速度。质量为 $m$ 的质点以速率 $v$ 做半径为 $R$ 的圆周运动时，合外力必须指向圆心，大小为
  \begin{equation}
    F=m\frac{v^2}{R}.
  \end{equation}
  这个时刻指向圆心的合力通常称为向心力；若没有这样的合力，质点将依照第一定律沿直线运动。
  \item 太阳系各行星的运动满足开普勒三定律。通过分析行星运动可知，行星所受的力时刻指向太阳；在这些结论基础上，牛顿进一步作出大胆概括，得到万有引力定律。任意两个物体之间都有指向对方的引力，其大小为
  \begin{equation}
    F=G\frac{m_1m_2}{r^2}.
  \end{equation}
  万有引力是牛顿力学最重要的成果之一，后面将专门讨论它的性质和应用。
  \item 当物体的运动受到约束时会受到约束力。与重力、弹簧弹力等不同，约束力的大小和方向往往不能在一开始就决定，而要由物体所受的其它力以及运动状态共同确定。平面的弹力、连杆内的弹力、绳子的拉力都是典型的约束力；相当一部分问题的复杂性，正来自需要在运动的不同时刻确定这些约束力。
  \item 自然界还广泛存在各种阻力，它们会降低机械效率，也会增加预测运动的难度。当两物体接触且存在相对运动或相对运动趋势时，会出现滑动摩擦力或静摩擦力。常见的滑动摩擦力通常与接触面压力 $N$ 成正比，方向与相对运动方向相反：
  \begin{equation}
    f=\mu N.
  \end{equation}
  其中比例系数 $\mu$ 是摩擦因数，与接触物体的性质有关。水中或空气中运动的物体也会受到阻力；流体阻力由流体性质、运动物体的形状和速度共同决定，通常难以简单定量描述，后面将在若干理想化条件下讨论。
\end{itemize}

\subsection{从物理对象到运动方程}

牛顿定律的应用不是把若干力机械地写进公式，而是先说明研究对象、参考系和作用过程。
把一个物体单独取出时，凡是来自系统边界以外的相互作用都要画在受力图上；物体内部各部分
之间的作用不应重复计入。重力由质量和重力场决定，弹簧弹力由形变量决定，支持力、绳张力
和静摩擦力却通常要与运动约束一起求出。因此受力图上可以先标未知量，但不能凭经验预先
认定它们的大小，更不能把“运动方向”误当成“合力方向”。

选定惯性系后，质点的运动方程为
\begin{equation}
 \frac{\dd\vec p}{\dd t}=\sum_i\vec F_i,
 \qquad\text{质量不变时为}\qquad
 m\ddot{\vec r}=\sum_i\vec F_i.
\end{equation}
这是一个矢量方程。若在直角坐标系中使用，就应在每个方向分别投影；若轨迹弯曲，沿切向与
法向分解往往更直接：切向分量改变速率，法向分量改变速度方向。得到加速度以后，还必须结合
初始位置和初速度积分，运动才被唯一确定。由此可见，同一受力方程配上不同初始条件，可以
描述完全不同的运动。

约束把不同物体的运动联系起来。不可伸长轻绳给出绳长不变；对于方向不转动的平面接触，
保持接触意味着物体与平面沿固定法向的加速度相同；纯滚动给出接触点相对速度为零。
曲面接触则不能直接照搬“法向加速度相同”。例如质点沿固定圆轨道运动时，
由 $\vec r\cdot\vec r=R^2$ 先求导得 $\vec r\cdot\vec v=0$，再求导得
$\vec r\cdot\vec a=-v^2$，因此即使轨道静止，质点仍有大小 $v^2/R$ 的向心加速度。
一般约束须先写几何关系，再对时间求导，不能遗漏曲率或转动带来的项。
可靠的列式顺序是：分别画受力图，选取便于投影的方向，
为每个独立物体列动力学方程，再用几何或运动学约束把未知量闭合。方程数量不足时，往往不是
“少列了一条牛顿定律”，而是遗漏了一个约束关系。

牛顿第三定律说明相互作用总是成对出现：$\vec F_{A\to B}=-\vec F_{B\to A}$。这两个力作用在
不同物体上，所以不能在单个物体的受力图中相互抵消；只有把 $A,B$ 作为一个整体研究时，它们
才成为内力并在系统合力中消去。第三定律还要求同时、共线，并与两个物体是否静止无关。
处理多物体问题时，先整体求加速度、再隔离局部求张力或支持力，正是利用了这一结构。


\begin{origexamplebox}{}\label{ch06:ex:s04-2-1}
一个质量为 $m$，沿平面上 $x$ 轴运动质点的位置随时间的关系为
\[
  x(t)=A\cos\omega t,
\]
求它所受到外力随位置的关系。

\begin{qualityfigure}
\begin{tikzpicture}[>=Latex,font=\small]
\draw[->](-3,0)--(3.2,0)node[right]{$x$};\fill(0,0)circle(2pt)node[below]{$O$};
\fill(2.2,0)circle(2pt)node[below]{$A$};\fill(1.2,0)circle(2pt)node[above]{$P(x)$};\node at(2.2,.9){$t=0:\ x=A$};
\draw[dashed](-2.2,-.4)--(-2.2,.4)node[above]{$-A$};\draw[<->](-2.2,-.8)--(2.2,-.8)node[midway,below]{$x=A\cos\omega t$};
\end{tikzpicture}
\figcap{沿 x 轴的余弦位置函数与振幅边界}\label{mech:06:1:q}
\end{qualityfigure}
\end{origexamplebox}

\Needspace{6\baselineskip}\begin{xesolution}题图~\ref{mech:06:1:q}中，$O$ 是平衡位置，$P$ 的坐标为有符号的 $x$。以下求二阶导数时沿这个固定坐标轴投影，再判断力的符号与 $x$ 的关系。

对位置函数连续求两次导数：
\[
  v_x=\frac{\dd x}{\dd t}=-A\omega\sin\omega t,
  \qquad
  a_x=\frac{\dd^2x}{\dd t^2}=-A\omega^2\cos\omega t=-\omega^2x.
\]
因此外力与位置的关系为
\[
  \boxed{F_x(x)=-m\omega^2x}.
\]
负号的几何意义如解图所示：
\soldiagTwoOne
当 $x>0$ 时力向左，当 $x<0$ 时力向右；合力始终指向平衡位置 $x=0$。
\end{xesolution}

\begin{origexamplebox}{}\label{ch06:ex:s04-2-2}
在一固定平面上，一个质点的运动由方程
\[
  x(t)=A\cos\omega t,
  \qquad
  y(t)=A\sin\omega t
\]
给出，其中 $A>0$、$\omega>0$，质点质量为 $m$。求它在运动过程中的合外力。

\begin{qualityfigure}
\begin{tikzpicture}[>=Latex,font=\small]
\draw[DeepBlue,thick] (0,0) circle(1.5);\fill (0,0) circle(2pt) node[below left]{$O$};
\fill (1.5,0) circle(2pt) node[below right]{$P$};\draw[thin] (0,0)--(1.5,0) node[midway,below]{$A$};
\draw[->,DeepBlue,thick] (1.5,.1)--(1.5,1.3) node[right]{$v$};
\end{tikzpicture}
\figcap{给定参数方程在平面中的圆周}\label{mech:06:2:q}
\end{qualityfigure}
\end{origexamplebox}

\Needspace{6\baselineskip}\begin{xesolution}题图~\ref{mech:06:2:q}中，$OP=A$ 是到原点的距离；由参数方程可知初始点在正 $x$ 轴上且速度向上。以下分别求 $x,y$ 的二阶导数，检查合力是否指向 $O$。

对两个坐标分别求二阶导数：
\[
  \ddot x=-\omega^2x,
  \qquad
  \ddot y=-\omega^2y.
\]
故
\[
  \boxed{\vec F=-m\omega^2(x\,\hat{\boldsymbol i}+y\,\hat{\boldsymbol j})
  =-m\omega^2\vec r}.
\]
又因为 $x^2+y^2=A^2$，质点做半径为 $A$ 的匀速圆周运动。速度为
\[
 \vec v=A\omega(-\sin\omega t\,\hat{\boldsymbol i}
 +\cos\omega t\,\hat{\boldsymbol j}).
\]
从 $t=0$ 的 $(A,0)$ 点出发时速度向上，故运动方向是逆时针。图中 $P$ 位于第一象限，
其速度沿切线指向左上，而合力沿半径指向 $O$：
\soldiagTwoTwo
合力时刻指向圆心，大小为
\[
  \boxed{F=m\omega^2A}.
\]
\end{xesolution}

\begin{origexamplebox}{}\label{ch06:ex:s04-2-3}
质量为 $m$ 的物体受何种合外力，可以使其按下列方程运动？
\[
  x(t)=A\cos\omega t,
  \qquad y(t)=A\sin\omega t,
  \qquad z(t)=q(t)
\]
其中 $A>0$、$\omega>0$，$q(t)$ 是给定的二阶可导位置函数。另讨论 $q(t)=ut$、$u$ 为常量的情形。

\begin{qualityfigure}
\begin{tikzpicture}[>=Latex,font=\small]
\draw[->](0,-.2)--(0,4)node[above]{$z$};\draw[->](0,0)--(3,-.7)node[right]{$x$};\draw[->](0,0)--(-1.2,-.7)node[left]{$y$};
\draw[dashed](0,0)ellipse(1.3 and .35);\draw[DeepBlue,thick,domain=0:720,samples=140]plot({1.3*cos(\x)},{.0042*\x+.35*sin(\x)});
\node[align=left]at(4,2.4){$x=A\cos\omega t$\\$y=A\sin\omega t$\\$z=q(t)$};
\node[below]at(1.3,0){$A$};\node at(4,.8){螺旋线示意对应 $q(t)=ut$};
\end{tikzpicture}
\figcap{水平圆投影和独立给定的竖直位置}\label{mech:06:3:q}
\end{qualityfigure}
\end{origexamplebox}

\Needspace{6\baselineskip}\begin{xesolution}题图~\ref{mech:06:3:q}中，水平投影半径为 $A$，而竖直坐标由 $q(t)$ 独立给定。图中螺旋线仅示意 $q(t)=ut$ 的特例，不能据此在一般情形中令 $q''(t)=0$。

三个方向的加速度为
\[
  \ddot x=-A\omega^2\cos\omega t,
  \qquad
  \ddot y=-A\omega^2\sin\omega t,
  \qquad
  \ddot z=q''(t).
\]
因此所需合外力为
\[
  \boxed{\vec F(t)=
  -m\omega^2\bigl[x(t)\hat{\boldsymbol i}+y(t)\hat{\boldsymbol j}\bigr]
  +mq''(t)\hat{\boldsymbol k}}.
\]
水平合力指向 $z$ 轴，竖直合力由 $q''(t)$ 决定；不能仅凭水平投影为圆便断定合力没有竖直分量。
若 $q(t)=ut$，则 $\dot z=u$、$\ddot z=0$，轨迹为等螺距螺旋线：$u>0$ 时上升，
$u<0$ 时下降，$u=0$ 时退化为圆。图取 $u>0$，$H$ 是与 $P$ 等高的轴上点，
$PH$ 位于水平面内；$P_0$ 是 $P$ 在 $xy$ 平面上的投影。
\Needspace{12\baselineskip}
\soldiagTwoThree
在这一特例中合外力没有竖直分量，大小恒为 $m\omega^2A$。每转一周历时 $2\pi/\omega$，
竖直位移为 $2\pi u/\omega$。一般情况下，合力大小则为
$m\sqrt{\omega^4A^2+[q''(t)]^2}$；两项的量纲均为加速度的平方。
\end{xesolution}

\begin{origexamplebox}{}\label{ch06:ex:s04-2-4}
在一个研究质点在特定介质中运动所受阻力的实验中发现，一个质量为 $m$、初速度为 $v_0$
的质点在阻力作用下减速，且速度随它在介质中前进的距离的变化关系为
\[
  v(x)=v_0-\alpha x,\qquad v_0>0,\quad\alpha>0.
\]
探讨它所受阻力的形式。

\begin{qualityfigure}
\begin{tikzpicture}[>=Latex,font=\small]
\draw[->] (0,0)--(6.2,0) node[right]{$x$};\fill (1,0) circle(2pt) node[below]{$P$};
\draw[->,DeepBlue,thick] (1,.5)--(2.6,.5) node[right]{$v>0$};
\draw[->,Gold,thick] (2.6,1.3)--(1,1.3) node[midway,above]{$a<0$};
\end{tikzpicture}
\figcap{正向减速运动的速度与加速度方向}\label{mech:06:4:q}
\end{qualityfigure}
\end{origexamplebox}

\Needspace{6\baselineskip}\begin{xesolution}题图~\ref{mech:06:4:q}中，质点仍沿 $+x$ 运动，给定关系 $v=v_0-\alpha x$ 对应 $0\le x<v_0/\alpha$ 的运动阶段。要由它求力，须使用 $a=v\,\dd v/\dd x$，而非把 $\dd v/\dd x$ 直接当成加速度。

在质点仍沿 $+x$ 方向运动时，利用链式求导
\[
  a=\frac{\dd v}{\dd t}
  =\frac{\dd v}{\dd x}\frac{\dd x}{\dd t}
  =(-\alpha)v.
\]
所以
\[
  \boxed{F_x=-m\alpha v=-m\alpha(v_0-\alpha x)}.
\]
力与运动方向的关系如解图所示：
\soldiagTwoFour
因此阻力的大小与瞬时速度成正比，方向与速度相反。左图是速度与位置的函数图，右图才是在
真实空间中对质点作出的受力示意，二者不能混为同一幅位置图。
进一步用初值 $v(0)=v_0$ 解运动方程：
\[
 \frac{\dd v}{\dd t}=-\alpha v,
 \qquad\int_{v_0}^{v(t)}\frac{\dd v}{v}=-\alpha\int_0^t\dd t,
 \qquad\ln\frac{v(t)}{v_0}=-\alpha t.
\]
指数化得到 $v(t)=v_0\e^{-\alpha t}$。再取 $x(0)=0$，积分得
\[
 x(t)=\int_0^t v_0\e^{-\alpha\tau}\dd\tau
 =\left[-\frac{v_0}{\alpha}\e^{-\alpha\tau}\right]_0^t.
\]
因此
\[
  v(t)=v_0\e^{-\alpha t},
  \qquad
  x(t)=\frac{v_0}{\alpha}\bigl(1-\e^{-\alpha t}\bigr),
\]
即质点逐渐逼近 $x=v_0/\alpha$，而不是在有限时间内突然停下。
有限的 $t\ge0$ 对应 $0\le x<v_0/\alpha$，不能将直线向右延长后的负速度段当作这次实验的后续运动。
$\alpha$ 的单位为 $\mathrm{s^{-1}}$，所以 $m\alpha v$ 的单位确为牛顿。
\end{xesolution}

力学中的许多分支都是从观察运动物体开始，进而寻找特定环境下相互作用力的特点；再用牛顿定律预测
后续运动，并与测量结果比较，最终形成可靠的结论并应用于实践。

在已知力的性质后，可以利用牛顿第二定律计算未来的运动。当质量为 $m$ 的物体在某一时刻的位置
$\vec r$、速度 $\vec v$ 和所受外力 $\vec F$ 均为已知时，其加速度为
\begin{equation}
  \vec a=\frac{\vec F}{m}.
\end{equation}
在极小的时间间隔 $\Delta t$ 里，近似认为速度和受力的变化可以忽略，则在 $\Delta t$ 后有
\begin{equation}
  \Delta\vec r\approx\vec v\,\Delta t,
  \qquad
  \Delta\vec v\approx\frac{\vec F}{m}\Delta t.
\end{equation}
知道新的位置和速度后，可以重新计算受力并不断重复这一过程，从而近似预测任意时刻的位置和速度。

由第二定律可见，在恒力作用下物体做匀加速运动：加速度的大小等于力与质量的比值，方向与力的方向
一致。匀加速运动在运动学中已经得到充分研究，因此，只要对物体进行完整而无遗漏的受力分析，就能
确定它的加速度，进而求出受力作用下的运动。

如果所受外力并非恒力，例如引力随物体位置变化，就必须借助微积分处理。事实上，微积分的出现正是
牛顿研究引力作用下物体运动时受到推动而发展起来的。下面从相对简单的问题出发，逐步分析外力作用下
运动的性质，把得到的结论和方法用于越来越复杂的系统，并由此形成研究动力学问题的一般步骤。

\begin{origexamplebox}{}\label{ch06:ex:s04-2-5}
一个质量为 $m$ 的质点在恒力 $F$ 作用下沿直线加速运动，经过一段时间，它的速度由 $v_0$
变成了 $v_1$，求外力作用的时间以及外力作用下质点的位移。

\begin{qualityfigure}
\begin{tikzpicture}[>=Latex,font=\small]
\draw[->](0,0)--(6.5,0)node[right]{$x$};\fill(0,0)circle(2pt)node[below]{$O$};
\fill(1,0)circle(2pt)node[below]{$m$};\draw[->,DeepBlue,thick](1,.4)--(2.4,.4)node[right]{$v_0$};\draw[->,Gold,thick](1,1.1)--(2.8,1.1)node[right]{$F$};
\fill(5,0)circle(2pt)node[below]{末态};\draw[->,DeepBlue](5,.4)--(6.4,.4)node[above]{$v_1$};
\end{tikzpicture}
\figcap{恒力作用下的初末速度与位移方向}\label{mech:06:5:q}
\end{qualityfigure}
\end{origexamplebox}

\Needspace{6\baselineskip}\begin{xesolution}题图~\ref{mech:06:5:q}中，正方向沿恒力选取，初末速度应理解为该方向的有符号分量。时间由速度变化决定，位移则由初末速度平方之差决定。

取恒力方向为正方向，则运动中的各量如解图所示：
\soldiagTwoFive
根据牛顿第二定律，加速度为
\[
  a=\frac{F}{m}.
\]
由 $v_1=v_0+at$ 可得力的作用时间
\[
  \boxed{t=\frac{m(v_1-v_0)}{F}}.
\]
再利用 $v_1^2-v_0^2=2as$，得位移
\[
  \boxed{s=\frac{m(v_1^2-v_0^2)}{2F}}
  \qquad
  \left(=\frac{v_0+v_1}{2}\,t\right).
\]
式中 $v_0$、$v_1$ 是沿选定正方向的速度分量；若仅给出速率，还需要结合运动方向判定符号。
\end{xesolution}

\begin{origexamplebox}{}\label{ch06:ex:s04-2-6}
质量不变的质点受到非零恒定合力，且初速度垂直于该力的分量不为零。
证明轨迹是开口指向恒力方向的抛物线，并讨论初速度与恒力平行的退化情形。

\begin{qualityfigure}
\begin{tikzpicture}[>=Latex,font=\small]
\draw[->](-.2,0)--(4,0)node[right]{$x$};\draw[->](0,-.2)--(0,3.4)node[above]{$y$};
\fill(0,0)circle(2pt)node[below left]{$O$};\draw[->,DeepBlue,thick](0,0)--(2,1)node[right]{$\vect v_0$};
\draw[->,Gold,thick](2.7,1.2)--(2.7,2.8)node[right]{$\vect F$};\draw[dashed](2,1)--(2,0)node[below]{$u$};\node[left]at(0,1){$w$};\draw[dashed](0,1)--(2,1);
\end{tikzpicture}
\figcap{沿恒力与垂直恒力的坐标分解}\label{mech:06:6:q}
\end{qualityfigure}
\end{origexamplebox}

\Needspace{6\baselineskip}\begin{xesolution}题图~\ref{mech:06:6:q}中，$y$ 轴沿恒力，$x$ 轴垂直于恒力；初速度的水平分量 $u$ 非零时，$x=ut$ 才能用来消去时间。

该命题在初速度对恒力垂直方向的分量非零时成立。若初速度与恒力平行，轨迹退化为直线，不应再称为非退化抛物线。

在运动平面内取 $y$ 轴沿恒力 $\vec F$ 的方向，$x$ 轴与之垂直。设初速度分量为 $v_{0x}=u\ne0$、$v_{0y}=w$，并把初始位置取为原点。由解图所示的坐标选择，
\soldiagTwoSix
运动方程为
\[
  x=ut,
  \qquad
  y=wt+\frac{F}{2m}t^2.
\]
由 $t=x/u$ 消去时间，得
\[
  \boxed{y=\frac wu x+\frac{F}{2mu^2}x^2}.
\]
这是抛物线，二次项的方向即 $\vec F$ 的方向，所以开口指向恒力。若 $u=0$，则质点始终沿 $\vec F$ 所在直线运动，是退化情形。
\end{xesolution}

\begin{origexamplebox}{}\label{ch06:ex:s04-2-7}
一个质量为 $m$ 的质点在一平面内运动，开始时刻其运动速度为 $v$。在此后的运动过程中，它时刻受到一个与速度方向垂直、指向前进方向右侧且大小不变的力，大小为 $F$。试根据牛顿第二定律给出该质点的运动轨迹。如果该质点做周期运动，请给出周期与已知量的关系。

\begin{qualityfigure}
\begin{tikzpicture}[>=Latex,font=\small]
\fill(0,0)circle(2pt)node[below left]{$m$};\draw[->,DeepBlue,very thick](0,0)--(3,0)node[right]{$\vect v$};
\draw[->,Gold,very thick](0,0)--(0,-2)node[below]{$\vect F$（前进方向右侧）};\draw(.3,0)--(.3,-.3)--(0,-.3);
\end{tikzpicture}
\figcap{初始速度及始终指向前进方向右侧的力}\label{mech:06:7:q}
\end{qualityfigure}
\end{origexamplebox}

\Needspace{6\baselineskip}\begin{xesolution}题图~\ref{mech:06:7:q}中，垂直标记表示合力没有切向分量；“右侧”是随瞬时前进方向改变的方向。先由此确定速率不变，再分析法向曲率。

先沿速度方向列切向方程。因为力始终与速度垂直，$F_t=0$，所以
\[
 m\frac{\dd v}{\dd t}=0,\qquad v(t)=v.
\]
速率不变。用功和能量的语言，同一事实也可表述为瞬时功率
\[
  \mathcal P=\vec F\mathbin{\boldsymbol\cdot}\vec v=0,
\]
为零，因而动能不变。力只改变速度方向，而且始终指向轨迹右侧，如解图所示：
\soldiagTwoSeven
这个恒定的法向力提供向心加速度：
\[
  \frac{F}{m}=\frac{v^2}{R}.
\]
因此曲率半径 $R=mv^2/F$ 不随时间改变。在固定平面中，转向始终相同且曲率恒定的轨迹是圆；
本题始终向右转，所以沿顺时针方向绕行。半径和周期分别为
\[
  \boxed{R=\frac{mv^2}{F}},
  \qquad
  \boxed{T=\frac{2\pi R}{v}=\frac{2\pi mv}{F}}.
\]
\end{xesolution}

\subsection{第三定律：相互作用}

第三定律看似平淡，而且不像前两个定律那样直接描述单个物体的运动，但它在力学分析中不可缺少；深入
考察三大定律之间的内在关系，也会发现第三定律扮演着十分重要的角色。它告诉我们任意两个物体之间
相互作用力的性质：大小相等、方向相反，并且在一条直线上。利用这一点，可以在已知部分受力时准确找出
多个物体之间的受力关系。

从某种意义上说，万有引力定律的建立也以第三定律所揭示的相互作用结构为基础。地球在太阳引力作用下
沿椭圆轨道运动，同时地球也以同样大小的力吸引太阳；只因太阳质量远大于地球，在相同大小的力作用下
获得的加速度很小。

\noindent\textbf{【思考】}事实上，牛顿第三定律给出了质量严格且具有可操作性的定义，你能够独立地找到它吗？

\begin{thinkanswer}
设两个隔离外界影响的物体只相互作用。第三定律给出 $\vec F_{12}=-\vec F_{21}$，第二定律则给出
\[
  m_1\vec a_1=-m_2\vec a_2.
\]
因此两物体的惯性质量之比可由同一次相互作用中的加速度大小之比操作性地定义：
\[
  \boxed{\frac{m_1}{m_2}=\frac{|\vec a_2|}{|\vec a_1|}}.
\]
再选定一个标准物体的质量作为单位，就可以通过比较加速度测定其它物体的质量。
\end{thinkanswer}

\begin{origexamplebox}{}\label{ch06:ex:s04-2-8}
如图所示，光滑水平面上两个质量分别为 $m_1$、$m_2$ 的物体 $A$、$B$ 由不可伸长且质量可忽略不计的绳连接，有一个大小和方向均不变的力 $F$ 作用于 $A$，求两物体的加速度。

\begin{center}\begin{minipage}{\linewidth}\centering
\begin{tikzpicture}[scale=0.72]
  \draw[thick] (-3.2,-0.7)--(3.35,-0.7);
  \draw[thick] (-2.45,-0.7) rectangle (-1.25,0.5);
  \draw[thick] (0.55,-0.7) rectangle (1.75,0.5);
  \draw[thick] (-1.25,-0.1)--(0.55,-0.1);
  \draw[-{Latex[length=2.2mm]},thick] (1.75,-0.1)--(2.85,-0.1) node[right] {$F$};
  \node at (-1.85,-0.1) {$m_2$};
  \node[above] at (-1.85,.55) {$B$};
  \node at (1.15,-0.1) {$m_1$};
  \node[above] at (1.15,.55) {$A$};
\end{tikzpicture}
\end{minipage}\end{center}
\bookcaption{轻绳连接的两物体在光滑水平面上共同加速。}\label{ch06:fig:s04-2-8}
\end{origexamplebox}

\Needspace{6\baselineskip}\begin{xesolution}
以下推导所用的解图为图~\ref{mech:legacy:ch06:s04-2-8}。

轻绳不可伸长，所以两物体具有相同的向右加速度 $a$。先把 $A$、$B$ 作为整体，再分离 $B$，如解图所示：
\soldiagTwoEight
对整体而言，绳张力是内力，水平方向的外力只有 $F$，故
\[
  F=(m_1+m_2)a.
\]
因此两物体的加速度均为
\[
  \boxed{a_A=a_B=\frac{F}{m_1+m_2}}.
\]
由解图右侧对 $B$ 列式，还可得绳张力
\[
  T=m_2a=\frac{m_2F}{m_1+m_2}.
\]
\end{xesolution}

\begin{origexamplebox}{}\label{ch06:ex:s04-2-9}
如图，质量分别为 $m$ 和 $M$ 的物体放在光滑水平面上，现有一恒定大小的外力 $F$ 作用于
$M$。两块初始相对静止，两物体间的最大静摩擦因数和动摩擦因数均取为 $\mu$。
在上块尚未滑出下块的阶段，求两个物体的加速度。

\begin{center}\begin{minipage}{\linewidth}\centering
\begin{tikzpicture}[scale=0.74]
  \draw[thick] (-3.0,-1.0)--(3.1,-1.0);
  \draw[thick] (-1.75,-1.0) rectangle (1.65,0.15);
  \draw[thick] (-0.55,0.15) rectangle (0.75,1.25);
  \draw[-{Latex[length=2.2mm]},thick] (-2.75,-0.42)--(-1.75,-0.42) node[midway,above] {$F$};
  \node at (0.10,0.70) {$m$};
  \node at (-0.05,-0.42) {$M$};
\end{tikzpicture}
\end{minipage}\end{center}
\bookcaption{上物块 $m$ 放在下物块 $M$ 上，外力作用于 $M$。}\label{ch06:fig:s04-2-9}
\end{origexamplebox}

\Needspace{6\baselineskip}\begin{xesolution}
以下推导所用的解图为图~\ref{mech:legacy:ch06:s04-2-9}。

若两块共同向右运动，设共同加速度大小为 $a$。此时摩擦力必须带动上块，方向和两种运动分支如解图所示：
\soldiagTwoNine
对整体列式得
\[
  a=\frac{F}{M+m}.
\]
上块所需静摩擦力为
\[
  f_{\rm req}=ma=\frac{mF}{M+m}.
\]
而其最大值为 $f_{max}=\mu mg$，因此不滑动的条件为
\[
  \boxed{F\le\mu(M+m)g}.
\]
在此范围内
\[
  \boxed{a_m=a_M=\frac{F}{M+m}}.
\]
若 $F>\mu(M+m)g$，则两块发生相对滑动，摩擦力大小为 $\mu mg$，上、下两块的加速度大小分别为
\[
  \boxed{a_m=\mu g},
  \qquad
  \boxed{a_M=\frac{F-\mu mg}{M}},
\]
方向都向右，且 $a_M>a_m$，所以下块相对上块向右滑动。
\end{xesolution}

\begin{origexamplebox}{}\label{ch06:ex:s04-2-10}
两个质量分别为 $m_1$、$m_2$ 的物体由一根质量可忽略且不可伸长的绳连接，绳绕过质量
和轴摩擦均可忽略的定滑轮，如图放置。水平面光滑，已知 $m_2>m_1$，求两者的加速度。

\begin{center}\begin{minipage}{\linewidth}\centering
\begin{tikzpicture}[scale=0.74]
  \draw[thick] (-3.0,1.00)--(0.60,1.00);
  \draw[thick] (-2.45,1.00) rectangle (-1.55,1.90);
  \node at (-2.00,1.45) {$m_1$};
  \draw[thick] (-1.55,1.90)--(1.15,1.90);
  \draw[thick] (1.15,1.90) arc[start angle=90,end angle=0,radius=0.55];
  \draw[thick] (1.70,1.35)--(1.70,-0.55);
  \draw[thick] (1.15,1.35) circle (0.55);
  \draw[thick] (1.25,-1.45) rectangle (2.15,-0.55);
  \node at (1.70,-1.00) {$m_2$};
\end{tikzpicture}
\end{minipage}\end{center}
\bookcaption{水平物块与悬挂物通过定滑轮相连。}\label{ch06:fig:s04-2-10}
\end{origexamplebox}

\Needspace{6\baselineskip}\begin{xesolution}
以下推导所用的解图为图~\ref{mech:legacy:ch06:s04-2-10}。

轻绳不可伸长，因此 $m_1$ 向右和 $m_2$ 向下的加速度大小相等，记为 $a$。无摩擦定滑轮与轻绳使绳两端张力均为 $T$，如解图所示：
\soldiagTwoTen
分别对两物体沿运动方向列牛顿第二定律：
\[
  T=m_1a,
  \qquad
  m_2g-T=m_2a.
\]
相加消去 $T$，得
\[
  \boxed{a=\frac{m_2g}{m_1+m_2}}.
\]
所以 $m_1$ 向右、$m_2$ 向下，加速度大小均为上式。绳张力为
\[
  T=m_1a=\frac{m_1m_2g}{m_1+m_2}.
\]
\end{xesolution}
}

{
\graphicspath{{../../../sample/xes-original/fifth-section/}}
\section{牛顿定律的建模方法与典型应用}
% Original problem figures and added solution diagrams for Section 2.2.

\tikzset{
  solbody/.style={draw=Ink,line width=0.8pt},
  solforce/.style={-{Latex[length=2.0mm]},draw=DeepBlue,line width=0.9pt},
  solgeom/.style={draw=TextGray,densely dashed,line width=0.55pt},
  sollabel/.style={font=\footnotesize,fill=none,inner sep=0.7pt,outer sep=1.2pt,text=Ink},
  solpanel/.style={font=\scriptsize\bfseries,text=TextGray}
}

\newcommand{\diagcaption}[2]{\figcap{#1：#2}}

\newcommand{\soldiagTwoEleven}{%
\begin{center}\begin{minipage}{\linewidth}\centering\begin{tikzpicture}[scale=.78]
  \fill[Ink] (0,0) circle (5pt) node[below left,sollabel] {地球};
  \draw[solbody] (0,0) circle (2.25);
  \fill[DeepBlue] (2.25,0) circle (2.6pt) node[right,sollabel] {月球};
  \draw[solforce] (2.25,0)--(.85,0) node[midway,above,sollabel] {$a_{\rm M}$};
  \draw[solgeom] (0,0)--(2.25,0) node[midway,below,sollabel] {$r$};
  \draw[solgeom,-{Latex[length=1.6mm]}] (2.25,0)--++(0,.95) node[above,sollabel] {$v$};
\end{tikzpicture}\diagcaption{解图}{把月球公转近似为匀速圆周运动，用 $a=4\pi^2r/T^2$ 计算向心加速度。}\end{minipage}\end{center}}

\newcommand{\probdiagTwoTwelve}{%
\begin{center}\begin{minipage}{\linewidth}\centering\begin{tikzpicture}[scale=.78]
  \def\slopeang{30}
  \coordinate (V) at (-3.2,-1.0);
  \coordinate (C) at ($(V)+(\slopeang:4.85)$);
  \coordinate (T) at ($(C)+(120:.45)$);
  \coordinate (E) at ($(C)+(0:.45)$);
  \coordinate (Q) at ($(E)+(0,-1.65)$);
  % The rope is tangent to the pulley at T.  A is the midpoint of the
  % block's right face, and the block's lower face lies on the incline.
  \coordinate (A) at ($(T)+(210:1.60)$);
  \coordinate (B) at ($(A)+(210:.42)$);
  \draw[solbody] (V)--(C)--(C |- V)--cycle;
  \draw[solbody] (C) circle (.45);
  \draw[solbody] (A)--(T) arc[start angle=120,end angle=0,radius=.45]--(Q);
  \begin{scope}[shift={(B)},rotate=30]\draw[solbody,fill=white] (-.42,-.45) rectangle (.42,.45);\end{scope}
  \node[sollabel] at (B) {$m_1$};
  \draw[solbody,fill=white] ($(Q)+(-.28,0)$) rectangle ($(Q)+(.28,-.83)$);
  \node[sollabel] at ($(Q)+(0,-.415)$) {$m_2$};
  \draw ( -2.48,-1.0) arc[start angle=0,end angle=30,radius=.72];
  \node[sollabel] at (-2.10,-.78) {$\alpha$};
\end{tikzpicture}\bookcaption{图示第二次布置：$m_1$ 在斜面上，$m_2$ 悬空。}\label{ch06:fig:s05-2-12}\end{minipage}\end{center}}

\newcommand{\soldiagTwoTwelve}{%
\begin{center}\begin{minipage}{\linewidth}\centering\begin{tikzpicture}[scale=.84]
\foreach \x/\title/\massA/\massB/\tension in {0/第一次/m_2/m_1/T,5.8/第二次/m_1/m_2/T'}{
 \begin{scope}[xshift=\x cm]
 \node[solpanel] at (0,1.9) {\title};
 \fill[Ink] (-1,0) circle (2pt);
 \node[sollabel,above left=7pt] at (-1,0) {$\massA$};
 \draw[solforce] (-1,0)--++(30:1.3) node[above right,sollabel] {$\tension$};
 \draw[solforce] (-1,0)--++(210:1.3) node[below left,sollabel] {$\massA g\sin\alpha$};
 \node[sollabel] at (-1,-1.3) {沿斜面分量};
 \fill[Ink] (1.5,0) circle (2pt);
 \node[sollabel,right=8pt] at (1.5,0) {$\massB$};
 \draw[solforce] (1.5,0)--++(0,1.2) node[above,sollabel] {$\tension$};
 \draw[solforce] (1.5,0)--++(0,-1.2) node[below,sollabel] {$\massB g$};
 \end{scope}
}
\end{tikzpicture}\diagcaption{解图}{左侧只列沿斜面的力分量，右侧列悬挂物的竖直受力；两次分别使用 $T,T'$。}\label{mech:legacy:ch06:s05-2-12}\end{minipage}\end{center}}

% 题图仅呈现题目的 A、斜槽、竖直参考线和两角；H/K/h/ell 属于解答构造。
\newcommand{\probdiagTwoThirteen}{%
\begin{center}\begin{minipage}{\linewidth}\centering\begin{tikzpicture}[scale=.84]
\def\al{30}\def\be{40}
 \coordinate (H) at (0,0); \coordinate (A) at (0,3);
 \pgfmathsetmacro{\slotlength}{3*cos(\al)/cos(\be-\al)}
 \coordinate (B) at ($(A)+({-90+\be}:\slotlength)$);
 \draw[solbody] (-2,{-2*tan(\al)})--(3,{3*tan(\al)});
 \draw[solgeom] (-2,{-2*tan(\al)})--++(1.5,0);
 \draw[draw=TextGray] (-1.35,{-2*tan(\al)}) arc[start angle=0,end angle=\al,radius=.65];
 \node[sollabel] at (-1.05,-.95) {$\alpha$};
 \draw[solgeom] (A)--(H);
 \draw[solbody] (A)--(B);
 \fill (A) circle(1.5pt) node[above,sollabel] {$A$};
 \draw[draw=TextGray] ($(A)+(0,-.95)$) arc[start angle=-90,end angle={-90+\be},radius=.95];
 \node[sollabel] at ($(A)+(-70:1.3)$) {$\beta$};
\end{tikzpicture}\bookcaption{斜槽从 $A$ 通向斜面，与竖直方向的夹角为 $\beta$；斜面倾角为 $\alpha$。}\label{ch06:fig:s05-2-13}\end{minipage}\end{center}}


\newcommand{\soldiagTwoThirteen}{%
\begin{center}\begin{minipage}{\linewidth}\centering\begin{tikzpicture}[scale=.82]
\def\al{30}\def\be{40}
 \coordinate (H) at (0,0); \coordinate (A) at (0,3);
 \pgfmathsetmacro{\slotlength}{3*cos(\al)/cos(\be-\al)}
 \coordinate (B) at ($(A)+({-90+\be}:\slotlength)$);
 \coordinate (K) at (H |- B);
 \draw[solbody] (-2,{-2*tan(\al)})--(3,{3*tan(\al)});
 \draw[solgeom] (-2,{-2*tan(\al)})--++(1.5,0);
 \draw[draw=TextGray] (-1.35,{-2*tan(\al)}) arc[start angle=0,end angle=\al,radius=.65];
 \node[sollabel] at (-1.05,-.95) {$\alpha$};
 \draw[solgeom] (A)--(H);
 \draw[curve] (A)--(B) node[pos=.48,right=4pt,sollabel] {$\ell$};
 \draw[solgeom] (K)--(B);
 \fill (A) circle(1.5pt) node[above,sollabel] {$A$};
 \fill (B) circle(1.5pt) node[below right,sollabel] {$B$};
 \node[sollabel,below right=7pt] at (H) {$H$};
 \node[sollabel,left=3pt] at (K) {$K$};
 \draw[<->,draw=TextGray] (-.65,0)--(-.65,3) node[midway,left=3pt,sollabel] {$h$};
 \draw[draw=TextGray] ($(A)+(0,-.95)$) arc[start angle=-90,end angle={-90+\be},radius=.95];
 \node[sollabel] at ($(A)+(-70:1.3)$) {$\beta$};
 \coordinate (Q) at ($(A)!.4!(B)$);
 \fill (4.7,2.5) circle(2pt);
 \draw[solforce] (4.7,2.5)--(4.7,1.2) node[midway,left=5pt,sollabel] {$mg$};
 \draw[solgeom,->] (4.7,2.5)--++(-50:1.3) node[right=4pt,sollabel] {沿槽};
 \node[sollabel] at (4.9,.65) {$a=g\cos\beta$};
\end{tikzpicture}\diagcaption{解图}{$AK=\ell\cos\beta$，$KB=\ell\sin\beta$，$HK=KB\tan\alpha$，故 $AH=AK+HK$。}\label{mech:legacy:ch06:s05-2-13}\end{minipage}\end{center}}

\newcommand{\probdiagTwoFourteen}{%
\begin{center}\begin{minipage}{\linewidth}\centering\begin{tikzpicture}[scale=.78]
  \draw[solbody] (-3.6,-.9)--(3.7,-.9);
  \draw[solbody,fill=white] (-2.5,-.20) rectangle (1.35,.70);
  \draw[solbody] (-1.85,-.55) circle (.35) (.70,-.55) circle (.35);
  \draw[solbody,fill=white] (.67,.70) rectangle (1.32,1.35);
  \node[sollabel] at (-.60,.25) {$M$}; \node[sollabel] at (1.00,1.03) {$m$};
  \draw[solgeom] (-2.5,.70)--(-2.5,1.85) (1.35,1.35)--(1.35,1.85);
  \draw[<->,draw=TextGray] (-2.5,1.70)--(1.35,1.70) node[midway,above=2pt,sollabel] {$L$};
  \draw[solforce] (1.35,.10)--++(1.25,0) node[right,sollabel] {$v_0$};
\end{tikzpicture}\bookcaption{小物块相对地面静止地放到车顶前缘。}\label{ch06:fig:s05-2-14}\end{minipage}\end{center}}

\newcommand{\soldiagTwoFourteen}{%
\begin{center}\begin{minipage}{\linewidth}\centering\begin{tikzpicture}[scale=.80]
  \node[solpanel] at (-2.15,1.35) {物块};
  \draw[solbody,fill=white] (-2.65,-.25) rectangle (-1.65,.75);
  \draw[solforce] (-2.15,.25)--++(1.15,0) node[right=5pt,sollabel] {$f$};
  \node[solpanel] at (2.10,1.35) {小车};
  \draw[solbody,fill=white] (1.55,-.25) rectangle (2.65,.75);
  \draw[solforce] (2.10,.25)--++(-1.15,0) node[left=5pt,sollabel] {$f$};
  \draw[solgeom,-{Latex[length=1.6mm]}] (-2.15,-.65)--++(.85,0) node[below,sollabel] {$a_m$};
  \draw[solgeom,-{Latex[length=1.6mm]}] (2.10,-.65)--++(-.85,0) node[below,sollabel] {$a_M$};
\end{tikzpicture}\diagcaption{解图}{只列水平受力，两力大小均为 $f=\mu mg$；摩擦使物块加速、小车减速。}\label{mech:legacy:ch06:s05-2-14}\end{minipage}\end{center}}

\newcommand{\soldiagTwoFifteen}{%
\begin{center}\begin{minipage}{\linewidth}\centering\begin{tikzpicture}[scale=.9]
 \draw[axis] (-.4,0)--(5.7,0) node[right,sollabel] {$x/\mathrm m$};
 \foreach \x/\t in {0/A,1/撤力点,5/B}{\fill (\x,0) circle(1.5pt) node[below=5pt,sollabel] {\t};}
 \draw[solforce] (0,.7)--(1,.7) node[midway,above=3pt,sollabel] {$F$};
 \draw[solforce] (3.4,-.7)--(2.2,-.7) node[midway,below=3pt,sollabel] {$\mu mg$};
 \draw[<->,draw=TextGray] (0,1.65)--(1,1.65) node[midway,above,sollabel] {$1\,\mathrm m$};
 \draw[<->,draw=TextGray] (1,1.65)--(5,1.65) node[midway,above,sollabel] {$4\,\mathrm m$};
 \node[sollabel] at (3,-1.35) {撤力后 $t=2.0\,\mathrm s$};
\end{tikzpicture}\diagcaption{解图}{撤力点距 $A$ 为 $1\,\mathrm m$，随后减速滑行 $4\,\mathrm m$ 到 $B$；两段按同一长度比例绘制。}\end{minipage}\end{center}}

\newcommand{\probdiagTwoSixteen}{%
\begin{center}\begin{minipage}{\linewidth}\centering\begin{tikzpicture}[scale=.78]
  \draw[solbody] (-3.4,-.75)--(3.5,-.75);
  \draw[solbody,fill=white] (-2.35,-.75) rectangle (2.45,.80);
  \draw[solbody] (0.25,-.75)--(1.35,.80);
  \node[sollabel] at (-1.45,-.05) {$A$}; \node[sollabel] at (1.75,-.05) {$B$};
  \draw[draw=TextGray] (.70,-.75) arc[start angle=0,end angle=54.6,radius=.45];
  \node[sollabel] at (.79,-.48) {$\theta$};
  \draw[solforce] (-3.35,0)--(-2.35,0) node[midway,above,sollabel] {$F$};
\end{tikzpicture}\bookcaption{两物块的相互接触面光滑，整体向右加速。}\label{ch06:fig:s05-2-16}\end{minipage}\end{center}}

\newcommand{\soldiagTwoSixteen}{%
\begin{center}\begin{minipage}{\linewidth}\centering\begin{tikzpicture}[scale=.80]
  \node[solpanel] at (0,2.45) {分离物块 $B$};
  \draw[solbody,fill=white] (-.55,-.45) rectangle (.55,.65);
  \draw[solforce] (0,.65)--++(0,1.0) node[above,sollabel] {$N_g$};
  \draw[solforce] (0,-.45)--++(0,-1.0) node[below,sollabel] {$m_Bg$};
  \draw[solforce] (-.15,.10)--++(-35:1.10) node[below right,sollabel] {$N$};
  \draw[solforce] (0,.10)--++(-1.20,0) node[left,sollabel] {$\mu N_g$};
  \draw[solgeom,-{Latex[length=1.6mm]}] (.75,.10)--++(.95,0) node[right,sollabel] {$a$};
\end{tikzpicture}\diagcaption{解图}{由斜面法向力的竖直分量和水平分量确定共同加速度与地面压力。}\label{mech:legacy:ch06:s05-2-16}\end{minipage}\end{center}}

\newcommand{\probdiagTwoSeventeen}{%
\begin{center}\begin{minipage}{\linewidth}\centering\begin{tikzpicture}[scale=.78]
  \draw[solbody] (-3.3,-.9)--(3.4,-.9);
  \draw[solbody,fill=white] (-1.9,-.9) rectangle (1.5,.25);
  \draw[solbody,fill=white] (-.85,.25) rectangle (.55,1.30);
  \node[sollabel] at (-.15,.77) {$m_1$}; \node[sollabel] at (-.20,-.32) {$m_2$};
  \draw[solforce] (.55,.78)--++(1.45,0) node[right,sollabel] {$F=kt$};
\end{tikzpicture}\bookcaption{上块受到随时间线性增大的水平拉力。}\label{ch06:fig:s05-2-17}\end{minipage}\end{center}}

\newcommand{\soldiagTwoSeventeen}{%
\begin{center}\begin{minipage}{\linewidth}\centering\begin{tikzpicture}[x=1.25cm,y=.78cm]
 \draw[axis] (0,0)--(4.5,0) node[right,sollabel] {$t$};
 \draw[axis] (0,0)--(0,5.2) node[above,sollabel] {$a$};
 \draw[curve] (0,0)--(2,2)--(3.45,4.9) node[right=3pt,sollabel] {$a_1$};
 \draw[curve,dashed] (2,2)--(4.1,2) node[right=3pt,sollabel] {$a_2$};
 \draw[solgeom] (2,0) node[below=3pt,sollabel] {$t_0$}--(2,2);
 \draw[solgeom] (0,2) node[left=3pt,sollabel] {$a_c$}--(2,2);
 \node[sollabel,anchor=west] at (.4,4.5) {$t<t_0$：$a_1=a_2$};
 \node[sollabel,anchor=west] at (.4,3.85) {$a_c=\mu m_1g/m_2$};
\end{tikzpicture}\diagcaption{解图}{以 $m_1=m_2$ 示意两段斜率关系；达到 $t_0$ 后上块加速度增长更快，下块加速度保持 $a_c$。}\label{mech:legacy:ch06:s05-2-17}\end{minipage}\end{center}}

\newcommand{\probdiagTwoEighteen}{%
\begin{center}\begin{minipage}{\linewidth}\centering\begin{tikzpicture}[scale=.9]
 \coordinate (O) at (0,0); \coordinate (A) at (52:2.6);
 \draw[solbody] (O)--(A) node[midway,below right=5pt,sollabel] {$L$};
 \fill (O) circle(1.5pt) node[below left,sollabel] {$O$};
 \fill (A) circle(3pt) node[above right=5pt,sollabel] {$A\ (m)$};
 \draw[solgeom] (O)--(0,3);
 \draw[draw=TextGray] (0,.85) arc[start angle=90,end angle=52,radius=.85];
 \node[sollabel] at (71:1.12) {$\theta$};
 \draw[->,draw=TextGray] (40:3.15) arc[start angle=40,end angle=10,radius=3.15];
 \node[sollabel] at (23:3.52) {$\omega$};
\end{tikzpicture}\bookcaption{杆长为 $L$；$\theta$ 自竖直向上方向顺时针量起，角速度指向纸内。}\label{ch06:fig:s05-2-18}\end{minipage}\end{center}}

\newcommand{\soldiagTwoEighteen}{%
\begin{center}\begin{minipage}{\linewidth}\centering\begin{tikzpicture}[scale=1]
 \coordinate (A) at (0,0); \coordinate (G) at (0,-1.3);
 \coordinate (Q) at ({-2.6*sin(38)},{-2.6*cos(38)});
 \draw[solforce] (A)--(G) node[midway,right=4pt,sollabel] {$m\vec g$};
 \draw[solforce] (G)--(Q) node[midway,below=5pt,sollabel] {$\vec R$};
 \draw[solforce] (A)--(Q) node[midway,above left=5pt,sollabel] {$\vec F_c$};
 \fill (A) circle(2pt) node[above right,sollabel] {$A$};
 \node[sollabel,align=left,anchor=west] at (1,-.8) {$\vec F_c=-m\omega^2L\hat e_r$\\[5pt]$\vec R+m\vec g=\vec F_c$};
\end{tikzpicture}\diagcaption{解图}{力矢量首尾相接：重力向下，合力 $\vec F_c$ 指向圆心；杆对质点的力 $\vec R$ 一般不沿杆。}\label{mech:legacy:ch06:s05-2-18}\end{minipage}\end{center}}

\newcommand{\soldiagTwoNineteen}{%
\begin{center}\begin{minipage}{\linewidth}\centering\begin{tikzpicture}[scale=.80]
  \coordinate (O) at (0,0); \coordinate (P) at (1.55,.15);
  \draw[solbody] (O) ellipse (2.5 and .78); \fill[Ink] (O) circle (1.4pt) node[below,sollabel] {$O$};
  \fill[Ink] (P) circle (3pt) node[above,sollabel] {$m$};
  \draw[solgeom] (O)--(P) node[midway,below,sollabel] {$r$};
  \draw[solforce] (P)--($(P)!0.78!(O)$) node[midway,above,sollabel] {$f_s$};
  \draw[->,draw=TextGray] (-.20,1.18) arc[start angle=105,end angle=30,radius=1.0];
  \node[sollabel] at (.8,1.10) {$\omega$};
\end{tikzpicture}\diagcaption{解图}{静摩擦力提供全部向心力，极限状态下 $f_s=\mu mg$。}\end{minipage}\end{center}}

\newcommand{\soldiagTwoTwenty}{%
\begin{center}\begin{minipage}{\linewidth}\centering\begin{tikzpicture}[scale=.86]
  \draw[solbody] (0,1.38) circle (.44);
  \draw[solbody] (-.44,1.38) arc[start angle=180,end angle=0,radius=.44];
  \draw[solbody] (-.44,1.38)--(-.44,-1.08) (.44,1.38)--(.44,-2.05);
  \draw[solbody,fill=white] (-.86,-1.90) rectangle (-.02,-1.08)
    node[midway,sollabel] {$m_1$};
  % 人物的双手直接抓住右侧绳，明确“同学相对绳运动”的几何关系。
  \draw[solbody] (.78,-1.02) circle (.18);
  \draw[solbody] (.78,-1.20)--(.78,-1.72)
    (.78,-1.33)--(.44,-1.20)
    (.78,-1.33)--(.44,-1.46)
    (.78,-1.72)--(.57,-2.02)
    (.78,-1.72)--(1.02,-2.02);
  \fill[Ink] (.44,-1.20) circle (1.3pt) (.44,-1.46) circle (1.3pt);
  \node[below right=2pt,sollabel] at (.78,-2.02) {同学 $m_2$};
  \draw[solforce] (-.44,-1.08)--++(0,.92) node[left=2pt,sollabel] {$T$};
  \draw[solforce] (.66,-1.38)--++(0,.88) node[pos=.52,left=3pt,sollabel] {$T$};
  \draw[solgeom,-{Latex[length=1.8mm]}] (1.72,-1.82)--++(0,1.18)
    node[right=3pt,sollabel] {$a_{\rm rel}$};
\end{tikzpicture}\diagcaption{解图}{先用绳长约束把同学的绝对加速度与重物加速度联系起来。}\end{minipage}\end{center}}

\newcommand{\probdiagTwoTwentyOne}{%
\begin{center}\begin{minipage}{\linewidth}\centering\begin{tikzpicture}[scale=.85]
 \draw[solbody] (-2,3.9)--(2,3.9);
 \draw[solbody] (0,3.9)--(0,2.5);
 \draw[solforce] (.35,3.15)--(.35,3.8) node[pos=.4,right=5pt,sollabel] {$T$};
 \draw[solbody] (0,2.5) circle(.85);
 \fill (0,2.5) circle(1.5pt);
 \node[sollabel] at (0,2.15) {$A$};
 \draw[solbody] (-.85,2.5) arc[start angle=180,end angle=0,radius=.85];
 \draw[solbody] (-.85,2.5)--(-.85,-.05) (.85,2.5)--(.85,.65);
 \node[sollabel,left=4pt] at (-.85,1.55) {$H$};
 \node[sollabel,right=4pt] at (.85,1.55) {$H$};
 \draw[solbody,fill=white] (.5,-.1) rectangle (1.2,.65) node[midway,sollabel] {$M_1$};
 \draw[solbody] (-.85,-.05) circle(.55);
 \fill (-.85,-.05) circle(1.5pt);
 \node[sollabel,left=5pt] at (-1.4,-.05) {$M_2$};
 \draw[solbody] (-1.4,-.05) arc[start angle=180,end angle=0,radius=.55];
 \draw[solbody] (-1.4,-.05)--(-1.4,-1.6) (-.3,-.05)--(-.3,-2.0);
 \node[sollabel,left=4pt] at (-1.4,-.85) {$T_1$};
 \node[sollabel,right=4pt] at (-.3,-.85) {$T_1$};
 \draw[solbody,fill=white] (-1.72,-2.25) rectangle (-1.08,-1.6) node[midway,sollabel] {$m_2$};
 \draw[solbody,fill=white] (-.62,-2.65) rectangle (.02,-2) node[midway,sollabel] {$m_1$};
\end{tikzpicture}\bookcaption{上绳接在动滑轮的轴心；顶端悬挂力为 $T$，绕定滑轮的绳张力为 $H$，下绳张力为 $T_1$。}\label{ch06:fig:s05-2-21}\end{minipage}\end{center}}

\newcommand{\soldiagTwoTwentyOne}{%
\begin{center}\begin{minipage}{\linewidth}\centering\begin{tikzpicture}[scale=.9]
 \node[solpanel] at (-2.4,2.5) {上绳：两端连接轴心与 $M_1$};
 \draw[solbody] (-2.4,1.4) circle(.6);
 \draw[solbody] (-3,1.4)--(-3,-.4) (-1.8,1.4)--(-1.8,-.15);
 \draw[solbody] (-3,-.4) circle(.35);
 \fill (-3,-.4) circle(1.5pt);
 \node[sollabel,left=3pt] at (-3.35,-.4) {$M_2$};
 \draw[solbody] (-2.1,-.8) rectangle (-1.5,-.15) node[midway,sollabel] {$M_1$};
 \draw[solgeom,->] (-3.8,1.1)--(-3.8,.2) node[midway,left,sollabel] {$a_p$};
 \draw[solgeom,->] (-1.1,1.1)--(-1.1,.2) node[midway,right,sollabel] {$a_{M_1}$};
 \node[sollabel] at (-2.4,-1.3) {$a_{M_1}+a_p=0$};
 \node[solpanel] at (3,2.5) {下绳：两端连接 $m_2,m_1$};
 \draw[solbody] (3,1.4) circle(.6);
 \draw[solbody] (2.4,1.4)--(2.4,-.15) (3.6,1.4)--(3.6,-.15);
 \draw[solbody] (2.1,-.8) rectangle (2.7,-.15) node[midway,sollabel] {$m_2$};
 \draw[solbody] (3.3,-.8) rectangle (3.9,-.15) node[midway,sollabel] {$m_1$};
 \draw[solgeom,->] (1.7,1.1)--(1.7,.2) node[midway,left,sollabel] {$a_2$};
 \draw[solgeom,->] (4.3,1.1)--(4.3,.2) node[midway,right,sollabel] {$a_1$};
 \draw[solgeom,->] (5.3,1.9)--(5.3,1) node[midway,right=5pt,sollabel] {$a_p$};
 \node[sollabel] at (3,-1.3) {$a_1+a_2=2a_p$};
\end{tikzpicture}\diagcaption{解图}{物体位置与题图保持一致；箭头表示向下取正，并不预先断言实际加速度方向。}\label{mech:legacy:ch06:s05-2-21}\end{minipage}\end{center}}

\newcommand{\probdiagTwoTwentyTwo}{%
\begin{center}\begin{minipage}{\linewidth}\centering\begin{tikzpicture}[scale=.84]
 \def\th{30}
 \coordinate (L) at (-3,-1); \coordinate (R) at (2.6,-1);
 \coordinate (P) at ($(L)+({5.6},{5.6*tan(\th)})$);
 \coordinate (S) at ($(L)!.58!(P)$);
 \coordinate (B) at ($(S)+({90+\th}:.44)$);
 \draw[solbody] (L)--(R)--(P)--cycle;
 \begin{scope}[shift={(B)},rotate=\th]\draw[solbody,fill=white] (-.45,-.44) rectangle (.45,.44);\end{scope}
 \node[sollabel] at (B) {$m$}; \node[sollabel] at (1.6,-.5) {$M$};
 \draw (L)++(.7,0) arc[start angle=0,end angle=\th,radius=.7];
 \node[sollabel] at ($(L)+(15:1.05)$) {$\theta$};
\end{tikzpicture}\bookcaption{物块底面与斜面重合；斜面可在光滑水平面上自由运动。}\label{ch06:fig:s05-2-22}\end{minipage}\end{center}}

\newcommand{\soldiagTwoTwentyTwo}{%
\begin{center}\begin{minipage}{\linewidth}\centering\begin{tikzpicture}[scale=.95]
 \node[solpanel] at (-2.5,2) {小物块受力};
 \coordinate (C) at (-2.5,0);
 \draw[solgeom] ($(C)+(210:1.6)$)--($(C)+(30:1.6)$);
 \fill (C) circle(2pt);
 \draw[solforce] (C)--++(120:1.5) node[above left=3pt,sollabel] {$N$};
 \draw[solforce] (C)--++(0,-1.5) node[below=3pt,sollabel] {$mg$};
 \node[solpanel] at (2,2) {地面系加速度分解};
 \coordinate (O) at (3,.7);
 \coordinate (Q) at ($(O)+({-2.4*cos(30)},-1.2)$);
 \coordinate (P) at ($(Q)+({1.2*cos(30)},0)$);
 \draw[solgeom,->] (O)--(Q) node[pos=.45,above left=5pt,sollabel] {$\ddot s\,\hat e_s$};
 \draw[solgeom,->] (Q)--(P) node[midway,below=5pt,sollabel] {$\ddot X\,\hat e_x$};
 \draw[solforce] (O)--(P) node[pos=.5,right=5pt,sollabel] {$\vec a_m$};
 \node[sollabel] at (2,-1.6) {$\vec a_m=\ddot s\,\hat e_s+\ddot X\,\hat e_x$};
\end{tikzpicture}\diagcaption{解图}{右图取 $M=m$ 示意，$\hat e_s$ 沿斜面向下；地面系加速度比斜面方向更陡，不沿斜面。}\label{mech:legacy:ch06:s05-2-22}\end{minipage}\end{center}}

\newcommand{\probdiagTwoTwentyThree}{%
\begin{center}\begin{minipage}{\linewidth}\centering\begin{tikzpicture}[scale=.85]
 \def\th{25}
 \coordinate (L) at (-2.5,-1); \coordinate (R) at (2.5,-1);
 \coordinate (P) at ($(L)+(0,{5*tan(\th)})$);
 \coordinate (C) at ($(P)!.42!(R)$);
 \draw[solbody] (L)--(R)--(P)--cycle;
 \draw[solbody,line width=2pt] (C)--++(0,2);
 \draw[solbody] ($(C)+(-.2,.6)$)--++(0,1.1) ($(C)+(.2,.6)$)--++(0,1.1);
 \node[sollabel,right=5pt] at ($(C)+(.2,1.4)$) {固定导轨};
 \node[sollabel,left=6pt] at ($(C)+(0,1.1)$) {$m$};
 \node[sollabel] at (.7,-.6) {$M$};
 \fill (C) circle(1.5pt) node[below right=3pt,sollabel] {$C$};
 \draw (R)++(155:.8) arc[start angle=155,end angle=180,radius=.8];
 \node[sollabel] at ($(R)+(167.5:1.15)$) {$\theta$};
 \draw[solforce] (-3.7,-.2)--(-2.55,-.2) node[midway,above,sollabel] {$F$};
 \draw[solgeom,->] (.5,-1.4)--(1.7,-1.4) node[right,sollabel] {$a$};
\end{tikzpicture}\bookcaption{杆的下端 $C$ 接触斜面；固定光滑导轨限制杆只能竖直运动，斜面倾角在右端标出。}\label{ch06:fig:s05-2-23}\end{minipage}\end{center}}

\newcommand{\soldiagTwoTwentyThree}{%
\begin{center}\begin{minipage}{\linewidth}\centering\begin{tikzpicture}[scale=.95]
 \node[solpanel] at (-2.2,2) {杆的受力};
 \fill (-2.2,0) circle(2pt);
 \draw[solforce] (-2.2,0)--++(65:1.4) node[above,sollabel] {$N$};
 \draw[solforce] (-2.2,0)--++(0,-1.4) node[below,sollabel] {$mg$};
 \draw[solforce] (-2.2,0)--++(-1.2,0) node[left,sollabel] {$R_h$};
 \draw[solgeom,->] (-1.1,-.4)--(-1.1,.7) node[right,sollabel] {$a\tan\theta$};
 \node[solpanel] at (2.5,2) {斜面的水平受力};
 \fill (2.5,0) circle(2pt);
 \draw[solforce] (2.5,0)--++(-1.4,0) node[left,sollabel] {$N\sin\theta$};
 \draw[solforce] (2.5,0)--++(1.4,0) node[right,sollabel] {$F$};
 \draw[solgeom,->] (2,-.8)--(3.1,-.8) node[right,sollabel] {$a$};
\end{tikzpicture}\diagcaption{解图}{左图保留固定导轨的水平约束力 $R_h$；右图只列斜面的水平力分量，二者的加速度方向不同。}\label{mech:legacy:ch06:s05-2-23}\end{minipage}\end{center}}


将牛顿三定律与受力对象的性质和特点有机地结合在一起，可以解决几乎所有物体在受力作用下的
运动问题。但是，随着相互作用物体的增加、系统本身复杂性的提升，预测运动将变得越来越复杂。
物体在受力作用下的运动问题将贯穿整个物理学。我们将从简单问题开始，掌握受力分析、运动学关系
与约束条件相结合的一般技术；随着学习不断推进，将能够解决越来越多、越来越复杂的问题。

\subsection{牛顿定律应用的一般次序}

复杂装置仍由若干简单物体和约束组成。第一步确定哪些物体具有相同大小的加速度，哪些方向
必须分别处理；第二步对每个研究对象画受力图；第三步选轴并列出分量方程；最后才代入运动学
关系。在忽略绳质量、滑轮转动惯量和轮轴摩擦，且绳上没有其他切向分布力的理想模型中，
同一根绳中的张力才可视为处处相等。绳长不变则给出绳各段位移、速度和加速度之间的几何约束；
有动滑轮时须把所有可变绳段计入，不能一概认为两端加速度大小相等。
绳或滑轮有质量时，张力通常不再相等，但只要绳仍不可伸长，绳长约束仍成立。

摩擦力尤其不能先写成 $\mu N$。静摩擦力在零到最大值之间自适应，只有由平衡或动力学方程
算出的需要值满足 $|f_s|\leq\mu_sN$，原先假定的相对静止才可能成立；若超过极限，物体发生
滑动，此时才用 $f_k=\mu_kN$，方向与接触面的相对滑动速度相反。在叠放物块、传送带和斜面
问题中，应先假定一种运动状态，算出所需摩擦，再用不等式检验，必要时分段讨论。

圆周运动也不对应一种新的力。“向心力”只是合力的法向分量，大小为
\begin{equation}
 F_n=m\frac{v^2}{R}=m\omega^2R.
\end{equation}
重力、支持力、张力或摩擦力中的某些分量共同提供它。竖直圆周运动中速率随高度改变，最低点
与最高点的受力方程不能混用；轻绳模型还要检查张力不能为负，轨道约束则要检查支持力不能
指向轨道不可能施力的一侧。这些条件常常决定物体能否完成整周运动。

万有引力问题把上述思路推广到距离变化的中心力。球对称天体外部可等效为质量集中在球心，
$F=GMm/r^2$；近地面小范围内才可把 $g$ 当作常量。月球的向心加速度与地表重力的比较说明，
落体与天体运动可以由同一种相互作用解释。计算以后还应检查量纲、方向和极限情形，避免只得到
形式正确却物理上不可能的根。


\begin{origexamplebox}{}\label{ch06:ex:s05-2-11}
地球的半径约为 $6.4\times10^3\,\mathrm{km}$，地表附近重力加速度为
$9.8\,\mathrm{m/s^2}$；通过天文观测发现月球的公转轨道平均半径约为
$3.8\times10^5\,\mathrm{km}$，公转周期约为 $27.3$ 天。传说牛顿在被苹果砸到以后开始
考虑导致地球表面物体下落和月球公转的是同一种力。你能否通过上面的数据给出月球公转
加速度与地球表面重力加速度的关系？

\begin{qualityfigure}
\begin{tikzpicture}[>=Latex,font=\small]
\draw[fill=PaleBlue](0,0)circle(.65);\node at(0,.95){地球};\fill(0,0)circle(1.5pt)node[below left]{$O$};
\draw[DeepBlue](0,0)circle(2.4);\fill(2.4,0)circle(2.5pt)node[right]{月球};
\draw[dashed](0,0)--(2.4,0)node[midway,above]{$r$};\draw(0,0)--(-.46,.46)node[above left=6pt]{$R_E$};
\node[align=left]at(5,1.1){$r=3.8\times10^5\,\mathrm{km}$\\$T=27.3$ 天};\node at(1,-3){示意不按实际比例绘制};
\end{tikzpicture}
\figcap{地表重力尺度与月球轨道尺度}\label{mech:06:11:q}
\end{qualityfigure}
\end{origexamplebox}

\Needspace{6\baselineskip}\begin{xesolution}题图~\ref{mech:06:11:q}中，$r$ 从地心 $O$ 量到月球，$R_E$ 是地球半径；比较平方反比规律时应使用这两个到地心的距离，而非月球离地表的高度。

把月球公转近似为半径为 $r$ 的匀速圆周运动，解图给出了向心加速度的方向：
\soldiagTwoEleven
取
\[
 r=3.8\times10^8\,\mathrm m,
 \qquad
 T=27.3\times86400\,\mathrm s,
\]
则
\[
 a_{\rm M}=\frac{4\pi^2r}{T^2}
 \approx2.70\times10^{-3}\,\mathrm{m/s^2}.
\]
因此
\[
 \boxed{\frac{a_{\rm M}}{g}\approx2.75\times10^{-4}
 \approx\frac1{3.63\times10^3}}.
\]
另一方面，$r/R_{\rm E}\approx59.4$，故
\[
 \left(\frac{R_{\rm E}}r\right)^2\approx2.84\times10^{-4},
\]
与上面的观测结果十分接近。这正是“引力强度随距离平方反比衰减”的重要线索。
\end{xesolution}

\begin{origexamplebox}{}\label{ch06:ex:s05-2-12}
质量分别为 $m_1$ 和 $m_2$ 的两个小物块用轻绳连接，绳跨过位于倾角
$\alpha=30^\circ$ 的光滑斜面顶端的轻滑轮，滑轮与转轴之间的摩擦不计，斜面固定在水平
桌面上，如图所示。第一次，$m_1$ 悬空，$m_2$ 放在斜面上，用 $t$ 表示 $m_2$ 自斜面
底端由静止开始运动至斜面顶端所需的时间。第二次，将 $m_1$ 和 $m_2$ 位置互换，使
$m_2$ 悬空，$m_1$ 放在斜面上，发现 $m_1$ 自斜面底端由静止开始运动至斜面顶端所需
的时间为 $t/3$。求 $m_1$ 与 $m_2$ 之比。
\probdiagTwoTwelve
\end{origexamplebox}

\Needspace{6\baselineskip}\begin{xesolution}
以下推导所用的解图为图~\ref{mech:legacy:ch06:s05-2-12}。

设斜面长度为 $s$。两次都是从静止出发并通过同一段距离，所以
\[
 s=\frac12a_1t^2
   =\frac12a_2\left(\frac t3\right)^2,
 \qquad\Rightarrow\qquad a_2=9a_1.
\]
沿绳方向列牛顿第二定律，受力关系如解图所示：
\soldiagTwoTwelve
第一次 $m_1$ 下落、$m_2$ 沿斜面上滑，
分别取两者的运动方向为正，绳张力记为 $T$，有
\[
 m_1g-T=m_1a_1,\qquad T-m_2g\sin\alpha=m_2a_1.
\]
两式相加，张力消去，得到
\[
 a_1=g\frac{m_1-m_2\sin\alpha}{m_1+m_2}.
\]
第二次 $m_2$ 下落、$m_1$ 沿斜面上滑，
张力改记为 $T'$，两物体分别满足
\[
 m_2g-T'=m_2a_2,\qquad T'-m_1g\sin\alpha=m_1a_2.
\]
相加可得
\[
 a_2=g\frac{m_2-m_1\sin\alpha}{m_1+m_2}.
\]
令 $x=m_1/m_2$，并代入 $\sin30^\circ=1/2$、$a_2=9a_1$：
\[
 \frac{1-x/2}{x+1}=9\frac{x-1/2}{x+1}.
\]
解得
\[
 \boxed{\frac{m_1}{m_2}=\frac{11}{19}}.
\]
该结果同时满足两次运动方向所要求的 $m_1>m_2/2$ 和 $m_2>m_1/2$。
\end{xesolution}

\begin{origexamplebox}{}\label{ch06:ex:s05-2-13}
一个质点自倾角为 $\alpha$ 的斜面上方定点 $A$，沿光滑斜槽从静止开始滑下。为了使质点
在最短时间到达斜面，求斜槽与竖直方向的夹角 $\beta$ 应等于多少。
\probdiagTwoThirteen
\end{origexamplebox}

\Needspace{6\baselineskip}\begin{xesolution}
以下推导所用的解图为图~\ref{mech:legacy:ch06:s05-2-13}。

取 $0<\alpha<\pi/2$，研究位于图示竖直平面内、向右下方倾斜的直槽。
设 $H$ 是 $A$ 正下方的斜面点，$AH=h$；槽端为 $B$，槽长 $AB=\ell$，
槽与竖直方向的夹角为 $0\le\beta<\pi/2$。图中 $BK$ 水平、$AK$ 竖直：
\Needspace{13\baselineskip}
\soldiagTwoThirteen
沿槽加速度为 $g\cos\beta$。由三角形 $AKB$，$BK=\ell\sin\beta$、$AK=\ell\cos\beta$；
又由斜面倾角，$KH=BK\tan\alpha$。利用 $AH=AK+KH$，有
\[
 \ell\cos\beta=h-\ell\sin\beta\tan\alpha,
 \qquad \ell(\cos\beta\cos\alpha+\sin\beta\sin\alpha)=h\cos\alpha.
\]
这说明竖直下降高度 $AK$ 随槽的方向改变，不能把它误当成固定的 $h$。槽长为
\[
 \ell=\frac{h\cos\alpha}{\cos(\beta-\alpha)}.
\]
又因 $\ell=\tfrac12g\cos\beta\,t^2$，所以
\[
 t^2=\frac{2h\cos\alpha}
 {g\cos\beta\cos(\beta-\alpha)}.
\]
令分母
$D=\cos\beta\cos(\beta-\alpha)$ 最大。利用积化和差，
\[
 D=\frac12\bigl[\cos\alpha+\cos(2\beta-\alpha)\bigr],
\]
显然在 $2\beta-\alpha=0$ 时取得最大值。因此
\[
  \boxed{\beta=\frac\alpha2}.
\]
它落在允许区间内。此时
\[
 t_{\min}^2=\frac{4h\cos\alpha}{g(1+\cos\alpha)}.
\]
当 $\alpha\to0$ 时，最优槽趋于竖直，$t_{\min}\to\sqrt{2h/g}$，与自由落下高度 $h$ 一致。
\end{xesolution}

在系统中包含摩擦力时要格外小心。两固体表面之间的滑动摩擦力正比于相互之间的压力，
比例系数称为摩擦系数；如果两物体保持相对静止但存在相对运动趋势，则由静摩擦力维持，
其大小介于零与最大静摩擦力之间。分析含摩擦的运动问题时，通常需要先判断相互接触的
物体间究竟有无相对运动；当速度或外力发生改变时，两种状态之间还可能发生转化。

\begin{origexamplebox}{}\label{ch06:ex:s05-2-14}
一质量为 $M$ 的平顶小车，以速度 $v_0$ 沿水平的光滑轨道作匀速直线运动。现将一质量为
$m$、相对地面初速度为零的小物块放置在车顶前缘。已知物块和车顶之间的滑动摩擦系数为 $\mu$，
物块按质点处理，忽略放置时的竖直冲击和车轮转动惯量。
若要求物块不会从车顶后缘掉下，则该车顶长度至少要多少？
\probdiagTwoFourteen
\end{origexamplebox}

\Needspace{6\baselineskip}\begin{xesolution}
以下推导所用的解图为图~\ref{mech:legacy:ch06:s05-2-14}。

物块相对地面初速度为零，小车初速度为 $v_0$。滑动阶段的受力如解图所示：
\soldiagTwoFourteen
物块受摩擦向右，$a_m=\mu g$；小车受摩擦向左，$a_M=-\mu mg/M$。在随小车平动的相对
坐标中，物块相对小车的初速度和加速度分别为
\[
 v_{\rm rel,0}=-v_0,
 \qquad
 a_{\rm rel}=a_m-a_M=\mu g\left(1+\frac mM\right).
\]
相对速度减为零前，物块相对车向后滑过的距离为
\[
 L_{\min}=\frac{v_0^2}{2a_{\rm rel}}
 =\boxed{\frac{Mv_0^2}{2\mu g(M+m)}}.
\]
到相对速度为零时两者取得共同速度；之后不再需要滑动摩擦，物块恰不会从后缘掉下。
达到共同速度所需时间为
\[
 t_* =\frac{Mv_0}{\mu g(M+m)},\qquad
 v_*=\mu g t_* =\frac{Mv_0}{M+m}.
\]
于是 $(M+m)v_*=Mv_0$，与系统水平动量不变相符。$M\gg m$ 时车速几乎不变，
$L_{\min}\to v_0^2/(2\mu g)$，也符合物块在匀速车面上减小相对速度的极限。
\end{xesolution}

\begin{origexamplebox}{}\label{ch06:ex:s05-2-15}
质量 $m=1.5\,\mathrm{kg}$ 的物块在水平恒力 $F$ 作用下从水平面上 $A$ 点由静止开始运动，
运动一段距离之后撤去该力，物块继续滑行 $t=2.0\,\mathrm s$ 后停在 $B$ 点。已知
$A,B$ 两点的距离 $x=5.0\,\mathrm m$，物块与水平面之间的动摩擦系数 $\mu=0.20$，
求恒力 $F$ 的大小。简单起见取 $g=10\,\mathrm{m/s^2}$。

\begin{qualityfigure}
\begin{tikzpicture}[>=Latex,font=\small]
\draw[thick](-.5,0)--(6.8,0);\draw[fill=PaleBlue](.2,0)rectangle(1.2,.65);\draw[->,DeepBlue](1.3,.3)--(2.3,.3)node[above]{$F$};
\fill(.7,0)circle(1pt)node[below]{$A$};\fill(3.2,0)circle(2pt)node[below]{撤力点 $C$};\fill(6.3,0)circle(2pt)node[above]{$B$（停止）};
\draw[<->](.7,-.7)--(6.3,-.7)node[midway,below]{$x=5.0\,\mathrm m$};\node at(4.8,1){撤力后滑行 $2.0\,\mathrm s$};\node at(3,-1.6){全过程接触面粗糙，$\mu=0.20$};
\end{tikzpicture}
\figcap{受恒力段与撤力后的滑行段}\label{mech:06:15:q}
\end{qualityfigure}
\end{origexamplebox}

\Needspace{6\baselineskip}\begin{xesolution}题图~\ref{mech:06:15:q}中，撤力点 $C$ 把总位移 $AB$ 分成两段，$CB$ 段的已知停车时间可先确定 $C$ 点速度，随后再倒推 $AC$ 段。

按“先处理撤力后的运动，再倒推受力阶段”的顺序分析，运动分段如解图所示：
\soldiagTwoFifteen
撤力后减速度大小为
\[
 a_2=\mu g=2.0\,\mathrm{m/s^2}.
\]
物块在 $2.0\,\mathrm s$ 后停止，故撤力瞬间速度
\[
 v=a_2t=4.0\,\mathrm{m/s},
\]
撤力后位移
\[
 x_2=\frac{v+0}{2}t=4.0\,\mathrm m.
\]
于是受力阶段位移 $x_1=5.0-4.0=1.0\,\mathrm m$。由 $v^2=2a_1x_1$ 得
\[
 a_1=8.0\,\mathrm{m/s^2}.
\]
受力阶段列牛顿第二定律：
\[
 F-\mu mg=ma_1,
\]
所以
\[
 \boxed{F=m(a_1+\mu g)=15\,\mathrm N}.
\]
\end{xesolution}

\begin{origexamplebox}{}\label{ch06:ex:s05-2-16}
如图所示，质量分别为 $m_A$ 和 $m_B$ 的两木块 $A,B$ 静止放在粗糙水平面上，两者与地面之间
的动摩擦因数均为 $\mu$；本题同时取最大静摩擦因数为 $\mu$。两木块的接触面是倾角为
$\theta$ 的光滑斜面。现对 $A$ 施加水平推力 $F$，要使两块由静止开始向右加速，
且 $A,B$ 间不发生相对滑动、均保持与地面接触，$\mu$ 和 $F$ 各应满足什么条件？
\probdiagTwoSixteen
\end{origexamplebox}

\Needspace{6\baselineskip}\begin{xesolution}
以下推导所用的解图为图~\ref{mech:legacy:ch06:s05-2-16}。

设共同向右加速度为 $a$。分离 $B$，接触面法向力 $N$ 垂直于斜面并指向右下方，受力如
解图所示：
\soldiagTwoSixteen
竖直方向没有加速度，故地面对 $B$ 的支持力
\[
 N_{gB}=m_Bg+N\cos\theta.
\]
水平方向列式：
\[
 N\sin\theta-\mu N_{gB}=m_Ba.
\]
要使 $N>0$ 成为可能，必须先有
\[
 \boxed{\mu<\tan\theta}.
\]
对 $A+B$ 整体，地面总压力为 $(m_A+m_B)g$，所以
\[
 a=\frac{F-\mu(m_A+m_B)g}{m_A+m_B}.
\]
另一方面，由 $B$ 的方程解得
\[
 N=\frac{m_B(a+\mu g)}{\sin\theta-\mu\cos\theta}.
\]
对 $A$ 的竖直方向，地面支持力
\[
 N_{gA}=m_Ag-N\cos\theta\ge0.
\]
代入上式和整体加速度，得到不翘起的上限
\[
 \boxed{F\le
 \frac{m_A(m_A+m_B)g}{m_B}\bigl(\tan\theta-\mu\bigr)}.
\]
启动前，整体的水平静摩擦力至多为 $\mu(N_{gA}+N_{gB})=\mu(m_A+m_B)g$。
因此要排除静止平衡，并在滑动后获得向右的加速度，必须有
\[
 \boxed{F>\mu(m_A+m_B)g}.
\]
要使上面的 $F$ 区间非空，$\mu$ 还必须满足更强的条件
\[
 \boxed{\mu<\frac{m_A}{m_A+m_B}\tan\theta}.
\]
因此最终条件是这一 $\mu$ 限制与上述 $F$ 的上下界同时成立。上界等号对应 $N_{gA}=0$，
即 $A$ 刚好不离地的临界状态；若要求地面对 $A$ 有正压力，应将上界改为严格不等号。
若最大静摩擦因数另有给定，启动条件须另行检验，不能直接用动摩擦因数代替。
\end{xesolution}

\begin{origexamplebox}{}\label{ch06:ex:s05-2-17}
如图所示，质量分别为 $m_1,m_2$ 的两个木块重叠后放在光滑的水平面上。两者之间的动
摩擦因数和最大静摩擦因数均为 $\mu$，现给 $m_1$ 施加随时间增大的力 $F=kt$，式中 $k$
是正常数。两块初始相对静止；只研究两者仍保持接触的阶段。试求
$m_1$ 和 $m_2$ 的加速度 $a_1,a_2$ 与时间的关系并画出此关系的图像。
\probdiagTwoSeventeen
\end{origexamplebox}

\Needspace{6\baselineskip}\begin{xesolution}
以下推导所用的解图为图~\ref{mech:legacy:ch06:s05-2-17}。

尚未滑动时两块共同加速：
\[
 a_1=a_2=\frac{kt}{m_1+m_2}.
\]
带动下块所需静摩擦力
\[
 f=m_2a_2=\frac{m_2kt}{m_1+m_2}.
\]
令 $f=\mu m_1g$，得到临界时刻
\[
 \boxed{t_0=\frac{\mu m_1g(m_1+m_2)}{km_2}}.
\]
当 $t>t_0$ 时两块相对滑动，摩擦力大小为 $\mu m_1g$，于是
\[
 \boxed{
 a_1(t)=\frac{kt-\mu m_1g}{m_1},
 \qquad
 a_2(t)=\frac{\mu m_1g}{m_2}}
 \qquad(t>t_0).
\]
完整的分段图像如解图：
\soldiagTwoSeventeen
在 $t=t_0$ 处，由于题设取最大静摩擦因数和动摩擦因数相同，两条加速度曲线连续；若两者
不相等，开始滑动时加速度会发生跳变。
\end{xesolution}

根据第一定律，不受力的质点将做匀速直线运动；其直接推论是，做曲线运动的质点必然
受到外力。在运动学中，加速度沿速度方向的分量衡量速率的变化率，垂直于速度的分量则
决定轨迹的曲率半径。若 $\theta$ 是加速度与速度方向的夹角，则
\begin{align}
 a\cos\theta&=\frac{\dd v}{\dd t}, \label{ch06:eq:s05-2-10}\\
 \rho&=\frac{v^2}{a\sin\theta}. \label{ch06:eq:s05-2-11}
\end{align}
牛顿第二定律给出了加速度与合力的正比例关系，因此这些运动学结论可以直接用于动力学。

\begin{origexamplebox}{}\label{ch06:ex:s05-2-18}
质量为 $m$ 的质点 $A$ 固定在长为 $L$、不可变形且质量可以忽略不计的轻刚杆一端，杆的
另一端由 $O$ 处的驱动装置约束。驱动装置强制杆以匀角速度 $\omega$ 在竖直平面内转动，
并可通过刚杆向质点传递沿杆和垂直于杆的约束力。求 $A$ 对杆作用力的大小和方向。
\probdiagTwoEighteen
\end{origexamplebox}

\Needspace{6\baselineskip}\begin{xesolution}
以下推导所用的解图为图~\ref{mech:legacy:ch06:s05-2-18}。

设杆与竖直向上的夹角为 $\theta$，从竖直向上沿顺时针方向计角。
在图示右半圆，$\hat{\boldsymbol e}_r=(\sin\theta,\cos\theta)$ 沿杆向外，
$\hat{\boldsymbol e}_t=(\cos\theta,-\sin\theta)$ 沿 $\theta$ 增大的方向。
质点作匀速圆周运动，所需合力为指向 $O$ 的
$m\omega^2L$。设杆对质点的力为 $\vec R$，则
\[
 \vec R+m\vec g=-m\omega^2L\,\hat{\boldsymbol e}_r,
\]
其中 $\hat{\boldsymbol e}_r$ 从 $O$ 指向 $A$。力的矢量关系如解图：
\soldiagTwoEighteen
故
\[
 \vec R=-m\omega^2L\,\hat{\boldsymbol e}_r-m\vec g.
\]
重力的径向分量为 $-mg\cos\theta$、切向分量为 $mg\sin\theta$。
匀速圆周运动的切向加速度为零，所以
\[
 R_r-mg\cos\theta=-m\omega^2L,\qquad R_t+mg\sin\theta=0.
\]
因此
\[
 R_r=m(g\cos\theta-\omega^2L),
 \qquad
 R_t=-mg\sin\theta.
\]
因此杆对质点的力大小为
\[
 R=m\sqrt{\omega^4L^2+g^2-2\omega^2Lg\cos\theta}.
\]
由牛顿第三定律，质点 $A$ 对杆的作用力为 $-\vec R$，大小相同、方向相反：
\[
 \boxed{F_{A\to\text{杆}}
 =m\sqrt{\omega^4L^2+g^2-2\omega^2Lg\cos\theta}}.
\]
其沿杆向外分量为 $m(\omega^2L-g\cos\theta)$，切向分量为 $mg\sin\theta$。
例如最高点 $\theta=0$ 时力的大小为 $m|\omega^2L-g|$，最低点 $\theta=\pi$ 时为
$m(\omega^2L+g)$。驱动装置通过杆补偿重力的切向作用，才能使角速度保持不变。
\end{xesolution}

\begin{origexamplebox}{}\label{ch06:ex:s05-2-19}
水平的粗糙圆盘以恒定角速度绕竖直中心轴转动。盘上放有一个质量为 $m$ 的物体，它与圆盘之间的最大静摩擦因数为 $\mu$，与圆心的
距离为 $r$。求该物体能够与圆盘保持相对静止时，圆盘绕中心转动的最大角速度
$\omega_0$。

\begin{qualityfigure}
\begin{tikzpicture}[>=Latex,font=\small]
\draw[fill=PaleBlue](0,0)circle(2);\fill(0,0)circle(2pt)node[below left]{$O$};\fill(1.4,0)circle(3pt)node[right]{$m$};
\draw[dashed](0,0)--(1.4,0)node[midway,below]{$r$};\draw[->,Gold](0,1.4)arc(90:170:1.4)node[midway,left]{$\omega$};
\node at(0,-2.6){俯视图；最大静摩擦因数 $\mu$};
\end{tikzpicture}
\figcap{圆盘俯视中的物体半径和角速度}\label{mech:06:19:q}
\end{qualityfigure}
\end{origexamplebox}

\Needspace{6\baselineskip}\begin{xesolution}题图~\ref{mech:06:19:q}中，$r$ 是物体到轴心 $O$ 的距离；相对圆盘静止并不表示对地静止，所需水平向心力必须由静摩擦承担。

物体随圆盘作圆周运动，水平方向只有静摩擦力能够提供向心力，方向如解图：
\soldiagTwoNineteen
所需向心力为 $m\omega^2r$，而
\[
 f_s\le f_{s,\max}=\mu mg.
\]
故
\[
 m\omega^2r\le\mu mg,
 \qquad
 \boxed{\omega_0=\sqrt{\frac{\mu g}{r}}}.
\]
质量 $m$ 消去，说明在该理想模型中临界角速度与物体质量无关。
\end{xesolution}

当一个动力学系统包含多个物体时，各物体的加速度大小和方向都可能不同，约束对确定它们
之间的加速度关系至关重要。在运动学中可以用几何、代数或微元法求速度约束；在动力学中，
限制物体运动的连接或表面还会施加约束力。求解这类系统时必须同时写出运动约束与牛顿定律。

\begin{origexamplebox}{}\label{ch06:ex:s05-2-20}
一个质量和摩擦均不计的定滑轮两端通过不可伸长的绳连接，绳子的一端系有质量为 $m_1$
的重物，另一端有一名质量为 $m_2$ 的同学。现该同学用力拉绳子，使他相对于绳子以加速度
$a$ 向上运动；规定向上为 $a$ 的正方向。先在绳保持绷紧的阶段求重物的加速度和绳中张力，
再求绳刚好失去张力的极限状态所对应的 $a$，并说明其物理含义。

\begin{qualityfigure}
\begin{tikzpicture}[>=Latex,font=\small]
\draw[thick](0,2.6)circle(.5);\draw[thick](0,3.1)--(0,3.6);\draw[thick](-.5,2.6)--(-.5,0)(.5,2.6)--(.5,-.6);
\draw[fill=PaleBlue](-.9,-.7)rectangle(-.1,0);\node at(-.5,-.35){$m_1$};\fill(.9,.8)circle(.16);\draw(.9,.64)--(.9,-.15)(.9,.3)--(.5,.7)(.9,-.15)--(.6,-.7)(.9,-.15)--(1.2,-.7);\node[right]at(1,.3){$m_2$};
\draw[->,Gold](1.8,-.4)--(1.8,1.2)node[midway,right]{$a$（相对绳）};\node at(-2.1,1){轻绳、轻滑轮};
\end{tikzpicture}
\figcap{重物、同学及相对绳的加速度}\label{mech:06:20:q}
\end{qualityfigure}
\end{origexamplebox}

\Needspace{6\baselineskip}\begin{xesolution}题图~\ref{mech:06:20:q}中，$a$ 的箭头表示同学相对他这一侧绳的加速度，而非相对地面的加速度。绷紧轻绳两侧的加速度反向，所以应先建立相对加速度约束再列动力学方程。

统一取竖直向上为正。设重物和同学相对地面的加速度分别为 $a_1,a_2$。重物一侧绳端的
加速度为 $a_1$，因此同学一侧绳的加速度为 $-a_1$；同学相对绳向上加速度为 $a$，因而
\[
 \boxed{a_2=a-a_1}.
\]
约束和受力如解图所示：
\soldiagTwoTwenty
对重物和同学分别列式：
\[
 T-m_1g=m_1a_1,
 \qquad
 T-m_2g=m_2a_2.
\]
代入 $a_2=a-a_1$，解得
\[
 \boxed{a_1=\frac{(m_2-m_1)g+m_2a}{m_1+m_2}},
 \qquad
 \boxed{T=\frac{m_1m_2(2g+a)}{m_1+m_2}}.
\]
这里的正负号是相对于“向上为正”的约定，且绳保持绷紧要求 $T\ge0$。

若要求 $T=0$，上式给出 $a=-2g$。这里 $a$ 以“相对绳向上”为正，因此实际含义是同学
相对绳以 $2g$ 的加速度向下运动；此时同学和重物都作自由落体，绳处于松弛临界。更直接
地说，同学应放开绳子并自由下落，而不是继续通过绷紧的绳主动上爬。

当 $T=0$ 时绳约束已经退化，不能再把“相对绳的加速度”当作由绷紧绳
唯一确定的运动量；因此最后一问应按自由落体的物理图景理解。
\end{xesolution}

\begin{origexamplebox}{}\label{ch06:ex:s05-2-21}
如图所示，一个质量为 $M_1$ 的物体挂在质量和摩擦均不计的定滑轮 $A$ 的一侧，另一侧
挂有一个质量为 $M_2$ 的动滑轮组件；它的两侧分别挂有质量为 $m_1,m_2$ 的两个物体。
两根绕轮绳轻且不可伸长，动滑轮轮缘光滑，绳可无摩擦地滑过轮缘，不计轮体转动的影响。
图中 $T$ 表示定滑轮上方悬挂线的张力，$T_1$ 表示下部绕轮绳的张力；上部绕轮绳的张力另记为 $H$。
求 $m_1,m_2$ 的加速度及 $T,T_1$，并给出 $H$。
\probdiagTwoTwentyOne
\end{origexamplebox}

\Needspace{6\baselineskip}\begin{xesolution}
以下推导所用的解图为图~\ref{mech:legacy:ch06:s05-2-21}。

统一取竖直向下为正，设动滑轮 $M_2$ 的加速度为 $a_p$，两个小物体的加速度为
$a_1,a_2$。两根轻绳分别给出
\[
 a_{M_1}=-a_p,
 \qquad
 a_1+a_2=2a_p,
\]
如解图所示：
\soldiagTwoTwentyOne
四个动力学方程为
\begin{align*}
 M_1g-H&=-M_1a_p,\\
 m_1g-T_1&=m_1a_1,\\
 m_2g-T_1&=m_2a_2,\\
 M_2g+2T_1-H&=M_2a_p.
\end{align*}
定滑轮本身质量可忽略、加速度为零，受到两段竖直绳向下各拉 $H$，所以
\[
 \boxed{T=2H}.
\]
不能将支承处的 $T$ 与绕轮绳的 $H$ 视为同一个力。下面逐步解出各量。
两个小物体的方程给出
\[
 a_1=g-\frac{T_1}{m_1},\qquad a_2=g-\frac{T_1}{m_2}.
\]
代入 $a_1+a_2=2a_p$，得
\[
 T_1=\frac{2m_1m_2}{m_1+m_2}(g-a_p).
\]
再用 $H=M_1(g+a_p)$ 代入动滑轮方程，整理为
\[
 \left(M_1+M_2+\frac{4m_1m_2}{m_1+m_2}\right)a_p
 =\left(M_2-M_1+\frac{4m_1m_2}{m_1+m_2}\right)g.
\]
记
\[
 D=(M_1+M_2)(m_1+m_2)+4m_1m_2.
\]
因此 $a_p=g[(M_2-M_1)(m_1+m_2)+4m_1m_2]/D$。回代到两个小物体的方程，得
\[
 \boxed{a_1=
 \frac{g\,[M_1m_1-3M_1m_2+M_2(m_1+m_2)+4m_1m_2]}{D}},
\]
\[
 \boxed{a_2=
 \frac{g\,[-3M_1m_1+M_1m_2+M_2(m_1+m_2)+4m_1m_2]}{D}},
\]
以及
\[
 \boxed{H=\frac{2M_1g\,[M_2(m_1+m_2)+4m_1m_2]}{D}},
\]
\[
 \boxed{T=\frac{4M_1g\,[M_2(m_1+m_2)+4m_1m_2]}{D}},\qquad
 \boxed{T_1=\frac{4M_1m_1m_2g}{D}}.
\]
加速度若为负，表示实际方向与预设的“向下”相反。
各张力均为质量乘重力加速度的量纲。若 $m_1=m_2=m$ 且 $M_1=M_2+2m$，
代回得 $a_p=a_1=a_2=0$，此时 $T_1=mg$、$H=M_1g$、$T=2M_1g$，
与整个装置的静力平衡相符。这里若改成有转动惯量且绳不打滑的轮，须分别设置下绳两侧张力并增加转动方程。
\end{xesolution}

\begin{origexamplebox}{}\label{ch06:ex:s05-2-22}
一质量为 $m$ 的物体放在质量为 $M$、倾角为 $\theta$ 的斜面上，两者一起放置于水平面
上。忽略一切摩擦，求两者的加速度。
\probdiagTwoTwentyTwo
\end{origexamplebox}

\Needspace{6\baselineskip}\begin{xesolution}
以下推导所用的解图为图~\ref{mech:legacy:ch06:s05-2-22}。

设斜面向右位移为 $X$，物块沿斜面向下位移为 $s$。在地面惯性系中
\[
 x_m=X-s\cos\theta,
 \qquad
 y_m=-s\sin\theta.
\]
因此小物块加速度及受力关系如解图：
\soldiagTwoTwentyTwo
沿斜面向下的单位矢量为 $\hat{\boldsymbol e}_s=(-\cos\theta,-\sin\theta)$，水平方向记为 $\hat{\boldsymbol e}_x=(1,0)$。
由坐标式求两次导数，物块对地加速度为
\[
 \vec a_m=(\ddot X-\ddot s\cos\theta,\,-\ddot s\sin\theta).
\]
它在 $\hat{\boldsymbol e}_s$ 上的投影是
\[
 \vec a_m\cdot\hat{\boldsymbol e}_s
 =-\ddot X\cos\theta+\ddot s(\cos^2\theta+\sin^2\theta)
 =\ddot s-\ddot X\cos\theta.
\]
支持力垂直斜面，投影为零；重力沿斜面的分量为 $mg\sin\theta$。所以
\[
 \ddot s-\ddot X\cos\theta=g\sin\theta.
\]
系统水平方向无外力，故
\[
 (M+m)\ddot X-m\ddot s\cos\theta=0.
\]
联立得到斜面向右的水平加速度
\[
 \boxed{a_M=\ddot X
 =\frac{mg\sin\theta\cos\theta}
 {M+m\sin^2\theta}}
\]
物块相对斜面沿斜面向下的加速度为
\[
 \boxed{a_{m/M}=\ddot s
 =\frac{(M+m)g\sin\theta}
 {M+m\sin^2\theta}}.
\]
若需要物块对地加速度，可写为
\[
\boxed{a_{mx}=-\frac{Mg\sin\theta\cos\theta}{M+m\sin^2\theta},
 \qquad
 a_{my}=-\frac{(M+m)g\sin^2\theta}{M+m\sin^2\theta}}.
\]
图中虚线表示相对加速度和斜面加速度，蓝色箭头表示两者相加后的对地加速度；图示取 $M=m$。
对任意有限的正质量，$|a_{my}/a_{mx}|=(1+m/M)\tan\theta$，因此对地加速度比斜面向下方向更陡。
当 $M\to\infty$ 时，$a_M\to0$、$\ddot s\to g\sin\theta$，还原为固定斜面。
\end{xesolution}

\begin{origexamplebox}{}\label{ch06:ex:s05-2-23}
一根质量为 $m$ 的杆被限制在竖直方向运动，它与一个倾角为 $\theta$、质量为 $M$ 的
斜面接触，接触点的摩擦可忽略不计；不考虑斜面与水平面的摩擦。求使斜面以加速度 $a$
向右运动所需要的水平外力 $F$ 的大小。
\probdiagTwoTwentyThree
\end{origexamplebox}

\Needspace{6\baselineskip}\begin{xesolution}
以下推导所用的解图为图~\ref{mech:legacy:ch06:s05-2-23}。

几何约束决定：斜面向右移动 $x$ 时，固定在水平位置的杆接触高度增加 $x\tan\theta$，
所以杆的竖直加速度为
\[
 a_y=a\tan\theta.
\]
杆与斜面的受力如解图所示。设接触面法向力大小为 $N$，分别对杆沿竖直方向、对斜面沿
水平方向列式：
\soldiagTwoTwentyThree
\begin{align*}
 N\cos\theta-mg&=ma\tan\theta,\\
 F-N\sin\theta&=Ma.
\end{align*}
由第一式得 $N\cos\theta=m(g+a\tan\theta)$，代入第二式，得到
\[
 \boxed{F=Ma+mg\tan\theta+ma\tan^2\theta}.
\]
图中导轨对杆的横向反力 $R_h=N\sin\theta$ 向左，使杆的水平加速度为零；它没有竖直分量，
所以不进入杆的竖直方程。导轨须光滑且固定于地面，不能把该反力漏掉后误认为杆会水平加速。
\end{xesolution}
}

{
\graphicspath{{../../../sample/xes-original/sixth-section/}}
\section{非惯性系与惯性力}
% Original problem figures and added solution diagrams for Section 2.3.
\tikzset{
  solbody/.style={draw=Ink,line width=0.8pt},
  solforce/.style={-{Latex[length=2.0mm]},draw=DeepBlue,line width=0.9pt},
  solgeom/.style={draw=TextGray,densely dashed,line width=0.55pt},
  sollabel/.style={font=\footnotesize,fill=none,inner sep=0.7pt,outer sep=1.2pt,text=Ink},
  solpanel/.style={font=\scriptsize\bfseries,text=TextGray}
}
\newcommand{\diagcaption}[2]{\figcap{#1：#2}}

\newcommand{\probdiagTwoTwentyFour}{%
\begin{center}\begin{minipage}{\linewidth}\centering\begin{tikzpicture}[scale=.86]
 \draw[solbody] (-2.7,-1.1) rectangle (2.5,1.7);
 \draw[solbody] (-1.75,-1.45) circle(.35) (1.45,-1.45) circle(.35);
 \coordinate (O) at (.4,1.7); \coordinate (B) at ($(O)+(-115:2)$);
 \draw[solbody] (O)--(B); \fill (B) circle(3pt);
 \node[sollabel,left=5pt] at (B) {$m$};
 \node[sollabel,above=3pt] at (O) {$O$};
 \draw[solgeom] (O)--++(0,-2.3);
 \draw[draw=TextGray] ($(O)+(-90:1)$) arc[start angle=-90,end angle=-115,radius=1];
 \node[sollabel] at ($(O)+(-102.5:1.4)$) {$\theta_1$};
 \draw[solforce] (1,2.35)--(2.2,2.35) node[right,sollabel] {$a$};
\end{tikzpicture}\bookcaption{车向右匀加速；平衡时细绳与竖直线夹角为 $\theta_1$。}\label{ch06:fig:s06-2-24}\end{minipage}\end{center}}

\newcommand{\soldiagTwoTwentyFour}{%
\begin{center}\begin{minipage}{\linewidth}\centering\begin{tikzpicture}[scale=.9]
 \node[solpanel] at (-2.4,2.7) {平衡时：三力矢量和为零};
 \fill (-2.4,.4) circle(2pt);
 \draw[solforce] (-2.4,.4)--++(0,-1.5) node[below,sollabel] {$mg$};
 \draw[solforce] (-2.4,.4)--++({-1.5*tan(25)},0) node[left,sollabel] {$ma$};
 \draw[solforce] (-2.4,.4)--++({1.5*tan(25)},1.5) node[above right,sollabel] {$T$};
 \node[solpanel] at (2.3,2.7) {从竖直位置静止释放};
 \coordinate (O) at (3.5,1.9);
 \coordinate (Bz) at ($(O)+(-90:2.6)$);
 \coordinate (Be) at ($(O)+(-115:2.6)$);
 \coordinate (Bt) at ($(O)+(-140:2.6)$);
 \draw[solgeom] (O)--(Bz);
 \draw[solgeom] (O)--(Be);
 \draw[solbody] (O)--(Bt);
 \draw[solgeom] (Bz) arc[start angle=-90,end angle=-140,radius=2.6];
 \foreach \P/\lab/\anc in {Bz/B_0/below right,Be/B_e/below,Bt/B_2/left}{
  \fill (\P) circle(2pt) node[\anc=4pt,sollabel] {$\lab$};
 }
 \node[sollabel,above right=3pt] at (O) {$O$};
 \draw[draw=TextGray] ($(O)+(-90:.72)$) arc[start angle=-90,end angle=-115,radius=.72];
 \node[sollabel] at ($(O)+(-102.5:1.02)$) {$\theta_1$};
 \draw[draw=TextGray] ($(O)+(-90:1.6)$) arc[start angle=-90,end angle=-140,radius=1.6];
 \node[sollabel] at ($(O)+(-122:1.95)$) {$\theta_2$};
 \draw[solforce] (4.65,1.2)--++(-115:1.25) node[right=5pt,sollabel] {$\vec g_{\rm eff}$};
\end{tikzpicture}\diagcaption{解图}{$B_0$ 为释放点，$B_e$ 为平衡点，$B_2$ 为另一端点；$OB_e$ 平行于有效重力，$\theta_2=2\theta_1$。}\label{mech:legacy:ch06:s06-2-24}\end{minipage}\end{center}}

\newcommand{\probdiagTwoTwentyFive}{%
\begin{center}\begin{minipage}{\linewidth}\centering\begin{tikzpicture}[scale=.76]
  \draw[solbody] (-1.35,-1.55) rectangle (1.35,1.55);
  \fill[Ink] (0,.45) circle (3.5pt);
  \draw[solforce] (1.85,.45)--++(0,-1.25) node[below,sollabel] {$g$};
\end{tikzpicture}\bookcaption{电梯与内部物体同时自由下落。}\label{ch06:fig:s06-2-25}\end{minipage}\end{center}}

\newcommand{\soldiagTwoTwentyFive}{%
\begin{center}\begin{minipage}{\linewidth}\centering\begin{tikzpicture}[scale=.78]
  \fill[Ink] (0,0) circle (3pt);
  \draw[solforce] (0,0)--++(0,-1.25) node[below,sollabel] {$mg$};
  \draw[solforce] (0,0)--++(0,1.25) node[above,sollabel] {$F_i=mg$};
  \node[sollabel] at (2.2,.2) {$\vec g_{\rm eff}=\vec g-\vec g=0$};
\end{tikzpicture}\diagcaption{解图}{在自由落体电梯系中，惯性力与重力抵消，局部呈现失重。}\label{mech:legacy:ch06:s06-2-25}\end{minipage}\end{center}}

\newcommand{\probdiagTwoTwentySix}{%
\begin{center}\begin{minipage}{\linewidth}\centering\begin{tikzpicture}[scale=.87]
 \def\th{22}
 \draw[solbody] (-2.8,-1.2) rectangle (2.8,1.8);
 \draw[solbody] (-1.85,-1.55) circle(.35) (1.65,-1.55) circle(.35);
 \coordinate (L) at (-1.9,-.85);
 \coordinate (R) at (2.3,-.85);
 \coordinate (P) at ($(L)+(4.2,{4.2*tan(\th)})$);
 \coordinate (S) at ($(L)+(\th:.45)$);
 \coordinate (B) at ($(S)+({90+\th}:.38)$);
 \draw[solbody] (L)--(R)--(P)--cycle;
 \begin{scope}[shift={(B)},rotate=\th]\draw[solbody,fill=white] (-.45,-.38) rectangle (.45,.38);\end{scope}
 \node[sollabel] at (B) {$m$};
 \draw (L)++(1.3,0) arc[start angle=0,end angle=\th,radius=1.3];
 \node[sollabel] at ($(L)+(11:1.68)$) {$\theta$};
 \draw[solforce] (-2.9,.5)--(-4,.5) node[left,sollabel] {$a$};
\end{tikzpicture}\bookcaption{汽车向左加速，物块从斜面底端释放；物块底面与斜面贴合。}\label{ch06:fig:s06-2-26}\end{minipage}\end{center}}

\newcommand{\soldiagTwoTwentySix}{%
\begin{center}\begin{minipage}{\linewidth}\centering\begin{tikzpicture}[scale=.95]
 \draw[solgeom] (202:2)--(22:2.4);
 \fill (0,0) circle(2pt);
 \draw[solforce] (0,0)--(1.6,0) node[right,sollabel] {$ma$};
 \draw[solforce] (0,0)--(0,-1.5) node[below,sollabel] {$mg$};
 \draw[solforce] (0,0)--(112:1.5) node[above left,sollabel] {$N$};
 \draw[solgeom,->] ($(0,0)+(112:.7)$)--++(22:1.7) node[above right=3pt,sollabel] {$a_{\parallel}$};
\end{tikzpicture}\diagcaption{解图}{$N$ 垂直斜面；沿斜面向上取正，$a_{\parallel}=a\cos\theta-g\sin\theta$。}\label{mech:legacy:ch06:s06-2-26}\end{minipage}\end{center}}

\newcommand{\probdiagTwoTwentySeven}{%
\begin{center}\begin{minipage}{\linewidth}\centering\begin{tikzpicture}[scale=.85]
 \def\al{26}
 \coordinate (L) at (-2.5,-1);
 \coordinate (R) at (2.5,-1);
 \coordinate (P) at ($(L)+(0,{5*tan(\al)})$);
 \coordinate (S) at ($(P)!.56!(R)$);
 \coordinate (B) at ($(S)+({90-\al}:.42)$);
 \draw[solbody] (L)--(R)--(P)--cycle;
 \begin{scope}[shift={(B)},rotate=-\al]\draw[solbody,fill=white] (-.43,-.42) rectangle (.43,.42);\end{scope}
 \node[sollabel] at (B) {$m$};
 \draw[solforce] ($(B)+({90-\al}:.42)$)--++({90-\al}:1) node[above right,sollabel] {$N$};
 \draw (R)++({180-\al}:.8) arc[start angle={180-\al},end angle=180,radius=.8];
 \node[sollabel] at ($(R)+(167:1.12)$) {$\alpha$};
 \draw[solforce] (-3.7,0)--(-2.55,0) node[midway,above,sollabel] {$F$};
\end{tikzpicture}\bookcaption{水平力推动斜面向右；支持力 $N$ 垂直于接触面。}\label{ch06:fig:s06-2-27}\end{minipage}\end{center}}

\newcommand{\soldiagTwoTwentySeven}{%
\begin{center}\begin{minipage}{\linewidth}\centering\begin{tikzpicture}[scale=1]
 \draw[solgeom] (154:1.8)--(-26:1.8);
 \fill (0,0) circle(2pt);
 \draw[solforce] (0,0)--(0,-1.5) node[below,sollabel] {$mg$};
 \draw[solforce] (0,0)--({-1.5*tan(26)},0) node[left,sollabel] {$ma$};
 \draw[solforce] (0,0)--({1.5*tan(26)},1.5) node[above right,sollabel] {$N$};
\end{tikzpicture}\diagcaption{解图}{虚线与题图斜面同向；平衡时 $N$ 正好抵消重力和惯性力的合力，并与斜面垂直。}\label{mech:legacy:ch06:s06-2-27}\end{minipage}\end{center}}

\newcommand{\probdiagTwoTwentyEight}{%
\begin{center}\begin{minipage}{\linewidth}\centering\begin{tikzpicture}[scale=.72]
  % The two slopes deliberately have visibly different inclinations, as in the
  % source figure: alpha is the gentler left slope and beta the steeper right one.
  \coordinate (L) at (-3.05,-1.05);
  \coordinate (P) at (.65,1.55);
  \coordinate (R) at (2.15,-1.05);
  \draw[solbody,fill=SoftGray] (L)--(P)--(R)--cycle;
  \draw[solbody] (-3.65,-1.05)--(3.25,-1.05);
  \foreach \x in {-3.55,-3.20,...,3.15}
    \draw[solgeom] (\x,-1.05)--++(-.16,-.18);

  % Each block is drawn in local slope coordinates; y=0 is its contact face.
  \coordinate (SL) at ($(L)!.48!(P)$);
  \pgfmathsetmacro{\leftang}{atan2(2.6,3.7)}
  \pgfmathsetmacro{\rightang}{-atan2(2.6,1.5)}
  \begin{scope}[shift={(SL)},rotate=\leftang]
    \draw[solbody,fill=white] (-.43,0) rectangle (.43,.72);
    \node[sollabel] at (0,.36) {$m_1$};
    \node[sollabel] at (0,.98) {$A$};
  \end{scope}
  \coordinate (SR) at ($(P)!.58!(R)$);
  \begin{scope}[shift={(SR)},rotate=\rightang]
    \draw[solbody,fill=white] (-.43,0) rectangle (.43,.72);
    \node[sollabel] at (0,.36) {$m_2$};
    \node[sollabel] at (0,.98) {$B$};
  \end{scope}

  \draw[solgeom] ([xshift=.43cm]L) arc[start angle=0,end angle=35.1,radius=.43];
  \draw[solgeom] ([shift={(120:.40)}]R) arc[start angle=120,end angle=180,radius=.40];
  \node[sollabel] at (-2.45,-.78) {$\alpha$};
  \node[sollabel] at (1.68,-.73) {$\beta$};
  \node[sollabel] at (-.05,-.25) {$M$};
\end{tikzpicture}\bookcaption{双斜面在光滑水平面上自由运动。}\label{ch06:fig:s06-2-28}\end{minipage}\end{center}}

\newcommand{\soldiagTwoTwentyEight}{%
\begin{center}\begin{minipage}{\linewidth}\centering\begin{tikzpicture}[scale=.75]
  \coordinate (L) at (-3.05,-.85);
  \coordinate (P) at (.65,1.75);
  \coordinate (R) at (2.15,-.85);
  \draw[solgeom] (L)--(P)--(R)--cycle;
  \draw[axis] (-3.60,-.85)--(3.45,-.85) node[right,sollabel] {$x$};

  \coordinate (SL) at ($(L)!.48!(P)$);
  \pgfmathsetmacro{\leftang}{atan2(2.6,3.7)}
  \pgfmathsetmacro{\rightang}{-atan2(2.6,1.5)}
  \begin{scope}[shift={(SL)},rotate=\leftang]
    \draw[solbody,fill=white] (-.35,0) rectangle (.35,.58);
    \node[sollabel] at (0,.29) {$m_1$};
    \draw[solforce] (-.10,1.02)--++(-1.15,0);
  \end{scope}
  \coordinate (SR) at ($(P)!.58!(R)$);
  \begin{scope}[shift={(SR)},rotate=\rightang]
    \draw[solbody,fill=white] (-.35,0) rectangle (.35,.58);
    \node[sollabel] at (0,.29) {$m_2$};
    \draw[solforce] (.10,1.02)--++(1.05,0);
  \end{scope}

  \node[sollabel] at (-2.5,1.65) {$s_1\ (+)$};
  \node[sollabel] at (3.12,1.35) {$s_2\ (+)$};
  \draw[solforce] (-.30,-1.22)--++(1.25,0)
    node[right,sollabel] {$\ddot X=A\ (+)$};
  \node[sollabel] at (-2.45,-.57) {$\alpha$};
  \node[sollabel] at (1.68,-.52) {$\beta$};
\end{tikzpicture}\diagcaption{解图}{取 $x$ 轴向右；$s_1,s_2$ 分别沿左右斜面向下为正，斜面加速度 $A=\ddot X$ 向右为正。}\label{mech:legacy:ch06:s06-2-28}\end{minipage}\end{center}}

\newcommand{\probdiagTwoTwentyNine}{%
\begin{center}\begin{minipage}{\linewidth}\centering\begin{tikzpicture}[scale=.76]
  \fill[TextGray!45] (-2.3,-1.0) rectangle (2.3,1.0); \draw[solbody] (-2.3,-1.0) rectangle (2.3,1.0);
  \fill[white,draw=Ink] (0,.15) circle (.22);
  \draw[solforce] (2.55,0)--++(1.0,0) node[right,sollabel] {$a$};
\end{tikzpicture}\bookcaption{装满水的水箱向前加速，水中有气泡。}\label{ch06:fig:s06-2-29}\end{minipage}\end{center}}

\newcommand{\soldiagTwoTwentyNine}{%
\begin{center}\begin{minipage}{\linewidth}\centering\begin{tikzpicture}[scale=.95]
 \foreach \b in {-.8,0,.8}{\draw[solgeom] (-2.1,{1.68+\b})--(2.1,{-1.68+\b});}
 \fill[white,draw=Ink] (0,0) circle(.14);
 \draw[solforce] (0,0)--(-1.12,-1.4) node[below left=4pt,sollabel] {$\vec g_{\rm eff}=(-a,-g)$};
 \draw[solforce] (0,0)--(1.12,1.4) node[above right=4pt,sollabel] {浮力};
 \node[sollabel,anchor=west] at (2.35,-1.1) {等压线};
 \draw[draw=TextGray] (2.3,-1.2)--(1.9,-1.52);
\end{tikzpicture}\diagcaption{解图}{等压线沿 $(g,-a)$ 方向，与有效重力 $(-a,-g)$ 垂直；压力向左下方增大，浮力指向右上方。}\label{mech:legacy:ch06:s06-2-29}\end{minipage}\end{center}}

% 题图保留杆、转轴、质点与给定量；半径投影及圆轨迹在解图中构造。
\newcommand{\probdiagTwoThirty}{%
\begin{center}\begin{minipage}{\linewidth}\centering\begin{tikzpicture}[scale=.95]
 \def\th{35}
 \coordinate (O) at (0,-1.1);
 \coordinate (P) at ($(O)+({90-\th}:2.6)$);
 \draw[solgeom] (O)--(0,2.6) node[above,sollabel] {转轴};
 \draw[solbody] (O)--++({90-\th}:3.6);
 \fill (O) circle(1.5pt) node[below left,sollabel] {$O$};
 \fill (P) circle(2.5pt) node[right=4pt,sollabel] {$P$};
 \draw[draw=TextGray] ($(O)+(0,.65)$) arc[start angle=90,end angle={90-\th},radius=.65];
 \node[sollabel] at ($(O)+(72.5:.94)$) {$\theta$};
 \draw[<->,draw=TextGray] ($(O)+(-35:.4)$)--($(P)+(-35:.4)$) node[midway,below right=4pt,sollabel] {$l$};
 \draw[->,draw=TextGray] (-.65,2.1) arc[start angle=180,end angle=350,x radius=.65,y radius=.17];
 \node[sollabel] at (.9,2.2) {$\omega$};
\end{tikzpicture}\bookcaption{质点在杆上距固定端 $O$ 为 $l$ 处，杆与竖直转轴夹角为 $\theta$。}\label{ch06:fig:s06-2-30}\end{minipage}\end{center}}


\newcommand{\soldiagTwoThirtyGeometry}{%
\begin{center}\begin{minipage}{\linewidth}\centering\begin{tikzpicture}[scale=.95]
 \def\th{35}
 \coordinate (O) at (0,-1.1);
 \coordinate (P) at ($(O)+({90-\th}:2.6)$);
 \coordinate (H) at (O |- P);
 \pgfmathsetmacro{\rad}{2.6*sin(\th)}
 \draw[solgeom] ($(H)+(-\rad,0)$) arc[start angle=180,end angle=0,x radius=\rad,y radius=.3];
 \draw[solbody] ($(H)+(-\rad,0)$) arc[start angle=180,end angle=360,x radius=\rad,y radius=.3];
 \draw[solgeom] (O)--(0,2.6) node[above,sollabel] {转轴};
 \draw[solbody] (O)--++({90-\th}:3.6);
 \fill (O) circle(1.5pt) node[below left,sollabel] {$O$};
 \fill (P) circle(2.5pt) node[right=4pt,sollabel] {$P$};
 \draw[solgeom] (H)--(P) node[midway,above=5pt,sollabel] {$r$};
 \draw[draw=TextGray] ($(O)+(0,.65)$) arc[start angle=90,end angle={90-\th},radius=.65];
 \node[sollabel] at ($(O)+(72.5:.94)$) {$\theta$};
 \draw[<->,draw=TextGray] ($(O)+(-35:.4)$)--($(P)+(-35:.4)$) node[midway,below right=4pt,sollabel] {$l$};
 \draw[->,draw=TextGray] (-.65,2.1) arc[start angle=180,end angle=350,x radius=.65,y radius=.17];
 \node[sollabel] at (.9,2.2) {$\omega$};
\end{tikzpicture}\diagcaption{解图}{$OP=l$，水平转动半径 $r=l\sin\theta$；圆轨道后半部以虚线表示。}\label{mech:legacy:ch06:s06-2-30}\end{minipage}\end{center}}

\newcommand{\soldiagTwoThirty}{%
\begin{center}\begin{minipage}{\linewidth}\centering\begin{tikzpicture}[scale=.95]
 \node[solpanel] at (-2.3,2.3) {旋转系中的受力};
 \coordinate (P) at (-2.3,0);
 \draw[solgeom] ($(P)+(235:1.8)$)--($(P)+(55:1.8)$);
 \fill (P) circle(2pt);
 \draw[solforce] (P)--++(1.5,0) node[right,sollabel] {$F_{\rm cf}$};
 \draw[solforce] (P)--++(0,-1.5) node[below,sollabel] {$mg$};
 \draw[solforce] (P)--++(145:1.35) node[above left,sollabel] {$N$};
 \draw[<->,draw=TextGray] ($(P)+(-35:.8)+(235:.4)$)--++(55:.8) node[midway,below right=6pt,sollabel] {$f_s$};
 \node[solpanel] at (2.9,2.3) {沿杆方向的投影};
 \coordinate (Q) at (2.9,0);
 \draw[solgeom] ($(Q)+(235:1.9)$)--($(Q)+(55:1.9)$);
 \fill (Q) circle(2pt);
 \draw[solforce] (Q)--++(55:1.3) node[pos=.65,right=8pt,sollabel] {$F_{{\rm cf},\parallel}$};
 \draw[solforce] (Q)--++(235:1.3) node[pos=.7,below right=7pt,sollabel] {$mg\cos\theta$};
\end{tikzpicture}\diagcaption{解图}{两图的杆方向相同；$N$ 垂直杆，$F_{{\rm cf},\parallel}=m\omega^2l\sin^2\theta$，$f_s$ 沿杆方向并可取正或负。}\end{minipage}\end{center}}

\newcommand{\soldiagTwoThirtyOne}{%
\begin{center}\begin{minipage}{\linewidth}\centering\begin{tikzpicture}[scale=.86]
  \draw[solbody] (0,0) circle (2.05);
  \fill[Ink] (0,0) circle (1.5pt) node[below left=3pt,sollabel] {转轴};
  \fill[DeepBlue] (2.05,0) circle (3.5pt) node[right=5pt,sollabel] {实验对象};
  \draw[solgeom] (0,0)--(2.05,0)
    node[pos=.48,below=4pt,sollabel] {$r=7\,\mathrm m$};
  \draw[solforce] (2.05,0)--(.62,0)
    node[pos=.58,above=4pt,sollabel] {$a_c=5g$};
  \draw[solgeom,-{Latex[length=1.7mm]}] (2.05,0)--++(0,1.05)
    node[above=3pt,sollabel] {$v$};
\end{tikzpicture}\diagcaption{解图}{离心机提供大小为 $5g$ 的水平向心加速度。}\end{minipage}\end{center}}


前面曾多次强调，牛顿定律只在惯性参考系中成立。于是自然会产生一个问题：相对于惯性系
做加速运动的非惯性参考系中，物体的运动是否就不能用类似的方式描述？严格地说，牛顿
定律不能直接用于非惯性系；但这并不妨碍我们先在惯性系中写出运动方程，再利用已知的
坐标变换，把它改写成相对于非惯性系的形式。这样做既保留了牛顿定律的适用条件，也能
方便地研究车厢、电梯和转盘中的运动。

考虑两个相对运动的参考系 $S$ 和 $S'$。其中 $S$ 是惯性参考系，$S'$ 相对于 $S$ 以
加速度 $\vec a_r$ 作平动。质量为 $m$ 的质点受到真实合力 $\vec F$，相对于 $S$ 的加速度
为 $\vec a$，相对于 $S'$ 的加速度为
\[
 \vec a'=\vec a-\vec a_r.
\]
由 $\vec F=m\vec a$ 得
\[
 m\vec a'=\vec F-m\vec a_r.
\]
因此，如果希望在 $S'$ 中仍采用“质量乘加速度等于力之和”的形式，就要引入惯性力
\begin{equation}
 \boxed{\vec F_i=-m\vec a_r}. \label{ch06:eq:s06-2-12}
\end{equation}
惯性力是参考系加速带来的表观力，没有施力物体，牛顿第三定律不适用于它。对于相对
非惯性系静止的物体，引入惯性力后，动力学问题常可转化为静力平衡问题。

上述结论并不限于单个质点。若非惯性系中有多个彼此作用的质点，就对每个质量为 $m_i$
的质点同时加入 $-m_i\vec a_r$。真实外力和质点间相互作用力仍按通常方式列出，惯性力
只负责补偿参考系的加速运动。于是每个质点都满足
\[
 m_i\vec a_i'=\vec F_{\mathrm{real},i}+\vec F_{\mathrm{inert},i},
 \qquad \vec F_{\mathrm{inert},i}=-m_i\vec a_r.
\]
这里 $\vec F_{\mathrm{real},i}$ 是第 $i$ 个质点所受真实合力，与惯性力分开记号。
由于所有质点所加惯性力都与各自质量成正比，许多相对车厢或电梯保持静止的问题便可按
平衡问题处理；但必须记住，这只是非惯性坐标中的等效写法，并没有产生新的相互作用。

\subsubsection*{平动非惯性系与等效重力场}

设参考系 $S'$ 相对惯性系作平动，原点加速度为 $\vec a_0$。若两系坐标轴不转动，则
$\vec a=\vec a'+\vec a_0$，于是
\begin{equation}
 m\vec a'=\sum\vec F-m\vec a_0.
\end{equation}
把 $-m\vec a_0$ 称为惯性力，便可在车厢、电梯等非惯性系中沿用平衡和动力学方程。惯性力
与质量成正比，没有施力物体，也没有第三定律的反作用力；它只是坐标系加速在运动方程中的
表现。若所有真实物体都同时受均匀重力 $m\vec g$，可把两者合并为
\begin{equation}
 \vec g_{\rm eff}=\vec g-\vec a_0.
\end{equation}
悬线、自由液面和小振动平衡位置都由有效重力方向决定。

这种局部等效有明确限度。自由落体电梯中，均匀重力与惯性力相消，舱内物体近似失重；但地球
重力随位置改变，不同点的自由落体加速度略有差异，仍会出现潮汐效应。因此自由落体系只能在
足够小的空间和足够短的时间内近似消除均匀重力对自由物体相对运动的影响。
按牛顿力学的定义，它仍是相对地面加速的非惯性系，不能仅因自由物体相对匀速就改变这个判定。
广义相对论中自由落体与局部惯性坐标的联系属于另一理论层次，这里不把两种定义混用。
判断均匀场近似是否可用，应比较被忽略的差异加速度与题目
要求的精度，而不能仅凭“物体正在自由下落”下结论。

若参考系还在转动，需要区分相对位置和绝对位置。设地面惯性系为 $S$，转动系为 $S'$，
$S'$ 原点的位置、速度、加速度分别为 $\vec R_0,\vec V_0,\vec a_0$；
$\vec r$ 是从 $S'$ 原点指向质点的位置矢量，$\vec v'$、$\vec a'$ 是在 $S'$ 中求得的相对速度、相对加速度。
转动系角速度为 $\vec\omega$，角加速度为 $\dot{\vec\omega}$。
对任意矢量 $\vec b$，两种求导满足
\[
 \left(\frac{\dd\vec b}{\dd t}\right)_S
 =\left(\frac{\dd\vec b}{\dd t}\right)_{S'}+\vec\omega\times\vec b.
\]
由绝对位置 $\vec R_0+\vec r$ 求导，得
\[
 \vec v=\vec V_0+\vec v'+\vec\omega\times\vec r.
\]
再求一次导数：$\vec v'$ 产生 $\vec a'+\vec\omega\times\vec v'$，而最后一项产生
$\dot{\vec\omega}\times\vec r+\vec\omega\times(\vec v'+\vec\omega\times\vec r)$。
两处相同的交叉项相加，得到
\begin{equation}
 \vec a=\vec a_0+\vec a'+2\vec\omega\times\vec v'
 +\vec\omega\times(\vec\omega\times\vec r)
 +\dot{\vec\omega}\times\vec r.
\end{equation}
五项依次为原点平动加速度、相对加速度、科里奥利项、随转动产生的向心项和角加速度项。
将除 $\vec a'$ 外的参考系运动项移到动力学方程右侧并乘 $-m$，就得到相应惯性力。
其中 $-m\vec\omega\times(\vec\omega\times\vec r)$ 指向远离转轴的方向，即离心力。
本节转杆问题中原点固定、$\vec\omega$ 恒定且质点相对杆静止，因而
$\vec a_0=0$、$\dot{\vec\omega}=0$、$\vec v'=0$，只需加入离心力。
转盘上运动物体的侧向偏转则还涉及科里奥利项，不能把所有非惯性效应都写成同一个 $-m\vec a_0$。


\begin{origexamplebox}{}\label{ch06:ex:s06-2-24}
很多时候可以近似认为地球表面是一个惯性参考系。现有一辆相对于地面以加速度 $a$ 作
匀加速直线运动的车，车内一个质量为 $m$ 的小球被悬挂在车顶。求小球相对车厢静止时，
悬线与竖直方向的夹角 $\theta_1$。悬线长为 $l$，轻且不可伸长；若将小球拉到竖直方向，
相对车厢由静止释放，求另一侧最高点与竖直方向的夹角 $\theta_2$。
全过程中车的加速度保持不变，忽略阻力及小球与车厢边界的碰撞。
\probdiagTwoTwentyFour
\end{origexamplebox}

\Needspace{16\baselineskip}
\Needspace{6\baselineskip}\begin{xesolution}
以下推导所用的解图为图~\ref{mech:legacy:ch06:s06-2-24}。

以车为参考系。车向右加速，因此小球受到大小为 $ma$、方向向左的惯性力。平衡位置的
受力和有效重力方向如解图所示：
\soldiagTwoTwentyFour
由水平、竖直方向平衡条件
\[
 T\sin\theta_1=ma,
 \qquad
 T\cos\theta_1=mg,
\]
得到
\[
 \boxed{\tan\theta_1=\frac ag}.
\]
有效重力大小为 $g_{\rm eff}=\sqrt{g^2+a^2}$，方向相对竖直向后偏 $\theta_1$。
图中 $B_0$ 是竖直释放点，$B_e$ 是平衡位置，$B_2$ 是另一转折点。设摆线相对通常竖直方向
向后偏转的角度为 $\theta$，它与有效重力的夹角为 $\theta-\theta_1$。
沿有效重力方向向下的坐标为 $l\cos(\theta-\theta_1)$，因此以悬点为势能零点，
与恒定有效重力对应的势能是
\[
 U=-mg_{\rm eff}l\cos(\theta-\theta_1).
\]
这是将重力和恒定惯性力共同计入的有效势能。车系中绳张力始终与相对速度垂直，不做功，
故小球相对车厢的动能 $mv'^2/2$ 与这一势能之和保持不变。
小球从 $\theta=0$ 处静止释放。另一侧最高点也是速度为零，能量守恒给出
\[
 \cos(\theta_2-\theta_1)=\cos\theta_1.
\]
除初始点外的另一个转折点为
\[
 \boxed{\theta_2=2\theta_1
 =2\arctan\frac ag}.
\]
图中这两个转折位置关于 $OB_e$ 对称，但并不关于通常的竖直线对称。
在摆动范围内张力为 $ml^{-1}v'^2+mg_{\rm eff}\cos(\theta-\theta_1)\ge mg>0$，
所以悬线保持绷紧。$a\to0$ 时两个角均趋于零，竖直释放的小球不再偏摆。
\end{xesolution}

\begin{origexamplebox}{}\label{ch06:ex:s06-2-25}
不计空气阻力，证明在地表附近自由下落的“电梯”中，仅受重力、彼此不相互作用且尚未碰壁的
自由物体相对于电梯均保持静止或作匀速直线运动。思考一下：在证明过程中忽略了什么，
得到了什么样的结论？这个电梯是一个惯性
参考系吗？
\probdiagTwoTwentyFive
\end{origexamplebox}

\Needspace{6\baselineskip}\begin{xesolution}
以下推导所用的解图为图~\ref{mech:legacy:ch06:s06-2-25}。

在地面惯性系的近似下，电梯与其中任意自由物体都只受重力，因而
\[
 \vec a_{\rm E}=\vec g,
 \qquad
 \vec a_{\rm P}=\vec g.
\]
故物体相对电梯的加速度为
\[
 \vec a_{\rm rel}=\vec a_{\rm P}-\vec a_{\rm E}=0.
\]
于是它相对电梯保持静止或作匀速直线运动。在电梯系中也可以用解图理解：向上的惯性力
$-m\vec g$ 恰好抵消向下的重力 $m\vec g$。
\soldiagTwoTwentyFive
这里忽略了地球自转、空气阻力以及重力随高度和水平位置的变化。尤其是重力场不完全均匀：
相距有限的物体具有微小的相对加速度，即潮汐效应。

在牛顿力学中，这个电梯相对于地面惯性系有加速度，因此不是惯性参考系；只是对足够小的
空间和足够短的时间，均匀重力与惯性力抵消，使自由物体呈现相对匀速的运动。这个结论也是
失重现象与等效原理的经典出发点。
\end{xesolution}

\begin{origexamplebox}{}\label{ch06:ex:s06-2-26}
如图所示，汽车内固定有一个倾角为 $\theta$ 的斜面，斜面表面光滑，底部有一个质量为
$m$ 的小物块。开始时系统静止，某时刻汽车突然开始向左以加速度 $a$ 作匀加速运动。
求使物块刚好具有沿斜面向上运动趋势的临界加速度；若要求物块从静止开始真正向上滑动，
$a$ 还应满足什么条件？
\probdiagTwoTwentySix
\end{origexamplebox}

\Needspace{6\baselineskip}\begin{xesolution}
以下推导所用的解图为图~\ref{mech:legacy:ch06:s06-2-26}。

在汽车参考系中，惯性力 $ma$ 水平向右。沿斜面方向分解重力与惯性力，受力关系如解图：
\soldiagTwoTwentySix
取沿斜面向上为正，则
\[
 ma_{\parallel}=ma\cos\theta-mg\sin\theta.
\]
使物块刚好具有向上运动趋势的临界条件为 $a_{\parallel}=0$，故
\[
 \boxed{a_{\min}=g\tan\theta}.
\]
在临界值处沿斜面合力为零，物块不会从静止状态自行上滑；若要真正开始向上运动，必须有
\[
 \boxed{a>g\tan\theta}.
\]
\end{xesolution}

\begin{origexamplebox}{}\label{ch06:ex:s06-2-27}
如图所示，在光滑水平面上有一光滑斜面，倾角为 $\alpha$。一质量为 $m$ 的物体放在斜面上，
在水平力 $F$ 推动下恰使物体 $m$ 与斜面无相对滑动。求斜面对物块弹力的大小。
\probdiagTwoTwentySeven
\end{origexamplebox}

\Needspace{6\baselineskip}\begin{xesolution}
以下推导所用的解图为图~\ref{mech:legacy:ch06:s06-2-27}。

以随斜面向右加速的参考系分析。物块相对斜面静止，受到重力、支持力和水平向左的惯性力，
三力平衡，如解图所示：
\soldiagTwoTwentySeven
沿斜面方向平衡要求
\[
 ma\cos\alpha=mg\sin\alpha,
 \qquad\Rightarrow\qquad a=g\tan\alpha.
\]
沿斜面法向平衡，或者直接在地面惯性系的竖直方向列式，得到
\[
 N\cos\alpha=mg.
\]
因此
\[
 \boxed{N=\frac{mg}{\cos\alpha}}.
\]
题设没有给出斜面质量，所以不能由现有条件求 $F$；但弹力由物块的竖直平衡唯一确定，
不需要知道 $F$ 和斜面质量。
\end{xesolution}

\begin{origexamplebox}{}\label{ch06:ex:s06-2-28}
如图所示，一个质量为 $M$、两边倾角分别为 $\alpha$ 和 $\beta$ 的斜面放在水平面上。
两个斜面上分别放有质量为 $m_1$ 和 $m_2$ 的物块 $A$ 和 $B$。忽略一切摩擦，求由静止
释放后所有物体的加速度。
\probdiagTwoTwentyEight
\end{origexamplebox}

\Needspace{6\baselineskip}\begin{xesolution}
以下推导所用的解图为图~\ref{mech:legacy:ch06:s06-2-28}。

取水平向右为正，设斜面的水平加速度为 $A=\ddot X$；令 $s_1,s_2$ 分别表示两物块相对
斜面沿各自斜面向下的位移。坐标和正方向如解图：
\soldiagTwoTwentyEight
地面系中的坐标关系为（省略不影响加速度的坐标常数）
\[
 x_1=X-s_1\cos\alpha,\quad y_1=-s_1\sin\alpha,
 \qquad x_2=X+s_2\cos\beta,\quad y_2=-s_2\sin\beta.
\]
两块分别受重力和垂直于各自斜面的支持力。以左块为例，将加速度沿左斜面向下投影，得到
\[
 (A-\ddot s_1\cos\alpha)(-\cos\alpha)
 +(-\ddot s_1\sin\alpha)(-\sin\alpha)
 =\ddot s_1-A\cos\alpha.
\]
支持力投影为零，重力投影为 $m_1g\sin\alpha$。右块同理，两个方程为
\[
 \ddot s_1-A\cos\alpha=g\sin\alpha,
 \qquad
 \ddot s_2+A\cos\beta=g\sin\beta.
\]
系统水平方向不受外力，因此
\[
 (M+m_1+m_2)A-m_1\ddot s_1\cos\alpha
 +m_2\ddot s_2\cos\beta=0.
\]
记
\[
 D=M+m_1\sin^2\alpha+m_2\sin^2\beta.
\]
联立解得斜面的加速度
\[
 \boxed{A=\frac{g\left(m_1\sin\alpha\cos\alpha
 -m_2\sin\beta\cos\beta\right)}{D}}.
\]
$A>0$ 表示斜面向右，$A<0$ 表示向左。两物块相对斜面的加速度分别为
\[
 \boxed{a_{1/\!M}=\ddot s_1=g\sin\alpha+A\cos\alpha},
\]
\[
 \boxed{a_{2/\!M}=\ddot s_2=g\sin\beta-A\cos\beta}.
\]
它们相对地面的加速度分量为
\[
 \boxed{
 a_{1x}=A-\ddot s_1\cos\alpha,
 \qquad
 a_{1y}=-\ddot s_1\sin\alpha},
\]
\[
 \boxed{
 a_{2x}=A+\ddot s_2\cos\beta,
 \qquad
 a_{2y}=-\ddot s_2\sin\beta}.
\]
这些式子给出了斜面和两个物块的全部对地加速度。
若 $m_1=m_2$ 且 $\alpha=\beta$，分子相消，$A=0$，两块沿固定斜面以 $g\sin\alpha$ 下滑。
若 $m_2\to0$，公式还原为单物块在活动斜面上的结果。两个法向力分别为
$N_1=m_1(g\cos\alpha-A\sin\alpha)$、$N_2=m_2(g\cos\beta+A\sin\beta)$。
例如代入后
\[
 N_1=\frac{m_1g[M\cos\alpha+m_2\sin\beta\sin(\alpha+\beta)]}{D}>0,
\]
$N_2$ 交换下标及角度可得同样结论，因此在 $M>0$、两倾角均为锐角时接触约束成立。
\end{xesolution}

\begin{origexamplebox}{}\label{ch06:ex:s06-2-29}
一个装满水的水箱中有一个气泡。当水箱以加速度 $a$ 向前加速运动时，这个气泡将
\underline{\hspace{3em}}（向前、向后）运动。
\probdiagTwoTwentyNine
\end{origexamplebox}

\Needspace{6\baselineskip}\begin{xesolution}
以下推导所用的解图为图~\ref{mech:legacy:ch06:s06-2-29}。

在水箱参考系中，液体受到向后的惯性力。把重力和惯性力合成，有效重力方向为“向下且
向后”。取 $x$ 向前、$y$ 向上，则 $\vec g_{\rm eff}=(-a,-g)$。
在液体随箱共同加速、相对流动的瞬态可忽略时，取宽 $\dd x$、横截面积 $S$ 的液体微元，
水平平衡要求 $-S\dd p-\rho S\dd x\,a=0$，即 $\partial p/\partial x=-\rho a$。
竖直方向同理有 $\partial p/\partial y=-\rho g$。所以
\[
 p=p_0-\rho ax-\rho gy.
\]
压强沿有效重力方向增大；在等压线上 $a\dd x+g\dd y=0$，故其斜率为 $-a/g$，
与有效重力垂直。图中箭头及等压线体现了这一关系：
\soldiagTwoTwentyNine
体积为 $V$ 的气泡受到周围水的合压力
$\vec F_b=-\rho_{\text{水}}V\vec g_{\rm eff}$；再加上气泡自身的重力和惯性力，
驱动力为 $(\rho_{\text{气}}-\rho_{\text{水}})V\vec g_{\rm eff}$。
气泡密度远小于水，所以它相对水箱向上且向前运动。这里只判断运动趋势，
实际加速度和速度还受周围液体运动及阻力影响。横向选择应填
\[
 \boxed{\text{向前}}.
\]
容易混淆之处是：致密小物体趋向有效重力方向，轻气泡却因浮力而向相反方向运动。
\end{xesolution}

如果物体相对于匀速转动的参考系保持静止，那么在惯性系中它作匀速圆周运动，所需合力指向转轴，大小为
\begin{equation}
 F=m\omega^2r \label{ch06:eq:s06-2-13}
\end{equation}
在转动系看来，可引入大小相同、方向沿径向向外的离心力，使相对
静止问题写成平衡形式。

\begin{origexamplebox}{}\label{ch06:ex:s06-2-30}
一根细杆与竖直方向的夹角为 $0<\theta<\pi/2$，绕通过其固定下端 $O$ 的竖直轴，
以恒定角速度 $\omega\ge0$ 匀速转动。杆上距 $O$ 为 $l>0$ 处套有一个质点。
先设接触无摩擦，质点能够与杆保持相对静止，求角速度
$\omega$。如果质点与杆之间有摩擦，且最大静摩擦因数为 $\mu$，求能够保持相对静止的角速度
$\omega$ 的范围。
\probdiagTwoThirty
\end{origexamplebox}

\Needspace{6\baselineskip}\begin{xesolution}
以下推导所用的解图为图~\ref{mech:legacy:ch06:s06-2-30}。

先作质点到竖直转轴的水平投影。杆长段 $OP=l$ 与竖直方向成 $\theta$，故转动半径为 $r=l\sin\theta$：
\soldiagTwoThirtyGeometry
在随杆转动的参考系中，质点相对静止，受到
重力、杆的支持力以及水平向外的离心力
\[
 F_{\rm cf}=m\omega^2l\sin\theta.
\]
沿杆方向的两个分量如解图所示：
\soldiagTwoThirty
无摩擦时沿杆方向平衡：
\[
 m\omega^2l\sin^2\theta=mg\cos\theta.
\]
因此
\[
 \boxed{\omega=\sqrt{\frac{g\cos\theta}{l\sin^2\theta}}}.
\]

有摩擦时，维持平衡所需摩擦力的大小为
\[
 |f|=m\left|g\cos\theta-\omega^2l\sin^2\theta\right|.
\]
重力与离心力在垂直于杆方向的合分量由支持力平衡，故
\[
 N=m\sin\theta\left(g+\omega^2l\cos\theta\right).
\]
静摩擦条件 $|f|\le\mu N$ 给出最稳妥的角速度范围表达式
\[
 \boxed{
 \left|g\cos\theta-\omega^2l\sin^2\theta\right|
 \le\mu\sin\theta\left(g+\omega^2l\cos\theta\right)}.
\]
取沿杆离开 $O$ 为正，所需有符号摩擦力为
$f_s=m(g\cos\theta-\omega^2l\sin^2\theta)$：转得慢时它阻止下滑，沿杆向上；
转得快时它阻止外滑，沿杆向下。拆开绝对值并令 $z=\omega^2\ge0$，得到
\begin{align*}
 zl(\sin^2\theta+\mu\sin\theta\cos\theta)&\ge g(\cos\theta-\mu\sin\theta),\\
 zl(\sin^2\theta-\mu\sin\theta\cos\theta)&\le g(\cos\theta+\mu\sin\theta).
\end{align*}
第一式中 $z$ 的系数总为正，因而对任何 $\mu\ge0$ 都给出下界
\[
 \boxed{\omega^2\ge
 \max\!\left\{0,
 \frac{g(\cos\theta-\mu\sin\theta)}
 {l(\sin^2\theta+\mu\sin\theta\cos\theta)}\right\}},
\]
第二式右端为正。若 $\mu<\tan\theta$，左端系数也为正，另有上界
\[
 \boxed{\omega^2\le
 \frac{g(\cos\theta+\mu\sin\theta)}
 {l(\sin^2\theta-\mu\sin\theta\cos\theta)}}.
\]
当 $\mu=\tan\theta$ 时第二式左端为零；当 $\mu>\tan\theta$ 时其左端非正。
两种情况都自动满足第二式，所以只需满足前述下界，理想模型下\textbf{没有有限的角速度上界}。
若 $\mu=0$，上下界重合为无摩擦结果；若 $\mu\ge\cot\theta$，下界为零，说明静止杆上也能靠摩擦托住质点。
\end{xesolution}

\clearpage
\begin{origexamplebox}{}\label{ch06:ex:s06-2-31}
航空医学研究中使用的离心机，是一个绕竖直轴转动的水平杆，杆的一端携带一个实验对象。
如果从转动中心到实验对象的距离为 $7\,\mathrm m$，离心机要转动多快，才能使乘坐者获得
大小为 $5g$ 的水平向心加速度？

\begin{qualityfigure}
\begin{tikzpicture}[>=Latex,font=\small]
\draw[dashed](0,0)circle(2.3);\draw[very thick](0,0)--(2.3,0)node[midway,below]{$r=7\,\mathrm m$};\fill(0,0)circle(2pt)node[below left]{$O$};\fill(2.3,0)circle(4pt)node[right]{实验对象};
\draw[->,Gold](.5,1.9)arc(75:150:2)node[midway,above]{$\omega$};\node at(1,-2.8){水平向心加速度要求 $5g$};
\end{tikzpicture}
\figcap{离心机的水平杆、轴心和实验对象}\label{mech:06:31:q}
\end{qualityfigure}
\end{origexamplebox}

\Needspace{6\baselineskip}\begin{xesolution}题图~\ref{mech:06:31:q}中，半径从转轴 $O$ 量到实验对象；$5g$ 明确指水平向心加速度，故直接用 $\omega^2r=5g$，不把竖直重力混入这一给定分量。

离心机提供的水平向心加速度满足
\[
 \omega^2r=5g.
\]
运动关系如解图所示：
\soldiagTwoThirtyOne
取 $r=7\,\mathrm m$、$g=9.8\,\mathrm{m/s^2}$，得
\[
 \boxed{\omega=\sqrt{\frac{5g}{r}}=\sqrt7
 \approx2.65\,\mathrm{rad/s}}.
\]
对应转动频率和转速为
\[
 f=\frac{\omega}{2\pi}\approx0.421\,\mathrm{Hz},
 \qquad
 \boxed{n\approx25.3\,\mathrm{r/min}},
\]
周期约为 $2.37\,\mathrm s$。

如果实际装置把“$5g$ 过载”定义为座椅支持力对应的总表观加速度，并计入竖直重力，条件将
改为
\[
 \sqrt{(\omega^2r)^2+g^2}=5g.
\]
这与题设明确给出的“水平向心加速度为 $5g$”是两个不同的指标。
\end{xesolution}
}

\endgroup

% MechanicsNumberedSolutionFigures20261023-v1
% MechanicsQualityRevision20261023-v1
% QualityRevision20261023
\begingroup
\tcbset{breakable/.default=false}
\chapter{动量、变质量系统与碰撞}

{
\graphicspath{{../../../sample/xes-original/seventh-section/}}
\section{动量定理}
% Original problem figures and added solution diagrams for Section 3.1.
\tikzset{
  solbody/.style={draw=Ink,line width=0.8pt},
  solforce/.style={-{Latex[length=2.0mm]},draw=DeepBlue,line width=0.9pt},
  solimpulse/.style={-{Latex[length=2.4mm]},draw=DeepBlue,line width=1.25pt},
  solgeom/.style={draw=TextGray,densely dashed,line width=0.55pt},
  solflow/.style={-{Latex[length=1.8mm]},draw=Blue,line width=0.75pt},
  sollabel/.style={font=\scriptsize,fill=none,inner sep=1pt},
  solpanel/.style={font=\scriptsize\bfseries,text=TextGray}
}
\newcommand{\diagcaption}[2]{%
 \refstepcounter{figure}\par\nopagebreak[4]\smallskip
 {\centering\scriptsize\color{TextGray}图~\thefigure\quad #2\par}\smallskip}

\newcommand{\probdiagThreeOne}{%
\begin{center}\begin{minipage}{.97\linewidth}\centering
\begin{tikzpicture}[scale=.78]
 \draw[solbody] (-2.8,1.2)--(0,1.2)--(0,-1.2)--(3,-1.2);
 \draw[DeepBlue,line width=1pt] (-1.8,1.22)--(.9,1.22);
 \draw[solbody,fill=SoftGray] (-.82,1.24) rectangle (-.08,1.85);
 \node[sollabel] at (-.45,1.55) {铁块};
 \draw[solflow] (.3,1.65)--(1.5,1.65) node[midway,above=3pt,sollabel] {$v$：抽纸速度};
 \node[sollabel] at (1.4,.96) {纸条};
 \draw[solbody,fill=SoftGray] (1.65,-1.2) rectangle (2.35,-.62);
 \node[sollabel] at (2,-.35) {$P$};
\end{tikzpicture}
\diagcaption{题图}{纸条从桌边抽出，铁块随后从桌边落下。}\label{ch07:fig:probdiagThreeOne}\end{minipage}\end{center}}

\newcommand{\soldiagThreeOne}{%
\begin{center}\begin{minipage}{.97\linewidth}\centering
\begin{tikzpicture}[scale=.78]
 \draw[solbody] (-1.1,2.2)--(0,2.2)--(0,0);
 \draw[axis] (0,0)--(6.7,0) node[right,sollabel] {$x$};
 \draw[curve,domain=0:1,samples=60] plot ({3*\x},{2.2*(1-\x*\x)});
 \draw[TextGray,line width=.9pt,domain=0:1,samples=60] plot ({5.5*\x},{2.2*(1-\x*\x)});
 \draw[solflow] (0,2.2)--(1.3,2.2) node[midway,above=4pt,sollabel] {$v_x$};
 \node[sollabel,below=4pt] at (3,0) {$O$：抽速 $2v$};
 \node[sollabel,below=4pt] at (5.5,0) {$P$：抽速 $v$};
 \node[sollabel] at (3.1,3.05) {$I_f=f\tau$；接触时间较短，离桌速度较小};
\end{tikzpicture}
\diagcaption{解图}{抽得越快，摩擦力作用越短，离桌速度和水平射程越小。}\label{ch07:fig:soldiagThreeOne}\end{minipage}\end{center}}

\newcommand{\soldiagThreeTwo}{%
\begin{center}\begin{minipage}{.97\linewidth}\centering
\begin{tikzpicture}[x=.42cm,y=.25cm]
  \draw[axis] (0,-5.5)--(0,9.7) node[above,sollabel] {$F/\mathrm N$};
  \draw[axis] (0,0)--(16.5,0) node[right,sollabel] {$t/\mathrm s$};
  \draw[curve] (0,8)--(5,8)--(5,5)--(9,5)--(9,0)--(15,0);
  \draw[draw=TextGray,line width=.9pt] (0,-4)--(15,-4);
  \draw[solgeom] (5,.5)--(5,8.5) (9,.5)--(9,5.5);
  \foreach \x in {5,9,15}{\draw[solgeom] (\x,-4.5)--(\x,-1.7);}
  \foreach \x/\lab in {5/5,9/9,15/15}{\draw (\x,.25)--(\x,-.25) node[below,sollabel] {\lab};}
  \node[sollabel] at (2.5,8.8) {$F_1=8$}; \node[sollabel] at (7,5.8) {$F_2=5$};
  \node[sollabel] at (12,-4.8) {$f=4$（向左）};
\end{tikzpicture}
\diagcaption{解图}{应用力的正冲量面积等于摩擦力的负冲量面积。}\label{ch07:fig:soldiagThreeTwo}\end{minipage}\end{center}}

\newcommand{\probdiagThreeThree}{%
\begin{center}\begin{minipage}{.97\linewidth}\centering
\begin{tikzpicture}[scale=.72]
  \draw[solbody,fill=white] (-.48,1.55) rectangle (.48,2.25); \node[sollabel] at (0,1.90) {$m_1$};
  \draw[decorate,decoration={coil,segment length=4pt,amplitude=3pt},line width=.75pt] (0,1.55)--(0,-.55);
  \draw[solbody,fill=white] (-.48,-1.25) rectangle (.48,-.55); \node[sollabel] at (0,-.90) {$m_2$};
  \draw[solforce] (1.0,1.9)--++(0,-.9) node[below,sollabel] {$v_1$};
  \node[sollabel] at (-1.25,.5) {$k,\ l_0$};
\end{tikzpicture}
\diagcaption{题图}{弹簧初态未知的二物体系统由静止竖直下落。}\label{ch07:fig:probdiagThreeThree}\end{minipage}\end{center}}

\newcommand{\soldiagThreeThree}{%
\begin{center}\begin{minipage}{.97\linewidth}\centering
\begin{tikzpicture}[scale=.78]
  \draw[solbody,rounded corners=1.5pt] (-1.45,-1.2) rectangle (1.45,1.25);
  \node[sollabel] at (0,.55) {$m_1+m_2+$ 弹簧};
  \draw[solimpulse] (0,.15)--++(0,-1.65) node[below,sollabel] {$(m_1+m_2)gt$};
  \node[sollabel,align=left] at (3.9,.2) {$P_f=m_1v_1+m_2v_2$\\（向下为正）};
  \node[sollabel] at (2.05,.2) {$\Longrightarrow$};
\end{tikzpicture}
\diagcaption{解图}{把弹簧力纳入系统内部，只保留重力的外冲量。}\label{ch07:fig:soldiagThreeThree}\end{minipage}\end{center}}

\newcommand{\probdiagThreeFour}{%
\begin{center}\begin{minipage}{.97\linewidth}\centering
\begin{tikzpicture}[scale=.78]
 \begin{scope}[rotate=25]
  \draw[solbody] (-3,0)--(3.2,0);
  \draw[solbody,fill=SoftGray] (-2.2,0) rectangle (-1.25,.7);
  \draw[solbody,fill=SoftGray] (.35,0) rectangle (1.3,.7);
  \draw[DeepBlue] (-1.25,.35)--(.35,.35);
  \node[sollabel,rotate=0] at (-1.725,.35) {$A$};
  \node[sollabel,rotate=0] at (.825,.35) {$B$};
  \draw[solforce] (1.3,.35)--(2.8,.35) node[midway,above=4pt,sollabel] {$F$};
 \end{scope}
\end{tikzpicture}
\diagcaption{题图}{$A$、$B$ 沿斜面上滑，经过 $t_1$ 后连接线断裂。}\label{ch07:fig:probdiagThreeFour}\end{minipage}\end{center}}

\newcommand{\soldiagThreeFour}{%
\begin{center}\begin{minipage}{.97\linewidth}\centering
\begin{tikzpicture}[x=1cm,y=.53cm]
 \draw[axis] (0,0)--(5.7,0) node[right,sollabel] {$t$};
 \draw[axis] (0,0)--(0,8) node[above,sollabel] {$v$};
 % Illustrative ratio m_A=m_B and t_2=t_1; all slopes share these parameters.
 \draw[curve] (0,0)--(2,2);
 \draw[curve] (2,2)--(4,0);
 \draw[TextGray,line width=.9pt] (2,2)--(4,8);
 \draw[solgeom] (0,2)--(2,2)--(2,0) (4,0)--(4,8);
 \node[sollabel,left=4pt] at (0,2) {$at_1$};
 \node[sollabel,below=4pt] at (2,0) {$t_1$};
 \node[sollabel,below=4pt] at (4,0) {$t_1+t_2$};
 \node[sollabel,anchor=west] at (4.15,7.3) {$B$：$a_B>a$};
 \node[sollabel,anchor=west] at (4.15,1.2) {$A$：减速至零};
\end{tikzpicture}
\diagcaption{解图}{断线后 $B$ 的斜率增大；图取 $m_A=m_B,\ t_2=t_1$ 显示相同定性关系。}\label{ch07:fig:soldiagThreeFour}\end{minipage}\end{center}}

\newcommand{\probdiagThreeFive}{%
\begin{center}\begin{minipage}{.97\linewidth}\centering
\begin{tikzpicture}[scale=.74]
  \draw[solbody] (-3.0,-.85)--(3.15,-.85);
  \draw[solbody,fill=white] (-2.15,-.84) rectangle (-1.3,-.12); \node[sollabel] at (-1.72,-.48) {$A$};
  \draw[solbody,fill=white] (1.35,-.84) rectangle (2.20,-.12); \node[sollabel] at (1.78,-.48) {$B$};
  \draw[solforce] (-1.30,-.48)--++(1.25,0) node[right,sollabel] {$v_0$};
  \node[sollabel] at (-1.72,-1.15) {$m_A=2\,\mathrm{kg}$};
  \node[sollabel] at (1.78,-1.15) {$m_B=1\,\mathrm{kg}$};
\end{tikzpicture}
\diagcaption{题图}{箱子 $A$ 在光滑面上冲向静止的箱子 $B$。}\label{ch07:fig:probdiagThreeFive}\end{minipage}\end{center}}

\newcommand{\soldiagThreeFive}{%
\begin{center}\begin{minipage}{.97\linewidth}\centering
\begin{tikzpicture}[scale=.8]
 \node[solpanel] at (-2.4,1.6) {碰前};
 \node[solpanel] at (2.4,1.6) {碰后};
 \draw[solbody,fill=SoftGray] (-3.5,-.35) rectangle (-2.6,.35);
 \draw[solbody,fill=SoftGray] (-1.8,-.35) rectangle (-.9,.35);
 \node[sollabel] at (-3.05,0) {$A$}; \node[sollabel] at (-1.35,0) {$B$};
 \draw[solbody,fill=SoftGray] (1,-.35) rectangle (1.9,.35);
 \draw[solbody,fill=SoftGray] (2.6,-.35) rectangle (3.5,.35);
 \node[sollabel] at (1.45,0) {$A$}; \node[sollabel] at (3.05,0) {$B$};
 \draw[solimpulse] (-3.5,.85)--(-.9,.85) node[midway,above=3pt,sollabel] {$P_i=m_Av_0$};
 \draw[solimpulse] (1,.85)--(3.6,.85) node[midway,above=3pt,sollabel] {$P_f=P_i$};
 \node[sollabel] at (0,-1.05) {以上箭头表示系统总动量；$I_x^{\rm ext}=0$};
\end{tikzpicture}
\diagcaption{解图}{碰撞内力冲量相消，系统总动量由碰前原样传到碰后。}\label{ch07:fig:soldiagThreeFive}\end{minipage}\end{center}}

\newcommand{\probdiagThreeSix}{%
\begin{center}\begin{minipage}{.97\linewidth}\centering
\begin{tikzpicture}[scale=.85]
 \coordinate (A) at (-2.4,.8); \coordinate (B) at (0,.8);
 \coordinate (C) at ($(B)+(-60:2.4)$);
 \draw[solbody] (A)--(B)--(C); \draw[solgeom] (B)--++(1.8,0);
 \draw[TextGray] ($(B)+(0:.65)$) arc (0:-60:.65);
 \draw[TextGray] ($(B)+(180:.5)$) arc (180:300:.5);
 \foreach \p in {A,B,C}{\fill[Ink] (\p) circle (2.5pt);}
 \node[sollabel,above=7pt] at (A) {$A$}; \node[sollabel,above=7pt] at (B) {$B$};
 \node[sollabel,below=6pt] at (C) {$C$};
 \node[sollabel] at (.95,.34) {$60^\circ$}; \node[sollabel] at (-.8,.05) {$120^\circ$};
 \draw[solimpulse] (A)--++(-1.1,0) node[midway,above=4pt,sollabel] {$I$};
\end{tikzpicture}
\diagcaption{题图}{冲量 $I$ 沿 $BA$ 方向；两根绳在冲击后仍不可伸长。}\label{ch07:fig:probdiagThreeSix}\end{minipage}\end{center}}

\newcommand{\soldiagThreeSix}{%
\begin{center}\begin{minipage}{.97\linewidth}\centering
\begin{tikzpicture}[scale=.85]
 \coordinate (A) at (-2.7,1); \coordinate (B) at (0,1);
 \coordinate (C) at ($(B)+(-60:2.7)$);
 \draw[solgeom] (A)--(B)--(C);
 \foreach \p in {A,B,C}{\fill[Ink] (\p) circle (2.4pt);}
 \node[sollabel,below=8pt] at (A) {$A$}; \node[sollabel,above=7pt] at (B) {$B$};
 \node[sollabel,below=7pt] at (C) {$C$};
 \draw[solimpulse] (A)--++(-1.05,0) node[midway,above=4pt,sollabel] {$I$};
 \draw[solimpulse] (A)--++(.95,0) node[midway,above=4pt,sollabel] {$J_1$};
 \draw[solimpulse] (B)--++(-.95,0) node[midway,above=4pt,sollabel] {$J_1$};
 \draw[solimpulse] (B)--++(-60:1) node[midway,right=5pt,sollabel] {$J_2$};
 \draw[solimpulse] (C)--++(120:1) node[midway,right=5pt,sollabel] {$J_2$};
 \draw[solflow] (2.4,.7)--++(-60:.85) node[right=4pt,sollabel] {$\hat{\boldsymbol u}$};
 \draw[solflow] (2.4,.7)--++(.95,0) node[right=3pt,sollabel] {$+x$};
\end{tikzpicture}
\diagcaption{解图}{先对三球分别列冲量方程，再用两条绳方向上的速度约束闭合方程组。}\label{ch07:fig:soldiagThreeSix}\end{minipage}\end{center}}

\newcommand{\soldiagThreeSeven}{%
\begin{center}\begin{minipage}{.97\linewidth}\centering
\begin{tikzpicture}[scale=.82]
 % g=1, v0=2, theta=45 degrees, wall at s0=2.5 in drawing units.
 \coordinate (Q) at (2.5,.9375);
 \draw[solbody] (-.3,0)--(2.5,0)--(2.5,2.8);
 \draw[curve,domain=0:2.5,samples=60] plot ({\x},{\x-\x*\x/4});
 \draw[solflow] (0,0)--(45:1.1) node[above left=3pt,sollabel] {$v_0$};
 \draw[TextGray] (.5,0) arc (0:45:.5); \node[sollabel] at (.85,.22) {$\theta$};
 \draw[<->,TextGray] (0,-.45)--(2.5,-.45) node[midway,below=3pt,sollabel] {$s_0$};
 \fill[Ink] (Q) circle (2pt) node[below left=5pt,sollabel] {$Q$};
 \draw[solflow] (2.5,1.05)--(2.5,2.3);
 \node[sollabel,anchor=west] at (2.7,2.25) {蹬后竖直上升};
 \begin{scope}[shift={(5.2,1.2)}]
  \fill[Ink] (0,0) circle (2.5pt);
  \draw[solimpulse] (0,0)--(-1.1,0) node[midway,below=4pt,sollabel] {$J_N$};
  \draw[solimpulse] (0,0)--(0,1.2) node[right=4pt,sollabel] {$J_f$};
  \node[sollabel] at (0,-1.1) {接触冲量，单独画出};
 \end{scope}
\end{tikzpicture}
\diagcaption{解图}{法向冲量消去水平速度，最大向上摩擦冲量继续增加竖直速度。}\label{ch07:fig:soldiagThreeSeven}\end{minipage}\end{center}}

\newcommand{\probdiagThreeEight}{%
\begin{center}\begin{minipage}{.97\linewidth}\centering
\begin{tikzpicture}[scale=.8]
 \coordinate (O) at (-.5,1.4); \coordinate (S) at (-2,-.4);
 \coordinate (A) at ({-.5+2.8*1.5/1.8},-1.4);
 \coordinate (B) at ({-.5+2.8*1.5/1.8+.7},-1.4);
 \draw[solbody] (-3,1.4)--(3.1,1.4) (-3,-1.4)--(3.1,-1.4);
 \draw[TextGray,line width=.8pt] (S)--(O)--(A);
 \draw[curve] (O)--(B);
 \draw[solflow] (S)--($(S)!.45!(O)$) node[midway,above left=3pt,sollabel] {$\vec v$};
 \draw[solgeom] (O)--(-.5,-1.4);
 \foreach \p in {O,S,A,B}{\fill[Ink] (\p) circle (2pt);}
 \node[sollabel,above=5pt] at (O) {$O$}; \node[sollabel,left=5pt] at (S) {$S$};
 \node[sollabel,below=5pt] at (A) {$A$}; \node[sollabel,below=5pt] at (B) {$B$};
 \draw[<->,TextGray] (-3.35,-1.4)--(-3.35,1.4) node[midway,left=4pt,sollabel] {$d$};
 \node[sollabel] at (0,2.35) {俯视；本图取摩擦冲量向右};
\end{tikzpicture}
\diagcaption{题图}{光滑墙落点为 $A$；有摩擦时落点为 $B$。}\label{ch07:fig:probdiagThreeEight}\end{minipage}\end{center}}

\newcommand{\soldiagThreeEight}{%
\begin{center}\begin{minipage}{.97\linewidth}\centering
\begin{tikzpicture}[scale=.8]
 \coordinate (O) at (-1,1.2); \coordinate (A) at (.2,-1.2); \coordinate (B) at (1.4,-1.2);
 \draw[solbody] (-2,1.2)--(2,1.2) (-2,-1.2)--(2,-1.2);
 \draw[TextGray] (O)--(A); \draw[curve] (O)--(B);
 \node[sollabel,above=5pt] at (O) {$O$};
 \node[sollabel,below left=5pt] at (A) {$A$}; \node[sollabel,below right=5pt] at (B) {$B$};
 \draw[solgeom] (A)--(.2,-1.9) (B)--(1.4,-1.9);
 \draw[<->,DeepBlue] (.2,-1.8)--(1.4,-1.8) node[midway,below=4pt,sollabel] {$AB=|\Delta v_t|d/v_n$};
 \begin{scope}[shift={(4,0)}]
  \fill[Ink] (0,.6) circle (2pt);
  \draw[solimpulse] (0,.6)--(0,-.55) node[midway,left=4pt,sollabel] {$J_N$};
  \draw[solimpulse] (0,.6)--(1.2,.6) node[midway,above=4pt,sollabel] {$J_t$};
  \node[sollabel,align=center] at (.35,-1.25) {持续滑动时\\$J_t=\mu J_N$};
 \end{scope}
\end{tikzpicture}
\diagcaption{解图}{左右两部分分别表示落点偏移和接触冲量，不能把冲量箭头当作轨迹。}\label{ch07:fig:soldiagThreeEight}\end{minipage}\end{center}}

\newcommand{\soldiagThreeNine}{%
\begin{center}\begin{minipage}{.97\linewidth}\centering
\begin{tikzpicture}[scale=.82]
 \draw[solbody,fill=SoftGray] (1.2,-1.1) rectangle (1.9,1.1);
 \draw[solflow] (-2.5,.5)--(1.2,.5) node[pos=.35,above=4pt,sollabel] {$v_{\rm in}=v$};
 \draw[solflow] (1.2,-.5)--(-2.5,-.5) node[pos=.65,below=4pt,sollabel] {$v_{\rm out}=-v$};
 \node[sollabel] at (-.1,1.75) {飞船系：碰前、碰后路线分行示意};
 \draw[solforce] (1.6,-1.6)--(3,-1.6) node[midway,below=4pt,sollabel] {$F=2\rho Sv^2$};
 \node[sollabel] at (2.55,.45) {迎流面};
\end{tikzpicture}
\diagcaption{解图}{飞船系中的入射与原路反弹分别列示；飞船所受反作用力沿来流方向。}\label{ch07:fig:soldiagThreeNine}\end{minipage}\end{center}}

\newcommand{\soldiagThreeTen}{%
\begin{center}\begin{minipage}{.97\linewidth}\centering
\begin{tikzpicture}[scale=.8]
 \coordinate (T) at (-1,0);
 \fill[SoftGray] (-4,1.2)--(T)--(-4,-1.2)--cycle;
 \draw[solgeom] (-4,-1.2) arc (-90:90:.25 and 1.2);
 \draw[solbody] (-4,1.2) arc (90:270:.25 and 1.2);
 \draw[solbody] (-4,1.2)--(T)--(-4,-1.2);
 \draw[solgeom] (-4.4,0)--(-.65,0);
 \foreach \y in {-.7,.7}{\draw[solflow] (-.1,\y)--({-1-abs(\y)/.4},\y);}
 \draw[<->,TextGray] (-4,-1.65)--(-1,-1.65) node[midway,below=4pt,sollabel] {$L$};
 \draw[<->,TextGray] (-4.55,0)--(-4.55,1.2) node[midway,left=3pt,sollabel] {$R$};
 \node[solpanel] at (-2.5,2.6) {空间外形：平行来流};
 \begin{scope}[shift={(1.4,0)}]
  \pgfmathsetmacro{\ang}{atan(.4)}
  \coordinate (P) at (3,0); \coordinate (Q) at (1.2,.72);
  \draw[solbody] (0,1.2)--(P)--(0,-1.2);
  \draw[solgeom] (-.2,0)--(3.5,0);
  \draw[solflow] (3.3,.72)--(Q) node[pos=.35,above=4pt,sollabel] {$\vec v_{\rm in}$};
  \draw[solflow] (Q)--++({180-2*\ang}:1.4) node[above left=3pt,sollabel] {$\vec v_{\rm out}$};
  \draw[solgeom,->] (Q)--++({90-\ang}:1.15) node[above right=3pt,sollabel] {$\hat{\boldsymbol n}$};
  \draw[TextGray] (2.35,0) arc (180:{180-\ang}:.65);
  \node[sollabel] at (2.15,.15) {$\alpha$};
  \node[sollabel] at (1.5,-1.7) {$R=L\tan\alpha$};
 \end{scope}
 \node[solpanel] at (2.9,2.6) {轴截面：法向分量反向};
\end{tikzpicture}
\diagcaption{解图}{左图说明 $L$、$R$ 与来流的空间关系；右图取包含轴线和碰撞点的截面，显示镜面反射。}\label{ch07:fig:soldiagThreeTen}\end{minipage}\end{center}}

\newcommand{\soldiagThreeEleven}{%
\begin{center}\begin{minipage}{.97\linewidth}\centering
\begin{tikzpicture}[x=5.3cm,y=5.7cm]
  \draw[axis] (0,0)--(1.13,0) node[right,sollabel] {$v/v_0$};
  \draw[axis] (0,0)--(0,.19) node[above,sollabel] {$P/(2\rho Sv_0^3)$};
  \draw[curve,domain=0:1,samples=70,smooth,variable=\x] plot ({\x},{\x*(1-\x)^2});
  \draw[solgeom] (1/3,0)--(1/3,4/27);
  \fill[DeepBlue] (1/3,4/27) circle (.55pt);
  \node[sollabel,below] at (1/3,0) {$1/3$};
  \node[sollabel,above right] at (1/3,4/27) {最大功率};
  \node[sollabel,below] at (1,0) {$1$};
\end{tikzpicture}
\diagcaption{解图}{无量纲功率 $x(1-x)^2$ 在 $x=v/v_0=1/3$ 处取得最大值。}\label{ch07:fig:soldiagThreeEleven}\end{minipage}\end{center}}

\newcommand{\soldiagThreeTwelve}{%
\begin{center}\begin{minipage}{.97\linewidth}\centering
\begin{tikzpicture}[scale=.8]
 \coordinate (O) at (-1.4,0); \coordinate (S) at (1.4,0);
 \draw[solgeom] (O) circle (2.8);
 \draw[solbody,fill=SoftGray] (O) circle (.45);
 \node[sollabel] at (-1.4,0) {太阳};
 \foreach \a in {-18,0,18}{\draw[solflow] ($(O)+(\a:.5)$)--($(O)+(\a:2.8)$);}
 \draw[solbody,line width=1.3pt] (1.4,-.4)--(1.4,.4);
 \draw[solforce] (1.4,0)--(2.65,0) node[midway,above=4pt,sollabel] {$F$};
 \node[sollabel,anchor=west] at (1.65,-.7) {小接收面 $S$，正入射};
 \draw[solgeom] (-1.4,-.5)--(-1.4,-1.4) (1.4,-.4)--(1.4,-1.4);
 \draw[<->,TextGray] (-1.4,-1.35)--(1.4,-1.35) node[midway,below=4pt,sollabel] {$r$};
 \node[sollabel] at (-1.3,3.3) {虚线球面的面积 $4\pi r^2$};
\end{tikzpicture}
\diagcaption{解图}{距离从太阳中心量起，接收面中心落在半径 $r$ 的球面上。}\label{ch07:fig:soldiagThreeTwelve}\end{minipage}\end{center}}


前面已经学习了牛顿运动定律。运动物体通常可以近似看成质点，并用给定时刻的速度描述
它的运动状态。然而，速度相同而质量不同的物体，其运动状态并不相同。为了同时刻画质量
和速度的影响，定义运动物体的\textbf{动量}为
\begin{equation}
 \vec p=m\vec v.
\end{equation}
日常经验也说明了质量不能被忽略：同样以约 $5\,\mathrm{m/s}$ 迎面而来，一个人和一辆
卡车所代表的运动状态显然完全不同。动量把质量和速度同时纳入同一个物理量中。

\textbf{动量是矢量。}\quad 质量是标量、速度是矢量，因此动量也是矢量，其方向与速度
方向相同。改变速度大小或改变速度方向，都会使动量发生变化。

\textbf{动量的变化。}\quad 若质点在 $t$ 时刻的动量为 $\vec p(t)$，在
$t+\Delta t$ 时刻的动量为 $\vec p(t+\Delta t)$，则这段时间内的动量变化量为
\begin{equation}
 \Delta\vec p=\vec p(t+\Delta t)-\vec p(t).
\end{equation}
这里的减法是矢量减法；即使速率不变，只要运动方向改变，$\Delta\vec p$ 也不为零。

\textbf{系统的合动量。}\quad 对于由 $N$ 个质点组成的系统，总动量定义为各质点动量
的矢量和：
\begin{equation}
 \vec P=\sum_{i=1}^{N}\vec p_i=\sum_{i=1}^{N}m_i\vec v_i.
\end{equation}

质量保持不变的物体受到合外力 $\vec F$ 时，由牛顿第二定律
\begin{equation}
 \vec F=m\vec a=m\frac{\dd\vec v}{\dd t}
 =\frac{\dd(m\vec v)}{\dd t}
 =\frac{\dd\vec p}{\dd t}.
\end{equation}
可见，力也可以表达为动量的变化率。对恒力作用的有限时间间隔，上式可写为
\begin{equation}
 \vec F\,\Delta t=\Delta\vec p.
\end{equation}
力与作用时间的乘积称为\textbf{冲量}，通常用 $\vec I$ 表示。对于变力，应把各微小时间
间隔内的冲量累积起来，于是一般形式的动量定理为
\begin{equation}
 \vec I=\int_{t_i}^{t_f}\vec F\,\dd t
 =\Delta\vec p=\vec p_f-\vec p_i.
\end{equation}
动量定理是牛顿运动定律的自然推论。它不要求知道力在每一时刻造成的加速度；只要能够
求出总冲量，就可以直接联系过程的初、末状态。对有些过程，它和逐时刻使用牛顿定律会
得到相同结果；而在碰撞等作用力复杂、作用时间极短的问题中，动量定理能显著简化处理，
甚至可以在不知道碰撞力随时间变化规律的情况下求出结果。

\subsection{单个物体的动力学}

当问题中只包含单个物体时，若过程持续时间或合力的总冲量容易确定，使用动量定理往往
可以省去分段求速度、位移的步骤。

\subsubsection*{冲量把短暂作用与动量变化联系起来}

牛顿第二定律写成 $\dd\vec p/\dd t=\vec F$ 后，在 $t_1$ 到 $t_2$ 上积分得
\begin{equation}
 \Delta\vec p=\vec p_2-\vec p_1
 =\int_{t_1}^{t_2}\vec F\,\dd t\equiv\vec J.
\end{equation}
冲量是力对时间的累积，方向由矢量积分决定；在 $F$--$t$ 图上，每个分量的冲量等于相应的
有向面积。碰撞力可能很大且随时间快速变化，但只要碰撞前后的动量可测，便无需知道力的每个
瞬时值。平均力 $\bar{\vec F}=\vec J/\Delta t$ 只保证总冲量相同，并不表示真实力恒定。

对质点系逐个应用动量定理并相加，成对内力因第三定律消去：
\begin{equation}
 \vec P_2-\vec P_1=\int_{t_1}^{t_2}\vec F_{\rm ext}\,\dd t.
\end{equation}
碰撞、爆炸等短过程里，重力等有限外力的冲量常比接触冲量小得多，可以在这一时间尺度上
忽略；这不是说外力不存在，而是做了量级比较。若某方向外冲量不可忽略，只有其余方向的动量
分量可能守恒。把过程时间无限放大或把很小的外力累积很久，原近似就会失效。

连续介质问题可用“单位时间通过控制面的动量”理解。质量流率为 $\dot m=\rho Av$ 的稳定射流，
若速度由 $\vec v_1$ 变为 $\vec v_2$，维持该变化所需合力为
\begin{equation}
 \vec F=\dot m(\vec v_2-\vec v_1).
\end{equation}
管壁所受力方向与此相反。弯管、喷嘴、雨滴落入车厢和水柱冲击挡板看似不同，本质都是在选定
时间内统计流入、流出和留在系统内的动量。画清控制面和速度方向，比死记某个压力公式可靠。


\begin{origexamplebox}{}\label{ch07:ex:s07-3-1}
如图，铁块压着一纸条放在水平桌面上。当以速度 $v$ 抽出纸条后，铁块掉在地上的 $P$ 点。
若以速度 $2v$ 抽出纸条，则铁块落地点 $O$ 在何处？
\begin{enumerate}[label=\Alph*.,leftmargin=2.2em]
 \item 仍在 $P$ 点
 \item $P$ 点左边
 \item $P$ 点右边不远处
 \item $P$ 点右边、原水平位移的两倍处
\end{enumerate}
图~\ref{ch07:fig:probdiagThreeOne}中，纸条从桌边抽出，铁块随后从桌边落下。
\probdiagThreeOne
\end{origexamplebox}

\Needspace{6\baselineskip}\begin{xesolution}
纸条相对铁块滑动时，铁块受到向右的滑动摩擦力。设纸条从铁块下方完全抽出所需的接触
时间为 $\tau$，则铁块在这段时间内获得的水平冲量为
\[
 I_x=f\tau,
 \qquad
 v_x=\frac{f\tau}{m}.
\]
抽纸速度增大，纸条与铁块的接触时间变短，铁块得到的摩擦冲量和离开桌面时的水平速度
都减小。桌面高度不变，铁块离开桌面后的下落时间相同，所以第二次的水平射程更小。
图~\ref{ch07:fig:soldiagThreeOne}中，抽得越快，摩擦力作用越短，离桌速度和水平射程越小。
\soldiagThreeOne
因此 $O$ 在 $P$ 点左边，选择
\[
 \boxed{\text{B}}.
\]
这里比较的是冲量而不是摩擦力大小；两次滑动摩擦力可以相同，但作用时间不同。
\end{xesolution}

\begin{origexamplebox}{}\label{ch07:ex:s07-3-2}
一个质量 $m=2\,\mathrm{kg}$ 的物体，在 $F_1=8\,\mathrm N$ 的水平力作用下从静止开始沿
水平面运动 $t_1=5\,\mathrm s$。随后外力变为 $F_2=5\,\mathrm N$，方向不变，物体又运动
$t_2=4\,\mathrm s$；撤去外力后，物体再经过 $t_3=6\,\mathrm s$ 停下。求物体所受的
摩擦力。

\begin{qualityfigure}
\begin{tikzpicture}[>=Latex,font=\small]
\draw[->](0,0)--(6.8,0)node[right]{$t/\mathrm s$};\draw[->](0,0)--(0,2.8)node[above]{$F/\mathrm N$};
\draw[DeepBlue,very thick](0,2.4)--(2,2.4)--(2,1.5)--(3.6,1.5)--(3.6,0)--(6,0);
\foreach\x/\t in {0/0,2/5,3.6/9,6/15}{\draw(\x,.08)--(\x,-.08)node[below]{$\t$};}
\node[left]at(0,2.4){$8$};\node[left]at(0,1.5){$5$};\node at(4.8,1.4){末速度为零};
\end{tikzpicture}
\figcap{两段推力和撤力阶段的时间边界}\label{mech:07:2:q}
\end{qualityfigure}
\end{origexamplebox}

\Needspace{6\baselineskip}\begin{xesolution}题图~\ref{mech:07:2:q}中，横轴把全过程分成 $5,4,6\,\mathrm s$ 三段。摩擦力在三个运动阶段都作负冲量，外加推力则仅在前两段存在。

取运动方向为正，设滑动摩擦力大小为 $f$。物体初、末速度都为零，所以全过程的总动量
变化为零。各力的冲量按时间分段相加：
\[
 F_1t_1+F_2t_2-f(t_1+t_2+t_3)=0.
\]
力与时间的对应关系见图~\ref{ch07:fig:soldiagThreeTwo}。

\soldiagThreeTwo
代入数据得
\[
 8\times5+5\times4-f(5+4+6)=0,
\]
故
\[
 \boxed{f=4\,\mathrm N}.
\]
在物体始终向右滑动的这 $15\,\mathrm s$ 内，摩擦力方向始终向左；撤去外力并不改变
摩擦力方向。
\end{xesolution}

\subsection{多个物体的运动}

若一个力学系统含有多个物体，每个物体除了受到系统外部的力，还受到其他物体施加的
内力。对第 $i$ 个物体，有
\begin{equation}
 \vec F_i=\vec F_i^{\rm ext}+\sum_{j\ne i}\vec f_{ij},
\end{equation}
其中 $\vec f_{ij}$ 表示第 $j$ 个物体对第 $i$ 个物体的力。由牛顿第三定律
\begin{equation}
 \vec f_{ij}=-\vec f_{ji}.
\end{equation}
在一段时间 $[t_i,t_f]$ 内，对第 $i$ 个物体应用动量定理：
\begin{equation}
 \int_{t_i}^{t_f}\vec F_i(t)\,\dd t=\Delta\vec p_i.
\end{equation}
对系统内全部质点求和，内力冲量两两抵消，得到系统的动量定理
\begin{equation}
 \int_{t_i}^{t_f}\left(\sum_{i=1}^{N}\vec F_i^{\rm ext}(t)\right)\dd t
 =\Delta\left(\sum_{i=1}^{N}\vec p_i\right).
\end{equation}
因此，系统总动量的变化只由外力的总冲量决定，与内力的具体形式无关。

\Needspace{8\baselineskip}
\begin{origexamplebox}{}\label{ch07:ex:s07-3-3}
两个质量分别为 $m_1$ 和 $m_2$ 的物体由劲度系数为 $k$、原长为 $l_0$ 的弹簧连接，
初始弹簧状态未知。初始时整个系统竖直放置，随后由静止释放。已知下落时间为 $t$ 时，
质量为 $m_1$ 的物体向下速度为 $v_1$，求另一物体的速度 $v_2$。
图~\ref{ch07:fig:probdiagThreeThree}中，弹簧初态未知的二物体系统由静止竖直下落。
\probdiagThreeThree
\end{origexamplebox}

\Needspace{6\baselineskip}\begin{xesolution}
以两物体（弹簧质量忽略）为系统，弹簧对两物体的作用力是一对内力，初始伸长量、劲度
系数和原长都不影响系统总动量。
取竖直向下为正，时间 $t$ 内外力总冲量只有重力冲量：
\[
 (m_1+m_2)gt=m_1v_1+m_2v_2.
\]
外力冲量与末态动量的对应关系见图~\ref{ch07:fig:soldiagThreeThree}。

\soldiagThreeThree
所以
\[
 \boxed{v_2=\frac{(m_1+m_2)gt-m_1v_1}{m_2}}.
\]
若计算结果为负，表示第二个物体此刻实际向上运动；因此这里的 $v_2$ 是以向下为正的
速度分量，而不只是速度大小。
\end{xesolution}

\Needspace{8\baselineskip}
\begin{origexamplebox}{}\label{ch07:ex:s07-3-4}
如图，$A$、$B$ 两小物体被平行于斜面的轻细线相连，均静止于斜面上。以平行于斜面
向上的恒力拉 $B$，使两物体同时由静止起以加速度 $a$ 沿斜面向上运动。经过时间 $t_1$，
细线突然被拉断；再经过时间 $t_2$，$A$ 上滑到最高点。已知 $A$、$B$ 的质量分别为
$m_A$ 和 $m_B$，断线后拉 $B$ 的恒力保持不变，求此时 $B$ 的速度。
图~\ref{ch07:fig:probdiagThreeFour}中，$A$、$B$ 沿斜面上滑，经过 $t_1$ 后连接线断裂。
\probdiagThreeFour
\end{origexamplebox}

\Needspace{6\baselineskip}\begin{xesolution}
取沿斜面向上为正。把除绳张力以外、分别作用于 $A$、$B$ 的恒定合力记为 $X_A$、$X_B$。
断线前两物体的共同加速度为 $a$，把两式相加可得
\begin{equation}\label{ch07:eq:connected-force}
 X_A+X_B=(m_A+m_B)a.
\end{equation}
断线瞬间 $A$ 的速度为 $at_1$；它在随后 $t_2$ 内减速到零，故
\begin{equation}\label{ch07:eq:broken-force}
 X_A t_2=m_A(0-at_1),
 \qquad
 X_A=-\frac{m_Aat_1}{t_2}.
\end{equation}
两阶段的速度变化见图~\ref{ch07:fig:soldiagThreeFour}。

\soldiagThreeFour
由式\eqref{ch07:eq:connected-force}、\eqref{ch07:eq:broken-force}得
\[
 X_B=(m_A+m_B)a+\frac{m_Aat_1}{t_2}.
\]
因此 $B$ 在断线后再运动 $t_2$ 的末速度为
\begin{align*}
 v_B&=at_1+\frac{X_B}{m_B}t_2\\
 &=\boxed{\frac{m_A+m_B}{m_B}\,a(t_1+t_2)}.
\end{align*}
推导只使用了两段运动中非绳力保持恒定这一条件，斜面倾角、拉力和阻力的具体数值都已
被消去。
特别地，断线后 $a_B=X_B/m_B>a$。因此解图中 $B$ 的速度折线必须比断线前更陡；
$A$ 的速度则由 $at_1$ 降至零。速度图只表示这一上滑阶段，不延伸到 $A$ 反向下滑以后。
\end{xesolution}

\Needspace{8\baselineskip}
\begin{origexamplebox}{}\label{ch07:ex:s07-3-5}
光滑水平面上，一个箱子 $A$ 的质量 $m_A=2\,\mathrm{kg}$，以初速度
$v_0=10\,\mathrm{m/s}$ 冲向静止的箱子 $B$，$B$ 的质量 $m_B=1\,\mathrm{kg}$。碰撞后：
\begin{enumerate}[label=\arabic*.,leftmargin=2.2em]
 \item 如果 $A$ 的速度变为 $5\,\mathrm{m/s}$，则 $B$ 的速度为多少？
 \item 如果 $A$、$B$ 之间涂有强力胶水，一碰就粘在一起，则合体后的速度为多少？
\end{enumerate}
图~\ref{ch07:fig:probdiagThreeFive}中，箱子 $A$ 在光滑面上冲向静止的箱子 $B$。
\probdiagThreeFive
\end{origexamplebox}

\Needspace{6\baselineskip}\begin{xesolution}
碰撞时间很短，系统水平方向外力冲量为零，因此两箱总动量在碰撞前后相同。
图~\ref{ch07:fig:soldiagThreeFive}中，碰撞内力冲量相消，系统总动量由碰前原样传到碰后。
\soldiagThreeFive
（1）取 $A$ 初速度方向为正：
\[
 m_Av_0=m_Av_A+m_Bv_B.
\]
代入数据，
\[
 2\times10=2\times5+1\times v_B,
 \qquad
 \boxed{v_B=10\,\mathrm{m/s}}.
\]

（2）粘合后两箱具有共同速度 $v$：
\[
 m_Av_0=(m_A+m_B)v,
\]
所以
\[
 \boxed{v=\frac{20}{3}\,\mathrm{m/s}}.
\]
两问都只用动量关系；“粘在一起”说明碰撞是完全非弹性的，并不意味着动能守恒。
\end{xesolution}

\subsection{冲击和碰撞}

若始终知道力的大小和方向，动量定理与牛顿运动定律当然可以给出相同结果。但在碰撞中，
作用时间极短，碰撞力极大且随时间的变化很复杂，逐时刻测量碰撞力十分困难；相比之下，
碰撞前后的速度却容易确定。因此，碰撞问题通常直接研究冲量和动量变化。

在碰撞或冲击的极短时间内，重力、普通摩擦力、弹簧力等有限外力的冲量与碰撞力冲量
相比可以忽略。相互碰撞的两个物体所受冲量大小相等、方向相反。根据系统的连接和约束，
还可以进一步确定冲量的方向以及碰撞后的瞬时速度关系。

\Needspace{8\baselineskip}
\begin{origexamplebox}{}\label{ch07:ex:s07-3-6}
如图，三个小球 $A$、$B$、$C$ 的质量均为 $m$，放在光滑水平桌面上。三球之间用质量
可忽略、不可伸长的绳 $AB$、$BC$ 相连，$\angle ABC=120^\circ$。现沿 $BA$ 绳方向给
$A$ 一个瞬时冲量 $I$，求 $C$ 球获得的速度。
图~\ref{ch07:fig:probdiagThreeSix}中，冲量 $I$ 沿 $BA$ 方向；两根绳在冲击后仍不可伸长。
\probdiagThreeSix
\end{origexamplebox}

\Needspace{6\baselineskip}\begin{xesolution}
设绳 $AB$、$BC$ 对各球传递的张力冲量大小分别为 $J_1$、$J_2$。取 $x$ 轴向右，令
$\hat{\boldsymbol u}$ 为从 $B$ 指向 $C$ 的单位矢量，则
$\hat{\boldsymbol x}\mathbin{\cdot}\hat{\boldsymbol u}=\cos60^\circ=1/2$。
各球的冲量—动量关系如解图所示：
图~\ref{ch07:fig:soldiagThreeSix}中，先对三球分别列冲量方程，再用两条绳方向上的速度约束闭合方程组。
\soldiagThreeSix
\[
 m\vec v_A=-I\hat{\boldsymbol x}+J_1\hat{\boldsymbol x},
\]
\[
 m\vec v_B=-J_1\hat{\boldsymbol x}+J_2\hat{\boldsymbol u},
 \qquad
 m\vec v_C=-J_2\hat{\boldsymbol u}.
\]
绳不可伸长，所以冲击后两条绳方向上的相对速度均为零：
\[
 (\vec v_B-\vec v_A)\mathbin{\cdot}\hat{\boldsymbol x}=0,
 \qquad
 (\vec v_C-\vec v_B)\mathbin{\cdot}\hat{\boldsymbol u}=0.
\]
代入上面的速度表达式，得到
\[
 I-2J_1+\frac{J_2}{2}=0,
 \qquad
 \frac{J_1}{2}-2J_2=0.
\]
从而
\[
 J_1=\frac{8I}{15},
 \qquad
 J_2=\frac{2I}{15}.
\]
$C$ 球沿 $CB$ 方向运动，速度大小为
\[
 \boxed{v_C=\frac{J_2}{m}=\frac{2I}{15m}}.
\]
这里不能把外冲量 $I$ 简单地在两根绳方向上作静力分解；冲量分配还必须同时满足不可
伸长绳的速度约束。
\end{xesolution}

\begin{origexamplebox}{}\label{ch07:ex:s07-3-7}
军训时，战士距墙的水平距离为 $s_0$，以速度 $v_0$ 起跳，再用脚蹬墙一次，使身体转为
竖直向上的运动以继续升高。墙面与鞋底之间的摩擦因数为 $\mu$，求能使人体重心总升高
最大的起跳角度。
采用主动蹬墙的理想冲量模型：短暂接触使水平速度降为零，向上切向冲量可达到
$J_f=\mu J_N$，接触期间重力冲量忽略。这里的 $\mu>0$ 是该模型的有效摩擦上限；
若实际接触不能达到上限，须另给切向冲量或鞋底运动规律。

\begin{qualityfigure}
\begin{tikzpicture}[>=Latex,font=\small]
\draw[thick](-.2,0)--(5,0);\draw[very thick](4,0)--(4,3.6)node[above]{竖直墙};
\fill(0,0)circle(2pt)node[below]{$O$};\draw[->,DeepBlue,thick](0,0)--(1.7,1.2)node[above]{$v_0$};\draw(.8,0)arc(0:35:.8)node[midway,right]{$\theta$};
\draw[<->](0,-.7)--(4,-.7)node[midway,below]{$s_0$};\fill(4,1.7)circle(2pt)node[left]{$P$};\draw[->,Gold](3.6,1.7)--(3.6,3)node[left]{$J_f$};\draw[->,Gold](4,1.7)--(2.8,1.7)node[below]{$J_N$};
\end{tikzpicture}
\figcap{到墙前的飞行与蹬墙冲量方向}\label{mech:07:7:q}
\end{qualityfigure}
\end{origexamplebox}

\Needspace{6\baselineskip}\begin{xesolution}题图~\ref{mech:07:7:q}中，从 $O$ 到墙上接触点 $P$ 的水平距离为 $s_0$；墙的法向冲量向左消除水平速度，向上的切向冲量增加竖直速度。

设起跳角为 $\theta$。从起跳点到墙的时间为
\[
 t=\frac{s_0}{v_0\cos\theta}.
\]
蹬墙前的水平、竖直速度分别为
\[
 v_x=v_0\cos\theta,
 \qquad
 v_y=v_0\sin\theta-\frac{gs_0}{v_0\cos\theta}.
\]
为使蹬墙后水平速度变为零，墙面对人的法向冲量为 $J_N=mv_x$。在题设主动蹬墙模型中，
取向上的切向冲量 $J_f=\mu J_N$，所以离墙后的竖直速度
\[
 v_y'=v_y+\mu v_x.
\]
冲量方向和两段上升过程见图~\ref{ch07:fig:soldiagThreeSeven}。

\soldiagThreeSeven
到墙时的高度为
\[
 h_1=s_0\tan\theta-\frac{gs_0^2}{2v_0^2\cos^2\theta}.
\]
离墙后还能上升 $h_2={v_y'}^2/(2g)$。把 $v_x,v_y$ 代入并整理，得到
\[
 h=h_1+h_2
 =\frac{v_0^2}{2g}(\sin\theta+\mu\cos\theta)^2-\mu s_0.
\]
先求不考虑到墙条件时的驻点。由于
\[
 \sin\theta+\mu\cos\theta
 \le\sqrt{1+\mu^2},
\]
等号在 $\tan\theta=1/\mu$ 时成立。因此候选角为
\[
 \theta_* =\arctan\frac{1}{\mu}.
\]
\textbf{可行性检验。}\quad 令 $\eta=gs_0/v_0^2$。必须同时满足
\[
 h_1\ge0\ \Longleftrightarrow\ \sin2\theta\ge\eta,
 \qquad
 v_y'>0\ \Longleftrightarrow\
 \sin\theta+\mu\cos\theta-\frac{\eta}{\cos\theta}>0.
\]
第一式表示到墙时尚未落地；若要求在离地后蹬墙，则改取 $h_1>0$。
把 $\theta_*$ 代入，这两个条件分别成为
\[
 \eta\le\frac{2\mu}{1+\mu^2},\qquad \eta<\mu.
\]
只有同时成立，$\theta_*$ 才是蹬墙后最高点的最优角，且
\[
 h_{\max}=\frac{v_0^2}{2g}(1+\mu^2)-\mu s_0.
\]

\textbf{候选角不可行时。}\quad 用 $z=\tan\theta>0$ 把条件化为两个二次不等式：
\[
 \eta z^2-2z+\eta\le0,
 \qquad
 \eta z^2-z+\eta-\mu<0.
\]
若 $\eta>1$，第一次抛体的最大水平射程也不足以到墙，不存在题设运动。
当 $0<\eta\le1$ 时，第一个不等式给出
\[
 z_-\le z\le z_+,
 \qquad z_\pm=\frac{1\pm\sqrt{1-\eta^2}}{\eta}.
\]
第二个不等式先要求 $D=1-4\eta(\eta-\mu)>0$，再给出
\[
 w_-<z<w_+,\qquad w_\pm=\frac{1\pm\sqrt D}{2\eta}.
\]
把这两个区间与 $z>0$ 取交集，便得到完整可行区间 $\mathcal I$；交集为空时不能完成
蹬墙后向上的运动。对非空区间，目标函数为
\[
 h(z)=\frac{v_0^2}{2g}\frac{(z+\mu)^2}{1+z^2}-\mu s_0,
 \qquad
 h'(z)=\frac{v_0^2}{g}\frac{(z+\mu)(1-\mu z)}{(1+z^2)^2}.
\]
它在 $z<1/\mu$ 时增加、$z>1/\mu$ 时减少。因此将 $1/\mu$ 限制在可行区间内即可：
若它在区间内就取它；若它在区间左侧就取最小可行端点；在右侧则取最大可行端点。
若所需端点因严格不等式被排除，只能无限逼近相应高度上限，并不存在达到上限的起跳角。
例如 $\mu=2,\eta=0.9$ 时，候选角约 $26.6^\circ$ 的射程不够，而 $45^\circ$ 满足两项条件，
不能仍把候选角当作答案。

\textbf{“全程最高”与“蹬墙后最高”的区别。}\quad 若人在到墙前已经越过抛物线顶点，
还应计入此前达到的高度 $v_0^2\sin^2\theta/(2g)$。全程最高值可统一写成
\[
 H(z)=\max\left\{h(z),\frac{v_0^2}{2g}\frac{z^2}{1+z^2}\right\}.
\]
若到墙前尚未到顶点，第二项虽是未走完抛物线的顶点高度，但 $v_y\ge0$ 且切向冲量向上，
此时 $h$ 不小于第二项，故这个写法仍成立。第二项随 $z$ 单调增加，因此全程优化只须比较
上述 $h$ 的最优值与可行区间最右端的第二项；同样保留端点可达性判断。
这区分了“蹬墙后的总升高”与字面意义的“整个运动曾达到的最高点”。
\end{xesolution}

\Needspace{8\baselineskip}
\begin{origexamplebox}{}\label{ch07:ex:s07-3-8}
在光滑水平面上有两面相距为 $d$ 的平行竖直墙。一个旋转的小球撞向其中一面墙。若墙面
光滑，弹性碰撞后小球垂直于墙的速度分量反向，平行于墙的分量不变，最终击中对面墙上的
$A$ 点。若球与墙之间的摩擦因数为 $\mu$，仍假设碰后垂直于墙的速度大小不变、方向反向，
证明：小球反弹后在另一面墙上的落点与 $A$ 点的距离，同小球初始位置、初速度方向均
无关。
主问限定碰撞期间接触点始终沿同一方向滑动，且使用库仑滑动摩擦模型；再讨论若提前
停止滑动，上述结论是否仍成立。水平面始终光滑，球的运动取俯视图。
图~\ref{ch07:fig:probdiagThreeEight}中，光滑墙落点为 $A$；有摩擦时落点为 $B$。
\probdiagThreeEight
\end{origexamplebox}

\Needspace{6\baselineskip}\begin{xesolution}
设碰墙前垂直、平行于墙的速度分量大小分别为 $v_n$、$v_t$。法向速度反向而大小不变，
所以法向冲量大小为
\[
 J_N=2mv_n.
\]
若接触点在整个碰撞中保持滑动，摩擦冲量大小为
\[
 J_t=\mu J_N=2\mu mv_n,
\]
从而平行于墙的质心速度分量改变量为
\[
 |\Delta v_t|=\frac{J_t}{m}=2\mu v_n.
\]
球从一面墙飞到另一面墙的时间只由法向运动决定：
\[
 T=\frac{d}{v_n}.
\]
与光滑墙情形相比，落点沿墙方向的偏移为
\[
 AB=|\Delta v_t|T=2\mu v_n\frac{d}{v_n}
 =\boxed{2\mu d}.
\]
图~\ref{ch07:fig:soldiagThreeEight}中，左右两部分分别表示落点偏移和接触冲量，不能把冲量箭头当作轨迹。
\soldiagThreeEight
初始位置和 $v_t$ 决定光滑墙情形的基准点 $A$，但附加偏移中的 $v_n$ 恰好消去，因此
$AB$ 与初始位置、初速度方向无关。旋转决定摩擦冲量的方向，故 $B$ 位于 $A$ 的哪一侧
由接触点的相对滑动方向决定。

若滑动提前停止，设小球半径为 $R$、绕球心转动惯量为 $\mathcal J$，取沿墙向右为正，
初始接触点相对墙的切向速度为 $w_0$。切向冲量 $J$ 对质心平动贡献 $J/m$，
对接触点的转动速度贡献 $R^2J/\mathcal J$，故
\[
 w=w_0+\left(\frac1m+\frac{R^2}{\mathcal J}\right)J.
\]
达到无滑动所需的冲量为
\[
 J_{\rm stop}=-\frac{w_0}{1/m+R^2/\mathcal J}.
\]
在球对平墙、法向通过球心且此后能够保持无滑动的简化模型中，
\[
 |J|=\min\left(2\mu m v_n,\frac{|w_0|}{1/m+R^2/\mathcal J}\right),
 \qquad AB=\frac{|J|d}{mv_n}.
\]
只有第一项起限制作用时，才恢复 $AB=2\mu d$；第二项起限制作用时，偏移还依赖初始
自转和平动。切向接触弹性等更细模型不在上述库仑冲量近似内。
\end{xesolution}

\subsection{应用到流体}

动量定理还可以用来处理流体或大量微观粒子对宏观物体的作用。设介质由大量静止粒子
组成，每个粒子质量为 $m$，单位体积内有 $n$ 个粒子；截面积为 $S$ 的长条形物体以速度
$v$ 沿轴线运动，并吸附迎面遇到的全部粒子。在很短的 $\Delta t$ 内，物体扫过体积
$Sv\Delta t$，吸附粒子的总质量为 $nmSv\Delta t=\rho Sv\Delta t$。这些粒子由静止变为
与物体同速，动量变化为
\begin{equation}
 \Delta p=(nSv\Delta t)mv=nmSv^2\Delta t
 =\rho Sv^2\Delta t,
\end{equation}
其中 $\rho=nm$ 是介质密度。由动量定理和牛顿第三定律，物体受到与运动方向相反的阻力
\begin{equation}
 F=\frac{\Delta p}{\Delta t}=\rho Sv^2.
\end{equation}
可见宏观力不仅取决于介质密度、迎风面积和速度，还取决于单个粒子碰撞后的运动方式。

\begin{origexamplebox}{}\label{ch07:ex:s07-3-9}
飞船以速率 $v$ 穿过在地面系静止、密度为 $\rho$ 的稀薄粒子介质，迎风平面面积为 $S$。粒子与飞船作正碰，碰后以相同的相对速率原路反弹，忽略单次碰撞对飞船速率的改变。求飞船受到的阻力。

\begin{qualityfigure}
\begin{tikzpicture}[>=Latex,font=\small]
\draw[->](-1,-1.5)--(6,-1.5)node[right]{$x$};\draw[->,Gold](.2,2.2)--(1.8,2.2)node[right]{$v$（飞船）};
\foreach\yy in {-.7,0,.7}{\fill(5,\yy)circle(1.5pt);\draw[->,DeepBlue](4.6,\yy)--(3.6,\yy);}
\node[align=center]at(4.4,1.6){船体系入射粒子\\相对速度 $-v$};\draw[very thick](.2,-1)--(.2,1);\node[right]at(.2,0){迎风面积 $S$};
\end{tikzpicture}
\figcap{飞船参考系中粒子迎面入射平面}\label{mech:07:9:q}
\end{qualityfigure}
\end{origexamplebox}

\Needspace{6\baselineskip}\begin{xesolution}题图~\ref{mech:07:9:q}中，粒子的入射方向相对飞船向后，图中箭头表示船体系中的速度；弹性原路反弹后应先变回介质原来静止的参考系，再求每粒子的动量变化。

在飞船参考系中，粒子碰前以速度 $v$ 迎面而来，弹性反弹后以同样大小的相对速度离开；
在介质原来静止的参考系中，粒子沿飞船运动方向的末速度为 $2v$。单位时间撞上的粒子
质量仍为 $\rho Sv$，故粒子每秒获得的动量为
\[
 \frac{\Delta p}{\Delta t}=(\rho Sv)(2v)=2\rho Sv^2.
\]
动量流和反作用力方向见图~\ref{ch07:fig:soldiagThreeNine}。

\soldiagThreeNine
所以飞船所受阻力大小为
\[
 \boxed{F=2\rho Sv^2}.
\]
弹性反弹使单个粒子的动量改变量成为完全吸附情形的两倍，因此阻力也加倍。
\end{xesolution}

\begin{origexamplebox}{}\label{ch07:ex:s07-3-10}
若飞船外形可近似为半顶角为 $\alpha$、轴向长度为 $L$ 的圆锥，在密度为 $\rho$ 的静止介质中以速率 $v$ 沿轴向前进，锥面光滑且粒子作弹性镜面反射，
求飞船受到的阻力。

\begin{qualityfigure}
\begin{tikzpicture}[>=Latex,font=\small]
\draw[->](-1,-1.5)--(6,-1.5)node[right]{$x$};\draw[->,Gold](.2,2.2)--(1.8,2.2)node[right]{$v$（飞船）};
\foreach\yy in {-.7,0,.7}{\fill(5,\yy)circle(1.5pt);\draw[->,DeepBlue](4.6,\yy)--(3.6,\yy);}
\node[align=center]at(4.4,1.6){船体系入射粒子\\相对速度 $-v$};\draw[very thick](2,0)--(0,1)--(0,-1)--cycle;\draw[dashed](0,0)--(2,0);\draw(1.2,0)arc(180:153.435:.8)node[midway,left]{$\alpha$};\draw[<->](0,-2.2)--(2,-2.2)node[midway,below]{$L$};
\end{tikzpicture}
\figcap{圆锥的轴向长度与半顶角}\label{mech:07:10:q}
\end{qualityfigure}
\end{origexamplebox}

\Needspace{6\baselineskip}\begin{xesolution}题图~\ref{mech:07:10:q}中，$L$ 沿轴线量取，半顶角 $\alpha$ 在母线与轴线之间。粒子在光滑锥面作镜面反射，动量改变量先沿局部法线计算，再投影到轴向。

令 $\alpha$ 为圆锥母线与轴线的夹角（半顶角），$L$ 为轴向长度，则底面半径
\[
 R=L\tan\alpha.
\]
在飞船参考系中，粒子以轴向相对速度 $v$ 撞上圆锥面并作镜面反射。把速度沿圆锥面法线
分解。取飞船前进方向为 $+x$，入射相对速度为 $-v\hat{\boldsymbol x}$，上表面单位外法线
的轴向分量为 $n_x=\sin\alpha$。镜面反射只使法向分量反号：
\[
 \vec v_{\rm out}=\vec v_{\rm in}
 -2(\vec v_{\rm in}\cdot\hat{\boldsymbol n})\hat{\boldsymbol n}.
\]
由 $\vec v_{\rm in}\cdot\hat{\boldsymbol n}=-v\sin\alpha$，粒子的轴向速度改变量为
\[
 \Delta v_x=2v\sin^2\alpha.
\]
几何分解见图~\ref{ch07:fig:soldiagThreeTen}。

\soldiagThreeTen
单位时间进入圆锥投影面积 $\pi R^2$ 的介质质量为 $\rho v\pi R^2$，故轴向动量变化率为
\[
 F=(\rho v\pi R^2)(2v\sin^2\alpha).
\]
因此飞船所受阻力大小为
\[
 \boxed{F=2\pi\rho L^2v^2\tan^2\alpha\sin^2\alpha
 =\frac{2\pi\rho L^2v^2\sin^4\alpha}{\cos^2\alpha}}.
\]
若题设把 $\alpha$ 定义为完整顶角，应在上述结果中用 $\alpha/2$ 代替 $\alpha$；若 $L$
表示母线长，则应改用 $R=L\sin\alpha$。因此使用结论前必须先核对角度与长度的几何定义。
\end{xesolution}

\begin{origexamplebox}{}\label{ch07:ex:s07-3-11}
一帆船在静水中顺风航行，风速为 $v_0$。假设帆面是完全弹性面，风与帆面垂直。船速
多大时，风供给船的功率最大？

\begin{qualityfigure}
\begin{tikzpicture}[>=Latex,font=\small]
\draw[thick](-.5,0)--(6,0);\draw[thick](1,.4)--(3,.4)--(2.6,0)--(1.4,0)--cycle;\draw[very thick](2,.4)--(2,2.6);
\draw[->,DeepBlue](0,1.6)--(1.5,1.6)node[midway,above]{$v_0$};\draw[->,Gold](3.2,.5)--(4.8,.5)node[above]{$v$};\node at(3.7,2.2){帆面垂直来风};\node at(2,-.7){静水；风与船同向};
\end{tikzpicture}
\figcap{同向的风速、船速与帆面}\label{mech:07:11:q}
\end{qualityfigure}
\end{origexamplebox}

\Needspace{6\baselineskip}\begin{xesolution}题图~\ref{mech:07:11:q}中，风和船同向，碰帆前的相对速率应为 $v_0-v$，而传给船的地面系功率是推力乘船速 $v$。两种速度不可混用。

设船速为 $v$，则风相对帆的速度为 $u=v_0-v$。由例~\ref{ch07:ex:s07-3-9} 的弹性反射模型，风对帆的推力为
\[
 F=2\rho S(v_0-v)^2.
\]
风对船的瞬时功率为
\[
 P=Fv=2\rho S\,v(v_0-v)^2.
\]
其变化趋势见图~\ref{ch07:fig:soldiagThreeEleven}。

\soldiagThreeEleven
求导得
\[
 \frac{\dd P}{\dd v}=2\rho S(v_0-v)(v_0-3v).
\]
在 $0<v<v_0$ 的物理区间内，内部极大值对应
\[
 \boxed{v=\frac{v_0}{3}}.
\]
这里求的是给定瞬时船速下的能量传递功率最大值；若要确定真实航行的稳定速度，还需要
加入水阻等动力学条件。
\end{xesolution}

\begin{origexamplebox}{}\label{ch07:ex:s07-3-12}
近代物理认为光线由光子构成，光子动量与能量满足 $E=pc$，其中 $c$ 为光速。太阳的总
辐射功率为 $P_S$，求距离太阳中心为 $r$ 处由太阳光冲击产生的辐射压强。
将太阳近似为各向同性点光源，接收面与入射光垂直，分别讨论完全吸收和原路完全反射。

\begin{qualityfigure}
\begin{tikzpicture}[>=Latex,font=\small]
\draw[DeepBlue](0,0)circle(3.5);\fill(0,0)circle(4pt)node[below left]{太阳 $P_S$};
\draw[dashed](0,0)--(3.5,0)node[midway,below]{$r$};\draw[very thick](3.5,-.8)--(3.5,.8)node[right]{接收面 $S$};
\draw[->,Gold](.4,.12)--(3.3,.12);\draw[->,Gold](.3,.4)--(2.6,1.2);\draw[->,Gold](.3,-.4)--(2.6,-1.2);\node at(2,-2.2){各向同性辐射；接收面与入射光垂直};
\end{tikzpicture}
\figcap{各向同性辐射球面与法向接收面}\label{mech:07:12:q}
\end{qualityfigure}
\end{origexamplebox}

\Needspace{6\baselineskip}\begin{xesolution}题图~\ref{mech:07:12:q}中，同一总功率通过半径 $r$ 的球面分摊，接收面 $S$ 与径向入射垂直；吸收与反射的区别进入光子沿该法线的动量改变量。

假设太阳各向同性辐射，则半径 $r$ 的球面上单位面积接收的功率（辐照度）为
\[
 \mathcal E=\frac{P_S}{4\pi r^2}.
\]
在时间 $\Delta t$ 内，面积 $S$ 吸收的光能为 $\Delta E=\mathcal E S\Delta t$。由
$p=E/c$，完全吸收时传给表面的动量为
\[
 \Delta p=\frac{\Delta E}{c}.
\]
能量流、动量流与受力的对应关系见图~\ref{ch07:fig:soldiagThreeTwelve}。

\soldiagThreeTwelve
因此太阳光产生的\textbf{辐射压强}为
\[
 \boxed{p_{\rm rad}=\frac{F}{S}
 =\frac{P_S}{4\pi r^2c}}.
\]
完全吸收、迎光面积为 $S$ 的物体所受力为
\[
 \boxed{F_{\rm abs}=\frac{P_SS}{4\pi r^2c}}.
\]
若光沿入射方向的反方向完全反射，单个光子的动量由 $+E/c$ 变为 $-E/c$，
表面得到的动量为 $2E/c$，因而
\[
 \boxed{p_{\rm refl}=\frac{P_S}{2\pi r^2c}},\qquad
 F_{\rm refl}=p_{\rm refl}S=\frac{P_SS}{2\pi r^2c}.
\]
主问是单位面积的受力，即压强，所以不需要另给总接收面积。引入面积 $S$ 只是为了通过
能量和动量逐步推导，并在需要时得到总力。量纲为
$[P_S/(r^2c)]=\mathrm{N/m^2}$；距离加倍时压强变为四分之一。
\end{xesolution}
}

{
\graphicspath{{../../../sample/xes-original/eighth-section/}}
\section{动量守恒}
% Original problem figures and added solution diagrams for Section 3.2.
\tikzset{
  solbody/.style={draw=Ink,line width=0.8pt},
  solforce/.style={-{Latex[length=2.0mm]},draw=DeepBlue,line width=0.9pt},
  solimpulse/.style={-{Latex[length=2.4mm]},draw=DeepBlue,line width=1.25pt},
  solgeom/.style={draw=TextGray,densely dashed,line width=0.55pt},
  solflow/.style={-{Latex[length=1.8mm]},draw=Blue,line width=0.75pt},
  sollabel/.style={font=\scriptsize,fill=none,inner sep=1pt},
  solpanel/.style={font=\scriptsize\bfseries,text=TextGray}
}
\newcommand{\diagcaption}[2]{%
 \refstepcounter{figure}\par\nopagebreak[4]\smallskip
 {\centering\scriptsize\color{TextGray}图~\thefigure\quad #2\par}\smallskip}
% Shared geometry for repeated momentum problems. Captions remain at call sites.
\newcommand{\momentumRain}{%
\begin{tikzpicture}[scale=.78]
  \fill[Blue!8] (-1.2,-.65) rectangle (1.2,-.12);
  \draw[solbody] (-1.2,.7)--(-1.2,-.65)--(1.2,-.65)--(1.2,.7);
  \foreach \x in {-.9,-.45,0,.45,.9}{\draw[solflow] (\x,1.55)--(\x,.86);}
  \node[sollabel] at (0,2.05) {$q=\rho S\Delta h/\tau$，雨水水平速度为零};
  \node[sollabel] at (0,-1.15) {$m(t)=M+qt$};
  \draw[solforce] (1.2,-.38)--(2.55,-.38) node[midway,below=3pt,sollabel] {$F$};
  \draw[solflow] (1.6,.52)--(2.8,.52) node[midway,above=3pt,sollabel] {$v$};
\end{tikzpicture}}

\newcommand{\momentumPickup}{%
\begin{tikzpicture}[scale=.8]
  \draw[solbody] (-2,0)--(2.5,0);
  \draw[DeepBlue,line width=1.2pt] (-1.8,0)--(-.15,0)
    .. controls (-.05,0) and (0,.05) .. (0,.15)--(0,2.4);
  \fill[Ink] (0,2.4) circle (2pt);
  \draw[solforce] (0,2.4)--(0,3.3) node[left=3pt,sollabel] {$F$};
  \draw[solforce] (-.65,1.65)--(-.65,.7) node[midway,left=3pt,sollabel] {$\lambda xg$};
  \draw[solflow] (.65,1.1)--(.65,2.1) node[midway,right=3pt,sollabel] {$v$};
  \draw[solgeom] (0,2.4)--(1.8,2.4);
  \draw[<->,draw=TextGray] (1.8,0)--(1.8,2.4) node[midway,right=3pt,sollabel] {$x$};
  \node[sollabel] at (-1.0,-.38) {静止绳段};
  \node[sollabel] at (0,-.95) {$\dot m_{\rm mov}=\lambda v$，新增动量率 $\lambda v^2$};
\end{tikzpicture}}

\newcommand{\momentumFallingRope}{%
\begin{tikzpicture}[scale=.78]
  \draw[solbody] (-.45,3)--(.45,3);
  \foreach \x in {-.4,-.2,0,.2,.4}{\draw[TextGray] (\x,3)--++(.12,.14);}
  \draw[solgeom] (0,3)--(0,2.1);
  \draw[DeepBlue,line width=1.1pt] (0,2.1)--(0,.28)
    .. controls (0,.08) and (-.15,.08) .. (-.3,.08);
  \draw[DeepBlue,line width=1.1pt,domain=-.3:1.6,samples=120]
    plot ({\x},{.08+.08*sin((\x+.3)*720)});
  \draw[solbody] (-.75,0)--(1.9,0);
  \draw[solgeom] (-.7,3)--(0,3) (-.7,2.1)--(0,2.1);
  \draw[<->,draw=TextGray] (-.7,2.1)--(-.7,3) node[midway,left=3pt,sollabel] {$x$};
  \draw[solflow] (.55,1.8)--(.55,.9) node[midway,right=3pt,sollabel] {$v$};
  \node[sollabel,anchor=west] at (.65,2.8) {虚线：上端已下降的距离};
  \node[sollabel,anchor=west] at (1.15,1.5) {$l-x$：仍在下落的绳长};
  \node[sollabel] at (.4,-.4) {秤盘；堆积绳段长度为 $x$};
\end{tikzpicture}}

\newcommand{\momentumRopeLoad}{%
\begin{tikzpicture}[scale=.78]
  \draw[DeepBlue,line width=1.1pt] (0,2)--(0,.28)
    .. controls (0,.08) and (-.15,.08) .. (-.3,.08);
  \draw[DeepBlue,line width=1.1pt,domain=-.3:1.8,samples=120]
    plot ({\x},{.08+.08*sin((\x+.3)*720)});
  \draw[solbody] (-.65,0)--(2.05,0);
  \draw[solflow] (.55,1.7)--(.55,.85) node[midway,right=3pt,sollabel] {$v$};
  \draw[solforce] (.45,0)--(.45,-1.0) node[below=3pt,sollabel] {$\lambda xg$};
  \draw[solforce] (1.6,0)--(1.6,-1.0) node[below=3pt,sollabel] {$\lambda v^2$};
  \node[sollabel,anchor=west,align=left] at (2.7,1.15)
    {秤盘承受的向下作用\\$N=\lambda xg+\lambda v^2$\\第二项来自使绳元停止};
\end{tikzpicture}}

\newcommand{\momentumAccretion}{%
\begin{tikzpicture}[scale=.8]
  \draw[solbody,fill=SoftGray] (-.8,0) rectangle (.8,.9);
  \draw[solforce] (-.2,.45)--(-.2,-.85) node[below=3pt,sollabel] {$mg$};
  \draw[solflow] (1.25,.7)--(1.25,-.25) node[midway,right=3pt,sollabel] {$v$};
  \foreach \x in {-1.3,-.85,-.4,.05,.5,.95,1.4}{\fill[TextGray] (\x,-1.55) circle (1.5pt);}
  \node[sollabel] at (0,-1.95) {地面系中静止的介质};
  \draw[<->,TextGray] (-.8,1.3)--(.8,1.3) node[midway,above=3pt,sollabel] {截面积 $S$};
  \node[sollabel,anchor=west,align=left] at (2.3,.8)
    {$\dot m=\rho Sv$\\$\dd(mv)/\dd t=mg$\\$m\dot v=mg-\rho Sv^2$};
\end{tikzpicture}}

\newcommand{\momentumDroplet}{%
\begin{tikzpicture}[scale=.8]
  \draw[solbody,fill=Blue!8] (0,.5) circle (.7);
  \fill[Ink] (0,.5) circle (1.5pt);
  \draw[solgeom] (0,.5)--(.7,.5) node[midway,above=4pt,sollabel] {$r$};
  \draw[solforce] (0,.5)--(0,-1.0) node[below=3pt,sollabel] {$mg$};
  \draw[solflow] (1.25,.75)--(1.25,-.25) node[midway,right=3pt,sollabel] {$v$};
  \foreach \x in {-1.3,-.85,-.4,.05,.5,.95,1.4}{\fill[TextGray] (\x,-1.65) circle (1.5pt);}
  \node[sollabel] at (0,-2.05) {地面系中静止的雾滴};
  \node[sollabel,anchor=west,align=left] at (2.3,.8)
    {$S=\pi r^2$\\$\dot m=\rho_f\pi r^2v$\\$\dot r=\rho_f v/(4\rho_w)$};
\end{tikzpicture}}


\newcommand{\soldiagThreeThirteen}{%
\begin{center}\begin{minipage}{.97\linewidth}\centering
\begin{tikzpicture}[scale=.76]
  \draw[solbody] (-2.8,-.75)--(2.8,-.75);
  \fill[TextGray!28,draw=Ink] (0,-.75) ellipse (2.1 and .48);
  \fill[Ink] (0,-.18) circle (2.8pt);
  \draw[solimpulse] (0,-.12)--++(0,1.35) node[above,sollabel] {$p=mv$};
  \draw[solimpulse] (0,-.78)--++(0,-.75) node[below,sollabel] {$P_{\rm E}=-p$};
  \node[sollabel] at (1.1,-.75) {地球};
  \node[sollabel,align=left] at (3.9,.45) {$mv+M_{\rm E}V_{\rm E}=0$\\内力交换动量};
\end{tikzpicture}
\diagcaption{解图}{人获得向上动量时，地球获得大小相等、方向相反的反冲动量。}\label{ch07:fig:soldiagThreeThirteen}\end{minipage}\end{center}}

\newcommand{\soldiagThreeFourteen}{%
\begin{center}\begin{minipage}{.97\linewidth}\centering
\begin{tikzpicture}[scale=.8]
 \draw[solbody] (-3.2,0)--(3.2,0);
 \foreach \x/\m in {-2.25/1,-.75/2,.75/3,2.25/4}{
  \draw[solbody,fill=SoftGray] (\x,.25) circle (.25);
  \node[sollabel] at (\x,.85) {$m_{\m}$};}
 \node[sollabel] at (0,1.6) {$\sum_i m_i\dot x_i=0\ \Longrightarrow\ \dot x_C=0$};
 \node[sollabel] at (0,-.45) {各球仅在水平面内运动};
\end{tikzpicture}
\diagcaption{解图}{水平总动量初值为零，故质心的水平速度始终为零。}\label{ch07:fig:soldiagThreeFourteen}\end{minipage}\end{center}}

\newcommand{\probdiagThreeFifteen}{%
\begin{center}\begin{minipage}{.97\linewidth}\centering
\begin{tikzpicture}[scale=.85]
 \pgfmathsetmacro{\ell}{3} \pgfmathsetmacro{\half}{40}
 \pgfmathsetmacro{\span}{2*\ell*sin(\half)}
 \pgfmathsetmacro{\xL}{-3*\span/8} \pgfmathsetmacro{\xH}{\span/8}
 \pgfmathsetmacro{\xR}{5*\span/8} \pgfmathsetmacro{\yH}{\ell*cos(\half)}
 \coordinate (L) at (\xL,0); \coordinate (H) at (\xH,\yH); \coordinate (R) at (\xR,0);
 \draw[axis] (-2.1,0)--(3.15,0) node[right,sollabel] {$x$};
 \draw[solbody] (L)--(H)--(R);
 \foreach \p in {L,H,R}{\fill[Ink] (\p) circle (2.5pt);}
 \node[sollabel,below left=6pt] at (L) {$L,2m$}; \node[sollabel,below right=6pt] at (R) {$R,m$};
 \node[sollabel,above=6pt] at (H) {$H,m$};
 \node[sollabel,left=7pt] at ($(L)!.5!(H)$) {$l$};
 \node[sollabel,right=7pt] at ($(R)!.5!(H)$) {$l$};
 \draw[solgeom] (H)--(\xH,0);
 \draw[TextGray] ($(H)+(-130:.65)$) arc (-130:-50:.65);
 \node[sollabel] at ({\xH+.35},{\yH-.95}) {$\theta$};
 \fill[DeepBlue] (0,{\yH/4}) circle (2pt) node[left=7pt,sollabel] {$C$};
 \draw[solgeom] (0,0)--(0,{\yH/4});
 \node[sollabel,below=6pt] at (0,0) {$O$};
 \draw[solgeom] (L)--(\xL,-.9) (R)--(\xR,-.9);
 \draw[<->,TextGray] (\xL,-.85)--(\xR,-.85) node[midway,below=4pt,sollabel] {$2l\sin(\theta/2)$};
 \node[sollabel] at (.5,-1.6) {$x_C=0$；$C$ 的高度随杆张开而改变};
\end{tikzpicture}
\diagcaption{题图}{两等长轻杆在竖直平面内张开，三个球的水平质心保持不变。}\label{ch07:fig:probdiagThreeFifteen}\end{minipage}\end{center}}

\newcommand{\soldiagThreeFifteen}{%
\begin{center}\begin{minipage}{.97\linewidth}\centering
\begin{tikzpicture}[scale=.85]
 \pgfmathsetmacro{\ell}{3} \pgfmathsetmacro{\half}{40}
 \pgfmathsetmacro{\span}{2*\ell*sin(\half)}
 \pgfmathsetmacro{\xL}{-3*\span/8} \pgfmathsetmacro{\xH}{\span/8}
 \pgfmathsetmacro{\xR}{5*\span/8} \pgfmathsetmacro{\yH}{\ell*cos(\half)}
 \coordinate (L) at (\xL,0); \coordinate (H) at (\xH,\yH); \coordinate (R) at (\xR,0);
 \draw[axis] (-2.1,0)--(3.15,0) node[right,sollabel] {$x$};
 \draw[solbody] (L)--(H)--(R);
 \foreach \p in {L,H,R}{\fill[Ink] (\p) circle (2.5pt);}
 \node[sollabel,below left=6pt] at (L) {$L,2m$}; \node[sollabel,below right=6pt] at (R) {$R,m$};
 \node[sollabel,above=6pt] at (H) {$H,m$};
 \node[sollabel,left=7pt] at ($(L)!.5!(H)$) {$l$};
 \node[sollabel,right=7pt] at ($(R)!.5!(H)$) {$l$};
 \draw[solgeom] (H)--(\xH,0);
 \draw[TextGray] ($(H)+(-130:.65)$) arc (-130:-50:.65);
 \node[sollabel] at ({\xH+.35},{\yH-.95}) {$\theta$};
 \fill[DeepBlue] (0,{\yH/4}) circle (2pt) node[left=7pt,sollabel] {$C$};
 \draw[solgeom] (0,0)--(0,{\yH/4});
 \node[sollabel,below=6pt] at (0,0) {$O$};
 \draw[solgeom] (L)--(\xL,-.9) (R)--(\xR,-.9);
 \draw[<->,TextGray] (\xL,-.85)--(\xR,-.85) node[midway,below=4pt,sollabel] {$2l\sin(\theta/2)$};
 \node[sollabel] at (.5,-1.6) {$x_C=0$；$C$ 的高度随杆张开而改变};
\end{tikzpicture}
\diagcaption{解图}{等长杆给出 $x_H=(x_L+x_R)/2$，质心约束给出 $2x_L+x_H+x_R=0$。}\label{ch07:fig:soldiagThreeFifteen}\end{minipage}\end{center}}

\newcommand{\probdiagThreeSixteen}{%
\begin{center}\begin{minipage}{.97\linewidth}\centering
\begin{tikzpicture}[scale=.73]
  \draw[solbody] (-2.5,-1.55)--(2.75,-1.55);
  \foreach \x in {-2.35,-2.0,...,2.6}{\draw[TextGray,line width=.35pt] (\x,-1.55)--++(-.15,-.18);}
  \draw[solgeom] (.35,-1.55)--(.35,2.25);
  \draw[solbody,fill=white] (.10,1.95) rectangle (.60,2.35); \node[sollabel,left=6pt] at (.1,2.15) {$M$};
  \draw[solbody,fill=white] (.10,.10) rectangle (.60,.50);
  \draw[solflow] (-2.25,.30)--(.05,.30) node[midway,above,sollabel] {$m,\ v$};
  \draw[<->,draw=TextGray] (1.05,-1.55)--(1.05,2.15) node[midway,right,sollabel] {$H$};
  \draw[<->,draw=TextGray] (1.72,-1.55)--(1.72,.30) node[midway,right,sollabel] {$h$};
\end{tikzpicture}
\diagcaption{题图}{物体下落至高度 $h$ 时被水平子弹击中并粘合。}\label{ch07:fig:probdiagThreeSixteen}\end{minipage}\end{center}}

\newcommand{\soldiagThreeSixteen}{%
\begin{center}\begin{minipage}{.97\linewidth}\centering
\begin{tikzpicture}[scale=.8]
 \node[solpanel] at (-2.5,3) {碰前动量};
 \fill[Ink] (-2.3,1.6) circle (2.5pt);
 \draw[solimpulse] (-3.8,1.6)--(-2.3,1.6) node[midway,above=5pt,sollabel] {$mv$};
 \draw[solimpulse] (-2.3,1.6)--(-2.3,.35) node[midway,left=5pt,sollabel] {$Mu$};
 \node[solpanel] at (2.4,3) {碰后：$V_x>0,\ V_y>0$（向下）};
 \coordinate (Q) at (.3,2.2);
 \draw[curve,domain=0:1,samples=70] plot ({.3+2.8*\x},{2.2-.9*\x-1.3*\x*\x});
 \draw[solflow] (Q)--(1.42,2.2) node[midway,above=5pt,sollabel] {$V_x$};
 \draw[solflow] (Q)--(.3,1.84) node[left=6pt,sollabel] {$V_y$};
 \draw[solgeom] (Q)--(.3,0);
 \draw[solbody] (-.1,0)--(3.7,0);
 \draw[<->,DeepBlue] (.3,-.45)--(3.1,-.45) node[midway,below=4pt,sollabel] {$l=V_xT$};
 \fill[Ink] (3.1,0) circle (2pt);
\end{tikzpicture}
\diagcaption{解图}{碰撞用二维动量守恒求 $V_x,V_y$，随后再用抛体运动求落点。}\label{ch07:fig:soldiagThreeSixteen}\end{minipage}\end{center}}

\newcommand{\soldiagThreeSeventeen}{%
\begin{center}\begin{minipage}{.97\linewidth}\centering
\begin{tikzpicture}[scale=.82]
 \coordinate (C) at (.5,0);
 \draw[decorate,decoration={coil,segment length=5pt,amplitude=3pt},line width=.8pt] (-2.5,0)--(2.5,0);
 \fill[Ink] (-2.5,0) circle (2.5pt); \fill[Ink] (2.5,0) circle (2.5pt);
 \node[sollabel] at (-2.5,.7) {$m_1$}; \node[sollabel] at (2.5,.7) {$m_2$};
 \draw[solgeom] (C)--(.5,.6); \node[sollabel] at (.5,1) {$C$：固定材料点};
 \draw[solgeom] (-2.5,-.2)--(-2.5,-.8) (C)--(.5,-.8) (2.5,-.2)--(2.5,-.8);
 \draw[<->,TextGray] (-2.5,-.65)--(.5,-.65) node[midway,below=4pt,sollabel] {$L_1$};
 \draw[<->,TextGray] (.5,-.65)--(2.5,-.65) node[midway,below=4pt,sollabel] {$L_2$};
 \node[sollabel] at (0,-1.4) {$m_1L_1=m_2L_2$};
\end{tikzpicture}
\diagcaption{解图}{零总动量使质心固定；固定点按反质量比分割弹簧瞬时长度。}\label{ch07:fig:soldiagThreeSeventeen}\end{minipage}\end{center}}

\newcommand{\probdiagThreeEighteen}{%
\begin{center}\begin{minipage}{.97\linewidth}\centering
\begin{tikzpicture}[scale=.72]
  % Draw the bowl as a closed rigid body, not just as an isolated circular arc.
  \path[fill=SoftGray]
    (-2.65,1.10)--(-2.65,-1.35)--(2.65,-1.35)--(2.65,1.10)--
    (1.95,1.10) arc[start angle=0,end angle=-180,radius=1.95]--cycle;
  \draw[solbody] (-2.65,1.10)--(-2.65,-1.35)--(2.65,-1.35)--(2.65,1.10);
  \draw[solbody] (-2.65,1.10)--(-1.95,1.10) (1.95,1.10)--(2.65,1.10);
  \draw[solbody] (-1.95,1.10) arc[start angle=180,end angle=360,radius=1.95];
  \draw[solbody] (-3.05,-1.35)--(3.05,-1.35);
  \foreach \x in {-2.9,-2.5,...,2.9}{\draw[TextGray,line width=.35pt] (\x,-1.35)--++(-.14,-.14);}
  \fill[Ink] (1.95,1.10) circle (4pt) node[above=7pt,sollabel] {$m$};
  \node[sollabel] at (2.16,.72) {$A$}; \node[sollabel] at (2.18,-.94) {$M$};
  % The endpoint lies on the inner circular wall, so this segment is visibly a radius.
  \draw[solgeom] (0,1.10)--({-1.95/sqrt(2)},{1.10-1.95/sqrt(2)}) node[midway,above left=4pt,sollabel] {$R$};
  \fill[Ink] (0,1.10) circle (1.7pt);
\end{tikzpicture}
\diagcaption{题图}{质点从可水平移动的半球形碗右边缘由静止释放。}\label{ch07:fig:probdiagThreeEighteen}\end{minipage}\end{center}}

\newcommand{\soldiagThreeEighteen}{%
\begin{center}\begin{minipage}{.97\linewidth}\centering
\begin{tikzpicture}[scale=.92]
 % Example mass ratio M=m. Circle and trajectory share R, theta and X.
 \pgfmathsetmacro{\rad}{2} \pgfmathsetmacro{\ang}{35}
 \pgfmathsetmacro{\X}{\rad/2*(1-sin(\ang))}
 \coordinate (C) at (\X,\rad);
 \coordinate (A) at ({\X+\rad*sin(\ang)},{\rad*(1-cos(\ang))});
 \draw[axis] (-2.4,0)--(3.1,0) node[right,sollabel] {$x$};
 \draw[axis] (0,-.45)--(0,3.1) node[above,sollabel] {$y$};
 \draw[solgeom] (-2,2) arc (180:360:2);
 \draw[solbody] ({\X-2},2) arc (180:360:2);
 \draw[curve,domain=90:35,samples=70] plot ({1+sin(\x)},{2*(1-cos(\x))});
 \draw[solgeom] (C)--(A) (C)--(\X,0);
 \draw[TextGray] ($(C)+(-90:.5)$) arc (-90:-55:.5);
 \node[sollabel] at ({\X+.25},1.25) {$\theta$};
 \node[sollabel,left=6pt] at ($(C)!.6!(A)$) {$R$};
 \fill[Ink] (A) circle (2.5pt) node[right=7pt,sollabel] {$A$};
 \fill[Ink] (2,2) circle (2pt) node[above right=5pt,sollabel] {$A_0$};
 \node[sollabel,below left=5pt] at (0,0) {$O$};
 \draw[solgeom] (0,2)--(0,2.6) (C)--(\X,2.6);
 \draw[<->,TextGray] (0,2.5)--(\X,2.5) node[midway,above=4pt,sollabel] {$X$};
 \node[sollabel] at (.2,-.9) {虚线为初始碗；实线为当前碗；蓝线为对地轨迹};
\end{tikzpicture}
\diagcaption{解图}{圆心、接触点和轨迹使用同一组参数；图取 $M=m$，一般关系见推导。}\label{ch07:fig:soldiagThreeEighteen}\end{minipage}\end{center}}

\newcommand{\soldiagThreeNineteen}{%
\begin{center}\begin{minipage}{.97\linewidth}\centering
\momentumRain
\diagcaption{解图}{雨水只增加质量，不带入水平动量；外力使总水平动量按 $Ft$ 增长。}\label{ch07:fig:soldiagThreeNineteen}\end{minipage}\end{center}}

\newcommand{\soldiagThreeTwenty}{%
\begin{center}\begin{minipage}{.97\linewidth}\centering
\momentumPickup
\diagcaption{解图}{提力同时承担已提起绳段的重量，并为新进入运动段的绳元提供动量。}\label{ch07:fig:soldiagThreeTwenty}\end{minipage}\end{center}}

\newcommand{\probdiagThreeTwentyOne}{%
\begin{center}\begin{minipage}{.97\linewidth}\centering
\momentumFallingRope
\diagcaption{题图}{柔绳上端释放后自由下落，绳元依次落到秤盘上并停止。}\label{ch07:fig:probdiagThreeTwentyOne}\end{minipage}\end{center}}

\newcommand{\soldiagThreeTwentyOne}{%
\begin{center}\begin{minipage}{.97\linewidth}\centering
\momentumRopeLoad
\diagcaption{解图}{秤盘既承受已落绳段的重量，也要连续吸收下落绳元的动量。}\label{ch07:fig:soldiagThreeTwentyOne}\end{minipage}\end{center}}

\newcommand{\soldiagThreeTwentyTwo}{%
\begin{center}\begin{minipage}{.97\linewidth}\centering
\momentumAccretion
\diagcaption{解图}{静止介质进入运动物体时要被加速，形成随 $v^2$ 增长的动量流项。}\label{ch07:fig:soldiagThreeTwentyTwo}\end{minipage}\end{center}}

\newcommand{\soldiagThreeTwentyThree}{%
\begin{center}\begin{minipage}{.97\linewidth}\centering
\momentumDroplet
\diagcaption{解图}{半径增长扩大迎风面积，质量增长又改变速度；三者通过动量流彼此耦合。}\label{ch07:fig:soldiagThreeTwentyThree}\end{minipage}\end{center}}


由动量定理可知，一个力学系统在某段过程中的总动量变化，等于外力在该过程中的总冲量，
与内力无关。因此，当系统所受合外力为零时，总动量不发生变化，这就是动量守恒。它是
动量定理的直接推论，却包含十分丰富的内容：不必知道系统内部各物体间作用力的大小和
变化过程，也能把运动前后的总体状态联系起来。在实际问题中，动量守恒常有以下三种
表现形式。

\textbf{孤立系统}\quad 不受任何外力作用的力学系统称为孤立系统。孤立系统的总动量
必然守恒。动量是矢量，因此在三维空间中，动量守恒可分别给出三个方向上的独立方程；
每一时刻的三个动量分量都分别等于相应的初始分量。

\textbf{分动量守恒}\quad 严格的孤立系统在现实中很少见，但若系统在某个特定方向上的
合外力为零，则该方向上的总动量分量守恒；其他方向的动量仍可以发生变化。

\textbf{碰撞瞬间的动量守恒}\quad 在冲击或碰撞的极短时间内，碰撞力远大于重力、普通
摩擦力等外力。这些有限外力的冲量与碰撞冲量相比可以忽略，因而碰撞前后的总动量近似
守恒。这里的“守恒”只针对碰撞持续的极短过程，不能把它延伸到碰后受外力作用的整个
运动阶段。

动量守恒虽然是动量定理的特殊情形，却能够把复杂内力全部消去，是处理多物体运动、
碰撞、质心运动和变质量吸附问题的基本工具。

\subsection{总动量守恒}

自然界中的物体总与外界存在相互作用，所以一个系统能否近似看成孤立系统，取决于所研究
问题的时间尺度和精度。例如，在某些问题中地球和月球可以近似看成孤立系统；但若太阳
引力的影响不可忽略，就必须扩大研究系统或计入太阳的外力冲量。一旦系统可近似视为孤立
系统，它的总动量三个分量均守恒，而与内力的形式、大小无关。换言之，“是否孤立”不是
脱离题目精度的绝对判断，而是针对指定对象、指定过程和指定精度作出的物理近似。

\subsubsection*{守恒条件、碰撞与质心系}

系统某方向的合外力始终为零，该方向的总动量才在整个过程中保持不变。若仅知道某段时间的
总外冲量为零，只能直接推出这段时间起点与终点的动量相等，中间可以变化。例如物块在光滑
水平面内碰撞，水平方向无外力；竖直方向虽有支持力，但支持力可与重力抵消，不能因为存在
支持力就断言竖直动量不守恒。列式前必须明确系统边界、方向及所研究的时间段。

动能定义为 $K=mv^2/2$；它是标量，不是动量的另一个名称。一维碰撞通常还需要恢复系数
\begin{equation}
 e=\frac{\text{分离相对速度}}{\text{接近相对速度}}
 =\frac{v_2-v_1}{u_1-u_2},\qquad 0\leq e\leq1,
\end{equation}
与动量守恒共同确定碰后速度。$e=1$ 为完全弹性碰撞，系统动能也守恒；$e=0$ 时两物体具有相同
末速度，是完全非弹性碰撞。一般碰撞中动能减少量转化为形变、内能和声能，不能把动能不守恒
误说成能量不守恒。

质心速度 $\vec V_C=\vec P/M$。把每个质点速度分解成 $\vec v_i=\vec V_C+\vec v_i'$，其中
$\vec v_i'$ 是相对质心的速度。由定义
\[
 \sum_i m_i\vec v_i'=\sum_i m_i\vec v_i-M\vec V_C=0.
\]
逐项展开速度平方，系统动能为
\[
 T=\sum_i\frac12m_i(V_C^2+2\vec V_C\cdot\vec v_i'+{v_i'}^2)
 =\frac12MV_C^2+\vec V_C\cdot\sum_i m_i\vec v_i'+\sum_i\frac12m_i{v_i'}^2.
\]
中间交叉项为零，故
\begin{equation}
 T=\frac12MV_C^2+\sum_i\frac12m_i{v_i'}^2.
\end{equation}
第一项是整体平动动能，第二项是质心系中的相对运动动能。无外力时 $\vec V_C$ 不变，碰撞只会
重新分配或耗散第二项。这个分解能直接给出：在总动量一定时，所有部分共同以质心速度运动时
总动能最小；完全非弹性碰撞的动能损失因此最大。

爆炸与非弹性碰撞的能量转化方向不同，但不意味着一般非弹性碰撞能够严格逆演。
若爆炸时间很短，可忽略外冲量，总动量仍守恒，但内部储存的化学能或弹性
势能会使动能增加。解题时先在方便的惯性系中写动量关系，再用相对速度或能量条件闭合；在质心
系中各碎片总动量始终为零，几何关系通常最清楚。


\begin{origexamplebox}{}\label{ch07:ex:s08-3-13}
地球和所有地球上的物体可以看成近似的孤立系统，理应满足动量守恒。然而一个人从地面
跳起时，他的动量发生了改变。如何解释这一现象？

\begin{qualityfigure}
\begin{tikzpicture}[>=Latex,font=\small]
\draw[thick](-2,0)..controls(-1,.3)and(1,.3)..(2,0);\node at(0,-.8){地球 $M_E$};
\fill(0,1.5)circle(.15);\draw(0,1.35)--(0,.7)(0,1.1)--(-.3,.9)(0,1.1)--(.3,.9)(0,.7)--(-.3,.4)(0,.7)--(.3,.4);\node[right]at(.6,1.3){人 $m$};
\draw[->,Gold](1.5,.7)--(1.5,1.9)node[right]{$v$};\draw[dashed](-2.5,-1.2)rectangle(2.5,2.3);\node at(0,2.7){研究系统：人和地球};
\end{tikzpicture}
\figcap{人和地球的系统边界}\label{mech:07:13:q}
\end{qualityfigure}
\end{origexamplebox}

\Needspace{6\baselineskip}\begin{xesolution}题图~\ref{mech:07:13:q}中，虚线把人和地球同时包括在系统内。地面给人的支持力在单独研究人时是外力，在这个完整系统内则属于内力。

只把人作为研究对象时，他受到地面的支持力冲量，因此人的动量当然可以改变。若把“人和
地球”作为完整系统，人与地面之间的作用力属于内力，系统总动量仍守恒。

设人起跳后获得向上的动量 $p=mv$。忽略外界冲量，地球必获得大小相等、方向相反的动量：
\[
 mv+M_{\rm E}V_{\rm E}=0,
 \qquad
 V_{\rm E}=-\frac{m}{M_{\rm E}}v.
\]
动量交换过程见图~\ref{ch07:fig:soldiagThreeThirteen}。

\soldiagThreeThirteen
由于地球质量 $M_{\rm E}$ 极大，它的反冲速度极小，日常观察中难以觉察；但地球的反冲
动量并不为零。人与地球之间的重力同样是系统内力，人在空中运动时总动量仍保持不变。
\end{xesolution}

\subsection{动量的分量守恒}

动量是矢量。若系统在某一方向上的合外力为零，则动量定理在该方向上给出
\[
 \Delta P_{\parallel}=I_{\parallel}^{\rm ext}=0,
\]
即该方向的总动量分量守恒。判断时必须先明确研究系统和方向；“系统受到外力”并不等于
“每个方向的动量都不守恒”。


\paragraph{全程守恒与首末动量相同的区别}
选定固定惯性系中的 $x$ 方向后，$\dd P_x/\dd t=F_x^{\rm ext}$。只有该方向合外力在整个过程中都为零，才能说 $P_x$ 全程守恒。若合外力并不始终为零，只是 $\int_{t_1}^{t_2}F_x^{\rm ext}\dd t=0$，则只能推出首末两时刻的分量相同，不能推出中间各时刻都不变。
图~\ref{quality:07:component}把两个相互作用物体一起选为系统，水平方向的相互作用力是内力；竖直方向的支持力和重力仍是外力。由哪个方向的外冲量可忽略来选择守恒分量，而不是看到碰撞就把两个方向同时写成守恒。若选的是随物体转动的轴，还会有基矢变化项，不能直接把固定轴的分量定理照搬过去。
\begin{qualityfigure}\begin{tikzpicture}[>=Latex]
 \draw[ground] (-.5,0)--(6.5,0);\draw (0,0) rectangle (1.5,1);\draw (4,0) rectangle (5.5,1);
 \draw[guide] (-.3,-.3) rectangle (5.8,1.35);
 \draw[vector] (1.5,.5)--node[above] {$F_{21}$}(2.7,.5);
 \draw[vector] (4,.5)--node[above] {$F_{12}$}(2.8,.5);
 \draw[->] (.75,1)--node[left] {$N_1$}(.75,2);
 \draw[->] (4.75,1)--node[right] {$N_2$}(4.75,2);
 \draw[->] (.75,.5)--node[left] {$m_1g$}(.75,-1);
 \draw[->] (4.75,.5)--node[right] {$m_2g$}(4.75,-1);
 \node[below] at (2.7,-1.3) {整体系统；光滑水平面上 $F_x^{\rm ext}=0$};
 \draw[axis] (6,.4)--(7,.4) node[right] {$x$};
\end{tikzpicture}
\figcap{内力在整体动量方程中相消；支持力与重力仍属于竖直方向的外力。}\label{quality:07:component}
\end{qualityfigure}


\begin{origexamplebox}{}\label{ch07:ex:s08-3-14}
水平面上有多个物体构成的系统，它们与水平面之间均无摩擦，且始终被约束在该水平面内运动。
除重力和水平面的支持力外，物体只受系统内部的相互作用力，不受其他外力。开始时所有物体
均处于静止状态。证明：在此后的运动过程中，系统质心位置保持不变。

\begin{qualityfigure}
\begin{tikzpicture}[>=Latex,font=\small]
\draw[thick](-.5,0)--(6,0);\fill(1,.15)circle(3pt)node[above]{$m_1$};\fill(4.5,.15)circle(3pt)node[above]{$m_2$};\node at(2.7,-.7){水平面光滑};\draw[->,Gold](1,.7)--(2.3,.7)node[above]{$\vec F_{12}$};\draw[->,Gold](4.5,.7)--(3.2,.7)node[above]{$\vec F_{21}$};\node at(3,1.7){所有初速度为零；仅有水平内力};
\end{tikzpicture}
\figcap{光滑平面内仅有系统内力的质点组}\label{mech:07:14:q}
\end{qualityfigure}
\end{origexamplebox}

\Needspace{6\baselineskip}\begin{xesolution}题图~\ref{mech:07:14:q}中，相互作用箭头成对出现，水平方向的总外力为零；所有初速度为零使初始总动量也为零，二者共同决定质心运动。

设各物体质量、水平坐标分别为 $m_i,x_i$，总质量为 $M=\sum_i m_i$。系统在水平方向
不受外力，且初始总动量为零，因此
\[
 P_x=\sum_i m_i\dot x_i=0.
\]
质心水平坐标为
\[
 x_C=\frac{\sum_i m_ix_i}{M},
\]
所以
\[
 \dot x_C=\frac{1}{M}\sum_i m_i\dot x_i
 =\frac{P_x}{M}=0.
\]
图~\ref{ch07:fig:soldiagThreeFourteen}中，水平总动量初值为零，故质心的水平速度始终为零。
\soldiagThreeFourteen
由此得到
\[
 \boxed{x_C=\text{常量}}.
\]
重力和支持力是竖直方向的外力，它们不影响水平方向的分动量守恒。题目所说的“质心位置”
应理解为水平面内的质心位置；竖直坐标由平面约束保持不变。
\end{xesolution}

\Needspace{8\baselineskip}
\begin{origexamplebox}{}\label{ch07:ex:s08-3-15}
如图，两根长度均为 $l$、质量可忽略的轻杆，一端通过质量为 $m$ 的球形铰链连接，另一端
分别接有质量为 $m$ 和 $2m$ 的小球。将两杆并拢，铰链向上放在桌上，然后轻敲一下使两
小球向两边滑动。两杆始终在竖直平面内，所有摩擦均忽略。建立图示坐标系，初始时铰链
位于 $x=0$，求两杆夹角为 $90^\circ$ 以及完全展开成水平状态时，质量为 $2m$ 的小球位置。
图~\ref{ch07:fig:probdiagThreeFifteen}中，两等长轻杆在竖直平面内张开，三个球的水平质心保持不变。
\probdiagThreeFifteen
\end{origexamplebox}

\Needspace{6\baselineskip}\begin{xesolution}
记左侧 $2m$ 小球、铰链和右侧 $m$ 小球的横坐标分别为 $x_L,x_H,x_R$。系统水平方向
无外力，且初始静止、质心在 $x=0$，所以始终有
\begin{equation}\label{ch07:eq:rods-com}
 2m x_L+m x_H+m x_R=0.
\end{equation}
两杆等长，铰链总在两个下端点连线的中垂线上，因此
\begin{equation}\label{ch07:eq:rods-midpoint}
 x_H=\frac{x_L+x_R}{2}.
\end{equation}
若两杆夹角为 $\theta$，两个下端点间距为
\begin{equation}\label{ch07:eq:rods-span}
 x_R-x_L=2l\sin\frac{\theta}{2}.
\end{equation}
几何关系和质心约束见图~\ref{ch07:fig:soldiagThreeFifteen}。

\soldiagThreeFifteen
把式\eqref{ch07:eq:rods-midpoint}代入式\eqref{ch07:eq:rods-com}并除以 $m/2$，得
$5x_L+3x_R=0$，即 $x_R=-5x_L/3$。再代入式\eqref{ch07:eq:rods-span}，有
$-8x_L/3=2l\sin(\theta/2)$，得到
\[
 x_L=-\frac{3l}{4}\sin\frac{\theta}{2}.
\]
因此，夹角为 $90^\circ$ 时
\[
 \boxed{x_L=-\frac{3\sqrt2}{8}l},
\]
完全展开成水平状态，即 $\theta=180^\circ$ 时
\[
 \boxed{x_L=-\frac34l}.
\]
负号表示质量为 $2m$ 的小球位于原点左侧。铰链并不保持在 $x=0$；由于左右端质量不同，
它会向较轻的右侧移动，以维持系统质心不动。
\end{xesolution}

\subsection{冲击、碰撞过程中的动量守恒}

碰撞时间极短时，重力等有限外力的冲量可以忽略，系统在碰撞前后的总动量近似守恒。
碰撞结束后的后续运动则仍须按正常的动力学规律处理，不能把“碰撞瞬间守恒”无限延长到
整个运动过程。

\Needspace{8\baselineskip}
\begin{origexamplebox}{}\label{ch07:ex:s08-3-16}
质量为 $M$ 的质点 $A$ 从距水平面高 $H$ 处由静止开始自由下落。当它下落到距地面高 $h$
处时，被一颗质量为 $m$、沿水平方向以速度 $v$ 飞行的子弹击中。碰撞时间极短，子弹最终
停留在 $A$ 内部。求 $A$ 的落地点与未被击中时下落点之间的距离 $l$。
图~\ref{ch07:fig:probdiagThreeSixteen}中，物体下落至高度 $h$ 时被水平子弹击中并粘合。
\probdiagThreeSixteen
\end{origexamplebox}

\Needspace{6\baselineskip}\begin{xesolution}
物体落到高度 $h$ 时，碰前竖直速度大小为
\[
 u=\sqrt{2g(H-h)}.
\]
在碰撞的极短时间内，重力冲量可以忽略。取水平向右、竖直向下为正，分别应用两方向的
动量守恒：
\[
 mv=(M+m)V_x,
 \qquad
 Mu=(M+m)V_y.
\]
所以粘合体碰后速度分量为
\[
 V_x=\frac{mv}{M+m},
 \qquad
 V_y=\frac{M}{M+m}\sqrt{2g(H-h)}.
\]
碰撞与后续抛体运动的衔接见图~\ref{ch07:fig:soldiagThreeSixteen}。

\soldiagThreeSixteen
设碰后到落地的时间为 $T$，则
\[
 h=V_yT+\frac12gT^2,
\]
取正根
\[
 T=\frac{\sqrt{V_y^2+2gh}-V_y}{g}.
\]
未被击中时物体沿原竖直线落地，因此所求距离就是碰后的水平位移：
\[
 \boxed{
 l=\frac{mv}{(M+m)g}
 \left[
 \sqrt{\frac{2M^2g(H-h)}{(M+m)^2}+2gh}
 -\frac{M}{M+m}\sqrt{2g(H-h)}
 \right]}.
\]
子弹嵌入说明碰撞是完全非弹性的，碰撞过程中机械能不守恒；本题只使用动量守恒确定
碰后初速度，再用运动学求落点。
\end{xesolution}

\section{质心系与动量分解}

考虑由 $N$ 个质点组成的系统，第 $i$ 个质点的质量和位置分别为 $m_i$、$\vec x_i$。
系统质心位置定义为
\begin{equation}
 \vec x_C=\frac{\sum_{i=1}^{N}m_i\vec x_i}{\sum_{i=1}^{N}m_i}.
\end{equation}
记总质量 $M=\sum_i m_i$，系统总动量为
\begin{equation}
 \vec P=\sum_{i=1}^{N}m_i\vec v_i
 =\frac{\dd}{\dd t}\left(\sum_{i=1}^{N}m_i\vec x_i\right)
 =M\frac{\dd\vec x_C}{\dd t}=M\vec v_C.
\end{equation}
因此，总动量也可以看成总质量与质心速度的乘积。若系统动量守恒，质心速度保持不变；
不受外力的系统质心或者静止，或者作匀速直线运动，具体由初始质心速度决定。

\begin{origexamplebox}{}\label{ch07:ex:s08-3-17}
两个质量分别为 $m_1,m_2$ 的质点由劲度系数为 $k$、原长为 $l_0$ 的轻弹簧连接。开始时
把弹簧拉长到 $l>l_0$ 后静止释放。在此后的运动中，弹簧上是否存在一个位置始终不变的点？
两质点在光滑水平直线上运动，弹簧均匀且质量忽略，不受其他水平外力。

\begin{qualityfigure}
\begin{tikzpicture}[>=Latex,font=\small]
\draw[thick](-.5,0)--(6,0);\fill(1,.15)circle(3pt)node[above]{$m_1$};\fill(4.5,.15)circle(3pt)node[above]{$m_2$};\node at(2.7,-.7){水平面光滑};\draw[decorate,decoration={coil,aspect=.5,segment length=5pt,amplitude=3pt}](1.1,.15)--(4.4,.15)node[midway,above=7pt]{$k,l_0$};\draw[<->](1,-1.3)--(4.5,-1.3)node[midway,below]{$l>l_0$（初始静止）};
\end{tikzpicture}
\figcap{两端质量、弹簧原长与初始伸长}\label{mech:07:17:q}
\end{qualityfigure}
\end{origexamplebox}

\Needspace{6\baselineskip}\begin{xesolution}题图~\ref{mech:07:17:q}中，两端质量决定质心在瞬时弹簧长度上的分点比例；在均匀理想弹簧作仿射伸缩的模型中，同一材料分点才可始终对应这处固定空间位置。

系统只受内力，且初始总动量为零，所以质心位置始终不变。设两个质点的位置为 $x_1,x_2$，
则固定点就是系统质心
\[
 x_C=\frac{m_1x_1+m_2x_2}{m_1+m_2}.
\]
把它写成沿弹簧从 $m_1$ 端起的分点形式：
\[
 x_C=x_1+\frac{m_2}{m_1+m_2}(x_2-x_1).
\]
图~\ref{ch07:fig:soldiagThreeSeventeen}中，零总动量使质心固定；固定点按反质量比分割弹簧瞬时长度。
\soldiagThreeSeventeen
因此，在理想轻弹簧各段均匀伸缩的模型下，弹簧上确有一点始终不动。该点把弹簧的瞬时
长度分成
\[
 \boxed{L_1:L_2=m_2:m_1}.
\]
这里“始终不变”指相对地面的位置不变。均匀无质量弹簧内部张力在同一时刻处处相同，
各段单位原长的伸长相同，所以同一材料点占总长度的比例不变。选取原长分数为
$m_2/(m_1+m_2)$ 的材料点，其位置恰等于 $x_C$，这才把固定质心与弹簧上的点联系起来。
若不假定均匀形变，只能由动量守恒断言质心位置固定，不能仅凭端点运动确认同一个材料点固定。
\end{xesolution}

\Needspace{8\baselineskip}
\begin{origexamplebox}{}\label{ch07:ex:s08-3-18}
质量为 $M$、半径为 $R$ 的球形碗放在光滑水平面上。质量为 $m$ 的质点 $A$ 开始时位于
碗的边缘，由静止释放。求 $A$ 在下落过程中相对地面的运动轨迹。
图~\ref{ch07:fig:probdiagThreeEighteen}中，质点从可水平移动的半球形碗右边缘由静止释放。
\probdiagThreeEighteen
\end{origexamplebox}

\Needspace{6\baselineskip}\begin{xesolution}
在地面惯性系中，取初始碗最低点为 $O$，$x$ 轴水平向右、$y$ 轴竖直向上；碗只作水平
平移，最低点的高度始终为零。设碗中心由 $(0,R)$ 移到 $(X,R)$，用 $\theta$ 表示从中心
指向质点的半径与竖直向下方向的夹角。质点相对碗的水平坐标和对地高度分别为
\[
 x_{A/M}=R\sin\theta,
 \qquad
 y_A=R(1-\cos\theta).
\]
释放点在右边缘，初态 $X=0$、$\theta=\pi/2$。系统水平方向质心不动，因此
\[
 MX+m(X+R\sin\theta)=mR.
\]
解得碗和质点的对地横坐标
\[
 X=\frac{mR}{M+m}(1-\sin\theta),
\]
\[
 x_A=X+R\sin\theta
 =\frac{R}{M+m}(m+M\sin\theta).
\]
参数与质心约束见图~\ref{ch07:fig:soldiagThreeEighteen}。

\soldiagThreeEighteen
消去 $\theta$，得到轨迹方程
\[
 \boxed{
 \left(\frac{(M+m)x_A-mR}{MR}\right)^2
 +\left(1-\frac{y_A}{R}\right)^2=1}.
\]
因此，质点相对地面的轨迹不是圆弧，而是椭圆的一段；其水平半轴为 $MR/(M+m)$，竖直
半轴为 $R$。碗质量 $M\to\infty$ 时，水平半轴趋于 $R$，轨迹才退化为近似固定圆弧。
\end{xesolution}

如果一个系统的运动变化完全由内力引起，那么系统质心的运动状态与先前相同。这一结论
不仅适用于弹簧和可移动曲面，也适用于行星系统等更复杂的多体系统。把太阳与行星作为系统，
在忽略外部天体引力的精度内，可近似使用总动量守恒。太阳的质量远大于各行星，因而它绕
系统质心的运动通常较小，入门计算常把太阳近似固定；但精细计算时不能把这个近似当作严格
事实。在质心系中 $M_\odot\vec v_\odot+\sum_i m_i\vec v_i=0$，行星动量改变时，
太阳的动量也随之改变，太阳并非完全不动。

\subsection{从封闭系统看质量交换}

汽车可能不断吸收雨水，运动物体也可能吸附周围介质。首先只做代数展开：
\begin{equation}
 \Delta(m\vec v)=\vec v\,\Delta m+m\,\Delta\vec v
 +\Delta m\,\Delta\vec v.
\end{equation}
忽略二阶小量并取连续极限，得到 $\dd(m\vec v)=\vec v\,\dd m+m\,\dd\vec v$。
这只是乘积的微分规则，\textbf{不能一般地将变质量物体的动量变化等同于外力冲量}。
设时间 $\dd t$ 内有质量 $\dd m>0$ 的物质以地面系速度 $\vec u_{\rm in}$ 被吸收。
选择“物体及将要被吸收的这部分物质”为同一批物质，动量定理给出
\[
 m\vec v+\vec u_{\rm in}\dd m+\vec F_{\rm ext}\dd t
 =(m+\dd m)(\vec v+\dd\vec v).
\]
展开、略去二阶项并除以 $\dd t$，得到
\begin{equation}\label{ch07:eq:accretion-balance}
 \frac{\dd(m\vec v)}{\dd t}
 =\vec F_{\rm ext}+\vec u_{\rm in}\frac{\dd m}{\dd t}.
\end{equation}
只有新物质在所研究方向上原来速度为零，右侧该方向的动量通量才为零。下面先研究这类
吸附问题，随后再推导同时包含流入和流出的通式。

\begin{origexamplebox}{}\label{ch07:ex:s08-3-19}
在雨中用恒力 $F$ 推一个初始质量为 $M$、截面积为 $S$ 的容器。一个静止在地面的同样
容器，单位时间内水位上升 $\Delta h$。不计雨水对容器前方的阻力，水的密度为 $\rho$。
用 $\tau$ 表示这里的一个单位时间，即静止容器在时长 $\tau$ 内水位上升 $\Delta h$。
容器从静止开始受推力，求经过 $3\tau$ 后的速度。

\begin{qualityfigure}
\begin{tikzpicture}[>=Latex,font=\small]
\draw[thick](0,1.8)--(0,0)--(3,0)--(3,1.8);\fill[PaleBlue](.05,.05)rectangle(2.95,.8);\draw[thick](-.5,-.2)--(6,-.2);
\foreach\x in {.4,1.2,2,2.8}{\draw[->,DeepBlue](\x,3)--(\x,2);}
\draw[->,Gold](3.2,.8)--(4.7,.8)node[right]{$F$};\node at(1.5,1.2){$M,S$};\node at(1.5,3.5){雨水密度 $\rho$};
\end{tikzpicture}
\figcap{竖直入雨的容器与水平推力}\label{mech:07:19:q}
\end{qualityfigure}
\end{origexamplebox}

\Needspace{6\baselineskip}\begin{xesolution}题图~\ref{mech:07:19:q}中，雨滴进入前没有水平速度，因此它们没有带入水平动量；容器内水量却不断增加，左侧必须对整个 $m(t)v(t)$ 求导。

雨水竖直下落，进入容器前的水平动量为零。雨水质量流率为
\[
 q=\frac{\rho S\Delta h}{\tau}.
\]
故运动容器在时刻 $t$ 的总质量为
\[
 m(t)=M+qt.
\]
取水平方向为正，把“容器加上已收集的雨水”作为变质量系统。外力的水平冲量全部转化为
系统水平动量：
\[
 \frac{\dd}{\dd t}[m(t)v(t)]=F.
\]
由初始静止条件积分得
\[
 [M+qt]v=Ft,
 \qquad
 v(t)=\frac{Ft}{M+qt}.
\]
质量流与水平动量的关系见图~\ref{ch07:fig:soldiagThreeNineteen}。

\soldiagThreeNineteen
在一般式中代入 $t=3\tau$ 和上述质量流率，得到
\[
 \boxed{v(3\tau)=\frac{3F\tau}{M+3\rho S\Delta h}}.
\]
若要求经过任意时长 $t$，则保留一般式
\[
 \boxed{v(t)=\frac{Ft}{M+(\rho S\Delta h/\tau)t}}.
\]
雨水虽没有带来水平动量，却不断增加需要加速的质量，所以速度小于质量恒为 $M$ 时的
$3F\tau/M$。分子 $F\tau$ 是冲量，除以质量后确为速度。若容器初速度为 $v_i$，
积分常数改为 $Mv_i$，于是 $v(t)=(Mv_i+Ft)/(M+qt)$；不能无说明地省略初始动量。
\end{xesolution}

\begin{origexamplebox}{}\label{ch07:ex:s08-3-20}
长度为 $l$、质量为 $m$ 的柔软绳放在水平桌面上。用手将绳的一端以恒定速度 $v$ 向上
提起。求提起高度为 $x$（$x<l$）时手的提力。

\begin{qualityfigure}
\begin{tikzpicture}[>=Latex,font=\small]
\draw[thick](-1,0)--(5,0);\draw[very thick](0,2.8)--(0,.2)..controls(.1,0)and(.3,0)..(.5,0)--(3.5,0);
\draw[->,DeepBlue](.4,2)--(.4,3.1)node[right]{$v$};\draw[<->](-.5,0)--(-.5,2.8)node[midway,left]{$x$};\node at(2,1){桌上绳段原来静止};
\end{tikzpicture}
\figcap{以恒速上提时绳元由静止进入运动段}\label{mech:07:20:q}
\end{qualityfigure}
\end{origexamplebox}

\Needspace{6\baselineskip}\begin{xesolution}题图~\ref{mech:07:20:q}中，桌面上的静止绳元依次从底部进入竖直运动段；手的提力既要托住这段绳的重量，又要提供新绳元的向上动量流率。

绳的线密度为 $\lambda=m/l$。当竖直段长度为 $x$ 时，运动部分质量为 $\lambda x$；绳以
速率 $v$ 从静止的桌面部分进入竖直运动部分，所以质量流率为
\[
 \frac{\dd m_{\rm mov}}{\dd t}=\lambda v.
\]
竖直运动部分的动量为
\[
 p=(\lambda x)v.
\]
由于 $\dot x=v$ 且 $v$ 恒定，
\[
 \frac{\dd p}{\dd t}=\lambda v^2.
\]
对竖直运动段列动量方程，手的提力减去该段重力等于动量增加率：
\[
 F-\lambda xg=\lambda v^2.
\]
图~\ref{ch07:fig:soldiagThreeTwenty}中，提力同时承担已提起绳段的重量，并为新进入运动段的绳元提供动量。
\soldiagThreeTwenty
因此
\[
 \boxed{F=\frac{m}{l}(gx+v^2)}.
\]
其中 $\lambda gx$ 用来平衡已提起绳段的重力，$\lambda v^2$ 用来把单位时间新提起的绳元
从静止加速到速度 $v$。
\end{xesolution}

\Needspace{8\baselineskip}
\begin{origexamplebox}{}\label{ch07:ex:s08-3-21}
如图，一根质量为 $m$、长为 $l$ 的柔软绳自由悬垂，下端恰与秤盘接触。某时刻放开柔软
绳上端，求台秤在任意时刻的读数以及读数的最大值。
图~\ref{ch07:fig:probdiagThreeTwentyOne}中，柔绳上端释放后自由下落，绳元依次落到秤盘上并停止。
\probdiagThreeTwentyOne
\end{origexamplebox}

\Needspace{6\baselineskip}\begin{xesolution}
设释放后绳上端下降距离为 $x$，落到秤盘上的绳长也为 $x$；尚在运动的竖直绳段作自由
落体，速度满足
\[
 v^2=2gx,
 \qquad
 x=\frac12gt^2
 \quad\left(0\le t<\sqrt{\frac{2l}{g}}\right).
\]
线密度 $\lambda=m/l$。秤盘读数由两部分组成：已经静止在盘上的绳段重力 $\lambda xg$，
以及单位时间使落下绳元停止所需的动量流 $\lambda v^2$。因此
\[
 N=\lambda xg+\lambda v^2
 =3\lambda gx.
\]
两项来源见图~\ref{ch07:fig:soldiagThreeTwentyOne}。

\soldiagThreeTwentyOne
在全部绳落完之前，
\[
 \boxed{N(t)=\frac{3mg^2t^2}{2l}},
 \qquad
 0\le t<\sqrt{\frac{2l}{g}}.
\]
最后一段绳刚要落到秤盘前，$x\to l$，读数达到
\[
 \boxed{N_{\max}=3mg}.
\]
理想模型中最后一段停止后，冲击过程结束，整根绳静止在秤盘上，读数回到 $mg$。因此
$3mg$ 是落绳过程中瞬时读数的上限，而不是最终静止读数。
\end{xesolution}

\begin{origexamplebox}{}\label{ch07:ex:s08-3-22}
在密度为 $\rho$ 的均匀介质中，一个截面积为 $S$ 的物体下落。物体能吸附所有与它相遇的
静止介质粒子，因而质量变化不可忽略。列出物体速度和质量随时间变化的方程，并定性分析
它的运动；不要求求解。

\begin{qualityfigure}
\begin{tikzpicture}[>=Latex,font=\small]
\foreach\x/\y in {-1/1,1/1.5,2.5/.5,-.5/-1.2,2/-1.7,3/1.5}{\fill[DeepBlue](\x,\y)circle(1.4pt);}
\draw[->,Gold](3.7,1.4)--(3.7,-1.4)node[below]{$y,v$（向下）};\draw[thick,fill=PaleBlue](0,-.5)rectangle(2,.5);\node at(1,0){$m$};\node at(1,-.9){截面积 $S$};\node at(1,2.2){介质粒子在地面系静止};
\end{tikzpicture}
\figcap{定截面下落物体对静止介质的扫掠}\label{mech:07:22:q}
\end{qualityfigure}
\end{origexamplebox}

\Needspace{6\baselineskip}\begin{xesolution}题图~\ref{mech:07:22:q}中，扫掠厚度在 $\dd t$ 内为 $v\dd t$，迎风面积为 $S$，故质量流率是 $\rho Sv$；介质入流的初始动量为零。

取竖直向下为正。物体在时间 $\dd t$ 内扫过体积 $Sv\dd t$，故吸附质量为
\[
 \dd m=\rho Sv\dd t,
 \qquad
 \boxed{\frac{\dd m}{\dd t}=\rho Sv}.
\]
介质粒子吸附前在地面参考系中静止，不携带初始动量；系统受到的外力只有重力，所以
\[
 \boxed{\frac{\dd(mv)}{\dd t}=mg}.
\]
展开得速度方程
\[
 \boxed{m\frac{\dd v}{\dd t}=mg-\rho Sv^2}.
\]
动量流关系见图~\ref{ch07:fig:soldiagThreeTwentyTwo}。

\soldiagThreeTwentyTwo
吸附静止粒子产生的 $\rho Sv^2$ 项使加速度不大于 $g$，但仅凭这一点不能断言速度总增加。
进一步取初始质量 $m_0$、初速度为零，并令 $k=\rho S$。令 $y=v^2$，在 $v>0$ 的时段消去时间：
\[
 \frac{\dd y}{\dd m}=\frac{2v\dot v}{\dot m}
 =\frac{2g}{k}-\frac{2y}{m},\qquad
 \frac{\dd}{\dd m}(m^2y)=\frac{2g}{k}m^2.
\]
从 $(m_0,0)$ 积分到 $(m,y)$，得 $m^2y=2g(m^3-m_0^3)/(3k)$，故
\[
 v^2=\frac{2g}{3k}\left(m-\frac{m_0^3}{m^2}\right).
\]
代回加速度式，$a=g/3+2g(m_0/m)^3/3>0$。所以对从静止出发的这一解，质量、速度均增加，
加速度从 $g$ 降低。当 $m$ 很大时，$v^2\sim2gm/(3k)$，速度并不趋于有限的终端值；加速度逐渐趋近
\[
 \boxed{a\to\frac{g}{3}}.
\]
这个极限补充了题设所要求的定性分析，并不是说物体作匀速运动。
若初始已有向下速度 $v_i$，则应加上初值项
$v^2=2gm/(3k)+m_0^2[v_i^2-2gm_0/(3k)]/m^2$。当 $kv_i^2>m_0g$ 时，初始阶段反而减速。
\end{xesolution}

\begin{origexamplebox}{}\label{ch07:ex:s08-3-23}
球状小水滴在静止且密度均匀的雾气中下落，并吸附全部遇到的水分子。水滴始终保持球形，
忽略空气阻力，重力加速度取恒定值 $g$。列出下落过程中水滴的质量、速度、半径所满足的
方程，并定性分析水滴运动；不要求求解运动方程。

\begin{qualityfigure}
\begin{tikzpicture}[>=Latex,font=\small]
\foreach\x/\y in {-1/1,1/1.5,2.5/.5,-.5/-1.2,2/-1.7,3/1.5}{\fill[DeepBlue](\x,\y)circle(1.4pt);}
\draw[->,Gold](3.7,1.4)--(3.7,-1.4)node[below]{$y,v$（向下）};\draw[thick,fill=PaleBlue](1,0)circle(.65);\draw(1,0)--(1.65,0)node[midway,above]{$r$};\node at(1,2.2){介质粒子在地面系静止};
\end{tikzpicture}
\figcap{球形水滴吸附静止雾滴的迎风截面}\label{mech:07:23:q}
\end{qualityfigure}
\end{origexamplebox}

\Needspace{6\baselineskip}\begin{xesolution}题图~\ref{mech:07:23:q}中，增长中的球半径为 $r$，迎风面积是 $\pi r^2$，质量却与 $r^3$ 成正比。质量流率与半径增长率须由这两个不同几何量联系。

记雾气中水分子的质量密度为 $\rho_f$，液态水的密度为 $\rho_w$，水滴半径、质量、下落
速度分别为 $r,m,v$。球形条件给出
\[
 \boxed{m=\frac43\pi\rho_w r^3}.
\]
水滴迎风截面积为 $\pi r^2$，所以吸附质量率为
\[
 \boxed{\frac{\dd m}{\dd t}=\rho_f\pi r^2v}.
\]
吸附前的雾滴静止，竖直动量方程仍为
\[
 \boxed{\frac{\dd(mv)}{\dd t}=mg}.
\]
将质量—半径关系代入质量增长方程，得到
\[
 \boxed{\frac{\dd r}{\dd t}=\frac{\rho_f}{4\rho_w}v}.
\]
展开动量方程还可写成
\[
 \boxed{
 \frac{\dd v}{\dd t}
 =g-\frac{3\rho_f}{4\rho_w r}v^2}.
\]
三个量之间的反馈关系见图~\ref{ch07:fig:soldiagThreeTwentyThree}。

\soldiagThreeTwentyThree
因此，在向下运动时，水滴下落越快，吸附质量越快；但速度是否增长还要看初值。
令 $\alpha=\rho_f/(4\rho_w)$，若初始半径为 $r_0$、初速度为零，令 $y=v^2$，由
$\dot r=\alpha v$ 和 $\dot v=g-3\alpha v^2/r$ 得到
\[
 \frac{\dd y}{\dd r}=\frac{2g}{\alpha}-\frac{6y}{r},\qquad
 \frac{\dd}{\dd r}(r^6y)=\frac{2g}{\alpha}r^6.
\]
积分给出 $r^6y=2g(r^7-r_0^7)/(7\alpha)$，即
\[
 v^2=\frac{2g}{7\alpha}\left(r-\frac{r_0^7}{r^6}\right).
\]
代回速度方程得 $a=g/7+6g(r_0/r)^7/7>0$。所以这一初值下半径、质量和速度都增加，
加速度从 $g$ 下降。长时间后 $v^2\propto r$，速度不会趋于有限终端值，而加速度趋近
\[
 \boxed{a\to\frac{g}{7}}.
\]
题设没有另给液态水密度，记液态水密度为常量 $\rho_w$，以便把质量与半径的关系写完整。
若初始速度为 $v_i$，积分常数改为 $r_0^6[v_i^2-2gr_0/(7\alpha)]$；当
$3\alpha v_i^2/r_0>g$ 时，水滴起初减速。$\alpha$ 无量纲，$gr/\alpha$ 的量纲为速度平方。
\end{xesolution}
}

{
\graphicspath{{../../../sample/xes-original/ninth-section/}}
\section{变质量系统}
% Original problem figures and added solution diagrams for Section 3.3.
\tikzset{
  solbody/.style={draw=Ink,line width=0.8pt},
  solforce/.style={-{Latex[length=2.0mm]},draw=DeepBlue,line width=0.9pt},
  solgeom/.style={draw=TextGray,densely dashed,line width=0.55pt},
  solflow/.style={-{Latex[length=1.8mm]},draw=Blue,line width=0.75pt},
  sollabel/.style={font=\scriptsize,fill=none,inner sep=1pt},
  solpanel/.style={font=\scriptsize\bfseries,text=TextGray}
}
\newcommand{\diagcaption}[2]{%
 \refstepcounter{figure}\par\nopagebreak[4]\smallskip
 {\centering\scriptsize\color{TextGray}图~\thefigure\quad #2\par}\smallskip}
% Shared geometry for repeated momentum problems. Captions remain at call sites.
\newcommand{\momentumRain}{%
\begin{tikzpicture}[scale=.78]
  \fill[Blue!8] (-1.2,-.65) rectangle (1.2,-.12);
  \draw[solbody] (-1.2,.7)--(-1.2,-.65)--(1.2,-.65)--(1.2,.7);
  \foreach \x in {-.9,-.45,0,.45,.9}{\draw[solflow] (\x,1.55)--(\x,.86);}
  \node[sollabel] at (0,2.05) {$q=\rho S\Delta h/\tau$，雨水水平速度为零};
  \node[sollabel] at (0,-1.15) {$m(t)=M+qt$};
  \draw[solforce] (1.2,-.38)--(2.55,-.38) node[midway,below=3pt,sollabel] {$F$};
  \draw[solflow] (1.6,.52)--(2.8,.52) node[midway,above=3pt,sollabel] {$v$};
\end{tikzpicture}}

\newcommand{\momentumPickup}{%
\begin{tikzpicture}[scale=.8]
  \draw[solbody] (-2,0)--(2.5,0);
  \draw[DeepBlue,line width=1.2pt] (-1.8,0)--(-.15,0)
    .. controls (-.05,0) and (0,.05) .. (0,.15)--(0,2.4);
  \fill[Ink] (0,2.4) circle (2pt);
  \draw[solforce] (0,2.4)--(0,3.3) node[left=3pt,sollabel] {$F$};
  \draw[solforce] (-.65,1.65)--(-.65,.7) node[midway,left=3pt,sollabel] {$\lambda xg$};
  \draw[solflow] (.65,1.1)--(.65,2.1) node[midway,right=3pt,sollabel] {$v$};
  \draw[solgeom] (0,2.4)--(1.8,2.4);
  \draw[<->,draw=TextGray] (1.8,0)--(1.8,2.4) node[midway,right=3pt,sollabel] {$x$};
  \node[sollabel] at (-1.0,-.38) {静止绳段};
  \node[sollabel] at (0,-.95) {$\dot m_{\rm mov}=\lambda v$，新增动量率 $\lambda v^2$};
\end{tikzpicture}}

\newcommand{\momentumFallingRope}{%
\begin{tikzpicture}[scale=.78]
  \draw[solbody] (-.45,3)--(.45,3);
  \foreach \x in {-.4,-.2,0,.2,.4}{\draw[TextGray] (\x,3)--++(.12,.14);}
  \draw[solgeom] (0,3)--(0,2.1);
  \draw[DeepBlue,line width=1.1pt] (0,2.1)--(0,.28)
    .. controls (0,.08) and (-.15,.08) .. (-.3,.08);
  \draw[DeepBlue,line width=1.1pt,domain=-.3:1.6,samples=120]
    plot ({\x},{.08+.08*sin((\x+.3)*720)});
  \draw[solbody] (-.75,0)--(1.9,0);
  \draw[solgeom] (-.7,3)--(0,3) (-.7,2.1)--(0,2.1);
  \draw[<->,draw=TextGray] (-.7,2.1)--(-.7,3) node[midway,left=3pt,sollabel] {$x$};
  \draw[solflow] (.55,1.8)--(.55,.9) node[midway,right=3pt,sollabel] {$v$};
  \node[sollabel,anchor=west] at (.65,2.8) {虚线：上端已下降的距离};
  \node[sollabel,anchor=west] at (1.15,1.5) {$l-x$：仍在下落的绳长};
  \node[sollabel] at (.4,-.4) {秤盘；堆积绳段长度为 $x$};
\end{tikzpicture}}

\newcommand{\momentumRopeLoad}{%
\begin{tikzpicture}[scale=.78]
  \draw[DeepBlue,line width=1.1pt] (0,2)--(0,.28)
    .. controls (0,.08) and (-.15,.08) .. (-.3,.08);
  \draw[DeepBlue,line width=1.1pt,domain=-.3:1.8,samples=120]
    plot ({\x},{.08+.08*sin((\x+.3)*720)});
  \draw[solbody] (-.65,0)--(2.05,0);
  \draw[solflow] (.55,1.7)--(.55,.85) node[midway,right=3pt,sollabel] {$v$};
  \draw[solforce] (.45,0)--(.45,-1.0) node[below=3pt,sollabel] {$\lambda xg$};
  \draw[solforce] (1.6,0)--(1.6,-1.0) node[below=3pt,sollabel] {$\lambda v^2$};
  \node[sollabel,anchor=west,align=left] at (2.7,1.15)
    {秤盘承受的向下作用\\$N=\lambda xg+\lambda v^2$\\第二项来自使绳元停止};
\end{tikzpicture}}

\newcommand{\momentumAccretion}{%
\begin{tikzpicture}[scale=.8]
  \draw[solbody,fill=SoftGray] (-.8,0) rectangle (.8,.9);
  \draw[solforce] (-.2,.45)--(-.2,-.85) node[below=3pt,sollabel] {$mg$};
  \draw[solflow] (1.25,.7)--(1.25,-.25) node[midway,right=3pt,sollabel] {$v$};
  \foreach \x in {-1.3,-.85,-.4,.05,.5,.95,1.4}{\fill[TextGray] (\x,-1.55) circle (1.5pt);}
  \node[sollabel] at (0,-1.95) {地面系中静止的介质};
  \draw[<->,TextGray] (-.8,1.3)--(.8,1.3) node[midway,above=3pt,sollabel] {截面积 $S$};
  \node[sollabel,anchor=west,align=left] at (2.3,.8)
    {$\dot m=\rho Sv$\\$\dd(mv)/\dd t=mg$\\$m\dot v=mg-\rho Sv^2$};
\end{tikzpicture}}

\newcommand{\momentumDroplet}{%
\begin{tikzpicture}[scale=.8]
  \draw[solbody,fill=Blue!8] (0,.5) circle (.7);
  \fill[Ink] (0,.5) circle (1.5pt);
  \draw[solgeom] (0,.5)--(.7,.5) node[midway,above=4pt,sollabel] {$r$};
  \draw[solforce] (0,.5)--(0,-1.0) node[below=3pt,sollabel] {$mg$};
  \draw[solflow] (1.25,.75)--(1.25,-.25) node[midway,right=3pt,sollabel] {$v$};
  \foreach \x in {-1.3,-.85,-.4,.05,.5,.95,1.4}{\fill[TextGray] (\x,-1.65) circle (1.5pt);}
  \node[sollabel] at (0,-2.05) {地面系中静止的雾滴};
  \node[sollabel,anchor=west,align=left] at (2.3,.8)
    {$S=\pi r^2$\\$\dot m=\rho_f\pi r^2v$\\$\dot r=\rho_f v/(4\rho_w)$};
\end{tikzpicture}}


\newcommand{\soldiagThreeTwentyFour}{%
\begin{center}\begin{minipage}{.97\linewidth}\centering
\momentumRain
\diagcaption{解图}{雨水只增加质量，不带入水平动量；外力使系统水平动量按 $Ft$ 增长。}\label{ch07:fig:soldiagThreeTwentyFour}\end{minipage}\end{center}}

\newcommand{\soldiagThreeTwentyFive}{%
\begin{center}\begin{minipage}{.97\linewidth}\centering
\momentumPickup
\diagcaption{解图}{提力既承担运动绳段的重量，也为新进入运动段的绳元提供动量。}\label{ch07:fig:soldiagThreeTwentyFive}\end{minipage}\end{center}}

\newcommand{\probdiagThreeTwentySix}{%
\begin{center}\begin{minipage}{.97\linewidth}\centering
\momentumFallingRope
\diagcaption{题图}{柔绳上端释放后下落，绳元依次落到秤盘上并停止。}\label{ch07:fig:probdiagThreeTwentySix}\end{minipage}\end{center}}

\newcommand{\soldiagThreeTwentySix}{%
\begin{center}\begin{minipage}{.97\linewidth}\centering
\momentumRopeLoad
\diagcaption{解图}{秤盘既承受已落绳段的重量，也连续吸收下落绳元的动量。}\label{ch07:fig:soldiagThreeTwentySix}\end{minipage}\end{center}}

\newcommand{\soldiagThreeTwentySeven}{%
\begin{center}\begin{minipage}{.97\linewidth}\centering
\momentumAccretion
\diagcaption{解图}{静止介质被吸收后要加速到 $v$，故出现随 $v^2$ 增长的动量流项。}\label{ch07:fig:soldiagThreeTwentySeven}\end{minipage}\end{center}}

\newcommand{\soldiagThreeTwentyEight}{%
\begin{center}\begin{minipage}{.97\linewidth}\centering
\momentumDroplet
\diagcaption{解图}{半径、迎风面积和质量共同增长，速度必须与它们联立求解。}\label{ch07:fig:soldiagThreeTwentyEight}\end{minipage}\end{center}}

\newcommand{\rocketMicroDiagram}{%
\begin{center}\begin{minipage}{.97\linewidth}\centering
\begin{tikzpicture}[scale=.77]
 \node[solpanel] at (-2.8,1.5) {喷气前}; \node[solpanel] at (2.6,1.5) {喷气后};
 \draw[solbody,fill=SoftGray] (-4,-.35) rectangle (-2,.35);
 \node[sollabel] at (-3,0) {$M$};
 \draw[solflow] (-2,0)--(-.7,0) node[midway,above=4pt,sollabel] {$v$};
 \draw[solbody,fill=SoftGray] (1.4,-.35) rectangle (3.4,.35);
 \node[sollabel] at (2.4,0) {$M+\Delta M$};
 \draw[solflow] (3.4,0)--(4.7,0) node[midway,above=4pt,sollabel] {$v+\Delta v$};
 \foreach \x in {-.2,.2,.6}{\fill[TextGray] (\x,0) circle (2pt);}
 \node[sollabel] at (.2,.75) {燃气 $-\Delta M$};
 \draw[solflow] (.65,-.95)--(-.65,-.95) node[midway,below=4pt,sollabel] {$-u$：相对火箭};
 \node[sollabel] at (.2,-1.8) {燃气对地速度为 $v-u$，可以向左、为零或向右};
\end{tikzpicture}
\diagcaption{示意图}{在地面系中，喷出燃气的速度是 $v-u$；$\Delta M<0$。}\label{ch07:fig:rocketMicroDiagram}\end{minipage}\end{center}}

\newcommand{\probdiagThreeTwentyNine}{%
\begin{center}\begin{minipage}{.97\linewidth}\centering
\begin{tikzpicture}[xscale=.82,yscale=.70]
  \draw[axis] (0,0)--(5.6,0) node[right,sollabel] {火箭质量};
  \draw[axis] (0,0)--(0,3.8) node[above,sollabel] {火箭速度};
  \draw[solgeom] (4.45,0)--(4.45,1.15)--(0,1.15);
  \fill[Ink] (4.45,1.15) circle (2.2pt);
  \node[sollabel,below] at (4.45,0) {$M$};
  \node[sollabel,left] at (0,1.15) {$v$};
  \node[sollabel,above right] at (4.45,1.15) {当前状态};
\end{tikzpicture}
\diagcaption{题图}{曲线应从当前状态 $(M,v)$ 出发，并标明燃烧过程的方向。}\label{ch07:fig:probdiagThreeTwentyNine}\end{minipage}\end{center}}

\newcommand{\soldiagThreeTwentyNine}{%
\begin{center}\begin{minipage}{.97\linewidth}\centering
\begin{tikzpicture}[scale=.82]
 \draw[axis] (0,0)--(5.2,0) node[right,sollabel] {$\mathcal M$};
 \draw[axis] (0,0)--(0,3.7) node[above,sollabel] {$V$};
 \draw[curve,domain=.6:4.4,samples=80] plot ({\x},{1+ln(4.4/\x)});
 \draw[curve,-{Latex[length=2mm]},domain=3.4:2.4,samples=20] plot ({\x},{1+ln(4.4/\x)});
 \draw[solgeom] (0,1)--(4.4,1)--(4.4,0);
 \node[sollabel,left=4pt] at (0,1) {$v$}; \node[sollabel,below=4pt] at (4.4,0) {$M$};
 \node[sollabel,anchor=west] at (2.5,3.2) {$V=v+u\ln(M/\mathcal M)$};
 \node[sollabel,anchor=west] at (2.5,2.6) {燃烧方向：沿曲线向左上方};
\end{tikzpicture}
\diagcaption{解图}{横轴质量减小时，火箭沿对数曲线向左上方运动。}\label{ch07:fig:soldiagThreeTwentyNine}\end{minipage}\end{center}}

\newcommand{\soldiagThreeThirty}{%
\begin{center}\begin{minipage}{.97\linewidth}\centering
\begin{tikzpicture}[scale=.82]
 \draw[axis] (0,0)--(5.8,0) node[right,sollabel] {$\mathcal M$};
 \draw[axis] (0,0)--(0,3.8) node[above,sollabel] {$u/\mathcal M$};
 \path[fill=Blue!12] (1.5,0)--plot[domain=1.5:4.6,samples=50] (\x,{2.1/\x})--(4.6,0)--cycle;
 \draw[curve,domain=.6:5.1,samples=80] plot (\x,{2.1/\x});
 \draw[solgeom] (1.5,0)--(1.5,1.4) (4.6,0)--(4.6,{2.1/4.6});
 \node[sollabel,below=4pt] at (1.5,0) {$M$}; \node[sollabel,below=4pt] at (4.6,0) {$M_0$};
 \node[sollabel] (L) at (3.8,2.5) {阴影面积 $u\ln(M_0/M)$};
 \draw[TextGray] (L.south)--(3,.4);
 \draw[solflow] (4.6,-.65)--(1.5,-.65) node[midway,below=4pt,sollabel] {质量减少的方向};
\end{tikzpicture}
\diagcaption{解图}{对 $\dd v=-u\,\dd M/M$ 积分，速度增量等于 $u/M$ 曲线下对应的面积。}\label{ch07:fig:soldiagThreeThirty}\end{minipage}\end{center}}


一辆行驶的汽车会不断消耗汽油，飞行中的飞机或火箭也有同样的特点：运动过程中，
研究对象的质量不断变化。处理这类问题时，必须同时计入速度变化和质量变化带来的
动量变化。变质量并不意味着动量定理失效；真正需要重新说明的是研究系统的边界——哪一
部分物质在当前时刻属于研究对象，哪些物质正从边界流入或流出，以及它们穿过边界前后的
速度分别是多少。

先只对主体动量作代数展开，不把它等同于外力冲量：
\begin{equation}
 \Delta\vec p=\Delta(m\vec v)
 =\Delta m\,\vec v+m\,\Delta\vec v+\Delta m\,\Delta\vec v.
\end{equation}
略去二阶小量并取连续极限，得到严格的微分关系
\[
 \dd\vec p=\vec v\,\dd m+m\,\dd\vec v.
\]
应用这一关系时还要注意：新进入系统或刚离开系统的物质，未必具有研究对象本身的速度。
这些物质所带入或带走的动量，正是变质量问题中容易遗漏的一项。

\subsubsection*{变质量系统的边界与动量通量}

火箭、落链和接雨小车的研究对象会随时间增减质量，不能直接把 $\vec F=m\vec a$ 中的
$m$ 换成 $m(t)$。动量定理适用于同一批物质；若研究对象在两个时刻包含的物质不同，
就要把跨越边界所带的动量单独记入。

设主体质量为 $m$、速度为 $\vec v$，时间 $\dd t$ 内吸收质量 $\dd m_{\rm in}>0$，
排出质量 $\dd m_{\rm out}>0$。流入、流出物质在同一个惯性系中的速度分别为
$\vec u_{\rm in},\vec u_{\rm out}$，主体质量改变
$\dd m=\dd m_{\rm in}-\dd m_{\rm out}$。把主体及这次交换的物质选为同一批物质：
\[
 m\vec v+\vec u_{\rm in}\dd m_{\rm in}+\vec F_{\rm ext}\dd t
 =(m+\dd m)(\vec v+\dd\vec v)+\vec u_{\rm out}\dd m_{\rm out}.
\]
这里 $\vec F_{\rm ext}$ 是连续极限下作用于主体的外力；微小交换质量所受有限外力的冲量
为二阶小量。展开、略去二阶项并除以 $\dd t$，得到
\begin{equation}\label{ch07:eq:mass-flux-balance}
 \frac{\dd(m\vec v)}{\dd t}
 =\vec F_{\rm ext}+\dot m_{\rm in}\vec u_{\rm in}
 -\dot m_{\rm out}\vec u_{\rm out}.
\end{equation}
左侧是主体动量增加率，右侧依次为外力、流入动量率、流出动量率。展开左侧乘积，得
\begin{equation}\label{ch07:eq:mass-acceleration}
 m\frac{\dd\vec v}{\dd t}
 =\vec F_{\rm ext}+\dot m_{\rm in}(\vec u_{\rm in}-\vec v)
 -\dot m_{\rm out}(\vec u_{\rm out}-\vec v).
\end{equation}
只有一股交换流时，也可用有符号的 $\dot m$（流入为正、流出为负）和相对主体速度
$\vec u_{\rm rel}=\vec u-\vec v$ 写成 $m\dot{\vec v}=\vec F_{\rm ext}+\vec u_{\rm rel}\dot m$。
不能把这个相对速度与前一式的惯性系绝对速度混用。动量通量的量纲为
$\mathrm{kg/s}\times\mathrm{m/s}=\mathrm N$，与外力一致。

理想火箭在无外力的一维运动中，以相对火箭大小为 $u$ 的速度向后喷气，满足
\begin{equation}
 m\,\dd v=-u\,\dd m,
 \qquad
 v-v_0=u\ln\frac{m_0}{m}.
\end{equation}
速度增量取决于质量比而不是燃料质量的简单差值。若有重力或阻力，还要把它们在燃烧时间内的
冲量计入；若喷气速度相对地面给出，必须先换成相对火箭速度。
火箭自身不是固定物质系统，即使无外力，其自身动量也不必守恒；守恒的是火箭加燃气的总动量。

链条落到秤盘、雨水落入车厢等问题还包含“捕获并加速新质量”的作用。流入物质在进入系统前
已有速度，主体必须提供冲量把它变为共同速度，于是支持力或阻力中出现 $\dot m\Delta v$ 项。
柔链问题还应区分移动段、静止段和弯折区，弯折区虽很短，却承担改变动量方向的有限合力。
检验答案时可看两个极限：质量流率趋于零时应回到定质量问题；相对流速趋于零时，动量通量项
也应消失。最后一种情形只说明无需额外改变主体速度，主体总动量仍可随它包含的质量改变。


\begin{origexamplebox}{变式：横向来雨的入流动量}\label{ch07:ex:s09-3-24}

一个初始质量为 $M$、水平开口面积为 $S$ 的容器，从静止开始受恒定水平力 $F>0$ 推动。
地面光滑，雨水以恒定的向下速度落入开口，且在地面系中还具有与 $F$ 同向的恒定水平速度 $u\ge0$。
雨水完全留在容器内；稳定质量流率记为 $q=\rho S\Delta h/\tau>0$，其中静止同样容器在
$\tau$ 时间内水位上升 $\Delta h$。忽略侧面撞雨及溢出。求 $t=3\tau$ 时的速度，并讨论长时间极限。

\begin{qualityfigure}
\begin{tikzpicture}[>=Latex,font=\small]
\draw[thick](0,1.8)--(0,0)--(3,0)--(3,1.8);\fill[PaleBlue](.05,.05)rectangle(2.95,.8);\draw[thick](-.5,-.2)--(6,-.2);
\foreach\x in {.4,1.2,2,2.8}{\draw[->,DeepBlue](\x,3)--({\x+.6},2);}
\draw[->,Gold](3.2,.8)--(4.7,.8)node[right]{$F$};\node at(1.5,1.2){$M,S$};\node at(1.5,3.5){雨水密度 $\rho$};\draw[->,Gold](3.5,2.7)--(5,2.7)node[right]{$u$（雨水水平速度）};
\end{tikzpicture}
\figcap{横向有风的雨水携带入流动量}\label{mech:07:24:q}
\end{qualityfigure}
\end{origexamplebox}

\Needspace{6\baselineskip}\begin{xesolution}
\begin{qualityfigure}
\begin{tikzpicture}[>=Latex,font=\small]
\draw[thick](0,1.8)--(0,0)--(3,0)--(3,1.8);\fill[PaleBlue](.05,.05)rectangle(2.95,.8);\draw[thick](-.5,-.2)--(6,-.2);
\foreach\x in {.4,1.2,2,2.8}{\draw[->,DeepBlue](\x,3)--({\x+.6},2);}
\draw[->,Gold](3.2,.8)--(4.7,.8)node[right]{$F$};\node at(1.5,1.2){$M,S$};\node at(1.5,3.5){雨水密度 $\rho$};\draw[->,Gold](3.5,2.7)--(5,2.7)node[right]{$u$（雨水水平速度）};\node[align=left]at(5,-1){质量流率 $q$\\带入水平动量流率 $qu$};
\end{tikzpicture}
\figcap{横向有风的雨水携带入流动量（解答用辅助图）}\label{mech:07:24:a}
\end{qualityfigure}
图~\ref{mech:07:24:a}中，雨水在地面系中带有水平速度 $u$，所以每秒不仅带入质量 $q$，还带入动量 $qu$。这是本变式与竖直入雨前例的关键差别。

取推力方向为正。容器与已收集的水构成变质量系统，$m(t)=M+qt$。
在 $\dd t$ 内进入的雨水质量为 $q\dd t$，它在进入前已有水平动量 $qu\dd t$。
将这部分带入动量和外力冲量同时计入，得
\[
 \frac{\dd}{\dd t}[(M+qt)v]=F+qu.
\]
初始动量为零，积分为 $(M+qt)v=(F+qu)t$，所以
\[
 \boxed{v(3\tau)=\frac{3\tau(F+qu)}{M+3q\tau}},\qquad
 \boxed{\lim_{t\to\infty}v=u+\frac Fq}.
\]
展开式也可写成 $(M+qt)\dot v=F+q(u-v)$：入流相对容器的速度差产生 $q(u-v)$ 项。
当容器超过雨水的水平速度时，该项反而使它减速，但恒力仍可使速度继续增加并渐近上述极限。
检查 $u=0$ 恢复竖直来雨前例；$q\to0$ 的有限时间极限恢复 $v=Ft/M$。
$qu$ 与 $F$ 同为力，分母为质量，所得速度量纲正确。极限仅在持续稳定入雨且容器不溢出的模型内成立。
\end{xesolution}

\begin{origexamplebox}{变式：线密度不均匀的提绳}\label{ch07:ex:s09-3-25}

长度为 $l$、总质量为 $m$ 的柔软绳放在水平桌面上。用手捏住一端，以恒定速度 $v>0$ 竖直提起。
从手持端沿绳的材料坐标记为 $s$，线密度为 $\lambda(s)=\lambda_0(1+s/l)$，$0\le s\le l$。
当手端高度为 $0<x<l$ 时，已提起的绳竖直，余绳仍静止，忽略弯折区长度和弹性。
求 $\lambda_0$、运动绳段质量与手的提力。新进入运动段的绳元按瞬时拾取模型处理。

\begin{qualityfigure}
\begin{tikzpicture}[>=Latex,font=\small]
\draw[thick](-1,0)--(5,0);\draw[very thick](0,2.8)--(0,.2)..controls(.1,0)and(.3,0)..(.5,0)--(3.5,0);
\draw[->,DeepBlue](.4,2)--(.4,3.1)node[right]{$v$};\draw[<->](-.5,0)--(-.5,2.8)node[midway,left]{$x$};\node at(2,1){桌上绳段原来静止};\node[align=left]at(3.5,2.4){材料坐标 $s$ 从手端起\\$\lambda(s)=\lambda_0(1+s/l)$};
\end{tikzpicture}
\figcap{线密度沿材料坐标增大的绳}\label{mech:07:25:q}
\end{qualityfigure}
\end{origexamplebox}

\Needspace{6\baselineskip}\begin{xesolution}
\begin{qualityfigure}
\begin{tikzpicture}[>=Latex,font=\small]
\draw[thick](-1,0)--(5,0);\draw[very thick](0,2.8)--(0,.2)..controls(.1,0)and(.3,0)..(.5,0)--(3.5,0);
\draw[->,DeepBlue](.4,2)--(.4,3.1)node[right]{$v$};\draw[<->](-.5,0)--(-.5,2.8)node[midway,left]{$x$};\node at(2,1){桌上绳段原来静止};\node[align=left]at(3.5,2.4){材料坐标 $s$ 从手端起\\$\lambda(s)=\lambda_0(1+s/l)$};\draw[->,Gold](4.5,.2)--(4.5,1.1)node[right]{$\dot m=\lambda(x)v$};\node at(2,-.8){运动段质量 $m_x=\int_0^x\lambda(s)\,\dd s$};
\end{tikzpicture}
\figcap{线密度沿材料坐标增大的绳（解答用辅助图）}\label{mech:07:25:a}
\end{qualityfigure}
图~\ref{mech:07:25:a}中，材料坐标从手端起算，当前进入运动段的绳元位于 $s=x$，其线密度为 $\lambda(x)$；已运动部分的质量则是从 $0$ 到 $x$ 的积分。

先用总质量归一化线密度：
\[
 m=\int_0^l\lambda_0(1+s/l)\,\dd s=\frac32\lambda_0l,
 \qquad \boxed{\lambda_0=\frac{2m}{3l}}.
\]
已运动的材料是 $0\le s\le x$，故运动段质量为
\[
 \boxed{m_x=\lambda_0\left(x+\frac{x^2}{2l}\right)}.
\]
拾取点对应 $s=x$，每秒进入的绳长为 $v$，因此
$\dot m_x=\lambda(x)v=\lambda_0(1+x/l)v$。
新绳元从零速度变为向上速度 $v$；运动段自身保持恒速，手力与重力的差供给这股动量流率：
\[
 F-m_xg=v\dot m_x=\lambda_0(1+x/l)v^2.
\]
于是
\[
 \boxed{F=\frac{2m}{3l}\left[g\left(x+\frac{x^2}{2l}\right)
 +(1+x/l)v^2\right]},\qquad 0<x<l.
\]
重力项包含运动段线密度的积分，拾取项使用当前位置的局部线密度，不能用同一个平均值替代二者。
若将分布改回常量 $\lambda=m/l$，立即恢复 $F=\lambda gx+\lambda v^2$。
两项量纲均为力；$x\to l^-$ 仍有拾取项，但最后一段完全离桌后质量流率变为零，恒速手力随之变为 $mg$。
\end{xesolution}

\Needspace{8\baselineskip}
\begin{origexamplebox}{}\label{ch07:ex:s09-3-26}
如图，一根质量为 $m$、长度为 $l$ 的均匀柔绳竖直悬挂，下端刚好与台秤秤盘接触。
释放绳的上端，求从释放开始台秤示数随时间的变化，并求台秤的最大示数。
图~\ref{ch07:fig:probdiagThreeTwentySix}中，柔绳上端释放后下落，绳元依次落到秤盘上并停止。
\probdiagThreeTwentySix
\end{origexamplebox}

\Needspace{6\baselineskip}\begin{xesolution}
设绳的线密度为 $\lambda=m/l$，释放后经过时间 $t$，绳的上端下落距离为 $x$。理想模型
中尚未落到秤盘上的绳段作自由落体，因而
\[
 x=\frac12gt^2,\qquad v=gt,\qquad v^2=2gx.
\]
秤盘的受力来源有两部分：
\[
 \underbrace{\lambda xg}_{\text{已停在盘上的绳重}}
 \quad\text{和}\quad
 \underbrace{\lambda v^2}_{\text{使下落绳元停止的动量流}}.
\]
图~\ref{ch07:fig:soldiagThreeTwentySix}中，秤盘既承受已落绳段的重量，也连续吸收下落绳元的动量。
\soldiagThreeTwentySix
所以在整根绳尚未全部落下时，
\[
 N=\lambda xg+\lambda v^2
 =3\lambda gx
 =\frac32\lambda g^2t^2.
\]
上端落到秤盘所需时间为
\[
 t_0=\sqrt{\frac{2l}{g}}.
\]
因此理想化的示数为
\[
 \boxed{
 N(t)=
 \begin{cases}
 \dfrac{3m g^2}{2l}\,t^2,
   &0\le t<\sqrt{\dfrac{2l}{g}},\\[6pt]
 mg,&t>\sqrt{\dfrac{2l}{g}}.
 \end{cases}}
\]
在全部落下前的一瞬间，
\[
 \boxed{N_{\max}=3mg}.
\]
随后动量流消失，稳定示数为 $mg$。理想模型在 $t=t_0$ 处给出突变；真实柔绳因弯曲、
碰撞持续时间和秤的响应时间有限，峰值会被平滑。
\end{xesolution}

\begin{origexamplebox}{变式：介质密度随深度增加的吸附下落}\label{ch07:ex:s09-3-27}

初始质量为 $m_0$、固定迎风面积为 $S$ 的物体从静止释放，竖直向下的深度记为 $y\ge0$。
介质粒子在地面系中静止，质量密度为 $\rho(y)=\rho_0(1+y/L)$，其中 $\rho_0,L>0$。
物体吸附所有遇到的粒子，重力加速度为常量 $g$，忽略其他作用。求质量 $m(y)$，写出并积分
速度平方的方程，比较深处运动与均匀介质前例的差别。

\begin{qualityfigure}
\begin{tikzpicture}[>=Latex,font=\small]
\foreach\x/\y in {-1/1,1/1.5,2.5/.5,-.5/-1.2,2/-1.7,3/1.5}{\fill[DeepBlue](\x,\y)circle(1.4pt);}
\draw[->,Gold](3.7,1.4)--(3.7,-1.4)node[below]{$y,v$（向下）};\draw[thick,fill=PaleBlue](0,-.5)rectangle(2,.5);\node at(1,0){$m$};\node at(1,-.9){截面积 $S$};\node at(1,2.2){介质粒子在地面系静止};\node at(1,-2.4){$\rho(y)=\rho_0(1+y/L)$};
\end{tikzpicture}
\figcap{介质密度随下落深度线性增加}\label{mech:07:27:q}
\end{qualityfigure}
\end{origexamplebox}

\Needspace{6\baselineskip}\begin{xesolution}
\begin{qualityfigure}
\begin{tikzpicture}[>=Latex,font=\small]
\foreach\x/\y in {-1/1,1/1.5,2.5/.5,-.5/-1.2,2/-1.7,3/1.5}{\fill[DeepBlue](\x,\y)circle(1.4pt);}
\draw[->,Gold](3.7,1.4)--(3.7,-1.4)node[below]{$y,v$（向下）};\draw[thick,fill=PaleBlue](0,-.5)rectangle(2,.5);\node at(1,0){$m$};\node at(1,-.9){截面积 $S$};\node at(1,2.2){介质粒子在地面系静止};\node at(1,-2.4){$\rho(y)=\rho_0(1+y/L)$};\node at(1,3){$\dot m=\rho(y)Sv$，入流初动量为零};
\end{tikzpicture}
\figcap{介质密度随下落深度线性增加（解答用辅助图）}\label{mech:07:27:a}
\end{qualityfigure}
图~\ref{mech:07:27:a}中，越深处介质越密，瞬时质量流率由当前位置的 $\rho(y)$ 决定，不能继续把 $\rho S$ 当作常量。图中 $y$ 从释放点向下量取。

以题图向下的 $y$ 为正。扫掠厚度为 $\dd y$ 时，吸附质量为
\[
 \dd m=\rho(y)S\,\dd y.
\]
令 $k=\rho_0S$，由 $m(0)=m_0$ 得
\[
 \boxed{m(y)=m_0+k\left(y+\frac{y^2}{2L}\right)},\qquad
 \dot m=k(1+y/L)v.
\]
入流粒子原本静止，故 $\dd(mv)/\dd t=mg$，即
$\dot v=g-(m'(y)/m)v^2$。在开始下落后令 $w=v^2$，用 $\dot v=v\,\dd v/\dd y$ 得
\[
 w'+2\frac{m'}m w=2g,\qquad (m^2w)'=2gm^2.
\]
结合 $w(0)=0$，得到含有完整初值的结果
\[
 \boxed{v^2(y)=\frac{2g}{m(y)^2}\int_0^y m(\xi)^2\,\dd\xi}.
\]
本题积分可直接展开为
\[
 \int_0^y m(\xi)^2\dd\xi
 =m_0^2y+m_0ky^2+\frac{k^2+m_0k/L}{3}y^3
 +\frac{k^2}{4L}y^4+\frac{k^2}{20L^2}y^5.
\]
在 $y\to0$ 时 $v^2\sim2gy$，恢复刚释放时的自由落体起步。
深处 $m\sim ky^2/(2L)$，所以 $v^2\sim2gy/5$，于是 $a=\tfrac12\dd(v^2)/\dd y\to g/5$。
这与均匀介质的 $g/3$ 不同，差别来自密度增长使质量随深度按平方增长。
令 $L\to\infty$ 可恢复均匀介质；积分量纲是质量平方乘长度，乘 $g/m^2$ 后为速度平方。
\end{xesolution}

\Needspace{8\baselineskip}
\begin{origexamplebox}{变式：捕获效率随半径变化的水滴}\label{ch07:ex:s09-3-28}

液态密度为 $\rho_w$ 的球形水滴在静止均匀雾中由静止下落，雾的质量密度为 $\rho_f$。
初始半径为 $r_0$，重力加速度为 $g$。与全部吸附的前例不同，水滴只捕获迎面雾滴的比例
$\eta(r)=r_*/r$，其中 $0<r_*\le r_0$，未捕获粒子直接穿过且不交换动量。
忽略空气阻力和蒸发，水滴始终为球形。列出运动方程，求 $v^2(r)$ 及长时间加速度的极限。

\begin{qualityfigure}
\begin{tikzpicture}[>=Latex,font=\small]
\foreach\x/\y in {-1/1,1/1.5,2.5/.5,-.5/-1.2,2/-1.7,3/1.5}{\fill[DeepBlue](\x,\y)circle(1.4pt);}
\draw[->,Gold](3.7,1.4)--(3.7,-1.4)node[below]{$y,v$（向下）};\draw[thick,fill=PaleBlue](1,0)circle(.65);\draw(1,0)--(1.65,0)node[midway,above]{$r$};\node at(1,2.2){介质粒子在地面系静止};\node at(1,-2.4){捕获比例 $\eta(r)=r_*/r$};
\end{tikzpicture}
\figcap{捕获效率随水滴半径减小}\label{mech:07:28:q}
\end{qualityfigure}
\end{origexamplebox}

\Needspace{6\baselineskip}\begin{xesolution}
\begin{qualityfigure}
\begin{tikzpicture}[>=Latex,font=\small]
\foreach\x/\y in {-1/1,1/1.5,2.5/.5,-.5/-1.2,2/-1.7,3/1.5}{\fill[DeepBlue](\x,\y)circle(1.4pt);}
\draw[->,Gold](3.7,1.4)--(3.7,-1.4)node[below]{$y,v$（向下）};\draw[thick,fill=PaleBlue](1,0)circle(.65);\draw(1,0)--(1.65,0)node[midway,above]{$r$};\node at(1,2.2){介质粒子在地面系静止};\node at(1,-2.4){捕获比例 $\eta(r)=r_*/r$};\node at(1,3){$\dot m=\eta(r)\rho_f\pi r^2v$};
\end{tikzpicture}
\figcap{捕获效率随水滴半径减小（解答用辅助图）}\label{mech:07:28:a}
\end{qualityfigure}
图~\ref{mech:07:28:a}中，迎风截面仍为 $\pi r^2$，但只有比例 $\eta(r)=r_*/r$ 被吸附。这个效率因子改变质量通量，随后也改变动量方程中的稀释项。

令 $b=\rho_fr_*/(4\rho_w)$，它的量纲是长度。球形质量关系不变，但吸附率含效率因子：
\[
 m=\frac43\pi\rho_wr^3,\qquad
 \dot m=\eta(r)\rho_f\pi r^2v=\rho_f\pi r_*rv.
\]
由质量式求导，得 $\dot r=bv/r$。被捕获雾滴入流前动量为零，所以
\[
 \frac{\dd(mv)}{\dd t}=mg,\qquad
 \boxed{\dot v=g-\frac{3b}{r^2}v^2},\qquad
 \boxed{\dot r=\frac br v}.
\]
在 $v>0$ 阶段令 $w=v^2$，消去时间：
\[
 \frac{\dd w}{\dd r}=\frac{2gr}{b}-\frac{6w}{r},
 \qquad \frac{\dd}{\dd r}(r^6w)=\frac{2g}{b}r^7.
\]
从 $(r_0,0)$ 积分，得到
\[
 \boxed{v^2(r)=\frac{g}{4b}\left(r^2-\frac{r_0^8}{r^6}\right)}.
\]
代回速度方程：
\[
 \boxed{a=\frac g4+\frac{3g}{4}\left(\frac{r_0}{r}\right)^8},
 \qquad \boxed{a\longrightarrow g/4}.
\]
初始 $r=r_0$ 时 $v=0,a=g$，之后半径增大、加速度下降但始终为正。
$r\ge r_0\ge r_*$ 保证 $0<\eta\le1$，因此捕获比例的模型自洽。
长时 $v^2\sim gr^2/(4b)$，从而 $\dot r\to\sqrt{gb}/2$、$v\sim gt/4$。
这里的 $g/4$ 与全捕获模型的 $g/7$ 不同，差别由半径相关效率改变了质量通量。
\end{xesolution}

火箭依靠从尾部喷出燃料燃烧产生的高温高压气体而获得巨大推力。实际升空过程十分复杂，
为简化分析，下面只研究火箭沿直线飞行的理想模型。忽略重力和其他外力，设火箭喷出的
燃气相对火箭向后速度大小为 $u$。在极短时间 $\Delta t$ 内，喷气前火箭的质量和速度为 $M,v$；
喷气后为 $M+\Delta M,v+\Delta v$，其中 $\Delta M<0$。喷出燃气的质量为
$-\Delta M$，在地面参考系中的速度为 $v-u$。
图~\ref{ch07:fig:rocketMicroDiagram}中，在地面系中，喷出燃气的速度是 $v-u$；$\Delta M<0$。
\rocketMicroDiagram
由喷气前后的总动量守恒，
\[
 Mv=(M+\Delta M)(v+\Delta v)-\Delta M(v-u).
\]
展开并消去 $Mv$：
\[
 M\Delta v+u\Delta M+\Delta M\Delta v=0.
\]
略去二阶小量 $\Delta M\Delta v$，得到
\begin{equation}
 \label{ch07:eq:rocket-micro}
 \Delta v=-u\frac{\Delta M}{M}.
\end{equation}
因为 $\Delta M<0$，所以 $\Delta v>0$：火箭向后喷出质量，自身向前加速。在连续极限下，
式 \eqref{ch07:eq:rocket-micro} 可写成
\[
 \dd v=-u\frac{\dd M}{M}.
\]

\Needspace{8\baselineskip}
\begin{origexamplebox}{}\label{ch07:ex:s09-3-29}
一枚火箭在某时刻的质量为 $M$、速度为 $v$，燃气相对火箭的喷出速度大小为 $u$。
忽略重力，请根据式 \eqref{ch07:eq:rocket-micro}，在下面坐标系中定性画出此后火箭速度与质量的关系。
图~\ref{ch07:fig:probdiagThreeTwentyNine}中，曲线应从当前状态 $(M,v)$ 出发，并标明燃烧过程的方向。
\probdiagThreeTwentyNine
\end{origexamplebox}

\Needspace{6\baselineskip}\begin{xesolution}
设此后某时刻火箭的剩余质量为 $\mathcal M$、速度为 $V$。由
\[
 \dd V=-u\frac{\dd\mathcal M}{\mathcal M}
\]
从当前状态 $(M,v)$ 积分到 $(\mathcal M,V)$：
\[
 V-v=-u\int_M^{\mathcal M}\frac{\dd m}{m}
 =u\ln\frac{M}{\mathcal M}.
\]
因此
\[
 \boxed{V(\mathcal M)=v+u\ln\frac{M}{\mathcal M}},
 \qquad 0<\mathcal M\le M.
\]
图~\ref{ch07:fig:soldiagThreeTwentyNine}中，横轴质量减小时，火箭沿对数曲线向左上方运动。
\soldiagThreeTwentyNine
曲线经过 $(M,v)$，并满足
\[
 \frac{\dd V}{\dd\mathcal M}=-\frac{u}{\mathcal M}<0,
 \qquad
 \frac{\dd^2V}{\dd\mathcal M^2}
 =\frac{u}{\mathcal M^2}>0.
\]
所以按“横轴质量、纵轴速度”作图时，曲线单调下降，且斜率随质量增大逐渐变平（曲线
开口向上）；实际燃烧过程的方向是沿曲线向左上方运动，即质量减小、速度增大。
$\mathcal M\to0$ 时速度发散只是理想模型的数学
外推，真实火箭不可能把结构质量全部喷掉。
\end{xesolution}

\Needspace{8\baselineskip}
\begin{origexamplebox}{}\label{ch07:ex:s09-3-30}
一枚单级火箭的初始质量为 $M_0$，燃料燃烧完毕后的质量为 $M$。燃气相对火箭的喷出速度
大小恒为 $u$，火箭由静止开始运动，忽略重力和其他阻力。证明火箭的末速度为
\[
 v=u\ln\frac{M_0}{M}.
\]

\begin{qualityfigure}
\begin{tikzpicture}[>=Latex,font=\small]
\draw[thick,fill=PaleBlue](0,-.5)--(2,-.5)--(2.8,0)--(2,.5)--(0,.5)--cycle;\node at(1,0){$M(t)$};
\draw[->,DeepBlue](2.8,0)--(4.2,0)node[right]{$v$};\draw[->,Gold](0,0)--(-1.7,0)node[left]{$u$（相对火箭）};
\node[align=center]at(1,-1.5){燃烧前 $M_0$，燃尽后 $M$\\无重力和其他外力};
\end{tikzpicture}
\figcap{火箭减质量与相对喷气速度}\label{mech:07:30:q}
\end{qualityfigure}
\end{origexamplebox}

\Needspace{6\baselineskip}\begin{xesolution}题图~\ref{mech:07:30:q}中，喷气速率 $u$ 是相对火箭的速率，方向与火箭速度相反。火箭质量递减，所以质量微元 $\dd M<0$，积分时必须保持这一符号约定。

由式 \eqref{ch07:eq:rocket-micro} 的连续形式，
\[
 \dd v=-u\frac{\dd M}{M}.
\]
火箭由静止开始，质量由 $M_0$ 减少到 $M$。对整个燃烧过程积分：
\[
 \int_0^v\dd v
 =-u\int_{M_0}^{M}\frac{\dd M}{M}.
\]
图~\ref{ch07:fig:soldiagThreeThirty}中，对 $\dd v=-u\,\dd M/M$ 积分，速度增量等于 $u/M$ 曲线下对应的面积。
\soldiagThreeThirty
于是
\[
 v=-u\left(\ln M-\ln M_0\right),
\]
即
\[
 \boxed{v=u\ln\frac{M_0}{M}}.
\]
这就是理想火箭方程。若定义质量比 $R=M_0/M$，则速度增量为
$\Delta v=u\ln R$：提高喷气速度 $u$ 能线性提高速度增量，而增大质量比的收益只有
对数增长。证明成立的条件是喷气相对速度恒定，并忽略重力、空气阻力等外力。
\end{xesolution}
}

{
\graphicspath{{../../../sample/xes-original/twelfth-section/}}
\section{碰撞}
% Problem figures and added solution diagrams for Section 4.3.
\tikzset{
  solbody/.style={draw=Ink,line width=0.8pt},
  solforce/.style={-{Latex[length=2.0mm]},draw=DeepBlue,line width=0.9pt},
  solgeom/.style={draw=TextGray,densely dashed,line width=0.55pt},
  solflow/.style={-{Latex[length=1.8mm]},draw=Blue,line width=0.75pt},
  sollabel/.style={font=\scriptsize,fill=none,inner sep=1pt},
  solpanel/.style={font=\scriptsize\bfseries,text=TextGray},
  solpath/.style={draw=Ink,line width=.75pt}
}
\newcommand{\diagcaption}[2]{%
 \refstepcounter{figure}\par\nopagebreak[4]\smallskip
 {\centering\scriptsize\color{TextGray}图~\thefigure\quad #2\par}\smallskip}

\newcommand{\probdiagFourFifteen}{%
\begin{center}\begin{minipage}{.97\linewidth}\centering
\begin{tikzpicture}[scale=.72]
  \draw[solbody] (-4.1,-.70)--(4.1,-.70);
  \draw[solbody,fill=SoftGray] (-3.15,-.70) rectangle (-2.20,.20);
  \draw[solbody,fill=SoftGray] (2.20,-.70) rectangle (3.15,.20);
  \draw[decorate,decoration={coil,segment length=5pt,amplitude=3pt},line width=.75pt]
    (-2.20,-.25)--(2.20,-.25);
  \node[sollabel] at (-2.68,.75) {$m_1$};
  \node[sollabel] at (2.68,.75) {$m_2$};
  \draw[<->,draw=DeepBlue] (-2.20,-1.05)--(2.20,-1.05)
    node[midway,below,sollabel] {$r=l<l_0$};
  \draw[solflow] (-2.70,1.10)--(-1.70,1.10) node[midway,above,sollabel] {$v_1$};
  \draw[solflow] (2.25,1.10)--(3.25,1.10) node[midway,above,sollabel] {$v_2$};
\end{tikzpicture}
\diagcaption{题意图}{弹簧只改变两物体的相对运动；整体质心速度保持不变。}\label{ch07:fig:probdiagFourFifteen}\end{minipage}\end{center}}

\newcommand{\soldiagFourFifteen}{%
\begin{center}\begin{minipage}{.97\linewidth}\centering
\begin{tikzpicture}[scale=.85]
 \draw[axis] (0,0)--(6.6,0) node[right,sollabel] {$r$};
 \draw[axis] (0,0)--(0,2.7) node[above,sollabel] {能量};
 \draw[curve,domain=.6:5.4,samples=80] plot ({\x},{.35*(\x-3)^2});
 \draw[TextGray,line width=.8pt] (.45,1.12)--(5.75,1.12);
 \draw[solgeom] (1.21115,0)--(1.21115,1.12) (4.78885,0)--(4.78885,1.12);
 \node[sollabel,above=4pt] at (5.7,1.12) {$E_{\rm rel}$};
 \node[sollabel] at (3,2.45) {$U(r)=\frac12k(r-l_0)^2$};
 \node[sollabel,below=4pt] at (1.21115,0) {$r_{\min}$};
 \node[sollabel,below=4pt] at (3,0) {$l_0$};
 \node[sollabel,below=4pt] at (4.78885,0) {$r_{\max}$};
\end{tikzpicture}
\diagcaption{解图}{相对势能关于原长位置对称，同一相对能量给出最近与最远两个转折点。}\label{ch07:fig:soldiagFourFifteen}\end{minipage}\end{center}}

\newcommand{\probdiagFourSixteen}{%
\begin{center}\begin{minipage}{.97\linewidth}\centering
\begin{tikzpicture}[scale=.72]
  \draw[solbody] (-1.75,-1.75)--(1.75,-1.75);
  \foreach \x in {-1.55,-1.25,...,1.55}{\draw[draw=TextGray,line width=.45pt] (\x,-1.75)--++(-.25,-.22);}
  \draw[solbody,fill=SoftGray] (-.55,-1.75) rectangle (.55,-.85);
  \draw[decorate,decoration={coil,segment length=5pt,amplitude=3pt},line width=.75pt]
    (0,-.85)--(0,1.10);
  \draw[solbody,fill=SoftGray] (-.55,1.10) rectangle (.55,2.00);
  \node[sollabel,right] at (.65,1.55) {$A,m_1$};
  \node[sollabel,right] at (.65,-1.30) {$B,m_2$};
  \draw[<->,draw=DeepBlue] (-.95,-.85)--(-.95,1.10) node[midway,left,sollabel] {$l<l_0$};
\end{tikzpicture}
\diagcaption{题意图}{在 $B$ 尚未离地的阶段，下端固定，只有上方物体 $A$ 运动。}\label{ch07:fig:probdiagFourSixteen}\end{minipage}\end{center}}

\newcommand{\soldiagFourSixteen}{%
\begin{center}\begin{minipage}{.97\linewidth}\centering
\begin{tikzpicture}[scale=.8]
 \node[solpanel] at (0,2.6) {下物块 $B$ 的临界受力图};
 \draw[solbody,fill=SoftGray] (-.4,-.35) rectangle (.4,.35);
 \draw[solbody] (-1.2,-.35)--(1.2,-.35);
 \draw[solforce] (0,0)--(0,1.7) node[above=4pt,sollabel] {$T=k(x-l_0)$};
 \draw[solforce] (0,0)--(0,-1.7) node[below=4pt,sollabel] {$m_2g$};
 \node[sollabel,right=5pt] at (.4,0) {$B$};
 \node[sollabel,anchor=west,align=left] at (2,1)
   {临界时 $N=0$\\$T=m_2g$\\$x_{\rm c}=l_0+m_2g/k$};
\end{tikzpicture}
\diagcaption{解图}{离地瞬间只保留弹力和重力；支持力降为零。}\label{ch07:fig:soldiagFourSixteen}\end{minipage}\end{center}}

\newcommand{\probdiagFourSeventeen}{%
\begin{center}\begin{minipage}{.97\linewidth}\centering
\begin{tikzpicture}[scale=.72]
  \draw[solbody] (-4.4,-.42)--(4.4,-.42);
  \shade[ball color=PaleBlue] (-2.25,0) circle (.42);
  \shade[ball color=SoftGray] (1.65,0) circle (.42);
  \node[sollabel,above] at (-2.25,.55) {$A,m_A$};
  \node[sollabel,above] at (1.65,.55) {$B,m_B$};
  \draw[solflow] (-1.72,0)--(-.50,0) node[midway,above,sollabel] {$v_0$};
  \node[sollabel] at (1.65,-.02) {$0$};
\end{tikzpicture}
\diagcaption{题意图}{一维完全弹性碰撞同时满足动量守恒和动能守恒。}\label{ch07:fig:probdiagFourSeventeen}\end{minipage}\end{center}}

\newcommand{\soldiagFourSeventeen}{%
\begin{center}\begin{minipage}{.97\linewidth}\centering
\begin{tikzpicture}[scale=.72]
  \node[solpanel] at (-2.60,1.10) {碰前};
  \draw[solflow] (-3.75,0)--(-1.45,0) node[midway,above,sollabel] {$v_0$};
  \node[solpanel] at (2.15,1.10) {碰后（$m_A>m_B$）};
  \draw[solflow] (.45,.25)--(1.55,.25) node[midway,above,sollabel] {$v_A$};
  \draw[solflow] (.45,-.45)--(3.60,-.45) node[midway,below,sollabel] {$v_B$};
\end{tikzpicture}
\diagcaption{解图}{非平凡碰撞后两速度都向右要求 $m_A>m_B$，且被撞球 $B$ 的速度更大。}\label{ch07:fig:soldiagFourSeventeen}\end{minipage}\end{center}}

\newcommand{\probdiagFourEighteen}{%
\begin{center}\begin{minipage}{.97\linewidth}\centering
\begin{tikzpicture}[scale=.9]
 \coordinate (A) at (0,0); \coordinate (B) at (.64,.48);
 \draw[solgeom] (-2.4,0)--(1.9,0);
 \draw[solbody,fill=SoftGray] (A) circle (.4);
 \draw[solbody,fill=SoftGray] (B) circle (.4);
 \draw[solgeom] (A)--(B);
 \node[sollabel,below=15pt] at (A) {$A$}; \node[sollabel,above=15pt] at (B) {$B$};
 \draw[solflow] (-2.1,0)--(-.5,0) node[midway,above=5pt,sollabel] {$\vec v_0$};
 \node[sollabel] at (.2,1.6) {俯视：接触瞬间；圆心连线为法线};
\end{tikzpicture}
\diagcaption{题意图}{等质量小球在光滑平面内发生极短的完全弹性碰撞。}\label{ch07:fig:probdiagFourEighteen}\end{minipage}\end{center}}

\newcommand{\soldiagFourEighteen}{%
\begin{center}\begin{minipage}{.97\linewidth}\centering
\begin{tikzpicture}[scale=.8]
 % Use the same laboratory directions as the contact figure: n=(.8,.6).
 \coordinate (O) at (0,1.92); \coordinate (A) at (1.44,0); \coordinate (B) at (4,1.92);
 \draw[solforce] (O)--(A) node[midway,below left=5pt,sollabel] {$\vec v_A$};
 \draw[solforce] (A)--(B) node[midway,below right=5pt,sollabel] {$\vec v_B$};
 \draw[solforce] (O)--(B) node[midway,above=5pt,sollabel] {$\vec v_0$};
 \draw[TextGray] (1.26,.24)--(1.5,.42)--(1.68,.18);
 \node[sollabel,anchor=west,align=left] at (4.8,1)
  {动量：首尾相接\\动能：交叉项为零\\两非零速度互相垂直};
\end{tikzpicture}
\diagcaption{解图}{沿用接触图的方向：$\vec v_0$ 向右，$\vec v_B$ 沿圆心连线，$\vec v_A\perp\vec v_B$。}\label{ch07:fig:soldiagFourEighteen}\end{minipage}\end{center}}

\newcommand{\probdiagFourNineteen}{%
\begin{center}\begin{minipage}{.97\linewidth}\centering
\begin{tikzpicture}[scale=.8]
 \begin{scope}[rotate=20]
  \draw[solbody] (-3,0)--(3,0);
  \draw[solbody,fill=SoftGray] (1.1,0) rectangle (2,.65);
  \draw[solgeom] (-2.2,0) rectangle (-1.3,.65);
  \draw[solgeom] (-1.75,.8)--(-1.75,1.55) (1.55,.8)--(1.55,1.55);
  \draw[<->,TextGray] (-1.75,1.45)--(1.55,1.45) node[midway,above=4pt,sollabel] {$L$};
  \draw[solforce] (1.55,.325)--(2.85,.325) node[above right=5pt,sollabel] {$f$};
  \draw[solflow] (.8,.75)--(-.4,.75) node[midway,above=4pt,sollabel] {下滑};
 \end{scope}
 \coordinate (Q) at ({-3*cos(20)},{-3*sin(20)});
 \draw[TextGray] (Q)--++(1.1,0);
 \draw[TextGray] ($(Q)+(.7,0)$) arc (0:20:.7);
 \node[sollabel] at ($(Q)+(1.1,.23)$) {$\theta$};
\end{tikzpicture}
\diagcaption{题意图}{沿斜面方向，重力分力做正功，滑动摩擦力做负功。}\label{ch07:fig:probdiagFourNineteen}\end{minipage}\end{center}}

\newcommand{\probdiagFourTwenty}{%
\begin{center}\begin{minipage}{.97\linewidth}\centering
\begin{tikzpicture}[scale=.78]
 \draw[solbody] (-3.8,0)--(3.8,0) (-3.6,0)--(-3.6,1.4);
 \draw[solgeom] (-.45,0) rectangle (.45,.8);
 \draw[decorate,decoration={coil,segment length=5pt,amplitude=3pt}] (-3.6,.4)--(1.55,.4);
 \draw[solbody,fill=SoftGray] (1.55,0) rectangle (2.45,.8);
 \draw[solgeom] (0,-.65)--(0,1.15) (2,-.65)--(2,1.15);
 \draw[<->,TextGray] (0,-.5)--(2,-.5) node[midway,below=4pt,sollabel] {$a$};
 \node[sollabel,above=5pt,anchor=south east] at (0,1.15) {$O$：自然长度时的中心};
 \draw[solflow] (2.6,1.65)--(1.4,1.65) node[midway,above=5pt,sollabel] {释放后向左};
\end{tikzpicture}
\diagcaption{题图}{弹簧在 $O$ 点为自然长度；物块从相对 $O$ 向右位移 $a$ 的位置由静止释放。}\label{ch07:fig:probdiagFourTwenty}\end{minipage}\end{center}}

\newcommand{\soldiagFourTwenty}{%
\begin{center}\begin{minipage}{.97\linewidth}\centering
\begin{tikzpicture}[scale=.67]
  \node[solpanel] at (-3.55,1.55) {运动中};
  \draw[axis] (-5.0,0)--(-2.15,0) node[right,sollabel] {$x$};
  \draw[solbody,fill=SoftGray] (-3.75,-.35) rectangle (-3.05,.35);
  \draw[solforce] (-3.40,.65)--(-4.45,.65) node[left,sollabel] {$F_s$};
  \draw[solforce] (-3.40,-.65)--(-2.35,-.65) node[right,sollabel] {$f_k$};
  \node[sollabel] at (-3.40,-1.02) {向左运动、$x>0$};
  \node[solpanel] at (1.95,1.55) {停止后（图示 $x>0$）};
  \draw[axis] (.35,0)--(3.85,0) node[right,sollabel] {$x$};
  \draw[solbody,fill=SoftGray] (1.60,-.35) rectangle (2.30,.35);
  \draw[solforce] (1.95,.65)--(.90,.65) node[left,sollabel] {$F_s$};
  \draw[solforce] (1.95,-.65)--(3.00,-.65) node[right,sollabel] {$f_s$};
  \node[sollabel] at (1.95,-1.02) {$f_s$ 按平衡需要取值};
\end{tikzpicture}
\diagcaption{解图}{滑动摩擦反向于速度；静摩擦则反向于物体停止位置处的弹簧力。}\label{ch07:fig:soldiagFourTwenty}\end{minipage}\end{center}}


动量关系之外，碰撞与弹簧问题还需要能量关系。这里先给出所需的基本定义；功和能量的
系统讨论在后续章节展开。质点的动能为 $E_k=mv^2/2$，力在位移上的功为
$W=\int\vec F\cdot\dd\vec r$。将 $m\dot{\vec v}=\vec F$ 与速度点乘并对时间积分，
得到 $\Delta E_k=W$。如果一类力的功只取决于起终点，就称为保守力，并定义其势能变化
$\Delta E_p=-W$；于是只受保守力时，$E_k+E_p$ 不变。近地面重力势能可写成 $mgy$，
理想弹簧势能为 $k(\ell-l_0)^2/2$，零点分别取 $y=0$ 和自然长度。

现在考察由多个质点构成的力学系统。
系统内每个质点除受外力外，还会受到其他质点施加的、满足牛顿第三定律的内力。作用力与
反作用力的矢量和为零，所以内力不能改变系统总动量；但两力作用点的位移通常不同，它们的
功并不必然相互抵消，因此内力完全可能改变系统的总动能。这正是动量问题与能量问题在处理
内力时的重要区别。

内力做功有两种性质截然不同的情形。若质点间的内力是保守力，内力功只与质点间的相对位置
有关；系统动能的减少会以内部势能的形式储存在相对位置中，并可在以后的运动中重新释放为
各质点的动能。若内力像滑动摩擦力一样，其功与相对运动的具体过程有关，那么机械能会转化
为内能等形式，在单纯的力学过程中不能自行完整地转回宏观动能。这里所谓“机械能损失”并不
是能量从自然界消失，而是能量形式发生了转化。

当系统的内力和外力均为保守力时，可把系统总机械能定义为所有质点的动能、外力对应的势能
以及系统内部质点之间的势能之和：
\[
 E=\sum_i E_{k,i}+E_{p,\text{外}}+E_{p,\text{内}},
\]
在势能没有显含时间、且无其他外部能量输入的模型中，运动过程中 $E$ 保持不变。

\begin{origexamplebox}{}\label{ch07:ex:s12-4-15}
光滑水平面上有一根无质量的弹簧，劲度系数为 $k$，原长为 $l_0$。它的两端分别连接质量
为 $m_1,m_2$ 的两个物体。
\begin{enumerate}[label=\arabic*.,leftmargin=2.2em]
 \item 初始时刻两物体的速度均为零，相距 $l<l_0$。求释放以后两物体之间的最大距离。
 \item 取向右为正方向，若初始时刻两物体的速度分别为 $v_1,v_2$，求此后运动过程中
 两物体的最近和最远距离。
\end{enumerate}
图~\ref{ch07:fig:probdiagFourFifteen}中，弹簧只改变两物体的相对运动；整体质心速度保持不变。
\probdiagFourFifteen
\end{origexamplebox}

\Needspace{6\baselineskip}\begin{xesolution}
记两物体之间的距离为 $r=x_2-x_1$，总质量和约化质量分别为
\[
 M=m_1+m_2,
 \qquad
 \mu_r=\frac{m_1m_2}{m_1+m_2}.
\]
质心速度
\[
 V=\frac{m_1v_1+m_2v_2}{M}
\]
保持不变。由 $\dot r=\dot x_2-\dot x_1$ 和质心速度的定义，先解出
\[
 \dot x_1=V-\frac{m_2}{M}\dot r,\qquad
 \dot x_2=V+\frac{m_1}{M}\dot r.
\]
代入两物体动能，交叉项系数 $-m_1m_2/M+m_2m_1/M=0$；相对速度平方项系数为
$(m_1m_2^2+m_2m_1^2)/(2M^2)=\mu_r/2$。因此
\[
 \frac12m_1\dot x_1^2+\frac12m_2\dot x_2^2
 =\frac12MV^2+\frac12\mu_r\dot r^{\,2},
\]
系统机械能可分成不变的质心平动能和相对运动能：
\[
 E=\frac12MV^2+\frac12\mu_r\dot r^{\,2}
   +\frac12k(r-l_0)^2.
\]

\textbf{(1)} 初始相对速度为零，弹簧的初始压缩量为
$\delta=l_0-l$。在最大距离处相对速度再次为零，因而
\[
 \frac12k(l-l_0)^2=\frac12k(r_{\max}-l_0)^2.
\]
除初始转折点 $r=l$ 外，另一转折点位于原长位置的对侧：
\[
 \boxed{r_{\max}=2l_0-l}.
\]

\textbf{(2)} 初始相对速度为 $\dot r_0=v_2-v_1$。相对运动的总能量为
\[
 E_{\rm rel}
 =\frac12\mu_r(v_2-v_1)^2+\frac12k(l-l_0)^2.
\]
最近和最远距离都是相对运动的转折点，满足 $\dot r=0$。令
\[
 A=\sqrt{(l-l_0)^2+\frac{\mu_r}{k}(v_2-v_1)^2},
\]
则
\[
 \boxed{r_{\min}=l_0-A},
 \qquad
 \boxed{r_{\max}=l_0+A}.
\]
图~\ref{ch07:fig:soldiagFourFifteen}中，相对势能关于原长位置对称，同一相对能量给出最近与最远两个转折点。
\soldiagFourFifteen
两物体的共同平动速度只改变质心运动，不改变它们之间的距离；所以结果只含相对速度
$v_2-v_1$。上式默认 $A\le l_0$ 且物体可视为质点；若模型中物体不能相互穿越，而计算
得到 $A>l_0$，则在弹簧长度降为零以前已发生接触，最近距离须由接触条件另行确定。
\end{xesolution}

\Needspace{8\baselineskip}
\begin{origexamplebox}{}\label{ch07:ex:s12-4-16}
质量分别为 $m_1,m_2$ 的两个物体 $A,B$ 用原长为 $l_0$、劲度系数为 $k$ 的轻弹簧连接。
初始时刻 $A$ 放在 $B$ 的上方且均静止，将弹簧压缩至 $l<l_0$。求系统释放以后 $B$
能够离开地面的条件。
图~\ref{ch07:fig:probdiagFourSixteen}中，在 $B$ 尚未离地的阶段，下端固定，只有上方物体 $A$ 运动。
\probdiagFourSixteen
\end{origexamplebox}

\Needspace{6\baselineskip}\begin{xesolution}
在 $B$ 尚未离地时，它受地面约束而保持静止，弹簧下端等效固定。设某时刻弹簧长度为
$x$，上方物体 $A$ 的速率为 $u$。以 $B$ 所在高度为重力势能零点，机械能守恒给出
\[
 \frac12m_1u^2+m_1gx+\frac12k(x-l_0)^2
 =m_1gl+\frac12k(l-l_0)^2.
\]
物体 $A$ 到达最高点时 $u=0$。去掉代表初始位置的根 $x=l$，另一转折位置为
\[
 x_{\max}=2l_0-l-\frac{2m_1g}{k}.
\]

当弹簧已经伸长时，它对 $B$ 的拉力向上。$B$ 刚要离地的临界条件是地面支持力
$N=0$，即
\[
 k(x-l_0)=m_2g,
 \qquad
 x_{\rm c}=l_0+\frac{m_2g}{k}.
\]
图~\ref{ch07:fig:soldiagFourSixteen}中，离地瞬间只保留弹力和重力；支持力降为零。
\soldiagFourSixteen
要使 $B$ 真正获得离地速度，$A$ 在其上升过程中必须越过这一临界位置，故
\[
 x_{\max}>x_{\rm c}.
\]
代入两式并整理，得到
\[
 \boxed{k(l_0-l)>(2m_1+m_2)g}.
\]
等号只表示 $A$ 恰在最高点时使 $N$ 瞬间降为零；随后弹簧拉力立即减小，$B$ 不会真正
上升。因此若“离开地面”要求产生向上的运动，应取严格不等号。
\end{xesolution}

\Needspace{8\baselineskip}
\begin{origexamplebox}{}\label{ch07:ex:s12-4-17}
在光滑平面上有两个球 $A,B$，质量分别为 $m_A,m_B$。开始时刻 $B$ 球静止，$A$ 球以
速度 $v_0$ 撞向 $B$。如果碰撞之后两球运动速度均与 $A$ 的初速度方向相同，且碰撞过程
极短、无机械能损失，求碰后两球的速度。
图~\ref{ch07:fig:probdiagFourSeventeen}中，一维完全弹性碰撞同时满足动量守恒和动能守恒。
\probdiagFourSeventeen
\end{origexamplebox}

\Needspace{6\baselineskip}\begin{xesolution}
取 $A$ 的初速度方向为正，设碰后速度分别为 $v_A,v_B$。碰撞时间极短，外力冲量可忽略，
故动量守恒；又因无机械能损失，动能守恒：
\[
 m_Av_0=m_Av_A+m_Bv_B,
\]
\[
 \frac12m_Av_0^2
 =\frac12m_Av_A^2+\frac12m_Bv_B^2.
\]
由动量式写成 $m_A(v_0-v_A)=m_Bv_B$，动能式则可因式分解为
\[
 m_A(v_0-v_A)(v_0+v_A)=m_Bv_B^2.
\]
排除“没有发生速度交换”的平凡情形后，两式相除得到
\[
 v_0+v_A=v_B.
\]
再与动量式联立，得到一维完全弹性碰撞公式
\[
 \boxed{v_A=\frac{m_A-m_B}{m_A+m_B}v_0},
 \qquad
 \boxed{v_B=\frac{2m_A}{m_A+m_B}v_0}.
\]
图~\ref{ch07:fig:soldiagFourSeventeen}中，非平凡碰撞后两速度都向右要求 $m_A>m_B$，且被撞球 $B$ 的速度更大。
\soldiagFourSeventeen
题设要求两球碰后都沿正方向运动，因此必须有 $m_A>m_B$；若 $m_A=m_B$，则 $A$ 停止、
$B$ 以 $v_0$ 前进，不能说 $A$ 的速度仍具有原方向。
\end{xesolution}

\Needspace{8\baselineskip}
\begin{origexamplebox}{}\label{ch07:ex:s12-4-18}
光滑平面上有两个质量相同的球 $A,B$。$A$ 球以某一初速度 $\vec v_0$ 撞向静止的 $B$，
碰撞时间极短且无机械能损失。求证：无论 $A$ 球以何种方式弹开，$B$ 球的速度方向总与
$A$ 球弹开的方向垂直。
图~\ref{ch07:fig:probdiagFourEighteen}中，等质量小球在光滑平面内发生极短的完全弹性碰撞。
\probdiagFourEighteen
\end{origexamplebox}

\Needspace{6\baselineskip}\begin{xesolution}
设两球质量均为 $m$，碰后速度分别为 $\vec v_A,\vec v_B$。在平面内必须使用矢量动量
守恒：
\begin{equation}\label{ch07:eq:equal-ball-momentum}
 \vec v_0=\vec v_A+\vec v_B.
\end{equation}
动能守恒给出
\begin{equation}\label{ch07:eq:equal-ball-energy}
 v_0^2=v_A^2+v_B^2.
\end{equation}
对式\eqref{ch07:eq:equal-ball-momentum}两边取矢量自身的点积：
\[
 v_0^2=v_A^2+v_B^2+2\vec v_A\cdot\vec v_B.
\]
与式\eqref{ch07:eq:equal-ball-energy}比较，得到
\[
 \vec v_A\cdot\vec v_B=0.
\]
因此在两球碰后速度均不为零的非退化碰撞中
\[
 \boxed{\vec v_A\perp\vec v_B}.
\]
图~\ref{ch07:fig:soldiagFourEighteen}沿用接触图的方向：$\vec v_0$ 向右，$\vec v_B$ 沿两球圆心连线指向右上方，$\vec v_A$ 指向右下方。将 $\vec v_A$ 与 $\vec v_B$ 首尾相接，合矢量就是 $\vec v_0$；动能关系使两条直角边互相垂直。
\soldiagFourEighteen
这也说明三条速度矢量构成以 $\vec v_0$ 为斜边的直角三角形。若其中一球碰后静止，则零
速度没有确定方向，应把它看作上述结论的退化极限。
\end{xesolution}

\subsection{机械能的损失：摩擦}

当物体所受的力中含有非保守力时，机械能不再守恒。对于单个质点从 $A$ 到 $B$ 的运动，
把外力分成保守力和非保守力，由动能定理得
\begin{equation}
 \frac12mv_B^2-\frac12mv_A^2=W_C+W_{NC},
\end{equation}
其中 $W_C,W_{NC}$ 分别表示保守力和非保守力的功。保守力做功等于势能的减少量：
\begin{equation}
 W_C=-E_p(B)+E_p(A).
\end{equation}
代入并整理可得
\begin{equation}\label{ch07:eq:nonconservative-work}
 W_{NC}
 =\left(\frac12mv_B^2+E_p(B)\right)
  -\left(\frac12mv_A^2+E_p(A)\right)
 =E(B)-E(A).
\end{equation}
也就是说，\textbf{非保守力的功等于系统机械能的变化量}。摩擦力通常做负功，使机械能
减小；其他非保守力也可能做正功，使物体的机械能增加。

\Needspace{8\baselineskip}
\begin{origexamplebox}{}\label{ch07:ex:s12-4-19}
在一个物体沿粗糙斜面下滑的过程中验证式\eqref{ch07:eq:nonconservative-work}。物体由静止出发，沿动摩擦因数为
$\mu$、倾角为 $\theta$ 的粗糙斜面下滑。求沿斜面下滑距离为 $L$ 时的速度。
图~\ref{ch07:fig:probdiagFourNineteen}中，沿斜面方向，重力分力做正功，滑动摩擦力做负功。
\probdiagFourNineteen
\end{origexamplebox}

\Needspace{6\baselineskip}\begin{xesolution}
\begin{qualityfigure}
\begin{tikzpicture}[>=Latex,font=\small]
\coordinate(A)at(0,2.6);\coordinate(B)at(5,0);\draw[very thick](A)--(B);\draw[dashed](0,0)--(B);
\fill(A)circle(2pt)node[above]{$A$};\fill(B)circle(2pt)node[right]{$B$};\draw[<->](-.3,2.2)--(4.7,-.4)node[midway,below]{$L$};
\draw[<->](-.8,0)--(-.8,2.6)node[midway,left]{$L\sin\theta$};\draw(4.1,0)arc(180:152.526:.9)node[midway,left]{$\theta$};
\draw[->,DeepBlue,thick](2,1.56)--(2.9,1.09)node[right]{$mg\sin\theta$};\draw[->,Gold,thick](2,1.56)--(1.1,2.03)node[left]{$f_k$};
\end{tikzpicture}
\figcap{粗糙斜面上的位移、高度差和功的符号（解答用辅助图）}\label{mech:07:119:a}
\end{qualityfigure}
图~\ref{mech:07:119:a}中，$A\to B$ 沿斜面向下，重力切向分量同向、摩擦力反向；高度差为 $L\sin\theta$，据此分别计算重力势能变化与摩擦功。

物体下降的竖直高度为 $L\sin\theta$，支持力为 $N=mg\cos\theta$，滑动摩擦力的功为
\[
 W_{NC}=W_f=-\mu mg\cos\theta\,L.
\]
取末位置为重力势能零点，则初、末机械能分别为
\[
 E_A=mgL\sin\theta,
 \qquad
 E_B=\frac12mv^2.
\]
代入式\eqref{ch07:eq:nonconservative-work}：
\[
 -\mu mgL\cos\theta
 =\frac12mv^2-mgL\sin\theta.
\]
因此
\[
 \boxed{v=\sqrt{2gL(\sin\theta-\mu\cos\theta)}}.
\]
物体若要由静止开始下滑且不断加速，须有 $\sin\theta>\mu\cos\theta$；若还要精确判断
能否启动，应另用静摩擦因数检验初始临界条件。
\end{xesolution}

\Needspace{8\baselineskip}
\begin{origexamplebox}{}\label{ch07:ex:s12-4-20}
如图所示，一轻弹簧一端固定在墙上，另一端连接一小物块。小物块放在摩擦因数为 $\mu$
的水平面上；弹簧处于自然状态时，小物块中心位于 $O$ 点。现用手将小物块从 $O$ 向右移动距离 $a>0$，
然后由静止释放，发现小物块开始向左移动。取最大静摩擦因数与滑动摩擦因数均为 $\mu$。
判断下列说法是否正确：
\begin{enumerate}[label=\Alph*.,leftmargin=2.2em]
 \item 小物块可能停在 $O$ 点；
 \item 小物块停止以后所受的摩擦力必不为零；
 \item 小物块无论停在 $O$ 点左边还是右边，停前所受摩擦力的方向和停止后所受摩擦力的
 方向，两者既可能相同也可能相反；
 \item 小物块在通过 $O$ 点后向右运动直到最远处的过程中，速度的大小总是减小；小物块
 由右边最远处回到 $O$ 点的过程中，速度大小总是增加。
\end{enumerate}
图~\ref{ch07:fig:probdiagFourTwenty}中，虚线物块表示弹簧自然长度时的位置；标注 $a$ 量取的是两个位置的中心间距。
\probdiagFourTwenty
\end{origexamplebox}

\Needspace{6\baselineskip}\begin{xesolution}
令 $x$ 为物块相对 $O$ 点的坐标，向右为正，并记
\[
 d=\frac{\mu mg}{k}.
\]
物块开始向左移动说明初始弹簧力足以克服最大静摩擦力；在题目用同一摩擦因数描述时，
至少有 $ka>\mu mg$。

\textbf{A．正确。}\quad 从 $x=a$ 向左运动到 $O$ 的过程中，滑动摩擦力做功
$-\mu mga$。若物块恰好在 $O$ 停下，则
\[
 \frac12ka^2-\mu mga=0,
\]
非零解为
\[
 a=2d=\frac{2\mu mg}{k}.
\]
此时 $O$ 点弹簧力为零，静摩擦力也可为零，物块能够停在该处。

\textbf{B．错误。}\quad A 项已经给出反例：物块停在 $O$ 点时，所需静摩擦力为零。

\textbf{C．正确。}\quad 物块从 $x=a$ 向左运动后的第一个停止位置 $x_1$ 满足
\[
 \frac12ka^2-\mu mg(a-x_1)=\frac12kx_1^2,
\]
除去初始位置，得到
\[
 x_1=2d-a.
\]
若 $d<a<2d$，则 $0<x_1<d$：物块停在 $O$ 点右侧，停前、停后的摩擦力均向右，两者
同向。若 $2d<a<3d$，则 $-d<x_1<0$：物块停在左侧，停前摩擦力向右，停后静摩擦力
向左，两者反向。

若 $a>3d$，物块到达 $x_1<-d$ 后还会反向。下一停止位置由同样的功—能关系得到
\[
 x_2=a-4d.
\]
当 $3d<a<4d$ 时，物块停在左侧且停前向右运动，停前、停后的摩擦力均向左，两者同向；
当 $4d<a<5d$ 时，物块停在右侧且停前向右运动，停前摩擦力向左、停后静摩擦力向右，
两者反向。因此，无论最后停在 $O$ 点左侧还是右侧，同向和反向都确有可能。

\textbf{D．错误。}\quad 物块从 $O$ 向右运动时，弹簧力与滑动摩擦力都向左，速度确实
一直减小。若物块能从右侧最远点返回，向左运动时弹簧力向左、摩擦力向右：在 $x>d$
时合力向左，速率增加；进入 $0<x<d$ 后合力向右，速率反而减小，并非直到 $O$ 都增加。
图~\ref{ch07:fig:soldiagFourTwenty}中，滑动摩擦反向于速度；静摩擦则反向于物体停止位置处的弹簧力。
\soldiagFourTwenty
所以正确选项为
\[
 \boxed{\text{A、C}}.
\]
本题中“摩擦力反向”不是由位置本身决定的：滑动时它反向于速度，静止后则按维持平衡所
需的方向和大小取值。
\end{xesolution}

对于更复杂过程中机械能怎样变化，也可以沿用“先区分保守力与非保守力，再考察各段功和
能量变化”的方法逐段分析。

\subsection{碰撞阶段、能量损失与摩擦生热}

碰撞过程持续时间很短，通常先用动量守恒联系碰撞前后，再用恢复系数或能量条件描述材料性质。
完全弹性碰撞同时守恒动量和动能；完全非弹性碰撞只表示碰后没有相对速度，并不意味着系统能量
消失。对质量 $m_1,m_2$ 的一维完全非弹性碰撞，动能损失为
\begin{equation*}
 \Delta K_{\rm loss}=\frac12\mu(u_1-u_2)^2,
 \qquad \mu=\frac{m_1m_2}{m_1+m_2},
\end{equation*}
可见损失只由碰前相对速度和约化质量决定。若碰撞后还有弹簧压缩、物块爬升或摆动，应把
“碰撞瞬间”和“碰后运动”分成两个阶段，分别使用适合的守恒关系。

滑动摩擦把宏观机械能转化为接触物体的内能。两个接触面相对滑动距离为 $s_{\rm rel}$ 时，产生的
热量为
\begin{equation*}
 Q=f_k s_{\rm rel}.
\end{equation*}
这里必须使用相对路程，而不是任一物体对地路程。若研究对象只取其中一个物体，摩擦功可以为正
也可以为负；若把两个物体及接触面一起纳入系统，摩擦内力总功等于机械能减少量的负值，而内能
增加量为正。区分“单个物体的摩擦功”和“系统的摩擦生热”，可以避免传送带、叠块等问题中常见的
符号错误。
}

\endgroup

% ThermalQualityRevision20261023-v1
% QualityRevision20261023
\begingroup
\tcbset{breakable/.default=false}
\chapter{功、能量和守恒定律}
% Editable TikZ: work, frames, and energy system boundaries.
\providecommand{\ThermalFigure}[3]{\begin{qualityfigure}
\begin{tikzpicture}[>=Latex,font=\small,x=.8cm,y=.8cm]#1\end{tikzpicture}
\figcap{#2}\label{#3}\end{qualityfigure}}
\newcommand{\ThQSlidingBoard}{\ThermalFigure{
 \draw[thick](-1,0)--(7,0);\draw[fill=DeepBlue!10,thick](0,0)rectangle(6,.5);
 \draw[fill=DeepBlue!20,thick](.7,.5)rectangle(2,1.4);\node at(1.35,.95){$A,m$};\node at(4,.25){$B,M$};
 \draw[->,thick](2.3,1)--(4,1)node[right]{$v_0$};\node[below]at(3,-.4){板与地面光滑；物块与板间摩擦因数 $\mu$};
}{初态物块向右滑上静止长板，两个物体的位移须分别计算。}{ch08:quality:board-q}}
\newcommand{\ThTrainBase}[1]{
 \draw[thick](0,0)rectangle(6,2);\draw[fill=DeepBlue!12](1,.2)rectangle(2,1.1);\node at(1.5,.65){$m$};
 \draw[->,thick](2,1)--(3.6,1)node[right]{$F$};\draw[->](6.5,1)--(8,1)node[right]{#1};
 \draw[<->](1,-.6)--node[below]{车内位移 $L$}(4.5,-.6);}
\newcommand{\ThQTrain}{\ThermalFigure{\ThTrainBase{$v_0$ 恒定}}{拉力及物块相对车厢的位移；地面系的作用点位移还包含车厢前进量。}{ch08:quality:train-q}}
\newcommand{\ThQAcceleratingTrain}{\ThermalFigure{\ThTrainBase{加速度 $a$}
 \node[above]at(3,2.4){初态物块相对车厢静止};
}{真实拉力在两参考系中相同；在加速车厢内列动力学方程需另处理惯性项。}{ch08:quality:accelerating-q}}
\newcommand{\ThRadialWorkBase}[1]{
 \fill(0,0)circle(3pt)node[below]{$O$（固定物体）};\draw[->](-.5,0)--(7,0)node[right]{$r$};
 \fill[DeepBlue](2,0)circle(2pt)node[above]{$A$};\fill[DeepBlue](5,0)circle(2pt)node[above]{$B$};
 \draw[<->](0,-.7)--node[below]{$r_A$}(2,-.7);\draw[<->](0,-1.5)--node[below]{$r_B$}(5,-1.5);
 \node[above]at(4,1){#1};}
\newcommand{\ThQInverseForce}{\ThermalFigure{\ThRadialWorkBase{吸引力大小 $k/r$}}{始末距离从同一固定力心量起；图示 $r_B>r_A$，公式也应适用于相反方向。}{ch08:quality:inverse-q}}
\newcommand{\ThQInverseSquareForce}{\ThermalFigure{\ThRadialWorkBase{吸引力大小 $k/r^2$}}{平方反比吸引力的始末位置；不要把路程与有符号径向位移混淆。}{ch08:quality:inverse-square-q}}
\newcommand{\ThQConservative}{\ThermalFigure{
 \draw[->](0,0)--(3,0)node[right]{$x$};\fill(1,0)circle(2pt);\draw[->,DeepBlue](1,0)--(2.5,0);
 \node[above,align=center]at(1.5,1){一维三种力\\$at,\ -\alpha v,\ -kx^3$};
 \begin{scope}[shift={(7,0)}]\fill(0,0)circle(2pt)node[below]{$O$};\fill(2,1)circle(2pt)node[right]{$P$};
 \draw[dashed](0,0)--node[above]{$r$}(2,1);\draw[->,DeepBlue,thick](2,1)--(1,.5)node[below right]{$\vect f(r)$};\end{scope}
}{分别辨认力依赖时间、速度还是位置；径向吸引力的适用域还须避开 $r=0$。}{ch08:quality:conservative-q}}
\newcommand{\ThStoppingBase}[1]{
 \draw[thick](-.5,0)--(7,0);\foreach\x in {0,.5,...,6.5}{\draw[TextGray](\x,0)--++(-.2,-.2);}
 \draw[fill=DeepBlue!12,thick](0,0)rectangle(1.3,1);\node at(.65,.5){$m$};
 \draw[->,thick](1.4,.6)--(3,.6)node[right]{$v_0$};
 \draw[dashed](5,0)rectangle(6.3,1);\node[above]at(5.65,1.1){停止};\draw[<->](0,-.7)--node[below]{$s$}(5,-.7);#1}
\newcommand{\ThQStopping}{\ThermalFigure{\ThStoppingBase{}}{同一标志点从初态到停下前的位移为 $s$。}{ch08:quality:stopping-q}}
\newcommand{\ThAStopping}{\ThermalFigure{\ThStoppingBase{
 \draw[->,DeepBlue,thick](.65,.5)--(-1,.5)node[left]{$f_k$};
 \draw[dashed](-1.5,-1.2)rectangle(7.2,1.6);\node[above]at(3,1.7){大系统边界：滑块＋地面};
}}{摩擦力对滑块做负功；扩大系统后，摩擦引起内部能量转化。}{ch08:quality:stopping-a}
图~\ref{ch08:quality:stopping-a}的摩擦力与向右位移相反，因此滑块所受摩擦功为 $-f_ks$。虚线把地面也纳入系统后，这一作用成为内部作用，减少的宏观动能转为内能；改变系统边界不会改变停车距离。
}
\newcommand{\ThARoughSpring}{\ThermalFigure{
 \draw[->](-.5,0)--(6,0)node[right]{$x$};\draw[->](0,0)--(0,3.4)node[above]{$K(x)$};
 \draw[DeepBlue,thick,domain=0:4,samples=70]plot(\x,{2.4-.15*\x^2-.05*(4-\x)});
 \fill(4,0)circle(1.5pt)node[below]{$a$};\node[right]at(4.4,2.8){示意：$\mu_kmg<ka/2$};
 \node[right]at(4.4,1.9){仅在实际滑动路径上取 $K\ge0$};
 \draw[->,thick](3.4,1.1)--(2.8,1.65);
}{用能量曲线检查运动能否走到原长位置；启动仍须先通过静摩擦条件。}{ch08:quality:rough-spring-a}
图~\ref{ch08:quality:rough-spring-a}把释放点 $x=a$ 的动能取为零。向左滑动时，弹簧释放能量但摩擦持续耗散，故 $K(x)=\tfrac12k(a^2-x^2)-\mu_kmg(a-x)$。只有在运动已启动的前提下，检查 $K(0)$ 才能区分到达与以非零速度通过 $O$；静摩擦条件不能从这条滑动曲线中省掉。
}
\newcommand{\ThBulletBase}[1]{
 \fill[DeepBlue](0,.6)circle(3pt)node[above]{$m$};\draw[->,thick](.3,.6)--(1.7,.6)node[midway,above]{$u$};
 \draw[thick](2.2,0)rectangle(4,1.2);\node at(3.1,.6){$M$};\draw[->](4.5,.6)--(5.4,.6);
 \draw[fill=DeepBlue!12,thick](6,0)rectangle(8,1.2);\fill(6.4,.6)circle(3pt);\node at(7.2,.6){$M+m$};#1}
\newcommand{\ThQBullet}{\ThermalFigure{\ThBulletBase{}}{子弹留在木块中，两者随后共同运动；水平外力冲量在碰撞时可忽略。}{ch08:quality:bullet-q}}
\newcommand{\ThABullet}{\ThermalFigure{\ThBulletBase{
 \draw[->,thick](8.3,.6)--(9.5,.6)node[right]{$v$};
 \draw[dashed](-.5,-.5)rectangle(4.3,1.7);\node[below]at(2,-.7){同一系统的碰前状态};
}}{动量守恒跨越碰撞；内能增加由碰前、碰后宏观动能差确定。}{ch08:quality:bullet-a}
图~\ref{ch08:quality:bullet-a}中虚线把子弹和木块一起选为系统，内部碰撞力的冲量成对抵消。右图两者共用速度 $v$，所以末动能必须使用总质量 $M+m$；不能只计算木块动能而漏掉嵌入的子弹。
}



{
\graphicspath{{../../../sample/xes-original/tenth-section/}}
\section{功、功率与动能定理}
% Technical geometry and labels share one coordinate system per panel.
\tikzset{solbody/.style={draw=Ink,line width=.8pt},solpath/.style={draw=Ink,line width=.8pt},
 solforce/.style={-{Latex[length=2mm]},draw=DeepBlue,line width=.9pt},
 solflow/.style={-{Latex[length=1.8mm]},draw=Blue,line width=.8pt},
 solgeom/.style={draw=TextGray,densely dashed,line width=.5pt},
 sollabel/.style={font=\scriptsize,fill=none,inner sep=1.5pt},
 solpanel/.style={font=\scriptsize\bfseries,text=TextGray},
 spring/.style={decorate,decoration={coil,segment length=4pt,amplitude=2.5pt},draw=Ink}}
\newcommand{\diagcaption}[2]{\refstepcounter{figure}\par\nopagebreak[4]\smallskip
 {\centering\scriptsize\color{TextGray}图~\thefigure\quad #2\par}\smallskip}
\newenvironment{workfig}{\begin{center}\begin{minipage}{.98\linewidth}\centering}{\end{minipage}\end{center}}

\newcommand{\workDerivationDiagram}{%
\begin{workfig}\begin{tikzpicture}[scale=.9]
 \node[solpanel] at (1.6,2.8) {(a) 轨迹与切向速度};
 \draw[solpath,domain=0:2.7,samples=60] plot (\x,{.25*\x*\x});
 \coordinate (A) at (.5,.0625); \coordinate (B) at (1.5,.5625);
 \fill (A) circle (1.6pt) node[below,sollabel] {$A$};
 \fill (B) circle (1.6pt);\node[sollabel] at (1.4,1) {$B$};
 \draw[solforce] (A)--++(1,.25) node[below right,sollabel] {$\vec v_A$};
 \draw[solforce] (B)--++(1,.75) node[right,sollabel] {$\vec v_B$};
 \draw[solforce] (A)--++(0,1.3) node[left,sollabel] {$\vec a_A$};
 \draw[TextGray] ($(A)+(14.036:.45)$) arc(14.036:90:.45);
 \node[sollabel] at ($(A)+(53:.68)$) {$\theta$};
 \begin{scope}[xshift=5cm]
 \node[solpanel] at (1.2,2.8) {(b) 精确的速度矢量关系};
 \coordinate (P) at (0,0); \coordinate (Q) at (2,.5); \coordinate (S) at (2,1.5);
 \draw[solforce] (P)--(Q) node[midway,below,sollabel] {$\vec v_A$};
 \draw[solforce] (P)--(S) node[midway,above left,sollabel] {$\vec v_B$};
 \draw[solforce] (Q)--(S) node[midway,right,sollabel] {$\Delta\vec v$};
 \node[sollabel] at (1,-.7) {$\vec v_B=\vec v_A+\Delta\vec v$};
 \end{scope}
\end{tikzpicture}\diagcaption{示意图}{速度沿轨迹切线；速度增量是两个速度矢量之差。}
\label{ch08:fig:workDerivationDiagram}\end{workfig}}

\newcommand{\probdiagFourOne}{%
\begin{workfig}\begin{tikzpicture}[xscale=.82,yscale=.72]
 \draw[solgeom](-.25,2.8)--(11.55,2.8);\draw[solgeom](-.25,0)--(11.55,0);
 \foreach \x in {0,2.35,3.875,5.6,6.9}{\node[sollabel,above]at(\x,2.8){$A$};}
 \node[sollabel,above right]at(8.425,2.8){$A$};\node[sollabel,above left]at(10.525,2.8){$A$};
 \foreach \x in {0,.4,1.725,3.2,4.95,7.075,8.95}{\node[sollabel,below]at(\x,0){$B$};}
 \draw[solpath](0,2.8)--(0,0);
 \draw[solpath](2.35,2.8)--(.4,0);
 \draw[solpath](3.875,2.8)--(3.2,.93)--(1.725,0);
 \draw[solpath](5.6,2.8)--(5.1,2.075)--(5.025,1.575)--(4.675,1.4)--(4.575,1.025)--(4.15,1.025)--(3.975,.525)--(3.2,0);
 \draw[solpath](6.9,2.8)..controls(7.2,1.8)and(5.8,.6)..(4.95,0);
 \draw[solpath](7.075,0)--(8.25,4.1)--(8.425,2.8);
 \draw[solpath](8.95,0)..controls(9.2,5.0)and(11.25,5.2)..(10.525,2.8);
 \draw[<->,TextGray](11.35,0)--(11.35,2.8)node[midway,right,sollabel]{$h$};
\end{tikzpicture}\diagcaption{题图}{七条路径的始末高度差都为 $h$。}\label{ch08:fig:probdiagFourOne}\end{workfig}}

\newcommand{\soldiagFourOne}{%
\begin{workfig}\begin{tikzpicture}[scale=.9]
 \draw[solgeom] (-2,1.5)--(3,1.5);\draw[solgeom] (-2,-1)--(3,-1);
 \draw[solpath] (-1.5,1.5)..controls(-1.5,-.3)and(2,.5)..(2,-1);
 \fill (-1.5,1.5) circle(1.7pt) node[above,sollabel] {$A:y_A$};
 \fill (2,-1) circle(1.7pt) node[below,sollabel] {$B:y_B$};
 \draw[solforce] (-1,0)--(-1,-.8) node[left,sollabel] {$m\vec g$};
 \draw[<->,DeepBlue] (3,-1)--(3,1.5) node[midway,right,sollabel] {$h=y_A-y_B$};
\end{tikzpicture}\diagcaption{解图}{重力功的积分只保留竖直坐标的始末差。}\label{ch08:fig:soldiagFourOne}\end{workfig}}

\newcommand{\probdiagFourTwo}{%
\begin{workfig}\begin{tikzpicture}[scale=.9]
 \draw[solbody] (-2.2,0)--(3.8,0);
 \draw[solbody,fill=SoftGray] (-.5,0) rectangle (.5,.7);\node[sollabel] at (0,.35) {$m$};
 \draw[solforce] (.5,.55)--++(30:2) node[right,sollabel] {$F$};
 \draw[solgeom] (.5,.55)--(1.9,.55);
 \draw[TextGray] (1.2,.55) arc(0:30:.7);\node[sollabel] at (1.55,.81) {$\theta$};
 \node[sollabel,above] at (-1.5,0) {$\mu$};
 \draw[solflow] (-.5,-.5)--(2.5,-.5) node[midway,below,sollabel] {$L$};
\end{tikzpicture}\diagcaption{题图}{力 $F$ 向右上方，物体沿水平面向右滑动。}\label{ch08:fig:probdiagFourTwo}\end{workfig}}

\newcommand{\soldiagFourTwo}{%
\begin{workfig}\begin{tikzpicture}[scale=.9]
 \draw[solbody] (-2,0)--(2.7,0);\draw[solbody,fill=SoftGray] (-.5,0) rectangle(.5,.8);
 \draw[solforce] (.5,.6)--++(30:1.7) node[right,sollabel] {$F$};
 \draw[solforce] (0,.8)--(0,2) node[above,sollabel] {$N$};
 \draw[solforce] (0,.4)--(0,-1.1) node[below,sollabel] {$mg$};
 \draw[solforce] (-.5,.15)--(-1.7,.15) node[above left,sollabel] {$f=\mu N$};
 \draw[solflow] (1,-.55)--(2.5,-.55) node[midway,below,sollabel] {$\Delta\vec r$};
 \node[sollabel,align=left,anchor=west] at (3.5,1) {$W_F=FL\cos\theta$\\[5pt]$W_f=-fL$\\[5pt]$W_N=W_G=0$};
\end{tikzpicture}\diagcaption{解图}{竖直方向合力为零；各力作功须与水平位移配对。}\label{ch08:fig:soldiagFourTwo}\end{workfig}}

\newcommand{\probdiagFourThree}{%
\begin{workfig}\begin{tikzpicture}[scale=.85]
 \coordinate(O) at (0,1.8);\coordinate(A) at (2.4,1.8);\coordinate(B) at (0,-1.4);
 \draw[spring] (O)--(A);\draw[spring] (O)--(B);
 \fill (O) circle(2pt) node[above left,sollabel] {$O$};\fill(A)circle(3pt)node[right,sollabel]{$A$};
 \fill(B)circle(3pt)node[below,sollabel]{$B$};
 \node[sollabel] at (1.2,2.15) {$l_0,\ k$};\node[sollabel] at (2.7,1.35) {$m$};
 \draw[solgeom] (A)--(B);\node[sollabel,align=left,anchor=west] at (2.1,-.2) {虚线仅连接\!始末位置\\不表示轨迹};
 \draw[solgeom] (-1.25,-.6)--(.3,-.6);
 \draw[<->,TextGray] (-.8,-.6)--(-.8,1.8)node[midway,left,sollabel]{$l_0$};
 \draw[<->,TextGray] (-.8,-1.4)--(-.8,-.6)node[midway,left,sollabel]{$l_0/n$};
\end{tikzpicture}\diagcaption{题图}{两条弹簧线表示先后两个状态，并非同时存在两根弹簧。}\label{ch08:fig:probdiagFourThree}\end{workfig}}

\newcommand{\soldiagFourThree}{%
\begin{workfig}\begin{tikzpicture}[scale=.9]
 \draw[solgeom] (0,1.5)--(6,1.5);\draw[solgeom] (0,-.8)--(6,-.8);
 \fill (1,1.5)circle(2pt)node[above,sollabel]{$A$};\fill (4,-.8)circle(2pt)node[below,sollabel]{$B$};
 \node[sollabel,align=left,anchor=west] at (1.5,2.15) {$x_A=0,\quad E_{k,A}=0$};
 \node[sollabel,align=left,anchor=west] at (.4,.35) {$x_B=l_0/n$\\[6pt]$W_{\rm s}=-k(l_0/n)^2/2$};
 \draw[<->,DeepBlue] (5.7,-.8)--(5.7,1.5)node[midway,right,sollabel]{$l_0(1+1/n)$};
\end{tikzpicture}\diagcaption{解图}{高度差决定重力功，弹簧形变量的平方差决定弹力功。}\label{ch08:fig:soldiagFourThree}\end{workfig}}

\newcommand{\probdiagFourFour}{%
\begin{workfig}\begin{tikzpicture}[scale=.9]
 \draw[solbody] (-3.5,0)--(4,0);\draw[solbody] (-3.3,0)--(-3.3,1.1);
 \draw[spring] (-3.3,.4)--(-1.3,.4);
 \draw[solbody,fill=SoftGray] (-1.3,0)rectangle(-.5,.8);\node[sollabel]at(-.9,.4){$A$};
 \draw[solbody] (-.5,0)rectangle(.3,.8);\node[sollabel]at(-.1,.4){$B$};
 \draw[solgeom] (1,0)--(1,1.25);\node[sollabel,above]at(1,1.25){原长连接点 $x=0$};
 \draw[solgeom] (-1.3,0)--(-1.3,-.65);
 \draw[<->,DeepBlue] (-1.3,-.5)--(1,-.5)node[midway,below,sollabel]{$x_1$};
 \draw[solflow] (2,.5)--(3.4,.5)node[right,sollabel]{$+x$};
\end{tikzpicture}\diagcaption{题图}{以 $A$ 左侧弹簧连接点为位置标志，初态为 $x=-x_1$。}\label{ch08:fig:probdiagFourFour}\end{workfig}}

\newcommand{\soldiagFourFour}{%
\begin{workfig}\begin{tikzpicture}[scale=.85]
 \foreach \y/\a/\b in {3.2/-1.8/-1,1.6/0/.8,0/1.1/3.0}{
 \draw[solbody](-3,\y)--(4,\y);\draw[solbody](-3,\y)--++(0,.7);
 \draw[spring](-3,{\y+.35})--(\a,{\y+.35});
 \draw[solbody,fill=SoftGray](\a,\y)rectangle++(.8,.7);\node[sollabel]at({\a+.4},{\y+.35}){$A$};
 \draw[solbody](\b,\y)rectangle++(.8,.7);\node[sollabel]at({\b+.4},{\y+.35}){$B$};}
 \draw[solgeom](0,-.5)--(0,4.15);\node[sollabel,above]at(0,4.15){$x=0$};
 \node[sollabel,anchor=east]at(-3.4,3.55){初始压缩};
 \node[sollabel,anchor=east]at(-3.4,1.95){原长处分离};
 \node[sollabel,anchor=east]at(-3.4,.35){$A$ 的最右端};
 \draw[solflow](.1,2.55)--(.7,2.55)node[midway,above,sollabel]{$V$};
 \draw[solflow](1,2.55)--(1.6,2.55)node[midway,above,sollabel]{$V$};
 \draw[solflow](3.1,1)--(3.9,1)node[midway,above,sollabel]{$V$};
 \node[sollabel]at(1.5,1){$v_A=0$};
 \draw[<->,DeepBlue](-1.8,-.5)--(0,-.5)node[midway,below,sollabel]{$x_1$};
 \draw[<->,DeepBlue](0,-.5)--(1.1,-.5)node[midway,below,sollabel]{$x_2$};
\end{tikzpicture}\diagcaption{解图}{三行使用同一原长基准；分离后 $B$ 匀速，$A$ 继续伸长弹簧至转折。}\label{ch08:fig:soldiagFourFour}\end{workfig}}

\newcommand{\soldiagFourFive}{%
\begin{workfig}\begin{tikzpicture}[scale=.9]
 \draw[solbody,fill=SoftGray](-2,0)rectangle(3,.55);\node[sollabel]at(2.3,.27){$B,M$};
 \draw[solbody](-.8,.55)rectangle(.2,1.25);\node[sollabel]at(-.3,.95){$A,m$};
 \draw[solforce](-.8,.7)--(-2,.7)node[above,sollabel]{$f_A$};
 \draw[solforce](3,.27)--(4.2,.27)node[right,sollabel]{$f_B$};
 \draw[solflow](-1,1.8)--(2,1.8)node[midway,above,sollabel]{$s_A$};
 \draw[solflow](-1,-.5)--(.5,-.5)node[midway,below,sollabel]{$s_B$};
 \draw[solgeom](.5,-.5)--(.5,-1.25);\draw[solgeom](2,1.8)--(2,1.3);\draw[solgeom](2,0)--(2,-1.25);
 \draw[<->,DeepBlue](.5,-1.2)--(2,-1.2)node[midway,below,sollabel]{$s_A-s_B$};
\end{tikzpicture}\diagcaption{解图}{$A$ 与板接触，摩擦力等大反向；位移箭头比较同一时间内的地面位移。}\label{ch08:fig:soldiagFourFive}\end{workfig}}

\newcommand{\soldiagFourSix}{%
\begin{workfig}\begin{tikzpicture}[scale=.9]
 \foreach \y/\len in {1.8/2.5,0/4.3}{
 \draw[solbody](0,\y)--(5,\y);\draw[solbody,fill=SoftGray](.3,\y)rectangle(1,{\y+.6});
 \draw[solforce](1,{\y+.3})--(2,{\y+.3})node[right,sollabel]{$F$};
 \draw[solflow](.3,{\y-.4})--(\len,{\y-.4});}
 \node[sollabel,anchor=east]at(-.4,2.1){火车系：$0\to u$};
 \node[sollabel,anchor=east]at(-.4,.3){地面系：$v_0\to v_0+u$};
 \node[sollabel,below]at(1.4,1.4){$L$};\node[sollabel,below]at(2.3,-.4){$L+v_0t$};
\end{tikzpicture}\diagcaption{解图}{同一物体、同一时间区间，在两个惯性系中配对比较位移和动能。}\label{ch08:fig:soldiagFourSix}\end{workfig}}

\newcommand{\soldiagFourSeven}{%
\begin{workfig}\begin{tikzpicture}[scale=.9]
 \draw[solbody](-2.5,0)--(3,0);\draw[solbody,fill=SoftGray](-.5,0)rectangle(.5,.8);
 \draw[solforce](.5,.4)--(2,.4)node[right,sollabel]{$F$};
 \draw[solforce](-.5,.4)--(-2,.4)node[left,sollabel]{$ma$};
 \draw[solflow](-1,1.5)--(1,1.5)node[midway,above,sollabel]{火车加速度 $a$};
 \draw[solflow](0,-.5)--(2,-.5)node[midway,below,sollabel]{$L$（$F>ma$）};
\end{tikzpicture}\diagcaption{解图}{惯性力大小 $ma$、方向向左；图示向前位移只对应 $F>ma$ 的分支。}\label{ch08:fig:soldiagFourSeven}\end{workfig}}

\newcommand{\movingPulleyDiagram}{%
\begin{workfig}\begin{tikzpicture}[scale=.9]
 \draw[solbody](.4,2.3)--(1.3,2.3);\draw[solbody](0,0)circle(.8);
 % Rope: left free end, lower semicircle, right fixed end; draw after pulley.
 \draw[DeepBlue,line width=1pt](-.8,1.4)--(-.8,0)arc(180:360:.8)--(.8,2.3);
 \fill(.8,2.3)circle(2pt);\node[sollabel,above]at(.8,2.3){固定端};
 \draw[solbody](0,0)--(0,-1.2);\fill(0,0)circle(2pt);\node[sollabel,above]at(0,.05){轮轴};
 \draw[solbody,fill=SoftGray](-.35,-1.85)rectangle(.35,-1.2);\node[sollabel]at(0,-1.5){$m$};
 \draw[solforce](-.8,1.4)--(-.8,2.35)node[left,sollabel]{$F$};
 \draw[solforce](0,-1.85)--(0,-2.55)node[right,sollabel]{$mg$};
 \draw[solflow](-1.65,.9)--(-1.65,1.9)node[midway,left,sollabel]{$\Delta y_e=L$};
 \draw[solflow](1.55,-.25)--(1.55,.25)node[midway,right,sollabel]{$\Delta y_p=L/2$};
 \node[sollabel,align=left,anchor=west]at(2.3,1.5){蓝线：绕轮缘的绳\\中线：连接轮轴与重物的挂杆};
\end{tikzpicture}\diagcaption{原理图}{位移箭头表示同一过程的上升量，不是某段绳的长度。}\label{ch08:fig:movingPulleyDiagram}\end{workfig}}

\newcommand{\soldiagFourEight}{%
\begin{workfig}\begin{tikzpicture}[scale=.9]
 \draw[solflow](0,0)--(7,0)node[right,sollabel]{$r$};\fill(0,0)circle(3pt)node[below,sollabel]{固定点};
 \fill(1.5,0)circle(2pt)node[below,sollabel]{$r_A$};\fill(5.5,0)circle(2pt)node[below,sollabel]{$r_B$};
 \fill(3.5,0)circle(2pt)node[below,sollabel]{$r$};
 \draw[solforce](3.5,.8)--(2.2,.8)node[midway,above,sollabel]{$\vec F$（向内）};
 \draw[solflow](3.5,-.7)--(4.2,-.7)node[midway,below,sollabel]{$\dd r>0$};
\end{tikzpicture}\diagcaption{解图}{图示向外的一小段位移；引力功微元为 $-(k/r)\dd r$。}\label{ch08:fig:soldiagFourEight}\end{workfig}}

\newcommand{\soldiagFourNine}{%
\begin{workfig}\begin{tikzpicture}[scale=.9]
 \draw[fill=Blue!12,draw=none](1.4,0)--plot[domain=1.4:3.7,samples=60](\x,{1.6/(\x*\x)})--(3.7,0)--cycle;
 \draw[solflow](0,0)--(5,0)node[right,sollabel]{$r$};\draw[solflow](0,0)--(0,3.6)node[above,sollabel]{$F=k/r^2$};
 \draw[solpath,domain=.7:4.7,samples=80]plot(\x,{1.6/(\x*\x)});
 \foreach \r/\name in {1.4/A,3.7/B}{\draw[solgeom](\r,0)--(\r,{1.6/(\r*\r)});\node[sollabel,below]at(\r,0){$r_\name$};}
 \node[sollabel,align=left,anchor=west]at(2,2.2){阴影面积 $=k(1/r_A-1/r_B)$\\[5pt]图示 $r_B>r_A$，故 $W<0$};
\end{tikzpicture}\diagcaption{解图}{阴影是力的大小对距离的积分；功的正负还须结合运动方向。}\label{ch08:fig:soldiagFourNine}\end{workfig}}


力是改变物体速度的原因。在外力作用下，物体的速度会发生变化。牛顿第二定律集中描述
力与加速度的瞬时关系；在此基础上得到的动量定理，则说明力对时间的积累——冲量——
怎样改变物体的动量。除了动量，还可以用另一个物理量描述运动物体的性质，这就是动能。
本节要建立的核心联系是：力在空间位移上的积累形成“功”，总功等于物体动能的变化。

以下在惯性系中讨论质量不变的质点，$\vec F_{\text{合}}$ 表示它受到的合力。
如图~\ref{ch08:fig:workDerivationDiagram}，质点在 $t$ 时刻位于 $A$，速度为 $\vec v_A$；
经过 $\Delta t$ 到达 $B$，速度为 $\vec v_B$。先定义速度增量
\[
\Delta\vec v=\vec v_B-\vec v_A.
\]
速度有方向，因此应先作矢量相加，再求模的平方，不能直接把速率相加。
\workDerivationDiagram
\begin{equation}
 v_B^2-v_A^2=2\vec v_A\cdot\Delta\vec v+|\Delta\vec v|^2.
 \label{ch08:eq:velocity-square}
\end{equation}
式~\eqref{ch08:eq:velocity-square} 是精确恒等式。若加速度随时间连续，短时间内
$\Delta\vec v=\vec a_A\Delta t+o(\Delta t)$，其中 $o(\Delta t)/\Delta t\to0$。
设 $\vec a_A$ 与 $\vec v_A$ 夹角为 $\theta$，则保留一阶项时
\begin{equation}
 v_B^2-v_A^2=2a_Av_A\cos\theta\,\Delta t+o(\Delta t).
\end{equation}
这里不能把 $|\Delta\vec v|^2$ 当作恒等于零；它在除以 $\Delta t$ 后取极限时才消失。
同样，$\Delta\vec r=\vec v_A\Delta t+o(\Delta t)$，$v_A\Delta t$ 只是小段位移大小的一阶近似。
两边除以 $\Delta t$ 并令 $\Delta t\to0$，再用牛顿第二定律，得到
\begin{equation}
 \frac{\dd}{\dd t}\left(\frac12mv^2\right)
 =m\vec v\cdot\frac{\dd\vec v}{\dd t}
 =\vec F_{\text{合}}\cdot\vec v.
 \label{ch08:eq:kinetic-rate}
\end{equation}
定义动能 $E_k=\tfrac12mv^2$。利用 $\dd\vec r=\vec v\,\dd t$，上式可写成
\begin{equation}
 \dd E_k=\vec F_{\text{合}}\cdot\dd\vec r.
 \label{ch08:eq:differential-work-energy}
\end{equation}
任意力 $\vec F$ 与实际位移微元的夹角为 $\theta$ 时，其功定义为
\begin{equation}
 \dd W=\vec F\cdot\dd\vec r=F\cos\theta\,\dd s,
 \qquad \dd s=|\dd\vec r|.
 \label{ch08:eq:work-element}
\end{equation}
功和动能都是标量，没有空间方向；动能非负，功则可以为正、负或零。
力沿运动方向的分量作正功，反向分量作负功，垂直于瞬时位移的分量不作功。
在有限过程中，把实际路径分段、求和并取分段长度趋于零的极限：
\begin{equation}
 W_{A\to B}=\lim_{\max\Delta s_i\to0}\sum_i\vec F_i\cdot\Delta\vec r_i
 =\int_A^B\vec F\cdot\dd\vec r.
 \label{ch08:eq:work-integral}
\end{equation}
将式~\eqref{ch08:eq:differential-work-energy} 沿同一过程积分，中间状态的动能增减相互衔接，得到
\begin{equation}
 W_{\text{合}}(A\to B)=E_{k,B}-E_{k,A}
 =\frac12mv_B^2-\frac12mv_A^2.
 \label{ch08:eq:work-energy}
\end{equation}
合力功等于各分力功的代数和。总功为正，
物体动能增加；总功为负，物体动能减小。直观地说，沿运动方向推物体通常使其速率增大，
逆着运动方向施力通常使其速率减小；但严格判断仍应以力与位移的夹角以及全过程功的
代数和为准。

\subsection{功和动能的确定}

应用动能定理可以使许多问题大为简化，但必须先明确研究对象、过程的初末状态，以及各力
所做功的正负。对恒力而言，只有力在位移方向上的分量做功；垂直于位移的力不做功。

\subsubsection*{恒力功、功率与动能定理的使用范围}

恒力的大小和方向都不变，可从式~\eqref{ch08:eq:work-integral} 的积分号内提出：
\[
 W=\vec F\cdot\int_A^B\dd\vec r
  =\vec F\cdot(\vec r_B-\vec r_A)=F|\vec r_B-\vec r_A|\cos\theta.
\]
这里用的是始末位置之间的位移，不是弯曲路径的长度。变力一般不能这样提出积分号，必须沿实际路径计算。
点积还说明只有切向分量作功；法向力即使很大也可能不作功，但它仍可改变速度方向。

功率表示作功的快慢。时间间隔 $\Delta t$ 内的平均功率为 $\overline P=W/\Delta t$；令时间间隔趋于零，得到瞬时功率
\begin{equation}
 P=\frac{\dd W}{\dd t}
 =\vec F\cdot\frac{\dd\vec r}{\dd t}
 =\vec F\cdot\vec v=F_t v.
 \label{ch08:eq:power}
\end{equation}
$F_t$ 是力沿速度方向的有符号分量。功的单位是焦耳，$1\,\mathrm J=1\,\mathrm{N\,m}$；功率单位是瓦特，
$1\,\mathrm W=1\,\mathrm{J/s}$。只有牵引力平行于速度且 $v>0$ 时，才能由此写成 $F=P/v$。
因此匀功率与恒牵引力不是同一条件：匀功率时速度升高，牵引力减小。不能把这个理想关系直接外推到静止启动瞬间；实际机器的牵引力有上限。

也可以在直角坐标中逐项验证式~\eqref{ch08:eq:kinetic-rate}。把牛顿第二定律的三个分量分别乘以相应速度分量，再相加：
\begin{align*}
 F_xv_x+F_yv_y+F_zv_z
 &=m\left(v_x\frac{\dd v_x}{\dd t}+v_y\frac{\dd v_y}{\dd t}+v_z\frac{\dd v_z}{\dd t}\right)\\
 &=\frac{\dd}{\dd t}\left[\frac m2(v_x^2+v_y^2+v_z^2)\right]
 =\frac{\dd E_k}{\dd t}.
\end{align*}
再乘 $\dd t$ 并积分，左边利用 $v_x\dd t=\dd x$ 等关系化为合力功，右边成为末动能减初动能。
所以动能定理是牛顿第二定律沿实际路径的积分形式，适于联系速率和位置；若还要求时间或速度方向，通常要配合运动学或分量方程。

对质点系，外力功与内力功共同决定总动能变化。两个物体间作用力等大反向，但其作用点位移一般不同，
所以内力总功未必为零；弹簧伸缩、碰撞形变和摩擦发热都体现了这一点。只有把内力功转化为相应势能或内能，
才能写成更完整的能量关系。参考系改变时速度、动能和功都会改变，但在同一惯性系中二者始终满足动能定理。

虚功方法研究满足约束的假想无穷小位移。理想约束力对允许的虚位移的总功为零，平衡时主动外力的总虚功也为零。
它把合力、合力矩平衡改写到约束允许的坐标中，常能避开未知支持力和绳张力；后文将用动滑轮说明。


\begin{origexamplebox}{}\label{ch08:ex:s10-4-1}
以下各图中 $A,B$ 两点的高度差均为 $h$。物体分别沿图示各条路径由 $A$ 到 $B$，
求每一过程中重力所做的功。
\probdiagFourOne
\end{origexamplebox}

\Needspace{6\baselineskip}\begin{xesolution}
图~\ref{ch08:fig:probdiagFourOne} 的各条路径具有同一高度差；在图~\ref{ch08:fig:soldiagFourOne} 中把端点记为 $y_A,y_B$。
取竖直向上为 $y$ 轴正方向。重力为 $\vec G=-mg\,\vec e_y$，任意一小段位移中重力做功
\[
 \dd W_G=\vec G\cdot\dd\vec r=-mg\,\dd y.
\]
从 $A$ 到 $B$ 累积：
\[
 W_G=-mg\int_{y_A}^{y_B}\dd y
 =mg(y_A-y_B).
\]
\soldiagFourOne
由于所有图中 $y_A-y_B=h$，无论路径是直线、折线、曲线，途中是否先升高或绕行，
均有
\[
 \boxed{W_G=mgh}.
\]
重力功只由初、末位置的高度差决定，与实际路径的形状和长度无关。
\end{xesolution}

\Needspace{8\baselineskip}
\begin{origexamplebox}{}\label{ch08:ex:s10-4-2}
质量为 $m$ 的物体在与水平方向夹角为 $\theta$ 的恒力 $F$ 作用下，沿摩擦系数为
$\mu$ 的水平面向右运动距离 $L$。求在此过程中恒力 $F$、摩擦力 $f$、重力 $mg$
以及支持力 $N$ 所做的功。若初速度为 $v_0$、方向向右，求物体最后的速度。
\probdiagFourTwo
\end{origexamplebox}

\Needspace{6\baselineskip}\begin{xesolution}
图~\ref{ch08:fig:probdiagFourTwo} 给出 $F$ 与水平线的夹角，图~\ref{ch08:fig:soldiagFourTwo} 分别标出四个力。
物体在竖直方向没有位移。若全过程保持与水平面接触，则
\[
 N+F\sin\theta-mg=0,
 \qquad
 N=mg-F\sin\theta.
\]
滑动摩擦力大小为
\[
 f=\mu N=\mu(mg-F\sin\theta),
\]
方向与位移相反。各力所做的功分别为
\[
 \boxed{W_F=FL\cos\theta},
 \qquad
 \boxed{W_f=-\mu(mg-F\sin\theta)L},
\]
\[
 \boxed{W_G=0},
 \qquad
 \boxed{W_N=0}.
\]
\soldiagFourTwo
对物体应用动能定理：
\[
 FL\cos\theta-\mu(mg-F\sin\theta)L
 =\frac12mv^2-\frac12mv_0^2.
\]
所以
\[
 \boxed{
 v=\sqrt{v_0^2+\frac{2L}{m}
 \left[F\cos\theta-\mu(mg-F\sin\theta)\right]}}.
\]
这一表达式要求 $N\ge0$，且物体确实向右滑完整段距离 $L$。若根号内在到达 $L$ 前
已降为零，物体会先停下，摩擦力方向和后续运动必须重新判断。
\end{xesolution}

\Needspace{8\baselineskip}
\begin{origexamplebox}{}\label{ch08:ex:s10-4-3}
如图，质量为 $m$ 的小球系在劲度系数为 $k$、原长为 $l_0$ 的轻弹簧一端，弹簧另一端
固定在 $O$ 点。开始时小球位于水平位置 $A$，弹簧处于原长。小球由 $A$ 自由释放，
下落到 $O$ 点正下方的 $B$ 点时，弹簧伸长量为 $l_0/n$，其中 $n>0$。已知小球能够到达该状态，求其速率。
\probdiagFourThree
\end{origexamplebox}

\Needspace{6\baselineskip}\begin{xesolution}
图~\ref{ch08:fig:probdiagFourThree} 标出两个状态，虚线只连接始末位置，不用于指定运动轨迹。
在 $A$ 点弹簧长度为原长 $l_0$，在 $B$ 点弹簧长度为
\[
 OB=l_0+\frac{l_0}{n}=l_0\left(1+\frac1n\right).
\]
因此小球下降的高度就是 $OB$。重力做功
\[
 W_G=mg\,l_0\left(1+\frac1n\right).
\]
令弹簧伸长量为 $x=l-l_0$，径向向外为正。弹力为 $-kx\vec e_r$，
而实际位移在径向上的投影为 $\dd l=\dd x$，所以即使轨迹弯曲，也有
\[
 \dd W_{\rm s}=(-kx\vec e_r)\cdot\dd\vec r=-kx\,\dd x.
\]
从 $x=0$ 积分到 $x=l_0/n$，得到
\[
 W_{\rm s}=-\int_0^{l_0/n}kx\,\dd x=-\frac12k\left(\frac{l_0}{n}\right)^2.
\]
图~\ref{ch08:fig:soldiagFourThree} 将高度差和形变量分开标出，分别用于计算重力功和弹力功。
\soldiagFourThree
小球由静止释放，由动能定理
\[
 mg\,l_0\left(1+\frac1n\right)
 -\frac12k\frac{l_0^2}{n^2}
 =\frac12mv^2.
\]
故
\[
 \boxed{
 v=\sqrt{2gl_0\left(1+\frac1n\right)
 -\frac{k l_0^2}{mn^2}}}.
\]
弹力虽然在曲线路径上方向不断改变，但其总功只由弹簧初、末形变量决定。
根号内各项量纲均为速度平方。它必须非负；题给的 $n$ 是实际到达状态的参数，不能把任意 $n$ 都当作必可实现的运动状态。
\end{xesolution}

\subsection{多个物体的情况}

动能定理不但适用于单个质点，也适用于由多个质点构成的力学系统。系统的总动能等于
各质点动能的代数和。与动量定理不同，系统内力的总功一般不为零；求系统总动能的变化
时，外力功和内力功都不能遗漏。例如枪和子弹的总动量在发射过程中可以守恒，但火药内力
做正功，使系统总动能明显增加。

\Needspace{8\baselineskip}
\begin{origexamplebox}{}\label{ch08:ex:s10-4-4}
如图，把弹簧的一端固定在墙上，另一端系物体 $A$。将弹簧压缩 $x_1$ 后，在 $A$ 的
右边紧靠放置物体 $B$，然后撤去外力。物体 $A,B$ 均在光滑水平面上，质量分别为
$m_A,m_B$，弹簧劲度系数为 $k$。求物体 $A$ 移动的最大距离。
\probdiagFourFour
\end{origexamplebox}

\Needspace{6\baselineskip}\begin{xesolution}
按图~\ref{ch08:fig:probdiagFourFour}，以弹簧连接点在原长时的位置为 $x=0$，伸长为正。
弹簧由压缩状态恢复时，$A$ 推着 $B$ 一起向右运动。在弹簧仍被压缩的阶段，两物体速度
相同，可看作总质量为 $m_A+m_B$ 的整体。此时 $x<0$，共同加速度及 $A$ 对 $B$ 的接触力为
\[
 a=\frac{-kx}{m_A+m_B},\qquad N=m_Ba=\frac{-m_Bkx}{m_A+m_B}>0.
\]
因此在恢复原长之前不会分离；到 $x=0$ 时 $N$ 才降为零。由动能定理
\[
 \frac12kx_1^2
 =\frac12(m_A+m_B)V^2,
\]
故
\[
 V=x_1\sqrt{\frac{k}{m_A+m_B}}.
\]
越过原长位置后，弹簧开始伸长并向左拉 $A$，而 $A$ 不能“拉住”右侧的 $B$，所以两物体
在弹簧原长处恰好分离。此后 $B$ 以速度 $V$ 匀速前进，$A$ 单独继续向右并拉伸弹簧。
\soldiagFourFour
图~\ref{ch08:fig:soldiagFourFour} 的三行分别表示初态、分离和 $A$ 到达最右端；三行使用相同的 $x=0$。
设 $A$ 在原长位置右侧继续移动的最大距离为 $x_2$，则
\[
 \frac12m_AV^2=\frac12kx_2^2,
\]
从而
\[
 x_2=x_1\sqrt{\frac{m_A}{m_A+m_B}}.
\]
物体 $A$ 从初始压缩位置到最右端的总路程为 $x_1+x_2$，所以
\[
 \boxed{
 s_{A,\max}=x_1\left(1+\sqrt{\frac{m_A}{m_A+m_B}}\right)}.
\]
若只问相对弹簧原长位置的最大伸长，则答案是 $x_2$；本题问的是 $A$ 从初始位置起移动的
最大距离，因此还要加上先前走过的 $x_1$。
\end{xesolution}

\Needspace{8\baselineskip}
\begin{origexamplebox}{}\label{ch08:ex:s10-4-5}
光滑水平面上有一块足够长、质量为 $M$、初始静止的板 $B$。某时刻质量为 $m$、
速度为 $v_0$ 的物体 $A$ 冲上板的上表面，两者间的动摩擦因数为 $\mu$。求从开始相对
滑动到相对速度变为零的过程中，摩擦力对两个物体所做的功各为多少；并求整个系统的
动能变化量。

\ThQSlidingBoard
\end{origexamplebox}

\Needspace{6\baselineskip}\begin{xesolution}
水平方向没有外力，达到共同速度 $V$ 时由动量守恒
\[
 mv_0=(m+M)V,
 \qquad
 V=\frac{m}{m+M}v_0.
\]
滑动期间，$A$ 的加速度为 $-\mu g$，板 $B$ 的加速度为 $\mu mg/M$。相对速度由
$v_0$ 减到零。相对加速度为两加速度之差，故
\[
 a_{A/B}=-\mu g-\frac{\mu mg}{M},\qquad
 0=v_0-\mu g\left(1+\frac mM\right)t.
\]
所需时间为
\[
 t=\frac{Mv_0}{\mu g(m+M)}.
\]
两物体相对地面的位移分别为
\[
 s_A=\frac{v_0+V}{2}\,t,
 \qquad
 s_B=\frac{V}{2}\,t.
\]
\soldiagFourFive
图~\ref{ch08:fig:soldiagFourFive} 中 $s_A>s_B$，两位移都相对地面计算。
摩擦力对 $A$ 方向向左，对 $B$ 方向向右，因此
\[
 \boxed{
 W_{f,A}=-\mu mgs_A
 =-\frac{mM(M+2m)}{2(m+M)^2}v_0^2},
\]
\[
 \boxed{
 W_{f,B}=+\mu mgs_B
 =\frac{m^2M}{2(m+M)^2}v_0^2}.
\]
系统动能变化为两摩擦功之和：
\[
 \boxed{
 \Delta E_k=W_{f,A}+W_{f,B}
 =-\frac{mM}{2(m+M)}v_0^2}.
\]
两个摩擦力大小相等、方向相反，但作用点相对地面的位移不同，所以它们的功不互相抵消。
损失的机械能等于摩擦力与相对滑动距离的乘积。也可分别核验
$W_{f,A}=\tfrac12m(V^2-v_0^2)$、$W_{f,B}=\tfrac12MV^2$，代入共同速度即得到同样结果。
这里取 $\mu>0$；减小 $\mu$ 会延长滑动时间，最终损失却不变，而 $\mu=0$ 时根本不会达到共同速度，不能套用这一终态。
\end{xesolution}

\section{功与动能的参考系依赖}

功依赖于力的作用点在所选参考系中的位移，动能依赖于物体在所选参考系中的速度，因此
二者都与参考系有关。只有在惯性参考系中，动能定理才能直接写成通常形式；同一过程可以
在不同惯性系中使用动能定理，但必须在同一个参考系内同时计算功和动能。

\begin{origexamplebox}{}\label{ch08:ex:s10-4-6}
一列火车以速度 $v_0$ 做匀速直线运动。火车上的人在光滑水平面上用恒力 $F$ 向前拉动质量为 $m$ 的物体，
使它相对火车向前移动距离 $L$。物体起初相对火车静止。先在火车参考系中求物体最后的
速度，再在相对地面静止的参考系中描述同一过程，并列出动能定理方程。

\ThQTrain
\end{origexamplebox}

\Needspace{6\baselineskip}\begin{xesolution}
火车做匀速直线运动，火车系与地面系都是惯性系。设物体相对火车的末速度为 $u$。在火车
系中，力的作用点位移为 $L$，动能定理给出
\[
 FL=\frac12mu^2,
 \qquad
 \boxed{u=\sqrt{\frac{2FL}{m}}}.
\]
设拉动持续时间为 $t$。由于 $u=Ft/m$，所以 $m u=Ft$。在地面系中，物体的初、末速度
分别为 $v_0$ 和 $v_0+u$；同一时间内，力的作用点相对地面的位移为 $L+v_0t$。
\soldiagFourSix
图~\ref{ch08:fig:soldiagFourSix} 的两行对应同一过程：地面位移等于车内位移加上火车的位移。
地面系中的动能定理为
\[
 \boxed{
 F(L+v_0t)
 =\frac12m(v_0+u)^2-\frac12mv_0^2}.
\]
展开右端并利用 $FL=\tfrac12mu^2$、$Ft=mu$，两边完全相等。不同惯性系中功和动能
变化的数值可以不同，但各自在本参考系内仍满足动能定理。
\end{xesolution}

\Needspace{8\baselineskip}
\begin{origexamplebox}{}\label{ch08:ex:s10-4-7}
同上，但火车以加速度 $a$ 做匀加速运动。在火车参考系和地面参考系中分别直接使用
动能定理，能否得到正确结果？若要使动能定理在火车参考系中也成立，应怎样处理？

\ThQAcceleratingTrain
\end{origexamplebox}

\Needspace{6\baselineskip}\begin{xesolution}
地面系近似为惯性系，因而可以直接使用动能定理。设过程开始时火车和物体的地面速度均为
$v_0$，持续时间为 $t$。物体只受水平拉力 $F$，故地面加速度为 $F/m$：
\[
 v_{\rm f}=v_0+\frac{F}{m}t,
 \qquad
 x_{\rm obj}=v_0t+\frac12\frac{F}{m}t^2.
\]
地面系中有
\[
 \boxed{
 Fx_{\rm obj}
 =\frac12m\left(v_0+\frac{F}{m}t\right)^2-\frac12mv_0^2}.
\]
火车系以加速度 $a$ 运动，不是惯性系。若仍只把 $F$ 的功写成 $FL$，将不能得到正确的
相对运动。要在火车系中使用动能定理，必须给物体加入方向与火车加速度相反的惯性力
\[
 \vec F_{\rm in}=-m\vec a.
\]
\soldiagFourSeven
图~\ref{ch08:fig:soldiagFourSeven} 中真实拉力向右、惯性力向左。
若拉力与火车加速度同向，且 $F>ma$，物体相对火车由静止向前移动距离 $L$，则
\[
 (F-ma)L=\frac12mu_{\rm rel}^2,
\]
所以
\[
 \boxed{u_{\rm rel}
 =\sqrt{2\left(\frac{F}{m}-a\right)L}}.
\]
若 $F=ma$，相对加速度为零，物体始终相对火车静止，不会完成 $L>0$ 的车内位移。
若 $F<ma$，物体相对火车向后运动；用有符号位移 $\xi<0$ 写
$(F-ma)\xi=\tfrac12mu_{\rm rel}^2$，两边仍为非负，不能把“向前移动 $L$”的根式硬套到此分支。
惯性力不是新的相互作用力，而是为了在非惯性系中继续使用牛顿定律和动能定理而引入的
等效力。
\end{xesolution}

\subsection{功与力的关系：虚功原理初步}

杠杆、轮轴、齿轮、定滑轮和动滑轮
等简单机械可以组合成复杂装置，把动力转化为所需要的形式。机械的一端连接动力源，动力
可以来自人力、蒸汽机或电动机；另一端是输出端，用来提升重物或产生其他运动。

机械不能创造能量，但可以储存能量。无摩擦、无其他损耗时，输入功与输出功的正确关系是
\[
 W_{\rm in}=W_{\rm out}+\Delta E_{\rm machine}.
\]
只有机械本身的动能、势能等不变时，才有输入功等于输出功；无摩擦本身并不能保证这一点。
在这个条件下，机械用较小的力克服较大的阻力，代价是动力作用点移动更长的距离，“省力”不意味着“省功”。

图~\ref{ch08:fig:movingPulleyDiagram} 为轻绳、轻滑轮构成的理想动滑轮，绳不可伸长，下方重物质量为 $m$。
取竖直向上为正，固定端高度为 $y_f$，自由端高度为 $y_e$，轮心高度为 $y_p$。
绳长为 $(y_f-y_p)+(y_e-y_p)+\pi R=\text{常量}$，所以
$\Delta y_e=2\Delta y_p$。自由端向上移动 $L$ 时，重物向上移动 $L/2$。
重物匀速上升，动能不变，因此拉力的功等于重力势能的增加：
\movingPulleyDiagram
\begin{equation}
 F L=mg\frac{L}{2}.
\end{equation}
消去共同的位移尺度 $L$，得到
\begin{equation}
 F=\frac12mg.
\end{equation}
如果装置处于平衡而并未实际运动，也可以假想一次符合绳长约束的微小位移，记作 $\delta y$。
它不是已经发生的位移，也不要求给它分配一段时间。理想约束下，平衡要求主动外力的总虚功为零：
\[
 F\delta y_e-mg\delta y_p=0,\qquad \delta y_e=2\delta y_p.
\]
由此同样得到 $F=mg/2$。这就是虚功原理在本例中的用法。
它把允许的位移关系与平衡条件结合起来，可避免逐个求出约束复杂系统中的支持力和内部绳张力。

\subsection{变力作功}

作用力不是恒力时，仍可使用动能定理，关键是正确计算变力的功。建立直角坐标系，把力写成
$(F_x,F_y,F_z)$，位移微元写成 $(\dd x,\dd y,\dd z)$，则功的微元为
\begin{equation}
 \dd W=F_x\dd x+F_y\dd y+F_z\dd z=\vec F\cdot\dd\vec r.
\end{equation}
在有限过程中，应把各微小功累积起来，即
\[
 W=\int_A^B\vec F\cdot\dd\vec r.
\]


\paragraph{积分沿真实路径，不能只代入首末位置}
上式的积分还包含路径信息。若已知运动参数 $\vect r(t)$，一般情形应写为
\[
 W=\int_{t_A}^{t_B}\vect F(\vect r(t),t)\cdot\dot{\vect r}(t)\dd t.
\]
图~\ref{quality:08:work}中 $\dd\vect r$ 沿轨迹切线，故只有 $F\cos\theta$ 进入功；法向分力即使改变运动方向，也不在这一微小位移上做功。若使用 $\int F_x(x)\dd x$，必须确认其他分量没有贡献，且已把沿该段路径的 $F_x$ 正确表示为 $x$ 的函数。
保守力的功可进一步化为势能差，但这需要力场具有相应的路径无关性，不能仅因题中给出了起点与终点就省去路径。例如滑动摩擦功是 $-\int\mu N\dd s$，积分用路程元 $\dd s\ge0$；往返运动的净位移可以为零，摩擦耗散却不为零。
\begin{qualityfigure}\begin{tikzpicture}[>=Latex]
 \draw[thick] (0,0) node[left] {$A$} .. controls (1,1.7) and (1.8,.55) .. (2.5,.55)
 .. controls (3.2,.55) and (4,.4) .. (5,1.4) node[right] {$B$};
 \fill (2.5,.55) circle (1.5pt);
 \draw[vector] (2.5,.55)--(4.1,.55) node[right] {$\dd\vect r$};
 \draw[vector,draw=TextGray] (2.5,.55)--(3.6,1.75) node[above] {$\vect F$};
 \node at (3.1,.95) {$\theta$};
 \node[below] at (2.5,-.5) {$\delta W=F\cos\theta\,\dd s$};
\end{tikzpicture}
\figcap{路径积分在每个位置取力与局部切向位移的点乘，箭头长度只作局部方向示意。}\label{quality:08:work}
\end{qualityfigure}


\Needspace{14\baselineskip}
\begin{origexamplebox}{}\label{ch08:ex:s10-4-8}
两个物体间的相互吸引力大小与它们之间的距离成反比。固定其中一个物体，求另一个物体
由相距 $r_A$ 运动到相距 $r_B$ 的过程中，该吸引力所做的功。

\ThQInverseForce
\end{origexamplebox}

\Needspace{6\baselineskip}\begin{xesolution}
设吸引力大小为
\[
 F=\frac{k}{r},
 \qquad k>0.
\]
取径向向外为正方向，吸引力的径向分量为 $F_r=-k/r$。因此
\[
 W_{A\to B}
 =\int_{r_A}^{r_B}F_r\,\dd r
 =-\int_{r_A}^{r_B}\frac{k}{r}\,\dd r.
\]
\soldiagFourEight
图~\ref{ch08:fig:soldiagFourEight} 区分有限的始末距离与当前位置上的径向位移微元。积分得到
\[
 \boxed{W_{A\to B}=k\ln\frac{r_A}{r_B}}.
\]
若题目所说“大小”专指绝对值，则为
\[
 \boxed{|W|=k\left|\ln\frac{r_A}{r_B}\right|}.
\]
当物体向内运动，即 $r_B<r_A$ 时，引力与位移同向，功为正；向外运动时功为负。
\end{xesolution}

\begin{origexamplebox}{}\label{ch08:ex:s10-4-9}
两个物体间的相互吸引力大小与它们之间距离的平方成反比。固定其中一个物体，求另一个
物体由相距 $r_A$ 运动到相距 $r_B$ 的过程中，该吸引力所做的功。

\ThQInverseSquareForce
\end{origexamplebox}

\Needspace{6\baselineskip}\begin{xesolution}
设吸引力大小为
\[
 F=\frac{k}{r^2},
 \qquad k>0.
\]
仍取径向向外为正，故 $F_r=-k/r^2$：
\[
 W_{A\to B}
 =-\int_{r_A}^{r_B}\frac{k}{r^2}\,\dd r
 =\left.\frac{k}{r}\right|_{r_A}^{r_B}.
\]
\soldiagFourNine
图~\ref{ch08:fig:soldiagFourNine} 画出 $r_B>r_A$ 时力的大小曲线：阴影面积为正，而向外运动的引力功为负。
所以功及其绝对值分别为
\[
\boxed{
 W_{A\to B}=k\left(\frac1{r_B}-\frac1{r_A}\right)},
\qquad
\boxed{
 |W|=k\left|\frac1{r_B}-\frac1{r_A}\right|}.
\]
符号判断与例~\ref{ch08:ex:s10-4-8} 相同：向内运动时引力做正功，向外运动时引力做负功。
\end{xesolution}
}

{
\graphicspath{{../../../sample/xes-original/eleventh-section/}}
\section{势能与机械能守恒}
\tikzset{solbody/.style={draw=Ink,line width=.8pt},solpath/.style={draw=Ink,line width=.8pt},
 solforce/.style={-{Latex[length=2mm]},draw=DeepBlue,line width=.9pt},
 solflow/.style={-{Latex[length=1.8mm]},draw=Blue,line width=.8pt},
 solgeom/.style={draw=TextGray,densely dashed,line width=.5pt},
 sollabel/.style={font=\scriptsize,fill=none,inner sep=1.5pt},
 solpanel/.style={font=\scriptsize\bfseries,text=TextGray},
 spring/.style={decorate,decoration={coil,segment length=4pt,amplitude=2.5pt},draw=Ink}}
\newcommand{\diagcaption}[2]{\refstepcounter{figure}\par\nopagebreak[4]\smallskip
 {\centering\scriptsize\color{TextGray}图~\thefigure\quad #2\par}\smallskip}
\newenvironment{energyfig}{\begin{center}\begin{minipage}{.98\linewidth}\centering}{\end{minipage}\end{center}}
% Generated by tools/generate_ch08_geometry.py
\newcommand{\screenedPotentialCurve}{\draw[solpath] (0.40000000,-1.26860129)--(0.42000000,-1.17739386)--(0.44000000,-1.09554935)--(0.46000000,-1.02178726)--(0.48000000,-0.95504559)--(0.50000000,-0.89443654)--(0.52000000,-0.83921253)--(0.54000000,-0.78873977)--(0.56000000,-0.74247769)--(0.58000000,-0.69996252)--(0.60000000,-0.66079430)--(0.62000000,-0.62462641)--(0.64000000,-0.59115703)--(0.66000000,-0.56012225)--(0.68000000,-0.53129036)--(0.70000000,-0.50445717)--(0.72000000,-0.47944208)--(0.74000000,-0.45608482)--(0.76000000,-0.43424272)--(0.78000000,-0.41378841)--(0.80000000,-0.39460782)--(0.82000000,-0.37659853)--(0.84000000,-0.35966839)--(0.86000000,-0.34373421)--(0.88000000,-0.32872077)--(0.90000000,-0.31455991)--(0.92000000,-0.30118970)--(0.94000000,-0.28855379)--(0.96000000,-0.27660079)--(0.98000000,-0.26528376)--(1.00000000,-0.25455974)--(1.02000000,-0.24438936)--(1.04000000,-0.23473647)--(1.06000000,-0.22556781)--(1.08000000,-0.21685277)--(1.10000000,-0.20856312)--(1.12000000,-0.20067278)--(1.14000000,-0.19315763)--(1.16000000,-0.18599535)--(1.18000000,-0.17916523)--(1.20000000,-0.17264806)--(1.22000000,-0.16642599)--(1.24000000,-0.16048241)--(1.26000000,-0.15480184)--(1.28000000,-0.14936987)--(1.30000000,-0.14417304)--(1.32000000,-0.13919879)--(1.34000000,-0.13443534)--(1.36000000,-0.12987172)--(1.38000000,-0.12549761)--(1.40000000,-0.12130336)--(1.42000000,-0.11727991)--(1.44000000,-0.11341874)--(1.46000000,-0.10971186)--(1.48000000,-0.10615176)--(1.50000000,-0.10273134)--(1.52000000,-0.09944396)--(1.54000000,-0.09628332)--(1.56000000,-0.09324349)--(1.58000000,-0.09031890)--(1.60000000,-0.08750425)--(1.62000000,-0.08479455)--(1.64000000,-0.08218509)--(1.66000000,-0.07967140)--(1.68000000,-0.07724926)--(1.70000000,-0.07491465)--(1.72000000,-0.07266379)--(1.74000000,-0.07049308)--(1.76000000,-0.06839910)--(1.78000000,-0.06637862)--(1.80000000,-0.06442855)--(1.82000000,-0.06254598)--(1.84000000,-0.06072812)--(1.86000000,-0.05897232)--(1.88000000,-0.05727608)--(1.90000000,-0.05563700)--(1.92000000,-0.05405280)--(1.94000000,-0.05252129)--(1.96000000,-0.05104042)--(1.98000000,-0.04960821)--(2.00000000,-0.04822277)--(2.02000000,-0.04688231)--(2.04000000,-0.04558510)--(2.06000000,-0.04432951)--(2.08000000,-0.04311398)--(2.10000000,-0.04193699)--(2.12000000,-0.04079713)--(2.14000000,-0.03969302)--(2.16000000,-0.03862336)--(2.18000000,-0.03758688)--(2.20000000,-0.03658240)--(2.22000000,-0.03560875)--(2.24000000,-0.03466484)--(2.26000000,-0.03374961)--(2.28000000,-0.03286205)--(2.30000000,-0.03200120)--(2.32000000,-0.03116611)--(2.34000000,-0.03035590)--(2.36000000,-0.02956970)--(2.38000000,-0.02880671)--(2.40000000,-0.02806612)--(2.42000000,-0.02734718)--(2.44000000,-0.02664916)--(2.46000000,-0.02597136)--(2.48000000,-0.02531311)--(2.50000000,-0.02467375)--(2.52000000,-0.02405267)--(2.54000000,-0.02344927)--(2.56000000,-0.02286296)--(2.58000000,-0.02229320)--(2.60000000,-0.02173945)--(2.62000000,-0.02120120)--(2.64000000,-0.02067795)--(2.66000000,-0.02016923)--(2.68000000,-0.01967457)--(2.70000000,-0.01919353)--(2.72000000,-0.01872569)--(2.74000000,-0.01827064)--(2.76000000,-0.01782797)--(2.78000000,-0.01739730)--(2.80000000,-0.01697828)--(2.82000000,-0.01657053)--(2.84000000,-0.01617371)--(2.86000000,-0.01578750)--(2.88000000,-0.01541157)--(2.90000000,-0.01504562)--(2.92000000,-0.01468934)--(2.94000000,-0.01434245)--(2.96000000,-0.01400466)--(2.98000000,-0.01367571)--(3.00000000,-0.01335534)--(3.02000000,-0.01304330)--(3.04000000,-0.01273934)--(3.06000000,-0.01244323)--(3.08000000,-0.01215474)--(3.10000000,-0.01187365)--(3.12000000,-0.01159975)--(3.14000000,-0.01133282)--(3.16000000,-0.01107269)--(3.18000000,-0.01081914)--(3.20000000,-0.01057199)--(3.22000000,-0.01033107)--(3.24000000,-0.01009619)--(3.26000000,-0.00986720)--(3.28000000,-0.00964392)--(3.30000000,-0.00942619)--(3.32000000,-0.00921386)--(3.34000000,-0.00900679)--(3.36000000,-0.00880483)--(3.38000000,-0.00860783)--(3.40000000,-0.00841566)--(3.42000000,-0.00822820)--(3.44000000,-0.00804531)--(3.46000000,-0.00786686)--(3.48000000,-0.00769275)--(3.50000000,-0.00752285)--(3.52000000,-0.00735705)--(3.54000000,-0.00719524)--(3.56000000,-0.00703731)--(3.58000000,-0.00688317)--(3.60000000,-0.00673271)--(3.62000000,-0.00658583)--(3.64000000,-0.00644244)--(3.66000000,-0.00630246)--(3.68000000,-0.00616578)--(3.70000000,-0.00603232)--(3.72000000,-0.00590201)--(3.74000000,-0.00577475)--(3.76000000,-0.00565048)--(3.78000000,-0.00552911)--(3.80000000,-0.00541056)--(3.82000000,-0.00529477)--(3.84000000,-0.00518167)--(3.86000000,-0.00507119)--(3.88000000,-0.00496326)--(3.90000000,-0.00485781)--(3.92000000,-0.00475479)--(3.94000000,-0.00465413)--(3.96000000,-0.00455578)--(3.98000000,-0.00445967)--(4.00000000,-0.00436575)--(4.02000000,-0.00427397)--(4.04000000,-0.00418427)--(4.06000000,-0.00409660)--(4.08000000,-0.00401091)--(4.10000000,-0.00392715)--(4.12000000,-0.00384528)--(4.14000000,-0.00376524)--(4.16000000,-0.00368700)--(4.18000000,-0.00361051)--(4.20000000,-0.00353573);}


\newcommand{\gravityPathsDiagram}{%
\begin{energyfig}\begin{tikzpicture}[scale=.9]
 \coordinate(A)at(2.7,1.5);\coordinate(B)at(-2.7,-1);
 \draw[solpath](A)--(B);
 \draw[solpath](A)..controls(1.2,2.3)and(-1.8,2.3)..(B);
 \draw[solpath](A)..controls(1,.9)and(-.8,1.2)..(B);
 \draw[solpath](A)..controls(2.4,-.3)and(.2,-1.6)..(B);
 \fill(A)circle(2pt)node[above right,sollabel]{$A,y_A$};\fill(B)circle(2pt)node[below left,sollabel]{$B,y_B$};
 \draw[solgeom](-3.2,1.5)--(3,1.5);\draw[solgeom](-3.2,-1)--(3,-1);
 \draw[<->,DeepBlue](-3.2,-1)--(-3.2,1.5)node[midway,left,sollabel]{$h$};
\end{tikzpicture}\diagcaption{示意图}{只有在物体可达且保持固定轨道约束时，才能比较这些路径的终态速率。}\label{ch08:fig:gravityPathsDiagram}\end{energyfig}}

\newcommand{\probdiagFourTen}{%
\begin{energyfig}\begin{tikzpicture}[scale=.85]
 \begin{scope}[xshift=-3.5cm,yshift=3.6cm]
 \node[solpanel]at(0,1.7){(1) 不脱离的光滑槽};
 \draw[solpath](-1.6,-.6)--(-.3,-.6)..controls(.9,-.6)and(.1,.8)..(.9,1.3);
 \draw[solpath](-1.6,-.72)--(-.3,-.72)..controls(1.03,-.72)and(.23,.72)..(1.03,1.3);
 \fill(-1.2,-.66)circle(2pt);\draw[solflow](-1.2,-.3)--(-.4,-.3)node[midway,above,sollabel]{$v$};
 \end{scope}
 \begin{scope}[xshift=3.5cm,yshift=3.6cm]
 \node[solpanel]at(0,1.7){(2) 斜抛};\draw[solgeom](-1.5,-.6)--(1.5,-.6);
 \fill(-1,-.6)circle(2pt);\draw[solforce](-1,-.6)--++(40:1.8)node[right,sollabel]{$v$};
 \draw[TextGray](-.4,-.6)arc(0:40:.6);\node[sollabel]at(-.15,-.35){$\theta$};
 \end{scope}
 \foreach \pos/\title in {-3.5/{(3) 珠子套在细圆环上},3.5/{(4) 质点沿圆柱内壁}}{
 \begin{scope}[xshift=\pos cm]
 \node[solpanel]at(0,1.75){\title};\draw[solbody](0,0)circle(1.1);
 \draw[solgeom](0,0)--(135:1.1)node[midway,left,sollabel]{$R$};
 \draw[solflow](.25,-1.4)--(1.15,-1.4)node[right,sollabel]{$v$};
 \end{scope}}
 \draw[DeepBlue,line width=1pt](-3.5,-1.1)circle(.12);
 \node[sollabel,below]at(-3.9,-1.3){套孔};\draw[TextGray](-3.9,-1.3)--(-3.62,-1.15);
 \fill(3.5,-1.1)circle(2.2pt);
\end{tikzpicture}\diagcaption{题图}{第三种约束可双向作用；第四种只能由内壁向圆心推压。}\label{ch08:fig:probdiagFourTen}\end{energyfig}}

\newcommand{\soldiagFourTen}{%
\begin{energyfig}\begin{tikzpicture}[scale=.9]
 \begin{scope}[xshift=-3.6cm]
 \node[solpanel]at(0,2.9){套环：不发生脱离};\draw[solbody](0,0)circle(1.35);
 \draw[DeepBlue,line width=1pt](45:1.35)circle(.13);
 \draw[solforce](45:1.35)--(45:.4)node[below right,sollabel]{$N$};
 \node[sollabel,align=center]at(0,-2){此处只画向内的一种状态\\其他状态允许反向};
 \end{scope}
 \begin{scope}[xshift=3.2cm]
 \def\rr{1.6}\def\cc{-.8}\pgfmathsetmacro\ss{sqrt(1-\cc*\cc)}\pgfmathsetmacro\uu{sqrt(-\cc*\rr)}
 \pgfmathsetmacro\ta{\uu*\ss}
 \coordinate(O)at(0,0);\coordinate(D)at({\rr*\ss},{-\rr*\cc});
 \node[solpanel]at(0,2.9){内壁：$D$ 点 $N=0$ 后脱离};
 \draw[solbody](O)circle(\rr);\fill(O)circle(1.5pt)node[below left,sollabel]{$O$};
 \draw[solgeom](0,-\rr)--(O)--(D);\fill(D)circle(2pt);\node[sollabel,anchor=west]at($(D)+(.25,.35)$){$D$};
 \draw[TextGray](0,-.43)arc(-90:53.1301:.43);\node[sollabel]at(.65,-.3){$\varphi_{\rm d}$};
 \draw[Blue,line width=.9pt,domain=0:{1.2*\ta},samples=80,variable=\t]
 plot({\rr*\ss+\uu*\cc*\t},{-\rr*\cc+\uu*\ss*\t-.5*\t*\t});
 \draw[solforce](D)--++({1.2*\cc},{1.2*\ss})node[above,sollabel]{$\vec u_{\rm d}$};
 \draw[solgeom](2.25,{-\rr*\cc})--(D);
 \draw[solforce](2.25,{-\rr*\cc})--++(0,{1.2*\ss})node[right,sollabel]{$u_{{\rm d},y}$};
 \draw[<->,TextGray](-2,-\rr)--(-2,{-\rr*\cc})node[midway,left,sollabel]{$R(1-c)$};
 \node[sollabel]at(0,-2){蓝线由抛体方程确定};
 \end{scope}
\end{tikzpicture}\diagcaption{解图}{右图取 $c=-0.8$ 展示脱离到首次最高点附近的运动，抛出切线与半径垂直。}\label{ch08:fig:soldiagFourTen}\end{energyfig}}

\newcommand{\probdiagFourEleven}{%
\begin{energyfig}\begin{tikzpicture}[scale=.9]
 \def\rr{2}\coordinate(O)at(0,0);\coordinate(D)at(45:\rr);
 \draw[solbody](-2.6,0)--(3,0);\draw[solbody](-\rr,0)arc(180:0:\rr);
 \fill(0,\rr)circle(2.5pt)node[above,sollabel]{初态质点 $m$};
 \fill(O)circle(1.5pt)node[below,sollabel]{$O$};\draw[solgeom](0,\rr)--(O)--(D);
 \fill(D)circle(1.5pt);\node[sollabel]at(1.1,.65){$R$};
 \draw[TextGray](45:.65)arc(45:90:.65);\node[sollabel]at(67.5:.92){$\theta$};
\end{tikzpicture}\diagcaption{题图}{初态位于圆面顶点；角度 $\theta$ 从竖直向上的半径量起。}\label{ch08:fig:probdiagFourEleven}\end{energyfig}}

\newcommand{\soldiagFourEleven}{%
\begin{energyfig}\begin{tikzpicture}[scale=1]
 \def\rr{2}\pgfmathsetmacro\ss{sqrt(5)/3}\def\cc{.6666666667}
 \pgfmathsetmacro\uu{sqrt(2*\rr/3)}\pgfmathsetmacro\ux{\uu*\cc}\pgfmathsetmacro\uy{-\uu*\ss}
 \pgfmathsetmacro\tf{\uy+sqrt(\uy*\uy+2*\rr*\cc)}\pgfmathsetmacro\xf{\rr*\ss+\ux*\tf}
 \coordinate(D)at({\rr*\ss},{\rr*\cc});
 \draw[solbody](-2.5,0)--(3.4,0);\draw[solbody](-\rr,0)arc(180:0:\rr);
 \draw[solgeom](0,\rr)--(0,0)--(D);\fill(0,0)circle(1.5pt)node[below left,sollabel]{$O$};
 \fill(D)circle(2pt)node[above right,sollabel]{$D$};
 \draw[Blue,line width=1pt,domain=0:\tf,samples=80,variable=\t]
 plot({\rr*\ss+\ux*\t},{\rr*\cc+\uy*\t-.5*\t*\t});
 \draw[solforce](D)--++({.95*\cc},{-.95*\ss})node[right,sollabel]{$\vec u$};
 \draw[solforce](D)--++(0,-.75)node[below left,sollabel]{$mg$};
 \draw[TextGray](41.8103:.5)arc(41.8103:90:.5);\node[sollabel]at(65:.8){$\theta$};
 \fill(\xf,0)circle(2pt)node[above right,sollabel]{落点};
 \draw[solgeom](0,0)--(0,-.55);\draw[solgeom](\xf,0)--(\xf,-.55);
 \draw[<->,DeepBlue](0,-.45)--(\xf,-.45)node[midway,below,sollabel]{$x$};
 \node[sollabel,align=left]at(-1.3,2.6){$D$ 点：$N=0$\\$\cos\theta=2/3$};
\end{tikzpicture}\diagcaption{解图}{脱离速度沿圆面切线，落点在半圆右端之外；脱离后仅受重力。}\label{ch08:fig:soldiagFourEleven}\end{energyfig}}

\newcommand{\probdiagFourTwelve}{%
\begin{energyfig}\begin{tikzpicture}[scale=.9]
 \def\rr{1.8}\draw[solbody](-3,-\rr)--(0,-\rr)arc(-90:90:\rr);
 \fill(-2.2,-\rr)circle(2pt)node[above,sollabel]{$m$};
 \draw[solflow](-1.9,{-\rr+.4})--(-.8,{-\rr+.4})node[midway,above,sollabel]{$v$};
 \fill(0,0)circle(1.5pt)node[left,sollabel]{$O$};
 \draw[solgeom](0,0)--(45:\rr)node[midway,above left,sollabel]{$R$};
 \node[sollabel,anchor=west]at(2.4,0){右半圆内壁};
\end{tikzpicture}\diagcaption{题图}{水平段在圆弧最低点相切，质点先沿右半圆上升。}\label{ch08:fig:probdiagFourTwelve}\end{energyfig}}

\newcommand{\soldiagFourTwelve}{%
\begin{energyfig}\begin{tikzpicture}[scale=1]
 \def\rr{1.9}\pgfmathsetmacro\cc{-1/sqrt(3)}\pgfmathsetmacro\ss{sqrt(2/3)}
 \pgfmathsetmacro\uu{sqrt(-\rr*\cc)}\pgfmathsetmacro\tf{-\rr*\ss/(\uu*\cc)}
 \coordinate(D)at({\rr*\ss},{-\rr*\cc});
 \draw[solbody](-2.2,-\rr)--(0,-\rr)arc(-90:90:\rr);
 \fill(0,0)circle(1.5pt)node[left,sollabel]{$O$};\fill(D)circle(2pt);\node[sollabel,anchor=west]at($(D)+(.3,.25)$){$D$};
 \draw[solgeom](0,-\rr)--(0,0)--(D);\node[sollabel]at(1,.3){$R$};
 \draw[TextGray](0,-.5)arc(-90:35.2644:.5);\node[sollabel]at(.75,-.35){$\varphi$};
 \draw[Blue,line width=1pt,domain=0:\tf,samples=100,variable=\t]
 plot({\rr*\ss+\uu*\cc*\t},{-\rr*\cc+\uu*\ss*\t-.5*\t*\t});
 \draw[solforce](D)--++({1.05*\cc},{1.05*\ss})node[above right,sollabel]{$\vec u$};
 \draw[solforce](D)--++(0,-.65)node[above left,sollabel]{$mg$};
 \node[sollabel,align=left,anchor=east]at(-.65,1.15){$D$ 点 $N=0$\\[4pt]$u_x<0,\ u_y>0$};
\end{tikzpicture}\diagcaption{解图}{脱离后轨迹沿切线进入圆内，水平坐标单调减小，随后经过 $O$。}\label{ch08:fig:soldiagFourTwelve}\end{energyfig}}

\newcommand{\soldiagFourThirteen}{%
\begin{energyfig}\begin{tikzpicture}[scale=.9]
 \begin{scope}[xshift=-3.6cm]
 \node[solpanel]at(0,2.6){(3) 四次势能};\draw[solflow](-1.6,0)--(1.8,0)node[right,sollabel]{$x$};
 \draw[solflow](0,-.3)--(0,2)node[above,sollabel]{$E_p$};
 \draw[solpath,domain=-1.25:1.25,samples=60]plot(\x,{.65*\x*\x*\x*\x});
 \node[sollabel]at(0,-.75){$x=0$ 稳定，$E_p''(0)=0$};
 \end{scope}
 \begin{scope}[xshift=1.7cm]
 \node[solpanel]at(2,2.6){(4) 由定积分确定的吸引势};
 \draw[solflow](0,0)--(4.5,0)node[right,sollabel]{$r$};
 \draw[solflow](0,-1.8)--(0,1.2)node[above,sollabel]{$E_p$};
 \screenedPotentialCurve
 \node[sollabel,align=left]at(2.3,1.35){$E_p'(r)=ke^{-\alpha r}/r^2>0$\\[5pt]$r\to\infty$ 时 $E_p\to0^-$};
 \end{scope}
\end{tikzpicture}\diagcaption{解图}{曲线纵坐标为势能；斜率的负值才是力，不能只凭曲线大致形状作判断。}\label{ch08:fig:soldiagFourThirteen}\end{energyfig}}

\newcommand{\probdiagFourFourteen}{%
\begin{energyfig}\begin{tikzpicture}[scale=.9]
 \begin{scope}[xshift=-3cm]
 \node[solpanel]at(0,2.3){(1)(2) 沿弹簧轴线运动};
 \draw[solbody](-2,0)--(1.8,0);\draw[solbody](-2,0)--(-2,1);
 \draw[spring](-2,.35)--(.5,.35);\draw[solbody,fill=SoftGray](.5,0)rectangle(1.3,.7);
 \node[sollabel,above left]at(-2,.35){$A$};\node[sollabel]at(.9,.35){$m$};
 \draw[solflow](.5,1.2)--(1.6,1.2)node[midway,above,sollabel]{$v_1$ 或 $v_2$};
 \end{scope}
 \begin{scope}[xshift=3.6cm]
 \node[solpanel]at(0,2.3){(3) 竖直运动的三个位置};
 \draw[solbody](-.7,1.7)--(.7,1.7);\node[sollabel,above left]at(0,1.7){$A$};
 \draw[spring](0,1.7)--(0,-1.5);\fill(0,-1.5)circle(3pt);\node[sollabel]at(.30,-1.88){$m$};
 \foreach \y/\txt in {.4/{原长 $x=0$},-.5/{平衡 $x=x_{\rm e}$},-1.5/{初态 $x=h$}}{
 \draw[solgeom](-.5,\y)--(.6,\y);\node[sollabel,anchor=west]at(.85,\y){\txt};}
 \draw[<->,TextGray](-1,-1.5)--(-1,.4)node[midway,left,sollabel]{$h$};
 \draw[solflow](-1.9,-1.5)--(-1.9,-.6)node[left,sollabel]{$v_3$};
 \end{scope}
\end{tikzpicture}\diagcaption{题图}{右图展示 $h>x_{\rm e}$ 的一个初态，解答同时适用于其他 $h$；速度方向与位移基准分开标注。}\label{ch08:fig:probdiagFourFourteen}\end{energyfig}}

\newcommand{\soldiagFourFourteen}{%
\begin{energyfig}\begin{tikzpicture}[scale=.9]
 \begin{scope}[xshift=-3.6cm]
 \node[solpanel]at(0,2.3){水平：平衡在原长位置};
 \draw[solflow](-1.8,0)--(1.9,0)node[right,sollabel]{$x$};\draw[solflow](0,-.2)--(0,1.7)node[above,sollabel]{$E_p$};
 \draw[solpath,domain=-1.5:1.5,samples=60]plot(\x,{.65*\x*\x});
 \node[sollabel,below]at(0,-.25){$0$};\end{scope}
 \begin{scope}[xshift=1.8cm]
 \node[solpanel]at(1.6,2.3){竖直：平衡位置向下移动};
 \draw[solflow](-.3,0)--(3.8,0)node[right,sollabel]{$x$};\draw[solflow](0,-.6)--(0,1.7)node[above,sollabel]{$E_p$};
 \draw[solpath,domain=-.2:3.3,samples=60]plot(\x,{.6*(\x-1.4)^2-.45});
 \draw[solgeom](1.4,-1)--(1.4,1.2);\node[sollabel,below]at(1.4,-1){$x_{\rm e}=mg/k$};
 \draw[solgeom](.2,.414)--(2.6,.414);\node[sollabel,anchor=west]at(2.75,.5){$E$};
 \draw[<->,DeepBlue](1.4,1.15)--(2.6,1.15)node[midway,above,sollabel]{$\mathcal A$};
 \draw[solgeom](2.6,.414)--(2.6,1.15);
 \end{scope}
\end{tikzpicture}\diagcaption{解图}{最低点斜率为零；右图的能量水平线与势阱相交于两个转折点。}\label{ch08:fig:soldiagFourFourteen}\end{energyfig}}


不同的力，做功的性质可能很不相同。一类是重力、理想弹簧弹力等力：它们所做的功只与
物体的初、末位置有关，与物体在力作用下经过的轨迹无关。另一类是滑动摩擦力等各种阻力：
它们所做的功与实际经过的路程有关。阻力会消耗机械能，工程设计中常设法减小它，但在真实
环境中又难以完全避免。当物体在无阻力的空间中运动，或者阻力小到可以忽略时，运动会表现
出一个十分清楚的共同性质。但“没有阻力”还不够：只有保守力作功、其他力不作功时，才有本节讨论的机械能守恒。

\subsection{重力作用下物体的运动}

如图~\ref{ch08:fig:gravityPathsDiagram}，地球表面附近有高度差为 $h$ 的两点 $A,B$。在两点之间设置若干形状
不同、固定不动的光滑轨道，让物体从 $A$ 点沿轨道运动。以下比较以物体确能到达 $B$、途中保持约束为前提；
若路径先升高，初动能还须足够越过最高处。轨道对物体的压力始终垂直于轨道切线，
因而不做功；重力功只取决于初、末位置的高度差：
\gravityPathsDiagram
\begin{equation}
 W_G=mgh .
\end{equation}
设所有物体在 $A$ 点的初速率均为 $v_A$，由动能定理
\begin{equation}
 mgh=\frac12mv_B^2-\frac12mv_A^2.
\end{equation}
移项得到
\begin{equation}
 \frac12mv_B^2=\frac12mv_A^2+mgh.\label{ch08:eq:gravity-kinetic}
\end{equation}
虽然物体在 $B$ 点的速度方向受到各自轨道的限制，但它们的速率和动能相同，与所走路径
无关。也就是说，到达 $B$ 点时的动能等于 $A$ 点的动能与从 $A$ 到 $B$ 的重力功之和；
轨道只改变速度方向，并不改变这个能量关系。

建立竖直向上的 $y$ 轴，设 $A,B$ 两点纵坐标为 $y_A,y_B$，则
\begin{equation}
 h=y_A-y_B.
\end{equation}
代入式~\eqref{ch08:eq:gravity-kinetic}：
\begin{equation}
 \frac12mv_B^2=\frac12mv_A^2+mg(y_A-y_B).
\end{equation}
进一步整理为
\begin{equation}
 \frac12mv_B^2+mgy_B
 =\frac12mv_A^2+mgy_A.
\end{equation}
定义物体在重力场中的重力势能为
\begin{equation}
 E_p(y)=mgy.\label{ch08:eq:gravity-potential}
\end{equation}
上述推理并不依赖 $A,B$ 两点处在某两个特殊高度，因此对运动中任意选取的两个位置都成立。
于是，当物体只受重力以及不做功的约束力时，动能与重力势能之和保持不变。定义机械能
\begin{equation}
 E=E_k+E_p=\frac12mv^2+mgy,
\end{equation}
则整个运动过程中 $E$ 为常量。势能零点可以任意选择；改变零点只会给所有势能同时加上
一个常数，不影响势能差和运动规律。

\begin{origexamplebox}{}\label{ch08:ex:s11-4-10}
质量均为 $m$ 的质点，初速率均为 $v$。轨道固定、无摩擦，不计空气阻力。
高度从初始位置算起，圆周运动均由最低点开始。求以下各情况下首次上升所能达到的最大高度：
\begin{enumerate}[label=\arabic*.,leftmargin=2.2em]
 \item 沿一条高度单调上升且无限高的弯曲光滑槽运动，槽保证质点不脱离；
 \item 以与水平面夹角为 $\theta$ 的方向斜向抛出，$0\le\theta\le\pi/2$；
 \item 套在半径为 $R$ 的光滑圆环上运动；
 \item 在半径为 $R$ 的光滑圆柱体内部运动。
\end{enumerate}
\probdiagFourTen
\end{origexamplebox}

\Needspace{6\baselineskip}\begin{xesolution}
图~\ref{ch08:fig:probdiagFourTen} 的第三种是套住细圆环的珠子，第四种是内壁单侧接触，不能把二者当作同一约束。
以下高度均从初始位置所在的水平面算起。

\textbf{(1) 无限高的光滑轨道}\quad 物体受到的支持力不做功。到达最高点时速率为零，
\[
 \frac12mv^2=mgH_1,
 \qquad
 \boxed{H_1=\frac{v^2}{2g}}.
\]

\textbf{(2) 斜抛}\quad 重力没有水平分量，所以 $v_x=v\cos\theta$ 保持不变。
最高点只有竖直分速度为零，并非总速度为零。将最高点仍有的水平动能扣除：
\[
 mgH_2=\frac12mv^2-\frac12m(v\cos\theta)^2=\frac12m(v\sin\theta)^2.
\]
因此
\[
 \frac12m(v\sin\theta)^2=mgH_2,
 \qquad
 \boxed{H_2=\frac{v^2\sin^2\theta}{2g}}.
\]

\textbf{(3) 套在圆环上}\quad 圆环可同时提供指向圆心或背离圆心的约束力，质点不会
脱离圆环。这是不同状态下可出现的两种方向，并非同时存在两个支持力。
若初始动能不足以到达环顶，它在中途反向；若能到达环顶，几何上的最大高度
就是 $2R$：
\[
 \boxed{H_3=\min\left(\frac{v^2}{2g},\,2R\right)}.
\]

临界情况 $v^2=4gR$ 下，理想质点无限接近环顶，到达时间趋于无穷；此时 $2R$ 严格说是高度上确界。

\textbf{(4) 在圆柱体内部}\quad 这一情形的约束力只能指向圆心，物体可能脱离内壁。
令质点沿内壁从最低点上升的圆心角为 $\varphi$，并记
\[
 q=\frac{v^2}{gR}.
\]
在接触阶段，能量守恒和径向方程分别为
\[
 u^2=v^2-2gR(1-\cos\varphi),
 \qquad
 N=m\left(\frac{u^2}{R}+g\cos\varphi\right).
\]
代入可得 $N=mg(q-2+3\cos\varphi)$。若 $q\le2$，速度先在下半圆降为零；
若 $q\ge5$，环顶仍有 $N=mg(q-5)\ge0$，能保持接触。
所以仅当 $2<q<5$ 时，物体在上半圆由 $N=0$ 脱离。令
\[
 c=\cos\varphi_{\rm d}=\frac{2-q}{3},
\]
则脱离速度满足 $u_{\rm d}^2=-gRc$。图~\ref{ch08:fig:soldiagFourTen} 的右图中，$D$ 为脱离点，
速度沿切线，竖直分量为 $u_{{\rm d},y}=u_{\rm d}\sin\varphi_{\rm d}$。之后作斜抛，
向上速度降为零前增加的高度为 $u_{{\rm d},y}^2/(2g)$。与脱离点高度相加，
得到
\[
 H_4
 =R(1-c)+\frac{u_{\rm d}^2(1-c^2)}{2g}
 =R\left(1-\frac32c+\frac12c^3\right).
\]
\soldiagFourTen
综上，
\[
 \boxed{
 H_4=
 \begin{cases}
 \dfrac{v^2}{2g},&0\le q\le2,\\[6pt]
 R\left(1-\dfrac32c+\dfrac12c^3\right),
   &2<q<5,\quad c=\dfrac{2-q}{3},\\[6pt]
 2R,&q\ge5.
 \end{cases}}
\]
第一段中质点在到达圆环侧面以前已经反向；中间一段先脱离再作抛体；最后一段始终接触
内壁并越过最高点。在 $q=2,5$ 处相邻表达式分别都给出 $R,2R$，结果连续，每项均为长度。
这里问首次上升；如果再碰到内壁，后续还与碰撞模型有关。只用能量写 $H=v^2/(2g)$ 会漏掉约束失效后的运动。
\end{xesolution}

\Needspace{8\baselineskip}
\begin{origexamplebox}{}\label{ch08:ex:s11-4-11}
一个质量为 $m$ 的质点从固定在水平面上、半径为 $R$ 的半圆柱形光滑表面的顶端由静止
开始下滑。求质点脱离圆柱体时，质点与圆心连线和竖直方向的夹角 $\theta$；并求落地点
与圆心的距离。理想的正顶点静止状态是非稳定平衡；“开始下滑”按微小扰动趋于零的极限理解。
\probdiagFourEleven
\end{origexamplebox}

\Needspace{6\baselineskip}\begin{xesolution}
图~\ref{ch08:fig:probdiagFourEleven} 中半圆的圆心在地面，顶点距圆心确为 $R$。
设质点在圆柱面上运动到圆心角 $\theta$ 时速率为 $u$。从顶端由静止下滑，机械能守恒：
\[
 \frac12mu^2=mgR(1-\cos\theta),
 \qquad
 u^2=2gR(1-\cos\theta).
\]
取指向圆心为径向正方向。质点在圆柱外表面运动时
\[
 mg\cos\theta-N=\frac{mu^2}{R}.
\]
脱离瞬间 $N=0$，所以
\[
 g\cos\theta=\frac{u^2}{R}.
\]
与能量式联立：
\[
 \cos\theta=2(1-\cos\theta),
 \qquad
 \boxed{\cos\theta=\frac23},
\]
即
\[
 \boxed{\theta=\arccos\frac23}.
\]
此时
\[
 \sin\theta=\frac{\sqrt5}{3},
 \qquad
 u=\sqrt{\frac{2gR}{3}}.
\]
\soldiagFourEleven
图~\ref{ch08:fig:soldiagFourEleven} 以圆心为原点、水平向右为 $x$ 轴、竖直向上为 $y$ 轴。
半径方向为 $(\sin\theta,\cos\theta)$，向右下方的单位切向量为 $(\cos\theta,-\sin\theta)$，二者点积为零。
所以脱离点坐标和速度分量为
\[
 x_0=\frac{\sqrt5}{3}R,\quad
 y_0=\frac23R,\quad
 u_x=\frac23\sqrt{\frac{2gR}{3}},\quad
 u_y=-\frac{\sqrt5}{3}\sqrt{\frac{2gR}{3}}.
\]
脱离后水平速度不变、竖直加速度为 $-g$。落地时 $y=0$，故
\[
 \frac12gt^2-u_yt-y_0=0,\qquad t=\frac{u_y+\sqrt{u_y^2+2gy_0}}{g}.
\]
另一根为负，不能表示脱离后的时刻。代入 $u_y$ 和 $y_0$，得到
\[
 t=\sqrt{\frac{R}{g}}\,
 \frac{\sqrt{46}-\sqrt{10}}{3\sqrt3}.
\]
因此落地点到圆心的水平距离为
\[
 \boxed{
 x=x_0+u_xt
 =\frac{5\sqrt5+4\sqrt{23}}{27}\,R}.
\]
结果约为 $1.12458R$，大于半圆右端的横坐标 $R$，量纲也是长度。
脱离后支持力消失，必须改用抛体运动，不能把圆周运动方程继续使用到落地。
\end{xesolution}

\Needspace{8\baselineskip}
\begin{origexamplebox}{}\label{ch08:ex:s11-4-12}
如图，光滑水平面上质量为 $m$ 的物体以初速度 $v$ 运动，冲向一个半径为 $R$ 的光滑
圆形表面。图~\ref{ch08:fig:probdiagFourTwelve} 的圆弧为固定的右半圆内壁，与水平面在最低点相切。
求使物体脱离圆弧后能够击中圆心 $O$ 的初速度 $v$。
\probdiagFourTwelve
\end{origexamplebox}

\Needspace{6\baselineskip}\begin{xesolution}
以圆心 $O$ 为原点，令物体从圆弧最低点开始运动，半径由竖直向下转过角
$\varphi$。接触点坐标为
\[
 x=R\sin\varphi,\qquad y=-R\cos\varphi.
\]
设脱离时速率为 $u$。物体在圆弧内侧运动，能量守恒和脱离条件为
\[
 u^2=v^2-2gR(1-\cos\varphi),
\qquad
 N=m\left(\frac{u^2}{R}+g\cos\varphi\right)=0.
\]
因此
\begin{equation}
 u^2=-gR\cos\varphi,\qquad v^2=gR(2-3\cos\varphi).
 \label{ch08:eq:circle-entry}
\end{equation}
脱离点在右上象限，故 $\cos\varphi<0$。切向速度分量为
\[
 u_x=u\cos\varphi,\qquad u_y=u\sin\varphi.
\]
若抛体经过圆心，则存在 $t>0$ 使
\[
 0=R\sin\varphi+u\cos\varphi\,t,
\]
\[
 0=-R\cos\varphi+u\sin\varphi\,t-\frac12gt^2.
\]
\soldiagFourTwelve
图~\ref{ch08:fig:soldiagFourTwelve} 中 $D$ 在右上方、速度指向左上方，抛出后水平分速度保持为负。
令 $c=\cos\varphi<0$、$s=\sin\varphi>0$。由横坐标方程得
\[
 t=-\frac{Rs}{uc}>0.
\]
代入纵坐标方程，除以 $R$，再用 $u^2=-gRc$：
\[
 0=-c-\frac{s^2}{c}-\frac{gRs^2}{2u^2c^2}
   =-c-\frac{s^2}{c}+\frac{s^2}{2c^3}.
\]
乘以 $2c^3$ 并代入 $s^2=1-c^2$，逐项整理：
\[
 0=-2c^4-2s^2c^2+s^2
 =-2c^4-2(1-c^2)c^2+1-c^2=1-3c^2.
\]
因此
\[
 1-3\cos^2\varphi=0.
\]
结合 $\cos\varphi<0$：
\[
 \cos\varphi=-\frac1{\sqrt3}.
\]
代入式~\eqref{ch08:eq:circle-entry}，得到
\[
 \boxed{v=\sqrt{(2+\sqrt3)\,gR}}.
\]
结果中 $v^2/(gR)=2+\sqrt3$ 在 $2$ 与 $5$ 之间，符合在上半圆脱离的条件；$\sqrt{gR}$ 的量纲为速度。
这里“击中圆心”发生在脱离后的抛体阶段；若不先判定脱离条件，不能确定正确的初速度。
\end{xesolution}

\section{保守力、势能与稳定性}

重力做功只与始末高度有关。具有同样性质的力称为保守力。保守力的关键性质可整理为：
\begin{enumerate}[leftmargin=2.2em]
 \item 保守力的大小和方向只由位置决定，不显含时间或速度。近地面重力、固定端理想弹簧
 弹力，以及以后学习的万有引力、静电场中的电场力，都是典型例子。
 \item 保守力做功只与初、末位置有关，与实际路径无关。重力功只取决于高度差，弹簧弹力
 的功只取决于始、末形变量。
 \item 由路径无关性立即得到：物体沿任意闭合路径运动一周，保守力的总功必为零。
 \item 对保守力可以定义只依赖位置的势能。势能不依赖物体的速度、运动时间或到达该位置
 的路径，并且
 \[
 W_{A\to B}=E_p(A)-E_p(B)=-\Delta E_p.
 \]
 势能的形式由相应的保守力决定，势能零点可以自由选取；真正影响运动的是势能差。
 \item 已知保守力随位置的关系并选定势能零点，可以求出各位置的势能；反过来，已知势能
 也可以求出力的大小和方向。空间中势能相等的点构成等势面，保守力与等势面垂直，并指向
 势能下降最快的方向。
 \item 若物体只受保守力和不做功的约束力作用，便可把机械能定义为动能与各项势能之和，
 且运动过程中机械能保持不变；这个常量可由运动中任意一个位置的动能和势能确定。
\end{enumerate}
重力势能已由式~\eqref{ch08:eq:gravity-potential} 给出。劲度系数为 $k$ 的理想弹簧，形变量为 $x$ 时，弹性势能为
\begin{equation}
 E_p(x)=\frac12kx^2.
\end{equation}
对于万有引力，若把无穷远处选作势能零点，两质量 $M,m$ 相距 $r$ 时的势能为
\begin{equation}
 E_p(r)=-\frac{GMm}{r}.
\end{equation}
势能的绝对值取决于零点选择，真正决定运动的是势能差。一维微小位移满足
$F_x\dd x=-\dd E_p$，故 $F_x=-\dd E_p/\dd x$：曲线斜率为正时力向负方向，斜率为负时力向正方向。
三维情况下，势能一般是 $E_p(x,y,z)$，全微分为
\[
 \dd E_p=\frac{\partial E_p}{\partial x}\dd x+
 \frac{\partial E_p}{\partial y}\dd y+\frac{\partial E_p}{\partial z}\dd z.
\]
例如 $\partial E_p/\partial x$ 表示保持 $y,z$ 不变时，势能随 $x$ 的变化率。
与 $\vec F\cdot\dd\vec r=-\dd E_p$ 逐项比较，三个分量分别是对应偏导数的负值。
把三个偏导数组成矢量，称为势能的梯度，记为
\[
 \nabla E_p=\left(\frac{\partial E_p}{\partial x},\frac{\partial E_p}{\partial y},\frac{\partial E_p}{\partial z}\right),
 \qquad \vec F=-\nabla E_p.
\]
在梯度非零的地方，梯度指向势能增加最快的方向；因此保守力与等势面垂直，指向势能降低最快的方向。
若物体只受保守力作用，则动能和相应势能之和守恒，即
\[
 E=E_k+E_p=\text{常量}.
\]

\Needspace{22\baselineskip}
\begin{origexamplebox}{}\label{ch08:ex:s11-4-13}
判断下列各力是否为保守力；若是，分析其势能函数的性质。各小题独立，系数取正值；
$a$ 在第(1)问是力随时间增加的系数，不是加速度，各题 $\alpha$ 的单位由所在公式决定。
\begin{enumerate}[label=\arabic*.,leftmargin=2.2em]
 \item 正比于时间的拉力 $f=a t$；
 \item 正比于速度、方向与速度相反的阻力 $f=-\alpha v$；
 \item 正比于坐标三次方的力 $f=-kx^3$；
 \item 随距离变化的吸引力
 \[
 f(r)=-k\frac{e^{-\alpha r}}{r^2}.
 \]
\end{enumerate}

\ThQConservative
\end{origexamplebox}

\Needspace{6\baselineskip}\begin{xesolution}
\textbf{(1)} $f=at$ 显含时间。同一位置在不同时刻受到的力不同，不能由单值的位置势能
确定，因而不是通常意义下的保守力。

\textbf{(2)} $f=-\alpha v$ 依赖速度。其瞬时功率为
\[
 P=fv=-\alpha v^2\le0,
\]
机械能随运动永久耗散，故不是保守力。

\textbf{(3)} $f=-kx^3$ 只依赖位置，在一维中为保守力。由
\[
 -\frac{\dd E_p}{\dd x}=-kx^3
\]
得到
\[
 \boxed{E_p(x)=\frac{k}{4}x^4+C}.
\]
若取 $E_p(0)=0$，则 $C=0$。势能在 $x=0$ 取得稳定极小值，关于原点对称。
但最低点的二阶导数为零，小振动不能直接当作具有非零固有频率的简谐振动。

\textbf{(4)} 这是只依赖距离的中心力，在 $r>0$ 的区域内为保守力。取
$E_p(\infty)=0$，可写成
\[
 \boxed{
 E_p(r)=-\int_r^\infty
 k\frac{e^{-\alpha s}}{s^2}\,\dd s}.
\]
上述定积分已经完整定义了势能，不必依赖特殊函数。若进一步引入指数积分函数
\[
 \operatorname{Ei}(-z)=-\int_z^\infty\frac{e^{-t}}{t}\,\dd t,\qquad z>0,
\]
注意 $\dd(e^{-\alpha s}/s)/\dd s=-\alpha e^{-\alpha s}/s-e^{-\alpha s}/s^2$，因此分部积分可得
\[
 \int_r^\infty\frac{e^{-\alpha s}}{s^2}\dd s
 =\frac{e^{-\alpha r}}r-\alpha\int_r^\infty\frac{e^{-\alpha s}}s\dd s
 =\frac{e^{-\alpha r}}r+\alpha\operatorname{Ei}(-\alpha r).
\]
于是也可表示为
\[
 E_p(r)
 =-k\left[
 \frac{e^{-\alpha r}}{r}
 +\alpha\,\operatorname{Ei}(-\alpha r)
 \right].
\]
\soldiagFourThirteen
图~\ref{ch08:fig:soldiagFourThirteen} 的右图用这个定积分作图，而不是 $-ke^{-\alpha r}/r$；后者求负导数会多出一项，不能代表题给的力。
由微积分基本定理，$E_p'(r)=ke^{-\alpha r}/r^2>0$，负导数恰为题给的引力。
它随 $r$ 单调增加：$r\to\infty$ 时趋于 $0$，$r\to0^+$ 时趋于 $-\infty$，对应始终
指向内侧的吸引力。前三项中，只有第 (3) 项具有普通的位置势能；第 (4) 项虽形式复杂，
但是否保守只取决于它是位置的单值中心力，而不取决于积分能否用初等函数表示。
\end{xesolution}

\Needspace{8\baselineskip}
\begin{origexamplebox}{}\label{ch08:ex:s11-4-14}
劲度系数为 $k$、自然长度为 $l_0$ 的弹簧一端固定在给定点 $A$，另一端连接质量为
$m$ 的物体。前两问限定物体沿弹簧轴线作一维运动，弹簧可拉伸也可压缩，忽略碰壁等几何限制。
求以下各情况下弹簧的最大伸长量：
\begin{enumerate}[label=\arabic*.,leftmargin=2.2em]
 \item 物体在与 $A$ 同一水平面的光滑平面上运动。初始时弹簧为原长，物体初速度为
 $v_1$；
 \item 仍在上述光滑水平面内运动。初始时弹簧压缩量为 $\delta l$，初速度为 $v_2$，
 方向指向平衡位置。若初速度方向背离平衡位置，结果如何？
 \item 物体在通过 $A$ 的竖直线上运动。初始时位于原长位置以下 $h$ 处，速度 $v_3$
 竖直向上。若初速度改为竖直向下，结果如何？
\end{enumerate}
\probdiagFourFourteen
\end{origexamplebox}

\Needspace{6\baselineskip}\begin{xesolution}
图~\ref{ch08:fig:probdiagFourFourteen} 将水平轴线和竖直轴线分开表示，固定点都记为 $A$。
以下用 $x$ 表示弹簧相对原长的形变量，规定伸长为正；“最大伸长量”就是运动中最大的
$x$。

\textbf{(1)} 水平面上没有重力势能变化。初态 $x=0$，在最大伸长处速度为零：
\[
\frac12mv_1^2=\frac12kx_{\max}^2,
\qquad
\boxed{x_{\max}=v_1\sqrt{\frac{m}{k}}}.
\]

\textbf{(2)} 初态为 $x=-\delta l$。在正向最大伸长处速度为零，故
\[
\frac12mv_2^2+\frac12k(\delta l)^2
=\frac12kx_{\max}^2,
\qquad
\boxed{
 x_{\max}=\sqrt{(\delta l)^2+\frac{m}{k}v_2^2}}.
\]
初速度若背离平衡位置，物体先到压缩侧的转折点，再向另一侧运动；由于总机械能不变，
最终的最大伸长量仍相同，只有到达各转折点的先后次序改变。

\textbf{(3)} 取竖直向下为 $x$ 正方向，忽略与原长 $l_0$ 有关的势能常数：
\[
E_p(x)=\frac12kx^2-mgx
=\frac12k\left(x-\frac{mg}{k}\right)^2+\text{常量}.
\]
竖直运动的平衡形变量和相对该位置的振幅分别为
\[
x_{\rm e}=\frac{mg}{k},
\qquad
\mathcal A=\sqrt{\left(h-\frac{mg}{k}\right)^2+\frac{m}{k}v_3^2}.
\]
\soldiagFourFourteen
这里 $\mathcal A$ 表示振幅，区别于固定点 $A$。它来自初态与转折点的能量比较：
\[
 \frac12mv_3^2+\frac12k(h-x_{\rm e})^2=\frac12k\mathcal A^2.
\]
图~\ref{ch08:fig:soldiagFourFourteen} 中势阱关于 $x_{\rm e}$ 对称，两个转折点为 $x_{\rm e}\pm\mathcal A$。
最大伸长位置为 $x_{\rm e}+\mathcal A$，所以
\[
\boxed{
 x_{\max}=\frac{mg}{k}
 +\sqrt{\left(h-\frac{mg}{k}\right)^2+\frac{m}{k}v_3^2}}.
\]
把初速度由向上改为向下，只改变速度符号，不改变机械能和振幅，故最大伸长量仍不变。
取 $v_3=0,h=x_{\rm e}$ 时振幅为零，物体保持平衡，结果给出 $x_{\max}=x_{\rm e}$。
若前两问不限制为沿轴线运动，非径向初速度会产生角动量；最大伸长处径向速度虽为零，切向速度却不一定为零。
此时需另用角动量守恒和包含切向动能的能量式，不能直接使用本题的一维答案。
\end{xesolution}

\subsection{势能方法的适用范围与稳定性}

势能函数也常记为 $U$。近地面重力势能 $mgy$、万有引力势能 $-GMm/r$、弹簧势能 $kx^2/2$
来自不同的力随位置的关系，不能跨越适用范围混用。尤其当高度与地球半径可比时，$g$ 随位置变化，$mgy$ 不再准确。

先选出已用势能 $U$ 表示的、与时间无关的保守相互作用；把所有尚未计入 $U$ 的力的功合记为 $W_{\rm other}$。
动能定理给出 $\Delta K=-\Delta U+W_{\rm other}$，所以
\begin{equation}
 \Delta(K+U)=W_{\rm other}.
 \label{ch08:eq:other-work}
\end{equation}
当 $W_{\rm other}=0$，机械能守恒。若还要把右边分为外力功和非保守内力功，两类必须互不重叠，
不能把同一个外部摩擦力计算两次。也不能既写重力功，又把同一重力势能变化列入左边。
在固定外场中可以把外场作用写入势能；例如研究近地面质点时使用 $mgy$，不必每次显式计算地球的运动。
但若场源也运动或相互作用显含时间，就须重新检查这个简化是否成立。

一维平衡点 $x_0$ 满足 $U'(x_0)=0$，即合力为零。若 $U''(x_0)>0$，势能在该点取局部极小，
微小偏离后出现指向平衡点的回复力，是稳定平衡；若 $U''(x_0)<0$，回复方向相反，平衡不稳定。
在二阶导数为正的平衡点附近展开：
\[
 U(x)\simeq U(x_0)+\frac12U''(x_0)(x-x_0)^2,
 \qquad F_x\simeq-U''(x_0)(x-x_0).
\]
与 $m\ddot x=-k_{\rm eff}(x-x_0)$ 比较可知，$k_{\rm eff}=U''(x_0)$，故小振动角频率
$\omega=\sqrt{U''(x_0)/m}$。如果 $U''(x_0)=0$，这个二阶判据不能定性，必须继续查看更高阶项；
例如 $U=kx^4/4$ 仍是稳定势阱，却不是上述线性回复模型。

给定一维总机械能 $E$，允许位置满足 $E\ge U(x)$；转折点满足 $E=U(x)$。在允许区间内
\[
 |v(x)|=\sqrt{\frac{2[E-U(x)]}{m}}.
\]
这只能先给出速率，运动方向由所处的往返分支确定。若需时间，在方向不变的一段上取
$v(x)=\pm\sqrt{2[E-U(x)]/m}$，再积分 $t_B-t_A=\int_{x_A}^{x_B}\dd x/v(x)$。
因此能量曲线既能判断可达区域，也能与运动方程配合求运动时间。
}

{
\graphicspath{{../../../sample/xes-original/thirteenth-section/}}
\section{能量的转化与守恒}
机械能守恒是有条件的，能量守恒则讨论系统中所有形式的能量。摩擦、非弹性碰撞或物体的塑性形变
会使宏观机械能减少，但这部分能量并没有消失，而是转入物体的内能、形变能、声能等其他形式。

能量并不只有动能和势能。物体的无规则热运动、电磁场、光辐射等也与能量有关；不同的能量形式可以在相互作用中转化。
例如摩擦使运动减慢并使接触体温度升高，电动机把电能转化为机械能，发光过程又可以把内部能量转化为辐射能。
这些形式不能各自孤立看待：研究一种形式减少了多少，还要追踪它传到哪个物体、变成哪种形式。

能量概念的适用范围也不限于具有确定位置、确定速度的经典质点。光没有静止质量，却能传递能量；微观系统即使不能
按经典轨道描述，也仍然可以讨论它与外界交换的能量。物理研究依靠可观测的相互作用，能量传递是把不同现象联系起来的重要线索，
但不应把“能量”与物体的质量或某一个可见运动简单等同。

人们对能量的认识并不是一次完成的。在研究热、机械运动以及新的微观现象时，曾对能量是否守恒提出过疑问。
面对表面上不相等的初末能量，不能仅凭一句“守恒”结束讨论，而要检验测量误差、系统边界，以及是否遗漏了热、辐射或其他携能对象。
将这些独立测得的能量传递纳入比较，才构成对守恒关系的实验检验。能量守恒有广泛的实验依据；本章在经典力学的适用范围内，
把它用于明确选定的系统，而不是用它替代对具体过程和相互作用的分析。

\subsection{先选系统，再写能量守恒关系}

“能量守恒”必须对一个明确的系统而言。同一个过程，选取的系统边界不同，方程中的“外力做功”和
“系统内部的能量转化”也不同。对一个无质量穿过边界的封闭系统，取“外界传入系统”为正，可写成
\begin{equation}
  \Delta E_{\rm sys}=W_{\rm ext}+Q,
  \label{ch08:eq:general-energy-account}
\end{equation}
其中 $W_{\rm ext}$ 是外界通过宏观功传入的能量，$Q$ 是由于温差传入的能量。在通常的短时间力学过程中，
若热传递可忽略，可取 $Q=0$。

对孤立系统，$W_{\rm ext}=Q=0$，于是
\begin{equation}
  \boxed{\Delta E_{\rm sys}=0}.
  \label{ch08:eq:isolated-energy}
\end{equation}
这里的 $E_{\rm sys}$ 不只是动能加势能，而是系统的总能量。在本章可用的层次上，常写为
\begin{equation}
  E_{\rm sys}=K+U+E_{\rm int}+E_{\rm other}.
  \label{ch08:eq:energy-components}
\end{equation}
其中 $K$ 是宏观动能，$U$ 是保守相互作用所对应的势能，$E_{\rm int}$ 可包括与无规则热运动、塑性形变等相关的内部能量。

图~\ref{ch08:fig:energy-forms} 表示孤立系统中能量形式可以改变，而总能量保持不变；箭头只是转化示例，不表示转化只能单向发生。
\begin{center}\begin{minipage}{.98\linewidth}\centering
\begin{tikzpicture}[node distance=11mm and 24mm]
  \node[draw=RuleGray,fill=PaleBlue,minimum width=25mm,minimum height=9mm] (mech) {机械能 $K+U$};
  \node[draw=RuleGray,fill=SoftGray,minimum width=23mm,minimum height=9mm,right=of mech] (int) {内能 $E_{\rm int}$};
  \node[draw=RuleGray,fill=SoftGray,minimum width=23mm,minimum height=9mm,right=of int] (other) {其他形式};
  \draw[-{Latex[length=2.2mm]},thick,draw=DeepBlue] (mech)--node[above,font=\scriptsize,fill=none]{摩擦、形变}(int);
  \draw[-{Latex[length=2.2mm]},thick,draw=TextGray] (int)--node[above,font=\small]{转化}(other);
  \node[below=7mm of int,font=\small\color{TextGray}] {总量不变，分项可改变};
\end{tikzpicture}
\refstepcounter{figure}\par\smallskip{\scriptsize\color{TextGray}图~\thefigure\quad 孤立系统内部的能量转化示例。\par}
\label{ch08:fig:energy-forms}
\end{minipage}\end{center}

\subsection{“机械能损失”不等于“能量损失”}

若把已选定的保守力全部纳入势能，机械能定义为
\[
  E_{\rm mech}=K+U.
\]
对只有这些保守力做功的过程，$\Delta E_{\rm mech}=0$。若有滑动摩擦、流体阻力等非保守作用，则一般有
\begin{equation}
  \Delta E_{\rm mech}=W_{\rm nc}.
  \label{ch08:eq:mechanical-energy-change}
\end{equation}
式 \eqref{ch08:eq:mechanical-energy-change} 描述机械能的变化，而不是说总能量减少了 $|W_{\rm nc}|$。

摩擦问题尤其需要先说清系统。若只选滑块为系统，地面对它的摩擦是外力；若选“滑块+地面”为系统，
这一接触作用变成内部作用，宏观机械能的减少对应系统内能的增加。不知道材料细节时，通常只能确定两接触体
内能增量的总和，不能凭空判定热量如何在两者之间分配。

\Needspace{14\baselineskip}
\begin{origexamplebox}{粗糙水平面上的停车距离}\label{ch08:ex:s13-4-21}
质量为 $m$ 的滑块以速度 $v_0$ 沿粗糙水平面滑动，动摩擦因数为 $\mu_k$。求滑块停下前的距离，
并分别用“滑块”和“滑块+地面”为系统说明能量变化。

\ThQStopping
\end{origexamplebox}

\Needspace{6\baselineskip}\begin{xesolution}
\ThAStopping

滑块竖直方向平衡，所以 $N=mg$，滑动摩擦力大小为 $f_k=\mu_kmg$。设停车距离为 $s$。

若只选滑块为系统，摩擦力是外力，由动能定理
\[
  0-\frac12mv_0^2=-\mu_kmgs,
\]
因而
\[
  \boxed{s=\frac{v_0^2}{2\mu_kg}}.
\]
若改选“滑块+地面”为系统，摩擦成为内部作用。忽略对外散热和声音携带的能量，总能量守恒给出
\[
  \frac12mv_0^2=\Delta E_{\rm int},
\]
即
\[
  \boxed{\Delta E_{\rm int}=\mu_kmgs=\frac12mv_0^2}.
\]
两种写法对同一过程的描述完全一致，只是系统边界不同。
\end{xesolution}

\subsection{含摩擦的弹簧系统：别忘了分段和启动条件}

\begin{origexamplebox}{弹簧能否把物块送到原长位置}\label{ch08:ex:s13-4-22}
水平粗糙面上，质量为 $m$ 的物块与劲度系数为 $k$ 的水平弹簧相连。弹簧原长位置记为 $O$。
以物块左侧弹簧连接点为位置标志，先从原长位置 $O$ 向右拉到 $x=a>0$ 处由静止释放。
已知最大静摩擦因数为 $\mu_s$，动摩擦因数为 $\mu_k$，$\mu_s\ge\mu_k\ge0$。
求物块能够启动、到达 $O$ 点和通过 $O$ 点的条件。

\begin{center}\begin{minipage}{.98\linewidth}\centering
\begin{tikzpicture}[scale=0.9,every node/.style={font=\scriptsize,fill=none}]
  \draw[thick] (-3.2,-0.8)--(3.4,-0.8);
  \foreach \x in {-3.0,-2.6,...,3.2}{\draw[thin,TextGray] (\x,-0.8)--++(-0.18,-0.18);}
  \draw[thick] (-3.0,0.15)--(-3.0,-0.8);
  \draw[decorate,decoration={coil,aspect=0.45,segment length=3.5mm,amplitude=2.2mm},thick]
    (-3.0,-0.3)--(1.8,-0.3);
  \draw[fill=PaleBlue,thick] (1.8,-0.8) rectangle (2.8,0.2);
  \node at (2.3,-.3) {$m$};
  \draw[guide] (0,-.8)--(0,-1.4) node[below]{$O:x=0$};
  \draw[guide] (1.8,-.8)--(1.8,-1.4);
  \draw[<->] (0,-1.2)--(1.8,-1.2) node[midway,above]{$a$};
  \draw[-{Latex}] (1.8,.7)--(.8,.7) node[midway,above] {释放后向左};
\end{tikzpicture}
\refstepcounter{figure}\par\smallskip{\scriptsize\color{TextGray}图~\thefigure\quad 物块底面与粗糙面接触；$a$ 从同一个连接点的原长位置量起。\par}
\label{ch08:fig:rough-spring}
\end{minipage}\end{center}
\end{origexamplebox}

\Needspace{6\baselineskip}\begin{xesolution}
\ThARoughSpring

图~\ref{ch08:fig:rough-spring} 中弹簧向左拉，启动前静摩擦力向右。释放瞬间弹簧力大小为 $ka$。
若 $ka\le\mu_smg$，静摩擦可以平衡它，物块可保持静止；不依赖额外扰动而启动的条件为
\begin{equation}
  \boxed{ka>\mu_smg}.
  \label{ch08:eq:spring-start}
\end{equation}
启动后，从 $x=a$ 向左运动到某一位置 $x$ 时，摩擦力所做的功为 $-\mu_kmg(a-x)$。对“物块+弹簧”写机械能变化关系，
\[
  \frac12mv^2+\frac12kx^2
  =\frac12ka^2-\mu_kmg(a-x).
\]
整理任一中间位置的速度：
\[
 mv^2=(a-x)\bigl[k(a+x)-2\mu_kmg\bigr].
\]
从 $a$ 到 $0$ 的途中，第一因子非负，第二因子的最小值出现在 $x=0$。
所以只要原点处的能量条件满足，途中不会先停下；不能只检查末态而不管过程可达性。
在 $O$ 点有 $x=0$，因而
\begin{equation}
  \frac12mv_O^2=\frac12ka^2-\mu_kmga.
  \label{ch08:eq:spring-at-origin}
\end{equation}
式 \eqref{ch08:eq:spring-at-origin} 右边非负时，物块才能到达 $O$；严格为正时才能以非零速度通过。因此
\[
  \boxed{a\ge\frac{2\mu_kmg}{k}\quad\text{可到达 }O},
\]
\[
  \boxed{a>\frac{2\mu_kmg}{k}\quad\text{可通过 }O}.
\]
完整答案必须同时满足启动条件 \eqref{ch08:eq:spring-start}。这也说明为什么不能把“摩擦因数”只记成一个统一的 $\mu$：
启动由 $\mu_s$ 控制，滑动阶段的能量损失则由 $\mu_k$ 控制。
\end{xesolution}

\Needspace{16\baselineskip}
\subsection{非弹性碰撞：动量守恒，动能不守恒}

\begin{origexamplebox}{子弹射入木块}\label{ch08:ex:s13-4-23}
质量为 $m$ 的子弹以速度 $u$ 水平射入质量为 $M$ 的静止木块并留在其中。碰撞时间极短，
可忽略水平外力冲量。求共同速度以及碰撞中增加的内能。

\ThQBullet
\end{origexamplebox}

\Needspace{6\baselineskip}\begin{xesolution}
\ThABullet

以“子弹+木块”为系统。碰撞时间内水平外力冲量可忽略，所以水平动量守恒。设碰后共同速度为 $v$，
\[
  m u=(M+m)v,\qquad \boxed{v=\frac{m}{M+m}u}.
\]
将共同速度代入碰撞前后的宏观动能差，逐项化简：
\begin{align*}
  \Delta E_{\rm int}
  &=\frac12mu^2-\frac12(M+m)v^2\\
  &=\frac12mu^2-\frac12(M+m)
    \left(\frac{m}{M+m}u\right)^2\\
  &=\boxed{\frac12\frac{mM}{m+M}u^2}.
\end{align*}
在忽略碰撞期间对外散热和声音携能的近似下，这一动能减少量转化为子弹与木块的内能和永久形变能。
不能在这一碰撞中同时强行使用“动量守恒”和“动能守恒”。
\end{xesolution}

\subsection{一个可执行的解题次序}\nopagebreak[4]

处理能量转化问题时，可以按以下顺序检查：
\begin{enumerate}[label=\arabic*.,leftmargin=2.2em]
  \item 用一句话写清系统包含哪些物体，以及初、末状态。
  \item 判断哪些相互作用纳入势能，不要把同一个保守力同时既写成功又写成势能变化。
  \item 列出跨过系统边界的功、热和其他能量传递。
  \item 若有摩擦、碰撞或形变，明确机械能减少后转入了哪一项，不要简写成“能量消失”。
  \item 最后检查结果的符号、量纲和极限情形。
\end{enumerate}
\nopagebreak[4]
能量方法的优势是跨越中间过程，直接联系两个状态；它的代价是通常不直接给出时间、轨迹和力的瞬时值。
竞赛题中若既问速度又问支持力或时间，常见的组合是：先用能量法求状态量，再用牛顿定律或冲量定理求局部过程量。
}

\endgroup

% ThermalQualityRevision20261023-v1
% QualityRevision20261023
\hypersetup{bookmarksdepth=3}
\begingroup
\tcbset{breakable/.default=false}
\chapter{角动量与刚体动力学}
% Pure problem geometry. Rich legacy force/process figures now appear in solutions.
\providecommand{\ThermalFigure}[3]{\begin{qualityfigure}
\begin{tikzpicture}[>=Latex,font=\small,x=.8cm,y=.8cm]#1\end{tikzpicture}
\figcap{#2}\label{#3}\end{qualityfigure}}
\newcommand{\ThVectorBase}[1]{
 \coordinate(A)at(0,0);\coordinate(C)at(3,1);\coordinate(B)at(3,3);
 \fill(A)circle(2pt)node[left]{$A$};\fill(C)circle(2pt)node[right]{$C$};\fill(B)circle(2pt)node[above]{$B$};
 \draw[->,thick](A)--(B)node[midway,above left]{$(1,1,1)$};\draw[->,DeepBlue,thick](A)--(C)node[midway,below]{$(1,1,0)$};
 \draw[dashed](B)--(C);#1}
\newcommand{\ThQVectors}{\ThermalFigure{\ThVectorBase{}}{两矢量共起点于 $A$；$C$ 是 $B$ 向 $z=0$ 平面的正交投影，图为斜投影。}{ch09:quality:vectors-q}}
\newcommand{\ThAVectors}{\ThermalFigure{\ThVectorBase{
 \node[right,align=left]at(4.7,1.8){$AC=\sqrt2$\\$BC=1$\\$AB=\sqrt3$};
}}{直角三角形 $ACB$ 提供独立的夹角核验；叉积方向垂直两矢量所在平面。}{ch09:quality:vectors-a}
图~\ref{ch09:quality:vectors-a}中的 $BC$ 沿 $z$ 方向，而 $AC$ 位于 $z=0$，故二者垂直。三角形给出 $\sin\theta=BC/AB$、$\cos\theta=AC/AB$，与分量内外积的结果互相核验；交换叉积次序只反转法向，不改变其大小。
}
\newcommand{\ThQArea}{\ThermalFigure{
 \fill[DeepBlue!10](0,0)--(4,.3)--(3.6,1.5)--cycle;
 \draw[->](0,0)--(4,.3)node[midway,below]{$\vect r$};\draw[->](0,0)--(3.6,1.5);
 \draw[->,thick](4,.3)--(3.6,1.5)node[right]{$\Delta\vect r$};\node[left]at(0,0){$O$};
}{质点在短时间内从一条位矢的末端移到另一条位矢的末端，扫过窄三角形。}{ch09:quality:area-q}}
\newcommand{\ThQInertias}{\ThermalFigure{
 \draw[very thick](-1.3,0)--(1.3,0);\draw[->](0,-.8)--(0,1.3);\node[below]at(0,-1){细杆 $l$，中垂轴};
 \begin{scope}[shift={(4,0)}]\draw[thick](0,0)circle(1);\draw[->](-1.5,0)--(1.5,0);\node[below]at(0,-1.3){薄圆盘 $R$，直径轴};\end{scope}
 \foreach\x/\lab in {8/薄球壳,12/实心球}{\begin{scope}[shift={(\x,0)}]\draw[thick](0,0)circle(1);\draw[dashed](0,0)ellipse(1 and .3);\draw[->](0,-1.4)--(0,1.4);\node[below]at(0,-1.7){\lab $R$};\end{scope}}
}{四种均匀分布均具有质量 $m$；箭头明确指定各问所求转轴。}{ch09:quality:inertias-q}}
\newcommand{\ThQConePendulum}{\ThermalFigure{
 \draw[thick](0,3)--(2,0)node[midway,right]{$l$};\draw[dashed](0,3)--(0,0)--(2,0);
 \draw[dashed](0,0)ellipse(2 and .5);\fill(0,3)circle(2pt)node[above]{$O$};\fill[DeepBlue](2,0)circle(3pt)node[right]{$m$};
 \draw(0,2.2)arc(270:304:.8)node[below]{$\theta$};\node[right]at(3.5,1.5){小球在水平圆周上运动};
}{圆锥摆的悬点、绳长与倾角；受力、角动量和角速度方向在解答中分别确定。}{ch09:quality:conical-q}}
\newcommand{\ThQDumbbell}{\ThermalFigure{
 \draw[thick](-1,0)--(5,0);\draw[very thick](0,0)--(3,2);\fill(0,0)circle(3pt)node[below]{$A,m$};\fill(3,2)circle(3pt)node[right]{$B,m$};
 \draw[->,thick](1.5,2.4)--(1.5,1.2)node[left]{$v_0$};\draw(1,0)arc(0:33.7:1)node[right]{$\theta$};\node[above left]at(1.5,1){$2l$};
}{碰前哑铃作竖直平动，只有下端 $A$ 先接触地面。}{ch09:quality:dumbbell-q}}
\newcommand{\ThQDumbbells}{\ThermalFigure{
 \draw[very thick](0,0)--(2.5,0) (2.5,1.8)--(5,1.8);
 \foreach\x/\y/\lab in {0/0/A,2.5/0/B,2.5/1.8/C,5/1.8/D}{\fill(\x,\y)circle(3pt)node[above right]{$\lab$};}
 \draw[->,thick](1,-1)--(1,.8)node[left]{$v_0$};\draw[dashed](2.5,.15)--(2.5,1.65);
 \node[right]at(6,1){碰撞线 $BC$ 垂直两杆；上杆初始静止};
}{两杆沿水平方向，每端质量为 $m$；杆长均为 $2l$。}{ch09:quality:dumbbells-q}}
\newcommand{\ThQSeesaws}{\ThermalFigure{
 \foreach\s/\a/\b in {0/A/B,7/B/C}{\begin{scope}[shift={(\s,0)}]
 \draw[very thick](-2,0)--(2,0);\draw[thick](-.4,-.7)--(0,0)--(.4,-.7);
 \fill(-2,0)circle(3pt)node[above]{$\a$};\fill(2,0)circle(3pt)node[above]{$\b$};\end{scope}}
 \draw[->](3,0)--(4,0);\node[above]at(3.5,1.2){第一次冲击后，B等高落回另一装置};
}{两次等臂弹射的先后关系；$C$ 表示小明。}{ch09:quality:seesaws-q}}
\newcommand{\ThQConeBead}{\ThermalFigure{
 \draw[thick](-2,3.46)--(0,0)--(2,3.46);\draw[dashed](0,0)--(0,4);
 \fill[DeepBlue](1.5,2.598)circle(3pt)node[right]{$P$};\draw[dashed](0,2.598)--(1.5,2.598);
 \node[above]at(.7,2.6){$r$};\node[DeepBlue]at(2.4,2.6){$\otimes\ v_0$};\node[left]at(0,0){$O$};
}{正三角形轴截面，初速度垂直该截面；力和有效势的分析在解答中展开。}{ch09:quality:bead-q}}
\newcommand{\ThQBrakingWheel}{\ThermalFigure{
 \draw[thick](0,0)circle(1.6);\fill(0,0)circle(2pt)node[below]{$O$};\fill(1,0)circle(2pt);
 \draw[->,thick](1,0)--(1,1.1)node[right]{$f$};\draw[->](140:1.9)arc(140:20:1.9)node[right]{$\omega_0$};
 \draw[<->](0,-.5)--node[below]{$r$}(1,-.5);\node[right]at(3,.3){均匀砂轮 $m,R$；固定轴};
}{力作用于圆心右侧且向上，初始转向为顺时针。}{ch09:quality:braking-q}}
\newcommand{\ThQWheels}{\ThermalFigure{
 \draw[thick](0,0)circle(1.5);\draw[thick](2.7,0)circle(1.2);
 \node[align=center]at(0,0){$m_1$\\$R_1$};\node[align=center]at(2.7,0){$m_2$\\$R_2$};
 \draw[->](150:1.8)arc(150:35:1.8);\draw[->]($(2.7,0)+(30:1.5)$)arc(30:145:1.5);
 \node[above]at(-.2,2.1){$\omega_{10}$};\node[above]at(3.2,1.7){$\omega_{20}$};
 \node[right]at(5,0){仅相对滑动时摩擦力大小为 $f$};
}{两固定轴砂轮的初始转向；接触处两个表面均向下运动，滑动方向由速度差决定。}{ch09:quality:wheels-q}}
\newcommand{\ThQBilliard}{\ThermalFigure{
 \draw[thick](-2,0)--(4,0);\draw[thick](0,1)circle(1);\draw[->,thick](0,1)--(2.5,1)node[right]{$v_0$};
 \node[above]at(0,2.4){初态 $\omega=0$};\node[below]at(1,-.4){粗糙水平面 $\mu$};
}{匀质球起初只有平动，自转由后续摩擦建立。}{ch09:quality:billiard-q}}
\newcommand{\ThQBilliards}{\ThermalFigure{
 \draw[thick](-2,0)--(6,0);\draw[thick](0,1)circle(1) (3,1)circle(1);
 \draw[->,thick](0,1)--(1.6,1)node[right]{$v_0$};\draw[->]($(0,1)+(140:1.2)$)arc(140:20:1.2);
 \node[above]at(0,2.5){白球 $\omega_0$};\node[above]at(3,2.5){红球静止};\node[below]at(2,-.5){球间无摩擦，地面粗糙};
}{碰前白球向右纯滚动，球间冲量沿两球心连线。}{ch09:quality:billiards-q}}
\newcommand{\ThQEggs}{\ThermalFigure{
 \foreach\s/\lab in {0/生,6/熟}{\begin{scope}[shift={(\s,0)}]
 \draw[thick](-1,2)--(4,0);\begin{scope}[shift={(1.4,2.12)},rotate=-21.8]\draw[thick](0,0)ellipse(.7 and 1);\node at(0,0){\lab};\end{scope}\end{scope}}
 \node[above]at(4.5,3.8){两次测试：同一坡角、相同释放高度};
}{两个面板分别表示生、熟鸡蛋的同条件测试，不把初始位置差异混入比较。}{ch09:quality:eggs-q}}
\newcommand{\ThQPulley}{\ThermalFigure{
 \draw[thick](0,2.5)circle(.8);\draw[thick](-.8,2.5)arc(180:0:.8)--(.8,.5) (-.8,2.5)--(-.8,0);
 \draw[thick](-1.15,-.7)rectangle(-.45,0);\draw[thick](.45,-.2)rectangle(1.15,.5);
 \node at(-.8,-.35){$m_1$};\node at(.8,.15){$m_2$};\fill(0,2.5)circle(2pt)node[right]{$O$};
 \draw[thick](0,2.5)--(0,4);\node[right]at(2,2.5){滑轮 $M,R$，$m_1>m_2$};
}{绳沿轮缘绕过固定轴滑轮，两端连接重物；无相对滑动。}{ch09:quality:pulley-q}}
\newcommand{\ThQRodWork}{\ThermalFigure{
 \draw[dashed](0,2.5)--(5,2.5);\draw[very thick](0,2.5)--(4,.2);
 \fill(0,2.5)circle(2pt)node[left]{$O$};\fill(2,1.35)circle(3pt)node[below left]{$2,m$};\fill(4,.2)circle(3pt)node[right]{$1,m$};
 \node[above left]at(1,1.925){$l$};\node[above left]at(3,.775){$l$};\draw(1.2,2.5)arc(0:-29.9:1.2)node[right]{$\theta$};
}{水平初态用虚线表示，当前杆从水平线向下转过 $\theta$。}{ch09:quality:rod-work-q}}
\newcommand{\ThQPercussion}{\ThermalFigure{
 \draw[very thick](0,3.5)--(0,0);\fill(0,3.5)circle(2pt)node[above]{$O$};\fill(0,1)circle(2pt)node[right]{$P$};
 \draw[->,thick](-2,1)--(0,1)node[midway,above]{$J$};\draw[<->](.8,0)--node[right]{$a$}(.8,1);
 \draw[<->](1.8,0)--node[right]{$L$}(1.8,3.5);\node[right]at(3,1.5){悬点冲量为零的打击位置待求};
}{打击点 $P$ 距下端 $a$，距铰点则为 $L-a$。}{ch09:quality:percussion-q}}
\newcommand{\ThQChimney}{\ThermalFigure{
 \draw[very thick](0,0)--(2.4,3.84);\draw[dashed](0,0)--(0,4.5);\fill(0,0)circle(2pt)node[below]{$O$};
 \draw(0,1)arc(90:58:1)node[above]{$\theta$};\node[above left]at(1.2,1.92){$L,m$};
 \draw[dashed](.55,1.5)--(1.1,1.15);\node[right]at(3,2){比较不同截面的剪力与弯矩};
}{烟囱按绕底端转动的均匀杆处理，倾角从铅垂方向量起。}{ch09:quality:chimney-q}}



{
\graphicspath{{../../../sample/xes-original/sixteenth-section/}{../../../sample/xes-original/sixteenth-section/assets/}}
\section{角动量、力矩与参考点}
% Chapter-local deterministic geometry. All text nodes have transparent fill.
\newcommand{\crossProductDiagram}{%
\begin{center}\begin{minipage}{\linewidth}\centering\begin{tikzpicture}[>=Latex,font=\small]
 \coordinate(O)at(0,0);\coordinate(A)at(3.5,0);\coordinate(B)at(1.2,1.8);
 \fill[PaleBlue](O)--(A)--($(A)+(B)$)--(B)--cycle;
 \draw[TextGray](A)--($(A)+(B)$)--(B);
 \draw[vector](O)--(A)node[below]{$\vect a$};
 \draw[vector](O)--(B)node[above left]{$\vect b$};
 \draw pic[draw,"$\theta$",angle radius=8mm,angle eccentricity=1.3]{angle=A--O--B};
 \node[DeepBlue,font=\Large]at(2.5,.8){$\odot$};
 \node[above]at(2.5,1.15){$\vect a\times\vect b$};
 \node[anchor=west,align=left]at(5.3,1){四指：$\vect a\longrightarrow\vect b$\\[5pt]拇指：出纸面};
\end{tikzpicture}\figcap{外积的面积与右手方向}\label{ch09:fig:cross}\end{minipage}\end{center}}

\newcommand{\particleAngularMomentumDiagram}{%
\begin{center}\begin{minipage}{\linewidth}\centering\begin{tikzpicture}[>=Latex,font=\small]
 \coordinate(O)at(0,0);\coordinate(P)at(3,1.4);\coordinate(V)at(5.3,1.4);
 \coordinate(R)at($(P)+(1.8,.84)$);
 \draw[vector](O)--(P)node[pos=.43,above left]{$\vect r$};
 \draw[vector](P)--(V)node[right]{$\vect p=m\vect v$};
 \draw[guide](P)--(R)node[above]{$\vect r$ 的平移线};
 \draw pic[draw,"$\varphi$",angle radius=12mm,angle eccentricity=1.3]{angle=V--P--R};
 \fill(O)circle(1.5pt)node[below]{$O$};\fill(P)circle(2pt)node[below=4pt]{$P$};
 \node[DeepBlue,font=\Large]at(.5,1.3){$\otimes$};\node[left]at(.2,1.3){$\vect L_O$};
 \node[anchor=west,align=left]at(7,1.2){$L_z=-yp_x<0$\\[5pt]角动量入纸面};
\end{tikzpicture}\figcap{夹角从同向平移后的两矢量之间量取}\label{ch09:fig:particle}\end{minipage}\end{center}}

\newcommand{\apsidesDiagram}{%
\begin{center}\begin{minipage}{\linewidth}\centering\begin{tikzpicture}[>=Latex,font=\small]
 \def\aa{3.4}\def\bb{2.3}\pgfmathsetmacro{\cc}{sqrt(\aa^2-\bb^2)}
 \coordinate(F)at(\cc,0);\coordinate(P)at(\aa,0);\coordinate(A)at(-\aa,0);
 \draw[curve](0,0)ellipse[x radius=\aa,y radius=\bb];
 \fill(F)circle(2.5pt)node[above left=3pt]{$M$};
 \fill(P)circle(2pt)node[right=3pt]{近点};\fill(A)circle(2pt)node[left=3pt]{远点};
 \draw[vector](P)--++(0,1.4)node[above]{$\vect v_{\rm p}$};
 \draw[vector](A)--++(0,-.65)node[left=4pt]{$\vect v_{\rm a}$};
 \draw[guide](F)--(\cc,-3.5)(P)--(\aa,-2.85)(A)--(-\aa,-3.5);
 \draw[<->](\cc,-2.7)--(\aa,-2.7)node[midway,below=3pt]{$r_{\rm p}$};
 \draw[<->](-\aa,-3.35)--(\cc,-3.35)node[midway,below]{$r_{\rm a}$};
 \node[anchor=west,align=left]at(4.7,-.2){两点均有 $\vect r\perp\vect v$\\[6pt]$r_{\rm p}v_{\rm p}=r_{\rm a}v_{\rm a}$\\[6pt]速度箭头仅表示方向};
\end{tikzpicture}\figcap{同一椭圆轨道的两个回归点}\label{ch09:fig:apsides}\end{minipage}\end{center}}

\newcommand{\keplerContextDiagram}{%
\begin{center}\begin{minipage}{\linewidth}\centering\begin{tikzpicture}[>=Latex,font=\small]
 \def\aa{4}\def\bb{2.4}\pgfmathsetmacro{\cc}{sqrt(\aa^2-\bb^2)}
 \coordinate(F)at(\cc,0);\coordinate(P)at(\aa,0);\coordinate(A)at(-\aa,0);
 \draw[curve](0,0)ellipse[x radius=\aa,y radius=\bb];
 \fill[DeepBlue](F)circle(3pt)node[above left=3pt]{太阳};
 \fill(P)circle(2pt)node[right=3pt]{近日点};\fill(A)circle(2pt)node[left=3pt]{远日点};
 \draw[vector](P)--++(0,1.35)node[above]{$v_{\rm near}$};
 \draw[vector](A)--++(0,-.65)node[left=4pt]{$v_{\rm far}$};
 \draw[guide](F)--(\cc,-3.55)(P)--(\aa,-2.9)(A)--(-\aa,-3.55);
 \draw[<->](\cc,-2.8)--(\aa,-2.8)node[midway,below=3pt]{$r_{\rm near}$};
 \draw[<->](-\aa,-3.4)--(\cc,-3.4)node[midway,below]{$r_{\rm far}$};
 \node[align=center]at(0,.5){$r_{\rm near}=a-c$\\[7pt]$r_{\rm far}=a+c$};
\end{tikzpicture}\figcap{日心距离均从同一个焦点量起；速度箭头不按速率比例绘制}\label{ch09:fig:kepler}\end{minipage}\end{center}}

\newcommand{\momentsContextDiagram}{%
\begin{center}\begin{minipage}{\linewidth}\centering\begin{tikzpicture}[>=Latex,font=\small]
 \foreach \xx/\mode in {0/0,4/1,8/2}{
  \begin{scope}[shift={(\xx,0)}]
   \ifnum\mode=2\fill[PaleBlue](-.8,-.35)rectangle(.8,.65);\fi
   \draw[thick](-.8,-.35)--(-.8,.65) (.8,-.35)--(.8,.65);
   \draw[thick](0,.65)ellipse(.8 and .25);
   \draw[thick](-.8,-.35)arc[start angle=180,end angle=360,x radius=.8,y radius=.25];
   \draw[guide](-.8,-.35)arc[start angle=180,end angle=0,x radius=.8,y radius=.25];
   \ifnum\mode=1\draw[thick](0,.65)ellipse(.4 and .125);\fi
   \draw[axis](0,-.7)--(0,1.3);
   \ifnum\mode=1
    \draw[->](0,.65)--(.8,.65)node[right]{$R_2$};
    \draw[->](0,.65)--(-.4,.65);
    \draw[TextGray](-.4,.65)--(-.85,1.05)node[above left]{$R_1$};
   \else\draw[->](0,.65)--(.8,.65)node[right]{$R$};\fi
   \ifnum\mode=0\node at(0,-1.15){薄壁圆筒};\node at(0,-1.85){$I=MR^2$};\fi
   \ifnum\mode=1
    \node at(0,-1.15){厚壁圆筒};\node at(0,-1.85){$I=\dfrac{M(R_1^2+R_2^2)}2$};\fi
   \ifnum\mode=2\node at(0,-1.15){实心圆柱};\node at(0,-1.85){$I=\dfrac12MR^2$};\fi
  \end{scope}}
 \begin{scope}[shift={(12,0)}]
  \draw[curve](-1,.15)--(1,.15);\draw[axis](0,-.35)--(0,1.2);
  \draw[guide](-1,.1)--(-1,-.65)(1,.1)--(1,-.65);
  \draw[<->](-1,-.55)--(1,-.55)node[midway,below]{$L$};
  \node at(0,-1.15){细杆绕中心};\node at(0,-1.85){$I=\dfrac1{12}ML^2$};
 \end{scope}
 \begin{scope}[shift={(0,-4.8)}]
  \draw[curve](-1,.15)--(1,.15);\draw[axis](-1,-.3)--(-1,1.2);
  \draw[guide](-1,.1)--(-1,-.65)(1,.1)--(1,-.65);
  \draw[<->](-1,-.55)--(1,-.55)node[midway,below]{$L$};
  \node at(0,-1.25){细杆绕端点};\node at(0,-2){$I=\dfrac13ML^2$};
 \end{scope}
 \begin{scope}[shift={(4,-4.8)}]
  \draw[fill=PaleBlue,thick](-1,-.1)--(.55,-.1)--(1,.65)--(-.55,.65)--cycle;
  \draw[axis](0,-.5)--(0,1.35);
  \node[below]at(-.25,-.2){$a$};\node[right]at(.85,.25){$b$};
  \node at(0,-1.25){矩形薄板，轴垂直板面};\node at(0,-2){$I=\dfrac{M(a^2+b^2)}{12}$};
 \end{scope}
 \foreach \xx/\mode in {8/0,12/1}{
  \begin{scope}[shift={(\xx,-4.8)}]
   \ifnum\mode=0\fill[PaleBlue](0,.25)circle(.72);\fi
   \draw[thick](0,.25)circle(.72);\draw[guide](0,.25)ellipse(.72 and .20);
   \draw[axis](0,-.7)--(0,1.4);\draw[->](0,.25)--(.72,.25)node[right]{$R$};
   \ifnum\mode=0\node at(0,-1.25){均匀实心球};\node at(0,-2){$I=\dfrac25MR^2$};
   \else\node at(0,-1.25){均匀薄球壳};\node at(0,-2){$I=\dfrac23MR^2$};\fi
  \end{scope}}
\end{tikzpicture}\figcap{八种均匀刚体及其指定转轴；圆柱、圆筒均绕中心轴}\label{ch09:fig:inertia}\end{minipage}\end{center}}

\newcommand{\gyroDiagram}{%
\begin{center}\begin{minipage}{\linewidth}\centering\begin{tikzpicture}[>=Latex,font=\small]
 \coordinate(O)at(0,0);\coordinate(C)at(1.7,2.4);\coordinate(Z)at(0,3.6);
 \draw[thick](-.7,0)--(.7,0);\draw[guide](O)--(Z)node[above]{$z$};
 \begin{scope}[shift={(C)},rotate=54.6888]
  \draw[fill=PaleBlue,thick](0,0)ellipse(.22 and .8);
 \end{scope}
 \draw[curve](O)--($(O)!1.26!(C)$);
 \draw[vector]($(O)!.85!(C)$)--($(O)!1.45!(C)$)node[above]{$\vect L$};
 \draw[vector](C)--++(0,-1.15)node[right=3pt]{$M\vect g$};
 \draw[vector](O)--++(-.28,.9)node[left]{$\vect n$};
 \fill(O)circle(2pt)node[below]{$O$};\fill(C)circle(1.5pt);\draw[thin,TextGray](C)--(2.55,2.75)node[above right]{$C$};
 \node[right]at(1.12,1.2){$r$};
 \draw pic[draw,"$\theta$",angle radius=10mm,angle eccentricity=1.3]{angle=C--O--Z};
 \node[DeepBlue,font=\Large]at(-1.25,1.6){$\otimes$};\node[above]at(-1.25,1.9){$\vect\tau$};
 \node[below,align=center]at(1,-.4){侧视：$\vect r$、重力在纸面内\\力矩入纸面};
 \begin{scope}[shift={(7,0)}]
  \coordinate(K)at(0,0);\coordinate(I)at({1.6*cos(35)},{2.7+.42*sin(35)});
  \coordinate(J)at({1.6*cos(70)},{2.7+.42*sin(70)});
  \draw[guide](K)--(0,3.4);\draw[TextGray](K)--(-1.6,2.7)(K)--(1.6,2.7);
  \draw[guide](0,2.7)ellipse(1.6 and .42);
  \draw[vector](K)--(I)node[pos=.55,right=3pt]{$\vect L_i$};
  \draw[vector](K)--(J)node[pos=.65,left=3pt]{$\vect L_f$};
  \draw[vector](I)--(J)node[midway,above=4pt]{$\Delta\vect L$};
  \draw[->,DeepBlue](1.85,2.7)arc[start angle=0,end angle=125,x radius=1.85,y radius=.55];
  \node[above]at(-1.2,3.3){$\omega_p$};
  \node[below,align=center]at(0,-.4){矢量端点的轨迹\\$\Delta\vect L=\vect L_f-\vect L_i$};
 \end{scope}
\end{tikzpicture}\figcap{陀螺受力与进动时的角动量矢量差}\label{ch09:fig:gyro}\end{minipage}\end{center}}

\newcommand{\areaVelocityDiagram}{%
\begin{center}\begin{minipage}{\linewidth}\centering\begin{tikzpicture}[>=Latex,font=\small]
 \coordinate(O)at(0,0);\coordinate(A)at(4,.25);\coordinate(B)at(3.7,1.65);
 \fill[PaleBlue](O)--(A)--(B)--cycle;
 \draw[vector](O)--(A)node[midway,below=3pt]{$\vect r$};
 \draw[vector](O)--(B)node[midway,above=4pt]{$\vect r+\dd\vect r$};
 \draw[vector](A)--(B)node[midway,right=7pt]{$\dd\vect r=\vect v\dd t$};
 \fill(O)circle(1.5pt)node[left]{$O$};\node[DeepBlue,font=\Large]at(2.7,.7){$\odot$};
 \node[anchor=west,align=left]at(6,.8){$\dd\vect S$ 沿出纸面方向\\[5pt]$\dd S=\tfrac12|\vect r\times\dd\vect r|$};
\end{tikzpicture}\figcap{面积微元与法向；小位移为便于识别而放大}\label{ch09:fig:area}\end{minipage}\end{center}}

\newcommand{\sphereBandDiagram}{%
\begin{center}\begin{minipage}{\linewidth}\centering\begin{tikzpicture}[>=Latex,font=\small]
 \def\rr{1.8}\def\th{55}
 \coordinate(O)at(0,0);\coordinate(P)at({\rr*sin(\th)},{\rr*cos(\th)});\coordinate(Z)at(0,2.8);
 \draw[thick](O)circle(\rr);\draw[axis](0,-2.1)--(Z)node[above]{$z$};
 \draw[DeepBlue]({-\rr*sin(52)},{\rr*cos(52)})--({\rr*sin(52)},{\rr*cos(52)});
 \draw[DeepBlue]({-\rr*sin(59)},{\rr*cos(59)})--({\rr*sin(59)},{\rr*cos(59)});
 \draw[guide](O)--(P)node[midway,below right]{$R$};
 \draw[guide](P)--({\rr*sin(\th)},2.3);
 \draw[<->](0,2.2)--({\rr*sin(\th)},2.2)node[midway,above=3pt]{$R\sin\theta$};
 \draw pic[draw,"$\theta$",angle radius=5.5mm,angle eccentricity=1.25]{angle=P--O--Z};
 \fill(P)circle(1.5pt)node[below right=8pt]{$P$};\node[below left]at(O){$O$};
 \draw[DeepBlue,line width=2pt]({\rr*sin(52)},{\rr*cos(52)})arc[start angle=38,end angle=31,radius=\rr];
 \draw[TextGray](P)--(2.4,1.55)node[right]{$R\dd\theta$（沿球面）};
 \node[anchor=west,align=left]at(3.3,-.3){球带面积：\\[5pt]$\dd A=(2\pi R\sin\theta)(R\dd\theta)$\\[8pt]球带质量：$\dd m=\sigma\dd A$};
\end{tikzpicture}\figcap{窄球带的轴截面；球带绕z轴一周，周长为$2\pi R\sin\theta$}\label{ch09:fig:sphere-band}\end{minipage}\end{center}}

\newcommand{\conicalPendulumDiagram}{%
\begin{center}\begin{minipage}{\linewidth}\centering\begin{tikzpicture}[>=Latex,font=\small]
 \coordinate(O)at(0,3);\coordinate(P)at(2,0);\coordinate(V)at(0,0);
 \draw[thick](-.7,3)--(.7,3);\draw[guide](O)--(V);
 \draw[curve](O)--(P)node[pos=.55,left=4pt]{$l$};
 \draw[vector](P)--($(P)!.4!(O)$)node[pos=.5,right=3pt]{$\vect T$};
 \draw[vector](P)--++(0,-1.2)node[right]{$m\vect g$};
 \fill(O)circle(1.5pt)node[above left=4pt]{$O$};\fill(P)circle(2pt)node[below left]{$P$};
 \draw[vector](O)--++(1.5,1)node[right]{$\vect L_O$};
 \draw[vector](-1,3.1)--++(0,1)node[above]{$\vect\omega$};
 \draw pic[draw,"$\theta$",angle radius=9mm,angle eccentricity=1.4]{angle=V--O--P};
 \draw[<->](0,-.4)--(2,-.4)node[midway,below]{$r=l\sin\theta$};
 \node[DeepBlue,font=\Large]at(2.65,.1){$\otimes$};\node[right]at(3,.1){$\vect v$};
 \node[anchor=west,align=left]at(4.4,2){$+x$ 向右，$+z$ 向上\\[5pt]$+y$ 入纸面\\[8pt]$\vect L_O\perp OP$};
\end{tikzpicture}\figcap{圆锥摆的同一瞬时位置、速度、受力与角动量}\label{ch09:fig:conical}\end{minipage}\end{center}}

\newcommand{\dumbbellGroundDiagram}{%
\begin{center}\begin{minipage}{\linewidth}\centering\begin{tikzpicture}[>=Latex,font=\small]
 \coordinate(A)at(0,0);\coordinate(B)at(3,2);\coordinate(H)at(4,0);
 \draw[thick](-1,0)--(4.5,0);\foreach\x in{-.8,-.4,...,4.4}{\draw[TextGray](\x,0)--++(-.13,-.15);}
 \draw[curve](A)--(B)node[midway,above left=3pt]{$2l$};
 \fill(A)circle(2.5pt)node[below=8pt]{$A,m$};\fill(B)circle(2.5pt)node[right=4pt]{$B,m$};
 \draw[vector](B)--++(0,-1.2)node[right]{$\vect v_0$};
 \draw[vector](-.6,1.8)--(-.6,.6)node[left]{$\vect v_0$};
 \draw pic[draw,"$\theta$",angle radius=11mm,angle eccentricity=1.3]{angle=H--A--B};
 \node[anchor=west,align=left]at(5,1.1){碰前：整体竖直平动\\[6pt]碰后：A粘住固定};
\end{tikzpicture}\figcap{哑铃下端A与地面的短时粘着碰撞}\label{ch09:fig:dumbbell-ground}\end{minipage}\end{center}}

\newcommand{\dumbbellCollisionDiagram}{%
\begin{center}\begin{minipage}{\linewidth}\centering\begin{tikzpicture}[>=Latex,font=\small]
 \draw[curve](0,0)--(2.5,0)(2.5,1.7)--(5,1.7);
 \foreach\x/\y/\n in{0/0/A,2.5/0/B,2.5/1.7/C,5/1.7/D}{\fill(\x,\y)circle(2.5pt)node[above right=3pt]{$\n$};}
 \draw[guide](2.5,.12)--(2.5,1.58);
 \draw[vector](1,-.5)--(1,.8)node[left]{$v_0$};
 \node[below]at(1.85,-.15){$2l$};\node[below]at(4,1.6){$2l$};
 \node[anchor=west,align=left]at(6,.7){碰撞前：B向C接近\\[5pt]BC垂直于两杆\\[5pt]每个端点质量均为m};
\end{tikzpicture}\figcap{两哑铃的碰前位置与碰撞线}\label{ch09:fig:dumbbell-collision}\end{minipage}\end{center}}

\newcommand{\seesawDiagram}{%
\begin{center}\begin{minipage}{\linewidth}\centering\begin{tikzpicture}[>=Latex,font=\small]
 \foreach\s/\leftname/\rightname/\vl/\vr in{0/A/B/u_A/v_B,7/B/C/v_B/v_C}{
  \begin{scope}[xshift=\s cm]
   \draw[curve](-2,0)--(2,0);\draw[thick](-.4,-.6)--(0,0)--(.4,-.6)--cycle;
   \fill(-2,0)circle(2.4pt)node[below=4pt]{$m_{\leftname}$};
   \fill(2,0)circle(2.4pt)node[below=4pt]{$m_{\rightname}$};
   \draw[vector](-2,1.9)--(-2,.45)node[midway,left]{$\vl$};
   \draw[vector](2,.15)--(2,1.6)node[midway,right]{$\vr$};
   \draw[guide](-2,-.9)--(2,-.9);\node[below]at(0,-.9){等臂、近水平的冲击位置};
  \end{scope}}
 \draw[->,TextGray](3,1)--(4,1);\node[above,align=center]at(3.5,1.1){随后\\等高落回};
\end{tikzpicture}\figcap{两次理想弹射的过程图；中间箭头不是空间飞行轨迹}\label{ch09:fig:seesaw}\end{minipage}\end{center}}

\newcommand{\funnelDiagram}{%
\begin{center}\begin{minipage}{\linewidth}\centering\begin{tikzpicture}[>=Latex,font=\small]
 \coordinate(O)at(0,0);\coordinate(V)at(0,4.1);
 \coordinate(P)at({3/sqrt(3)},3);
 \draw[curve]({-3.8/sqrt(3)},3.8)--(O)--({3.8/sqrt(3)},3.8);
 \draw[guide](O)--(V);\fill(O)circle(1.5pt)node[below]{$O$};
 \fill(P)circle(2.4pt)node[right=4pt]{$P,m$};
 \draw[vector](P)--++({-1.15*cos(30)},{1.15*sin(30)})node[left=3pt]{$\vect N$};
 \draw[vector](P)--++(0,-1.1)node[right=3pt]{$m\vect g$};
 \draw[guide](0,3)--(P);
 \draw[<->](0,2.65)--({3/sqrt(3)},2.65)node[midway,below=3pt]{$r$};
 \draw[guide](O)--(-2.8,0)(0,3)--(-2.8,3);
 \draw[<->](-2.65,0)--(-2.65,3)node[midway,left]{$h$};
 \draw pic[draw,"$\theta$",angle radius=10mm,angle eccentricity=1.3]{angle=P--O--V};
 \node[DeepBlue,font=\Large]at(2.7,3.3){$\otimes$};\node[right]at(3.05,3.3){$\vect v_\varphi$};
 \node[anchor=west,align=left]at(3.9,1.6){轴截面内不画切向速度箭头\\[6pt]$\theta=30^\circ$，$r=h\tan\theta$};
\end{tikzpicture}\figcap{圆锥漏斗轴截面、接触受力与入纸面的圆周速度}\label{ch09:fig:funnel}\end{minipage}\end{center}}


在运动学中，角速度和角加速度描述转动怎样进行；在动力学中，还需要知道转动状态怎样受力而改变。
地球的自转、收拢双臂的滑冰者和高速旋转的陀螺，都涉及这一问题。相应的物理量是角动量，
改变角动量的是力矩。零外力矩保证角动量不变，却不保证任意刚体的角速度方向都不变。

\subsection{外积与转动方向}
内积把两个矢量联系成标量，适合描述投影、功和功率；外积又称叉积，适合描述有方向的转动。
对非零矢量 $\vect a,\vect b$，取它们同起点时的夹角 $0\le\theta\le\pi$，有
\begin{equation}
 |\vect a\times\vect b|=ab\sin\theta.
\end{equation}
其大小是以两矢量为邻边的平行四边形面积，方向垂直于该平面。右手四指沿小于或等于
$180^\circ$ 的转角从 $\vect a$ 弯向 $\vect b$，拇指所指就是外积方向。平行或反平行时
外积为零，不再给零矢量指定方向。图~\ref{ch09:fig:cross}中从向右的 $\vect a$ 转到右上的
$\vect b$ 为逆时针，故外积出纸面，以 $\odot$ 表示；$\otimes$ 表示入纸面。
\crossProductDiagram
外积满足反交换律和分配律，但一般不满足结合律：
\begin{align}
 \vect a\times\vect b&=-\vect b\times\vect a,\qquad \vect a\times\vect a=\vect0,\\
 (\vect a+\vect b)\times\vect c&=\vect a\times\vect c+\vect b\times\vect c,\\
 (\vect a\times\vect b)\times\vect c&=(\vect a\cdot\vect c)\vect b-(\vect b\cdot\vect c)\vect a,\\
 \vect a\times(\vect b\times\vect c)&=(\vect a\cdot\vect c)\vect b-(\vect a\cdot\vect b)\vect c.
\end{align}
在右手直角坐标系中，分量计算为
\begin{equation}
 \vect a\times\vect b=
 (a_yb_z-a_zb_y,\ a_zb_x-a_xb_z,\ a_xb_y-a_yb_x).
\end{equation}
例如纸面为 $xy$ 平面，向右、向上分别为正方向，$+z$ 就是出纸面。这个约定下
$(\vect r\times\vect p)_z=xp_y-yp_x$，正值对应逆时针趋势，负值对应顺时针趋势。
这是右手定则的代数形式，不能在一道题中途更换正方向。

位置、速度、动量属于极矢量，角速度、力矩和角动量属于轴矢量。以空间反演为例，
极矢量 $\vect r,\vect p$ 都变号，而 $(-\vect r)\times(-\vect p)=\vect r\times\vect p$ 不变。
因此外积的变换性质不同于普通位置矢量；例如轴矢量与极矢量的外积
$\vect v=\vect\omega\times\vect r$ 仍是极矢量。

\subsection{角动量、力矩与参考点}
先在一个惯性系内取固定点 $O$。质量恒定为 $m$ 的质点位置为 $\vect r=\overrightarrow{OP}$，
速度为 $\vect v$，动量为 $\vect p=m\vect v$。定义
\begin{equation}
 \vect L_O=\vect r\times\vect p,\qquad L_O=mrv\sin\varphi=p\,d,
\end{equation}
其中 $\varphi$ 是 $\vect r$ 与 $\vect v$ 的夹角，$d$ 是 $O$ 到动量作用线的垂距，
不是 $O$ 到质点的距离。图~\ref{ch09:fig:particle}把 $\vect r$ 平移到 $P$ 再量角，
两矢量的夹角为锐角；$p_y=0,y>0$，故 $L_z=-yp_x<0$，角动量入纸面。
\particleAngularMomentumDiagram
若质点沿过 $O$ 的直线径向运动，则 $d=0$、$L_O=0$；若以恒定速度沿不经过 $O$ 的直线运动，
$d$ 不变，角动量也不变。角动量必须附带参考点，同一动量对不同点的角动量一般不同。

由于 $O$ 固定，$\dd\vect r/\dd t=\vect v$。对定义求导，并利用牛顿第二定律：
\begin{align}
 \frac{\dd\vect L_O}{\dd t}
 &=\frac{\dd\vect r}{\dd t}\times m\vect v+
   \vect r\times\frac{\dd(m\vect v)}{\dd t}\\
 &=\vect v\times m\vect v+\vect r\times\vect F
 =\vect r\times\vect F.
\end{align}
第一项为零，因为两个矢量平行。定义对同一点的力矩 $\vect M_O=\vect r\times\vect F$，得
\begin{equation}
 \boxed{\frac{\dd\vect L_O}{\dd t}=\vect M_O}.\label{ch09:eq:angular-law}
\end{equation}
力矩大小是力乘以相应垂直力臂，平行于 $\vect r$ 的力分量不贡献力矩。合力非零，合力矩仍可为零。
角动量单位为 $\mathrm{kg\,m^2/s}$，力矩单位为 $\mathrm{N\,m}=\mathrm{kg\,m^2/s^2}$，
二者相差一个时间因子。积分形式为
\begin{equation}
 \vect L_O(t_2)-\vect L_O(t_1)=\int_{t_1}^{t_2}\vect M_O\,\dd t.
\end{equation}
右侧称为冲量矩。碰撞时外力瞬时很大不一定破坏近似守恒，但必须确认该外力的冲量矩可忽略，
或其作用线通过所取矩心；不能仅凭“时间短”就一概使用角动量守恒。

\subsection{质点系、质心分解及移动矩心}
对质量不变的质点系，设总质量为 $M$，总动量为 $\vect P$，质心位置、速度分别为
$\vect R_C,\vect V_C$。令 $\vect r_i=\vect R_C+\vect r_i'$、
$\vect v_i=\vect V_C+\vect v_i'$。质心定义给出
\begin{equation}
 \sum_i m_i\vect r_i'=\vect0,\qquad
 \sum_i m_i\vect v_i'=\frac{\dd}{\dd t}\sum_i m_i\vect r_i'=\vect0.
\end{equation}
逐项展开角动量求和，而不是把混合项略去：
\begin{align}
 \vect L_O&=\sum_i m_i(\vect R_C+\vect r_i')\times(\vect V_C+\vect v_i')\\
 &=M\vect R_C\times\vect V_C
 +\vect R_C\times\sum_i m_i\vect v_i'
 +\left(\sum_i m_i\vect r_i'\right)\times\vect V_C
 +\sum_i m_i\vect r_i'\times\vect v_i'\\
 &=\boxed{M\vect R_C\times\vect V_C+\vect L_C}.
\end{align}
第一项称轨道角动量，第二项是相对质心的内部角动量。以地球为例，对太阳取矩时，
公转与自转分别对应这两项。此分解不要求系统是刚体：粒子团、爆炸碎片和可形变物体同样适用。

设 $\vect F_{ij}$ 是 $j$ 对 $i$ 的内力，若 $\vect F_{ji}=-\vect F_{ij}$ 且内力沿两质点连线，
则一对内力的力矩和为
\begin{equation}
 \vect r_i\times\vect F_{ij}+\vect r_j\times\vect F_{ji}
 =(\vect r_i-\vect r_j)\times\vect F_{ij}=\vect0.
\end{equation}
因此总角动量的变化只由合外力矩决定。若其全部分量为零，则角动量矢量守恒；
若只有某一固定方向的力矩分量为零，则只能断定那个方向的角动量分量守恒。
轨道角动量的变化率是 $\vect R_C\times\vect F_{\rm ext}$，即把合外力作用于质心后
对 $O$ 的力矩；自转角动量的变化率是各外力对质心的合力矩。两项相加才是各外力对O的总力矩，
不能把总力矩全部算进轨道部分。

\textbf{换固定矩心。}若 $O,O'$ 都固定在同一惯性系中，$\vect a=\overrightarrow{OO'}$，则
\begin{align}
 \vect L_{O'}&=\vect L_O-\vect a\times\vect P,\\
 \vect M_{O'}&=\vect M_O-\vect a\times\vect F_{\rm ext}.
\end{align}
总动量为零时角动量与固定参考点选择无关；合外力为零时外力矩与固定参考点选择无关。
一般情形两者均随矩心改变。选支点是为消去支点力矩，并不意味着其他外力矩自动为零。

\textbf{矩心移动。}若 $O$ 以速度 $\vect v_O$ 移动，仍用惯性系中的绝对速度定义
$\vect L_O^{\rm abs}=\sum_i(\vect r_i-\vect R_O)\times m_i\vect v_i$，则
\begin{equation}
 \frac{\dd\vect L_O^{\rm abs}}{\dd t}
 =\sum_i(\vect v_i-\vect v_O)\times m_i\vect v_i+\vect M_O
 =\vect M_O-\vect v_O\times\vect P.
\end{equation}
所以即使矩心匀速运动，也不能不加条件地沿用固定点公式。
若改用相对速度定义 $\vect L_O^{\rm rel}$，并让参考轴只平移、不旋转，则
\begin{equation}
 \frac{\dd\vect L_O^{\rm rel}}{\dd t}
 =\vect M_O-\left[\sum_i m_i(\vect r_i-\vect R_O)\right]\times\vect a_O.
\end{equation}
两种公式对应不同的角动量定义，不应混用。特别取 $O=C$ 时，括号中的量为零。
在只随质心平移的坐标系里，也可把每个质点的惯性力写成 $-m_i\vect a_C$，其对质心的合力矩
\begin{equation}
 \sum_i\vect r_i'\times(-m_i\vect a_C)
 =-\left(\sum_i m_i\vect r_i'\right)\times\vect a_C=\vect0.
\end{equation}
这才是质心即使加速，仍有 $\dd\vect L_C/\dd t=\vect M_C^{\rm ext}$ 的原因。

\section{角动量守恒与有心力运动}
力始终沿质点与固定力心的连线，称为有心力。它对力心的力矩为零，角动量守恒；
若 $\vect L\ne\vect0$，$\vect r,\vect v$ 始终垂直于固定的 $\vect L$，轨迹处于一个固定平面。
$L=0$ 对应径向运动，应单独处理。平面极坐标中取沿运动的方位角 $\varphi$，有
\begin{equation}
 \vect v=\dot r\,\uvec e_r+r\dot\varphi\,\uvec e_\varphi,
 \qquad L=mr^2\dot\varphi.
\end{equation}
窄扇形面积为 $\dd A=\tfrac12r^2\dd\varphi$，故面积速度
\begin{equation}
 \frac{\dd A}{\dd t}=\frac12r^2\dot\varphi=\frac{L}{2m}.
\end{equation}
这对任何固定中心力成立，不只适用于平方反比引力；面积矢量的直接推导见后面的面积速度例题。

质点到力心的距离达到极值的位置称回归点，此时 $\dot r=0$，速度只有横向分量。
若有两个回归点 $r_1,r_2$，便有 $r_1v_1=r_2v_2$；回归点不是合力为零的点。
图~\ref{ch09:fig:apsides}中两处速度都与日心连线垂直，近点半径小，速率反而大。
\apsidesDiagram
若中心力又是保守力，势能为 $U(r)$，把速度分解代入机械能：
\begin{align}
 E&=\frac12m\dot r^2+\frac12mr^2\dot\varphi^2+U(r)\\
 &=\frac12m\dot r^2+U_{\rm eff}(r),
 \qquad U_{\rm eff}(r)=U(r)+\frac{L^2}{2mr^2}.
\end{align}
有效势中的第二项不是新增的真实相互作用，而是角动量约束在径向方程中的体现。
可达范围必须有 $E\ge U_{\rm eff}$，回归点满足等号。圆轨道还要求
$U'_{\rm eff}(r_0)=0$；若 $U''_{\rm eff}(r_0)>0$，小径向扰动稳定，负值则不稳定，
二阶导数为零时不能仅凭二阶判据下结论。能量给速度大小，角动量把横向速度与半径联系起来，二者不能互相替代。

以 $M\gg m$ 的万有引力为例，$U=-GMm/r$。在回归点代入 $v=L/(mr)$，令 $u=1/r$：
\begin{equation}
 \frac{L^2}{2m}u^2-GMm\,u-E=0,
 \quad
 u_\pm=\frac{GMm^2}{L^2}
 \left(1\pm\sqrt{1+\frac{2EL^2}{G^2M^2m^3}}\right).
\end{equation}
这里已假定 $L\ne0$。允许能量须满足判别式非负：
$E\ge E_{\min}=-G^2M^2m^3/(2L^2)$。
当 $E_{\min}<E<0$ 时，两根均正，对应椭圆的近、远日点；$E=E_{\min}$ 时重根对应圆轨道。
$E=0$ 时一个根为零，即远端趋向无穷远；$E>0$ 时仅一个根为正，只存在一个有限距离的回归点。
图~\ref{ch09:fig:kepler}专画束缚椭圆情形，太阳位于焦点而非中心，两个半径都从太阳量起。
\keplerContextDiagram

\subsection{中心力轨道方程与开普勒定律}
进一步把时间消去，可以直接求轨道形状。设 $h=L/m=r^2\dot\varphi\ne0$，令 $u=1/r$，
撇号在本小节表示对 $\varphi$ 求导。由 $\dot\varphi=hu^2$，
\begin{equation}
 \dot r=-\frac{u'}{u^2}\dot\varphi=-hu',\qquad
 \ddot r=-h u''\dot\varphi=-h^2u^2u'',\qquad
 r\dot\varphi^2=h^2u^3.
\end{equation}
把它们代入径向牛顿方程 $m(\ddot r-r\dot\varphi^2)=F(r)$，得到比内方程（Binet equation）
\begin{equation}
 u''+u=-\frac{F(1/u)}{mh^2u^2}.
\end{equation}
引力 $F=-GMmu^2$ 使右侧成为常量 $GM/h^2$，故
\begin{equation}
 u=\frac{GM}{h^2}+A\cos\varphi+B\sin\varphi
 =\frac1p[1+e\cos(\varphi-\varphi_0)],
 \quad p=\frac{h^2}{GM},\quad e=p\sqrt{A^2+B^2}.
\end{equation}
其中 $p$ 称半通径，$e\ge0$ 是离心率，$\varphi_0$ 给出近点方位；于是
$r=p/[1+e\cos(\varphi-\varphi_0)]$。$0\le e<1$ 是椭圆，$e=1$ 是抛物线，
$e>1$ 是双曲线，圆是 $e=0$ 的特殊情形。各类均须限制 $r>0$ 的方位范围。
椭圆半长轴 $a$、半短轴 $b$、焦距 $c$ 满足
$c=ae=\sqrt{a^2-b^2}$，$p=a(1-e^2)=b^2/a$。

对椭圆，$r_{\rm p}=a(1-e)$、$r_{\rm a}=a(1+e)$。
把近点半径及 $L^2=GMm^2p$ 代入能量，逐步化简得
\begin{align}
 E&=\frac{GMm\,a(1-e^2)}{2a^2(1-e)^2}-\frac{GMm}{a(1-e)}\\
 &=\frac{GMm}{2a(1-e)}[(1+e)-2]
 =-\frac{GMm}{2a}.
\end{align}
一周扫过面积为 $\pi ab$，面积速度为 $h/2$，所以
\begin{equation}
 T=\frac{2\pi ab}{h},\qquad
 T^2=\frac{4\pi^2a^2b^2}{GMp}=\frac{4\pi^2a^3}{GM}.
\end{equation}
这分别给出轨道能量与开普勒第三定律。$a$ 是半长轴，不是近点距离；同一中心天体周围的
椭圆轨道满足 $T^2\propto a^3$，与面积速度恒定一起把几何、时间与动力学联系起来。

\section{转动惯量、移轴定理与进动}
刚体各质元相对位置不变，绕固定轴转动时有共同角速度。对轴向角动量，有
\begin{equation}
 L_{\rm axis}=I_{\rm axis}\omega,\qquad
 I_{\rm axis}=\sum_i m_i r_{i\perp}^2=\int r_\perp^2\dd m.
\end{equation}
$r_\perp$ 是到轴的垂距。仅当该轴是相应的主惯量轴时，才可把整个角动量写成
$\vect L=I\vect\omega$；一般情况下还有垂直于该轴的角动量分量。
在三维中可写 $\vect L_C=\mathbf I_C\vect\omega$，$\mathbf I_C$ 称惯量张量，
其作用不一定只是乘一个数，也可能改变方向。后面用两质点刚体具体说明。

图~\ref{ch09:fig:inertia}列出八种均匀物体及指定转轴。薄圆筒的质量全在半径R处，故为 $MR^2$；
实心圆柱的一部分质量更靠近轴，故为 $MR^2/2$。形状、尺寸含义和轴位置必须同时核对。
\momentsContextDiagram
例如边长 $a,b$ 的矩形薄板，取过中心且垂直板面的轴。面密度为 $M/(ab)$，有
\begin{align}
 I&=\frac{M}{ab}\int_{-a/2}^{a/2}\int_{-b/2}^{b/2}(x^2+y^2)\,\dd y\,\dd x\\
 &=\frac{M}{ab}\left(b\frac{a^3}{12}+a\frac{b^3}{12}\right)
 =\frac{M(a^2+b^2)}{12}.
\end{align}
量纲为质量乘长度平方，$b\to0$ 时回到细杆过中心垂轴的结果。

设目标轴与过质心的轴平行，轴间距为d。质元到质心轴的横向矢量为 $\vect\rho_i$，
到目标轴的横向矢量为 $\vect\rho_i+\vect d$。展开平方得到平行轴定理：
\begin{align}
 I&=\sum_i m_i|\vect\rho_i+\vect d|^2\\
 &=I_C+2\vect d\cdot\sum_i m_i\vect\rho_i+Md^2
 =\boxed{I_C+Md^2}.
\end{align}
中间项由质心定义消失。这解释了同一细杆绕端点的惯量从 $ML^2/12$ 增加到 $ML^2/3$。
两轴必须平行；复合体各部分先移到同一轴再相加，挖孔则减去被移除部分关于同一轴的惯量。

若所有质元在同一薄平面内，x、y轴在板面内且互相垂直，z轴过同一点垂直板面，则
\begin{equation}
 I_z=\int(x^2+y^2)\dd m=\int x^2\dd m+\int y^2\dd m=I_y+I_x.
\end{equation}
这就是垂直轴定理。它不要求原点必须是质心，但要求三轴共点、两两垂直及物体为平面薄片，
不能直接套用于有厚度的三维物体。

\subsection{角动量守恒与陀螺进动}
系统总角动量守恒，不要求每个部分的角动量分别守恒。收拢双臂使转动惯量减小，
在外力矩可忽略且绕同一主轴的条件下，$I\omega=L$ 不变而转速增加。
此时 $E_k=L^2/(2I)$ 反而增加，能量来自肌肉做功；角动量守恒不等于机械能守恒。
图~\ref{ch09:fig:cosmic}的飓风与天体吸积结构说明旋转在不同尺度都很常见。
分析它们仍要检查外力矩、摩擦、角动量输运及物质交换，不能仅凭旋涡外形断言整个系统严格守恒。
\historyfigure[0.78\linewidth]{cosmic-angular-momentum-redraw.png}{飓风与天体吸积结构中的旋转\label{ch09:fig:cosmic}}

图~\ref{ch09:fig:gyro}中，对称陀螺固定支点为O，质心为C，$OC=r$，质量为M；
自转轴对竖直方向的倾角为$\theta$，高速自转角速度为 $\omega_0$。若自转角动量占主导，
$L\simeq I_0\omega_0$ 沿自转轴。重力对O的力矩大小 $Mgr\sin\theta$，垂直于
$\overrightarrow{OC}$ 与重力所在平面；支点反力通过O，不贡献对O的力矩。
\gyroDiagram
稳定进动时$\theta$近似不变，$\vect L$ 的端点沿水平圆周移动，圆半径为 $L\sin\theta$。
两根等长矢量 $\vect L_i,\vect L_f$ 的差连接其端点，若转过小方位角 $\dd\phi$，
\begin{equation}
 |\dd\vect L|=L\sin\theta\,\dd\phi,
 \qquad \left|\frac{\dd\vect L}{\dd t}\right|=L\sin\theta\,\omega_p.
\end{equation}
与重力矩相等，且 $\sin\theta\ne0$，得到
\begin{equation}
 Mgr\sin\theta=I_0\omega_0\sin\theta\,\omega_p,
 \qquad \boxed{\omega_p=\frac{Mgr}{I_0\omega_0}}.
\end{equation}
进动主要改变L的方向。上述式子忽略了进动本身对角动量的贡献，并要求自转足够快、倾角稳定、
支点固定；自转较慢时可能出现明显章动，即倾角随时间振荡。$\theta$=0时轴没有可辨认的绕竖直进动圆，
不能用约去零因子的步骤推断可观测的进动。式中r始终是支点到质心的距离，不是陀螺轮体半径。


\Needspace{8\baselineskip}
\begin{origexamplebox}{【统统算出来】}\label{ch09:ex:s16-14-1}
已知 $\overrightarrow{AB}=(1,1,1)$，$\overrightarrow{AC}=(1,1,0)$，求：
\begin{enumerate}[label=(\arabic*),leftmargin=2.2em]
  \item $\overrightarrow{AB}\times\overrightarrow{AC}$；
  \item $\overrightarrow{AB}\cdot\overrightarrow{AC}$；
  \item $\overrightarrow{AB}$ 与 $\overrightarrow{AC}$ 夹角的正弦；
  \item $\overrightarrow{AB}$ 与 $\overrightarrow{AC}$ 夹角的余弦；
  \item 验证 $\sin^2\theta+\cos^2\theta=1$。
\end{enumerate}

\ThQVectors
\end{origexamplebox}

\Needspace{6\baselineskip}\begin{xesolution}
\ThAVectors

记 $\vect a=(1,1,1)$、$\vect b=(1,1,0)$。直接按分量计算：
\[
 \vect a\times\vect b
 =\begin{vmatrix}
   \vect i&\vect j&\vect k\\
   1&1&1\\
   1&1&0
  \end{vmatrix}
 =\boxed{(-1,1,0)}.
\]
而
\[
 \vect a\cdot\vect b=1+1+0=\boxed2,
 \qquad
 |\vect a|=\sqrt3,\quad |\vect b|=\sqrt2.
\]
由内、外积的几何意义，
\[
 \sin\theta=\frac{|\vect a\times\vect b|}{|\vect a||\vect b|}
 =\frac{\sqrt2}{\sqrt6}=\boxed{\frac1{\sqrt3}},
\]
\[
 \cos\theta=\frac{\vect a\cdot\vect b}{|\vect a||\vect b|}
 =\frac2{\sqrt6}=\boxed{\sqrt{\frac23}}.
\]
最后
\[
 \sin^2\theta+\cos^2\theta=\frac13+\frac23=1,
\]
验证完毕。由于内积为正，夹角是锐角；因此这里的正弦、余弦均取正值。
\end{xesolution}

\Needspace{8\baselineskip}
\begin{origexamplebox}{【面积速度】}\label{ch09:ex:s16-14-2}
学习角动量之前，小明定义了“面积速度”：在参考点看来，一个质点与参考点的连线在运动
过程中会扫过一定面积。按照右手螺旋方向把面积定义为面积矢量 $\vect S$，面积速度矢量即
单位时间内扫过的面积矢量。试说明面积速度矢量与角动量的关系。


\ThQArea
\end{origexamplebox}

\Needspace{6\baselineskip}\begin{xesolution}
\areaVelocityDiagram

在短时间 $\dd t$ 内，质点位移为 $\dd\vect r=\vect v\dd t$。
图~\ref{ch09:fig:area}中窄三角形的有向面积为
\[
 \dd\vect S=\frac12\vect r\times\dd\vect r
 =\frac12\vect r\times\vect v\,\dd t.
\]
所以面积速度矢量
\[
 \boxed{\frac{\dd\vect S}{\dd t}
 =\frac12\vect r\times\vect v
 =\frac{\vect L}{2m}}.
\]
面积速度与角动量方向相同，大小相差常数 $2m$。于是，有心力作用下角动量守恒等价于
面积速度恒定，这正是开普勒第二定律的几何本质。
\end{xesolution}

\Needspace{8\baselineskip}
\begin{origexamplebox}{【继续算出来】}\label{ch09:ex:s16-14-3}
小明研究刚体的角动量，需要先计算转动惯量。以下刚体质量均匀分布，质量均为 $m$。求：
\begin{enumerate}[label=(\arabic*),leftmargin=2.2em]
  \item 长为 $l$ 的细杆绕其中垂线的转动惯量；
  \item 半径为 $R$ 的圆盘绕一条直径的转动惯量；
  \item 半径为 $R$ 的空心球壳绕一条直径的转动惯量；
  \item 半径为 $R$ 的实心球体绕一条直径的转动惯量。
\end{enumerate}

\ThQInertias
\end{origexamplebox}

\Needspace{6\baselineskip}\begin{xesolution}
四问都从 $I=\int r_\perp^2\dd m$ 出发；$r_\perp$ 必须是到所求转轴的垂距，而不一定是到原点的距离。

\textbf{(1) 细杆。}\quad 取杆中点为原点，线密度 $\lambda=m/l$：
\[
 I=\lambda\int_{-l/2}^{l/2}x^2\dd x
 =\frac ml\left[\frac{x^3}{3}\right]_{-l/2}^{l/2}
 =\frac ml\frac{l^3}{12}=\boxed{\frac1{12}ml^2}.
\]

\textbf{(2) 圆盘。}\quad 先算垂直盘面过圆心轴的转动惯量。面密度
$\sigma=m/(\pi R^2)$，半径为 $r$ 的圆环质量 $\dd m=2\pi r\sigma\dd r$，故
\[
 I_z=\int_0^Rr^2(2\pi r\sigma\dd r)=\frac12mR^2.
\]
对平面薄片使用垂直轴定理 $I_z=I_x+I_y$，又由对称性 $I_x=I_y$，所以绕直径
\[
 \boxed{I_x=\frac14mR^2}.
\]

\textbf{(3) 空心球壳。}\quad 这里指厚度可忽略的均匀薄球壳。取z为所求直径轴，
极角$\theta$从正z方向量起。图~\ref{ch09:fig:sphere-band}中球带的环半径为 $R\sin\theta$，
沿球面的带宽为 $R\dd\theta$，二者相乘再乘 $2\pi$ 得球带面积；它不是投影到水平面后的面积。
\sphereBandDiagram
面密度 $\sigma=m/(4\pi R^2)$，因而
\[
 \dd m=2\pi R^2\sigma\sin\theta\dd\theta,
 \qquad r_\perp=R\sin\theta.
\]
于是
\[
 I=2\pi\sigma R^4\int_0^\pi\sin^3\theta\dd\theta.
\]
令 $u=\cos\theta$，$\dd u=-\sin\theta\dd\theta$，上下限由1变为−1：
\[
 \int_0^\pi\sin^3\theta\dd\theta
 =\int_{-1}^{1}(1-u^2)\dd u
 =\left[u-\frac{u^3}{3}\right]_{-1}^{1}=\frac43.
\]
代入面密度，$I=2\pi[m/(4\pi R^2)]R^4(4/3)=\boxed{2mR^2/3}$。

\textbf{(4) 实心球。}\quad 把球体看成半径为 $r$、厚度 $\dd r$ 的球壳。设体密度
$\rho=3m/(4\pi R^3)$，球壳质量为 $\dd m=4\pi r^2\rho\dd r$，其转动惯量为
$\dd I=(2/3)r^2\dd m$，故
\[
 I=\frac{8\pi\rho}{3}\int_0^Rr^4\dd r
 =\frac{8\pi\rho R^5}{15}
 =\boxed{\frac25mR^2}.
\]
四个结果都有质量乘长度平方的量纲；同质量、同外半径下，实心球的质量更靠近轴，
其惯量小于薄球壳，$2/5<2/3$ 与几何分布一致。
\end{xesolution}

\Needspace{8\baselineskip}
\begin{origexamplebox}{【小明的疑惑】}\label{ch09:ex:s16-14-4}
质量为 $m$ 的质点被长为 $l$ 的绳子拉着，在水平面内做圆周运动，绳子与竖直方向夹角
$\theta$。
\begin{enumerate}[label=(\arabic*),leftmargin=2.2em]
  \item 小球的速度大小为多大？
  \item 小球旋转的角速度 $\vect\omega$ 的大小和方向？
  \item 相对于悬挂点，小球角动量 $\vect L$ 的方向？
  \item 角动量与角速度似乎不在同一方向上，这样对吗？
  \item 验证这个运动是否满足角动量定理。
\end{enumerate}


\ThQConePendulum
\end{origexamplebox}

\Needspace{6\baselineskip}\begin{xesolution}
\conicalPendulumDiagram

取 $0<\theta<\pi/2$，绳不可伸长且小球可视为质点。图~\ref{ch09:fig:conical}中圆周半径
$r=l\sin\theta$。绳的张力 $T$ 的竖直分量平衡重力，水平分量提供向心力：
\[
 T\cos\theta=mg,
 \qquad
 T\sin\theta=\frac{mv^2}{l\sin\theta}.
\]
相除得
\[
 \boxed{v=\sqrt{gl\,\frac{\sin^2\theta}{\cos\theta}}}.
\]
小球绕竖直轴转动，故
\[
 \boxed{\omega=\frac{v}{l\sin\theta}
 =\sqrt{\frac{g}{l\cos\theta}}},
\]
$\vect\omega$ 沿竖直轴，方向由圆周运动方向按右手螺旋法则确定。

再看相对于悬点 $O$ 的角动量。与图一致，取 $+x$ 向右、$+z$ 向上、$+y$ 入纸面，
小球此刻在O的右下方，速度沿 $+y$，因此 $\vect\omega$ 向上。此时
\[
 \vect r=(l\sin\theta,0,-l\cos\theta),\qquad
 \vect v=(0,v,0).
\]
所以
\[
 \boxed{\vect L=m\vect r\times\vect v
 =(mlv\cos\theta,0,mlv\sin\theta)}.
\]
$\vect L$ 与绳、竖直轴共面，且垂直于绳；它一般并不与竖直的 $\vect\omega$ 平行。
这并不矛盾：只有绕对称主轴定轴转动时才能简单写成 $\vect L=I\vect\omega$，质点作圆锥摆
并不满足该条件。

最后验证角动量定理。张力沿 $\vect r$，对悬点力矩为零；只有重力产生力矩：
\[
 \vect M_O=\vect r\times m\vect g=(0,mgl\sin\theta,0).
\]
上面的 $\vect L$ 随小球以角速度 $\omega$ 绕竖直轴转动。其水平分量大小为
$L_h=mlv\cos\theta$，因而
\[
 \left|\frac{\dd\vect L}{\dd t}\right|
 =\omega L_h
 =\frac{v}{l\sin\theta}(mlv\cos\theta)
 =\frac{mv^2\cos\theta}{\sin\theta}
 =mgl\sin\theta,
\]
其方向为 $+y$：绕 $+z$ 转动使L的正x水平分量转向 $+y$。
也可直接用 $\dot{\vect L}=\vect\omega\times\vect L$
得到 $(0,\omega mlv\cos\theta,0)$，故大小和方向都与 $\vect M_O$ 一致。
当$\theta$趋于零时v趋于零，轨道缩为一点；角速度公式的有限极限不代表静止点仍有可见的圆周轨迹。
\end{xesolution}

\Needspace{8\baselineskip}
\begin{origexamplebox}{【哑铃落地】}\label{ch09:ex:s16-14-5}
两端各为质量m质点的哑铃由长 $2l$ 的轻杆连接，以速度 $v_0$ 竖直向下纯平动，杆与地面
夹角为 $0\le\theta\le\pi/2$。下端A首先碰到沙地，并在极短碰撞中粘住、固定在该点；
另一端B不同时碰地。忽略碰撞期间重力冲量。用角动量方法求碰后瞬间角速度、B的速度及机械能损失。
这也是把碰撞与能量方法用于哑铃落地模型的另一种解法。


\ThQDumbbell
\end{origexamplebox}

\Needspace{6\baselineskip}\begin{xesolution}
\dumbbellGroundDiagram

图~\ref{ch09:fig:dumbbell-ground}标出了A、B及碰前速度。对固定接触点A取矩，
沙地冲量的力矩为零，故碰撞前后关于A的角动量守恒。取顺时针为正，B的动量产生正角动量。
$\theta$=0的情形应理解为A先接触的极限；若两端同时接触，则是另一种约束，不能套本解。

碰前上端 $B$ 的动量为 $m\vect v_0$，下端 $A$ 相对自身无力臂，所以
\[
 L_A^-=m v_0(2l\cos\theta).
\]
碰后哑铃绕 $A$ 转动，转动惯量只有上端质点贡献：
\[
 I_A=m(2l)^2=4ml^2.
\]
因此
\[
 m v_0(2l\cos\theta)=4ml^2\omega,
 \qquad
 \boxed{\omega=\frac{v_0\cos\theta}{2l}}.
\]
碰后瞬间上端速度大小为
\[
 \boxed{v_B'=2l\omega=v_0\cos\theta},
\]
方向垂直于杆。碰撞损失的机械能为
\[
 \Delta E
 =mv_0^2-\frac12I_A\omega^2
 =\boxed{mv_0^2\left(1-\frac12\cos^2\theta\right)}.
\]
这里不能对整个哑铃应用线动量守恒，因为沙地给出的外冲量不可忽略；但以碰撞点取矩后，
该冲量的力矩冲量为零，角动量守恒恰好可用。
\end{xesolution}

\Needspace{8\baselineskip}
\begin{origexamplebox}{【哑铃的碰撞】}\label{ch09:ex:s16-14-6}
两个相同哑铃进行完全弹性碰撞：一个初始静止，另一个以速度 $v_0$ 撞来。左侧哑铃的右端
重物与右侧哑铃的左端重物发生对心正碰，且碰撞线垂直于两根轻杆。每个端点重物的质量均为
$m$。求：
\begin{enumerate}[label=(\arabic*),leftmargin=2.2em]
  \item 碰后四个重物的运动速度；
  \item 计算接触两端点的恢复系数。
\end{enumerate}


\ThQDumbbells
\end{origexamplebox}

\Needspace{6\baselineskip}\begin{xesolution}
\dumbbellCollisionDiagram

令每根杆长为 $2l$。只讨论沿碰撞线方向的速度；图~\ref{ch09:fig:dumbbell-collision}画的是碰撞前，
运动哑铃纯平动，四个端点依次记为
$A,B,C,D$，其中 $B,C$ 接触。设碰撞冲量大小为 $J$，方向使 $B$ 减速、$C$ 加速。

对左哑铃，总质量 $2m$，质心速度改变量为 $-J/(2m)$；关于质心的角冲量为 $-Jl$，而
$I_{\rm c}=2ml^2$，故角速度改变量为 $-J/(2ml)$。于是
\[
 v_A'=v_0,\qquad v_B'=v_0-\frac{J}{m}.
\]
同理对右哑铃，
\[
 v_C'=\frac{J}{m},\qquad v_D'=0.
\]
这组结果清楚地显示：冲量只改变两个接触端点的瞬时速度，远端点在碰后瞬间仍保持原速度。

完全弹性碰撞还要求系统动能守恒：
\[
 \frac12m(v_0^2+v_0^2)
 =\frac12m\left[v_0^2+\left(v_0-\frac Jm\right)^2
 +\left(\frac Jm\right)^2+0^2\right].
\]
除去无碰撞的 $J=0$ 解，得到 $J=mv_0$，因而
\[
 \boxed{v_A'=v_0,\quad v_B'=0,\quad v_C'=v_0,\quad v_D'=0}.
\]
接触点的接近速度为 $v_B-v_C=v_0$，分离速度为 $v_C'-v_B'=v_0$，所以
\[
 \boxed{e=\frac{v_C'-v_B'}{v_B-v_C}=1}.
\]
这也与题设“完全弹性碰撞”一致。
\end{xesolution}

\Needspace{8\baselineskip}
\begin{origexamplebox}{【跷跷板跳高】}\label{ch09:ex:s16-14-7}
小明想跳得更高，以便扣篮。他把两个轻质、等臂跷跷板顺次放在篮球场上。A 同学质量
$90\,\mathrm{kg}$，B 同学质量 $80\,\mathrm{kg}$，小明质量 $60\,\mathrm{kg}$。A 从
$50\,\mathrm{cm}$ 高处落在第一块跷跷板上，把 B 弹起；B 随后落在第二块跷跷板上，把
小明弹起。所有碰撞均为弹性碰撞。按短时、等臂的理想冲击模型，只计算各人的竖直运动：
跷跷板在碰撞位置近于水平，B两次接触点等高，忽略转移到第二装置所需的水平运动对竖直运动的影响，
两次弹射后人与板均可分离。小明相对于起跳位置的重心能上升多高？


\ThQSeesaws
\end{origexamplebox}

\Needspace{6\baselineskip}\begin{xesolution}
\seesawDiagram

A 落到跷跷板时的速率为
\[
 u_A=\sqrt{2gh_0},\qquad h_0=0.50\,\mathrm m.
\]
图~\ref{ch09:fig:seesaw}用两个面板表示两次冲击，面板间的箭头只表示过程先后，不是飞行轨迹。
设第一次接触中板给A、B的竖直向上冲量大小分别为 $J_A,J_B$。对轻质板的固定支点取冲量矩，
等臂且板的转动惯量忽略，故 $J_A=J_B=J$。令A碰后向下速度为 $u_A'$、B向上速度为 $v_B$：
\[
 m_A(u_A'-u_A)=-J,\qquad m_Bv_B=J.
\]
消去J得 $m_Au_A=m_Au_A'+m_Bv_B$。完全弹性的理想传递还给出
\[
 \frac12m_Au_A^2=\frac12m_Au_A'^2+\frac12m_Bv_B^2.
\]
由第一式 $m_A(u_A-u_A')=m_Bv_B$，代入第二式并除去无传递的零解，得到
$u_A+u_A'=v_B$。再与动量形式的关系联立，第一次弹射速度为
\[
 v_B=\frac{2m_A}{m_A+m_B}u_A.
\]
B 上升后再落回同一高度，撞击第二块跷跷板前速率仍为 $v_B$。第二次弹性碰撞后，小明的
起跳速度为
\[
 v_C=\frac{2m_B}{m_B+m_C}v_B
 =\frac{2m_B}{m_B+m_C}\frac{2m_A}{m_A+m_B}u_A.
\]
故小明相对于起跳位置的上升高度
\begin{align*}
 h&=\frac{v_C^2}{2g}
 =h_0\left(
 \frac{2m_A}{m_A+m_B}
 \frac{2m_B}{m_B+m_C}
 \right)^2\\
 &=0.50\left(\frac{180}{170}\cdot\frac{160}{140}\right)^2\mathrm m
 \approx\boxed{0.732\,\mathrm m}.
\end{align*}
结果是在题设理想传递与等高条件下的重心上升量。轻板两端的瞬时几何约束不意味着A、B分离后的
自由速度相等；正是冲量关系与弹性条件共同给出结果。若B第二次落点低于第一次起跳点 $\Delta h$，
应把第二次入射速率改为 $\sqrt{v_B^2+2g\Delta h}$，不能仍用 $v_B$。
式中速度与 $\sqrt{gh_0}$ 同量纲，全部质量同比放大时高度不变；实际耗散、板质量和主动蹬跳须另建模型。
\end{xesolution}

\Needspace{8\baselineskip}
\begin{origexamplebox}{【漏斗与黄豆】}\label{ch09:ex:s16-14-8}
一个光滑漏斗内部为圆锥面，轴截面是正三角形。小明把一粒黄豆放在漏斗内，给它合适的
水平初速度 $v_0$，黄豆在漏斗中作匀速圆周运动。
\begin{enumerate}[label=(\arabic*),leftmargin=2.2em]
  \item 黄豆转圈位置与圆心的距离为多大？
  \item 若把初速度加倍，且黄豆不会飞出或漏下，最高、最低分别能达到多高？
  \item 若把初速度减半，且黄豆不会飞出或漏下，最高、最低分别能达到多高？
\end{enumerate}


\ThQConeBead
\end{origexamplebox}

\Needspace{6\baselineskip}\begin{xesolution}
\funnelDiagram

图~\ref{ch09:fig:funnel}以O为圆锥顶点，h竖直向上量；P处的初速度垂直轴截面，不能沿半径。
圆锥轴截面是正三角形，所以母线与竖直轴夹角 $\theta=30^\circ$，半径与高度满足
$r=h\tan\theta=h/\sqrt3$。

\textbf{(1)} 黄豆作水平匀速圆周运动。支持力 $N$ 的竖直分量平衡重力，水平分量提供向心力：
\[
 N\sin\theta=mg,
 \qquad
 N\cos\theta=\frac{mv_0^2}{r_0}.
\]
所以
\[
 \frac{v_0^2}{r_0}=g\cot\theta,
 \qquad r_0=\frac{v_0^2\tan\theta}{g}=\boxed{\frac{v_0^2}{\sqrt3g}},
 \qquad
 h_0=r_0\cot\theta=\boxed{\frac{v_0^2}{g}}.
\]

\textbf{(2)(3)} 改变初速度后，黄豆不再作定高圆周运动，但关于圆锥轴的角动量和机械能仍守恒。
令新初速度为 $u=kv_0$，其中加倍时 $k=2$，减半时 $k=1/2$。初态位于 $r_0$，在最高、
最低回归点处径向速度为零。轴向角动量守恒给出
\[
 mur_0=mvr,
 \qquad v=\frac{ur_0}{r}.
\]
机械能守恒，且 $h=\sqrt3r$：
\[
 \frac12mu^2+mg\sqrt3r_0
 =\frac12m\frac{u^2r_0^2}{r^2}+mg\sqrt3r.
\]
令 $x=r/r_0$，再用原匀速圆周条件 $v_0^2=\sqrt3gr_0=gh_0$，把能量等式除以 $mv_0^2$：
\[
 \frac{k^2}{2}+1=\frac{k^2}{2x^2}+x.
\]
乘以 $2x^2$ 并移项，化为
\[
 2x^3-(k^2+2)x^2+k^2
 =(x-1)\bigl(2x^2-k^2x-k^2\bigr)=0.
\]
$x=1$ 是改变速度的瞬间所在回归点，另一个正根为
\[
 x_* =\frac{k^2+k\sqrt{k^2+8}}{4}.
\]
因高度与半径成正比，$h_*=x_*h_0$。

当 $k=2$ 时，$x_*=1+\sqrt3>1$，初始位置是最低点：
\[
 \boxed{h_{\min}=h_0=\frac{v_0^2}{g}},
 \qquad
 \boxed{h_{\max}=(1+\sqrt3)h_0
 =\frac{(1+\sqrt3)v_0^2}{g}}.
\]
\end{xesolution}

\begin{solutioncontinuation}
当 $k=1/2$ 时，$x_*=(1+\sqrt{33})/16<1$，初始位置是最高点：
\[
 \boxed{h_{\max}=h_0=\frac{v_0^2}{g}},
 \qquad
 \boxed{h_{\min}=\frac{1+\sqrt{33}}{16}h_0
 =\frac{(1+\sqrt{33})v_0^2}{16g}}.
\]
图中的 $h$ 是从圆锥顶点所在水平面量起的竖直高度；两种结论都需要满足题设“不飞出、
不漏下”的边界条件。若k=1，方程化为 $(x-1)^2(2x+1)=0$，唯一正回归半径仍为 $r_0$，
恢复原匀速圆周；这也是对全部系数的检验。各高度都与 $v_0^2/g$ 同量纲。
光滑圆锥面的法向方程还给出 $N=m(g\sin\theta+v_\varphi^2\cos\theta/r)>0$，
说明在r为正的这段运动中无需负支持力；实际漏斗的口缘与底孔仍须按题设检查。
\end{solutioncontinuation}
}

{
\graphicspath{{../../../sample/xes-original/seventeenth-section/}}
\section{刚体的运动描述}
% Chapter-local figures. Text remains transparent; geometric coordinates are shared.
\newcommand{\eulerDiagram}{%
\begin{center}\begin{minipage}{\linewidth}\centering\begin{tikzpicture}[>=Latex,font=\small]
 \begin{scope}
  \coordinate(O)at(0,0);\draw[guide](O)circle(1.25);
  \draw[axis](-.6,0)--(1.6,0)node[right]{$x$};\draw[axis](0,-.6)--(0,1.6)node[above]{$y$};
  \draw[vector](O)--(35:1.65)node[above right]{$N$};
  \draw[->,DeepBlue](.7,0)arc[start angle=0,end angle=35,radius=.7];\node at(18:1){$\alpha$};
  \node[DeepBlue,font=\large]at(O){$\odot$};\node[below left=4pt]at(O){$z$};
  \node at(0,2.25){(1) 绕z正转$\alpha$};\node[align=center]at(0,-1.8){沿$-z$方向看向原点\\$z$朝向读者};
 \end{scope}
 \begin{scope}[xshift=4.7cm]
  \coordinate(O)at(0,0);\draw[TextGray](-1.2,0)--(1.2,0)node[right]{$xy$};
  \draw[DeepBlue](145:1.25)--(-35:1.25)node[right]{$XY$};
  \draw[axis](O)--(0,1.7)node[above]{$z$};\draw[vector](O)--(55:1.7)node[above right]{$Z$};
  \draw[->,DeepBlue](90:.65)arc[start angle=90,end angle=55,radius=.65];\node at(72.5:1){$\beta$};
  \node[DeepBlue,font=\large]at(O){$\otimes$};\node[below left=4pt]at(O){$N$};
  \node at(0,2.25){(2) 绕N正转$\beta$};\node[align=center]at(0,-1.8){沿$+N$方向看\\$N$指向纸内};
 \end{scope}
 \begin{scope}[xshift=9.4cm]
  \coordinate(O)at(0,0);\draw[guide](O)circle(1.25);
  \draw[guide](-1.3,0)--(1.6,0)node[right]{$N$};
  \draw[vector](O)--(35:1.65)node[above right]{$X$};\draw[vector](O)--(125:1.65)node[above left]{$Y$};
  \draw[->,DeepBlue](.7,0)arc[start angle=0,end angle=35,radius=.7];\node at(18:1){$\gamma$};
  \node[DeepBlue,font=\large]at(O){$\odot$};\node[below left=4pt]at(O){$Z$};
  \node at(0,2.25){(3) 绕Z正转$\gamma$};\node[align=center]at(0,-1.8){沿$-Z$方向看向原点\\$Z$朝向读者};
 \end{scope}
\end{tikzpicture}\figcap{欧拉角的三步定义与各面板视向}\label{ch09:fig:euler}\end{minipage}\end{center}}

\newcommand{\rollingKinematicsDiagram}{%
\begin{center}\begin{minipage}{\linewidth}\centering\begin{tikzpicture}[>=Latex,font=\small]
 \def\rr{1.1}\coordinate(C)at(0,\rr);
 \draw[thick](-1.8,0)--(2.4,0);\draw[curve](C)circle(\rr);
 \fill(C)circle(1.5pt)node[below left]{$C$};\draw[vector](C)--++(1.9,0)node[right]{$\vect v_C$};
 \draw[->,DeepBlue]($(C)+(155:1.38)$)arc[start angle=155,end angle=-25,radius=1.38];
 \node[left]at(-1.4,2){$\omega$};\node[below]at(0,-.2){平动与顺时针转动};
 \begin{scope}[xshift=6.8cm]
  \coordinate(C)at(0,\rr);\coordinate(P)at(0,0);\coordinate(Top)at(0,2.2);
  \coordinate(Q)at({\rr*cos(145)},{\rr+\rr*sin(145)});
  \coordinate(S)at({\rr*cos(315)},{\rr+\rr*sin(315)});
  \draw[thick](-1.8,0)--(3,0);\draw[curve](C)circle(\rr);
  \foreach\p in{Q,Top,S}{\draw[guide](P)--(\p);\fill(\p)circle(1.5pt);}
  \fill(P)circle(2pt)node[below=4pt]{$P$};\fill(C)circle(1.5pt)node[left]{$C$};
  \node[left=4pt]at(Q){$Q$};\node[above=7pt]at(Top){$P'$};
  \draw[TextGray](S)--(1.45,.72)node[right]{$S$};
  \draw[vector](C)--++(.66,0)node[above]{$\vect v_C$};
  \draw[vector](Top)--++(1.32,0)node[right]{$\vect v_{P'}$};
  \draw[vector](Q)--++({.6*(\rr+\rr*sin(145))},{-.6*\rr*cos(145)});
  \draw[TextGray](-.38,2)--(-1.1,2.6)node[above left]{$\vect v_Q$};
  \draw[vector](S)--++({.6*(\rr+\rr*sin(315))},{-.6*\rr*cos(315)})node[below right=5pt]{$\vect v_S$};
 \end{scope}
\end{tikzpicture}\figcap{固定水平面纯滚动；右图所有速度使用同一比例且垂直于从P出发的连线}\label{ch09:fig:rolling}\end{minipage}\end{center}}

\newcommand{\nonparallelAngularDiagram}{%
\begin{center}\begin{minipage}{\linewidth}\centering\begin{tikzpicture}[>=Latex,font=\small]
 \coordinate(O)at(0,0);\coordinate(A)at(1.25,1.7);\coordinate(B)at(-1.25,-1.7);\coordinate(Z)at(0,2.7);
 \draw[axis](-2,0)--(2.3,0)node[right]{$x$};\draw[axis](0,-2.1)--(Z)node[above]{$z$};
 \draw[curve](B)--(A);\fill(A)circle(2.5pt)node[above=5pt]{$A,m$};\fill(B)circle(2.5pt)node[below=5pt]{$B,m$};
 \fill(O)circle(1.5pt)node[below right=4pt]{$O$};
 \draw[vector](-.32,.1)--(-.32,2.3)node[left]{$\vect\omega$};
 \draw[vector](O)--(-1.7,1.25)node[above left]{$\vect L_O$};
 \node[DeepBlue,font=\Large]at(1.75,1.5){$\otimes$};\node[right]at(2.05,1.5){$\vect v_A$};
 \node[DeepBlue,font=\Large]at(-1.75,-1.45){$\odot$};\node[left]at(-2.05,-1.45){$\vect v_B$};
 \node[right=3pt]at(.55,.75){$l$};
 \draw pic[draw,"$\alpha$",angle radius=9mm,angle eccentricity=1.3]{angle=A--O--Z};
 \node[anchor=west,align=left]at(4,.2){$+y$入纸面\\[6pt]$\vect L_O\perp AB$\\[6pt]$\vect L_O$一般不平行于$\vect\omega$};
\end{tikzpicture}\figcap{两质点斜杆绕竖直轴转动，轴向角动量不等于整个角动量}\label{ch09:fig:nonparallel}\end{minipage}\end{center}}

\newcommand{\inclineRollingDiagram}{%
\begin{center}\begin{minipage}{\linewidth}\centering\begin{tikzpicture}[>=Latex,font=\small]
 \def\tt{25}\def\rr{.82}\coordinate(A)at(0,{5.4*tan(\tt)});\coordinate(B)at(5.4,0);
 \coordinate(P)at(2.8,{2.6*tan(\tt)});\coordinate(C)at($(P)+({\rr*sin(\tt)},{\rr*cos(\tt)})$);
 \draw[thick](0,0)--(A)--(B)--cycle;\draw[curve](C)circle(\rr);
 \draw[guide](P)--(C);
 \draw[vector](P)--++({-1.05*cos(\tt)},{1.05*sin(\tt)})node[above=9pt]{$\vect f$};
 \draw[vector](P)--++({1.9*sin(\tt)},{1.9*cos(\tt)})node[above right=5pt]{$\vect N$};
 \draw[vector](C)--++({1.7*cos(\tt)},{-1.7*sin(\tt)})node[right=5pt]{$\vect a_C$};
 \draw[vector](C)--++(0,-1.1)node[below right=5pt]{$M\vect g$};
 \draw[->,DeepBlue]($(C)+(160:1.05)$)arc[start angle=160,end angle=95,radius=1.05];
 \node[above left=4pt]at($(C)+(150:1.15)$){$\beta$};
 \fill(C)circle(1.5pt)node[above left=2pt]{$C$};\fill(P)circle(1.5pt)node[below left=4pt]{$P$};
 \coordinate(H)at(4.3,0);\draw pic[draw,"$\theta$",angle radius=8mm,angle eccentricity=1.3]{angle=A--B--H};
 \begin{scope}[shift={(8,1.8)}]
  \draw[thick](0,0)circle(.55);\draw[guide](0,0)ellipse(.55 and .17);
  \node at(0,-.95){实心球};\node at(0,-1.65){$I=\frac25MR^2$};
 \end{scope}
 \begin{scope}[shift={(10.7,1.8)}]
  \fill[PaleBlue](0,-.58)rectangle(.65,.58);
  \draw[fill=PaleBlue,thick](.65,0)ellipse(.28 and .58);
  \fill[PaleBlue](0,-.58)rectangle(.65,.58);
  \draw[thick](0,-.58)--(.65,-.58)(0,.58)--(.65,.58);
  \draw[fill=PaleBlue,thick](0,0)ellipse(.28 and .58);
  \node at(.15,-.95){实心圆柱};\node at(.15,-1.65){$I=\frac12MR^2$};
 \end{scope}
 \begin{scope}[shift={(13.4,1.8)}]
  \draw[guide](.65,0)ellipse(.28 and .58);
  \draw[thick](.65,-.58)arc[start angle=-90,end angle=90,x radius=.28,y radius=.58];
  \draw[thick](0,-.58)--(.65,-.58)(0,.58)--(.65,.58);
  \draw[fill=white,thick](0,0)ellipse(.28 and .58);
  \draw[thick](0,0)ellipse(.23 and .51);
  \node at(.15,-.95){薄壁圆筒};\node at(.15,-1.65){$I=MR^2$};
 \end{scope}
\end{tikzpicture}\figcap{圆形刚体沿斜面纯滚动；右侧分辨三种不同质量分布的刚体}\label{ch09:fig:incline}\end{minipage}\end{center}}

\newcommand{\brakeWheelDiagram}{%
\begin{center}\begin{minipage}{\linewidth}\centering\begin{tikzpicture}[>=Latex,font=\small]
 \coordinate(O)at(0,0);\coordinate(P)at(1,0);\draw[curve](O)circle(1.5);
 \draw[->,DeepBlue] (130:1.8)arc[start angle=130,end angle=-25,radius=1.8];\node[above]at(0,1.9){$\omega_0$（顺时针）};
 \draw[vector](P)--++(0,1.1)node[left=5pt]{$\vect f$};
 \draw[thick](O)--(225:1.5)node[midway,above left=5pt]{$R$};
 \draw[guide](O)--(0,-.4)(P)--(1,-.4);\draw[<->](0,-.3)--(1,-.3)node[midway,below]{$r$};
 \fill(O)circle(1.5pt)node[above left]{$O$};\fill(P)circle(1.5pt)node[right]{$P$};
 \node[anchor=west,align=left]at(3,0){力矩逆时针\\[6pt]与初始转动相反};
\end{tikzpicture}\figcap{固定轴砂轮的制动力矩}\label{ch09:fig:brake}\end{minipage}\end{center}}

\newcommand{\twoWheelsDiagram}{%
\begin{center}\begin{minipage}{\linewidth}\centering\begin{tikzpicture}[>=Latex,font=\small]
 \def\ra{1.45}\def\rb{1.15}\coordinate(Ca)at(0,0);\coordinate(Cb)at(2.6,0);\coordinate(P)at(1.45,0);
 \draw[curve](Ca)circle(\ra);\draw[curve](Cb)circle(\rb);
 \node[align=center]at(Ca){$m_1$\\$R_1$};\node[align=center]at(Cb){$m_2$\\$R_2$};
 \draw[->,DeepBlue]($(Ca)+(155:1.7)$)arc[start angle=155,end angle=40,radius=1.7];
 \draw[->,DeepBlue]($(Cb)+(25:1.4)$)arc[start angle=25,end angle=140,radius=1.4];
 \node[above]at(-.5,1.9){$\omega_{10}$};\node[above]at(3.1,1.6){$\omega_{20}$};
 \fill(P)circle(2pt);\draw[TextGray](P)--(1.45,-1.6)node[below]{$P$};
 \draw[guide](P)--(5,0);\node[above]at(5.7,1.6){接触处放大，$u_{10}>u_{20}$};
 \draw[thick](5.7,-1.1)--(5.7,1.1);
 \draw[vector](5.4,0)--(5.4,.85)node[left]{$\vect f_1$};
 \draw[vector](6,0)--(6,-.85)node[right]{$\vect f_2$};
 \draw[->,TextGray](4.7,.3)--(4.7,-.5)node[left]{$u_1$};
 \draw[->,TextGray](6.7,.3)--(6.7,-.5)node[right]{$u_2$};
\end{tikzpicture}\figcap{两个固定轴砂轮的接触与摩擦；另一相对速度分支的力方向反转}\label{ch09:fig:two-wheels}\end{minipage}\end{center}}

\newcommand{\billiardRollingDiagram}{%
\begin{center}\begin{minipage}{\linewidth}\centering\begin{tikzpicture}[>=Latex,font=\small]
 \foreach\s in{0,5.5}{\draw[thick](\s-1.4,0)--(\s+2.1,0);\draw[curve](\s,.85)circle(.85);}
 \draw[vector](0,.85)--(1.7,.85)node[right]{$v_0$};\node at(0,2.35){初态：$\omega=0$};
 \draw[vector](5.5,.85)--(7.2,.85)node[right]{$v(t)$};
 \draw[vector](5.5,0)--(4.3,0)node[below=4pt]{$f=\mu mg$};
 \draw[->,DeepBlue](5.5,.85)++(145:1.05)arc[start angle=145,end angle=20,radius=1.05];
 \node at(5.5,2.35){滑动阶段：$\omega(t)>0$};
\end{tikzpicture}\figcap{台球初态与摩擦建立转动的后续状态}\label{ch09:fig:billiard}\end{minipage}\end{center}}

\newcommand{\ballsCollisionDiagram}{%
\begin{center}\begin{minipage}{\linewidth}\centering\begin{tikzpicture}[>=Latex,font=\small]
 \draw[thick](-1.7,0)--(4,0);\draw[curve](0,1)circle(1);\draw[curve](2,1)circle(1);
 \node[below,align=center]at(0,-.25){白球\\$m,R$};\node[below,align=center]at(2,-.25){红球\\$m,R$};
 \draw[vector](0,1)--(.75,1)node[pos=.45,above=4pt]{$v_0$};
 \draw[->,DeepBlue](0,1)++(140:1.2)arc[start angle=140,end angle=45,radius=1.2];
 \node[above]at(-.8,2.1){$\omega_0$};\node[align=center]at(2,1){静止\\不自转};
 \node[anchor=west,align=left]at(4.5,1){碰撞前\\[6pt]白球：$v_0=\omega_0R$};
\end{tikzpicture}\figcap{相同台球无摩擦正碰的初始状态}\label{ch09:fig:balls}\end{minipage}\end{center}}

\newcommand{\bounceSpinDiagram}{%
\begin{center}\begin{minipage}{\linewidth}\centering\begin{tikzpicture}[>=Latex,font=\small]
 \draw[thick](-1.7,0)--(1.7,0);\coordinate(Ca)at(0,1.7);\draw[curve](Ca)circle(.8);
 \draw[vector](Ca)--++(0,-1.1)node[right=3pt]{$v_0$};
 \draw[->,DeepBlue]($(Ca)+(20:1.02)$)arc[start angle=20,end angle=150,radius=1.02];\node[above]at(0,2.8){$\omega_0$};
 \node at(0,-.55){碰前：竖直下落并自旋};
 \begin{scope}[xshift=5.2cm]
  \draw[thick](-1.7,0)--(2,0);\draw[curve](0,.8)circle(.8);\coordinate(P)at(0,0);
  \draw[vector](P)--(-1.3,0)node[below=5pt]{$\vect J_t$};
  \draw[vector](P)--(0,1.25)node[right=3pt]{$\vect J_n$};
  \fill(P)circle(1.6pt)node[below right=3pt]{$P$};
  \node[above]at(0,2.1){接触冲量图};\node at(0,-.75){球与地面相切于P};
 \end{scope}
 \begin{scope}[xshift=9.9cm]
  \coordinate(O)at(0,.7);\draw[guide](O)--(0,2.4);\draw[vector](O)--(-.85,2.15)node[above left]{$\vect v'$};
  \coordinate(V)at(0,2.4);\coordinate(W)at(-.85,2.15);
  \draw pic[draw,"$\alpha$",angle radius=8mm,angle eccentricity=1.35]{angle=V--O--W};
  \node at(0,-.55){弹起方向向左偏};
 \end{scope}
\end{tikzpicture}\figcap{旋转篮球的碰前、接触和碰后三个不同阶段}\label{ch09:fig:bounce-analysis}\end{minipage}\end{center}}

\newcommand{\eggRollingDiagram}{%
\begin{center}\begin{minipage}{\linewidth}\centering\begin{tikzpicture}[>=Latex,font=\small]
 \foreach\s/\mode in{0/0,6.6/1}{
  \begin{scope}[xshift=\s cm]
   \def\tt{22}\coordinate(P)at(2,1.1);
   \coordinate(C)at($(P)+({.78*sin(\tt)},{.78*cos(\tt)})$);
   \draw[thick]($(P)+({-2*cos(\tt)},{2*sin(\tt)})$)--($(P)+({2.8*cos(\tt)},{-2.8*sin(\tt)})$);
   \begin{scope}[shift={(C)},rotate=-22]
    \draw[fill=PaleBlue,thick](0,0)ellipse(.58 and .78);
    \ifnum\mode=0\draw[fill=DeepBlue!12,thick](0,0)ellipse(.36 and .51);
    \else\draw[DeepBlue](0,-.1)ellipse(.35 and .43);\draw[TextGray](-.32,.1)..controls(-.12,.24)and(.06,-.02)..(.32,.12);\fi
   \end{scope}
   \fill(P)circle(1.3pt)node[below left=3pt]{$P$};
   \draw[guide]($(C)+(-1.3,0)$)--($(C)+(1.3,0)$);
   \draw[vector]($(C)+(.9,.2)$)--++({1.1*cos(\tt)},{-1.1*sin(\tt)})node[right]{$\vect v_C$};
   \draw[->,DeepBlue]($(C)+(150:1.05)$)arc[start angle=150,end angle=65,radius=1.05];
   \node[above]at($(C)+(0,1.2)$){向右下滚时外壳顺时针转};
   \ifnum\mode=0\node[align=center]at(2,-.45){熟鸡蛋\\整体近似共同转动};
   \else\node[align=center]at(2,-.45){生鸡蛋\\内部液体可相对外壳运动};\fi
  \end{scope}}
\end{tikzpicture}\figcap{相同释放高度的两种内部结构；剖面仅用于区分凝固与液体}\label{ch09:fig:egg}\end{minipage}\end{center}}

\newcommand{\massivePulleyDiagram}{%
\begin{center}\begin{minipage}{\linewidth}\centering\begin{tikzpicture}[>=Latex,font=\small]
 \coordinate(O)at(0,2.5);\def\rr{.8}
 \draw[thick](-1.8,4)--(1.8,4);\draw[thick](0,4)--(O);\draw[curve](O)circle(\rr);
 \fill(O)circle(2pt)node[right=4pt]{$O$};
 \draw[TextGray](-.6,3.03)--(-1.2,2.9)node[left]{$M,R$};
 \draw[thick](-.8,.65)--(-.8,2.5)arc[start angle=180,end angle=0,radius=.8]--(.8,1.05);
 \draw[thick](-1.1,0)rectangle(-.5,.65);\node at(-.8,.325){$m_1$};
 \draw[thick](.5,.4)rectangle(1.1,1.05);\node at(.8,.725){$m_2$};
 \draw[vector](-1.55,1.1)--(-1.55,.1)node[left]{$\vect a_1$};
 \draw[vector](1.55,.4)--(1.55,1.4)node[right]{$\vect a_2$};
 \node[anchor=west,align=left]at(3,2.1){轴O固定\\[6pt]绳不伸长、与轮缘不滑\\[6pt]$a=\beta R$};
\end{tikzpicture}\figcap{重滑轮的轮轴、绳路和两端加速度}\label{ch09:fig:pulley}\end{minipage}\end{center}}

\newcommand{\torqueWorkDiagram}{%
\begin{center}\begin{minipage}{\linewidth}\centering\begin{tikzpicture}[>=Latex,font=\small]
 \coordinate(O)at(0,2.7);\coordinate(H)at(4.7,2.7);\coordinate(P)at(4,.4);\coordinate(C)at($(O)!.5!(P)$);
 \draw[thick](-.8,2.7)--(.8,2.7);\draw[guide](O)--(H);
 \draw[curve](O)--(P);\fill(O)circle(1.8pt)node[above left]{$O$};
 \fill(C)circle(2.5pt)node[below left=4pt]{$m_2=m$};\fill(P)circle(2.5pt)node[right=4pt]{$m_1=m$};
 \node[below left=5pt]at($(O)!.25!(P)$){$l$};\node[above right=3pt]at($(O)!.75!(P)$){$l$};
 \draw pic[draw,"$\theta$",angle radius=13mm,angle eccentricity=1.3]{angle=P--O--H};
\end{tikzpicture}\figcap{轻杆从水平位置向下转过$\theta$，两质点分别下降$2l\sin\theta$和$l\sin\theta$}\label{ch09:fig:torque-work}\end{minipage}\end{center}}

\newcommand{\percussionDiagram}{%
\begin{center}\begin{minipage}{\linewidth}\centering\begin{tikzpicture}[>=Latex,font=\small]
 \coordinate(O)at(0,3.8);\coordinate(P)at(0,1.3);\coordinate(B)at(0,0);
 \draw[thick](-.8,3.8)--(.8,3.8);\draw[curve](O)--(B);
 \fill(O)circle(2pt)node[above right]{$O$};\fill(P)circle(1.6pt)node[below right=4pt]{$P$};
 \draw[vector](-2.2,1.3)--(P)node[pos=.28,above]{$\vect J$};
 \draw[guide](O)--(1.3,3.8)(B)--(1.3,0);
 \draw[guide](O)--(-.8,3.8)(B)--(-.8,0);
 \draw[<->](-.65,1.3)--(-.65,3.8)node[midway,left=4pt]{$x=L-a$};
 \draw[<->](-.65,0)--(-.65,1.3)node[midway,left=4pt]{$a$};
 \draw[<->](1.2,0)--(1.2,3.8)node[midway,right]{$L$};
 \node[anchor=west]at(3,1.9){支点冲量为零时求P的位置};
\end{tikzpicture}\figcap{固定铰点细杆的打击位置与长度定义}\label{ch09:fig:percussion}\end{minipage}\end{center}}

\newcommand{\supportedRodDiagram}{%
\begin{center}\begin{minipage}{\linewidth}\centering\begin{tikzpicture}[>=Latex,font=\small]
 \draw[curve](0,1)--(4,1);\foreach\x in{0,4}{\draw[thick](\x-.3,.45)--(\x,1)--(\x+.3,.45)--cycle;}
 \draw[vector](0,1)--(0,2)node[above]{$N_A=mg/2$};\draw[vector](4,1)--(4,2)node[above]{$N_B=mg/2$};
 \draw[vector](2,1)--(2,-.05)node[below]{$mg$};\node[above]at(2,1.1){$C$};
 \node[below]at(0,.4){$A$};\node[below]at(4,.4){$B$};\node at(2,2.65){整杆};
 \begin{scope}[xshift=7cm]
  \draw[curve](0,1)--(2,1);\draw[thick](-.3,.45)--(0,1)--(.3,.45)--cycle;
  \draw[vector](0,1)--(0,2)node[above]{$N_A$};\draw[vector](1,1)--(1,-.05)node[below]{$mg/2$};
  \draw[guide](2,.4)--(2,1.6);\node[above]at(2,1.75){中截面C};
  \draw[->,DeepBlue](2,1)++(-65:.55)arc[start angle=-65,end angle=200,radius=.55];
  \node[right]at(2.8,.9){$M_{\rm int}$};\node[right]at(2.8,.35){$Q=0$};
  \draw[<->](0,-.7)--(1,-.7)node[midway,below]{$L/4$};\draw[<->](1,-.7)--(2,-.7)node[midway,below]{$L/4$};
  \node at(1,2.65){左半杆};
 \end{scope}
\end{tikzpicture}\figcap{端点支撑杆的整体与左半杆；中截面仅有非零弯矩}\label{ch09:fig:rod-analysis}\end{minipage}\end{center}}

\newcommand{\chimneyDiagram}{%
\begin{center}\begin{minipage}{\linewidth}\centering\begin{tikzpicture}[>=Latex,font=\small]
 \def\tt{32}\def\ll{4.5}\def\xx{1.45}
 \coordinate(O)at(0,0);\coordinate(T)at({\ll*sin(\tt)},{\ll*cos(\tt)});
 \coordinate(P)at({\xx*sin(\tt)},{\xx*cos(\tt)});\coordinate(G)at($(P)!.5!(T)$);\coordinate(V)at(0,4.5);
 \draw[guide](O)--(V);\draw[curve,line width=1.4pt](O)--(T);
 \fill(O)circle(2pt)node[below]{$O$};\fill(P)circle(1.5pt)node[below right=4pt]{$P$};\fill(G)circle(1.5pt)node[right=4pt]{$G$};
 \draw[vector](P)--++({-1.3*cos(\tt)},{1.3*sin(\tt)})node[above left]{$-Q\uvec e_t$};
 \draw[vector](G)--++(0,-.85)node[right=5pt]{$m_s\vect g$};
 \draw[->,DeepBlue]($(P)+(80:.65)$)arc[start angle=80,end angle=-110,radius=.65];
 \node[right=4pt]at($(P)+(2,-.2)$){$M$（正向）};
 \draw[TextGray]($(P)+(1.85,-.2)$)--($(P)+(.65,0)$);
 \draw[->,TextGray](3.4,2.6)--++({.8*cos(\tt)},{-.8*sin(\tt)})node[right]{$\uvec e_t$};
 \draw pic[draw,"$\theta$",angle radius=8mm,angle eccentricity=1.3]{angle=T--O--V};
 \draw[<->]($(O)+(1,-.63)$)--($(P)+(1,-.63)$)node[midway,below right]{$x$};
 \draw[<->]($(P)+(-.6,.38)$)--($(T)+(-.6,.38)$)node[midway,above left]{$s$};
 \node[anchor=west,align=left]at(5,2.4){研究P以上一段\\[6pt]G为该段质心\\[6pt]$\beta$、M以顺时针为正\\[6pt]Q为有符号量};
\end{tikzpicture}\figcap{烟囱上段的剪力、弯矩、重力与统一方向约定}\label{ch09:fig:chimney}\end{minipage}\end{center}}


\newcommand{\bounceSpinProblemDiagram}{%
\begin{center}\begin{minipage}{\linewidth}\centering\begin{tikzpicture}[>=Latex,font=\small]
 \draw[thick](-1.7,0)--(1.7,0);\coordinate(Ca)at(0,1.7);\draw[curve](Ca)circle(.8);
 \draw[vector](Ca)--++(0,-1.1)node[right=3pt]{$v_0$};
 \draw[->,DeepBlue]($(Ca)+(20:1.02)$)arc[start angle=20,end angle=150,radius=1.02];
 \node[above]at(0,2.8){$\omega_0$};\node[right]at(1.2,1.7){$m,R$};
\end{tikzpicture}\figcap{篮球落地前的平动与自旋方向}\label{ch09:fig:bounce}\end{minipage}\end{center}}

\newcommand{\supportedRodProblemDiagram}{%
\begin{center}\begin{minipage}{\linewidth}\centering\begin{tikzpicture}[>=Latex,font=\small]
 \draw[curve](0,1)--(6,1);\foreach\x in{0,6}{\draw[thick](\x-.3,.45)--(\x,1)--(\x+.3,.45)--cycle;}
 \node[below]at(0,.4){$A$};\node[below]at(6,.4){$B$};\node[above]at(3,1.08){均匀杆，质量 $m$};
 \draw[<->](0,-.45)--(6,-.45)node[midway,below]{$L$};
\end{tikzpicture}\figcap{均匀杆由两端支点水平支撑}\label{ch09:fig:rod}\end{minipage}\end{center}}


刚体是忽略形变的理想质点系：任意两质点之间的距离在运动中保持不变。真实固体受力后都会
发生形变，只有形变相对物体尺度很小、且对所研究运动的影响可忽略时，才采用刚体模型。
它不是“绝对坚硬”的真实材料；局部压缩、弹性波、内部振动及断裂需要另行考虑。

\subsection{位置、姿态与自由度}
质心的三个坐标只确定位置，不确定朝向。例如一本书可以质心不动而转动。自由三维刚体需要
三个位置参数和三个姿态参数，共六个构型自由度。固定一点后剩三个转动自由度；固定一条轴后
剩一个；平面运动有两个平移和一个转动自由度。无滑动接触还在速度层面增加约束，不能重复计数。

欧拉角（Euler angles）是一组描述姿态的坐标。取原点都在质心的两组右手轴：xyz的方向
保持相对惯性系不变，XYZ与刚体固连。前者若随质心加速平移，其原点并非惯性系中的固定点，
“方向固定”不等于“整个坐标系是惯性系”。采用 $z$--$N$--$Z$ 转动约定：
\begin{enumerate}[label=(\arabic*),leftmargin=2.4em]
 \item 绕z转$\alpha$，把x转到节线N；N是xy与XY平面的交线，取向由转动约定给出。
 \item 绕N按右手规则转$\beta$，把z转到Z。
 \item 绕Z转$\gamma$，把N转到固连轴X。
\end{enumerate}
图~\ref{ch09:fig:euler}把三次转动分开画；第二面板N指向纸内，因而正$\beta$在该观察方向下是顺时针。
三面板的视向不同，不能把它们的纸面左右直接当成同一组三维坐标。
\eulerDiagram
通常约定 $0\le\alpha,\gamma<2\pi$、$0\le\beta\le\pi$。
当$\beta$为0或$\pi$时，两平面重合，节线不唯一，$\alpha$与$\gamma$不能独立确定；这是坐标退化，不是刚体失去姿态。
因此欧拉角只在排除退化并固定角度范围后局部唯一。位置和姿态确定\textbf{构型}；要确定
动力学状态，还须给质心速度和角速度，或等价的广义速度。

有限转动一般不可交换，先绕一轴再绕另一轴与倒序结果不同。无穷小转动在一阶近似下才可按
矢量相加。故实际角速度应为
\begin{equation}
 \vect\omega=\dot\alpha\,\uvec e_z+\dot\beta\,\uvec e_N+\dot\gamma\,\uvec e_Z,
\end{equation}
右侧三根轴不通常互相垂直，应投影到同一组轴再求分量；三个角的导数不是三个固定直角分量。
定轴或平面运动只需一个角$\varphi$时，才可直接以有符号量写 $\omega=\dot\varphi$、$\beta=\ddot\varphi$。

\subsection{刚体速度、加速度与瞬心}
平动中$\omega$=0，各点速度相同，任一固连线段始终平行于自身；定轴转动中固连轴不动，
各点绕该轴作圆周运动；平面平行运动则包含质心平动及绕质心的转动，车轮和滑杆是常见模型。
对刚体任意两点A、B，有
\begin{equation}
 \vect v_B=\vect v_A+\vect\omega\times\overrightarrow{AB}.\label{ch09:eq:rigid-velocity}
\end{equation}
与 $\overrightarrow{AB}$ 点乘，可得
$(\vect v_B-\vect v_A)\cdot\overrightarrow{AB}=0$，即相对速度在连线方向上的分量为零。
这也等于对固定长度条件 $\overrightarrow{AB}\cdot\overrightarrow{AB}=\text{常量}$ 求导。

平面运动且$\omega$不为零时，速度场可等效为绕瞬时速度为零的点P转动。若已知两点速度方向，
过两点作各自速度的垂线，其非平行交点给出瞬心；平行等速的特殊情形应回到速度方程判断，
不能强求一个有限交点。瞬心可以位于刚体之外的延伸速度场中，不一定是实际材料点。
对非零速度纯平动，$\omega$=0，不存在有限距离的瞬时静止点。

图~\ref{ch09:fig:rolling}中车轮向右纯滚，接触地面的P瞬时静止，故任一点Q满足
$\vect v_Q=\vect\omega\times\overrightarrow{PQ}$。
速度必须垂直PQ，且 $v_Q/\omega=PQ$。顶点 $P'$ 到P的距离为2R，故顶点速率为质心速率的两倍。
\rollingKinematicsDiagram
瞬心不是一个始终固定的铰点。下一时刻轮缘接触地面的是另一个材料点，所以不能把整个过程
当成绕同一P的定轴转动。对同一对材料点，式~\eqref{ch09:eq:rigid-velocity}求导给出
\begin{equation}
 \vect a_B=\vect a_A+\vect\beta\times\overrightarrow{AB}
 +\vect\omega\times(\vect\omega\times\overrightarrow{AB}).
\end{equation}
第二项是切向项，第三项是法向项。P瞬时速度为零，不代表其加速度为零。
例如水平面上以恒定质心速度纯滚的圆轮，接地点材料点有竖直向上的加速度 $\omega^2R$。
因此瞬心法可用于速度，不能未经推导就用 $\sum M_P=I_P\beta$ 代替质心动力学方程。

\section{刚体动力学与纯滚动}
将各质元速度相加，利用相对质心速度的质量加权和为零，得到
\begin{equation}
 \vect P=M\vect v_C,\qquad \sum\vect F_{\rm ext}=M\vect a_C.
\end{equation}
同一外力无论作用在何处，对质心加速度的贡献相同；但对转动的作用由其作用线决定。
角动量则为 $\vect L_O=\overrightarrow{OC}\times M\vect v_C+\vect L_C$。
前一节已说明 $\vect L_C$ 不一定与 $\vect\omega$ 平行，图~\ref{ch09:fig:nonparallel}给出具体反例。
\nonparallelAngularDiagram
设轻杆两端各质量m，质心O到各端距离l，杆与竖直轴夹角$\alpha$。图中 $+x$ 向右、$+z$ 向上，
为右手系，$+y$ 指向纸内。若 $\vect\omega=(0,0,\omega)$，上端A满足
\begin{equation}
 \vect r_A=(l\sin\alpha,0,l\cos\alpha),\qquad
 \vect v_A=(0,\omega l\sin\alpha,0).
\end{equation}
下端B的位置、速度均相反，所以两端对O的角动量相同，求和为
\begin{equation}
 \vect L_O=2m\vect r_A\times\vect v_A
 =2ml^2\omega(-\sin\alpha\cos\alpha,0,\sin^2\alpha).
\end{equation}
当 $0<\alpha<\pi/2$，它既有负x分量又有正z分量，不与竖直角速度同向。
其轴向分量仍是 $L_z=(2ml^2\sin^2\alpha)\omega=I_z\omega$。
这说明标量轴向关系可以正确，而向量式 $\vect L=I\vect\omega$ 却不成立。

一般关系为 $\vect L_C=\mathbf I_C\vect\omega$。仅沿主惯量轴转动时，角动量与角速度
平行，才能用一个标量惯量描述整个向量关系；球对称物体每个过中心方向都是主轴。
一般自由三维转动即使外力矩为零，L保持固定，刚体的轴仍可能摆动。

对平面刚体，最常用的方程组是
\begin{equation}
 \sum F_x=Ma_{Cx},\qquad \sum F_y=Ma_{Cy},\qquad \sum M_C=I_C\beta.
\end{equation}
最后一式只取垂直运动平面的分量，转轴、力矩和$\beta$使用同一正方向。若物体确绕惯性系中的固定轴
转动，也可用关于该轴的惯量写 $\sum M=I\beta$。对任意运动点则须使用前一节的参考点变换式，
不能因为某点瞬时静止就视为始终固定。摩擦、支持力通常是联立方程的未知量，不应先凭直觉定值。

\subsection{斜面纯滚动与摩擦判据}
设质量M、外半径R、绕质心轴惯量I的圆形刚体从静止沿倾角$\theta$的固定斜面下滚，
无外加驱动转矩。沿坡向下及顺时针为正，先假设静摩擦沿坡向上。
图~\ref{ch09:fig:incline}中P为接触点，$CP=R$ 且CP垂直斜面；N沿CP，重力通过C。
\inclineRollingDiagram
法向没有质心加速度，切向和转动方程分别为
\begin{equation}
 N=Mg\cos\theta,\qquad Mg\sin\theta-f=Ma_C,\qquad fR=I\beta.
\end{equation}
若不滑动，$a_C=\beta R$，因此 $f=Ia_C/R^2$，代回切向方程：
\begin{equation}
 Mg\sin\theta=(M+I/R^2)a_C,
 \qquad \boxed{a_C=\frac{g\sin\theta}{1+I/(MR^2)}}.
\end{equation}
继而求出实际所需的静摩擦
\begin{equation}
 f=\frac{I}{I+MR^2}Mg\sin\theta,\qquad
 \boxed{\mu_s\ge\frac fN=\frac{I}{I+MR^2}\tan\theta}.
\end{equation}
只有满足不等式，纯滚动假设才自洽。否则改用动摩擦 $\mu_kN$，并撤去 $a_C=\beta R$，
分别求平动与转动，直到接触点相对速度归零后再判断能否维持纯滚动。

令 $\kappa=I/(MR^2)$，图中实心球、实心圆柱、薄壁圆筒分别有
$\kappa=2/5,1/2,1$，相应加速度为 $5g\sin\theta/7,2g\sin\theta/3,g\sin\theta/2$。
在同一坡面、相同质心沿坡路程、同样从静止出发且均不滑时，实心球先到，薄壁圆筒后到。
总质量与绝对尺寸从比值中约去，决定先后的是质量分布。

静摩擦方向并非永远向上坡：加外转矩、改变作用力位置或让接触面运动，都可能使其反向或为零。
应先判断没有摩擦时接触点的相对滑动趋势，再用方程的正负校验。固定面上纯滚时
$P_f=\vect f\cdot\vect v_P=0$，但摩擦仍同时改变平动和转动动能；这两部分的功率相加为零。
运动传送带上的接触点对地速度不为零，静摩擦对物体可以做非零功。

\section{刚体的功、能量与碰撞}
任一质元速度分解为 $\vect v_i=\vect v_C+\vect\omega\times\vect r_i'$。
平方并求和，混合项含 $\sum_i m_i\vect r_i'=0$，因此
\begin{equation}
 E_k=\frac12Mv_C^2+\frac12\sum_i m_i|\vect\omega\times\vect r_i'|^2.
\end{equation}
平面刚体的角速度沿同一法向，第二项成为 $I_C\omega^2/2$；一般三维情形为
$\vect\omega\cdot\mathbf I_C\vect\omega/2$。不能把一般三维运动也一律套单个I。

外力作用点速度不同，应先逐力计算功率。利用标量三重积恒等式，
\begin{align}
 P&=\sum_i\vect F_i\cdot(\vect v_C+\vect\omega\times\vect r_i')\\
 &=\left(\sum_i\vect F_i\right)\cdot\vect v_C+
 \left(\sum_i\vect r_i'\times\vect F_i\right)\cdot\vect\omega\\
 &=\vect F_{\rm ext}\cdot\vect v_C+\vect M_C\cdot\vect\omega.
\end{align}
第一项给质心平动动能变化率，第二项给转动动能变化率。同一外力可同时参与两部分，
不能在计算功率时任意把作用点搬到质心。理想刚性约束的成对内力不对整体做净功，
但能在部分物体之间传递能量；若含主动内部驱动或可变形自由度，需另计相应功和储能。

绕固定轴转动，作用点的切向位移为 $r_\perp\dd\varphi$，故
\begin{equation}
 \dd W=M_{\rm axis}\dd\varphi,\qquad P=M_{\rm axis}\omega,
 \qquad W=\int_{\varphi_1}^{\varphi_2}M_{\rm axis}(\varphi)\dd\varphi=\Delta E_k.
\end{equation}
恒力矩且I不变才对应恒角加速度；力臂随姿态改变时应积分，不能简单用初始力矩乘总转角。
对纯滚动下坡，也可令势能减少等于两部分动能之和；它与动力学结果相符，且说明把物体当作
仅平动的质点会高估下坡速度和加速度。

\subsection{碰撞、打击与截面内力}
短时碰撞应同时考虑线冲量和冲量矩，对固定矩心O有
\begin{equation}
 M(\vect v_C^+-\vect v_C^-)=\vect J_{\rm ext},\qquad
 \vect L_O^+-\vect L_O^-=\int\vect M_O^{\rm ext}\dd t.
\end{equation}
若只有某一接触冲量J，对质心还有
$\Delta\vect L_C=\overrightarrow{CP}\times\vect J$；平面情况下可再写成I乘角速度改变量。
选择固定支点或碰撞点取矩，常能消去该处冲量矩，但不一定能消去其余外冲量矩。
打击中心就是平动与转动响应恰好配合，使固定支点不需提供冲量的打击位置；并非简单的几何中点。

碰撞通常伴随形变、声、热，动量或角动量守恒不保证动能守恒。仅有法向恢复系数等于1时，
还须查看切向摩擦是否耗能，不能直接宣称全部动能不变。
生鸡蛋内部液体引入相对流动自由度，不是单一角速度刚体；既要考虑流体参与转动的程度，
也要考虑耗散，后面的生熟鸡蛋题说明这一模型边界。

刚体假设各点距离不变，不意味着内部作用为零。杆或烟囱靠轴力、剪力和弯矩维持整体运动。
分析方法是先用整体方程求运动，再在待考察位置截开，对其中一段补齐截面相互作用，分别列力
与力矩方程。具体题里某项可以为零；只有实际非零的内力分布才用于比较断裂位置。
能量法适合连接状态、求角速度，动力学适合求瞬时反力及加速度，两者配合并以量纲、方向和
极限情况相互校验，不能仅靠整体能量判断哪个截面最危险。


\Needspace{8\baselineskip}
\begin{origexamplebox}{【砂轮制动】}\label{ch09:ex:s17-15-1}
一个质量为 $m$、半径为 $R$ 的匀质砂轮以角速度 $\omega_0$ 旋转。在距离圆心为 $r$ 处
有恒力 $f$ 作用，方向如图。求随后砂轮的运动。


\ThQBrakingWheel
\end{origexamplebox}

\Needspace{6\baselineskip}\begin{xesolution}
\brakeWheelDiagram

图~\ref{ch09:fig:brake}中砂轮初始顺时针转动，f在圆心右侧向上，故力矩逆时针。
取顺时针为正，匀质砂轮视为均匀圆盘，绕中心轴的转动惯量
\[
 I=\frac12mR^2.
\]
力对轴的力矩大小为 $fr$，方向与初始转动相反，所以角加速度为常量
\[
 \beta=-\frac{fr}{I}=-\frac{2fr}{mR^2}.
\]
于是制动阶段
\[
 \boxed{\omega(t)=\omega_0-\frac{2fr}{mR^2}t},
 \qquad
 \boxed{\varphi(t)=\omega_0t-\frac{fr}{mR^2}t^2}.
\]
砂轮在
\[
 \boxed{t_s=\frac{mR^2\omega_0}{2fr}}
\]
时停止，停止前转过的角度为
\[
 \boxed{\varphi_s=\frac{mR^2\omega_0^2}{4fr}}.
\]
若题意中的恒力在停止后仍保持原方向，砂轮会反向转动；“制动”通常只讨论首次停转之前。
\end{xesolution}

\Needspace{8\baselineskip}
\begin{origexamplebox}{【两个砂轮】}\label{ch09:ex:s17-15-2}
两个匀质砂轮质量分别为 $m_1,m_2$，半径分别为 $R_1,R_2$，初始分别以角速度
$\omega_{10},\omega_{20}$ 按图示方向旋转。$t=0$ 时两轮接触，此后接触处有恒定摩擦力
$f$（仅在相对滑动期间取此大小，两轮轴固定且忽略其他阻力矩）。求：
\begin{enumerate}[label=(\arabic*),leftmargin=2.2em]
  \item 任意时刻二者的角速度；
  \item 这个问题是否有类似的质点动力学问题相对应？
\end{enumerate}


\ThQWheels
\end{origexamplebox}

\Needspace{6\baselineskip}\begin{xesolution}
\twoWheelsDiagram

图~\ref{ch09:fig:two-wheels}的第一轮以顺时针为正、第二轮以逆时针为正；两接触表面初速度均向下。
令 $u_{10}=R_1\omega_{10}$、$u_{20}=R_2\omega_{20}$，以向下为正的初始相对速度
$s_0=u_{10}-u_{20}$。若 $s_0>0$，摩擦使第一轮减速、第二轮加速，图中力箭头对应此分支；
$s_0<0$ 时两摩擦箭头都反向。以 $\varepsilon=\operatorname{sgn}(s_0)$ 统一记号，
$I_i=\tfrac12m_iR_i^2$，持续滑动期间
\[
 \boxed{\omega_1(t)=\omega_{10}-\varepsilon\frac{fR_1}{I_1}t
 =\omega_{10}-\varepsilon\frac{2f}{m_1R_1}t},
\]
\[
 \boxed{\omega_2(t)=\omega_{20}+\varepsilon\frac{fR_2}{I_2}t
 =\omega_{20}+\varepsilon\frac{2f}{m_2R_2}t}.
\]
两接触表面不再相对滑动的条件为 $R_1\omega_1=R_2\omega_2$；因此滑动终止时刻
\[
 \boxed{t_*=
 \frac{|R_1\omega_{10}-R_2\omega_{20}|}
 {f(R_1^2/I_1+R_2^2/I_2)}}.
\]
上述含t的式子只用于 $0\le t\le t_*$。此后两轮无其他阻力矩，维持无相对滑动所需的静摩擦为零，
各自以最终角速度匀速转动。若初始 $s_0=0$，则 $t_*=0$，直接处于这一状态，不能强加大小f的动摩擦。

它对应于两个质点在恒定相互摩擦力下相对滑动的问题。映射关系是：切向线速度
$u_i=R_i\omega_i$ 对应质点速度，而“等效质量”
\[
 m_i^{\rm eff}=\frac{I_i}{R_i^2}
\]
对应质点质量。记 $q_i=I_i/R_i^2$，两轮切向方程为 $q_1\dot u_1=-\varepsilon f$、
$q_2\dot u_2=\varepsilon f$，相加得 $q_1u_1+q_2u_2$ 不变。于是最终共同轮缘速度及角速度为
\[
 u_* =\frac{q_1u_{10}+q_2u_{20}}{q_1+q_2},\qquad
 \omega_{1*}=\frac{u_*}{R_1},\quad \omega_{2*}=\frac{u_*}{R_2}.
\]
这不是说固定轴承系统的总线动量守恒，而是切向方程与两个相互作用质点的方程具有相同形式。
共同速度位于两初始轮缘速度之间；机械能损失
\[
 \Delta E=\frac12\frac{q_1q_2}{q_1+q_2}(u_{10}-u_{20})^2\ge0
\]
可由初末转动动能相减得到，量纲为能量且 $s_0=0$ 时为零。
\end{xesolution}

\Needspace{8\baselineskip}
\begin{origexamplebox}{【滚动的台球】}\label{ch09:ex:s17-15-3}
质量为 $m$、半径为 $R$ 的匀质台球刚被击出时速度为 $v_0$，没有旋转。已知台球与地面的
滑动摩擦因数为 $\mu$，求台球此后的运动。


\ThQBilliard
\end{origexamplebox}

\Needspace{6\baselineskip}\begin{xesolution}
\billiardRollingDiagram

图~\ref{ch09:fig:billiard}分别画出初态和滑动阶段，避免把 $\omega=0$ 的初态与后续转动混为一图。
初始接触点相对地面向前滑，滑动摩擦力大小为 $f=\mu mg$、方向向后。匀质实心球
$I=\tfrac25mR^2$。在滑动阶段
\[
 ma=-\mu mg,
 \qquad I\beta=\mu mgR.
\]
取使球进入向前纯滚动的角速度为正，得到
\[
 v(t)=v_0-\mu gt,
 \qquad
 \omega(t)=\frac{\mu mgR}{I}t=\frac{5\mu g}{2R}t.
\]
当 $v=\omega R$ 时开始无滑动滚动：
\[
 v_0-\mu gt_* =\frac52\mu gt_*,
 \qquad
 \boxed{t_* =\frac{2v_0}{7\mu g}}.
\]
相应速度和角速度为
\[
 \boxed{v_* =\frac57v_0},
 \qquad
 \boxed{\omega_* =\frac{5v_0}{7R}}.
\]
若忽略滚动阻力，$t>t_*$ 后台球以这两个恒定值作无滑动滚动。滑动阶段质心前进距离为
\[
 \boxed{x_* =v_0t_* -\frac12\mu gt_*^2
 =\frac{12v_0^2}{49\mu g}}.
\]
最终总动能为 $\tfrac12m(5v_0/7)^2+\tfrac12(2mR^2/5)[5v_0/(7R)]^2=5mv_0^2/14$，
比初值 $mv_0^2/2$ 少 $mv_0^2/7$。损失来自接触点滑动，不能用质心路程直接乘摩擦力代表总损失。
\end{xesolution}

\Needspace{8\baselineskip}
\begin{origexamplebox}{【台球碰撞】}\label{ch09:ex:s17-15-4}
质量均为 $m$、半径均为 $R$ 的匀质白球和红球发生无摩擦正碰。碰前白球以速度 $v_0$
向前运动，同时具有角速度 $\omega_0$，且 $v_0=\omega_0R$。红球起初静止且不自转。
两球所在的同一水平面与球的动摩擦因数均为$\mu$，忽略滚动阻力；球间碰撞极短，可忽略该期间的地面冲量。
球间碰撞完全弹性。求碰撞后两个
球的运动，并判断白球是否可能再次撞到红球。


\ThQBilliards
\end{origexamplebox}

\Needspace{6\baselineskip}\begin{xesolution}
\ballsCollisionDiagram

图~\ref{ch09:fig:balls}标出碰前状态，取向右与顺时针为正，碰后另以t=0计时。
球与球之间无摩擦，所以碰撞冲量沿两球心连线，不对任一球心产生力矩。因而碰撞过程中
两球角速度不变：
\[
 \omega_{\mathrm W}'=\omega_0,
 \qquad \omega_{\mathrm R}'=0.
\]
沿球心连线方向，这是相同质量的一维完全弹性碰撞，二球交换平动速度：
\[
 v_{\mathrm W}'=0,
 \qquad v_{\mathrm R}'=v_0.
\]
碰后白球仍自转，接触地面的点相对地面向后滑，所以地面对白球的滑动摩擦力向前；红球只
平动不转动，摩擦力向后。分别应用例~\ref{ch09:ex:s17-15-3} 的滑动阶段公式：
\[
 v_{\mathrm W}(t)=\mu gt,
 \qquad
 \omega_{\mathrm W}(t)=\omega_0-\frac{5\mu g}{2R}t,
\]
\[
 v_{\mathrm R}(t)=v_0-\mu gt,
 \qquad
 \omega_{\mathrm R}(t)=\frac{5\mu g}{2R}t.
\]
两球分别达到纯滚动的时刻相同：
\[
 t_{\mathrm W}=t_{\mathrm R}=\frac{2v_0}{7\mu g}.
\]
且最终纯滚动速度为
\[
 \boxed{v_{\mathrm W}=\frac27v_0},
 \qquad
 \boxed{v_{\mathrm R}=\frac57v_0}.
\]
在整个滑动阶段 $v_{\mathrm R}-v_{\mathrm W}=v_0-2\mu gt>0$；两球进入纯滚动后红球仍更快，
所以间距单调增大，\boxed{\text{白球不会再次撞到红球}}。
更明确地，滑动阶段两球心间距减去2R的净间隙为
\[
 d(t)=\int_0^t(v_{\mathrm R}-v_{\mathrm W})\dd t
 =v_0t-\mu gt^2,\qquad 0\le t\le t_*.
\]
因此 $d(t_*)=10v_0^2/(49\mu g)>0$；之后间隙以 $3v_0/7$ 的恒定速率增大。
当$\mu$趋零时，滑动调节时间趋无穷，但球间碰后红球前进、白球只自转，仍不会追及。
\end{xesolution}

\Needspace{8\baselineskip}
\begin{origexamplebox}{【小明拍球】}\label{ch09:ex:s17-15-5}
小明把一个篮球快速旋转地拍下。落地前篮球速度大小为 $v_0$、方向竖直向下，角速度为
$\omega_0$，转向如图。篮球视为质量 $m$、半径 $R$ 的空心球壳，与地面的滑动摩擦因数为
$\mu$，采用法向恢复系数为1、切向库仑摩擦的理想碰撞模型；一旦切向相对速度归零，
接触不再滑动，不考虑切向弹性回弹。求篮球弹起后的速度方向。
\bounceSpinProblemDiagram
\end{origexamplebox}

\Needspace{6\baselineskip}\begin{xesolution}
图~\ref{ch09:fig:bounce-analysis}将碰前状态、接触冲量与碰后方向分开。碰撞极短，重力冲量可忽略。
\bounceSpinDiagram
法向恢复系数为1意味着竖直速度反向而大小不变，因此
\[
 v_y'=v_0 \quad\text{（竖直向上）}.
\]
图中的逆时针自旋使落地点相对地面向右滑，地面对球的切向冲量 $J_t$ 向左。
先假设整个接触阶段都保持向右滑动，此时动摩擦冲量为
\[
 J_t=\mu J_n.
\]
法向冲量由竖直动量改变给出
\[
 J_n=m(v_0-(-v_0))=2mv_0,
\]
故水平速度
\[
 \boxed{v_x'=-\frac{J_t}{m}=-2\mu v_0}.
\]
于是弹起方向相对竖直线向左偏，偏角 $\alpha$ 满足
\[
 \boxed{\tan\alpha=\frac{|v_x'|}{v_y'}=2\mu}.
\]
同时空心球壳 $I=\tfrac23mR^2$，其自旋改变为
\[
 \omega' =\omega_0-\frac{J_tR}{I}
 =\omega_0-\frac{3\mu v_0}{R}.
\]
若碰撞全程保持滑动，即
\[
 \omega_0R>5\mu v_0.
\]
则以上结果成立。这个临界条件来自碰后接触点速度
$v_x'+\omega'R=\omega_0R-5\mu v_0$ 仍为正。
若 $\omega_0R\leq5\mu v_0$，摩擦足以使接触点在分离前或恰在分离时停止滑动。设所需切向
冲量为 $J_t$，由
\[
 v_x'+\omega'R
 =-\frac{J_t}{m}+\left(\omega_0-\frac{J_tR}{I}\right)R=0
\]
得
\[
 J_t=\frac{2}{5}m\omega_0R,
 \qquad
 v_x'=-\frac25\omega_0R.
\]
因此完整的弹起方向为
\[
 \boxed{
 \tan\alpha=
 \begin{cases}
  2\mu, & \omega_0R>5\mu v_0,\\[2pt]
  \dfrac{2\omega_0R}{5v_0}, & \omega_0R\leq5\mu v_0.
 \end{cases}}
\]
两种情形都向图中摩擦冲量所指的一侧偏转；第二种情形中 $\omega'=2\omega_0/5$，
与 $v_x'=-2\omega_0R/5$ 正好使接触点相对速度为零。
临界点两式相同；$\mu$=0时水平冲量为零，仍竖直弹起；$\omega_0=0$ 时也无侧向偏转。
法向动能恢复不等于总动能恢复。切向平动与转动的能量改变为
\[
 \Delta K_t=-J_t\omega_0R+\frac{J_t^2}{2}\left(\frac1m+\frac{R^2}{I}\right)\le0,
\]
在本题允许的冲量范围内非正，说明摩擦耗能与题设法向弹性条件相容。
\end{xesolution}

\Needspace{8\baselineskip}
\begin{origexamplebox}{【生熟鸡蛋】}\label{ch09:ex:s17-15-6}
把不同鸡蛋放在同一个斜面上向下滚，如何利用物理方法区分生鸡蛋和熟鸡蛋？


\ThQEggs
\end{origexamplebox}

\Needspace{6\baselineskip}\begin{xesolution}
\eggRollingDiagram

图~\ref{ch09:fig:egg}把两只鸡蛋放在相同坡角的两个面板中，接触位置与释放高度相同；
用剖面区分内部凝固与内部液体。比较前还要尽量控制外形、质量、放置朝向和坡面的粗糙程度。
熟鸡蛋内部凝固，外壳与内部近似共同转动；生鸡蛋内部液体可以相对外壳流动，不能给全部质量
指定同一角速度。这个内部响应差异才是判断依据，而不是预先认定某一种一定先到。

先用圆形、质心位于几何中心的模型解释转动惯量的影响。如果外形在滚动截面近似为圆，
且凝固整体保持纯滚动，刚体近似给出
\[
 a_{\rm cooked}=\frac{g\sin\theta}{1+I/(mR^2)}.
\]
真实鸡蛋不是圆球，其接触半径和质心高度随姿态改变，上式只能作为圆形近似，不能精确描述
每个瞬间。生鸡蛋还同时存在两种不同影响：内部液体未充分随壳转动时，参与转动的有效惯量可能较小；
液体相对流动又通过黏滞作用耗散能量。两种影响不能只取其一。

例如取圆形壳内液体几乎不受黏性转矩的理想极限，平动质量为总质量m，参与转动的惯量主要是
外壳的 $I_{\rm shell}$，则可能有
\[
 a_{\rm liquid}\simeq\frac{g\sin\theta}{1+I_{\rm shell}/(mR^2)}.
\]
它反而可大于凝固整体的加速度。这不是对所有真实鸡蛋的预测，而是说明“存在黏滞耗能”
不能单独推出“生鸡蛋一定更慢”。在较强耦合极限，液体逐渐跟随外壳，又更接近共同转动；
真实结果还受形状和滚动姿态影响。

因此，斜面实验可记录滚动时间、姿态变化以及短暂改变外壳运动后的响应，并与已知生熟样品比较；
只给本题的条件，不能仅凭一次到达先后作无条件判别。为检验内部是否仍在运动，可另在水平桌面
将蛋旋转，短暂按停外壳后立即松手：生蛋内部液体若仍在转，会通过黏滞作用重新带动外壳；
熟蛋整体已停住，通常不会自行恢复转动。按住过久会让内部也停止，所以必须说明“短暂”。
这个补充实验利用内部相对运动，不能反过来当作斜面滚动谁先到底的证明。
\end{xesolution}

\Needspace{8\baselineskip}
\begin{origexamplebox}{【重滑轮】}\label{ch09:ex:s17-15-7}
两个重物挂在滑轮两端，绳与定滑轮间不相对滑动。定滑轮质量为 $M$、半径为 $R$，质量均匀
分布。求松手时 $m_1$ 的加速度，设 $m_1>m_2$。


\ThQPulley
\end{origexamplebox}

\Needspace{6\baselineskip}\begin{xesolution}
\massivePulleyDiagram

图~\ref{ch09:fig:pulley}中绳绕轮缘且连到两重物的中心吊点，支架连接固定轮轴。
设 $m_1$ 向下、$m_2$ 向上加速度大小均为 $a$，两侧张力分别为 $T_1,T_2$：
\[
 m_1g-T_1=m_1a,
 \qquad
 T_2-m_2g=m_2a.
\]
均匀圆盘滑轮的转动惯量 $I=\tfrac12MR^2$。绳不打滑给出 $\beta=a/R$，滑轮转动方程为
\[
 (T_1-T_2)R=I\frac aR.
\]
联立消去两张力：
\[
 (m_1-m_2)g
 =\left(m_1+m_2+\frac I{R^2}\right)a.
\]
因此
\[
 \boxed{a=\frac{(m_1-m_2)g}{m_1+m_2+M/2}},
\]
方向为 $m_1$ 向下。滑轮有转动惯性，所以加速度小于把滑轮视为轻滑轮时的结果。
\end{xesolution}

\Needspace{8\baselineskip}
\begin{origexamplebox}{【力矩做功】}\label{ch09:ex:s17-15-8}
两个质量均为 $m$ 的质点固定在长为 $2l$ 的轻杆上，一个在最远端，另一个在中点。杆的
另一端铰接在天花板上。初始时杆水平，由静止释放。求杆与水平呈 $\theta$ 角时：
\begin{enumerate}[label=(\arabic*),leftmargin=2.2em]
  \item 杆的角速度；
  \item 1 号、2 号质点的动能；
  \item 1 号、2 号质点重力势能的减小量；
  \item 说明本题与力矩做功的关系。
\end{enumerate}


\ThQRodWork
\end{origexamplebox}

\Needspace{6\baselineskip}\begin{xesolution}
\torqueWorkDiagram

图~\ref{ch09:fig:torque-work}中$\theta$从水平参考线向下量，取第一次下落的 $0\le\theta\le\pi/2$。
两质点到转轴的距离分别为 $2l$ 与 $l$，系统转动惯量
\[
 I=m(2l)^2+ml^2=5ml^2.
\]
杆由水平转到向下倾斜 $\theta$ 时，两质点分别下降 $2l\sin\theta$ 与 $l\sin\theta$，
故机械能守恒给出
\[
 \frac12I\omega^2=mg(2l\sin\theta)+mg(l\sin\theta)=3mgl\sin\theta.
\]
所以
\[
 \boxed{\omega=\sqrt{\frac{6g\sin\theta}{5l}}}.
\]
两质点动能分别为
\[
 \boxed{E_{k1}=\frac12m(2l\omega)^2
 =\frac{12}{5}mgl\sin\theta},
\]
\[
 \boxed{E_{k2}=\frac12m(l\omega)^2
 =\frac35mgl\sin\theta}.
\]
两者重力势能减小量分别为
\[
 \boxed{-\Delta E_{p1}=2mgl\sin\theta},
 \qquad
 \boxed{-\Delta E_{p2}=mgl\sin\theta}.
\]
个别质点的动能增加并不等于它自己的重力势能减少，因为轻杆的内力在两质点之间传递能量；
但对整个刚体，支点不做功，总动能增加等于总重力势能减少。

从力矩做功看，重力对支点的合力矩为
\[
 M(\phi)=3mgl\cos\phi,
\]
故
\[
 W=\int_0^\theta M(\phi)\dd\phi
 =3mgl\sin\theta=\Delta E_k.
\]
这正是转动形式的动能定理。
\end{xesolution}

\Needspace{8\baselineskip}
\begin{origexamplebox}{【打击中心】}\label{ch09:ex:s17-15-9}
一根质量均匀、质量为 $m$、长度为 $L$ 的杆悬挂在天花板上。用冲量 $J$ 水平打击杆上距
下端 $a$ 的位置，发现悬挂端没有受到冲量。求打击位置。


\ThQPercussion
\end{origexamplebox}

\Needspace{6\baselineskip}\begin{xesolution}
\percussionDiagram

图~\ref{ch09:fig:percussion}中O是固定铰点，杆初始静止；打击点P距悬点为 $x=L-a$。
悬点冲量为零，杆质心获得的线动量完全来自打击冲量：
\[
 mv_C=J.
\]
对悬点取矩，悬点冲量没有力矩，角动量定理给出
\[
 I_O\omega=Jx,
 \qquad I_O=\frac13mL^2.
\]
刚体绕悬点转动，质心速度还满足
\[
 v_C=\omega\frac L2.
\]
代入前两式：
\[
 \frac Jm=\frac{Jx}{I_O}\frac L2,
\]
从而
\[
 \boxed{x=\frac{2L}{3}},
 \qquad
 \boxed{a=L-x=\frac L3}.
\]
即打击点距悬点 $2L/3$（距下端 $L/3$）。该点称为打击中心；在此处横向打击，悬点瞬间
不承受冲量。
\end{xesolution}

\Needspace{8\baselineskip}
\begin{origexamplebox}{【小明的疑惑】}\label{ch09:ex:s17-15-10}
一根质量m、长度L、质量均匀分布的棍靠两端A、B的支撑点水平悬空放置。小明只分析左半根棍时，发现外力矩似乎
不能平衡。这是怎么回事？
\supportedRodProblemDiagram
\end{origexamplebox}

\Needspace{6\baselineskip}\begin{xesolution}
图~\ref{ch09:fig:rod-analysis}左图以整杆为研究对象，对称性与竖直力平衡给出两端支持力各为 $mg/2$，
合力与对任一点的合力矩均为零。
\supportedRodDiagram

但若只截取左半杆，除了左支点支持力 $mg/2$ 与左半杆重力 $mg/2$ 之外，截面还受到右半杆
施加的\textbf{内力和内力矩}。在本题中，中截面的水平、竖直力平衡分别给出轴力为零及
\[
 N_A-\frac{mg}{2}+Q=0\quad\Longrightarrow\quad Q=0.
\]
所以右图没有画一个实际非零的剪力箭头。两竖直外力大小相等、方向相反，但作用线不同，构成力偶。设杆长
为 $L$，左半杆重心距左支点 $L/4$，其外力偶矩大小为
\[
 M_{\rm ext}=\frac{mg}{2}\frac L4=\frac{mgL}{8}.
\]
以逆时针为正，左支持力与左半杆重力对中截面的合外力矩为
$-(mg/2)(L/2)+(mg/2)(L/4)=-mgL/8$。为保持静止，中截面必须提供逆时针弯矩，大小
\[
 \boxed{M_{\rm int}=\frac{mgL}{8}}
\]
的反向弯矩。

疑惑来自把刚体截开后只补了截面内力，却漏掉截面内力矩。对能承受弯曲的杆件，截面相互
作用一般包括轴力、剪力和弯矩；只有把它们全部画出，局部受力与力矩才能平衡。
\end{xesolution}

\Needspace{8\baselineskip}
\begin{origexamplebox}{【烟囱倒地】}\label{ch09:ex:s17-15-11}
把烟囱视为质量 $m$、长度 $L$ 的均匀硬杆，绕底端倾倒，杆与竖直方向夹角为 $\theta$。
用两种模型判断哪里最容易断裂：
\begin{enumerate}[label=(\arabic*),leftmargin=2.2em]
  \item 烟囱内部会有垂直于烟囱的剪力，剪力最大处最先断裂；
  \item 烟囱内部会提供弯矩，弯矩最大处最先断裂。
\end{enumerate}


\ThQChimney
\end{origexamplebox}

\Needspace{6\baselineskip}\begin{xesolution}
\chimneyDiagram

图~\ref{ch09:fig:chimney}中 $x$ 从固定基底O沿杆向上量，研究截面P以上的一段；
规定切向正方向 $\uvec e_t$ 为倾倒的顺时针方向，角速度、角加速度也以顺时针为正。
线密度 $\lambda=m/L$。整杆绕底端转动惯量
$I_O=mL^2/3$。重力矩给出角加速度
\[
 \beta=\frac{mg(L/2)\sin\theta}{mL^2/3}
 =\frac{3g\sin\theta}{2L}.
\]
从竖直静止位置以任意小扰动启动的极限下，机械能守恒给出
\[
 \frac12I_O\omega^2=mg\frac L2(1-\cos\theta),
 \qquad
 \omega^2=\frac{3g}{L}(1-\cos\theta).
\]

在 $x$ 处截开，研究上段，长度 $s=L-x$、质量 $m_s=ms/L$，质心位于
$r_c=(L+x)/2$。沿垂直杆方向列质心运动方程。重力的垂杆分量为
$m_sg\sin\theta$，质心切向加速度为 $\beta r_c$。定义截面对上段的剪力为
$-Q\uvec e_t$，Q是有符号量，图中箭头表示Q的正向。因此
$m_sg\sin\theta-Q=m_s\beta r_c$，得
\begin{align*}
 Q(x)&=m_s\bigl(g\sin\theta-\beta r_c\bigr)\\
 &=\frac{mg\sin\theta}{4L^2}(L-x)(L-3x).
\end{align*}
Q为正时剪力与倾倒切向相反，为负时实际方向反转；剪力\textbf{大小}是 $|Q|$。令 $y=x/L$，$|Q|$ 正比于
$|(1-y)(1-3y)|$；在 $0\le y\le1$ 上最大值出现在 $y=0$，所以按模型一
\[
 \boxed{\text{基底附近最容易发生剪切断裂}}.
\]

再对上段质心取矩。重力通过上段质心，不产生力矩；截面剪力对上段质心产生力矩。为使上段
具有顺时针角加速度$\beta$，截面还必须提供弯矩。关于上段质心G，剪力 $-Q\uvec e_t$ 的
顺时针力矩是 $Qs/2$。定义截面弯矩的有符号量M以顺时针为正，则
\[
 M+Q\frac s2=I_c\beta,\qquad I_c=\frac{m_ss^2}{12}.
\]
移项并代入，可得
\begin{align*}
 M(x)&=I_c\beta-Q(x)\frac s2\\
 &=\frac{mg\sin\theta}{8L^2}(L-x)^3
   -\frac{mg\sin\theta}{8L^2}(L-x)^2(L-3x)\\
 &=\frac{mg\sin\theta}{4L^2}\,x(L-x)^2.
\end{align*}
\end{xesolution}

\begin{solutioncontinuation}
\enlargethispage{\baselineskip}
令 $y=x/L$，需最大化 $y(1-y)^2$：
\[
 \frac{\dd}{\dd y}\bigl[y(1-y)^2\bigr]
 =(1-y)(1-3y)=0.
\]
非端点极值为 $y=1/3$，所以按模型二
\[
 \boxed{x=\frac L3\quad\text{处弯矩最大，最容易弯折断裂}}.
\]
对 $0<\theta<\pi$、$0<x<L$，M为正，与图中的顺时针约定一致；两端M均为零。
Q具有力的量纲，M具有力乘长度的量纲，正负关系也满足上段转动方程。
值得注意的是，$M(x)$ 中的 $\sin\theta$ 只改变总体大小，不改变最大位置；以上结果针对题目
给定的理想均匀刚杆模型。
\end{solutioncontinuation}
}

\endgroup

\hypersetup{bookmarksdepth=2}

% ThermalQualityRevision20261023-v1
\hypersetup{bookmarksdepth=3}
\begingroup
\tcbset{breakable/.default=false}
\chapter{简谐振动、阻尼、受迫振动和共振}
% Editable TikZ, including frequency-response and weak-nonlinearity analysis.
\providecommand{\ThermalFigure}[3]{\begin{qualityfigure}
\begin{tikzpicture}[>=Latex,font=\small,x=.8cm,y=.8cm]#1\end{tikzpicture}
\figcap{#2}\label{#3}\end{qualityfigure}}
\newcommand{\ThQVerticalSpring}{\ThermalFigure{
 \draw[thick](-1,3)--(1,3);\draw[decorate,decoration={coil,segment length=3mm,amplitude=2mm},thick](0,3)--(0,.5);
 \draw[fill=DeepBlue!15,thick](-.4,-.2)rectangle(.4,.5);\node at(0,.15){$m$};
 \draw[dashed](-1.1,1)--(1,1)node[right]{原长末端};\draw[<->](1.5,1)--node[right]{伸长量}(1.5,.5);
 \node[left]at(-.5,1.9){$k$};\draw[->](3,1.8)--(3,.5)node[right]{$g$};
}{竖直轻弹簧与小球；须区分从原长量起的伸长与从平衡位置量起的位移。}{ch10:quality:vertical-q}}
\newcommand{\ThQPhysicalPendulum}{\ThermalFigure{
 \fill[DeepBlue!10](.5,1.3)ellipse(1 and 1.8);\draw[thick](.5,1.3)ellipse(1 and 1.8);
 \coordinate(O)at(0,3);\coordinate(C)at(.6,.9);\draw[thick](O)--(C);
 \fill(O)circle(2pt)node[above]{$O$};\fill(C)circle(2pt)node[right]{$C$};
 \draw[dashed](0,3)--(0,.2);\draw[<->](-.5,2.9)--node[left]{$a$}(.1,.8);
 \draw(0,2.2)arc(270:286:.8)node[below]{$\theta$};\node[right,align=left]at(3,1.7){绕 $O$ 转动惯量 $I$\\质心 $C$ 位于轴下方};
}{任意刚体的复摆；小角度是质心连线与向下铅垂线的夹角。}{ch10:quality:pendulum-q}}
\newcommand{\ThBandwidthBase}[1]{
 \draw[->](0,0)--(7,0)node[right]{$\omega$};\draw[->](0,0)--(0,3.2)node[above]{$\overline P$};
 \draw[DeepBlue,thick,domain=.4:6.6,samples=100]plot(\x,{2.8/(1+(\x-3.5)^2/.64)});
 \draw[dashed](0,1.4)--(6.5,1.4);\node[left]at(0,1.4){$\overline P_{\max}/2$};
 \draw[dashed](2.7,0)--(2.7,1.4) (4.3,0)--(4.3,1.4);
 \node[below]at(2.7,0){$\omega_1$};\node[below]at(4.3,0){$\omega_2$};#1}
\newcommand{\ThQBandwidth}{\ThermalFigure{\ThBandwidthBase{}}{半功率带宽由功率曲线与半高线的两个交点定义；图为弱阻尼近峰示意。}{ch10:quality:bandwidth-q}}
\newcommand{\ThABandwidth}{\ThermalFigure{\ThBandwidthBase{
 \draw[<->](2.7,-.7)--node[below]{$\Delta\omega$}(4.3,-.7);\draw[dashed](3.5,0)--(3.5,2.8);\node[above]at(3.5,2.9){功率峰 $\omega_0$};
}}{先对功率施加半高条件，再比较两个正频率根。}{ch10:quality:bandwidth-a}
图~\ref{ch10:quality:bandwidth-a}中的水平线表示功率减半，并非位移振幅减半。对于线性黏性阻尼模型，它给出 $|k-m\omega^2|=b\omega$；左右交点取绝对值的不同分支，故两个二次方程的一次项符号相反。
}
\newcommand{\ThHardeningBase}[1]{
 \draw[->](-2.7,0)--(2.7,0)node[right]{$x$};\draw[->](0,-2.3)--(0,2.3)node[above]{$F$};
 \draw[TextGray,dashed](-2,1)--(2,-1);
 \draw[DeepBlue,thick,domain=-2:2,samples=70]plot(\x,{-.5*\x-.12*\x^3});
 \node[left]at(-2,2){$-kx-\lambda x^3$};\node[right]at(2,-.8){$-kx$};#1}
\newcommand{\ThQHardening}{\ThermalFigure{\ThHardeningBase{}}{硬化回复力比线性模型更强；初态为正转折点 $x=A$。}{ch10:quality:hardening-q}}
\newcommand{\ThAHardening}{\ThermalFigure{
 \draw[->](-2.6,0)--(2.6,0)node[right]{$x$};\draw[->](0,0)--(0,3.1)node[above]{$U$};
 \draw[TextGray,dashed,domain=-2:2,samples=70]plot(\x,{.3*(\x)*(\x)});
 \draw[DeepBlue,thick,domain=-2:2,samples=70]plot(\x,{.3*(\x)*(\x)+.075*(\x)*(\x)*(\x)*(\x)});
 \draw[dashed](-2,2.4)--(2,2.4);\node[below]at(-2,0){$-A$};\node[below]at(2,0){$A$};
 \node[right,align=left]at(3,1.8){实线：$U=kx^2/2+\lambda x^4/4$\\虚线：线性弹簧势能};
}{正四次势能使势阱在大位移处更陡，转折点仍由总能量与势能相等确定。}{ch10:quality:hardening-a}
图~\ref{ch10:quality:hardening-a}中的正四次项使回复作用增强，预期频率随振幅增大而升高。作一阶估计时，$\cos^3\Omega t=(3\cos\Omega t+\cos3\Omega t)/4$ 的基频部分改变频率，三倍频部分只在同阶修正波形；这给出相对频移的系数 $3/8$。
}
\newcommand{\ThQSphereCharges}{\ThermalFigure{
 \draw[thick](0,0)circle(2);\draw[dashed](0,0)ellipse(2 and .55);
 \foreach \x/\y/\n in {1.5/1.323/1,-1.8/.872/2,.5/-1.936/3}{\fill[DeepBlue](\x,\y)circle(3pt)node[right]{$\n:Q,m$};}
 \draw[->](0,0)--(-1,-1.732)node[midway,right]{$R$};\node[right]at(4,0){各粒子限于球面；初图不预设平衡形状};
}{约束球面上的三个同号点电荷；平衡构形及简正模由势能极小与二阶展开求出。}{ch10:quality:charges-q}}
\newcommand{\ThQSpringPendulum}{\ThermalFigure{
 \coordinate(S)at(0,3);\coordinate(O)at(0,0);\coordinate(P)at(1.5,-.8);
 \draw[thick](-.7,3)--(.7,3);\draw[dashed](S)--(O);
 \draw[decorate,decoration={coil,segment length=3mm,amplitude=2mm},thick](S)--(P);
 \fill[DeepBlue](P)circle(3pt)node[right]{$m$};\fill(O)circle(1pt)node[above left]{$O$（平衡）};
 \draw[->](O)--(2.7,0)node[right]{$x$};\draw[->](O)--(0,-1.7)node[below]{$y$};
 \draw[dashed](P)--(1.5,0) (P)--(0,-.8);\node[left]at(0,1.5){$l_e$};
}{竖直平衡长度为 $l_e$，原长为 $l_0$；展开坐标从平衡位置而非原长末端量起。}{ch10:quality:spring-pendulum-q}}



% Chapter-local, deterministic geometry. Text nodes never have opaque fill.
\newenvironment{chTenFigure}{\par\begin{center}\begin{minipage}{.98\linewidth}\centering}{\end{minipage}\end{center}\par}
\tikzset{chTenLabel/.style={font=\small,fill=none,inner sep=3pt},chTenArrow/.style={-{Latex[length=2mm]},DeepBlue,thick}}
\pgfplotsset{chTenPlot/.style={axis lines=left,tick label style={font=\scriptsize},label style={font=\small,fill=none},legend style={draw=none,fill=none,font=\scriptsize},title style={font=\small},grid=none}}

\newcommand{\chTenPhaseFigure}{\begin{chTenFigure}
\begin{tikzpicture}[x=1cm,y=1cm,>=Latex,every node/.style=chTenLabel]
\draw[->](-3.4,0)--(3.4,0)node[right]{$x$};\draw[->](0,-2.3)--(0,2.35)node[above]{$v$};
\draw[DeepBlue,thick,domain=0:360,samples=121]plot({2.6*cos(\x)},{1.55*sin(\x)});
\draw[chTenArrow,domain=75:45,samples=16]plot({2.6*cos(\x)},{1.55*sin(\x)});
\draw[chTenArrow,domain=245:215,samples=16]plot({2.6*cos(\x)},{1.55*sin(\x)});
\node[below right]at(2.68,-.08){$A$};\node[below left]at(-2.68,-.08){$-A$};
\node[above left]at(-.08,1.63){$A\omega_0$};\node[below right]at(.08,-1.63){$-A\omega_0$};
\node at(4.1,1.1){\shortstack{$x=A\cos\psi$\\$v=-A\omega_0\sin\psi$}};
\end{tikzpicture}
\figcap{无阻尼相轨迹沿顺时针方向运动；在顶部$v>0$，状态点向右移动。}\label{ch10:fig:phase}
\end{chTenFigure}}

\newcommand{\chTenVerticalFigure}{\begin{chTenFigure}
\begin{tikzpicture}[every node/.style=chTenLabel,>=Latex]
\draw[ground](-1,0)--(6.5,0);
\foreach \p/\h/\txt in {0/1.65/原长参照,2.5/2.4/平衡位置,5/3.05/瞬时位置}{
 \draw[spring](\p,0)--(\p,-\h);\fill[DeepBlue](\p,-\h)circle(3pt);
 \node[below]at(\p,-3.55){\txt};}
\draw[guide](0,-1.65)--(6.3,-1.65);\draw[guide](2.5,-2.4)--(6.3,-2.4);
\draw[<->](3.45,-1.65)--(3.45,-2.4)node[midway,right]{$y_0$};
\draw[<->](6,-2.4)--(6,-3.05)node[midway,right]{$x$};
\node[right]at(5.1,-3.1){$m$};
\end{tikzpicture}\figcap{伸长$y=y_0+x$；辅助水平线只用于比较三个位置。}\label{ch10:fig:vertical-spring}\end{chTenFigure}}

\newcommand{\chTenPendulumFigure}{\begin{chTenFigure}
\begin{tikzpicture}[every node/.style=chTenLabel,>=Latex]
\coordinate(O)at(0,0);\coordinate(C)at({2.5*sin(28)},{-2.5*cos(28)});
\draw[guide](O)--(0,-.18) (0,-.55)--(0,-3.3);\draw[thick](O)--(C);
\draw[DeepBlue,thick,rotate around={28:(O)}](-.48,-.7)rectangle(.48,-2.95);
\fill(O)circle(2pt)node[above left]{$O$};\fill(C)circle(2pt)node[right]{$C$};
\draw[chTenArrow](C)--++(0,-1.0)node[right]{$Mg$};
\draw(0,-.5)arc(-90:-62:.5);\node[left]at(-.55,-.25){$\theta$};\draw[TextGray](-.5,-.25)--(.1,-.49);
\draw[<->](-.78,-.415)--($(C)+(-.78,-.415)$)node[midway,left]{$a$};\node[right]at(2.1,-.7){转轴垂直纸面};
\end{tikzpicture}\figcap{复摆绕固定轴$O$转动，$OC=a$；重力作用于质心$C$。}\label{ch10:fig:physical-pendulum}\end{chTenFigure}}

\newcommand{\chTenDampingFigure}{\begin{chTenFigure}
\begin{tikzpicture}
\begin{axis}[chTenPlot,width=.31\linewidth,height=4.7cm,xmin=0,xmax=9,ymin=-.7,ymax=1.1,xlabel={$\omega_0t$},ylabel={$x/x_0$},title={欠阻尼：$\gamma/\omega_0=.2$}]
\addplot[DeepBlue,thick,domain=0:9,samples=140]{exp(-.2*x)*(cos(deg(sqrt(.96)*x))+.2/sqrt(.96)*sin(deg(sqrt(.96)*x)))};\end{axis}
\begin{axis}[chTenPlot,at={(.335\linewidth,0)},anchor=south west,width=.31\linewidth,height=4.7cm,xmin=0,xmax=9,ymin=-.7,ymax=1.1,xlabel={$\omega_0t$},title={临界：$\gamma/\omega_0=1$}]
\addplot[DeepBlue,thick,domain=0:9,samples=100]{(1+x)*exp(-x)};\end{axis}
\begin{axis}[chTenPlot,at={(.67\linewidth,0)},anchor=south west,width=.31\linewidth,height=4.7cm,xmin=0,xmax=9,ymin=-.7,ymax=1.1,xlabel={$\omega_0t$},title={过阻尼：$\gamma/\omega_0=2$}]
\addplot[DeepBlue,thick,domain=0:9,samples=100]{.5*(1+2/sqrt(3))*exp((-2+sqrt(3))*x)+.5*(1-2/sqrt(3))*exp((-2-sqrt(3))*x)};\end{axis}
\end{tikzpicture}\figcap{同一初始位移、零初速度下的三类响应；坐标标度完全相同。}\label{ch10:fig:damping-types}\end{chTenFigure}}

\newcommand{\chTenEnvelopeFigure}{\begin{chTenFigure}
\begin{tikzpicture}\begin{axis}[chTenPlot,width=.86\linewidth,height=5.4cm,xmin=0,xmax=30,ymin=-1.1,ymax=1.2,xlabel={$\omega_0t$},ylabel={$x/A_0$},legend pos=north east]
\addplot[DeepBlue,thick,domain=0:30,samples=350]{exp(-.08*x)*cos(deg(sqrt(1-.08^2)*x))};\addlegendentry{位移}
\addplot[TextGray,dashed,domain=0:30]{exp(-.08*x)};\addlegendentry{$\pm\mathrm e^{-\gamma t}$}
\addplot[TextGray,dashed,domain=0:30,forget plot]{-exp(-.08*x)};
\end{axis}\end{tikzpicture}\figcap{弱阻尼示例$\gamma/\omega_0=0.08$；虚线给振幅包络，不是另一条运动轨迹。}\label{ch10:fig:damping-envelope}\end{chTenFigure}}

\newcommand{\chTenDampedPhaseFigure}{\begin{chTenFigure}
\begin{tikzpicture}
\begin{axis}[chTenPlot,width=.46\linewidth,height=6cm,axis equal image,xmin=-1.2,xmax=1.2,ymin=-1.2,ymax=1.2,xlabel={$x/A_0$},ylabel={$v/(A_0\omega_0)$},title={真实相轨迹：$\gamma/\omega_0=.2$}]
\addplot[DeepBlue,thick,domain=0:28,samples=450]({exp(-.2*x)*cos(deg(sqrt(.96)*x))},{-exp(-.2*x)*(.2*cos(deg(sqrt(.96)*x))+sqrt(.96)*sin(deg(sqrt(.96)*x))});
\addplot[DeepBlue,->,thick,domain=.5:1,samples=12]({exp(-.2*x)*cos(deg(sqrt(.96)*x))},{-exp(-.2*x)*(.2*cos(deg(sqrt(.96)*x))+sqrt(.96)*sin(deg(sqrt(.96)*x))});\end{axis}
\begin{axis}[chTenPlot,at={(.52\linewidth,0)},anchor=south west,width=.46\linewidth,height=6cm,axis equal image,xmin=-1.2,xmax=1.2,ymin=-1.2,ymax=1.2,xlabel={$X/A_0$},ylabel={$Y/A_0$},title={扣除包络：$\gamma/\omega_0=.5$}]
\addplot[TextGray,dashed,domain=-1:1]{-x};\addplot[DeepBlue,thick,domain=0:360,samples=180]({cos(x)},{-.5*cos(x)-sqrt(.75)*sin(x)});\end{axis}
\end{tikzpicture}\figcap{左图螺旋收缩；右图$Y=\mathrm e^{\gamma t}v/\omega_0$，虚线为长轴方向。左右阻尼值分别标明。}\label{ch10:fig:damped-phase}\end{chTenFigure}}

\newcommand{\chTenResonanceFigure}{\begin{chTenFigure}
\begin{tikzpicture}
\begin{axis}[chTenPlot,width=.46\linewidth,height=6cm,xmin=0,xmax=2,ymin=0,ymax=21,xlabel={$\omega/\omega_0$},ylabel={$Ak/F_0$},legend pos=north east]
\addplot[Ink,dashdotted,thick,domain=0:2,samples=500]{1/sqrt((1-x^2)^2+(.05*x)^2)};\addlegendentry{$Q=20$}
\addplot[DeepBlue,thick,domain=0:2,samples=350]{1/sqrt((1-x^2)^2+(.1*x)^2)};\addlegendentry{$Q=10$}
\addplot[Blue,dashed,thick,domain=0:2,samples=220]{1/sqrt((1-x^2)^2+(.2*x)^2)};\addlegendentry{$Q=5$}
\addplot[TextGray,densely dotted,thick,domain=0:2,samples=120]{1/sqrt((1-x^2)^2+(2*x)^2)};\addlegendentry{$Q=0.5$}\end{axis}
\begin{axis}[chTenPlot,at={(.52\linewidth,0)},anchor=south west,width=.46\linewidth,height=6cm,xmin=0,xmax=2,ymin=0,ymax=3.3,xlabel={$\omega/\omega_0$},ylabel={$\delta$},ytick={0,1.570796,3.141593},yticklabels={$0$,$\pi/2$,$\pi$}]
\addplot[Ink,dashdotted,thick,domain=0:2,samples=500]{rad(atan2(.05*x,1-x^2))};
\addplot[DeepBlue,thick,domain=0:2,samples=350]{rad(atan2(.1*x,1-x^2))};
\addplot[Blue,dashed,thick,domain=0:2,samples=220]{rad(atan2(.2*x,1-x^2))};
\addplot[TextGray,densely dotted,thick,domain=0:2,samples=120]{rad(atan2(2*x,1-x^2))};\end{axis}
\end{tikzpicture}\figcap{同组参数的幅频和相频曲线；较强阻尼$Q=0.5$没有正频率的位移峰。}\label{ch10:fig:resonance}\end{chTenFigure}}

\newcommand{\chTenProfileFigure}{\begin{chTenFigure}
\begin{tikzpicture}\begin{axis}[chTenPlot,width=.8\linewidth,height=5.6cm,xmin=-4,xmax=4,ymin=0,ymax=1.12,xlabel={$(z-z_0)/\Gamma$},ylabel={相对峰高},legend pos=north east]
\addplot[DeepBlue,thick,domain=-4:4,samples=160]{1/(1+x^2)};\addlegendentry{洛伦兹型}
\addplot[TextGray,dashed,thick,domain=-4:4,samples=160]{exp(-ln(2)*x^2)};\addlegendentry{高斯型}
\draw[guide](axis cs:-1,0)--(axis cs:-1,.5)--(axis cs:1,.5)--(axis cs:1,0);
\draw[<->](axis cs:-1,.18)--(axis cs:1,.18)node[midway,below,font=\small,fill=none]{全宽$2\Gamma$};
\end{axis}\end{tikzpicture}\figcap{相同峰高、相同半高全宽时的两种线型；峰外长尾明显不同。}\label{ch10:fig:profiles}\end{chTenFigure}}

\newcommand{\chTenCoupledPendulaFigure}{\begin{chTenFigure}
\begin{tikzpicture}[every node/.style=chTenLabel,>=Latex]
\coordinate(A)at(0,0);\coordinate(B)at(4,0);\coordinate(P)at({2.6*sin(18)},{-2.6*cos(18)});\coordinate(Q)at({4+2.6*sin(-12)},{-2.6*cos(-12)});
\draw[ground](-.7,0)--(4.7,0);\draw[thick](A)--(P) (B)--(Q);\draw[spring](P)--(Q);
\draw[guide](A)--(0,-.4) (0,-1)--(0,-3) (B)--(4,-.4) (4,-1)--(4,-3);\fill(P)circle(3pt)node[below left]{$m$};\fill(Q)circle(3pt)node[below right]{$m$};
\node[right]at(.7,-1.45){$l$};\node[left]at(3.4,-1.45){$l$};
\draw(0,-.7)arc(-90:-72:.7);\node[left]at(-.4,-.62){$\theta_1$};\draw[TextGray](-.38,-.62)--(.11,-.69);
\draw(4,-.7)arc(-90:-102:.7);\node[right]at(4.4,-.62){$\theta_2$};\draw[TextGray](4.38,-.62)--(3.93,-.69);
\draw[<->](0,.5)--(4,.5)node[midway,above]{$d$};\node[above]at(2.15,-2.15){$k_c$};
\end{tikzpicture}\figcap{两球间弹簧实际连接在球上；图中第二个角为负，两角正向约定仍一致。}\label{ch10:fig:coupled-pendula}\end{chTenFigure}}

\newcommand{\chTenCoupledMassFigure}{\begin{chTenFigure}
\begin{tikzpicture}[every node/.style=chTenLabel,>=Latex]
\draw[ground](-5,-.6)--(-5,.6) (5,-.6)--(5,.6);
\draw[spring](-5,0)--(-3.2,0) (-1.8,0)--(1.8,0) (3.2,0)--(5,0);
\draw[thick](-3.2,-.5)rectangle(-1.8,.5);\draw[thick](1.8,-.5)rectangle(3.2,.5);
\node at(-2.5,0){$m$};\node at(2.5,0){$m$};
\node at(-4.1,.65){$k$};\node at(0,.65){$k_c$};\node at(4.1,.65){$k$};
\draw[guide](-2.5,-.6)--(-2.5,-1.1) (2.5,-.6)--(2.5,-1.1);
\draw[chTenArrow](-2.5,-.85)--(-1.6,-.85)node[right]{$x_1$};\draw[chTenArrow](2.5,-.85)--(3.4,-.85)node[right]{$x_2$};
\end{tikzpicture}\figcap{图示平衡构形；短虚线为坐标零点，箭头只规定位移正向，不表示此刻的实际位移。}\label{ch10:fig:coupled-masses}\end{chTenFigure}}

\newcommand{\chTenDoublePendulumFigure}{\begin{chTenFigure}
\begin{tikzpicture}[every node/.style=chTenLabel,>=Latex]
\coordinate(O)at(0,0);\coordinate(P)at({2.2*sin(23)},{-2.2*cos(23)});\coordinate(Q)at($(P)+({2.2*sin(-18)},{-2.2*cos(-18)})$);
\draw[guide](O)--(0,-.4) (0,-1)--(0,-2.5) (P)--++(0,-.4) ($(P)+(0,-1)$)--++(0,-1.3);\draw[thick](O)--(P)--(Q);
\fill(O)circle(2pt)node[above]{$O$};\fill(P)circle(3pt)node[right]{$m_1=m$};\fill(Q)circle(3pt)node[below left]{$m_2=m$};
\draw(0,-.7)arc(-90:-67:.7);\node[left]at(-.5,-.62){$\theta_1$};\draw[TextGray](-.48,-.62)--(.14,-.68);
\draw($(P)+(0,-.7)$)arc(-90:-108:.7);\node[right]at($(P)+(.55,-.62)$){$\theta_2$};\draw[TextGray]($(P)+(.5,-.62)$)--($(P)+(-.11,-.69)$);
\node[right]at(.78,-1.35){$l$};\node[left]at($(P)+(-.65,-1.5)$){$l$};
\end{tikzpicture}\figcap{串接双摆。第二杆角从上球处的竖直虚线量起。}\label{ch10:fig:double-pendulum}\end{chTenFigure}}

\newcommand{\chTenDoubleModesFigure}{\begin{chTenFigure}
\begin{tikzpicture}[every node/.style=chTenLabel,>=Latex]
\foreach \dx/\sg/\txt in {0/1/低频模,5/-1/高频模}{\begin{scope}[xshift=\dx cm]
\coordinate(O)at(0,0);\coordinate(P)at({2*sin(12)},{-2*cos(12)});\coordinate(Q)at($(P)+({2*sin(\sg*sqrt(2)*12)},{-2*cos(\sg*sqrt(2)*12)})$);
\draw[guide](O)--(0,-3.95) (P)--++(0,-1.9);\draw[thick](O)--(P)--(Q);\fill(P)circle(3pt);\fill(Q)circle(3pt);
\node[above]at(0,0){\txt};\node[right]at($(P)+(.15,0)$){$\theta_1$};\node[below]at($(Q)+(0,-.2)$){$\theta_2=\ifnum\sg<0-\fi\sqrt2\theta_1$};
\draw[chTenArrow]($(O)+(0,-.65)$)arc(-90:-72:.65);
\draw[chTenArrow]($(P)+(0,-.65)$)arc(-90:{-90+\sg*24}:.65);
\end{scope}}
\end{tikzpicture}\figcap{双摆的两个模态；低频同向、高频反向，下杆角振幅均较大。}\label{ch10:fig:double-modes}\end{chTenFigure}}

\newcommand{\chTenRingBase}{
\draw[TextGray](0,0)circle(1.4);
\foreach \j in {0,...,5}{\pgfmathsetmacro{\ang}{60*\j}\fill[DeepBlue](\ang:1.4)circle(2.2pt);\ifnum\j=0\node[font=\scriptsize,fill=none]at(1.75,.23){$\j$};\else\ifnum\j=3\node[font=\scriptsize,fill=none]at(-1.75,.23){$\j$};\else\node[font=\scriptsize,fill=none]at(\ang:1.75){$\j$};\fi\fi}}
\newcommand{\chTenBenzeneModelFigure}{\begin{chTenFigure}
\begin{minipage}{.46\linewidth}\centering\includegraphics[width=.96\linewidth]{assets/ch10/benzene-model.png}\par\small 苯：六个碳、六个氢\end{minipage}\hfill
\begin{minipage}{.46\linewidth}\centering\begin{tikzpicture}[every node/.style=chTenLabel]
\chTenRingBase
\foreach \j in {0,...,5}{\draw[chTenArrow]({60*\j}:1.4)--++({60*\j+90}:.48);}
\node at(0,-2.25){$x_j$均沿逆时针切向};
\end{tikzpicture}\end{minipage}
\figcap{真实分子的立体结构与六坐标简化模型；左图淡蓝区域表示离域电子云的空间示意，不是额外原子。}\label{ch10:fig:benzene-model}\end{chTenFigure}}

\newcommand{\chTenRingSymmetryFigure}{\begin{chTenFigure}
\begin{tikzpicture}[every node/.style=chTenLabel,>=Latex]
\foreach \dx/\kind in {0/0,4.3/1,8.6/2}{\begin{scope}[xshift=\dx cm,scale=.9]
\chTenRingBase
\ifnum\kind=0\draw[chTenArrow](15:.85)arc(15:75:.85);\node at(0,-2.3){$C_6:60^\circ$};
\else\ifnum\kind=1\draw[guide](30:2)--(210:2);\node at(0,-2.3){$C_2:180^\circ$};
\else\draw[guide](-2,0)--(2,0);\node at(0,-2.3){$C'_2:180^\circ$};\fi\fi
\end{scope}}\end{tikzpicture}\figcap{三种对称操作；后二图虚线是环面内的旋转轴，而不是粒子的运动轨迹。}\label{ch10:fig:ring-symmetry}\end{chTenFigure}}

\newcommand{\chTenOneRingMode}[2]{\begin{tikzpicture}[scale=.75,every node/.style=chTenLabel]
\chTenRingBase
\foreach \v[count=\j from 0] in {#1}{\ifnum\v=0\else\draw[chTenArrow]({60*\j}:1.4)--++({60*\j+90}:{.58*\v});\fi}
\node[font=\small]at(0,-2.45){#2};\end{tikzpicture}}
\newcommand{\chTenRingModesFigure}{\begin{chTenFigure}
\chTenOneRingMode{1,1,0,-1,-1,0}{$m\omega^2/k=1$}\hspace{.8cm}
\chTenOneRingMode{0,1,1,0,-1,-1}{同频旋转伙伴}\par\medskip
\chTenOneRingMode{1,-1,0,1,-1,0}{$m\omega^2/k=3$}\hspace{.8cm}
\chTenOneRingMode{0,1,-1,0,1,-1}{同频旋转伙伴}
\figcap{两组二重简并实振型。相邻点的箭头必须按各自切向比较，不能只按纸面左右比较。}\label{ch10:fig:ring-modes}\end{chTenFigure}}

\newcommand{\chTenMonoChainFigure}{\begin{chTenFigure}\begin{tikzpicture}[every node/.style=chTenLabel,>=Latex]
\foreach \j in {0,...,4}{\fill[DeepBlue](2.1*\j,0)circle(4pt);\node[above]at(2.1*\j,.3){$m$};}
\foreach \j in {0,...,3}{\draw[spring](2.1*\j+.14,0)--(2.1*\j+1.96,0);}
\node[left]at(-.3,0){$\cdots$};\node[right]at(8.7,0){$\cdots$};
\foreach \j/\txt in {1/{j-1},2/{j},3/{j+1}}{\draw[chTenArrow](2.1*\j,-.55)--++(.6,0)node[below]{$u_{\txt}$};}
\draw[<->](2.1,.85)--(4.2,.85)node[midway,above]{$a$};
\end{tikzpicture}\figcap{单原子链的平衡格距$a$及向右位移坐标。}\label{ch10:fig:mono-chain}\end{chTenFigure}}

\newcommand{\chTenMonoDispersionFigure}{\begin{chTenFigure}\begin{tikzpicture}\begin{axis}[chTenPlot,width=.72\linewidth,height=5.5cm,xmin=-1.1,xmax=1.1,ymin=0,ymax=2.25,xlabel={$qa/\pi$},ylabel={$\omega/\sqrt{k/m}$},xtick={-1,0,1},ytick={0,1,2}]
\addplot[DeepBlue,thick,domain=-1:1,samples=201]{2*abs(sin(90*x))};\addplot[TextGray,dashed,domain=-1:1]{2};
\draw[guide](axis cs:-1,0)--(axis cs:-1,2);\draw[guide](axis cs:1,0)--(axis cs:1,2);
\end{axis}\end{tikzpicture}\figcap{单原子链色散；两端达到截止频率，斜率降为零。}\label{ch10:fig:mono-dispersion}\end{chTenFigure}}

\newcommand{\chTenDiatomicChainFigure}{\begin{chTenFigure}\begin{tikzpicture}[every node/.style=chTenLabel,>=Latex]
\foreach \j in {0,...,5}{\pgfmathtruncatemacro{\odd}{mod(\j,2)}\ifnum\odd=0\fill[DeepBlue](1.7*\j,0)circle(5pt);\node[above]at(1.7*\j,.35){$M$};\else\draw[DeepBlue,thick](1.7*\j,0)circle(3pt);\node[above]at(1.7*\j,.35){$m$};\fi}
\foreach \j in {0,...,4}{\draw[spring](1.7*\j+.18,0)--(1.7*\j+1.52,0);}
\draw[<->](3.4,1.45)--(6.8,1.45)node[midway,above]{晶胞$a$};
\draw[<->](3.4,-1.4)--(5.1,-1.4)node[midway,below]{$a/2$};
\draw[chTenArrow](3.4,-.55)--++(.6,0)node[midway,below]{$u_j$};\draw[chTenArrow](5.1,-.55)--++(.6,0)node[midway,below]{$v_j$};
\draw[guide](2.55,-1)rectangle(5.95,.95);
\end{tikzpicture}\figcap{虚框圈出同一晶胞内的一重一轻两原子；相邻距离与晶胞长度不混用。}\label{ch10:fig:diatomic-chain}\end{chTenFigure}}

\newcommand{\chTenDiatomicDispersionFigure}{\begin{chTenFigure}\begin{tikzpicture}\begin{axis}[chTenPlot,width=.8\linewidth,height=6cm,xmin=-1.1,xmax=1.42,ymin=0,ymax=3.5,xlabel={$qa/\pi$},ylabel={$\omega/\sqrt{k/M}$},xtick={-1,0,1},legend pos=south east]
\addplot[DeepBlue,thick,domain=-1:1,samples=150]{sqrt(5-sqrt(25-16*sin(90*x)^2))};\addlegendentry{声学支}
\addplot[TextGray,thick,domain=-1:1,samples=150]{sqrt(5+sqrt(25-16*sin(90*x)^2))};\addlegendentry{光学支}
\draw[guide](axis cs:-1,1.414214)--(axis cs:1.1,1.414214);\draw[guide](axis cs:-1,2.828427)--(axis cs:1.1,2.828427);
\draw[<->](axis cs:1.18,1.414214)--(axis cs:1.18,2.828427)node[midway,right,font=\small,fill=none]{禁带};
\end{axis}\end{tikzpicture}\figcap{双原子链示例$M=4m$。水平虚线对应$\sqrt{2k/M}$与$\sqrt{2k/m}$；区中心上支为$\sqrt{2k(1/M+1/m)}$。}\label{ch10:fig:diatomic-dispersion}\end{chTenFigure}}
\newcommand{\chTenFrictionFigure}{\begin{chTenFigure}\begin{tikzpicture}[every node/.style=chTenLabel,>=Latex]
\draw[ground](-.5,0)--(8,0);\draw[thick](0,0)rectangle(1.25,.8);\draw[thick](6.1,0)rectangle(7.35,.8);
\node at(.625,.4){$A,M$};\node at(6.725,.4){$B,m$};\draw[spring](1.25,.4)--(6.1,.4);
\draw[chTenArrow](.2,1.1)--(1.2,1.1)node[above]{$v_A$};\draw[chTenArrow](7.1,1.1)--(6.1,1.1)node[above]{$v_B$};
\draw[<->](.625,-.55)--(6.725,-.55)node[midway,below]{$r=x_B-x_A$};
\draw[->](8.2,.1)--++(.7,0)node[right]{$x$};
\node[anchor=west]at(0,-1.7){第一阶段：$f_A$向左，$f_B$向右。};
\node[anchor=west]at(0,-2.35){第二阶段：$v_A=0$，$f_A$向右；$v_B<0$，$f_B$仍向右。};
\end{tikzpicture}\figcap{物块形状仅为识别符号，计算中按质点取中心间距；两阶段摩擦方向分别标明。}\label{ch10:fig:friction}\end{chTenFigure}}

\newcommand{\chTenGovernorFigure}{\begin{chTenFigure}\begin{tikzpicture}[every node/.style=chTenLabel,>=Latex]
\coordinate(T)at(0,2);\coordinate(B)at(0,-2);\coordinate(P)at(2,0);\coordinate(Q)at(-2,0);
\draw[guide](0,-3.1)--(0,3.0);\draw[thick](T)--(P)--(B)--(Q)--cycle;
\draw[thick](-.17,1.78)rectangle(.17,2.22);\draw[thick](-.22,-2.23)rectangle(.22,-1.77);
\fill[DeepBlue](P)circle(4pt);\fill[DeepBlue](Q)circle(4pt);
\node[below left]at(-2.13,-.1){$m$};\node[right]at(2.13,0){$m$};\node[left]at(-.35,-2){$M$};\node[above left]at(-.2,2.23){固定套筒};
\draw[guide](.2,2)--(3.35,2) (.2,-2)--(3.35,-2);\draw[<->](3.15,-2)--(3.15,2)node[midway,right]{$y$};
\draw[<->](0,0)--(2,0)node[midway,above]{$\rho$};
\node[above right]at(1.12,1.1){$l$};\node[below right]at(1.12,-1.1){$l$};
\draw[guide](-.2,2)--(-3.15,2) (-2.15,0)--(-3.15,0);
\draw[<->](-3,0)--(-3,2)node[midway,left]{$y/2$};
\draw[chTenArrow,domain=210:345,samples=30]plot({.75*cos(\x)},{-2.75+.2*sin(\x)});\node[right]at(.9,-2.75){$\omega$};
\end{tikzpicture}\figcap{飞球调速器的瞬时子午截面；四根铰臂等长，球与两套筒的高度差均为$y/2$。}\label{ch10:fig:governor}\end{chTenFigure}}

\newcommand{\chTenChainFigure}{\begin{chTenFigure}\begin{tikzpicture}[every node/.style=chTenLabel,>=Latex]
\draw[ground](-1.2,4.5)--(6.5,4.5);\draw[ground](-1.2,0)--(6.5,0);
\draw[spring](0,4.5)--(0,2.5);\fill(0,2.5)circle(2pt)node[left]{$P$};
\foreach \j in {0,...,11}{\draw[DeepBlue,thick](0,{.125+.2*\j})ellipse(.085 and .125);}
\draw[DeepBlue,thick](0,2.45)--(0,2.5);
\draw[DeepBlue,thick](0,.19)to[out=-90,in=180](.3,.12)--(1.7,.12)to[out=0,in=-90](1.85,.25)to[out=90,in=0](1.6,.37)--(.6,.37);
\draw[guide](.2,2.5)--(3.5,2.5);\draw[<->](3.2,0)--(3.2,2.5)node[midway,right]{$x$};\draw[<->](3.2,2.5)--(3.2,4.5)node[midway,right]{$s=H-x$};
\draw[<->](6,0)--(6,4.5)node[midway,right]{$H$};
\node[right]at(.45,1.5){移动部分$\lambda x$};\node[below]at(.9,-.2){地板上的静止余链};
\draw[chTenArrow](-.75,1.05)--(-.75,2.05)node[midway,left]{$\dot x>0$};
\end{tikzpicture}\figcap{零原长弹簧、悬垂链与地板余链；箭头示上移拾取阶段，下移时发生沉积。}\label{ch10:fig:chain}\end{chTenFigure}}

\newcommand{\chTenWedgeFigure}{\begin{chTenFigure}\begin{tikzpicture}[every node/.style=chTenLabel,>=Latex,scale=.86]
\draw[ground](-4.1,0)--(4.1,0);
\draw[thick](-3,0)--(-3,3)--(0,0)--cycle;\draw[thick](0,0)--(3,3)--(3,0)--cycle;
\fill[DeepBlue](-3,3)circle(3pt);\fill[DeepBlue](3,3)circle(3pt);\draw[thick,DeepBlue](-3,3)--(3,3);
\node[above]at(0,3.15){初始杆长$L_0$};\node at(-2.1,.6){$m$};\node at(2.1,.6){$m$};
\draw(.85,0)arc(0:45:.85);\node at(1.25,.42){$45^\circ$};
\node at(0,-.65){初始构形：球按质点表示};
\begin{scope}[xshift=7.3cm]
\draw[ground](-1,0)--(4.6,0);\draw[guide](0,0)--(3,3);
\draw[thick](.8,0)--(3.8,3)--(3.8,0)--cycle;
\coordinate(P)at(2.65,1.85);\fill[DeepBlue](P)circle(3pt)node[below right]{$P$};
\draw[DeepBlue,thick](-.9,1.85)--(P);\node[above]at(.2,1.95){弹性杆};
\draw[chTenArrow](P)--++(-.7,.7);\node[left]at(1.6,2.55){$N$};
\draw[guide](3.8,3)--(2.9,3.9) (P)--(1.75,2.75);
\draw[<->](2.9,3.9)--(1.75,2.75)node[midway,above left]{$s$};
\draw[guide](3,3)--(3,3.8) (3.8,3)--(3.8,3.8);\draw[chTenArrow](3,3.6)--(3.8,3.6)node[midway,above]{$X$};
\draw[guide](3,-.2)--(3,-.8) (2.65,-.2)--(2.65,-.8);
\draw[chTenArrow](3,-.6)--(2.65,-.6);\node[left]at(2.45,-.6){$q=X-s/\sqrt2$};
\end{scope}\end{tikzpicture}\figcap{左图为严格$45^\circ$初始构形；右图放大右半边的运动几何，虚斜线为初始斜面。$N$垂直当前斜面。}\label{ch10:fig:wedges}\end{chTenFigure}}

\newcommand{\chTenChargeMode}[2]{\begin{tikzpicture}[scale=.6,every node/.style=chTenLabel]
\draw[TextGray](0,0)circle(1.5);
\foreach \v[count=\j from 0] in {#1}{\fill[DeepBlue]({120*\j}:1.5)circle(3pt);\ifnum\v=0\else\draw[chTenArrow]({120*\j}:1.5)--++({120*\j+90}:{.35*\v});\fi}
\node at(0,-2.2){#2};\end{tikzpicture}}
\newcommand{\chTenChargesFigure}{\begin{chTenFigure}
\begin{tikzpicture}[every node/.style=chTenLabel,>=Latex]
\draw[TextGray](0,0)circle(1.7);
\draw[DeepBlue,thick](0:1.7)--(120:1.7)--(240:1.7)--cycle;
\foreach \j in {0,1,2}{\fill[DeepBlue]({120*\j}:1.7)circle(3pt);}
\draw[->](0,0)--(2.5,0)node[right]{$x$};\draw[->](0,0)--(0,2.4)node[above]{$y$};
\node at(-.38,.55){$\odot\ z$};\node[above right]at(1.7,.1){1};\node[above left]at(120:1.8){2};\node[below left]at(240:1.8){3};\node[below]at(0,-2.0){大圆正视，$z$轴朝向读者};
\end{tikzpicture}\hspace{.4cm}
\chTenChargeMode{1,-1,0}{$\vect a\propto(1,-1,0)$}\hspace{.2cm}
\chTenChargeMode{1,1,-2}{$\vect a\propto(1,1,-2)$}\par
\begin{tikzpicture}[every node/.style=chTenLabel]
\foreach \x in {0,1.7,3.4}{\draw[DeepBlue,thick](\x,0)circle(.17);\fill(\x,0)circle(1pt);}
\node[right]at(4,0){$z_1=z_2=z_3>0$：三球均出大圆平面};
\end{tikzpicture}
\figcap{平衡构形、两个大圆内振型及一个离面振型。各圆均从右侧点1起逆时针编号；切向箭头表示$a_i$。}\label{ch10:fig:charges}\end{chTenFigure}}

\newcommand{\chTenSpringPendulumFigure}{\begin{chTenFigure}\begin{tikzpicture}[every node/.style=chTenLabel,>=Latex]
\coordinate(O)at(0,3.6);\coordinate(E)at(0,.7);\coordinate(P)at(2,-.3);
\draw[ground](-.8,3.6)--(.8,3.6);\draw[guide](O)--(E);\draw[spring](O)--(P);\fill(P)circle(3pt)node[below right]{$P,m$};
\draw[->](E)--(3.6,.7)node[right]{$x$};\draw[->](E)--(0,-1.2)node[below]{$y$};
\fill(E)circle(2pt)node[left]{$E$};\node[above]at(O){$O$};
\draw[guide](P)--(2,.7);\draw[<->](-.55,.7)--(-.55,3.6)node[midway,left]{$L$};
\draw[<->](2.8,-.3)--(2.8,.7)node[midway,right]{$y$};\node[right]at(1.35,1.75){$r$};
\node at(2.6,2.9){$OP^2=x^2+(L+y)^2$};
\end{tikzpicture}\figcap{平衡点$E$与瞬时位置$P$；$y$向下为正，弹簧长度$r$从悬点$O$量起。}\label{ch10:fig:spring-pendulum}\end{chTenFigure}}

% 初始装置图仅含题设；分析量与受力另见解答图。
\newcommand{\chTenFrictionProblemFigure}{\begin{chTenFigure}\begin{tikzpicture}[every node/.style=chTenLabel,>=Latex]
\draw[ground](-.5,0)--(8,0);\draw[thick](0,0)rectangle(1.25,.8);\draw[thick](6.1,0)rectangle(7.35,.8);
\node at(.625,.4){$A,M$};\node at(6.725,.4){$B,m$};\draw[spring](1.25,.4)--(6.1,.4);
\node[above]at(3.675,.8){$k,\ l_0$};
\draw[guide](.625,-.1)--(.625,-.8) (6.725,-.1)--(6.725,-.8);
\draw[<->](.625,-.6)--(6.725,-.6)node[midway,below]{$l$};\node[right]at(8,0){$\mu$};
\end{tikzpicture}\figcap{由静止释放的初始状态；物块按质点建模，间距按两中心计。}\label{ch10:fig:friction-problem}\end{chTenFigure}}

\newcommand{\chTenGovernorProblemFigure}{\begin{chTenFigure}\begin{tikzpicture}[every node/.style=chTenLabel,>=Latex,scale=.8]
\coordinate(T)at(0,2);\coordinate(B)at(0,-2);\coordinate(P)at(2,0);\coordinate(Q)at(-2,0);
\draw[guide](0,-3.1)--(0,3);\draw[thick](T)--(P)--(B)--(Q)--cycle;
\draw[thick](-.17,1.78)rectangle(.17,2.22);\draw[thick](-.22,-2.23)rectangle(.22,-1.77);
\fill[DeepBlue](P)circle(4pt);\fill[DeepBlue](Q)circle(4pt);
\node[left]at(-2.15,0){$m$};\node[right]at(2.15,0){$m$};
\node[above left]at(-.2,2.23){固定套筒};\node[left]at(-.4,-2){$M$};\node[right]at(.4,-2){可沿轴滑动};
\node[above right]at(1.12,1.1){$l$};\node[below right]at(1.12,-1.1){$l$};
\node[above left]at(-1.12,1.1){$l$};\node[below left]at(-1.12,-1.1){$l$};
\draw[chTenArrow,domain=210:345,samples=30]plot({.75*cos(\x)},{-2.75+.2*sin(\x)});\node[right]at(.9,-2.75){$\omega_0$};
\end{tikzpicture}\figcap{调速器的装置截面：两球、四根等长轻杆与上、下套筒。}\label{ch10:fig:governor-problem}\end{chTenFigure}}

\newcommand{\chTenChainProblemFigure}{\begin{chTenFigure}\begin{tikzpicture}[every node/.style=chTenLabel,>=Latex,scale=.8]
\draw[ground](-1.2,4.5)--(4.5,4.5);\draw[ground](-1.2,0)--(4.5,0);
\draw[spring](0,4.5)--(0,2.5);\fill(0,2.5)circle(2pt);
\foreach \j in {0,...,11}{\draw[DeepBlue,thick](0,{.125+.2*\j})ellipse(.085 and .125);}
\draw[DeepBlue,thick](0,2.45)--(0,2.5);
\draw[DeepBlue,thick](0,.19)to[out=-90,in=180](.3,.12)--(1.7,.12)to[out=0,in=-90](1.85,.25)to[out=90,in=0](1.6,.37)--(.6,.37);
\node[right]at(.4,3.5){$k$，原长为零};\node[right]at(.4,1.6){线密度 $\lambda$};\node[below]at(.9,-.2){地板上的余链};
\draw[<->](4,0)--(4,4.5)node[midway,right]{$H$};
\end{tikzpicture}\figcap{悬于零原长弹簧下的软链，下部平放在地板上。}\label{ch10:fig:chain-problem}\end{chTenFigure}}

\newcommand{\chTenWedgeProblemFigure}{\begin{chTenFigure}\begin{tikzpicture}[every node/.style=chTenLabel,>=Latex,scale=.8]
\draw[ground](-4.1,0)--(4.1,0);
\draw[thick](-3,0)--(-3,3)--(0,0)--cycle;\draw[thick](0,0)--(3,3)--(3,0)--cycle;
\fill[DeepBlue](-3,3)circle(3pt);\fill[DeepBlue](3,3)circle(3pt);\draw[thick,DeepBlue](-3,3)--(3,3);
\node[above]at(0,3.15){$L_0,\ k$};\node[above left]at(-3,3){$m$};\node[above right]at(3,3){$m$};
\node at(-2.1,.6){$m$};\node at(2.1,.6){$m$};
\draw(.85,0)arc(0:45:.85);\node at(1.25,.42){$45^\circ$};
\draw(-.85,0)arc(180:135:.85);\node at(-1.25,.42){$45^\circ$};
\end{tikzpicture}\figcap{两球在斜面顶端由静止释放，水平弹性轻杆初始处于自由长度。}\label{ch10:fig:wedges-problem}\end{chTenFigure}}

\section{简谐振动的描述与判据}\label{ch10:sec:shm}
\subsection{振幅、相位与初始条件}
振动是物理量在平衡值附近反复变化的过程。弹簧振子的位移、交流电路中的电荷以及分子中原子的偏移，虽然所指对象不同，却可以服从相同的数学方程。最简单的周期振动是简谐振动（simple harmonic oscillation）：
\begin{equation}\label{ch10:eq:definition}
 Q(t)=Q_0+A\cos(\omega t+\varphi),\qquad A\geq0,\quad\omega>0.
\end{equation}
这里$Q_0$是平衡值，$A$是振幅，即相对平衡值的最大偏离量；$\psi=\omega t+\varphi$是相位，表示振动进行到一个周期中的哪一阶段；$\varphi$是$t=0$时的初相位。角频率$\omega$的单位为$\mathrm{rad/s}$，相位增加$2\pi$完成一次振动，故周期与频率分别为
\[
 T=\frac{2\pi}{\omega},\qquad f=\frac1T=\frac{\omega}{2\pi}.
\]
同一个位移通常在一周期内出现两次，此时运动方向相反。因此只有位移还不能完全确定状态，还需要速度，或者等价的运动方向信息。以下以平衡位置为$x=0$的质点为例，直接对位移求导：
\begin{align}
 x&=A\cos\psi,\label{ch10:eq:shm-x}\\
 v=\dot x&=-A\omega\sin\psi,\label{ch10:eq:shm-v}\\
 a=\ddot x&=-A\omega^2\cos\psi=-\omega^2x.\label{ch10:eq:shm-a}
\end{align}
速度振幅为$A\omega$，加速度振幅为$A\omega^2$。当$x>0$时$a<0$，当$x<0$时$a>0$，故加速度总指向平衡位置；这不是说物体总向平衡位置运动，而是说物体背离平衡位置时速率减小、接近平衡位置时速率增大。

若初始条件为$x(0)=x_0$、$v(0)=v_0$，则
\[
 x_0=A\cos\varphi,\qquad -\frac{v_0}{\omega}=A\sin\varphi.
\]
平方后相加，得到
\begin{equation}\label{ch10:eq:initial}
 A=\sqrt{x_0^2+\frac{v_0^2}{\omega^2}},\qquad
 \cos\varphi=\frac{x_0}{A},\quad \sin\varphi=-\frac{v_0}{A\omega}.
\end{equation}
用余弦确定角的候选值，再用正弦的符号选象限，便能避免只用反三角函数主值产生的歧义。也可记为
$\varphi=\operatorname{atan2}(-v_0/\omega,x_0)$，其中$\operatorname{atan2}(Y,X)$表示由横分量$X$、纵分量$Y$确定象限的角。$A=0$对应静止，初相位无物理意义。

更一般地，若已知$x/A=c$，则$\psi=\pm\arccos c+2n\pi$；若已知$-v/(A\omega)=s$，则$\psi=\arcsin s+2n\pi$或$\pi-\arcsin s+2n\pi$。两条条件应同时满足。求两个状态间的历时，要沿实际运动方向连续增加相位，不能每次把相位跳回主值区间：
\begin{equation}\label{ch10:eq:phase-time}
 \Delta t=\frac{\psi_2-\psi_1}{\omega}.
\end{equation}
例如从$x=A,v=0$到$x=0,v<0$，相位从0增加到$\pi/2$，历时$T/4$；到$x=0,v>0$则要增加到$3\pi/2$，历时$3T/4$。

\subsection{复数表示为什么方便}
把平衡值先减去，令$q=Q-Q_0$。利用欧拉公式$\mathrm e^{\ii\psi}=\cos\psi+\ii\sin\psi$，引入辅助复变量
\[
 \widetilde q(t)=A\mathrm e^{\ii(\omega t+\varphi)}
 =\widetilde A\mathrm e^{\ii\omega t},\qquad
 \widetilde A=A\mathrm e^{\ii\varphi},\qquad q=\Re\widetilde q.
\]
$\widetilde A$称复振幅，它的模是振幅，辐角是初相位。复数本身不意味着物体运动到一个“虚位置”；它是同时保存两个正交相位分量的计算工具。由于求导与取实部可交换，
\[
 \frac{\dd\widetilde q}{\dd t}=\ii\omega\widetilde q,
 \qquad \frac{\dd^2\widetilde q}{\dd t^2}=-\omega^2\widetilde q,
 \qquad \dot q=\Re\dot{\widetilde q}.
\]
于是线性微分方程中的求导可以转为乘$\ii\omega$。线性叠加同样可以先在复数域计算再取实部。但平方、乘积等非线性运算不能这样直接交换：通常
$(\Re\widetilde q)^2\ne\Re(\widetilde q^{2})$。例如计算动能时必须使用真实速度的平方，而不是复速度平方的实部。

\subsection{回复判据、能量判据与相图}
若一个坐标的动力学方程为
\begin{equation}\label{ch10:eq:linear-test}
 \ddot x=-\omega_0^2x,\qquad\omega_0^2>0,
\end{equation}
其中$\omega_0$不随$x,t$变化，则它作简谐振动。平衡位置不在零点时，先用$x=q-q_0$表示相对平衡位置的偏移。若$\ddot x=+\lambda^2x$，则是指数型不稳定运动，不能仅凭“加速度正比于位移”判断为振动。

将式\eqref{ch10:eq:linear-test}乘$\dot x$，逐项写成全导数：
\[
 \dot x\ddot x+\omega_0^2x\dot x=0
 \quad\Longrightarrow\quad
 \frac{\dd}{\dd t}\left(\frac12\dot x^2+\frac12\omega_0^2x^2\right)=0.
\]
结合初始条件，可写为
\begin{equation}\label{ch10:eq:ellipse}
 v^2+\omega_0^2x^2=\omega_0^2A^2,
 \qquad \frac{x^2}{A^2}+\frac{v^2}{A^2\omega_0^2}=1.
\end{equation}
这就是位移和速度大小“此消彼长”的关系。在质量$m$、劲度系数$k=m\omega_0^2$的振子中，它等价于能量守恒：
\begin{equation}\label{ch10:eq:energy}
 E=\frac12mv^2+\frac12kx^2=\frac12kA^2.
\end{equation}
平衡点$x=0$处动能最大；两端$x=\pm A$处动能为零、势能最大。

反过来，在光滑、非静止的实际动力学轨道上，若式\eqref{ch10:eq:ellipse}恒成立，在$v\ne0$的区间求导并约去$v$可得式\eqref{ch10:eq:linear-test}，再以加速度连续性延拓到转向点。也可在单调运动的一段上积分
\[
 \frac{\dd x}{\sqrt{A^2-x^2}}=\pm\omega_0\dd t,
 \qquad \arcsin(x/A)=\pm\omega_0t+C,
\]
并在转向点接上正确的速度分支。单独一条能量等式并不允许任意安排在端点停留多久；这还需要原动力学和光滑性条件。恒定的$x=\pm A$也满足该代数等式，却不是回复力非零时的动力学解。

\chTenPhaseFigure
图\ref{ch10:fig:phase}把一个运动状态画成相平面中的一点。顶部$x=0,v>0$时点必须向右移动；右端$x=A,v=0$时加速度为负，点向下移动，故轨迹是顺时针方向。横纵轴表示不同的物理量，纸面椭圆的长短还与标度有关。用$X=x/A$、$V=v/(A\omega_0)$归一化后，轨迹才明确成为$X^2+V^2=1$的单位圆；不能仅凭一幅图外观看起来像圆就断言量纲有误。

\begin{origexamplebox}{竖直弹簧的平衡位置与周期}\label{ch10:ex:vertical-spring}
劲度系数为$k$的轻弹簧竖直悬挂，末端连接质量为$m$的小球。求平衡伸长量及小球在平衡位置附近作小振动的周期。运动过程中弹簧保持在线性弹性范围内。

\ThQVerticalSpring
\end{origexamplebox}
\Needspace{6\baselineskip}\begin{xesolution}
以弹簧原长时球的位置为伸长量$y$的零点，向下为正。球所受重力为$+mg$，弹簧力为$-ky$，故
\[
 m\ddot y=mg-ky.
\]
平衡要求加速度为零，因此$ky_0=mg$。图\ref{ch10:fig:vertical-spring}中的$x$不是弹簧总伸长，而是相对平衡位置的位移：$x=y-y_0$。由于$y_0$为常量，$\ddot y=\ddot x$，代入得
\[
 m\ddot x=mg-k(y_0+x)=-kx.
\]
\chTenVerticalFigure
比较回复方程，$\omega_0^2=k/m$，所以
\[
 \boxed{y_0=mg/k},\qquad \boxed{T=2\pi\sqrt{m/k}}.
\]
重力被平衡伸长抵消，并未改变线性回复系数。$m/k$的单位为秒平方，周期量纲正确。如果振幅过大使弹簧失去线性、碰到限位或连接松弛，便不能继续使用这同一方程。
\end{xesolution}

\section{稳定平衡附近的小振动}\label{ch10:sec:small}
\subsection{先找平衡，再确定回复系数}
严格简谐要求回复关系精确线性；小振动则是非线性系统在稳定平衡附近的近似。对一个无显含时间约束、无耗散的一自由度系统，设广义坐标为$q$，势能为$U(q)$。平衡点$q_0$满足$U'(q_0)=0$。取$\eta=q-q_0$，泰勒展开给出
\[
 U(q)=U(q_0)+\frac12U''(q_0)\eta^2
 +\frac1{3!}U'''(q_0)\eta^3+\cdots.
\]
当$U''(q_0)>0$且振幅足够小，高阶项相对二阶项很小，势能可近似为抛物线。动能一般写作$T_{\rm kin}=\tfrac12M(q)\dot q^2$；展开$M(q)=M(q_0)+M'(q_0)\eta+\cdots$，其中$\eta\dot\eta^2$已经是三阶小量，二阶近似只留下
\[
 T_{\rm kin}\simeq\frac12M_0\dot\eta^2,\qquad M_0=M(q_0)>0.
\]
于是二阶能量为两个平方项之和，求导得
\begin{equation}\label{ch10:eq:small-frequency}
 M_0\ddot\eta+U''(q_0)\eta=0,
 \qquad \omega_0=\sqrt{\frac{U''(q_0)}{M_0}}.
\end{equation}
$M_0$是对坐标$q$而言的惯性系数，不一定具有质量的量纲；若$q$是角度，它就是转动惯量。若$U''<0$，该平衡不稳定；若$U''=0$，必须继续检查高阶项。例如$U\propto\eta^4$可以稳定，却没有与振幅无关的非零小振动频率。

\subsection{单摆：把小角度近似落到能量上}
长$l$的单摆取偏离竖直向下方向的角$\theta$为坐标。小球速率为$l|\dot\theta|$，故
\[
 T_{\rm kin}=\frac12ml^2\dot\theta^2,\qquad
 U=mgl(1-\cos\theta).
\]
在$\theta=0$处，$U'=mgl\sin\theta=0$，$U''=mgl>0$，所以
\[
 \omega_0^2=\frac{mgl}{ml^2}=\frac gl,\qquad
 T_0=2\pi\sqrt{l/g}.
\]
这里的判定过程依次是：写能量、找稳定平衡、保留二阶项、比较惯性和回复系数。它对弹性约束、转动约束和复杂机构的小振动都适用。

周期为什么只在小摆角下近似不变？若最大摆角为$\theta_m$，能量守恒给
\[
 \dot\theta^2=\frac{2g}{l}(\cos\theta-\cos\theta_m),\quad
 T=4\sqrt{\frac l{2g}}\int_0^{\theta_m}
 \frac{\dd\theta}{\sqrt{\cos\theta-\cos\theta_m}}.
\]
令$\sin(\theta/2)=s\sin u$，$s=\sin(\theta_m/2)$，可得
\[
 T=4\sqrt{\frac lg}\int_0^{\pi/2}\frac{\dd u}{\sqrt{1-s^2\sin^2u}}
 =T_0\left(1+\frac{s^2}{4}+\cdots\right)
 =T_0\left(1+\frac{\theta_m^2}{16}+\cdots\right).
\]
第二个等号使用$(1-z)^{-1/2}=1+z/2+\cdots$及$\int_0^{\pi/2}\sin^2u\,\dd u=\pi/4$。当$\theta_m=10^\circ=\pi/18\simeq0.1745\,\mathrm{rad}$，首项相对修正约$0.00190$，即$0.19\%$。它给出了“小”到底有多小的可检验尺度。

\begin{origexamplebox}{复摆的小振动}\label{ch10:ex:physical-pendulum}
质量为$M$、质心到水平固定转轴的距离为$a>0$的刚体可以绕该轴无摩擦转动，绕轴转动惯量为$I$。求质心位于轴下方的平衡位置附近的小振动周期。

\ThQPhysicalPendulum
\end{origexamplebox}
\Needspace{6\baselineskip}\begin{xesolution}
取图\ref{ch10:fig:physical-pendulum}中向右偏转的$\theta$为正。重力对轴$O$的力矩使偏角减小，其有符号值为$-Mga\sin\theta$，因此精确转动方程是
\[
 I\ddot\theta=-Mga\sin\theta.
\]
\chTenPendulumFigure
小角度以弧度表示，$\sin\theta\simeq\theta$，得到$I\ddot\theta+Mga\theta=0$。也可由$U=Mga(1-\cos\theta)\simeq Mga\theta^2/2$及$T_{\rm kin}=I\dot\theta^2/2$读出相同的两个系数。故
\[
 \boxed{\omega_0=\sqrt{Mga/I}},\qquad
 \boxed{T=2\pi\sqrt{I/(Mga)}}.
\]
当所有质量集中在距轴$a$处，$I=Ma^2$，结果退化为长$a$的单摆；一般刚体还含绕质心的转动惯量，不能直接用$Ma^2$替代$I$。若$a\to0$，重力回复力矩消失，周期趋于无穷，与模型一致。
\end{xesolution}

\section{阻尼振动}\label{ch10:sec:damping}
\subsection{三种阻尼情形及其通解}
阻尼表示运动中的耗散。以下采用黏性阻力模型$F_d=-b\dot x$，其中$b>0$，单位为$\mathrm{kg/s}$。这种线性关系只适用于相应速度和介质条件，不能替代干摩擦或任意流体阻力。运动方程为
\begin{equation}\label{ch10:eq:damping}
 m\ddot x+b\dot x+kx=0,
 \qquad\ddot x+2\gamma\dot x+\omega_0^2x=0,
 \quad\gamma=\frac{b}{2m},\quad\omega_0=\sqrt{k/m}.
\end{equation}
这里$\gamma$和$\omega_0$都是频率量纲，但$\gamma$控制衰减，不是振荡角频率。设$x=\mathrm e^{rt}$，除去非零指数，特征方程和根为
\[
 r^2+2\gamma r+\omega_0^2=0,
 \qquad r_\pm=-\gamma\pm\sqrt{\gamma^2-\omega_0^2}.
\]
\paragraph{欠阻尼：$0<\gamma<\omega_0$。}
令$\omega_d=\sqrt{\omega_0^2-\gamma^2}$。共轭根$-\gamma\pm\ii\omega_d$给出两个实独立解：
\begin{equation}\label{ch10:eq:under}
 x=\mathrm e^{-\gamma t}(C\cos\omega_dt+D\sin\omega_dt)
 =A_0\mathrm e^{-\gamma t}\cos(\omega_dt+\varphi).
\end{equation}
若$x(0)=x_0,\dot x(0)=v_0$，则$C=x_0$，由$v_0=-\gamma C+\omega_dD$得$D=(v_0+\gamma x_0)/\omega_d$。准周期$T_d=2\pi/\omega_d$大于无阻尼周期；严格说$x(t)$不再周期重复，而是振荡与衰减的乘积。

\paragraph{临界阻尼：$\gamma=\omega_0$。}
此时两个根相同，不能把两个相同的指数当成两条独立解。直接令$x=\mathrm e^{-\gamma t}y$代入方程，得到$\ddot y=0$，故
\begin{equation}\label{ch10:eq:critical}
 x=(C+Dt)\mathrm e^{-\omega_0t}
 =[x_0+(v_0+\omega_0x_0)t]\mathrm e^{-\omega_0t}.
\end{equation}
也可在欠阻尼解中用$\sin\omega_dt/\omega_d\to t$求极限；初速度对应的正弦项系数随$\omega_d\to0$变化，不能只把正弦置零。

\paragraph{过阻尼：$\gamma>\omega_0$。}
记$\kappa=\sqrt{\gamma^2-\omega_0^2}$，则两个根均为负实数：
\begin{equation}\label{ch10:eq:over}
 x=C\mathrm e^{(-\gamma+\kappa)t}+D\mathrm e^{(-\gamma-\kappa)t}.
\end{equation}
初值给$C+D=x_0$、$\kappa(C-D)=v_0+\gamma x_0$，也可合写为
\[
 x=\mathrm e^{-\gamma t}\left[x_0\cosh\kappa t+
 \frac{v_0+\gamma x_0}{\kappa}\sinh\kappa t\right].
\]
它没有反复过零的振荡，但初速度很大时仍可能穿过平衡位置一次；“非振荡”不能误读为“绝不越过平衡点”。

\chTenDampingFigure
图\ref{ch10:fig:damping-types}三幅图均取$x_0>0,v_0=0$，只改变$\gamma/\omega_0$。欠阻尼有过冲；临界阻尼不反复振荡；过阻尼中的慢指数使返回延迟。其慢衰减率为
\[
 \gamma-\sqrt{\gamma^2-\omega_0^2}
 =\frac{\omega_0^2}{\gamma+\sqrt{\gamma^2-\omega_0^2}},
\]
随$\gamma$增大而减小。对于相同$m,k$、从同一偏离位置静止释放，以及用长期收敛尺度比较的情形，临界阻尼是欠阻尼与非振荡返回之间的分界，具有最快的非振荡返回特性。不能脱离初始条件、允许过冲量和“返回”的容差定义，宣称它在所有指标下都最快。临界衰减的特征时间是$1/\omega_0$量级，并非一次完整的无阻尼周期。

\subsection{包络、能量与品质因数}
欠阻尼位移被两条包络$\pm A_0\mathrm e^{-\gamma t}$限制，振幅尺度减少到$1/\mathrm e$的时间为$\tau_A=1/\gamma$。相邻同向极值相隔$T_d$，极值绝对值之比恰为$\mathrm e^{\gamma T_d}$，其对数$\gamma T_d$称对数衰减率。图\ref{ch10:fig:damping-envelope}同时画出了位移及两条包络。
\chTenEnvelopeFigure
机械能仍定义为$E=m\dot x^2/2+kx^2/2$，但已不守恒。将运动方程乘$\dot x$，可得精确结论
\begin{equation}\label{ch10:eq:dissipation}
 \frac{\dd E}{\dd t}=\dot x(m\ddot x+kx)=-b\dot x^2\leq0.
\end{equation}
令$\psi=\omega_dt+\varphi$，精确求导得
\[
 \dot x=-A_0\mathrm e^{-\gamma t}(\gamma\cos\psi+\omega_d\sin\psi),
\]
所以瞬时能量为
\[
 E(t)=\frac{mA_0^2}{2}\mathrm e^{-2\gamma t}
 \left[(\gamma\cos\psi+\omega_d\sin\psi)^2+\omega_0^2\cos^2\psi\right].
\]
方括号一般随相位变化，因此不能把瞬时能量直接写成固定常数乘$\mathrm e^{-2\gamma t}$。当$\gamma/\omega_0\ll1$，一周期内包络变化很小，将慢变包络暂时固定再平均，得到
\[
 \overline E\simeq\frac12m\omega_0^2A^2(t),\qquad
 \overline P_{\rm loss}\simeq\frac12b\omega_0^2A^2(t),
 \qquad \frac{\dd\overline E}{\dd t}\simeq-2\gamma\overline E.
\]
这里$P_{\rm loss}>0$是损失功率的大小。于是平均能量近似按$\mathrm e^{-2\gamma t}$衰减，$\tau_E\simeq1/(2\gamma)$。品质因数（quality factor）可定义为本线性模型的参数
\begin{equation}\label{ch10:eq:quality}
 Q=\frac{m\omega_0}{b}=\frac{\omega_0}{2\gamma}.
\end{equation}
在弱阻尼下，它还有能量意义：
\[
 Q\simeq\frac{\omega_0\overline E}{\overline P_{\rm loss}}
 \simeq2\pi\frac{\overline E}{|\Delta E|_{\text{一周期}}}.
\]
第二个比值要求一周期损耗占比较小。$Q$大表示每次振动损耗较少，而不是完全没有耗散。

\subsection{阻尼相图与损耗相位}
相图是以坐标和速度（或动量）为轴画出的状态轨迹。因为$x$和$v$都衰减，欠阻尼相轨迹沿顺时针向原点收缩，见图\ref{ch10:fig:damped-phase}左图；它不是一组互不连接的椭圆。
\chTenDampedPhaseFigure
为分开“包络收缩”和“相位偏离”，定义$X=\mathrm e^{\gamma t}x$、$V=\mathrm e^{\gamma t}v$，注意$V\ne\dot X$。由$X=A_0\cos\psi$及$V=-\gamma A_0\cos\psi-\omega_d A_0\sin\psi$消去相位：
\begin{equation}\label{ch10:eq:tilted}
 X^2+\frac{(V+\gamma X)^2}{\omega_d^2}=A_0^2.
\end{equation}
这条曲线含交叉项$XV$。令$Y=V/\omega_0$并利用$\omega_d^2+\gamma^2=\omega_0^2$，得
\[
 X^2+Y^2+2\frac\gamma{\omega_0}XY
 =A_0^2\left(1-\frac{\gamma^2}{\omega_0^2}\right).
\]
在同一长度标度的$X,Y$图上，其主轴沿$Y=\pm X$；阻尼改变轴长比，而不是使这个特定归一化下的主轴角逐渐增大。若采用不同横纵标度，纸面倾角还会改变。图中的椭圆倾角不能直接当作损耗相位。

速度还可写为
\[
 v=-\omega_0A_0\mathrm e^{-\gamma t}\sin(\psi+\eta),\qquad
 \sin\eta=\frac\gamma{\omega_0},\quad
 \cos\eta=\frac{\omega_d}{\omega_0},\quad0<\eta<\frac\pi2.
\]
因此扣除衰减包络后，速度相对位移的相位偏离无阻尼的$\pi/2$，多出$\eta$。在此欠阻尼自由衰减约定下，可称$\eta$为损耗角，并有$\sin\eta=1/(2Q)$。不同学科也用“损耗角”表示受迫响应中的材料损耗相位，定义未必相同，使用时应说明；下文用$\delta$专指位移相对外力的滞后角。

\section{受迫振动、共振与谱线}\label{ch10:sec:driven}
\subsection{暂态与稳态：从复振幅到真实位移}
给振子施加$F(t)=F_0\cos\omega t$，其中$F_0>0,\omega>0$。线性方程为
\begin{equation}\label{ch10:eq:forced}
 m\ddot x+b\dot x+kx=F_0\cos\omega t.
\end{equation}
为什么先研究单一频率？满足通常傅里叶展开条件的周期力可以分解为常量和不同频率的正弦项，线性系统对各项的响应相加便是总响应；非周期驱动力在适当条件下可用傅里叶积分。这不是说任意力都等于一个正弦力。

先将右端暂写成$F_0\mathrm e^{\ii\omega t}$，设特解$\widetilde x=\widetilde X\mathrm e^{\ii\omega t}$。求导一次、两次分别带来$\ii\omega$、$-\omega^2$，因此
\[
 (k-m\omega^2+\ii b\omega)\widetilde X=F_0,
 \quad\widetilde X=\frac{F_0}{k-m\omega^2+\ii b\omega}
 =\frac{F_0(k-m\omega^2-\ii b\omega)}{(k-m\omega^2)^2+b^2\omega^2}.
\]
写$\widetilde X=A\mathrm e^{-\ii\delta}$，取实部得稳态$x_s=A\cos(\omega t-\delta)$，其中
\begin{align}
 A&=\frac{F_0}{\sqrt{(k-m\omega^2)^2+b^2\omega^2}},\label{ch10:eq:response-A}\\
 \cos\delta&=\frac{k-m\omega^2}{\sqrt{(k-m\omega^2)^2+b^2\omega^2}},\qquad
 \sin\delta=\frac{b\omega}{\sqrt{(k-m\omega^2)^2+b^2\omega^2}}.\label{ch10:eq:response-phase}
\end{align}
对$b>0$，必须取$0<\delta<\pi$，即$\delta=\operatorname{atan2}(b\omega,k-m\omega^2)$。只写$\tan\delta=b\omega/(k-m\omega^2)$会在$\omega>\omega_0$时丢失象限；在$\omega=\omega_0$，$\delta=\pi/2$而非“无定义”。

完整解为$x=x_h+x_s$，其中$x_h$是上一节相应阻尼类型的齐次解，用来满足初始位移和速度。例如欠阻尼时从静止开始施力，$x(0)=v(0)=0$，可写
\[
 x_h=\mathrm e^{-\gamma t}(C\cos\omega_dt+D\sin\omega_dt),\quad
 C=-A\cos\delta,\quad
 D=\frac{-A\omega\sin\delta+\gamma C}{\omega_d}.
\]
经过若干衰减时间，$x_h$才可忽略；稳态频率由外力决定，不是自由衰减频率$\omega_d$。在启动初期，两个频率的叠加还可能出现拍状变化。

\subsection{幅频、相频与共振条件}
令$r=\omega/\omega_0$，则
\[
 \frac{Ak}{F_0}=\frac1{\sqrt{(1-r^2)^2+r^2/Q^2}},\qquad
 \delta=\operatorname{atan2}(r/Q,1-r^2).
\]
图\ref{ch10:fig:resonance}按同一组$Q$画出幅频和相频两部分。实线和虚线不仅表示不同峰高，也给出相位经过$\pi/2$的陡缓变化。
\chTenResonanceFigure
当$\omega\to0$，弹性项控制响应，$A\to F_0/k,\delta\to0$，振子近似随外力准静态移动。当$\omega$很高，需同时有$\omega\gg\omega_0$和$\omega\gg b/m$，惯性项才压倒其余项，此时$A\simeq F_0/(m\omega^2)$、$\delta\simeq\pi$。在$\omega=\omega_0$，弹性项与惯性项在复表示中抵消，阻尼项平衡外力，$A=F_0/(b\omega_0)=QF_0/k$。

位移峰值需要最小化$A$的分母平方$D(\omega)$。求导：
\[
 D=m^2(\omega_0^2-\omega^2)^2+b^2\omega^2,
 \quad D'=4m^2\omega(\omega^2-\omega_0^2+2\gamma^2).
\]
所以只有$0<\gamma<\omega_0/\sqrt2$时，存在正频率的位移共振峰：
\[
 \omega_r=\sqrt{\omega_0^2-2\gamma^2},\qquad
 A_{\max}=\frac{F_0}{2m\gamma\sqrt{\omega_0^2-\gamma^2}}
 =\frac{QF_0/k}{\sqrt{1-1/(4Q^2)}}.
\]
若$\gamma\geq\omega_0/\sqrt2$，$A$在$\omega\geq0$上的最大值位于$\omega=0$，没有非零频率的位移共振峰。弱阻尼$Q\gg1$时$\omega_r\simeq\omega_0$、$A_{\max}\simeq QF_0/k$，这给出$Q$作为共振放大倍数的近似意义。

\subsection{吸收功率与无阻尼共振}
稳态平均输入功率必须用真实力乘真实速度。因$v=-A\omega\sin(\omega t-\delta)$，展开正弦并平均，交叉项平均为零、$\overline{\cos^2\omega t}=1/2$，得到
\begin{equation}\label{ch10:eq:power}
 \overline P=\overline{Fv}=\frac12F_0A\omega\sin\delta
 =\frac12b\omega^2A^2
 =\frac{F_0^2b\omega^2}{2[(k-m\omega^2)^2+b^2\omega^2]}.
\end{equation}
它恰等于$\overline{bv^2}$，所以稳态没有能量逐周期累积。把分母除以$\omega^2$，可看出功率在$k/\omega-m\omega=0$即$\omega=\omega_0$处最大：
\[
 \overline P_{\max}=\frac{F_0^2}{2b},\qquad
 (\omega A)_{\max}=\frac{F_0}{b}.
\]
因此速度振幅峰、吸收功率峰与$\delta=\pi/2$的位置一致；位移峰通常与它们不同。在$\omega_0$处，外力与速度同相，做功最有效；此时$\overline E=m\omega_0^2A^2/2$，从而$\omega_0\overline E/\overline P=Q$。

如果$b=0$，不能把有阻尼的峰值公式直接用于$\omega=\omega_0$。此时没有有限振幅的周期稳态。取共振特解$x_p=Ct\sin\omega_0t$，则
\[
 \ddot x_p+\omega_0^2x_p=2C\omega_0\cos\omega_0t,
 \quad C=\frac{F_0}{2m\omega_0}.
\]
由静止起振时$x=F_0t\sin\omega_0t/(2m\omega_0)$，包络随$t$增长；无限增长只属于无阻尼且始终线性的理想模型，实际会受耗散、非线性或结构极限限制。

\subsection{线型与半高全宽}
洛伦兹线型（Lorentzian profile）在概率论中对应柯西分布。中心为$z_0$、半宽为$\Gamma>0$的归一化形式是
\[
 L(z)=\frac1\pi\frac{\Gamma}{(z-z_0)^2+\Gamma^2},\qquad
 \int_{-\infty}^{\infty}L(z)\,\dd z=1.
\]
积分用$u=(z-z_0)/\Gamma$化为$\pi^{-1}[\arctan u]_{-\infty}^{\infty}=1$。分子中的$\Gamma$不可遗漏，否则量纲和归一化均不对。高斯线型（Gaussian profile）为
\[
 G(z)=\frac1{\sqrt{2\pi}\sigma}\exp[-(z-z_0)^2/(2\sigma^2)].
\]
比较图\ref{ch10:fig:profiles}时，为单看形状，先各自除以峰高，并令二者半高全宽相同。洛伦兹线型在$|z-z_0|=\Gamma$时降到一半，故全宽$2\Gamma$；高斯线型则为$2\sqrt{2\ln2}\,\sigma$。洛伦兹尾部按距离平方的倒数衰减，高斯尾部衰减更快。
\chTenProfileFigure
一般受迫振子的$A^2(\omega)$并非严格洛伦兹函数，因为分母含$\omega^4$。只有弱阻尼、$|\Delta|=|\omega-\omega_0|\ll\omega_0$时，$\omega_0^2-\omega^2\simeq-2\omega_0\Delta$，于是
\[
 A^2(\omega)\simeq\frac{F_0^2}{4m^2\omega_0^2}\frac1{\Delta^2+\gamma^2}.
\]
半高必须说明测的是$A$、$A^2$还是功率。下题求功率峰的带宽，再说明何时可把它用于位移强度谱。

\begin{origexamplebox}{共振半功率带宽}\label{ch10:ex:bandwidth}
弱阻尼受迫振子在共振附近的两个半功率频率为$\omega_1<\omega_2$。证明$Q\simeq\omega_0/(\omega_2-\omega_1)$。

\ThQBandwidth
\end{origexamplebox}
\Needspace{6\baselineskip}\begin{xesolution}
\ThABandwidth

由式\eqref{ch10:eq:power}，将功率除以峰值，得
\[
 \frac{\overline P}{\overline P_{\max}}
 =\frac{b^2\omega^2}{(k-m\omega^2)^2+b^2\omega^2}.
\]
半功率条件是该比值等于$1/2$，交叉相乘并移项：
\[
 2b^2\omega^2=(k-m\omega^2)^2+b^2\omega^2
 \quad\Longrightarrow\quad|k-m\omega^2|=b\omega.
\]
低频侧$k-m\omega_1^2>0$，高频侧相反，所以分别解
\[
 \omega_1^2+2\gamma\omega_1-\omega_0^2=0,
 \qquad\omega_2^2-2\gamma\omega_2-\omega_0^2=0.
\]
各取正根：
\[
 \omega_1=\sqrt{\omega_0^2+\gamma^2}-\gamma,
 \qquad\omega_2=\sqrt{\omega_0^2+\gamma^2}+\gamma.
\]
因此$\Delta\omega=2\gamma=b/m$，即
\[
 \boxed{Q=\frac{\omega_0}{\omega_2-\omega_1}}.
\]
对这一线性黏性阻尼模型，\emph{吸收功率}的等式实际上是精确的。弱阻尼时两个半功率点近似为$\omega_0\pm\gamma$，且在窄峰内$\omega^2$近似不变，$\overline P\propto\omega^2A^2$才近似与$A^2$成正比；这时也可把$2\gamma$作为位移强度谱的近似半高全宽。$\gamma$和带宽都是角频率，$Q$无量纲；$b\to0^+$时峰越来越窄，但恰好$b=0$的共振需使用上一小节的增长解。
\end{xesolution}

\section{多自由度小振动与简正模（选读）}\label{ch10:sec:modes}
\subsection{惯性矩阵和弹性矩阵从哪里来}
一个单摆只需一个角度，但双摆需要两个角度。具有$n$个独立广义坐标的系统，在平衡构形附近令$q_i$表示第$i$个坐标的微小偏移，并组成列向量$\vect q$。在光滑势能、无耗散、约束不显含时间的条件下，平衡使一次项消失，二阶近似为
\begin{equation}\label{ch10:eq:quadratic}
 T_{\rm kin}=\frac12\sum_{i,j}M_{ij}\dot q_i\dot q_j
 =\frac12\dot{\vect q}^{\mathsf T}\mathbf M\dot{\vect q},\qquad
 U-U_0=\frac12\sum_{i,j}K_{ij}q_iq_j
 =\frac12\vect q^{\mathsf T}\mathbf K\vect q.
\end{equation}
上标$\mathsf T$表示转置；$\mathbf M$是惯性矩阵，$\mathbf K$是弹性矩阵。其元素分别为动能对速度的二阶偏导和势能在平衡点的二阶偏导。混合偏导可交换，故两矩阵对称。即使最初把交叉项写成不对称系数，也可替换为对称部分，因为$\vect q^{\mathsf T}(\mathbf B-\mathbf B^{\mathsf T})\vect q=0$。

例如$U=aq_1^2/2+cq_1q_2+dq_2^2/2$的矩阵是$\bigl(\begin{smallmatrix}a&c\\c&d\end{smallmatrix}\bigr)$，不是把两个非对角元都写成$2c$。求$\partial U/\partial q_i$得到$\sum_jK_{ij}q_j$，所以拉格朗日方程给出
\begin{equation}\label{ch10:eq:matrix-motion}
 \mathbf M\ddot{\vect q}+\mathbf K\vect q=0.
\end{equation}

对于真正独立、非奇异的物理坐标，非零广义速度应具有正动能，即$\mathbf M$正定。若出现零惯性方向，要检查是否引入了冗余坐标、忽略了必要惯量，或在该平衡点使用了奇异坐标；不能直接求逆继续套用。势能矩阵则有三种重要情况：正定保证严格局部稳定；出现负方向意味着不稳定；有零方向时需看对称性或更高阶势能。零频整体转动可以是中性运动，$U\propto q^4$也可能稳定，而$U\propto-q^4$不稳定。

\subsection{两个模型：耦合摆与双摆}
先看图\ref{ch10:fig:coupled-pendula}：两根摆长均为$l$，摆球均为$m$，固定悬点水平相距$d$，摆球间轻弹簧原长恰为$d$、劲度系数为$k_c$。两角都从各自竖直向下方向向右量为正，不以图中左右位置改变角度正向。
\chTenCoupledPendulaFigure
两球坐标差的水平分量为$d+l\sin\theta_2-l\sin\theta_1$，竖直分量为$l\cos\theta_1-l\cos\theta_2$，故弹簧长度
\[
 d'=\sqrt{(d+l\sin\theta_2-l\sin\theta_1)^2+
 l^2(\cos\theta_1-\cos\theta_2)^2}.
\]
小角下$d'-d=l(\theta_2-\theta_1)+O(\theta^2)$，于是
\begin{align*}
 T_{\rm kin}&=\frac12ml^2(\dot\theta_1^2+\dot\theta_2^2),\\
 U&=mgl(2-\cos\theta_1-\cos\theta_2)+\frac12k_c(d'-d)^2\\
 &\simeq\frac12mgl(\theta_1^2+\theta_2^2)
 +\frac12k_cl^2(\theta_2-\theta_1)^2.
\end{align*}
记$\epsilon=k_cl/(mg)$，得
\[
 \mathbf M=ml^2\begin{pmatrix}1&0\\0&1\end{pmatrix},\qquad
 \mathbf K=mgl\begin{pmatrix}1+\epsilon&-\epsilon\\-\epsilon&1+\epsilon\end{pmatrix}.
\]
这里耦合出现在势能中。交换两摆不改变方程，故先试同相$\theta_1=\theta_2$与反相$\theta_1=-\theta_2$：
\[
 \ddot\theta_1=-\frac gl[(1+\epsilon)\theta_1-\epsilon\theta_2],
 \quad\omega_S^2=\frac gl,
 \qquad\omega_A^2=\frac gl(1+2\epsilon).
\]
同相时弹簧一阶伸长为零；反相时两端位移反向，伸长加倍。取$Q_S=(\theta_1+\theta_2)/2$、$Q_A=(\theta_1-\theta_2)/2$，直接相加、相减方程也得到两条独立简谐方程。任意初始状态是这两种模式的叠加，不必始终具有左右对称形状。

同一数学结构也适用于两物块弹簧模型，见图\ref{ch10:fig:coupled-masses}。两端弹簧均为$k$，中间弹簧为$k_c$，$x_1,x_2$从各自平衡位置向右量：
\[
 m\ddot x_1=-(k+k_c)x_1+k_cx_2,\qquad
 m\ddot x_2=k_cx_1-(k+k_c)x_2.
\]
因此同相、反相频率分别为$\sqrt{k/m}$和$\sqrt{(k+2k_c)/m}$。中间弹簧的伸长量为$x_2-x_1$，不是$x_1+x_2$。
\chTenCoupledMassFigure

再看图\ref{ch10:fig:double-pendulum}的等质量、等杆长双摆。两角均相对竖直方向，而不是第二杆相对第一杆的夹角。上球坐标$(l\sin\theta_1,-l\cos\theta_1)$，下球再加$(l\sin\theta_2,-l\cos\theta_2)$，故下球速率平方为
\[
 v_2^2=l^2\dot\theta_1^2+l^2\dot\theta_2^2
 +2l^2\dot\theta_1\dot\theta_2\cos(\theta_1-\theta_2).
\]
\chTenDoublePendulumFigure
将两球动能相加，重力势能相加，小角近似给
\begin{align*}
 T_{\rm kin}&\simeq\frac12ml^2(2\dot\theta_1^2+2\dot\theta_1\dot\theta_2+\dot\theta_2^2),\\
 U&=mgl[2(1-\cos\theta_1)+(1-\cos\theta_2)]
 \simeq\frac12mgl(2\theta_1^2+\theta_2^2),\\
 \mathbf M&=ml^2\begin{pmatrix}2&1\\1&1\end{pmatrix},\qquad
 \mathbf K=mgl\begin{pmatrix}2&0\\0&1\end{pmatrix}.
\end{align*}
这里势能没有交叉项，耦合却在动能中。仅把一种能量配方并不够。例如令$u=\theta_1+\theta_2/2,w=\theta_2$，则
\[
 T_{\rm kin}=ml^2\dot u^2+\frac14ml^2\dot w^2,
 \quad U=mgl\left(u^2-uw+\frac34w^2\right).
\]
动能虽已化为平方和，势能反而出现交叉项；真正目标是用\emph{同一个}坐标变换同时对角化两者。

\subsection{同时对角化、正交性和简正坐标}
一般二次型可以逐变量配方。例如$f_{11}\ne0$时，把$q_1$换成$q_1'=q_1+\sum_{j>1}f_{1j}q_j/f_{11}$，便可将含$q_1$的项写成$f_{11}(q_1')^2/2$和其余变量的二次型；继续处理余项即可消去交叉项。零主元时需要先换坐标。这个过程叫二次型的对角化，但它未保证同时处理两个二次型。

令旧坐标$\vect q=\mathbf B\vect Q$。代入式\eqref{ch10:eq:quadratic}，得到新矩阵
\[
 \mathbf M'=\mathbf B^{\mathsf T}\mathbf M\mathbf B,
 \qquad\mathbf K'=\mathbf B^{\mathsf T}\mathbf K\mathbf B.
\]
这称合同变换，与$\mathbf B^{-1}\mathbf K\mathbf B$那种相似变换不能混淆。如何选$\mathbf B$？先设所有坐标以同频率运动，$\vect q=\vect a\cos(\omega t+\phi)$，代入动力学方程：
\begin{equation}\label{ch10:eq:eigen}
 (\mathbf K-\omega^2\mathbf M)\vect a=0,
 \qquad\det(\mathbf K-\omega^2\mathbf M)=0.
\end{equation}
前式是广义本征值问题，后式是存在非零振幅向量的条件。但“行列式有$n$个根”还不足以保证$n$个独立本征向量，下面补上保证完备性的论证。

$\mathbf M$实对称正定，故有实对称正定平方根$\mathbf M^{1/2}$和逆$\mathbf M^{-1/2}$。令$\vect u=\mathbf M^{1/2}\vect a$，广义本征方程化为
\[
 \underbrace{\mathbf M^{-1/2}\mathbf K\mathbf M^{-1/2}}_{\mathbf D}
 \vect u=\lambda\vect u,\qquad\lambda=\omega^2.
\]
$\mathbf D$为实对称矩阵，所以存在$n$个实特征值及一组完备正交特征向量；重根子空间内也能选择正交基。令这些单位向量为$\vect u_i$，取$\vect a_i=\mathbf M^{-1/2}\vect u_i$，便有
\[
 \vect a_i^{\mathsf T}\mathbf M\vect a_j=\delta_{ij},\qquad
 \vect a_i^{\mathsf T}\mathbf K\vect a_j=\lambda_j\delta_{ij}.
\]
这既证明存在性，也给出构造。不同频率的质量正交性还可直接验证：由$\mathbf K\vect a_j=\lambda_j\mathbf M\vect a_j$及矩阵对称性，
\[
 \lambda_j\vect a_i^{\mathsf T}\mathbf M\vect a_j
 =\vect a_i^{\mathsf T}\mathbf K\vect a_j
 =\lambda_i\vect a_i^{\mathsf T}\mathbf M\vect a_j,
\]
所以$\lambda_i\ne\lambda_j$时该内积为零。若不把向量归一化，可记$M_i=\vect a_i^{\mathsf T}\mathbf M\vect a_i>0$，$K_i=\lambda_iM_i$；取$\mathbf B=(\vect a_1,\ldots,\vect a_n)$后，
\[
 T_{\rm kin}=\frac12\sum_i M_i\dot Q_i^2,
 \qquad U-U_0=\frac12\sum_i K_iQ_i^2,
 \qquad M_i\ddot Q_i+K_iQ_i=0.
\]
$Q_i$称简正坐标，$\omega_i=\sqrt{K_i/M_i}$称简正频率；只有一个$Q_i$被激发的运动称简正模（normal mode）。$\lambda_i>0$时$Q_i=C_i\cos(\omega_it+\phi_i)$；$\lambda_i=0$时$Q_i=C_i+D_it$；$\lambda_i<0$时出现增长/衰减指数。因此“$\mathbf M$正定、$\mathbf K$对称”只保证能解耦，并不保证每个坐标都振动。

初始条件也可明确投影：$Q_i(0)=\vect a_i^{\mathsf T}\mathbf M\vect q(0)/M_i$，$\dot Q_i(0)=\vect a_i^{\mathsf T}\mathbf M\dot{\vect q}(0)/M_i$。分别解每个标量方程，再用$\vect q=\sum_i\vect a_iQ_i$还原运动。原坐标通常含多个频率，并不各自简谐。

\subsection{完整求出双摆的两个模态}
在式\eqref{ch10:eq:eigen}中提出公共因子，令$\alpha=g/l$、$\lambda=\omega^2$，得到
\[
 \det\begin{pmatrix}2\alpha-2\lambda&-\lambda\\-\lambda&\alpha-\lambda\end{pmatrix}
 =2(\alpha-\lambda)^2-\lambda^2=0.
\]
展开为$\lambda^2-4\alpha\lambda+2\alpha^2=0$，故
\[
 \lambda_-=(2-\sqrt2)g/l,\qquad
 \lambda_+=(2+\sqrt2)g/l.
\]
第二行给$a_2/a_1=\lambda/(\alpha-\lambda)$，分别为$\sqrt2$和$-\sqrt2$。可选$\vect a_-=(1,\sqrt2)^{\mathsf T}$、$\vect a_+=(1,-\sqrt2)^{\mathsf T}$。所以同向摆动的低频模中，下杆角振幅是上杆的$\sqrt2$倍；高频模两杆反向，绝对比仍为$\sqrt2$，不是简单的等幅反相。

令
\[
 \theta_1=Q_-+Q_+,\quad\theta_2=\sqrt2(Q_--Q_+),
 \qquad Q_\pm=\frac{\theta_1}{2}\mp\frac{\theta_2}{2\sqrt2}.
\]
这里最后一式中$Q_+$取减号、$Q_-$取加号。代入可验算
\begin{align*}
 T_{\rm kin}&=\frac12[2(2+\sqrt2)ml^2]\dot Q_-^2
 +\frac12[2(2-\sqrt2)ml^2]\dot Q_+^2,\\
 U&=\frac12(4mgl)Q_-^2+\frac12(4mgl)Q_+^2.
\end{align*}
分别用弹性系数除惯性系数，恰好得到上面的$\lambda_\pm$，完成独立校验。图\ref{ch10:fig:double-modes}画的是各模的一个瞬时构形，杆长始终为$l$，箭头表示角运动的相对方向而不是独立平移。
\chTenDoubleModesFigure

\subsection{六原子环：从对称操作到六个模式}
图\ref{ch10:fig:benzene-model}左为苯分子的立体模型，右为简化的六原子环。真实苯含六个碳和六个氢；在右图的简化中，每个碳氢单元只当作质量$m$的质点，被限制在固定圆环上作切向小位移$x_j$，相邻单元用有效劲度$k$耦合。它只保留六个坐标，不能代替全部分子振动。
\chTenBenzeneModelFigure
给六点按逆时针编号$j=0,\ldots,5$，所有$x_j$的正向均为逆时针切向，循环约定$x_{j+6}=x_j$。有效能量和方程为
\[
 T_{\rm kin}=\frac12m\sum_{j=0}^5\dot x_j^2,\qquad
 U=\frac12k\sum_{j=0}^5(x_{j+1}-x_j)^2,
 \qquad m\ddot x_j=k(x_{j+1}+x_{j-1}-2x_j).
\]
这里$k$定义的是相邻切向坐标差的弹性系数，不把画在弦上的普通伸缩弹簧刚度不加换算地当成$k$。

对称操作的意义是使能量函数和动力学规律不变，而非要求任意时刻的实际位形都对称。图\ref{ch10:fig:ring-symmetry}依次画出绕环轴转$60^\circ$的$C_6$、绕对边中点连线转$180^\circ$的$C_2$和绕对顶点连线转$180^\circ$的$C'_2$。后两种在环面投影中表现为反射，切向正向也随之反转。以图中轴为例，其作用可写为
\[
 (C_6x)_j=x_{j-1},\qquad (C_2x)_j=-x_{1-j},\qquad
 (C'_2x)_j=-x_{-j},
\]
下标按模6理解。每个操作只是对$\sum\dot x_j^2$、$\sum(x_{j+1}-x_j)^2$重新编号，所以能量不变；单独交换任意两点一般会改变相邻关系，因而不一定是对称操作。
\chTenRingSymmetryFigure
如果$x_j(t)$是一个解，对它作对称操作仍得到解且不改变频率。初始条件却可以不对称，因此一般解不必保持某种对称形状；应区分“方程的对称性”和“所选状态的对称性”，不能把一切不对称初值都等同于严格意义的自发对称性破缺。

设$\mathbf C=C_6$，有$\mathbf C^6=\mathbf I$，其本征值只能是六次单位根。对任意向量$\vect a$，令$w=\mathrm e^{2\pi\ii/6}$，构造
\[
 \vect a^{(n)}=\sum_{s=0}^5w^{-ns}\mathbf C^s\vect a,
 \qquad\mathbf C\vect a^{(n)}=w^n\vect a^{(n)}.
\]
等号由把求和指标平移一位并用$w^{6n}=1$得到；某个投影可以为零，不能把零向量当模式。由于动力学矩阵与$\mathbf C$可交换，同一频率的子空间在$\mathbf C$作用下保持不变，可在其中选择这些投影。反射操作的平方为恒等，其本征值则只能为$\pm1$。

更直接地，设$x_j=A\exp[\ii(2\pi nj/6-\omega t)]$，代入每个质点的方程，约去公共因子得
\[
 -m\omega_n^2=k(\mathrm e^{2\pi\ii n/6}+\mathrm e^{-2\pi\ii n/6}-2),
 \quad\omega_n^2=\frac{4k}{m}\sin^2\frac{n\pi}{6}.
\]
所以频率平方为$0,k/m,3k/m,4k/m$，简并度依次为1、2、2、1。$n$与$6-n$是一对共轭行波，实部和虚部给两个独立实驻波模式。零频模是整体沿环匀速转动，不是周期无穷大的普通正弦振动。

也可不用先猜行波，而从两种二重轴的正负特征构造实模式。为方便分类，在对顶点轴这一类中选穿过$j=2,5$的轴，记其操作为$D_2$，于是$(D_2x)_j=-x_{4-j}$。它与上面的$C_2$轴垂直，两操作可交换。联立$C_2\vect x=\sigma\vect x$和$D_2\vect x=\tau\vect x$，其中$\sigma,\tau=\pm1$，可以选出
\begin{center}\small
\begin{tabular}{ccl}
\toprule $\sigma$&$\tau$&一个代表向量\\\midrule
$-1$&$-1$&$(1,1,1,1,1,1)$\\
$+1$&$-1$&$(1,-1,1,-1,1,-1)$\\
$-1$&$+1$&$(1,1,0,-1,-1,0)$\\
$+1$&$+1$&$(1,-1,0,1,-1,0)$\\\bottomrule
\end{tabular}\end{center}
将第一分量代回$m\ddot x_0=k(x_1+x_5-2x_0)$：均匀向量给右端为零；交替向量给$-4kx_0$；后两个向量分别给$-kx_0$与$-3kx_0$。其余分量也须代入验证，例如后两种的静止点$j=2,5$，两邻点位移相反，合回复力确为零，因而可以一直静止。

前两种向量各自也是$C_6$的本征向量，分别给本征值$+1$和$-1$。后两种经$C_6$作用却变成不成比例的新向量，仍有同一频率，故必须补入各自的简并伙伴。单凭二重轴的两个正负号，不能把每个代表态都当成独立的单重频率。

用$+$、$-$、0表示相同绝对值的正、负和零位移，下列六个实向量张成全部坐标空间：
\begin{center}\small
\begin{tabular}{cll}
\toprule
$m\omega^2/k$&代表向量&独立模式数\\\midrule
0&$(1,1,1,1,1,1)$&1\\
1&$(1,1,0,-1,-1,0)$，$(0,1,1,0,-1,-1)$&2\\
3&$(1,-1,0,1,-1,0)$，$(0,1,-1,0,1,-1)$&2\\
4&$(1,-1,1,-1,1,-1)$&1\\\bottomrule
\end{tabular}\end{center}
这组简并向量可再正交化，但不是必须先正交才能验证模式。两种二重模的代表向量分别满足
\[
 \begin{aligned}
 \vect a&=(1,1,0,-1,-1,0),&\mathbf C^2\vect a&=\mathbf C\vect a-\vect a,\\
 \vect b&=(1,-1,0,1,-1,0),&\mathbf C^2\vect b&=-\mathbf C\vect b-\vect b.
 \end{aligned}
\]
两种关系符号不同，分别对应单位根实部$1/2$与$-1/2$。图\ref{ch10:fig:ring-modes}画出两组代表态及旋转后的伙伴，箭头严格沿所在点的圆周切线；不动点明确保留。
\chTenRingModesFigure
交替向量经$C_6$后变号，故其$C_6$本征值为$-1$；均匀向量的本征值为$+1$。对于二重简并，某一个实向量通常会被旋转为同频子空间中的另一向量，而不是自身乘一个实数。对称操作的乘法结构及其在线性空间中的作用，分别属于群论和群表示论；上述六环是它们在分子与固体振动中的一个具体入口。

\subsection{真实分子的振型与光谱}
一个非线性的12原子分子有$3\times12=36$个位置自由度。除去3个整体平动和3个整体转动，剩30个内部振动自由度。苯的对称性使其中10对两两简并，故可分成20组基频。图\ref{ch10:fig:benzene-modes-a}至图\ref{ch10:fig:benzene-modes-d}逐组列出相应振型；每个双重组有两幅独立图，不能把30个振动自由度误数成20个。

读图时，深色大圆为碳，浅色小圆为氢；平面内箭头表示一个选定相位的主要位移方向，$\odot$表示出纸面，$\otimes$表示入纸面。半个周期后所有位移同时反向。这里给出振型的定性示意，箭头不按本征矢量的分量比例绘制，不能用箭头长度计算质心位移或动能；精确振型还须由力常数矩阵求出。波数$\widetilde\nu$的单位为$\mathrm{cm^{-1}}$，与频率的关系为$\nu=c\widetilde\nu$、$\omega=2\pi\nu$；模型参数、同位素和测量条件会影响数值，不能把图上近似基频当成无误差常数。A、B表示单重，E表示二重；g、u表示相应反演对称性，进一步的光谱选择定则需结合电偶极矩或极化率变化。
图\ref{ch10:fig:benzene-modes-a}中的点、叉可用来区分相邻原子的同向和反向离面运动；图\ref{ch10:fig:benzene-modes-b}中，$1014\,\mathrm{cm^{-1}}$一组的各原子同时向环心收拢，而$994\,\mathrm{cm^{-1}}$一组的相邻单元交替伸缩。图\ref{ch10:fig:benzene-modes-c}中$1332\,\mathrm{cm^{-1}}$一组的碳、氢切向运动相反，不能看成整个分子刚性转动。图\ref{ch10:fig:benzene-modes-d}中较高频的位移以碳氢键方向上的伸缩为主；双重组的两幅图应分别看作两个独立振动坐标。
% Source figure 4.11, physical page 93: twenty groups, thirty panels.
% Arrows are qualitative normal-pattern directions, not computed amplitudes.
\newcommand{\chTenMolecule}[4]{%
\begin{tikzpicture}[scale=.55,>=Latex]
\foreach \j in {0,...,5}{\draw[TextGray]({60*\j}:1)--({60*(\j+1)}:1);\draw[TextGray]({60*\j}:1)--({60*\j}:1.7);}
\foreach \j in {0,...,5}{\draw[DeepBlue,thick]({60*\j}:1)circle(.10);\draw[TextGray]({60*\j}:1.7)circle(.075);}
\ifnum#1=0
 \foreach \v[count=\j from 0] in {#2}{\ifnum\v=0\else
 \node[font=\tiny,inner sep=0,fill=none]at({60*\j}:1){\ifnum\v>0$\cdot$\else$\times$\fi};\fi}
 \foreach \v[count=\j from 0] in {#3}{\ifnum\v=0\else
 \node[font=\tiny,inner sep=0,fill=none]at({60*\j}:1.7){\ifnum\v>0$\cdot$\else$\times$\fi};\fi}
\else
 \foreach \ang/\mag[count=\j from 0] in {#2}{\ifdim\mag pt=0pt\else\draw[DeepBlue,-{Latex[length=.8mm]},line width=.45pt]({60*\j}:1)++(\ang:.12)--++(\ang:{.45*max(.6,\mag)});\fi}
 \foreach \ang/\mag[count=\j from 0] in {#3}{\ifdim\mag pt=0pt\else\draw[DeepBlue,-{Latex[length=.8mm]},line width=.45pt]({60*\j}:1.7)++(\ang:.09)--++(\ang:{.42*max(.6,\mag)});\fi}
\fi
\path[use as bounding box](-2.5,-2.4)rectangle(2.5,2.4);
\end{tikzpicture}%
}
\newcommand{\chTenModeGroup}[2]{\begin{minipage}[t]{.47\linewidth}\centering #2\par\smallskip{\small #1}\end{minipage}}
\newcommand{\chTenZeroVectors}{0/0,0/0,0/0,0/0,0/0,0/0}
\begin{chTenFigure}
\chTenModeGroup{$395\,\mathrm{cm^{-1}}\ (E_{2u})$}{
\chTenMolecule{0}{-1,1,1,-1,1,1}{-1,1,1,-1,1,1}{}
\chTenMolecule{0}{1,1,-1,1,1,-1}{1,1,-1,1,1,-1}{}}
\hfill
\chTenModeGroup{$599\,\mathrm{cm^{-1}}\ (E_{2g})$}{
\chTenMolecule{1}{0/1,240/.6,300/.6,180/1,60/.6,120/.6}{0/1,240/.6,300/.6,180/1,60/.6,120/.6}{}
\chTenMolecule{1}{0/0,60/1,300/1,0/0,240/1,120/1}{0/0,90/1,270/1,0/0,270/1,90/1}{}}
\par\medskip
\chTenModeGroup{$677\,\mathrm{cm^{-1}}\ (A_{2u})$}{
\chTenMolecule{0}{1,1,1,1,1,1}{-1,-1,-1,-1,-1,-1}{}}
\hfill
\chTenModeGroup{$700\,\mathrm{cm^{-1}}\ (B_{2g})$}{
\chTenMolecule{0}{1,-1,1,-1,1,-1}{1,-1,1,-1,1,-1}{}}
\par\medskip
\chTenModeGroup{$838\,\mathrm{cm^{-1}}\ (E_{1g})$}{
\chTenMolecule{0}{-1,-1,-1,1,1,1}{1,1,1,-1,-1,-1}{}
\chTenMolecule{0}{1,1,-1,-1,-1,1}{-1,-1,1,1,1,-1}{}}
\figcap{苯分子振型（一）：五组、八个独立模式。离面模式用圆内点、叉标示；未标箭头的方向不表示该原子永远静止。}\label{ch10:fig:benzene-modes-a}
\end{chTenFigure}

\begin{chTenFigure}
\chTenModeGroup{$954\,\mathrm{cm^{-1}}\ (E_{2u})$}{
\chTenMolecule{0}{-1,1,-1,-1,1,-1}{1,-1,1,1,-1,1}{}
\chTenMolecule{0}{-1,1,1,-1,1,1}{1,-1,-1,1,-1,-1}{}}
\hfill
\chTenModeGroup{$981\,\mathrm{cm^{-1}}\ (B_{2g})$}{
\chTenMolecule{0}{1,-1,1,-1,1,-1}{-1,1,-1,1,-1,1}{}}
\par\medskip
\chTenModeGroup{$994\,\mathrm{cm^{-1}}\ (B_{2u})$}{
\chTenMolecule{1}{0/1,240/1,120/1,0/1,240/1,120/1}{0/1,240/1,120/1,0/1,240/1,120/1}{}}
\hfill
\chTenModeGroup{$1014\,\mathrm{cm^{-1}}\ (A_{1g})$}{
\chTenMolecule{1}{180/1,240/1,300/1,0/1,60/1,120/1}{180/.6,240/.6,300/.6,0/.6,60/.6,120/.6}{}}
\par\medskip
\chTenModeGroup{$1048\,\mathrm{cm^{-1}}\ (E_{1u})$}{
\chTenMolecule{1}{0/.7,0/0,0/0,0/.7,0/0,0/0}{180/1,315/1,45/1,180/1,315/1,45/1}{}
\chTenMolecule{1}{0/0,240/.6,300/.6,0/0,240/.6,300/.6}{90/1,180/1,0/1,90/1,180/1,0/1}{}}
\hfill
\chTenModeGroup{$1137\,\mathrm{cm^{-1}}\ (B_{1u})$}{
\chTenMolecule{1}{0/0,0/0,0/0,0/0,0/0,0/0}{90/1,330/1,210/1,90/1,330/1,210/1}{}}
\figcap{苯分子振型（二）：六组、八个独立模式。不同基频的同名对称类型仍是不同模。}\label{ch10:fig:benzene-modes-b}
\end{chTenFigure}

\begin{chTenFigure}
\chTenModeGroup{$1162\,\mathrm{cm^{-1}}\ (E_{2g})$}{
\chTenMolecule{1}{0/0,0/0,0/0,0/0,0/0,0/0}{0/0,330/1,210/1,0/0,150/1,30/1}{}
\chTenMolecule{1}{0/0,0/0,0/0,0/0,0/0,0/0}{270/1,150/1,150/1,90/1,330/1,330/1}{}}
\hfill
\chTenModeGroup{$1332\,\mathrm{cm^{-1}}\ (A_{2g})$}{
\chTenMolecule{1}{270/.12,330/.12,30/.12,90/.12,150/.12,210/.12}{90/1,150/1,210/1,270/1,330/1,30/1}{}}
\par\medskip
\chTenModeGroup{$1376\,\mathrm{cm^{-1}}\ (B_{1u})$}{
\chTenMolecule{1}{90/1,330/1,210/1,90/1,330/1,210/1}{0/0,0/0,0/0,0/0,0/0,0/0}{}}
\hfill
\chTenModeGroup{$1477\,\mathrm{cm^{-1}}\ (E_{1u})$}{
\chTenMolecule{1}{180/.4,180/1,0/1,180/.4,180/1,0/1}{180/1,315/1,45/1,180/1,315/1,45/1}{}
\chTenMolecule{1}{270/1,0/0,0/0,270/1,0/0,0/0}{90/1,0/0,0/0,90/1,0/0,0/0}{}}
\par\medskip
\chTenModeGroup{$1612\,\mathrm{cm^{-1}}\ (E_{2g})$}{
\chTenMolecule{1}{0/1,240/1,300/1,180/1,60/1,120/1}{270/1,0/0,180/1,90/1,0/0,0/1}{}
\chTenMolecule{1}{0/0,180/1,0/1,0/0,0/1,180/1}{0/0,150/1,30/1,0/0,330/1,210/1}{}}
\figcap{苯分子振型（三）：五组、八个独立模式。此组以平面内弯曲和环骨架形变为主。}\label{ch10:fig:benzene-modes-c}
\end{chTenFigure}

\begin{chTenFigure}
\chTenModeGroup{$3080\,\mathrm{cm^{-1}}\ (B_{2u})$}{
\chTenMolecule{1}{0/0,0/0,0/0,0/0,0/0,0/0}{0/1,240/1,120/1,0/1,240/1,120/1}{}}
\hfill
\chTenModeGroup{$3092\,\mathrm{cm^{-1}}\ (E_{2g})$}{
\chTenMolecule{1}{0/0,0/0,0/0,0/0,0/0,0/0}{0/0,240/1,120/1,0/0,60/1,300/1}{}
\chTenMolecule{1}{0/0,0/0,0/0,0/0,0/0,0/0}{0/1,240/.6,300/.6,180/1,60/.6,120/.6}{}}
\par\medskip
\chTenModeGroup{$3110\,\mathrm{cm^{-1}}\ (E_{1u})$}{
\chTenMolecule{1}{0/0,0/0,0/0,0/0,0/0,0/0}{0/0,240/1,300/1,0/0,240/1,300/1}{}
\chTenMolecule{1}{0/0,0/0,0/0,0/0,0/0,0/0}{0/1,0/0,0/0,0/1,0/0,0/0}{}}
\hfill
\chTenModeGroup{$3121\,\mathrm{cm^{-1}}\ (A_{1g})$}{
\chTenMolecule{1}{0/0,0/0,0/0,0/0,0/0,0/0}{0/1,60/1,120/1,180/1,240/1,300/1}{}}
\figcap{苯分子振型（四）：四组、六个独立模式，以碳氢键伸缩为主。四张图合计20组、30个模式；未画出的很小位移不可据此当成精确零。}\label{ch10:fig:benzene-modes-d}
\end{chTenFigure}

例如对称伸缩时，若分子的电偶极矩不随该坐标改变，即使外光频率相等，也未必产生电偶极红外吸收；有些模主要通过拉曼散射观察。频率相合是共振的条件之一，不是“每个模都必有同样强的红外谱线”。

经验势能或微观计算给出平衡点的力常数矩阵，经质量加权对角化便得到$\omega_i$和振型。量子描述中，每个近似独立模的能级为
\[
 E_{n_i}=\left(n_i+\frac12\right)\hbar\omega_i,\qquad n_i=0,1,2,\ldots.
\]
相邻能级相差$\hbar\omega_i$，故允许的单量子跃迁满足$h\nu=\hbar\omega_i$，即$\nu=\omega_i/(2\pi)$。这把经典简正频率与分子光谱联系起来；非谐性还会改变能级间隔并引入组合频率等现象。

\section{非线性振动与摄动（选读）}\label{ch10:sec:nonlinear}
保留势能的高阶项后，回复力不再线性，叠加原理也不能原样使用。一个保留左右对称性的例子是杜芬方程（Duffing equation）
\[
 \ddot x+\omega_0^2x+\varepsilon\alpha x^3=0,
 \qquad U=\frac12m\omega_0^2x^2+\frac14m\varepsilon\alpha x^4.
\]
$\varepsilon$用来标记小修正阶数，要求$|\varepsilon\alpha A^2|\ll\omega_0^2$。若仍把频率固定为$\omega_0$，摄动方程会产生$t\sin\omega_0t$这样的长期项，时间稍长就失去一致性。为避免它，令$\tau=\Omega t$，$\Omega=\omega_0+\varepsilon\Omega_1+\cdots$，$x=A\cos\tau+\varepsilon x_1(\tau)+\cdots$。

因为$\ddot x=\Omega^2\dd^2x/\dd\tau^2$，收集一阶项得到
\[
 \omega_0^2(x_1''+x_1)=2\omega_0\Omega_1 A\cos\tau-\alpha A^3\cos^3\tau.
\]
利用$\cos^3\tau=(3\cos\tau+\cos3\tau)/4$，右端的基频系数是$2\omega_0\Omega_1A-3\alpha A^3/4$。令它为零消去共振驱动，得到
\[
 \Omega_1=\frac{3\alpha A^2}{8\omega_0},\qquad
 \boxed{\Omega\simeq\omega_0\left(1+\frac{3\varepsilon\alpha A^2}{8\omega_0^2}\right)}.
\]
余下$x_1''+x_1=-\alpha A^3\cos3\tau/(4\omega_0^2)$，特解为$\alpha A^3\cos3\tau/(32\omega_0^2)$。若把$A$定义为$t=0$的转向位移，还应加齐次项使$x_1(0)=0$，从而
\[
 x\simeq A\cos\tau+
 \frac{\varepsilon\alpha A^3}{32\omega_0^2}(\cos3\tau-\cos\tau).
\]
这明确显示三次谐波。$\varepsilon\alpha>0$为硬化，振幅增大时频率升高；$\varepsilon\alpha<0$为软化，只能在平衡附近和模型有效范围讨论，不能外推到负四次势能导致的远处失稳区域。令非线性参数趋于零，以上结果恢复简谐振动。

另一类现象来自参数调制，例如
\[
 \ddot x+\omega_0^2[1+h\cos(\Omega t)]x=0,\qquad |h|\ll1.
\]
外界改变的是回复系数，而不是直接加同频外力。在$\Omega=2\omega_0$附近，设$x=a(t)\cos\omega_0t+b(t)\sin\omega_0t$，$a,b$缓变，忽略其二阶时间导数。保留右端乘积中的基频项，得到
\[
 2\omega_0\dot b=-\frac12h\omega_0^2a,
 \qquad -2\omega_0\dot a=\frac12h\omega_0^2b.
\]
即$\dot a=-h\omega_0b/4$、$\dot b=-h\omega_0a/4$，两种包络组合有指数率$\pm h\omega_0/4$；一般初始微扰包含增长分量。这是主参数共振的近似机制。若另有弱黏性阻尼，增长须克服$\gamma$，在精确调谐的一阶近似下要求$|h|\omega_0/4>\gamma$；频率失谐也会缩小不稳定区。严格零位移、零速度仍是齐次方程的零解，不会在没有扰动时凭空起振。
\subsection{小修正为什么会在长时间后变得重要}
摄动不是把某个小项直接删去，而是在一个明确的小参数下，追踪它对主运动的首阶影响。以本节的三次回复项为例，在振幅 $A$ 附近，非线性加速度与线性加速度的比为 $|\varepsilon\alpha|A^2/\omega_0^2$。只有这个无量纲比值很小，才允许把运动看成简谐振动加小修正。

如果固定原频率写 $x=x_0+\varepsilon x_1$，其中 $x_0=A\cos\omega_0t$，一阶方程会是
\[
 \ddot x_1+\omega_0^2x_1
 =-\frac{3\alpha A^3}{4}\cos\omega_0t
  -\frac{\alpha A^3}{4}\cos3\omega_0t.
\]
第一项恰以系统本征频率驱动修正量。验证
\[
 \left(\frac{\dd^2}{\dd t^2}+\omega_0^2\right)
 [t\sin\omega_0t]=2\omega_0\cos\omega_0t,
\]
便知其中包含 $-3\alpha A^3t\sin\omega_0t/(8\omega_0)$。即使 $\varepsilon$ 很小，乘上足够长的时间后也可能不再小，这就是长期项。
它不表示保守振子真的无界增长，而是说明选错了比较基准。若真实频率为 $\omega_0+\delta\omega$，则
\[
 A\cos[(\omega_0+\delta\omega)t]
 \simeq A\cos\omega_0t-A\delta\omega\,t\sin\omega_0t.
\]
相位误差积累就产生同样形式的长期项。把频率也纳入待求量，是把这部分相位变化吸收到主运动中，而不是人为删除真实物理效应。

\begin{origexamplebox}{弱硬化弹簧的频率与适用范围}\label{ch10:ex:teaching-hardening}
质点质量为 $m$，在光滑直线上受到回复力 $F=-kx-\lambda x^3$，$k,\lambda>0$。从 $x=A$ 静止释放，且 $\lambda A^2/k\ll1$。求振动频率的首阶修正；若要求频率相对简谐值的修正不超过 $1\%$，振幅应满足什么近似条件？

\ThQHardening
\end{origexamplebox}
\Needspace{6\baselineskip}\begin{xesolution}
\ThAHardening

由牛顿定律，$\ddot x+(k/m)x+(\lambda/m)x^3=0$，所以 $\omega_0=\sqrt{k/m}$，而本节的一阶系数对应 $\varepsilon\alpha=\lambda/m$。代入已推导的消长期项条件，得到
\[
 \Omega\simeq\sqrt{\frac km}\left(1+\frac{3\lambda A^2}{8k}\right),
 \qquad T\simeq2\pi\sqrt{\frac mk}\left(1-\frac{3\lambda A^2}{8k}\right).
\]
周期使用 $(1+q)^{-1}\simeq1-q$，不能照抄频率的正号。回复力在较大位移处比线性弹簧更强，频率升高、周期缩短与物理判断相符。
首阶相对修正为 $3\lambda A^2/(8k)$，故近似要求
\[
 A\lesssim\sqrt{\frac{0.08k}{3\lambda}}.
\]
该边界使 $\lambda A^2/k\lesssim0.0267$，与小参数假设一致，但这是首阶误差估计，不是包含全部高阶项的严格误差界。由 $[\lambda]=\mathrm{N/m^3}$ 可知 $k/\lambda$ 为长度平方；$A\to0$ 时恢复线性结果。
\end{xesolution}

\subsection{参数共振的能量从哪里来}
受迫振动是外界施加随时间变化的力；参数共振则是外界改变弹簧劲度、摆长等系统参数。对 $k(t)=m\omega_0^2[1+h\cos(\Omega t)]$，瞬时机械能写作
\[
 E=\frac12m\dot x^2+\frac12k(t)x^2.
\]
对时间求导并代入 $m\ddot x+k(t)x=0$，得到
\[
 \dot E=\dot x[m\ddot x+k(t)x]+\frac12\dot k(t)x^2
 =\frac12\dot k(t)x^2.
\]
所以能量由调节参数的外界装置提供，并未违反守恒。若在位移大时增大劲度、位移小时减小劲度，平均上可以输入能量。一次振动有两次位移绝对值极大，这提示主参数共振出现在约两倍本征频率处，但精确增长率仍须由缓变振幅方程计算。
加入阻力 $-2m\gamma\dot x$ 后，还要减去 $2m\gamma\dot x^2$ 的耗散功率。只有平均输入克服耗散，振幅才增长；超出小振幅范围后，非线性、参数调制范围及装置反馈也会限制增长，线性指数解不能无限外推。


\section{格波、声学支与光学支（选读）}\label{ch10:sec:lattice}
\subsection{由原子小振动理解固体中的波}
晶体中的原子实携带主要质量，在平衡格点附近通过相互作用共同运动。把相互作用在平衡附近线性化，就得到很多耦合振子；一个位置的扰动可以传到相邻位置。相位有组织的长波运动体现为宏观声波，多种模式按热统计规律被占据则体现为热运动。并非每个原子都是一个独立振子，真正独立的是简正模。

声速与热导率都受晶格影响，但二者不能简单互相推出：声速与弹性和密度有关，热导率还取决于热容、散射和平均自由程，金属中电子常有重要贡献。例如金刚石兼有较高声速和热导率，而晶体方向、温度、纯度及缺陷会改变数值和排序。讨论典型量级时，固体声速约$10^3$—$10^4\,\mathrm{m/s}$；热导率则跨越很宽范围，不能据“声速更大”断言“一定更导热”。

量子化的晶格简正模称声子（phonon）。爱因斯坦固体模型先用同频量子振子解释经典热容失效，德拜模型进一步考虑低频声学模及其频率分布。下文先用经典力学求出这些模式；量子占据和统计平均在热学及近代物理中继续使用。

\subsection{单原子链：运动方程、色散与截止}
图\ref{ch10:fig:mono-chain}中，相邻平衡位置相距$a$，各质量为$m$，相邻弹簧劲度为$k$。$u_j$表示第$j$个原子从自身平衡位置沿链向右的位移，不是它的绝对位置。右、左弹簧的伸长分别为$u_{j+1}-u_j$、$u_j-u_{j-1}$，所以
\begin{equation}\label{ch10:eq:mono-motion}
 m\ddot u_j=k(u_{j+1}+u_{j-1}-2u_j).
\end{equation}
\chTenMonoChainFigure
平移一个格距后方程不变，故设相邻原子的复振幅相差相位$qa$：
\[
 u_j=\Re\{A\mathrm e^{\ii(qja-\omega t)}\}.
\]
$q$是一维波矢，单位为$\mathrm{m^{-1}}$；这里取$-\omega t$约定，所以$q>0$的长波沿右方传播。代入方程，约去公共指数：
\[
 -m\omega^2=k(\mathrm e^{\ii qa}+\mathrm e^{-\ii qa}-2)
 =-4k\sin^2(qa/2),
\]
从而色散关系为
\begin{equation}\label{ch10:eq:mono-dispersion}
 \omega(q)=2\sqrt{k/m}\,|\sin(qa/2)|.
\end{equation}
因为$j$是整数，$q$与$q+2\pi/a$给完全相同的格点位移。只需取一个不重复区间，例如$-\pi/a<q\leq\pi/a$，称第一布里渊区（Brillouin zone）。图\ref{ch10:fig:mono-dispersion}标出区中心和两端；两端代表相同的格点相位关系，不应算成两种独立状态。
\chTenMonoDispersionFigure
最大频率$\omega_c=2\sqrt{k/m}$称截止频率。为了说明“截止”不是凭图猜测，可改设振幅比为$\lambda$，$u_j\propto\lambda^j\mathrm e^{-\ii\omega t}$，则
\[
 \lambda+\lambda^{-1}=2-\frac{m\omega^2}{k}.
\]
在允许频带内可取$|\lambda|=1$，即传播型相位因子。当$\omega>\omega_c$，两根为负实数且互为倒数，其中一根绝对值小于1。半无限链受边界策动时，可选向内部衰减的根，产生交替符号的倏逝响应；不能在整个无限链中接受另一侧无界增长的解作为有限能量体波。恰在截止点$\lambda=-1$，相邻原子反相、群速度为零，仍是允许的边界模，不是“根本不存在振动”。

\subsection{长波、短波与传播速度}
在$q>0$的第一布里渊区内，相速度和群速度分别为
\[
 v_p=\frac\omega q=\frac{2\sqrt{k/m}\sin(qa/2)}q,
 \qquad v_g=\frac{\dd\omega}{\dd q}=a\sqrt{k/m}\cos(qa/2).
\]
相速度描述相位面的运动，群速度描述窄带波包包络的传播，在此无耗散链模型中也与能量传播有关。长波$|qa|\ll1$时，$\sin(qa/2)\simeq qa/2$，故$\omega\simeq a\sqrt{k/m}|q|$，二者大小都趋于$c=a\sqrt{k/m}$。连续近似中，单位长度质量$\rho_1=m/a$、轴向刚度$C_1=ka$，恰有$c=\sqrt{C_1/\rho_1}$，与连续弹性波一致。

短波极限$|q|a=\pi$对应最短独立波长$2a$，相邻粒子反相，$v_g=0$。若在扩展区间继续增加$q$，色散曲线斜率可以反号；此时波包传播方向应看$\dd\omega/\dd q$，不能只看$q$符号。区边界并非通常晶体学中称为$\Gamma$的点；$\Gamma$通常表示区中心$q=0$。

\subsection{双原子链：一个晶胞内有两个坐标}
图\ref{ch10:fig:diatomic-chain}中质量$M,m$交替排列，假设$M\geq m$。每条最近邻弹簧劲度为$k$，\emph{晶胞长度}定义为$a$，相邻原子间距则为$a/2$。设第$j$胞中重、轻原子的位移分别为$u_j,v_j$，则
\begin{align}
 M\ddot u_j&=k(v_j+v_{j-1}-2u_j),\label{ch10:eq:diatomic-M}\\
 m\ddot v_j&=k(u_j+u_{j+1}-2v_j).\label{ch10:eq:diatomic-m}
\end{align}
\chTenDiatomicChainFigure
取$u_j=A\mathrm e^{\ii(qja-\omega t)}$、$v_j=B\mathrm e^{\ii(qja-\omega t)}$的实部；晶胞内部的位置相位并入复振幅$B$，不能再重复加一次。代入得
\[
 \begin{pmatrix}
 2k-M\omega^2&-k(1+\mathrm e^{-\ii qa})\\
 -k(1+\mathrm e^{\ii qa})&2k-m\omega^2
 \end{pmatrix}\begin{pmatrix}A\\B\end{pmatrix}=0.
\]
行列式为零，使用$(1+\mathrm e^{\ii qa})(1+\mathrm e^{-\ii qa})=4\cos^2(qa/2)$，得
\[
 Mm\omega^4-2k(M+m)\omega^2+4k^2\sin^2(qa/2)=0.
\]
这是关于$\omega^2$的二次方程，两根为
\begin{equation}\label{ch10:eq:diatomic-dispersion}
 \omega_\pm^2=k\left(\frac1M+\frac1m\right)
 \pm k\sqrt{\left(\frac1M+\frac1m\right)^2-
 \frac4{Mm}\sin^2\frac{qa}{2}}.
\end{equation}
负号支为声学支，正号支为光学支。式中根号项与前项同为$\mathrm{s^{-2}}$，最外层结果是频率平方，不能漏掉最后开方。

\subsection{两支的端点、带隙及轻质量极限}
在$q=0$，下支$\omega_-=0$且$A=B$，整个晶体平移；上支
$\omega_+^2=2k(1/M+1/m)$，第一行给$MA+mB=0$，两原子反向且胞内质心不动。在区边界$qa=\pi$，耦合矩阵非对角元为零，故下支$\omega_-^2=2k/M$时只有重原子运动，上支$\omega_+^2=2k/m$时只有轻原子运动。图\ref{ch10:fig:diatomic-dispersion}把这些端点与带隙一并标出。
\chTenDiatomicDispersionFigure
当$M>m$，$\sqrt{2k/M}<\omega<\sqrt{2k/m}$之间没有无限周期链的传播型体模，称频率禁带；边界受迫问题仍可有倏逝响应。当$M=m$时禁带闭合，两条支在缩小的布里渊区边界相接，实质上是单原子链色散折叠后的表示，并非多出新的物理自由度。

令$\beta=m/M\ll1$。低频支中轻原子的惯性可忽略，它近似始终位于两侧重原子位移的平均位置，相当于两根$k$弹簧串联成$k/2$，故
\[
 \omega_-\simeq\sqrt{2k/M}\,|\sin(qa/2)|.
\]
高频支主要由轻原子振动，展开精确根式到$\beta$一阶，得
\[
 \omega_+^2\simeq\frac{2k}{m}\left[1+\beta\cos^2(qa/2)\right],
 \qquad\omega_+\simeq\sqrt{2k/m}\sqrt{1+\beta\cos^2(qa/2)}.
\]
低频的长波波速更一般地为$c=a\sqrt{k/[2(M+m)]}$；$m\ll M$时恢复刚才的串联模型。高频$q=0$时相邻晶胞同相但胞内两原子反向，因此“晶胞之间没有相位差”并不等于“没有内部振动”。在两种原子带相反有效电荷的晶体中，该运动可与光场耦合，这说明“光学支”的名称；并非所有晶体的每一光学模都有相同红外活性。

声学支的低频模和日常声音相连，光学支则是多原子晶胞内部振动的延续。每个允许$q$对应两个独立模；$N$个晶胞在周期边界条件下总计$2N$个自由度，恰与$2N$个原子的一维坐标数相等。

\subsection{声子频带与电子能带的联系和区别}
色散的“带”首先是经典周期结构允许的模式频率集合。把每个格波量子化后，一个模每增加一个声子，能量增加$\hbar\omega(q)$；在热平衡中各模的占据不同，其总能量和温度依赖决定晶格热容等性质。热传导还需考虑这些激发被缺陷、边界和其他激发散射的过程，单靠无耗散谐链并不能给出真实材料的有限热导率。

周期势中的电子也有由波矢标记的允许能量带和禁带，但它来自电子波动方程，不能把格波的$\omega(q)$直接当成电子能量。电子服从费米统计，声子作为集体激发服从玻色统计且数目通常不守恒。电子在振动晶格中运动，二者的耦合既影响输运和能量交换，也构成进一步讨论电子—声子相互作用的基础。共同点是周期性组织了本征模式；粒子性质、动力学方程和占据规律仍须分别处理。


\Needspace{18\baselineskip}
\section{综合问题：振动系统的建模}

\begin{origexamplebox}{摩擦面上的两物块与弹簧}\label{ch10:ex:friction}
质量分别为 $M=3.6\,\mathrm{kg}$ 和 $m=1.8\,\mathrm{kg}$ 的物块 $A,B$，连接在原长
$l_0=0.30\,\mathrm m$、劲度系数 $k=24\,\mathrm{N/m}$ 的弹簧两端。两物块与水平面间的动、静摩擦因数均取
$\mu=0.30$，$g=10\,\mathrm{m/s^2}$。开始时弹簧长 $l=1.20\,\mathrm m$，系统由静止释放。
求两物块相碰前瞬间的速度和从释放到相碰的时间。物块尺寸忽略，碰撞发生在弹簧长度降为零时。

初始装置见图~\ref{ch10:fig:friction-problem}。
\chTenFrictionProblemFigure
\end{origexamplebox}
\Needspace{6\baselineskip}\begin{xesolution}
取水平向右为正，图\ref{ch10:fig:friction}中$A$在左、$B$在右，$r=x_B-x_A$。初始弹力$k(l-l_0)=21.6\,\mathrm N$，大于两者最大静摩擦力$10.8\,\mathrm N$、$5.4\,\mathrm N$，故释放后$A$向右、$B$向左。
\chTenFrictionFigure
第一阶段两者均滑动。摩擦力对$A$向左、对$B$向右；有符号弹力用$k(r-l_0)$表示，即使弹簧后来压缩也适用：
\[
 M\ddot x_A=k(r-l_0)-\mu Mg,\qquad
 m\ddot x_B=-k(r-l_0)+\mu mg.
\]
令
\[
 \omega^2=k\left(\frac1M+\frac1m\right)=20\,\mathrm{s^{-2}}.
\]
由两物块方程相减，得到
\[
 \ddot r+\omega^2\left(r-l_0-\frac{2\mu g}{\omega^2}\right)=0.
\]
两运动方程相加，质心加速度为
\[
 \dot V=\frac{\mu g(m-M)}{M+m}=-1.00\,\mathrm{m/s^2},
 \qquad V=-(1.00\,\mathrm{m/s^2})t.
\]
相对运动平衡值为$r_e=l_0+2\mu g/\omega^2=0.60\,\mathrm m$。由$r(0)=1.20\,\mathrm m,\dot r(0)=0$得到（下式长度数值以米计）
\[
 r=0.60+0.60\cos\omega t,\qquad
 \dot r=-0.60\omega\sin\omega t.
\]
所以
\[
 v_A=V-\frac{m}{M+m}\dot r
 =-(1.00\,\mathrm{m/s^2})t+(0.20\,\mathrm m)\omega\sin\theta,
 \qquad \theta=\omega t.
\]
$A$ 先于碰撞停下；首次非零停下时刻由
\[
 \theta_1=4\sin\theta_1
\]
确定。零根是释放瞬间，必须排除；函数$\theta-4\sin\theta$在$2.4$与$2.5$间变号，且此区间单调，二分迭代得到$\theta_1=2.47458$。于是
\[
 t_1=0.55333\,\mathrm s,\qquad
 r_1=0.12860\,\mathrm m,\qquad
 \dot r_1=-1.6600\,\mathrm{m/s}.
\]

此时$r_1>0$，尚未相碰，$v_B=\dot r_1<0$。第二阶段弹簧对$A$向左，静摩擦力应向右，不再是第一阶段的方向。到相碰为止弹簧最大压缩力不超过 $kl_0=7.2\,\mathrm N$，
小于 $\mu Mg=10.8\,\mathrm N$，故静摩擦足以使 $A$ 保持静止。$B$ 单独向左运动，满足
\[
 \ddot r+\omega_B^2(r-r_B)=0,\qquad
 \omega_B=\sqrt{\frac{k}{m}},\qquad
 r_B=l_0+\frac{\mu mg}{k}=0.525\,\mathrm m.
\]
令第二阶段时间为 $\tau$，则
\[
 r=r_B+(r_1-r_B)\cos\omega_B\tau
 +\frac{\dot r_1}{\omega_B}\sin\omega_B\tau.
\]
取$r=0$的最小正根，得$\tau=0.09282\,\mathrm s$。对上式逐项求导，
\[
 v_B=-(r_1-r_B)\omega_B\sin\omega_B\tau+\dot r_1\cos\omega_B\tau,
\]
得到碰撞前
\[
 \boxed{v_A=0},\qquad
 \boxed{v_B=-1.0843\,\mathrm{m/s}},\qquad
 \boxed{t=t_1+\tau=0.64615\,\mathrm s}.
\]
负号表示向左。能量校验可由两物块位移求摩擦功：$A$滑动路程约$0.20405\,\mathrm m$，$B$至相碰的总路程约$0.99595\,\mathrm m$。初始弹性势能$9.72\,\mathrm J$等于末态弹性能$1.08\,\mathrm J$、$B$动能约$1.05816\,\mathrm J$与摩擦损耗约$7.58184\,\mathrm J$之和。若忽略$A$中途停下及静摩擦阶段，便会把相对运动按同一方程错误外推。
\end{xesolution}

\begin{origexamplebox}{飞球调速器的小振动}\label{ch10:ex:governor}
两个质量均为 $m$ 的小球由四根长为 $l$ 的轻杆连接上、下套筒；上套筒固定，下套筒质量为 $M$，可沿竖直轴无摩擦滑动。
装置绕轴以恒定角速度 $\omega_0$ 转动，忽略杆与套筒的转动惯量。求：（1）平衡时两套筒距离；
（2）在该平衡位置附近的竖直小振动角频率；（3）若轴改为自由转动、总角动量守恒，小振动角频率如何改变？

初始装置见图~\ref{ch10:fig:governor-problem}。
\chTenGovernorProblemFigure
\end{origexamplebox}
\Needspace{6\baselineskip}\begin{xesolution}
图\ref{ch10:fig:governor}画通过两球与转轴的瞬时竖直截面。每球位于两套筒中间高度，离转轴半径为$\rho$。取向下的套筒间距为$y$，由直角三角形得
\[
 \rho=\sqrt{l^2-\frac{y^2}{4}}.
\]
\chTenGovernorFigure
由$\rho^2+y^2/4=l^2$求导，得$\dot\rho=-y\dot y/(4\rho)$；每球的竖直速度为$\dot y/2$。两球在子午面内的径向、竖直动能加上下套筒动能为
\[
 T_y=m\left(\dot\rho^2+\frac14\dot y^2\right)+\frac12M\dot y^2
 =\frac12\left(M+\frac m2+\frac{my^2}{8\rho^2}\right)\dot y^2.
\]
利用$4\rho^2+y^2=4l^2$，化为
\[
 T_y=\frac12\left(M+\frac{ml^2}{2\rho^2}\right)\dot y^2.
\]
恒角速度由外部驱动维持时，两球转动动能为 $m\rho^2\omega_0^2$，重力势能为
$-(M+m)gy$。约化拉格朗日量为$T_y+m\rho^2\omega_0^2-U_g$，因此用于$y$运动的有效势能为
\[
 U_\omega=-(M+m)gy-m\omega_0^2(l^2-y^2/4).
\]
以$U_\omega'=0$求平衡，等价于
\[
 -\frac12m\omega_0^2y_0+(M+m)g=0,\qquad
 \boxed{y_0=\frac{2(M+m)g}{m\omega_0^2}}.
\]
存在该构形还要求$0<y_0<2l$，即$\omega_0^2>(M+m)g/(ml)$。记
\[
 \rho_0^2=l^2-\frac{y_0^2}{4},\qquad
 M_{\rm eff}=M+\frac{ml^2}{2\rho_0^2},
\]
由$U_\omega''=m\omega_0^2/2$读出回复系数。令$y=y_0+\eta$，变惯性引起的$\dot\eta^2$项为二阶，线性方程只保留平衡处惯性：$M_{\rm eff}\ddot\eta+(m\omega_0^2/2)\eta=0$，所以
\[
 \boxed{\Omega_{\omega}=\omega_0
 \sqrt{\frac{m}{2M_{\rm eff}}}}.
\]

轴自由时，总角动量 $J=2m\rho^2\omega$ 守恒。消去循环坐标后，有效势能多出
$J^2/(4m\rho^2)$，即$U_J=-(M+m)gy+J^2/(4m\rho^2)$。令$D=\rho^2=l^2-y^2/4$，逐次求导得
\[
 U_J'=-(M+m)g+\frac{J^2y}{8mD^2},\qquad
 U_J''=\frac{J^2}{8mD^2}\left(1+\frac{y^2}{D}\right).
\]
比较同一平衡构形、平衡角速度同为$\omega_0$时，$J=2mD_0\omega_0$，一阶平衡条件仍相同，而二阶曲率为
\[
 \frac{m\omega_0^2}{2}\left(1+\frac{y_0^2}{\rho_0^2}\right),
\]
故
\[
 \boxed{\Omega_J=\omega_0
 \sqrt{\frac{m}{2M_{\rm eff}}
 \left(1+\frac{y_0^2}{\rho_0^2}\right)}}.
\]
套筒向下运动会减小$\rho$，角动量守恒使转速增大，因此自由轴情形附加回复更强。两结果根号内均无量纲。自由轴时可用$T_y+U_J$守恒；恒速驱动时装置可能与驱动机构交换功，不能混用能量条件。
\end{xesolution}

\Needspace{14\baselineskip}
\begin{origexamplebox}{落链与零原长弹簧}\label{ch10:ex:chain}
原长为零、劲度系数为 $k$ 的弹簧悬挂线密度为 $\lambda$ 的长软链；链的下部平放在地板上。
天花板与地板距离为 $H$，重力加速度为 $g$。（1）求平衡时弹簧长和悬垂链长；（2）把悬垂链长从平衡值改变
$\delta_0$ 后静止释放，求小振动角频率；（3）计及链节在地板处非弹性拾取、沉积造成的能量损失，求弱衰减下的振幅包络。
初始装置见图~\ref{ch10:fig:chain-problem}。
\chTenChainProblemFigure
\end{origexamplebox}
\Needspace{6\baselineskip}\begin{xesolution}
图\ref{ch10:fig:chain}中$x$是悬垂链长，也是连接点离地高度，向上增加；弹簧长为$H-x$。$\dot x>0$时拾取地板上的静链，$\dot x<0$时运动链沉积。假设链足够长，地板上始终有静止余链，链节不反弹。
\chTenChainFigure
平衡时
\[
 k(H-x_0)=\lambda gx_0,\qquad
 \boxed{x_0=\frac{kH}{k+\lambda g}},\qquad
 \boxed{s_0=H-x_0=\frac{\lambda gH}{k+\lambda g}}.
\]
以地板为重力势能零面，悬垂链的重心在$x/2$，故
\[
 U(x)=\frac12k(H-x)^2+\frac12\lambda gx^2
 =U(x_0)+\frac12(k+\lambda g)(x-x_0)^2.
\]
上移时，进入移动部分的链节初速度为零，动量平衡为$\dd(\lambda x\dot x)/\dd t=k(H-x)-\lambda gx$，即
\[
 \lambda x\ddot x=-(k+\lambda g)(x-x_0)-\lambda\dot x^2\quad(\dot x>0).
\]
下移时，离开移动部分的链节带走动量，质量流项抵消求导中的$\lambda\dot x^2$，得到
\[
 \lambda x\ddot x=-(k+\lambda g)(x-x_0)\quad(\dot x<0).
\]
取$q=x-x_0$。小振动线性阶中，$\dot q^2$和$q\ddot q$为二阶，两式共同化为$\lambda x_0\ddot q+(k+\lambda g)q=0$。运动质量近似为$\lambda x_0$，故
\[
 \boxed{\omega=\sqrt{\frac{k+\lambda g}{\lambda x_0}}
 =\frac{k+\lambda g}{\sqrt{\lambda kH}}}.
\]

地板处单位时间转入或转出运动部分的链质量为$\lambda|\dot q|$；完全非弹性的速度改变每次损失
$\tfrac12\lambda|\dot q|^3\dd t$。令一个周期内 $q=A\cos\omega t$，则
\[
 \Delta E=-\frac12\lambda A^3\omega^2
 \int_0^{2\pi}|\sin u|^3\dd u
 =-\frac43\lambda A^3\omega^2.
\]
这里$\int_0^{2\pi}|\sin u|^3\dd u=4\int_0^{\pi/2}\sin^3u\dd u$，令$z=\cos u$，后一积分为$\int_0^1(1-z^2)\dd z=2/3$。损耗率也可通过对$E=\lambda x\dot x^2/2+U$求导、分别代入两条分段方程验证为$\dot E=-\lambda|\dot x|^3/2$。
又振动能$E-U(x_0)\simeq\tfrac12\lambda x_0\omega^2A^2$、$T=2\pi/\omega$，故$\lambda x_0\omega^2A\dot A\simeq\Delta E/T$，平均化得到
\[
 \frac{\dd A}{\dd t}=-\frac{2\omega}{3\pi x_0}A^2.
\]
积分$\dd(1/A)/\dd t=2\omega/(3\pi x_0)$，若以$\delta_0>0$表示初始偏离的大小，则
\[
 \boxed{A(t)=\frac{\delta_0}
 {1+\dfrac{2\omega\delta_0}{3\pi x_0}t}}.
\]
若$\delta_0$带符号，公式中改用$|\delta_0|$。包络按反比例式而非指数衰减，因为耗散率与$|\dot q|^3$成正比。要求$A\ll x_0$、$A\ll H-x_0$且每周期包络变化很小，连接点不到地板或天花板，余链未被提完。$\lambda g$与$k$同为$\mathrm{N/m}$，所得$\omega$为频率量纲。
\end{xesolution}

\Needspace{14\baselineskip}
\begin{origexamplebox}{弹性轻杆连接的两球与可分离斜劈}\label{ch10:ex:wedges}
两个完全相同的 $45^\circ$ 斜劈质量均为 $m$，放在光滑水平面上；两个质量均为 $m$ 的小球分别从两斜面的最高点静止释放。
两球由水平弹性轻杆连接，杆的自由长度为 $L_0$、劲度系数为 $k$，初始恰为自由长度。斜劈可以彼此分离，
球与斜面光滑接触。释放瞬间不计为一次“达到极小”；已知支持力第二次达到极小时两球恰好到达地面。
求支持力第二次达到极大时弹性杆的长度。

初始装置见图~\ref{ch10:fig:wedges-problem}。
\chTenWedgeProblemFigure
\end{origexamplebox}
\Needspace{6\baselineskip}\begin{xesolution}
图\ref{ch10:fig:wedges}两侧的斜面均严格为$45^\circ$。只看右半边，令斜劈向外位移为$X$，球沿斜面向左下滑$s$，并定义球相对初始位置的水平位移
\[
 q=X-s\cos45^\circ.
\]
\chTenWedgeFigure
杆长为$L=L_0+2q$。每劈速度为$\dot X=\dot q+\dot s/\sqrt2$，每球水平速度为$\dot q$、向下速度为$\dot s/\sqrt2$。把两劈和两球的动能相加，得到
\[
 T=m(\dot q+\dot s/\sqrt2)^2+m(\dot q^2+\dot s^2/2).
\]
杆形变为$2q$，两球下降均为$s/\sqrt2$，因此整理后
\[
 T=m\left(2\dot q^2+\sqrt2\dot q\dot s+\dot s^2\right),\qquad
 U=2kq^2-\sqrt2mgs.
\]
分别对$q,s$用拉格朗日方程，得到
\[
 4m\ddot q+\sqrt2m\ddot s+4kq=0,\qquad
 \sqrt2m\ddot q+2m\ddot s-\sqrt2mg=0.
\]
第二式给$\ddot s=(g-\ddot q)/\sqrt2$，代入第一式消去$s$：
\[
 3\ddot q+\frac{4k}{m}q+g=0.
\]
令
\[
 \omega=\sqrt{\frac{4k}{3m}},\qquad A=\frac{mg}{4k},
\]
利用 $q(0)=\dot q(0)=0$ 得
\[
 q=A(\cos\omega t-1),\qquad
 L=L_0+\frac{mg}{2k}(\cos\omega t-1).
\]
右球的竖直向上方程为$N/\sqrt2-mg=-m\ddot s/\sqrt2$，其中杆力水平，不贡献竖直分量。故
\[
 N=\sqrt2mg-m\ddot s=\frac m{\sqrt2}(g+\ddot q),
 \qquad\ddot q=-\frac g3\cos\omega t.
\]
代入得到单侧支持力
\[
 N=\frac{mg}{\sqrt2}\left(1-\frac13\cos\omega t\right).
\]
释放后不把 $t=0$ 计作一次“达到极小”，则第二次极大在 $\omega t=3\pi$，早于第二次极小
$\omega t=4\pi$；题给落地条件保证上述接触解在该时刻仍有效。此时 $\cos\omega t=-1$，所以
\[
 \boxed{L=L_0-\frac{mg}{k}}.
\]
全过程$N>0$，没有因负支持力提前脱离；$mg/k$是长度量纲。物理上还需$L_0>mg/k$，否则杆先缩到零长，模型失效。将两条拉格朗日方程分别乘$\dot q,\dot s$相加，可验$\dd(T+U)/\dd t=0$，斜劈运动已计入总动能，整个系统机械能守恒。
\end{xesolution}

\begin{origexamplebox}{球面上三个等量点电荷的简正模}\label{ch10:ex:charges}
三个质量均为 $m$、电荷量均为 $Q$ 的点电荷被限制在半径为 $R$ 的球面上，只能沿切平面运动，
相互作用为 $KQ^2/r^2$，其中 $K$ 为库仑常量。（1）描述稳定平衡构形；（2）数出转动和振动自由度；
（3）求小振动的不同角频率及简并度。

作为计数参照，非线性三原子水分子共有9个位置坐标，扣除整体平动3个、转动3个，剩3个内部振动：键角弯曲、两氢氧键对称伸缩和反对称伸缩。两种伸缩频率通常较接近，弯曲频率较低。本题同样要区分内部形变和整体转动。

\ThQSphereCharges
\end{origexamplebox}
\Needspace{6\baselineskip}\begin{xesolution}
令$d_{ij}=|\vect r_i-\vect r_j|$，由$|\vect r_i|=R$得
\[
 \sum_{i<j}d_{ij}^2=9R^2-|\vect r_1+\vect r_2+\vect r_3|^2\leq9R^2.
\]
函数$z^{-1/2}$在$z>0$严格凸，故
\[
 \sum_{i<j}\frac1{d_{ij}}\geq
 \frac3{\sqrt{\tfrac13\sum d_{ij}^2}}\geq\frac{\sqrt3}{R}.
\]
两步取等要求三边相等且$\sum_i\vect r_i=0$，因此三个电荷位于某一大圆上，夹角均为$120^\circ$，形成边长$\sqrt3R$的正三角形。这证明全局最低势能，而非仅凭“三对距离尽可能大”猜测。图\ref{ch10:fig:charges}取大圆为$xy$平面。
\chTenChargesFigure
每个电荷有两个切向自由度，共6个；球面限制了整体平移，整体构形在三维空间有3个转动自由度，
余下3个为振动自由度。

取平衡方位角$\phi_i=0,2\pi/3,4\pi/3$，沿大圆的小角位移为$a_i$，垂直大圆的无量纲位移为$z_i$。球面坐标可写为
\[
 \frac{\vect r_i}{R}=(\sqrt{1-z_i^2}\cos(\phi_i+a_i),
 \sqrt{1-z_i^2}\sin(\phi_i+a_i),z_i).
\]
对一对电荷记$\Delta=\phi_i-\phi_j$、$\delta=a_i-a_j$，点积展开为
\[
 c_{ij}=\frac{\vect r_i\cdot\vect r_j}{R^2}
 =-\frac12-\sin\Delta\,\delta+\frac14\delta^2
 +\frac14(z_i^2+z_j^2)+z_iz_j+O(3).
\]
设$c_{ij}=-1/2+\epsilon_{ij}$，则
\[
 \frac R{d_{ij}}=[2(1-c_{ij})]^{-1/2}
 =\frac1{\sqrt3}\left(1+\frac{\epsilon_{ij}}3+\frac{\epsilon_{ij}^2}{6}\right)+O(3).
\]
二阶$\delta^2$系数为$[1/12+\sin^2\Delta/6]/\sqrt3=5/(24\sqrt3)$。三对相加，一次项相消，势能
\[
 U=\sum_{i<j}\frac{KQ^2}{|\vect r_i-\vect r_j|}
\]
展开到二阶为
\[
 U=\frac{KQ^2}{R}\left[\sqrt3+\frac12\vect a^{\mathsf T}\mathbf H_a\vect a
 +\frac12\vect z^{\mathsf T}\mathbf H_z\vect z\right]+O(3),
\]
\[
 \mathbf H_a=\frac5{12\sqrt3}\begin{pmatrix}2&-1&-1\\-1&2&-1\\-1&-1&2\end{pmatrix},\quad
 \mathbf H_z=\frac1{3\sqrt3}\begin{pmatrix}1&1&1\\1&1&1\\1&1&1\end{pmatrix}.
\]
这些矩阵已提出$KQ^2/R$，其元素和本征值均无量纲。$\mathbf H_a(1,1,1)^{\mathsf T}=0$对应绕$z$轴整体转动；在$\sum a_i=0$的子空间，$\mathbf H_a\vect a=5\vect a/(4\sqrt3)$，可选$(1,-1,0)$、$(1,1,-2)$两个独立振型。$\mathbf H_z(1,1,1)^{\mathsf T}=(1,1,1)^{\mathsf T}/\sqrt3$，对应三球同向离开大圆；$\sum z_i=0$的两个零模则是绕平面内轴整体转动。因此非零本征值为
\[
 \lambda_1=\frac{\sqrt3}{3}\quad(1\text{重}),\qquad
 \lambda_2=\frac{5\sqrt3}{12}\quad(2\text{重}).
\]
而二阶动能为 $\tfrac12mR^2\sum_i(\dot a_i^2+\dot z_i^2)$。故三个振动模的频率为
\[
 \boxed{\omega_1^2=\frac{KQ^2}{mR^3}\frac{\sqrt3}{3}}\quad(1\text{重}),
 \qquad
 \boxed{\omega_2^2=\frac{KQ^2}{mR^3}\frac{5\sqrt3}{12}}\quad(2\text{重}).
\]
共3个零模与3个振动模，和自由度计数一致，非零二阶系数为正，故除整体取向外稳定。量纲检查：$KQ^2/R^3$是力除以长度，再除质量为$\mathrm{s^{-2}}$；电荷趋零时回复作用消失，频率也趋零。
\end{xesolution}

\Needspace{16\baselineskip}
\begin{origexamplebox}{弹簧摆势能的三阶展开与内共振}\label{ch10:ex:spring-pendulum}
原长为 $l_0$、劲度系数为 $k$ 的轻弹簧悬挂质量为 $m$ 的小球。以平衡位置为原点，$x$ 为水平位移，
$y$ 为竖直向下位移。把势能写到三阶小量
\[
 V=\frac12m(\omega_1^2x^2+\omega_2^2y^2)
 +ax^3+bx^2y+cxy^2+dy^3,
\]
求全部系数，并讨论最低阶非线性耦合产生的内共振条件。

\ThQSpringPendulum
\end{origexamplebox}
\Needspace{6\baselineskip}\begin{xesolution}
平衡长度
\[
 L=l_0+\frac{mg}{k}.
\]
图\ref{ch10:fig:spring-pendulum}中球到悬点距离为$r=\sqrt{x^2+(L+y)^2}$。势能（略去常量）为
\[
 V=\frac12k(r-l_0)^2-mgy.
\]
\chTenSpringPendulumFigure
当$|x|,|y|\ll L$，先提出正的$L+y$，逐阶展开：
\[
 r=(L+y)\sqrt{1+\frac{x^2}{(L+y)^2}}
 =L+y+\frac{x^2}{2(L+y)}+O(4).
\]
再对$1/(L+y)$用几何级数，得
\[
 r=L+y+\frac{x^2}{2L}-\frac{x^2y}{2L^2}+O(4)
\]
$O(4)$表示$x,y$总次数至少4，相应$L$的幂补足量纲。令$e=L-l_0=mg/k$，平方并保留至三阶：
\[
 (r-l_0)^2=e^2+2ey+y^2+\frac eLx^2
 +\left(\frac1L-\frac e{L^2}\right)x^2y+O(4).
\]
乘$k/2$、减$mgy$后，线性项因$ke=mg$消失，利用$L-e=l_0$得到
\[
 V=\frac12\frac{mg}{L}x^2+\frac12ky^2
 +\frac{kl_0}{2L^2}x^2y+O(4).
\]
因此
\[
 \boxed{\omega_1^2=\frac gL},\qquad
 \boxed{\omega_2^2=\frac km},
\]
\[
 \boxed{a=c=d=0},\qquad
 \boxed{b=\frac{kl_0}{2L^2}}.
\]
$x\to-x$的对称性排除了$x^3,xy^2$；$x=0$时竖直弹簧恰好线性，故无$y^3$。$b$单位为能量/长度三次方。由$F_x=-\partial V/\partial x,F_y=-\partial V/\partial y$得
\[
 m\ddot x+\frac{mg}{L}x+2bxy=0,\qquad
 m\ddot y+ky+bx^2=0.
\]
若先取$x\simeq A\cos\omega_1t$，则$x^2=A^2(1+\cos2\omega_1t)/2$，所以$x^2y$项使竖直模受到平均偏移作用及频率约$2\omega_1$的驱动，同时$2byx$使水平模受到竖直振动的参数调制；
最低阶强内共振条件为
\[
 \boxed{\omega_2\simeq2\omega_1},\qquad
 \frac{k}{m}\simeq\frac{4g}{L}.
\]
结合$L=l_0+mg/k$，精确线性频率比为2时有$kl_0=3mg$。远离该条件时，三阶项主要造成小的频率修正和高次谐波；接近共振时，两模可显著交换能量，但没有外界无限供能。将两方程乘对应速度相加，可验在所保留阶数内$\dd(T+V)/\dd t=0$。特别地$l_0=0$时$b=0$，这一阶的耦合消失，不能只看频率关系便断言相同机制仍存在。
\end{xesolution}

\endgroup

\hypersetup{bookmarksdepth=2}

% ThermalQualityRevision20261023-v1
\hypersetup{bookmarksdepth=3}
\begingroup
\renewcommand{\Needspace}[1]{\par\begingroup
\dimen0=#1\relax\vskip0pt plus\dimen0\penalty-100
\vskip0pt plus-\dimen0\vskip\dimen0\penalty9999\vskip-\dimen0
\vskip0pt\endgroup}
\tcbset{breakable/.default=false}
\chapter{有心力、万有引力与天体运动}
% Editable TikZ: use the force centre, not the centre of an ellipse, for distances.
\providecommand{\ThermalFigure}[3]{\begin{qualityfigure}
\begin{tikzpicture}[>=Latex,font=\small,x=.8cm,y=.8cm]#1\end{tikzpicture}
\figcap{#2}\label{#3}\end{qualityfigure}}
\newcommand{\ThRadialFallBase}[1]{
 \draw[->](-.5,0)--(7,0)node[right]{$r$};\fill(0,0)circle(4pt)node[below]{$M$};
 \fill[DeepBlue](2.5,0)circle(2pt)node[above]{$r_1$};\fill[DeepBlue](6,0)circle(2pt)node[above]{$r_0$};
 \draw[->,thick](5.5,.7)--(3.9,.7)node[midway,above]{下落};#1}
\newcommand{\ThQRadialFall}{\ThermalFigure{\ThRadialFallBase{}}{从外侧静止释放，径向坐标随时间减小，始末距离都为正。}{ch11:quality:radial-q}}
\newcommand{\ThARadialFall}{\ThermalFigure{\ThRadialFallBase{
 \draw[<->](3,-.8)--node[below]{$|\dd r|$}(3.7,-.8);\node[below]at(4.5,-1.6){$\dot r<0,\quad\dd t=-\dd r/|\dot r|$};
}}{径向速度取负分支，时间元为正；积分上下限须与此相配。}{ch11:quality:radial-a}
图~\ref{ch11:quality:radial-a}中位移箭头朝向力心，故能量方程求出的速率须配 $\dot r=-\sqrt{2GM(1/r-1/r_0)}$。时间积分若交换为从 $r_1$ 到 $r_0$，负号才随之消失；这一符号核对避免得到负的下落时间。
}
\newcommand{\ThOrbitBase}[1]{
 \fill[DeepBlue!12](0,0)circle(1.4);\draw[thick](0,0)circle(1.4);\node at(0,0){地球};
 \draw[dashed](0,0)circle(1.65);\fill[DeepBlue](1.65,0)circle(2pt)node[right]{$P$};
 \draw[->](0,0)--(140:1.4)node[midway,above]{$R$};#1}
\newcommand{\ThQOrbit}{\ThermalFigure{\ThOrbitBase{
 \node[right]at(3,1){近地轨道：高度相对 $R$ 可忽略};
}}{圆轨道速度和逃逸速度对应不同的运动要求；图中间隙为示意。}{ch11:quality:orbit-q}}
\newcommand{\ThAOrbit}{\ThermalFigure{\ThOrbitBase{
 \draw[->,thick](1.65,0)--(1.65,1.5)node[right]{$v_c$};
 \draw[->,DeepBlue,thick](1.65,0)--(.5,0)node[above]{$F_g$};
 \node[right,align=left]at(3,0){圆轨道：引力提供向心力\\逃逸临界：无穷远速度趋零};
}}{用 $GM=gR^2$ 把引力参数换为题给地表量；两速度不能混用。}{ch11:quality:orbit-a}
图~\ref{ch11:quality:orbit-a}中的圆周条件是 $mv_c^2/R=GMm/R^2$。逃逸临界态则以无穷远为势能零点并令总机械能为零；二者分别由力与能量决定，故逃逸速度比圆轨道速度大 $\sqrt2$ 倍。
}
\newcommand{\ThApsidesBase}[1]{
 \draw[thick](0,0)ellipse(3 and 1.8);\fill(2.4,0)circle(3pt)node[above left]{$S$};
 \fill(3,0)circle(2pt)node[right]{$P$};\fill(-3,0)circle(2pt)node[left]{$A$};
 \draw[<->](2.4,-2.2)--node[below]{$r_p$}(3,-2.2);\draw[<->](-3,-3)--node[below]{$r_a$}(2.4,-3);#1}
\newcommand{\ThQApsides}{\ThermalFigure{\ThApsidesBase{}}{近点 $P$、远点 $A$ 的距离从焦点力心 $S$ 量起；椭圆中心不作为距离原点。}{ch11:quality:apsides-q}}
\newcommand{\ThTransferBase}[1]{
 \draw[TextGray](0,0)circle(1.2);\draw[TextGray](0,0)circle(3.2);
 \draw[DeepBlue,thick](-1,0)ellipse(2.2 and 1.96);\fill(0,0)circle(3pt)node[below right]{$S$};
 \fill(1.2,0)circle(2pt)node[below right]{$P$};\fill(-3.2,0)circle(2pt)node[left]{$A$};
 \draw[->](0,0)--(50:1.2)node[right]{$r_1$};\draw[->](0,0)--(120:3.2)node[above]{$r_2$};#1}
\newcommand{\ThQTransfer}{\ThermalFigure{\ThTransferBase{}}{转移椭圆与内、外圆轨道分别相切于 $P,A$；力心为它们共同的焦点 $S$。}{ch11:quality:transfer-q}}
\newcommand{\ThATransfer}{\ThermalFigure{\ThTransferBase{
 \draw[->,thick](1.2,0)--(1.2,1.3)node[right]{$\Delta v_1$};
 \draw[->,thick](-3.2,0)--(-3.2,-1)node[left]{$\Delta v_2$};
 \draw[->,DeepBlue,thick](.35,1.55)--(-.55,1.9);
}}{两次脉冲均沿当地顺行切向；无推力段从近点飞到远点，历时半个椭圆周期。}{ch11:quality:transfer-a}
图~\ref{ch11:quality:transfer-a}中的第一脉冲把内圆速度提高到转移椭圆近点速度；第二脉冲把较低的远点速度提高到外圆速度。两处箭头在纸面方向相反，却都沿顺行方向，因此两次速度增量的大小均为正。
}
\newcommand{\ThMassLossBase}[1]{
 \draw[dashed](-1.2,0)circle(1.2);\draw[dashed](-1.2,0)circle(3);
 \fill[DeepBlue](-2.4,0)circle(5pt)node[below]{$M$};\fill(1.8,0)circle(3pt)node[below]{$m$};
 \fill(-1.2,0)circle(1.5pt)node[below]{$C$};\draw[<->](-2.4,-1)--node[below]{$d$}(1.8,-1);
 \draw[->,thick](-2.4,0)--(-2.4,-1.5);\draw[->,thick](1.8,0)--(1.8,1.5);#1}
\newcommand{\ThQMassLoss}{\ThermalFigure{\ThMassLossBase{
 \foreach\a in {30,90,150,210,270,330}{\draw[->,TextGray](-2.4,0)--++(\a:.8);}
 \node[right]at(3,0){$M$ 瞬时各向同性抛出 $\Delta M$};
}}{爆发前绕共同质心 $C$ 转动；无反冲使残余两星瞬时速度连续。}{ch11:quality:mass-loss-q}}
\newcommand{\ThAMassLoss}{\ThermalFigure{
 \fill[DeepBlue](0,0)circle(4pt)node[below]{$M-\Delta M$};\fill(5,0)circle(3pt)node[below]{$m$};
 \draw[<->](0,-.8)--node[below]{$d$（瞬时不变）}(5,-.8);
 \draw[->,thick](5,0)--(5,1.8)node[right]{$v_{\rm rel}$（瞬时不变）};
 \node[above]at(2.5,2.4){相对运动的引力参数改变为 $G(M+m-\Delta M)$};
}{爆发后转入残余双星的相对坐标，束缚条件由相对比能量为负判断。}{ch11:quality:mass-loss-a}
图~\ref{ch11:quality:mass-loss-a}固定的是爆发瞬间的相对距离与相对速度；爆发前的质心位置和分量动能不应继续当作残余系统的质心量。用 $v_{\rm rel}^2=G(M+m)/d$ 与新的引力参数比较，才能直接判断新轨道是否束缚。
}
\newcommand{\ThLRLBase}[1]{
 \fill(0,0)circle(3pt)node[below]{$O$};\draw[->](0,0)--(5,0)node[right]{极轴};
 \fill[DeepBlue](2.4,1.7)circle(2pt)node[right]{$P$};\draw[->](0,0)--(2.4,1.7)node[midway,above left]{$\vect r$};
 \draw(.8,0)arc(0:35.3:.8)node[right]{$\theta$};#1}
\newcommand{\ThQLRL}{\ThermalFigure{\ThLRLBase{
 \draw[->,thick](2.4,1.7)--(1.3,2.5)node[above]{$\vect p$};
}}{平面有心力运动中的位置、动量与极角；待证明守恒矢量可用于选取极轴。}{ch11:quality:lrl-q}}
\newcommand{\ThALRL}{\ThermalFigure{\ThLRLBase{
 \draw[->,DeepBlue,thick](0,0)--(2,0)node[below]{$\vect e$};
 \draw[->,thick](0,0)--(4.8,3.4)node[right]{$\uvec r$ 的方向};
}}{将固定的偏心率矢量作为极轴，点乘径向单位矢量引入 $e\cos\theta$。}{ch11:quality:lrl-a}
图~\ref{ch11:quality:lrl-a}中的 $\theta$ 从常矢量 $\vect e$ 量起，故 $\vect e\cdot\uvec r=e\cos\theta$。另一方面，$\vect L$ 垂直轨道平面；利用三重积把 $(\vect p\times\vect L)\cdot\uvec r$ 化为 $L^2/r$，便把矢量守恒式转成轨道方程。圆轨道 $e=0$ 时极轴可任取。
}
\newcommand{\ThQDust}{\ThermalFigure{
 \draw[thick](0,0)circle(.7);\node at(0,0){$M_s$};\draw[->](0,0)--(-45:.7)node[below]{$r_s$};
 \draw[->,DeepBlue,thick](-4,1.8)--(-1.8,1.8)node[midway,above]{$v$};\draw[dashed](-4,0)--(0,0);
 \draw[<->](-3.7,0)--node[left]{$b$}(-3.7,1.8);\node[right]at(2,1.6){远处平行来流，密度 $\rho$};
}{飞船静止系中的入射尘埃；冲量参数 $b$ 从远处渐近线量起，不等于最近距离。}{ch11:quality:dust-q}}
\newcommand{\ThQPrecession}{\ThermalFigure{
 \draw[thick](0,0)circle(2);\fill(0,0)circle(3pt)node[below]{$M$};\fill[DeepBlue](2,0)circle(2pt)node[above]{$m$};
 \draw[->](2,0)--(.8,0)node[above]{$F_g$};\draw[->,DeepBlue](2,0)--(3,0);
 \node[DeepBlue]at(3.6,.7){$F_1=kmR$};\draw[thin,TextGray](3,0)--(3.4,.45);
 \draw[<->](0,-.7)--node[below]{$R$}(2,-.7);\node[align=left,anchor=west]at(5,1.3){近圆轨道上\\存在微小径向扰动};
}{径向附加力朝外；圆周角运动与径向小振动须分开求频率。}{ch11:quality:precession-q}}
\newcommand{\ThAPrecession}{\ThermalFigure{
 \draw[TextGray,domain=0:360,samples=100]plot({2*cos(\x)/(1+.18*cos(\x))},{2*sin(\x)/(1+.18*cos(\x))});
 \draw[DeepBlue,thick,domain=0:360,samples=100]plot({2*cos(\x+25)/(1+.18*cos(\x))},{2*sin(\x+25)/(1+.18*cos(\x))});
 \fill(0,0)circle(2pt)node[below]{$O$};\draw[dashed](0,0)--(1.695,0);\draw[dashed](0,0)--(25:1.695);
 \draw[->](1,0)arc(0:25:1);\node[right]at(2.8,.5){$\Delta\varpi=2\pi(\Omega/\kappa_r-1)$};
}{两次近日点方向的示意叠图，不把缓慢进动轨道误认为一条闭合椭圆。}{ch11:quality:precession-a}
图~\ref{ch11:quality:precession-a}中一次近日点到下一次近日点用的是径向周期 $2\pi/\kappa_r$，这期间方位角增加 $2\pi\Omega/\kappa_r$。超出 $2\pi$ 的部分就是顺行进动角；因此进动角速度取两频率之差，而不是把附加力简单代入开普勒周期。
}
\newcommand{\ThQOrbitDrag}{\ThermalFigure{
 \draw[thick](0,0)ellipse(3 and 1.8);\fill(2.4,0)circle(3pt)node[below]{$M$};\fill[DeepBlue](0,1.8)circle(2pt)node[above]{$m$};
 \draw[->,thick](0,1.8)--(-1.4,1.8)node[above]{$\vect v$};\draw[->,DeepBlue,thick](0,1.8)--(1.4,1.8)node[above]{$-\gamma\vect v$};
}{弱阻力始终与瞬时速度反向；研究对象是很多周期后的缓慢平均演化。}{ch11:quality:drag-q}}
\newcommand{\ThAOrbitDrag}{\ThermalFigure{
 \draw[thick](-1.8,0)ellipse(3 and 2.4);\draw[DeepBlue,thick](-1.08,0)ellipse(1.8 and 1.44);
 \fill(0,0)circle(3pt)node[below]{$S$};\draw[->](-4,1.6)--(-2.7,1.05);
 \node[right,align=left]at(2,1){同一焦点 $S$\\周期平均下 $a$ 减小、$e$ 保持不变};
}{以力心为中心按比例缩放的相继密切椭圆；实际轨道不在一个周期内突然跳变。}{ch11:quality:drag-a}
图~\ref{ch11:quality:drag-a}中的两条椭圆共用焦点 $S$，两个半轴按同一因子缩小。证明这一点需要同时追踪 $E$ 与 $L$：由它们的平均变化检验 $e^2=1+2EL^2/(G^2M^2m^3)$ 不变，不能只由能量降低就断言偏心率不变。
}



第9章已建立角动量守恒、有效势与有心力轨道的基本方法。本章先作必要回顾，再重点讨论轨道初始条件、二体运动及天体运动中的具体应用。

\section{有心力运动的守恒量}
\label{ch11:sec:central}

把力心固定在惯性系原点 $O$。若质点所受合力始终沿 $O$ 与质点的连线，就称它为有心力。
先研究大小只依赖距离、且不显含时间的情形：
\begin{equation}\label{ch11:eq:central-force}
 \vect F=F(r)\uvec r,\qquad r=|\vect r|>0.
\end{equation}
这里 $F(r)$ 是有符号的径向分量：$F>0$ 向外，$F<0$ 向内。
“沿径向”保证力矩为零；“只依赖 $r$、不显含时间”才使这里的机械能也守恒。
例如随时间改变的径向力仍可使角动量守恒，却未必有守恒的机械能。
真实的两颗星都会运动，暂时固定力心是质量悬殊时的近似；后面将用相对坐标解除这个限制。

对原点的角动量为 $\vect L=m\vect r\times\vect v$。由于叉乘两个平行矢量为零，
\begin{equation}
 \frac{\dd\vect L}{\dd t}=\vect r\times\vect F=0.
\end{equation}
若 $\vect L\ne0$，任意时刻都有 $\vect r\cdot\vect L=0$，所以轨道在一个通过力心、垂直于
$\vect L$ 的固定平面内。若 $L=0$，速度与位置矢量平行，运动退化为沿直线的径向运动；
这一情形不能套用后面含 $1/L$ 或 $1/h$ 的轨道公式。

在非退化轨道平面内取极坐标 $(r,\theta)$，选择 $\theta$ 的正向使角动量大小 $L>0$，记
\begin{equation}\label{ch11:eq:h}
 h=\frac Lm=r^2\dot\theta>0.
\end{equation}
$h$ 是单位质量角动量，量纲为 $\mathrm{m^2/s}$。极坐标的两个单位矢量也随时间转动，故
\[
 \dot{\uvec r}=\dot\theta\uvec\theta,\qquad
 \dot{\uvec\theta}=-\dot\theta\uvec r,
\]
\[
 \vect v=\dot r\uvec r+r\dot\theta\uvec\theta,\qquad
 \vect a=(\ddot r-r\dot\theta^2)\uvec r
 +(r\ddot\theta+2\dot r\dot\theta)\uvec\theta.
\]
牛顿第二定律的径向、切向分量分别是
\begin{equation}\label{ch11:eq:polar-newton}
 m(\ddot r-r\dot\theta^2)=F(r),\qquad
 r\ddot\theta+2\dot r\dot\theta=0.
\end{equation}
将后一式乘以 $r$，左端正是 $\dd(r^2\dot\theta)/\dd t$，再次得到式\eqref{ch11:eq:h}。
在短时间 $\dd t$ 内位置矢量扫过面积 $\dd A=\tfrac12r^2\dd\theta$，因此
\begin{equation}\label{ch11:eq:area-law}
 \frac{\dd A}{\dd t}=\frac12r^2\dot\theta=\frac h2.
\end{equation}
这就是开普勒（Kepler）第二定律的动力学来源；等面积定律适用于有心力，并非平方反比引力独有。

再取势能 $U(r)$ 满足 $F(r)=-\dd U/\dd r$。两个速度分量正交，机械能为
\begin{equation}\label{ch11:eq:energy-original}
 E=\frac12m(\dot r^2+r^2\dot\theta^2)+U(r).
\end{equation}
利用式\eqref{ch11:eq:polar-newton}可得
$\dot E=m\vect v\cdot\vect a+U'(r)\dot r=F(r)\dot r-F(r)\dot r=0$。
对于这种没有速度相关势的模型，哈密顿量（Hamiltonian）就是上述机械能。
给定初始位置与速度，$E$ 和 $\vect L$ 就由初态确定，以后不再改变。

\subsection{束缚运动与散射}
同一组动力学方程可以回答两类不同的问题。束缚运动关心粒子是否始终停留在有限距离范围内、
是否稳定以及多久返回；散射则关心从远方入射的粒子经过力心附近后朝哪个方向离开。
图\ref{ch11:fig:bound-scatter}把这两种初始条件并列画出。

\begin{center}\begin{minipage}{\linewidth}\centering
\begin{tikzpicture}[>=Latex,scale=.84,font=\small]
 \begin{scope}[shift={(-3.8,0)}]
  \draw[curve] (-1,0) ellipse (2 and 1.732);
  \draw[curve,->,domain=45:67,samples=30] plot ({-1+2*cos(\x)},{1.732*sin(\x)});
  \fill (0,0) circle (1.8pt) node[below left=3pt] {$O$};
  \fill (1,0) circle (1.8pt) node[below right=3pt] {$P_0$};
  \draw[guide] (0,0)--(1,0) node[midway,below=5pt] {$r_0$};
  \draw[vector] (1,0)--(1,1.05) node[right=4pt] {$v_0$};
  \node at (-.7,-2.85) {(a) 有限区域内的束缚运动};
 \end{scope}
 \begin{scope}[shift={(3.3,0)},scale=.78]
  \pgfmathsetmacro{\aa}{.35}\pgfmathsetmacro{\bb}{1}
  \pgfmathsetmacro{\cc}{sqrt(\aa*\aa+\bb*\bb)}
  \pgfmathsetmacro{\cs}{\aa/\cc}\pgfmathsetmacro{\sn}{\bb/\cc}
  \draw[guide] (-3.8,1)--(3.6,1);
  \draw[curve,domain=2.0:-1.7,samples=180,variable=\u]
   plot ({\cs*(\cc-\aa*cosh(\u))-\sn*\bb*sinh(\u)},
         {\sn*(\cc-\aa*cosh(\u))+\cs*\bb*sinh(\u)});
  \fill (0,0) circle (2pt) node[below left=3pt] {$O$};
  \draw[guide] (-3.8,0)--(-2.6,0);
  \draw[<->] (-3.25,0)--(-3.25,1) node[midway,left=4pt] {$b$};
  \draw[vector] (-3.4,1)--(-2.1,1) node[midway,above=5pt] {$\vect v_\infty$};
  
  \begin{scope}[shift={(.3,-1.6)}]
   \draw[->] (0,0)--(1.35,0) node[right] {$\vect v_\infty$};
   \draw[->,DeepBlue] (0,0)--(-38.58:1.35) node[below right] {$\vect v'_\infty$};
   \draw (.58,0) arc(0:-38.58:.58);
   \node at (-19.29:.83) {$\chi$};
  \end{scope}
  \node at (-.2,-3.65) {(b) 远处速度方向的偏折};
 \end{scope}
\end{tikzpicture}
\figcap{两类有心力问题。左图的 $r_0,v_0$ 决定初始能量与角动量；右图的 $b$ 是力心到入射渐近线的垂距。}
\label{ch11:fig:bound-scatter}
\end{minipage}\end{center}

在图\ref{ch11:fig:bound-scatter}(a)中，初速度垂直于半径，初始径向速度为零；
它可能是近点或远点，却未必是圆轨道。只有径向向心条件也满足才是圆轨道。
纯粹斥力不能在无其他约束的空间维持这样的束缚轨道；有引力也不自动保证束缚或稳定。
在图(b)中，$b$ 称为冲量参数，也称瞄准距离；远处速度为 $v_\infty$ 时 $h=bv_\infty$。
散射角是远处入射、出射\emph{速度}的夹角，不是两位置矢量的夹角。
引力和斥力均可造成散射。

描述远处自由入射时，通常要求力及选定的势能随 $r\to\infty$ 衰减，并检查远场渐近运动能否定义。
“短程力”与“长程力”还涉及衰减速率及具体散射理论；平方反比引力和库仑力通常称为长程力，
但仍能研究其经典散射，不能把“必须短程”当作这里的普遍条件。
量子散射、光线在引力场中的偏折也延续了这一问题思路，不过其运动方程不能不加条件地照搬经典质点模型。

\section{有效势能与径向稳定性}
\label{ch11:sec:effective}
把 $\dot\theta=h/r^2$ 代入式\eqref{ch11:eq:energy-original}，得
\begin{equation}\label{ch11:eq:effective}
 E=\frac12m\dot r^2+U_{\rm eff}(r),\qquad
 U_{\rm eff}(r)=U(r)+\frac{mh^2}{2r^2}=U(r)+\frac{L^2}{2mr^2}.
\end{equation}
这样，二维轨道的径向变化等价于一维粒子在有效势中的运动。后一项来自横向动能，称为离心势垒；
它不是新添的实际相互作用，改变轨道角动量也会改变这条有效势曲线。
由径向方程同样得到 $m\ddot r=-U_{\rm eff}'(r)$。
允许区域满足 $E\ge U_{\rm eff}$；若等号成立且 $U_{\rm eff}'\ne0$，径向速度为零并转向。
圆轨道则必须同时满足
\begin{equation}\label{ch11:eq:circle-condition}
 U_{\rm eff}'(r_c)=0,\qquad E=U_{\rm eff}(r_c),\qquad \dot r=0.
\end{equation}

图\ref{ch11:fig:effective}取引力势 $U=-GMm/r$。此时
$r_c=h^2/(GM)$、$E_c=U_{\rm eff}(r_c)=-GMm/(2r_c)$。
令 $x=r/r_c$，便有 $U_{\rm eff}/|E_c|=x^{-2}-2x^{-1}$，纵横轴均为无量纲量。

\begin{center}\begin{minipage}{\linewidth}\centering
\begin{tikzpicture}
\begin{axis}[width=.80\linewidth,height=6.5cm,axis lines=left,
 xlabel={$r/r_c$},ylabel={$U_{\rm eff}/|E_c|$},xmin=.35,xmax=4.35,ymin=-1.45,ymax=2.0,
 xtick=\empty,ytick={-1,0,1},clip=false,
 every axis label/.style={fill=none},xlabel style={at={(axis description cs:1,-.04)},anchor=north},
 ylabel style={at={(axis description cs:-.04,1)},anchor=south}]
 \addplot[curve,domain=.42:4.25,samples=180]{1/x^2-2/x};
 \addplot[guide] coordinates {(.35,0)(4.25,0)};
 \addplot[guide] coordinates {(.35,-.5)(4.25,-.5)};
 \draw[guide] (axis cs:.585786,-1.45)--(axis cs:.585786,-.5);
 \draw[guide] (axis cs:3.414214,-1.45)--(axis cs:3.414214,-.5);
 \fill[DeepBlue] (axis cs:1,-1) circle (2pt);
 \draw[thin] (axis cs:1.1,-1.06)--(axis cs:1.65,-1.23);
 \node[anchor=west] at (axis cs:1.7,-1.25) {稳定圆轨道};
 \node[anchor=south] at (axis cs:2.8,-.47) {$E/|E_c|=-1/2$};
 \node[anchor=north] at (axis cs:.585786,-1.45) {$r_-/r_c$};
 \node[anchor=north] at (axis cs:3.414214,-1.45) {$r_+/r_c$};
\end{axis}
\end{tikzpicture}
\figcap{同一角动量下的引力有效势。水平能量线与曲线的两个交点确定近点 $r_-$ 与远点 $r_+$，最低点对应圆轨道。}
\label{ch11:fig:effective}
\end{minipage}\end{center}

图中的束缚粒子在 $r_-$ 与 $r_+$ 之间径向往返，但同时 $\dot\theta=h/r^2$ 始终为正；
所以粒子不是沿一条半径来回走。有效势先回答允许区和稳定性，精确的轨道形状仍需求 $r(\theta)$。
在圆轨道附近令 $r=r_c+\xi$，小扰动满足
\[
 m\ddot\xi=-U_{\rm eff}''(r_c)\xi+\order{\xi^2}.
\]
$U_{\rm eff}''>0$ 时是径向小振动，$U_{\rm eff}''<0$ 时扰动指数增长；
若二阶导数为零，不能立即宣布稳定，还需检查更高阶或完整方程。

\subsection{幂律力的三类有效势}
\label{ch11:sec:power}
取 $F=-kr^n$、$k>0$，这里 $n$ 是\emph{力}的指数，而不是势的指数。积分得
\[
 U(r)=\frac{k}{n+1}r^{n+1}+C\quad(n\ne-1),\qquad
 U(r)=k\ln(r/r_*)+C\quad(n=-1).
\]
参考长度 $r_*$ 只改变势能零点。$n\ne-3$ 时，由圆轨道条件有
\[
 kr_c^n=\frac{mh^2}{r_c^3},\qquad
 r_c=\left(\frac{mh^2}{k}\right)^{1/(n+3)}.
\]
对有效势再求一次导数，并用圆轨道条件消去 $k$ 或 $h$：
\begin{equation}\label{ch11:eq:power-stability}
 U_{\rm eff}''(r_c)=knr_c^{n-1}+\frac{3mh^2}{r_c^4}
 =\frac{mh^2}{r_c^4}(n+3).
\end{equation}
所以稳定圆轨道要求 $n>-3$。这不是“引力越强越稳定”：若引力在靠近力心时比 $r^{-3}$ 增长得更快，
小的向内偏离会使引力增长超过维持原角动量所需的离心项，从而继续向内偏离。

\begin{center}\begin{minipage}{\linewidth}\centering
\begin{tikzpicture}[font=\small]
 \foreach \xx/\nn/\expr in {0/-4/{.5/x^2-1/(3*x^3)},4.65/-2/{.5/x^2-1/x},9.3/1/{.5/x^2+.5*x^2}}{
  \begin{axis}[at={(\xx cm,0)},anchor=south west,width=4.6cm,height=4.7cm,
   axis lines=middle,xmin=.3,xmax=3.4,ymin=-1.1,ymax=2.2,ticks=none,
   xlabel={$r$},ylabel={$U_{\rm eff}$},clip=true,
   xlabel style={at={(axis description cs:1.02,.333333)},anchor=west},
   ylabel style={at={(axis description cs:0,1.05)},anchor=south}]
   \addplot[curve,domain=.42:3.2,samples=180]{\expr};
   \ifnum\nn=-4
    \fill (axis cs:1,.166667) circle (1.7pt);
    \node[anchor=south west] at (axis cs:1.1,.28) {$P$：不稳定};
   \else\ifnum\nn=-2
    \addplot[guide] coordinates {(.612574113,-.3)(2.72075922,-.3)};
    \fill (axis cs:1,-.5) circle (1.7pt);
    \node[anchor=north west] at (axis cs:1.1,-.60) {$P$：稳定};
   \else
    \addplot[guide] coordinates {(.618,1.5)(1.618,1.5)};
    \fill (axis cs:1,1) circle (1.7pt);
    \node[anchor=north west] at (axis cs:1.1,.87) {$P$：稳定};
   \fi\fi
  \end{axis}
 }
 \node at (2.2,-.55) {(a) $n<-3$，取 $n=-4$};
 \node at (6.85,-.55) {(b) $-3<n<-1$，取 $n=-2$};
 \node at (11.5,-.55) {(c) $n>-1$，取 $n=1$};
\end{tikzpicture}
\figcap{三类幂律引力的有效势。各面板使用各自的无量纲单位；能量线用于显示势阱内的允许区，不表示三种系统的能量相同。}
\label{ch11:fig:power-potentials}
\end{minipage}\end{center}

图\ref{ch11:fig:power-potentials}(a)中圆轨道处是势垒顶而非势阱底；
向两侧扰动可导致向力心坠落或向远处逃离。图(b)的有效势在无穷远趋于零，
负能量的适当范围可束缚，非负能量则可通向无穷远。图(c)的势能在远处无限增大，
有限能量运动的外侧也受到限制。$n=-1$ 的对数势同样有内侧离心势垒和远处上升的势能，
应直接用对数式讨论，不能代入含 $1/(n+1)$ 的公式。

当 $n=-3$ 时，
\[
 U_{\rm eff}=\frac{mh^2-k}{2r^2}+C.
\]
若 $mh^2=k$，有效势为常数，各个半径都可在 $\dot r=0$ 时维持圆周运动，
但没有径向回复力；加入微小径向速度便会漂移，不是稳定势阱中的振动。
若 $mh^2\ne k$，有效势单调，不存在圆轨道。这个临界情形应与一般幂律分开。

\section{从径向运动求时间与轨道}
\label{ch11:sec:quadrature}
将有效势能量式解出径向速度，必须保留方向分支：
\[
 \dot r=\pm\sqrt{\frac2m[E-U_{\rm eff}(r)]}.
\]
出向段取正号、入向段取负号，转向点处连接两支。
从 $r_a$ 到 $r_b$ 的一段运动，其时间可按沿途方向取
\[
 \Delta t=\int_{r_a}^{r_b}\frac{\dd r}{\dot r}.
\]
对于图\ref{ch11:fig:effective}中两个简单转向点之间的束缚运动，径向周期为
\begin{equation}\label{ch11:eq:radial-period}
 T_r=2\int_{r_-}^{r_+}\frac{\dd r}{\sqrt{2[E-U_{\rm eff}(r)]/m}}.
\end{equation}
因从近点到远点只占一个径向周期的一半，前面有因子2。
端点虽使分母为零，但简单转向点附近积分形如 $\int\dd r/\sqrt{|r-r_\pm|}$，通常收敛。
恰好圆轨道的两个端点重合，应改用小振动极限，而不是直接把积分上下限设成同一点。

用 $\dot\theta=h/r^2$ 除以 $\dot r$，得到
\begin{equation}\label{ch11:eq:angle-integral}
 \theta_b-\theta_a=\pm\int_{r_a}^{r_b}
 \frac{h\,\dd r}{r^2\sqrt{2[E-U_{\rm eff}(r)]/m}}.
\end{equation}
这里正负号也由出向/入向段决定。一般位置之间的积分端点是这两个位置，
\emph{不是必须取转向点}。只有求近点到远点夹角、完整径向周期或散射全过程时，才按问题选择转向点及无穷远。
例如每个径向周期内转过的方位角为
\[
 \Delta\theta_r=2\int_{r_-}^{r_+}
 \frac{h\,\dd r}{r^2\sqrt{2[E-U_{\rm eff}(r)]/m}}.
\]
若 $\Delta\theta_r/(2\pi)$ 是有理数，径向状态和方位角可以在整数个周期后同时重复；
若为无理数，轨道一般不闭合。径向周期和绕力心转一圈所需时间不能先验地当成同一个量。

\subsection{散射角与冲量参数}
远方取 $U(\infty)=0$，则 $E=mv_\infty^2/2$、$h=bv_\infty$。
设散射轨道有一个最近距离 $r_m$，由 $E=U_{\rm eff}(r_m)$ 求出。
从最近点到无穷远的方位角变化为
\[
 \Phi=\int_{r_m}^{\infty}
 \frac{b\,\dd r}{r^2\sqrt{1-b^2/r^2-2U(r)/(mv_\infty^2)}}.
\]
不受力的直线轨道有 $\Phi=\pi/2$。对于不绕转多圈的通常单次散射，
远处速度的偏折角为 $\chi=|2\Phi-\pi|$；引力向力心弯折、斥力背离力心弯折。
若轨道绕力心多圈或有多个转向结构，需保留累积方位角，并按实际入出射方向取夹角，不能丢掉绕转次数。

\subsection{比内方程的换元}
轨道形状用 $r(\theta)$ 描述比先求 $r(t),\theta(t)$ 再消去时间更直接。
令 $u(\theta)=1/r$，以下撇号表示对 $\theta$ 求导。首先
\[
 \dot\theta=hu^2,\qquad
 \frac{\dd r}{\dd\theta}=-\frac{u'}{u^2},\qquad
 \frac{\dd^2r}{\dd\theta^2}=-\frac{u''}{u^2}+\frac{2(u')^2}{u^3}.
\]
由链式法则分别得到
\[
 \dot r=-\frac{u'}{u^2}\,hu^2=-hu',\qquad
 \ddot r=-hu''\dot\theta=-h^2u^2u'',\qquad
 r\dot\theta^2=h^2u^3.
\]
代入径向牛顿方程，整理为比内方程（Binet equation）：
\begin{equation}\label{ch11:eq:binet}
 u''+u=-\frac{F(1/u)}{mh^2u^2}
       =-\frac{mF(1/u)}{L^2u^2}.
\end{equation}
左边和右边的量纲均为长度的倒数。这一换元把时间消去，但不改变轨道的物理条件；
径向退化运动 $h=0$ 仍应直接用时间方程处理。

\subsection{积分上下限和运动分支如何选择}
上面的积分式解决三个不同问题：时间积分用来问“多久到达”，角度积分用来问“到达何处”，比内方程用来直接求轨道形状。列式之前，先由 $E\ge U_{\rm eff}(r)$ 找允许区间，再用初始径向速度选入向或出向分支。不能因根号约定为正，就把正在靠近力心的 $\dot r$ 也取正。

\Needspace{16\baselineskip}\begin{origexamplebox}{从静止径向落向引力中心的时间}\label{ch11:ex:teaching-radial-fall}
质量可忽略的质点从距固定点质量 $M$ 为 $r_0$ 处静止释放，仅受万有引力。求降到 $r_1$ 的时间，其中 $0<r_1<r_0$，并讨论形式上的 $r_1\to0$ 极限。把引力源看作点质量。

\ThQRadialFall
\end{origexamplebox}
\Needspace{6\baselineskip}\begin{xesolution}
\ThARadialFall

初速度为零，故角动量为零，运动是纯径向的，不可使用含 $1/L$ 的轨道公式。记 $\mu=GM$，单位质量机械能守恒给出
\[
 \tfrac12\dot r^2-\frac\mu r=-\frac\mu{r_0}.
\]
质点向内落，选负根 $\dot r=-\sqrt{2\mu(1/r-1/r_0)}$。为使时间积分为正，写成
\[
 t=\int_{r_1}^{r_0}\frac{\dd r}{\sqrt{2\mu(1/r-1/r_0)}}.
\]
令 $r=r_0\cos^2u$，$u_1=\arccos\sqrt{r_1/r_0}$。则 $|\dd r|=2r_0\sin u\cos u\,\dd u$，分母为 $\sqrt{2\mu/r_0}\tan u$，故
\[
 t=\sqrt{\frac{2r_0^3}{\mu}}\int_0^{u_1}\cos^2u\,\dd u
 =\sqrt{\frac{r_0^3}{2\mu}}\left(u_1+\sin u_1\cos u_1\right).
\]
回到半径变量，括号中的第二项为 $\sqrt{(r_1/r_0)(1-r_1/r_0)}$。当 $r_1\to r_0$，时间趋零；当 $r_1\to0$，形式极限为 $\pi\sqrt{r_0^3/(8GM)}$，量纲为时间。
现实中有半径的天体会先与质点接触；点质量模型在中心还有奇点，因此有限的形式时间不意味着质点能无条件穿过中心继续沿同一公式运动。
\end{xesolution}

\section{平方反比力与万有引力二体问题}
\label{ch11:sec:kepler}
经典平方反比力问题是天体力学的重要起点。统一写成 $F=-m\kappa/r^2$，
则引力的 $\kappa=GM>0$；库仑相互作用可取
$\kappa=-Qq/(4\pi\varepsilon_0m)$，同号电荷给 $\kappa<0$。
有些文献用 $\sqrt{GM}$ 作为引力强度参数；本章始终直接使用 $GM$，避免平方与符号混淆。
对 $\kappa>0$，式\eqref{ch11:eq:binet}成为常系数方程
\[
 u''+u=\frac{\kappa}{h^2}.
\]
特解为常数，齐次解为正弦与余弦的组合。将两积分常数合并为振幅和初相位：
\begin{equation}\label{ch11:eq:conic}
 u=\frac{\kappa}{h^2}[1+e\cos(\theta-\theta_0)],\qquad
 r=\frac{p}{1+e\cos(\theta-\theta_0)},\quad p=\frac{h^2}{\kappa}.
\end{equation}
$e\ge0$ 是偏心率，$p$ 为半通径，$\theta_0$ 取近点方向；$e=0$ 时是圆，近点方向不唯一。
物理轨道只取分母为正的区域。$0<e<1$ 为椭圆，$e=1$ 为抛物线，$e>1$ 为双曲线。
这称为开普勒问题。对斥力 $\kappa<0$，更方便把解写成
$u=-|\kappa|/h^2+C\cos(\theta-\theta_0)$，取 $u>0$ 的双曲线支；
不能仍把 $p=h^2/\kappa$ 当成正长度并沿用引力的全部范围。

\subsection{有效势积分给出能量与偏心率}
比内方程容易给出圆锥曲线，却还没有把 $e$ 与初始能量联系起来。
令比机械能 $\varepsilon=E/m$，将 $\dot r=-hu'$ 代入引力能量式：
\[
 \varepsilon=\tfrac12h^2(u')^2+\tfrac12h^2u^2-\kappa u.
\]
配方而不是省略交叉项，得
\[
 (u')^2+\left(u-\frac{\kappa}{h^2}\right)^2
 =\frac{\kappa^2}{h^4}+\frac{2\varepsilon}{h^2}.
\]
记右边平方根为 $C$，则
\[
 \frac{\dd(u-\kappa/h^2)}{\sqrt{C^2-(u-\kappa/h^2)^2}}=\pm\dd\theta.
\]
积分并把常相位吸收进 $\theta_0$，恢复式\eqref{ch11:eq:conic}，同时得到
\begin{equation}\label{ch11:eq:energy-eccentricity}
 e^2=1+\frac{2\varepsilon h^2}{\kappa^2}
     =1+\frac{2EL^2}{m^3\kappa^2}.
\end{equation}
因此负、零、正能量分别对应椭圆、抛物线、双曲线；这依赖引力势零点取在无穷远。
对于椭圆，近远点之和为 $2a$，
\[
 r_p=\frac p{1+e}=a(1-e),\qquad
 r_a=\frac p{1-e}=a(1+e),\qquad p=a(1-e^2).
\]
代回能量关系得
\begin{equation}\label{ch11:eq:vis-viva}
 E=-\frac{m\kappa}{2a},\qquad
 v^2=\kappa\left(\frac2r-\frac1a\right).
\end{equation}
后一式称为活力公式。对于双曲线若用正的几何半实轴 $a_h=p/(e^2-1)$，
则 $E=+m\kappa/(2a_h)$；抛物线用 $E=0$ 单独处理。

于是给定初始 $r_0,v_{r0},v_{\theta0}$ 后，可以依次计算
\[
 h=r_0v_{\theta0},\qquad
 \varepsilon=\tfrac12(v_{r0}^2+v_{\theta0}^2)-\frac\kappa{r_0},\qquad
 p=\frac{h^2}{\kappa},\quad e=\sqrt{1+\frac{2\varepsilon h^2}{\kappa^2}}.
\]
当 $e>0$ 时，由
$e\cos(\theta-\theta_0)=p/r_0-1$ 和
$e\sin(\theta-\theta_0)=hv_{r0}/\kappa$ 共同确定初相位象限，
不能只用反余弦丢掉入向与出向之别。

\subsection{两颗星都运动时怎样使用上述公式}
设两颗星质量为 $m_1,m_2$，没有外力，定义
\[
 M=m_1+m_2,\quad
 \vect R=\frac{m_1\vect r_1+m_2\vect r_2}{M},\quad
 \vect r=\vect r_1-\vect r_2,\quad \mu=\frac{m_1m_2}{M}.
\]
$\mu$ 称为约化质量。相加两个牛顿方程得 $M\ddot{\vect R}=0$；相减则得
\[
 \ddot{\vect r}=-GM\frac{\vect r}{r^3},\qquad
 \mu\ddot{\vect r}=-Gm_1m_2\frac{\vect r}{r^3}.
\]
因此前述固定中心方程仍适用，但运动的是\emph{两星的相对位置}，引力参数为 $GM$。
由 $\vect r_1=\vect R+(m_2/M)\vect r$、$\vect r_2=\vect R-(m_1/M)\vect r$，展开并相消交叉项：
\[
 K=\tfrac12M|\dot{\vect R}|^2+\tfrac12\mu|\dot{\vect r}|^2,\qquad
 E_{\rm rel}=\tfrac12\mu|\dot{\vect r}|^2-\frac{Gm_1m_2}{r}.
\]
相对角动量为 $\vect L_{\rm rel}=\mu\vect r\times\dot{\vect r}$，故 $h=L_{\rm rel}/\mu$。
在质心系，两颗星的轨道与相对轨道相似、比例分别为 $m_2/M$ 与 $m_1/M$，且位于质心两侧。
只有 $m_2\ll m_1$ 时才能近似把 $GM$ 换为较大星体的 $Gm_1$。

\subsection{圆轨道、逃逸速度与周期}
圆轨道由 $mv_c^2/r=GMm/r^2$ 得
\[
 v_c=\sqrt{GM/r},\qquad T_c=\frac{2\pi r}{v_c}=2\pi\sqrt{r^3/(GM)}.
\]
若从距离 $r$ 处恰好逃到无穷远且末速为零，则初态总能量为零：
$mv_e^2/2-GMm/r=0$，所以 $v_e=\sqrt{2GM/r}=\sqrt2v_c$。
对椭圆不能直接用周长除以某处速度，因为速度不恒定。设短半轴为 $b_e=a\sqrt{1-e^2}$，
用式\eqref{ch11:eq:area-law}的面积速度除去椭圆面积：
\[
 T=\frac{\pi ab_e}{h/2}=\frac{2\pi ab_e}{\sqrt{GMp}}
 =2\pi\sqrt{\frac{a^3}{GM}}.
\]
故开普勒第三定律为
\begin{equation}\label{ch11:eq:third-law}
 \boxed{T^2=\frac{4\pi^2a^3}{G(m_1+m_2)}}.
\end{equation}
其中 $a$ 是相对椭圆的半长轴；若给某一颗星相对质心的半长轴，必须先按质量比换算。

\Needspace{16\baselineskip}\begin{origexamplebox}{近地圆轨道和逃逸速度}\label{ch11:ex:01}
忽略大气和地球自转，地球半径为 $R$、地表重力加速度为 $g$。求近地圆轨道速度、周期和第二宇宙速度。

\ThQOrbit
\end{origexamplebox}
\Needspace{6\baselineskip}\begin{xesolution}
\ThAOrbit

地表重力满足 $mg=GMm/R^2$，消去试探物体质量得 $GM=gR^2$。
近地圆轨道径向方程为 $mv_c^2/R=mg$，因此
\[
 \boxed{v_c=\sqrt{gR}},\qquad
 \boxed{T=\frac{2\pi R}{v_c}=2\pi\sqrt{\frac Rg}}.
\]
恰好逃逸是能量阈值而非“重力突然消失”：途中仍受引力，动能逐渐减小。
令初态机械能等于无穷远静止时的零能量，
\[
 \frac12mv_e^2-\frac{GMm}{R}=0
 \quad\Longrightarrow\quad\boxed{v_e=\sqrt{2gR}}.
\]
$gR$ 的量纲为速度平方、$R/g$ 为时间平方。取 $R=6.37\times10^6\,\mathrm m$、
$g=9.81\,\mathrm{m/s^2}$，分别约为 $7.9\,\mathrm{km/s}$、$84.4\,\mathrm{min}$ 和 $11.2\,\mathrm{km/s}$。
这是理想近地极限；实际卫星不能贴地飞行，需避开地形与大气阻力。
\end{xesolution}

\Needspace{16\baselineskip}\begin{origexamplebox}{椭圆轨道远近点速度}\label{ch11:ex:02}
卫星绕质量为 $M$ 的天体作椭圆运动，近点和远点距离分别为 $r_p,r_a$。求两点速度之比，并给出速度本身。

\ThQApsides
\end{origexamplebox}
\Needspace{6\baselineskip}\begin{xesolution}
图\ref{ch11:fig:apsis}中力心在椭圆的焦点而非中心。两拱点有 $\dot r=0$，速度均垂直于位置矢量。
\begin{center}\begin{minipage}{\linewidth}\centering
\begin{tikzpicture}[>=Latex,scale=.78,font=\small]
 \draw[curve] (0,0) ellipse (3 and 1.8);
 \coordinate (S) at (2.4,0);
 \fill (S) circle (2pt) node[above left=4pt] {$S$};
 \fill (3,0) circle (2pt) node[below right=4pt] {$P$};
 \fill (-3,0) circle (2pt) node[below left=4pt] {$A$};
 \draw[vector] (3,0)--(3,1.3) node[right=4pt] {$v_p$};
 \draw[vector] (-3,0)--(-3,-.7) node[left=4pt] {$v_a$};
 \draw[guide] (-3,0)--(3,0);
 \draw[<->] (2.4,-2.2)--(3,-2.2) node[midway,below=5pt] {$r_p$};
 \draw[<->] (-3,-2.2)--(2.4,-2.2) node[midway,below=5pt] {$r_a$};
\end{tikzpicture}
\figcap{近点 $P$ 与远点 $A$ 的速度沿同一绕行方向。两距离之和等于椭圆长轴 $2a$。}\label{ch11:fig:apsis}
\end{minipage}\end{center}
角动量守恒给出 $r_pv_p=r_av_a$，故 $v_p/v_a=r_a/r_p$。
这还没有给出速度的绝对大小。利用两点机械能相等，
\[
 \tfrac12v_p^2-\frac{GM}{r_p}
 =\tfrac12v_a^2-\frac{GM}{r_a}.
\]
代入 $v_a=(r_p/r_a)v_p$，在非圆情形消去 $r_a-r_p$，
\[
 v_p^2\frac{(r_a-r_p)(r_a+r_p)}{2r_a^2}
 =GM\frac{r_a-r_p}{r_pr_a}.
\]
于是
\[
 \boxed{v_p=\sqrt{\frac{2GMr_a}{r_p(r_p+r_a)}}},\qquad
 \boxed{v_a=\sqrt{\frac{2GMr_p}{r_a(r_p+r_a)}}},\qquad
 \boxed{\frac{v_p}{v_a}=\frac{r_a}{r_p}}.
\]
也可把图中 $a=(r_p+r_a)/2$ 代入活力公式\eqref{ch11:eq:vis-viva}复核。
圆轨道 $r_p=r_a=r$ 虽不能在上一步除以零，但由连续极限或向心方程都得 $v=\sqrt{GM/r}$。
\end{xesolution}

\section{轨道能量、变速与拱线进动}
椭圆总机械能 $E=-GMm/(2a)$ 只依赖半长轴，角动量则还区分形状。
在某点瞬时变速时，位置和势能未变，能量改变量来自动能，角动量改变量来自速度的切向分量。
所以只给“增加多少能量”通常不足以确定一条新轨道，还要给变速方向。

从半径 $r_0$ 的\emph{圆轨道}出发，若只作沿原切向的瞬时变速且不反向，则 $\dot r=0$ 仍成立。
新速度 $v$ 满足
\[
 \ddot r=\frac{v^2}{r_0}-\frac{GM}{r_0^2}.
\]
若 $v_c<v<v_e$，该点是新椭圆近点；若 $0<v<v_c$，该点是远点；$v=v_c$ 保持圆轨道。
$v=v_e$ 是抛物轨道、$v>v_e$ 是双曲轨道，都不再是束缚椭圆。
若 $v=0$，则退化为径向坠落。若从一般椭圆上 $\dot r\ne0$ 的位置变速，
仅改变切向速度不会消去径向速度，该位置不能因此成为近点或远点。

\subsection{一次变速先确定新轨道，再讨论怎样接续}
瞬时变速只改变速度，不改变该瞬间的位置。此后若发动机关闭，卫星沿由新能量、新角动量确定的轨道运动；不能把“速度变大”直接翻译成“立即进入更大的圆”。较高圆轨道的圆周速度虽然更低，从较低圆轨道转移过去却通常要先顺行加速，因为首先要提高轨道能量。

\Needspace{16\baselineskip}\begin{origexamplebox}{两个共面圆轨道之间的双脉冲转移}\label{ch11:ex:teaching-transfer}
卫星原在半径 $r_1$ 的圆轨道上，拟转入同平面、同绕行方向、半径 $r_2>r_1$ 的圆轨道。采用一条近点为 $r_1$、远点为 $r_2$ 的椭圆转移轨道，在两端各施加一次沿切向的瞬时速度增量。忽略卫星质量及推进过程的有限持续时间，求两次增量与转移时间。

\ThQTransfer
\end{origexamplebox}
\Needspace{6\baselineskip}\begin{xesolution}
\ThATransfer

图~\ref{ch11:fig:apsis} 的 $P,A$ 分别对应转移椭圆的近点、远点，故半长轴 $a=(r_1+r_2)/2$。记 $\mu=GM$。转移过程中单位质量能量为 $-\mu/(2a)$，在任一点有
\[
 \frac{v^2}{2}-\frac\mu r=-\frac\mu{2a},
 \qquad v^2=\mu\left(\frac2r-\frac1a\right).
\]
这先由能量守恒给出端点速度，尚未把它们误认为圆周速度。在 $P$ 点：
\[
 v_P=\sqrt{\frac{2\mu r_2}{r_1(r_1+r_2)}},\qquad
 \Delta v_1=v_P-\sqrt{\frac\mu{r_1}}>0.
\]
第一次顺行加速，把圆轨道切换为转移椭圆。飞到 $A$ 点时：
\[
 v_A=\sqrt{\frac{2\mu r_1}{r_2(r_1+r_2)}},\qquad
 \Delta v_2=\sqrt{\frac\mu{r_2}}-v_A>0.
\]
第二次仍是顺行加速，因为转移椭圆远点速度比当地圆周速度更小。两端满足 $r_1v_P=r_2v_A$，核验了无推力飞行段的角动量守恒；两次脉冲处卫星自身角动量则改变，不能跨过脉冲继续当作守恒。
从近点到远点用半个椭圆周期，因此
\[
 t_{\rm tr}=\pi\sqrt{\frac{a^3}{\mu}}
 =\pi\sqrt{\frac{(r_1+r_2)^3}{8\mu}}.
\]
$r_2\to r_1$ 时两次速度增量均趋零；形式上的飞行时间趋于原圆轨道半周期，表示仍按“飞到对面”的约定计时，而不是说留在原轨道必须等待半周期。结果的速度与时间量纲分别由 $\sqrt{\mu/r}$、$\sqrt{r^3/\mu}$ 给出。
\end{xesolution}

\begingroup\tcbset{breakable/.default=true}\begin{starread}{近圆轨道的径向频率与拱线}
\label{ch11:sec:precession}
设圆轨道平均角频率为 $\Omega=h/r_c^2$，径向小振动频率
$\kappa_r=\sqrt{U_{\rm eff}''(r_c)/m}$。每个径向周期转过
$\Delta\theta_r\simeq2\pi\Omega/\kappa_r$。
对于力指数为 $n$ 的幂律引力，式\eqref{ch11:eq:power-stability}给
$\kappa_r^2=(n+3)\Omega^2$。
平方反比引力 $n=-2$ 有 $\kappa_r=\Omega$，相邻近日点方向相同；
各向同性谐振子 $n=1$ 则有 $\kappa_r=2\Omega$，相邻最近点相差 $\pi$，
是同一条固定椭圆短轴上的另一端，并不是长轴持续进动。
因此“两个频率不相等”不能直接推出拱线转动，必须检查方位角及方向的等价关系。

在接近开普勒势的弱扰动中，$\kappa_r$ 接近 $\Omega$，每径向周期超过 $2\pi$ 的小角
$\delta\varpi=2\pi(\Omega/\kappa_r-1)$ 对应平均进动速率
$\dot\varpi=\delta\varpi/(2\pi/\kappa_r)=\Omega-\kappa_r$。
这个公式的使用背景是近圆、近开普勒轨道，不能不加解释地套到所有幂律势。
在通常光滑有心势的条件下，使所有束缚轨道都闭合的平方反比引力和各向同性谐振子是两个重要特例。
\end{starread}\endgroup

\Needspace{22\baselineskip}
\section{天体运动综合问题}

\Needspace{16\baselineskip}\begin{origexamplebox}{哈雷彗星的轨道状态}
\label{ch11:ex:03}
哈雷彗星绕太阳沿椭圆轨道逆时针运动，周期为 $T=76.1\,\mathrm y$。1986年经过近日点
$P_0$ 时，日彗距离为 $r_0=0.590\,\mathrm{AU}$；此后到达 $P$ 点，位置矢量
$\overrightarrow{SP}$ 与近日点方向 $\overrightarrow{SP_0}$ 的夹角为 $\theta=72^\circ$。
取 $1\,\mathrm{AU}=1.5\times10^{11}\,\mathrm m$。求此时的日彗距离、速度大小和速度方向
（以速度矢量与射线 $SP_0$ 的逆时针夹角表示）。


\begin{center}\begin{minipage}{\linewidth}\centering
\begin{tikzpicture}[scale=.90,>=Latex]
  \coordinate (S) at (0,0);
  \draw[axis,->] (0,0)--(4.3,0) node[right] {近日点方向};
  \draw[curve,samples=180,smooth,variable=\a,domain=-104:104]
    plot ({2.0*1.16062*cos(\a)/(1+.96715*cos(\a))},
          {2.0*1.16062*sin(\a)/(1+.96715*cos(\a))});
  \fill (S) circle (2.6pt) node[below left] {太阳 $S$};
  \coordinate (P0) at (1.18,0);
  \fill[DeepBlue] (P0) circle (2.2pt) node[below right] {$P_0$};
  \coordinate (P) at ({1.78712*cos(72)},{1.78712*sin(72)});
  \draw[guide] (S)--(P);
  \fill[DeepBlue] (P) circle (2.4pt) node[above right] {$P$};
  \draw[->,TextGray] (.70,0) arc[start angle=0,end angle=72,radius=.70];
  \node at (36:.95) {$\theta$};
\end{tikzpicture}
\figcap{题设轨道的近日点附近示意。$S$ 为太阳，$P_0$ 为近日点，$\theta=72^\circ$。}
\label{ch11:fig:halley}
\end{minipage}\end{center}
\end{origexamplebox}
\Needspace{6\baselineskip}\begin{xesolution}
图\ref{ch11:fig:halley-analysis}中 $SP_0$ 是位置角的零方向，速度沿轨道切线。
\begin{center}\begin{minipage}{\linewidth}\centering
\begin{tikzpicture}[scale=.90,>=Latex]
  \coordinate (S) at (0,0);
  \draw[axis,->] (0,0)--(4.3,0) node[right] {近日点方向};
  \draw[curve,samples=180,smooth,variable=\a,domain=-104:104]
    plot ({2.0*1.16062*cos(\a)/(1+.96715*cos(\a))},
          {2.0*1.16062*sin(\a)/(1+.96715*cos(\a))});
  \fill (S) circle (2.6pt) node[below left] {太阳 $S$};
  \coordinate (P0) at (1.18,0);
  \fill[DeepBlue] (P0) circle (2.2pt) node[below right] {$P_0$};
  \coordinate (P) at ({1.78712*cos(72)},{1.78712*sin(72)});
  \draw[guide] (S)--(P);
  \fill[DeepBlue] (P) circle (2.4pt) node[above right] {$P$};
  \draw[->,DeepBlue,very thick] (P)--++(126.70:1.25) node[above left] {$\vect v$};
  \draw[->,TextGray] (.70,0) arc[start angle=0,end angle=72,radius=.70];
  \node at (36:.95) {$\theta$};
\end{tikzpicture}
\figcap{近日点方向取作极轴。图中 $\theta$ 是位置角，而速度方向还要由径向、切向两个分量共同确定。}
\label{ch11:fig:halley-analysis}
\end{minipage}\end{center}
以地球轨道为单位使用开普勒第三定律，太阳质量占主导时有
\[
 \left(\frac{T}{1\,\mathrm y}\right)^2
 =\left(\frac{a}{1\,\mathrm{AU}}\right)^3,
 \qquad
 a=(76.1)^{2/3}\,\mathrm{AU}=17.9579\,\mathrm{AU}.
\]
近日点满足 $r_0=a(1-e)$，所以
\[
 e=1-\frac{r_0}{a}=0.967145,
 \qquad
 p=a(1-e^2)=1.16062\,\mathrm{AU}.
\]
图中 $SP_0$ 已选作真近点角的零方向，故在 $\theta=72^\circ$ 处
\[
 \boxed{r=\frac{p}{1+e\cos\theta}=0.893562\,\mathrm{AU}}
 \approx1.34\times10^{11}\,\mathrm m.
\]

在 $\mathrm{AU}$、年为单位时，太阳引力参数为
$\mu_\odot=GM_\odot=4\pi^2\,\mathrm{AU^3/y^2}$。由活力公式
\[
 v=\sqrt{\mu_\odot\left(\frac2r-\frac1a\right)}
   =9.28243\,\mathrm{AU/y},
\]
取一年为 $365.25$ 天，把 $\mathrm{AU/y}$ 换算成 $\mathrm{km/s}$，即
\[
 \boxed{v=44.0\,\mathrm{km/s}}.
\]
只给速度大小还不能确定图中的箭头。令单位质量角动量
$h=r^2\dot\theta=\sqrt{\mu_\odot p}$。对轨道方程求导，
\[
 \frac{\dd r}{\dd\theta}=\frac{pe\sin\theta}{(1+e\cos\theta)^2},\qquad
 \dot r=\frac{\dd r}{\dd\theta}\frac h{r^2}=\frac h p e\sin\theta,\qquad
 r\dot\theta=\frac h r.
\]
利用 $h/p=\mu_\odot/h$，沿径向外和逆时针切向的分量分别为
\[
 v_r=\frac{\mu_\odot}{h}e\sin\theta=5.36456\,\mathrm{AU/y},
 \qquad
 v_\theta=\frac{\mu_\odot}{h}(1+e\cos\theta)=7.57529\,\mathrm{AU/y}.
\]
因 $v_r>0$，彗星正在远离太阳。速度相对径向 $SP$ 再逆时针偏转
$\arctan(v_\theta/v_r)=54.695^\circ$，所以相对 $SP_0$ 的方向角为
\[
 \boxed{\psi=72^\circ+54.695^\circ=126.695^\circ}.
\]
代回 $v^2=v_r^2+v_\theta^2$ 可恢复前述活力公式数值，方向与大小相互校验。
\end{xesolution}

\Needspace{16\baselineskip}\begin{origexamplebox}{双星系统中的瞬时对称失质量}
\label{ch11:ex:04}
质量分别为 $M$ 和 $m$（$M>m$）的两颗星相距 $d$，在相互引力作用下绕共同质心作圆周运动。
某次超新星爆发使质量较大的星体瞬时、各向同性地失去质量 $\Delta M$；抛出物不碰撞另一颗星，
且爆发不对残余星体施加反冲。求爆发后两颗残余星体仍保持束缚的条件。

\ThQMassLoss
\end{origexamplebox}
\Needspace{6\baselineskip}\begin{xesolution}
\ThAMassLoss

爆发前相对运动的圆轨道速度满足
\[
 v_{\rm rel}^2=\frac{G(M+m)}{d}.
\]
爆发是瞬时的，位置与两颗残余星体的速度均来不及改变；爆发后的引力质量和却变为
$M-\Delta M+m$。于是单位约化质量的相对轨道能量为
\begin{align*}
 \varepsilon'
 &=\frac12v_{\rm rel}^2-\frac{G(M-\Delta M+m)}{d}\\
 &=\frac{G}{2d}\,[2\Delta M-(M+m)].
\end{align*}
实际失质量还应满足 $0\le\Delta M<M$。残余双星保持束缚等价于 $\varepsilon'<0$，故
\[
 \boxed{\Delta M<\frac{M+m}{2}}.
\]
等号对应抛物线临界逃逸；超过该值为双曲线分离。这个判据依赖“瞬时、对称、无反冲”，若爆发给残余星体一个踢速度，
还必须把该速度矢量加入爆发瞬间的相对速度后重新计算能量。
\end{xesolution}

\Needspace{16\baselineskip}\begin{origexamplebox}{拉普拉斯--龙格--楞次矢量}
\label{ch11:ex:05}
质量为 $m$ 的质点在质量为 $M$ 的固定中心产生的平方反比引力中运动。证明
拉普拉斯--龙格--楞次矢量（Laplace--Runge--Lenz vector）
\[
 \vect e=\frac{\vect p\times\vect L}{GMm^2}-\frac{\vect r}{r}
\]
守恒，并由此推出轨道的圆锥曲线方程
\[
 r=\frac{p_0}{1+e\cos\theta}.
\]

\ThQLRL
\end{origexamplebox}
\Needspace{6\baselineskip}\begin{xesolution}
\ThALRL

这里 $\vect p=m\vect v$、$\vect L=\vect r\times\vect p$，并记
$\uvec r=\vect r/r$。引力为
\[
 \vect F=-\frac{GMm}{r^3}\vect r,
\]
所以 $\dot{\vect L}=\vect r\times\vect F=0$。另一方面，利用三重积恒等式，
\begin{align*}
 \frac{\dd}{\dd t}(\vect p\times\vect L)
 &=\vect F\times\vect L\\
 &=-\frac{GMm}{r^3}\,
   \vect r\times(\vect r\times\vect p)\\
 &=GMm\left(\frac{\vect p}{r}
 -\frac{\vect r(\vect r\cdot\vect p)}{r^3}\right)
 =GMm^2\frac{\dd\uvec r}{\dd t}.
\end{align*}
因此 $\dot{\vect e}=0$，所给矢量确为常矢量。

取极轴沿 $\vect e$，令 $\theta$ 为 $\vect e$ 与 $\vect r$ 的夹角。将守恒式点乘 $\uvec r$，并注意
\[
 (\vect p\times\vect L)\cdot\uvec r
 =\frac1r\vect L\cdot(\vect r\times\vect p)=\frac{L^2}{r},
\]
便有
\[
 e\cos\theta=\frac{L^2}{GMm^2r}-1.
\]
整理得
\[
 \boxed{r=\frac{p_0}{1+e\cos\theta}},
 \qquad
 \boxed{p_0=\frac{L^2}{GMm^2}}.
\]
这里取 $L\ne0$。$0<e<1,e=1,e>1$ 分别给出椭圆、抛物线、双曲线；
矢量的模就是轨道偏心率，方向指向近心点。$e=0$ 时轨道为圆，$\vect e=0$，
不能沿零矢量选极轴，任取平面内固定方向即可。$L=0$ 的纯径向运动直接用径向能量方程，
不将 $p_0=0$ 代入非退化轨道公式。$\vect e$ 无量纲，$p_0$ 具有长度量纲。
\end{xesolution}

\Needspace{16\baselineskip}\begin{origexamplebox}{星际尘埃造成的引力阻力}
\label{ch11:ex:06}
宇宙飞船质量为 $M_s$、几何半径为 $r_s$，以速度 $v$ 穿过密度为 $\rho$ 的均匀稀薄尘埃。
尘埃微粒质量远小于飞船质量，远处可视为平行来流；只有冲量参数不大于 $R_d$ 的尘埃计入环境范围。

（1）先把飞船视为固定的引力中心，求尘埃被散射到 $\theta$ 至 $\theta+\dd\theta$ 时的微分散射截面
$\dd\sigma/\dd\Omega$；（2）假设直接撞击的尘埃被飞船吸收，并且引力散射均为小角度，求直接碰撞阻力与
掠过飞船的尘埃造成的引力阻力。

\ThQDust
\end{origexamplebox}
\Needspace{6\baselineskip}\begin{xesolution}
设尘埃的冲量参数为 $b$，即远处入射直线与引力中心的距离，而不是最近距离。
远处单位质量能量为 $v^2/2$，单位质量角动量为 $bv$。代入偏心率公式得到
\[
 e=\sqrt{1+\frac{b^2v^4}{G^2M_s^2}}=\sqrt{1+\frac{b^2}{a_g^2}},
 \qquad a_g=\frac{GM_s}{v^2}.
\]
在无穷远处 $1+e\cos\nu_\infty=0$；入射与出射速度的夹角为
$\theta=2\nu_\infty-\pi$，故 $\sin(\theta/2)=1/e$，从而
\[
 \tan\frac\theta2=\frac{a_g}{b},
 \qquad
 a_g=\frac{GM_s}{v^2}.
\]
图\ref{ch11:fig:scattering}左侧实线是双曲线轨迹，右侧单独比较远处的速度矢量，
不把近处尚在变化的切线方向当作最终散射方向。
\begin{center}\begin{minipage}{\linewidth}\centering
\begin{tikzpicture}[scale=.87,>=Latex]
  \pgfmathsetmacro{\ag}{.35}\pgfmathsetmacro{\cc}{sqrt(1+\ag*\ag)}
  \pgfmathsetmacro{\cs}{\ag/\cc}\pgfmathsetmacro{\sn}{1/\cc}
  \pgfmathsetmacro{\chi}{2*atan(\ag)}
  \draw[guide] (-4.1,0)--(3.4,0);
  \draw[guide] (-4.1,1)--(-.2,1);
  \draw[<->] (-3.9,0)--(-3.9,1) node[midway,left] {$b$};
  \draw[curve,domain=2.05:-1.8,samples=180,variable=\u,->]
   plot ({\cs*(\cc-\ag*cosh(\u))-\sn*sinh(\u)},
         {\sn*(\cc-\ag*cosh(\u))+\cs*sinh(\u)});
  \filldraw[draw=DeepBlue,fill=DeepBlue!12] (0,0) circle(.14);
  \node[below] at (0,-.18) {$M_s$};
  \draw[->,DeepBlue] (-3.4,1.4)--(-2.2,1.4) node[midway,above] {远处入射};
  \node at (-.3,-1.65) {轨迹：$a_g/b=0.35$};
  \begin{scope}[shift={(5,1)}]
   \draw[vector] (0,0)--(2.3,0) node[right] {$\vect v$};
   \draw[vector] (0,0)--(-\chi:2.3) node[below right] {$\vect v'$};
   \draw[->] (.8,0) arc(0:-\chi:.8);
   \node at (-.5*\chi:1.14) {$\theta$};
   \draw[guide] (-\chi:2.3)--({2.3*cos(\chi)},0);
   \node[above] at (1.0,.20) {$v'=v$};
   \node at (1.15,-2.45) {纵向变化 $v-v\cos\theta$};
  \end{scope}
\end{tikzpicture}
\figcap{引力双曲线与远处速度夹角。右图把远处速度平移到同一起点比较，两者大小相等。}
\label{ch11:fig:scattering}
\end{minipage}\end{center}

由于 $b=a_g\cot(\theta/2)$、$\dd\Omega=2\pi\sin\theta\,\dd\theta$，轴对称散射给出
\begin{align*}
 \frac{\dd\sigma}{\dd\Omega}
 &=\frac{b}{\sin\theta}\left|\frac{\dd b}{\dd\theta}\right|\\
 &=\boxed{\left(\frac{GM_s}{2v^2}\right)^2
 \csc^4\frac\theta2}.
\end{align*}
这里使用了 $|\dd b/\dd\theta|=(a_g/2)\csc^2(\theta/2)$ 以及
$\sin\theta=2\sin(\theta/2)\cos(\theta/2)$。尘埃质量在速度与能量关系中约去，
因此截面与单个尘粒质量无关；结果量纲为面积。

几何半径与远处碰撞截面半径不同。恰好碰到表面的轨道在最近处速度为 $b_{\rm hit}v/r_s$，
由能量守恒
\[
 \frac{v^2}{2}=\frac{b_{\rm hit}^2v^2}{2r_s^2}-\frac{GM_s}{r_s}
 \quad\Longrightarrow\quad b_{\rm hit}^2=r_s^2+2a_gr_s.
\]
所以直接撞击的质量流率为 $\rho v\pi b_{\rm hit}^2$。在“吸收并停到飞船参考系”的假设下，每单位质量损失纵向动量 $v$，故
\[
 F_{\rm hit}=\pi b_{\rm hit}^2\rho v^2
 =\pi\rho v^2(r_s^2+2a_gr_s)\simeq\boxed{\pi r_s^2\rho v^2}.
\]
末一步与题设的小角近似一致，要求 $a_g\ll r_s$，在这一阶忽略引力聚焦对直接碰撞的修正。
冲量参数在 $b$ 到 $b+\dd b$ 的掠射尘埃，单位时间通过的质量为
$\rho v\,2\pi b\dd b$，纵向速度损失为 $v(1-\cos\theta)$。又
\[
 1-\cos\theta=\frac{2a_g^2}{b^2+a_g^2},
\]
所以排除直接碰撞区域后
\begin{align*}
 F_{\rm grav}
 &=\rho v^2\int_{b_{\rm hit}}^{R_d}2\pi b(1-\cos\theta)\,\dd b\\
 &=2\pi\rho v^2a_g^2
 \ln\!\frac{R_d^2+a_g^2}{b_{\rm hit}^2+a_g^2}.
\end{align*}
当所有散射角都很小，即 $a_g\ll r_s<R_d$ 时，化为
\[
 \boxed{F_{\rm grav}\simeq
 \frac{4\pi\rho G^2M_s^2}{v^2}\ln\frac{R_d}{r_s}}.
\]
总阻力为 $F_{\rm hit}+F_{\rm grav}$，方向与飞船相对尘埃的速度相反。若尘埃发生镜面反弹、烧蚀或强引力聚焦，
直接碰撞项的系数和积分下限都需相应修正；上述精确积分要求 $R_d>b_{\rm hit}$，小角结果还用到 $b_{\rm hit}\simeq r_s$。
碰撞区不能再计入掠射阻力，否则会重复计算。
\end{xesolution}

\Needspace{16\baselineskip}\begin{origexamplebox}{弱线性径向附加力造成的近日点进动}
\label{ch11:ex:07}
除太阳质量 $M$ 的平方反比引力外，质量为 $m$ 的行星还受到很小的径向外力
$\vect F_1=kmr\,\uvec r$，其中 $k>0$ 为常量。设近圆轨道半径为 $R$，求近日点进动角速度。

\ThQPrecession
\end{origexamplebox}
\Needspace{6\baselineskip}\begin{xesolution}
\ThAPrecession

附加力向外，故单位质量势能可写为
\[
 \Phi(r)=-\frac{GM}{r}-\frac12kr^2.
\]
在半径 $R$ 的圆轨道上，向心条件给出实际平均角频率
\[
 \boxed{\Omega^2=\frac{GM}{R^3}-k}.
\]
为求近日点进动，必须把径向小振动频率与 $\Omega$ 区分。令单位质量角动量为 $h=R^2\Omega$，
有效势为 $\Phi_{\rm eff}=\Phi+h^2/(2r^2)$。逐次求导得
\[
 \Phi_{\rm eff}'=\frac{GM}{r^2}-kr-\frac{h^2}{r^3},\qquad
 \Phi_{\rm eff}''=-\frac{2GM}{r^3}-k+\frac{3h^2}{r^4}.
\]
圆轨道条件使 $\Phi_{\rm eff}'(R)=0$ 且 $h^2=GMR-kR^4$。
写 $r=R+x$，将 $\ddot r=-\Phi_{\rm eff}'(r)$ 保留到 $x$ 的一阶，得到
\[
 \ddot x+\kappa_r^2x=0,
 \qquad
 \kappa_r^2=\Phi_{\rm eff}''(R)
 =\frac{GM}{R^3}-4k.
\]
因此一次径向周期内方位角略多于 $2\pi$，近日点沿运动方向进动，其角速度为
\[
 \boxed{\dot\varpi=\Omega-\kappa_r
 =\sqrt{\frac{GM}{R^3}-k}
 -\sqrt{\frac{GM}{R^3}-4k}}.
\]
记无扰动平均运动 $n=\sqrt{GM/R^3}$。在 $k\ll n^2$ 时，
$\Omega=n-k/(2n)+O(k^2/n^3)$，
$\kappa_r=n-2k/n+O(k^2/n^3)$，相减得
\[
 \boxed{\dot\varpi\simeq\frac{3k}{2n}
 =\frac{3k}{2}\sqrt{\frac{R^3}{GM}}}.
\]
稳定近圆运动还要求 $GM/R^3>4k$。当 $k\to0$ 时 $\Omega=\kappa_r$、进动消失，极限正确。
\end{xesolution}

\Needspace{16\baselineskip}\begin{origexamplebox}{弱速度阻力下椭圆轨道的长期演化}
\label{ch11:ex:08}
质量为 $m$ 的质点在固定质量 $M$ 的平方反比引力下沿椭圆运动，同时受到很小的阻力
$\vect f=-\gamma\vect v$。取时间间隔 $\dd t$ 远长于轨道周期、又短到轨道根数的相对变化仍很小。
已知机械能和角动量的大小分别为 $E<0,L$。求周期平均后的 $\dd E,\dd L$，并证明轨道只按比例缩小、偏心率不变。

\ThQOrbitDrag
\end{origexamplebox}
\Needspace{6\baselineskip}\begin{xesolution}
\ThAOrbitDrag

阻力的瞬时功率为
\[
 \dot E=\vect f\cdot\vect v=-\gamma v^2.
\]
一周内先忽略微小阻力对轨道形状的改变。沿这一闭合开普勒轨道，
\[
 \frac{\dd}{\dd t}(\vect r\cdot\vect v)
 =v^2+\vect r\cdot\dot{\vect v}=v^2-\frac{GM}{r}.
\]
一周后 $\vect r\cdot\vect v$ 回到原值，故左端周期平均为零，
$\langle v^2\rangle=GM\langle1/r\rangle$。代入
$E=m\langle v^2\rangle/2-GMm\langle1/r\rangle$，
得 $E=-m\langle v^2\rangle/2$。于是长期平均为
\[
 \boxed{\dd E=\frac{2\gamma}{m}E\,\dd t}.
\]
因为 $E<0$，右端为负，轨道能量变得更负。阻力矩则直接给出
\[
 \dot{\vect L}=\vect r\times(-\gamma\vect v)
 =-\frac\gamma m\vect L,
 \qquad
 \boxed{\dd L=-\frac\gamma mL\,\dd t}.
\]

开普勒椭圆的偏心率满足
\[
 e^2=1+\frac{2EL^2}{G^2M^2m^3}.
\]
在当前精度下
\[
 \frac{\dd(EL^2)}{EL^2}
 =\frac{\dd E}{E}+2\frac{\dd L}{L}
 =\frac{2\gamma}{m}\dd t-\frac{2\gamma}{m}\dd t=0,
\]
所以
\[
 \boxed{\dd e=0}.
\]
另一方面 $E=-GMm/(2a)$，从而
\[
 \frac{\dd a}{a}=-\frac{\dd E}{E}
 =-\frac{2\gamma}{m}\dd t.
\]
半长轴按比例减小而偏心率不变，因此近、远点距离 $a(1\mp e)$ 以相同比例收缩。这个结论是周期平均结果，
单个轨道周期内的瞬时轨道根数仍会有小幅振荡。

把半长轴方程积分，得到
\[
 \boxed{a(t)=a_0\exp\!\left(-\frac{2\gamma t}{m}\right)}.
\]
由开普勒第三定律 $T\propto a^{3/2}$，轨道周期随时间按
\[
 \boxed{T(t)=T_0\exp\!\left(-\frac{3\gamma t}{m}\right)}
\]
缩短。这个指数规律并不意味着质点在有限时间撞到中心；在该平均模型中 $a$ 只在
$t\to\infty$ 时趋于零。

还可直接核对能量损失率：阻力功率为
$\dot E=\vect F_d\cdot\vect v=-\gamma v^2$。椭圆轨道一周平均满足
$\langle v^2\rangle=GM/a$，于是
\[
 \langle\dot E\rangle=-\frac{\gamma GM}{a},
\]
与 $E=-GMm/(2a)$ 及 $\dot a/a=-2\gamma/m$ 完全一致。
上述推导要求一周内轨道根数变化很小，即 $\gamma T/m\ll1$；阻尼过强时不能再把运动看成一串缓慢收缩的开普勒椭圆。
\end{xesolution}

\endgroup

\hypersetup{bookmarksdepth=2}

% ThermalQualityRevision20261023-v1
% QualityRevision20261023
\hypersetup{bookmarksdepth=3}
\begingroup
\renewcommand{\Needspace}[1]{\par\begingroup
\dimen0=#1\relax\vskip0pt plus\dimen0\penalty-100
\vskip0pt plus-\dimen0\vskip\dimen0\penalty9999\vskip-\dimen0
\vskip0pt\endgroup}
\tcbset{breakable/.default=false}
\chapter{弹性体与机械波}
% Editable TikZ: boundary conditions, propagation and energy transport.
\providecommand{\ThermalFigure}[3]{\begin{qualityfigure}
\begin{tikzpicture}[>=Latex,font=\small,x=.8cm,y=.8cm]#1\end{tikzpicture}
\figcap{#2}\label{#3}\end{qualityfigure}}
\newcommand{\ThQSelfWeight}{\ThermalFigure{
 \draw[very thick](-1,3)--(1,3);\draw[fill=DeepBlue!12,thick](-.25,3)rectangle(.25,0);
 \draw[<->](1,0)--node[right]{$L$}(1,3);\draw[->](2.5,2.7)--(2.5,1.5)node[right]{$g$};
 \node[right,align=left]at(3.5,1.7){均匀杆：$A,\rho,Y$\\上端固定，下端自由};
}{自重加载的竖直细杆；下端自由意味着该处轴向张力为零。}{ch12:quality:self-weight-q}}
\newcommand{\ThStringBase}[1]{
 \draw[very thick](0,-.7)--(0,.7) (6,-.7)--(6,.7);\draw[thick](0,0)--(6,0);
 \draw[<->](0,-1)--node[below]{$L$}(6,-1);\node[above]at(3,.8){均匀弦，线密度 $\mu$};#1}
\newcommand{\ThQStringFrequency}{\ThermalFigure{\ThStringBase{\draw[->](0,0)--(-1,0)node[left]{$\mathcal T$};\draw[->](6,0)--(7,0)node[right]{$\mathcal T$};}}{两端固定、长度与线密度保持不变，改变张力而不改变边界条件。}{ch12:quality:string-frequency-q}}
\newcommand{\ThAStringFrequency}{\ThermalFigure{\ThStringBase{
 \draw[DeepBlue,thick,domain=0:6,samples=70]plot(\x,{.65*sin(30*\x)});
 \node[above]at(3,1.5){基模：两端为节点，中间为一个波腹};
}}{基模中弦长恒为半个波长，升频来自波速变化。}{ch12:quality:string-frequency-a}
图~\ref{ch12:quality:string-frequency-a}的两端节点固定了 $\lambda_1=2L$。因此 $f_1=c/(2L)$，而 $c=\sqrt{\mathcal T/\mu}$；频率翻倍要求波速翻倍、张力变为四倍。不能误把增加张力等同于基模变成二次谐波。
}
\newcommand{\ThQPulse}{\ThermalFigure{
 \draw[->](-4,0)--(4,0)node[right]{$x$};\draw[->](0,-.2)--(0,2)node[above]{$y$};
 \draw[DeepBlue,thick,domain=-3.8:3.8,samples=90]plot(\x,{1.4*exp(-(\x)*(\x))});
 \node[right]at(4,1.3){$t=0:\ y=F(x),\ y_t=0$};
}{局域初始弦形，释放时每个质元的横向速度都为零。}{ch12:quality:pulse-q}}
\newcommand{\ThAPulse}{\ThermalFigure{
 \draw[->](-5,0)--(5,0)node[right]{$x$};
 \draw[DeepBlue,thick,domain=-4.8:4.8,samples=100]plot(\x,{.7*exp(-(\x-2.2)*(\x-2.2))+.7*exp(-(\x+2.2)*(\x+2.2))});
 \draw[->,thick](-2.2,1.1)--(-3.4,1.1)node[left]{$c$};\draw[->,thick](2.2,1.1)--(3.4,1.1)node[right]{$c$};
 \node[below]at(-2.2,-.3){$F(x+ct)/2$};\node[below]at(2.2,-.3){$F(x-ct)/2$};
}{两半幅行波相背离开；相遇重叠时仍按位移代数相加。}{ch12:quality:pulse-a}
图~\ref{ch12:quality:pulse-a}的左右两列波在 $t=0$ 重合，位移相加而速度相消。释放后两列行波各自同时具有动能与势能，所以能量核查应与初态的纯势能比较，不能只把振幅平方简单相加。
}
\newcommand{\ThQBeats}{\ThermalFigure{
 \foreach \x/\lab in {0/标准音叉,4/待测音叉}{\draw[thick](\x,2)--(\x,1)arc(180:360:.5)--(\x+1,2);\draw[thick](\x+.5,.5)--(\x+.5,-.2);\node[below]at(\x+.5,-.4){\lab};}
 \node[above]at(.5,2.4){$440\,\mathrm{Hz}$};\node[above]at(4.5,2.4){频率 $f$};
 \node[right,align=left]at(6,1){初始拍频 $3\,\mathrm{Hz}$\\加蜡后拍频 $1\,\mathrm{Hz}$};
}{两音叉共同发声；加蜡只改变待测音叉的有效振动质量。}{ch12:quality:beats-q}}
\newcommand{\ThABeats}{\ThermalFigure{
 \draw[->](0,0)--(8,0)node[right]{$f/\mathrm{Hz}$};
 \foreach\x/\lab in {1/437,4/440,7/443}{\draw(\x,-.15)--(\x,.15)node[above]{$\lab$};}
 \draw[->,DeepBlue,thick](7,-.7)--(5,-.7)node[midway,below]{趋近标准};
 \draw[->,thick](1,-.7)--(0,-.7)node[below]{远离};
}{加蜡使待测频率下降，只有高频候选会先靠近 $440\,\mathrm{Hz}$。}{ch12:quality:beats-a}
图~\ref{ch12:quality:beats-a}保留初始的两个频率候选，再沿“加质量使频率降低”的方向移动。低频候选向左会使频差增大，与拍频下降矛盾；因此可排除它，不能只凭 $|f-440|=3$ 直接选正号。
}
\newcommand{\ThQDoppler}{\ThermalFigure{
 \draw[thick](0,0)rectangle(2,1);\foreach\x in {.4,1.6}{\draw(\x,-.15)circle(.2);}
 \draw[->,thick](2.2,.5)--(4.2,.5)node[midway,above]{$v_s=20\,\mathrm{m/s}$};
 \fill(7.5,.5)circle(3pt)node[above]{静止观察者};\node[below]at(1,-.7){汽笛 $f_0=600\,\mathrm{Hz}$};
 \node[below]at(6,-.7){空气静止，$c=340\,\mathrm{m/s}$};
}{列车接近观察者；通过后驶离，声源的本征周期保持不变。}{ch12:quality:doppler-q}}
\newcommand{\ThADoppler}{\ThermalFigure{
 \draw[->](-1,0)--(7,0)node[right]{$x$};\fill(0,0)circle(2pt)node[below]{$S_0$};\fill(1.2,0)circle(2pt)node[below]{$S_1$};
 \draw[<->](0,-.8)--node[below]{$v_sT_s$}(1.2,-.8);\draw[<->](0,1.2)--node[above]{$cT_s$}(5.5,1.2);
 \draw[DeepBlue,thick](5.5,-.4)--(5.5,.7);\draw[DeepBlue,thick](1.2,-.4)--(1.2,.7);
 \draw[<->](1.2,.5)--node[above]{$\lambda_{\rm ahead}$}(5.5,.5);
}{相隔一个本征周期发射的两波面：前方间距等于声波前进量减去声源前进量。}{ch12:quality:doppler-a}
图~\ref{ch12:quality:doppler-a}取第一波面发出后经过 $T_s=1/f_0$ 的时刻，此时第二波面从 $S_1$ 刚发出。前方波长为 $(c-v_s)T_s$，后方则为 $(c+v_s)T_s$；静止观察者再用当地声速除以波长求频率，所以两种情形的分母符号不同。
}
\newcommand{\ThQSoundSources}{\ThermalFigure{
 \foreach\x in {0,2,4,6}{\foreach\y in {0,1.5}{\draw[thick](\x,\y)rectangle(\x+.7,\y+.7);\draw(\x+.35,\y+.35)circle(.2);}}
 \fill[DeepBlue](9,1.1)circle(3pt)node[above]{$P$};\node[right]at(10,1.1){相位互不相关};
 \node[below]at(3.5,-.5){每台单独在 $P$ 产生 $70\,\mathrm{dB}$};
}{八台独立声源在同一观察点叠加，给定的是各自单独的声强级。}{ch12:quality:sources-q}}
\newcommand{\ThASoundSources}{\ThermalFigure{
 \draw[thick](0,0)rectangle(8,1);\foreach\x in {1,...,7}{\draw(\x,0)--(\x,1);}
 \node[above]at(4,1.4){$I_{\rm tot}=I_1+\cdots+I_8$};\node[below]at(4,-.4){每格为单台声源的平均强度 $I_1$};
}{互不相干时相加强度；强度相加完成后再取对数。}{ch12:quality:sources-a}
图~\ref{ch12:quality:sources-a}中的八格表示强度的八份相加。相位互不相关使长期平均交叉项为零，才有 $I_{\rm tot}=8I_1$；对数刻度上增加的是 $10\log_{10}8$，不能把八个分贝数相加。
}
\newcommand{\ThSoundDistanceBase}[1]{
 \fill(0,0)circle(3pt)node[above]{甲};\fill(2,0)circle(2pt)node[above]{$M$};\fill[DeepBlue](4,0)circle(3pt)node[right]{$P$};
 \fill(4,2.2)circle(3pt)node[above]{乙};\draw[<->](0,-.6)--node[below]{$2\,\mathrm m$}(2,-.6);
 \draw[<->](0,-1.3)--node[below]{$4\,\mathrm m$}(4,-1.3);\draw[dashed](4,.2)--(4,2);#1}
\newcommand{\ThQSoundDistance}{\ThermalFigure{\ThSoundDistanceBase{
 \node[right]at(5.5,1.3){甲在 $M$ 单独：$80\,\mathrm{dB}$};\node[right]at(5.5,.2){乙在 $P$ 单独：$70\,\mathrm{dB}$};
}}{甲的测量点 $M$ 与待求观察点 $P$ 不同，乙的给定数据已经对应 $P$。}{ch12:quality:distance-q}}
\newcommand{\ThASoundDistance}{\ThermalFigure{\ThSoundDistanceBase{
 \draw[DeepBlue,dashed](0,0)circle(2);\draw[->,DeepBlue](2,.6)--(3.5,.6);
 \node[right]at(5.5,.7){先把甲换算到 $P$，再加乙的强度};
}}{同一声源传播距离加倍，球面面积变为四倍。}{ch12:quality:distance-a}
图~\ref{ch12:quality:distance-a}中甲的功率穿过以它为中心的球面。由 $4\pi r^2I$ 不变，先将 $M$ 点强度除以四得到 $P$ 点的甲源贡献；随后与乙在同一点的贡献相加，避免把不同位置的分贝数据直接相加。
}
\newcommand{\ThQDeepWave}{\ThermalFigure{
 \draw[->](0,0)--(6,0)node[right]{$k$};\draw[->](0,0)--(0,3.2)node[above]{$\omega$};
 \draw[DeepBlue,thick,domain=.05:5.5,samples=80]plot(\x,{1.1*sqrt(\x)});\node[right]at(6,2){$\omega^2=gk$};
}{深水、重力主导范围内的正频率色散曲线。}{ch12:quality:deep-q}}
\newcommand{\ThADeepWave}{\ThermalFigure{
 \draw[->](0,0)--(6,0)node[right]{$k$};\draw[->](0,0)--(0,3.5)node[above]{$\omega$};
 \draw[DeepBlue,thick,domain=.05:5.5,samples=80]plot(\x,{1.1*sqrt(\x)});
 \draw[thick](0,0)--(4.5,2.858);\draw[TextGray,thick](.3,1.048)--(5.4,2.668);
 \fill(3,1.905)circle(2pt)node[above left]{$(k_0,\omega_0)$};
 \node[right,align=left]at(6,1){割线斜率：$v_p=\omega_0/k_0$\\切线斜率：$v_g=\dd\omega/\dd k$};
}{相速度取原点割线斜率，群速度取同一点的切线斜率。}{ch12:quality:deep-a}
图~\ref{ch12:quality:deep-a}中根号曲线的切线比原点割线平缓，直接反映 $v_g<v_p$。对 $\omega=\sqrt{gk}$ 求导得到恰好一半的斜率，因此波峰穿过以较慢群速度运动的包络；两种速度不能由同一个 $\omega/k$ 重复计算。
}
\newcommand{\ThCutoffBase}[1]{
 \draw[->](0,0)--(4.5,0)node[right]{$c_0k/\omega_c$};\draw[->](0,0)--(0,4)node[above]{$\omega/\omega_c$};
 \draw[DeepBlue,thick,domain=0:3.6,samples=80]plot(\x,{sqrt(1+\x^2)});
 \node[left]at(0,1){$1$};\fill(1.732,2)circle(2pt)node[above left]{$\omega=2\omega_c$};#1}
\newcommand{\ThQCutoff}{\ThermalFigure{\ThCutoffBase{}}{色散曲线在 $k=0$ 有非零截止频率；采用正波数、正频率支。}{ch12:quality:cutoff-q}}
\newcommand{\ThACutoff}{\ThermalFigure{\ThCutoffBase{
 \draw[thick](0,0)--(3,3.464);\draw[TextGray,thick](.4,.846)--(3.4,3.444);
 \node[right]at(5,2){割线与切线分别给出相速度和群速度};
}}{在同一点比较两种斜率；接近截止时割线变陡、切线变平。}{ch12:quality:cutoff-a}
图~\ref{ch12:quality:cutoff-a}中的原点割线随 $k\to0^+$ 越来越陡，因为纵截距仍为 $\omega_c$；曲线自身的切线却逐渐水平。这解释了相速度发散而群速度趋零，两者的不同极限来自不同的几何定义。
}
\newcommand{\ThQClosedPipe}{\ThermalFigure{
 \draw[thick](0,-.6)--(6,-.6) (0,.6)--(6,.6);\draw[very thick](0,-.6)--(0,.6);
 \draw[<->](0,-1.2)--node[below]{$L=0.85\,\mathrm m$}(6,-1.2);
 \node[above]at(0,1){闭端};\node[above]at(6,1){开口};\node[right]at(7,0){$c=340\,\mathrm{m/s}$};
}{一端闭、一端开；边界条件应针对质元位移或压强分别说明。}{ch12:quality:closed-q}}
\newcommand{\ThAClosedPipe}{\ThermalFigure{
 \foreach \j/\yy in {1/3,3/1,5/-1}{\begin{scope}[shift={(0,\yy)}]
 \draw[dashed](0,0)--(5,0);\draw[very thick](0,-.7)--(0,.7);
 \draw[DeepBlue,thick,domain=0:5,samples=90]plot(\x,{.6*sin(\j*18*\x)});
 \node[right]at(5.4,0){$L=\j\lambda/4$};\end{scope}}
}{最低三个质元位移驻波：闭端为节点、开口为波腹，只出现奇数个四分之一波长。}{ch12:quality:closed-a}
图~\ref{ch12:quality:closed-a}中开口端的位移不为零，闭端位移为零；这与压强驻波的节点、波腹恰好相反。故允许的波数满足 $kL=(2n-1)\pi/2$，不能把两端同为节点的弦公式直接搬来。
}
\newcommand{\ThResonanceBase}[1]{
 \foreach \x/\y/\lab in {0/2.4/1,4/.7/2}{
 \draw[thick](\x,0)--(\x,3.2) (\x+1,0)--(\x+1,3.2);\draw[DeepBlue,thick](\x,\y)--(\x+1,\y);
 \fill[DeepBlue!15](\x,0)rectangle(\x+1,\y);\draw[<->](\x+1.5,\y)--node[right]{$L_{\lab}$}(\x+1.5,3.2);}
 \node[above]at(2.5,3.8){同一音叉 $f=500\,\mathrm{Hz}$};#1}
\newcommand{\ThQResonanceLength}{\ThermalFigure{\ThResonanceBase{}}{调节水面得到连续两次共振；长度从水面量到实际管口。}{ch12:quality:resonance-q}}
\newcommand{\ThAResonanceLength}{\ThermalFigure{\ThResonanceBase{
 \draw[dashed](-.2,3.5)--(1.2,3.5) (3.8,3.5)--(5.2,3.5);
 \draw[<->](-.6,3.2)--node[left]{$e$}(-.6,3.5);\draw[<->](3.4,3.2)--node[left]{$e$}(3.4,3.5);
}}{两次共振使用相同端部修正，等效开口均在实际管口外侧距离 $e$；图中修正间隔为辨认而放大。}{ch12:quality:resonance-a}
图~\ref{ch12:quality:resonance-a}的两条虚线表示等效开口，所以基模满足 $L_1+e=\lambda/4$，下一允许模式满足 $L_2+e=3\lambda/4$。两式相减消去未知 $e$，先求声速，再回代求修正量；不能先假定 $e=0$。
}
\newcommand{\ThQStringInterface}{\ThermalFigure{
 \draw[thick](-4,0)--(0,0);\draw[line width=2pt](0,0)--(4,0);\fill(0,0)circle(2pt);
 \draw[->,DeepBlue](-3,1)--(-1,1)node[midway,above]{入射};\node[below]at(-2,-.3){$\mu_1$};\node[below]at(2,-.3){$\mu_2=4\mu_1$};
 \node[above]at(0,2){两侧张力相同};
}{两根弦在无质量结点连接，波从轻弦传播到重弦。}{ch12:quality:interface-q}}
\newcommand{\ThAStringInterface}{\ThermalFigure{
 \draw[thick](-4,0)--(0,0);\draw[line width=2pt](0,0)--(4,0);\draw[dashed](0,-1.3)--(0,1.8);
 \draw[->,DeepBlue](-3,1)--(-1,1)node[midway,above]{$A_i$};\draw[<-,thick](-3,-.8)--(-1,-.8)node[midway,below]{$A_r$};
 \draw[->,DeepBlue](1,1)--(3,1)node[midway,above]{$A_t$};
 \node[below,align=center]at(0,-1.7){位移连续：$A_i+A_r=A_t$\\横向力连续：$Z_1(A_i-A_r)=Z_2A_t$};
}{反射波传播方向相反，在斜率边界条件中带相反符号。}{ch12:quality:interface-a}
图~\ref{ch12:quality:interface-a}把反射箭头指向左，因而横向力条件中出现 $A_i-A_r$，而位移条件仍是 $A_i+A_r$。求出振幅比后，透射能流还须乘阻抗比 $Z_2/Z_1$；只计算 $t^2$ 会漏掉重弦的能流权重。
}
\newcommand{\ThQMembrane}{\ThermalFigure{
 \draw[thick](0,0)circle(2);\draw[DeepBlue](0,0)circle(.7);\draw[DeepBlue](0,0)circle(1.2);\draw[DeepBlue](0,0)circle(1.7);
 \draw[->](0,0)--(2.8,0)node[right]{$r$};\draw[->,thick](.8,.8)--(1.6,1.6)node[right]{外行};
 \node[right,align=left]at(4,0){面密度 $\mu$\\单位长度张力 $\sigma$};
}{轴对称圆膜外行波的俯视图，波面为同心圆；环半径是传播变量。}{ch12:quality:membrane-q}}



\section{弹性形变、应力与材料定律}
\subsection{弹性体、物质坐标与小形变}

弹性（elasticity）指撤去使物体形变的作用后，物体能够恢复原有构形的性质；
“完全弹性”不等于“任意大形变都服从线性胡克定律”。固体可以储存剪切弹性能，
流体则主要以黏性应力抵抗形变速率。黏弹性材料的应力还与加载历史有关；
塑性形变可以留下不可恢复的部分。非牛顿流体、黏弹性与塑性属于流变学
（rheology）研究的范围，不能把三者简单等同。本章先研究小形变、可逆、线性弹性体。

用参考构形中某材料点的位置 $\vect R=(X,Y,Z)$ 给它作永久标记。
该点在时刻 $t$ 的实际位置为 $\vect r(\vect R,t)$，位移为
\begin{equation}
 \vect u(\vect R,t)=\vect r(\vect R,t)-\vect R.
 \label{ch12:eq:displacement}
\end{equation}
$\vect R$ 是“哪一个材料点”，$\vect r$ 是“它现在在哪里”；不能把位移场误认为速度场。
若相邻点的原间隔为 $\dd\vect R$，泰勒展开到一阶给出
\begin{equation}
 \dd u_i=\sum_{j=1}^{3}\frac{\partial u_i}{\partial X_j}\dd X_j,\qquad
 \dd\vect r=(\mathsf I+\mathsf H)\dd\vect R,\quad H_{ij}=\frac{\partial u_i}{\partial X_j}.
 \label{ch12:eq:gradient}
\end{equation}
这里 $X_1=X,X_2=Y,X_3=Z$，$\mathsf H$ 是九个分量组成的线性映射；
本章采用“行是位移分量、列是求导坐标”的矩阵约定。整体平移改变 $\vect u$，
却不改变相邻点的间隔，因而不能算作形变。

任何矩阵都有唯一分解
\[
 \mathsf H=\frac{\mathsf H+\mathsf H^{\mathsf T}}2+
            \frac{\mathsf H-\mathsf H^{\mathsf T}}2
          \equiv\boldsymbol\varepsilon+\mathsf W.
\]
上标 $\mathsf T$ 表示转置。$\boldsymbol\varepsilon$ 对称，$\mathsf W$ 反对称。
例如刚体发生小转角 $\delta\vect\phi$ 时，
$\dd\vect u=\delta\vect\phi\times\dd\vect R$，相应矩阵为
\[
 \mathsf W=
 \begin{pmatrix}
 0&-\delta\phi_z&\delta\phi_y\\
 \delta\phi_z&0&-\delta\phi_x\\
 -\delta\phi_y&\delta\phi_x&0
 \end{pmatrix}.
\]
它只转动微元，不改变微元各边长度和夹角。反过来，由
\[
 |\dd\vect r|^2-|\dd\vect R|^2
 =2\dd\vect R\cdot\boldsymbol\varepsilon\,\dd\vect R+O(H^2)
\]
可见，一阶长度变化只取决于对称部分。因此定义小应变张量（strain tensor）
\begin{equation}
 \varepsilon_{ij}=\frac12\left(\frac{\partial u_i}{\partial X_j}
                           +\frac{\partial u_j}{\partial X_i}\right),\qquad
 \boldsymbol\varepsilon=
 \begin{pmatrix}
 \varepsilon_{xx}&\gamma_{xy}/2&\gamma_{xz}/2\\
 \gamma_{xy}/2&\varepsilon_{yy}&\gamma_{yz}/2\\
 \gamma_{xz}/2&\gamma_{yz}/2&\varepsilon_{zz}
 \end{pmatrix}.
 \label{ch12:eq:strain}
\end{equation}
这一分解适用于小位移梯度；大角度刚体转动不能直接用截断的一阶应变判断。

\subsubsection{正应变与剪应变的几何意义}

沿 $X$ 方向原长为 $L$ 的小线段，伸长为 $\Delta L$，正应变为
$\varepsilon_{xx}=\Delta L/L=\partial u_x/\partial X$。
正号表示伸长，负号表示缩短。两个原来正交的单位线段，变形后方向分别近似为
$\uvec x+\partial_X\vect u$ 和 $\uvec y+\partial_Y\vect u$。
它们点积到一阶为 $\partial_Yu_x+\partial_Xu_y$。
若夹角从 $\pi/2$ 变成 $\pi/2-\alpha$，因
$\cos(\pi/2-\alpha)=\sin\alpha\simeq\alpha$，故
\begin{equation}
 \gamma_{xy}=2\varepsilon_{xy}
 =\frac{\partial u_x}{\partial Y}+\frac{\partial u_y}{\partial X}
 \simeq\alpha.
 \label{ch12:eq:shear-angle}
\end{equation}
$\gamma_{xy}$ 是工程剪应变，不能把它与张量分量 $\varepsilon_{xy}$ 混用。
例如 $u_x=\gamma Y,u_y=0$ 时，上边相对下边横移 $\gamma H$，
$\tan\alpha=\gamma\simeq\alpha$。正应变、剪应变均无量纲。
图\ref{ch12:fig:strain}还说明：同一受拉材料内，微元取向不同，看到的正应变、剪应变分量也不同。

\begin{center}\begin{minipage}{\linewidth}\centering
\begin{tikzpicture}[>=Latex,scale=.87]
 \draw[guide] (0,0) rectangle(4,1.25);
 \draw[thick] (-.35,.12) rectangle(4.35,1.13);
 \draw[vector] (-.35,.63)--(-1.3,.63);
 \draw[vector] (4.35,.63)--(5.3,.63);
 \draw[thick,DeepBlue] (.5,.4) rectangle(1.3,.85);
 \draw[thick,DeepBlue] (2.5,.25)--(3.15,.63)--(2.5,1)--(1.85,.63)--cycle;
 \node at (2,-.55) {同一受拉杆内的不同取向微元};
 \begin{scope}[shift={(7.1,0)}]
  \draw[guide] (0,0) rectangle(1.3,1.25);
  \draw[thick,DeepBlue] (0,0) rectangle(1.8,1.05);
  \draw[<->] (0,-.4)--(1.3,-.4) node[midway,below] {$L$};
  \node at (.9,1.75) {正应变};
 \end{scope}
 \begin{scope}[shift={(10.5,0)}]
  \draw[guide] (0,0) rectangle(1.3,1.25);
  \draw[thick,DeepBlue] (0,0)--(1.3,0)--(1.8,1.25)--(.5,1.25)--cycle;
  \draw[<->] (0,1.65)--(.5,1.65);
  \node[above] at (.25,1.65) {$\gamma H$};
  \draw (0,.65) arc(90:68.20:.65);
  \node[left] at (-.12,.55) {$\alpha$};
  \node at (.9,-.8) {剪应变};
 \end{scope}
\end{tikzpicture}
\figcap{受拉材料中的正应变与剪应变。虚线为参考构形；右侧将微元平移、转动对齐后比较长度和直角的变化。}
\label{ch12:fig:strain}
\end{minipage}\end{center}

\subsection{有向截面、应力与弹性定律}

\subsubsection{为什么内部受力需要应力张量}

只指定物体内部一点，尚不能指定“谁对谁施力”。需在该点切出法向为 $\vect n$ 的小面，
取 $\vect n$ 指向所研究部分的外侧，以外部对内部的面力除以面积，定义面力密度
$\vect t(\vect n)=\dd\vect F/\dd S$，单位为 $\mathrm{N/m^2}$，即帕斯卡。
同一点不同方向的截面通常有不同的面力；所以内部应力（stress）不是一个矢量。

面力与有向面积是线性关系，可以用小四面体证明：斜面的面积为 $\dd S$，
其三个坐标面投影面积分别为 $n_j\dd S$。体积力及惯性力是微元边长的三次量，
面力是二次量；微元趋于零时，主阶平衡给出
$\vect t(\vect n)=\sum_j n_j\vect t(\uvec e_j)$。
负的 $n_j$ 由反向截面的作用反作用关系 $\vect t(-\vect n)=-\vect t(\vect n)$ 处理。
因此定义
\begin{equation}
 t_i(\vect n)=\sum_j\sigma_{ij}n_j,\qquad
 \dd\vect F=\boldsymbol\sigma\,\vect n\,\dd S.
 \label{ch12:eq:traction}
\end{equation}
$\sigma_{ij}$ 表示法向沿第 $j$ 轴的面上，第 $i$ 方向的面力分量。
$\sigma_{nn}=\vect n\cdot\boldsymbol\sigma\vect n$ 为正应力，余下的面内分量为剪应力，
如图\ref{ch12:fig:stress}所示。拉伸正应力取正；流体压强 $p>0$ 对应 $\boldsymbol\sigma=-p\mathsf I$。

\begin{center}\begin{minipage}{\linewidth}\centering
\begin{tikzpicture}[>=Latex,scale=.93]
 \begin{scope}
  \draw[thick] (0,0)--(2.5,0)--(3.25,.8)--(.75,.8)--cycle;
  \coordinate(P) at(1.6,.4);
  \draw[vector] (P)--++(0,1.5) node[above] {$\vect t=\sigma_{nn}\vect n$};
  \draw[->,TextGray] (.55,.32)--(.55,1.15) node[left] {$\vect n$};
  \node at (1.6,-.6) {正应力：垂直截面};
 \end{scope}
 \begin{scope}[shift={(6,0)}]
  \draw[thick] (0,0)--(2.5,0)--(3.25,.8)--(.75,.8)--cycle;
  \coordinate(Q) at(1.35,.35);
  \draw[vector] (Q)--++(1.9,0) node[right] {$\vect t_\parallel$};
  \draw[->,TextGray] (Q)--++(0,1.2) node[above] {$\vect n$};
  \node at (1.6,-.6) {剪应力：平行截面};
 \end{scope}
\end{tikzpicture}
\figcap{正应力与剪应力是同一有向截面面力的垂直、平行分量。箭头起于截面内部，面力单位为力除以面积。}
\label{ch12:fig:stress}
\end{minipage}\end{center}

在没有微观偶力矩的普通连续介质中，应力还满足 $\sigma_{ij}=\sigma_{ji}$。
例如考察尺寸 $\dd x,\dd y,\dd z$ 的小长方体绕 $z$ 轴的力矩：
$x$ 法向两面上的 $y$ 向剪力合力矩为 $\sigma_{yx}\dd x\dd y\dd z$，
$y$ 法向两面上的 $x$ 向剪力合力矩为 $-\sigma_{xy}\dd x\dd y\dd z$。
正应力不贡献这一主阶力矩。微元转动惯量比体积多两个长度幂次，
故角动量方程的最低阶要求 $\sigma_{yx}-\sigma_{xy}=0$，其余两对同理。

\subsubsection{胡克定律、各向同性与独立弹性常数}

在指定坐标系中，应力映射也可显式写成
\[
 \begin{pmatrix}\dd F_x\\\dd F_y\\\dd F_z\end{pmatrix}
 =
 \begin{pmatrix}
 \sigma_{xx}&\sigma_{xy}&\sigma_{xz}\\
 \sigma_{yx}&\sigma_{yy}&\sigma_{yz}\\
 \sigma_{zx}&\sigma_{zy}&\sigma_{zz}
 \end{pmatrix}
 \begin{pmatrix}n_x\\n_y\\n_z\end{pmatrix}\dd S.
\]
对法向沿 $+y$ 的截面，法向分量为 $\sigma_{yy}\dd S$，
切向分量为 $\sigma_{xy}\dd S$、$\sigma_{zy}\dd S$，不能漏掉三维中的第二个切向方向。
小形变范围内的一般线性本构关系写成
$\sigma_{ij}=\sum_{k,l}C_{ijkl}\varepsilon_{kl}$。
四个指标各有三个取值，形式上有 $3^4=81$ 项。
应力和应变分别对称，使独立分量由 $9\times9$ 降为 $6\times6=36$。
若形变完全可逆、存在弹性能密度 $w(\boldsymbol\varepsilon)$，则
$\sigma_{ij}=\partial w/\partial\varepsilon_{ij}$（非对角独立分量要一致计数），
混合二阶导数相等给出 $C_{ijkl}=C_{klij}$，因此最多有 $6\times7/2=21$ 个独立常数。
这是弹性对称性和能量条件共同给出的结论，并非所有材料一开始只有两个常数。

各向同性意味着旋转坐标后本构关系不变。各向同性四阶张量只能由克罗内克符号
$\delta_{ij}$ 的三种两两配对组成；再用前述对称性，便只剩两个常数：
\begin{equation}
 \sigma_{ij}=\lambda_e\,\varepsilon_{kk}\delta_{ij}+2G\varepsilon_{ij},
 \qquad \varepsilon_{kk}=\varepsilon_{xx}+\varepsilon_{yy}+\varepsilon_{zz}.
 \label{ch12:eq:lame}
\end{equation}
这里重复指标 $k$ 表示对三个方向求和；$\delta_{ij}$ 在 $i=j$ 时为1，其余为0。
$\lambda_e,G$ 称为拉梅常数（Lamé constants）；用 $\lambda_e$ 的下标区别后文波长 $\lambda$。
在线性范围，$\varepsilon_{kk}=\Delta V/V$，因为
$(1+\varepsilon_{xx})(1+\varepsilon_{yy})(1+\varepsilon_{zz})-1$
的一阶项正是三者之和。

图\ref{ch12:fig:moduli}给出三种实验操作。纯剪切中
$\sigma_{xy}=2G\varepsilon_{xy}=G\gamma_{xy}$，故 $G$ 就是剪切模量。
静水压缩时 $\sigma_{xx}=\sigma_{yy}=\sigma_{zz}=-p$，取式\eqref{ch12:eq:lame}的迹得到
\begin{equation}
 -3p=(3\lambda_e+2G)\frac{\Delta V}{V},\qquad
 K=-V\frac{\dd p}{\dd V}=\lambda_e+\frac{2G}{3}.
 \label{ch12:eq:bulk}
\end{equation}
压强增加时体积减小，负号使稳定材料的体积模量 $K$ 为正。

\begin{center}\begin{minipage}{\linewidth}\centering
\begin{tikzpicture}[>=Latex,scale=.87]
 \draw[guide] (0,0) rectangle(2,1.5);
 \draw[thick,DeepBlue] (0,0)--(2,0)--(2.5,1.5)--(.5,1.5)--cycle;
 \draw[vector] (.7,1.9)--(2,1.9) node[right] {$\tau$};
 \draw[vector] (1.4,-.4)--(.1,-.4);
 \node at (1.2,-1) {剪切：$\tau=G\gamma$};
 \begin{scope}[shift={(5,0)}]
  \draw[guide] (0,0) rectangle(2,1.5);
  \draw[thick,DeepBlue] (.2,.15) rectangle(1.8,1.35);
  \draw[vector] (1,2.1)--(1,1.4);
  \draw[vector] (1,-.6)--(1,.1);
  \draw[vector] (-.7,.75)--(.15,.75);
  \draw[vector] (2.7,.75)--(1.85,.75);
  \node at (1,-1) {体压缩：$K$};
 \end{scope}
 \begin{scope}[shift={(10,0)}]
  \draw[guide] (0,0) rectangle(2,1.5);
  \draw[thick,DeepBlue] (-.3,.2) rectangle(2.3,1.3);
  \draw[vector] (-.3,.75)--(-1,.75);
  \draw[vector] (2.3,.75)--(3,.75);
  \node at (1,1.9) {横向收缩};
  \node at (1,-1) {单轴拉伸：$Y,\nu$};
 \end{scope}
\end{tikzpicture}
\figcap{剪切、均匀压缩、单轴拉伸。虚线是未变形轮廓；中图箭头都朝内，右图轴向伸长并伴随横向收缩。}
\label{ch12:fig:moduli}
\end{minipage}\end{center}

通常实验的参考构形已经处于大气压下；在其上增加小载荷时，
可先扣除共同的大气压及原有压缩，研究增量应力与增量应变。
弹性模量的数值依材料而异，一些常见固体可达 $10^{10}\,\mathrm{Pa}$ 或更高量级。
单轴拉伸时侧面自由，因此 $\sigma_{yy}=\sigma_{zz}=0$。
记 $\varepsilon_{xx}=\varepsilon_\parallel$，
$\varepsilon_{yy}=\varepsilon_{zz}=\varepsilon_\perp$。横向本构式给出
\[
 0=\lambda_e(\varepsilon_\parallel+2\varepsilon_\perp)+2G\varepsilon_\perp,
 \qquad
 \nu=-\frac{\varepsilon_\perp}{\varepsilon_\parallel}
     =\frac{\lambda_e}{2(\lambda_e+G)}.
\]
$\nu$ 为泊松比（Poisson's ratio）。正泊松比材料受拉后横向收缩；也存在负泊松比材料。
再将 $\varepsilon_\perp=-\nu\varepsilon_\parallel$ 代回轴向应力，
\begin{equation}
 \sigma_{xx}=Y\varepsilon_\parallel,\qquad
 Y=\frac{G(3\lambda_e+2G)}{\lambda_e+G}
  =2G(1+\nu)=3K(1-2\nu).
 \label{ch12:eq:young}
\end{equation}
$Y$ 是杨氏模量（Young's modulus），$Y,G,K,\lambda_e$ 的量纲均为压强。
稳定各向同性材料要求剪切与体压缩能量均为正，即 $G>0,K>0$，从而 $-1<\nu<1/2$。
对于原长 $L$、截面积 $A$ 的均匀细杆，
$F/A=Y\Delta L/L$，故整体等效劲度系数 $k=YA/L$；
材料的刚硬程度与整个物体的劲度不是同一个物理量。

\Needspace{16\baselineskip}\begin{origexamplebox}{自重引起的伸长}
\label{ch12:ex:01}
长为 $L$、截面积为 $A$、密度为 $\rho$、杨氏模量为 $Y$ 的均匀细杆竖直悬挂，
求自重引起的总伸长。取小形变，$L,A,\rho$ 均指参考构形的量。

\ThQSelfWeight
\end{origexamplebox}
\Needspace{6\baselineskip}\begin{xesolution}
从杆的自由下端向上量原构形坐标 $x$。图\ref{ch12:fig:selfweight}右侧取出截面以下的部分：
它长 $x$、质量 $\rho Ax$，上方截面的拉力平衡其全部重量，因此
\[
 F(x)=\rho Agx,\qquad \varepsilon(x)=\frac{F(x)}{YA}=\frac{\rho gx}{Y}.
\]
\begin{center}\begin{minipage}{\linewidth}\centering
\begin{tikzpicture}[>=Latex,scale=.85]
 \draw[thick] (-.9,3.5)--(.9,3.5);
 \draw[thick] (-.25,0) rectangle(.25,3.5);
 \draw[guide] (-.6,1.7)--(.7,1.7);
 \draw[<->] (-.7,0)--(-.7,1.7) node[midway,left] {$x$};
 \draw[<->] (.8,0)--(.8,3.5) node[midway,right] {$L$};
 \node[below] at (0,0) {自由端};
 \begin{scope}[shift={(4,0)}]
  \draw[thick] (-.25,0) rectangle(.25,1.7);
  \draw[vector] (0,1.7)--(0,2.7) node[above] {$F(x)$};
  \draw[vector] (0,.85)--(0,-.5) node[below] {$\rho Axg$};
  \node[right] at (.5,1) {截面以下的杆};
 \end{scope}
\end{tikzpicture}
\figcap{自重拉伸的坐标与受力。右图的向上拉力作用于截面，向下重力作用于所取部分。}
\label{ch12:fig:selfweight}
\end{minipage}\end{center}
微段原长 $\dd x$，伸长为 $\varepsilon(x)\dd x$，所以
\[
 \boxed{\Delta L=\int_0^L\frac{\rho gx}{Y}\dd x=\frac{\rho gL^2}{2Y}}.
\]
$F$ 从下端的零增至上端的 $\rho AgL$，故平均应变恰为最大值的一半。
截面积消去，是因为增大截面同时增大承重和抗拉能力。$\rho gL/Y$ 无量纲，
$\Delta L$ 是长度；小形变还要求 $\rho gL/Y\ll1$。非均匀杆须重算每一截面以下的重量。
\end{xesolution}

\section{一维机械波的建立与传播}
\subsection{细杆纵波与弹性能}

设侧面自由的均匀细杆截面积为 $A$，纵向位移为 $u(x,t)$。
$x$ 标记原构形位置；图\ref{ch12:fig:rod-element}中的两端实际位置分别为
$x+u(x,t)$、$x+\dd x+u(x+\dd x,t)$，因此
\[
 \dd L-\dd x=u(x+\dd x,t)-u(x,t)=u_x\dd x,\qquad\varepsilon=u_x.
\]
用下标 $x,t$ 表示偏导数。微元右端受力向右为 $A\sigma(x+\dd x)$，
左端受力向左为 $A\sigma(x)$，其差为 $A\sigma_x\dd x$，而质量为 $\rho A\dd x$。
于是
\begin{equation}
 \rho A\dd x\,u_{tt}=A\sigma_x\dd x,\qquad
 u_{tt}=\frac{Y}{\rho}u_{xx},\qquad c_{\rm rod}=\sqrt{\frac{Y}{\rho}}.
 \label{ch12:eq:rod-wave}
\end{equation}
这里已用 $\sigma=Yu_x$。该一维杆波近似还要求波长远大于横向尺寸，
横向惯性可以忽略，侧面能够随泊松效应自由伸缩。

\begin{center}\begin{minipage}{\linewidth}\centering
\begin{tikzpicture}[>=Latex,scale=.88]
 \draw[thick] (-.5,0)--(8,0);
 \draw[thick] (1,-.2)--(1,.2) (4,-.2)--(4,.2);
 \node[above] at (1,.2) {$x$}; \node[above] at (4,.2) {$x+\dd x$};
 \draw[<->] (1,-.5)--(4,-.5) node[midway,below] {$\dd x$};
 \draw[thick,DeepBlue] (-.5,-2)--(8,-2);
 \draw[thick,DeepBlue] (1.5,-2.2)--(1.5,-1.8) (5,-2.2)--(5,-1.8);
 \draw[->,TextGray] (1,-1.3)--(1.5,-1.3) node[midway,above] {$u(x)$};
 \draw[->,TextGray] (4,-1.3)--(5,-1.3) node[midway,above] {$u(x+\dd x)$};
 \node[below] at (3.25,-2.2) {$\dd L=\dd x+u_x\dd x$};
 \draw[thick] (1.5,-4.3) rectangle(5,-3.6);
 \draw[vector] (1.5,-3.95)--(-.3,-3.95) node[above] {$A\sigma(x)$};
 \draw[vector] (5,-3.95)--(7.2,-3.95) node[above] {$A\sigma(x+\dd x)$};
 \node[below] at (3.25,-4.4) {$\dd m=\rho A\dd x$};
\end{tikzpicture}
\figcap{细杆微元的三层图：原位置、变形后位置、取出微元后的两端受力。位移差决定应变，应力差决定加速度。}
\label{ch12:fig:rod-element}
\end{minipage}\end{center}

单位参考体积的动能密度为 $\rho u_t^2/2$。准静态地把应变从零增加到 $\varepsilon$，
外力每单位体积所做的功为
$\int_0^\varepsilon \sigma(\eta)\dd\eta=\int_0^\varepsilon Y\eta\dd\eta=Y\varepsilon^2/2$。
因此总能量密度为
\begin{equation}
 w=\frac12\rho u_t^2+\frac12Yu_x^2.
 \label{ch12:eq:rod-energy}
\end{equation}
$\sigma\,\dd\varepsilon$ 量纲是 $\mathrm{J/m^3}$；不能把它当作整根杆的能量。

\subsection{弦的横波、两个偏振与能流}

张力为 $T$、线密度为 $\mu$ 的柔弦沿 $x$ 轴伸展，两个横向位移分别为
$\xi(x,t),\eta(x,t)$，空间位置为 $(x,\xi,\eta)$。
相邻点的切向线元与单位切向量为
\[
 \dd\vect r=(1,\xi_x,\eta_x)\dd x,\qquad
 \vect t=\frac{(1,\xi_x,\eta_x)}{\sqrt{1+\xi_x^2+\eta_x^2}}
 \simeq(1,\xi_x,\eta_x).
\]
小斜率 $|\xi_x|,|\eta_x|\ll1$ 时，张力大小到一阶视为恒定。
图\ref{ch12:fig:string-element}中微元受力为右端 $T\vect t(x+\dd x)$ 与左端 $-T\vect t(x)$；
两者的横向差给出
\begin{equation}
 \mu\,\xi_{tt}=T\xi_{xx},\qquad
 \mu\,\eta_{tt}=T\eta_{xx},\qquad c_{\rm string}=\sqrt{\frac T\mu}.
 \label{ch12:eq:string-wave}
\end{equation}
这里 $T$ 是力，单位 $\mathrm N$；$\mu$ 是质量除以长度，故 $T/\mu$ 的量纲为速度平方。
两个横向自由度具有相同波速，不表示弦元沿传播方向随波前进。

\begin{center}\begin{minipage}{\linewidth}\centering
\begin{tikzpicture}[x={(1cm,0cm)},y={(-.5cm,-.42cm)},z={(0cm,1cm)},>=Latex]
 \draw[axis] (0,0,0)--(6.4,0,0) node[right] {$x$};
 \draw[axis] (0,0,0)--(0,2.2,0) node[below left] {$y\ (\xi)$};
 \draw[axis] (0,0,0)--(0,0,1.9) node[above] {$z\ (\eta)$};
 \draw[curve,domain=.25:5.8,samples=90,variable=\q]
  plot ({\q},{.65*sin(50*\q)},{.5*cos(50*\q)});
 \coordinate(P) at(3,{.65*sin(150)},{.5*cos(150)});
 \coordinate(Q) at(3.7,{.65*sin(185)},{.5*cos(185)});
 \fill (P) circle(1.7pt); \fill(Q) circle(1.7pt);
 \draw[vector] (P)--++(-1.1,{-.65*.872665*cos(150)*1.1},{.5*.872665*sin(150)*1.1});
 \draw[vector] (Q)--++(1.1,{.65*.872665*cos(185)*1.1},{-.5*.872665*sin(185)*1.1});
 \node[above] at (2.25,0,.35) {$-T\vect t(x)$};
 \node[below] at (5,0,-1.1) {$T\vect t(x+\dd x)$};
 \draw[guide] (P)--(3,0,0); \draw[guide] (Q)--(3.7,0,0);
 \draw[<->] (3,0,1.1)--(3.7,0,1.1) node[midway,above] {$\dd x$};
\end{tikzpicture}
\figcap{右手坐标系中的弦微元。实线由同一空间曲线生成，两个力箭头沿其端点切线、背离所取微元。}
\label{ch12:fig:string-element}
\end{minipage}\end{center}

弦元长度
$\dd s=\sqrt{1+\xi_x^2+\eta_x^2}\dd x
\simeq[1+(\xi_x^2+\eta_x^2)/2]\dd x$。
在恒定预张力近似下，多出的势能为 $T(\dd s-\dd x)$，所以单位长度的能量为
\begin{equation}
 \mathcal E=\frac12\mu(\xi_t^2+\eta_t^2)+\frac12T(\xi_x^2+\eta_x^2).
 \label{ch12:eq:string-energy}
\end{equation}
以一个偏振 $y$ 为例，截面左侧通过张力对右侧做功的功率为
$P=-Ty_xy_t$。右行波 $y=A\cos(kx-\omega t)$ 满足
$y_x=-kA\sin(kx-\omega t)$、$y_t=\omega A\sin(kx-\omega t)$，
因此 $P=Tk\omega A^2\sin^2(kx-\omega t)\ge0$。周期平均得到
\begin{equation}
 \overline P=\frac12Tk\omega A^2
 =\frac12\mu c_{\rm string}\omega^2A^2,\qquad
 Z=\frac T{c_{\rm string}}=\mu c_{\rm string}=\sqrt{\mu T}.
 \label{ch12:eq:string-power}
\end{equation}
$Z$ 为弦的特征阻抗，单位为 $\mathrm{kg/s}$。线性、无色散近似下，
波速由介质与预应力决定，不由振幅和频率决定；固定频率下能流与振幅平方成正比。

\subsection{一维波动方程、行波与驻波}

对 $u_{tt}=c^2u_{xx}$，令特征坐标 $\zeta=x-ct,\ \chi=x+ct$。
链式法则给出
$\partial_x=\partial_\zeta+\partial_\chi$、
$\partial_t=-c\partial_\zeta+c\partial_\chi$，因此
\[
 u_{tt}-c^2u_{xx}=-4c^2u_{\zeta\chi}=0.
\]
先对 $\chi$ 积分，$u_\zeta$ 只依赖 $\zeta$；再对 $\zeta$ 积分，得到
\begin{equation}
 u(x,t)=f(x-ct)+g(x+ct).
 \label{ch12:eq:dalembert}
\end{equation}
第一项保持形状向右平移，第二项向左平移。若给定
$u(x,0)=F(x)$、$u_t(x,0)=V(x)$，则
$f+g=F$、$-cf'+cg'=V$，解这两式并积分可得
\[
 u(x,t)=\frac{F(x-ct)+F(x+ct)}2+
 \frac1{2c}\int_{x-ct}^{x+ct}V(s)\dd s.
\]
这说明求波的运动还需初始位移与初始速度，只有波动方程本身不足以唯一确定波形。

简谐右行波为 $u=A\cos(kx-\omega t+\varphi)$，
$\omega=ck,\ k=2\pi/\lambda,\ f=\omega/(2\pi),\ T_{\rm per}=1/f$。
$kx-\omega t+\varphi$ 为相位。同一相位点满足 $k\dd x-\omega\dd t=0$，
故速度为 $c=\omega/k$。图\ref{ch12:fig:travel-standing}左侧的瞬时零位移点会移动，
不能叫作驻波的“波节”。

两列等幅、同频、反向行波可组成
$u=2A\sin kx\cos\omega t$。其空间振幅为 $|2A\sin kx|$，
$x=n\pi/k$ 始终不动，是波节；两节之间振幅最大的地方是波腹。
固定端要求 $u=0$，反射位移反相；自由端要求横向力 $Tu_x=0$，反射位移同相。
两端固定、间距 $L$ 时 $\sin kL=0$，因此
\begin{equation}
 k_n=\frac{n\pi}{L},\quad \lambda_n=\frac{2L}{n},\quad
 f_n=\frac{nc}{2L},\qquad n=1,2,\ldots.
 \label{ch12:eq:fixed-modes}
\end{equation}

\begin{center}\begin{minipage}{\linewidth}\centering
\begin{tikzpicture}[>=Latex,x=.85cm,y=.75cm]
 \draw[axis] (-.3,0)--(6.6,0) node[right] {$x$};
 \draw[curve,domain=0:6,samples=100] plot(\x,{sin(60*\x)});
 \draw[guide,domain=0:6,samples=100] plot(\x,{sin(60*(\x-.65))});
 \draw[vector] (2.6,1.5)--(4,1.5) node[midway,above] {$c$};
 \node[above] at (1.5,1.05) {波峰};
 \node[below] at (4.5,-1.05) {波谷};
 \draw[<->] (1.5,-1.8)--(4.5,-1.8) node[midway,below] {$\lambda/2$};
 \node at (3,-2.8) {行波：实、虚线为不同时刻};
 \begin{scope}[shift={(9,0)}]
  \draw[axis] (-.3,0)--(6.6,0) node[right] {$x$};
  \draw[curve,domain=0:6,samples=100] plot(\x,{sin(60*\x)});
  \draw[guide,domain=0:6,samples=100] plot(\x,{-.5*sin(60*\x)});
  \foreach \p in {0,3,6}{\fill(\p,0)circle(2pt);}
  \node[above] at (1.5,1.05) {波腹};
  \draw[thin,TextGray] (3,-.15)--(3,-1.4);
  \node[below] at (3,-1.55) {波节};
  \node at (3,-2.8) {驻波：波节位置不随时间改变};
 \end{scope}
\end{tikzpicture}
\figcap{行波与驻波的区别。左侧虚线是平移后的波形；右侧虚线仅改变各点的振幅比例，波节保持不动。}
\label{ch12:fig:travel-standing}
\end{minipage}\end{center}

\Needspace{16\baselineskip}\begin{origexamplebox}{弦的基频与张力}
\label{ch12:ex:02}
长 $L$、线密度 $\mu$ 的弦两端固定，基频为 $f_1$。
若长度、线密度不变而基频提高为 $2f_1$，张力应变为原来的多少倍？

\ThQStringFrequency
\end{origexamplebox}
\Needspace{6\baselineskip}\begin{xesolution}
\ThAStringFrequency

基频对应两端各一个波节、中间一个波腹，由式\eqref{ch12:eq:fixed-modes}，
$f_1=c/(2L)=\sqrt{T/\mu}/(2L)$。两次状态相除，
\[
 \frac{2f_1}{f_1}=\frac{\sqrt{T'/\mu}}{\sqrt{T/\mu}}
 =\sqrt{\frac{T'}T},\qquad \boxed{T'=4T}.
\]
频率加倍先要求波速加倍，而张力与波速平方成正比，所以不是两倍。
若张力变化使弦的长度、线密度明显改变，就必须同时计入这些变化。
\end{xesolution}

两个横向分量可以有不同相位。令
$\xi=A\cos(kx-\omega t)$、
$\eta=B\cos(kx-\omega t+\delta)$。
在固定材料点处，$\delta=0,\pi$ 时运动轨迹为直线，称线偏振；
$A=B,\delta=\pm\pi/2$ 时轨迹为圆，称圆偏振；一般情形为椭圆偏振。
这是介质点在垂直传播方向平面内的运动轨迹，不是波的传播轨迹。
\subsubsection{波动方程不等于完整的运动答案}
质点运动必须给初位置、初速度；连续介质则要对每一个位置给初位移、初速度。波动方程只规定变化规律，初始条件才决定是哪一次运动，边界条件又决定传播到端点后的反射方式。
对无界弦，右行项 $f(x-ct)$ 在 $x$ 增加 $c\Delta t$、时间增加 $\Delta t$ 后保持相同的自变量，所以波形向右；左行项 $g(x+ct)$ 同理向左。固定材料点只随时间上下振动，不能把波形的水平平移理解成弦本身向右滑动。

\Needspace{16\baselineskip}\begin{origexamplebox}{松手后分裂成两份的弦脉冲}\label{ch12:ex:teaching-initial-pulse}
一根足够长的均匀弦，波速为 $c$。在 $t=0$ 时拉成位移形状 $F(x)$ 并从静止释放，即 $y(x,0)=F(x)$、$y_t(x,0)=0$；$F$ 为局域的平滑函数。求后续波形，并解释为什么两列波的振幅各为原来的一半，能量却没有丢失。

\ThQPulse
\end{origexamplebox}
\Needspace{6\baselineskip}\begin{xesolution}
\ThAPulse

写 $y=f(x-ct)+g(x+ct)$。初位移给 $f(x)+g(x)=F(x)$；初速度给 $-cf'(x)+cg'(x)=0$。后式说明 $f-g$ 为常量。局域条件使远处位移趋零，常量可取零，故
\[
 y(x,t)=\frac12F(x-ct)+\frac12F(x+ct).
\]
每一份以速度 $c$ 向相反方向传播。代回 $t=0$，两半重叠恢复 $F$，两者的速度贡献反号，确实使初速度为零。
设弦张力为 $\mathcal T$、线密度为 $\mu$，满足 $\mathcal T=\mu c^2$。初始全是弹性势能：
\[
 E_0=\frac{\mathcal T}{2}\int_{-\infty}^{\infty}[F'(x)]^2\dd x.
\]
对右行半幅脉冲，$y_x=F'/2$、$y_t=-cF'/2$，动能与势能相等，其总能量为
\[
 E_+=\int\left(\frac{\mu c^2}{8}[F']^2+
 \frac{\mathcal T}{8}[F']^2\right)\dd x=\frac12E_0.
\]
左行脉冲同样带走 $E_0/2$。把“一半振幅对应四分之一能量”直接套在初始静止弦形上，会漏掉行波中的动能；初始状态与单向行波的动能条件不同。该解还要求端点反射尚未返回所研究区域。
\end{xesolution}

\section{声波、拍频与多普勒效应}
\subsection{流体中的声波与绝热声速}

静态流体不能维持剪切应力，在无黏理想流体模型中 $G=0$，只剩压缩纵波。
设平衡密度为 $\rho_0$，小纵位移为 $u(x,t)$。质量守恒给出
$\delta\rho=-\rho_0 u_x$：伸长区域密度减小。
小压强变化为
$\delta p=(\partial p/\partial\rho)_{\rm process}\delta\rho=-K u_x$。
两端压强差形成力 $-\partial_x\delta p$，所以
\[
 \rho_0u_{tt}=-\partial_x\delta p=K u_{xx},\qquad
 c^2=\frac K{\rho_0}=\left(\frac{\partial p}{\partial\rho}\right)_{\rm process}.
\]
这里“过程”必须说明，否则 $K$ 不唯一。普通空气中声波的压缩、稀疏较快，
在一周期内、一个波长尺度上来不及通过热传导抹平温度差；忽略热耗散时取局部可逆绝热过程。
它不是说气体温度不变，而是每一物质微元与邻域的净换热可忽略。

理想气体满足 $pV^\gamma=\text{常量}$，其中 $\gamma=C_p/C_V$。
对固定质量微元求微分：
$\dd p/p=-\gamma\,\dd V/V$，于是
\begin{equation}
 K_s=-V\left(\frac{\partial p}{\partial V}\right)_s=\gamma p,\qquad
 c=\sqrt{\frac{\gamma p}{\rho_0}}
  =\sqrt{\frac{\gamma RT_0}{M_{\rm mol}}}.
 \label{ch12:eq:sound-speed}
\end{equation}
最后一步用 $pV=nRT_0$ 与 $\rho_0=nM_{\rm mol}/V$，
$M_{\rm mol}$ 为摩尔质量而不是单个分子质量。若使用单分子质量 $m_0$，
则应写 $c=\sqrt{\gamma k_{\rm B}T_0/m_0}$。

牛顿（Newton）采用等温压缩关系，会得到 $K_T=p$；
拉普拉斯（Laplace）的绝热修正使预测声速乘以 $\sqrt\gamma$。
例如 $T_0=293\,\mathrm K$、$M_{\rm mol}=0.029\,\mathrm{kg/mol}$、$\gamma=1.4$ 时，
等温模型给出约 $290\,\mathrm{m/s}$，绝热模型给出约 $343\,\mathrm{m/s}$。
这是在同一温度、同一气体参数下的比较，不能拿不同条件的历史近似数值混作精确测量。



\subsection{拍频与近频率叠加}

同方向、同振幅而角频率接近的两列简谐波在固定位置可写成
\begin{align}
 y&=A\cos(\omega_1t+\phi_1)+A\cos(\omega_2t+\phi_2)\\
 &=2A\cos\!\left(\frac{\Delta\omega}{2}t+\frac{\Delta\phi}{2}\right)
 \cos\!\left(\bar\omega t+\bar\phi\right),
\end{align}
其中 $\Delta\omega=\omega_1-\omega_2$、$\bar\omega=(\omega_1+\omega_2)/2$。
快速振动的包络振幅缓慢变化。相邻两次响度极大的时间间隔对应包络绝对值的周期，故拍频
\begin{equation}
 f_b=|f_1-f_2|.
\end{equation}
拍频不是第三个独立本征频率，而是线性叠加后能量或响度包络的变化频率。

\begin{center}\begin{minipage}{\linewidth}\centering
\begin{tikzpicture}
\begin{axis}[width=.82\linewidth,height=5.0cm,axis lines=middle,
  xlabel={$t$},ylabel={$y$},ylabel style={at={(axis description cs:-.025,.5)},anchor=east},xmin=0,xmax=12.6,ymin=-2.3,ymax=3.2,ticks=none,
  every axis label/.style={fill=none}]
  \addplot[DeepBlue,very thick,domain=0:12.6,samples=650]{2*cos(deg(.5*x))*cos(deg(7*x))};
  \addplot[TextGray,dashed,domain=0:12.6,samples=180]{2*abs(cos(deg(.5*x)))};
  \addplot[TextGray,dashed,domain=0:12.6,samples=180]{-2*abs(cos(deg(.5*x)))};
  \draw[<->] (axis cs:0,2.55)--(axis cs:6.283185,2.55) node[midway,above] {$T_b=1/|f_1-f_2|$};
\end{axis}
\end{tikzpicture}
\figcap{两近频率简谐振动叠加形成拍。虚线是振幅包络，不是介质中另一列波。}\label{ch12:fig:beats}
\end{minipage}\end{center}
图~\ref{ch12:fig:beats} 中相邻两个包络极大之间就是一个拍周期，快速载波在此期间振动许多次。

\Needspace{16\baselineskip}\begin{origexamplebox}{用拍频校准音叉}
\label{ch12:ex:03}
标准音叉频率为 $440\,\mathrm{Hz}$，待测音叉与它同时发声时每秒听到3次响度增强。在待测音叉上粘少量蜡后，拍频降为 $1\,\mathrm{Hz}$。求待测音叉原频率。

\ThQBeats
\end{origexamplebox}
\Needspace{6\baselineskip}\begin{xesolution}
\ThABeats

初始只有两个候选：
\[
 f=440\pm3\,\mathrm{Hz},
\]
即 $437\,\mathrm{Hz}$ 或 $443\,\mathrm{Hz}$。粘蜡增加振动系统的有效质量，使待测频率降低。若原来是
$443\,\mathrm{Hz}$，降低后会接近 $440\,\mathrm{Hz}$，拍频可从3降到1；若原来是 $437\,\mathrm{Hz}$，继续降低只会使频差增大。
因此
\[
 \boxed{f=443\,\mathrm{Hz}}.
\]
方向判断是这道题不可省略的一步；单凭初始拍频不能确定高低。
\end{xesolution}

\subsection{多普勒效应}

声波相对介质传播。这里讨论均匀静止介质中的共线运动、声源亚声速 $|v_s|<c$，所有速度均相对介质测量。设介质中声速为 $c$，观察者以速度 $v_o$ 向声源运动，声源以速度 $v_s$ 向观察者运动，则
\begin{equation}
 f'=f\frac{c+v_o}{c-v_s}.
\end{equation}
若实际运动方向相反，应给相应速度取负号。分子来自观察者遇到波峰的速率变化，分母来自移动声源改变波长。
对于真空光，不能套用“相对介质速度”的声学公式，应使用相对论多普勒关系。

\begin{center}\begin{minipage}{\linewidth}\centering
\begin{tikzpicture}[scale=.85,>=Latex]
  \foreach \n in {1,2,3,4}{
   \draw[TextGray] ({-.245*\n},0) circle({.7*\n});
   \fill[TextGray] ({-.245*\n},0) circle(1pt);
  }
  \fill[DeepBlue] (0,0) circle(2.8pt);
  \node[below] at (0,-.12) {$S$};
  \fill (3.9,0) circle(2.8pt) node[right=4pt] {$O$};
  \draw[vector] (.1,-3.3)--(1.3,-3.3) node[midway,below] {$v_s$};
  \draw[vector] (3.8,-1.1)--(2.6,-1.1) node[midway,below] {$v_o$};
  \draw[vector] (.3,3.1)--(1.7,3.1) node[midway,above] {$c$};
  \draw[<->] (1.365,-.30)--(1.82,-.30);
  \draw[guide] (1.60,-.35)--(2.6,-2.35);
  \node[right] at (2.7,-2.45) {$\lambda_{\text{前}}=(c-v_s)/f$};
  \draw[<->] (-3.78,-.35)--(-2.835,-.35);
  \node[below] at (-3.3,-.45) {$\lambda_{\text{后}}$};
  \node at (-1.5,-4.3) {圆心是历次发射位置，非当前声源位置};
\end{tikzpicture}
\figcap{声源和观察者相向运动的符号约定。波面相对介质以声速 $c$ 传播。}\label{ch12:fig:doppler}
\end{minipage}\end{center}
图~\ref{ch12:fig:doppler}中，第 $n$ 个波面发出于 $-n/f$ 时，
圆心为 $x=-nv_s/f$、半径为 $nc/f$。相邻前方波峰间距为
$(c-v_s)/f$，观察者迎着波峰运动，单位时间跨过 $(c+v_o)/\lambda_{\text{前}}$ 个波峰，
这就推出上面的频率公式。图固定了正号方向：$v_s,v_o>0$ 都使接收频率升高。

\Needspace{16\baselineskip}\begin{origexamplebox}{迎面驶来的鸣笛列车}
\label{ch12:ex:04}
列车以 $20\,\mathrm{m/s}$ 接近静止观察者，汽笛本征频率 $600\,\mathrm{Hz}$，空气中声速
$340\,\mathrm{m/s}$。求接近时和驶离后的听到频率，忽略风。

\ThQDoppler
\end{origexamplebox}
\Needspace{6\baselineskip}\begin{xesolution}
\ThADoppler

观察者静止，$v_o=0$。接近时声源向观察者，
\[
 f_{\rm app}=600\frac{340}{340-20}
 =\boxed{637.5\,\mathrm{Hz}}.
\]
驶离时 $v_s$ 取负，
\[
 f_{\rm rec}=600\frac{340}{340+20}
 =\boxed{566.7\,\mathrm{Hz}}.
\]
接近频率高、驶离频率低，符合波面前方被压缩、后方被拉长的图像。分式分子、分母都是速度，故频率单位保持不变。
\end{xesolution}

\section{声能传播、声强与声强级}
\subsection{声强与声强级}

声强 $I$ 是单位面积上平均能流功率，单位为 $\mathrm{W/m^2}$。点声源在均匀无耗介质中各向同性辐射时
\begin{equation}
 I(r)=\frac{P}{4\pi r^2}.
\end{equation}
声强级定义为
\begin{equation}
 L_I=10\log_{10}\frac I{I_0}\quad(\mathrm{dB}),
 \qquad I_0=10^{-12}\,\mathrm{W/m^2}.
\end{equation}
分贝是对数比值；多个互不相干声源应先相加强度，再取对数，不能直接相加分贝数。


\paragraph{距离规律与叠加规律都有条件}
球面面积 $4\pi r^2$ 来自以点声源为中心的同心球面。只有声功率没有被吸收、没有其他声源或反射波叠加，且辐射近似各向同性时，穿过这些球面的总功率才相同，从而得到平方反比律。靠近有限尺寸振源、具有强方向性的声束或封闭房间的混响，均需另行分析。
对相同阻抗中沿同一方向传播的相干平面波，强度可写成
\[
 I=I_1+I_2+2\sqrt{I_1I_2}\cos\phi.
\]
只有相位差在测量时间内充分变化、交叉项平均为零时，才可取 $I=I_1+I_2$。两个等强、不相干声源合成后的声强级增量是 $10\log_{10}2\simeq3.01\,\mathrm{dB}$；这个数来自强度倍增，不是把原来的分贝数再乘2。声压级则用 $20\log_{10}(p_{\rm rms}/p_0)$，指数前的20来自强度与声压平方的关系，使用时还须给出参考声压和介质条件。


\Needspace{16\baselineskip}\begin{origexamplebox}{相同声源的声强级叠加}
\label{ch12:ex:05}
一台机器在某点产生 $70\,\mathrm{dB}$ 的声强级。若有8台相同机器独立运行，且它们在该点的声波相位互不相关，求总声强级。

\ThQSoundSources
\end{origexamplebox}
\Needspace{6\baselineskip}\begin{xesolution}
\ThASoundSources

每台强度相同，非相干时总强度为 $I_{\rm tot}=8I$。因此
\[
 L_{\rm tot}-L_1=10\log_{10}8\approx9.03\,\mathrm{dB}.
\]
故
\[
 \boxed{L_{\rm tot}\approx79.0\,\mathrm{dB}}.
\]
不是 $8\times70\,\mathrm{dB}$。若只增加到两台，相同推理给出约 $3.01\,\mathrm{dB}$ 的增量。
\end{xesolution}

\subsubsection{声强与声压、质元速度怎样联系}
声音传播时，空气质元在平衡位置附近往复运动，并不是从声源一路飞到听者。取沿 $+x$ 传播的平面简谐声波，位移为
\[
 \xi(x,t)=\xi_m\cos(kx-\omega t),\qquad \omega=ck.
\]
质元速度是对时间求偏导，而波速是波形的传播速度，二者不能混淆：
\[
 v_{\rm part}=\partial_t\xi=\omega\xi_m\sin(kx-\omega t).
\]
小形变的体压缩量为 $-\partial_x\xi$；用绝热体积模量 $B_s=\rho_0c^2$，声压扰动为
\[
 p'=-B_s\partial_x\xi
 =\rho_0c\omega\xi_m\sin(kx-\omega t)
 =\rho_0c\,v_{\rm part}.
\]
这里 $p'$ 是总压强相对平衡值的微小变化，可以为正或负；$\rho_0$ 为平衡密度。对于这列右行波，声压与质元速度同相。截面上每单位面积传递的瞬时功率为 $p'v_{\rm part}$，取一个周期的平均，得到
\[
 I=\langle p'v_{\rm part}\rangle
 =\frac12\rho_0c\omega^2\xi_m^2
 =\frac{p_m^2}{2\rho_0c}
 =\frac{p_{\rm rms}^2}{\rho_0c}.
\]
因正弦平方的周期平均为 $1/2$，有效声压 $p_{\rm rms}=p_m/\sqrt2$。所以声强与振幅的平方成正比，而不是与振幅本身成正比。上述声压与声强关系适用于单向平面行波；在驻波中声压和速度的时空关系不同，不能照搬。

\subsubsection{先算强度，再换成分贝}
球面面积随距离平方增加。在无吸收且各向同性辐射的条件下，每个球面通过的总功率都为 $P$，所以 $4\pi r^2I=P$。这不是能量丢失，而是同一功率分散到更大的面积。距离从 $r_1$ 变成 $r_2$ 时，
\[
 L_I(r_2)-L_I(r_1)
 =10\log_{10}\frac{I(r_2)}{I(r_1)}
 =-20\log_{10}\frac{r_2}{r_1}.
\]
距离加倍，声强减为四分之一，声强级下降约 $6.02\,\mathrm{dB}$。这与“两个独立相同声源使声强级增加约 $3.01\,\mathrm{dB}$”并不矛盾，前者改变传播面积，后者改变总入射强度。

\Needspace{16\baselineskip}\begin{origexamplebox}{不同距离、不同声强级的两个独立声源}\label{ch12:ex:teaching-sound-distance}
两个可视为点声源的机器各向同性发声。在距机器甲 $2\,\mathrm m$ 处，它单独产生 $80\,\mathrm{dB}$ 的声强级；机器乙在观察点单独产生 $70\,\mathrm{dB}$。观察点距甲 $4\,\mathrm m$。忽略反射和吸收，二者互不相干。求同时工作时观察点的总声强级。

\ThQSoundDistance
\end{origexamplebox}
\Needspace{6\baselineskip}\begin{xesolution}
\ThASoundDistance

先把甲的测量位置换到观察点。距离加倍，强度减为四分之一，故
\[
 I_A=\frac14I_0 10^{80/10}=2.5\times10^{-5}\,\mathrm{W/m^2}.
\]
乙在该点的强度为 $I_B=I_0 10^{70/10}=10^{-5}\,\mathrm{W/m^2}$。互不相干时，平均交叉项为零，才有 $I=I_A+I_B$，所以
\[
 L=10\log_{10}(3.5\times10^7)\approx75.44\,\mathrm{dB}.
\]
结果略高于甲单独在该点的 $73.98\,\mathrm{dB}$，不是两个分贝数的算术和。若乙关闭，保留 $I_A$ 即回到甲单独的结果。
\end{xesolution}

若两列同频同向平面声波保持固定相位差 $\delta$，应先叠加声压，再平方取平均，得到
\[
 I=I_1+I_2+2\sqrt{I_1I_2}\cos\delta.
\]
等强同相时局部强度为 $4I_1$，反相时可以相消；若相位不断随机变化，平均 $\cos\delta=0$，才恢复非相干相加。驻波、不同传播方向和复杂声场还应计算矢量能流，不能把这条同向行波公式当作任意声场的通式。

\section{波包、色散与传播速度}
\subsection{波包、相速度与群速度}

把相近波数的简谐波叠加可形成局域波包。单个相位点按
\begin{equation}
 v_p=\frac\omega k
\end{equation}
传播，而窄带包络按
\begin{equation}
 v_g=\frac{\dd\omega}{\dd k}
\end{equation}
传播。无色散介质中 $\omega=ck$，二者相等；有色散时二者通常不同。若波包带宽太大或介质强吸收，单一群速度未必能完整描述包络变形。

\begin{center}\begin{minipage}{\linewidth}\centering
\begin{tikzpicture}
\begin{axis}[width=.82\linewidth,height=4.9cm,axis lines=middle,
  xlabel={$x$},ylabel={$y$},ylabel style={at={(axis description cs:-.025,.5)},anchor=east},xmin=-6,xmax=6,ymin=-1.5,ymax=1.7,ticks=none,
  every axis label/.style={fill=none}]
  \addplot[DeepBlue,very thick,domain=-6:6,samples=650]{exp(-x^2/8)*cos(deg(5*x))};
  \addplot[TextGray,dashed,domain=-6:6,samples=220]{exp(-x^2/8)};
  \addplot[TextGray,dashed,domain=-6:6,samples=220]{-exp(-x^2/8)};
  \draw[vector] (axis cs:1.2,1.35)--(axis cs:3.5,1.35) node[midway,above,clearlabel]{$v_g$};
  \draw[component] (axis cs:2.0,-.85)--(axis cs:3.4,-.85) node[midway,below,clearlabel]{$v_p$};
\end{axis}
\end{tikzpicture}
\figcap{窄带波包。载波相位在包络内移动，包络和相位的速度可以不同。}\label{ch12:fig:wave-packet}
\end{minipage}\end{center}
图~\ref{ch12:fig:wave-packet} 中虚线包络的位置随 $v_g$ 移动，而某个波峰按 $v_p$ 穿过包络。

\Needspace{16\baselineskip}\begin{origexamplebox}{深水重力波的相速度与群速度}
\label{ch12:ex:06}
深水表面重力波满足 $kh\gg1$（$h$ 为水深），并处在重力远强于表面张力的波长范围。忽略黏性，色散关系 $\omega^2=gk$。求相速度和群速度的关系，并判断一个窄带波包中能量包络相对波峰怎样运动。

\ThQDeepWave
\end{origexamplebox}
\Needspace{6\baselineskip}\begin{xesolution}
\ThADeepWave

取正频率支 $\omega=\sqrt{gk}$，
\[
 v_p=\frac\omega k=\sqrt{\frac gk},
 \qquad
 v_g=\frac{\dd\omega}{\dd k}=\frac12\sqrt{\frac gk}.
\]
所以
\[
 \boxed{v_g=\frac12v_p}.
\]
波峰以 $v_p$ 前进，会从波包后部生成、穿过包络并在前部消失；能量包络只以一半速度传播。
量纲上 $g/k$ 为 $L^2T^{-2}$，开方后是速度。若水深有限，色散关系改变，比例不再恒为 $1/2$。
\end{xesolution}

\subsubsection{从两个相近波数推导包络速度}
相速度跟踪一个确定相位，例如某一波峰。对 $kx-\omega t=\text{常量}$ 求时间导数，$k\dd x/\dd t-\omega=0$，所以 $v_p=\omega/k$。这并未回答一组有限长度波动的整体轮廓怎样移动。

先把问题简化为两个等振幅、同向的近频率行波：
\[
 y=A\cos(k_1x-\omega_1t)+A\cos(k_2x-\omega_2t).
\]
定义 $\bar k=(k_1+k_2)/2$、$\bar\omega=(\omega_1+\omega_2)/2$、$\Delta k=k_1-k_2$、$\Delta\omega=\omega_1-\omega_2$。和差化积给出
\[
 y=2A\cos\frac{\Delta kx-\Delta\omega t}{2}
       \cos(\bar kx-\bar\omega t).
\]
第二项快速振荡，第一项缓慢调节振幅。追踪一个包络极大，就要保持 $\Delta kx-\Delta\omega t$ 不变，于是
\[
 v_{\rm env}=\frac{\Delta\omega}{\Delta k}.
\]
当两波数都趋近中心波数 $k_0$，差商趋于色散曲线的斜率，才得到 $v_g=\omega'(k_0)$。因此相速度对应从原点到色散曲线上一点的割线斜率，群速度对应该点的切线斜率；只有线性无截距的 $\omega=ck$ 才在所有波数上都相同。

两个单色波的包络在整个空间周期重复，还不是真正局域的波包。若叠加一段窄波数范围，在 $k_0$ 附近展开
\[
 \omega(k)=\omega_0+v_g(k-k_0)+\tfrac12\omega''(k_0)(k-k_0)^2+\cdots,
\]
只保留常数项和一次项时，各分量的慢相位共同依赖 $x-v_gt$，局域包络近似刚性平移。二次项使不同波数获得不同额外相位，波包会展宽或变形；忽略它要求 $|\omega''(k_0)|(\Delta k)^2t\ll1$。图~\ref{ch12:fig:wave-packet} 中虚线轮廓对应这种窄带包络，内部细密波峰对应载波，不能把两个箭头解释为同一个质元的两种速度。

\Needspace{16\baselineskip}\begin{origexamplebox}{带截止频率的波包}\label{ch12:ex:teaching-cutoff}
某无耗一维线性系统的色散关系为 $\omega^2=\omega_c^2+c_0^2k^2$，其中 $\omega_c>0$、$c_0>0$ 均为常量。取沿 $+x$ 传播的窄带波包，中心频率为 $2\omega_c$。求相速度、群速度及二者之积，并讨论接近截止频率时的变化。

\ThQCutoff
\end{origexamplebox}
\Needspace{6\baselineskip}\begin{xesolution}
\ThACutoff

取 $k>0$、$\omega>0$，在中心频率处有 $k_0=\sqrt3\omega_c/c_0$。先由相位定义计算
\[
 v_p=\frac{2\omega_c}{\sqrt3\omega_c/c_0}=\frac{2c_0}{\sqrt3}.
\]
对色散关系求导，$2\omega\,\dd\omega/\dd k=2c_0^2k$，因而
\[
 v_g=\frac{c_0^2k_0}{\omega_0}=\frac{\sqrt3c_0}{2},
 \qquad v_pv_g=c_0^2.
\]
中心频率从上方趋于 $\omega_c$ 时，$k_0\to0$，$v_p\to\infty$ 而 $v_g\to0$。发散的是等相位点的形式速度，不能据此断言信息以无限速度传递。对越来越靠近截止的波包，还须使带宽足够窄，才能继续使用单一群速度近似。
\end{xesolution}

\section{边界条件、驻波与反射透射}
\subsection{管中驻波}

开口端附近声压扰动近似为零，是声压节点、空气位移腹；刚性闭口端空气位移为零，是位移节点、声压腹。
两端开口或两端闭口的理想管长 $L$ 时
\begin{equation}
 f_n=\frac{nc}{2L},\qquad n=1,2,3,\ldots.
\end{equation}
一端闭口、一端开口时只允许奇次谐波：
\begin{equation}
 f_m=\frac{(2m-1)c}{4L},\qquad m=1,2,3,\ldots.
\end{equation}
真实开口的位移腹略在管口之外，可用端部修正增加有效长度。

\begin{center}\begin{minipage}{\linewidth}\centering
\begin{tikzpicture}[x=1.1cm,y=.8cm,>=Latex]
  \draw[very thick] (0,-1.2)--(0,1.2);
  \draw[thick] (0,1.2)--(6,1.2);
  \draw[thick] (0,-1.2)--(6,-1.2);
  \draw[DeepBlue,very thick,domain=0:6,samples=180] plot (\x,{.85*sin(45*\x)});
  \draw[guide] (0,0)--(6,0);
  \node[above] at (0,1.4) {闭口};
  \node[above] at (6,1.4) {开口};
  \draw[<->,TextGray] (0,-1.7)--(6,-1.7) node[midway,below]{$L=3\lambda/4$};
\end{tikzpicture}
\figcap{一端闭口一端开口管的第二个允许位移模态。曲线表示纵向位移随位置的分布，并非空气真的沿曲线弯曲；对应 $L=3\lambda/4$。}\label{ch12:fig:pipe-mode}
\end{minipage}\end{center}
图~\ref{ch12:fig:pipe-mode} 的闭口是位移节点、开口是位移腹，因而管长只能容纳奇数个四分之一波长。

\Needspace{16\baselineskip}\begin{origexamplebox}{闭管的共振频率判别}
\label{ch12:ex:07}
一端闭口、一端开口的空气柱长 $L=0.85\,\mathrm m$，声速 $c=340\,\mathrm{m/s}$。求最低三个共振频率，并判断 $300\,\mathrm{Hz}$ 是否为共振频率。

\ThQClosedPipe
\end{origexamplebox}
\Needspace{6\baselineskip}\begin{xesolution}
\ThAClosedPipe

基频
\[
 f_1=\frac c{4L}=\frac{340}{4\times0.85}=100\,\mathrm{Hz}.
\]
闭管只取奇数倍，故最低三项为
\[
 \boxed{100\,\mathrm{Hz},\quad300\,\mathrm{Hz},\quad500\,\mathrm{Hz}}.
\]
因此 $300\,\mathrm{Hz}$ 是第二个允许模态，而不是“二次谐波”；按全部整数倍编号时它是第三谐波。
频率公式中 $c/L$ 的量纲为 $\mathrm{s^{-1}}$。
\end{xesolution}

\subsubsection{先决定边界约束的是位移还是声压}
令 $x$ 从闭口端指向开口端，空气纵向位移为 $\xi(x,t)$。闭口处空气不能穿过刚性壁，故 $\xi(0,t)=0$；开口处与外界相通，理想近似下声压扰动 $p'(L,t)=0$，并不是绝对压强为零。
对一个两端原本相距 $\dd x$ 的空气微元，实际长度增量为 $\xi_x\dd x$。质量守恒给出密度变化 $\delta\rho=-\rho_0\xi_x$，结合 $p'=c^2\delta\rho$，得到
\[
 p'=-\rho_0c^2\xi_x.
\]
因此开口的边界条件是 $\xi_x(L,t)=0$，与闭口的 $\xi=0$ 不同。对简谐驻波写 $\xi=X(x)\cos\omega t$，波动方程化为 $X''+k^2X=0$，通解为 $C\sin kx+D\cos kx$。闭口条件消去 $D$，开口条件给
\[
 X'(L)=Ck\cos kL=0,
 \quad kL=\frac{(2m-1)\pi}{2}.
\]
再用 $\omega=ck$，就得到闭管只允许奇次谐波的频率公式。图~\ref{ch12:fig:pipe-mode} 选的是 $m=2$：$kL=3\pi/2$，故 $L=3\lambda/4$。图中的曲线是“各位置的纵向位移值”，不是空气微元走过的弯曲路径。

同一驻波的声压空间因子正比于 $\cos kx$，而位移因子为 $\sin kx$。所以位移波节处为声压波腹，位移波腹处为声压波节；不能在同一位置把两者都标成波节。两端闭口给 $\sin kL=0$；两端开口可从 $X=C\cos kx$ 出发，也得到 $kL=n\pi$。零频率的整体平移不是这里讨论的声学共振模态。

\Needspace{16\baselineskip}\begin{origexamplebox}{由连续共振长度测声速与端部修正}\label{ch12:ex:teaching-resonance-length}
频率 $f=500\,\mathrm{Hz}$ 的音叉放在竖直细管口附近，管下部水面形成近似刚性闭端。调节水面，测得连续两个共振空气柱长度为 $L_1=0.160\,\mathrm m$、$L_2=0.500\,\mathrm m$。设两次开口端修正量相同，记为 $e$；且第一次为基模。求声速和 $e$。

\ThQResonanceLength
\end{origexamplebox}
\Needspace{6\baselineskip}\begin{xesolution}
\ThAResonanceLength

开口位移腹稍在管外，所以有效长度是 $L+e$。第一、第二允许模态分别满足
\[
 L_1+e=\frac\lambda4,\qquad L_2+e=\frac{3\lambda}4.
\]
两式相减可消去不能直接量取的 $e$：$L_2-L_1=\lambda/2$。因此
\[
 \lambda=0.680\,\mathrm m,\qquad
 c=f\lambda=340\,\mathrm{m/s},\qquad
 e=\lambda/4-L_1=0.010\,\mathrm m.
\]
修正量为正，与位移腹在管外一致。若不知道第一次是哪一级，仍能由相邻共振间隔求 $c$，却不能唯一确定 $e$；这说明“第一次为基模”不是多余条件。
\end{xesolution}

\subsection{介质边界的反射与透射}

两根张力相同而线密度不同的弦在 $x=0$ 连接。令特征阻抗 $Z_j=\sqrt{T\mu_j}$，位移振幅满足边界处位移连续和横向力连续。
用同一时间因子，写
$y_1=A_i e^{\ii(\omega t-k_1x)}+A_r e^{\ii(\omega t+k_1x)}$、
$y_2=A_t e^{\ii(\omega t-k_2x)}$，实际取实部。
$x=0$ 处位移连续要求 $A_i+A_r=A_t$，横向力连续要求
$k_1(A_i-A_r)=k_2A_t$。因 $Tk_j=\omega Z_j$，
第二式为 $Z_1(A_i-A_r)=Z_2A_t$。
将 $A_t=A_i+A_r$ 代入并移项，得到 $(Z_1-Z_2)A_i=(Z_1+Z_2)A_r$，
因此位移振幅系数为
\begin{equation}
 r=\frac{A_r}{A_i}=\frac{Z_1-Z_2}{Z_1+Z_2},
 \qquad
 t=\frac{A_t}{A_i}=\frac{2Z_1}{Z_1+Z_2}.
\end{equation}
功率系数为
\begin{equation}
 \mathcal R=r^2,
 \qquad
 \mathcal T=\frac{Z_2}{Z_1}t^2
 =\frac{4Z_1Z_2}{(Z_1+Z_2)^2},
 \qquad \mathcal R+\mathcal T=1.
\end{equation}
位移振幅可以大于入射振幅，仍不代表能量超过入射能量；功率还含阻抗因子。

\begin{center}\begin{minipage}{\linewidth}\centering
\begin{tikzpicture}[scale=.9,>=Latex]
  \draw[very thick,DeepBlue] (-4,0)--(0,0);
  \draw[very thick,TextGray] (0,0)--(4,0);
  \draw[vector] (-3,1.0)--(-.4,1.0) node[midway,above]{$A_i$};
  \draw[vector] (-.4,-.9)--(-3,-.9) node[midway,below]{$A_r$};
  \draw[vector] (.35,1.0)--(3,1.0) node[midway,above]{$A_t$};
  \draw[guide] (0,-1.4)--(0,2.2);
  \node[clearlabel] at (-2,2.15) {$Z_1$};
  \node[clearlabel] at (2,2.15) {$Z_2$};
  \node[point,label={[clearlabel]below right:连接点}] at (0,0) {};
\end{tikzpicture}
\figcap{弦连接点处的入射、反射与透射方向。连接点位移和横向力必须同时连续。}\label{ch12:fig:boundary}
\end{minipage}\end{center}
按图~\ref{ch12:fig:boundary} 的传播方向，反射波的波数分量变号，这正是横向力边界式中出现 $A_i-A_r$ 的原因。

\Needspace{16\baselineskip}\begin{origexamplebox}{轻弦到重弦的反射相位与能量}
\label{ch12:ex:08}
同一张力下两弦线密度满足 $\mu_2=4\mu_1$，波从弦1入射到弦2。求位移振幅反射、透射系数和能量反射、透射率，并判断反射是否反相。

\ThQStringInterface
\end{origexamplebox}
\Needspace{6\baselineskip}\begin{xesolution}
\ThAStringInterface

$Z\propto\sqrt\mu$，所以 $Z_2=2Z_1$。于是
\[
 r=\frac{Z_1-Z_2}{Z_1+Z_2}=-\frac13,
 \qquad
 t=\frac{2Z_1}{Z_1+Z_2}=\frac23.
\]
$r<0$ 表示位移反射波反相。能量系数
\[
 \boxed{\mathcal R=r^2=\frac19},
 \qquad
 \boxed{\mathcal T=\frac{Z_2}{Z_1}t^2
 =2\times\frac49=\frac89}.
\]
二者之和为1，验证无耗边界的能量守恒。若 $Z_2\to\infty$，$r\to-1$，回到固定端反射；若 $Z_2=Z_1$，反射消失。
\end{xesolution}

\section{二维与三维弹性波（选读）}
\subsection{张紧薄膜的二维波动（选读）}

膜的两个空间自变量为 $x,y$，横向位移只有一个分量 $z(x,y,t)$。
设参考面密度为 $\mu_s$，各向同性的单位长度张力为 $\sigma_s$。
$\sigma_s$ 的单位为 $\mathrm{N/m}$，不同于弦张力 $T$ 的 $\mathrm N$，
也不同于三维应力的 $\mathrm{N/m^2}$。这一模型取预张力恒定、小斜率，忽略弯曲刚度。

取膜上一块区域 $D$，从 $+z$ 方向看，其边界按逆时针方向绕行。
图\ref{ch12:fig:membrane}中的边界线元与向上单位法向到一阶为
\[
 \dd\vect r=(\dd x,\dd y,z_x\dd x+z_y\dd y),\qquad
 \vect n\simeq(-z_x,-z_y,1).
\]
周围膜对这块膜的力沿膜面、垂直边界，并指向区域外，所以
\[
 \dd\vect F=\sigma_s\,\dd\vect r\times\vect n
 =\sigma_s(\dd y,-\dd x,z_x\dd y-z_y\dd x)+O(|\nabla z|^2).
\]
检验方向：在区域右边界，边界线元沿 $+y$，而 $\dd\vect r\times\vect n$ 的水平分量沿 $+x$，
确实朝外，不能把叉积次序颠倒。环绕闭曲线时 $\oint\dd x=\oint\dd y=0$，
故一阶水平合力消失。竖直合力用格林公式变成
\begin{align}
 F_z&=\sigma_s\oint_{\partial D}(z_x\dd y-z_y\dd x)\\
 &=\sigma_s\iint_D(z_{xx}+z_{yy})\dd x\dd y.
 \label{ch12:eq:membrane-force}
\end{align}
这里所用公式为 $\oint(P\dd x+Q\dd y)=\iint(Q_x-P_y)\dd x\dd y$，
取 $P=-z_y,Q=z_x$。将任意区域的质量乘以加速度，与合力相等，就得到局部方程
\begin{equation}
 z_{tt}=\frac{\sigma_s}{\mu_s}(z_{xx}+z_{yy}),\qquad
 c_{\rm membrane}=\sqrt{\frac{\sigma_s}{\mu_s}}.
 \label{ch12:eq:membrane-wave}
\end{equation}

\begin{center}\begin{minipage}{\linewidth}\centering
\begin{tikzpicture}[x={(1cm,0cm)},y={(-.45cm,-.80cm)},z={(0cm,1cm)},>=Latex]
 \draw[axis] (0,0,0)--(5.1,0,0) node[right] {$x$};
 \draw[axis] (0,0,0)--(0,4.3,0) node[below left] {$y$};
 \draw[axis] (0,0,0)--(0,0,2.4) node[above] {$z$};
 \foreach \q in {0,1,2,3}{
  \draw[TextGray!45,domain=0:4,samples=45,variable=\s]
    plot({\s},{\q},{.12*\s+.07*\q+.045*\s*\q});
 }
 \foreach \q in {0,1,2,3,4}{
  \draw[TextGray!45,domain=0:3,samples=40,variable=\s]
    plot({\q},{\s},{.12*\q+.07*\s+.045*\q*\s});
 }
 \draw[DeepBlue,thick] (1,1,.235)--(3,1,.565)--(3,2,.77)--(1,2,.35)--cycle;
 \draw[vector] (3,1.25,.61625)--(3,1.8,.729);
 \coordinate(P) at(3,1.5,.6675);
 \draw[vector] (P)--++({1.25*(1+.205*.205)},{-1.25*.1875*.205},{1.25*.1875});
 \draw[vector,TextGray] (P)--++(-.1875,-.205,1);
 \node[right=4pt] at (4.4,1.5,.93) {$\dd\vect F$};
 \node[above=4pt] at (2.8125,1.295,1.75) {$\vect n$};
 \draw[thin,TextGray] (3.04,1.78,.733)--(3.9,3.2,.50);
 \node[below right=3pt] at (3.9,3.2,.50) {$\dd\vect r$};
 \draw[thin,TextGray] (1.5,1.55,.393125)--(.1,2.8,.45);
 \node[left=4pt] at (.1,2.8,.45) {$D$};
\end{tikzpicture}
\figcap{右手坐标系中的膜面及有向边界。曲面网格、边界、切线、法向和外向膜力均由同一曲面 $z=0.12x+0.07y+0.045xy$ 生成。}
\label{ch12:fig:membrane}
\end{minipage}\end{center}

面积元由两个切向量的叉积求得：
$\dd S=\sqrt{1+z_x^2+z_y^2}\dd x\dd y$。
在小斜率近似下，面积增加为
$\dd S-\dd x\dd y\simeq(z_x^2+z_y^2)\dd x\dd y/2$。
恒定单位长度张力所对应的势能增加等于 $\sigma_s$ 乘以面积增加，
故单位参考面积的能量为
\begin{equation}
 \mathcal E_s=\frac12\mu_s z_t^2+\frac12\sigma_s(z_x^2+z_y^2).
 \label{ch12:eq:membrane-energy}
\end{equation}
这里常数平面面积的能量可以选为零，并不意味着可以省略形变面积的变化。

无限均匀膜上的平面简谐波可写为
$z=\operatorname{Re}\{A e^{\ii(\omega t-k_xx-k_yy)}\}$，
代入波动方程给出
$\omega^2=c_{\rm membrane}^2(k_x^2+k_y^2)$。
频率决定 $|\vect k|$，$k_x:k_y$ 及其符号决定传播方向，
复数 $A$ 的模和幅角分别决定振幅与初相位；这些是不同的自由度。

矩形膜的边长为 $a,b$，四边固定时试取
$z=X(x)Y(y)\cos(\omega t+\phi)$。边界要求 $X(0)=X(a)=Y(0)=Y(b)=0$，
因此 $X\propto\sin(m\pi x/a)$、$Y\propto\sin(n\pi y/b)$，代回方程得到
\begin{equation}
 z_{mn}\propto\sin\frac{m\pi x}{a}\sin\frac{n\pi y}{b}\cos(\omega_{mn}t+\phi),
 \quad
 \omega_{mn}=\pi\sqrt{\frac{\sigma_s}{\mu_s}}
 \sqrt{\frac{m^2}{a^2}+\frac{n^2}{b^2}}.
 \label{ch12:eq:membrane-modes}
\end{equation}
$m,n$ 均为正整数。不同的空间方向各带一个模态指标，若 $a=b$，
交换 $m,n$ 不改变频率，可能出现简并。圆膜需用极坐标和贝塞尔函数
（Bessel functions）；允许频率由形状与边界共同决定，而不是只由材料决定。

\subsection{三维弹性体与纵横体波（选读）}

\subsubsection{从应力散度到运动方程}

对任意物质体积 $D$，外向面力合力的第 $i$ 分量为
\[
 F_i=\oint_{\partial D}\sigma_{ij}n_j\dd S
 =\iiint_D\frac{\partial\sigma_{ij}}{\partial x_j}\dd V.
\]
第二步对每个分量使用高斯定理。若体积力密度为 $f_i$，单位 $\mathrm{N/m^3}$，
牛顿第二定律给出
\begin{equation}
 \rho u_{i,tt}=\partial_j\sigma_{ij}+f_i.
 \label{ch12:eq:stress-div}
\end{equation}
小形变时，在这一阶精度内可将参考坐标导数与空间坐标导数等同。
均匀各向同性介质中，代入式\eqref{ch12:eq:lame}，逐项求导：
\[
 \partial_j\sigma_{ij}
 =\lambda_e\partial_i(\partial_k u_k)
   +G\partial_j(\partial_i u_j+\partial_j u_i)
 =(\lambda_e+G)\partial_i(\nabla\cdot\vect u)+G\nabla^2u_i.
\]
无体积力时得到纳维—柯西方程（Navier--Cauchy equation）
\begin{equation}
 \rho\vect u_{tt}=(\lambda_e+G)\nabla(\nabla\cdot\vect u)+G\nabla^2\vect u
 =(\lambda_e+2G)\nabla(\nabla\cdot\vect u)-G\nabla\times(\nabla\times\vect u).
 \label{ch12:eq:navier}
\end{equation}
末式用到 $\nabla^2\vect u=\nabla(\nabla\cdot\vect u)-\nabla\times(\nabla\times\vect u)$。
为看清耦合，第一分量完全展开为
\[
 \rho u_{x,tt}=(\lambda_e+2G)u_{x,xx}
 +G(u_{x,yy}+u_{x,zz})
 +(\lambda_e+G)(u_{y,xy}+u_{z,xz});
\]
另外两式由 $x,y,z$ 循环替换得到。横向二阶导数不能从展开式中漏去。

\subsubsection{三维能量密度}

从零应变沿直线路径增加到 $\boldsymbol\varepsilon$，应力随路径参数也线性增加，
所以单位体积弹性能为
\begin{equation}
 w_{\rm el}=\frac12\sum_{i,j}\sigma_{ij}\varepsilon_{ij}
 =\frac{\lambda_e}{2}(\operatorname{tr}\boldsymbol\varepsilon)^2
   +G\sum_{i,j}\varepsilon_{ij}^2.
 \label{ch12:eq:elastic-energy}
\end{equation}
双重求和中的 $\varepsilon_{xy}^2$ 与 $\varepsilon_{yx}^2$ 都要计入；
纯剪切时便得到 $w_{\rm el}=G\gamma_{xy}^2/2$，与 $\tau=G\gamma_{xy}$ 的做功一致。
加上动能密度 $\rho|\vect u_t|^2/2$ 就得到总机械能密度。
把应变分成体积部分和无迹部分，还可写成
\[
 w_{\rm el}=\frac K2(\operatorname{tr}\boldsymbol\varepsilon)^2
 +G\left|\boldsymbol\varepsilon-\frac13
 (\operatorname{tr}\boldsymbol\varepsilon)\mathsf I\right|^2,
\]
清楚显示 $K>0,G>0$ 是稳定性条件。

\subsubsection{平面波的三种偏振与不同波速}

取平面波 $\vect u=\vect A e^{\ii(\vect k\cdot\vect r-\omega t)}$ 的实部。
代入式\eqref{ch12:eq:navier}得到
\[
 \rho\omega^2\vect A=(\lambda_e+G)\vect k(\vect k\cdot\vect A)+Gk^2\vect A.
\]
将振幅唯一分成
$\vect A_P=(\vect A\cdot\hat{\vect k})\hat{\vect k}$ 与
$\vect A_S=\vect A-\vect A_P$。
前者与 $\vect k$ 平行，满足 $\vect k\times\vect A_P=0$；
后者与 $\vect k$ 垂直，满足 $\vect k\cdot\vect A_S=0$。
分别投影原式：
\begin{equation}
 c_P^2=\frac{\omega_P^2}{k^2}=\frac{\lambda_e+2G}{\rho}
      =\frac{K+4G/3}{\rho},\qquad
 c_S^2=\frac{\omega_S^2}{k^2}=\frac G\rho.
 \label{ch12:eq:body-speeds}
\end{equation}
一个传播方向对应一种纵偏振、两种互相独立的横偏振，如图\ref{ch12:fig:polarization}。
对一般波场，在适当边界条件下可作无旋与无散的赫尔姆霍兹分解；
对运动方程分别取散度与旋度，也得到
$(\nabla\cdot\vect u)_{tt}=c_P^2\nabla^2(\nabla\cdot\vect u)$、
$(\nabla\times\vect u)_{tt}=c_S^2\nabla^2(\nabla\times\vect u)$。
因此上述结论不只是某一个特殊方向的平面波巧合。

\begin{center}\begin{minipage}{\linewidth}\centering
\begin{tikzpicture}[>=Latex]
 \foreach \q/\name in {0/纵偏振,4.5/横偏振一,9/横偏振二}{
 \begin{scope}[shift={(\q,0)}]
  \draw[axis] (0,0)--(2.7,0) node[right] {$\vect k$};
  \node at (1.2,-1) {\name};
 \end{scope}}
 \draw[<->,DeepBlue,very thick] (.6,.5)--(1.9,.5) node[midway,above] {$\vect A_P$};
 \draw[<->,DeepBlue,very thick] (5.7,-.6)--(5.7,.9) node[above] {$\vect A_{S1}$};
 \draw[DeepBlue,thick] (10.2,.4) circle(.16);
 \fill[DeepBlue] (10.2,.4) circle(1.6pt);
 \node[above] at (10.2,.75) {$\vect A_{S2}$：出入纸面};
\end{tikzpicture}
\figcap{沿同一方向传播的三种独立体波偏振。箭头与点圆符号表示材料点振动方向，不是另外的传播方向。}
\label{ch12:fig:polarization}
\end{minipage}\end{center}

细杆纵波的 $\sqrt{Y/\rho}$ 不能替换为体纵波的 $c_P$：
前者侧面可以横向收缩，后者平面纵波伴随的横向约束使有效纵向模量为 $\lambda_e+2G$。
稳定各向同性材料满足 $c_P>c_S$。地震中先到达的压缩体波称P波（primary wave），
较慢的剪切体波称S波（secondary wave）。液态外核不能传递剪切体波；
这处边界称古登堡不连续面（Gutenberg discontinuity），核幔边界深约 $2900\,\mathrm{km}$，波的折射、反射和模式转换为认识地球内部提供依据。
S波在边界可转换为能穿过液体的压缩波，不能将“液体不传S波”说成“所有S波能量均原样反射”。


\section{弹性体与波动综合问题}

\Needspace{16\baselineskip}\begin{origexamplebox}{绕端点旋转的均匀弹性绳}
\label{ch12:ex:09}
一根自然长度为 $L_0$、质量为 $M$ 的均匀弹性绳，以角速度 $\omega$ 绕通过一端 $O$、且与绳垂直的轴匀速旋转，
忽略重力。绳的轴向刚度为 $EA=kL_0$，即同一根绳在均匀受拉时的整体劲度系数为 $k$。
求稳定旋转时的绳长、弹性势能和动能。


\begin{center}\begin{minipage}{\linewidth}\centering
\begin{tikzpicture}[scale=.92,>=Latex]
  \draw[axis] (0,-2.0)--(0,2.1);
  \fill (0,0) circle (2.5pt) node[below left] {$O$};
  \draw[curve,very thick] (0,0)--(5.3,0) node[right] {自由端};
  \draw[->,DeepBlue,thick] (-.75,1.12) arc[start angle=150,end angle=-145,x radius=.75,y radius=.32];
  \node[left] at (-.72,.75) {$\omega$};
  \draw[<->,guide] (0,-.65)--(5.3,-.65) node[midway,below] {$L$};
  \node[above] at (2.65,0) {$L_0,\ M,\ k$};
\end{tikzpicture}
\figcap{弹性绳绕过端点 $O$ 的垂直轴旋转。$L_0$ 为自然长度，$L$ 为待求的旋转绳长。}
\label{ch12:fig:rotating-rope}
\end{minipage}\end{center}
\end{origexamplebox}
\Needspace{6\baselineskip}\begin{xesolution}
图\ref{ch12:fig:rotating-rope}中，整根绳沿径向伸展，但各处的伸长不相同。


以自然构形中的坐标 $s\in[0,L_0]$ 标记绳元，令其稳定旋转时的半径为 $r(s)$。
自然线密度为 $\lambda_0=M/L_0$，局部小形变张力为
\[
 T(s)=EA\left(\frac{\dd r}{\dd s}-1\right).
\]
图\ref{ch12:fig:rotating-element}把自然微元与其实际半径分开标出。
\begin{center}\begin{minipage}{\linewidth}\centering
\begin{tikzpicture}[>=Latex]
 \draw[guide] (0,0)--(7,0); \fill(0,0)circle(2pt) node[left] {$O$};
 \draw[very thick,DeepBlue] (3,0)--(4,0);
 \draw[vector] (3,0)--(1.7,0) node[midway,above] {$T(s)$};
 \draw[vector] (4,0)--(5.2,0) node[midway,above] {$T(s+\dd s)$};
 \draw[<->] (0,-.6)--(3,-.6) node[midway,below] {$r(s)$};
 \node at (3.5,1.0) {$\dd m=(M/L_0)\dd s$};
 \draw[vector] (3.5,-1.35)--(2.2,-1.35) node[left] {$\omega^2r$};
\end{tikzpicture}
\figcap{旋转绳微元的径向受力。内侧张力较大，两端张力的差产生向心加速度。}
\label{ch12:fig:rotating-element}
\end{minipage}\end{center}
长度为 $\dd s$ 的绳元所需向心力为 $\lambda_0\omega^2r\,\dd s$；由相邻截面的张力差得到
\[
 \frac{\dd T}{\dd s}=-\lambda_0\omega^2r.
\]
代入 $EA=kL_0$，并定义
\[
 \alpha=\omega\sqrt{\frac Mk},
 \qquad q=\frac\alpha{L_0},
\]
便有
\[
 r''+q^2r=0.
\]
边界条件是转轴端 $r(0)=0$，自由端无张力 $T(L_0)=0$，即 $r'(L_0)=1$。因此
\[
 r(s)=\frac{\sin(qs)}{q\cos\alpha},
 \qquad
 \boxed{L=r(L_0)=L_0\frac{\tan\alpha}{\alpha}}.
\]
满足正伸长且无内部折返的主解支为 $0\le\alpha<\pi/2$，但线性小应变还要求最大应变 $\sec\alpha-1\ll1$；当 $\alpha\to\pi/2$ 时线性弹性模型预言长度发散，实际上材料会先离开线性范围。

弹性势能必须按非均匀应变积分：
\begin{align*}
 U
 &=\frac{EA}{2}\int_0^{L_0}(r'-1)^2\,\dd s\\
 &=\frac{kL_0^2}{2}\int_0^1\left(\frac{\cos(\alpha z)}{\cos\alpha}-1\right)^2\dd z\\
 &=\boxed{\frac{kL_0^2}{2}
 \left(1+\frac{1}{2\cos^2\alpha}
 -\frac{3\tan\alpha}{2\alpha}\right)}.
\end{align*}
动能为各绳元转动动能之和：
\begin{align*}
 K
 &=\frac12\lambda_0\omega^2\int_0^{L_0}r^2(s)\,\dd s\\
 &=\frac{kL_0^2}{2\cos^2\alpha}\int_0^1\sin^2(\alpha z)\dd z\\
 &=\boxed{\frac{kL_0^2}{4\cos^2\alpha}
 \left(1-\frac{\sin2\alpha}{2\alpha}\right)}.
\end{align*}
上述积分使用
$\int_0^1\cos(\alpha z)\dd z=\sin\alpha/\alpha$、
$\int_0^1\cos^2(\alpha z)\dd z=1/2+\sin2\alpha/(4\alpha)$；
$\sin^2$ 的积分把末项改为负号。小角速度极限下 $L=L_0(1+\alpha^2/3+\cdots)$、
$K\simeq M\omega^2L_0^2/6$，后者正是未伸长均匀细绳绕端点转动的动能，因而结果通过极限检验。
\end{xesolution}

\Needspace{16\baselineskip}\begin{origexamplebox}{弦上点质量的反射与透射}
\label{ch12:ex:10}
一根无限长弦的线密度为 $\lambda$、张力为 $T$。从 $x<0$ 远处入射行波
\[
 y_i=A\cos(\omega t-kx).
\]
在 $x=0$ 处附着一个可随弦横向运动、质量为 $m$ 的质点，弦在质点两侧材料相同。求：（1）$k$ 与 $\omega$ 的关系；
（2）位移振幅的反射、透射系数，以及能量反射率、透射率。


\begin{center}\begin{minipage}{\linewidth}\centering
\begin{tikzpicture}[>=Latex,scale=.8]
 \draw[axis] (-3.5,0)--(3.6,0) node[right] {$x$};
 \draw[axis] (0,-.5)--(0,2) node[above] {$y$};
 \draw[curve] (-3,.25) .. controls (-1.8,.25) and (-1.5,1.1) .. (0,1.1)
 .. controls (1.5,1.1) and (1.8,1.5) .. (3,1.75);
 \filldraw[fill=DeepBlue!20,draw=DeepBlue] (0,1.1) circle(.12);
 \node[above left] at (-.12,1.15) {$m$};
 \node[below left] at (0,0) {$O$};
 \draw[->] (-3,2)--(-1.5,2) node[midway,above] {入射方向};
\end{tikzpicture}
\figcap{原题装置示意：质点附着在 $x=0$ 处，并随弦横向运动。曲线仅示意弦的瞬时形状。}
\label{ch12:fig:point-mass}
\end{minipage}\end{center}
\end{origexamplebox}
\Needspace{6\baselineskip}\begin{xesolution}
先按题设的持续单频激励画出图\ref{ch12:fig:point-mass-analysis}。
\begin{center}\begin{minipage}{\linewidth}\centering
\begin{tikzpicture}[>=Latex,x=.9cm,y=.85cm]
 \draw[axis] (-4.5,0)--(4.6,0) node[right] {$x$};
 \draw[curve,domain=-4:0,samples=100] plot(\x,{.7*(cos(90*\x)+sin(90*\x))});
 \draw[curve,domain=0:4,samples=100] plot(\x,{.7*(cos(90*\x)-sin(90*\x))});
 \filldraw[fill=DeepBlue!20,draw=DeepBlue] (0,.7) circle(.15);
 \node[above right] at (.15,.85) {$m$};
 \draw[vector] (-3,1.6)--(-1.5,1.6) node[midway,above] {入射};
 \draw[vector] (-1.5,-1.7)--(-3,-1.7) node[midway,below] {反射};
 \draw[vector] (1.5,1.6)--(3,1.6) node[midway,above] {透射};
 \node at (0,-2.6) {$\beta=1,\ t=0$：实线为两侧总位移};
\end{tikzpicture}
\figcap{持续单频波在点质量处的一个瞬间。左侧入射与反射相加，右侧为透射波；两侧在质点处位移连续而斜率不同。}
\label{ch12:fig:point-mass-analysis}
\end{minipage}\end{center}

弦的波动方程给出
\[
 \boxed{k=\omega\sqrt{\frac\lambda T}},
 \qquad c=\frac\omega k=\sqrt{\frac T\lambda}.
\]
用复振幅并取时间因子 $e^{\ii\omega t}$，写成
\[
 y_-=(Ae^{-\ii kx}+Be^{\ii kx})e^{\ii\omega t},
 \qquad
 y_+=Ce^{-\ii kx}e^{\ii\omega t}.
\]
实际位移取复表达式的实部。质点处位移连续，所以 $A+B=C$。
图\ref{ch12:fig:point-mass-analysis}取 $\beta=1$，即 $B/A=(-1-\ii)/2$、$C/A=(1-\ii)/2$，
用同一波数、同一时刻绘制，不能把三列单频波画成互不相连的孤立脉冲。
右侧张力向右，其竖直分量为 $Ty_x(0^+)$；左侧张力向左，其竖直分量为 $-Ty_x(0^-)$。
质点受到的横向合力为
$T[y_x(0^+)-y_x(0^-)]$，牛顿第二定律给出
\[
 T[y_x(0^+)-y_x(0^-)]=m y_{tt}(0).
\]
在 $x=0$，$y_x(0^+)=-\ii kC$、
$y_x(0^-)=-\ii kA+\ii kB$，省略共同时间因子。
用 $C=A+B$，左端为 $-2\ii kTB$，右端为 $-m\omega^2C$，故
\[
 2\ii kT B=m\omega^2C.
\]
令
\[
 \beta=\frac{m\omega^2}{2kT}
 =\frac{m\omega}{2\sqrt{T\lambda}},
\]
则位移振幅系数为
\[
 \boxed{t\equiv\frac CA=\frac1{1+\ii\beta}},
 \qquad
 \boxed{r\equiv\frac BA=-\frac{\ii\beta}{1+\ii\beta}}.
\]
两侧阻抗相同，故能流比就是振幅模平方：
\[
 \boxed{\mathcal R=|r|^2=\frac{\beta^2}{1+\beta^2}},
 \qquad
 \boxed{\mathcal T=|t|^2=\frac1{1+\beta^2}},
 \qquad \mathcal R+\mathcal T=1.
\]
低频时质点几乎随弦运动，$\mathcal T\to1$；高频时惯性使连接点近似固定，$r\to-1$，反射发生相位反转。
\end{xesolution}

\Needspace{16\baselineskip}\begin{origexamplebox}{圆膜上的轴对称外行波}
\label{ch12:ex:11}
面密度为 $\mu$、单位长度张力为 $\sigma$ 的均匀张紧圆膜上，有一列从中心向外传播的轴对称小振幅波。
远离中心处写成
\[
 z(r,t)=A(r)\cos(\omega t-kr).
\]
求 $\omega/k$ 和振幅 $A(r)$ 的渐近形式，并用环形微元的牛顿第二定律说明这一形式为什么只在 $kr\gg1$ 时成立。

\ThQMembrane
\end{origexamplebox}
\Needspace{6\baselineskip}\begin{xesolution}
半径 $r$、宽度 $\dd r$ 的环形微元质量为 $\mu2\pi r\dd r$。
图\ref{ch12:fig:ring-element}中内、外边界张力朝相反的径向；
竖直分量分别为 $-2\pi r\sigma z_r(r)$ 与 $2\pi(r+\dd r)\sigma z_r(r+\dd r)$。
\begin{center}\begin{minipage}{\linewidth}\centering
\begin{tikzpicture}[>=Latex]
 \draw[thick] (0,0) circle(1.1); \draw[thick] (0,0) circle(1.45);
 \draw[->] (0,0)--(135:1.1) node[midway,above right] {$r$};
 \draw[vector] (1.45,0)--(2.1,0); \draw[vector] (1.1,0)--(.5,0);
 \node at (0,-2) {俯视：环内侧与外侧的力反向};
 \begin{scope}[shift={(5,0)}]
  \draw[curve,domain=0:3,samples=60] plot(\x,{.14*(\x-1)^2});
  \draw[vector] (1,0)--(.1,0) node[above] {内侧};
  \draw[vector] (2,.14)--(2.9,.392) node[above] {外侧};
  \draw[guide] (1,0)--(1,-.7) (2,.14)--(2,-.7);
  \draw[<->] (1,-.6)--(2,-.6) node[midway,below] {$\dd r$};
  \node at (1.5,-2) {径向截面：力沿局部切线};
 \end{scope}
\end{tikzpicture}
\figcap{圆膜的环形微元与径向剖面。外边界比内边界周长大，求合力时不能只对斜率作差。}
\label{ch12:fig:ring-element}
\end{minipage}\end{center}
两项相减并展开到 $\dd r$ 一阶，竖直合力为
\[
 2\pi\sigma\frac{\partial}{\partial r}
 \left(r\frac{\partial z}{\partial r}\right)\dd r.
\]
牛顿第二定律因此给出径向波动方程
\[
 z_{tt}=c^2\left(z_{rr}+\frac1r z_r\right),
 \qquad
 \boxed{c=\sqrt{\frac\sigma\mu}=\frac\omega k}.
\]
把复数形式 $z=A(r)e^{\ii(kr-\omega t)}$ 代入，并使用 $\omega=ck$，得到
\[
 A''+\left(2\ii k+\frac1r\right)A'
 +\frac{\ii k}{r}A=0.
\]
当 $kr\gg1$ 时，振幅在一个波长内变化很慢，可忽略 $A''$ 与 $A'/r$ 相对 $2\ii kA'$ 的贡献，于是
\[
 2A'+\frac Ar\simeq0,
 \qquad
 \boxed{A(r)=\frac{C}{\sqrt r}}.
\]
用 $A=C r^{-1/2}$ 回查：$A'=-A/(2r)$、$A''=3A/(4r^2)$，
被忽略项量级为 $A/r^2$，保留项量级为 $kA/r$，比值为 $O((kr)^{-1})$；
所以必须有 $kr\gg1$。同一结论也来自能流守恒：圆周长正比于 $r$，而单位长度平均能流正比于 $A^2$，故 $rA^2$ 为常量。
精确外行解是零阶汉克尔函数（Hankel function）$H_0^{(1)}(kr)$，其大宗量渐近式确为
$r^{-1/2}e^{\ii(kr-\pi/4)}$；靠近 $r=0$ 时上述渐近式发散，必须用完整的贝塞尔函数（Bessel function）并结合中心激励条件。
\end{xesolution}

\Needspace{16\baselineskip}\begin{origexamplebox}{音叉驱动弦的驻波}
\label{ch12:ex:12}
线密度为 $0.15\,\mathrm{g/cm}$ 的弦，一端连接频率为 $50\,\mathrm{Hz}$、
沿弦横向振动的音叉，另一端跨过光滑小滑轮并悬挂配重；音叉到滑轮的水平弦长为
$l=1.00\,\mathrm m$。讨论弱端位移、弦内部共振振幅明显大于叉股振幅的通常实验情形，
忽略滑轮半径相对 $l$ 的修正。

（1）共振时音叉连接点近似看作波节还是波腹？
（2）水平弦段有一、二、三个波腹时，所需配重质量分别为多少？
\begin{center}\begin{minipage}{\linewidth}\centering
\begin{tikzpicture}[>=Latex,scale=.83]
 \draw[thick] (-4.5,-.4)--(-4.05,-.4);
 \draw[thick] (-2.8,.25)--(-3.4,.25) arc(90:270:.65)--(-2.8,-1.05);
 \draw[curve,very thick] (-2.8,.25)--(3,.25) arc(90:0:.30)--(3.3,-1.7);
 \draw[thick] (3,-.05) circle(.30);
 \draw[thick,fill=DeepBlue!8] (2.9,-1.7) rectangle(3.7,-2.35);
 \node at (3.3,-2.02) {$m$};
 \draw[<->] (-2.8,.9)--(3,.9) node[midway,above] {$l$};
 \node[left] at (-3.6,-1.2) {音叉};
\end{tikzpicture}
\figcap{音叉、水平弦、光滑滑轮和配重。水平弦接在叉股末端，并与滑轮顶部相切。}
\label{ch12:fig:fork}
\end{minipage}\end{center}
\end{origexamplebox}
\Needspace{6\baselineskip}\begin{xesolution}
连接点有位移，不等于它就是波腹。设叉股位移振幅为 $a_d$，
弦的稳态空间振幅为 $U(x)$。给定位移的边界条件是
$U(0)=a_d,\ U(l)=0$，并不是自由端的 $U'(0)=0$。
由 $U''+k^2U=0$ 先满足右端边界，写 $U=C\sin[k(l-x)]$，
再用左端边界得到
\[
 U(x)=a_d\frac{\sin[k(l-x)]}{\sin(kl)}.
\]
当 $kl$ 接近 $n\pi$ 时，分母很小，内部振幅显著大于 $a_d$；
真实阻尼和音叉耦合使共振振幅有限。因端点振幅相对内部波腹很小，
共振的两端均近似波节，而不是严格声称运动的叉股位移恒为零。
图\ref{ch12:fig:fork-analysis}画清了叉股末端连接与弦在滑轮上的相切关系。

\begin{center}\begin{minipage}{\linewidth}\centering
\begin{tikzpicture}[>=Latex,scale=.83]
 \draw[thick] (-4.5,-.4)--(-4.05,-.4);
 \draw[thick] (-2.8,.25)--(-3.4,.25) arc(90:270:.65)--(-2.8,-1.05);
 \draw[curve,very thick] (-2.8,.25)--(3,.25) arc(90:0:.30)--(3.3,-1.7);
 \draw[thick] (3,-.05) circle(.30);
 \draw[thick,fill=DeepBlue!8] (2.9,-1.7) rectangle(3.7,-2.35);
 \node at (3.3,-2.02) {$m$};
 \draw[<->] (-2.8,.9)--(3,.9) node[midway,above] {$l$};
 \node[left] at (-3.6,-1.2) {音叉};
 \node[above] at (-2.8,1.55) {近似波节};
 \node[above] at (3,1.55) {近似波节};
 \foreach \n in {1,2,3}{
  \begin{scope}[shift={(-2.8,{-2.9-1.1*\n})}]
   \draw[guide] (0,0)--(5.8,0);
   \draw[curve,domain=0:5.8,samples=100] plot(\x,{.37*sin(\n*180*\x/5.8)});
   \fill(0,0)circle(1.5pt); \fill(5.8,0)circle(1.5pt);
   \node[left] at (-.4,0) {$n=\n$};
  \end{scope}}
\end{tikzpicture}
\figcap{音叉横向驱动弦及前三个共振模态。水平弦接在叉股末端；到滑轮顶部相切，再沿圆弧转向竖直。}
\label{ch12:fig:fork-analysis}
\end{minipage}\end{center}

$n$ 个波腹对应 $l=n\lambda_n/2$，故
\[
 f=\frac{n}{2l}\sqrt{\frac{T_n}{\mu}},\quad
 \mu=0.15\,\mathrm{g/cm}=0.015\,\mathrm{kg/m},\quad
 T_n=\mu\left(\frac{2lf}{n}\right)^2=\frac{150}{n^2}\,\mathrm N.
\]
配重静止时 $T_n=m_ng$，因此
\[
\begin{array}{c|ccc}
 n&1&2&3\\ \hline
 T_n/\mathrm N&150&37.5&16.667\\
 m_n/\mathrm{kg}\ (g=9.8\,\mathrm{m/s^2})&15.306&3.827&1.701\\
 m_n/\mathrm{kg}\ (g=10\,\mathrm{m/s^2})&15.0&3.75&1.667
\end{array}
\]
同一 $f,l$ 下波腹越多，波长越短、波速越小，所需张力按 $1/n^2$ 减小。
“强驱动”本身不会把给定位移端变成自由端；若叉股与弦强烈耦合，
应联立叉股动力学，而非自行换用奇数四分之一波长公式。
\end{xesolution}

\Needspace{16\baselineskip}\begin{origexamplebox}{高空绳索上的水平移动}
\label{ch12:ex:13}
两悬点 $A,B$ 的坐标为 $(-4l,0)$、$(4l,0)$，取 $y$ 轴竖直向下。
重为 $G$ 的演员初始骑独轮车、未踩动而平衡在一根轻绳上，位置为 $(0,3l)$。
他蹬轮缓慢向右移动水平距离 $3l$，初末静止。分别求终点坐标及所做的功：

（1）绳不可伸长；（2）绳为均匀的零原长胡克弹性绳。

采用轻绳、理想无滑动驱动接触，轮半径相对 $l$ 可忽略；
忽略平动、转动动能及耗散，但保留驱动力矩引起的两侧张力差。
\begin{center}\begin{minipage}{\linewidth}\centering
\begin{tikzpicture}[>=Latex,scale=.62]
 \draw[axis] (-5,0)--(5.2,0) node[right] {$x$};
 \draw[axis] (0,.5)--(0,-4.2) node[below] {$y$};
 \draw[curve,thick] (-4,0)--(0,-3)--(4,0);
 \fill (-4,0) circle(2pt) node[above] {$A(-4l,0)$};
 \fill (4,0) circle(2pt) node[above] {$B(4l,0)$};
 \filldraw[fill=DeepBlue!10] (0,-3) circle(.13);
 \node[below left] at (-.18,-3.1) {$(0,3l)$};
 \draw[<->] (-.65,0)--(-.65,-3) node[midway,left] {$3l$};
 \draw[->] (.35,-3)--(3,-3) node[midway,below] {$3l$};
\end{tikzpicture}
\figcap{两悬点间连续绳的初态。$y$ 轴竖直向下，箭头表示题设的水平位移。}
\label{ch12:fig:rope-initial}
\end{minipage}\end{center}
\end{origexamplebox}
\Needspace{6\baselineskip}\begin{xesolution}
把人和轮合起来研究。准静态平动平衡要求绳的两个张力与重力合力为零；
驱动轮能够传递力矩，所以不能另外强加 $T_L=T_R$。
令 $a=4l$，当前位置 $P=(x,y)$，两段实际绳长为
\[
 L_L=\sqrt{(a+x)^2+y^2},\qquad L_R=\sqrt{(a-x)^2+y^2}.
\]
图\ref{ch12:fig:rope-branches}分别画出两种绳在终态仍与悬点相连的形状。

\begin{center}\begin{minipage}{\linewidth}\centering
\begin{tikzpicture}[>=Latex,scale=.62]
 \foreach \sh/\yf/\name in {0/2.4/不可伸长绳,11/1.3125/连续弹性绳}{
 \begin{scope}[shift={(\sh,0)}]
  \coordinate(A)at(-4,0);\coordinate(B)at(4,0);
  \draw[guide] (A)--(0,-3)--(B);
  \draw[curve] (A)--(3,-\yf)--(B);
  \fill(A)circle(2pt)node[above]{$A$};
  \fill(B)circle(2pt)node[above]{$B$};
  \fill(0,-3)circle(2pt); \node[below] at(-.3,-3.2){$(0,3l)$};
  \fill(3,-\yf)circle(2pt); \node[right] at(3.25,-\yf-.25){$P_f$};
  \node at (0,1.2) {\name};
  \draw[<->] (0,-4.25)--(3,-4.25)node[midway,below]{$3l$};
 \end{scope}}
\end{tikzpicture}
\figcap{连续绳的初末构形。虚线为初态，实线为终态；两幅图都保留同一根绳的两个支撑段。}
\label{ch12:fig:rope-branches}
\end{minipage}\end{center}

（1）初态总长 $2\sqrt{(4l)^2+(3l)^2}=10l$，运动时
$L_L+L_R=10l$。这是焦点为 $A,B$、半长轴 $5l$、半短轴 $3l$ 的椭圆：
\[
 \frac{x^2}{25l^2}+\frac{y^2}{9l^2}=1.
\]
在 $x_f=3l$ 处取悬点下方的正根，得到
\[
 \boxed{(x_f,y_f)=\left(3l,\frac{12}{5}l\right)}.
\]
不可伸长绳不储存新的弹性能，初末动能均为零。
演员上升 $3l-12l/5=3l/5$，其输出的机械功为
$\boxed{W_1=3Gl/5}$。

（2）必须跟踪连续绳两侧的材料份额，而不是把它改成两根固定劲度的弹簧。
设整根绳的柔度为 $C_0$（劲度倒数，单位 $\mathrm{m/N}$）。
某一材料段占总材料份额 $q$，其柔度为 $qC_0$；零原长胡克定律给出
$L_L=qC_0T_L,\ L_R=(1-q)C_0T_R$。
因此无论有多少材料转移到轮的另一侧，
\begin{equation}
 \frac{L_L}{T_L}+\frac{L_R}{T_R}=C_0.
 \label{ch12:eq:rope-compliance}
\end{equation}
初态左右对称，$T_0(3/5)\times2=G$，所以
$T_0=5G/6,\ C_0=10l/T_0=12l/G$。

当前位置的水平、竖直平衡分别为
\[
 T_L\frac{a+x}{L_L}=T_R\frac{a-x}{L_R},\qquad
 T_L\frac y{L_L}+T_R\frac y{L_R}=G.
\]
令 $U=T_L/L_L,V=T_R/L_R$，这两式成为
$(a+x)U=(a-x)V,\ U+V=G/y$，解得
\[
 U=\frac{G(a-x)}{2ay},\qquad V=\frac{G(a+x)}{2ay}.
\]
代入式\eqref{ch12:eq:rope-compliance}，
\[
 C_0=\frac1U+\frac1V
 =\frac{4a^2y}{G(a^2-x^2)},\qquad
 y=\frac{GC_0}{4}\left(1-\frac{x^2}{a^2}\right)
 =3l\left(1-\frac{x^2}{16l^2}\right).
\]
这次是抛物线而不是椭圆，因总实际绳长可以改变。在 $x=3l$ 处
\[
 \boxed{(x_f,y_f)=\left(3l,\frac{21}{16}l\right)}.
\]
回查材料份额 $q=(L_L/T_L)/C_0=(a+x)/(2a)$，终态 $q=7/8$，
仍在 $0<q<1$ 内。两侧张力都为正，模型在本题位移范围内自洽。

零原长线性弹性段的储能为 $TL/2$。利用上面的 $U,V$，
\begin{align*}
 E_{\rm el}&=\frac12(T_LL_L+T_RL_R)\\
 &=\frac{G}{4ay}\left[(a-x)L_L^2+(a+x)L_R^2\right]\\
 &=\frac{G(a^2-x^2+y^2)}{2y}
 =\frac G2\left(\frac{4a^2}{GC_0}+y\right).
\end{align*}
因此 $\Delta E_{\rm el}=G(y_f-y_0)/2$，而重力势能变化为
$-G(y_f-y_0)$。由能量守恒，
\[
 \boxed{W_2=\Delta E_{\rm el}+\Delta E_g
 =\frac G2(y_0-y_f)=\frac{27}{32}Gl}.
\]
演员上升的同时弹性能减少，故所做功小于单独的重力势能增加，这并不违背守恒。

若另行改成“两根相同、材料段不交换的独立零原长弹性绳，并由外部水平力拖动”，
每根劲度 $k=G/(6l)$，合力为 $(-2kx,-2ky)$；
竖直平衡给出 $y=3l$，外部做功
$\int_0^{3l}2kx\,\dd x=3Gl/2$。
这个结果对应另一装置，不是上面的连续绳驱动模型。
若轮与绳完全无摩擦，蹬轮又不能传递驱动力矩，也不能直接沿用本题的驱动过程。
\end{xesolution}

\Needspace{16\baselineskip}\begin{origexamplebox}{门轴粘滑与木门啸叫的简化模型}
\label{ch12:ex:14}
把门轴承视为半径 $a=1.0\,\mathrm{cm}$、高 $h=1.0\,\mathrm{cm}$ 的实心圆柱。
门重 $W=100\,\mathrm N$ 为它对水平支承面的法向载荷，压力均匀，
静摩擦因数 $\mu_f=0.30$，材料剪切模量 $G_s=0.50\,\mathrm{GPa}$。
门带动顶端以恒定角速度 $\omega$ 转动。采用集中参数近似：圆柱为扭转弹簧，下端固连刚性端盘；
端盘以接触面可传递的总摩擦力矩为粘滑阈值，忽略接触区的局部部分滑移。
忽略滑动时间，求自激粘滑的重复角频率与 $\omega$ 的关系；
先保留每次释放后的剩余扭转角，再讨论近似释放到零的情形。
\begin{center}\begin{minipage}{\linewidth}\centering
\begin{tikzpicture}[>=Latex,scale=.86]
 \draw[thick] (-1.2,.1)--(-1.7,-.55)--(1.7,-.55)--(2.2,.1)--(1,.1);
 \draw[thick] (-1.2,.1)--(-1,.1);
 \draw[thick] (-1,0)--(-1,2.1) (1,0)--(1,2.1);
 \draw[thick] (0,2.1) ellipse(1 and .3);
 \draw[thick] (-1,0) arc(180:360:1 and .3);
 \draw[guide] (-1,0) arc(180:0:1 and .3);
 \draw[thick] (-1,0)--(-1,-.18) (1,0)--(1,-.18);
 \draw[thick] (-1,-.18) arc(180:360:1 and .3);
 \draw[TextGray,thin] (-.75,-.35)--(-1.8,-.9);
 \node[left,font=\small] at (-1.8,-.9) {刚性端盘};
 \draw[guide] (0,0)--(0,2.1);
 \draw[<->] (1.55,0)--(1.55,2.1) node[midway,right] {$h$};
 \draw[->] (0,2.1)--(1,2.1) node[midway,above] {$a$};
 \draw[->,DeepBlue] (-.9,2.75) arc(140:30:1.17 and .35);
 \node[above] at (0,3.15) {$\omega$};
\end{tikzpicture}
\figcap{门轴承的几何与驱动。圆柱下端连接刚性端盘，端盘与支承面构成等效摩擦接触。}
\label{ch12:fig:hinge}
\end{minipage}\end{center}
\end{origexamplebox}
\Needspace{6\baselineskip}\begin{xesolution}
图\ref{ch12:fig:hinge-analysis}用两个截面的相对径向角表示扭转 $\phi$。
轴线保持直线；表面母线的倾斜反映剪切角，而不是把轴线倾斜当作扭转。
半径 $r$ 处上下截面的相对切向位移为 $r\phi$，
剪应变 $\gamma(r)=r\phi/h$，剪应力为 $G_s r\phi/h$。
面积环 $\dd A=2\pi r\dd r$ 对轴的力矩为 $r(G_s r\phi/h)\dd A$，积分得到
\[
 \tau_{\rm el}=\frac{G_s\phi}{h}\int_0^a2\pi r^3\dd r
 =K_\phi\phi,\qquad K_\phi=\frac{G_s\pi a^4}{2h}.
\]
\begin{center}\begin{minipage}{\linewidth}\centering
\begin{tikzpicture}[>=Latex,scale=.86]
 \draw[thick] (-1.2,.1)--(-1.7,-.55)--(1.7,-.55)--(2.2,.1)--(1,.1);
 \draw[thick] (-1.2,.1)--(-1,.1);
 \draw[thick] (-1,0)--(-1,2.1) (1,0)--(1,2.1);
 \draw[thick] (0,2.1) ellipse(1 and .3);
 \draw[thick] (-1,0) arc(180:360:1 and .3);
 \draw[guide] (-1,0) arc(180:0:1 and .3);
 \draw[thick] (-1,0)--(-1,-.18) (1,0)--(1,-.18);
 \draw[thick] (-1,-.18) arc(180:360:1 and .3);
 \draw[TextGray,thin] (-.75,-.35)--(-1.8,-.9);
 \node[left,font=\small] at (-1.8,-.9) {刚性端盘};
 \draw[guide] (0,0)--(0,2.1);
 \draw[<->] (1.55,0)--(1.55,2.1) node[midway,right] {$h$};
 \draw[->] (0,2.1)--(1,2.1) node[midway,above] {$a$};
 \draw[->,DeepBlue] (-.9,2.75) arc(140:30:1.17 and .35);
 \node[above] at (0,3.15) {$\omega$};
 \begin{scope}[shift={(5.1,1.1)}]
  \draw[thick] (0,0) circle(1);
  \draw[guide] (0,0)--(1,0);
  \draw[thick,DeepBlue] (0,0)--(40:1);
  \draw[->] (.55,0) arc(0:40:.55);
  \node[right] at (.63,.25) {$\phi$};
  \node[align=center,text width=3cm] at (0,-1.65) {两截面径向标记\\的相对转角};
 \end{scope}
 \begin{scope}[shift={(9.2,0)}]
  \draw[guide] (0,0)--(0,2.1);
  \draw[thick,DeepBlue] (0,0)--(.7,2.1);
  \draw[<->] (0,2.5)--(.7,2.5);
  \node[above] at (.35,2.6) {$a\phi$};
  \node[align=center,text width=3.3cm] at (.45,-.85) {展开表面\\$\gamma(a)\simeq a\phi/h$};
 \end{scope}
\end{tikzpicture}
\figcap{门轴承的圆柱、截面转角和展开表面。$\phi$ 是绕轴转角，$\gamma$ 是表面剪切角，两者含义不同。}
\label{ch12:fig:hinge-analysis}
\end{minipage}\end{center}

刚性端盘的接触压力为 $p=W/(\pi a^2)$。在所采用的总力矩阈值模型中，
临界时各环的摩擦应力均取极限值 $\mu_fp$，于是
\[
 \tau_s=\int_0^a r\,\mu_fp(2\pi r\dd r)=\frac23\mu_fWa,
 \qquad\phi_{\max}=\frac{\tau_s}{K_\phi}.
\]
这个阈值属于刚性端盘的等效模型。若圆柱弹性底面直接接触支承面，并要求每个局部都严格粘着，
则 $\tau(r)=G_s r\phi/h$ 随半径增加，外缘会先达到 $\mu_fp$：
\[
 \phi_{\rm edge}=\frac{\mu_fWh}{\pi G_sa^3},\qquad
 K_\phi\phi_{\rm edge}=\frac12\mu_fWa.
\]
这只标志局部滑移开始，不能直接当作整个底面的滑动阈值；后续还需接触应力重分布模型。
以下数值均限于题干明确采用的刚性端盘总力矩近似。

粘着时端盘静止，顶面转动使相对角以 $\omega$ 增加。
若一次滑动结束后剩余角为 $\phi_{\min}$，
\[
 t_{\rm stick}=\frac{\phi_{\max}-\phi_{\min}}{\omega},\quad
 f_{\rm rep}=\frac{\omega}{\phi_{\max}-\phi_{\min}},\quad
 \boxed{\Omega_{\rm rep}=\frac{2\pi\omega}{\phi_{\max}-\phi_{\min}}}.
\]
“滑动很快”只说明其持续时间短，并不能推出剩余扭转必为零。
释放后重新粘着还要求 $|\phi_{\min}|\le\phi_{\max}$，否则弹性力矩仍超过最大静摩擦力矩。
设 $q=\phi_{\min}/\phi_{\max}$、$-1\le q<1$，则
\[
 \boxed{\Omega_{\rm rep}=\frac{3\pi^2G_sa^3}{2\mu_fWh(1-q)}\,\omega}.
\]
若采用释放后近似无扭转的闭合模型 $q=0$，代入数值得
\[
 K_\phi=7.85\times10^2\,\mathrm{N\,m/rad},\quad
 \tau_s=0.200\,\mathrm{N\,m},\quad
 \phi_{\max}=2.55\times10^{-4}\,\mathrm{rad},
 \qquad
 \boxed{\Omega_{\rm rep}\simeq2.47\times10^4\,\omega}.
\]
例如 $\omega=0.10\,\mathrm{rad/s}$，该理想模型给出
$t_{\rm stick}=2.55\times10^{-3}\,\mathrm s$、
$f_{\rm rep}\simeq393\,\mathrm{Hz}$。是否真能忽略滑动时间，需要另验接触动力学。

这里求的是粘滑事件的重复频率，不是材料的自由扭转固有频率。
$W$ 是门施加的载荷，不能令轴承质量等于 $W/g$。
若另给一个确切的等效振型、有效转动惯量 $I_{\rm eff}$，
才可建立单自由度近似
$I_{\rm eff}\ddot\phi+K_\phi\phi=0$，得到
$\Omega_0=\sqrt{K_\phi/I_{\rm eff}}$。
快慢分离需 $t_{\rm stick}$ 远大于实际释放时间；
若用 $1/\Omega_0$ 估计释放时间，就还需 $t_{\rm stick}\Omega_0\gg1$。
没有 $I_{\rm eff}$ 及滑动摩擦规律，不能凭门重给这个条件添上数值。

在 $q=0$ 模型下，一次储能为
\[
 U_c=\frac12K_\phi\phi_{\max}^2
 =\frac{\tau_s^2}{2K_\phi}=2.55\times10^{-5}\,\mathrm J.
\]
若释放时这部分能量耗散，平均输入功率为
$\overline P=U_c/t_{\rm stick}=\tau_s\omega/2$，
与粘着期间力矩从零线性增长到 $\tau_s$ 的平均功率相符。
一般 $q$ 时，释放能量为
$K_\phi(\phi_{\max}^2-\phi_{\min}^2)/2$；
相应平均功率为 $K_\phi(\phi_{\max}+\phi_{\min})\omega/2$。
实际声能只是耗散的一部分，其余主要变成接触区内能。
$G_sa^3/(Wh)$ 无量纲，故重复角频率与驱动角速度量纲相同。
\end{xesolution}

\endgroup

\hypersetup{bookmarksdepth=2}

% ThermalQualityRevision20261023-v1
% QualityRevision20261023
\hypersetup{bookmarksdepth=3}
\begingroup
\renewcommand{\Needspace}[1]{\par\begingroup
\dimen0=#1\relax\vskip0pt plus\dimen0\penalty-100
\vskip0pt plus-\dimen0\vskip\dimen0\penalty9999\vskip-\dimen0
\vskip0pt\endgroup}
\tcbset{breakable/.default=false}
\chapter{流体静力学与动力学}
% Editable TikZ: control volumes, free bodies, and local integration elements.
\providecommand{\ThermalFigure}[3]{\begin{qualityfigure}
\begin{tikzpicture}[>=Latex,font=\small,x=.8cm,y=.8cm]#1\end{tikzpicture}
\figcap{#2}\label{#3}\end{qualityfigure}}
\newcommand{\ThFlowElement}[1]{
 \draw[thick](0,.5)--(2.5,.5)--(4,1)--(6.5,1);
 \draw[thick](0,-.5)--(2.5,-.5)--(4,-1)--(6.5,-1);
 \fill[DeepBlue!18](.6,-.5)rectangle(1.8,.5);\fill[DeepBlue!18](4.5,-1)rectangle(5.1,1);
 \draw[<->](.6,-1.5)--node[below]{$v_A\Delta t$}(1.8,-1.5);
 \draw[<->](4.5,-1.5)--node[below]{$v_B\Delta t$}(5.1,-1.5);
 \node[above]at(1.2,.7){$S_A$};\node[above]at(4.8,1.2){$S_B$};#1}
\newcommand{\ThAFlowVolume}{\ThermalFigure{\ThFlowElement{}}{在同一时间内经过两个截面的流体柱，长度分别为当地流速乘时间。}{ch13:quality:volume-a}
图~\ref{ch13:quality:volume-a}中的阴影柱体不是同一时刻被跟踪的两个质点，而是相同时间内穿过截面的体积。稳态、不可压缩且无侧向支流保证控制体没有净体积积累，因而两柱体积相等。若右侧截面变大，右柱必须变短，正对应流速减小。
}
\newcommand{\ThAUniformFlow}{\ThermalFigure{
 \draw[thick](0,-.5)rectangle(7,.5);\draw[dashed](1,-.5)--(1,.5) (5,-.5)--(5,.5);
 \draw[->,DeepBlue,thick](1,0)--(2.2,0)node[above]{$v_1$};\draw[->,DeepBlue,thick](5,0)--(6.2,0)node[above]{$v_2$};
 \node[below]at(3.5,-.8){控制体：等截面直管的两截面之间};
}{一维稳态模型中的相同截面积与流量；不是一般三维速度场的结论。}{ch13:quality:uniform-a}
图~\ref{ch13:quality:uniform-a}的侧壁不漏流，任取两端均有相同的截面积 $S$。若 $v_2>v_1$，每秒流出体积将大于流入体积，控制体内就不能保持处处充满且密度不变，与题设矛盾。
}
\newcommand{\ThAFallingFlow}{\ThermalFigure{
 \draw[DeepBlue,thick](-.8,0)..controls(-.5,-1)and(-.4,-2)..(-.35,-3);
 \draw[DeepBlue,thick](.8,0)..controls(.5,-1)and(.4,-2)..(.35,-3);
 \draw[thick](-.8,0)--(.8,0) (-.38,-2.5)--(.38,-2.5);
 \draw[<->](-1.6,0)--node[left]{$h$}(-1.6,-2.5);
 \draw[->](1.3,0)--(1.3,-.7)node[right]{$v$};\draw[->](1.3,-2.5)--(1.3,-3.6)node[right]{$v(h)$};
 \node[right]at(2.7,-1.2){$S\,v=S(h)v(h)$};
}{先用同一水微元的自由下落求速度，再用稳态连续性比较两个截面。}{ch13:quality:falling-a}
图~\ref{ch13:quality:falling-a}将两个步骤分开：高度差 $h$ 决定微元增加的动能，两个横截面则用于守恒单位时间的流体体积。流速增加时截面必须缩小，不能把两处面积也设为相同。
}
\newcommand{\ThAFlowMass}{\ThermalFigure{\ThFlowElement{
 \node[right]at(7,1){$\Delta m_A=\rho_AS_Av_A\Delta t$};
 \node[right]at(7,-.1){$\Delta m_B=\rho_BS_Bv_B\Delta t$};
}}{可压缩流的两体积柱可以不等大，但通过的质量须相同。}{ch13:quality:mass-a}
图~\ref{ch13:quality:mass-a}中每一柱体还须乘当地密度才得到质量。控制体内部的质量既不累积也不凭空产生，因此应令 $\Delta m_A=\Delta m_B$；只有同密度时才可进一步令两柱体积相等。
}
\newcommand{\ThDilutingBase}[1]{
 \draw[thick](-.6,0)--(-.6,-3) (.6,0)--(.6,-3);
 \draw[dashed](-1,0)--(1,0) (-1,-2.5)--(1,-2.5);
 \node[above]at(0,.2){$S,\rho,v$};\node[below]at(0,-3.2){$S,\rho(h),v(h)$};
 \draw[<->](-1.4,0)--node[left]{$h$}(-1.4,-2.5);
 \foreach \y in {-.2,-.4,-.6,-1.4,-1.9,-2.5}{\fill[DeepBlue](0,\y)circle(1.5pt);}#1}
\newcommand{\ThQDilutingColumn}{\ThermalFigure{\ThDilutingBase{}}{题设人为固定横截面积；分子沿竖直方向自由下落。}{ch13:quality:diluting-q}}
\newcommand{\ThADilutingColumn}{\ThermalFigure{\ThDilutingBase{
 \draw[->](1.5,-.1)--(1.5,-.7)node[right]{$v$};\draw[->](1.5,-2)--(1.5,-3)node[right]{$v(h)$};
}}{下方流速更大而面积固定，分子间的纵向间距增大。}{ch13:quality:diluting-a}
图~\ref{ch13:quality:diluting-a}中粒子点的疏密表示题设模型的密度变化。两个截面共用同一个 $S$，连续性要求 $\rho(h)v(h)=\rho v$；这里降低的是密度，不能照搬上一题缩小截面的真实水束模型。
}
\newcommand{\ThLayerBase}[1]{
 \fill[DeepBlue!10](0,1.5)rectangle(3,3);\fill[DeepBlue!25](0,0)rectangle(3,1.5);
 \draw[thick](0,3.3)--(0,0)--(3,0)--(3,3.3);\draw[DeepBlue](0,3)--(3,3) (0,1.5)--(3,1.5);
 \node at(1.5,2.3){$\rho_1,h_1$};\node at(1.5,.75){$\rho_2,h_2$};
 \draw[->](3.6,1.5)--(3.6,.2)node[right]{$x$};#1}
\newcommand{\ThQLayers}{\ThermalFigure{\ThLayerBase{\node[right]at(4.5,2.5){壁宽 $b$ 垂直纸面};}}{界面是下层深度坐标的起点，外侧大气压在净压力中抵消。}{ch13:quality:layers-q}}
\newcommand{\ThALayers}{\ThermalFigure{\ThLayerBase{
 \draw[DeepBlue,thick](3,1.5)--(4.2,1.5)--(5,0)--(3,0);
 \foreach \y in {.2,.65,1.1}{\draw[->,DeepBlue](3,\y)--({5-.8*\y/1.5},\y);}
 \node[right]at(5.3,.8){梯形表压分布};
}}{下层侧壁表压等于常量项与线性项之和，合力矩须逐条积分。}{ch13:quality:layers-a}
图~\ref{ch13:quality:layers-a}中下层顶部已经有上层液柱贡献的表压，故压力图是梯形而非从零开始的三角形。深度为 $x$ 的窄条面积为 $b\dd x$，其关于界面的力矩为 $x\Delta p(x)b\dd x$，由此得到合力的作用点。
}
\newcommand{\ThFloatBase}[1]{
 \fill[DeepBlue!10](-1,-1.5)rectangle(4,.4);\draw[DeepBlue](-1,.4)--(4,.4);
 \draw[thick](.7,-.8)rectangle(2.3,1.8);\node at(1.5,1.3){$m,S$};
 \draw[<->](2.9,-.8)--node[right]{$h_0$}(2.9,.4);#1}
\newcommand{\ThQFloat}{\ThermalFigure{\ThFloatBase{}}{直立柱体的平衡浸深；小振动中顶面仍露出、底面仍浸没。}{ch13:quality:float-q}}
\newcommand{\ThAFloat}{\ThermalFigure{\ThFloatBase{
 \draw[dashed](.7,-1.1)--(2.3,-1.1);\draw[->](.2,-.8)--(.2,-1.1)node[left]{$x$};
 \draw[->,thick](1.5,.1)--(1.5,-1.5)node[below]{$mg$};
 \draw[->,DeepBlue,thick](1.5,.1)--(1.5,1)node[right]{$F_b$};
}}{向下位移增加排液体积 $Sx$，额外浮力朝向上方。}{ch13:quality:float-a}
图~\ref{ch13:quality:float-a}将 $x$ 取向下为正；额外排液体积为 $Sx$，因此浮力增量的有向分量是 $-\rho gSx$。先用平衡条件消去 $mg-\rho gSh_0$，才留下关于平衡位置的振动方程。
}
\newcommand{\ThQOutlet}{\ThermalFigure{
 \draw[thick](0,3)--(0,0)--(3,0)--(3,.8) (3,1.1)--(3,3);\draw[DeepBlue](0,2.5)--(3,2.5);
 \node[left]at(0,2.5){$A$};\node[right]at(3.1,.95){$B$};\draw[->](3.1,.95)--(4.8,.95)node[right]{$v$};
 \draw[<->](5.5,.95)--node[right]{$h$}(5.5,2.5);
}{自由液面 $A$ 与出口 $B$ 均接大气，容器横截面远大于孔口。}{ch13:quality:outlet-q}}
\newcommand{\ThJetBase}[1]{
 \draw[DeepBlue,thick](0,.3)--(4,.3) (0,-.3)--(4,-.3);\draw[very thick](4,-1.6)--(4,1.6);
 \draw[->,thick](.6,0)--(2,0)node[midway,above]{$v$};\node[below]at(1,-.6){$A,\rho$};#1}
\newcommand{\ThQJet}{\ThermalFigure{\ThJetBase{\node[right]at(4.5,1){静止大平板};}}{自由射流垂直撞击平板后沿板面流走。}{ch13:quality:jet-q}}
\newcommand{\ThQMovingPlate}{\ThermalFigure{\ThJetBase{
 \draw[->,thick](4.5,1)--(6,1)node[right]{$u$};\node[right]at(4.5,-.8){$0\le u<v$};
}}{平板沿来流方向匀速后退，截面流速 $v$ 仍相对地面给定。}{ch13:quality:moving-q}}
\newcommand{\ThAMovingPlate}{\ThermalFigure{
 \draw[dashed](-.5,-1)rectangle(4.5,1);\draw[very thick](3.5,-1.4)--(3.5,1.4);
 \draw[->,DeepBlue,thick](0,0)--(2,0)node[midway,above]{$v-u$};
 \draw[->,DeepBlue,thick](3.5,.1)--(3.5,1.4);\draw[->,DeepBlue,thick](3.5,-.1)--(3.5,-1.4);
 \node[right,align=left]at(5,0){板系控制体静止\\出流法向速度为零};
}{在匀速平板的惯性系中计算撞击质量率和法向动量变化。}{ch13:quality:moving-a}
图~\ref{ch13:quality:moving-a}的虚线控制体随板一起平移，水以相对速度 $v-u$ 穿过入口。因此流入质量率和每单位质量损失的法向动量都含 $v-u$，合力中出现它的平方；向地面系中的平板传功时再乘板速 $u$。
}
\newcommand{\ThQSettling}{\ThermalFigure{
 \draw[thick](-2,2)--(-2,-2)--(2,-2)--(2,2);\draw[DeepBlue](-2,1.6)--(2,1.6);
 \draw[thick](0,0)circle(.55);\draw[->](0,0)--node[above]{$R$}(.55,0);
 \draw[->,DeepBlue,thick](1,0)--(1,-1.2)node[below]{$v$};
 \node[right,align=left]at(3,0){球 $\rho_s$\\液体 $\rho_f,\eta$};
}{小球远离容器壁的缓慢下沉模型；终端状态由合力为零确定。}{ch13:quality:settling-q}}
\newcommand{\ThDropletBase}[1]{
 \draw[DeepBlue,thick](0,0)circle(.65);\draw[DeepBlue,thick](2,0)circle(.65);
 \node[below]at(0,-.8){$R$};\node[below]at(2,-.8){$R$};\draw[->,thick](3.2,0)--(4.2,0);
 \draw[DeepBlue,thick](5.7,0)circle(.82);\node[below]at(5.7,-1){#1};}
\newcommand{\ThQDroplets}{\ThermalFigure{\ThDropletBase{$R_f$ 待求}}{两滴同种液体合并为一滴，体积守恒而界面面积改变。}{ch13:quality:droplets-q}}
\newcommand{\ThADroplets}{\ThermalFigure{\ThDropletBase{$R_f=2^{1/3}R$}
 \node[above]at(1,1.3){$A_i=8\pi R^2$};\node[above]at(5.7,1.3){$A_f=4\pi R_f^2$};
}{先由体积相加求半径，再比较新旧界面面积。}{ch13:quality:droplets-a}
图~\ref{ch13:quality:droplets-a}中的两个初始界面都计入 $A_i$，末态只剩一个界面。界面能的减少为 $\gamma(A_i-A_f)$；这些能量转移到流动和内能，不能解释为总能量消失。
}
\newcommand{\ThACoupledMercury}{\ThermalFigure{
 \draw[thick](0,-.4)rectangle(7,.4);\draw[line width=3pt](1,-.4)--(1,.4) (6,-.4)--(6,.4);
 \draw[->,DeepBlue,thick](1,1)--(2,1)node[right]{$q_1$};\draw[->,DeepBlue,thick](6,1)--(5,1)node[left]{$q_2$};
 \node at(3.5,0){绝热空气 $\bar p+\delta p$};
 \draw[<->](1,-1)--node[below]{$\ell-q_1-q_2$}(6,-1);
}{抽出两个水银内侧界面的位移，以空气柱的压缩量建立耦合项。}{ch13:quality:coupled-a}
图~\ref{ch13:quality:coupled-a}把两个位移都取向空气内部为正，所以气柱缩短量为 $q_1+q_2$。同一个 $\delta p$ 同时作用于两段水银，但两段惯性分别由各自总长度 $L_1,L_2$ 决定，不能只用竖直高度 $h$ 作为质量长度。
}
\newcommand{\ThACouette}{\ThermalFigure{
 \draw[thick](0,0)circle(1);\draw[thick](0,0)circle(2.3);\draw[DeepBlue,dashed](0,0)circle(1.6);
 \draw[->](0,0)--node[above]{$r$}(1.6,0);\fill(1.6,0)circle(1.5pt);
 \draw[->,DeepBlue,thick](1.6,0)--(1.6,1.2)node[right]{$v_\theta=r\omega$};
 \node[right,align=left]at(3,0){圆柱面面积 $2\pi r l$\\切应力 $\tau_{r\theta}$\\力矩再乘力臂 $r$};
}{在半径 $r$ 的同轴圆柱面上传递黏性力矩，不能只令切应力恒定。}{ch13:quality:couette-a}
图~\ref{ch13:quality:couette-a}的虚线圆柱面传递的总力矩为 $2\pi l r^2\tau_{r\theta}$。稳态无角动量积累，要求该乘积与 $r$ 无关；柱坐标中的 $-v_\theta/r$ 项保证刚体转动 $v_\theta\propto r$ 时没有虚假的剪切应力。
}
\newcommand{\ThAHangar}{\ThermalFigure{
 \draw[thick](-2,0)arc(180:0:2);\draw[thick](-2,0)--(2,0);\fill(0,0)circle(1pt)node[below]{$O$};
 \draw[->](0,0)--(1.414,1.414)node[midway,below right]{$R$};
 \draw[->,DeepBlue,thick](1.414,1.414)--(2.15,2.15)node[right]{$\Delta p\,\dd A$};
 \draw[->](.65,0)arc(0:45:.65)node[right]{$\theta$};\node[left]at(-2,0){$A$};
}{曲面微元的净压力沿外法线；竖直分量另乘 $\sin\theta$。}{ch13:quality:hangar-a}
图~\ref{ch13:quality:hangar-a}中微元面积为 $\dd A=lR\dd\theta$。压力差先由门口 $A$ 的停滞压确定库内压强，再与曲面外压比较；左右对称使水平分量相消，竖直积分还应包含法线的投影因子 $\sin\theta$。
}
\newcommand{\ThAFunnel}{\ThermalFigure{
 \draw[thick](-2,3)--(-.12,0) (.12,0)--(2,3);\draw[DeepBlue](-1.35,2)--(1.35,2);
 \draw[dashed](0,0)--(0,3.2);\draw[<->](2.5,0)--node[right]{$h$}(2.5,2);
 \draw[<->](-1.35,2.5)--node[above]{$2Rh/H$}(1.35,2.5);
 \draw[->,DeepBlue](0,0)--(0,-.8)node[right]{$Q$};
}{尖锥近似下自由液面的半径随水深变化，积分时不能保持水面面积不变。}{ch13:quality:funnel-a}
图~\ref{ch13:quality:funnel-a}中的相似三角形给出自由液面半径 $Rh/H$；液面下降 $\dd h<0$ 排出的体积是 $-A(h)\dd h$。将它等于小孔流出体积 $Q\dd t$，即可把流速公式转为排液时间积分。
}
\newcommand{\ThAShore}{\ThermalFigure{
 \draw[thick](-1,0)--(6,0)node[right]{岸线 $y=0$};\draw[->](0,0)--(0,3.4)node[above]{$y$};
 \draw[DeepBlue,thick](0,3)..controls(3,3)and(4,1.5)..(4.4,.3);\fill(0,3)circle(1.5pt)node[left]{$B$};
 \fill(3.2,1.9)circle(1.5pt)node[above right]{$P$};\draw[dashed](3.2,1.9)--(3.2,.6);
 \draw[->,thick](3.2,1.9)--(3.9,1.1);\draw(3.2,1.3)arc(270:311:.6)node[below]{$\theta$};
}{局部传播方向与朝岸法线 $-y$ 的夹角为 $\theta$；图示只标方向，不代替精确摆线。}{ch13:quality:shore-a}
图~\ref{ch13:quality:shore-a}中射线朝岸传播，所以 $\dd y<0$ 而 $\dd x>0$，这决定 $\dd x/\dd y=-\tan\theta$ 的负号。沿岸平移对称给出 $\sin\theta/c(y)$ 不变，初点 $B$ 的水平切向又给出其常量值。
}
\newcommand{\ThAHose}{\ThermalFigure{
 \coordinate(O)at(0,0);\coordinate(E)at(-.8,-2.3);\coordinate(P)at(1.5,-3.1);
 \draw[very thick](O)--(E)--(P);\fill(O)circle(2pt)node[above]{$O$};\node[left]at(E){$E$};\node[below]at(P){$P$};
 \draw[->,DeepBlue,thick](P)--(2.7,-3.5)node[right]{$u\vect e_2$};
 \fill(-.025,-1.925)circle(1.5pt)node[left]{$C$};\draw[->,thick](-.025,-1.925)--(-.025,-3.5)node[below]{$Mg$};
 \draw[dashed](0,0)--(1.5,-3.1);\node[right]at(2,-1){关于 $O$ 比较重力矩与出流角动量通量};
}{管和管内水组成控制体；出口动量通量关于接头 $O$ 有非零力矩。}{ch13:quality:hose-a}
图~\ref{ch13:quality:hose-a}的 $C$ 表示管与管内水的共同质心，两段均匀且等长时它位于两段中点的平均位置。入口位于 $O$，关于 $O$ 的入口角动量为零；出口 $P$ 则必须保留位置矢量与出流动量的叉乘。
}



{
\graphicspath{{../../../sample/xes-original/kinematics-supplement/}{../../../sample/xes-original/kinematics-supplement/assets/}}
\section{流线、流管与质量流率}
\newcommand{\SphereStreamlinesDiagram}{%
\begin{center}\begin{minipage}{\linewidth}\centering\begin{tikzpicture}[scale=.85,>=Latex]
 \shade[ball color=TextGray!35] (0,0) circle (1.2);
 \draw[Blue,thick] plot[smooth] coordinates {(-4.000000,0.354765) (-3.955556,0.354930) (-3.911111,0.355102) (-3.866667,0.355283) (-3.822222,0.355472) (-3.777778,0.355670) (-3.733333,0.355878) (-3.688889,0.356097) (-3.644444,0.356326) (-3.600000,0.356568) (-3.555556,0.356822) (-3.511111,0.357089) (-3.466667,0.357371) (-3.422222,0.357668) (-3.377778,0.357981) (-3.333333,0.358312) (-3.288889,0.358662) (-3.244444,0.359032) (-3.200000,0.359424) (-3.155556,0.359839) (-3.111111,0.360279) (-3.066667,0.360747) (-3.022222,0.361244) (-2.977778,0.361772) (-2.933333,0.362335) (-2.888889,0.362935) (-2.844444,0.363575) (-2.800000,0.364259) (-2.755556,0.364992) (-2.711111,0.365776) (-2.666667,0.366618) (-2.622222,0.367522) (-2.577778,0.368496) (-2.533333,0.369545) (-2.488889,0.370677) (-2.444444,0.371902) (-2.400000,0.373228) (-2.355556,0.374668) (-2.311111,0.376235) (-2.266667,0.377943) (-2.222222,0.379809) (-2.177778,0.381853) (-2.133333,0.384097) (-2.088889,0.386568) (-2.044444,0.389297) (-2.000000,0.392320) (-1.955556,0.395681) (-1.911111,0.399429) (-1.866667,0.403626) (-1.822222,0.408346) (-1.777778,0.413675) (-1.733333,0.419722) (-1.688889,0.426615) (-1.644444,0.434514) (-1.600000,0.443613) (-1.555556,0.454151) (-1.511111,0.466414) (-1.466667,0.480748) (-1.422222,0.497548) (-1.377778,0.517251) (-1.333333,0.540285) (-1.288889,0.566991) (-1.244444,0.597501) (-1.200000,0.631616) (-1.155556,0.668746) (-1.111111,0.707979) (-1.066667,0.748256) (-1.022222,0.788579) (-0.977778,0.828132) (-0.933333,0.866331) (-0.888889,0.902794) (-0.844444,0.937303) (-0.800000,0.969748) (-0.755556,1.000093) (-0.711111,1.028348) (-0.666667,1.054547) (-0.622222,1.078741) (-0.577778,1.100985) (-0.533333,1.121337) (-0.488889,1.139853) (-0.444444,1.156584) (-0.400000,1.171580) (-0.355556,1.184883) (-0.311111,1.196533) (-0.266667,1.206563) (-0.222222,1.215001) (-0.177778,1.221873) (-0.133333,1.227198) (-0.088889,1.230991) (-0.044444,1.233262) (0.000000,1.234019) (0.044444,1.233262) (0.088889,1.230991) (0.133333,1.227198) (0.177778,1.221873) (0.222222,1.215001) (0.266667,1.206563) (0.311111,1.196533) (0.355556,1.184883) (0.400000,1.171580) (0.444444,1.156584) (0.488889,1.139853) (0.533333,1.121337) (0.577778,1.100985) (0.622222,1.078741) (0.666667,1.054547) (0.711111,1.028348) (0.755556,1.000093) (0.800000,0.969748) (0.844444,0.937303) (0.888889,0.902794) (0.933333,0.866331) (0.977778,0.828132) (1.022222,0.788579) (1.066667,0.748256) (1.111111,0.707979) (1.155556,0.668746) (1.200000,0.631616) (1.244444,0.597501) (1.288889,0.566991) (1.333333,0.540285) (1.377778,0.517251) (1.422222,0.497548) (1.466667,0.480748) (1.511111,0.466414) (1.555556,0.454151) (1.600000,0.443613) (1.644444,0.434514) (1.688889,0.426615) (1.733333,0.419722) (1.777778,0.413675) (1.822222,0.408346) (1.866667,0.403626) (1.911111,0.399429) (1.955556,0.395681) (2.000000,0.392320) (2.044444,0.389297) (2.088889,0.386568) (2.133333,0.384097) (2.177778,0.381853) (2.222222,0.379809) (2.266667,0.377943) (2.311111,0.376235) (2.355556,0.374668) (2.400000,0.373228) (2.444444,0.371902) (2.488889,0.370677) (2.533333,0.369545) (2.577778,0.368496) (2.622222,0.367522) (2.666667,0.366618) (2.711111,0.365776) (2.755556,0.364992) (2.800000,0.364259) (2.844444,0.363575) (2.888889,0.362935) (2.933333,0.362335) (2.977778,0.361772) (3.022222,0.361244) (3.066667,0.360747) (3.111111,0.360279) (3.155556,0.359839) (3.200000,0.359424) (3.244444,0.359032) (3.288889,0.358662) (3.333333,0.358312) (3.377778,0.357981) (3.422222,0.357668) (3.466667,0.357371) (3.511111,0.357089) (3.555556,0.356822) (3.600000,0.356568) (3.644444,0.356326) (3.688889,0.356097) (3.733333,0.355878) (3.777778,0.355670) (3.822222,0.355472) (3.866667,0.355283) (3.911111,0.355102) (3.955556,0.354930) (4.000000,0.354765)};
\draw[Blue,thick,->] (2.444444,0.371902)--(2.711111,0.365776);
\draw[Blue,thick] plot[smooth] coordinates {(-4.000000,-0.354765) (-3.955556,-0.354930) (-3.911111,-0.355102) (-3.866667,-0.355283) (-3.822222,-0.355472) (-3.777778,-0.355670) (-3.733333,-0.355878) (-3.688889,-0.356097) (-3.644444,-0.356326) (-3.600000,-0.356568) (-3.555556,-0.356822) (-3.511111,-0.357089) (-3.466667,-0.357371) (-3.422222,-0.357668) (-3.377778,-0.357981) (-3.333333,-0.358312) (-3.288889,-0.358662) (-3.244444,-0.359032) (-3.200000,-0.359424) (-3.155556,-0.359839) (-3.111111,-0.360279) (-3.066667,-0.360747) (-3.022222,-0.361244) (-2.977778,-0.361772) (-2.933333,-0.362335) (-2.888889,-0.362935) (-2.844444,-0.363575) (-2.800000,-0.364259) (-2.755556,-0.364992) (-2.711111,-0.365776) (-2.666667,-0.366618) (-2.622222,-0.367522) (-2.577778,-0.368496) (-2.533333,-0.369545) (-2.488889,-0.370677) (-2.444444,-0.371902) (-2.400000,-0.373228) (-2.355556,-0.374668) (-2.311111,-0.376235) (-2.266667,-0.377943) (-2.222222,-0.379809) (-2.177778,-0.381853) (-2.133333,-0.384097) (-2.088889,-0.386568) (-2.044444,-0.389297) (-2.000000,-0.392320) (-1.955556,-0.395681) (-1.911111,-0.399429) (-1.866667,-0.403626) (-1.822222,-0.408346) (-1.777778,-0.413675) (-1.733333,-0.419722) (-1.688889,-0.426615) (-1.644444,-0.434514) (-1.600000,-0.443613) (-1.555556,-0.454151) (-1.511111,-0.466414) (-1.466667,-0.480748) (-1.422222,-0.497548) (-1.377778,-0.517251) (-1.333333,-0.540285) (-1.288889,-0.566991) (-1.244444,-0.597501) (-1.200000,-0.631616) (-1.155556,-0.668746) (-1.111111,-0.707979) (-1.066667,-0.748256) (-1.022222,-0.788579) (-0.977778,-0.828132) (-0.933333,-0.866331) (-0.888889,-0.902794) (-0.844444,-0.937303) (-0.800000,-0.969748) (-0.755556,-1.000093) (-0.711111,-1.028348) (-0.666667,-1.054547) (-0.622222,-1.078741) (-0.577778,-1.100985) (-0.533333,-1.121337) (-0.488889,-1.139853) (-0.444444,-1.156584) (-0.400000,-1.171580) (-0.355556,-1.184883) (-0.311111,-1.196533) (-0.266667,-1.206563) (-0.222222,-1.215001) (-0.177778,-1.221873) (-0.133333,-1.227198) (-0.088889,-1.230991) (-0.044444,-1.233262) (0.000000,-1.234019) (0.044444,-1.233262) (0.088889,-1.230991) (0.133333,-1.227198) (0.177778,-1.221873) (0.222222,-1.215001) (0.266667,-1.206563) (0.311111,-1.196533) (0.355556,-1.184883) (0.400000,-1.171580) (0.444444,-1.156584) (0.488889,-1.139853) (0.533333,-1.121337) (0.577778,-1.100985) (0.622222,-1.078741) (0.666667,-1.054547) (0.711111,-1.028348) (0.755556,-1.000093) (0.800000,-0.969748) (0.844444,-0.937303) (0.888889,-0.902794) (0.933333,-0.866331) (0.977778,-0.828132) (1.022222,-0.788579) (1.066667,-0.748256) (1.111111,-0.707979) (1.155556,-0.668746) (1.200000,-0.631616) (1.244444,-0.597501) (1.288889,-0.566991) (1.333333,-0.540285) (1.377778,-0.517251) (1.422222,-0.497548) (1.466667,-0.480748) (1.511111,-0.466414) (1.555556,-0.454151) (1.600000,-0.443613) (1.644444,-0.434514) (1.688889,-0.426615) (1.733333,-0.419722) (1.777778,-0.413675) (1.822222,-0.408346) (1.866667,-0.403626) (1.911111,-0.399429) (1.955556,-0.395681) (2.000000,-0.392320) (2.044444,-0.389297) (2.088889,-0.386568) (2.133333,-0.384097) (2.177778,-0.381853) (2.222222,-0.379809) (2.266667,-0.377943) (2.311111,-0.376235) (2.355556,-0.374668) (2.400000,-0.373228) (2.444444,-0.371902) (2.488889,-0.370677) (2.533333,-0.369545) (2.577778,-0.368496) (2.622222,-0.367522) (2.666667,-0.366618) (2.711111,-0.365776) (2.755556,-0.364992) (2.800000,-0.364259) (2.844444,-0.363575) (2.888889,-0.362935) (2.933333,-0.362335) (2.977778,-0.361772) (3.022222,-0.361244) (3.066667,-0.360747) (3.111111,-0.360279) (3.155556,-0.359839) (3.200000,-0.359424) (3.244444,-0.359032) (3.288889,-0.358662) (3.333333,-0.358312) (3.377778,-0.357981) (3.422222,-0.357668) (3.466667,-0.357371) (3.511111,-0.357089) (3.555556,-0.356822) (3.600000,-0.356568) (3.644444,-0.356326) (3.688889,-0.356097) (3.733333,-0.355878) (3.777778,-0.355670) (3.822222,-0.355472) (3.866667,-0.355283) (3.911111,-0.355102) (3.955556,-0.354930) (4.000000,-0.354765)};
\draw[Blue,thick,->] (2.444444,-0.371902)--(2.711111,-0.365776);
\draw[Blue,thick] plot[smooth] coordinates {(-4.000000,0.759789) (-3.955556,0.760117) (-3.911111,0.760460) (-3.866667,0.760818) (-3.822222,0.761193) (-3.777778,0.761585) (-3.733333,0.761996) (-3.688889,0.762426) (-3.644444,0.762877) (-3.600000,0.763350) (-3.555556,0.763847) (-3.511111,0.764369) (-3.466667,0.764917) (-3.422222,0.765493) (-3.377778,0.766100) (-3.333333,0.766738) (-3.288889,0.767411) (-3.244444,0.768120) (-3.200000,0.768868) (-3.155556,0.769657) (-3.111111,0.770491) (-3.066667,0.771373) (-3.022222,0.772307) (-2.977778,0.773295) (-2.933333,0.774342) (-2.888889,0.775453) (-2.844444,0.776632) (-2.800000,0.777885) (-2.755556,0.779218) (-2.711111,0.780637) (-2.666667,0.782148) (-2.622222,0.783760) (-2.577778,0.785481) (-2.533333,0.787320) (-2.488889,0.789287) (-2.444444,0.791394) (-2.400000,0.793653) (-2.355556,0.796078) (-2.311111,0.798683) (-2.266667,0.801485) (-2.222222,0.804502) (-2.177778,0.807756) (-2.133333,0.811267) (-2.088889,0.815062) (-2.044444,0.819167) (-2.000000,0.823613) (-1.955556,0.828432) (-1.911111,0.833662) (-1.866667,0.839340) (-1.822222,0.845511) (-1.777778,0.852221) (-1.733333,0.859519) (-1.688889,0.867457) (-1.644444,0.876089) (-1.600000,0.885468) (-1.555556,0.895648) (-1.511111,0.906679) (-1.466667,0.918602) (-1.422222,0.931452) (-1.377778,0.945247) (-1.333333,0.959987) (-1.288889,0.975652) (-1.244444,0.992195) (-1.200000,1.009543) (-1.155556,1.027597) (-1.111111,1.046235) (-1.066667,1.065313) (-1.022222,1.084677) (-0.977778,1.104165) (-0.933333,1.123613) (-0.888889,1.142866) (-0.844444,1.161777) (-0.800000,1.180214) (-0.755556,1.198058) (-0.711111,1.215207) (-0.666667,1.231576) (-0.622222,1.247091) (-0.577778,1.261692) (-0.533333,1.275330) (-0.488889,1.287967) (-0.444444,1.299572) (-0.400000,1.310119) (-0.355556,1.319591) (-0.311111,1.327973) (-0.266667,1.335253) (-0.222222,1.341424) (-0.177778,1.346479) (-0.133333,1.350415) (-0.088889,1.353228) (-0.044444,1.354917) (0.000000,1.355480) (0.044444,1.354917) (0.088889,1.353228) (0.133333,1.350415) (0.177778,1.346479) (0.222222,1.341424) (0.266667,1.335253) (0.311111,1.327973) (0.355556,1.319591) (0.400000,1.310119) (0.444444,1.299572) (0.488889,1.287967) (0.533333,1.275330) (0.577778,1.261692) (0.622222,1.247091) (0.666667,1.231576) (0.711111,1.215207) (0.755556,1.198058) (0.800000,1.180214) (0.844444,1.161777) (0.888889,1.142866) (0.933333,1.123613) (0.977778,1.104165) (1.022222,1.084677) (1.066667,1.065313) (1.111111,1.046235) (1.155556,1.027597) (1.200000,1.009543) (1.244444,0.992195) (1.288889,0.975652) (1.333333,0.959987) (1.377778,0.945247) (1.422222,0.931452) (1.466667,0.918602) (1.511111,0.906679) (1.555556,0.895648) (1.600000,0.885468) (1.644444,0.876089) (1.688889,0.867457) (1.733333,0.859519) (1.777778,0.852221) (1.822222,0.845511) (1.866667,0.839340) (1.911111,0.833662) (1.955556,0.828432) (2.000000,0.823613) (2.044444,0.819167) (2.088889,0.815062) (2.133333,0.811267) (2.177778,0.807756) (2.222222,0.804502) (2.266667,0.801485) (2.311111,0.798683) (2.355556,0.796078) (2.400000,0.793653) (2.444444,0.791394) (2.488889,0.789287) (2.533333,0.787320) (2.577778,0.785481) (2.622222,0.783760) (2.666667,0.782148) (2.711111,0.780637) (2.755556,0.779218) (2.800000,0.777885) (2.844444,0.776632) (2.888889,0.775453) (2.933333,0.774342) (2.977778,0.773295) (3.022222,0.772307) (3.066667,0.771373) (3.111111,0.770491) (3.155556,0.769657) (3.200000,0.768868) (3.244444,0.768120) (3.288889,0.767411) (3.333333,0.766738) (3.377778,0.766100) (3.422222,0.765493) (3.466667,0.764917) (3.511111,0.764369) (3.555556,0.763847) (3.600000,0.763350) (3.644444,0.762877) (3.688889,0.762426) (3.733333,0.761996) (3.777778,0.761585) (3.822222,0.761193) (3.866667,0.760818) (3.911111,0.760460) (3.955556,0.760117) (4.000000,0.759789)};
\draw[Blue,thick,->] (2.444444,0.791394)--(2.711111,0.780637);
\draw[Blue,thick] plot[smooth] coordinates {(-4.000000,-0.759789) (-3.955556,-0.760117) (-3.911111,-0.760460) (-3.866667,-0.760818) (-3.822222,-0.761193) (-3.777778,-0.761585) (-3.733333,-0.761996) (-3.688889,-0.762426) (-3.644444,-0.762877) (-3.600000,-0.763350) (-3.555556,-0.763847) (-3.511111,-0.764369) (-3.466667,-0.764917) (-3.422222,-0.765493) (-3.377778,-0.766100) (-3.333333,-0.766738) (-3.288889,-0.767411) (-3.244444,-0.768120) (-3.200000,-0.768868) (-3.155556,-0.769657) (-3.111111,-0.770491) (-3.066667,-0.771373) (-3.022222,-0.772307) (-2.977778,-0.773295) (-2.933333,-0.774342) (-2.888889,-0.775453) (-2.844444,-0.776632) (-2.800000,-0.777885) (-2.755556,-0.779218) (-2.711111,-0.780637) (-2.666667,-0.782148) (-2.622222,-0.783760) (-2.577778,-0.785481) (-2.533333,-0.787320) (-2.488889,-0.789287) (-2.444444,-0.791394) (-2.400000,-0.793653) (-2.355556,-0.796078) (-2.311111,-0.798683) (-2.266667,-0.801485) (-2.222222,-0.804502) (-2.177778,-0.807756) (-2.133333,-0.811267) (-2.088889,-0.815062) (-2.044444,-0.819167) (-2.000000,-0.823613) (-1.955556,-0.828432) (-1.911111,-0.833662) (-1.866667,-0.839340) (-1.822222,-0.845511) (-1.777778,-0.852221) (-1.733333,-0.859519) (-1.688889,-0.867457) (-1.644444,-0.876089) (-1.600000,-0.885468) (-1.555556,-0.895648) (-1.511111,-0.906679) (-1.466667,-0.918602) (-1.422222,-0.931452) (-1.377778,-0.945247) (-1.333333,-0.959987) (-1.288889,-0.975652) (-1.244444,-0.992195) (-1.200000,-1.009543) (-1.155556,-1.027597) (-1.111111,-1.046235) (-1.066667,-1.065313) (-1.022222,-1.084677) (-0.977778,-1.104165) (-0.933333,-1.123613) (-0.888889,-1.142866) (-0.844444,-1.161777) (-0.800000,-1.180214) (-0.755556,-1.198058) (-0.711111,-1.215207) (-0.666667,-1.231576) (-0.622222,-1.247091) (-0.577778,-1.261692) (-0.533333,-1.275330) (-0.488889,-1.287967) (-0.444444,-1.299572) (-0.400000,-1.310119) (-0.355556,-1.319591) (-0.311111,-1.327973) (-0.266667,-1.335253) (-0.222222,-1.341424) (-0.177778,-1.346479) (-0.133333,-1.350415) (-0.088889,-1.353228) (-0.044444,-1.354917) (0.000000,-1.355480) (0.044444,-1.354917) (0.088889,-1.353228) (0.133333,-1.350415) (0.177778,-1.346479) (0.222222,-1.341424) (0.266667,-1.335253) (0.311111,-1.327973) (0.355556,-1.319591) (0.400000,-1.310119) (0.444444,-1.299572) (0.488889,-1.287967) (0.533333,-1.275330) (0.577778,-1.261692) (0.622222,-1.247091) (0.666667,-1.231576) (0.711111,-1.215207) (0.755556,-1.198058) (0.800000,-1.180214) (0.844444,-1.161777) (0.888889,-1.142866) (0.933333,-1.123613) (0.977778,-1.104165) (1.022222,-1.084677) (1.066667,-1.065313) (1.111111,-1.046235) (1.155556,-1.027597) (1.200000,-1.009543) (1.244444,-0.992195) (1.288889,-0.975652) (1.333333,-0.959987) (1.377778,-0.945247) (1.422222,-0.931452) (1.466667,-0.918602) (1.511111,-0.906679) (1.555556,-0.895648) (1.600000,-0.885468) (1.644444,-0.876089) (1.688889,-0.867457) (1.733333,-0.859519) (1.777778,-0.852221) (1.822222,-0.845511) (1.866667,-0.839340) (1.911111,-0.833662) (1.955556,-0.828432) (2.000000,-0.823613) (2.044444,-0.819167) (2.088889,-0.815062) (2.133333,-0.811267) (2.177778,-0.807756) (2.222222,-0.804502) (2.266667,-0.801485) (2.311111,-0.798683) (2.355556,-0.796078) (2.400000,-0.793653) (2.444444,-0.791394) (2.488889,-0.789287) (2.533333,-0.787320) (2.577778,-0.785481) (2.622222,-0.783760) (2.666667,-0.782148) (2.711111,-0.780637) (2.755556,-0.779218) (2.800000,-0.777885) (2.844444,-0.776632) (2.888889,-0.775453) (2.933333,-0.774342) (2.977778,-0.773295) (3.022222,-0.772307) (3.066667,-0.771373) (3.111111,-0.770491) (3.155556,-0.769657) (3.200000,-0.768868) (3.244444,-0.768120) (3.288889,-0.767411) (3.333333,-0.766738) (3.377778,-0.766100) (3.422222,-0.765493) (3.466667,-0.764917) (3.511111,-0.764369) (3.555556,-0.763847) (3.600000,-0.763350) (3.644444,-0.762877) (3.688889,-0.762426) (3.733333,-0.761996) (3.777778,-0.761585) (3.822222,-0.761193) (3.866667,-0.760818) (3.911111,-0.760460) (3.955556,-0.760117) (4.000000,-0.759789)};
\draw[Blue,thick,->] (2.444444,-0.791394)--(2.711111,-0.780637);
\draw[Blue,thick] plot[smooth] coordinates {(-4.000000,1.214450) (-3.955556,1.214907) (-3.911111,1.215383) (-3.866667,1.215879) (-3.822222,1.216396) (-3.777778,1.216935) (-3.733333,1.217498) (-3.688889,1.218085) (-3.644444,1.218699) (-3.600000,1.219339) (-3.555556,1.220009) (-3.511111,1.220709) (-3.466667,1.221441) (-3.422222,1.222207) (-3.377778,1.223009) (-3.333333,1.223849) (-3.288889,1.224730) (-3.244444,1.225652) (-3.200000,1.226620) (-3.155556,1.227635) (-3.111111,1.228701) (-3.066667,1.229820) (-3.022222,1.230995) (-2.977778,1.232231) (-2.933333,1.233531) (-2.888889,1.234898) (-2.844444,1.236338) (-2.800000,1.237854) (-2.755556,1.239450) (-2.711111,1.241134) (-2.666667,1.242909) (-2.622222,1.244781) (-2.577778,1.246757) (-2.533333,1.248844) (-2.488889,1.251047) (-2.444444,1.253376) (-2.400000,1.255837) (-2.355556,1.258439) (-2.311111,1.261191) (-2.266667,1.264102) (-2.222222,1.267182) (-2.177778,1.270441) (-2.133333,1.273890) (-2.088889,1.277540) (-2.044444,1.281403) (-2.000000,1.285489) (-1.955556,1.289811) (-1.911111,1.294382) (-1.866667,1.299212) (-1.822222,1.304313) (-1.777778,1.309697) (-1.733333,1.315374) (-1.688889,1.321353) (-1.644444,1.327641) (-1.600000,1.334246) (-1.555556,1.341170) (-1.511111,1.348415) (-1.466667,1.355979) (-1.422222,1.363858) (-1.377778,1.372040) (-1.333333,1.380514) (-1.288889,1.389261) (-1.244444,1.398259) (-1.200000,1.407481) (-1.155556,1.416894) (-1.111111,1.426464) (-1.066667,1.436149) (-1.022222,1.445908) (-0.977778,1.455693) (-0.933333,1.465457) (-0.888889,1.475150) (-0.844444,1.484722) (-0.800000,1.494124) (-0.755556,1.503306) (-0.711111,1.512221) (-0.666667,1.520822) (-0.622222,1.529066) (-0.577778,1.536914) (-0.533333,1.544325) (-0.488889,1.551267) (-0.444444,1.557706) (-0.400000,1.563615) (-0.355556,1.568967) (-0.311111,1.573741) (-0.266667,1.577916) (-0.222222,1.581477) (-0.177778,1.584409) (-0.133333,1.586700) (-0.088889,1.588343) (-0.044444,1.589332) (0.000000,1.589662) (0.044444,1.589332) (0.088889,1.588343) (0.133333,1.586700) (0.177778,1.584409) (0.222222,1.581477) (0.266667,1.577916) (0.311111,1.573741) (0.355556,1.568967) (0.400000,1.563615) (0.444444,1.557706) (0.488889,1.551267) (0.533333,1.544325) (0.577778,1.536914) (0.622222,1.529066) (0.666667,1.520822) (0.711111,1.512221) (0.755556,1.503306) (0.800000,1.494124) (0.844444,1.484722) (0.888889,1.475150) (0.933333,1.465457) (0.977778,1.455693) (1.022222,1.445908) (1.066667,1.436149) (1.111111,1.426464) (1.155556,1.416894) (1.200000,1.407481) (1.244444,1.398259) (1.288889,1.389261) (1.333333,1.380514) (1.377778,1.372040) (1.422222,1.363858) (1.466667,1.355979) (1.511111,1.348415) (1.555556,1.341170) (1.600000,1.334246) (1.644444,1.327641) (1.688889,1.321353) (1.733333,1.315374) (1.777778,1.309697) (1.822222,1.304313) (1.866667,1.299212) (1.911111,1.294382) (1.955556,1.289811) (2.000000,1.285489) (2.044444,1.281403) (2.088889,1.277540) (2.133333,1.273890) (2.177778,1.270441) (2.222222,1.267182) (2.266667,1.264102) (2.311111,1.261191) (2.355556,1.258439) (2.400000,1.255837) (2.444444,1.253376) (2.488889,1.251047) (2.533333,1.248844) (2.577778,1.246757) (2.622222,1.244781) (2.666667,1.242909) (2.711111,1.241134) (2.755556,1.239450) (2.800000,1.237854) (2.844444,1.236338) (2.888889,1.234898) (2.933333,1.233531) (2.977778,1.232231) (3.022222,1.230995) (3.066667,1.229820) (3.111111,1.228701) (3.155556,1.227635) (3.200000,1.226620) (3.244444,1.225652) (3.288889,1.224730) (3.333333,1.223849) (3.377778,1.223009) (3.422222,1.222207) (3.466667,1.221441) (3.511111,1.220709) (3.555556,1.220009) (3.600000,1.219339) (3.644444,1.218699) (3.688889,1.218085) (3.733333,1.217498) (3.777778,1.216935) (3.822222,1.216396) (3.866667,1.215879) (3.911111,1.215383) (3.955556,1.214907) (4.000000,1.214450)};
\draw[Blue,thick,->] (2.444444,1.253376)--(2.711111,1.241134);
\draw[Blue,thick] plot[smooth] coordinates {(-4.000000,-1.214450) (-3.955556,-1.214907) (-3.911111,-1.215383) (-3.866667,-1.215879) (-3.822222,-1.216396) (-3.777778,-1.216935) (-3.733333,-1.217498) (-3.688889,-1.218085) (-3.644444,-1.218699) (-3.600000,-1.219339) (-3.555556,-1.220009) (-3.511111,-1.220709) (-3.466667,-1.221441) (-3.422222,-1.222207) (-3.377778,-1.223009) (-3.333333,-1.223849) (-3.288889,-1.224730) (-3.244444,-1.225652) (-3.200000,-1.226620) (-3.155556,-1.227635) (-3.111111,-1.228701) (-3.066667,-1.229820) (-3.022222,-1.230995) (-2.977778,-1.232231) (-2.933333,-1.233531) (-2.888889,-1.234898) (-2.844444,-1.236338) (-2.800000,-1.237854) (-2.755556,-1.239450) (-2.711111,-1.241134) (-2.666667,-1.242909) (-2.622222,-1.244781) (-2.577778,-1.246757) (-2.533333,-1.248844) (-2.488889,-1.251047) (-2.444444,-1.253376) (-2.400000,-1.255837) (-2.355556,-1.258439) (-2.311111,-1.261191) (-2.266667,-1.264102) (-2.222222,-1.267182) (-2.177778,-1.270441) (-2.133333,-1.273890) (-2.088889,-1.277540) (-2.044444,-1.281403) (-2.000000,-1.285489) (-1.955556,-1.289811) (-1.911111,-1.294382) (-1.866667,-1.299212) (-1.822222,-1.304313) (-1.777778,-1.309697) (-1.733333,-1.315374) (-1.688889,-1.321353) (-1.644444,-1.327641) (-1.600000,-1.334246) (-1.555556,-1.341170) (-1.511111,-1.348415) (-1.466667,-1.355979) (-1.422222,-1.363858) (-1.377778,-1.372040) (-1.333333,-1.380514) (-1.288889,-1.389261) (-1.244444,-1.398259) (-1.200000,-1.407481) (-1.155556,-1.416894) (-1.111111,-1.426464) (-1.066667,-1.436149) (-1.022222,-1.445908) (-0.977778,-1.455693) (-0.933333,-1.465457) (-0.888889,-1.475150) (-0.844444,-1.484722) (-0.800000,-1.494124) (-0.755556,-1.503306) (-0.711111,-1.512221) (-0.666667,-1.520822) (-0.622222,-1.529066) (-0.577778,-1.536914) (-0.533333,-1.544325) (-0.488889,-1.551267) (-0.444444,-1.557706) (-0.400000,-1.563615) (-0.355556,-1.568967) (-0.311111,-1.573741) (-0.266667,-1.577916) (-0.222222,-1.581477) (-0.177778,-1.584409) (-0.133333,-1.586700) (-0.088889,-1.588343) (-0.044444,-1.589332) (0.000000,-1.589662) (0.044444,-1.589332) (0.088889,-1.588343) (0.133333,-1.586700) (0.177778,-1.584409) (0.222222,-1.581477) (0.266667,-1.577916) (0.311111,-1.573741) (0.355556,-1.568967) (0.400000,-1.563615) (0.444444,-1.557706) (0.488889,-1.551267) (0.533333,-1.544325) (0.577778,-1.536914) (0.622222,-1.529066) (0.666667,-1.520822) (0.711111,-1.512221) (0.755556,-1.503306) (0.800000,-1.494124) (0.844444,-1.484722) (0.888889,-1.475150) (0.933333,-1.465457) (0.977778,-1.455693) (1.022222,-1.445908) (1.066667,-1.436149) (1.111111,-1.426464) (1.155556,-1.416894) (1.200000,-1.407481) (1.244444,-1.398259) (1.288889,-1.389261) (1.333333,-1.380514) (1.377778,-1.372040) (1.422222,-1.363858) (1.466667,-1.355979) (1.511111,-1.348415) (1.555556,-1.341170) (1.600000,-1.334246) (1.644444,-1.327641) (1.688889,-1.321353) (1.733333,-1.315374) (1.777778,-1.309697) (1.822222,-1.304313) (1.866667,-1.299212) (1.911111,-1.294382) (1.955556,-1.289811) (2.000000,-1.285489) (2.044444,-1.281403) (2.088889,-1.277540) (2.133333,-1.273890) (2.177778,-1.270441) (2.222222,-1.267182) (2.266667,-1.264102) (2.311111,-1.261191) (2.355556,-1.258439) (2.400000,-1.255837) (2.444444,-1.253376) (2.488889,-1.251047) (2.533333,-1.248844) (2.577778,-1.246757) (2.622222,-1.244781) (2.666667,-1.242909) (2.711111,-1.241134) (2.755556,-1.239450) (2.800000,-1.237854) (2.844444,-1.236338) (2.888889,-1.234898) (2.933333,-1.233531) (2.977778,-1.232231) (3.022222,-1.230995) (3.066667,-1.229820) (3.111111,-1.228701) (3.155556,-1.227635) (3.200000,-1.226620) (3.244444,-1.225652) (3.288889,-1.224730) (3.333333,-1.223849) (3.377778,-1.223009) (3.422222,-1.222207) (3.466667,-1.221441) (3.511111,-1.220709) (3.555556,-1.220009) (3.600000,-1.219339) (3.644444,-1.218699) (3.688889,-1.218085) (3.733333,-1.217498) (3.777778,-1.216935) (3.822222,-1.216396) (3.866667,-1.215879) (3.911111,-1.215383) (3.955556,-1.214907) (4.000000,-1.214450)};
\draw[Blue,thick,->] (2.444444,-1.253376)--(2.711111,-1.241134);
\draw[Blue,thick] plot[smooth] coordinates {(-4.000000,1.768360) (-3.955556,1.768887) (-3.911111,1.769433) (-3.866667,1.769999) (-3.822222,1.770586) (-3.777778,1.771195) (-3.733333,1.771827) (-3.688889,1.772483) (-3.644444,1.773164) (-3.600000,1.773870) (-3.555556,1.774604) (-3.511111,1.775367) (-3.466667,1.776159) (-3.422222,1.776982) (-3.377778,1.777837) (-3.333333,1.778726) (-3.288889,1.779651) (-3.244444,1.780612) (-3.200000,1.781612) (-3.155556,1.782652) (-3.111111,1.783734) (-3.066667,1.784860) (-3.022222,1.786032) (-2.977778,1.787252) (-2.933333,1.788522) (-2.888889,1.789844) (-2.844444,1.791220) (-2.800000,1.792653) (-2.755556,1.794145) (-2.711111,1.795699) (-2.666667,1.797317) (-2.622222,1.799002) (-2.577778,1.800757) (-2.533333,1.802583) (-2.488889,1.804485) (-2.444444,1.806464) (-2.400000,1.808524) (-2.355556,1.810668) (-2.311111,1.812898) (-2.266667,1.815217) (-2.222222,1.817627) (-2.177778,1.820132) (-2.133333,1.822734) (-2.088889,1.825436) (-2.044444,1.828239) (-2.000000,1.831145) (-1.955556,1.834157) (-1.911111,1.837274) (-1.866667,1.840500) (-1.822222,1.843833) (-1.777778,1.847274) (-1.733333,1.850822) (-1.688889,1.854477) (-1.644444,1.858236) (-1.600000,1.862097) (-1.555556,1.866056) (-1.511111,1.870110) (-1.466667,1.874252) (-1.422222,1.878478) (-1.377778,1.882779) (-1.333333,1.887149) (-1.288889,1.891577) (-1.244444,1.896054) (-1.200000,1.900569) (-1.155556,1.905109) (-1.111111,1.909662) (-1.066667,1.914215) (-1.022222,1.918751) (-0.977778,1.923258) (-0.933333,1.927717) (-0.888889,1.932115) (-0.844444,1.936433) (-0.800000,1.940655) (-0.755556,1.944765) (-0.711111,1.948746) (-0.666667,1.952581) (-0.622222,1.956254) (-0.577778,1.959750) (-0.533333,1.963053) (-0.488889,1.966149) (-0.444444,1.969024) (-0.400000,1.971666) (-0.355556,1.974062) (-0.311111,1.976203) (-0.266667,1.978078) (-0.222222,1.979679) (-0.177778,1.980998) (-0.133333,1.982031) (-0.088889,1.982772) (-0.044444,1.983218) (0.000000,1.983367) (0.044444,1.983218) (0.088889,1.982772) (0.133333,1.982031) (0.177778,1.980998) (0.222222,1.979679) (0.266667,1.978078) (0.311111,1.976203) (0.355556,1.974062) (0.400000,1.971666) (0.444444,1.969024) (0.488889,1.966149) (0.533333,1.963053) (0.577778,1.959750) (0.622222,1.956254) (0.666667,1.952581) (0.711111,1.948746) (0.755556,1.944765) (0.800000,1.940655) (0.844444,1.936433) (0.888889,1.932115) (0.933333,1.927717) (0.977778,1.923258) (1.022222,1.918751) (1.066667,1.914215) (1.111111,1.909662) (1.155556,1.905109) (1.200000,1.900569) (1.244444,1.896054) (1.288889,1.891577) (1.333333,1.887149) (1.377778,1.882779) (1.422222,1.878478) (1.466667,1.874252) (1.511111,1.870110) (1.555556,1.866056) (1.600000,1.862097) (1.644444,1.858236) (1.688889,1.854477) (1.733333,1.850822) (1.777778,1.847274) (1.822222,1.843833) (1.866667,1.840500) (1.911111,1.837274) (1.955556,1.834157) (2.000000,1.831145) (2.044444,1.828239) (2.088889,1.825436) (2.133333,1.822734) (2.177778,1.820132) (2.222222,1.817627) (2.266667,1.815217) (2.311111,1.812898) (2.355556,1.810668) (2.400000,1.808524) (2.444444,1.806464) (2.488889,1.804485) (2.533333,1.802583) (2.577778,1.800757) (2.622222,1.799002) (2.666667,1.797317) (2.711111,1.795699) (2.755556,1.794145) (2.800000,1.792653) (2.844444,1.791220) (2.888889,1.789844) (2.933333,1.788522) (2.977778,1.787252) (3.022222,1.786032) (3.066667,1.784860) (3.111111,1.783734) (3.155556,1.782652) (3.200000,1.781612) (3.244444,1.780612) (3.288889,1.779651) (3.333333,1.778726) (3.377778,1.777837) (3.422222,1.776982) (3.466667,1.776159) (3.511111,1.775367) (3.555556,1.774604) (3.600000,1.773870) (3.644444,1.773164) (3.688889,1.772483) (3.733333,1.771827) (3.777778,1.771195) (3.822222,1.770586) (3.866667,1.769999) (3.911111,1.769433) (3.955556,1.768887) (4.000000,1.768360)};
\draw[Blue,thick,->] (2.444444,1.806464)--(2.711111,1.795699);
\draw[Blue,thick] plot[smooth] coordinates {(-4.000000,-1.768360) (-3.955556,-1.768887) (-3.911111,-1.769433) (-3.866667,-1.769999) (-3.822222,-1.770586) (-3.777778,-1.771195) (-3.733333,-1.771827) (-3.688889,-1.772483) (-3.644444,-1.773164) (-3.600000,-1.773870) (-3.555556,-1.774604) (-3.511111,-1.775367) (-3.466667,-1.776159) (-3.422222,-1.776982) (-3.377778,-1.777837) (-3.333333,-1.778726) (-3.288889,-1.779651) (-3.244444,-1.780612) (-3.200000,-1.781612) (-3.155556,-1.782652) (-3.111111,-1.783734) (-3.066667,-1.784860) (-3.022222,-1.786032) (-2.977778,-1.787252) (-2.933333,-1.788522) (-2.888889,-1.789844) (-2.844444,-1.791220) (-2.800000,-1.792653) (-2.755556,-1.794145) (-2.711111,-1.795699) (-2.666667,-1.797317) (-2.622222,-1.799002) (-2.577778,-1.800757) (-2.533333,-1.802583) (-2.488889,-1.804485) (-2.444444,-1.806464) (-2.400000,-1.808524) (-2.355556,-1.810668) (-2.311111,-1.812898) (-2.266667,-1.815217) (-2.222222,-1.817627) (-2.177778,-1.820132) (-2.133333,-1.822734) (-2.088889,-1.825436) (-2.044444,-1.828239) (-2.000000,-1.831145) (-1.955556,-1.834157) (-1.911111,-1.837274) (-1.866667,-1.840500) (-1.822222,-1.843833) (-1.777778,-1.847274) (-1.733333,-1.850822) (-1.688889,-1.854477) (-1.644444,-1.858236) (-1.600000,-1.862097) (-1.555556,-1.866056) (-1.511111,-1.870110) (-1.466667,-1.874252) (-1.422222,-1.878478) (-1.377778,-1.882779) (-1.333333,-1.887149) (-1.288889,-1.891577) (-1.244444,-1.896054) (-1.200000,-1.900569) (-1.155556,-1.905109) (-1.111111,-1.909662) (-1.066667,-1.914215) (-1.022222,-1.918751) (-0.977778,-1.923258) (-0.933333,-1.927717) (-0.888889,-1.932115) (-0.844444,-1.936433) (-0.800000,-1.940655) (-0.755556,-1.944765) (-0.711111,-1.948746) (-0.666667,-1.952581) (-0.622222,-1.956254) (-0.577778,-1.959750) (-0.533333,-1.963053) (-0.488889,-1.966149) (-0.444444,-1.969024) (-0.400000,-1.971666) (-0.355556,-1.974062) (-0.311111,-1.976203) (-0.266667,-1.978078) (-0.222222,-1.979679) (-0.177778,-1.980998) (-0.133333,-1.982031) (-0.088889,-1.982772) (-0.044444,-1.983218) (0.000000,-1.983367) (0.044444,-1.983218) (0.088889,-1.982772) (0.133333,-1.982031) (0.177778,-1.980998) (0.222222,-1.979679) (0.266667,-1.978078) (0.311111,-1.976203) (0.355556,-1.974062) (0.400000,-1.971666) (0.444444,-1.969024) (0.488889,-1.966149) (0.533333,-1.963053) (0.577778,-1.959750) (0.622222,-1.956254) (0.666667,-1.952581) (0.711111,-1.948746) (0.755556,-1.944765) (0.800000,-1.940655) (0.844444,-1.936433) (0.888889,-1.932115) (0.933333,-1.927717) (0.977778,-1.923258) (1.022222,-1.918751) (1.066667,-1.914215) (1.111111,-1.909662) (1.155556,-1.905109) (1.200000,-1.900569) (1.244444,-1.896054) (1.288889,-1.891577) (1.333333,-1.887149) (1.377778,-1.882779) (1.422222,-1.878478) (1.466667,-1.874252) (1.511111,-1.870110) (1.555556,-1.866056) (1.600000,-1.862097) (1.644444,-1.858236) (1.688889,-1.854477) (1.733333,-1.850822) (1.777778,-1.847274) (1.822222,-1.843833) (1.866667,-1.840500) (1.911111,-1.837274) (1.955556,-1.834157) (2.000000,-1.831145) (2.044444,-1.828239) (2.088889,-1.825436) (2.133333,-1.822734) (2.177778,-1.820132) (2.222222,-1.817627) (2.266667,-1.815217) (2.311111,-1.812898) (2.355556,-1.810668) (2.400000,-1.808524) (2.444444,-1.806464) (2.488889,-1.804485) (2.533333,-1.802583) (2.577778,-1.800757) (2.622222,-1.799002) (2.666667,-1.797317) (2.711111,-1.795699) (2.755556,-1.794145) (2.800000,-1.792653) (2.844444,-1.791220) (2.888889,-1.789844) (2.933333,-1.788522) (2.977778,-1.787252) (3.022222,-1.786032) (3.066667,-1.784860) (3.111111,-1.783734) (3.155556,-1.782652) (3.200000,-1.781612) (3.244444,-1.780612) (3.288889,-1.779651) (3.333333,-1.778726) (3.377778,-1.777837) (3.422222,-1.776982) (3.466667,-1.776159) (3.511111,-1.775367) (3.555556,-1.774604) (3.600000,-1.773870) (3.644444,-1.773164) (3.688889,-1.772483) (3.733333,-1.771827) (3.777778,-1.771195) (3.822222,-1.770586) (3.866667,-1.769999) (3.911111,-1.769433) (3.955556,-1.768887) (4.000000,-1.768360)};
\draw[Blue,thick,->] (2.444444,-1.806464)--(2.711111,-1.795699);

 \draw[Blue,thick,->] (-4,0)--(-1.25,0);\draw[Blue,thick,->] (1.25,0)--(4,0);
\end{tikzpicture}\figcap{圆球绕流的子午面流线。不同流线绕过球面，不在球顶或球底合并。}\label{ch13:fig:k10-2-23}\end{minipage}\end{center}}

\newcommand{\StreamtubeDiagram}{%
\begin{center}\begin{minipage}{\linewidth}\centering\begin{tikzpicture}[scale=.95,>=Latex]
 % Center c(t)=(t,0,.1t^2); section basis e_y and (-.2t,0,1)/sqrt(1+.04t^2).
 \foreach \ang in {0,60,120,180,240,300}{
  \draw[Blue,thick,domain=0:5,samples=75,variable=\t]
   plot ({\t+(.70-.055*\t)*(.45*cos(\ang)-.2*\t*sin(\ang)/sqrt(1+.04*\t*\t))},
         {.1*\t*\t+(.70-.055*\t)*(.22*cos(\ang)+sin(\ang)/sqrt(1+.04*\t*\t))});}
 \foreach \t in {0,5}{
  \draw[Ink,thick,domain=0:360,samples=80,variable=\a]
   plot ({\t+(.70-.055*\t)*(.45*cos(\a)-.2*\t*sin(\a)/sqrt(1+.04*\t*\t))},
         {.1*\t*\t+(.70-.055*\t)*(.22*cos(\a)+sin(\a)/sqrt(1+.04*\t*\t))});}
 \foreach \t/\lab in {0.7/1,2.5/2,4.2/3}{
  \draw[->,DeepBlue,very thick] (\t,{.1*\t*\t})--++(.55,{.11*\t});
  \draw[TextGray,thin] (\t,{.1*\t*\t})--(\t,{.1*\t*\t-1.0});
  \node[below,font=\small] at (\t,{.1*\t*\t-1.07}) {$\vect v_{\lab}$};}
 \node[left] at (-.5,0) {$C_1$};\node[left] at (-.5,-.45) {$S_A$};
 \node[right] at (5.55,2.8) {$C_2$};\node[right] at (5.55,2.35) {$S_B$};
\end{tikzpicture}\figcap{闭合截线 $C_1,C_2$ 围出横截面，侧壁由流线构成。箭头沿局部速度方向；图为弯曲流管的透视示意。}\label{ch13:fig:k10-2-24}\end{minipage}\end{center}}

\newcommand{\PipeContinuityDiagram}{%
\begin{center}\begin{minipage}{\linewidth}\centering\begin{tikzpicture}[scale=.74,>=Latex]
 \draw[very thick] (-4,.55)--(-1.5,.55)--(1.2,1.2)--(4,1.2);
 \draw[very thick] (-4,-.55)--(-1.5,-.55)--(1.2,-1.2)--(4,-1.2);
 \draw[vector] (-2.9,0)--(-1.65,0);
 \draw[TextGray,thin] (-2.3,-.12)--(-2.3,-.95);
 \node[below] at (-2.3,-1.0) {$v_A$};
 \draw[vector] (1.65,0)--(2.4,0) node[midway,below=4pt]{$v_B$};
 \draw[DeepBlue] (-3.25,-.55)--(-3.25,.55) node[above=5pt]{$S_A$};
 \draw[DeepBlue] (3.0,-1.2)--(3.0,1.2) node[above=5pt]{$S_B$};
\end{tikzpicture}\figcap{管道的两个横截面及流速方向，$v_B$ 为待求量。}\label{ch13:fig:k10-ex-2-97}\end{minipage}\end{center}}

\newcommand{\UniformPipeDiagram}{%
\begin{center}\begin{minipage}{\linewidth}\centering\begin{tikzpicture}[scale=.8,>=Latex]
 \draw[thick] (-4,.65)--(4,.65) (-4,-.65)--(4,-.65);
 \foreach \xx/\lab in {-2.8/1,2.0/2}{
  \draw[dashed,TextGray] (\xx,-.65)--(\xx,.65) node[above=5pt]{$S$};
  \draw[->,DeepBlue,thick] (\xx+.2,0)--(\xx+1.25,0) node[midway,below=5pt]{$v_{\lab}$};}
\end{tikzpicture}\figcap{同一圆柱管的两个等截面；在一维稳态模型中比较 $Sv_1$ 与 $Sv_2$。}\label{ch13:fig:k10-ex-2-98}\end{minipage}\end{center}}

\newcommand{\CompressiblePipeDiagram}{%
\begin{center}\begin{minipage}{\linewidth}\centering\begin{tikzpicture}[scale=.8,>=Latex]
 \draw[thick] (-4,.6)--(-1.5,.6)--(1.3,1.3)--(4,1.3);
 \draw[thick] (-4,-.6)--(-1.5,-.6)--(1.3,-1.3)--(4,-1.3);
 \foreach \xx/\rr/\lab in {-3/.6/A,3/1.3/B}{
  \draw[DeepBlue,dashed] (\xx,-\rr)--(\xx,\rr) node[above=6pt]{$S_{\lab}$};
  \node[below=6pt] at (\xx,-\rr) {$\rho_{\lab}$};}
 \draw[->,DeepBlue,thick] (-2.65,0)--(-1.65,0);
 \draw[TextGray,thin] (-1.9,.1)--(-1.9,1.05);
 \node[above] at (-1.9,1.1) {$v_A$};
 \draw[->,DeepBlue,thick] (1.55,0)--(2.5,0) node[midway,below=5pt]{$v_B$};
\end{tikzpicture}\figcap{两截面的面积、密度及流速符号，$v_B$ 为待求量。}\label{ch13:fig:k10-ex-2-100}\end{minipage}\end{center}}

\newcommand{\FallingJetDiagram}{%
\begin{center}\begin{minipage}{\linewidth}\centering\begin{tikzpicture}[scale=.85,>=Latex]
 \draw[thick] (-1.8,1)--(-1.8,0)--(-.50,0) (.50,0)--(1.8,0)--(1.8,1);
 \foreach \sgn in {-1,1}{\draw[DeepBlue,thick,domain=0:3.7,samples=90,variable=\s]
  plot ({\sgn*.5/(1+1.2*\s)^.25},{-\s});}
 \draw[TextGray,dashed] (-.75,0)--(1.3,0) (-.75,-3)--(1.3,-3);
 \draw[<->,TextGray] (-1.1,0)--(-1.1,-3) node[midway,left]{$h$};
 \draw[DeepBlue,dashed] (-.5,0)--(.5,0);
 \draw[DeepBlue,thick] ({-.5/4.6^.25},-3)--({.5/4.6^.25},-3);
 \node[above] at (0,.12) {$S$};\node[left] at (-.55,-3.25) {$S(h)$};
 \draw[->,DeepBlue,thick] (1.3,0)--(1.3,-.65) node[right]{$v$};
 \draw[->,DeepBlue,thick] (1.3,-3)--(1.3,-4.1) node[right]{$v(h)$};
\end{tikzpicture}\figcap{高度差 $h$ 的末端就是截面 $S(h)$；两速度箭头分别从对应截面起画。}\label{ch13:fig:k10-ex-2-99}\end{minipage}\end{center}}


\label{ch13:sec:flow-intro}

水和空气没有固定形状，受力后可以流动，统称流体。流体运动极其复杂，本节先讨论稳定流动：空间中任一给定点的速度大小和方向都不随时间改变。稳定并不表示某个流体微元的速度沿途不变；微元运动到不同位置时，会取得当地不同的流速。

流体由大量分子组成，逐一跟踪每个分子的运动通常既不可行也没有必要。稳定流动改为考察“空间中某一点此刻允许多大的速度、指向哪里”：只要这个速度场不随时间变化，后来经过同一点的流体微元就会取得相同速度。某个微元沿途经过不同位置时仍可能不断加速、减速或转弯，所以“稳定”描述的是空间分布不随时间变，并不等于每个微元作匀速直线运动。

水在通常条件下密度变化很小，可理想化为不可压缩流体；空气的密度较易随压强和速度改变，是可压缩流体。对稳定流动，给出空间每点的速度矢量便给出流速分布。沿速度场逐步跟踪微元可得到流线。

不可压缩同样是一种近似：它表示同一流体微元在运动中密度近似不变，并不是说流体“不能改变形状”。气体通常更容易压缩，完整描述稳定气流时除了速度场，还须给出各点密度。后面的连续性方程会分别说明体积流率和质量流率在两种模型中的区别。

\SphereStreamlinesDiagram

在速度场连续光滑且速度非零的区域，流线具有三条基本性质：流线在每一点与当地速度相切；每一点通过一条局部唯一的流线；两条不同流线不能相交。停滞点的速度为零，不能由零矢量指定切线方向，需作为例外单独讨论。

流线可以用逐步推进的方法构造：从一个起点出发，沿当地速度方向前进很短距离，到达新位置后改沿新位置的速度方向继续前进；令每一步趋于无限小，得到的光滑曲线就是流线。图~\ref{ch13:fig:k10-2-23}中的每条线都应在圆球附近平滑转弯而不穿过球体。若两条流线相交，交点处会出现两个切线方向，与稳定速度场在每点只有一个速度矢量矛盾。

取一个与流动方向垂直的面积 $S$，若截面上流速均为 $v$，单位时间通过的体积为 $Sv$。由一圈流线围成的管状区域称为流线管，其侧壁没有流体穿过。

在短时间 $\Delta t$ 内，穿过截面的流体近似占据底面积 $S$、长度 $v\Delta t$ 的柱体，体积为 $Sv\Delta t$，因此体积流率为 $Sv$。流线管的侧壁处处由流线组成，速度与侧壁相切，没有法向分量，所以流体不会从侧壁进出；流入和流出只能通过两端截面。

\StreamtubeDiagram

图~\ref{ch13:fig:k10-2-24}的两截面属于同一流线管。不可压缩流体稳定流过同一流线管时，单位时间通过各截面的体积相等。截面积增大则流速降低，截面积减小则流速增大。

\Needspace{16\baselineskip}\begin{origexamplebox}{不可压缩流体的连续性方程}\label{ch13:ex:k10-2-97}
不可压缩流体在管中稳定流动，每个垂直截面上流速均匀。截面 $A$ 的面积为 $S_A$、流速为 $v_A$；截面 $B$ 的面积为 $S_B$。求 $v_B$。

\PipeContinuityDiagram
\end{origexamplebox}
\Needspace{6\baselineskip}\begin{xesolution}
\ThAFlowVolume

图~\ref{ch13:fig:k10-ex-2-97}标出截面 $S_A,S_B$ 与对应速度。取相同时间 $\Delta t$，通过两截面的体积分别为 $S_Av_A\Delta t$、$S_Bv_B\Delta t$。稳态且不可压缩时两者相等：
\[
S_Av_A=S_Bv_B.
\]
所以
\[
\boxed{v_B=\frac{S_A}{S_B}v_A}.
\]
量纲检查：面积比无量纲，结果仍为速度。
\end{xesolution}

\Needspace{16\baselineskip}\begin{origexamplebox}{等截面管内能否稳定加速}\label{ch13:ex:k10-2-98}
圆柱形管道处处充满不可压缩流体。流体能否在管中作稳定的加速运动？采用沿直管轴线流动、各横截面流速均匀、无支流的一维模型。

\UniformPipeDiagram
\end{origexamplebox}
\Needspace{6\baselineskip}\begin{xesolution}
\ThAUniformFlow

图~\ref{ch13:fig:k10-ex-2-98}中圆柱管的截面积 $S$ 处处相同。连续性方程要求任意两截面
\[
Sv_1=Sv_2,
\]
故 $v_1=v_2$。在“不可压缩、稳定、处处充满且无支流”的条件下，流速不能沿管长增加，因而
\[
\boxed{\text{不能作沿程稳定加速运动}}.
\]
若流体可压缩、管道有泄漏或流动随时间变化，结论可能改变；这些都超出了题设。
若不作一维近似，连续性只保证截面平均流速相同；截面内部仍可发生速度重新分布。因此不能由“等截面”一句话推出任何微元的局部加速度都为零。
\end{xesolution}

\Needspace{16\baselineskip}\begin{origexamplebox}{下落水束的截面积}\label{ch13:ex:k10-2-99}
极大容器底部有面积为 $S$ 的圆孔，水刚离孔时速度为 $v$。每个水分子随后自由落体，水束不散开。求孔下方 $h$ 处水流截面积 $S(h)$。

\FallingJetDiagram
\end{origexamplebox}
\Needspace{6\baselineskip}\begin{xesolution}
\ThAFallingFlow

图~\ref{ch13:fig:k10-ex-2-99}中，$h$ 的下端、$S(h)$ 与 $v(h)$ 标在同一横截面。自由下落的速度由
\[
v(h)^2=v^2+2gh
\]
给出，故 $v(h)=\sqrt{v^2+2gh}$。水不可压缩，体积流率不变：
\[
Sv=S(h)v(h).
\]
因此
\[
\boxed{S(h)=\frac{Sv}{\sqrt{v^2+2gh}}}.
\]
随 $h$ 增大，水束加速而截面收细，乘积 $S(h)v(h)$ 保持为 $Sv$。
\end{xesolution}

对于可压缩流体，稳定流动中真正守恒的是质量流率而非体积流率。若没有物质产生或消失，对任意闭合曲面，单位时间流入的质量必须等于流出的质量，否则曲面内的质量会随时间积累或减少。这就是连续性条件。

若截面上密度和速度近似均匀，短时间通过的质量为 $\rho Sv\Delta t$，质量流率为 $\rho Sv$。不可压缩情形中同一流线管内 $\rho$ 不变，才可约去密度得到 $Sv$ 恒定；可压缩流体则可能用密度变化补偿截面积或速度变化。连续性只表达物质守恒，不自动给出压强和速度的全部关系，解决更完整的流体问题还需要后续动力学方程。

连续性方程的适用对象必须选成同一股稳定流。若管道存在支路，应把所有入口和出口的质量流率带符号相加，而不能只随意比较两个截面；若流动随时间变化，控制体内还可能发生质量积累。对题图的单一流线管，侧壁无穿流、内部又无源汇，所以任意两个横截面的流率才相等。

自由水束的两个例题展示了同一守恒条件的两种实现方式。真实水近似不可压缩，密度保持不变，下落加速后必须缩小截面积；若人为规定截面积不变而又让各微元自由加速，连续性就要求密度降低。二者不能同时把截面积和密度都固定，否则 $v$ 增大将使质量流率凭空增大，违反稳定流动的物质守恒。

定常流动中，流线与单个微元的轨迹重合；非定常流动中一般不重合。若速度场随时间改变，某一时刻画出的流线仍处处与当时速度相切，但一个微元随后看到的速度场已经变化，它实际走出的迹线可能与这条瞬时流线不同。本节始终假定稳定流动，因此可以用流线同时表示当地速度方向和微元运动轨迹；超出这一条件时应区分两者。

截面若不垂直于速度，穿过它的体积流率应使用法向分量 $v_n$，即 $Sv_n$，而不是总速率 $Sv$。若截面上速度或密度不均匀，还需把 $\rho\vect v\cdot\vect n$ 在整个截面上积分。本节例题把截面选为垂直于流动且各处均匀，才得到简洁的 $Sv$ 或 $\rho Sv$；这些条件也是公式的一部分，不能从图中省略。

\Needspace{16\baselineskip}\begin{origexamplebox}{可压缩流体的质量连续性}\label{ch13:ex:k10-2-100}
流体在管道中稳定流动，且同一垂直截面上性质均匀。$A$ 处截面积、密度、流速分别为 $S_A,\rho_A,v_A$；$B$ 处为 $S_B,\rho_B,v_B$。求 $v_B$。

\CompressiblePipeDiagram
\end{origexamplebox}
\Needspace{6\baselineskip}\begin{xesolution}
\ThAFlowMass

图~\ref{ch13:fig:k10-ex-2-100}将两侧密度分别标为 $\rho_A,\rho_B$。单位时间通过截面的质量为 $\rho Sv$，不能只比较体积。质量连续性给出
\[
\rho_AS_Av_A=\rho_BS_Bv_B,
\]
所以
\[
\boxed{v_B=\frac{\rho_AS_A}{\rho_BS_B}v_A}.
\]
当 $\rho_A=\rho_B$ 时退化为例~\ref{ch13:ex:k10-2-97}的不可压缩结果。
\end{xesolution}

\Needspace{16\baselineskip}\begin{origexamplebox}{等截面自由落体液柱的密度}\label{ch13:ex:k10-2-101}
极大容器中液体初始密度为 $\rho$，从面积为 $S$ 的孔以速度 $v$ 流出。忽略分子间相互作用和表面张力，设下落过程中水束截面积始终为 $S$，每个分子作自由落体。求密度 $\rho(h)$。

\ThQDilutingColumn
\end{origexamplebox}
\Needspace{6\baselineskip}\begin{xesolution}
\ThADilutingColumn

下落 $h$ 后速度仍为
\[
v(h)=\sqrt{v^2+2gh}.
\]
稳态质量流率在孔口和高度 $h$ 处相等：
\[
\rho Sv=\rho(h)S v(h).
\]
截面积约去，得到
\[
\boxed{\rho(h)=\frac{\rho v}{\sqrt{v^2+2gh}}}.
\]
这里人为固定了截面积，因此加速必须由密度降低来满足质量连续性；这正与例~\ref{ch13:ex:k10-2-99}中“密度不变、截面积减小”互为对照。
\end{xesolution}
}

\newcommand{\FluidBernoulliDiagram}{%
\begin{center}\begin{minipage}{\linewidth}\centering\begin{tikzpicture}[scale=.9,>=Latex]
 \draw[->,TextGray] (-2,-2.6)--(-2,3.4) node[above]{$z$};
 \draw[TextGray,dashed] (-2,-2.3)--(7.8,-2.3);
 \begin{scope}[rotate=16]
  \draw[thick] (0,.85)--(6,.5) (0,-.85)--(6,-.5);
  \draw[TextGray,dashed] (0,0)--(6,0);
  \draw[DeepBlue,thick] (0,0) ellipse[x radius=.20,y radius=.85];
  \draw[DeepBlue,thick] (6,0) ellipse[x radius=.117647,y radius=.5];
  \draw[TextGray,dashed] (.35,0) ellipse[x radius=.195294,y radius=.830];
  \draw[TextGray,dashed] (7.0115,0) ellipse[x radius=.117647,y radius=.5];
  \draw[->,DeepBlue,thick] (-1.1,0)--(0,0);
  \draw[->,DeepBlue,thick] (7.4,0)--(6,0);
  \node[above] at (-.9,.15) {$p_1A_1$};\node[above right] at (7.0,.6) {$p_2A_2$};
  \draw[TextGray,thin] (0,-.87)--(0,-1.28) (.35,-.87)--(.35,-1.28);
  \draw[->,TextGray] (-.55,-1.15)--(0,-1.15);
  \draw[->,TextGray] (.90,-1.15)--(.35,-1.15);
  \node[below] at (.175,-1.3) {$\dd s_1$};
  \draw[<->,TextGray] (6,-.85)--(7.0115,-.85);
  \node[below] at (6.5,-1.0) {$\dd s_2$};
  \node[above] at (0,1.02) {$A_1,\,v_1$};\node[above] at (5.9,.72) {$A_2,\,v_2$};
  \coordinate (FluidA) at (0,0);\coordinate (FluidB) at (6,0);
 \end{scope}
 \draw[<->,TextGray] (2,-2.3)--(2,0) node[midway,right]{$z_1$};
 \draw[TextGray,dotted] (FluidA)--(2,0);
 \draw[<->,TextGray] (5,-2.3)--(5,{6*sin(16)}) node[midway,left]{$z_2$};
 \draw[TextGray,dotted] (FluidB)--(5,{6*sin(16)});
\end{tikzpicture}\figcap{流管两端压力做功。相同时间内移入、移出的质量相等；高度从同一水平基准面量起。}\label{ch13:fig:bernoulli}
\end{minipage}\end{center}}

\newcommand{\FluidShearDiagram}{%
\begin{center}\begin{minipage}{\linewidth}\centering\begin{tikzpicture}[scale=.85,>=Latex]
 \foreach \yy in {0,1.8}{\draw[Ink,thick] (0,\yy)--(4.2,\yy)--(5.15,\yy+.55)--(.95,\yy+.55)--cycle;}
 \foreach \yy in {.15,.60,1.05,1.50}{\draw[->,DeepBlue] (1,\yy)--({1.65+.65*\yy},\yy);}
 \draw[->,DeepBlue,very thick] (2,2.08)--(4.25,2.08);
 \node[right] at (5.5,2.08) {$v_x+\dd v_x$};
 \draw[->,DeepBlue,very thick] (2,.27)--(3.4,.27) node[right=4pt]{$v_x$};
 \draw[<->,TextGray] (-.45,0)--(-.45,1.8) node[midway,left]{$\dd y$};
 \draw[->,Ink] (-1.3,0)--(-1.3,2.7) node[above]{$y$};
 \draw[->,Ink] (0,-.8)--(2,-.8) node[right]{$x$};
 \node[above] at (1.6,2.5) {上层对下层：向右拉};
 \node[below] at (3.1,-1.05) {$\tau_{xy}=\dd F_x/\dd S=\eta\,\dd v_x/\dd y$};
\end{tikzpicture}\figcap{相距 $\dd y$ 的流层与切向速度梯度。上层流得较快，黏性作用减小两层速度差；反作用力方向相反。}\label{ch13:fig:shear}
\end{minipage}\end{center}}

\newcommand{\FluidOutletDiagram}{%
\begin{center}\begin{minipage}{\linewidth}\centering\begin{tikzpicture}[scale=.78,>=Latex]
 \fill[Blue!12] (0,0) rectangle (4,3);
 \draw[thick] (0,3.6)--(0,0)--(4,0)--(4,.7) (4,1.05)--(4,3.6);
 \draw[Blue,thick] (0,3)--(4,3);
 \draw[Blue,thick] (4,.7)--(5.5,.5) (4,1.05)--(5.5,.85);
 \fill (1,3) circle(1.6pt) node[above]{$A$};
 \fill (4,.875) circle(1.6pt);
 \draw[TextGray,thin] (4,.875)--(3.4,1.3);
 \node[above left] at (3.4,1.3) {$B$};
 \draw[->,DeepBlue,thick] (4.7,.67)--(5.6,.67) node[right]{$v$};
 \draw[<->,TextGray] (-.6,.875)--(-.6,3) node[midway,left]{$h$};
 \draw[dotted,TextGray] (-.6,.875)--(3.7,.875);
 \node[above] at (2,4.0) {$p_A=p_B=p_0$};
\end{tikzpicture}\figcap{取自由液面 $A$ 与小孔出口 $B$ 比较压强、速度和高度。}\label{ch13:fig:outlet}\end{minipage}\end{center}}

\newcommand{\FluidPlateDiagram}{%
\begin{center}\begin{minipage}{\linewidth}\centering\begin{tikzpicture}[scale=.85,>=Latex]
 \draw[Ink,very thick] (4,-2)--(4,2);
 \draw[Blue,thick] (0,-.25)--(3.5,-.25)..controls(3.8,-.25)and(3.8,-1)..(3.8,-1.65);
 \draw[Blue,thick] (0,.25)--(3.5,.25)..controls(3.8,.25)and(3.8,1)..(3.8,1.65);
 \draw[->,DeepBlue,thick] (.5,0)--(2,0) node[midway,above=8pt]{$v$};
 \draw[->,DeepBlue,thick] (3.65,.7)--(3.65,1.65);
 \draw[->,DeepBlue,thick] (3.65,-.7)--(3.65,-1.65);
 \draw[dashed,TextGray] (2.3,-1.75) rectangle (4.3,1.75);
 \draw[->,Ink,thick] (4.7,0)--(6,0) node[right]{$F_{\text{水对板}}$};
 \node[left] at (0,0) {$A$};\node[above] at (3.3,2.2) {虚线为控制体};
 \draw[->,TextGray] (0,-1.6)--(1.4,-1.6) node[right]{$x$};
\end{tikzpicture}\figcap{来流有正向 $x$ 动量，流出水沿板面，法向动量为零。}\label{ch13:fig:jet-plate}\end{minipage}\end{center}}

\newcommand{\FluidSettlingDiagram}{%
\begin{center}\begin{minipage}{\linewidth}\centering\begin{tikzpicture}[scale=.9,>=Latex]
 \draw[thick] (0,0) circle(.62);
 \draw[->,Ink,thick] (0,-.62)--(0,-2) node[right]{$mg$};
 \draw[->,DeepBlue,thick] (-.26,{sqrt(.62*.62-.26*.26)})--(-.26,1.6) node[above left]{$F_{\!B}$};
 \draw[->,DeepBlue,thick] (.26,{sqrt(.62*.62-.26*.26)})--(.26,2.2) node[above right]{$6\pi\eta Rv_t$};
 \draw[->,TextGray,thick] (1.8,.6)--(1.8,-.7) node[right]{$v_t$};
 \node[right] at (2.2,1.6) {$F_{\!B}=\frac43\pi R^3\rho_f g$};
\end{tikzpicture}\figcap{下沉小球的受力。终端状态速度向下，浮力和黏性阻力均向上。}\label{ch13:fig:settling}\end{minipage}\end{center}}

\section{流体运动学与物质导数}\label{ch13:sec:material}

流体由大量不断更换邻居的质元组成。拉格朗日描述（Lagrangian description）给每个质元贴上初始标签 $\vect a$，记录其位置 $\vect r(\vect a,t)$；欧拉描述（Eulerian description）则在空间各点记录速度、密度和压强场
\begin{equation}
 \vect v=\vect v(\vect r,t),\qquad \rho=\rho(\vect r,t),\qquad p=p(\vect r,t).
\end{equation}
两种描述说的是同一群质元。固定初始标签 $\vect a$，对其位置求时间导数，就得到该质元此刻的速度：
\[
 \vect v\bigl(\vect r(\vect a,t),t\bigr)
 =\left.\frac{\partial\vect r(\vect a,t)}{\partial t}\right|_{\vect a}.
\]
左边先找到质元当前所在的位置，再读取空间速度场；右边始终跟踪同一个质元。标签 $\vect a$ 不能在求导时随意改变。
竞赛流体题通常使用欧拉描述，但研究某个质元所经历的变化时必须沿它的轨迹求导。若标量场为 $f(\vect r,t)$，物质导数（material derivative）为
\begin{equation}
 \frac{\mathrm D f}{\mathrm Dt}
 =\frac{\partial f}{\partial t}+\vect v\cdot\nabla f.
\end{equation}
这是多元函数链式法则的直接结果。把沿质元轨迹读到的量写成
$f(x(t),y(t),z(t),t)$，则
\[
 \frac{\dd f}{\dd t}=\frac{\partial f}{\partial t}
 +\frac{\partial f}{\partial x}\frac{\dd x}{\dd t}
 +\frac{\partial f}{\partial y}\frac{\dd y}{\dd t}
 +\frac{\partial f}{\partial z}\frac{\dd z}{\dd t}.
\]
代入 $\dd x/\dd t=v_x$ 等关系，后三项才合写成 $\vect v\cdot\nabla f$。
第一项描述固定位置随时间的变化，后三项描述质元移动到不同位置而读到的变化；二者都不能遗漏。
因此流体质元的加速度不是单纯的 $\partial\vect v/\partial t$，而是
\begin{equation}
 \vect a=\frac{\mathrm D\vect v}{\mathrm Dt}
 =\frac{\partial\vect v}{\partial t}+(\vect v\cdot\nabla)\vect v.
\end{equation}
第二项是对流加速度：即使某点的速度场不随时间改变，质元进入速度不同的区域时仍可有加速度。

迹线是一个质元随时间走过的轨迹，满足 $\dd\vect r/\dd t=\vect v(\vect r,t)$；流线是在同一时刻处处与速度相切的曲线。定常流动中二者重合；非定常流动中一般不重合，非定常流动时不能把某一时刻画出的流线当成质元未来轨迹。
对固定时刻 $t_0$，流线可用参数 $s$ 表示为 $\dd\vect r/\dd s=\lambda(s)\vect v(\vect r,t_0)$。
在相应分量非零处，也常写成
\[
 \frac{\dd x}{v_x(x,y,z,t_0)}=\frac{\dd y}{v_y(x,y,z,t_0)}
 =\frac{\dd z}{v_z(x,y,z,t_0)}.
\]
某分量为零时应使用参数形式，不能除以零。前面的流管图中，$C_1,C_2$ 之间的侧壁由同一组流线组成，速度沿侧壁而不穿过它。

为比较相邻质元，令两位置相差 $\delta\vect r$，作一阶展开：
\[
 v_i(\vect r+\delta\vect r)=v_i(\vect r)
 +\sum_j\frac{\partial v_i}{\partial x_j}\delta x_j+O(|\delta\vect r|^2).
\]
第一项给出共同平动。剩余的速度梯度可分解成对称和反对称两部分：
\begin{equation}
 \frac{\partial v_i}{\partial x_j}
 =D_{ij}+\Omega_{ij},
 \quad
 D_{ij}=\frac12\left(\frac{\partial v_i}{\partial x_j}+\frac{\partial v_j}{\partial x_i}\right),
 \quad
 \Omega_{ij}=\frac12\left(\frac{\partial v_i}{\partial x_j}-\frac{\partial v_j}{\partial x_i}\right).
\end{equation}
$D_{ij}$ 描述局部拉伸和剪切速率，$\Omega_{ij}$ 描述刚体式转动；涡量 $\vect\omega=\nabla\times\vect v$，局部流体微团的角速度为 $\vect\omega/2$。不可压缩条件是 $\nabla\cdot\vect v=\operatorname{tr}\mathbf D=0$，表示质元体积不随运动改变，并不等价于“速度处处相同”。
具体地，$\boldsymbol\Omega\,\delta\vect r=(\vect\omega/2)\times\delta\vect r$，所以它只改变相邻质元连线的方向。
连线长度平方的变化率则为
\[
 \frac{\dd}{\dd t}|\delta\vect r|^2
 =2\,\delta\vect r\cdot\mathbf D\,\delta\vect r,
\]
反对称部分的贡献为零。故 $\mathbf D$ 才描述变形快慢；它的量纲是 $\mathrm{s^{-1}}$，不同于弹性体中无量纲的应变。

\section{静止流体}\label{ch13:sec:statics}

静止流体不能承受持续剪切应力：一旦持续受切向力，普通流体会不断变形，而不是保持静止。
任取内部一个面，其表面力只能沿法线；如果不同方向的正应力不等，换一个倾斜截面便会出现切向分量，与静止矛盾。
所以静止流体的应力各向同性。取拉伸应力为正，应力张量为
\[
 \boldsymbol\sigma=-p\mathbf I.
\]
体积 $V$ 内流体所受周围流体的压力合力为
\[
 \vect F_p=-\oint_{\partial V}p\vect n\,\dd S
 =-\int_V\nabla p\,\dd V,
\]
其中 $\vect n$ 是向外单位法向，第二步对各分量使用高斯定理。因此单位体积的压力合力为 $-\nabla p$。
若外体积力密度为 $\vect f$，静力平衡要求 $\nabla p=\vect f$。这还意味着在光滑单连通区域内
$\nabla\times\vect f=0$；注意 $\vect f$ 是单位体积的力，不能与单位质量外力混淆。
在重力场中 $\vect f=\rho\vect g$，取竖直向上为 $z$，薄层受力平衡给出
\begin{equation}
 \frac{\dd p}{\dd z}=-\rho g.
\end{equation}
不可压缩液体密度恒定时，深度增加 $h$，压强增加 $\rho gh$。同一连通静止液体中，同一水平面压强相等；
平直界面或忽略表面张力时，不同液体交界面两侧压强连续；曲面界面还可能存在后文讨论的拉普拉斯压差。
物体浸入液体后，液体对其表面的压力合力等于把相同体积液体移走所失去的压力合力：
\[
 \vect F_B=-\int_{V_{\text{排}}}\rho\vect g\,\dd V.
\]
故浮力等于排开流体的重量、方向与重力相反。密度均匀时，其作用线通过排开体积的几何中心；密度不均匀时应取排开流体的质量中心。

\subsection{从一小柱液体看压强随高度的变化}
压强不是向下的力，而是作用在所选小面上的法向力除以面积。液体内部同一点的压强是标量，受压面的朝向不同，压力方向随之改变。因此不能因为重力向下，就把侧壁受到的液体压力也画成向下。

取截面积为 $S$、底面高度为 $z$、顶面高度为 $z+\dd z$ 的竖直小柱。周围液体在底面给它向上的力 $p(z)S$，在顶面给它向下的力 $p(z+\dd z)S$；侧面压力没有竖直分量。向上为正，平衡式为
\[
 p(z)S-p(z+\dd z)S-\rho gS\dd z=0.
\]
减去相同项、除以 $S\dd z$ 并取极限，才得到 $\dd p/\dd z=-\rho g$。负号表示向上压强降低；若改用向下的深度 $h$，则 $\dd p/\dd h=\rho g$。对密度恒定的液体，从自由面 $p=p_a$ 积分到深度 $h$，有
\[
 p(h)-p_a=\int_0^h\rho g\,\dd h'=\rho gh.
\]
$p$ 是绝对压强，$p-p_a$ 是相对大气的表压。自由面与大气相通并不意味着 $p=0$，而是表压为零。比较受压面两侧的合力时，才可能把相同的大气压贡献抵消。

\subsection{不同液层与受压面的合力}
密度分层时，应在每层分别累加压强增量。在同一连通静止液体中取水平小柱，水平外体积力为零，压强差必须为零；这就是等高处等压的力学依据。若两点属于互不连通的液体，或者有水平加速度，就不能仅凭高度相同认定压强相等。

\Needspace{16\baselineskip}\begin{origexamplebox}{两层液体中的侧壁压力}\label{ch13:ex:teaching-layered}
一个敞口容器中，上层液体密度为 $\rho_1$、厚度为 $h_1$，下层密度为 $\rho_2>\rho_1$、厚度为 $h_2$。竖直侧壁宽度为 $b$，外侧为大气。求液体作用在侧壁下层部分上的净压力及其作用点深度，忽略界面张力。

\ThQLayers
\end{origexamplebox}
\Needspace{6\baselineskip}\begin{xesolution}
\ThALayers

从两液体界面向下量 $x$，$0\le x\le h_2$。上层已贡献表压 $\rho_1gh_1$，下层再贡献 $\rho_2gx$，故
\[
 \Delta p(x)=\rho_1gh_1+\rho_2gx,
 \qquad \dd F=\Delta p(x)b\,\dd x.
\]
每一水平窄条所受净压力均沿侧壁法线向外，因而可以直接相加其大小：
\[
 F=bg\left(\rho_1h_1h_2+\frac12\rho_2h_2^2\right).
\]
求作用点不能直接取几何中心，因为下方压强更大。设合力位于界面下方 $x_c$，对界面处水平轴取矩，有
\[
 Fx_c=\int_0^{h_2}x\,\dd F
 =bg\left(\frac12\rho_1h_1h_2^2+\frac13\rho_2h_2^3\right),
\]
所以
\[
 x_c=\frac{\rho_1h_1h_2^2/2+\rho_2h_2^3/3}
 {\rho_1h_1h_2+\rho_2h_2^2/2}.
\]
距自由面的深度是 $h_1+x_c$。没有上层时 $x_c=2h_2/3$；若上层压力远大于下层压强变化，分布接近均匀，$x_c\to h_2/2$。两极限分别对应三角形和近似矩形压强分布，且 $\rho gbh^2$ 的量纲确为力。
\end{xesolution}

\subsection{浮力来自压力差，稳定性还取决于力矩}
设想把浸没物体占据的空间重新填满周围液体。静止时，这部分液体的重量被周围压力合力平衡；只要边界上的压强分布相同，把它换成物体后，压力合力也不变。这一替换论证给出浮力，而不是另外假设一种向上的力。
对于跨越两层液体的物体，应分别计算 $\rho_1gV_1+\rho_2gV_2$，不能把总排水体积都乘同一个密度。浮力与重力相等只保证竖直合力为零；物体稍微倾斜后是否产生回复力矩，还需比较重心与浮力作用线，不能仅据“浮着”判断稳定。

\Needspace{16\baselineskip}\begin{origexamplebox}{浮体的竖直小振动}\label{ch13:ex:teaching-buoyancy}
质量为 $m$、水平截面积恒为 $S$ 的直立柱体，在密度 $\rho$ 的液体中平衡漂浮。只允许竖直运动，忽略流体附加惯性、阻力和表面张力，且小振动中底面始终浸没、顶面始终露出。求平衡浸深及小振动周期。

\ThQFloat
\end{origexamplebox}
\Needspace{6\baselineskip}\begin{xesolution}
\ThAFloat

平衡时 $mg=\rho gSh_0$，所以 $h_0=m/(\rho S)$。取柱体向下的位移 $x$ 为正；此时浸深为 $h_0+x$，浮力向上且增大：
\[
 m\ddot x=mg-\rho gS(h_0+x)=-\rho gSx.
\]
平衡项消去后，剩下的正是指向平衡位置的回复力。因此
\[
 \omega^2=\frac{\rho gS}{m},\qquad
 T=2\pi\sqrt{\frac{m}{\rho gS}}.
\]
若 $x<0$，浮力减少，合力向下，也把柱体推回平衡。$\rho gS$ 的单位为 $\mathrm{N/m}$，相当于等效劲度系数。真实液体随物体一起加速会增加有效惯性，故本题结果依赖题设的忽略条件；也不能把它用于已经完全浸没、排液体积不随 $x$ 变化的柱体。
\end{xesolution}

\section{连续性方程与伯努利方程}\label{ch13:sec:bernoulli}

一般质量守恒写成
\begin{equation}
 \frac{\partial\rho}{\partial t}+\nabla\cdot(\rho\vect v)=0.
\end{equation}
对固定控制体，质量的增加率等于净流入率：
\[
 \frac{\dd}{\dd t}\int_V\rho\,\dd V
 =-\oint_{\partial V}\rho\vect v\cdot\vect n\,\dd S.
\]
把右边化成体积分，利用控制体可以任意选取，才得到上面的局部式。
再展开 $\nabla\cdot(\rho\vect v)=\vect v\cdot\nabla\rho+\rho\nabla\cdot\vect v$，可写成
\[
 \boxed{\frac{\mathrm D\rho}{\mathrm Dt}+\rho\nabla\cdot\vect v=0}.
\]
这说明散度正负控制质元体积的胀缩；质量不变时密度随之降低或升高。不可压缩流体满足
$\mathrm D\rho/\mathrm Dt=0$，因而 $\nabla\cdot\vect v=0$；不同质元原本可以有不同密度。
对稳定的一维流管，质量流率 $\dot m=\rho Av$ 沿程不变；不可压缩时化为 $Av=\text{常量}$。

无黏、不可压缩流体沿同一条流线作稳定流动时，欧拉方程（Euler equation）沿流线积分得到
\begin{equation}
 p+\frac12\rho v^2+\rho gz=\text{常量}.
\end{equation}
三项分别是单位体积的压强能、动能和重力势能。伯努利方程（Bernoulli equation）不能跨越泵、涡轮、强黏性耗散区随意使用，
也不能在旋涡流中默认所有流线具有同一常量。
下面把推导写开。对无黏质元应用牛顿定律，压力合力密度为 $-\nabla p$，故
\[
 \rho\frac{\mathrm D\vect v}{\mathrm Dt}=-\nabla p+\rho\vect g.
\]
在定常流中，沿流线取弧长 $s$，切向加速度为 $v\,\dd v/\dd s$。沿切向投影得到
\[
 \rho v\frac{\dd v}{\dd s}=-\frac{\dd p}{\dd s}-\rho g\frac{\dd z}{\dd s},
 \qquad \dd(v^2/2)+\frac{\dd p}{\rho}+g\,\dd z=0.
\]
密度为常数时逐项积分，就是伯努利方程。
\FluidBernoulliDiagram
图~\ref{ch13:fig:bernoulli}也给出功—能解释。相同短时间内，两端移过的质量相等：
$\dd m=\rho A_1\dd s_1=\rho A_2\dd s_2$。
两端压力总功为
\[
 \dd W_p=p_1A_1\dd s_1-p_2A_2\dd s_2
 =\dd m\left(\frac{p_1}{\rho}-\frac{p_2}{\rho}\right).
\]
它等于动能与重力势能的增加
$\dd m[(v_2^2-v_1^2)/2+g(z_2-z_1)]$。
出口压力做负功，正是式中压强项不能漏掉的原因。

更一般地，对定常、无黏、正压流体，若密度可写成压强的单值函数 $\rho=\rho(p)$，沿流线有
\begin{equation}
 \frac{v^2}{2}+gz+\int^p\frac{\dd p'}{\rho(p')}=\text{常量}.
\end{equation}
不可压缩时 $\rho$ 为常量，积分项化为 $p/\rho$；可压缩气流中密度和内能通常与温度耦合，不能仍把 $p+\rho v^2/2+\rho gz$ 当作普适守恒式。
若是定常、无黏、绝热且没有激波等不可逆过程的流动，单位质量内能 $u$ 满足
\[
 \dd u=-p\,\dd(1/\rho)=\frac{p}{\rho^2}\dd\rho.
\]
定义单位质量焓 $h_{\rm sp}=u+p/\rho$，则
\[
 \dd h_{\rm sp}=\dd u+\frac{\dd p}{\rho}-\frac p{\rho^2}\dd\rho
 =\frac{\dd p}{\rho}.
\]
代回沿流线的动力学式，得到
\[
 \boxed{h_{\rm sp}+\frac{v^2}{2}+gz=\text{沿该流线的常量}}.
\]
这里焓含有内能，不能只保留 $p/\rho$。例如理想气体的内能随温度改变，气流加速可能伴随温度降低。
受热驱动的自然对流还涉及密度变化和浮力；液体密度变化虽小，也不能一概删去其浮力效应。
常用近似是在质量守恒中取无散速度场，而在重力项中保留小密度差；这与“所有动力学都与温度无关”不同。

\Needspace{16\baselineskip}\begin{origexamplebox}{小孔流出速度}\label{ch13:ex:outlet}
大水箱自由液面到侧壁小孔的高度差为 $h$，水箱横截面远大于小孔。忽略黏性，求流出速度。

\ThQOutlet
\end{origexamplebox}
\Needspace{6\baselineskip}\begin{xesolution}
在自由液面与小孔出口之间使用伯努利方程。两点都与大气相通，压强项相消；大水箱自由液面速度近似为零：
\FluidOutletDiagram
图~\ref{ch13:fig:outlet}中 $z_A-z_B=h$，$v_A\simeq0$，于是
\[
 p_0+\rho gh=p_0+\frac12\rho v^2.
\]
所以
\[
 \boxed{v=\sqrt{2gh}}.
\]
这与物体自由下落高度 $h$ 的速度相同，但流体微元的真实轨迹和压强变化并不等同于一个孤立质点自由落体。
\end{xesolution}

\section{动量流与流体作用力}\label{ch13:sec:momentum}

稳定流动中，控制体所受外力等于流出动量流率减去流入动量流率：
\begin{equation}
 \sum\vect F=\sum_{\rm out}\dot m\vect v-
 \sum_{\rm in}\dot m\vect v.
\end{equation}
这里的外力包括压强、重力和固体壁面对流体的力。求流体对管壁的力时，最后还要用牛顿第三定律反向。

\Needspace{16\baselineskip}\begin{origexamplebox}{水射流正撞平板}\label{ch13:ex:jet-plate}
密度为 $\rho$、截面积为 $A$、速度为 $v$ 的稳定水射流垂直撞击静止大平板，撞后水沿板面向四周流走。忽略重力，求水对板的平均力。

\ThQJet
\end{origexamplebox}
\Needspace{6\baselineskip}\begin{xesolution}
取法向为 $x$ 轴。质量流率 $\dot m=\rho Av$。流入法向速度为 $v$，流出水流沿板面，法向速度为零。
\FluidPlateDiagram
图~\ref{ch13:fig:jet-plate}中虚线包围转向区；图右箭头是水对板的力，列控制体方程时应使用它的反作用力。
平板对水的力为
\[
 F_{\text{板对水},x}=\dot m(0-v)=-\rho Av^2.
\]
所以水对板的力方向沿来流方向，大小
\[
 \boxed{F=\rho Av^2}.
\]
量纲为 $\mathrm{kg\,m^{-3}}\cdot\mathrm{m^2}\cdot\mathrm{m^2s^{-2}}=\mathrm N$。
\end{xesolution}

\subsection{为什么稳态流动仍然需要外力}
稳态表示固定位置的流速不随时间变化，不表示经过这里的每一份流体都没有加速度。图~\ref{ch13:fig:jet-plate} 中，虚线区域内的水总量可以不变，但新进入的水不断丢失法向动量，这就需要平板持续施力。
在短时间 $\dd t$ 内，流入质量为 $\dot m\dd t$，流入动量为 $\dot m\vect v_{\rm in}\dd t$；流出端同理。稳态时控制体内的动量不积累，所以外力冲量必须补足出入动量之差。除以 $\dd t$ 得到上面的动量流方程。多出口时必须先按矢量相加，不能直接相加速率。

对管内流动，在入口截面，外部流体把所研究流体推入控制体；在出口截面，外部流体阻碍它流出。两个压力方向均指向控制体内部。若求管道与环境之间的净作用力，使用表压时还须把管外大气压的抵消计入，不能只把绝对压强随意置零。

\Needspace{16\baselineskip}\begin{origexamplebox}{射流推动匀速后退的平板}\label{ch13:ex:teaching-moving-plate}
沿水平 $+x$ 方向的细水射流速度为 $v$、截面积为 $A$、密度为 $\rho$。一块足够大的平板沿同方向以恒速 $u$ 后退，$0\le u<v$。水撞板后沿板面流走，相对平板没有法向速度。忽略重力及空气作用，求水对板的力和传给板的功率。

\ThQMovingPlate
\end{origexamplebox}
\Needspace{6\baselineskip}\begin{xesolution}
\ThAMovingPlate

以匀速平板为参考系，它仍是惯性系。对应图~\ref{ch13:fig:jet-plate} 的控制体随板平移，入射相对速度为 $v-u$，故每秒真正撞板的质量为
\[
 \dot m_{\rm hit}=\rho A(v-u).
\]
这不是固定截面的 $\rho Av$：一部分水还在追赶平板。水相对板的法向速度从 $v-u$ 变成零，每单位质量丢失法向动量 $v-u$。由反作用关系，板受到沿 $+x$ 的力
\[
 F=\dot m_{\rm hit}(v-u)=\rho A(v-u)^2.
\]
也可在地面系看同一份水，撞前法向速度为 $v$，撞后为 $u$，其动量变化给出同一个结果。力对平板的功率必须使用地面系板速：
\[
 P=Fu=\rho Au(v-u)^2.
\]
在 $0<u<v$ 内，$\dd P/\dd u=\rho A(v-u)(v-3u)$，故最大功率出现在 $u=v/3$，大小为 $4\rho Av^3/27$。$u=0$ 时有力却不向静止板传机械功；$u\to v$ 时几乎没有水追上平板，力和功率均趋零。这也说明力最大与功率最大不是同一个条件。
\end{xesolution}

\section{黏性、层流与相似准则}\label{ch13:sec:viscosity}

牛顿流体（Newtonian fluid）的剪切应力与速度梯度成正比：
\begin{equation}
 \tau=\eta\frac{\partial v_x}{\partial y},
\end{equation}
其中 $\eta$ 为动力黏度。半径 $R$、长度 $L$ 的圆管内充分发展层流满足哈根--泊肃叶定律（Hagen--Poiseuille law）
\begin{equation}
 Q=\frac{\pi R^4}{8\eta L}\Delta p.
\end{equation}
\FluidShearDiagram
图~\ref{ch13:fig:shear}中，上下两流层沿 $x$ 流动，速度只随法向坐标 $y$ 改变。
速度差除以层间距是剪切速率；乘以动力黏度才是单位面积的切向力。
$\eta$ 的单位为 $\mathrm{Pa\,s}$，另定义运动黏度
\[
 \boxed{\nu=\eta/\rho},\qquad [\nu]=\mathrm{m^2/s}.
\]
推导圆管流量时，设 $G=\Delta p/L>0$，距轴 $r$ 处的定常速度为 $v_z(r)$。沿轴的力平衡给出
$\eta r^{-1}\dd[r\,\dd v_z/\dd r]/\dd r=-G$。
积分一次为 $r\,\dd v_z/\dd r=-Gr^2/(2\eta)+C$；轴心速度梯度有限要求 $C=0$。
再积分并用管壁无滑移条件 $v_z(R)=0$：
\[
 v_z(r)=\frac{G}{4\eta}(R^2-r^2),\qquad
 Q=\int_0^R v_z(r)\,2\pi r\dd r
 =\frac{\pi GR^4}{8\eta}.
\]
这一步同时说明充分发展、直圆管、不可压缩牛顿流体和无滑移边界等条件为何重要。
流量对半径的四次方依赖说明细管尺寸误差会产生很大影响。雷诺数（Reynolds number）
\begin{equation}
 \mathrm{Re}=\frac{\rho v\ell}{\eta}
\end{equation}
比较惯性与黏性效应；其临界值依赖几何和扰动，不能把某个数值当成所有流动的普适分界。

低雷诺数下，半径 $R$ 的球以速度 $v$ 在黏性流体中缓慢运动，斯托克斯阻力（Stokes drag）为 $6\pi\eta Rv$。

\Needspace{16\baselineskip}\begin{origexamplebox}{小球在黏性液体中的终端速度}\label{ch13:ex:settling}
密度为 $\rho_s$、半径为 $R$ 的小球在密度 $\rho_f$、黏度 $\eta$ 的液体中下落。假设满足斯托克斯阻力，求终端速度。

\ThQSettling
\end{origexamplebox}
\Needspace{6\baselineskip}\begin{xesolution}
终端状态合力为零。向下重力减去浮力等于向上黏性阻力：
\FluidSettlingDiagram
图~\ref{ch13:fig:settling}按下沉情形画出三个力；向下为正，平衡式为
\[
 \frac43\pi R^3(\rho_s-\rho_f)g=6\pi\eta Rv_t.
\]
因此
\[
 \boxed{v_t=\frac{2(\rho_s-\rho_f)gR^2}{9\eta}}.
\]
该式要求 $\rho_s>\rho_f$ 且雷诺数足够小；若小球上浮，速度方向反向，大小用密度差绝对值。
解出后可以代回 $\mathrm{Re}=2\rho_fR|v_t|/\eta$ 检查是否远小于1。
这里特征长度取球直径 $2R$；若检验不满足，就必须更换阻力模型，而不是继续相信代数上求出的速度。
\end{xesolution}

\section{表面张力}\label{ch13:sec:surface}

表面张力 $\gamma$ 既可理解为沿界面单位长度的拉力，也可理解为增加单位面积界面所需的可逆功。
曲面两侧压强差满足拉普拉斯公式；球形液滴只有一个界面，肥皂泡有内外两个界面：
\begin{equation}
 \Delta p_{\rm drop}=\frac{2\gamma}{R},\qquad
 \Delta p_{\rm bubble}=\frac{4\gamma}{R}.
\end{equation}
半径越小，维持曲率所需压强差越大。毛细管中接触角为 $\theta$ 时，液面升降高度为
\begin{equation}
 h=\frac{2\gamma\cos\theta}{\rho gr}.
\end{equation}

\subsection{同一个系数为什么既是力除以长度，又是功除以面积}
设薄液膜有一条可滑动的直边，长度为 $\ell$，准静态向外移动 $\dd x$。液膜有前后两个界面，面积总增量为 $\dd A=2\ell\dd x$，外力做功为 $F\dd x$。在温度、组成不变，且过程可逆的条件下，$F\dd x=\gamma\dd A$，所以 $F=2\gamma\ell$。若只有一个界面，相应的因子2消失。
因此 $\mathrm{N/m}=\mathrm{J/m^2}$ 并非偶然的单位巧合，而是力做功与面积变化之间的关系。液膜沿界面拉紧边界，不是处处沿固定空间方向施力。

\subsection{球面压差与毛细升降的推导}
先看单界面的球形液滴。令半径准静态增加 $\dd R$，体积和面积的变化为
\[
 \dd V=4\pi R^2\dd R,\qquad \dd A=8\pi R\dd R.
\]
内外压强差做功为 $(p_{\rm in}-p_{\rm out})\dd V$，用于增加界面自由能 $\gamma\dd A$。约去公共因子后，$p_{\rm in}-p_{\rm out}=2\gamma/R$。薄壁肥皂泡有两个近似同半径的界面，面积变化加倍，压差才成为 $4\gamma/R$；“两个界面”不是把半径也乘2。

再看半径为 $r$ 的竖直毛细管，管外液面宽广。接触角 $\theta$ 从液体内部量取，圆形接触线长 $2\pi r$，界面张力的竖直合力为 $2\pi r\gamma\cos\theta$。在弯月面体积相对液柱体积可忽略、气体密度可忽略的近似下，它平衡高出外液面的液柱重量：
\[
 2\pi r\gamma\cos\theta=\rho g\pi r^2h.
\]
于是得到前式。$h$ 是相对管外液面的有符号高度：润湿时 $\theta<\pi/2$、$h>0$；不润湿时 $\theta>\pi/2$、$h<0$，表示液面下降。忽略弯月面体积通常要求管半径远小于毛细长度 $\sqrt{\gamma/(\rho g)}$。接触角变化、管壁污染或液柱尚在运动时，不能原样套用静态结果。

\Needspace{16\baselineskip}\begin{origexamplebox}{液滴合并时的界面能变化}\label{ch13:ex:teaching-droplets}
两滴同种不可压缩液体均为半径 $R$ 的球形液滴，合并后形成一个球形液滴。设表面张力为常量 $\gamma$，忽略重力。求合并后的半径、界面能减少量及新液滴的内外压差。界面能的减少是否等于总能量的消失？

\ThQDroplets
\end{origexamplebox}
\Needspace{6\baselineskip}\begin{xesolution}
\ThADroplets

先守恒的是体积而不是表面积。设新半径为 $R_f$，则
\[
 \frac43\pi R_f^3=2\frac43\pi R^3,
 \qquad R_f=2^{1/3}R.
\]
初、末总表面积分别为 $8\pi R^2$ 和 $4\pi 2^{2/3}R^2$。因此界面能减少
\[
 \Delta E_s=4\pi\gamma R^2(2-2^{2/3})>0.
\]
这部分能量可先转化为液滴内部流动、形变振动的能量，再由黏性耗散转化为内能；没有凭空消失。若过程中温度明显改变，常量 $\gamma$ 的近似还须重查。
新液滴压差为 $\Delta p_f=2\gamma/(2^{1/3}R)$，比原压差小。半径增大、曲率减小与压差下降一致；$\gamma R^2$ 的量纲是能量，$\gamma/R$ 的量纲是压强。
\end{xesolution}

\section{流体运动方程（选读）}\label{ch13:sec:equations-reading}
\begingroup\tcbset{breakable/.default=true}\begin{starread}{欧拉方程与纳维--斯托克斯方程}
无黏流体的局部动力学方程为
\[
 \rho\left(\frac{\partial\vect v}{\partial t}+\vect v\cdot\nabla\vect v\right)
 =-\nabla p+\rho\vect g.
\]
加入牛顿黏性应力后得到不可压缩纳维--斯托克斯方程（Navier--Stokes equation）
\[
 \rho\left(\frac{\partial\vect v}{\partial t}+\vect v\cdot\nabla\vect v\right)
 =-\nabla p+\eta\nabla^2\vect v+\rho\vect g,
 \qquad \nabla\cdot\vect v=0.
\]
非线性对流项使涡旋、边界层和湍流问题远比一维管流复杂。竞赛计算若采用伯努利或泊肃叶模型，
必须先说明忽略或保留了哪些项。

对可压缩牛顿流体，黏性应力的常用形式为
\[
 \boldsymbol\tau=\eta\left[\nabla\vect v+(\nabla\vect v)^{\mathsf T}
 -\frac23(\nabla\cdot\vect v)\mathbf 1\right]+\zeta(\nabla\cdot\vect v)\mathbf 1,
\]
其中 $\eta$ 是剪切黏度，$\zeta$ 是体黏度。不可压缩且 $\eta$ 均匀时才简化为动量方程中的 $\eta\nabla^2\vect v$。这说明泊肃叶流、低雷诺数绕流和可压缩声流采用的黏性模型虽同源，保留的项并不相同。
其来由是把总应力分为 $\boldsymbol\sigma=-p\mathbf I+\boldsymbol\tau$。
不含局部力偶的微小流体立方体满足角动量平衡，要求 $\tau_{xy}=\tau_{yx}$ 等，因此应力张量对称。
形变率还可分成无迹部分与纯体积变化：
\[
 \mathbf D=\left(\mathbf D-\frac13\nabla\cdot\vect v\,\mathbf I\right)
 +\frac13\nabla\cdot\vect v\,\mathbf I.
\]
各向同性牛顿流体的应力分别与这两部分成正比，所以出现剪切黏度和体黏度两个系数。
热力学压强 $p$ 与总正应力的平均值在可压缩黏性流动中未必相同，不能把包含体黏度项的全部 $\boldsymbol\tau$ 都称为无迹偏应力。
对任意小控制体，表面应力的合力密度为 $\nabla\cdot\boldsymbol\sigma$，故一般动量式为
\[
 \rho\frac{\mathrm D\vect v}{\mathrm Dt}=-\nabla p+\nabla\cdot\boldsymbol\tau+\rho\vect g.
\]
若 $\eta,\zeta$ 均匀，逐项求散度得到
\[
 \nabla\cdot\boldsymbol\tau=\eta\nabla^2\vect v+
 \left(\zeta+\frac\eta3\right)\nabla(\nabla\cdot\vect v).
\]
再用不可压缩条件，后一项才消失。方程虽然形式紧凑，三维非线性流动的存在性、正则性以及湍流描述仍涉及深刻数学问题；写出方程不等于已经求出了任意流动。
\end{starread}\endgroup

\newcommand{\FluidCoupledHgDiagram}{%
\begin{center}\begin{minipage}{\linewidth}\centering\begin{tikzpicture}[scale=.66,>=Latex]
 \fill[TextGray!30] (-.16,2.4)--(.16,2.4)--(.16,.16)--(.8,.16)--(.8,-.16)--(-.16,-.16)--cycle;
 \fill[TextGray!30] (10.8,-.16)--(14.56,-.16)--(14.56,2.4)--(14.24,2.4)--(14.24,.16)--(10.8,.16)--cycle;
 \draw[thick] (-.16,3.1)--(-.16,-.16)--(14.56,-.16)--(14.56,3.1);
 \draw[thick] (.16,3.1)--(.16,.16)--(14.24,.16)--(14.24,3.1);
 \draw[DeepBlue,thick] (-.16,2.4)--(.16,2.4) (14.24,2.4)--(14.56,2.4) (.8,-.16)--(.8,.16) (10.8,-.16)--(10.8,.16);
 \node at (5.8,.65) {封闭空气};
 \node[above right] at (.35,2.6) {$p_0$};\node[above left] at (14.05,2.6) {$p_0$};
 \draw[<->,TextGray] (-.9,0)--(-.9,2.4) node[midway,left]{$24\,\mathrm{cm}$};
 \draw[<->,TextGray] (15.2,0)--(15.2,2.4) node[midway,right]{$24\,\mathrm{cm}$};
 \foreach \xx in {0,.8,10.8,14.4}{\draw[dotted,TextGray] (\xx,-.22)--(\xx,-1.15);}
 \draw[<->,TextGray] (0,-.75)--(.8,-.75) node[midway,below=6pt]{$8\,\mathrm{cm}$};
 \draw[<->,TextGray] (.8,-.75)--(10.8,-.75) node[midway,below=6pt]{$\ell=100\,\mathrm{cm}$};
 \draw[<->,TextGray] (10.8,-.75)--(14.4,-.75) node[midway,below=6pt]{$36\,\mathrm{cm}$};
\end{tikzpicture}\figcap{平衡时两段水银柱之间封住空气。尺寸按管道中心线计，左右开口均通大气。}\label{ch13:fig:coupled-hg}\end{minipage}\end{center}}

\newcommand{\FluidCouetteDiagram}{%
\begin{center}\begin{minipage}{\linewidth}\centering\begin{tikzpicture}[scale=.9,>=Latex]
 \draw[thick] (0,0) circle(2.2) (0,0) circle(.85);
 \fill (0,0) circle(1.5pt) node[below left]{$O$};
 \draw[->,DeepBlue,thick] ({1.12*cos(20)},{1.12*sin(20)}) arc(20:85:1.12);
 \node[right] at (.65,1.22) {$\omega_0$};
 \draw[<->,TextGray] (0,0)--(-.85,0);
 \draw[TextGray,thin] (-.5,0)--(-1.1,-.5);
 \node[left] at (-1.1,-.5) {$R_1$};
 \draw[<->,TextGray] (0,0)--({2.2*cos(125)},{2.2*sin(125)}) node[midway,left=5pt]{$R_2$};
 \draw[<->,TextGray] (0,0)--(1.45,0);
 \node at (1.2,-.45) {$r$};
 \fill (1.45,0) circle(1.5pt) node[below right]{$P$};
 \draw[->,DeepBlue,thick] (1.45,0)--(1.45,.95) node[right]{$v_\theta$};
 \node[right] at (2.6,0) {外筒固定};
 \node[below] at (0,-2.55) {$z$ 轴竖直向上，图为水平截面};
\end{tikzpicture}\figcap{同轴圆筒的水平截面。内筒以 $\omega_0$ 转动，外筒固定；$P$ 为间隙中的考察点。}\label{ch13:fig:couette}\end{minipage}\end{center}}

\newcommand{\FluidHangarDiagram}{%
\begin{center}\begin{minipage}{\linewidth}\centering\begin{tikzpicture}[scale=.85,>=Latex]
 \draw[Blue,thick] plot[smooth] coordinates {(-4.000000,0.357197) (-3.955556,0.360639) (-3.911111,0.364264) (-3.866667,0.368088) (-3.822222,0.372125) (-3.777778,0.376394) (-3.733333,0.380912) (-3.688889,0.385701) (-3.644444,0.390785) (-3.600000,0.396189) (-3.555556,0.401943) (-3.511111,0.408078) (-3.466667,0.414633) (-3.422222,0.421646) (-3.377778,0.429165) (-3.333333,0.437240) (-3.288889,0.445931) (-3.244444,0.455303) (-3.200000,0.465432) (-3.155556,0.476401) (-3.111111,0.488308) (-3.066667,0.501261) (-3.022222,0.515383) (-2.977778,0.530813) (-2.933333,0.547705) (-2.888889,0.566233) (-2.844444,0.586584) (-2.800000,0.608962) (-2.755556,0.633578) (-2.711111,0.660646) (-2.666667,0.690366) (-2.622222,0.722907) (-2.577778,0.758385) (-2.533333,0.796837) (-2.488889,0.838195) (-2.444444,0.882276) (-2.400000,0.928775) (-2.355556,0.977284) (-2.311111,1.027322) (-2.266667,1.078380) (-2.222222,1.129955) (-2.177778,1.181586) (-2.133333,1.232872) (-2.088889,1.283486) (-2.044444,1.333169) (-2.000000,1.381725) (-1.955556,1.429015) (-1.911111,1.474943) (-1.866667,1.519449) (-1.822222,1.562501) (-1.777778,1.604088) (-1.733333,1.644211) (-1.688889,1.682886) (-1.644444,1.720133) (-1.600000,1.755979) (-1.555556,1.790455) (-1.511111,1.823591) (-1.466667,1.855421) (-1.422222,1.885976) (-1.377778,1.915288) (-1.333333,1.943390) (-1.288889,1.970311) (-1.244444,1.996081) (-1.200000,2.020726) (-1.155556,2.044275) (-1.111111,2.066750) (-1.066667,2.088177) (-1.022222,2.108577) (-0.977778,2.127972) (-0.933333,2.146380) (-0.888889,2.163821) (-0.844444,2.180311) (-0.800000,2.195867) (-0.755556,2.210503) (-0.711111,2.224234) (-0.666667,2.237072) (-0.622222,2.249029) (-0.577778,2.260117) (-0.533333,2.270344) (-0.488889,2.279721) (-0.444444,2.288256) (-0.400000,2.295955) (-0.355556,2.302827) (-0.311111,2.308876) (-0.266667,2.314109) (-0.222222,2.318528) (-0.177778,2.322140) (-0.133333,2.324945) (-0.088889,2.326947) (-0.044444,2.328148) (0.000000,2.328548) (0.044444,2.328148) (0.088889,2.326947) (0.133333,2.324945) (0.177778,2.322140) (0.222222,2.318528) (0.266667,2.314109) (0.311111,2.308876) (0.355556,2.302827) (0.400000,2.295955) (0.444444,2.288256) (0.488889,2.279721) (0.533333,2.270344) (0.577778,2.260117) (0.622222,2.249029) (0.666667,2.237072) (0.711111,2.224234) (0.755556,2.210503) (0.800000,2.195867) (0.844444,2.180311) (0.888889,2.163821) (0.933333,2.146380) (0.977778,2.127972) (1.022222,2.108577) (1.066667,2.088177) (1.111111,2.066750) (1.155556,2.044275) (1.200000,2.020726) (1.244444,1.996081) (1.288889,1.970311) (1.333333,1.943390) (1.377778,1.915288) (1.422222,1.885976) (1.466667,1.855421) (1.511111,1.823591) (1.555556,1.790455) (1.600000,1.755979) (1.644444,1.720133) (1.688889,1.682886) (1.733333,1.644211) (1.777778,1.604088) (1.822222,1.562501) (1.866667,1.519449) (1.911111,1.474943) (1.955556,1.429015) (2.000000,1.381725) (2.044444,1.333169) (2.088889,1.283486) (2.133333,1.232872) (2.177778,1.181586) (2.222222,1.129955) (2.266667,1.078380) (2.311111,1.027322) (2.355556,0.977284) (2.400000,0.928775) (2.444444,0.882276) (2.488889,0.838195) (2.533333,0.796837) (2.577778,0.758385) (2.622222,0.722907) (2.666667,0.690366) (2.711111,0.660646) (2.755556,0.633578) (2.800000,0.608962) (2.844444,0.586584) (2.888889,0.566233) (2.933333,0.547705) (2.977778,0.530813) (3.022222,0.515383) (3.066667,0.501261) (3.111111,0.488308) (3.155556,0.476401) (3.200000,0.465432) (3.244444,0.455303) (3.288889,0.445931) (3.333333,0.437240) (3.377778,0.429165) (3.422222,0.421646) (3.466667,0.414633) (3.511111,0.408078) (3.555556,0.401943) (3.600000,0.396189) (3.644444,0.390785) (3.688889,0.385701) (3.733333,0.380912) (3.777778,0.376394) (3.822222,0.372125) (3.866667,0.368088) (3.911111,0.364264) (3.955556,0.360639) (4.000000,0.357197)};
\draw[Blue,thick,->] (-2.444444,0.882276)--(-2.711111,0.660646);
\draw[Blue,thick] plot[smooth] coordinates {(-4.000000,0.912334) (-3.955556,0.919978) (-3.911111,0.927966) (-3.866667,0.936318) (-3.822222,0.945053) (-3.777778,0.954194) (-3.733333,0.963765) (-3.688889,0.973790) (-3.644444,0.984296) (-3.600000,0.995310) (-3.555556,1.006861) (-3.511111,1.018979) (-3.466667,1.031695) (-3.422222,1.045042) (-3.377778,1.059051) (-3.333333,1.073757) (-3.288889,1.089193) (-3.244444,1.105393) (-3.200000,1.122390) (-3.155556,1.140213) (-3.111111,1.158893) (-3.066667,1.178455) (-3.022222,1.198921) (-2.977778,1.220308) (-2.933333,1.242627) (-2.888889,1.265881) (-2.844444,1.290064) (-2.800000,1.315164) (-2.755556,1.341157) (-2.711111,1.368008) (-2.666667,1.395672) (-2.622222,1.424096) (-2.577778,1.453213) (-2.533333,1.482953) (-2.488889,1.513233) (-2.444444,1.543967) (-2.400000,1.575067) (-2.355556,1.606438) (-2.311111,1.637990) (-2.266667,1.669629) (-2.222222,1.701267) (-2.177778,1.732819) (-2.133333,1.764206) (-2.088889,1.795353) (-2.044444,1.826193) (-2.000000,1.856662) (-1.955556,1.886706) (-1.911111,1.916274) (-1.866667,1.945325) (-1.822222,1.973819) (-1.777778,2.001725) (-1.733333,2.029014) (-1.688889,2.055662) (-1.644444,2.081650) (-1.600000,2.106962) (-1.555556,2.131584) (-1.511111,2.155506) (-1.466667,2.178720) (-1.422222,2.201219) (-1.377778,2.223000) (-1.333333,2.244059) (-1.288889,2.264395) (-1.244444,2.284007) (-1.200000,2.302895) (-1.155556,2.321062) (-1.111111,2.338507) (-1.066667,2.355235) (-1.022222,2.371246) (-0.977778,2.386544) (-0.933333,2.401133) (-0.888889,2.415015) (-0.844444,2.428194) (-0.800000,2.440672) (-0.755556,2.452455) (-0.711111,2.463543) (-0.666667,2.473942) (-0.622222,2.483654) (-0.577778,2.492683) (-0.533333,2.501030) (-0.488889,2.508699) (-0.444444,2.515692) (-0.400000,2.522012) (-0.355556,2.527661) (-0.311111,2.532640) (-0.266667,2.536952) (-0.222222,2.540598) (-0.177778,2.543580) (-0.133333,2.545898) (-0.088889,2.547553) (-0.044444,2.548545) (0.000000,2.548876) (0.044444,2.548545) (0.088889,2.547553) (0.133333,2.545898) (0.177778,2.543580) (0.222222,2.540598) (0.266667,2.536952) (0.311111,2.532640) (0.355556,2.527661) (0.400000,2.522012) (0.444444,2.515692) (0.488889,2.508699) (0.533333,2.501030) (0.577778,2.492683) (0.622222,2.483654) (0.666667,2.473942) (0.711111,2.463543) (0.755556,2.452455) (0.800000,2.440672) (0.844444,2.428194) (0.888889,2.415015) (0.933333,2.401133) (0.977778,2.386544) (1.022222,2.371246) (1.066667,2.355235) (1.111111,2.338507) (1.155556,2.321062) (1.200000,2.302895) (1.244444,2.284007) (1.288889,2.264395) (1.333333,2.244059) (1.377778,2.223000) (1.422222,2.201219) (1.466667,2.178720) (1.511111,2.155506) (1.555556,2.131584) (1.600000,2.106962) (1.644444,2.081650) (1.688889,2.055662) (1.733333,2.029014) (1.777778,2.001725) (1.822222,1.973819) (1.866667,1.945325) (1.911111,1.916274) (1.955556,1.886706) (2.000000,1.856662) (2.044444,1.826193) (2.088889,1.795353) (2.133333,1.764206) (2.177778,1.732819) (2.222222,1.701267) (2.266667,1.669629) (2.311111,1.637990) (2.355556,1.606438) (2.400000,1.575067) (2.444444,1.543967) (2.488889,1.513233) (2.533333,1.482953) (2.577778,1.453213) (2.622222,1.424096) (2.666667,1.395672) (2.711111,1.368008) (2.755556,1.341157) (2.800000,1.315164) (2.844444,1.290064) (2.888889,1.265881) (2.933333,1.242627) (2.977778,1.220308) (3.022222,1.198921) (3.066667,1.178455) (3.111111,1.158893) (3.155556,1.140213) (3.200000,1.122390) (3.244444,1.105393) (3.288889,1.089193) (3.333333,1.073757) (3.377778,1.059051) (3.422222,1.045042) (3.466667,1.031695) (3.511111,1.018979) (3.555556,1.006861) (3.600000,0.995310) (3.644444,0.984296) (3.688889,0.973790) (3.733333,0.963765) (3.777778,0.954194) (3.822222,0.945053) (3.866667,0.936318) (3.911111,0.927966) (3.955556,0.919978) (4.000000,0.912334)};
\draw[Blue,thick,->] (-2.444444,1.543967)--(-2.711111,1.368008);
\draw[Blue,thick] plot[smooth] coordinates {(-4.000000,1.497260) (-3.955556,1.507103) (-3.911111,1.517277) (-3.866667,1.527793) (-3.822222,1.538663) (-3.777778,1.549898) (-3.733333,1.561508) (-3.688889,1.573505) (-3.644444,1.585897) (-3.600000,1.598697) (-3.555556,1.611912) (-3.511111,1.625551) (-3.466667,1.639623) (-3.422222,1.654135) (-3.377778,1.669092) (-3.333333,1.684499) (-3.288889,1.700360) (-3.244444,1.716675) (-3.200000,1.733445) (-3.155556,1.750667) (-3.111111,1.768338) (-3.066667,1.786450) (-3.022222,1.804994) (-2.977778,1.823960) (-2.933333,1.843334) (-2.888889,1.863099) (-2.844444,1.883236) (-2.800000,1.903724) (-2.755556,1.924539) (-2.711111,1.945656) (-2.666667,1.967047) (-2.622222,1.988681) (-2.577778,2.010528) (-2.533333,2.032554) (-2.488889,2.054726) (-2.444444,2.077010) (-2.400000,2.099370) (-2.355556,2.121772) (-2.311111,2.144179) (-2.266667,2.166559) (-2.222222,2.188876) (-2.177778,2.211096) (-2.133333,2.233189) (-2.088889,2.255123) (-2.044444,2.276867) (-2.000000,2.298394) (-1.955556,2.319676) (-1.911111,2.340688) (-1.866667,2.361407) (-1.822222,2.381809) (-1.777778,2.401874) (-1.733333,2.421582) (-1.688889,2.440916) (-1.644444,2.459860) (-1.600000,2.478398) (-1.555556,2.496516) (-1.511111,2.514202) (-1.466667,2.531444) (-1.422222,2.548233) (-1.377778,2.564558) (-1.333333,2.580411) (-1.288889,2.595785) (-1.244444,2.610673) (-1.200000,2.625069) (-1.155556,2.638968) (-1.111111,2.652365) (-1.066667,2.665255) (-1.022222,2.677636) (-0.977778,2.689503) (-0.933333,2.700854) (-0.888889,2.711687) (-0.844444,2.722000) (-0.800000,2.731790) (-0.755556,2.741056) (-0.711111,2.749797) (-0.666667,2.758012) (-0.622222,2.765700) (-0.577778,2.772859) (-0.533333,2.779490) (-0.488889,2.785592) (-0.444444,2.791164) (-0.400000,2.796206) (-0.355556,2.800718) (-0.311111,2.804700) (-0.266667,2.808151) (-0.222222,2.811071) (-0.177778,2.813460) (-0.133333,2.815319) (-0.088889,2.816646) (-0.044444,2.817443) (0.000000,2.817708) (0.044444,2.817443) (0.088889,2.816646) (0.133333,2.815319) (0.177778,2.813460) (0.222222,2.811071) (0.266667,2.808151) (0.311111,2.804700) (0.355556,2.800718) (0.400000,2.796206) (0.444444,2.791164) (0.488889,2.785592) (0.533333,2.779490) (0.577778,2.772859) (0.622222,2.765700) (0.666667,2.758012) (0.711111,2.749797) (0.755556,2.741056) (0.800000,2.731790) (0.844444,2.722000) (0.888889,2.711687) (0.933333,2.700854) (0.977778,2.689503) (1.022222,2.677636) (1.066667,2.665255) (1.111111,2.652365) (1.155556,2.638968) (1.200000,2.625069) (1.244444,2.610673) (1.288889,2.595785) (1.333333,2.580411) (1.377778,2.564558) (1.422222,2.548233) (1.466667,2.531444) (1.511111,2.514202) (1.555556,2.496516) (1.600000,2.478398) (1.644444,2.459860) (1.688889,2.440916) (1.733333,2.421582) (1.777778,2.401874) (1.822222,2.381809) (1.866667,2.361407) (1.911111,2.340688) (1.955556,2.319676) (2.000000,2.298394) (2.044444,2.276867) (2.088889,2.255123) (2.133333,2.233189) (2.177778,2.211096) (2.222222,2.188876) (2.266667,2.166559) (2.311111,2.144179) (2.355556,2.121772) (2.400000,2.099370) (2.444444,2.077010) (2.488889,2.054726) (2.533333,2.032554) (2.577778,2.010528) (2.622222,1.988681) (2.666667,1.967047) (2.711111,1.945656) (2.755556,1.924539) (2.800000,1.903724) (2.844444,1.883236) (2.888889,1.863099) (2.933333,1.843334) (2.977778,1.823960) (3.022222,1.804994) (3.066667,1.786450) (3.111111,1.768338) (3.155556,1.750667) (3.200000,1.733445) (3.244444,1.716675) (3.288889,1.700360) (3.333333,1.684499) (3.377778,1.669092) (3.422222,1.654135) (3.466667,1.639623) (3.511111,1.625551) (3.555556,1.611912) (3.600000,1.598697) (3.644444,1.585897) (3.688889,1.573505) (3.733333,1.561508) (3.777778,1.549898) (3.822222,1.538663) (3.866667,1.527793) (3.911111,1.517277) (3.955556,1.507103) (4.000000,1.497260)};
\draw[Blue,thick,->] (-2.444444,2.077010)--(-2.711111,1.945656);
\draw[Blue,thick] plot[smooth] coordinates {(-4.000000,2.155276) (-3.955556,2.165382) (-3.911111,2.175721) (-3.866667,2.186297) (-3.822222,2.197111) (-3.777778,2.208165) (-3.733333,2.219461) (-3.688889,2.231000) (-3.644444,2.242783) (-3.600000,2.254809) (-3.555556,2.267079) (-3.511111,2.279591) (-3.466667,2.292343) (-3.422222,2.305333) (-3.377778,2.318559) (-3.333333,2.332016) (-3.288889,2.345699) (-3.244444,2.359605) (-3.200000,2.373725) (-3.155556,2.388055) (-3.111111,2.402585) (-3.066667,2.417308) (-3.022222,2.432213) (-2.977778,2.447292) (-2.933333,2.462534) (-2.888889,2.477926) (-2.844444,2.493456) (-2.800000,2.509112) (-2.755556,2.524880) (-2.711111,2.540746) (-2.666667,2.556696) (-2.622222,2.572714) (-2.577778,2.588785) (-2.533333,2.604893) (-2.488889,2.621022) (-2.444444,2.637156) (-2.400000,2.653280) (-2.355556,2.669376) (-2.311111,2.685428) (-2.266667,2.701420) (-2.222222,2.717335) (-2.177778,2.733159) (-2.133333,2.748876) (-2.088889,2.764469) (-2.044444,2.779924) (-2.000000,2.795227) (-1.955556,2.810362) (-1.911111,2.825317) (-1.866667,2.840078) (-1.822222,2.854631) (-1.777778,2.868966) (-1.733333,2.883068) (-1.688889,2.896928) (-1.644444,2.910534) (-1.600000,2.923876) (-1.555556,2.936944) (-1.511111,2.949729) (-1.466667,2.962221) (-1.422222,2.974412) (-1.377778,2.986295) (-1.333333,2.997861) (-1.288889,3.009103) (-1.244444,3.020015) (-1.200000,3.030591) (-1.155556,3.040824) (-1.111111,3.050709) (-1.066667,3.060242) (-1.022222,3.069416) (-0.977778,3.078229) (-0.933333,3.086674) (-0.888889,3.094750) (-0.844444,3.102451) (-0.800000,3.109776) (-0.755556,3.116719) (-0.711111,3.123280) (-0.666667,3.129455) (-0.622222,3.135242) (-0.577778,3.140638) (-0.533333,3.145643) (-0.488889,3.150253) (-0.444444,3.154467) (-0.400000,3.158284) (-0.355556,3.161703) (-0.311111,3.164722) (-0.266667,3.167340) (-0.222222,3.169557) (-0.177778,3.171372) (-0.133333,3.172785) (-0.088889,3.173794) (-0.044444,3.174399) (0.000000,3.174601) (0.044444,3.174399) (0.088889,3.173794) (0.133333,3.172785) (0.177778,3.171372) (0.222222,3.169557) (0.266667,3.167340) (0.311111,3.164722) (0.355556,3.161703) (0.400000,3.158284) (0.444444,3.154467) (0.488889,3.150253) (0.533333,3.145643) (0.577778,3.140638) (0.622222,3.135242) (0.666667,3.129455) (0.711111,3.123280) (0.755556,3.116719) (0.800000,3.109776) (0.844444,3.102451) (0.888889,3.094750) (0.933333,3.086674) (0.977778,3.078229) (1.022222,3.069416) (1.066667,3.060242) (1.111111,3.050709) (1.155556,3.040824) (1.200000,3.030591) (1.244444,3.020015) (1.288889,3.009103) (1.333333,2.997861) (1.377778,2.986295) (1.422222,2.974412) (1.466667,2.962221) (1.511111,2.949729) (1.555556,2.936944) (1.600000,2.923876) (1.644444,2.910534) (1.688889,2.896928) (1.733333,2.883068) (1.777778,2.868966) (1.822222,2.854631) (1.866667,2.840078) (1.911111,2.825317) (1.955556,2.810362) (2.000000,2.795227) (2.044444,2.779924) (2.088889,2.764469) (2.133333,2.748876) (2.177778,2.733159) (2.222222,2.717335) (2.266667,2.701420) (2.311111,2.685428) (2.355556,2.669376) (2.400000,2.653280) (2.444444,2.637156) (2.488889,2.621022) (2.533333,2.604893) (2.577778,2.588785) (2.622222,2.572714) (2.666667,2.556696) (2.711111,2.540746) (2.755556,2.524880) (2.800000,2.509112) (2.844444,2.493456) (2.888889,2.477926) (2.933333,2.462534) (2.977778,2.447292) (3.022222,2.432213) (3.066667,2.417308) (3.111111,2.402585) (3.155556,2.388055) (3.200000,2.373725) (3.244444,2.359605) (3.288889,2.345699) (3.333333,2.332016) (3.377778,2.318559) (3.422222,2.305333) (3.466667,2.292343) (3.511111,2.279591) (3.555556,2.267079) (3.600000,2.254809) (3.644444,2.242783) (3.688889,2.231000) (3.733333,2.219461) (3.777778,2.208165) (3.822222,2.197111) (3.866667,2.186297) (3.911111,2.175721) (3.955556,2.165382) (4.000000,2.155276)};
\draw[Blue,thick,->] (-2.444444,2.637156)--(-2.711111,2.540746);

 \draw[thick] (-2.2,0) arc(180:0:2.2);
 \draw[->,TextGray] (-4.2,0)--(4.5,0) node[right]{$x$};
 \fill (-2.2,0) circle(1.6pt);\node[below] at (-2.2,-.15) {$A$（门口）};
 \fill (0,0) circle(1.6pt) node[below]{$O$};
 \coordinate (B) at ({2.2*cos(55)},{2.2*sin(55)});
 \draw[TextGray] (0,0)--(B);\fill (B) circle(1.6pt);
 \draw[TextGray,thin] (.8,1.7)--(B);
 \node[left] at (.7,1.7) {$B$};\node[left] at (.55,.95) {$R$};
 \draw[TextGray] (.65,0) arc(0:55:.65);\node at (.92,.35) {$\theta$};
 \node[above left] at (-2.4,4.1) {来流向左};
\end{tikzpicture}\figcap{半圆柱机库的横截面与题设流线。$A$ 为敞开的门口，$R$ 为半径，$\theta$ 从正 $x$ 方向量起。}\label{ch13:fig:hangar}\end{minipage}\end{center}}

\newcommand{\FluidFunnelDiagram}{%
\begin{center}\begin{minipage}{\linewidth}\centering\begin{tikzpicture}[scale=.85,>=Latex]
 \def\aa{.16}\def\RR{2.4}\def\HH{3.4}\def\hh{2.65}
 \pgfmathsetmacro{\rh}{\aa+(\RR-\aa)*\hh/\HH}
 \fill[Blue!12] (-\aa,0)--(-\rh,\hh)--(\rh,\hh)--(\aa,0)--cycle;
 \draw[thick] (-\RR,\HH)--(-\aa,0) (\aa,0)--(\RR,\HH);
 \draw[DeepBlue,thick] (-\rh,\hh)--(\rh,\hh);
 \draw[DeepBlue,thick] (-\aa,0)--(-\aa,-1) (\aa,0)--(\aa,-1);
 \draw[->,DeepBlue,thick] (0,-.4)--(0,-1.25) node[midway,right=7pt]{$Q$};
 \draw[<->,TextGray] (-3,0)--(-3,\HH) node[midway,left]{$H$};
 \draw[<->,TextGray] (2.9,0)--(2.9,\hh) node[midway,right]{$h$};
 \draw[dotted,TextGray] (\aa,0)--(2.9,0) (\rh,\hh)--(2.9,\hh);
 \draw[<->,TextGray] (0,3.85)--(\RR,3.85) node[midway,above]{$R$};
 \draw[TextGray,thin] (-\aa,-1.25)--(-\aa,-1.68) (\aa,-1.25)--(\aa,-1.68);
 \draw[->,TextGray] (-.9,-1.55)--(-\aa,-1.55);
 \draw[->,TextGray] (.9,-1.55)--(\aa,-1.55);
 \node[below] at (0,-1.65) {有限开口 $2a$};
\end{tikzpicture}\figcap{圆锥形漏斗与连通的有限出口。$H$、$h$ 均从出口量起，图示排液过程中的一个状态。}\label{ch13:fig:funnel}\end{minipage}\end{center}}

\newcommand{\FluidShoreDiagram}{%
\begin{center}\begin{minipage}{\linewidth}\centering\begin{tikzpicture}[scale=.9,>=Latex]
 \begin{scope}[shift={(-5.6,1.4)}]
  \fill[Blue!12] (0,0)--(4.2,0)--(4.2,-1.2)--cycle;
  \draw[DeepBlue,thick] (0,0)--(4.2,0);
  \draw[Ink,thick] (-.7,.2)--(4.2,-1.2);
  \draw[TextGray] (1.1,0) arc(0:-15.9454:1.1);
  \node at (1.4,-.22) {$\alpha$};
  \fill (0,0) circle(1.5pt) node[above] {岸};
  \node[above] at (3.1,.15) {水面};
  \node[below] at (3.1,-1.3) {岸坡剖面};
 \end{scope}
 \begin{scope}[shift={(.2,0)}]
  \fill[Blue!8] (0,0) rectangle (3.6,2.7);
  \draw[->,TextGray] (0,0)--(3.9,0) node[right]{$x$};
  \draw[->,TextGray] (0,0)--(0,3) node[above]{$y$};
  \draw[DeepBlue,very thick,domain=0:180,samples=120,variable=\p]
   plot ({pi*\p/180+sin(\p)},{1+cos(\p)});
  \fill (0,2) circle(1.8pt) node[left]{$B(0,y_0)$};
  \fill (pi,0) circle(1.8pt) node[below]{$A$};
  \draw[->,DeepBlue,thick] (0,2)--(.8,2) node[midway,above=5pt]{$\vect c_0$};
  \node[below] at (1,-.4) {岸线 $y=0$};
 \end{scope}
\end{tikzpicture}\figcap{左为岸坡剖面，水只在水面与坡底之间；右为原题所给的弯向岸边的射线示意，$B$ 在离岸坐标轴上。}\label{ch13:fig:shore}\end{minipage}\end{center}}

\newcommand{\FluidHoseDiagram}{%
\begin{center}\begin{minipage}{\linewidth}\centering\begin{tikzpicture}[scale=.95,>=Latex]
 \def\ang{25}\def\LL{2.65}
 \coordinate (O) at (0,0);
 \coordinate (E) at ({-\LL*sin(\ang)},{-\LL*cos(\ang)});
 \coordinate (P) at ({\LL*(cos(\ang)-sin(\ang))},{-\LL*(cos(\ang)+sin(\ang))});
 \draw[dashed,TextGray] (O)--(0,-\LL)--(\LL,-\LL);
 \draw[DeepBlue,line width=2pt] (O)--(E)--(P);
 \foreach \pt in {O,E,P}{\fill (\pt) circle(1.7pt);}
 \node[above] at (O) {$O$};\node[left=5pt] at (E) {$E$};\node[below=7pt] at (P) {$P$};
 \draw[->,Ink] ($(O)!.25!(E)$)--($(O)!.65!(E)$);
 \draw[->,Ink] ($(E)!.25!(P)$)--($(E)!.65!(P)$);
 \draw[TextGray] (0,-.75) arc(-90:{-90-\ang}:.75);
 \node[right] at (.15,-.95) {$\theta$};
 \coordinate (d1) at ({-.6*cos(\ang)},{.6*sin(\ang)});
 \coordinate (d2) at ({-.55*sin(\ang)},{-.55*cos(\ang)});
 \draw[<->,TextGray] ($(O)+(d1)$)--($(E)+(d1)$) node[midway,left=5pt]{$L$};
 \draw[<->,TextGray] ($(E)+(d2)$)--($(P)+(d2)$) node[midway,below=5pt]{$L$};
 \draw[->,Ink] (2.6,-.2)--(2.6,-1.4) node[right]{$\vect g$};
 \draw[->,DeepBlue,thick] (P)--++({.85*cos(\ang)},{-.85*sin(\ang)}) node[right]{出流方向};
\end{tikzpicture}\figcap{两段等长管在 $E$ 处保持直角；实线仅示意未知偏角，虚线为题设的角度参考构形。}\label{ch13:fig:hose}\end{minipage}\end{center}}

\newcommand{\FluidHoseAnalysisDiagram}{%
\begin{center}\begin{minipage}{\linewidth}\centering\begin{tikzpicture}[scale=.95,>=Latex]
 \def\ang{31.0809622}\def\LL{2.65}
 \coordinate (O) at (0,0);
 \coordinate (E) at ({-\LL*sin(\ang)},{-\LL*cos(\ang)});
 \coordinate (P) at ({\LL*(cos(\ang)-sin(\ang))},{-\LL*(cos(\ang)+sin(\ang))});
 \draw[dashed,TextGray] (O)--(0,-\LL)--(\LL,-\LL);
 \draw[DeepBlue,line width=2pt] (O)--(E)--(P);
 \foreach \pt in {O,E,P}{\fill (\pt) circle(1.7pt);}
 \node[above] at (O) {$O$};\node[left=5pt] at (E) {$E$};\node[below=7pt] at (P) {$P$};
 \draw[->,Ink] ($(O)!.25!(E)$)--($(O)!.65!(E)$);
 \draw[->,Ink] ($(E)!.25!(P)$)--($(E)!.65!(P)$);
 \draw[TextGray] (0,-.75) arc(-90:{-90-\ang}:.75);
 \node[right] at (.15,-.95) {$\theta=31.1^\circ$};
 \coordinate (d1) at ({-.6*cos(\ang)},{.6*sin(\ang)});
 \coordinate (d2) at ({-.55*sin(\ang)},{-.55*cos(\ang)});
 \draw[<->,TextGray] ($(O)+(d1)$)--($(E)+(d1)$) node[midway,left=5pt]{$L$};
 \draw[<->,TextGray] ($(E)+(d2)$)--($(P)+(d2)$) node[midway,below=5pt]{$L$};
 \draw[->,Ink] (2.6,-.2)--(2.6,-1.4) node[right]{$\vect g$};
 \draw[->,DeepBlue,thick] (P)--++({.85*cos(\ang)},{-.85*sin(\ang)}) node[right]{$u=Q/A$};
\end{tikzpicture}\figcap{两段等长管在 $E$ 处保持直角；实线按求得的稳定角度绘制，虚线为比较用初始构形。}\label{ch13:fig:hose-analysis}\end{minipage}\end{center}}

\section{流体动力学综合问题}\label{ch13:sec:problems}

\Needspace{16\baselineskip}\begin{origexamplebox}{两段水银柱与封闭空气的耦合振动}\label{ch13:ex:coupled-hg}
等截面的 U 形管内封有两段水银柱，中间夹着一段空气柱。平衡时，两侧开口处的水银自由面均高出
水平空气柱 $h=24\,\mathrm{cm}$；左、右水银柱沿管道的总长度分别为
$L_1=24\,\mathrm{cm}+8\,\mathrm{cm}$、$L_2=24\,\mathrm{cm}+36\,\mathrm{cm}$，
空气柱长 $\ell=100\,\mathrm{cm}$。大气压等于 $76\,\mathrm{cmHg}$，取
$g=9.8\,\mathrm{m/s^2}$。忽略阻力，把空气视为双原子理想气体并认为微振动过程绝热，求两个本征频率。

\FluidCoupledHgDiagram
\end{origexamplebox}
\Needspace{6\baselineskip}\begin{xesolution}
\ThACoupledMercury

图~\ref{ch13:fig:coupled-hg}中，水平气柱中心线到两个自由面的竖直高度均为 $h$。
平衡时空气压强等于大气压再加一段高度为 $h$ 的水银柱压强：
\[
 \bar p=p_0+\rho gh=\rho g(0.76+0.24)\,\mathrm m=\rho g\ell .
\]
这里 $\rho$ 是水银密度。定义 $q_1$、$q_2$ 分别为左、右水银内侧界面向空气柱内的位移，
即左界面向右为正，右界面向左为正；对应外侧自由面均下降同样距离。
因此空气体积的一阶相对变化为
\[
 \frac{\delta V}{V_0}=-\frac{q_1+q_2}{\ell}.
\]
绝热关系 $pV^\gamma=\text{常量}$ 给出
\[
 \delta p=\gamma\bar p\frac{q_1+q_2}{\ell}.
\]

以左水银柱为例。自由面下降 $q_1$ 使重力压头减少 $\rho gq_1$；空气压强增加 $\delta p$，也把水银柱推离空气。
沿 $q_1$ 正向写牛顿第二定律，管截面积约去后得到
\[
 L_1\ddot q_1+gq_1+\frac{\delta p}{\rho}=0.
\]
右侧同理。利用 $\bar p/(\rho\ell)=g$，两式成为
\begin{align*}
 L_1\ddot q_1+gq_1+\gamma g(q_1+q_2)&=0,\\
 L_2\ddot q_2+gq_2+\gamma g(q_1+q_2)&=0.
\end{align*}
令 $q_i=Q_i e^{\ii\omega t}$，非零振幅要求
\[
 \begin{vmatrix}
 g(1+\gamma)-\omega^2L_1&\gamma g\\
 \gamma g&g(1+\gamma)-\omega^2L_2
 \end{vmatrix}=0.
\]
记 $x=\omega^2/g$，便有
\[
 L_1L_2x^2-(1+\gamma)(L_1+L_2)x+(1+2\gamma)=0.
\]
代入 $L_1=0.32\,\mathrm m$、$L_2=0.60\,\mathrm m$、$\gamma=7/5$，得到
\[
 x_-=2.1071\,\mathrm{m^{-1}},\qquad x_+=9.3929\,\mathrm{m^{-1}},
\]
所以
\[
 \boxed{\omega_-=4.54\,\mathrm{rad/s},\quad \omega_+=9.59\,\mathrm{rad/s}},
 \qquad
 \boxed{f_-=0.723\,\mathrm{Hz},\quad f_+=1.53\,\mathrm{Hz}}.
\]
低频模中两段水银大体反向运动，空气压缩较小；高频模中两者大体同时压缩空气。
若振动足够慢、空气近似等温，只需把 $\gamma$ 改为 $1$，会得到
$f_-=0.719\,\mathrm{Hz}$、$f_+=1.36\,\mathrm{Hz}$。
\end{xesolution}

\Needspace{16\baselineskip}\begin{origexamplebox}{同轴圆筒间的黏性旋转流}\label{ch13:ex:couette}
半径为 $R_1$ 与 $R_2$（$R_2>R_1$）的两个长圆筒竖直同轴放置，间隙充满密度为 $\rho$、动力黏度为 $\eta$ 的
不可压缩牛顿流体。内筒以恒定角速度 $\omega_0$ 转动，外筒静止，流动稳定且轴对称。

（1）求角速度分布 $\omega(r)$；（2）先求同一高度的压强分布，再补出竖直方向的变化；（3）说明柱坐标中的剪切应力公式
\[
 \tau_{r\theta}=\eta\left(\frac1r\frac{\partial v_r}{\partial\theta}
 +\frac{\partial v_\theta}{\partial r}-\frac{v_\theta}{r}\right)
\]
为何与直角坐标中的形式不同。

\FluidCouetteDiagram
\end{origexamplebox}
\Needspace{6\baselineskip}\begin{xesolution}
\ThACouette

稳定轴对称流满足 $v_r=v_z=0$、$v_\theta=r\omega(r)$。
图~\ref{ch13:fig:couette}中 $P$ 的切向速度为 $v_\theta$，半径为 $r$；所画水平截面上各点的高度相同。将速度代入切应力公式，得到
\[
 \tau_{r\theta}=\eta\left(\frac{\dd v_\theta}{\dd r}-\frac{v_\theta}{r}\right)
 =\eta r\frac{\dd\omega}{\dd r}.
\]
任取半径为 $r$、长度为 $l$ 的圆柱面，黏性力矩为
$2\pi rl\,\tau_{r\theta}\,r$。稳定时各圆柱面传递的力矩相同，故
$r^2\tau_{r\theta}$ 为常量，即 $r^3\omega'(r)$ 为常量。积分并使用
$\omega(R_1)=\omega_0$、$\omega(R_2)=0$，得
\[
 \boxed{\omega(r)=\frac{\omega_0R_1^2}{R_2^2-R_1^2}
 \left(\frac{R_2^2}{r^2}-1\right)}.
\]
这一步可写为 $r^3\omega'=K$，因此 $\omega=C_1+C_2/r^2$。外壁条件给 $C_1=-C_2/R_2^2$，内壁条件再给
$C_2=\omega_0R_1^2R_2^2/(R_2^2-R_1^2)$。两边界均满足，且 $\omega_0=0$ 时全流体静止。

径向速度虽为零，流体微元仍有向心加速度 $-v_\theta^2/r$。径向欧拉方程因此给出
\[
 \frac{\dd p}{\dd r}=\rho\frac{v_\theta^2}{r}.
\]
令 $C=\omega_0R_1^2/(R_2^2-R_1^2)$，由
$v_\theta=C(R_2^2/r-r)$ 积分可得
\[
 \boxed{p(r)-p(R_1)=\rho C^2\left[
 \frac{r^2-R_1^2}{2}-2R_2^2\ln\frac r{R_1}
 -\frac{R_2^4}{2}\left(\frac1{r^2}-\frac1{R_1^2}\right)
 \right]}.
\]
因 $\dd p/\dd r=\rho v_\theta^2/r\ge0$，压强必向外增大，这也直接检查了积分式的符号。
上式只比较同一高度的两点。取竖直向上为 $z$，由于 $v_z=0$，竖直平衡还要求
$\partial p/\partial z=-\rho g$。若把方括号中的径向差记为 $\Delta p_r(r)$，完整压强场为
\[
 p(r,z)=p(R_1,z_0)+\Delta p_r(r)-\rho g(z-z_0).
\]
在不转动极限，径向差消失，只剩熟悉的静水压分布。

柱坐标基矢 $\vect e_r,\vect e_\theta$ 随方位角改变，且方位方向的实际弧长是 $r\,\dd\theta$，所以：
（i）$r^{-1}\partial v_r/\partial\theta$ 来自弧长尺度；
（ii）$-v_\theta/r$ 来自基矢方向随位置转动；
（iii）圆柱面传递力矩时还多一个力臂 $r$，因此守恒的是 $r^2\tau_{r\theta}$，而非平面流中的恒定剪切应力。
同样，曲线运动还产生直角平行流中没有的向心压强梯度。
\end{xesolution}

\Needspace{16\baselineskip}\begin{origexamplebox}{半圆柱形机库受到的风力}\label{ch13:ex:hangar}
长度为 $l$、横截面半径为 $R$ 的半圆柱形机库，处在垂直于柱轴的均匀风中。空气密度为 $\rho$，忽略重力与黏性，外部速度势为
\[
 \phi=-v_0\left(r+\frac{R^2}{r}\right)\cos\theta .
\]
地面处的门口 $A$ 敞开，使库内空气可视为静止且压强等于 $A$ 点外部压强。求风对机库曲面部分的合力。

\FluidHangarDiagram
\end{origexamplebox}
\Needspace{6\baselineskip}\begin{xesolution}
\ThAHangar

图~\ref{ch13:fig:hangar}中，$\theta$ 从正 $x$ 方向量起，门口 $A$ 对应 $\pi$。
由 $\vect v=\nabla\phi$，在曲面 $r=R$ 上有
\[
 v_r=0,\qquad v_\theta=\frac1r\frac{\partial\phi}{\partial\theta}
 =2v_0\sin\theta.
\]
外部沿流线使用伯努利方程：
\[
 p_{\rm out}(\theta)=p_\infty+\frac12\rho v_0^2
 \left(1-4\sin^2\theta\right).
\]
门口 $A$ 对应 $\theta=\pi$，是停滞点，故库内静止空气压强为
\[
 p_{\rm in}=p_A=p_\infty+\frac12\rho v_0^2.
\]
于是曲面内外压强差为
\[
 \Delta p=p_{\rm in}-p_{\rm out}=2\rho v_0^2\sin^2\theta.
\]
曲面面积元为 $\dd S=lR\,\dd\theta$，外法线
$\vect n=(\cos\theta,\sin\theta)$。左右对称使水平分量相消，而竖直分量为
\begin{align*}
 F_y&=2\rho v_0^2lR\int_0^\pi\sin^3\theta\,\dd\theta\\
 &=\boxed{\frac83\rho v_0^2Rl},\qquad F_x=0.
\end{align*}
合力竖直向上。若门关闭且库内压强不是停滞点压强，这个结论会随内部压强改变，不能直接套用。
\end{xesolution}

\Needspace{16\baselineskip}\begin{origexamplebox}{圆锥形漏斗排液}\label{ch13:ex:funnel}
一个尖端向下的圆锥形漏斗，高为 $H$、顶面半径为 $R$，初始装满水。底部有半径为 $a$ 的小孔，打开后自由流出。
忽略自由液面速度和黏性，求水面从高度 $H$ 降到 $H/2$ 所需的时间，并说明何时必须考虑黏性。

\FluidFunnelDiagram
\end{origexamplebox}
\Needspace{6\baselineskip}\begin{xesolution}
\ThAFunnel

图~\ref{ch13:fig:funnel}保留有限出口。先取 $a\ll R$ 的尖锥近似，忽略小孔截断锥尖带来的几何修正。
由相似三角形，水深为 $h$ 时的自由液面面积为
\[
 A(h)=\pi\left(\frac{Rh}{H}\right)^2.
\]
小孔处托里拆利速度为 $v=\sqrt{2gh}$，流量
$Q=\pi a^2\sqrt{2gh}$。体积守恒 $A(h)\dd h/\dd t=-Q$，故
\[
 \dd t=-\frac{R^2}{a^2H^2\sqrt{2g}}h^{3/2}\,\dd h.
\]
从 $H$ 积分到 $H/2$：
\[
 \boxed{t=\frac{2R^2}{5a^2}\sqrt{\frac{H}{2g}}
 \left(1-2^{-5/2}\right)}.
\]
量纲由 $\sqrt{H/g}$ 给出时间。这里忽略黏性的条件不是“小孔很小”，而是小孔足够短、雷诺数足够大，流出损失可用接近 $1$ 的流量系数描述。
实际锐边小孔若流量系数为 $C_d<1$，时间应除以 $C_d$；若出口是一根长细管，则黏性压降占主导，必须改用泊肃叶流模型。
若不忽略有限孔径，令 $H$ 为出口到顶面的真实高度，截面半径应为
$r(h)=a+(R-a)h/H$。设 $b=(R-a)/H$，则把同一个体积守恒式中的面积换成 $\pi(a+bh)^2$，得到
\[
 t=\frac1{\sqrt{2g}}\left[
 2\sqrt h+\frac{4b}{3a}h^{3/2}+\frac{2b^2}{5a^2}h^{5/2}
 \right]_{H/2}^{H}.
\]
当 $a/R\to0$ 时，最后一项占主导，恢复前面的尖锥结果。有限开口与理想模型的关系因此是可检验的近似，不是把管口画成封闭尖点。
\end{xesolution}

\Needspace{16\baselineskip}\begin{origexamplebox}{浅水波向岸边弯折的轨迹}\label{ch13:ex:shore}
平直海岸的水深随离岸距离线性增加。取 $x$ 轴沿岸、$y$ 轴垂直海岸并指向海面，岸线为 $y=0$，水深
$h(y)=y\tan\alpha$。浅水波速满足 $c=\sqrt{gh}=C\sqrt y$。一列波包从
$B=(0,y_0)$ 出发，初始传播方向平行于岸线。求波包射线，并证明它是一段摆线。

\FluidShoreDiagram
\end{origexamplebox}
\Needspace{6\baselineskip}\begin{xesolution}
\ThAShore

介质只随 $y$ 变化，沿岸方向平移对称。令 $\theta$ 为射线与指向岸边的法线方向（$-y$ 方向）的夹角，
图~\ref{ch13:fig:shore}中 $B$ 在 $y$ 轴上，故初始沿岸坐标为零。
射线折射不变量为
\[
 \frac{\sin\theta}{c(y)}=\text{常量}.
\]
在 $B$ 点射线平行于岸线，$\theta_0=\pi/2$，故
\[
 \sin\theta=\frac{c(y)}{c(y_0)}=\sqrt{\frac y{y_0}}.
\]
沿传播方向 $y$ 减小、$x$ 增加，因而
\[
 \frac{\dd x}{\dd y}=-\tan\theta.
\]
作参数
\[
 \varphi=\pi-2\theta,\qquad 0\le\varphi\le\pi,
\]
由 $y/y_0=\sin^2\theta=(1+\cos\varphi)/2$，有
\[
 \frac{\dd y}{\dd\varphi}=-\frac{y_0}{2}\sin\varphi,
 \qquad \tan\theta=\cot(\varphi/2).
\]
所以
\[
 \frac{\dd x}{\dd\varphi}
 =\left(-\cot\frac\varphi2\right)\left(-\frac{y_0}2\sin\varphi\right)
 =\frac{y_0}2(1+\cos\varphi).
\]
从0积分到 $\varphi$，并取 $x(0)=0$，得到
\[
 \boxed{x=\frac{y_0}{2}(\varphi+\sin\varphi)},\qquad
 \boxed{y=\frac{y_0}{2}(1+\cos\varphi)}.
\]
这就是一段上下翻转的摆线。$\varphi=0$ 时位于 $B=(0,y_0)$，$\varphi=\pi$ 时到达
$A=(\pi y_0/2,0)$。此时 $\theta=0$，即射线垂直岸线到达，而不是斜着撞岸。
最后一句是线性浅水波射线模型的极限描述；实际近岸波幅可能不再远小于水深，出现破碎等现象，不能把理想轨迹继续用于这些区域。
\end{xesolution}

\Needspace{16\baselineskip}\begin{origexamplebox}{自由转动直角水管的稳定偏角}\label{ch13:ex:hose}
一根均匀直角软管的内截面积为 $A=2.0\,\mathrm{cm^2}$，软管自身质量为 $m_h=100\,\mathrm g$，
两段管长均为 $L=40\,\mathrm{cm}$，转动时两段形状不变且始终保持直角。进水端在固定接头 $O$ 处可在竖直平面内自由转动；水从第一段流向弯头，再由第二段射出，
体积流量为 $Q=0.30\,\mathrm{L/s}$。取水的密度 $\rho=1.0\times10^3\,\mathrm{kg/m^3}$、
$g=10\,\mathrm{m/s^2}$，忽略黏性和出口之外水流的重量，求稳定时软管相对“第一段竖直向下、第二段水平向右”构形向后转过的角度 $\theta$。

\FluidHoseDiagram
\end{origexamplebox}
\Needspace{6\baselineskip}\begin{xesolution}
\ThAHose

图~\ref{ch13:fig:hose}中 $O$ 是转轴，$E$ 是直角弯头，$P$ 是出水口。两段长度相等且转动中保持垂直。
出口平均速度为
\[
 u=\frac QA=\frac{3.0\times10^{-4}}{2.0\times10^{-4}}
 =1.50\,\mathrm{m/s}.
\]
管内水的质量不能遗漏：
\[
 m_w=\rho A(2L)=0.160\,\mathrm{kg},\qquad
 M=m_h+m_w=0.260\,\mathrm{kg}.
\]
取 $x$ 向右、$y$ 向上。顺时针转过 $\theta$ 后，第一、第二段的单位方向分别为
\[
 \vect e_1=(-\sin\theta,-\cos\theta),\qquad
 \vect e_2=(\cos\theta,-\sin\theta).
\]
软管与其中水的线密度都沿两段均匀分布，整体质心位置为
\[
 \vect r_C=\frac{3L}{4}\vect e_1+\frac L4\vect e_2,\qquad
 x_C=\frac L4(\cos\theta-3\sin\theta).
\]

以软管连同管内的稳态水流为控制体，对 $O$ 点列角动量矩方程。入口就在 $O$，没有角动量矩；出口位置为
$L(\vect e_1+\vect e_2)$，速度为 $u\vect e_2$，故流出的角动量流率为
\[
 \rho Q\,[L(\vect e_1+\vect e_2)]\times(u\vect e_2)
 =\rho QuL\,\vect e_z.
\]
接头反力对 $O$ 无力矩，入口压强力作用线通过 $O$，出口取大气表压为零。因此外力矩只剩总重力矩：
\[
 Mg\frac L4(3\sin\theta-\cos\theta)=\rho QuL.
\]
约去 $L$ 并代入数据，
\[
 3\sin\theta-\cos\theta
 =\frac{4\rho Qu}{Mg}=0.6923.
\]
写成
$\sqrt{10}\sin[\theta-\arctan(1/3)]=0.6923$，取图示的小于 $90^\circ$ 的稳定构形，得到
\[
 \boxed{\theta=31.1^\circ}.
\]
\FluidHoseAnalysisDiagram

这个角度必须大于 $\arctan(1/3)=18.4^\circ$，此后质心位于转轴左侧，重力才提供与流出角动量同向的正力矩；
若误把管内 $0.160\,\mathrm{kg}$ 的水重漏掉，会得到约 $53.1^\circ$，与题给截面积信息的作用不相容。

再作三项核验。第一，等式两边都已除以 $L$：左边 $Mg$ 与右边 $\rho Qu$ 的量纲同为力，
所以无量纲方程没有漏掉长度因子。第二，令
\[
 f(\theta)=3\sin\theta-\cos\theta.
\]
在题图允许的 $0<\theta<90^\circ$ 内，$f'(\theta)=3\cos\theta+\sin\theta>0$，
故 $f$ 单调增加；又有 $f(0)=-1$、$f(90^\circ)=3$，右端 $0.6923$ 落在该区间内，
因此这里恰有一个几何上可接受的解，不会因反三角函数的另一分支产生歧义。

第三，若把流量逐渐降到零，平衡式退化为
$3\sin\theta-\cos\theta=0$，即 $\theta\to\arctan(1/3)$。
这正是直角管连同管内静水仅在重力作用下使质心落到转轴正下方的构形；随着 $Q$ 增大，
出口角动量流率 $\rho Q^2L/A$ 增大，所需重力矩随之增大，因而 $\theta$ 应继续增大。
这个极限与单调性均和求得的 $31.1^\circ$ 一致。
\end{xesolution}

\endgroup

\hypersetup{bookmarksdepth=2}

\hypersetup{bookmarksdepth=3}
\begingroup
\renewcommand{\Needspace}[1]{\par\begingroup
\dimen0=#1\relax\vskip0pt plus\dimen0\penalty-100
\vskip0pt plus-\dimen0\vskip\dimen0\penalty9999\vskip-\dimen0
\vskip0pt\endgroup}
\tcbset{breakable/.default=false}
\chapter{温度、物态、理想气体与分子模型}

{
\graphicspath{{/Users/dongxingxi/Documents/物理竞赛一轮资料/sample/xes-original/twenty-eighth-section/}{/Users/dongxingxi/Documents/物理竞赛一轮资料/sample/xes-original/twenty-eighth-section/assets/}}

% Teaching diagrams. Labels are transparent; geometry is shared where applicable.
\newcommand{\boltzmannRayleighContext}{%
\begin{center}\begin{minipage}{\linewidth}\centering
\includegraphics[width=.9\linewidth]{boltzmann-rayleigh-context.png}
\figcap{左：玻尔兹曼（Ludwig Boltzmann）与分子运动的统计图像；右：瑞利勋爵（Lord Rayleigh）与精密气体测量。}
\label{ch14:fig:history}\end{minipage}\end{center}}
\newcommand{\molecularMicroscopyFigure}{%
\begin{center}\begin{minipage}{\linewidth}\centering
\includegraphics[width=.72\linewidth]{figure19-1-redraw.png}
\figcap{分子尺度成像示意。左、右分别表示局部分子结构与规则排列的分子；此图不表示实验测量数据。}\label{ch14:fig:s28-19-1}
\end{minipage}\end{center}}

\newcommand{\thermometerMeasurementDiagram}{%
\begin{center}\begin{minipage}{\linewidth}\centering
\begin{tikzpicture}[x=1cm,y=1cm,>=Latex]
\draw[axis] (0,0)--(6.6,0) node[right]{$t$};
\draw[axis] (0,0)--(0,3.4) node[above]{$T$};
\draw[DeepBlue,thick,domain=0:6,samples=65] plot(\x,{1.6+1.3*exp(-\x)});
\draw[Ink,thick,domain=0:6,samples=65] plot(\x,{1.6-1.1*exp(-\x)});
\draw[TextGray,dashed] (0,1.6)--(6.2,1.6);
\node[left] at(0,2.9){$T_{\text{物},0}$}; \node[left] at(0,.5){$T_{\text{计},0}$};
\node[above,DeepBlue] at(3,2){被测物降温};
\node[below] at(3,1.1){温度计升温};
\node[right] at(6.3,1.6){$T_f$};
\end{tikzpicture}
\figcap{接触后的温度变化：在隔热的整体中，两者从不同初温趋于共同温度 $T_f$。}\label{ch14:fig:thermometry}
\end{minipage}\end{center}}

\newcommand{\molecularEquilibriumDiagram}{%
\begin{center}\begin{minipage}{\linewidth}\centering
\begin{tikzpicture}[x=.8cm,y=.8cm,>=Latex]
\foreach \offset/\title in {0/{较高温度},4.6/{较低温度},10.2/{共同温度}}{
\draw[Ink,thick] (\offset,0) rectangle ++(3.8,2.8);
\node[above] at(\offset+1.9,3){\title};}
\foreach \xx/\yy/\ang in {.5/.5/35,1.8/.8/110,2.6/1.6/190,.8/2.1/320}{
\fill[Ink] (\xx,\yy) circle(1.5pt);
\draw[->,Ink,thick] (\xx,\yy)--++(\ang:.85);
\fill[Ink] (\xx+4.6,\yy) circle(1.5pt);
\draw[->,Ink,thick] (\xx+4.6,\yy)--++(\ang:.38);
\fill[DeepBlue] (\xx+10.2,\yy) circle(1.5pt);
\draw[->,DeepBlue,thick] (\xx+10.2,\yy)--++(\ang:.6);}
\draw[->,DeepBlue,thick] (8.65,1.4)--(9.95,1.4);
\node[below,align=center] at(4.2,-.2){接触并交换能量};
\node[below,align=center] at(12.1,-.2){单分子能量仍涨落};
\end{tikzpicture}
\figcap{同种气体的热平衡示意：箭头表示典型热运动速度，不表示共同的宏观流速。}\label{ch14:fig:equilibrium}
\end{minipage}\end{center}}

% q=r/r0, U/epsilon=q^-12-2q^-6; exact roots of U/epsilon=E.
\newcommand{\ThermalPotentialAxes}[1]{%
\pgfmathsetmacro{\qleft}{(1+sqrt(1+#1))^(-1/6)}
\pgfmathsetmacro{\qright}{(1-sqrt(1+#1))^(-1/6)}
\draw[axis] (.76,0)--(1.93,0) node[right]{$r/r_0$};
\draw[axis] (.76,-1.2)--(.76,1.0) node[above]{$U/\varepsilon$};
\begin{scope}\clip (.79,-1.08) rectangle (1.87,.88);
\draw[DeepBlue,very thick,domain=.80:1.86,samples=180,smooth] plot(\x,{(\x)^(-12)-2*(\x)^(-6)});
\end{scope}
\draw[TextGray,dashed] (.8,#1)--(1.83,#1);
\fill[Ink] (\qleft,#1) circle(.4pt) (\qright,#1) circle(.4pt);
\draw[TextGray,dashed] (1,-1)--(1,0);
\node[above] at(1,0){$1$};
\node[left] at(.76,-1){$-1$};
\node[right,font=\small] at(1.84,#1){$-1+E_T/\varepsilon$};}
\newcommand{\anharmonicPotentialDiagram}{%
\begin{center}\begin{minipage}{\linewidth}\centering
\begin{tikzpicture}[x=5.5cm,y=1.65cm,>=Latex]
\ThermalPotentialAxes{-.5}
\draw[TextGray,dashed] (\qleft,-.5)--(\qleft,-1.38);
\draw[TextGray,dashed] (\qright,-.5)--(\qright,-1.38);
\draw[<->,DeepBlue] (\qleft,-1.3)--(\qright,-1.3);
\node[below,font=\small] at({(\qleft+\qright)/2},-1.32){$r_2-r_1=2A_T$};
\node[font=\small] at(.86,-.18){$r_1$};
\draw[TextGray,thin] (.88,-.23)--(\qleft,-.5);
\node[font=\small] at(1.4,-.98){$r_2$};
\draw[TextGray,thin] (1.37,-.9)--(\qright,-.5);
\end{tikzpicture}
\figcap{非对称势阱与同一能量水平的两个真实交点。横坐标按 $r_0$ 归一化；$A_T$ 只表示两转折点间距的一半。}\label{ch14:fig:anharmonic}
\end{minipage}\end{center}}

\newcommand{\gasWallCollisionDiagram}{%
\begin{center}\begin{minipage}{\linewidth}\centering
\begin{tikzpicture}[x=1cm,y=1cm,>=Latex]
\fill[DeepBlue!10] (1,0) rectangle (4,2.4);
\draw[Ink,line width=1.2pt] (4,-.4)--(4,2.9);
\draw[TextGray,dashed] (1,0)--(4,0) (1,2.4)--(4,2.4) (1,0)--(1,2.4);
\foreach \xx/\yy in {1.3/.5,2.3/1.8,2.7/.9}{
\fill[Ink] (\xx,\yy) circle(1.5pt); \draw[->,DeepBlue] (\xx,\yy)--++(.8,0);}
\draw[<->] (1,-.65)--node[below]{$v_x\Delta t$}(4,-.65);
\node[above] at(2.5,3.2){朝向壁面的候选分子层};
\node[align=left,anchor=west] at(4.5,1.2){壁面面积 $S$\\层体积 $Sv_x\Delta t$\\单次壁面冲量 $2mv_x$\\此处 $v_x>0$};
\end{tikzpicture}
\figcap{给定正法向速度的分子在 $\Delta t$ 内能够到壁面的薄层；图为法向剖面。}\label{ch14:fig:collision-layer}
\end{minipage}\end{center}}

\newcommand{\traditionalBarometerDiagram}{%
\begin{center}\begin{minipage}{\linewidth}\centering
\begin{tikzpicture}[x=.85cm,y=.85cm,>=Latex]
\ThermalBarometerBase{3.5}{4.6}
\node[anchor=west,align=left] at(3,2.1){管顶残余气体压强 $p_g$\\[4pt]$p_{\rm atm}=\rho_{\rm Hg}gH+p_g$\\[4pt]真空近似：$p_g\ll p_{\rm atm}$};
\end{tikzpicture}
\figcap{水银气压计的压力平衡。高度 $H$ 从汞槽自由液面量到管内液面，不从槽底起算。}\label{ch14:fig:barometer-principle}
\end{minipage}\end{center}}

\newcommand{\rayleighApparatusDiagram}{%
\begin{center}\begin{minipage}{\linewidth}\centering
\begin{tikzpicture}[x=.9cm,y=.9cm,>=Latex]
\foreach \sx/\lab in {0/B,5/C}{
\begin{scope}[shift={(\sx,0)}]
\fill[DeepBlue!22] (-.65,1.8)--(-.25,1.65)--(-.25,.7)--(.25,.7)--(.25,1.65)--(.65,1.8)--(.65,0)--(-.65,0)--cycle;
\draw[Ink,thick] (-.65,2.6)--(-.65,0)--(.65,0)--(.65,2.6);
\draw[Ink,thick] (-.25,3)--(-.25,.7)--(.25,.7)--(.25,3);
\draw[DeepBlue,thick] (-.65,1.8)..controls(-.45,1.8)and(-.3,1.72)..(-.25,1.65);
\draw[DeepBlue,thick] (.25,1.65)..controls(.3,1.72)and(.45,1.8)..(.65,1.8);
\draw[TextGray,dashed] (-1,.7)--(1,.7);
\node[below] at(0,-.3){$\lab$ 处的局部剖面};
\end{scope}}
\node[align=center] at(2.5,2.2){同形玻璃元件\\同种液体\\相同接触状态};
\draw[<->,DeepBlue] (1.3,.7)--(3.7,.7);
\end{tikzpicture}
\figcap{两处同形叉状汞面与标志的比较原理。放大图把两处移到同一基准高度，以比较几何形状；实际装置的高度差见前图。}\label{ch14:fig:rayleigh-forks}
\end{minipage}\end{center}}

\newcommand{\heatFeelingDiagram}{%
\begin{center}\begin{minipage}{\linewidth}\centering
\begin{tikzpicture}[x=.9cm,y=.9cm,>=Latex]
\foreach \sx/\title/\desc in {0/{冷热水后摸温水}/{初态与热流方向},4/{湿冷环境}/{多种换热因素},8/{泡脚后打喷嚏}/{单次观察不定因},12/{高温环境}/{不能据此认定安全}}{
\node[align=center] at(\sx,1){\title};
\node[align=center,font=\small,TextGray] at(\sx,0){\desc};}
\end{tikzpicture}
\figcap{四个情境需要分别判断：热感觉、换热过程与安全防护不是同一问题。}\label{ch14:fig:heat-feeling}
\end{minipage}\end{center}}

\newcommand{\surfacePhenomenaDiagram}{%
\begin{center}\begin{minipage}{\linewidth}\centering
\begin{tikzpicture}[x=.85cm,y=.85cm,>=Latex]
\begin{scope}[shift={(-4.8,0)}]
\fill[DeepBlue!15] (-1,-1)--(-1,.3)..controls(-.55,1.1)and(.55,1.1)..(1,.3)--(1,-1)--cycle;
\draw[Ink,thick] (-1,-1.2)--(-1,1.8) (1,-1.2)--(1,1.8);
\draw[DeepBlue,thick] (-1,.3)..controls(-.55,1.1)and(.55,1.1)..(1,.3);
\draw[TextGray,dashed] (0,.9)--(2,.9);\node[right,font=\small] at(2,.9){水平视线};
\node[above] at(0,2){凸液面};
\end{scope}
\begin{scope}
\draw[Ink,thick] (-.22,1.6)--(-.22,.65) (.22,1.6)--(.22,.65);
\fill[DeepBlue!15,draw=DeepBlue,thick] (-.22,.65)..controls(-.2,.25)and(-.75,-.15)..(-.65,-.55)..controls(-.5,-1.3)and(.5,-1.3)..(.65,-.55)..controls(.75,-.15)and(.2,.25)..(.22,.65)--cycle;
\draw[->,Ink] (0,-.5)--(0,-1.8) node[below]{$mg$};
\draw[->,DeepBlue] (-.22,.6)--(-.22,1.4);
\draw[->,DeepBlue] (.22,.6)--(.22,1.4);
\node[right,font=\small] at(.6,1.2){张力竖直分量};
\node[above] at(0,2){悬在管口的液滴};
\end{scope}
\begin{scope}[shift={(5.6,0)}]
\draw[Ink,line width=2pt] (-1.05,1.5)--(.75,-1.2);
\draw[DeepBlue,line width=1.8pt] (-2,1.2)..controls(-1.4,1.15)and(-1.1,1.4)..(-.95,1.17)--(.77,-1.4);
\draw[->,DeepBlue] (-.4,.05)--(.1,-.7);
\draw[Ink,thick] (-.05,-.6)--(-.05,-1.8)--(2,-1.8)--(2,-.6);
\node[above] at(0,2){液流贴着玻璃棒一侧};
\end{scope}
\end{tikzpicture}
\figcap{凸弯月面、与管口相连的液滴及沿玻璃棒一侧流下的液流。蓝色液流不穿过玻璃棒。}\label{ch14:fig:surface-phenomena}
\end{minipage}\end{center}}

% Openings in the circular bulbs are the same intersections used by the tube walls.
\newcommand{\ThermalBulbOutline}{%
\pgfmathsetmacro{\ang}{asin(.15/.8)}\pgfmathsetmacro{\dx}{sqrt(.8*.8-.15*.15)}
\draw[Ink,thick] ({-2+\dx},.15) arc[start angle=\ang,end angle={360-\ang},radius=.8];
\draw[Ink,thick] ({2-\dx},-.15) arc[start angle={180+\ang},end angle={540-\ang},radius=.8];
\draw[Ink,thick] ({-2+\dx},.15)--({2-\dx},.15) ({-2+\dx},-.15)--({2-\dx},-.15);}
\newcommand{\bulbThermometerDiagram}{%
\begin{center}\begin{minipage}{\linewidth}\centering
\begin{tikzpicture}[x=.85cm,y=.85cm,>=Latex]
\begin{scope}[shift={(-4.2,0)}]
\ThermalBulbOutline
\fill[DeepBlue] (-.09,-.15) rectangle (.09,.15);
\node[above] at(0,1.1){水平：活塞居中};
\node at(-2,0){$V_0$};\node at(2,0){$V_0$};
\end{scope}
\begin{scope}[shift={(2.5,0)}]
\begin{scope}[rotate=90]\ThermalBulbOutline\end{scope}
\fill[DeepBlue] (-.15,-.54) rectangle (.15,-.36);
\draw[TextGray,dashed] (-.7,0)--(.9,0);
\draw[{Latex[length=1.6pt,width=2pt]}-{Latex[length=1.6pt,width=2pt]}] (.75,0)--(.75,-.45);\node[right] at(.9,-.23){$h$};
\node[right] at(1,2){$V_0+Sh,\ p_{\text{上}}$};
\node[right] at(1,-2){$V_0-Sh,\ p_{\text{下}}$};
\draw[->,Ink] (-.65,-.45)--(-.65,-1.1) node[left]{$mg$};
\node[above] at(0,3){竖直：活塞下移};
\end{scope}
\end{tikzpicture}
\figcap{两球通过细管真正连通；活塞只在直管段内滑动，密封整个管截面。}\label{ch14:fig:bulb-thermometer}
\end{minipage}\end{center}}

\newcommand{\torricelliDiagram}{%
\begin{center}\begin{minipage}{\linewidth}\centering
\begin{tikzpicture}[x=.85cm,y=.85cm,>=Latex]
\foreach \sx/\height/\idx/\gasLen in {-3.2/3.7/1/8,3.2/3.6/2/10}{
\begin{scope}[shift={(\sx,0)}]
\ThermalBarometerBase{\height}{4.1}
\draw[{Latex[length=1.6pt,width=2pt]}-{Latex[length=1.6pt,width=2pt]}] (-.9,\height)--(-.9,4.1);
\node[above left,font=\small] at(-.9,4.1){$\ell_{\idx}=\gasLen\,\mathrm{cm}$};
\node[below] at(0,-1){状态 \idx};
\end{scope}}
\end{tikzpicture}
\figcap{同一固定管顶、同一汞槽液面基准。左：$H_1=74\,\mathrm{cm}$；右：$H_2=72\,\mathrm{cm}$，残余气柱由8厘米增至10厘米。}\label{ch14:fig:torricelli}
\end{minipage}\end{center}}

\newcommand{\heatedCylinderDiagram}{%
\begin{center}\begin{minipage}{\linewidth}\centering
\begin{tikzpicture}[x=1cm,y=1cm,>=Latex]
\begin{axis}[width=9cm,height=6cm,xmin=.97,xmax=1.14,ymin=.87,ymax=1.14,
xlabel={$V/V_0$},ylabel={$p/p_0$},xtick={1,1.1},ytick={.9,1,1.1},grid=major,
scaled ticks=false,tick label style={/pgf/number format/fixed},clip=false]
\addplot[DeepBlue,very thick,->] coordinates{(1,.9)(1,1)};
\addplot[DeepBlue,very thick,->] coordinates{(1,1)(1.1,1)};
\addplot[DeepBlue,very thick,->] coordinates{(1.1,1)(1.1,1.1)};
\node[left] at(axis cs:1,.95){定容};
\node[above] at(axis cs:1.05,1){定压};
\node[right] at(axis cs:1.1,1.05){定容};
\end{axis}
\end{tikzpicture}
\figcap{缓慢加热的三个阶段。坐标只显示 $V/V_0$ 与 $p/p_0$ 的局部范围，刻度在各自轴上均为线性。}\label{ch14:fig:heated-path}
\end{minipage}\end{center}}

\newcommand{\sandPistonDiagram}{%
\begin{center}\begin{minipage}{\linewidth}\centering
\begin{tikzpicture}[x=.85cm,y=.85cm,>=Latex]
\foreach \sx/\height/\idx/\sand in {-4.2/1.5/0/1,0/2.25/1/1,4.2/2.7/2/0}{
\begin{scope}[shift={(\sx,0)}]
\draw[Ink,thick] (-1,3.6)--(-1,0)--(1,0)--(1,3.6);
\draw[Ink,line width=2pt] (-1,1.48)--(-.8,1.48) (.8,1.48)--(1,1.48);
\filldraw[fill=SoftGray,draw=Ink] (-1,\height) rectangle (1,\height+.14);
\ifnum\sand=1\fill[TextGray!45] (-.78,\height+.14)--(0,\height+.65)--(.78,\height+.14)--cycle;\fi
\draw[<->] (-1.35,0)--node[left]{$H_{\idx}$}(-1.35,\height);
\node at(0,.65){气体};
\node[below,align=center,font=\small] at(0,-.3){\ifnum\idx=0 初态及起动前\else\ifnum\idx=1 砂未取走\else 砂已取尽\fi\fi};
\end{scope}}
\end{tikzpicture}
\figcap{三个状态的几何比例：$H_1=1.5H_0$，$H_2=1.8H_0$；活塞密封气缸，铁砂直接承托在活塞上。}\label{ch14:fig:sand-states}
\end{minipage}\end{center}}

\newcommand{\mercurySpillDiagram}{%
\begin{center}\begin{minipage}{\linewidth}\centering
\begin{tikzpicture}[x=.85cm,y=.85cm,>=Latex]
\foreach \sx/\air/\hg/\title in {-4/1.75/.7/{初态},0/2.8/.7/{刚到管口},4/3.5/0/{全部溢出}}{
\begin{scope}[shift={(\sx,0)}]
\fill[DeepBlue!30] (-.3,\air) rectangle (.3,\air+\hg);
\draw[Ink,thick] (-.3,3.5)--(-.3,0)--(.3,0)--(.3,3.5);
\node[below] at(0,-.2){\title};
\end{scope}}
\begin{scope}[shift={(-3,-5.7)}]
\begin{axis}[width=8cm,height=4.5cm,xmin=0,xmax=20,ymin=450,ymax=490,
xtick={0,12,20},ytick={456,480,484},xlabel={$x/\mathrm{cm}$},ylabel={$T/\mathrm K$},
scaled ticks=false,grid=major,xmajorgrids=false,clip=false]
\addplot[DeepBlue,thick,domain=0:20,samples=60]{(76+x)*(100-x)/16};
\addplot[only marks,mark=*] coordinates{(12,484)};
\node[below] at(axis cs:12,478){最高点};
\draw[thin,Ink] (axis cs:12,479)--(axis cs:12,483);
\draw[->,Ink] (axis cs:19,461)--(axis cs:15,461);
\node[above,font=\small] at(axis cs:17,461){溢出方向};
\end{axis}
\end{scope}
\end{tikzpicture}
\figcap{水银到达管口后，剩余长度 $x$ 从20厘米减至0；所需温度在 $x=12\,\mathrm{cm}$ 处最大。}\label{ch14:fig:mercury-spill}
\end{minipage}\end{center}}

\newcommand{\lennardJonesDiagram}{%
\begin{center}\begin{minipage}{\linewidth}\centering
\begin{tikzpicture}[x=5.5cm,y=1.6cm,>=Latex]
\ThermalPotentialAxes{-.8}
\draw[TextGray,densely dashed,domain=.86:1.15,samples=60] plot(\x,{-1+36*(\x-1)^2});
\node[anchor=west,font=\small,align=left] at(1.45,.65){实线：完整势能\\虚线：局部二阶近似};
\end{tikzpicture}
\figcap{势阱最低点为 $(r_0,-\varepsilon)$；二阶近似关于 $r_0$ 对称，而三次项改变左右两侧的形状。}\label{ch14:fig:lennard-jones}
\end{minipage}\end{center}}

% Source figure roles retained; quantitative geometry independently constructed.
\newcommand{\figureNineteenTwoFaithful}{%
\begin{center}\begin{minipage}{\linewidth}\centering
\includegraphics[width=.9\linewidth]{thermometer-contact.png}
\figcap{左侧温度计尚未接触水；右侧温包浸入水中且不接触容器底，等待热平衡后读数。}\label{ch14:fig:s28-19-2}
\end{minipage}\end{center}}

\newcommand{\figureNineteenThreeFaithful}{%
\begin{center}\begin{minipage}{\linewidth}\centering
\begin{tikzpicture}[x=5.5cm,y=1.6cm,>=Latex]
\ThermalPotentialAxes{-.5}
\pgfmathsetmacro{\qa}{(1+sqrt(.2))^(-1/6)}
\pgfmathsetmacro{\qb}{(1-sqrt(.2))^(-1/6)}
\draw[TextGray,dashed] (.8,-.8)--(1.8,-.8);
\fill[Ink] (\qa,-.8) circle(.4pt) (\qb,-.8) circle(.4pt);
\draw[DeepBlue,->] (\qa,-1.37)--(\qleft,-1.37);
\draw[DeepBlue,->] (\qb,-1.37)--(\qright,-1.37);
\draw[TextGray,dotted] (\qleft,-.5)--(\qleft,-1.35) (\qright,-.5)--(\qright,-1.35);
\node[below,font=\small] at(1.14,-1.4){能量升高时，右端移动更多};
\node[right,font=\small] at(1.84,-.8){较低能量};
\end{tikzpicture}
\figcap{非对称势阱中的两组转折点。平均间距必须按运动时间或统计权重计算，不能把左右转折点的中点直接当作平均值。}\label{ch14:fig:s28-19-3}
\end{minipage}\end{center}}

\newcommand{\ThermalBarometerBase}[2]{%
\fill[DeepBlue!23] (-1.5,-.65) rectangle (1.5,0);
\fill[DeepBlue!23] (-.28,-.5) rectangle (.28,#1);
\draw[Ink,thick] (-1.5,.35)--(-1.5,-.65)--(1.5,-.65)--(1.5,.35);
\draw[Ink,thick] (-.28,-.4)--(-.28,#2)--(.28,#2)--(.28,-.4);
\draw[DeepBlue,thick] (-1.5,0)--(-.28,0) (.28,0)--(1.5,0) (-.28,#1)--(.28,#1);
\draw[TextGray,dashed] (.35,#1)--(1.05,#1);
\draw[<->] (.9,0)--node[right]{$H$}(.9,#1);}
\newcommand{\figureNineteenFourFaithful}{%
\begin{center}\begin{minipage}{\linewidth}\centering
\begin{tikzpicture}[x=.85cm,y=.85cm,>=Latex]
\ThermalBarometerBase{3.5}{4.6}
\node[font=\small,anchor=west] at(.7,4.1){近似真空};
\draw[Ink,thin] (.3,4.05)--(.6,4.1);
\draw[Ink] (2,0)--(2,4.2);
\foreach \height in {0,.5,...,4}{\draw[Ink] (2,\height)--(2.15,\height);}
\draw[Ink,thin] (2.15,1.8)--(2.5,1.8);
\node[right] at(2.5,1.8){读数尺};
\node[left] at(-1.55,-.3){汞槽};
\end{tikzpicture}
\figcap{传统水银气压计。管口浸入汞槽，管内汞柱与槽中水银连续连通。}\label{ch14:fig:s28-19-4}
\end{minipage}\end{center}}

\newcommand{\figureNineteenFiveFaithful}{%
\begin{center}\begin{minipage}{\linewidth}\centering
\begin{tikzpicture}[x=.85cm,y=.85cm,>=Latex]
% Same fork-shaped chambers at different elevations. Gas spaces are not joined.
\foreach \yy/\lab in {0/B,4.3/C}{
\begin{scope}[shift={(1,\yy)}]
\fill[DeepBlue!22] (-.65,1.5)--(-.23,1.38)--(-.23,.65)--(.23,.65)--(.23,1.38)--(.65,1.5)--(.65,0)--(-.65,0)--cycle;
\draw[Ink,thick] (-.65,2.1)--(-.65,0)--(.65,0)--(.65,2.1);
\draw[Ink,thick] (-.65,2.1)--(-.23,2.1) (.23,2.1)--(.65,2.1);
\draw[Ink,thick] (-.23,2.35)--(-.23,.65)--(.23,.65)--(.23,2.35);
\draw[DeepBlue] (-.65,1.5)..controls(-.4,1.5)and(-.25,1.4)..(-.23,1.38);
\draw[DeepBlue] (.23,1.38)..controls(.25,1.4)and(.4,1.5)..(.65,1.5);
\node[left] at(-.95,.65){$\lab$};
\end{scope}}
% Hydraulic line, shown as a connection, distinct from measurement chambers.
\draw[DeepBlue,line width=1.5pt] (1.6,.2)--(3,.2)--(3,4.5)--(1.6,4.5);
\node[right,align=left,font=\small] at(3.2,3.6){连通汞路\\可调储槽};
\draw[DeepBlue] (3,1.1)--(4,1.1);
\draw[Ink,thick] (3.65,1.5)--(3.65,.7)--(4.5,.7)--(4.5,1.5);
\draw[DeepBlue,line width=2pt] (3.67,1.1)--(4.48,1.1);
% Sample gas to lower chamber; upper gas has its own vacuum valve.
\draw[Ink,thick] (-3.1,1.6) circle(.65);
\draw[Ink,thick] (-2.45,1.6)--(.35,1.9);
\node[below,align=center,font=\small] at(-3.1,.6){待测气体球\\低温恒温环境};
\draw[Ink,thick] (.35,6.2)--(-.7,6.2)--(-.7,7);
\draw[Ink] (-.9,6.6)--(-.5,6.9) (-.9,6.9)--(-.5,6.6);
\node[left] at(-.9,6.75){$D$};
\node[left,font=\small] at(-.9,7.2){抽气、检验真空};
\draw[TextGray,dashed] (1.3,.65)--(2.5,.65) (1.3,4.95)--(2.5,4.95);
\draw[<->] (2.3,.65)--node[left]{$h$}(2.3,4.95);
\end{tikzpicture}
\figcap{瑞利精密气压计的测量结构示意。$B,C$ 为同形叉状端面，$D$ 位于真空参考气路；蓝线表示汞路，黑线表示气路。}\label{ch14:fig:s28-19-5}
\end{minipage}\end{center}}

\newcommand{\figureNineteenSixFaithful}{%
\begin{center}\begin{minipage}{\linewidth}\centering
\begin{tikzpicture}[x={(.9cm,0cm)},y={(.38cm,.24cm)},z={(0cm,.9cm)},>=Latex]
\begin{scope}[shift={(-4,0)}]
\draw[TextGray,dashed] (0,2.5,0)--(0,0,0) (0,2.5,0)--(2.5,2.5,0) (0,2.5,0)--(0,2.5,2.5);
\draw[Ink,thick] (0,0,0)--(2.5,0,0)--(2.5,0,2.5)--(0,0,2.5)--cycle;
\draw[Ink,thick] (2.5,0,0)--(2.5,2.5,0)--(2.5,2.5,2.5)--(2.5,0,2.5);
\draw[Ink,thick] (0,0,2.5)--(0,2.5,2.5)--(2.5,2.5,2.5);
\draw[->] (0,0,0)--(3.2,0,0) node[right]{$x$};
\draw[->] (2.5,0,0)--(2.5,3.4,0) node[right]{$y$};
\draw[->] (0,0,0)--(0,0,3.3) node[above]{$z$};
\fill[DeepBlue] (.7,.7,.8) circle(1.6pt);
\draw[->,DeepBlue,thick] (.7,.7,.8)--(1.9,1.1,1.8) node[above]{$\vect v$};
\draw[->,Ink] (.7,.7,.8)--(1.9,.7,.8);
\draw[TextGray,thin] (1.3,.7,.8)--(-.35,0,.55);
\node[left=3pt] at(-.35,0,.55){$v_x$};
\node[below] at(1.25,0,-.55){边长 $d$；右手坐标系};
\end{scope}
\begin{scope}[x=1cm,y=1cm,shift={(2.3,0)}]
\draw[Ink,thick] (3.3,-.2)--(3.3,3.7);
\draw[->,DeepBlue,thick] (.3,.3)--(3.3,1.8);
\draw[->,DeepBlue,thick] (3.3,1.8)--(.3,3.3);
\draw[TextGray,dashed] (1.8,1.8)--(3.3,1.8);
\node[below right] at(1.1,.7){$\vect v=(v_x,v_z)$};
\node[above right] at(.3,3.3){$\vect v'=(-v_x,v_z)$};
\fill[Ink] (3.3,1.8) circle(1.5pt);
\node[right] at(3.5,1.8){$Q$};
\node[below,font=\small] at(1.5,-.5){$xz$ 剖面，$v_y$ 未画出};
\end{scope}
\end{tikzpicture}
\figcap{左为立方容器中的速度分解；右为垂直于 $x$ 轴的壁面反射，法向速度反号、切向速度不变，两线与法线夹角相等。}\label{ch14:fig:s28-19-6}
\end{minipage}\end{center}}

\newcommand{\exampleNineteenSixFaithful}{%
\begin{center}\begin{minipage}{\linewidth}\centering
\begin{tikzpicture}[x=1cm,y=1cm,>=Latex]
\draw[Ink,thick] (0,1.4)--(0,0)--(4.5,0) (0,1.4)--(4.5,1.4);
\draw[Ink,line width=3pt] (2.88,0)--(2.88,.22) (3.42,1.4)--(3.42,1.18);
\filldraw[fill=SoftGray,draw=Ink] (2.95,0) rectangle (3.05,1.4);
\draw[TextGray,dashed] (3.3,0)--(3.3,1.4);
\draw[->] (2.5,1.9)--(3,1.9);
\draw[->] (3.8,1.9)--(3.3,1.9);
\node[above,align=center,font=\small] at(3.15,2){$AB$ 间容积\\$0.1V_0$};
\draw[TextGray,thin] (3,-.08)--(2.6,-.5);
\draw[TextGray,thin] (3.3,-.08)--(3.7,-.5);
\node[below] at(2.6,-.5){$B$};\node[below] at(3.7,-.5){$A$};
\node at(1.4,.7){$V_0,\ 0.9p_0$};
\node at(4,.7){$p_0$};
\begin{scope}[shift={(6.1,-.2)}]
\begin{axis}[width=5.8cm,height=4.3cm,xmin=.97,xmax=1.14,ymin=.87,ymax=1.24,
xlabel={$V/V_0$},ylabel={$p/p_0$},xtick={1,1.1},ytick={.9,1,1.1,1.2},grid=major,
scaled ticks=false,tick label style={font=\scriptsize,/pgf/number format/fixed}]
\end{axis}
\end{scope}
\end{tikzpicture}
\figcap{左：活塞初始在 $B$ 处；两个挡块限制左、右移动。右：待画过程线的局部 $p$--$V$ 坐标，坐标轴均为线性刻度。}\label{ch14:fig:heated-apparatus}
\end{minipage}\end{center}}

\newcommand{\figureNineteenSevenFaithful}{%
\begin{center}\begin{minipage}{\linewidth}\centering
\begin{tikzpicture}[x=.9cm,y=.9cm,>=Latex]
\draw[Ink,thick] (-1.2,3.5)--(-1.2,0)--(1.2,0)--(1.2,3.5);
\draw[Ink,line width=2pt] (-1.2,1.48)--(-1,1.48) (1,1.48)--(1.2,1.48);
\filldraw[fill=SoftGray,draw=Ink] (-1.2,1.5) rectangle (1.2,1.68);
\fill[TextGray!45] (-1,1.68)--(0,2.35)--(1,1.68)--cycle;
\node at(0,.7){$p_0,T_0$};
\draw[<->] (-1.65,0)--node[left]{$H_0$}(-1.65,1.5);
\node[right] at(1.7,1.55){卡环承托活塞};
\draw[TextGray] (1.1,1.47)--(1.65,1.55);
\node[right] at(1.7,2.2){铁砂};
\end{tikzpicture}
\figcap{初态气体压强等于大气压；卡环承担铁砂的重量，活塞与缸壁之间没有漏气间隙。}\label{ch14:fig:s28-19-7}
\end{minipage}\end{center}}

\newcommand{\figureNineteenEightFaithful}{%
\begin{center}\begin{minipage}{\linewidth}\centering
\begin{tikzpicture}[x=.9cm,y=.9cm,>=Latex]
\fill[DeepBlue!30] (-.3,2) rectangle (.3,2.8);
\draw[Ink,thick] (-.3,4)--(-.3,0)--(.3,0)--(.3,4);
\draw[<->] (.8,0)--node[right]{$a=50\,\mathrm{cm}$}(.8,2);
\draw[<->] (.8,2)--node[right]{$h=20\,\mathrm{cm}$}(.8,2.8);
\draw[<->] (-2.8,0)--node[left]{$L=100\,\mathrm{cm}$}(-2.8,4);
\draw[Ink,thin] (0,.8)--(0,-.25);
\node[below] at(0,-.35){封闭空气};
\node[above] at(0,4.2){开口，外压 $p_0$};
\end{tikzpicture}
\figcap{管内空气、汞柱及管口的位置按 $50:20:100$ 的长度比例绘制。}\label{ch14:fig:s28-19-8}
\end{minipage}\end{center}}

\newcommand{\ThermalUTubeState}[1]{%
\pgfmathsetmacro{\leftlevel}{2+#1}
\pgfmathsetmacro{\rightlevel}{2-#1}
\fill[DeepBlue!18] (-1.7,\leftlevel)--(-1.1,\leftlevel)--(-1.1,0) arc(180:360:1.1)--(1.1,\rightlevel)--(1.7,\rightlevel)--(1.7,0) arc(0:-180:1.7)--cycle;
\draw[Ink,thick] (-1.7,4)--(-1.7,0) arc(180:360:1.7)--(1.7,4)--(1.1,4)--(1.1,0) arc(360:180:1.1)--(-1.1,4)--cycle;
\draw[DeepBlue,thick] (-1.7,\leftlevel)--(-1.1,\leftlevel) (1.1,\rightlevel)--(1.7,\rightlevel);
\draw[<->] (-2.1,\leftlevel)--(-2.1,4);
\draw[<->] (2.1,\rightlevel)--(2.1,4);}
\newcommand{\closedUTubeDiagram}{%
\begin{center}\begin{minipage}{\linewidth}\centering
\begin{tikzpicture}[x=.75cm,y=.75cm,>=Latex]
\begin{scope}[shift={(-3.7,0)}]
\ThermalUTubeState{0}
\node[left] at(-2.1,3){$L$};\node[right] at(2.1,3){$L$};
\node[above] at(0,4.9){初态：两端封闭};
\node[above,font=\small] at(-1.4,4.15){$p_0,T_0$};
\node[above,font=\small] at(1.4,4.15){$p_0,T_0$};
\end{scope}
\begin{scope}[shift={(3.7,0)}]
\ThermalUTubeState{.4}
\node[left] at(-2.1,3.2){$L-x$};\node[right] at(2.1,2.8){$L+x$};
\draw[TextGray,dashed] (-1.1,2.4)--(.4,2.4) (.4,1.6)--(1.1,1.6);
\draw[<->] (.25,1.6)--node[left]{$2x$}(.25,2.4);
\node[above] at(0,4.9){末态：右侧受热};
\node[above,font=\small] at(-1.4,4.15){$p_L,T_0$};
\node[above,font=\small] at(1.4,4.15){$p_R,T_R$};
\end{scope}
\end{tikzpicture}
\figcap{水连续充满两腿及弯管；加热后左液面上升、右液面下降，气体仍各自封闭。}\label{ch14:fig:closed-u}
\end{minipage}\end{center}}

% Ideal narrow L pipe. Filled mercury is clipped to its same constant-width boundary.
\newcommand{\ThermalLPipe}[1]{%
\begin{scope}
\clip (-2,-.12)--(2.62,-.12)--(2.62,2)--(2.38,2)--(2.38,.12)--(-2,.12)--cycle;
\pgfmathsetmacro{\tailx}{min(#1,2.5)}
\pgfmathsetmacro{\taily}{max(0,#1-2.5)}
\pgfmathsetmacro{\headx}{min(#1+1,2.5)}
\pgfmathsetmacro{\heady}{min(2,max(0,#1+1-2.5))}
\pgfmathparse{#1<2.5?1:0}
\ifnum\pgfmathresult=1
\pgfmathparse{#1+1<2.5?1:0}
\ifnum\pgfmathresult=1
\fill[DeepBlue!45] (\tailx,-.12) rectangle (\headx,.12);
\else
\fill[DeepBlue!45] (\tailx,-.12)--(2.62,-.12)--(2.62,\heady)--(2.38,\heady)--(2.38,.12)--(\tailx,.12)--cycle;
\fi
\else
\fill[DeepBlue!45] (2.38,\taily) rectangle (2.62,\heady);
\fi
\end{scope}
\draw[Ink,thick] (-2,.12)--(2.38,.12)--(2.38,2);
\draw[Ink,thick] (-2,.12)--(-2,-.12)--(2.62,-.12)--(2.62,2);
\draw[TextGray,dashed] (0,-.45)--(0,.55);
\node[below,font=\scriptsize] at(0,-.48){$s=0$};}
\newcommand{\lPipeInitialDiagram}{%
\begin{center}\begin{minipage}{\linewidth}\centering
\begin{tikzpicture}[x=1cm,y=1cm,>=Latex]
\ThermalLPipe{0}
\draw[<->] (-2,-1.15)--node[below]{$a$}(0,-1.15);
\draw[<->] (0,-1.15)--node[below]{$b$}(1,-1.15);
\draw[<->] (0,-1.75)--node[below]{$c$}(2.5,-1.75);
\draw[<->] (3.1,0)--node[right]{$d$}(3.1,2);
\node[above] at(-1,.3){封闭气体};
\node[above] at(.5,.3){水银};
\node[above] at(2.5,2.2){出口 $s=c+d$};
\end{tikzpicture}
\figcap{初始位置与沿管中心线坐标。弯头体积和曲率半径相对直管长度忽略不计。}\label{ch14:fig:l-initial}
\end{minipage}\end{center}}
\newcommand{\lPipeStagesDiagram}{%
\begin{center}\begin{minipage}{\linewidth}\centering
\begin{tikzpicture}[x=.9cm,y=.9cm,>=Latex]
\foreach \sx/\sy/\xx/\title in {0/0/.5/{水平段},7/0/2/{跨越弯头},0/-4/2.8/{竖直段},7/-4/4/{逐渐流出}}{
\begin{scope}[shift={(\sx,\sy)}]
\ThermalLPipe{\xx}
\node[below] at(.3,-1){\title};
\fill[Ink] ({min(\xx,2.5)},{max(0,\xx-2.5)}) circle(1pt);
\end{scope}}
\end{tikzpicture}
\figcap{四阶段水银位置。黑点为水银后端，其沿程坐标等于 $x$；流出阶段管内只保留长度 $c+d-x$ 的水银。}\label{ch14:fig:l-stages}
\end{minipage}\end{center}}


热物理研究热、温度以及物质宏观状态的变化。日常生活中经常使用温度，但把温度作为基本物理量写入
定量规律，是热学区别于纯力学的重要特征。十九世纪中叶以前，热学与动力学长期沿着两条相对独立的
道路发展，能量守恒也多半只在特定力学系统中使用。焦耳等人的实验逐渐表明，伴随热现象发生的
能量转移与机械功没有本质区别，热只是能量转移的一种形式。此后，分子运动论和统计物理把宏观温度
同微观粒子的无规则运动联系起来。连接两种尺度的基本常量之一是玻尔兹曼常数
\begin{equation}
  \boxed{k_B=\SI{1.380649e-23}{J.K^{-1}}}.
\end{equation}
它的量纲是“能量/温度”，意味着 $k_BT$ 正是温度 $T$ 所对应的微观能量尺度。

从历史次序看，这一理论框架并不是一开始就完整出现的。人们先在热机、量热和气体实验中积累经验，
再逐步确认机械功与热量能够相互转化；此后近半个世纪中，分子运动论、玻尔兹曼的统计方法以及对
微观粒子存在性的实验检验相互推进，才使“热和温度具有微观基础”成为可以定量使用的物理理论。
所以本讲既从宏观可测量的温度、压强出发，也反复回到分子的运动和碰撞，说明两种描述如何对应。

图\ref{ch14:fig:history}把统计理论与精密气体实验两条研究路径并列展示。
\boltzmannRayleighContext

\section{分子运动论}\label{ch14:sec:molecules}

化学课程往往从原子、分子出发，热物理也可以采用同样的路径：宏观物质由大量微观粒子构成，粒子
遵循力学和电磁学的基本规律，而温度、压强、密度等宏观量则是大量粒子共同运动的统计结果。这样的
观点避开了早期“热质”等历史概念，直接把可观察的热现象放进微观动力学框架中。

这里可以作一个很鲜明的课程比较：化学课一开始就引入原子、分子，而初等物理通常从牛顿
力学开始。如果物理也完全按现代基本理论的层次起步，第一课就不得不谈相对论、量子论、粒子物理和
宇宙学标准模型，这显然超出了现阶段所能掌握的范围。热学却有一个便利之处：我们不必重新经历关于
“热质”等古老概念的全部曲折，只要承认物质由微粒构成，便能较快切入热现象的本质。这里采用分子
运动论，不是省略宏观实验事实，而是用一个统一的微观图景解释这些事实。

十九世纪八十年代，原子、分子虽然已经成为重要的科学概念，但它们是否真实存在，仍曾引起激烈争论。
爱因斯坦（Albert Einstein）对布朗运动的理论分析给出了可检验的粒子图像；让·佩兰（Jean Perrin） 随后的实验测量使这一图像得到
有力支持：悬浮在液体中的微粒不会全部沉到底部，而会在液体分子不断碰撞下持续作无规则运动。现代的
原子力显微镜和扫描隧道显微镜甚至可以在适当条件下直接获得分子尺度的图像。

\molecularMicroscopyFigure

今天我们知道，分子是普通宏观物质的重要组成单元，但不是自然界最基本的粒子。在本课程讨论的温度
和能量范围内，通常把分子的内部结构视为稳定；若发生化学反应，可以理解为原子组合方式改变，而不是
把热运动直接当成任意的粒子裂变过程。

佩兰的工作之所以重要，不只是“看见一些小颗粒在动”。如果液体是完全连续、内部没有微粒的介质，
悬浮颗粒长期不规则运动及其随颗粒大小、温度变化的统计规律就很难得到统一解释；在分子图景下，它们
恰好是大量液体分子从各个方向撞击颗粒而产生的可见结果。图\ref{ch14:fig:s28-19-1}帮助理解分子尺度成像的外观，并不是一份实验记录。
真实原子力显微镜探测针尖与样品之间的力或力梯度，扫描隧道显微镜探测隧穿电流；它们的衬度
还受仪器模式和样品电子性质影响，不能一概看成几何高度。两类技术都能把研究尺度推进到分子量级。

分子有几个共同特点。第一，尺度很小，典型线度约为 $\SI{0.1}{nm}$ 到 $\SI{1}{nm}$，单个分子的质量也
很小，所以宏观物体总含有极多分子。第二，同种分子具有全同性；除非能连续追踪系统中每个粒子的轨迹，
原则上无法给同种分子贴上永久标签。第三，构成物质的分子并不静止：气体分子可在较大空间内自由运动
并频繁碰撞；液体分子彼此靠近，却仍能不断交换近邻；固体原子虽有近似固定的平衡位置，仍会在其附近
振动并通过相互作用交换能量。

“全同性”在大量粒子问题中尤其值得注意。在通常的少体力学中，我们可以把第一个、第二个物体分别
编号并追踪轨迹；同种分子却没有可测量的永久身份差别。除非能够从始至终追踪系统中每一个分子，否则
交换了位置的两个同种分子在物理上不可分辨。正是粒子数巨大而又彼此全同，使宏观热系统的描述方式与
只研究几个物体的动力学有很大差别，也使统计方法成为必需。

分子运动论的基本假设是：每一个分子的运动仍遵循基本物理定律。只是在面对约 $10^{23}$ 个粒子时，
逐一求轨迹既不现实，也没有必要；真正需要的是从微观规律推导密度、弹性、声速、热导率、电导率、
热膨胀和折射率等宏观量。统计规律并不否定单个粒子的动力学，而是在巨量自由度下提取稳定的整体性质。

因此，“每个分子遵守力学定律”不是一句无法检验的口号。以微观相互作用和粒子统计为起点，人们确实
能够计算物质的密度、弹性、声速、热导率、电导率、热膨胀系数以及透明介质的折射率，并同实验比较。
热物理的任务，就是说明为什么单个分子的剧烈涨落不会使宏观物体的温度和压强同样剧烈地跳动，以及
哪些条件下这些平均量能够稳定、可重复地测量。

\section{温度、热平衡与能量均分}\label{ch14:sec:temperature}

测量温度时，需要让温度计与被测物充分接触，并等待温度计示数稳定。两个物体接触后，若有净能量从
一方流向另一方，就说它们温度不同；若在宏观上没有净能量交换，就说二者处于热平衡。热平衡关系具有
传递性，这使我们能够用温度计比较不同物体的温度。
这里的“没有净能量交换”专指通过允许的热接触不再发生净传热；两个温度不同、却被绝热壁隔开的
物体，也可以暂时没有热量交换，不能因此判定热平衡。
热平衡的传递性通常称为热力学第零定律：若 $A$、$B$ 分别与温度计 $C$ 热平衡，
则 $A$、$B$ 也彼此热平衡。温标还需要选定可测的测温量和标定规则，不能把“热平衡”本身当成数值刻度。

例如，暂把温度计和被测物的热容量分别视为常数 $C_t,C_b$，初温分别为 $T_t,T_b$，
接触过程中整体绝热且不做功。把二者一同选为系统，能量守恒为
\[
 C_t(T_f-T_t)+C_b(T_f-T_b)=0,
 \qquad T_f=\frac{C_tT_t+C_bT_b}{C_t+C_b}.
\]
因此测量对被测物的扰动是
\[
 T_f-T_b=\frac{C_t}{C_t+C_b}(T_t-T_b).
\]
只有 $C_t\ll C_b$ 时，温度计达到平衡的读数才接近被测物接触前的温度。
“等到示数稳定”排除了尚未达到平衡的误差，却不能消除温度计本身吸、放热造成的扰动。

图\ref{ch14:fig:s28-19-2}比较未接触与正确浸入的两种放置方式；只有右图的温包与水充分接触，示数才可用来判断水温。
\figureNineteenTwoFaithful

图\ref{ch14:fig:thermometry}中，初温较低的温度计吸收能量，而被测物略微冷却。
只在温度计热容量相对很小时，共同终温 $T_f$ 才近似等于被测物原来的温度。
\thermometerMeasurementDiagram

以气体为例，分子之间不断发生碰撞。能量较高与能量较低的分子相撞时，单次碰撞的能量转移方向并不
确定，但在大量碰撞的统计意义下，极端高能或极端低能状态都不占多数。若追踪单个分子，它的能量会
随时间不断涨落；若观察足够多分子，平均量却可以稳定下来。两个初始平均能量不同的气体接触后，碰撞
会使能量重新分配，最后达到共同的热平衡状态。

更具体地说，两分子碰撞满足能量、动量守恒，但这并不保证每次都是“高能分子把能量交给低能分子”。
碰撞结果还与质量、入射方向和相位有关。能够成立的是概率陈述：在一个能量很高的分子与普通分子频繁
碰撞的长时间过程中，前者更可能损失一部分能量；能量很低的分子则更可能从周围获得能量。若真的追踪
一个分子并记录它的能量--时间曲线，曲线会不断上下波动，偶尔偏离平均值，却不大可能永久停留在极大
或极小值。温度对应的是这种统计分布的尺度，而不是某个指定分子在某一瞬间的动能。

图\ref{ch14:fig:equilibrium}中的箭头比较同种分子热运动的典型大小；箭头方向杂乱，
表示平均流速为零，而不是所有分子停止运动。
\molecularEquilibriumDiagram

大量分子经过频繁碰撞后能够充分交换原先所具有的能量，最后达到稳定状态。此时气体均匀充满可用空间，
不存在整体宏观流动，密度、压强等状态量也不再随时间系统地改变，这就是热平衡。若把两份各自平衡的
气体混合，而它们原先每个分子的平均平动能相同，混合后的碰撞不会引起持续的净能量转移；若平均动能
不同，则原先平均动能较高的一方会在统计上失去能量，另一方获得能量，直至两者达到共同温度。这正与
“热量从高温物体传向低温物体”的宏观经验相对应。

温度与分子无规则运动的平均能量有直接关系。经典能量均分定理指出：在经典统计成立且系统处于温度
$T$ 的热平衡状态时，每一个能写成二次型的自由度，平均具有
\begin{equation}
  \boxed{\langle E\rangle=\frac12k_BT}.
\end{equation}
绝对温度与摄氏温度的关系是
\begin{equation}
  \boxed{T=t_{\rm C}+273.15},
\end{equation}
其中 $T$ 用 K 表示，$t_{\rm C}$ 用 $\si{\degreeCelsius}$ 表示。温差的数值满足
$\Delta T=\Delta t_{\rm C}$，但物态方程中必须使用绝对温度。

平动动能、转动动能以及近似简谐振子的动能和势能，都是速度、角速度或形变量的二次型；在相应自由度
被充分激发时，它们可按能量均分定理分配能量。单原子理想气体分子有三个平动自由度，因此
\begin{equation}
  \left\langle E_{\rm tr}\right\rangle=\frac32k_BT,
  \qquad
  v_{\rm rms}=\sqrt{\langle v^2\rangle}=\sqrt{\frac{3k_BT}{m}}.
\end{equation}
例如在 $T=\SI{300}{K}$ 时，氧分子的均方根速率约为数百米每秒，远大于通常风速，并与声速处于同一
数量级。宏观空气没有整体高速流动，是因为各分子速度方向无规则，矢量平均速度接近零，而
$\langle v^2\rangle$ 并不为零。

在 $T=\SI{300}{K}$ 时，每个二次型自由度的平均能量为
\begin{equation}
  \frac12k_BT\simeq\SI{2.07e-21}{J},
\end{equation}
三个平动自由度合计
\begin{equation}
  \left\langle E_{\rm tr}\right\rangle
  =\frac32k_BT\simeq\SI{6.21e-21}{J}.
\end{equation}
取氧分子摩尔质量 $\mu=\SI{32e-3}{kg.mol^{-1}}$，单分子质量 $m=\mu/N_A$，可得
\begin{equation}
  v_{\rm rms}=\sqrt{\frac{3k_BT}{m}}\simeq\SI{484}{m.s^{-1}}.
\end{equation}
这一定量结果补足了“数百米每秒”的数量级判断。它与声速同量级并非偶然，后续研究气体中的声传播时，
声速正是由分子的热运动尺度和气体的可压缩性共同决定。

能量均分定理有适用范围。它要求相应自由度能够用经典物理描述并在热平衡中被充分激发。金属中的电子、
低温下冻结的分子振动自由度等对象需要量子统计，不能把 $\tfrac12k_BT$ 不加判断地套给所有自由度。

金属中的电子提供了一个反例：若在室温下机械地使用经典能量均分，所得电子典型速率约为
$10^5\,\mathrm{m/s}$；实际金属中电子的特征速率可达到约 $10^6\,\mathrm{m/s}$。这一数量级差别不是
实验粗糙造成的，而是电子服从量子统计的直接表现。能量均分定理非常有用，但“先判断是否经典、相应
自由度是否已被激发”是使用它不可省略的一步。

\section{热膨胀}\label{ch14:sec:expansion}

几乎所有物质在温度升高时都会发生尺寸变化。许多温度计正是利用液体的热胀冷缩来显示温度。固体中的
原子不能自由远离近邻；原子间相互作用在平衡距离 $r_0$ 附近形成势阱。绝对零度的经典图像中，原子
处在势能最低点；温度升高后，原子获得热运动能量，在平衡位置附近振动。

若势阱关于 $r_0$ 完全对称，振动幅度虽然增加，平均位置仍等于 $r_0$，不会产生一阶热膨胀。真实原子间
势能在短程排斥侧很陡，在长程吸引侧较缓，因此右侧转折点比左侧移动得更多，平均原子间距随温度升高
而增加。这种微观的不对称性正是通常固体热膨胀的来源。

图\ref{ch14:fig:anharmonic}的两个转折点 $r_1,r_2$ 是同一能量水平与势能曲线的交点。
图\ref{ch14:fig:s28-19-3}进一步比较两个能量水平。右侧变宽是理解热膨胀的线索，
但平均间距 $\langle r\rangle$ 需要按粒子在各位置停留的时间或平衡概率加权，不能只取两端的中点。
\anharmonicPotentialDiagram

\figureNineteenThreeFaithful

设固体在参考温度 $T_0$ 时长度为 $L_0$，温度改变 $\Delta T=T-T_0$ 后长度改变 $\Delta L$。线膨胀系数
定义为
\begin{equation}
  \boxed{\alpha=\frac{1}{L_0}\frac{\Delta L}{\Delta T}},
  \qquad
  \boxed{\Delta L=\alpha L_0\Delta T},
  \qquad
  \boxed{L=L_0(1+\alpha\Delta T)}.
\end{equation}
这里必须保留初始长度 $L_0$，以保证伸长量具有正确量纲。对于各向同性固体，在改变量很小时，面积和体积
膨胀系数分别近似为 $2\alpha$ 和 $3\alpha$。

线膨胀系数严格说来可以随温度缓慢变化。上面的线性关系是在温区不太宽、总相对伸长很小的条件下，把
$\alpha(T)$ 近似看成常量后得到的。若温区很宽，应写成
\begin{equation}
  \frac{\mathrm dL}{L}=\alpha(T)\,\mathrm dT,
  \qquad
  \ln\frac{L}{L_0}=\int_{T_0}^{T}\alpha(T')\,\mathrm dT'.
\end{equation}
表~\ref{ch14:tab:s28-19-1}所列数值是室温附近的典型实验值，用于估算时还应注意材料成分、热处理状态以及各向异性。

\begin{center}\begin{minipage}{\linewidth}\centering
\small
\begin{tabular}{lclc}
\toprule
材料 & $\alpha/\si{K^{-1}}$ & 材料 & $\alpha/\si{K^{-1}}$\\
\midrule
铝 & $24\times10^{-6}$ & 黄铜、青铜 & $19\times10^{-6}$\\
耐热玻璃 & $3.2\times10^{-6}$ & 普通玻璃 & $9\times10^{-6}$\\
铅 & $29\times10^{-6}$ & 钢 & $11\times10^{-6}$\\
铜 & $17\times10^{-6}$ & 混凝土 & $12\times10^{-6}$\\
\bottomrule
\end{tabular}
\booktablecaption{室温附近若干常见材料的线膨胀系数。}\label{ch14:tab:s28-19-1}
\end{minipage}\end{center}

热膨胀产生的约束力可以很大，桥梁、轨道、管路和精密仪器都必须预留温度变化造成的尺寸改变。纪尧姆（Guillaume）
因研究镍铁合金获得1920年诺贝尔物理学奖。含镍约 $35.4\%$ 的因瓦合金线膨胀系数可低至约
$0.9\times10^{-6}\,\si{K^{-1}}$，因此适合制作精密仪器的结构件。

工程中真正危险的往往不是自由膨胀，而是“想膨胀却被约束”。若一根杆完全不能伸长，温度变化产生的
热应变 $\alpha\Delta T$ 会转化为应力，在线弹性范围内约为 $\sigma=E\alpha\Delta T$。由于杨氏
模量 $E$ 很大，即使 $\alpha\Delta T$ 很小也可能产生显著内力。桥梁伸缩缝、铁轨间隙和管道补偿器
都是对这一事实的工程回应；因瓦合金的价值则在于从材料本身减小热应变。

\subsection{自由膨胀量怎样转化为约束力}
先区分两种应变。温度变化引起的自由热应变是 $\alpha\Delta T$；外力引起的轴向弹性应变是 $\sigma/Y$，其中 $Y$ 为杨氏模量，规定拉应力为正。在小形变、线弹性和温度均匀的近似下，总应变为
\[
 \varepsilon=\alpha\Delta T+\frac{\sigma}{Y}.
\]
固定两端不是令每一项都为零，而是要求总伸长为零。于是均匀杆升温后有 $\sigma=-Y\alpha\Delta T$，负号表示受压。前面所说的热应力量值 $Y\alpha\Delta T$，应与这里的有符号应力区分。
若杆可以自由移动，则约束力为零，伸长才等于自由膨胀量；若已经屈服、脱离支座或发生屈曲，就要重新建模。

\Needspace{180pt}
\begin{origexamplebox}{两种串联材料受热后的共同轴力}\label{ch14:ex:teaching-composite}
两根细杆首尾牢固连接，长度为 $L_1,L_2$，截面积同为 $A$，杨氏模量为 $Y_1,Y_2$，线膨胀系数为 $\alpha_1,\alpha_2$。初温时恰好无应力地安装在两个固定支座间。现均匀升温 $\Delta T>0$，求轴力大小和两杆各自的长度变化。假定保持直线，不发生屈曲；两个外端支座固定，两杆在连接处无相对滑移，但连接点可以整体移动。
\end{origexamplebox}
\Needspace{65pt}
\begin{xesolution}
取拉力 $F$ 为正。连接处没有额外轴向外力，静力平衡要求两杆承受同一轴力，而不是承受相同伸长。分别有
\[
 \Delta L_i=\alpha_iL_i\Delta T+\frac{FL_i}{AY_i},\qquad i=1,2.
\]
固定的是两杆总长，因此 $\Delta L_1+\Delta L_2=0$，得到
\[
 F=-A\Delta T\frac{\alpha_1L_1+\alpha_2L_2}{L_1/Y_1+L_2/Y_2}.
\]
若两膨胀系数均为正，$F<0$，两杆受压，压力大小为 $|F|$。将同一个 $F$ 代回第一式，并将分母同乘 $Y_1Y_2$，得
\[
 \Delta L_1=\frac{L_1L_2(\alpha_1Y_1-\alpha_2Y_2)}{L_1Y_2+L_2Y_1}\Delta T,
 \qquad \Delta L_2=-\Delta L_1.
\]
若第一根杆位于左边，则连接点向右的位移正是 $\Delta L_1$。连接点可以移动，所以通常并不是 $\Delta L_1=\Delta L_2=0$；一段略伸长时另一段必须等量缩短。仅当 $\alpha_1Y_1=\alpha_2Y_2$ 时，两段才都不改变长度。
分母的单位为长度除以压强，所得 $F$ 的单位为力；若两种材料相同，就回到 $F=-AY\alpha\Delta T$，与总长度无关。若仅一段非常柔软，更多的弹性压缩会集中在该段，这也解释了不能简单平均两个杨氏模量的原因。
\end{xesolution}

\section{理想气体物态方程}\label{ch14:sec:ideal-gas}

物态方程用宏观状态参量之间的函数关系描述物质。较高温度、较低密度的气体具有最简单的近似：忽略
分子本身体积和分子间相互作用，把它处理为理想气体。若容器体积为 $V$，气体压强为 $P$、绝对温度为
$T$、物质的量为 $\nu$，则
\begin{equation}
  \boxed{PV=\nu RT},
  \qquad R\approx\SI{8.314462618}{J.mol^{-1}.K^{-1}}.
\end{equation}

这一简洁方程来自很长时期的气体实验。若容器中气体总质量为 $M$、平均摩尔质量为 $\mu$，则
$\nu=M/\mu$，物态方程也可写成
\begin{equation}
  PV=\frac{M}{\mu}RT.
\end{equation}
实验上分别改变压强、体积和温度，经过一个多世纪的积累才逐渐确认：在足够稀薄、远离液化条件时，
不同气体都趋近同一种形式。所谓“理想气体”因此是一种受实验支持的极限模型，而不是声称真实分子
没有体积、没有相互作用。
对同一份封闭气体，在没有泄漏和化学反应时 $\nu$ 不变，任意两个平衡状态满足
\begin{equation}
  \boxed{\frac{P_1V_1}{T_1}=\frac{P_2V_2}{T_2}}.
\end{equation}
使用这个比例式时，压强和体积可采用任意彼此一致的单位，但温度仍必须使用 K。

这里有两种常见使用情形。第一，对给定容器内的一份封闭气体，只要过程没有泄漏、化学反应或物质交换，
物质的量 $\nu$ 就不变，因而 $PV/T$ 保持常量。比例式中同一种物理量的单位会在分子、分母间相消，所以
可以直接用 $\mathrm{cmHg}$、升等方便单位。第二，对通常大气这类开放环境，无法指定“整团空气”的
固定体积和总物质的量，就应改用单位体积内的质量，即密度形式，而不能生硬套用一份封闭气体的比例式。

若研究开放环境中的气体，没有一个固定的总体积。把 $\nu=M/\mu$ 与密度 $\rho=M/V$ 代入物态方程，可得
\begin{equation}
  \boxed{P=\frac{\rho RT}{\mu}},
  \qquad
  \boxed{\rho=\frac{P\mu}{RT}},
\end{equation}
其中 $\mu$ 是摩尔质量。气体越远离液化条件，理想气体近似通常越好；温度较低或密度较大时，分子有限
体积和相互作用会带来明显偏差。

“离液化状态越远，越像理想气体”是一个实用判断。升高温度使热运动更容易压过分子间吸引，降低密度
则减小分子彼此靠近的机会；相反，低温高密度下，分子有限体积和相互作用都不可忽略，需要更精细的
状态方程。竞赛题若直接称“理想气体”，意味着题目已经授权使用上述近似；若讨论真实气体实验精度，
则必须重新检查偏差是否小到可以忽略。

\section{混合气体与氩气的发现}\label{ch14:sec:mixtures}

空气主要由氮气和氧气组成，并含少量其他气体。在理想气体近似下，第 $i$ 种成分满足
\begin{equation}
  P_iV=\nu_iRT=\frac{m_i}{\mu_i}RT,
  \qquad
  P_i=\frac{\rho_iRT}{\mu_i}.
\end{equation}
道尔顿分压定律给出
\begin{equation}
  \boxed{P=\sum_iP_i=RT\sum_i\frac{\rho_i}{\mu_i}}.
\end{equation}
因此，若能准确测出气体压强和密度，就有机会识别混合物中少量而未知的成分。

传统水银气压计通过水银柱高度确定压强，但残余气体、弯月面端点和管壁不规则折射都会形成系统误差。

传统装置的基本原理很直接：汞槽液面受大气压作用，长管上端若接近真空，则汞柱静止时有
$p_{\rm atm}\simeq\rho_{\rm Hg}gH$。它很适合日常测量，却未必满足精密研究。第一，很难仅凭外观确认
汞柱上方是否达到足够可靠的真空；第二，汞面弯月线使“上端”和“下端”究竟取哪一点缺少完全一致的
判据；第三，玻璃管壁的不规则折射会让看到的位置偏离真实位置。这些不是多测几次就自动消失的随机
误差，而是需要改变装置结构才能压低的系统误差。

图\ref{ch14:fig:s28-19-4}中，管口始终浸在汞槽内，管内与槽中的水银连续连通。
\figureNineteenFourFaithful

图\ref{ch14:fig:barometer-principle}标明了高度 $H$ 的两端及残余气体压强 $p_g$。
\traditionalBarometerDiagram

瑞利改进了气体压强测量装置。他特别设计出形状一致的叉形液面，使两端弯月面的表面张力效应能够
相互抵消；可靠的阀门和参考端用于检查真空状态。这样，真正需要比较的是两个明确标志位置之间的高度差，
而不是凭肉眼猜测弯月面的“顶端”。

图\ref{ch14:fig:s28-19-5}中 $C$ 处的结构把汞柱限制成与 $B$ 处相同的叉形，使两个端面的毛细和表面张力效应尽量同形；
测量 $C$、$B$ 处活塞或标志之间的高度差时，共同误差便可相消。$D$ 处阀门用于检验参考端真空是否
可靠，其余储槽、连接管和调节部件则用来缓慢改变液面位置。精密仪器的核心不是“画得更复杂”，而是
把每个难以判断的端点改造成可重复比较的量。

\figureNineteenFiveFaithful

图\ref{ch14:fig:rayleigh-forks}放大比较 $B,C$ 两处的同形接触结构。
只有两处液体、润湿状态及几何相同，毛细修正才可近似作为共同项抵消。
\rayleighApparatusDiagram

通过化合物分解得到的氮气密度约为 $\SI{1.2511}{g.L^{-1}}$，而从空气中除去已知活泼成分后所得“氮气”
密度约为 $\SI{1.2572}{g.L^{-1}}$。差值远大于实验误差，说明后者并非纯氮。设未知重气体在混合气体中
的分压占比为 $\alpha$，同温同压下氮气和未知气体的密度分别为 $\rho_1,\rho_2$，则混合气体密度
可以从各成分质量逐步推出。若给定体积 $V$ 中两种气体质量分别为 $m_1,m_2$，摩尔质量分别为
$\mu_1,\mu_2$，它们的分压满足
\begin{equation}
  P_1V=\frac{m_1}{\mu_1}RT,
  \qquad
  P_2V=\frac{m_2}{\mu_2}RT.
\end{equation}
由 $P=P_1+P_2$ 和 $\rho'=(m_1+m_2)/V$ 得
\begin{equation}
  \rho'=\frac{P_1\mu_1+P_2\mu_2}{RT}.
\end{equation}
令未知重气体的分压 $P_2=\alpha P$，而同温同压下纯氮和未知气体密度分别为
$\rho_1=P\mu_1/(RT)$、$\rho_2=P\mu_2/(RT)$，便得到
\begin{equation}
  \rho'=(1-\alpha)\rho_1+\alpha\rho_2,
\end{equation}
因而
\begin{equation}
  \Delta\rho=\rho'-\rho_1=\alpha(\rho_2-\rho_1),
  \qquad
  \rho_2=\rho_1+\frac{\Delta\rho}{\alpha}.
\end{equation}
瑞利勋爵与拉姆齐（William Ramsay）通过进一步去除空气中的已知成分并测定体积、分压，得到 $\alpha\approx0.0104$，
这里的 $\alpha$ 是除去氧等成分后残余混合气体中的比例，不是全部空气中的比例。
这些实验确认了一种新气体——氩。氩是18号元素，摩尔质量约为
$\SI{39.948}{g.mol^{-1}}$，主要以单原子形式存在，在干燥空气中的体积分数约为 $0.934\%$。瑞利
因气体密度研究及氩的发现获得1904年诺贝尔物理学奖。

这个故事还要追溯到1785年。卡文迪许曾不断去除空气中的已知活泼成分，得到极少量残余气体，却没有
把它认定为新元素。一个多世纪后，瑞利重新采用类似分离思路：如果残余气体近似为纯未知成分，
便可通过分离前后的体积变化确定它在原混合气体中的分压比例 $\alpha$。问题的代数并不复杂，困难在于
$\Delta\rho$ 和 $\alpha$ 都很小，任何弯月面读数、残余气体或温度控制误差都会被放大。正因有精密、
可复现的压强和密度测量，密度差才不能再被当作偶然误差。

在 $0\,{}^\circ\mathrm C$、$1\,\mathrm{atm}$ 条件下，氩气密度约为 $\SI{1.784}{g.L^{-1}}$。须注意，上列历史近似数据并不是一套足以重现现代精确密度的自洽数据：
直接用 $\Delta\rho=0.0061\,\mathrm{g/L}$ 和 $\alpha=0.0104$ 代入，会得到
$\rho_2\simeq1.84\,\mathrm{g/L}$，不能把它写成精确算得 $1.784\,\mathrm{g/L}$。
精密比较必须采用同一温压条件及同一批实验得到的比例。
氩化学性质稳定，今天广泛用于惰性保护气氛、科研
实验、医疗和照明等领域。这里的重点不是记住几个年份，而是看清一条实验推理链：先发现超过
误差范围的反常密度，再改进仪器、排除系统误差，最后用混合气体物态方程把微小差异转化为未知成分的
密度和摩尔质量。

\subsection{分压按分子数分配，不按质量分配}
各组分处在同一温度、同一可用体积时，由 $p_iV=n_iRT$ 和 $pV=(\sum_i n_i)RT$ 相除，得到
\[
 x_i=\frac{n_i}{n},\qquad \frac{p_i}{p}=x_i,\qquad
 \overline\mu=\frac{\sum_i n_i\mu_i}{\sum_i n_i}=\sum_i x_i\mu_i.
\]
这里 $x_i$ 是摩尔分数，$\overline\mu$ 是混合气体的平均摩尔质量。若已知的是质量分数 $w_i=m_i/\sum_jm_j$，先用 $n_i=m_i/\mu_i$ 换算，才有
\[
 \frac1{\overline\mu}=\sum_i\frac{w_i}{\mu_i}.
\]
因此，质量相等的轻、重气体并不贡献相等分压。用密度反推成分之前，必须先说明给出的是哪一种比例。

\Needspace{180pt}
\begin{origexamplebox}{等质量氦、氮混合物的分压与密度}\label{ch14:ex:teaching-mixture}
质量各为 $m$ 的氦气和氮气装在同一容器中，达到温度 $T$、总压强 $p$ 的平衡态，均按理想气体处理。取摩尔质量分别为 $4$ 和 $28\,\mathrm{g/mol}$。求两者分压、混合气体平均摩尔质量及密度。
\end{origexamplebox}
\Needspace{65pt}
\begin{xesolution}
相同质量下，物质的量之比为
\[
 n_{\rm He}:n_{\rm N_2}=(m/4):(m/28)=7:1.
\]
所以氦、氮的摩尔分数分别为 $7/8,1/8$，分压分别为 $7p/8,p/8$，并非各为 $p/2$。平均摩尔质量既可按摩尔分数加权，也可用总质量除以总物质的量：
\[
 \overline\mu=\frac78\times4+\frac18\times28
 =7\,\mathrm{g/mol},\qquad
 \rho=\frac{p\overline\mu}{RT}.
\]
最后一式使用SI单位时须代入 $\overline\mu=7\times10^{-3}\,\mathrm{kg/mol}$。平均值位于两种摩尔质量之间且更靠近氦，是因为氦的分子数更多，而不是因为两种质量不相等。总压由两分压相加恢复为 $p$，总密度也等于各组分密度之和。
\end{xesolution}

\section{理想气体压强的微观推导}\label{ch14:sec:pressure}

设体积为 $V$ 的容器内有 $N$ 个同种分子，每个质量为 $m$。取一块垂直于 $x$ 轴、面积为 $S$ 的壁面。
某分子以 $x$ 方向速度分量 $v_x>0$ 撞向壁面并完全弹性反射，分子的 $x$ 动量由 $mv_x$ 变为
$-mv_x$，故壁面获得冲量 $2mv_x$。

\figureNineteenSixFaithful

图\ref{ch14:fig:s28-19-6}右侧的 $Q$ 点为碰撞点，入射、反射线具有相等的法向夹角。
图\ref{ch14:fig:collision-layer}另画出参与计数的薄层，避免把三维反射轨迹与计数体积混在一起。
\gasWallCollisionDiagram

在时间 $\Delta t$ 内，能到达该壁面的同类分子位于底面积 $S$、长度 $v_x\Delta t$ 的薄柱中；还要乘
$1/2$，因为热平衡时约有一半分子朝向该壁面。把所有分子的贡献相加，壁面得到的总冲量为
\begin{equation}
  \Delta p_x=\frac12\sum_{i=1}^N(2mv_{ix})
  \frac{Sv_{ix}\Delta t}{V}
  =\frac{mS\Delta t}{V}\sum_{i=1}^Nv_{ix}^2.
\end{equation}

这里的几何因子可以直接从图\ref{ch14:fig:collision-layer}读出：速度的法向分量为 $v_x$，所以在 $\Delta t$ 内，距离壁面不超过
$v_x\Delta t$ 的分子才可能撞到所取面积 $S$；相应候选区域体积为 $Sv_x\Delta t$，占整个容积的比例
为 $Sv_x\Delta t/V$。统计时若把 $v_x^2$ 相同的正、负方向分子都算进求和，就必须再乘 $1/2$，因为
背离该壁面运动的那一半分子不会在这段时间撞上它。这个 $1/2$ 不是碰撞损失能量，而是方向计数修正。
于是平均力 $F=\Delta p_x/\Delta t$，压强为
\begin{equation}
  \boxed{P=\frac FS=\frac{mN}{V}\overline{v_x^2}
  =n_Nm\overline{v_x^2}},
\end{equation}
其中 $n_N=N/V$ 是分子数密度。热平衡气体各向同性，故
\begin{equation}
  \overline{v_x^2}=\overline{v_y^2}=\overline{v_z^2}
  =\frac13\overline{v^2}.
\end{equation}

分子动量的改变来自壁面对分子的力；按牛顿第三定律，分子同时把大小相等、方向相反的冲量交给
壁面。单位时间总冲量就是壁面所受平均力，再除以面积才得到压强。热平衡时没有优先方向，因此三个
直角坐标方向的速度平方平均值相同。注意这里相等的是 $\overline{v_x^2},\overline{v_y^2},
\overline{v_z^2}$，而不是要求每个分子在每一瞬间都有 $v_x=v_y=v_z$。
所以
\begin{equation}
  \boxed{P=\frac13n_Nm\overline{v^2}
  =\frac23n_N\left\langle\frac12mv^2\right\rangle
  =\frac23u_{\rm tr}},
\end{equation}
其中 $u_{\rm tr}$ 是单位体积内的平均平动动能。再用
$\langle E_{\rm tr}\rangle=3k_BT/2$，得到
\begin{equation}
  \boxed{P=n_Nk_BT},
  \qquad
  \boxed{PV=Nk_BT=\frac{N}{N_A}N_Ak_BT=\nu RT}.
\end{equation}
这里 $\nu=N/N_A$ 是物质的量（单位 $\mathrm{mol}$），$n_N=N/V$ 是数密度（单位 $\mathrm{m^{-3}}$），二者不可混用。
因此普适气体常数满足 $R=N_Ak_B$。宏观物态方程并不是孤立的经验公式，而是大量分子碰壁冲量的统计
结果。

\Needspace{8\baselineskip}
\Needspace{16\baselineskip}\section{综合例题}
\begin{origexamplebox}{【人体膨胀】}\label{ch14:ex:s28-19-1}
已知人体的线膨胀系数为 $\alpha$。那么冬天时小明在室内（$25\,{}^\circ\mathrm C$）比室外
（$0\,{}^\circ\mathrm C$）体积大多少？小明的密度膨胀系数为多少？
\end{origexamplebox}

\Needspace{6\baselineskip}\begin{xesolution}
先明确模型：环境温度从 $0\,{}^\circ\mathrm C$ 变到 $25\,{}^\circ\mathrm C$，
并不意味着真实人体各部分也改变同样的温度；人体有体温调节，不能据此算出实际体积变化。
若按题意作理想化计算，须另外假定物体整体温度确实改变 $25\,\mathrm K$、质量不变、
膨胀系数恒定且各向同性。以下求的是这一模型的结果。设 $0\,{}^\circ\mathrm C$ 时三个方向的特征长度分别为
$a_0,b_0,c_0$，升高温度 $\Delta T$ 后，三个方向都按同一个线膨胀系数变化：
\begin{equation}
  a=a_0(1+\alpha\Delta T),\quad
  b=b_0(1+\alpha\Delta T),\quad
  c=c_0(1+\alpha\Delta T).
\end{equation}
于是
\begin{equation}
  \frac{V}{V_0}=(1+\alpha\Delta T)^3
  \simeq 1+3\alpha\Delta T,
\end{equation}
其中略去了 $\alpha^2$ 及更高阶小量。取 $\Delta T=25\,\mathrm K$，得到
\begin{equation}
  \boxed{\frac{V_{25}-V_0}{V_0}\simeq(75\,\mathrm K)\alpha},
  \qquad
  \boxed{V_{25}-V_0\simeq(75\,\mathrm K)\alpha V_0}.
\end{equation}

若人体质量近似不随温度改变，$\rho=m/V$。因此
\begin{equation}
  \frac{\rho}{\rho_0}=\frac{1}{1+3\alpha\Delta T}
  \simeq1-3\alpha\Delta T.
\end{equation}
若把体积膨胀系数定义为 $\beta_V=(1/V)(\mathrm dV/\mathrm dT)$，则
$\boxed{\beta_V\simeq3\alpha}$；若题目所说的“密度膨胀系数”是
$\beta_\rho=(1/\rho)(\mathrm d\rho/\mathrm dT)$，则
$\boxed{\beta_\rho\simeq-3\alpha}$。负号表示升温时密度减小。
\end{xesolution}

\Needspace{10\baselineskip}
\Needspace{16\baselineskip}\begin{origexamplebox}{【“真实”感受】}\label{ch14:ex:s28-19-2}
小明为了研究“物体的冷热程度”做了以下实验。试通过以下实验说明，小明感受到的冷热程度究竟是什么。

（1）一只手放在冷水里，另一只手放在热水里，然后同时拿出后立刻放在温水里。

（2）小明发现南方的冬天很冷，即便大多数日子天气预报报道的气温是比北方高的。

（3）小明给妈妈打热水泡脚，妈妈把脚伸进水盆后打了个喷嚏。

（4）小明发现有的叔叔需要进入到 $60\,{}^\circ\mathrm C$ 以上的环境工作，进去之前要先吃点辣椒，
再穿上棉被。
\end{origexamplebox}

\Needspace{6\baselineskip}\begin{xesolution}
皮肤并不是直接读取环境的温度数值。冷、热感受首先取决于皮肤与外界之间的热流大小和方向，同时还
受到感受器适应、血管舒缩、出汗、蒸发以及神经反射的影响。因此，相同温度的物体可以带来不同感觉，
不同温度的环境也可能带来相似感觉。

图\ref{ch14:fig:heat-feeling}把四类观察分开列出；前三者不能推出第四种做法安全。
\heatFeelingDiagram

（1）放过冷水的手温度较低，进入温水后热量由水流入手，感觉“热”；放过热水的手温度较高，进入
同一盆温水后热量由手流向水，感觉“冷”。温水的客观温度只有一个，感觉不同来自皮肤初态以及热流
方向不同。这也说明手感不能代替温度计。

（2）体感温度不只由气温决定。空气湿度、风速、日照、衣物含水程度、室内取暖条件以及人体活动量
都会改变对流、蒸发和辐射换热。南方冬季在高湿、衣物或墙体容易潮湿且室内缺少持续供暖的情形下，
皮肤失热状况可能比气温数值显示的更明显；“南方一定比北方冷”并不是普遍物理定律，必须结合具体
湿度、风速和生活环境判断。

（3）泡脚与喷嚏先后发生，只说明时间上相邻，不能由一次观察确定因果。
是否存在温度刺激之外的因素，题目没有提供足够证据。这个事例不能用一个热学公式解释喷嚏机制，
也不能把喷嚏当成温度计读数；可以确定的是，单次生理反应不足以判断水温或人体净吸热量。

（4）当环境温度高于体表温度时，热量可能从环境流向人体。棉被在这里不是“生热”，而是用低导热、
低对流的空气层短时减小某些外来热流；但厚重衣物也会阻碍散热和汗液蒸发。
因此不能单凭隔热一项就判断净热量减少，更不能把吃辣椒当成可靠的防热措施。
进入高温作业环境需要专门的风险评估、工程控制、合适防护和休息制度；题述做法不能作为安全操作方法，
不应模仿。这个例子在物理上要说明的是：隔热同时影响热量进入与排出，必须比较全部热流。
\end{xesolution}

\Needspace{10\baselineskip}
\Needspace{16\baselineskip}\begin{origexamplebox}{【学习化学的小明】}\label{ch14:ex:s28-19-3}
小明开始学习化学了。

（1）老师教小明用量筒时说，视线要与凹液面相平，读数才能准确。可是有一次小明在称量液体时，
竟然发现没有凹液面，反而是凸液面。请说出可能的原因，这样是否会造成测量不准确？

（2）使用滴定管时，小明发现液滴半天也不滴下去，请说出可能的原因。

（3）老师说在烧杯中倒水时，需要“引流”。这是什么意思？
\end{origexamplebox}

\Needspace{6\baselineskip}\begin{xesolution}
这三个现象都与液体的表面张力以及液体对固体表面的润湿性有关。

图\ref{ch14:fig:surface-phenomena}左图的最高点与视线同高；中图液滴颈部仍接触管口；
右图液流沿玻璃棒的一侧进入接收容器。这三种接触关系正对应三个子问。
\surfacePhenomenaDiagram

（1）水能润湿洁净玻璃，玻璃壁附近液面被拉高，形成凹液面；水银等不能润湿玻璃的液体，接触角
大于 $90^\circ$，玻璃壁附近液面反而降低，形成凸液面。量筒刻度仍然可以使用，但读数基准要改为凸
液面的最高点，并保持视线水平，避免视差。若仍机械地按凹液面的最低点读数，当然会产生系统误差；
此外，量筒若不是按这种液体校准，附着残留和弯月面形状也会带来额外误差。

（2）滴定管尖端的液滴受到向下的重力，同时液滴颈部周围的表面张力提供向上的合力。液滴很小时，
$mg$ 尚不足以克服表面张力，便会悬在管口。继续流入液体使液滴变大，或轻触接收容器内壁改变接触线，
液滴才可能落下。因此“暂时不滴”不一定意味着管口堵塞；只有在旋塞已打开、液柱压差充分而仍无液体
流出时，才需要检查气泡、旋塞和尖嘴堵塞。

（3）所谓引流，是让玻璃棒一端靠在烧杯嘴，另一端贴近接收容器内壁，使水沿玻璃棒连续流下。水对
洁净玻璃的附着与润湿使液流贴着玻璃棒运动，从而减小飞溅并便于控制落点。玻璃棒只起导流作用，不能
把它当作搅动液流的杠杆去敲碰容器。
\end{xesolution}

\Needspace{10\baselineskip}
\Needspace{16\baselineskip}\begin{origexamplebox}{【简易温度计】}\label{ch14:ex:s28-19-4}
细玻璃管连接着两个相同玻璃球，管壁光滑，管内有一个质量为 $m$ 的活塞，组成一个实验装置。
当装置水平放置时，活塞恰在管中点。问竖直放置时活塞的位移 $h$ 满足什么条件？可否根据这一原理
制作简易气温计？
\end{origexamplebox}

\Needspace{6\baselineskip}\begin{xesolution}
设细管横截面积为 $S$。装置水平、活塞居中时，两边封闭气体的压强相同；玻璃球相同，若两边温度
也相同，则两边气体的物质的量相同，记为 $\nu$，各自体积记为 $V_0$。

图\ref{ch14:fig:bulb-thermometer}标出活塞中点基准及下移量 $h$。
设两球之间可供活塞移动的直管有效长度为 $\ell_{\text{管}}$，还必须检查活塞是否留在直管内。
\bulbThermometerDiagram

竖直放置后，活塞在重力作用下向下移动 $h$。下腔体积减为 $V_0-Sh$，上腔体积增为 $V_0+Sh$。
设共同温度为 $T$，两侧气体压强分别为
\begin{equation}
  p_{\text{下}}=\frac{\nu RT}{V_0-Sh},\qquad
  p_{\text{上}}=\frac{\nu RT}{V_0+Sh}.
\end{equation}
活塞静止时，向上的压强差作用力与重力平衡：
\begin{equation}
  (p_{\text{下}}-p_{\text{上}})S=mg.
\end{equation}
代入物态方程，得到位移必须满足
\begin{equation}
  \boxed{\frac{2\nu RTS^2h}{V_0^2-S^2h^2}=mg},
  \qquad 0<h<\min\left(\frac{V_0}{S},\frac{\ell_{\text{管}}}2\right).
\end{equation}
反解温度可得
\begin{equation}
  \boxed{T(h)=\frac{mg\bigl(V_0^2-S^2h^2\bigr)}{2\nu RS^2h}}.
\end{equation}
因此，只要先用两个已知温度标定 $h$ 与 $T$ 的对应关系，原理上可以做成气体温度计。不过它并不是
等刻度温度计；实际灵敏度还会受到活塞摩擦、漏气、管径不均匀、玻璃热膨胀、两球温度不完全相同等
因素影响。题设的“管壁光滑”正是为了排除最显著的静摩擦误差。
\end{xesolution}

\Needspace{10\baselineskip}
\Needspace{16\baselineskip}\begin{origexamplebox}{【托里拆利实验】}\label{ch14:ex:s28-19-5}
有一水银气压计，因水银柱上的玻璃管内有微量空气，致使读数不准确。当大气压力为
$76\,\mathrm{cmHg}$ 时，此气压计读数为 $74\,\mathrm{cm}$，此时其上空气柱长为
$8.0\,\mathrm{cm}$。若气压计读数为 $72.0\,\mathrm{cm}$ 时，则正确的大气压力为多少？
（不考虑温度变化）
\end{origexamplebox}

\Needspace{6\baselineskip}\begin{xesolution}
先写带量纲的关系 $p_{\rm atm}=\rho_{\rm Hg}gH+p_g$。
若引入汞柱压强水头 $b_{\rm atm}=p_{\rm atm}/(\rho_{\rm Hg}g)$、$b_g=p_g/(\rho_{\rm Hg}g)$，
各个 $b$ 都是与压强等效的汞柱高度，单位为长度。管顶若是真空，$b_g=0$，于是
$b_{\rm atm}=H$；现在残余气体对汞柱上表面施压，压力平衡可等价写成
\begin{equation}
  b_{\rm atm}=H+b_g.
\end{equation}
例如，压强为 $76\,\mathrm{cmHg}$，就对应 $b_{\rm atm}=76\,\mathrm{cm}$；
不能直接把以帕为单位的压强与以米为单位的高度相加。

图\ref{ch14:fig:torricelli}采用足够大的汞槽，使两状态的槽液面高度近似不变。
\torricelliDiagram

第一状态中
\begin{equation}
  b_{g1}=b_{\rm atm,1}-H_1=(76-74)\,\mathrm{cm}=2\,\mathrm{cm},
  \qquad \ell_1=8.0\,\mathrm{cm}.
\end{equation}
第二次读数由 $74\,\mathrm{cm}$ 降到 $72\,\mathrm{cm}$，在管顶固定、管径均匀的条件下，水银柱
上端下降 $2.0\,\mathrm{cm}$，故残余空气柱长度变为
\begin{equation}
  \ell_2=8.0+2.0=10.0\,\mathrm{cm}.
\end{equation}
温度不变，且残余气体量不变，由波义耳定律（Boyle's law）
\begin{equation}
  \begin{gathered}
  p_{g1}\ell_1=p_{g2}\ell_2
  \quad\Longrightarrow\quad b_{g1}\ell_1=b_{g2}\ell_2,
  \\
  b_{g2}=2\,\mathrm{cm}\times\frac{8}{10}=1.6\,\mathrm{cm}.
  \end{gathered}
\end{equation}
所以第二状态的真实大气压为
\begin{equation}
  \begin{gathered}
  b_{\rm atm,2}=H_2+b_{g2}=(72.0+1.6)\,\mathrm{cm}=73.6\,\mathrm{cm},
  \\
  \boxed{p_{\rm atm,2}=73.6\,\mathrm{cmHg}}.
  \end{gathered}
\end{equation}
图中的 $H$ 是水银柱高度，$\ell$ 是管顶残余气柱长度；二者此消彼长，不能把气柱长度仍误取为
$8.0\,\mathrm{cm}$。
\end{xesolution}

\Needspace{11\baselineskip}
\Needspace{16\baselineskip}\begin{origexamplebox}{【加热气体】}\label{ch14:ex:s28-19-6}
如图所示，水平放置的汽缸内壁光滑，活塞厚度不计，在 $A$、$B$ 两处设有限制装置，使活塞只能在
$A$、$B$ 之间运动，$B$ 左面汽缸的容积为 $V_0$，$A$、$B$ 之间的容积为 $0.1V_0$。开始时活塞
在 $B$ 处，缸内气体的压强为 $0.9p_0$（$p_0$ 为大气压强），温度为 $297\,\mathrm K$，现缓慢
加热汽缸内气体，直至 $399.3\,\mathrm K$。

（1）活塞刚离开 $B$ 处时的温度 $T_B$；

（2）缸内气体最后的压强 $p'$；

（3）画出整个过程的 $p$--$V$ 图像。
\exampleNineteenSixFaithful
\end{origexamplebox}

\Needspace{6\baselineskip}\begin{xesolution}
活塞经历三个阶段：在 $B$ 处被挡住时定容升压；离开 $B$ 后在大气压下向 $A$ 移动，近似定压膨胀；
到达 $A$ 后再次被挡住，继续定容升压。图\ref{ch14:fig:heated-apparatus}中两个挡块分别限制向左、向右移动，
图\ref{ch14:fig:heated-path}把三个阶段与 $p$--$V$ 图中的三段折线一一对应。

\heatedCylinderDiagram

第一阶段 $V=V_0$。活塞刚能离开 $B$ 时，缸内压强升到外界大气压 $p_0$，故
\begin{equation}
  \frac{0.9p_0}{297}=\frac{p_0}{T_B},\qquad
  \boxed{T_B=330\,\mathrm K}.
\end{equation}

第二阶段 $p=p_0$，活塞从 $B$ 移到 $A$，体积由 $V_0$ 增至 $1.1V_0$。到达 $A$ 时
\begin{equation}
  \frac{V_0}{330}=\frac{1.1V_0}{T_A},\qquad
  T_A=363\,\mathrm K.
\end{equation}

第三阶段 $V=1.1V_0$，由 $363\,\mathrm K$ 升到 $399.3\,\mathrm K$：
\begin{equation}
  \frac{p_0}{363}=\frac{p'}{399.3},\qquad
  \boxed{p'=1.10p_0}.
\end{equation}
$p$--$V$ 图依次经过
\begin{equation}
  (V_0,0.9p_0)\longrightarrow(V_0,p_0)
  \longrightarrow(1.1V_0,p_0)\longrightarrow(1.1V_0,1.1p_0),
\end{equation}
即“竖直—水平—竖直”三段，分别对应定容、定压、定容过程。
\end{xesolution}

\Needspace{11\baselineskip}
\Needspace{16\baselineskip}\begin{origexamplebox}{【各种过程】}\label{ch14:ex:s28-19-7}
一个质量可不计的活塞将一定量的理想气体封闭在上端开口的直立圆筒形气缸内，活塞上堆放着铁砂，
如图所示。最初活塞搁置在气缸内壁的固定卡环上，气体柱的高度为 $H_0$，压强等于大气压强 $p_0$。
现对气体缓慢加热，当气体温度升高了 $\Delta T=60\,\mathrm K$ 时，活塞（及铁砂）开始离开卡环而
上升；继续加热直到气柱高度为 $H_1=1.5H_0$。此后，在维持温度不变的条件下逐渐取走铁砂，直到
铁砂全部取走时，气柱高度变为 $H_2=1.8H_0$。求此时气体的温度（不计活塞与气缸之间的摩擦）。
\figureNineteenSevenFaithful
\end{origexamplebox}

\Needspace{6\baselineskip}\begin{xesolution}
设初温为 $T_0$，活塞刚离开卡环时气体压强为 $P_1$。气缸横截面积不变，所以气体体积与气柱高度
成正比。图\ref{ch14:fig:s28-19-7}给出初态卡环与活塞的接触关系，
图\ref{ch14:fig:sand-states}按实际比例画出 $H_0,H_1,H_2$，分别对应三个状态中的气体体积。

\sandPistonDiagram

最后铁砂全部取走，活塞质量又可忽略，故气体最终压强等于大气压 $p_0$。取走铁砂的过程保持温度
不变，于是
\begin{equation}
  P_1H_1=p_0H_2,\qquad
  \frac{P_1}{p_0}=\frac{H_2}{H_1}=\frac{1.8}{1.5}=1.2.
\end{equation}
最初到活塞刚起动前，卡环使气体体积保持不变；温度从 $T_0$ 升到 $T_0+60$，压强从 $p_0$ 升到
$P_1$：
\begin{equation}
  \frac{P_1}{p_0}=\frac{T_0+60}{T_0}=1.2,
\end{equation}
解得
\begin{equation}
  \boxed{T_0=300\,\mathrm K}.
\end{equation}
活塞离开卡环后，铁砂质量尚未改变，外压恒定为 $P_1$，所以从 $H_0$ 膨胀至 $H_1=1.5H_0$ 的
过程是定压过程：
\begin{equation}
  \frac{T_1}{T_0+60}=\frac{H_1}{H_0}=1.5.
\end{equation}
因此在气柱达到 $H_1$ 时以及随后恒温取砂直至末态，气体温度均为
\begin{equation}
  \boxed{T_1=1.5\times360=540\,\mathrm K}.
\end{equation}
这里最容易颠倒的是等温取砂阶段的压强比：体积由 $1.5H_0$ 增到 $1.8H_0$，压强应由 $1.2p_0$
降到 $p_0$。
\end{xesolution}

\Needspace{11\baselineskip}
\Needspace{16\baselineskip}\begin{origexamplebox}{【气体膨胀】}\label{ch14:ex:s28-19-8}
如图所示，开口向上、粗细均匀的玻璃管长 $L=100\,\mathrm{cm}$，管内有一段高
$h=20\,\mathrm{cm}$ 的水银柱，封闭着长 $a=50\,\mathrm{cm}$ 的空气柱，大气压强
$p_0=76\,\mathrm{cmHg}$，温度 $t_0=27\,{}^\circ\mathrm C$。求温度至少升到多高时，
可使水银全部溢出？
\figureNineteenEightFaithful
\end{origexamplebox}

\Needspace{6\baselineskip}\begin{xesolution}
这是一个“先膨胀到管口，再边溢出边变化”的过程。仅求水银柱上端刚到管口的温度还不够；水银开始
溢出后，剩余汞柱变短，气体压强下降，但气体可占体积继续增大，两种作用竞争会形成一个温度势垒。

\mercurySpillDiagram

图\ref{ch14:fig:s28-19-8}给出初始长度；
图\ref{ch14:fig:mercury-spill}中 $x$ 表示剩余汞柱的高度，不是活塞位移。
为避免把压强与长度相加，引入 $b_0=p_0/(\rho_{\rm Hg}g)=76\,\mathrm{cm}$。
初态封闭气体压强为
\begin{equation}
  p_1=\rho_{\rm Hg}g(b_0+h)=96\,\mathrm{cmHg},\qquad T_0=27+273=300\,\mathrm K,
\end{equation}
这里沿用题目计算精度，将 $273.15$ 近似取作 $273$，故 $T_0\simeq300\,\mathrm K$。
气柱长度为 $a=50\,\mathrm{cm}$。水银柱上端刚到管口时，汞柱仍高 $20\,\mathrm{cm}$，气柱长度
变为 $L-h=80\,\mathrm{cm}$，由理想气体状态方程得到
\begin{equation}
  T_{\text{口}}=T_0\frac{96\times80}{96\times50}=480\,\mathrm K.
\end{equation}
但此时还没有“全部溢出”。令溢出过程中管内剩余汞柱高度为 $x$。它的下端距管底 $L-x$，所以封闭
气柱长为 $L-x$；汞柱平衡又给出气体压强 $p=\rho_{\rm Hg}g(b_0+x)$。于是准静态过程所需温度为
\begin{equation}
  T(x)=T_0\frac{(b_0+x)(L-x)}{(b_0+h)a},\qquad 0\le x\le h.
\end{equation}
要使水银能连续溢出，外界温度至少应达到该函数在区间 $[0,h]$ 上的最大值。令
\begin{equation}
  f(x)=(b_0+x)(L-x),\qquad
  f'(x)=L-b_0-2x,
\end{equation}
得到最大值位置
\begin{equation}
  x_*=\frac{L-b_0}{2}=12\,\mathrm{cm}.
\end{equation}
因此
\begin{equation}
  T_{\min}=300\times\frac{(76+12)(100-12)}{96\times50}
  =484\,\mathrm K,
\end{equation}
即
\begin{equation}
  \boxed{t_{\min}\simeq484-273=211\,{}^\circ\mathrm C}.
\end{equation}
图中的曲线在 $x=12\,\mathrm{cm}$ 处达到最高点：加热到 $480\,\mathrm K$ 只能使水银刚到管口，
还必须跨过 $484\,\mathrm K$ 的最高点，越过极大值后，在不降低加热温度的条件下，不再有静态平衡阻止剩余水银流出。
此后的实际运动未必继续准静态；极大值用于判断能否越过整条平衡路径。
\end{xesolution}

\Needspace{18\baselineskip}
\Needspace{16\baselineskip}\begin{origexamplebox}{【定性分析】}\label{ch14:ex:s28-19-9}
等截面U形管的两端封闭，两腿竖直。初态液面距弯管上缘的高度至少为 $L$，以保证以下过程中液面均在竖直段内。管内水柱把左右两团理想气体隔开。初态两液面等高，两侧气柱长度均为
$L$，压强和温度均为 $p_0,T_0$。现缓慢加热右侧气体至 $T_R>T_0$，左侧气体始终保持
$T_0$。水的密度为 $\rho$，忽略毛细现象。判断两液面的移动方向，并写出末态位移所满足的方程；
说明物理解是否唯一。


\closedUTubeDiagram
\end{origexamplebox}

\Needspace{6\baselineskip}\begin{xesolution}
图\ref{ch14:fig:closed-u}两端都封闭，水柱在弯管处连续连通。右侧升温后，若液面尚未来得及移动，其压强先增大，所以水柱被向左推：右液面下降、左液面上升。
设两液面各移动 $x$。截面积相同保证两侧体积变化量相等，因而末态气柱长度分别为 $L-x$、$L+x$。
由理想气体状态方程
\[
 p_L=\frac{p_0L}{L-x},\qquad
 p_R=\frac{p_0L}{L+x}\frac{T_R}{T_0}.
\]
右液面比左液面低 $2x$，静水压关系为
\[
 \boxed{\frac{p_0L(T_R/T_0)}{L+x}-\frac{p_0L}{L-x}=2\rho gx},
 \qquad 0<x<L.
\]
令左端减右端为 $F(x)$。有
\[
 F(0)=p_0\!\left(\frac{T_R}{T_0}-1\right)>0,\qquad
 \lim_{x\to L^-}F(x)=-\infty,
\]
并且
\[
 F'(x)=-\frac{p_0L(T_R/T_0)}{(L+x)^2}
       -\frac{p_0L}{(L-x)^2}-2\rho g<0.
\]
所以 $F(x)$ 在 $(0,L)$ 内严格递减且只过零一次，移动方向和末态位移均唯一。
当 $T_R\to T_0^+$ 时 $F(0)\to0$，唯一解趋于 $x=0$；这与没有温差便不移动一致。
题设的直腿长度条件保证 $0<x<L$ 内右液面不会提前进入弯头。
\end{xesolution}

\Needspace{22\baselineskip}
\Needspace{16\baselineskip}\begin{origexamplebox}{【L 型管】}\label{ch14:ex:s28-19-10}
粗细均匀的L形玻璃管左端封闭、上端开口。沿管中心线取坐标 $s$：初态封闭气体占据
$-a\le s\le0$，长为 $b$ 的水银柱占据 $0\le s\le b$；弯头在 $s=c$，竖直出口终点在
$s=c+d$。已知 $0<b<c$、$d\ge b$，初态气体温度为 $T_0$、大气压为 $p_0$，管截面积为
$S$，弯头曲率半径和容积相对各直管段忽略不计。缓慢改变气体温度，使水银准静态地向出口移动。求水银能够全部流出时气体必须达到的最低温度。


\lPipeInitialDiagram
\end{origexamplebox}

\Needspace{6\baselineskip}\begin{xesolution}
图\ref{ch14:fig:l-initial}规定全部长度与坐标；图\ref{ch14:fig:l-stages}分别画出水银的四种位置。
黑点标明后端，前端在尚未流出时位于沿程坐标 $x+b$。
\lPipeStagesDiagram
设水银柱后端移动了 $x$。封闭气体体积为 $V(x)=S(a+x)$。水银柱两端的竖直高度差记为
$h(x)$；它随水银柱位置分四段：
\[
h(x)=
\begin{cases}
0,&0\le x\le c-b,\\
x+b-c,&c-b\le x\le c,\\
b,&c\le x\le c+d-b,\\
c+d-x,&c+d-b\le x\le c+d.
\end{cases}
\]
准静态平衡要求气体压强为 $p(x)=p_0+\rho_{\rm Hg}g h(x)$。初态 $p(0)=p_0$、
$V(0)=Sa$，故沿着平衡路径所需温度为
\[
 T(x)=T_0\frac{p(x)V(x)}{p_0Sa}
 =T_0\frac{[p_0+\rho_{\rm Hg}g h(x)](a+x)}{p_0a}.
\]
要连续越过全过程，至少应达到 $T(x)$ 在 $0\le x\le c+d$ 上的最大值。前三段中压强不降而
体积增加，所以温度单调增加。最后一段令
\[
 \Phi(x)=[p_0+\rho_{\rm Hg}g(c+d-x)](a+x),
\]
逐项求导，
\[
 \Phi'(x)=p_0+\rho_{\rm Hg}g(c+d-a-2x),\qquad
 \Phi''(x)=-2\rho_{\rm Hg}g<0.
\]
所以最后一段是开口向下的抛物线，其唯一驻点为
\[
 x_*=\frac{p_0/(\rho_{\rm Hg}g)+c+d-a}{2}.
\]
把它截取到最后一段的允许区间：
\[
 \widehat x=\min\!\left\{c+d,
 \max\!\left[c+d-b,x_*\right]\right\}.
\]
于是所求最低温度为
\[
 \boxed{T_{\min}=T_0
 \frac{[p_0+\rho_{\rm Hg}g(c+d-\widehat x)](a+\widehat x)}{p_0a}}.
\]
若 $x_*\le c+d-b$，最后一段一开始便下降，最大值在它的起点；
若 $c+d-b<x_*<c+d$，最大值在内部驻点；
若 $x_*\ge c+d$，最后一段始终上升，最大值在全部流出的终点。
因而截取公式包含这三种情形。量纲检查中
$pV/(p_0V_0)$ 无量纲，故结果仍为温度。
\end{xesolution}

\Needspace{12\baselineskip}
\Needspace{16\baselineskip}\begin{origexamplebox}{【一维原子链】}\label{ch14:ex:s28-19-11}
如果认为单原子链中每个原子感受到的平均势能为莱纳德--琼斯势
\begin{equation}
  U(r)=4\varepsilon\left[\left(\frac{\sigma}{r}\right)^{12}
  -\left(\frac{\sigma}{r}\right)^6\right],
\end{equation}
其中 $\varepsilon$、$\sigma$ 为常数，$r$ 为所有原子之间的统一间距。求：

（1）零温下平衡态每原子结合能、晶格常数；

（2）把原子势能在平衡位置展开到二阶项和三阶项，求单原子在温度为 $T$ 时的振动幅度、原子链的
形变模量和热膨胀系数。
\end{origexamplebox}

\Needspace{6\baselineskip}\begin{xesolution}
先令 $U'(r_0)=0$。对势能求导：
\begin{equation}
  U'(r)=4\varepsilon\left(-\frac{12\sigma^{12}}{r^{13}}
  +\frac{6\sigma^6}{r^7}\right).
\end{equation}
于是 $r_0^6=2\sigma^6$，即
\begin{equation}
  \boxed{r_0=2^{1/6}\sigma}.
\end{equation}
把它代回势能，得到 $U(r_0)=-\varepsilon$。因此零温下的晶格常数为 $r_0$，从势阱底把一个原子
移到无穷远所需的结合能大小为
\begin{equation}
  \boxed{E_b=\varepsilon}.
\end{equation}

图\ref{ch14:fig:lennard-jones}标出真实最低点，并将局部二阶近似与完整势能比较。
\lennardJonesDiagram

令 $x=r-r_0$。在平衡位置作泰勒展开：
\begin{equation}
  U(r)=-\varepsilon+\frac12kx^2+\frac16\lambda x^3+O(x^4),
\end{equation}
二、三阶导数分别为
\[
 U''(r)=4\varepsilon\left(\frac{156\sigma^{12}}{r^{14}}-\frac{42\sigma^6}{r^8}\right),\qquad
 U'''(r)=4\varepsilon\left(-\frac{2184\sigma^{12}}{r^{15}}+\frac{336\sigma^6}{r^9}\right).
\]
利用 $\sigma^6/r_0^6=1/2$、$\sigma^{12}/r_0^{12}=1/4$，可分别提出 $r_0^{-2}$、$r_0^{-3}$，
得到 $4(156/4-42/2)=72$、$4(-2184/4+336/2)=-1512$。因此
\begin{equation}
  \boxed{k=U''(r_0)=\frac{72\varepsilon}{r_0^2}},
  \qquad
  \boxed{\lambda=U'''(r_0)=-\frac{1512\varepsilon}{r_0^3}}.
\end{equation}
二阶近似是简谐振子。其势能这个二次型自由度平均分得 $k_BT/2$，故
\begin{equation}
  \frac12k\langle x^2\rangle=\frac12k_BT,
  \qquad
  \boxed{x_{\rm rms}=r_0\sqrt{\frac{k_BT}{72\varepsilon}}}.
\end{equation}
若把“振动幅度”理解为平均总振动能 $k_BT$ 对应的经典转折幅度，则
\begin{equation}
  \frac12kA^2=k_BT,\qquad
  \boxed{A=r_0\sqrt{\frac{k_BT}{36\varepsilon}}}.
\end{equation}

设原子链应变为 $\eta=x/r_0$。每个原子的弹性能为
\begin{equation}
  \frac12kx^2=\frac12kr_0^2\eta^2.
\end{equation}
除以每个原子在链上占据的长度 $r_0$，与一维弹性能密度
$u=\tfrac12Y_{1\rm D}\eta^2$ 比较，得到
\begin{equation}
  \boxed{Y_{1\rm D}=kr_0=\frac{72\varepsilon}{r_0}}.
\end{equation}
这里 $Y_{1\rm D}$ 的量纲是力；若要得到通常三维杨氏模量，还需给定原子链的有效横截面积。

三次项 $\lambda<0$ 使势阱右缓左陡。热膨胀所需的是平均间距，
而不是左右转折点的几何中点。以下在粒子被约束于平衡距离附近、经典统计有效、振幅较小的条件下计算；
可用的温度窗口为 $\hbar\sqrt{k/m}\ll k_BT\ll\varepsilon$。下文积分是局部势阱的渐近近似，
不把粒子在无穷远自由解离的状态算入这条原子链的平衡分布。

令 $\beta=1/(k_BT)$，零阶高斯权重为 $w_0(x)=\exp(-\beta kx^2/2)$。
三次项小，玻尔兹曼权重可展开为
\[
 e^{-\beta(kx^2/2+\lambda x^3/6)}
 \simeq w_0(x)\left(1-\frac{\beta\lambda x^3}{6}\right).
\]
平均位移是归一化后的积分之比：
\[
 \langle x\rangle=
 \frac{\int xw_0(x)(1-\beta\lambda x^3/6)\,\mathrm dx}
 {\int w_0(x)(1-\beta\lambda x^3/6)\,\mathrm dx}.
\]
在局部小振动近似中将积分限延伸为 $(-\infty,\infty)$；奇函数积分为零，
所以分子中零阶的 $\int xw_0\,\mathrm dx=0$，分母中的一阶修正
$\int x^3w_0\,\mathrm dx=0$。于是
\[
 \langle x\rangle=-\frac{\beta\lambda}{6}\langle x^4\rangle_0.
\]
四阶矩也可直接推导。记 $a=\beta k/2$、$I(a)=\int e^{-ax^2}\mathrm dx=\sqrt{\pi/a}$，
对 $a$ 连续求两次导数：
\[
 I'(a)=-\int x^2e^{-ax^2}\mathrm dx,\qquad
 I''(a)=\int x^4e^{-ax^2}\mathrm dx=\frac{3\sqrt\pi}{4a^{5/2}}.
\]
因而 $\langle x^4\rangle_0=I''/I=3/(4a^2)=3(k_BT/k)^2$，最终
\[
 \langle x\rangle=-\frac{\lambda k_BT}{2k^2}>0,\qquad
 \boxed{\alpha=\frac1{r_0}\frac{\mathrm d\langle x\rangle}{\mathrm dT}
 =\frac{7k_B}{48\varepsilon}}.
\]
其量纲为 $\mathrm K^{-1}$；$\lambda<0$ 给出正膨胀，与图中的右缓左陡相符。
若令三次项为零，平均位移与这一阶热膨胀都为零。

还可比较一个容易混淆的几何估计。把同一能量的转折点写成
$x_\pm=\pm A+\delta$，其中 $kA^2/2=k_BT$。代入势能并保留一阶非谐项，
\[
 \frac{k}{2}(\pm A+\delta)^2+\frac{\lambda}{6}(\pm A+\delta)^3
 \simeq\frac{kA^2}{2}\pm kA\delta\pm\frac{\lambda A^3}{6},
\]
故 $\delta=-\lambda A^2/(6k)=-\lambda k_BT/(3k^2)$。
中点随温度变化所给出的系数为 $7k_B/(72\varepsilon)$，它不是上面所求的热膨胀系数。
原因在于粒子并不在两转折点之间匀速运动；必须使用时间或统计权重，不能把中点法与统计法并列成两种正确答案。
\end{xesolution}
}

\endgroup

\hypersetup{bookmarksdepth=2}

\hypersetup{bookmarksdepth=3}
\begingroup
\renewcommand{\Needspace}[1]{\par\begingroup
\dimen0=#1\relax\vskip0pt plus\dimen0\penalty-100
\vskip0pt plus-\dimen0\vskip\dimen0\penalty9999\vskip-\dimen0
\vskip0pt\endgroup}
\tcbset{breakable/.default=false}
\chapter{热力学第一定律}

{
\graphicspath{{/Users/dongxingxi/Documents/物理竞赛一轮资料/sample/xes-original/twenty-ninth-section/}{/Users/dongxingxi/Documents/物理竞赛一轮资料/sample/xes-original/twenty-ninth-section/assets/}}

% Diagrams for Lesson 20: the first law of thermodynamics.

\newcommand{\heatTransferMapDiagram}{%
\begin{center}
\begin{tikzpicture}[x=1cm,y=1cm,>=Latex]
  \node[draw=DeepBlue,fill=PaleBlue,thick,rounded corners=2pt,minimum width=2.6cm,minimum height=1.0cm] (sys) at (0,0) {热力学系统};
  \node[draw=RuleGray,fill=SoftGray,thick,rounded corners=2pt,minimum width=2.6cm,minimum height=1.0cm] (cond) at (-5.1,1.8) {直接接触（传导）};
  \node[draw=RuleGray,fill=SoftGray,thick,rounded corners=2pt,minimum width=2.6cm,minimum height=1.0cm] (conv) at (-5.1,-1.8) {流体整体运动（对流）};
  \node[draw=RuleGray,fill=SoftGray,thick,rounded corners=2pt,minimum width=2.6cm,minimum height=1.0cm] (rad) at (5.1,1.8) {电磁波（辐射）};
  \node[draw=RuleGray,fill=SoftGray,thick,rounded corners=2pt,minimum width=2.6cm,minimum height=1.0cm] (work) at (5.1,-1.8) {边界位移（功）};
  \draw[<->,DeepBlue,very thick] (cond)--node[above,sloped,font=\small] {$Q$，需要介质}(sys);
  \draw[<->,DeepBlue,very thick] (conv)--node[below,sloped,font=\small] {$Q$，需要介质}(sys);
  \draw[<->,DeepBlue,very thick] (rad)--node[above,sloped,font=\small] {$Q$，可穿过真空}(sys);
  \draw[<->,DeepBlue,very thick] (work)--node[below,sloped,font=\small] {$W$}(sys);
\end{tikzpicture}
\end{center}\nopagebreak[2]}

\newcommand{\pistonWorkDiagram}{%
\begin{center}
\begin{tikzpicture}[x=1cm,y=1cm,>=Latex]
  \begin{scope}[shift={(-4.8,0)}]
    \draw[Ink,very thick] (-1.35,-1.6)--(-1.35,1.7) (1.35,-1.6)--(1.35,1.7) (-1.35,-1.6)--(1.35,-1.6);
    \draw[fill=SoftGray,draw=Ink,thick] (-1.2,0.25) rectangle (1.2,0.55);
    \draw[Ink,very thick] (0,0.55)--(0,2.15);
    \draw[->,DeepBlue,very thick] (0,2.2)--(0,2.95) node[right] {$\dd y$};
    \draw[->,DeepBlue,very thick] (-0.8,0.1)--(-0.8,0.95) node[left] {$F=PA$};
    \node at (0,-0.55) {$P,V$};
    \node[font=\small] at (0,-2.1) {$\dd V=A\,\dd y$};
  \end{scope}
  \begin{scope}[shift={(2.2,-1.55)},x=1.15cm,y=0.95cm]
    \draw[axis] (0,0)--(5.4,0) node[right] {$V$};
    \draw[axis] (0,0)--(0,4.7) node[above] {$P$};
    \draw[fill=PaleBlue,draw=none] (1,3.8) .. controls (2.0,2.3) and (3.4,1.35) .. (4.6,1.0)--(4.6,0)--(1,0)--cycle;
    \draw[DeepBlue,very thick,postaction={decorate},decoration={markings,mark=at position .55 with {\arrow{Latex}}}]
      (0.65,4.25) .. controls (1.7,2.25) and (3.4,1.25) .. (5.0,0.9);
    \fill[Ink] (1,3.8) circle (2pt) node[above right] {$i$};
    \fill[Ink] (4.6,1.0) circle (2pt) node[above right] {$f$};
    \draw[guide] (1,0)--(1,3.8) (4.6,0)--(4.6,1.0);
    \node[DeepBlue,font=\small] at (2.65,0.65) {$W^{\rm out}=\displaystyle\int_{V_i}^{V_f}P\,\dd V$};
  \end{scope}
\end{tikzpicture}
\end{center}\nopagebreak[2]}

\newcommand{\processComparisonDiagram}{%
\begin{center}
\begin{tikzpicture}[x=1cm,y=1cm,>=Latex]
  \foreach \x/\title/\line/\note in {
    -5.25/{等容}/{V=\mathrm{const}}/{W=0},
    -1.75/{等压}/{P=\mathrm{const}}/{Q=n c_p\Delta T},
    1.75/{等温}/{T=\mathrm{const}}/{\Delta U=0},
    5.25/{绝热}/{Q=0}/{PV^\gamma=\mathrm{const}}
  }{
    \draw[fill=SoftGray,draw=RuleGray,rounded corners=1pt] (\x-1.45,-1.05) rectangle (\x+1.45,1.05);
    \node[font=\bfseries] at (\x,0.55) {\title};
    \node[font=\small] at (\x,0.0) {$\line$};
    \node[font=\small,DeepBlue] at (\x,-0.55) {$\note$};
  }
\end{tikzpicture}
\end{center}\nopagebreak[2]}

\newcommand{\adiabaticCurvesDiagram}{%
\begin{center}
\begin{tikzpicture}[x=1.25cm,y=1.0cm,>=Latex]
  \draw[axis] (0,0)--(6.2,0) node[right] {$V$};
  \draw[axis] (0,0)--(0,5.0) node[above] {$P$};
  \draw[DeepBlue,very thick,domain=1:5.7,samples=120,smooth] plot (\x,{4.5/\x});
  \draw[Ink,very thick,domain=1:5.7,samples=120,smooth] plot (\x,{4.5/(\x^1.4)});
  \fill[Ink] (1,4.5) circle (2pt) node[above right] {同一初态};
  \node[DeepBlue,anchor=west] at (4.2,1.15) {等温 $PV=\mathrm{const}$};
  \node[Ink,anchor=west] at (3.8,0.42) {绝热 $PV^\gamma=\mathrm{const}$};
  \draw[->,TextGray,thick] (2.0,2.2)--(3.3,1.1) node[midway,above right,font=\small] {膨胀降温};
  \node[align=left,anchor=west,font=\small] at (7.0,2.6)
    {$TV^{\gamma-1}=\mathrm{const}$\\[3pt]
     $\displaystyle\frac{T^\gamma}{P^{\gamma-1}}=\mathrm{const}$\\[3pt]
     $P=\mathrm{const}\cdot\rho^\gamma$};
\end{tikzpicture}
\end{center}\nopagebreak[2]}

\newcommand{\convectionCriterionDiagram}{%
\begin{center}
\begin{tikzpicture}[x=1cm,y=1cm,>=Latex]
  \draw[axis] (-5.3,-2.0)--(-5.3,2.7) node[above] {$h$};
  \fill[PaleBlue] (-5.1,-1.75) rectangle (-1.25,2.25);
  \draw[DeepBlue,very thick,->] (-3.15,-1.1)--(-3.15,1.25) node[midway,right] {气团上升};
  \node at (-3.15,-1.45) {$T$};
  \node at (-3.15,1.65) {$T+\dd T_1$};
  \draw[<->,TextGray] (-1.7,-1.1)--node[right] {$\dd h$}(-1.7,1.25);
  \node[align=left,anchor=west] at (0.1,1.4)
    {上升气团：$\displaystyle \dd T_1=-\frac{\mu g}{c_p}\,\dd h$};
  \node[align=left,anchor=west] at (0.1,0.35)
    {环境气体：$\displaystyle \dd T_2=\frac{\dd T}{\dd h}\,\dd h$};
  \node[draw=DeepBlue,fill=PaleBlue,thick,rounded corners=2pt,align=center] at (4.0,-1.15)
    {$T+\dd T_1>T+\dd T_2$\\[3pt]$\displaystyle\frac{\dd T}{\dd h}<-\frac{\mu g}{c_p}$};
\end{tikzpicture}
\end{center}\nopagebreak[2]}

\newcommand{\threeCyclesDiagram}{%
\begin{center}
\begin{tikzpicture}[x=1cm,y=1cm,>=Latex]
  \foreach \sx/\title/\kind in {-5/{矩形循环}/0,0/{上三角循环}/1,5/{下三角循环}/2}{
    \begin{scope}[shift={(\sx,-1.55)},x=1.05cm,y=0.95cm]
      \draw[axis] (0,0)--(3.8,0) node[right] {$V$};
      \draw[axis] (0,0)--(0,3.8) node[above] {$P$};
      \draw[guide] (0,0.8)--(3.25,0.8) (0,3.05)--(3.25,3.05) (0.8,0)--(0.8,3.45) (3.1,0)--(3.1,3.45);
      \ifnum\kind=0
        \draw[DeepBlue,very thick,
          postaction={decorate},decoration={markings,
            mark=at position .14 with {\arrow{Latex}},
            mark=at position .39 with {\arrow{Latex}},
            mark=at position .64 with {\arrow{Latex}},
            mark=at position .89 with {\arrow{Latex}}}]
          (0.8,0.8)--(0.8,3.05)--(3.1,3.05)--(3.1,0.8)--cycle;
      \else\ifnum\kind=1
        \draw[DeepBlue,very thick,
          postaction={decorate},decoration={markings,
            mark=at position .17 with {\arrow{Latex}},
            mark=at position .48 with {\arrow{Latex}},
            mark=at position .82 with {\arrow{Latex}}}]
          (0.8,0.8)--(0.8,3.05)--(3.1,3.05)--cycle;
      \else
        \draw[DeepBlue,very thick,
          postaction={decorate},decoration={markings,
            mark=at position .18 with {\arrow{Latex}},
            mark=at position .51 with {\arrow{Latex}},
            mark=at position .83 with {\arrow{Latex}}}]
          (0.8,0.8)--(3.1,3.05)--(3.1,0.8)--cycle;
      \fi\fi
      \node[font=\scriptsize,left] at (0,0.8) {$P_1$};
      \node[font=\scriptsize,left] at (0,3.05) {$P_2$};
      \node[font=\scriptsize,below] at (0.8,0) {$V_1$};
      \node[font=\scriptsize,below] at (3.1,0) {$V_2$};
    \end{scope}
    \node[font=\bfseries] at (\sx,2.55) {\title};
  }
\end{tikzpicture}
\end{center}\nopagebreak[2]}

\newcommand{\standardHeatCyclesDiagram}{%
\begin{center}
\begin{tikzpicture}[x=1cm,y=1cm,>=Latex]
  \foreach \sx/\title/\leftlab/\rightlab in {-5/{卡诺}/{T_h}/{T_c},0/{奥托}/{V_1}/{V_2},5/{布雷顿}/{P_1}/{P_2}}{
    \begin{scope}[shift={(\sx,-1.4)},x=1cm,y=0.85cm]
      \draw[axis] (0,0)--(3.7,0) node[right] {$V$};
      \draw[axis] (0,0)--(0,4.0) node[above] {$P$};
      \ifdim\sx pt=-5pt
        \draw[DeepBlue,very thick,postaction={decorate},decoration={markings,mark=at position .2 with {\arrow{Latex}},mark=at position .7 with {\arrow{Latex}}}]
          (0.8,3.2) .. controls (1.5,2.7) and (2.2,2.35) .. (3.0,2.2)
          .. controls (2.75,1.4) and (2.45,1.0) .. (2.1,0.8)
          .. controls (1.55,0.95) and (1.1,1.15) .. (0.75,1.45)
          .. controls (0.7,2.05) and (0.72,2.65) .. cycle;
        \node[font=\scriptsize] at (2.0,3.0) {$T_h$}; \node[font=\scriptsize] at (1.55,0.45) {$T_c$};
      \else\ifdim\sx pt=0pt
        \draw[DeepBlue,very thick,->] (0.85,0.8)--(0.85,3.2);
        \draw[DeepBlue,very thick] (0.85,3.2) .. controls (1.45,2.35) and (2.25,1.65) .. (3.0,1.35);
        \draw[DeepBlue,very thick,->] (3.0,1.35)--(3.0,0.55);
        \draw[DeepBlue,very thick,->] (3.0,0.55) .. controls (2.2,0.75) and (1.45,1.0) .. (0.85,0.8);
        \node[font=\scriptsize,below] at (0.85,0) {$V_{\min}$}; \node[font=\scriptsize,below] at (3.0,0) {$V_{\max}$};
      \else
        \draw[DeepBlue,very thick,->] (0.75,3.05)--(2.7,2.35);
        \draw[DeepBlue,very thick,->] (2.7,2.35) .. controls (3.05,1.5) and (2.8,1.0) .. (2.3,0.75);
        \draw[DeepBlue,very thick,->] (2.3,0.75)--(0.7,1.2);
        \draw[DeepBlue,very thick,->] (0.7,1.2) .. controls (0.4,2.0) and (0.45,2.6) .. (0.75,3.05);
        \node[font=\scriptsize] at (1.7,2.95) {$P_h$}; \node[font=\scriptsize] at (1.55,0.55) {$P_l$};
      \fi\fi
    \end{scope}
    \node[font=\bfseries] at (\sx,2.45) {\title};
  }
\end{tikzpicture}
\end{center}\nopagebreak[2]}

\newcommand{\variableHeatCapacityDiagram}{%
\begin{center}
\begin{tikzpicture}[x=1.1cm,y=1.0cm,>=Latex]
  \draw[axis] (0,0)--(6.2,0) node[right] {$V$};
  \draw[axis] (0,0)--(0,4.8) node[above] {$P$};
  \draw[DeepBlue,very thick] (0.8,3.9)--(5.2,0.75);
  \draw[guide] (3.0,0)--(3.0,2.33);
  \node[below,font=\small] at (3.0,0) {$V_0/2$};
  \node[DeepBlue,above right] at (1.15,3.5) {$\displaystyle\frac{P}{P_0}+\frac{V}{V_0}=1$};
  \node[align=left,anchor=west,font=\small] at (6.7,2.7)
    {$\displaystyle T(V)=\frac{P_0}{R}\!\left(V-\frac{V^2}{V_0}\right)$\\[5pt]
     $V=V_0/2$ 处 $\dd T/\dd V=0$，\\过程热容出现发散。};
\end{tikzpicture}
\end{center}\nopagebreak[2]}

\newcommand{\isothermalAtmosphereDiagram}{%
\begin{center}
\begin{tikzpicture}[x=1cm,y=1cm,>=Latex]
  \shade[bottom color=DeepBlue!22,top color=white] (-5.7,-2.0) rectangle (-1.8,2.7);
  \draw[axis] (-6.0,-2.0)--(-6.0,3.0) node[above] {$h$};
  \foreach \y/\lab in {-1.5/{P_0,n_0},0/{P_0\e^{-h/H}},1.6/{P_0\e^{-2h/H}}}{
    \draw[TextGray] (-6.0,\y)--(-5.7,\y); \node[font=\small,anchor=west] at (-5.3,\y) {$\lab$};
  }
  \node[draw=DeepBlue,fill=PaleBlue,thick,rounded corners=2pt,align=center] at (3.0,0.4)
    {$\displaystyle P(h)=P_0\e^{-h/H}$\\[5pt]
     $\displaystyle n(h)=n_0\e^{-h/H}$\\[5pt]
     $\displaystyle H=\frac{RT}{\mu g}$};
\end{tikzpicture}
\end{center}\nopagebreak[2]}

\newcommand{\capillaryBubbleDiagram}{%
\begin{center}
\begin{tikzpicture}[x=1cm,y=1cm,>=Latex]
  \begin{scope}[shift={(-4.3,0)}]
    \fill[DeepBlue!12] (-2.2,-1.8) rectangle (2.2,-0.5);
    \draw[DeepBlue,very thick] (-2.2,-0.5)--(2.2,-0.5);
    \draw[Ink,very thick] (-0.45,-1.8)--(-0.45,2.6) (0.45,-1.8)--(0.45,2.6);
    \fill[DeepBlue!18] (-0.42,-1.8) rectangle (0.42,1.35);
    \draw[DeepBlue,very thick] (-0.42,1.35) arc[start angle=180,end angle=0,x radius=0.42,y radius=0.24];
    \draw[<->,TextGray] (0.9,-0.5)--node[right] {$h$}(0.9,1.35);
    \node[font=\small] at (0,-2.2) {$h=2\sigma/(\rho gR)$};
  \end{scope}
  \begin{scope}[shift={(3.7,0)}]
    \shade[ball color=DeepBlue!18,opacity=0.55] (0,0.2) circle (1.65);
    \draw[DeepBlue,very thick] (0,0.2) circle (1.65);
    \draw[DeepBlue,thick] (0,0.2) circle (1.53);
    \draw[->,Ink,thick] (0,0.2)--(1.17,1.37) node[midway,above left] {$R$};
    \node at (0,0.2) {$p_{\rm in}$};
    \node[font=\small] at (0,-2.2) {$p_{\rm in}-p_0=4\sigma/R$};
  \end{scope}
\end{tikzpicture}
\end{center}\nopagebreak[2]}

\newcommand{\cycleExampleDiagram}{%
\begin{center}
\begin{tikzpicture}[x=1.2cm,y=0.9cm,>=Latex]
  \draw[axis] (0,0)--(6.2,0) node[right] {$V$};
  \draw[axis] (0,0)--(0,5.2) node[above] {$P$};
  \draw[guide] (0,4.2)--(1.25,4.2);
  \draw[guide] (0,1.05)--(5.1,1.05);
  \draw[guide] (1.25,0)--(1.25,4.2);
  \node[left,font=\small] at (0,4.2) {$P_1$};
  \node[left,font=\small] at (0,1.05) {$P_1/4$};
  \node[below,font=\small] at (1.25,0) {$V_1$};
  \draw[DeepBlue,very thick,->] (1.25,4.2)--(1.25,1.05);
  \draw[DeepBlue,very thick,->] (1.25,1.05)--(5.0,1.05);
  \draw[DeepBlue,very thick,postaction={decorate},decoration={markings,mark=at position .52 with {\arrow{Latex}}}]
    (5.0,1.05) .. controls (3.4,1.3) and (1.8,2.35) .. (1.25,4.2);
  \fill[Ink] (1.25,4.2) circle (2pt) node[above right] {$a$};
  \fill[Ink] (1.25,1.05) circle (2pt) node[above right] {$b$};
  \fill[Ink] (5.0,1.05) circle (2pt) node[below right] {$c$};
  \node[font=\small,DeepBlue] at (3.45,2.4) {$PV=P_1V_1$};
  \node[align=left,anchor=west,font=\small] at (7.0,2.6)
    {$V_c=4V_1$\\[3pt]
     $a\to b$：等容降温\\
     $b\to c$：等压吸热\\
     $c\to a$：等温压缩};
\end{tikzpicture}
\end{center}\nopagebreak[2]}

\newcommand{\straightProcessDiagram}{%
\begin{center}
\begin{tikzpicture}[x=1.2cm,y=1.0cm,>=Latex]
  \draw[axis] (0,0)--(5.8,0) node[right] {$V$};
  \draw[axis] (0,0)--(0,4.8) node[above] {$P$};
  \draw[guide] (0,3.8)--(1.2,3.8);
  \draw[guide] (0,1.9)--(4.6,1.9);
  \node[left,font=\small] at (0,3.8) {$P_A$};
  \node[left,font=\small] at (0,1.9) {$P_A/2$};
  \draw[DeepBlue,very thick,->] (1.2,3.8)--(4.6,1.9);
  \fill[Ink] (1.2,3.8) circle (2pt) node[above right] {$A$};
  \fill[Ink] (4.6,1.9) circle (2pt) node[below right] {$B$};
  \draw[guide] (1.2,0)--(1.2,3.8);
  \draw[guide] (4.6,0)--(4.6,1.9);
  \node[below,font=\small] at (1.2,0) {$V_A$};
  \node[below,font=\small] at (4.6,0) {$V_B$};
  \node[align=left,anchor=west,font=\small] at (6.7,2.7)
    {$\displaystyle W^{\rm out}=\frac{P_A+P_B}{2}(V_B-V_A)$\\[5pt]
     $\displaystyle \Delta U=\frac32(P_BV_B-P_AV_A)$};
\end{tikzpicture}
\end{center}\nopagebreak[2]}

\newcommand{\liftingPistonDiagram}{%
\begin{center}
\begin{tikzpicture}[x=1cm,y=1cm,>=Latex]
  \draw[Ink,very thick] (-4.7,-2.0)--(-4.7,3.2) (-1.7,-2.0)--(-1.7,3.2) (-4.7,-2.0)--(-1.7,-2.0);
  \draw[fill=SoftGray,draw=Ink,thick] (-4.55,0.15) rectangle (-1.85,0.48);
  \draw[->,DeepBlue,very thick] (-3.2,0.5)--(-3.2,1.65) node[right] {缓慢上提};
  \node at (-3.2,-0.7) {$p,T$};
  \draw[->,Ink,thick] (-3.2,-1.95)--(-3.2,-2.8) node[right] {$Mg$};
  \draw[->,TextGray,thick] (-4.0,-2.0)--(-4.0,-1.15) node[midway,left] {$p_0A$};
  \draw[->,TextGray,thick] (-2.5,-1.15)--(-2.5,-2.0) node[midway,right] {$pA$};
  \node[draw=DeepBlue,fill=PaleBlue,thick,rounded corners=2pt,align=left,anchor=west] at (0.2,0.2)
    {刚离地：$(p_0-p)A=Mg$\\[4pt]
     绝热膨胀：$\displaystyle\frac{T^\gamma}{p^{\gamma-1}}=\mathrm{const}$\\[4pt]
     $\displaystyle T_f=T_0\left(1-\frac{Mg}{p_0A}\right)^{(\gamma-1)/\gamma}$};
\end{tikzpicture}
\end{center}\nopagebreak[2]}

% Original vector redraws replacing all figures cropped from the source PDF.
% The physical relationships are retained, while composition, line work,
% colour, labels and visual hierarchy are newly designed for this handout.


\newcommand{\perpetualMotionRedraw}{%
\begin{center}
\begin{tikzpicture}[x=1cm,y=1cm,>=Latex]
  \foreach \cx/\title in {-5.1/{偏心重锤轮},0/{循环配重带},5.1/{单向摆锤轮}}{
    \draw[fill=SoftGray,draw=RuleGray,rounded corners=2pt] (\cx-2.15,-2.35) rectangle (\cx+2.15,2.65);
    \node[font=\small\bfseries] at (\cx,2.25) {\title};
  }
  % Deliberately new overbalanced-wheel construction.
  \begin{scope}[shift={(-5.1,0)}]
    \draw[Ink,very thick] (-1.55,-1.85)--(1.55,-1.85) (0,-1.85)--(0,-1.15);
    \draw[Ink,very thick] (0,0.15) circle (1.3);
    \foreach \a in {0,45,...,315}{
      \draw[TextGray,thick] (0,0.15)--(\a:1.05);
      \pgfmathsetmacro{\rr}{ifthenelse(cos(\a)>0,1.3,0.78)}
      \fill[DeepBlue] (\a:\rr) ++(0,0.15) circle (2.5pt);
    }
    \draw[->,DeepBlue,very thick] (1.55,0.75) arc[start angle=25,end angle=-70,radius=1.05];
    \node[font=\scriptsize,align=center] at (0,-2.08) {设想：右侧力矩始终较大};
  \end{scope}
  % New closed-loop chain construction.
  \begin{scope}[shift={(0,0)}]
    \draw[Ink,very thick,rounded corners=12pt] (-1.15,-1.45) rectangle (1.15,1.35);
    \foreach \y in {-1.05,-0.45,0.15,0.75}{
      \draw[fill=DeepBlue!16,draw=DeepBlue,thick] (-1.35,\y-0.17) rectangle (-0.93,\y+0.17);
      \draw[fill=DeepBlue!35,draw=DeepBlue,thick] (0.86,\y-0.27) rectangle (1.44,\y+0.27);
    }
    \draw[->,DeepBlue,very thick] (1.72,0.55)--(1.72,-0.5);
    \draw[->,TextGray,thick] (-1.72,-0.5)--(-1.72,0.55);
    \node[font=\scriptsize,align=center] at (0,-2.08) {设想：重块一侧持续下落};
  \end{scope}
  % New one-way pendulum construction.
  \begin{scope}[shift={(5.1,0)}]
    \draw[Ink,very thick] (-1.55,-1.85)--(1.55,-1.85) (0,-1.85)--(0,-1.05);
    \draw[Ink,very thick] (0,0.18) circle (1.18);
    \foreach \a in {10,50,...,330}{
      \draw[TextGray,thick] (0,0.18)--(\a:1.05);
      \draw[Ink,thick] (\a:1.05) ++(0,0.18)--++(0,-0.34);
      \fill[DeepBlue] (\a:1.05) ++(0,-0.16) circle (2.7pt);
    }
    \draw[->,DeepBlue,very thick] (-1.48,0.85) arc[start angle=155,end angle=245,radius=1.08];
    \node[font=\scriptsize,align=center] at (0,-2.08) {设想：摆锤只在一侧做功};
  \end{scope}
  \node[draw=DeepBlue,fill=PaleBlue,thick,rounded corners=2pt,font=\small,align=center]
    at (0,-3.05) {共同漏洞：装置回到原状态时，若无外界能量输入，就不可能留下正的净功};
\end{tikzpicture}
\end{center}\nopagebreak[2]}

\newcommand{\thermalPlumeRedraw}{%
\begin{center}
\begin{tikzpicture}[x=1cm,y=1cm,>=Latex]
  \shade[bottom color=DeepBlue!8,top color=white] (-4.9,-2.35) rectangle (5.0,2.75);
  % Newly drawn stylised hand and forearm.
  \path[fill=orange!28,draw=Ink,thick]
    (-4.65,-2.1) .. controls (-3.6,-1.55) and (-2.65,-1.05) .. (-1.65,-0.8)
    .. controls (-1.1,-0.65) and (-0.62,-0.62) .. (-0.18,-0.38)
    .. controls (0.18,-0.18) and (0.08,0.13) .. (-0.22,0.18)
    .. controls (-0.7,0.22) and (-1.12,0.08) .. (-1.58,0.04)
    .. controls (-1.05,0.28) and (-0.62,0.48) .. (-0.23,0.75)
    .. controls (0.13,1.0) and (-0.05,1.28) .. (-0.38,1.24)
    .. controls (-1.02,1.15) and (-1.56,0.76) .. (-2.08,0.55)
    .. controls (-2.7,0.30) and (-3.05,0.05) .. (-3.55,-0.35)
    .. controls (-4.05,-0.78) and (-4.35,-1.35) .. (-4.65,-2.1) -- cycle;
  \foreach \x/\h/\bend in {-2.8/2.0/0.35,-2.15/2.55/-0.25,-1.45/2.15/0.30,-0.78/1.7/-0.2}{
    \draw[orange!85!red,very thick,->]
      (\x,0.35) .. controls (\x+\bend,0.9) and (\x-\bend,1.45) .. (\x+0.15,\h);
  }
  \foreach \x/\y in {-3.2/0.25,-2.65/1.0,-2.05/1.55,-1.4/1.0,-0.9/1.55,0.05/0.95}{
    \fill[orange!80!red] (\x,\y) circle (1.5pt);
  }
  \node[draw=DeepBlue,fill=PaleBlue,rounded corners=2pt,align=left,anchor=west] at (1.0,1.25)
    {手掌加热邻近空气\\热空气密度减小并上升};
  \draw[->,DeepBlue,thick] (1.15,0.85) .. controls (0.35,0.65) and (-0.15,1.1) .. (-0.55,1.45);
  \node[font=\small,TextGray,anchor=west] at (1.0,-0.15) {上升暖流携带热量：对流};
  \draw[very thick,orange!85!red] (1.0,-0.75)--(2.0,-0.75) node[right,font=\small,Ink] {温度较高};
  \draw[very thick,DeepBlue!55] (1.0,-1.22)--(2.0,-1.22) node[right,font=\small,Ink] {环境空气};
\end{tikzpicture}
\end{center}\nopagebreak[2]}

\newcommand{\stellarConvectionRedraw}{%
\begin{center}
\begin{tikzpicture}[x=1cm,y=1cm,>=Latex]
  \begin{scope}[shift={(-3.9,0)}]
    \shade[ball color=orange!55] (0,0) circle (1.75);
    \clip (0,0) circle (1.75);
    \foreach \x/\y/\r in {-1.0/0.8/0.45,-0.25/1.05/0.40,0.65/0.8/0.48,1.15/0.1/0.42,0.6/-0.65/0.5,-0.35/-0.75/0.48,-1.15/-0.25/0.42,-0.25/0.1/0.55}{
      \draw[orange!85!black,opacity=.7,thick] (\x,\y) circle (\r);
      \fill[yellow!75!white,opacity=.65] (\x,\y) circle (0.11);
    }
  \end{scope}
  \draw[Ink,thick] (-3.9,0) circle (1.75);
  \node[font=\small\bfseries] at (-3.9,-2.2) {太阳表面：较细的对流颗粒};
  \begin{scope}[shift={(3.35,0)}]
    \shade[ball color=red!65!orange] (0,0) circle (2.45);
    \clip (0,0) circle (2.45);
    \foreach \x/\y/\r in {-1.3/1.2/0.62,-0.15/1.4/0.68,1.25/1.0/0.72,1.55/-0.25/0.72,0.75/-1.2/0.78,-0.55/-1.35/0.72,-1.55/-0.5/0.65,-0.35/0.15/0.92,0.95/0.2/0.72}{
      \draw[red!70!black,opacity=.72,very thick] (\x,\y) circle (\r);
      \fill[orange!80!yellow,opacity=.65] (\x,\y) circle (0.16);
    }
  \end{scope}
  \draw[Ink,thick] (3.35,0) circle (2.45);
  \node[font=\small\bfseries] at (3.35,-2.85) {红巨星表面：尺度更大的对流胞};
  \foreach \cx/\cy in {-4.5/-0.5,-3.4/0.35,2.4/-0.5,3.6/0.55,4.2/-0.85}{
    \draw[->,DeepBlue,thick] (\cx,\cy-0.25)--(\cx,\cy+0.35);
  }
  \node[draw=DeepBlue,fill=PaleBlue,rounded corners=2pt,font=\small] at (0,-3.55)
    {亮区：热物质上升\qquad 暗边界：冷物质下沉};
\end{tikzpicture}
\end{center}\nopagebreak[2]}

\newcommand{\thermalRadiationRedraw}{%
\begin{center}
\begin{tikzpicture}[x=1cm,y=1cm,>=Latex]
  \draw[Ink,thick,rounded corners=3pt,fill=SoftGray] (-5.4,-1.45) rectangle (5.4,1.45);
  \shade[left color=red!85!black,middle color=orange!85,right color=DeepBlue!35]
    (-4.85,-0.6) rectangle (4.85,0.6);
  \draw[Ink,very thick] (-4.85,-0.6) rectangle (4.85,0.6);
  \node[white,font=\bfseries] at (-3.65,0) {高温区};
  \node[Ink,font=\bfseries] at (3.65,0) {低温区};
  \foreach \x/\s in {-4.4/1.0,-3.55/.85,-2.7/.7}{
    \draw[orange!90!red,very thick,decorate,decoration={snake,amplitude=1.3mm,segment length=4mm}]
      (\x,0.78)--(\x,0.78+\s);
  }
  \foreach \x/\s in {2.9/.42,3.75/.32,4.45/.25}{
    \draw[DeepBlue!70,thick,decorate,decoration={snake,amplitude=.8mm,segment length=4mm}]
      (\x,0.78)--(\x,0.78+\s);
  }
  \node[font=\small,anchor=east] at (-2.2,-1.95) {$j_{\rm rad}=\sigma T^4$：温度升高，辐射急剧增强};
  \draw[->,DeepBlue,thick] (0,-2.0)--(3.8,-2.0) node[right,font=\small] {温度降低};
\end{tikzpicture}
\end{center}\nopagebreak[2]}

\newcommand{\heatCapacityRedraw}{%
\begin{center}
\begin{tikzpicture}[x=1cm,y=1cm,>=Latex]
  \draw[axis] (0,0)--(11.2,0) node[right] {$T/\mathrm K$（对数刻度）};
  \draw[axis] (0,0)--(0,5.6) node[above] {$c_v$};
  \foreach \x/\lab in {0.8/{10},2.8/{50},4.3/{100},5.9/{500},7.4/{10^3},9.2/{5\!\times\!10^3},10.7/{10^4}}{
    \draw[TextGray] (\x,0.08)--(\x,-0.08) node[below,font=\scriptsize] {$\lab$};
  }
  \foreach \y/\lab in {1.35/{\frac32R},3.15/{\frac52R},4.95/{\frac72R}}{
    \draw[guide] (0,\y)--(10.75,\y);
    \node[left,font=\small] at (0,\y) {$\lab$};
  }
  \draw[DeepBlue,very thick]
    (0.7,1.35)--(2.0,1.35)
    .. controls (3.0,1.35) and (3.15,3.15) .. (4.65,3.15)
    --(6.45,3.15)
    .. controls (7.5,3.15) and (7.7,4.95) .. (9.35,4.95)
    --(10.65,4.95);
  \draw[<->,Ink,thick] (1.3,0.18)--(1.3,1.28) node[midway,right,font=\small] {平动};
  \draw[<->,Ink,thick] (5.25,1.43)--(5.25,3.07) node[midway,right,font=\small] {转动};
  \draw[<->,Ink,thick] (9.9,3.23)--(9.9,4.87) node[midway,right,font=\small] {振动};
  \node[draw=DeepBlue,fill=PaleBlue,rounded corners=2pt,font=\small,align=left] at (5.55,5.15)
    {温度升高后，更高能级对应的自由度才逐步参与能量均分};
\end{tikzpicture}
\end{center}\nopagebreak[2]}

\newcommand{\pressureLayerRedraw}{%
\begin{center}
\begin{tikzpicture}[x=1cm,y=1cm,>=Latex]
  \shade[bottom color=DeepBlue!16,top color=DeepBlue!4] (-4.2,-1.55) rectangle (1.2,2.15);
  \draw[Ink,very thick] (-4.2,-1.55)--(1.2,-1.55);
  \draw[fill=DeepBlue!20,draw=DeepBlue,thick] (-3.7,0.2) rectangle (0.7,0.72);
  \draw[<->,Ink,thick] (-3.1,-1.48)--node[left] {$h$}(-3.1,0.18);
  \draw[<->,DeepBlue,thick] (1.65,0.2)--node[right] {$\dd h$}(1.65,0.72);
  \node[font=\small] at (-1.5,-0.55) {$T,\ P,\ \rho$};
  \node[font=\small] at (-1.5,1.15) {$T+\dd T,\ P+\dd P$};
  \draw[->,DeepBlue,very thick] (-0.2,0.45)--(-0.2,-0.7) node[midway,right] {$\rho gA\,\dd h$};
  \node[draw=DeepBlue,fill=PaleBlue,rounded corners=2pt,align=left,anchor=west] at (3.05,0.3)
    {薄层静力平衡：\\[3pt]$A\,\dd P=-\rho gA\,\dd h$\\[3pt]$\boxed{\dd P=-\rho g\,\dd h}$};
\end{tikzpicture}
\end{center}\nopagebreak[2]}

\newcommand{\atmosphereProfileRedraw}{%
\begin{center}
\begin{tikzpicture}[x=1cm,y=0.82cm,>=Latex]
  \fill[DeepBlue!10] (-5.2,-2.25) rectangle (5.6,-1.55);
  \fill[DeepBlue!6]  (-5.2,-1.55) rectangle (5.6,-0.25);
  \fill[orange!9]    (-5.2,-0.25) rectangle (5.6,2.0);
  \fill[DeepBlue!5]  (-5.2,2.0) rectangle (5.6,3.45);
  \fill[violet!7]    (-5.2,3.45) rectangle (5.6,5.15);
  \foreach \y/\h in {-2.25/0,-1.55/10,-0.25/20,2.0/50,3.45/80,5.15/100}{
    \draw[TextGray!65] (-5.2,\y)--(5.6,\y);
    \node[left,font=\scriptsize] at (-5.25,\y) {$\h$};
  }
  \draw[axis] (-5.2,-2.25)--(-5.2,5.55) node[above] {$h/\mathrm{km}$};
  \draw[axis] (-4.65,-2.25)--(0.8,-2.25) node[right] {$T$};
  \draw[DeepBlue,very thick]
    (-2.0,-2.25)--(-3.3,-1.55)--(-2.45,-0.25)--(-3.15,2.0)--(-1.65,3.45)--(-0.55,5.15);
  \foreach \x/\y in {-2.0/-2.25,-3.3/-1.55,-2.45/-0.25,-3.15/2.0,-1.65/3.45,-0.55/5.15}{\fill[DeepBlue] (\x,\y) circle (1.5pt);}
  \node[font=\scriptsize\bfseries,anchor=east] at (5.35,-1.88) {对流层};
  \node[font=\scriptsize\bfseries,anchor=east] at (5.35,-0.72) {平流层};
  \node[font=\scriptsize\bfseries,anchor=east] at (5.35,0.85) {中间层};
  \node[font=\scriptsize\bfseries,anchor=east] at (5.35,3.72) {热层};
  % Ground, mountains and clouds.
  \fill[Ink!70] (2.6,-2.25)--(3.45,-1.63)--(3.9,-1.98)--(4.45,-1.5)--(5.05,-2.25)--cycle;
  \foreach \x/\y in {1.15/-1.86,2.05/-1.78}{
    \fill[white,draw=TextGray] (\x,\y) circle (.18) (\x+.22,\y+.05) circle (.23) (\x+.45,\y) circle (.18);
  }
  % Airliner near the tropopause.
  \draw[Ink,very thick] (1.45,-1.02)--(2.65,-0.9);
  \draw[Ink,thick] (1.95,-.97)--(2.25,-.65) (1.95,-.97)--(2.28,-1.18);
  \node[font=\scriptsize,anchor=west] at (2.8,-0.9) {民航巡航};
  % Newly composed aurora ribbons.
  \draw[green!55!DeepBlue,very thick,decorate,decoration={snake,amplitude=1.1mm,segment length=5mm}]
    (1.25,4.48)--(4.15,4.68);
  \draw[violet!55!DeepBlue,thick,decorate,decoration={snake,amplitude=.9mm,segment length=5mm}]
    (1.45,4.05)--(3.9,4.22);
  \node[font=\scriptsize,anchor=west] at (4.25,4.72) {极光};
\end{tikzpicture}
\end{center}\nopagebreak[2]}

\newcommand{\cycleQuestionRedraw}{%
\begin{center}
\begin{tikzpicture}[x=1.05cm,y=.82cm,>=Latex]
  \draw[axis] (0,0)--(5.7,0) node[right] {$V$};
  \draw[axis] (0,0)--(0,4.8) node[above] {$P$};
  \draw[guide] (0,3.9)--(1.2,3.9);
  \draw[guide] (0,1.0)--(5.0,1.0);
  \node[left,font=\small] at (0,3.9) {$P_1$};
  \node[left,font=\small] at (0,1.0) {$P_1/4$};
  \draw[DeepBlue,very thick,->] (1.2,3.9)--(1.2,1.0);
  \draw[DeepBlue,very thick,->] (1.2,1.0)--(4.8,1.0);
  \draw[DeepBlue,very thick,postaction={decorate},decoration={markings,mark=at position .52 with {\arrow{Latex}}}]
    (4.8,1.0) .. controls (3.25,1.22) and (1.75,2.25) .. (1.2,3.9);
  \fill[Ink] (1.2,3.9) circle (2pt) node[above right] {$a$};
  \fill[Ink] (1.2,1.0) circle (2pt) node[below left] {$b$};
  \fill[Ink] (4.8,1.0) circle (2pt) node[below right] {$c$};
  \draw[guide] (1.2,0)--(1.2,3.9);
  \node[below,font=\small] at (1.2,0) {$V_1$};
  \node[font=\small,DeepBlue] at (3.25,2.2) {$c\to a$ 等温};
\end{tikzpicture}
\end{center}\nopagebreak[2]}

\newcommand{\straightQuestionRedraw}{%
\begin{center}
\begin{tikzpicture}[x=1cm,y=1cm,>=Latex]
  \draw[axis] (0,0)--(5.8,0) node[right] {$V$};
  \draw[axis] (0,0)--(0,4.7) node[above] {$P$};
  \draw[DeepBlue,very thick,->] (1.35,3.7)--(4.55,1.55);
  \fill[Ink] (1.35,3.7) circle (2pt) node[above left] {$A$};
  \fill[Ink] (4.55,1.55) circle (2pt) node[below right] {$B$};
  \node[draw=DeepBlue,fill=PaleBlue,rounded corners=2pt,font=\small,anchor=west] at (6.45,2.65)
    {$P_B/P_A=1/2$};
\end{tikzpicture}
\end{center}\nopagebreak[2]}

\newcommand{\liftingPistonQuestionRedraw}{%
\begin{center}
\begin{tikzpicture}[x=1cm,y=1cm,>=Latex]
  \draw[Ink,very thick] (-2.15,-1.9)--(-2.15,3.1) (2.15,-1.9)--(2.15,3.1) (-2.15,-1.9)--(2.15,-1.9);
  \draw[fill=SoftGray,draw=Ink,thick] (-2.0,0.15) rectangle (2.0,0.48);
  \fill[DeepBlue!10] (-2.0,-1.75) rectangle (2.0,0.15);
  \draw[->,DeepBlue,very thick] (0,0.55)--(0,1.8) node[right] {缓慢上提};
  \draw[Ink,very thick] (-2.45,-1.9)--(2.45,-1.9);
  \foreach \x in {-2.3,-1.7,...,2.3}{\draw[TextGray] (\x,-1.9)--++(-.25,-.25);}
  \node at (0,-0.7) {$p_0,T_0$};
  \node[font=\small,anchor=west] at (2.45,1.45) {绝热气缸质量 $M$};
  \node[font=\small,anchor=west] at (2.45,0.3) {无质量、无摩擦活塞};
  \node[draw=DeepBlue,fill=PaleBlue,rounded corners=2pt,font=\small,align=left,anchor=west] at (5.4,-0.7)
    {截面积 $A$\\外界压强 $p_0$\\刚离地时求 $T_f$};
\end{tikzpicture}
\end{center}\nopagebreak[2]}

% High-fidelity second pass: these drawings keep the source figure's object
% inventory and physical layout while using newly authored geometry.
\newcommand{\perpetualMotionFaithful}{%
\begin{center}
\begin{tikzpicture}[x=1cm,y=1cm,>=Latex]
  \foreach \cx/\lab in {-5.0/{装置 A},0/{装置 B},5.0/{装置 C}}{
    \draw[fill=SoftGray,draw=RuleGray,rounded corners=2pt] (\cx-2.1,-2.5) rectangle (\cx+2.1,2.75);
    \node[font=\small\bfseries] at (\cx,2.36) {\lab};
  }
  % A: rotor on a two-post stand, with curved ball channels.
  \begin{scope}[shift={(-5.0,0)}]
    \draw[Ink,very thick] (-1.55,-2.0)--(1.55,-2.0);
    \draw[Ink,very thick] (-0.72,-2.0)--(-0.38,-1.1) (0.72,-2.0)--(0.38,-1.1);
    \draw[fill=SoftGray,draw=Ink,very thick] (0,0.25) circle (1.35);
    \begin{scope}[shift={(0,.25)}]
    \foreach \a in {0,60,...,300}{
      \draw[DeepBlue,thick,rounded corners=5pt]
        (\a:0.30) .. controls (\a+18:0.72) and (\a+12:1.03) .. (\a:1.17);
      \fill[Ink] (\a:1.12) circle (2.2pt);
    }
    \end{scope}
    \fill[Ink] (0,0.25) circle (3pt);
    \draw[->,DeepBlue,very thick] (1.55,0.75) arc[start angle=20,end angle=-70,radius=1.05];
    \node[font=\scriptsize,align=center] at (0,-2.25) {带弯曲滚道的偏心转轮};
  \end{scope}
  % B: capsule-shaped endless carrier with blocks and a side guide.
  \begin{scope}[shift={(0,0)}]
    \draw[Ink,very thick,rounded corners=17pt] (-1.02,-1.62) rectangle (1.02,1.62);
    \fill[Ink] (0,1.05) circle (2.5pt) (0,-1.05) circle (2.5pt);
    \foreach \y in {-1.15,-.55,.05,.65,1.25}{
      \draw[fill=DeepBlue!10,draw=Ink,thick] (-1.28,\y-.17) rectangle (-.86,\y+.17);
      \draw[fill=DeepBlue!18,draw=Ink,thick] (.86,\y-.17) rectangle (1.28,\y+.17);
    }
    \draw[fill=TextGray!18,draw=Ink,thick]
      (1.45,-1.72)--(1.95,-1.72)--(1.95,1.35)--(1.68,1.58)--(1.43,1.23)--cycle;
    \draw[->,DeepBlue,very thick] (1.55,.65)--(1.55,-.35);
    \node[font=\scriptsize,align=center] at (0,-2.25) {沿封闭轨道循环的配重块};
  \end{scope}
  % C: perforated hub, radial rods and hanging end masses.
  \begin{scope}[shift={(5.0,0)}]
    \draw[Ink,very thick] (-1.55,-2.0)--(1.55,-2.0) (0,-2.0)--(0,-1.15);
    \draw[fill=SoftGray,draw=Ink,very thick] (0,0.2) circle (1.05);
    \foreach \a in {0,72,...,288}{
      \fill[white,draw=TextGray,thick] (\a:0.57) ++(0,0.2) circle (.22);
    }
    \begin{scope}[shift={(0,.2)}]
    \foreach \a in {-80,-50,...,250}{
      \draw[Ink,thick] (\a:0.95)--(\a:1.58);
      \fill[DeepBlue] (\a:1.58) circle (2.8pt);
    }
    \end{scope}
    \fill[Ink] (0,0.2) circle (3pt);
    \draw[->,DeepBlue,very thick] (-1.35,0.85) arc[start angle=155,end angle=245,radius=1.0];
    \node[font=\scriptsize,align=center] at (0,-2.25) {带径向摆锤的开孔转轮};
  \end{scope}
  \node[draw=DeepBlue,fill=PaleBlue,thick,rounded corners=2pt,font=\small,align=center]
    at (0,-3.08) {若装置完成一个循环并回到原状态，内部各部分的势能变化也必须一并结算};
\end{tikzpicture}
\end{center}\nopagebreak[2]}

\newcommand{\pistonWorkFaithful}{%
\begin{center}
\begin{tikzpicture}[x=1cm,y=1cm,>=Latex]
  \foreach \sx/\top/\gas/\name in {-6.0/.15/{V}/{初态},-2.9/.65/{V+\dd V}/{末态}}{
    \begin{scope}[shift={(\sx,0)}]
      \draw[Ink,very thick] (-1.05,-1.7)--(-1.05,2.05) (1.05,-1.7)--(1.05,2.05) (-1.05,-1.7)--(1.05,-1.7);
      \fill[DeepBlue!9] (-.95,-1.58) rectangle (.95,\top-.16);
      \draw[fill=SoftGray,draw=Ink,thick] (-.92,\top-.16) rectangle (.92,\top+.16);
      \draw[Ink,very thick] (0,\top+.16)--(0,2.45);
      \node at (0,-.7) {$\gas$};
      \node[font=\scriptsize] at (0,-2.05) {\name};
    \end{scope}
  }
  \draw[<->,DeepBlue,thick] (-1.45,.15)--node[right,font=\small] {$\dd y$}(-1.45,.65);
  \draw[->,DeepBlue,very thick] (-2.9,2.5)--(-2.9,3.08);
  \node[font=\small] at (-4.45,-2.55) {$\dd V=A\,\dd y$};
  \begin{scope}[shift={(1.0,-1.6)},x=1.05cm,y=.92cm]
    \draw[axis] (0,0)--(5.5,0) node[right] {$V$};
    \draw[axis] (0,0)--(0,4.85) node[above] {$P$};
    \fill[PaleBlue] (1.0,3.9) .. controls (2.0,2.55) and (3.6,1.35) .. (4.75,.9)--(4.75,0)--(1.0,0)--cycle;
    \draw[DeepBlue,very thick,postaction={decorate},decoration={markings,mark=at position .55 with {\arrow{Latex}}}]
      (.7,4.35) .. controls (1.85,2.45) and (3.45,1.25) .. (5.0,.82);
    \fill[Ink] (1.0,3.9) circle (2pt) node[above right] {$i$};
    \fill[Ink] (4.75,.9) circle (2pt) node[above right] {$f$};
    \draw[guide] (1.0,0)--(1.0,3.9) (4.75,0)--(4.75,.9);
    \node[font=\small,DeepBlue] at (2.8,.62) {$W^{\rm out}=\displaystyle\int_{V_i}^{V_f}P\,\dd V$};
  \end{scope}
\end{tikzpicture}
\end{center}\nopagebreak[2]}

\newcommand{\atmosphereProfileFaithful}{%
\begin{center}
\begin{tikzpicture}[x=1cm,y=.86cm,>=Latex]
  \fill[DeepBlue!9] (-5.8,0) rectangle (5.65,.62);
  \fill[DeepBlue!5] (-5.8,.62) rectangle (5.65,1.25);
  \fill[orange!8] (-5.8,1.25) rectangle (5.65,3.1);
  \fill[DeepBlue!5] (-5.8,3.1) rectangle (5.65,4.95);
  \fill[violet!8] (-5.8,4.95) rectangle (5.65,6.2);
  \foreach \y in {0,.62,1.25,3.1,4.95,6.2}{\draw[TextGray!55] (-5.8,\y)--(5.65,\y);}
  % Pressure scale retained from the source information.
  \node[rotate=90,font=\scriptsize\bfseries] at (-6.95,3.1) {压强（相对海平面）};
  \foreach \y/\p in {0/{1},.62/{10^{-1}},1.85/{10^{-2}},3.1/{10^{-3}},4.05/{10^{-4}},4.95/{10^{-5}},5.58/{10^{-6}},6.2/{10^{-7}}}{
    \node[left,font=\scriptsize] at (-5.92,\y) {$\p$};
  }
  % Central altitude scale, as in the historical infographic.
  \draw[axis] (.65,0)--(.65,6.55) node[above] {$h/\mathrm{km}$};
  \foreach \y/\h in {0/0,.62/10,1.25/20,3.1/50,4.95/80,5.58/90,6.2/100}{
    \draw[TextGray] (.57,\y)--(.73,\y);
    \node[right,font=\scriptsize] at (.78,\y) {$\h$};
  }
  % Newly plotted, source-faithful zig-zag temperature profile.
  \draw[DeepBlue,very thick]
    (-1.35,0)--(-3.25,.62)--(-3.05,1.25)--(-1.65,3.1)--(-3.75,4.95)--(-2.55,5.58)--(-1.75,6.2);
  \foreach \x/\y in {-1.35/0,-3.25/.62,-3.05/1.25,-1.65/3.1,-3.75/4.95,-2.55/5.58,-1.75/6.2}{\fill[DeepBlue] (\x,\y) circle (1.5pt);}
  \draw[axis] (-4.9,0)--(-.55,0) node[right] {$T$};
  \node[font=\scriptsize] at (-4.15,-.3) {低温};
  \node[font=\scriptsize] at (-1.15,-.3) {高温};
  % Numbered layers and concise heating explanations.
  \foreach \x/\y/\n in {-4.65/.32/4,-4.65/2.15/3,-4.65/3.95/2,-4.65/5.55/1}{
    \draw[Ink,thick,fill=white] (\x,\y) circle (.18);
    \node[font=\tiny] at (\x,\y) {\n};
  }
  \node[font=\scriptsize,align=left,anchor=west] at (1.65,2.5)
    {低层大气主要吸收\\温暖地表发出的辐射};
  \node[font=\scriptsize,align=left,anchor=west] at (1.65,5.2)
    {高层稀薄气体直接\\吸收太阳高能辐射};
  % Clouds, mountains, low aircraft and upper-space vehicle.
  \fill[Ink!65] (2.8,0)--(3.55,.5)--(4.0,.15)--(4.65,.72)--(5.3,0)--cycle;
  \foreach \x in {1.55,2.25}{\fill[white,draw=TextGray] (\x,.4) circle (.12) (\x+.16,.44) circle (.17) (\x+.32,.4) circle (.12);}
  \draw[Ink,thick] (3.0,.78)--(4.05,.88) (3.45,.82)--(3.72,1.08) (3.45,.82)--(3.78,.6);
  \draw[Ink,thick] (3.75,5.9)--(4.75,6.02) (4.15,5.95)--(4.42,6.22) (4.15,5.95)--(4.45,5.72);
  \draw[green!55!DeepBlue,very thick,decorate,decoration={snake,amplitude=1mm,segment length=5mm}]
    (1.65,5.72)--(3.45,5.9);
  \node[font=\scriptsize,anchor=west] at (1.65,5.98) {极光};
\end{tikzpicture}
\end{center}\nopagebreak[2]}

% Parameter-based technical figures, revised 2026-09-15.
\newcommand{\stellarInteriorRevision}{\begin{center}\begin{tikzpicture}[font=\scriptsize]
\foreach \sx/\outer/\inner/\core/\title in {0/1.7/1.19/.43/太阳,7/2.1/.68/.23/红巨星}{
\begin{scope}[shift={(\sx,0)}]
\fill[orange!18](0,0)circle(\outer);\fill[yellow!18](0,0)circle(\inner);\fill[orange!70](0,0)circle(\core);
\draw[Ink](0,0)circle(\outer);\draw[TextGray](0,0)circle(\inner);
\pgfmathsetmacro{\midradius}{(\outer+\inner)/2}
\pgfmathsetmacro{\cellradius}{(\outer-\inner)*.38}
\foreach \angle in {35,95,155,215,275,335}{
\draw[DeepBlue,thick,->] (\angle:\midradius)++(0:\cellradius)arc(0:305:\cellradius);}
\node at(0,2.6){\title};
\draw[TextGray](0,0)--(.25,-2.6);\node[below]at(.25,-2.6){核心};
\draw[TextGray](20:{(\core+\inner)/2})--(2.45,-1.2);\node[anchor=west]at(2.45,-1.2){辐射区};
\draw[TextGray](40:\midradius)--(2.45,1.25);\node[anchor=west]at(2.45,1.25){对流包层};
\end{scope}}\end{tikzpicture}
\par\smallskip\includegraphics[width=.63\linewidth]{figure20-3-redraw-v2.png}
\par\small 表面对流细节：太阳的细小颗粒与红巨星的较大对流胞
\end{center}\nopagebreak[2]}

\renewcommand{\atmosphereProfileFaithful}{\begin{center}\begin{tikzpicture}[x=1cm,y=.065cm,font=\scriptsize]
\foreach \lo/\hi/\col in {0/11/DeepBlue!8,11/50/orange!8,50/86/DeepBlue!5,86/100/violet!8}{\fill[\col](0,\lo)rectangle(11,\hi);}
\draw[axis](0,0)--(0,108)node[above]{$h/\mathrm{km}$};
\foreach \h in {0,20,40,60,80,100}{\draw(-.08,\h)--(.08,\h);\node[left]at(0,\h){$\h$};}
\foreach \h in {11,50,86}{\draw[guide](0,\h)--(11,\h);}
% Celsius axis: x=(t+100)/30. Altitudes and temperatures are representative.
\draw[axis](0,0)--(7.1,0)node[right]{$t/{}^\circ\mathrm C$};
\foreach \t in {-100,-50,0,50,100}{\draw({(\t+100)/30},1)--({(\t+100)/30},-1)node[below]{$\t$};}
\draw[DeepBlue,thick]plot coordinates {(3.833,0)(1.45,11)(1.45,20)(1.85,32)(3.25,47)(3.25,51)(1.383,71)(.46,86)(.46,91)(.731,100)};
\node[anchor=west]at(7.3,5.5){对流层};\node[anchor=west]at(7.3,31){平流层};
\node[anchor=west]at(7.3,68){中间层};\node[anchor=west]at(7.3,95){热层下部};
\node[anchor=west]at(11.4,109){$p/p_0$（约值）};
\foreach \h/\p in {0/1,20/{5.4\!\times\!10^{-2}},50/{7.9\!\times\!10^{-4}},80/{1.05\!\times\!10^{-5}},100/{3.2\!\times\!10^{-7}}}{\node[anchor=west]at(11.4,\h){$\p$};}
\end{tikzpicture}
\par\medskip\includegraphics[width=.94\linewidth]{atmosphere-scenes.png}
\par\small 云与山地\hfill 民航巡航\hfill 高空飞行器\hfill 极光
\par\smallskip\small 下方情景不按高度比例；上方折线为代表性分层模型，实际温度随地点、季节和太阳活动改变。
\end{center}\nopagebreak[2]}
\tikzset{cyclepath/.style={DeepBlue,thick,postaction={decorate},decoration={markings,mark=at position .55 with {\arrow{Latex}}}},revlabel/.style={font=\small,fill=none,inner sep=2pt}}
\renewcommand{\standardHeatCyclesDiagram}{%
\begin{center}\begin{tikzpicture}[x=1cm,y=.86cm,font=\small]
\foreach \s/\name in {0/卡诺循环,5/奥托循环,10/布雷顿循环}{
\begin{scope}[shift={(\s,0)}]
\draw[axis](0,0)--(4.1,0) node[right]{$V$};\draw[axis](0,0)--(0,4.65) node[above]{$P$};
\node at (1.8,5.1){\name};
\ifnum\s=0
% gamma=5/3, Th/Tc=(3/2)^(2/3), V1=.8,V2=2,V3=3,V4=1.2.
\pgfmathsetmacro{\cold}{3.2/pow(1.5,2/3)}
\draw[cyclepath,domain=.8:2,samples=65]plot(\x,{3.2/\x});
\draw[cyclepath,domain=2:3,samples=65]plot(\x,{3.2*pow(2,2/3)/pow(\x,5/3)});
\draw[cyclepath,domain=3:1.2,samples=65]plot(\x,{\cold/\x});
\draw[cyclepath,domain=1.2:.8,samples=65]plot(\x,{3.2*pow(.8,2/3)/pow(\x,5/3)});
\foreach \x/\y/\n/\a in {.8/4/1/left,2/1.6/2/above right,3/{\cold/3}/3/right,1.2/{\cold/1.2}/4/below left}{\fill({\x},{\y})circle(1.4pt)node[\a]{$\n$};}
\node at (2.65,3.15){$T_h$};\draw[thin,TextGray](2.4,2.95)--(1.45,{3.2/1.45});
\node at (2,-.65){$1\to2\to3\to4\to1$};
\else\ifnum\s=5
% Vmin=1,Vmax=3, gamma=1.4; lower adiabat K=1.8, upper K=4.
\draw[cyclepath,domain=3:1,samples=65]plot(\x,{1.8/pow(\x,1.4)});
\draw[cyclepath](1,1.8)--(1,4);
\draw[cyclepath,domain=1:3,samples=65]plot(\x,{4/pow(\x,1.4)});
\draw[cyclepath](3,{4/pow(3,1.4)})--(3,{1.8/pow(3,1.4)});
\foreach \x/\y/\n/\a in {3/{1.8/pow(3,1.4)}/1/below right,1/1.8/2/left,1/4/3/above,3/{4/pow(3,1.4)}/4/above right}{\fill({\x},{\y})circle(1.4pt)node[\a]{$\n$};}
\draw[guide](1,0)--(1,1.6);\node[below]at(1,0){$V_{\min}$};\node[below,yshift=-12pt]at(3,0){$V_{\max}$};
\else
% Pl=1,Ph=3,V2=.7,V3=1.5, gamma=1.4.
\pgfmathsetmacro{\ratio}{pow(3,1/1.4)}
\draw[cyclepath,domain={.7*\ratio}:.7,samples=65]plot(\x,{3*pow(.7/\x,1.4)});
\draw[cyclepath](.7,3)--(1.5,3);
\draw[cyclepath,domain=1.5:{1.5*\ratio},samples=65]plot(\x,{3*pow(1.5/\x,1.4)});
\draw[cyclepath]({1.5*\ratio},1)--({.7*\ratio},1);
\foreach \x/\y/\n/\a in {{.7*\ratio}/1/1/below left,.7/3/2/above left,1.5/3/3/above right,{1.5*\ratio}/1/4/right}{\fill({\x},{\y})circle(1.4pt)node[\a]{$\n$};}
\draw[guide](0,3)--(.55,3);\node[left]at(0,3){$P_h$};\draw[guide](0,1)--(1.4,1);\node[left]at(0,1){$P_l$};
\fi\fi\end{scope}}
\end{tikzpicture}\end{center}\nopagebreak[2]}

\renewcommand{\adiabaticCurvesDiagram}{\begin{center}\begin{tikzpicture}[x=1cm,y=.85cm,font=\small]
\draw[axis](0,0)--(6,0)node[right]{$V$};\draw[axis](0,0)--(0,5.1)node[above]{$P$};
\draw[cyclepath,domain=1:5.6,samples=90]plot(\x,{4.5/\x});
\draw[Ink,thick,domain=1:5.6,samples=90]plot(\x,{4.5/pow(\x,1.4)});
\fill(1,4.5)circle(1.5pt)node[above right]{共同初态};
\node[anchor=west]at(6.6,3.4){等温：$PV=\text{常量}$};
\node[anchor=west]at(6.6,2.2){绝热：$PV^\gamma=\text{常量}$};
\draw[TextGray](5.3,{4.5/5.3})--(6.3,3.4);
\draw[TextGray](5.3,{4.5/pow(5.3,1.4)})--(6.1,.35)--(6.3,2.2);
\end{tikzpicture}\end{center}\nopagebreak[2]}

\renewcommand{\convectionCriterionDiagram}{\begin{center}\begin{tikzpicture}[font=\small]
\draw[axis](0,0)--(0,4)node[above]{$h$};
\draw[DeepBlue,thick](2,.65)ellipse(.7 and .3);\draw[DeepBlue,thick](2,3.1)ellipse(.9 and .38);
\draw[->,DeepBlue,thick](2,1.15)--(2,2.55);\node[left]at(1.75,1.9){气团上升};
\node at(2,.65){$T$};\node at(2,3.1){$T+\dd T_1$};
\draw[<->](3.3,.65)--node[right]{$\dd h$}(3.3,3.1);
\node[anchor=west,align=left]at(5,2.85){气团：$\dd T_1=-\mu g\,\dd h/c_p$\\[6pt]环境：$\dd T_2=T'(h)\,\dd h$};
\node[anchor=west]at(5,1){继续上升须有 $\dd T_1>\dd T_2$};
\end{tikzpicture}\end{center}\nopagebreak[2]}

\renewcommand{\variableHeatCapacityDiagram}{\begin{center}\begin{tikzpicture}[font=\small]
\begin{scope}[x=4cm,y=3cm]
\draw[axis](0,0)--(1.2,0)node[right]{$V$};\draw[axis](0,0)--(0,1.2)node[above]{$P$};
\draw[DeepBlue,thick](0,1)--(1,0);\draw[guide](.5,0)--(.5,.5)--(0,.5);
\node[left]at(0,1){$P_0$};\node[below]at(1,0){$V_0$};\node[below]at(.5,0){$V_0/2$};
\node at(.65,.8){$P=P_0(1-V/V_0)$};\end{scope}
\begin{scope}[shift={(7,0)},x=4cm,y=3cm]
\draw[axis](0,0)--(1.2,0)node[right]{$V$};\draw[axis](0,0)--(0,1.2)node[above]{$T$};
\draw[DeepBlue,thick,domain=0:1,samples=70]plot(\x,{4*\x*(1-\x)});
\draw[guide](.5,0)--(.5,1);\fill(.5,1)circle(.4pt);
\node[above]at(.5,1){$T_{\max}=P_0V_0/(4R)$};\node[below]at(.5,0){$V_0/2$};\node[below]at(1,0){$V_0$};
\end{scope}\end{tikzpicture}\end{center}\nopagebreak[2]}

\newcommand{\capillaryRevisionDiagram}{\begin{center}\begin{tikzpicture}[font=\small]
\fill[DeepBlue!10](-2,-.7)rectangle(2,-1.8);
\fill[DeepBlue!10](-.5,-1.8)--(-.5,1.6)arc(180:360:.5)--(.5,-1.8)--cycle;
\draw[DeepBlue,thick](-2,-.7)--(-.5,-.7) (.5,-.7)--(2,-.7);
\draw[Ink,thick](-.5,-1.8)--(-.5,2.5) (.5,-1.8)--(.5,2.5);
\draw[DeepBlue,thick](-.5,1.6)arc(180:360:.5);
\draw[->,DeepBlue,thick](-.5,1.6)--(-.5,2.3);\draw[->,DeepBlue,thick](.5,1.6)--(.5,2.3);
\node[anchor=east]at(-.8,2.3){沿接触线向上};
\draw[<->](1,-.7)--node[right]{$h$}(1,1.1);\draw[guide](.15,1.1)--(1.2,1.1);
\draw[<->](0,-.2)--(.5,-.2);\node[above left]at(.25,-.16){$R$};
\node[anchor=west,align=left]at(3.2,1.4){完全润湿：$\theta=0$\\[5pt]接触处液面切线平行管壁\\[5pt]周界长度：$2\pi R$};
\node[anchor=west,align=left]at(3.2,-.7){$F_\sigma=2\pi R\sigma$\\[5pt]忽略弯月面体积修正，$R\ll h$};
\end{tikzpicture}\end{center}\nopagebreak[2]}

\newcommand{\bubbleRevisionDiagram}{\begin{center}\begin{tikzpicture}[font=\small]
\draw[DeepBlue,thick](0,0)circle(1.6);\draw[DeepBlue,thick](0,0)circle(1.5);
\draw[guide](0,0)circle(1.85);\fill(0,0)circle(1pt);
\draw[->](0,0)--node[above left]{$R$}(45:1.6);\node at(-.5,-.5){$p_R$};
\draw[->,thick](1.6,0)--(2.2,0)node[right]{$x>0$};\node at(-2.1,1.4){$p_0$};
\node[anchor=west,align=left]at(3.8,1){实线：内、外两个气液界面\\[5pt]虚线：扰动后半径 $r=R+x$\\[5pt]总面积 $A_s=8\pi r^2$};
\node[anchor=west]at(3.8,-1){液膜厚度仅为示意，不按比例};
\end{tikzpicture}\end{center}\nopagebreak[2]}

\renewcommand{\cycleQuestionRedraw}{\begin{center}\begin{tikzpicture}[x=1.1cm,y=.85cm,font=\small]
\draw[axis](0,0)--(5.1,0)node[right]{$V$};\draw[axis](0,0)--(0,4.8)node[above]{$P$};
\draw[guide](0,4)--(1,4)--(1,0) (0,1)--(4,1)--(4,0);
\draw[cyclepath](1,4)--(1,1);\draw[cyclepath](1,1)--(4,1);
\draw[cyclepath,domain=4:1,samples=95]plot(\x,{4/\x});
\node[left]at(0,4){$P_1$};\node[left]at(0,1){$P_1/4$};\node[below]at(1,0){$V_1$};\node[below]at(4,0){$4V_1$};
\foreach \x/\y/\n/\a in {1/4/a/above right,1/1/b/below left,4/1/c/right}{\fill(\x,\y)circle(1.5pt)node[\a]{$\n$};}
\node at(3.5,3){$PV=P_1V_1$};\draw[TextGray](3.1,2.8)--(2,2);
\end{tikzpicture}\end{center}\nopagebreak[2]}
\renewcommand{\cycleExampleDiagram}{\cycleQuestionRedraw}
\renewcommand{\straightQuestionRedraw}{\begin{center}\begin{tikzpicture}[font=\small]
\draw[axis](0,0)--(5.7,0)node[right]{$V$};\draw[axis](0,0)--(0,4.7)node[above]{$P$};
\draw[guide](0,4)--(1.4,4)--(1.4,0) (0,2)--(4.6,2)--(4.6,0);
\draw[cyclepath](1.4,4)--(4.6,2);\fill(1.4,4)circle(1.5pt)node[above]{$A$};\fill(4.6,2)circle(1.5pt)node[above right]{$B$};
\node[left]at(0,4){$P_A$};\node[left]at(0,2){$P_A/2$};\node[below]at(1.4,0){$V_A$};\node[below]at(4.6,0){$V_B$};
\end{tikzpicture}\end{center}\nopagebreak[2]}
\renewcommand{\straightProcessDiagram}{\straightQuestionRedraw}

\renewcommand{\liftingPistonQuestionRedraw}{\begin{center}\begin{tikzpicture}[font=\small]
\foreach \sx/\height/\lab/\gas in {0/1.4/初态/{p_0,T_0},5/3.8/最后密封的瞬间/{p_f,T_f}}{
\begin{scope}[shift={(\sx,0)}]
\fill[DeepBlue!9](-1,0)rectangle(1,\height);
\draw[thick](-1,3.8)--(-1,0)--(1,0)--(1,3.8);
\draw[fill=SoftGray,thick](-1,\height)rectangle(1,{\height+.18});
\draw[thick](0,{\height+.18})--(0,{\height+.85});
\ifnum\sx=0
\node[above]at(0,{\height+1.35}){$F=0$};
\else
\draw[->,DeepBlue,thick](0,{\height+.9})--(0,{\height+1.45});
\node[right]at(0,{\height+1.2}){$F$};
\fi
\node at(0,.7){$\gas$};
\draw[thick](-1.6,-.4)--(1.6,-.4);\draw[thick](-.8,0)--(-.8,-.4) (.8,0)--(.8,-.4);
\draw[->](0,-.32)--(0,0);\node[below]at(0,-.55){底部通大气};\node at(0,5.6){\lab};
\end{scope}}
\node[anchor=west,align=left]at(8,2.2){缸壁与活塞密封接触\\[5pt]缸体质量 $M$、截面积 $A$\\[5pt]终态活塞下缘位于开口};
\end{tikzpicture}\end{center}\nopagebreak[2]}
\renewcommand{\liftingPistonDiagram}{\begin{center}\begin{tikzpicture}[font=\small]
\node at(0,2.6){活塞};\draw[thick,fill=SoftGray](-1.2,-.1)rectangle(1.2,.1);
\draw[->,thick](-.6,.1)--(-.6,1.6)node[above]{$F$};\draw[->,thick](.6,-1.5)--(.6,-.1)node[midway,right]{$pA$};
\draw[->,thick](.6,1.6)--(.6,.1)node[midway,right]{$p_0A$};
\begin{scope}[shift={(6,0)}]\node at(0,2.6){缸体（刚离地，$N=0$）};
\draw[thick](-1,1.6)--(-1,0)--(1,0)--(1,1.6);
\draw[->,thick](-.5,-1.5)--(-.5,0)node[midway,left]{$p_0A$};
\draw[->,thick](.5,1.4)--(.5,0)node[midway,left]{$p_fA$};
\draw[->,thick](1,0)--(1,-1.5)node[right]{$Mg$};\end{scope}
\end{tikzpicture}\end{center}\nopagebreak[2]}

\renewcommand{\heatTransferMapDiagram}{\begin{center}\begin{tikzpicture}[font=\small]
\node[draw=DeepBlue,minimum width=2.4cm,minimum height=.8cm](sys)at(0,0){热力学系统};
\node[draw=RuleGray,align=center,minimum width=3.5cm,minimum height=1.1cm](cond)at(-5,1.8){直接接触（传导）\\需要介质};
\node[draw=RuleGray,align=center,minimum width=3.5cm,minimum height=1.1cm](conv)at(-5,-1.8){流体运动（对流）\\需要介质};
\node[draw=RuleGray,align=center,minimum width=3.5cm,minimum height=1.1cm](rad)at(5,1.8){电磁波（辐射）\\可穿过真空};
\node[draw=RuleGray,align=center,minimum width=3.5cm,minimum height=1.1cm](work)at(5,-1.8){边界位移（功）};
\draw[<->,DeepBlue,thick](cond.east)--node[above=4pt,sloped]{$Q$}(sys.north west);
\draw[<->,DeepBlue,thick](conv.east)--node[below=4pt,sloped]{$Q$}(sys.south west);
\draw[<->,DeepBlue,thick](rad.west)--node[above=4pt,sloped]{$Q$}(sys.north east);
\draw[<->,DeepBlue,thick](work.west)--node[below=4pt,sloped]{$W$}(sys.south east);
\end{tikzpicture}\end{center}\nopagebreak[2]}

\renewcommand{\isothermalAtmosphereDiagram}{\begin{center}\begin{tikzpicture}[font=\small]
\shade[bottom color=DeepBlue!22,top color=white](0,0)rectangle(4,4.8);
\draw[axis](-.4,0)--(-.4,5)node[above]{$h$};
\foreach \y/\alt/\lab in {0/0/{P_0,n_{N0}},2/H/{P_0/\e,n_{N0}/\e},4/{2H}/{P_0/\e^2,n_{N0}/\e^2}}{
\draw(-.5,\y)--(-.25,\y);\node[left]at(-.5,\y){$\alt$};\node[anchor=west]at(.3,{\y+.22}){$\lab$};}
\node[anchor=west,align=left]at(5.3,2.5){$P(h)=P_0\e^{-h/H}$\\[8pt]$n_N(h)=n_{N0}\e^{-h/H}$\\[8pt]$H=RT/(\mu g)$};
\end{tikzpicture}\end{center}\nopagebreak[2]}

\renewcommand{\heatCapacityRedraw}{\begin{center}\begin{tikzpicture}[x=3.2cm,y=1.05cm,font=\small]
\draw[axis](0,0)--(3.25,0)node[right]{$T/\mathrm K$};\draw[axis](0,0)--(0,4.3)node[above]{$c_v/R$};
\foreach \x/\t in {0/10,1/100,2/1000,3/10000}{\draw(\x,.06)--(\x,-.06)node[below]{$\t$};}
\foreach \y/\t in {1.5/{3/2},2.5/{5/2},3.5/{7/2}}{\draw[guide](0,\y)--(3.1,\y);\node[left]at(0,\y){$\t$};}
\draw[DeepBlue,thick,domain=0:3,samples=140]plot(\x,{1.5+1/(1+exp(min(8,max(-8,-15*(\x-.72)))))+1/(1+exp(min(8,max(-8,-15*(\x-2.25)))))});
\node at(.35,1){平动};\node at(1.5,2){再加转动};\node at(2.65,3){再加振动};
\node[anchor=west]at(.15,4.15){横轴为对数刻度；激发温区依分子种类而异};
\end{tikzpicture}\end{center}\nopagebreak[2]}

\renewcommand{\pistonWorkFaithful}{\begin{center}\begin{tikzpicture}[font=\small]
\foreach \sx/\top/\gas in {0/1.3/V,3/1.8/{V+\dd V}}{
\begin{scope}[shift={(\sx,0)}]\fill[DeepBlue!9](-.85,0)rectangle(.85,\top);
\draw[thick](-.85,2.4)--(-.85,0)--(.85,0)--(.85,2.4);
\draw[thick,fill=SoftGray](-.85,\top)rectangle(.85,{\top+.18});
\draw[thick](0,{\top+.18})--(0,{\top+1.25});\node at(0,.6){$\gas$};\end{scope}}
\draw[guide](3.85,1.3)--(4.25,1.3);\draw[<->](4.2,1.3)--node[right]{$\dd y$}(4.2,1.8);
\node at(1.5,-.6){$\dd V=A\dd y$};
\begin{scope}[shift={(6,0)},x=1.1cm,y=.8cm]
\draw[axis](0,0)--(5.3,0)node[right]{$V$};\draw[axis](0,0)--(0,4.9)node[above]{$P$};
\fill[PaleBlue](1,0)--plot[domain=1:4.5,samples=100](\x,{4/\x})--(4.5,0)--cycle;
\draw[cyclepath,domain=.9:4.8,samples=100]plot(\x,{4/\x});
\foreach \x/\s in {1/i,4.5/f}{\draw[guide](\x,0)--(\x,{4/\x});\fill(\x,{4/\x})circle(1.3pt)node[above right]{$\s$};\node[below]at(\x,0){$V_\s$};}
\node at(2.5,-1.1){$W^{\rm out}=\int_{V_i}^{V_f}P\,\dd V$};\end{scope}
\end{tikzpicture}\end{center}\nopagebreak[2]}


在热力学第一定律建立以前，人们曾反复尝试制造一种不需要任何能量供应、却能持续对外做功的机器，
也就是所谓的第一类永动机。所有这类尝试的失败，使人们逐渐认识到：必定存在一条普遍规律，使这种
机器从根本上不可能实现。今天我们知道，这条规律就是能量的转化和守恒定律。它是十九世纪科学发展中
最重要的成就之一；热力学第一定律则是能量守恒在包含热现象的系统中的定量陈述。

对于一个既不与外界交换物质、也不交换能量的孤立系统，总能量始终保持不变；若系统与环境之间存在
物质或能量交换，那么特定过程前后系统总能量的改变量，必然等于净输入系统的总能量。这里的总能量
包括内能、整体动能及整体势能；流入、流出的物质也能携带能量，不能只计算热量和边界功。
这一定律把一种新的宏观能量形式——内能——同热量、功以及其他可能的能量交换方式联系起来。

\par\noindent\begin{minipage}{\linewidth}
\perpetualMotionFaithful
\diagramcaption{三类第一类永动机设想。转轮、循环配重和径向摆锤的结构不同，却都企图在没有净能量输入时持续输出功。}\label{ch15:fig:s29-20-1}
\end{minipage}\par\medskip

图~\ref{ch15:fig:s29-20-1}中的装置并不是因为“某个齿轮比例没有调好”才失败。只要装置完成一个循环并回到原状态，内能
的净变化就是零；若又没有从外界吸收能量，便不可能有正的净功输出。后文将把这一判断写成可计算的
方程。

\section{内能}

热力学中的内能（internal energy）定义为包含在给定系统中的总能量，扣除系统质心整体运动带来的
动能，以及系统在外场中的整体势能。内能由系统内部微观粒子的运动状态和相互作用决定。在最简单的
经典图像中，它可以理解为微观粒子热运动的动能与粒子间相互作用势能之和。

一般情形下，系统内能的绝对值很难直接确定，因为分子之间存在复杂的相互作用。幸运的是，在经典物理
中，真正进入可观测规律的通常是内能的改变 $\Delta U$。这与海面航行的类比相似：许多问题并不需要
知道海水总量或海底的绝对深度，只需要知道海面高度在过程中的变化。内能的转化过程仍由底层的动力学
规律支配，但热力学把大量微观自由度压缩为少数宏观状态量。

只有在某些简单模型中，内能可以直接写出。例如单原子理想气体：构成气体的原子在所研究的过程中内部
结构不变，气体处于绝对温度 $T$ 的热平衡状态，分子数为 $N$，物质的量为 $n=N/N_A$。理想气体忽略
分子间相互作用，因而不存在分子间势能；每个单原子分子有三个平动自由度。由能量均分定理，每个分子
平均具有 $3k_BT/2$ 的平动动能，所以全部分子的内能为
\begin{equation}
  \boxed{U=\frac32Nk_BT=\frac32nRT}. \label{ch15:eq:s29-20-1}
\end{equation}
由此可见，在热平衡下，一定量单原子理想气体的内能只与温度有关。

以氧气、氮气为代表的双原子气体稍复杂。除质心的三个平动自由度外，分子还可以转动，也可能沿两原子
连线方向振动。若不加判断地使用经典能量均分定理，会得到高温极限下的 $U=7nRT/2$；实际测量却显示，
常温下振动自由度并没有充分激发，气体的行为更像只有平动和转动自由度。这个问题长期困扰经典物理，
直到量子理论建立后，人们才理解：相应能级间隔若远大于 $k_BT$，该自由度会被冻结；温度升高后，
振动贡献才逐渐显现。

对于不能忽略分子间势能的物质，内能通常不仅与温度有关，也会与体积有关，因为体积改变会改变分子
间距和相互作用势能。尽管内能的具体表达式可能很复杂，但有一点极其重要：当系统的宏观平衡状态完全
给定时，其内能也随之确定。内能是状态函数；系统以何种过程来到同一个状态，不会改变该状态的内能。

\section{热交换：传导、对流与辐射}

当系统与外界环境存在温度差时，就具备发生热量交换的条件。约定系统从状态1到状态2的过程中，从外界
净输入热量记作 $Q_{1\to2}$。若 $Q>0$，能量以热的形式净流入系统；若
$Q<0$，系统净放热。后文第一定律中的 $Q^{\rm in}$ 也表示这一带符号的净热量，即 $Q^{\rm in}=Q$。
计算循环效率时，另把各吸热段的热量总和记为正量 $Q_h$，各放热段排出的热量总和记为正量 $Q_c$，则
\begin{equation*}
  Q_{\rm cyc}=Q_h-Q_c.
\end{equation*}
这两个正的累计热量一般不相等，也不是互为负数。
热量像功一样，是过程量而不是状态量：不能说“某个状态含有多少热量”，只能说某一过程中传递了多少
热量。

温度差只创造了热交换的条件，能量真正传递还需要具体机制。常见机制包括直接接触传导、流体对流和
热辐射。两个温度不同的固体直接接触，或通过某种介质接触时，热端粒子的平均能量较高；相互作用使
能量在大量碰撞中向冷端传递。宏观上热端降温、冷端升温，这就是热传导。在最简单的近似下，热流率
与温差成正比，比例系数还取决于材料和几何尺寸，所以不同介质的传热能力差别显著。

\par\noindent\begin{minipage}{\linewidth}
\redrawimage[0.62\linewidth]{figure20-2-redraw-v3.png}
\diagramcaption{手掌附近的对流羽流示意。空气受热、密度减小并上升；颜色表示相对冷热，不代表实测温度。}\label{ch15:fig:s29-20-2}
\end{minipage}\par\medskip
图~\ref{ch15:fig:s29-20-2}中，手掌是局部热源；上方羽流逐渐扩展，表明热空气上升时还会卷入周围空气。

气体只靠分子扩散传导热量的能力通常不强；当介质是重力场中的流体，并在竖直方向存在足够大的温度梯度
时，会出现效率更高的对流。以地表附近的空气为例，随高度增加，大气压强和密度逐渐减小。若低处一团
较热气体受到扰动向上运动，到达高处后密度仍小于周围气体，它将继续受到向上的浮力，并把所携带的
热量送往高处。对流一旦形成，流体的整体运动会迅速交换各区域的能量，趋向减小温差。

\par\noindent\begin{minipage}{\linewidth}
\stellarInteriorRevision
\diagramcaption{太阳与红巨星的内部剖面及表面对流细节。箭头表示对流区物质的循环运动；两星大小与内部层厚不按同一比例。}\label{ch15:fig:s29-20-3}
\end{minipage}\par\medskip
图~\ref{ch15:fig:s29-20-3}左侧太阳的核心产生能量，能量先穿过辐射区，再经外层对流区向外传递；右侧红巨星具有扩展的
对流包层和相对紧凑的内部区域。下方局部图中，亮区对应较热物质上涌，暗边界对应较冷物质下沉。
这说明对流不仅发生在地球空气中，也能成为恒星内部的重要输能方式；不能把表面纹理直接当作内部结构。

传导和对流都需要物质介质，热辐射则可以穿过真空。太阳与地球之间的能量交换就是典型例子。所有温度
高于绝对零度的物体都会发出电磁辐射，辐射带走能量；温度越高，单位面积的总辐射功率越大。不同物质
的发射能力不同。吸收能力越强的表面，在同一温度下通常也具有更强的发射能力。能够吸收全部入射辐射
的理想物体称为黑体（black body）。黑体温度为 $T$ 时，单位表面积的总辐射功率为
\begin{equation}
  \boxed{j_{\rm rad}=\sigma T^4}, \label{ch15:eq:s29-20-2}
\end{equation}
其中 $\sigma\approx\SI{5.670374419e-8}{W.m^{-2}.K^{-4}}$ 是斯蒂芬--玻尔兹曼常数。四次方关系
说明，温度升高时热辐射能力会急剧增强。

\par\noindent\begin{minipage}{\linewidth}
\redrawimage[0.72\linewidth]{figure20-4-redraw-v2.png}
\diagramcaption{手部热辐射强弱示意。掌心、指尖、手腕和边缘的温度不同；图中颜色不是实验测量值。}\label{ch15:fig:s29-20-4}
\end{minipage}\par\medskip
图~\ref{ch15:fig:s29-20-4}的各部位即使没有接触远处探测器，也能通过电磁辐射交换能量。

\par\noindent\begin{minipage}{\linewidth}
\heatTransferMapDiagram
\figcap{热传递与边界功的路径。热辐射不需要物质介质，功对应边界受力和位移。}\label{ch15:fig:heat-transfer-map}
\end{minipage}\par\medskip
图\ref{ch15:fig:heat-transfer-map}中的 $Q$ 与 $W$ 都跨越系统边界，却对应不同的能量交换机制。

上述三种热交换都由温度差驱动，但路径不同：传导靠微观碰撞逐步传能，对流还伴随流体团块的宏观位移，
辐射则以电磁波携带能量。判断题目中的热传递方式时，不能只看“有没有接触”，还要看是否发生整体流动、
是否允许真空传播，以及过程所经历的时间尺度。

\subsection{从温差到热流率：串联热阻}
热量 $Q$ 的单位是焦耳，热流率 $\dot Q$ 的单位是瓦。对于温度只沿 $x$ 方向变化的均匀平板，傅里叶定律为
\[
 \dot Q=-\lambda A\frac{\dd T}{\dd x},
\]
其中 $\lambda$ 是热导率，单位为 $\mathrm{W/(m\,K)}$，$A$ 是垂直热流方向的面积。负号表示能量自发流向温度降低的方向。无内部热源的稳态中，各截面的 $\dot Q$ 相同；若 $\lambda$ 为常数，对厚度 $L$ 积分得
\[
 T_h-T_c=\dot Q\frac{L}{\lambda A}=\dot Q\,\mathcal R_{\rm th}.
\]
定义的热阻 $\mathcal R_{\rm th}=L/(\lambda A)$ 以 $\mathrm{K/W}$ 为单位。串联层的温差相加、热流率相同，因此热阻相加；并联通道则温差相同、热流率相加。这里借用的是相同的线性关系，不是说热流就是电流。图\ref{ch15:fig:heat-transfer-map}中传导路径若有多种介质，就要沿该路径逐段计算。

\Needspace{180pt}
\begin{origexamplebox}{复合平板中的界面温度}\label{ch15:ex:teaching-conduction}
面积为 $A$ 的两块平板紧贴串联，厚度为 $L_1,L_2$，热导率为 $\lambda_1,\lambda_2$。第一块外侧恒为 $T_h$，第二块外侧恒为 $T_c<T_h$。忽略接触热阻、侧面散热和内部热源，求稳态热流率及两板接触面的温度 $T_i$。
\end{origexamplebox}
\Needspace{65pt}
\begin{xesolution}
取由热端指向冷端为正，两块板满足
\[
 \dot Q=\frac{T_h-T_i}{\mathcal R_1}
       =\frac{T_i-T_c}{\mathcal R_2},\qquad
 \mathcal R_j=\frac{L_j}{\lambda_jA}.
\]
两段温降之和为总温降，所以
\[
 \dot Q=\frac{T_h-T_c}{\mathcal R_1+\mathcal R_2},\qquad
 T_i=\frac{\mathcal R_2T_h+\mathcal R_1T_c}{\mathcal R_1+\mathcal R_2}.
\]
较大热阻承担较大温降。例如 $\mathcal R_1\gg\mathcal R_2$ 时，第一层几乎承担全部温差，$T_i$ 接近冷端而不是热端。两段热流率相等保证界面没有能量积累；若初始温度还在变化，就不能仅凭这一稳态公式求整个升温过程。热量本身还要乘经过的时间：稳态持续 $\Delta t$ 时，$Q=\dot Q\Delta t$。
\end{xesolution}

\subsection{辐射发出功率与净散热功率}
物体在发射的同时也吸收环境辐射。若灰体完全被温度均匀的巨大环境包围，相关温区内发射率与吸收率可用同一常数 $\epsilon$ 近似，且忽略其他换热，则净辐射散热率为
\[
 P_{\rm net}=\epsilon\sigma A(T^4-T_0^4).
\]
不是 $\epsilon\sigma AT^4$。当 $T=T_0$ 时，发射并未停止，只是发出与吸收相等。
这与图\ref{ch15:fig:s29-20-4}中各部位仍持续发射辐射的情景相容。若视场内有多个温度不同的物体，就不能把它们任意替换为一个算术平均温度。

\Needspace{180pt}
\begin{origexamplebox}{电加热物体的辐射平衡与回稳时间}\label{ch15:ex:teaching-radiation}
热容量为常数 $C$ 的物体内部温度均匀、表面积为 $A$，被温度 $T_0$ 的巨大环境包围。恒定电功率 $P$ 全部转化为物体内能，只通过上述灰体辐射模型交换热量。求平衡温度，并求偏离平衡很小时温差的衰减时间常数。
\end{origexamplebox}
\Needspace{65pt}
\begin{xesolution}
以物体为系统，功率形式的第一定律为
\[
 C\frac{\dd T}{\dd t}=P-\epsilon\sigma A(T^4-T_0^4).
\]
平衡时左右能流相等，故
\[
 T_e=\left(T_0^4+\frac{P}{\epsilon\sigma A}\right)^{1/4}.
\]
令 $T=T_e+\vartheta$、$|\vartheta|\ll T_e$，展开 $T^4=T_e^4+4T_e^3\vartheta+O(\vartheta^2)$，平衡项抵消，得到
\[
 \frac{\dd\vartheta}{\dd t}=-\frac{\vartheta}{\tau},\qquad
 \vartheta(t)=\vartheta(0)e^{-t/\tau},\qquad
 \tau=\frac{C}{4\epsilon\sigma AT_e^3}.
\]
温度偏高时散热大于供热，偏低时反之，所以平衡稳定。分母是热流率对温度的变化率，单位为 $\mathrm{W/K}$，故 $\tau$ 的单位是秒。该指数衰减只适用于小偏离；很大温差时须保留四次方的非线性项。
\end{xesolution}

\section{机械功与准静态过程}

当系统与环境之间存在压强差，并且系统边界允许移动时，就会发生另一类能量交换——机械功。系统膨胀
时可对环境做功；系统被压缩时，环境对系统做功。

设圆柱形气缸横截面积为 $A$，封闭气体压强为 $P$。气体作用在活塞上的力为 $F=PA$。活塞向外移动
极小距离 $\dd y$ 时，体积增量 $\dd V=A\dd y$，气体对外所做的元功为
\begin{equation*}
  \dd W^{\rm out}=F\dd y=P\,\dd V.
\end{equation*}
这一结果与容器的具体形状无关。若可逆过程中压强随体积按 $P(V)$ 改变，从 $V_i$ 变到 $V_f$ 的对外功
就是 $P$--$V$ 图像下的有向面积：
\begin{equation}
  \boxed{W^{\rm out}_{i\to f}=\int_{V_i}^{V_f}P(V)\,\dd V}. \label{ch15:eq:s29-20-3}
\end{equation}

\par\noindent\begin{minipage}{\linewidth}
\pistonWorkFaithful
\diagramcaption{活塞发生微小位移时，气体体积增加 $\dd V=A\dd y$；同一过程在 $P$--$V$ 图中曲线下的面积就是气体对外功。}\label{ch15:fig:s29-20-5}
\end{minipage}\par\medskip
图~\ref{ch15:fig:s29-20-5}右图的 $i,f$ 与横轴上的 $V_i,V_f$ 一一对应，阴影上边界就是所积分的同一条压强曲线。

若改用环境对系统所做的功，符号相反：$\dd W^{\rm in}=-P\dd V$。但把边界处的功直接写成
$P\dd V$ 隐含一个条件：外界压强与系统压强在全过程中几乎相等，只保留足以推动边界移动的无穷小差别。
因此过程必须足够缓慢，使系统在每一时刻都接近平衡态；这样的理想过程称为准静态过程。若再忽略摩擦、
黏滞、有限温差传热等耗散，准静态过程可进一步视为可逆过程。

真实压缩或膨胀总会偏离这一理想极限，但竞赛题若给出“缓慢”“光滑”“无摩擦”等条件，通常正是在
授权使用准静态、可逆近似。非准静态过程仍满足能量守恒，只是系统内部可能没有统一压强，不能直接用
系统平衡压强代入式~(\ref{ch15:eq:s29-20-3})计算边界功。

\section{热力学第一定律}

热交换和机械功是系统与环境交换能量的两种最常见方式。本章主要讨论不交换物质的封闭系统。
一般应写 $\Delta(U+E_k+E_p)=Q+W^{\rm in}$。若系统从状态1经任意过程到达状态2，
整体动能和外场势能的变化可忽略，只考虑热量和功两种交换形式，取“流入系统为正”，则可简化为
\begin{equation}
  \boxed{U_2-U_1=Q^{\rm in}_{1\to2}+W^{\rm in}_{1\to2}}. \label{ch15:eq:s29-20-4}
\end{equation}
若改用气体对外做功 $W^{\rm out}=-W^{\rm in}$，同一规律等价地写成
\begin{equation*}
  \boxed{Q^{\rm in}=\Delta U+W^{\rm out}}.
\end{equation*}
后一个形式在气体过程题中更常用：吸收的热量一部分增加内能，另一部分转化为气体对外功。

热力学第一定律是能量守恒的直接体现，因此对可逆、不可逆过程都成立。可逆性的作用不是决定能量是否
守恒，而是决定能否仅凭一条平衡态路径计算功和热量。内能 $U$ 是状态函数，$\Delta U$ 只取决于始末
状态；热量 $Q$ 与功 $W$ 是过程量，同样的两个状态之间，沿不同路径的数值可以不同。

若系统经历一个循环又回到初态，则 $\Delta U_{\rm cyc}=0$。因此
\begin{equation*}
  0=W^{\rm in}_{\rm cyc}+Q^{\rm in}_{\rm cyc},
  \qquad
  \boxed{W^{\rm out}_{\rm cyc}=Q^{\rm in}_{\rm cyc}}.
\end{equation*}
整个循环对外的净功必等于同一循环从外界吸收的净热量。若循环不从外部吸收任何能量，净功只能为零；
任何企图无能量输入而持续对外做功的第一类永动机都不可能存在。

\begin{warningbox}{先统一功的正方向}
有的教材把“气体对外做功”记为正，有的教材把“外界对气体做功”记为正。两种约定都可以，但不能在
同一推导中混用。本讲解答统一采用 $Q>0$ 表示气体吸热，$W^{\rm out}>0$ 表示气体对外做功，故
$Q=\Delta U+W^{\rm out}$。
\end{warningbox}

\section{理想气体的等容、等压与等温过程}

热力学第一定律配合理想气体状态方程，可以进一步研究具体过程。理想气体的内能只与温度有关；在温区
不太宽、摩尔热容可近似为常量时，$n$ 摩尔气体有
\begin{equation*}
  \dd U=n c_v\,\dd T,
\end{equation*}
其中 $c_v$ 是定容摩尔热容。微分式也适用于随温度变化的 $c_v(T)$；但有限温差时一般应算
\[
 \Delta U=n\int_{T_1}^{T_2}c_v(T)\,\dd T.
\]
只有在该温区内 $c_v$ 近似不变，才可把它提出积分得到 $n c_v\Delta T$。

理想气体经历等容过程时，$\dd V=0$，因而不发生体积功。吸收的热量全部表现为内能增加：
\begin{equation*}
  \boxed{Q_V=n c_v\Delta T}.
\end{equation*}
单原子理想气体在经典适用范围内 $c_v=3R/2$。双原子气体的热容随温度改变：低温时部分转动自由度
冻结，常温时主要有平动和转动贡献，高温时振动自由度逐渐激发。

\par\noindent\begin{minipage}{\linewidth}
\heatCapacityRedraw
\diagramcaption{双原子气体热容的自由度激发示意。三个平台为 $3R/2,5R/2,7R/2$，横轴各十倍温区等宽；不计解离、电子激发和相变。}\label{ch15:fig:s29-20-6}
\end{minipage}\par\medskip
图~\ref{ch15:fig:s29-20-6}中每增加一个平台，表示又有一组自由度开始贡献热容，不是气体发生了不连续的温度跳变。

等压过程比等容过程多了膨胀功。温度增加 $\Delta T$ 时，理想气体状态方程给出
$P\Delta V=nR\Delta T$，于是
\begin{equation*}
  Q_P=n c_v\Delta T+P\Delta V
  =n(c_v+R)\Delta T=n c_p\Delta T.
\end{equation*}
因此理想气体满足迈耶关系
\begin{equation*}
  \boxed{c_p=c_v+R}.
\end{equation*}
$c_p$ 是一摩尔气体在等压过程中升高一开尔文所需的热量。它比 $c_v$ 大一个 $R$，正是因为等压升温
时还要为膨胀功提供能量。

若气体全过程同恒温热库接触并始终保持温度 $T$，便经历等温过程。对理想气体，$U$ 只与 $T$ 有关，
所以等温过程 $\Delta U=0$。气体等温膨胀时吸收的热量完全转化为对外功。$n$ 摩尔理想气体从 $V_1$
等温膨胀到 $V_2$ 时
\begin{equation}
  \boxed{Q^{\rm in}=W^{\rm out}
  =\int_{V_1}^{V_2}P\,\dd V
  =nRT\ln\!\frac{V_2}{V_1}}. \label{ch15:eq:s29-20-5}
\end{equation}
理想气体具有“等温即内能不变”的简单性质。真实物质的内能通常还依赖体积，不能把这条结论直接
推广到液体、固体或高密度真实气体。

\par\noindent\begin{minipage}{\linewidth}
\processComparisonDiagram
\figcap{四类常用气体过程的约束。等压热量式和绝热幂律分别使用常热容与可逆绝热近似。}\label{ch15:fig:process-constraints}
\end{minipage}\par\medskip
图\ref{ch15:fig:process-constraints}应按已知条件逐项选用；例如 $Q=0$ 本身并不足以推出 $PV^\gamma$ 恒定。

\subsection{同一对初末态，不一定有同样的功和热}

以内能只依赖温度的理想气体为例，设初态为 $(T_0,V_1)$，末态为 $(T_0,V_2)$，且 $V_2>V_1$。
若气体接触温度为 $T_0$ 的热库，作可逆等温膨胀，则
\[
 \Delta U=0,\qquad W^{\rm out}=nRT_0\ln\frac{V_2}{V_1},\qquad Q=W^{\rm out}>0.
\]
若改在绝热刚性容器中撤去隔板，让气体向真空自由膨胀到同一体积，则外压为零，
整个容器外壁也不移动，所以 $W^{\rm out}=0$；绝热给出 $Q=0$，于是仍有 $\Delta U=0$。
理想气体的平衡末温仍为 $T_0$，两条路径连接相同的初末态，但功、热不同。
自由膨胀的中间过程并非一串均匀平衡态，不能把平衡物态方程 $p=nRT_0/V$
当作边界上的实际外压去积分。写体积功之前必须说明使用哪一侧的压强，以及边界是否准静态。

同理，$Q=0$ 不能单独推出 $TV^{\gamma-1}$ 恒定；刚才的自由膨胀就是反例。
内能的改变量由状态决定，热和功则由实际能量交换的方式决定，这是后续所有过程题选式的依据。

\section{绝热过程}

温差出现后，真正传递有限热量需要一定时间。现实中的快速过程虽然可能存在明显温差，但在过程持续的
短时间内，热交换相对于功和内能改变很小，可以近似为绝热过程。绝热过程定义为系统与环境之间没有
热量交换，即 $Q=0$。若没有其他能量交换，绝热膨胀对外做的功完全来自内能降低；绝热压缩时，外界
做功则转化为内能增加。
快速并不是绝热的定义：隔热良好的系统也可以缓慢地发生绝热变化。若要同时使用“绝热”和“准静态”
近似，过程时间应长于系统内部恢复力学平衡的时间，又短于显著漏热的时间。

考虑 $n$ 摩尔理想气体从平衡状态 $(P,V,T)$ 作无穷小可逆绝热变化。由
\begin{equation*}
  T=\frac{PV}{nR},
  \qquad
  \dd T=\frac{P\dd V+V\dd P}{nR},
\end{equation*}
以及第一定律 $\dd U+P\dd V=n c_v\dd T+P\dd V=0$，先代入温度微分：
\[
 \frac{c_v}{R}(P\dd V+V\dd P)+P\dd V=0,
 \qquad (c_v+R)P\dd V+c_v V\dd P=0.
\]
利用 $c_p=c_v+R$，再除以 $c_vPV$，得到
\begin{equation}
  \boxed{\gamma\frac{\dd V}{V}+\frac{\dd P}{P}=0},
  \qquad
  \boxed{\gamma=\frac{c_p}{c_v}}. \label{ch15:eq:s29-20-6}
\end{equation}
$\gamma$ 称为绝热指数。以下幂律公式还要求 $c_v,c_p$ 在所研究温区内为常数。对微分式从初态积分至末态：
\[
 \int_{P_1}^{P_2}\frac{\dd P}{P}+\gamma\int_{V_1}^{V_2}\frac{\dd V}{V}=0,
 \qquad \ln\frac{P_2}{P_1}+\gamma\ln\frac{V_2}{V_1}=0.
\]
把对数相加并取指数，得 $P_2V_2^\gamma=P_1V_1^\gamma$，即泊松（Poisson）绝热方程
\begin{equation}
  \boxed{PV^\gamma=\mathrm{const}}. \label{ch15:eq:s29-20-7}
\end{equation}

利用 $P=nRT/V$，有 $PV^\gamma=nRTV^{\gamma-1}$；定量气体的 $nR$ 不变，故
\begin{equation}
  \boxed{TV^{\gamma-1}=\mathrm{const}}. \label{ch15:eq:s29-20-8}
\end{equation}
这一形式适合研究封闭气体在热机中的体积、温度变化。若代入 $V=nRT/P$，则
$PV^\gamma=(nR)^\gamma T^\gamma/P^{\gamma-1}$，所以
\begin{equation}
  \boxed{\frac{T^\gamma}{P^{\gamma-1}}=\mathrm{const}}. \label{ch15:eq:s29-20-9}
\end{equation}
压强和温度都是强度量，式~(\ref{ch15:eq:s29-20-9})不显含气体总量，适合跟踪开放环境中不与外界混合的
同一小气团，并不意味着可随意用于有质量流入、流出的系统。再以密度
$\rho$ 代替体积，可写成
\begin{equation}
  \boxed{P=\mathrm{const}\cdot\rho^\gamma}. \label{ch15:eq:s29-20-10}
\end{equation}
压强与密度的关系是描述气体物态及其变化的直接工具。
这里每条绝热线的常数取决于气体状态。对固定物质的量、常热容理想气体，下一章的熵公式给出
\[
 S_2-S_1=n c_v\ln\frac{P_2V_2^\gamma}{P_1V_1^\gamma}.
\]
所以同一条可逆绝热线也是等熵线；这个常数不是焓。若热容明显随温度变化，须改用
$\int_{T_1}^{T_2}c_v(T)\dd T/T+R\ln(V_2/V_1)=0$，不能套固定指数的幂律。

\par\noindent\begin{minipage}{\linewidth}
\adiabaticCurvesDiagram
\figcap{从同一初态出发，理想气体等温线与可逆绝热线的比较。}\label{ch15:fig:adiabatic-comparison}
\end{minipage}\par\medskip
图\ref{ch15:fig:adiabatic-comparison}中，在任一相同的末体积处，下支的压强较低，对应膨胀降温的绝热过程。

在相同初态出发的膨胀过程中，绝热线比等温线下降得更陡。等温膨胀可以从热库持续吸热来补偿对外功；
绝热膨胀没有热量补充，只能消耗内能，温度随体积增大而下降。这也是水汽和空气团快速上升时降温的
基本原因。

\subsection{给定过程曲线时，先由约束求温度变化}
常见等压、等温与可逆绝热过程不是全部可能路径。若定量理想气体按准静态路径 $pV^{\kappa}=K$ 变化，$\kappa$ 为给定过程指数，则由 $pV=nRT$ 得
\[
 TV^{\kappa-1}=\mathrm{const},\qquad
 \frac{\dd T}{T}=(1-\kappa)\frac{\dd V}{V}.
\]
$\kappa$ 是路径参数，不能一看到幂律便将它写成气体的 $\gamma=c_p/c_v$。
设 $c_v$ 为常数，当 $\kappa\ne1$ 时，第一定律逐项给出
\[
 \delta Q=nc_v\dd T+p\dd V
 =n\left(c_v+\frac{R}{1-\kappa}\right)\dd T.
\]
括号内是沿该路径定义的摩尔过程热容，不是新的物性常数。等压的 $\kappa=0$ 恢复 $c_p$；可逆绝热的 $\kappa=\gamma$ 使其为零。$\kappa=1$ 时温度不变却可以传热，必须直接算功，不能把上式发散理解为吸收无限热量。

\Needspace{180pt}
\begin{origexamplebox}{温度下降的膨胀是否一定吸热}\label{ch15:ex:teaching-polytropic}
一摩尔单原子理想气体从 $(p_1,V_1,T_1)$ 沿准静态路径 $pV^2=\mathrm{const}$ 膨胀至 $V_2=2V_1$。只做体积功，求末温、对外功、内能变化和热量，并解释热量的符号。
\end{origexamplebox}
\Needspace{65pt}
\begin{xesolution}
由路径和状态方程，$TV=\mathrm{const}$，所以 $T_2=T_1/2$。实际压力曲线给出
\[
 W^{\rm out}=\int_{V_1}^{2V_1}\frac{p_1V_1^2}{V^2}\dd V
 =\frac12p_1V_1=\frac12RT_1>0.
\]
单原子气体内能变化为
\[
 \Delta U=\frac32R(T_2-T_1)=-\frac34RT_1,
 \qquad Q=\Delta U+W^{\rm out}=-\frac14RT_1.
\]
气体边膨胀边放热：内能减少量比所做功更大，多减少的部分排给环境。也可用过程热容 $c=3R/2-R=R/2$，得到 $Q=c(T_2-T_1)$，与直接积分一致。与图\ref{ch15:fig:adiabatic-comparison}的两条路径比较，本题 $\kappa=2>\gamma=5/3$，膨胀降温比同初态可逆绝热更快，因而必须再向外放热。
\end{xesolution}

\section{重力场中形成对流的条件}

绝热方程的一个重要应用，是判断重力场中的理想气体能否形成对流。受扰动气团来到新高度后，会随外界
压强变化而膨胀或压缩。大气中高处压强较小，气团上升时迅速膨胀；若上升足够快，气团与周围空气来不及
交换显著热量，便可近似为绝热过程。

设大气具有已知的稳定温度分布 $T(h)$。在高度 $h$ 与 $h+\dd h$ 之间，静力平衡给出
\begin{equation*}
  \dd P=-\rho g\,\dd h=-\frac{\mu gP}{RT}\,\dd h,
\end{equation*}
其中 $\mu$ 为气体平均摩尔质量。把同一小气团从 $h$ 绝热抬升到 $h+\dd h$，由式~(\ref{ch15:eq:s29-20-9})可得它的
温度改变
\begin{equation*}
  \dd T_1=\frac{\gamma-1}{\gamma}\frac{T}{P}\dd P
  =-\frac{\gamma-1}{\gamma}\frac{\mu g}{R}\dd h
  =-\frac{\mu g}{c_p}\dd h.
\end{equation*}
而环境中原本位于 $h+\dd h$ 处的空气温度为
\begin{equation*}
  T_2=T+\frac{\dd T}{\dd h}\dd h.
\end{equation*}
两团气体在同一高度具有近似相同压强和成分；温度较高者密度较小。要使受扰动上升后的气团比周围空气
更热、更轻，从而继续上升，必须满足
\begin{equation}
  \boxed{\frac{\dd T}{\dd h}<-\frac{\mu g}{c_p}}. \label{ch15:eq:s29-20-11}
\end{equation}
因此，只要求低处比高处热还不够；环境温度随高度降低的速率必须大于气团的绝热降温率。

\par\noindent\begin{minipage}{\linewidth}
\pressureLayerRedraw
\diagramcaption{取横截面积为 $A$、厚度为 $\dd h$ 的水平气层，由竖直方向受力平衡得到 $\dd P=-\rho g\dd h$。}\label{ch15:fig:s29-20-7}
\end{minipage}\par\medskip
图~\ref{ch15:fig:s29-20-7}中上、下表面压强差支撑薄层的重量；$\dd h>0$ 时 $\dd P<0$。

\par\noindent\begin{minipage}{\linewidth}
\convectionCriterionDiagram
\figcap{受扰动气团与环境温度的比较。两者来到同一高度后，压强近似相等。}\label{ch15:fig:convection-displacement}
\end{minipage}\par\medskip
图\ref{ch15:fig:convection-displacement}中，应比较上方气团的 $T+\dd T_1$ 与该高度环境的 $T+\dd T_2$，
不能只比较气团上升前后的温度。

\par\noindent\begin{minipage}{\linewidth}
\atmosphereProfileFaithful
\diagramcaption{地球大气的代表性温度、压强与分层。下方列出相关自然情景和飞行器，二者不是同一比例的图。}\label{ch15:fig:s29-20-8}
\end{minipage}\par\medskip
图~\ref{ch15:fig:s29-20-8}上方的温度和压强对应同一个高度刻度；下方情景只帮助辨认物体，并不按高度比例排布。

从地表向上，对流层的平均温度通常随高度降低，但“随高度降低”本身不保证干对流不稳定。
例如约 $6.5\,\mathrm{K/km}$ 的平均温度递减率，小于干空气约 $9.8\,\mathrm{K/km}$ 的绝热递减率；
地表局部强加热或水汽凝结才可能使情况改变。再向上进入平流层，
温度可能随高度增加，对流受到抑制。云层多位于对流活跃区域，民航巡航高度则常接近平流层下部，以避开
低层大气中较强的天气扰动。水汽、二氧化碳等吸收地表长波辐射；平流层臭氧吸收太阳紫外线，
热层稀薄气体吸收更高能辐射。极光主要来自带电粒子激发高层气体，并不是空气热辐射的直接显示。
真实大气还受到水汽凝结、辐射吸收和地表加热等影响，式~(\ref{ch15:eq:s29-20-11})给出的是
干燥理想气体的基本判据。



\Needspace{16\baselineskip}\section{综合例题}
\begin{origexamplebox}{【正方形循环】}\label{ch15:ex:s29-20-1}
小明制作了一个气缸，封住了一些理想气体，并让气体在 $P$--$V$ 图中完成三种顺时针循环。三种循环
具有相同的边界压强 $P_1,P_2$ 和边界体积 $V_1,V_2$：

（1）矩形循环；（2）上三角循环；（3）下三角循环；（4）比较这三种循环以及不同大小循环的效率。

下图给出三种循环的形状与运行方向。设气体的定容摩尔
热容为 $c_v$，并在所考察温区内近似不变。
\par\noindent\begin{minipage}{\linewidth}
\threeCyclesDiagram
\figcap{矩形、上三角与下三角循环。三个图使用相同的 $P_1,P_2,V_1,V_2$。}\label{ch15:fig:three-geometric-cycles}
\end{minipage}\par\medskip
\end{origexamplebox}\par\nopagebreak[4]

\par\nopagebreak[4]\begin{xesolution}
记
\begin{equation*}
  \Delta P=P_2-P_1,\qquad \Delta V=V_2-V_1,
  \qquad \alpha=\frac{c_v}{R}.
\end{equation*}
对任意过程，理想气体的温度满足 $nRT=PV$，所以
\begin{equation*}
  \Delta U=n c_v\Delta T=\alpha\Delta(PV),
  \qquad Q=\Delta U+W^{\rm out}.
\end{equation*}

图\ref{ch15:fig:three-geometric-cycles}中的箭头均围出正的循环功，但各段吸热量并不相同。

（1）矩形循环的净功就是矩形面积：
\begin{equation*}
  W_{\text{矩}}=\Delta P\,\Delta V.
\end{equation*}
吸热发生在左侧等容升压段与上侧等压膨胀段。两段吸热量分别为
\begin{equation*}
  Q_{V_1}=\alpha V_1\Delta P,
  \qquad
  Q_{P_2}=(\alpha+1)P_2\Delta V.
\end{equation*}
因此
\begin{equation}
  \boxed{\eta_{\text{矩}}=
  \frac{\Delta P\Delta V}
  {\alpha V_1\Delta P+(\alpha+1)P_2\Delta V}}.
\end{equation}

（2）上三角循环的顺时针面积为矩形的一半。它的吸热过程仍是完整的左侧等容升压段和上侧等压膨胀段；
从右上角沿斜线回到左下角时，$P,V,T$ 都下降，同时外界对气体做功，该段为放热。因此
\begin{equation}
  \boxed{\eta_{\rm upper}=
  \frac{\tfrac12\Delta P\Delta V}
  {\alpha V_1\Delta P+(\alpha+1)P_2\Delta V}}.
\end{equation}

（3）下三角循环从左下角沿直线膨胀到右上角时吸热，其余等容降压、等压压缩两段都放热。直线段对外功
等于梯形面积，故吸热量为
\begin{align*}
  Q_{\rm diag}
  &=\alpha(P_2V_2-P_1V_1)
    +\frac{P_1+P_2}{2}(V_2-V_1).
\end{align*}
于是
\begin{equation}
  \boxed{\eta_{\rm lower}=
  \frac{\tfrac12\Delta P\Delta V}
  {\alpha(P_2V_2-P_1V_1)+\tfrac12(P_1+P_2)\Delta V}}.
\end{equation}

为了比较，令矩形循环的吸热量为 $Q_{\text{矩}}$、净功为 $W_{\text{矩}}$。下三角斜边吸热量可化为
\begin{equation*}
  Q_{\rm diag}=Q_{\text{矩}}-\frac12W_{\text{矩}}.
\end{equation*}
又因为这里的矩形热机有 $0<W_{\text{矩}}<Q_{\text{矩}}$，所以
\begin{equation*}
  \boxed{\eta_{\text{矩}}>\eta_{\rm lower}>\eta_{\rm upper}}.
\end{equation*}

（4）循环画得“面积更大”并不自动意味着效率更高。效率只依赖无量纲比值
$P_2/P_1$、$V_2/V_1$ 和 $c_v/R$。若把循环关于原点相似放大，使这两个比值不变，功和吸热量按同一
比例增大，效率不变；若固定 $P_1,V_1$ 而改变 $P_2,V_2$，则必须代入上述公式具体比较。功率、每循环
输出功和效率是三个不同概念。
\end{xesolution}


\Needspace{16\baselineskip}\begin{origexamplebox}{【新式气缸】}\label{ch15:ex:s29-20-2}
小明制作了一种新式气缸，气体在其中完成下列循环：

（1）两绝热、两等温循环；（2）两绝热、两等体循环；（3）两绝热、两等压循环。

试用各状态点的温度、体积或压强，以最简单形式表示循环效率。取定量、定组成的理想气体，
摩尔热容在整个温区内近似不变，所有绝热段均可逆。下图依次为卡诺、奥托和布雷顿循环。
\par\noindent\begin{minipage}{\linewidth}
\standardHeatCyclesDiagram
\figcap{三种循环的状态点与方向。均按 $1\to2\to3\to4\to1$ 运行；各曲线由同一气体的状态方程确定。}\label{ch15:fig:three-standard-cycles}
\end{minipage}\par\medskip
\end{origexamplebox}\par\nopagebreak[4]

\par\nopagebreak[4]\begin{xesolution}
三种循环都用绝热段连接吸热、放热过程。绝热段本身 $Q=0$，关键是分别计算另外两段的吸、放热量。

在图\ref{ch15:fig:three-standard-cycles}中，卡诺循环的 $1\to2$ 为高温等温膨胀；奥托与布雷顿循环的
$1\to2$ 则为绝热压缩。下面的温度下标严格沿用各图的状态号。

（1）两条等温线温度分别为 $T_h,T_c$。高温等温膨胀吸热，低温等温压缩放热。由两条绝热线
$TV^{\gamma-1}=\mathrm{const}$ 有
\[
 T_hV_2^{\gamma-1}=T_cV_3^{\gamma-1},\qquad
 T_hV_1^{\gamma-1}=T_cV_4^{\gamma-1}.
\]
相除得 $(V_2/V_1)^{\gamma-1}=(V_3/V_4)^{\gamma-1}$。
记两边的体积比为 $r>1$，两等温段的吸、放热量大小因此为
\begin{equation*}
  Q_h=nRT_h\ln r,\qquad Q_c=nRT_c\ln r.
\end{equation*}
故卡诺循环效率为
\begin{equation}
  \boxed{\eta_{\rm C}=1-\frac{T_c}{T_h}}.
\end{equation}

（2）两条等容线位于 $V_{\min}$、$V_{\max}$，定义压缩比
$r=V_{\max}/V_{\min}>1$。等容吸、放热量为
\begin{equation*}
  Q_h=n c_v(T_3-T_2),\qquad Q_c=n c_v(T_4-T_1).
\end{equation*}
两条绝热线给出
\begin{equation*}
  \frac{T_2}{T_1}=\frac{T_3}{T_4}=r^{\gamma-1},
\end{equation*}
分别写成 $T_2=r^{\gamma-1}T_1$ 和 $T_3=r^{\gamma-1}T_4$，相减得
$T_3-T_2=r^{\gamma-1}(T_4-T_1)$。代入热量比并约去 $nc_v$，
得到 $Q_c/Q_h=1/r^{\gamma-1}$，奥托循环效率为
\begin{equation}
  \boxed{\eta_{\rm O}=1-\frac{1}{r^{\gamma-1}}}.
\end{equation}

（3）两条等压线压强分别为 $P_h,P_l$，定义压强比 $r_p=P_h/P_l>1$。等压吸、放热量为
\begin{equation*}
  Q_h=n c_p(T_3-T_2),\qquad Q_c=n c_p(T_4-T_1).
\end{equation*}
由绝热关系 $T^\gamma/P^{\gamma-1}=\mathrm{const}$，
\begin{equation*}
  \frac{T_2}{T_1}=\frac{T_3}{T_4}=r_p^{(\gamma-1)/\gamma}.
\end{equation*}
令 $b=r_p^{(\gamma-1)/\gamma}$，则 $T_2=bT_1,T_3=bT_4$，所以
$T_3-T_2=b(T_4-T_1)$。约去热量比中的 $nc_p$，得到布雷顿循环效率
\begin{equation}
  \boxed{\eta_{\rm B}=1-\frac{1}{r_p^{(\gamma-1)/\gamma}}}.
\end{equation}
三式的共同结构都是 $\eta=1-Q_c/Q_h$；区别在于吸、放热段受何种状态量约束。
这里 $Q_h,Q_c$ 均为正的热量大小。效率无量纲，且在温差、压缩比或压强比
分别趋于零、1、1时趋于零；同一气体、同一可逆绝热条件必须贯穿整个推导。
\end{xesolution}


\Needspace{16\baselineskip}\begin{origexamplebox}{【变化的热容】}\label{ch15:ex:s29-20-3}
小明发现，气体在一段过程中对应的过程热容未必保持为常数。为验证这一点，他设计了一摩尔理想气体
经历的过程
\begin{equation*}
  \frac{P}{P_0}+\frac{V}{V_0}=1.
\end{equation*}
设气体的定容摩尔热容为 $c_v$，求整个过程中摩尔过程热容随体积的变化关系。
\end{origexamplebox}\par\nopagebreak[4]

\par\nopagebreak[4]\begin{xesolution}
过程热容定义为 $C=\delta Q/\dd T$。它是过程属性，不等同于气体本身的 $c_v$ 或 $c_p$。
由 $P>0,V>0$，本题的物理区间为 $0<V<V_0$；两截距只是边界，不是气体可取的常温有限状态。由过程方程
\begin{equation*}
  P=P_0\left(1-\frac{V}{V_0}\right),
\end{equation*}
一摩尔理想气体的温度为
\begin{equation*}
  T(V)=\frac{PV}{R}
  =\frac{P_0}{R}\left(V-\frac{V^2}{V_0}\right),
\end{equation*}
从而
\begin{equation*}
  \frac{\dd T}{\dd V}=\frac{P_0}{R}\left(1-\frac{2V}{V_0}\right).
\end{equation*}

\par\noindent\begin{minipage}{\linewidth}
\variableHeatCapacityDiagram
\figcap{线性 $P(V)$ 与相应的 $T(V)$。温度抛物线的顶点严格位于两个体积截距的中间。}\label{ch15:fig:variable-heat}
\end{minipage}\par\medskip
图\ref{ch15:fig:variable-heat}中左侧压强随体积线性下降，右侧温度却先升后降，这正是过程热容变号的原因。

由第一定律
\begin{equation*}
  \delta Q=c_v\dd T+P\dd V,
\end{equation*}
所以
\begin{align*}
  C(V)
  &=c_v+P\frac{\dd V}{\dd T}\\
  &=c_v+R\frac{1-V/V_0}{1-2V/V_0}.
\end{align*}
即
\begin{equation}
  \boxed{C(V)=c_v+R\frac{V_0-V}{V_0-2V}}.
\end{equation}
在 $V=V_0/2$ 处，温度达到最大值，温度的一阶变化为零，但热量的一阶变化仍有
$\delta Q=P\dd V=(P_0/2)\dd V\ne0$，所以以 $\delta Q/\dd T$
定义的过程热容发散。发散不是气体储存了无穷能量，而是“用温度变化量作分母”在温度转折点失去良好
定义。令 $z=V/V_0$、$\alpha=c_v/R>0$，则
\[
 \frac{C}{R}=\frac{(\alpha+1)-(2\alpha+1)z}{1-2z},\qquad
 z_* =\frac{\alpha+1}{2\alpha+1}\in(1/2,1).
\]
因此 $0<z<1/2$ 时 $C>0$；$1/2<z<z_*$ 时 $C<0$；$z=z_*$ 时 $C=0$；
$z_*<z<1$ 时 $C>0$。膨胀经过负热容区时仍吸热，却因对外做功更多而降温。
在零点，内能减少恰好补偿膨胀功，元热量为零。$C$ 的单位始终与 $R,c_v$ 相同。
\end{xesolution}


\Needspace{16\baselineskip}\begin{origexamplebox}{【恒定的热容】}\label{ch15:ex:s29-20-4}
小明猜想，可能存在一些过程使过程热容保持不变，例如常见的等容、等压、等温和绝热过程。但逐一枚举
容易遗漏。能否得到“过程热容为常量”的统一过程方程？
\end{origexamplebox}\par\nopagebreak[4]

\par\nopagebreak[4]\begin{xesolution}
设一摩尔、定组成理想气体沿准静态过程的摩尔过程热容为有限常量 $C$，并取常数 $c_v,c_p$。由定义与第一定律
\begin{equation*}
  C\dd T=c_v\dd T+P\dd V,
\end{equation*}
即
\begin{equation*}
  (C-c_v)\dd T=P\dd V.
\end{equation*}
又由 $RT=PV$，
\begin{equation*}
  R\dd T=P\dd V+V\dd P.
\end{equation*}
消去 $\dd T$，得到
\begin{equation*}
  (C-c_v)V\dd P+(C-c_p)P\dd V=0.
\end{equation*}
在 $C\ne c_v$ 时两边除以 $PV(C-c_v)$ 并积分：
\begin{equation*}
  \frac{\dd P}{P}+k\frac{\dd V}{V}=0,
  \qquad
  k=\frac{C-c_p}{C-c_v},
\end{equation*}
所以统一的过程方程是
\begin{equation}
  \boxed{PV^k=\mathrm{const}},
  \qquad
  \boxed{k=\frac{C-c_p}{C-c_v}},
\end{equation}
或反解为
\begin{equation}
  \boxed{C=c_v+\frac{R}{1-k}}.
\end{equation}
这就是多方过程。它包含若干重要极限：
\begin{center}
\small
\begin{tabular}{c c c}
\toprule
过程 & 多方指数 $k$ & 过程热容 $C$\\
\midrule
等压 & $0$ & $c_p$\\
等温 & $1$ & 形式上发散（$\dd T=0$）\\
绝热 & $\gamma$ & $0$\\
等容 & $k\to\infty$ & $c_v$\\
\bottomrule
\end{tabular}
\end{center}\nopagebreak[2]
特别地，$C=c_v$ 时不能除以 $C-c_v$，应回到 $(C-c_v)\dd T=P\dd V$，直接得到 $\dd V=0$。
$k=1$ 的等温过程通常有 $\delta Q\ne0$、$\dd T=0$，并没有一个有限的过程热容，
只能列为 $|C|\to\infty$ 的极限。因此有限恒热容过程除等容这一单列情形外，构成
$PV^k=\mathrm{const}$ 的一整族过程；不能把“无穷大”当作与 $c_p$ 一样的有限常数。
\end{xesolution}


\Needspace{16\baselineskip}\begin{origexamplebox}{【等温大气】}\label{ch15:ex:s29-20-5}
小明研究一个简化的大气模型。假设大气经过长时间热接触后处处具有相同温度 $T$，气体平均摩尔质量
为 $\mu$。考虑气体密度随压强改变，求压强和分子数密度随高度 $h$ 的变化关系。
\end{origexamplebox}\par\nopagebreak[4]

\par\nopagebreak[4]\begin{xesolution}
取高度 $h$ 处厚度为 $\dd h$ 的水平气层。静力平衡要求
\begin{equation*}
  \dd P=-\rho g\,\dd h.
\end{equation*}
等温理想气体满足
\begin{equation*}
  \rho=\frac{\mu P}{RT},
\end{equation*}
故
\begin{equation*}
  \frac{\dd P}{P}=-\frac{\mu g}{RT}\dd h.
\end{equation*}
取地面 $h=0$ 处压强为 $P_0$，积分得到
\begin{equation}
  \boxed{P(h)=P_0\exp\!\left(-\frac{\mu gh}{RT}\right)}.
\end{equation}
分子数密度 $n_N=P/(k_BT)$，所以
\begin{equation}
  \boxed{n_N(h)=n_{N0}\exp\!\left(-\frac{\mu gh}{RT}\right)}.
\end{equation}

\par\noindent\begin{minipage}{\linewidth}
\isothermalAtmosphereDiagram
\figcap{等温大气的指数分布。标高 $H$ 决定压强和数密度衰减的快慢。}\label{ch15:fig:isothermal-atmosphere}
\end{minipage}\par\medskip
图\ref{ch15:fig:isothermal-atmosphere}的指数关系要求温度、成分与重力加速度近似不随高度变化。

定义等温标高
\begin{equation*}
  \boxed{H=\frac{RT}{\mu g}},
\end{equation*}
则压强、质量密度和分子数密度都按 $\e^{-h/H}$ 衰减。真实大气并不等温，重力加速度和气体成分也会随
高度改变；本题给出的是最基本的指数大气模型。
\end{xesolution}


\Needspace{16\baselineskip}\begin{origexamplebox}{【毛细现象】}\label{ch15:ex:s29-20-6}
小明把一根内半径为 $R$ 的玻璃管插入水面。假设水完全润湿玻璃，表面张力系数为 $\sigma$，水的密度
为 $\rho$。水面会上升多高？小明认为大树吸水、蒸腾也是同样机制，那么仅靠毛细作用给
$\SI{10}{m}$ 高处供水，需要多细的管半径？
\end{origexamplebox}\par\nopagebreak[4]

\par\nopagebreak[4]\begin{xesolution}
完全润湿时接触角 $\theta=0$。沿管内圆周，表面张力的竖直合力为
\begin{equation*}
  F_\sigma=2\pi R\sigma\cos\theta=2\pi R\sigma.
\end{equation*}
上升水柱的重力为
\begin{equation*}
  G=\rho g(\pi R^2h).
\end{equation*}
平衡时二者相等，故
\begin{equation}
  \boxed{h=\frac{2\sigma}{\rho gR}}.
\end{equation}

\par\noindent\begin{minipage}{\linewidth}
\capillaryRevisionDiagram
\figcap{完全润湿的毛细管。弯月面中心低、壁边高，壁边表面张力沿切线向上。}\label{ch15:fig:capillary-wetting}
\end{minipage}\par\medskip
图\ref{ch15:fig:capillary-wetting}的 $h$ 从外部水平水面量起。这里忽略弯月面附近体积相对长水柱的微小修正；
接触角为零使周界张力全用于托住水柱，量纲也给出 $[\sigma/(\rho gR)]=\mathrm m$。

若要求 $h=\SI{10}{m}$，则
\begin{equation*}
  \boxed{R=\frac{2\sigma}{\rho gh}}.
\end{equation*}
取室温水 $\sigma\approx\SI{0.073}{N.m^{-1}}$、$\rho\approx\SI{1.0e3}{kg.m^{-3}}$，得到
\begin{equation*}
  R\approx\frac{2\times0.073}{10^3\times9.8\times10}
  \approx\SI{1.5e-6}{m}.
\end{equation*}
这说明微米级管道的毛细压强足以对应十米水柱。但真实树木的远距离输水并非单靠静态毛细平衡：叶片
蒸腾在木质部水柱中产生负压，水分子的内聚力和导管结构共同维持连续水柱；毛细作用主要帮助润湿和
稳定细小导管。
\end{xesolution}


\Needspace{16\baselineskip}\begin{origexamplebox}{【吹泡泡】}\label{ch15:ex:s29-20-7}
小明拿出泡泡水吹泡泡。设泡泡水的表面张力系数为 $\sigma$，泡泡为半径 $R$ 的正球。

（1）若不计液膜厚度，泡泡内部压强为多大？

（2）若平衡时液膜的质量面密度为 $\sigma_0$，忽略空气质量，把空气看成双原子理想气体。在小的径向
扰动下，设气泡仍能维持球形，判断平衡是否稳定。
\end{origexamplebox}\par\nopagebreak[4]

\par\nopagebreak[4]\begin{xesolution}
（1）肥皂泡有内、外两个气液界面，每个球面都产生 $2\sigma/R$ 的拉普拉斯压强差，因此总压强差为
\begin{equation*}
  p_{\rm in}-p_0=\frac{4\sigma}{R}.
\end{equation*}
所以
\begin{equation}
  \boxed{p_{\rm in}=p_0+\frac{4\sigma}{R}}.
\end{equation}
注意：只有一个界面的液滴压强差是 $2\sigma/R$；肥皂泡有两个界面，不能漏掉因子2。

\par\noindent\begin{minipage}{\linewidth}
\bubbleRevisionDiagram
\figcap{肥皂泡的两界面与径向扰动。$p_R$ 是平衡时内部压强；虚线表示均匀向外位移。}\label{ch15:fig:bubble-radius}
\end{minipage}\par\medskip
也可直接对图\ref{ch15:fig:bubble-radius}作虚功计算：$\dd V=4\pi r^2\dd r$，
两界面的面积增量为 $\dd A_s=16\pi r\dd r$，平衡要求
$(p_R-p_0)\dd V=\sigma\dd A_s$，由此再次得到因子4。

（2）设半径受扰动变为 $r=R+x$，$|x|\ll R$。小振幅并不自动意味着变化很快：
若传热足够快，气体接近等温；若振动周期远短于漏热时间、又足以维持内部压强近似均匀，则接近绝热。
先统一用多方指数 $\kappa$ 表示这两种极限：等温为1，常温双原子气体的可逆绝热近似为 $7/5$。
取表面张力 $\sigma$ 在扰动中不变，忽略气体惯性、膜的黏性耗散和外部空气附加质量。记 $p_R=p_{\rm in}$，于是
\begin{equation*}
  p(r)r^{3\kappa}=p_R R^{3\kappa},
  \qquad
  p'(R)=-\frac{3\kappa p_R}{R}.
\end{equation*}
球面对称时所有真实力的矢量合力为零，不能把它作为整个泡的平移合力。这里计算的是与半径 $r$
共轭的广义力，即虚功中 $\dd r$ 的系数：压强贡献为 $4\pi r^2[p(r)-p_0]$。
两个界面的表面张力功对应有效势 $E_s=8\pi\sigma r^2$，向内的广义力为
\begin{equation*}
  -\frac{\dd E_s}{\dd r}=-16\pi\sigma r.
\end{equation*}
故总径向广义力为
\begin{equation*}
  F(r)=4\pi r^2[p(r)-p_0]-16\pi\sigma r.
\end{equation*}
在 $r=R$ 处利用平衡条件 $p_{\rm in}-p_0=4\sigma/R$，得到
\begin{equation*}
  F'(R)=4\pi\bigl(4\sigma-3\kappa Rp_R\bigr).
\end{equation*}
液膜质量为 $m_f=4\pi R^2\sigma_0$。线性化运动方程写成
\begin{equation*}
  m_f\ddot x=F'(R)x=-m_f\omega^2x,
\end{equation*}
其中
\begin{equation}
  \boxed{\omega^2=
  \frac{3\kappa Rp_R-4\sigma}{R^2\sigma_0}}.
\end{equation}
代入 $p_R=p_0+4\sigma/R$，分子为
\begin{equation*}
  3\kappa Rp_0+(12\kappa-4)\sigma>0\qquad(\kappa=1\ \text{或}\ 7/5).
\end{equation*}
所以 $\omega^2>0$，小的球对称径向扰动会产生回复力，
\begin{equation*}
  \boxed{\text{该平衡对题设的小径向扰动稳定。}}
\end{equation*}
这里判断的是理想化的球对称模式；真实泡膜还可能因排液、蒸发、膜厚不均或非球形扰动而破裂。
由 $[\sigma]=\mathrm{N/m}$、$[\sigma_0]=\mathrm{kg/m^2}$ 可验得 $[\omega^2]=\mathrm{s^{-2}}$。
一般传热速率介于两极限之间时，还可能出现相位滞后和阻尼，题设不足以指定唯一振动频率；但上述两极限都稳定。
\end{xesolution}


\Needspace{16\baselineskip}\begin{origexamplebox}{【华氏温标】}\label{ch15:ex:s29-20-8}
美国常用的华氏温标把标准大气压下冰的熔点定为 $32\,{}^\circ\mathrm F$，水的沸点定为
$212\,{}^\circ\mathrm F$。

（1）导出华氏温度与摄氏温度之间的换算关系；（2）给出绝对零度的华氏温度。
\end{origexamplebox}\par\nopagebreak[4]

\par\nopagebreak[4]\begin{xesolution}
摄氏温标中冰点与沸点相差 $100$ 度，华氏温标中相差 $212-32=180$ 度，因此同一温差满足
\begin{equation*}
  \Delta t_F=\frac{180}{100}\Delta t_C=\frac95\Delta t_C.
\end{equation*}
再用 $t_C=0$ 对应 $t_F=32$，得到
\begin{equation}
  \boxed{t_F=\frac95t_C+32},
  \qquad
  \boxed{t_C=\frac59(t_F-32)}.
\end{equation}
绝对零度 $t_C=-273.15\,{}^\circ\mathrm C$，所以
\begin{equation*}
  t_F=\frac95(-273.15)+32=-459.67\,{}^\circ\mathrm F.
\end{equation*}
即
\begin{equation}
  \boxed{0\,\mathrm K=-459.67\,{}^\circ\mathrm F}.
\end{equation}
温差换算只乘 $9/5$，绝对温度读数换算还必须处理零点偏移；这是两类公式最容易混淆之处。
\end{xesolution}


\Needspace{16\baselineskip}\begin{origexamplebox}{【循环过程】}\label{ch15:ex:s29-20-9}
如图所示，一定量单原子理想气体从初态 $a(P_1,V_1)$ 出发，先经等容过程到状态 $b$，再经等压过程
到状态 $c$，最后经等温过程回到 $a$，完成一个循环。图中 $P_b=P_c=P_1/4$。已知一摩尔气体温度
改变 $\SI{1}{K}$ 时内能改变 $c_v=3R/2$。

（1）$c\to a$ 吸热还是放热，热量多少？（2）$a\to b$ 吸热还是放热，热量多少？

（3）$b\to c$ 吸热还是放热，热量多少？（4）整个循环的净功和净吸放热量各为多少？
\par\noindent\begin{minipage}{\linewidth}
\cycleQuestionRedraw
\diagramcaption{$a\to b\to c\to a$ 循环。箭头给出过程方向，$c\to a$ 为等温压缩。}\label{ch15:fig:s29-20-9}
\end{minipage}\par\medskip
图~\ref{ch15:fig:s29-20-9}中 $a,c$ 在同一条双曲线上，而 $b,c$ 具有相同压强。
\end{origexamplebox}\par\nopagebreak[4]

\par\nopagebreak[4]\begin{xesolution}
过程 $c\to a$ 等温，所以 $P_cV_c=P_aV_a$。由 $P_c=P_1/4$ 得
\begin{equation*}
  \frac{P_1}{4}V_c=P_1V_1,
  \qquad
  \boxed{V_c=4V_1}.
\end{equation*}
设气体物质的量为 $n$，并统一取 $Q>0$ 为吸热、$W^{\rm out}>0$ 为气体对外做功。

\par\noindent\begin{minipage}{\linewidth}
\cycleExampleDiagram
\figcap{循环计算所用状态。$a,b$ 等体积，$b,c$ 等压，$a,c$ 等温。}\label{ch15:fig:cycle-solution}
\end{minipage}\par\medskip
以下逐段热量与图\ref{ch15:fig:cycle-solution}的三条有向路径对应。

（1）$c\to a$ 为温度 $T_a=P_1V_1/(nR)$ 的等温压缩，$\Delta U_{ca}=0$，
\begin{equation*}
  Q_{ca}=W^{\rm out}_{ca}
  =nRT_a\ln\frac{V_a}{V_c}
  =P_1V_1\ln\frac14.
\end{equation*}
因此
\begin{equation}
  \boxed{Q_{ca}=-P_1V_1\ln4},
\end{equation}
即气体放热 $P_1V_1\ln4$。

（2）$a\to b$ 等容，且 $T_b=T_a/4$。所以
\begin{equation*}
  W^{\rm out}_{ab}=0,
  \qquad
  Q_{ab}=\Delta U_{ab}
  =\frac32nR\left(\frac{T_a}{4}-T_a\right)
  =-\frac98P_1V_1.
\end{equation*}
故
\begin{equation}
  \boxed{Q_{ab}=-\frac98P_1V_1},
\end{equation}
即放热 $9P_1V_1/8$。

（3）$b\to c$ 等压，温度由 $T_a/4$ 升到 $T_a$。单原子气体 $c_p=c_v+R=5R/2$，所以
\begin{equation*}
  Q_{bc}=nc_p(T_c-T_b)
  =\frac52nR\cdot\frac34T_a
  =\frac{15}{8}P_1V_1.
\end{equation*}
即
\begin{equation}
  \boxed{Q_{bc}=\frac{15}{8}P_1V_1},
\end{equation}
气体吸热。

（4）净功可由各段相加。等容段功为零，等压段功为
\begin{equation*}
  W^{\rm out}_{bc}=\frac{P_1}{4}(4V_1-V_1)=\frac34P_1V_1,
\end{equation*}
等温压缩段功为 $-P_1V_1\ln4$。因此
\begin{equation}
  \boxed{W^{\rm out}_{\rm cyc}=P_1V_1\left(\frac34-\ln4\right)<0}.
\end{equation}
循环内能改变量为零，故净热量与净功相等：
\begin{equation}
  \boxed{Q_{\rm cyc}=P_1V_1\left(\frac34-\ln4\right)<0}.
\end{equation}
所以该循环是逆时针循环，外界对气体做净功
$P_1V_1(\ln4-3/4)$，气体向外净放出同样多的热量。
\end{xesolution}


\Needspace{16\baselineskip}\begin{origexamplebox}{【直线过程】}\label{ch15:ex:s29-20-10}
如图所示，一摩尔理想气体从 $P$--$V$ 图中的状态 $A$ 出发，沿缓慢直线过程到达状态 $B$。已知
$P_B/P_A=1/2$。若要使全过程的最终结果是气体从外界吸收热量，状态 $B$ 与状态 $A$ 的体积比应满足
什么条件？已知该理想气体每摩尔内能为 $U=3RT/2$。
\par\noindent\begin{minipage}{\linewidth}
\straightQuestionRedraw
\diagramcaption{状态 $A$ 到状态 $B$ 的准静态直线过程；解答中的详图进一步补出端点压强、体积及面积关系。}\label{ch15:fig:s29-20-10}
\end{minipage}\par\medskip
图~\ref{ch15:fig:s29-20-10}中 $B$ 的纵坐标为 $A$ 的一半，这个比例将同时用于内能变化和梯形面积。
\end{origexamplebox}\par\nopagebreak[4]

\par\nopagebreak[4]\begin{xesolution}
令
\begin{equation*}
  x=\frac{V_B}{V_A},
  \qquad P_B=\frac12P_A.
\end{equation*}
直线过程的对外功等于图下梯形面积：
\begin{equation*}
  W^{\rm out}
  =\frac{P_A+P_B}{2}(V_B-V_A)
  =\frac34P_AV_A(x-1).
\end{equation*}

\par\noindent\begin{minipage}{\linewidth}
\straightProcessDiagram
\figcap{直线过程的两个端点及轴上投影。积分区间是 $V_A$ 到 $V_B$。}\label{ch15:fig:straight-solution}
\end{minipage}\par\medskip
图\ref{ch15:fig:straight-solution}中直线下方的梯形面积给功，端点的 $PV$ 差给内能变化，二者要分别计算。

由 $U=3RT/2$ 以及一摩尔理想气体 $RT=PV$，
\begin{align*}
  \Delta U
  &=\frac32(P_BV_B-P_AV_A)\\
  &=\frac34P_AV_A(x-2).
\end{align*}
因此吸热量
\begin{align*}
  Q&=\Delta U+W^{\rm out}\\
   &=\frac34P_AV_A[(x-2)+(x-1)]\\
   &=\frac34P_AV_A(2x-3).
\end{align*}
要求气体净吸热即 $Q>0$，得到
\begin{equation}
  \boxed{\frac{V_B}{V_A}>\frac32}.
\end{equation}
若 $1<x<3/2$，气体虽然膨胀并对外做功，但温度、内能下降得更多，仍会净放热；因此“膨胀”不等于
“吸热”。
\end{xesolution}


\Needspace{16\baselineskip}\begin{origexamplebox}{【提升活塞】}\label{ch15:ex:s29-20-11}
如图所示，在质量为 $M$、内部横截面积为 $A$ 的竖直绝热气缸中，活塞封闭了一定量初温为 $T_0$ 的
理想气体。活塞也绝热，活塞质量及活塞与气缸之间的摩擦均可忽略。大气压强为 $p_0$，重力加速度为
$g$。气缸由不密封的小支脚支承，大气能作用于缸底外侧；忽略气体质量，取恒定绝热指数
$\gamma=c_p/c_v$，并满足 $Mg<p_0A$。现把活塞缓慢上提；当活塞到达气缸开口处、气体尚未泄出时，
气缸恰好离开支承面。求此时气体温度。
\par\noindent\begin{minipage}{\linewidth}
\liftingPistonQuestionRedraw
\diagramcaption{初态与刚离地终态。终态取活塞仍密封气体的极限位置；缸底外侧受到大气压。}\label{ch15:fig:s29-20-11}
\end{minipage}\par\medskip
图~\ref{ch15:fig:s29-20-11}右侧取开口前最后密封的状态；活塞再向上、气体泄出后就不能继续使用同一封闭系统的绝热方程。
\end{origexamplebox}\par\nopagebreak[4]

\par\nopagebreak[4]\begin{xesolution}
初始时活塞没有质量，又未受额外拉力，故气体压强为 $p_i=p_0$。上提过程缓慢，使每一时刻都可视为
平衡态；气缸、活塞又绝热，所以封闭气体经历可逆绝热膨胀。

设某时刻气体压强为 $p$，外界对活塞的上拉力为 $F$。质量可忽略的活塞受力平衡给出
\begin{equation*}
  F+pA=p_0A,
  \qquad
  F=(p_0-p)A.
\end{equation*}
还必须单独检查缸体，不能说低压气体把缸底“向上拉”。缸底内侧气体给缸体向下的力 $pA$，
外侧大气给向上的力 $p_0A$，支承力为 $N$。竖直平衡给出
\[
 N+p_0A=pA+Mg,\qquad N=Mg-(p_0-p)A.
\]
随着气体膨胀，$p$ 下降，支承力逐渐减小。刚离地时 $N=0$，于是
\begin{equation*}
  \boxed{p_f=p_0-\frac{Mg}{A}}.
\end{equation*}

\par\noindent\begin{minipage}{\linewidth}
\liftingPistonDiagram
\figcap{活塞与缸体分别受力。力箭头的作用对象不同，不能直接在一张受力图中混加。}\label{ch15:fig:lift-forces}
\end{minipage}\par\medskip
图\ref{ch15:fig:lift-forces}左图给 $F=(p_0-p)A$，右图在 $N=0$ 时给 $(p_0-p_f)A=Mg$，因此才有 $F_f=Mg$。

由绝热关系 $T^\gamma/P^{\gamma-1}=\mathrm{const}$，
\begin{equation*}
  \frac{T_f^\gamma}{p_f^{\gamma-1}}
  =\frac{T_0^\gamma}{p_0^{\gamma-1}}.
\end{equation*}
因此
\begin{equation}
  \boxed{T_f=T_0\left(\frac{p_f}{p_0}\right)^{(\gamma-1)/\gamma}
  =T_0\left(1-\frac{Mg}{p_0A}\right)^{(\gamma-1)/\gamma}}.
\end{equation}
要使末态压强为正并能准静态达到，题设参数必须满足 $Mg<p_0A$。上提过程中气体膨胀并对外做功，却
没有从环境吸热，因此内能和温度都降低，结果中的括号确实小于1。
\end{xesolution}
}

\endgroup

\hypersetup{bookmarksdepth=2}

% ThermalQualityRevision20261023-v1
% QualityRevision20261023
\hypersetup{bookmarksdepth=3}
\begingroup
\renewcommand{\Needspace}[1]{\par\begingroup
\dimen0=#1\relax\vskip0pt plus\dimen0\penalty-100
\vskip0pt plus-\dimen0\vskip\dimen0\penalty9999\vskip-\dimen0
\vskip0pt\endgroup}
\tcbset{breakable/.default=false}
\chapter{热力学第二定律}
% Editable TikZ and prose: state entropy versus actual process.
\providecommand{\ThermalFigure}[3]{\begin{qualityfigure}
\begin{tikzpicture}[>=Latex,font=\small,x=.8cm,y=.8cm]#1\end{tikzpicture}
\figcap{#2}\label{#3}\end{qualityfigure}}
\newcommand{\ThQTwoPaths}{\ThermalFigure{
 \draw[thick](0,0)rectangle(4,2);\draw[line width=2pt](2,0)--(2,2);
 \node at(1,1){$T_0,V_0$};\node at(3,1){真空};
 \draw[->,thick](4.5,1)--(5.5,1);\draw[thick](6,0)rectangle(10,2);\node at(8,1){$T_f,2V_0$};
 \node[below,align=center]at(5,-.4){甲：热库 $T_0$，可逆膨胀\quad 乙：隔热，撤去隔板};
}{同一初态到体积加倍的两种实现方式；末温及熵产生由过程分别判断。}{ch16:quality:paths-q}}
\newcommand{\ThATwoPaths}{\ThermalFigure{
 \node[draw,minimum width=2.4cm,minimum height=1cm] (A)at(0,0){$(T_0,V_0)$};
 \node[draw,minimum width=2.4cm,minimum height=1cm] (B)at(7,0){$(T_0,2V_0)$};
 \draw[->,thick,DeepBlue](A.north)..controls(2,2)and(5,2)..(B.north)node[midway,above]{甲：可逆等温，$Q_{\rm rev}$};
 \draw[->,thick](A.south)..controls(2,-2)and(5,-2)..(B.south)node[midway,below]{乙：自由膨胀，$Q=0$};
}{气体的熵变由相同端点决定，总熵产生还须计入热库。}{ch16:quality:paths-a}
图~\ref{ch16:quality:paths-a}的两条箭头连接同一对平衡态，所以气体的 $\Delta S$ 相同。乙过程不能用真实热量 $Q=0$ 代替上方辅助可逆路径的 $Q_{\rm rev}$；两种路径的环境熵变不同，才使总熵产生不同。
}
\newcommand{\ThFiniteWorkBase}[1]{
 \node[draw,minimum width=2.4cm,minimum height=.9cm](H)at(0,2){$C_1,T_h$};
 \node[draw,minimum width=2.4cm,minimum height=.9cm](L)at(0,-2){$C_2,T_c$};
 \node[draw,circle,minimum size=1cm](E)at(0,0){热机};
 \draw[->,thick](H)--node[right]{$\delta Q_h$}(E);\draw[->,thick](E)--node[right]{$\delta Q_c$}(L);
 \draw[->,DeepBlue,thick](E)--(3,0)node[right]{$\delta W$};#1}
\newcommand{\ThQFiniteWork}{\ThermalFigure{\ThFiniteWorkBase{
 \node[right,align=left]at(4,1){两端是有限物体\\温度均随取热、放热改变};
}}{整体隔热，只向外输出功；两物体最终达到共同温度。}{ch16:quality:finite-q}}
\newcommand{\ThAFiniteWork}{\ThermalFigure{
 \draw[->](0,0)--(6,0)node[right]{过程进度};\draw[->](0,0)--(0,3.6)node[above]{$T$};
 \draw[DeepBlue,thick](.3,3.1)..controls(2,2.3)and(4,1.8)..(5.5,1.8);
 \draw[thick](.3,.6)..controls(2,1.3)and(4,1.8)..(5.5,1.8);
 \node[left]at(0,3.1){$T_h$};\node[left]at(0,.6){$T_c$};\node[right]at(5.5,1.8){$T_*$（可逆极限）};
}{热物体降温，冷物体升温；可逆运行给出能达到的最低共同终温。}{ch16:quality:finite-a}
图~\ref{ch16:quality:finite-a}中的共同终温不是预先给定的热库温度。终温越低，留在两物体中的总内能越少，输出功越大；因此先用总熵不减确定终温下界，再由能量守恒求最大功。
}
\newcommand{\ThCoolingBase}[1]{
 \node[draw,minimum width=2.6cm,minimum height=.9cm](C)at(0,0){物体 $C,T$};
 \node[draw,circle,minimum size=1cm](R)at(4,0){制冷机};
 \node[draw,minimum width=2.6cm,minimum height=.9cm](H)at(8,0){环境 $T_0$};
 \draw[->,thick](C)--node[above]{$\delta Q_c$}(R);\draw[->,thick](R)--node[above]{$\delta Q_h$}(H);
 \draw[->,DeepBlue,thick](4,-1.7)--(R)node[midway,right]{$\delta W$};#1}
\newcommand{\ThQCooling}{\ThermalFigure{\ThCoolingBase{}}{制冷机循环复原，有限物体从 $T_0$ 冷却至 $T_1$。}{ch16:quality:cooling-q}}
\newcommand{\ThACooling}{\ThermalFigure{\ThCoolingBase{
 \node[below]at(0,-.9){$\dd T<0,\quad\delta Q_c=-C\dd T>0$};
 \node[above]at(4,1){可逆元循环：$\delta Q_h/T_0=\delta Q_c/T$};
}}{按物体瞬时温度逐次计算取热与元功，再沿降温过程积分。}{ch16:quality:cooling-a}
图~\ref{ch16:quality:cooling-a}中的热量箭头均取正的大小。物体每降低 $\dd T<0$ 放出 $-C\dd T$，环境收到的热量还包含输入功；可逆熵平衡给出 $\delta W=\delta Q_c(T_0/T-1)$，随 $T$ 改变，不能把终态制冷系数用于整个过程。
}
\newcommand{\ThQCarnot}{\ThermalFigure{
 \node[draw](H)at(0,2){热库 $T_1$};\node[draw,circle](E)at(0,0){循环工质};\node[draw](C)at(0,-2){热库 $T_2$};
 \draw[->,thick](H)--node[right]{$Q_1$}(E);\draw[->,thick](E)--node[right]{$Q_2$}(C);
 \draw[->,DeepBlue,thick](E)--(3,0)node[right]{$W$};\node[right]at(4,1.5){$T_1>T_2$};
}{只与两个热库交换热量的循环热机；先约定热量大小与方向。}{ch16:quality:carnot-q}}
\newcommand{\ThQSecondLaw}{\ThermalFigure{
 \node[draw](H)at(0,1.5){单一热库};\node[draw,circle](E)at(0,0){循环装置};
 \draw[->](H)--node[right]{$Q$}(E);\draw[->](E)--(2.3,0)node[right]{$W$};
 \node[below]at(0,-1){开尔文情境};
 \node[draw](C)at(6,0){低温热库 $T_c$};\node[draw](T)at(6,2){高温热库 $T_h$};
 \draw[->,thick](C)--node[right]{$Q$}(T);\node[below]at(6,-1){克劳修斯情境};
}{两种假设过程：完全把单库热量变成功，或不附带其他变化地逆温差传热。}{ch16:quality:second-law-q}}
\newcommand{\ThQIrreversibility}{\ThermalFigure{
 \draw[->](0,0)--(4.7,0)node[right]{$V$};\draw[->](0,0)--(0,3)node[above]{$p$};
 \draw[DeepBlue,thick,domain=1:4,samples=60]plot(\x,{2.6/\x});
 \draw[thick,domain=1:4,samples=60]plot(\x,{2.6/\x^1.5});
 \fill(1,2.6)circle(1.5pt)node[above right]{$A$};
 \node[right]at(4.8,2.2){等温与可逆绝热};
 \node[right]at(4.8,1.2){还需分别讨论有限温差传热、快速膨胀};
}{状态曲线的几何性质与过程能否自然发生属于不同问题，需要分别核查条件。}{ch16:quality:irreversibility-q}}
\newcommand{\ThQThreeBodies}{\ThermalFigure{
 \foreach \x/\lab/\temp in {0/1/300,3/2/300,6/3/100}{
 \draw[thick](\x,0)rectangle(2+\x,1.5);\node at(\x+1,1){物体 $\lab$};\node at(\x+1,.45){$c,\ \temp\,\mathrm K$};}
 \draw[dashed](-.4,-.4)rectangle(8.4,2);\node[above]at(4,2.1){外界不供热、不做功；内部可用循环热机};
}{三个物体组成隔离整体，最高温度由总能量与总熵共同限制。}{ch16:quality:three-q}}
\newcommand{\ThQSurfaceEngine}{\ThermalFigure{
 \draw[thick](0,0)rectangle(4,2);\fill[DeepBlue!12](0,0)rectangle(2.7,2);
 \draw[line width=2pt](2.7,0)--(2.7,2);\draw[->](2.7,2.5)--(3.7,2.5)node[right]{$\dd x$};
 \draw[<->](-.5,0)--node[left]{$\ell$}(-.5,2);\node at(1.3,1){表面 $A$};
 \draw[->,DeepBlue,thick](2.8,1)--(4.7,1)node[right]{维持准静态的外力};
 \node[right,align=left]at(5.4,1.4){界面张力 $\sigma(T)$\\温度由热库控制};
}{可移动边界改变界面面积；外力功与表面张力功的符号应随面积变化判断。}{ch16:quality:surface-q}}



{
\graphicspath{{../../../sample/xes-original/thirtieth-section/}{../../../sample/xes-original/thirtieth-section/assets/}}

% Original vector redraws for Lesson 21.  Text labels deliberately have no
% rectangular background; spacing is achieved by placement, not masking.

\tikzset{
  thermobody/.style={draw=Ink,thick},
  thermoarrow/.style={-{Latex[length=2.2mm]},draw=DeepBlue,very thick},
  thermoguide/.style={draw=TextGray,densely dashed},
  thermolabel/.style={font=\scriptsize,text=Ink,fill=none,inner sep=0pt,outer sep=.8pt}
}


\newcommand{\pistonstage}[4]{%
  \begin{scope}[shift={(#1,0)}]
    \draw[thermobody] (-0.62,-1.20)--(-0.62,1.15) (0.62,-1.20)--(0.62,1.15);
    \draw[thermobody] (-0.62,1.15)--(0.62,1.15);
    \draw[fill=SoftGray,draw=Ink,thick] (-0.54,#2) rectangle (0.54,{#2+0.18});
    \coordinate (Pin) at (0,#2);
    \coordinate (Cen) at (0,-1.78);
    \coordinate (Crank) at ($(Cen)+(#3:0.52)$);
    \draw[Ink,thick] (Pin)--(Crank);
    \draw[Ink,thick] (Cen) circle (0.52);
    \draw[Ink,thick] (Cen)--(Crank);
    \fill[Ink] (Cen) circle (1.1pt) (Crank) circle (1.1pt);
    \node[thermolabel,anchor=north] at (0,-2.42) {#4};
  \end{scope}
}

\newcommand{\ottoCycleFigure}{%
\begin{center}
\begin{tikzpicture}[x=.92cm,y=.92cm]
  \pistonstage{-6.8}{0.05}{220}{进气}
  \pistonstage{-5.05}{0.52}{300}{压缩}
  \pistonstage{-3.30}{0.72}{80}{点火}
  \pistonstage{-1.55}{-0.10}{155}{做功}
  \pistonstage{0.20}{0.18}{235}{排气}
  \draw[thermoarrow] (-7.85,0.85)--(-7.10,0.85);
  \node[thermolabel,above] at (-7.48,0.85) {进气};
  \draw[thermoarrow] (0.53,0.85)--(1.25,0.85);
  \node[thermolabel,above] at (0.88,0.85) {排气};
  \draw[thermoarrow] (-6.8,1.85)--(-6.8,1.25);
  \node[thermolabel,above right] at (-6.8,1.85) {火花塞};
  \draw[DeepBlue,thick] (-3.53,1.22)--(-3.30,1.57)--(-3.10,1.22)
    (-3.45,1.35)--(-3.28,1.80)--(-3.18,1.37);

  \begin{scope}[shift={(2.35,-2.15)},x=.9cm,y=.75cm]
    \draw[axis] (0,0)--(5.4,0) node[right] {$V$};
    \draw[axis] (0,0)--(0,4.9) node[above] {$P$};
    \coordinate (O) at (.55,.85);
    \coordinate (A) at (4.55,.85);
    \coordinate (B) at (.90,2.35);
    \coordinate (C) at (.90,4.15);
    \coordinate (D) at (4.55,1.75);
    \draw[DeepBlue,very thick,postaction={decorate},decoration={markings,
      mark=at position .35 with {\arrow{Latex}},mark=at position .78 with {\arrow{Latex}}}]
      (O)--(A)--(B)--(C) .. controls (1.85,3.55) and (3.20,2.35) .. (D)--(A)--(O);
    \foreach \p/\lab/\pos in {O/O/below left,A/A/below right,B/B/left,C/C/above left,D/D/above right}{
      \fill[Ink] (\p) circle (1.7pt) node[thermolabel,\pos] {$\lab$};
    }
    \draw[thermoguide] (.90,0)--(.90,4.15) (4.55,0)--(4.55,1.75);
    \node[thermolabel,below] at (.90,0) {$V_2$};
    \node[thermolabel,below] at (4.55,0) {$V_1$};
    \draw[thermoarrow] (1.25,2.75)--(1.25,3.65);
    \node[thermolabel,right] at (1.25,3.20) {$Q_h$};
    \draw[thermoarrow] (4.85,1.70)--(5.30,1.35);
    \node[thermolabel,right] at (5.18,1.48) {$Q_c$};
  \end{scope}
\end{tikzpicture}
\end{center}\nopagebreak[2]}

\newcommand{\carnotCycleFigure}{%
\begin{center}
\begin{tikzpicture}[x=1.08cm,y=.95cm]
  \draw[axis] (0,0)--(7.4,0) node[right] {$V$};
  \draw[axis] (0,0)--(0,5.0) node[above] {$P$};
  \draw[TextGray,thick,domain=.8:6.7,samples=100] plot (\x,{4.55/(\x^.43)});
  \draw[TextGray,thick,domain=.8:6.7,samples=100] plot (\x,{2.30/(\x^.43)});
  \coordinate (P1) at (1.35,4.00);
  \coordinate (P2) at (4.35,2.48);
  \coordinate (P3) at (5.55,1.10);
  \coordinate (P4) at (2.05,1.70);
  \draw[DeepBlue,very thick,postaction={decorate},decoration={markings,
    mark=at position .12 with {\arrow{Latex}},mark=at position .38 with {\arrow{Latex}},
    mark=at position .63 with {\arrow{Latex}},mark=at position .88 with {\arrow{Latex}}}]
    (P1) .. controls (2.25,3.38) and (3.35,2.76) .. (P2)
    .. controls (4.75,1.83) and (5.10,1.38) .. (P3)
    .. controls (4.35,1.20) and (3.00,1.45) .. (P4)
    .. controls (1.90,2.62) and (1.62,3.36) .. cycle;
  \foreach \p/\lab/\pos in {P1/1/above left,P2/2/above right,P3/3/below right,P4/4/below left}{
    \fill[Ink] (\p) circle (1.8pt) node[thermolabel,\pos] {$\lab$};
  }
  \node[thermolabel,anchor=west] at (5.55,3.15) {$T_1$ 等温线};
  \node[thermolabel,anchor=west] at (4.85,.55) {$T_2$ 等温线};
  \node[thermolabel] at (3.0,3.18) {$Q_1$};
  \node[thermolabel] at (3.75,1.15) {$Q_2$};
  \node[thermolabel,DeepBlue] at (4.05,4.35) {$T_1>T_2$};
\end{tikzpicture}
\end{center}\nopagebreak[2]}

\newcommand{\heatdevicepanel}[5]{%
  \begin{scope}[shift={(#1,0)}]
    \draw[thermobody,fill=SoftGray] (-1.40,2.15) rectangle (1.40,2.62);
    \draw[thermobody,fill=SoftGray] (-1.40,-2.62) rectangle (1.40,-2.15);
    \node[thermolabel] at (0,2.38) {高温热库 $T_h$};
    \node[thermolabel] at (0,-2.38) {低温热库 $T_c$};
    \draw[thermobody,fill=PaleBlue] (0,0) circle (.64);
    \node[thermolabel] at (0,0) {#2};
    #3 #4 #5
  \end{scope}
}

\newcommand{\heatDevicesFigure}{%
\begin{center}
\begin{tikzpicture}[x=1cm,y=.82cm]
  \heatdevicepanel{-5.7}{热机}{\draw[thermoarrow] (0,2.12)--(0,.70) node[midway,right,thermolabel] {$Q_h$};}{\draw[thermoarrow] (0,-.70)--(0,-2.12) node[midway,right,thermolabel] {$Q_c$};}{\draw[thermoarrow] (.68,0)--(1.55,0) node[midway,above,thermolabel] {$W$};}
  \heatdevicepanel{-1.9}{“完美”热机}{\draw[thermoarrow] (0,2.12)--(0,.70) node[midway,right,thermolabel] {$Q_h$};}{}{\draw[thermoarrow] (.68,0)--(1.55,0);\node[thermolabel] at (1.15,-1.05) {$W=Q_h$};}
  \heatdevicepanel{1.9}{制冷机}{\draw[thermoarrow] (0,.70)--(0,2.12) node[midway,right,thermolabel] {$Q_h$};}{\draw[thermoarrow] (0,-2.12)--(0,-.70) node[midway,right,thermolabel] {$Q_c$};}{\draw[thermoarrow] (1.55,0)--(.68,0) node[midway,above,thermolabel] {$W$};}
  \heatdevicepanel{5.7}{“完美”制冷机}{\draw[thermoarrow] (0,.70)--(0,2.12);}{\draw[thermoarrow] (0,-2.12)--(0,-.70);}{}
  \node[thermolabel] at (-5.7,-3.05) {(a) 真实热机};
  \node[thermolabel] at (-1.9,-3.05) {(b) 第二类永动机};
  \node[thermolabel] at (1.9,-3.05) {(c) 真实制冷机};
  \node[thermolabel] at (5.7,-3.05) {(d) 无功制冷机};
\end{tikzpicture}
\end{center}\nopagebreak[2]}

\newcommand{\portraitOne}[5]{%
  \begin{scope}[shift={(#1,0)}]
    \draw[fill=Ink,draw=Ink,rounded corners=1pt] (-1.62,-2.15) rectangle (1.62,2.25);
    \fill[#2] (-1.18,-2.10) .. controls (-1.02,-1.00) and (-.75,-.63) .. (-.48,-.50)
      --(.48,-.50) .. controls (.82,-.72) and (1.07,-1.15) .. (1.18,-2.10)--cycle;
    \fill[SoftGray] (0,.35) ellipse (.68 and 1.02);
    #3
    \draw[Ink,thick] (-.30,.47)--(-.08,.50) (.10,.50)--(.32,.47);
    \fill[Ink] (-.18,.49) circle (.04) (.18,.49) circle (.04);
    \draw[Ink,thin] (0,.37)--(-.05,.05)--(.06,.00);
    \draw[Ink,thin] (-.24,-.17) .. controls (0,-.28) and (.18,-.22) .. (.28,-.14);
    #4
    \node[font=\small\bfseries,anchor=north] at (0,-2.38) {#5};
  \end{scope}
}

\newcommand{\foundersPortraitFigure}{%
\begin{center}
\begin{tikzpicture}
  \portraitOne{-4.6}{TextGray}{\fill[Ink] (-.70,.82) .. controls (-.48,1.58) and (.55,1.58) .. (.72,.74) .. controls (.28,1.10) and (-.25,1.15) .. cycle;}{\draw[SoftGray,very thick] (-.40,-.62)--(-.15,-.88)--(0,-.70)--(.18,-.88)--(.42,-.62);}{萨迪·卡诺}
  \portraitOne{0}{RuleGray}{\fill[SoftGray] (-.60,.92) arc[start angle=180,end angle=0,x radius=.60,y radius=.50];}{\fill[SoftGray] (-.52,-.18) .. controls (-.62,-.95) and (-.35,-1.28) .. (0,-1.48) .. controls (.36,-1.25) and (.62,-.92) .. (.50,-.18) .. controls (.18,-.42) and (-.18,-.42) .. cycle;}{开尔文勋爵}
  \portraitOne{4.6}{TextGray}{\fill[Ink] (-.66,.88) .. controls (-.38,1.42) and (.48,1.34) .. (.68,.70) .. controls (.30,1.02) and (-.15,1.02) .. cycle;}{\draw[Ink,very thick] (-.34,-.12) .. controls (0,-.02) and (.18,-.04) .. (.36,-.12);}{鲁道夫·克劳修斯}
\end{tikzpicture}
\end{center}\nopagebreak[2]}

\newcommand{\entropyPathFigure}{%
\begin{center}
\begin{tikzpicture}[x=1cm,y=.9cm]
  \begin{scope}[shift={(-3.8,-1.4)}]
    \draw[axis] (0,0)--(5.0,0) node[right] {$V$};
    \draw[axis] (0,0)--(0,4.3) node[above] {$P$};
    \coordinate (A) at (1.0,1.0);
    \coordinate (B) at (1.0,3.55);
    \coordinate (C) at (4.25,.85);
    \draw[DeepBlue,very thick,->] (A)--(B);
    \draw[DeepBlue,very thick,->] (B) .. controls (2.0,2.95) and (3.0,1.70) .. (C);
    \fill[Ink] (A) circle (1.7pt) node[thermolabel,left] {$(V_0,T_0)$};
    \fill[Ink] (B) circle (1.7pt) node[thermolabel,above left] {$(V_0,T)$};
    \fill[Ink] (C) circle (1.7pt) node[thermolabel,right] {$(V,T)$};
    \node[thermolabel,anchor=west] at (1.15,2.25) {等容升温};
    \node[thermolabel,rotate=-28] at (2.90,2.05) {等温膨胀};
  \end{scope}
  \node[draw=DeepBlue,rounded corners=2pt,thick,align=left,inner sep=7pt] at (4.0,.2)
    {$\displaystyle \Delta S=n c_v\ln\frac{T}{T_0}+nR\ln\frac{V}{V_0}$\\[5pt]
     熵是状态量：可任选一条可逆路径计算。};
\end{tikzpicture}
\end{center}\nopagebreak[2]}

\newcommand{\boltzmannMemorialFigure}{%
\begin{center}
\begin{tikzpicture}[x=1cm,y=1cm]
  \fill[TextGray!28] (-4.6,-2.5) rectangle (4.6,1.65);
  \fill[TextGray!28] (-4.6,1.60) arc[start angle=180,end angle=0,x radius=4.6,y radius=2.0]--cycle;
  \draw[Ink,thick] (-4.6,-2.5)--(-4.6,1.6) arc[start angle=180,end angle=0,x radius=4.6,y radius=2.0]--(4.6,-2.5)--cycle;
  \node[font=\large\bfseries] at (0,2.25) {$S=k_B\ln\Omega$};
  \fill[Ink] (-1.35,-2.48) .. controls (-1.05,-1.35) and (-.78,-.78) .. (-.48,-.55)
    --(.48,-.55) .. controls (.82,-.86) and (1.12,-1.45) .. (1.35,-2.48)--cycle;
  \fill[SoftGray] (0,.10) ellipse (.80 and 1.12);
  \fill[Ink] (-.75,.54) .. controls (-.42,1.36) and (.52,1.36) .. (.78,.48)
    .. controls (.36,.84) and (-.20,.88) .. cycle;
  \draw[Ink,thick] (-.40,.25)--(-.12,.28) (.12,.28)--(.40,.25);
  \fill[Ink] (-.22,.26) circle (.045) (.22,.26) circle (.045);
  \draw[Ink,thin] (0,.15)--(-.05,-.18)--(.06,-.15);
  \fill[SoftGray] (-.55,-.20) .. controls (-.62,-1.02) and (-.30,-1.35) .. (0,-1.52)
    .. controls (.33,-1.35) and (.64,-1.00) .. (.55,-.20) .. controls (.16,-.45) and (-.16,-.45) .. cycle;
  \draw[thermoguide] (-3.8,-1.95)--(-2.6,-.25)--(-3.3,.70);
  \draw[thermoguide] (3.8,-1.95)--(2.6,-.25)--(3.3,.70);
\end{tikzpicture}
\end{center}\nopagebreak[2]}

\newcommand{\carnotProofDiagram}{%
\begin{center}
\begin{tikzpicture}[x=1cm,y=1cm]
  \node[draw=Ink,fill=SoftGray,minimum width=3.0cm,minimum height=.65cm] (H) at (0,2.0) {高温热库 $T_h$};
  \node[draw=DeepBlue,fill=PaleBlue,minimum width=2.2cm,minimum height=.8cm] (E) at (0,0) {可逆热机};
  \node[draw=Ink,fill=SoftGray,minimum width=3.0cm,minimum height=.65cm] (C) at (0,-2.0) {低温热库 $T_c$};
  \draw[thermoarrow] (H)--node[right,thermolabel] {$Q_h$}(E);
  \draw[thermoarrow] (E)--node[right,thermolabel] {$Q_c$}(C);
  \draw[thermoarrow] (E)--(2.25,0) node[right,thermolabel] {$W=Q_h-Q_c$};
  \node[draw=DeepBlue,rounded corners=2pt,align=left,inner sep=6pt] at (5.4,0)
    {$\Delta S_{\rm cyc}=0$\\[2pt]
     $\displaystyle \frac{Q_h}{T_h}-\frac{Q_c}{T_c}=0$\\[4pt]
     $\displaystyle \eta=1-\frac{T_c}{T_h}$};
\end{tikzpicture}
\end{center}\nopagebreak[2]}

\newcommand{\secondLawEntropyDiagram}{%
\begin{center}
\begin{tikzpicture}[x=1cm,y=1cm]
  \node[draw=RuleGray,rounded corners=2pt,minimum width=5.4cm,minimum height=2.25cm,align=center] at (-3.5,0)
    {单一热库取热并全部变成功\\[4pt]$\Delta S_{\rm tot}=-Q/T<0$\\[2pt]\color{DeepBlue}不可能};
  \node[draw=RuleGray,rounded corners=2pt,minimum width=5.4cm,minimum height=2.25cm,align=center] at (3.5,0)
    {无补偿地把热从 $T_c$ 送往 $T_h$\\[4pt]
     $\displaystyle\Delta S_{\rm tot}=\frac{Q}{T_h}-\frac{Q}{T_c}<0$\\[2pt]\color{DeepBlue}不可能};
\end{tikzpicture}
\end{center}\nopagebreak[2]}

\newcommand{\irreversibleProcessesDiagram}{%
\begin{center}
\begin{tikzpicture}[x=1cm,y=1cm]
  \begin{scope}[shift={(-4.3,0)}]
    \draw[thermobody] (-2,-1.2) rectangle (2,1.2);
    \draw[thermobody] (0,-1.2)--(0,1.2);
    \fill[PaleBlue] (-2,-1.2) rectangle (0,1.2);
    \node at (-1,0) {气体};
    \draw[thermoarrow] (-.25,0)--(1.35,0);
    \node[thermolabel] at (0,-1.55) {快速自由膨胀：$\Delta S=nR\ln(V_f/V_i)>0$};
  \end{scope}
  \begin{scope}[shift={(4.3,0)}]
    \draw[thermobody,fill=SoftGray] (-2,-1.1) rectangle (-.35,1.1);
    \draw[thermobody,fill=PaleBlue] (.35,-1.1) rectangle (2,1.1);
    \node at (-1.18,0) {$T_h$};
    \node at (1.18,0) {$T_c$};
    \draw[thermoarrow] (-.2,0)--(.2,0);
    \node[thermolabel] at (0,-1.55) {$\Delta S=Q(1/T_c-1/T_h)>0$};
  \end{scope}
\end{tikzpicture}
\end{center}\nopagebreak[2]}

\newcommand{\threeBodiesDiagram}{%
\begin{center}
\begin{tikzpicture}[x=1cm,y=1cm]
  \foreach \x/\t/\lab in {-4/300/{物体1},0/300/{物体2},4/100/{物体3}}{
    \draw[thermobody,fill=SoftGray] (\x-1.05,-.55) rectangle (\x+1.05,.55);
    \node at (\x,.15) {\lab};
    \node[thermolabel] at (\x,-.25) {$\t\,\mathrm K$};
  }
  \draw[thermoarrow,<->] (-2.85,0)--(-1.15,0);
  \draw[thermoarrow,<->] (1.15,0)--(2.85,0);
  \node[thermolabel,above] at (-2.0,0) {可逆热机};
  \node[thermolabel,above] at (2.0,0) {可逆热机};
  \node[draw=DeepBlue,rounded corners=2pt,align=center,inner sep=6pt] at (0,-2.05)
    {$T_{\max}+2T=700\,\mathrm K$\\[2pt]$T_{\max}T^2=300^2\!\times100\,\mathrm K^3$};
\end{tikzpicture}
\end{center}\nopagebreak[2]}

\newcommand{\surfaceEngineDiagram}{%
\begin{center}
\begin{tikzpicture}[x=1cm,y=1cm]
  \begin{scope}[shift={(-3.4,0)}]
    \draw[axis] (0,0)--(5.0,0) node[right] {$A$};
    \draw[axis] (0,0)--(0,4.0) node[above] {$\sigma$};
    \coordinate (A) at (1.0,1.45);
    \coordinate (B) at (4.1,1.45);
    \coordinate (C) at (4.1,3.05);
    \coordinate (D) at (1.0,3.05);
    \draw[DeepBlue,very thick,postaction={decorate},decoration={markings,
      mark=at position .12 with {\arrow{Latex}},mark=at position .37 with {\arrow{Latex}},
      mark=at position .62 with {\arrow{Latex}},mark=at position .87 with {\arrow{Latex}}}]
      (A)--(B)--(C)--(D)--cycle;
    \node[thermolabel,left] at (0,1.45) {$\sigma_h$};
    \node[thermolabel,left] at (0,3.05) {$\sigma_c$};
    \node[thermolabel,below] at (1.0,0) {$A_1$};
    \node[thermolabel,below] at (4.1,0) {$A_2$};
    \node[thermolabel] at (2.55,2.22) {$W_{\rm out}$};
  \end{scope}
  \node[draw=DeepBlue,rounded corners=2pt,align=left,inner sep=6pt] at (4.4,.3)
    {$S=-A\,\mathrm d\sigma/\mathrm dT$\\[3pt]
     $Q_h=T_h\Delta S$\\[3pt]
     $Q_c=T_c\Delta S$\\[3pt]
     $\eta=1-T_c/T_h$};
\end{tikzpicture}
\end{center}\nopagebreak[2]}

% Shared coordinates for mechanical contact and thermodynamic paths.
\tikzset{revcycle/.style={DeepBlue,thick,postaction={decorate},decoration={markings,mark=at position .55 with {\arrow{Latex}}}}}
\newcommand{\pistonStageRevision}[4]{\begin{scope}[shift={(#1,0)}]
\pgfmathsetmacro{\pinY}{-2+.4*cos(#2)+sqrt(1.5*1.5-.4*.4*sin(#2)*sin(#2))}
\coordinate(pin)at(0,\pinY);\coordinate(crank)at({.4*sin(#2)},{-2+.4*cos(#2)});
\fill[DeepBlue!7](-.62,\pinY+.12)rectangle(.62,.75);
\draw[thick](-.62,-1.2)--(-.62,.75) (.62,-1.2)--(.62,.75) (-.25,.75)--(.25,.75);
\ifnum#4=1 \draw[thick](-.62,.55)--(-.25,.4);\else\draw[thick](-.62,.75)--(-.25,.75);\fi
\ifnum#4=5 \draw[thick](.25,.4)--(.62,.55);\else\draw[thick](.25,.75)--(.62,.75);\fi
\draw[thick,fill=SoftGray](-.62,\pinY-.06)rectangle(.62,\pinY+.12);
\draw[thick](pin)--(crank)--(0,-2);\draw[thick](0,-2)circle(.4);\fill(pin)circle(1.2pt)(crank)circle(1.2pt)(0,-2)circle(1.2pt);
\draw[thick](0,1.03)--(0,.59);\draw[thick](-.09,.64)--(.09,.64);
\ifnum#4=3 \draw[orange!85!black,thick](-.14,.52)--(.05,.38)--(-.06,.24)--(.13,.12);\fi
\ifnum#4=1 \draw[->,DeepBlue,thick](-1.05,.95)--(-.45,.95)--(-.45,.62);\fi
\ifnum#4=5 \draw[->,DeepBlue,thick](.45,.62)--(.45,.95)--(1.05,.95);\fi
\ifnum#4=3\else
\ifnum#4=2 \draw[->](.87,-.8)--(.87,-.1);\else
\ifnum#4=5 \draw[->](.87,-.8)--(.87,-.1);\else\draw[->](.87,-.1)--(.87,-.8);\fi\fi\fi
\node[font=\small]at(0,-2.8){#3};\end{scope}}
\renewcommand{\ottoCycleFigure}{\begin{center}\begin{tikzpicture}[font=\small]
\pistonStageRevision{0}{90}{进气（口开）}{1}
\pistonStageRevision{2.65}{270}{压缩}{2}
\pistonStageRevision{5.3}{0}{点火}{3}
\pistonStageRevision{7.95}{90}{做功}{4}
\pistonStageRevision{10.6}{270}{排气（口开）}{5}
\node[anchor=west]at(6.5,1.45){火花塞};\draw[TextGray](6.35,1.42)--(5.3,1.03);
\node[font=\scriptsize]at(5.3,-3.35){相同缸径、曲柄半径与连杆长度；同一位置可属于不同冲程}
;\end{tikzpicture}
\par\medskip\begin{tikzpicture}[x=1.6cm,y=.9cm,font=\small]
\pgfmathsetmacro{\pa}{1.8*pow(1.2/3,1.4)}\pgfmathsetmacro{\pd}{4*pow(1.2/3,1.4)}
\draw[axis](0,0)--(4,0)node[right]{$V$};\draw[axis](0,0)--(0,4.8)node[above]{$P$};
\draw[revcycle,domain=3:1.2,samples=90]plot(\x,{1.8*pow(1.2/\x,1.4)});
\draw[revcycle](1.2,1.8)--(1.2,4);
\draw[revcycle,domain=1.2:3,samples=90]plot(\x,{4*pow(1.2/\x,1.4)});
\draw[revcycle](3,\pd)--(3,\pa);
\draw[guide](1.2,0)--(1.2,1.65) (3,0)--(3,\pa);
\draw[TextGray,dashed,->](1.2,\pa)--(2.5,\pa);\draw[TextGray,->](2.4,\pa-.12)--(1.6,\pa-.12);
\foreach \x/\y/\n/\a in {3/\pa/A/below right,1.2/1.8/B/left,1.2/4/C/above,3/\pd/D/above right,1.2/\pa/O/left}{\fill(\x,\y)circle(1.2pt)node[\a]{$\n$};}
\node[below]at(1.2,0){$V_2$};\node[below]at(3,0){$V_1$};
\node[anchor=west,align=left]at(4.6,3.2){蓝线：封闭空气标准循环\\[6pt]$A\to B$、$C\to D$：绝热\\[6pt]$B\to C$：等容吸热 $Q_h$\\[6pt]$D\to A$：等容放热 $Q_c$};
\node[anchor=west,align=left,font=\scriptsize]at(4.6,.65){灰色往返线仅提示真实进、排气\\不是蓝色闭式循环的一部分};
\end{tikzpicture}\end{center}\nopagebreak[2]}

\renewcommand{\carnotCycleFigure}{\begin{center}\begin{tikzpicture}[x=1.45cm,y=1cm,font=\small]
\pgfmathsetmacro{\cold}{4/pow(1.5,2/3)}
\draw[axis](0,0)--(4.4,0)node[right]{$V$};\draw[axis](0,0)--(0,4.9)node[above]{$P$};
\draw[TextGray,domain=.9:4.15,samples=100]plot(\x,{4/\x});
\draw[TextGray,domain=.9:4.15,samples=100]plot(\x,{\cold/\x});
\draw[revcycle,domain=1:2.4,samples=85]plot(\x,{4/\x});
\draw[revcycle,domain=2.4:3.6,samples=85]plot(\x,{4*pow(2.4,2/3)/pow(\x,5/3)});
\draw[revcycle,domain=3.6:1.5,samples=85]plot(\x,{\cold/\x});
\draw[revcycle,domain=1.5:1,samples=85]plot(\x,{4/pow(\x,5/3)});
\foreach \x/\y/\n/\a in {1/4/1/above left,2.4/{4/2.4}/2/above right,3.6/{\cold/3.6}/3/below,1.5/{\cold/1.5}/4/below left}{\fill(\x,\y)circle(1.4pt)node[\a]{$\n$};}
\node[anchor=west,align=left]at(5,3.4){高温等温：$PV=nRT_1$\\[5pt]$1\to2$ 吸热 $Q_1$};
\node[anchor=west,align=left]at(5,1.6){低温等温：$PV=nRT_2$\\[5pt]$3\to4$ 放热 $Q_2$};
\draw[TextGray](4.1,{4/4.1})--(4.55,3.4)--(4.8,3.4);
\draw[TextGray](4.1,{\cold/4.1})--(4.65,1.6)--(4.8,1.6);
\end{tikzpicture}\end{center}\nopagebreak[2]}

\renewcommand{\heatdevicepanel}[5]{\begin{scope}[shift={(#1,0)}]
\draw[thermobody,fill=SoftGray](-1.4,2.15)rectangle(1.4,2.62);\draw[thermobody,fill=SoftGray](-1.4,-2.62)rectangle(1.4,-2.15);
\node[thermolabel]at(0,2.38){高温热库 $T_h$};\node[thermolabel]at(0,-2.38){低温热库 $T_c$};
\draw[thermobody](-.66,-.65)rectangle(.66,.65);\node[thermolabel,align=center,text width=1.1cm]at(0,0){#2};#3 #4 #5\end{scope}}
\renewcommand{\foundersPortraitFigure}{\begin{center}\includegraphics[width=.96\linewidth]{thermodynamics-founders.png}\par\smallskip
\makebox[.31\linewidth]{萨迪·卡诺}\makebox[.31\linewidth]{开尔文}\makebox[.31\linewidth]{鲁道夫·克劳修斯}\end{center}\nopagebreak[2]}
\renewcommand{\boltzmannMemorialFigure}{\begin{center}\begin{tikzpicture}
\node[inner sep=0]at(0,0){\includegraphics[width=6.2cm]{boltzmann-memorial.png}};
\node[font=\large]at(0,2.32){$S=k_B\ln\Omega$};\end{tikzpicture}\end{center}\nopagebreak[2]}

\renewcommand{\entropyPathFigure}{\begin{center}\begin{tikzpicture}[font=\small]
\draw[axis](0,0)--(5,0)node[right]{$V$};\draw[axis](0,0)--(0,4.8)node[above]{$P$};
\draw[revcycle](1,1.6)--(1,4);\draw[revcycle,domain=1:4,samples=90]plot(\x,{4/\x});
\foreach \x/\y/\n/\a in {1/1.6/A/left,1/4/B/above,4/1/C/right}{\fill(\x,\y)circle(1.5pt)node[\a]{$\n$};}
\node[anchor=west,align=left]at(6,3.45){$A:(V_0,T_0)$\\[5pt]$B:(V_0,T)$\\[5pt]$C:(V,T)$};
\node[anchor=west,align=left]at(6,1.25){$A\to B$：等容\\[5pt]$B\to C$：等温，$PV=nRT$};
\end{tikzpicture}\par\smallskip
$\displaystyle\Delta S=n c_v\ln\frac{T}{T_0}+nR\ln\frac{V}{V_0}$
\end{center}\nopagebreak[2]}

\renewcommand{\carnotProofDiagram}{\begin{center}\begin{tikzpicture}[font=\small]
\node[draw,minimum width=3cm,minimum height=.6cm](H)at(0,2){高温热库 $T_1$};
\node[draw,minimum width=2.2cm,minimum height=.8cm](E)at(0,0){可逆热机};
\node[draw,minimum width=3cm,minimum height=.6cm](C)at(0,-2){低温热库 $T_2$};
\draw[thermoarrow](H)--node[right]{$Q_1$}(E);\draw[thermoarrow](E)--node[right]{$Q_2$}(C);
\draw[thermoarrow](E.east)--(3.1,0);\node[above]at(2.2,.15){$W$};
\node[anchor=west,align=left]at(4,0){$W=Q_1-Q_2$\\[8pt]$Q_1/T_1=Q_2/T_2$\\[8pt]$\eta=1-T_2/T_1$};
\end{tikzpicture}\end{center}\nopagebreak[2]}

\renewcommand{\irreversibleProcessesDiagram}{\begin{center}\begin{tikzpicture}[font=\small]
\foreach \sx/\phase in {0/0,5/1}{\begin{scope}[shift={(\sx,0)}]
\ifnum\phase=0 \fill[PaleBlue](-1.6,-1)rectangle(0,1);\draw[thick](0,-1)--(0,1);\node at(-.8,0){气体};\node at(.8,0){真空};
\else\fill[PaleBlue](-1.6,-1)rectangle(1.6,1);\draw[guide](0,-1)--(0,1);\node at(.8,0){气体};\fi
\draw[thick](-1.6,-1)rectangle(1.6,1);\end{scope}}
\draw[->,DeepBlue,thick](1.85,0)--(3.15,0)node[midway,above]{撤去隔板};
\node at(2.5,-1.55){绝热自由膨胀：$Q=W=0$，$\Delta S=nR\ln(V_f/V_i)>0$};
\draw[thick](9,-1)rectangle(10.5,1);\draw[thick](11.3,-1)rectangle(12.8,1);
\node at(9.75,0){$T_h$};\node at(12.05,0){$T_c$};\draw[thermoarrow](10.6,0)--(11.2,0);
\node[align=center]at(10.9,-1.6){两恒温热库间传热\\$\Delta S=Q(1/T_c-1/T_h)>0$};
\end{tikzpicture}\end{center}\nopagebreak[2]}

\renewcommand{\threeBodiesDiagram}{\begin{center}\begin{tikzpicture}[font=\small]
\node[draw,minimum width=2.7cm,minimum height=.7cm](Y)at(0,2.2){供热体 $Y$};
\node[draw,minimum width=2.7cm,minimum height=.7cm](X)at(6,2.2){目标体 $X$};
\node[draw,minimum width=2.7cm,minimum height=.7cm](E)at(0,0){可逆热机};
\node[draw,minimum width=2.7cm,minimum height=.7cm](P)at(6,0){可逆热泵};
\node[draw,minimum width=3cm,minimum height=.7cm](Z)at(3,-2.4){低温体 $Z$};
\draw[thermoarrow](Y)--node[left]{$-c\dd Y$}(E);\draw[thermoarrow](P)--node[right]{$c\dd X$}(X);
\draw[thermoarrow](E.east)--node[above]{$\delta W$}(P.west);
\draw[thermoarrow](E.south)--(0,-2.4)--(Z.west);
\draw[thermoarrow](Z.east)--(6,-2.4)--(P.south);
\node[left,align=right]at(-.1,-1.3){热机向 $Z$ 放热};\node[right,align=left]at(6.1,-1.3){热泵从 $Z$ 吸热};
\node at(3,-3.4){初态 $(X,Y,Z)=(300,300,100)\,\mathrm K$；极限末态 $(400,150,150)\,\mathrm K$};
\end{tikzpicture}\end{center}\nopagebreak[2]}

\renewcommand{\surfaceEngineDiagram}{\begin{center}\begin{tikzpicture}[x=1.5cm,y=.9cm,font=\small]
\draw[axis](0,0)--(3.35,0)node[right]{$A$};\draw[axis](0,0)--(0,4.9)node[above]{$\sigma$};
\fill[PaleBlue](.8,1)--(1.4,1)--plot[domain=1.4:2.8,samples=70](\x,{5-pow(2.8/\x,2)})--(1.6,4)--plot[domain=1.6:.8,samples=70](\x,{5-pow(1.6/\x,2)})--cycle;
\draw[revcycle](.8,1)--(1.4,1);\draw[revcycle,domain=1.4:2.8,samples=70]plot(\x,{5-pow(2.8/\x,2)});
\draw[revcycle](2.8,4)--(1.6,4);\draw[revcycle,domain=1.6:.8,samples=70]plot(\x,{5-pow(1.6/\x,2)});
\foreach \x/\y/\n/\a in {.8/1/1/left,1.4/1/2/below right,2.8/4/3/right,1.6/4/4/above left}{\fill(\x,\y)circle(1.4pt)node[\a]{$\n$};}
\node[left]at(0,1){$\sigma_h$};\node[left]at(0,4){$\sigma_c$};
\node[anchor=west,align=left]at(4.2,3.4){$\sigma(T)=\sigma_*-aT^2$\\[6pt]等熵线：$AT=\text{常量}$\\[6pt]$1\to2$：高温扩面\\[6pt]$3\to4$：低温缩面};
\node[anchor=west]at(4.2,1){$W_{\rm out}=-\oint\sigma\dd A>0$};
\end{tikzpicture}\end{center}\nopagebreak[2]}


从某种意义上讲，热力学的发展与人们制造和改进热机的努力是分不开的。最初，人们设想制造一种
不需要能量输入却能够持续工作的热机，但热力学第一定律告诉我们，这样的努力没有任何希望。退而求
其次，人们又希望得到效率为百分之百的热机，即把吸收的热量全部转化为机械功。从能量守恒的角度看，
这种设想并不直接违背热力学第一定律：能量输入与输出可以相等，但这还没有回答过程是否能够发生。
第二定律额外指出，循环装置不能以“只从单一热库取热并把它全部变成功”为唯一效果。
常见气体热机确有膨胀做功、压缩复原和向外放热的过程，但这只能帮助理解具体模型，不能证明一般定律。
例如液体表面缩小时也能对外做功，后面将研究这种不同的工质。

人们在追求尽可能高效的热机时，逐渐发现这一事实背后隐藏着一条普遍的热力学原理，这就是今天的
热力学第二定律。热机、可逆性以及熵，将是本讲讨论的主题。

\section{热力学循环的效率}

热机在工作时循环运行：工质从某个状态出发，经历一系列过程后再次回到初始状态，完成一个热力学循环。
对在两个热库之间工作的普通热机，工质从高温热库吸热，还向低温热库放热。把整个循环
从外界吸收的总热量记为 $Q^{\rm in}$，把向外排出的总热量记为 $Q^{\rm out}$。由于循环结束后工质回到
初态，内能改变量为零。这里 $Q^{\rm in},Q^{\rm out}$ 分别是累计吸热、累计放热的正量，
也记为 $Q_h,Q_c$；净输入热量为二者之差，不沿用上一章单段过程中带符号的 $Q^{\rm in}=Q$。
热力学第一定律给出循环的净功
\begin{equation}
  \boxed{W^{\rm out}=Q^{\rm in}-Q^{\rm out}}.
\end{equation}
若工质仅以体积功与外界交换机械能，且过程可逆，则净功为 $\oint P\dd V$，等于 $P$--$V$ 图上
围成的有向面积，顺时针为正。对于电功、表面功等形式，须使用相应的功表达式，不能一概用 $P$--$V$ 面积。
只要排出的热量不为零，热机就没有把吸收的全部热量变成功。定义热机
效率为对外净功与吸热量的比值：
\begin{equation}
  \boxed{\eta=\frac{W^{\rm out}}{Q^{\rm in}}
  =1-\frac{Q^{\rm out}}{Q^{\rm in}}}.
\end{equation}

最常见的内燃机可用奥托热机近似描述。图~\ref{ch16:fig:s30-21-1}把活塞的工作过程与相应的 $P$--$V$ 循环放在一起。
从活塞运动看，一个循环可分为进气、压缩、点火、做功和排气几个环节。首先关闭排气口、打开进气口，
空气和汽油的混合气体被吸入气缸；自然吸气发动机的进气压强接近一个大气压，涡轮增压发动机则更高。
活塞到达最下端后进气口关闭，活塞因惯性向上运动并压缩混合气体。压缩过程通常很快，虽然气体温度
升高，但与气缸的热量交换可以忽略，因而可近似为绝热压缩。气体被完全压缩后，火花塞点燃高压油气
混合物，燃烧很快，活塞来不及移动，气体体积近似不变而温度、压强急剧升高。随后高温高压气体推动
活塞向下运动，进入做功冲程；这一快速膨胀过程也可近似为绝热膨胀。活塞到达最下端后排气口打开，
废气被排出，装置完成一个循环。

\par\noindent\begin{minipage}{\linewidth}
\ottoCycleFigure
\diagramcaption{上方是四冲程发动机的五个动作阶段，下方是空气标准奥托循环。点火单列显示，不是额外增加一个活塞冲程。}\label{ch16:fig:s30-21-1}
\end{minipage}\par\medskip

最小、最大气缸体积分别为 $V_2,V_1$。先看真实进排气：$O\to A$ 提示吸气，压强近似不变，
但缸内气体量不断增加，并非定量气体等压膨胀；$A\to O$ 提示排气，气体不断离开气缸。
在真实做功冲程后打开排气门时，废气也会先因内外压差排出。这些开放过程伴随物质携能，
不能套用一份封闭气体的 $nc_v\Delta T$ 来计算排出能量。

为了求一个便于比较的理论效率，现在换用\textbf{空气标准模型}：用同一份定量、定组成的理想气体
代替不断更新的混合气；用可逆等容吸热代替燃烧，用可逆等容放热代替排气散热的效果。
蓝色闭式循环只有四段：$A\to B$ 绝热压缩，$B\to C$ 等容吸热，$C\to D$ 绝热膨胀，
$D\to A$ 等容放热。灰色 $O\leftrightarrow A$ 不是这份封闭工质实际走过的两段。

在这一理想化循环中，只有 $B\to C$ 的等容过程吸热。设气缸内气体可视为总物质的量为 $n$、定容
摩尔热容 $c_v$ 在整个温区内近似恒定的理想气体，并忽略摩擦、漏热和燃烧成分变化，则
\[
  Q_h=n c_v(T_C-T_B).
\]
类似地，只有 $D\to A$ 的等容过程向外排热，排热量为
\[
  Q_c=n c_v(T_D-T_A).
\]
所以
\[
  \eta=1-\frac{Q_c}{Q_h}
  =1-\frac{T_D-T_A}{T_C-T_B}.
\]
再利用两段绝热过程
\[
  T_A V_1^{\gamma-1}=T_B V_2^{\gamma-1},
  \qquad
  T_D V_1^{\gamma-1}=T_C V_2^{\gamma-1},
\]
令压缩比 $r=V_1/V_2>1$，两式分别给出 $T_B=r^{\gamma-1}T_A$、
$T_C=r^{\gamma-1}T_D$。相减而不是直接约去温度，可得
\[
 T_C-T_B=r^{\gamma-1}(T_D-T_A),\qquad
 \frac{Q_c}{Q_h}=\frac{T_D-T_A}{T_C-T_B}=r^{1-\gamma}.
\]
于是得到奥托循环效率与发动机几何参数的关系
\begin{equation}
  \boxed{\eta=1-\left(\frac{V_2}{V_1}\right)^{\gamma-1}}.
\end{equation}
这说明在理想模型中，压缩比 $V_1/V_2$ 越大，循环效率越高；真实发动机还会受到爆震、传热、摩擦和
材料强度等条件限制。

\section{卡诺循环}

在所有热力学循环中，对理论发展最重要的是卡诺设计的卡诺热机。它来自人们寻找最高可能效率
热机的努力；当理论上确定这种热机的效率以后，人们才认识到其背后隐藏着一条基本的热力学定律。
以下示范理想气体循环的计算时，仍取定量气体、恒定 $c_v,c_p$；卡诺定理本身并不限于这一模型。

\par\noindent\begin{minipage}{\linewidth}
\carnotCycleFigure
\diagramcaption{卡诺循环的 $P$--$V$ 图。$1\to2$、$3\to4$ 为等温过程，$2\to3$、$4\to1$ 为绝热过程。}\label{ch16:fig:s30-21-2}
\end{minipage}\par\medskip

如图~\ref{ch16:fig:s30-21-2}所示，卡诺循环由两段绝热过程和两段等温过程组成，两条等温线的温度分别为 $T_1>T_2$。
正是两个热库之间存在温差，才有可能把一部分热能转化为机械功。循环从状态1出发：工质与温度为 $T_1$
的高温热库接触，持续吸热并作等温膨胀，工质对外做功；到达状态2后与热库脱离，继续绝热膨胀，直到
温度降到 $T_2$，来到状态3；然后工质与低温热库接触，受到外界做功而作等温压缩，同时向低温热库
排出热量 $Q_2$；到达状态4后再与热库脱离，作绝热压缩，最后恰好回到初态1。

卡诺可以证明，无论采用何种工作物质，可逆卡诺循环都具有相同效率。先以理想气体为例。设热机
内有 $n$ 摩尔理想气体，高温等温膨胀吸入的热量与低温等温压缩排出的热量分别为
\[
  Q_1=nRT_1\ln\frac{V_2}{V_1},
  \qquad
  Q_2=nRT_2\ln\frac{V_3}{V_4}.
\]
两段绝热过程满足
\[
  T_1V_2^{\gamma-1}=T_2V_3^{\gamma-1},
  \qquad
  T_1V_1^{\gamma-1}=T_2V_4^{\gamma-1}.
\]
两式相除得到 $V_2/V_1=V_3/V_4$。于是效率表达式中所有体积因子恰好抵消，最终得到
\begin{equation}
  \boxed{\eta_{\rm Carnot}=1-\frac{T_2}{T_1}}.
\end{equation}
令人惊讶的是，卡诺热机的效率只与两个热库的热力学温度有关，与热机的几何结构无关；这与奥托
热机效率依赖压缩比形成鲜明对比。

\section{热力学第二定律的早期陈述形式}

卡诺关于热机效率的研究使人们逐渐意识到，有一条基本规律阻止人们制造效率为 $100\%$ 的热机。
这种机器并不直接违反热力学第一定律，却从未能够实现。图~\ref{ch16:fig:s30-21-3}(a)表示真实热机：它从高温热库吸收
$Q_h$，对外做功 $W=Q_h-Q_c$，并向低温热库排出 $Q_c$。图~\ref{ch16:fig:s30-21-3}(b)所示的“完美”热机则只从单一
热库取热并把它全部变成功。开尔文勋爵彻底断绝了这种幻想，他指出：

\begin{quote}\small
不存在这样的装置：当它完成一个循环时，能从一个物体取出热量并把它转换为功，而不发生其他变化
（至少要在其他物体中发生变化）。
\end{quote}

历史上这种机器被称为第二类永动机，因此热力学第二定律的开尔文表述也可理解为：第二类永动机不可能
制成。

\par\noindent\begin{minipage}{\linewidth}
\heatDevicesFigure
\diagramcaption{热机和制冷机能量流。由左至右为普通热机、第二类永动机、普通制冷机和无功制冷机；两种“完美”设想不允许实现。}\label{ch16:fig:s30-21-3}
\end{minipage}\par\medskip

生活经验还指出，热量会自发地从高温物体传向低温物体。冰箱、空调试图把热量从低温部分搬运到高温
环境，必须依赖外部能量供应。在力学中，物体可以凭惯性短时间沿受力的反方向运动，例如竖直上抛物体
仍继续上升，或弹簧振子在回复力作用下仍远离平衡位置；不能据此类比出一种可持续的无补偿逆向传热。
在这里讨论的通常宏观传热中，若不付出其他代价，热量只能自发地从高温传向低温。克劳修斯把这一事实表述为：

\begin{quote}\small
不存在这样的装置：当它完成一个循环时，能使热量从较低温度的热库传到较高温度的热库，而没有补偿
变化。
\end{quote}

从高温向低温的传热过程可以自然发生，而它的逆过程不能独立发生，因此克劳修斯表述也意味着：有限
温差下的热传导是一种不可逆过程。开尔文与克劳修斯的经典英文表述分别为：
\begin{quote}\footnotesize
It is impossible, by means of inanimate material agency, to derive mechanical effect from any portion of
matter by cooling it below the temperature of the coldest of the surrounding objects.

Heat can never pass from a colder to a warmer body without some other change, connected therewith,
occurring at the same time.
\end{quote}

\par\noindent\begin{minipage}{\linewidth}
\foundersPortraitFigure
\diagramcaption{三位热力学奠基人：卡诺研究热机的理想效率，开尔文和克劳修斯提出第二定律的重要表述。}\label{ch16:fig:s30-21-4}
\end{minipage}\par\medskip
图~\ref{ch16:fig:s30-21-4}按卡诺、开尔文、克劳修斯的顺序排列，分别对应前述热机研究及两种早期陈述。

\section{熵与热力学第二定律}

在得到热力学第二定律的早期陈述以后，人们逐渐认识到：对于任意热力学系统，除压强、密度、温度、
内能等状态变量外，还存在另一个状态变量——熵（entropy），记为 $S$。熵从出现之初就带有神秘色彩，
尤其是由熵增出发对宇宙结局的推测；正因为如此，更应准确理解熵的定义和适用条件，避免在不正确的
条件下滥用相关定律。

熵首先是状态变量。与内能一样，它完全由热力学系统的状态决定，而与系统以什么过程到达该状态无关。
状态变量必然同组成系统的微观粒子的性质有关：压强是粒子碰撞容器表面的宏观效果；在经典图像中，
温度由粒子的平均能量决定；内能可理解为全部微观粒子的能量总和。熵也由微观状态决定，玻尔兹曼后来
揭示了这一联系。

当热力学系统经历某个过程时，熵可能改变。对于可逆过程中的任意状态，都可以定义系统温度。若系统在
温度为 $T$ 的平衡态吸收一小份热量 $\delta Q_{\rm rev}$，且热量足够小，使温度在这一微元过程中可视
为不变，则熵的改变为
\begin{equation}
  \boxed{\dd S=\frac{\delta Q_{\rm rev}}{T}}.
\end{equation}
可逆吸热时熵增加，可逆放热时熵减小；相同热量在较低温度下传递，引起的熵变更大。可逆绝热过程没有
热量交换，因而也是等熵过程。与内能相似，经典热力学首先给出的是熵在过程中的改变，而不是任意系统
熵的绝对数值。

不可逆过程同样会导致熵变。由于熵是状态变量，虽然不能把不可逆路径上的每一中间状态都当作平衡态，
仍可另选一条连接同一初、末态的可逆路径来计算 $\Delta S$。对孤立系统，任意自发过程只可能使熵增加；
在理想可逆极限中熵保持不变，孤立系统中的过程不能使总熵减小。这就是熵增原理。其直接推论是：孤立
系统的热平衡态是熵最大的状态。

在熵概念提出的早期，克劳修斯曾把这一结论直接应用于整个宇宙，设想宇宙最终达到处处均匀、温度相同
的热平衡状态，所有细节全部消失。今天我们知道，不能把一个简单孤立热力学系统的结论不加条件地套用
到演化的宇宙；宇宙中丰富结构的形成还受到引力、膨胀和非平衡过程等因素控制。

\subsection{准静态为什么还不够}

准静态只表示变化足够缓慢，使系统能用一串近似平衡态描述；可逆还要求系统和环境都能在不留下
其他变化的条件下复原。例如，缓慢拖动有摩擦的活塞可以近似准静态，但摩擦仍把机械能转为内能，
单纯反向拖动不能把这些内能全部恢复为原来的机械能。有限温差传热也有同样的问题。
因此，可逆传热须令两侧温差趋于零，可逆体积功须令内外压强差趋于零，并排除摩擦等耗散。
这是一种极限模型，不是要求实际装置在有限时间内严格达到它。

\subsection{熵流与熵产生}

对不交换物质的系统，热量可以从不同温度的环境输入。设第 $j$ 个恒温热库温度为 $T_j$，
在全过程向系统传入的带符号热量为 $Q_j$。热库自身的熵变为 $-Q_j/T_j$。
把系统和这些热库一起考察，第二定律要求总熵变非负，因而
\begin{equation}
 \Delta S_{\rm sys}-\sum_j\frac{Q_j}{T_j}=S_{\rm gen}\ge0.
\end{equation}
这里 $S_{\rm gen}$ 是系统连同有关传热边界的总熵产生；它不是新的能量。
等号对应整个过程可逆。如果环境温度变化，应把 $Q_j/T_j$ 换成相应的积分，不能随意选一个平均温度。
对循环，$\Delta S_{\rm sys}=0$，于是
\begin{equation}
 \sum_j\frac{Q_j}{T_j}\le0.
\end{equation}
这就是以恒温热库分段表示的克劳修斯不等式。循环可以把吸热和放热的熵交换相互抵消，
却不能用负的熵产生去抵消耗散。

例如，两热库间直接传递热量 $Q>0$，高温为 $T_h$、低温为 $T_c<T_h$。
高温热库熵变 $-Q/T_h$，低温热库熵变 $Q/T_c$，故
\[
 S_{\rm gen}=Q\left(\frac1{T_c}-\frac1{T_h}\right)>0.
\]
总能量没有损失，但可以利用的温差减小了。若 $T_h\to T_c$，单位热量引起的熵产生趋于零，
这正是可逆传热必须取无穷小温差的原因。

\subsection{由熵不等式判断热机与制冷机}

在 $T_h>T_c>0$ 的两个热库之间，热机每循环从高温热库取热 $Q_h>0$，向低温热库排热 $Q_c>0$。
工质循环复原，环境总熵变为
\[
 -\frac{Q_h}{T_h}+\frac{Q_c}{T_c}\ge0
 \quad\Longrightarrow\quad
 \frac{Q_c}{Q_h}\ge\frac{T_c}{T_h}.
\]
结合 $W=Q_h-Q_c$，得到
\begin{equation}
 \eta=1-\frac{Q_c}{Q_h}\le1-\frac{T_c}{T_h}.
\end{equation}
这次推导没有使用气体物态方程，所以卡诺效率是任意工质的上限，而不只是理想气体计算的巧合。
若装置的某一部分不可逆，通常有严格小于号。

制冷机的有用效果是从低温热库取出 $Q_c$，外界输入功 $W>0$，向高温热库排出
$Q_h=Q_c+W$。此时热库熵变满足
\[
 \frac{Q_h}{T_h}-\frac{Q_c}{T_c}\ge0
 \quad\Longrightarrow\quad
 W\ge Q_c\frac{T_h-T_c}{T_c}.
\]
制冷系数和供热系数分别为
\begin{equation}
 \varepsilon_{\rm ref}=\frac{Q_c}{W}\le\frac{T_c}{T_h-T_c},\qquad
 \varepsilon_{\rm hp}=\frac{Q_h}{W}=1+\varepsilon_{\rm ref}
 \le\frac{T_h}{T_h-T_c}.
\end{equation}
二者可以大于1，因为它们比较的是搬运热量与输入功，而不是把输入热量转为功的比例。
例如 $T_h=300\,\mathrm K,T_c=270\,\mathrm K$ 时，可逆制冷系数为9；搬走
$900\,\mathrm J$ 的热量至少需要 $100\,\mathrm J$ 的功，向高温环境排热为
$1000\,\mathrm J$。由 $1000/300=900/270$ 可同时核验熵平衡。


\section{具有简单物态方程系统的熵}

在绝大多数情形下，热力学系统的熵不容易直接求出；但对物态较简单的系统，某些过程中的熵变可以精确
或近似计算。最简单的系统之一是热库。它有极大的热容量，无论吸收还是放出所研究的有限热量，其温度
都可以近似认为不变。因此，温度为 $T$ 的热库吸收热量 $Q$ 时熵变就是 $Q/T$。

另一类系统以固体和液体为代表。通常吸、放热会引起温度变化，虽然物质体积也随温度改变，但固体、液体
的体积变化往往很小。若忽略体积变化，就同时忽略了系统与环境通过体积功发生的能量交换。设系统比热
为 $c$，并在所研究温区内近似为常数。质量为 $m$ 的物质在温度为 $T$ 时可逆吸收
$\dd Q=mc\,\dd T$，因而
\[
  \dd S=\frac{\dd Q}{T}=mc\frac{\dd T}{T}.
\]
从参考温度 $T_0$ 积分到任意温度 $T$，得到
\begin{equation}
  \boxed{S(T)-S(T_0)=mc\ln\frac{T}{T_0}}.
\end{equation}

理想气体的熵不仅依赖温度，还依赖体积。设有 $n$ 摩尔理想气体，参考状态为 $(V_0,T_0)$，要计算任意
状态 $(V,T)$ 相对参考状态的熵变，先假设 $c_v$ 在温区内为常数。由于熵是状态量，可以任意设计一条连接两状态的可逆路径。图~\ref{ch16:fig:s30-21-5}
选择两步路径：先作等容过程 $(V_0,T_0)\to(V_0,T)$，再作等温过程 $(V_0,T)\to(V,T)$。第一步熵变为
\[
  S(V_0,T)-S(V_0,T_0)=n c_v\ln\frac{T}{T_0};
\]
第二步中理想气体的吸热量等于对外功，故
\[
  S(V,T)-S(V_0,T)=nR\ln\frac{V}{V_0}.
\]
两式相加即得
\begin{equation}
  \boxed{S(V,T)-S(V_0,T_0)
  =n c_v\ln\frac{T}{T_0}+nR\ln\frac{V}{V_0}}.
\end{equation}

\par\noindent\begin{minipage}{\linewidth}
\entropyPathFigure
\diagramcaption{理想气体任意两状态间的熵差。图中用等容升温加等温膨胀的可逆路径代替任意真实路径。}\label{ch16:fig:s30-21-5}
\end{minipage}\par\medskip
图~\ref{ch16:fig:s30-21-5}的 $A\to B$ 改变温度而不改体积，$B\to C$ 改变体积而不改温度；因此两个对数项各有清楚的来源。
若实际过程是降温或压缩，相应积分上限小于下限，对数自动给出负号，并不要求实际过程也按图运行。
若 $c_v$ 随温度变化，应保留积分
\[
 \Delta S=n\int_{T_0}^{T}\frac{c_v(T')}{T'}\,\dd T'+nR\ln\frac{V}{V_0}.
\]

\subsection{熵变用初末态计算，熵产生还要看实际环境}
选一条可逆辅助路径，是为了计算状态函数的差，不是把真实过程宣称为可逆。图\ref{ch16:fig:s30-21-5}中的 $A\to B\to C$ 正是这种计算工具。系统熵变、热库熵变和二者的总熵变须分开；系统自己可以减熵，只要包括有关环境后的总熵不减。

\Needspace{180pt}
\begin{origexamplebox}{同末态的两种膨胀，熵产生为何不同}\label{ch16:ex:teaching-two-paths}
$n$ 摩尔理想气体初态为 $(T_0,V_0)$。过程甲是在 $T_0$ 热库作用下可逆等温膨胀至 $2V_0$；过程乙是在隔热刚性容器中向真空自由膨胀至 $2V_0$。求气体、有关环境的熵变和总熵产生。忽略容器热容量、整体动能及势能的变化。

\ThQTwoPaths
\end{origexamplebox}
\Needspace{65pt}
\begin{xesolution}
\ThATwoPaths

乙过程既不传热也不向外做功，内能不变，理想气体末温仍为 $T_0$。所以甲、乙的气体初末态完全相同，均有
\[
 \Delta S_{\rm gas}=nc_v\ln\frac{T_0}{T_0}+nR\ln2=nR\ln2.
\]
甲从热库吸热 $Q=nRT_0\ln2$，热库熵变 $-Q/T_0=-nR\ln2$。接收功的理想机械装置不增加熵，故甲的总熵产生为零。
乙没有向环境传热，环境熵变为零，因此
\[
 S_{{\rm gen},\text{甲}}=0,\qquad
 S_{{\rm gen},\text{乙}}=nR\ln2>0.
\]
乙中 $Q=0$ 却有 $\Delta S_{\rm gas}>0$，并不违反熵定义，因为 $\Delta S=\int\delta Q_{\rm rev}/T$ 中的热量属于辅助可逆路径，而非真实自由膨胀的热量。若末体积趋于初体积，两者熵变都趋零；$nR$ 的单位为 $\mathrm{J/K}$，对数中的体积比无量纲。
\end{xesolution}

\subsection{有限热源不能始终当作恒温热库}
实际物体一边供热一边降温，热机每一小步可达到的效率也在改变。因此不能选一个平均温度代入卡诺效率再乘总热量。解决有限热源问题的顺序是：先明确装置是否循环复原，再写总能量和总熵条件，最后判断上界能否由可逆过程逼近。

\Needspace{180pt}
\begin{origexamplebox}{两物体之间能够取出的最大功}\label{ch16:ex:teaching-finite-work}
两个物体的恒定热容量为 $C_1,C_2$，初温为 $T_h>T_c>0$。在它们之间运行热机，直到两者达到共同终温，热机最后恢复原态。整体与外界不交换热量，但可以输出机械功；忽略体积变化和其他能量改变。求最多能输出多少功，并与两物体直接接触比较。

\ThQFiniteWork
\end{origexamplebox}
\Needspace{65pt}
\begin{xesolution}
\ThAFiniteWork

若共同终温为 $T_f$，热机不储存净能量，输出功为
\[
 W=C_1T_h+C_2T_c-(C_1+C_2)T_f.
\]
能量零点抵消，只涉及温度差。总熵不能减少：
\[
 C_1\ln\frac{T_f}{T_h}+C_2\ln\frac{T_f}{T_c}\ge0.
\]
把第一项的对数拆成 $\ln(T_f/T_c)-\ln(T_h/T_c)$，得到
\[
 (C_1+C_2)\ln\frac{T_f}{T_c}\ge C_1\ln\frac{T_h}{T_c},\qquad
 T_f\ge T_*=T_c\left(\frac{T_h}{T_c}\right)^{C_1/(C_1+C_2)}.
\]
所有对数和幂函数的自变量均无量纲。由于 $W$ 随 $T_f$ 增加而减小，最大功为
\[
 W_{\max}=C_1(T_h-T_*)+C_2(T_c-T_*).
\]
可在两物体的瞬时温度之间运行无穷小可逆循环，使总熵保持不变，直到温差趋零，所以该上界可以在可逆极限下逼近。图\ref{ch16:fig:s30-21-3}中的普通热机能流方向仍适用，但两端温度在不断变化。
直接接触时不输出功，终温是加权算术平均 $T_d=(C_1T_h+C_2T_c)/(C_1+C_2)\ge T_*$。若 $C_1=C_2=C$，则
\[
 T_*=\sqrt{T_hT_c},\qquad W_{\max}=C(\sqrt{T_h}-\sqrt{T_c})^2.
\]
两初温相等时最大功为零。直接传热没有损失总能量，却损失了温差所能提供的做功能力；不能把“能量守恒”误读为“有用功不变”。
\end{xesolution}

\Needspace{180pt}
\begin{origexamplebox}{将有限物体冷却到环境温度以下所需的最少功}\label{ch16:ex:teaching-cooling-work}
热容量为常数 $C$ 的物体初温与环境热库同为 $T_0$。用循环复原的制冷机将物体冷却到 $T_1<T_0$，环境保持 $T_0$。求最少输入功。忽略体积变化及漏热，假定此温区无相变。

\ThQCooling
\end{origexamplebox}
\Needspace{65pt}
\begin{xesolution}
\ThACooling

在物体瞬时温度 $T$ 下再降低 $\dd T<0$，须取出热量 $\delta Q_c=-C\dd T>0$。可逆制冷系数为 $T/(T_0-T)$，所以最小元功为
\[
 \delta W_{\min}=(-C\dd T)\frac{T_0-T}{T}.
\]
系数随温度改变，必须积分：
\[
 W_{\min}=C\int_{T_1}^{T_0}\left(\frac{T_0}{T}-1\right)\dd T
 =C\left[T_0\ln\frac{T_0}{T_1}-(T_0-T_1)\right].
\]
物体总计放出 $Q_c=C(T_0-T_1)$，环境得到 $Q_h=Q_c+W_{\min}=CT_0\ln(T_0/T_1)$。因此物体熵变 $C\ln(T_1/T_0)$ 恰被环境的 $Q_h/T_0$ 抵消，验证了可逆极限。
令 $x=(T_0-T_1)/T_0$，小温差时 $W_{\min}\simeq CT_0x^2/2>0$，而不是与温差一次方成正比。真实有限速率制冷还会产生熵，所需功更大；本模型不能跨越热容变化或相变而直接外推至绝对零度。
\end{xesolution}

\section{熵的微观解释}

宏观量内能可以解释为微观粒子动能与势能的总和。熵也最初来自宏观热过程的研究，但随着分子运动理论
的发展，人们逐渐认识到它的微观含义。玻尔兹曼最早系统提出基于分子运动理论的熵解释。虽然这一理论
起初饱受争议，但后来被广泛接受；在主流学界尚未普遍接受玻尔兹曼理论时，普朗克和爱因斯坦已经
基于其统计思想建立早期量子理论。

\par\noindent\begin{minipage}{\linewidth}
\boltzmannMemorialFigure
\diagramcaption{玻尔兹曼纪念碑上的熵公式，把宏观熵与微观状态的可能数目联系起来。}\label{ch16:fig:s30-21-6}
\end{minipage}\par\medskip
图~\ref{ch16:fig:s30-21-6}中纪念碑上方的公式，是接下来计数解释的核心。

今天通常把玻尔兹曼公式写成
\begin{equation}
  \boxed{S=k_B\ln\Omega},
\end{equation}
其中 $k_B$ 是玻尔兹曼常数，$\Omega$ 是给定宏观状态所对应的微观状态数。一个宏观状态允许的微观
排布可能性越多，熵就越大；这与把熵理解为系统“无序程度”的传统图像接近，但微观状态数的说法更准确。

先看气体分布这一例子。把容器分为等大的左、右两半，暂只记录每个粒子位于哪一半，而不记录具体位置和速度。
给 $N$ 个粒子临时编号，是为了计数各粒子的左右选择，不是把编号当作新的物理自由度。
若左侧有 $N_1$ 个，右侧有 $N-N_1$ 个，则先从全部粒子中挑出左侧粒子的有序序列，有
$N(N-1)\cdots(N-N_1+1)$ 种；同一组左侧粒子的内部排列被重复计了 $N_1!$ 次，故应除去：
\[
 \Omega_{\rm LR}(N_1)=\frac{N(N-1)\cdots(N-N_1+1)}{N_1!}
 =\frac{N!}{N_1!(N-N_1)!}.
\]
例如 $N=4$ 时，全在左侧仅1种，两边各2个有6种。一般有
$\Omega_{\rm LR}(N_1+1)/\Omega_{\rm LR}(N_1)=(N-N_1)/(N_1+1)$，
所以计数先增后减，在左右数目最接近时最大。每个粒子左右等可能且彼此独立时，
概率为 $\Omega_{\rm LR}/2^N$。这解释了为什么宏观气体倾向于分散到整个容器。
$\Omega_{\rm LR}$ 只计算这一粗粒化的分布因素，并不是包含全部位置、动量的完整微观状态数。

同一组合计数也可用于顺磁性物质。某些原子具有非零磁矩，其性质有点像指南针；在外磁场作用下，简化模型中每个磁矩
只有沿磁场或逆磁场两个方向。设系统含 $N$ 个原子，其中 $N_1$ 个磁矩向上，剩下
$N_2=N-N_1$ 个向下。宏观测量只能确认上下数目的差，不能确定究竟哪一些原子向上。于是，在 $N$ 个
原子中选出 $N_1$ 个向上的微观排布数为
\begin{equation}
  \boxed{\Omega=\frac{N!}{N_1!(N-N_1)!}}.
\end{equation}
在没有外场时，$N_1=N/2$ 附近的微观状态数最大，这正是通常情况下顺磁性物质不显示宏观磁性的统计
原因。

中学阶段通常不会进行复杂的微观熵运算，但必须注意微观世界与宏观世界之间巨大的数目差异。设一个人
吃下一份标称化学能为 $\SI{40}{J}$ 的食物，体温为 $\SI{310}{K}$，能否把人体熵增算成 $40/310$？
不能：进食是物质流入和化学反应，不等于向恒温封闭系统传入40焦耳的热量；还要考虑物质携熵、散热和代谢。
为比较数目大小，改考察一个明确的理想模型：温度恒为 $\SI{310}{K}$ 的热库净吸收
$\SI{40}{J}$ 热量。这个热库的熵变才是
\[
  \Delta S=\frac{Q}{T}
  =\frac{\SI{40}{J}}{\SI{310}{K}}
  \approx\SI{0.13}{J.K^{-1}}.
\]
另一方面，玻尔兹曼公式给出
\[
  \Delta S=k_B\ln\frac{\Omega_f}{\Omega_i}.
\]
联立两式，取 $k_B=\SI{1.380649e-23}{J.K^{-1}}$，得到
\begin{equation}
  \boxed{\ln\frac{\Omega_f}{\Omega_i}
  =\frac{40/310}{1.380649\times10^{-23}}
  \approx9.35\times10^{21}}.
\end{equation}
这里先报告对数，避免把数量级粗略相近的巨大指数误写成状态数之比的数值近似。按这一统计描述，
该比值可写为 $\exp(9.35\times10^{21})$ 的量级说明；宏观上很小的热量也可对应极大的微观可能数目变化。


\Needspace{12\baselineskip}\section{综合例题}
\begin{origexamplebox}{【卡诺热机】}\label{ch16:ex:s30-21-1}
利用熵增原理证明卡诺定理：任意工作物质作卡诺循环的效率都只与高、低温热库的温度有关，
与工质的数量、物态方程均无关，效率均为
\[
  \eta=1-\frac{T_2}{T_1},
\]
其中 $T_1,T_2$ 分别为高、低温热库的温度。

\ThQCarnot
\end{origexamplebox}\par\nopagebreak[4]

\par\nopagebreak[4]\begin{xesolution}
设热机在一个可逆循环中从温度为 $T_1$ 的高温热库吸收热量 $Q_1$，向温度为 $T_2$ 的低温热库放出
热量 $Q_2$。工质经历完整循环后回到初态，作为状态函数的熵没有净改变。

高温热库失去热量，熵变为 $-Q_1/T_1$；低温热库得到热量，熵变为 $Q_2/T_2$。卡诺循环是可逆
循环，把工质、两热库和只储存机械能的理想功库一并作为孤立整体；功库熵不变，所以总熵变为零，因此
\[
  -\frac{Q_1}{T_1}+\frac{Q_2}{T_2}=0,
  \qquad
  \frac{Q_2}{Q_1}=\frac{T_2}{T_1}.
\]

\par\noindent\begin{minipage}{\linewidth}
\carnotProofDiagram
\figcap{卡诺热机的能量流与熵平衡。$Q_1,Q_2,W$ 均表示图中箭头方向上的正量。}\label{ch16:fig:carnot-proof}
\end{minipage}\par\medskip
图\ref{ch16:fig:carnot-proof}中热量的大小不同，但可逆情况下从高温热库取走的熵等于送给低温热库的熵。

热机一个循环的对外净功为 $W=Q_1-Q_2$，故
\[
  \boxed{\eta=\frac{W}{Q_1}
  =1-\frac{Q_2}{Q_1}
  =1-\frac{T_2}{T_1}}.
\]
推导中没有使用工质的物态方程、热容量或物质的量，只使用了“循环后工质熵不变”和“整个过程可逆”。
所以任意工质作卡诺循环都具有同一效率。若过程不可逆，总熵会增加，向低温热库排出的热量更多，
实际效率必低于卡诺效率。
\end{xesolution}


\Needspace{12\baselineskip}\begin{origexamplebox}{【熵增原理】}\label{ch16:ex:s30-21-2}
利用熵增原理，重新对热力学第二定律的开尔文和克劳修斯陈述给出定量说明。

\ThQSecondLaw
\end{origexamplebox}\par\nopagebreak[4]

\par\nopagebreak[4]\begin{xesolution}
\textbf{开尔文表述。} 假设一台循环装置从温度为 $T$ 的单一热库吸收热量 $Q$，并把它全部变为功。
装置完成循环后熵变为零；热库失去热量，熵变为
\[
  \Delta S_{\rm reservoir}=-\frac{Q}{T}<0.
\]
还需把接收机械功的理想功库纳入整体，例如无摩擦地被提高的重物；它只增加机械能，熵不变。
于是孤立的“装置、热库加功库”总熵减少，违背熵增原理。因此不可能在没有其他变化的条件下，从单一热库
吸热并把热量全部转化为功。

\textbf{克劳修斯表述。} 假设一台循环装置不消耗功，却把热量 $Q$ 从温度 $T_c$ 的低温热库搬到温度
$T_h$ 的高温热库，其中 $T_h>T_c$。两热库的总熵变为
\[
  \Delta S_{\rm tot}
  =\frac{Q}{T_h}-\frac{Q}{T_c}<0.
\]
这同样违背熵增原理。因此把热量从低温处送往高温处，必须伴随外界做功或其他补偿变化。

\par\noindent\begin{minipage}{\linewidth}
\secondLawEntropyDiagram
\figcap{两种被禁止的唯一效果及其负熵变。比较时都须把装置和所有环境变化计入。}\label{ch16:fig:second-law-comparison}
\end{minipage}\par\medskip
图\ref{ch16:fig:second-law-comparison}左侧的热机必须已完成循环；单次等温膨胀并不属于这个禁止条件。

两个陈述表面上分别谈“热变功”和“逆向传热”，本质上都排除了使孤立系统总熵减少的循环装置。
\end{xesolution}


\Needspace{12\baselineskip}\begin{origexamplebox}{【统统证出来】}\label{ch16:ex:s30-21-3}
请利用热力学第二定律辨析下列说法，指出所需条件，并给出正确结论：
\begin{enumerate}[label=\alph*.]
  \item 理想气体的等温膨胀过程是不可能自然发生的；
  \item 对本讲的定量理想气体，可逆绝热线与等温线不能相交两次；
  \item 对固定组成、稳定单相的简单可压缩系统，两条不同可逆绝热线不能相交；
  \item 有限温差热传递过程不可逆；
  \item 快速膨胀过程不可逆。
\end{enumerate}

\ThQIrreversibility
\end{origexamplebox}\par\nopagebreak[4]

\par\nopagebreak[4]\begin{xesolution}
这组判断的共同方法是：若假定某过程可以在没有其他影响的情况下逆转，就尝试拼出第二类永动机，或
直接算出总熵减少。

\textbf{a. 等温膨胀。} 这一说法本身不成立，不能给“自然发生”附加原本没有的限制来证明它。
理想气体在恒温环境中推动活塞缓慢膨胀可以发生，单个过程并不违反第二定律。此时 $\Delta U=0$，对外功
\[
  W=nRT\ln\frac{V_2}{V_1}>0
\]
必须完全来自热库供热。工质的体积已经改变，这不是“装置循环复原后只留下热变功”这一被禁止的效果。
它不能在既绝热又输出功的条件下仍维持等温，这首先可由第一定律判断。另一个更直接的反例是撤去
隔板，让理想气体向真空自由膨胀，则 $Q=W=\Delta U=0$，温度确实不变，而且过程会自发发生；这正是
不可逆等温自由膨胀，不能与可逆等温膨胀混为一谈。

\par\noindent\begin{minipage}{\linewidth}
\irreversibleProcessesDiagram
\figcap{左侧分别画出撤去隔板前、后的气体状态；右侧为两恒温热库间的有限温差传热。}\label{ch16:fig:irreversibility}
\end{minipage}\par\medskip
图\ref{ch16:fig:irreversibility}中左图初态有真空区，末态虚线只标隔板原位置，不再是阻挡气体的实体壁；
右图的传热将在第d问中计算总熵变。

\textbf{b. 绝热线与等温线的交点。} 对本讲的理想气体，固定温度下
$S(T,V)=\text{常量}+nR\ln V$ 随 $V$ 严格增加，而同一可逆绝热线有固定熵，故只能选出一个 $V$。
也可代入 $P=nRT/V$ 与 $PV^\gamma=K$，得 $nRTV^{\gamma-1}=K$，在 $\gamma>1,V>0$ 时解唯一。
若强行假设两次相交并能围成非零功的循环，就可用一段等温线与一段绝热线构造单热库循环，违反第二定律。
对于一般稳定单相系统，还需可逆绝热线上的温度随体积单调，例如通常热膨胀系数为正的介质。
仅有正热容不能保证这个单调性；存在反常热膨胀时，不能把“两次相交不可能”无条件推广。

\textbf{c. 两条绝热线不能相交。} 可逆绝热线就是等熵线。若两条绝热线在某平衡状态相交，它们在该点
具有相同的熵；对于给定组成、单相稳定且物态方程光滑的简单系统，$c_v>0$ 使
$(\partial S/\partial T)_V=nc_v/T>0$，故固定熵和体积可唯一确定温度。通过该点的可逆绝热路径由微分方程唯一确定，因此
两条曲线实际是同一条绝热线，而不是两条不同曲线。若把它们看作不同曲线，还可以在两曲线间补上一小段
等温线，构造出上一问所排除的单热库循环。

\textbf{d. 有限温差传热不可逆。} 先设两个恒温热库，热量 $Q$ 从 $T_h$ 传到 $T_c$，$T_h>T_c$。总熵变
\[
  \Delta S_{\rm tot}
  =-\frac{Q}{T_h}+\frac{Q}{T_c}
  =Q\left(\frac1{T_c}-\frac1{T_h}\right)>0.
\]
若要完全逆转，就必须使孤立系统熵减少，故不能自发实现。只有在传热温差趋于零的可逆极限中，熵产生
才趋于零。
若两个有限物体的温度会变化，则应对微小传热量积分：
$\Delta S_{\rm tot}=\int(1/T_c-1/T_h)\,\delta Q>0$；不能把初始温度直接代入有限热量公式。

\textbf{e. 快速膨胀。} “快”本身不是可直接代入熵公式的物理条件。通常快速膨胀来不及维持平衡，
伴随有限压差、黏滞、冲击或最终弛豫时，会产生熵，因而不可逆；若只给一个短时间而没有机制，证明并不充分。
以绝热容器中撤去隔板、气体向真空自由膨胀为明确模型：气体由 $V_i$ 扩展到 $V_f$，温度不变，但
\[
  \Delta S=nR\ln\frac{V_f}{V_i}>0.
\]
若想让气体与环境都无变化地回到初态，就必须使总熵减少，所以过程不可逆。

\end{xesolution}


\Needspace{12\baselineskip}\begin{origexamplebox}{【热传递】}\label{ch16:ex:s30-21-4}
三个热容量均为 $c$ 的物体，已知 $T_1=T_2=\SI{300}{K}$，$T_3=\SI{100}{K}$。若外界不做功也不供热，
利用热机工作于三物体之间，使三物中某一物体温度升高，求它所能达到的最高温度；此时另外两物体的
温度各为多少？

\ThQThreeBodies
\end{origexamplebox}\par\nopagebreak[4]

\par\nopagebreak[4]\begin{xesolution}
先证明任何允许过程的温度上界，再证明这个上界可由可逆装置趋近，不能只用两条等式猜测最优状态。
取最终最高温度为 $X$，另外两体温度为 $Y,Z$，均以开尔文计。热机完成循环，不储存净能量，
则第一定律和熵增原理分别给出
\[
 X+Y+Z=700,\qquad
 \Delta S=c\ln\frac{XYZ}{300^2\times100}\ge0.
\]
所以 $XYZ\ge9.0\times10^6$。又由 $(Y-Z)^2\ge0$，
\[
 YZ\le\left(\frac{Y+Z}{2}\right)^2=\left(\frac{700-X}{2}\right)^2,
 \quad X(700-X)^2\ge36.0\times10^6.
\]
因为 $X$ 是最高温度，$X\ge700/3$。函数 $f(X)=X(700-X)^2$ 在该区间内有
$f'(X)=(700-X)(700-3X)\le0$；而 $f(400)=36.0\times10^6$，故必有 $X\le400$。
等号同时要求总熵不变及 $Y=Z$。这才说明可逆性和另外两体等温是上界成立的必要条件。

设目标物体的最高温度为 $T_{\max}$。三个物体的热容量均为常量 $c$，系统与外界既不交换热量也不交换
功，因此总能量守恒：
\[
  c(T_{\max}+2T)=c(300+300+100),
  \qquad T_{\max}+2T=700.
\]
可逆极限下总熵不变：
\[
  c\ln\frac{T_{\max}}{300}
  +c\ln\frac{T}{300}
  +c\ln\frac{T}{100}=0.
\]
于是
\[
  T_{\max}T^2=300\times300\times100.
\]

\par\noindent\begin{minipage}{\linewidth}
\threeBodiesDiagram
\figcap{两台可逆装置在三物体之间协同工作：热机的功全部驱动热泵，不从外界取得净功。}\label{ch16:fig:three-body-engine}
\end{minipage}\par\medskip

联立 $T_{\max}=700-2T$，得到
\[
  (700-2T)T^2=9.0\times10^6.
\]
物理解为 $T=\SI{150}{K}$，从而
\[
  \boxed{T_{\max}=\SI{400}{K}},
  \qquad
  \boxed{T_{\rm other}=\SI{150}{K},\ \SI{150}{K}}.
\]
方程还有 $T=\SI{300}{K},T_{\max}=\SI{100}{K}$ 这一根，它只是把初态中三物体的编号重新排列，并非
所求的最高温度状态。
最后验证可实现性。图\ref{ch16:fig:three-body-engine}中，令供热体 $Y$ 与低温体 $Z$ 之间的可逆热机输出功，
恰好驱动从 $Z$ 向目标体 $X$ 输热的可逆热泵。每小步都令两机完成循环，故没有外界净功输入。
热机从 $Y$ 取出 $-c\dd Y>0$，输出功
$\delta W=(-c\dd Y)(1-Z/Y)$；热泵向 $X$ 送入 $c\dd X$，需要功
$\delta W=c\dd X(1-Z/X)$。二者相等，等价于
\[
 \frac{\dd X}{X}+\frac{\dd Y}{Y}+\frac{\dd Z}{Z}=0,
 \qquad \dd X+\dd Y+\dd Z=0.
\]
沿这一可逆路径，$XYZ=9.0\times10^6$ 且 $X+Y+Z=700$ 始终成立。对 $300\le X<400$，
另外两个正根是
\[
 Y,Z=\frac{700-X\pm\sqrt{(700-X)^2-36.0\times10^6/X}}{2},
\]
并满足 $X\ge Y>Z$，热机与热泵均可按图运行。随着 $X\to400$，$Y,Z\to150$，
两者的温差趋于零，达到可逆极限；实际有限速率装置只能逼近这一最高温度。
所有对数都作用于温度比，温度积的两边单位都是 $\mathrm K^3$。
\end{xesolution}


\Needspace{12\baselineskip}\begin{origexamplebox}{【液体表面能】}\label{ch16:ex:s30-21-5}
小明学习了卡诺热机后，想把液体表面作为工作物质。液体表面能够对外界物体施加表面张力，因此
表面具有能量，并且可以得到
\[
  U=\left(\sigma-T\frac{\dd\sigma}{\dd T}\right)A,
\]
其中 $\sigma$ 为液体的表面张力系数，$A$ 为液体表面积。请研究利用液体表面作为工作物质的热机，并
求出其效率。

\ThQSurfaceEngine
\end{origexamplebox}\par\nopagebreak[4]

\par\nopagebreak[4]\begin{xesolution}
先约定 $A$ 是所有被计入的气液界面的总面积；若是肥皂膜，须把两面都计入，不能再额外乘2。
本题 $U$ 指随界面增加而产生的表面过剩内能，不包括整块液体的体相内能。
可逆扩展表面时外界对表面做元功 $\delta W^{\rm in}=\sigma\dd A$，因而
\[
 \dd U=T\dd S+\sigma\dd A,\qquad
 \delta W^{\rm out}=-\sigma\dd A.
\]
定义亥姆霍兹（Helmholtz）自由能 $F=U-TS$，有 $\dd F=-S\dd T+\sigma\dd A$。
在忽略曲率效应、$\sigma$ 仅随温度变化时，表面过剩自由能为
\[
  F=\sigma(T)A.
\]
又因为 $F=U-TS$，把题给的 $U$ 代入可得表面熵
\[
  S=\frac{U-F}{T}
  =-A\frac{\dd\sigma}{\dd T}.
\]
普通液体的表面张力通常随温度升高而减小，故 $\dd\sigma/\dd T<0$；在给定温度下扩大表面会增加熵，
并需要从热库吸热。

让液体表面工作在温度分别为 $T_h,T_c$ 的两个热库之间，组成两段等温过程和两段等熵过程。在高温
$T_h$ 下，表面积从 $A_1$ 可逆增加到 $A_2$，表面吸热
\[
  Q_h=T_h\Delta S,
  \qquad
  \Delta S=-\left.\frac{\dd\sigma}{\dd T}\right|_{T_h}(A_2-A_1)>0.
\]
然后沿等熵过程降到 $T_c$，在低温下使表面积减小并向低温热库放热。由于两条等温段的熵变量大小相等，
放出的热量为
\[
  Q_c=T_c\Delta S.
\]
所以一个循环的净功与效率为
\[
  W_{\rm out}=Q_h-Q_c=(T_h-T_c)\Delta S,
  \qquad
  \boxed{\eta=1-\frac{T_c}{T_h}}.
\]

\par\noindent\begin{minipage}{\linewidth}
\surfaceEngineDiagram
\figcap{表面工质的可逆循环。上下为等温段，两条弯曲边为等熵段；面积与温度都随等熵过程改变。}\label{ch16:fig:surface-cycle}
\end{minipage}\par\medskip

还需确定等熵线，不能直接把它画成等面积直线。对 $S=-A\sigma'(T)$ 微分得
\[
 \dd S=-\sigma'(T)\dd A-A\sigma''(T)\dd T=0,
 \qquad A\sigma'(T)=\text{常量}.
\]
一般等熵过程既改变面积也改变温度。图\ref{ch16:fig:surface-cycle}采用一个明确的非退化模型
$\sigma=\sigma_*-aT^2$，其中 $a>0$，温区内 $\sigma>0$。这时
$S=2aAT$、$U=(\sigma_*+aT^2)A$、$C_A=2aAT>0$，等熵线为 $AT=\text{常量}$。
高温等温段 $1\to2$ 的面积为 $A_1\to A_2$；冷却到低温的状态3满足
$A_3=A_2T_h/T_c$；低温压缩到状态4时 $A_4=A_1T_h/T_c$，再沿等熵线回到状态1。
由此 $Q_h=2aT_h^2(A_2-A_1)$、$Q_c=2aT_hT_c(A_2-A_1)$，再次给出相同效率。

如果只在很窄温区用线性表面张力近似，形式上 $\sigma'$ 为常量，等熵线会变成 $A=\text{常量}$，
循环呈矩形，形式功为
\[
  W_{\rm out}=(\sigma_c-\sigma_h)(A_2-A_1).
\]
但仅含表面项时这还导致 $C_A=-TA\sigma''=0$，是退化近似，不能据此描述真实液膜如何绝热升降温。
若计入体相自由能 $F_0(T)$，总熵变成 $S_0(T)-A\sigma'(T)$，等熵线必须重新求，通常不再竖直。
无论哪一可靠模型，只要可逆且工作在两个热库之间，效率仍为 $1-T_c/T_h$。
在 $\sigma$--$A$ 图中，本题是逆时针循环，因 $W_{\rm out}=-\oint\sigma\dd A$，其净功为正；
这与气体的 $W_{\rm out}=\oint P\dd V$ 符号相反，不能只凭“循环方向”判断正负。
\end{xesolution}
}

\endgroup

\hypersetup{bookmarksdepth=2}

% ThermalQualityRevision20261023-v1
% QualityRevision20261023
\hypersetup{bookmarksdepth=3}
\begingroup
\renewcommand{\Needspace}[1]{\par\begingroup
\dimen0=#1\relax\vskip0pt plus\dimen0\penalty-100
\vskip0pt plus-\dimen0\vskip\dimen0\penalty9999\vskip-\dimen0
\vskip0pt\endgroup}
\chapter{相变与表面现象}
% Editable TikZ for phase equilibrium and mixed gases.
\providecommand{\ThermalFigure}[3]{\begin{qualityfigure}
\begin{tikzpicture}[>=Latex,font=\small,x=.8cm,y=.8cm]#1\end{tikzpicture}
\figcap{#2}\label{#3}\end{qualityfigure}}
\newcommand{\ThQCondensation}{\ThermalFigure{
 \draw[thick](0,2.2)--(0,0)--(4,0)--(4,2.2);\draw[line width=2pt](0,1.7)--(4,1.7);
 \node[align=center] at(2,.9){干空气 $p_{a0}$\\水蒸气 $\phi p_s$};
 \draw[->,thick](4.7,2)--(4.7,1)node[right]{等温压缩};
 \node[right,align=left]at(6,1.5){初始 $V_0$、无液水\\压缩后 $V=V_0/r$};
}{先按各组分分压判断饱和；总压强不是水蒸气的饱和判据。}{ch17:quality:condensation-q}}
\newcommand{\ThACondensation}{\ThermalFigure{
 \draw[->](0,0)--(6.3,0)node[right]{$r$};\draw[->](0,0)--(0,3)node[above]{$p_v/p_s$};
 \draw[DeepBlue,thick](1,.6)--(3,2)--(6,2);\draw[dashed](3,0)--(3,2) (0,2)--(6,2);
 \node[below]at(1,0){$1$};\node[below]at(3,0){$1/\phi$};\node[left]at(0,2){$1$};
 \node[above]at(1.8,1.5){无液水};\node[above]at(4.6,2.1){气液共存};
}{凝结前蒸气分压随压缩倍数增大，凝结后被饱和蒸气压限制。}{ch17:quality:condensation-a}
图~\ref{ch17:quality:condensation-a}中的折点满足 $p_v=p_s$。折点右侧气相水的物质的量减少，才使分压保持不变；空气分压没有这一平台，因而总压强仍随 $r$ 增长。
}
\newcommand{\ThBoilingBase}[1]{
 \draw[->](0,0)--(5,0)node[right]{$T$};\draw[->](0,0)--(0,3.3)node[above]{$p_s$};
 \draw[DeepBlue,thick,domain=.5:4.6,samples=60]plot(\x,{.15*exp(.6*\x)});
 \draw[dashed](0,2.2)--(4.5,2.2);\draw[dashed](0,1.1)--(4.5,1.1);
 \node[left]at(0,2.2){$p_1$};\node[left]at(0,1.1){$p_1/2$};#1}
\newcommand{\ThQBoiling}{\ThermalFigure{\ThBoilingBase{
 \node[right]at(5,2){已知 $p_s(T_1)=p_1$};
}}{降低外压后，沸腾条件仍为饱和蒸气压等于外压；曲线仅作定性示意。}{ch17:quality:boiling-q}}
\newcommand{\ThABoiling}{\ThermalFigure{\ThBoilingBase{
 \draw[dashed](4.476,0)--(4.476,2.2) (3.321,0)--(3.321,1.1);
 \node[below]at(4.476,0){$T_1$};\node[below]at(3.321,0){$T_2$};
 \draw[->,thick](4.4,2.8)--(3.4,2.8);
}}{在正斜率共存线上，压强下降对应沸点下降；温度并不减半。}{ch17:quality:boiling-a}
图~\ref{ch17:quality:boiling-a}中两条水平线与同一共存线相交。积分克拉珀龙方程应在这两个交点之间进行，使用 $\ln(p_2/p_1)$ 与倒温度差，而不能把压强比直接当成温度比。
}
\newcommand{\ThQMercury}{\ThermalFigure{
 \draw[thick](6,0)--(0,0)--(0,.6)--(6,.6);\fill[TextGray!40](3.4,0)rectangle(5.4,.6);
 \node[above]at(1.7,.8){空气＋汞蒸气，$17\,\mathrm{cm}$};\node[below]at(4.4,-.1){汞柱 $10\,\mathrm{cm}$};
 \draw[thick](8,-1)--(8,3)--(8.6,3)--(8.6,-1);\fill[TextGray!40](8,0)rectangle(8.6,1);
 \node[right]at(9,2){封闭气柱 $20\,\mathrm{cm}$};\node[right]at(9,.5){汞柱 $10\,\mathrm{cm}$};
 \node[right]at(9,-.7){下端与大气相通};
}{水平与封闭端朝上的竖直情形；两图分别按各自尺度画出，竖直方向由题给膨胀事实辨明。}{ch17:quality:mercury-q}}
\newcommand{\ThVaporizationBase}[1]{
 \node[draw,minimum width=2.2cm,minimum height=1cm](L)at(0,0){液态苯 $m$};
 \node[draw,minimum width=2.6cm,minimum height=1.6cm](G)at(6,0){蒸气 $m$};
 \draw[->,thick](L)--node[above]{$p,T$ 不变}(G);
 \draw[->,DeepBlue](2.4,-1.6)--(2.4,-.2)node[midway,left]{$Q$};#1}
\newcommand{\ThQVaporization}{\ThermalFigure{\ThVaporizationBase{}}{同一质量的苯等压汽化；气相体积远大于液相体积。}{ch17:quality:vaporization-q}}
\newcommand{\ThAVaporization}{\ThermalFigure{\ThVaporizationBase{
 \draw[->,thick](G.east)--(9,0)node[right]{$W=p\Delta V$};
 \node[below]at(6,-1.4){吸热分为内能增加与对外做功};
}}{把同一质量的物质作为系统；系统边界随汽化膨胀而移动。}{ch17:quality:vaporization-a}
图~\ref{ch17:quality:vaporization-a}中的出功箭头表示克服外压的膨胀功。因此 $mL$ 是焓的增加量，求内能增加必须扣除 $p(V_g-V_l)$；忽略 $V_l$ 后该项为 $nRT$。
}
\newcommand{\ThQWetExpansion}{\ThermalFigure{
 \foreach \x/\w/\lab in {0/2/初态,5/3/水刚消失}{
 \draw[thick](\x,2)--(\x,0)--(\x+\w,0)--(\x+\w,2);\draw[line width=2pt](\x,1.5)--(\x+\w,1.5);
 \node[above]at(\x+.5*\w,2){\lab};}
 \fill[DeepBlue!30](.2,.1)rectangle(1.8,.3);
 \node[below]at(1,-.3){$V,\ 3.0\,\mathrm{atm}$};\node[below]at(6.5,-.3){$2V,\ 2.0\,\mathrm{atm}$};
 \draw[->](2.6,.8)--(4.4,.8)node[midway,above]{等温};
}{液水尚存与刚好消失是两个不同约束；干空气物质的量始终不变。}{ch17:quality:wet-expansion-q}}
\newcommand{\ThQTwoVessels}{\ThermalFigure{
 \draw[thick](0,0)rectangle(3,2);\draw[thick](6,0)rectangle(9,2);
 \draw[thick](3,.8)--(6,.8) (3,1.2)--(6,1.2);
 \node at(1.5,1){容器甲，$V$};\node at(7.5,1){容器乙，$V$};
 \node[below]at(1.5,-.3){保持 $100\,{}^\circ\mathrm C$};\node[below]at(7.5,-.3){冷却至 $0\,{}^\circ\mathrm C$};
 \node[above]at(4.5,2.6){初态均为 $1\,\mathrm{atm}$、$100\,{}^\circ\mathrm C$、$50\%$ 相对湿度};
}{细管允许组分重新分配，温度分别由两个外界维持。}{ch17:quality:vessels-q}}
\newcommand{\ThQCloud}{\ThermalFigure{
 \draw[thick](-1,0)--(5,0);\draw[->](0,0)--(0,3.8)node[above]{$z$};
 \draw[DeepBlue,thick](2,.5)circle(.5);\draw[DeepBlue,thick](2,2.6)circle(.8);
 \draw[->,thick](2,1.2)--(2,1.7)node[right]{未凝结气团上升};
 \node[right]at(3.3,.5){地面 $p_0,T_0$};\node[right]at(3.3,2.7){与周围同压，绝热降温};
 \draw[dashed](.3,3.4)--(5,3.4)node[right]{凝结起点待定};
}{干绝热模型仅用于到达饱和之前；凝结后需重新计入潜热。}{ch17:quality:cloud-q}}



{
\graphicspath{{../../../sample/xes-original/thirty-first-section/}}

% All annotations are transparent: labels never carry an opaque rectangular mask.
\tikzset{
  xeslabel/.style={font=\small,text=Ink,fill=none,inner sep=0.6pt},
  xessmall/.style={font=\footnotesize,text=TextGray,fill=none,inner sep=0.5pt},
  xesline/.style={draw=Ink,line width=0.85pt},
  xesblue/.style={draw=DeepBlue,line width=1.05pt},
  xesarrow/.style={-{Latex[length=2.2mm,width=1.4mm]},draw=DeepBlue,line width=0.9pt},
  forcearrow/.style={-{Latex[length=2.3mm,width=1.5mm]},draw=DeepBlue,line width=1.0pt}
}


\newcommand{\phaseDiagramWater}{%
\begin{center}\begin{minipage}{\linewidth}\centering
\begin{tikzpicture}[x=1.03cm,y=0.73cm]
  \draw[xesline,-{Latex[length=2mm]}] (0,0)--(8.7,0) node[xeslabel,below left=1pt] {$T$};
  \draw[xesline,-{Latex[length=2mm]}] (0,0)--(0,6.2) node[xeslabel,below right=1pt] {$p$};
  \coordinate (A) at (2.55,1.55);
  \coordinate (C) at (7.05,4.35);
  \draw[xesblue] (0.25,0.35) .. controls (0.95,0.7) and (1.75,1.15) .. (A);
  \draw[xesblue] (A) .. controls (4.0,1.9) and (5.65,3.0) .. (C);
  \draw[xesblue] (A) .. controls (2.25,3.15) and (1.92,4.75) .. (1.7,5.9);
  \fill[DeepBlue] (A) circle (1.4pt) (C) circle (1.4pt);
  \node[xeslabel,anchor=south east] at ($(A)+(-0.08,0.05)$) {$A$};
  \node[xeslabel,anchor=south west] at ($(C)+(0.08,0.04)$) {$C$};
  \node[xessmall,align=center] at (1.35,3.35) {固态};
  \node[xessmall,align=center] at (4.1,3.8) {液态};
  \node[xessmall,align=center] at (5.55,1.15) {气态};
  \node[xessmall] at (4.75,3.25) {汽化曲线};
  \node[xessmall] at (.5,-.75) {升华曲线};
  \draw[RuleGray,->] (.5,-.5)--(.85,.67);
  \node[xessmall,rotate=79] at (2.8,4.65) {熔化曲线};
  \draw[RuleGray,densely dashed] (A)--(2.55,0) node[xessmall,below,align=center] {$0.01\,{}^{\circ}\mathrm C$};
  \draw[RuleGray,densely dashed] (A)--(0,1.55) node[xessmall,left] {$611.73\,\mathrm{Pa}$};
  \draw[RuleGray,densely dashed] (C)--(7.05,0) node[xessmall,below] {$374\,{}^{\circ}\mathrm C$};
  \draw[RuleGray,densely dashed] (C)--(0,4.35) node[xessmall,left,align=right] {$2.2\times10^7\,\mathrm{Pa}$};
  \node[xessmall,anchor=south] at (7.05,4.95) {临界点};
  \node[xessmall,anchor=west] at (2.72,1.18) {三相点};

  % microscopic phase sketches, placed away from the curves
  \begin{scope}[shift={(8.0,1.1)},scale=0.58]
    \draw[RuleGray,rounded corners=2pt] (-0.75,-0.75) rectangle (0.75,0.75);
    \foreach \x in {-0.48,0,0.48}\foreach \y in {-0.48,0,0.48}{\fill[DeepBlue!70] (\x,\y) circle (1.1pt);}
    \node[xessmall] at (0,-1.05) {固};
  \end{scope}
  \begin{scope}[shift={(8.0,2.45)},scale=0.58]
    \draw[RuleGray,rounded corners=2pt] (-0.75,-0.75) rectangle (0.75,0.75);
    \foreach \p in {(-.5,-.35),(-.05,-.48),(.42,-.25),(-.4,.15),(.08,.06),(.5,.3),(-.15,.48)}{\fill[DeepBlue!70] \p circle (1.1pt);}
    \node[xessmall] at (0,-1.05) {液};
  \end{scope}
  \begin{scope}[shift={(8.0,3.8)},scale=0.58]
    \draw[RuleGray,rounded corners=2pt] (-0.75,-0.75) rectangle (0.75,0.75);
    \foreach \p in {(-.53,-.48),(.48,-.31),(-.18,.08),(.45,.5),(-.5,.43)}{\fill[DeepBlue!70] \p circle (1.1pt);}
    \node[xessmall] at (0,-1.05) {气};
  \end{scope}
\end{tikzpicture}
\end{minipage}\end{center}}


\newcommand{\criticalApparatusFigure}{%
\begin{center}\begin{minipage}{\linewidth}\centering
\includegraphics[width=.83\linewidth]{../../../sample/xes-original/restructured-complete/assets/ch17/critical-cell.png}\par\smallskip
\begin{tikzpicture}[x=1cm,y=1cm]
  \foreach \dx in {0,4,8}{
    \draw[Ink,line width=1pt,rounded corners=7pt] (\dx,0) rectangle +(2.5,3.35);
    \draw[RuleGray,line width=2.2pt] (\dx+.28,.25)--(\dx+.28,3.1);
  }
  % below critical: two visibly separated phases
  \fill[DeepBlue!12] (.08,.08) rectangle (2.42,1.48);
  \draw[DeepBlue,line width=1.1pt] (.08,1.48)--(2.42,1.48);
  \foreach \p in {(.45,.45),(.82,.72),(1.18,.38),(1.62,.72),(2.05,.42),(.55,1.05),(1.45,1.08),(2.1,1.1)}{\fill[DeepBlue!65] \p circle (1.2pt);}
  \foreach \p in {(.65,2.1),(1.75,2.65),(2.05,2.0)}{\fill[DeepBlue!50] \p circle (1.0pt);}
  \node[xessmall,align=center] at (1.25,-.48) {$T<T_c$\\清晰气液界面};
  % critical opalescence: fluctuating cloudy density
  \foreach \p/\r in {(4.45,.45)/.18,(4.8,1.0)/.25,(5.35,.62)/.3,(5.9,1.25)/.22,(4.55,1.75)/.35,(5.25,1.65)/.2,(5.8,2.1)/.32,(4.85,2.55)/.28,(5.65,2.85)/.22}{
    \shade[inner color=DeepBlue!28,outer color=white] \p circle (\r);
  }
  \draw[DeepBlue!60,densely dotted] (4.15,1.5) .. controls (4.8,1.7) and (5.5,1.25) .. (6.35,1.55);
  \node[xessmall,align=center] at (5.25,-.48) {$T\approx T_c$\\临界乳光};
  % above critical: uniform supercritical fluid
  \fill[DeepBlue!8] (8.08,.08) rectangle (10.42,3.19);
  \foreach \p in {(8.45,.45),(8.85,.85),(9.35,.42),(9.9,.75),(8.55,1.35),(9.15,1.55),(9.8,1.25),(8.7,2.1),(9.35,2.25),(10,2.05),(8.45,2.75),(9.05,2.95),(9.75,2.7)}{\fill[DeepBlue!45] \p circle (1.05pt);}
  \node[xessmall,align=center] at (9.25,-.48) {$T>T_c,\ p>p_c$\\均匀超临界流体};
\end{tikzpicture}
\end{minipage}\end{center}}

% Quantitative diagrams: shared formulas determine points and curves.
\newcommand{\vdWIsotherms}{%
\begin{center}\begin{minipage}{\linewidth}\centering\begin{tikzpicture}
\begin{axis}[width=.82\linewidth,height=6cm,axis lines=left,xmin=.48,xmax=3.5,ymin=0,ymax=2.7,
xlabel={$v'=V/V_c$},ylabel={$p'=p/p_c$},xtick={1,2,3},ytick={1,2},clip=true]
\foreach \t in {.85,.9,1,1.08,1.16,1.24}{
\addplot[DeepBlue,domain=.48:3.3,samples=250]{8*\t/(3*x-1)-3/x^2};}
\addplot[only marks,mark=*,mark size=1.3pt] coordinates{(1,1)};
\node[above right] at (axis cs:1.02,1.04) {$C$};
\draw[dashed,TextGray] (axis cs:1,0)--(axis cs:1,1);
\node[align=left,anchor=north east,font=\small] at (axis cs:3.35,2.5)
{$t'=0.85,\ 0.90,\ 1.00,$\\$1.08,\ 1.16,\ 1.24$\\从下向上};
\end{axis}\end{tikzpicture}\end{minipage}\end{center}}

\newcommand{\criticalPointFigure}{%
\begin{center}\begin{minipage}{\linewidth}\centering\begin{tikzpicture}
\begin{axis}[width=6cm,height=5.3cm,axis lines=left,xmin=.55,xmax=2.3,ymin=0,ymax=2.4,
xlabel={$v'$},ylabel={$p'$},xtick={1,2},ytick={1,2}]
\addplot[DeepBlue,thick,domain=.56:2.2,samples=250]{8/(3*x-1)-3/x^2};
\addplot[only marks,mark=*] coordinates{(1,1)};
\draw[dashed,TextGray] (axis cs:.8,1)--(axis cs:1.3,1);
\node[above right] at (axis cs:1.07,1.1) {$C$};
\end{axis}
\begin{scope}[xshift=7cm,x=.88cm,y=.82cm]
\draw[axis] (0,0)--(5.7,0) node[right]{$T$};
\draw[axis] (0,0)--(0,5.5) node[above]{$p$};
\draw[DeepBlue,thick,domain=.6:4,samples=100] plot(\x,{.2+.2*\x*\x});
\coordinate (P) at (1,1.5);\coordinate (M) at (2.54951,1.5);\coordinate (Q) at (5,1.5);
\coordinate (A) at (1,4.6);\coordinate (B) at (5,4.6);\coordinate (C) at (4,3.4);
\draw[->,Ink] (P)--(M);\draw[->,Ink] (M)--(Q);
\draw[->,DeepBlue,dashed] (P)--(A);\draw[->,DeepBlue,dashed] (A)--(B);\draw[->,DeepBlue,dashed] (B)--(Q);
\foreach \q in {P,M,Q,A,B,C}{\fill (\q) circle(1.2pt);}
\foreach \q in {P,M,Q}{\node[below=4pt] at (\q){$\q$};}
\foreach \q in {A,B,C}{\node[above=4pt] at (\q){$\q$};}
\node[font=\small] at (2.2,3.7){液相};\node[font=\small] at (4.7,.55){气相};
\end{scope}
\node[font=\small] at (2.5,-1) {(a) 临界等温线};
\node[font=\small] at (9,-1) {(b) 穿越与绕行路径};
\end{tikzpicture}\end{minipage}\end{center}}

\newcommand{\maxwellAreaFigure}{%
\begin{center}\begin{minipage}{\linewidth}\centering\begin{tikzpicture}
\begin{axis}[width=.85\linewidth,height=6cm,axis lines=left,xmin=.48,xmax=2.9,ymin=.25,ymax=1.08,
xlabel={$v'$},ylabel={$p'$},ylabel style={at={(axis description cs:0,1)},anchor=south east},xtick={.603402,2.348842},xticklabels={$v_l'$,$v_g'$},ytick={.64699835},yticklabels={$p_T'$}]
\addplot[draw=none,fill=DeepBlue!15,domain=.603401903:1.090526633,samples=120]
{7.2/(3*x-1)-3/x^2} -- (axis cs:1.090526633,.646998352)--(axis cs:.603401903,.646998352)\closedcycle;
\addplot[draw=none,fill=Gold!20,domain=1.090526633:2.348842376,samples=120]
{7.2/(3*x-1)-3/x^2} -- (axis cs:2.348842376,.646998352)--(axis cs:1.090526633,.646998352)\closedcycle;
\addplot[DeepBlue,thick,domain=.54:2.8,samples=300]{7.2/(3*x-1)-3/x^2};
\draw[Ink,thick] (axis cs:.603401903,.646998352)--(axis cs:2.348842376,.646998352);
\addplot[only marks,mark=*,mark size=1.2pt] coordinates{(.603402,.646998)(.718597,.419843)(1.090527,.646998)(1.528505,.724013)(2.348842,.646998)};
\node[above left] at (axis cs:.603402,.646998){$M$};
\node[below] at (axis cs:.718597,.419843){$A$};
\node[above left] at (axis cs:1.090527,.646998){$P$};
\node[above] at (axis cs:1.528505,.724013){$B$};
\node[above right] at (axis cs:2.348842,.646998){$N$};
\node at (axis cs:.78,.58){$S_1$};\node at (axis cs:2.05,.83){$S_2$};
\draw[->,TextGray] (axis cs:2.05,.79)--(axis cs:1.85,.681);
\node[font=\small] at (axis cs:2.2,.96){$t'=0.9,\ S_1=S_2$};
\end{axis}\end{tikzpicture}\end{minipage}\end{center}}

\newcommand{\contrailFigure}{%
\begin{center}\begin{minipage}{\linewidth}\centering\includegraphics[width=.88\linewidth]{../../../sample/xes-original/restructured-complete/assets/ch17/contrail.png}\end{minipage}\end{center}}

\newcommand{\mercuryTubeFigure}{%
\begin{center}\begin{minipage}{\linewidth}\centering\begin{tikzpicture}[x=.17cm,y=.17cm]
\draw[Ink,thick] (0,5)--(0,10)--(31,10) (0,5)--(31,5);
\fill[DeepBlue!25] (17,5.1) rectangle(27,9.9);
\draw[DeepBlue] (17,5)--(17,10) (27,5)--(27,10);
\draw[<->] (0,2)--(17,2) node[midway,below,font=\small]{$17\,\mathrm{cm}$};
\draw[<->] (17,13)--(27,13) node[midway,above,font=\small]{$10\,\mathrm{cm}$};
\node[font=\small] at (22,7.5){Hg};
\node[font=\small,align=center] at (8.5,18){封闭气体\\空气与饱和汞蒸气};
\draw[->,TextGray] (8.5,15)--(8.5,9);
\draw[Ink,thick] (47,0)--(47,33)--(52,33)--(52,0);
\fill[DeepBlue!25] (47.1,3) rectangle(51.9,13);
\draw[DeepBlue] (47,3)--(52,3) (47,13)--(52,13);
\node[font=\small] at (49.5,8){Hg};
\draw[<->] (56,13)--(56,33) node[midway,right,font=\small]{$20\,\mathrm{cm}$};
\draw[<->] (56,3)--(56,13) node[midway,right,font=\small]{$10\,\mathrm{cm}$};
\node[font=\small,align=center] at (37,26){同温度的\\封闭气体};
\draw[->,TextGray] (44,24)--(48,24);
\node[below,font=\small] at (49.5,0){开口，外接大气};
\node[above,font=\small] at (49.5,34){封闭端};
\end{tikzpicture}\end{minipage}\end{center}}

\newcommand{\mixedGasPVFigure}{%
\begin{center}\begin{minipage}{\linewidth}\centering\begin{tikzpicture}
\begin{axis}[width=.8\linewidth,height=5.7cm,axis lines=left,xmin=0,xmax=4.6,ymin=0,ymax=5.7,
xlabel={$x=V'/V$},ylabel={$p/\mathrm{atm}$},xtick={.5,1,2,4},ytick={1,2,3,5}]
\addplot[DeepBlue,thick,domain=.45:2,samples=150]{1+2/x};
\addplot[DeepBlue,thick,dashed,domain=2:4.4,samples=120]{4/x};
\addplot[only marks,mark=*] coordinates{(.5,5)(1,3)(2,2)(4,1)};
\draw[dotted] (axis cs:2,0)--(axis cs:2,2);
\node[align=left,font=\small] at (axis cs:3.2,4.3){$x<2$：液体仍在\\$x\ge2$：水全汽化};
\node[above right,font=\small] at (axis cs:2,2.1){液体刚好消失};
\end{axis}\end{tikzpicture}\end{minipage}\end{center}}

\newcommand{\pressureCookerQuestionFigure}{%
\begin{center}\begin{minipage}{\linewidth}\centering\begin{tikzpicture}
\begin{scope}[x=.6cm,y=.6cm]
\fill[DeepBlue!15] (.05,.05) rectangle(3.95,1.7);
\draw[thick] (0,4)--(0,0)--(4,0)--(4,4);
\draw[thick] (-.2,4)--(1.7,4)--(1.7,4.6) (2.3,4.6)--(2.3,4)--(4.2,4);
\draw[thick,fill=gray!15] (1.4,4.65) rectangle(2.6,5.15);
\draw[DeepBlue] (0,1.7)--(4,1.7);
\node[font=\small,above] at (2,5.3){限压阀 $S$};
\node[font=\small] at (2,2.6){气相};
\node[font=\small] at (2,.8){液态水};
\end{scope}
\begin{axis}[at={(3.6cm,0)},anchor=south west,width=8cm,height=5.0cm,axis lines=left,
xmin=0,xmax=130,ymin=0,ymax=260,xlabel={$t/{}^\circ\mathrm C$},ylabel={$p_w/\mathrm{kPa}$},
xtick={10,30,50,70,90,100,120},ytick={0,50,100,150,200,250},grid=both,
minor tick num=1,grid style={draw=RuleGray!70},tick label style={font=\scriptsize}]
\addplot[DeepBlue,thick] table[x=t,y=ps]{assets/ch17/saturation.dat};
\end{axis}\end{tikzpicture}\end{minipage}\end{center}}

\newcommand{\pressureCookerFigure}{%
\begin{center}\begin{minipage}{\linewidth}\centering\begin{tikzpicture}
\begin{scope}[x=.7cm,y=.7cm]
\fill[DeepBlue!15] (.05,.05) rectangle(3.95,1.7);
\draw[thick] (0,4)--(0,0)--(4,0)--(4,4);
\draw[thick] (-.2,4)--(1.7,4)--(1.7,4.6) (2.3,4.6)--(2.3,4)--(4.2,4);
\draw[thick,fill=gray!15] (1.4,4.65) rectangle(2.6,5.15);
\draw[DeepBlue] (0,1.7)--(4,1.7);
\draw[->,DeepBlue] (2,3.2)--(2,4.5);
\draw[->,DeepBlue] (1.7,4.63)--(.7,4.63);
\node[font=\small,above] at (2,5.3){限压阀 $S$};
\node[font=\small,align=center] at (2,2.4){空气与水蒸气};
\node[font=\small] at (2,.8){液态水};
\end{scope}
\begin{axis}[at={(4.2cm,0)},anchor=south west,width=8cm,height=5.8cm,axis lines=left,
xmin=0,xmax=130,ymin=0,ymax=340,xlabel={$t/{}^\circ\mathrm C$},ylabel={$p/\mathrm{kPa}$},
xtick={10,90,120},ytick={100,200,300},legend style={at={(.02,.98)},anchor=north west,font=\scriptsize,draw=none,fill=none}]
\addplot[DeepBlue,thick] table[x=t,y=ps]{assets/ch17/saturation.dat};\addlegendentry{$p_w$}
\addplot[TextGray,dashed] table[x=t,y=pa]{assets/ch17/saturation.dat};\addlegendentry{$p_a$（封闭外推）}
\addplot[Ink,thick,restrict x to domain=0:90.174448744] table[x=t,y=total]{assets/ch17/saturation.dat};\addlegendentry{$p_a+p_w$（开启前）}
\addplot[Ink,dashed,restrict x to domain=90.174448744:130] table[x=t,y=total]{assets/ch17/saturation.dat};
\addplot[Ink,thick,domain=90.174448744:120]{198.671424787};
\addplot[dotted,domain=0:128]{198.671424787};
\addplot[only marks,mark=*] coordinates{(90.174448744,198.671424787)(120,198.671424787)};
\draw[dotted] (axis cs:90.174448744,0)--(axis cs:90.174448744,198.671424787)
(axis cs:120,0)--(axis cs:120,198.671424787);
\node[font=\small,anchor=north] at (axis cs:109,191){$p_*$};
\end{axis}\end{tikzpicture}\end{minipage}\end{center}}

\newcommand{\twoVesselsFigure}{%
\begin{center}\begin{minipage}{\linewidth}\centering\begin{tikzpicture}[x=1cm,y=1cm]
\fill[DeepBlue!12] (7.22,.02) rectangle(10.38,.95);
\draw[thick] (3.2,1.5)--(3.2,0)--(0,0)--(0,3.45)--(3.2,3.45)--(3.2,1.9);
\draw[thick] (7.2,1.5)--(7.2,0)--(10.4,0)--(10.4,3.45)--(7.2,3.45)--(7.2,1.9);
\draw[thick] (3.2,1.5)--(7.2,1.5) (3.2,1.9)--(7.2,1.9);
\draw[DeepBlue] (7.2,.95)--(10.4,.95);
\node[font=\small,above] at (5.2,2.15){连通细管};
\node[align=center] at (1.6,2.6){热端\\$100{}^\circ\mathrm C$};
\node[align=center] at (8.8,2.6){冷端\\$0{}^\circ\mathrm C$};
\node[font=\small] at (8.8,.45){凝结水};
\draw[->,DeepBlue] (4.2,.9)--(6.2,.9);
\node[font=\small,below] at (5.2,.8){水蒸气向冷端迁移};
\draw[TextGray] (7,-.2)--(10.6,-.2)--(10.6,1.2) (7,-.2)--(7,1.2);
\node[font=\small,below] at (8.8,-.25){外部冰水浴};
\end{tikzpicture}\end{minipage}\end{center}}

\newcommand{\cloudFormationFigure}{%
\begin{center}\begin{minipage}{\linewidth}\centering\begin{tikzpicture}
\begin{scope}[x=.8cm,y=.8cm]
\draw[axis] (0,0)--(0,5.5) node[above]{$z$};
\foreach \y/\r in {.8/.28,2/.36,3.2/.44,4.4/.52}{\draw[DeepBlue] (1.5,\y) circle(\r);}
\draw[->,DeepBlue] (2.4,.7)--(2.4,4.6);
\draw[dashed,TextGray] (0,4.4)--(3.8,4.4);
\node[font=\small,align=center] at (1.8,5.5){云底 $z_d$\\首次饱和};
\node[font=\small,align=center] at (1.8,-.5){上升膨胀、降温};
\end{scope}
\begin{axis}[at={(4.2cm,0)},anchor=south west,width=8cm,height=5.5cm,axis lines=left,
xmin=269,xmax=290,ymin=.25,ymax=1.12,xlabel={$T/\mathrm K$},ylabel={$p_v/p_s(T_0)$},
xtick={270,278.475,288},xticklabels={270,$T_d$,288},ytick={.4,.6,.8,1}]
\addplot[DeepBlue,thick,domain=270:288,samples=200]{exp(44000/8.314*(1/288-1/x))};
\addplot[Ink,thick,domain=278.474927:288,samples=120]{.6*(x/288)^3.5};
\addplot[Ink,dashed,domain=270:278.474927,samples=80]{.6*(x/288)^3.5};
\addplot[only marks,mark=*] coordinates{(278.474927,.533370375)};
\draw[dotted] (axis cs:278.474927,.25)--(axis cs:278.474927,.533370375);
\node[font=\small] at (axis cs:280,.99){$p_s(T)$};
\node[font=\small] at (axis cs:285,.44){$p_v(T)$};
\draw[->,TextGray] (axis cs:287,.67)--(axis cs:282,.62);
\end{axis}\end{tikzpicture}\end{minipage}\end{center}}


在前面的讨论中，我们把气体看作大量自由运动的分子组成的系统，并用理想气体物态方程描述它们。
这一模型在气体密度较低、分子间距离较大时十分成功，但它有一个明显的局限：理想气体分子之间完全
没有相互作用，也没有自身大小，因此无论怎样降温、压缩，理想气体都不会自动变成液体或固体。真实物质
却可以在固、液、气三种物态之间转化。要理解这些相变，就必须把分子的有限大小和分子间作用力考虑进来。

从微观上看，物质的物态由两种趋势的竞争决定。热运动倾向于使分子分散、排列无序；分子间吸引作用
倾向于把分子聚集起来。当温度较低或密度较大时，分子间作用更重要，物质容易成为液体或固体；当温度
较高、密度较小时，热运动占优势，物质表现为气体。固体熔化成液体的过程叫熔化，液体凝固成固体的过程
叫凝固；液体变成气体叫汽化，气体变成液体叫液化；固体也可以不经过液态直接变成气态，这一过程叫升华。

\section{相图}

物质在不同温度和压强下稳定存在的物态，可以用相图表示。图~\ref{ch17:fig:s31-22-1}是水的 $p$--$T$ 相图。图中三条曲线
分别是固液共存的熔化曲线、液气共存的汽化曲线和固气共存的升华曲线。曲线上每一点都表示两种物态能够
平衡共存；曲线把平面分为固态区、液态区和气态区。跨过一条相界线时，物质发生相变。

\par\noindent\begin{minipage}{\linewidth}
\phaseDiagramWater
\diagramcaption{水的 $p$--$T$ 相图及三种物态的微观示意。曲线、三相点和临界点的位置关系与水的相图特征一致。}\label{ch17:fig:s31-22-1}
\end{minipage}\par\medskip

三条相界线交于点 $A$，称为三相点。在三相点，固态、液态和气态能够同时平衡共存。水的三相点温度为
$0.01\,{}^{\circ}\mathrm C$，压强为 $611.73\,\mathrm{Pa}$。液气共存曲线并不无限延伸，它终止于点 $C$，
称为临界点。水的临界温度约为 $374\,{}^{\circ}\mathrm C$，临界压强约为
$2.2\times10^7\,\mathrm{Pa}$。温度高于临界温度以后，不论怎样等温压缩，都不会出现清晰的气液分界面。

相变时系统需要吸收或放出热量。例如在一定压强下，固体达到熔点后继续吸热，温度可以保持不变而逐渐
熔化；液体达到沸点后继续吸热，温度也可以保持不变而逐渐汽化。单位质量物质完成某种相变所需的热量
称为相应的潜热。相变潜热说明：即使宏观温度不变，分子的空间结构和分子间势能仍会发生变化。

\subsection{相变时为什么吸热而不升温}

对固定压强下的两相平衡，温度由共存条件确定。输入热量主要用于改变两相比例，
并不必然增大每个分子的平均平动能。设单位质量潜热为 $L$，转化质量为 $m$，则
$Q=mL$。但 $L$ 不是单位质量内能变化：若只有体积功，第一定律给出
\begin{equation}
 \Delta u=L-p(v_2-v_1),
\end{equation}
其中 $u$ 为比内能，$v$ 为比体积，单位分别为 $\mathrm{J/kg}$、$\mathrm{m^3/kg}$。
把 $h=u+pv$ 定义为比焓，在同一压强下就有 $L=h_2-h_1$。
气化时 $v_g-v_l>0$，吸热的一部分成为向外膨胀功；因此通常 $\Delta u<L$。
若过程包含显著电功、表面功或动能变化，就必须另列这些能量项，不能直接套用这一简式。

\subsection{两相比例与杠杆关系}

在给定的气液共存温度下，令总质量为 $m$，其中气相质量为 $m_g$，气相质量分数
$x=m_g/m$。忽略界面体积，两相总体积满足
\[
 V=(m-m_g)v_l+m_gv_g,
 \qquad \bar v=\frac Vm=(1-x)v_l+xv_g.
\]
移项得
\begin{equation}
 x=\frac{\bar v-v_l}{v_g-v_l},\qquad 0\le x\le1.
\end{equation}
所以等温两相区的体积变化不是各相任意改变密度，而主要是气相、液相比例改变。
当 $\bar v=v_l$ 时全部为液体；当 $\bar v=v_g$ 时刚好全部成为蒸气。
这也解释了后面含水空气题中，液水尚存与液水恰好消失之后必须使用不同压强表达式。

\subsection{共存曲线的斜率：克拉珀龙方程}

对固定组成、仅有体积功的纯物质，可逆基本关系为
$\dd U=T\dd S-p\dd V$。定义 $G=U+pV-TS$，逐项微分并约去对应项：
\[
 \dd G=(T\dd S-p\dd V)+p\dd V+V\dd p-T\dd S-S\dd T
 =V\dd p-S\dd T.
\]
对同一种物质的单位质量，$g=G/m$ 满足 $\dd g=v\dd p-s\dd T$。
两相平衡要求 $g_1=g_2$；沿共存曲线移动一小步后仍然相等，因此
\[
 v_1\dd p-s_1\dd T=v_2\dd p-s_2\dd T,
 \qquad (v_2-v_1)\dd p=(s_2-s_1)\dd T.
\]
相变可逆且等温时 $L=T(s_2-s_1)$，于是
\begin{equation}
 \frac{\dd p}{\dd T}=\frac{L}{T(v_2-v_1)}.
\end{equation}
两边的单位均为 $\mathrm{Pa/K}$。汽化时 $L>0,v_g>v_l$，共存曲线斜率为正；
普通条件下冰融化为水时体积反而减小，故水的固液共存线斜率为负。
这给出了相图中倾斜方向的物理依据，不应靠记忆图形判断。

在远离临界点的汽化过程中，若 $v_g\gg v_l$、蒸气近似理想气体，
$v_g=RT/(Mp)$，其中 $M$ 是摩尔质量，则
\[
 \frac{\dd\ln p_s}{\dd T}=\frac{ML}{RT^2}.
\]
再在不太宽的温区内把 $L$ 视为常数，从 $T_1$ 积分到 $T_2$ 得
\begin{equation}
 \ln\frac{p_s(T_2)}{p_s(T_1)}
 =\frac{ML}{R}\left(\frac1{T_1}-\frac1{T_2}\right).
\end{equation}
升温时右端为正，饱和蒸气压升高；临界点附近两相比体积接近，不能继续忽略 $v_l$。


\subsection{先判断有没有液相，再决定蒸气分压}
饱和蒸气压 $p_s(T)$ 是两相平衡条件，不是所有蒸气在该温度下必有的压强。若蒸气总量太少，没有液体留下，蒸气分压应由实际物质的量和体积计算；只有液相存在且达到平衡时，分压才被 $p_s(T)$ 固定。
混有不凝结气体时，理想混合近似下被固定的也是蒸气\emph{分压}，不是总压强。图\ref{ch17:fig:s31-22-5}的水平共存线针对纯物质；含空气的容器即使正在凝结，总压强也未必保持不变。

\Needspace{180pt}
\begin{origexamplebox}{湿空气等温压缩的凝结起点}\label{ch17:ex:teaching-condensation}
温度 $T$ 下，体积为 $V_0$ 的密闭容器内有理想干空气和水蒸气，初始无液水。干空气分压为 $p_{a0}$，水蒸气分压为 $\phi p_s(T)$，其中 $0<\phi<1$。缓慢等温压缩至 $V=V_0/r$、$r\ge1$。假定气液平衡及时建立，忽略液水体积及空气在液水中的溶解，求开始凝结时的 $r$、总压强和凝结的水的物质的量。

\ThQCondensation
\end{origexamplebox}
\Needspace{65pt}
\begin{xesolution}
\ThACondensation

凝结前两种气体各自的物质的量不变，分压都随压缩倍数增长：$p_a=rp_{a0}$、$p_v=r\phi p_s$。达到饱和的条件为 $r\phi=1$，故临界压缩倍数为 $r_*=1/\phi$。
凝结后，水分子从气相转入液相，不能继续对气相水蒸气单独使用定量气体的玻意耳关系。于是
\[
 p(r)=rp_{a0}+\begin{cases}r\phi p_s,&1\le r\le1/\phi,\\p_s,&r\ge1/\phi.\end{cases}
\]
初始水总物质的量为 $n_{w0}=\phi p_sV_0/(RT)$；凝结后气相水为 $n_v=p_sV_0/(rRT)$，所以
\[
 n_l=\begin{cases}0,&1\le r\le1/\phi,\\
 \displaystyle\frac{p_sV_0}{RT}\left(\phi-\frac1r\right),&r\ge1/\phi.
 \end{cases}
\]
临界点两段总压强连续、凝结量从零开始且非负。气液两相的水相加仍为 $n_{w0}$，而干空气始终满足 $p_aV=p_{a0}V_0$。总压强在凝结后仍增长，增长来自空气分压。若压缩到液体体积不再可忽略，就必须用扣除液体后的气相空间重算，本式不再适用。
\end{xesolution}

\subsection{从共存线斜率到沸点变化}
沸腾需考虑液体内部气泡的机械平衡。在液柱静压和曲率附加压强均可忽略的平面近似下，沸点满足 $p_s(T_b)=p_{\rm ext}$；它并非脱离外压的固定温度。
前面的共存线斜率确定了改变外压后沸点的变化方向：汽化潜热为正、蒸气比体积更大，所以减压使沸点降低。

\Needspace{180pt}
\begin{origexamplebox}{减压到一半时的沸点}\label{ch17:ex:teaching-boiling}
某纯液体在外压 $p_1$ 下的沸点为 $T_1=360\,\mathrm K$。在所讨论温区内，汽化摩尔潜热近似为 $L_m=40.0\,\mathrm{kJ/mol}$，蒸气服从理想气体关系，液相比体积可忽略。求外压降为 $p_1/2$ 时的沸点，取 $R=8.314\,\mathrm{J/(mol\,K)}$。

\ThQBoiling
\end{origexamplebox}
\Needspace{65pt}
\begin{xesolution}
\ThABoiling

将摩尔潜热直接代入共存线积分式，不能再乘一次摩尔质量：
\[
 \ln\frac{p_2}{p_1}=\frac{L_m}{R}\left(\frac1{T_1}-\frac1{T_2}\right).
\]
本题左边为 $-\ln2$，移项得
\[
 \frac1{T_2}=\frac1{360\,\mathrm K}+\frac{R\ln2}{L_m},
 \qquad T_2\simeq342\,\mathrm K.
\]
减压后温度下降，与共存曲线正斜率一致；$R/L_m$ 的单位为 $\mathrm{K^{-1}}$。若压强变化很小，可线性化为 $\Delta T\simeq RT_1^2\ln(p_2/p_1)/L_m$；本题用倒温度形式直接求解，保留了这一模型积分结果的非线性。近临界区、潜热显著变化或液柱很高时，须恢复被忽略的条件。
\end{xesolution}

\section{范德瓦耳斯方程}

理想气体模型忽略了分子的大小和相互作用。范德瓦耳斯用两个修正把这两种效应纳入物态方程。第一，
分子本身占有体积，分子质心能够活动的有效体积小于容器的几何体积。若每个分子对应的排除体积参数为
$b$，有效体积写成 $V-Nb$。第二，容器内部的分子受到周围分子近似对称的吸引，而靠近器壁的分子受到
的吸引不对称，撞击器壁前会被向内拉回，因此实测压强小于没有吸引力时的压强。单位体积内分子数为
$N/V$，这一修正与分子对数成正比，量级为 $a(N/V)^2$。于是得到
\begin{equation}
  \boxed{\left[p+a\left(\frac{N}{V}\right)^2\right](V-Nb)=Nk_B T}.
  \label{ch17:eq:s31-22-1}
\end{equation}
这里 $a$ 衡量分子间吸引作用，$b$ 衡量分子的有限大小。若以物质的量 $n$ 表示，同一个方程也常写为
\[
  \left(p+a_m\frac{n^2}{V^2}\right)(V-nb_m)=nRT,
\]
此时 $a_m=N_A^2a$、$b_m=N_Ab$ 是摩尔参数，单位分别为
$\mathrm{Pa\,m^6\,mol^{-2}}$ 和 $\mathrm{m^3\,mol^{-1}}$，$N_A$ 为阿伏伽德罗常量。
表~\ref{ch17:tab:s31-22-1}中的 $a,b$ 按传统表头表示这两个摩尔参数，不能与分子形式的数值混用。
排除体积参数也不是分子自身几何体积：硬球模型中单粒子参数是球自身体积的4倍，来自两球球心的不可重叠区域及成对计数。
表~\ref{ch17:tab:s31-22-1}列出若干物质的经验参数；不同物质的参数不同，反映了分子大小和吸引作用
强弱的差异。

\begin{table}[htbp]
  \centering
  \caption{若干物质的范德瓦耳斯参数}\label{ch17:tab:s31-22-1}
  \renewcommand{\arraystretch}{1.12}
  \begin{tabular}{ccc}
    \toprule
    物质 & $a/(\mathrm{Pa\,m^6\,mol^{-2}})$ & $b/(10^{-3}\,\mathrm{m^3\,mol^{-1}})$ \\
    \midrule
    $\mathrm H_2$   & 0.01945 & 0.022 \\
    $\mathrm H_2O$ & 0.56539 & 0.031 \\
    $\mathrm N_2$  & 0.13882 & 0.039 \\
    $\mathrm O_2$  & 0.13983 & 0.032 \\
    $\mathrm{CO}_2$& 0.37186 & 0.043 \\
    \bottomrule
  \end{tabular}
\end{table}

范德瓦耳斯方程仍然是近似模型，但它已经能同时描述稀薄气体、液体的有限体积以及气液相变的主要
特征。它的重要性不只在于提高数值精度，更在于用一个简单方程把气态和液态放进同一幅物理图景中。

\section{气液相变}

把式~(\ref{ch17:eq:s31-22-1})改写为
\[
  p=\frac{Nk_BT}{V-Nb}-a\left(\frac NV\right)^2.
\]
在密度很低、$Nb/V\ll1$ 时，可以把第一项按 $Nb/V$ 展开：
\begin{equation}
  \boxed{p=\frac{Nk_BT}{V}
  +\left(\frac NV\right)^2(k_BT b-a)+\cdots }.
  \label{ch17:eq:s31-22-2}
\end{equation}
第一项就是理想气体压强；第二项同时包含有限分子体积造成的正修正和分子吸引造成的负修正。当温度较高
时，$k_BT b$ 项更重要，有限体积使压强偏大；温度较低时，$a$ 项可能占优势，吸引作用使压强偏小。

图~\ref{ch17:fig:s31-22-2}给出了范德瓦耳斯方程在不同温度下的等温线。温度高于某一特征温度时，等温线随体积增加
单调下降，与普通气体相似；温度降低后，等温线上出现一段波状回环。临界温度 $T_c$ 恰好把这两种情形
分开，临界等温线在点 $C$ 处出现水平拐点。图中从下到上的温度依次为
$T_1=0.85T_c,T_2=0.90T_c,T_c,T_3=1.08T_c,T_4=1.16T_c,T_5=1.24T_c$。
图中先用无量纲比值 $t'=T/T_c$ 标出这六个温度；坐标 $v'=V/V_c$、$p'=p/p_c$
也分别用临界值作单位，临界参数的求法将在下一节推导。

\par\noindent\begin{minipage}{\linewidth}
\vdWIsotherms
\diagramcaption{范德瓦耳斯方程的等温线。低于临界温度时出现回环，$T_c$ 等温线在临界点出现水平拐点。}\label{ch17:fig:s31-22-2}
\end{minipage}\par\medskip

方程给出的整段回环并不全是均匀物质可以稳定经历的状态。真实物质在低于临界温度时等温压缩，达到
饱和蒸气状态后会开始液化；在一段有限体积范围内，液体和蒸气共同存在，压强保持在该温度的饱和蒸气压。
随着体积继续减小，蒸气所占比例下降、液体所占比例上升，直到物质全部变为液体。如何用理论等温线确定
这段恒压相变线，将在后面的麦克斯韦等面积法中讨论。

\section{临界点}

临界点是气液共存曲线的终点。图~\ref{ch17:fig:s31-22-3}(a)中，临界等温线在 $C$ 点具有水平拐点，所以
\[
  \left(\frac{\partial p}{\partial V}\right)_{T_c}=0,
  \qquad
  \left(\frac{\partial^2 p}{\partial V^2}\right)_{T_c}=0.
\]
把范德瓦耳斯方程逐次求导，两个条件为
\[ -\frac{Nk_BT_c}{(V_c-Nb)^2}+\frac{2aN^2}{V_c^3}=0,
\qquad \frac{2Nk_BT_c}{(V_c-Nb)^3}-\frac{6aN^2}{V_c^4}=0.\]
分别移项后相除，得 $2/(V_c-Nb)=3/V_c$，因此 $V_c=3Nb$。
代回第一式，$Nk_BT_c=2aN^2(2Nb)^2/(3Nb)^3=8aN/(27b)$；
最后代回状态方程即可求出临界压强。于是
\begin{equation}
  \boxed{p_c=\frac{a}{27b^2},\qquad
  T_c=\frac{8a}{27k_Bb},\qquad V_c=3Nb}.
  \label{ch17:eq:s31-22-3}
\end{equation}

\par\noindent\begin{minipage}{\linewidth}
\criticalPointFigure
\diagramcaption{临界点的两种表示。(a) 临界等温线的水平拐点；(b) $p$--$T$ 相图中气液共存线止于临界点。}\label{ch17:fig:s31-22-3}
\end{minipage}\par\medskip

图~\ref{ch17:fig:s31-22-3}(b)区分了两条路径：水平的 $P\to M\to Q$ 在 $M$ 穿过气液共存线，会发生相变；
另一条 $P\to A\to B\to Q$ 先增压，再进入高于临界温度的区域，绕过临界点而不穿过共存线。
沿后一条路径，物质的密度可以连续变化，途中不会出现突变
的气液界面。由此可见，气体和液体并不是两种在任何条件下都截然分开的物态；在临界点以上，二者可以
连续过渡。

在透明高压容器中观察临界过程，会看到图~\ref{ch17:fig:s31-22-4}所示的现象。低于临界温度时，容器内有清晰的液面；接近
临界温度时，液体与蒸气的密度趋于相同，密度涨落对光产生强烈散射，系统呈现乳白色浑浊，称为临界乳光；
继续升温，使温度和压强都超过临界值后，界面消失，容器中成为均匀的超临界流体。
仅有 $T>T_c$ 并不足以称为超临界流体，还要有 $p>p_c$；若温度虽高而压强很低，仍是低密度气态。
下图表示选定充装量下经过临界区域的三个阶段，
使界面、密度涨落和均匀状态可以直接比较。

\par\noindent\begin{minipage}{\linewidth}
\criticalApparatusFigure
\diagramcaption{临界点附近的观察情景与状态示意。左：清晰气液界面；中：临界乳光；右：均匀超临界流体。}\label{ch17:fig:s31-22-4}
\end{minipage}\par\medskip

临界点附近还有一个重要特征：等温压缩系数
\[
  \kappa_T=-\frac1V\left(\frac{\partial V}{\partial p}\right)_T
\]
变得很大，并在理想临界模型中发散。这意味着很小的压强变化就能引起很大的密度变化，也解释了临界点
附近强烈的密度涨落。

把压强、体积和温度分别用临界值作单位，定义约化变量
\begin{equation}
  \boxed{p'=\frac{p}{p_c},\qquad v'=\frac{V}{V_c},\qquad t'=\frac{T}{T_c}}.
  \label{ch17:eq:s31-22-4}
\end{equation}
代入时，$p+aN^2/V^2=[a/(27b^2)](p'+3/v'^2)$，
$V-Nb=Nb(3v'-1)$，$Nk_BT=8aNt'/(27b)$。
两边约去共同因子 $aN/(27b)$，参数 $a,b$ 完全消失，得到
\begin{equation}
  \boxed{\left(p'+\frac{3}{v'^2}\right)(3v'-1)=8t'}.
  \label{ch17:eq:s31-22-5}
\end{equation}
这说明在范德瓦耳斯模型中，不同物质只要用各自临界参数作单位，就具有同一形式的物态方程和相似
的相变行为，这一结论称为对应态定律。真实物质不完全服从简单的范德瓦耳斯方程，但约化变量和对应
态思想仍然非常有用。

\section{低于临界温度的等温线}

在图~\ref{ch17:fig:s31-22-5}中，蓝色曲线是某一 $T<T_c$ 时范德瓦耳斯方程给出的等温线。真实物质从体积较小的液态
$M$ 出发等温膨胀时，不会沿曲线走过 $M\to A\to P\to B\to N$ 的全部波状部分，而是沿水平直线
$M\to P\to N$ 发生气液共存相变。$M$ 是饱和液体，$N$ 是饱和蒸气，水平线压强 $p_T$ 就是温度 $T$
下的饱和蒸气压。

\par\noindent\begin{minipage}{\linewidth}
\maxwellAreaFigure
\diagramcaption{麦克斯韦等面积作图。水平相变线的位置由上下两块面积相等确定。}\label{ch17:fig:s31-22-5}
\end{minipage}\par\medskip

水平线的位置不能任意选择。设饱和液体和饱和蒸气的体积分别为 $V_l,V_g$。
令亥姆霍兹自由能 $F=U-TS$，在固定温度和粒子数下 $\dd F=-p\dd V$。
对模型的均匀态自由能作解析延拓，得到
\[F(V_g)-F(V_l)=-\int_{V_l}^{V_g}p(V)\dd V.\]
延拓曲线中的不稳定段不是实际可走过的稳定准静态过程。
平衡共存要求两相的吉布斯自由能 $G=F+p_TV$ 相等（比较同样粒子数），
即 $F(V_g)-F(V_l)=-p_T(V_g-V_l)$。二式比较就得到等面积关系。
也可从内能是状态函数理解这个形式积分：模型延拓给出
\[
  \Delta_1U=U(N)-U(M)
  =T\,[S(N)-S(M)]-\int_{V_l}^{V_g}p(V)\,\dd V.
\]
若沿真实的恒压相变线从 $M$ 到 $N$，同一对初、末态的内能改变量为
\[
  \Delta_2U=T\,[S(N)-S(M)]-p_T(V_g-V_l).
\]
内能是状态函数，$\Delta_1U=\Delta_2U$，所以
\begin{equation}
  \boxed{\int_{V_l}^{V_g}\bigl[p(V)-p_T\bigr]\,\dd V=0}.
  \label{ch17:eq:s31-22-6}
\end{equation}
也就是说，水平线上方的面积 $S_2$ 必须等于水平线下方的面积 $S_1$。这就是麦克斯韦等面积法。它把
理论方程中不稳定的回环部分替换为真实的恒压气液共存线，同时保证相变两端具有相同的热力学势。

\section{过冷气体与过热液体}

虽然平衡态相变线由麦克斯韦等面积法确定，但实际系统不一定恰好在到达饱和状态时立即相变。液滴的
形成需要一个凝结核，气泡的形成也需要一个成核中心；如果容器非常洁净，系统中缺少尘埃、划痕、自由
表面等成核条件，蒸气可以越过饱和点而暂时保持气态，液体也可以越过沸点而暂时保持液态。这些状态分别
称为过冷蒸气（或过饱和蒸气）和过热液体，它们是亚稳态。
在图~\ref{ch17:fig:s31-22-5}中，$M\to A$ 是过热液体分支，$N\to B$ 是过饱和蒸气分支。
它们尚满足 $(\partial p/\partial V)_T<0$，即等温压缩系数为正；但相分离后的自由能更低。
极值点 $A,B$ 是失稳界限（spinodal）；其间 $A\to P\to B$ 有
$(\partial p/\partial V)_T>0$，压缩系数为负，连局部力学稳定性也不满足。

亚稳态并非真正稳定。一旦出现凝结核或气泡核，相变会迅速发生。过冷蒸气遇到尘埃、离子或飞机尾气中
的微粒，就可能凝结成小液滴或冰晶。飞机在高空留下的白色尾迹，常常就是尾气中的水蒸气在低温空气中
凝结、冻结形成的。云的形成也与空气上升冷却并达到水蒸气饱和有关。

图~\ref{ch17:fig:s31-22-6}中，尾迹从发动机后方开始，随混合扩散逐渐展宽。它并非水蒸气本身的颜色，而是凝结、冻结后的微小颗粒散射光所形成的可见轨迹。
\par\noindent\begin{minipage}{\linewidth}
\contrailFigure
\diagramcaption{飞机尾迹的形成情景。可见尾迹来自水蒸气在微粒上凝结或凝华。}\label{ch17:fig:s31-22-6}
\end{minipage}\par\medskip

亚稳态对粒子探测也十分重要。云室利用过饱和蒸气：带电粒子通过时沿途产生离子，水滴在离子上凝结，
于是粒子轨迹变成可见雾迹。气泡室则利用过热液体：带电粒子经过时触发微小气泡，形成可见径迹。云室
和气泡室曾经是发现基本粒子的重要工具，与这些技术有关的工作获得过诺贝尔物理学奖。它们共同利用了
同一原理：把系统准备在接近失稳的亚稳态，让微小扰动被迅速放大成可见的宏观结构。

\section{表面张力、附加压强与润湿}

\subsection{先数界面，再列受力}

把液体表面面积记为 $A$。在温度和组成给定时，缓慢增加面积所需的可逆功为
$\delta W_{\rm on}=\sigma\dd A$，$\sigma$ 的单位既可写成 $\mathrm{J/m^2}$，
也可写成 $\mathrm{N/m}$。面积功与表面内能并不总是同一件事：若 $\sigma$ 随温度改变，
拉伸时还可能吸、放热，相关熵项已在第二定律的表面工质题中讨论。

考虑半径 $R$ 的球形液滴，或忽略膜厚的球形肥皂泡。沿赤道切开，取上半部作为受力对象。
表面张力在赤道边缘沿球面切线作用，其竖直分量向下。图\ref{ch17:ref:fig-laplace-cut}
用竖直剖面标明切口；三个向上的箭头表示压差在水平投影面上的合力方向，
并不表示压强只作用在三个点上。

\begin{center}\begin{minipage}{\linewidth}\centering
\begin{tikzpicture}[x=1cm,y=1cm,>=Latex,font=\small]
 \draw[Ink,thick] (-1.5,0) arc[start angle=180,end angle=0,radius=1.5];
 \draw[TextGray,dashed] (-1.5,0)--(1.5,0);
 \fill[Ink] (0,0) circle (1pt);
 \draw[<->,TextGray] (0,-.25)--(1.5,-.25);
 \node[below] at (.75,-.25) {$R$};
 \foreach \x in {-.75,0,.75}{\draw[->,DeepBlue,thick] (\x,.10)--(\x,.82);}
 \draw[->,DeepBlue,thick] (-1.5,0)--(-1.5,-.72);
 \draw[->,DeepBlue,thick] (1.5,0)--(1.5,-.72);
 \node[anchor=east] at (-1.72,-.35) {边缘张力};
 \node[anchor=west] at (1.75,.75) {$\Delta p\,\pi R^2$ 向上};
 \node[anchor=west] at (1.75,-.35) {$k\sigma\,2\pi R$ 向下};
 \node at (0,1.85) {上半球受力的竖直剖面};
\end{tikzpicture}
\diagramcaption{球面沿赤道截开。$k$ 是界面数；每个界面在赤道各处都贡献向下的切向拉力。}\label{ch17:ref:fig-laplace-cut}
\end{minipage}\end{center}

压差的合力等于 $\Delta p\,\pi R^2$，不是压差乘以半球曲面面积。
若有 $k$ 个具有相同张力系数、半径近似相同的界面，边缘张力合力为 $k\sigma\,2\pi R$。
平衡条件给出
\begin{equation}
 \Delta p\,\pi R^2=k\sigma\,2\pi R
 \quad\Longrightarrow\quad
 \Delta p=\frac{2k\sigma}{R}.
\end{equation}
液滴只有一个液气界面，$k=1$；薄肥皂泡有内外两个液气界面，$k=2$，故分别得到
$2\sigma/R$ 和 $4\sigma/R$。液体内部单个气泡通常也只有一个气液界面，不能因名称中有“泡”
便取 $k=2$。这些公式忽略重力造成的明显形变；$R$ 增大时附加压强减小，平面极限趋于零。

\subsection{接触角如何决定液面升降}

接触角 $\theta$ 在液体一侧量取，是固体壁面与液气界面切平面之间的夹角。
圆管内半径为 $r$，液体密度为 $\rho$，两侧上方气压相同。管壁一周的表面张力合力
竖直分量为 $2\pi r\sigma\cos\theta$。当管径足够小、液面弯曲部分的体积可相对
$\pi r^2h$ 忽略时，与高出外液面的液柱重力平衡：
\[
 2\pi r\sigma\cos\theta=\rho g\pi r^2h,
 \qquad h=\frac{2\sigma\cos\theta}{\rho gr}.
\]
$h$ 是有符号的液面高度差。润湿时 $\theta<90^\circ$，$h>0$，液面升高；
不润湿时 $\theta>90^\circ$，$h<0$，液面降低。若先假定表面张力必向上，再把负号丢掉，
就会把凸液面画成错误的上升液柱。接触角、凹凸形状与液面升降必须同时核对。


\Needspace{10\baselineskip}
\section{综合例题}
\begin{origexamplebox}{【饱和蒸气】}\label{ch17:ex:s31-22-1}
一端封闭的均匀管子内有空气和水银蒸气，被 $10\,\mathrm{cm}$ 长的水银柱封闭（不论管子如何放置）。
当管子水平放置时，空气柱稳定长度为 $17\,\mathrm{cm}$；管子竖直放置时，空气柱稳定长度为
$20\,\mathrm{cm}$。两次均在同一温度下达到平衡。若大气压强为 $76\,\mathrm{cmHg}$，求水银蒸气的饱和蒸气压。

\ThQMercury
\end{origexamplebox}

\begin{xesolution}
水银蒸气与液态水银共存，所以只要温度不变，它的分压始终等于饱和蒸气压 $p_s$；能使用波义耳定律的
是管内那部分物质的量固定的空气，而不是“空气与水银蒸气”的混合气体整体。

水平放置时，水银柱两侧没有高度差，封闭气体总压强等于大气压强 $76\,\mathrm{cmHg}$，故空气分压为
$76-p_s$。竖直放置并使封闭端在上时，下方大气对水银柱向上作用，上方封闭气体的压力与重力向下。
由 $p_{\rm atm}A=p_{\rm trap}A+\rho_{\rm Hg}g\ell A$，其总压强为
$76-10=66\,\mathrm{cmHg}$，空气分压为 $66-p_s$。两种放置方式与各段长度的对应关系如下图。

\par\noindent\begin{minipage}{\linewidth}
\mercuryTubeFigure
\diagramcaption{水平和竖直放置时的封闭气柱与水银柱。}\label{ch17:ref:fig-mercury}
\end{minipage}\par\medskip

图~\ref{ch17:ref:fig-mercury} 中竖直气柱长20厘米、水银柱长10厘米，图上的共同长度比例与题设一致。
管子截面积不变，空气体积与空气柱长度成正比。对空气列波义耳定律：
\[
  (76-p_s)\times17=(66-p_s)\times20.
\]
解得
\[
  1292-17p_s=1320-20p_s,
  \qquad
  \boxed{p_s=\frac{28}{3}\,\mathrm{cmHg}\approx9.33\,\mathrm{cmHg}}.
\]
若把竖直图误画成“开口在上、封闭端在下”，封闭气体压强就会变成 $86\,\mathrm{cmHg}$，与题给长度
由 $17\,\mathrm{cm}$ 增至 $20\,\mathrm{cm}$ 的事实矛盾；因此图中的上下位置也是方程的一部分。
\end{xesolution}

\Needspace{9\baselineskip}
\begin{origexamplebox}{【内能变化】}\label{ch17:ex:s31-22-2}
在 $80.2\,{}^{\circ}\mathrm C$ 和 $1\,\mathrm{atm}$ 下，液体苯（$\mathrm{C_6H_6}$）蒸发为苯蒸气。
已知在此条件下苯的汽化热 $L=94.4\,\mathrm{cal/g}$，苯蒸气符合理想气体物态方程，且蒸气体积远大于
液体体积。求内能的改变。

\ThQVaporization
\end{origexamplebox}

\begin{xesolution}
\ThAVaporization

取质量为 $m$ 的苯为系统。等压汽化过程中系统吸收的热量是 $Q=mL$。热力学第一定律给出
\[
  Q=\Delta U+W,
  \qquad W=p(V_g-V_l).
\]
因为 $V_g\gg V_l$，可忽略液体体积；苯蒸气视为理想气体，若物质的量为 $n=m/M$，则
\[
  W\approx pV_g=nRT=\frac{m}{M}RT.
\]
所以单位质量苯汽化时的内能增量为
\[
  \frac{\Delta U}{m}=L-\frac{RT}{M}.
\]
取 $T=80.2+273.15=353.35\,\mathrm K$，$M=78.1\,\mathrm{g/mol}$，
$R=1.987\,\mathrm{cal/(mol\,K)}$，得到
\[
  \boxed{\frac{\Delta U}{m}
  =94.4-\frac{1.987\times353.35}{78.1}
  \approx85.4\,\mathrm{cal/g}}.
\]
相应的摩尔内能增量为
\[
  \boxed{\Delta U_m=ML-RT
  \approx6.67\times10^3\,\mathrm{cal/mol}
  \approx27.9\,\mathrm{kJ/mol}}.
\]
汽化热并不全部变成内能，其中约 $RT$ 用来克服外界压强完成膨胀功。
\end{xesolution}

\Needspace{11\baselineskip}
\begin{origexamplebox}{【气体膨胀】}\label{ch17:ex:s31-22-3}
一汽缸的初始体积为 $V$，其中盛有 $2\,\mathrm{mol}$ 空气和少量水（水的体积可以忽略）。平衡时气体
总压强为 $3.0\,\mathrm{atm}$。经等温膨胀使体积变为 $2V$，在膨胀结束时水刚好全部消失，此时总压强
为 $2.0\,\mathrm{atm}$。试计算：
\begin{enumerate}[label=\arabic*.]
  \item 汽缸中气体的温度；
  \item 若继续作等温压缩，使体积变为 $V/2$，汽缸中气体的总压强；
  \item 若继续作等温膨胀，使体积变为 $4V$，汽缸中气体的总压强。
\end{enumerate}

\ThQWetExpansion
\end{origexamplebox}

\begin{xesolution}
在液态水尚未完全消失时，水蒸气分压恒等于该温度的饱和蒸气压 $p_s$；空气物质的量始终是
$2\,\mathrm{mol}$，记为 $n_a$，其分压服从理想气体方程。初态和体积为 $2V$ 的末态分别满足
\[
  3.0\,\mathrm{atm}=\frac{n_aRT}{V}+p_s,
  \qquad
  2.0\,\mathrm{atm}=\frac{n_aRT}{2V}+p_s,
\]
两式相减得到
\[
  \frac{n_aRT}{2V}=1.0\,\mathrm{atm},
  \qquad p_s=1.0\,\mathrm{atm}.
\]

\par\noindent\begin{minipage}{\linewidth}
\mixedGasPVFigure
\diagramcaption{含水空气的分段等温过程。各状态点由分压方程确定。}\label{ch17:ref:fig-mixed}
\end{minipage}\par\medskip

因此温度为
\[
  \boxed{T\approx373.15\,\mathrm K\quad(t\approx100\,{}^\circ\mathrm C)}.
\]
这是由水在 $1\,\mathrm{atm}$ 饱和蒸气压下的温度确定的，不需要预先给出体积。
要注意：这里是水蒸气的分压，不是初态的总压强。反过来可求
\[V=\frac{n_aRT}{2.0\,\mathrm{atm}}\approx30.6\,\mathrm L.\]

在体积 $2V$ 处，水刚好全部成为蒸气，水蒸气的物质的量为
\[
  n_w=\frac{p_s(2V)}{RT}=2\,\mathrm{mol}.
\]

继续压缩到 $V/2$ 时，部分水蒸气凝结，故水蒸气分压仍为 $1.0\,\mathrm{atm}$。
空气分压为 $p_a=n_aRT/(V/2)=4.0\,\mathrm{atm}$，所以总压强为
\[
  \boxed{p_{V/2}=p_a+p_s=5.0\,\mathrm{atm}}.
\]

从 $2V$ 继续膨胀时已经没有液态水，空气和水蒸气的总物质的量固定为 $2+2=4\,\mathrm{mol}$。在
$4V$ 处
\[
  \boxed{p_{4V}=\frac{(n_a+n_w)RT}{4V}=1.0\,\mathrm{atm}}.
\]
图~\ref{ch17:ref:fig-mixed} 中蓝色实线表示 $V'\le2V$ 时的 $p=p_s+n_aRT/V'$，虚线表示 $V'\ge2V$ 时的 $p=(n_a+n_w)RT/V'$；两段在
$V'=2V$ 处连续相接。
\end{xesolution}

\Needspace{11\baselineskip}
\begin{origexamplebox}{【高压锅】}\label{ch17:ex:s31-22-4}
考虑下述理想化压力锅模型：先使锅内水沸腾并排除空气，再由限压阀 $S$ 密封，此时可以
认为锅内只有水的饱和蒸气，空气已经全部排除；然后继续加热，直到压力阀被锅内水蒸气顶起，锅内达到
设计温度。现有一压力锅，按正确方法加热能达到的预期温度为 $120\,{}^{\circ}\mathrm C$，大气压为
$1.01\times10^5\,\mathrm{Pa}$。

若改为在加热前、锅内温度为 $10\,{}^{\circ}\mathrm C$ 时就封闭压力阀，那么加热
到压力阀刚被顶起时，锅内水的温度是多少？若继续加热，锅内水的温度最高可达多少？假设空气不溶于水。
已知水的饱和蒸气压 $p_w(t)$ 与摄氏温度的关系如图所示。始终保留液态水；本题模型不作为实际器具的操作说明。
\par\noindent\begin{minipage}{\linewidth}
\pressureCookerQuestionFigure
\diagramcaption{限压阀与水的饱和蒸气压曲线；曲线给出题目所需的读数条件。}\label{ch17:ref:fig-cooker-given}
\end{minipage}\par\medskip
\end{origexamplebox}

\begin{xesolution}
压力阀是否顶起取决于锅内外压强差。正确使用时锅内只有饱和水蒸气，并在 $120\,{}^{\circ}\mathrm C$
顶起，因此阀门对应的锅内绝对压强阈值就是
\[
  p_*=p_w(120\,{}^{\circ}\mathrm C)\approx1.99\times10^5\,\mathrm{Pa}.
\]

错误使用时，$10\,{}^{\circ}\mathrm C$ 就把阀门装上，锅内同时封入空气和少量水蒸气。设锅内容积不变，
初始空气分压为
\[
  p_{a0}=p_{\rm atm}-p_w(10\,{}^{\circ}\mathrm C).
\]
在阀门开启前空气物质的量不变，故
\[
  p_a(t)=p_{a0}\frac{t+273.15}{283.15}.
\]
锅内始终有水，水蒸气分压为 $p_w(t)$。首次顶起条件是
\[
  p_w(t)+\bigl[p_{\rm atm}-p_w(10)\bigr]
  \frac{t+273.15}{283.15}=p_w(120).
\]

\par\noindent\begin{minipage}{\linewidth}
\pressureCookerFigure
\diagramcaption{限压装置及首次开启的压强条件。}\label{ch17:ref:fig-cooker}
\end{minipage}\par\medskip

由曲线读数 $p_w(10)\approx1.2\times10^3\,\mathrm{Pa}$、
$p_w(120)\approx1.99\times10^5\,\mathrm{Pa}$，联立上式，即在图~\ref{ch17:ref:fig-cooker} 中找总压强曲线与水平阈值线的交点，得到
\[
  \boxed{t_{\rm lift}\approx90\,{}^{\circ}\mathrm C}.
\]
此时饱和蒸气压约 $7.0\times10^4\,\mathrm{Pa}$，空气分压约 $1.29\times10^5\,\mathrm{Pa}$，二者之和
已经达到阀门阈值，所以水尚未达到 $100\,{}^{\circ}\mathrm C$，阀门便会被顶起。

阀门开启后，锅内总压强被限制在 $p_*$ 附近，空气与蒸气的混合气不断逸出。继续升温时 $p_w(t)$ 增加，
空气分压 $p_a=p_*-p_w(t)$ 相应减小；到 $120\,{}^{\circ}\mathrm C$ 时，$p_w=p_*$，锅内空气已趋于排尽。
再供热只会使水猛烈沸腾并从阀门排出蒸气，因而
\[
  \boxed{t_{\max}=120\,{}^{\circ}\mathrm C}.
\]
图中的细虚线表示封闭空气分压随绝对温度线性增加，蓝线表示饱和蒸气压；二者相加得到灰色总压强曲线。
总压强曲线与 $p_*$ 水平线的交点决定首次开启温度。超过交点后，图中虚线仅表示假设阀门仍未开启的外推，
实际压强受限于水平的 $p_*$。$120\,{}^\circ\mathrm C$ 是模型中空气排尽后的上限，前提是液态水尚存。
\end{xesolution}

\par\medskip
\begin{origexamplebox}{【奇怪的过程】}\label{ch17:ex:s31-22-5}
用不导热细管连接的两个相同容器里装有压强为 $1\,\mathrm{atm}$、相对湿度为 $50\%$、温度为
$100\,{}^{\circ}\mathrm C$ 的空气。现将其中一个容器浸在 $0\,{}^{\circ}\mathrm C$ 的冰中，另一个容器
仍维持在 $100\,{}^{\circ}\mathrm C$。忽略凝结水所占体积，并采用普通气体的扩散平衡近似，
不考虑细管内温差引起的热扩散效应。问系统的压强改变为多大？每一容器中的相对湿度为多少？已知
$0\,{}^{\circ}\mathrm C$ 时水的饱和蒸气压为 $4.6\,\mathrm{mmHg}$。

\ThQTwoVessels
\end{origexamplebox}

\begin{xesolution}
在 $100\,{}^{\circ}\mathrm C$ 时水的饱和蒸气压约为 $1\,\mathrm{atm}$，所以初态相对湿度 $50\%$
对应的水蒸气分压和干空气分压均为 $0.5\,\mathrm{atm}$。设每个容器体积为 $V$。

冷端降到 $0\,{}^{\circ}\mathrm C$ 后，水蒸气在那里凝结，整个连通系统中的水蒸气分压被限制为
\[
  p_v=p_s(0)=4.6\,\mathrm{mmHg}=0.00605\,\mathrm{atm}.
\]
因为细管连通，两端同一种气体的分压相等；但两端温度不同，同样分压所对应的物质的量不同。

\par\noindent\begin{minipage}{\linewidth}
\twoVesselsFigure
\diagramcaption{冷热容器通过开放管口连通；凝结水的厚度只用于辨识液相，不按比例，计算中忽略其体积。}\label{ch17:ref:fig-vessels}
\end{minipage}\par\medskip
图~\ref{ch17:ref:fig-vessels} 中连通管位于气相区域，冷端凝结水在管口下方，两端分别维持给定温度。

干空气的总物质的量守恒。初态两容器中共有
\[
  n_a=2\frac{(0.5p_{\rm atm})V}{R(373\,\mathrm K)}
  =\frac{p_{\rm atm}V}{R(373\,\mathrm K)}.
\]
末态设干空气共同分压为 $p_a$，则
\[
  n_a=\frac{p_aV}{R(373\,\mathrm K)}
      +\frac{p_aV}{R(273\,\mathrm K)}.
\]
两式相等，解得
\[
  p_a=p_{\rm atm}\frac{273}{273+373}
  \approx0.4226\,\mathrm{atm}.
\]
因此系统末态总压强为
\[
  p_f=p_a+p_v\approx0.4226+0.00605=0.4287\,\mathrm{atm}
  \approx326\,\mathrm{mmHg}.
\]
相对初态的压强改变为
\[
  \boxed{\Delta p=p_f-p_i\approx-0.571\,\mathrm{atm}
  \approx-434\,\mathrm{mmHg}\approx-57.9\,\mathrm{kPa}}.
\]

冷容器内有凝结水，水蒸气达到饱和，所以
\[
  \boxed{\varphi_{0{}^\circ\mathrm C}=100\%}.
\]
热容器仍为 $100\,{}^{\circ}\mathrm C$，其饱和蒸气压为 $760\,\mathrm{mmHg}$，实际水蒸气分压只有
$4.6\,\mathrm{mmHg}$，故
\[
  \boxed{\varphi_{100{}^\circ\mathrm C}
  =\frac{4.6}{760}\times100\%\approx0.61\%}.
\]
在忽略热扩散、细管温差效应并采用普通气体扩散平衡的模型中，细管使两端分压相同，却不使两端相对湿度相同；相对湿度还取决于当地温度对应的饱和蒸气压。
\end{xesolution}

\Needspace{12\baselineskip}
\begin{origexamplebox}{【云的形成】}\label{ch17:ex:s31-22-6}
以下是一个形成云的简单模型。热空气带着地球表面的水蒸气上升，当升到一定高度时，水蒸气达到饱和并
开始凝结成水滴，云就产生了。假设空气是由氧和氮构成的理想气体，比例为 $21:79$；令 $g$ 为重力加速度，
忽略它随高度的变化。凝结前采用静力平衡、恒定组成且比熵处处相同的干绝热大气柱模型；水蒸气含量很少，
对空气平均摩尔质量和绝热系数的影响可忽略。上升气团与周围压强近似相等，凝结前作可逆绝热变化。
以下温度递减率与压强分布只用于开始凝结以前；不得将干绝热关系直接用于凝结放热后的阶段。在这些假设下可证明
\[
  p=p_0\left(\frac{T_0-\Gamma z}{T_0}\right)^\alpha,
\]
其中 $p_0,T_0$ 分别是海平面 $z=0$ 处的压强和温度，$\Gamma>0$ 是温度随高度的递减率。
\begin{enumerate}[label=\arabic*.]
  \item 用 $\gamma,R,g,m_a$ 给出 $\Gamma$，其中 $\gamma$ 为绝热系数，$m_a$ 为空气的摩尔质量；
  \item 估计高度每升高 $1\,\mathrm{km}$ 时温度的改变量；
  \item 用 $\gamma$ 表示指数 $\alpha$；
  \item 在该模型中大气层的总厚度为多少；
  \item 饱和蒸气压随温度的变化满足克拉珀龙方程
  \[
    \frac{\dd p_s}{\dd T}=\frac{L}{T(v_2-v_1)},
  \]
  其中 $L$ 为摩尔相变潜热（单位 $\mathrm{J/mol}$），$v_1,v_2$ 是两相的摩尔体积。试在液体体积远小于蒸气体积的极限下求
  饱和蒸气压与温度的关系；
  \item 找出水蒸气开始凝结的条件。
\end{enumerate}

\ThQCloud
\end{origexamplebox}

\begin{xesolution}
\textbf{1. 温度递减率。} 采用静力平衡、恒定组成及等熵的大气柱近似，气团与周围压强近似相等。
这比单说“一个气团绝热上升”更强；以下线性递减率是这一模型的结果。
高度为 $z$ 处空气密度为
\[
  \rho=\frac{m_a p}{RT}.
\]
静力平衡给出
\[
  \frac{\dd p}{\dd z}=-\rho g=-\frac{m_ag}{RT}p.
\]
另一方面，理想气体的可逆绝热关系为
\[
  p\propto T^{\gamma/(\gamma-1)}.
\]
对数微分后与静力平衡式比较，并利用 $\dd T/\dd z=-\Gamma$，得到
\[
  \boxed{\Gamma=\frac{\gamma-1}{\gamma}\frac{m_ag}{R}}.
\]

\textbf{2. 数值估计。} 空气平均摩尔质量
\[
  m_a\approx0.21\times32+0.79\times28=28.8\,\mathrm{g/mol}.
\]
取 $\gamma\approx1.40$、$g\approx9.8\,\mathrm{m/s^2}$，有
\[
  \Gamma\approx9.7\times10^{-3}\,\mathrm{K/m}
  \approx9.7\,\mathrm{K/km}.
\]
所以在这个“干绝热”模型中，高度每升高 $1\,\mathrm{km}$，温度约降低
\[
  \boxed{\Delta T\approx-9.7\,\mathrm K}.
\]

\textbf{3. 压强指数。} 由绝热关系直接得到
\[
  \boxed{\alpha=\frac{\gamma}{\gamma-1}}.
\]
也可把题给表达式对 $z$ 求导，与静力平衡式比较，得到 $\alpha\Gamma=m_ag/R$，再代入上一问的
$\Gamma$，结果相同。

\textbf{4. 模型大气厚度。} 题给线性温度模型在 $T=T_0-\Gamma z=0$ 时终止，因此
\[
  \boxed{H=\frac{T_0}{\Gamma}}.
\]
若取 $T_0\approx288\,\mathrm K$，则 $H\approx29.7\,\mathrm{km}$。真实大气当然不会在这个高度突然
终止；这一结果只表明“恒定干绝热递减率”不能无限向上外推。

\textbf{5. 饱和蒸气压。} 取汽化潜热 $L_m>0$，用 $v_g\gg v_l$ 写克拉珀龙方程：
\[
  \frac{\dd p_s}{\dd T}
  =\frac{L_m}{T(v_g-v_l)}
  \approx\frac{L_m}{Tv_g}.
\]
饱和蒸气近似为理想气体，$v_g=RT/p_s$，故
\[
  \frac{1}{p_s}\frac{\dd p_s}{\dd T}
  =\frac{L_m}{RT^2}.
\]
若在所研究温区把 $L_m$ 近似为常量，从参考温度 $T_0$ 积分到 $T$：
\[
  \boxed{p_s(T)=p_s(T_0)
  \exp\!\left[\frac{L_m}{R}\left(\frac1{T_0}-\frac1T\right)\right]}.
\]
这就是克劳修斯--克拉珀龙近似关系，它表明饱和蒸气压随温度下降得非常快。

\textbf{6. 凝结条件。} 在开始凝结以前，水蒸气与干空气一起绝热上升，水蒸气摩尔分数 $x_v$ 近似
不变，因而
\[
  p_v(z)=x_vp(z)=p_{v0}\left(\frac{T}{T_0}\right)^\alpha,
  \qquad T=T_0-\Gamma z.
\]
水蒸气首次凝结的条件是实际分压等于该温度的饱和蒸气压：
\[
  \boxed{p_{v0}\left(\frac{T_d}{T_0}\right)^\alpha=p_s(T_d)}.
\]
若海平面相对湿度为 $\varphi_0$，则 $p_{v0}=\varphi_0p_s(T_0)$，条件可写成
\[
  \boxed{\varphi_0\left(\frac{T_d}{T_0}\right)^\alpha
  =\exp\!\left[\frac{L_m}{R}
  \left(\frac1{T_0}-\frac1{T_d}\right)\right]}.
\]
求出 $T_d$ 后，云底高度为
\[
  \boxed{z_d=\frac{T_0-T_d}{\Gamma}}.
\]

\par\noindent\begin{minipage}{\linewidth}
\par\noindent\begin{minipage}{\linewidth}
\cloudFormationFigure
\diagramcaption{未凝结气团上升膨胀；实际分压与饱和蒸气压相等时到达云底。}\label{ch17:ref:fig-cloud}
\end{minipage}\par\medskip
\end{minipage}\par\smallskip

图~\ref{ch17:ref:fig-cloud} 右侧以 $T_0=288\,\mathrm K,\varphi_0=0.6,\alpha=3.5,L_m=44\,\mathrm{kJ/mol}$ 为一组示例参数，
空气团随上升冷却从右向左运动；两曲线第一次相交时开始凝结，
左图对应的高度就是云底。图形把“压强降低”与“饱和蒸气压下降得更快”这两个同时发生的效应放在一起。
凝结以后释放潜热，气相水的物质的量也改变，不能继续沿未凝结分支外推。
\end{xesolution}
}

\endgroup

\hypersetup{bookmarksdepth=2}

\hypersetup{bookmarksdepth=3}
\begingroup
\renewcommand{\Needspace}[1]{\par\begingroup
\dimen0=#1\relax\vskip0pt plus\dimen0\penalty-100
\vskip0pt plus-\dimen0\vskip\dimen0\penalty9999\vskip-\dimen0
\vskip0pt\endgroup}
\chapter{统计物理基础}

\section{概率、随机变量与涨落}

一次随机试验的全部互斥且完备的结果组成样本空间（sample space）。事件是样本空间的子集；互斥事件满足概率加法，
独立事件 $A,B$ 满足
\begin{equation}
 P(A\cap B)=P(A)P(B).
\end{equation}
离散随机变量 $X$ 的期望、方差分别为
\begin{equation}
 \langle X\rangle=\sum_i p_iX_i,
 \qquad
 \operatorname{Var}(X)=\left\langle(X-\langle X\rangle)^2\right\rangle.
\end{equation}
连续变量只需把求和改成对归一化概率密度的积分。若 $X_1,\ldots,X_N$ 相互独立且同分布，则
\[
 \operatorname{Var}\!\left(\sum_{i=1}^NX_i\right)=N\operatorname{Var}(X_1),
\]
而总量的平均值正比于 $N$，所以相对涨落通常按 $N^{-1/2}$ 衰减。这是宏观量看起来稳定、微观量却始终涨落的数学原因。

随机变量变换时不能只把 $x$ 换成 $y(x)$。若一个 $y$ 对应若干个根 $x_j$，概率密度应写成
\begin{equation}
 p_Y(y)=\sum_jp_X(x_j)\left|\frac{\dd x_j}{\dd y}\right|.
\end{equation}
绝对值是变量变换的雅可比因子；求速率分布时出现的 $4\pi v^2$，正是三维速度空间球壳的几何雅可比。
\subsection{从事件到条件概率}

基本事件是一次试验不可再细分的结果；不同基本事件互斥，所有基本事件的并集是样本空间。
概率满足 $P(A)\ge0$、$P(\Omega)=1$。若 $A,B$ 互斥，$P(A\cup B)=P(A)+P(B)$；
一般情形则须扣除重复计算的交集：$P(A\cup B)=P(A)+P(B)-P(A\cap B)$。
有限个基本事件等可能时，概率等于有利结果数除以总结果数，这叫古典概率模型。
例如公平骰子掷出偶数的概率是 $3/6=1/2$。连续区域内均匀取点时，概率等于长度、面积或体积之比，
这是几何概率模型；必须先说明对哪一种几何测度均匀。
连续样本空间的事件须可测。例如在 $[0,1]$ 上均匀取点，“取到有理数”是可测事件，概率为零，
因为有理数可数且每个点的概率为零；概率为零不等于集合为空。某些不可测集合则不能在该测度下直接赋予概率。
多次试验的联合结果是一个数组 $(\omega_1,\ldots,\omega_n)$，联合样本空间是各空间的直积；
只有另加独立性假设，联合概率才是各分量概率的乘积。

已经知道事件 $A$ 发生，相当于把样本空间缩小到 $A$，并重新归一化。只要 $P(A)>0$，
\begin{equation}
 P(B\mid A)=\frac{P(A\cap B)}{P(A)},\qquad
 P(A\cap B)=P(A)P(B\mid A).
\end{equation}
已知骰子点数为偶数，再问点数是否为6，概率变为 $(1/6)/(1/2)=1/3$，不是原先的 $1/6$。
独立的含义是知道 $A$ 不改变 $B$ 的概率，即 $P(B\mid A)=P(B)$，从而得到乘法公式。
“互斥”与“独立”不同：两个正概率的互斥事件不可能独立，因为一个发生时另一个的条件概率是零。
独立性对双方是对称的，也传给补事件。例如
$P(A\cap B^c)=P(A)-P(A\cap B)=P(A)[1-P(B)]$，
同理 $P(A^c\cap B)=P(A^c)P(B)$、$P(A^c\cap B^c)=P(A^c)P(B^c)$。
随机变量是样本空间上的实值函数 $X(\omega)$；同一个数值可对应多个基本事件，其概率须将这些事件合并计算。

对于连续变量，$f_X(x)\dd x$ 是落在微小区间内的概率，而 $f_X(x)$ 本身不是概率；它的单位是 $X$ 的单位的倒数。
\[P(a<X<b)=\int_a^bf_X(x)\dd x,\qquad
f_X(x)\ge0,\qquad\int f_X(x)\dd x=1.\]
密度可以大于1，但任一区间的概率不能大于1。两个变量的联合密度 $f_{XY}$ 与边缘密度的关系为
\[f_X(x)=\int f_{XY}(x,y)\dd y,\qquad
f_Y(y)=\int f_{XY}(x,y)\dd x.\]
独立时联合密度才可分解为 $f_X(x)f_Y(y)$。离散情形对应 $p_i=\sum_jp_{ij}$、$q_j=\sum_ip_{ij}$。

\subsection{变量变换为什么要累加不同原像}

设无量纲速度分量 $u$ 服从标准高斯密度
$f_u(u)=(2\pi)^{-1/2}\exp(-u^2/2)$，令无量纲动能 $e=u^2/2$。
任意 $e>0$ 都有两个原像 $u_+=\sqrt{2e}$、$u_-=-\sqrt{2e}$。
两个速度区间都贡献同一能量区间的概率，因此
\begin{align*}
 f_e(e)&=f_u(u_+)\left|\frac{\dd u_+}{\dd e}\right|
       +f_u(u_-)\left|\frac{\dd u_-}{\dd e}\right|\\
 &=2\frac{e^{-e}}{\sqrt{2\pi}}\frac1{\sqrt{2e}}
 =\frac{e^{-e}}{\sqrt{\pi e}},\qquad e>0.
\end{align*}
代换 $e=s^2$ 检验归一化：
\[\int_0^\infty f_e(e)\dd e
=\frac2{\sqrt\pi}\int_0^\infty e^{-s^2}\dd s=1.\]
密度在零点附近发散但可积，并不表示零能量具有有限点概率。若只取正根，积分会错成 $1/2$。
一般变换公式在各单调分支上使用，导数为零的端点通过极限处理。

\subsection{样本平均与宏观稳定性}

设独立同分布变量的均值为 $\mu$、有限方差为 $\sigma^2$，样本平均为 $\bar X=N^{-1}\sum_iX_i$。
线性性给出 $\langle\bar X\rangle=\mu$；独立性使交叉协方差为零，于是
\[\operatorname{Var}(\bar X)=\frac1{N^2}\sum_i\operatorname{Var}(X_i)
=\frac{\sigma^2}{N}.\]
对于任何固定容许误差 $\delta>0$，切比雪夫不等式给出
\[P(|\bar X-\mu|\ge\delta)\le\frac{\sigma^2}{N\delta^2}\longrightarrow0.\]
这证明了有限方差情形的弱大数定律：样本平均依概率趋向总体平均，而不是每一次有限试验必然完全相等。
期望或方差不存在时不能直接套用本推导。例如柯西密度 $1/[\pi(1+x^2)]$ 的通常期望积分不收敛，
不能因图形关于零对称就把其期望当作零。若变量相关，方差中还须保留 $2\sum_{i<j}\operatorname{Cov}(X_i,X_j)$。


\subsection{由等分容器过渡到任意体积分区}
“左、右各一半”只是几何概率的一种特殊情形。粒子处在某区域的概率，先由允许位置的测度决定；再用独立试验规则计算粒子数的分布。固定总数并不表示各区域粒子数固定。

\Needspace{180pt}
\begin{origexamplebox}{体积比为四比一的两区域粒子数}\label{ch18:ex:teaching-asymmetric-binomial}
均匀平衡的稀薄理想气体含 $N$ 个粒子。把容器在几何上分成体积为 $4V/5$ 和 $V/5$ 的两区域，并无阻止粒子穿过的实体隔板。忽略粒子位置间的关联，求较小区域粒子数 $n$ 的分布、平均值和标准差，并解释宏观密度为何仍近似稳定。
\end{origexamplebox}
\Needspace{65pt}
\begin{xesolution}
每个粒子在小区域的概率为 $q=1/5$。选定哪 $n$ 个粒子进入小区域有 $\binom Nn$ 种组合，每种组合的概率为 $q^n(1-q)^{N-n}$，所以
\[
 P(n)=\binom Nn(1/5)^n(4/5)^{N-n},\qquad 0\le n\le N.
\]
二项式定理保证全部概率之和为1。令 $X_i$ 表示第 $i$ 个粒子是否在小区域，$X_i=1$ 的概率为 $q$。由于 $X_i^2=X_i$，有 $\langle X_i\rangle=q$、$\operatorname{Var}(X_i)=q(1-q)$。独立性消去协方差，故
\[
 \langle n\rangle=N/5,\qquad \operatorname{Var}(n)=4N/25,
 \qquad \sigma_n=2\sqrt N/5,\qquad
 \frac{\sigma_n}{\langle n\rangle}=\frac2{\sqrt N}.
\]
粒子数的绝对起伏随 $N$ 增大，却比平均数增长得慢，所以相对密度起伏变小。大区域粒子数为 $N-n$，与小区域完全反相关；不能把两区域总粒子数的方差写成两个方差直接相加，否则会误判固定的 $N$ 也在涨落。这里的临时粒子编号只用于稀薄经典模型的位置计数，不是为全同粒子添加新的可观测身份。
\end{xesolution}

\section{微观状态、宏观状态与熵}

热力学只用 $p,V,T,N$ 等少数宏观量描述系统；统计物理则把宏观规律解释为大量微观自由度的概率规律。
一个宏观状态通常对应极多微观状态。若孤立系统在给定能量壳内的可及微观状态等概率，某宏观状态出现的概率
正比于它包含的微观状态数 $\Omega$。玻尔兹曼熵（Boltzmann entropy）定义为
\begin{equation}
 S=k_B\ln\Omega.
\end{equation}
两个系统各自的能量等宏观约束已分别给定、允许状态可独立组合时，微观状态数相乘，$\Omega_{12}=\Omega_1\Omega_2$，所以熵相加；对数正是把乘法计数变成
可加宏观量的函数。
仅知道总能量时，还须对能量在两系统间的分配求和，不能不加条件地直接相乘。

当各微观状态的概率并不相等时，更一般的吉布斯熵（Gibbs entropy）为
\begin{equation}
 S_G=-k_B\sum_i p_i\ln p_i.
\end{equation}
若能量壳内的 $\Omega$ 个状态等概率，$p_i=1/\Omega$，上式立即退化为 $S=k_B\ln\Omega$。
因此“状态数熵”和“概率分布的信息熵”并非两套无关定义，而是同一结构在等概率分布下的特例。
\subsection{信息量、条件熵与相关性}

一个事件越罕见，观察到它所获得的信息越多。要求独立事件的信息可加、确定事件的信息为零，
即 $I(pq)=I(p)+I(q)$、$I(1)=0$；再要求连续及随概率递减，得到 $I(p)=-c\ln p$，其中 $c>0$。
例如令 $p=e^{-x}$，则 $J(x)=I(e^{-x})$ 满足连续加法方程 $J(x+y)=J(x)+J(y)$，故 $J(x)=cx$。
取 $c=1$，平均信息量是
\[H(X)=\sum_i p_i I(p_i)=-\sum_i p_i\ln p_i.\]
自然对数给出的信息单位为奈特；若采用以2为底的对数，单位为比特。
公平骰子的熵是 $\ln6$；已知为偶数后的熵是 $\ln3$，所获得的信息是 $\ln2$。

把一个结果进一步细分为 $(i,j)$，写 $p_{ij}=p_i p(j\mid i)$，代入定义：
\begin{align*}
 H(X,Y)&=-\sum_{ij}p_{ij}\ln[p_i p(j\mid i)]\\
 &= -\sum_i p_i\ln p_i
    -\sum_i p_i\sum_jp(j\mid i)\ln p(j\mid i)\\
 &=H(X)+H(Y\mid X).
\end{align*}
这里 $H(Y\mid X)$ 是先知道 $X$ 后剩余不确定性的加权平均。独立时 $p(j\mid i)=q_j$，
因而 $H(X,Y)=H(X)+H(Y)$。

一般情况下不能直接相加。令 $q_{ij}=p_iq_j$，利用 $-\ln x\ge1-x$，有
\[\sum_{p_{ij}>0}p_{ij}\ln\frac{p_{ij}}{q_{ij}}
\ge\sum_{p_{ij}>0}(p_{ij}-q_{ij})
=1-\sum_{p_{ij}>0}q_{ij}\ge0.\]
零概率项以连续极限解释；若 $p_{ij}=0$ 的位置上 $q_{ij}>0$，上述不等式仍成立。
左边等于 $H(X)+H(Y)-H(X,Y)$，即互信息。所以联合熵不大于两边之和，等号当且仅当独立。
相关性会减少额外信息，但协方差为零并不足以证明独立；不能把两者混淆。
例如用独立、等方差的零均值高斯变量 $X,Y$ 模拟随机光场的两个正交分振幅，
与 $x$ 轴夹角 $\theta$ 方向的分振幅为 $Z=X\cos\theta+Y\sin\theta$。
它与 $X$ 的相关系数为 $\cos\theta$：$\theta=\pi/2$ 时独立；锐角时有关联。
对于这个高斯模型，$I(X{:}Z)=-\tfrac12\ln(1-\cos^2\theta)$，可由二维高斯熵与两个边缘熵相减得到。
同方向的连续、无测量噪声变量完全相同，互信息发散；若按有限分辨率离散读数，同方向的互信息才等于该离散变量的熵。
这说明讨论光场相关性时应交代分布模型和测量分辨率，不能把连续与离散结论直接混用。

\subsection{连续分布与最大熵的条件}

连续变量可形式地写微分熵 $h=-\int f\ln f\,\dd x$，但量纲与坐标选择必须明确。
可以先选定参考尺度把变量无量纲化，或用参考测度写相对熵；不能对带单位的密度直接取对数后比较不同单位下的数值。
对固定的无量纲坐标，给定均值 $\mu$ 和方差 $\sigma^2>0$，令
\[g(x)=\frac1{\sqrt{2\pi}\sigma}\exp\!\left[-\frac{(x-\mu)^2}{2\sigma^2}\right].\]
任意具有相同均值、方差的密度 $f$ 满足 $D(f\Vert g)=\int f\ln(f/g)\dd x\ge0$。
因 $\ln g$ 只含常数和 $(x-\mu)^2$，由约束可算出
\[D(f\Vert g)=-h(f)+\frac12\ln(2\pi\sigma^2)+\frac12\ge0.\]
所以 $h(f)\le\tfrac12\ln(2\pi e\sigma^2)$，且仅在 $f=g$ 时取等号。
高斯分布是这些约束下的最大熵分布，不是任何约束下都必然出现的分布。

均匀角分布也须考虑测度：球面上均匀取方向时，面积元为
$\dd A=r^2\sin\theta\,\dd\theta\dd\phi$，故极角密度为 $f_\theta(\theta)=\tfrac12\sin\theta$，
不是在 $0\le\theta\le\pi$ 上均匀。常数由 $\int_0^\pi\sin\theta\dd\theta=2$ 确定。
这是几何概率、变量变换和统计物理中速度空间球壳因子的共同思想。


\Needspace{22\baselineskip}
\begin{origexamplebox}{二项分布与最概然宏观态}
$N$ 个彼此独立、可区分的粒子，每个粒子以相同概率处在容器左半或右半。左侧有 $n$ 个粒子的微观状态数是多少？哪个 $n$ 最可能？
\end{origexamplebox}
\begin{xesolution}
选择哪 $n$ 个粒子位于左侧，有
\[
 \Omega(n)=\binom Nn,
 \qquad P(n)=2^{-N}\binom Nn.
\]
相邻两项比值
\[
 \frac{P(n+1)}{P(n)}=\frac{N-n}{n+1}
\]
在 $n<(N-1)/2$ 时大于1，大于该界限时小于1，等于时比值为1。
因此偶数 $N$ 只有峰值 $n=N/2$；奇数 $N$ 有两个等高峰 $n=(N-1)/2,(N+1)/2$。
令每个粒子处于左侧的指示变量为 $X_i=0$ 或1，则 $n=\sum_iX_i$，
$\langle X_i\rangle=1/2$、$\operatorname{Var}(X_i)=1/4$。独立性给出
\[\langle n\rangle=N/2,\qquad \operatorname{Var}(n)=N/4,\qquad \sigma_n=\sqrt N/2.\]
这是任意 $N$ 的精确结果。以总粒子数为分母时 $\sigma_n/N=1/(2\sqrt N)$；
以左侧平均数为分母时 $\sigma_n/\langle n\rangle=1/\sqrt N$。两者都随粒子数增大而消失。
\end{xesolution}

\begin{center}\begin{minipage}{\linewidth}\centering
\begin{tikzpicture}
\begin{axis}[width=.74\linewidth,height=4.7cm,ybar,bar width=5.5pt,
  axis lines=left,xlabel={$n$},ylabel={$P(n)$},xmin=-.5,xmax=20.5,ymin=0,ymax=.19,
  xtick={0,5,10,15,20},ytick={0,.05,.10,.15},
  every axis label/.style={fill=none},xlabel style={yshift=-5pt}]
  \addplot[fill=DeepBlue!35,draw=DeepBlue] coordinates{
   (0,.000001)(1,.000019)(2,.000181)(3,.001087)(4,.004621)(5,.014786)
   (6,.036964)(7,.073929)(8,.120134)(9,.160179)(10,.176197)
   (11,.160179)(12,.120134)(13,.073929)(14,.036964)(15,.014786)
   (16,.004621)(17,.001087)(18,.000181)(19,.000019)(20,.000001)};
\end{axis}
\end{tikzpicture}
\figcap{$N=20$ 时二项分布集中在 $n=N/2$ 附近；粒子数增大后绝对宽度增长，而相对宽度缩小。}\label{ch18:fig:binomial}
\end{minipage}\end{center}
图~\ref{ch18:fig:binomial} 把例题中的最概然值和 $\sqrt N$ 量级涨落直接画了出来。

\begingroup\tcbset{breakable/.default=true}\begin{starread}{相空间、系综与刘维尔定理}
具有 $f$ 个经典自由度的微观状态可表示为相空间（phase space）中的一点
\[
 \Gamma=(q_1,\ldots,q_f;p_1,\ldots,p_f).
\]
系综（ensemble）不是许多粒子，而是许多份具有相同宏观条件、微观状态各异的系统副本；其状态由相空间概率密度
$\rho_\Gamma(\Gamma,t)$ 描述。哈密顿运动保持相空间体积；若 $H$ 为哈密顿量（Hamiltonian），概率密度满足刘维尔方程（Liouville equation）
\[
 \frac{\partial\rho_\Gamma}{\partial t}+\{\rho_\Gamma,H\}=0,
\]
其中哈密顿量 $H(q,p)$ 是用正则坐标和动量表示的系统能量，花括号表示泊松括号（\mbox{Poisson bracket}）：
\[\{A,H\}=\sum_{j=1}^{f}\left(\frac{\partial A}{\partial q_j}\frac{\partial H}{\partial p_j}
-\frac{\partial A}{\partial p_j}\frac{\partial H}{\partial q_j}\right).\]
由哈密顿方程 $\dot q_j=\partial H/\partial p_j$、$\dot p_j=-\partial H/\partial q_j$，
相空间流速的散度为 $\sum_j(H_{p_jq_j}-H_{q_jp_j})=0$；将其代入概率连续性方程即得上式。
这里要求哈密顿量足够光滑。体积守恒并不自动证明遍历性，也不自动导出粗粒化熵增。
孤立平衡系统在给定能量壳上的等概率原理，是微正则系综的起点。它说的是相空间体积元等概率，而不是每个位置坐标等概率；
例如单摆在转折点附近停留更久，位置分布并不均匀，但相轨道上的正则相空间测度仍与哈密顿演化相容。

若联合分布为 $p_{ij}$、边缘分布为 $p_i,q_j$，互信息（mutual information）
\[
 I(X{:}Y)=\sum_{ij}p_{ij}\ln\frac{p_{ij}}{p_iq_j}\ge0
\]
衡量知道一个变量后对另一个变量减少了多少不确定性。独立时 $p_{ij}=p_iq_j$，互信息为零，联合信息熵等于两边之和。
\end{starread}\endgroup

\section{正则分布与玻尔兹曼因子}

把小系统与温度为 $T$ 的巨大热库接触，总能量固定。小系统处在能量 $E_i$ 的微观状态时，热库剩余能量为
$E_{\rm tot}-E_i$。利用 $\partial S_R/\partial E_R=1/T$ 展开热库状态数，
\[
 \ln\Omega_R(E_{\rm tot}-E_i)
 \simeq\ln\Omega_R(E_{\rm tot})-\frac{E_i}{k_BT},
\]
便得到概率
\begin{equation}
 P_i=\frac{e^{-\beta E_i}}{Z},\qquad
 \beta=\frac1{k_BT},\qquad
 Z=\sum_i e^{-\beta E_i}.
\end{equation}
$Z$ 是配分函数（partition function）。任意物理量的平均值为 $\langle A\rangle=\sum_iP_iA_i$，平均能量可写成
\begin{equation}
 \langle E\rangle=-\frac{\partial\ln Z}{\partial\beta}.
\end{equation}
玻尔兹曼因子比较的是能量差：
\begin{equation}
 \frac{P_2}{P_1}=e^{-(E_2-E_1)/(k_BT)}.
\end{equation}
温度升高时高能态相对概率增加，但并不意味着所有微观状态概率完全相等；只有 $k_BT$ 远大于相关能级差时才近似如此。

\subsection{从归一化到平均能量}

玻尔兹曼因子 $e^{-\beta E_i}$ 只给出相对权重。若所有权重都乘同一常数，物理概率不应改变；
除以权重之和 $Z$ 才使 $\sum_iP_i=1$。这里下标 $i$ 表示\emph{一个微观状态}。
若能量为 $E_j$ 的能级对应 $g_j$ 个不同状态，落在整个能级上的概率应为
\[
 P(E_j)=\frac{g_je^{-\beta E_j}}{\sum_kg_ke^{-\beta E_k}}.
\]
例如基态只有1个状态、能量高出 $\varepsilon$ 的激发能级有3个状态，
两能级的总布居比为 $3e^{-\beta\varepsilon}$，不是 $e^{-\beta\varepsilon}$。
高温极限下每个微观状态近似等概率，但两个能级的总布居比趋于3，说明“能级等概率”与“状态等概率”不同。

在体积等外参量固定、各 $E_i$ 不随 $\beta$ 变化时，把配分函数逐项求导：
\[
 \frac{\partial Z}{\partial\beta}=-\sum_iE_ie^{-\beta E_i},\qquad
 -\frac1Z\frac{\partial Z}{\partial\beta}
 =\sum_iE_i\frac{e^{-\beta E_i}}Z=\langle E\rangle.
\]
这就解释了为什么对 $\ln Z$ 求导可以得到平均能量。再求一次导数有
\[
 \frac{\partial^2\ln Z}{\partial\beta^2}
 =\frac{\sum_iE_i^2e^{-\beta E_i}}Z
 -\left(\frac{\sum_iE_ie^{-\beta E_i}}Z\right)^2
 =\langle E^2\rangle-\langle E\rangle^2.
\]
利用 $\dd\beta/\dd T=-1/(k_BT^2)$，便得到
\[
 C_V=\left(\frac{\partial\langle E\rangle}{\partial T}\right)_V
 =\frac{\operatorname{Var}(E)}{k_BT^2}\ge0.
\]
因此热容与能量涨落的联系来自同一个归一化分布，而不是另加的经验公式。
这里的非负性针对这些条件下的正则平衡系统，不能不加条件推广到所有长程相互作用的孤立系统。

\begin{origexamplebox}{二能级系统的平均能量与热容}
$N$ 个独立、可区分或局域化的粒子各有能量 $0$ 和 $\varepsilon>0$ 两个非简并能级，与温度 $T>0$ 的热库平衡。求平均总能量和热容，并判断低温和高温极限。
\end{origexamplebox}
\begin{xesolution}
单粒子配分函数 $z=1+e^{-\beta\varepsilon}$，处于激发态的概率为
\[
 p_e=\frac{e^{-\beta\varepsilon}}{1+e^{-\beta\varepsilon}}
 =\frac1{e^{\beta\varepsilon}+1}.
\]
所以
\[
 \boxed{U=\frac{N\varepsilon}{e^{\varepsilon/(k_BT)}+1}}.
\]
$T\to0$ 时 $U\to0$；$T\to\infty$ 时两个能级等概率，$U\to N\varepsilon/2$。
令 $x=\varepsilon/(k_BT)$，则 $\dd x/\dd T=-x/T$，固定 $N,\varepsilon$ 时
\[ C=\frac{\dd U}{\dd T}
=-\frac{N\varepsilon e^x}{(e^x+1)^2}\left(-\frac{x}{T}\right).\]
利用 $\varepsilon x/T=k_Bx^2$，热容为
\[
 C=Nk_B\left(\frac{\varepsilon}{k_BT}\right)^2
 \frac{e^{\varepsilon/(k_BT)}}{[1+e^{\varepsilon/(k_BT)}]^2}
\]
在高低温两端都趋于零，中间出现峰值。高温能量饱和说明有限能级系统不能无限吸热。
\end{xesolution}

\begin{center}\begin{minipage}{\linewidth}\centering
\begin{tikzpicture}[x=1.25cm,y=1.2cm]
  \draw[axis] (-.25,0)--(4.7,0) node[right]{$E$};
  \foreach \x in {0,1,2,3,4}{
    \pgfmathsetmacro{\h}{2.5*(1-.6)*pow(.6,\x)}
    \draw[DeepBlue,very thick] (\x,0)--(\x,\h);
    \fill[DeepBlue] (\x,\h) circle (1.2pt);
  }
  \node[clearlabel,above] at (0,1) {$P_0=0.4$};
  \node[clearlabel,below=5pt] at (2,0) {$E_i=i\varepsilon$};
  \node[clearlabel] at (2.7,1.05) {$q=e^{-\beta\varepsilon}=0.6$};
  \node at (4.5,.25) {$\cdots$};
\end{tikzpicture}
\figcap{等间隔能级的归一化概率 $P_n=(1-q)q^n$，图取 $q=0.6$，省略号表示后续能级。}\label{ch18:fig:boltzmann-population}
\end{minipage}\end{center}
图~\ref{ch18:fig:boltzmann-population} 的每根竖线对应一个离散能级，图中画的是归一化概率
$P_n=q^n/Z=(1-q)q^n$，所有高度之和为1。下面求配分函数时，求和的则是未归一化权重 $q^n$，
不能把图中的概率再相加当成 $Z$。

\begin{origexamplebox}{等间隔能级的配分函数}
一个粒子具有能级 $E_n=n\varepsilon$（$n=0,1,2,\ldots$），各能级非简并。求单粒子配分函数和平均能量，并检验高、低温极限。
\end{origexamplebox}
\begin{xesolution}
令 $q=e^{-\beta\varepsilon}$。因 $T>0,\varepsilon>0$，有 $0<q<1$。配分函数是收敛的等比级数：
\[
 Z=\sum_{n=0}^\infty q^n=\frac1{1-q}.
\]
有限和为 $(1-q^{M+1})/(1-q)$，取 $M\to\infty$ 即得上式。
又因 $\partial q/\partial\beta=-\varepsilon q$，有
$\partial\ln Z/\partial\beta=-\varepsilon q/(1-q)$。
于是
\[
 U=-\frac{\partial\ln Z}{\partial\beta}
 =\frac{\varepsilon}{e^{\beta\varepsilon}-1}.
\]
即
\[
 \boxed{U=\frac{\varepsilon}{e^{\varepsilon/(k_BT)}-1}}.
\]
低温 $k_BT\ll\varepsilon$ 时 $U\sim\varepsilon e^{-\varepsilon/(k_BT)}\to0$；高温时用
$e^x-1\sim x$ 得 $U/(k_BT)\to1$，即 $U\sim k_BT$，与一维经典谐振子两个二次项的均分结果一致。
更精细地，$U=k_BT-\varepsilon/2+O(\varepsilon^2/(k_BT))$，故不能把渐近等价理解为两者差值趋于零。
$Z$ 无量纲，$U$ 的量纲由 $\varepsilon$ 给出。
\end{xesolution}

\Needspace{180pt}
\begin{origexamplebox}{简并度能否抵消较高能量的不利因素}\label{ch18:ex:teaching-degeneracy}
一个系统只有两组能级：能量为0的基能级有1个微观状态，能量为 $\varepsilon>0$ 的激发能级有3个微观状态。系统与正温度热库平衡。求激发能级的总概率、两组总布居相等的温度及平均能量，并讨论高、低温极限。
\end{origexamplebox}
\Needspace{65pt}
\begin{xesolution}
每个激发微观状态的权重是 $e^{-\beta\varepsilon}$，三个不同状态的权重须相加：
\[
 Z=1+3e^{-\beta\varepsilon},\qquad
 P_e=\frac{3e^{-\beta\varepsilon}}{1+3e^{-\beta\varepsilon}},\qquad
 P_g=\frac1{1+3e^{-\beta\varepsilon}}.
\]
两组等布居要求 $3e^{-\varepsilon/(k_BT)}=1$，因此 $T_*=\varepsilon/(k_B\ln3)$。平均能量不是两个能级的算术平均，而是
\[
 U=0\cdot P_g+\varepsilon P_e=\varepsilon P_e.
\]
当 $T\to0^+$ 时，$P_e\to0$、$U\to0$；当 $T\to\infty$ 时，四个微观状态等概率，$P_e\to3/4$、$U\to3\varepsilon/4$，并非 $\varepsilon/2$。即使激发组总概率超过基组，任一单独激发态的概率仍低于唯一基态；这不是逐态布居反转，而是简并状态数带来的总数优势。
\end{xesolution}

\section{能量均分定理}

经典正则分布中，哈密顿量里的每个独立二次项平均贡献 $k_BT/2$。例如
$p_x^2/(2m)$ 是一个二次项，三维单原子理想气体每个分子平均平动能为
\begin{equation}
 \left\langle\frac{p_x^2+p_y^2+p_z^2}{2m}\right\rangle=\frac32k_BT,
\end{equation}
从而 $U=3Nk_BT/2$、$C_V=3Nk_B/2$。一个经典一维简谐振子含动能和势能两个二次项，平均能量为 $k_BT$。

均分定理要求经典连续近似有效并达到热平衡。若量子能级间隔远大于 $k_BT$，相应自由度会冻结；这正是双原子气体
振动热容在常温下小于经典预测的原因。

更一般地，对正则分布中满足边界项消失的相空间变量 $x$，分部积分给出
\begin{equation}
 \left\langle x\frac{\partial H}{\partial x}\right\rangle=k_BT.
\end{equation}
若 $H$ 中含一个二次项 $ax^2$，左边来自该项的贡献就是 $2\langle ax^2\rangle$，故
$\langle ax^2\rangle=k_BT/2$。这也说明均分结论依赖“二次型”和“经典连续积分”两个条件。

可先用一个二次项看清积分过程。取 $E_x=ax^2$、$a>0$，$x$ 可遍历整个实轴。
其归一化积分和平均能量分别为
\[
 I(\beta)=\int_{-\infty}^{\infty}e^{-\beta ax^2}\dd x
 =\sqrt{\frac{\pi}{\beta a}},\qquad
 \langle E_x\rangle=-\frac{\partial\ln I}{\partial\beta}
 =\frac1{2\beta}=\frac12k_BT.
\]
对一维谐振子，$E=p^2/(2m)+kx^2/2$，两个积分可以分离，故动能和势能各贡献一份
$k_BT/2$。不能只数“振动方向”就漏掉其中一项；也不能把任意一个角坐标都当成可在整个实轴上
独立积分的二次变量。分部积分的完整条件和非二次项的反例，在本章后面的连续相空间计算中继续展开。

\section{麦克斯韦速率分布}

平衡理想气体的速度三个分量相互独立，每个分量服从高斯分布（Gaussian distribution）。
令 $a=m/(2k_BT)$。一维动能为 $mv_x^2/2$，故玻尔兹曼因子给出
\[f_x(v_x)=\sqrt{a/\pi}\,e^{-av_x^2},\qquad
F(\boldsymbol v)=(a/\pi)^{3/2}e^{-a(v_x^2+v_y^2+v_z^2)}.\]
高斯积分 $\int_{-\infty}^{\infty}e^{-av_x^2}\dd v_x=\sqrt{\pi/a}$ 确定了归一化系数。
速率在 $v$ 到 $v+\dd v$ 内的速度向量组成半径 $v$、厚度 $\dd v$ 的球壳，体积为 $4\pi v^2\dd v$，
所以须将三维密度乘以球壳体积，而不是令 $v_x=v$ 就结束。
麦克斯韦速率分布（Maxwell speed distribution）为
\begin{equation}
 f(v)=4\pi\left(\frac{m}{2\pi k_BT}\right)^{3/2}
 v^2e^{-mv^2/(2k_BT)},\qquad v\ge0,
\end{equation}
并满足 $\int_0^\infty f(v)\dd v=1$。最概然速率、平均速率和方均根速率分别为
\begin{equation}
 v_{\rm mp}=\sqrt{\frac{2k_BT}{m}},\qquad
 \bar v=\sqrt{\frac{8k_BT}{\pi m}},\qquad
 v_{\rm rms}=\sqrt{\frac{3k_BT}{m}}.
\end{equation}
对 $v>0$，$\dd\ln f/\dd v=2/v-2av$，令其为零得 $v_{\rm mp}=1/\sqrt a$。
计算平均值时，代换 $u=av^2$ 给出 $\int_0^\infty v^3e^{-av^2}\dd v=1/(2a^2)$；
而高斯积分对 $a$ 求导可得 $\int_0^\infty v^4e^{-av^2}\dd v=3\sqrt\pi/(8a^{5/2})$。
分别乘归一化系数，得到 $\bar v=2/\sqrt{\pi a}$ 和 $\langle v^2\rangle=3/(2a)$。

\begin{center}\begin{minipage}{\linewidth}\centering
\begin{tikzpicture}
\begin{axis}[width=.74\linewidth,height=4.8cm,axis lines=left,xlabel={$u=v/\sqrt{2k_BT_1/m}$},ylabel={$f_u(u)$},
  xmin=0,xmax=4,ymin=0,ymax=1.05,ticks=none,every axis label/.style={fill=none}]
  \addplot[DeepBlue,very thick,domain=0:4,samples=150]{4/sqrt(pi)*x^2*exp(-x^2)};
  \addplot[TextGray,dashed,domain=0:4,samples=150]{4/sqrt(pi)*pow(.5,1.5)*x^2*exp(-.5*x^2)};
  \node at (axis cs:1.0,.96) {$T_1$};
  \node at (axis cs:2.8,.8) {$T_2=2T_1$};
  \draw[->,TextGray] (axis cs:2.8,.71)--(axis cs:2.45,.24);
\end{axis}
\end{tikzpicture}
\figcap{以同一速率单位比较两个温度；两条概率密度在完整区间 $0\le u<\infty$ 上的积分均为1。}\label{ch18:fig:speed}
\end{minipage}\end{center}
图~\ref{ch18:fig:speed} 中高温曲线的峰向右移且降低，横轴使用共同单位，不能分别重标横轴后比较峰的位置。

\begin{origexamplebox}{氮分子的方均根速率}
把氮气视为理想气体，摩尔质量 $M=28.0\,\mathrm{g/mol}$。求 $T=300\,\mathrm K$ 时的方均根速率。
\end{origexamplebox}
\begin{xesolution}
用单分子质量公式也可计算，但摩尔形式更直接：
\[
 v_{\rm rms}=\sqrt{\frac{3RT}{M}}
 =\sqrt{\frac{3\times8.314\times300}{28.0\times10^{-3}}}
 \approx5.17\times10^2\,\mathrm{m/s}.
\]
所以 $\boxed{v_{\rm rms}\approx517\,\mathrm{m/s}}$。必须把摩尔质量换成 $\mathrm{kg/mol}$；否则数值会错约 $\sqrt{1000}$ 倍。
\end{xesolution}

\begin{origexamplebox}{麦克斯韦分布中的越阈概率}
质量为 $m$ 的理想气体分子处于温度 $T$。写出速率大于给定阈值 $v_0$ 的概率积分，并在
$v_0\gg\sqrt{k_BT/m}$ 时判断其主要数量级。
\end{origexamplebox}
\begin{xesolution}
由归一化速率分布，所求概率为
\[
 P(v>v_0)=4\pi\left(\frac{m}{2\pi k_BT}\right)^{3/2}
 \int_{v_0}^{\infty}v^2e^{-mv^2/(2k_BT)}\dd v.
\]
令 $a=m/(2k_BT)$，分部积分一次：
\[
 \int_{v_0}^{\infty}v^2e^{-av^2}\dd v
 =\frac{v_0e^{-av_0^2}}{2a}
 +\frac1{2a}\int_{v_0}^{\infty}e^{-av^2}\dd v.
\]
令 $x_0=\sqrt a\,v_0$，定义余误差函数
\[\operatorname{erfc}(x_0)=\frac2{\sqrt\pi}\int_{x_0}^{\infty}e^{-x^2}\dd x.\]
将归一化系数乘回上式，得到
\[
 \boxed{P(v>v_0)=\operatorname{erfc}(x_0)+\frac{2x_0}{\sqrt\pi}e^{-x_0^2}}.
\]
再次分部积分可得 $\operatorname{erfc}(x_0)\sim e^{-x_0^2}/(\sqrt\pi x_0)$，因此
\[P(v>v_0)=\frac{2x_0}{\sqrt\pi}e^{-x_0^2}
\left[1+\frac1{2x_0^2}+O(x_0^{-4})\right]\quad(x_0\gg1).\]
各项无量纲。$v_0=0$ 时精确式为1，$v_0\to\infty$ 时趋于0；大阈值近似不能用于 $v_0=0$。
\end{xesolution}

\section{熵增的统计含义}

孤立宏观系统从低概率宏观态走向高概率宏观态，是因为后者包含压倒性更多的微观状态。熵增不是微观力学定律失去时间反演性，
而是典型初态与粗粒化描述下的概率结论。对有限系统，短时熵减涨落并非逻辑上不可能，但概率随粒子数呈指数压低。

\subsection{理想气体混合熵}

两种可区分理想气体在相同温度、压强下分别占有体积 $V_1,V_2$，撤去隔板后各自自由膨胀到
$V=V_1+V_2$。因为理想气体内能只依赖温度，等温混合的熵变为
\begin{equation}
 \Delta S_{\rm mix}
 =n_1R\ln\frac{V}{V_1}+n_2R\ln\frac{V}{V_2}
 =-R\sum_i n_i\ln x_i,
\end{equation}
其中 $x_i=n_i/(n_1+n_2)$ 是摩尔分数。若两侧其实是完全相同的气体，交换粒子不产生新的可分辨微观状态，
不能沿用“可区分组分”的计数；此时去掉隔板的熵变为零。

\begin{center}\begin{minipage}{\linewidth}\centering
\begin{tikzpicture}[scale=.92]
  \draw[ground] (0,0) rectangle (6,2.7);
  \draw[very thick] (3,0)--(3,2.7);
  \foreach \p in {(0.5,.5),(.9,1.2),(1.6,.7),(.55,2.1),(1.45,1.8)}
    \fill[DeepBlue] \p circle (2pt);
  \foreach \p in {(3.4,.55),(3.6,1.1),(4.7,.65),(5.4,1.7),(4.4,2.25)}
    \draw[TextGray,thick] \p circle (2pt);
  \node[clearlabel,above,align=center,font=\footnotesize] at (1.5,2.7) {组分1\\实心点};
  \node[clearlabel,above,align=center,font=\footnotesize] at (4.5,2.7) {组分2\\空心点};
  \draw[vector] (6.45,1.35)--(7.6,1.35);
  \begin{scope}[xshift=8.2cm]
    \draw[ground] (0,0) rectangle (6,2.7);
    \foreach \p in {(0.5,.5),(.9,1.9),(2.6,2.1),(3.5,.55),(5.5,.8)}
      \fill[DeepBlue] \p circle (2pt);
    \foreach \p in {(.55,1.3),(2.1,1.45),(3.1,.95),(4.9,.5),(5.3,2.05)}
      \draw[TextGray,thick] \p circle (2pt);
  \end{scope}
\end{tikzpicture}
\figcap{可区分理想气体混合。隔板撤去后每种组分可占据的体积增大，因而产生混合熵。}\label{ch18:fig:mixing}
\end{minipage}\end{center}
图~\ref{ch18:fig:mixing} 用两种点型区分组分；若点完全不可区分，不能再把左右来源当成新的微观标记。

\begin{origexamplebox}{两种等量理想气体的混合熵}
各有 $1.00\,\mathrm{mol}$ 的两种不同理想气体，在同温同压下由隔板分开，各占体积 $V$。撤去隔板并等温混合，求熵变；若两侧是同一种气体，结果怎样？
\end{origexamplebox}
\begin{xesolution}
两种气体各自从 $V$ 自由扩展到 $2V$，所以
\[
 \Delta S=R\ln2+R\ln2=\boxed{2R\ln2\approx11.5\,\mathrm{J/K}}.
\]
熵变为正，单位是 $\mathrm{J/K}$。若两侧是温度、压强和成分完全相同的同一种气体，撤板前后宏观状态相同；同种粒子的置换不构成新的可分辨状态，故
\[
 \boxed{\Delta S=0}.
\]
把同种气体也人为贴上“左、右来源标签”会造成虚假的吉布斯佯谬。
\end{xesolution}

\subsection{能量涨落与热容}

固定 $N,V$ 且能级 $E_i$ 不显含温度时，$Z'= -\sum_iE_ie^{-\beta E_i}$、
$Z''=\sum_iE_i^2e^{-\beta E_i}$，所以
\[\frac{\partial^2\ln Z}{\partial\beta^2}
=\frac{Z''}Z-\left(\frac{Z'}Z\right)^2=\langle E^2\rangle-\langle E\rangle^2.\]
又因 $U=-\partial\ln Z/\partial\beta$、$\partial\beta/\partial T=-1/(k_BT^2)$，
\[ C_V=\frac{\partial U}{\partial T}
=\frac{\operatorname{Var}(E)}{k_BT^2}.\]
整理得到
\begin{equation}
 \operatorname{Var}(E)=\langle E^2\rangle-\langle E\rangle^2
 =k_BT^2C_V.
\end{equation}
对普通宏观系统，$E\propto N$、$C_V\propto N$，所以 $\sigma_E/E\propto N^{-1/2}$。
热容非负对应正则平衡态能量方差非负；若某些长程相互作用孤立系统出现负热容，不能直接把这条正则系综结论照搬过去。

\begin{center}\begin{minipage}{\linewidth}\centering
\begin{tikzpicture}
\begin{axis}[width=.72\linewidth,height=4.5cm,axis lines=left,
  xlabel={$x=(E-U)/\sigma_E$},ylabel={$f_x(x)$},xmin=-4,xmax=4,ymin=0,ymax=.55,ticks=none,
  every axis label/.style={fill=none}]
  \addplot[DeepBlue,very thick,domain=-4:4,samples=180]{exp(-x^2/2)/sqrt(2*pi)};
  \draw[guide] (axis cs:0,0)--(axis cs:0,.399);
  \node[clearlabel] at (axis cs:0,.49) {$E=U$};
  \draw[guide] (axis cs:-1,.02)--(axis cs:-1,.30) (axis cs:1,.02)--(axis cs:1,.30);
  \draw[<->,TextGray] (axis cs:-1,.08)--(axis cs:1,.08);
  \node[clearlabel] at (axis cs:2.5,.16) {$\Delta E=2\sigma_E$};
  \draw[TextGray,->] (axis cs:2,.13)--(axis cs:1,.08);
\end{axis}
\end{tikzpicture}
\figcap{平均能量附近的高斯近似，横轴以 $\sigma_E$ 为单位；$\sigma_E^2=k_BT^2C_V$。}\label{ch18:fig:energy-fluctuation}
\end{minipage}\end{center}
图~\ref{ch18:fig:energy-fluctuation} 标出的 $2\sigma_E$ 是典型宽度，下面的例题将其与平均能量比较。

\Needspace{15\baselineskip}
\begin{origexamplebox}{能量相对涨落的粒子数标度}
经典单原子理想气体含 $N$ 个粒子，处于温度 $T$ 的正则系综。求能量标准差与平均能量之比。
\end{origexamplebox}
\begin{xesolution}
理想气体
\[
 U=\frac32Nk_BT,
 \qquad C_V=\frac32Nk_B.
\]
由涨落公式
\[
 \sigma_E=\sqrt{k_BT^2C_V}=k_BT\sqrt{\frac{3N}{2}}.
\]
因此
\[
 \boxed{\frac{\sigma_E}{U}=\sqrt{\frac{2}{3N}}}.
\]
该比值无量纲并随 $N^{-1/2}$ 消失；$N\to\infty$ 时热力学能量成为相对稳定的宏观量。对少粒子系统，涨落不可忽略。
\end{xesolution}

\section{分子通量、压强与统计假设}

\subsection{穿过小孔的分子并不是无偏样本}

设数密度为 $n_N=N/V$，取孔面法线为 $z$ 轴。速度 $v_z>0$ 的分子在时间 $\dd t$ 内，
只有位于厚度 $v_z\dd t$ 内的部分能够到达面积 $\dd A$。因此单位面积、单位时间的穿越数为
\[J=n_N\int_{v_z>0}v_z F(\boldsymbol v)\,\dd^3v
=n_N\sqrt{\frac{k_BT}{2\pi m}}=\frac14 n_N\bar v.\]
第一步中的 $v_z$ 权重说明较快且较接近法向运动的分子更容易穿孔；不能直接用容器内的速率密度代表逸出束的速率密度。
几何上，对于各向同性分布，前半球上的平均通量因子是
\[\frac1{4\pi}\int_0^{2\pi}\dd\phi\int_0^{\pi/2}\cos\theta\sin\theta\dd\theta=\frac14.\]
小孔逸出公式适用于孔尺度小于平均自由程、孔外回流可忽略的分子流模型。

分子在壁面弹性反射时，每次传递法向动量 $2mv_z$。将其乘碰撞通量并积分，得到
\[p=2mn_N\int_{v_z>0}v_z^2F\,\dd^3v
=mn_N\langle v_z^2\rangle
=\frac13mn_N\langle v^2\rangle=\frac23n_N\langle\varepsilon_k\rangle.\]
结合 $\langle\varepsilon_k\rangle=3k_BT/2$，就恢复 $p=n_Nk_BT$。
对各向同性光子气，速率固定为 $c$、单光子动量大小为 $\varepsilon/c$，同样的角积分给出
数通量 $J_\gamma=n_\gamma c/4$、能量通量 $j_E=uc/4$ 和压强 $p=u/3$，其中 $u$ 为能量密度。
平行光束不满足各向同性条件，不能直接套用 $p=u/3$。

\subsection{把穿孔权重用于能量计算}
分子从小孔逸出是一个条件抽样过程：容器内随机选取一个分子，与随机选取一个实际穿孔事件，不具有同样的速度分布。前面的通量积分中 $v_z$ 权重既影响穿孔数，也影响每个逸出分子平均带走的能量。

\Needspace{180pt}
\begin{origexamplebox}{小孔逸出为何使绝热容器中的气体冷却}\label{ch18:ex:teaching-effusion}
刚性绝热容器中有稀薄单原子理想气体，粒子数和温度初值为 $N_0,T_0$。通过足够小的孔缓慢向真空逸出；逸出之间气体有足够时间恢复麦克斯韦分布。忽略孔边作用和回流。求一个逸出分子的平均动能，以及容器内剩余气体的温度随 $N$ 的关系。
\end{origexamplebox}
\Needspace{65pt}
\begin{xesolution}
取孔面外法线为 $z$ 轴。麦克斯韦分布三个速度分量独立；穿孔权重只含 $v_z>0$，故两个切向分量仍各贡献 $k_BT/2$。令 $a=m/(2k_BT)$，法向分量的加权二阶矩为
\[
 \langle v_z^2\rangle_{\rm hit}
 =\frac{\int_0^\infty v_z^3e^{-av_z^2}\dd v_z}
        {\int_0^\infty v_ze^{-av_z^2}\dd v_z}
 =\frac{1/(2a^2)}{1/(2a)}=\frac1a.
\]
两积分用 $u=av_z^2$ 即可得到。法向平均动能为 $m/(2a)=k_BT$，因此每个逸出分子平均携带
\[
 \overline E_{\rm out}=\frac12k_BT+\frac12k_BT+k_BT=2k_BT,
\]
大于容器内分子的 $3k_BT/2$。快分子更易逸出，剩余分子因而冷却。
容器边界固定且不传热，但有物质携能离开，不能写成 $\dd U=0$。在 $\dd N<0$ 的一步中，流出的粒子数为 $-\dd N$，能量变化为 $\dd U=2k_BT\dd N$。而
\[
 \dd U=\frac32k_B(N\dd T+T\dd N).
\]
两式相等得到 $3N\dd T=T\dd N$，积分为
\[
 \ln\frac{T}{T_0}=\frac13\ln\frac{N}{N_0},\qquad
 T=T_0\left(\frac{N}{N_0}\right)^{1/3}.
\]
$N$ 减小时 $T$ 减小，初值和符号均相符。这不是封闭气体的可逆绝热膨胀，不能套 $TV^{\gamma-1}=\mathrm{const}$；本题体积没有变化，变的是粒子数。若漏气过快、孔尺度超过平均自由程或来不及再平衡，穿孔抽样及此冷却律都须重新判断。
\end{xesolution}

\subsection{动力学与概率描述之间的条件}

单个微观态按确定性方程运动，不等于我们知道其初始条件。统计描述来自微观初态的不完全信息以及宏观分辨率的限制，
也可包括外界随机作用。系综是系统副本的概率集合；时间平均与系综平均是否相等，还需要动力学条件，不能由定义保证。
对于存在混沌的多自由度系统，初态和外界扰动的微小误差可能迅速放大；但不是所有多自由度运动都必然混沌。
理想的时间反演会让经典轨迹倒走，实际实验中反转误差和环境扰动使这一过程极难实现。
量子衰变则提供不同的随机模型：在恒定衰变率 $\lambda$ 近似下，生存概率满足
$\dd P/\dd t=-\lambda P$，所以 $P(t)=e^{-\lambda t}$。
若初始有 $N_0$ 个独立粒子，未衰变数均值为 $N_0P(t)$、方差为 $N_0P(t)[1-P(t)]$。
平均规律是指数衰减，不表示每个样本在每一时刻的粒子数都正好等于均值。
孤立系统细粒化吉布斯熵在哈密顿演化下保持不变；宏观熵增涉及粗粒化和特殊初始条件。

在稀薄气体的玻尔兹曼方程中，分子混沌假设把碰撞前两粒子分布近似为单粒子分布的乘积。
记一次碰撞前后乘积分别为 $b=f_1f_2$、$a=f'_1f'_2$。微观碰撞的正反对称性允许把
$H=\int f\ln f\,\dd^3r\dd^3v$ 的碰撞贡献整理成
\[\frac{\dd H}{\dd t}=-\frac14\int B\,(a-b)\ln\frac ab\,\dd\Lambda\le0.\]
这里 $B\ge0$ 是碰撞权重，$\dd\Lambda$ 表示对位置、两入射速度及散射方向积分；边界不向外输运熵。
因对数单调，$(a-b)(\ln a-\ln b)\ge0$，所以符号确定。该结果称为H定理，不能脱离稀薄气体和碰撞前无关联近似来理解。
平衡时所有允许碰撞满足 $f_1f_2=f'_1f'_2$，即 $\ln f_1+\ln f_2$ 在碰撞中守恒。
对普通弹性气体，它由粒子数、动量和能量这些碰撞不变量组成；在气体静止参考系中得到 $f\propto e^{-mv^2/(2k_BT)}$。

\section{从配分函数连接功、热与熵}

\subsection{连续相空间平均与非二次能量项}

若系统近似独立，$H=\sum_ih_i$，经典正则权重可分解。
这里的近独立可来自忽略相互作用，也可来自用平均场代替其他粒子的作用；后一种做法还要满足平均场的自洽条件，不能保证强关联系统适用。
单个自由度在连续情形的平均为
\[\langle A\rangle=
\frac{\int A(q,p)e^{-\beta h(q,p)}\,\dd q\dd p}
{\int e^{-\beta h(q,p)}\,\dd q\dd p}.\]
相空间测度的共同常数在平均值比式中抵消；计算绝对熵时则须保留它，例如以普朗克常数规定相空间单元。
能量均分的分部积分具体为
\[\int x\frac{\partial H}{\partial x}e^{-\beta H}\dd x
=-\frac1\beta\int x\,\partial_x(e^{-\beta H})\dd x
=\frac1\beta\int e^{-\beta H}\dd x,\]
最后一步要求边界项 $xe^{-\beta H}$ 消失。其他坐标积分后除以配分函数，即得前面的均分恒等式。

二次项之外不能机械分配 $k_BT/2$。例如在 $z\ge0$ 的均匀重力场中，势能为 $mgz$，
\[Z_z=\int_0^\infty e^{-\beta mgz}\dd z=\frac1{\beta mg},
\qquad \langle mgz\rangle=-\partial_\beta\ln Z_z=k_BT.\]
离散的两能级系统应使用求和；谐振子的经典势能和动能各为二次项，才各贡献 $k_BT/2$。
对于大小固定为 $d$ 的经典偶极子，在场强 $\mathcal E$ 中势能为 $-d\mathcal E\cos\theta$。
令 $x=\beta d\mathcal E$，其取向配分积分为
\[Z_{\rm rot}=2\pi\int_0^\pi e^{x\cos\theta}\sin\theta\dd\theta
=4\pi\frac{\sinh x}{x}.\]
对 $x$ 求对数导数得到 $\langle\cos\theta\rangle=\coth x-1/x$，从而
$\langle U_{\rm rot}\rangle=-d\mathcal E(\coth x-1/x)$。弱场时 $\coth x-1/x\simeq x/3$，
强场时趋于1，分别对应小偏向和接近完全取向。这里没有理由直接使用一个二次项的均分值。

\subsection{能级变化与布居变化}

系统平均能量为 $U=\sum_iP_iE_i$，所以
\begin{equation}
 \dd U=\sum_iP_i\dd E_i+\sum_iE_i\dd P_i.
\end{equation}
第一项是在概率不变时外参量改变能级的贡献，第二项是在能级给定时布居改变的贡献。
在弱耦合、准静态、仅有相应外参量功的热力学描述中，取外界对系统做功为正，
\[\delta W_{\rm on}=\sum_iP_i\dd E_i,\qquad
\delta Q=\sum_iE_i\dd P_i,\qquad \dd U=\delta W_{\rm on}+\delta Q.\]
若外参量为体积，$p=-\sum_iP_i\partial E_i/\partial V$，便得到 $\delta W_{\rm on}=-p\dd V$。
对于连续相空间，只需在固定的相空间测度下将上述求和换成积分。

正则概率满足 $\ln P_i=-\beta E_i-\ln Z$，因此
\[S=-k_B\sum_iP_i\ln P_i=k_B\ln Z+\frac UT.\]
又因 $\sum_i\dd P_i=0$，沿平衡态序列
\[\dd S=-k_B\sum_i(\ln P_i+1)\dd P_i
=k_B\beta\sum_iE_i\dd P_i=\frac{\delta Q_{\rm rev}}T.\]
从而 $F=U-TS=-k_BT\ln Z$，并与 $\dd F=-S\dd T-p\dd V$ 相容（固定粒子数）。
不可逆实际过程一般满足熵不等式，不能把任意传热都写成 $T\dd S$。
“绝热”仅表示不传热；快速改变外参量仍可能改变微观布居。只有满足相应绝热演化条件时才可进一步谈布居保持。


\Needspace{12\baselineskip}
\begingroup\tcbset{breakable/.default=true}\begin{starread}{玻色--爱因斯坦统计与费米--狄拉克统计}
不可分辨量子粒子需按交换对称性分类。平均占据数分别为
\[
 \bar n_B(\varepsilon)=\frac1{e^{(\varepsilon-\mu)/(k_BT)}-1},\qquad
 \bar n_F(\varepsilon)=\frac1{e^{(\varepsilon-\mu)/(k_BT)}+1}.
\]
式中负号对应玻色--爱因斯坦（Bose--Einstein）统计，正号对应费米--狄拉克（Fermi--Dirac）统计；
化学势 $\mu=(\partial F/\partial N)_{T,V}$ 描述固定温度、体积下增加粒子所需的自由能变化，其量纲与 $\varepsilon$ 相同。
后者服从泡利不相容原理（Pauli exclusion principle）。二者的经典极限均为麦克斯韦--玻尔兹曼分布；
平衡光子气有 $\mu=0$，低温电子气出现费米简并。
\end{starread}\endgroup

\endgroup

\hypersetup{bookmarksdepth=2}

\begingroup
\tcbset{breakable/.default=false,after={\par\nobreak\smallskip}}
\RenewDocumentEnvironment{xesolution}{+b}{%
  \par\nobreak\smallskip\nobreak\begingroup\small
  \clubpenalties 2 10000 10000\displaywidowpenalties 2 10000 10000\noindent
  {\bfseries\color{DeepBlue}解答}\quad #1\par\endgroup\smallskip
}{}
\begingroup
\renewcommand{\Needspace}[1]{\par\begingroup
\dimen0=#1\relax\vskip0pt plus\dimen0\penalty-100
\vskip0pt plus-\dimen0\vskip\dimen0\penalty9999\vskip-\dimen0
\vskip0pt\endgroup}
\tcbset{breakable/.default=false}
\chapter{静电场与电势}

\graphicspath{{/Users/dongxingxi/Documents/物理竞赛一轮资料/sample/xes-original/thirty-second-section/}}
\section{电荷、库仑定律与场的叠加}

% All text annotations are transparent. No label uses a white backing rectangle.
\tikzset{
  xeslabel/.style={font=\small,text=Ink,fill=none,inner sep=0.6pt},
  xessmall/.style={font=\footnotesize,text=TextGray,fill=none,inner sep=0.45pt},
  xestiny/.style={font=\scriptsize,text=TextGray,fill=none,inner sep=0.35pt},
  xesline/.style={draw=Ink,line width=0.85pt},
  xesblue/.style={draw=DeepBlue,line width=1.0pt},
  xesarrow/.style={-{Latex[length=2.15mm,width=1.35mm]},draw=DeepBlue,line width=0.9pt},
  fieldline/.style={draw=DeepBlue,line width=0.82pt,postaction={decorate},
    decoration={markings,mark=at position .64 with {\arrow{Latex[length=1.65mm,width=1.1mm]}}}},
  negfield/.style={draw=DeepBlue,line width=0.82pt,postaction={decorate},
    decoration={markings,mark=at position .36 with {\arrowreversed{Latex[length=1.65mm,width=1.1mm]}}}}
}


\newcommand{\matterScaleFigure}{%
\begin{center}
\begin{tikzpicture}[x=1cm,y=1cm]
  
  \node[xessmall,anchor=north west] at (.25,3.86) {尺度逐级放大};
  % atom
  \begin{scope}[shift={(1.55,1.85)}]
    \draw[DeepBlue!70,line width=.7pt,rotate=18] (0,0) ellipse (1.08 and .38);
    \draw[DeepBlue!70,line width=.7pt,rotate=-42] (0,0) ellipse (1.08 and .38);
    \draw[DeepBlue!70,line width=.7pt,rotate=78] (0,0) ellipse (1.08 and .38);
    \fill[Gold!80!Ink] (0,0) circle (2.8pt);
    \foreach \p in {(1.01,.13),(-.77,.46),(-.24,-.87)}{\fill[DeepBlue] \p circle (1.8pt);}
    \node[xeslabel] at (0,-1.45) {原子};
    \node[xestiny] at (0,-1.83) {$\sim10^{-8}\,\mathrm{cm}$};
  \end{scope}
  \draw[xesarrow] (2.75,1.86)--(3.55,1.86);
  % nucleus
  \begin{scope}[shift={(4.7,1.85)}]
    \foreach \x/\y/\c in {-.32/.18/DeepBlue,.18/.32/Gold,-.12/-.3/Gold,.38/-.18/DeepBlue}{
      \fill[\c!75] (\x,\y) circle (7pt);\draw[Ink!45,line width=.3pt] (\x,\y) circle (7pt);}
    \node[xeslabel] at (0,-1.45) {原子核};
    \node[xestiny] at (0,-1.83) {$10^{-13}$--$10^{-12}\,\mathrm{cm}$};
  \end{scope}
  \draw[xesarrow] (5.85,1.86)--(6.65,1.86);
  % nucleon
  \begin{scope}[shift={(7.6,1.85)}]
    \shade[ball color=DeepBlue!52] (0,0) circle (.63);
    \node[xeslabel] at (0,-1.45) {质子或中子};
    \node[xestiny] at (0,-1.83) {$\sim10^{-13}\,\mathrm{cm}$};
  \end{scope}
  \draw[xesarrow] (8.65,1.86)--(9.35,1.86);
  % quarks
  \begin{scope}[shift={(10.2,1.85)}]
    \fill[DeepBlue!78] (-.42,.24) circle (4pt);
    \fill[Gold!88] (.42,.24) circle (4pt);
    \fill[DeepBlue!78] (0,-.43) circle (4pt);
    \draw[RuleGray,line width=.65pt,densely dashed] (-.42,.24)--(.42,.24)--(0,-.43)--cycle;
    \node[xeslabel] at (0,-1.45) {夸克};
    \node[xestiny] at (0,-1.83) {$<10^{-16}\,\mathrm{cm}$};
  \end{scope}
  \draw[RuleGray,-{Latex[length=1.6mm]},line width=.55pt] (1.0,.32)--(.75,-.2);
  \node[xestiny,anchor=north] at (.72,-.24) {电子尺度 $<10^{-16}\,\mathrm{cm}$};
\end{tikzpicture}
\end{center}}

\newcommand{\coulombHistoryFigure}{%
\begin{center}
\begin{minipage}{.21\linewidth}\centering
\includegraphics[width=\linewidth]{/Users/dongxingxi/Documents/物理竞赛一轮资料/sample/xes-original/restructured-complete/assets/revision-ch19-ch20/coulomb.png}\par\small 库仑
\end{minipage}\hfill
\begin{minipage}{.76\linewidth}\centering
\begin{tikzpicture}[scale=.83]
 \draw[xesline] (-2,3.3)--(2,3.3);
 \draw[xesline] (0,3.3)--(0,.65);
 \draw[xesline] (-1.65,.65)--(1.65,.65);
 \fill[Ink] (-1.65,.65) circle (3pt);\fill[DeepBlue] (1.65,.65) circle (4pt);
 \fill[Gold] (2.3,.65) circle (4pt);\draw[xesline] (2.3,.65)--(2.3,-.5);
 \node[xessmall,above] at (1.65,.82) {$A$};\node[xessmall,above] at (-1.65,.82) {$D$};
 \node[xessmall,right] at (2.45,.65) {$B$};\node[xessmall,left] at (-.18,2.2) {悬丝};
 \node[xessmall,below] at (0,-.7) {侧视：两球同高};
 \begin{scope}[xshift=6.1cm]
  \coordinate(O) at (0,1);\coordinate(A) at (1.8,1.55);\coordinate(D) at (-1.8,.45);\coordinate(B) at (1.9,.6);
  \draw[RuleGray,dashed] (-2.1,1)--(2.1,1);
  \draw[xesline] (D)--(A);\fill[Ink] (O) circle (2pt);
  \fill[DeepBlue] (A) circle (4pt);\fill[Ink] (D) circle (3pt);\fill[Gold] (B) circle (4pt);
  \draw[xesarrow] (A)--($(A)+(-.1,.95)$);
  \draw[RuleGray,dotted] (B)--(A);
  \draw[RuleGray] (.7,1) arc[start angle=0,end angle=17,radius=.7];
  \node[xessmall] at (.96,1.13) {$\theta$};
  \node[xessmall,below] at (O) {$O$};\node[xessmall,left] at (D) {$D$};
  \node[xessmall,right] at ($(A)+(.1,.1)$) {$A$};\node[xessmall,right] at (B) {$B$};
  \node[xessmall,above] at (1.7,2.55) {$\boldsymbol F_{BA}$};
  \node[xessmall,below] at (0,-.7) {俯视：$O$ 为悬丝转轴};
 \end{scope}
\end{tikzpicture}\end{minipage}
\end{center}}


\newcommand{\fieldLineGallery}{%
\begin{center}
\includegraphics[width=.95\linewidth]{/Users/dongxingxi/Documents/物理竞赛一轮资料/sample/xes-original/restructured-complete/assets/revision-ch19-ch20/field-gallery.pdf}
\end{center}}


\newcommand{\gaussLawFigure}{%
\begin{center}
\begin{tikzpicture}[x=1cm,y=1cm]
  \node[inner sep=0pt] at (1.1,2.05) {\includegraphics[height=2.7cm]{/Users/dongxingxi/Documents/物理竞赛一轮资料/sample/xes-original/restructured-complete/assets/revision-ch19-ch20/gauss.png}};
  \node[xessmall] at (1.1,.38) {高斯};
  % local surface element
  \begin{scope}[shift={(4.25,2.05)}]
    \fill[DeepBlue!8] (-1,-.42)--(.75,-.72)--(1.05,.38)--(-.7,.7)--cycle;
    \draw[xesblue] (-1,-.42)--(.75,-.72)--(1.05,.38)--(-.7,.7)--cycle;
    \coordinate (P) at (0,.0);
    \fill[Ink] (P) circle (1.1pt);
    \draw[xesarrow] (P)--(.25,1.2);
    \draw[xesarrow] (P)--(.85,.83);
    \node[xessmall,anchor=west] at (.29,1.12) {$\dd\boldsymbol A_i$};
    \node[xessmall,anchor=west] at (.9,.77) {$\boldsymbol E_i$};
    \draw[RuleGray,line width=.55pt] (.07,.36) arc[start angle=79,end angle=44,radius=.43];
    \node[xestiny] at (.39,.74) {$\theta_i$};
    \node[xessmall] at (0,-1.8) {(a) 面元与外法线};
  \end{scope}
  % enclosed charge
  \begin{scope}[shift={(7.3,2.05)}]
    \draw[Ink,line width=.75pt] (0,0) ellipse (1.1 and .8);
    \foreach \a in {10,55,100,145,190,235,280,325}{
      \draw[fieldline] (-.15,.06)--++(\a:1.55);}
    \fill[DeepBlue] (-.15,.06) circle (5.0pt);\node[xestiny,text=white] at (-.15,.06) {$+$};
    \node[xessmall] at (0,-1.8) {(b) 电荷在闭面内};
  \end{scope}
  % external charge
  \begin{scope}[shift={(10.35,2.05)}]
    \draw[Ink,line width=.75pt] (0,0) ellipse (1.0 and .8);
    \fill[DeepBlue] (-1.65,.1) circle (5.0pt);\node[xestiny,text=white] at (-1.65,.1) {$+$};
    \draw[fieldline] (-1.58,.114)--(1.25,.68);
    \draw[fieldline] (-1.58,.1)--(1.3,.1);
    \draw[fieldline] (-1.58,.086)--(1.25,-.48);
    \node[xessmall] at (0,-1.8) {(c) 电荷在闭面外};
  \end{scope}
\end{tikzpicture}
\end{center}}

\newcommand{\differentialCubeFigure}{%
\begin{center}
\begin{tikzpicture}[x=1cm,y=1cm]
 \coordinate(A) at (0,0);\coordinate(B) at (4,0);\coordinate(C) at (5.25,1.25);\coordinate(D) at (1.25,1.25);
 \coordinate(E) at (0,3.1);\coordinate(F) at (4,3.1);\coordinate(G) at (5.25,4.35);\coordinate(H) at (1.25,4.35);
 \fill[DeepBlue!5] (A)--(B)--(F)--(E)--cycle;
 \fill[DeepBlue!8] (B)--(C)--(G)--(F)--cycle;
 \fill[DeepBlue!11] (E)--(F)--(G)--(H)--cycle;
 \draw[xesline] (A)--(B)--(C)--(G)--(H)--(E)--cycle (B)--(F)--(G) (E)--(F);
 \draw[RuleGray,dashed] (A)--(D)--(C) (D)--(H);
 \node[xessmall] at (1,-.32) {$\Delta x$};
 \node[xessmall,rotate=45] at (4.9,.35) {$\Delta y$};
 \node[xessmall] at (-.4,1.0) {$\Delta z$};
 \draw[xesarrow,dashed] (.625,2.175)--(-1.35,2.175);
 \node[xessmall,anchor=east,align=right] at (-1.45,2.175) {面1\\$E_x(x)$};
 \draw[xesarrow] (4.625,2.175)--(6.15,2.175);
 \node[xessmall,anchor=west,align=left] at (6.25,2.175) {面2\\$E_x(x+\Delta x)$};
 \draw[xesarrow] (2,1.55)--(.8,.35);
 \node[xessmall,anchor=east,align=right] at (-.25,.25) {面3：$E_y(y)$};
 \draw[RuleGray] (-.1,.25)--(.67,.32);
 \draw[xesarrow,dashed] (3.25,2.8)--(4.45,4.0);
 \node[xessmall,anchor=west,align=left] at (5.45,4.2) {面4\\$E_y(y+\Delta y)$};
 \draw[RuleGray] (4.5,4.03)--(5.3,4.2);
 \draw[xesarrow,dashed] (2.625,.625)--(2.625,-.75);
 \node[xessmall,anchor=north] at (2.625,-.82) {面5：$E_z(z)$};
 \draw[xesarrow] (2.625,3.725)--(2.625,5.15);
 \node[xessmall,anchor=south] at (2.625,5.22) {面6：$E_z(z+\Delta z)$};
 \node[xessmall,anchor=east] at (-.35,-.62) {$(x,y,z)$};
 \draw[RuleGray] (-.3,-.5)--(A);
\end{tikzpicture}
\end{center}}


\newcommand{\extendedChargeFigure}{%
\begin{center}
\begin{tikzpicture}[x=1cm,y=1cm]
  \draw[Ink,line width=1.1pt] (0,-2.15)--(0,2.15);
  \coordinate (P) at (4,0);\fill[Ink] (P) circle (1.3pt);\node[xessmall,anchor=west] at (4.12,0) {$P$};
  \draw[RuleGray,densely dashed] (0,0)--(P);
  \draw[DeepBlue,line width=.8pt] (0,2.15)--(P)--(0,-2.15);
  \draw[<->,RuleGray,line width=.6pt] (0,-2.6)--(4,-2.6);\node[xessmall,below] at (2,-2.6) {$a$};
  \draw[DeepBlue,line width=.65pt] (3.25,0) arc[start angle=180,end angle=151.7,radius=.75];
  \draw[DeepBlue,line width=.65pt] (3.25,0) arc[start angle=180,end angle=208.3,radius=.75];
  \node[xessmall] at (3.08,.34) {$\theta_1$};\node[xessmall] at (3.08,-.34) {$\theta_2$};
  \fill[DeepBlue] (0,1.05) circle (1.6pt);\node[xestiny,anchor=east] at (-.12,1.05) {$\dd q=\lambda\dd y$};
  \draw[xesarrow] (P)--++(1.2,-.315);\node[xessmall,anchor=north] at (5.05,-.55) {$\dd\boldsymbol E$};
  \node[xessmall,align=left,anchor=west] at (6.0,1.3) {同一几何关系用于有限直线、\\无限长条带和圆盘轴线积分；\\不同对象只改变微元的选取。};
  \draw[RuleGray,-{Latex[length=1.5mm]},line width=.6pt] (5.85,1.22)--(3.5,.45);
\end{tikzpicture}
\end{center}}

\newcommand{\unequalChargesFigure}{%
\begin{center}
\includegraphics[width=.76\linewidth]{/Users/dongxingxi/Documents/物理竞赛一轮资料/sample/xes-original/restructured-complete/assets/revision-ch19-ch20/unequal.pdf}
\end{center}}


\newcommand{\cavityFigure}{%
\begin{center}
\begin{tikzpicture}[x=1cm,y=1cm]
 \fill[DeepBlue!8] (0,0) circle (2.5);\draw[xesline] (0,0) circle (2.5);
 \fill[white] (.9,.65) circle (1);\draw[xesblue] (.9,.65) circle (1);
 \coordinate(O) at (0,0);\coordinate(C) at (.9,.65);\coordinate(P) at (1.5,.05);
 \foreach \p in {O,C,P}{\fill[Ink] (\p) circle (1.5pt);}
 \node[xessmall] at (-.23,-.22) {$O$};\node[xessmall] at (.88,.93) {$C$};\node[xessmall] at (1.72,-.12) {$P$};
 \draw[xesarrow] (O)--(C);\draw[xesarrow] (O)--(P);\draw[xesarrow] (C)--(P);
 \node[xessmall] at (.25,.58) {$\boldsymbol a$};\node[xessmall] at (.7,-.22) {$\boldsymbol r$};\node[xessmall] at (1.62,.56) {$\boldsymbol r'$};
 \node[xessmall,align=left,anchor=west] at (4,1) {完整均匀球：$+\rho$\\叠加的负球：$-\rho$};
 \node[xessmall,align=left,anchor=west] at (4,-.55) {空腔内 $\boldsymbol r-\boldsymbol r'=\boldsymbol a$，\\因而叠加场为匀强场。};
\end{tikzpicture}
\end{center}}


\newcommand{\dipoleFarFigure}{%
\begin{center}
\includegraphics[width=.95\linewidth]{/Users/dongxingxi/Documents/物理竞赛一轮资料/sample/xes-original/restructured-complete/assets/revision-ch19-ch20/dipole-far.pdf}
\end{center}}


\newcommand{\quadrupoleFarFigure}{%
\begin{center}
\begin{tikzpicture}[x=1cm,y=1cm]
  \fill[Ink] (-1.5,0) circle (4pt);\node[xessmall,below=4pt] at (-1.5,0) {$-q$};
  \fill[DeepBlue] (0,0) circle (5pt);\node[xessmall,below=5pt] at (0,0) {$+2q$};
  \fill[Ink] (1.5,0) circle (4pt);\node[xessmall,below=4pt] at (1.5,0) {$-q$};
  \draw[<->,RuleGray,line width=.55pt] (-1.5,-.7)--(0,-.7);\node[xessmall,below] at (-.75,-.7) {$l$};
  \draw[<->,RuleGray,line width=.55pt] (0,-.7)--(1.5,-.7);\node[xessmall,below] at (.75,-.7) {$l$};
  \coordinate (P) at (4.6,2.2);\fill[Ink] (P) circle (1.4pt);\node[xessmall,anchor=west] at ($(P)+(.08,0)$) {$P$};
  \draw[RuleGray,densely dashed] (0,0)--(P);\node[xessmall,above left] at (2.25,1.1) {$r$};
  \draw[DeepBlue,line width=.6pt] (1.0,0) arc[start angle=0,end angle=25.6,radius=1.0];\node[xessmall] at (1.24,.28) {$\theta$};
  \node[inner sep=0pt] at (8.55,1.25) {\includegraphics[width=4.4cm]{/Users/dongxingxi/Documents/物理竞赛一轮资料/sample/xes-original/restructured-complete/assets/revision-ch19-ch20/quadrupole.pdf}};
\end{tikzpicture}
\end{center}}

\newcommand{\chargedSlabFigure}{%
\begin{center}
\begin{tikzpicture}[x=1cm,y=1cm]
  \begin{scope}
    \fill[DeepBlue!10] (-1.35,-2.0) rectangle (1.35,2.0);
    \draw[DeepBlue,line width=.9pt] (-1.35,-2)--(-1.35,2) (1.35,-2)--(1.35,2);
    \foreach \x in {-1,-.5,0,.5,1}\foreach \y in {-1.55,-.75,.05,.85,1.65}{\node[xestiny,text=DeepBlue] at (\x,\y) {$+$};}
    \draw[xesline,-{Latex[length=2mm]}] (-3.25,0)--(3.25,0) node[xeslabel,below left] {$x$};
    \node[xessmall,below=3pt,fill=white,inner sep=1pt] at (-1.35,0) {$-d$};\node[xessmall,below=3pt,fill=white,inner sep=1pt] at (1.35,0) {$d$};
    \draw[xesarrow] (0,0)--(.8,0);\node[xessmall,above] at (.45,0) {$\boldsymbol v_0$};
    \fill[Ink] (0,0) circle (3pt);\node[xessmall,above left=3pt] at (0,0) {$-q$};
    \node[xessmall] at (0,-2.45) {(a) 带电无限厚板};
  \end{scope}
  \begin{scope}[shift={(7.3,0)},xscale=1.05,yscale=.85]
    \draw[xesline,-{Latex[length=2mm]}] (-3.2,0)--(3.2,0) node[xeslabel,below left] {$x$};
    \draw[xesline,-{Latex[length=2mm]}] (0,-2.35)--(0,2.35) node[xeslabel,below right] {$E_x$};
    \draw[RuleGray,densely dashed] (-1.35,-2.0)--(-1.35,2.0) (1.35,-2.0)--(1.35,2.0);
    \draw[xesblue] (-3,-1.5)--(-1.35,-1.5)--(1.35,1.5)--(3,1.5);
    \node[xessmall,below=2pt] at (-1.35,0) {$-d$};\node[xessmall,below=2pt] at (1.35,0) {$d$};
    \node[xessmall,anchor=west] at (1.55,1.94) {$\rho d/\varepsilon_0$};
    \node[xessmall] at (0,-2.78) {(b) 电场的分段线性关系};
  \end{scope}
\end{tikzpicture}
\end{center}}

\newcommand{\electronLensFigure}{%
\begin{center}
\begin{tikzpicture}[x=1cm,y=1cm]
  \draw[RuleGray,densely dashed] (-4,0)--(4,0);
  \fill[DeepBlue!12] (0,0) circle (8pt);\node[xeslabel] at (0,0) {$+$};
  \node[xessmall,below=9pt] at (0,0) {正线电荷 $\lambda$（垂直纸面）};
  \draw[Ink,line width=.9pt,postaction={decorate},decoration={markings,mark=at position .62 with {\arrow{Latex[length=2mm,width=1.25mm]}}}] (-4,1.55)--(4,1.55);
  \node[xessmall,anchor=east] at (-4.08,1.55) {电子};
  \draw[<->,RuleGray,line width=.55pt] (-2.8,0)--(-2.8,1.55);\node[xessmall,anchor=east] at (-2.88,.78) {$b$};
  \coordinate (A) at (-1.35,1.55);
  \draw[xesarrow] (A)--++(0,-1.0);\node[xessmall,anchor=west] at ($(A)+(0,-.5)$) {$\boldsymbol F_\perp$};
  \draw[xesarrow] (-3.15,1.55)--(-2.15,1.55);\node[xessmall,above] at (-2.65,1.55) {$\boldsymbol v$};
  \draw[DeepBlue,line width=1pt] (4,1.55) .. controls (4.65,1.49) and (5.15,1.36) .. (5.8,1.08);
  \draw[RuleGray,densely dashed] (4,1.55)--(5.85,1.55);
  \draw[DeepBlue,line width=.6pt] (4.75,1.55) arc[start angle=0,end angle=-14,radius=.75];
  \node[xessmall] at (5.02,1.37) {$\alpha$};
  \node[xessmall,align=left,anchor=west] at (6.35,1.0) {弱偏转近似中，计算沿未偏转直线\\的横向冲量；结果与瞄准距离 $b$ 无关。};
\end{tikzpicture}
\end{center}}

% Pure question geometry; analysis figure retained separately.
\newcommand{\emFiniteLineQuestion}{%
\begin{center}
\begin{tikzpicture}[x=1cm,y=1cm]
  \draw[Ink,line width=1.1pt] (0,-2.15)--(0,2.15);
  \coordinate (P) at (4,0);\fill[Ink] (P) circle (1.3pt);\node[xessmall,anchor=west] at (4.12,0) {$P$};
  \draw[RuleGray,densely dashed] (0,0)--(P);
  \draw[DeepBlue,line width=.8pt] (0,2.15)--(P)--(0,-2.15);
  \draw[<->,RuleGray,line width=.6pt] (0,-2.6)--(4,-2.6);\node[xessmall,below] at (2,-2.6) {$a$};
  \draw[DeepBlue,line width=.65pt] (3.25,0) arc[start angle=180,end angle=151.7,radius=.75];
  \draw[DeepBlue,line width=.65pt] (3.25,0) arc[start angle=180,end angle=208.3,radius=.75];
  \node[xessmall] at (3.08,.34) {$\theta_1$};\node[xessmall] at (3.08,-.34) {$\theta_2$};
\end{tikzpicture}
\end{center}}

% Pure question geometry; analysis figure retained separately.
\newcommand{\emQuadrupoleQuestion}{%
\begin{center}
\begin{tikzpicture}[x=1cm,y=1cm]
  \fill[Ink] (-1.5,0) circle (4pt);\node[xessmall,below=4pt] at (-1.5,0) {$-q$};
  \fill[DeepBlue] (0,0) circle (5pt);\node[xessmall,below=5pt] at (0,0) {$+2q$};
  \fill[Ink] (1.5,0) circle (4pt);\node[xessmall,below=4pt] at (1.5,0) {$-q$};
  \draw[<->,RuleGray,line width=.55pt] (-1.5,-.7)--(0,-.7);\node[xessmall,below] at (-.75,-.7) {$l$};
  \draw[<->,RuleGray,line width=.55pt] (0,-.7)--(1.5,-.7);\node[xessmall,below] at (.75,-.7) {$l$};
  \coordinate (P) at (4.6,2.2);\fill[Ink] (P) circle (1.4pt);\node[xessmall,anchor=west] at ($(P)+(.08,0)$) {$P$};
  \draw[RuleGray,densely dashed] (0,0)--(P);\node[xessmall,above left] at (2.25,1.1) {$r$};
  \draw[DeepBlue,line width=.6pt] (1.0,0) arc[start angle=0,end angle=25.6,radius=1.0];\node[xessmall] at (1.24,.28) {$\theta$};
\end{tikzpicture}
\end{center}}

% Pure question geometry; analysis figure retained separately.
\newcommand{\emSlabQuestion}{%
\begin{center}
\begin{tikzpicture}[x=1cm,y=1cm]
  \begin{scope}
    \fill[DeepBlue!10] (-1.35,-2.0) rectangle (1.35,2.0);
    \draw[DeepBlue,line width=.9pt] (-1.35,-2)--(-1.35,2) (1.35,-2)--(1.35,2);
    \foreach \x in {-1,-.5,0,.5,1}\foreach \y in {-1.55,-.75,.05,.85,1.65}{\node[xestiny,text=DeepBlue] at (\x,\y) {$+$};}
    \draw[xesline,-{Latex[length=2mm]}] (-3.25,0)--(3.25,0) node[xeslabel,below left] {$x$};
    \node[xessmall,below=3pt,fill=white,inner sep=1pt] at (-1.35,0) {$-d$};\node[xessmall,below=3pt,fill=white,inner sep=1pt] at (1.35,0) {$d$};
    \draw[xesarrow] (0,0)--(.8,0);\node[xessmall,above] at (.45,0) {$\boldsymbol v_0$};
    \fill[Ink] (0,0) circle (3pt);\node[xessmall,above left=3pt] at (0,0) {$-q$};
    \node[xessmall] at (0,-2.45) {(a) 带电无限厚板};
  \end{scope}
  \end{tikzpicture}
\end{center}}


电磁相互作用是自然界的基本相互作用之一。我们日常接触到的摩擦力、支持力、弹力以及绝大多数化学作用，
从微观上说都与带电粒子之间的电磁作用有关。法拉第把“场”引入物理学：电荷并不是隔着空间瞬时作用于
另一个电荷，而是在周围空间建立电场，另一个电荷再与所在位置的电场发生作用。麦克斯韦把电与磁统一起来，
并预言电磁波；量子力学建立以后，量子电动力学又给出了电磁相互作用的量子描述。电磁学因此既是经典
物理的重要组成部分，也是相对论和量子理论发展的出发点之一。


\subsection{物质结构和电荷}


通常物质由分子组成，分子由原子组成。原子的典型尺度约为 $10^{-10}\,\mathrm m$，内部包含尺度约为
$1$--$10\,\mathrm{fm}$（$10^{-15}$--$10^{-14}\,\mathrm m$，随核的质量数而变）的原子核和核外电子。原子核由质子与中子构成，质子和中子统称核子；核子内部还有
夸克和胶子等更基本的自由度。图~\ref{ch19:fig:s32-1-1}把这些层级关系放在同一条尺度链上，图中的相对大小只是
示意，尺度数字才是定量信息。

\par\noindent\begin{minipage}{\linewidth}
\matterScaleFigure
\diagramcaption{从原子、原子核到核子和夸克的结构层级。箭头表示继续放大观察，并不表示实物相互分离。}\label{ch19:fig:s32-1-1}
\end{minipage}\par\medskip

电荷具有三条基本性质。

第一，\textbf{可加性}。一个系统的总电荷等于各部分电荷的代数和。第二，\textbf{守恒性}。在孤立系统中，
电荷既不会凭空产生，也不会凭空消失；成对产生或湮灭的正、负电荷总量相等。第三，\textbf{量子化}。
自然界可单独观测到的物体所带电荷通常是元电荷 $e$ 的整数倍。电子带电 $-e$，质子带电 $+e$，其中
\begin{equation}
  \boxed{e=1.602\,176\,634\times10^{-19}\,\mathrm C}.
\end{equation}
夸克及反夸克的电荷可为 $\pm e/3$ 或 $\pm2e/3$，但夸克在通常条件下不会以自由粒子的形式单独出现。


% EM teaching enrichment 20260927
\paragraph{从电荷守恒到可测量的电子数}
摩擦起电通常是电子转移，并非把中性物质中的电荷“创造”出来。若两个原先中性的绝缘物体只彼此交换电子，最后必有 $Q_A+Q_B=0$；若还与地面相通，就应把大地计入系统。
例如物体带电 $-3.20\,\mathrm{nC}$，多出的电子数约为
\[
 N=\frac{|Q|}{e}=\frac{3.20\times10^{-9}}{1.60\times10^{-19}}=2.00\times10^{10}.
\]
电子数很大，所以宏观电荷密度可以连续处理；这不取消微观电荷量子化。解题先选封闭的电荷交换系统，再列守恒式，最后判断是否需要离散电子数。


\subsection{库仑定律}


同号电荷相斥，异号电荷相吸。真空中两个静止点电荷 $Q_1,Q_2$ 相距 $r$ 时，相互作用力的大小为
\begin{equation}
  \boxed{F=\frac{1}{4\pi\varepsilon_0}\frac{|Q_1Q_2|}{r^2}}.
\end{equation}
若保留 $Q_1Q_2$ 的正负号，则还必须同时规定径向单位矢量的方向；只写“大小”时应取绝对值。真空介电
常量为 $\varepsilon_0$，常用的静电力常量为
\begin{equation}
  \boxed{\frac{1}{4\pi\varepsilon_0}
  \approx8.99\times10^9
  \,\mathrm{N\,m^2\,C^{-2}}}.
\end{equation}

库仑定律的直接形式要求：电荷可视为点电荷，电荷相对静止或速度远小于光速，并且周围介质可视为
真空。若带电体的线度远小于彼此间距，就可以把其电荷集中在一点；若带电体形状、大小不可忽略，就要把
它分割成电荷元后作矢量积分。图~\ref{ch19:fig:s32-1-2}左侧是库仑肖像，中间显示悬丝与横杆，右侧是与横杆同高的水平俯视图。固定带电球 $B$ 与活动带电球 $A$ 处在同一水平面；$O$ 是悬丝转轴，$D$ 是平衡物。$A$ 受到沿 $BA$ 的斥力，只有该力相对 $O$ 的力矩使悬丝扭转。扭丝转过的角度反映
静电力矩，从而可以检验力随距离的平方反比关系。

\par\noindent\begin{minipage}{\linewidth}
\coulombHistoryFigure
\diagramcaption{库仑与扭秤。活动球 $A$、固定球 $B$ 与横杆同高；俯视图的虚线为横杆零位。}\label{ch19:fig:s32-1-2}
\end{minipage}\par\medskip

静电力与万有引力都服从平方反比律，因此两体问题的轨道都可以是圆锥曲线。引力永远相吸；静电力既可
相吸也可相斥。同号电荷散射时出现的是双曲线的另一支，这正是作用方向不同带来的几何差别。

把相距无穷远处的相互作用势能取为零，两个点电荷的静电势能为
\begin{equation}
  \boxed{E_p=\frac{1}{4\pi\varepsilon_0}\frac{Q_1Q_2}{r}}.
\end{equation}
同号电荷的势能为正，异号电荷的势能为负。以两个质子为例，库仑斥力与万有引力之比为
\begin{equation}
  \frac{F_C}{F_G}
  =\frac{e^2}{4\pi\varepsilon_0Gm_p^2}
  \approx1.24\times10^{36}.
\end{equation}
可见在原子和分子尺度，重力通常完全可以忽略。原子核中的质子却能结合在一起，是因为短程强相互作用
在核尺度上压倒了库仑斥力。


% EM teaching enrichment 20260927
\paragraph{带符号矢量式与对称位置的稳定性}
令源电荷位于 $\boldsymbol r_i$，试探电荷位于 $\boldsymbol r$，统一写成
\[
 \boldsymbol F_i=kqQ_i\frac{\boldsymbol r-\boldsymbol r_i}{|\boldsymbol r-\boldsymbol r_i|^3},
 \qquad k=\frac1{4\pi\varepsilon_0}.
\]
这样位移矢量始终从源指向观察点，吸引或排斥由 $qQ_i$ 的符号给出。取两个固定正电荷 $Q$ 位于 $x=\pm a$，小正电荷 $q$ 在原点受力为零。沿轴小移 $x$ 后
\[
 F_x=kQq\left[\frac1{(a+x)^2}-\frac1{(a-x)^2}\right]
 \simeq-\frac{4kQq}{a^3}x.
\]
轴向有回复力；但垂直轴小移 $y$ 后 $F_y\simeq2kQq y/a^3$，反而远离原点。因此“合力为零”不等于三维稳定平衡，也不能只检验一个方向。


\subsection{电场}


若空间中已经存在电场，把电荷量为 $q$ 的小试探电荷放在某点，它受到的电场力定义为
\begin{equation}
  \boxed{\boldsymbol F=q\boldsymbol E}.
\end{equation}
所以 $\boldsymbol E=\boldsymbol F/q$。电场是空间本身的属性，并不依赖我们是否真的放入试探电荷；试探
电荷必须足够小，才不会显著改变原有电荷分布。一般情况下电场是位置和时间的函数
$\boldsymbol E(x,y,z,t)$。

电场线是一组有方向的曲线，曲线上每一点的切线方向与该点电场方向相同；线越密的地方通常表示场越强。
场线只是描述矢量场的几何工具，不是带电粒子的实际轨迹，也不是空间中真实存在的“细线”。\footnote{电场线
与流体力学中的流线相似：都用曲线的切向表示一个矢量场，但流线也不一定等同于流体质点的完整轨迹。}
图~\ref{ch19:fig:s32-1-3}前两幅是静电场，第三幅表示电荷振荡时场的变化不能瞬时传到远处，远处场线会出现弯折并形成向外
传播的电磁扰动。

\par\noindent\begin{minipage}{\linewidth}
\fieldLineGallery
\diagramcaption{三种电场线结构。(a) 异号电荷；(b) 四极电荷分布；(c) 振荡电荷的瞬时场线。箭头均表示电场方向。}\label{ch19:fig:s32-1-3}
\end{minipage}\par\medskip

点电荷 $Q$ 在位置矢量 $\boldsymbol r$ 处产生的电场为
\begin{equation}
  \boxed{\boldsymbol E(\boldsymbol r)
  =\frac{1}{4\pi\varepsilon_0}\frac{Q}{r^3}\boldsymbol r
  =\frac{1}{4\pi\varepsilon_0}\frac{Q}{r^2}\hat{\boldsymbol r}}.
\end{equation}
$Q>0$ 时场沿径向向外，$Q<0$ 时场沿径向向内。多个电荷产生的总场等于各电荷产生电场的矢量和，这就是
电场的叠加原理。


% EM teaching enrichment 20260927
\paragraph{连续分布：先定电荷元，再消去分量}
线、面、体分布分别用 $\dd q=\lambda\dd l$、$\sigma\dd S$、$\rho\dd V$，并积分
$\dd\boldsymbol E=k\,\dd q\,\boldsymbol R/R^3$。这里 $\boldsymbol R$ 从电荷元指向观察点，不能与电荷元的位置矢量混淆。
以半径 $a$、总电荷 $Q$ 的均匀细圆环为例，轴线上距圆心 $z$ 处的横向场两两抵消，而每个元的轴向投影含因子 $z/\sqrt{a^2+z^2}$，故
\[
 E_z=\int\frac{kz\,\dd q}{(a^2+z^2)^{3/2}}
 =\frac{kQz}{(a^2+z^2)^{3/2}}.
\]
$z=0$ 时对称性给出零场，$z\gg a$ 时恢复点电荷场。先写“哪些分量为何相消”，可避免把矢量大小直接相加。


% EM structure and teaching revision 20260928
\subsection{连续带电体积分的建模步骤}
求场时先说明源点与观察点，不要把电荷元的位置直接当成场矢量的方向。若观察点为 $\boldsymbol r$，源点为 $\boldsymbol r'$，则
\[
 \dd\boldsymbol E=\frac{1}{4\pi\varepsilon_0}
 \frac{\boldsymbol r-\boldsymbol r'}{|\boldsymbol r-\boldsymbol r'|^3}\dd q.
\]
线、面、体分布分别选 $\dd q=\lambda\dd l$、$\sigma\dd S$、$\rho\dd V$。先用对称性消去积分后必为零的分量，再计算剩下的分量；只有对同一观察点求得的矢量才可以叠加。
以半径 $R$、总电荷 $Q$ 的均匀圆环为例，轴上距中心 $z$ 处各电荷元的横向分量抵消，轴向分量都含投影因子 $z/\sqrt{R^2+z^2}$，于是
\[
 E_z=\frac{Qz}{4\pi\varepsilon_0(R^2+z^2)^{3/2}}.
\]
当 $|z|\gg R$ 时退回点电荷场；在中心附近 $E_z\simeq Qz/(4\pi\varepsilon_0R^3)$。若 $Q>0$，正试探电荷沿轴偏离中心后受力继续向外，负试探电荷则沿轴受到回复力。这个结论只判断轴向稳定性，不能据此认定三维空间中的稳定平衡。
检查积分时还应核对量纲、远场主项和特殊对称点。对有限带电体，远场通常先由总电荷决定；总电荷为零时才需要进一步保留偶极或更高阶贡献。


\subsection{专题例题与方法应用}

\begin{origexamplebox}{【统统算出来】}\label{ch19:ex:s32-1-1}
小明认为电场的高斯定理有点神奇，不是很相信高斯定理的结果，于是希望利用微积分的方式求解一些
对称系统的电场。另外，利用微积分也可以计算一些不是无穷的系统的电场。请帮助小明计算下列均匀带电
系统产生的电场：
\begin{enumerate}[label=\arabic*.]
  \item 无限长直导线外、距导线 $a$ 的一点；
  \item 有限长直导线中点附近、距导线 $a$ 的一点，点到两端的连线与垂线所成角分别为
  $\theta_1,\theta_2$，如图；
  \item 无限大均匀带电平面外、距平面 $a$ 的一点；
  \item 只有一个方向无限延伸的均匀带电平面（无限长条带）外、距平面 $a$ 的一点，两侧边缘对应的角为
  $\theta_1,\theta_2$；
  \item 均匀带电圆盘中轴线上、距圆盘中心 $a$ 的一点，圆盘边缘对轴线上该点的半张角为 $\theta$。
\end{enumerate}
直线电荷的线密度记为 $\lambda$，平面和圆盘的面密度记为 $\sigma$。
\par\noindent\begin{minipage}{\linewidth}
\emFiniteLineQuestion
\figcap{有限直线的端点与观察点；长条带的截面使用同样的端点角约定。}\label{ch19:fig:em-question-1}
\end{minipage}\par\medskip
\end{origexamplebox}\par\nopagebreak[4]

\begin{xesolution}
五个对象都可以拆成“已经知道电场的微元”，再利用对称性消去不需要的分量。图中把有限直线画成竖直，
$P$ 到直线的垂距为 $a$；对长条带和圆盘，仍可沿用相同的截面角度关系。

\par\noindent\begin{minipage}{\linewidth}
\extendedChargeFigure
\diagramcaption{有限长直线的积分几何。$\theta_1,\theta_2$ 分别对应垂足上方和下方的端点。}\label{ch19:fig:s32-1-6}
\end{minipage}\par\medskip

\textbf{1. 无限长直导线。} 先直接积分。取导线沿 $y$ 轴，$P=(a,0)$，电荷元
$\dd q=\lambda\dd y$。电场的横向分量为
\[
  \dd E_x=\frac{1}{4\pi\varepsilon_0}
  \frac{\lambda a\,\dd y}{(a^2+y^2)^{3/2}}.
\]
上下两个电荷元的纵向分量互相抵消，所以
\[
  E=\frac{\lambda}{4\pi\varepsilon_0}
  \int_{-\infty}^{\infty}\frac{a\,\dd y}{(a^2+y^2)^{3/2}}
  =\boxed{\frac{\lambda}{2\pi\varepsilon_0a}}.
\]
$\lambda>0$ 时方向垂直导线向外，$\lambda<0$ 时方向向内。也可用半径 $a$、长度 $L$ 的同轴圆柱
高斯面立即得到 $E(2\pi aL)=\lambda L/\varepsilon_0$。

\textbf{2. 有限长直导线。} 规定从垂线向上、向下量得的端点角分别为 $\theta_1,\theta_2$。
对应端点为 $y_1=a\tan\theta_1$、$y_2=-a\tan\theta_2$。沿用图~\ref{ch19:fig:s32-1-6}的坐标，
\[
 \dd E_x=k\lambda\frac{a\,\dd y}{(a^2+y^2)^{3/2}},\qquad
 \dd E_y=-k\lambda\frac{y\,\dd y}{(a^2+y^2)^{3/2}},\quad k=\frac1{4\pi\varepsilon_0}.
\]
负号来自电荷元指向 $P$ 的纵向位移为 $-y$。令 $y=a\tan u$，则 $\dd y=a\sec^2u\dd u$，有
\[
 E_x=\frac{k\lambda}{a}\int_{-\theta_2}^{\theta_1}\cos u\,\dd u,
 \qquad E_y=-\frac{k\lambda}{a}\int_{-\theta_2}^{\theta_1}\sin u\,\dd u.
\]
上下限代入后，垂直导线方向和平行导线方向的分量分别为
\[
  \boxed{E_\perp=\frac{\lambda}{4\pi\varepsilon_0a}
  (\sin\theta_1+\sin\theta_2)},
  \qquad
  \boxed{E_\parallel=\frac{\lambda}{4\pi\varepsilon_0a}
  (\cos\theta_1-\cos\theta_2)}.
\]
这里 $E_\parallel$ 的正方向取向上；若换用向下为正，第二式整体变号。导线关于垂足对称时
$\theta_1=\theta_2=\theta$，平行分量为零，
\[
  \boxed{E=\frac{\lambda}{2\pi\varepsilon_0a}\sin\theta}.
\]
让两端趋向无穷远，$\theta\to\pi/2$，就回到第1问。

\textbf{3. 无限大平面。} 为直接积分，以 $P$ 在平面上的垂足为圆心，将平面分成半径 $s$、宽度 $\dd s$ 的圆环。
环电荷 $\dd q=2\pi\sigma s\dd s$，到 $P$ 的距离为 $\sqrt{a^2+s^2}$。环内横向场相消，法向分量相加：
\[
 E_\perp=\frac{\sigma a}{2\varepsilon_0}\int_0^\infty
 \frac{s\,\dd s}{(a^2+s^2)^{3/2}}
 =\frac{\sigma a}{2\varepsilon_0}
 \left[-\frac1{\sqrt{a^2+s^2}}\right]_0^\infty
 =\frac{\sigma}{2\varepsilon_0}\quad(a>0).
\]
这个结果由库仑积分独立得到，接着用高斯定律校验。取横跨平面的薄圆柱高斯面。两个端面上的场大小相等、方向分别向两侧，
侧面通量为零。设端面面积为 $A$，则
\[
  2EA=\frac{\sigma A}{\varepsilon_0},
  \qquad
  \boxed{E=\frac{\sigma}{2\varepsilon_0}}.
\]
电场与距离 $a$ 无关。正电面两侧的场都背离平面，负电面两侧的场都指向平面。

\textbf{4. 无限长条带。} 令条带位于 $z=0$ 平面，沿 $y$ 方向无限长；横向坐标范围为
$-b_1\le x\le b_2$，观察点为 $(0,0,a)$。把宽度为 $\dd x$ 的窄条看成线电荷
$\dd\lambda=\sigma\dd x$。由第1问，窄条在 $P$ 处产生
\[
  \dd E=\frac{\sigma\dd x}{2\pi\varepsilon_0\sqrt{a^2+x^2}}.
\]
令 $\theta_1=\arctan(b_1/a)$、$\theta_2=\arctan(b_2/a)$，电荷元到 $P$ 的方向正比于 $(-x,0,a)$，故
\[
 E_z=\frac{\sigma}{2\pi\varepsilon_0}\int_{-b_1}^{b_2}\frac{a\,\dd x}{a^2+x^2}
 =\frac{\sigma}{2\pi\varepsilon_0}\left[\arctan\frac xa\right]_{-b_1}^{b_2},
\]
\[
 E_x=-\frac{\sigma}{2\pi\varepsilon_0}\int_{-b_1}^{b_2}\frac{x\,\dd x}{a^2+x^2}
 =-\frac{\sigma}{4\pi\varepsilon_0}\ln\frac{a^2+b_2^2}{a^2+b_1^2}.
\]
再用 $\cos\theta_i=a/\sqrt{a^2+b_i^2}$，得到
\[
  \boxed{E_z=\frac{\sigma}{2\pi\varepsilon_0}(\theta_1+\theta_2)},
  \qquad
  \boxed{E_x=\frac{\sigma}{2\pi\varepsilon_0}
  \ln\frac{\cos\theta_2}{\cos\theta_1}}.
\]
对称条带有 $\theta_1=\theta_2=\theta$，横向分量消失，
$E_z=\sigma\theta/(\pi\varepsilon_0)$。宽度趋于无穷时 $\theta\to\pi/2$，又回到无限大平面的
$\sigma/(2\varepsilon_0)$。

\textbf{5. 均匀带电圆盘。} 把半径为 $R$ 的圆盘分割成半径 $r$、宽度 $\dd r$ 的同心圆环。
环上横向分量互相抵消，轴向分量为
\[
  \dd E_z=\frac{1}{4\pi\varepsilon_0}
  \frac{a\,(\sigma 2\pi r\dd r)}{(a^2+r^2)^{3/2}}.
\]
从 $0$ 积分到 $R$：
\[
  E_z=\frac{\sigma}{2\varepsilon_0}
  \left(1-\frac{a}{\sqrt{a^2+R^2}}\right).
\]
因 $\cos\theta=a/\sqrt{a^2+R^2}$，所以
\[
  \boxed{E_z=\frac{\sigma}{2\varepsilon_0}(1-\cos\theta)}.
\]
当 $R\to\infty$ 时 $\theta\to\pi/2$，结果连续过渡到无限大平面；当 $a\gg R$ 时，展开后得到
$E\approx(1/4\pi\varepsilon_0)(\sigma\pi R^2/a^2)$，即远处只看见总电荷。
\end{xesolution}



\begin{origexamplebox}{【奇妙的电场线】}\label{ch19:ex:s32-1-2}
小明发现两个点电荷的电场线很有趣，想通过作图描绘电场线的形状；可是小明希望画得准确一些。已知
$+Q$、$-q$ 分别位于 $(-a_0,0)$、$(a_0,0)$ 处，且 $Q>q>0$。系统绕二者连线旋转对称，因此只需先画
包含连线的平面内的电场线。回答下列问题：
\begin{enumerate}[label=\arabic*.]
  \item 是否有从 $+Q$ 出发的电场线延伸到无穷远？是否有从无穷远出发、终止于 $-q$ 的电场线？
  \item 零场点在哪里？该点附近的电场线有什么特征？
  \item 远离两个电荷处的电场具有什么特征？
  \item 写出包含 $x$ 轴的平面内的电场线方程。
\end{enumerate}
\end{origexamplebox}\par\nopagebreak[4]

\begin{xesolution}
总电荷为 $Q-q>0$。因此 $+Q$ 发出的场线中，数量与 $q$ 对应的那一部分终止于 $-q$，其余部分必须
延伸到无穷远；$+Q$ 足以提供全部终止于 $-q$ 的场线，所以不需要从无穷远引入场线。图中实线为场线，
场线由下文的第一积分计算。通向无穷远与终止于 $-q$ 的场线之间存在分界结构，它趋近零场点；零场点处方向未定义，不能把它当普通连接点。

\unequalChargesFigure

零场点只能在两个电荷连线的外侧。由于 $Q>q$，它位于较小电荷 $-q$ 的右侧。设零场点坐标为
$x_0>a_0$，两场反向且大小相等：
\[
  \frac{Q}{(x_0+a_0)^2}=\frac{q}{(x_0-a_0)^2}.
\]
取正平方根并整理：
\[
  \sqrt Q(x_0-a_0)=\sqrt q(x_0+a_0),
  \qquad
  \boxed{x_0=a_0\frac{\sqrt Q+\sqrt q}{\sqrt Q-\sqrt q}}.
\]
在零场点附近，沿 $x$ 轴从左向右穿过该点，电场由向左变为向右；垂直于 $x$ 轴的小位移则受到指向
轴线的电场。局部展开可写成
\[
  \boldsymbol E\approx\kappa\left[(x-x_0)\boldsymbol e_x
  -\frac{y}{2}\boldsymbol e_y-\frac{z}{2}\boldsymbol e_z\right],\qquad\kappa>0.
\]
所以它是一个鞍点，不是场线可以平滑穿过的普通点；轴线和分界面构成局部的分离结构。

当 $r\gg a_0$ 时，最低阶电场只取决于总电荷：
\[
  \boxed{\boldsymbol E\approx\frac{1}{4\pi\varepsilon_0}
  \frac{Q-q}{r^2}\hat{\boldsymbol r}}.
\]
因此远处场线近似从一个总电荷为 $Q-q$ 的正点电荷径向向外；再下一阶才看出偶极修正。

最后求电场线方程。取包含 $x$ 轴的平面，令点的坐标为 $(x,y)$，到两电荷的距离分别为
\[
  r_1=\sqrt{(x+a_0)^2+y^2},\qquad r_2=\sqrt{(x-a_0)^2+y^2}.
\]
轴对称点电荷场的流函数可以取为电荷量乘以相应极角的余弦。叠加后，电场线是下式的等值线：
\[
  \boxed{Q\frac{x+a_0}{r_1}-q\frac{x-a_0}{r_2}=C}.
\]
这一第一积分可以直接核验。令 $F=Q(x+a_0)/r_1-q(x-a_0)/r_2$，
逐项求导得到 $F_x=yE_y/k$、$F_y=-yE_x/k$。沿电场线的切向位移满足
$(\dd x,\dd y)\parallel(E_x,E_y)$，所以
\[
 \dd F=F_x\dd x+F_y\dd y=\frac yk(E_y\dd x-E_x\dd y)=0.
\]
故 $F=C$ 确为场线方程，而不是单凭外观猜出的曲线。改变常数 $C$ 就得到图中的不同场线；绕 $x$ 轴旋转这些曲线，可得到完整的三维场线曲面。
\end{xesolution}











\section{高斯定律与对称带电体}

\subsection{电场高斯定律——积分形式}


对于面积为 $A$ 的平面，若匀强电场 $\boldsymbol E$ 与面积外法线的夹角为 $\theta$，穿过该面的电通量为
\begin{equation}
  \boxed{\Phi_E=EA\cos\theta\equiv\boldsymbol E\cdot\boldsymbol A}.
\end{equation}
其中面积矢量 $\boldsymbol A=A\boldsymbol n$ 的方向规定为法线方向。对于任意曲面，把它划分成足够小的面元，
每个面元上电场近似不变，总通量为
\[
  \Phi_E=\int_S\boldsymbol E\cdot\dd\boldsymbol A.
\]
若 $S$ 是闭合曲面，面元方向统一取向外，高斯定律给出
\begin{equation}
  \boxed{\oint_S\boldsymbol E\cdot\dd\boldsymbol A
  =\frac{Q_{\mathrm{in}}}{\varepsilon_0}}.
\end{equation}
闭面外的电荷虽然会改变闭面上各点的电场，却不会改变穿入、穿出闭面的净通量：每一束从外部进入的场线
还会从另一处离开。闭面内正电荷是场线的源，负电荷是场线的汇；在无电荷区域，场线不能凭空开始或终止。

\par\noindent\begin{minipage}{\linewidth}
\gaussLawFigure
\diagramcaption{高斯定律的几何意义。(a) 面元的法向面积矢量；(b) 内部电荷产生非零净通量；(c) 外部电荷的穿入与穿出相抵。}\label{ch19:fig:s32-1-4}
\end{minipage}\par\medskip

图~\ref{ch19:fig:s32-1-4}可以帮助理解通量，却不是高斯定律的独立证明，因为若先用点电荷的库仑场画出这些场线，
推理本身已经使用了库仑定律。更根本地说，高斯定律是电磁场的基本规律之一；即使电荷正在运动，
对任意时刻的电场仍有同样的积分关系。只靠“数场线”不能代替定量证明：场线条数是人为选取的，
真正不依赖画法的量是曲面积分 $\oint\boldsymbol E\cdot\dd\boldsymbol A$。


% EM teaching enrichment 20260927
\paragraph{何时能由通量直接求场强}
高斯定律总成立，但只有对称性使选定面上的场强可提出积分号时，才成为简便的求场方法。球对称选同心球面；无限长轴对称选同轴圆柱；无限均匀平面选跨面的柱形盒。
例如半径 $a$、均匀体电荷密度 $\rho$ 的绝缘实心球，半径 $r<a$ 的高斯面只包围 $4\pi\rho r^3/3$，于是
\[
 E_r=\frac{\rho r}{3\varepsilon_0}\quad(r<a),\qquad
 E_r=\frac{\rho a^3}{3\varepsilon_0r^2}\quad(r>a).
\]
两式在表面连续，球心为零；球内不能把总电荷全部代入 $kQ/r^2$。若给球外再放一个点电荷，闭面的总通量仍只取决于内部电荷，但面上总场已不再均匀，不能再用 $E\,4\pi r^2$ 表示通量。


\subsection{电场高斯定律——微分形式}


把一个有限区域分成许多很小的长方体。相邻体积元的公共面上，一个体积元的外法线恰好是另一个体积元
外法线的反方向，所以内部公共面的通量两两抵消，只剩整个区域外边界的净通量。反过来考察图~\ref{ch19:fig:s32-1-5}中的一个
小长方体，就能从积分形式得到局域的微分形式。

\par\noindent\begin{minipage}{\linewidth}
\differentialCubeFigure
\diagramcaption{位于 $(x,y,z)$ 附近的微小长方体。每一对相对表面的外法线方向相反，图中箭头画的是相应通量项的方向。}\label{ch19:fig:s32-1-5}
\end{minipage}\par\medskip

以垂直于 $x$ 轴的两面为例，右面流出的通量减去左面流入的通量为
\[
  \bigl[E_x(x+\Delta x)-E_x(x)\bigr]\Delta y\Delta z
  \approx\frac{\partial E_x}{\partial x}\Delta x\Delta y\Delta z.
\]
$y,z$ 两个方向完全类似。体积元内电荷为
$\rho\Delta x\Delta y\Delta z$。代入高斯定律并约去体积，得到
\begin{equation}
  \boxed{\frac{\partial E_x}{\partial x}
  +\frac{\partial E_y}{\partial y}
  +\frac{\partial E_z}{\partial z}
  =\frac{\rho}{\varepsilon_0}}
  \qquad\text{即}\qquad
  \boxed{\boldsymbol\nabla\cdot\boldsymbol E=\frac{\rho}{\varepsilon_0}}.
\end{equation}
这说明电荷密度决定电场在该点的“发散程度”。积分形式联系一个闭面与其包围的总电荷，微分形式则把
规律落实到每一个空间点；二者是同一条物理定律的两种表述。

本章随后反复使用六个核心概念：元电荷、电荷可视为点电荷的条件、库仑平方反比律、电场强度
$\boldsymbol E=\boldsymbol F/q$、电场线以及高斯定律
$\oint\boldsymbol E\cdot\dd\boldsymbol S=\sum Q/\varepsilon_0$。下面的例题将这些概念放入连续带电体、
多极场、带空腔物体和带电粒子运动中。

% EM teaching enrichment 20260927
\paragraph{散度给局部电荷，边界给整体场}
若某一区域 $\boldsymbol E=\alpha x\boldsymbol e_x+\beta y\boldsymbol e_y$，则
$\rho=\varepsilon_0(\alpha+\beta)$。其中 $\alpha,\beta$ 的单位为 $\mathrm{V/m^2}$。
当 $\beta=-\alpha$ 时，区域无电荷而场仍可非零。这正说明“没有场源”不等于“没有场”，外部电荷和边界仍能建立电场。
静电问题还要满足 $\nabla\times\boldsymbol E=0$，再给定边界条件。只给散度不足以重建矢量场；后面导体问题中的唯一性定理就是把这些信息合在一起。




% EM structure and teaching revision 20260928
\subsection{分段高斯面与表面跳变}
高斯面必须是闭合曲面；曲面上场强是否相等，要由电荷分布的对称性判断。以半径 $R$、体电荷密度 $\rho$ 的均匀实心球为例，半径 $r$ 的同心球面给出
\[
 E(r)=\begin{cases}\rho r/(3\varepsilon_0),&r<R,\\
 \rho R^3/(3\varepsilon_0r^2),&r>R.
 \end{cases}
\]
这里 $E$ 表示沿径向外的有符号分量。球内包围电荷随 $r^3$ 增加，不能把全体电荷代入球内的高斯定律。由于表面没有额外的电荷薄层，$E$ 在 $r=R$ 连续；其导数可以不连续。
若表面另有总面电荷密度 $\sigma_s$（包含自由与束缚电荷），在界面两侧取薄柱形高斯面，令厚度趋于零，侧面的通量消失，得到
\[
 E_{n,\mathrm{out}}-E_{n,\mathrm{in}}=\frac{\sigma_s}{\varepsilon_0}.
\]
若只给定自由面电荷密度 $\sigma_f$，介质界面的条件应写为 $D_{n,\mathrm{out}}-D_{n,\mathrm{in}}=\sigma_f$，不能把 $\sigma_f$ 直接当作上式的总面电荷。切向分量则由静电场的环量为零而连续。这两条边界关系不要求整个问题球对称，因而也能用于导体表面、带电平面和分层介质的局部检查。
不要把“闭合面净通量为零”误读为“面上处处无场”：外部电荷产生的场可以穿入又穿出同一曲面，净贡献恰好抵消。


\subsection{专题例题与方法应用}

\begin{origexamplebox}{【球壳空洞】}\label{ch19:ex:s32-1-3}
一个绝缘球体内有一个球形空洞，不空的地方均匀带有固定的体电荷密度 $\rho$。另比较：若它换成达到静电平衡的金属且空洞内无电荷，会有什么不同？
\begin{enumerate}[label=\arabic*.]
  \item 小明想起来似乎万有引力那一章中也计算过类似模型。您还记得答案吗？
  \item 这两个模型有什么区别？
\end{enumerate}
\end{origexamplebox}\par\nopagebreak[4]

\begin{xesolution}
设完整大球球心为 $O$，空洞球心为 $C$，从 $O$ 指向 $C$ 的位移为 $\boldsymbol a$。空洞内任一点
$P$ 相对 $O,C$ 的位置矢量分别为 $\boldsymbol r,\boldsymbol r'$，于是
$\boldsymbol r=\boldsymbol a+\boldsymbol r'$。

\cavityFigure

把“带空洞的均匀带电体”叠加成：一个体电荷密度为 $+\rho$ 的完整大球，再在空洞位置放一个体电荷密度
为 $-\rho$ 的小球。均匀带电实心球内部的电场由高斯定律得到
\[
  \boldsymbol E_{\rm sphere}=\frac{\rho}{3\varepsilon_0}\boldsymbol r.
\]
因此空洞内
\[
  \boldsymbol E
  =\frac{\rho}{3\varepsilon_0}\boldsymbol r
  -\frac{\rho}{3\varepsilon_0}\boldsymbol r'
  =\boxed{\frac{\rho}{3\varepsilon_0}\boldsymbol a}.
\]
结果与 $P$ 在空洞中的位置无关，空洞内是匀强电场；$\rho>0$ 时方向从大球球心指向空洞球心。

引力模型的叠加步骤完全相同，但引力永远吸引。若质量密度为 $\rho_m$，空洞内引力场为
\[
  \boxed{\boldsymbol g=-\frac{4\pi G\rho_m}{3}\boldsymbol a}.
\]
负号说明方向与电荷密度为正时的电场相反。

“均匀体电荷”与“处于静电平衡的金属”不能混为一谈。若物体是真正的导体，且达到静电平衡，自由电荷
会迁移到表面，金属内部有 $\boldsymbol E=0$，不能维持非零、均匀的体电荷密度。若空洞中没有另放电荷，
空洞内也有 $\boldsymbol E=0$。所以第一问的叠加结果适用于绝缘体中固定的均匀体电荷；金属静电平衡模型则由等势边界条件决定，不能沿用均匀体电荷假设。
\end{xesolution}



\begin{origexamplebox}{【无穷大厚板】}\label{ch19:ex:s32-1-6}
如图，一片区域在两个方向上无限大，另一个方向只有 $-d$ 到 $+d$ 内有均匀正电荷分布，体电荷密度为
$\rho>0$。初始时，质量为 $m$、电荷为 $-q$ 的点电荷从原点以初速度 $v_0$ 向右运动，求运动周期。
\par\noindent\begin{minipage}{\linewidth}
\emSlabQuestion
\figcap{厚板边界及负电荷的初始位置和速度。}\label{ch19:fig:em-question-6}
\end{minipage}\par\medskip
\end{origexamplebox}\par\nopagebreak[4]

\begin{xesolution}
先用关于中面镜面对称的柱形高斯面求场。厚板内、外电场分别为
\[
  \boxed{E_x=\frac{\rho x}{\varepsilon_0}\quad(|x|\le d)},
  \qquad
  \boxed{E_x=\frac{\rho d}{\varepsilon_0}\operatorname{sgn}x\quad(|x|>d)}.
\]
图(b)中斜线必须穿过原点，并在 $x=\pm d$ 与两侧水平线连续；这三个位置关系正是后面分段运动的依据。

\chargedSlabFigure

在板内，负电荷受到指向中面的回复力
\[
  F_x=-qE_x=-\frac{q\rho}{\varepsilon_0}x=-m\omega^2x,
  \qquad
  \boxed{\omega=\sqrt{\frac{q\rho}{m\varepsilon_0}}}.
\]
若简谐振幅 $A=v_0/\omega$ 不超过 $d$，粒子始终留在板内，周期直接为
\[
  \boxed{T=\frac{2\pi}{\omega}}
  \qquad (v_0\le\omega d).
\]

若 $v_0>\omega d$，粒子会穿出厚板。在 $0\le x\le d$ 内由
$x=(v_0/\omega)\sin\omega t$ 得到从中面运动到边界的时间
\[
  t_1=\frac1\omega\arcsin\frac{\omega d}{v_0},
\]
到达边界的速度为
\[
  v_d=\sqrt{v_0^2-\omega^2d^2}.
\]
板外回复加速度大小恒为
\[
  a_{\rm out}=\frac{q\rho d}{m\varepsilon_0}=\omega^2d,
\]
故从边界到右侧转折点的时间为
\[
  t_2=\frac{v_d}{\omega^2d}.
\]
从中面向右到右侧转折点是整个周期的四分之一，所以
\[
  \boxed{T=4\left[
  \frac1\omega\arcsin\frac{\omega d}{v_0}
  +\frac{\sqrt{v_0^2-\omega^2d^2}}{\omega^2d}
  \right]}
  \qquad(v_0>\omega d).
\]
在临界速度 $v_0=\omega d$ 处，第二项为零、反正弦为 $\pi/2$，两段答案连续接成
$2\pi/\omega$。
\end{xesolution}





\graphicspath{{/Users/dongxingxi/Documents/物理竞赛一轮资料/sample/xes-original/thirty-third-section/}}
\section{电势与场势关系}

% Text annotations are transparent; no label carries an opaque backing box.
\tikzset{
  xeslabel/.style={font=\small,text=Ink,fill=none,inner sep=.6pt},
  xessmall/.style={font=\footnotesize,text=TextGray,fill=none,inner sep=.45pt},
  xestiny/.style={font=\scriptsize,text=TextGray,fill=none,inner sep=.35pt},
  xesline/.style={draw=Ink,line width=.85pt},
  xesblue/.style={draw=DeepBlue,line width=1pt},
  xesarrow/.style={-{Latex[length=2.15mm,width=1.35mm]},draw=DeepBlue,line width=.9pt},
  fieldline/.style={draw=DeepBlue,line width=.82pt,postaction={decorate},
    decoration={markings,mark=at position .62 with {\arrow{Latex[length=1.65mm,width=1.1mm]}}}},
  eqline/.style={draw=RuleGray,line width=.66pt,densely dashed}
}


\newcommand{\equipotentialGallery}{%
\begin{center}
\includegraphics[width=.95\linewidth]{/Users/dongxingxi/Documents/物理竞赛一轮资料/sample/xes-original/restructured-complete/assets/revision-ch19-ch20/equipotentials.pdf}
\end{center}}


\newcommand{\dipolePotentialGeometry}{%
\begin{center}
\begin{tikzpicture}[x=1cm,y=1cm]
  \coordinate (C) at (0,0);\coordinate (P) at (5.3,2.35);
  \coordinate (A) at (0,1.0);\coordinate (B) at (0,-1.0);
  \fill[DeepBlue] (A) circle (5pt);\node[xeslabel,anchor=east] at (-.16,1) {$+q$};
  \fill[Ink] (B) circle (5pt);\node[xeslabel,anchor=east] at (-.16,-1) {$-q$};
  \fill[Ink] (C) circle (1.3pt);\node[xessmall,anchor=east] at (-.12,0) {$C$};
  \fill[Ink] (P) circle (1.5pt);\node[xeslabel,anchor=west] at (5.42,2.35) {$P$};
  \draw[DeepBlue,line width=.75pt] (A)--(P) node[xessmall,pos=.53,above] {$r_+$};
  \draw[DeepBlue,line width=.75pt] (B)--(P) node[xessmall,pos=.53,below] {$r_-$};
  \draw[RuleGray,densely dashed] (C)--(P) node[xessmall,pos=.52,below] {$r$};
  \draw[<->,RuleGray,line width=.55pt] (.35,0)--(.35,1);\node[xessmall,anchor=west] at (.42,.5) {$a$};
  \draw[<->,RuleGray,line width=.55pt] (.35,-1)--(.35,0);\node[xessmall,anchor=west] at (.42,-.5) {$a$};
  \draw[DeepBlue,line width=.6pt] (0,1.15) arc[start angle=90,end angle=23.9,radius=1.15];
  \node[xessmall] at (.72,.75) {$\theta$};
  \node[xessmall,align=left,anchor=west] at (7.1,.6) {$r\gg a$ 时\\$r_- -r_+\approx2a\cos\theta$};
  \draw[RuleGray,-{Latex[length=1.5mm]},line width=.55pt] (7.0,.55)--(4.35,1.92);
\end{tikzpicture}
\end{center}}

\newcommand{\zeroEquipotentialFigure}{%
\begin{center}
\includegraphics[width=.76\linewidth]{/Users/dongxingxi/Documents/物理竞赛一轮资料/sample/xes-original/restructured-complete/assets/revision-ch19-ch20/zero-potential.pdf}
\end{center}}


\newcommand{\ionizationEnergyFigure}{%
\begin{center}
\includegraphics[width=.95\linewidth]{/Users/dongxingxi/Documents/物理竞赛一轮资料/sample/xes-original/restructured-complete/assets/revision-ch19-ch20/ionization.pdf}
\end{center}}


\newcommand{\alphaDecayFigure}{%
\begin{center}
\begin{tikzpicture}[x=1cm,y=1cm]
  \begin{scope}[shift={(1.8,0)},scale=.9]
    \foreach \x/\y/\c in {-1/.1/DeepBlue,-.65/.65/Gold,-.25/.2/Gold,.15/.72/DeepBlue,.55/.18/Gold,.9/.62/DeepBlue,-.8/-.55/Gold,-.3/-.55/DeepBlue,.2/-.35/Gold,.7/-.48/DeepBlue,0/0/DeepBlue}{\shade[ball color=\c!65] (\x,\y) circle (.28);}
  \end{scope}
  \draw[DeepBlue,line width=.65pt,dashed] (3.0,.28)--(5.35,1.15);
  \draw[DeepBlue,line width=.65pt,dashed] (3.0,-.02)--(5.35,.82);
  \begin{scope}[shift={(5.75,1.0)}]
    \shade[ball color=Gold!70] (-.2,.18) circle (.27);\shade[ball color=DeepBlue!70] (.2,.18) circle (.27);
    \shade[ball color=DeepBlue!70] (-.2,-.18) circle (.27);\shade[ball color=Gold!70] (.2,-.18) circle (.27);
  \end{scope}
  \draw[xesarrow] (5.1,.98)--(6.7,1.58);
  \node[xessmall,anchor=west] at (6.85,1.62) {$\alpha$ 粒子（$^4\mathrm{He}$ 核）};
  \node[xessmall,align=left,anchor=west] at (6.85,.35) {母核质量数减 4，\\核电荷数减 2。};
\end{tikzpicture}
\end{center}}

\newcommand{\acceleratorTunnelFigure}{%
\begin{center}
\includegraphics[width=.76\linewidth]{/Users/dongxingxi/Documents/物理竞赛一轮资料/sample/xes-original/restructured-complete/assets/revision-ch19-ch20/accelerator.png}
\end{center}}


\newcommand{\lightningFigure}{%
\begin{center}
\includegraphics[width=.76\linewidth]{/Users/dongxingxi/Documents/物理竞赛一轮资料/sample/xes-original/restructured-complete/assets/revision-ch19-ch20/lightning.png}
\end{center}}


\newcommand{\potentialMethodFigure}{%
\begin{center}
\begin{tikzpicture}[x=1cm,y=1cm]
  \begin{scope}
    \draw[xesline,-{Latex[length=1.8mm]}] (0,0)--(3.3,0) node[xessmall,below left] {$s$};
    \draw[xesline,-{Latex[length=1.8mm]}] (0,0)--(0,2.6) node[xessmall,below right] {$\varphi$};
    \draw[xesblue] (.35,2.35)..controls(.8,1.65)and(1.65,.75)..(3.05,.25);
    \node[xestiny] at (1.8,-.42) {有限带电体：可取 $\varphi(\infty)=0$};
  \end{scope}
  \begin{scope}[shift={(4.4,0)}]
    \draw[xesline,-{Latex[length=1.8mm]}] (0,0)--(3.3,0) node[xessmall,below left] {$s$};
    \draw[xesline,-{Latex[length=1.8mm]}] (0,0)--(0,2.6) node[xessmall,below right] {$\varphi$};
    \draw[xesblue] (.28,2.35)..controls(.85,1.65)and(1.7,1.02)..(3.1,.48);
    \draw[RuleGray,densely dashed] (.55,0)--(.55,2.1) (2.55,0)--(2.55,.62);
    \node[xestiny] at (1.8,-.42) {无限直线：只定义电势差};
  \end{scope}
  \begin{scope}[shift={(8.8,0)}]
    \draw[xesline,-{Latex[length=1.8mm]}] (0,0)--(3.3,0) node[xessmall,below left] {$z$};
    \draw[xesline,-{Latex[length=1.8mm]}] (0,0)--(0,2.6) node[xessmall,below right] {$\varphi$};
    \draw[xesblue] (0,2.25)--(3.05,.25);
    \node[xestiny] at (1.8,-.42) {无限平面：$\varphi$ 随距离线性变化};
  \end{scope}
\end{tikzpicture}
\end{center}}

\newcommand{\chargedSphereEnergyFigure}{%
\begin{center}
\begin{tikzpicture}[x=1cm,y=1cm]
  \begin{scope}[shift={(1.5,1.45)}]
    \draw[Ink,line width=1pt] (0,0) circle (1.05);
    \foreach \a in {0,30,...,330}{\node[xestiny,text=DeepBlue] at (\a:.93) {$+$};}
    \node[xessmall] at (0,-1.48) {导体球壳：电荷在表面};
  \end{scope}
  \begin{scope}[shift={(5.0,1.45)}]
    \fill[DeepBlue!7] (0,0) circle (1.05);\draw[Ink,line width=1pt] (0,0) circle (1.05);
    \foreach \x/\y in {-0.6/0,-.3/.5,0/.15,.42/.55,.65/.05,-.5/-.48,.05/-.55,.5/-.4}{\node[xestiny,text=DeepBlue] at (\x,\y) {$+$};}
    \node[xessmall] at (0,-1.48) {绝缘实心球：电荷在体内};
  \end{scope}
  \begin{scope}[shift={(7.5,.05)},xscale=.9,yscale=.8]
    \draw[xesline,-{Latex[length=1.8mm]}] (0,0)--(4.2,0) node[xessmall,below left] {$r$};
    \draw[xesline,-{Latex[length=1.8mm]}] (0,0)--(0,3.55) node[xessmall,below right] {$E$};
    \draw[RuleGray,densely dashed] (1.55,0)--(1.55,3.15);\node[xessmall,below] at (1.55,0) {$R$};
    \draw[DeepBlue,line width=1pt] (0,0)--(1.55,2.75);
    \draw[DeepBlue,line width=1pt] (1.55,2.75)..controls(2.15,1.5)and(3.05,.65)..(4,.35);
    \draw[Ink,line width=.85pt] (0,0)--(1.55,0);
    \draw[Ink,line width=.85pt] (1.55,2.75)..controls(2.15,1.5)and(3.05,.65)..(4,.35);
    \node[xestiny,anchor=west] at (2.5,2.45) {蓝：实心球};
    \node[xestiny,anchor=west] at (2.5,2.08) {黑：球壳};
  \end{scope}
\end{tikzpicture}
\end{center}}

\newcommand{\chargeAssemblyFigure}{%
\begin{center}
\begin{tikzpicture}[x=1cm,y=1cm]
  \foreach \x/\y/\i in {0/0/1,1.7/.5/2,3.2/-.25/3,4.5/1.0/4,5.8/-.1/5,.7/2.0/6,2.4/2.2/7,4.0/2.7/8,5.7/2.1/9,7.1/1.0/10}{
    \fill[DeepBlue!75] (\x,\y) circle (4pt);\node[xestiny,above=4pt] at (\x,\y) {$Q_{\i}$};}
  \draw[RuleGray,densely dashed] (0,0)--(3.2,-.25) node[xessmall,pos=.52,below] {$r_{13}$};
  \draw[RuleGray,densely dashed] (1.7,.5)--(5.7,2.1) node[xessmall,pos=.55,above] {$r_{29}$};
  \draw[xesarrow] (8.3,1.2)--(7.25,1.05);\node[xessmall,anchor=west,align=left] at (8.45,1.2) {最后移入的电荷只与\\已经就位的电荷计一次相互作用};
\end{tikzpicture}
\end{center}}

\newcommand{\apolloniusSphereFigure}{%
\begin{center}
\includegraphics[width=.76\linewidth]{/Users/dongxingxi/Documents/物理竞赛一轮资料/sample/xes-original/restructured-complete/assets/revision-ch19-ch20/apollonius-example.pdf}
\end{center}}


\newcommand{\twoLinesFamiliesFigure}{%
\begin{center}
\includegraphics[width=.76\linewidth]{/Users/dongxingxi/Documents/物理竞赛一轮资料/sample/xes-original/restructured-complete/assets/revision-ch19-ch20/two-lines.pdf}
\end{center}}


\newcommand{\finiteRodFamiliesFigure}{%
\begin{center}
\includegraphics[width=.76\linewidth]{/Users/dongxingxi/Documents/物理竞赛一轮资料/sample/xes-original/restructured-complete/assets/revision-ch19-ch20/finite-rod.pdf}
\end{center}}


\newcommand{\repellingSpheresFigure}{%
\begin{center}
\begin{tikzpicture}[x=1cm,y=1cm]
  \coordinate (S) at (0,3.5);\coordinate (B) at (0,0);\coordinate (A) at ({-3.5*sin(40)},{3.5*(1-cos(40))});
  \draw[Ink,line width=.9pt] (-.7,3.5)--(.7,3.5);\foreach \x in {-.6,-.3,0,.3,.6}{\draw[RuleGray,line width=.5pt] (\x,3.5)--++(.18,.2);}
  \draw[Ink,line width=1pt] (S)--(A);
  \fill[DeepBlue] (A) circle (5pt);\node[xeslabel,anchor=east] at ($(A)+(-.14,0)$) {$A$};
  \fill[DeepBlue] (B) circle (5pt);\node[xeslabel,anchor=west] at ($(B)+(.14,0)$) {$B$};
  \draw[RuleGray,densely dashed] (S)--(B);
  \draw[RuleGray] (0,2.7) arc[start angle=270,end angle=230,radius=.8];
  \node[xessmall] at (-.4,2.47) {$\theta$};
  \node[xessmall,anchor=south east] at (-1.35,2.1) {$L$};
  \draw[<->,RuleGray,line width=.55pt] (.65,0)--(.65,{3.5*(1-cos(40))});\node[xessmall,anchor=west] at (.73,.4) {$h$};
  \draw[RuleGray,densely dashed] (A)--(B) node[xessmall,pos=.54,above] {$r$};
  \draw[xesarrow] (A)--++(0,-1.15);\node[xessmall,anchor=north] at ($(A)+(0,-1.18)$) {$m\boldsymbol g$};
  \draw[xesarrow] (A)--($(A)!.38!(S)$);\node[xessmall,anchor=south east] at (-1.5,1.75) {$\boldsymbol T$};
  \draw[xesarrow] (A)--($(A)!-.48!(B)$);\node[xessmall,anchor=south east] at (-3.15,1.15) {$\boldsymbol F_e$};
  \node[xessmall,align=left,anchor=west] at (2.1,1.7) {$r=2L\sin(\theta/2)$\\$h=2L\sin^2(\theta/2)$\\因此 $h=r\sin(\theta/2)$};
\end{tikzpicture}
\end{center}}

\newcommand{\bohrAtomFigure}{%
\begin{center}
\begin{tikzpicture}[x=1cm,y=1cm]
  \fill[Gold!80!Ink] (0,0) circle (7pt);\node[xestiny,text=white] at (0,0) {$+$};
  \draw[DeepBlue,line width=1pt] (0,0) circle (2.1);
  \fill[Ink] (2.1,0) circle (5pt);\node[xeslabel,anchor=west] at (2.27,0) {$e^-$};
  \draw[xesarrow] (2.1,0)--(2.1,1.15);\node[xessmall,anchor=west] at (2.25,.75) {$\boldsymbol v$};
  \draw[xesarrow] (2.0,0)--(.85,0);\node[xessmall,above] at (1.4,0) {$\boldsymbol F_C$};
  \draw[<->,RuleGray,line width=.55pt] (0,-.42)--(2.1,-.42);\node[xessmall,below] at (1.05,-.42) {$r$};
  \node[xessmall,align=left,anchor=west] at (4.45,.65) {圆周运动条件：$mv^2/r=ke^2/r^2$\\总能量：$E=\frac12mv^2-ke^2/r$};
\end{tikzpicture}
\end{center}}

\newcommand{\rutherfordFigure}{%
\begin{center}
\includegraphics[width=.76\linewidth]{/Users/dongxingxi/Documents/物理竞赛一轮资料/sample/xes-original/restructured-complete/assets/revision-ch19-ch20/rutherford.pdf}
\end{center}}


% Source question conditions restored, analysis retained separately.
\newcommand{\emRepellingQuestion}{%
\begin{center}
\begin{tikzpicture}[x=1cm,y=1cm]
  \coordinate (S) at (0,3.5);\coordinate (B) at (0,0);\coordinate (A) at ({-3.5*sin(40)},{3.5*(1-cos(40))});
  \draw[Ink,line width=.9pt] (-.7,3.5)--(.7,3.5);\foreach \x in {-.6,-.3,0,.3,.6}{\draw[RuleGray,line width=.5pt] (\x,3.5)--++(.18,.2);}
  \draw[Ink,line width=1pt] (S)--(A);
  \fill[DeepBlue] (A) circle (5pt);\node[xeslabel,anchor=east] at ($(A)+(-.14,0)$) {$A$};
  \fill[DeepBlue] (B) circle (5pt);\node[xeslabel,anchor=west] at ($(B)+(.14,0)$) {$B$};
  \draw[RuleGray,densely dashed] (S)--(B);
  \draw[RuleGray] (0,2.7) arc[start angle=270,end angle=230,radius=.8];
  \node[xessmall] at (-.4,2.47) {$\theta$};
  \node[xessmall,anchor=south east] at (-1.35,2.1) {$L$};
  \draw[<->,RuleGray,line width=.55pt] (.65,0)--(.65,{3.5*(1-cos(40))});\node[xessmall,anchor=west] at (.73,.4) {$h$};
  \draw[RuleGray,densely dashed] (A)--(B) node[xessmall,pos=.54,above] {$r$};
\end{tikzpicture}
\end{center}}


本章前文从带电体受力的角度认识电场。物理中许多问题都可以从“力”和“能量交换”两个方面考察，电场也不
例外：除了用各点的电场强度描述电场，还可以研究给定电荷在各点具有的能量。后一种描述引出电势能和
电势。

把电荷量为 $q$ 的试探电荷放进一个既定电场，并假定它不足以改变场源电荷的位置，试探电荷受到的力为
$\boldsymbol F=q\boldsymbol E$。当它的位置改变时，电场力会做功。静电场中的电场力是保守力，因此可以
定义电势能。由于电场力处处正比于 $q$，电势能也必定正比于 $q$：
\begin{equation}
  \boxed{E_p=q\varphi}.
\end{equation}
比例因子 $\varphi$ 只由场源和空间位置决定，称为该点的\textbf{电势}。它等于单位正电荷在该点具有的
电势能。电场强度是矢量，电势是标量；二者描述的是同一个静电场。


\subsection{电势与电场}


在直角坐标系中，电势写成标量函数 $\varphi(x,y,z)$，试探电荷的势能为
$E_p(x,y,z)=q\varphi(x,y,z)$。保守力分量等于势能对相应坐标导数的负值：
\[
  F_x=-\frac{\partial E_p}{\partial x},\qquad
  F_y=-\frac{\partial E_p}{\partial y},\qquad
  F_z=-\frac{\partial E_p}{\partial z}.
\]
再与 $F_x=qE_x,F_y=qE_y,F_z=qE_z$ 比较，得到
\begin{equation}
  \boxed{E_x=-\frac{\partial\varphi}{\partial x},\qquad
  E_y=-\frac{\partial\varphi}{\partial y},\qquad
  E_z=-\frac{\partial\varphi}{\partial z}}
  \qquad\Longleftrightarrow\qquad
  \boxed{\boldsymbol E=-\boldsymbol\nabla\varphi}.
\end{equation}
负号表示电场总是指向电势降低最快的方向。知道电势函数后，对三个坐标求偏导即可得到电场。

反过来，若电场 $\boldsymbol E(x,y,z)$ 已知，电荷从 $A$ 沿任意路径运动到 $B$，电场力做功
\[
  W_{A\to B}=q\int_A^B\boldsymbol E\cdot\dd\boldsymbol l.
\]
保守力做功等于势能减少量，因此
\begin{equation}
  \boxed{\varphi(B)-\varphi(A)
  =-\int_A^B\boldsymbol E\cdot\dd\boldsymbol l}.
\end{equation}
静电力是保守力，积分结果只取决于起点和终点，与路径无关。沿电场方向移动时
$\boldsymbol E\cdot\dd\boldsymbol l>0$，电势降低；逆电场方向移动时电势升高。


% EM teaching enrichment 20260927
\paragraph{势的零点与平衡点}
给电势加常数不改变电场，也不改变电势差。若沿一维约束轨道有
$\varphi(x)=\varphi_0+ax^2/2$，则 $E_x=-ax$，电荷运动方程为 $m\ddot x=-qa x$。
$qa>0$ 时小振动角频率为 $\sqrt{qa/m}$；$qa<0$ 时原点不稳定。因此稳定性取决于 $q\varphi$ 的曲率，负电荷在电势极大处反而可能具有势能极小值。本结论限于给定的一维约束，不能直接推出三维自由空间中的稳定悬浮。


\subsection{叠加原理}


把无穷远选为电势零点，静止点电荷 $Q$ 在距它 $r$ 处产生的电势为
\begin{equation}
  \boxed{\varphi(r)=\frac{1}{4\pi\varepsilon_0}\frac{Q}{r}}.
\end{equation}
在少数延伸到无穷远的带电系统中，这一零点选择并不可行，只能谈两个有限位置之间的电势差。对于有限的
孤立带电体，通常都取 $\varphi(\infty)=0$。

多个点电荷的电场可以矢量叠加，它们的电势则可以直接作代数和。空间中一点到各电荷 $Q_n$ 的距离为
$r_n$ 时
\begin{equation}
  \boxed{\varphi=\frac{1}{4\pi\varepsilon_0}
  \sum_n\frac{Q_n}{r_n}}.
\end{equation}
若电荷连续分布，把求和改为积分即可。因为标量叠加通常比矢量叠加简单，复杂带电体常采用“先求电势，
再用 $\boldsymbol E=-\boldsymbol\nabla\varphi$ 求电场”的路线。

图中的电偶极子由 $+q,-q$ 构成，两电荷到中心 $C$ 的距离均为 $a$。任意点 $P$ 到两电荷的距离分别为
$r_+,r_-$，到中心的距离为 $r$，$\boldsymbol r$ 与偶极矩方向的夹角为 $\theta$。由叠加原理
\[
  \varphi(P)=\frac{q}{4\pi\varepsilon_0}
  \left(\frac1{r_+}-\frac1{r_-}\right).
\]
当 $r\gg a$ 时，$r_- -r_+\approx2a\cos\theta$。定义偶极矩
$\boldsymbol p$ 由负电荷指向正电荷，大小 $p=2qa$，则
\begin{equation}
  \boxed{\varphi(P)=\frac{1}{4\pi\varepsilon_0}
  \frac{p\cos\theta}{r^2}
  =\frac{1}{4\pi\varepsilon_0}\frac{\boldsymbol p\cdot\boldsymbol r}{r^3}}.
\end{equation}

\par\noindent\begin{minipage}{\linewidth}
\dipolePotentialGeometry
\diagramcaption{电偶极子远场近似中的距离和角度关系。$\boldsymbol p$ 沿 $-q$ 指向 $+q$。}
\end{minipage}\par\medskip

另一个重要例子是半径为 $R$、总电荷为 $Q$ 的均匀带电球壳。它与均匀质量球壳的引力势问题具有完全相同
的平方反比结构。球壳外部的电势等同于把全部电荷集中在球心，内部电势处处等于表面电势：
\begin{equation}
  \boxed{\varphi(r)=
  \begin{cases}
    \dfrac{1}{4\pi\varepsilon_0}\dfrac{Q}{r},&r>R,\\[6pt]
    \dfrac{1}{4\pi\varepsilon_0}\dfrac{Q}{R},&r<R.
  \end{cases}}
\end{equation}
球壳内部电势虽不一定为零，但其空间导数为零，因此内部电场为零。


% EM teaching enrichment 20260927
\paragraph{先算势再求场：求导时保持什么不变}
对均匀细圆环，轴线上电势为 $\varphi(z)=kQ/\sqrt{a^2+z^2}$，求导立即得到
$E_z=-\dd\varphi/\dd z=kQz/(a^2+z^2)^{3/2}$，其中 $Q,a$ 保持不变。
但只知道轴线上的电势函数，并不能直接对横坐标求导得到轴外电场；轴外场必须从完整空间函数或额外对称关系获得。
无限长线电荷则给出
\[
 \varphi(r)-\varphi(r_0)=-\frac{\lambda}{2\pi\varepsilon_0}\ln\frac r{r_0}.
\]
令 $r_0\to\infty$ 会发散，所以应保留有限参考半径。场强是否随距离趋零，与无穷远电势积分是否收敛，是两件不同的事。


\subsection{等势面}


把空间中电势相同的点连接起来，得到的曲面称为等势面。它与地形图中的等高线相似，使标量函数
$\varphi(x,y,z)$ 具有直观的几何表示。图~\ref{ch19:fig:s33-2-1}给出匀强电场、正点电荷和电偶极子的等势面截线；实线带箭头
表示电场线，灰色虚线表示等势面在纸面内的截线。

\par\noindent\begin{minipage}{\linewidth}
\equipotentialGallery
\diagramcaption{三种典型静电场中的等势面与电场线。两族曲线在每个交点处互相垂直。}\label{ch19:fig:s33-2-1}
\end{minipage}\par\medskip

等势面和电场线具有三条直接性质。第一，不同电势的等势面不能相交，电场线也不能在电场非零处相交；
否则同一点会出现两个电势值或两个电场方向。第二，电场线处处垂直于等势面，因为电荷沿等势面移动时
$\dd\varphi=0$，电场力不做功。第三，若用相同电势差画相邻等势面，等势面越密的地方场强越大，因为
$|\boldsymbol E|$ 是电势沿法线方向的空间变化率。

两个不等量异号点电荷 $+Q,-q$ 产生的电势为
\[
  \varphi(P)=\frac{1}{4\pi\varepsilon_0}
  \left(\frac{Q}{r_+}-\frac{q}{r_-}\right).
\]
零电势面满足
\[
  \frac{r_+}{r_-}=\frac{Q}{q}.
\]
到两个定点的距离之比为常数的点构成 Apollonius 球面。因此当 $Q\ne q$ 时，零电势面是一个包围电荷
绝对值较小者的球面；当 $Q=q$ 时，它退化为两电荷连线的中垂面。

图~\ref{ch19:fig:s33-2-2} 取 $Q=2q$；圆周上 $P$ 到正、负电荷的距离分别为 $r_+,r_-$，满足 $r_+=2r_-$，故两项电势正好相消。

\par\noindent\begin{minipage}{\linewidth}
\zeroEquipotentialFigure
\diagramcaption{不等量异号点电荷的零电势面截面。球面包围绝对值较小的电荷。}\label{ch19:fig:s33-2-2}
\end{minipage}\par\medskip


% EM teaching enrichment 20260927
\paragraph{等势线读数不能代替粒子轨迹}
沿局部法线，相邻等势面差为 $\Delta\varphi$、间距为 $\Delta n$ 时，
$E_n\simeq-\Delta\varphi/\Delta n$。必须比较相同电势差的线距；随意多画几条线不会使真实场变强。
例如两相邻等势面差 $20\,\mathrm V$、法向间距 $2.0\,\mathrm{mm}$，局部场强约为 $10^4\,\mathrm{V/m}$。
电荷沿等势面移动时电场不做功，但可存在垂直于速度的电场力；若没有合适约束，粒子通常不会一直留在该面上。


% EM structure and teaching revision 20260928
\subsection{由场求势：分段积分与参考点}
电势差由路径积分定义。若场具有径向对称性，可写为
\[
 \varphi(r)-\varphi(r_0)=-\int_{r_0}^{r}E_r(s)\dd s.
\]
积分路径穿过不同场区时必须分段，并让同一界面的电势值衔接。半径 $R$、电荷 $Q$ 的均匀实心球取无穷远为零，球外 $\varphi=kQ/r$，球内则为
\[
 \varphi(r)=\varphi(R)+\int_r^R\frac{kQs}{R^3}\dd s
 =\frac{kQ}{2R}\left(3-\frac{r^2}{R^2}\right),\qquad k=\frac1{4\pi\varepsilon_0}.
\]
中心电势有限；当 $Q>0$ 时它大于球面电势，而中心场强为零。若 $Q<0$，电势高低关系反向，但中心场强仍为零。这说明电势的数值与其空间变化率是不同问题；“场为零”只能推出所在连通无场区域内电势为常数。
无限长线电荷的场按 $1/r$ 衰减，到无穷远的积分不收敛，不能仍机械地选 $\varphi(\infty)=0$。应指定有限参考半径 $r_0$，得到
$\varphi(r)-\varphi(r_0)=-\lambda\ln(r/r_0)/(2\pi\varepsilon_0)$。绝对零点的选择不影响电势差、力和能量变化。


\subsection{专题例题与方法应用}

\begin{origexamplebox}{【电偶极子】}\label{ch19:ex:s32-1-4}
相距 $l$ 的两个点电荷 $-q$、$+q$ 构成电偶极子。定义电偶极矩
$\boldsymbol p=q\boldsymbol l$，方向由负电荷指向正电荷。求远处 $r\gg l$ 的电场，并研究其电场线。
\begin{center}\begin{tikzpicture}[>=Latex,scale=.8]

 \fill ( -1,0) circle(3pt) node[below]{$-q$};\fill[DeepBlue](1,0)circle(3pt)node[below]{$+q$};
 \draw[<->](-1,-.65)--(1,-.65)node[midway,below]{$l$};
 \draw[->,DeepBlue](0,0)--(2.1,0)node[right]{$\boldsymbol p$};
 \fill(0,0)circle(1pt)node[below]{$O$};\draw[dashed](0,0)--(4.0,2.2)node[midway,above left]{$r$};
 \fill(4.0,2.2)circle(1.5pt)node[right]{$P$};
 \draw(.9,0)arc[start angle=0,end angle=28.81,radius=.9];\node at(1.18,.3){$\theta$};

\end{tikzpicture}\end{center}
\figcap{电偶极子的电荷间距与远处观察方向，电场待求。}\label{ch19:fig:em-source-thirty-second-section-4}
\end{origexamplebox}\par\nopagebreak[4]

\begin{xesolution}
把坐标原点取在两电荷中点，极轴沿 $\boldsymbol p$。观察点的位置矢量为 $\boldsymbol r$，它与
$\boldsymbol p$ 的夹角为 $\theta$。在 $l/r$ 的一阶精度下，电势为
\[
  \varphi(\boldsymbol r)
  =\frac{1}{4\pi\varepsilon_0}\frac{\boldsymbol p\cdot\hat{\boldsymbol r}}{r^2}
  =\frac{1}{4\pi\varepsilon_0}\frac{p\cos\theta}{r^2}.
\]
由 $\boldsymbol E=-\boldsymbol\nabla\varphi$ 得到球坐标分量
\[
  E_r=-\frac{\partial\varphi}{\partial r}
  =\frac{1}{4\pi\varepsilon_0}\frac{2p\cos\theta}{r^3},
  \qquad
  E_\theta=-\frac1r\frac{\partial\varphi}{\partial\theta}
  =\frac{1}{4\pi\varepsilon_0}\frac{p\sin\theta}{r^3}.
\]
所以矢量形式为
\[
  \boxed{\boldsymbol E
  =\frac{1}{4\pi\varepsilon_0r^3}
  \left[3(\boldsymbol p\cdot\hat{\boldsymbol r})\hat{\boldsymbol r}-\boldsymbol p\right]}.
\]
场强大小和相对径向的偏角 $\beta$ 分别满足
\[
  \boxed{E=\frac{p}{4\pi\varepsilon_0r^3}\sqrt{1+3\cos^2\theta}},
  \qquad
  \tan\beta=\frac{E_\theta}{E_r}=\frac12\tan\theta.
\]

图~\ref{ch19:fig:s32-1-7} 左侧把场分解为径向和极角方向，右侧以同一偶极矩方向展示远场线族。

\par\noindent\begin{minipage}{\linewidth}
\dipoleFarFigure
\diagramcaption{电偶极子的远场分量与典型场线。$\boldsymbol p$ 的方向由负电荷指向正电荷。}\label{ch19:fig:s32-1-7}
\end{minipage}\par\medskip

场线的切线与电场平行，故在极坐标中
\[
  \frac{\dd r}{r\dd\theta}=\frac{E_r}{E_\theta}=2\cot\theta.
\]
积分得到
\[
  \ln r=2\ln\sin\theta+\text{常量},
  \qquad
  \boxed{r=C\sin^2\theta}.
\]
图右侧按这个方程画出理想点偶极子的场线，偶极矩集中于原点；对左侧有限间距的真实电荷，这个近似只用于 $r\gg l$，不延伸至两电荷附近。远场按 $r^{-3}$ 衰减，比有净电荷的
$r^{-2}$ 更快。
\end{xesolution}
\begin{origexamplebox}{【直线电四极子】}\label{ch19:ex:s32-1-5}
如图，三个电荷依次为 $-q,+2q,-q$，相邻两电荷相距 $l$（两端相距 $2l$）。求距离该电四极子 $r\gg l$ 处的电场强度大小
和方向，并求在这一近似下的电场线方程。
\par\noindent\begin{minipage}{\linewidth}
\emQuadrupoleQuestion
\figcap{三个点电荷、相邻间距及远处观察方向。}\label{ch19:fig:em-question-5}
\end{minipage}\par\medskip
\end{origexamplebox}\par\nopagebreak[4]

\begin{xesolution}
总电荷为零，偶极矩也为零，所以远场的最低非零项是电四极项。取极轴沿电荷排列方向，观察方向与极轴
夹角为 $\theta$。把两端电荷的距离倒数按 $l/r$ 展开到二阶：
\[
  \frac1{|\boldsymbol r\mp l\boldsymbol e_x|}
  =\frac1r\left[1\pm\frac lr\cos\theta
  +\frac{l^2}{2r^2}(3\cos^2\theta-1)+\cdots\right].
\]
与中心 $+2q$ 的电势相加后，零阶和一阶项都抵消，得到
\[
  \boxed{\varphi(r,\theta)=\frac{1}{4\pi\varepsilon_0}
  \frac{ql^2}{r^3}(1-3\cos^2\theta)}.
\]
于是
\[
  \boxed{E_r=\frac{3ql^2}{4\pi\varepsilon_0r^4}(1-3\cos^2\theta)},
  \qquad
  \boxed{E_\theta=-\frac{6ql^2}{4\pi\varepsilon_0r^4}\sin\theta\cos\theta}.
\]
场强大小为
\[
  \boxed{E=\frac{3ql^2}{4\pi\varepsilon_0r^4}
  \sqrt{1-2\cos^2\theta+5\cos^4\theta}}.
\]

图~\ref{ch19:fig:s32-1-8} 左侧的三个电荷共线，相邻间距为 $l$；观察角 $\theta$ 从这条轴量起，右侧场线由三个电荷共同产生。

\par\noindent\begin{minipage}{\linewidth}
\quadrupoleFarFigure
\diagramcaption{直线电四极子的电荷位置、远场观察方向与场线四瓣结构。}\label{ch19:fig:s32-1-8}
\end{minipage}\par\medskip

场线方程满足
\[
  \frac{\dd r}{r\dd\theta}
  =\frac{E_r}{E_\theta}
  =-\frac{1-3\cos^2\theta}{2\sin\theta\cos\theta}.
\]
令 $u=\cos\theta$，积分得
\[
  \ln r=\frac12\ln|\cos\theta|+\ln|\sin\theta|+\text{常量},
\]
所以每一瓣中的场线可写为
\[
  \boxed{r=C|\sin\theta|\sqrt{|\cos\theta|}}.
\]
它具有关于轴线和中垂面的对称性。四极场按 $r^{-4}$ 衰减，比偶极场的 $r^{-3}$ 又快一阶。
\end{xesolution}
\begin{origexamplebox}{【电势算电场】}\label{ch19:ex:s33-2-1}
小明发现电势和电场有非常明确的关系，并且电势是标量，应该更方便计算。请利用先求电势再求电场的方式，
求解本章前文中诸多问题的电场。您一定会发现有的问题这样求解似乎不太正确，思考为什么以及对应的解决办法。
\end{origexamplebox}\par\nopagebreak[4]

\begin{xesolution}
基本路线是先由
\[
  \varphi(\boldsymbol r)=\frac{1}{4\pi\varepsilon_0}
  \int\frac{\dd q}{|\boldsymbol r-\boldsymbol r'|}
\]
求标量电势，再作 $\boldsymbol E=-\boldsymbol\nabla\varphi$。困难不在求导，而在于：对无限延伸的带电体，
把无穷远取作零电势会使电势积分发散。

\potentialMethodFigure

\textbf{1. 点电荷。} 对 $\varphi=kQ/r$ 直接求梯度：
\[
  \boldsymbol E=-\frac{\dd\varphi}{\dd r}\hat{\boldsymbol r}
  =\boxed{\frac{1}{4\pi\varepsilon_0}\frac{Q}{r^2}\hat{\boldsymbol r}}.
\]

\textbf{2. 有限长直线。} 令导线沿 $y$ 轴，从 $-L_2$ 延伸到 $L_1$，线密度为 $\lambda$；观察点在
$(a,0)$。电势为
\[
  \varphi(a)=\frac{\lambda}{4\pi\varepsilon_0}
  \int_{-L_2}^{L_1}\frac{\dd y}{\sqrt{a^2+y^2}}
  =\frac{\lambda}{4\pi\varepsilon_0}
  \left[\operatorname{arsinh}\frac{L_1}{a}
  +\operatorname{arsinh}\frac{L_2}{a}\right].
\]
对称导线 $L_1=L_2=L$ 时
\[
  E_\perp=-\frac{\dd\varphi}{\dd a}
  =\frac{\lambda}{2\pi\varepsilon_0a}
  \frac{L}{\sqrt{a^2+L^2}}
  =\boxed{\frac{\lambda}{2\pi\varepsilon_0a}\sin\theta},
\]
与本章前文的矢量积分一致。非对称时不能只对 $a$ 求导就认为得到了全部电场。
在一般观察点 $(x,y)$、$x>0$ 写成
\[
 \varphi(x,y)=k\lambda\left[\operatorname{arsinh}\frac{L_1-y}{x}
 +\operatorname{arsinh}\frac{L_2+y}{x}\right].
\]
于是先对 $x,y$ 分别求偏导，再令 $(x,y)=(a,0)$：
\[
 E_x=\frac{k\lambda}{a}\left(\frac{L_1}{\sqrt{a^2+L_1^2}}+
 \frac{L_2}{\sqrt{a^2+L_2^2}}\right),\quad
 E_y=k\lambda\left(\frac1{\sqrt{a^2+L_1^2}}-\frac1{\sqrt{a^2+L_2^2}}\right).
\]
这也恢复了非对称有限导线的两个分量。

\textbf{3. 无限长直线。} 若直接令 $L\to\infty$，上式中的 $\varphi$ 发散；原因是无限直线的电场只按
$1/a$ 衰减，从有限位置积分到无穷远得到对数发散。解决办法是选择一个有限参考距离 $a_0$：
\[
  \boxed{\varphi(a)-\varphi(a_0)
  =-\frac{\lambda}{2\pi\varepsilon_0}\ln\frac{a}{a_0}}.
\]
求导仍给出
\[
  \boxed{E_a=-\frac{\dd\varphi}{\dd a}
  =\frac{\lambda}{2\pi\varepsilon_0a}}.
\]
所以“绝对电势发散”不妨碍有限位置间的电势差和电场有确定值。

\textbf{4. 无限大平面。} 设平面在 $z=0$，则
\[
  \boxed{\varphi(z)-\varphi(0)=-\frac{\sigma}{2\varepsilon_0}|z|},
\]
从而
\[
  \boxed{E_z=-\frac{\dd\varphi}{\dd z}
  =\frac{\sigma}{2\varepsilon_0}\operatorname{sgn}z}.
\]
这里也不能令 $\varphi(\infty)=0$，因为匀强电场的电势沿距离线性变化。

\textbf{5. 均匀圆盘。} 半径为 $R$、面密度为 $\sigma$ 的圆盘，在轴线上 $z>0$ 处的电势为
\[
  \varphi(z)=\frac{1}{4\pi\varepsilon_0}
  \int_0^R\frac{\sigma2\pi r\dd r}{\sqrt{z^2+r^2}}
  =\frac{\sigma}{2\varepsilon_0}
  \left(\sqrt{z^2+R^2}-z\right).
\]
因此
\[
  \boxed{E_z=-\frac{\dd\varphi}{\dd z}
  =\frac{\sigma}{2\varepsilon_0}
  \left(1-\frac{z}{\sqrt{z^2+R^2}}\right)}.
\]
\textbf{6. 无限长条带。} 条带位于 $z=0$，沿 $y$ 方向无限，宽度范围 $-b_1\le u\le b_2$。
取有限参考点 $P_0=(0,0,z_0)$，$z_0>0$。窄条 $\dd\lambda=\sigma\dd u$ 是无限线电荷，因而
\[
 \varphi(x,z)-\varphi(0,z_0)
 =-\frac{\sigma}{4\pi\varepsilon_0}\int_{-b_1}^{b_2}
 \ln\frac{(x-u)^2+z^2}{u^2+z_0^2}\,\dd u.
\]
分母固定参考电势，且比值无量纲。对 $x,z$ 求偏导，积分中的参考项导数为零：
\[
 E_x=\frac{\sigma}{2\pi\varepsilon_0}\int_{-b_1}^{b_2}
 \frac{x-u}{(x-u)^2+z^2}\,\dd u,\qquad
 E_z=\frac{\sigma}{2\pi\varepsilon_0}\int_{-b_1}^{b_2}
 \frac{z}{(x-u)^2+z^2}\,\dd u.
\]
在 $P=(0,0,a)$ 处得到
\[
 E_x=\frac{\sigma}{4\pi\varepsilon_0}\ln\frac{a^2+b_1^2}{a^2+b_2^2},\qquad
 E_z=\frac{\sigma}{2\pi\varepsilon_0}\left(\arctan\frac{b_1}{a}+\arctan\frac{b_2}{a}\right),
\]
与直接积分的条带结果相同；当 $b_1=b_2$ 时横向分量确为零。

有限系统可以取无穷远为零点；无限直线、无限条带和无限平面则应先计算电势差，或先保留有限尺寸、求导
得到电场后再取无限大极限。
\end{xesolution}



\section{电势应用与放电现象}

\subsection{电子伏特}


电荷量为 $q$ 的粒子从电势 $\varphi_A$ 运动到 $\varphi_B$，电场做功
\[
  W=q(\varphi_A-\varphi_B)=qU,
\]
其中 $U=\varphi_A-\varphi_B$ 是沿运动方向的电势降。电子通过 $1\,\mathrm V$ 的电势差获得的能量定义为
$1\,\mathrm{eV}$：
\begin{equation}
  \boxed{1\,\mathrm{eV}=1.602\,176\,634\times10^{-19}\,\mathrm J}.
\end{equation}
电子伏特特别适合描述原子、原子核和高能粒子的能量。由 $k_BT=1\,\mathrm{eV}$ 可得
\begin{equation}
  \boxed{T\approx1.2\times10^4\,\mathrm K}.
\end{equation}
因此做数量级估算时，可以把“约一万开尔文”与“热运动典型能量约 $1\,\mathrm{eV}$”联系起来。

中性原子中的电子被原子核的库仑吸引束缚。外来粒子若传给电子足够能量，电子就会脱离原子，原子
成为正离子，这个最低能量称为第一电离能。图~\ref{ch19:fig:s33-2-3}显示其随原子序数的周期性：每个周期的碱金属处较低，
稀有气体处较高；绝大多数原子的第一电离能在 $5$--$25\,\mathrm{eV}$ 的尺度内。

\par\noindent\begin{minipage}{\linewidth}
\ionizationEnergyFigure
\diagramcaption{原子的第一电离能随原子序数的周期性变化。覆盖 $Z=1$--$108$，符号区分元素类别；空心点表示含理论或估计标识的数据。}\label{ch19:fig:s33-2-3}
\end{minipage}\par\medskip

质量数较大的原子核可能发生 $\alpha$ 衰变，放出由两个质子和两个中子组成的氦核。母核的质量数减少 4，
核电荷数减少 2；典型 $\alpha$ 粒子能量约为 $5\,\mathrm{MeV}$。

\alphaDecayFigure

以 $5\,\mathrm{MeV}$ 的 $\alpha$ 粒子正碰金核为例，金的核电荷数 $Z=79$。若金核近似不动，最近距离
$r_m$ 由动能完全转化为库仑势能确定：
\[
  5\,\mathrm{MeV}
  =\frac{1}{4\pi\varepsilon_0}\frac{2\times79e^2}{r_m},
  \qquad r_m\approx4.5\times10^{-14}\,\mathrm m.
\]
在如此小的距离上仍观察到平方反比散射，说明原子的正电荷集中在远小于原子尺度的原子核中。

为了研究更高能量下的粒子相互作用，人们建造粒子加速器。大型电子--正电子对撞机曾把单束电子加速到
$104.5\,\mathrm{GeV}$，使对撞质心能量达到 $209\,\mathrm{GeV}$；大型强子对撞机的设计单束质子能量达到
$7\,\mathrm{TeV}$。把巨大能量集中在极小空间区域，可以产生并探测更重的粒子。宇宙线中还存在能量超过
$10^{20}\,\mathrm{eV}$ 的粒子，远高于人造加速器能量。

图~\ref{ch19:fig:s33-2-4} 中的蓝色磁体沿隧道弯曲排列，束流在内部真空管中传播；磁场用于约束轨道，粒子的能量由加速电场提供。

\par\noindent\begin{minipage}{\linewidth}
\acceleratorTunnelFigure
\diagramcaption{大型强子对撞机环形隧道中的超导磁体和束流管结构示意。环形主隧道周长约 $26.7\,\mathrm{km}$。}\label{ch19:fig:s33-2-4}
\end{minipage}\par\medskip


% EM teaching enrichment 20260927
\paragraph{先判能量变化，再换速度}
电子由静止从低电势点进入高 $U$ 的电势点，$q=-e$，故
$\Delta K=-q\Delta\varphi=eU>0$。若 $U=100\,\mathrm V$，其动能为 $100\,\mathrm{eV}$，非相对论速率为
\[
 v=\sqrt{\frac{2eU}{m_e}}\simeq5.93\times10^6\,\mathrm{m/s}.
\]
电子伏特是能量单位，不能直接代入以千克为质量单位的速度公式，须先换焦耳。若所得 $eU$ 与 $m_ec^2$ 同量级，应改用 $eU=(\gamma-1)m_ec^2$，而非继续使用 $mv^2/2$。


\subsection{气体放电}


气体通常近似绝缘，但过强电场会使它击穿。气体中少量自由电子先被电场加速；若电子在两次碰撞之间获得
的能量足以激发原子，原子退激时便可能发光；碰撞能量超过电离阈值时，还可能电离原子并产生新的电子。
超过阈值不意味着每次碰撞都电离，碰撞过程具有一定概率。电离产生的载流子足以补偿复合、扩散等损失时，
放电才能持续。辉光放电同样包含电离，并常依靠阴极二次电子发射补充电子，不能说成“只激发、不电离”。
辉光、弧光、火花是不同的放电状态，区别还取决于气压、电流、气体温度、电极及供电条件；弧光通常具有
较大的持续电流和明显加热，火花则是击穿时的快速放电。图~\ref{ch19:fig:s33-2-5}中的闪电是大尺度气体击穿。

\par\noindent\begin{minipage}{\linewidth}
\lightningFigure
\diagramcaption{闪电中的分支放电通道。亮通道是空气被击穿后形成的高温等离子体。}\label{ch19:fig:s33-2-5}
\end{minipage}\par\medskip

典型原子的电离能约为 $10\,\mathrm{eV}$。若电子与原子的碰撞截面取
$\sigma\sim10^{-20}\,\mathrm{m^2}$，理想气体数密度为 $n=P/(k_BT)$，则平均自由程估计为
\[
  l\sim\frac1{n\sigma}\sim\frac{k_BT}{P\sigma}.
\]
在 $T=300\,\mathrm K$、$P=1\,\mathrm{atm}$ 时，$l\sim4\times10^{-6}\,\mathrm m$。若电子要在一个
自由程内获得约 $10\,\mathrm{eV}$，所需场强量级为
\[
  E\sim\frac{10\,\mathrm V}{4\times10^{-6}\,\mathrm m}
  \approx2.5\times10^6\,\mathrm{V/m}.
\]
真实击穿场强还受气体种类、压强、湿度、电极曲率和初始电离等因素影响，但这一估算说明为什么尖端附近
的强局域电场容易触发放电。


% EM teaching enrichment 20260927
\paragraph{碰撞尺度估算的用途与限制}
把电离所需能量记为 $W_i$，单次自由飞行的获得能量约为 $eE\ell$。条件 $eE\ell\sim W_i$ 只检验一次碰撞是否可能有效，不能独自决定能否形成自持放电。
在固定温度、固定碰撞截面的粗略模型中 $\ell\propto1/P$，故相似的单次加速条件对应近似固定的 $E/P$。实际有限电极间隙还要能容纳足够多次增殖碰撞；气压极低时虽自由程变长，碰撞次数却减少。因此不能由该局部估算推断“气压越低必然越容易击穿”。


\subsection{专题例题与方法应用}

\begin{origexamplebox}{【零等势面】}\label{ch19:ex:s33-2-4}
小明将两个点电荷 $+Q,-q$ 分别置于 $(a,0,0)$ 和原点。求零电势等势面的方程（以无穷远为电势零点）。
\end{origexamplebox}\par\nopagebreak[4]

\begin{xesolution}
空间点 $P(x,y,z)$ 到 $+Q$ 和 $-q$ 的距离分别为
\[
  r_+=\sqrt{(x-a)^2+y^2+z^2},
  \qquad r_-=\sqrt{x^2+y^2+z^2}.
\]
零电势条件为
\[
  \frac{Q}{r_+}=\frac{q}{r_-}.
\]
令 $\mu=Q/q$，平方并整理：
\[
  \mu^2(x^2+y^2+z^2)=(x-a)^2+y^2+z^2.
\]
当 $\mu\ne1$ 时，配方得到球面
\[
  \boxed{\left(x+\frac{a}{\mu^2-1}\right)^2+y^2+z^2
  =\left(\frac{a\mu}{|\mu^2-1|}\right)^2}.
\]

\apolloniusSphereFigure

若 $Q>q$，球心在负 $x$ 轴，球面包围位于原点的较小电荷 $-q$；若 $Q<q$，则包围 $+Q$。特别地，
$Q=q$ 时不能除以 $\mu^2-1$，原方程直接化为
\[
  \boxed{x=\frac a2},
\]
即两电荷连线的中垂面。
\end{xesolution}



\begin{origexamplebox}{【无限长导线】}\label{ch19:ex:s33-2-5}
两根无限长直导线沿 $z$ 方向，带电线密度分别为 $+\lambda$ 和 $-\lambda$，与 $xOy$ 平面的交点分别为
$(a,0)$ 和 $(-a,0)$。求平面内任意等势线的方程及电场线方程。
\end{origexamplebox}\par\nopagebreak[4]

\begin{xesolution}
设 $P(x,y)$ 到正、负线电荷的距离为
\[
  r_+=\sqrt{(x-a)^2+y^2},\qquad
  r_-=\sqrt{(x+a)^2+y^2}.
\]
单根无限线电荷的绝对电势发散，但一正一负两根导线的发散常数相消。取中垂面 $x=0$ 为零电势，可写成
\[
  \boxed{\varphi(x,y)=\frac{\lambda}{2\pi\varepsilon_0}
  \ln\frac{r_-}{r_+}}.
\]
令 $C=\exp(2\pi\varepsilon_0\varphi_0/\lambda)>0$，等势线满足
\[
  \boxed{\frac{r_-}{r_+}=C},
\]
即
\[
  \boxed{(x+a)^2+y^2=C^2\bigl[(x-a)^2+y^2\bigr]}.
\]
$C\ne1$ 时是圆，$C=1$ 时退化为直线 $x=0$。

电场线与等势线正交。令
\[
  \beta=\arg\bigl[(x-a)+\mathrm i y\bigr]
  -\arg\bigl[(x+a)+\mathrm i y\bigr]
\]
为点 $P$ 看两导线所成的有向角，则 $\beta=$ 常量给出电场线。消去角度：
\[
  \boxed{x^2+y^2-a^2=2ay\cot\beta},
\]
或
\[
  \boxed{x^2+(y-a\cot\beta)^2=a^2\csc^2\beta}.
\]
这是一族通过 $(\pm a,0)$ 的圆；$y=0$ 是需单独保留的直线分支。
为何有向角不变？令 $u=\ln(r_-/r_+)$，直接求导可得
$\beta_x=u_y$、$\beta_y=-u_x$。而 $\boldsymbol E=-(\lambda/2\pi\varepsilon_0)\nabla u$，
所以 $\nabla\beta\cdot\boldsymbol E=0$，沿场线 $\beta$ 不变。
等势圆的圆心 $(c,0)$ 和半径 $R_c$ 满足 $c^2-R_c^2=a^2$；
场线圆的圆心 $(0,d)$ 和半径 $R_d$ 满足 $R_d^2-d^2=a^2$。
两圆相交于 $P=(x,y)$ 时，相应半径的点积为
$ (x-c,y)\cdot(x,y-d)=x^2+y^2-cx-dy=0$，证明两族曲线正交。

\twoLinesFamiliesFigure

图中 $+\lambda$ 位于右侧，$-\lambda$ 位于左侧，箭头从正电荷指向负电荷；实线等势圆与虚线场线圆处处正交。三维空间中，两类截线沿 $z$ 方向平移，分别形成等势柱面和场线
柱面。
\end{xesolution}



\begin{origexamplebox}{【有限长导线】}\label{ch19:ex:s33-2-6}
一根有限长直导线带电线密度 $\lambda$，位于 $(a,0,0)$ 到 $(-a,0,0)$ 之间。系统绕 $x$ 轴旋转对称。
求包含 $x$ 轴的平面内任意等势线的方程及电场线方程。
\end{origexamplebox}\par\nopagebreak[4]

\begin{xesolution}
在包含导线的子午面内，以 $s=\sqrt{y^2+z^2}$ 表示点到 $x$ 轴的距离，并定义点到两个端点的距离
\[
  r_+=\sqrt{(x-a)^2+s^2},
  \qquad r_-=\sqrt{(x+a)^2+s^2}.
\]
直接积分得到
\[
  \varphi(x,s)=\frac{\lambda}{4\pi\varepsilon_0}
  \int_{-a}^{a}\frac{\dd x'}{\sqrt{(x-x')^2+s^2}}
  =\boxed{\frac{\lambda}{4\pi\varepsilon_0}
  \ln\frac{r_++r_-+2a}{r_++r_--2a}}.
\]
因此 $\varphi=$ 常量等价于
\[
  \boxed{r_++r_-=\text{常量}},
\]
在子午面内是一族以导线两端为焦点的椭圆，绕 $x$ 轴旋转后成为长旋转椭球面。

引入长旋转椭圆坐标
\[
  \xi=\frac{r_++r_-}{2a},
  \qquad \eta=\frac{r_--r_+}{2a}.
\]
由两焦点距离可反解出
\[
 x=a\xi\eta,\qquad s=a\sqrt{(\xi^2-1)(1-\eta^2)},\quad \xi\ge1,\quad -1\le\eta\le1.
\]
两条坐标切向量的点积为
\[
 \frac{\partial(x,s)}{\partial\xi}\cdot\frac{\partial(x,s)}{\partial\eta}
 =a^2\xi\eta-
 a^2\frac{\xi\sqrt{1-\eta^2}}{\sqrt{\xi^2-1}}
 \frac{\eta\sqrt{\xi^2-1}}{\sqrt{1-\eta^2}}=0.
\]
因此它们正交。电势只依赖 $\xi$，场垂直于 $\xi=$ 常量的椭圆，
沿 $\eta=$ 常量的坐标线，即
\[
  \boxed{r_--r_+=\text{常量}}.
\]
它在子午面内是一族共焦双曲线，绕轴旋转后成为双曲面。

\finiteRodFamiliesFigure

图中粗黑线表示 $-a\le x\le a$ 的均匀带电导线（$\lambda>0$）；实线是等势椭圆，虚线箭头从导线不同位置向外。椭圆与双曲线在每个交点正交。导线两端是共同焦点，不是椭圆或双曲线的圆心；这一位置关系是正确画图的
关键。
\end{xesolution}



\section{静电能量与带电粒子运动}

\subsection{静电势能}


必须区分“一个试探电荷在外电场中的势能”与“构造整个带电系统所需的总静电势能”。前者是
$E_p=q\varphi$；后者是把全部电荷从无穷远逐个搬到指定位置时，外力所做功的总和。

设系统由 $N$ 个点电荷 $Q_i$ 构成，$r_{ij}$ 为 $Q_i,Q_j$ 的距离。把每一对电荷的相互作用能只计算一次，
得到
\begin{equation}
  \boxed{U=\sum_{i=1}^{N}\sum_{j=1}^{i-1}
  \frac{1}{4\pi\varepsilon_0}\frac{Q_iQ_j}{r_{ij}}
  =\frac12\sum_{i=1}^{N}Q_i\varphi_i'}.
\end{equation}
$\varphi_i'$ 是除 $Q_i$ 自身以外的其他电荷在 $Q_i$ 处产生的电势。第二种写法前面的 $1/2$ 不能遗漏，
因为直接求和 $Q_i\varphi_i'$ 会把每一对相互作用计算两次；点电荷自身的无穷大自能并不包含在式中。

对总电荷为 $Q$、半径为 $R$ 的均匀带电球壳，表面电势为
$Q/(4\pi\varepsilon_0R)$，所以
\begin{equation}
  \label{ch19:ref:shell-energy}\boxed{U=\frac12Q\varphi(R)
  =\frac{1}{2}\frac{1}{4\pi\varepsilon_0}\frac{Q^2}{R}}
  =\boxed{\frac{Q^2}{8\pi\varepsilon_0R}}.
\end{equation}
半径越小，静电势能越高。压缩球壳时必须克服表面电荷之间的斥力，外力所做的功就储存在静电系统中。


% EM teaching enrichment 20260927
\paragraph{用逐个装配核对二分之一}
三个相同点电荷 $q$ 放在边长 $a$ 的等边三角形顶点。搬入第一个不做相互作用功，第二个需 $kq^2/a$，第三个需 $2kq^2/a$，所以总能为 $3kq^2/a$。
若只把第三个移到无穷远、其余固定，电场对它做功 $2kq^2/a$，此处应使用第三个在既定外场中的势能 $q\varphi'$，不能再乘 $1/2$。
“整个系统能量按电荷对计数”与“一个粒子在不变外场中的能量”是两种不同问题，辨明过程比记忆系数更可靠。


\subsection{电场能量}


静电能不仅可以看成电荷对之间的相互作用能，也可以看成分布在电场中的能量。真空中电场强度为 $E$ 的
位置具有能量密度
\begin{equation}
  \boxed{w_E=\frac12\varepsilon_0E^2}.
\end{equation}
从无场状态建立电场必须输入能量。例如把原本紧贴的等量异号电荷拉开，外力做功形成电场；移动电荷时，
场的分布和储存的能量也随之改变。

均匀带电球壳内部 $E=0$，外部电场为点电荷场。球外距离 $r$ 处的能量密度是
\[
  w_E(r)=\frac12\varepsilon_0
  \left(\frac{1}{4\pi\varepsilon_0}\frac{Q}{r^2}\right)^2.
\]
把它对球壳外部全空间积分：
\begin{equation}
  \boxed{U=\int_R^\infty w_E(r)\,4\pi r^2\dd r}.
\end{equation}
代入后得到 $U=Q^2/(8\pi\varepsilon_0R)$，与式\eqref{ch19:ref:shell-energy}完全一致。电荷对求和强调相互作用，能量密度积分
强调能量保存在空间中的场里；后者可以推广到随时间变化的电磁场。

本章使用的核心概念依次是：电势能、电势、等势面、静电系统的相互作用能以及电场能量密度。计算时必须
同时检查电势零点是否能取在无穷远、求和是否重复计算电荷对，以及所求的是单个电荷的势能还是整个系统
的总能量。

% EM teaching enrichment 20260927
\paragraph{总场平方中的交叉项}
两组固定电荷产生 $\boldsymbol E_1,\boldsymbol E_2$ 时，场能含
\[
 \frac{\varepsilon_0}{2}(\boldsymbol E_1+\boldsymbol E_2)^2
 =\frac{\varepsilon_0E_1^2}{2}+\frac{\varepsilon_0E_2^2}{2}
 +\varepsilon_0\boldsymbol E_1\cdot\boldsymbol E_2.
\]
前两项是各自自能，交叉项给出相互作用能；它可以为负，虽然完整场能密度非负。
移动刚性带电体而不改变自身分布时，自能不变，求能量差可只算相互作用项。点电荷模型的自能发散，不能把含自场的 $q\varphi$ 当成有限相互作用势能。




% EM structure and teaching revision 20260928
\subsection{组装能、相互作用能与自能}
把若干点电荷从无穷远逐个移入，先前已放置的电荷对后来电荷做功。每对电荷只计一次，因此
\[
 W=\sum_{i<j}\frac{kq_iq_j}{r_{ij}}
 =\frac12\sum_i q_i\varphi_{\mathrm{other}}(\boldsymbol r_i).
\]
右式的电势必须排除该点电荷自身的发散项。两个同号电荷由很远处靠近时 $W$ 增大，需要外力做正功；异号电荷靠近则相反。
连续分布可写为 $W=\tfrac12\int\rho\varphi\dd V$。对于半径 $R$ 的均匀实心带电球，也可逐层组装：已放好的内球电荷为 $q(r)=Qr^3/R^3$，新薄壳为 $\dd q=3Qr^2\dd r/R^3$，故
\[
 W=\int_0^R\frac{kq(r)}r\dd q=\frac{3kQ^2}{5R}.
\]
这包含整个连续球的自能；若题目仅问两部分之间的相互作用能，则应减去各部分各自的自能。把各区域的场直接平方相加会漏掉交叉项。
用于粒子运动时，先界定哪些带电体固定、哪些能运动。单粒子在固定外场中可用 $K+q\varphi=\mathrm{const}$；若场源也运动，则须把场源的动能和体系相互作用能一起计入。


\subsection{专题例题与方法应用}

\begin{origexamplebox}{【电子透镜】}\label{ch19:ex:s32-1-7}
在电子显微镜中，一束原先通过电势差 $U_0$ 加速、具有一定能量的电子，经过一根与电子束原来方向垂直的
绷紧带电细金属线。该线带正电，电荷线密度为 $\lambda$；金属线产生的电场相对较弱，因此电子轨道只稍受
影响。入射电子的瞄准参数为 $b$，求电子通过带电金属线以后运动方向的偏转角。
\end{origexamplebox}\par\nopagebreak[4]

\begin{xesolution}
取 $b>0$ 为瞄准距离的大小，并采用非相对论、细线及弱偏转近似。电子的电荷为 $-e$，进入相互作用区前，加速电势差把电势能转化为动能：
\[
  \frac12mv^2=eU_0.
\]
正线电荷的场沿径向向外，电子受力指向线电荷。弱偏转近似中，可以沿原来的水平直线计算横向冲量。
令电子在 $x$ 方向运动，线电荷在原点垂直纸面，未偏转轨迹为 $y=b$，则 $x=vt$。

\electronLensFigure

线电荷场强为
\[
  E(r)=\frac{\lambda}{2\pi\varepsilon_0r},
  \qquad r=\sqrt{x^2+b^2}.
\]
横向力的大小为
\[
  |F_y|=eE\frac{b}{r}
  =\frac{e\lambda}{2\pi\varepsilon_0}
  \frac{b}{x^2+b^2}.
\]
于是横向动量改变量
\[
  |\Delta p_y|
  =\int_{-\infty}^{\infty}|F_y|\dd t
  =\frac{e\lambda}{2\pi\varepsilon_0v}
  \int_{-\infty}^{\infty}\frac{b\,\dd x}{x^2+b^2}
  =\frac{e\lambda}{2\varepsilon_0v}.
\]
偏转角很小时 $\alpha\approx|\Delta p_y|/(mv)$，代入 $mv^2=2eU_0$：
\[
  \boxed{\alpha=\frac{e\lambda}{2\varepsilon_0mv^2}
  =\frac{\lambda}{4\varepsilon_0U_0}}.
\]
在这一近似下，偏转角与电子质量、电荷量以及瞄准距离 $b$ 都无关，只由线电荷密度和加速电压决定。
$b$ 的消失来自两个效应的抵消：靠得更近时横向力更大，但强作用区域的通过时间更短。

这里把积分限延至无穷，是对有限装置中横向冲量的近似，不是说无限线电荷的对数电势在无穷远有有限值。
若实际主要相互作用区为 $-X\le x\le X$，$X\gg b$，同一恒速近似给出
\[
 |\Delta p_y|=\frac{e\lambda}{\pi\varepsilon_0v}\arctan\frac Xb,
 \qquad \alpha_X=\frac{\lambda}{2\pi\varepsilon_0U_0}\arctan\frac Xb.
\]
令 $X/b$ 很大才回到上面的结果。同时须满足
$[\lambda/(2\pi\varepsilon_0U_0)]\ln(X/b)\ll1$，使纵向动能变化也远小于初始动能；
细线半径须远小于 $b$，线长须足以近似为无限线。
若这些条件不满足，应给定实际场区和入射条件再求轨迹，不能将该常数偏转角公式直接外推。
\end{xesolution}
\begin{origexamplebox}{【均匀带电球】}\label{ch19:ex:s33-2-2}
小明给一个金属球壳均匀充电，得到一个半径为 $R$、电量为 $Q$ 的均匀带电球壳。试求：
\begin{enumerate}[label=\arabic*.]
  \item 空间中总的电场能量；
  \item 必须用力压住球壳才能保证其不会四散飞走。需要多大的压强？
  \item 若是均匀带电的绝缘实心球，答案有什么变化？
  \item 为什么常见的是均匀带电的“金属球壳”和“绝缘球体”？
\end{enumerate}
\end{origexamplebox}\par\nopagebreak[4]

\begin{xesolution}
金属球壳与绝缘实心球的电场分布如下。黑线是球壳：内部场为零，表面外侧跃变到
$Q/(4\pi\varepsilon_0R^2)$；蓝线是均匀带电实心球：内部场从球心开始线性增加，并在表面与外部点电荷场
连续相接。

\chargedSphereEnergyFigure

\textbf{1. 金属球壳的能量。} 球外
$E=Q/(4\pi\varepsilon_0r^2)$，球内 $E=0$，所以
\[
  U=\int_R^\infty\frac12\varepsilon_0E^2\,4\pi r^2\dd r
  =\boxed{\frac{Q^2}{8\pi\varepsilon_0R}}.
\]

\textbf{2. 维持球壳的压强。} 设球壳半径发生准静态改变，电荷量保持 $Q$ 不变。电斥力对应的向外广义力
为
\[
  F_R=-\frac{\dd U}{\dd R}
  =\frac{Q^2}{8\pi\varepsilon_0R^2}.
\]
除以球面积 $4\pi R^2$，得到需要由材料或外界提供的向内压强
\[
  \boxed{p_e=\frac{F_R}{4\pi R^2}
  =\frac{Q^2}{32\pi^2\varepsilon_0R^4}}
  =\boxed{\frac{\sigma_s^2}{2\varepsilon_0}}
  =\boxed{\frac12\varepsilon_0E_{\rm out}^2},
\]
其中 $\sigma_s=Q/(4\pi R^2)$。不能用 $\sigma_sE_{\rm out}$，因为面电荷受到的是去掉自身奇异场后的平均
场 $(E_{\rm out}+E_{\rm in})/2=E_{\rm out}/2$。

\textbf{3. 均匀带电绝缘实心球。} 球内体电荷密度
$\rho=3Q/(4\pi R^3)$，由高斯定律
\[
  E(r)=\frac{Qr}{4\pi\varepsilon_0R^3}\quad(r<R),
  \qquad
  E(r)=\frac{Q}{4\pi\varepsilon_0r^2}\quad(r>R).
\]
内、外能量相加：
\[
  U_{\rm in}=\frac{Q^2}{40\pi\varepsilon_0R},
  \qquad U_{\rm out}=\frac{Q^2}{8\pi\varepsilon_0R},
\]
故
\[
  \boxed{U_{\rm solid}=\frac{3Q^2}{20\pi\varepsilon_0R}
  =\frac35\frac{1}{4\pi\varepsilon_0}\frac{Q^2}{R}}.
\]
绝缘球内部存在体电荷力，维持它的机械应力随半径变化，不能只用一个表面压强描述。若只问把整个球作
相似压缩所需的平均广义压强，则
\[
  \bar p=-\frac{\dd U}{\dd V}
  =\frac{3Q^2}{80\pi^2\varepsilon_0R^4}.
\]

\textbf{4. 材料区别。} 导体达到静电平衡后，内部自由电荷会迁移到外表面，因此均匀带电导体自然表现为
“带电球壳”；绝缘体中的电荷可以被冻结在体内，才可能形成均匀体电荷分布。几何上当然也能制作绝缘空壳，
但“均匀带电金属实心球”在静电平衡下不可能保持均匀体电荷。
\end{xesolution}



\begin{origexamplebox}{【相互作用能】}\label{ch19:ex:s33-2-3}
小明尝试构建一个由 10 个点电荷组成的体系，电荷记为 $Q_i$，$i=1,2,\ldots,10$；任意两电荷之间的
距离为 $r_{ij}$。为了把这 10 个电荷从无穷远搬到对应位置，小明需要做功多少？
\end{origexamplebox}\par\nopagebreak[4]

\begin{xesolution}
按任意顺序逐个移入电荷。第一个电荷移入时没有其他电荷，外力做功为零；第二个电荷只与第一个发生相互
作用；第 $i$ 个电荷移入时，只需计算它与已经就位的 $i-1$ 个电荷的作用。这样每一对电荷恰好出现一次。

\chargeAssemblyFigure

外力在准静态组装过程中所做的总功等于系统静电势能：
\[
  \boxed{W_{\rm ext}=U
  =\frac{1}{4\pi\varepsilon_0}
  \sum_{i=1}^{10}\sum_{j=1}^{i-1}\frac{Q_iQ_j}{r_{ij}}}.
\]
也可写成
\[
  \boxed{U=\frac12\sum_{i=1}^{10}Q_i\varphi_i'},
  \qquad
  \varphi_i'=\frac{1}{4\pi\varepsilon_0}
  \sum_{j\ne i}\frac{Q_j}{r_{ij}}.
\]
若异号电荷使某些项为负，外力总功也可能为负，表示组装过程中静电力向外界释放能量。答案与移入顺序
无关，因为静电力是保守力。
\end{xesolution}



\begin{origexamplebox}{【同性相斥】}\label{ch19:ex:s33-2-7}
一个可视为点电荷的小球 $A$，质量为 $m$、电量为 $+q$，用长度未知的轻绳悬挂。
另一小球 $B$ 也带电 $+q$，初始足够远，以两球初始相互作用能为零的极限处理，$A$ 初始竖直静止。
小明缓慢移动 $B$ 并排斥 $A$，使 $A$ 始终近似平衡；最终 $B$ 占据 $A$ 原来的最低位置，$A$ 上升 $h$。
忽略两球尺寸变化及 $B$ 的重力功（所求手的功仅计克服静电相互作用的部分），求小明做的功和电场能量增量。
\par\noindent\begin{minipage}{\linewidth}
\emRepellingQuestion
\figcap{悬点、轻绳及两球的最终位置；虚线标出原来的最低位置。}\label{ch19:fig:em-source-thirty-third-section-7}
\end{minipage}\par\medskip
\end{origexamplebox}\par\nopagebreak[4]

\begin{xesolution}
取初态 $B$ 在足够远处、$A$ 竖直悬挂，末态缓慢达到平衡。设绳长为 $L$，末态绳与竖直方向夹角为
$\theta$，两球距离为 $r$。$A,B$ 都在以悬点为圆心、半径为 $L$ 的圆周上，所以
\[
  r=2L\sin\frac\theta2,
  \qquad h=L(1-\cos\theta)=2L\sin^2\frac\theta2.
\]
因此
\[
  \boxed{h=r\sin\frac\theta2}.
\]

图~\ref{ch19:fig:s33-2-6} 中两段悬线位置均长 $L$，夹角为 $\theta$；$A,B$ 的连线决定电斥力方向，$h$ 是两球的高度差。

\par\noindent\begin{minipage}{\linewidth}
\repellingSpheresFigure
\diagramcaption{末态几何与 $A$ 球受力。$B$ 位于 $A$ 原来的最低位置，三力方向由实际位置唯一确定。}\label{ch19:fig:s33-2-6}
\end{minipage}\par\medskip

沿 $A$ 球运动圆周的切向列平衡条件。电斥力
$F_e=kq^2/r^2$，其切向分量为 $F_e\cos(\theta/2)$；重力切向分量为 $mg\sin\theta$：
\[
  F_e\cos\frac\theta2=mg\sin\theta
  =2mg\sin\frac\theta2\cos\frac\theta2.
\]
故
\[
  F_e=2mg\sin\frac\theta2=\frac{2mgh}{r}.
\]
再代入 $F_e=kq^2/r^2$：
\[
  \frac{kq^2}{r}=2mgh.
\]
两球相互作用的电场能量从零增至
\[
  \boxed{\Delta U_E=\frac{kq^2}{r}=2mgh}.
\]
$A$ 的重力势能增加 $mgh$，绳的拉力不做功；准静态过程中小明所做功等于总势能增加：
\[
  \boxed{W_{\text{小明}}=\Delta U_E+\Delta U_g=3mgh}.
\]
答案与未知绳长 $L$ 无关，正是因为平衡条件和圆周几何把 $L,\theta,r$ 全部消去了。
\end{xesolution}



\begin{origexamplebox}{【玻尔氢原子】}\label{ch19:ex:s33-2-8}
在玻尔氢原子模型中，电子在原子核外以光速的 $1/137$ 绕原子核作圆周运动。这个氢原子的总能量是多少？
为什么是负能量？
\end{origexamplebox}\par\nopagebreak[4]

\begin{xesolution}
氢核带电 $+e$，电子带电 $-e$。库仑吸引力提供向心力：
\[
  \frac{m_ev^2}{r}=\frac{1}{4\pi\varepsilon_0}\frac{e^2}{r^2}.
\]
因此
\[
  \frac{1}{4\pi\varepsilon_0}\frac{e^2}{r}=m_ev^2.
\]

\bohrAtomFigure

电子动能和电势能分别为
\[
  K=\frac12m_ev^2,
  \qquad U=-\frac{1}{4\pi\varepsilon_0}\frac{e^2}{r}=-m_ev^2.
\]
所以总能量
\[
  E=K+U=-\frac12m_ev^2.
\]
代入 $v=c/137$：
\[
  \boxed{E=-\frac{m_ec^2}{2\times137^2}
  \approx-13.6\,\mathrm{eV}}.
\]
同时可得轨道半径 $r\approx5.29\times10^{-11}\,\mathrm m$。总能量为负，是因为把相距无穷远、静止的
质子和电子规定为零能量；氢原子是束缚态，要输入 $13.6\,\mathrm{eV}$ 才能把电子完全移到无穷远。
\end{xesolution}



\begin{origexamplebox}{【$\alpha$ 粒子散射】}\label{ch19:ex:s33-2-9}
小明用一个 $\alpha$ 粒子打向带 $Z$ 个单位正电荷的金原子核。
\begin{enumerate}[label=\arabic*.]
  \item 若认为金原子核不动，$\alpha$ 粒子的轨迹如何随初始状态变化？
  \item 若有很多 $\alpha$ 粒子同时入射，出射的微分散射截面是多少？
  \item 历史上谁第一次做了这个实验？
\end{enumerate}
\end{origexamplebox}\par\nopagebreak[4]

\begin{xesolution}
记 $\alpha$ 粒子电荷为 $+2e$，入射动能为 $E_k=\tfrac12mv_\infty^2$，瞄准参数为 $b$。把库仑
排斥势写成
\[
  V(r)=\frac{\kappa}{r},
  \qquad
  \kappa=\frac{1}{4\pi\varepsilon_0}(2e)(Ze).
\]

\rutherfordFigure
图中三条曲线使用同一 $\kappa/(2E_k)=0.7$ 和同一核位置，长度以同一个任意单位表示。
虚线是入射与出射渐近线，$b$ 是入射渐近线到核的垂距；蓝线的 $b$ 最小，因而偏转最大。


\textbf{1. 轨迹。} 能量与角动量守恒，角动量为 $L=mv_\infty b$。排斥平方反比力的轨迹是以原子核为
焦点的双曲线。为补出角度关系，令 $u=1/r$，以最近点的径向为极角 $\psi=0$。
比奈方程为 $u''+u=-m\kappa/L^2$，其对称解为
\[
 u=\frac{m\kappa}{L^2}(\epsilon\cos\psi-1).
\]
将最近点 $u'(0)=0$ 代入能量式
$E_k=(L^2/2m)(u'^2+u^2)+\kappa u$，得到
$\epsilon^2=1+2E_kL^2/(m\kappa^2)=1+(2E_kb/\kappa)^2$。
两条渐近线处 $u=0$，故 $\cos\psi_\infty=1/\epsilon$。
速度入射方向与入射位置矢量反向，因而偏转角
$\Theta=\pi-2\psi_\infty$，有
$\sin(\Theta/2)=1/\epsilon$。消去 $\epsilon$ 后，散射角 $\Theta$ 与瞄准参数满足
\[
  \boxed{b=\frac{\kappa}{2E_k}\cot\frac\Theta2},
  \qquad
  \boxed{\tan\frac\Theta2=\frac{\kappa}{2E_kb}}.
\]
$b$ 越小或入射能量越低，偏转越强；$b=0$ 时发生正碰并反向散射，最近距离
\[
  \boxed{r_{\min}=\frac{\kappa}{E_k}}.
\]
双曲线离心率为
\[
  \epsilon=\sqrt{1+\left(\frac{2E_kb}{\kappa}\right)^2}>1.
\]

\textbf{2. 卢瑟福微分截面。} 入射束中瞄准参数在 $b$ 到 $b+\dd b$ 的粒子占据环带面积
$\dd\sigma=2\pi b|\dd b|$；散射到 $\Theta$ 到 $\Theta+\dd\Theta$ 的立体角为
$\dd\Omega=2\pi\sin\Theta\dd\Theta$。令 $B_0=\kappa/(2E_k)$，则
\[
 \left|\frac{\dd b}{\dd\Theta}\right|=\frac{B_0}{2}\csc^2\frac\Theta2,
 \qquad \frac{\dd\sigma}{\dd\Omega}=
 \frac{b}{\sin\Theta}\left|\frac{\dd b}{\dd\Theta}\right|.
\]
代入 $b=B_0\cot(\Theta/2)$ 和 $\sin\Theta=2\sin(\Theta/2)\cos(\Theta/2)$，得到
\[
  \boxed{\frac{\dd\sigma}{\dd\Omega}
  =\left(\frac{\kappa}{4E_k}\right)^2
  \frac{1}{\sin^4(\Theta/2)}}.
\]
若关心散射角大于 $\Theta_0$ 的总截面，则
\[
  \boxed{\sigma(\Theta>\Theta_0)
  =\pi\left(\frac{\kappa}{2E_k}\right)^2
  \cot^2\frac{\Theta_0}{2}}.
\]

\textbf{3. 历史。} 1909 年，汉斯·盖革和恩斯特·马斯登在恩斯特·卢瑟福指导下完成金箔
$\alpha$ 粒子散射实验；卢瑟福于 1911 年用核式原子模型解释大角散射。严格说，实际执行实验并最先
报告结果的是盖革与马斯登，提出决定性理论解释的是卢瑟福。
\end{xesolution}

\endgroup

\endgroup

\begingroup
\tcbset{breakable/.default=false,after={\par\nobreak\smallskip}}
\RenewDocumentEnvironment{xesolution}{+b}{%
  \par\nobreak\smallskip\nobreak\begingroup\small
  \clubpenalties 2 10000 10000\displaywidowpenalties 2 10000 10000\noindent
  {\bfseries\color{DeepBlue}解答}\quad #1\par\endgroup\smallskip
}{}
\begingroup
\renewcommand{\Needspace}[1]{\par\begingroup
\dimen0=#1\relax\vskip0pt plus\dimen0\penalty-100
\vskip0pt plus-\dimen0\vskip\dimen0\penalty9999\vskip-\dimen0
\vskip0pt\endgroup}
\tcbset{breakable/.default=false}
\chapter{导体、电介质与电容器}

{
\graphicspath{{/Users/dongxingxi/Documents/物理竞赛一轮资料/sample/xes-original/thirty-fourth-section/}{/Users/dongxingxi/Documents/物理竞赛一轮资料/sample/xes-original/thirty-fourth-section/assets/}}
\section{导体平衡、屏蔽与电荷响应}

本节先用等势与电荷守恒确定导体的约束。后面的异号导体题还会用到静电势的最大值原理：无电荷区域的电势满足拉普拉斯方程，非恒定解不能在区域内部取得严格极大或极小值，极值由边界控制。这为判断导体电势与表面电荷符号提供工具；唯一性定理及电像构造将在下一节系统展开。

% All annotations have transparent backgrounds so labels never mask curves.
\tikzset{
  xeslabel/.style={font=\small,text=Ink,fill=none,inner sep=.55pt},
  xessmall/.style={font=\footnotesize,text=TextGray,fill=none,inner sep=.4pt},
  xestiny/.style={font=\scriptsize,text=TextGray,fill=none,inner sep=.3pt},
  xesline/.style={draw=Ink,line width=.85pt},
  xesblue/.style={draw=DeepBlue,line width=1pt},
  xesarrow/.style={-{Latex[length=2.1mm,width=1.3mm]},draw=DeepBlue,line width=.9pt},
  fieldline/.style={draw=DeepBlue,line width=.82pt,postaction={decorate},
    decoration={markings,mark=at position .58 with {\arrow{Latex[length=1.65mm,width=1.05mm]}}}},
  eqline/.style={draw=RuleGray,line width=.66pt,densely dashed}
}


\newcommand{\conductorFieldFigure}{%
\begin{center}
\includegraphics[width=.76\linewidth]{/Users/dongxingxi/Documents/物理竞赛一轮资料/sample/xes-original/restructured-complete/assets/revision-ch19-ch20/conductor-pair.pdf}
\end{center}}


\newcommand{\inductionSequenceFigure}{%
\begin{center}
\begin{tikzpicture}[x=1cm,y=1cm]
  \foreach \x/\tag in {0/(a),2.7/(b),5.4/(c),8.1/(d),10.8/(e)}{
    \begin{scope}[shift={(\x,0)}]
      \fill[DeepBlue!6] (0,0) circle (.67);\draw[xesblue] (0,0) circle (.67);
      \node[xestiny] at (0,-1.0) {\tag};
    \end{scope}}
  % (a) neutral
  \foreach \a in {0,60,...,300}{\node[xestiny,text=DeepBlue] at (\a:.48) {$+$};\node[xestiny,text=Ink] at (\a+30:.35) {$-$};}
  % negative rod and polarized spheres in b--d
  \foreach \x in {2.7,5.4,8.1}{
    \fill[Ink!15,rounded corners=2pt] (\x-1.35,-.18) rectangle (\x-.82,.18);
    \draw[Ink,line width=.65pt,rounded corners=2pt] (\x-1.35,-.18) rectangle (\x-.82,.18);
    \foreach \dx in {-1.24,-1.08,-.92}{\node[xestiny,text=Ink] at (\x+\dx,0) {$-$};}
    \foreach \a in {125,155,185,215}{\node[xestiny,text=DeepBlue] at ($(\x,0)+(\a:.52)$) {$+$};}}
  \foreach \a in {-55,-20,20,55}{\node[xestiny,text=Ink] at ($(2.7,0)+(\a:.52)$) {$-$};}
  % ground lead in c
  \draw[xesline] (5.92,-.32)--(6.35,-.32)--(6.35,-.62);
  \draw[xesline] (6.08,-.62)--(6.62,-.62) (6.17,-.73)--(6.53,-.73) (6.27,-.84)--(6.43,-.84);
  % (e) retained positive charge redistributes
  \foreach \a in {0,45,...,315}{\node[xestiny,text=DeepBlue] at ($(10.8,0)+(\a:.51)$) {$+$};}
  \node[xessmall,align=center] at (5.4,-1.42) {靠近负电体 $\rightarrow$ 接地 $\rightarrow$ 先断地线 $\rightarrow$ 移走负电体};
\end{tikzpicture}
\end{center}}

\newcommand{\faradayCageFigure}{%
\begin{center}
  \includegraphics[width=.78\linewidth]{faraday-cage-built-in-v2.png}

  \vspace{.25em}
  {\footnotesize\color{TextGray}
    紫白色放电只接触金属网笼外表面；笼内人员所在区域没有放电通道。}
\end{center}}

\newcommand{\discretePotentialFigure}{%
\begin{center}
\begin{tikzpicture}[x=.95cm,y=1cm]
  \draw[xesline,-{Latex[length=1.8mm]}] (-.4,0)--(11.8,0) node[xeslabel,below left] {$x$};
  \fill[DeepBlue!24] (0,-.16) rectangle (10.5,.16);
  \foreach \i in {0,...,10}{\fill[DeepBlue] (\i*1.05,0) circle (2.1pt);}
  \node[xeslabel,below=6pt] at (0,0) {$A$};\node[xeslabel,below=6pt] at (10.5,0) {$B$};
  \foreach \i/\lab in {0/0,1/1,2/2,4/{n-1},5/n,6/{n+1},8/{N-2},9/{N-1},10/N}{
    \node[xestiny,above=5pt] at (\i*1.05,0) {$\varphi_{\lab}$};}
  \draw[<->,Gold!85,line width=.7pt] (5.25,-.45)--(6.3,-.45);
  \node[xestiny,text=Gold!80!black] at (5.78,-.68) {$a$};
\end{tikzpicture}
\end{center}}

\newcommand{\groundedSphereImageFigure}{%
\begin{center}
\includegraphics[width=.76\linewidth]{/Users/dongxingxi/Documents/物理竞赛一轮资料/sample/xes-original/restructured-complete/assets/revision-ch19-ch20/grounded-sphere.pdf}
\end{center}}


\newcommand{\ionosphereFigure}{%
\begin{center}
  \includegraphics[width=.82\linewidth]{ionosphere-communication-built-in-v2.png}

  \vspace{.25em}
  \begin{tikzpicture}[baseline=-.55ex]
    \draw[orange!88!black,line width=1.25pt,-{Latex[length=1.8mm]}] (0,0)--(1.25,0);
    \node[xessmall,anchor=west] at (1.4,0) {$\omega<\omega_p$：在电离层返回地面};
    \draw[cyan!70!DeepBlue,line width=1.25pt,-{Latex[length=1.8mm]}] (6.0,0)--(7.25,0);
    \node[xessmall,anchor=west] at (7.4,0) {$\omega>\omega_p$：穿透电离层到达卫星};
  \end{tikzpicture}
\end{center}}

\newcommand{\plasmaSlabFigure}{%
\begin{center}
\begin{tikzpicture}[x=1cm,y=1cm]
  \fill[DeepBlue!8] (0,0) rectangle (9.2,3.5);\draw[xesblue] (0,0) rectangle (9.2,3.5);
  \fill[DeepBlue!30] (0,0) rectangle (.65,3.5);
  \fill[Gold!25] (8.55,0) rectangle (9.2,3.5);
  \foreach \x in {1.0,1.7,...,8.2}{\foreach \y in {.55,1.25,1.95,2.65}{\node[xestiny,text=RuleGray] at (\x,\y) {$+$};}}
  \node[xessmall,text=Ink] at (.33,1.75) {$-\sigma$};
  \node[xessmall,text=Ink] at (8.88,1.75) {$+\sigma$};
  \foreach \y in {.65,1.35,2.05,2.75}{\draw[xesarrow] (8.3,\y)--(.9,\y);}
  \node[xeslabel] at (4.6,-.55) {$\boldsymbol E'=-\dfrac{nq}{\varepsilon_0}x\,\hat{\boldsymbol x}$};
  \draw[<->,Gold!80,line width=.65pt] (8.55,3.83)--(9.2,3.83);\node[xestiny,text=Gold!80!black] at (8.88,4.05) {$x$};
  \node[xessmall,anchor=west] at (9.55,2.55) {载流子整体右移};
  \node[xessmall,anchor=west] at (9.55,1.78) {$\boldsymbol E'$ 指向左};
\end{tikzpicture}
\end{center}}

\newcommand{\twoConductorPatchFigure}{%
\begin{center}
\begin{tikzpicture}[x=1cm,y=1cm]
  \fill[DeepBlue!7] (-2,0) circle (1.1);\draw[xesblue] (-2,0) circle (1.1);
  \fill[DeepBlue!7] (2,0) circle (.85);\draw[xesblue] (2,0) circle (.85);
  \foreach \a in {110,145,180,215,250}{\node[xestiny,text=DeepBlue] at ($(-2,0)+(\a:.92)$) {$+$};}
  \foreach \a in {-55,-20,20,55}{\node[xestiny,text=DeepBlue] at ($(-2,0)+(\a:.92)$) {$+$};}
  \foreach \a in {135,165,195,225}{\node[xestiny,text=Ink] at ($(2,0)+(\a:.68)$) {$-$};}
  \foreach \a in {-48,-16,16,48}{\node[xestiny,text=DeepBlue] at ($(2,0)+(\a:.68)$) {$+$};}
  \draw[RuleGray,-{Latex[length=1.55mm]},line width=.6pt] (.82,.35)--(.25,.35);
  \node[xessmall,anchor=south] at (.53,.42) {感应增强};
  \node[xessmall] at (-2,-1.48) {$A:+Q$，表面电荷保持同号};
  \node[xessmall] at (2,-1.48) {$B:-q$，远侧可能出现正斑块};
\end{tikzpicture}
\end{center}}

\newcommand{\chargedIsolatedSphereImageFigure}{%
\begin{center}
\begin{tikzpicture}[x=1cm,y=1cm]
  \fill[DeepBlue!7] (0,0) circle (1.45);\draw[xesblue] (0,0) circle (1.45);
  \fill[Ink] (4.0,0) circle (5pt);\node[xeslabel,above=5pt] at (4,0) {$+q$};
  \fill[Ink] ({1.45*1.45/4},0) circle (4.5pt);
  \fill[DeepBlue] (0,0) circle (4.5pt);
  \node[xessmall,anchor=east] at (-1.72,-.68) {$q_0=q+qR/a$};
  \draw[RuleGray,-{Latex[length=1.45mm]},line width=.5pt] (-1.62,-.62)--(-.12,-.08);
  \node[xessmall,anchor=west,align=left] at (1.45,.76) {$q'=-qR/a$\\$b=R^2/a$};
  \draw[RuleGray,-{Latex[length=1.45mm]},line width=.5pt] (1.35,.68)--(.64,.1);
  \draw[<->,Gold!80,line width=.65pt] (0,-1.78)--(4,-1.78);\node[xestiny,text=Gold!80!black] at (2,-2.03) {$a$};
  \draw[-{Latex[length=1.7mm]},DeepBlue,line width=.8pt] (3.85,.4)--(3.1,.4);\node[xessmall,above] at (3.48,.43) {吸引};
  \draw[-{Latex[length=1.7mm]},Gold!80!black,line width=.8pt] (4.15,-.4)--(4.85,-.4);\node[xessmall,below] at (4.5,-.43) {斥力};
\end{tikzpicture}
\end{center}}

\newcommand{\cavityImageFigure}{%
\begin{center}
\begin{tikzpicture}[x=1cm,y=1cm]
  \draw[xesblue,line width=1.15pt] (0,0) circle (1.75);
  \fill[Ink] (-.78,0) circle (5pt);\node[xeslabel,below=5pt] at (-.78,0) {$-q$};
  \fill[DeepBlue] ({-1.75*1.75/.78},0) circle (5pt);\node[xeslabel,below=5pt] at ({-1.75*1.75/.78},0) {$+qR/d$};
  \fill[Ink] (0,0) circle (1.6pt);\node[xessmall,below=4pt] at (0,0) {$O$};
  \draw[<->,Gold!80,line width=.65pt] (-.78,.45)--(0,.45);\node[xestiny,text=Gold!80!black] at (-.39,.67) {$d$};
  \draw[<->,Gold!80,line width=.65pt] ({-1.75*1.75/.78},-.55)--(0,-.55);\node[xestiny,text=Gold!80!black] at (-1.97,-.78) {$R^2/d$};
  \draw[xesline] (1.75,0)--(2.35,0)--(2.35,-.55);
  \draw[xesline] (2.08,-.55)--(2.62,-.55) (2.17,-.68)--(2.53,-.68) (2.27,-.81)--(2.43,-.81);
  \node[xessmall,align=left,anchor=west] at (3.25,.48) {像电荷在球外，\\只用于球腔内部};
\end{tikzpicture}
\end{center}}

\newcommand{\cornerImagesFigure}{%
\begin{center}
\begin{tikzpicture}[x=1cm,y=1cm]
  \fill[DeepBlue!5] (0,0) rectangle (4.8,4.0);
  \draw[xesblue,line width=1.2pt] (-4.8,0)--(4.8,0) (0,-4.0)--(0,4.0);
  \fill[DeepBlue] (2.3,1.65) circle (5pt);\node[xeslabel,above right] at (2.3,1.65) {$+q$};
  \fill[Ink] (-2.3,1.65) circle (5pt);\node[xeslabel,above left] at (-2.3,1.65) {$-q$};
  \fill[Ink] (2.3,-1.65) circle (5pt);\node[xeslabel,below right] at (2.3,-1.65) {$-q$};
  \fill[DeepBlue] (-2.3,-1.65) circle (5pt);\node[xeslabel,below left] at (-2.3,-1.65) {$+q$};
  \draw[RuleGray,densely dashed] (2.3,1.65)--(-2.3,1.65) (2.3,1.65)--(2.3,-1.65);
  \draw[<->,Gold!80,line width=.6pt] (0,2.55)--(2.3,2.55);\node[xestiny,text=Gold!80!black] at (1.15,2.78) {$a$};
  \draw[<->,Gold!80,line width=.6pt] (3.15,0)--(3.15,1.65);\node[xestiny,text=Gold!80!black] at (3.38,.83) {$b$};
  \node[xessmall,anchor=west] at (5.35,2.25) {第一象限为真实区域};
  \node[xessmall,anchor=west] at (5.35,1.48) {每跨过一块接地板，像电荷变号};
\end{tikzpicture}
\end{center}}

\newcommand{\hemisphereImagesFigure}{%
\begin{center}
\begin{tikzpicture}[x=1cm,y=1cm]
  \draw[xesblue,line width=1.1pt] (-5,0)--(-1.45,0) (1.45,0)--(5,0);
  \draw[xesblue,line width=1.1pt] (-1.45,0) arc[start angle=180,end angle=0,radius=1.45];
  \fill[DeepBlue] (0,3.25) circle (5pt);\node[xeslabel,above=4pt] at (0,3.25) {$+q$};
  \fill[Ink] (0,-3.25) circle (4.5pt);\node[xeslabel,below=4pt] at (0,-3.25) {$-q$};
  \fill[Ink] (0,{1.45*1.45/3.25}) circle (4.5pt);\node[xeslabel,right=5pt] at (0,{1.45*1.45/3.25}) {$-qR/a$};
  \fill[DeepBlue] (0,{-1.45*1.45/3.25}) circle (4.5pt);\node[xeslabel,right=5pt] at (0,{-1.45*1.45/3.25}) {$+qR/a$};
  \draw[<->,Gold!80,line width=.6pt] (-.42,0)--(-.42,3.25);\node[xestiny,text=Gold!80!black] at (-.68,1.63) {$a$};
  \node[xessmall,anchor=west] at (5.55,1.35) {四电荷同时使平面与半球面电势为零};
  \node[xessmall,anchor=west] at (5.55,.55) {三个像电荷均位于真实区域之外};
\end{tikzpicture}
\end{center}}

\newcommand{\planeEnergyFigure}{%
\begin{center}
\begin{tikzpicture}[x=1cm,y=1cm]
  \begin{scope}[shift={(-3.2,0)}]
    \draw[xesblue,line width=1.05pt] (-2.15,0)--(2.15,0);
    \foreach \x in {-2.0,-1.5,...,2.0}{\draw[RuleGray,line width=.4pt] (\x,0)--(\x-.18,-.22);}
    \draw[xesarrow] (0,1.05)--(0,2.18);
    \draw[DeepBlue,line width=.65pt,densely dashed,-{Latex[length=1.55mm]}] (0,-1.05)--(0,-2.18);
    \draw[DeepBlue,line width=.7pt] (0,1.0) circle (4.2pt);\node[xestiny,anchor=west] at (.15,1.0) {$a$};
    \fill[DeepBlue] (0,2.3) circle (4.5pt);\node[xeslabel,anchor=west] at (.15,2.3) {$+q$};
    \draw[Ink,line width=.7pt] (0,-1.0) circle (4.2pt);\node[xestiny,anchor=west] at (.15,-1.0) {$-a$};
    \fill[Ink] (0,-2.3) circle (4.5pt);\node[xeslabel,anchor=west] at (.15,-2.3) {$-q$（像）};
    \node[xessmall,align=center] at (0,3.0) {缓慢移动\\真实电荷与感应分布同步调整};
  \end{scope}
  \begin{scope}[shift={(3.55,0)}]
    \draw[xesblue,line width=1.05pt] (-2.15,0)--(2.15,0);
    \foreach \x in {-2.0,-1.5,...,2.0}{\draw[RuleGray,line width=.4pt] (\x,0)--(\x-.18,-.22);}
    \draw[xesarrow] (0,1.05)--(0,2.18);
    \draw[DeepBlue,line width=.7pt] (0,1.0) circle (4.2pt);\node[xestiny,anchor=west] at (.15,1.0) {$a$};
    \fill[DeepBlue] (0,2.3) circle (4.5pt);\node[xeslabel,anchor=west] at (.15,2.3) {$+q$};
    \fill[Ink] (0,-1.0) circle (4.5pt);\node[xeslabel,anchor=west] at (.15,-1.0) {$-q$（初始像）};
    \node[xessmall,align=center] at (0,3.0) {瞬时移动\\初始感应分布暂时不变};
  \end{scope}
\end{tikzpicture}
\end{center}}

\newcommand{\uniformFieldSphereFigure}{%
\begin{center}
\includegraphics[width=.76\linewidth]{/Users/dongxingxi/Documents/物理竞赛一轮资料/sample/xes-original/restructured-complete/assets/revision-ch19-ch20/uniform-sphere.pdf}
\end{center}}


% Source question conditions restored, analysis retained separately.
\newcommand{\emCornerQuestion}{%
\begin{center}
\begin{tikzpicture}[x=1cm,y=1cm]
  \fill[DeepBlue!5] (0,0) rectangle (4.8,4.0);
  \draw[xesblue,line width=1.2pt] (0,0)--(4.8,0) (0,0)--(0,4.0);
  \fill[DeepBlue] (2.3,1.65) circle (5pt);\node[xeslabel,above right] at (2.3,1.65) {$+q$};
  \draw[<->,Gold!80,line width=.6pt] (0,2.55)--(2.3,2.55);\node[xestiny,text=Gold!80!black] at (1.15,2.78) {$a$};
  \draw[<->,Gold!80,line width=.6pt] (3.15,0)--(3.15,1.65);\node[xestiny,text=Gold!80!black] at (3.38,.83) {$b$};
\end{tikzpicture}
\end{center}}

% Source question conditions restored, analysis retained separately.
\newcommand{\emHemisphereQuestion}{%
\begin{center}
\begin{tikzpicture}[x=1cm,y=1cm]
  \draw[xesblue,line width=1.1pt] (-5,0)--(-1.45,0) (1.45,0)--(5,0);
  \draw[xesblue,line width=1.1pt] (-1.45,0) arc[start angle=180,end angle=0,radius=1.45];
  \draw[xesline] (0,0)--(35:1.45) node[midway,below right]{$R$};
  \fill[DeepBlue] (0,3.25) circle (5pt);\node[xeslabel,above=4pt] at (0,3.25) {$+q$};
  \draw[<->,Gold!80,line width=.6pt] (-.42,0)--(-.42,3.25);\node[xestiny,text=Gold!80!black] at (-.68,1.63) {$a$};
\end{tikzpicture}
\end{center}}


第19章分别从力和能量两个角度考察了电场，引入了描述电场的矢量场
$\boldsymbol E$ 与标量场 $\varphi$。电场强度可以用于一般电磁问题，电势则专用于静电场。电场线和
等势面虽然是图形化的表达，却能把许多空间关系直接呈现出来。此前主要讨论电场对少数点电荷的作用，
本章开始研究电场与宏观物质的相互作用。

所有宏观物质都由原子构成，原子由带正电的原子核与核外电子组成。把物质置于外电场中，正、负电荷会
受到方向相反的作用；物质内部电荷能否移动，决定了它对外电场的响应方式。


\subsection{导体和绝缘体}


孤立原子中的电子受原子核束缚。许多原子聚集成固体以后，每个电子不仅受到最近原子核的作用，也受到
周围大量原子的影响。某些距离原子核较远的电子不再局限于单个原子，而可以在整个固体中运动。准确描述
这一现象需要量子理论；在宏观静电学中，可以把导体理解为含有大量可移动载流子的物质，把绝缘体理解为
电荷在外场作用下仍基本束缚在原处的物质。

金属是典型导体。向金属传递一些电荷后，电荷不会固定在接触点附近，而会在极短时间内重新分布到导体
表面。要让导体长期带电，必须用绝缘体把它与其他导电区域隔开，否则电荷会迅速流失。玻璃、橡胶和干燥
木材是常见绝缘体；它们经过摩擦后，电荷可以长时间停留在局部区域。

上述“自由电子”和“完全束缚电子”的图像是简化模型。真实材料的导电性还随温度、杂质、场强和频率
变化；绝缘体在足够强的电场下也会击穿。但这一模型已经足以建立静电平衡的主要规律。

应当进一步说明，这里所谓“自由”并不是电子完全不受原子核和晶格作用。电子在固体中始终处于所有离子
实产生的周期势场里，还会与晶格振动、杂质和其他电子碰撞。导体与绝缘体的本质差别，是材料中是否存在
能够在很小外场下改变占据状态并形成宏观电流的载流子。金属中的载流子数目很大，所以外场一出现，表面
电荷就能在极短时间内完成重新分布；绝缘体中的电荷只能产生很小的束缚位移，通常表现为极化而不是持续
导电。半导体处在二者之间，其载流子数目可被温度、光照或掺杂显著改变。

导体得到的多余电荷会遍布它的可达表面，而不会停在最初接触的位置。这里的“表面”包括外表面，也可能
包括含有内部电荷的空腔表面；究竟分布到哪一部分，要由静电平衡条件共同决定。若导体与另一个导体直接
接触，它们将构成一个新的整体等势体，电荷继续移动，直到整个连通导体达到共同电势。因此实验中常用
绝缘支架、绝缘手柄或真空把带电导体与环境隔离。


\subsection{静电平衡}


把导体放入预先存在的外电场，内部载流子立即受力并发生定向移动。载流子到达导体边界后不能继续离开，
于是表面出现净电荷。这些表面电荷又产生电场，并与外电场叠加。只要导体内部总电场不为零，电荷分布就
还会继续改变；最终导体内部处处无场，宏观状态不再改变，这就是\textbf{静电平衡}。移动后留在表面的净
电荷称为\textbf{感应电荷}。

图~\ref{ch20:fig:s34-3-1}以一个带正电导体和一个净电荷为零的导体为例。中性导体靠近带电体的一侧出现负感应电荷，远侧
出现正感应电荷；电场线从正电荷出发，终止于负电荷或延伸到无穷远。导体表面本身是等势面，导体内部
没有电场线。

\par\noindent\begin{minipage}{\linewidth}
\conductorFieldFigure
\diagramcaption{静电平衡状态下两个导体的电场线、等势线与感应电荷分布。以两球为例，实线为场线、虚线为等势线；电场线垂直球面，延伸到画框外的线并非在框边终止。}\label{ch20:fig:s34-3-1}
\end{minipage}\par\medskip

根据形成静电平衡的过程，可以得到以下性质。

\begin{enumerate}[label=\arabic*.]
  \item \textbf{导体材料内部电场处处为零。}如果某处 $\boldsymbol E\ne0$，可移动电荷仍会受力，电荷分布
  还会变化，因而不可能处于静电平衡。

  \item \textbf{导体材料内部没有净体电荷。}在导体内部取任意闭合高斯面，因为面上
  $\boldsymbol E=0$，电通量为零，所包围的净电荷也必为零。因此导体的多余电荷只能分布在表面。

  \item \textbf{整个导体是等势体，表面是等势面。}导体内部任意两点间的电势差为
  $-\int\boldsymbol E\cdot\dd\boldsymbol l=0$。把导体连接到大地称为接地；大地可吸收或提供大量电荷，
  通常规定接地导体的电势为零。

  \item \textbf{导体表面外侧的电场垂直于表面。}若存在切向分量，表面电荷会沿导体表面流动。取一个
  横跨导体表面的极薄柱形高斯面，导体内侧电场为零，得到表面外侧法向电场与面电荷密度的关系
  \begin{equation}
    \boxed{E_n=\frac{\sigma}{\varepsilon_0}}.
  \end{equation}
  其中 $E_n$ 以导体外法线方向为正。$\sigma>0$ 时电场指向导体外，$\sigma<0$ 时指向导体内。
\end{enumerate}

图~\ref{ch20:fig:s34-3-2}给出用静电感应使导体带正电的完整顺序。起初导体呈电中性；把带负电物体靠近导体，导体左侧聚集
正电荷、右侧聚集负电荷；保持负电体不动并将导体接地，电子受排斥流入大地，导体上只剩净正电荷；先
断开地线，导体上的正电荷被保留下来；最后移走负电体，正电荷重新分布到整个外表面。顺序不能颠倒：若
先移走带电体再断开地线，电荷会经地线回到大地，导体仍为中性。

\par\noindent\begin{minipage}{\linewidth}
\inductionSequenceFigure
\diagramcaption{一个完整的静电感应带电过程。图中负电棒始终位于导体左侧，接地只发生在步骤 (c)。}\label{ch20:fig:s34-3-2}
\end{minipage}\par\medskip

感应电荷的总量和局部分布必须分开理解。未接地的中性导体在感应过程中总电荷始终为零，但近端、远端可
分别出现等量异号电荷；接地以后，地线允许电荷进入或离开，导体总电荷不再固定。断开地线只固定导体的
总电荷，并不固定表面电荷分布；当外部带电体继续移动时，表面电荷仍会沿导体重新排列。正因如此，静电
感应问题不能只画几个“正、负号”就结束，还必须同时检查总电荷守恒、导体电势条件和表面法向电场方向。

接地也不能简单理解为“把所有电荷都导走”。它真正固定的是导体电势。为了使电势达到零，大地可能吸收
电子，也可能向导体提供电子，方向取决于外部电荷的符号和几何位置。图~\ref{ch20:fig:s34-3-2}中负电棒排斥电子，所以电子
从导体流向大地；若靠近的是正电棒，过程中的电荷流向和最后带电符号都将反转。

静电平衡时，同一导体表面不可能有一条电场线从正电荷区域出发、再终止于该导体的负电荷区域。因为沿
电场线电势必须下降，而同一导体各点电势相等。把这一结论用于内部无电荷的封闭导体壳：空腔内表面任意
两点之间不能由电场线连接；又由包围空腔的高斯面可知内表面总电荷为零，于是空腔内表面处处无电荷，
空腔内电场为零。这就是静电屏蔽。

法拉第笼是静电屏蔽的典型应用。完整金属外壳可以阻止外部静电场进入内部；网状金属笼对变化较慢、波长
远大于网孔的电磁场也有明显屏蔽作用。屏蔽结构可保护敏感仪器免受射频干扰，也可以把无线电发射设备的
干扰限制在内部。遭遇雷击或静电放电时，电流主要沿导体笼外表面传导，内部人员和设备受到保护。

图~\ref{ch20:fig:s34-3-3} 的金属网构成连通屏蔽壳；在上述静电条件下，壳内无电荷的空腔与外界电场隔离。

\par\noindent\begin{minipage}{\linewidth}
\faradayCageFigure
\diagramcaption{法拉第笼的屏蔽情景。理想静电极限下腔内无场；雷击瞬态的安全还依赖连接、接地及屏蔽完整性。}\label{ch20:fig:s34-3-3}
\end{minipage}\par\medskip


% EM teaching enrichment 20260927
\paragraph{带电空腔：把内表面与外表面分开}
中性孤立导体壳内放入电荷 $q$，在导体材料内取包围空腔的高斯面，因 $E=0$，有
$q+Q_{\rm inner}=0$。再由壳的净电荷守恒得 $Q_{\rm outer}=q$。这两步只决定总量；内电荷偏心时内表面电荷一般不均匀。
若外表面为球面、外界无其他电荷，外部解由球形等势边界确定为球心点电荷 $q$ 的场，与空腔内电荷偏心方向无关。
若壳接地且外界无其他电荷，外表面可以不带电，外部场为零；此时壳总电荷已由地线改变，不能再使用原来的中性约束。


\subsection{等离子体频率}


导体最终达到静电平衡，不仅因为内部存在可移动载流子，还因为载流子的运动会耗散能量。竖直弹簧下突然
挂上重物，若没有阻力，重物不会停在新的静力平衡位置，而会围绕它持续振动；只有阻尼耗散机械能，重物
才会最终静止。导体中的载流子也一样：若完全没有碰撞和电阻，它们不会停在屏蔽外场的电荷分布上，而会
围绕电中性位置集体振荡。

把弹簧的过程说得更具体些：重物刚挂上时，原来的弹簧长度不是新的静力平衡位置。重力先使重物向下加速，
弹簧伸长后恢复力增大；重物经过平衡位置时仍有速度，于是继续向下，随后被拉回。没有空气阻力和弹簧内耗
时，这一往复过程永不停止。导体中载流子的惯性对应重物质量，表面电荷产生的恢复场对应弹簧恢复力，碰撞
和电阻则对应阻尼。通常金属中的碰撞时间很短，因此宏观观察只看到迅速建立的静电平衡，而看不到长时间
的自由振荡。

设导体中可移动载流子的质量为 $m$、电荷为 $q$、数密度为 $n$。外电场 $E_0$ 沿 $+x$ 方向突然加上，
载流子整体位移 $x$ 后，样品两侧出现面电荷，大小为 $\sigma=nqx$。这些面电荷在内部产生恢复场
\[
  E'=-\frac{nq}{\varepsilon_0}x.
\]

图中以正载流子向右位移表示正负面电荷的几何关系。
若载流子是电子，$q=-e$，在向右外场下平衡位移 $x_0<0$，电子向左；
两端面的符号仍应从乘积 $nqx$ 判断，不能仅由位移方向猜测。
\plasmaSlabFigure

忽略碰撞和阻尼，每个载流子的运动方程为
\[
  qE_0-\frac{nq^2}{\varepsilon_0}x=m\ddot x.
\]
整理为
\begin{equation}
  \boxed{\ddot x+\frac{nq^2}{m\varepsilon_0}x=\frac{qE_0}{m}}.
\end{equation}
这是围绕平衡位置 $x_0=\varepsilon_0E_0/(nq)$ 的简谐振动，固有角频率为
\begin{equation}
  \boxed{\omega_p=\left(\frac{nq^2}{m\varepsilon_0}\right)^{1/2}}.
\end{equation}

若平均阻力为 $-m\nu\dot x$，令 $y=x-x_0$，则
\[
 \ddot y+\nu\dot y+\omega_p^2y=0,\qquad
 s_\pm=\frac{-\nu\pm\sqrt{\nu^2-4\omega_p^2}}2.
\]
当 $0<\nu<2\omega_p$ 时为衰减振荡，振幅包络为 $e^{-\nu t/2}$；
$\nu=2\omega_p$ 时为临界阻尼；$\nu>2\omega_p$ 时为非振荡的过阻尼趋近。
特别是 $\nu\gg\omega_p$ 时，慢根 $s_+\simeq-\omega_p^2/\nu$，
弛豫时间约为 $\nu/\omega_p^2$，所以不能笼统说“碰撞越强，平衡越快”。
弱阻尼而外场交变时，在适当驱动下可于 $\omega\approx\omega_p$ 附近出现强响应。因此“静电平衡”和“等离子体振荡”并不矛盾：前者是有耗散、长时间极限，后者是惯性
不可忽略时的动力学行为。

等离子体是离子和电子组成、能够表现出集体行为的带电粒子体系。离子质量远大于电子，最简单近似下可把
离子视为静止背景，只考虑电子相对离子的振荡；此时带符号电荷为 $q=-e$、$m=m_e$，频率中的 $q^2=e^2$，$\omega_p$ 主要由电子数密度决定。
电子偏离电中性位置后，电荷分离产生恢复场；电子越过平衡位置时仍有惯性，于是形成等离子体振荡。

等离子体频率公式还说明等离子体频率只由载流子数密度、电荷量和质量决定，与样品在理想无限均匀近似下的横向
尺寸无关。电子的质量很小，所以电子等离子体频率通常很高；离子也可以振荡，但其特征频率低得多。在
有限等离子体中还会出现边界效应、温度效应和不同传播方向的多种波模，等离子体频率公式给出的是最基本的均匀纵向
集体振荡尺度。

地球约 $60\,\mathrm{km}$ 以上的高层大气受到太阳高能辐射和宇宙线作用而部分电离，形成电离层。入射电磁波的频率低于电离层的等离
子体频率时，电子能够集体响应并使电磁波被反射；频率足够高时，电磁波可以穿过电离层。短波通信可借助
电离层一次或多次反射跨越地平线，而卫星通信需要采用能够穿透电离层的较高频率。

图~\ref{ch20:fig:s34-3-6} 中橙色路径表示被电离层折回的通信波，蓝色路径穿过电离层到达卫星；两条路径对应不同的频率条件。

\par\noindent\begin{minipage}{\linewidth}
\ionosphereFigure
\diagramcaption{电离层反射实现远距离通信。图中同时区分低频反射与高频穿透。}\label{ch20:fig:s34-3-6}
\end{minipage}\par\medskip

现实导体中的碰撞使载流子振荡衰减，系统才会趋于静电平衡；在高频电磁场或等离子体中，惯性和集体振荡
则不能忽略。静电平衡与等离子体频率看似分属静力学和动力学，实际是同一套电荷重新分布机制在有、无
耗散两种极限下的表现。

% EM teaching enrichment 20260927
\paragraph{把振荡模型接回电荷弛豫}
在强碰撞极限，线性碰撞模型给出（电导率的微观推导见第21章） $\sigma=nq^2/(m\nu)$，从而
\[
 \frac{\nu}{\omega_p^2}=\frac{\varepsilon_0}{\sigma}.
\]
这正是忽略惯性后，由电荷连续性与 $\boldsymbol j=\sigma\boldsymbol E$ 得到的体电荷弛豫时间：
$\partial_t\rho=-\nabla\cdot\boldsymbol j=-\sigma\rho/\varepsilon_0$。
二者相符须满足强碰撞近似，不能把 $\varepsilon_0/\sigma$ 生硬解释为任意金属中电子惯性运动的真实完成时间。静电模型的使用应比较观测时间与材料响应时间，而不只是看电源是否标为直流。




% EM structure and teaching revision 20260928
\subsection{接地、孤立与相连导体的条件选择}
导体内部场为零只给出导体等势，还不足以决定它的电势值或总电荷。孤立导体给定总电荷 $Q$，其电势待求；接地导体给定电势，电荷可以经地线流入流出；导线连接的多个导体电势相等，但通常只守恒它们的总电荷。
例如两个相距很远的金属球，半径为 $R_1,R_2$，忽略相互影响。用细导线连接后
\[
 \frac{Q_1}{R_1}=\frac{Q_2}{R_2},\qquad Q_1+Q_2=Q,
 \qquad Q_i=\frac{QR_i}{R_1+R_2}.
\]
球面场强却满足 $E_1/E_2=R_2/R_1$；等势不意味着等电荷或等场强。若两球靠得很近，各自电势还包含对方电荷的贡献，上述孤立球近似便不再成立。
在有空腔的问题中，先用导体材料内部的高斯面确定内表面总感应电荷，再按接地或孤立条件确定外表面电荷。偏心电荷通常使内表面分布不均匀，却不妨碍用高斯定律求其总量。把“总电荷已知”与“分布已知”分开，是避免误用球对称公式的关键。


\subsection{专题例题与方法应用}

\begin{origexamplebox}{【异号导体】}\label{ch20:ex:s34-3-1}
小明使导体 $A,B$ 分别带电 $+Q,-q$，若 $Q>q$。由于二者会产生静电感应，问导体 $A$ 上是否有可能
出现负电荷，导体 $B$ 上是否有可能出现正电荷？若 $Q=q$，结果有无变化？
\end{origexamplebox}\par\nopagebreak[4]

\begin{xesolution}
题目只给了净电荷，没有给导体形状、尺寸和间距，因此应先区分“必然结论”和“取决于几何的结论”。

静电平衡时每个导体各有一个常电势 $V_A,V_B$。在导体外部，电势满足拉普拉斯方程。若某导体处在整个
外部区域的最高电势，它的表面向外离开时电势只能降低，因此
$E_n=-\partial\varphi/\partial n\ge0$，由 $\sigma=\varepsilon_0E_n$ 可知该导体表面不可能出现负电荷。

当 $Q>q$ 时，系统总电荷 $Q-q>0$，正电导体 $A$ 是两导体中的高电势者；否则若净负电的 $B$ 处在最高
电势，最大值原理将迫使它表面处处 $\sigma\ge0$，与其净电荷为负矛盾。因此
\[
  \boxed{A\text{ 表面不出现负电荷}}.
\]
$B$ 的净电荷为负，所以它至少有一部分表面带负电；但若 $A$ 的影响足够强，$B$ 远离 $A$ 的一侧可以出现
正电荷斑块。因此
\[
  \boxed{B\text{ 表面可能出现正电荷}\quad\text{但并非必然}}.
\]

\twoConductorPatchFigure

当 $Q=q$ 时系统总电荷为零。只有两个导体时，$V_A>0>V_B$：$A$ 是电势最大边界，$B$ 是电势最小边界。
最大值原理和最小值原理分别给出
\[
  \boxed{\sigma_A\ge0,\qquad \sigma_B\le0}.
\]
也就是说，$A$ 不出现负电荷，$B$ 不出现正电荷；净电通量从 $A$ 流出并由 $B$ 流入。尖角处面电荷密度
可以很大，但符号不变。以上结论默认只有这两个导体、无额外外场，也不讨论封闭空腔内另置电荷的特殊拓扑。
\end{xesolution}



\begin{origexamplebox}{【同号导体】}\label{ch20:ex:s34-3-2}
小明使导体 $A,B$ 分别带电 $+Q,+q$，若 $Q>q$。由于二者会产生静电感应，问导体 $A$ 上是否有可能
出现负电荷，导体 $B$ 上是否有可能出现负电荷？若 $Q=q$，结果有无变化？若有更多的球（比如 100 个），
能得出什么定性结论？
\end{origexamplebox}\par\nopagebreak[4]

\begin{xesolution}
两个导体净电荷都为正，并不意味着每个表面点都一定带正电。决定局部面电荷符号的是该点外侧法向电场，
而决定法向电场的是两个导体的电势和几何形状。

设 $V_A,V_B>0$。电势较高的导体是外部区域的最大电势边界，表面必有 $\sigma\ge0$；电势较低的导体
面对高电势导体时，离开其表面电势可能先升高，于是 $E_n<0$，局部可以出现负电荷。不过它的净电荷仍为
正，所以其余表面必须有更多正电荷。

仅由 $Q>q$ 不能判断 $V_A$ 与 $V_B$ 的大小，因为电势还取决于导体电容，也就是尺寸、形状和相互位置。
因此对任意导体，能确定的是
\[
  \boxed{\text{电势较高者不出现负电荷；电势较低者可能出现负电荷斑块。}}
\]
若 $A,B$ 是半径相同、环境对称的两球，则 $Q>q$ 会给出 $V_A>V_B$；此时 $A$ 表面全为正电荷，$B$
在两球很近时可能于迎向 $A$ 的一面出现负电荷。

对一般导体，即使 $Q=q$，尺寸、形状不同仍可使 $V_A\ne V_B$，所以电势较低者仍可能出现负电荷斑块；等量电荷本身并不排除这种现象。

当两个相同球带等量正电 $Q=q$ 且放置对称时，$V_A=V_B$。把两球视作处于同一电势的复合边界，外部
电势从正值降到无穷远的零，因此两球表面均有 $\sigma\ge0$，球间也没有从一个球终止到另一个球的电场线。

推广到 100 个导体：至少有一个处于最高电势的导体，其表面不可能出现负电荷；较低电势导体可能出现
局部负电荷。若 100 个完全相同的球在对称环境中带等量电荷，使它们电势相同，则所有球表面均不出现负
电荷。若几何和电荷量不对称，不能仅凭“每个球净带正电”推出每处面电荷都为正。
\end{xesolution}



\begin{origexamplebox}{【凭空相吸】}\label{ch20:ex:s34-3-3}
小明梦到两个金属球本身都不带电。由于小扰动，某一个球体内的电荷分布发生了一些变化，进而两个球互相
静电感应、局部带电并相互吸引。醒来后小明觉得这个现象不对，却说不清哪里不对。请解释。
\end{origexamplebox}\par\nopagebreak[4]

\begin{xesolution}
两个球都不带净电荷，且没有外加电荷和外场。电势函数
\[
  \varphi(\boldsymbol r)=0
\]
满足全部条件：球外无体电荷、无穷远电势为零、两个球各为等势体且各自总电荷为零。唯一性定理说明这就
是唯一的静电平衡解，所以
\[
  \boxed{\boldsymbol E=0,\qquad \sigma_A=\sigma_B=0}.
\]

从能量角度看也一样。任何自发的正负电荷分离都会建立非零电场，增加
$\int\varepsilon_0E^2/2\,\dd V>0$。系统没有外界能量来源，不会从能量最低的无场状态自动跳到更高能量
状态。短暂的热涨落或量子涨落会很快衰减，不能形成持续的经典静电吸引。

现实中的两个“中性物体相吸”可能来自外部带电体诱导、接触电势、残余电荷，或微观尺度的色散力；这些
都增加了题目梦境中没有给出的物理条件，并不是两个理想中性导体凭空产生稳定静电感应。
\end{xesolution}



\section{唯一性定理与电像法}

\subsection{唯一性定理}


定性规律可以判断静电平衡的许多特征；要确定空间中电势和电场的具体分布，还需要静电场的唯一性定理。
由
$\boldsymbol E=-\boldsymbol\nabla\varphi$ 与 $\boldsymbol\nabla\cdot\boldsymbol E=\rho/\varepsilon_0$
联立，得到
\begin{equation}
  \boxed{\frac{\partial^2\varphi}{\partial x^2}
  +\frac{\partial^2\varphi}{\partial y^2}
  +\frac{\partial^2\varphi}{\partial z^2}
  =-\frac{\rho}{\varepsilon_0}}.
\end{equation}
这就是泊松方程。无电荷区域中 $\rho=0$，它退化为拉普拉斯方程
$\nabla^2\varphi=0$。

中学阶段不需要直接求解一般的三维偏微分方程，但离散化可以直观说明为什么必须给出边界条件。设电荷和
电势只随 $x$ 变化，则
\begin{equation}
  \boxed{\frac{\dd^2\varphi}{\dd x^2}=-\frac{\rho(x)}{\varepsilon_0}}.
\end{equation}

离散化的思想与测量城市温度很相似。温度本来是位置的连续函数，但实际不可能在每一个数学点都放置温度
计；只要在足够密的观测站取得温度，再利用相邻点之间的变化，就能近似恢复整座城市的温度分布。电势也
一样：先在有限个节点上求 $\varphi_n$，节点越密，离散结果越接近连续函数。泊松方程的离散式不是
额外的物理假设，而是连续二阶导数的差分近似。
在 $A,B$ 之间取间距均为 $a$ 的离散点，点 $n$ 处的电势为 $\varphi_n$。其右、左两侧一阶导数分别近似为
$\bigl(\varphi_{n+1}-\varphi_n\bigr)/a$ 与
$\bigl(\varphi_n-\varphi_{n-1}\bigr)/a$，因此二阶导数的离散形式为
\begin{equation}
  \boxed{\frac{\varphi_{n+1}+\varphi_{n-1}-2\varphi_n}{a^2}
  =-\frac{\rho_n}{\varepsilon_0}}.
\end{equation}

图~\ref{ch20:fig:s34-3-4} 中节点 $n$ 与两侧相邻节点均相距 $a$，所以其二阶差分恰由上式三项构成；$A,B$ 的电势不再作为未知量。

\par\noindent\begin{minipage}{\linewidth}
\discretePotentialFigure
\diagramcaption{把 $A,B$ 之间连续的电势分布离散化。相邻节点间距均为 $a$，端点电势是边界条件。}\label{ch20:fig:s34-3-4}
\end{minipage}\par\medskip

对每一个内部节点都列一个方程，就得到包含大量未知电势的线性方程组。内部节点的电荷密度必须已知；对
最靠近两端的节点，方程还会包含端点 $A,B$ 的电势，因此边界电势也必须给定。令 $a\to0$，离散点无限
细密，便回到连续的泊松方程。

静电场唯一性定理可以表述为：给定区域内的电荷分布，并在边界上给定适当条件，区域内的电势和电场便
唯一确定。常见边界条件有三类：边界电势已知；边界法向电场已知；一部分边界给电势、另一部分边界给
法向电场。导体是等势体，所以“某个导体电势已知”本身就是一种边界条件；若只给导体的总电荷，还需用
总电荷约束共同确定导体电势。

为什么内部电荷和边界条件缺一不可，可以直接从差分方程看出。对 $n=2,3,\ldots,N-2$，方程右侧由
$\rho_n$ 给定；但对最靠近左端的 $n=1$，方程包含端点未知量 $\varphi_0$，对最靠近右端的
$n=N-1$，方程包含 $\varphi_N$。若不给 $A,B$ 的电势，未知数就比独立方程多，整体可以发生不受限制的
改变。三维问题虽然节点更多、邻接关系更复杂，逻辑完全相同：内部的泊松方程规定局部曲率，边界
条件负责把所有局部关系闭合成唯一的全局解。

若给定的是边界法向场而非边界电势，由高斯定律还必须满足总通量与区域内总电荷相容。纯诺依曼
边界条件只决定电势到一个任意加法常量，但这个常量不影响电场；混合边界条件则把两类信息分别施加在
不同边界部分。这三种分支条件，正是狄利克雷、诺依曼和混合边界条件在静电学中的对应。

唯一性定理的力量在于：只要构造出一个满足相同区域电荷分布和相同边界条件的电势函数，就不必再去求
复杂的感应电荷分布；这个构造出来的解就是唯一的物理解。


\subsection{电像法}


静电平衡问题通常同时包含给定电荷和导体表面上未知的感应电荷。直接积分感应电荷非常困难。若能在所
研究区域之外放置若干假想点电荷，使它们与真实电荷共同产生的电势恰好满足导体边界条件，那么由唯一性
定理可知，假想电荷在研究区域内的作用与真实感应电荷完全等效。这就是\textbf{电像法}。像电荷不是真实
电荷，只能放在研究区域之外，不能把它当成物体的一部分去计算受力。

电像法的经典例子是：半径为 $R$ 的接地导体球，球心为 $C$；球外轴线上距球心 $L>R$ 处放一点电荷
$+Q$。在球内同一轴线上、距球心 $l$ 处放一个像电荷 $-q$。只要整个球面电势为零，球外电场就与真实
问题完全相同。取轴线上球面的近、远两个点 $a,b$，有
\[
  \frac{1}{4\pi\varepsilon_0}
  \left(\frac{Q}{L-R}-\frac{q}{R-l}\right)=0,
  \qquad
  \frac{1}{4\pi\varepsilon_0}
  \left(\frac{Q}{L+R}-\frac{q}{R+l}\right)=0.
\]
解得
\begin{equation}
  \boxed{q=\frac RLQ,\qquad l=\frac RL R=\frac{R^2}{L}}.
\end{equation}

这里只在球面上选取两个点列方程，看起来似乎不足以保证整个球面电势都为零。关键在于所得参数同时满足
$Ll=R^2$ 和 $q/Q=R/L$。对球面任意点 $P$，三角形几何给出
$r_1/r_2=L/R$，从而 $Q/r_1=q/r_2$；所以任意 $P$ 的电势都为零。近、远两点的方程只是求出候选参数，
任意点的距离比关系才完成对整个边界条件的验证。电像法中的每一次构造都应做这样的全边界检查，不能只
验证图上少数特殊点。

威廉·汤姆孙（后来获封开尔文勋爵）系统发展了这一方法。它之所以巧妙，是因为我们不必知道每一
小块球面上到底聚集了多少感应电荷，只需找到一个位于禁入区域内的简单点电荷系统，使球外电势满足完全
相同的方程和边界条件。像电荷的位置落在导体内部，那里不是所求球外区域，因此像电荷产生的奇点不会
污染真实解。

图~\ref{ch20:fig:s34-3-5} 中 $C$ 是球心，$P$ 严格位于球面；$CP=R$，真实电荷与像电荷到 $C$ 的距离分别为 $L,l$。

\par\noindent\begin{minipage}{\linewidth}
\groundedSphereImageFigure
\diagramcaption{接地导体球的电像法几何。真实电荷为 $+Q$，像电荷为 $-q$，且 $Ll=R^2$。}\label{ch20:fig:s34-3-5}
\end{minipage}\par\medskip

对球面上任意点 $P$，几何关系保证
$r_1/r_2=L/R=Q/q$，所以 $Q/r_1-q/r_2=0$；球面确为零等势面。像电荷的
位置和大小确定后，外电荷受到的吸引力等于它与像电荷之间的库仑力：
\[
  F=\frac{1}{4\pi\varepsilon_0}\frac{Qq}{(L-l)^2}.
\]
外电荷与接地球组成系统的静电势能不是两个真实点电荷势能，而是
\[
  U=\frac12Q\varphi_{\rm ind}(Q)
  =-\frac12\frac{1}{4\pi\varepsilon_0}\frac{Qq}{L-l}.
\]
因子 $1/2$ 来自导体感应电荷随真实电荷逐步建立；若把像电荷误当真实电荷，这个因子会被遗漏。

把真实电荷与像电荷的场在球外叠加，在球面处取法向分量，再用 $\sigma=\varepsilon_0 E_n$，得到距真实电荷为 $r_1$ 的球面
点上的感应电荷密度
\begin{equation}
  \boxed{\sigma=-\frac{1}{4\pi}\frac{Q(L^2-R^2)}{Rr_1^3}}.
\end{equation}
这里已经约去 $\varepsilon_0$；结果的单位为电荷面密度。球面各处 $\sigma<0$，且离 $Q$ 越近，$r_1$
越小，负感应电荷密度的绝对值越大。


% EM teaching enrichment 20260927
\paragraph{先用平面模型检验力与能量}
电荷 $q$ 距无限接地平面为 $a$，像电荷为镜像位置的 $-q$，二者相距 $2a$。
真实电荷受力大小 $kq^2/(4a^2)$，指向平面；系统与感应有关的势能却为
\[
 U(a)=\frac12q\varphi_{\rm ind}=-\frac{kq^2}{4a},\qquad
 F_a=-\frac{\dd U}{\dd a}=-\frac{kq^2}{4a^2}.
\]
求力时像电荷场可直接使用；求位置导数时像的位置和感应分布随 $a$ 一起改变。像法只在导体外替代感应场，像电荷所在处不是一个真实可访问的受力粒子。


% EM structure and teaching revision 20260928
\subsection{电像构造的验收与能量系数}
电像法的替代电荷只能放在不求解的区域。构造完成后依次检查：真实电荷位置的奇性、导体表面的等势条件、远处或其他边界条件。满足这些条件后，唯一性定理才保证指定区域内的电势正确；替代区域内部的场没有对应的物理意义。
点电荷 $q$ 距无限接地平面 $a$，像电荷为平面另一侧的 $-q$。真实电荷所受吸力大小为 $kq^2/(4a^2)$，缓慢从远处移到 $a$ 时体系能量变化为
\[
 W(a)=\int_{\infty}^{a}\frac{kq^2}{4s^2}\dd s
 =-\frac{kq^2}{4a}.
\]
这里外力沿离开平面的方向，故外力沿坐标 $s$ 的分量为正。若把真电荷和像电荷误当成两个独立的真实点电荷，会写出 $-kq^2/(2a)$，多出一倍。
两种处理的区别在于真实导体表面的感应电荷会随 $q$ 的移动重新分布。在线性接地边界下，组装给出 $W=\tfrac12q\varphi_{\mathrm{ind}}$；有另加外场或固定偏置时，外加场能与感应能应分别计账，不能把所有电势项都一律乘以二分之一。


\subsection{专题例题与方法应用}

\begin{origexamplebox}{【同性相吸】}\label{ch20:ex:s34-3-4}
把一点电荷 $+q$ 放在距球心 $a$ 处，旁边是半径为 $R$、净电荷也为 $+q$ 的孤立金属球，$a>R$。竟然
可能出现同性相吸。请定量解释。
\end{origexamplebox}\par\nopagebreak[4]

\begin{xesolution}
孤立带电球与接地球的差别是：除了保证球面等势，还必须保证球上总电荷为 $+q$。先放置接地球所需的像
电荷
\[
  q'=-\frac{qR}{a},\qquad b=\frac{R^2}{a},
\]
再在球心放一个像电荷
\[
  q_0=q-q'=q+\frac{qR}{a}.
\]
球心电荷在球面上只增加一个常量电势，不破坏“球面等势”；两个像电荷的代数和为球的真实净电荷 $+q$。

\chargedIsolatedSphereImageFigure

取远离球心的方向为正。真实电荷受到球心像电荷的斥力，以及偏心负像电荷的吸力：
\[
  F=k\frac{qq_0}{a^2}-k\frac{q|q'|}{(a-b)^2}
  =kq^2\left[\frac1{a^2}+\frac{R}{a^3}
  -\frac{Ra}{(a^2-R^2)^2}\right],
  \qquad k=\frac1{4\pi\varepsilon_0}.
\]
前两项表示净正电荷带来的斥力，最后一项表示近侧强烈集中的负感应电荷带来的吸力。令 $x=a/R$，临界条件
$F=0$ 化为
\[
  (x-1)^2(x+1)^3=x^4.
\]
大于 1 的解为黄金比 $x=(1+\sqrt5)/2$。因此
\[
  \boxed{R<a<\frac{1+\sqrt5}{2}R\quad\Longrightarrow\quad F<0\quad\text{两者相吸}},
\]
而 $a>(1+\sqrt5)R/2$ 时净作用为斥力。距离很近时，负感应电荷集中在球的近端，其吸力随间隙减小增长得
比远处分散正电荷的斥力快，所以即使两物体净电荷同号，仍能相吸。
\end{xesolution}



\begin{origexamplebox}{【球内屏蔽】}\label{ch20:ex:s34-3-5}
用一个内半径为 $R$、外半径为 $R_o>R$ 的同心接地金属球壳把一点电荷 $-q$ 封在球腔内，点电荷距球心为 $x$。求：
\begin{enumerate}[label=(\arabic*)]
  \item 空间中的电场分布；
  \item 感应电荷的面密度分布；
  \item 若球壳不接地，而是净电荷为零，结果怎样变化；
  \item 为什么电荷在壳外时有“静电屏蔽”，在壳内时却不能使腔内无场？
\end{enumerate}
\begin{center}\begin{tikzpicture}[>=Latex,scale=.8]

 \fill[SoftGray,even odd rule](0,0)circle(2.1)(0,0)circle(1.5);
 \draw[thick](0,0)circle(2.1)(0,0)circle(1.5);
 \fill(-.7,0)circle(3pt)node[below=3pt]{$-q$};\fill(0,0)circle(1pt)node[below right]{$O$};
 \draw[<->](-.7,.45)--(0,.45)node[midway,above]{$x$};
 \draw(0,0)--(60:1.5)node[midway,right]{$R$};\draw(0,0)--(240:2.1)node[pos=.45,right=3pt]{$R_o$};
 \draw(2.1,0)--(3,0)--(3,-.65);\draw(2.65,-.65)--(3.35,-.65)(2.75,-.8)--(3.25,-.8)(2.9,-.95)--(3.1,-.95);

\end{tikzpicture}\end{center}
\figcap{接地同心金属球壳与腔内偏心点电荷。}\label{ch20:fig:em-source-thirty-fourth-section-5}
\end{origexamplebox}\par\nopagebreak[4]

\begin{xesolution}
记 $d=x<R$，取 $z$ 轴从球心指向真实电荷，真实电荷的位置矢量为 $\boldsymbol d$。为使半径 $R$ 的内球面
电势为零，在同一轴线上球外放像电荷
\[
  q'=+\frac{qR}{d},\qquad
  \boldsymbol d'=\frac{R^2}{d^2}\boldsymbol d.
\]

图~\ref{ch20:fig:s34-3-7} 中真实电荷到球心距离为 $d<R$；像点距离为 $R^2/d>R$，因而不在要求电场的球腔内。

\par\noindent\begin{minipage}{\linewidth}
\cavityImageFigure
\diagramcaption{球腔内点电荷的像位于球面之外。像电荷不在所求球腔区域内。}\label{ch20:fig:s34-3-7}
\end{minipage}\par\medskip

\textbf{(1) 接地球壳的电场。} 球腔 $r<R$ 内的电势为
\[
  \boxed{\varphi(\boldsymbol r)=\frac1{4\pi\varepsilon_0}
  \left[-\frac q{|\boldsymbol r-\boldsymbol d|}
  +\frac{qR/d}{|\boldsymbol r-\boldsymbol d'|}\right]},
\]
电场为 $\boldsymbol E=-\boldsymbol\nabla\varphi$。导体材料内部 $\boldsymbol E=0$。球壳接地后，外表面
没有净电荷，外部边界电势和无穷远电势均为零，由唯一性定理
\[
  \boxed{\boldsymbol E_{\rm outside}=0}.
\]

以上像电荷写法取 $0<d<R$。若真实电荷恰在球心，取 $d\to0$ 的极限即可：
球腔内 $\varphi=-kq/r+kq/R$，$\boldsymbol E=-kq\hat{\boldsymbol r}/r^2$；
内表面 $\sigma_{\rm in}=q/(4\pi R^2)$，避免把发散的像电荷坐标当作真实奇点。

\textbf{(2) 内表面面电荷密度。} 令 $\theta$ 为球面点与 $+z$ 轴的夹角，该点到真实电荷的距离为
$s=(R^2+d^2-2Rd\cos\theta)^{1/2}$。由球腔侧法向场得到
\[
  \boxed{\sigma_{\rm in}(\theta)=
  \frac{q(R^2-d^2)}{4\pi R\bigl(R^2+d^2-2Rd\cos\theta\bigr)^{3/2}}}.
\]
真实电荷为负，所以内壁感应电荷处处为正；离电荷最近的 $\theta=0$ 处最大。积分可得内表面总感应电荷
$+q$。接地时外表面电荷为零，地面提供了这 $+q$ 电荷。

\textbf{(3) 球壳孤立且净电荷为零。} 导体材料内部仍须无场，因此内表面的电荷分布和上式完全不变，总量
仍为 $+q$。为使球壳总电荷为零，外表面必须带 $-q$。外表面为半径 $R_o$ 的球面，则这部分电荷均匀
分布，外部电场等效于球心的 $-q$：
\[
  \boldsymbol E_{\rm outside}=-\frac1{4\pi\varepsilon_0}
  \frac q{r^2}\hat{\boldsymbol r}.
\]
球腔内电场不变，只是电势整体增加常量 $-q/(4\pi\varepsilon_0R_o)$。

\textbf{(4) 屏蔽的方向性。} 外部电荷只改变导体外表面的电荷分布；空腔边界仍为等势面，且腔内无真实
电荷，唯一解是腔内无场。内部真实电荷却直接出现在空腔的泊松方程中，边界为等势并不能把它产生的
腔内电场消去，只能迫使内表面出现感应电荷。导体材料仍屏蔽了“穿过金属”的场；接地时外部无场，孤立
时外表面补偿电荷还会产生外场。
\end{xesolution}



\begin{origexamplebox}{【能量】}\label{ch20:ex:s34-3-6}
计算前文接地导体球外放点电荷的问题，以及例~\ref{ch20:ex:s34-3-5} 接地球腔内放点电荷的问题中，系统的有限静电势能。
\end{origexamplebox}\par\nopagebreak[4]

\begin{xesolution}
导体感应问题的有限相互作用能统一写成
\[
  U=\frac12q_{\rm real}\varphi_{\rm ind}(\boldsymbol r_{\rm real}).
\]
这里 $\varphi_{\rm ind}$ 是全部感应电荷在真实电荷处产生的电势。因感应电荷随真实电荷逐渐建立，不能把
像电荷当成另一个独立真实电荷；前面的 $1/2$ 必须保留。

\textbf{1. 球外点电荷。} $+Q$ 距接地球心 $L>R$，像电荷为 $-QR/L$，像位置为 $R^2/L$。因此
\[
  \boxed{U_{\rm out}=-\frac1{4\pi\varepsilon_0}
  \frac{Q^2R}{2(L^2-R^2)}}.
\]
把它对 $L$ 求负导数，恰好得到例题正文中真实电荷所受的吸力。

\textbf{2. 球腔内点电荷。} $-q$ 距球心 $d<R$，像电荷为 $+qR/d$，二者间距为
$(R^2-d^2)/d$，故
\[
  \boxed{U_{\rm in}=-\frac1{4\pi\varepsilon_0}
  \frac{q^2R}{2(R^2-d^2)}}.
\]
当电荷靠近球面时 $d\to R$，能量趋向负无穷，与局部近似为点电荷靠近接地无限平面的结果一致。

若例~\ref{ch20:ex:s34-3-5}中的球壳改为孤立且净电荷为零，外表面 $-q$ 的场还带来有限能量修正
$+q^2/(8\pi\varepsilon_0R_o)$，所以
\[
  U_{\rm neutral}=U_{\rm in}+\frac{q^2}{8\pi\varepsilon_0R_o}.
\]
这里比较的是扣除了点电荷自身无穷大自能后的相互作用能。
\end{xesolution}



\begin{origexamplebox}{【电像法】}\label{ch20:ex:s34-3-7}
两个互相垂直、半无限大的接地金属板围成空间的四分之一。以交线在截面内的交点为原点，两板分别为
$x=0$ 与 $y=0$，在第一象限 $(a,b)$ 处放一点电荷 $+q$。求第一象限中的电场。
\par\noindent\begin{minipage}{\linewidth}
\emCornerQuestion
\figcap{相互垂直的两接地导体板与真实点电荷；像电荷待构造。}\label{ch20:fig:em-source-thirty-fourth-section-7}
\end{minipage}\par\medskip
\end{origexamplebox}\par\nopagebreak[4]

\begin{xesolution}
先关于 $x=0$ 镜像，放置 $-q$ 于 $(-a,b)$，使 $x=0$ 面电势为零；再关于 $y=0$ 镜像，放置 $-q$ 于
$(a,-b)$。但这两个像电荷会破坏另一块板的零电势条件，因此还必须在 $(-a,-b)$ 放置 $+q$。四电荷关于
每块板都是等量异号镜像，故两块板上电势同时为零。

图~\ref{ch20:fig:s34-3-8} 的真实求解域是第一象限，另三个点都是像点；逐块检查镜像的相反符号即可验证两板零电势。

\par\noindent\begin{minipage}{\linewidth}
\cornerImagesFigure
\diagramcaption{两块正交接地板的三个像电荷。只有第一象限的 $+q$ 是真实电荷。}\label{ch20:fig:s34-3-8}
\end{minipage}\par\medskip

在第一象限取观察点 $\boldsymbol r=(x,y,z)$，令
\[
  \boldsymbol r_1=(a,b,0),\quad
  \boldsymbol r_2=(-a,b,0),\quad
  \boldsymbol r_3=(a,-b,0),\quad
  \boldsymbol r_4=(-a,-b,0).
\]
电势与电场分别为
\[
  \boxed{\varphi(\boldsymbol r)=\frac q{4\pi\varepsilon_0}
  \left(\frac1{|\boldsymbol r-\boldsymbol r_1|}
  -\frac1{|\boldsymbol r-\boldsymbol r_2|}
  -\frac1{|\boldsymbol r-\boldsymbol r_3|}
  +\frac1{|\boldsymbol r-\boldsymbol r_4|}\right)},
\]
\[
  \boxed{\boldsymbol E(\boldsymbol r)=\frac q{4\pi\varepsilon_0}
  \left(\frac{\boldsymbol r-\boldsymbol r_1}{|\boldsymbol r-\boldsymbol r_1|^3}
  -\frac{\boldsymbol r-\boldsymbol r_2}{|\boldsymbol r-\boldsymbol r_2|^3}
  -\frac{\boldsymbol r-\boldsymbol r_3}{|\boldsymbol r-\boldsymbol r_3|^3}
  +\frac{\boldsymbol r-\boldsymbol r_4}{|\boldsymbol r-\boldsymbol r_4|^3}\right)}.
\]
三个像电荷都位于第一象限之外，所以没有在真实区域内引入虚假奇点，构造合法。
\end{xesolution}



\begin{origexamplebox}{【半球和平面】}\label{ch20:ex:s34-3-8}
一个无限大接地平面上方凸起一个半径为 $R$ 的半球，球心在平面内。在球心正上方 $a>R$ 处放置一点
电荷 $+q$，求上方空间中的电场。若半球改为向下凹入平面的半球形凹槽，为什么同样的有限像电荷构造不能
直接使用？
\par\noindent\begin{minipage}{\linewidth}
\emHemisphereQuestion
\figcap{接地导体平面上的半球凸起及轴线上真实点电荷。}\label{ch20:fig:em-source-thirty-fourth-section-8}
\end{minipage}\par\medskip
\end{origexamplebox}\par\nopagebreak[4]

\begin{xesolution}
取平面为 $z=0$，真实电荷位于 $z=a$。要同时使平面和半径 $R$ 的球面电势为零，需要三枚像电荷：
\[
  -q\text{ 位于 }z=-a,\qquad
  -\frac{qR}{a}\text{ 位于 }z=\frac{R^2}{a},\qquad
  +\frac{qR}{a}\text{ 位于 }z=-\frac{R^2}{a}.
\]
前两枚 $\pm q$ 关于平面反号镜像，后两枚 $\mp qR/a$ 也关于平面反号镜像，所以平面电势为零；在球面上，
$+q$ 与上方球内像电荷相消，$-q$ 与下方球内像电荷相消，所以半球面电势也为零。

图~\ref{ch20:fig:s34-3-9} 中四个点都在对称轴上；关于平面成对反射使平面电势为零，关于球面的像变换使半球面电势为零。

\par\noindent\begin{minipage}{\linewidth}
\hemisphereImagesFigure
\diagramcaption{接地平面上凸半球的三像电荷构造。粗蓝线是实际导体边界。}\label{ch20:fig:s34-3-9}
\end{minipage}\par\medskip

对上半空间且半球外的观察点，电势为四个点电荷电势的代数和，电场为其负梯度：
\[
  \boxed{\varphi(\boldsymbol r)=\frac1{4\pi\varepsilon_0}
  \left[
  \frac q{|\boldsymbol r-a\hat{\boldsymbol z}|}
  -\frac q{|\boldsymbol r+a\hat{\boldsymbol z}|}
  -\frac{qR/a}{|\boldsymbol r-(R^2/a)\hat{\boldsymbol z}|}
  +\frac{qR/a}{|\boldsymbol r+(R^2/a)\hat{\boldsymbol z}|}
  \right]},
  \qquad \boldsymbol E=-\boldsymbol\nabla\varphi.
\]

若半球向下凹，球内上方点 $z=R^2/a$ 属于真实可达空间，而不是导体内部。把像电荷放在那里会在真实
区域制造一个并不存在的奇点，违反电像法的基本要求。此时边界虽仍由平面和半球面组成，但简单的有限
反射不再封闭，必须改用其他坐标方法、积分方程或数值方法求解。
\end{xesolution}



\begin{origexamplebox}{【能量相关】}\label{ch20:ex:s34-3-9}
一个无穷大的接地金属板，一侧距板 $a$ 处放置一点电荷 $+q$。
\begin{enumerate}[label=(\arabic*)]
  \item 系统电势能为多大？
  \item 缓慢将点电荷拉至距板 $2a$ 处，需要做功多少？
  \item 瞬间将点电荷移至距板 $2a$ 处，需要做功多少？
  \item 后两问为何不同？差值对应的能量去了哪里？
\end{enumerate}
\end{origexamplebox}\par\nopagebreak[4]

\begin{xesolution}
这里“瞬间移动”采用明确的冻结感应电荷近似：在搬运阶段保持原有面电荷分布，
只计算真实电荷在该固定静电场中的势能变化；不把这一步解释为真实电荷的无限速运动。
若实际问极快运动时外力的总功，还必须给出运动过程，并计入惯性和电磁辐射，
仅凭初末位置不能决定这些附加项。
接地平面的像电荷为位于板另一侧对称位置的 $-q$。真实电荷与像电荷距离为 $2a$。

\planeEnergyFigure

\textbf{(1) 初始能量。} 导体感应问题保留 $1/2$ 因子：
\[
  \boxed{U(a)=\frac12q\left(-\frac{kq}{2a}\right)
  =-\frac{kq^2}{4a}=-\frac{q^2}{16\pi\varepsilon_0a}}.
\]

\textbf{(2) 缓慢移动。} 感应电荷在每一时刻都跟随真实电荷调整，外力做功等于系统能量增量：
\[
  \boxed{W_{\rm slow}=U(2a)-U(a)
  =-\frac{kq^2}{8a}+\frac{kq^2}{4a}
  =\frac{kq^2}{8a}}.
\]

\textbf{(3) 瞬时移动。} 移动发生得足够快时，板上的初始感应电荷分布来不及改变。在这个极短阶段，它的
外部作用仍等效于固定在 $z=-a$ 的像电荷 $-q$。真实电荷从 $z=a$ 跳到 $z=2a$，它与固定感应场的距离
从 $2a$ 变为 $3a$，故外力做功为
\[
  \boxed{W_{\rm fast}=q\left(-\frac{kq}{3a}+\frac{kq}{2a}\right)
  =\frac{kq^2}{6a}}.
\]

\textbf{(4) 能量去向。} 瞬移完成后，旧感应电荷分布不是新位置对应的静电平衡分布。表面电荷重新流动，
最终达到 $U(2a)$。瞬时过程比准静态过程多输入
\[
  \boxed{W_{\rm fast}-W_{\rm slow}=\frac{kq^2}{24a}}.
\]
这部分能量在感应电荷的暂态运动中转化为焦耳热，或以电磁辐射形式向外传出；两者的份额取决于暂态动力学，不能由静电模型确定；它没有保留在最终静电场中。
若理想导体完全无耗散，电荷会持续振荡，不能立即到达新的静电平衡，这正与前文等离子体振荡的讨论呼应。
\end{xesolution}



\begin{origexamplebox}{【匀强电场中的导体球】}\label{ch20:ex:s34-3-10}
把半径为 $R$ 的中性导体球放置在原本沿 $+z$ 方向的匀强电场
$\boldsymbol E_0=E_0\hat{\boldsymbol z}$ 中，求空间中的电势、电场及球面感应电荷分布。
\end{origexamplebox}\par\nopagebreak[4]

\begin{xesolution}
无导体时可取电势 $\varphi_0=-E_0r\cos\theta$。导体球外需要补上一个位于球心、沿 $+z$ 方向的偶极子
势 $A\cos\theta/r^2$。令球面 $r=R$ 为等势面：
\[
  -E_0R\cos\theta+\frac{A\cos\theta}{R^2}=\text{常量}.
\]
球不带净电荷，所以没有 $1/r$ 单极项；取球内常电势为零，得到 $A=E_0R^3$。于是
\[
  \boxed{\varphi(r,\theta)=
  \begin{cases}
    -E_0\left(r-\dfrac{R^3}{r^2}\right)\cos\theta,&r\ge R,\\[6pt]
    0,&r<R.
  \end{cases}}
\]
球外电场的球坐标分量为
\[
  \boxed{E_r=E_0\left(1+2\frac{R^3}{r^3}\right)\cos\theta,
  \qquad
  E_\theta=-E_0\left(1-\frac{R^3}{r^3}\right)\sin\theta},
\]
球内 $\boldsymbol E=0$。在 $r=R$，$E_\theta=0$，所以场线严格垂直球面；法向场为
$E_r(R)=3E_0\cos\theta$，故
\[
  \boxed{\sigma(\theta)=3\varepsilon_0E_0\cos\theta}.
\]
北半球 $\cos\theta>0$ 带正感应电荷，南半球带负感应电荷，总感应电荷为零。远处附加项等同于偶极矩
\[
  \boxed{\boldsymbol p=4\pi\varepsilon_0R^3\boldsymbol E_0}.
\]

\par\noindent\begin{minipage}{\linewidth}
\uniformFieldSphereFigure
\diagramcaption{匀强电场中的中性导体球。实线为电场线，虚线为等势线；球内无场，球面切向场为零。}
\end{minipage}\par\medskip

取图中外场箭头方向为 $+z$ 轴，极角 $\theta$ 也从这一方向量起。场线在
球面两极处最密，恰对应 $|\sigma|$ 最大；赤道处球面场为零；离开球面后，赤道外侧的场逐渐恢复，场线绕过导体而不进入其内部。
\end{xesolution}
}

{
\graphicspath{{/Users/dongxingxi/Documents/物理竞赛一轮资料/sample/xes-original/thirty-fifth-section/}{/Users/dongxingxi/Documents/物理竞赛一轮资料/sample/xes-original/thirty-fifth-section/assets/}}
\section{电容的定义与导体组}

% Labels are deliberately transparent: no white boxes may cover geometry.
\tikzset{
  xeslabel/.style={font=\small,text=Ink,fill=none,inner sep=.55pt},
  xessmall/.style={font=\footnotesize,text=TextGray,fill=none,inner sep=.4pt},
  xestiny/.style={font=\scriptsize,text=TextGray,fill=none,inner sep=.3pt},
  xesline/.style={draw=Ink,line width=.85pt},
  xesblue/.style={draw=DeepBlue,line width=1pt},
  xesarrow/.style={-{Latex[length=2.1mm,width=1.3mm]},draw=DeepBlue,line width=.9pt},
  fieldline/.style={draw=DeepBlue,line width=.82pt,postaction={decorate},
    decoration={markings,mark=at position .57 with {\arrow{Latex[length=1.65mm,width=1.05mm]}}}},
  circuit/.style={draw=Ink,line width=.9pt,line cap=round,line join=round}
}


\newcommand{\capacitorFieldFigure}{%
\begin{center}
\includegraphics[width=.76\linewidth]{/Users/dongxingxi/Documents/物理竞赛一轮资料/sample/xes-original/restructured-complete/assets/revision-ch19-ch20/capacitor-pair.pdf}
\end{center}}


\newcommand{\parallelPlateFigure}{%
\begin{center}
\begin{tikzpicture}[x=1cm,y=1cm]
  % geometry panel
  \begin{scope}[shift={(-4.45,0)}]
    \fill[Gold!18] (-1.5,-1.25)--(-.72,-.85)--(-.72,1.25)--(-1.5,.85)--cycle;
    \draw[xesblue] (-1.5,-1.25)--(-.72,-.85)--(-.72,1.25)--(-1.5,.85)--cycle;
    \fill[DeepBlue!10] (.4,-1.25)--(1.18,-.85)--(1.18,1.25)--(.4,.85)--cycle;
    \draw[xesblue] (.4,-1.25)--(1.18,-.85)--(1.18,1.25)--(.4,.85)--cycle;
    \node[xeslabel,anchor=east] at (-1.58,.65) {$+Q$};
    \node[xeslabel,anchor=west] at (1.26,.65) {$-Q$};
    \draw[<->,Gold!80!black,line width=.65pt] (-.68,-1.55)--(.37,-1.55);
    \node[xestiny,text=Gold!70!black] at (-.15,-1.8) {$d$};
    \node[xessmall] at (0,1.65) {面积 $A$，板间距 $d$};
  \end{scope}
  \node[inner sep=0pt] at (3,0) {\includegraphics[width=6.7cm]{/Users/dongxingxi/Documents/物理竞赛一轮资料/sample/xes-original/restructured-complete/assets/revision-ch19-ch20/plate-edges.pdf}};

\end{tikzpicture}
\end{center}}

\newcommand{\coaxialCapacitorFigure}{%
\begin{center}
\begin{tikzpicture}[x=1cm,y=1cm]
  \begin{scope}[shift={(-3.0,0)}]
    \fill[DeepBlue!9] (-1.45,-1.55) rectangle (1.45,1.25);
    \draw[xesblue] (-1.45,-1.55)--(-1.45,1.25) (1.45,-1.55)--(1.45,1.25);
    \draw[xesblue] (0,1.25) ellipse (1.45 and .42);
    \draw[xesblue] (-1.45,-1.55) arc[start angle=180,end angle=360,x radius=1.45,y radius=.42];
    \fill[Gold!28] (-.62,-1.55) rectangle (.62,1.25);
    \draw[Gold!80!black,line width=.9pt] (-.62,-1.55)--(-.62,1.25) (.62,-1.55)--(.62,1.25);
    \draw[Gold!80!black,line width=.9pt] (0,1.25) ellipse (.62 and .2);
    \node[xessmall] at (0,-2.1) {同轴圆柱，长度 $h$};
    \draw[<->,Gold!75!black,line width=.6pt] (1.8,-1.55)--(1.8,1.25);
    \node[xestiny,text=Gold!70!black,anchor=west] at (1.95,-.15) {$h$};
  \end{scope}
  \begin{scope}[shift={(3.2,0)}]
    \fill[DeepBlue!8] (0,0) circle (1.65);
    
    \fill[Gold!26] (0,0) circle (.56);
    \draw[xesblue] (0,0) circle (1.65);
    \draw[Gold!80!black,line width=.9pt] (0,0) circle (.56);
    \foreach \a in {15,45,...,345}{\draw[fieldline] (\a:.62)--(\a:1.55);}
    \draw[Ink,line width=.55pt] (0,0)--(-.56,0);\node[xestiny,above] at (-.28,.05) {$a$};
    \draw[Ink,line width=.55pt] (0,0)--(1.65,0);\node[xestiny,below] at (1.12,-.04) {$b$};
    \node[xessmall] at (0,-2.1) {横截面：$a<r<b$};
  \end{scope}
\end{tikzpicture}
\end{center}}

\newcommand{\capacitorParallelFigure}{%
\begin{center}
\begin{tikzpicture}[x=1cm,y=1cm]
  \draw[circuit] (-3,-1.7)--(-3,1.7)--(-1.0,1.7);
  \draw[circuit] (1.0,1.7)--(3,1.7)--(3,-1.7);
  \draw[circuit] (-3,0)--(-1.0,0) (1.0,0)--(3,0);
  \draw[circuit] (-3,-1.7)--(-.45,-1.7) (.45,-1.7)--(3,-1.7);
  \draw[circuit,line width=1.4pt] (-1.0,1.42)--(-1.0,1.98) (1.0,1.42)--(1.0,1.98);
  \draw[circuit,line width=1.4pt] (-1.0,-.28)--(-1.0,.28) (1.0,-.28)--(1.0,.28);
  \draw[circuit,line width=1.4pt] (-.45,-2.08)--(-.45,-1.32) (.45,-1.94)--(.45,-1.46);
  \node[xeslabel,above] at (0,1.98) {$C_1$};
  \node[xeslabel,above] at (0,.28) {$C_2$};
  \node[xeslabel,below] at (0,-2.08) {$U$};
  \node[xessmall] at (0,1.15) {$Q_1=C_1U$};
  \node[xessmall] at (0,-.58) {$Q_2=C_2U$};
  \node[xessmall,align=left,anchor=west] at (4.0,.35) {两端电压相同\\总电荷相加};
\end{tikzpicture}
\end{center}}

\newcommand{\capacitorSeriesFigure}{%
\begin{center}
\begin{tikzpicture}[x=1cm,y=1cm]
  \draw[circuit] (-3,-1.45)--(-3,1.45)--(-1.55,1.45);
  \draw[circuit] (-.75,1.45)--(.75,1.45);
  \draw[circuit] (1.55,1.45)--(3,1.45)--(3,-1.45);
  \draw[circuit] (-3,-1.45)--(-.45,-1.45) (.45,-1.45)--(3,-1.45);
  \draw[circuit,line width=1.4pt] (-1.55,1.17)--(-1.55,1.73) (-.75,1.17)--(-.75,1.73);
  \draw[circuit,line width=1.4pt] (.75,1.17)--(.75,1.73) (1.55,1.17)--(1.55,1.73);
  \draw[circuit,line width=1.4pt] (-.45,-1.83)--(-.45,-1.07) (.45,-1.69)--(.45,-1.21);
  \node[xeslabel,above] at (-1.15,1.76) {$C_1$};
  \node[xeslabel,above] at (1.15,1.76) {$C_2$};
  \node[xeslabel,below] at (0,-1.87) {$U$};
  \node[xessmall,align=center] at (0,.55) {中间导体净电荷为零\\两电容电荷量均为 $Q$};
  \node[xessmall] at (0,-.45) {$U=Q/C_1+Q/C_2$};
\end{tikzpicture}
\end{center}}

\newcommand{\greenSphereFigure}{%
\begin{center}
\begin{tikzpicture}[x=1cm,y=1cm]
  \fill[Gold!16] (-1.3,0) circle (1.25);\draw[xesblue] (-1.3,0) circle (1.25);
  \fill[Ink] (3.35,0) circle (2.2pt);\node[xeslabel,above=4pt] at (3.35,0) {$P$};
  \fill[Ink] (-1.3,0) circle (1.7pt);
  \draw[xesline] (-1.3,0)--(3.35,0);
  \node[xeslabel,above] at (1.0,0) {$r$};
  \node[xessmall,align=left,anchor=west] at (4.25,.35) {状态 A：球带单位电荷\\状态 B：$P$ 点放电荷 $q$};
\end{tikzpicture}
\end{center}}

\newcommand{\greenTwoConductorFigure}{%
\begin{center}
\begin{tikzpicture}[x=1cm,y=1cm]
  \fill[DeepBlue!10] (-2.6,0) ellipse (1.3 and .95);\draw[xesblue] (-2.6,0) ellipse (1.3 and .95);
  \fill[Gold!16] (2.4,.1) ellipse (1.05 and 1.18);\draw[xesblue] (2.4,.1) ellipse (1.05 and 1.18);
  \node[xeslabel,align=center] at (-2.6,0) {$Q_1$\\$\varphi_1$};
  \node[xeslabel,align=center] at (2.4,.1) {$Q_2$\\$\varphi_2$};
  \draw[<->,Gold!75!black,line width=.65pt] (-1.15,0)--(1.2,0);
  \node[xessmall,text=Gold!70!black,above] at (0,.05) {静电感应耦合};
  \node[xessmall] at (0,-1.55) {另一状态用 $Q'_i,\varphi'_i$ 表示，几何形状保持不变};
\end{tikzpicture}
\end{center}}

\newcommand{\dielectricMeterFigure}{%
\begin{center}
\includegraphics[width=.93\linewidth]{/Users/dongxingxi/Documents/物理竞赛一轮资料/sample/xes-original/restructured-complete/assets/revision-ch19-ch20/dielectric-meter.png}
\end{center}}


\newcommand{\polarizationFigure}{%
\begin{center}
\begin{tikzpicture}[x=1cm,y=1cm]
  \begin{scope}[shift={(-3.6,0)}]
    \draw[xesblue] (-2.1,-1.45) rectangle (2.1,1.45);
    \foreach \x in {-1.45,0,1.45}{\foreach \y in {-0.85,0,.85}{
      \fill[DeepBlue!20] (\x-.24,\y) ellipse (.38 and .2);
      \fill[Gold!35] (\x+.24,\y) ellipse (.38 and .2);
      \node[xestiny,text=Ink] at (\x-.28,\y) {$-$};
      \node[xestiny,text=Ink] at (\x+.28,\y) {$+$};}}
    \draw[xesarrow] (-1.8,1.9)--(1.8,1.9);
    \node[xessmall,above] at (0,1.93) {$\boldsymbol E_0$};
    \node[xessmall] at (0,-1.8) {微观：偶极矩沿外场排列};
  \end{scope}
  \begin{scope}[shift={(3.6,0)}]
    \draw[xesblue] (-2.1,-1.45) rectangle (2.1,1.45);
    \foreach \y in {-1.1,-.55,0,.55,1.1}{
      \node[xestiny,text=Ink] at (-1.82,\y) {$-$};
      \node[xestiny,text=Ink] at (1.82,\y) {$+$};
      \draw[-{Latex[length=1.65mm]},Gold!80!black,line width=.72pt] (1.35,\y)--(-1.35,\y);}
    \node[xessmall,text=Gold!70!black] at (0,.28) {$\boldsymbol E_p$};
    \node[xessmall] at (0,-1.8) {宏观：束缚面电荷产生反向极化场};
    \node[xestiny,anchor=east] at (-2.25,0) {$-\sigma_p$};
    \node[xestiny,anchor=west] at (2.25,0) {$+\sigma_p$};
  \end{scope}
\end{tikzpicture}
\end{center}}

\newcommand{\exampleParallelFigure}{%
\begin{center}
\begin{tikzpicture}[x=1cm,y=1cm]
  \draw[xesblue,line width=1.25pt] (-3,.75)--(3,.75);
  \draw[xesblue,line width=1.25pt] (-3,-.75)--(3,-.75);
  \draw[<->,Gold!75!black,line width=.65pt] (3.35,-.75)--(3.35,.75);
  \node[xestiny,text=Gold!70!black,anchor=west] at (3.52,0) {$d$};
  \draw[<->,Gold!75!black,line width=.65pt] (-3,-1.18)--(3,-1.18);
  \node[xestiny,text=Gold!70!black] at (0,-1.42) {$l$};
  \node[xessmall] at (0,1.15) {$S=l^2$};
  \foreach \x in {-2.5,-1.5,...,2.5}{\draw[fieldline] (\x,.62)--(\x,-.62);}
\end{tikzpicture}
\end{center}}

\newcommand{\exampleSphericalFigure}{%
\begin{center}
\begin{tikzpicture}[x=1cm,y=1cm]
  
  \fill[Gold!16] (0,0) circle (.9);
  \draw[xesblue] (0,0) circle (1.7);
  \draw[Gold!80!black,line width=.9pt] (0,0) circle (.9);
  \draw[xesline] (0,0)--(.9,0) node[xestiny,pos=.62,above] {$R_1$};
  \draw[xesline] (0,0)--(1.7,0) node[xestiny,pos=.82,below] {$R_2$};
  \foreach \a in {15,45,...,345}{\draw[fieldline] (\a:.97)--(\a:1.62);}
\end{tikzpicture}
\end{center}}

\newcommand{\touchingSpheresFigure}{%
\begin{center}
\begin{tikzpicture}[x=1cm,y=1cm]
  \fill[DeepBlue!7] (-1.25,0) circle (1.25);\draw[xesblue] (-1.25,0) circle (1.25);
  \fill[DeepBlue!7] (1.25,0) circle (1.25);\draw[xesblue] (1.25,0) circle (1.25);
  \draw[xesline] (1.25,0)--(2.13,.88);
  \node[xeslabel,above] at (1.72,.48) {$R$};
  \node[xessmall] at (0,-1.65) {两球接触并处于同一电势 $V$};
\end{tikzpicture}
\end{center}}

\newcommand{\touchingSphereImagesFigure}{%
\begin{center}
\begin{tikzpicture}[x=1cm,y=1cm]
  \draw[xesblue] (-1.55,0) circle (1.55);\draw[xesblue] (1.55,0) circle (1.55);
  \draw[RuleGray,densely dashed,line width=.55pt] (-3.2,0)--(3.2,0);
  \foreach \x/\sgn in {-1.55/DeepBlue,-.78/Ink,-.52/DeepBlue,-.39/Ink,
                         1.55/DeepBlue,.78/Ink,.52/DeepBlue,.39/Ink}{
    \fill[\sgn] (\x,0) circle (3.1pt);}
  % 接触点附近的像电荷十分密集；用错层引线避免 q_2,q_3,q_4 互相覆盖。
  \node[xestiny,above=3pt] at (-1.55,0) {$q_1$};
  \draw[RuleGray,line width=.45pt] (-.78,-.05)--(-.90,-.34)
    node[xestiny,below] {$q_2$};
  \draw[RuleGray,line width=.45pt] (-.52,.05)--(-.70,.38)
    node[xestiny,above] {$q_3$};
  \draw[RuleGray,line width=.45pt] (-.39,-.05)--(-.48,-.64)
    node[xestiny,below] {$q_4$};
  \node[xestiny,above=3pt] at (1.55,0) {$q_1$};
  \draw[RuleGray,line width=.45pt] (.78,-.05)--(.90,-.34)
    node[xestiny,below] {$q_2$};
  \draw[RuleGray,line width=.45pt] (.52,.05)--(.70,.38)
    node[xestiny,above] {$q_3$};
  \draw[RuleGray,line width=.45pt] (.39,-.05)--(.48,-.64)
    node[xestiny,below] {$q_4$};
  \node[xessmall] at (0,-1.95) {$x_n=R/n,\quad q_n=4\pi\varepsilon_0RV\,(-1)^{n-1}/n$};
  \node[xessmall] at (0,-2.43) {两侧电荷序列都向接触点聚集};
\end{tikzpicture}
\end{center}}

\newcommand{\twoCylindersFigure}{%
\begin{center}
\includegraphics[width=.76\linewidth]{/Users/dongxingxi/Documents/物理竞赛一轮资料/sample/xes-original/restructured-complete/assets/revision-ch19-ch20/two-cylinders.pdf}
\end{center}}


\newcommand{\modifiedCapacitorFigure}{%
\begin{center}
\begin{tikzpicture}[x=1cm,y=1cm]
  \begin{scope}[shift={(-3.5,0)}]
    \draw[xesblue,line width=1.1pt] (-2,1)--(2,1) (-2,-1)--(2,-1);
    \fill[Gold!22,pattern=north east lines,pattern color=Gold!70!black] (-2,-.9) rectangle (0,.9);
    \draw[RuleGray,line width=.55pt] (0,-.9)--(0,.9);
    \node[xessmall] at (-1,0) {$\varepsilon_r$};\node[xessmall] at (1,0) {真空};
    \node[xessmall] at (0,-1.42) {(a) 占一半面积：两个电容并联};
  \end{scope}
  \begin{scope}[shift={(3.5,0)}]
    \draw[xesblue,line width=1.1pt] (-2,1)--(2,1) (-2,-1)--(2,-1);
    \fill[Gold!22,pattern=north east lines,pattern color=Gold!70!black] (-2,0) rectangle (2,.9);
    \draw[RuleGray,line width=.55pt] (-2,0)--(2,0);
    \node[xessmall] at (0,.45) {$\varepsilon_r$};\node[xessmall] at (0,-.48) {真空};
    \node[xessmall] at (0,-1.42) {(b) 占一半厚度：两个电容串联};
  \end{scope}
\end{tikzpicture}
\end{center}}

\newcommand{\dielectricInsertionFigure}{%
\begin{center}
\begin{tikzpicture}[x=1cm,y=1cm]
  \draw[xesblue,line width=1.2pt] (-4,1)--(4,1) (-4,-1)--(4,-1);
  \fill[Gold!22,pattern=north east lines,pattern color=Gold!70!black] (-4,-.9) rectangle (1.0,.9);
  \draw[Gold!80!black,line width=.75pt] (1.0,-.9)--(1.0,.9);
  \draw[<->,Gold!75!black,line width=.65pt] (-4,1.38)--(1.0,1.38);
  \node[xestiny,text=Gold!70!black] at (-1.5,1.62) {$x$};
  \draw[<->,Gold!75!black,line width=.65pt] (1.0,1.38)--(4,1.38);
  \node[xestiny,text=Gold!70!black] at (2.5,1.62) {$a-x$};
  \draw[<->,Gold!75!black,line width=.65pt] (4.35,-1)--(4.35,1);
  \node[xestiny,text=Gold!70!black,anchor=west] at (4.5,0) {$d$};
  \draw[xesarrow] (.6,0)--(1.75,0);\node[xessmall,above] at (1.18,.05) {$F$};
  \node[xessmall] at (-1.5,0) {$\varepsilon_r$};
  \node[xessmall] at (0,-1.48) {极板宽度为 $b$，总面积 $S=ab$};
\end{tikzpicture}
\end{center}}

\newcommand{\yCapacitorFigure}{%
\begin{center}
\begin{tikzpicture}[x=1cm,y=1cm]
  \coordinate (O) at (0,0);
  \draw[circuit] (O)--(0,.62);
  \draw[circuit,line width=1.35pt] (-.34,.62)--(.34,.62) (-.34,1.02)--(.34,1.02);
  \draw[circuit] (0,1.02)--(0,1.85);
  \node[xeslabel,anchor=west] at (.45,.82) {$C_1$};\node[xeslabel,above] at (0,1.85) {$\varphi_1$};
  \begin{scope}[rotate=120]
    \draw[circuit] (0,0)--(0,.72);
    \draw[circuit,line width=1.35pt] (-.34,.72)--(.34,.72) (-.34,1.12)--(.34,1.12);
    \draw[circuit] (0,1.12)--(0,1.95);
  \end{scope}
  \node[xeslabel] at (-1.85,-1.18) {$\varphi_2$};\node[xeslabel] at (-1.00,-1.02) {$C_2$};
  \begin{scope}[rotate=-120]
    \draw[circuit] (0,0)--(0,.72);
    \draw[circuit,line width=1.35pt] (-.34,.72)--(.34,.72) (-.34,1.12)--(.34,1.12);
    \draw[circuit] (0,1.12)--(0,1.95);
  \end{scope}
  \node[xeslabel] at (1.85,-1.18) {$\varphi_3$};\node[xeslabel] at (1.00,-1.02) {$C_3$};
  \fill[DeepBlue] (O) circle (2.3pt);\node[xessmall,anchor=west] at (.18,.12) {浮置中性节点，电势 $V$};
\end{tikzpicture}
\end{center}}

\newcommand{\dielectricSphereFieldFigure}{%
\begin{center}
\includegraphics[width=.76\linewidth]{/Users/dongxingxi/Documents/物理竞赛一轮资料/sample/xes-original/restructured-complete/assets/revision-ch19-ch20/dielectric-sphere.pdf}
\end{center}}


\newcommand{\polarizationEnergyFigure}{%
\begin{center}
\begin{tikzpicture}[x=1cm,y=1cm]
  \begin{scope}[shift={(-3.3,0)}]
    \draw[xesblue] (-1.7,-1.0) rectangle (1.7,1.0);
    \foreach \x in {-1.15,0,1.15}{\draw[decorate,decoration={zigzag,segment length=3pt,amplitude=2pt},Gold!75!black,line width=.7pt] (\x-.35,0)--(\x+.35,0);}
    \foreach \x in {-1.15,0,1.15}{\node[xestiny] at (\x-.48,0) {$-$};\node[xestiny] at (\x+.48,0) {$+$};}
    \node[xessmall] at (0,-1.38) {无外场：微观“弹簧”未拉伸};
  \end{scope}
  \begin{scope}[shift={(3.3,0)}]
    \draw[xesblue] (-1.7,-1.0) rectangle (1.7,1.0);
    \foreach \x in {-1.15,0,1.15}{\draw[decorate,decoration={zigzag,segment length=3pt,amplitude=2pt},Gold!75!black,line width=.7pt] (\x-.55,0)--(\x+.55,0);}
    \foreach \x in {-1.15,0,1.15}{\node[xestiny] at (\x-.7,0) {$-$};\node[xestiny] at (\x+.7,0) {$+$};}
    \draw[xesarrow] (-1.45,.68)--(1.45,.68);\node[xessmall,above] at (0,.71) {$\boldsymbol E$};
    \node[xessmall] at (0,-1.38) {有外场：极化自由度也储存能量};
  \end{scope}
\end{tikzpicture}
\end{center}}

\newcommand{\chargingEnergyFigure}{%
\begin{center}
\begin{tikzpicture}[x=1cm,y=1cm]
  \node[draw=DeepBlue,line width=.9pt,rounded corners=2pt,minimum width=2.3cm,minimum height=1.0cm,text=Ink,fill=DeepBlue!6] (S) at (-3.6,0) {电源 $C\mathcal E^2$};
  \node[draw=DeepBlue,line width=.9pt,rounded corners=2pt,minimum width=2.4cm,minimum height=1.0cm,text=Ink,fill=DeepBlue!6] (C) at (0,1.15) {电容场能 $\tfrac12C\mathcal E^2$};
  \node[draw=Gold!80!black,line width=.9pt,rounded corners=2pt,minimum width=2.4cm,minimum height=1.0cm,text=Ink,fill=Gold!12] (H) at (0,-1.15) {焦耳热 $\tfrac12C\mathcal E^2$};
  \draw[xesarrow] (S.east)--(C.west);\draw[-{Latex[length=2.1mm]},Gold!80!black,line width=.9pt] (S.east)--(H.west);
  \node[xessmall,align=left,anchor=west] at (2.1,.25) {$R>0$ 时分配比例与 $R$ 无关；\\$R\to0^+$ 只使过程更快。};
\end{tikzpicture}
\end{center}}

\newcommand{\earthCapacitorFigure}{%
\begin{center}
\begin{tikzpicture}[x=1cm,y=1cm]
  \fill[DeepBlue!16] (0,0) circle (2.05);\draw[xesblue] (0,0) circle (2.05);
  \draw[Gold!80!black,line width=3.3pt,opacity=.7] (0,0) circle (2.34);
  \node[xessmall] at (0,0) {地球};
  \node[xestiny] at (0,-.45) {半径 $R_\oplus$，电荷 $-Q$};
  \draw[<->,Gold!75!black,line width=.65pt] (2.05,0)--(2.34,0);
  \node[xestiny,text=Gold!70!black,anchor=west] at (2.48,0) {$d=5\,\mathrm{km}$};
  \node[xessmall,anchor=west] at (3.35,.62) {等效大气上边界：$+Q$};
  \draw[RuleGray,-{Latex[length=1.5mm]},line width=.55pt] (3.2,.55)--(2.25,.35);
  \node[xessmall,anchor=west] at (3.35,-.35) {$C=4\pi\varepsilon_0\dfrac{R_\oplus(R_\oplus+d)}{d}$};
\end{tikzpicture}
\end{center}}

% Pure question geometry; analysis figure retained separately.
\newcommand{\emParallelQuestion}{%
\begin{center}
\begin{tikzpicture}[x=1cm,y=1cm]
  \draw[xesblue,line width=1.25pt] (-3,.75)--(3,.75);
  \draw[xesblue,line width=1.25pt] (-3,-.75)--(3,-.75);
  \draw[<->,Gold!75!black,line width=.65pt] (3.35,-.75)--(3.35,.75);
  \node[xestiny,text=Gold!70!black,anchor=west] at (3.52,0) {$d$};
  \draw[<->,Gold!75!black,line width=.65pt] (-3,-1.18)--(3,-1.18);
  \node[xestiny,text=Gold!70!black] at (0,-1.42) {$l$};
  \node[xessmall] at (0,1.15) {$S=l^2$};
\end{tikzpicture}
\end{center}}

% Pure question geometry; analysis figure retained separately.
\newcommand{\emSphericalQuestion}{%
\begin{center}
\begin{tikzpicture}[x=1cm,y=1cm]
  
  \fill[Gold!16] (0,0) circle (.9);
  \draw[xesblue] (0,0) circle (1.7);
  \draw[Gold!80!black,line width=.9pt] (0,0) circle (.9);
  \draw[xesline] (0,0)--(.9,0) node[xestiny,pos=.62,above] {$R_1$};
  \draw[xesline] (0,0)--(1.7,0) node[xestiny,pos=.82,below] {$R_2$};
\end{tikzpicture}
\end{center}}

% Pure question geometry; analysis figure retained separately.
\newcommand{\emModifiedQuestion}{%
\begin{center}
\begin{tikzpicture}[x=1cm,y=1cm]
  \path[use as bounding box] (-5.55,-1.7) rectangle (5.55,1.15);
  \begin{scope}[shift={(-3.5,0)}]
    \fill[Gold!22,pattern=north east lines,pattern color=Gold!70!black] (-2,-.9) rectangle (0,.9);
    \draw[RuleGray,line width=.55pt] (0,-.9)--(0,.9);
    \node[xessmall] at (-1,0) {$\varepsilon_r$};\node[xessmall] at (1,0) {真空};
    \draw[xesblue,line width=1.1pt] (-2,1)--(2,1) (-2,-1)--(2,-1);
    \node[xessmall] at (0,-1.42) {(a) 占一半面积};
  \end{scope}
  \begin{scope}[shift={(3.5,0)}]
    \fill[Gold!22,pattern=north east lines,pattern color=Gold!70!black] (-2,0) rectangle (2,.9);
    \draw[RuleGray,line width=.55pt] (-2,0)--(2,0);
    \node[xessmall] at (0,.45) {$\varepsilon_r$};\node[xessmall] at (0,-.48) {真空};
    \draw[xesblue,line width=1.1pt] (-2,1)--(2,1) (-2,-1)--(2,-1);
    \node[xessmall] at (0,-1.42) {(b) 占一半厚度};
  \end{scope}
\end{tikzpicture}
\end{center}}

% Pure question geometry; analysis figure retained separately.
\newcommand{\emInsertionQuestion}{%
\begin{center}
\begin{tikzpicture}[x=1cm,y=1cm]
  \draw[xesblue,line width=1.2pt] (-4,1)--(4,1) (-4,-1)--(4,-1);
  \fill[Gold!22,pattern=north east lines,pattern color=Gold!70!black] (-4,-.9) rectangle (1.0,.9);
  \draw[Gold!80!black,line width=.75pt] (1.0,-.9)--(1.0,.9);
  \draw[<->,Gold!75!black,line width=.65pt] (-4,1.38)--(1.0,1.38);
  \node[xestiny,text=Gold!70!black] at (-1.5,1.62) {$x$};
  \draw[<->,Gold!75!black,line width=.65pt] (1.0,1.38)--(4,1.38);
  \node[xestiny,text=Gold!70!black] at (2.5,1.62) {$a-x$};
  \draw[<->,Gold!75!black,line width=.65pt] (4.35,-1)--(4.35,1);
  \node[xestiny,text=Gold!70!black,anchor=west] at (4.5,0) {$d$};
  \node[xessmall] at (-1.5,0) {$\varepsilon_r$};
  \node[xessmall] at (0,-1.48) {极板宽度为 $b$，总面积 $S=ab$};
\end{tikzpicture}
\end{center}}


本章前文讨论了导体在静电平衡状态下的性质：导体内部电场为零，整个导体是等势体，电荷分布在表面。
当几个导体共同存在时，每个导体都可以用其电荷和电势描述；改变一个导体的电荷，还会通过静电感应
改变其他导体的表面电荷分布与电势。本章把这种相互作用定量化。最简单的量是孤立导体的电容以及由
两个导体构成的电容器的电容；在此基础上，还要讨论一般导体组、介质极化和电位移矢量。


\subsection{孤立导体的电容}


先设想空间中有许多导体，其中只有一个导体不接地，其余导体均与大地相连，电势固定为零。给这个
孤立导体电荷 $Q$ 后，它达到静电平衡并具有电势 $U$。由静电场的唯一性和叠加原理可知：导体的几何
形状、位置及周围接地边界不变时，把所有电荷同时放大 $k$ 倍，电势也放大 $k$ 倍。因此 $Q$ 与 $U$
成正比，写成
\begin{equation}
  \boxed{Q=CU}.
\end{equation}
比例系数 $C$ 称为该孤立导体的电容。它只由导体的形状、大小、位置以及周围边界决定，与本次给它多少
电荷无关。在国际单位制中电容的单位是法拉，$1\,\mathrm F=1\,\mathrm C/\mathrm V$；实际孤立导体的
电容通常很小，常用 $\mu\mathrm F$、$\mathrm{nF}$ 和 $\mathrm{pF}$。

半径为 $R$ 的孤立导体球外没有其他导体。球带电 $Q$ 时，球面外电场与位于球心的点电荷相同，取无穷远
电势为零，球面电势为
\[
  U=\int_R^{\infty}\frac{Q}{4\pi\varepsilon_0r^2}\,\mathrm dr
   =\frac{Q}{4\pi\varepsilon_0R}.
\]
代入定义得到
\begin{equation}
  \boxed{C=\frac QU=4\pi\varepsilon_0R}.
\end{equation}
可见球半径越大，达到相同电势所需的电荷越多，电容也越大。例如 $R=1\,\mathrm m$ 时，电容仍只有
约 $111\,\mathrm{pF}$。

若孤立导体的电容为 $C$，把它从不带电缓慢充到电荷 $Q$。中间状态带电 $q$ 时电势为 $q/C$，再搬入
$\mathrm dq$ 所做的功为 $(q/C)\mathrm dq$。总静电能为
\begin{equation}
  \boxed{W=\int_0^Q\frac qC\,\mathrm dq
  =\frac{Q^2}{2C}=\frac12QU=\frac12CU^2}.
\end{equation}
这些能量保存在导体带电所建立的电场中。由于静电场能量密度非负，电容也必为正数。


\subsection{电容器的电容}


取两个彼此绝缘的导体，分别带 $+Q$ 和 $-Q$。若它们的放置使电场主要限制在两者之间，这一导体组
称为电容器，两个导体称为两极。带正电的一极为正极，另一极为负极。

图~\ref{ch20:fig:s35-4-1} 以两球为例，正球上的场线沿外法线发出，负球上的场线沿内法线进入；两球各自等势，但电势一般不同。

\par\noindent\begin{minipage}{\linewidth}
\capacitorFieldFigure
\diagramcaption{双极电容器的两球示例。两极带等量异号电荷，电场线由正极指向负极。}\label{ch20:fig:s35-4-1}
\end{minipage}\par\medskip

由唯一性定理，两极电荷量 $Q$ 与电势差 $U=\varphi_+-\varphi_-$ 成正比：
\begin{equation}
  \boxed{Q=CU}.
\end{equation}
这里的 $C$ 是电容器的电容量。孤立导体与双极电容器的两式 $Q=CU$ 形式相同，但含义不同：孤立导体的 $U$ 是相对无穷远
或接地边界的电势，双极电容器的 $U$ 是两极之间的电势差。电容器的电容由两极的大小、形状、相对位置
以及两极间介质决定。

最简单的是平行板电容器。两块面积均为 $A$ 的平面导体板正对放置，板间距为 $d$。真实电场在中央近似
均匀，在边缘向外弯曲；当板的线度远大于 $d$ 时，边缘场只占很小部分，可以忽略。

图~\ref{ch20:fig:s35-4-2} 左侧标出面积 $A$ 和间距 $d$；右侧场线在板间中央近乎平行，只有边缘附近明显弯曲，因此下文匀强近似要求极板尺度远大于 $d$。

\par\noindent\begin{minipage}{\linewidth}
\parallelPlateFigure
\diagramcaption{平行板电容器的几何参数与实际电场。右图为无限延伸方向的截面模型：上板带正电、下板带负电；中央近似匀强，边缘场由同一电势边界求得。}\label{ch20:fig:s35-4-2}
\end{minipage}\par\medskip

忽略边缘效应时，每块板的面电荷密度为 $\sigma=Q/A$。单个无限大带电面的场强为
$\sigma/(2\varepsilon_0)$，两板之间同向叠加、外部反向抵消，所以
\[
  E=\frac{\sigma}{\varepsilon_0}=\frac{Q}{\varepsilon_0A},\qquad
  U=Ed=\frac{Qd}{\varepsilon_0A}.
\]
因此
\begin{equation}
  \boxed{C=\varepsilon_0\frac Ad}.
\end{equation}
板面积越大，相同电势差下能够存储的电荷越多；板距越大，相同电荷产生的电势差越大。因此 $C$ 正比于
$A$、反比于 $d$。

另一种常见结构是同轴圆柱电容器。内、外导体半径分别为 $a,b$，共同长度为 $h$。忽略两端边缘场，
在 $a<r<b$ 处取半径 $r$、长度 $h$ 的同轴圆柱高斯面。若内柱带 $+Q$，则
\[
  E(r)\,2\pi rh=\frac{Q}{\varepsilon_0},\qquad
  E(r)=\frac{Q}{2\pi\varepsilon_0h}\frac1r.
\]

图~\ref{ch20:fig:s35-4-3} 的横截面中，积分面半径必须满足 $a<r<b$；沿径向从 $a$ 积到 $b$，才能得到两导体间的电势差。

\par\noindent\begin{minipage}{\linewidth}
\coaxialCapacitorFigure
\diagramcaption{同轴圆柱形电容器。电场只存在于 $a<r<b$，并沿径向分布。}\label{ch20:fig:s35-4-3}
\end{minipage}\par\medskip

两极之间的电势差为
\[
  U=\int_a^bE(r)\,\mathrm dr
   =\frac{Q}{2\pi\varepsilon_0h}\ln\frac ba,
\]
所以
\begin{equation}
  \boxed{C=\frac{2\pi\varepsilon_0h}{\ln(b/a)}}.
\end{equation}
它正比于长度 $h$。外半径 $b$ 固定时，$a$ 越接近 $b$，$\ln(b/a)$ 越小，电容迅速增大。
若保持 $Q$ 不变，电势差减小，而内表面场 $Q/(2\pi\varepsilon_0ha)$ 反而随 $a$ 增大而减小；
若保持电压不变，场强的变化须代入 $Q=CU$ 后另作判断。公式假定圆柱足够长；实际电容器尺寸有限时，两端边缘场会修正结果。


\subsection{导体组的一般理论}


设有 $N$ 个形状与位置固定的导体。某一静电平衡状态下，第 $i$ 个导体的电荷和电势分别为
$Q_i,\varphi_i$；保持几何结构不变，再建立另一状态 $Q'_i,\varphi'_i$。由格林定理可得静电倒易
关系
\begin{equation}
  \boxed{\sum_iQ_i\varphi'_i=\sum_iQ'_i\varphi_i}.
\end{equation}
它把一个状态的电荷与另一个状态的电势交叉相乘再求和，两种交叉方式所得结果完全相同。
其条件是几何结构、介质分布和线性互易的介电响应在两状态中相同，无穷远边界项消失。
在导体外区域使用 $\boldsymbol E=-\nabla\varphi$、$\boldsymbol D=\varepsilon\boldsymbol E$，
有 $\boldsymbol D\cdot\nabla\varphi'=\boldsymbol D'\cdot\nabla\varphi$。
区域内无自由体电荷时，
\[
 \nabla\cdot(\varphi'\boldsymbol D-\varphi\boldsymbol D')=0.
\]
积分并用散度定理。区域在导体表面的外法线指入导体，故该处通量分别为 $-Q_i$、$-Q_i'$；
又因各导体等势，边界积分便给出
$-\sum_i\varphi_i'Q_i+\sum_i\varphi_iQ_i'=0$。
若区域另有固定点电荷，还须把它们的 $q\varphi'$ 与 $q'\varphi$ 项一起计入。

先看一个导体球与球外点 $P$。状态 A 中球带单位电荷，$P$ 点不放电荷；状态 B 中球不带净电荷，
$P$ 点放电荷 $q$。状态 A 在距离球心 $r$ 的 $P$ 点产生电势
$1/(4\pi\varepsilon_0r)$。把两状态代入倒易关系，状态 B 中球的电势立即得到
\[
  \varphi=\frac{q}{4\pi\varepsilon_0r}.
\]
这说明外点电荷使中性导体球发生复杂的表面电荷重排，但球的共同电势仍可不求分布而直接得到。

图~\ref{ch20:fig:s35-4-6} 中距离 $r$ 从球心量到 $P$；交换球上电荷与 $P$ 处电荷时，球的位置和半径均保持不变。

\par\noindent\begin{minipage}{\linewidth}
\greenSphereFigure
\diagramcaption{格林倒易定理的一个例子。两种状态共用同一导体球和同一点 $P$。}\label{ch20:fig:s35-4-6}
\end{minipage}\par\medskip

再取两个任意导体。状态 A 的量为 $Q_1,Q_2,\varphi_1,\varphi_2$，状态 B 的量带撇号。倒易关系成为
\begin{equation}
  \boxed{Q_1\varphi'_1+Q_2\varphi'_2
   =Q'_1\varphi_1+Q'_2\varphi_2}.
\end{equation}

图~\ref{ch20:fig:s35-4-7} 将两导体的电荷与电势一一配对；倒易关系比较的是同一几何下的两个状态，而不是把两导体互换位置。

\par\noindent\begin{minipage}{\linewidth}
\greenTwoConductorFigure
\diagramcaption{格林倒易定理的另一个例子。两导体形状和位置固定，只改变各自电荷。}\label{ch20:fig:s35-4-7}
\end{minipage}\par\medskip

几个特殊选择很有用：若状态 A 中仅导体 1 带电 $Q$，状态 B 中仅导体 2 带相同电荷 $Q$，则
$\varphi'_1=\varphi_2$，说明相互感应的影响具有对称性；若两状态中分别把一个导体接地、另一个保持
同一电势，可得两种情况下接地导体的感应总电荷相同。还可以选取单位电势状态，构造导体组的电容系数
矩阵，并由倒易定理证明矩阵是对称的。
还可直接处理一边定电荷、一边接地的混合条件：设状态 A 为 $Q_1=Q$、$\varphi_2=0$，
状态 B 为 $\varphi_1'=0$、$Q_2'=Q'$，则
\[
 Q_2\varphi_2'=Q_1'\varphi_1.
\]
这里两侧分别是“感应电荷乘另一状态的非零电势”，不能误删感应电荷项。
它与分别指定净电荷或指定电势的前两种选择，均由同一个两导体倒易式直接得到。
另一个混合选择更适合用于感应电荷的计算：状态 A 令 $\varphi_1=V_*$、$Q_2=0$；
状态 B 令 $\varphi_1'=0$、$Q_2'=Q_*$。代入倒易式后左侧两项均为零，故
\[
 0=Q_1'V_*+Q_*\varphi_2,\qquad
 \boxed{\frac{Q_1'}{Q_*}=-\frac{\varphi_2}{V_*}}.
\]
也就是说，在导体 2 中性时把导体 1 升到给定电势所得的导体 2 电势，
可反求导体 1 接地、向导体 2 加入给定电荷时导体 1 的感应电荷。
采用参考量 $V_*,Q_*$ 后，两边都是无量纲比值，不能直接把有不同单位的电荷与电势写成相等。


% EM structure and teaching revision 20260928
\subsection{从几何求电容的完整流程}
电容不是给定电荷本身的性质，而是导体几何、介质和边界条件共同决定的比例系数。两端电容器依次选定自由电荷 $+Q,-Q$，求介质中的场，再沿从正极到负极的路径算电压，最后消去 $Q$。
例如长为 $\ell$、内外半径 $a,b$ 的同轴圆柱，忽略端部效应，介电常数为 $\varepsilon$，则
\[
 E_r(r)=\frac{Q}{2\pi\varepsilon\ell r},\quad
 U=\int_a^b E_r\dd r=\frac{Q\ln(b/a)}{2\pi\varepsilon\ell},\quad
 C=\frac{2\pi\varepsilon\ell}{\ln(b/a)}.
\]
若间隙 $d=b-a\ll a$，利用 $\ln(1+d/a)\simeq d/a$ 得
$C\simeq\varepsilon(2\pi a\ell)/d$，正好退回展开为平行板后的近似。
多导体问题中，必须先说明哪两部分作为端口、其余导体是接地还是悬浮。一般关系为 $Q_i=\sum_j C_{ij}V_j$；改变悬浮导体的约束，相当于消去一组不同的未知电势，得到的端口等效电容也可能不同。


\subsection{专题例题与方法应用}

\begin{origexamplebox}{【统统算出来】}\label{ch20:ex:s35-4-1}
如图，最简单的双极电容为平行板电容器、球电容器和柱电容器，各参数如图。且平行板边长
$l\gg d$，柱电容器长度远大于极板间距。试利用高斯定理求解这三个电容器的电容。
\begin{center}
\begin{minipage}{.53\linewidth}\emParallelQuestion\end{minipage}
\begin{minipage}{.40\linewidth}\emSphericalQuestion\end{minipage}
\end{center}
\figcap{平行板与同心球形电容器的几何尺寸。}\label{ch20:fig:em-capacitors-question}
\begin{center}\begin{tikzpicture}[>=Latex,scale=.7]

 \draw[thick](0,0)circle(1.35)(0,0)circle(.55);
 \draw(0,0)--(140:.55)node[midway,above]{$a$};\draw(0,0)--(-35:1.35)node[midway,below]{$b$};
 \node[align=left,anchor=west]at(2.2,.25){同轴圆柱横截面\\轴向长度 $h$\\内、外导体半径 $a,b$};

\end{tikzpicture}\end{center}
\figcap{柱形电容器的横截面与长度定义。}\label{ch20:fig:em-source-thirty-fifth-section-1}
\end{origexamplebox}\par\nopagebreak[4]

图~\ref{ch20:fig:s35-4-10} 的板间距为 $d$，板面边长为 $l$，故用来包围极板电荷的高斯面有效面积为 $S=l^2$。

\par\noindent\begin{minipage}{\linewidth}
\exampleParallelFigure
\diagramcaption{例~\ref{ch20:ex:s35-4-1}的平行板电容器，$S=l^2$。}\label{ch20:fig:s35-4-10}
\end{minipage}\par\medskip

图~\ref{ch20:fig:s35-4-11} 中场仅分布于 $R_1<r<R_2$；两导体之间的电压须沿半径方向积分。

\par\noindent\begin{minipage}{\linewidth}
\exampleSphericalFigure
\diagramcaption{例~\ref{ch20:ex:s35-4-1}的同心球形电容器，内、外半径分别为 $R_1,R_2$。}\label{ch20:fig:s35-4-11}
\end{minipage}\par\medskip

\begin{xesolution}
三种结构都具有足够的对称性。统一方法是：先由高斯定理求 $E$，再沿场线积分求电势差 $U$，
最后使用 $C=Q/U$。

\textbf{(1) 平行板电容器。} 忽略边缘效应，两极之间
\[
 E=\frac{Q}{\varepsilon_0S},\qquad U=Ed=\frac{Qd}{\varepsilon_0S},
\]
故
\[
 \boxed{C_{\rm p}=\varepsilon_0\frac Sd=\varepsilon_0\frac{l^2}{d}}.
\]

\textbf{(2) 同心球形电容器。} 在 $R_1<r<R_2$ 取半径为 $r$ 的球面，
\[
 E(r)4\pi r^2=\frac Q{\varepsilon_0},\qquad
 E(r)=\frac{Q}{4\pi\varepsilon_0r^2}.
\]
所以
\[
 U=\int_{R_1}^{R_2}E(r)\,\mathrm dr
 =\frac{Q}{4\pi\varepsilon_0}\left(\frac1{R_1}-\frac1{R_2}\right),
\]
从而
\[
 \boxed{C_{\rm s}=4\pi\varepsilon_0\frac{R_1R_2}{R_2-R_1}}.
\]
令 $R_2\to\infty$，自然退化为孤立球的 $C=4\pi\varepsilon_0R_1$。

\textbf{(3) 同轴圆柱电容器。} 按题图取内、外半径 $a,b$、轴向长度 $h$，忽略端部效应，
\[
 E(r)2\pi rh=\frac Q{\varepsilon_0},\qquad
 U=\int_{a}^{b}\frac{Q}{2\pi\varepsilon_0hr}\,\mathrm dr
 =\frac{Q}{2\pi\varepsilon_0h}\ln\frac{b}{a}.
\]
因此
\[
 \boxed{C_{\rm c}=\frac{2\pi\varepsilon_0h}{\ln(b/a)}}.
\]
这里每一步都使用了实际存在电场的区域：平行板在两板之间，球形在两球面之间，圆柱形在两同轴柱面之间。
\end{xesolution}



\begin{origexamplebox}{【两个导体球】}\label{ch20:ex:s35-4-2}
两个导体球紧挨着，半径均为 $R$。试求这两个球作为一个导体的单极电容。
\par\noindent\begin{minipage}{\linewidth}
\touchingSpheresFigure
\diagramcaption{例~\ref{ch20:ex:s35-4-2}。两球相切并电连接，整个导体处于同一电势。}\label{ch20:fig:s35-4-12}
\end{minipage}\par\medskip
\end{origexamplebox}\par\nopagebreak[4]

图~\ref{ch20:fig:s35-4-12} 的接触点是两球的公共点；在计算电容时，两球应视为同一等势导体，而非两个独立电容器。


\begin{xesolution}
令两球共同电势为 $V$，接触点为坐标原点，两球心位于 $x=\pm R$。可用一列像电荷表示球外电势：
在每个球内、距接触点 $R/n$ 的轴线上放置
\[
 q_n=4\pi\varepsilon_0RV\frac{(-1)^{n-1}}n,
 \qquad n=1,2,\ldots
\]
下面验证这个序列为何成立。令 $K=4\pi\varepsilon_0RV$。
在右球球心 $x=R$ 放 $q_1=K$，其自身在右球面产生常电势 $V$。
左侧第 $n$ 枚电荷在 $x=-R/n$，距右球心为 $R+R/n$；关于右球反演，
\[
 q_{n+1}=-\frac{R}{R+R/n}q_n=-\frac n{n+1}q_n,\qquad
 x_{n+1}=R-\frac{R^2}{R+R/n}=\frac{R}{n+1}.
\]
因此左侧第 $n$ 枚与右侧第 $n+1$ 枚在右球面电势相消，只留下右球心的 $V$。
左右互换同理，故两个球面都是电势 $V$。接触点处取连续极限。
这两列像电荷产生的等势面恰包含两个相切球面。像电荷只是计算工具；它们的总和等于真实导体所带总电荷。

\par\noindent\begin{minipage}{\linewidth}
\touchingSphereImagesFigure
\diagramcaption{解图：相切双球的像电荷序列。两列交错符号的像电荷向接触点聚集。}
\end{minipage}\par\medskip

一个球对应的总电荷为
\[
 Q_1=\sum_{n=1}^{\infty}q_n
 =4\pi\varepsilon_0RV\left(1-\frac12+\frac13-\frac14+\cdots\right)
 =4\pi\varepsilon_0RV\ln2.
\]
两球对称，所以 $Q=2Q_1$。因此组合导体的单极电容为
\[
 \boxed{C=\frac QV=8\pi\varepsilon_0R\ln2}.
\]
该值小于两个相距很远的孤立球电容之和 $8\pi\varepsilon_0R$；接触后彼此的正电势抬高，使达到同一
电势所需的总电荷略少。
\end{xesolution}



\begin{origexamplebox}{【两个圆柱】}\label{ch20:ex:s35-4-3}
两根无穷长平行圆柱导体的半径均为 $R$，圆心距为 $L>2R$；两柱分别带等量异号线电荷。
试求这一双极电容器单位长度的电容。
\begin{center}\begin{tikzpicture}[>=Latex,scale=.8]

 \draw[thick](-2.1,0)circle(1)(2.1,0)circle(1);
 \fill(-2.1,0)circle(1pt);\fill(2.1,0)circle(1pt);
 \draw(-2.1,0)--++(55:1)node[midway,above left]{$R$};\draw(2.1,0)--++(55:1)node[midway,above left]{$R$};
 \draw[dashed](-2.1,0)--(-2.1,-1.6)(2.1,0)--(2.1,-1.6);
 \draw[<->](-2.1,-1.5)--(2.1,-1.5)node[midway,below]{$L$};
 \node at(-2.1,1.45){$-\lambda$};\node at(2.1,1.45){$+\lambda$};

\end{tikzpicture}\end{center}
\figcap{无穷长平行导体圆柱的横截面，半径与圆心距如图。}\label{ch20:fig:em-source-thirty-fifth-section-3}
\end{origexamplebox}\par\nopagebreak[4]

\par\noindent\begin{minipage}{\linewidth}
\twoCylindersFigure
\diagramcaption{例~\ref{ch20:ex:s35-4-3}。两圆柱截面是双极坐标中的两条等坐标圆。}\label{ch20:fig:s35-4-13}
\end{minipage}\par\medskip

\begin{xesolution}
先说明为什么圆柱截面可由两条虚拟线电荷构造。令 $d=L/2$、$a=\sqrt{d^2-R^2}$，
在图~\ref{ch20:fig:s35-4-13}中 $x=\pm a$ 放虚拟线电荷 $\pm\lambda$。它们位于各自圆柱内部，
只在柱外使用对应电势 $\varphi=(\lambda/2\pi\varepsilon_0)\ln(r_-/r_+)$。
取 $\xi=\ln(r_-/r_+)$ 为常量，配方得到圆心 $a\coth\xi$、半径 $a/|\sinh\xi|$；
因此选择 $\cosh\xi_0=d/R$ 就恰好得到两条真实圆柱边界。
每柱外围的高斯面包围一根虚拟线电荷，故计算出的 $\lambda$ 就是真实单位长度总电荷。

这是二维静电问题。采用双极坐标 $(\xi,\eta)$，其尺度因子为
\[
 h_\xi=h_\eta=\frac a{\cosh\xi-\cos\eta}.
\]
两个圆柱面可取为 $\xi=\pm\xi_0$。由双极坐标几何关系
\[
 R=\frac a{\sinh\xi_0},\qquad
 \frac L2=a\coth\xi_0,
\]
消去 $a$ 得
\[
 \cosh\xi_0=\frac{L}{2R},\qquad
 \xi_0=\operatorname{arcosh}\frac{L}{2R}.
\]
令两柱电势分别为 $+U/2$ 和 $-U/2$。满足拉普拉斯方程及边界条件的电势是
\[
 \varphi(\xi)=\frac{U}{2\xi_0}\xi.
\]
在 $\xi=\xi_0$ 的圆柱面上，法向电场大小
$E_n=(1/h_\xi)\,\partial\varphi/\partial\xi$，弧元为
$\mathrm ds=h_\eta\mathrm d\eta$。因此单位长度电荷量
\[
 \lambda=\varepsilon_0\int_0^{2\pi}E_n\,\mathrm ds
 =\varepsilon_0\int_0^{2\pi}\frac{U}{2\xi_0}\,\mathrm d\eta
 =\frac{\pi\varepsilon_0U}{\xi_0}.
\]
单位长度电容为
\[
 \boxed{C'=\frac\lambda U
 =\frac{\pi\varepsilon_0}{\operatorname{arcosh}(L/2R)}
 =\frac{\pi\varepsilon_0}
 {\ln\!\left[\dfrac{L+\sqrt{L^2-4R^2}}{2R}\right]}}.
\]
当 $L\gg R$ 时，$C'\simeq\pi\varepsilon_0/\ln(L/R)$。
\end{xesolution}



\section{电介质、极化与界面条件}

\subsection{电介质与相对介电常数}


两极之间是真空的电容器称为真空电容器。实际装置往往在两极之间填充橡胶、玻璃、石蜡、纸等不导电
物质，这些绝缘材料称为电介质。它们不允许自由电荷像导体中那样长距离迁移，因此不能产生导体式的
静电屏蔽；但内部正负电荷可以发生微小相对位移，从而改变电场。

对于均匀、线性、各向同性介质，在几何结构不变且全部有效电场区域被填满的近似下，把同一种电介质完全填满不同形状的真空电容器，电容都按相同倍数增大。这个只由材料决定的
倍数称为相对介电常数 $\varepsilon_r$。若原真空电容为 $C_0$，完全填充后
\begin{equation}
  \boxed{C=\varepsilon_rC_0}.
\end{equation}

图~\ref{ch20:fig:s35-4-8}中先充电再断开电源，接入高输入阻抗静电计；要求静电计输入电容远小于待测电容，
且观测时间远短于漏电时间常数。这样电容器自由电荷近似不变。
右图的介质紧贴两块极板并覆盖全部正对面积，不能留有额外空气层；右侧表针较小的偏转表示 $U$ 降低。
\par\noindent\begin{minipage}{\linewidth}
\dielectricMeterFigure
\diagramcaption{插入电介质使电容增大的实验情景。插入过程中电容器与外界绝缘，$Q_0$ 不变，因此高输入阻抗静电计的电压读数降低。}\label{ch20:fig:s35-4-8}
\end{minipage}\par\medskip

\begin{table}[htbp]
\centering
\small
\booktablecaption{室温下典型电介质的相对介电常数与击穿场强}\label{ch20:tab:s35-4-1}
\begin{tabular}{lcc}
\toprule
材料 & 相对介电常数 $\varepsilon_r$ & 击穿场强 $E_{\rm b}/(\mathrm{V\,m^{-1}})$\\
\midrule
干燥空气 & $1.00059$ & $3\times10^6$\\
石英 & $3.78$ & $24\times10^6$\\
纸 & $3.7$ & $16\times10^6$\\
玻璃 & $5.6$ & $14\times10^6$\\
水 & $80$ & --\\
橡胶 & $6.7$ & $12\times10^6$\\
尼龙 & $3.4$ & $14\times10^6$\\
聚苯乙烯 & $2.56$ & $24\times10^6$\\
\bottomrule
\end{tabular}
\end{table}

所有表中材料都有 $\varepsilon_r>1$。在自由电荷 $Q$ 固定、介质把电场区域完全填满时，电容增大
$\varepsilon_r$ 倍，电势差和电场相应变为真空时的 $1/\varepsilon_r$。电场的减小来自介质内部原子或
分子的极化，而不是自由电荷被导走。


% EM teaching enrichment 20260927
\paragraph{部分填充为何不能直接乘一个常数}
面积 $A$、间距 $d$ 的平行板中，介质沿电场方向占厚度 $t$，余下是真空，忽略边缘效应。各层自由电荷对应的 $D_n$ 相同，电压分别相加，等效为串联：
\[
 C=\frac{\varepsilon_0A}{d-t+t/\varepsilon_r}.
\]
若介质贯穿板距、只占正对面积 $A_1$，其余面积 $A-A_1$ 为真空，两区电压相同，等效为并联：
\[
 C=\frac{\varepsilon_0}{d}\bigl[\varepsilon_rA_1+(A-A_1)\bigr].
\]
取 $t=0,d$ 或 $A_1=0,A$ 均应回到真空与全填充极限。判断串并联应沿场线看各区共享电压还是共享法向自由电荷。


\subsection{介质极化与电位移矢量}


经典图像中，可把绝缘体看成许多中性原子或分子。外电场作用时，正、负电荷中心发生微小分离，形成
电偶极子；原来具有永久偶极矩的分子也会出现取向。偶极矩总体趋向外场方向，其产生的附加电场却在
介质内部反向抵消一部分外场，因此实际电场减小。

图~\ref{ch20:fig:s35-4-9} 左侧每个偶极矩由负端指向正端，右侧由束缚面电荷产生的场则指向左，因而抵消一部分向右的外加场。

\par\noindent\begin{minipage}{\linewidth}
\polarizationFigure
\diagramcaption{外场作用下电介质的极化。左图是微观偶极矩取向，右图是等效束缚面电荷及其反向极化场。}\label{ch20:fig:s35-4-9}
\end{minipage}\par\medskip

若每个分子的平均电偶极矩为矢量 $\boldsymbol p$，单位体积内分子数为 $N$，均匀极化时定义极化强度
\begin{equation}
  \boxed{\boldsymbol P=N\boldsymbol p}.
\end{equation}
$\boldsymbol P$ 是单位体积内电偶极矩的矢量和。均匀极化体内部的偶极电荷相互抵消，只在边界留下束缚
面电荷；一般情况下束缚体电荷密度为 $\rho_{\rm b}=-\boldsymbol\nabla\cdot\boldsymbol P$，束缚面电荷
密度为 $\sigma_{\rm b}=\boldsymbol P\cdot\hat{\boldsymbol n}$。

为把自由电荷和极化束缚电荷分开，引入电位移矢量
\begin{equation}
  \boxed{\boldsymbol D=\varepsilon_0\boldsymbol E+\boldsymbol P}.
\end{equation}
它是对材料响应进行记账的辅助场，并不是另一种独立的基本相互作用。由
$\boldsymbol\nabla\cdot\boldsymbol E=(\rho_{\rm f}+\rho_{\rm b})/\varepsilon_0$ 和
$\rho_{\rm b}=-\boldsymbol\nabla\cdot\boldsymbol P$ 可得
\begin{equation}
  \boxed{\oint_S\boldsymbol D\cdot\mathrm d\boldsymbol S=Q_{\rm f}},
\end{equation}
其中 $Q_{\rm f}$ 只包括曲面内的自由电荷，不把介质极化产生的束缚电荷重复计算。

外场不太强时，许多电介质的响应近似线性。单个原子的诱导偶极矩可写为
\begin{equation}
  \boxed{\boldsymbol p=\alpha\boldsymbol E_{\rm loc}},
\end{equation}
其中 $\alpha$ 为原子极化率，$\boldsymbol E_{\rm loc}$ 是原子所在位置的局域场。宏观各向同性线性介质中
$\boldsymbol P=\varepsilon_0\chi_e\boldsymbol E$，于是
\begin{equation}
  \boxed{\boldsymbol D=\varepsilon_0\boldsymbol E+\boldsymbol P
  =\varepsilon_r\varepsilon_0\boldsymbol E},
\end{equation}
并有 $\varepsilon_r=1+\chi_e$。因此相对介电常数反映了微观粒子在外场中的可极化程度。

若考虑球形空穴给出的局域场修正，对于数密度为 $N$、极化率为 $\alpha$ 的非极性稀薄介质，可得到
克劳修斯--莫索蒂关系
\begin{equation}
  \boxed{\alpha=\frac{3\varepsilon_0}{N}\,
  \frac{\varepsilon_r-1}{\varepsilon_r+2}}.
\end{equation}
它说明可以由宏观测得的 $\varepsilon_r$ 反推微观极化率。关系式假设各向同性、线性、粒子间关联较弱；
对强极性液体、晶体各向异性或强场情形，需要更精细的模型。

% EM teaching enrichment 20260927
\paragraph{界面条件把宏观场与束缚电荷接起来}
取法线从介质1指向介质2，薄盒高斯面给出
$D_{2n}-D_{1n}=\sigma_f$；静电场的小矩形环路给出 $E_{2t}=E_{1t}$。
无自由面电荷时连续的是 $D_n$，一般不是 $E_n$。
例如场垂直于真空与介质界面，$D_n=D$，则 $E_{\rm vac}=D/\varepsilon_0$、$E_{\rm diel}=D/(\varepsilon_r\varepsilon_0)$。
介质中 $P=D(1-1/\varepsilon_r)$，端面束缚电荷为 $\pm P$。这组关系同时说明了自由电荷不变、场减弱与束缚电荷出现的原因；不能把束缚电荷再次加进 $\oint\boldsymbol D\cdot\dd\boldsymbol S$ 的右侧。




% EM structure and teaching revision 20260928
\subsection{分层介质：同电压还是同自由电荷}
对线性、各向同性、无明显边缘效应的平板结构，介质分层方向决定等效连接。若两层沿电场方向依次排列，厚度为 $d_1,d_2$，界面无自由电荷，则法向位移 $D=Q/S$ 在两层中相同，电压为
\[
 U=E_1d_1+E_2d_2=\frac QS\left(\frac{d_1}{\varepsilon_1}+\frac{d_2}{\varepsilon_2}\right),\qquad
 C=\frac S{d_1/\varepsilon_1+d_2/\varepsilon_2}.
\]
若介质左右并排、共同跨接两块极板，则两部分电压相同，电荷相加，得到
$C=(\varepsilon_1S_1+\varepsilon_2S_2)/d$。不能仅凭“各占一半”选择串联或并联。
取 $d_1=d_2=d/2$、$\varepsilon_1=\varepsilon_0$、$\varepsilon_2=4\varepsilon_0$，沿场分层得到 $C=1.6\varepsilon_0S/d$，等面积并排则为 $2.5\varepsilon_0S/d$。低介电常数层中的场更强，也更可能先达到击穿阈值。
上述计算假定介质响应线性且几何保持不变。若要求介质受力，还要说明电源是否连接，并把电源的能量交换纳入功的计算。


\subsection{专题例题与方法应用}

\begin{origexamplebox}{【改造电容器】}\label{ch20:ex:s35-4-4}
如图，一个平行板电容器的极板面积为 $S$、间距为 $d$。

(1) 在左边一半插入相对介电常数为 $\varepsilon_r$ 的介质，电容变为多大？变大还是变小？

(2) 在上边一半插入相对介电常数为 $\varepsilon_r$ 的介质，电容变为多大？变大还是变小？
\emModifiedQuestion\figcap{介质分别占一半正对面积或一半板间厚度。}\label{ch20:fig:em-modified-question}
\end{origexamplebox}\par\nopagebreak[4]

图~\ref{ch20:fig:s35-4-14} 左图的两部分共有相同板间电压；右图的两层沿电场方向依次排列，分别对应并联与串联。

\par\noindent\begin{minipage}{\linewidth}
\modifiedCapacitorFigure
\diagramcaption{例~\ref{ch20:ex:s35-4-4}。左图按面积分区，应视为并联；右图按厚度分层，应视为串联。}\label{ch20:fig:s35-4-14}
\end{minipage}\par\medskip

\begin{xesolution}
原电容为 $C_0=\varepsilon_0S/d$。

\textbf{(1) 按面积分区。} 两个区域的电势差均为同一个 $U$，相当于并联：
\[
 C_1=\varepsilon_r\varepsilon_0\frac{S/2}{d},\qquad
 C_2=\varepsilon_0\frac{S/2}{d}.
\]
故
\[
 \boxed{C=C_1+C_2=\frac{\varepsilon_r+1}{2}C_0}.
\]

\textbf{(2) 按厚度分层。} 电场依次穿过厚度均为 $d/2$ 的介质层和真空层，自由面电荷相同，相当于串联：
\[
 C_1=\varepsilon_r\varepsilon_0\frac S{d/2}=2\varepsilon_rC_0,
 \qquad C_2=\varepsilon_0\frac S{d/2}=2C_0.
\]
所以
\[
 \boxed{C=\frac{C_1C_2}{C_1+C_2}
 =\frac{2\varepsilon_r}{\varepsilon_r+1}C_0}.
\]
因为 $\varepsilon_r>1$，两种改造都使电容增大；但两者通常不相等。判断“并联还是串联”的依据不是
图画在左边还是上边，而是各区域共享电压还是依次承受电压。
\end{xesolution}



\begin{origexamplebox}{【匀强电场中的电介质球】}\label{ch20:ex:s35-4-6}
半径为 $R$、相对介电常数为 $\varepsilon_r$ 的电介质球放在匀强电场
$\boldsymbol E_0=E_0\hat{\boldsymbol z}$ 中，求空间中的电场。
\end{origexamplebox}\par\nopagebreak[4]

\par\noindent\begin{minipage}{\linewidth}
\dielectricSphereFieldFigure
\diagramcaption{解图：$\varepsilon_r=4$ 时的电介质球，$+z$ 沿外场。线表示场的方向，线密度不用于跨介质比较 $E$ 的大小；界面处切向 $E$、法向 $D$ 连续。}
\end{minipage}\par\medskip

\begin{xesolution}
问题关于 $z$ 轴对称，且远处电势应趋于 $-E_0r\cos\theta$。球内电势在原点必须有限，球外附加项
必须在无穷远消失，因此设
\[
 \varphi_{\rm in}=-Ar\cos\theta,
 \qquad
 \varphi_{\rm out}=-E_0r\cos\theta+\frac{B\cos\theta}{r^2}.
\]
球面 $r=R$ 上没有自由面电荷，故电势连续、法向电位移连续：
\[
 \varphi_{\rm in}=\varphi_{\rm out},\qquad
 \varepsilon_r\varepsilon_0E_{r,\rm in}=\varepsilon_0E_{r,\rm out}.
\]
比较两侧 $\cos\theta$ 的系数，第一条边界条件给出 $A=E_0-B/R^3$；
由 $E_r=-\partial\varphi/\partial r$，第二条给出
$\varepsilon_r A=E_0+2B/R^3$。将第一式代入第二式：
\[
 (\varepsilon_r+2)\frac B{R^3}=(\varepsilon_r-1)E_0.
\]
再回代即得
\[
 A=\frac{3E_0}{\varepsilon_r+2},\qquad
 B=\frac{\varepsilon_r-1}{\varepsilon_r+2}E_0R^3.
\]
所以球内电场为
\[
 \boxed{\boldsymbol E_{\rm in}=\frac{3}{\varepsilon_r+2}\boldsymbol E_0}.
\]
球外电场分量为
\[
 \boxed{\begin{aligned}
 E_r&=E_0\cos\theta+
 2\frac{\varepsilon_r-1}{\varepsilon_r+2}E_0\frac{R^3}{r^3}\cos\theta,\\
 E_\theta&=-E_0\sin\theta+
 \frac{\varepsilon_r-1}{\varepsilon_r+2}E_0\frac{R^3}{r^3}\sin\theta.
 \end{aligned}}
\]
附加项正是位于球心、方向沿 $\boldsymbol E_0$ 的偶极子场。等效偶极矩为
\[
 \boldsymbol p_{\rm eff}=4\pi\varepsilon_0R^3
 \frac{\varepsilon_r-1}{\varepsilon_r+2}\boldsymbol E_0.
\]
\end{xesolution}



\section{电容网络、储能与充放电}

\subsection{电容的连接}


多个电容器构成的系统，在外部端口看来可以由一个等效电容代替。并联时，所有电容器接在相同的两个
节点之间，因此两端电势差相同。两个电容 $C_1,C_2$ 并接到电压 $U$ 上，有
$Q_1=C_1U,Q_2=C_2U$；电源搬运的总电荷是
\[
  Q=Q_1+Q_2=(C_1+C_2)U.
\]

图~\ref{ch20:fig:s35-4-4} 的两支路连接同一对节点，所以每个电容器的电压均为 $U$，而电源提供的电荷量等于各支路电荷量之和。

\par\noindent\begin{minipage}{\linewidth}
\capacitorParallelFigure
\diagramcaption{电容并联。各支路电压相同，等效电容由总电荷与端电压之比确定。}\label{ch20:fig:s35-4-4}
\end{minipage}\par\medskip

于是多个电容并联时
\begin{equation}
  \boxed{C_{\mathrm{eq}}=C_1+C_2+\cdots}.
\end{equation}

串联时，相邻电容器之间的连接导体与外界绝缘。如果初始不带净电荷，它在充电后仍保持净电荷为零；
因此相邻两块极板带等量异号电荷，各电容器的电荷量大小均为 $Q$。总电压为
\[
  U=U_1+U_2+\cdots=Q\left(\frac1{C_1}+\frac1{C_2}+\cdots\right).
\]

图~\ref{ch20:fig:s35-4-5} 中两电容器之间的导线与两片内侧极板组成孤立中性导体，因此内侧两片极板的电荷等量异号。

\par\noindent\begin{minipage}{\linewidth}
\capacitorSeriesFigure
\diagramcaption{电容串联。中间浮置导体净电荷守恒，故各电容器的电荷量大小相同。}\label{ch20:fig:s35-4-5}
\end{minipage}\par\medskip

所以
\begin{equation}
  \boxed{\frac1{C_{\mathrm{eq}}}=\frac1{C_1}+\frac1{C_2}+\cdots}.
\end{equation}
并联增加可存储电荷的有效面积，等效电容增大；串联增加等效极板距离，等效电容小于任一个串联电容。
电容的串、并联规律恰与电阻的串、并联规律相反，而形式上与弹簧劲度系数的连接规律相似。


% EM teaching enrichment 20260927
\paragraph{中间导体预先带电时}
两电容 $C_1,C_2$ 串接，左端电势 $U$、右端为零，中间浮置导体预带净电荷 $q_0$。设中点电势为 $V$，该导体相邻两片极板的电荷相加为
\[
 C_1(V-U)+C_2V=q_0,
 \qquad V=\frac{C_1U+q_0}{C_1+C_2}.
\]
只有 $q_0=0$ 才恢复通常的等电量串联关系。断开或重接电容后，应先圈出仍然孤立的导体并守恒其净电荷，再写每只电容的电压关系。


\subsection{电容储能}


容量为 $C$ 的电容器最初不带电。当两极已有电荷量 $q$ 时，电势差为 $u=q/C$。从负极向正极再搬运
电量 $\mathrm dq$，外界需要输入
\begin{equation}
  \mathrm dW=u\,\mathrm dq=\frac qC\,\mathrm dq.
\end{equation}
把电荷从 $0$ 积分到 $Q$，得到
\begin{equation}
  \boxed{W=\int_0^Q\frac qC\,\mathrm dq
  =\frac{Q^2}{2C}=\frac12QU=\frac12CU^2}.
\end{equation}
该能量以电场形式储存在两极之间。使用三个等价形式时必须说明控制量：若 $Q$ 固定，改变几何形状时
宜用 $Q^2/(2C)$；若接在理想恒压源上，电源也会交换能量，单看电容场能不足以直接判断受力。


% EM teaching enrichment 20260927
\paragraph{恒压与恒电荷：同一吸力的两种能量账}
平行板间距取 $d$，增大 $d$ 的方向为广义坐标正向。固定 $Q$ 时
$W=Q^2d/(2\varepsilon_0A)$，故
$F_d=-\partial W/\partial d=-Q^2/(2\varepsilon_0A)$，负号表示吸引。
固定 $U$ 时，电源功 $\dd W_s=U\dd Q=U^2\dd C$，场能改变为 $\dd W=U^2\dd C/2$，两者之差才是电力机械功。因此
\[
 F_d=\frac12U^2\frac{\dd C}{\dd d}
 =-\frac{\varepsilon_0AU^2}{2d^2}.
\]
两种表达用 $Q=CU$ 可相互转化。也可由表面电荷只受“另一块板”所生场得到压力 $p=\varepsilon_0E^2/2$，直接用总场 $\sigma/\varepsilon_0$ 乘本板电荷会把力算大一倍。


% EM structure and teaching revision 20260928
\subsection{充电过程为何不能只算末态储能}
电容通过电阻 $R$ 接到恒压源 $U_0$，初始电压为零。由 $U_0=Ri+u_C$ 与 $i=C\dd u_C/\dd t$ 得
\[
 u_C=U_0(1-e^{-t/(RC)}),\qquad i=\frac{U_0}{R}e^{-t/(RC)}.
\]
电源总供能为 $\int_0^\infty U_0i\dd t=CU_0^2$，电容末态只储存 $CU_0^2/2$，另半转化为焦耳热。减小 $R$ 会增大峰值电流、缩短过程，却不改变理想模型中总焦耳热。
两个初始电压为 $U_1,U_2$ 的电容 $C_1,C_2$ 同名端相连，且体系无外接电源时，总电荷守恒给出
\[
 U_f=\frac{C_1U_1+C_2U_2}{C_1+C_2},\qquad
 W_i-W_f=\frac12\frac{C_1C_2}{C_1+C_2}(U_1-U_2)^2.
\]
能量减少量不会因为连接线画成理想导线就自动消失。实际过程还可能涉及导线电感、电磁辐射及介质损耗；若严格设为无损LC系统，则会振荡而不自行停在上述平衡末态。先说明终态能否建立，再使用末态电荷约束。


\subsection{专题例题与方法应用}

\begin{origexamplebox}{【奇妙的吸引力】}\label{ch20:ex:s35-4-5}
把相对介电常数为 $\varepsilon_r$ 的电介质部分插入平行板电容器。极板长为 $a$、宽为 $b$、间距为
$d$，介质插入长度为 $x$，如图。利用虚功法求介质受到的力：

(1) 电容器电荷 $Q$ 恒定；

(2) 电容器接在电压为 $U$ 的理想电源上。

两问均用电容器内的电场 $E$ 表示结果。
\emInsertionQuestion\figcap{介质插入长度与电容器尺寸；受力方向待求。}\label{ch20:fig:em-insertion-question}
\end{origexamplebox}\par\nopagebreak[4]

图~\ref{ch20:fig:s35-4-15} 中 $x$ 向右增加，表示介质插入更深；因此由能量求得的正 $F$ 指向右。

\par\noindent\begin{minipage}{\linewidth}
\dielectricInsertionFigure
\diagramcaption{例~\ref{ch20:ex:s35-4-5}。介质区和真空区并联，电场方向竖直，吸引力沿插入方向。}\label{ch20:fig:s35-4-15}
\end{minipage}\par\medskip

\begin{xesolution}
忽略边缘场。插入部分与未插入部分共享相同板间电压，因此
\[
 C(x)=\varepsilon_r\varepsilon_0\frac{bx}{d}
      +\varepsilon_0\frac{b(a-x)}d
 =\frac{\varepsilon_0b}{d}\bigl[a+(\varepsilon_r-1)x\bigr],
\]
从而
\[
 \frac{\mathrm dC}{\mathrm dx}
 =\frac{\varepsilon_0b}{d}(\varepsilon_r-1)>0.
\]

\textbf{(1) 定电荷。} 电容场能 $W=Q^2/(2C)$，虚位移 $\mathrm dx$ 时
\[
 F=-\frac{\mathrm dW}{\mathrm dx}
 =\frac{Q^2}{2C^2}\frac{\mathrm dC}{\mathrm dx}.
\]
又 $Q/C=U=Ed$，故
\[
 \boxed{F=\frac12\varepsilon_0(\varepsilon_r-1)b d\,E^2}.
\]

\textbf{(2) 定电压。} 电容器本身的场能是 $CU^2/2$，但介质移动时电源还要搬运电荷
$\mathrm dQ=U\,\mathrm dC$。把电容器与电源作为整体，其与机械坐标有关的势能是
\[
 \mathcal G=\frac12CU^2-UQ=-\frac12CU^2.
\]
因此
\[
 F=-\frac{\mathrm d\mathcal G}{\mathrm dx}
 =\frac12U^2\frac{\mathrm dC}{\mathrm dx}
 =\boxed{\frac12\varepsilon_0(\varepsilon_r-1)b d\,E^2}.
\]
两问用瞬时场强 $E$ 表示时形式相同，方向都指向增大 $x$ 的方向，即介质被吸入电容器。定电压情形若
遗漏电源做功，会错误地得到相反的力。
\end{xesolution}



\begin{origexamplebox}{【能量怎么多了】}\label{ch20:ex:s35-4-7}
线性电介质中的电场能量密度为
$u=\tfrac12\varepsilon_r\varepsilon_0E^2$，而真空中在相同场强下只有
$u_0=\tfrac12\varepsilon_0E^2$。试提出一个物理图像解释这个现象。
\end{origexamplebox}\par\nopagebreak[4]

\par\noindent\begin{minipage}{\linewidth}
\polarizationEnergyFigure
\diagramcaption{解图：极化可类比为微观正负电荷之间的“弹簧”被拉伸；建立同样场强时还要为极化自由度输入能量。}
\end{minipage}\par\medskip

\begin{xesolution}
线性介质中
\[
 \boldsymbol D=\varepsilon_0\boldsymbol E+\boldsymbol P,
 \qquad
 u=\frac12\boldsymbol E\cdot\boldsymbol D
 =\frac12\varepsilon_0E^2+\frac12\boldsymbol E\cdot\boldsymbol P.
\]
第一项是真空电场能；第二项是建立极化时外界对束缚正、负电荷做功所对应的能量。可以把原子内部的
正负电荷中心想成由微观弹性约束连接：外场把它们拉开，形成偶极矩，同时在材料的极化自由度中储存
能量。因此在\emph{相同实际场强 $E$} 下，介质比真空多出
\[
 u-u_0=\frac12\boldsymbol E\cdot\boldsymbol P
 =\frac12(\varepsilon_r-1)\varepsilon_0E^2.
\]
这并不意味着插入介质后能量一定增加。若自由电荷 $Q$ 固定，插入介质使 $E$ 减小，电容器总能量
$Q^2/(2C)$ 反而降低，差值转化为机械功；若电压固定，电源继续输送电荷，场能增加，增加的能量来自
电源。必须先说明固定的是 $E$、$Q$ 还是 $U$。
\end{xesolution}



\begin{origexamplebox}{【充电过程】}\label{ch20:ex:s35-4-8}
用电动势为 $\mathcal E$ 的电源给电容 $C$ 充电。

(1) 若电路中有电阻 $R$，直至充电完成，电源做功多少，发热多少？

(2) 若电阻 $R\to0^+$，电源做功多少，发热多少？

(3) 如果电路中完全没有电阻，答案会有什么变化？
\end{origexamplebox}\par\nopagebreak[4]

\par\noindent\begin{minipage}{\linewidth}
\chargingEnergyFigure
\diagramcaption{解图：理想恒压源经任意正电阻给电容充电时，电源做功的一半进入电场，另一半成为焦耳热。}
\end{minipage}\par\medskip

\begin{xesolution}
充电完成后 $Q=C\mathcal E$。电源每搬运电量 $\mathrm dq$ 做功
$\mathcal E\,\mathrm dq$，所以总功
\[
 W_{\rm s}=\int_0^Q\mathcal E\,\mathrm dq
 =\mathcal EQ=C\mathcal E^2.
\]
电容器最终储能
\[
 W_C=\frac12C\mathcal E^2.
\]
由能量守恒，电阻发热
\[
 \boxed{Q_R=W_{\rm s}-W_C=\frac12C\mathcal E^2}.
\]

也可直接积分验证。$RC$ 电路电流
$i(t)=(\mathcal E/R)e^{-t/(RC)}$，因此
\[
 Q_R=\int_0^\infty i^2R\,\mathrm dt
 =\frac{\mathcal E^2}{R}\int_0^\infty e^{-2t/(RC)}\,\mathrm dt
 =\frac12C\mathcal E^2.
\]
结果与 $R$ 的大小无关。故 $R\to0^+$ 时，充电时间趋于零、瞬时功率趋于无穷，但积分后的电源功和
发热仍分别为 $C\mathcal E^2$ 与 $C\mathcal E^2/2$。

若令 $R$ \emph{严格等于零}，同时又把电感、辐射和电源内阻全部忽略，理想模型在接通瞬间要求电容
电压跳变并产生冲激电流，问题本身不再自洽。真实电路总有电感或辐射损耗；若只保留电感 $L$ 而没有
任何耗散，系统会发生 $LC$ 振荡，电容也不会不可逆地停在“充电完成”状态。因此第(2)问的极限不能
简单替代第(3)问的严格零电阻模型。
\end{xesolution}



\begin{origexamplebox}{【电容充电】}\label{ch20:ex:s35-4-9}
如图，三个电容器初始均不带电，连接成 $Y$ 型，三条外支路的电势分别为
$\varphi_1,\varphi_2,\varphi_3$。求三个电容器外侧极板所带电荷。
\par\noindent\begin{minipage}{\linewidth}
\yCapacitorFigure
\diagramcaption{例~\ref{ch20:ex:s35-4-9}。三个电容器的内侧极板连接成一个浮置中性节点。}\label{ch20:fig:s35-4-16}
\end{minipage}\par\medskip
\end{origexamplebox}\par\nopagebreak[4]

图~\ref{ch20:fig:s35-4-16} 的三个内侧极板连接在同一节点上，电势均为 $V$；该节点与外界不导通，所以其总电荷保持为零。


\begin{xesolution}
设中心浮置节点电势为 $V$。第 $i$ 个电容器外侧极板的电荷为
\[
 q_i=C_i(\varphi_i-V).
\]
中心节点由三个内侧极板和连接导线组成，初始净电荷为零，又与外界绝缘，所以始终满足
\[
 \sum_{i=1}^3C_i(V-\varphi_i)=0.
\]
由此
\[
 \boxed{V=\frac{C_1\varphi_1+C_2\varphi_2+C_3\varphi_3}
 {C_1+C_2+C_3}}.
\]
代回得
\[
 \boxed{q_i=C_i\left(\varphi_i-
 \frac{\sum_jC_j\varphi_j}{\sum_jC_j}\right)}.
\]
显式写出例如
\[
 q_1=\frac{C_1}{C_1+C_2+C_3}
 \bigl[C_2(\varphi_1-\varphi_2)+C_3(\varphi_1-\varphi_3)\bigr],
\]
其余两式循环置换指标即可。可检查 $q_1+q_2+q_3=0$，这与整个三电容系统初始总电荷为零一致。
\end{xesolution}



\begin{origexamplebox}{【电容串并联】}\label{ch20:ex:s35-4-10}
证明两个电容器串联或并联时都可以等效成一个电容器，求出等效公式，并讨论更多电容器串、并联时
是否也可以等效。
\end{origexamplebox}\par\nopagebreak[4]

\begin{xesolution}
等效的含义是：从外部两个端口看，在任意端电压 $U$ 下，系统从电源取得的总电荷 $Q$ 与某个单电容
$C_{\rm eq}$ 完全相同，即 $Q=C_{\rm eq}U$。

\textbf{并联：} 各电容两端电压均为 $U$，总电荷
\[
 Q=\sum_iQ_i=\sum_iC_iU,
\]
故
\[
 \boxed{C_{\rm eq}=\sum_iC_i}.
\]

\textbf{串联：} 每个中间浮置节点初始净电荷为零，因此各电容器电荷量大小均为 $Q$。总电压
\[
 U=\sum_iU_i=Q\sum_i\frac1{C_i},
\]
故
\[
 \boxed{\frac1{C_{\rm eq}}=\sum_i\frac1{C_i}}.
\]
结论可直接推广到任意多个纯并联或纯串联电容。由串、并联组成的两端网络可逐级化简；更一般的线性
两端电容网络，只要内部节点初始无净电荷，也存在一个等效电容，但未必能只用一次串并联公式求出，
需要先列内部节点电荷守恒方程。若网络有三个以上独立外端口，则通常需要电容系数矩阵，不能用一个
标量电容完整描述。
\end{xesolution}



\begin{origexamplebox}{【超大电容】}\label{ch20:ex:s35-4-11}
开尔文曾把大气层构造为一个电容器模型：地球表面是一个极板，带有
$5\times10^5\,\mathrm C$ 的负电荷；把距地表 $5\,\mathrm{km}$ 的同心球面作为带正电的等效外极板，两极之间的大气层厚度为 $5\,\mathrm{km}$。

(1) 求这个球形电容器的电容；

(2) 求地球表面的电场能量密度以及球形电容器的总能量；

(3) 已知空气电阻率为 $3\times10^{13}\,\Omega\!\cdot\!\mathrm m$，求两极间大气层的电阻；

(4) 大气电容与电阻构成 $RC$ 放电回路，求时间常数；

(5) 若平均一次雷电向地面补充 $25\,\mathrm C$ 电荷，为维持上述电荷量，全球每天平均需要多少次雷电？
\end{origexamplebox}\par\nopagebreak[4]

\par\noindent\begin{minipage}{\linewidth}
\earthCapacitorFigure
\diagramcaption{解图：地球表面与高空等效层构成同心球形电容器；大气同时提供漏电电阻。}
\end{minipage}\par\medskip

\begin{xesolution}
取地球半径 $R_\oplus=6.4\times10^6\,\mathrm m$，大气等效厚度
$d=5.0\times10^3\,\mathrm m$，外层半径 $R_\oplus+d$，电荷量大小
$Q=5.0\times10^5\,\mathrm C$。

\textbf{(1) 电容。} 同心球形电容公式给出
\[
 C=4\pi\varepsilon_0\frac{R_\oplus(R_\oplus+d)}d
 \simeq4\pi\varepsilon_0\frac{R_\oplus^2}{d}
 =\boxed{0.912\,\mathrm F}.
\]

\textbf{(2) 能量密度与总能量。} 地表场强
\[
 E(R_\oplus)=\frac{Q}{4\pi\varepsilon_0R_\oplus^2}
 \simeq1.10\times10^2\,\mathrm{V/m}.
\]
故地表能量密度
\[
 \boxed{u(R_\oplus)=\frac12\varepsilon_0E^2
 \simeq5.33\times10^{-8}\,\mathrm{J/m^3}}.
\]
总能量
\[
 \boxed{W=\frac{Q^2}{2C}\simeq1.37\times10^{11}\,\mathrm J}.
\]
相应电势差约为 $U=Q/C\simeq5.48\times10^5\,\mathrm V$。

\textbf{(3) 大气电阻。} 径向电流通过半径 $r$ 处面积 $4\pi r^2$，微元电阻
$\mathrm dR=\rho\,\mathrm dr/(4\pi r^2)$，因此
\[
 R_{\rm air}=\frac\rho{4\pi}
 \left(\frac1{R_\oplus}-\frac1{R_\oplus+d}\right)
 =\boxed{2.91\times10^2\,\Omega}.
\]

\textbf{(4) 时间常数。} 两个球壳使用完全相同的几何因子，故
\[
 \tau=R_{\rm air}C=\rho\varepsilon_0
 =\boxed{2.66\times10^2\,\mathrm s}\approx4.43\,\mathrm{min}.
\]

\textbf{(5) 雷电次数。} 为维持近似恒定的 $Q$，平均补充电流需抵消漏电流
$I=Q/\tau$。一天需要补充的电荷为 $Q(86400\,\mathrm s)/\tau$，所以
\[
 N=\frac{Q}{\tau}\frac{86400}{25}
 \simeq\boxed{6.5\times10^6\ \text{次/天}}.
\]
这个数值是题设均匀电阻率和固定等效高度模型的预测，不能单凭它认定真实全球雷电次数或验证大气电路。
真实电导率随高度变化，雷暴供电机制与单次输送电荷也有分布；本模型用于说明电容储能与漏电并存时的收支关系。
\end{xesolution}
}

\endgroup

\endgroup

\begingroup
\tcbset{breakable/.default=false,after={\par\nobreak\smallskip}}
\RenewDocumentEnvironment{xesolution}{+b}{%
  \par\nobreak\smallskip\nobreak\begingroup\small
  \clubpenalties 2 10000 10000\displaywidowpenalties 2 10000 10000\noindent
  {\bfseries\color{DeepBlue}解答}\quad #1\par\endgroup\smallskip
}{}
\begingroup
\renewcommand{\Needspace}[1]{\par\begingroup
\dimen0=#1\relax\vskip0pt plus\dimen0\penalty-100
\vskip0pt plus-\dimen0\vskip\dimen0\penalty9999\vskip-\dimen0
\vskip0pt\endgroup}
\tcbset{breakable/.default=false}
\chapter{恒定电流与复杂电路}

{
\graphicspath{{/Users/dongxingxi/Documents/物理竞赛一轮资料/sample/xes-original/thirty-sixth-section/}{/Users/dongxingxi/Documents/物理竞赛一轮资料/sample/xes-original/thirty-sixth-section/assets/}}
\section{电流密度与微观输运}

% All labels are transparent. Rectangles in circuit figures represent actual elements,
% not label backgrounds.
\tikzset{
  xeslabel/.style={font=\small,text=Ink,fill=none,inner sep=.55pt},
  xessmall/.style={font=\footnotesize,text=TextGray,fill=none,inner sep=.4pt},
  xestiny/.style={font=\scriptsize,text=TextGray,fill=none,inner sep=.3pt},
  xesline/.style={draw=Ink,line width=.85pt,line cap=round,line join=round},
  xesblue/.style={draw=DeepBlue,line width=1pt,line cap=round,line join=round},
  xesarrow/.style={-{Latex[length=2.1mm,width=1.3mm]},draw=DeepBlue,line width=.9pt},
  grayarrow/.style={-{Latex[length=1.7mm,width=1.1mm]},draw=TextGray,line width=.65pt},
  circuit/.style={draw=Ink,line width=.9pt,line cap=round,line join=round},
  element/.style={draw=Ink,fill=PaleBlue,minimum width=.66cm,minimum height=.30cm,
    inner sep=0pt,line width=.75pt},
  nodepoint/.style={circle,fill=DeepBlue,draw=none,inner sep=1.8pt},
  driftpath/.style={draw=DeepBlue!75,line width=.65pt}
}


\newcommand{\ivCharacteristicFigure}{%
\begin{center}
\begin{tikzpicture}[x=1cm,y=1cm]
  \begin{scope}[shift={(-3.2,0)}]
    \draw[circuit] (-2.15,0)--(-.75,0);
    \draw[circuit,fill=DeepBlue!7] (-.75,-.42) rectangle (.75,.42);
    \draw[circuit] (.75,0)--(2.15,0);
    \fill[white,draw=Ink,line width=.8pt] (-2.15,0) circle (2pt);
    \fill[white,draw=Ink,line width=.8pt] (2.15,0) circle (2pt);
    \draw[xesarrow] (-1.95,.62)--(-.95,.62);
    \node[xeslabel,above] at (-1.45,.65) {$I$};
    \node[xeslabel,below] at (0,-.48) {$V$};
    \node[xeslabel,below] at (-2.15,-.12) {$+$};
    \node[xeslabel,below] at (2.15,-.12) {$-$};
    \node[xessmall] at (0,-1.15) {二端元件与人为规定的测试方向};
  \end{scope}
  \begin{scope}[shift={(3.25,0)}]
    \draw[xesline,-{Latex[length=1.8mm]}] (-2.05,0)--(2.05,0) node[xeslabel,right] {$V$};
    \draw[xesline,-{Latex[length=1.8mm]}] (0,-1.55)--(0,1.65) node[xeslabel,above] {$I$};
    \draw[xesblue] (-1.65,-1.22)..controls(-1.20,-1.12)and(-.92,-.72)..(-.62,-.50)
      ..controls(-.20,-.18)and(.38,.24)..(.73,.53)
      ..controls(1.05,.79)and(1.28,1.13)..(1.43,1.48);
    \node[xessmall] at (0,-1.90) {伏安关系 $I=I(V)$};
  \end{scope}
\end{tikzpicture}
\end{center}}

\newcommand{\ohmHistoryFigure}{%
\begin{center}
  \includegraphics[width=.88\linewidth]{ohm-portrait-and-stamp-built-in-v1.png}
  \par\smallskip
  \begin{minipage}[t]{.43\linewidth}\centering\small\color{TextGray}
    (a) 格奥尔格·西蒙·欧姆（1789--1854）的历史肖像。
  \end{minipage}\hfill
  \begin{minipage}[t]{.43\linewidth}\centering\small\color{TextGray}
    (b) 以电阻和 $U=RI$ 为主题的纪念邮票。
  \end{minipage}
  \figcap{欧姆肖像与电阻定律纪念邮票。}\label{ch21:fig:s36-5-2}
\end{center}}

\newcommand{\cylindricalResistorFigure}{%
\begin{center}
\begin{tikzpicture}[x=1cm,y=1cm]
  \fill[DeepBlue!8] (-2.4,-.55) rectangle (2.4,.55);
  \draw[xesblue] (-2.4,-.55)--(2.4,-.55) (2.4,.55)--(-2.4,.55);
  \draw[xesblue] (-2.4,0) ellipse (.26 and .55);
  \draw[xesblue] (2.4,.55) arc[start angle=90,end angle=-90,x radius=.26,y radius=.55];
  \draw[xesblue,dashed] (2.4,-.55) arc[start angle=-90,end angle=-270,x radius=.26,y radius=.55];
  \draw[<->,Gold!70!black,line width=.65pt] (-2.35,.95)--(2.35,.95);
  \node[xestiny,text=Gold!70!black] at (0,1.18) {$l$};
  \draw[xesarrow] (-3.15,0)--(-2.48,0);
  \node[xeslabel,above] at (-2.82,.06) {$I$};
  \node[xeslabel] at (0,0) {$\rho$};
  \draw[RuleGray,line width=.55pt,-{Latex[length=1.5mm]}] (-3.45,-.85)--(-2.58,-.32);
  \node[xessmall,anchor=east] at (-3.50,-.88) {截面积 $A$};
  \node[xeslabel,below] at (-2.40,-.65) {$+$};
  \node[xeslabel,below] at (2.40,-.65) {$-$};
\end{tikzpicture}
\end{center}}

\newcommand{\ohmicIVFigure}{%
\begin{center}
\begin{tikzpicture}[x=1cm,y=1cm]
  \draw[xesline,-{Latex[length=1.8mm]}] (-2.15,0)--(2.15,0) node[xeslabel,right] {$V$};
  \draw[xesline,-{Latex[length=1.8mm]}] (0,-1.45)--(0,1.65) node[xeslabel,above] {$I$};
  \draw[xesblue] (-1.65,-1.25)--(1.72,1.30);
  \draw[RuleGray,densely dashed] (1.10,0)--(1.10,.83)--(0,.83);
  \node[xessmall,anchor=west] at (1.18,.58) {斜率 $1/R$};
  \node[xestiny,below left] at (0,0) {$O$};
\end{tikzpicture}
\end{center}}

\newcommand{\drudeFigure}{%
\begin{center}
\begin{tikzpicture}[x=1cm,y=1cm]
  \fill[DeepBlue!4] (-4.2,-1.45) rectangle (4.2,1.45);
  \foreach \x/\y in {-3.6/1.0,-2.4/.95,-1.15/1.05,.2/.95,1.45/1.08,2.75/.92,3.70/1.10,
                        -3.1/.05,-1.85/.15,-.55/-.02,.85/.12,2.15/.0,3.35/.15,
                        -3.65/-.95,-2.35/-1.02,-1.0/-.90,.35/-1.05,1.65/-.92,2.95/-1.06,3.82/-.92}{
    \fill[Gold!78] (\x,\y) circle (3.8pt);\draw[Gold!60!black,line width=.45pt] (\x,\y) circle (3.8pt);}
  \draw[driftpath] (-4.0,-.58)..controls(-3.55,-.20)and(-3.30,.45)..(-3.05,.10)
    ..controls(-2.70,-.30)and(-2.25,-.42)..(-2.02,.02)
    ..controls(-1.75,.50)and(-1.18,.35)..(-.96,-.05)
    ..controls(-.65,-.62)and(-.15,-.42)..(.10,.02)
    ..controls(.38,.55)and(.95,.46)..(1.18,.05)
    ..controls(1.45,-.42)and(2.0,-.52)..(2.28,-.04)
    ..controls(2.55,.42)and(3.12,.30)..(3.95,.58);
  \foreach \x/\y in {-3.80/-.48,-3.08/.10,-2.03/.02,-.96/-.05,.10/.02,1.18/.05,2.28/-.04,3.78/.49}{
    \fill[DeepBlue] (\x,\y) circle (2pt);}
  \draw[xesarrow] (-3.45,-1.82)--(-1.55,-1.82);
  \node[xessmall,below] at (-2.50,-1.88) {电子平均漂移 $\boldsymbol v_d$};
  \draw[xesarrow] (3.45,-1.82)--(1.55,-1.82);
  \node[xessmall,below] at (2.50,-1.88) {电场与电流 $\boldsymbol E,\boldsymbol j$};
  \node[xessmall] at (0,1.82) {高速无规则运动上叠加很小的定向漂移；每次碰撞近似“重置”漂移速度};
\end{tikzpicture}
\end{center}}

\newcommand{\gasDischargeFigure}{%
\begin{center}
\begin{tikzpicture}[x=.82cm,y=.78cm]
  \draw[xesline,-{Latex[length=1.8mm]}] (0,0)--(10.1,0) node[xessmall,right] {电流 $I$（对数标度）};
  \draw[xesline,-{Latex[length=1.8mm]}] (0,0)--(0,5.55) node[xessmall,above] {电压 $U$};
  \foreach \x/\lab in {1/{10^{-15}},3.7/{10^{-10}},6.2/{10^{-5}},8.7/{10^{0}}}{
    \draw[RuleGray] (\x,0)--(\x,-.10);\node[xestiny,below] at (\x,-.12) {$\lab$};}
  \foreach \y/\lab in {1/{200},2/{400},3/{600},4/{800},5/{1000}}{
    \draw[RuleGray] (0,\y)--(-.10,\y);\node[xestiny,anchor=east] at (-.14,\y) {$\lab$};}
  \draw[RuleGray,densely dashed] (5.72,0)--(5.72,5.15) (8.55,0)--(8.55,5.15);
  \node[xessmall] at (2.80,5.20) {暗放电};
  \node[xessmall] at (7.15,5.20) {辉光放电};
  \node[xessmall] at (9.25,5.20) {电弧};
  \draw[xesblue] (.55,.42)..controls(.95,.42)and(1.05,2.65)..(1.34,2.95)
    ..controls(1.58,3.12)and(4.8,3.20)..(5.70,3.26)
    ..controls(6.18,3.22)and(6.62,2.20)..(7.23,1.64)
    ..controls(7.55,1.34)and(7.96,1.50)..(8.18,2.20)
    ..controls(8.38,2.90)and(8.43,4.72)..(8.55,5.05);
  \draw[Gold!85!black,line width=1.05pt] (8.55,5.05)..controls(8.67,2.40)and(8.72,.72)..(9.05,.62)
    ..controls(9.36,.55)and(9.72,.66)..(9.92,.74);
  \foreach \x/\y/\lab in {.58/.45/{A'},3.10/3.16/{C},5.68/3.26/{D},6.15/2.92/{E},6.85/1.95/{F},7.45/1.46/{G},8.12/2.22/{H},8.55/5.05/{I},9.08/.62/{J},9.85/.74/{K}}{
    \fill[Ink] (\x,\y) circle (1.7pt);\node[xestiny,above=2pt] at (\x,\y) {$\lab$};}
  \fill[Ink] (1.35,2.95) circle (1.7pt);\node[xestiny,below left=1pt] at (1.35,2.95) {$B'$};
  \fill[Ink] (1.65,3.08) circle (1.7pt);\node[xestiny,above right=1pt] at (1.65,3.08) {$B$};
\end{tikzpicture}
\end{center}}

\newcommand{\seriesParallelResistorsFigure}{%
\begin{center}
\begin{tikzpicture}[x=1cm,y=1cm]
  \begin{scope}[shift={(-3.55,0)}]
    \draw[circuit] (-2.05,0)--(-1.35,0);
    \draw[circuit,decorate,decoration={zigzag,segment length=2.8mm,amplitude=1.05mm}] (-1.35,0)--(-.10,0);
    \draw[circuit] (-.10,0)--(.35,0);
    \draw[circuit,decorate,decoration={zigzag,segment length=2.8mm,amplitude=1.05mm}] (.35,0)--(1.60,0);
    \draw[circuit] (1.60,0)--(2.05,0);
    \node[xeslabel,above] at (-.72,.18) {$R_1$};\node[xeslabel,above] at (.98,.18) {$R_2$};
    \node[xessmall] at (0,-.72) {串联：电流相同};
  \end{scope}
  \begin{scope}[shift={(3.55,0)}]
    \draw[circuit] (-2.05,0)--(-1.55,0) (-1.55,0)--(-1.55,1.0) (-1.55,0)--(-1.55,-1.0);
    \draw[circuit,decorate,decoration={zigzag,segment length=2.8mm,amplitude=1.05mm}] (-1.55,1.0)--(1.55,1.0);
    \draw[circuit,decorate,decoration={zigzag,segment length=2.8mm,amplitude=1.05mm}] (-1.55,-1.0)--(1.55,-1.0);
    \draw[circuit] (1.55,1.0)--(1.55,0) (1.55,-1.0)--(1.55,0) (1.55,0)--(2.05,0);
    \node[xeslabel,above] at (0,1.12) {$R_1$};\node[xeslabel,below] at (0,-1.12) {$R_2$};
    \node[xessmall] at (0,-1.68) {并联：电压相同};
  \end{scope}
\end{tikzpicture}
\end{center}}

\newcommand{\networkTermsFigure}{%
\begin{center}
\begin{tikzpicture}[x=1cm,y=1cm]
  \coordinate (A) at (0,2.0);\coordinate (B) at (-2.0,.62);\coordinate (C) at (-1.25,-1.72);
  \coordinate (D) at (1.25,-1.72);\coordinate (E) at (2.0,.62);\coordinate (O) at (0,0);
  \draw[circuit] (A)--(B)--(C)--(D)--(E)--cycle;
  \draw[circuit] (A)--(O)--(B) (O)--(C) (O)--(D) (O)--(E);
  \foreach \P in {A,B,C,D,E,O}{\node[nodepoint] at (\P) {};}
  \foreach \X/\Y in {A/B,B/C,C/D,D/E,E/A,A/O,O/B,O/C,O/D,O/E}{
    \path (\X)--(\Y) node[element,pos=.48,sloped] {};}
  \draw[Gold!80!black,densely dashed,line width=.7pt,rounded corners=7pt]
    (-.15,1.48)--(-.65,.58)--(0,0)--(1.05,.35)--(.85,1.20)--cycle;
  \node[xessmall,text=Gold!70!black] at (.60,.68) {回路};
  \draw[grayarrow] (-3.45,1.65)--(-2.05,.68);\node[xeslabel,anchor=east] at (-3.52,1.65) {节点};
  \draw[grayarrow] (-3.45,-.25)--($(B)!.60!(C)$);\node[xeslabel,anchor=east] at (-3.52,-.25) {支路};
  \draw[grayarrow] (3.55,1.35)--($(E)!.48!(A)$);\node[xeslabel,anchor=west] at (3.62,1.35) {元件};
  \node[xessmall] at (0,-2.32) {节点由导线连续区域确定；两个节点之间是一条支路};
\end{tikzpicture}
\end{center}}

\newcommand{\kclKvlFigure}{%
\begin{center}
\begin{tikzpicture}[x=.95cm,y=.95cm]
  \begin{scope}[shift={(-3.9,0)}]
    \coordinate (N1) at (0,1.85);\coordinate (N2) at (-1.70,-1.15);\coordinate (N4) at (1.70,-1.15);\coordinate (N3) at (0,-.05);
    \draw[circuit] (N1)--(N2)--(N4)--cycle (N1)--(N3)--(N2) (N3)--(N4);
    \foreach \P in {N1,N2,N3,N4}{\node[nodepoint] at (\P) {};}
    \foreach \X/\Y in {N1/N2,N1/N4,N2/N4,N1/N3,N3/N2,N3/N4}{\path (\X)--(\Y) node[element,pos=.48,sloped] {};}
    \draw[xesarrow] ($(N1)!.09!(N2)$)--($(N1)!.29!(N2)$);
    \node[xestiny,anchor=east] at (-.47,1.18) {$i_1$};
    \draw[xesarrow] ($(N1)!.08!(N3)$)--($(N1)!.28!(N3)$);
    \node[xestiny,anchor=west] at (.12,1.50) {$i_2$};
    \draw[xesarrow] ($(N1)!.09!(N4)$)--($(N1)!.29!(N4)$);
    \node[xestiny,anchor=west] at (.47,1.18) {$i_3$};
    \draw[xesarrow] ($(N3)!.08!(N2)$)--($(N3)!.29!(N2)$);
    \node[xestiny,anchor=east] at (-.30,-.22) {$i_4$};
    \draw[xesarrow] ($(N3)!.08!(N4)$)--($(N3)!.29!(N4)$);
    \node[xestiny,anchor=west] at (.30,-.22) {$i_5$};
    \draw[xesarrow] ($(N2)!.07!(N4)$)--($(N2)!.29!(N4)$);
    \node[xestiny,below] at (-1.05,-1.20) {$i_6$};
    \node[xeslabel,above] at (N1) {$N_1$};\node[xeslabel,below left] at (N2) {$N_2$};
    \node[xeslabel,above=2pt] at (N3) {$N_3$};\node[xeslabel,below right] at (N4) {$N_4$};
    \node[xessmall] at (0,-2.05) {KCL：先规定支路电流方向};
  \end{scope}
  \begin{scope}[shift={(3.9,0)}]
    \coordinate (N1) at (0,1.85);\coordinate (N2) at (-1.70,-1.15);\coordinate (N4) at (1.70,-1.15);\coordinate (N3) at (0,-.05);
    \draw[circuit] (N1)--(N2)--(N4)--cycle (N1)--(N3)--(N2) (N3)--(N4);
    \foreach \P in {N1,N2,N3,N4}{\node[nodepoint] at (\P) {};}
    \foreach \X/\Y in {N1/N2,N1/N4,N2/N4,N1/N3,N3/N2,N3/N4}{\path (\X)--(\Y) node[element,pos=.48,sloped] {};}
    \node[xestiny,anchor=east] at (-1.13,.63) {$v_1$};
    \node[xestiny,anchor=west] at (1.13,.63) {$v_3$};
    \node[xestiny,anchor=west] at (.18,1.03) {$v_2$};
    \node[xestiny,anchor=east] at (-.72,-.30) {$v_4$};
    \node[xestiny,anchor=west] at (.72,-.30) {$v_5$};
    \node[xestiny,below] at (0,-1.46) {$v_6$};
    \node[xestiny] at (-.39,.88) {$+$};\node[xestiny] at (-.94,-.08) {$-$};
    \node[xestiny] at (-.18,1.25) {$+$};\node[xestiny] at (-.18,.61) {$-$};
    \node[xestiny] at (.39,.88) {$+$};\node[xestiny] at (.94,-.08) {$-$};
    \node[xestiny] at (-.39,-.44) {$+$};\node[xestiny] at (-.94,-.80) {$-$};
    \node[xestiny] at (.39,-.44) {$+$};\node[xestiny] at (.94,-.80) {$-$};
    \node[xestiny] at (-.61,-.98) {$+$};\node[xestiny] at (.61,-.98) {$-$};
    \draw[Gold!80!black,-{Latex[length=1.5mm]},line width=.65pt]
      (-.66,.58) arc[start angle=120,end angle=430,radius=.20];
    \node[xestiny,text=Gold!70!black] at (-.58,.20) {$L_1$};
    \draw[Gold!80!black,-{Latex[length=1.5mm]},line width=.65pt]
      (.46,.58) arc[start angle=60,end angle=-250,radius=.20];
    \node[xestiny,text=Gold!70!black] at (.58,.20) {$L_2$};
    \draw[Gold!80!black,-{Latex[length=1.5mm]},line width=.65pt]
      (-.20,-.70) arc[start angle=170,end angle=-140,radius=.20];
    \node[xestiny,text=Gold!70!black] at (0,-.78) {$L_3$};
    \draw[Gold!80!black,-{Latex[length=1.7mm]},line width=.75pt]
      (0,2.28) arc[start angle=90,end angle=430,x radius=2.25,y radius=2.02];
    \node[xestiny,text=Gold!70!black] at (1.92,1.52) {$L_4$};
    \node[xessmall] at (0,-2.05) {KVL：沿所选回路累计支路电压};
  \end{scope}
\end{tikzpicture}
\end{center}}

\newcommand{\nodeMethodFigure}{%
\begin{center}
\begin{tikzpicture}[x=1cm,y=1cm]
  \coordinate (a) at (-2.4,1.35);\coordinate (b) at (1.25,1.35);\coordinate (c) at (1.25,-1.35);
  \draw[circuit] (a)--(-1.25,1.35);
  \draw[circuit,decorate,decoration={zigzag,segment length=3mm,amplitude=1.05mm}] (-1.25,1.35)--(.15,1.35);
  \draw[circuit] (.15,1.35)--(b);
  \draw[circuit] (b)--(1.25,.40);
  \draw[circuit,decorate,decoration={zigzag,segment length=3mm,amplitude=1.05mm}] (1.25,.40)--(1.25,-.78);
  \draw[circuit] (1.25,-.78)--(c)--(-2.4,-1.35);
  \draw[circuit] (-2.4,-1.35)--(-2.4,-.34);
  \draw[circuit] (-2.63,-.34)--(-2.17,-.34) (-2.78,.02)--(-2.02,.02);
  \draw[circuit] (-2.4,.02)--(a);
  \draw[circuit] (b)--(3.0,1.35)--(3.0,.34);
  \draw[circuit] (3.0,-.34)--(3.0,-1.35)--(c);
  \draw[circuit] (3.0,0) circle (.34);
  \draw[xesarrow] (3.0,-.20)--(3.0,.20);
  \foreach \P in {a,b,c}{\node[nodepoint] at (\P) {};}
  \node[xeslabel,above] at (a) {$a$};\node[xeslabel,above] at (b) {$b$};\node[xeslabel,below] at (c) {$c$};
  \node[xeslabel,above] at (-.55,1.48) {$R_1$};\node[xeslabel,right] at (1.42,-.18) {$R_2$};
  \node[xeslabel,left] at (-2.82,-.15) {$E_0$};\node[xeslabel,right] at (3.40,0) {$I_0$};
  \draw[circuit] (c)--(1.25,-1.65) (.92,-1.65)--(1.58,-1.65) (1.02,-1.82)--(1.48,-1.82) (1.12,-1.98)--(1.38,-1.98);
  \node[xessmall] at (0,-2.42) {取 $c$ 为参考地后，只需求节点 $b$ 的电势};
\end{tikzpicture}
\end{center}}

\newcommand{\tolmanRingFigure}{%
\begin{center}
\begin{tikzpicture}[x=1cm,y=1cm]
  \draw[xesblue,line width=1.2pt] (0,0) ellipse (2.55 and 1.12);
  \draw[xesblue,line width=.75pt] (0,0) ellipse (2.30 and .92);
  \draw[xesline] (0,-1.55)--(0,1.55);
  \draw[xesarrow] (-.55,1.45) arc[start angle=160,end angle=25,x radius=.65,y radius=.30];
  \node[xeslabel,above] at (0,1.52) {$\omega$};
  \draw[xesarrow] (2.08,.52)..controls(1.55,.85)and(.95,1.02)..(.35,1.08);
  \node[xeslabel,above] at (1.25,.98) {$I(t)$};
  \draw[RuleGray,dashed] (0,0)--(2.40,.37);\node[xeslabel,above] at (1.42,.27) {$r$};
  \node[xessmall] at (0,-1.98) {晶格骤停后，电子因惯性暂时保持切向速度};
\end{tikzpicture}
\end{center}}

\newcommand{\transientCurrentFigure}{%
\begin{center}
\begin{tikzpicture}[x=1cm,y=1cm]
  \draw[xesline,-{Latex[length=1.8mm]}] (0,0)--(5.2,0) node[xeslabel,right] {$t$};
  \draw[xesline,-{Latex[length=1.8mm]}] (0,0)--(0,2.55) node[xeslabel,above] {$I$};
  \draw[RuleGray,densely dashed] (0,2.0)--(4.75,2.0);
  \node[xessmall,anchor=east] at (-.12,2.0) {$U/R$};
  \draw[xesblue] plot[domain=0:4.7,samples=80] (\x,{2*(1-exp(-\x/1.0))});
  \draw[<->,Gold!75!black,line width=.65pt] (0,-.42)--(1.0,-.42);
  \node[xestiny,text=Gold!70!black,below] at (.5,-.46) {$\tau$};
\end{tikzpicture}
\end{center}}

\newcommand{\rotatingChargedRingFigure}{%
\begin{center}
\begin{tikzpicture}[x=1cm,y=1cm]
  \draw[xesblue,line width=1.2pt] (0,0) ellipse (2.55 and 1.03);
  \draw[xesblue,line width=.72pt] (0,0) ellipse (2.30 and .85);
  \draw[xesline] (0,-1.50)--(0,1.50);
  \draw[xesarrow] (-.55,1.40) arc[start angle=160,end angle=25,x radius=.65,y radius=.30];
  \node[xeslabel,above] at (0,1.48) {$\omega$};
  \draw[xesarrow] (.48,1.02)..controls(1.1,.96)and(1.65,.83)..(2.15,.50);
  \node[xeslabel,above] at (1.37,.95) {$I$};
  \draw[RuleGray,dashed] (0,0)--(2.38,.35);\node[xeslabel,above] at (1.42,.24) {$R$};
  \draw[grayarrow] (-3.05,.62)--(-2.34,.51);\node[xeslabel,anchor=east] at (-3.12,.62) {$Q$};
  \node[xeslabel,below left] at (0,0) {$O$};
\end{tikzpicture}
\end{center}}

\newcommand{\vacuumDiodeFigure}{%
\begin{center}
\begin{tikzpicture}[x=1cm,y=1cm]
  \draw[xesblue,line width=1.25pt] (-2.35,-.75)--(-2.35,.75) (2.35,-.75)--(2.35,.75);
  \draw[<->,Gold!75!black,line width=.65pt] (-2.30,1.52)--(2.30,1.52);
  \node[xestiny,text=Gold!70!black] at (0,1.76) {$L$};
  \node[xeslabel,above] at (-2.35,.82) {阴极 $C$};\node[xeslabel,above] at (2.35,.82) {阳极 $A$};
  \draw[xesarrow] (-1.75,0)--(1.75,0);\node[xessmall,above] at (0,.05) {电子运动};
  \draw[grayarrow] (2.95,.45)--(2.40,.45);\node[xeslabel,anchor=west] at (3.02,.45) {$I$};
  \draw[xesline,-{Latex[length=1.6mm]}] (-2.60,-1.05)--(2.80,-1.05) node[xeslabel,right] {$x$};
  \node[xestiny,below] at (-2.35,-1.06) {$0$};\node[xestiny,below] at (2.35,-1.06) {$L$};
\end{tikzpicture}
\end{center}}

\newcommand{\childLangmuirFigure}{%
\begin{center}
\begin{tikzpicture}[x=1cm,y=1cm]
  \draw[xesline,-{Latex[length=1.7mm]}] (0,0)--(5.15,0) node[xeslabel,right] {$x$};
  \draw[xesline,-{Latex[length=1.7mm]}] (0,0)--(0,2.65) node[xeslabel,above] {$\varphi$};
  \draw[xesblue] plot[domain=0:4.7,samples=80] (\x,{2.25*(\x/4.7)^(4/3)});
  \draw[RuleGray,densely dashed] (4.7,0)--(4.7,2.25)--(0,2.25);
  \node[xestiny,below] at (4.7,0) {$L$};\node[xestiny,anchor=east] at (-.10,2.25) {$U$};
  \node[xessmall,anchor=west] at (1.00,1.82) {$\varphi(x)=U(x/L)^{4/3}$};
\end{tikzpicture}
\end{center}}

\newcommand{\conductingShellFigure}{%
\begin{center}
\begin{tikzpicture}[x=1cm,y=1cm]
  \fill[DeepBlue!8] (0,0) circle (1.85);\fill[white] (0,0) circle (.82);
  \draw[xesblue] (0,0) circle (1.85) (0,0) circle (.82);
  \draw[xesline] (0,0)--(.58,-.58);\node[xeslabel,below right] at (.62,-.62) {$a$};
  \draw[xesline] (0,0)--(0,1.85);\node[xeslabel,left] at (0,1.25) {$b$};
  \draw[xesarrow] (1.05,.25)--(1.62,.39);\node[xessmall,anchor=west] at (1.72,.44) {$\boldsymbol j\parallel\boldsymbol E$};
  \node[xessmall] at (0,-2.30) {$\sigma(r)=K E(r)$，电流在任意同心球面上相同};
\end{tikzpicture}
\end{center}}

\newcommand{\ionChamberFigure}{%
\begin{center}
\begin{tikzpicture}[x=1cm,y=1cm]
  \draw[xesblue,line width=1.2pt] (-3.5,1.45)--(3.5,1.45) (-3.5,-1.45)--(3.5,-1.45);
  \draw[xesline] (-3.5,1.45)--(-3.5,-1.45) (3.5,1.45)--(3.5,-1.45);
  \node[xeslabel,above] at (-1.90,1.48) {阴极};\node[xeslabel,below] at (-1.90,-1.48) {阳极};
  \draw[circuit] (1.90,1.45)--(1.90,2.05)--(3.15,2.05);\node[xeslabel,right] at (3.15,2.05) {$-U_0$};
  \draw[circuit] (1.90,-1.45)--(1.90,-2.05);
  \draw[circuit,decorate,decoration={zigzag,segment length=2.7mm,amplitude=1.0mm}] (1.90,-2.05)--(1.90,-2.95);
  \draw[circuit] (1.90,-2.95)--(1.90,-3.20);
  \draw[circuit] (1.55,-3.20)--(2.25,-3.20) (1.66,-3.37)--(2.14,-3.37) (1.77,-3.54)--(2.03,-3.54);
  \node[xeslabel,right] at (2.10,-2.50) {$R$};\draw[circuit] (1.90,-1.70)--(3.20,-1.70);
  \node[xeslabel,right] at (3.20,-1.70) {信号};
  \draw[xesline,-{Latex[length=1.7mm]}] (-4.65,-1.45)--(-4.65,2.05) node[xeslabel,above] {$x$};
  \draw[<->,Gold!75!black,line width=.65pt] (-3.78,-1.42)--(-3.78,1.42);
  \node[xestiny,text=Gold!70!black,anchor=east] at (-3.92,0) {$D$};
  \draw[<->,Gold!75!black,line width=.65pt] (-3.25,-1.40)--(-3.25,-.15);
  \node[xestiny,text=Gold!70!black,anchor=east] at (-3.36,-.80) {$x_0$};
  \draw[xesarrow] (-2.85,-.14)--(2.85,-.14);\node[xessmall,above] at (1.5,.75) {入射粒子径迹};
  \foreach \x in {-2.3,-1.65,-.95,-.25,.45,1.15,1.85,2.50}{
    \node[xestiny,text=Gold!70!black] at (\x,.02) {$+$};\node[xestiny,text=DeepBlue] at (\x,-.29) {$-$};}
  \draw[grayarrow] (-.42,.08)--(-.42,.78);\node[xestiny,anchor=west] at (-.34,.63) {离子 $v^+$};
  \draw[grayarrow] (.42,-.32)--(.42,-.98);\node[xestiny,anchor=west] at (.50,-.88) {电子 $v^-$};
  \node[xessmall] at (0,2.58) {气密容器内充工作气体；电场把电子与正离子分别引向两极};
\end{tikzpicture}
\end{center}}

\newcommand{\ionSignalFigure}{%
\begin{center}
\begin{tikzpicture}[x=1cm,y=1cm]
  \draw[xesline,-{Latex[length=1.7mm]}] (0,0)--(6.1,0) node[xeslabel,right] {$t$};
  \draw[xesline,-{Latex[length=1.7mm]}] (0,0)--(0,2.65) node[xeslabel,above] {$|\Delta U|$};
  \draw[xesblue] (0,0)--(1.25,1.25)--(4.70,2.25);
  \draw[RuleGray,densely dashed] (1.25,0)--(1.25,1.25) (4.70,0)--(4.70,2.25);
  \node[xestiny,below] at (1.25,0) {$t_e=x_0/v^-$};
  \node[xestiny,below] at (4.70,0) {$t_i=(D-x_0)/v^+$};
  \node[xessmall] at (3.10,2.55) {电子收集完后，离子继续形成较慢斜坡};
\end{tikzpicture}
\end{center}}

\newcommand{\ringResistorFigure}{%
\begin{center}
\begin{tikzpicture}[x=1cm,y=1cm]
  \draw[xesblue,line width=1.1pt] (0,0) circle (1.55);
  \coordinate (a) at (-1.55,0);\coordinate (b) at ({1.55*cos(-35)},{1.55*sin(-35)});
  \node[nodepoint] at (a) {};\node[nodepoint] at (b) {};
  \draw[circuit] (a)--(-2.65,0);\draw[circuit] (b)--($(b)+(1.0,-.70)$);
  \draw[RuleGray,dashed] (0,0)--(a) (0,0)--(b);
  \node[xeslabel,above left] at (a) {$a$};\node[xeslabel,below right] at (b) {$b$};
  \node[xeslabel,below left] at (0,0) {$O$};\node[xeslabel] at (-.45,.90) {$R_1$};\node[xeslabel] at (.72,-.80) {$R_2$};
  \node[xessmall] at (0,-2.10) {$R_1+R_2=R$，两段圆弧在端点 $a,b$ 之间并联};
\end{tikzpicture}
\end{center}}

\newcommand{\dualSourceFigure}{%
\begin{center}
\begin{tikzpicture}[x=1cm,y=1cm]
  \coordinate (L) at (-2.65,0);\coordinate (R) at (2.65,0);
  \draw[circuit] (L)--(-1.50,0);
  \draw[circuit,decorate,decoration={zigzag,segment length=3mm,amplitude=1.05mm}] (-1.50,0)--(1.20,0);
  \draw[circuit] (1.20,0)--(R);
  \node[xeslabel,above] at (0,0.25) {$R$};
  \draw[circuit] (L)--(-2.65,-1.20)--(-.55,-1.20);
  \draw[circuit] (-.55,-1.60)--(-.55,-.80) (-.15,-1.48)--(-.15,-.92);
  \draw[circuit] (-.15,-1.20)--(.35,-1.20);
  \draw[circuit,decorate,decoration={zigzag,segment length=2.8mm,amplitude=.9mm}] (.35,-1.20)--(1.70,-1.20);
  \draw[circuit] (1.70,-1.20)--(2.65,-1.20)--(R);
  \node[xeslabel,below] at (-.35,-1.62) {$\varepsilon_1$};\node[xeslabel,above] at (1.02,-1.03) {$r_1$};
  \draw[circuit] (L)--(-2.65,-2.40)--(-.55,-2.40);
  \draw[circuit] (-.55,-2.80)--(-.55,-2.00) (-.15,-2.68)--(-.15,-2.12);
  \draw[circuit] (-.15,-2.40)--(.35,-2.40);
  \draw[circuit,decorate,decoration={zigzag,segment length=2.8mm,amplitude=.9mm}] (.35,-2.40)--(1.70,-2.40);
  \draw[circuit] (1.70,-2.40)--(2.65,-2.40)--(R);
  \node[xeslabel,below] at (-.35,-2.82) {$\varepsilon_2$};\node[xeslabel,above] at (1.02,-2.23) {$r_2$};
  \node[nodepoint] at (L) {};\node[nodepoint] at (R) {};
  \draw[xesarrow] (-1.9,.45)--(-.85,.45);\node[xeslabel,above] at (-1.38,.48) {$I$};
  \node[xessmall] at (0,-3.30) {两个实际电源同极性并联后向负载 $R$ 供电};
\end{tikzpicture}
\end{center}}

% Source question conditions restored, analysis retained separately.
\newcommand{\emRotatingRingQuestion}{%
\begin{center}
\begin{tikzpicture}[x=1cm,y=1cm]
  \draw[xesblue,line width=1.2pt] (0,0) ellipse (2.55 and 1.03);
  \draw[xesblue,line width=.72pt] (0,0) ellipse (2.30 and .85);
  \draw[xesline] (0,-1.50)--(0,1.50);
  \draw[xesarrow] (-.55,1.40) arc[start angle=160,end angle=25,x radius=.65,y radius=.30];
  \node[xeslabel,above] at (0,1.48) {$\omega$};
  \draw[RuleGray,dashed] (0,0)--(2.38,.35);\node[xeslabel,above] at (1.42,.24) {$R$};
  \draw[grayarrow] (-3.05,.62)--(-2.34,.51);\node[xeslabel,anchor=east] at (-3.12,.62) {$Q$};
  \node[xeslabel,below left] at (0,0) {$O$};
\end{tikzpicture}
\end{center}}

% Source question conditions restored, analysis retained separately.
\newcommand{\emConductingShellQuestion}{%
\begin{center}
\begin{tikzpicture}[x=1cm,y=1cm]
  \fill[DeepBlue!8] (0,0) circle (1.85);\fill[white] (0,0) circle (.82);
  \draw[xesblue] (0,0) circle (1.85) (0,0) circle (.82);
  \draw[xesline] (0,0)--(.58,-.58);\node[xeslabel,below right] at (.62,-.62) {$a$};
  \draw[xesline] (0,0)--(0,1.85);\node[xeslabel,left] at (0,1.25) {$b$};
\end{tikzpicture}
\end{center}}

% Source question conditions restored, analysis retained separately.
\newcommand{\emIonChamberQuestion}{%
\begin{center}
\begin{tikzpicture}[x=1cm,y=1cm]
  \draw[xesblue,line width=1.2pt] (-3.5,1.45)--(3.5,1.45) (-3.5,-1.45)--(3.5,-1.45);
  \draw[xesline] (-3.5,1.45)--(-3.5,-1.45) (3.5,1.45)--(3.5,-1.45);
  \node[xeslabel,above] at (-1.90,1.48) {阴极};\node[xeslabel,below] at (-1.90,-1.48) {阳极};
  \draw[circuit] (1.90,1.45)--(1.90,2.05)--(3.15,2.05);\node[xeslabel,right] at (3.15,2.05) {$-U_0$};
  \draw[circuit] (1.90,-1.45)--(1.90,-2.05);
  \draw[circuit,decorate,decoration={zigzag,segment length=2.7mm,amplitude=1.0mm}] (1.90,-2.05)--(1.90,-2.95);
  \draw[circuit] (1.90,-2.95)--(1.90,-3.20);
  \draw[circuit] (1.55,-3.20)--(2.25,-3.20) (1.66,-3.37)--(2.14,-3.37) (1.77,-3.54)--(2.03,-3.54);
  \node[xeslabel,right] at (2.10,-2.50) {$R$};\draw[circuit] (1.90,-1.70)--(3.20,-1.70);
  \node[xeslabel,right] at (3.20,-1.70) {信号};
  \draw[xesline,-{Latex[length=1.7mm]}] (-4.65,-1.45)--(-4.65,2.05) node[xeslabel,above] {$x$};
  \draw[<->,Gold!75!black,line width=.65pt] (-3.78,-1.42)--(-3.78,1.42);
  \node[xestiny,text=Gold!70!black,anchor=east] at (-3.92,0) {$D$};
  \draw[<->,Gold!75!black,line width=.65pt] (-3.25,-1.40)--(-3.25,-.15);
  \node[xestiny,text=Gold!70!black,anchor=east] at (-3.36,-.80) {$x_0$};
  \draw[xesarrow] (-2.85,-.14)--(2.85,-.14);\node[xessmall,above] at (1.5,.75) {入射粒子径迹};
  \foreach \x in {-2.3,-1.65,-.95,-.25,.45,1.15,1.85,2.50}{
    \node[xestiny,text=Gold!70!black] at (\x,.02) {$+$};\node[xestiny,text=DeepBlue] at (\x,-.29) {$-$};}
\end{tikzpicture}
\end{center}}

% Source question conditions restored, analysis retained separately.
\newcommand{\emRingResistorQuestion}{%
\begin{center}
\begin{tikzpicture}[x=1cm,y=1cm]
  \draw[xesblue,line width=1.1pt] (0,0) circle (1.55);
  \coordinate (a) at (-1.55,0);\coordinate (b) at ({1.55*cos(-35)},{1.55*sin(-35)});
  \node[nodepoint] at (a) {};\node[nodepoint] at (b) {};
  \draw[circuit] (a)--(-2.65,0);\draw[circuit] (b)--($(b)+(1.0,-.70)$);
  \draw[RuleGray,dashed] (0,0)--(a) (0,0)--(b);
  \node[xeslabel,above left] at (a) {$a$};\node[xeslabel,below right] at (b) {$b$};
  \node[xeslabel,below left] at (0,0) {$O$};\node[xeslabel] at (-.45,.90) {$R_1$};\node[xeslabel] at (.72,-.80) {$R_2$};
\end{tikzpicture}
\end{center}}


在最一般的意义下，电流是大量带电粒子沿确定方向发生整体定向运动的宏观表现。电能的输运和使用
已经成为现代社会的基础；如果所有电力驱动的设备突然停止工作，生活与生产会立即受到全面影响。
本章从电流的定义出发，讨论物质导电的经典图像、元件的伏安特性、欧姆定律（Ohm's law）以及一般直流网络的
基尔霍夫定律（Kirchhoff's circuit laws）。本章后续内容还将在这些基础上研究更复杂的电阻网络。

导体和绝缘体的关键差别，在于物质内部是否存在能够在外场作用下长距离迁移的带电粒子，即载流子。
金属中的载流子主要是电子；电解液中正、负离子都可参与导电；半导体中则可能同时存在电子和空穴。
载流子的浓度、质量、电荷量和碰撞性质，以及温度、外电场、光照等外部条件，共同决定材料的导电性。
必须强调：金属导电的精确理论属于量子物理。下面的经典模型不能从第一性原理给出全部材料参数，
但能把宏观规律与载流子的运动联系起来，并说明经典图像在哪些条件下有效、又会在何处失效。


\subsection{电流强度和电流密度}


电流并不只存在于接有电源的闭合回路中；只要带电粒子发生有方向的输运，就形成电流。取一个有向截面，
若在时间间隔 $\Delta t$ 内按规定正方向通过它的净电量为 $\Delta Q$，平均电流强度定义为
\begin{equation}
  \boxed{I=\frac{\Delta Q}{\Delta t}}.
\end{equation}
当时间间隔趋于零，瞬时电流是通过电量 $Q(t)$ 对时间的导数：
\begin{equation}
  \boxed{I(t)=\lim_{\Delta t\to0}\frac{\Delta Q}{\Delta t}
  =\frac{\mathrm dQ}{\mathrm dt}}.
\end{equation}
一般情况下 $I$ 随时间变化；若它不随时间变化，称为稳恒电流。反过来，已知电流随时间的关系后，
从 $t_1$ 到 $t_2$ 通过截面的净电量为
\begin{equation}
  \boxed{Q(t_1\to t_2)=\int_{t_1}^{t_2}I(t)\,\mathrm dt}.
\end{equation}
对于恒定电流，上式退化为 $Q=I(t_2-t_1)$。电流的国际单位是安培，
$1\,\mathrm A=1\,\mathrm C\,\mathrm s^{-1}$。

设导线横截面积为 $A$，单位体积内有 $n$ 个同种载流子，每个载流子的电量为 $q$，沿导线的平均定向
漂移速度为 $v_d$。在 $\Delta t$ 内能够穿过截面的载流子原来位于长为 $v_d\Delta t$ 的柱体中，
其数目为 $nAv_d\Delta t$，因而通过的电量是 $nqAv_d\Delta t$。代入定义得到
\begin{equation}
  \boxed{I=nqAv_d}.
\end{equation}
式中 $q$ 是带符号的电荷量，$v_d$ 是沿所选截面法向的速度分量。对电子，电荷与沿常规电流方向的
速度分量均为负，乘积为正。若只讨论大小，则写 $|I|=n|q|A|v_d|$，不可只对电荷取绝对值。

真实导线中的漂移速度通常很小。以横截面积 $A=1\,\mathrm{mm^2}$ 的纯银导线为例，若每个银原子提供
一个导电电子，则 $n\simeq6\times10^{28}\,\mathrm{m^{-3}}$。通有 $1\,\mathrm A$ 电流时
\begin{equation}
  \boxed{v_d=\frac{I}{neA}\approx1.0\times10^{-4}\,\mathrm{m\,s^{-1}}}.
\end{equation}
这个速度比金属中电子本身的无规则运动速度小许多个数量级。因此“电子像小球一样沿导线迅速跑到用电器”
并不是正确图像：电路中建立电场的过程很快，但单个电子的净漂移极慢。
金属电子的无规则运动速率不能用经典能量均分定理准确求得，必须由量子统计决定；其典型高速电子可达到
约光速的百分之一，而上面算出的净漂移仍只有 $10^{-4}\,\mathrm{m\,s^{-1}}$ 量级。

若截面较大，不同位置的载流子漂移可能并不相同，单一的 $I$ 不能描述局部流动。定义电流密度矢量
\begin{equation}
  \boxed{\boldsymbol j=nq\boldsymbol v_d}.
\end{equation}
它的方向按正电荷运动方向规定。对于电子，$q=-e$，所以 $\boldsymbol j$ 与电子漂移速度方向相反；
若有多种载流子同时运动，总电流密度是各类载流子贡献之和。穿过任意有向曲面 $S$ 的电流为
\[
  I=\iint_S\boldsymbol j\cdot\mathrm d\boldsymbol S.
\]
这正像流体的体积流量等于速度场在截面上的通量，只不过电流计算的是电荷通量。
例如河水中央常比近岸处流得快，必须先描述各点流速才能求总流量；非均匀导体中的电流也完全类似。


\subsection{导电的微观机制——简单模型}


先构造一个纯粹唯象的“黏滞阻力”模型。设载流子定向运动时受到与漂移速度成正比、方向相反的平均阻力
$\boldsymbol F_r=-\alpha\boldsymbol v_d$。恒定电场中达到稳定漂移后，电场力与阻力平衡：
\[
  q\boldsymbol E-\alpha\boldsymbol v_d=0,
  \qquad \boldsymbol v_d=\frac q\alpha\boldsymbol E.
\]
于是
\[
  \boldsymbol j=nq\boldsymbol v_d=\frac{nq^2}{\alpha}\boldsymbol E.
\]
与微观欧姆定律比较可得
\begin{equation}
  \boxed{\sigma=\frac{nq^2}{\alpha}}.
\end{equation}
这个模型把极复杂的散射过程压缩成一个平均阻力系数。它虽然不是真实金属中电子运动的精确图像，却能
说明为什么稳定状态下 $\boldsymbol j$ 与 $\boldsymbol E$ 同向且成正比。

它还提示了暂态过程：电压突然改变时，载流子具有惯性，漂移速度不能无延迟地跳变。若保留质量 $m$，
运动方程为 $m\dot{\boldsymbol v}_d=q\boldsymbol E-\alpha\boldsymbol v_d$，电流会在特征时间
$\tau=m/\alpha$ 内趋近新的稳态值。通常金属中的 $\tau$ 极短，所以普通低频电路把这一过程近似为
瞬时完成；例~\ref{ch21:ex:s36-5-2} 将把它定量写出。


% EM teaching enrichment 20260927 supplementary

\paragraph{从建立过程看稳定电流}
这个模型也给出电场接通后的响应。初始漂移速度为零时，
\[
 m\frac{\dd\vect v_d}{\dd t}=q\vect E-\alpha\vect v_d,
 \qquad \vect v_d(t)=\frac{q\vect E}{\alpha}(1-\mathrm e^{-t/\tau}),\quad \tau=\frac m\alpha.
\]
乘以 $nq$ 后得到 $\vect j(t)=\sigma\vect E(1-\mathrm e^{-t/\tau})$，其中 $\sigma=nq^2/\alpha$。
达到 $t=\tau$ 时电流约为最终值的 $63\%$，约数个 $\tau$ 后才可用稳态欧姆定律。
负载流子的漂移方向与电场相反，但 $q\vect v_d$ 与电场同向，故电导率仍为正。
这一时间尺度描述载流子的局部松弛，不能把它当作信号沿整根导线传播的时间。


\subsection{电子的扩散和漂移}


金属中的电子并非只做缓慢定向运动。它们本来就在晶体中高速无规则运动，外电场只在这种运动上叠加
一个很小的平均漂移。电子在两次散射之间近似自由运动，受电场加速；散射会把先前积累的定向速度随机化，
大量电子的统计平均才形成稳定电流。

\par\noindent\begin{minipage}{\linewidth}
\drudeFigure
\diagramcaption{导体内电子的无规则运动与定向漂移。蓝色曲线表示一条示意轨迹，金色圆点表示晶格散射中心。}\label{ch21:fig:s36-5-6}
\end{minipage}\par\medskip

设自由电子数密度为 $n$、电子电荷量的绝对值为 $e$、质量为 $m$，两次有效散射之间的平均时间为
$\tau$。在弱电场中，电子由电场获得的平均漂移速度大小约为 $v_d=eE\tau/m$。因此
\begin{equation}
  \boxed{\boldsymbol j=\sigma\boldsymbol E,\qquad
  \sigma=\frac{ne^2\tau}{m}}.
\end{equation}
这就是经典德鲁德模型给出的电导率。若从一次碰撞后零定向速度开始计算到下一次碰撞，简单的匀加速
平均会出现 $1/2$；但真实碰撞时间服从分布，玻尔兹曼方程的弛豫时间近似给出上式的标准系数。

该模型也帮助理解金属电阻率随温度升高而增大的现象。金属中载流子浓度 $n$ 大致不变；升温增强晶格
振动，使散射更频繁、$\tau$ 变小，于是 $\sigma$ 减小而 $\rho$ 增大。半导体则不同：升温会显著增加
载流子浓度，$n$ 增大的效果常超过 $\tau$ 减小的效果，因此其电阻通常随温度升高而减小。
在较粗略的经典估计中，也可暂把平均自由程 $\ell$ 看作不变：温度升高使电子典型速率增大，
$\tau\sim\ell/\bar v$ 因而减小，仍得到同一趋势；但金属中更准确的解释需要量子统计与晶格散射理论。

推导上述电导率公式时隐含了几项条件：电子在一次自由飞行中由外场获得的定向速度远小于其无规则运动速度；
散射机制和平均碰撞时间不被外场显著改变；散射足够频繁；载流子浓度不随外场改变。任何一项遭到破坏，
伏安关系都可能偏离直线。稀薄气体在强电场中的放电就是典型例子：电子在两次碰撞之间获得很大能量，
可激发或电离气体分子，使载流子浓度发生剧烈变化，依次出现暗放电、辉光放电和电弧区。

\par\noindent\begin{minipage}{\linewidth}
\gasDischargeFigure
\diagramcaption{气体放电的典型电压--电流曲线。横轴为对数电流标度，各区段明显不满足线性欧姆定律。}\label{ch21:fig:s36-5-7}
\end{minipage}\par\medskip

经典模型仍不能精确描述金属电子，但能把实测电导率换算成有意义的微观尺度。以银为例，
$\sigma\simeq6.3\times10^7\,\mathrm{S\,m^{-1}}$，取
$n\simeq6.0\times10^{28}\,\mathrm{m^{-3}}$，由 $\sigma=ne^2\tau/m$ 得到
\begin{equation}
  \boxed{\tau=\frac{m\sigma}{ne^2}\approx3.8\times10^{-14}\,\mathrm s}.
\end{equation}
若电子典型无规则运动速率取 $\bar v\sim10^6\,\mathrm{m\,s^{-1}}$，平均自由程约为
\begin{equation}
  \boxed{\ell\sim\bar v\tau\approx3.8\times10^{-8}\,\mathrm m}.
\end{equation}
它远大于晶格中原子间的典型距离，说明电子并非每经过一个原子就发生一次经典小球式碰撞；这一事实正是
后来建立能带理论和量子输运理论的重要线索。例如取原子间距约 $0.3\,\mathrm{nm}$，
上述 $38\,\mathrm{nm}$ 自由程约为其 $127$ 倍。这与朴素的“逐原子碰撞”图像相差很大。


% EM teaching enrichment 20260927
\paragraph{浓度不均匀时还有扩散流}
无规则热运动在均匀密度下彼此抵消；若载流子数密度 $n$ 随位置变化，粒子通量可近似写成
$\boldsymbol\Gamma_{\rm diff}=-D\nabla n$，相应电流为 $q\boldsymbol\Gamma_{\rm diff}$。加上漂移项有
\[
 \boldsymbol j=nq\boldsymbol v_d-qD\nabla n.
\]
这里 $D$ 是扩散系数，单位 $\mathrm{m^2/s}$。对电子 $q=-e$，扩散电流方向与电子扩散方向相反。
稳态零电流可以来自漂移与扩散相互抵消，而非各项分别为零。因此含浓度梯度的区域不能只据 $\boldsymbol j=0$ 就断言 $\boldsymbol E=0$；前面的纯欧姆关系省略了这一驱动力。


% EM structure and teaching revision 20260928
\subsection{电流连续性与载流子建立时间}
宏观电流密度由各类载流子的定向运动叠加：$\boldsymbol j=\sum_s n_sq_s\boldsymbol v_s$。电子的平均漂移方向与常规电流方向相反，但其随机热运动通常远快于漂移；随机速度在平衡状态下的平均值为零。
取简单碰撞模型，电子电荷为 $-e$、平均弛豫时间为 $\tau$，恒定外场接通后
\[
 m\frac{\dd\boldsymbol v_d}{\dd t}=-e\boldsymbol E-\frac m\tau\boldsymbol v_d,
 \qquad\boldsymbol v_d(t)=-\frac{e\tau}{m}\boldsymbol E(1-e^{-t/\tau}).
\]
长时间后得到 $\boldsymbol j=ne^2\tau\boldsymbol E/m$。这说明欧姆关系是该线性模型的稳定响应，而非任何材料、任何时间下都成立的定义。
局部电荷守恒写为 $\partial\rho/\partial t+\nabla\cdot\boldsymbol j=0$。稳定电流满足 $\nabla\cdot\boldsymbol j=0$，但不要求电流密度处处相同；在变截面导体中 $j=I/S$，较细处的场和单位体积发热都可能更大。
实际建立过程还受电磁扰动传播和电路储能影响，因此微观碰撞时间不能直接替代整条电路的RC或LR时间常数。


\subsection{专题例题与方法应用}

\begin{origexamplebox}{【惯性维持的电流】}\label{ch21:ex:s36-5-1}
1917 年，托尔曼和斯图尔特的实验直接证明了电子在金属中参与导电。他们让一个铜环以很高的恒定
角速度绕主轴转动，此时环中没有相对晶格的电流；再让铜环突然停止，发现环内出现短暂电流。求从铜环
停止开始流过任一横截面的总电量。结果用环半径 $r$、初始角速度 $\omega$、电子电量绝对值 $e$、
电子质量 $m$ 和铜的电导率 $\sigma$ 表示。记导线横截面积为 $S$。
\end{origexamplebox}\par\nopagebreak[4]

\par\noindent\begin{minipage}{\linewidth}
\tolmanRingFigure
\diagramcaption{解图：铜环晶格骤停后，电子因惯性保持原切向速度；常规电流方向与电子运动方向相反。}
\end{minipage}\par\medskip

\begin{xesolution}
停止前，自由电子随晶格共同转动，其切向速度为 $r\omega$；相对于晶格没有漂移，所以没有电流。晶格
突然停止时，电子速度不能突变，刚停止后的相对漂移速度仍为 $r\omega$。用德鲁德弛豫模型描述碰撞：
\[
  m\frac{\mathrm dv}{\mathrm dt}=-\frac m\tau v,
  \qquad v(0)=r\omega,
\]
故
\[
  v(t)=r\omega e^{-t/\tau}.
\]
电子数密度为 $n$，于是常规电流大小
\[
  |I(t)|=neS\,r\omega e^{-t/\tau}.
\]
从 $t=0$ 积分到电流完全衰减：
\[
  |Q|=\int_0^\infty |I(t)|\,\mathrm dt=neSr\omega\tau.
\]
再利用德鲁德关系 $\sigma=ne^2\tau/m$ 消去 $n$ 与 $\tau$，得到
\[
  \boxed{|Q|=\frac{m\sigma S r\omega}{e}}.
\]
若把转动方向规定为正，由于电子带负电，常规电流以及带符号的 $Q$ 均沿反方向。
总电量与横截面积 $S$ 成正比；单位横截面积通过的电量为
$|Q|/S=m\sigma r\omega/e$。
\end{xesolution}



\begin{origexamplebox}{【暂态电流】}\label{ch21:ex:s36-5-2}
在一个阻值为 $R$ 的圆柱形电阻两端突然加上恒定电压 $U$。考虑载流子的惯性，求通过电阻的电流
随时间的变化关系 $I(t)$。
\end{origexamplebox}\par\nopagebreak[4]

\begin{xesolution}
设电子漂移速度的大小为 $v_d$，取常规电流方向为正。德鲁德模型的平均运动方程可写成
\[
  m\frac{\mathrm dv_d}{\mathrm dt}+\frac m\tau v_d=eE.
\]
圆柱导体内 $E=U/l$，而 $I=neSv_d$。把方程乘以 $neS$，并利用稳态时
$U/R=neS(eE\tau/m)$，可得电流满足
\[
  \tau\frac{\mathrm dI}{\mathrm dt}+I=\frac UR.
\]
突然加压前 $I(0)=0$，解这个一阶方程：
\[
  \boxed{I(t)=\frac UR\left(1-e^{-t/\tau}\right),\qquad t\ge0}.
\]

\par\noindent\begin{minipage}{\linewidth}
\transientCurrentFigure
\diagramcaption{解图：电流以时间常数 $\tau$ 从零趋近稳态值 $U/R$；普通电路中 $\tau$ 极短。}
\end{minipage}\par\medskip

欧姆定律 $I=U/R$ 是 $t\gg\tau$ 后的稳态结果。若理想化地令 $\tau\to0$，才会得到“加压瞬间电流
立刻跳到 $U/R$”的集中参数模型。
\end{xesolution}



\begin{origexamplebox}{【转动的带电环】}\label{ch21:ex:s36-5-3}
一个半径为 $R$、总电荷为 $Q$ 的刚性均匀带电圆环，绕其对称轴以角速度 $\omega$ 匀速转动。
这一过程可以等效为圆形导线中的稳恒电流，求等效电流 $I$。
\par\noindent\begin{minipage}{\linewidth}
\emRotatingRingQuestion
\figcap{刚性带电圆环及给定转动方向，等效电流待求。}\label{ch21:fig:em-source-thirty-sixth-section-3}
\end{minipage}\par\medskip
\end{origexamplebox}\par\nopagebreak[4]

\par\noindent\begin{minipage}{\linewidth}
\rotatingChargedRingFigure
\diagramcaption{，图示取 $Q>0$。圆环每转一周，总电荷 $Q$ 就完整通过任一固定截面一次。}\label{ch21:fig:s36-5-12}
\end{minipage}\par\medskip

\begin{xesolution}
圆环转动周期为
\[
  T=\frac{2\pi}{\omega}.
\]
在一个周期内，环上全部电荷恰好通过任意固定的径向截面一次。由电流定义
\[
  \boxed{I=\frac QT=\frac{Q\omega}{2\pi}}.
\]
若 $Q>0$，电流方向与转动方向相同；若 $Q<0$，常规电流方向相反。半径 $R$ 不出现在结果中，因为
半径增大虽然使线速度增大，却同时使一周路程同比增大，周期仍只由 $\omega$ 决定。
\end{xesolution}



\begin{origexamplebox}{【真空二极管】}\label{ch21:ex:s36-5-4}
真空二极管由相对放置的阴极和阳极平板构成，两极板面积均为 $S$，距离为 $L$。稳定工作时有电流
$I$ 通过。阴极发射充分，工作于零初速的空间电荷限制状态，忽略边缘效应。取阴极处 $x=0$ 且电势为零，阳极处 $x=L$、电势为 $U>0$。依次回答：

(1) 写出两极间电势分布 $\varphi(x)$ 满足的方程；

(2) 假设电子从阴极由静止出发，求电势为 $\varphi(x)$ 处的电子速度；

(3) 利用稳恒电流与位置无关，写出电子数密度 $n(x)$；

(4) 消去 $n(x),v(x)$，得到关于 $\varphi(x)$ 的微分方程及边界条件；

(5) 积分证明查尔德--朗缪尔定律 $I=KU^{3/2}$，并求 $K$。
\par\noindent\begin{minipage}{\linewidth}
\vacuumDiodeFigure
\diagramcaption{电子由阴极射向阳极，图中 $I$ 表示相反方向的常规电流大小。}\label{ch21:fig:s36-5-13}
\end{minipage}\par\medskip
\end{origexamplebox}\par\nopagebreak[4]


\begin{xesolution}
电子云带负电，会改变板间电势，这正是空间电荷限制电流的来源。下面把 $I$ 作为电流大小。

\textbf{(1) 泊松方程。} 电子电荷密度 $\rho_q=-en(x)$，一维泊松方程为
\[
  \frac{\mathrm d^2\varphi}{\mathrm dx^2}
  =-\frac{\rho_q}{\varepsilon_0}=\frac{en(x)}{\varepsilon_0}.
\]

\textbf{(2) 能量关系。} 电子从阴极静止出发，电势能减少量转化为动能：
\[
  \frac12mv^2=e\varphi,
  \qquad v(x)=\sqrt{\frac{2e\varphi(x)}m}.
\]

\textbf{(3) 稳恒电流。} 任意截面上单位时间通过的电荷量相同，故
\[
  I=en(x)v(x)S,
  \qquad n(x)=\frac{I}{eSv(x)}.
\]

\textbf{(4) 消元与边界条件。} 代回泊松方程：
\[
  \frac{\mathrm d^2\varphi}{\mathrm dx^2}
  =A\varphi^{-1/2},
  \qquad
  A=\frac{I}{\varepsilon_0S}\sqrt{\frac m{2e}}.
\]
两板电势给出 $\varphi(0)=0,\ \varphi(L)=U$。空间电荷限制的极限状态下，阴极表面电场刚好为零，
所以还应取 $\varphi'(0)=0$。注意 $E_x=-\varphi'$，电子受力 $F_x=-eE_x=e\varphi'$：
若 $E_x(0)>0$，零初速电子受到朝阴极的力，不能向阳极出发；若 $E_x(0)<0$，电子可被拉向阳极，
但在发射不足时电流尚未达到空间电荷限制值。发射充分时，负空间电荷削弱阴极附近的加速场，
极限恰为 $E_x(0)=0$。因此这里的零场边界来自题设工作状态，不是任意稳恒电流的通用条件。

\textbf{(5) 积分。} 方程两边乘 $\varphi'$：
\[
  \varphi''\varphi'=A\varphi^{-1/2}\varphi'.
\]
从阴极积分并用 $\varphi(0)=\varphi'(0)=0$，得到
\[
  \frac12(\varphi')^2=2A\sqrt\varphi,
  \qquad \varphi'=2\sqrt A\,\varphi^{1/4}.
\]
再积分一次：
\[
  \frac43\varphi^{3/4}=2\sqrt A\,x,
  \qquad
  \varphi(x)=\left(\frac32\sqrt A\,x\right)^{4/3}.
\]
用 $\varphi(L)=U$ 消去 $A$，得到
\[
  \boxed{\varphi(x)=U\left(\frac xL\right)^{4/3}}.
\]

\par\noindent\begin{minipage}{\linewidth}
\childLangmuirFigure
\diagramcaption{解图：空间电荷限制状态下的电势曲线。阴极处切线水平，随后按 $x^{4/3}$ 增长。}
\end{minipage}\par\medskip

最后将 $A=(4/9)U^{3/2}/L^2$ 与其定义比较：
\[
  \boxed{I=\frac49\varepsilon_0S\sqrt{\frac{2e}{m}}
  \frac{U^{3/2}}{L^2}=KU^{3/2}},
  \qquad
  \boxed{K=\frac49\frac{\varepsilon_0S}{L^2}\sqrt{\frac{2e}{m}}}.
\]
\end{xesolution}



\section{伏安特性、电阻与测量}

\subsection{元件的伏安特性}


电路由元件按一定方式连接而成。对电路设计者而言，常可暂时不追究元件内部结构，而只研究端口处可测量
的电压和电流。最简单的是二端元件。图中先人为规定端口电压 $V$ 的正负方向和电流 $I$ 的测试方向；
这种规定不等于元件本身一定具有固定正、负极，计算出的负号只说明真实方向与规定方向相反。

\par\noindent\begin{minipage}{\linewidth}
\ivCharacteristicFigure
\diagramcaption{二端元件及其 $I$--$V$ 特性。标出的正、负号与箭头只是统一的测试约定。}\label{ch21:fig:s36-5-1}
\end{minipage}\par\medskip

若元件稳定工作时，同一个端电压总对应同一个端电流，就可用函数 $I=I(V)$ 描述它的外部性能，称为
伏安关系；把它画在 $I$--$V$ 平面上，便得到伏安特性曲线。不同元件的曲线可以完全不同，线性电阻
只是其中最简单的一类。使用这种二端描述还隐含两个条件：一是所研究的回路区域内没有随时间变化的磁通，
否则端电压会受感应电场影响而与路径有关；二是元件内部不能持续积累净电荷，否则流入与流出的电流并不
相等。满足这些条件后，元件的伏安关系加上连接规律就足以决定整个稳恒电路的行为。


% EM teaching enrichment 20260927
\paragraph{工作点与两种电阻}
电动势 $\mathcal E$、内阻 $r$ 的电源驱动非线性元件时，电源端口满足 $U=\mathcal E-rI$。它与元件曲线 $I=I(U)$ 的交点才是实际工作点，不能把电源电动势直接当作元件端电压。
在该点，$R_{\rm dc}=U/I$ 是静态电阻，$r_d=\dd U/\dd I$ 是小变化对应的微分电阻。若元件满足 $U=aI+bI^3$，则两者分别为 $a+bI^2$ 与 $a+3bI^2$；它们仅在线性情形相同。微分关系只能描述工作点附近的小扰动。


\subsection{欧姆定律}


欧姆定律给出一类导体两端电压与其中电流的实验关系。它并不直接解释导电的微观机制，却在广泛的材料
和工作条件下极其准确，因而成为电路分析中使用最频繁的规律之一。格奥尔格·西蒙·欧姆在十九世纪系统
研究了这一关系；后来以他的姓氏命名电阻单位欧姆，并用纪念邮票表现电阻和公式。

\ohmHistoryFigure

在温度和其他外部条件不变时，若导体两端电压为 $U$、通过电流为 $I$，欧姆定律写成
\begin{equation}
  \boxed{U=IR,\qquad I=\frac UR}.
\end{equation}
比例系数 $R$ 是电阻，单位为欧姆，$1\,\Omega=1\,\mathrm V/\mathrm A$。对长度为 $l$、横截面积为
$A$ 的均匀圆柱形导体，实验进一步给出
\begin{equation}
  \boxed{R=\rho\frac lA},
\end{equation}
其中 $\rho$ 是材料在给定状态下的电阻率。长度越长，载流子沿途经历的散射越多；截面积越大，可并行
输运电荷的通道越多，因而分别出现 $R\propto l$ 与 $R\propto1/A$。

\par\noindent\begin{minipage}{\linewidth}
\cylindricalResistorFigure
\diagramcaption{长度为 $l$、截面积为 $A$ 的均匀圆柱形电阻。电流方向与端电压测试方向一致。}\label{ch21:fig:s36-5-4}
\end{minipage}\par\medskip

理想电阻的伏安关系是一条过原点的直线。若纵轴为 $I$、横轴为 $V$，直线斜率为 $1/R$；曲线越陡，
相同电压产生的电流越大，电阻越小。

\par\noindent\begin{minipage}{\linewidth}
\ohmicIVFigure
\diagramcaption{理想电阻的 $I$--$V$ 特性。直线斜率是电导 $1/R$。}\label{ch21:fig:s36-5-5}
\end{minipage}\par\medskip

把宏观欧姆定律写成局部形式。对图~\ref{ch21:fig:s36-5-4} 的均匀导体，内部电场 $E=U/l$，电流密度 $j=I/A$。
由 $U=IR$ 与 $R=\rho l/A$ 得
\begin{equation}
  \boxed{\boldsymbol j=\frac1\rho\boldsymbol E=\sigma\boldsymbol E},
\end{equation}
其中 $\sigma=1/\rho$ 为电导率。它不含导体的长度和截面积，直接描述材料局部的导电能力。上式只在
各向同性、线性且工作状态不改变的欧姆材料中成立；对各向异性材料，$\sigma$ 一般应写成张量。


% EM teaching enrichment 20260927
\paragraph{非均匀导体的电阻积分与功率检验}
缓慢变截面的细导体可近似看成许多串接薄片，故
$\dd R=\rho(x)\dd x/A(x)$，$R=\int\rho(x)\dd x/A(x)$。若长度 $l$ 内截面积由 $A_0$ 线性增至 $2A_0$、电阻率不变，得到
\[
 R=\frac{\rho l}{A_0}\int_0^1\frac{\dd u}{1+u}
 =\frac{\rho l}{A_0}\ln2.
\]
这里采用一维电流近似；骤变截面附近的三维电流拥挤不能忽略时，应另求电流场。
局部耗散功率密度为 $\boldsymbol j\cdot\boldsymbol E=j^2/\sigma$，对体积积分应回到 $UI=I^2R$。单位和总功率是检查非均匀导体积分的有效方法。


% EM structure and teaching revision 20260928
\subsection{测量加载与电功率的局部核对}
理想电压表不取电流，理想电流表不产生压降；实际仪表会改变被测状态。电阻 $R_x$ 与内阻 $R_V$ 的电压表并联时，若电流表测的是总电流，则
\[
 R_{\text{读}}=\frac UI=\frac{R_xR_V}{R_x+R_V}<R_x.
\]
相反，若电压表跨接 $R_x$ 与电流表内阻 $R_A$ 的串联组合，而电流表只测流过 $R_x$ 的电流，则 $R_{\text{读}}=R_x+R_A$。先画出仪表端点，再分析电流究竟流过哪些支路，误差方向便不需死记。
场描述给出局部功率密度 $p=\boldsymbol j\cdot\boldsymbol E$。对欧姆介质，$p=j^2/\sigma=\sigma E^2\ge0$；积分后与端口功率 $UI$ 相等。以长直均匀电阻为例，$j=I/S$、$E=U/\ell$，故 $\int p\dd V=UI$。
这里的非负性只适用于耗散电阻。电源支路可以输出功率，其端口 $UI$ 的符号取决于是否采用电流流入电压正端的关联参考方向。


\subsection{专题例题与方法应用}

\begin{origexamplebox}{【导电厚球壳】}\label{ch21:ex:s36-5-5}
两个同心金属球壳的内、外半径分别为 $a,b$，其间充满一种特殊导电材料。该材料在各点的电导率
与当地电场大小成正比，即 $\sigma=KE$。若两球壳之间的电势差为 $U$，求两壳之间的电流。
\par\noindent\begin{minipage}{\linewidth}
\emConductingShellQuestion
\figcap{导电厚球壳的内外半径；两球面之间加给定电压。}\label{ch21:fig:em-source-thirty-sixth-section-5}
\end{minipage}\par\medskip
\end{origexamplebox}\par\nopagebreak[4]

\par\noindent\begin{minipage}{\linewidth}
\conductingShellFigure
\diagramcaption{球对称性保证电流密度和电场均沿径向，且通过任一同心球面的电流相同。}\label{ch21:fig:s36-5-14}
\end{minipage}\par\medskip

\begin{xesolution}
稳恒状态下，任意半径 $r$ 的同心球面都通过同一电流 $I$。因此
\[
  j(r)=\frac{I}{4\pi r^2}.
\]
材料的局部伏安关系不是线性的，而是
\[
  j=\sigma E=KE^2.
\]
两式联立：
\[
  E(r)=\sqrt{\frac{I}{4\pi K}}\frac1r.
\]
内外球壳的电势差为电场沿径向的积分：
\[
  U=\int_a^bE(r)\,\mathrm dr
  =\sqrt{\frac{I}{4\pi K}}\ln\frac ba.
\]
于是
\[
  \boxed{I=\frac{4\pi K}{[\ln(b/a)]^2}U^2}.
\]
这里 $I\propto U^2$，明确说明该材料不满足线性欧姆定律。若交换两球壳电势高低，电流方向反向；
上式给出其大小。
\end{xesolution}



\begin{origexamplebox}{【电离计数器】}\label{ch21:ex:s36-5-6}
最简单的电离室由一对面积为 $S$、间距为 $D$ 的平行板电极构成，置于充有混合气体的气密容器中。
阴极电势为 $-U_0$，阳极通过负载电阻 $R$ 接地。带电粒子平行极板入射，在距阳极 $x_0$ 处沿径迹
产生 $N$ 对电子--正离子；产生一对粒子所需平均能量 $W$ 由气体决定。假设全部电荷均被收集，不发生
次级电离和电子俘获，电子、正离子漂移速率分别为 $v^-,v^+$。求：

(1) 两极间电场强度；(2) 电离室电容 $C$；(3) $R$ 很大时信号电压 $\Delta U(t)$；
(4) 当 $v^-\gg v^+$ 时信号的特征；(5) $R$ 的作用不可忽略时的 $\Delta U(t)$。
\par\noindent\begin{minipage}{\linewidth}
\emIonChamberQuestion
\figcap{平行板电离室的装置、入射径迹与外接电路；竖直侧壁表示绝缘容器边界。}\label{ch21:fig:em-source-thirty-sixth-section-6}
\end{minipage}\par\medskip
\end{origexamplebox}\par\pagebreak[1]

\par\noindent\begin{minipage}{\linewidth}
\ionChamberFigure
\diagramcaption{题图：平行板电离室。$x$ 从阳极指向阴极；电子向阳极、正离子向阴极漂移并共同形成感应信号。}
\end{minipage}\par\medskip

\begin{xesolution}
若入射粒子在气体中沉积能量 $E_{\rm dep}$，则理想情况下 $N=E_{\rm dep}/W$；本题已直接给出 $N$。
以下取阳极对地电压变化为 $\Delta U$，并采用小信号条件 $Ne/C\ll U_0$，使收集过程中的场强和漂移速率
可近似不变。信号大小记为 $u=-\Delta U\ge0$，避免把实际负脉冲误画成正电压。

\textbf{(1) 电场。} 忽略边缘效应，板间电场均匀，大小
\[
  \boxed{E=\frac{U_0}{D}},
\]
方向由阳极指向阴极。电子受力方向与电场相反，故移向阳极；正离子移向阴极。

\textbf{(2) 电容。} 两平行板构成电容器：
\[
  \boxed{C=\frac{\varepsilon_0S}{D}}.
\]

\textbf{(3) $R$ 很大。} 先把两板电势保持不变。位于 $x$ 的电荷 $q$ 在阳极上感应的总电荷为
$Q_A=-q(1-x/D)$。可由叠加与静电互易关系得到：令阳极为单位电势、阴极为零，板间权重电势为
$\psi_A(x)=1-x/D$，于是 $Q_A+q\psi_A(x)=0$。也可将很薄的均匀面电荷放在两板间，
联立两侧场强跳变和总电压不变条件，得到同一关系。
电子向 $-x$ 移动、正离子向 $+x$ 移动时，保持阳极接地所需的感应正电荷都增加，增量为
$e|\Delta x|/D$。若阳极近似隔离，这部分正电荷不能由地补入，阳极电压因而降低。
在时刻 $t$，电子已经移动
$\min(v^-t,x_0)$，正离子已经移动 $\min(v^+t,D-x_0)$，所以总感应电荷大小
\[
  Q_{\rm sig}(t)=\frac{Ne}{D}
  \left[\min(v^-t,x_0)+\min(v^+t,D-x_0)\right].
\]
$R$ 极大时可忽略泄放，信号电压大小为
\[
  \boxed{|\Delta U(t)|=\frac{Ne}{CD}
  \left[\min(v^-t,x_0)+\min(v^+t,D-x_0)\right]}.
\]
最终两类电荷都到达电极，括号内等于 $D$，故 $|\Delta U(\infty)|=Ne/C$，与产生位置 $x_0$ 无关。

\textbf{(4) 快、慢两段。} 定义电子收集时间和离子收集时间
\[
  t_e=\frac{x_0}{v^-},\qquad t_i=\frac{D-x_0}{v^+}.
\]
在常见的 $t_e<t_i$ 情形，信号大小先以较大斜率快速上升；电子收集完成后，只剩离子造成的缓慢上升，
在 $t_i$ 达到平台。精确条件是 $x_0<Dv^-/(v^-+v^+)$；仅有 $v^-\gg v^+$ 不能排除极靠近阴极的
$t_i<t_e$ 情形。后一情形先失去离子分量，剩下电子分量直至 $t_e$。两种情况都由上式的两个最小值统一描述，
实际电压为 $\Delta U=-u$，最终平台为 $-Ne/C$。

\par\noindent\begin{minipage}{\linewidth}
\ionSignalFigure
\diagramcaption{解图：忽略泄放时的信号大小。图中把常见的 $t_e<t_i$ 情形画成快速前沿加缓慢离子尾。}
\end{minipage}\par\medskip

\textbf{(5) 有限负载电阻。} 读出节点相当于电容 $C$ 与电阻 $R$ 并联，信号满足
\[
  C\frac{\mathrm d u}{\mathrm dt}+\frac{u}{R}=i_s(t),\qquad \Delta U=-u.
\]
这里 $i_s$ 表示使阳极电势降低的信号电流大小，是两个矩形脉冲之和（$t\ge0$）：
\[
  i_s(t)=\frac{Ne}{D}\left[v^-H(t_e-t)+v^+H(t_i-t)\right],
\]
其中 $H$ 为阶跃函数。令 $\tau_R=RC$。对于某个持续时间为 $T_k$、幅值为
$I_k=Ne v_k/D$ 的分量，零初值解为
\[
  u_k(t)=
  \begin{cases}
    RI_k\left(1-e^{-t/\tau_R}\right),&0\le t\le T_k,\\[2mm]
    RI_k\left(1-e^{-T_k/\tau_R}\right)e^{-(t-T_k)/\tau_R},&t>T_k.
  \end{cases}
\]
分别取 $(v_k,T_k)=(v^-,t_e)$ 与 $(v^+,t_i)$，得到 $\Delta U(t)=-u_e(t)-u_i(t)$。当
$RC\gg\max(t_e,t_i)$ 时，在收集时间内展开指数，退化为前面近似无泄放的折线；当 $RC$ 较小时，信号在收集过程中
已明显经电阻衰减，收集结束后则按 $e^{-t/(RC)}$ 回到零。
\end{xesolution}



\section{电路守恒与节点方程}

\subsection{简单电阻网络}


电阻最基本的连接方式是串联和并联。串联支路中各电阻通过同一电流 $I$，总电压是各段电压之和；
并联支路中各电阻接在同一对节点之间，端电压相同，总电流是各支路电流之和。

\par\noindent\begin{minipage}{\linewidth}
\seriesParallelResistorsFigure
\diagramcaption{电阻的串联与并联。判断连接方式的依据是节点关系，而不是图形看起来是否横排或竖排。}\label{ch21:fig:s36-5-8}
\end{minipage}\par\medskip

对 $n$ 个串联电阻，$U_i=IR_i$，所以
\[
  U=U_1+\cdots+U_n=I(R_1+\cdots+R_n).
\]
等效电阻为
\begin{equation}
  \boxed{R_{\rm s}=R_1+R_2+\cdots+R_n}.
\end{equation}
对 $n$ 个并联电阻，公共端电压为 $U$，第 $i$ 条支路电流为
\begin{equation}
  \boxed{I_i=\frac{U}{R_i}}.
\end{equation}
总电流
$I=\sum_i I_i=U\sum_i(1/R_i)$，故等效电阻满足
\begin{equation}
  \boxed{\frac1{R_{\rm p}}=\frac1{R_1}+\frac1{R_2}+\cdots+\frac1{R_n}}.
\end{equation}
串联等效电阻大于任一分电阻，并联等效电阻小于任一分电阻。若网络能够逐级辨认出纯串联或纯并联，
即可反复用上述串并联公式化简；但含桥路、多个电源或交叉耦合支路的网络通常不能这样处理。


% EM teaching enrichment 20260927
\paragraph{分压、分流及仪表加载}
两电阻串联时 $U_2=UR_2/(R_1+R_2)$；两电阻并联、总电流为 $I$ 时 $I_1=IR_2/(R_1+R_2)$。分流分子是另一支路的电阻，来源是先求共同电压。
若用内阻 $R_V$ 的电压表测 $R_2$，被测支路已变为 $R_2\parallel R_V$，读数是
\[
 U_V=U\frac{R_2R_V/(R_2+R_V)}{R_1+R_2R_V/(R_2+R_V)}.
\]
只有 $R_V\gg R_2$ 且加载误差可忽略时才恢复理想分压。先把仪表画成实际元件，再讨论它的极限，可防止“测量不影响电路”的未经检验假设。


\subsection{电阻网络定理与基本术语}


一般电路是由元件按拓扑关系连接成的网络。图中的每个方框都代表一个元件，它可以是电阻、电源或更复杂
的二端器件。两个或更多元件的连接点称为节点；同一段理想导线上的各点电势相同，属于同一个节点。
相邻两节点之间包含元件的一段称为支路。若一个节点只连接两条支路，通常只是串接中的平凡节点；三条及
以上支路相接时会发生电流分流。沿若干支路从一个节点出发，最后回到原节点的闭合路径称为回路。

\par\noindent\begin{minipage}{\linewidth}
\networkTermsFigure
\diagramcaption{电路分析中的元件、节点、支路和回路。虚线只标示所选回路，不是导线。}\label{ch21:fig:s36-5-9}
\end{minipage}\par\medskip


% EM teaching enrichment 20260927
\paragraph{画图之前先识别节点}
所有用无电阻导线直接连通的点标为同一节点；交叉但未连接的导线不合并。随后把每个元件记成“节点 $a$ 到节点 $b$ 的支路”，图形折叠或拉伸便不会改变电路。
含 $N$ 个节点、$b$ 条支路的连通网络中，选 $N-1$ 条无闭环支路可连接全部节点，称为生成树。每补入一条其余支路，恰增加一个独立回路，因此独立回路数为 $b-N+1$。这是方程数量的来源，也能用来检查是否把交叉线误当节点。


\subsection{基尔霍夫电流定律（KCL）}


基尔霍夫电流定律是电荷守恒在集中参数电路中的表达：任一节点流入电流之和等于流出电流之和；
等价地，按统一符号约定计入后，节点各支路电流的代数和为零。未知支路电流的方向可以预先任意假设，
若解得负值，说明真实方向与假设方向相反。

\par\noindent\begin{minipage}{\linewidth}
\kclKvlFigure
\diagramcaption{左图为 KCL 的支路电流约定，右图为同一网络的支路电压和回路方向。所有字母标注均不覆盖元件。}\label{ch21:fig:s36-5-10}
\end{minipage}\par\medskip

按左图箭头，对四个节点分别可写为
\[
\begin{array}{ll}
N_1:&0=-i_1-i_2-i_3,\\
N_2:&0=i_1+i_4-i_6,\\
N_3:&0=i_2-i_4-i_5,\\
N_4:&0=i_3+i_5+i_6.
\end{array}
\]
最后一式中的 $i_6$ 由底边支路方向确定。把四式相加得到恒等式 $0=0$，因此含 $N$ 个节点的连通网络
只有 $N-1$ 个独立 KCL 方程。少写会丢失约束，多写则只是重复已有信息。


% EM teaching enrichment 20260927
\paragraph{节点暂时储存电荷时的完整守恒式}
用闭面包围一个小节点，流入减流出等于 $\dd Q_{\rm node}/\dd t$。稳恒状态右边为零，才得到常用 KCL。
在集中参数模型中，若节点对地有不可忽略电容 $C$，就应把 $C\dd V/\dd t$ 作为一条支路电流列入。例如通过电阻 $R$ 给该节点充电时，
\[
 \frac{\mathcal E-V}{R}=C\frac{\dd V}{\dd t}.
\]
常用节点法不是否定电荷积累，而是把积累明确归入电容元件。先选电流正向再写等式，解出的负数不应被擅自改成正数。


\subsection{基尔霍夫电压定律（KVL）}


基尔霍夫电压定律是稳恒电路能量守恒的等价表达：沿任一闭合回路，各支路电压的代数和为零。使用时
先规定每个支路电压的正方向，再选定绕行方向；沿支路从正端走向负端记为一个符号，反向走则记为相反符号。
对图~\ref{ch21:fig:s36-5-10} 右侧标出的四个回路，可写成
\[
\begin{array}{ll}
L_1:&0=-v_1+v_2+v_4,\\
L_2:&0=-v_2+v_3-v_5,\\
L_3:&0=-v_4+v_5-v_6,\\
L_4:&0=v_1+v_6-v_3.
\end{array}
\]
四式相加同样得到恒等式，所以它们并不独立。前三个小回路已经覆盖全部支路，外围回路 $L_4$ 可由前三式
线性组合得到。对连通网络，独立回路数等于“支路数 $-$ 节点数 $+1$”。KCL、KVL 再配合每个元件的
伏安关系，原则上可以求出全部支路电流和元件电压。


% EM teaching enrichment 20260927
\paragraph{一圈电压与一圈能量如何对应}
对电动势 $\mathcal E$、内阻 $r$、外电阻 $R$ 的闭合回路，沿电流绕行，源内非静电作用提升每单位电荷的能量 $\mathcal E$，电阻段合计降低 $I(R+r)$，故
$\mathcal E-Ir-IR=0$。
乘以 $I$ 得电源输出功率 $\mathcal EI=I^2r+I^2R$，同时核对了符号与能量。
若回路有变化磁通，应加入感应电动势；不能一面把感应效应遗漏，一面仍以纯静电势差的闭合和为零处理全部电场。


\subsection{节点电压法}


直接把所有支路电流和元件电压都列为未知量，哪怕小网络也会产生很大的方程组。节点电压法先选一个节点
作为参考地，把其电势定为零；随后只把其余独立节点电势作为未知量。支路电压由两端节点电势之差自动给出，
因此 KVL 自动满足；再对未知节点写 KCL，就得到数量最少的一组方程。

\par\noindent\begin{minipage}{\linewidth}
\nodeMethodFigure
\diagramcaption{节点电压法示例。取 $c$ 为参考地，理想电压源固定 $a$ 点电势，只剩 $b$ 点电势未知。}\label{ch21:fig:s36-5-11}
\end{minipage}\par\medskip

在图中令 $\varphi_c=0$，电压源使 $\varphi_a=E_0$，设 $\varphi_b=e$。以流入 $b$ 点为正：
从 $a$ 经 $R_1$ 流入的电流是 $(E_0-e)/R_1$，电流源注入 $I_0$，从 $b$ 经 $R_2$ 流出的电流是
$e/R_2$。KCL 给出
\[
  \frac{E_0-e}{R_1}+I_0-\frac e{R_2}=0.
\]
解得
\[
  \boxed{e=\frac{E_0/R_1+I_0}{1/R_1+1/R_2}}.
\]
节点电势一旦求出，每个元件两端的电压和各支路电流便可由其伏安关系依次得到。参考节点原则上可以任意
选择；经验上，选与较多理想电压源相连的节点作参考地，往往能减少辅助未知量，使代数最简。

% EM teaching enrichment 20260927
\paragraph{悬浮电压源与解后的核对}
若理想电压源位于两个未知节点 $a,b$ 之间，不能把其电流写成“电压除以零电阻”。可把两节点合成超节点，对流向外界的支路写 KCL，再补 $V_a-V_b=\mathcal E$。
例如 $a,b$ 分别经 $R_a,R_b$ 接地，则
$V_a/R_a+V_b/R_b=0$，与电压约束联立得
$V_a=\mathcal E R_a/(R_a+R_b)$、$V_b=-\mathcal E R_b/(R_a+R_b)$。
负节点电势是相对参考点而言，并非异常。解后应检查各节点电流守恒，以及全部元件吸收功率的代数和为零。




% EM structure and teaching revision 20260928
\subsection{把拓扑、方程与功率检查分开}
设连通电路有 $N$ 个节点和 $B$ 条支路。固定一个参考电势后，可独立写出 $N-1$ 个节点电流方程；若改用支路电流未知量，则还需 $B-N+1$ 个独立回路方程及元件关系。逐个面孔写回路虽方便，但非平面画法中的线条交叉未必表示新节点。
例如节点电势 $V$ 经 $R_1$ 接到已知电势 $E$、经 $R_2$ 接地，同时有恒流源 $I_s$ 向该节点注入，则
\[
 \frac{V-E}{R_1}+\frac V{R_2}-I_s=0,
 \qquad V=\frac{E/R_1+I_s}{1/R_1+1/R_2}.
\]
令 $I_s=0$ 应退回分压公式；令 $R_2\to\infty$ 应得到 $V=E+I_sR_1$。这类极限检查比单纯把计算结果代回原方程更容易发现参考方向选错。
若理想电压源连接两个未知电势节点，把两节点作为超节点写总电流平衡，再补电压差约束。解完后分别计算各电阻消耗和各电源输出，检查全电路净功率为零；这一步同时检验数值、支路方向和正负号。


\subsection{专题例题与方法应用}

\begin{origexamplebox}{【圆环电阻】}\label{ch21:ex:s36-5-7}
一段总电阻为 $R$ 的均匀电阻丝弯成圆环。在环上固定一点 $a$，另一接点 $b$ 可沿电阻丝移动。
先在 $a,b$ 间接恒定电压 $U$，问 $b$ 在何处时整条电阻丝的发热功率最小，最小值是多少？若改接
电流为 $I$ 的恒流源，结论又如何？
\par\noindent\begin{minipage}{\linewidth}
\emRingResistorQuestion
\figcap{均匀电阻圆环及接点；两段弧的电阻记为 $R_1,R_2$。}\label{ch21:fig:em-source-thirty-sixth-section-7}
\end{minipage}\par\medskip
\end{origexamplebox}\par\nopagebreak[4]

\par\noindent\begin{minipage}{\linewidth}
\ringResistorFigure
\diagramcaption{题图：接点 $a,b$ 把圆环分成电阻为 $R_1,R_2$ 的两段，两段端点相同，故彼此并联。}
\end{minipage}\par\medskip

\begin{xesolution}
两段弧的电阻之和固定：$R_1+R_2=R$。它们并联后的等效电阻为
\[
  R_{\rm eq}=\frac{R_1R_2}{R_1+R_2}=\frac{R_1R_2}{R}.
\]

\textbf{恒压源。} 总功率
\[
  P_U=\frac{U^2}{R_{\rm eq}}=\frac{U^2R}{R_1R_2}.
\]
在和固定时，乘积满足 $R_1R_2\le(R/2)^2$；要使功率最小，应使乘积最大，故
\[
  R_1=R_2=\frac R2.
\]
均匀电阻丝上这意味着 $b$ 位于 $a$ 的直径对面，且
\[
  \boxed{P_{U,\min}=\frac{4U^2}{R}}.
\]

\textbf{恒流源。} 此时
\[
  P_I=I^2R_{\rm eq}=I^2\frac{R_1R_2}{R}.
\]
要使功率最小，应使 $R_1R_2$ 尽可能小，即让 $b$ 趋近 $a$，一段电阻趋近零，因而
\[
  \boxed{P_{I,\min}=0}
\]
（若规定 $a,b$ 必须是不同点，则 $0$ 是下确界，只能无限接近而不能取到）。相反，$b$ 位于直径对面
时 $R_{\rm eq}$ 最大，恒流情况下反而得到最大功率 $I^2R/4$。这两种结论相反，根源是恒压时
$P=U^2/R_{\rm eq}$，恒流时 $P=I^2R_{\rm eq}$。
\end{xesolution}



\begin{origexamplebox}{【双电源供电】}\label{ch21:ex:s36-5-8}
两个电动势和内阻分别为 $(\varepsilon_1,r_1)$、$(\varepsilon_2,r_2)$ 的实际电源同极性并联，
共同连接负载电阻 $R$。

(1) 利用基尔霍夫方程求通过 $R$ 的电流；

(2) 对结果变形，说明它怎样给出关于电源叠加和戴维南等效的启发。
\par\noindent\begin{minipage}{\linewidth}
\dualSourceFigure
\diagramcaption{每个实际电源均由理想电动势与内阻串联构成，三条支路跨接同一对节点。}\label{ch21:fig:s36-5-15}
\end{minipage}\par\medskip
\end{origexamplebox}\par\nopagebreak[4]


\begin{xesolution}
记左、右公共节点电势差为 $V=V_L-V_R$，并把三条支路电流都暂按从左向右取正。负载支路电流为
$V/R$；对第 $k$ 个电源支路，按图示极性有 $V=\varepsilon_k+i_kr_k$，所以
$i_k=(V-\varepsilon_k)/r_k$。在左节点写 KCL：
\[
  \frac VR+\frac{V-\varepsilon_1}{r_1}
  +\frac{V-\varepsilon_2}{r_2}=0.
\]
解出 $V$ 并除以 $R$，负载电流为
\[
  \boxed{I=\frac VR
  =\frac{\varepsilon_1r_2+\varepsilon_2r_1}
  {r_1r_2+R(r_1+r_2)}}.
\]

把两个电源及其内阻看成一个二端网络。定义
\[
  \varepsilon_{\rm eq}
  =\frac{\varepsilon_1/r_1+\varepsilon_2/r_2}{1/r_1+1/r_2}
  =\frac{\varepsilon_1r_2+\varepsilon_2r_1}{r_1+r_2},
  \qquad
  r_{\rm eq}=\frac{r_1r_2}{r_1+r_2}.
\]
原结果立即化成
\[
  \boxed{I=\frac{\varepsilon_{\rm eq}}{R+r_{\rm eq}}}.
\]
这说明并联电源网络从负载端看，等效为一个电动势 $\varepsilon_{\rm eq}$ 与内阻 $r_{\rm eq}$ 串联的
实际电源，正是戴维南等效的形式。同时，$V$ 对 $\varepsilon_1,\varepsilon_2$ 是线性的，每个电源的
贡献可以分别计算再相加，体现了线性网络的叠加思想。若两电源极性相反，只需把其中一个
$\varepsilon_k$ 按符号取负；公式仍然成立。
\end{xesolution}
}

{
\graphicspath{{/Users/dongxingxi/Documents/物理竞赛一轮资料/sample/xes-original/thirty-seventh-section/}}
\section{等效电源、对称性与电桥}

% Technical figures are redrawn as vectors.  Every text node is transparent;
% rectangles are used only for actual circuit elements or explanatory regions.
\tikzset{
  xeslabel/.style={font=\small,text=Ink,fill=none,inner sep=.45pt},
  xessmall/.style={font=\footnotesize,text=TextGray,fill=none,inner sep=.35pt},
  xestiny/.style={font=\scriptsize,text=TextGray,fill=none,inner sep=.25pt},
  circuit/.style={draw=Ink,line width=.82pt,line cap=round,line join=round},
  activewire/.style={draw=DeepBlue,line width=1.05pt,line cap=round,line join=round},
  hiddenwire/.style={draw=TextGray,densely dashed,line width=.72pt,line cap=round},
  solutionwire/.style={draw=DeepBlue!78,line width=.9pt,line cap=round,line join=round},
  currentarrow/.style={-{Latex[length=1.8mm,width=1.2mm]},draw=DeepBlue,line width=.75pt},
  nodepoint/.style={circle,fill=Ink,draw=none,inner sep=1.45pt},
  bluenode/.style={circle,fill=DeepBlue,draw=white,line width=.3pt,inner sep=1.7pt},
  goldnode/.style={circle,fill=Gold!85!black,draw=white,line width=.3pt,inner sep=1.7pt},
  port/.style={circle,fill=white,draw=Ink,line width=.75pt,inner sep=1.6pt},
  sourcecircle/.style={circle,draw=Ink,line width=.75pt,minimum size=.54cm,inner sep=0pt},
  networkblob/.style={draw=DeepBlue!75,fill=DeepBlue!4,line width=.8pt},
  eqbox/.style={draw=none,fill=none,inner sep=2pt}
}


\newcommand{\zigres}[2]{%
  \draw[circuit,decorate,decoration={zigzag,segment length=2.25mm,amplitude=.72mm}] (#1)--(#2);}

\newcommand{\drawvoltage}[2]{%
  \draw[circuit] (#1)--++(0,.26);
  \draw[circuit] ($(#1)+(-.16,.26)$)--($(#1)+(.16,.26)$);
  \draw[circuit] ($(#1)+(-.27,.39)$)--($(#1)+(.27,.39)$);
  \draw[circuit] ($(#1)+(0,.39)$)--(#2);}

\newcommand{\superpositionFigure}{%
\begin{center}
\begin{tikzpicture}[x=1cm,y=1cm]
  \begin{scope}[shift={(-3.25,0)}]
    \coordinate (a) at (-1.7,1.05);\coordinate (b) at (.15,1.05);\coordinate (c) at (.15,-1.0);
    \coordinate (d) at (-1.7,-1.0);\coordinate (r) at (1.6,0);
    \draw[circuit] (a)--(-1.7,.32);
    \draw[circuit] (-1.93,.18)--(-1.47,.18);
    \draw[circuit] (-1.82,.02)--(-1.58,.02);
    \draw[circuit] (-1.7,.02)--(d);
    \zigres{a}{b}\zigres{b}{c}\draw[circuit] (c)--(d);
    \draw[circuit] (b)--(1.55,1.05)--(1.55,.55);
    \draw[circuit] (c)--(1.55,-1.0)--(1.55,-.50);
    \node[nodepoint] at (a) {};\node[nodepoint] at (b) {};\node[nodepoint] at (c) {};
    \node[xeslabel,above] at (a) {$a$};\node[xeslabel,above] at (b) {$b$};\node[xeslabel,below] at (c) {$c$};
    \node[xestiny,above] at (-.77,1.28) {$R_1$};\node[xestiny,right] at (.40,.03) {$R_2$};
    \node[xestiny,left] at (-1.98,-.02) {$E_0$};
    \node[xessmall] at (0,-1.55) {只保留电压源：电流源断路};
  \end{scope}
  \begin{scope}[shift={(3.25,0)}]
    \coordinate (a) at (-1.7,1.05);\coordinate (b) at (.15,1.05);\coordinate (c) at (.15,-1.0);
    \draw[circuit] (a)--(-1.7,-1.0)--(c);
    \zigres{a}{b}\zigres{b}{c}
    \draw[circuit] (b)--(1.45,1.05)--(1.45,.36);
    \draw[circuit] (c)--(1.45,-1.0)--(1.45,-.36);
    \node[sourcecircle] at (1.45,0) {};
    \draw[currentarrow] (1.45,-.22)--(1.45,.22);
    \node[xestiny,right] at (1.73,0) {$I_0$};
    \node[nodepoint] at (a) {};\node[nodepoint] at (b) {};\node[nodepoint] at (c) {};
    \node[xeslabel,above] at (a) {$a$};\node[xeslabel,above] at (b) {$b$};\node[xeslabel,below] at (c) {$c$};
    \node[xestiny,above] at (-.77,1.28) {$R_1$};\node[xestiny,right] at (.40,.03) {$R_2$};
    \node[xessmall] at (0,-1.55) {只保留电流源：电压源短路};
  \end{scope}
\end{tikzpicture}
\end{center}}

\newcommand{\networkBlob}[1]{%
  \draw[networkblob] plot[smooth cycle,tension=.7] coordinates
    {(-1.05,.72) (-.55,1.05) (.10,.85) (.82,1.02) (1.08,.35) (.92,-.38) (.34,-.86) (-.40,-.78) (-1.04,-.25)};
  \zigres{-0.65,0.42}{0.20,0.42}
  \node[sourcecircle,minimum size=.42cm] at (-.45,-.35) {};
  \draw[currentarrow] (-.45,-.50)--(-.45,-.20);
  \draw[circuit] (.42,-.55)--(.42,-.18);\draw[circuit] (.24,-.18)--(.60,-.18);
  \draw[circuit] (.15,-.05)--(.69,-.05);\draw[circuit] (.42,-.05)--(.42,.15);
  \coordinate (#1u) at (1.72,.34);\coordinate (#1d) at (1.72,-.34);
  \draw[circuit] (1.00,.34)--(#1u);\draw[circuit] (1.00,-.34)--(#1d);
  \node[port] at (#1u) {};\node[port] at (#1d) {};
}

\newcommand{\theveninDerivationFigure}{%
\begin{center}
\begin{tikzpicture}[x=.82cm,y=.82cm]
  \begin{scope}[shift={(-6.2,0)}]
    \networkBlob{pA}
    \draw[currentarrow] (2.72,.34)--(2.05,.34);
    \node[xestiny,above] at (2.38,.43) {$i$};
    \node[xestiny,right] at (1.92,.25) {$+$};\node[xestiny,right] at (1.92,-.43) {$-$};
    \node[xestiny] at (2.28,-.12) {$v$};\node[xessmall] at (.25,-1.43) {(a) 原有源二端网络};
  \end{scope}
  \begin{scope}[shift={(-2.05,0)}]
    \networkBlob{pB}
    \node[sourcecircle] at (2.45,0) {};
    \draw[currentarrow] (2.45,-.22)--(2.45,.22);
    \draw[circuit] (pBu)--(2.45,.34)--(2.45,.27) (pBd)--(2.45,-.34)--(2.45,-.27);
    \draw[currentarrow] (2.05,.70)--(1.42,.70);
    \node[xestiny,above] at (1.74,.78) {$i_t$};
    \node[xestiny,right] at (2.78,.02) {$i_t$};\node[xestiny,below] at (2.45,-.43) {$v_t$};
    \node[xessmall] at (.45,-1.43) {(b) 接入测试电流源};
  \end{scope}
  \begin{scope}[shift={(2.45,0)}]
    \draw[networkblob] plot[smooth cycle,tension=.7] coordinates
      {(-1.05,.72) (-.4,1.0) (.28,.83) (.9,.98) (1.08,.30) (.88,-.55) (.14,-.85) (-.72,-.65) (-1.05,-.05)};
    \zigres{-0.62,0.35}{0.28,0.35}
    \draw[circuit] (-.55,-.35)--(.56,-.35);
    \draw[circuit] (1.0,.34)--(1.72,.34) (1.0,-.34)--(1.72,-.34);
    \node[port] at (1.72,.34) {};\node[port] at (1.72,-.34) {};
    \node[sourcecircle] at (2.45,0) {};\draw[currentarrow] (2.45,-.22)--(2.45,.22);
    \draw[circuit] (1.72,.34)--(2.45,.34)--(2.45,.27) (1.72,-.34)--(2.45,-.34)--(2.45,-.27);
    \node[xestiny,right] at (2.82,.02) {$i_t$};\node[xestiny,below] at (2.45,-.43) {$v_a=i_tr$};
    \node[xessmall] at (.55,-1.43) {(c) 内部独立源置零};
  \end{scope}
  \begin{scope}[shift={(6.6,0)}]
    \networkBlob{pD}
    \draw[circuit] (pDu)--++(.48,0);\draw[circuit] (pDd)--++(.48,0);
    \node[xestiny,right] at (1.92,.25) {$+$};\node[xestiny,right] at (1.92,-.43) {$-$};
    \node[xestiny,below] at (2.10,-.48) {$v_b$};\node[xestiny,above] at (2.25,.48) {$i_t=0$};
    \node[xessmall] at (.35,-1.43) {(d) 端口开路};
  \end{scope}
\end{tikzpicture}
\end{center}}

\newcommand{\sourceEquivalenceFigure}{%
\begin{center}
\begin{tikzpicture}[x=1cm,y=1cm]
  \begin{scope}[shift={(-4.65,0)}]
    \networkBlob{qA}
    \node[xessmall] at (.30,-1.28) {线性有源网络};
  \end{scope}
  \node[xeslabel] at (-1.95,0) {$\Longleftrightarrow$};
  \begin{scope}[shift={(0,0)}]
    \coordinate (u) at (-1.1,.72);\coordinate (d) at (-1.1,-.72);
    \node[sourcecircle] at (-1.1,0) {};
    \node[xestiny] at (-1.1,.12) {$+$};\node[xestiny] at (-1.1,-.12) {$-$};
    \draw[circuit] (u)--(-1.1,.27) (d)--(-1.1,-.27);
    \draw[circuit] (u)--(-.42,.72);\zigres{-0.42,.72}{.70,.72}\draw[circuit] (.70,.72)--(1.30,.72);
    \draw[circuit] (d)--(1.30,-.72);\node[port] at (1.30,.72) {};\node[port] at (1.30,-.72) {};
    \node[xestiny,below] at (-1.10,-.34) {$V_{\rm Th}$};\node[xestiny,above] at (.14,.95) {$R_{\rm Th}$};
    \node[xessmall] at (0,-1.28) {戴维南形式};
  \end{scope}
  \node[xeslabel] at (2.15,0) {$\Longleftrightarrow$};
  \begin{scope}[shift={(4.35,0)}]
    \draw[circuit] (-1.25,.72)--(1.25,.72) (-1.25,-.72)--(1.25,-.72);
    \draw[circuit] (-.85,.72)--(-.85,.27) (-.85,-.72)--(-.85,-.27);
    \node[sourcecircle] at (-.85,0) {};\draw[currentarrow] (-.85,-.22)--(-.85,.22);
    \zigres{.25,.72}{.25,-.72}
    \node[port] at (1.25,.72) {};\node[port] at (1.25,-.72) {};
    \node[xestiny,left] at (-1.15,0) {$I_{\rm N}$};\node[xestiny,right] at (.52,0) {$R_{\rm N}$};
    \node[xessmall] at (0,-1.28) {诺顿形式};
  \end{scope}
\end{tikzpicture}
\end{center}}

\newcommand{\bridgeFigure}{%
\begin{center}
\begin{tikzpicture}[x=1cm,y=1cm]
  \coordinate (L) at (-3,0);\coordinate (T) at (0,1.25);\coordinate (D) at (0,-1.25);\coordinate (R) at (3,0);
  \zigres{L}{T}\zigres{T}{R}\zigres{L}{D}\zigres{D}{R}\zigres{T}{D}
  \draw[circuit] (L)--(-3,-2.0)--(-.55,-2.0);
  \draw[circuit] (.55,-2.0)--(3,-2.0)--(R);
  \draw[circuit] (-.55,-2.0)--(-.18,-2.0);\draw[circuit] (.18,-2.0)--(.55,-2.0);
  \draw[circuit] (-.18,-1.72)--(-.18,-2.28);\draw[circuit] (.18,-1.83)--(.18,-2.17);
  \foreach \P in {L,T,D,R}{\node[nodepoint] at (\P) {};}
  \node[xestiny,above left] at (-1.55,.70) {$R_1$};\node[xestiny,above right] at (1.55,.70) {$R_2$};
  \node[xestiny,below left] at (-1.55,-.70) {$R_3$};\node[xestiny,below right] at (1.55,-.70) {$R_4$};
  \node[xestiny,right] at (.25,0) {$R_5$};\node[xestiny,below] at (0,-2.28) {$\mathcal E$};
  \draw[currentarrow] (.43,.54)--(.43,-.54);\node[xestiny,right] at (.54,.02) {$I_5=0$};
\end{tikzpicture}
\end{center}}

\newcommand{\deltaStarFigure}{%
\begin{center}
\begin{tikzpicture}[x=1cm,y=1cm]
  \begin{scope}[shift={(-3.4,0)}]
    \coordinate (A) at (0,1.75);\coordinate (B) at (-1.75,-1.25);\coordinate (C) at (1.75,-1.25);
    \zigres{A}{B}\zigres{A}{C}\zigres{B}{C}
    \foreach \P in {A,B,C}{\node[port] at (\P) {};}
    \node[xeslabel,above] at (A) {$A$};\node[xeslabel,below left] at (B) {$B$};\node[xeslabel,below right] at (C) {$C$};
    \node[xestiny,left] at (-1.05,.38) {$R_{AB}$};\node[xestiny,right] at (1.05,.38) {$R_{AC}$};
    \node[xestiny,below] at (0,-1.48) {$R_{BC}$};\node[xessmall] at (0,-2.05) {$\Delta$ 连接};
  \end{scope}
  \node[xeslabel] at (0,0) {$\Longleftrightarrow$};
  \begin{scope}[shift={(3.4,0)}]
    \coordinate (A) at (0,1.75);\coordinate (B) at (-1.75,-1.25);\coordinate (C) at (1.75,-1.25);\coordinate (O) at (0,0);
    \zigres{A}{O}\zigres{B}{O}\zigres{C}{O}
    \foreach \P in {A,B,C}{\node[port] at (\P) {};}
    \node[nodepoint] at (O) {};
    \node[xeslabel,above] at (A) {$A$};\node[xeslabel,below left] at (B) {$B$};\node[xeslabel,below right] at (C) {$C$};
    \node[xestiny,right] at (.28,.93) {$R_a$};\node[xestiny,above left] at (-.95,-.50) {$R_b$};
    \node[xestiny,above right] at (.95,-.50) {$R_c$};\node[xestiny,above right] at (O) {$\varphi_0$};
    \node[xessmall] at (0,-2.05) {$Y$ 连接};
  \end{scope}
\end{tikzpicture}
\end{center}}

\newcommand{\sourceTransformExampleFigure}{%
\begin{center}
\begin{tikzpicture}[x=1cm,y=1cm]
  \begin{scope}[shift={(-3.2,0)}]
    \node[sourcecircle] at (-1.55,0) {};
    \node[xestiny] at (-1.55,.12) {$+$};\node[xestiny] at (-1.55,-.12) {$-$};
    \draw[circuit] (-1.55,.27)--(-1.55,.90)--(-.85,.90);
    \zigres{-0.85,.90}{.35,.90}\draw[circuit] (.35,.90)--(1.40,.90);
    \draw[circuit] (-1.55,-.27)--(-1.55,-.90)--(1.40,-.90);
    \node[port] at (1.40,.90) {};\node[port] at (1.40,-.90) {};
    \draw[currentarrow] (.92,1.24)--(1.38,1.24);
    \node[xestiny,above] at (1.15,1.28) {$I$};
    \node[xestiny,left] at (-1.88,0) {$E$};\node[xestiny,above] at (-.25,1.12) {$r$};
    \node[xestiny,right] at (1.60,.68) {$+$};\node[xestiny,right] at (1.60,-1.02) {$-$};
    \node[xestiny] at (1.85,0) {$U$};\node[xessmall] at (0,-1.45) {电压源形式};
  \end{scope}
  \node[xeslabel] at (0,0) {$\Longleftrightarrow$};
  \begin{scope}[shift={(3.2,0)}]
    \draw[circuit] (-1.45,.90)--(1.40,.90) (-1.45,-.90)--(1.40,-.90);
    \draw[circuit] (-.85,.90)--(-.85,.27) (-.85,-.90)--(-.85,-.27);
    \node[sourcecircle] at (-.85,0) {};\draw[currentarrow] (-.85,-.22)--(-.85,.22);
    \zigres{.20,.90}{.20,-.90}
    \node[port] at (1.40,.90) {};\node[port] at (1.40,-.90) {};
    \draw[currentarrow] (.92,1.24)--(1.38,1.24);
    \node[xestiny,above] at (1.15,1.28) {$I$};
    \node[xestiny,left] at (-1.17,0) {$E/r$};\node[xestiny,right] at (.47,0) {$r$};
    \node[xestiny,right] at (1.60,.68) {$+$};\node[xestiny,right] at (1.60,-1.02) {$-$};
    \node[xestiny] at (1.85,0) {$U$};\node[xessmall] at (0,-1.45) {电流源形式};
  \end{scope}
\end{tikzpicture}
\end{center}}

\newcommand{\gridProblemFigure}{%
\begin{center}
\begin{tikzpicture}[x=1.05cm,y=1.05cm]
  \foreach \i in {0,1,2}{\draw[circuit] (0,\i)--(2,\i);\draw[circuit] (\i,0)--(\i,2);}
  \foreach \i in {0,1,2}{\foreach \j in {0,1,2}{\node[nodepoint] at (\i,\j) {};}}
  \node[xeslabel,below left] at (0,0) {$A$};\node[xeslabel,below right] at (2,0) {$B$};
  \node[xeslabel,above right] at (2,2) {$C$};
  \draw[<->,DeepBlue!75,line width=.6pt] (0,-.55)--(1,-.55);
  \node[xestiny,below] at (.5,-.60) {$R$};
\end{tikzpicture}
\end{center}}

\newcommand{\potentialNode}[3]{%
  \node[circle,fill=#3,draw=white,line width=.3pt,inner sep=1.55pt] at (#1) {};
  \node[xestiny,text=Ink,anchor=south west] at ($(#1)+(.08,.08)$) {$#2$};}

\newcommand{\gridPotentialDiagonalFigure}{%
\begin{center}
\begin{tikzpicture}[x=1.2cm,y=1.2cm]
  \foreach \i in {0,1,2}{\draw[circuit] (0,\i)--(2,\i);\draw[circuit] (\i,0)--(\i,2);}
  \potentialNode{0,0}{1}{Gold!70}
  \potentialNode{1,0}{\frac23}{DeepBlue!28}
  \potentialNode{2,0}{\frac12}{RuleGray!55}
  \potentialNode{0,1}{\frac23}{DeepBlue!28}
  \potentialNode{1,1}{\frac12}{RuleGray!55}
  \potentialNode{2,1}{\frac13}{DeepBlue!12}
  \potentialNode{0,2}{\frac12}{RuleGray!55}
  \potentialNode{1,2}{\frac13}{DeepBlue!12}
  \potentialNode{2,2}{0}{Ink!20}
  \node[xestiny,below left] at (0,0) {$A$};\node[xestiny,below right=3pt] at (2,2) {$C$};
\end{tikzpicture}
\end{center}}

\newcommand{\gridPotentialAdjacentFigure}{%
\begin{center}
\begin{tikzpicture}[x=1.2cm,y=1.2cm]
  \foreach \i in {0,1,2}{\draw[circuit] (0,\i)--(2,\i);\draw[circuit] (\i,0)--(\i,2);}
  \potentialNode{0,0}{1}{Gold!70}
  \potentialNode{1,0}{\frac12}{RuleGray!55}
  \potentialNode{2,0}{0}{Ink!20}
  \potentialNode{0,1}{\frac7{10}}{DeepBlue!30}
  \potentialNode{1,1}{\frac12}{RuleGray!55}
  \potentialNode{2,1}{\frac3{10}}{DeepBlue!12}
  \potentialNode{0,2}{\frac35}{DeepBlue!23}
  \potentialNode{1,2}{\frac12}{RuleGray!55}
  \potentialNode{2,2}{\frac25}{DeepBlue!16}
  \node[xestiny,below left] at (0,0) {$A$};\node[xestiny,below right] at (2,0) {$B$};
\end{tikzpicture}
\end{center}}

\newcommand{\tetrahedronProblemFigure}{%
\begin{center}
\begin{tikzpicture}[x=1cm,y=1cm]
  \coordinate (T) at (0,2.0);\coordinate (L) at (-2.0,.15);\coordinate (R) at (2.0,.15);\coordinate (B) at (0,-1.65);
  \draw[circuit] (T)--(L)--(B)--(R)--(T) (L)--(R) (T)--(B);
  \foreach \P in {T,L,R,B}{\node[bluenode] at (\P) {};}
  \node[xestiny,above] at (T) {$A$};\node[xestiny,left] at (L) {$C$};
  \node[xestiny,right] at (R) {$D$};\node[xestiny,below] at (B) {$B$};
  \node[xessmall] at (0,-2.15) {交叉点只是投影重叠，不是电路节点};
\end{tikzpicture}
\end{center}}

\newcommand{\completeGraphProblemFigure}{%
\begin{center}
\begin{tikzpicture}[x=1cm,y=1cm]
  \foreach \i in {0,...,4}{\coordinate (P\i) at ({90+72*\i}:1.75);}
  \foreach \i in {0,...,4}{\foreach \j in {0,...,4}{\ifnum\i<\j\draw[circuit] (P\i)--(P\j);\fi}}
  \foreach \i in {0,...,4}{\node[bluenode] at (P\i) {};}
  \node[xestiny,above] at (P0) {$N=5$};
\end{tikzpicture}
\end{center}}

\newcommand{\completeGraphReductionFigure}{%
\begin{center}
\begin{tikzpicture}[x=1cm,y=1cm]
  \coordinate (A) at (-3,0);\coordinate (B) at (3,0);
  \node[goldnode] at (A) {};\node[goldnode] at (B) {};
  \node[xeslabel,left] at (A) {$A:V$};\node[xeslabel,right] at (B) {$B:0$};
  \draw[activewire] (A)--(B);
  \foreach \y in {1.25,.42,-.42,-1.25}{
    \coordinate (M) at (0,\y);\draw[solutionwire] (A)--(M)--(B);\node[bluenode] at (M) {};}
  \draw[RuleGray,densely dashed,rounded corners=3pt] (-.45,1.60) rectangle (.45,-1.60);
  \node[xessmall,anchor=west] at (.58,1.42) {$N-2$ 个等势点：$V/2$};
  \node[xestiny] at (-1.5,1.55) {$R$};\node[xestiny] at (1.5,1.55) {$R$};
  \node[xessmall] at (0,-1.95) {每条支路段均为 $R$；包括直接连接 $A,B$ 的一段};
\end{tikzpicture}
\end{center}}

\newcommand{\octagonProblemFigure}{%
\begin{center}
\begin{tikzpicture}[x=1cm,y=1cm]
  \foreach \i in {0,...,7}{\coordinate (O\i) at ({22.5+45*\i}:2.0);}
  \foreach \i in {0,...,7}{
    \foreach \j in {0,...,7}{
      \ifnum\i<\j
        \pgfmathtruncatemacro{\d}{mod(\j-\i,8)}
        \ifnum\d=4\else\draw[circuit,line width=.55pt] (O\i)--(O\j);\fi
      \fi}}
  \foreach \i in {0,...,7}{\node[bluenode] at (O\i) {};}
  \node[xestiny,right] at (O0) {$A$};\node[xestiny,left] at (O4) {$B$};
  \draw[DeepBlue!70,densely dashed,line width=.7pt] (O0)--(O4);
  \node[xessmall,below] at (0,-2.38) {蓝色虚线直径不接电阻；其余任意顶点对均有 $R$};
\end{tikzpicture}
\end{center}}

\newcommand{\octagonReductionFigure}{%
\begin{center}
\begin{tikzpicture}[x=1cm,y=1cm]
  \coordinate (A) at (-3.2,0);\coordinate (B) at (3.2,0);
  \node[goldnode] at (A) {};\node[goldnode] at (B) {};
  \node[xeslabel,left] at (A) {$A$};\node[xeslabel,right] at (B) {$B$};
  \foreach \y in {1.55,.95,.32,-.32,-.95,-1.55}{
    \coordinate (M) at (0,\y);\draw[solutionwire] (A)--(M)--(B);\node[bluenode] at (M) {};}
  \node[xestiny] at (-1.62,1.75) {$R$};\node[xestiny] at (1.62,1.75) {$R$};
  \draw[RuleGray,densely dashed,rounded corners=3pt] (-.45,1.82) rectangle (.45,-1.82);
  \node[xessmall] at (0,-2.12) {六个中间点等势};
  \draw[DeepBlue!70,densely dashed] (A)--(B);
  \node[xessmall] at (0,-2.62) {虚线只表示直径位置，$A,B$ 间无直接支路};
\end{tikzpicture}
\end{center}}

\newcommand{\cubeCoordinates}{%
  \coordinate (A) at (-1.35,-1.05);\coordinate (B) at (1.05,-1.05);
  \coordinate (E) at (-1.35,1.10);\coordinate (F) at (1.05,1.10);
  \coordinate (D) at (-.45,-.32);\coordinate (C) at (1.95,-.32);
  \coordinate (H) at (-.45,1.83);\coordinate (G) at (1.95,1.83);}

\newcommand{\cubeEdges}{%
  \draw[circuit] (A)--(B)--(F)--(E)--cycle;
  \draw[circuit] (B)--(C)--(G)--(F) (E)--(H)--(G);
  \draw[hiddenwire] (A)--(D)--(C) (D)--(H);}

\newcommand{\cubeProblemFigure}{%
\begin{center}
\begin{tikzpicture}[x=1cm,y=1cm]
  \cubeCoordinates\cubeEdges
  \foreach \P in {A,B,C,D,E,F,G,H}{\node[bluenode] at (\P) {};}
  \node[xeslabel,below left] at (A) {$A$};\node[xeslabel,below right] at (B) {$B$};
  \node[xeslabel,right] at (C) {$C$};\node[xeslabel,above right] at (G) {$C_1$};
  \node[xessmall] at (.30,-1.62) {虚线棱在立方体后方，所有投影交叉处均不另设节点};
\end{tikzpicture}
\end{center}}

\newcommand{\cubePotentialFigure}{%
\begin{center}
\begin{tikzpicture}[x=.69cm,y=.69cm]
  \begin{scope}[shift={(-6.0,0)}]
    \cubeCoordinates\cubeEdges
    \potentialNode{A}{1}{Gold!72}\potentialNode{B}{0}{Ink!20}
    \potentialNode{E}{\frac9{14}}{DeepBlue!28}\potentialNode{D}{\frac9{14}}{DeepBlue!28}
    \potentialNode{H}{\frac47}{DeepBlue!23}\potentialNode{F}{\frac5{14}}{DeepBlue!14}
    \potentialNode{C}{\frac5{14}}{DeepBlue!14}\potentialNode{G}{\frac37}{DeepBlue!18}
    \node[xestiny] at (.30,-1.70) {(1) 邻点 $A\to B$};
  \end{scope}
  \begin{scope}[shift={(0,0)}]
    \cubeCoordinates\cubeEdges
    \potentialNode{A}{1}{Gold!72}\potentialNode{C}{0}{Ink!20}
    \potentialNode{E}{\frac23}{DeepBlue!30}\potentialNode{D}{\frac12}{RuleGray!55}
    \potentialNode{H}{\frac12}{RuleGray!55}\potentialNode{B}{\frac12}{RuleGray!55}
    \potentialNode{F}{\frac12}{RuleGray!55}\potentialNode{G}{\frac13}{DeepBlue!13}
    \node[xestiny] at (.30,-1.70) {(2) 面对角 $A\to C$};
  \end{scope}
  \begin{scope}[shift={(6.0,0)}]
    \cubeCoordinates\cubeEdges
    \potentialNode{A}{1}{Gold!72}\potentialNode{G}{0}{Ink!20}
    \potentialNode{B}{\frac35}{DeepBlue!27}\potentialNode{D}{\frac35}{DeepBlue!27}
    \potentialNode{E}{\frac35}{DeepBlue!27}\potentialNode{C}{\frac25}{DeepBlue!15}
    \potentialNode{F}{\frac25}{DeepBlue!15}\potentialNode{H}{\frac25}{DeepBlue!15}
    \node[xestiny] at (.30,-1.70) {(3) 体对角 $A\to G$};
  \end{scope}
\end{tikzpicture}
\end{center}}

\newcommand{\infiniteLadderProblemFigure}{%
\begin{center}
\begin{tikzpicture}[x=1cm,y=1cm]
  \coordinate (a) at (-3.7,.75);\coordinate (b) at (-3.7,-.75);
  \node[port] at (a) {};\node[port] at (b) {};
  \node[xeslabel,left] at (a) {$a$};\node[xeslabel,left] at (b) {$b$};
  \draw[circuit] (a)--(-3.45,.75) (b)--(-3.45,-.75);
  \foreach \i in {0,1,2,3}{
    \pgfmathsetmacro{\xa}{-3.45+1.65*\i}\pgfmathsetmacro{\xb}{-2.05+1.65*\i}
    \zigres{\xa,.75}{\xb,.75}\zigres{\xa,-.75}{\xb,-.75}
    \zigres{\xb,.75}{\xb,-.75}
    \ifnum\i<3\draw[circuit] (\xb,.75)--({\xb+.25},.75) (\xb,-.75)--({\xb+.25},-.75);\fi}
  \node[xestiny,above] at (-2.75,1.02) {$R$};
  \draw[circuit,densely dotted] (2.90,.75)--(4.10,.75) (2.90,-.75)--(4.10,-.75);
  \node[xeslabel,right] at (4.10,0) {$\cdots$};
\end{tikzpicture}
\end{center}}

\newcommand{\ladderRecursionFigure}{%
\begin{center}
\begin{tikzpicture}[x=1cm,y=1cm]
  \coordinate (a) at (-3.4,.85);\coordinate (b) at (-3.4,-.85);
  \node[port] at (a) {};\node[port] at (b) {};\node[xeslabel,left] at (a) {$a$};\node[xeslabel,left] at (b) {$b$};
  \draw[circuit] (a)--(-3.15,.85) (b)--(-3.15,-.85);
  \zigres{-3.15,.85}{-1.65,.85}\zigres{-3.15,-.85}{-1.65,-.85}
  \node[xestiny,above] at (-2.40,1.12) {$R$};\node[xestiny,below] at (-2.40,-1.12) {$R$};
  \coordinate (u) at (-1.45,.85);\coordinate (d) at (-1.45,-.85);
  \draw[circuit] (-1.65,.85)--(u) (-1.65,-.85)--(d);
  \zigres{u}{d}\node[xestiny,left] at (-1.73,0) {$R$};
  \draw[circuit] (u)--(-.65,.85) (d)--(-.65,-.85);
  \draw[DeepBlue!70,dashed,rounded corners=3pt,line width=.8pt] (-.65,1.25) rectangle (2.65,-1.25);
  \node[xeslabel] at (1.0,0) {无限尾部 $\equiv X$};
  \draw[circuit] (2.65,.85)--(3.15,.85) (2.65,-.85)--(3.15,-.85);
  \node[port] at (3.15,.85) {};\node[port] at (3.15,-.85) {};
  \node[xessmall] at (0,-1.72) {$X=R+\bigl(R\parallel X\bigr)+R$};
\end{tikzpicture}
\end{center}}

\newcommand{\threeRailProblemFigure}{%
\begin{center}
\begin{tikzpicture}[x=1cm,y=1cm]
  \foreach \r/\lab in {1.35/A,0/B,-1.35/C}{
    \node[port] at (-4.05,\r) {};\node[xeslabel,left] at (-4.05,\r) {$\lab$};
    \draw[circuit] (-4.05,\r)--(-3.82,\r);
    \foreach \i in {0,1,2,3}{
      \pgfmathsetmacro{\xa}{-3.82+1.65*\i}\pgfmathsetmacro{\xb}{-2.48+1.65*\i}
      \zigres{\xa,\r}{\xb,\r}
      \ifnum\i<3\draw[circuit] (\xb,\r)--({\xb+.31},\r);\fi}}
  \foreach \i in {0,1,2,3}{
    \pgfmathsetmacro{\x}{-2.48+1.65*\i}
    \zigres{\x,1.35}{\x,0}\zigres{\x,0}{\x,-1.35}}
  \draw[circuit,densely dotted] (2.47,1.35)--(3.55,1.35) (2.47,0)--(3.55,0) (2.47,-1.35)--(3.55,-1.35);
  \node[xeslabel,right] at (3.55,0) {$\cdots$};
  \node[xestiny,above] at (-3.15,1.65) {$R$};\node[xestiny,right] at (-2.20,.68) {$R$};
\end{tikzpicture}
\end{center}}

\newcommand{\threeRailModesFigure}{%
\begin{center}
\begin{tikzpicture}[x=1cm,y=1cm]
  \begin{scope}[shift={(-3.1,0)}]
    \foreach \y in {1,0,-1}{\node[bluenode] at (0,\y) {};}
    \draw[circuit] (0,1)--(0,0)--(0,-1);
    \draw[currentarrow] (.35,0)--(.35,1.0);\draw[currentarrow] (.35,0)--(.35,-1.0);
    \node[xeslabel,right] at (.58,1) {$+1$};\node[xeslabel,right] at (.58,0) {$0$};\node[xeslabel,right] at (.58,-1) {$-1$};
    \node[xessmall] at (0,-1.52) {$\boldsymbol e_1=(1,0,-1)$，$\mu=1$};
  \end{scope}
  \begin{scope}[shift={(3.1,0)}]
    \foreach \y in {1,0,-1}{\node[bluenode] at (0,\y) {};}
    \draw[circuit] (0,1)--(0,0)--(0,-1);
    \draw[currentarrow] (.35,0)--(.35,.78);\draw[currentarrow] (.35,.0)--(.35,-.78);
    \draw[currentarrow] (-.35,.85)--(-.35,-.85);
    \node[xeslabel,right] at (.58,1) {$+1$};\node[xeslabel,right] at (.58,0) {$-2$};\node[xeslabel,right] at (.58,-1) {$+1$};
    \node[xessmall] at (0,-1.52) {$\boldsymbol e_3=(1,-2,1)$，$\mu=3$};
  \end{scope}
\end{tikzpicture}
\end{center}}

\newcommand{\nestedTriangleProblemFigure}{%
\begin{center}
\begin{tikzpicture}[x=1cm,y=1cm]
  \coordinate (O) at (0,0);
  \foreach \k in {0,...,6}{
    \pgfmathsetmacro{\rad}{2.55/(2^\k)}
    \pgfmathsetmacro{\ang}{90+180*mod(\k,2)}
    \draw[circuit,line width=.72pt]
      ($(O)+(\ang:\rad)$)--($(O)+({\ang+120}:\rad)$)--($(O)+({\ang+240}:\rad)$)--cycle;}
  \node[xeslabel,below left] at ($(O)+(210:2.55)$) {$A$};
  \node[xeslabel,below right] at ($(O)+(330:2.55)$) {$B$};
  \node[xeslabel,above] at ($(O)+(90:2.55)$) {$C$};
\end{tikzpicture}
\end{center}}

\newcommand{\nestedTriangleRecursionFigure}{%
\begin{center}
\begin{tikzpicture}[x=1cm,y=1cm]
  \coordinate (A) at (-2.55,-1.45);\coordinate (B) at (2.55,-1.45);\coordinate (C) at (0,2.97);
  \coordinate (Mab) at (0,-1.45);\coordinate (Mbc) at (1.275,.76);\coordinate (Mca) at (-1.275,.76);
  \draw[solutionwire] (A)--(Mab)--(B) (B)--(Mbc)--(C) (C)--(Mca)--(A);
  \draw[activewire] (Mab)--(Mbc)--(Mca)--cycle;
  \foreach \P in {A,B,C}{\node[goldnode] at (\P) {};}
  \foreach \P in {Mab,Mbc,Mca}{\node[bluenode] at (\P) {};}
  \node[xeslabel,below left] at (A) {$A$};\node[xeslabel,below right] at (B) {$B$};\node[xeslabel,above] at (C) {$C$};
  \node[xestiny,below] at (-1.28,-1.55) {$r/2$};\node[xestiny,below] at (1.28,-1.55) {$r/2$};
  \node[xestiny] at (0,.06) {内网等效支路 $kr/2$};
  \node[xessmall] at (0,-2.05) {蓝色三角形不是新增金属丝，而是对全部后续嵌套网络的三端等效};
\end{tikzpicture}
\end{center}}

\newcommand{\sierpinskiProblemFigure}{%
\begin{center}
\begin{tikzpicture}[x=1cm,y=1cm]
  \begin{scope}[shift={(-5.2,0)}]
    \fill[DeepBlue!48] (-1.45,-1.0)--(1.45,-1.0)--(0,1.51)--cycle;
    \draw[circuit] (-1.45,-1.0)--(1.45,-1.0)--(0,1.51)--cycle;
    \node[xessmall] at (0,-1.48) {初始};
  \end{scope}
  \begin{scope}[shift={(0,0)}]
    \foreach \x/\y in {-0.725/-1.0,0.725/-1.0,0/.255}{
      \fill[DeepBlue!48] (\x-.725,\y)--(\x+.725,\y)--(\x,{\y+1.255})--cycle;
      \draw[circuit] (\x-.725,\y)--(\x+.725,\y)--(\x,{\y+1.255})--cycle;}
    \node[xessmall] at (0,-1.48) {第一阶};
  \end{scope}
  \begin{scope}[shift={(5.2,0)}]
    \foreach \X/\Y in {-1.0875/-1.0,-.3625/-1.0,.3625/-1.0,1.0875/-1.0,-.725/-.3725,.725/-.3725,-.3625/.255,.3625/.255,0/.8825}{
      \fill[DeepBlue!48] (\X-.3625,\Y)--(\X+.3625,\Y)--(\X,{\Y+.6275})--cycle;
      \draw[circuit,line width=.6pt] (\X-.3625,\Y)--(\X+.3625,\Y)--(\X,{\Y+.6275})--cycle;}
    \node[xessmall] at (0,-1.48) {第二阶};
  \end{scope}
  \node[xestiny] at (-6.75,-1.10) {$A$};\node[xestiny] at (-3.65,-1.10) {$B$};\node[xestiny] at (-5.20,1.68) {$C$};
\end{tikzpicture}
\end{center}}

\newcommand{\sierpinskiResistanceFigure}{%
\begin{center}
\begin{tikzpicture}[x=1cm,y=1cm]
  \coordinate (A) at (-3.0,-1.25);\coordinate (B) at (3.0,-1.25);\coordinate (C) at (0,3.95);
  \coordinate (Mab) at (0,-1.25);\coordinate (Mbc) at (1.50,1.35);\coordinate (Mca) at (-1.50,1.35);
  \foreach \P/\Q/\R in {A/Mab/Mca,Mab/B/Mbc,Mca/Mbc/C}{
    \draw[solutionwire,fill=DeepBlue!5] (\P)--(\Q)--(\R)--cycle;}
  \foreach \P in {A,B,C,Mab,Mbc,Mca}{\node[bluenode] at (\P) {};}
  \node[xeslabel,below left] at (A) {$A$};\node[xeslabel,below right] at (B) {$B$};\node[xeslabel,above] at (C) {$C$};
  \node[xessmall] at (-1.5,-1.95) {缩小的同型三端网};
  \draw[->,TextGray,line width=.6pt] (-1.5,-1.70)--(-1.5,-.55);
  \node[eqbox,anchor=west] at (3.65,1.35) {$\displaystyle
    R_{n+1}=\frac53R_n\quad(\text{单段电阻不缩放})$};
  \node[eqbox,anchor=west] at (3.65,.20) {$\displaystyle
    r_{n+1}=\frac12r_n\quad(\text{真实长度缩放})$};
  \node[eqbox,anchor=west] at (3.65,-.95) {$\displaystyle
    R_{n+1}^{\rm phys}=\frac56R_n^{\rm phys}$};
\end{tikzpicture}
\end{center}}


原则上，基尔霍夫定律足以处理绝大多数集中参数电路；但网络一旦变复杂，直接设置每条支路电流、
列出全部节点和回路方程，会迅速产生规模很大的线性方程组。这样的做法可靠，却冗长、难以检查，也很
容易在符号上出错。因此，对只含线性元件的网络，应优先利用线性叠加、等效电源、对称性、电桥平衡、
星形--三角形变换和自相似性，把网络先化简，再在必要时使用基尔霍夫方程。对于计算机，求解大规模
线性方程组并不困难；对人工竞赛解题而言，找到结构才是关键。


\subsection{叠加原理}


节点电压法、支路电流法和基尔霍夫方程不仅适用于电阻，也可与任意元件的伏安关系联立。若元件是
非线性的，得到的通常是非线性方程组，叠加原理一般不成立；若网络只含线性电阻和独立电源，方程组对
电源产生的非齐次项具有可加性。于是多个独立电源同时工作时，任一支路电流或任一节点电势，等于各个
电源单独作用时对应结果的代数和。

考察某一个独立电源的贡献时，其余独立电源要“置零”，同时保留它们的内阻和原有连接关系：

\begin{itemize}
  \item 理想电压源置零，即把其电压改为零，等效为短路；
  \item 理想电流源置零，即把其电流改为零，等效为断路。
\end{itemize}

下图把本章前面含电压源 $E_0$、电流源 $I_0$ 和两个电阻的电路分成两个子问题。仅保留电压源时，
$R_1,R_2$ 串联，流过 $R_2$ 的电流为
\[
  I_1=\frac{E_0}{R_1+R_2}.
\]
仅保留电流源时，$R_1,R_2$ 并联，按分流关系，流过 $R_2$ 的电流为
\[
  I_2=\frac{R_1}{R_1+R_2}I_0.
\]
若图示两部分在 $R_2$ 中产生的电流方向相同，则原电路中
$I_{R_2}=I_1+I_2$；若方向相反，就应按所选正方向作代数相加。

\par\noindent\begin{minipage}{\linewidth}
\superpositionFigure
\diagramcaption{叠加原理。左图只保留电压源，右图只保留电流源；被置零的电压源短路，被置零的电流源断路。}\label{ch21:fig:s37-6-1}
\end{minipage}\par\medskip

叠加的是电压和电流，不是功率。由于 $P=I^2R$ 或 $P=U^2/R$ 是二次量，通常不能把各电源单独作用时
的功率直接相加。


\subsection{等效电源定理}


线性叠加还有两个重要推论：戴维南（Thevenin）等效电压源定理和诺顿（Norton）等效电流源定理。
无论内部结构多复杂，一个线性有源二端网络对外只表现为端电压与端电流之间的一条直线关系；只要确定
这条直线的截距和斜率，就不必再追踪网络内部的每条支路。

对图~\ref{ch21:fig:s37-6-2}(a) 的二端网络，规定测试电流 $i_t$ 流入上端口，端电压为 $v_t$。在端口接入可调理想电流源，
得到图~\ref{ch21:fig:s37-6-2}(b)。由叠加原理，$v_t$ 可拆成两部分：

\begin{enumerate}
  \item 保留测试电流源，把网络内部所有独立源置零。此时内部成为纯电阻网络，端口等效电阻记为 $r$，
  因而图~\ref{ch21:fig:s37-6-2}(c) 中 $v_a=i_t r$；
  \item 断开测试电流源，恢复网络内部全部独立源。端口电流为零，图~\ref{ch21:fig:s37-6-2}(d) 中的端电压 $v_b$ 就是
  开路电压。
\end{enumerate}

所以
\[
  v_t=v_a+v_b=i_t r+v_b.
\]
这正是一个理想电压源与内阻串联时的端口关系。

\par\noindent\begin{minipage}{\linewidth}
\theveninDerivationFigure
\diagramcaption{戴维南定理的推导。四幅图的端口方向完全一致，第三幅只保留测试源，第四幅只保留网络内部独立源。}\label{ch21:fig:s37-6-2}
\end{minipage}\par\medskip

具体地说，任意线性有源二端网络都可等效为开路电压 $V_{\rm Th}$ 与电阻 $R_{\rm Th}$ 串联：
$V_{\rm Th}$ 由原网络端口开路时的电压得到；$R_{\rm Th}$ 由内部所有独立源置零后、从端口向里看见的
等效电阻得到。若网络中还含受控源，受控源不能置零，应在端口加测试源求 $R_{\rm Th}=v_t/i_t$。

同一网络也可等效为短路电流 $I_{\rm N}$ 与电阻 $R_{\rm N}$ 并联。$I_{\rm N}$ 是端口短路时的电流，
$R_{\rm N}$ 与 $R_{\rm Th}$ 相同。两种等效电路在端口上满足
\begin{equation}
  \boxed{R_{\rm Th}=R_{\rm N},\qquad V_{\rm Th}=I_{\rm N}R_{\rm N}.}
\end{equation}

\par\noindent\begin{minipage}{\linewidth}
\sourceEquivalenceFigure
\diagramcaption{同一个线性有源二端网络的戴维南形式与诺顿形式。等效只保证端口外部特性相同。}\label{ch21:fig:s37-6-3}
\end{minipage}\par\medskip

“等效”并不意味着三种电路内部的电流和功率分布相同；它只表示接到同一个负载后，端口电压和负载电流
完全一致。


% EM teaching enrichment 20260927
\paragraph{负载功率为何在阻值相等时最大}
当等效内阻 $R_{\rm Th}>0$、外接纯电阻 $R_L>0$ 时，
\[
 P_L=\frac{V_{\rm Th}^2R_L}{(R_{\rm Th}+R_L)^2},\qquad
 \frac{\dd P_L}{\dd R_L}
 =\frac{V_{\rm Th}^2(R_{\rm Th}-R_L)}{(R_{\rm Th}+R_L)^3}.
\]
故 $R_L=R_{\rm Th}$ 时最大，$P_{L,\max}=V_{\rm Th}^2/(4R_{\rm Th})$。这里最大化的是负载功率，不是输电效率；在这个简单串联等效模型中，匹配时负载与内阻耗散相等，效率仅为一半。端口等效也不能据此还原复杂原网络内部每个元件的耗散。


\subsection{对称性}


若交换网络中的两个节点后，所有连接、元件参数以及外加电源条件都保持不变，那么这两个节点不可区分，
电势必然相等。等势节点之间没有电流，可以按需要用理想导线相连，也可以把连接它们的支路删去。使用
对称性时，必须把“电路本身的对称”与“端口激励的对称”同时检查；只看几何外形相似而忽略电流从哪里
流入、从哪里流出，是复杂电阻网络中最常见的错误之一。


% EM teaching enrichment 20260927
\paragraph{等势结论来自同一个边值问题}
把端口 $A$ 固定为 $U$、$B$ 固定为零。若一个变换保持两个端口各自不动，并把等阻支路一一对应，变换前后的节点电势都满足同一组 KCL；由电阻网络解的唯一性，相互交换的节点必等势。
若变换交换 $A,B$，一般得到的是 $V_{i'}=U-V_i$，而非 $V_{i'}=V_i$。只有处在该变换不动位置的节点才能推出 $V=U/2$。
例如所有棱阻值为 $R$ 的正方体，从一对体对角顶点通电，靠近入端的三个顶点等势，靠近出端的三个也等势，三层等效为
\[
 R_{AB}=\frac R3+\frac R6+\frac R3=\frac{5R}{6}.
\]
若把出端换到相邻顶点，原先的三点分组不再保持端口条件，不能照搬上述结果。


\subsection{电桥平衡}


图~\ref{ch21:fig:s37-6-4} 中，电源跨接在左右两端。若桥臂满足
\[
  \boxed{R_1R_4=R_2R_3},
\]
则上、下两个中点电势相等，$R_5$ 中电流为零，这就是电桥平衡。此时 $R_5$ 既可看作没有电流的断路支路
而删去，也可利用两端等势把它短接；两种处理对其余电路的电流、电压都没有影响。要注意，“断路”说的是
$R_5$ 中无电流，“短路”说的是它的两个端点等势，二者描述的是同一平衡状态的不同侧面。

\par\noindent\begin{minipage}{\linewidth}
\bridgeFigure
\diagramcaption{平衡电桥。$R_5$ 的上下端分别是两个分压支路的中点；平衡条件下二者等势。}\label{ch21:fig:s37-6-4}
\end{minipage}\par\medskip


% EM teaching enrichment 20260927
\paragraph{先断开桥支路，比较两个分压}
设图中上支路由左至右为 $R_1,R_2$，下支路为 $R_3,R_4$，左右端电势分别为 $U,0$。暂断开 $R_5$，两中点电势为
\[
 V_{\rm up}=U\frac{R_2}{R_1+R_2},\qquad
 V_{\rm down}=U\frac{R_4}{R_3+R_4}.
\]
相等即给出 $R_1R_4=R_2R_3$。接回任意有限 $R_5$ 不产生桥电流，因而原分压保持成立；若不相等，接回后两个分压都会受加载影响。
还可把桥端口外的网络化成戴维南源，其内阻为
$(R_1\parallel R_2)+(R_3\parallel R_4)$，桥电流便等于开路中点电压差除以该内阻与 $R_5$ 之和。这给出了偏离平衡时可继续计算的路线。


% EM structure and teaching revision 20260928
\subsection{端口等效的测定与适用边界}
线性有源二端网络可在端口写成 $U=U_{\mathrm{oc}}-IR_{\mathrm{th}}$，其中 $I$ 取网络向负载输出的方向。开路电压确定截距，短路电流在存在且有限时给出 $R_{\mathrm{th}}=U_{\mathrm{oc}}/I_{\mathrm{sc}}$。
只有独立源的电路，求等效电阻时才可把理想独立电压源短接、理想独立电流源断开；实际源内阻仍应保留。含受控源时必须保留控制关系，并在端口加测试电压或测试电流来求比值。
对 $R_{\mathrm{th}}>0$ 的被动负载 $R_L$，
\[
 P_L=\frac{U_{\mathrm{oc}}^2R_L}{(R_{\mathrm{th}}+R_L)^2},\qquad
 \frac{\dd P_L}{\dd R_L}=\frac{U_{\mathrm{oc}}^2(R_{\mathrm{th}}-R_L)}{(R_{\mathrm{th}}+R_L)^3}.
\]
最大功率对应阻值匹配，但此时等效内阻与负载各消耗一半功率。若追求输电效率，应使线路电阻远小于负载的等效阻值，目标与最大功率传输不同。
网络对称性也必须连同端口激励一起判断：图形上对称的两个节点，在不对称接入电源后未必等势。


\subsection{专题例题与方法应用}

\begin{origexamplebox}{【电流源】}\label{ch21:ex:s37-6-1}
小明利用一个电流源向外供电，但是小红却误以为是一个电压源。小明发现无法说服小红，即电流源和
电压源是等价的。试帮助小明证明这个结论，并利用这个关系求解电压源并联的公式。
\end{origexamplebox}\par\nopagebreak[4]

\par\noindent\begin{minipage}{\linewidth}
\sourceTransformExampleFigure
\diagramcaption{解图：电压源串联内阻与电流源并联内阻的端口方向。两图对同一负载给出相同的 $U,I$。}
\end{minipage}\par\medskip

\begin{xesolution}
这里的“等价”指\textbf{带内阻电源的端口等价}，并不是说孤立的理想电压源与孤立的理想电流源可以
无条件互换。设实际电压源的电动势为 $E$、串联内阻为 $r$，负载端电压为 $U$，向负载输出的电流为
$I$。由 KVL，端口关系为
\[
  U=E-Ir,
  \qquad I=\frac Er-\frac Ur.
\]
若改用输出电流 $I_{\rm N}=E/r$ 的理想电流源，并在端口并联同一个电阻 $r$，KCL 给出
\[
  I=I_{\rm N}-\frac Ur=\frac Er-\frac Ur,
\]
与上式完全相同。因此
\[
  \boxed{E\ \text{串联}\ r
  \quad\Longleftrightarrow\quad
  I_{\rm N}=\frac Er\ \text{并联}\ r}.
\]

现在把 $n$ 个电压源 $(E_i,r_i)$ 同极性并联。先把每一支路变成诺顿形式：电流源 $E_i/r_i$ 与
$r_i$ 并联。并联电流源的电流相加，并联内阻的电导相加，故
\[
  I_{\rm N,eq}=\sum_{i=1}^n\frac{E_i}{r_i},
  \qquad
  \frac1{r_{\rm eq}}=\sum_{i=1}^n\frac1{r_i}.
\]
再变回电压源，得到
\[
  \boxed{r_{\rm eq}=\left(\sum_{i=1}^n\frac1{r_i}\right)^{-1}},
  \qquad
  \boxed{E_{\rm eq}=r_{\rm eq}\sum_{i=1}^n\frac{E_i}{r_i}}.
\]
对两个电源，公式化为
\[
  \boxed{r_{\rm eq}=\frac{r_1r_2}{r_1+r_2},\qquad
  E_{\rm eq}=\frac{E_1r_2+E_2r_1}{r_1+r_2}}.
\]
若某支路极性相反，应把相应 $E_i$ 取负。两个电动势不同、内阻又被理想化为零的电压源不能直接并联，
因为模型会要求无限大的环流；真实电源的内阻正是公式成立所必需的物理条件。
\end{xesolution}



\begin{origexamplebox}{【田电阻】}\label{ch21:ex:s37-6-2}
有一个田字格的电阻网络，每条边电阻为 $R$，试求：

(1) 对角线两个点 $A,C$ 的等效电阻；

(2) 相邻的两个顶点 $A,B$ 的等效电阻。
\par\noindent\begin{minipage}{\linewidth}
\gridProblemFigure
\diagramcaption{九个交点都是节点，十二个相邻节点间的线段各代表一个电阻 $R$。}\label{ch21:fig:s37-6-6}
\end{minipage}\par\medskip
\end{origexamplebox}\par\nopagebreak[4]


\begin{xesolution}
\textbf{(1) 对角端点 $A,C$。} 令 $\varphi_A=V,\ \varphi_C=0$。关于对角线 $AC$ 反射，成对节点
等势；再结合绕中心旋转 $180^\circ$ 后交换 $A,C$，可知中心点和另两个外角的电势均为 $V/2$。
与 $A$ 相邻的两个边中点等势，设其电势为 $x$。其中任一点连接 $A$ 和两个 $V/2$ 节点，KCL 给出
\[
  \frac{x-V}{R}+2\frac{x-V/2}{R}=0,
  \qquad x=\frac{2V}{3}.
\]
从 $A$ 流出的总电流是
\[
  I=2\frac{V-x}{R}=\frac{2V}{3R},
\]
故
\[
  \boxed{R_{AC}=\frac VI=\frac32R}.
\]

\par\noindent\begin{minipage}{\linewidth}
\gridPotentialDiagonalFigure
\diagramcaption{解图（1）：把 $A,C$ 的电势分别归一化为 $1,0$。同色节点等势，标注放在网格空隙内而不遮挡支路。}
\end{minipage}\par\medskip

\textbf{(2) 相邻外角 $A,B$。} 仍令 $\varphi_A=V,\ \varphi_B=0$。左右反射同时交换两个端点，
所以竖直中线上的三个节点均为 $V/2$。设左边中点电势为 $u$、左上角为 $w$，其右侧镜像节点的
电势分别为 $V-u,V-w$。左上角和左边中点的 KCL 为
\[
  2w-u-\frac V2=0,
  \qquad
  3u-V-w-\frac V2=0.
\]
解得 $u=7V/10,\ w=3V/5$。从 $A$ 流出的电流
\[
  I=\frac{V-V/2}{R}+\frac{V-7V/10}{R}=\frac{4V}{5R},
\]
因此
\[
  \boxed{R_{AB}=\frac54R}.
\]

\par\noindent\begin{minipage}{\linewidth}
\gridPotentialAdjacentFigure
\diagramcaption{解图（2）：相邻端点激励不再具有对角情形的全部对称性；图中的九个电势已经逐点满足 KCL。}
\end{minipage}\par\medskip
\end{xesolution}



\begin{origexamplebox}{【正四面体电阻】}\label{ch21:ex:s37-6-3}
有一个正四面体框架的电阻网络，每条边电阻为 $R$，试求任意两点间的等效电阻。
\par\noindent\begin{minipage}{\linewidth}
\tetrahedronProblemFigure
\diagramcaption{四个顶点、六条棱均完整画出；交叉处没有节点。}\label{ch21:fig:s37-6-7}
\end{minipage}\par\medskip
\end{origexamplebox}\par\nopagebreak[4]


\begin{xesolution}
从任意两顶点 $A,B$ 接入电源。其余两顶点 $C,D$ 在保持 $A,B$ 不动的对称变换下可以交换，所以
$\varphi_C=\varphi_D$，连接 $C,D$ 的电阻中没有电流，可删去。此后 $A$ 到 $B$ 之间有三条并联支路：
一条是直接的 $R$，另两条分别是 $A-C-B$ 和 $A-D-B$，每条为 $2R$。于是
\[
  \frac1{R_{AB}}=\frac1R+\frac1{2R}+\frac1{2R}=\frac2R,
  \qquad
  \boxed{R_{AB}=\frac R2}.
\]
这个结论与所选顶点无关，因为正四面体的任意两顶点在网络拓扑上完全等价。
\end{xesolution}



\begin{origexamplebox}{【$N$ 边形电阻】}\label{ch21:ex:s37-6-4}
有一个 $N$ 边形电阻网络（图中 $N=5$），每条边、每条对角线都用阻值为 $R$ 的电阻相连，试求
任意两点间的等效电阻。
\par\noindent\begin{minipage}{\linewidth}
\completeGraphProblemFigure
\diagramcaption{图示为 $K_5$；一般情形是任意两顶点之间都有且只有一个电阻 $R$。}\label{ch21:fig:s37-6-8}
\end{minipage}\par\medskip
\end{origexamplebox}\par\nopagebreak[4]


\begin{xesolution}
取所求端点为 $A,B$，加电压 $V$。其余 $N-2$ 个顶点在保持 $A,B$ 不变的任意置换下都等价，故电势
相同。它们之间虽有大量电阻，但每一条两端等势，电流均为零。由交换 $A,B$ 的反对称性，这些中间点
电势都是 $V/2$。

从 $A$ 流出的电流包括直接流向 $B$ 的 $V/R$，以及流向 $N-2$ 个中间点的
$(N-2)V/(2R)$：
\[
  I=\frac VR+(N-2)\frac{V}{2R}=\frac{NV}{2R}.
\]
所以
\[
  \boxed{R_{AB}=\frac{2R}{N}}.
\]
代入 $N=4$ 就得到上一题正四面体的 $R/2$；二者本质上分别是完全图 $K_N$ 与 $K_4$。

\par\noindent\begin{minipage}{\linewidth}
\completeGraphReductionFigure
\diagramcaption{解图：固定端点 $A,B$ 后，其余 $N-2$ 点全部等势。虚线框表示把这些点视为同一电势层，而非新增导线。}
\end{minipage}\par\medskip
\end{xesolution}



\begin{origexamplebox}{【变种八边形电阻】}\label{ch21:ex:s37-6-5}
有一个八边形电阻网络，除最长的对角线外，每条边、每条对角线都用阻值为 $R$ 的电阻相连，试求相距
最远的两点间的等效电阻。
\par\noindent\begin{minipage}{\linewidth}
\octagonProblemFigure
\diagramcaption{八边形中唯独四条直径不接电阻；任意其他两顶点之间均接一个 $R$。}\label{ch21:fig:s37-6-9}
\end{minipage}\par\medskip
\end{origexamplebox}\par\nopagebreak[4]


\begin{xesolution}
相距最远的端点 $A,B$ 互为对顶点，它们之间恰好没有直径支路。设端电势为 $V,0$，试令其余六点均为
$V/2$：任一中间点与 $A,B$ 相连的两支路电流相互抵消，与其他中间点相连的支路电流为零，所以各点
KCL均满足。正电阻网络在端电势给定后内部电势解唯一，故这就是实际电势，不需要声称任意交换顶点都
保持网络不变。每个中间点都通过一个 $R$ 接到 $A$，并通过另一个 $R$ 接到
$B$；中间点之间的支路两端等势，均无电流。

网络于是化成六条相同的 $A$--中间点--$B$ 支路并联，每条阻值 $2R$：
\[
  \boxed{R_{AB}=\frac{2R}{6}=\frac R3}.
\]
这里不能误加一条 $A,B$ 间的直接电阻；“最长对角线除外”正是本题与完全图 $K_8$ 的区别。

\par\noindent\begin{minipage}{\linewidth}
\octagonReductionFigure
\diagramcaption{解图：六个蓝色中间点等势，内部连线均不载流；真正参与端口电流的是六条长度为两个电阻的并联通路。}
\end{minipage}\par\medskip
\end{xesolution}



\begin{origexamplebox}{【正方体电阻】}\label{ch21:ex:s37-6-6}
有一个立方体电阻网络，每条棱电阻为 $R$。求：

(1) 相邻两顶点；

(2) 面对角线两顶点；

(3) 体对角线两顶点的等效电阻。
\par\noindent\begin{minipage}{\linewidth}
\cubeProblemFigure
\diagramcaption{实线和虚线共同构成立方体的十二条棱；虚线只表示被遮挡的三条棱。}\label{ch21:fig:s37-6-10}
\end{minipage}\par\medskip
\end{origexamplebox}\par\nopagebreak[4]


\begin{xesolution}
统一令输入端电势为 $V$、输出端为 $0$。下图列出三种接法中所有顶点的归一化电势，数值由对称性合并
节点后写 KCL 得到；把每个内部节点的相邻电势代入，都可直接验证流入等于流出。

\par\noindent\begin{minipage}{\linewidth}
\cubePotentialFigure
\diagramcaption{解图：三种端点选择的归一化电势。每个数值紧邻顶点外侧，虚线棱与实线棱的空间关系保持一致。}
\end{minipage}\par\medskip

\textbf{(1) 相邻顶点。} 用三位坐标 $000$ 至 $111$ 标记立方体顶点，一条棱只改变其中一位。
取 $000$ 为输入端、$100$ 为输出端。交换后两坐标可知 $010,001$ 等势，记为 $x$；
反转第一坐标同时交换端电势，可知 $110,101$ 电势为 $V-x$。另记 $011$ 电势为 $y$，则 $111$ 为 $V-y$。
在 $010$ 和 $011$ 两点分别列KCL：
\[
 (x-V)+(x-(V-x))+(x-y)=0,\qquad
 2(y-x)+(y-(V-y))=0.
\]
整理为 $4x-y=2V,\ 4y-2x=V$，消去 $y$ 得 $14x=9V$，再得 $y=4V/7$。
因此输入端除直达输出端的一条棱外，另两条棱通向电势均为 $9V/14$ 的节点，
所以
\[
  I=\frac VR+2\frac{V-9V/14}{R}=\frac{12V}{7R},
  \qquad
  \boxed{R_{\text{邻}}=\frac{7R}{12}}.
\]

\textbf{(2) 面对角线顶点。} 取 $000$ 输入、$110$ 输出。交换前两坐标以及同时反转这两位，
可知 $100,010,101,011$ 电势均为 $V/2$。令 $001$ 的电势为 $x$，它的邻点为 $000,101,011$，故
\[
 (x-V)+2(x-V/2)=0,\qquad x=2V/3.
\]
输入端的两个同面邻点电势均为 $V/2$，第三个邻点电势为 $2V/3$，故
\[
  I=2\frac{V-V/2}{R}+\frac{V-2V/3}{R}=\frac{4V}{3R},
  \qquad
  \boxed{R_{\text{面}}=\frac{3R}{4}}.
\]

\textbf{(3) 体对角线顶点。} 输入端的三个邻点等势，设为 $x$；输出端的三个邻点由反对称性为 $V-x$。
任取一个 $x$ 节点，它连接输入端和两个 $V-x$ 节点，故
\[
  \frac{x-V}{R}+2\frac{x-(V-x)}R=0,
  \qquad x=\frac{3V}{5}.
\]
于是
\[
  I=3\frac{V-3V/5}{R}=\frac{6V}{5R},
  \qquad
  \boxed{R_{\text{体}}=\frac{5R}{6}}.
\]
三种结果依次增大，但这不是由“几何距离”直接决定的，而是由端点选择后整个网络允许的并联通道决定。
\end{xesolution}



\section{多端、无限与自相似网络}

\subsection{电流分布法}


求两点间等效电阻时，常规做法是在两点间加电压 $U$，求出总电流 $I$，再用 $R_{\rm eq}=U/I$。电流
分布法则反过来：从一点灌入给定电流 $I$，从另一点抽出同样的电流，先利用对称性和 KCL 确定各支路
如何分流，再求两点间电势差 $U$，最后用 $R_{\rm eq}=U/I$。

这种方法同样服从叠加。例如，要得到“电流从 $A$ 进、从 $B$ 出”的分布，可以把“从 $A$ 进、从
$C$ 出”和“从 $C$ 进、从 $B$ 出”两种分布按相同电流强度相加。中间节点 $C$ 的注入与抽出相互抵消，
剩下的正是所求端口条件。


% EM teaching enrichment 20260927
\paragraph{分流必须同时满足回路电压条件}
只写 KCL 可以构造许多看似守恒的支路电流，却未必对应真实电势。得到候选分流后，还应检查任意独立回路的 $\sum IR=0$。
一个实用核验是功率：若从 $A$ 注入 $I$、从 $B$ 抽出 $I$，则
\[
 I(V_A-V_B)=\sum_{\text{各电阻支路}}R_kI_k^2,
 \qquad R_{AB}=\frac{\sum_kR_kI_k^2}{I^2}.
\]
每条支路只计一次。对三角形三个等阻 $R$、端口取任意两顶点，直连支路分得 $2I/3$，绕行两阻各为 $I/3$，代入得 $R_{AB}=2R/3$，与 $R\parallel2R$ 相符。


\subsection{多端纯电阻网络}


只有 $N$ 个端口能与外界连接、内部只含电阻的电路，称为 $N$ 端纯电阻网络。一般的无源线性 $N$ 端
电阻网络，可以用每两端口之间各接一个适当电阻的完全图表示。端口若具有对称性，若干等效电阻可能相等；
若无对称性，完全图中的各支路阻值一般不同。

三端网络是最重要的特例。它既可写成三角形（$\Delta$）连接，也可写成星形（$Y$）连接。设端口电势
为 $\varphi_a,\varphi_b,\varphi_c$，星形中心电势为 $\varphi_0$，三条星形支路为 $R_a,R_b,R_c$。
中心节点满足 KCL：
\begin{equation}
  \frac{\varphi_a-\varphi_0}{R_a}
  +\frac{\varphi_b-\varphi_0}{R_b}
  +\frac{\varphi_c-\varphi_0}{R_c}=0.
\end{equation}
解得
\begin{equation}
  \varphi_0=
  \frac{\dfrac{\varphi_a}{R_a}+\dfrac{\varphi_b}{R_b}+\dfrac{\varphi_c}{R_c}}
       {\dfrac1{R_a}+\dfrac1{R_b}+\dfrac1{R_c}}
  =\frac{\varphi_aR_bR_c+\varphi_bR_cR_a+\varphi_cR_aR_b}
  {R_aR_b+R_bR_c+R_cR_a}.
\end{equation}
由端口 $A$ 流入星形网络的电流为
\begin{equation}
  I_a=\frac{\varphi_a-\varphi_0}{R_a}
  =\frac{\varphi_a(R_b+R_c)-\varphi_bR_c-\varphi_cR_b}
  {R_aR_b+R_bR_c+R_cR_a}
  =\frac{\varphi_a-\varphi_b}{R_{AB}}
  +\frac{\varphi_a-\varphi_c}{R_{AC}}.
\end{equation}
最后一个等号是三角形连接的端口电流。令
\[
  S=R_aR_b+R_bR_c+R_cR_a,
\]
比较 $\varphi_a,\varphi_b,\varphi_c$ 的系数，并循环置换端口，可得星形变三角形：
\[
  \boxed{R_{AB}=\frac{S}{R_c},\qquad
  R_{AC}=\frac{S}{R_b},\qquad
  R_{BC}=\frac{S}{R_a}}.
\]
反解得到三角形变星形：
\[
\boxed{
  R_a=\frac{R_{AB}R_{AC}}{R_{AB}+R_{AC}+R_{BC}},\quad
  R_b=\frac{R_{AB}R_{BC}}{R_{AB}+R_{AC}+R_{BC}},\quad
  R_c=\frac{R_{AC}R_{BC}}{R_{AB}+R_{AC}+R_{BC}}.}
\]

\par\noindent\begin{minipage}{\linewidth}
\deltaStarFigure
\diagramcaption{同一三端网络的三角形连接与星形连接。对应关系由端口字母决定，不应靠电阻在图中的朝向记忆。}\label{ch21:fig:s37-6-5}
\end{minipage}\par\medskip

星形--三角形变换能把许多表面上既不串联也不并联的网络改写成可逐级化简的形式；但在高度对称或
自相似网络中，先用等势、分流和递归往往更简洁。

% EM teaching enrichment 20260927
\paragraph{用三个端口状态检验星三角变换}
若三角形三边均为 $R$，对应星形每臂为 $R/3$。只在 $A,B$ 之间通电、$C$ 外端开路时，星形等效为 $2R/3$，恰与三角形的 $R\parallel2R$ 相同。
但若把 $C$ 接地，它就不再是开路条件，不能把通向 $C$ 的星形臂删掉。三端等效保证任意端电势下的三端电流均相同，比只匹配某一对端口的一个电阻值更强。
复杂网络宜先识别可局部替换的三端块，替换时保留端口名字；随后重新判断串、并联关系，避免按图形朝向误配公式分母。




% EM structure and teaching revision 20260928
\subsection{无限网络的递推、选根与收敛}
自相似结构允许用整体的等效量表示去掉首段后剩下的部分。设每段先串联一个 $R$，再把一个 $R$ 对地分流，余下网络等效仍为 $X$，则
\[
 X=R+\frac{RX}{R+X},\qquad X^2-RX-R^2=0.
\]
正电阻网络取 $X=(1+\sqrt5)R/2$。负根虽满足代数方程，却不满足被动端口的非负耗散条件。
更可靠的定义来自有限截断：从某个非负末端电阻 $X_0$ 出发，令 $X_{n+1}=R+RX_n/(R+X_n)$。一步后有 $R\le X_n\le2R$，且该区间内导数
$f'(X)=R^2/(R+X)^2\le1/4$，所以相邻近似间的差不断缩小，极限确为所选正根。
有多个端口时，一个标量等效电阻往往不够。应保留端口电势和注入电流的线性关系，或者利用对称性把端口激励分解成独立模式。未经证明就把三端网络替换成两端电阻，会丢失约束，使看似正确的递推求出错误结果。


\subsection{专题例题与方法应用}

\begin{origexamplebox}{【无穷网络】}\label{ch21:ex:s37-6-7}
求图示无限电阻网络中 $a,b$ 两点间的等效电阻。图中每个电阻的阻值均为 $R$。
\par\noindent\begin{minipage}{\linewidth}
\infiniteLadderProblemFigure
\diagramcaption{上下两条横向支路和每一根竖直梯档都为 $R$，网络向右无限延伸。}\label{ch21:fig:s37-6-11}
\end{minipage}\par\medskip
\end{origexamplebox}\par\nopagebreak[4]


\begin{xesolution}
设整个无限梯形网络的等效电阻为 $X$。删去最左边一格后，余下部分与原网络完全相同，因此从第一对
竖直节点向右看进去仍是 $X$。在这两个节点之间，第一根竖直电阻 $R$ 与余下网络 $X$ 并联；它们外侧
还各串联一个横向电阻 $R$。所以
\[
  X=2R+\frac{RX}{R+X}.
\]

\par\noindent\begin{minipage}{\linewidth}
\ladderRecursionFigure
\diagramcaption{解图：虚线框内的无限尾部仍等效为 $X$。两只横向 $R$ 串在 $R\parallel X$ 的两侧。}
\end{minipage}\par\medskip

令 $x=X/R$，有
\[
  x=2+\frac{x}{1+x},
  \qquad x^2-2x-2=0.
\]
电阻必须为正，舍去 $x=1-\sqrt3$，得到
\[
  \boxed{R_{ab}=X=(1+\sqrt3)R}.
\]
回代可见 $X>2R$，与“入口处必经上下两个 $R$，中间还有有限并联电阻”的结构相符。
\end{xesolution}



\begin{origexamplebox}{【三端无穷网络】}\label{ch21:ex:s37-6-8}
求图示无限电阻网络中 $A,C$ 两点和 $A,B$ 两点之间的等效电阻，已知图中每个电阻值均为 $R$。
\par\noindent\begin{minipage}{\linewidth}
\threeRailProblemFigure
\diagramcaption{题图示意：三条横向链向右无限延伸，相邻两行在每一列以竖直电阻连接；左端为 $A,B,C$。}
\end{minipage}\par\medskip
\end{origexamplebox}\par\nopagebreak[4]


\begin{xesolution}
这个网络有三个端口，不能只用一个标量电阻表示尾部。利用上下反射对称，把每一列的三个电势分解成
彼此独立的三种横向模式：
\[
  \boldsymbol e_0=(1,1,1),\qquad
  \boldsymbol e_1=(1,0,-1),\qquad
  \boldsymbol e_3=(1,-2,1).
\]
这些模式可直接从KCL导出。令一列的电势向量为 $\boldsymbol V=(V_A,V_B,V_C)^{\mathsf T}$，
流向本列两根竖直电阻的电流向量为
\[
 \boldsymbol I_{\rm vertical}=\frac1R
 \begin{pmatrix}1&-1&0\\-1&2&-1\\0&-1&1\end{pmatrix}\boldsymbol V
 =\frac1R L\boldsymbol V.
\]
逐项相乘得到 $L\boldsymbol e_0=0$、$L\boldsymbol e_1=\boldsymbol e_1$、
$L\boldsymbol e_3=3\boldsymbol e_3$。横向三根电阻同为 $R$，不会混合这些正交模式，
因此后两种模式分别看到无量纲横向电导 $\mu=1,3$。
端口一注入一抽出、总电流为零，不激发沿链流动的共同模式；各列相同的共同电势只表示参考电势的选择。

\par\noindent\begin{minipage}{\linewidth}
\threeRailModesFigure
\diagramcaption{解图：三行节点的两个有效电势模式。箭头长度表示相对电势（或注入电流）分量，不是实际支路。}
\end{minipage}\par\medskip

设某一模式从边界看入的无量纲输入电导为 $y_\mu$，真实电导为 $y_\mu/R$。越过最左边的横向电阻后，
节点处同时接有横向电导 $\mu/R$ 和一个完全相同的无限尾部 $y_\mu/R$。串并联关系给出
\[
  y_\mu=\frac{\mu+y_\mu}{1+\mu+y_\mu},
\]
所以
\[
  y_\mu^2+\mu y_\mu-\mu=0,
  \qquad
  y_1=\frac{\sqrt5-1}{2},
  \qquad
  y_3=\frac{\sqrt{21}-3}{2}.
\]

为说明端口电阻如何组合，令单位模态向量为 $\boldsymbol u_j=\boldsymbol e_j/|\boldsymbol e_j|$，
注入电流为 $I\boldsymbol b$。其模态电流是 $I(\boldsymbol b\cdot\boldsymbol u_j)$，模态电压是
$IR(\boldsymbol b\cdot\boldsymbol u_j)/y_j$。端电压为 $\boldsymbol b\cdot\boldsymbol V$，故
\[
 R_{\rm port}=R\sum_{j=1,3}\frac{(\boldsymbol b\cdot\boldsymbol u_j)^2}{y_j}.
\]
对 $A,C$ 注入与抽出的电流向量正比于 $(1,0,-1)=\boldsymbol e_1$，只激发 $\mu=1$ 模式。该向量
长度平方为 $2$，故
\[
  \boxed{R_{AC}=\frac{2R}{y_1}=(1+\sqrt5)R}.
\]

对 $A,B$，向量 $(1,-1,0)$ 在归一化的 $\boldsymbol e_1,\boldsymbol e_3$ 上的投影平方分别为
$1/2,3/2$。两个正交模式的电压降相加：
\[
  R_{AB}=R\left(\frac{1/2}{y_1}+\frac{3/2}{y_3}\right)
  =\boxed{\frac{4+\sqrt5+\sqrt{21}}4R}.
\]
数值上 $R_{AB}\approx2.70466R$、$R_{AC}\approx3.23607R$。这也符合端点距离直觉：$A,C$ 分处
最上、最下两行，等效电阻较大。
\end{xesolution}



\begin{origexamplebox}{【嵌套网络】}\label{ch21:ex:s37-6-9}
如图所示，框架由同种金属丝制成，单位长度的电阻为 $\rho$。一连串内接等边三角形的数目可认为趋向
无穷，取 $AB$ 边长为 $a$，以后每个三角形的边长依次减小一半。求框架上 $A,B$ 两点间的电阻。
\par\noindent\begin{minipage}{\linewidth}
\nestedTriangleProblemFigure
\diagramcaption{每个新三角形的三个顶点是上一三角形的三边中点，方向逐级交替。}\label{ch21:fig:s37-6-12}
\end{minipage}\par\medskip
\end{origexamplebox}\par\nopagebreak[4]


\begin{xesolution}
任意一层以后仍是一个关于三个端点完全对称的三端网络，所以可把边长为 $a$ 的完整无限网络等效成
三角形连接，设每一条等效支路为
\[
  x(a)=k\rho a,
\]
其中无量纲常数 $k$ 与尺度无关。记 $r=\rho a$。外三角形的每条边被中点分成两个 $r/2$；中点围成的
无限内网边长为 $a/2$，其三角形等效支路就是 $kr/2$。

\par\noindent\begin{minipage}{\linewidth}
\nestedTriangleRecursionFigure
\diagramcaption{解图：把中点以内的无限嵌套部分替换成三条 $kr/2$。外层六段金属丝各为 $r/2$。}
\end{minipage}\par\medskip

在图中令 $A,B$ 电势为 $V,0$。反射对称性与电势反对称性给出 $C$ 及 $AB$ 中点电势为 $V/2$；
令 $CA$ 中点电势为 $u$，则 $CB$ 中点为 $V-u$。在 $CA$ 中点列KCL：
\[
 \frac{u-V}{r/2}+\frac{u-V/2}{r/2}
 +\frac{u-V/2}{kr/2}+\frac{u-(V-u)}{kr/2}=0.
\]
乘 $kr/2$ 后得 $(2k+3)u=3(k+1)V/2$。从 $A$ 流入的总电流为
\[
 I=\frac{V-V/2}{r/2}+\frac{V-u}{r/2}
 =\frac{3(k+2)V}{(2k+3)r}.
\]
所以 $A,B$ 间等效电阻
\[
  R_{AB}=\frac{r(2k+3)}{3(k+2)}.
\]
另一方面，若从完整三端网络的等效三角形看，$A,B$ 之间是一条 $kr$ 与另两条串联的 $2kr$ 并联，
因此
\[
  R_{AB}=kr\parallel 2kr=\frac{2kr}{3}.
\]
令两式相等：
\[
  \frac{2k}{3}=\frac{2k+3}{3(k+2)},
  \qquad 2k^2+2k-3=0.
\]
取正根 $k=(\sqrt7-1)/2$，最终
\[
  \boxed{R_{AB}=\frac{\sqrt7-1}{3}\,\rho a}.
\]
结果小于只保留外三角形时的 $2\rho a/3$，因为不断加入的内接金属丝提供了额外并联通路。
\end{xesolution}



\begin{origexamplebox}{【分形电阻】}\label{ch21:ex:s37-6-10}
如图，考虑大三角形 $ABC$，取每条边的中点并连接，得到第二阶图形，其中含有三个正立小正三角形。
把中间的倒立三角形镂空，便得到一阶“镂垫”。对每一个保留的小三角形重复相同操作，得到更高阶图形。
如此无穷进行，得到谢尔宾斯基三角形，也称谢尔宾斯基镂垫。讨论下列性质：

(1) 最终镂垫的总面积是多少？

(2) 最终镂垫的总边长是多少（内、外边界均计入）？

(3) 若所有边都由单位长度电阻为 $\lambda$ 的电阻丝构成，$AB$ 两点的等效电阻为多少？

(4) 讨论上述结果。
\par\noindent\begin{minipage}{\linewidth}
\sierpinskiProblemFigure
\diagramcaption{前三阶图形由同一中点规则生成；白色倒三角形是被镂空的部分。}\label{ch21:fig:s37-6-13}
\end{minipage}\par\medskip
\end{origexamplebox}\par\nopagebreak[4]


\begin{xesolution}
设最初等边三角形边长为 $a$、面积为 $S_0=\sqrt3a^2/4$。

\textbf{(1) 面积。} 每次操作把一个实心三角形分成四个全等小三角形，只保留其中三个。因此第 $n$ 次
操作后
\[
  S_n=S_0\left(\frac34\right)^n.
\]
令 $n\to\infty$，得到
\[
  \boxed{S_\infty=0}.
\]

\textbf{(2) 总边长。} 第 $n$ 阶有 $3^n$ 个最小正三角形，每个边长 $a/2^n$；这些三角形彼此只在
顶点相接，不共享完整边。把每条几何线段只计一次，总边长为
\[
  L_n=3^n\cdot3\frac{a}{2^n}
  =3a\left(\frac32\right)^n.
\]
所以
\[
  \boxed{L_\infty=\infty}.
\]

\textbf{(3) 等效电阻。} 第 $n$ 阶的最小边长为 $a/2^n$，每段电阻
\[
  r_n=\lambda\frac{a}{2^n}.
\]
若所有最小边的电阻都暂记为 $r$，三角形初阶的端点电阻为 $2r/3$。下面推导增加一阶的比例。
把三个小三角形拼接成一阶网，九条边均为 $r$。令 $A,B$ 电势为 $V,0$，则 $C$ 与 $AB$ 中点为 $V/2$；
$CA,CB$ 中点分别为 $u,V-u$。前一中点连接 $A,C$ 及另外两个中点，KCL给出
\[
 (u-V)+(u-V/2)+(u-V/2)+(u-(V-u))=0,
 \qquad u=3V/5.
\]
所以 $I=(V-V/2)/r+(V-3V/5)/r=9V/(10r)$，一阶端电阻为 $10r/9$。
与初阶 $2r/3$ 相比，比例为 $5/3$。任意更高阶的子网都可等效为对称三端三角形，因而同一比例逐阶适用。
因此
\[
  R_n=\frac23\left(\frac53\right)^n r_n
  =\frac{2\lambda a}{3}\left(\frac56\right)^n.
\]

\par\noindent\begin{minipage}{\linewidth}
\sierpinskiResistanceFigure
\diagramcaption{解图：每一阶由三个缩小到 $1/2$ 的同型三端网络拼成；拓扑使等效电阻乘 $5/3$，长度缩放又使单段电阻乘 $1/2$。}
\end{minipage}\par\medskip

取极限得
\[
  \boxed{R_{AB}=0}.
\]
这是理想无限阶模型的极限，不表示有限粗细的真实金属丝会有严格为零的电阻。

\textbf{(4) 讨论。} 这个分形同时具有“面积趋于零、总边长趋于无穷、端点电阻趋于零”三种看似矛盾
的性质。它们分别测量二维占据、线性边界和电流并联通道，不应互相替代。面积每阶乘 $3/4$；总边长
每阶乘 $3/2$；电阻则是拓扑因子 $5/3$ 与几何缩放 $1/2$ 的乘积，即每阶乘 $5/6$。真实导线有有限
线径、接触电阻和最小结构尺度，递归在某一阶必然停止，因而真实面积、长度和电阻都保持有限且非零。
\end{xesolution}
}

\endgroup

\endgroup

\begingroup
\tcbset{breakable/.default=false,after={\par\nobreak\smallskip}}
\RenewDocumentEnvironment{xesolution}{+b}{%
  \par\nobreak\smallskip\nobreak\begingroup\small
  \clubpenalties 2 10000 10000\displaywidowpenalties 2 10000 10000\noindent
  {\bfseries\color{DeepBlue}解答}\quad #1\par\endgroup\smallskip
}{}
\begingroup
\renewcommand{\Needspace}[1]{\par\begingroup
\dimen0=#1\relax\vskip0pt plus\dimen0\penalty-100
\vskip0pt plus-\dimen0\vskip\dimen0\penalty9999\vskip-\dimen0
\vskip0pt\endgroup}
\tcbset{breakable/.default=false}
\chapter{洛伦兹力与带电粒子运动}

{
\graphicspath{{/Users/dongxingxi/Documents/物理竞赛一轮资料/sample/xes-original/thirty-eighth-section/}{/Users/dongxingxi/Documents/物理竞赛一轮资料/sample/xes-original/thirty-eighth-section/assets/}}
\section{洛伦兹力与回旋轨道}

% Restrained diagrams for Lesson 7.  Every arrow follows F=q v x B.

\newcommand{\fieldInto}{\mathord{\otimes}}
\newcommand{\fieldOut}{\mathord{\odot}}

\newcommand{\lorentzDirectionDiagram}{%
\begin{center}
\begin{tikzpicture}[scale=0.86,>=Latex]
  \begin{scope}[xshift=-5.2cm]
    \draw[rounded corners=4pt,draw=RuleGray] (-2.1,-1.55) rectangle (2.1,1.55);
    \fill[DeepBlue] (0,0) circle (3pt) node[below left] {$+q$};
    \draw[vector] (0,0)--(1.55,0) node[right] {$\vect v$};
    \draw[vector] (0,0)--(0,1.18) node[above] {$\vect F$};
    \node[font=\Large] at (-1.25,0.65) {$\fieldInto$};
    \node[TextGray] at (0,-1.9) {$q>0$，$\vect B$ 入纸面};
  \end{scope}
  \begin{scope}
    \draw[rounded corners=4pt,draw=RuleGray] (-2.1,-1.55) rectangle (2.1,1.55);
    \fill[Ink] (0,0) circle (3pt) node[above left] {$-q$};
    \draw[vector] (0,0)--(1.55,0) node[right] {$\vect v$};
    \draw[vector] (0,0)--(0,-1.18) node[below] {$\vect F$};
    \node[font=\Large] at (-1.25,0.65) {$\fieldInto$};
    \node[TextGray] at (0,-1.9) {$q<0$，受力反向};
  \end{scope}
  \begin{scope}[xshift=5.2cm]
    \draw[rounded corners=4pt,draw=RuleGray] (-2.1,-1.55) rectangle (2.1,1.55);
    \fill[DeepBlue] (0,0) circle (3pt);
    \node at (0.35,-0.48) {$+q$};
    \draw[guide,dashed] (0,0)--(-0.72,-0.52)--(-0.72,0.78)--(0,0);
    \draw[vector] (0,0)--(1.42,0) node[right] {$\vect v$};
    \draw[vector] (0,0)--(-0.72,-0.52) node[below left] {$\vect B$};
    \draw[vector] (0,0)--(0,1.30) node[above] {$\vect F$};
    \node[TextGray,font=\small] at (0,-1.9) {三维正交示意（纸面上为投影）};
  \end{scope}
\end{tikzpicture}

{\small\color{TextGray}$\vect F=q\vect v\times\vect B$；先按正电荷定向，负电荷再反向}
\end{center}}

\newcommand{\gyroHelixDiagram}{%
\begin{center}
\begin{tikzpicture}[scale=0.82,>=Latex]
  \begin{scope}[xshift=-4.4cm]
    \draw[curve] (0,0) circle (1.55);
    \node[font=\Large] at (0,0) {$\fieldOut$};
    \fill[DeepBlue] (0,1.55) circle (2.6pt) node[above left] {$q>0$};
    \draw[vector] (0,1.55)--(1.25,1.55) node[right] {$\vect v_\perp$};
    \draw[vector] (0,1.55)--(0,0.48) node[right] {$\vect F$};
    \draw[<->] (0,0)--(1.55,0);
    \node[above=2pt,inner sep=0.5pt] at (0.78,0) {$r_{\rm c}$};
    \draw[->,DeepBlue,thick] (1.85,0.35) arc[start angle=10,end angle=-70,radius=0.55];
    \node[TextGray] at (0,-2.05) {$v_\parallel=0$：匀速圆周运动};
  \end{scope}
  \begin{scope}[xshift=3.3cm]
    \draw[axis] (-0.5,0)--(7.0,0) node[right] {$\vect B$};
    \draw[guide,dashed] (0,0)--(6.2,0);
    % Oblique projection of a true helix.  The cosine term supplies the
    % front/back displacement, so the path forms visible loops rather than
    % being mistaken for a transverse sine wave.
    \draw[DeepBlue,very thick,domain=0:1440,samples=440,smooth,variable=\t]
      plot ({0.0042*\t+0.70*cos(\t)},{0.78*sin(\t)});
    \fill[DeepBlue] (0,0) circle (2.4pt);
    \draw[vector] (0,0)--(1.4,0) node[below] {$\vect v_\parallel$};
    \draw[vector] (0,0)--(0,1.0) node[left] {$\vect v_\perp$};
    \node[TextGray] at (3.1,-1.55) {$v_\parallel\ne0$：沿磁场轴线作螺旋运动};
  \end{scope}
\end{tikzpicture}
\end{center}}

\newcommand{\particleTracksContext}{%
\historyfigure[0.76\linewidth]{particle-tracks-built-in-v1.png}%
{云雾室与气泡室径迹。曲率给出动量信息，弯曲方向结合磁场方向可判定电荷符号。}\label{ch22:fig:s38-7-3}}

\newcommand{\cyclotronDiagram}{%
\begin{center}
\begin{tikzpicture}[scale=0.88,>=Latex]
  \fill[PaleBlue] (-3.5,-2.5) arc[start angle=270,end angle=90,radius=2.5] -- (-0.25,2.5)--(-0.25,-2.5)--cycle;
  \fill[SoftGray] (3.5,-2.5) arc[start angle=-90,end angle=90,radius=2.5] -- (0.25,2.5)--(0.25,-2.5)--cycle;
  \draw[DeepBlue,thick] (-3.5,-2.5) arc[start angle=270,end angle=90,radius=2.5] -- (-0.25,2.5)--(-0.25,-2.5)--cycle;
  \draw[DeepBlue,thick] (3.5,-2.5) arc[start angle=-90,end angle=90,radius=2.5] -- (0.25,2.5)--(0.25,-2.5)--cycle;
  \node at (-1.7,0) {D}; \node at (1.7,0) {D};
  \draw[<->,Ink,thick] (-0.18,2.9)--(0.18,2.9) node[midway,above=3pt] {交变电压};
  \node[font=\Large] at (4.7,1.65) {$\fieldOut$};
  \node[right] at (5.05,1.65) {$\vect B$};
  \draw[DeepBlue,very thick,domain=0:900,samples=240,smooth]
    plot ({0.0028*\x*cos(\x)},{0.0028*\x*sin(\x)});
  \draw[vector] (0.1,0)--(0.9,0) node[below right] {$q\Delta U$};
  \node[TextGray,align=left,anchor=west] at (4.35,-0.15)
    {$\omega_{\rm rf}=\omega_{\rm c}=|q|B/m$\\[2pt]
     每过狭缝一次，动能增加 $|q|\Delta U$};
\end{tikzpicture}

{\small\color{TextGray}回旋加速器俯视图：D 形盒内由磁场弯转，狭缝处由交变电场加速}
\end{center}}

\newcommand{\zeemanSpectrumDiagram}{%
\begin{center}
\begin{tikzpicture}[x=1cm,y=1cm,>=Latex]
  \node[inner sep=0] at (-4.2,0.4) {\includegraphics[width=5.0cm]{lorentz-zeeman-built-in-v1.png}};
  \node[TextGray] at (-5.35,-2.05) {洛伦兹};
  \node[TextGray] at (-3.05,-2.05) {塞曼};
  \begin{scope}[xshift=0.2cm]
    \node[TextGray] at (2.2,1.8) {无磁场};
    \node[TextGray,font=\small] at (2.2,-2.8) {锌谱线 $\lambda_0\simeq468.0\,\mathrm{nm}$};
    \draw[Ink,very thick] (2.2,-0.8)--(2.2,1.35);
    \node[TextGray] at (6.2,1.8) {加磁场};
    \draw[DeepBlue,very thick] (5.25,-0.8)--(5.25,1.35);
    \draw[Ink,very thick] (6.2,-0.8)--(6.2,1.35);
    \draw[DeepBlue,very thick] (7.15,-0.8)--(7.15,1.35);
    \draw[<->] (5.25,-1.2)--(6.2,-1.2);
    \draw[<->] (6.2,-1.2)--(7.15,-1.2);
    \node[below=2pt,inner sep=0.5pt] at (5.73,-1.2) {$\Delta\omega$};
    \node[below=2pt,inner sep=0.5pt] at (6.68,-1.2) {$\Delta\omega$};
    \node[TextGray,align=center] at (4.7,-2.05)
      {$\omega_0$ 与 $\omega_0\pm |q|B/(2m)$\\弱场下形成对称三线分裂};
  \end{scope}
\end{tikzpicture}
\end{center}}

\newcommand{\driftDiagram}{%
\begin{center}
\begin{tikzpicture}[x=0.95cm,y=0.95cm,>=Latex]
  \draw[axis] (-0.4,0)--(9.4,0) node[right] {$\vect v_D$};
  \draw[DeepBlue,very thick,domain=0:720,samples=220,smooth]
    plot ({0.012*\x-0.72*sin(\x)},{0.72*(1-cos(\x))});
  \draw[TextGray,dashed,domain=0:720,samples=220,smooth]
    plot ({0.012*\x},{0.72});
  \foreach \x in {1.2,5.5,9.0}{\node[font=\large] at (\x,2.2) {$\fieldInto$};}
  \draw[vector] (0.8,-0.55)--(2.4,-0.55) node[midway,below] {$\vect F\times\vect B$};
  \node[TextGray,align=left,anchor=west] at (9.7,0.9)
    {$\displaystyle\vect v_D=\frac{\vect F\times\vect B}{qB^2}$\\导向中心匀速平移};
\end{tikzpicture}
\end{center}}

\newcommand{\splitProductsDiagram}{%
\begin{center}
\begin{tikzpicture}[scale=0.82,>=Latex]
  \begin{scope}[xshift=-4.2cm]
    \node[font=\large] at (-2.55,1.05) {$\fieldInto$};
    \node[left=2pt] at (-2.25,1.05) {$\vect B$};
    \fill[DeepBlue] (0,2.0) circle (2.5pt);
    \node[above=8pt,inner sep=0.5pt] at (0,2.0) {$P_0$};
    \draw[DeepBlue,very thick] (0,2.0) arc[start angle=90,end angle=450,radius=1.0];
    \draw[Ink,very thick] (0,2.0) arc[start angle=90,end angle=-270,radius=2.0];
    \draw[vector] (0,2.0)--(-1.25,2.0);
    \draw[vector] (0,2.0)--(1.25,2.0);
    \node[above=2pt,anchor=south east,inner sep=0.5pt] at (-0.42,2.0) {$+q,m,\ v_0$};
    \node[above=2pt,anchor=south west,inner sep=0.5pt] at (0.42,2.0) {$-q,2m,\ v_0$};
    \node[TextGray,align=center,font=\small] at (0,-2.45)
      {碰撞前：半径 $r,2r$\\共同返回周期 $2T$};
  \end{scope}
  \begin{scope}[xshift=4.2cm]
    \fill[DeepBlue] (0,-1.1) circle (2.5pt) node[below left] {$P_0$};
    \draw[DeepBlue,very thick] (0,-1.1) arc[start angle=-90,end angle=630,radius=2.0];
    \draw[Ink,very thick] (0,-1.1) arc[start angle=-90,end angle=-450,radius=0.8];
    \draw[vector] (0,-1.1)--(1.55,-1.1);
    \draw[vector] (0,-1.1)--(-1.0,-1.1);
    \node[below=2pt,anchor=north west,inner sep=0.5pt] at (0.45,-1.1) {$5v_0/3$};
    \node[below=2pt,anchor=north east,inner sep=0.5pt] at (-0.35,-1.1) {$v_0/3$};
    \node[TextGray,align=center,font=\small] at (0,-2.45)
      {第一次碰撞后：半径 $5r/3,2r/3$\\再过 $2T$ 恢复初态};
  \end{scope}
\end{tikzpicture}
\end{center}}

\newcommand{\beamFocusDiagram}{%
\begin{center}
\begin{tikzpicture}[x=0.95cm,y=0.95cm,>=Latex]
  \coordinate (BeamO) at (0,0);
  \coordinate (BeamZ) at (1.2,0);
  \coordinate (BeamV) at (1.35,0.45);
  \draw[axis] (0,0)--(9.0,0) node[right] {$z$};
  \draw[thick] (8,-2.0)--(8,2.0) node[above] {荧光屏};
  \draw[DeepBlue,very thick,domain=0:720,samples=220,smooth]
    plot ({0.0111*\x},{0.75*sin(\x)});
  \draw[TextGray,dashed] (0,0)--(8,0);
  \draw[vector] (0,0)--(1.35,0.45) node[above] {$\vect v$};
  \draw pic[draw,"$\theta$",angle radius=8mm,angle eccentricity=1.25]
    {angle=BeamZ--BeamO--BeamV};
  \draw[<->] (0,-1.45)--(8,-1.45);
  \node[below=2pt,inner sep=0.5pt] at (4,-1.45) {$L$};
  \draw[<->,DeepBlue] (8.35,0)--node[right] {$\rho_{\max}$}(8.35,0.72);
\end{tikzpicture}

{\small\color{TextGray}纵向磁场使横向速度转过相位 $\omega_{\rm c}L/v$；整数圈处重新聚焦}
\end{center}}

\newcommand{\magneticBilliardDiagram}{%
\begin{center}
\begin{tikzpicture}[scale=0.86,>=Latex]
  \coordinate (BucketO) at (0,0);
  \draw[curve] (0,0) circle (3.0);
  \coordinate (P) at (0,-3);
  \coordinate (C) at (-0.866,-3.5);
  \coordinate (Qhit) at (-1.396,-2.652);
  \coordinate (Rin) at (0,-1.8);
  \coordinate (Vtip) at (-0.65,-1.7);
  \draw[DeepBlue,very thick] (P) arc[start angle=30,end angle=122,radius=1.0];
  \draw[guide] (0,0)--(P) (0,0)--(-2.52,-1.62);
  \draw[guide] (C)--(P);
  \fill[DeepBlue] (P) circle (2.5pt) node[below] {小孔};
  \fill[Ink] (C) circle (2.2pt) node[below left] {回旋圆心};
  \draw[vector] (P)--++(-0.65,1.3) node[left] {$\vect v$};
  \draw pic[draw,"$\theta$",angle radius=7mm,angle eccentricity=1.25]
    {angle=Rin--P--Vtip};
  \draw pic[draw,"$\delta$",angle radius=10mm,angle eccentricity=1.25]
    {angle=Qhit--BucketO--P};
  \draw[<->] (C)--(P);
  \node[right=2pt,inner sep=0.5pt] at ($(C)!0.52!(P)$) {$r_{\rm c}$};
  \node[TextGray,align=left,anchor=west] at (4.15,0.4)
    {每次反射保持入射角；\\相邻碰撞点的圆心角恒为 $\delta$。\\[2pt]
     $\delta/(2\pi)\in\mathbb Q$ 时才会返回小孔。};
  \node[font=\large] at (1.9,1.9) {$\fieldInto$};
\end{tikzpicture}
\end{center}}

\newcommand{\velocitySelectorDiagram}{%
\begin{center}
\begin{tikzpicture}[scale=0.9,>=Latex]
  \draw[thick] (-4,1.65)--(4,1.65) (-4,-1.65)--(4,-1.65);
  \foreach \x in {-3.5,-2.5,...,3.5}{
    \node at (\x,1.95) {$+$}; \node at (\x,-1.95) {$-$};
    \node[font=\large,TextGray] at (\x,0.75) {$\fieldInto$};
  }
  \fill[DeepBlue] (-2.2,0) circle (3pt) node[below left] {$+q$};
  \draw[vector] (-2.2,0)--(-0.5,0) node[right] {$\vect v$};
  \draw[vector] (-2.2,0)--(-2.2,1.1) node[left] {$qvB$};
  \draw[vector] (-2.2,0)--(-2.2,-1.1) node[left] {$qE$};
  \draw[vector] (2.8,1.1)--(2.8,-0.7) node[midway,right] {$\vect E$};
  \node[TextGray] at (0,-2.55) {只有 $v=E/B$ 的粒子不偏转；电荷符号改变时两力同时反向};
\end{tikzpicture}
\end{center}}

\newcommand{\crossedFieldCycloidDiagram}{%
\begin{center}
\begin{tikzpicture}[x=1cm,y=1cm,>=Latex]
  \draw[thick] (-0.3,2.3)--(9.0,2.3) (-0.3,-0.1)--(9.0,-0.1);
  \foreach \x in {0.2,1.2,...,8.7}{\node at (\x,2.6) {$+$}; \node at (\x,-0.4) {$-$};}
  \draw[DeepBlue,very thick,domain=0:720,samples=220,smooth]
    plot ({0.012*\x-0.72*sin(\x)},{2.3-0.72*(1-cos(\x))});
  \foreach \x in {1.2,4.5,7.8}{\node[font=\large,TextGray] at (\x,1.25) {$\fieldInto$};}
  \draw[<->] (8.7,2.3)--node[right] {$h_{\max}$}(8.7,0.86);
  \node[TextGray] at (4.4,-0.9) {$h_{\max}=2mE/(|q|B^2)$；轨迹为沿 $\vect E\times\vect B$ 方向推进的摆线};
\end{tikzpicture}
\end{center}}

\newcommand{\gravityDriftDiagram}{%
\begin{center}
\begin{tikzpicture}[x=0.92cm,y=0.92cm,>=Latex]
  \fill[SoftGray] (-0.5,-0.35) rectangle (9.4,-0.75);
  \draw[thick] (-0.5,-0.35)--(9.4,-0.35) node[right] {地面};
  \draw[DeepBlue,very thick,domain=0:620,samples=210,smooth]
    plot ({8.0-0.013*\x+0.45*sin(\x)},{2.7-0.62*(1-cos(\x))-0.28*sin(\x)});
  \fill[DeepBlue] (8.0,2.7) circle (2.5pt) node[above] {初始位置};
  \draw[vector] (8.0,2.7)--(8.0,1.7) node[right] {$\vect v_0$};
  \draw[vector] (8.0,3.0)--(6.5,3.0) node[midway,above] {$\vect v_D$};
  \draw[<->] (8.75,-0.35)--node[right] {$H$}(8.75,2.7);
  \node[font=\large] at (5.2,2.9) {$\fieldOut$};
\end{tikzpicture}
\end{center}}

\newcommand{\alternatingLayersDiagram}{%
\begin{center}
\begin{tikzpicture}[x=1.0cm,y=0.82cm,>=Latex]
  \foreach \k in {0,2,4}{\fill[PaleBlue] (0,\k) rectangle (8,\k+1);}
  \foreach \k in {1,3,5}{\fill[SoftGray] (0,\k) rectangle (8,\k+1);}
  \foreach \k in {0,...,6}{\draw[RuleGray] (0,\k)--(8,\k);}
  \node[anchor=east] at (-0.2,0.5) {电场层};
  \node[anchor=east] at (-0.2,1.5) {第1磁场层};
  \node[anchor=east] at (-0.2,2.5) {电场层};
  \node[anchor=east] at (-0.2,3.5) {第2磁场层};
  \node[anchor=east] at (-0.2,4.5) {电场层};
  \node[anchor=east] at (-0.2,5.5) {第3磁场层};
  \foreach \y in {0.5,2.5,4.5}{\draw[vector] (7.4,\y-0.25)--(7.4,\y+0.28) node[above] {$\vect E$};}
  \foreach \y in {1.5,3.5,5.5}{\node[font=\large] at (7.4,\y) {$\fieldInto$};}
  \draw[DeepBlue,very thick] (1.0,0)--(1.0,1)
    .. controls (1.0,1.45) and (0.65,1.75) .. (0.55,2.0)
    --(0.55,3.0)
    .. controls (0.55,3.45) and (0.0,3.85) .. (0.15,4.0)
    --(0.15,5.0)
    .. controls (0.15,5.55) and (-0.15,5.75) .. (0.05,5.95);
  \draw[<->] (8.35,0)--node[right] {$d$}(8.35,1);
  \node[TextGray] at (4.0,-0.55) {穿过第 $k$ 个磁场层后，$v_x=-k\omega_{\rm c}d$};
\end{tikzpicture}
\end{center}}

\newcommand{\hallDiagram}{%
\begin{center}
\begin{tikzpicture}[scale=0.9,>=Latex]
  \coordinate (A) at (-3,-1); \coordinate (B) at (2,-1);
  \coordinate (C) at (3.5,0.2); \coordinate (D) at (-1.5,0.2);
  \coordinate (At) at (-3,1); \coordinate (Bt) at (2,1);
  \coordinate (Ct) at (3.5,2.2); \coordinate (Dt) at (-1.5,2.2);
  \fill[PaleBlue] (A)--(B)--(C)--(D)--cycle;
  \fill[SoftGray] (At)--(Bt)--(Ct)--(Dt)--cycle;
  \draw[thick] (A)--(B)--(C)--(D)--cycle (At)--(Bt)--(Ct)--(Dt)--cycle
    (A)--(At) (B)--(Bt) (C)--(Ct) (D)--(Dt);
  \draw[vector] (-4.2,0.5)--(-2.0,0.5) node[midway,above] {$I\parallel x$};
  \draw[vector] (0.2,-2.0)--(2.0,-0.55) node[midway,below right] {$\vect B\parallel y$};
  \draw[vector] (2.6,1.85)--(2.6,0.15) node[midway,right] {$\vect E_H\parallel -z$};
  \node[DeepBlue] at (-0.2,1.78) {$+$};
  \node[Ink] at (-0.2,-0.58) {$-$};
  \draw[<->] (-3.35,-1)--node[left] {$a$}(-3.35,1);
  \draw[<->] (-3,-1.35)--node[below] {$c$}(2,-1.35);
  \draw[<->] (2.35,-1.15)--node[below right] {$b$}(3.75,0.0);
\end{tikzpicture}

{\small\color{TextGray}霍尔电场抵消载流子的横向洛伦兹力；电势差的符号可判定载流子电性}
\end{center}}

\newcommand{\corbinoDiagram}{%
\begin{center}
\begin{tikzpicture}[scale=0.92,>=Latex]
  \fill[SoftGray,even odd rule] (0,0) circle (2.8) (0,0) circle (1.25);
  \draw[curve] (0,0) circle (2.8) (0,0) circle (1.25);
  \foreach \a in {0,45,...,315}{
    \draw[vector] ({1.35*cos(\a)},{1.35*sin(\a)})--({2.35*cos(\a)},{2.35*sin(\a)});
  }
  \draw[DeepBlue,thick,->] (2.0,0) arc[start angle=0,end angle=55,radius=2.0];
  \node[DeepBlue] at (2.7,1.55) {$j_\varphi$};
  \draw[<->] (0,0)--(1.25,0);
  \draw[<->] (0,0)--(-2.8,0);
  \node[above=2pt,inner sep=0.5pt] at (0.63,0) {$a$};
  \node[below=2pt,inner sep=0.5pt] at (-1.4,0) {$b$};
  \node[font=\Large] at (0,0) {$\fieldOut$};
  \node[TextGray,align=left,anchor=west] at (4.0,0.5)
    {图示电子载流子，$q_c<0$；\\无磁场：$\vect j\parallel\vect E_r$；\\有轴向磁场：电流产生方位分量，\\径向电导降为 $\sigma_0/(1+\mu^2B^2)$。};
\end{tikzpicture}
\end{center}}

\newcommand{\epicycleDiagram}{%
\begin{center}
\begin{tikzpicture}[scale=0.9,>=Latex]
  \draw[guide] (0,0) circle (2.2);
  \draw[curve] (0.45,0) circle (2.2);
  \fill[Ink] (0,0) circle (2pt) node[below left] {$O$};
  \fill[DeepBlue] (0.45,0) circle (2pt) node[below right] {$C$};
  \draw[<->] (0,0)--(0.45,0);
  \node[above=3pt,inner sep=0.5pt] at (0.225,0) {$\delta$};
  \draw[DeepBlue,thick,->] (2.65,0) arc[start angle=0,end angle=-70,radius=2.2];
  \node[TextGray,align=left,anchor=west] at (4.0,0.3)
    {$C$ 为移动后的导向中心；\\径向小扰动只把回旋圆心略微移开；\\相对原点的半径以 $\omega_{\rm c}$ 周期起伏。};
\end{tikzpicture}
\end{center}}

\newcommand{\radialFieldOrbitDiagram}{%
\begin{center}
\begin{tikzpicture}[scale=0.9,>=Latex]
  \draw[curve] (0,0) circle (1.8);
  \node[font=\Large] at (0,0) {$\fieldInto$};
  \node[TextGray,right] at (0.3,0) {$\vect B$};
  \fill[DeepBlue] (0,1.8) circle (2.6pt) node[above right] {$+q$};
  \draw[vector] (0,1.8)--(-1.25,1.8) node[left] {$\vect v$};
  \draw[vector] (0,1.8)--(0,0.45) node[right] {$\vect F_B$};
  \draw[vector] (0,1.8)--(0,2.65) node[right] {$\vect F_E$};
  \draw[guide] (-2.4,0)--(-.2,0) (-2.4,1.8)--(-.2,1.8);
  \draw[<->] (-2.4,0)--(-2.4,1.8) node[midway,left]{$R$};
  \draw[DeepBlue,thick,->] (-2.15,0.25) arc[start angle=175,end angle=245,radius=0.65];
  \node[TextGray,align=left,anchor=west] at (3.1,0.35)
    {$B$ 入纸面、$+q$ 逆时针时磁力向内；\\向外电场力抵消一部分磁向心力。};
\end{tikzpicture}
\end{center}}

\newcommand{\dampedSpiralDiagram}{%
\begin{center}
\begin{tikzpicture}[scale=0.9,>=Latex]
  \coordinate (S) at (-1.4,0);
  \coordinate (P) at (0.6,0);
  \draw[DeepBlue,very thick,domain=0:920,samples=260,smooth]
    plot ({0.6-2.0*exp(-0.0032*\x)*cos(\x)},{-2.0*exp(-0.0032*\x)*sin(\x)});
  \fill[DeepBlue] (S) circle (2.5pt) node[below] {初始位置};
  \fill[Ink] (P) circle (2.8pt) node[right] {$\vect r_\infty$};
  \draw[vector] (S)--++(1.35,0) node[above] {$\vect v_0$};
  \node[TextGray,align=left,anchor=west] at (3.2,0)
    {速度按 $e^{-\gamma t/m}$ 衰减；\\轨迹绕有限终点作圆形收敛螺旋。};
\end{tikzpicture}
\end{center}}

\newcommand{\dampedDriftDiagram}{%
\begin{center}
\begin{tikzpicture}[x=0.95cm,y=0.95cm,>=Latex]
  \draw[axis] (-0.3,0)--(9.0,0) node[right] {$x$};
  \draw[axis] (0,-1.2)--(0,2.4) node[above] {$y$};
  \draw[guide] (0,0)--(8.7,1.55);
  \draw[DeepBlue,very thick,domain=0:620,samples=240,smooth]
    plot ({0.013*\x+0.72*(1-exp(-0.006*\x))*sin(\x)},
          {0.0025*\x+0.75*exp(-0.006*\x)*sin(\x)});
  \draw[vector] (6.0,1.08)--(7.6,1.36) node[above] {$\vect v_\infty$};
  \node[TextGray] at (4.3,-1.55) {回旋瞬态指数衰减，长期运动趋于匀速斜向漂移};
\end{tikzpicture}
\end{center}}

\newcommand{\magneticMomentDiagram}{%
\begin{center}
\begin{tikzpicture}[scale=0.88,>=Latex]
  \begin{scope}[xshift=-4.2cm]
    \draw[curve] (0,0) circle (1.6);
    \fill[DeepBlue] (0,1.6) circle (2.6pt) node[above] {$q$};
    \draw[vector] (0,1.6)--(1.2,1.6) node[right] {$\vect v_\perp$};
    \node[font=\Large] at (-0.32,0.25) {$\fieldOut$};
    \node[font=\Large] at (0.32,-0.25) {$\fieldInto$};
    \node[TextGray,left] at (-0.62,0.25) {$\vect B$};
    \node[TextGray,right] at (0.62,-0.25) {$\vect\mu$};
    \node[TextGray] at (0,-2.25) {$\mu=mv_\perp^2/(2B)=E_\perp/B$};
  \end{scope}
  \begin{scope}[xshift=3.2cm]
    \draw[DeepBlue,thick] (-3,-2.0) .. controls (-1,-1.1) and (1,-1.1) .. (3,-2.0);
    \draw[DeepBlue,thick] (-3,2.0) .. controls (-1,1.1) and (1,1.1) .. (3,2.0);
    \foreach \y in {-1.4,-0.7,0,0.7,1.4}{\draw[->,TextGray] (-2.8,\y)--(2.8,{0.65*\y});}
    \draw[DeepBlue,very thick,domain=0:720,samples=180,smooth]
      plot ({-2.4+0.0067*\x},{0.55*sin(\x)*(1-0.00065*\x)});
    \draw[vector] (-2.3,0)--(-0.8,0) node[below] {$v_\parallel$};
    \node[TextGray,align=center] at (0,-2.55)
      {$\rho/L_B\ll1$ 时，$\mu$ 为绝热不变量；\\进入强场区时 $v_\perp$ 增大、$v_\parallel$ 减小。};
  \end{scope}
\end{tikzpicture}
\end{center}}

% All curves below share the initial conditions used in the adjacent solution.
\newcommand{\rTwoDataPath}{/Users/dongxingxi/Documents/物理竞赛一轮资料/sample/xes-original/thirty-eighth-section/assets/trajectory-data}
\newenvironment{rTwoFigure}{\par\noindent\begin{minipage}{\linewidth}\centering}{\end{minipage}\par\medskip}
\newcommand{\cyclotronExtractionDiagram}{\begin{rTwoFigure}
\begin{tikzpicture}[x=.9cm,y=.9cm,>=Latex]
 \begin{scope}[xshift=-3.4cm]
  \draw[thick](-.2,.7)--(-2,.7)--(-2,-.7)--(-.2,-.7);
  \draw[thick](.2,.7)--(2,.7)--(2,-.7)--(.2,-.7);
  \foreach \x in {-1.2,0,1.2}{\draw[vector](\x,1.3)--(\x,-1.3);}
  \draw[DeepBlue,densely dashed](-1.9,0)--(1.9,0);
  \node at(0,1.7){D 形电极的侧剖面};
  \node[align=center,text width=4.8cm,font=\footnotesize] at(0,-1.9){束流在中空腔内\\$\boldsymbol B$垂直轨道平面};
 \end{scope}
 \begin{scope}[xshift=3.2cm]
  \draw[guide](-2,0)..controls(-.6,0)and(0,.8)..(1,1.3);
  \draw[DeepBlue,thick,->](-2,0)..controls(-.5,0)and(.2,-.3)..(2,-.3);
  \draw[thick](-.9,-.8)--(.3,-.8);
  \node[below=3pt,font=\footnotesize]at(-.3,-.8){负偏转板};
  \draw[thin](1,-.55)--(1,.25);
  \draw[thin](1,.25)--(1,1.6)node[above]{薄铝窗};
  \draw[very thick](2,-.75)--(2,.2)node[above]{靶};
  \node[above] at(-1.1,.4){虚线：未引出时};
  \node[align=center,text width=4.8cm,font=\footnotesize] at(0,-1.9){实线：正离子被引向窗口\\再轰击靶};
 \end{scope}
\end{tikzpicture}
\end{rTwoFigure}}
\renewcommand{\magneticMomentDiagram}{\begin{rTwoFigure}
\begin{tikzpicture}[scale=.84,>=Latex]
\begin{scope}[xshift=-3cm]
 \draw[curve](0,0)circle(1.25);\fill(0,1.25)circle(2pt)node[above]{$+q$};
 \draw[vector](0,1.25)--(1,1.25)node[right]{$v_\perp$};
 \node at(-.2,.35){$\fieldOut\ B$};\node at(.2,-.35){$\fieldInto\ \mu$};
 \node at(0,-1.7){$\mu=E_\perp/B$};
\end{scope}
\begin{scope}[xshift=.8cm]
 \foreach \a in {-1.4,-.9,0,.9,1.4}{
  \draw[TextGray,->,domain=0:6,samples=60,variable=\t]plot(\t,{\a/sqrt(1+.15*\t)});
 }
 \draw[DeepBlue,thick]plot file {\rTwoDataPath/mirror.dat};
 \node[above]at(3,1.6){向右磁场逐渐增强};
 \node[align=center]at(3,-2){回旋半径减小，$v_\perp$增大\\$v_\parallel$减小；图为缓变近似投影};
\end{scope}
\end{tikzpicture}\end{rTwoFigure}}
\renewcommand{\gyroHelixDiagram}{\begin{rTwoFigure}
\begin{tikzpicture}[scale=.76,>=Latex]
\begin{scope}[xshift=-3cm]
 \draw[curve] (0,0) circle(1.4);
 \node at(0,0){$\fieldOut\ B$};
 \foreach \a in {90,210,330}{
  \fill[DeepBlue] (\a:1.4) circle(2pt);
  \draw[vector] (\a:1.4)--++({.65*sin(\a)},{-.65*cos(\a)});
  \draw[vector,TextGray] (\a:1.4)--(\a:.8);
 }
 \node[above] at(0,1.55){$q>0$};
 \node[below,align=center] at(0,-1.7){切向箭头：速度\\向心箭头：磁力};
\end{scope}
\begin{scope}[xshift=1cm]
 \draw[axis] (-.4,0)--(6.6,0) node[right]{$\boldsymbol B$};
 \draw[DeepBlue,thick,domain=0:1440,samples=450,variable=\t]
 plot({.0042*\t+.70*(cos(\t)-1)},{.78*sin(\t)});
 \fill[DeepBlue] (0,0) circle(2pt) node[below left]{$P$};
 \draw[vector] (0,0)--(1,0);
 \draw[thin] (.8,0)--(1.1,-.95)--(2.3,-.95) node[right]{$v_\parallel$};
 \draw[vector] (0,0)--(0,1.12) node[above]{$v_\perp$};
 \draw[<->] (0,-1.2)--(1.512,-1.2) node[midway,below]{$h$};
 \node at(3,-2){螺旋线的斜投影};
\end{scope}
\end{tikzpicture}
\end{rTwoFigure}}

\renewcommand{\driftDiagram}{\begin{rTwoFigure}
\begin{tikzpicture}[x=.70cm,y=.65cm,>=Latex]
 \foreach \row/\sgn/\name in {0/1/{质子},1/-1/{电子}}{
  \begin{scope}[yshift={-3.4*\row cm}]
   \draw[DeepBlue,thick,domain=0:720,samples=240,variable=\t]
    plot({.50*(pi*\t/180-sin(\t))},{\sgn*.50*(1-cos(\t))});
   \draw[->] (1.8,-1.4)--(3.5,-1.4);\node at(2.6,-1.8){$\boldsymbol v_E$};
   \node[anchor=east]at(-.3,0){\name};
  \end{scope}}
 \node at(3,2){$\boldsymbol E$向上，$\boldsymbol B$出纸面};
 \begin{scope}[xshift=7cm]
 \foreach \row/\sgn/\name in {0/1/{质子},1/-1/{电子}}{
  \begin{scope}[yshift={-3.4*\row cm}]
   \draw[DeepBlue,thick,domain=0:360,samples=200,variable=\t]
    plot({3-\sgn*.45*(pi*\t/180-sin(\t))},{-.45*(1-cos(\t))});
   \draw[->] ({3+\sgn*.7},-1.4)--({3-\sgn*.7},-1.4);
   \node at(3,-1.8){$\boldsymbol v_g$};
  \end{scope}}
 \node at(3,2){$\boldsymbol g$向下，$\boldsymbol B$出纸面};
 \end{scope}
\end{tikzpicture}
\par\smallskip{\small 各图分别归一化长度尺度，不比较回旋半径大小；电漂移同向，重力漂移随电荷改号而反向。}
\end{rTwoFigure}}

\renewcommand{\cyclotronDiagram}{\begin{rTwoFigure}
\begin{tikzpicture}[scale=.9,>=Latex]
 \filldraw[fill=PaleBlue,draw=DeepBlue] (-.12,-2) arc(-90:-270:2)--cycle;
 \filldraw[fill=SoftGray,draw=DeepBlue] (.12,-2) arc(-90:90:2)--cycle;
 \node at(-1.65,.9){D};\node at(1.65,.9){D};
 \node at(1.45,1.4){$\fieldOut\ B$};
 \def\rTwoY{0}
 \foreach \k in {1,...,8}{
  \pgfmathsetmacro{\rr}{.42*sqrt(\k)}
  \ifodd\k
   \draw[DeepBlue,thick] (-.12,\rTwoY)--(.12,\rTwoY);
   \draw[DeepBlue,thick] (.12,\rTwoY) arc(90:-90:\rr);
   \pgfmathparse{\rTwoY-2*\rr}\xdef\rTwoY{\pgfmathresult}
  \else
   \draw[DeepBlue,thick] (.12,\rTwoY)--(-.12,\rTwoY);
   \draw[DeepBlue,thick] (-.12,\rTwoY) arc(-90:-270:\rr);
   \pgfmathparse{\rTwoY+2*\rr}\xdef\rTwoY{\pgfmathresult}
  \fi
 }
 \fill (0,0) circle(2pt);
 \draw[->] (.12,.07)--(.43,-.10);
 \draw[thin] (-.12,-2)--(-.5,-2.55)--(-.5,-2.9);
 \draw[thin] (.12,-2)--(.5,-2.55)--(.5,-2.9);
 \node[draw,inner xsep=6pt,minimum height=.6cm,anchor=north] at(0,-2.9){高频振荡器};
 \node[anchor=west,align=left] at(2.7,.6){盒内：定半径圆弧\\[3pt]过缝：动能增加 $|q|\Delta U$\\[3pt]$r_k\propto\sqrt{k}$};
 \node[anchor=west,align=left] at(2.7,-1.9){窄缝用小间隔表示；\\实际加速轨迹连续。};
 \draw[thin] (0,2.15)--(0,2.5) node[above]{加速缝};
\end{tikzpicture}
\par\smallskip{\small 俯视图中磁场出纸面，正粒子顺时针弯转；半径只在过缝加速后增大。}
\end{rTwoFigure}}

\renewcommand{\splitProductsDiagram}{\begin{rTwoFigure}
\begin{tikzpicture}[scale=.9,>=Latex]
\begin{scope}[xshift=-3.6cm]
 \draw[DeepBlue,thick] (0,0) arc(90:330:1.3);
 \draw[Ink,dashed,thick] (0,0) arc(90:-30:1.3);
 \coordinate(P) at(0,0);\coordinate(C) at(0,-1.3);
 \coordinate(Q) at({sqrt(3)*1.3/2},{-1.5*1.3});
 \draw[guide] (P)--(C)--(Q);
 \fill(P) circle(2pt) node[above=12pt]{$P_0$};
 \fill(Q) circle(2pt) node[right]{$P_1$};
 \fill(C) circle(2pt) node[left]{$C_0$};
 \draw[vector] (P)--(-1.2,0) node[above]{$2v_0$};
 \draw[vector] (P)--(.6,0) node[above right]{$v_0$};
 \node at(-1.8,-1.2){$\fieldInto\ B$};
 \node[align=center] at(0,-3.5){实线：$+q,m$；虚线：$-q,2m$\\首段同圆反向，$\Delta t=2T/3$};
\end{scope}
\begin{scope}[xshift=3.1cm,yshift=-1.3cm,scale=.9]
 \foreach \k in {0,...,5}{
  \coordinate(P\k) at({1.3*cos(180+60*\k)},{1.3*sin(180+60*\k)});
  \fill(P\k) circle(2pt);
  \node at({1.65*cos(180+60*\k)},{1.65*sin(180+60*\k)}){$P_{\k}$};
 }
 \draw[DeepBlue,->] (P0)--(P1);\draw[guide] (P1)--(P2)--(P3)--(P4)--(P5)--(P0);
 \node[align=center] at(0,-2.5){六个碰撞点（连线不是轨迹）\\$P_6=P_0$，总周期 $4T$};
\end{scope}
\end{tikzpicture}
\end{rTwoFigure}}

\renewcommand{\beamFocusDiagram}{\begin{rTwoFigure}
\begin{tikzpicture}[x=.8cm,y=.8cm,>=Latex]
 \foreach \k/\phase/\lab in {0/270/{未聚焦},1/360/{首个近轴焦点}}{
 \begin{scope}[yshift={-3.8*\k cm}]
 \draw[axis] (0,0)--(8.8,0) node[right]{$z$};
 \draw[thick] (8,-1.4)--(8,1.4) node[above]{屏};
 \foreach \amp in {-.48,-.24,.24,.48}{
  \draw[DeepBlue,domain=0:8,samples=150,variable=\t] plot(\t,{\amp*(1-cos(\phase*\t/8))});
 }
 \fill(0,0) circle(2pt) node[below left]{$O$};
 \draw[<->] (0,-1.3)--(8,-1.3) node[midway,below]{$L$};
 \node[anchor=west] at(9,.6){\lab};
 \ifnum\k=0
 \draw[<->] (8.35,-.48)--(8.35,.48);
 \node[anchor=west] at(9,-.3){投影包络};
 \fi
 \end{scope}}
\end{tikzpicture}
\par\smallskip{\small 同一组近轴螺旋束在含轴平面上的投影；屏上最大径向位移需由两个横向分量合成。}
\end{rTwoFigure}}

\renewcommand{\magneticBilliardDiagram}{\begin{rTwoFigure}
\begin{tikzpicture}[scale=1.05,>=Latex]
 \def\R{2.3}\def\rc{1}\def\ang{30}
 \pgfmathsetmacro{\del}{2*atan2(cos(\ang),\R+sin(\ang))}
 \coordinate(O) at(0,0);\coordinate(P) at(0,-\R);
 \coordinate(C) at({-cos(\ang)},{-\R-sin(\ang)});
 \coordinate(Q) at({-\R*sin(\del)},{-\R*cos(\del)});
 \draw[curve] (O) circle(\R);
 \draw[DeepBlue,very thick] (P) arc(30:{150-\del}:1);
 \draw[guide] (P)--(O)--(Q) (C)--(P) (C)--(Q);
 \draw[vector] (P)--++(-.5,.866) node[above left]{$\boldsymbol v$};
 \fill(P) circle(2pt) node[below right]{$P$（小孔）};
 \fill(Q) circle(2pt) node[left]{$Q$};\fill(C) circle(2pt) node[below left]{$C$};
 \fill(O) circle(2pt) node[above]{$O$};
 \draw (0,-1.86) arc(90:120:.44) node[midway,above]{$\theta$};
 \draw (0,-.65) arc(-90:{-90-\del}:.65);
 \node at(-.26,-.89){$\delta$};
 \node[anchor=west,align=left] at(2.8,.4){$OP=OQ=R$\\$CP=CQ=r_{\rm c}$\\$CP\perp\boldsymbol v$\\[4pt]$\fieldInto\ B,\quad q>0$};
\end{tikzpicture}
\end{rTwoFigure}}

\renewcommand{\crossedFieldCycloidDiagram}{\begin{rTwoFigure}
\begin{tikzpicture}[x=.85cm,y=.85cm,>=Latex]
 \draw[thick] (-.3,2.3)--(9.5,2.3) (-.3,-.1)--(9.5,-.1);
 \foreach \x in {0,1,...,9}{\node at(\x,2.65){$+$};\node at(\x,-.4){$-$};}
 \draw[DeepBlue,very thick,domain=0:720,samples=300,variable=\t]
 plot({.72*(pi*\t/180-sin(\t))},{2.3-.72*(1-cos(\t))});
 \fill(0,2.3) circle(2pt);\node[above left] at(0,2.8){$v(0)=0$};
 \node at(4.5,1.5){$\fieldInto\ B$};
 \draw[<->] (9.3,2.3)--(9.3,.86) node[midway,right]{$h_{\max}$};
\end{tikzpicture}
\par\smallskip{\small 正电荷从上板静止释放；每个尖点处速度为零，$h_{\max}=2mE/(qB^2)$。}
\end{rTwoFigure}}

\renewcommand{\gravityDriftDiagram}{\begin{rTwoFigure}
\begin{tikzpicture}[x=1cm,y=1cm,>=Latex]
 \draw[thick] (-.2,0)--(8.8,0) node[right]{地面};
 \draw[DeepBlue,thick,domain=0:720,samples=300,variable=\t]
 plot({7+.3*(sin(\t)-pi*\t/180)+(cos(\t)-1)},{2+.3*(cos(\t)-1)-sin(\t)});
 \fill(7,2) circle(2pt) node[above right]{初始位置};
 \draw[vector] (7,2)--(7,1.3) node[right]{$v_0$};
 \draw[vector] (5.8,2.5)--(4.7,2.5) node[midway,above]{$v_D$};
 \draw[<->] (8.2,0)--(8.2,2) node[midway,right]{$H$};
 \node at(2.5,3.15){$q>0,\quad\fieldOut\ B$};
\end{tikzpicture}
\par\smallskip{\small 本图取不触地的一组参数；开始时切线竖直向下，回旋中心向左漂移。}
\end{rTwoFigure}}

\renewcommand{\alternatingLayersDiagram}{\begin{rTwoFigure}
\begin{tikzpicture}[x=.68cm,y=.68cm,>=Latex]
 \foreach \k in {0,2,4}{\fill[PaleBlue] (-1,\k) rectangle(8,\k+1);}
 \foreach \k in {1,3,5}{\fill[SoftGray] (-1,\k) rectangle(8,\k+1);}
 \foreach \k in {0,...,6}{\draw[RuleGray] (-1,\k)--(8,\k);}
 \draw[DeepBlue,very thick] plot file {\rTwoDataPath/layers.dat};
 \foreach \y in {.5,2.5,4.5}{\draw[vector](7,\y-.2)--(7,\y+.2);\node[right] at(7.2,\y){$E$};}
 \foreach \y/\n in {1.5/1,3.5/2,5.5/3}{\node at(7,\y){$\fieldInto$};\node[left]at(-1,\y){第\n 磁场层};}
 \draw[<->] (8.5,0)--(8.5,1) node[midway,right]{$d$};
\end{tikzpicture}
\par\smallskip{\small 图取 $\Lambda=2.5$，在第三磁场层内转向；水平速度不会在界面处突变。}
\end{rTwoFigure}}

\renewcommand{\dampedSpiralDiagram}{\begin{rTwoFigure}
\begin{tikzpicture}[scale=1.05,>=Latex]
 \draw[axis] (-1,0)--(3,0) node[right]{$x$};\draw[axis] (0,-.3)--(0,3.7) node[above]{$y$};
 \draw[DeepBlue,thick,domain=0:1260,samples=450,variable=\t]
 plot({2/(1+.22^2)*(.22+exp(-.22*pi*\t/180)*(-.22*cos(\t)+sin(\t)))},
 {2/(1+.22^2)*(1+exp(-.22*pi*\t/180)*(-.22*sin(\t)-cos(\t)))});
 \fill(0,0) circle(2pt) node[below left]{$O$};
 \fill({.44/(1+.22^2)},{2/(1+.22^2)}) circle(2pt);
 \draw[thin] ({.44/(1+.22^2)},{2/(1+.22^2)})--(2.8,2.1) node[right]{$\boldsymbol r_\infty$};
 \draw[vector] (0,0)--(1,0) node[below]{$v_0$};
 \node[anchor=west,align=left] at(3.8,1){$\fieldInto\ B,\quad q>0$\\初始速度沿 $+x$\\终点有 $y_\infty>0$};
\end{tikzpicture}
\end{rTwoFigure}}

\renewcommand{\dampedDriftDiagram}{\begin{rTwoFigure}
\begin{tikzpicture}[x=.7cm,y=.7cm,>=Latex]
 \pgfmathsetmacro{\ux}{1/(1+.25^2)}\pgfmathsetmacro{\uy}{.25/(1+.25^2)}
 \pgfmathsetmacro{\cr}{((3-\ux)*.25-\uy)/(1+.25^2)}
 \pgfmathsetmacro{\ci}{(-(3-\ux)-.25*\uy)/(1+.25^2)}
 \draw[axis] (-.3,0)--(12.6,0) node[right]{$x$};\draw[axis](0,-3)--(0,2)node[above]{$y$};
 \draw[guide] (0,{\ci-\uy*\cr/\ux})--(12,{\ci+\uy*(12-\cr)/\ux});
 \draw[DeepBlue,thick,domain=0:720,samples=350,variable=\t]
 plot({\ux*pi*\t/180+\cr*(1-exp(-.25*pi*\t/180)*cos(\t))-\ci*exp(-.25*pi*\t/180)*sin(\t)},
 {\uy*pi*\t/180+\ci*(1-exp(-.25*pi*\t/180)*cos(\t))+\cr*exp(-.25*pi*\t/180)*sin(\t)});
 \draw[vector](0,0)--(1,0)node[above]{$v_0$};
 \node[anchor=west]at(6,-2.5){$E\parallel+y,\quad\fieldOut\ B$};
 \draw[vector](9,.8)--(10.5,1.175)node[above]{$\boldsymbol v_\infty$};
\end{tikzpicture}
\par\smallskip{\small 虚线是长期直线渐近线；两个方向的回旋瞬态均指数衰减。}
\end{rTwoFigure}}


磁场与电场都是矢量场，但它们对点电荷的作用方式不同。静电场可以直接改变粒子的动能；磁场只对
运动电荷产生与速度垂直的力。这个“始终垂直”看似只是方向信息，实际上决定了回旋、螺旋、漂移、
速度选择、霍尔效应以及粒子探测器中几乎全部的轨迹结构。

人们很早就认识磁体，却曾长期把电与磁看成两类彼此独立的自然现象。电流磁效应和电磁感应相继被发现
以后，运动电荷、电流、磁场和变化电场才逐渐被纳入统一的电磁理论。这一统一不只是概念整理：发电机、
电动机、远距离输电和通信技术由此建立，成为第二次工业革命的重要物理基础。后面的讨论将从运动电荷
在磁场中的受力出发，再过渡到电流的受力、电流产生的磁场以及变化电磁场。

磁现象最直观的入口是磁铁：磁体有南、北两极，同名磁极相斥、异名磁极相吸。把磁铁置于铺有铁屑的
薄板下方，铁屑会沿局部磁场方向排列，显出一组从一极附近弯向另一极附近的曲线。图中每一点的切线
方向给出该处磁场方向；这些曲线只是磁场的可视化，不是空间中真实存在的细线。

\historyfigure[0.72\linewidth]{magnet-field-built-in-v1.png}%
{马蹄形磁铁上方铁屑形成的磁场分布；每条曲线在局部的切向给出磁场方向。}\label{ch22:fig:s38-7-1}


\subsection{磁场与洛伦兹力（Lorentz force）}


在给定惯性系中，电量为 $q$、速度为 $\vect v$ 的点粒子处在磁感应强度为 $\vect B$ 的磁场中，
所受磁力为
\begin{equation}
  \boxed{\vect F_B=q\vect v\times\vect B},
  \qquad F_B=|q|vB\sin\theta .
\end{equation}
方向先按右手定则判断 $\vect v\times\vect B$，再由电荷正负决定是否反向。磁力同时垂直于
$\vect v$ 与 $\vect B$。因此
\begin{equation}
  \vect F_B\cdot\vect v=0,
  \qquad \frac{\dd}{\dd t}\left(\frac12mv^2\right)=0.
\end{equation}
纯磁场不改变速率与动能，只改变速度方向；所谓“磁场加速粒子”通常是磁场负责弯转，真正增加动能的
仍是电场。

\lorentzDirectionDiagram

国际单位制中 $B$ 的单位是特斯拉（T），磁通量
\begin{equation}
  \Phi_B=\int_S\vect B\cdot\dd\vect S
\end{equation}
的单位是韦伯（Wb），且 $\SI{1}{T}=\SI{1}{Wb.m^{-2}}$。常用换算为
$\SI{1}{T}=10^4\,\mathrm G$。高斯在一些工程和天体物理资料中仍很常见，因为日常磁场若用特斯拉表示
往往是很小的数。永久磁体可产生数特斯拉量级的磁场，脉冲装置能够短时达到更高数值；中子星表面则可
达到远超实验室磁场的极端量级。因此，讨论磁场强弱时还必须同时说明它是稳态场、脉冲场还是天体环境。


% EM teaching enrichment 20260927
\paragraph{把右手判断写成分量核对}
若 $\boldsymbol v=v\boldsymbol e_x$、$\boldsymbol B=B\boldsymbol e_z$，则
$\boldsymbol v\times\boldsymbol B=-vB\boldsymbol e_y$。正电荷向 $-y$ 偏，电子向 $+y$ 偏。
同时有 $\dd K/\dd t=q\boldsymbol E\cdot\boldsymbol v$：即使电、磁场同时存在，瞬时改变动能的仍只有电场项。轨迹弯曲只能说明存在横向加速度，不能单独证明粒子能量增加。


\subsection{匀强磁场中的回旋与螺旋运动}


把速度分解为平行、垂直于磁场的两部分：
\begin{equation}
  \vect v=\vect v_\parallel+\vect v_\perp .
\end{equation}
$\vect v_\parallel$ 不受磁力影响而保持不变；$\vect v_\perp$ 在垂直于磁场的平面内作匀速圆周运动。
由 $|q|v_\perp B=mv_\perp^2/r_{\rm c}$ 得
\begin{equation}
  \boxed{r_{\rm c}=\frac{mv_\perp}{|q|B}=\frac{p_\perp}{|q|B}},
  \qquad
  \boxed{\omega_{\rm c}=\frac{|q|B}{m}},
\end{equation}
\begin{equation}
  \boxed{T_{\rm c}=\frac{2\pi m}{|q|B}},
  \qquad
  f_{\rm c}=\frac{|q|B}{2\pi m}.
\end{equation}
经典非相对论范围内，回旋角频率与速率无关。若 $v_\parallel\ne0$，圆周运动同时沿磁场方向平移，
形成螺旋线；每一圈沿轴线前进的螺距为
\begin{equation}
  h=v_\parallel T_{\rm c}=\frac{2\pi m v_\parallel}{|q|B}.
\end{equation}

\gyroHelixDiagram

轨迹曲率还能反推出粒子动量 $p_\perp=|q|Br_{\rm c}$。若粒子穿过物质后损失能量，曲率半径减小；
结合弯曲方向、磁场方向和已知运动方向，还可判定电荷符号。这是云雾室、气泡室以及现代径迹探测器的
基本判读逻辑。

1932 年 Anderson 在置于磁场中的云雾室内加入一块铅板：粒子穿过铅板会损失能量，因而板后轨迹的
曲率变大。这样不仅能由弯曲方向判断电荷符号，还能由板前、板后的半径变化辨认粒子的运动方向。照片中
出现的粒子带正电，但射程和电离能力又不像质子，于是被认定为狄拉克预言的电子反粒子——正电子。

气泡室把同一思想推广到更高能粒子反应。已知外磁场后，轨迹的切向给出速度方向，曲率给出横向动量，
顶点位置显示粒子在哪里发生散射或衰变。早期分析需要人工检查海量照片；欧洲核子研究中心确认弱中性流
时，候选事例就是从约一百四十万张照片中筛出的少数事件。格拉泽因发明气泡室获得 1960 年诺贝尔
物理学奖。这些历史事例说明图像并非装饰，而是把电荷、动量和相互作用类型转化为可测量证据的仪器。

\particleTracksContext


% EM teaching enrichment 20260927
\paragraph{螺距与半径测出的是不同速度分量}
给定同一粒子的轨迹半径 $r_c$、螺距 $h$ 和场强 $B$，可依次求
\[
 v_\perp=\frac{|q|Br_c}{m},\qquad
 |v_\parallel|=\frac{|q|Bh}{2\pi m},\qquad
 \tan\alpha=\frac{2\pi r_c}{h},
\]
其中 $\alpha$ 是速度与磁场轴的锐夹角，$h$ 取每圈轴向位移的大小。总动能由两分量平方和得到，而不是用螺旋半径反推总速率。
$h\to0$ 回到平面圆周；$r_c\to0$ 回到沿场直线。轨迹的几何螺旋若没有运动方向信息，仍不足以独自确定电荷正负。


\subsection{回旋方向与运动方程的完整解}


仅记半径和周期还不够，许多题目需要同时判断圆心位置、转向和离开磁场的边界点。设匀强磁场
$\vect B=B\uvec z$，则横向分量满足
\begin{equation}
 m\dot v_x=qBv_y,
 \qquad
 m\dot v_y=-qBv_x.
\end{equation}
定义带符号回旋角频率 $\Omega=qB/m$，解可写为
\begin{align}
 v_x(t)&=v_{x0}\cos\Omega t+v_{y0}\sin\Omega t,\notag\\
 v_y(t)&=v_{y0}\cos\Omega t-v_{x0}\sin\Omega t.
\end{align}
正电荷在 $B$ 指向纸面外时顺时针转动，负电荷方向相反。积分一次得到导向中心
\begin{equation}
 X=x+\frac{mv_y}{qB},
 \qquad
 Y=y-\frac{mv_x}{qB},
\end{equation}
在纯匀强磁场中保持不变。由初始位置和初速度先求 $(X,Y)$，再画半径
$r_c=mv_\perp/(|q|B)$ 的圆，比凭弯曲方向猜圆心更可靠。

粒子从有界磁场区域射出时，轨迹只取完整圆的一段。圆弧所对圆心角 $\Delta\phi$ 同时决定飞行
时间和弦长：
\begin{equation}
 t=\frac{\Delta\phi}{\omega_c},
 \qquad
 \ell_{\rm chord}=2r_c\sin\frac{\Delta\phi}{2}.
\end{equation}
处理矩形、圆形或扇形磁场边界时，应先用几何求圆心角，再求时间，不能把穿越时间简单写成
直线长度除以速率。


% EM teaching enrichment 20260927 supplementary

\paragraph{先定圆心，再选真正的出射交点}
例如 $q>0$，粒子从原点以 $v_0\uvec x$ 进入 $\vect B=B\uvec z$，有圆心 $(0,-r_c)$，
\[
 x=r_c\sin\Omega t,\qquad y=r_c(\cos\Omega t-1).
\]
若磁场为条带 $0\le x\le a$ 且 $0<a<r_c$，首次碰到远边界的时间为
$t=\arcsin(a/r_c)/\Omega$，此时 $v_y<0$。同一个圆与 $x=a$ 的另一个交点对应更晚时刻，
粒子早已离开磁场，不能继续沿原圆推算。若 $a>r_c$，轨迹到不了远边界，而在半周后返回 $x=0$。
边界问题的核心是把代数交点按时间排序，并逐一检查是否仍处在磁场内部。


\subsection{经典公式的相对论边界}


速度接近光速时，动量应写成 $\vect p=\gamma m\vect v$。回旋半径仍满足
\begin{equation}
 r_c=\frac{p_\perp}{|q|B},
\end{equation}
但角频率变为
\begin{equation}
 \omega_c=\frac{|q|B}{\gamma m}.
\end{equation}
随着能量升高，$\gamma$ 增大，固定射频的经典回旋加速器逐渐失谐。同步回旋加速器通过改变频率，
同步加速器则同时调节磁场和射频，使高能粒子保持在设计轨道上。竞赛题若未说明相对论效应，通常
默认 $v\ll c$；若给出的能量与静能 $mc^2$ 同量级，就必须改用相对论动量和能量关系。

% EM teaching enrichment 20260927
\paragraph{以动能为已知量时的连续表达}
总能量 $\mathcal E=K+mc^2$ 与动量满足 $\mathcal E^2=p^2c^2+m^2c^4$，故垂直入射时
\[
 r_c=\frac{\sqrt{K(K+2mc^2)}}{|q|Bc}.
\]
$K\ll mc^2$ 时恢复 $\sqrt{2mK}/(|q|B)$；$K\gg mc^2$ 时近似 $K/(|q|Bc)$。低能区半径随 $\sqrt K$ 增大，高能区近似随 $K$ 增大，是判断题中经典近似能否继续使用的直观尺度检验。



% EM structure and teaching revision 20260928
\subsection{从实验数据判断模型范围}
粒子垂直磁场运动时，测得曲率半径直接给出 $p=|q|Br$。低速区用 $K=p^2/(2m)$，而相对论关系为
\[
 K=\sqrt{m^2c^4+p^2c^2}-mc^2.
\]
把根式按 $p/(mc)$ 展开，首项才是经典动能。因此“半径公式仍成立”不代表经典能量和周期公式仍成立；高能粒子的回旋角频率为 $|q|B/(\gamma m)$。
若粒子由电压差大小 $U$ 加速且初始静止，经典模型先给出 $K=|q|U$，再检查 $|q|U\ll mc^2$。对于电子，静能约为 $511\,\mathrm{keV}$：几千电子伏特的加速通常可作低速近似，而达到数百千电子伏特时不能再忽略相对论修正。
有界磁场问题还要检查几何可达性。先确定圆心、半径和真实绕行方向，再比较候选交点的到达时间；切线接触边界与穿越边界不同。若粒子在到达某个代数交点之前已离开磁区，后续轨迹就必须换用新区的运动方程。


\subsection{专题例题与方法应用}

\begin{origexamplebox}{【基本计算】}\label{ch22:ex:s38-7-1}
质量为 $m$、电量为 $q$ 的粒子由静止经过电压 $U$ 加速，随后垂直注入磁感应强度为 $B$ 的匀强
磁场。求其回旋半径，忽略相对论效应。
\end{origexamplebox}\par\nopagebreak[4]

\begin{xesolution}
加速电场所做的功转化为动能。用 $Q=|q|$ 表示电荷量的大小，则
\[
  QU=\frac{p^2}{2m},
  \qquad p=\sqrt{2mQU}.
\]
进入磁场后 $p=QBr$，所以
\begin{equation}
  \boxed{r=\frac{\sqrt{2mQU}}{QB}
  =\sqrt{\frac{2mU}{QB^2}}}.
\end{equation}
式中使用的是使粒子获得能量的电势差大小；电荷正负只改变弯曲方向，不改变半径。
\end{xesolution}



\begin{origexamplebox}{【比例关系】}\label{ch22:ex:s38-7-2}
电子、质子、$\alpha$ 粒子都经过电压 $U$ 加速，再垂直注入同一匀强磁场 $B$。求三者的回旋半径
与周期之比。
\end{origexamplebox}\par\nopagebreak[4]

\begin{xesolution}
相同加速电压下
\[
  r=\sqrt{\frac{2mU}{QB^2}}\propto\sqrt{\frac mQ},
  \qquad T=\frac{2\pi m}{QB}\propto\frac mQ.
\]
因此
\begin{equation}
  \boxed{r_e:r_p:r_\alpha
  =\sqrt{m_e}:\sqrt{m_p}:\sqrt{\frac{m_\alpha}{2}}},
\end{equation}
\begin{equation}
  \boxed{T_e:T_p:T_\alpha=m_e:m_p:\frac{m_\alpha}{2}}.
\end{equation}
若取 $m_\alpha\simeq4m_p$，则
$r_e:r_p:r_\alpha=\sqrt{m_e}:\sqrt{m_p}:\sqrt{2m_p}$，
$T_e:T_p:T_\alpha=m_e:m_p:2m_p$。
\end{xesolution}



\begin{origexamplebox}{【统统算出来】}\label{ch22:ex:s38-7-3}
质量为 $m$、电量为 $q$ 的粒子以速率 $v$ 垂直进入匀强磁场 $B$。忽略重力、电磁辐射与相对论
效应，求周期和频率、动能和动量、关于回旋圆心的角动量、等效电流、轨道所围磁通量以及轨道磁矩。
\end{origexamplebox}\par\nopagebreak[4]

\begin{xesolution}
记 $Q=|q|$。由 $r=mv/(QB)$、$\omega_{\rm c}=QB/m$ 依次得到
\begin{align}
  T&=\frac{2\pi m}{QB},
  & f&=\frac{QB}{2\pi m},\\
  E_k&=\frac12mv^2,
  & p&=mv,\\
  L&=mvr=\frac{m^2v^2}{QB},
  & |I|&=\frac{Q}{T}=\frac{Q^2B}{2\pi m},\\
  \Phi_B&=B\pi r^2=\frac{\pi m^2v^2}{Q^2B},
  & |\mu|&=|I|\pi r^2=\frac{mv^2}{2B}=\frac{E_k}{B}.
\end{align}
方向关系为
\begin{equation}
  \boxed{\vect\mu=\frac{q}{2m}\vect L}.
\end{equation}
对由磁力维持的回旋运动，$\vect\mu$ 总与 $\vect B$ 反向。磁通量与角动量的大小还满足
$\Phi_B=\pi L/Q$。
\end{xesolution}



\begin{origexamplebox}{【磁场中的裂变产物】}\label{ch22:ex:s38-7-4}
匀强磁场 $\vect B$ 垂直纸面向里。一个静止的中性粒子裂变成两个粒子：质量分别为
$m,2m$，电量分别为 $+q,-q$，前者以 $2v_0$ 向左、后者以 $v_0$ 向右运动。二者相碰时作完全弹性碰撞且不
交换电量。定量描述后续运动。

\begin{center}
\begin{tikzpicture}[>=Latex,scale=.85]
 \fill[DeepBlue] (-.25,0) circle(2.5pt);\fill[Ink] (.25,0) circle(2.5pt);
 \draw[vector] (-.4,0)--(-2.7,0) node[midway,above]{$2v_0$};
 \draw[vector] (.4,0)--(1.55,0) node[midway,above]{$v_0$};
 \node[below] at (-.8,-.3){$m,+q$};\node[below] at (.8,-.3){$2m,-q$};
 \node at (3,1){$\otimes\ B$};
 \node[font=\small] at (0,-1.1){同一裂变点，两粒子仅为辨认而稍分开};
\end{tikzpicture}
\end{center}
\figcap{裂变后的瞬时速度、质量与电荷。}\label{ch22:fig:em22-question-4}
\end{origexamplebox}\par\nopagebreak[4]

\begin{xesolution}
\textbf{初态与第一次相遇。} 取向右为 $+x$、向上为 $+y$，裂变点为 $P_0$。
初始总动量 $m(-2v_0)+2m(v_0)=0$，与静止母粒子相符。两个磁力均向下，回旋半径分别为
\[
 r_1=\frac{m(2v_0)}{qB},\qquad r_2=\frac{2mv_0}{qB},\qquad
 r_1=r_2\equiv r,\quad \omega=\frac{qB}{m},\quad T=\frac{2\pi}{\omega}.
\]
因此图中两粒子从同一点沿同一个圆反向运动，正粒子角速度为 $+\omega$，负粒子为 $-\omega/2$。
首次相遇需相对转角为 $2\pi$，即
\[
 (\omega+\omega/2)\Delta t=2\pi,\qquad \boxed{\Delta t=2T/3}.
\]
此时正粒子逆时针转 $4\pi/3$，负粒子顺时针转 $2\pi/3$，到达图中的 $P_1$。

\textbf{碰撞。} 两速度沿同一圆的切线反向，大小仍为 $2v_0,v_0$。碰撞时间极短，磁场冲量可忽略。
以正粒子碰前速度方向为正，一维完全弹性碰撞公式给出
\[
 u_1'=\frac{m-2m}{3m}(2v_0)+\frac{4m}{3m}(-v_0)=-2v_0,
 \qquad
 u_2'=\frac{2m}{3m}(2v_0)+\frac{m}{3m}(-v_0)=v_0.
\]
两粒子各自反向而不改速率，总动量仍为零、总动能仍为 $3mv_0^2$。
碰后共用的新圆心为 $C_1=2P_1-C_0$：半径大小不变，但磁力因速度反向而反向。

\splitProductsDiagram

\textbf{重复过程。} 每次都在同一个半径为 $r$ 的圆上反向追遇，故相邻碰撞间隔始终为 $2T/3$。
令 $\mathcal R(\beta)$ 为平面逆时针转角 $\beta$ 的旋转，取 $P_0=(0,0)$，则可递推写出全部碰撞点与
第 $k$ 次碰撞后的速度（$k=0$表示裂变后）：
\[
 \boldsymbol v_{1,k}=\mathcal R(k\pi/3)(-2v_0,0),\qquad
 \boldsymbol v_{2,k}=-\boldsymbol v_{1,k}/2,
\]
\[
 P_{k+1}-P_k=\mathcal R(k\pi/3)
 \left(\frac{\sqrt3 r}{2},-\frac{3r}{2}\right).
\]
后一个式子来自首段转过 $4\pi/3$ 后的弦向量；下一段出发方向比上一段转过
$4\pi/3+\pi=7\pi/3$，模 $2\pi$ 即 $\pi/3$。
因此六条等长弦依次转 $60^\circ$，六个碰撞位置构成正六边形；六次碰撞后位置和两粒子的速度同时恢复。
完整周期为 $6(2T/3)=\boxed{4T}$。这个 $4T$ 来自六段同半径轨道，不能理解成两组内切圆之间往返。
任一段的连续轨迹也已确定：令 $J(x,y)=(-y,x)$，$C_k=P_k+J\boldsymbol v_{1,k}/\omega$，则在
$0\le s\le2T/3$ 内，两粒子位置分别为
\[
 \boldsymbol r_1=C_k+\mathcal R(\omega s)(P_k-C_k),\qquad
 \boldsymbol r_2=C_k+\mathcal R(-\omega s/2)(P_k-C_k).
\]

\end{xesolution}



\begin{origexamplebox}{【自聚焦】}\label{ch22:ex:s38-7-5}
质量为 $m$、电量大小为 $Q$ 的粒子经电压 $U$ 加速后大致沿 $z$ 方向射出，方向与 $z$ 轴的最大
夹角为 $\theta\ll1$，出射口在 $z=0$。在 $z=L$ 放置荧光屏：
\begin{enumerate}[label=(\arabic*),leftmargin=2.2em]
  \item 无磁场时亮斑多大？
  \item 加沿 $z$ 方向的匀强磁场 $B$ 后亮斑多大？怎样选择 $B$ 使亮斑最小？
\end{enumerate}
\end{origexamplebox}\par\nopagebreak[4]

\begin{xesolution}
加速后的速率
\[
  v=\sqrt{\frac{2QU}{m}}.
\]
无磁场时，最大横向速度为 $v\theta$，飞行时间为 $L/v$，故亮斑半径
\begin{equation}
  \boxed{\rho_0=L\theta}
\end{equation}
（若题目把“大小”理解为直径，则再乘 $2$）。

加纵向磁场后，横向速度以 $\omega_{\rm c}=QB/m$ 转动，而纵向速度近似仍为 $v$。设初始横向
速度为复数 $u_\perp(0)$，则
\[
  \Delta r_\perp(t)=\int_0^t u_\perp(0)\e^{\mathrm i\omega_{\rm c}t'}\dd t',
\]
其大小为
\[
  |\Delta r_\perp|=\frac{2u_\perp(0)}{\omega_{\rm c}}
  \left|\sin\frac{\omega_{\rm c}t}{2}\right|.
\]
取 $u_\perp(0)\le v\theta$、$t=L/v$，亮斑半径为
\begin{equation}
  \boxed{\rho_B=\frac{2mv\theta}{QB}
  \left|\sin\frac{QBL}{2mv}\right|}.
\end{equation}
若
\begin{equation}
  \frac{QBL}{mv}=2\pi n,
  \qquad
  \boxed{B_n=\frac{2\pi n m v}{QL}},\quad n=1,2,\ldots,
\end{equation}
所有近轴粒子的横向速度恰转整数圈，屏上重新聚焦；最小非零磁场取 $n=1$。

\beamFocusDiagram
\end{xesolution}



\begin{origexamplebox}{【请君入瓮】}\label{ch22:ex:s38-7-6}
半径为 $R$ 的圆柱形桶内有垂直横截面的匀强磁场 $B$。质量为 $m$、电量大小为 $Q$ 的粒子从桶壁
小孔以速率 $v$ 入射，和桶壁均作完全弹性碰撞。先讨论沿半径入射，再讨论入射方向与向内半径成
$\theta$ 的情形：粒子何时能从原小孔出去？

\begin{center}
\begin{tikzpicture}[>=Latex,scale=.75]
 \foreach \x/\ang/\lab in {0/90/径向入射,6/120/斜入射}{
  \begin{scope}[shift={(\x,0)}]
   \draw[thick] (272:1.7) arc[start angle=272,end angle=628,radius=1.7];
   \fill (0,0) circle(1.5pt) node[above left]{$O$};
   \draw[guide] (0,-1.7)--(0,0);\node[right=2pt] at (0,-.95){$R$};
   \draw[vector] (0,-1.7)--++(\ang:1.45) node[pos=.6,left=3pt]{$\boldsymbol v$};
   \node at (.9,.7){$\otimes B$};\node[below] at (0,-1.85){\lab};
   \ifnum\ang=120
    \draw[thin] (0,-.92) arc[start angle=90,end angle=120,radius=.78];
    \node at (-.35,-.55){$\theta$};
   \fi
  \end{scope}}
\end{tikzpicture}
\end{center}
\figcap{从桶壁小孔径向入射或以角度 $\theta$ 斜入射。}\label{ch22:fig:em22-question-6}
\end{origexamplebox}\par\nopagebreak[4]

\begin{xesolution}
粒子的回旋半径为 $r_{\rm c}=mv/(QB)$。圆壁反射只改变速度的法向分量，速率和入射角不变；圆形
边界又有转动对称性，所以相邻两次碰壁点在桶心所张的圆心角恒为 $\delta$。粒子能够回到小孔的
充要条件是存在正整数 $N,k$ 使
\begin{equation}
  N\delta=2\pi k,
  \qquad\text{即}\qquad \boxed{\frac{\delta}{2\pi}\in\mathbb Q}.
\end{equation}

沿半径入射时，桶心、回旋圆心和碰撞点构成直角三角形，几何上
\begin{equation}
  \delta=2\arctan\frac{r_{\rm c}}R.
\end{equation}
故可写成
\begin{equation}
  \boxed{\frac{r_{\rm c}}R=\tan\frac{k\pi}{N}},
  \qquad 0<k<\frac N2.
\end{equation}

一般入射时，把 $\theta$ 规定为从向内半径朝洛伦兹偏转侧量取的有向角。两圆交点几何给出
\begin{equation}
  \boxed{\tan\frac\delta2=
  \frac{r_{\rm c}\cos\theta}{R+r_{\rm c}\sin\theta}}.
\end{equation}
具体地，取小孔 $P=(0,-R)$，正电荷、磁场入纸面，速度方向为 $(-\sin\theta,\cos\theta)$。
图中回旋圆心为 $C=(-r_{\rm c}\cos\theta,-R-r_{\rm c}\sin\theta)$。
另一个交点记为 $Q=(-R\sin\delta,-R\cos\delta)$。用 $|Q-C|^2=|P-C|^2$ 消去半径，得
\[
 r_{\rm c}\cos\theta\sin\delta
 =(R+r_{\rm c}\sin\theta)(1-\cos\delta).
\]
再用半角公式即得上式。对 $|\theta|<\pi/2$，用
$\delta=2\operatorname{atan2}(r_{\rm c}\cos\theta,R+r_{\rm c}\sin\theta)$ 选取 $0<\delta<2\pi$，
避免分母为负时丢失角度分支。图中虚线 $OP,OQ$ 正是定义圆心角的两条半径。
若从另一侧倾斜入射，只需令 $\theta\to-\theta$。把由此算出的 $\delta$ 代入
$\delta/(2\pi)\in\mathbb Q$，便得到一般逃出条件。若该比为无理数，碰撞点会在圆周上稠密分布，
但在理想的零尺寸小孔模型中永不精确返回。

\magneticBilliardDiagram
\end{xesolution}



\section{交叉场、漂移与载流子输运}

\subsection{匀强横向力场中的导向中心漂移}


粒子除磁力外再受与 $\vect B$ 垂直的恒力 $\vect F$：
\begin{equation}
  m\frac{\dd\vect v}{\dd t}=q\vect v\times\vect B+\vect F.
\end{equation}
令 $\vect v=\vect v_D+\vect u$，选取恒定的 $\vect v_D$ 使
$q\vect v_D\times\vect B+\vect F=0$，则 $\vect u$ 仍满足纯回旋运动方程。由此
\begin{equation}
  \boxed{\vect v_D=\frac{\vect F\times\vect B}{qB^2}}.
\end{equation}
若 $\vect F=q\vect E$，得到与电荷符号、质量都无关的
\begin{equation}
  \boxed{\vect v_E=\frac{\vect E\times\vect B}{B^2}};
\end{equation}
若 $\vect F=m\vect g$，则
\begin{equation}
  \boxed{\vect v_g=\frac{m\vect g\times\vect B}{qB^2}},
\end{equation}
正、负电荷的重力漂移方向相反，因而可能造成电荷分离。

\driftDiagram


% EM teaching enrichment 20260927 supplementary

\paragraph{漂移公式的方向与适用条件}
由 $q\vect v_D\times\vect B+\vect F=0$ 出发，再与 $\vect B$ 作叉乘。
因为选择 $\vect v_D\perp\vect B$ 且 $\vect F\perp\vect B$，三重积给出
$-qB^2\vect v_D+\vect F\times\vect B=0$，从而得到上述漂移速度。
若 $\vect B=B\uvec z$、$\vect F=F\uvec x$，则
\[
 \vect v_D=-\frac{F}{qB}\uvec y.
\]
正、负电荷沿相反方向漂移；换成电力 $\vect F=q\vect E$ 后，电荷符号消去，漂移方向相同。
在漂移参考系中只剩纯磁场回旋，因此实验室轨迹是圆周运动与匀速平移的叠加，并非沿外力方向持续加速。
若外力含沿磁场分量，该分量仍使粒子沿场加速，须另列纵向方程；外力或磁场随空间显著变化时，这个恒定漂移分解不再直接成立。


\subsection{交叉电磁场、速度选择与漂移}


若 $\vect E\perp\vect B$，满足
\begin{equation}
 \vect E+\vect v\times\vect B=\vect0
\end{equation}
的粒子不发生横向偏转。速度选择器通常令粒子束、$\vect E$、$\vect B$ 两两垂直，选出的速率为
\begin{equation}
 \boxed{v=\frac EB}.
\end{equation}
方向必须使电力与磁力相反；改变电荷正负会同时反转两种力，因而同一速度方向的正、负电荷都可
满足无偏转条件。

对任意初速度，运动可分成以 $\vect v_E=\vect E\times\vect B/B^2$ 匀速移动的导向中心和绕
导向中心的回旋。因为 $\vect v_E$ 与 $q,m$ 无关，理想均匀交叉场不会靠这一漂移分离正负电荷。
若外力不是 $q\vect E$，一般漂移
\begin{equation}
 \vect v_D=\frac{\vect F\times\vect B}{qB^2}
\end{equation}
会随电荷符号改变，才可能造成电荷分离。

霍尔效应正是磁力造成横向电荷积累，积累建立的霍尔电场最终与磁力平衡。若载流子漂移速度为
$\vect v_d$，稳态有
\begin{equation}
 q(\vect E_H+\vect v_d\times\vect B)=0.
\end{equation}
对宽度 $w$、厚度 $d$ 的矩形导体，电流 $I=n|q|v_dwd$，霍尔电压大小为
\begin{equation}
 \boxed{|U_H|=\frac{IB}{n|q|d}}.
\end{equation}
霍尔电压极性可以判断主要载流子的电荷符号，大小则可测载流子浓度。


% EM teaching enrichment 20260927
\paragraph{选择速度之后如何测质量}
先用交叉场选出 $v=E/B_s$，再进入只有垂直匀强磁场 $B_a$ 的分析区，半径满足
\[
 r=\frac{mv}{|q|B_a},\qquad
 \frac m{|q|}=\frac{rB_aB_s}{E}.
\]
这里选速区与分析区的磁场分别记为 $B_s,B_a$，不能未经题设认定为同一个场强。测得半径给出质荷比；若电荷态未知，仅靠这一组数据不能唯一确定质量。同速且同电荷态的同位素才可由半径直接区分。


% EM structure and teaching revision 20260928
\subsection{霍尔平衡与载流子浓度}
稳定载流子横向不再加速时，$q(\boldsymbol E_H+\boldsymbol v_d\times\boldsymbol B)=0$。电荷符号在这一步约去，但把漂移速度换为常规电流 $I=nqv_dwd$ 时，符号又通过 $q$ 进入测量关系。
取电流沿 $+x$、磁场沿 $+z$、宽度沿 $y$，定义 $U_H=V(-w/2)-V(w/2)$，有
\[
 E_y=v_dB,\qquad U_H=E_yw=\frac{IB}{nqd}.
\]
宽度消去，厚度 $d$ 留在分母。测出电压的符号可以判断单一载流子模型中的电荷符号，测出大小可以求浓度。
例如 $I=1\,\mathrm A$、$B=0.50\,\mathrm T$、$d=0.10\,\mathrm{mm}$，若 $|U_H|=1.0\,\mathrm{mV}$、$|q|=e$，则 $n\simeq3.1\times10^{25}\,\mathrm{m^{-3}}$。把毫米误作米会造成三个数量级的错误。
普通霍尔条形样品可在侧边积累电荷，圆环几何却没有同样的开路侧边；科尔比诺圆盘必须同时求径向和环向的运动。两种装置不能仅靠替换样品宽度互相套用公式。


\subsection{专题例题与方法应用}

\begin{origexamplebox}{【速度选择器】}\label{ch22:ex:s38-7-7}
一束 $\alpha$ 粒子向右射出，速率各不相同。设计装置只让速率为 $v$ 的粒子沿直线通过。
\end{origexamplebox}\par\nopagebreak[4]

\begin{xesolution}
让匀强电场、匀强磁场和粒子速度两两垂直，并使电场力与磁力反向。直线通过的条件为
\begin{equation}
  Q E=QvB,
  \qquad \boxed{v=\frac EB}.
\end{equation}
速率偏大的粒子磁力占优，偏小的粒子电场力占优，二者向相反方向偏转。该关系与粒子质量、电量大小
无关；若电荷改号，两种力同时反向，仍可选择相同速率。

\velocitySelectorDiagram
\end{xesolution}



\begin{origexamplebox}{【漂移运动】}\label{ch22:ex:s38-7-8}
平行板电容器电压为 $U$、板间距为 $d$，板间有平行极板且垂直纸面的匀强磁场 $B$。质量为 $m$、
正电量为 $Q$ 的粒子从正极板附近由静止释放。若要求粒子不能到达负极板，$B$ 至少多大？

\begin{center}
\begin{tikzpicture}[>=Latex,scale=.8]
 \draw[very thick] (-3,1.2)--(3,1.2) (-3,-1.2)--(3,-1.2);
 \foreach \x in {-2.7,-1.8,...,2.7}{\node[above] at (\x,1.2){$+$};\node[below] at (\x,-1.2){$-$};}
 \foreach \x in {-2,-.5,1,2.5}{\node at (\x,0){$\otimes$};}
 \draw[<->] (3.45,-1.2)--(3.45,1.2) node[midway,right]{$d$};
 \fill[DeepBlue] (-1,1.03) circle(2.5pt);\node[below right=2pt] at (-1,.85){$m,+Q;\ v_0=0$};
 \draw[vector] (-2.7,.7)--(-2.7,-.7) node[midway,left]{$\boldsymbol E$};
 \node at (.4,-.6){$B$ 入纸面};\node[font=\small] at (0,-1.95){两板电势差为 $U$};
\end{tikzpicture}
\end{center}
\figcap{正电荷从正极板附近静止释放，轨迹待求。}\label{ch22:fig:em22-question-8}
\end{origexamplebox}\par\nopagebreak[4]

\begin{xesolution}
板间电场大小 $E=U/d$。取 $\vect E$ 向下、$\vect B$ 入纸面，则导向中心以
$v_D=E/B$ 水平漂移，粒子相对导向中心以 $\omega_{\rm c}=QB/m$ 回旋。从静止初态积分可得摆线
\[
  x=\frac{v_D}{\omega_{\rm c}}(\omega_{\rm c}t-\sin\omega_{\rm c}t),
  \qquad
  y=\frac{v_D}{\omega_{\rm c}}(1-\cos\omega_{\rm c}t),
\]
其中 $y$ 沿电场方向由正极板量起。最大下降距离
\begin{equation}
  y_{\max}=\frac{2v_D}{\omega_{\rm c}}
  =\frac{2mE}{QB^2}=\frac{2mU}{QB^2d}.
\end{equation}
临界条件 $y_{\max}=d$，故
\begin{equation}
  \boxed{B_{\min}=\frac1d\sqrt{\frac{2mU}{Q}}}.
\end{equation}
若要求轨迹严格不接触负极板，应取 $B>B_{\min}$。

\crossedFieldCycloidDiagram
\end{xesolution}



\begin{origexamplebox}{【花式漂移】}\label{ch22:ex:s38-7-9}
在地球赤道附近，把地面近似为平面、地磁场近似为匀强场。质量为 $m$、电量为 $q$ 的粒子从地面
上方 $H$ 处以速率 $v_0$ 竖直向下运动，且 $H\ll R_\oplus$。求粒子不与地面相碰的条件，以及任意
时刻它到初始位置的距离。
\end{origexamplebox}\par\nopagebreak[4]

\begin{xesolution}
取 $y$ 轴竖直向上，初始位置为 $(0,H)$，地面为 $y=0$；令地磁场沿 $+z$ 方向，并记有号回旋频率
\[
  \Omega=\frac{qB}{m}.
\]
运动方程为
\[
  \dot v_x=\Omega v_y,
  \qquad
  \dot v_y=-\Omega v_x-g,
  \qquad v_x(0)=0,\quad v_y(0)=-v_0.
\]
把恒定的重力漂移 $v_{D,x}=-g/\Omega$ 分离出来，积分得到
\begin{align}
  x(t)&=-\frac g\Omega t+\frac g{\Omega^2}\sin\Omega t
       +\frac{v_0}{\Omega}(\cos\Omega t-1),\\
  y(t)&=H+\frac g{\Omega^2}(\cos\Omega t-1)
       -\frac{v_0}{\Omega}\sin\Omega t.
\end{align}
电荷改号时 $x$ 方向反转，但竖直下降的最大值不变。令
$\omega_{\rm c}=|\Omega|$，一个回旋周期内相对初始高度的最大下降量是
\begin{equation}
  \Delta y_{\max}=\frac g{\omega_{\rm c}^2}
  +\sqrt{\left(\frac g{\omega_{\rm c}^2}\right)^2
  +\left(\frac{v_0}{\omega_{\rm c}}\right)^2}.
\end{equation}
因此严格不碰地的条件为
\begin{equation}
  \boxed{H>\frac g{\omega_{\rm c}^2}
  +\sqrt{\frac{g^2}{\omega_{\rm c}^4}+\frac{v_0^2}{\omega_{\rm c}^2}}}.
\end{equation}
任意时刻到初始位置的距离则是
\begin{equation}
  \boxed{s(t)=\sqrt{x(t)^2+[y(t)-H]^2}},
\end{equation}
其中 $x(t),y(t)$ 由上式给出。第一项造成持续的水平漂移，余下项为有界回旋。

\gravityDriftDiagram
\end{xesolution}



\begin{origexamplebox}{【交替电磁场】}\label{ch22:ex:s38-7-10}
空间沿竖直方向交替分布电场层与磁场层，每层厚度均为 $d$。电场 $\vect E$ 向上，磁场 $\vect B$
垂直纸面向里。质量为 $m$、正电量为 $q$ 的粒子从第一电场层底部由静止释放。它最多能到达第几层
磁场？
\end{origexamplebox}\par\nopagebreak[4]

\begin{xesolution}
电场每帮助粒子向上通过一层，动能增加 $qEd$。所以进入第 $k$ 个磁场层时
\begin{equation}
  v_k^2=\frac{2kqEd}{m}.
\end{equation}
令 $\omega=qB/m$。在磁场层内
\[
  \dot v_x=-\omega v_y,
  \qquad \dot y=v_y,
\]
故
\begin{equation}
  \frac{\dd}{\dd t}(v_x+\omega y)=0.
\end{equation}
粒子每向上穿过一个厚度 $d$ 的磁场层，$v_x$ 就减少 $\omega d$。因此进入第 $k$ 层前已穿过
$k-1$ 个磁场层，$v_{x,\mathrm{in}}=-(k-1)\omega d$；若能穿出第 $k$ 层，则
\[
  v_{x,\mathrm{out}}=-k\omega d.
\]
磁场不改变总速率，所以穿出条件是 $v_k^2-v_{x,\mathrm{out}}^2\ge0$：
\begin{equation}
  2k\frac{qE}{m}d-k^2\left(\frac{qB}{m}\right)^2d^2\ge0,
\end{equation}
即
\begin{equation}
  k\le \Lambda,
  \qquad \Lambda=\frac{2mE}{qB^2d}.
\end{equation}
它能完整穿过的磁场层数为 $\lfloor\Lambda\rfloor$，随后还会进入下一电场层并到达下一磁场层，
但不能再穿出。因此最多到达
\begin{equation}
  \boxed{N_{\max}=\lfloor\Lambda\rfloor+1
  =\left\lfloor\frac{2mE}{qB^2d}\right\rfloor+1}
\end{equation}
层磁场（临界擦边时约定按上方电场层处理）。若 $\Lambda$ 恰为整数，边界处为法向分量 $v_y=0$，
而切向分量 $v_x=-k\omega d\ne0$。理想阶跃场在界面的取值必须规定：若按上方电场处理，粒子会重新
向上加速，仍得上述式子；若按磁场侧处理，则在此切触后返回，不能把“到达界面”等同于“穿过界面”。
真实过渡区应按实际场分布求解，非整数 $\Lambda$ 的结论不受这个临界约定影响。
图中界面两侧位置与速度连续，后续电场段不能再画成竖直线。

\alternatingLayersDiagram
\end{xesolution}



\begin{origexamplebox}{【霍尔效应】}\label{ch22:ex:s38-7-11}
长方体导体沿 $x$ 方向通电流 $I$，长度、沿磁场的宽度和竖直厚度分别为 $c,b,a$。载流子数密度为
$n$，每个载流子带电 $+q$，外加磁场沿 $y$ 方向。稳定后上下表面的霍尔电势差是多少？若载流子
带电 $-q$，结果怎样？说明其用途，并判断理想单载流子模型中纵向电阻是否改变。

\begin{center}
\begin{tikzpicture}[>=Latex,scale=.67]
 \coordinate(A)at(-3,-1);\coordinate(B)at(2,-1);\coordinate(C)at(3.5,.2);\coordinate(D)at(-1.5,.2);
 \coordinate(At)at(-3,1);\coordinate(Bt)at(2,1);\coordinate(Ct)at(3.5,2.2);\coordinate(Dt)at(-1.5,2.2);
 \fill[SoftGray] (At)--(Bt)--(Ct)--(Dt)--cycle;
 \draw[thick] (A)--(B)--(C)--(Ct)--(Dt)--(At)--cycle (At)--(Bt)--(Ct) (B)--(Bt);
 \draw[guide] (A)--(D)--(C) (D)--(Dt);
 \draw[vector] (-4.4,0)--(-3,0) node[midway,above]{$I\parallel x$};
 \draw[vector] (1.8,-2.3)--(3.3,-1.1) node[below right]{$\boldsymbol B\parallel y$};
 \draw[<->] (-3.45,-1)--(-3.45,1) node[pos=.15,left]{$a$};
 \draw[<->] (-3,-1.45)--(2,-1.45) node[midway,below]{$c$};
 \draw[<->] (2.35,-1.15)--(3.85,.05) node[midway,right]{$b$};
 \draw (0,1.6)--(0,3)--(5,3)--(5,1.0);
 \draw (0,-1)--(0,-2.6)--(5,-2.6)--(5,.2);
 \draw (5,.6) circle(.4);\node at(5,.6){V};
 \node[right,font=\small] at(5.55,.6){上下表面电势差};
\end{tikzpicture}
\end{center}
\figcap{霍尔测量的导体尺寸、外磁场与电流方向；V 表示高内阻电压表。}\label{ch22:fig:em22-question-11}
\end{origexamplebox}\par\nopagebreak[4]

\begin{xesolution}
垂直于电流方向的横截面积为 $ab$，故电流密度和漂移速率满足
\begin{equation}
  j_x=\frac{I}{ab}=nqv_d,
  \qquad v_d=\frac{I}{nqab}.
\end{equation}
载流子被磁力推向上下表面，电荷积累建立霍尔电场。稳定时横向电流为零，
\[
  qE_H=qv_dB,
  \qquad E_H=v_dB.
\]
上下表面相距 $a$，所以
\begin{equation}
  \boxed{|U_H|=E_Ha=\frac{IB}{nqb}}.
\end{equation}
为给出极性，按题设取常规电流沿 $+x$、磁场沿 $+y$，以竖直向上为 $+z$，组成右手坐标系。
正载流子受到 $q\boldsymbol v_d\times\boldsymbol B$ 的方向为 $+z$，故稳态霍尔电场向下，上表面电势较高。明确规定
$U_H=V_{\text{上}}-V_{\text{下}}$ 后，有
\[
 \boxed{U_H=+\frac{IB}{nqb}\quad(\text{正载流子})}.
\]
若载流子改为 $-q$，在给定常规电流方向下漂移速度也反向，但电荷符号同时改变；最终霍尔电场与
原来反向，$U_H$ 的符号改变、大小不变。因此霍尔电压既可测磁场，也可判定主要载流子正负并测量
载流子浓度。

在理想均匀、单载流子德鲁德模型中，横向霍尔电场恰好抵消横向磁力，纵向关系仍为
$E_x=\rho j_x$，故四端测得的纵向电阻不变。真实材料若有多种载流子、复杂能带或几何混入，则可能
出现磁阻。

\hallDiagram
\end{xesolution}



\begin{origexamplebox}{【科尔比诺】}\label{ch22:ex:s38-7-12}
内、外半径分别为 $a,b$ 的圆环形导体电阻率为 $\rho$、载流子迁移率为 $\mu$，轴向长度为 $L$，且
$L\gg a,b$。内外圆柱面间加电压 $U$。求无磁场时的电流；再加轴向磁场 $B$ 后电流如何改变，并说明
这种结构的用途。

\begin{center}
\begin{tikzpicture}[>=Latex,scale=.75]
 \fill[SoftGray,even odd rule] (0,0)circle(1.9) (0,0)circle(.8);
 \draw[thick] (0,0)circle(1.9) (0,0)circle(.8);
 \draw[thin] (0,0)--(30:.8) node[midway,above]{$a$};
 \draw[thin] (0,0)--(210:1.9) node[midway,below]{$b$};
 \node at (-1.05,1.0){$\odot B$};
 \draw (.8,0)--(3,0)--(3,-2.5)--(1.7,-2.5);
 \draw (0,-1.9)--(0,-2.5)--(1.3,-2.5);
 \draw[thick] (1.3,-2.68)--(1.3,-2.32) (1.7,-2.8)--(1.7,-2.2);
 \node[below] at(1.5,-2.85){$U$};
 \node[align=left,anchor=west,font=\small] at(3.7,.7){圆柱横截面\\轴向长度 $L$\\电阻率 $\rho$，迁移率 $\mu$};
\end{tikzpicture}
\end{center}
\figcap{同轴圆柱导体的横截面。引线绝缘，仅端点接触内、外电极，电压加在两圆柱面之间。}\label{ch22:fig:em22-question-12}
\end{origexamplebox}\par\nopagebreak[4]

\begin{xesolution}
半径 $r$ 处电流通过的面积为 $2\pi rL$，微元电阻
\[
  \dd R=\frac{\rho\,\dd r}{2\pi rL}.
\]
积分得
\begin{equation}
  R_0=\frac{\rho}{2\pi L}\ln\frac ba,
  \qquad
  \boxed{I_0=\frac{2\pi LU}{\rho\ln(b/a)}}.
\end{equation}

加轴向磁场后，径向漂移会受到方位方向的洛伦兹力。科尔比诺圆环中不能靠静电荷建立一个沿整圈
同向、线积分非零的静电场来完全抵消这种偏转，因此出现方位电流，径向电导降低。单载流子德鲁德
模型可以直接求出这一电导。取 $B$ 沿 $+z$，载流子带符号电荷为 $q_c$，质量为 $m_c$，弛豫时间为 $\tau$。
稳态平均速度满足 $0=q_c(\boldsymbol E+\boldsymbol v\times\boldsymbol B)-m_c\boldsymbol v/\tau$。
轴对称静电场的 $E_\varphi=0$，于是
\[
 v_r=\frac{q_c\tau}{m_c}(E_r+Bv_\varphi),\qquad
 v_\varphi=-\frac{q_c\tau B}{m_c}v_r.
\]
把第二式代入第一式，令 $\mu=|q_c|\tau/m_c$、$\sigma_0=nq_c^2\tau/m_c=1/\rho$，得到
\[
 v_r=\frac{q_c\tau E_r/m_c}{1+\mu^2B^2},\qquad
 j_r=\frac{\sigma_0 E_r}{1+\mu^2B^2},\qquad
 j_\varphi=-\frac{q_c\tau B}{m_c}j_r.
\]
因此
\begin{equation}
  \sigma_{rr}(B)=\frac{\sigma_0}{1+(\mu B)^2}.
\end{equation}
所以
\begin{equation}
  \boxed{R(B)=R_0[1+(\mu B)^2]},
  \qquad
  \boxed{I(B)=\frac{I_0}{1+(\mu B)^2}}.
\end{equation}
这种没有横向边界积累霍尔电荷的圆环几何可以直接测量磁阻，进而求迁移率或作为磁场传感结构。

\corbinoDiagram
\end{xesolution}



\section{缓变场、扰动与阻尼运动}

\subsection{缓变磁场中的绝热近似}


若磁场在一个回旋半径内变化很小、在一个回旋周期内变化很慢，粒子仍可在局部看成回旋，第一
绝热不变量近似守恒：
\begin{equation}
 \boxed{\frac{p_\perp^2}{2mB}=\mu=\text{常量}}.
\end{equation}
粒子进入更强磁场时，横向动能随 $B$ 增大；总动能若不变，平行动能便减小。当平行速度降到零，
粒子会被反射，形成磁镜。这个机制说明磁场虽然不做功，仍能在平行和垂直两个运动分量之间重新
分配动能。

轴对称缓变磁场还可使偏离轴线的粒子受到平均聚焦作用。实际粒子束装置常把均匀场中的回旋、
场边缘的方向变化和多段磁场组合起来，实现选速、聚焦与动量分析。使用局部回旋近似时必须检查
$r_c$ 是否远小于磁场变化尺度；若二者同量级，简单的导向中心图像不再可靠。


% EM teaching enrichment 20260927
\paragraph{磁镜反射条件的逐步求法}
在静态磁场且无电场做功时，总动能 $K$ 不变。设起点场强为 $B_0$、速度与场夹角为 $\alpha_0$，则
$\mu=K\sin^2\alpha_0/B_0$，沿磁力线到场强 $B$ 处有
\[
 K_\parallel=K-\mu B
 =K\left(1-\frac B{B_0}\sin^2\alpha_0\right).
\]
反射发生在 $K_\parallel=0$，即 $B_m=B_0/\sin^2\alpha_0$。若装置最大场强为 $B_{\max}$，能被反射须满足
$\sin^2\alpha_0\ge B_0/B_{\max}$；过于沿场的粒子会从损失锥逃逸。这个判据同时依赖能量守恒与绝热不变量，两者缺一不可。


\subsection{线性阻尼与小扰动的求解框架}
取匀强磁场 $\boldsymbol B=B\boldsymbol e_z$，并在横向运动中加入线性阻力 $-m\nu\boldsymbol v$。这里 $\nu>0$ 的单位为时间的倒数。令 $w=v_x+\mathrm i v_y$、$\Omega=qB/m$，运动方程化为
\[
 \dot w+(\nu+\mathrm i\Omega)w=\frac q m(E_x+\mathrm iE_y).
\]
无电场时 $w(t)=w_0e^{-\nu t}e^{-\mathrm i\Omega t}$：速度方向仍以带符号回旋频率旋转，而速率指数衰减，动能按 $e^{-2\nu t}$ 衰减。此时瞬时曲率半径随速率变小；轨迹并非围绕一个固定圆心的同心圆，位置应再对速度积分。

恒定横向电场下，先求稳定速度
\[
 w_\infty=\frac{q(E_x+\mathrm iE_y)}{m(\nu+\mathrm i\Omega)},\qquad
 w(t)=w_\infty+(w_0-w_\infty)e^{-(\nu+\mathrm i\Omega)t}.
\]
取 $E_y=0$，稳定速度包含沿电场的耗散分量和横向分量：
\[
 v_{x,\infty}=\frac{qE_x\nu}{m(\nu^2+\Omega^2)},\qquad
 v_{y,\infty}=-\frac{qE_x\Omega}{m(\nu^2+\Omega^2)}.
\]
令 $\nu\to0$，稳定特解退回 $v_x=0,v_y=-E_x/B$；但无阻尼时初值引起的回旋不会衰减，所以这不表示所有粒子都自动趋向匀速漂移。令 $B\to0$，则回到电场与阻力平衡的漂移速度。

若外场随位置缓慢变化，小扰动法要围绕指定参考轨道展开，把 $\boldsymbol r=\boldsymbol r_0+\delta\boldsymbol r$、$\boldsymbol v=\boldsymbol v_0+\delta\boldsymbol v$ 代入，减去参考轨道方程，只保留扰动的一次项。常数项没有消去通常意味着参考轨道本身不满足原方程；丢弃二次项之前也应说明扰动位移相对场变化尺度足够小。稳定性由线性解是否增长判断，不能只凭某个瞬时力的方向下结论。

例题若把阻力系数记为 $\gamma$、写成 $-\gamma\boldsymbol v$，则本小节的衰减率应取 $\nu=\gamma/m$；这两种系数的单位不同。

% EM structure and teaching revision 20260928
\subsection{磁镜条件与绝热近似的两重检查}
无电场且磁场随位置缓慢变化时，总动能保持不变，而第一绝热不变量
$\mu=m v_\perp^2/(2B)$ 近似保持。由初始磁场 $B_0$ 与俯仰角 $\alpha_0$ 得
\[
 v_\parallel^2=v^2\left(1-\frac{B}{B_0}\sin^2\alpha_0\right),\qquad
 B_{\mathrm{turn}}=\frac{B_0}{\sin^2\alpha_0}.
\]
磁镜反射发生在纵向速度降到零时；总速率没有下降，只是平行与垂直分量重新分配。若沿途最大场为 $B_{\max}$，反射要求 $\sin^2\alpha_0\ge B_0/B_{\max}$。
这个判据还要满足绝热条件：一回旋半径内场的相对变化很小，典型地 $r_c|\nabla B|/B\ll1$；还应检查粒子沿场前进时一回旋周期内所经历的场变化，例如 $|v_\parallel\partial_\parallel B|/(|\Omega|B)\ll1$，其中 $|\Omega|=|q|B/m$。仅有横向回旋半径很小，对 $v_\parallel\gg v_\perp$ 的粒子并不充分。若场本身随时间变化，也需回旋周期远小于场的变化时间。磁场在一个锐利界面突然改变时，应按分区轨迹匹配，不能直接宣称磁矩跨界不变。
若存在碰撞、辐射或强阻尼，总动能和磁矩可能同时改变。解题时把精确守恒量、近似不变量和纯几何约束分别列出，才知道哪一步需要额外假设。


\subsection{专题例题与方法应用}

\begin{origexamplebox}{【磁场中的小扰动】}\label{ch22:ex:s38-7-13}
柱形区域内有垂直纸面的匀强磁场 $B$。质量为 $m$、电量为 $+q$ 的粒子在半径 $R$ 处绕圆心作匀速
圆周运动。受到很小的径向扰动后，径向运动的周期是多少？
\end{origexamplebox}\par\nopagebreak[4]

\begin{xesolution}
匀强磁场中任意平面轨迹都是回旋圆。小的径向速度或位移扰动不会改变回旋频率，只会使新回旋圆的
圆心相对原圆心产生一个小位移 $\delta$。若新圆半径仍记作 $R+\order{\delta}$，则相对原圆心的距离
\[
  r(t)=\sqrt{R^2+\delta^2+2R\delta\cos(\omega_{\rm c}t+\varphi_0)}
  \simeq R+\delta\cos(\omega_{\rm c}t+\varphi_0).
\]
故径向小振动角频率就是回旋角频率：
\begin{equation}
  \boxed{\omega_r=\omega_{\rm c}=\frac{qB}{m}},
  \qquad
  \boxed{T_r=\frac{2\pi m}{qB}}.
\end{equation}

\epicycleDiagram
\end{xesolution}



\begin{origexamplebox}{【电磁场中的小扰动】}\label{ch22:ex:s38-7-14}
柱形区域内有垂直纸面向里的匀强磁场 $B$，并有径向向外的电场
$E=\alpha/r$，$\alpha>0$。质量为 $m$、电量为 $+q$ 的粒子在半径 $R$ 处作逆时针匀速圆周运动。求径向小扰动的
周期。

\begin{center}
\begin{tikzpicture}[>=Latex,scale=.8]
 \draw[thick] (0,0)circle(1.7);\node at(0,.35){$\otimes B$};
 \foreach \a in {0,60,120,180,240,300}{\draw[vector] (\a:.8)--(\a:1.35);}
 \node[right] at(2.2,.2){$E=\alpha/r$，径向向外};
 \fill[DeepBlue] (0,1.7)circle(2pt) node[above right]{$m,+q$};
 \draw[vector] (0,1.7)--(-1.2,1.7) node[left]{$\boldsymbol v$};
 \draw[->] (0,0)--(0,-1.7) node[pos=.25,right=2pt]{$R$};
\end{tikzpicture}
\end{center}
\figcap{径向向外电场与给定的逆时针圆周初态。}\label{ch22:fig:em22-question-14}
\end{origexamplebox}\par\nopagebreak[4]

\begin{xesolution}
先判断转向：图中正电荷在圆顶向左运动，$B$入纸面，故磁力向内、电场力向外。
若改成顺时针，二力都向外，就不能维持圆周运动。参数 $\alpha=Er$ 的单位为伏特，电场力的系数为 $q\alpha$。

记 $\omega_B=qB/m>0$，圆周运动角速度为 $\Omega>0$。磁力向内，径向力平衡为
\begin{equation}
  qBR\Omega=mR\Omega^2+\frac{q\alpha}R,
\end{equation}
即
\begin{equation}
  \Omega^2-\omega_B\Omega+\frac{q\alpha}{mR^2}=0.
\end{equation}
圆轨道存在需 $\omega_B^2\ge4q\alpha/(mR^2)$。

在轴对称场中正则角动量守恒，可写为
\begin{equation}
  r^2\left(\dot\varphi-\frac{\omega_B}{2}\right)
  =R^2\left(\Omega-\frac{\omega_B}{2}\right).
\end{equation}
令 $r=R+x$、$|x|\ll R$，则
\[
  \delta\dot\varphi=-\frac{2}{R}
  \left(\Omega-\frac{\omega_B}{2}\right)x.
\]
径向方程为 $\ddot r=r\dot\varphi^2-\omega_Br\dot\varphi+q\alpha/(mr)$。
记 $a=q\alpha/m$，在 $R,\Omega$ 附近展开，并减去圆轨道的零阶平衡式：
\[
 \ddot x=(\Omega^2-\omega_B\Omega-a/R^2)x
       +R(2\Omega-\omega_B)\delta\dot\varphi.
\]
由 $\Omega^2-\omega_B\Omega=-a/R^2$ 和上面的角速度变化式，系数变为
$-2a/R^2-(2\Omega-\omega_B)^2$。再用 $(2\Omega-\omega_B)^2=\omega_B^2-4a/R^2$，得到
\begin{equation}
  \ddot x+\left(\omega_B^2-\frac{2q\alpha}{mR^2}\right)x=0.
\end{equation}
因此
\begin{equation}
  \boxed{\omega_r=\sqrt{\omega_B^2-\frac{2q\alpha}{mR^2}}},
  \qquad
  \boxed{T_r=\frac{2\pi}{\sqrt{(qB/m)^2-2q\alpha/(mR^2)}}}.
\end{equation}
圆轨道的存在条件已经保证根号内为正。$q\alpha/(mR^2)$ 的量纲为时间的负二次方；
令 $\alpha\to0$，即回到纯磁场的 $\omega_r=\omega_B$。

\radialFieldOrbitDiagram
\end{xesolution}



\begin{origexamplebox}{【磁场中的阻尼运动】}\label{ch22:ex:s38-7-15}
匀强磁场 $B$ 垂直纸面向里。正电粒子 $m,+q$ 从原点以 $v_0$ 沿 $+x$ 方向运动，同时受到线性阻力
$\vect f=-\gamma\vect v$。求任意时刻的位置及最终停止位置。
\end{origexamplebox}\par\nopagebreak[4]

\begin{xesolution}
记
\[
  a=\frac\gamma m,
  \qquad \omega=\frac{qB}{m}.
\]
运动方程为
\[
  \dot v_x=-av_x-\omega v_y,
  \qquad
  \dot v_y=\omega v_x-av_y.
\]
令 $u=v_x+\mathrm i v_y$，则
\[
  \dot u=(-a+\mathrm i\omega)u,
  \qquad u(0)=v_0,
\]
所以
\begin{equation}
  v_x=v_0\e^{-at}\cos\omega t,
  \qquad
  v_y=v_0\e^{-at}\sin\omega t.
\end{equation}
下图的初始切线沿 $+x$，磁力使轨迹向上弯转；标出的极限点将在积分后确定。
\dampedSpiralDiagram
再积分一次：
\begin{align}
  x(t)&=\frac{v_0}{a^2+\omega^2}
  \left[a+\e^{-at}(-a\cos\omega t+\omega\sin\omega t)\right],\\
  y(t)&=\frac{v_0}{a^2+\omega^2}
  \left[\omega+\e^{-at}(-a\sin\omega t-\omega\cos\omega t)\right].
\end{align}
令 $t\to\infty$，最终位置为
\begin{equation}
  \boxed{(x_\infty,y_\infty)=
  \frac{v_0}{a^2+\omega^2}(a,\omega)}.
\end{equation}
速度指数衰减而方向继续转动，故轨迹不是无限延伸的螺旋，而是收敛到有限终点。

\end{xesolution}



\begin{origexamplebox}{【电磁场中的阻尼运动】}\label{ch22:ex:s38-7-16}
匀强磁场 $B$ 垂直纸面向外，匀强电场 $E$ 沿 $+y$ 方向。正电粒子 $m,+q$ 从原点以 $v_0$ 沿
$+x$ 方向运动，并受线性阻力 $\vect f=-\gamma\vect v$。求任意时刻的位置并概括运动规律。
\end{origexamplebox}\par\nopagebreak[4]

\begin{xesolution}
令
\[
  a=\frac\gamma m,
  \qquad \omega=\frac{qB}{m},
  \qquad b=\frac{qE}{m},
  \qquad u=v_x+\mathrm i v_y.
\]
分量方程合并为
\begin{equation}
  \dot u=-(a+\mathrm i\omega)u+\mathrm i b.
\end{equation}
稳态速度为
\begin{equation}
  u_\infty=\frac{\mathrm i b}{a+\mathrm i\omega}
  =\frac{b\omega}{a^2+\omega^2}
  +\mathrm i\frac{ba}{a^2+\omega^2}.
\end{equation}
因此完整速度和位置（令 $z=x+\mathrm i y$）为
\begin{equation}
  \boxed{u(t)=u_\infty+(v_0-u_\infty)\e^{-(a+\mathrm i\omega)t}},
\end{equation}
\begin{equation}
  \boxed{z(t)=u_\infty t+
  \frac{v_0-u_\infty}{a+\mathrm i\omega}
  \left[1-\e^{-(a+\mathrm i\omega)t}\right]}.
\end{equation}
这已经同时给出 $x(t)=\Re z(t)$ 与 $y(t)=\Im z(t)$。运动由两部分组成：以
$\omega$ 回转、以 $e^{-at}$ 衰减的瞬态，加上恒定漂移
\begin{equation}
  \boxed{\vect v_\infty=
  \frac{q^2}{\gamma^2+q^2B^2}
  \left(BE\,\uvec x+\frac{\gamma E}{q}\,\uvec y\right)}.
\end{equation}
当 $\gamma\to0$ 时恢复纯 $\vect E\times\vect B$ 漂移 $v_x=E/B$；阻力存在时还出现沿电场方向的
速度分量。长期轨迹趋于沿 $\vect v_\infty$ 的直线。

\dampedDriftDiagram
\end{xesolution}



\section{磁矩、频移与加速装置}

\subsection{磁场中的角动量与轨道磁矩}


电荷作圆周运动等效为环电流。周期 $T=2\pi/\omega_c$，电流为 $I=q/T$，轨道面积为
$\pi r_c^2$，故磁矩
\begin{equation}
 \vect\mu=I\vect S.
\end{equation}
它与机械轨道角动量 $\vect L=m\vect r\times\vect v$ 的关系为
\begin{equation}
 \boxed{\vect\mu=\frac q{2m}\vect L}.
\end{equation}
负电荷的磁矩与角动量方向相反。外磁场中的相互作用能为
\begin{equation}
 U=-\vect\mu\cdot\vect B,
\end{equation}
非均匀磁场还会对磁矩产生指向强场或弱场的净力。正常塞曼效应的经典频移正与轨道磁矩
和外场的耦合有关。

磁场存在时应区分机械动量 $m\vect v$ 与正则动量。若用矢势
$\vect B=\nabla\times\vect A$，拉格朗日量为
\begin{equation}
 \mathcal L=\frac12mv^2+q\vect A\cdot\vect v-q\varphi,
\end{equation}
相应正则动量为 $\vect p=m\vect v+q\vect A$。多数高中竞赛轨迹题只需机械动量，但涉及对称性、
广义动量或量子磁效应时，两者不能混用。


% EM teaching enrichment 20260927 supplementary

\paragraph{磁矩与磁场的夹角不能只看电荷符号}
纯匀强磁场中的回旋满足 $r_c=mv_\perp/(|q|B)$。取沿真实运动绕行方向的面积矢量，
环电流定义直接给出 $\vect\mu=(q/2m)\vect L$，其中角动量对回旋中心计算。
正电荷与负电荷的回旋方向相反，但它们形成的常规环电流方向相同，磁矩均与外磁场相反。
其大小为
\[
 |\vect\mu|=\frac{|q|v_\perp r_c}{2}=\frac{mv_\perp^2}{2B}=\frac{K_\perp}{B}.
\]
前面绝热不变量中的 $\mu$ 指这个正的大小。
对于固定磁矩的回路可用 $U=-\vect\mu\cdot\vect B$ 讨论外场作用；
回旋粒子的轨道和磁矩会随场改变，不能把相互作用能再机械地加到已给出的回旋动能上重复计数。


\subsection{正常塞曼效应的经典图像}


经典模型把原子中的电子近似成各向同性谐振子，固有角频率为 $\omega_0$。令磁场沿 $z$ 轴，电子在
$z$ 方向的运动不受磁力影响，频率仍为 $\omega_0$；在垂直于磁场的平面内，两种相反旋向的圆偏振
模式可由分量方程求出。电子电荷为 $-e$，记 $\omega_c=eB/m_e>0$，有
\[
 \ddot x=-\omega_0^2x-\omega_c\dot y,\qquad
 \ddot y=-\omega_0^2y+\omega_c\dot x,\qquad
 \ddot z=-\omega_0^2z.
\]
取圆偏振试探解 $x=A\cos\omega t,\ y=sA\sin\omega t$，其中 $s=\pm1$ 表示两种旋向。
代入第一式得 $-\omega^2=-\omega_0^2-s\omega_c\omega$；第二式给出同一条件。于是两种模式满足
\begin{equation}
  \omega^2\mp\omega_{\rm c}\omega-\omega_0^2=0,
  \qquad \omega_{\rm c}=\frac{eB}{m_e}.
\end{equation}
故正频率为
\begin{equation}
  \omega_\pm=\sqrt{\omega_0^2+\frac{\omega_{\rm c}^2}{4}}
  \pm\frac{\omega_{\rm c}}2
  \simeq \omega_0\pm\frac{eB}{2m_e}
\end{equation}
（末式为弱场近似）。于是原来的一条谱线分裂为等间隔三条。真实原子光谱需要量子力学解释，但这个
经典推导准确抓住了正常塞曼分裂与 $B$ 成正比的结构。

塞曼效应是近代物理建立前经典电子模型少有的定量成功之一。无外场时的一条谱线，在磁场中分裂为
中央不变的一条和分列两侧的两条，间隔随 $B$ 增大。洛伦兹把物质中的束缚电子看成具有线性回复力的
振子：沿磁场方向的振动不受磁力影响，垂直方向的两个相反旋向则分别被磁力“加快”和“减慢”。洛伦兹
和塞曼因相关工作共同获得 1902 年诺贝尔物理学奖；今天对反常分裂、精细结构和选择定则的解释
必须使用量子力学，但这里的经典模型仍清楚展示了方向与频移的来源。

\zeemanSpectrumDiagram


\subsection{回旋加速器}


回旋加速器用两个相互绝缘的 D 形盒提供近似无电场的回旋区，在狭缝处施加交变电压。若射频角频率
满足
\begin{equation}
  \omega_{\rm rf}=\omega_{\rm c}=\frac{|q|B}{m},
\end{equation}
粒子每次通过狭缝时电场方向都恰好帮助它加速。若装置的最大轨道半径为 $R$，非相对论极限下最大
动能为
\begin{equation}
  \boxed{E_{k,\max}=\frac{q^2B^2R^2}{2m}}.
\end{equation}
能量升高后相对论质量效应使回旋频率下降，固定频率的电场将逐渐失谐，这是经典回旋加速器的限制。

劳伦斯因回旋加速器的发明与发展获得 1939 年诺贝尔物理学奖。实际装置在 D 形盒上下布置大型
磁极或线圈，在盒间狭缝接入高频振荡器；粒子在盒内只弯转，在狭缝处才被电场加速。
达到约 $730\,\mathrm{MeV}$ 质子能量的历史大型装置采用同步回旋方案，不能直接套用上述固定频率的
非相对论模型。高能区需通过改变射频频率、调整磁场或采用同步加速方案补偿相对论失谐。

\cyclotronDiagram

图中的俯视光滑圆弧对应下面装置剖视图中两个中空的 D 形电极内部。上下磁极与铜色励磁线圈
提供近似垂直于电极平面的磁场；两个绝缘馈电端分别连接高频电源的两端。电极置于真空腔内，
粒子达到外缘后由引出结构送入右侧束流管。剖视图强调部件空间关系，轨迹计算采用上面的俯视模型。
\begin{rTwoFigure}
\includegraphics[width=.86\linewidth]{/Users/dongxingxi/Documents/物理竞赛一轮资料/sample/xes-original/thirty-eighth-section/assets/cyclotron-apparatus-built-in.png}
\captionof{figure}{回旋加速器的装置结构：磁极、励磁线圈、真空腔与 D 形电极。}
\label{ch22:fig:cyclotron-apparatus}
\end{rTwoFigure}
中空结构和束流引出细节见下图。左侧是经过加速缝的侧剖面，虚线轨道位于电极腔内，
不是穿过金属。右侧的负偏转板吸引正离子，使其偏离继续回旋的虚线路径，进入轰击室；
薄铝窗把真空区与外侧隔开，同时允许足够高能的离子通过，靶放在引出束流的路径上。
气体入口用于给中心离子源供气，真空系统负责维持束流运动区的低气压。
\cyclotronExtractionDiagram


% EM teaching enrichment 20260927
\paragraph{逐次过缝的能量账}
若每次过缝获得近似相同的能量 $|q|U_g$，从动能 $K_0$ 起，过 $n$ 次缝后
$K_n=K_0+n|q|U_g$，对应盒内半径
\[
 r_n=\frac{\sqrt{2m(K_0+n|q|U_g)}}{|q|B}.
\]
所以等量增能并不对应等距增大半径，外圈间隔反而逐渐变小。一整圈通常过缝两次，所需时间近似 $nT_c/2$；起始和引出的位置还可能各差一段圆弧。
恒频同步要求每个半圈的时间相同，并忽略过缝时间。把“电压每周期的变化”和“粒子每次到缝的相位”对应起来，比只背共振频率更有用。


% EM structure and teaching revision 20260928
\subsection{回旋加速的同步与失谐}
理想回旋加速器的匀强磁场只负责弯转，交变电场在电极缝隙中做功。非相对论粒子每半周回到缝隙，时间为 $\pi m/(|q|B)$，所以高频电压必须在这段时间内反号，驱动频率等于回旋频率。
若每次过缝时电压大小近似为 $U_g$，且粒子总在加速相位通过，则每次获得动能 $|q|U_g$。经过 $N$ 次过缝，忽略初能与损失，有
\[
 K_N=N|q|U_g,\qquad r_N=\sqrt{\frac{2mNU_g}{|q|B^2}}.
\]
因而半径随 $\sqrt N$ 增长，轨道间距逐渐变小。半径达到装置外缘 $R$ 时，经典最大动能为 $q^2B^2R^2/(2m)$，它同时受到装置尺寸与磁场强度限制。
高能区中 $\gamma$ 增大，回旋周期随之变长，恒频电源与粒子到缝时间失配。此时即使磁场强度和电极电压不变，粒子也可能在不合适的相位减速；调整驱动频率或磁场才可能恢复同步。


\subsection{专题例题与方法应用}

\begin{origexamplebox}{【运动电荷的轨道磁矩】}\label{ch22:ex:s38-7-17}
平面电流环的磁矩定义为 $\vect\mu=I\vect S$。
\begin{enumerate}[label=(\arabic*),leftmargin=2.2em]
  \item 求带电粒子在匀强磁场中回旋运动的等效磁矩；
  \item 在沿轴线缓慢变化的磁场中，若粒子回旋很多圈后磁场才明显改变，写出局部轨道磁矩；
  \item 证明该磁矩在上述缓变条件下近似守恒。
\end{enumerate}
\end{origexamplebox}\par\nopagebreak[4]

\begin{xesolution}
\textbf{(1) 匀强场。}
令 $v_\perp$ 为垂直磁场的速度，回旋半径 $\rho=mv_\perp/(Q B)$，周期
$T=2\pi m/(QB)$。等效电流大小 $Q/T$，所以
\begin{equation}
  \boxed{\mu=\frac{Q}{T}\pi\rho^2
  =\frac{mv_\perp^2}{2B}=\frac{E_\perp}{B}}.
\end{equation}
矢量方向由真实电流方向按右手定则确定。对无论正负电荷的自由回旋，$\vect\mu$ 都与局部
$\vect B$ 反向；同时仍有 $\vect\mu=(q/2m)\vect L$。

\textbf{(2) 缓变轴向场。}
设磁场沿导向中心弧长 $s$ 缓慢变化，局部回旋半径为 $\rho$，场变化尺度
$L_B=B/|\dd B/\dd s|$。当
\begin{equation}
  \frac{\rho}{L_B}\ll1,
  \qquad
  \frac1{\omega_{\rm c}B}\left|\frac{\dd B}{\dd t}\right|\ll1,
\end{equation}
一个周期内磁场近似不变，可以逐点定义
\begin{equation}
  \boxed{\mu(s)=\frac{mv_\perp^2(s)}{2B(s)}}.
\end{equation}
下图左侧给出回旋电流与磁矩的方向；右侧沿磁场增强方向，轨道半径逐渐减小，
平行速度同时减小。下面从磁力分量与能量关系推导这一变化。
\magneticMomentDiagram

\textbf{(3) 绝热守恒。}
先给出不依赖作用量定理的证明。取局部磁轴为 $z$，轴对称场在轴附近满足
$\nabla\cdot\boldsymbol B=0$，故 $B_r\simeq-(r/2)\,\dd B_z/\dd z$。
一次回旋内把 $B_z$ 视为常量，带符号电荷为 $q$ 的粒子有 $v_\varphi=-qB\rho/m$。
利用 $\boldsymbol e_\varphi\times\boldsymbol e_r=-\boldsymbol e_z$，轴向平均磁力为
\[
 F_\parallel=-qv_\varphi B_r
 =-\frac{q^2B\rho^2}{2m}\frac{\dd B}{\dd s}
 =-\frac{mv_\perp^2}{2B}\frac{\dd B}{\dd s}=-\mu\frac{\dd B}{\dd s}.
\]
磁场不做总功，$E_\perp+E_\parallel$保持不变。因此
\[
 \frac{\dd E_\perp}{\dd t}=-F_\parallel v_\parallel
 =\frac{E_\perp}{B}\frac{\dd B}{\dd s}\frac{\dd s}{\dd t}
 =\frac{E_\perp}{B}\frac{\dd B}{\dd t},\qquad
 \frac{\dd}{\dd t}\frac{E_\perp}{B}=0
\]
在所取缓变近似的精度内成立。这解释了横向与纵向动能的交换，而不把磁矩守恒当作证明的前提。

也可用正则作用量表述。此处必须取正则动量 $\boldsymbol p_{\rm can}=m\boldsymbol v+q\boldsymbol A$，
不是半径公式里的机械动量。采用 $\boldsymbol A=\boldsymbol B\times\boldsymbol r/2$，沿实际回旋方向积分，
机械项为 $4\pi E_\perp/\omega_c$，磁矢势项为其负的一半，故
\begin{equation}
  J=\oint\vect p_{{\rm can},\perp}\cdot\dd\vect r_\perp
  =\frac{2\pi E_\perp}{\omega_{\rm c}}
  =\frac{2\pi m}{Q}\frac{E_\perp}{B}
\end{equation}
是绝热不变量，因此
\begin{equation}
  \boxed{\frac{\dd}{\dd t}\left(\frac{E_\perp}{B}\right)\simeq0},
  \qquad \boxed{\mu\simeq\text{常量}}.
\end{equation}
等价地，导向中心沿磁场方向受到平均镜力
\begin{equation}
  F_\parallel=-\mu\frac{\dd B}{\dd s}.
\end{equation}
粒子进入强场区时，$E_\perp=\mu B$ 增加；因磁场不做总功，总动能守恒，故平行动能相应减少，
甚至在 $v_\parallel=0$ 处反射，形成磁镜。若磁场在一圈内变化显著，上述绝热守恒便会失效。

\end{xesolution}
}

\endgroup

\endgroup

% EMOpticsQuality20261023
% QualityRevision20261023
\begingroup
\tcbset{breakable/.default=false,after={\par\nobreak\smallskip}}
\RenewDocumentEnvironment{xesolution}{+b}{%
  \par\nobreak\smallskip\nobreak\begingroup\small
  \clubpenalties 2 10000 10000\displaywidowpenalties 2 10000 10000\noindent
  \LectureSolutionHeading #1\par\endgroup\smallskip
}{}
\chapter{磁场、安培力与安培定律}

{
\graphicspath{{../../../sample/xes-original/thirty-ninth-section/}{../../../sample/xes-original/thirty-ninth-section/assets/}}
\section{电流元叠加与典型磁场}

% Restrained diagrams for Lesson 8.  Vector directions and integration geometry
% are kept explicit so that every figure can be read together with its proof.

\newcommand{\fieldInto}{\mathord{\otimes}}
\newcommand{\fieldOut}{\mathord{\odot}}

\newcommand{\ampereForceDiagram}{%
\begin{center}
\begin{tikzpicture}[scale=0.9,>=Latex]
  \fill[SoftGray] (-4.6,-0.38) rectangle (1.0,0.38);
  \draw[Ink,thick] (-4.6,-0.38) rectangle (1.0,0.38);
  \draw[vector] (-3.9,0)--(-1.9,0) node[midway,below] {$\vect v_d$};
  \fill[DeepBlue] (-3.9,0) circle (3pt) node[left=4pt] {$q$};
  \draw[<->] (-4.35,-0.7)--(0.75,-0.7);
  \node[below=2pt,inner sep=0.5pt] at (-1.8,-0.7) {$\dd\ell$};
  \foreach \x in {-2.8,-1.8,-.8}{
    \node[font=\Large,DeepBlue] at (\x,1.0) {$\fieldInto$};
  }
  \node[DeepBlue] at (0.65,1.0) {$\vect B$};
  \draw[vector] (-3.9,0)--(-3.9,1.45) node[left] {$\dd\vect F$};
  \node[align=left,anchor=west] at (2.0,0.75)
    {$I=nqAv_d$，$\dd N=nA\,\dd\ell$\\[4pt]
     $\displaystyle\dd\vect F=\dd N(q\vect v_d\times\vect B)
       =I\,\dd\vect\ell\times\vect B$};
\end{tikzpicture}

{\small\color{TextGray}安培力就是导线中载流子所受洛伦兹力的宏观合力}
\end{center}}

\newcommand{\bentWireResultDiagram}{%
\begin{center}
\begin{tikzpicture}[x=1cm,y=1cm,>=Latex]
  \foreach \x in {-0.3,0.7,...,7.7}{\draw[DeepBlue!45,->] (\x,-0.7)--++(0,3.4);}
  \node[DeepBlue] at (8.3,2.6) {$\vect B$};
  \draw[Ink,very thick] (0,0)
    .. controls (1.2,1.8) and (2.1,-0.35) .. (3.2,1.15)
    .. controls (4.3,2.65) and (5.9,0.55) .. (7.2,2.0);
  \fill[DeepBlue] (0,0) circle (2.5pt) node[below] {$a$};
  \fill[DeepBlue] (7.2,2.0) circle (2.5pt) node[above] {$b$};
  \draw[Ink,very thick,postaction={decorate},decoration={markings,mark=at position .45 with {\arrow{Latex}}}]
    (0,0).. controls (1.2,1.8) and (2.1,-0.35) .. (3.2,1.15);
  \node[anchor=south east] at (.6,.85) {$I$};
  \draw[vector,DeepBlue] (0,0)--(7.2,2.0) node[midway,below] {$\vect L_{ab}$};
  \node[TextGray,align=center] at (3.7,-1.2)
    {$\displaystyle\vect F=I\!\int_a^b\dd\vect\ell\times\vect B
      =I\vect L_{ab}\times\vect B$；闭合回路在匀强磁场中合力为零};
\end{tikzpicture}
\end{center}}

\newcommand{\loopTorqueDiagram}{%
\begin{center}
\begin{tikzpicture}[scale=0.9,>=Latex]
  \coordinate (O) at (0,0);
  \draw[DeepBlue,very thick] (-2.5,-1.1)--(1.2,-1.1)--(2.5,1.1)--(-1.2,1.1)--cycle;
  \draw[vector] (-0.5,-1.1)--(0.8,-1.1) node[below] {$I$};
  \node[font=\Large] at (-2.5,-0.25) {$\fieldOut$};
  \node[left=4pt] at (-2.5,-0.25) {$\vect F_1$};
  \node[font=\Large] at (2.5,0.25) {$\fieldInto$};
  \node[right=4pt] at (2.5,0.25) {$\vect F_2$};
  \draw[vector] (0,0)--(0.8,1.7) node[above] {$\vect\mu=I\vect S$};
  \draw[vector] (0,0)--(2.35,0) node[right] {$\vect B$};
  \draw[->,DeepBlue,thick] (1.0,1.15) arc[start angle=55,end angle=5,radius=1.3];
  \node[DeepBlue,font=\small,anchor=south west] at (1.62,1.42) {转动趋势};
  \node[font=\Large,DeepBlue] at (1.25,-0.28) {$\fieldInto$};
  \node[DeepBlue,right=2pt] at (1.25,-0.28) {$\vect\tau$};
  \node[TextGray,align=left,anchor=west] at (4.1,0.4)
    {$\displaystyle\vect\tau=\vect\mu\times\vect B$\\
     $\displaystyle U=-\vect\mu\cdot\vect B$\\[2pt]
     $\vect\mu\parallel\vect B$ 为稳定平衡};
\end{tikzpicture}
\end{center}}

\newcommand{\motorHelmholtzContext}{%
\historyfigure[0.9\linewidth]{motor-helmholtz-built-in-v1.png}%
{左：换向器使线圈力矩在每半周后保持驱动方向；右：两同轴线圈的间距等于半径时，中部磁场最均匀。}}

\newcommand{\ampereLawDiagram}{%
\begin{center}
\begin{tikzpicture}[scale=0.9,>=Latex]
  \draw[curve] (0,0) ellipse (3.3 and 1.65);
  \draw[->,DeepBlue,very thick] ({3.3*cos(16)},{1.65*sin(16)}) arc[start angle=16,end angle=78,x radius=3.3,y radius=1.65];
  \node[DeepBlue] at (2.75,1.65) {环路正向};
  \foreach \p/\lab in {(-1.6,0.35)/I_1,(0.15,-0.35)/I_2,(1.55,0.5)/I_3}{
    \node[font=\Large] at \p {$\fieldOut$};
    \node[below=4pt] at \p {$\lab$};
  }
  \node[font=\Large] at (4.7,0.1) {$\fieldInto$};
  \node[below=4pt] at (4.7,0.1) {$I_4$（环路外）};
  \node[align=left,anchor=west] at (6.0,0.5)
    {$\displaystyle\oint_C\vect B\cdot\dd\vect\ell=\mu_0 I_{\rm enc}$\\[4pt]
     环路外电流改变局部 $\vect B$，\\但不改变该环量};
\end{tikzpicture}
\end{center}}

\newcommand{\wireSolenoidDiagram}{%
\begin{center}
\begin{tikzpicture}[scale=0.88,>=Latex]
  \begin{scope}[xshift=-4.5cm]
    \draw[Ink,very thick] (0,-2.1)--(0,2.1);
    \draw[vector] (0,-1.4)--(0,-0.1) node[left] {$I$};
    \foreach \r in {0.7,1.35,2.0}{
      \draw[DeepBlue,thick,->] (\r,0) arc[start angle=0,end angle=145,x radius=\r,y radius=0.42*\r];
      \draw[DeepBlue,thick] (-\r,0) arc[start angle=180,end angle=360,x radius=\r,y radius=0.42*\r];
    }
    \draw[<->] (0,-2.5)--(2.0,-2.5);
    \node[below=2pt,inner sep=0.5pt] at (1.0,-2.5) {$r$};
    \node[TextGray] at (0,-3.05) {$B=\mu_0I/(2\pi r)$};
  \end{scope}
  \begin{scope}[xshift=3.4cm]
    \foreach \x in {0,0.45,...,5.4}{
      \draw[Ink,thick] (\x,-1.35) ellipse (0.25 and 1.35);
    }
    \draw[DeepBlue,->,thick] (-0.6,0.65)--(6.0,0.65);
    \draw[DeepBlue,->,thick] (-0.6,0)--(6.0,0);
    \draw[DeepBlue,->,thick] (-0.6,-0.65)--(6.0,-0.65);
    \draw[guide,dashed] (-0.4,-2.0) rectangle (3.7,2.0);
    \node[TextGray,align=center] at (2.7,2.45)
      {理想无限长密绕螺线管};
    \node[TextGray,align=center] at (2.7,-3.05)
      {$B_{\rm in}=\mu_0nI$，$B_{\rm out}=0$};
  \end{scope}
\end{tikzpicture}
\end{center}}

\newcommand{\biotSavartDiagram}{%
\begin{center}
\begin{tikzpicture}[scale=0.92,>=Latex]
  \draw[Ink,very thick] (-4,-1.2) .. controls (-2.3,1.8) and (-0.7,-0.4) .. (0.8,1.2);
  \coordinate (Q) at (-1.9,0.75);
  \coordinate (P) at (4.2,1.55);
  \fill[DeepBlue] (Q) circle (2.4pt);
  \fill[Ink] (P) circle (2.4pt) node[right] {$P$};
  \draw[vector] (Q)--++(1.25,0.25) node[above] {$I\dd\vect\ell$};
  \draw[vector,DeepBlue] (Q)--(P) node[midway,below] {$\vect r$};
  \node[font=\Large,DeepBlue] at ($(P)+(0,0.85)$) {$\fieldOut$};
  \node[right=3pt] at ($(P)+(0,0.85)$) {$\dd\vect B$};
  \node[align=left,anchor=west] at (-0.1,-1.2)
    {$\displaystyle\dd\vect B=\frac{\mu_0}{4\pi}
      \frac{I\,\dd\vect\ell\times\vect r}{r^3}$\\[5pt]
     稳恒电流必须闭合；最终对整个回路积分};
\end{tikzpicture}
\end{center}}

\newcommand{\loopAxisDiagram}{%
\begin{center}
\begin{tikzpicture}[scale=0.9,>=Latex]
  \draw[axis] (-1.0,0)--(8.2,0) node[right] {$z$};
  \draw[DeepBlue,very thick] (0,0) ellipse (0.8 and 2.05);
  \draw[vector] (0,-2.05) arc[start angle=-90,end angle=-15,x radius=0.8,y radius=2.05] node[right] {$I$};
  \fill[Ink] (0,0) circle (2.2pt) node[below left] {$O$};
  \fill[DeepBlue] (5.4,0) circle (2.5pt) node[below] {$P$};
  \coordinate (Q) at (0,1.45);
  \draw[vector,DeepBlue] (Q)--(5.4,0) node[midway,above] {$\boldsymbol\rho$};
  \draw[vector] (5.4,0)--++(1.25,0.8) node[right] {$\dd\vect B$};
  \draw[vector] (5.4,0)--++(1.45,0) node[right] {$\dd\vect B_z$};
  \draw[<->] (0,-2.45)--(5.4,-2.45);
  \draw[<->] (-0.25,0)--(-0.25,2.05);
  \node[below=2pt,inner sep=0.5pt] at (2.7,-2.45) {$z$};
  \node[left=3pt,inner sep=0.5pt] at (-0.25,1.03) {$R$};
  \node[TextGray,align=left,anchor=west] at (8.9,0.4)
    {横向分量成对抵消，\\仅轴向分量相加：\\[3pt]
     $\displaystyle B_z=\frac{\mu_0IR^2}{2(R^2+z^2)^{3/2}}$};
\end{tikzpicture}
\end{center}}

\newcommand{\helmholtzDiagram}{%
\begin{center}
\begin{tikzpicture}[scale=0.86,>=Latex]
  \draw[DeepBlue,very thick] (-2.5,0) ellipse (0.75 and 2.25);
  \draw[DeepBlue,very thick] (2.5,0) ellipse (0.75 and 2.25);
  \draw[axis] (-4.3,0)--(4.3,0) node[right] {$x$};
  \draw[<->] (-2.5,2.65)--(2.5,2.65);
  \draw[<->] (-2.5,0)--(-2.5,2.25);
  \node[above=2pt,inner sep=0.5pt] at (0,2.65) {$2D$};
  \node[left=3pt,inner sep=0.5pt] at (-2.5,1.13) {$R$};
  \draw[vector] (-2.55,-2.15) arc[start angle=-92,end angle=-18,x radius=0.75,y radius=2.25] node[right] {$NI$};
  \begin{scope}[xshift=7.1cm,yshift=-0.2cm]
    \draw[axis] (-2.8,0)--(2.8,0) node[right] {$x$};
    \draw[axis] (0,-0.2)--(0,2.6) node[above] {$B$};
    \draw[DeepBlue,very thick,domain=-2.5:2.5,samples=100]
      plot (\x,{1.55-0.03*\x*\x*\x*\x});
    \draw[TextGray,dashed] (-2.4,1.55)--(2.4,1.55);
    \node[TextGray] at (0,-0.6) {$2D=R$ 时 $B''(0)=0$};
  \end{scope}
\end{tikzpicture}
\end{center}}

\newcommand{\magneticGaussDiagram}{%
\begin{center}
\begin{tikzpicture}[scale=0.9,>=Latex]
  \shade[ball color=SoftGray,opacity=0.65] (0,0) ellipse (3.0 and 1.65);
  \draw[curve] (0,0) ellipse (3.0 and 1.65);
  \draw[Ink,very thick] (-0.55,0) rectangle (0.55,2.0);
  \node at (0,1.62) {N}; \node at (0,0.38) {S};
  \foreach \x in {-1.9,-1.25,-0.6,0.6,1.25,1.9}{
    \draw[DeepBlue,thick,->] (\x,1.25) .. controls (1.15*\x,2.25) and (1.15*\x,-1.0) .. (\x,-1.25);
  }
  \node[align=left,anchor=west] at (4.4,0.6)
    {磁感线没有起点或终点；\\穿出闭合面的通量必等于穿入量：\\[5pt]
     $\displaystyle\boxed{\oint_S\vect B\cdot\dd\vect S=0}$};
\end{tikzpicture}
\end{center}}

\newcommand{\finiteWireDiagram}{%
\begin{center}
\begin{tikzpicture}[scale=0.95,>=Latex]
  \coordinate (WireU) at (0,2.8);
  \coordinate (WireD) at (0,-2.6);
  \coordinate (WireH) at (0,0);
  \draw[Ink,very thick] (0,-2.6)--(0,2.8);
  \draw[vector] (0,-2.1)--(0,-0.8) node[left] {$I$};
  \coordinate (P) at (4.3,0);
  \fill[DeepBlue] (P) circle (2.6pt) node[right] {$P$};
  \draw[guide] (0,0)--(P) node[midway,below] {$r$};
  \draw[guide] (0,2.8)--(P) (0,-2.6)--(P);
  \draw pic[draw,"$\theta_1$",angle radius=9mm,angle eccentricity=1.25]
    {angle=WireU--P--WireH};
  \draw pic[draw,"$\theta_2$",angle radius=7mm,angle eccentricity=1.3]
    {angle=WireH--P--WireD};
  \node[font=\Large] at (5.35,0.8) {$\fieldInto$};
  \node[TextGray,align=left,anchor=west] at (6.0,-0.4)
    {$\displaystyle B=\frac{\mu_0I}{4\pi r}
       (\sin\theta_1+\sin\theta_2)$\\[3pt]
     无限长极限：$\theta_1=\theta_2=\pi/2$};
\end{tikzpicture}
\end{center}}

\newcommand{\currentSheetDiagram}{%
\begin{center}
\begin{tikzpicture}[scale=0.9,>=Latex]
  \coordinate (SheetL) at (-4.1,0);
  \coordinate (SheetR) at (3.2,0);
  \coordinate (SheetH) at (0,0);
  \draw[Ink,very thick] (SheetL)--(SheetR);
  \foreach \x in {-3.8,-3.0,...,2.6}{\node[font=\large] at (\x,-0.28) {$\fieldInto$};}
  \node[below] at (3.9,-0.65) {面电流 $j$};
  \coordinate (P) at (0,3.0);
  \fill[DeepBlue] (P) circle (2.5pt) node[above] {$A$};
  \draw[guide] (P)--(0,0);\node[left=4pt] at (0,.65) {$a$};
  \draw[guide] (P)--(-4.1,0) (P)--(3.2,0);
  \draw pic[draw,"$\theta_1$",angle radius=8mm,angle eccentricity=1.3]
    {angle=SheetL--P--SheetH};
  \draw pic[draw,"$\theta_2$",angle radius=11mm,angle eccentricity=1.25]
    {angle=SheetH--P--SheetR};
  \draw[vector] (P)--++(1.55,0) node[above] {$B_x$};
  \draw[vector] (P)--++(0,-1.8) node[below right=3pt] {$B_y$};
  \node[TextGray,align=left,anchor=west] at (5.3,1.4)
    {每条窄条等效为无限长直电流；\\左右分量积分后，非对称时还会有 $B_y$。\\[3pt]
     无限宽极限：$B=\mu_0j/2$};
\end{tikzpicture}
\end{center}}

\newcommand{\idealHelixDiagram}{%
\begin{center}
\begin{tikzpicture}[scale=0.85,>=Latex]
  \begin{scope}[xshift=-4.1cm]
    \foreach \y in {-2,-1.2,...,2}{\draw[DeepBlue,thick] (0,\y) ellipse (1.45 and 0.34);}
    \draw[DeepBlue,thick,->] ({1.45*cos(230)},{-2+.34*sin(230)}) arc(230:310:1.45 and .34);
    \node[below] at (0,-2.34) {$I$};
    \node[TextGray] at (0,-3.25) {“只旋不进”：纯方位面电流 $K_\varphi=nI$};
  \end{scope}
  \begin{scope}[xshift=4.1cm]
    \foreach \t in {-720,-360,0,360}{
      \draw[DeepBlue,dashed,thick,domain=\t:{\t+180},samples=61]
        plot ({1.45*cos(\x)},{0.0027*\x+.34*sin(\x)});
      \draw[DeepBlue,very thick,domain={\t+180}:{\t+360},samples=61,
        postaction={decorate},decoration={markings,mark=at position .55 with {\arrow{Latex}}}]
        plot ({1.45*cos(\x)},{0.0027*\x+.34*sin(\x)});}
    \draw[guide,dashed] (0,-2.2)--(0,2.2);
    \node[below] at (.2,-.7) {$I$};
    \draw[vector] (2.0,-0.9)--(2.0,0.7) node[right] {$K_z$};
    \node[TextGray] at (0,-3.25) {真实螺旋：另有轴向面电流 $K_z=I/(2\pi R)$};
  \end{scope}
\end{tikzpicture}
\end{center}}

\newcommand{\resistorGridForceDiagram}{%
\begin{center}
\begin{tikzpicture}[scale=0.86,>=Latex]
  \foreach \x in {0,1.5,3.0}{\draw[Ink,thick] (\x,0)--(\x,3.0);}
  \foreach \y in {0,1.5,3.0}{\draw[Ink,thick] (0,\y)--(3.0,\y);}
  \fill[DeepBlue] (0,0) circle (2.7pt) node[below right=4pt] {入};
  \fill[DeepBlue] (3,3) circle (2.7pt) node[below right=4pt] {出};
  \draw[vector,DeepBlue] (0,0)--(3,3) node[pos=.7,above left=3pt] {$\vect L$};
  \draw[Ink,thick] (-1,-.8)--(0,0);
  \draw[vector] (-.85,-.68)--(-.3,-.24) node[above left] {$I$};
  \draw[Ink,thick] (3,3)--(4,3.8);
  \draw[vector] (3.25,3.2)--(3.8,3.64) node[above left] {$I$};
  \begin{scope}[xshift=6.2cm,yshift=0.6cm]
    \coordinate (GridO) at (0,0);
    \coordinate (GridN) at (0,2.7);
    \coordinate (GridB) at (1.65,1.65);
    \draw[axis] (0,0)--(2.7,0) node[right] {$\vect e_1$};
    \draw[axis] (0,0)--(0,2.7) node[above] {$\vect n$};
    \draw[vector] (0,0)--(1.65,1.65) node[right] {$\vect B$};
    \node[below=10pt] at (1,0) {$\vect e_1$--$\vect n$ 平面侧视};
    \draw pic[draw,"$45^\circ$",angle radius=8mm,angle eccentricity=1.35]
      {angle=GridB--GridO--GridN};
    \node[TextGray,align=left,anchor=west] at (3.5,1.2)
      {$\displaystyle\sum_e I_e\vect\ell_e=I\vect L$\\[3pt]
       $\displaystyle\vect F=I\vect L\times\vect B$};
  \end{scope}
\end{tikzpicture}
\end{center}}

\newcommand{\currentElementsDiagram}{%
\begin{center}
\begin{tikzpicture}[scale=0.92,>=Latex]
  \coordinate (A) at (-3.6,-1.2); \coordinate (B) at (2.5,1.2);
  \draw[vector] (A)--++(1.4,1.0) node[above] {$I_1\dd\vect\ell_1$};
  \draw[vector] (B)--++(-1.0,1.35) node[above] {$I_2\dd\vect\ell_2$};
  \draw[vector,DeepBlue] (A)--(B) node[midway,below] {$\vect r$};
  \draw[vector] (A)--++(-0.7,1.35) node[left] {$\dd\vect F_{1\leftarrow2}$};
  \draw[vector] (B)--++(1.0,-0.9) node[right] {$\dd\vect F_{2\leftarrow1}$};
  \node[TextGray,align=center] at (-0.5,-2.4)
    {一般两电流元的两力既不等大反向，也不共线；孤立电流元不是可独立存在的物体};
\end{tikzpicture}
\end{center}}

\newcommand{\magneticLensDiagram}{%
\begin{center}
\begin{tikzpicture}[x=0.82cm,y=0.82cm,>=Latex]
  \draw[axis] (-7.0,0)--(7.4,0) node[right] {$z$};
  \draw[DeepBlue,very thick] (0,0) ellipse (0.35 and 2.25);
  \draw[vector] (0,-2.2) arc[start angle=-90,end angle=-15,x radius=0.35,y radius=2.25] node[right] {$I$};
  \draw[thick] (-5.6,-1.4)--(-5.6,1.4) node[above] {物面};
  \draw[thick] (5.0,-1.4)--(5.0,1.4) node[above] {像面};
  \fill[Ink] (-5.6,0) circle (2.3pt);
  \fill[Ink] (5.0,0) circle (2.3pt);
  \draw[DeepBlue,very thick] (-5.6,0)--(0,1.2)--(5.0,0);
  \draw[DeepBlue,very thick,dashed] (-5.6,0)--(0,-0.95)--(5.0,0);
  \draw[<->] (-5.6,-1.75)--(0,-1.75);
  \draw[<->] (0,-1.75)--(5.0,-1.75);
  \node[below=2pt,inner sep=0.5pt] at (-2.8,-1.75) {$u$};
  \node[below=2pt,inner sep=0.5pt] at (2.5,-1.75) {$v$};
  \draw[->,DeepBlue,thick] (0.75,2.35) arc[start angle=10,end angle=300,radius=0.45];
  \node[DeepBlue] at (1.65,2.1) {拉莫尔转角 $\phi_L$};
  \node[TextGray,align=left,anchor=west] at (7.8,0.7)
    {$\displaystyle\frac1u+\frac1v=\frac1f$\\[4pt]
     $\displaystyle\frac1f=\frac{q^2}{4m^2v_0^2}\int B_z^2\dd z$\\[4pt]
     $\displaystyle\phi_L=\frac{q}{2mv_0}\int B_z\dd z$};
\end{tikzpicture}
\end{center}}

\newcommand{\spinPrecessionDiagram}{%
\begin{center}
\begin{tikzpicture}[scale=0.72,>=Latex]
  \draw[axis] (0,-2.2)--(0,3.0) node[above] {$\vect B_0$};
  \draw[DeepBlue,thick,dashed] (0,1.0) ellipse (2.2 and 0.65);
  \draw[vector] (0,0)--(1.85,1.65) node[right] {$\vect s\parallel\vect\mu$};
  \draw[->,DeepBlue,very thick] (1.45,1.55) arc[start angle=30,end angle=150,x radius=1.35,y radius=0.48];
  \node[DeepBlue] at (0,2.25) {$\omega_0=\gamma B_0$};
  \begin{scope}[xshift=6.4cm]
    \draw[axis] (0,-2.2)--(0,3.0) node[above] {$\vect B_0$};
    \draw[vector,DeepBlue] (0,0)--(2.0,0) node[right] {$\vect B_1$};
    \draw[vector] (0,0)--(0,2.15) node[left] {初始 $\vect s$};
    \draw[->,Ink,very thick] (0,2.15) arc[start angle=90,end angle=-80,radius=2.15];
    \node[TextGray,align=center] at (0,-2.8)
      {随动系：自旋以 $\Omega_R=\gamma B_1$ 绕 $\vect B_1$ 翻转};
  \end{scope}
\end{tikzpicture}
\end{center}}

% Geometry- and sign-checked replacements, September 2026.
\newcommand{\straightConductorDiagram}{%
\begin{center}\begin{tikzpicture}[scale=.85,>=Latex]
 \draw[Ink,very thick](-2,0)--(2,0);
 \fill(-2,0)circle(2pt)node[left]{$a$};\fill(2,0)circle(2pt)node[right]{$b$};
 \draw[vector](-1.4,0)--(-.4,0)node[midway,below]{$I$};
 \draw[vector](0,0)--(0,1.6)node[above]{$\vect F$};
 \foreach \x in {-1.5,-.7,.7,1.5}{\node[DeepBlue]at(\x,.85){$\otimes$};}
 \node[DeepBlue]at(2.25,.85){$\vect B$};
 \draw[<->](-2,-.8)--(2,-.8)node[midway,below]{$L$};
 \node[anchor=west]at(3.2,.2){$\vect F=I\vect L_{ab}\times\vect B$，$F=IBL$};
\end{tikzpicture}\end{center}}
\newcommand{\finiteCoilFieldDiagram}{%
\begin{center}\begin{tikzpicture}[scale=.7,>=Latex]
 \foreach \sg in {-1,1}{\foreach \a/\b/\h in {.2/2.6/3.8,.6/1.85/2.8}{
  \draw[DeepBlue,thick,postaction={decorate},decoration={markings,mark=at position .5 with {\arrow{Latex}}}]
   ({\sg*\a},0)..controls({\sg*\a},\h)and({\sg*\b},\h)..({\sg*\b},0)
   ..controls({\sg*\b},-\h)and({\sg*\a},-\h)..cycle;}}
 \foreach \y in {-1.3,-.975,...,1.3}{\draw[Ink](0,\y)ellipse(1.15 and .25);}
 \draw[vector](0,-.8)--(0,3.25)node[above]{$\vect B$};
 \draw[Ink,very thick,->]({1.15*cos(230)},{-1.3+.25*sin(230)})arc(230:310:1.15 and .25);
 \node[below]at(0,-1.65){$I$};
 \node[align=left,anchor=west]at(3.4,.2){内场与外部回流场相接\\箭头表明场线方向\\中部匀强只是一种近似};
\end{tikzpicture}\end{center}}
\renewcommand{\loopTorqueDiagram}{%
\begin{center}\begin{tikzpicture}[scale=.82,>=Latex]
 \begin{scope}[xshift=-3.8cm]
  \coordinate (A) at (1,-1.2);\coordinate (C) at (-1,1.2);
  \draw[Ink,thick] (1.65,-.875)--(-.35,1.525)--(-1.65,.875)--(.35,-1.525)--cycle;
  \draw[guide,dashed] (-1.3,-.65)--(1.3,.65) node[right]{转轴};
  \draw[vector] (.35,-1.525)--(1.3,-1.05);
  \node[below left] at (.6,-1.5){$I$};
  \draw[vector] (A)--++(0,-1.2) node[right]{$\vect F_A$};
  \draw[vector] (C)--++(0,1.2) node[left]{$\vect F_C$};
  \node[right] at (1.65,-1.3){A};\node[above left=3pt] at (C){C};
  \draw[vector] (0,0)--(1.5,1.25) node[above]{$\vect\mu$};
  \draw[vector] (-2.0,3.0)--(1.7,3.0) node[right]{$\vect B$};
  \node at (0,-3.35){沿轴两边的力构成力偶};
 \end{scope}
 \begin{scope}[xshift=3.8cm]
  \draw[Ink,very thick] (1,-1.2)--(-1,1.2);
  \node[right] at (1,-1.2){A};\node[left] at (-1,1.2){C};
  \draw[vector] (1,-1.2)--(1,-2.5) node[right]{$IbB$};
  \draw[vector] (-1,1.2)--(-1,2.5) node[left]{$IbB$};
  \draw[guide,dashed] (-1,-1.3)--(-1,-2.0);
  \draw[<->] (-1,-1.85)--(1,-1.85) node[midway,above]{$a\sin\theta$};
  \draw[vector] (0,0)--(1.8,0) node[right]{$\vect B$};
  \draw[vector] (0,0)--(1.2,1.0) node[above]{$\vect n$};
  \draw (.65,0) arc(0:39.806:.65);\node at (.88,.27){$\theta$};
  \node at (0,-3.35){沿转轴观察的侧视图};
 \end{scope}
\end{tikzpicture}\end{center}}

\renewcommand{\wireSolenoidDiagram}{%
\begin{center}\begin{tikzpicture}[scale=.78,>=Latex]
 \begin{scope}[xshift=-4.4cm]
  \foreach \r in {.65,1.25,1.85}{
   \draw[DeepBlue,dashed] (\r,0) arc(0:180:{\r} and {.36*\r});
   \draw[DeepBlue,thick] (-\r,0) arc(180:360:{\r} and {.36*\r});
   \draw[DeepBlue,->] (270:{\r} and {.36*\r}) arc(270:310:{\r} and {.36*\r});}
  \draw[Ink,thick] (0,-1.9)--(0,2.0);\draw[vector] (0,.65)--(0,1.65) node[right]{$I$};
  \draw[<->] (0,-2.15)--(1.85,-2.15) node[midway,below]{$r$};
  \node at (0,-2.95){$B=\mu_0I/(2\pi r)$};
 \end{scope}
 \begin{scope}[xshift=.0cm]
  \foreach \x in {0,.45,...,5.4}{\draw[Ink] (\x,0) ellipse(.18 and 1.0);}
  \foreach \y in {-.5,0,.5}{\draw[DeepBlue,->,thick] (-.3,\y)--(5.75,\y);}
  \draw[guide,dashed] (.7,.2) rectangle(4.7,1.85);
  \node[above] at (2.7,1.9){外部段：$B=0$};
  \draw[<->] (.7,-1.45)--(4.7,-1.45) node[midway,below]{$\ell$};
  \node at (2.7,-2.95){内部段：$B=\mu_0nI$};
  \node[above] at (2.7,2.6){无限长模型（仅画一段）};
 \end{scope}
\end{tikzpicture}\end{center}}

\renewcommand{\biotSavartDiagram}{%
\begin{center}\begin{tikzpicture}[scale=.9,>=Latex]
 \draw[Ink,thick,domain=-2:2,samples=50] plot(\x,{.18*\x*\x});
 \coordinate(Q)at(0,0);\coordinate(P)at(3.5,1.6);
 \fill(Q)circle(2pt)node[below left=3pt]{$Q$};\fill(P)circle(2pt)node[right]{$P$};
 \draw[vector](Q)--(1.1,0)node[below right=3pt]{$I\dd\vect\ell$};
 \draw[vector](Q)--(P)node[midway,above]{$\vect r$};
 \node[DeepBlue]at(3.5,2.35){$\odot\ \dd\vect B$};
 \node[align=left,anchor=west]at(5.0,.8){$Q$ 在导线上\\$\dd\vect\ell$ 沿切线\\$\dd\vect\ell\times\vect r$ 指向纸外};
\end{tikzpicture}\end{center}}

\renewcommand{\loopAxisDiagram}{%
\begin{center}\begin{tikzpicture}[scale=.86,>=Latex]
 \draw[axis](-1,0)--(6,0)node[right]{$z$};
 \draw[DeepBlue,thick](0,0)ellipse(.6 and 1.6);
 \coordinate(Q)at(0,1.6);\coordinate(Qp)at(0,-1.6);\coordinate(P)at(4,0);
 \fill(Q)circle(2pt)node[above right]{$Q$};\fill(Qp)circle(2pt)node[below right]{$Q'$};
 \fill(P)circle(2pt)node[below left=4pt]{$P$};\node[below left]at(0,0){$O$};
 \draw[vector](Q)--(-.8,1.6)node[above]{$I\dd\vect\ell$};
 \draw[vector](Qp)--(.8,-1.6);
 \draw[guide](Q)--(P)node[midway,above]{$\boldsymbol\rho$};\draw[guide](Qp)--(P);
 \draw[vector](P)--(4.64,1.6)node[above]{$\dd\vect B_Q$};
 \draw[vector](P)--(4.64,-1.6)node[below]{$\dd\vect B_{Q'}$};
 \draw[guide,dashed](4.64,-1.6)--(4.64,1.6);
 \draw[<->](-1.25,0)--(-1.25,1.6)node[midway,left]{$R$};
 \draw[<->](0,-2.2)--(4,-2.2)node[midway,below]{$z$};
 \node[align=left,anchor=west]at(6.5,.1){对径元横向分量相消\\轴向分量相加\\$\rho=\sqrt{R^2+z^2}$};
\end{tikzpicture}\end{center}}

\renewcommand{\helmholtzDiagram}{%
\begin{center}\begin{tikzpicture}[scale=.85,>=Latex]
 \begin{scope}[xshift=-3.8cm]
  \draw[axis](-2.6,0)--(2.7,0)node[right]{$x$};
  \foreach \x in {-1,1}{\draw[Ink,thick](\x,0)ellipse(.32 and 1.9);
   \draw[DeepBlue,->]({\x+.32},0)arc(0:70:.32 and 1.9);}
  \draw[<->](-1,2.3)--(1,2.3)node[midway,above]{$2D$};
  \draw[<->](-1.6,0)--(-1.6,1.9)node[midway,left]{$R$};
  \node at(0,-2.4){两个线圈同向通电};
 \end{scope}
 \begin{scope}[xshift=3.8cm,x=2.4cm,y=2.2cm]
  \draw[axis](-1,0)--(1.15,0)node[right]{$x/R$};
  \draw[axis](0,0)--(0,1.9)node[above]{$BR/(\mu_0NI)$};
  \foreach \d/\col/\sty in {.4/DeepBlue/solid,.5/Ink/solid,.6/TextGray/dashed}{
   \draw[\col,\sty,thick,domain=-.95:.95,samples=90]plot(\x,{.5*((1+(\x-\d)^2)^(-1.5)+(1+(\x+\d)^2)^(-1.5))});}
  \node[anchor=west,font=\small,DeepBlue]at(.28,1.62){$D=0.4R$};
  \node[anchor=west,font=\small]at(.28,1.38){$D=0.5R$};
  \node[anchor=west,font=\small,TextGray]at(.28,1.14){$D=0.6R$};
 \end{scope}
\end{tikzpicture}\end{center}}

\renewcommand{\currentElementsDiagram}{%
\begin{center}\begin{tikzpicture}[scale=.9,>=Latex]
 \coordinate(A)at(0,0);\coordinate(B)at(4,2);
 \fill(A)circle(2pt)node[below]{$1$};\fill(B)circle(2pt)node[below right=4pt]{$2$};
 \draw[vector](A)--(0,1.5)node[above]{$I_1\dd\vect\ell_1$};
 \draw[vector](B)--(5.5,2)node[right]{$I_2\dd\vect\ell_2$};
 \draw[vector](B)--(A)node[midway,below right=4pt]{$\vect r=\vect r_1-\vect r_2$};
 \draw[vector](A)--(-1.6,0)node[above]{$\dd\vect F_{1\leftarrow2}$};
 \draw[vector](B)--(4,5.2)node[above]{$\dd\vect F_{2\leftarrow1}$};
 \node[align=left,anchor=west]at(6.6,.2){两力各垂直本地电流元\\一般不等大反向\\本图两元均在纸面内};
\end{tikzpicture}\end{center}}

\renewcommand{\magneticLensDiagram}{%
\begin{center}\begin{tikzpicture}[scale=.82,>=Latex]
 \draw[axis](-4.6,0)--(4.8,0)node[right]{$z$};
 \draw[Ink](-4,-.9)--(-4,2.0)node[above]{物面};\draw[Ink](4,-.9)--(4,2.0)node[above]{像面};
 \draw[DeepBlue,thick](-4,0)--(0,1.1)--(4,0);
 \draw[DeepBlue,dashed,thick](-4,0)--(0,-1.1)--(4,0);
 \draw[Ink,thick](0,0)ellipse(.25 and 1.7);
 \draw[<->](-4,-2)--(0,-2)node[midway,below]{$u$};
 \draw[<->](0,-2)--(4,-2)node[midway,below]{$v$};
 \node[align=left,anchor=west]at(5.3,.25){$\Delta\xi'=-\xi/f$\\[4pt]
 $\displaystyle\phi_L=-\frac{q}{2mv_0}\int B_z\dd z$};
 \node[above]at(0,2.1){拉莫尔随动平面内的射线};
\end{tikzpicture}\end{center}}

\renewcommand{\spinPrecessionDiagram}{%
\begin{center}\begin{tikzpicture}[scale=.78,>=Latex]
 \begin{scope}[xshift=-3.5cm]
  \draw[axis](0,-.5)--(0,3.1)node[above]{$\vect B_0$};
  \draw[DeepBlue,dashed](0,1.5)ellipse(1.7 and .48);
  \draw[vector](0,0)--({1.7*cos(35)},{1.5+.48*sin(35)})node[right]{$\vect s$};
  \draw[vector](-2.2,1.8)--(-2.2,.6)node[midway,left]{$\vect\omega_0$};
  \node[align=center]at(0,-1.25){实验室系\\$\vect\omega_0=-\gamma\vect B_0$};
 \end{scope}
 \begin{scope}[xshift=3.5cm]
  \draw[axis](-.5,0)--(2.8,0)node[right]{$y$};
  \draw[axis](0,-2)--(0,2.7)node[above]{$z$};
  \node[DeepBlue,anchor=east]at(-.3,-.55){$\odot\ B_1$};
  \draw[vector](0,0)--(0,1.7)node[left]{初始 $\vect s$};
  \draw[vector](0,0)--(1.7,0)node[above right]{$\pi/2$ 脉冲后};
  \draw[vector](0,0)--(0,-1.7)node[left]{$\pi$ 脉冲后};
  \draw[DeepBlue,->,thick](0,1.7)arc(90:-90:1.7);
  \node[align=center]at(.2,-2.7){随动系，沿 $-\vect B_1$ 观察\\转动平面为 $yz$ 面，$s_x=0$};
 \end{scope}
\end{tikzpicture}\end{center}}

\renewcommand{\magneticGaussDiagram}{%
\begin{center}\begin{tikzpicture}[scale=.8,>=Latex]
 \draw[Ink,thick](0,0)rectangle(.8,2.4);\node at(.4,1.95){N};\node at(.4,.45){S};
 \foreach \sg in {-1,1}{\foreach \a/\b/\h in {.08/2.1/3.8,.2/1.4/3.4,.32/.9/3.0}{
  \draw[DeepBlue,postaction={decorate},decoration={markings,mark=at position .5 with {\arrow{Latex}}}]
   ({.4+\sg*\a},2.4)..controls({.4+\sg*\a},\h)and({.4+\sg*\b},\h)..({.4+\sg*\b},1.2)
   ..controls({.4+\sg*\b},{2.4-\h})and({.4+\sg*\a},{2.4-\h})..({.4+\sg*\a},0);}}
 \draw[DeepBlue,->](.4,.8)--(.4,1.6);
 \draw[guide,dashed](-.4,.75)rectangle(1.2,2.7);
 \node[align=left,anchor=west]at(3.7,1.1){闭合面 $S$\\穿出与穿入磁通相消\\磁体内部由S极指向N极};
\end{tikzpicture}\end{center}}

\newcommand{\handSolenoidContext}{%
 \historyfigure[.84\linewidth]{right-hand-solenoid.png}{左为右手握持规则的手势，右为空心密绕螺线管的实物结构。}}

\newcommand{\fluxPatchDiagram}{%
\begin{center}\begin{tikzpicture}[scale=.78,>=Latex]
 \draw[Ink,thick]plot[smooth cycle]coordinates{(-2,-1)(-2.3,.6)(-.8,1.1)(1.3,.7)(2.2,-.5)(.3,-1.2)};
 \fill[Gold!35](-.35,-.15)--(.2,-.3)--(.65,.15)--(.1,.3)--cycle;
 \draw[vector](.1,0)--(.1,1.8)node[above]{$\dd\vect S$};
 \draw[vector](.1,0)--(1.75,1.4)node[right]{$\vect B$};
 \draw(.1,.6)arc(90:40:.6);\node at(.35,.58){$\theta$};
 \node[below]at(.1,-.45){$\dd S$};
 \node[align=left,anchor=west]at(3.2,.3){$\dd\Phi_B=B\dd S\cos\theta$\\法向与曲面局部切平面垂直};
\end{tikzpicture}\end{center}}

\newcommand{\motorStagesDiagram}{%
\begin{center}\begin{tikzpicture}[scale=.67,>=Latex]
 \foreach \ang/\xx/\lab in {25/0/1,75/4.2/2,115/8.4/3,165/12.6/4}{
  \begin{scope}[xshift=\xx cm]
   \draw[guide,->](-1.6,1.8)--(1.6,1.8)node[midway,above]{$B$};
   \coordinate(A)at({1.25*cos(\ang)},{1.25*sin(\ang)});
   \coordinate(C)at({-1.25*cos(\ang)},{-1.25*sin(\ang)});
   \draw[Ink,thick](A)--(C);\draw[Ink](0,0)circle(.27);
   \draw[Ink,very thick](-.5,0)--(-.27,0) (.27,0)--(.5,0);
   \ifnum\ang<90
    \node at(A){$\odot$};\node at(C){$\otimes$};
    \draw[vector](A)--++(0,.9);\draw[vector](C)--++(0,-.9);
   \else
    \node at(A){$\otimes$};\node at(C){$\odot$};
    \draw[vector](A)--++(0,-.9);\draw[vector](C)--++(0,.9);
   \fi
   \draw[DeepBlue,->](.9,0)arc(0:55:.9);
   \node at(0,-2.35){阶段\lab};
  \end{scope}}
\end{tikzpicture}\end{center}}

\newcommand{\chargeMonopoleComparison}{%
\begin{center}\begin{tikzpicture}[scale=.65,>=Latex]
 \foreach \yy/\f/\signa/\signb in {2.6/E/+/-, -2.6/B/N/S}{
  \foreach \xx/\sg/\dr in {0/\signa/1,4.5/\signb/-1}{
   \begin{scope}[shift={(\xx,\yy)}]
    \draw[Ink](0,0)circle(.32);\node at(0,0){$\sg$};
    \foreach \a in {0,45,...,315}{\ifnum\dr=1
     \draw[DeepBlue,->](\a:.42)--(\a:1.45);
     \else\draw[DeepBlue,->](\a:1.45)--(\a:.42);\fi}
    \node[above]at(0,1.65){$\vect\f$};
   \end{scope}}
  \begin{scope}[shift={(9.6,\yy)}]
   \draw[vector](-1.6,0)--(1.9,0)node[right]{$\vect v$};
   \node[above]at(0,1.55){\ifdim\yy pt>0pt 正电荷\else 假想正磁荷\fi};
   \foreach \x in {-.9,0,.9}{
    \draw[DeepBlue,dashed] (\x,1.35)arc(90:270:.23 and 1.35);
    \ifdim\yy pt>0pt
     \draw[DeepBlue,->](\x,1.35)arc(90:-90:.23 and 1.35);
    \else
     \draw[DeepBlue,->](\x,-1.35)arc(270:450:.23 and 1.35);
    \fi}
   \node[below]at(0,-1.6){\ifdim\yy pt>0pt 环绕的磁场 $\vect B$\else 假想环绕电场 $\vect E$\fi};
  \end{scope}}
 \node[align=center]at(4.8,-5.8){下排仅是假设磁荷存在的类比；普通磁铁的N、S极不是孤立磁荷。};
\end{tikzpicture}\end{center}}


第22章从单个带电粒子的洛伦兹力出发研究轨道；本章把许多载流子的受力合并起来，得到通电
导线的安培力（Ampere force），再反过来研究稳恒电流怎样产生磁场。两条主线分别回答“磁场怎样作用于电流”
和“电流怎样产生磁场”，最终由磁矩、安培环路定律（Ampere circuital law）与毕奥--萨伐尔定律（Biot--Savart law）连接起来。

需要特别提醒，这个讲授顺序并不是历史发现顺序。实验上，人们先看到的是通电导线在磁场中受力，
以及电流会在周围产生磁场；洛伦兹力则是在后来用带电粒子的微观运动解释电流与磁场现象时逐渐形成
的。现在先学过单个粒子的受力，再把大量载流子的力相加，能够直接看见安培力的来源，也能避免把
“电流元受力公式”当成毫无来由的新定律。电动机依靠的正是这种分布在整段载流导线上的力及其力矩，
它把电能持续转化为机械能，对现代工业的意义也来自这里。

磁场与静电场还有一个根本差别：迄今没有发现可以像点电荷那样单独存在的磁单极子。经典磁场的来源是
运动电荷形成的电流以及微观粒子的磁矩；磁感线没有磁荷式的起点、终点，可以闭合，也可以延伸至无穷远。
本章先研究普通电流的磁场，后面再借助相同的磁偶极图像理解物质磁性与粒子自旋的磁效应。

从结构上看，本章有两组容易混淆但方向相反的问题。已知外磁场，求载流导线受力，要使用
$\dd\vect F=I\dd\vect\ell\times\vect B$；已知电流分布，求它在空间中产生的磁场，则使用安培
环路定律或毕奥--萨伐尔定律。前者的结果是力或力矩，后者的结果是磁感应强度。实际问题常把二者
串联起来：先由一个回路的电流求磁场，再由该磁场求另一个回路的受力。每一步都应写清哪一部分是
“场源”，哪一部分是“受力对象”，否则极易把同一根导线产生的自场直接代回自身受力公式而产生歧义。


\subsection{毕奥--萨伐尔定律与圆电流}


安培环路定律给出磁场的环量，却不总能直接算出每一点的场。对任意稳恒闭合电流，
毕奥--萨伐尔定律给出电流元在场点 $P$ 产生的磁场
\begin{equation}
  \boxed{\dd\vect B=\frac{\mu_0}{4\pi}
  \frac{I\,\dd\vect\ell\times\vect r}{r^3}},
\end{equation}
其中 $\vect r$ 从电流元指向场点。必须沿整个闭合回路积分：孤立的稳恒电流元并不存在。

毕奥与萨伐尔在1820年前后从实验中总结出这一关系。它在形式上类似库仑定律：源的贡献随距离
平方反比衰减，并可叠加；但方向结构不同，$\dd\vect B$ 同时垂直于电流元和由电流元指向场点的矢量。
分母写作 $r^3$ 是为了同分子的矢量 $\vect r$ 配合，大小仍为
$\dd B=(\mu_0/4\pi)I\dd\ell\sin\theta/r^2$。使用时必须先统一 $\dd\vect\ell$、$\vect r$ 的方向，
再沿完整回路积分，不能把一个虚构的孤立电流元当成可以单独存在的磁场源。

毕奥--萨伐尔定律适用于给定稳恒电流分布。它具有叠加性：若空间中有多个回路，总磁场等于各回路
积分结果的矢量和。计算时可以按四步进行：先写电流元 $I\dd\vect\ell$ 的方向；再写从源点到场点的
$\vect r$；用叉积确定 $\dd\vect B$；最后只把不被对称性抵消的分量积分。常见错误是把 $\vect r$
反向写成从场点指向源点，这会使整个磁场方向反号；或只积分 $\dd B$ 的大小而忘记不同电流元的矢量
方向并不相同。

图~\ref{ch23:fig:review-biotsavartdiagram-1}中，源点Q位于导线上，线元取Q处切线；从Q到P的矢量与线元叉乘指向纸外。
\par\noindent\begin{minipage}{\linewidth}
\biotSavartDiagram
\figcap{源点、切线和磁场叉积方向\label{ch23:fig:review-biotsavartdiagram-1}}
\end{minipage}\par\medskip

半径为 $R$ 的圆电流在轴线上距圆心 $z$ 处，横向分量成对抵消，只留下轴向分量：
\begin{equation}
  \boxed{B_z(z)=\frac{\mu_0IR^2}{2(R^2+z^2)^{3/2}}}.
\end{equation}
远场 $|z|\gg R$ 时，令 $\mu=I\pi R^2$，有
\begin{equation}
  B_z\simeq\frac{\mu_0}{4\pi}\frac{2\mu}{|z|^3},
\end{equation}
这与磁偶极子的轴向远场一致。

轴线上每个电流元到场点距离相同。把单个 $\dd\vect B$ 分解后，垂直轴线的分量总能在圆环对径位置找到
大小相等、方向相反的伙伴，只有轴向分量叠加。圆心处退化为 $B=\mu_0I/(2R)$；远处则按 $1/|z|^3$
衰减，比点电荷的 $1/r^2$ 更快，这正反映了磁场没有单极子项，最先出现的是磁偶极项。

轴上结果还可做三项检验。第一，量纲上 $\mu_0I/R$ 才具有磁感应强度量纲；第二，$z=0$ 时所有
电流元贡献方向完全一致，应得到圆心场 $\mu_0I/(2R)$；第三，$z\to\infty$ 时场趋于零且关于 $z$
是偶函数。若推导结果在轴线两侧变号，或远场仍按 $1/z^2$ 衰减，就说明轴向投影或对径抵消处理有误。
远场只保留 $\mu=I\pi R^2$，意味着观察距离足够大时无法辨认回路半径和电流各自是多少，只能辨认
它们组成的磁矩。

图~\ref{ch23:fig:review-loopaxisdiagram-1}中，Q与对径点 $Q\prime$ 同在圆环上，到P距离相同；两个场的横向分量相消，轴向分量相加。
\par\noindent\begin{minipage}{\linewidth}
\loopAxisDiagram
\figcap{圆电流对径线元的轴向叠加\label{ch23:fig:review-loopaxisdiagram-1}}
\end{minipage}\par\medskip

两个半径均为 $R$、每个有 $N$ 匝的同轴线圈，圆心位于 $x=\pm D$，通以同向电流 $I$。轴线上
的合场为
\begin{equation}
  B(x)=\frac{\mu_0NIR^2}{2}
  \left\{
  \frac{1}{[(x-D)^2+R^2]^{3/2}}
  +\frac{1}{[(x+D)^2+R^2]^{3/2}}
  \right\}.
\end{equation}
对称性自动给出 $B'(0)=0$；再令 $B''(0)=0$，得到 $D=R/2$，即两线圈间距 $2D=R$。
这就是亥姆霍兹线圈能在中心产生高均匀度磁场的原因。

这里两线圈的中心分别在 $x=\pm D$，所以实际间距是 $2D$。对称性使中心的一阶梯度自动为零；进一步
令二阶导数也为零，得到间距等于线圈半径。这样中心附近场强的首个非零空间变化被推到更高阶，在一个
有限区域内都比单线圈均匀。实验室常用亥姆霍兹线圈校准磁传感器、补偿地磁场或为带电粒子实验提供
近似匀强磁场；均匀性来自成对线圈的几何消差，而不是把线圈简单做大。

在中心附近作 Taylor 展开更能看清“均匀”的含义。由于 $B(x)=B(-x)$，所有奇次项自动为零；一般
线圈间距下首个变化项是 $x^2$。当 $2D=R$ 时二次项也消失，展开从 $x^4$ 项开始，因而小范围内的
相对变化远小于单线圈。中心磁场为
\begin{equation}
  B(0)=\left(\frac45\right)^{3/2}\frac{\mu_0NI}{R}.
\end{equation}
若两线圈电流不等、半径不等或安装轴线偏斜，原本由对称性消去的一阶梯度会重新出现，所以真实装置
的均匀度还依赖绕组一致性和几何校准。

图~\ref{ch23:fig:review-helmholtzdiagram-1}中，两圆心分别在 $\pm D$；三条曲线均由同一个轴场表达式计算，$D=0.5R$ 的中心二阶曲率为零。
\par\noindent\begin{minipage}{\linewidth}
\helmholtzDiagram
\figcap{亥姆霍兹线圈间距与三条轴场曲线\label{ch23:fig:review-helmholtzdiagram-1}}
\end{minipage}\par\medskip

\motorHelmholtzContext


% EM teaching enrichment 20260927
\paragraph{圆弧与径向连接段的分工}
半径 $a$、张角 $\alpha$（弧度）的圆弧载电流 $I$，求圆心处磁场。每个线元与源点至圆心方向垂直，贡献方向相同，故
\[
 B_{\rm arc}=\frac{\mu_0I}{4\pi a^2}\int\dd l
 =\frac{\mu_0I\alpha}{4\pi a}.
\]
若连接线段沿经过圆心的径向，因 $\dd\boldsymbol l\parallel\boldsymbol r$，这些段对圆心的贡献为零。
回路仍须闭合；圆弧公式只是总积分的一部分。对两个半径 $a,b$ 的同心弧组成的扇环回路，两弧绕向相反，圆心磁场大小为
$\mu_0I\alpha|1/a-1/b|/(4\pi)$，方向由较小半径那段决定。


% EM structure and teaching revision 20260928
\subsection{有限直导线的积分与端角约定}
令导线沿 $z$ 轴延伸，观察点到轴的垂距为 $a>0$。电流元到观察点的距离为 $\sqrt{a^2+z^2}$，各元产生的磁场方向相同，故
\[
 B=\frac{\mu_0I}{4\pi}\int_{z_1}^{z_2}\frac{a\dd z}{(a^2+z^2)^{3/2}}
 =\frac{\mu_0I}{4\pi a}\left[\frac z{\sqrt{a^2+z^2}}\right]_{z_1}^{z_2}.
\]
若端角从观察点到垂足的方向量起，$z/\sqrt{a^2+z^2}$ 对应带符号的正弦；若从导线方向量起，则应换为余弦。不同教材的端角公式看似不同，往往只是角度定义不同。
令两端趋向正负无穷，结果为 $\mu_0I/(2\pi a)$；令一端在垂足、另一端趋向无穷，结果减半。观察点位于垂足以外时，两端贡献是差而不是两个无符号正弦的和。
对折线或圆弧回路，各段都必须对同一个观察点算矢量再叠加。导线与源点到观察点的方向平行时，该段在毕奥--萨伐尔积分中的贡献为零，但这不意味着该段电流本身不存在。


\subsection{专题例题与方法应用}

\begin{origexamplebox}{【长直电流线】}\label{ch23:ex:s39-8-1}
一根有限长直导线中通有电流 $I$。场点 $P$ 到导线的垂直距离为 $r$，从 $P$ 看向导线两端与
垂线的夹角分别为 $\theta_1,\theta_2$。利用毕奥--萨伐尔定律求 $P$ 点的磁感应强度；再令
导线无限长，并与安培环路定律比较。

% Pure source-question geometry: no derived forces, currents or solution curves.
%
\begin{center}
\begin{tikzpicture}[scale=0.95,>=Latex]
  \coordinate (WireU) at (0,2.8);
  \coordinate (WireD) at (0,-2.6);
  \coordinate (WireH) at (0,0);
  \draw[Ink,very thick] (0,-2.6)--(0,2.8);
  \draw[vector] (0,-2.1)--(0,-0.8) node[left] {$I$};
  \coordinate (P) at (4.3,0);
  \fill[DeepBlue] (P) circle (2.6pt) node[right] {$P$};
  \draw[guide] (0,0)--(P) node[midway,below] {$r$};
  \draw[guide] (0,2.8)--(P) (0,-2.6)--(P);
  \draw pic[draw,"$\theta_1$",angle radius=9mm,angle eccentricity=1.25]
    {angle=WireU--P--WireH};
  \draw pic[draw,"$\theta_2$",angle radius=7mm,angle eccentricity=1.3]
    {angle=WireH--P--WireD};
  \end{tikzpicture}\end{center}
\figcap{有限长直电流及观察点的端点角。\label{ch23:fig:em23-question-1}}
\end{origexamplebox}

\begin{xesolution}
取导线沿 $y$ 轴，垂足为 $y=0$。长度元 $\dd y$ 到 $P$ 的距离为
$\rho=\sqrt{r^2+y^2}$，且 $\dd\vect\ell\times\vect\rho$ 的大小为 $r\,\dd y$。于是
\[
  \dd B=\frac{\mu_0I}{4\pi}\frac{r\,\dd y}{(r^2+y^2)^{3/2}}.
\]
由 $y=r\tan\theta$，或直接积分，得到
\begin{equation}
  \boxed{B=\frac{\mu_0I}{4\pi r}
  (\sin\theta_1+\sin\theta_2)}.
\end{equation}
方向由右手螺旋法则决定；图示电流向上且 $P$ 在导线右侧时，$\vect B$ 指向纸内。

当导线两端都趋于无穷远，$\theta_1=\theta_2=\pi/2$，故
\[
  \boxed{B=\frac{\mu_0I}{2\pi r}},
\]
与以导线为中心、半径为 $r$ 的圆形安培环路所得结果完全一致。

图~\ref{ch23:fig:review-finitewirediagram-1}中，两个端点角以垂线为基准；场点P在电流右侧，右手规则给出纸内磁场。
\par\noindent\begin{minipage}{\linewidth}
\finiteWireDiagram
\figcap{有限直电流的两端角度\label{ch23:fig:review-finitewirediagram-1}}
\end{minipage}\par\medskip
\end{xesolution}



\begin{origexamplebox}{【半无限大电流面】}\label{ch23:ex:s39-8-2}
一个方向无限长、横向宽度有限的导体面内通有均匀面电流，单位横向长度所分配的电流为 $j$。
场点 $A$ 到导体面的距离为 $a$，从 $A$ 看向左右边缘与垂线的夹角为 $\theta_1,\theta_2$。
求 $A$ 点的磁感应强度；再令导体面横向也无限宽，并与安培环路定律比较。

% Pure source-question geometry: no derived forces, currents or solution curves.
%
\begin{center}
\begin{tikzpicture}[scale=0.9,>=Latex]
  \coordinate (SheetL) at (-4.1,0);
  \coordinate (SheetR) at (3.2,0);
  \coordinate (SheetH) at (0,0);
  \draw[Ink,very thick] (SheetL)--(SheetR);
  \foreach \x in {-3.8,-3.0,...,2.6}{\node[font=\large] at (\x,-0.28) {$\fieldInto$};}
  \node[below] at (3.9,-0.65) {面电流 $j$};
  \coordinate (P) at (0,3.0);
  \fill[DeepBlue] (P) circle (2.5pt) node[above] {$A$};
  \draw[guide] (P)--(0,0);\node[left=4pt] at (0,.65) {$a$};
  \draw[guide] (P)--(-4.1,0) (P)--(3.2,0);
  \draw pic[draw,"$\theta_1$",angle radius=8mm,angle eccentricity=1.3]
    {angle=SheetL--P--SheetH};
  \draw pic[draw,"$\theta_2$",angle radius=11mm,angle eccentricity=1.25]
    {angle=SheetH--P--SheetR};
  \end{tikzpicture}\end{center}
\figcap{沿纸面法线流动的均匀面电流，横向边缘由两端点角给定。\label{ch23:fig:em23-question-2}}
\end{origexamplebox}

\begin{xesolution}
把导体面切成宽为 $\dd x$ 的无限长直导线，每条电流为 $\dd I=j\,\dd x$。取导体面沿 $x$ 轴，
$A=(0,a)$，电流指向纸内。由无限长直导线的磁场可写成
\[
  \dd\vect B=\frac{\mu_0j}{2\pi}
  \frac{a\,\vect e_x+x\,\vect e_y}{a^2+x^2}\,\dd x.
\]
若左、右宽度分别为 $a\tan\theta_1,a\tan\theta_2$，积分得到
\begin{align}
  B_x&=\frac{\mu_0j}{2\pi}(\theta_1+\theta_2),\\
  B_y&=\frac{\mu_0j}{2\pi}
  \ln\!\left(\frac{\cos\theta_1}{\cos\theta_2}\right).
\end{align}
非对称有限面会留下法向分量；若场点在有限面的中垂面上，$\theta_1=\theta_2$，则 $B_y=0$。
取以垂足为中心、左右等宽同时扩展的无限面极限，先由对称性得 $B_y=0$，再令 $\theta_1=\theta_2\to\pi/2$，故
\begin{equation}
  \boxed{\vect B=\frac{\mu_0j}{2}\vect e_x}.
\end{equation}
取一条跨过电流面的矩形安培环路，上、下两边各贡献 $B\ell$，包围电流为 $j\ell$，于是
$2B\ell=\mu_0j\ell$，与积分结果相同。

图~\ref{ch23:fig:review-currentsheetdiagram-1}中，左侧宽于右侧，因此积分的法向分量向下；水平分量由各窄条同向叠加。
\par\noindent\begin{minipage}{\linewidth}
\currentSheetDiagram
\figcap{有限宽电流带的分量积分\label{ch23:fig:review-currentsheetdiagram-1}}
\end{minipage}\par\medskip
\end{xesolution}



\begin{origexamplebox}{【电流环】}\label{ch23:ex:s39-8-3}
半径为 $R$ 的圆形电流环中通有电流 $I$。利用毕奥--萨伐尔定律求轴线上距圆心 $z$ 处的
磁感应强度。

% Pure source-question geometry: no derived forces, currents or solution curves.
\begin{center}\begin{tikzpicture}[scale=.8,>=Latex]
 \draw[axis](-1,0)--(6,0)node[right]{$z$};
 \draw[DeepBlue,thick](0,0)ellipse(.6 and 1.6);
 \draw[vector](0,-1.6)arc[start angle=-90,end angle=-15,x radius=.6,y radius=1.6]node[right]{$I$};
 \fill(4,0)circle(2pt)node[above]{$P$};\node[below left]at(0,0){$O$};
 \draw[<->](-1,0)--(-1,1.6)node[midway,left]{$R$};
 \draw[<->](0,-2.1)--(4,-2.1)node[midway,below]{$z$};
\end{tikzpicture}\end{center}
\figcap{圆电流与轴上观察点。\label{ch23:fig:em23-question-3}}
\end{origexamplebox}

\begin{xesolution}
轴上场点到每个电流元的距离都为 $\rho=\sqrt{R^2+z^2}$。由于
$\dd\vect\ell\perp\vect\rho$，每个电流元产生的场大小为
\[
  \dd B=\frac{\mu_0I}{4\pi}\frac{\dd\ell}{R^2+z^2}.
\]
把 $\dd\vect B$ 分解为轴向与横向分量。横向分量与圆环对径位置的贡献抵消；轴向投影因子为
$R/\rho$，故
\[
  B_z=\frac{\mu_0I}{4\pi}
  \frac{R}{(R^2+z^2)^{3/2}}\oint\dd\ell.
\]
代入 $\oint\dd\ell=2\pi R$，得到
\begin{equation}
  \boxed{B_z(z)=\frac{\mu_0IR^2}{2(R^2+z^2)^{3/2}}}.
\end{equation}
圆心处 $B(0)=\mu_0I/(2R)$；方向由电流的右手螺旋法则决定。

图~\ref{ch23:fig:review-loopaxisdiagram-2}中，Q与对径点 $Q\prime$ 同在圆环上，到P距离相同；两个场的横向分量相消，轴向分量相加。
\par\noindent\begin{minipage}{\linewidth}
\loopAxisDiagram
\figcap{圆电流对径线元的轴向叠加\label{ch23:fig:review-loopaxisdiagram-2}}
\end{minipage}\par\medskip
\end{xesolution}



\section{对称性、安培环路与磁通}

\subsection{安培环路定律}


对任意闭合曲线 $C$ 规定绕行正向，再用右手螺旋法则确定其所围表面的正法向。稳恒电流产生的
磁场满足
\begin{equation}
  \boxed{\oint_C\vect B\cdot\dd\vect\ell=\mu_0 I_{\rm enc}},
  \qquad
  \mu_0\simeq4\pi\times10^{-7}\,\mathrm{H/m}.
\end{equation}
$I_{\rm enc}$ 是穿过该表面的电流代数和：由负面穿向正面取正，反向取负。环路外的电流虽然会
改变曲线上各点的 $\vect B$，却不改变整个环量。

环路 $C$ 不一定是真实导线，它是为了积分而选取的数学闭合曲线；它所围的曲面也可有许多种形状。
对稳恒电流而言，电流处处连续，同一边界的不同跨越面所截获的净电流相同，因此右端定义明确。选环路
的原则不是“让它围住尽量多的电流”，而是利用对称性使 $\vect B\cdot\dd\vect\ell$ 在各段容易积分。
若对称性不足，环路定律仍给出正确的总环量，却不能单独决定曲线上每一点的 $B$；这时必须改用
毕奥--萨伐尔积分或结合其他条件。

闭合曲线的正向可以任意选定，但一旦选定，所跨曲面的正法向必须用右手规则与之绑定：右手四指沿环路
正向弯曲，拇指即为正法向。从曲面负面穿向正面的电流记正，反向记负。磁场环量
$\oint\vect B\cdot\dd\vect\ell$ 可与力沿闭合路径所做的功类比理解：它把每一点磁场在切线方向的
分量累加起来。环路外电流对局部磁场并非没有贡献，只是对整个闭合线积分的净贡献为零；这与闭合
高斯面外的电荷会改变面上电场、却不改变总电通量的结构相似。

也要看清环量定律与通量定律的差别。电场高斯定律在闭合曲面上积分法向分量，反映电荷造成的源汇；
稳恒磁场的安培定律在闭合曲线上积分切向分量，反映电流造成的环流。一个是“穿出多少”，一个是
“绕行多少”。因此不能把 $\oint\vect B\cdot\dd\vect\ell$ 误写成磁通，也不能用包围电流除以面积
直接得到 $B$。只有当所选环路上的磁场方向和大小由对称性确定后，环量才可化成 $B$ 乘有效长度。

图~\ref{ch23:fig:review-amperelawdiagram-1}中，环路正向决定纸外电流为正；外部电流虽改变局部场，却不计入所穿净电流。
\par\noindent\begin{minipage}{\linewidth}
\ampereLawDiagram
\figcap{有向安培环路及穿面电流\label{ch23:fig:review-amperelawdiagram-1}}
\end{minipage}\par\medskip

环路定律最适合高对称电流分布。无限长直导线周围的磁场沿同心圆切向，故
\begin{equation}
  B(2\pi r)=\mu_0I,
  \qquad
  \boxed{B=\frac{\mu_0I}{2\pi r}}.
\end{equation}
对单位长度有 $n$ 匝的理想无限长密绕螺线管，取一段长为 $\ell$ 的矩形环路，得到
\begin{equation}
  B\ell=\mu_0(n\ell)I,
  \qquad
  \boxed{B_{\rm in}=\mu_0nI},
  \qquad B_{\rm out}=0.
\end{equation}
有限螺线管外部并非严格无场，端部磁感线会弯回；长度远大于半径时，中部才可用上述理想结果。

无限长直电流的圆柱对称性保证：同一半径圆周上 $B$ 大小相同且处处沿切线，因而环路积分才可化为
$B(2\pi r)$。如果没有这种对称性，环路定律仍然成立，却不能仅靠它求出局部磁场。

方向也可由对称性与右手规则共同确定。把右手拇指指向正电流方向，弯曲的四指给出磁场环绕方向；在
垂直导线的平面内，磁感线是同心圆。场强随 $1/r$ 衰减而不是 $1/r^2$，原因是环量
$2\pi rB=\mu_0I$ 固定，而环路周长正比于 $r$。不能用圆柱侧面积解释这一规律：磁场沿侧面切向，
穿过柱侧面的磁通为零。无限长线电荷的电场虽然也按 $1/r$ 衰减，却来自高斯通量
$E(2\pi r\ell)=\lambda\ell/\varepsilon_0$；一个用切向环量，一个用法向通量，理由不同。

密绕螺线管的磁感线分布与一根均匀磁化的条形磁铁十分相似：内部磁感线近似平行而密集，外部从一端
绕回另一端。管越长、线圈越密，中部越接近匀强磁场。理想无限长模型用矩形环路跨过线圈侧面，外部
一段的场取零，内部一段给出 $B\ell$；跨越面穿过 $n\ell$ 匝、每匝电流为 $I$，于是得到
$B=\mu_0nI$。有限线圈端部的漏场不是“误差画出来”，而是磁感线闭合所必需的部分。

螺线管外场取零是无限长极限的结论。严格说来，单独由安培环路定律首先得到的是内外场之差
\begin{equation}
  B_{\rm in}-B_{\rm out}=\mu_0nI.
\end{equation}
再利用无限长模型的对称性以及远离电流面时不能凭空保留额外的匀强场，才得到 $B_{\rm out}=0$。
在有限螺线管中，外部磁感线从一端绕到另一端，内部场也随轴向位置变化；只有离端部足够远的中心区域，
才能近似使用 $\mu_0nI$。把有限模型直接当作处处匀强，会在端部受力和磁通计算中造成明显错误。

图~\ref{ch23:fig:review-wiresolenoiddiagram-1}中，右图矩形环路的下边在管内、上边在管外，内部场线没有越出管体；左图虚线表示圆周后半段。
\par\noindent\begin{minipage}{\linewidth}
\wireSolenoidDiagram
\figcap{直电流与理想螺线管的磁场\label{ch23:fig:review-wiresolenoiddiagram-1}}
\end{minipage}\par\medskip

图~\ref{ch23:fig:finite-coil-field}中，线圈内部向上的磁场经两端接到外部向下的回流场；有限长线圈外部不为零。
\par\noindent\begin{minipage}{\linewidth}
\finiteCoilFieldDiagram
\figcap{有限长螺线管的内部场与外部回流场\label{ch23:fig:finite-coil-field}}
\end{minipage}\par\medskip
\handSolenoidContext
右手图中的四指弯曲方向可代表圆电流方向，拇指即为磁矩和轴上磁场方向；若把拇指改指直导线电流，
四指则给出磁场环绕方向。旁边实物是空心密绕线圈，两个接线端对应同一根连续绕线的两端，
不能把绕组误当成许多个相互独立的闭合圆环。


% EM teaching enrichment 20260927
\paragraph{有限半径导线：包围的电流不是处处相同}
半径 $a$ 的无限长直导线中电流密度均匀，总电流 $I$。在 $r<a$ 处，环路只包围
$I_{\mathrm{enc}}=I r^2/a^2$，故
\[
 B(r)=\begin{cases}\mu_0Ir/(2\pi a^2),&r<a,\\
 \mu_0I/(2\pi r),&r\ge a.
 \end{cases}
\]
内部场随 $r$ 线性增加，外部按 $1/r$ 下降，表面连续。若电流只在表面流动，内部包围电流为零，答案就不同。
这种分段计算与均匀带电实心球的高斯法相似，但电场计算包围电荷并除以面积，磁场计算穿面电流并除以周长，不能交换两套积分量。


\subsection{磁通量与磁场高斯定律}


穿过曲面 $S$ 的磁通量为
\begin{equation}
  \Phi_B=\int_S\vect B\cdot\dd\vect S.
\end{equation}
经典电磁学中没有磁荷源汇，磁场满足 $\nabla\cdot\vect B=0$；磁感线可闭合或延伸至无穷远。因此任意闭合曲面满足
\begin{equation}
  \boxed{\oint_S\vect B\cdot\dd\vect S=0}.
\end{equation}
这条定律并不说闭合面上处处没有磁场，而是说所有穿出通量与穿入通量严格抵消。

磁通量的单位是 weber（Wb），$1\,\mathrm{Wb}=1\,\mathrm{T\,m^2}$。对平面匀强磁场，若面积法向与
磁场夹角为 $\theta$，有 $\Phi_B=BS\cos\theta$；对弯曲面或非匀强场则必须逐点使用
$\dd\Phi_B=\vect B\cdot\dd\vect S$。开曲面的磁通一般不为零；改变边界、位置或方向可以改变磁通，
但固定同一边界时，仅改变跨接曲面的形状不会改变磁通。设 $S_1,S_2$ 按同一边界正向定向，
$S_1$ 与反向的 $S_2$ 拼成闭合面，则 $\Phi_1-\Phi_2=0$。图中小面元的箭头是法向，不是曲面内切向，
其贡献为 $B\dd S\cos\theta$。即使闭合面包住一块条形磁铁，北极附近穿出的磁通也必在南极
附近等量穿入，因为磁感线在磁体内部继续闭合。

一根磁感线若穿入闭合面，最终还必须从别处穿出，所以净磁通为零。这个结论允许闭合面上有很强、很不
均匀的磁场，也允许面内存在电流和磁性物质；它排除的是“磁感线在面内凭空起始或终止”。因此磁场
高斯定律与安培环路定律描述的是两个不同的几何方面：前者约束磁场的源汇结构，后者把磁场的
环流与穿面电流联系起来。

如果将来发现真正的磁单极子，磁场高斯定律的右端就需要像电场高斯定律那样出现磁荷项；但在
目前经典电磁学与实验事实的范围内，右端取零。这里的“未发现”是物理事实陈述，而不是单靠数学逻辑
证明磁单极子永远不可能存在。竞赛计算采用现行麦克斯韦理论，因此应把
$\nabla\cdot\vect B=0$ 或闭合面总磁通为零作为基本约束。

这条约束不仅适用于稳恒电流，也适用于随时间变化的电磁场。电流如何变化会改变磁场的大小和分布，
却不会使磁感线突然出现孤立起点或终点；因此任何时刻对任意闭合面的磁通积分仍然为零。

图~\ref{ch23:fig:review-magneticgaussdiagram-1}中，内部场从S指向N，外部场回到S；闭合面有穿入也有穿出，不是面上处处无场。
\par\noindent\begin{minipage}{\linewidth}
\magneticGaussDiagram
\figcap{磁体内部与外部的磁感线\label{ch23:fig:review-magneticgaussdiagram-1}}
\end{minipage}\par\medskip
图~\ref{ch23:fig:review-fluxpatchdiagram-1}中，小面元面积为 $\dd S$，向量 $\dd\vect S$ 沿局部法向，与B的夹角进入 $B\dd S\cos\theta$。
\par\noindent\begin{minipage}{\linewidth}
\fluxPatchDiagram
\figcap{面积微元与局部法向\label{ch23:fig:review-fluxpatchdiagram-1}}
\end{minipage}\par\medskip
图~\ref{ch23:fig:review-chargemonopolecomparison-1}中，上排给出电荷场，下排只是假设磁荷时的类比；两种源汇结构的差别对应两条高斯定律的右端。
\par\noindent\begin{minipage}{\linewidth}
\chargeMonopoleComparison
\figcap{电荷场与假想磁荷场的比较\label{ch23:fig:review-chargemonopolecomparison-1}}
\end{minipage}\par\medskip
比较图上排是电荷产生的场，下排是假设存在孤立磁荷时的对应图像，并非已经观测到磁单极子。
普通磁铁的N、S极不能拆成两个孤立源；切开磁铁仍得到同时带有两极的磁体。



% EM teaching enrichment 20260927
\paragraph{零净磁通与非零单面磁通}
匀强场穿过边长 $a$ 的立方体，若场沿一条棱方向，两个相对面的磁通分别为 $+Ba^2$ 与 $-Ba^2$，其余四面为零，总和为零。把立方体旋转后，各面磁通改变而总和仍为零。
对开口线圈选取跨越面计算的磁通通常不为零；“磁场无源”限制的是闭合曲面，并非任意曲面。
若两张曲面共享同一个边界并取相容法向，它们拼合成闭面，于是两张面的磁通相等。下一章使用回路磁通时，可以自由选择便于积分的跨越面，依据正是这条规律。


% EM structure and teaching revision 20260928
\subsection{为何有环路定律仍要判断对称性}
静磁场满足 $\oint\boldsymbol B\cdot\dd\boldsymbol l=\mu_0I_{\mathrm{enc}}$。只有在所选路径上磁场方向和大小受到足够对称性约束时，才能从一个积分方程直接求出局部场。
无限长圆柱导体半径为 $R$、电流密度均匀，则圆形安培回路给出
\[
 B_\phi(r)=\begin{cases}\mu_0Ir/(2\pi R^2),&r<R,\\
 \mu_0I/(2\pi r),&r>R.
 \end{cases}
\]
导体内部包围的是部分电流，故磁场随 $r$ 增大；外部才是反比律。若同一圆柱中的电流密度不均匀且不具有轴对称性，环路定律仍成立，但不能把 $B$ 直接提到积分号外。
磁场的散度为零则约束另一个积分：任意闭合面的净磁通量为零。它不要求任意开曲面的磁通量为零，也不禁止磁力线穿过该曲面。计算线圈磁通时必须指定开曲面、面积正向和匝数；线圈的磁链是各匝磁通之和，不一定等于一匝磁通乘统一匝数。


\subsection{专题例题与方法应用}

\begin{origexamplebox}{【密绕螺线管】}\label{ch23:ex:s39-8-4}
“密绕”螺线管可理想化为一系列相互平行的圆电流环连续堆叠到无穷远。设半径为 $R$，单位长度
有 $n$ 匝，每匝电流为 $I$。

(1) 求轴线上的磁感应强度；

(2) 求任意位置的磁感应强度；

(3) 定性比较“只旋不进”的理想密绕螺线管与真实“既旋又进”的螺旋导线。

% Pure source-question geometry: no derived forces, currents or solution curves.
%
\begin{center}
\begin{tikzpicture}[scale=0.85,>=Latex]
  \begin{scope}[xshift=-4.1cm]
    \foreach \y in {-2,-1.2,...,2}{\draw[DeepBlue,thick] (0,\y) ellipse (1.45 and 0.34);}
    \draw[DeepBlue,thick,->] ({1.45*cos(230)},{-2+.34*sin(230)}) arc(230:310:1.45 and .34);
    \node[below] at (0,-2.34) {$I$};
  \end{scope}
  \begin{scope}[xshift=4.1cm]
    \foreach \t in {-720,-360,0,360}{
      \draw[DeepBlue,dashed,thick,domain=\t:{\t+180},samples=61]
        plot ({1.45*cos(\x)},{0.0027*\x+.34*sin(\x)});
      \draw[DeepBlue,very thick,domain={\t+180}:{\t+360},samples=61,
        postaction={decorate},decoration={markings,mark=at position .55 with {\arrow{Latex}}}]
        plot ({1.45*cos(\x)},{0.0027*\x+.34*sin(\x)});}
    \draw[guide,dashed] (0,-2.2)--(0,2.2);
    \node[below] at (.2,-.7) {$I$};
  \end{scope}
\node[font=\small] at(-4.1,-3){理想圆匝排列};\node[font=\small] at(4.1,-3){连续螺旋绕线};
\end{tikzpicture}
\end{center}
\figcap{圆匝近似与连续螺旋绕组的区别。\label{ch23:fig:em23-question-4}}
\end{origexamplebox}

\begin{xesolution}
\textbf{(1) 轴线积分。} 位于 $z'$ 到 $z'+\dd z'$ 内的线圈数为 $n\dd z'$，它在轴上 $z$ 处
产生
\[
  \dd B_z=\frac{\mu_0nIR^2}{2[R^2+(z-z')^2]^{3/2}}\,\dd z'.
\]
令 $u=z-z'$，对 $z'\in(-\infty,\infty)$ 积分：
\[
  B_z=\frac{\mu_0nIR^2}{2}
  \int_{-\infty}^{\infty}\frac{\dd u}{(R^2+u^2)^{3/2}}
  =\boxed{\mu_0nI}.
\]

\textbf{(2) 全空间。} 无限长、连续密绕模型具有平移和轴对称性。取跨过圆柱面的矩形安培
环路，结合远处磁场趋零可得
\begin{equation}
  \boxed{\vect B=
  \begin{cases}
    \mu_0nI\,\vect e_z,&r<R,\\
    \vect0,&r>R.
  \end{cases}}
\end{equation}
理想电流面上磁场发生跃变；若取对称平均值，表面场为 $\mu_0nI/2$。

\textbf{(3) 真实螺旋的附加场。} “只旋不进”的模型只有方位面电流
$K_\varphi=nI$。真实单根螺旋还把电流沿轴向输运，连续化以后另有
\[
  K_z=\frac{I}{2\pi R},
\]
它会像轴向直电流那样产生方位磁场；外接回线和有限端部还会带来漏磁场。密绕条件
$nR\gg1$ 下，附加场与主轴向场之比约为 $1/(2\pi nR)$，理想模型才足够准确。

图~\ref{ch23:fig:review-idealhelixdiagram-1}中，左侧仅含方位电流，右侧绕线还沿轴推进；后者的轴向电流决定附加方位场。
\par\noindent\begin{minipage}{\linewidth}
\idealHelixDiagram
\figcap{圆环叠加近似与有轴向推进的绕线\label{ch23:fig:review-idealhelixdiagram-1}}
\end{minipage}\par\medskip
\end{xesolution}



\section{安培力与载流回路的合力}

\subsection{电流元受力与闭合回路}


设导线的载流子数密度为 $n$、电量为 $q$、平均漂移速度为 $\vect v_d$，导线截面积为 $A$。
电流满足 $I=nqAv_d$。长度为 $\dd\ell$ 的一小段导线中有 $nA\dd\ell$ 个载流子，把它们所受的
洛伦兹力相加，得到电流元所受安培力
\begin{equation}
  \boxed{\dd\vect F=I\,\dd\vect\ell\times\vect B}.
\end{equation}
这里 $\dd\vect\ell$ 沿约定的正电流方向。即使金属中的真实载流子是负电子，这个写法仍然成立：
电子电量和漂移方向同时反号，乘积给出的仍是常规电流方向。

这个推导还解释了为什么公式中出现的是电流 $I$ 而不是单个载流子的速度。长度元内载流子数与
$\dd\ell$ 成正比，单个载流子的磁力与 $v_d$ 成正比；二者相乘得到 $nqAv_d\dd\ell=I\dd\ell$。
若导线与磁场夹角为 $\alpha$，受力大小为
\begin{equation}
  \dd F=IB\dd\ell\sin\alpha.
\end{equation}
方向由 $\dd\vect\ell$ 叉乘 $\vect B$ 决定。公式给出的是磁场对整段导线中载流子及其约束作用的
宏观合力；金属晶格通过电磁约束承受载流子传来的横向力，所以实验中移动的是整根导线，而不是只有
电子在导线内部横向堆积。

图~\ref{ch23:fig:review-ampereforcediagram-1}中，正载流子的漂移方向向右，磁场向纸内，磁力向上；三者方向用于确定电流元的受力。
\par\noindent\begin{minipage}{\linewidth}
\ampereForceDiagram
\figcap{载流子与电流元受力\label{ch23:fig:review-ampereforcediagram-1}}
\end{minipage}\par\medskip

若长度为 $L$ 的直导线垂直于匀强磁场，合力大小为 $F=IBL$。更一般地，匀强磁场中从 $a$ 到
$b$ 的任意弯曲导线满足
\begin{equation}
  \vect F=I\int_a^b\dd\vect\ell\times\vect B
  =I\vect L_{ab}\times\vect B,
  \qquad \vect L_{ab}=\vect r_b-\vect r_a.
\end{equation}
因此，导线的弯曲细节不影响合力；完全置于匀强磁场中的闭合回路因为
$\oint\dd\vect\ell=\vect 0$，合力为零。不过力分布并不为零，回路仍可能受到拉伸、压缩和力矩。

上式之所以能把曲线积分化成端点位移，关键是 $\vect B$ 在积分范围内处处相同：
\begin{equation}
 I\int_a^b\dd\vect\ell\times\vect B
 =I\left(\int_a^b\dd\vect\ell\right)\times\vect B
 =I(\vect r_b-\vect r_a)\times\vect B.
\end{equation}
若磁场随位置改变，$\vect B(\vect r)$ 不能提出积分号，弯曲路径上各段受力的加权不同，合力便不再只由
端点决定。对于闭合线框，也只有在整条回路都处于同一匀强磁场时才能直接令端点位移为零；若只有一部分
伸入磁场，应该把“有场区内那一段”的起点和终点代入，不能因线框本身闭合就误判总力为零。

“合力为零”只对完全处于匀强磁场中的闭合回路成立。安培力密集地作用在各段导线上，各段力可以
彼此抵消合力，却仍可能使柔软线圈扩张、压缩或改变形状。若磁场不均匀，闭合回路一般还会受到平移合力。
把电流方向按右手螺旋法则对应到面积法向后，磁矩与局部磁场同向的回路趋向磁场较强处，反向的回路则
趋向磁场较弱处。这一性质是理解顺磁、抗磁以及不均匀磁场中磁性物质受力的起点。

对尺度很小、可视为点磁偶极子的线圈，不均匀场中的势能仍近似为
$U=-\vect\mu\cdot\vect B(\vect r)$，于是平移力可写成
\begin{equation}
  \vect F=-\nabla U=\nabla(\vect\mu\cdot\vect B).
\end{equation}
若 $\vect\mu$ 已被约束为与局部磁场同向，线圈趋向 $B$ 较大的区域；若反向，则趋向较小区域。
这里既有使磁矩转向的力矩，也有由场梯度造成的平移力。匀强场中梯度为零，所以闭合小线圈只有力矩、
没有净平移力；非匀强场中两种效应可以同时存在。

图~\ref{ch23:fig:review-bentwireresultdiagram-1}中，端点位移 $\vect L_{ab}$ 从a指向b；只有匀强场中才可把共同的磁场移到积分号外。
\par\noindent\begin{minipage}{\linewidth}
\bentWireResultDiagram
\figcap{弯曲导线合力的端点等效\label{ch23:fig:review-bentwireresultdiagram-1}}
\end{minipage}\par\medskip
图~\ref{ch23:fig:straight-conductor}把同一结果用于直导线：从a到b的电流向右、磁场向纸内，故合力向上。
\par\noindent\begin{minipage}{\linewidth}
\straightConductorDiagram
\figcap{垂直匀强磁场的直导线受力\label{ch23:fig:straight-conductor}}
\end{minipage}\par\medskip


% EM teaching enrichment 20260927
\paragraph{随堂推演：半个线框进入磁场}
边长 $a,b$ 的矩形线框沿水平方向跨在匀强磁场边界上，磁场垂直纸面，只有一条长为 $b$ 的竖边完全在有场区。两段水平边上的力方向相反、长度相同，合力抵消；竖边留下大小 $IBb$ 的水平力。
不能因线框闭合就判为零，也不能把总周长 $2(a+b)$ 代入 $IBl$。若线框整体完全进入同一匀强场，两竖边力才相互抵消。导线若运动并产生感应电流，电流方向还应由下一章的磁通变化确定。


% EM structure and teaching revision 20260928
\subsection{磁力做功与电源供能如何相容}
单个电荷的磁力垂直于它的总速度，不做功；载流导体却可以受到安培力并获得机械功。这两句不矛盾，因为导体平动、载流子的相对漂移以及维持电流的电源共同参与能量交换。
长为 $\ell$ 的导体棒在匀强磁场中垂直于棒和磁场运动，取电流方向使磁力推动棒，大小为 $F=BI\ell$。棒的运动产生反电动势 $\mathcal E_b=B\ell v$，因而
\[
 P_{\mathrm{mech}}=Fv=\mathcal E_b I.
\]
若电源电动势为 $\mathcal E_0$、总电阻为 $R$，稳定电流满足
$\mathcal E_0=IR+\mathcal E_b$，电源功率分为焦耳热与机械输出：
$\mathcal E_0I=I^2R+Fv$。
外力反过来拖动棒时，机械输入又可变为电能，进入电磁感应问题。判断功率方向应统一电流和速度的参考方向，不应只因出现磁力就断言“体系没有机械能变化”。


\subsection{专题例题与方法应用}

\begin{origexamplebox}{【电路安培力】}\label{ch23:ex:s39-8-5}
一个“田”字形电阻网络置于匀强磁场中。网络每一小段的电阻均为 $R$；电流 $I$ 从左下端流入，
从右上端流出。磁场位于网络法向与某一条边张成的平面内，并与法向成 $45^\circ$。求网络受到的
总安培力。

% Pure source-question geometry: no derived forces, currents or solution curves.
%
\begin{center}
\begin{tikzpicture}[scale=0.86,>=Latex]
  \foreach \x in {0,1.5,3.0}{\draw[Ink,thick] (\x,0)--(\x,3.0);}
  \foreach \y in {0,1.5,3.0}{\draw[Ink,thick] (0,\y)--(3.0,\y);}
  \fill[DeepBlue] (0,0) circle (2.7pt) node[below right=4pt] {入};
  \fill[DeepBlue] (3,3) circle (2.7pt) node[below right=4pt] {出};
  \draw[Ink,thick] (-1,-.8)--(0,0);
  \draw[vector] (-.85,-.68)--(-.3,-.24) node[above left] {$I$};
  \draw[Ink,thick] (3,3)--(4,3.8);
  \draw[vector] (3.25,3.2)--(3.8,3.64) node[above left] {$I$};
  \begin{scope}[xshift=6.2cm,yshift=0.6cm]
    \coordinate (GridO) at (0,0);
    \coordinate (GridN) at (0,2.7);
    \coordinate (GridB) at (1.65,1.65);
    \draw[axis] (0,0)--(2.7,0) node[right] {$\vect e_1$};
    \draw[axis] (0,0)--(0,2.7) node[above] {$\vect n$};
    \draw[vector] (0,0)--(1.65,1.65) node[right] {$\vect B$};
    \node[below=10pt] at (1,0) {$\vect e_1$--$\vect n$ 平面侧视};
    \draw pic[draw,"$45^\circ$",angle radius=8mm,angle eccentricity=1.35]
      {angle=GridB--GridO--GridN};
    \end{scope}\end{tikzpicture}\end{center}
\figcap{方格网络的通流端点，以及磁场和板面法线的夹角。\label{ch23:fig:em23-question-5}}
\end{origexamplebox}

\begin{xesolution}
不需要先求每条支路的电流。对每条有向边 $e$，记电流和位移为 $I_e,\vect\ell_e$，则
\[
  \vect F=\left(\sum_e I_e\vect\ell_e\right)\times\vect B.
\]
把 $\vect\ell_e=\vect r_{e,+}-\vect r_{e,-}$ 展开；任一内部结点前面的系数等于“流入电流减流出
电流”，由 KCL 为零，只剩入口、出口两端。因此
\begin{equation}
  \boxed{\vect F=I(\vect r_{\rm out}-\vect r_{\rm in})\times\vect B
  =I\vect L\times\vect B}.
\end{equation}
结果与内部电阻及分流细节无关。

若每个小正方形边长为 $a$，取 $\vect e_1,\vect e_2$ 沿两组边、
$\vect n=\vect e_1\times\vect e_2$，并令磁场在网络面内的投影沿 $\vect e_1$，则
\[
  \vect L=2a(\vect e_1+\vect e_2),
  \qquad
  \vect B=\frac{B}{\sqrt2}(\vect n+\vect e_1).
\]
从而
\begin{equation}
  \boxed{\vect F=\sqrt2IaB(\vect e_1-\vect e_2-\vect n)},
  \qquad
  \boxed{F=\sqrt6\,IaB}.
\end{equation}
若磁场投影沿另一组边，只需交换 $\vect e_1,\vect e_2$；力的大小不变。

图~\ref{ch23:fig:review-resistorgridforcediagram-1}中，入口和出口位移决定总合力，内部结点的流入流出项成对抵消。
\par\noindent\begin{minipage}{\linewidth}
\resistorGridForceDiagram
\figcap{电流网络合力的端点等效\label{ch23:fig:review-resistorgridforcediagram-1}}
\end{minipage}\par\medskip
\end{xesolution}



\begin{origexamplebox}{【牛顿第三定律】}\label{ch23:ex:s39-8-6}
在力学中，两个质点间的相互作用力通常被写成大小相等、方向相反且共线的作用力与反作用力。电磁相互
作用通过空间中的场传播，若只截取两段电流元而忽略维持电流的完整回路与场的动量，结论未必仍具有
相同的简单形式。下面先比较孤立写出的两个电流元，再回到两个完整闭合回路。

(1) 考虑两个电流元 $I_1\dd\vect\ell_1$ 和 $I_2\dd\vect\ell_2$，写出二者之间的受力并判断
是否等大反向；

(2) 再考虑两个闭合电流回路，证明两个回路受到的总相互作用力大小相等、方向相反。注意这并未
证明两力作用线重合。

% Pure source-question geometry: no derived forces, currents or solution curves.
%
\begin{center}\begin{tikzpicture}[scale=.9,>=Latex]
 \coordinate(A)at(0,0);\coordinate(B)at(4,2);
 \fill(A)circle(2pt)node[below]{$1$};\fill(B)circle(2pt)node[below right=4pt]{$2$};
 \draw[vector](A)--(0,1.5)node[above]{$I_1\dd\vect\ell_1$};
 \draw[vector](B)--(5.5,2)node[right]{$I_2\dd\vect\ell_2$};
 \draw[vector](B)--(A)node[midway,below right=4pt]{$\vect r=\vect r_1-\vect r_2$};
 \end{tikzpicture}\end{center}
\figcap{两电流元与从电流元2指向电流元1的位置差。\label{ch23:fig:em23-question-6}}
\end{origexamplebox}

\begin{xesolution}
\textbf{(1) 电流元。} 令 $\vect r=\vect r_1-\vect r_2$ 从电流元 2 指向电流元 1，并记
$C=\mu_0I_1I_2/(4\pi r^3)$。由毕奥--萨伐尔定律和安培力公式，
\begin{align}
  \dd\vect F_{1\leftarrow2}
  &=C\,\dd\vect\ell_1\times(\dd\vect\ell_2\times\vect r)\\
  &=C\left[\dd\vect\ell_2(\dd\vect\ell_1\cdot\vect r)
  -\vect r(\dd\vect\ell_1\cdot\dd\vect\ell_2)\right],\\
  \dd\vect F_{2\leftarrow1}
  &=-C\,\dd\vect\ell_2\times(\dd\vect\ell_1\times\vect r)\\
  &=C\left[-\dd\vect\ell_1(\dd\vect\ell_2\cdot\vect r)
  +\vect r(\dd\vect\ell_1\cdot\dd\vect\ell_2)\right].
\end{align}
两式一般并不互为相反数。这里没有违背力学定律，因为孤立稳恒电流元本来就不是可独立存在的
封闭物体；维持电流的导线、载流子约束和电磁场动量都不能删去。

图~\ref{ch23:fig:review-currentelementsdiagram-1}中，位置矢量从2指向1，两力分别由各自的双重叉积计算，均垂直于本地电流元。
\par\noindent\begin{minipage}{\linewidth}
\currentElementsDiagram
\figcap{两电流元的相互作用力\label{ch23:fig:review-currentelementsdiagram-1}}
\end{minipage}\par\medskip

\textbf{(2) 闭合回路。} 两回路的互感能量可写成
\begin{equation}
  U_{12}=\frac{\mu_0I_1I_2}{4\pi}
  \oint_{C_1}\!\oint_{C_2}
  \frac{\dd\vect\ell_1\cdot\dd\vect\ell_2}{|\vect r_1-\vect r_2|}.
\end{equation}
它只依赖两回路的相对位置。若分别把两回路整体平移 $\vect a_1,\vect a_2$，共同平移不改变
$U_{12}$，所以
\[
  \nabla_{\!a_1}U_{12}+\nabla_{\!a_2}U_{12}=0.
\]
固定电流条件下，机械力为 $\vect F_i=+\nabla_{\!a_i}U_{12}$，这里是正号。
为看清这一点，令互感增加 $\dd M$：两电源维持电流所输入的能量合计为
$I_1\dd(MI_2)+I_2\dd(MI_1)=2I_1I_2\dd M$；磁场互能增加 $I_1I_2\dd M$，
剩余 $I_1I_2\dd M$ 转为机械功。因此不能把磁场互能直接当成无电源参与的机械势能并取负梯度。
两个平移梯度之和为零，故
\begin{equation}
  \boxed{\vect F_{1\leftarrow2}+\vect F_{2\leftarrow1}=\vect0}.
\end{equation}
这只证明总力等大反向；要讨论力的作用线与总力矩，还必须同时使用角动量守恒。
\end{xesolution}



\section{磁矩、进动与磁场应用}

\subsection{力矩、磁矩与能量}


面积为 $S$ 的平面电流环置于匀强磁场中。若 $\vect S$ 按右手螺旋法则定向，则两侧安培力
构成力偶，力矩为
\begin{equation}
  \vect\tau=I\vect S\times\vect B.
\end{equation}
定义回路的磁偶极矩
\begin{equation}
  \boxed{\vect\mu=I\vect S},
\end{equation}
便有
\begin{equation}
  \boxed{\vect\tau=\vect\mu\times\vect B},
  \qquad
  \boxed{U=-\vect\mu\cdot\vect B}.
\end{equation}
$\vect\mu$ 与 $\vect B$ 同向时势能最低。普通线圈受到指向稳定平衡位置的回复力矩；直流电动机利用电刷和换向器
在每半周改变线圈中的电流方向，使电磁力矩持续沿同一转动方向做功。

以矩形线框为例，取转轴沿线框一组边，且垂直于磁场。图中长为 $b$、平行转轴的两边A、C
所受力各为 $IbB$，等大反向但不共线，构成力偶；另一组边所受力沿转轴方向，共线相消。
若面积法向与磁场夹角为 $\theta$，两力作用线的距离为 $a\sin\theta$，故力矩大小为
$\tau=IbB\,a\sin\theta=IBS\sin\theta$，其中 $S=ab$。
推广到任意平面回路，只要用回路包围的有向面积 $\vect S$ 代替矩形面积，结果不变。
无阻尼时，线框一般会越过平衡位置并继续运动；只有存在耗散时才可能最终停在势能最低位置。电动机要实现连续转动，
必须在越过平衡位置后及时反转线圈电流；换向器所做的正是这一步，而电刷负责把外部直流电源持续接入
转动的线圈。

势能公式可以从力矩做功得到。设外界缓慢地把线框法向从 $\theta$ 转到
$\theta+\dd\theta$。以增大 $\theta$ 为正，磁力矩的带符号分量是
$\tau_\theta=-IBS\sin\theta$，不能与非负的力矩大小混用。因此
$\dd W_B=\tau_\theta\dd\theta=-IBS\sin\theta\dd\theta$，势能变化满足
\begin{equation}
  \dd U=-\dd W_B=IBS\sin\theta\,\dd\theta.
\end{equation}
积分得 $U=-IBS\cos\theta+C$。把 $\theta=\pi/2$ 选作零势能即有
$U=-\vect\mu\cdot\vect B$。在 $\theta=0$ 附近，微小偏转会产生使其返回的力矩，是稳定平衡；
$\theta=\pi$ 虽然力矩也为零，却处于势能最高点，是不稳定平衡。因而“力矩为零”本身不足以判定稳定性，
还必须比较势能或考察微小偏转后的力矩方向。

把完整回路作为基本对象还有一个理由：稳恒电流不可能在一小段导线的两端无端产生和消失，孤立电流元
并不存在。磁矩 $\vect\mu=I\vect S$ 把回路电流与几何面积结合为一个整体量，国际单位为
$\mathrm{A\,m^2}$。在远离回路、看不清其形状细节时，不同小回路的磁场都由这一磁矩主导，正如远处
不同小电荷分布可由低阶多极矩描述。

经典图像中，物质内部的原子电流与电子自旋都可以贡献微观磁矩。大量微观磁矩若方向杂乱，宏观磁性
可能相互抵消；在外场或内部相互作用下出现一定程度的定向排列，材料便表现出可测的磁化。这里不讨论
具体材料理论，但电流环的力矩、势能和远场已经提供了理解微观磁矩响应外场的基本语言。

若线圈有 $N$ 匝且各匝面积、方向相同，总磁矩为 $\vect\mu=NI\vect S$，力矩和势能也同时乘以 $N$。
对于不完全平面的闭合回路，仍可定义矢量面积
\begin{equation}
  \vect S=\frac12\oint\limits_C\vect r\times\dd\vect\ell,
\end{equation}
在匀强磁场中仍有 $\vect\tau=I\vect S\times\vect B$。这说明磁矩不是只对矩形或圆形线框才有意义；
它把任意小闭合回路对外磁场的最低阶响应统一成同一个矢量。

图~\ref{ch23:fig:review-looptorquediagram-1}中，A、C两侧导线平行转轴，其力不共线；右侧视图的距离 $a\sin\theta$ 直接给出力矩。
\par\noindent\begin{minipage}{\linewidth}
\loopTorqueDiagram
\figcap{矩形线框的力偶与力臂\label{ch23:fig:review-looptorquediagram-1}}
\end{minipage}\par\medskip

图~\ref{ch23:fig:review-motorstagesdiagram-1}中，阶段2到3之间电流换向，导线电流的点、叉和受力箭头随之同时改变。
\par\noindent\begin{minipage}{\linewidth}
\motorStagesDiagram
\figcap{直流电动机的四个转动阶段\label{ch23:fig:review-motorstagesdiagram-1}}
\end{minipage}\par\medskip
四阶段图是沿转轴看线框的截面，两端的圆点、叉号分别表示沿轴向纸外、纸内的电流，
上下箭头由 $I\dd\vect\ell\times\vect B$ 确定。阶段2到3之间，导线转过中性位置，
换向器使同一根导线中的电流反向，力矩继续沿图示逆时针方向。图心短线表示电刷与换向器接触位置，
它们不是直接短接线框两端的导线。恰在死点处力矩为零，连续转动还依靠惯性或多组线圈。


% EM teaching enrichment 20260927
\paragraph{小转角振动把力矩变成运动方程}
磁矩大小 $\mu$ 固定的刚性线圈可绕一轴自由转动，转动惯量为 $J$，转轴垂直于匀强磁场。以磁矩平行磁场的位置为 $\theta=0$，有
\[
 J\ddot\theta=-\mu B\sin\theta\simeq-\mu B\theta,
 \qquad \omega=\sqrt{\frac{\mu B}{J}}.
\]
也可由 $U=-\mu B\cos\theta\simeq-\mu B+\mu B\theta^2/2$ 认出回复系数。
该模型要求转动时 $\mu$ 近似不变，例如由电流源维持线圈电流。若是闭合低阻线圈且转动导致显著感应电流，必须同时列电路方程，不能把磁矩当常量。


% EM structure and teaching revision 20260928
\subsection{转动稳定性、固有磁矩与进动}
平面小线圈磁矩为 $\boldsymbol\mu=IS\boldsymbol n$，在匀强外场中净力为零而力矩为 $\boldsymbol M=\boldsymbol\mu\times\boldsymbol B$。若电流大小恒定，转轴固定且垂直于磁场、磁矩始终在垂直转轴的平面内，转角 $\theta$ 也就是磁矩与磁场的夹角，势能为 $U=-\mu B\cos\theta$。
平行取向附近 $U\simeq-\mu B+\mu B\theta^2/2$，转动惯量为 $J$ 时
\[
 J\ddot\theta=-\mu B\theta,\qquad \omega_0=\sqrt{\frac{\mu B}{J}}.
\]
反平行取向对应势能极大值，微小偏转会继续增大。如果电流并非由外源维持，而是在线圈运动中感生出来，则 $\mu$ 也会变化，不能直接照搬恒磁矩势能模型。
对于 $\boldsymbol\mu=\gamma\boldsymbol L$ 且角动量主要沿自转轴的快速转动体系，$\dd\boldsymbol L/\dd t=\gamma\boldsymbol L\times\boldsymbol B$ 给出进动。力矩垂直于角动量时主要改变方向，而非其大小；其近似前提是自转远快于进动，且外力矩不足以在一个自转周期内显著改变姿态。


\subsection{专题例题与方法应用}

\begin{origexamplebox}{【电流环力矩】}\label{ch23:ex:s39-8-7}
一个闭合电流环置于匀强磁场 $\vect B$ 中。

由于回路所受合力为零，关于不同参考点计算的总力矩应当相同。题目要求从各电流元的力矩积分出发，
验证这一结论，并把结果化为磁矩与磁场的叉积；不能只引用正文公式作为证明。

(1) 任取参考点 $O$，写出各电流元对 $O$ 的力矩积分；

(2) 用积分形式表示 $\vect\tau=I\vect S\times\vect B$；

(3) 证明两种表达式对任意形状的闭合电流环都相等。

% Pure source-question geometry: no derived forces, currents or solution curves.
\begin{center}\begin{tikzpicture}[scale=.78,>=Latex]
 \draw[thick] plot[smooth cycle]coordinates{(-2,0)(-1,1.4)(1.5,1.1)(2,-.5)(0,-1.2)};
 \draw[vector](1.5,1.1)--(1.85,.7)node[right]{$I$};
 \fill(-3,-1.6)circle(2pt)node[left]{$O$};
 \draw[guide](-3,-1.6)--(-1,1.4)node[midway,left]{$\boldsymbol r$};
 \draw[vector](3,-.7)--(3,1.3)node[above]{$\boldsymbol B$};
\end{tikzpicture}\end{center}
\figcap{闭合载流回路、任意参考点与外磁场。\label{ch23:fig:em23-question-7}}
\end{origexamplebox}

\begin{xesolution}
以 $O$ 为原点，电流元位置为 $\vect r$。直接由受力定义，
\begin{equation}
  \vect\tau_O=I\oint_C\vect r\times
  (\dd\vect\ell\times\vect B).
\end{equation}
任意闭合空间曲线的矢量面积定义为
\begin{equation}
  \vect S=\frac12\oint_C\vect r\times\dd\vect\ell.
\end{equation}
它对平面回路就是“面积乘法向”，故待证的第二种形式为
\[
  \vect\tau=I\vect S\times\vect B.
\]

利用三重积展开：
\[
  \vect r\times(\dd\vect\ell\times\vect B)
  =\dd\vect\ell(\vect r\cdot\vect B)
  -\vect B(\vect r\cdot\dd\vect\ell).
\]
第二项沿闭合回路的积分为
$\vect B\oint\dd(r^2/2)=0$。为求第一项，记
$\vect J=\oint\dd\vect\ell(\vect r\cdot\vect B)$。因 $\vect B$ 为常矢量，闭合回路上全微分积分为零：
\[
  \vect0=\oint\dd[\vect r(\vect r\cdot\vect B)]
  =\vect J+\oint\vect r(\dd\vect\ell\cdot\vect B).
\]
再将矢量面积与 $\vect B$ 叉乘，逐项展开：
\[
  \vect S\times\vect B
  =\frac12\oint[(\vect r\times\dd\vect\ell)\times\vect B]
  =\frac12\left[\vect J-\oint\vect r(\dd\vect\ell\cdot\vect B)\right]
  =\vect J.
\]
由此得到所需恒等式
\[
  \oint_C\dd\vect\ell(\vect r\cdot\vect B)=\vect S\times\vect B.
\]
因此
\begin{equation}
  \boxed{\vect\tau_O=I\vect S\times\vect B}.
\end{equation}
又因匀强磁场中闭合回路的合力
$I\oint\dd\vect\ell\times\vect B=0$，换参考点时增加的
$-\vect a\times\vect F$ 为零，故力矩确实与参考点无关。

图~\ref{ch23:fig:review-looptorquediagram-2}中，A、C两侧导线平行转轴，其力不共线；右侧视图的距离 $a\sin\theta$ 直接给出力矩。
\par\noindent\begin{minipage}{\linewidth}
\loopTorqueDiagram
\figcap{矩形线框的力偶与力臂\label{ch23:fig:review-looptorquediagram-2}}
\end{minipage}\par\medskip
\end{xesolution}



% FinalReviewQuestionSplit20261024 【电子显微镜】
\tcbset{breakable/.default=true}
\begin{origexamplebox}{【电子显微镜】}\label{ch23:ex:s39-8-8}
半径为 $R$ 的平面圆电流环中通有电流 $I$。质量为 $m$、电量为 $+q$ 的粒子从轴上距线圈
$u\gg R$ 的物点近轴入射，轴向速率记为 $v_0$，入射角满足 $|\theta|u\ll R$。
考虑非相对论弱聚焦情形 $|q|\mu_0|I|/(mv_0)\ll1$；当存在远离线圈的实像时，像距记为 $v\gg R$。

采用“脉冲近似”：圆环磁场随轴向距离快速衰减，粒子在远处几乎不受磁场作用；粒子的轴向速度较大，
通过线圈附近强场区的时间很短，相互作用期间可忽略横向位置的改变，只累计横向速度的微小变化。
在近轴条件下，不同初始小角度的粒子都应在同一像点附近会聚。

(1) 在薄透镜（脉冲）近似下证明
\[
  \frac1u+\frac1v=\frac1f,
\]
并求焦距；

(2) 求像绕轴线产生的角度偏转。

\begin{center}\begin{minipage}{.96\linewidth}\centering
\begin{tikzpicture}[>=Latex,font=\small,every node/.style={fill=none},scale=.88]

\draw[->] (-2,0)--(4.5,0) node[right]{$z$};
\draw[thick] (0,0) ellipse (.25 and 1.2);\node[above] at (0,1.2){圆线圈 $R,I$};
\fill (-1.8,0) circle (1.6pt) node[below]{物点};
\draw[DeepBlue,->] (-1.8,0)--(-.6,.25) node[midway,above]{$v_0,\theta$};
\draw[<->] (-1.8,-.7)--(0,-.7) node[midway,below]{$u$};

\end{tikzpicture}
\figcap{题图：物点、圆线圈和近轴初始射线；成像位置由磁场聚焦决定。\label{ch23:quality:q8}}
\end{minipage}\end{center}
\end{origexamplebox}
\tcbset{breakable/.default=false}

\par\pagebreak[0]
\begin{xesolution}
焦距必须依赖粒子的轴向动量 $p=mv_0$；若未指定束流速率，就应保留这一参数。
取轴场为 $B_z(z)$。轴对称条件下磁场无散给出
$r^{-1}\partial_r(rB_r)+\partial_z B_z=0$。
近轴时令 $B_r=c(z)r+O(r^3)$，代入得 $2c+B_z'=0$，于是
\begin{equation}
  B_r=-\frac r2\frac{\dd B_z}{\dd z}+O(r^3).
\end{equation}
以 $z$ 为自变量，撇号表示对 $z$ 求导，且 $\dot z\simeq v_0$。
由 $B_x=-xB_z'/2$、$B_y=-yB_z'/2$，洛伦兹方程的两个横向分量为
\[
 x''=\frac q{mv_0}\left(B_z y'+\frac{B_z'}2y\right),\qquad
 y''=-\frac q{mv_0}\left(B_z x'+\frac{B_z'}2x\right).
\]
令横向复坐标 $w=x+\mathrm i y$，将第二式乘 $\mathrm i$ 后与第一式相加，得到
\begin{equation}
  w''+2\mathrm i\kappa w'+\mathrm i\kappa'w=0,
  \qquad
  \kappa(z)=\frac{qB_z(z)}{2mv_0}.
\end{equation}
作拉莫尔旋转
\[
  w(z)=\xi(z)\exp\!\left[-\mathrm i\int^z\kappa(s)\dd s\right],
\]
令 $K'=\kappa$，则 $w'=\e^{-\mathrm iK}(\xi'-\mathrm i\kappa\xi)$，
$w''=\e^{-\mathrm iK}(\xi''-2\mathrm i\kappa\xi'-\mathrm i\kappa'\xi-\kappa^2\xi)$。
代回后虚数项逐一抵消，方程化成
\begin{equation}
  \xi''+\kappa^2(z)\xi=0.
\end{equation}
薄透镜区内横向位置几乎不变，斜率突变量为
\[
  \Delta\xi'=-\xi\int_{-\infty}^{\infty}\kappa^2(z)\dd z
  \equiv-\frac{\xi}{f}.
\]
对高度为 $\xi$ 的一条近轴射线，透镜前斜率为 $\xi/u$，后斜率为 $-\xi/v$；
故 $\Delta\xi'=-\xi(1/u+1/v)$。与上式对比，得到
\begin{equation}
  \boxed{\frac1u+\frac1v=\frac1f},
  \qquad
  \boxed{\frac1f=\frac{q^2}{4m^2v_0^2}
  \int_{-\infty}^{\infty}B_z^2(z)\dd z}.
\end{equation}

单圆环轴场为
\[
  B_z(z)=\frac{\mu_0IR^2}{2(R^2+z^2)^{3/2}},
\]
令 $z=R\tan t$，则下面积分化为
$R^{-5}\int_{-\pi/2}^{\pi/2}\cos^4t\,\dd t$；使用
$\cos^4t=(3+4\cos2t+\cos4t)/8$，得到
\[
  \int_{-\infty}^{\infty}\frac{\dd z}{(R^2+z^2)^3}
  =\frac{3\pi}{8R^5}.
\]
所以
\begin{equation}
  \boxed{\frac1f=\frac{3\pi q^2\mu_0^2I^2}{128m^2v_0^2R}},
  \qquad
  \boxed{f=\frac{128m^2v_0^2R}{3\pi q^2\mu_0^2I^2}}.
\end{equation}

拉莫尔变换本身还使整个像绕轴转过有符号角
\begin{equation}
  \phi_L=-\frac{q}{2mv_0}\int_{-\infty}^{\infty}B_z(z)\dd z.
\end{equation}
对单圆环 $\int B_z\dd z=\mu_0I$，故
\begin{equation}
  \boxed{\phi_L=-\frac{q\mu_0I}{2mv_0}}.
\end{equation}
符号由电荷和电流方向决定；若比较带方向标记的实物与实像，还要另计普通实像的倒转 $\pi$。

图~\ref{ch23:fig:review-magneticlensdiagram-1}中，高度为 $\xi$ 的射线前后斜率为 $\xi/u$ 与 $-\xi/v$；图在拉莫尔随动平面内表示聚焦。
\par\noindent\begin{minipage}{\linewidth}
\magneticLensDiagram
\figcap{近轴磁透镜的斜率变化\label{ch23:fig:review-magneticlensdiagram-1}}
\end{minipage}\par\medskip
\end{xesolution}



\begin{origexamplebox}{【自旋进动】}\label{ch23:ex:s39-8-9}
把质子的自旋角动量记作 $\vect s$，其磁矩写成
\[
  \vect\mu=\frac{ge}{m_p}\vect s,
  \qquad g=2.79.
\]
其中 $g=2.79$ 是由质子内部结构决定的无量纲因子。自旋是粒子的内禀量子属性，本题不讨论其量子起源，
只把 $\vect s$ 与 $\vect\mu$ 当作方向可能随外场变化、大小保持不变的经典矢量，研究磁力矩造成的
进动以及共振横向磁场造成的翻转。

(1) 证明质子在匀强磁场中发生进动，并求进动角速度；

(2) 当 $B_0=\SI{1}{T}$、$m_p=1.67\times10^{-27}\,\mathrm{kg}$、
$e=1.60\times10^{-19}\,\mathrm C$ 时，求进动频率；

(3) 求同频电磁波的真空波长并判断所属波段；

(4) 再施加一个垂直于 $\vect B_0$、并以同频率旋转的弱磁场，定性描述自旋的变化。

\begin{center}\begin{minipage}{.96\linewidth}\centering
\begin{tikzpicture}[>=Latex,font=\small,every node/.style={fill=none},scale=.88]

\draw[->,thick] (0,0)--(0,2) node[above]{$\boldsymbol B_0$};
\draw[DeepBlue,->,thick] (0,0)--(1.2,1.3) node[right]{$\boldsymbol s$};
\draw[dashed] (1.2,0)--(1.2,1.3);
\draw[->] (2.2,0)--(3.4,0) node[right]{$\boldsymbol B_1(t)\perp\boldsymbol B_0$};

\end{tikzpicture}
\figcap{题图：初始自旋与纵向恒定场，以及待施加的横向旋转弱场。\label{ch23:quality:q9}}
\end{minipage}\end{center}
\end{origexamplebox}

\begin{xesolution}
令回磁比
\[
  \gamma=\frac{ge}{m_p}.
\]
磁场对磁矩的力矩等于自旋角动量变化率：
\begin{equation}
  \frac{\dd\vect s}{\dd t}=\vect\mu\times\vect B_0
  =\gamma\vect s\times\vect B_0.
\end{equation}
右端同时垂直于 $\vect s$ 和 $\vect B_0$，故 $s$ 及其在磁场方向的投影都不变，只有方位角连续
变化。写成 $\dd\vect s/\dd t=\vect\omega_0\times\vect s$，得到
\begin{equation}
  \boxed{\vect\omega_0=-\gamma\vect B_0},
  \qquad
  \boxed{\omega_0=\frac{geB_0}{m_p}}.
\end{equation}
负号表示对正质子，进动角速度矢量与外磁场反向；频率大小与自旋倾角无关。

代入数值：
\begin{align}
  \omega_0&=2.79\times\frac{1.60\times10^{-19}}{1.67\times10^{-27}}
  =2.67\times10^8\ \mathrm{rad/s},\\
  f_0&=\frac{\omega_0}{2\pi}=\boxed{4.25\times10^7\ \mathrm{Hz}}
  =\boxed{42.5\ \mathrm{MHz}}.
\end{align}
同频电磁波在真空中的波长为
\begin{equation}
  \boxed{\lambda=\frac{c}{f_0}\simeq7.05\ \mathrm m},
\end{equation}
属于射电波（甚高频波段），不是微波、红外线或可见光。

横向旋转场的作用，用随动系中的图像最直观：

图~\ref{ch23:fig:review-spinprecessiondiagram-1}中，右图B1垂直纸面，转动发生在yz平面，整个过程中 $s_x=0$，不会转成平行B1。
\par\noindent\begin{minipage}{\linewidth}
\spinPrecessionDiagram
\figcap{进动与随动系中的自旋翻转\label{ch23:fig:review-spinprecessiondiagram-1}}
\end{minipage}\par\medskip

最后加横向旋转场 $\vect B_1$，令它与自旋进动同频、同向，即旋转角速度矢量为
$\vect\omega_0=-\gamma\vect B_0$。只有频率大小相同而旋向相反，并不满足这里的共振条件。在该转动参考系
中，$\vect B_1$ 变为静止，而 $\vect B_0$ 的进动被参考系转动抵消；于是自旋以
$\boxed{\Omega_R=\gamma B_1}$ 绕横向 $\vect B_1$ 作较慢的拉比转动。若初始
$\vect s\parallel\vect B_0$，持续时间
$\pi/(2\Omega_R)$ 的脉冲把自旋转到横向，持续 $\pi/\Omega_R$ 则把它翻到与 $\vect B_0$ 反向。

\end{xesolution}
}

\endgroup

\begingroup
\tcbset{breakable/.default=false,after={\par\nobreak\smallskip}}
\RenewDocumentEnvironment{xesolution}{+b}{%
  \par\nobreak\smallskip\nobreak\begingroup\small
  \clubpenalties 2 10000 10000\displaywidowpenalties 2 10000 10000\noindent
  {\bfseries\color{DeepBlue}解答}\quad #1\par\endgroup\smallskip
}{}
\chapter{电磁感应}

{
\graphicspath{{/Users/dongxingxi/Documents/物理竞赛一轮资料/sample/xes-original/fortieth-section/}{/Users/dongxingxi/Documents/物理竞赛一轮资料/sample/xes-original/fortieth-section/assets/}}
\section{磁通法、楞次定律与发电机}

% Diagrams for Lesson 9: electromagnetic induction.
\newcommand{\faradayHistoryContext}{%
  \historyfigure[0.93\linewidth]{faraday-ring-built-in-v1.png}%
  {左：法拉第肖像；右：铁环上的双线圈、检流计与电池构成早期电磁感应实验}}

\newcommand{\movingMagnetExperimentFigure}{%
\begin{center}
\resizebox{0.98\linewidth}{!}{%
\begin{tikzpicture}[x=1cm,y=1cm,>=Latex]
  \node[inner sep=0] at (0,0)
    {\includegraphics[width=14cm]{moving-magnet-experiment-built-in-v1.png}};
  \node[anchor=east,align=center] at (-7.35,2.35)
    {(a) 靠近\\$\vect v_M\leftarrow$};
  \node[anchor=east,align=center] at (-7.35,0)
    {(b) 静止\\$\vect v_M=\vect0$};
  \node[anchor=east,align=center] at (-7.35,-2.35)
    {(c) 远离\\$\vect v_M\rightarrow$};
\end{tikzpicture}}%
\figcap{磁铁靠近时检流计向一侧偏转；停住时回零；远离时反向偏转。}\label{ch24:fig:s40-9-2}
\end{center}}

\newcommand{\fluxOrientationDiagram}{%
\begin{center}
\begin{tikzpicture}[scale=0.88,>=Latex]
  \draw[curve] (0,0) ellipse (2.7 and 1.25);
  \draw[->,DeepBlue,very thick] (2.55,0.35) arc[start angle=16,end angle=105,x radius=2.7,y radius=1.25];
  \node[DeepBlue] at (2.5,2.0) {回路正向};
  \draw[vector] (0,0)--(0,2.15) node[above] {$\vect n$};
  \foreach \x in {-1.7,-0.8,0.15,1.1,1.9}{
    \draw[component] (\x,-1.65)--(\x,1.65);
  }
  \node[DeepBlue] at (2.35,0.0) {$\vect B$};
  \node[align=left,anchor=west] at (3.7,0.35)
    {$\displaystyle\Phi_B=\int_S\vect B\cdot\dd\vect S$\\[5pt]
     $\displaystyle\mathcal E=-\frac{\dd\Phi_B}{\dd t}$\\[5pt]
     负号表示感应效应阻碍磁通的改变};
\end{tikzpicture}
\end{center}}

\newcommand{\motionalRodDiagram}{%
\begin{center}
\begin{tikzpicture}[x=1cm,y=1cm,>=Latex]
  \draw[Ink,very thick] (-3,1.5)--(3,1.5);
  \draw[Ink,very thick] (-3,-1.5)--(3,-1.5);
  \draw[Ink,very thick] (-3,-1.5)--(-3,-0.55)
    decorate[decoration={zigzag,segment length=4pt,amplitude=2pt}] {-- (-3,0.55)}--(-3,1.5);
  \node[anchor=west] at (-2.78,0) {$R$};
  \draw[DeepBlue,very thick] (0,-1.65)--(0,1.65);
  \draw[vector] (0.2,0.35)--(1.75,0.35) node[right] {$\vect v$};
  \draw[vector] (-0.18,-0.25)--(-1.45,-0.25) node[left] {$\vect F_B$};
  \draw[vector] (0.18,-0.85)--(1.45,-0.85) node[right] {$\vect F_{\rm ext}$};
  \draw[<->] (-3.55,-1.5)--node[left] {$\ell$}(-3.55,1.5);
  \foreach \x in {-2.1,-1.3,0.7,1.5}{
    \foreach \y in {-1.2,1.15}{\node[TextGray] at (\x,\y) {$\times$};}
  }
  \node[DeepBlue] at (2.55,1.0) {$\vect B$};
  \draw[vector] (0,-0.82)--(0,0.82) node[left] {$I$};
  \draw[->,DeepBlue,thick] (-0.55,1.5)--(-1.65,1.5) node[midway,above] {$I$};
  \draw[->,DeepBlue,thick] (-1.65,-1.5)--(-0.55,-1.5);
  \node[align=left,anchor=west] at (3.4,0.3)
    {$\displaystyle\mathcal E_{\rm m}=\int(\vect v\times\vect B)\cdot\dd\vect\ell$\\[5pt]
     $\displaystyle\mathcal E_{\rm m}=B\ell v$ （三者两两垂直）\\[5pt]
     $\displaystyle I=\frac{B\ell v}{R},\quad F_B=BI\ell$\\[4pt]
     $q>0:\ q(\vect v\times\vect B)$ 沿棒向上};
\end{tikzpicture}
\end{center}}

\newcommand{\movingLoopDiagram}{%
\begin{center}
\begin{tikzpicture}[scale=0.9,>=Latex]
  \draw[curve,postaction={decorate},decoration={markings,mark=at position .38 with {\arrow{Latex}}}]
    plot[smooth cycle] coordinates{(-2,-0.7)(-1,1.2)(0.8,1.0)(2.0,0)(1.1,-1.1)(-0.8,-1.3)};
  \node[DeepBlue,anchor=south west] at (0.45,1.12) {$C$ 正向};
  \draw[vector] (-1.2,-0.5)--(-0.25,0.15) node[above] {$\vect v$};
  \foreach \p in {(-1.5,0.3),(-0.5,0.7),(0.45,0.35),(1.05,-0.45),(-0.2,-0.75)}
    {\node[TextGray] at \p {$\times$};}
  \node[DeepBlue] at (2.5,-0.4) {$\vect B$};
  \node[align=left,anchor=west] at (3.4,0.2)
   {$\displaystyle\mathcal E=\oint_C(\vect E+\vect v\times\vect B)\cdot\dd\vect\ell$\\[5pt]
    $\displaystyle\mathcal E=-\frac{\dd}{\dd t}\int_{S(t)}\vect B\cdot\dd\vect S$\\[5pt]
    图示特例：匀强 $\vect B$ 中刚性平移，$\mathcal E=0$};
\end{tikzpicture}
\end{center}}

\newcommand{\vortexFieldDiagram}{%
\begin{center}
\begin{tikzpicture}[scale=0.88,>=Latex]
  \draw[Ink,thick] (0,0) circle (2.0);
  \foreach \r in {0.65,1.25,2.65,3.35}{
    \draw[DeepBlue,thick,->] (\r,0) arc[start angle=0,end angle=180,radius=\r];
    \draw[DeepBlue,thick] (-\r,0) arc[start angle=180,end angle=360,radius=\r];
  }
  \foreach \x in {-1.1,0,1.1}{\foreach \y in {-1.0,0,1.0}{
    \pgfmathparse{\x*\x+\y*\y<3.5?1:0}\ifnum\pgfmathresult=1\node[TextGray] at (\x,\y) {$\times$};\fi}}
  \draw[<->] (0,-3.85)--node[below] {$R$}(2.0,-3.85);
  \node[align=left,anchor=west] at (4.1,0.25)
    {$\displaystyle E_\varphi(r)=-\frac r2\dot B,\quad r\le R$\\[6pt]
     $\displaystyle E_\varphi(r)=-\frac{R^2}{2r}\dot B,\quad r\ge R$\\[4pt]
     图示 $\dot B>0$ 向纸内；$\varphi$ 正向为顺时针\\感生电场沿逆时针};
\end{tikzpicture}
\end{center}}

\newcommand{\selfInductanceDiagram}{%
\begin{center}
\begin{tikzpicture}[scale=0.86,>=Latex]
  \foreach \x in {-2.4,-2.0,...,2.4}{\draw[Ink,thick] (\x,0) ellipse (0.22 and 1.45);}
  \draw[DeepBlue,->,very thick] (-3.0,0)--(3.0,0) node[right] {$\vect B$};
  \draw[Ink,very thick,->] ({2.4+.22*cos(310)},{1.45*sin(310)}) arc(310:380:.22 and 1.45);
  \node[right] at (2.7,-.6) {$I$};
  \draw[guide,dashed] (-2.6,-1.75) rectangle (2.6,1.75);
  \node[align=left,anchor=west] at (4.0,0.25)
    {$\displaystyle\Lambda=N\Phi=LI$\\[5pt]
     $\displaystyle\mathcal E_L=-\frac{\dd(LI)}{\dd t}$\\[5pt]
     $\displaystyle W_L=\frac12LI^2$};
\end{tikzpicture}
\end{center}}

\newcommand{\rlInductorContext}{%
\begin{center}
\resizebox{0.96\linewidth}{!}{%
\begin{tikzpicture}[x=1cm,y=1cm,>=Latex,line cap=round,line join=round]
  \node[inner sep=0] at (-3.75,0)
    {\includegraphics[width=5.4cm]{rl-inductor-built-in-v1.png}};
  \begin{scope}[shift={(0.1,0)}]
    \draw[Ink,thick] (0,1.55)--(1.35,1.55);
    \draw[Ink,thick,decorate,decoration={zigzag,segment length=5pt,amplitude=2.5pt}]
      (1.35,1.55)--(3.15,1.55);
    \draw[Ink,thick] (3.15,1.55)--(5.1,1.55)--(5.1,0.72);
    \foreach \yy in {0.54,0.18,-0.18,-0.54}{
      \draw[Ink,thick] (5.1,{\yy+0.18}) arc[start angle=90,end angle=-90,x radius=0.25,y radius=0.18];
    }
    \draw[Ink,thick] (5.1,-0.72)--(5.1,-1.55)--(3.20,-1.55);
    \fill[white,draw=Ink,thick] (3.20,-1.55) circle (2.2pt);
    \fill[white,draw=Ink,thick] (2.25,-1.55) circle (2.2pt);
    \draw[Ink,thick] (2.25,-1.55)--(3.03,-1.10);
    \draw[Ink,thick] (2.25,-1.55)--(0,-1.55)--(0,-0.30);
    \draw[Ink,thick] (0,0.30)--(0,1.55);
    \draw[Ink,very thick] (-0.38,0.30)--(0.38,0.30);
    \draw[Ink,very thick] (-0.23,-0.30)--(0.23,-0.30);
    \node[left] at (-0.42,0.30) {$+$};
    \node[left] at (-0.30,-0.30) {$-$};
    \node[left] at (-0.48,0) {$\mathcal E$};
    \node[above] at (2.25,1.72) {$R$};
    \node[right] at (5.38,0) {$L$};
    \node[above left=3pt] at (5.12,0.72) {$a$};
    \node[below left=3pt] at (5.12,-0.72) {$b$};
    \node[below] at (2.67,-1.78) {$S$};
    \draw[->,DeepBlue,very thick] (0.55,1.90)--(1.35,1.90) node[midway,above] {$I$};
  \end{scope}
\end{tikzpicture}}%
\figcap{左为环形电感实物结构；右为串联 $RL$ 回路。闭合 $S$ 后，电流沿箭头方向逐渐建立。}\label{ch24:fig:s40-9-9}
\end{center}}

\newcommand{\mutualInductanceDiagram}{%
\begin{center}
\begin{tikzpicture}[scale=0.9,>=Latex]
  \draw[curve] (-2.0,-0.25) ellipse (2.3 and 1.1);
  \draw[curve] (2.0,0.45) ellipse (2.0 and 0.95);
  \draw[->,DeepBlue,very thick] (-0.1,0.15) arc[start angle=10,end angle=70,x radius=2.3,y radius=1.1];
  \draw[->,DeepBlue,very thick] (3.75,0.7) arc[start angle=15,end angle=70,x radius=2.0,y radius=0.95];
  \node at (-2.0,-0.25) {$1$}; \node at (2.0,0.45) {$2$};
  \node[DeepBlue] at (-0.1,1.25) {$I_1$}; \node[DeepBlue] at (3.75,1.4) {$I_2$};
  \foreach \x in {-0.7,0.1,0.9}{\draw[component] (\x,-1.5) to[bend left=18] (\x+0.9,1.8);}
  \node[align=left,anchor=west] at (5.0,0.45)
   {$\Phi_2=M_{12}I_1,\quad\Phi_1=M_{21}I_2$\\[5pt]
    $\boxed{M_{12}=M_{21}}$\\[5pt]
    $W=\tfrac12L_1I_1^2+MI_1I_2+\tfrac12L_2I_2^2$};
\end{tikzpicture}
\end{center}}

\newcommand{\concentricLoopsDiagram}{%
\begin{center}\begin{tikzpicture}[scale=0.78,>=Latex]
  \draw[Ink,very thick] (0,0) circle (2.0); \draw[DeepBlue,very thick] (0,0) circle (1.1);
  \draw[->,Ink,thick] (1.85,0.75) arc[start angle=22,end angle=-45,radius=2.0];
  \draw[->,DeepBlue,thick] (0.95,0.55) arc[start angle=30,end angle=120,radius=1.1];
  \node[Ink,anchor=west] at (2.65,0.4) {原电流}; \node[DeepBlue] at (0,1.55) {感应电流};
\end{tikzpicture}\end{center}}

\newcommand{\magnetThreeLoopsDiagram}{%
\begin{center}\begin{tikzpicture}[scale=0.8,>=Latex]
  \draw[DeepBlue,dashed,very thick] (0,1.25) arc(90:270:.33 and 1.25);
  \draw[fill=SoftGray] (-1.8,-0.65) rectangle (1.8,0.65);
  \node at (-1.25,0) {$S$}; \node at (1.25,0) {$N$};
  \foreach \x in {-3.4,3.4}{\draw[DeepBlue,very thick] (\x,0) ellipse (0.33 and 1.25);}
  \draw[DeepBlue,very thick] (0,-1.25) arc(270:450:.33 and 1.25);
  \draw[vector] (-1.7,1.1)--(-2.7,1.1) node[midway,above=4pt] {$\vect v_M$};
  \draw[vector] (-3.4,-1.55)--(-4.45,-1.55) node[midway,below] {$\vect a_1$};
  \node[anchor=north] at (0,-1.43) {$t=0:\ \vect a_2=0$};
  \draw[vector,dashed] (0,-2.23)--(-1.15,-2.23)
    node[midway,below] {左移后 $\vect a_2$};
  \draw[vector] (3.4,-1.55)--(2.35,-1.55) node[midway,below] {$\vect a_3$};
  \node[TextGray,align=center] at (-3.4,2.1) {极外近轴\\收缩};
  \node[TextGray,align=center] at (0,3.5) {初始无趋势\\近中部左移后收缩};
  \node[TextGray,align=center] at (3.4,2.1) {极外近轴\\扩张};
\end{tikzpicture}\end{center}}

\newcommand{\rodResistorDiagram}{%
\begin{center}\begin{tikzpicture}[scale=0.78,>=Latex]
  \draw[Ink,very thick] (-3,1.4)--(3,1.4); \draw[Ink,very thick] (-3,-1.4)--(3,-1.4);
  \draw[Ink,very thick] (-2.7,-1.4)--(-2.7,-0.45) decorate[decoration={zigzag,segment length=4pt,amplitude=2pt}] {-- (-2.7,0.45)}--(-2.7,1.4);
  \node[left=8pt] at (-2.7,0) {$R$};
  \draw[DeepBlue,very thick] (0.8,-1.55)--(0.8,1.55); \node[above] at (0.8,1.55) {$m$};
  \draw[vector] (1.0,0)--(2.4,0) node[right] {$v$};
  \draw[<->] (-3.7,-1.4)--node[left] {$L$}(-3.7,1.4);
  \foreach \x in {-1.7,-0.7,0.3,1.4,2.3}{\foreach \y in {-0.75,0.75}{\node[TextGray] at (\x,\y) {$\times$};}}
  \node[align=left,anchor=west] at (3.7,0.2)
    {$\displaystyle I=\frac{BLv}{R}$\\[4pt]
     $\displaystyle m\dot v=-\frac{B^2L^2}{R}v$};
\end{tikzpicture}\end{center}}

\newcommand{\twoRodDiagram}{%
\begin{center}\begin{tikzpicture}[scale=0.78,>=Latex]
  \draw[Ink,very thick] (-3,1.4)--(3,1.4); \draw[Ink,very thick] (-3,-1.4)--(3,-1.4);
  \draw[DeepBlue,very thick] (-1.5,-1.55)--(-1.5,1.55); \draw[DeepBlue,very thick] (1.5,-1.55)--(1.5,1.55);
  \node[above] at (-1.5,1.62) {$m,R$}; \node[above] at (1.5,1.62) {$m,R$};
  \draw[vector] (1.7,0)--(2.9,0) node[right] {$v_0$};
  \foreach \x in {-2.4,-0.5,0.5,2.4}{\foreach \y in {-0.7,0.7}{\node[TextGray] at (\x,\y) {$\times$};}}
  \draw[<->] (-3.25,-1.4)--node[left] {$L$}(-3.25,1.4);
  \node[align=left,anchor=west] at (3.7,0.1)
   {$v_{\rm cm}=v_0/2$\\[4pt]$u=v_2-v_1=v_0\e^{-t/\tau}$};
\end{tikzpicture}\end{center}}

\newcommand{\unequalRodDiagram}{%
\begin{center}\begin{tikzpicture}[x=0.85cm,y=0.85cm,>=Latex]
  \draw[Ink,very thick] (-5,1.0)--(-1,1.0); \draw[Ink,very thick] (-5,-1.0)--(-1,-1.0);
  \draw[DeepBlue,very thick] (-4.0,-1.15)--(-4.0,1.15) node[above] {$m_1$};
  \draw[Ink,dotted,very thick] (-1,0)--(1,0);
  \draw[Ink,very thick] (1,1.8)--(5,1.8); \draw[Ink,very thick] (1,-1.8)--(5,-1.8);
  \draw[DeepBlue,very thick] (4.0,-1.95)--(4.0,1.95) node[above] {$m_2$};
  \draw[vector] (4.2,0)--(5.5,0) node[right] {$v_0$};
  \node at (-2.2,0) {$B_1,L_1$}; \node at (2.2,0) {$B_2,L_2$};
  \foreach \x in {-4.7,-3.2,-1.7,1.5,3.0,4.6}{\node[TextGray] at (\x,0.55) {$\times$};}
\end{tikzpicture}\end{center}}

\newcommand{\rodCapacitorDiagram}{%
\begin{center}\begin{tikzpicture}[scale=0.78,>=Latex]
  \draw[Ink,very thick] (-3,1.4)--(3,1.4); \draw[Ink,very thick] (-3,-1.4)--(3,-1.4);
  \draw[Ink,very thick] (-2.7,-1.4)--(-2.7,-0.25); \draw[Ink,very thick] (-3.05,-0.25)--(-2.35,-0.25);
  \draw[Ink,very thick] (-3.05,0.25)--(-2.35,0.25); \draw[Ink,very thick] (-2.7,0.25)--(-2.7,1.4);
  \node[left] at (-2.7,0) {$C$};
  \draw[DeepBlue,very thick] (0.8,-1.55)--(0.8,1.55) node[above] {$m$};
  \draw[vector] (1.0,0)--(2.4,0) node[right] {$v_0$};
  \foreach \x in {-1.7,-0.7,0.2,1.5,2.4}{\foreach \y in {-0.7,0.7}{\node[TextGray] at (\x,\y) {$\times$};}}
  \node[align=left,anchor=west] at (3.5,0.2)
    {$q_\infty=CBLv_\infty$\\[4pt]$m(v-v_0)=-BLq$};
\end{tikzpicture}\end{center}}

\newcommand{\rodInductorDiagram}{%
\begin{center}\begin{tikzpicture}[scale=0.78,>=Latex]
  \draw[Ink,very thick] (-3,1.4)--(3,1.4); \draw[Ink,very thick] (-3,-1.4)--(3,-1.4);
  \draw[Ink,very thick] (-2.7,-1.4)--(-2.7,-0.6);
  \draw[Ink,thick] (-2.7,-0.6) decorate[decoration={coil,aspect=0.65,segment length=5pt,amplitude=5pt}] {-- (-2.7,0.6)};
  \draw[Ink,very thick] (-2.7,0.6)--(-2.7,1.4); \node[left] at (-2.9,0) {$\mathcal L$};
  \draw[DeepBlue,very thick] (0.8,-1.55)--(0.8,1.55) node[above] {$m$};
  \draw[vector] (1.0,0)--(2.4,0) node[right] {$v_0$};
  \foreach \x in {-1.7,-0.7,0.2,1.5,2.4}{\foreach \y in {-0.7,0.7}{\node[TextGray] at (\x,\y) {$\times$};}}
  \node[align=left,anchor=west] at (3.5,0.2)
    {$\mathcal L\dot I=B\ell v$\\[4pt]$m\dot v=-B\ell I$};
\end{tikzpicture}\end{center}}

\newcommand{\rodSourceDiagram}{%
\begin{center}\begin{tikzpicture}[scale=0.78,>=Latex]
  \draw[Ink,very thick] (-3,1.4)--(3,1.4); \draw[Ink,very thick] (-3,-1.4)--(3,-1.4);
  \draw[Ink,very thick] (-2.7,-1.4)--(-2.7,-0.55);
  \draw[Ink,thick] (-2.7,-0.55) decorate[decoration={zigzag,segment length=4pt,amplitude=2pt}] {-- (-2.7,0.1)};
  \draw[Ink,very thick] (-3.05,0.25)--(-2.35,0.25); \draw[Ink,thick] (-2.9,0.55)--(-2.5,0.55);
  \draw[Ink,very thick] (-2.7,0.55)--(-2.7,1.4); \node[left] at (-2.85,-0.25) {$R$}; \node[left] at (-2.9,0.5) {$\mathcal E$};
  \draw[DeepBlue,very thick] (0.8,-1.55)--(0.8,1.55) node[above] {$m$};
  \draw[vector] (1.0,0)--(2.4,0) node[right] {$v$};
  \foreach \x in {-1.7,-0.7,0.2,1.5,2.4}{\foreach \y in {-0.7,0.7}{\node[TextGray] at (\x,\y) {$\times$};}}
\end{tikzpicture}\end{center}}

\newcommand{\equivalentCircuitDiagram}{%
\begin{center}\begin{tikzpicture}[scale=0.82,>=Latex]
  \draw[Ink,very thick] (-4,0)--(-2.7,0);
  \draw[Ink,thick] (-2.7,0) decorate[decoration={zigzag,segment length=5pt,amplitude=2.5pt}] {-- (-0.9,0)};
  \node[above] at (-1.8,0.25) {$R$};
  \draw[Ink,very thick] (-0.9,0)--(0.0,0);
  \draw[Ink,very thick] (0,-0.75)--(0,0.75); \draw[Ink,very thick] (0.45,-0.75)--(0.45,0.75);
  \node[above] at (0.22,0.9) {$C_m=m/(B^2L^2)$};
  \draw[Ink,very thick] (0.45,0)--(2.0,0);
  \node[align=left,anchor=west] at (3.0,0.2)
   {$U=RI+U_m$\\[4pt]$I=C_m\dot U_m$\\[4pt]$\tfrac12mv^2=\tfrac12C_mU_m^2$};
\end{tikzpicture}\end{center}}

\newcommand{\chargedRingDiagram}{%
\begin{center}\begin{tikzpicture}[scale=0.82,>=Latex]
  \draw[Ink,very thick] (0,0) circle (2.4);
  \fill[DeepBlue] (2.4,0) circle (3pt);
  \node[anchor=south west] at (2.5,0.12) {$m,+q$};
  \draw[->,DeepBlue,very thick] (2.0,1.15) arc[start angle=30,end angle=105,radius=2.4];
  \node[DeepBlue] at (2.8,1.5) {$\dot\theta$};
  \foreach \x in {-1.1,0,1.1}{\foreach \y in {-1.1,0,1.1}{\node[TextGray] at (\x,\y) {$\times$};}}
  \node at (0,0) {$B(t)$};
  \draw[vector] (2.32,-0.32)--(1.05,-0.32)
    node[midway,below] {$a_r=-R\dot\theta^2$};
  \draw[axis] (2.48,0.32)--(3.75,0.32)
    node[right,align=left] {径向正向\\$N_r,\ qR\dot\theta B$ 均取代数值};
\end{tikzpicture}\end{center}}

\newcommand{\chordParticleDiagram}{%
\begin{center}\begin{tikzpicture}[scale=0.8,>=Latex]
  \draw[Ink,thick] (0,0) circle (2.5);
  \draw[Ink,very thick] ({-sqrt(2.5*2.5-1.25*1.25)},-1.25)--({sqrt(2.5*2.5-1.25*1.25)},-1.25);
  \fill[DeepBlue] (0,-1.25) circle (3pt) node[below] {$m,+q$};
  \draw[<->] (0,0)--node[right] {$h$}(0,-1.25);
  \draw[axis] (-2.2,-1.75)--(2.4,-1.75) node[right] {$x$};
  \node[TextGray] at (-1.4,1.0) {$\times$}; \node[TextGray] at (0,1.0) {$\times$}; \node[TextGray] at (1.4,1.0) {$\times$};
  \node at (0,1.65) {$B_0\sin\omega t$};
  \node[align=left,anchor=west] at (3.5,0)
    {$\displaystyle m\ddot x=\frac{qh}{2}\dot B$\\[5pt]
     $\displaystyle x=\frac{qhB_0}{2m\omega}(1-\cos\omega t)$};
\end{tikzpicture}\end{center}}

\newcommand{\chordEmfDiagram}{%
\begin{center}\begin{tikzpicture}[scale=0.82,>=Latex]
  \draw[Ink,thick] (0,0) circle (2.5);
  \draw[Ink,very thick] (-3.2,-1.2)--(4.3,-1.2);
  \coordinate (A) at ({-sqrt(2.5*2.5-1.2*1.2)},-1.2); \coordinate (C) at ({sqrt(2.5*2.5-1.2*1.2)},-1.2); \coordinate (D) at (4.1,-1.2);
  \fill[DeepBlue] (A) circle (2.5pt) node[below] {$A$}; \fill[DeepBlue] (C) circle (2.5pt) node[below] {$C$}; \fill[Ink] (D) circle (2.5pt) node[below] {$D$};
  \draw[<->] ({sqrt(2.5*2.5-1.2*1.2)},-2)--(4.1,-2) node[midway,below] {$\ell$};
  \draw[<->] (0,0)--node[right] {$h$}(0,-1.2);
  \foreach \r in {0.8,1.6,3.2}{\draw[DeepBlue,thick,->] (\r,0) arc[start angle=0,end angle=105,radius=\r];}
  \node[TextGray] at (-1.1,.2) {$\dot B=k$};
\end{tikzpicture}\end{center}}

\newcommand{\circularVoltmetersDiagram}{%
\begin{center}\begin{tikzpicture}[scale=0.76,>=Latex]
  \draw[Ink,very thick] (0,0) circle (2.4);
  \coordinate (A) at (150:2.4); \coordinate (C) at (30:2.4);
  \fill[DeepBlue] (A) circle (2.5pt) node[left] {$A$}; \fill[DeepBlue] (C) circle (2.5pt) node[right] {$C$};
  \draw[Ink,thick] (A) to[out=115,in=65,looseness=1.05]
    node[midway,circle,draw,fill=white,inner sep=2pt] {$V_1$} (C);
  \draw[Ink,thick] (A) .. controls (-3.2,-2.8) and (3.2,-2.8) ..
    node[pos=0.5,circle,draw,fill=white,inner sep=2pt] {$V_2$} (C);
  \draw[Ink,thick] (A)--(-0.35,0.2);
  \draw[Ink,thick] (0.35,0.2)--(C);
  \node[circle,draw,fill=white,inner sep=2pt] at (0,0.2) {$V_3$};
  \foreach \x in {-1,0,1}{\node[TextGray] at (\x,-0.65) {$\times$};}
  \node at (0,-1.15) {$\dot B=k$};
\end{tikzpicture}\end{center}}

\newcommand{\squareVoltmetersDiagram}{%
\begin{center}\begin{tikzpicture}[scale=0.76,>=Latex]
  \draw[Ink,very thick] (-2,-2) rectangle (2,2);
  \draw[Ink,thick] (-2,2.55)--(2,2.55); \node[circle,draw,fill=white,inner sep=2pt] at (0,2.55) {$V_1$};
  \draw[Ink,thick] (2.55,-2)--(2.55,2); \node[circle,draw,fill=white,inner sep=2pt] at (2.55,0) {$V_2$};
  \draw[Ink,thick] (-2,-2.55)--(2,-2.55); \node[circle,draw,fill=white,inner sep=2pt] at (0,-2.55) {$V_3$};
  \draw[Ink,thick] (-2.55,-2)--(-2.55,2); \node[circle,draw,fill=white,inner sep=2pt] at (-2.55,0) {$V_4$};
  \node[left] at (-2,0.75) {$R_1$}; \node[right] at (2,0.75) {$R_2$};
  \node[TextGray] at (0,0) {$\times\ B(t)$}; \draw[<->] (-2,-2.85)--node[below] {$a$}(2,-2.85);
\end{tikzpicture}\end{center}}

\newcommand{\gridCurrentDiagram}{%
\begin{center}\begin{tikzpicture}[scale=0.88,>=Latex]
  \draw[Ink,very thick] (-2,-2) grid[step=2] (2,2);
  \draw[->,DeepBlue,very thick] (-2,2)--(0,2);
  \draw[->,DeepBlue,very thick] (0,2)--(2,2);
  \draw[->,DeepBlue,very thick] (2,2)--(2,0);
  \draw[->,DeepBlue,very thick] (2,0)--(2,-2);
  \draw[->,DeepBlue,very thick] (2,-2)--(0,-2);
  \draw[->,DeepBlue,very thick] (0,-2)--(-2,-2);
  \draw[->,DeepBlue,very thick] (-2,-2)--(-2,0);
  \draw[->,DeepBlue,very thick] (-2,0)--(-2,2);
  \node[DeepBlue] at (2.6,0) {$I_0$};
  \node[TextGray] at (0,0.35) {内部支路}; \node[TextGray] at (0,-0.35) {$I=0$};
  \draw[<->] (-2,-2.5)--node[below] {$2a$}(2,-2.5);
\end{tikzpicture}\end{center}}

\newcommand{\solenoidSplitDiagram}{%
\begin{center}\begin{tikzpicture}[scale=0.78,>=Latex]
  \foreach \x in {-4.0,-3.6,...,4.0}{\draw[Ink,thick] (\x,0) ellipse (0.2 and 1.15);}
  \draw[guide,dashed] (0,-1.55)--(0,1.55);
  \draw[<->] (-4.2,-1.65)--node[below] {$l_1$}(0,-1.65);
  \draw[<->] (0,-1.65)--node[below] {$l_2$}(4.2,-1.65);
  \draw[vector] (-3.7,0)--(3.7,0) node[right] {$B=\mu_0nI$};
  \node[align=left,anchor=west] at (5.0,0.15)
    {$L=\mu_0n^2A(l_1+l_2)$\\[4pt]$=L_1+L_2$ （理想长螺线管）};
\end{tikzpicture}\end{center}}

\newcommand{\mutualLoopsProofDiagram}{%
\begin{center}\begin{tikzpicture}[scale=0.82,>=Latex]
  \draw[curve] (-2.6,0) ellipse (1.35 and 0.9); \draw[curve] (2.6,0.4) ellipse (1.25 and 0.8);
  \draw[vector] (-1.5,0.35)--(1.9,0.7) node[midway,above] {$\vect r_{12}$};
  \draw[vector] (-3.2,-0.65)--(-2.2,-0.85) node[below] {$\dd\vect\ell_1$};
  \draw[vector] (2.0,-0.15)--(3.0,-0.25) node[below] {$\dd\vect\ell_2$};
  \node[align=left,anchor=west] at (4.6,0.25)
   {$\displaystyle M_{12}=\frac{\mu_0}{4\pi}\oint_{C_1}\!\oint_{C_2}
   \frac{\dd\vect\ell_1\cdot\dd\vect\ell_2}{r_{12}}$\\[6pt]
    交换 $1,2$ 后积分核不变};
\end{tikzpicture}\end{center}}

\newcommand{\rlcDiagram}[1]{%
\begin{center}\begin{tikzpicture}[scale=0.70,>=Latex]
  \draw[Ink,very thick] (-4,0.9)--(-3.2,0.9);
  \draw[Ink,very thick] (-3.2,0.3)--(-3.2,1.5); \draw[Ink,very thick] (-2.75,0.5)--(-2.75,1.3);
  \node[above] at (-3.0,1.5) {$C,\ U$};
  \draw[Ink,very thick] (-2.75,0.9)--(-1.7,0.9);
  \draw[Ink,thick] (-1.7,0.9) decorate[decoration={coil,aspect=0.65,segment length=5pt,amplitude=5pt}] {-- (0.3,0.9)};
  \node[above] at (-0.7,1.45) {$L$};
  \draw[Ink,thick] (0.3,0.9) decorate[decoration={zigzag,segment length=5pt,amplitude=2.5pt}] {-- (2.2,0.9)};
  \node[above] at (1.25,1.25) {$R$};
  \draw[Ink,very thick] (2.2,0.9)--(3.0,0.9)--(3.0,-0.9);
  \ifnum#1=1
    \draw[Ink,very thick] (3.0,-0.9)--(0.25,-0.9);
    \draw[Ink,very thick] (-0.2,-0.9)--(-4,-0.9)--(-4,0.9);
    \draw[Ink,very thick] (-0.2,-1.4)--(-0.2,-0.4); \draw[Ink,very thick] (0.25,-1.2)--(0.25,-0.6);
    \node[below] at (0,-1.35) {$\mathcal E$};
  \else
    \draw[Ink,very thick] (3.0,-0.9)--(-4,-0.9)--(-4,0.9);
  \fi
  \draw[->,DeepBlue,very thick] (-1.6,0.38)--(0.8,0.38) node[midway,above] {$I(t)$};
\end{tikzpicture}\end{center}}

\newcommand{\travelingFieldTrainDiagram}{%
\begin{center}\begin{tikzpicture}[scale=0.72,>=Latex]
  \draw[Ink,very thick] (-4,1.25)--(4,1.25); \draw[Ink,very thick] (-4,-1.25)--(4,-1.25);
  \draw[DeepBlue,very thick,fill=PaleBlue] (-1.4,-1.25) rectangle (1.4,1.25);
  \node at (0,0) {$m$}; \draw[<->] (-1.4,-1.62)--node[below] {$a$}(1.4,-1.62);
  \draw[<->] (-3.4,-1.25)--node[left] {$L$}(-3.4,1.25);
  \draw[vector] (1.55,0.45)--(3.0,0.45) node[right] {$v$};
  \foreach \x in {-3.2,-2.4,-0.8,0,0.8,2.4,3.2}{\node[TextGray] at (\x,0.75) {$\times$};}
  \draw[->,TextGray,thick] (-3.5,-0.65)--(3.5,-0.65)
    node[right] {$v_\phi$};
  \node[align=left,anchor=west] at (4.6,0.2)
    {$B=B_0\cos(kx-\omega t)$\\[4pt]$v_\phi=\omega/k$};
\end{tikzpicture}\end{center}}

% Faithful independent redraw of the generator diagram from Lesson 9.
\newcommand{\generatorFigureFaithful}{%
\begin{center}
\resizebox{0.94\linewidth}{!}{%
\begin{tikzpicture}[x=0.90cm,y=0.90cm,>=Latex,line cap=round,line join=round]
  % left panel: rotating rectangular coil between pole pieces, slip rings and lamp
  \begin{scope}[shift={(-4.4,0)}]
    \draw[fill=SoftGray,draw=Ink,thick] (-3.8,0.45)--(-1.65,1.55)--(-0.85,1.05)--(-2.95,-0.05)--cycle;
    \draw[fill=SoftGray,draw=Ink,thick] (0.85,1.05)--(3.05,2.15)--(3.80,1.62)--(1.63,0.52)--cycle;
    \node[font=\bfseries] at (-2.45,0.73) {$N$}; \node[font=\bfseries] at (2.35,1.35) {$S$};
    \draw[DeepBlue,line width=1.6pt] (-1.28,-0.72)--(-0.78,1.22)--(1.30,0.70)--(0.80,-1.24)--cycle;
    \draw[Ink,thick] (-0.25,-2.12)--(-0.25,1.00);
    \draw[Ink,thick] (-0.82,-1.75) ellipse (0.58 and 0.22);
    \draw[Ink,thick] (0.28,-1.92) ellipse (0.58 and 0.22);
    \draw[DeepBlue,thick,->] (-1.70,-1.95) arc[start angle=190,end angle=350,x radius=1.45,y radius=0.55];
    \node[below] at (-0.25,-2.43) {$\omega$};
    \draw[Ink,thick] (-1.36,-1.68)--(-2.30,-2.60)--(-2.30,-3.05);
    \draw[Ink,thick] (0.85,-1.85)--(2.25,-2.60)--(2.25,-3.05);
    \draw[Ink,thick] (-2.30,-3.05)--(-1.05,-3.05) (1.05,-3.05)--(2.25,-3.05);
    \draw[Ink,thick] (-1.05,-3.05) .. controls (-0.95,-2.15) and (0.95,-2.15) .. (1.05,-3.05);
    \draw[Ink,thick] (-0.72,-2.86) arc[start angle=200,end angle=-20,x radius=0.72,y radius=0.45];
    \node[font=\scriptsize] at (0,-3.55) {滑环、电刷与负载};
  \end{scope}
  % right panel: coil normal in uniform field
  \begin{scope}[shift={(5.0,0.25)}]
    \foreach \y in {-2.0,-1.45,...,2.0}{\draw[TextGray,thick,->] (-3.45,\y)--(3.35,\y);}
    \node[above,font=\bfseries] at (0,2.18) {$\vect B$};
    \begin{scope}[rotate=37]
      \foreach \s in {-0.18,-0.09,0,0.09,0.18}{\draw[Ink,thick] (-1.55+\s,-1.05)--(1.55+\s,-1.05)--(1.55+\s,1.05)--(-1.55+\s,1.05)--cycle;}
    \end{scope}
    \draw[->,DeepBlue,very thick] (0,0)--(2.05,1.58) node[right] {法线 $\vect n$};
    \draw[TextGray,dashed] (0,0)--(2.65,0);
    \draw[DeepBlue,thick,->] (1.05,0) arc[start angle=0,end angle=37,radius=1.05];
    \node[DeepBlue] at (1.25,0.48) {$\theta$};
    \draw[DeepBlue,thick,->] (-1.58,-1.25) arc[start angle=220,end angle=305,x radius=0.88,y radius=0.50];
    \node[below,font=\scriptsize] at (0,-2.48) {$\Phi_B=NAB\cos\theta$};
  \end{scope}
\end{tikzpicture}}%
\figcap{磁极间转动线圈、滑环与负载；右图给出磁场、线圈面积法向和夹角。}\label{ch24:fig:s40-9-4}
\end{center}}

% Explicit circuit topology and geometric references; no opaque label backgrounds.
\renewcommand{\generatorFigureFaithful}{%
\begin{center}\begin{tikzpicture}[scale=.83,>=Latex]
 \begin{scope}[xshift=-3.6cm]
  \draw[Ink,thick](-.18,1.2)--(-1,1.2)--(-1,3.1)--(1,3.1)--(1,1.2)--(.18,1.2);
  \draw[Ink,thick](-.18,1.2)--(-.5,0);
  \draw[Ink,thick](.18,1.2)--(.75,.9)--(.75,-.85)--(.5,-.85);
  \draw[guide,dashed](0,-1.3)--(0,3.6)node[above]{转轴};
  \foreach \y in {0,-.85}{\draw[DeepBlue,thick](0,\y)ellipse(.5 and .18);}
  \draw[Ink,line width=2pt](-.75,0)--(-.5,0) (.5,-.85)--(.95,-.85);
  \draw[Ink,thick](-.75,0)--(-1.5,0)--(-1.5,-2)--(-.5,-2);
  \draw[Ink,thick](.95,-.85)--(1.5,-.85)--(1.5,-2)--(.5,-2);
  \draw[Ink,thick](-.5,-2.18)rectangle(.5,-1.82);\node[below]at(0,-2.25){负载};
  \node[left]at(-1.6,0){电刷1};\node[right]at(1.55,-.85){电刷2};
  \node[right]at(.95,.35){滑环1};\node[anchor=north]at(0,-1.07){滑环2};
  \node at(-1.8,2.3){N};\node at(1.8,2.3){S};
  \draw[DeepBlue,->](-.7,3.55)arc(140:30:.9 and .3);
  \node[align=center]at(0,-3.65){两个滑环同轴且彼此绝缘\\线框两端经电刷接入负载};
 \end{scope}
 \begin{scope}[xshift=3.9cm]
  \draw[Ink,very thick](-1,2.7)--(1,.3);
  \draw[vector](0,1.5)--(2.3,1.5)node[right]{$\vect B$};
  \draw[vector](0,1.5)--(1.2,2.5)node[above]{$\vect n$};
  \draw(.65,1.5)arc(0:39.806:.65);\node at(.91,1.79){$\theta$};
  \node[align=center,anchor=north](generatorview)at(.4,-.4){沿转轴观察\\斜线为线框平面的侧影\\法线与侧影垂直};
  \node[align=center,anchor=north]at([yshift=-5mm]generatorview.south){单匝：$\Phi_B=BS\cos\theta$\\总磁链：$\Lambda=N\Phi_B$};
 \end{scope}
\end{tikzpicture}
\figcap{转动线框、共轴滑环与外部回路；右侧用沿转轴的视图确定面积法向。}
\end{center}}

\renewcommand{\unequalRodDiagram}{%
\begin{center}\begin{tikzpicture}[scale=.82,>=Latex]
 \draw[Ink,thick](-5,1)--(-1,1)--(1,1.8)--(5,1.8);
 \draw[Ink,thick](-5,-1)--(-1,-1)--(1,-1.8)--(5,-1.8);
 \draw[DeepBlue,very thick](-4,-1.15)--(-4,1.15)node[above]{$m_1$};
 \draw[DeepBlue,very thick](4,-1.95)--(4,1.95)node[above]{$m_2$};
 \draw[vector](4,-.5)--(4,.6)node[left]{$I$};\draw[vector](-4,.5)--(-4,-.5)node[left]{$I$};
 \draw[vector](4.25,0)--(5.4,0)node[right]{$v_2$};
 \draw[vector](-3.7,0)--(-2.7,0)node[above]{$v_1$};
 \node at(-2,-.5){$B_1,L_1$};\node at(2,-.6){$B_2,L_2$};
 \foreach \x in {-3,-2,2,3}{\node[TextGray]at(\x,.65){$\times$};}
 \node[align=center]at(0,-2.6){上下两根固定过渡导线分别连接；图示电流以右棒向上为正。};
\end{tikzpicture}\end{center}}

\renewcommand{\rodSourceDiagram}{%
\begin{center}\begin{tikzpicture}[scale=.8,>=Latex]
 \draw[Ink,thick](-3,1.5)--(3,1.5) (-3,-1.5)--(3,-1.5);
 \draw[Ink,thick](-2.7,-1.5)--(-2.7,-.8);
 \draw[Ink,thick](-2.9,-.8)rectangle(-2.5,-.15);
 \draw[Ink,thick](-2.7,-.15)--(-2.7,.25);
 \draw[Ink,thick](-2.95,.25)--(-2.45,.25) (-3.1,.65)--(-2.3,.65);
 \draw[Ink,thick](-2.7,.65)--(-2.7,1.5);
 \node[left]at(-3,-.5){$R$};\node[left]at(-3.15,.45){$\mathcal E$};
 \node[right]at(-2.35,.25){$-$};\node[right]at(-2.25,.65){$+$};
 \draw[DeepBlue,very thick](.8,-1.65)--(.8,1.65)node[above]{$m$};
 \draw[vector](.8,.7)--(.8,-.6)node[right]{$I$};
 \draw[vector](1.2,0)--(2.6,0)node[right]{$v$};
 \foreach \x in {-1.6,-.6,1.8}{\foreach \y in {-.9,.9}{\node[TextGray]at(\x,\y){$\times$};}}
\end{tikzpicture}\end{center}}

\renewcommand{\squareVoltmetersDiagram}{%
\begin{center}\begin{tikzpicture}[scale=.82,>=Latex]
 \draw[Ink,thick](-2,1.65)--(2,1.65)--(2,.45) (2,-.45)--(2,-1.65)--(-2,-1.65)--(-2,-.45) (-2,.45)--(-2,1.65);
 \draw[Ink,thick](-2.15,-.45)rectangle(-1.85,.45) (1.85,-.45)rectangle(2.15,.45);
 \node[right]at(-1.7,.0){$R_1$};\node[left]at(1.7,0){$R_2$};
 \node[circle,draw,minimum size=.7cm,inner sep=1pt](Vone)at(0,2.35){$V_1$};
 \node[circle,draw,minimum size=.7cm,inner sep=1pt](Vtwo)at(2.8,0){$V_2$};
 \node[circle,draw,minimum size=.7cm,inner sep=1pt](Vthree)at(0,-2.35){$V_3$};
 \node[circle,draw,minimum size=.7cm,inner sep=1pt](Vfour)at(-2.8,0){$V_4$};
 \draw(-2,1.65)--(-2,2.35)--(Vone.west) (Vone.east)--(2,2.35)--(2,1.65);
 \draw(2,1.65)--(2.8,1.65)--(Vtwo.north) (Vtwo.south)--(2.8,-1.65)--(2,-1.65);
 \draw(-2,-1.65)--(-2,-2.35)--(Vthree.west) (Vthree.east)--(2,-2.35)--(2,-1.65);
 \draw(-2,1.65)--(-2.8,1.65)--(Vfour.north) (Vfour.south)--(-2.8,-1.65)--(-2,-1.65);
 \foreach \x in {-2,2}{\foreach \y in {-1.65,1.65}{\fill(\x,\y)circle(2pt);}}
 \node[TextGray]at(0,0){$\otimes\ B(t)$};
 \draw[<->](-2,-3)--(2,-3)node[midway,below]{$a$};
 \node[align=left,anchor=west]at(4,.2){仪表、引线均与顶点相接\\为辨认路径放大引线间距\\计算取紧贴极限，附加面积趋零};
\end{tikzpicture}\end{center}}

\renewcommand{\circularVoltmetersDiagram}{%
\begin{center}\begin{tikzpicture}[scale=.8,>=Latex]
 \draw[Ink,thick](0,0)circle(2.4);
 \coordinate(A)at(150:2.4);\coordinate(C)at(30:2.4);
 \fill(A)circle(2pt)node[left]{$A$};\fill(C)circle(2pt)node[right]{$C$};
 \node[circle,draw,minimum size=.7cm,inner sep=1pt](Vone)at(0,3.15){$V_1$};
 \node[circle,draw,minimum size=.7cm,inner sep=1pt](Vtwo)at(0,-3.15){$V_2$};
 \draw(A)--(150:3.15)arc(150:100:3.15)--(Vone.west);
 \draw(Vone.east)--(80:3.15)arc(80:30:3.15)--(C);
 \draw(A)--(150:3.15)arc(150:260:3.15)--(Vtwo.west);
 \draw(Vtwo.east)--(280:3.15)arc(280:390:3.15)--(C);
 \draw[Ink](0,0)circle(.38);\node at(0,0){$V_3$};
 \draw(A)--(150:.38) (30:.38)--(C);
 \draw[guide](30:.8)arc(30:150:.8);\node at(0,1.12){$2\pi/3$};
 \node[TextGray]at(0,-1.2){$\otimes\quad \dot B=k$};
 \node[align=left,anchor=west]at(3.8,.3){$V_3$ 的接线沿 $AO,CO$\\仪表视为点状，$O$ 为圆心\\$V_1,V_2$ 引线取紧贴弧线极限\\图中间距只为辨认路径};
\end{tikzpicture}\end{center}}

\renewcommand{\gridCurrentDiagram}{%
\begin{center}\begin{tikzpicture}[scale=.85,>=Latex]
 \draw[Ink,thick](-2,-2)grid[step=2](2,2);
 \draw[DeepBlue,->,very thick](2,2)--(0,2);\draw[DeepBlue,->,very thick](0,2)--(-2,2);
 \draw[DeepBlue,->,very thick](-2,2)--(-2,0);\draw[DeepBlue,->,very thick](-2,0)--(-2,-2);
 \draw[DeepBlue,->,very thick](-2,-2)--(0,-2);\draw[DeepBlue,->,very thick](0,-2)--(2,-2);
 \draw[DeepBlue,->,very thick](2,-2)--(2,0);\draw[DeepBlue,->,very thick](2,0)--(2,2);
 \node at(-1,1){$\otimes$};\node at(1,1){$\otimes$};\node at(-1,-1){$\otimes$};\node at(1,-1){$\otimes$};
 \node[align=left,anchor=west]at(3,.4){$k>0$：外周逆时针\\每段电流 $i=ka^2/(2r)$\\内部四段电流为零};
 \draw[<->](-2,-2.5)--(2,-2.5)node[midway,below]{$2a$};
\end{tikzpicture}\end{center}}

\renewcommand{\rlcDiagram}[1]{%
\begin{center}\begin{tikzpicture}[scale=.9,>=Latex]
 \draw[Ink,thick](-4,1)--(-3.2,1) (-2.75,1)--(-1.8,1);
 \draw[Ink,thick](-3.2,.5)--(-3.2,1.5) (-2.75,.5)--(-2.75,1.5);
 \node[above]at(-3,1.75){$C,\ U$};
 \ifnum#1=1\node[below]at(-3.25,.45){$+$};\node[below]at(-2.7,.45){$-$};
 \else\node[below]at(-3.25,.45){$-$};\node[below]at(-2.7,.45){$+$};\fi
 \foreach \xx in {-1.8,-1.3,-.8,-.3}{\draw[Ink,thick](\xx,1)arc(180:0:.25 and .3);}
 \draw[Ink,thick](.2,1)--(.65,1);\draw[Ink,thick](.65,.8)rectangle(1.85,1.2);
 \draw[Ink,thick](1.85,1)--(3,1)--(3,-1)--(1.2,-1);
 \node[above]at(-.8,1.45){$L$};\node[above]at(1.25,1.4){$R$};
 \fill(1.2,-1)circle(1.6pt);\fill(.3,-1)circle(1.6pt);
 \draw[Ink,thick](.3,-1)--(1.2,-1);\node[below]at(.75,-1.15){$S$（已闭合）};
 \ifnum#1=1
  \draw[Ink,thick](.3,-1)--(-1,-1) (-1.4,-1)--(-4,-1)--(-4,1);
  \draw[Ink,thick](-1.4,-1.5)--(-1.4,-.5) (-1,-1.3)--(-1,-.7);
  \node[below]at(-1.2,-1.65){$\mathcal E$};\node[above]at(-1.55,-.5){$+$};
 \else\draw[Ink,thick](.3,-1)--(-4,-1)--(-4,1);\fi
 \draw[vector](-1.7,0)--(.15,0)node[midway,above]{$I(t)$};
\end{tikzpicture}\end{center}}

\renewcommand{\chargedRingDiagram}{%
\begin{center}\begin{tikzpicture}[scale=.82,>=Latex]
 \draw[Ink,thick](0,0)circle(2.1);\fill(2.1,0)circle(2.4pt)node[above right]{$m,+q$};
 \draw[guide,dashed](0,0)--(2.1,0);\node[above]at(.95,0){$R$};
 \draw[DeepBlue,->](0:2.45)arc(0:-60:2.45);\node[right]at(2.5,-1.3){$\theta$ 正向};
 \draw[vector](2.1,0)--(2.1,-1.2)node[left]{$\vect e_\theta$};
 \draw[vector](2.1,0)--(3.5,0)node[right]{$\vect e_r$};
 \node[TextGray]at(-.5,.9){$\otimes\ B(t)$};
 \node[align=left,anchor=west]at(4.5,.4){$N_r>0$ 向外\\$N_r<0$ 向内\\图示 $\theta=0$ 的基矢};
\end{tikzpicture}\end{center}}

\renewcommand{\travelingFieldTrainDiagram}{%
\begin{center}\begin{tikzpicture}[scale=.8,>=Latex]
 \draw[Ink,thick](-4,1.4)--(4,1.4) (-4,-1.4)--(4,-1.4);
 \draw[DeepBlue,very thick](-1.2,-1.15)rectangle(1.2,1.15);
 \foreach \x in {-3.3,-2.5,2.5,3.3}{\node[TextGray]at(\x,.55){$\odot$};}
 \foreach \x in {-.7,0,.7}{\node[TextGray]at(\x,.55){$\otimes$};}
 \draw[guide,dashed](-1.9,-1.3)--(-1.9,1.3) (1.9,-1.3)--(1.9,1.3);
 \node at(0,-.35){$m$};\draw[vector](1.45,0)--(3,0)node[right]{$v$};
 \draw[vector](-1.8,2)--(1.8,2)node[right]{$v_\phi=\omega/k$};
 \draw[<->](-1.2,-1.9)--(1.2,-1.9)node[midway,below]{$a$};
 \draw[<->](-4.4,-1.15)--(-4.4,1.15)node[midway,left]{$L$};
 \node[align=center]at(0,-2.9){某时刻的正负场区，虚线处 $B=0$；导轨只作绝缘机械支承。};
\end{tikzpicture}\end{center}}

\renewcommand{\solenoidSplitDiagram}{%
\begin{center}\begin{tikzpicture}[scale=.75,>=Latex]
 \foreach \x in {-3.6,-3.2,...,3.6}{\draw[Ink](\x,0)ellipse(.18 and 1.0);}
 \draw[guide,dashed](0,-1.4)--(0,1.5);\draw[vector](-3.4,0)--(3.4,0);
 \draw[<->](-3.8,-1.5)--(0,-1.5)node[midway,below]{$l_1$};
 \draw[<->](0,-1.5)--(3.8,-1.5)node[midway,below]{$l_2$};
 \node[align=left,anchor=west]at(4.5,.2){保持位置、同向串联：\\$L=L_1+L_2+2M$\\两子段均很长，才可近似略去端部互感。};
\end{tikzpicture}\end{center}}

\renewcommand{\mutualLoopsProofDiagram}{%
\begin{center}\begin{tikzpicture}[scale=.82,>=Latex]
 \draw[Ink,thick](-2.5,0)circle(1.2);\draw[Ink,thick](2.5,.4)circle(1);
 \coordinate(A)at(-1.3,0);\coordinate(B)at(1.5,.4);
 \fill(A)circle(2pt);\fill(B)circle(2pt);
 \draw[vector](A)--(-1.3,1)node[right]{$\dd\vect\ell_1$};
 \draw[vector](B)--(1.5,-.6)node[right]{$\dd\vect\ell_2$};
 \draw[vector](A)--(B)node[midway,above]{$\vect r_{12}$};
 \node at(-2.5,0){$C_1$};\node at(2.5,.4){$C_2$};
 \node[align=left,anchor=west]at(4.2,.2){源点在各自回路上\\线元沿回路切线\\$r_{12}=|\vect r_2-\vect r_1|$};
\end{tikzpicture}\end{center}}

\renewcommand{\mutualInductanceDiagram}{%
\begin{center}\begin{tikzpicture}[scale=.8,>=Latex]
 \foreach \x/\lab in {-2/1,2/2}{
  \draw[Ink,thick](\x,0)ellipse(.45 and 1.5);
  \draw[DeepBlue,->]({\x+.45},0)arc(0:65:.45 and 1.5);
  \node[below]at(\x,-1.7){回路\lab};}
 \draw[vector](-3.2,0)--(3.2,0)node[right]{$\vect B_1$};
 \node[align=left,anchor=west]at(4.4,.2){同轴同向时互感为正\\$\Phi_{2\leftarrow1}=MI_1$\\$\Phi_{1\leftarrow2}=MI_2$};
\end{tikzpicture}\end{center}}

\newcommand{\selfFieldLinesDiagram}{%
\begin{center}\begin{tikzpicture}[scale=.66,>=Latex]
 \foreach \xx in {-3.8,3.8}{\begin{scope}[xshift=\xx cm]
  \foreach \sg in {-1,1}{\foreach \a/\b/\h in {.2/2.6/3.8,.6/1.85/2.8}{
   \draw[DeepBlue,thick,postaction={decorate},decoration={markings,mark=at position .5 with {\arrow{Latex}}}]
    ({\sg*\a},0)..controls({\sg*\a},\h)and({\sg*\b},\h)..({\sg*\b},0)
    ..controls({\sg*\b},-\h)and({\sg*\a},-\h)..cycle;}}
  \draw[vector](0,-.8)--(0,3.25)node[above]{$B$};
  \ifdim\xx pt<0pt
   \draw[Ink,thick](0,0)ellipse(1.45 and .38);
   \draw[Ink,->](230:1.45 and .38)arc(230:310:1.45 and .38);
   \node at(0,-3.65){单匝圆电流的自生磁场};
  \else
   \foreach \y in {-1.3,-.975,...,1.3}{\draw[Ink,thick](0,\y)ellipse(1.15 and .25);}
   \node at(0,-3.65){有限长同向绕组的回流场};
  \fi
 \end{scope}}
\end{tikzpicture}\end{center}}


第23章研究了稳恒电流与磁场；本章放开“稳恒”限制：回路运动、磁场随时间变化，或两者
同时发生时，磁通量的变化会在回路中产生感应电动势。法拉第电磁感应定律（Faraday's law）给出数值，楞次定律（Lenz's law）
把能量守恒写进了方向。

历史上，电流产生磁场以及磁场作用于电流先被发现。由此自然产生一个反向问题：既然运动电荷能产生
磁效应，磁场是否也能在适当条件下产生电效应？法拉第经过长期实验确认，导体在磁场中运动可以使
回路出现电流；随后又发现，即使导体固定，只要磁场大小或方向变化，同样可以产生电动势。前一类常称
动生电动势，后一类常称感生电动势。两类机制的微观解释不同，却都由回路磁通量的变化统一描述。

这一发现既有应用价值，也有理论意义。它使发电机成为可能，从而把水流、蒸汽轮机或其他动力装置的
机械能转成电能；它还首次明确表明，变化磁场能在空间中产生一种不同于静电场的电场。经典电磁学中，
电场既可由电荷产生，也可由变化磁场产生。麦克斯韦后来又补上“变化电场产生磁场”，两种场便能相互
维持并以电磁波传播，电、磁与光由此进入同一理论框架。

法拉第的实验装置可以用两个彼此绝缘、绕在同一铁环上的线圈来理解。一个线圈接电池，另一个线圈
接检流计：接通电源的瞬间，检流计发生一次偏转；电流稳定后，偏转消失；断开电源时，检流计又向相反
方向偏转。这说明感应效应不取决于“附近有没有电流”，而取决于磁场是否正在变化。把永久磁铁插入或
拔出线圈，也能观察到相同规律：运动越快，磁通变化越快，感应电动势越大；磁铁停住后，即使线圈中仍
有很强磁场，也不会维持感应电流。由这些现象抽象出来的核心量不是磁场本身，而是穿过回路的磁通量。

\faradayHistoryContext


\subsection{法拉第--楞次定律}


对闭合回路 $C$ 选定正向，再用右手螺旋法则确定跨越面 $S$ 的正法向。磁通量与感应电动势为
\begin{equation}
  \Phi_B=\int_S\vect B\cdot\dd\vect S,
  \qquad
  \boxed{\mathcal E=-\frac{\dd\Phi_B}{\dd t}}.
\end{equation}
负号不是额外的计算规则，它表示感应电流所产生的磁效应总是阻碍原磁通的变化。若是 $N$ 匝
线圈，各匝穿过的磁通都为 $\Phi_B$，则磁链 $\Lambda=N\Phi_B$，且 $\mathcal E=-\dd\Lambda/\dd t$。

图~\ref{ch24:fig:s40-9-2} 把同一装置的三个时刻并列出来。上排磁铁靠近线圈，磁通变化率非零，检流计偏转；中排磁铁
停住，虽然线圈中仍有磁场，但磁通不再变化，指针回到零位；下排磁铁远离，磁通变化方向反转，因而
电流和指针偏转方向也随之反转。这里真正决定电动势的是运动造成的磁通变化率，不是磁铁与线圈的距离本身。

\movingMagnetExperimentFigure

对匀强磁场中的平面回路，磁通可写为
\begin{equation}
  \Phi_B=BS\cos\theta.
\end{equation}
因此其时间导数包含三类贡献：
\begin{equation}
 \frac{\dd\Phi_B}{\dd t}
 =S\cos\theta\,\frac{\dd B}{\dd t}
 +B\cos\theta\,\frac{\dd S}{\dd t}
 -BS\sin\theta\,\frac{\dd\theta}{\dd t}.
\end{equation}
第一项对应磁场大小变化，第二项对应回路面积或形状变化，第三项对应线圈转动。一般非匀强场中不能
把 $B$ 提到积分号外，但“求总磁通的完整时间变化率”这一原则不变。若多匝线圈各匝磁通不同，应逐匝
相加形成磁链，而不能机械地用 $N$ 乘某一匝的磁通。

对固定回路，电动势是沿回路一周非静电驱动力对单位正电荷的线积分；沿正向为正，反之为负。
对运动导线，应使用 $\oint(\vect E+\vect v_{\rm wire}\times\vect B)\cdot\dd\vect\ell$，
其中 $\vect v_{\rm wire}$ 是导线元速度，不是载流子的总速度。这是瞬时回路的驱动量；
磁力对载流子实际运动的总功仍为零，输出能量来自维持导线运动的机械功。
回路正向一旦选定，曲面正法向就由右手规则确定。由于磁场
高斯定律给出任意闭合面净磁通为零，只要两个曲面具有同一个边界，它们的磁通量相同，因此无需再为
同一回路指定唯一的“计算曲面”。这保证法拉第定律不会因跨越面的形状不同而含糊。

电动势的单位是 volt，但它与两点间静电势差并非总是同一个概念。静电场中，单位电荷从一点移到另一点
所受电场力做功只与端点有关，沿闭合回路一周总功为零；感生电场中，沿闭合回路一周的功可以非零，
因此不能给整条回路上的每一点分配一个全局单值电势来替代环量。电压表在具体接线下仍会给出读数，但
在时变磁场问题中，读数还可能依赖表笔引线所围成的路径，必须把测量回路本身纳入磁通分析。

磁通量可以因三种因素改变：磁场大小变化、回路面积或形状变化、回路法向与磁场夹角变化。法拉第
定律只问总磁通量怎样变化，并不预先区分是哪一种因素造成；真正需要区分的是驱动电荷的局部力来自
$\vect v\times\vect B$，还是来自变化磁场产生的涡旋电场。

图~\ref{ch24:fig:review-fluxorientationdiagram-1}中，环路正向和法向按右手规则绑定，负号相对于这套共同约定解释。
\par\noindent\begin{minipage}{\linewidth}
\fluxOrientationDiagram
\figcap{回路正向与磁通法向}\label{ch24:fig:review-fluxorientationdiagram-1}
\end{minipage}\par\medskip

“阻碍变化”要与“阻碍原磁场”区分。原磁通增大时，感应磁场与它反向；原磁通减小时，感应磁场反而
与它同向。因此对方向题，先判断原磁通的正负和增减，再确定感应磁场，最后用右手螺旋法则找电流。

负号也可由能量守恒检验。若磁铁靠近线圈时，感应电流产生的磁场反而帮助磁铁继续靠近，磁铁就会在
没有外界做功的情况下越动越快，同时线圈中还产生电能，这将造成能量无源增加。真实情况是靠近时受到
排斥、远离时受到吸引，外力必须克服电磁阻力做功，所做功转化为电能和焦耳热。楞次定律不是一条
与法拉第定律并列的独立大小公式，而是法拉第定律负号所表达的方向与能量约束。

实际判向时可以固定四步：先选回路正向并由右手规则确定面积正法向；再写出 $\Phi_B$ 的正负；判断
$\dd\Phi_B/\dd t$；最后由 $\mathcal E=-\dd\Phi_B/\dd t$ 判断正电动势方向。若只凭“增反减同”口诀
而不先确定原磁场穿面的方向，在立体线圈、翻转线圈或多磁场叠加问题中很容易把方向判反。


% EM teaching enrichment 20260927
\paragraph{磁通为零的瞬间仍可有电动势}
固定面积 $A$ 的单匝回路处于 $B(t)=B_0\sin\omega t$ 的垂直匀强磁场中，取场的正方向为面法向，则
$\Phi=B_0A\sin\omega t$，$\mathcal E=-\omega B_0A\cos\omega t$。
$t=0$ 时磁通为零而电动势大小最大；$t=\pi/(2\omega)$ 时磁通最大而电动势为零。决定感应的是变化率，不是磁通本身。
若回路只含恒定电阻 $R$、自感可忽略，某段过程的有向通过电量为
\[
 \Delta q=\int I\dd t=-\frac{\Delta\Phi}{R}.
\]
它只取决于首末磁通；有显著自感时还要计入 $L\Delta I$，不能无条件沿用。


\subsection{交流发电机}


发电机是电磁感应的直接应用。最简单的交流发电机让 $N$ 匝、面积为 $A$ 的线圈在匀强磁场 $B$ 中以
角速度 $\omega$ 转动。若 $t=0$ 时面积法向与磁场同向，磁链为
\begin{equation}
  \Lambda(t)=NAB\cos\omega t,
\end{equation}
于是
\begin{equation}
  \boxed{\mathcal E(t)=NAB\omega\sin\omega t
  \equiv\mathcal E_0\sin\omega t}.
\end{equation}
若外电路近似纯电阻 $R$，电流 $I(t)=\mathcal E_0\sin\omega t/R$，方向每半周反转，因此称为交流。

当线圈法向与 $\vect B$ 平行时，磁通量达到极值，但瞬时变化率为零，所以电动势为零；当线圈平面与
$\vect B$ 平行、法向垂直于磁场时，磁通为零却变化最快，电动势达到峰值。这一“量为零而变化率最大”
的关系是理解正弦、余弦相差四分之一周期的关键。若初始相位不同，只会把
$\sin\omega t$ 改为带相位常数的正弦函数，峰值 $\mathcal E_0=NAB\omega$ 不变。

若关心交流在电阻上的平均发热效果，可引入有效值
\begin{equation}
  \mathcal E_{\rm rms}=\frac{\mathcal E_0}{\sqrt2},
  \qquad I_{\rm rms}=\frac{I_0}{\sqrt2}.
\end{equation}
于是周期平均功率仍写成 $\langle P\rangle=I_{\rm rms}^2R$。有效值不是瞬时电动势的算术平均；
正弦电动势一周期平均为零，但平方平均不为零，仍能持续向电阻输送能量。

\generatorFigureFaithful

水力发电把水的重力势能经涡轮转为转轴机械能；火力与核电站先以热源产生蒸汽，再驱动涡轮。能量来源
虽然不同，末级的电磁转换都依靠线圈与磁场的相对转动。负载接通后出现的电磁力矩总是反抗驱动转矩，
因此发出的电功率最终来自原动机做功，而不是磁场无中生有地提供能量。

交流发电机通常用两个彼此绝缘的滑环把转动线圈两端接到静止电刷，使外电路得到方向周期反转的电压。
若改用分裂换向器，可以把线圈内部的交变电动势整流为外电路中方向不变、大小脉动的电动势。发电机与
电动机的结构有很强的互逆性：外力驱动线圈转动时输出电能；给线圈输入电流时，安培力矩又能输出
机械运动。负载越大，电流越大，反向电磁转矩也越大，原动机必须提供更大转矩才能维持同一角速度。

若回路同时运动且存在感生电场，完整形式为
\begin{equation}
  \boxed{\mathcal E=\oint_{C(t)}(\vect E+\vect v\times\vect B)\cdot\dd\vect\ell
  =-\frac{\dd}{\dd t}\int_{S(t)}\vect B\cdot\dd\vect S}.
\end{equation}
左边第一项是变化磁场产生的感生电场，第二项是回路运动产生的动生效应；右边则把它们统一为“跟着
回路计算的磁通变化率”。

这里的时间导数不是只在固定空间区域上对 $B$ 求偏导，而是对随回路一起运动、变形的曲面求总导数。
因此即使 $\partial\vect B/\partial t=0$，曲面边界运动仍可使磁通变化并产生动生电动势；反之，即使
回路完全固定，$\partial\vect B/\partial t\ne0$ 也会产生感生电动势。把这两个来源分开有助于解释
局部受力，把它们合在同一磁通定律中则更适合计算整个回路。

若导体静止而磁体移动，导线中的自由电荷初始并没有随导体切割磁感线，不能用“导体内电荷的
$q\vect v\times\vect B$”作为唯一解释；法拉第定律仍预言回路中有电动势，说明运动或变化的磁场
在导体所在位置建立了电场。换到随带电粒子运动的参考系，可以把低速情形下的磁力效果等价地写成
\begin{equation}
  \vect E'\simeq\vect v\times\vect B=-\vect v_B\times\vect B,
\end{equation}
其中 $\vect v_B$ 是磁场源相对粒子的速度。这种推理能给出低速方向和数量级，但电场、磁场在不同惯性系
之间如何严格变换属于狭义相对论；因此它只能作为低速极限的物理说明，不能替代完整的洛伦兹变换。

“磁铁动、线圈不动”与“线圈动、磁铁不动”在实验中都能产生电流，但经典局部解释看起来不同：前者
强调变化磁场产生电场，后者强调运动电荷受到磁力。相对论统一了这种差别，因为电场和磁场不是彼此
绝对独立的实体，而是同一电磁场在不同惯性系中的分量。中学范围不需要计算完整变换，却应保留一个
重要结论：不同参考系中的电荷、电流、场和能量须按相应变换关联，能量数值不必相同；经变换后的物理预言必须相容，不能因选择“谁在运动”而得到互相矛盾
的实验预言。

图~\ref{ch24:fig:review-movingloopdiagram-1}中，图示各点速度相同，面积与法向不变，正负的局部动生贡献沿全回路相消。
\par\noindent\begin{minipage}{\linewidth}
\movingLoopDiagram
\figcap{匀强场中闭合回路的刚性平移}\label{ch24:fig:review-movingloopdiagram-1}
\end{minipage}\par\medskip


% EM teaching enrichment 20260927
\paragraph{从磁通相位到输入功率}
面积 $A$、匝数 $N$ 的线圈以角速度 $\omega$ 转动，若其法向与磁场夹角为 $\omega t$，则磁链为 $NBA\cos\omega t$，电动势为 $NBA\omega\sin\omega t$。
接纯电阻总负载 $R$ 并忽略自感后，周期平均输出为
\[
 \overline P=\frac{(NBA\omega)^2}{2R}.
\]
理想无机械损耗时，外界平均驱动力矩为 $\overline P/\omega$。线圈转到磁通最大的位置，瞬时电动势反而为零；输出电压的峰值由磁通变化最快的位置决定。加入电感后，电流不再与电动势同相，平均功率须按下一章的相位关系计算。


% EM structure and teaching revision 20260928
\subsection{磁通变化的三种来源与符号流程}
均匀磁场中的平面回路有 $\Phi=BS\cos\theta$。磁场大小、回路面积和取向都可能变化，因此
\[
 \mathcal E=-\frac{\dd\Phi}{\dd t}
 =-S\cos\theta\,\dot B-B\cos\theta\,\dot S
 +BS\sin\theta\,\dot\theta.
\]
这个展开要求回路仍可由一个统一面积和法向描述；一般变形回路应回到曲面积分。
例如回路面积为 $S_0$，在恒定磁场中以 $\theta=\omega t$ 转动，则 $\mathcal E=BS_0\omega\sin\omega t$。磁通最大时电动势为零，磁通经过零值时电动势幅值最大，二者不能按“谁大谁强”判断。
解题先任选回路绕行正向，再由右手关系确定面法向，接着写带符号磁通，最后用负导数得到沿所选绕行方向的电动势。结果为负表示实际方向与参考方向相反。楞次定律可作物理检查：感应电流的磁效应反抗磁通变化，不是无条件反抗外磁场本身。


\subsection{专题例题与方法应用}

\begin{origexamplebox}{【大小圆环】}\label{ch24:ex:s40-9-1}
两个同心导电圆环放置在同一平面上，回路电阻大小均为 $r$。

(1) 若大圆环中突然通入顺时针电流，求小圆环中感生电流的方向；

(2) 若小圆环中突然通入顺时针电流，求大圆环中感生电流的方向。

% Pure source-question geometry: no derived forces, currents or solution curves.
\begin{center}\begin{tikzpicture}[scale=.68,>=Latex]
 \foreach \x/\case in {0/1,7/2}{\begin{scope}[shift={(\x,0)}]
 \draw[thick](0,0)circle(2);\draw[DeepBlue,thick](0,0)circle(1.1);
 \ifnum\case=1\draw[vector](2,0)arc[start angle=0,end angle=-65,radius=2];\node at(0,2.5){大环电流突然建立};
 \else\draw[vector](1.1,0)arc[start angle=0,end angle=-65,radius=1.1];\node at(0,2.5){小环电流突然建立};\fi
 \node[below]at(0,-2.3){两环电阻均为 $r$};\end{scope}}
\end{tikzpicture}\end{center}
\figcap{两种激励情形；箭头仅表示题设中主动建立的顺时针电流。}\label{ch24:fig:em24-question-1}
\end{origexamplebox}

\begin{xesolution}
规定垂直纸面向外为磁通正方向。

(1) 大环的顺时针电流使穿过小环的向里磁通从零增大。小环要用向外的磁场阻碍这个增量，故其感生电流
为逆时针。

(2) 互感系数具有互易性，$M_{12}=M_{21}$。小环的顺时针电流同样使大环中的向里磁通增大，所以大环
中的感生电流也为逆时针。结果与谁是“大环”无关，只由磁通的变化方向决定。

图~\ref{ch24:fig:review-concentricloopsdiagram-1}中，所激励环的顺时针电流增加时，另一环的感应电流为逆时针；交换激励环仍须重算所穿磁通。
\par\noindent\begin{minipage}{\linewidth}
\concentricLoopsDiagram
\figcap{同心大小圆环的感应方向}\label{ch24:fig:review-concentricloopsdiagram-1}
\end{minipage}\par\medskip
\end{xesolution}



\begin{origexamplebox}{【奇怪的圆环】}\label{ch24:ex:s40-9-2}
一个条形磁铁的周围有三个导体圆环，分别位于 $S$ 极左侧、磁铁中部和 $N$ 极右侧，三环与磁铁同轴。
若磁铁突然向左运动，判断三个圆环的加速度方向，以及它们各自有扩张还是收缩的趋势。

% Pure source-question geometry: no derived forces, currents or solution curves.
%
\begin{center}\begin{tikzpicture}[scale=0.8,>=Latex]
  \draw[DeepBlue,dashed,very thick] (0,1.25) arc(90:270:.33 and 1.25);
  \draw[fill=SoftGray] (-1.8,-0.65) rectangle (1.8,0.65);
  \node at (-1.25,0) {$S$}; \node at (1.25,0) {$N$};
  \foreach \x in {-3.4,3.4}{\draw[DeepBlue,very thick] (\x,0) ellipse (0.33 and 1.25);}
  \draw[DeepBlue,very thick] (0,-1.25) arc(270:450:.33 and 1.25);
  \draw[vector] (-1.7,1.1)--(-2.7,1.1) node[midway,above=4pt] {$\vect v_M$};
  \end{tikzpicture}\end{center}
\figcap{磁铁开始左移，三个导体圆环的运动与形变待求。}\label{ch24:fig:em24-question-2}
\end{origexamplebox}

\begin{xesolution}
轴向力与径向张力必须分别判断，不能用环内总磁通的方向代替环线所在处的局部磁场方向。
忽略自感、初始环静止时，设磁铁向左速率为 $V>0$，以磁铁中心为原点、向右为 $z$ 正向，
记半径固定的环所穿正向磁通为 $\Phi(z)$、电阻为 $R_c$，则
\[
 I=-\frac{V}{R_c}\frac{\dd\Phi}{\dd z},\qquad
 F_z=I\frac{\dd\Phi}{\dd z}=-\frac V{R_c}\left(\frac{\dd\Phi}{\dd z}\right)^2\le0.
\]
所以非对称位置的圆环初始均向左加速；若保留自感，电流需要建立，严格初始瞬间可仍为零。

左环面对正在接近的 $S$ 极，穿过它的磁通增大。感应电流使圆环与接近的磁极相斥，因此左环向左加速；其感应磁矩
与当地轴向外磁场反向。在环线处的场仍朝右的通常近轴极外位置，径向安培力向内，故收缩。

右环面对正在远离的 $N$ 极，原磁通减小。它用同向磁矩维持磁通，与磁铁相吸，所以向左加速；
在环线处轴向场朝右时，径向力向外，故扩张。若环很大，线处的轴向场可能反向，就不能沿用这个径向结论。

中环在完全对称的初始位置上有 $\partial\Phi/\partial x=0$，所以理想化的瞬间感应电流和加速度都可为零。磁铁稍微左移后，
通过中环的磁通减小，中环会向左追随磁铁。在磁铁外侧的近赤道区域，局部轴向场与磁矩相反，
因此中环稍后通常是收缩，不是扩张。

可用磁偶极模型检验这一点。设磁矩大小为 $\mu_M$、环半径为 $a$，则
\[
 \Phi(z)=\frac{\mu_0\mu_Ma^2}{2(a^2+z^2)^{3/2}},\quad
 B_z(a,z)=\frac{\mu_0\mu_M(2z^2-a^2)}{4\pi(a^2+z^2)^{5/2}}.
\]
对 $z>0$ 有 $\Phi'<0$，感应电流为正；在 $0<z<a/\sqrt2$ 时局部 $B_z<0$，
每段径向力 $\dd F_a=I B_z\dd\ell<0$，即收缩；更远处才变为扩张。
这也说明没有环半径和位置等几何条件时，应给出条件分支，而不能把三个径向趋势当作普遍结论。

图~\ref{ch24:fig:review-magnetthreeloopsdiagram-1}中，轴向箭头描述对相对运动的阻碍；径向标记只适用于图示近轴极外和近中部的局部场条件。
\par\noindent\begin{minipage}{\linewidth}
\magnetThreeLoopsDiagram
\figcap{磁铁左移时的三个圆环}\label{ch24:fig:review-magnetthreeloopsdiagram-1}
\end{minipage}\par\medskip
\end{xesolution}



\section{动生电动势与机电耦合}

\subsection{动生电动势与运动回路}


导体以速度 $\vect v$ 在磁场中运动时，导体内电荷受到 $q\vect v\times\vect B$ 的磁力。对一段运动导线，
动生电动势为
\begin{equation}
  \boxed{\mathcal E_{\rm m}=\int(\vect v\times\vect B)\cdot\dd\vect\ell}.
\end{equation}
长为 $\ell$ 的直棒以速率 $v$ 垂直切割匀强磁场时，$\mathcal E_{\rm m}=B\ell v$。它是棒两端的非静电力做功
与单位电荷的比，不能简单说成“磁力对电荷做功”；能量来自维持导体运动的外力或导体的动能。

对直导体棒，若 $\vect v,\vect B$ 与棒方向彼此垂直，正电荷受力方向由 $\vect v\times\vect B$ 决定，
并在棒的一端积累；另一端留下异号电荷。开路时，电荷分离建立静电场，直到
$q\vect E+q\vect v\times\vect B=\vect0$，此后不再有持续电流，但两端仍有电势差 $B\ell v$。
闭合回路后，电荷能够沿外电路流动，棒相当于电源。若三者不完全垂直，只有
$\vect v\times\vect B$ 沿棒方向的分量贡献电动势，所以一般式必须保留点积，不能总写成 $B\ell v$。

以长度为 $\ell$ 的导体棒在匀强磁场中匀速滑动为例，若回路总电阻为 $R$，则
\begin{equation}
  I=\frac{B\ell v}{R},\qquad
  F_B=BI\ell=\frac{B^2\ell^2}{R}v.
\end{equation}
安培力阻碍棒的运动。维持匀速所需外力功率
\begin{equation}
  P_{\rm ext}=F_Bv=\frac{B^2\ell^2v^2}{R}
\end{equation}
恰好等于电阻热功率
\begin{equation}
  P_R=I^2R=\frac{B^2\ell^2v^2}{R}.
\end{equation}
所以磁力负责改变载流子的运动方向和建立电荷分离，真正输入系统的能量来自外力。若撤去外力，棒的
动能会在电磁阻尼作用下转化为焦耳热；这正是楞次定律在能量上的含义。

还可以直接由磁通验证。棒在平行导轨上移动距离 $\dd x=v\dd t$，回路面积增加
$\dd S=\ell v\dd t$，于是 $|\dd\Phi_B/\dd t|=B\ell v$。洛伦兹力积分与磁通变化率给出同一结果，
但它们强调的层次不同：前者说明棒内电荷怎样受力，后者不必追踪载流子，直接从整个回路的几何变化
求电动势。遇到多根棒同时运动或回路形状复杂时，用完整的运动回路积分通常更稳妥。

图~\ref{ch24:fig:review-motionalroddiagram-1}中，棒向右、B向纸内，正电荷磁力沿棒向上，因此环流逆时针，棒受到向左的阻力。
\par\noindent\begin{minipage}{\linewidth}
\motionalRodDiagram
\figcap{运动导体棒与闭合电阻回路}\label{ch24:fig:review-motionalroddiagram-1}
\end{minipage}\par\medskip


% EM teaching enrichment 20260927
\paragraph{导体棒的制动：电路与运动必须联立}
长度 $l$、质量 $m$ 的棒沿水平导轨滑动，闭合回路总电阻 $R$ 恒定、磁场 $B$ 垂直导轨面，忽略自感和摩擦。棒速 $v$ 时
$\mathcal E=Blv$，$I=Blv/R$，磁力阻碍滑动，因此
\[
 m\dot v=-\frac{B^2l^2}{R}v,\qquad
 v=v_0\exp\left(-\frac{B^2l^2}{mR}t\right).
\]
瞬时机械损失率 $-\dd(mv^2/2)/\dd t$ 恰等于 $I^2R$，总焦耳热为初动能。
若要求恒速，就必须有外力补偿磁阻力；磁场只是使机械功与电能交换的耦合条件，不能被当作持续供能的能源。


% EM structure and teaching revision 20260928
\subsection{机电耦合：两条方程与一份能量账}
导体棒长 $\ell$、质量 $m$，在匀强磁场中沿导轨运动。忽略回路自感、取总电阻 $R$，则感应电流大小为 $I=B\ell v/R$，磁力反向阻碍运动，故
\[
 m\dot v=F_{\mathrm{ext}}-\frac{B^2\ell^2}{R}v.
\]
恒定外力下，终端速度 $v_\infty=F_{\mathrm{ext}}R/(B^2\ell^2)$，趋近它的时间常数为 $\tau_m=mR/(B^2\ell^2)$。当回路电感 $L$ 不可忽略且近似不随棒的位置变化时（例如固定外接电感占主导），取电感电压关系 $u_L=L\dot I$，必须同时求解
\[
 L\dot I+RI=B\ell v,\qquad m\dot v=F_{\mathrm{ext}}-B\ell I.
\]
若 $L/R\ll\tau_m$，电流才可近似及时跟随速度，退回第一式。把两式分别乘以 $I$ 和 $v$ 后相加，得到
\[
 F_{\mathrm{ext}}v=\frac{\dd}{\dd t}\left(\frac12mv^2+\frac12LI^2\right)+I^2R.
\]
自感电压关系及其能量将在本章最后一节推导。若运动使 $L(x)$ 显著变化，还会产生相应的电动势和自感磁力，不能继续使用这里的常量 $L$ 模型。这条能量式能检验耦合项符号：机械能转出的 $B\ell Iv$ 必须恰好进入电路，不能在两边同时作为损耗。接入电容、电源或增加第二根可动棒时，仍按同样顺序增列各自的状态变量与约束。


\subsection{专题例题与方法应用}

\begin{origexamplebox}{【小棒加电阻】}\label{ch24:ex:s40-9-3}
一根导体棒横跨在宽度为 $L$ 的光滑导轨上，导轨间有垂直纸面向里的匀强磁场 $B$。棒的质量为 $m$，左侧电阻为 $R$，
导轨和棒的电阻不计。初始时棒以速度 $v_0$ 向右运动。

(1) 求任意时刻导体棒的速度；

(2) 求从开始运动到 $t\to\infty$ 通过电阻的总电量。

% Pure source-question geometry: no derived forces, currents or solution curves.
%
\begin{center}\begin{tikzpicture}[scale=0.78,>=Latex]
  \draw[Ink,very thick] (-3,1.4)--(3,1.4); \draw[Ink,very thick] (-3,-1.4)--(3,-1.4);
  \draw[Ink,very thick] (-2.7,-1.4)--(-2.7,-0.45) decorate[decoration={zigzag,segment length=4pt,amplitude=2pt}] {-- (-2.7,0.45)}--(-2.7,1.4);
  \node[left=8pt] at (-2.7,0) {$R$};
  \draw[DeepBlue,very thick] (0.8,-1.55)--(0.8,1.55); \node[above] at (0.8,1.55) {$m$};
  \draw[vector] (1.0,0)--(2.4,0) node[right] {$v_0$};
  \draw[<->] (-3.7,-1.4)--node[left] {$L$}(-3.7,1.4);
  \foreach \x in {-1.7,-0.7,0.3,1.4,2.3}{\foreach \y in {-0.75,0.75}{\node[TextGray] at (\x,\y) {$\times$};}}
  \end{tikzpicture}\end{center}
\figcap{棒与定值电阻组成的闭合导轨回路。}\label{ch24:fig:em24-question-3}
\end{origexamplebox}

\begin{xesolution}
棒向右运动时回路面积增大，向里磁通增大。感应电流的安培力必然向左。当速度为 $v$ 时，
\begin{equation}
  \mathcal E=BLv,
  \qquad I=\frac{BL}{R}v,
  \qquad m\dot v=-BLI=-\frac{B^2L^2}{R}v.
\end{equation}
令
\begin{equation}
  \tau=\frac{mR}{B^2L^2},
\end{equation}
得
\begin{equation}
  \boxed{v(t)=v_0\e^{-t/\tau}}.
\end{equation}
总电量可以直接积分，也可用冲量消去时间：
\begin{equation}
  q_{\rm tot}=\int_0^\infty I\dd t
  =\frac{BL}{R}v_0\tau
  =\boxed{\frac{mv_0}{BL}}.
\end{equation}
结果与 $R$ 无关：电阻越大，瞬时电流越小，但衰减时间按同样比例变长。

图~\ref{ch24:fig:review-rodresistordiagram-1}中，棒横跨上下导轨，左侧R闭合回路；速度减小导致电流和阻力同时减小。
\par\noindent\begin{minipage}{\linewidth}
\rodResistorDiagram
\figcap{导体棒经电阻耗散动能}\label{ch24:fig:review-rodresistordiagram-1}
\end{minipage}\par\medskip
\end{xesolution}



\begin{origexamplebox}{【小棒加小棒】}\label{ch24:ex:s40-9-4}
两根导体棒横跨在宽度为 $L$ 的光滑导轨上，导轨间有垂直纸面向里的匀强磁场 $B$。两棒质量均为 $m$、电阻均为 $R$，
导轨电阻不计。右棒初速为 $v_0$ 向右，左棒初始静止。

(1) 求任意时刻左右两棒的速度；

(2) 求两棒从开始运动到 $t\to\infty$ 的相对位移。

% Pure source-question geometry: no derived forces, currents or solution curves.
%
\begin{center}\begin{tikzpicture}[scale=0.78,>=Latex]
  \draw[Ink,very thick] (-3,1.4)--(3,1.4); \draw[Ink,very thick] (-3,-1.4)--(3,-1.4);
  \draw[DeepBlue,very thick] (-1.5,-1.55)--(-1.5,1.55); \draw[DeepBlue,very thick] (1.5,-1.55)--(1.5,1.55);
  \node[above] at (-1.5,1.62) {$m,R$}; \node[above] at (1.5,1.62) {$m,R$};
  \draw[vector] (1.7,0)--(2.9,0) node[right] {$v_0$};
  \foreach \x in {-2.4,-0.5,0.5,2.4}{\foreach \y in {-0.7,0.7}{\node[TextGray] at (\x,\y) {$\times$};}}
  \draw[<->] (-3.25,-1.4)--node[left] {$L$}(-3.25,1.4);
  \end{tikzpicture}\end{center}
\figcap{两根相同导体棒，右棒具有初速度。}\label{ch24:fig:em24-question-4}
\end{origexamplebox}

\begin{xesolution}
记左、右棒速度为 $v_1,v_2$，相对速度 $u=v_2-v_1$。回路电动势只与面积变化率有关：
\begin{equation}
  \mathcal E=BLu,
  \qquad I=\frac{BL}{2R}u.
\end{equation}
左棒受力向右，右棒受力向左，故
\begin{equation}
  m\dot v_1=BLI,
  \qquad m\dot v_2=-BLI.
\end{equation}
因而质心速度不变，$V_{\rm cm}=v_0/2$，而
\begin{equation}
  \dot u=-\frac{B^2L^2}{mR}u.
\end{equation}
令 $\tau=mR/(B^2L^2)$，得
\begin{equation}
  \boxed{v_1=\frac{v_0}{2}\left(1-\e^{-t/\tau}\right)},
  \qquad
  \boxed{v_2=\frac{v_0}{2}\left(1+\e^{-t/\tau}\right)}.
\end{equation}
两棒的相对位移为
\begin{equation}
  \boxed{\Delta x_{\rm rel}=\int_0^\infty u\dd t
  =v_0\tau=\frac{mRv_0}{B^2L^2}}.
\end{equation}

图~\ref{ch24:fig:review-tworoddiagram-1}中，两个运动边界的电动势相减，驱动量取决于 $v_2-v_1$，不是任一棒的绝对速度。
\par\noindent\begin{minipage}{\linewidth}
\twoRodDiagram
\figcap{同一导轨上的两根导体棒}\label{ch24:fig:review-tworoddiagram-1}
\end{minipage}\par\medskip
\end{xesolution}



\begin{origexamplebox}{【小棒再加小棒】}\label{ch24:ex:s40-9-5}
两根导体棒分别横跨在左右两组导轨上，左侧导轨宽度为 $L_1$、磁场为 $B_1$，右侧导轨宽度为 $L_2$、磁场为 $B_2$，
中间用很长的过渡导轨连接。两棒质量分别为 $m_1,m_2$，回路总电阻为 $R$。右棒初速为 $v_0$ 向右，左棒初始静止。

(1) 求任意时刻两棒的速度；

(2) 找出运动中的一个守恒量。

% Pure source-question geometry: no derived forces, currents or solution curves.
%
\begin{center}\begin{tikzpicture}[scale=.82,>=Latex]
 \draw[Ink,thick](-5,1)--(-1,1)--(1,1.8)--(5,1.8);
 \draw[Ink,thick](-5,-1)--(-1,-1)--(1,-1.8)--(5,-1.8);
 \draw[DeepBlue,very thick](-4,-1.15)--(-4,1.15)node[above]{$m_1$};
 \draw[DeepBlue,very thick](4,-1.95)--(4,1.95)node[above]{$m_2$};
 \draw[vector](4.25,0)--(5.4,0)node[right]{$v_0$};
 \node at(-2,-.5){$B_1,L_1$};\node at(2,-.6){$B_2,L_2$};
 \foreach \x in {-3,-2,2,3}{\node[TextGray]at(\x,.65){$\times$};}
 \end{tikzpicture}\end{center}
\figcap{不同棒长与场强的两导轨段，由固定导线相连。}\label{ch24:fig:em24-question-5}
\end{origexamplebox}

\begin{xesolution}
记 $k_1=B_1L_1$、$k_2=B_2L_2$。按图中回路方向取正，瞬时感应电动势为
\begin{equation}
  \mathcal E=k_2v_2-k_1v_1,
  \qquad I=\frac{k_2v_2-k_1v_1}{R}.
\end{equation}
两棒的方程是
\begin{equation}
  m_1\dot v_1=k_1I,
  \qquad m_2\dot v_2=-k_2I.
\end{equation}
于是
\begin{equation}
  \frac{\dd}{\dd t}\left(\frac{m_1v_1}{k_1}+\frac{m_2v_2}{k_2}\right)=0,
\end{equation}
所求守恒量可取为
\begin{equation}
  \boxed{\frac{m_1v_1}{B_1L_1}+\frac{m_2v_2}{B_2L_2}
  =\frac{m_2v_0}{B_2L_2}}.
\end{equation}
再令 $y=k_2v_2-k_1v_1$，有
\begin{equation}
  \dot y=-\frac1R\left(\frac{k_1^2}{m_1}+\frac{k_2^2}{m_2}\right)y,
  \qquad
  y(0)=k_2v_0.
\end{equation}
定义
\begin{equation}
  D=k_1^2m_2+k_2^2m_1,
  \qquad \tau=\frac{Rm_1m_2}{D},
\end{equation}
积分得
\begin{align}
  \boxed{v_1(t)}&=\boxed{\frac{k_1k_2m_2}{D}v_0\left(1-\e^{-t/\tau}\right)},\\
  \boxed{v_2(t)}&=\boxed{\frac{k_1^2m_2+k_2^2m_1\e^{-t/\tau}}{D}v_0}.
\end{align}
当 $t\to\infty$ 时 $k_2v_2=k_1v_1$，也就是总动生电动势降为零。

图~\ref{ch24:fig:review-unequalroddiagram-1}中，上下两根过渡导线分别连通；右棒正电流向上、左棒向下，故两棒的安培力方向相反。
\par\noindent\begin{minipage}{\linewidth}
\unequalRodDiagram
\figcap{不等宽导轨之间的闭合连接}\label{ch24:fig:review-unequalroddiagram-1}
\end{minipage}\par\medskip
\end{xesolution}



\begin{origexamplebox}{【小棒加电容】}\label{ch24:ex:s40-9-6}
一根质量为 $m$ 的导体棒横跨在宽度为 $L$ 的光滑导轨上，导轨间有垂直纸面向里的匀强磁场 $B$。回路中串联电容 $C$，
所有电阻均忽略。棒初速为 $v_0$ 向右，电容初始不带电。描述 $t=0$ 后导体棒的运动。

% Pure source-question geometry: no derived forces, currents or solution curves.
%
\begin{center}\begin{tikzpicture}[scale=0.78,>=Latex]
  \draw[Ink,very thick] (-3,1.4)--(3,1.4); \draw[Ink,very thick] (-3,-1.4)--(3,-1.4);
  \draw[Ink,very thick] (-2.7,-1.4)--(-2.7,-0.25); \draw[Ink,very thick] (-3.05,-0.25)--(-2.35,-0.25);
  \draw[Ink,very thick] (-3.05,0.25)--(-2.35,0.25); \draw[Ink,very thick] (-2.7,0.25)--(-2.7,1.4);
  \node[left] at (-2.7,0) {$C$};
  \draw[DeepBlue,very thick] (0.8,-1.55)--(0.8,1.55) node[above] {$m$};
  \draw[vector] (1.0,0)--(2.4,0) node[right] {$v_0$};
  \foreach \x in {-1.7,-0.7,0.2,1.5,2.4}{\foreach \y in {-0.7,0.7}{\node[TextGray] at (\x,\y) {$\times$};}}
  \end{tikzpicture}\end{center}
\figcap{导体棒与初始未充电电容器相连。}\label{ch24:fig:em24-question-6}
\end{origexamplebox}

\begin{xesolution}
零电阻与“初始电容不带电、棒却有非零速度”不能在普通时间尺度上同时满足；闭合回路后会出现一个极短的充电暂态。
设暂态后电容带电量为 $q$、棒速为 $v$，则回路方程和棒的冲量方程分别为
\begin{equation}
  \frac qC=BLv,
  \qquad m(v-v_0)=-BLq.
\end{equation}
解得
\begin{equation}
  \boxed{v=\frac{m}{m+CB^2L^2}v_0},
  \qquad
  \boxed{q=\frac{mCBL}{m+CB^2L^2}v_0}.
\end{equation}
上述状态是忽略自感、先保留微小电阻 $R_s>0$ 并等待暂态结束后的终态，随后电流为零，棒匀速运动。
为明确过程，令 $k_m=BL$，则
\[
 R_s\dot q+q/C=k_mv,\qquad mv+k_mq=mv_0.
\]
消去 $v$ 得 $R_s\dot q+(1/C+k_m^2/m)q=k_mv_0$。若 $q(0)=0$，则
\[
 q(t)=q_\infty(1-\e^{-t/\tau_s}),\quad
 \tau_s=\frac{R_smC}{m+Ck_m^2},\quad
 v(t)=v_0-\frac{k_m}{m}q(t).
\]
因此先令 $t\gg\tau_s$ 再取 $R_s\to0^+$，终态不变，暂态持续时间趋零。
能量检验给出
\[
 Q=\frac12mv_0^2-\frac12mv_\infty^2-\frac{q_\infty^2}{2C}
 =\frac{mCk_m^2v_0^2}{2(m+Ck_m^2)}>0.
\]
这个耗散不随上述顺序下的 $R_s\to0^+$ 而消失；直接令电阻、自感都严格为零会得到奇异初值问题。
若保留电感却使电阻严格为零，系统一般振荡，不能说它必然进入同一恒速终态。

图~\ref{ch24:fig:review-rodcapacitordiagram-1}中，电容最终电压等于棒的动生电动势；暂态需按文中微小电阻模型理解。
\par\noindent\begin{minipage}{\linewidth}
\rodCapacitorDiagram
\figcap{机械等效电容与真实电容的电荷重分配}\label{ch24:fig:review-rodcapacitordiagram-1}
\end{minipage}\par\medskip
\end{xesolution}



\begin{origexamplebox}{【小棒加电感】}\label{ch24:ex:s40-9-7}
一根质量为 $m$ 的导体棒横跨在宽度为 $\ell$ 的光滑导轨上，导轨间有垂直纸面向里的匀强磁场 $B$。回路中串联自感为 $\mathcal L$
的理想电感，忽略回路其余部分的自感与电阻，$\mathcal L$ 不随棒的位置变化。棒初速为 $v_0$ 向右，电感初始无电流。描述 $t=0$ 后导体棒的运动。

% Pure source-question geometry: no derived forces, currents or solution curves.
%
\begin{center}\begin{tikzpicture}[scale=0.78,>=Latex]
  \draw[Ink,very thick] (-3,1.4)--(3,1.4); \draw[Ink,very thick] (-3,-1.4)--(3,-1.4);
  \draw[Ink,very thick] (-2.7,-1.4)--(-2.7,-0.6);
  \draw[Ink,thick] (-2.7,-0.6) decorate[decoration={coil,aspect=0.65,segment length=5pt,amplitude=5pt}] {-- (-2.7,0.6)};
  \draw[Ink,very thick] (-2.7,0.6)--(-2.7,1.4); \node[left] at (-2.9,0) {$\mathcal L$};
  \draw[DeepBlue,very thick] (0.8,-1.55)--(0.8,1.55) node[above] {$m$};
  \draw[vector] (1.0,0)--(2.4,0) node[right] {$v_0$};
  \foreach \x in {-1.7,-0.7,0.2,1.5,2.4}{\foreach \y in {-0.7,0.7}{\node[TextGray] at (\x,\y) {$\times$};}}
  \end{tikzpicture}\end{center}
\figcap{导体棒与电感相连，电感初始电流为零。}\label{ch24:fig:em24-question-7}
\end{origexamplebox}

\begin{xesolution}
取感应电流 $I$ 的正方向使它对棒产生向左的力。电路与力学方程为
\begin{equation}
  \mathcal L\dot I=B\ell v,
  \qquad m\dot v=-B\ell I.
\end{equation}
对第二式再求导并代入第一式，得
\begin{equation}
  \ddot v+\Omega^2v=0,
  \qquad
  \Omega=\frac{B\ell}{\sqrt{m\mathcal L}}.
\end{equation}
用 $v(0)=v_0$、$I(0)=0$ 得
\begin{equation}
  \boxed{v(t)=v_0\cos\Omega t},
  \qquad
  \boxed{I(t)=v_0\sqrt{\frac m{\mathcal L}}\sin\Omega t},
\end{equation}
\begin{equation}
  \boxed{x(t)-x(0)=\frac{v_0}{\Omega}\sin\Omega t}.
\end{equation}
棒的动能和电感磁能周期性交换：
\begin{equation}
  \frac12mv^2+\frac12\mathcal LI^2=\frac12mv_0^2.
\end{equation}

图~\ref{ch24:fig:review-rodinductordiagram-1}中，线圈接入回路后，电感电流和棒速度连续变化，动能与磁能相互转化。
\par\noindent\begin{minipage}{\linewidth}
\rodInductorDiagram
\figcap{导体棒与电感的能量交换}\label{ch24:fig:review-rodinductordiagram-1}
\end{minipage}\par\medskip
\end{xesolution}



\begin{origexamplebox}{【小棒加电源】}\label{ch24:ex:s40-9-8}
一根质量为 $m$ 的导体棒横跨在宽度为 $\ell$ 的光滑导轨上，导轨间有垂直纸面向里的匀强磁场 $B$。回路总电阻为 $R$，
并接入恒定电源 $\mathcal E$，电源极性使静止棒受到向右的安培力。棒初速为 $v_0$ 向右。描述 $t=0$ 后导体棒的运动。

% Pure source-question geometry: no derived forces, currents or solution curves.
%
\begin{center}\begin{tikzpicture}[scale=.8,>=Latex]
 \draw[Ink,thick](-3,1.5)--(3,1.5) (-3,-1.5)--(3,-1.5);
 \draw[Ink,thick](-2.7,-1.5)--(-2.7,-.8);
 \draw[Ink,thick](-2.9,-.8)rectangle(-2.5,-.15);
 \draw[Ink,thick](-2.7,-.15)--(-2.7,.25);
 \draw[Ink,thick](-2.95,.25)--(-2.45,.25) (-3.1,.65)--(-2.3,.65);
 \draw[Ink,thick](-2.7,.65)--(-2.7,1.5);
 \node[left]at(-3,-.5){$R$};\node[left]at(-3.15,.45){$\mathcal E$};
 \node[right]at(-2.35,.25){$-$};\node[right]at(-2.25,.65){$+$};
 \draw[DeepBlue,very thick](.8,-1.65)--(.8,1.65)node[above]{$m$};
 \draw[vector](1.2,0)--(2.6,0)node[right]{$v_0$};
 \foreach \x in {-1.6,-.6,1.8}{\foreach \y in {-.9,.9}{\node[TextGray]at(\x,\y){$\times$};}}
\end{tikzpicture}\end{center}
\figcap{电源、定值电阻和运动导体棒的实际连接。}\label{ch24:fig:em24-question-8}
\end{origexamplebox}

\begin{xesolution}
取棒中向下为电流正方向。按右手定则，正电流在向纸内的磁场中产生向右的安培力；
实际电流的正负还取决于棒速，不能仅由电源极性判断。按此约定，电流和运动方程为
\begin{equation}
  I=\frac{\mathcal E-B\ell v}{R},
  \qquad
  m\dot v=B\ell I.
\end{equation}
运动电动势 $B\ell v$ 是电源的反电动势。定义
\begin{equation}
  v_\infty=\frac{\mathcal E}{B\ell},
  \qquad \tau=\frac{mR}{B^2\ell^2},
\end{equation}
得
\begin{equation}
  \boxed{v(t)=v_\infty+(v_0-v_\infty)\e^{-t/\tau}}.
\end{equation}
若 $v_0<v_\infty$，电源使棒加速；若 $v_0>v_\infty$，反电动势大于电源，电流反向，棒减速；两种情况都趋向同一终端速度。

图~\ref{ch24:fig:review-rodsourcediagram-1}中，按所标电源极性，正电流在棒中向下；当动生电动势大于电源时，电流可反向。
\par\noindent\begin{minipage}{\linewidth}
\rodSourceDiagram
\figcap{电源驱动的导体棒}\label{ch24:fig:review-rodsourcediagram-1}
\end{minipage}\par\medskip
\end{xesolution}



\begin{origexamplebox}{【小棒大结局】}\label{ch24:ex:s40-9-9}
一根质量为 $m$、电阻为 $R$ 的导体棒横跨在宽度为 $L$ 的光滑导轨上，导轨间有垂直纸面向里的匀强磁场 $B$。

(1) 若 $R=0$，推导棒两端电压 $U$ 与流过棒的电流 $I$ 的关系，辨认其等效元件与参数；

(2) 若 $R\ne0$，再求 $U$--$I$ 关系与等效元件；

(3) 用该结论重新理解前面的电阻、电容、电感和电源问题。

% Pure source-question geometry: no derived forces, currents or solution curves.
\begin{center}\begin{tikzpicture}[scale=.78,>=Latex]
 \draw[thick](-3,1.4)--(3,1.4) (-3,-1.4)--(3,-1.4);
 \draw[DeepBlue,very thick](.8,-1.55)--(.8,1.55)node[above]{$m,R$};
 \draw[vector](1.1,0)--(2.4,0)node[right]{$v$};
 \draw[vector](-2.7,1.4)--(-1.7,1.4)node[above]{$I$};
 \node[left]at(-3,1.4){$+$};\node[left]at(-3,-1.4){$-$};
 \draw[<->](-3.7,-1.4)--(-3.7,1.4)node[midway,left]{$U$};
 \draw[<->](3.4,-1.4)--(3.4,1.4)node[midway,right]{$L$};
 \foreach \x in {-1.6,-.5,1.8}{\foreach \y in {-.8,.8}{\node at(\x,\y){$\otimes$};}}
 \node at(-.7,0){$B$};
\end{tikzpicture}\end{center}
\figcap{可运动导体棒的电学端口；箭头和正负号定义关联参考方向。}\label{ch24:fig:em24-question-9}
\end{origexamplebox}

\begin{xesolution}
把棒的动生电动势记为 $U_m=BLv$，取正电流的安培力使棒向右，则
\begin{equation}
  m\dot v=BLI.
\end{equation}
因而
\begin{equation}
  \dot U_m=BL\dot v=\frac{B^2L^2}{m}I,
  \qquad
  \boxed{I=C_m\dot U_m},
  \qquad
  \boxed{C_m=\frac{m}{B^2L^2}}.
\end{equation}
当 $R=0$ 时，$U=U_m$，所以这根可运动导体棒在电学端口上等效为电容 $C_m$。其“电场能”正好就是棒的动能：
\begin{equation}
  \frac12C_mU_m^2=\frac12\frac{m}{B^2L^2}(BLv)^2=\frac12mv^2.
\end{equation}
当棒有电阻 $R$ 时，端口方程为
\begin{equation}
  \boxed{U=RI+U_m},
  \qquad I=C_m\dot U_m,
\end{equation}
因而等效为电阻 $R$ 与电容 $C_m$ 串联，电容初电压为 $U_m(0)=BLv_0$。

前面诸问题因此统一为基本电路：例~\ref{ch24:ex:s40-9-3} 是 $RC_m$ 放电；例~\ref{ch24:ex:s40-9-6} 是机械等效电容与真实电容的瞬态电荷重分配；
例~\ref{ch24:ex:s40-9-7} 是 $\mathcal LC_m$ 振荡；例~\ref{ch24:ex:s40-9-8} 是电源对 $RC_m$ 的充电。

图~\ref{ch24:fig:review-equivalentcircuitdiagram-1}中，机械端口电压 $U_m=BLv$ 与电流关系相当于电容，图中R与该电容串联。
\par\noindent\begin{minipage}{\linewidth}
\equivalentCircuitDiagram
\figcap{运动导体棒的电学端口等效}\label{ch24:fig:review-equivalentcircuitdiagram-1}
\end{minipage}\par\medskip
\end{xesolution}



\section{感生电场、加速与电压测量}

\subsection{感生电场}


导体固定而磁场随时间变化时，仍可有非零电动势。此时驱动电荷的是空间中的涡旋电场：
\begin{equation}
  \boxed{\oint_C\vect E\cdot\dd\vect\ell
  =-\frac{\dd}{\dd t}\int_S\vect B\cdot\dd\vect S},
  \qquad
  \boxed{\nabla\times\vect E=-\frac{\partial\vect B}{\partial t}}.
\end{equation}
感生电场的环量不为零，因而它不能在全空间由单值电势表示。在半径 $R$ 的圆柱区域内有轴向匀强磁场
$B(t)$ 时，由轴对称性得
\begin{equation}
  \boxed{E_\varphi(r)=-\frac r2\dot B\quad(r\le R)},
  \qquad
  \boxed{E_\varphi(r)=-\frac{R^2}{2r}\dot B\quad(r\ge R)}.
\end{equation}
一根径向导线上 $\vect E\perp\dd\vect\ell$，沿它没有感生电动势；但沿圆周或一般非径向路径则可能有非零线积分。

静电场与感生电场的差别可以用场线图像概括。静电场线从正电荷出发、终止于负电荷，在无电荷区域的
闭合环量为零；感生电场线则围绕变磁场区域闭合，没有能够给整个空间定义单值电势的起点和终点。
若在空间中放入导体，导体内自由电荷会重新分布并叠加静电场，所以实际导体中的总电场可能与纯粹由
法拉第定律算出的涡旋场不同；但沿闭合回路的静电场环量仍为零，不会改变总感应电动势。

在轴对称系统中，若半径 $r$ 的圆周内磁场处处相同，则
$2\pi rE_\varphi=-\pi r^2\dot B$，得到内区公式。外区圆周只包围半径 $R$ 的变磁场区域，磁通量为
$\pi R^2B$，因此出现 $E_\varphi\propto1/r$ 的外区结果。电场线是闭合圆，环量非零；它不能像静电场
那样写成全空间单值电势的负梯度。图中沿径向的导线电压为零，并不意味着空间中没有电场，而是该路径
处处与涡旋电场垂直。

外区虽没有磁场随时间变化，仍然可以存在感生电场，这是“局部场方程”与“磁通计算”需要共同理解的
地方。圆周本身位于 $r>R$，但它包围的任一跨越面都穿过内区变磁场，故环量不为零。不能把法拉第
定律误读成“只有 $\dot B\ne0$ 的点才有电场”；更准确地说，$\nabla\times\vect E$ 在变磁场区非零，
这一涡旋结构决定了整个周围空间的电场分布。

图~\ref{ch24:fig:review-vortexfielddiagram-1}中，磁场向纸内增强时，电场逆时针；半径R是内外电场公式的分界，不是电场消失的位置。
\par\noindent\begin{minipage}{\linewidth}
\vortexFieldDiagram
\figcap{柱形变磁场内外的涡旋电场}\label{ch24:fig:review-vortexfielddiagram-1}
\end{minipage}\par\medskip


% EM teaching enrichment 20260927
\paragraph{有场区外也可以存在涡旋电场}
轴向磁场只分布在半径 $a$ 的圆柱区域内，内部 $B(t)$ 均匀。取逆时针环路为正，$B$ 向纸外为正，由对称性
\[
 E_\varphi(r)=\begin{cases}
 -\dfrac r2\dot B,&r<a,\\[3pt]
 -\dfrac{a^2}{2r}\dot B,&r>a.
 \end{cases}
\]
两式在 $a$ 处连续。外部观察点虽有 $B=0$，其环路仍包围变化磁通，所以 $E$ 不为零。
此时把静止电荷放在圆柱外也会受到电场力；感生电场不要求那里先有导线。若加上一条闭合金属环，环内电流由环量与电路响应共同决定。


\subsection{电子感应加速器}


电子感应加速器（betatron）把涡旋电场用于加速电子。设电子沿半径 $r$ 的固定圆轨道运动，轨道以内
磁场的面积平均值为 $\overline B(t)$。法拉第定律给出圆周切向电场
\begin{equation}
  2\pi rE_\varphi=-\pi r^2\frac{\dd\overline B}{\dd t},
  \qquad
  E_\varphi=-\frac r2\frac{\dd\overline B}{\dd t}.
\end{equation}
以电子加速的绕行方向为正，电子电荷为 $-e$。切向方程为
$\dd p/\dd t=er\dot{\overline B}/2$。从注入时刻 $t_0$ 积分，先保留初值：
\begin{equation}
 p(t)-p(t_0)=\frac{er}{2}[\overline B(t)-\overline B(t_0)].
\end{equation}
若注入条件满足 $p(t_0)=er\overline B(t_0)/2$，才可简化为 $p=er\overline B/2$，
在非相对论情形下就是
\begin{equation}
  v=\frac{er\overline B}{2m}.
\end{equation}
另一方面，轨道处的轴向局部磁场 $B_{\rm orb}$ 必须提供向心力，
$mv^2/r=evB_{\rm orb}$，即 $v=erB_{\rm orb}/m$。比较两式可得保持圆轨道半径的条件
\begin{equation}
  \boxed{B_{\rm orb}=\frac12\overline B}.
\end{equation}
更一般地，固定半径要求 $\overline B-2B_{\rm orb}$ 为常量；常用的注入匹配使常量为零，
于是轨道处的局部磁场始终等于轨道面积平均磁场的一半。这一“二比一条件”同时保证涡旋
电场负责加速、局部磁场负责弯转。

这里要区分三个量：轨道半径 $r$，轨道处的轴向局部磁场 $B_{\rm orb}$，以及整个轨道圆面积上的平均磁场
$\overline B$。法拉第定律关心圆内总磁通，所以出现 $\overline B$；圆周运动的瞬时向心力只由电子
所在位置的局部场决定，所以出现 $B_{\rm orb}$。这里下标orb指轨道，不是径向分量。若把两者混成同一个均匀场，就会得到
$B=\overline B/2$ 的自相矛盾等式。实际 betatron 必须专门设计磁极形状，使空间分布满足平均场与轨道
场的二比一关系，同时还要提供径向、竖直方向的聚焦稳定性。

\historyfigure[0.72\linewidth]{betatron-built-in-v1.png}
{电子感应加速器及其发明者克斯特。}\label{ch24:fig:s40-9-6}

电子感应加速器曾可把电子加速到数百 MeV；能量继续提高时，同步辐射损耗快速增加，其他大型加速器
更有优势。此类装置后来大量用于产生 X 射线和放射治疗。它展示了感生电场并非只在导线里“等效存在”，
而是在真空轨道上也能直接对带电粒子做功。

上面的 $mv^2/r=evB_{\rm orb}$ 是非相对论写法。电子能量升高后，应把机械动量写成 $p=\gamma mv$，轨道
关系改为 $p=erB_{\rm orb}$；法拉第电场仍沿切向增加电子动量。随着速度接近 $c$，继续输入的能量主要增加
$\gamma$ 而不是显著增加速度，同时圆周加速运动带来的辐射损耗迅速增强。这些效应限制了感应加速器
可达到的能量，也说明“速度越来越大”不能在高能区继续使用牛顿形式外推。


% EM teaching enrichment 20260927
\paragraph{检验一个候选磁场分布}
设设计轨道半径为 $r_0$，轨道内轴向磁场取
$B_z(r,t)=B_c(t)[1-\alpha r^2/r_0^2]$。其面积平均值为
\[
 \overline B=\frac2{r_0^2}\int_0^{r_0}B_z(r,t)r\dd r
 =B_c(1-\alpha/2),
\]
而轨道处 $B_{\rm orb}=B_c(1-\alpha)$。采用零积分常量的注入匹配，要求 $\overline B=2B_{\rm orb}$，得 $\alpha=2/3$。
这个计算只核对固定半径的加速条件，未核对径向、轴向稳定性；满足二比一关系不能自动保证粒子束聚焦。


% EM structure and teaching revision 20260928
\subsection{非保守电场下的电压表读数}
静电场中，两点电势差与路径无关；变化磁场产生的感生电场满足 $\oint\boldsymbol E\cdot\dd\boldsymbol l=-\dot\Phi_B$，其沿两条不同连接路径的线积分之差一般不为零。
设路径 $\Gamma_1,\Gamma_2$ 都从 $A$ 到 $B$，按 $\Gamma_1$ 去、沿 $\Gamma_2$ 返回组成闭合回路，则
\[
 \int_{\Gamma_1}\boldsymbol E\cdot\dd\boldsymbol l-
 \int_{\Gamma_2}\boldsymbol E\cdot\dd\boldsymbol l=-\frac{\dd\Phi_B}{\dd t}.
\]
因此读电压表时，除了接点 $A,B$，还要给出引线围成的路径以及是否穿过变化磁通区域。把引线绞合或贴近布置，是减小附加回路面积与感应误差的一种方式。
计算可把各段电阻上的电压降和感应电动势按指定回路统一列入方程。应避免给整个空间强行指定一个能代表全部电场的单值静电势；只有感应贡献可忽略或已明确划入集中电动势元件时，普通节点电势法才可直接使用。
测量模型还假定电压表内阻足够大、对原回路电流影响可忽略；若不满足此条件，表支路本身也须参与电路求解。


\subsection{专题例题与方法应用}

\begin{origexamplebox}{【小环的响应】}\label{ch24:ex:s40-9-10}
半径为 $R$ 的固定绝缘大圆环与匀强磁场垂直，一个质量为 $m$、电荷为 $+q$ 的小环套在大环上，可无摩擦滑动。规定绕“垂直纸面
向里”的轴按右手螺旋法则为角坐标 $\theta$ 正向，初始 $\theta=0$ 且静止。

(1) $B(t)=\alpha t$；

(2) $B(t)=B_0\sin\omega t$；

(3) $B(t)=B_0\cos\omega t$。

分别求小环的角速度、位置以及大环对小环的径向支持力。
\par\noindent\begin{minipage}{\linewidth}
\chargedRingDiagram
\figcap{小环的角向和径向符号约定}\label{ch24:fig:review-chargedringdiagram-1}
\end{minipage}\par\medskip
\end{origexamplebox}

\begin{xesolution}
半径 $R$ 处的感生电场为
\begin{equation}
  E_\theta=-\frac R2\dot B.
\end{equation}
切向方程 $mR\ddot\theta=qE_\theta$ 给出一个对任意 $B(t)$ 都成立的通式：
\begin{equation}
  \boxed{\dot\theta(t)=-\frac q{2m}\,[B(t)-B(0)]},
  \qquad
  \boxed{\theta(t)=-\frac q{2m}\int_0^t[B(t')-B(0)]\dd t'}.
\end{equation}
取径向向外为正。磁力的径向分量为 $qR\dot\theta B$，径向加速度为 $-R\dot\theta^2$，故大环的有符号支持力为
\begin{equation}
  \boxed{N_r=-mR\dot\theta^2-qR\dot\theta B}.
\end{equation}
$N_r>0$ 表示支持力向外，$N_r<0$ 表示向内。

三种情况依次为
\begin{align}
  B=\alpha t:\quad
  &\boxed{\dot\theta=-\frac{q\alpha}{2m}t},
  &\boxed{\theta=-\frac{q\alpha}{4m}t^2};\\
  B=B_0\sin\omega t:\quad
  &\boxed{\dot\theta=-\frac{qB_0}{2m}\sin\omega t},
  &\boxed{\theta=\frac{qB_0}{2m\omega}(\cos\omega t-1)};\\
  B=B_0\cos\omega t:\quad
  &\boxed{\dot\theta=\frac{qB_0}{2m}(1-\cos\omega t)},
  &\boxed{\theta=\frac{qB_0}{2m}\left(t-\frac{\sin\omega t}{\omega}\right)}.
\end{align}
对第 (1) 问，将 $\dot\theta=-q\alpha t/(2m)$ 代入，向心项为
$-q^2R\alpha^2t^2/(4m)$，磁力移项为 $q^2R\alpha^2t^2/(2m)$，相加得
\begin{equation}
  \boxed{N_r=\frac{q^2R\alpha^2t^2}{4m}},
\end{equation}
方向向外：此时磁力向内且大于所需向心力，大环必须向外拉住小环。
对第 (2) 问同理得
\[
 N_r=-\frac{q^2RB_0^2}{4m}\sin^2\omega t
      +\frac{q^2RB_0^2}{2m}\sin^2\omega t
 =\boxed{\frac{q^2RB_0^2}{4m}\sin^2\omega t}.
\]
它向外或为零。对第 (3) 问，令 $c=\cos\omega t$，得到
\[
 N_r=-\frac{q^2RB_0^2}{4m}[(1-c)^2+2c(1-c)]
 =\boxed{-\frac{q^2RB_0^2}{4m}\sin^2\omega t}.
\]
它向内或为零。两者不同的原因是初始磁场分别为0与 $B_0$，角动量的积分常数不同，
不能只凭某时刻 $B$ 的大小判断速度。三问均满足初始速度为零；$N_r$ 的量纲为力。

图~\ref{ch24:fig:review-chargedringdiagram-1}中，顺时针为角坐标正向，径向向外为正；两个力的代数值按同一径向约定相加。

\end{xesolution}



\begin{origexamplebox}{【貌似简谐振动】}\label{ch24:ex:s40-9-11}
半径为 $R$ 的圆柱形区域内有垂直纸面向里的匀强磁场
$B=B_0\sin\omega t$。在距离圆心 $h$ 的一条弦上有光滑直轨，带电质点 $m,+q$ 被约束在轨道中，初始位于弦中点且静止。

(1) 求粒子的运动规律和任意时刻的位置；

(2) 求保证粒子不离开磁场区域的参数条件。

% Pure source-question geometry: no derived forces, currents or solution curves.
%
\begin{center}\begin{tikzpicture}[scale=0.8,>=Latex]
  \draw[Ink,thick] (0,0) circle (2.5);
  \draw[Ink,very thick] ({-sqrt(2.5*2.5-1.25*1.25)},-1.25)--({sqrt(2.5*2.5-1.25*1.25)},-1.25);
  \fill[DeepBlue] (0,-1.25) circle (3pt) node[below] {$m,+q$};
  \draw[<->] (0,0)--node[right] {$h$}(0,-1.25);
  \draw[axis] (-2.2,-1.75)--(2.4,-1.75) node[right] {$x$};
  \node[TextGray] at (-1.4,1.0) {$\times$}; \node[TextGray] at (0,1.0) {$\times$}; \node[TextGray] at (1.4,1.0) {$\times$};
  \node at (0,1.65) {$B_0\sin\omega t$};
  \end{tikzpicture}\end{center}
\figcap{正电荷位于弦中点，圆形区域内磁场按给定规律变化。}\label{ch24:fig:em24-question-11}
\end{origexamplebox}

\begin{xesolution}
取轨道方向向右为 $x$ 轴，原点在弦中点。对于圆柱内的匀强时变磁场，感生电场可写为
\begin{equation}
  \vect E=-\frac12\dot B_z\,\uvec z\times\vect r.
\end{equation}
磁场指向纸内，在圆心下方距离 $h$ 的弦上，电场沿轨道的分量与 $x$ 无关：
\begin{equation}
  E_x=\frac h2\dot B=\frac{hB_0\omega}{2}\cos\omega t.
\end{equation}
磁力与轨道方向垂直，由轨道的支持力平衡，因而只有电场力改变 $x$ 向运动：
\begin{equation}
  m\ddot x=\frac{qhB_0\omega}{2}\cos\omega t.
\end{equation}
由 $x(0)=0,\dot x(0)=0$ 得
\begin{equation}
  \boxed{\dot x(t)=\frac{qhB_0}{2m}\sin\omega t},
  \qquad
  \boxed{x(t)=\frac{qhB_0}{2m\omega}(1-\cos\omega t)}.
\end{equation}
这是绕平衡位置 $x_c=qhB_0/(2m\omega)$ 的简谐形运动，但驱动来自时变磁场，并非恢复力。弦的半长为
$s=\sqrt{R^2-h^2}$，而 $|x|$ 的最大值为 $|q|hB_0/(m\omega)$，故条件为
\begin{equation}
  \boxed{\frac{|q|hB_0}{m\omega}\le\sqrt{R^2-h^2}}.
\end{equation}

图~\ref{ch24:fig:review-chordparticlediagram-1}中，轨道在圆心下方h处，端点由同一个圆与直线的交点确定；粒子位移不能越过弦端。
\par\noindent\begin{minipage}{\linewidth}
\chordParticleDiagram
\figcap{弦上粒子的约束与磁场边界}\label{ch24:fig:review-chordparticlediagram-1}
\end{minipage}\par\medskip
\end{xesolution}



\begin{origexamplebox}{【感生电动势】}\label{ch24:ex:s40-9-12}
半径为 $R$ 的圆柱区域内有垂直纸面向里的匀强磁场
$B=B_0+kt$。距圆心 $h$ 处有一条直导线 $AD$，与圆的两个交点为 $A,C$，$D$ 位于 $C$ 右侧。记 $CD=\ell$，求沿 $A\to C$、$C\to D$ 两段的感生电动势，并讨论 $D$ 趋于无穷远的特例。

% Pure source-question geometry: no derived forces, currents or solution curves.
%
\begin{center}\begin{tikzpicture}[scale=0.82,>=Latex]
  \draw[Ink,thick] (0,0) circle (2.5);
  \draw[Ink,very thick] (-3.2,-1.2)--(4.3,-1.2);
  \coordinate (A) at ({-sqrt(2.5*2.5-1.2*1.2)},-1.2); \coordinate (C) at ({sqrt(2.5*2.5-1.2*1.2)},-1.2); \coordinate (D) at (4.1,-1.2);
  \fill[DeepBlue] (A) circle (2.5pt) node[below] {$A$}; \fill[DeepBlue] (C) circle (2.5pt) node[below] {$C$}; \fill[Ink] (D) circle (2.5pt) node[below] {$D$};
  \draw[<->] ({sqrt(2.5*2.5-1.2*1.2)},-2)--(4.1,-2) node[midway,below] {$\ell$};
  \draw[<->] (0,0)--node[right] {$h$}(0,-1.2);
  \node[TextGray] at (-1.1,.2) {$\dot B=k$};
\end{tikzpicture}\end{center}
\figcap{弦、观察端点和外延直线；感生电场分布待求。}\label{ch24:fig:em24-question-12}
\end{origexamplebox}

\begin{xesolution}
先取 $0<h<R$，令 $s=\sqrt{R^2-h^2}$，以垂足为坐标原点，导线向右为积分正向；则 $x_A=-s$、$x_C=s$、$x_D=s+\ell$。磁场向里增大时，感生电场从读者视角看为逆时针，因而在圆心下方的导线上其分量向右。

圆内电场 $E_\varphi=kr/2$。在水平线 $y=-h$ 上，其 $x$ 分量为常量 $kh/2$，故
\begin{equation}
  \boxed{\mathcal E_{AC}=\int_{-s}^{s}\frac{kh}{2}\dd x
  =kh\sqrt{R^2-h^2}}.
\end{equation}
圆外电场 $E_\varphi=kR^2/(2r)$，而水平分量为
\begin{equation}
  E_x=\frac{kR^2h}{2(x^2+h^2)}.
\end{equation}
因此对于有限端点 $D$，
\begin{align}
  \mathcal E_{CD}
  &=\frac{kR^2h}{2}\int_s^{s+\ell}\frac{\dd x}{x^2+h^2}\\
  &=\boxed{\frac{kR^2}{2}\left[\arctan\frac{s+\ell}{h}-\arctan\frac{s}{h}\right]}.
\end{align}
只有再令 $\ell\to\infty$，才得到
\[
 \boxed{\mathcal E_{C\infty}=\frac{kR^2}{2}\left(\frac\pi2-\arctan\frac{s}{h}\right)
 =\frac{kR^2}{2}\arcsin\frac hR}.
\]
当 $\ell=0$ 时第二段电动势为零；当 $h=0$ 时两段上的电场沿线分量都为零，可由几何直接判断。
若改变积分方向，电动势整体变号。
后一段虽然完全位于无磁场区域，感生电场却并不为零；这正是涡旋电场与“局部磁力”图像的区别。

图~\ref{ch24:fig:review-chordemfdiagram-1}中，A、C为圆与导线的交点，D为有限观察端点；圆外仍需积分涡旋电场。
\par\noindent\begin{minipage}{\linewidth}
\chordEmfDiagram
\figcap{跨越磁场边界的直导线}\label{ch24:fig:review-chordemfdiagram-1}
\end{minipage}\par\medskip
\end{xesolution}



\begin{origexamplebox}{【这是测的啥】}\label{ch24:ex:s40-9-13}
半径为 $R$ 的圆柱区域内有垂直纸面向里的匀强磁场 $B=B_0+kt$。边缘处放置一根不均匀的圆形电阻线，两端点为 $A,C$。
劣弧 $AC$ 的圆心角为 $2\pi/3$、电阻为 $r_0/4$；优弧 $AC$ 的圆心角为 $4\pi/3$、电阻为 $r_0/3$。电压表 $V_1,V_2$ 分别紧贴劣弧和优弧连接 $A,C$；
$V_3$ 的两根导线从 $A,C$ 沿径向通到圆心。求三表示数并解释它们测量的物理量。
\par\noindent\begin{minipage}{\linewidth}
\circularVoltmetersDiagram
\figcap{圆形电阻线与三种测量路径}\label{ch24:fig:review-circularvoltmetersdiagram-1}
\end{minipage}\par\medskip
\end{origexamplebox}

\begin{xesolution}
总感应电动势的大小为
\begin{equation}
  \mathcal E_0=k\pi R^2.
\end{equation}
电阻线中的电流为
\begin{equation}
  I=\frac{\mathcal E_0}{r_0/4+r_0/3}
  =\frac{12\mathcal E_0}{7r_0}.
\end{equation}
$V_1,V_2$ 的导线各自紧贴被测电阻弧，电压表与电阻弧围成的面积可趋于零，所以它们读出的是相应电阻上的欧姆压降：
\begin{equation}
  \boxed{V_1=I\frac{r_0}{4}=\frac37\mathcal E_0},
  \qquad
  \boxed{V_2=I\frac{r_0}{3}=\frac47\mathcal E_0}.
\end{equation}
$V_3$ 的两根引线都沿径向，引线上 $\vect E_{\rm ind}\cdot\dd\vect\ell=0$。劣弧和两根径向线围成的扇形面积为 $\pi R^2/3$，其感应电动势为 $\mathcal E_0/3$。因此
\begin{equation}
  \boxed{V_3=\left|I\frac{r_0}{4}-\frac{\mathcal E_0}{3}\right|
  =\frac{2}{21}\mathcal E_0}.
\end{equation}
$V_3$ 测得的是由该引线路径定义的端点读数，它不等于任意一段电阻压降。时变磁场中 $\oint\vect E\cdot\dd\vect\ell\ne0$，“两点间电压”会依赖电压表引线的路径。

图~\ref{ch24:fig:review-circularvoltmetersdiagram-1}中，V1、V2的引线分别紧贴两段弧，V3沿两条半径；不同闭合路径对应不同包围面积。

\end{xesolution}



\begin{origexamplebox}{【这也可以测】}\label{ch24:ex:s40-9-14}
边长为 $a$ 的正方形导线框放在垂直纸面向里的匀强磁场 $B=B_0+kt$ 中。左、右两边电阻分别为 $R_1,R_2$，上、下两边电阻忽略。
电压表 $V_1,V_2,V_3,V_4$ 分别紧贴上、右、下、左四边连接。求各表示数并解释其物理意义。
\par\noindent\begin{minipage}{\linewidth}
\squareVoltmetersDiagram
\figcap{四只电压表的真实连接点}\label{ch24:fig:review-squarevoltmetersdiagram-1}
\end{minipage}\par\medskip
\end{origexamplebox}

\begin{xesolution}
线框的感应电动势大小为 $\mathcal E_0=ka^2$，电流大小为
\begin{equation}
  I=\frac{ka^2}{R_1+R_2}.
\end{equation}
紧贴某一边的电压表与该边围成的面积趋于零，因此其读数就是该边的欧姆压降。上下两边无电阻，故
\begin{equation}
  \boxed{V_1=0},
  \qquad
  \boxed{V_2=\frac{R_2}{R_1+R_2}ka^2},
\end{equation}
\begin{equation}
  \boxed{V_3=0},
  \qquad
  \boxed{V_4=\frac{R_1}{R_1+R_2}ka^2}.
\end{equation}
虽然 $V_1$ 和 $V_3$ 都连在左右两侧导体之间，但它们的引线路径不同。在涡旋电场中，不能脱离测量路径把“两点电压”当成唯一的标量。

图~\ref{ch24:fig:review-squarevoltmetersdiagram-1}中，所有引线都接到顶点，表与被测边的间距仅为辨认而放大，计算取其所围附加面积趋零。

\end{xesolution}



\begin{origexamplebox}{【奇妙的田字格】}\label{ch24:ex:s40-9-15}
一个外框边长为 $2a$ 的“田”字格电阻网络放在垂直纸面向里的匀强磁场 $B=B_0+kt$ 中。每根导线的单位长度电阻为 $r/a$，因而每段长为 $a$ 的小边电阻为 $r$。求各支路电流。

% Pure source-question geometry: no derived forces, currents or solution curves.
%
\begin{center}\begin{tikzpicture}[scale=.85,>=Latex]
 \draw[Ink,thick](-2,-2)grid[step=2](2,2);
 \node at(-1,1){$\otimes$};\node at(1,1){$\otimes$};\node at(-1,-1){$\otimes$};\node at(1,-1){$\otimes$};
 \draw[<->](-2,-2.5)--(2,-2.5)node[midway,below]{$2a$};\node[right]at(2.6,.2){$\dot B=k$};
\end{tikzpicture}\end{center}
\figcap{四个方格中的垂直磁场按相同变化率改变。}\label{ch24:fig:em24-question-15}
\end{origexamplebox}

\begin{xesolution}
四个边长为 $a$ 的小方格所围面积相同，每个网孔的感应电动势都是
\begin{equation}
  \mathcal E_c=ka^2,
\end{equation}
且方向都为逆时针。由四重旋转对称性，四个网孔电流必然相等，记为 $i$。每一条内部十字支路都是两个相邻网孔的公用边，两个等大网孔电流在其上方向相反，故内部支路的实际电流为零。

对任一小网孔写基尔霍夫方程，只有位于外边界的两段电阻产生压降：
\begin{equation}
  2ri=ka^2.
\end{equation}
因此
\begin{equation}
  \boxed{i=\frac{ka^2}{2r}}.
\end{equation}
外周八段小边均通有大小 $ka^2/(2r)$ 的逆时针电流，中央十字的四段电流均为零。

图~\ref{ch24:fig:review-gridcurrentdiagram-1}中，向纸内磁场增强时，外周逆时针电流产生向纸外的反抗磁场，内部支路电流为零。
\par\noindent\begin{minipage}{\linewidth}
\gridCurrentDiagram
\figcap{田字格的外周环流}\label{ch24:fig:review-gridcurrentdiagram-1}
\end{minipage}\par\medskip
\end{xesolution}



\begin{origexamplebox}{【万能的基尔霍夫】}\label{ch24:ex:s40-9-16}
证明：任意复杂的、仅含电阻和理想电源的线性网络放入匀强时变磁场 $B=B_0+kt$ 中，只要所有支路参数和各独立回路所围面积均已知，就能求出所有支路电流；并证明此时仍满足电源叠加定理。
\end{origexamplebox}

\begin{xesolution}
为网络选取 $n$ 个独立网孔电流，组成列向量 $\vect i$。记 $\vect R$ 为由支路电阻组成的网孔电阻矩阵，$\vect e$ 为理想电源在各网孔中的代数和，
$\vect S=(S_1,\ldots,S_n)^{\mathsf T}$ 为有符号面积列向量。法拉第--基尔霍夫方程可写成
\begin{equation}
  \boxed{\vect R\vect i=\vect e-k\vect S}.
\end{equation}
对于非平面网络，“网孔”应改用一般的独立回路，$S_j$ 为其有向面积在磁场方向的投影。
设支路--回路矩阵为 $\vect A$，支路电阻矩阵为 $\vect D=\operatorname{diag}(R_e)$，
则 $\vect R=\vect A^{\mathsf T}\vect D\vect A$，且
$\vect i^{\mathsf T}\vect R\vect i=\sum_eR_eI_e^2\ge0$。
这只能无条件保证半正定。若每个非零独立环流都经过正电阻，才有严格正定及可逆性。所以在此条件下
\begin{equation}
  \boxed{\vect i=\vect R^{-1}\vect e-k\vect R^{-1}\vect S}.
\end{equation}
再用网孔--支路关联矩阵即可恢复所有支路电流。上式的第一项是普通电源单独作用的响应，第二项是时变磁场等效电源单独作用的响应，两者线性相加，故叠加定理仍然成立。
若存在纯理想电源或零电阻回路，必须先检查约束：回路电动势不相容时无有限稳态解，
相容时也可能不能唯一确定电流。例如两只等电压理想电源并联，其分担电流不由理想模型唯一决定。
因此题设的“任意”须理解为非退化、约束相容的线性网络；退化情况不能直接写 $\vect R^{-1}$。
\end{xesolution}



\section{自感、互感与动态回路}

\subsection{自感、电感的暂态与磁场能量}


给定线性回路中，由自身电流 $I$ 产生的磁链与电流成正比：
\begin{equation}
  \boxed{\Lambda=LI},
  \qquad
  \boxed{\mathcal E_L=-\frac{\dd(LI)}{\dd t}}.
\end{equation}
几何形状不变且介质线性时 $L$ 为常量，$\mathcal E_L=-L\dot I$。对长度 $l$、截面积 $A$、单位长度
$n$ 匝的理想长螺线管，
\begin{equation}
  \boxed{L=\mu_0n^2Al}.
\end{equation}
电感不阻碍已有电流，而是阻碍电流的改变。串联 $RL$ 回路接上恒定电源 $\mathcal E_0$ 后，
\begin{equation}
  L\dot I+RI=\mathcal E_0,
  \qquad
  I(t)=\frac{\mathcal E_0}{R}\left(1-\e^{-t/\tau_L}\right),
  \qquad \tau_L=\frac LR.
\end{equation}
下列螺线管公式采用真空或非磁性芯近似。电感储能与真空磁场能量密度分别为
\begin{equation}
  \boxed{W_L=\frac12LI^2},
  \qquad
  \boxed{u_B=\frac{B^2}{2\mu_0}}.
\end{equation}
若为无磁滞、可逆的线性介质，能量密度为 $u_B=\vect B\cdot\vect H/2$；
在线性各向同性介质中等于 $B^2/(2\mu)$，不能无条件仍用 $\mu_0$。

图~\ref{ch24:fig:review-selffieldlinesdiagram-1}中，单环和多匝线圈都有内部场与外部回流场，同向绕组的内场贡献相加。
\par\noindent\begin{minipage}{\linewidth}
\selfFieldLinesDiagram
\figcap{单环与有限螺线管的自生场}\label{ch24:fig:review-selffieldlinesdiagram-1}
\end{minipage}\par\medskip
图中的单环电流按右手规则产生向上的轴场，场线从上侧绕回下侧；多匝同向绕组使内部磁场叠加，
外部回流场仍存在。计算自感时，必须计算这些自生磁场穿过本回路各匝的磁链，而不是把外部场另行遗漏。

自感系数的单位是 henry（H）。由 $\mathcal E_L=-L\dot I$ 可见，
$1\,\mathrm H=1\,\mathrm{V\,s/A}$；由 $\Lambda=LI$ 又可写成
$1\,\mathrm H=1\,\mathrm{Wb/A}$。同一回路中电流增加时，自感电动势方向与电流方向相反；电流减小时，
自感电动势与原电流同向，试图维持它。所谓“电感阻碍电流”是不准确的：稳恒直流通过理想电感时
$\dot I=0$，自感电动势为零；它阻碍的是电流的变化率。

自感是所有闭合导体回路的普遍属性，大小由回路形状、匝数、尺寸和周围介质决定；在线性介质中与电流
大小无关。若几何结构本身随时间变化，则 $L$ 也可能变化，必须使用完整形式
$\mathcal E_L=-\dd(LI)/\dd t$，不能直接简写为 $-L\dot I$。

对总匝数 $N=nl$ 的长螺线管，内部 $B=\mu_0NI/l$。每匝磁通为 $BA$，总磁链
\begin{equation}
  \Lambda=NBA=\frac{\mu_0N^2A}{l}I,
\end{equation}
所以
\begin{equation}
  L=\frac{\mu_0N^2A}{l}=\mu_0n^2Al.
\end{equation}
匝数加倍会使磁场贡献和磁链计数各增加一倍，因此自感按 $N^2$ 增长。

接通 $RL$ 回路时，$I(0)=0$，电感两端最初承担几乎全部电源电压；随后电流按时间常量 $L/R$ 上升，
最终趋于 $\mathcal E_0/R$，此时电流不再变化，理想电感压降趋零。断开含大电感的电路时，电感会试图
维持原电流并产生很高反向电动势，这既可被利用，也可能造成开关打火或绝缘损坏。

图~\ref{ch24:fig:s40-9-9} 右侧把电源、开关、电阻和电感的连接关系完整画出。开关闭合后，沿图示电流方向列回路方程为
$\mathcal E_0=RI+L\dot I$；端点 $a,b$ 的电压为 $U_{ab}=L\dot I$。因此初始瞬间 $U_{ab}=\mathcal E_0$，
稳态时 $U_{ab}\to0$，这两端极限也能直接检查解析解是否合理。

\rlInductorContext

若把已经达到稳恒电流 $I_0$ 的 $RL$ 支路从电源断开，但仍让电阻与电感构成闭合回路，方程变为
\begin{equation}
  L\dot I+RI=0,\qquad I(t)=I_0\e^{-tR/L}.
\end{equation}
上升和衰减具有同一时间常量 $\tau_L=L/R$。在一个时间常量后，衰减电流剩下 $1/\e$，上升电流达到
终值的 $1-1/\e$。初始电流不能瞬间突变，否则 $\dot I$ 无穷大，需要无限自感电动势；实际电路中的
寄生电容、击穿和开关过程会限制电压峰值。

储能式可以直接由功率推出。给电感建立电流时，外电源克服自感电动势所提供的功率为
$P=UI=LI\dot I$，从零电流积分到 $I$：
\begin{equation}
  W_L=\int_0^I LI'\,\dd I'=\frac12LI^2.
\end{equation}
当电流减小时，这部分能量会重新流回电路或在电阻中耗散。能量并不是“藏在线圈铜线里”，而是分布在
线圈电流建立的磁场中；对全空间积分 $B^2/(2\mu_0)$ 应与 $LI^2/2$ 相等。反过来，若某个回路磁场
容易计算，也可先算总磁场能量，再由 $W=LI^2/2$ 反求自感。

以长螺线管作检验，内部体积为 $Al$，场强 $B=\mu_0nI$，所以
\begin{equation}
  \int u_B\,\dd V
  =\frac{(\mu_0nI)^2}{2\mu_0}Al
  =\frac12(\mu_0n^2Al)I^2
  =\frac12LI^2.
\end{equation}
这既验证了能量密度，也说明理想长螺线管的能量主要储存在管内近似匀强磁场区域。有限线圈还会有外部
漏场，对应的能量同样必须计入，实际自感也会因此偏离无限长近似。

图~\ref{ch24:fig:review-selfinductancediagram-1}中，内部磁场穿过每一匝，磁链要将各匝通量相加，才得到 $\Lambda=LI$。
\par\noindent\begin{minipage}{\linewidth}
\selfInductanceDiagram
\figcap{螺线管的自生磁链}\label{ch24:fig:review-selfinductancediagram-1}
\end{minipage}\par\medskip


% EM teaching enrichment 20260927
\paragraph{开关题先取连续量，再取最终状态}
有限端电压下电感电流连续：$I_L(0^+)=I_L(0^-)$；有限支路电流下电容电压连续：$U_C(0^+)=U_C(0^-)$。它们是从构成方程积分得到的初值条件。
例如初始电流 $I_0$ 的线圈只经电阻 $R$ 放电，热量
\[
 Q_R=\int_0^\infty RI_0^2e^{-2Rt/L}\dd t=\frac12LI_0^2,
\]
与初磁能相符。初始端电压大小为 $RI_0$，增加 $R$ 会加快衰减并提高电压，而不是使电流在开关瞬间消失。
理想“直接开路”不给电流出路，有限电压的简单 $RL$ 模型便不闭合，应考虑实际寄生电容或放电通道。


\subsection{互感与互易性}


回路 1 中的电流会在回路 2 中建立磁链，反之亦然。选定两个回路的正向后，定义
\begin{equation}
  \Lambda_2=M_{12}I_1,
  \qquad
  \Lambda_1=M_{21}I_2.
\end{equation}
相应的感应电动势为
\begin{equation}
  \mathcal E_2=-M_{12}\frac{\dd I_1}{\dd t},
  \qquad
  \mathcal E_1=-M_{21}\frac{\dd I_2}{\dd t},
\end{equation}
其中假定几何位置不变、互感为常量。若两线圈相对移动，必须保留
$-\dd(MI)/\dd t$ 的完整形式，运动本身也会改变磁链。
对线性、静止的回路系统，互易定理给出
\begin{equation}
  \boxed{M_{12}=M_{21}\equiv M}.
\end{equation}
$M$ 的正负取决于两回路正向的选取，但对称性不变。两回路共同产生的磁场能量为
\begin{equation}
  \boxed{W=\frac12L_1I_1^2+MI_1I_2+\frac12L_2I_2^2}.
\end{equation}
由能量必须非负还可得耦合界 $M^2\le L_1L_2$。

两个线圈不必有金属接触也能互相影响：回路1电流变化，使穿过回路2的磁通变化，从而在回路2中产生
感应电动势。互感的绝对值由两回路形状、距离、相对朝向以及介质决定；正负只反映两个回路正方向怎样
约定。若把其中一个线圈的正方向反过来，电流和磁通的符号定义随之改变，$M$ 也反号，但实际电磁现象
没有改变。

对两根细导线闭合回路，互感还可写成诺依曼形式
\begin{equation}
  M=\frac{\mu_0}{4\pi}\oint_{C_1}\!\oint_{C_2}
  \frac{\dd\vect\ell_1\cdot\dd\vect\ell_2}{|\vect r_1-\vect r_2|}.
\end{equation}
交换下标 $1,2$ 后，被积式和积分区域不变，因而直观看出 $M_{12}=M_{21}$。这个对称性不是说两个
回路中产生的磁通一定相等，而是说“回路1单位电流在回路2建立的磁链”与反向过程的比例系数相等。
若匝数不同，磁链已经包含各匝磁通的总和，互易性仍然成立。

互易性 $M_{12}=M_{21}$ 是线性静止系统的普遍结论，不要求两个线圈相同。能量中的交叉项 $MI_1I_2$
表示两组场在空间中的重叠：同向加强时可以为正，反向削弱时可以为负，但总磁场能量必须非负，这便
要求 $|M|\le\sqrt{L_1L_2}$。定义耦合系数 $k=M/\sqrt{L_1L_2}$ 后，总有 $|k|\le1$；理想变压器
希望 $|k|$ 尽可能接近1，而需要隔离的电路则会尽量减小互感。

能量不等式可以把耦合强弱限制得更直观。把总能量看成关于 $I_2$ 的二次式，若
$M^2>L_1L_2$，总能找到某个 $I_2/I_1$ 使二次式为负，这与磁场能量非负矛盾。所以
$|k|\le1$ 不是经验规定，而是能量正定性的结果。$|k|=1$ 表示两个线圈的全部磁通完全耦合；真实线圈
通常有漏磁，故 $|k|<1$。

变压器正是互感的典型应用：原线圈中的交流使铁芯磁通变化，副线圈因共享磁链而产生电动势。理想紧耦合
情况下，电压比近似等于匝数比；实际装置还要考虑漏磁、绕组电阻、铁芯损耗与有限励磁电流。本章只用
互感说明“无金属接触也能传递电磁能量”的机制，具体交流电路分析还需结合阻抗和功率关系。

判定互感电动势方向时仍应回到楞次定律：先判断原线圈电流变化使副线圈中的磁通怎样变化，再令副线圈
电流产生的磁场阻碍这一变化。所谓同名端标记只是把这一判向结果固定在器件符号上，不能脱离线圈绕向
和参考电流方向孤立使用。
绕向或端点定义改变后，符号会改变，但实际端电压与电流响应必须保持一致。
这也是互感计算的基本自检。

图~\ref{ch24:fig:review-mutualinductancediagram-1}中，两个同向电流的正法向一致，回路1的场穿过回路2，反向激励时交换源和受体。
\par\noindent\begin{minipage}{\linewidth}
\mutualInductanceDiagram
\figcap{两个同轴回路的互感}\label{ch24:fig:review-mutualinductancediagram-1}
\end{minipage}\par\medskip


% EM teaching enrichment 20260927
\paragraph{串联耦合线圈与正负号}
两个线圈串联后流过同一电流 $I$。若所选回路正向使互感项为正，总磁能为
$\frac12(L_1+L_2+2M)I^2$，所以等效电感为 $L_1+L_2+2M$。
反接一个线圈，互感项反号，等效电感为 $L_1+L_2-2M$。由 $|M|\le\sqrt{L_1L_2}$ 可检验等效电感不为负。
例如 $L_1=L_2=L$、$M=kL$ 且 $0\le k\le1$，两种接法分别给出 $2L(1+k)$、$2L(1-k)$。反接后只在理想完全耦合 $k=1$ 的极限才为零；真实漏磁使其仍保留正电感。



% EM structure and teaching revision 20260928
\subsection{开关初值、稳态与自然响应}
理想电感满足 $u_L=L\dot I$，所以在无冲激电压时电流连续；理想电容满足 $I_C=C\dot U_C$，所以在无冲激电流时电压连续。开关刚动作后的初值应从动作前的储能状态继承，不能立即套直流稳态的“电感短路、电容开路”。
串联RLC无源放电，取电容电荷 $Q$ 为状态量，电流方向与 $\dot Q$ 一致，则
\[
 L\ddot Q+R\dot Q+\frac QC=0,
 \qquad s_{1,2}=-\frac R{2L}\pm\sqrt{\frac{R^2}{4L^2}-\frac1{LC}}.
\]
当 $R<2\sqrt{L/C}$ 时为衰减振荡；等号为临界阻尼；较大电阻时两个实负根分别控制快、慢衰减。二阶方程需要两个初值，例如 $Q(0)$ 与 $I(0)$。
若外接恒定电源，先求长期特解，再叠加满足初值的自然响应。检查总能量
$W=Q^2/(2C)+LI^2/2$，无源时应有 $\dot W=-RI^2\le0$。若推导得到无源电阻使振幅持续增长，应首先检查电流参考方向和感应电动势的符号。


\subsection{专题例题与方法应用}

\begin{origexamplebox}{【自感线圈】}\label{ch24:ex:s40-9-17}
一个理想密绕长螺线管的长度为 $l$、半径为 $R$、单位长度匝数为 $n$。

(1) 求其自感；

(2) 若把它看成左右两部分，整体自感是否等于两部分自感之和？

% Pure source-question geometry: no derived forces, currents or solution curves.
%
\begin{center}\begin{tikzpicture}[scale=.75,>=Latex]
 \foreach \x in {-3.6,-3.2,...,3.6}{\draw[Ink](\x,0)ellipse(.18 and 1.0);}
 \draw[guide,dashed](0,-1.4)--(0,1.5);
 \draw[<->](-3.8,-1.5)--(0,-1.5)node[midway,below]{$l_1$};
 \draw[<->](0,-1.5)--(3.8,-1.5)node[midway,below]{$l_2$};
 \end{tikzpicture}\end{center}
\figcap{沿同一长螺线管把绕组分为两段，位置保持不变。}\label{ch24:fig:em24-question-17}
\end{origexamplebox}

\begin{xesolution}
螺线管内部磁场为 $B=\mu_0nI$，截面积 $A=\pi R^2$，总匝数 $N=nl$。磁链为
\begin{equation}
  \Lambda=N\Phi=(nl)(\mu_0nI)(\pi R^2).
\end{equation}
所以
\begin{equation}
  \boxed{L=\frac\Lambda I=\mu_0n^2\pi R^2l}.
\end{equation}
若左右两段长度为 $l_1,l_2$，且 $l_1+l_2=l$，先保持两段原有位置和绕向。
同一电流 $I$ 依次通过两段，磁能为
\[
 W=\frac12L_1I^2+MI^2+\frac12L_2I^2
   =\frac12(L_1+L_2+2M)I^2,
\]
所以一般关系是 $L=L_1+L_2+2M$，而不是无条件简单相加。
若两个子段都满足 $l_1,l_2\gg R$，并在长管主阶近似中忽略分界附近的交链端场，则
\begin{equation}
  L_1+L_2\simeq\mu_0n^2\pi R^2(l_1+l_2)\simeq L.
\end{equation}
因此“等于”只指这一近似的主阶结果。整体长不代表任意切出的短段也长；如果一段长度与半径相当，
不能对它继续使用长管公式。单纯把一根线圈在账面上分成两段不会改变其几何或磁场，互感项仍需计入；
实际切开并移远则改变了互感，是另一个装置。

图~\ref{ch24:fig:review-solenoidsplitdiagram-1}中，分界只是对同一线圈的划分；保持位置时总自感还包括两段互感的两倍。
\par\noindent\begin{minipage}{\linewidth}
\solenoidSplitDiagram
\figcap{同向串联的两段线圈}\label{ch24:fig:review-solenoidsplitdiagram-1}
\end{minipage}\par\medskip
\end{xesolution}



\begin{origexamplebox}{【互感线圈】}\label{ch24:ex:s40-9-18}
两个单匝线圈放置在任意位置，分别通入电流 $I_1,I_2$。证明：1 号线圈产生的磁场穿过 2 号线圈的磁通量与 $I_1$ 之比，等于 2 号线圈产生的磁场穿过 1 号线圈的磁通量与 $I_2$ 之比，即 $M_{12}=M_{21}$。
% Question contract: only the two loops and prescribed current reference directions.
\begin{center}\begin{tikzpicture}[scale=.82,>=Latex]
 \draw[Ink,thick](-2.5,0)circle(1.2);\draw[Ink,thick](2.5,.4)circle(1);
 \draw[vector](-1.3,0)arc[start angle=0,end angle=65,radius=1.2];
 \draw[vector](3.5,.4)arc[start angle=0,end angle=65,radius=1];
 \node at(-2.5,0){线圈1};\node at(2.5,.4){线圈2};
 \node[right]at(-1.35,.95){$I_1$};\node[right]at(3.45,1.2){$I_2$};
\end{tikzpicture}\end{center}
\figcap{两个单匝线圈及各自的电流参考方向。}\label{ch24:fig:em24-question-18}

\end{origexamplebox}

\begin{xesolution}
\par\noindent\begin{minipage}{\linewidth}
\mutualLoopsProofDiagram
\figcap{互感双重积分的源点和场点}\label{ch24:fig:review-mutualloopsproofdiagram-1}
\end{minipage}\par\medskip

使用稳恒线电流的矢势
\begin{equation}
  \vect A_1(\vect r_2)=\frac{\mu_0I_1}{4\pi}
  \oint_{C_1}\frac{\dd\vect\ell_1}{r_{12}}.
\end{equation}
由斯托克斯定理，1 号线圈在 2 号线圈中建立的磁通为
\begin{align}
  \Phi_{2\leftarrow1}
  &=\oint_{C_2}\vect A_1\cdot\dd\vect\ell_2\\
  &=\frac{\mu_0I_1}{4\pi}\oint_{C_2}\!\oint_{C_1}
  \frac{\dd\vect\ell_1\cdot\dd\vect\ell_2}{r_{12}}.
\end{align}
因此
\begin{equation}
  M_{12}=\frac{\Phi_{2\leftarrow1}}{I_1}
  =\frac{\mu_0}{4\pi}\oint_{C_1}\!\oint_{C_2}
  \frac{\dd\vect\ell_1\cdot\dd\vect\ell_2}{r_{12}}.
\end{equation}
交换积分次序及下标 $1,2$，点积和距离 $r_{12}$ 均不变，所以
\begin{equation}
  \boxed{M_{12}=M_{21}}.
\end{equation}
两回路的形状、大小和相对位置可以任意；以上证明采用真空中的稳恒线电流近似。
普通线性互易介质也有互易性，但不能把真空积分核无条件照搬到任意介质。

图~\ref{ch24:fig:review-mutualloopsproofdiagram-1}中，两个线元落在各自回路上并取切向，交换下标不改变点积及两点距离。

\end{xesolution}



\begin{origexamplebox}{【RLC 回路】}\label{ch24:ex:s40-9-19}
电阻 $R$、电感 $L$、电容 $C$ 串联成回路。初始电流为零，电容初电压为 $U$。$t=0$ 闭合开关，求任意时刻的放电电流。
\par\noindent\begin{minipage}{\linewidth}
\rlcDiagram{0}
\figcap{RLC回路及初始电压参考极性}\label{ch24:fig:review-rlcdiagram-1}
\end{minipage}\par\medskip
\end{origexamplebox}

\begin{xesolution}
记电容器带电量为 $q$，并把放电方向的电流取为 $I=-\dot q$。则
\begin{equation}
  L\ddot q+R\dot q+\frac qC=0,
  \qquad q(0)=CU,
  \qquad \dot q(0)=0.
\end{equation}
特征根为
\begin{equation}
  s_\pm=-\frac{R}{2L}\pm
  \sqrt{\left(\frac{R}{2L}\right)^2-\frac1{LC}}.
\end{equation}
对 $s_+\ne s_-$，电流的统一表达式为
\begin{equation}
  \boxed{I(t)=\frac UL\,
  \frac{\e^{s_+t}-\e^{s_-t}}{s_+-s_-}}.
\end{equation}
当 $R<2\sqrt{L/C}$ 时，记
\begin{equation}
  \alpha=\frac{R}{2L},
  \qquad
  \omega_d=\sqrt{\frac1{LC}-\alpha^2},
\end{equation}
则欠阻尼电流为
\begin{equation}
  \boxed{I(t)=\frac{U}{L\omega_d}\e^{-\alpha t}\sin\omega_dt}.
\end{equation}
临界阻尼 $R=2\sqrt{L/C}$ 时，
\begin{equation}
  \boxed{I(t)=\frac{U}{L}\,t\e^{-Rt/(2L)}}.
\end{equation}
过阻尼情形直接使用实根形式。电流的初值为零是电感电流不能突变的结果，而非“电容还没开始放电”。

图~\ref{ch24:fig:review-rlcdiagram-1}中，图示开关已闭合，C为等长双极板，R、L使用独立元件符号；电流正向与初始电压极性共同决定方程。

\end{xesolution}



\begin{origexamplebox}{【电源凑热闹】}\label{ch24:ex:s40-9-20}
电阻 $R$、电感 $L$、电容 $C$ 与恒定电源 $\mathcal E$ 串联成回路。初始电流为零，电容初电压为 $U$。$t=0$ 闭合开关，求任意时刻的电流。
\par\noindent\begin{minipage}{\linewidth}
\rlcDiagram{1}
\figcap{RLC回路及初始电压参考极性}\label{ch24:fig:review-rlcdiagram-2}
\end{minipage}\par\medskip
\end{origexamplebox}

\begin{xesolution}
取电源充电方向为电流正方向，令 $q$ 为电容正极板电量，$I=\dot q$。回路方程为
\begin{equation}
  L\ddot q+R\dot q+\frac qC=\mathcal E,
  \qquad q(0)=CU,
  \qquad \dot q(0)=0.
\end{equation}
把 $q=C\mathcal E+y$ 代入，$y$ 满足与上题完全相同的齐次方程，只是初值由 $CU$ 换成 $C(U-\mathcal E)$。因此
\begin{equation}
  \boxed{I(t)=\frac{\mathcal E-U}{L}\,
  \frac{\e^{s_+t}-\e^{s_-t}}{s_+-s_-}},
\end{equation}
其中 $s_\pm$ 与上题相同。欠阻尼时
\begin{equation}
  \boxed{I(t)=\frac{\mathcal E-U}{L\omega_d}
  \e^{-\alpha t}\sin\omega_dt}.
\end{equation}
若 $U=\mathcal E$，初始电容电压已与电源平衡，回路不产生电流；若 $U>\mathcal E$，上式给出负电流，即电容反向向电源一侧放电。
其中 $U$ 按图中极性定义为有符号初电压。临界阻尼时两个根重合，不能直接代入上面的商式；
利用 $\lim_{s_+\to s_-}(\e^{s_+t}-\e^{s_-t})/(s_+-s_-)=t\e^{s_-t}$，得到
\begin{equation}
 \boxed{I(t)=\frac{\mathcal E-U}{L}t\e^{-Rt/(2L)}},\qquad R=2\sqrt{L/C}.
\end{equation}
过阻尼时用两个不同实根的通式。各分支均满足 $I(0)=0$、$\dot I(0)=(\mathcal E-U)/L$。

图~\ref{ch24:fig:review-rlcdiagram-2}中，图示开关已闭合，C为等长双极板，R、L使用独立元件符号；电流正向与初始电压极性共同决定方程。

\end{xesolution}



\begin{origexamplebox}{【磁悬浮列车】}\label{ch24:ex:s40-9-21}
把列车等效为一个长为 $a$、宽为 $L$、总电阻为 $R_c$ 且自感可忽略的闭合导体框，质量为 $m$。导轨中存在行波磁场
\begin{equation*}
  B(x,t)=B_0\cos(kx-\omega t),
\end{equation*}
方向垂直导体框平面。

(1) 若某时刻车体左边坐标为 $x$、向右速度为 $v$，求瞬时电磁动力；

(2) 解释这种行波磁场为何能驱动列车。
% Question contract: given frame, speed and field law; no phase-speed result or zero-field construction.
\begin{center}\begin{tikzpicture}[scale=.8,>=Latex]
 \draw[Ink,thick](-4,1.4)--(4,1.4) (-4,-1.4)--(4,-1.4);
 \draw[DeepBlue,very thick](-1.2,-1.15)rectangle(1.2,1.15);
 \node at(0,.25){$m,\ R_c$};
 \draw[vector](1.45,0)--(3,0)node[right]{$v$};
 \draw[<->](-1.2,-1.9)--(1.2,-1.9)node[midway,below]{$a$};
 \draw[<->](-4.4,-1.15)--(-4.4,1.15)node[midway,left]{$L$};
 \node at(0,2.1){$B(x,t)=B_0\cos(kx-\omega t)$，垂直框面};
\end{tikzpicture}\end{center}
\figcap{闭合导体框与绝缘机械支承；框长、框宽及瞬时速度为已知量。}\label{ch24:fig:em24-question-21}

\end{origexamplebox}

\begin{xesolution}
\par\noindent\begin{minipage}{\linewidth}
\travelingFieldTrainDiagram
\figcap{行波磁场与绝缘机械支承}\label{ch24:fig:review-travelingfieldtraindiagram-1}
\end{minipage}\par\medskip

令
\begin{equation}
  \xi=k\left(x+\frac a2\right)-\omega t,
  \qquad
  \delta=\frac{ka}{2},
  \qquad
  v_\phi=\frac\omega k.
\end{equation}
穿过导体框的磁通为
\begin{align}
  \Phi(x,t)
  &=L\int_x^{x+a}B_0\cos(kx'-\omega t)\dd x'\\
  &=\frac{2LB_0}{k}\sin\delta\cos\xi.
\end{align}
跟随车体求时间导数，得总感应电动势
\begin{equation}
  \mathcal E=-\frac{\dd\Phi}{\dd t}
  =2LB_0\sin\delta\sin\xi\,(v-v_\phi).
\end{equation}
忽略自感时 $I=\mathcal E/R_c$。对回路做微小平移，磁场力可写为 $F_x=I\,\partial\Phi/\partial x$，而
\begin{equation}
  \frac{\partial\Phi}{\partial x}=-2LB_0\sin\delta\sin\xi.
\end{equation}
因此
\begin{equation}
  \boxed{F_x=\frac{4L^2B_0^2}{R_c}\sin^2\!\delta\,\sin^2\!\xi\,(v_\phi-v)}.
\end{equation}
若速度在一次相位循环中变化很小，可将其近似看成常量。车体看到的相位变化率是
$\dot\xi=kv-\omega$，应在滑差周期 $T_s=2\pi/|\omega-kv|$ 上取平均，而非固定位置的周期 $2\pi/\omega$。
当 $v\ne v_\phi$ 时，在此前提下 $\langle\sin^2\xi\rangle=1/2$，故
\begin{equation}
  \boxed{\langle F_x\rangle=\frac{2L^2B_0^2}{R_c}
  \sin^2\!\frac{ka}{2}\,(v_\phi-v)}.
\end{equation}
当列车慢于行波相速 $v_\phi$ 且耦合非零时，平均力向右；快于行波时平均力反向。
若 $\sin(ka/2)=0$，框长恰跨整数个波长，总磁通恒为零，驱动力对任意速度均为零。
同步时 $v=v_\phi$，磁通也不随车体改变，直接由瞬时式得到零电动势和零力，无须使用发散的滑差周期。

若保留自感 $L_c$，令 $\Omega=\omega-kv$、$K=2LB_0\sin(ka/2)$，在恒速周期稳态下
$L_c\dot I+R_cI=-K\Omega\sin\xi/k$，$F_x=-KI\sin\xi$。
正弦激励的同相响应被因子 $R_c/(R_c^2+\Omega^2L_c^2)$ 缩减，因而
\[
 \langle F_x\rangle=
 \frac{K^2R_c}{2(R_c^2+\Omega^2L_c^2)}(v_\phi-v).
\]
在非零耦合、$R_c>0$ 的稳态平均意义下，方向仍趋向同步，但不能把这个结论套用于任意瞬态。
图中的导轨只作绝缘机械支承，不与闭合导体框构成额外电路；正负磁场区随相位向右移动。
本模型解释电磁推进，不负责解释列车的竖直悬浮力。

图~\ref{ch24:fig:review-travelingfieldtraindiagram-1}中，点、叉区表示某一时刻磁场的相反方向，虚线为零场位置，场图样以相速移动。

\end{xesolution}
}

\endgroup

\begingroup
\tcbset{breakable/.default=false,after={\par\nobreak\smallskip}}
\RenewDocumentEnvironment{xesolution}{+b}{%
  \par\nobreak\smallskip\nobreak\begingroup\small
  \clubpenalties 2 10000 10000\displaywidowpenalties 2 10000 10000\noindent
  {\bfseries\color{DeepBlue}解答}\quad #1\par\endgroup\smallskip
}{}
\chapter{交流电、相量与RLC电路}

\section{周期信号与有效值}

\subsection{拟稳交流电路的适用范围}


把交流电路当作集中参数电路，隐含了“电磁作用传播得足够快”的条件。若电路的最大线度为 $\ell$，电磁扰动在导线中的传播速度为 $v_p$，则传播时间
\begin{equation}
 \tau_p=\frac{\ell}{v_p}
\end{equation}
应远小于交流周期 $T=2\pi/\omega$，即
\begin{equation}
 \omega\tau_p\ll1.
\end{equation}
各元件自身的尺寸也应满足相同条件。此时同一支路各处的电流可视为同时变化，基尔霍夫定律可以逐时刻使用，称为拟稳电路。若 $\ell$ 已能和波长 $\lambda=v_p/f$ 相比，就必须考虑传播时延和分布参数，不能再把一段长导线等同于一个无长度的理想元件。
还须忽略连接导线的分布电阻、电感与对地电容，或把它们明确计入模型；否则即使传播时延很小，导线两端也可能有不可忽略的电压差。以下只研究线性元件的周期稳态，接通电源后的暂态另行处理。


% EM teaching enrichment 20260927
\paragraph{用相位延迟作数量级判断}
传播时间造成的相位差约为 $\Delta\phi=\omega\ell/v_p$。例如取 $\ell=1\,\mathrm m$、$v_p=2.0\times10^8\,\mathrm{m/s}$，在 $f=50\,\mathrm{Hz}$ 时相位差约 $1.6\times10^{-6}$ 弧度；在 $f=100\,\mathrm{MHz}$ 时则约为 $\pi$，已经不能说整段导线“同时”处于同一交流状态。
这一判断不依赖电压大小。短而高频的线路和长而低频的线路，都应比较 $\omega\ell/v_p$，不能只根据“交流”两个字决定是否使用相量电路模型。
此外，相量只描述某一频率的稳态。开关闭合后的自然响应由初值决定，可能在若干时间常数内不可忽略。


\subsection{正弦交流电与有效值}


正弦电压写作
\begin{equation}
 u(t)=U_m\cos(\omega t+\varphi_u).
\end{equation}
交流表通常给出有效值（均方根值）
\begin{equation}
 U=\sqrt{\frac1T\int_0^T u^2(t)\dd t}=\frac{U_m}{\sqrt2},
 \qquad I=\frac{I_m}{\sqrt2}.
\end{equation}
只有纯正弦波才可直接使用峰值除以 $\sqrt2$；任意周期波形应回到均方根定义。


% EM teaching enrichment 20260927
\paragraph{有效值来自相同热效应}
电阻 $R$ 中一周期发热为 $\int_0^T i^2R\dd t$。定义一个直流 $I$ 在同一时间产生相同热量，即
$I^2RT=\int_0^Ti^2R\dd t$，便得到均方根定义；它既不是瞬时值的平均，也不是绝对值的平均。
例如在每周期中通电比例为 $D$ 的矩形脉冲，通电时电流 $I_0$、其余为零，则
\[
 I_{\rm avg}=DI_0,\qquad I_{\rm rms}=I_0\sqrt D.
\]
取 $D=1/4$，平均值为 $I_0/4$，有效值却为 $I_0/2$。
\subsection{非正弦波形的有效值}


\begin{center}
\begin{tikzpicture}[x=.85cm,y=.9cm,>=Latex]
  \draw[axis] (-.2,0)--(8.2,0) node[right]{$t$};
  \draw[axis] (0,-1.7)--(0,1.8) node[above]{$u$};
  \draw[DeepBlue,very thick] (0,1)--(2,1)--(2,-1)--(4,-1)--(4,1)--(6,1)--(6,-1)--(8,-1);
  \draw[guide] (0,1)--(8,1);
  \draw[guide] (0,-1)--(8,-1);
  \node[clearlabel,left] at (0,1) {$U_0$};
  \node[clearlabel,left] at (0,-1) {$-U_0$};
  \node[clearlabel,below] at (4,-1.12) {$T$};
\end{tikzpicture}
\figcap{双极性方波。有效值由平方平均定义，不能根据波形“正负抵消”判断。}\label{ch25:fig:square-wave}
\end{center}
图~\ref{ch25:fig:square-wave} 中正、负平台对平均值相消，却对 $u^2$ 作相同的正贡献。

\begin{origexamplebox}{方波与三角波的有效值}
周期电压在一个周期内先以 $+U_0$ 保持半周期，再以 $-U_0$ 保持半周期。（1）求平均值与有效值；
（2）若改为在 $-U_0$ 与 $+U_0$ 之间匀速往返的对称三角波，求有效值。
\end{origexamplebox}
\begin{xesolution}
方波的正负面积相消，平均值为零；但
\[
 U_{\rm rms}^2=\frac1T\int_0^T u^2\dd t=U_0^2,
 \qquad \boxed{U_{\rm rms}=U_0}.
\]
三角波在半周期内可写为 $u=-U_0+4U_0t/T$，由对称性
\begin{align*}
 U_{\rm rms}^2
 &=\frac{2U_0^2}{T}\int_0^{T/2}\left(1-\frac{8t}{T}+\frac{16t^2}{T^2}\right)\dd t\\
 &=\frac{2U_0^2}{T}\left[t-\frac{4t^2}{T}+\frac{16t^3}{3T^2}\right]_0^{T/2}
 =\frac{2U_0^2}{T}\left(\frac T2-T+\frac{2T}{3}\right)=\frac{U_0^2}{3}.
\end{align*}
故 $\boxed{U_{\rm rms}=U_0/\sqrt3}$。两种结果都具有电压单位；把峰值一律除以 $\sqrt2$ 只对正弦波成立。
\end{xesolution}


% EM teaching enrichment 20260927 supplementary

\paragraph{用同峰值比较检验结果}
峰值相同时，方波几乎始终停在最大幅值，故有效值最大；三角波有较多时间位于较小幅值，
有效值比正弦波更低。三者依次为 $U_0,U_0/\sqrt2,U_0/\sqrt3$。
若波形叠加直流偏置 $U_d$，先写 $u=U_d+u_{\rm ac}$ 且 $\langle u_{\rm ac}\rangle=0$，展开平方可得
$U_{\rm rms}^2=U_d^2+U_{{\rm ac,rms}}^2$；有效值不能直接相加。


% EM structure and teaching revision 20260928
\subsection{混合频率信号的有效值与适用范围}
设周期信号含直流和不同整数倍频率的谐波：
\[
 u(t)=U_d+\sum_{n=1}^{N}U_{mn}\cos(n\omega t+\phi_n).
\]
在共同周期内，直流与各正弦项的乘积平均为零，不同频率项之间的乘积平均也为零，因此
\[
 U_{\mathrm{rms}}^2=U_d^2+\frac12\sum_{n=1}^{N}U_{mn}^2.
\]
同频项却不能按平方直接相加，必须先考虑相位。两项幅值同为 $A$ 的同频电压，相位相同会合成峰值 $2A$，反相则完全抵消。
例如 $u=3+4\cos\omega t+2\sin2\omega t$（单位为伏），平均值为 $3\,\mathrm V$，有效值为 $\sqrt{19}\,\mathrm V$。若接在线性电阻 $R$ 上，平均发热功率为 $19/R$；若接在频率相关网络上，则须分别求各谐波支路响应后再合成实际波形。
相量法一次只处理一个频率。把多频信号压缩成一个“总相量”，会丢掉各频率面对的不同感抗、容抗及传播条件。



\section{相量、阻抗与串并联网络}

\subsection{电阻、电感与电容的相量关系}


本章统一使用\textbf{有效值相量}，而不是峰值相量。设电流
$i(t)=\sqrt2 I\cos(\omega t+\varphi_i)$，定义
\begin{equation}
 \widetilde U=Ue^{\ii\varphi_u},\qquad
 \widetilde I=Ie^{\ii\varphi_i},\qquad
 u(t)=\sqrt2\Re(\widetilde Ue^{\ii\omega t}),\quad
 i(t)=\sqrt2\Re(\widetilde Ie^{\ii\omega t}).
\end{equation}
$\ii^2=-1$ 是虚数单位，不能与瞬时电流 $i(t)$ 混淆。相量的模是有效值；公用的时间因子不再写入相量。
采用关联参考方向，即电流从元件电压的正端流入。电阻满足 $u=Ri$，故 $\widetilde U=R\widetilde I$；
电感满足 $u=L\dd i/\dd t$，对 $e^{\ii\omega t}$ 求导等于乘 $\ii\omega$，故 $\widetilde U=\ii\omega L\widetilde I$；
电容满足 $i=C\dd u/\dd t$，故 $\widetilde I=\ii\omega C\widetilde U$。因此三种理想元件的阻抗为
\begin{equation}
 Z_R=R,\qquad Z_L=\ii\omega L,\qquad Z_C=\frac1{\ii\omega C}=-\frac{\ii}{\omega C}.
\end{equation}
一般阻抗可写为 $Z=\widetilde U/\widetilde I=R+\ii X=|Z|e^{\ii\varphi}$，其中
$|Z|=\sqrt{R^2+X^2}$，$X$ 是电抗，$\varphi=\varphi_u-\varphi_i=\arg Z$。
$X>0$ 称感性，$X<0$ 称容性；容抗的大小是 $1/(\omega C)$，不是 $\omega C$。
因此电阻中电压电流同相；电感电压超前电流 $\pi/2$；电容电压滞后电流 $\pi/2$。
这些相位关系必须和所选时间因子约定一致。若改用 $e^{-\ii\omega t}$，阻抗虚部符号会同时改变，物理结果不变。

\begin{center}
\begin{minipage}{0.78\linewidth}
\centering
\begin{tikzpicture}[scale=.9,>=Latex]
  \draw[axis] (-.3,0)--(5,0) node[right]{$\Re$};
  \draw[axis] (0,-2.1)--(0,2.3) node[above]{$\Im$};
  \draw[vector] (0,0)--(4.1,0) node[below]{$\widetilde U_R$};
  \draw[vector] (0,0)--(0,1.8) node[left]{$\widetilde U_L$};
  \draw[component] (0,0)--(0,-1.05) node[left]{$\widetilde U_C$};
  \draw[vector] (0,0)--(4.1,.75) node[above right]{$\widetilde U$};
  \draw[guide] (4.1,0)--(4.1,.75);
\end{tikzpicture}
\figcap{以电流相量为实轴的串联RLC电压相量。}\label{ch25:fig:series-phasors}
\end{minipage}
\end{center}
图~\ref{ch25:fig:series-phasors} 中纵向净分量为 $U_L-U_C$，横向分量为 $U_R$；总电压是相量和，不能把三个电压有效值直接代数相加。图中箭头表示复平面上的相量，不是空间中的电场方向。


% EM teaching enrichment 20260927
\paragraph{从相量答案还原真实波形}
若某支路 $Z=3+4\ii\,\Omega$，端电压为 $u=10\sqrt2\cos\omega t\,\mathrm V$，取电压相量为实数 $10\,\mathrm V$，则
\[
 \widetilde I=\frac{10}{3+4\ii}=1.2-1.6\ii\,\mathrm A
 =2e^{-\ii\arctan(4/3)}\,\mathrm A.
\]
因而 $i(t)=2\sqrt2\cos[\omega t-\arctan(4/3)]\,\mathrm A$。负相角表示电流滞后电压，并不表示电流有效值为负。
计算时先统一频率、有效值和参考相位，再在复平面解网络，最后取实部还原。不能把相量的虚部当成“另一股实际电流”，也不能把不同频率的相量放在同一张图上直接相加。


\subsection{导纳与并联支路}


并联电路用导纳更直接。定义
\begin{equation}
 Y=\frac1Z=G+\ii B,
 \qquad \widetilde I=Y\widetilde U,
\end{equation}
其中 $G$ 为电导，$B$ 为电纳。并联各支路电压相同，总导纳等于支路导纳之和。对理想电阻、电感和电容，
\begin{equation}
 Y_R=\frac1R,
 \qquad Y_L=-\frac{\ii}{\omega L},
 \qquad Y_C=\ii\omega C.
\end{equation}
导纳也可写成 $Y=|Y|e^{-\ii\varphi}$。若 $Z=R+\ii X$，乘以分母的共轭得到
$G=R/(R^2+X^2)$、$B=-X/(R^2+X^2)$，所以导纳的相角与阻抗相角相反；$G,B$ 的单位均为西门子。

\begin{center}
\begin{tikzpicture}[scale=.9,>=Latex]
  \draw[axis] (-.3,0)--(4.5,0) node[right]{$\Re$};
  \draw[axis] (0,-2.2)--(0,2.4) node[above]{$\Im$};
  \draw[vector] (0,0)--(3.4,0) node[below]{$\widetilde I_R$};
  \draw[vector] (0,0)--(0,1.9) node[left]{$\widetilde I_C$};
  \draw[vector] (0,0)--(0,-1.25) node[left]{$\widetilde I_L$};
  \draw[vector] (0,0)--(3.4,.65) node[above right]{$\widetilde I$};
  \draw[guide] (3.4,0)--(3.4,.65);
\end{tikzpicture}
\figcap{以电压相量为实轴的并联RLC电流相量。}\label{ch25:fig:parallel-phasors}
\end{center}
图~\ref{ch25:fig:parallel-phasors} 中总电流的虚部为 $I_C-I_L$；相量的竖直分量相消时，电源电流只剩与电压同相的 $I_R$。


% EM teaching enrichment 20260927
\paragraph{支路电流很大，总电流仍可很小}
在某一频率，电阻支路 $R=10\,\Omega$、理想感抗 $X_L=5\,\Omega$、理想容抗大小 $X_C=5\,\Omega$ 并联，端电压有效值 $10\,\mathrm V$。以电压为参考，三支路电流相量分别为
$1$、$-2\ii$、$+2\ii\,\mathrm A$，总电流只有 $1\,\mathrm A$。
电源提供有功功率 $10\,\mathrm W$；电感与电容支路仍交换能量，并非没有电流。若把并联各支路电流有效值直接相加成 $5\,\mathrm A$，就忽略了相位。
这个算例也解释了并联谐振时电源电流最小与内部环流较大为何可以同时发生。


\subsection{复阻抗网络}




\begin{origexamplebox}{两支路并联的等效阻抗}
图~\ref{ch25:fig:impedance-network} 中，上支路为 $R_1$ 与 $L$ 串联，下支路为 $R_2$ 与 $C$ 串联。写出端口等效阻抗，并给出端口呈纯电阻的条件。
\begin{center}
\begin{tikzpicture}[scale=.92,>=Latex]
  \draw[ground] (0,0)--(.7,0)--(.7,1.25)--(1.4,1.25);
  \draw[ground] (.7,0)--(.7,-1.25)--(1.4,-1.25);
  \foreach \y/\name in {1.25/R_1,-1.25/R_2}{
    \draw[thick] (1.4,\y-.18) rectangle (2.8,\y+.18);
    \node[above] at (2.1,\y+.23) {$\name$};
    \draw[ground] (2.8,\y)--(3.5,\y);
  }
  \draw[thick] (3.5,1.25)
    .. controls (3.5,1.7) and (4,1.7) .. (4,1.25)
    .. controls (4,1.7) and (4.5,1.7) .. (4.5,1.25)
    .. controls (4.5,1.7) and (5,1.7) .. (5,1.25)
    .. controls (5,1.7) and (5.5,1.7) .. (5.5,1.25);
  \node[above] at (4.5,1.75) {$L$};
  \draw[ground] (3.5,-1.25)--(4.25,-1.25);
  \draw[thick] (4.25,-1.65)--(4.25,-.85) (4.55,-1.65)--(4.55,-.85);
  \node[below] at (4.4,-1.7) {$C$};
  \draw[ground] (4.55,-1.25)--(6.2,-1.25)--(6.2,1.25)--(5.5,1.25);
  \draw[ground] (6.2,0)--(7,0);
  \fill (0,0) circle (2pt) node[left] {$A$};
  \fill (7,0) circle (2pt) node[right] {$B$};
  \foreach \x in {.7,6.2}{\fill (\x,0) circle (1.5pt);}
\end{tikzpicture}
\figcap{两条串联支路连接于同一端口 $A,B$。}\label{ch25:fig:impedance-network}
\end{center}
\end{origexamplebox}
\begin{xesolution}
两支路阻抗为
\[
 Z_1=R_1+\ii\omega L,
 \qquad Z_2=R_2-\frac{\ii}{\omega C}.
\]
并联应相加导纳：
\[
 \boxed{Z_{AB}=\left(\frac1{Z_1}+\frac1{Z_2}\right)^{-1}}
 =\frac{Z_1Z_2}{Z_1+Z_2}.
\]
端口电流与电压同相要求总导纳虚部为零。由
\[
 \Im\!\left(\frac1{Z_1}+\frac1{Z_2}\right)
 =-\frac{\omega L}{R_1^2+(\omega L)^2}
 +\frac{1/(\omega C)}{R_2^2+[1/(\omega C)]^2}
\]
得到纯电阻条件
\[
 \boxed{\frac{\omega L}{R_1^2+(\omega L)^2}
 =\frac{1/(\omega C)}{R_2^2+[1/(\omega C)]^2}}.
\]
当 $R_1,R_2\ge0$ 且 $L,C>0$ 时，交叉相乘并约去非零的 $\omega$，还可将频率条件写成
\[
 L-CR_1^2=\omega^2LC(L-CR_2^2).
\]
若 $L-CR_2^2\ne0$，得到
\[
 \boxed{\omega^2=\frac{L-CR_1^2}{LC(L-CR_2^2)}}.
\]
右端必须严格为正才有非零实频率。若分母为零而分子不为零，无解；若两者都为零，
即 $R_1=R_2=\sqrt{L/C}$，则在任意非零频率都有 $Z_{AB}=\sqrt{L/C}$，不再只有一个特殊频率。
若至少一条支路有正电阻，则总电导为正，上式确实对应有限的纯电阻端口。
若 $R_1=R_2=0$，上式仍给出 $\omega^2LC=1$，但此时总导纳连实部也为零，即 $Y=0$、端口阻抗趋于无穷大，不能称为有限的纯电阻。
这时非零的电容、电感支路电流相消，电源电流为零。各导纳项单位都是西门子。
\end{xesolution}


% EM teaching enrichment 20260927 supplementary

\paragraph{把复数结果还原成可测量量}
算出 $Z_{AB}=R_{\rm eq}+\ii X_{\rm eq}$ 后，电流有效值为 $I=U/|Z_{AB}|$，电压相对于电流的相角为
$\arg Z_{AB}$，平均功率为 $I^2R_{\rm eq}$。这里的等效实部一般随频率变化，不能视为一只固定物理电阻。
也可分别算 $\widetilde I_1=\widetilde U/Z_1$、$\widetilde I_2=\widetilde U/Z_2$，检查相量和是否等于端口电流，
并核对 $I_1^2R_1+I_2^2R_2$ 是否等于总有功功率。这给出了独立于复杂代数化简的验算方法。


% EM structure and teaching revision 20260928
\subsection{并联RC网络的相位与频率极限}
例如电阻 $R$ 与电容 $C$ 并联，端电压相量为 $\widetilde U$，有
\[
 \widetilde I=\left(\frac1R+\ii\omega C\right)\widetilde U,
 \quad |I|=U\sqrt{R^{-2}+\omega^2C^2},
 \quad \arg\widetilde I-\arg\widetilde U=\arctan(\omega CR).
\]
电流超前电压；低频时电容支路趋于开路，结果退回 $I=U/R$；频率升高后无功支路电流增大，但电阻平均消耗仍为 $U^2/R$。
若取 $R=1.0\,\mathrm{k}\Omega$、$C=1.0\,\mu\mathrm F$，两支路电流幅值相等发生在 $\omega RC=1$，即 $f\simeq159\,\mathrm{Hz}$，此时总电流超前电压 $45^\circ$，总电流幅值是单个支路的 $\sqrt2$ 倍，而非两倍。
所有理想元件的低频、高频极限都是检查网络答案的有效工具，但极限下还要注意是否出现理想电源短路或无阻尼共振等使模型失去有限稳态的情形。



\section{谐振、品质因数与频率响应}

\subsection{串联RLC电路与共振}


串联电路总阻抗
\begin{equation}
 Z=R+\ii\left(\omega L-\frac1{\omega C}\right),
\end{equation}
设 $R,L,C>0$，电源有效值 $U$ 不随频率改变。电流有效值和电压超前电流的相位差为
\begin{equation}
 I=\frac{U}{\sqrt{R^2+(\omega L-1/\omega C)^2}},
 \qquad
 \tan\varphi=\frac{\omega L-1/(\omega C)}R.
\end{equation}
当
\begin{equation}
 \omega_0=\frac1{\sqrt{LC}}
\end{equation}
时发生串联谐振，总阻抗最小为 $R$，电流最大并与电源电压同相。此时 $U_L$ 与 $U_C$ 大小相等、相位相反，
各自可以远大于电源电压，这称为电压谐振，并不违反能量守恒：电感磁能与电容电能在周期内往返交换。

对上述恒参数串联模型，品质因数和半功率带宽满足精确关系
\begin{equation}
 Q=\frac{\omega_0L}{R}=\frac1{\omega_0CR},
\qquad \Delta\omega=\frac{R}{L},\qquad Q=\frac{\omega_0}{\Delta\omega}.
\end{equation}
这里 $Q=U_L/U_R$ 是谐振时的电压放大倍数。只有把两个半功率频率近似写成
$\omega_0\pm\Delta\omega/2$ 时，才额外要求 $Q\gg1$。后面的带宽例题给出不作弱损耗近似的证明。


% EM teaching enrichment 20260927
\paragraph{品质因数的能量含义}
谐振时电流有效值为 $I$，电感最大磁能为 $LI^2$；电容最大电能相同，总储能在两者间交换。
电阻每周期损失能量 $I^2R(2\pi/\omega_0)$，所以
\[
 2\pi\frac{\text{储能}}{\text{每周期耗能}}
 =\frac{\omega_0L}{R}=Q.
\]
高 $Q$ 表示每周期损失比例小，因而稳态可建立较强内部储能，但建立与衰减也较慢。电流峰值、相位过零与电容电压的峰值不必在所有模型中出现在同一频率；使用“共振频率”时应先说明最大化哪个响应量。

\begin{origexamplebox}{串联RLC的谐振参数}
$R=20\,\Omega$、$L=0.20\,\mathrm H$、$C=50\,\mu\mathrm F$ 串联接正弦电源，有效值 $U=10\,\mathrm V$。求谐振角频率、谐振电流和谐振时电感电压。
\end{origexamplebox}
\begin{xesolution}
\[
 \omega_0=\frac1{\sqrt{LC}}
 =\frac1{\sqrt{0.20\times50\times10^{-6}}}
 \approx316\,\mathrm{rad/s}.
\]
谐振时总阻抗为 $R$，故
\[
 I_0=\frac UR=0.50\,\mathrm A.
\]
电感电抗 $X_L=\omega_0L\approx63.2\,\Omega$，所以
\[
 \boxed{U_L=I_0X_L\approx31.6\,\mathrm V}.
\]
它大于电源电压，但电容电压同样为 $31.6\,\mathrm V$ 且相位相反，二者相量相消。
\end{xesolution}


\subsection{并联RLC电路与共振}


对电阻、电感和电容三个理想支路并联的电路，总导纳为
\begin{equation}
 Y=\frac1R+\ii\left(\omega C-\frac1{\omega L}\right).
\end{equation}
当 $\omega=\omega_0=1/\sqrt{LC}$ 时，总导纳的虚部为零，电源电流与电压同相，且大小取最小值 $I_0=U/R$。此时
\begin{equation}
 I_C=\omega_0CU,
 \qquad I_L=\frac{U}{\omega_0L},
 \qquad I_C=I_L,
\end{equation}
两条无功支路之间可以有很大的循环电流，这称为电流谐振；电源只需补偿电阻损耗。

真实线圈含串联电阻，电容也有损耗。此时“总电流与电压同相”的相位谐振频率，与总电流极小或端口阻抗极大的振幅谐振频率一般不严格相同。只有损耗足够小时，两者才都近似为 $1/\sqrt{LC}$；因此不能不加条件地把所有并联谐振公式照搬到实际电路。

\begin{center}
\begin{tikzpicture}[scale=.78,>=Latex]
 \begin{scope}
  \draw[thick] (0,.4)--(0,2)--(1,2) (2,2)--(3,2)--(3,1.35);
  \draw[thick] (1,1.8) rectangle (2,2.2);\node[above] at (1.5,2.25) {$R$};
  \draw[thick] (2.65,1.35)--(3.35,1.35) (2.65,1.1)--(3.35,1.1);
  \node[right] at (3.4,1.22) {$C$};
  \draw[thick] (3,1.1)--(3,.7)
    .. controls (3.6,.7) and (3.6,.2) .. (3,.2)
    .. controls (3.6,.2) and (3.6,-.3) .. (3,-.3)
    .. controls (3.6,-.3) and (3.6,-.8) .. (3,-.8)
    --(3,-1.5)--(0,-1.5)--(0,-.4);
  \node[right] at (3.55,-.1) {$L$};
  \draw (0,0) circle (.4);\node at (0,0) {$\sim$};\node[left] at (-.45,0) {$U$};
  \node at (1.5,-2.1) {串联：$|Z|$ 在谐振时最小};
 \end{scope}
 \begin{scope}[shift={(7,0)}]
  \draw[thick] (0,.4)--(0,2)--(4.8,2)--(4.8,1.5);
  \draw[thick] (0,-.4)--(0,-1.5)--(4.8,-1.5)--(4.8,-.8);
  \draw[thick] (2.4,2)--(2.4,1.5);
  \foreach \x/\name in {2.4/R_C,4.8/R_L}{
   \draw[thick] (\x-.18,.5) rectangle (\x+.18,1.5);
   \node[right] at (\x+.25,1) {$\name$};
  }
  \draw[thick] (2.4,.5)--(2.4,0);
  \draw[thick] (2.05,0)--(2.75,0) (2.05,-.25)--(2.75,-.25);
  \draw[thick] (2.4,-.25)--(2.4,-1.5);\node[right] at (2.85,-.15) {$C$};
  \draw[thick] (4.8,.5)--(4.8,.4)
   .. controls (5.35,.4) and (5.35,0) .. (4.8,0)
   .. controls (5.35,0) and (5.35,-.4) .. (4.8,-.4)
   .. controls (5.35,-.4) and (5.35,-.8) .. (4.8,-.8);
  \node[right] at (5.3,-.3) {$L$};
  \draw (0,0) circle (.4);\node at (0,0) {$\sim$};\node[left] at (-.45,0) {$U$};
  \foreach \y in {-1.5,2}{\fill (2.4,\y) circle (1.5pt);}
  \node at (2.4,-2.1) {有损并联：相位与振幅谐振须区分};
 \end{scope}
\end{tikzpicture}
\figcap{串联与有损并联谐振电路的对照。}\label{ch25:fig:resonance-circuits}
\end{center}
图~\ref{ch25:fig:resonance-circuits} 右图中 $R_C$ 与 $R_L$ 分别表示电容、电感支路的串联损耗，其总导纳为
\[
 Y=\frac1{R_C-\ii/(\omega C)}+\frac1{R_L+\ii\omega L}.
\]
先令 $\Im Y=0$ 求相位谐振，再根据需要对 $|Y|$ 求极小值；不能直接把左图的 $R$ 搬到右图。


% EM teaching enrichment 20260927
\paragraph{同一器件换连接方式，控制量也变了}
在理想并联模型中，电源固定端电压，电感与电容的电流虚部相消，谐振时总阻抗为并联电阻 $R$，电源电流最小。串联模型固定的是共同电流关系，各电抗电压相消，谐振时总阻抗为串联电阻 $R$，电源电流最大。
对有损并联支路，电导本身也可能随频率变化。因此 $\Im Y=0$ 只保证电源电流与电压同相；最小 $|Y|$ 还要求
$\dd(G^2+B^2)/\dd\omega=0$。若 $B=0$ 处 $G\,\dd G/\dd\omega\ne0$，它就不是电流最小点。后面的实际线圈例题应沿这两个问题分别计算。


\subsection{谐振曲线与半功率带宽}


\begin{center}
\begin{tikzpicture}
\begin{axis}[width=.74\linewidth,height=5.3cm,axis lines=left,
  xlabel={$\omega/\omega_0$},ylabel={$\overline P/P_{\max}$},
  xmin=0,xmax=2.2,ymin=0,ymax=1.15,
  xtick={1},ytick={.5,1},every axis label/.style={fill=none}]
  \addplot[DeepBlue,very thick,domain=.05:2.2,samples=320]{1/(1+16*(x-1/x)^2)};
  \addplot[TextGray,dashed] coordinates{(.05,.5)(2.2,.5)};
  \draw[TextGray,dotted] (axis cs:.8827822185,0)--(axis cs:.8827822185,.5);
  \draw[TextGray,dotted] (axis cs:1.1327822185,0)--(axis cs:1.1327822185,.5);
  \draw[<->,TextGray] (axis cs:.8827822185,.25)--(axis cs:1.1327822185,.25);
  \node[anchor=west] at (axis cs:1.32,.86) {$Q=4$};
  \node[anchor=west] at (axis cs:1.32,.72) {$\Delta\omega/\omega_0=1/4$};
\end{axis}
\end{tikzpicture}
\figcap{串联RLC的归一化平均功率曲线。半功率两频率之差定义带宽。}\label{ch25:fig:bandwidth}
\end{center}
图~\ref{ch25:fig:bandwidth} 中水平虚线与曲线的两个交点定义 $\omega_1,\omega_2$。

\begin{origexamplebox}{半功率频率的精确关系}
串联RLC接恒定有效值正弦电压。证明两个半功率角频率 $\omega_1<\omega_2$ 满足
$\omega_2-\omega_1=R/L$ 和 $\omega_1\omega_2=1/(LC)$。
\end{origexamplebox}
\begin{xesolution}
平均功率
\[
 \overline P=\frac{U^2R}{R^2+(\omega L-1/\omega C)^2},
 \qquad P_{\max}=\frac{U^2}{R}.
\]
半功率条件给出
\[
 \left|\omega L-\frac1{\omega C}\right|=R.
\]
低频侧取负号、高频侧取正号：
\[
 L\omega_1^2+R\omega_1-\frac1C=0,
 \qquad
 L\omega_2^2-R\omega_2-\frac1C=0.
\]
分别取两个二次方程的正根，令 $D=\sqrt{R^2+4L/C}$，得到
\[
 \omega_1=\frac{D-R}{2L},\qquad \omega_2=\frac{D+R}{2L}.
\]
相减时 $D$ 消去，得到 $(2R)/(2L)=R/L$；相乘时利用平方差，得到
$(D^2-R^2)/(4L^2)=(4L/C)/(4L^2)=1/(LC)$。因此
\[
 \boxed{\omega_2-\omega_1=\frac RL},
 \qquad \boxed{\omega_1\omega_2=\omega_0^2=\frac1{LC}}.
\]
第一式量纲为 $\Omega/\mathrm H=\mathrm{s^{-1}}$。弱损耗时两频率近似关于 $\omega_0$ 对称，精确地则是几何平均为 $\omega_0$。
\end{xesolution}


% EM teaching enrichment 20260927 supplementary

\paragraph{读图时区分半功率与半振幅}
电阻固定时 $P\propto I^2$，所以半功率点的电流是峰值的 $1/\sqrt2$，并非一半。
由 $\omega_1\omega_2=\omega_0^2$ 可知两个频率以几何平均意义围绕 $\omega_0$，在线性频率轴上并不严格对称。
例如 $Q=4$ 时 $\Delta\omega/\omega_0=1/4$，两端约为 $0.883\omega_0$ 与 $1.133\omega_0$，与图中交点相符。
只有 $Q\gg1$ 的窄带极限，才可用 $\omega_{1,2}\approx\omega_0\mp\Delta\omega/2$。


\subsection{实际线圈与并联谐振}




\begin{origexamplebox}{实际线圈并联电容的相位谐振}
图~\ref{ch25:fig:real-coil} 中线圈等效为 $r$ 与 $L$ 串联，再与理想电容 $C$ 并联。求端口总导纳虚部为零的角频率，并给出存在实谐振频率的条件。
\begin{center}
\begin{tikzpicture}[scale=.95]
  \draw[ground] (0,0)--(0,2.8)--(1,2.8);
  \draw[ground] (0,0)--(1,0);
  \draw[ground] (6,0)--(7,0)--(7,2.8)--(6,2.8);
  \draw[thick] (1,2.8)
   .. controls (1,3.25) and (1.5,3.25) .. (1.5,2.8)
   .. controls (1.5,3.25) and (2,3.25) .. (2,2.8)
   .. controls (2,3.25) and (2.5,3.25) .. (2.5,2.8)
   .. controls (2.5,3.25) and (3,3.25) .. (3,2.8);
  \node[clearlabel,above] at (2,3.3) {$L$};
  \draw[thick] (3,2.6) rectangle (4.5,3.0); \node[clearlabel] at (3.75,2.8) {$r$};
  \draw[ground] (4.5,2.8)--(6,2.8);
  \draw[ground] (1,0)--(3.2,0);
  \draw[thick] (3.2,-.45)--(3.2,.45); \draw[thick] (3.55,-.45)--(3.55,.45);
  \node[clearlabel,below] at (3.38,-.45) {$C$};
  \draw[ground] (3.55,0)--(6,0);
  \node[point,label={[clearlabel]left:$A$}] at (0,1.4) {};
  \node[point,label={[clearlabel]right:$B$}] at (7,1.4) {};
\end{tikzpicture}
\figcap{含串联损耗电阻 $r$ 的实际线圈与理想电容并联。}\label{ch25:fig:real-coil}
\end{center}
\end{origexamplebox}
\begin{xesolution}
线圈支路导纳
\[
 Y_L=\frac1{r+\ii\omega L}
 =\frac{r-\ii\omega L}{r^2+\omega^2L^2},
\]
总导纳
\[
 Y=Y_L+\ii\omega C.
\]
令虚部为零：
\[
 \omega C-\frac{\omega L}{r^2+\omega^2L^2}=0.
\]
排除直流 $\omega=0$ 后，
\[
 \boxed{\omega_r^2=\frac1{LC}-\frac{r^2}{L^2}}.
\]
存在非零实频率要求 $\boxed{r^2<L/C}$。当 $r\to0$ 时回到 $1/\sqrt{LC}$；损耗增大使相位谐振频率下降。
当 $r^2\ge L/C$ 时不存在非零的相位谐振频率，这不等于证明阻抗幅值的峰也消失；峰值位置应另对 $|Z(\omega)|$ 求极值。
在 $r>0$ 的相位谐振点，总电导为 $rC/L$，故端口阻抗 $L/(rC)$ 为正实数。
\end{xesolution}


% EM teaching enrichment 20260927 supplementary

\paragraph{解出频率后仍要判断是否存在}
相位谐振必须有正实角频率。若求出的 $\omega^2$ 为负，说明给定损耗下容性支路无法在任何正频率处
恰好抵消线圈支路的电纳；不能把负数取绝对值后继续求根。
反过来，当线圈损耗趋于零，结果应回到理想并联条件 $\omega^2LC=1$。
判断相位谐振是否存在、求谐振频率、求电源电流最小点是三个不同步骤，题目要求哪个就完成哪个，不能相互替代。


% EM structure and teaching revision 20260928
\subsection{相位谐振、幅值峰值与参数依赖}
串联理想RLC接恒幅电压源时，$|I|=U/\sqrt{R^2+(\omega L-1/(\omega C))^2}$，电抗相消使电流最大且输入电压、电流同相。此处阻尼只来自与频率无关的串联电阻，所以两种谐振判据重合。
含有损耗线圈的并联电路则未必如此。线圈为 $r+\ii\omega L$，与电容并联时
\[
 Y=\frac{r}{r^2+\omega^2L^2}
 +\ii\left(\omega C-\frac{\omega L}{r^2+\omega^2L^2}\right).
\]
令虚部为零可求相位谐振频率，但实部本身也随频率变化，因此 $|Y|$ 的极小值不必出现在同一点。若问电源电流最小，应直接极小化 $|Y|$；若问电压电流同相，才使用 $\Im Y=0$。
比较谐振曲线还必须固定激励方式。恒压源下看电流响应，恒流源下看电压响应，峰谷可能互换。品质因数很高时一些频率差别很小，可以作近似，但应先给出近似参数，例如 $r\ll\omega L$，再说明忽略的是哪一阶修正。



\section{交流功率、互感与变压器}

\subsection{交流功率}


令时间原点使电流初相为零，则 $i=\sqrt2 I\cos\omega t$、
$u=\sqrt2 U\cos(\omega t+\varphi)$。用积化和差公式，瞬时功率为
\begin{equation}
 p(t)=u(t)i(t)=UI[\cos\varphi+\cos(2\omega t+\varphi)].
\end{equation}
第二项在整数个周期内的平均值为零，所以周期平均功率为
\begin{equation}
 \overline P=UI\cos\varphi.
\end{equation}
$\cos\varphi$ 称功率因数。理想电感和电容的平均功率为零，但瞬时功率并非始终为零：它们在部分周期吸收能量，
另一些时段把能量送回电源。电阻消耗的平均功率为 $I^2R$。

复功率可记为 $\widetilde S=\widetilde U\widetilde I^*=P+\ii Q_r$，其中 $P$ 为有功功率，$Q_r$ 为无功功率；
采用有效值相量时这里没有 $1/2$ 因子。展开相角可得
\begin{equation}
 P=UI\cos\varphi=I^2R,\qquad Q_r=UI\sin\varphi=I^2X,
 \qquad |\widetilde S|=UI=I^2|Z|=\sqrt{P^2+Q_r^2}.
\end{equation}
$|\widetilde S|$ 称视在功率。三者具有相同物理量纲，但为区分用途，有功功率用瓦特（W），无功功率用乏（var），视在功率用伏安（VA）。
感性负载 $Q_r>0$，容性负载 $Q_r<0$；无功功率不表示每周期永久消耗的能量。
也可只讨论其大小 $|Q_r|=P\tan|\varphi|$。对被动元件定义损耗角
$\epsilon=\pi/2-|\varphi|$，则 $\tan\epsilon=R/|X|=P/|Q_r|$；理想电抗元件 $\epsilon=0$，纯电阻对应其 $\epsilon\to\pi/2$ 极限。


% EM teaching enrichment 20260927
\paragraph{用同相与正交分量理解补偿}
对感性负载，以电压为参考，把电流分成同相的 $I_P=P/U$ 与滞后正交分量的大小 $I_Q=Q_r/U>0$，即 $\widetilde I=I_P-\ii I_Q$。理想并联电容提供相反的正交电流，可降低电源电流，却不改变原负载在相同电压下消耗的有功功率。
若 $P=800\,\mathrm W$、$U=200\,\mathrm V$、原功率因数为 $0.8$，则电流为 $5\,\mathrm A$，其中同相分量 $4\,\mathrm A$、正交分量 $3\,\mathrm A$。完全补偿后电源电流降至 $4\,\mathrm A$。
同一线路电阻上的损耗由 $25R$ 降到 $16R$；减少的是上游输电损耗，不能说负载所需的 $800\,\mathrm W$ 被电容“节省”了。


\subsection{变压器与输电}


理想变压器由同一交变磁通穿过两组线圈，法拉第定律（Faraday's law）给出
\begin{equation}
 \frac{U_1}{U_2}=\frac{N_1}{N_2},
 \qquad U_1I_1=U_2I_2.
\end{equation}
升高输电电压可在传输相同功率时降低电流，从而把线路损耗 $I^2R$ 降低。真实变压器还存在绕组电阻、漏磁、
涡流和磁滞损耗；直流稳态不能在副边维持感应电压。


% EM teaching enrichment 20260927
\paragraph{把负载折算到原边}
令匝数比 $a=N_1/N_2$，副边接负载 $Z_L$。理想变压器的电压幅值比为 $a$，电流幅值比为 $1/a$，按一致端口参考方向，原边所见阻抗为
\[
 Z_{\rm in}=a^2 Z_L.
\]
匝数比平方源于一次电压变换与一次电流变换，不能只把阻抗乘 $a$。
若传输有功功率 $P$、功率因数为 $\cos\varphi$，输电电流 $I=P/(U\cos\varphi)$，线路总电阻 $R_{\rm line}$ 上损耗为
$P^2R_{\rm line}/(U^2\cos^2\varphi)$。提高电压或改善功率因数都可降低电流，但工程装置的绝缘、损耗与额定值仍限定实际工作范围。

\begin{origexamplebox}{功率因数校正}
一线性感性负载在角频率 $\omega$、有效值 $U$ 的正弦电源上消耗有功功率 $P>0$，原功率因数为 $\cos\varphi$，其中 $0<\varphi<\pi/2$。在负载两端并联理想电容，使总电流与电压同相。求所需电容。
\end{origexamplebox}
\begin{xesolution}
以电压相量为实轴，负载电流为 $\widetilde I_{\rm load}=I\cos\varphi-\ii I\sin\varphi$。
由于 $P=UI\cos\varphi$，同相分量为 $I_P=P/U$，滞后无功分量的大小为
\[
 I_Q=I_P\tan\varphi=\frac PU\tan\varphi.
\]
电容电流超前电压 $\pi/2$，大小 $I_C=\omega CU$。校正到同相要求 $I_C=I_Q$，所以
\[
 \boxed{C=\frac{P\tan\varphi}{\omega U^2}}.
\]
这与图~\ref{ch25:fig:parallel-phasors} 中用向上的容性分量抵消向下的感性分量相同。
并联电容不改变负载端电压和有功功率；电源电流由 $P/(U\cos\varphi)$ 降至 $P/U$。
$\varphi\to0$ 时 $C\to0$；$P/(\omega U^2)$ 的单位为 $\mathrm{A\,s/V=F}$。
\end{xesolution}


\subsection{互感}





\begin{origexamplebox}{互感线圈中的感应电压}
两线圈互感大小为 $M>0$，端点和绕向如图~\ref{ch25:fig:mutual-inductance}。电流
$i_1=I_0\cos\omega t$ 从线圈1同名端流入，线圈2开路。求 $u_2=V_a-V_b$，并作量纲检查。
\begin{center}
\begin{tikzpicture}[scale=.9,>=Latex]
  \foreach \x/\name in {0/N_1,4/N_2}{
    \draw[thick] (\x,1.75)--(\x,1)
      .. controls (\x+.65,1) and (\x+.65,.5) .. (\x,.5)
      .. controls (\x+.65,.5) and (\x+.65,0) .. (\x,0)
      .. controls (\x+.65,0) and (\x+.65,-.5) .. (\x,-.5)
      .. controls (\x+.65,-.5) and (\x+.65,-1) .. (\x,-1)
      --(\x,-1.75);
    \fill (\x,1.75) circle (2pt);
    \fill (\x,-1.75) circle (2pt);
    \fill[DeepBlue] (\x-.2,.9) circle (2.3pt);
    \node[anchor=east] at (\x-.55,0) {$\name$};
  }
  \draw[very thick,TextGray] (2,-1.2)--(2,1.2) (2.2,-1.2)--(2.2,1.2);
  \node at (2.1,1.7) {$M>0$};
  \draw[->,DeepBlue,thick] (0,1.55)--(0,1.12);
  \node[anchor=east] at (-.15,1.5) {$i_1$};
  \node[right] at (4,1.75) {$a$（同名端）};
  \node[right] at (4,-1.75) {$b$};
  \draw[<->,TextGray] (6.3,-1.5)--(6.3,1.5);
  \node[right] at (6.3,0) {$u_2=V_a-V_b$};
\end{tikzpicture}
\figcap{开路互感线圈。圆点为同名端，$i_1$ 从线圈1的同名端流入；线圈2端口不接负载。}\label{ch25:fig:mutual-inductance}
\end{center}
\end{origexamplebox}
\begin{xesolution}
图~\ref{ch25:fig:mutual-inductance} 的点号规定两绕组的相对绕向：流入线圈1同名端的电流增大时，线圈2同名端 $a$ 的电势高于 $b$。两根中间线表示共同磁芯，不是导线连接。

按同名端与关联参考方向，线圈2的互感磁链为 $\Lambda_{21}=Mi_1$，端口电压为
\[
 u_2=\frac{\dd\Lambda_{21}}{\dd t}
 =M\frac{\dd i_1}{\dd t}
 =\boxed{-\omega MI_0\sin\omega t}.
\]
法拉第定律中的感应电动势为 $-\dd\Lambda_{21}/\dd t$；沿线圈从 $a$ 走向 $b$ 的内部感应电动势，与端口电压降 $V_a-V_b$ 的参考方向相反。
因此不能把感应电动势的负号不加区分地搬到关联方向的端电压公式中。

例如 $0<t<\pi/\omega$ 时 $i_1$ 正在减小，$\dd i_1/\dd t<0$，故 $a$ 端比 $b$ 端电势低；
$t=0$ 时电流虽最大，但变化率为零，$u_2=0$。交换 $a,b$ 或交换一组绕组的同名端，结果整体反号。
幅值 $\omega MI_0$ 的单位为 $\mathrm{s^{-1}\,H\,A=V}$。
次级开路 $i_2=0$，所以输出功率 $u_2i_2=0$；有感应电压不等于已向负载输送能量。
\end{xesolution}


% EM teaching enrichment 20260927 supplementary

\paragraph{用初始变化方向核对符号}
取正电流从线圈1同名端流入，按图中端口定义，开路线圈2有 $u_2=M\,\dd i_1/\dd t$。
在 $t=0$ 处 $i_1$ 达正峰值而变化率为零，故感应电压为零；在 $0<t<\pi/(2\omega)$ 内电流下降，$u_2<0$。
若交换电压表两根引线，电压符号反转而幅值不变；若只把互感大小 $M$ 改成负数而不改变参考方向，反而会重复处理极性。


\subsection{变压器带负载}


\begin{center}
\begin{tikzpicture}[scale=.9,>=Latex]
  \draw[very thick,TextGray] (3.1,-1.3)--(3.1,1.3) (3.3,-1.3)--(3.3,1.3);
  \foreach \x/\sgn in {2.4/1,4/-1}{
    \draw[thick] (\x,1)
      .. controls (\x+\sgn*.6,1) and (\x+\sgn*.6,.5) .. (\x,.5)
      .. controls (\x+\sgn*.6,.5) and (\x+\sgn*.6,0) .. (\x,0)
      .. controls (\x+\sgn*.6,0) and (\x+\sgn*.6,-.5) .. (\x,-.5)
      .. controls (\x+\sgn*.6,-.5) and (\x+\sgn*.6,-1) .. (\x,-1);
    \fill[DeepBlue] (\x-\sgn*.2,.9) circle (2.2pt);
  }
  \draw[ground] (2.4,1)--(2.4,1.6)--(.2,1.6)--(.2,.4);
  \draw[ground] (2.4,-1)--(2.4,-1.6)--(.2,-1.6)--(.2,-.4);
  \draw[thick] (.2,0) circle (.4);
  \node at (.2,0) {$\sim$};
  \node[left] at (-.25,0) {$U_1$};
  \draw[ground] (4,1)--(4,1.6)--(6.2,1.6)--(6.2,.55);
  \draw[ground] (4,-1)--(4,-1.6)--(6.2,-1.6)--(6.2,-.55);
  \draw[thick] (6,-.55) rectangle (6.4,.55);
  \node[right] at (6.5,0) {$R_L$};
  \draw[->,DeepBlue] (1,1.6)--(1.8,1.6) node[midway,above] {$I_1$};
  \draw[->,DeepBlue] (4.5,1.6)--(5.3,1.6) node[midway,above] {$I_2$};
  \node[left] at (2.05,0) {$N_1$};
  \node[right] at (4.3,0) {$N_2$};
  \node at (5.1,-.55) {$U_2$};
\end{tikzpicture}
\figcap{理想变压器带电阻负载。负载阻抗折算到原边时乘匝数比的平方。}\label{ch25:fig:transformer-load}
\end{center}
图~\ref{ch25:fig:transformer-load} 中绕组共享同一主磁通，电压均以上方同名端为正。$I_1$ 流入原边同名端，$I_2$ 从副边同名端流出并送往负载；忽略励磁电流时 $N_1I_1=N_2I_2$。因此电压比等于匝数比，电流大小比为其倒数。

\begin{origexamplebox}{负载阻抗的原边折算}
理想变压器匝数比 $N_1:N_2=n:1$，副边接纯电阻 $R_L$。求原边输入阻抗；若原边电压有效值为 $U_1$，求负载功率。
\end{origexamplebox}
\begin{xesolution}
理想变压器满足
\[
 \frac{U_1}{U_2}=n,
 \qquad \frac{I_1}{I_2}=\frac1n.
\]
因此原边看到的阻抗
\[
 Z_{\rm in}=\frac{U_1}{I_1}
 =\frac{nU_2}{I_2/n}=\boxed{n^2R_L}.
\]
副边电压 $U_2=U_1/n$，负载功率
\[
 \boxed{P_L=\frac{U_2^2}{R_L}=\frac{U_1^2}{n^2R_L}}.
\]
它也等于原边输入功率 $U_1^2/Z_{\rm in}$，验证能量守恒。$n$ 无量纲，故折算阻抗仍以欧姆计。
\end{xesolution}

% EM teaching enrichment 20260927 supplementary

\paragraph{通过两侧功率核对折算}
当原边端电压保持不变且匝数比固定时，副边电压由匝比确定，电流随负载阻抗变化。
纯电阻负载增大时，两侧电流都减小；理想空载极限两侧负载电流均趋于零。
可用 $U_1^2/R_{\rm in}=U_2^2/R_L$ 直接检查折算电阻中的匝比是否平方。
真实变压器空载仍需励磁电流，这不在理想变压器的负载折算模型内，不能用空载电流不为零来否定上述理想关系。


% EM structure and teaching revision 20260928
\subsection{复功率守恒与无功补偿}
对线性稳态网络，逐支路的复功率相加可同时检验有功与无功的端口平衡。以下沿用前文的 $Q_r$ 表示无功功率，避免与电荷量混淆。
在感性负载两端并接电容，保持电压有效值不变时，电容提供 $Q_{r,C}=-\omega CU^2$。若负载从功率因数角 $\varphi_1$ 改善到 $\varphi_2$，所需
\[
 C=\frac{P(\tan\varphi_1-\tan\varphi_2)}{\omega U^2}.
\]
补偿减少电源与线路中的无功电流，不会把负载原来需要的有功功率消除。过度补偿会使输入变成容性，也不能用“电流越小越好”取代具体设计条件。
理想变压器保持有功和无功的复功率关系，负载折算必须乘匝数比的平方；电压按匝数比变化，而电流反比变化。若结果让无损变压器无端产生实功率，应检查电流参考方向或折算比是否倒置。



\endgroup

\begingroup
\tcbset{breakable/.default=false,after={\par\nobreak\smallskip}}
\RenewDocumentEnvironment{xesolution}{+b}{%
  \par\nobreak\smallskip\nobreak\begingroup\small
  \clubpenalties 2 10000 10000\displaywidowpenalties 2 10000 10000\noindent
  {\bfseries\color{DeepBlue}解答}\quad #1\par\endgroup\smallskip
}{}
\chapter{电磁场、电磁波与超导体}

{
\graphicspath{{/Users/dongxingxi/Documents/物理竞赛一轮资料/sample/xes-original/forty-first-section/}{/Users/dongxingxi/Documents/物理竞赛一轮资料/sample/xes-original/forty-first-section/assets/}}
\section{位移电流与麦克斯韦方程组}

% Diagrams for Lesson 10: electromagnetic fields and superconductors.
\newcommand{\maxwellOnnesContext}{%
  \historyfigure{maxwell-onnes-built-in-v1.png}%
  {左：麦克斯韦与电磁场理论；右：卡末林·昂内斯与低温超导实验。人物图用于建立历史情境，公式与物理结构见后文示意图。\label{ch26:ref:review-history-context}}}

\newcommand{\radialCurrentSphereDiagram}{%
\begin{center}
\begin{tikzpicture}[scale=.88,>=Latex]
 \draw[thick] (0,0) circle (1.7);
 \fill (0,0) circle (2.5pt);
 \foreach \a in {0,45,...,315}{\draw[vector] (\a:.8)--(\a:2.15);}
 \foreach \a in {22.5,67.5,...,337.5}{\draw[->,dashed,TextGray] (\a:.85)--(\a:1.95);}
 \node at (0,-.25) {$Q(t)$};
 \node[right] at (2.2,0) {$\vect E$};
 \node[above] at (.95,2) {$\vect j$};
 \begin{scope}[shift={(.9,.75)},rotate=25]
  \draw[DeepBlue] (0,0) ellipse (.42 and .22);
  \draw[->,DeepBlue] (.42,0) arc[start angle=0,end angle=95,x radius=.42,y radius=.22];
 \end{scope}
 \node[above] at (.8,1.08) {$\Gamma$};
 \node[below] at (0,-2.25) {球对称分布};
 \fill (5,0) circle (2pt) node[above] {$P$};
 \draw[vector] (5,-.1)--(6.7,-.1) node[right] {$\vect j$};
 \draw[->,TextGray,thick] (5,-.7)--(3.3,-.7) ;
 \node at (4.2,-1.2) {$\partial\vect D/\partial t$};
 \node[align=center] at (5,1.3) {同一点 $P$ 的两种源\\等大、反向};
 \node at (5,-2) {$\vect j+\partial\vect D/\partial t=0$};
\end{tikzpicture}
\figcap{径向泄漏电荷。右侧单独展示局部源的抵消。}\label{ch26:fig:s41-10-2}
\end{center}}

\newcommand{\displacementCurrentDiagram}{%
\begin{center}
\begin{tikzpicture}[scale=.9,>=Latex]
 % Two meridian views make the identical boundary and the deformed surface explicit.
 \foreach \xx/\tag in {0/S_1,7/S_2}{
  \begin{scope}[shift={(\xx,0)}]
   \draw[very thick] (-2.7,0)--(-.65,0) (.65,0)--(2.25,0);
   \draw[line width=2pt] (-.65,-1.1)--(-.65,1.1) (.65,-1.1)--(.65,1.1);
   \draw[->,DeepBlue,thick] (-2.6,.4)--(-1.9,.4) node[midway,above] {$I$};
   \fill[DeepBlue] (-1.7,.65) circle (2.2pt) (-1.7,-.65) circle (2.2pt);
   \node[left,DeepBlue] at (-1.4,.95) {$C$};
   \node[below] at (-.65,-1.8) {$+Q$};\node[below] at (.65,-1.8) {$-Q$};
  \end{scope}
 }
 \draw[DeepBlue,thick] (-1.7,-.65)--(-1.7,.65);
 \node[DeepBlue] at (-.6,-2.8) {$S_1$：平圆盘截线};
 \begin{scope}[shift={(7,0)}]
  \draw[DeepBlue,thick] (-1.7,.65)
   .. controls (-1.6,1.45) and (-1.2,1.7) .. (-.8,1.7)
   .. controls (-.25,1.7) and (0,1.45) .. (0,0);
  \draw[DeepBlue,thick] (-1.7,-.65)
   .. controls (-1.6,-1.45) and (-1.2,-1.7) .. (-.8,-1.7)
   .. controls (-.25,-1.7) and (0,-1.45) .. (0,0);
  \draw[->,TextGray] (-.45,.25)--(.45,.25);
  \node[right] at (1.1,.7) {$\partial\vect D/\partial t$};
  \node[DeepBlue] at (-.5,-2.8) {$S_2$：绕过左板后封口于间隙};
 \end{scope}
\end{tikzpicture}
\par\smallskip{\small\color{TextGray}同一轴截面的两种跨越面。蓝色剖面绕水平轴旋转成曲面，两个蓝点旋转成同一圆形边界 $C$；$S_2$ 不穿过极板或导线。}
\figcap{充电电容器的两个同边界跨越面。}\label{ch26:ref:review-surface-choice}
\end{center}}

\newcommand{\continuityDiagram}{%
\begin{center}
\begin{tikzpicture}[>=Latex]
 \node[draw=RuleGray,align=center,text width=8cm,inner sep=7pt] (a) at (0,1.7)
 {$\nabla\times\vect B=\mu_0\vect j$\\取散度只能得到 $\nabla\cdot\vect j=0$};
 \node[draw=DeepBlue,align=center,text width=8cm,inner sep=7pt] (b) at (0,-.3)
 {$\nabla\times\vect B=\mu_0(\vect j+\partial\vect D/\partial t)$\\
  配合 $\nabla\cdot\vect D=\rho$，得到 $\nabla\cdot\vect j=-\partial\rho/\partial t$};
 \draw[->,DeepBlue,thick] (a.south)--(b.north) node[midway,right] {补位移电流};
\end{tikzpicture}
\figcap{位移电流使环路定律与电荷连续性相容（非磁介质）。}\label{ch26:ref:review-continuity}
\end{center}}

\newcommand{\maxwellLightPanel}{%
\begin{center}
\begin{tikzpicture}[x=1cm,y=1cm,>=Latex]
  \node[draw=RuleGray,rounded corners=2pt,text width=8.8cm,
        minimum height=0.92cm,align=center,inner sep=5pt] (laws)
    at (0,1.75) {高斯定律、法拉第定律与安培--麦克斯韦定律};
  \node[draw=DeepBlue,rounded corners=2pt,text width=8.8cm,
        minimum height=1.08cm,align=center,inner sep=5pt] (wave)
    at (0,0) {在真空无源区联立，得到横向电磁波\\[-1pt]
      $\displaystyle c=\frac1{\sqrt{\mu_0\varepsilon_0}}$};
  \node[draw=DeepBlue,rounded corners=2pt,text width=8.8cm,
        minimum height=0.92cm,align=center,inner sep=5pt] (light)
    at (0,-1.82) {该速度在数值上等于真空光速 $c$：光是电磁波};
  \draw[->,DeepBlue,very thick] (laws.south)--(wave.north);
  \draw[->,DeepBlue,very thick] (wave.south)--(light.north);
\end{tikzpicture}
\figcap{由麦克斯韦方程到光的电磁理论。}\label{ch26:ref:review-maxwell-light}
\end{center}}

\newcommand{\planeWaveThreeDDiagram}{%
\begin{center}
\begin{tikzpicture}[x={(1.05cm,0cm)},y={(0cm,1cm)},z={(-0.46cm,-0.34cm)},>=Latex]
  \draw[axis] (0,0,0)--(8.25,0,0) node[right] {$x$（传播方向）};
  \draw[axis] (0,0,0)--(0,2.05,0) node[above] {$y$};
  \draw[axis] (0,0,0)--(0,0,2.4) node[below left] {$z$（朝向读者）};
  \draw[DeepBlue,very thick,domain=0:7.5,samples=120,smooth,variable=\t]
    plot ({\t},{1.18*sin(72*\t)},0);
  \draw[TextGray,very thick,domain=0:7.5,samples=120,smooth,variable=\t]
    plot ({\t},0,{1.05*sin(72*\t)});
  \foreach \t in {0.65,3.15,5.65}{
    \pgfmathsetmacro{\ey}{1.18*sin(72*\t)}
    \pgfmathsetmacro{\bz}{1.05*sin(72*\t)}
    \draw[->,DeepBlue,thick] (\t,0,0)--(\t,\ey,0);
    \draw[->,TextGray,thick] (\t,0,0)--(\t,0,\bz);
  }
  \node[DeepBlue,anchor=south west,inner sep=0.5pt] at (0.95,1.55,0.15) {$\vect E\parallel\uvec y$};
  \node[TextGray,anchor=south west,inner sep=0.5pt] at (.3,-1.45,0) {$\vect B\parallel\uvec z$};
  \draw[->,DeepBlue,line width=1.2pt] (5.55,1.72,0)--(7.4,1.72,0)
    node[midway,above] {$\vect S\parallel\vect k$};
  \node[DeepBlue,anchor=north east,inner sep=0.5pt] at (7.4,-1.72,0)
    {$\vect E\times\vect B\parallel +\uvec x$};
\end{tikzpicture}
\figcap{平面波的三维投影。电、磁场同相且与传播方向构成右手系。}\label{ch26:ref:review-wave-3d}
\end{center}}

\newcommand{\poyntingBalanceDiagram}{%
\begin{center}
\begin{tikzpicture}[>=Latex]
 \draw[draw=RuleGray,fill=SoftGray] (-2,-1.15) rectangle (2,1.15);
 \node at (0,.75) {体积 $V$};
 \draw[vector] (-1.2,0)--(.65,0) node[right] {$\vect E$};
 \draw[->,TextGray,thick] (-1.2,-.65)--(.65,-.65) node[right] {$\vect j$};
 \draw[vector] (2,.45)--(3.25,.45) node[right] {$\vect S\cdot\dd\vect A$};
 \draw[vector] (-2,0)--(-3,0);
 \node[align=center] at (0,-2)
 {场能减少率 $=$ 向外净能流 $+$ 对物质做功率\\[5pt]
 $\displaystyle-\frac{\dd}{\dd t}\int_Vu\,\dd V
 =\oint_{\partial V}\vect S\cdot\dd\vect A+\int_V\vect j\cdot\vect E\,\dd V$};
\end{tikzpicture}
\figcap{电磁场的能量流出与对物质做功。}\label{ch26:ref:review-energy-balance}
\end{center}}

\newcommand{\radiationPatternDiagram}{%
\begin{center}
\begin{tikzpicture}[scale=.83,>=Latex]
 \draw[axis] (0,-2.4)--(0,2.6) node[above] {$\vect a$};
 \draw[guide] (-2.9,0)--(2.9,0);
 \draw[DeepBlue,very thick] plot[domain=0:360,samples=240]
  ({2.6*sin(\x)^3},{2.6*cos(\x)*sin(\x)^2});
 \fill (0,0) circle (2.5pt);
 \node at (.35,-1.4) {$q$};
 \draw[->,TextGray] (0,0)--(40:2.7) node[above left] {$\uvec n$};
 \draw[->,DeepBlue,thick] (2.068,1.736)--(2.486,1.238);
 \node[above] at (3.4,2.6) {$\vect E_{\rm rad}\ (q>0)$};
 \draw (0,.7) arc[start angle=90,end angle=40,radius=.7];
 \node at (.3,1.3) {$\theta$};
 \node[align=left,anchor=west] at (3.4,.3)
 {蓝色径向长度 $\propto\dd P/\dd\Omega$\\[5pt]
  $\dd P/\dd\Omega\propto\sin^2\theta$\\[5pt]
  沿加速度轴为零，横向最强\\[5pt]
  曲线不是电荷轨迹或波前};
\end{tikzpicture}
\figcap{非相对论电荷的辐射功率角分布。}\label{ch26:ref:review-radiation-pattern}
\end{center}}

\newcommand{\radiationWavefrontDiagram}{%
\begin{center}
\begin{tikzpicture}[scale=.8,>=Latex]
 \draw[axis] (-4,0)--(4.1,0) node[right] {$r$};
 \draw[axis,dashed] (0,-3.6)--(0,3.7) node[above] {振动轴};
 \fill (0,0) circle (2.5pt);
 \node at (-.33,-.28) {$q$};
 \foreach \rad/\sgn in {.85/1,1.65/-1,2.45/1,3.25/-1}{
  \draw[dashed,TextGray] (0,0) circle (\rad);
  \foreach \a in {0,30,60,120,150,180,210,240,300,330}{
   \pgfmathsetmacro{\dx}{.32*\sgn*sin(\a)*cos(\a)/\rad}
   \pgfmathsetmacro{\dy}{-.32*\sgn*cos(\a)^2/\rad}
   \draw[->,DeepBlue,thick] ({\rad*cos(\a)-\dx},{\rad*sin(\a)-\dy})
    --({\rad*cos(\a)+\dx},{\rad*sin(\a)+\dy});
  }
 }
 \draw[<->,TextGray] (.85,-3.85)--(1.65,-3.85) node[midway,below] {$\lambda/2$};
 \draw[guide] (.85,0)--(.85,-3.85) (1.65,0)--(1.65,-3.85);
 \draw[vector] (3.3,.4)--(4.15,.4);
 \node[above] at (3.95,.5) {$\vect S$};
 \node[align=left,anchor=west] at (4.8,.2)
 {$q>0$；同一时刻，相邻圆相差半波长\\[5pt]
  箭头为横向电场，不是磁感线\\[5pt]
  大小同时体现 $\sin\theta$ 与 $1/r$\\[5pt]
  沿振动轴的电场箭头为零};
\end{tikzpicture}
\figcap{振荡电荷的远区辐射。相隔半波长的场反向，相隔一波长的场同向。}\label{ch26:fig:s41-10-7}
\end{center}}

\newcommand{\thomsonScatteringDiagram}{%
\begin{center}
\begin{tikzpicture}[scale=.85,>=Latex]
 \draw[axis] (-4,0)--(4,0) node[right] {$x$};
 \draw[DeepBlue,thick] plot[domain=-3.9:-1.3,samples=70] (\x,{.45*sin(180*\x)});
 \draw[vector] (-3.7,1.2)--(-2.3,1.2) node[midway,above] {入射波};
 \fill (0,0) circle (2.5pt);
 \node[anchor=east] at (-.45,-.28) {$e^-$};
 \draw[vector] (0,0)--(0,1.3) node[above] {$\vect E$};
 \draw[->,TextGray,thick] (0,0)--(0,-1.3) node[below] {$\vect a$};
 \foreach \a in {-50,0,50,130,180,230}{
  \draw[->,TextGray] (\a:.7)--(\a:2);
 }
 \node[align=left,anchor=west] at (4.6,.2)
 {选取 $E_z>0$ 的瞬间\\[4pt]
  $a_z=-eE_z/m<0$\\[4pt]
  横向辐射强，振动轴上无辐射\\[4pt]
  $\sigma_T=(8\pi/3)r_e^2$};
\end{tikzpicture}
\figcap{电子受场驱动的瞬时方向与再辐射。}\label{ch26:ref:review-thomson}
\end{center}}

\newcommand{\superconductivityTimeline}{%
\begin{center}
\begin{tikzpicture}[x=0.86cm,y=1cm,>=Latex]
  \draw[axis] (-0.2,0)--(13.3,0) node[right] {时间};
  \foreach \x/\year/\desc in {
    0.7/1911/{Hg，$4.2\,$K},
    3.55/1933/{迈斯纳效应},
    6.4/1987/{铜氧化物，$T_c>77\,$K},
    9.25/2001/{MgB$_2$，$39\,$K},
    12.1/2015/{高压硫氢化物}}
  {
    \fill[DeepBlue] (\x,0) circle (2.6pt);
    \draw[TextGray] (\x,0.12)--(\x,0.7);
    \node[align=center,anchor=south,font=\scriptsize,text width=2.25cm] at (\x,0.72) {\year\\\desc};
  }
  \node[TextGray,align=left,anchor=north west,font=\small,text width=13.5cm] at (0.65,-0.38)
    {这是一条概念史线索，不表示“临界温度越高，材料越容易应用”；临界磁场、电流密度与制备条件同样关键。};
\end{tikzpicture}
\figcap{超导研究的历史节点。}\label{ch26:ref:review-superconductivity-timeline}
\end{center}}

\newcommand{\meissnerComparisonDiagram}{%
\begin{center}
\begin{tikzpicture}[scale=.82,>=Latex]
 \foreach \cx/\title in {-3.6/{理想导体：冻结初始场},3.6/{迈斯纳态：排斥内部场}}{
  \draw[thick,fill=PaleBlue] (\cx,0) circle (1.15);
  \node[above] at (\cx,2.8) {\title};
 }
 \foreach \x in {-1.5,-1,-.5,0,.5,1,1.5}{
  \draw[->,TextGray] ({-3.6+\x},-2.4)--({-3.6+\x},2.4);
 }
 % Analytic field lines outside a perfectly diamagnetic cylinder:
 % psi/B0 = x*(1-a^2/(x^2+y^2)); solve by bisection before exporting.
 \draw[->,DeepBlue,thick] plot coordinates {(3.444431,-2.400000) (3.442849,-2.360000) (3.441152,-2.320000) (3.439326,-2.280000) (3.437357,-2.240000) (3.435228,-2.200000) (3.432922,-2.160000) (3.430416,-2.120000) (3.427686,-2.080000) (3.424701,-2.040000) (3.421426,-2.000000) (3.417821,-1.960000) (3.413835,-1.920000) (3.409411,-1.880000) (3.404476,-1.840000) (3.398945,-1.800000) (3.392710,-1.760000) (3.385640,-1.720000) (3.377572,-1.680000) (3.368301,-1.640000) (3.357572,-1.600000) (3.345062,-1.560000) (3.330368,-1.520000) (3.312995,-1.480000) (3.292355,-1.440000) (3.267795,-1.400000) (3.238686,-1.360000) (3.204596,-1.320000) (3.165532,-1.280000) (3.122131,-1.240000) (3.075631,-1.200000) (3.027572,-1.160000) (2.979413,-1.120000) (2.932282,-1.080000) (2.886918,-1.040000) (2.843741,-1.000000) (2.802945,-0.960000) (2.764579,-0.920000) (2.728616,-0.880000) (2.694980,-0.840000) (2.663579,-0.800000) (2.634310,-0.760000) (2.607074,-0.720000) (2.581775,-0.680000) (2.558325,-0.640000) (2.536643,-0.600000) (2.516654,-0.560000) (2.498293,-0.520000) (2.481502,-0.480000) (2.466226,-0.440000) (2.452420,-0.400000) (2.440042,-0.360000) (2.429057,-0.320000) (2.419434,-0.280000) (2.411145,-0.240000) (2.404168,-0.200000) (2.398485,-0.160000) (2.394080,-0.120000) (2.390942,-0.080000) (2.389062,-0.040000) (2.388436,0.000000) (2.389062,0.040000) (2.390942,0.080000) (2.394080,0.120000) (2.398485,0.160000) (2.404168,0.200000) (2.411145,0.240000) (2.419434,0.280000) (2.429057,0.320000) (2.440042,0.360000) (2.452420,0.400000) (2.466226,0.440000) (2.481502,0.480000) (2.498293,0.520000) (2.516654,0.560000) (2.536643,0.600000) (2.558325,0.640000) (2.581775,0.680000) (2.607074,0.720000) (2.634310,0.760000) (2.663579,0.800000) (2.694980,0.840000) (2.728616,0.880000) (2.764579,0.920000) (2.802945,0.960000) (2.843741,1.000000) (2.886918,1.040000) (2.932282,1.080000) (2.979413,1.120000) (3.027572,1.160000) (3.075631,1.200000) (3.122131,1.240000) (3.165532,1.280000) (3.204596,1.320000) (3.238686,1.360000) (3.267795,1.400000) (3.292355,1.440000) (3.312995,1.480000) (3.330368,1.520000) (3.345062,1.560000) (3.357572,1.600000) (3.368301,1.640000) (3.377572,1.680000) (3.385640,1.720000) (3.392710,1.760000) (3.398945,1.800000) (3.404476,1.840000) (3.409411,1.880000) (3.413835,1.920000) (3.417821,1.960000) (3.421426,2.000000) (3.424701,2.040000) (3.427686,2.080000) (3.430416,2.120000) (3.432922,2.160000) (3.435228,2.200000) (3.437357,2.240000) (3.439326,2.280000) (3.441152,2.320000) (3.442849,2.360000) (3.444431,2.400000)};
\draw[->,DeepBlue,thick] plot coordinates {(3.150238,-2.400000) (3.146057,-2.360000) (3.141599,-2.320000) (3.136839,-2.280000) (3.131749,-2.240000) (3.126299,-2.200000) (3.120454,-2.160000) (3.114177,-2.120000) (3.107425,-2.080000) (3.100151,-2.040000) (3.092305,-2.000000) (3.083829,-1.960000) (3.074661,-1.920000) (3.064732,-1.880000) (3.053970,-1.840000) (3.042296,-1.800000) (3.029628,-1.760000) (3.015881,-1.720000) (3.000973,-1.680000) (2.984824,-1.640000) (2.967367,-1.600000) (2.948549,-1.560000) (2.928341,-1.520000) (2.906748,-1.480000) (2.883810,-1.440000) (2.859611,-1.400000) (2.834279,-1.360000) (2.807984,-1.320000) (2.780925,-1.280000) (2.753328,-1.240000) (2.725425,-1.200000) (2.697447,-1.160000) (2.669610,-1.120000) (2.642113,-1.080000) (2.615124,-1.040000) (2.588791,-1.000000) (2.563231,-0.960000) (2.538540,-0.920000) (2.514790,-0.880000) (2.492037,-0.840000) (2.470321,-0.800000) (2.449671,-0.760000) (2.430105,-0.720000) (2.411633,-0.680000) (2.394261,-0.640000) (2.377988,-0.600000) (2.362812,-0.560000) (2.348729,-0.520000) (2.335730,-0.480000) (2.323809,-0.440000) (2.312959,-0.400000) (2.303170,-0.360000) (2.294435,-0.320000) (2.286747,-0.280000) (2.280099,-0.240000) (2.274484,-0.200000) (2.269898,-0.160000) (2.266335,-0.120000) (2.263793,-0.080000) (2.262269,-0.040000) (2.261761,0.000000) (2.262269,0.040000) (2.263793,0.080000) (2.266335,0.120000) (2.269898,0.160000) (2.274484,0.200000) (2.280099,0.240000) (2.286747,0.280000) (2.294435,0.320000) (2.303170,0.360000) (2.312959,0.400000) (2.323809,0.440000) (2.335730,0.480000) (2.348729,0.520000) (2.362812,0.560000) (2.377988,0.600000) (2.394261,0.640000) (2.411633,0.680000) (2.430105,0.720000) (2.449671,0.760000) (2.470321,0.800000) (2.492037,0.840000) (2.514790,0.880000) (2.538540,0.920000) (2.563231,0.960000) (2.588791,1.000000) (2.615124,1.040000) (2.642113,1.080000) (2.669610,1.120000) (2.697447,1.160000) (2.725425,1.200000) (2.753328,1.240000) (2.780925,1.280000) (2.807984,1.320000) (2.834279,1.360000) (2.859611,1.400000) (2.883810,1.440000) (2.906748,1.480000) (2.928341,1.520000) (2.948549,1.560000) (2.967367,1.600000) (2.984824,1.640000) (3.000973,1.680000) (3.015881,1.720000) (3.029628,1.760000) (3.042296,1.800000) (3.053970,1.840000) (3.064732,1.880000) (3.074661,1.920000) (3.083829,1.960000) (3.092305,2.000000) (3.100151,2.040000) (3.107425,2.080000) (3.114177,2.120000) (3.120454,2.160000) (3.126299,2.200000) (3.131749,2.240000) (3.136839,2.280000) (3.141599,2.320000) (3.146057,2.360000) (3.150238,2.400000)};
\draw[->,DeepBlue,thick] plot coordinates {(2.722243,-2.400000) (2.715907,-2.360000) (2.709255,-2.320000) (2.702267,-2.280000) (2.694927,-2.240000) (2.687217,-2.200000) (2.679117,-2.160000) (2.670608,-2.120000) (2.661672,-2.080000) (2.652291,-2.040000) (2.642447,-2.000000) (2.632124,-1.960000) (2.621306,-1.920000) (2.609981,-1.880000) (2.598140,-1.840000) (2.585776,-1.800000) (2.572887,-1.760000) (2.559475,-1.720000) (2.545547,-1.680000) (2.531118,-1.640000) (2.516207,-1.600000) (2.500841,-1.560000) (2.485053,-1.520000) (2.468883,-1.480000) (2.452376,-1.440000) (2.435584,-1.400000) (2.418564,-1.360000) (2.401374,-1.320000) (2.384079,-1.280000) (2.366742,-1.240000) (2.349427,-1.200000) (2.332200,-1.160000) (2.315123,-1.120000) (2.298255,-1.080000) (2.281653,-1.040000) (2.265373,-1.000000) (2.249462,-0.960000) (2.233967,-0.920000) (2.218929,-0.880000) (2.204386,-0.840000) (2.190371,-0.800000) (2.176913,-0.760000) (2.164040,-0.720000) (2.151774,-0.680000) (2.140134,-0.640000) (2.129139,-0.600000) (2.118804,-0.560000) (2.109141,-0.520000) (2.100160,-0.480000) (2.091872,-0.440000) (2.084285,-0.400000) (2.077404,-0.360000) (2.071235,-0.320000) (2.065783,-0.280000) (2.061051,-0.240000) (2.057043,-0.200000) (2.053760,-0.160000) (2.051205,-0.120000) (2.049380,-0.080000) (2.048284,-0.040000) (2.047918,0.000000) (2.048284,0.040000) (2.049380,0.080000) (2.051205,0.120000) (2.053760,0.160000) (2.057043,0.200000) (2.061051,0.240000) (2.065783,0.280000) (2.071235,0.320000) (2.077404,0.360000) (2.084285,0.400000) (2.091872,0.440000) (2.100160,0.480000) (2.109141,0.520000) (2.118804,0.560000) (2.129139,0.600000) (2.140134,0.640000) (2.151774,0.680000) (2.164040,0.720000) (2.176913,0.760000) (2.190371,0.800000) (2.204386,0.840000) (2.218929,0.880000) (2.233967,0.920000) (2.249462,0.960000) (2.265373,1.000000) (2.281653,1.040000) (2.298255,1.080000) (2.315123,1.120000) (2.332200,1.160000) (2.349427,1.200000) (2.366742,1.240000) (2.384079,1.280000) (2.401374,1.320000) (2.418564,1.360000) (2.435584,1.400000) (2.452376,1.440000) (2.468883,1.480000) (2.485053,1.520000) (2.500841,1.560000) (2.516207,1.600000) (2.531118,1.640000) (2.545547,1.680000) (2.559475,1.720000) (2.572887,1.760000) (2.585776,1.800000) (2.598140,1.840000) (2.609981,1.880000) (2.621306,1.920000) (2.632124,1.960000) (2.642447,2.000000) (2.652291,2.040000) (2.661672,2.080000) (2.670608,2.120000) (2.679117,2.160000) (2.687217,2.200000) (2.694927,2.240000) (2.702267,2.280000) (2.709255,2.320000) (2.715907,2.360000) (2.722243,2.400000)};
\draw[->,DeepBlue,thick] plot coordinates {(2.266006,-2.400000) (2.259327,-2.360000) (2.252420,-2.320000) (2.245280,-2.280000) (2.237902,-2.240000) (2.230283,-2.200000) (2.222420,-2.160000) (2.214309,-2.120000) (2.205949,-2.080000) (2.197340,-2.040000) (2.188482,-2.000000) (2.179378,-1.960000) (2.170029,-1.920000) (2.160441,-1.880000) (2.150620,-1.840000) (2.140574,-1.800000) (2.130313,-1.760000) (2.119847,-1.720000) (2.109191,-1.680000) (2.098360,-1.640000) (2.087369,-1.600000) (2.076239,-1.560000) (2.064988,-1.520000) (2.053639,-1.480000) (2.042214,-1.440000) (2.030738,-1.400000) (2.019236,-1.360000) (2.007734,-1.320000) (1.996258,-1.280000) (1.984836,-1.240000) (1.973493,-1.200000) (1.962259,-1.160000) (1.951158,-1.120000) (1.940218,-1.080000) (1.929463,-1.040000) (1.918920,-1.000000) (1.908612,-0.960000) (1.898563,-0.920000) (1.888793,-0.880000) (1.879325,-0.840000) (1.870179,-0.800000) (1.861372,-0.760000) (1.852922,-0.720000) (1.844845,-0.680000) (1.837157,-0.640000) (1.829872,-0.600000) (1.823002,-0.560000) (1.816558,-0.520000) (1.810552,-0.480000) (1.804994,-0.440000) (1.799891,-0.400000) (1.795252,-0.360000) (1.791084,-0.320000) (1.787392,-0.280000) (1.784183,-0.240000) (1.781459,-0.200000) (1.779226,-0.160000) (1.777486,-0.120000) (1.776242,-0.080000) (1.775494,-0.040000) (1.775245,0.000000) (1.775494,0.040000) (1.776242,0.080000) (1.777486,0.120000) (1.779226,0.160000) (1.781459,0.200000) (1.784183,0.240000) (1.787392,0.280000) (1.791084,0.320000) (1.795252,0.360000) (1.799891,0.400000) (1.804994,0.440000) (1.810552,0.480000) (1.816558,0.520000) (1.823002,0.560000) (1.829872,0.600000) (1.837157,0.640000) (1.844845,0.680000) (1.852922,0.720000) (1.861372,0.760000) (1.870179,0.800000) (1.879325,0.840000) (1.888793,0.880000) (1.898563,0.920000) (1.908612,0.960000) (1.918920,1.000000) (1.929463,1.040000) (1.940218,1.080000) (1.951158,1.120000) (1.962259,1.160000) (1.973493,1.200000) (1.984836,1.240000) (1.996258,1.280000) (2.007734,1.320000) (2.019236,1.360000) (2.030738,1.400000) (2.042214,1.440000) (2.053639,1.480000) (2.064988,1.520000) (2.076239,1.560000) (2.087369,1.600000) (2.098360,1.640000) (2.109191,1.680000) (2.119847,1.720000) (2.130313,1.760000) (2.140574,1.800000) (2.150620,1.840000) (2.160441,1.880000) (2.170029,1.920000) (2.179378,1.960000) (2.188482,2.000000) (2.197340,2.040000) (2.205949,2.080000) (2.214309,2.120000) (2.222420,2.160000) (2.230283,2.200000) (2.237902,2.240000) (2.245280,2.280000) (2.252420,2.320000) (2.259327,2.360000) (2.266006,2.400000)};
\draw[->,DeepBlue,thick] plot coordinates {(1.784889,-2.400000) (1.778969,-2.360000) (1.772920,-2.320000) (1.766741,-2.280000) (1.760433,-2.240000) (1.753999,-2.200000) (1.747439,-2.160000) (1.740757,-2.120000) (1.733955,-2.080000) (1.727037,-2.040000) (1.720007,-2.000000) (1.712870,-1.960000) (1.705630,-1.920000) (1.698294,-1.880000) (1.690869,-1.840000) (1.683362,-1.800000) (1.675780,-1.760000) (1.668132,-1.720000) (1.660427,-1.680000) (1.652674,-1.640000) (1.644884,-1.600000) (1.637068,-1.560000) (1.629237,-1.520000) (1.621402,-1.480000) (1.613576,-1.440000) (1.605771,-1.400000) (1.597999,-1.360000) (1.590273,-1.320000) (1.582608,-1.280000) (1.575015,-1.240000) (1.567509,-1.200000) (1.560102,-1.160000) (1.552807,-1.120000) (1.545639,-1.080000) (1.538609,-1.040000) (1.531731,-1.000000) (1.525017,-0.960000) (1.518480,-0.920000) (1.512130,-0.880000) (1.505981,-0.840000) (1.500041,-0.800000) (1.494324,-0.760000) (1.488837,-0.720000) (1.483592,-0.680000) (1.478598,-0.640000) (1.473862,-0.600000) (1.469394,-0.560000) (1.465201,-0.520000) (1.461290,-0.480000) (1.457668,-0.440000) (1.454341,-0.400000) (1.451313,-0.360000) (1.448592,-0.320000) (1.446180,-0.280000) (1.444081,-0.240000) (1.442300,-0.200000) (1.440839,-0.160000) (1.439700,-0.120000) (1.438885,-0.080000) (1.438395,-0.040000) (1.438232,0.000000) (1.438395,0.040000) (1.438885,0.080000) (1.439700,0.120000) (1.440839,0.160000) (1.442300,0.200000) (1.444081,0.240000) (1.446180,0.280000) (1.448592,0.320000) (1.451313,0.360000) (1.454341,0.400000) (1.457668,0.440000) (1.461290,0.480000) (1.465201,0.520000) (1.469394,0.560000) (1.473862,0.600000) (1.478598,0.640000) (1.483592,0.680000) (1.488837,0.720000) (1.494324,0.760000) (1.500041,0.800000) (1.505981,0.840000) (1.512130,0.880000) (1.518480,0.920000) (1.525017,0.960000) (1.531731,1.000000) (1.538609,1.040000) (1.545639,1.080000) (1.552807,1.120000) (1.560102,1.160000) (1.567509,1.200000) (1.575015,1.240000) (1.582608,1.280000) (1.590273,1.320000) (1.597999,1.360000) (1.605771,1.400000) (1.613576,1.440000) (1.621402,1.480000) (1.629237,1.520000) (1.637068,1.560000) (1.644884,1.600000) (1.652674,1.640000) (1.660427,1.680000) (1.668132,1.720000) (1.675780,1.760000) (1.683362,1.800000) (1.690869,1.840000) (1.698294,1.880000) (1.705630,1.920000) (1.712870,1.960000) (1.720007,2.000000) (1.727037,2.040000) (1.733955,2.080000) (1.740757,2.120000) (1.747439,2.160000) (1.753999,2.200000) (1.760433,2.240000) (1.766741,2.280000) (1.772920,2.320000) (1.778969,2.360000) (1.784889,2.400000)};
\draw[->,DeepBlue,thick] plot coordinates {(3.755569,-2.400000) (3.757151,-2.360000) (3.758848,-2.320000) (3.760674,-2.280000) (3.762643,-2.240000) (3.764772,-2.200000) (3.767078,-2.160000) (3.769584,-2.120000) (3.772314,-2.080000) (3.775299,-2.040000) (3.778574,-2.000000) (3.782179,-1.960000) (3.786165,-1.920000) (3.790589,-1.880000) (3.795524,-1.840000) (3.801055,-1.800000) (3.807290,-1.760000) (3.814360,-1.720000) (3.822428,-1.680000) (3.831699,-1.640000) (3.842428,-1.600000) (3.854938,-1.560000) (3.869632,-1.520000) (3.887005,-1.480000) (3.907645,-1.440000) (3.932205,-1.400000) (3.961314,-1.360000) (3.995404,-1.320000) (4.034468,-1.280000) (4.077869,-1.240000) (4.124369,-1.200000) (4.172428,-1.160000) (4.220587,-1.120000) (4.267718,-1.080000) (4.313082,-1.040000) (4.356259,-1.000000) (4.397055,-0.960000) (4.435421,-0.920000) (4.471384,-0.880000) (4.505020,-0.840000) (4.536421,-0.800000) (4.565690,-0.760000) (4.592926,-0.720000) (4.618225,-0.680000) (4.641675,-0.640000) (4.663357,-0.600000) (4.683346,-0.560000) (4.701707,-0.520000) (4.718498,-0.480000) (4.733774,-0.440000) (4.747580,-0.400000) (4.759958,-0.360000) (4.770943,-0.320000) (4.780566,-0.280000) (4.788855,-0.240000) (4.795832,-0.200000) (4.801515,-0.160000) (4.805920,-0.120000) (4.809058,-0.080000) (4.810938,-0.040000) (4.811564,0.000000) (4.810938,0.040000) (4.809058,0.080000) (4.805920,0.120000) (4.801515,0.160000) (4.795832,0.200000) (4.788855,0.240000) (4.780566,0.280000) (4.770943,0.320000) (4.759958,0.360000) (4.747580,0.400000) (4.733774,0.440000) (4.718498,0.480000) (4.701707,0.520000) (4.683346,0.560000) (4.663357,0.600000) (4.641675,0.640000) (4.618225,0.680000) (4.592926,0.720000) (4.565690,0.760000) (4.536421,0.800000) (4.505020,0.840000) (4.471384,0.880000) (4.435421,0.920000) (4.397055,0.960000) (4.356259,1.000000) (4.313082,1.040000) (4.267718,1.080000) (4.220587,1.120000) (4.172428,1.160000) (4.124369,1.200000) (4.077869,1.240000) (4.034468,1.280000) (3.995404,1.320000) (3.961314,1.360000) (3.932205,1.400000) (3.907645,1.440000) (3.887005,1.480000) (3.869632,1.520000) (3.854938,1.560000) (3.842428,1.600000) (3.831699,1.640000) (3.822428,1.680000) (3.814360,1.720000) (3.807290,1.760000) (3.801055,1.800000) (3.795524,1.840000) (3.790589,1.880000) (3.786165,1.920000) (3.782179,1.960000) (3.778574,2.000000) (3.775299,2.040000) (3.772314,2.080000) (3.769584,2.120000) (3.767078,2.160000) (3.764772,2.200000) (3.762643,2.240000) (3.760674,2.280000) (3.758848,2.320000) (3.757151,2.360000) (3.755569,2.400000)};
\draw[->,DeepBlue,thick] plot coordinates {(4.049762,-2.400000) (4.053943,-2.360000) (4.058401,-2.320000) (4.063161,-2.280000) (4.068251,-2.240000) (4.073701,-2.200000) (4.079546,-2.160000) (4.085823,-2.120000) (4.092575,-2.080000) (4.099849,-2.040000) (4.107695,-2.000000) (4.116171,-1.960000) (4.125339,-1.920000) (4.135268,-1.880000) (4.146030,-1.840000) (4.157704,-1.800000) (4.170372,-1.760000) (4.184119,-1.720000) (4.199027,-1.680000) (4.215176,-1.640000) (4.232633,-1.600000) (4.251451,-1.560000) (4.271659,-1.520000) (4.293252,-1.480000) (4.316190,-1.440000) (4.340389,-1.400000) (4.365721,-1.360000) (4.392016,-1.320000) (4.419075,-1.280000) (4.446672,-1.240000) (4.474575,-1.200000) (4.502553,-1.160000) (4.530390,-1.120000) (4.557887,-1.080000) (4.584876,-1.040000) (4.611209,-1.000000) (4.636769,-0.960000) (4.661460,-0.920000) (4.685210,-0.880000) (4.707963,-0.840000) (4.729679,-0.800000) (4.750329,-0.760000) (4.769895,-0.720000) (4.788367,-0.680000) (4.805739,-0.640000) (4.822012,-0.600000) (4.837188,-0.560000) (4.851271,-0.520000) (4.864270,-0.480000) (4.876191,-0.440000) (4.887041,-0.400000) (4.896830,-0.360000) (4.905565,-0.320000) (4.913253,-0.280000) (4.919901,-0.240000) (4.925516,-0.200000) (4.930102,-0.160000) (4.933665,-0.120000) (4.936207,-0.080000) (4.937731,-0.040000) (4.938239,0.000000) (4.937731,0.040000) (4.936207,0.080000) (4.933665,0.120000) (4.930102,0.160000) (4.925516,0.200000) (4.919901,0.240000) (4.913253,0.280000) (4.905565,0.320000) (4.896830,0.360000) (4.887041,0.400000) (4.876191,0.440000) (4.864270,0.480000) (4.851271,0.520000) (4.837188,0.560000) (4.822012,0.600000) (4.805739,0.640000) (4.788367,0.680000) (4.769895,0.720000) (4.750329,0.760000) (4.729679,0.800000) (4.707963,0.840000) (4.685210,0.880000) (4.661460,0.920000) (4.636769,0.960000) (4.611209,1.000000) (4.584876,1.040000) (4.557887,1.080000) (4.530390,1.120000) (4.502553,1.160000) (4.474575,1.200000) (4.446672,1.240000) (4.419075,1.280000) (4.392016,1.320000) (4.365721,1.360000) (4.340389,1.400000) (4.316190,1.440000) (4.293252,1.480000) (4.271659,1.520000) (4.251451,1.560000) (4.232633,1.600000) (4.215176,1.640000) (4.199027,1.680000) (4.184119,1.720000) (4.170372,1.760000) (4.157704,1.800000) (4.146030,1.840000) (4.135268,1.880000) (4.125339,1.920000) (4.116171,1.960000) (4.107695,2.000000) (4.099849,2.040000) (4.092575,2.080000) (4.085823,2.120000) (4.079546,2.160000) (4.073701,2.200000) (4.068251,2.240000) (4.063161,2.280000) (4.058401,2.320000) (4.053943,2.360000) (4.049762,2.400000)};
\draw[->,DeepBlue,thick] plot coordinates {(4.477757,-2.400000) (4.484093,-2.360000) (4.490745,-2.320000) (4.497733,-2.280000) (4.505073,-2.240000) (4.512783,-2.200000) (4.520883,-2.160000) (4.529392,-2.120000) (4.538328,-2.080000) (4.547709,-2.040000) (4.557553,-2.000000) (4.567876,-1.960000) (4.578694,-1.920000) (4.590019,-1.880000) (4.601860,-1.840000) (4.614224,-1.800000) (4.627113,-1.760000) (4.640525,-1.720000) (4.654453,-1.680000) (4.668882,-1.640000) (4.683793,-1.600000) (4.699159,-1.560000) (4.714947,-1.520000) (4.731117,-1.480000) (4.747624,-1.440000) (4.764416,-1.400000) (4.781436,-1.360000) (4.798626,-1.320000) (4.815921,-1.280000) (4.833258,-1.240000) (4.850573,-1.200000) (4.867800,-1.160000) (4.884877,-1.120000) (4.901745,-1.080000) (4.918347,-1.040000) (4.934627,-1.000000) (4.950538,-0.960000) (4.966033,-0.920000) (4.981071,-0.880000) (4.995614,-0.840000) (5.009629,-0.800000) (5.023087,-0.760000) (5.035960,-0.720000) (5.048226,-0.680000) (5.059866,-0.640000) (5.070861,-0.600000) (5.081196,-0.560000) (5.090859,-0.520000) (5.099840,-0.480000) (5.108128,-0.440000) (5.115715,-0.400000) (5.122596,-0.360000) (5.128765,-0.320000) (5.134217,-0.280000) (5.138949,-0.240000) (5.142957,-0.200000) (5.146240,-0.160000) (5.148795,-0.120000) (5.150620,-0.080000) (5.151716,-0.040000) (5.152082,0.000000) (5.151716,0.040000) (5.150620,0.080000) (5.148795,0.120000) (5.146240,0.160000) (5.142957,0.200000) (5.138949,0.240000) (5.134217,0.280000) (5.128765,0.320000) (5.122596,0.360000) (5.115715,0.400000) (5.108128,0.440000) (5.099840,0.480000) (5.090859,0.520000) (5.081196,0.560000) (5.070861,0.600000) (5.059866,0.640000) (5.048226,0.680000) (5.035960,0.720000) (5.023087,0.760000) (5.009629,0.800000) (4.995614,0.840000) (4.981071,0.880000) (4.966033,0.920000) (4.950538,0.960000) (4.934627,1.000000) (4.918347,1.040000) (4.901745,1.080000) (4.884877,1.120000) (4.867800,1.160000) (4.850573,1.200000) (4.833258,1.240000) (4.815921,1.280000) (4.798626,1.320000) (4.781436,1.360000) (4.764416,1.400000) (4.747624,1.440000) (4.731117,1.480000) (4.714947,1.520000) (4.699159,1.560000) (4.683793,1.600000) (4.668882,1.640000) (4.654453,1.680000) (4.640525,1.720000) (4.627113,1.760000) (4.614224,1.800000) (4.601860,1.840000) (4.590019,1.880000) (4.578694,1.920000) (4.567876,1.960000) (4.557553,2.000000) (4.547709,2.040000) (4.538328,2.080000) (4.529392,2.120000) (4.520883,2.160000) (4.512783,2.200000) (4.505073,2.240000) (4.497733,2.280000) (4.490745,2.320000) (4.484093,2.360000) (4.477757,2.400000)};
\draw[->,DeepBlue,thick] plot coordinates {(4.933994,-2.400000) (4.940673,-2.360000) (4.947580,-2.320000) (4.954720,-2.280000) (4.962098,-2.240000) (4.969717,-2.200000) (4.977580,-2.160000) (4.985691,-2.120000) (4.994051,-2.080000) (5.002660,-2.040000) (5.011518,-2.000000) (5.020622,-1.960000) (5.029971,-1.920000) (5.039559,-1.880000) (5.049380,-1.840000) (5.059426,-1.800000) (5.069687,-1.760000) (5.080153,-1.720000) (5.090809,-1.680000) (5.101640,-1.640000) (5.112631,-1.600000) (5.123761,-1.560000) (5.135012,-1.520000) (5.146361,-1.480000) (5.157786,-1.440000) (5.169262,-1.400000) (5.180764,-1.360000) (5.192266,-1.320000) (5.203742,-1.280000) (5.215164,-1.240000) (5.226507,-1.200000) (5.237741,-1.160000) (5.248842,-1.120000) (5.259782,-1.080000) (5.270537,-1.040000) (5.281080,-1.000000) (5.291388,-0.960000) (5.301437,-0.920000) (5.311207,-0.880000) (5.320675,-0.840000) (5.329821,-0.800000) (5.338628,-0.760000) (5.347078,-0.720000) (5.355155,-0.680000) (5.362843,-0.640000) (5.370128,-0.600000) (5.376998,-0.560000) (5.383442,-0.520000) (5.389448,-0.480000) (5.395006,-0.440000) (5.400109,-0.400000) (5.404748,-0.360000) (5.408916,-0.320000) (5.412608,-0.280000) (5.415817,-0.240000) (5.418541,-0.200000) (5.420774,-0.160000) (5.422514,-0.120000) (5.423758,-0.080000) (5.424506,-0.040000) (5.424755,0.000000) (5.424506,0.040000) (5.423758,0.080000) (5.422514,0.120000) (5.420774,0.160000) (5.418541,0.200000) (5.415817,0.240000) (5.412608,0.280000) (5.408916,0.320000) (5.404748,0.360000) (5.400109,0.400000) (5.395006,0.440000) (5.389448,0.480000) (5.383442,0.520000) (5.376998,0.560000) (5.370128,0.600000) (5.362843,0.640000) (5.355155,0.680000) (5.347078,0.720000) (5.338628,0.760000) (5.329821,0.800000) (5.320675,0.840000) (5.311207,0.880000) (5.301437,0.920000) (5.291388,0.960000) (5.281080,1.000000) (5.270537,1.040000) (5.259782,1.080000) (5.248842,1.120000) (5.237741,1.160000) (5.226507,1.200000) (5.215164,1.240000) (5.203742,1.280000) (5.192266,1.320000) (5.180764,1.360000) (5.169262,1.400000) (5.157786,1.440000) (5.146361,1.480000) (5.135012,1.520000) (5.123761,1.560000) (5.112631,1.600000) (5.101640,1.640000) (5.090809,1.680000) (5.080153,1.720000) (5.069687,1.760000) (5.059426,1.800000) (5.049380,1.840000) (5.039559,1.880000) (5.029971,1.920000) (5.020622,1.960000) (5.011518,2.000000) (5.002660,2.040000) (4.994051,2.080000) (4.985691,2.120000) (4.977580,2.160000) (4.969717,2.200000) (4.962098,2.240000) (4.954720,2.280000) (4.947580,2.320000) (4.940673,2.360000) (4.933994,2.400000)};
\draw[->,DeepBlue,thick] plot coordinates {(5.415111,-2.400000) (5.421031,-2.360000) (5.427080,-2.320000) (5.433259,-2.280000) (5.439567,-2.240000) (5.446001,-2.200000) (5.452561,-2.160000) (5.459243,-2.120000) (5.466045,-2.080000) (5.472963,-2.040000) (5.479993,-2.000000) (5.487130,-1.960000) (5.494370,-1.920000) (5.501706,-1.880000) (5.509131,-1.840000) (5.516638,-1.800000) (5.524220,-1.760000) (5.531868,-1.720000) (5.539573,-1.680000) (5.547326,-1.640000) (5.555116,-1.600000) (5.562932,-1.560000) (5.570763,-1.520000) (5.578598,-1.480000) (5.586424,-1.440000) (5.594229,-1.400000) (5.602001,-1.360000) (5.609727,-1.320000) (5.617392,-1.280000) (5.624985,-1.240000) (5.632491,-1.200000) (5.639898,-1.160000) (5.647193,-1.120000) (5.654361,-1.080000) (5.661391,-1.040000) (5.668269,-1.000000) (5.674983,-0.960000) (5.681520,-0.920000) (5.687870,-0.880000) (5.694019,-0.840000) (5.699959,-0.800000) (5.705676,-0.760000) (5.711163,-0.720000) (5.716408,-0.680000) (5.721402,-0.640000) (5.726138,-0.600000) (5.730606,-0.560000) (5.734799,-0.520000) (5.738710,-0.480000) (5.742332,-0.440000) (5.745659,-0.400000) (5.748687,-0.360000) (5.751408,-0.320000) (5.753820,-0.280000) (5.755919,-0.240000) (5.757700,-0.200000) (5.759161,-0.160000) (5.760300,-0.120000) (5.761115,-0.080000) (5.761605,-0.040000) (5.761768,0.000000) (5.761605,0.040000) (5.761115,0.080000) (5.760300,0.120000) (5.759161,0.160000) (5.757700,0.200000) (5.755919,0.240000) (5.753820,0.280000) (5.751408,0.320000) (5.748687,0.360000) (5.745659,0.400000) (5.742332,0.440000) (5.738710,0.480000) (5.734799,0.520000) (5.730606,0.560000) (5.726138,0.600000) (5.721402,0.640000) (5.716408,0.680000) (5.711163,0.720000) (5.705676,0.760000) (5.699959,0.800000) (5.694019,0.840000) (5.687870,0.880000) (5.681520,0.920000) (5.674983,0.960000) (5.668269,1.000000) (5.661391,1.040000) (5.654361,1.080000) (5.647193,1.120000) (5.639898,1.160000) (5.632491,1.200000) (5.624985,1.240000) (5.617392,1.280000) (5.609727,1.320000) (5.602001,1.360000) (5.594229,1.400000) (5.586424,1.440000) (5.578598,1.480000) (5.570763,1.520000) (5.562932,1.560000) (5.555116,1.600000) (5.547326,1.640000) (5.539573,1.680000) (5.531868,1.720000) (5.524220,1.760000) (5.516638,1.800000) (5.509131,1.840000) (5.501706,1.880000) (5.494370,1.920000) (5.487130,1.960000) (5.479993,2.000000) (5.472963,2.040000) (5.466045,2.080000) (5.459243,2.120000) (5.452561,2.160000) (5.446001,2.200000) (5.439567,2.240000) (5.433259,2.280000) (5.427080,2.320000) (5.421031,2.360000) (5.415111,2.400000)};
 \node[align=center,anchor=north] at (-3.6,-2.85) {$E=0$ 只要求初始场不变\\这里先加场、后变为理想导体};
 \node[align=center,anchor=north] at (3.6,-2.85) {平衡迈斯纳态，$B\approx0$\\屏蔽电流使磁场绕开样品};
\end{tikzpicture}
\figcap{先加磁场后冷却：理想导体与迈斯纳态的区别（长圆柱截面模型）。}\label{ch26:ref:review-meissner}
\end{center}}

\newcommand{\levitationDiagram}{%
\begin{center}
\begin{tikzpicture}[scale=.88,>=Latex]
 \draw[fill=SoftGray,thick] (-1.5,1.2) rectangle (1.5,2.05);
 \draw[thick] (0,1.2)--(0,2.05);
 \node at (-.75,1.63) {$N$};\node at (.75,1.63) {$S$};
 \draw[fill=PaleBlue,thick] (-2.2,-.6)--(-2.2,-1)
   arc[start angle=180,end angle=360,x radius=2.2,y radius=.5]--(2.2,-.6);
 \draw[fill=PaleBlue,thick] (0,-.6) ellipse (2.2 and .5);
 \node at (0,-.6) {超导圆盘};
 \draw[guide] (1.5,1.63)--(3.8,1.63);
 \fill (4.5,1.63) circle (2.5pt);
 \draw[vector] (4.5,1.63)--(4.5,2.85) node[right] {$F_{\rm m}$};
 \draw[->,TextGray,thick] (4.5,1.63)--(4.5,.41) node[right] {$mg$};
 \node at (4.5,-.1) {磁体受力图};
 \node[align=center,anchor=north] at (0,-2.05) {屏蔽电流改变磁场，磁体受向上磁力\\平衡时 $F_{\rm m}=mg$；稳定性还需考虑几何与钉扎};
\end{tikzpicture}
\figcap{超导磁悬浮的装置位置与磁体受力。}\label{ch26:ref:review-levitation}
\end{center}}

\newcommand{\superconductingRingDiagram}{%
\begin{center}
\begin{tikzpicture}[scale=.82,>=Latex]
 \foreach \cx/\title in {-4/{建立初始状态},0/{缓慢减小外场},4/{外场已撤去}}{
  \draw[DeepBlue,very thick] (\cx,0) circle (1.1);
  \node[align=center] at (\cx,-1.65) {\title};
  \foreach \dx in {-.4,0,.4}{
   \draw[TextGray] (\cx+\dx,0) circle (.09);
   \fill[TextGray] (\cx+\dx,0) circle (.025);
  }
 }
 \foreach \cx in {0,4}{
  \draw[->,DeepBlue,very thick] (\cx+1.4,0)
   arc[start angle=0,end angle=140,radius=1.4];
  \node[DeepBlue,above] at (\cx,1.55) {$I>0$};
 }
 \draw[->,TextGray] (-2.5,0)--(-1.65,0);
 \draw[->,TextGray] (1.7,0)--(2.5,0);
 \node[align=center,anchor=north,font=\small] at (0,-2.2)
 {俯视：$\odot$ 为出纸面正磁通；外磁通减少，逆时针环流产生同向自感磁通\\[5pt]
  左图为初始总场，中图为外场与自感场之和，右图仅为自感场};
\end{tikzpicture}
\figcap{宏观超导环的外场撤去过程。}\label{ch26:ref:review-ring-hold}
\end{center}}

% Historical comparison; values are approximate, not a precision dataset.
\newcommand{\superconductivityChartFaithful}{%
\begin{center}
\resizebox{.98\linewidth}{!}{%
\begin{tikzpicture}[x=.019cm,y=.62cm,>=Latex]
\begin{scope}[shift={(0,0)}]
\draw[axis] (0,.5)--(230,.5) node[right] {$T_c/\mathrm K$};
\foreach \t in {0,50,100,150,200}{
 \draw[TextGray] (\t,.5)--(\t,.65) node[above,font=\scriptsize] {$\t$};
 \draw[TextGray!30,dotted] (\t,.45)--(\t,-20.5);
}
\node[anchor=east,font=\scriptsize,inner sep=2pt] at (0,0) {Hg (1911)};
\draw[DeepBlue,line width=1pt] (0,0)--(4.2,0);
\fill[DeepBlue] (4.2,0) circle (1.8pt);
\node[anchor=west,font=\scriptsize] at (7.2,0) {4.2};
\node[anchor=east,font=\scriptsize,inner sep=2pt] at (0,-1) {Pb (1913)};
\draw[DeepBlue,line width=1pt] (0,-1)--(7.2,-1);
\fill[DeepBlue] (7.2,-1) circle (1.8pt);
\node[anchor=west,font=\scriptsize] at (10.2,-1) {7.2};
\node[anchor=east,font=\scriptsize,inner sep=2pt] at (0,-2) {Nb (1930)};
\draw[DeepBlue,line width=1pt] (0,-2)--(9.2,-2);
\fill[DeepBlue] (9.2,-2) circle (1.8pt);
\node[anchor=west,font=\scriptsize] at (12.2,-2) {9.2};
\node[anchor=east,font=\scriptsize,inner sep=2pt] at (0,-3) {NbN (1941)};
\draw[DeepBlue,line width=1pt] (0,-3)--(16,-3);
\fill[DeepBlue] (16,-3) circle (1.8pt);
\node[anchor=west,font=\scriptsize] at (19,-3) {16};
\node[anchor=east,font=\scriptsize,inner sep=2pt] at (0,-4) {V$_3$Si (1953)};
\draw[DeepBlue,line width=1pt] (0,-4)--(17,-4);
\fill[DeepBlue] (17,-4) circle (1.8pt);
\node[anchor=west,font=\scriptsize] at (20,-4) {17};
\node[anchor=east,font=\scriptsize,inner sep=2pt] at (0,-5) {Nb$_3$Sn (1954)};
\draw[DeepBlue,line width=1pt] (0,-5)--(18,-5);
\fill[DeepBlue] (18,-5) circle (1.8pt);
\node[anchor=west,font=\scriptsize] at (21,-5) {18};
\node[anchor=east,font=\scriptsize,inner sep=2pt] at (0,-6) {Nb$_3$Ge (1973)};
\draw[DeepBlue,line width=1pt] (0,-6)--(23,-6);
\fill[DeepBlue] (23,-6) circle (1.8pt);
\node[anchor=west,font=\scriptsize] at (26,-6) {23};
\node[anchor=east,font=\scriptsize,inner sep=2pt] at (0,-7) {CeCu$_2$Si$_2$ (1979)};
\draw[DeepBlue,line width=1pt] (0,-7)--(0.6,-7);
\fill[DeepBlue] (0.6,-7) circle (1.8pt);
\node[anchor=west,font=\scriptsize] at (3.6,-7) {0.6};
\node[anchor=east,font=\scriptsize,inner sep=2pt] at (0,-8) {UBe$_{13}$ (1983)};
\draw[DeepBlue,line width=1pt] (0,-8)--(0.9,-8);
\fill[DeepBlue] (0.9,-8) circle (1.8pt);
\node[anchor=west,font=\scriptsize] at (3.9,-8) {0.9};
\node[anchor=east,font=\scriptsize,inner sep=2pt] at (0,-9) {UPt$_3$ (1984)};
\draw[DeepBlue,line width=1pt] (0,-9)--(0.5,-9);
\fill[DeepBlue] (0.5,-9) circle (1.8pt);
\node[anchor=west,font=\scriptsize] at (3.5,-9) {0.5};
\node[anchor=east,font=\scriptsize,inner sep=2pt] at (0,-10) {UPd$_2$Al$_3$ (1991)};
\draw[DeepBlue,line width=1pt] (0,-10)--(2,-10);
\fill[DeepBlue] (2,-10) circle (1.8pt);
\node[anchor=west,font=\scriptsize] at (5,-10) {2};
\node[anchor=east,font=\scriptsize,inner sep=2pt] at (0,-11) {CeCoIn$_5$ (2001)};
\draw[DeepBlue,line width=1pt] (0,-11)--(2.3,-11);
\fill[DeepBlue] (2.3,-11) circle (1.8pt);
\node[anchor=west,font=\scriptsize] at (5.3,-11) {2.3};
\node[anchor=east,font=\scriptsize,inner sep=2pt] at (0,-12) {PuCoGa$_5$ (2002)};
\draw[DeepBlue,line width=1pt] (0,-12)--(18.5,-12);
\fill[DeepBlue] (18.5,-12) circle (1.8pt);
\node[anchor=west,font=\scriptsize] at (21.5,-12) {18.5};
\node[anchor=east,font=\scriptsize,inner sep=2pt] at (0,-13) {PuRhGa$_5$ (2004)};
\draw[DeepBlue,line width=1pt] (0,-13)--(9,-13);
\fill[DeepBlue] (9,-13) circle (1.8pt);
\node[anchor=west,font=\scriptsize] at (12,-13) {9};
\node[anchor=east,font=\scriptsize,inner sep=2pt] at (0,-14) {LaBaCuO (1986)};
\draw[DeepBlue,line width=1pt] (0,-14)--(30,-14);
\fill[DeepBlue] (30,-14) circle (1.8pt);
\node[anchor=west,font=\scriptsize] at (33,-14) {30};
\node[anchor=east,font=\scriptsize,inner sep=2pt] at (0,-15) {LaSrCuO (1987)};
\draw[DeepBlue,line width=1pt] (0,-15)--(36,-15);
\fill[DeepBlue] (36,-15) circle (1.8pt);
\node[anchor=west,font=\scriptsize] at (39,-15) {36};
\node[anchor=east,font=\scriptsize,inner sep=2pt] at (0,-16) {YBaCuO (1987)};
\draw[DeepBlue,line width=1pt] (0,-16)--(93,-16);
\fill[DeepBlue] (93,-16) circle (1.8pt);
\node[anchor=west,font=\scriptsize] at (96,-16) {93};
\node[anchor=east,font=\scriptsize,inner sep=2pt] at (0,-17) {BiSrCaCuO (1988)};
\draw[DeepBlue,line width=1pt] (0,-17)--(110,-17);
\fill[DeepBlue] (110,-17) circle (1.8pt);
\node[anchor=west,font=\scriptsize] at (113,-17) {110};
\node[anchor=east,font=\scriptsize,inner sep=2pt] at (0,-18) {TlBaCaCuO (1988)};
\draw[DeepBlue,line width=1pt] (0,-18)--(125,-18);
\fill[DeepBlue] (125,-18) circle (1.8pt);
\node[anchor=west,font=\scriptsize] at (128,-18) {125};
\node[anchor=east,font=\scriptsize,inner sep=2pt] at (0,-19) {HgBaCaCuO (1993)};
\draw[DeepBlue,line width=1pt] (0,-19)--(133,-19);
\fill[DeepBlue] (133,-19) circle (1.8pt);
\node[anchor=west,font=\scriptsize] at (136,-19) {133};
\node[anchor=east,font=\scriptsize,inner sep=2pt] at (0,-20) {HgBaCaCuO，30 GPa};
\draw[DeepBlue,line width=1pt] (0,-20)--(164,-20);
\fill[DeepBlue] (164,-20) circle (1.8pt);
\node[anchor=west,font=\scriptsize] at (167,-20) {164};
\end{scope}
\begin{scope}[shift={(470,0)}]
\draw[axis] (0,.5)--(230,.5) node[right] {$T_c/\mathrm K$};
\foreach \t in {0,50,100,150,200}{
 \draw[TextGray] (\t,.5)--(\t,.65) node[above,font=\scriptsize] {$\t$};
 \draw[TextGray!30,dotted] (\t,.45)--(\t,-20.5);
}
\node[anchor=east,font=\scriptsize,inner sep=2pt] at (0,0) {HgTlBaCaCuO (1995)};
\draw[DeepBlue,line width=1pt] (0,0)--(138,0);
\fill[DeepBlue] (138,0) circle (1.8pt);
\node[anchor=west,font=\scriptsize] at (141,0) {138};
\node[anchor=east,font=\scriptsize,inner sep=2pt] at (0,-1) {BKBO (1988)};
\draw[DeepBlue,line width=1pt] (0,-1)--(30,-1);
\fill[DeepBlue] (30,-1) circle (1.8pt);
\node[anchor=west,font=\scriptsize] at (33,-1) {30};
\node[anchor=east,font=\scriptsize,inner sep=2pt] at (0,-2) {YPd$_2$B$_2$C (1994)};
\draw[DeepBlue,line width=1pt] (0,-2)--(23,-2);
\fill[DeepBlue] (23,-2) circle (1.8pt);
\node[anchor=west,font=\scriptsize] at (26,-2) {23};
\node[anchor=east,font=\scriptsize,inner sep=2pt] at (0,-3) {MgB$_2$ (2001)};
\draw[DeepBlue,line width=1pt] (0,-3)--(39,-3);
\fill[DeepBlue] (39,-3) circle (1.8pt);
\node[anchor=west,font=\scriptsize] at (42,-3) {39};
\node[anchor=east,font=\scriptsize,inner sep=2pt] at (0,-4) {Li，33 GPa (2002)};
\draw[DeepBlue,line width=1pt] (0,-4)--(16,-4);
\fill[DeepBlue] (16,-4) circle (1.8pt);
\node[anchor=west,font=\scriptsize] at (19,-4) {16};
\node[anchor=east,font=\scriptsize,inner sep=2pt] at (0,-5) {K$_3$C$_{60}$ (1991)};
\draw[DeepBlue,line width=1pt] (0,-5)--(18,-5);
\fill[DeepBlue] (18,-5) circle (1.8pt);
\node[anchor=west,font=\scriptsize] at (21,-5) {18};
\node[anchor=east,font=\scriptsize,inner sep=2pt] at (0,-6) {RbCsC$_{60}$ (1991)};
\draw[DeepBlue,line width=1pt] (0,-6)--(33,-6);
\fill[DeepBlue] (33,-6) circle (1.8pt);
\node[anchor=west,font=\scriptsize] at (36,-6) {33};
\node[anchor=east,font=\scriptsize,inner sep=2pt] at (0,-7) {Cs$_3$C$_{60}$，1.4 GPa};
\draw[DeepBlue,line width=1pt] (0,-7)--(38,-7);
\fill[DeepBlue] (38,-7) circle (1.8pt);
\node[anchor=west,font=\scriptsize] at (41,-7) {38};
\node[anchor=east,font=\scriptsize,inner sep=2pt] at (0,-8) {CNT低温报道 (2001)};
\draw[DeepBlue,line width=1pt] (0,-8)--(0.55,-8);
\fill[DeepBlue] (0.55,-8) circle (1.8pt);
\node[anchor=west,font=\scriptsize] at (3.55,-8) {0.55};
\node[anchor=east,font=\scriptsize,inner sep=2pt] at (0,-9) {CNT高温报道 (2006)};
\draw[DeepBlue,line width=1pt] (0,-9)--(12,-9);
\fill[DeepBlue] (12,-9) circle (1.8pt);
\node[anchor=west,font=\scriptsize] at (15,-9) {12};
\node[anchor=east,font=\scriptsize,inner sep=2pt] at (0,-10) {掺硼金刚石 (2004)};
\draw[DeepBlue,line width=1pt] (0,-10)--(4,-10);
\fill[DeepBlue] (4,-10) circle (1.8pt);
\node[anchor=west,font=\scriptsize] at (7,-10) {4};
\node[anchor=east,font=\scriptsize,inner sep=2pt] at (0,-11) {LaOFeP (2006)};
\draw[DeepBlue,line width=1pt] (0,-11)--(4,-11);
\fill[DeepBlue] (4,-11) circle (1.8pt);
\node[anchor=west,font=\scriptsize] at (7,-11) {4};
\node[anchor=east,font=\scriptsize,inner sep=2pt] at (0,-12) {LaOFFeAs (2008)};
\draw[DeepBlue,line width=1pt] (0,-12)--(26,-12);
\fill[DeepBlue] (26,-12) circle (1.8pt);
\node[anchor=west,font=\scriptsize] at (29,-12) {26};
\node[anchor=east,font=\scriptsize,inner sep=2pt] at (0,-13) {铁基材料高温区 (2008)};
\draw[DeepBlue,line width=1pt] (0,-13)--(43,-13);
\fill[DeepBlue] (43,-13) circle (1.8pt);
\node[anchor=west,font=\scriptsize] at (46,-13) {43};
\node[anchor=east,font=\scriptsize,inner sep=2pt] at (0,-14) {掺杂SrFFeAs (2009)};
\draw[DeepBlue,line width=1pt] (0,-14)--(56,-14);
\fill[DeepBlue] (56,-14) circle (1.8pt);
\node[anchor=west,font=\scriptsize] at (59,-14) {56};
\node[anchor=east,font=\scriptsize,inner sep=2pt] at (0,-15) {FeSe单层报道 (2015)};
\draw[DeepBlue,line width=1pt] (0,-15)--(109,-15);
\fill[DeepBlue] (109,-15) circle (1.8pt);
\node[anchor=west,font=\scriptsize] at (112,-15) {109};
\node[anchor=east,font=\scriptsize,inner sep=2pt] at (0,-16) {硫氢化物，155 GPa};
\draw[DeepBlue,line width=1pt] (0,-16)--(203,-16);
\fill[DeepBlue] (203,-16) circle (1.8pt);
\node[anchor=west,font=\scriptsize] at (206,-16) {203};
\node[anchor=east,font=\scriptsize,inner sep=2pt] at (0,-17) {NbC（补充）};
\draw[DeepBlue,line width=1pt] (0,-17)--(11,-17);
\fill[DeepBlue] (11,-17) circle (1.8pt);
\node[anchor=west,font=\scriptsize] at (14,-17) {11};
\node[anchor=east,font=\scriptsize,inner sep=2pt] at (0,-18) {Nb$_3$Al（补充）};
\draw[DeepBlue,line width=1pt] (0,-18)--(18,-18);
\fill[DeepBlue] (18,-18) circle (1.8pt);
\node[anchor=west,font=\scriptsize] at (21,-18) {18};
\node[anchor=east,font=\scriptsize,inner sep=2pt] at (0,-19) {CaC$_6$（补充）};
\draw[DeepBlue,line width=1pt] (0,-19)--(11.5,-19);
\fill[DeepBlue] (11.5,-19) circle (1.8pt);
\node[anchor=west,font=\scriptsize] at (14.5,-19) {11.5};
\end{scope}
\end{tikzpicture}}
\figcap{不同超导材料族的历史温区。年代与温度为约值；高压、掺杂和单层条件不能忽略，CNT与单层FeSe列的是历史报道温区，不能直接等同为常压块材零电阻温度。}\label{ch26:fig:s41-10-9}
\end{center}
液态冷媒提供的温区约为：液氦 $4.2\,\mathrm K$、液氢 $20\,\mathrm K$、液氮 $77\,\mathrm K$、液态四氟化碳 $145\,\mathrm K$。
铜氧化物和铁基材料的简式表示材料族，不表示未经掺杂的任意同名样品都具有图示 $T_c$。
图中 MgB$_2$、富勒烯、重费米子、碳材料及铁基体系分布于不同温区，说明新材料的发现并不总是刷新最高温度纪录；高压硫氢化物还需与常压材料分开评价。
}


电与磁最初分别由静电学、稳恒电流和电磁感应来描述。图~\ref{ch26:ref:review-maxwell-portrait} 中的麦克斯韦（James Clerk Maxwell）的关键工作，是看出这些规律必须组成
一个彼此相容的整体；而卡末林·昂内斯（Heike Kamerlingh Onnes）在低温实验中发现的超导现象，则把电磁场与物质的相互作用推到了一个
极端而清晰的情形：电阻可以消失，磁通却不能只按“零电阻导体”来理解。图~\ref{ch26:ref:review-history-context} 将这两条历史线索并列，左侧是场的理论，右侧是场与低温物质的实验。

电磁感应定律揭示了变化的磁场能够产生电场。由此自然会追问：变化的电场能否产生磁场？若原有的
安培定律只能处理稳恒电流，它在含有随时间变化的电荷与电场时是否仍然完备？麦克斯韦在尚无直接
实验依据的情况下，从理论自洽性出发补上位移电流项，使电场和磁场能够彼此激发。由此得到的电磁波
传播速度与当时测得的光速一致，他进一步提出光的电磁理论，经典电磁学的基本框架遂告建立。

这一过程值得特别注意：麦克斯韦并不是把彼此无关的四条实验定律简单抄在一起，而是要求它们同时
满足电荷守恒、空间对称性和有限传播速度。静电场的高斯定律告诉我们电荷能够成为电场的源，磁场
高斯定律则表明磁感线没有已经被发现的起点和终点；法拉第定律说明变化磁场会在周围激起涡旋电场。
唯独旧的安培定律只容纳稳恒传导电流，在时间变化问题中显得不对称。位移电流修正消除了这种
不对称，并使方程组能够在没有电荷、没有导线的真空中仍然具有随时间传播的非零解。电磁波并非额外
塞进理论的假设，而是完备方程组自身推出来的结果。

从学习方法上说，本章的难点不在记住更多公式，而在分清三种层次。第一层是实验事实，例如电荷产生
电场、电流产生磁场以及电磁感应；第二层是把事实压缩为局部方程，并检查这些方程能否同时成立；第三层
才是从方程推出波速、能流和辐射。只记“变化电场产生磁场”会漏掉为什么必须有这一项；只记四条方程
又会看不出电磁波从何而来。后文的每一步都围绕“源—场—传播—能量”这条线索展开。

\maxwellOnnesContext

\historyfigure[0.27\linewidth]{maxwell-built-in-v1.png}%
{麦克斯韦（James Clerk Maxwell，1831--1879），历史人物肖像。\label{ch26:ref:review-maxwell-portrait}}\label{ch26:fig:s41-10-1}


\subsection{位移电流：补上安培定律的缺口}


本段限定真空或磁导率近似为 $\mu_0$ 的非磁介质，$\vect j$ 为自由传导电流密度，$\rho$ 为自由电荷密度。含磁化介质时应使用 $\vect H$，不能仍把任意介质中的 $\vect B$ 直接除以 $\mu_0$。
稳恒电流产生磁场时，安培环路定律为
\begin{equation}
  \oint_C\vect B\cdot\dd\vect l=\mu_0 I_{\rm c},
\end{equation}
其中 $I_{\rm c}=\int_S\vect j\cdot\dd\vect S$ 是穿过跨越面 $S$ 的传导电流。若电流处处稳恒，选择
同一边界 $C$ 的不同跨越面，所得 $I_{\rm c}$ 相同，公式没有歧义。

先看第一个矛盾。设一块均匀、弱导电的球形介质中心原有电荷，电荷随时间沿径向向外泄漏。
电流密度具有严格的球对称性。若在球内任取一条闭合环路，按原来的安培环路定律，只要跨越面上
有净电流穿过，环路周围似乎就应当出现环向磁场。然而球对称系统中不存在一个在每一点都非零、又能
选出固定切向的矢量场；任何非零环向磁场都会破坏球对称性。因此，单独用传导电流作磁场源会导出
相互冲突的结论。矛盾的根源是：电荷流出时，球内电荷量和由高斯定律决定的电场也同时在变化。

把这个论证写得更定量一些。若半径 $r$ 的球面内电荷为 $Q_r(t)$，球对称性给出
\begin{equation}
  4\pi r^2D_r=Q_r(t).
\end{equation}
设向外传导电流为 $I_r$，电荷守恒要求 $I_r=-\dd Q_r/\dd t$；另一方面，穿过同一球面的位移电流为
\begin{equation}
  I_{\rm d}=\frac{\dd}{\dd t}(4\pi r^2D_r)=\frac{\dd Q_r}{\dd t}=-I_r.
\end{equation}
球对称性还把这个积分结论加强为每一点的关系：
\[
 j_r=\frac{I_r}{4\pi r^2},\qquad
 \frac{\partial D_r}{\partial t}=-\frac{I_r}{4\pi r^2},
 \qquad \vect j+\frac{\partial\vect D}{\partial t}=0.
\]
因此不是只在闭合面上“总量抵消”，而是局部的两种源也相消。
无外加磁场时，对称性只允许径向磁场；再由任意同心球面磁通为零，可得 $B_r=0$。于是球对称解确实为 $\vect B=0$。

图~\ref{ch26:fig:s41-10-2} 中实线箭头表示径向电场，灰色虚线箭头表示径向传导电流。任选闭合环路 $\Gamma$ 都不能在
球对称条件下挑出唯一的环向磁场；图右把同一点的 $\vect j$ 与 $\partial\vect D/\partial t$ 单独放大，两者大小相等、方向相反。

\radialCurrentSphereDiagram

给电容器充电时，问题暴露出来：取 $S_1$ 穿过导线，有 $I_{\rm c}=I$；取 $S_2$ 从两极板之间鼓出，
它不切断任何导线，因而 $I_{\rm c}=0$。两个面有同一边界，却会给出两个不同的磁场环量。

这里真正参与定律的是“同一闭合回路的任意跨越面”。环路 $C$ 可以想成一根铁丝围成的圈，$S_1$ 与
$S_2$ 则像两张边缘都粘在铁丝上的橡皮膜。一张膜穿过导线，另一张膜向电容器间隙鼓起；二者形状不同，
边界却完全相同。若右端只计算传导电流，左端同一个 $\oint_C\vect B\cdot\dd\vect l$ 就会得到两个值，
这不是近似误差，而是定律内部失去定义。位移电流项的任务正是让右端也只依赖边界，而不依赖所选曲面。

图~\ref{ch26:ref:review-surface-choice} 中两种蓝色剖面具有相同边界；从左图变为右图时，跨越面绕开左极板而不切断导线。

\displacementCurrentDiagram

极板间虽然没有自由电荷穿越，电位移通量却在改变。麦克斯韦定义位移电流
\begin{equation}
  I_{\rm d}=\frac{\dd}{\dd t}\int_S\vect D\cdot\dd\vect S,
\end{equation}
一般介质中的形式为 $\oint_C\vect H\cdot\dd\vect l=\int_S(\vect j+\partial\vect D/\partial t)\cdot\dd\vect S$。在本段的非磁介质条件下，它写为
\begin{equation}
  \boxed{\oint_C\vect B\cdot\dd\vect l
  =\mu_0\left(\int_S\vect j\cdot\dd\vect S
  +\frac{\dd}{\dd t}\int_S\vect D\cdot\dd\vect S\right)}.
\end{equation}
真空中 $\vect D=\varepsilon_0\vect E$。平行板电容器忽略边缘效应时，板间
$E=Q/(\varepsilon_0 A)$，于是
\begin{equation}
  I_{\rm d}=\varepsilon_0\frac{\dd(EA)}{\dd t}
  =\frac{\dd Q}{\dd t}=I.
\end{equation}
因此 $S_1$ 上的传导电流与 $S_2$ 上的位移电流恰好相等，磁场环量不再依赖跨越面的选取。

若在极板之间取半径为 $r$、与极板同轴的圆形环路，并假设 $r$ 小于极板半径，则该环路所围面积中的
位移电流为 $I_{{\rm d},r}=I\pi r^2/A$。由圆柱对称性，环路上的磁感应强度大小处处相同且沿切向，故
\begin{equation}
  B(2\pi r)=\mu_0I\frac{\pi r^2}{A},
  \qquad B=\frac{\mu_0Ir}{2A}.
\end{equation}
在极板投影之外，所包围的总位移电流已等于 $I$，于是 $B=\mu_0I/(2\pi r)$。这表明充电间隙中的
磁场并非抽象的“修补量”：它同样可由环路定律给出具体空间分布，并与导线周围磁场连续衔接。

回看弱导电球的例子，穿过闭合面的向外传导电流使内部电荷减少，与此同时电通量也减少。传导电流与
位移电流的代数和满足连续性，因而不会再凭空选出一个破坏球对称性的磁场方向。两个表面上看不同的
矛盾，其共同点正是都遗漏了变化电场的磁效应。

这个修正并非只为解决电容器的特殊矛盾。局部电荷守恒要求
\begin{equation}
  \nabla\cdot\vect j=-\frac{\partial\rho}{\partial t}.
\end{equation}
若仍用 $\nabla\times\vect B=\mu_0\vect j$，两边取散度便得到
$0=\mu_0\nabla\cdot\vect j$，它只允许电荷密度不随时间变化。加入位移电流后，利用
$\nabla\cdot\vect D=\rho$，有
\begin{equation}
  0=\mu_0\left(\nabla\cdot\vect j+
  \frac{\partial}{\partial t}\nabla\cdot\vect D\right)
  =\mu_0\left(\nabla\cdot\vect j+\frac{\partial\rho}{\partial t}\right),
\end{equation}
恰好还原连续性方程。

图~\ref{ch26:ref:review-continuity} 从上向下补入位移电流后，散度方程恢复一般的电荷守恒。

\continuityDiagram

这里的“位移电流”不是极板间真的有电荷横渡真空；它是变化电场对磁场的贡献，单位与电流相同。
这个名称带有明显的历史色彩。麦克斯韦所处的时代仍常把电磁场设想为某种介质的状态，认为外场会使
假想的真空介质或真实电介质发生极化，极化状态的变化仿佛形成一种电流。今天已经不再需要这种“以太”
图景，但“位移电流”这一名称以及电位移矢量 $\vect D$ 被保留下来。现代观点强调的不是机械介质真的
流动，而是 $\partial\vect D/\partial t$ 与传导电流一样，是涡旋磁场的源。

位移电流与传导电流还应在物理机制上区分。金属导线中的传导电流伴随自由载流子的定向运动，可以在
截面两侧统计净通过的电荷；真空电容间隙中的位移电流则没有这样的电荷输运。二者之所以能写在同一个
括号中，是因为它们对磁场环量产生同样的源效应，并共同保证总电流连续。不能据此断言“真空中也有
电子从一块极板跳到另一块极板”，也不能把位移电流当作某种额外漏电。


% EM teaching enrichment 20260927
\paragraph{位移电流不是电子穿过真空间隙}
平行板充电时，导线中传导电流改变极板自由电荷，板间的变化电场对应位移电流。二者在磁场环量关系中以相同单位出现，却描述不同机制。
对于均匀线性介质，适用的宏观项为 $\partial\boldsymbol D/\partial t$，包括真空项与极化变化；若介质还漏电，则传导项也要保留。
检查一个跨越面时应分别问：有没有自由电荷穿面？穿面的电位移通量是否变化？只凭“这里没有导线”不能把安培--麦克斯韦方程的右边判为零。


\subsection{麦克斯韦方程组}


在真空的宏观描述中，电磁场的四条基本方程可写成
\begin{align}
  \nabla\cdot\vect E&=\frac{\rho}{\varepsilon_0},
  &\oint_{\partial V}\vect E\cdot\dd\vect A&=\frac{Q_V}{\varepsilon_0},
  \\
  \nabla\cdot\vect B&=0,
  &\oint_{\partial V}\vect B\cdot\dd\vect A&=0,
  \\
  \nabla\times\vect E&=-\frac{\partial\vect B}{\partial t},
  &\oint_C\vect E\cdot\dd\vect l&=-\frac{\dd}{\dd t}\int_S\vect B\cdot\dd\vect S,
  \\
  \nabla\times\vect B&=\mu_0\vect j+\mu_0\varepsilon_0
  \frac{\partial\vect E}{\partial t},
  &\oint_C\vect B\cdot\dd\vect l&=\mu_0 I_{\rm c}+\mu_0\varepsilon_0
  \frac{\dd}{\dd t}\int_S\vect E\cdot\dd\vect S.
\end{align}
左列是微分形式，描述一个点附近的场；右列是积分形式，直接联系闭合曲面或闭合回路上的整体量。
四条方程分别表达：电荷是电场的源；没有观测到孤立磁荷；变化的磁场产生涡旋电场；传导电流与
变化的电场共同产生涡旋磁场。

积分形式适合直接分析具有对称性的有限区域，微分形式则把定律落实到空间中的每一点。例如
$\nabla\cdot\vect E=\rho/\varepsilon_0$ 在直角坐标系中就是
\begin{equation}
  \frac{\partial E_x}{\partial x}+\frac{\partial E_y}{\partial y}
  +\frac{\partial E_z}{\partial z}=\frac{\rho}{\varepsilon_0},
\end{equation}
而法拉第定律包含三个按 $x\to y\to z\to x$ 循环轮换得到的分量式，例如
\begin{equation}
  \frac{\partial E_z}{\partial y}-\frac{\partial E_y}{\partial z}
  =-\frac{\partial B_x}{\partial t}.
\end{equation}
安培--麦克斯韦定律同样给出三个分量关系，例如
\begin{equation}
  \frac{\partial B_z}{\partial y}-\frac{\partial B_y}{\partial z}
  =\mu_0j_x+\mu_0\varepsilon_0\frac{\partial E_x}{\partial t}.
\end{equation}
这些式子可由积分定律分别作用在无限小的高斯盒或矩形回路上得到。中学竞赛很少要求逐分量求解
完整方程组，但看清分量间的循环结构，有助于理解后面的平面波推导不是另立的新规律。

完整写出两个旋度方程的分量结构，可以看到这种循环关系：
\begin{align}
 \frac{\partial E_z}{\partial y}-\frac{\partial E_y}{\partial z}
   &=-\frac{\partial B_x}{\partial t},\\
 \frac{\partial E_x}{\partial z}-\frac{\partial E_z}{\partial x}
   &=-\frac{\partial B_y}{\partial t},\\
 \frac{\partial E_y}{\partial x}-\frac{\partial E_x}{\partial y}
   &=-\frac{\partial B_z}{\partial t},\\
 \frac{\partial B_z}{\partial y}-\frac{\partial B_y}{\partial z}
   &=\mu_0j_x+\mu_0\varepsilon_0\frac{\partial E_x}{\partial t},\\
 \frac{\partial B_x}{\partial z}-\frac{\partial B_z}{\partial x}
   &=\mu_0j_y+\mu_0\varepsilon_0\frac{\partial E_y}{\partial t},\\
 \frac{\partial B_y}{\partial x}-\frac{\partial B_x}{\partial y}
   &=\mu_0j_z+\mu_0\varepsilon_0\frac{\partial E_z}{\partial t}.
\end{align}
每一个旋度分量都描述“垂直于某一小面的环流密度”。这也解释了为什么平面波中只需选择一对互相
垂直的电、磁分量，另外几条分量式就自动化为 $0=0$。

方程组还需要洛伦兹力定律
\begin{equation}
  \boxed{\vect F=q(\vect E+\vect v\times\vect B)}
\end{equation}
才能说明场如何作用于带电粒子。麦克斯韦方程决定“场怎样产生和传播”，洛伦兹力决定“场怎样推动
物质”，二者分工不同。

还应避免把“场的方程”和“粒子的运动方程”混为一谈。给定 $\rho,\vect j$ 与边界条件，麦克斯韦方程
求出 $\vect E,\vect B$；再把所得场代入洛伦兹力，结合牛顿定律，才能求粒子的轨迹。反过来，
粒子运动又会改变电荷、电流分布并重新产生场。完整的电磁动力学是这种相互耦合，而不是先验地把场
当成不会受带电粒子影响的固定背景。

图~\ref{ch26:ref:review-maxwell-light} 中连接三个方框的箭头分别表示联立求解和比较传播速度，不是对实验定律的简单罗列。

\maxwellLightPanel


% EM teaching enrichment 20260927
\paragraph{四条方程与材料关系各管什么}
电的通量方程约束电荷源，磁的通量方程排除磁单极源；两条旋度方程描述变化场及电流产生的环流。它们需与电荷运动方程、初值、边界条件共同使用，不是四个互不相关的算式。
例如对安培--麦克斯韦方程取散度，左侧为零，代入电的高斯定律后得到
$\partial\rho/\partial t+\nabla\cdot\boldsymbol J=0$。若删去位移电流项，就会错误地强迫任何问题都满足 $\nabla\cdot\boldsymbol J=0$，无法描述极板正在积累电荷的过程。
介质中还需给出 $\boldsymbol D(\boldsymbol E)$、$\boldsymbol B(\boldsymbol H)$ 等材料关系；这些关系不是所有材料都相同的基本场方程。


\subsection{位移电流的定量应用}


真空中的经典电磁场满足
\begin{align}
 \nabla\cdot\vect E&=\frac{\rho}{\varepsilon_0},&
 \nabla\cdot\vect B&=0,\\
 \nabla\times\vect E&=-\frac{\partial\vect B}{\partial t},&
 \nabla\times\vect B&=\mu_0\vect J+\mu_0\varepsilon_0
 \frac{\partial\vect E}{\partial t}.
\end{align}
最后一式中的 $\varepsilon_0\partial\vect E/\partial t$ 是位移电流密度。它保证电荷守恒与安培环路定律相容：
对充电电容器，无论选取跨过导线还是穿过极板间隙的曲面，环路积分都给出同一磁场。


下面从充电电流出发，确定极板间变化的电场如何产生磁场。

\begin{origexamplebox}{平行板电容器中的位移电流}
真空圆形平行板电容器极板半径为 $R$，板间距远小于 $R$，在拟稳近似下忽略边缘效应，以恒定电流 $I$ 充电。沿轴线观察时，板间电场沿出纸面方向增大，如图~\ref{ch26:fig:displacement-question}。求板间距轴线 $r<R$ 处的磁感应强度，并说明径向跨过 $r=R$ 时结果如何衔接。
% Question contract: show plate radius and observation point only; B and the 安培 loop belong in the solution.
\begin{center}
\begin{tikzpicture}[scale=.9,>=Latex]
 \draw[very thick] (0,0) circle (1.8);
 \foreach \x/\y in {0/0,.55/.55,-.55/.55,.55/-.55,-.55/-.55,1.35/0,-1.35/0,0/1.35,0/-1.35}{
   \draw[TextGray] (\x,\y) circle (.07);\fill[TextGray] (\x,\y) circle (.018);
 }
 \draw[TextGray] (0,0)--(-1.485,1.018) node[above left] {$R$};
 \draw[TextGray] (0,0)--(.743,-.743) node[midway,below left] {$r$};
 \fill[DeepBlue] (.743,-.743) circle (2pt) node[right] {$P$};
 \node[below] at (0,-2) {沿板间轴线观察};
 \node[align=left,anchor=west] at (2.3,.3)
  {$\odot$：电场沿出纸面方向增大\\[4pt]
   圆周：极板边缘的投影\\[4pt]
   $P$ 点的磁场待求};
\end{tikzpicture}
\figcap{极板半径、观察点位置及电场变化方向。}\label{ch26:fig:displacement-question}
\end{center}
\end{origexamplebox}
\begin{xesolution}
图~\ref{ch26:fig:displacement-current} 给出解答中选取的安培环路及相应磁场方向。
\begin{center}
\begin{tikzpicture}[scale=.9,>=Latex]
 \draw[very thick] (0,0) circle (1.8);
 \draw[DeepBlue,dashed,thick] (0,0) circle (1.05);
 \foreach \x/\y in {0/0,.55/.55,-.55/.55,.55/-.55,-.55/-.55,1.35/0,-1.35/0,0/1.35,0/-1.35}{
   \draw[TextGray] (\x,\y) circle (.07);\fill[TextGray] (\x,\y) circle (.018);
 }
 \draw[->,DeepBlue,thick] (1.05,0) arc[start angle=0,end angle=70,radius=1.05];
 \draw[TextGray] (0,0)--(-1.485,1.018);
 \node[above left] at (-1.485,1.018) {$R$};
 \draw[TextGray] (0,0)--(.743,-.743);
 \node[below left] at (.44,-.38) {$r$};
 \node[right,DeepBlue] at (1.08,.72) {$\vect B$};
 \node[below] at (0,-2) {沿板间轴线观察};
 \node[align=left,anchor=west] at (2.7,.5)
  {$\odot$：$\partial\vect E/\partial t$ 出纸面\\[4pt]
   虚线圆：安培环路 $C$\\[4pt]
   实线圆：极板边缘的投影\\[4pt]
   $I_{\rm d}(r)=I\,r^2/R^2$};
\end{tikzpicture}
\figcap{充电电容器中的位移电流。环路 $C$ 围成的曲面穿过极板间隙时，由变化电场补足传导电流。}\label{ch26:fig:displacement-current}
\end{center}

由 $Q=\varepsilon_0E\pi R^2$ 和 $\dd Q/\dd t=I$，有 $\varepsilon_0\partial E/\partial t=I/(\pi R^2)$。板间电场近似均匀，总位移电流等于传导电流 $I$。半径为 $r$ 的安培环路所包围的电场面积占总极板面积的 $r^2/R^2$，故
\[
 I_d(r)=I\frac{r^2}{R^2}.
\]
由安培--麦克斯韦定律
\[
 B(2\pi r)=\mu_0I_d(r),
 \qquad
 \boxed{B(r)=\frac{\mu_0Ir}{2\pi R^2}\quad(r<R)}.
\]
在同一横截面上，$r\ge R$ 的环路包围全部位移电流，故 $B(r)=\mu_0I/(2\pi r)$。两支表达式在 $r=R$ 相等，验证径向连续性；这不是把板间观察点直接移到任意一段真实导线附近。
图中的位移电流出纸面，右手拇指朝向读者，四指给出逆时针磁场方向。$r\to0$ 时 $B\to0$，轴线上不应存在优选切向；$\mu_0Ir/R^2$ 的单位为特斯拉。
\end{xesolution}


% EM structure and teaching revision 20260928
\subsection{用连续性检验位移电流项}
对安培--麦克斯韦方程取散度，利用旋度的散度为零，可得
\[
 0=\nabla\cdot\boldsymbol j+\frac{\partial}{\partial t}(\nabla\cdot\boldsymbol D)
 =\nabla\cdot\boldsymbol j+\frac{\partial\rho_f}{\partial t}.
\]
这正是自由电荷的连续性方程。位移电流不是为了给真空人为安排真实载流子，而是使变化电场、磁场与电荷守恒能够在同一组场方程中成立。
对正在充电的理想平行板电容器，取同一条闭合环路，可以选穿过导线的曲面，也可以选绕开导线、穿过极板间隙的曲面。前者包围传导电流，后者包围电位移通量的变化率，两者应给出相同环量。
若只取半径 $r$ 小于极板半径 $a$ 的同心圆回路，并忽略边缘场，极板间 $D$ 均匀，故包围的位移电流为总充电电流的 $r^2/a^2$。磁场因而随 $r$ 线性增加，而不是处处按整根导线的 $1/r$ 规律变化。曲面选择与包围比例，是这一类题目最容易混淆的两个步骤。


\section{电磁波、能流与辐射压}

\subsection{真空中的电磁波}


在没有电荷和传导电流的真空区域，$\rho=0,\vect j=0$。设一列平面波沿 $+x$ 方向传播，场具有共同相位
\begin{equation}
  \vect E=\vect E_0\cos(kx-\omega t),
  \qquad
  \vect B=\vect B_0\cos(kx-\omega t).
\end{equation}
由 $\nabla\cdot\vect E=0$ 得 $-kE_{0x}\sin(kx-\omega t)=0$，故 $E_{0x}=0$；同理由
$\nabla\cdot\vect B=0$ 得 $B_{0x}=0$。因此源自由空间中的平面电磁波是横波。

不妨选 $\vect E=E_0\cos(kx-\omega t)\uvec y$，并设
$\vect B=B_0\cos(kx-\omega t)\uvec z$。法拉第定律的 $z$ 分量给出
\begin{equation}
  \frac{\partial E_y}{\partial x}=-\frac{\partial B_z}{\partial t}
  \quad\Longrightarrow\quad
  kE_0=\omega B_0.
\end{equation}
安培--麦克斯韦定律的 $y$ 分量为
\begin{equation}
  -\frac{\partial B_z}{\partial x}
  =\mu_0\varepsilon_0\frac{\partial E_y}{\partial t}
  \quad\Longrightarrow\quad
  kB_0=\mu_0\varepsilon_0\omega E_0.
\end{equation}
消去振幅便得到
\begin{equation}
  \boxed{\frac{\omega}{k}=\frac{1}{\sqrt{\mu_0\varepsilon_0}}\equiv c},
  \qquad
  \boxed{E_0=cB_0}.
\end{equation}
这个由电磁常量预言的速度等于真空光速，是“光就是电磁波”的核心证据。更一般地，三者满足
\begin{equation}
  \boxed{\vect B=\frac{1}{\omega}\vect k\times\vect E},
  \qquad
  \vect E\perp\vect B,\quad \vect E\perp\vect k,
  \quad \vect B\perp\vect k.
\end{equation}

推导中场的振幅 $E_0,B_0$ 在方程两边被约去，这并不意味着振幅为零，而是说明真空常量只决定传播
速度及两种场的振幅比，实际振幅由波源强弱决定。频率同样由波源的振荡频率决定；一旦频率 $f$ 给定，
真空中的波长便不能任意选择，而必须满足
\begin{equation}
  \lambda=\frac{c}{f},\qquad k=\frac{2\pi}{\lambda},\qquad \omega=2\pi f.
\end{equation}
传播方向的垂直平面有两个彼此独立的偏振方向。任意线偏振可以分解成这两个方向的叠加；若两分量有
相位差，则可形成椭圆偏振或圆偏振。对这里的真空平面波，$\vect E$ 与 $\vect B$ 同相，三者按
$\vect E,\vect B,\vect k$ 的次序构成右手系。

也可以不先猜传播速度，而直接从麦克斯韦方程消去磁场。对法拉第定律取旋度，使用恒等式
$\nabla\times(\nabla\times\vect E)=\nabla(\nabla\cdot\vect E)-\nabla^2\vect E$；在真空无源区
$\nabla\cdot\vect E=0$，再代入安培--麦克斯韦定律，得到
\begin{equation}
  \nabla^2\vect E-\mu_0\varepsilon_0\frac{\partial^2\vect E}{\partial t^2}=0.
\end{equation}
磁场同理满足完全相同的波动方程。与标准形式
$\nabla^2\psi-(1/v^2)\partial^2\psi/\partial t^2=0$ 对照，立即读出
$v=1/\sqrt{\mu_0\varepsilon_0}$。这一推导说明速度结论并不依赖于选取沿 $x$ 传播、沿 $y$ 偏振的
特殊坐标；特殊坐标只是把矢量方程变成容易计算的分量式。

把平面简谐波代回波动方程可得 $\omega^2=c^2k^2$。正频率支对应沿 $+\vect k$ 方向传播，改变相位
符号则得到反向传播的波。波动方程是线性的：若两列波分别是解，它们的和仍是解，因此干涉、驻波和
偏振叠加都自然成立。所谓“振幅任意”只是在给定线性真空方程时，方程本身不替波源规定振幅；实际
系统中的天线功率、边界条件和能量守恒仍会决定可实现的场强。类似地，频率可由波源选择，但波速固定，
于是不同频率对应不同波长。可见光只是电磁波谱的一小段，改变频率并不改变其电磁波本性。

图~\ref{ch26:ref:review-wave-3d} 中电场位于蓝色平面，磁场位于灰色平面；其方向在下文逐项说明。

\planeWaveThreeDDiagram

图中蓝色曲线位于 $x$--$y$ 平面，灰色曲线位于 $x$--$z$ 平面；同一 $x$ 处的两根场矢量同相。
右手从 $\vect E$ 旋向 $\vect B$，拇指所指就是传播方向。图上的斜向 $z$ 轴只是纸面投影，不表示
磁场真的“倾斜”到传播方向。若 $\vect E$ 始终沿固定横向振动，这列波为线偏振波；若两个横向分量
存在相位差，电场矢量端点可描出椭圆或圆。


\subsection{场的能量与坡印廷矢量}


真空中电场与磁场的能量密度分别为
\begin{equation}
  u_E=\frac12\varepsilon_0E^2,
  \qquad
  u_B=\frac{B^2}{2\mu_0},
  \qquad u=u_E+u_B.
\end{equation}
为了找出能量怎样流动，将 $\vect E$ 点乘安培--麦克斯韦方程、将 $\vect B/\mu_0$ 点乘法拉第方程，
除去公共常数后，两个点乘式分别为
\begin{align}
 \frac1{\mu_0}\vect E\cdot(\nabla\times\vect B)
   &=\vect j\cdot\vect E+\frac{\partial}{\partial t}
      \left(\frac{\varepsilon_0E^2}{2}\right),\\
 \frac1{\mu_0}\vect B\cdot(\nabla\times\vect E)
   &=-\frac{\partial}{\partial t}\left(\frac{B^2}{2\mu_0}\right).
\end{align}
用第二式减去第一式，再利用恒等式
\begin{equation}
  \nabla\cdot(\vect E\times\vect B)
  =\vect B\cdot(\nabla\times\vect E)-
  \vect E\cdot(\nabla\times\vect B),
\end{equation}
可得
\begin{equation}
  \boxed{\frac{\partial u}{\partial t}+\nabla\cdot\vect S=-\vect j\cdot\vect E},
  \qquad
  \boxed{\vect S=\frac{1}{\mu_0}\vect E\times\vect B}.
\end{equation}
这就是电磁场的局部能量守恒式。$\vect S$ 称为坡印廷矢量，它的方向是能量传播方向，大小表示
单位时间穿过单位垂直面积的能量。

静电场和稳恒磁场可以把能量储存在空间中；当电场与磁场共同组成传播的电磁波时，能量还能够从一处
流向另一处。太阳并不需要把物质一路输送到地球，太阳辐射携带的能量正是由电磁场穿过空间传来的。
因此坡印廷矢量不是只为配合波动图形而引入的方向箭头，它给出了能量运输的定量描述。对于平面波，
$\vect E\times\vect B$ 指向传播方向，这也反过来帮助判断两场矢量的相对取向。

对有限体积 $V$ 积分，局部守恒式还可以写成
\begin{equation}
  \frac{\dd}{\dd t}\int_Vu\,\dd V
  =-\oint_{\partial V}\vect S\cdot\dd\vect A-
  \int_V\vect j\cdot\vect E\,\dd V.
\end{equation}
左端是体内电磁场能量的增加率；第一个负号项表示由边界流出的电磁能；第二项是场对带电物质做功的
功率。若 $\vect j\cdot\vect E>0$，场把能量交给物质，例如电阻发热；若它为负，外力驱动电荷并向场
输送能量，例如发射天线。$\vect S$ 的单位为 $\si{W/m^2}$，与“单位面积上的功率”这一物理解释一致。

图~\ref{ch26:ref:review-energy-balance} 中边界向外的箭头对应面积分，体内同向的电场、电流对应正的做功项。

\poyntingBalanceDiagram

对上面的真空平面波，$E=cB$ 且 $c^2=1/(\mu_0\varepsilon_0)$，所以
\begin{equation}
  u_E=\frac12\varepsilon_0E^2
  =\frac{B^2}{2\mu_0}=u_B.
\end{equation}
电、磁两部分能量在每一时刻都相等。瞬时能流密度为 $S=\varepsilon_0cE^2$，一个周期内的平均光强为
\begin{equation}
  \boxed{\langle S\rangle=\frac12\varepsilon_0cE_0^2
  =\frac{cB_0^2}{2\mu_0}}.
\end{equation}


% EM teaching enrichment 20260927
\paragraph{积分形式中的三个能量去向}
对固定体积 $V$，坡印廷定理写成
\[
 -\oint_{\partial V}\boldsymbol S\cdot\dd\boldsymbol A
 =\frac{\dd}{\dd t}\int_Vu\dd V+\int_V\boldsymbol J\cdot\boldsymbol E\dd V.
\]
左边是从边界净流入的功率，右边分别为内部场能增加率与场对物质的做功率。若是稳恒欧姆导体，场能不变，流入能量最终转化为焦耳热；若是真空充电电容器间隙，没有传导耗散，流入主要增加场能。
对平行板半径 $R$、间距 $d$ 的拟稳充电模型，边缘附近 $B=\mu_0I/(2\pi R)$，侧面流入功率的大小为
$(EB/\mu_0)(2\pi Rd)=IEd=IU$，与电路功率一致。此处忽略边缘修正，只用来核对能量流向和总量。


\subsection{平面波方向与相位判读}


在无电荷、无电流的真空区域，对法拉第定律再取旋度并使用
$\nabla\times(\nabla\times\vect E)=\nabla(\nabla\cdot\vect E)-\nabla^2\vect E$，得到
\begin{equation}
 \nabla^2\vect E-\mu_0\varepsilon_0\frac{\partial^2\vect E}{\partial t^2}=0.
\end{equation}
磁场满足同样的波动方程，传播速度为
\begin{equation}
 c=\frac1{\sqrt{\mu_0\varepsilon_0}}.
\end{equation}
平面波中 $\vect E$、$\vect B$ 和传播方向 $\uvec k$ 两两垂直，并满足
\begin{equation}
 \vect B=\frac1c\uvec k\times\vect E,
 \qquad B=\frac Ec.
\end{equation}

\begin{center}
\begin{tikzpicture}[x=.85cm,y=.78cm,>=Latex]
 \foreach \y/\lab in {1.65/{E_x/E_0},-1.65/{cB_y/E_0}}{
  \draw[axis] (0,\y)--(8.3,\y) node[right] {$z$};
  \draw[axis] (0,\y-1.25)--(0,\y+1.3);
  \node[anchor=east] at (-.15,\y+.8) {$\lab$};
  \draw[DeepBlue,very thick] plot[domain=0:8,samples=180] (\x,{\y+sin(90*\x)});
 }
 \foreach \x in {1,3,5,7}{\draw[guide] (\x,-2.8)--(\x,2.8);}
 \draw[->,TextGray,thick] (5,3.2)--(7,3.2) node[midway,above] {传播方向 $+\uvec z$};
\end{tikzpicture}
\figcap{同一时刻的电、磁场分量，分两行绘出以比较相位。}\label{ch26:fig:wave-phase}
\end{center}
图~\ref{ch26:fig:wave-phase} 中两行的波峰、波谷和零点位于相同的 $z$，因此两场同相。
纵坐标分别为归一化分量，而非两个空间坐标；实际方向为 $\vect E\parallel\uvec x$、
$\vect B\parallel\uvec y$，由 $\uvec x\times\uvec y=\uvec z$ 确定向前的能流。


% EM teaching enrichment 20260927
\paragraph{从一个电场分量补齐整束平面波}
给定 $\boldsymbol E=E_0\cos(kz-\omega t)\boldsymbol e_x$，等相位条件 $kz-\omega t=\text{常量}$ 表明波沿 $+z$ 传播。于是
\[
 \boldsymbol B=\frac{E_0}{c}\cos(kz-\omega t)\boldsymbol e_y,
 \qquad \omega=ck.
\]
电场与磁场同时反向时，$\boldsymbol E\times\boldsymbol B$ 仍沿 $+z$。若把相位改为 $kz+\omega t$，传播方向反向，磁场的相对方向也必须反向。
只画两条正弦曲线还不足以说明空间关系；应同时给出坐标轴、传播方向以及叉积检验。波形图纵坐标表示场分量，不是粒子在空间中的上下振动轨迹。


\subsection{能量、动量与辐射压}


电磁场能量密度与坡印廷矢量为
\begin{equation}
 u=\frac12\varepsilon_0E^2+\frac{B^2}{2\mu_0},
 \qquad \vect S=\frac1{\mu_0}\vect E\times\vect B.
\end{equation}
$\vect S$ 给出单位面积单位时间通过的能量及其方向。真空简谐平面波的平均光强
\begin{equation}
 I=\langle S\rangle=\frac12c\varepsilon_0E_0^2.
\end{equation}
正入射完全吸收表面的平均辐射压为 $I/c$，正入射理想反射表面为 $2I/c$；斜入射还需取法向动量变化。

\begin{origexamplebox}{激光束的场强和辐射压}
功率 $P$ 均匀分布在截面积 $A$ 的连续激光束中，将光束近似为单色线偏振平面波。求真空中的电场峰值；若光束正入射并被完全吸收，求平均压力。
\end{origexamplebox}
\begin{xesolution}
光强 $I=P/A$。由 $I=c\varepsilon_0E_0^2/2$，
\[
 \boxed{E_0=\sqrt{\frac{2P}{Ac\varepsilon_0}}}.
\]
峰值电场与有效值电场相差 $\sqrt2$，上式中的因子 $2$ 来自对 $\cos^2\omega t$ 取平均。
在短时间 $\Delta t$ 内，到达面积 $A$ 的能量为 $IA\Delta t$，携带向前的动量为 $IA\Delta t/c$。
完全吸收时这部分动量全部交给物体，故 $F=(IA\Delta t/c)/\Delta t=IA/c$，再除以面积得到压力
\[
 \boxed{p_{\rm rad}=\frac{P}{Ac}}.
\]
若改为理想镜面反射，法向动量反向，压力加倍。$P/(Ac)$ 的单位为 $\mathrm{N/m^2}$；
保持面积不变而使功率趋零，场强和压力也都趋零。
\end{xesolution}

\begin{center}
\begin{tikzpicture}[scale=.92,>=Latex]
  \draw[very thick] (5,-1.5)--(5,1.5);
  \foreach \y in {-1,-.5,0,.5,1}{\draw[vector] (.3,\y)--(3.5,\y);}
  \draw[->,TextGray,very thick] (5,0)--(6.5,0) node[right]{$\vect F$};
  \draw[->,DeepBlue,thick] (3.8,.85)--(4.8,.85) node[midway,above]{$\vect p_\gamma$};
  \node[clearlabel,above] at (5,1.65) {吸收面};
  \node[clearlabel] at (1.9,-1.35) {均匀光束，光强 $I$};
\end{tikzpicture}
\figcap{光束被吸收时把传播方向的动量交给物体；压力等于单位面积单位时间的法向动量通量。}\label{ch26:fig:radiation-pressure}
\end{center}
图~\ref{ch26:fig:radiation-pressure} 中力箭头与入射传播方向同向；若改成反射，光子法向动量反向，力增大一倍。

\begin{origexamplebox}{太阳辐射对完全反射薄膜的推力}
面积为 $A$ 的平面薄膜正对光强为 $I$ 的光束，薄膜作理想镜面反射且光束完全覆盖薄膜。求推力；若光束与薄膜法线夹角为 $\theta$，求法向力。
\end{origexamplebox}
\begin{xesolution}
正入射反射使单位能量携带的法向动量从 $+1/c$ 变为 $-1/c$，变化量为 $2/c$，所以
\[
 \boxed{F=\frac{2IA}{c}}.
\]
斜入射时，设 $0\le\theta\le\pi/2$，薄膜在垂直光束方向的投影面积为 $A\cos\theta$，整片薄膜截获的功率为 $IA\cos\theta$，每单位能量的法向动量变化为
$2\cos\theta/c$，故
\[
 \boxed{F_n=\frac{2IA}{c}\cos^2\theta}.
\]
$\theta=0$ 回到正入射结果，$\theta\to\pi/2$ 时截获功率和法向动量都趋零。$IA/c$ 的单位为牛顿。
\end{xesolution}


% EM structure and teaching revision 20260928
\subsection{场幅、平均强度与总功率的区别}
真空平面波中 $B_m=E_m/c$，电场与磁场能量密度相等，瞬时能流方向为 $\boldsymbol E\times\boldsymbol B$。对于单频正弦波，时间平均强度为
\[
 \langle S\rangle=\frac12\varepsilon_0cE_m^2
 =\varepsilon_0cE_{\mathrm{rms}}^2.
\]
峰值与有效值混用会造成二倍误差。总功率还需对横截面积分：$P=\int\langle\boldsymbol S\rangle\cdot\dd\boldsymbol A$。只有波束在所选截面上近似均匀并正入射时，才可简写 $P=\langle S\rangle A$。
设均匀波束以角 $\theta$ 入射平面，强度按垂直波束方向的单位面积定义，则表面每单位面积截获的能量通量为 $\langle S\rangle\cos\theta$。若镜面理想反射，法向动量反转再提供一个 $\cos\theta$，因而法向辐射压为
$p_n=2\langle S\rangle\cos^2\theta/c$。
斜入射问题若同时漏掉投影面积和动量投影，结果可能在正入射时正确，却在掠入射极限仍给出有限法向压力。用 $\theta\to\pi/2$ 作检查，能直接识别这类错误。


\section{辐射、散射与介质传播}

\subsection{加速电荷的辐射}


静止电荷产生静电场；匀速运动的电荷在固定观察点也可产生随时间变化的场，但在真空中不产生向远区输送净能量的辐射项。电磁作用始终以有限速度传播，不能把“有时间变化”和“有辐射”等同。加速度使场出现按 $1/r$ 衰减的辐射项。在非相对论条件 $v\ll c$ 下，距离电荷很远处的辐射电场大小为
\begin{equation}
  \boxed{E_{\rm rad}(r,t)=\frac{1}{4\pi\varepsilon_0}
  \frac{|q|\,a(t-r/c)\sin\theta}{c^2r}}.
\end{equation}
其中 $t-r/c$ 是推迟时刻，$\theta$ 是观察方向与加速度方向的夹角。这里用 $E_{\rm rad}$ 而不用
$E_r$：辐射场实际上垂直于传播方向，不是沿半径方向的场分量。
$a$ 在大小公式中表示加速度的模；方向应由矢量式判断：
\[
 \vect E_{\rm rad}
 =\frac{q}{4\pi\varepsilon_0c^2r}\,
 \uvec n\times(\uvec n\times\vect a)
 \quad\text{（取推迟时刻）}.
\]
利用三重积，$\uvec n\times(\uvec n\times\vect a)=
\uvec n(\uvec n\cdot\vect a)-\vect a$，恰好是加速度横向分量的反向；
电荷为负时场方向再反转。

辐射场随 $1/r$ 衰减，而库仑近场随 $1/r^2$ 衰减；所以在足够远处，辐射项反而占主导。远区有
$B_{\rm rad}=E_{\rm rad}/c$，于是单位立体角内的辐射功率为
\begin{equation}
  \frac{\dd P}{\dd\Omega}=r^2S
  =\frac{q^2a^2\sin^2\theta}{16\pi^2\varepsilon_0c^3}.
\end{equation}

图~\ref{ch26:ref:review-radiation-pattern} 中角度 $\theta$ 从加速度轴量起，蓝色曲线的径向长度代表功率而不是电场幅值。
图中另画出的 $\vect E_{\rm rad}$ 箭头对应 $q>0$：它垂直于观察方向 $\uvec n$，
沿 $\uvec n\times(\uvec n\times\vect a)$ 指向右下方；若电荷为负，电场方向反转而功率角分布不变。

\radiationPatternDiagram

沿加速度方向 $\theta=0,\pi$ 时不辐射，垂直于加速度时最强。对全空间积分，利用
\begin{equation}
  \int\sin^2\theta\,\dd\Omega
  =2\pi\int_0^\pi\sin^3\theta\,\dd\theta=\frac{8\pi}{3},
\end{equation}
得到拉莫尔公式
\begin{equation}
  \boxed{P=\frac{q^2a^2}{6\pi\varepsilon_0c^3}}.
\end{equation}
它说明辐射功率正比于加速度平方；这里的公式只适用于非相对论点电荷，不能直接套用于
$v$ 接近 $c$ 的情形。

图~\ref{ch26:fig:s41-10-7} 用一组相隔半波长的波前表示振荡电荷向外传递的场。每个蓝箭头都沿波前切向，因而与能流
方向垂直；相邻半波带的箭头反向，对应简谐场的相位翻转。沿加速度轴箭头趋于零，与
$E_{\rm rad}\propto\sin\theta$ 一致。

\radiationWavefrontDiagram

辐射图样常画成绕加速度轴的“甜甜圈”。图中的径向长度表示该方向的相对辐射强弱，并不是电荷真的
喷出一圈物质。把 $E_{\rm rad}$ 代入 $S=\varepsilon_0cE_{\rm rad}^2$，其中 $E_{\rm rad}^2\propto1/r^2$；
再乘球面面积元 $r^2\dd\Omega$，距离因子正好消去。因此无论在多大的球面上统计，只要球面包围波源，
在对应同一推迟时刻的波前上，穿出的总辐射功率都相同。若加速度随时间变化，不同半径处同一观察时刻所收到的是不同发射时刻的能量，不能直接要求瞬时功率相等。这是能量从波源持续向外传播而不是在空间中凭空损失的表现。

推迟时刻 $t-r/c$ 是有限传播速度的直接结果：观察者接收到的不是电荷“此刻”的运动信息，而是场从
电荷传播到观察点之前的状态，正如地球上看到的太阳是约八分钟以前的太阳。拉莫尔公式曾是经典
电磁学的重大成功，却也暴露出经典原子模型的根本困难：若电子像行星一样绕原子核转动，它始终有
向心加速度，便应不断辐射并损失机械能，轨道逐渐缩小直至落入原子核。真实原子却稳定存在，这个冲突
不能靠经典电磁学自身解决，后来成为量子理论产生的重要背景之一。

还要分清近场与辐射场。靠近电荷时，$1/r^2$ 的库仑型项以及与速度有关的感应项可能占主导，
它们储存的能量可在波源附近往返交换；远离波源后，按 $1/r$ 衰减的辐射项成为主导，并携带净能量
离开系统。因此“有时变电场”不等于“观察点处一定只有辐射场”，判断远区辐射必须同时看距离依赖和
能流方向。拉莫尔公式计算的是已经离开波源、穿过远处闭合面的总功率。


% EM teaching enrichment 20260927
\paragraph{近场与辐射场的距离检验}
库仑型近场随 $1/r^2$ 衰减，非相对论辐射场随 $1/r$ 衰减。远区能流与场强平方成正比，球面面积又正比于 $r^2$，因此只有 $1/r$ 的场项能在远处留下有限的总辐射功率。
沿加速度轴方向，辐射场的横向投影为零；垂直轴方向最强。图中的 $\sin^2\theta$ 表示每单位立体角的功率分布，不是电场振幅，振幅对应 $\sin\theta$。
若同一非相对论点电荷的加速度大小加倍，拉莫尔功率变为四倍；恒速直线运动的加速度为零，不因观察点处近场随时间改变就产生同样的辐射项。


\subsection{汤姆孙散射截面}


考虑沿 $+x$ 传播、沿 $z$ 偏振的电磁波，对初始平均静止的自由电子作弱场、非相对论近似
\begin{equation}
  E_z=E_0\cos(kx-\omega t).
\end{equation}
电子的振动速度幅值约为 $v_0=eE_0/(m\omega)$。要求 $eE_0/(m\omega c)\ll1$，既保证 $v_0\ll c$，也保证位移幅值 $v_0/\omega$ 远小于 $c/\omega$，可把电子所在区域的电场看成均匀。忽略磁力相对于电力的
$v/c$ 阶小量，电子加速度为
\begin{equation}
  a_z=-\frac{eE_0}{m}\cos\omega t.
\end{equation}
把它代入拉莫尔公式并对周期取平均，因 $\langle\cos^2\omega t\rangle=1/2$，有
\begin{equation}
  \langle P_e\rangle=\frac{e^4E_0^2}{12\pi\varepsilon_0m^2c^3}.
\end{equation}
入射波的平均能流密度是 $\langle S_{\rm in}\rangle=\varepsilon_0cE_0^2/2$。散射截面定义为
“总散射功率除以入射能流密度”，于是
\begin{align}
  \sigma_T
  &=\frac{\langle P_e\rangle}{\langle S_{\rm in}\rangle}
  =\frac{e^4}{6\pi\varepsilon_0^2m^2c^4}
  \\
  &=\boxed{\frac{8\pi}{3}\left(\frac{1}{4\pi\varepsilon_0}
  \frac{e^2}{mc^2}\right)^2}
  =\frac{8\pi}{3}r_e^2
  \approx\SI{6.65e-29}{m^2},
\end{align}
其中 $r_e=e^2/(4\pi\varepsilon_0mc^2)$ 称为经典电子半径。截面具有面积量纲，但它不意味着电子是
半径 $\sqrt{\sigma_T/\pi}$ 的实心小球；它衡量的是电子从入射波中截获并再辐射能量的能力。

在这一低频、非相对论经典近似中，$E_0^2$ 在散射功率与入射能流之比中消去，结果也不显含频率，
所以汤姆孙截面与入射波强度、频率均无关。这个结论有明确的适用边界：频率高到必须考虑电子
反冲和量子效应时（光子能量不再满足 $\hbar\omega\ll mc^2$），需要使用更一般的散射理论，不能仍把自由电子当成被均匀电场驱动的经典点电荷。这里还假定电子不受原子束缚；束缚电子在低频、近共振等区域一般不满足这一频率无关结论。
图中的电子电荷为 $-e$，所以瞬时加速度与电场反向；散射强弱取决于 $a^2$，不因这种反向而改变。

散射截面的定义可以想成：若入射能流密度为 $S_{\rm in}$，把总散射功率写成
$P_{\rm sc}=S_{\rm in}\sigma$，那么 $\sigma$ 就是一个等效受光面积。它允许我们不追踪每一束散射波，
而用一个具有面积量纲的量比较相互作用强弱。对一个自由电子，偶极辐射在电子加速度方向上最弱、在
垂直方向上最强；把各方向功率全部积分才得到上面的 $\sigma_T$。经典电子半径 $r_e$ 也不是测得的
电子几何边界，而是由 $e,m,c,\varepsilon_0$ 组合出来的电磁作用长度尺度。

图~\ref{ch26:ref:review-thomson} 中向下的加速度与向上的瞬时电场反向；两者沿同一振动轴，散射功率取决于其横向投影。

\thomsonScatteringDiagram


% EM teaching enrichment 20260927
\paragraph{截面的面积单位与物理含义}
散射截面定义为总散射功率除以入射光强，因此单位是面积，但它不是把电子当作半径固定的小硬球得到的几何面积。
经典弱场模型中，电子加速度振幅 $a_0=eE_0/m_e$，平均辐射功率正比于 $a_0^2$，入射光强正比于 $E_0^2$，相除后振幅抵消，得到与入射强度无关的截面。
若频率接近束缚电子的共振，或光子反冲与电子静能不可忽略，自由、非相对论电子的这套推导便不足以使用；它不能无条件推广为所有物质、所有频率的散射规律。


\subsection{介质中的传播与色散}


均匀、线性、各向同性且无耗的介质中 $\vect D=\varepsilon\vect E$、$\vect B=\mu\vect H$，相速度
$v=1/\sqrt{\varepsilon\mu}$。实际介质的 $\varepsilon$ 随频率变化，因此折射率
$n(\omega)=c/v_p$ 具有色散。波包包络以群速度
\begin{equation}
 v_g=\frac{\dd\omega}{\dd k}
\end{equation}
传播；在强色散和吸收区，群速度不能不加分析地解释为信息或能量速度。

法向入射到两个无耗、非磁介质的平面波，在界面处切向 $E,H$ 连续。定义波阻抗
\begin{equation}
 Z=\sqrt{\frac{\mu}{\varepsilon}},
\end{equation}
电场振幅反射系数和透射系数为
\begin{equation}
 r_E=\frac{Z_2-Z_1}{Z_2+Z_1},
 \qquad
 t_E=\frac{2Z_2}{Z_1+Z_2}.
\end{equation}
这些振幅系数来自两个边界方程。统一以切向电场的同一正方向记有符号振幅，
则 $E_i+E_r=E_t$；由于反射波传播反向，其磁场对该正方向的贡献反号，故
\[
 \frac{E_i-E_r}{Z_1}=\frac{E_t}{Z_2}.
\]
代入 $E_t=E_i+E_r$，整理为 $(Z_2-Z_1)E_i=(Z_1+Z_2)E_r$，
得到 $r_E$，再用 $t_E=1+r_E$ 得到透射系数。
平均能流为 $E_0^2/(2Z)$，所以
\[
 R=|r_E|^2,\qquad T=\frac{Z_1}{Z_2}|t_E|^2
   =\frac{4Z_1Z_2}{(Z_1+Z_2)^2},\qquad R+T=1.
\]
最后一步是将分子 $(Z_2-Z_1)^2+4Z_1Z_2$ 合并为 $(Z_1+Z_2)^2$。

\begin{center}
\begin{tikzpicture}[scale=.9,>=Latex]
  \fill[DeepBlue!5] (4,-2)--(8,-2)--(8,2)--(4,2)--cycle;
  \draw[very thick] (4,-2)--(4,2);
  \draw[vector] (.5,.6)--(3.65,.6) node[midway,above]{$\vect k_i$};
  \draw[vector] (3.55,-.5)--(.7,-.5) node[midway,below]{$\vect k_r$};
  \draw[vector] (4.35,.6)--(7.5,.6) node[midway,above]{$\vect k_t$};
  \node[clearlabel] at (2,1.55) {$Z_1$};
  \node[clearlabel] at (6,1.55) {$Z_2$};
  \node[clearlabel,below] at (4,-2) {界面};
\end{tikzpicture}
\figcap{法向入射的反射与透射。箭头表示传播方向；反射电场的正负由振幅系数 $r_E$ 决定。}\label{ch26:fig:interface}
\end{center}
图~\ref{ch26:fig:interface} 把三个传播方向分开；$r_E<0$ 表示电场相位翻转，并不把反射箭头改画成向右。


% EM teaching enrichment 20260927
\paragraph{相速度与群速度怎样从同一色散关系得到}
在窄带、弱吸收条件下写 $k=n(\omega)\omega/c$，则
\[
 v_p=\frac c{n(\omega)},\qquad
 v_g=\left(\frac{\dd k}{\dd\omega}\right)^{-1}
 =\frac c{n+\omega\,\dd n/\dd\omega}.
\]
无色散时两者相同；有色散时，单色波峰的速度与包络速度不同。若在某一频段近似 $n=n_0+a\omega$，则 $v_g=c/(n_0+2a\omega)$，并非简单的 $c/n$。
这个线性近似只在所给窄频段使用，不能外推到任意频率。强吸收引起波包重塑时，包络峰值运动也不能直接当作新信息的传播速度。

\begin{origexamplebox}{介质界面的能流守恒}
真空中的平面波法向入射到折射率为 $n$、磁导率近似为 $\mu_0$ 的无耗介质。设 $n>0$。求电场振幅反射系数、功率反射率与透射率，验证能流守恒，并检验 $n=1$ 与 $n\to\infty$ 的极限。
\end{origexamplebox}
\begin{xesolution}
界面没有自由面电流时，切向电场和磁场连续。把入射电场振幅归一化为1，
反射、透射振幅分别为 $r_E,t_E$，则
\[
 1+r_E=t_E,\qquad \frac{1-r_E}{Z_0}=\frac{t_E}{Z_0/n}.
\]
第二式中的负号来自反射波传播方向反转，不能把反射磁场仍按入射方向相加。
消去 $t_E$，即 $1-r_E=n(1+r_E)$。非磁介质波阻抗 $Z=Z_0/n$，于是
\[
 r_E=\frac{Z_2-Z_1}{Z_2+Z_1}
 =\frac{1-n}{1+n}.
\]
因此
\[
 \boxed{R=|r_E|^2=\left(\frac{n-1}{n+1}\right)^2}.
\]
透射振幅系数 $t_E=1+r_E=2/(1+n)$，所以
\[
 T=n|t_E|^2=\frac{4n}{(1+n)^2},\qquad
 R+T=\frac{(n-1)^2+4n}{(n+1)^2}=1.
\]
只有 $n>1$ 时 $r_E<0$，反射电场相位翻转；$0<n<1$ 时不翻转。
$n=1$ 时 $R=0,T=1$；在本无耗实折射率模型的 $n\to\infty$ 极限中 $R\to1,T\to0$。
后一极限不等于对真实吸收金属的描述。$R,T$ 均无量纲且位于 $[0,1]$。
\end{xesolution}


% EM structure and teaching revision 20260928
\subsection{平面波界面的振幅与能流核对}
两种无耗散、非磁性介质折射率为 $n_1,n_2$，平面波正入射，取同一切向电场正向。电、磁场的切向边界条件给出
\[
 E_i+E_r=E_t,\qquad n_1(E_i-E_r)=n_2E_t.
\]
反射波传播方向反转，故磁场方程中的反射项带负号。解得电场振幅系数
$r=(n_1-n_2)/(n_1+n_2)$、$t=2n_1/(n_1+n_2)$。负的 $r$ 表示电场相位反转，不表示能量反向增加。
能流与 $nE^2$ 成正比，所以反射率与透射率为
\[
 \mathcal R=|r|^2,\qquad \mathcal T=\frac{n_2}{n_1}|t|^2,
 \qquad \mathcal R+\mathcal T=1.
\]
例如空气近似 $n_1=1$、另一介质 $n_2=1.5$，有 $r=-0.2$、$t=0.8$，反射能量占 $4\%$，透射占 $96\%$。把 $t^2=0.64$ 当成透射率会误判能量不守恒。
这些公式使用正入射、线性、无耗散且磁导率相同的前提。斜入射须区分偏振；存在吸收时还要计入介质内的能量沉积。


\section{超导磁响应与磁通量子化}

\subsection{超导：零电阻还不够}


金属电阻率一般随温度降低而下降，因为晶格热振动减弱，传导电子受到的散射减少。随着低温技术发展，
人们自然追问电阻会趋向一个非零常数，还是能够真正消失。1908 年，荷兰物理学家卡末林·昂内斯（图~\ref{ch26:ref:review-onnes-portrait}）
成功液化氦，获得约 $\SI{4.2}{K}$ 的低温条件；1911 年他研究汞的电阻时，发现温度降到这一范围后，
电阻突然跌至实验分辨率以下，超导态由此被发现。

\historyfigure[0.23\linewidth]{onnes-built-in-v1.png}%
{卡末林·昂内斯（Heike Kamerlingh Onnes），低温物理和超导现象的先驱；历史人物肖像。\label{ch26:ref:review-onnes-portrait}}\label{ch26:fig:s41-10-8}

后来人们在多种元素、合金和铜氧化物中发现超导转变。材料进入超导态的温度称为临界温度 $T_c$；
实际超导还受临界磁场和临界电流密度限制，不能只凭 $T_c$ 判断应用性能。从 1911 年的汞到 1973 年前后，
传统材料的 $T_c$ 提升十分缓慢；1986 年贝德诺尔茨（Johannes Georg Bednorz）与缪勒（Karl Alexander M\"uller）发现铜氧化物高温超导体，随后材料的
临界温度迅速越过液氮沸点 $\SI{77}{K}$。下图以分材料族的历史数据展示不同年代的临界温度，反映了
超导材料探索中长期缓慢推进与关键发现后快速跃升并存的过程。

图~\ref{ch26:ref:review-superconductivity-timeline} 串起零电阻的发现、迈斯纳效应和新材料温区，图中的年份用于定位各项发现。

\superconductivityTimeline

\superconductivityChartFaithful

1987 年，吴茂昆、朱经武等人的钇钡铜氧材料使转变温度达到约 $90\,$K，超过液氮沸点。随后铋系、铊系和汞系铜氧化物又推动了高温区的探索。早期有过常规超导临界温度约 $30\,$K 的经验估计，但它不是 BCS 理论给出的普适严格上限，不能把后来的高温发现理解为违反一条已证明的上限定理。

零电阻最直接的后果，是闭合超导环中的电流可以长期存在而不发生焦耳耗散。但若只把超导体看作
电阻率严格为零的普通导体，仍解释不了它最重要的磁学性质。

“电阻测不到”与“电阻严格为零”在实验上并不完全等价，因为任何电压表都有分辨率。超导持续电流
实验提供了另一种证据：在闭合超导环中建立电流后，若电阻哪怕很小，磁能也应按一定时间尺度衰减；
实际观测到电流可以在极长时间内保持，由衰减上限反推的电阻率远低于普通良导体。零电阻意味着
$P_J=I^2R$ 不再产生焦耳热，却不意味着超导装置可以承受任意电流和磁场。温度、外磁场或电流密度
超过相应临界条件时，材料会回到有耗散的正常态。

超导临界温度的发展也说明实验条件与材料体系同样重要。早期研究依赖液氦低温，能探索的材料范围和
温区都很有限；传统金属与合金的转变温度长期只缓慢提高。铜氧化物高温超导出现后，液氮温区的实验
和应用成为可能，材料的晶体结构、载流子浓度及制备工艺也成为研究重点。图~\ref{ch26:fig:s41-10-9} 将材料名称、年代和实验条件列在各行，用横向长度表示临界温度，避免低温点拥挤。读图时应比较同一时代的数量级与跃迁趋势，
而不宜只盯住某一个最高点；图中高压条件下的数据也不能与常压工程材料直接等同。


% EM teaching enrichment 20260927
\paragraph{持续电流能给出多小的电阻上限}
把环近似为自感 $L$ 与极小电阻 $R$ 的闭合回路，则 $I(t)/I_0=e^{-Rt/L}$。若观察时间 $t$ 内电流相对下降不超过 $\delta$，便有
\[
 R\le-\frac Lt\ln(1-\delta)\simeq\frac{L\delta}{t}\quad(\delta\ll1).
\]
这是一条由测量精度得到的上限，不是由有限精度直接证明一个数学量严格为零。它说明长时间稳定的持续电流为何是灵敏的实验检验，同时也提示应把“测不到变化”与“没有任何变化”在表述上分开。


\subsection{理想导体与迈斯纳效应}


在理想导体内部，$\vect E=0$。法拉第定律给出
\begin{equation}
  \frac{\dd\Phi_B}{\dd t}=-\oint_C\vect E\cdot\dd\vect l=0.
\end{equation}
因此理想导体只保证磁通不再改变：降温前内部若已有磁场，进入零电阻状态后可以把这部分磁通“冻结”
下来。超导体却表现出更强的性质——在给定温度、外场条件下进入平衡迈斯纳态时，体内磁场被排斥到表面附近，这就是迈斯纳效应。对第二类超导体，需要外场低于相应下临界场；处于混合态时有量子化涡旋进入，不能笼统断言所有超导状态的体内场都为零。

两种操作顺序最能看出差别。若先把无磁场中的样品降至零电阻态、再施加外场，理想导体中的感应电流
会抵消磁通变化，体内可以保持 $B=0$；这个结果尚不能区分理想导体和超导体。若先在正常态样品中
施加磁场，再降温穿过 $T_c$，理想导体模型预言原有磁通应被冻结在体内。1933 年迈斯纳与奥森菲尔德
测量锡等样品的磁场分布，却发现无论先加磁场还是先降温，只要样品进入超导态且外场低于临界条件，
在平衡迈斯纳态下体内磁场都会被排斥。因而迈斯纳效应是一种平衡态磁学性质，不只是“电阻为零导致磁通不变”。

从磁场边界条件看，体内近似 $B=0$ 必须由表面附近的屏蔽电流来维持。这个电流不是普通电阻回路中
由持续电场推动的电流，而是超导凝聚态对外磁场的无耗散响应。宏观图常把磁感线画成在几何表面突然
转弯，真实情况则是在伦敦穿透深度的尺度上连续衰减。样品尺寸远大于穿透深度时，把体内看成
$B=0$ 是很好的竞赛近似；研究薄膜或表面效应时，就必须保留这段有限厚度。

图~\ref{ch26:ref:review-meissner} 中左图保留先前的内部磁场，右图外部场线绕开样品；蓝色场线彼此不相交，并满足样品表面无磁通穿入的条件。

\meissnerComparisonDiagram

图中左侧的磁力线穿过材料，表示理想导体保留初始磁通；右侧磁力线绕开超导体，表示迈斯纳排磁。
真实超导体的磁场并非在数学意义上于表面瞬间降为零，而是在称为伦敦穿透深度的薄层内衰减；
在本章的宏观近似中，可写成超导体体内 $B\approx0$。

因此，“完美导电性”和“完全抗磁性”回答的是两个不同问题。前者回答电场能否在体内维持以及电流是否
耗散；后者回答给定温度和外磁场后，平衡态磁通怎样分布。理想导体的最终磁场依赖进入零电阻状态前的
历史，超导体的迈斯纳态则在相同外部条件下趋向同一排磁结果。正是这种与路径无关的排磁性，使超导体
不能被简单等同为电导率无限大的普通金属。

把磁体放在超导体上方时，表面屏蔽电流改变空间磁场分布，可产生向上的磁力。若磁力与重力平衡，
磁体能够悬浮；图~\ref{ch26:ref:review-levitation-scene} 中磁体与低温样品间的间隙体现这一状态。实验中常见的稳定悬浮还与第二类超导体的磁通钉扎有关。

图~\ref{ch26:ref:review-levitation} 的右侧是磁体的受力图：向上的磁力与向下的重力都作用于该磁体，
平衡时二者大小相等。下方超导圆盘受到的力不属于这个受力分析对象。

\levitationDiagram

\historyfigure[0.32\linewidth]{levitation-built-in-v1.png}%
{超导磁悬浮实验：磁场与屏蔽电流共同产生向上的电磁力。\label{ch26:ref:review-levitation-scene}}


\subsection{超导环中的磁通保持}


以下采用宏观厚环、忽略动能电感的近似，并要求全过程不失超、无相位滑移且量子数保持不变。设超导环的几何自感为 $L$，外界穿过环的磁通为 $\Phi_{\rm ext}$，环中持续电流 $I$ 自身贡献磁链 $LI$。
环路电阻为零时，感应电动势满足
\begin{equation}
  L\frac{\dd I}{\dd t}=-\frac{\dd\Phi_{\rm ext}}{\dd t},
\end{equation}
所以
\begin{equation}
  \boxed{\frac{\dd}{\dd t}\bigl(\Phi_{\rm ext}+LI\bigr)=0}.
\end{equation}
先在外磁场中把圆环冷却到超导态，再撤去外场，$\Phi_{\rm ext}$ 减小；环中便出现持续电流，使 $LI$
增加，从而补偿总磁通的变化。

若圆环建立该超导状态时的总磁通为 $\Phi_{\rm ini}$，以后在上述适用条件内缓慢改变外磁场，都有
\begin{equation}
  I=\frac{\Phi_{\rm ini}-\Phi_{\rm ext}}{L}.
\end{equation}
外磁通减小时电流按楞次定律方向增大，外磁通增大时电流反向调整。这里保持的是“外磁通加自感磁通”
的总和，不是说外磁场本身被冻结不变。对于尺度足够大的经典问题，这一磁通守恒关系直接来自零电阻
环路的法拉第定律；在微观超导理论中还会出现磁通量子化，那是比本章宏观模型更进一步的内容。

图~\ref{ch26:ref:review-ring-hold} 中外场与自感场均取出纸面为正；逆时针持续电流按右手定则补偿逐渐撤去的外磁通。

\superconductingRingDiagram

这个宏观推导抓住了固定状态分支中总磁链保持的物理图像。$\Phi_{\rm ini}$ 是初始总磁通，不是后文定义的磁通量子 $\Phi_0=h/(2e)$。后面的磁通量子化一节进一步区分总磁通和磁通子，并讨论固定状态与最低能量状态的区别。

% EM teaching enrichment 20260927
\paragraph{撤场以后能量从何而来}
在固定状态分支中，若起初环流为零、外磁通为 $\Phi_i$，随后外磁通降到零，则
$I_f=\Phi_i/L$，最终环内能量为 $\Phi_i^2/(2L)$。
这不是零电阻环无中生有地产生能量；撤场装置与环之间通过互感交换功，必须把外磁场源及其驱动过程一并计入能量账。
零电阻保证没有 $I^2R$ 耗散，却不保证环的磁能不变。若改变几何结构而使 $L$ 改变，也应对 $\Phi_{\rm ext}+LI$ 整体求导，不能继续把 $L$ 提到导数外。



\subsection{超导环的量子态选择}


超导体进入迈斯纳态时，体内磁感应强度被排斥；含孔的超导环还允许磁通以离散方式被俘获。对电荷为 $2e$ 的库珀对，磁通量子为
\begin{equation}
 \Phi_0=\frac{h}{2e}.
\end{equation}
在壁厚远大于穿透深度、环路位于超导体内部的理想情形，磁通子（fluxoid）量子化可近似写成
$\Phi=n\Phi_0$。薄环或沿积分路径的超电流密度不可忽略时，量子化的是
\begin{equation}
 \Phi+\mu_0\lambda_L^2\oint_C\vect j_s\cdot\dd\vect l=n\Phi_0.
\end{equation}
其中 $\lambda_L$ 为伦敦穿透深度，$\vect j_s$ 为超电流密度，$C$ 位于超导材料内，磁通 $\Phi$ 的正法向与 $C$ 按右手规则配对。
第二项反映载流子运动；只有选取深入厚壁、$\vect j_s\approx0$ 的路径时才能略去。
不能把任意几何区域的磁通都机械地取整。用集中参数描述薄环时，可将这一项写成 $L_kI$，
并用总电感 $L=L_g+L_k$ 代替单纯的几何自感。

\begin{center}
\begin{tikzpicture}[scale=.9,>=Latex]
 \fill[DeepBlue!12,even odd rule] (0,0) circle (1.7) (0,0) circle (.9);
 \draw[thick] (0,0) circle (1.7) (0,0) circle (.9);
 \foreach \x/\y in {0/0,.45/0,-.45/0,0/.45,0/-.45}{
   \draw[DeepBlue] (\x,\y) circle (.09);\fill[DeepBlue] (\x,\y) circle (.025);
 }
 \draw[->,DeepBlue,very thick] (25:1.3) arc[start angle=25,end angle=150,radius=1.3];
 \node[above] at (0,1.9) {$I>0$：逆时针};
 \node[align=left,anchor=west] at (2.4,.3)
  {俯视图：$\odot$ 表示出纸面磁通\\[5pt]
   取出纸面为正，$I>0$ 产生正自感磁通\\[5pt]
   $I<0$ 时环流和自感场同时反向};
\end{tikzpicture}
\figcap{含孔超导环中的俘获磁通。环流方向调节总磁通子满足量子化条件。}\label{ch26:fig:flux-quantization}
\end{center}
图~\ref{ch26:fig:flux-quantization} 中持续电流产生的自感磁通与外磁通相加，正是下题模型 $\Phi_{\rm ext}+LI=n\Phi_0$ 的几何含义。

\begin{origexamplebox}{超导环的最邻近磁通态}
超导环外加磁通为 $\Phi_{\rm ext}$，总电感为 $L$（厚环可忽略动能电感）。假设不同量子态之间能够充分弛豫，并忽略其他态间能量差。用简化模型
$\Phi_{\rm ext}+LI=n\Phi_0$，相应电感能量为 $U_n=LI^2/2$。求给定 $\Phi_{\rm ext}$ 时的平衡量子数选择规则，并求相邻两态简并的位置。
\end{origexamplebox}
\begin{xesolution}
第 $n$ 态的持续电流和能量为
\[
 I_n=\frac{n\Phi_0-\Phi_{\rm ext}}{L},
 \qquad
 U_n=\frac{(n\Phi_0-\Phi_{\rm ext})^2}{2L}.
\]
平衡态选择使 $|n\Phi_0-\Phi_{\rm ext}|$ 最小的整数 $n$，即
\[
 \boxed{n=\operatorname*{argmin}_{m\in\mathbb Z}|m-\Phi_{\rm ext}/\Phi_0|}.
\]
$n$ 与 $n+1$ 两态简并时
\[
 (n\Phi_0-\Phi_{\rm ext})^2=[(n+1)\Phi_0-\Phi_{\rm ext}]^2,
\]
两边相减，得
\[
 [(n+1)-n]\Phi_0[(2n+1)\Phi_0-2\Phi_{\rm ext}]=0.
\]
因 $\Phi_0\ne0$，故
\[
 \boxed{\Phi_{\rm ext}=\left(n+\frac12\right)\Phi_0}.
\]
$\Phi/L$ 的单位为安培，能量式单位为焦耳。远离半整数点时最低能态唯一；在简并点两种环流分别为 $-\Phi_0/(2L)$ 和 $+\Phi_0/(2L)$，能量均为 $\Phi_0^2/(8L)$。
上述规则选择全局最低能量态，并不保证实际圆环随外场改变就立即跳到该态。若量子态转换需要跨越相位滑移势垒，圆环可滞留在固定 $n$ 的亚稳态；此时应沿同一个 $I_n(\Phi_{\rm ext})$ 分支计算，直到发生允许的状态转换。
\end{xesolution}

\begin{starread}{波导、驻波与边界条件}
理想导体表面切向电场为零。空心、单连通截面的金属波导只允许满足边界条件的离散横向模式，并存在截止频率。
这里不包括同轴线等双导体结构；后者可以支持无截止频率的横电磁（TEM）模式。
对真空填充、理想无耗的矩形波导，横电（TE）与横磁（TM）模式具有不同场结构；传播模式满足 $\omega^2=c^2k_z^2+\omega_c^2$，故 $v_p=\omega/k_z>c$、$v_g=c^2k_z/\omega<c$，并不允许超光速传信。
谐振腔把传播方向也离散化，品质因数由单位周期储能与损耗之比决定。电磁模式量子化后，每个模式成为频率为
$\omega$ 的量子谐振子，这是黑体辐射与光子概念的基础之一。
\end{starread}

% EM structure and teaching revision 20260928
\subsection{磁通保持、磁通子与材料临界条件}
理想导体由零电阻导出的结论是既有磁通不随时间改变；超导体的平衡磁响应还包括迈斯纳效应。二者对“先加磁场后降温”和“先降温后加磁场”的响应不能简单等同。
闭合超导环的量子条件约束的是磁通子。若把几何电感与动能电感之和记为 $L_{\mathrm{tot}}$，外磁通为 $\Phi_{\mathrm{ext}}$，简化模型给出
\[
 \Phi_{\mathrm{ext}}+L_{\mathrm{tot}}I=n\Phi_0,\qquad
 W_n=\frac{(n\Phi_0-\Phi_{\mathrm{ext}})^2}{2L_{\mathrm{tot}}},\quad
 \Phi_0=\frac{h}{2e}.
\]
相邻两条能量抛物线在半整数外磁通处相交，比较 $W_n$ 可确定热力学最低能态。但实际扫场时，能垒会使环停留在亚稳态，未必立即跳到最低的一支。
由某个整数 $n$ 算得电流之后，还须核对临界电流、临界磁场和温度条件；若超出材料允许范围，零电阻与给定电感的模型就不能继续使用。把“允许的量子态”“最低能态”与“实验能到达的状态”分开，才不会把整数取整规则误作所有情形下的实际演化规律。
}

\endgroup

\begingroup\tcbset{breakable/.default=false}
\renewcommand{\Needspace}[1]{\par\begingroup
\dimen0=#1\relax\vskip0pt plus\dimen0\penalty-100
\vskip0pt plus-\dimen0\vskip\dimen0\penalty9999\vskip-\dimen0
\vskip0pt\endgroup}
\chapter{几何光学：光线、反射和折射}


\graphicspath{{/Users/dongxingxi/Documents/物理竞赛一轮资料/sample/xes-original/forty-second-section/}{/Users/dongxingxi/Documents/物理竞赛一轮资料/sample/xes-original/forty-second-section/assets/}}


\section{光线模型、波面与光学发展}


\newcommand{\AtmosphereWavefrontPanel}{%
\begin{center}\begin{minipage}{\linewidth}\centering
\begin{tikzpicture}[>=Latex,font=\small,every node/.style={fill=none}]
\draw[thick] plot coordinates {(2.8670322,-1.4924825) (2.8279646,-1.4378127) (2.7878547,-1.3839031) (2.7467173,-1.3307734) (2.7045675,-1.2784432) (2.6614209,-1.2269318) (2.6172933,-1.1762583) (2.5722011,-1.1264412) (2.5261608,-1.0774989) (2.4791894,-1.0294495) (2.4313043,-0.9823107) (2.3825230,-0.9360999) (2.3328636,-0.8908340) (2.2823444,-0.8465298) (2.2309840,-0.8032037) (2.1788012,-0.7608715) (2.1258154,-0.7195488) (2.0720461,-0.6792510) (2.0175131,-0.6399929) (1.9622365,-0.6017888) (1.9062366,-0.5646530) (1.8495342,-0.5285991) (1.7921500,-0.4936404) (1.7341053,-0.4597897) (1.6754215,-0.4270596) (1.6161201,-0.3954621) (1.5562231,-0.3650088) (1.4957525,-0.3357111) (1.4347306,-0.3075796) (1.3731799,-0.2806248) (1.3111231,-0.2548565) (1.2485830,-0.2302843) (1.1855827,-0.2069173) (1.1221455,-0.1847640) (1.0582946,-0.1638327) (0.9940537,-0.1441309) (0.9294464,-0.1256661) (0.8644965,-0.1084449) (0.7992280,-0.0924738) (0.7336650,-0.0777587) (0.6678315,-0.0643049) (0.6017519,-0.0521174) (0.5354504,-0.0412007) (0.4689516,-0.0315588) (0.4022800,-0.0231953) (0.3354601,-0.0161133) (0.2685166,-0.0103154) (0.2014741,-0.0058036) (0.1343573,-0.0025798) (0.0671910,-0.0006450) (0.0000000,0.0000000) (-0.0671910,-0.0006450) (-0.1343573,-0.0025798) (-0.2014741,-0.0058036) (-0.2685166,-0.0103154) (-0.3354601,-0.0161133) (-0.4022800,-0.0231953) (-0.4689516,-0.0315588) (-0.5354504,-0.0412007) (-0.6017519,-0.0521174) (-0.6678315,-0.0643049) (-0.7336650,-0.0777587) (-0.7992280,-0.0924738) (-0.8644965,-0.1084449) (-0.9294464,-0.1256661) (-0.9940537,-0.1441309) (-1.0582946,-0.1638327) (-1.1221455,-0.1847640) (-1.1855827,-0.2069173) (-1.2485830,-0.2302843) (-1.3111231,-0.2548565) (-1.3731799,-0.2806248) (-1.4347306,-0.3075796) (-1.4957525,-0.3357111) (-1.5562231,-0.3650088) (-1.6161201,-0.3954621) (-1.6754215,-0.4270596) (-1.7341053,-0.4597897) (-1.7921500,-0.4936404) (-1.8495342,-0.5285991) (-1.9062366,-0.5646530) (-1.9622365,-0.6017888) (-2.0175131,-0.6399929) (-2.0720461,-0.6792510) (-2.1258154,-0.7195488) (-2.1788012,-0.7608715) (-2.2309840,-0.8032037) (-2.2823444,-0.8465298) (-2.3328636,-0.8908340) (-2.3825230,-0.9360999) (-2.4313043,-0.9823107) (-2.4791894,-1.0294495) (-2.5261608,-1.0774989) (-2.5722011,-1.1264412) (-2.6172933,-1.1762583) (-2.6614209,-1.2269318) (-2.7045675,-1.2784432) (-2.7467173,-1.3307734) (-2.7878547,-1.3839031) (-2.8279646,-1.4378127) (-2.8670322,-1.4924825)};
\draw[thick] plot coordinates {(3.1127778,-1.3204095) (3.0703616,-1.2610538) (3.0268137,-1.2025234) (2.9821502,-1.1448397) (2.9363876,-1.0880240) (2.8895427,-1.0320974) (2.8416327,-0.9770804) (2.7926755,-0.9229933) (2.7426889,-0.8698559) (2.6916914,-0.8176880) (2.6397018,-0.7665088) (2.5867393,-0.7163370) (2.5328234,-0.6671912) (2.4779739,-0.6190895) (2.4222112,-0.5720497) (2.3655556,-0.5260890) (2.3080282,-0.4812244) (2.2496501,-0.4374725) (2.1904428,-0.3948494) (2.1304282,-0.3533707) (2.0696283,-0.3130518) (2.0080657,-0.2739076) (1.9457629,-0.2359524) (1.8827429,-0.1992002) (1.8190291,-0.1636647) (1.7546447,-0.1293588) (1.6896137,-0.0962953) (1.6239599,-0.0644863) (1.5577075,-0.0339436) (1.4908810,-0.0046783) (1.4235051,0.0232986) (1.3556044,0.0499770) (1.2872041,0.0753469) (1.2183294,0.0993991) (1.1490056,0.1221245) (1.0792583,0.1435150) (1.0091132,0.1635625) (0.9385962,0.1822598) (0.8677333,0.1995998) (0.7965505,0.2155763) (0.7250742,0.2301833) (0.6533306,0.2434154) (0.5813462,0.2552678) (0.5091475,0.2657362) (0.4367612,0.2748165) (0.3642139,0.2825056) (0.2915323,0.2888005) (0.2187433,0.2936989) (0.1458737,0.2971991) (0.0729503,0.2992997) (0.0000000,0.3000000) (-0.0729503,0.2992997) (-0.1458737,0.2971991) (-0.2187433,0.2936989) (-0.2915323,0.2888005) (-0.3642139,0.2825056) (-0.4367612,0.2748165) (-0.5091475,0.2657362) (-0.5813462,0.2552678) (-0.6533306,0.2434154) (-0.7250742,0.2301833) (-0.7965505,0.2155763) (-0.8677333,0.1995998) (-0.9385962,0.1822598) (-1.0091132,0.1635625) (-1.0792583,0.1435150) (-1.1490056,0.1221245) (-1.2183294,0.0993991) (-1.2872041,0.0753469) (-1.3556044,0.0499770) (-1.4235051,0.0232986) (-1.4908810,-0.0046783) (-1.5577075,-0.0339436) (-1.6239599,-0.0644863) (-1.6896137,-0.0962953) (-1.7546447,-0.1293588) (-1.8190291,-0.1636647) (-1.8827429,-0.1992002) (-1.9457629,-0.2359524) (-2.0080657,-0.2739076) (-2.0696283,-0.3130518) (-2.1304282,-0.3533707) (-2.1904428,-0.3948494) (-2.2496501,-0.4374725) (-2.3080282,-0.4812244) (-2.3655556,-0.5260890) (-2.4222112,-0.5720497) (-2.4779739,-0.6190895) (-2.5328234,-0.6671912) (-2.5867393,-0.7163370) (-2.6397018,-0.7665088) (-2.6916914,-0.8176880) (-2.7426889,-0.8698559) (-2.7926755,-0.9229933) (-2.8416327,-0.9770804) (-2.8895427,-1.0320974) (-2.9363876,-1.0880240) (-2.9821502,-1.1448397) (-3.0268137,-1.2025234) (-3.0703616,-1.2610538) (-3.1127778,-1.3204095)};
\draw[thick] plot coordinates {(3.3585234,-1.1483366) (3.3127585,-1.0842949) (3.2657727,-1.0211436) (3.2175831,-0.9589060) (3.1682077,-0.8976049) (3.1176645,-0.8372630) (3.0659722,-0.7779025) (3.0131498,-0.7195454) (2.9592169,-0.6622130) (2.9041933,-0.6059266) (2.8480993,-0.5507068) (2.7909556,-0.4965741) (2.7327831,-0.4435484) (2.6736035,-0.3916492) (2.6134384,-0.3408957) (2.5523100,-0.2913066) (2.4902409,-0.2429001) (2.4272540,-0.1956940) (2.3633725,-0.1497059) (2.2986199,-0.1049526) (2.2330200,-0.0614507) (2.1665972,-0.0192161) (2.0993757,0.0217356) (2.0313805,0.0613892) (1.9626366,0.0997302) (1.8931693,0.1367444) (1.8230042,0.1724182) (1.7521672,0.2067385) (1.6806844,0.2396925) (1.6085822,0.2712681) (1.5358870,0.3014538) (1.4626258,0.3302383) (1.3888255,0.3576112) (1.3145133,0.3835621) (1.2397166,0.4080817) (1.1644629,0.4311609) (1.0887801,0.4527912) (1.0126960,0.4729645) (0.9362386,0.4916735) (0.8594361,0.5089113) (0.7823169,0.5246715) (0.7049093,0.5389482) (0.6272419,0.5517364) (0.5493434,0.5630311) (0.4712423,0.5728283) (0.3929676,0.5811244) (0.3145480,0.5879163) (0.2360125,0.5932014) (0.1573900,0.5969780) (0.0787095,0.5992444) (0.0000000,0.6000000) (-0.0787095,0.5992444) (-0.1573900,0.5969780) (-0.2360125,0.5932014) (-0.3145480,0.5879163) (-0.3929676,0.5811244) (-0.4712423,0.5728283) (-0.5493434,0.5630311) (-0.6272419,0.5517364) (-0.7049093,0.5389482) (-0.7823169,0.5246715) (-0.8594361,0.5089113) (-0.9362386,0.4916735) (-1.0126960,0.4729645) (-1.0887801,0.4527912) (-1.1644629,0.4311609) (-1.2397166,0.4080817) (-1.3145133,0.3835621) (-1.3888255,0.3576112) (-1.4626258,0.3302383) (-1.5358870,0.3014538) (-1.6085822,0.2712681) (-1.6806844,0.2396925) (-1.7521672,0.2067385) (-1.8230042,0.1724182) (-1.8931693,0.1367444) (-1.9626366,0.0997302) (-2.0313805,0.0613892) (-2.0993757,0.0217356) (-2.1665972,-0.0192161) (-2.2330200,-0.0614507) (-2.2986199,-0.1049526) (-2.3633725,-0.1497059) (-2.4272540,-0.1956940) (-2.4902409,-0.2429001) (-2.5523100,-0.2913066) (-2.6134384,-0.3408957) (-2.6736035,-0.3916492) (-2.7327831,-0.4435484) (-2.7909556,-0.4965741) (-2.8480993,-0.5507068) (-2.9041933,-0.6059266) (-2.9592169,-0.6622130) (-3.0131498,-0.7195454) (-3.0659722,-0.7779025) (-3.1176645,-0.8372630) (-3.1682077,-0.8976049) (-3.2175831,-0.9589060) (-3.2657727,-1.0211436) (-3.3127585,-1.0842949) (-3.3585234,-1.1483366)};
\fill (0.0000000,0.0000000) circle (1.2pt);
\node[below left,fill=none,inner sep=2pt] at (0.0000000,0.0000000) {$O$};
\draw[gray,dashed,thin] (-3.0000000,0.0000000)--(3.0000000,0.0000000);
\node[right,fill=none,inner sep=2pt] at (3.2000000,-0.2000000) {局地水平};
\draw[DeepBlue,thick,<-] plot coordinates {(0.0000000,0.0000000) (0.0350000,0.0155563) (0.0700000,0.0307340) (0.1050000,0.0455462) (0.1400000,0.0600054) (0.1750000,0.0741238) (0.2100000,0.0879130) (0.2450000,0.1013845) (0.2800000,0.1145492) (0.3150000,0.1274176) (0.3500000,0.1399998) (0.3850000,0.1523058) (0.4200000,0.1643450) (0.4550000,0.1761267) (0.4900000,0.1876596) (0.5250000,0.1989523) (0.5600000,0.2100131) (0.5950000,0.2208500) (0.6300000,0.2314706) (0.6650000,0.2418825) (0.7000000,0.2520927) (0.7350000,0.2621083) (0.7700000,0.2719358) (0.8050000,0.2815819) (0.8400000,0.2910526) (0.8750000,0.3003542) (0.9100000,0.3094922) (0.9450000,0.3184725) (0.9800000,0.3273004) (1.0150000,0.3359812) (1.0500000,0.3445199) (1.0850000,0.3529214) (1.1200000,0.3611905) (1.1550000,0.3693316) (1.1900000,0.3773492) (1.2250000,0.3852475) (1.2600000,0.3930307) (1.2950000,0.4007027) (1.3300000,0.4082673) (1.3650000,0.4157282) (1.4000000,0.4230890) (1.4350000,0.4303531) (1.4700000,0.4375238) (1.5050000,0.4446045) (1.5400000,0.4515980) (1.5750000,0.4585076) (1.6100000,0.4653360) (1.6450000,0.4720860) (1.6800000,0.4787603) (1.7150000,0.4853616) (1.7500000,0.4918923) (1.7850000,0.4983549) (1.8200000,0.5047518) (1.8550000,0.5110851) (1.8900000,0.5173570) (1.9250000,0.5235698) (1.9600000,0.5297253) (1.9950000,0.5358256) (2.0300000,0.5418726) (2.0650000,0.5478681) (2.1000000,0.5538139) (2.1350000,0.5597117) (2.1700000,0.5655632) (2.2050000,0.5713698) (2.2400000,0.5771333) (2.2750000,0.5828550) (2.3100000,0.5885364) (2.3450000,0.5941789) (2.3800000,0.5997838) (2.4150000,0.6053525) (2.4500000,0.6108861) (2.4850000,0.6163858) (2.5200000,0.6218529) (2.5550000,0.6272885) (2.5900000,0.6326936) (2.6250000,0.6380693) (2.6600000,0.6434166) (2.6950000,0.6487364) (2.7300000,0.6540298) (2.7650000,0.6592976) (2.8000000,0.6645408) (2.8350000,0.6697601) (2.8700000,0.6749563) (2.9050000,0.6801304) (2.9400000,0.6852830) (2.9750000,0.6904148) (3.0100000,0.6955267) (3.0450000,0.7006192) (3.0800000,0.7056930) (3.1150000,0.7107489) (3.1500000,0.7157873) (3.1850000,0.7208090) (3.2200000,0.7258144) (3.2550000,0.7308041) (3.2900000,0.7357788) (3.3250000,0.7407388) (3.3600000,0.7456847) (3.3950000,0.7506170) (3.4300000,0.7555362) (3.4650000,0.7604426) (3.5000000,0.7653368)};
\draw[gray,dashed,thin] (0.0000000,0.0000000)--(3.5000000,1.5750000);
\draw[gray,dashed,thin] (0.0000000,0.0000000)--(3.5000000,0.4550000);
\node[right,fill=none,inner sep=2pt] at (3.5000000,1.5800000) {视方向};
\node[right,fill=none,inner sep=2pt] at (3.5000000,0.3800000) {真实方向};
\node[below,fill=none,inner sep=2pt] at (-1.5000000,-0.4000000) {地表};
\begin{scope}[shift={(-4.0000000,2.5000000)}]
\draw[thick] plot coordinates {(0.2892544,-0.3447200) (0.2952266,-0.3396193) (0.3011088,-0.3344152) (0.3068993,-0.3291092) (0.3125963,-0.3237029) (0.3181981,-0.3181981) (0.3237029,-0.3125963) (0.3291092,-0.3068993) (0.3344152,-0.3011088) (0.3396193,-0.2952266) (0.3447200,-0.2892544) (0.3497157,-0.2831942) (0.3546048,-0.2770477) (0.3593860,-0.2708168) (0.3640576,-0.2645034) (0.3686184,-0.2581094) (0.3730669,-0.2516368) (0.3774018,-0.2450876) (0.3816216,-0.2384637) (0.3857253,-0.2317671) (0.3897114,-0.2250000) (0.3935789,-0.2181643) (0.3973264,-0.2112622) (0.4009529,-0.2042957) (0.4044573,-0.1972670) (0.4078385,-0.1901782) (0.4110955,-0.1830315) (0.4142272,-0.1758290) (0.4172327,-0.1685730) (0.4201112,-0.1612656) (0.4228617,-0.1539091) (0.4254834,-0.1465057) (0.4279754,-0.1390576) (0.4303371,-0.1315673) (0.4325678,-0.1240368) (0.4346666,-0.1164686) (0.4366331,-0.1088649) (0.4384665,-0.1012280) (0.4401664,-0.0935603) (0.4417322,-0.0858640) (0.4431635,-0.0781417) (0.4444598,-0.0703955) (0.4456206,-0.0626279) (0.4466458,-0.0548412) (0.4475349,-0.0470378) (0.4482876,-0.0392201) (0.4489038,-0.0313904) (0.4493833,-0.0235512) (0.4497259,-0.0157048) (0.4499315,-0.0078536) (0.4500000,0.0000000) (0.4499315,0.0078536) (0.4497259,0.0157048) (0.4493833,0.0235512) (0.4489038,0.0313904) (0.4482876,0.0392201) (0.4475349,0.0470378) (0.4466458,0.0548412) (0.4456206,0.0626279) (0.4444598,0.0703955) (0.4431635,0.0781417) (0.4417322,0.0858640) (0.4401664,0.0935603) (0.4384665,0.1012280) (0.4366331,0.1088649) (0.4346666,0.1164686) (0.4325678,0.1240368) (0.4303371,0.1315673) (0.4279754,0.1390576) (0.4254834,0.1465057) (0.4228617,0.1539091) (0.4201112,0.1612656) (0.4172327,0.1685730) (0.4142272,0.1758290) (0.4110955,0.1830315) (0.4078385,0.1901782) (0.4044573,0.1972670) (0.4009529,0.2042957) (0.3973264,0.2112622) (0.3935789,0.2181643) (0.3897114,0.2250000) (0.3857253,0.2317671) (0.3816216,0.2384637) (0.3774018,0.2450876) (0.3730669,0.2516368) (0.3686184,0.2581094) (0.3640576,0.2645034) (0.3593860,0.2708168) (0.3546048,0.2770477) (0.3497157,0.2831942) (0.3447200,0.2892544) (0.3396193,0.2952266) (0.3344152,0.3011088) (0.3291092,0.3068993) (0.3237029,0.3125963) (0.3181981,0.3181981) (0.3125963,0.3237029) (0.3068993,0.3291092) (0.3011088,0.3344152) (0.2952266,0.3396193) (0.2892544,0.3447200)};
\draw[thick] plot coordinates {(0.5463695,-0.6511378) (0.5576502,-0.6415031) (0.5687610,-0.6316731) (0.5796986,-0.6216506) (0.5904596,-0.6114388) (0.6010408,-0.6010408) (0.6114388,-0.5904596) (0.6216506,-0.5796986) (0.6316731,-0.5687610) (0.6415031,-0.5576502) (0.6511378,-0.5463695) (0.6605741,-0.5349223) (0.6698091,-0.5233123) (0.6788402,-0.5115428) (0.6876644,-0.4996175) (0.6962792,-0.4875400) (0.7046819,-0.4753140) (0.7128700,-0.4629432) (0.7208409,-0.4504314) (0.7285922,-0.4377824) (0.7361216,-0.4250000) (0.7434268,-0.4120882) (0.7505055,-0.3990508) (0.7573555,-0.3858919) (0.7639749,-0.3726155) (0.7703616,-0.3592255) (0.7765136,-0.3457261) (0.7824291,-0.3321215) (0.7881063,-0.3184156) (0.7935434,-0.3046128) (0.7987387,-0.2907171) (0.8036908,-0.2767329) (0.8083980,-0.2626644) (0.8128590,-0.2485159) (0.8170724,-0.2342918) (0.8210370,-0.2199962) (0.8247514,-0.2056336) (0.8282146,-0.1912084) (0.8314255,-0.1767249) (0.8343831,-0.1621876) (0.8370866,-0.1476010) (0.8395351,-0.1329693) (0.8417279,-0.1182971) (0.8436642,-0.1035889) (0.8453436,-0.0888492) (0.8467655,-0.0740824) (0.8479294,-0.0592930) (0.8488351,-0.0444856) (0.8494822,-0.0296646) (0.8498705,-0.0148345) (0.8500000,0.0000000) (0.8498705,0.0148345) (0.8494822,0.0296646) (0.8488351,0.0444856) (0.8479294,0.0592930) (0.8467655,0.0740824) (0.8453436,0.0888492) (0.8436642,0.1035889) (0.8417279,0.1182971) (0.8395351,0.1329693) (0.8370866,0.1476010) (0.8343831,0.1621876) (0.8314255,0.1767249) (0.8282146,0.1912084) (0.8247514,0.2056336) (0.8210370,0.2199962) (0.8170724,0.2342918) (0.8128590,0.2485159) (0.8083980,0.2626644) (0.8036908,0.2767329) (0.7987387,0.2907171) (0.7935434,0.3046128) (0.7881063,0.3184156) (0.7824291,0.3321215) (0.7765136,0.3457261) (0.7703616,0.3592255) (0.7639749,0.3726155) (0.7573555,0.3858919) (0.7505055,0.3990508) (0.7434268,0.4120882) (0.7361216,0.4250000) (0.7285922,0.4377824) (0.7208409,0.4504314) (0.7128700,0.4629432) (0.7046819,0.4753140) (0.6962792,0.4875400) (0.6876644,0.4996175) (0.6788402,0.5115428) (0.6698091,0.5233123) (0.6605741,0.5349223) (0.6511378,0.5463695) (0.6415031,0.5576502) (0.6316731,0.5687610) (0.6216506,0.5796986) (0.6114388,0.5904596) (0.6010408,0.6010408) (0.5904596,0.6114388) (0.5796986,0.6216506) (0.5687610,0.6316731) (0.5576502,0.6415031) (0.5463695,0.6511378)};
\draw[thick] plot coordinates {(0.8034845,-0.9575556) (0.8200738,-0.9433870) (0.8364133,-0.9289310) (0.8524980,-0.9141921) (0.8683230,-0.8991748) (0.8838835,-0.8838835) (0.8991748,-0.8683230) (0.9141921,-0.8524980) (0.9289310,-0.8364133) (0.9433870,-0.8200738) (0.9575556,-0.8034845) (0.9714325,-0.7866505) (0.9850134,-0.7695768) (0.9982944,-0.7522688) (1.0112712,-0.7347316) (1.0239401,-0.7169705) (1.0362970,-0.6989911) (1.0483382,-0.6807988) (1.0600601,-0.6623991) (1.0714591,-0.6437976) (1.0825318,-0.6250000) (1.0932746,-0.6060120) (1.1036845,-0.5868395) (1.1137582,-0.5674881) (1.1234926,-0.5479639) (1.1328847,-0.5282728) (1.1419318,-0.5084208) (1.1506311,-0.4884139) (1.1589798,-0.4682582) (1.1669755,-0.4479599) (1.1746158,-0.4275252) (1.1818982,-0.4069602) (1.1888206,-0.3862712) (1.1953809,-0.3654646) (1.2015771,-0.3445467) (1.2074073,-0.3235238) (1.2128697,-0.3024024) (1.2179626,-0.2811888) (1.2226845,-0.2598896) (1.2270340,-0.2385112) (1.2310097,-0.2170602) (1.2346104,-0.1955431) (1.2378351,-0.1739664) (1.2406827,-0.1523367) (1.2431524,-0.1306606) (1.2452434,-0.1089447) (1.2469551,-0.0871956) (1.2482869,-0.0654199) (1.2492385,-0.0436244) (1.2498096,-0.0218155) (1.2500000,0.0000000) (1.2498096,0.0218155) (1.2492385,0.0436244) (1.2482869,0.0654199) (1.2469551,0.0871956) (1.2452434,0.1089447) (1.2431524,0.1306606) (1.2406827,0.1523367) (1.2378351,0.1739664) (1.2346104,0.1955431) (1.2310097,0.2170602) (1.2270340,0.2385112) (1.2226845,0.2598896) (1.2179626,0.2811888) (1.2128697,0.3024024) (1.2074073,0.3235238) (1.2015771,0.3445467) (1.1953809,0.3654646) (1.1888206,0.3862712) (1.1818982,0.4069602) (1.1746158,0.4275252) (1.1669755,0.4479599) (1.1589798,0.4682582) (1.1506311,0.4884139) (1.1419318,0.5084208) (1.1328847,0.5282728) (1.1234926,0.5479639) (1.1137582,0.5674881) (1.1036845,0.5868395) (1.0932746,0.6060120) (1.0825318,0.6250000) (1.0714591,0.6437976) (1.0600601,0.6623991) (1.0483382,0.6807988) (1.0362970,0.6989911) (1.0239401,0.7169705) (1.0112712,0.7347316) (0.9982944,0.7522688) (0.9850134,0.7695768) (0.9714325,0.7866505) (0.9575556,0.8034845) (0.9433870,0.8200738) (0.9289310,0.8364133) (0.9141921,0.8524980) (0.8991748,0.8683230) (0.8838835,0.8838835) (0.8683230,0.8991748) (0.8524980,0.9141921) (0.8364133,0.9289310) (0.8200738,0.9433870) (0.8034845,0.9575556)};
\draw[DeepBlue,thick,->] (0.0000000,0.0000000)--(1.6500000,0.0000000);
\node[below,fill=none,inner sep=2pt] at (0.8000000,-1.1500000) {同心波面与法线};
\end{scope}
\end{tikzpicture}
\figcap{观测点 $O$ 在地表；按到眼光线的切线反向看，太阳视方向高于真实方向。左侧为同心波面。}\label{ch27:fig:s42-19-1}
\end{minipage}\end{center}
}

\newcommand{\ReflectionPanel}{%
\begin{center}\begin{minipage}{\linewidth}\centering
\begin{tikzpicture}[>=Latex,font=\small,every node/.style={fill=none}]
\begin{scope}[shift={(-4.1000000,0.0000000)}]
\draw[DeepBlue,thick,->] (-1.4000000,1.4000000)--(0.0000000,0.0000000);
\draw[DeepBlue,thick,->] (0.0000000,0.0000000)--(1.4000000,1.4000000);
\draw[thick] (-1.8000000,0.0000000)--(1.8000000,0.0000000);
\draw[gray,dashed,thin] (0.0000000,-0.2000000)--(0.0000000,1.8000000);
\node[above,fill=none,inner sep=2pt] at (-0.4000000,0.8000000) {$i$};
\node[above,fill=none,inner sep=2pt] at (0.4000000,0.8000000) {$r$};
\node[below,fill=none,inner sep=2pt] at (0.0000000,-0.3500000) {反射角等于入射角};
\end{scope}
\begin{scope}[shift={(0.3000000,0.0000000)}]
\draw[thick] (-1.9336619,-0.4762222)--(0.1336619,0.2762222);
\draw[thick] (-0.5547819,1.4332399)--(1.5125419,0.6807955);
\draw[DeepBlue,thick,->] (-2.4000000,-0.1000000)--(-0.9000000,-0.1000000);
\draw[DeepBlue,thick,->] (-0.9000000,-0.1000000)--(0.4788800,1.0570177);
\draw[DeepBlue,thick,->] (0.4788800,1.0570177)--(0.7219874,-0.3217132);
\node[below,fill=none,inner sep=2pt] at (0.0000000,-0.7000000) {两镜夹角 $40^\circ$};
\node[below,fill=none,inner sep=2pt] at (0.0000000,-1.1000000) {方向转过 $80^\circ$};
\end{scope}
\begin{scope}[shift={(4.7000000,0.0000000)}]
\draw[thick] (-1.2000000,1.3000000)--(-1.2000000,-1.2000000)--(1.6000000,-1.2000000);
\draw[DeepBlue,thick,->] (1.5000000,1.0000000)--(-1.2000000,-0.3500000);
\draw[DeepBlue,thick,->] (-1.2000000,-0.3500000)--(0.5000000,-1.2000000);
\draw[DeepBlue,thick,->] (0.5000000,-1.2000000)--(2.0000000,-0.4500000);
\node[below,fill=none,inner sep=2pt] at (0.1500000,-1.5000000) {直角反射器};
\end{scope}
\end{tikzpicture}
\figcap{单镜、双镜与直角反射器。各段在镜面交点处等角反射。}\label{ch27:fig:s42-19-3}
\end{minipage}\end{center}
}

\newcommand{\RefractionPrismPanel}{%
\begin{center}\begin{minipage}{\linewidth}\centering
\begin{tikzpicture}[>=Latex,font=\small,every node/.style={fill=none}]
\begin{scope}[shift={(-3.2000000,0.0000000)}]
\draw[thick] (-2.0000000,0.0000000)--(2.0000000,0.0000000);
\draw[gray,dashed,thin] (0.0000000,-1.8000000)--(0.0000000,1.8000000);
\draw[DeepBlue,thick,->] (-1.5000000,1.5000000)--(0.0000000,0.0000000);
\draw[DeepBlue,thick,->] (0.0000000,0.0000000)--(0.8017840,-1.5000000);
\node[above,fill=none,inner sep=2pt] at (-1.6000000,0.4000000) {$n_1=1$};
\node[above,fill=none,inner sep=2pt] at (-1.5000000,-1.1000000) {$n_2=1.5$};
\node[above,fill=none,inner sep=2pt] at (-0.3500000,0.9000000) {$i$};
\node[right,fill=none,inner sep=2pt] at (0.7500000,-0.8500000) {$r$};
\draw[gray,dashed,thin] (0.7200000,-0.8500000)--(0.2000000,-0.7000000);
\end{scope}
\begin{scope}[shift={(3.0000000,0.0000000)}]
\draw[thick] (-1.5000000,-1.0000000)--(0.0000000,1.5980760)--(1.5000000,-1.0000000)--(-1.5000000,-1.0000000);
\draw[DeepBlue,thick,->] (-2.4167961,-0.3128935)--(-0.8000000,0.2124354);
\draw[gray,dashed,thin] (-1.2763140,0.4874354)--(-0.1937822,-0.1375646);
\draw[red!65!black,thick,->] (-0.8000000,0.2124354)--(0.8048823,0.2039789);
\draw[red!65!black,thick,->] (0.8048823,0.2039789)--(2.1743453,-0.2725409);
\draw[gray,dashed,thin] (0.3718696,-0.0460211)--(1.3677988,0.5289789);
\draw[Blue,thick,->] (-0.8000000,0.2124354)--(0.8153120,0.1859141);
\draw[Blue,thick,->] (0.8153120,0.1859141)--(2.1609251,-0.3542946);
\node[above,fill=none,inner sep=2pt] at (0.0000000,1.0000000) {$A$};
\node[below,fill=none,inner sep=2pt] at (0.0000000,-1.3000000) {各颜色只在界面改变方向};
\end{scope}
\end{tikzpicture}
\figcap{平界面折射与棱镜色散；虚线为界面法线。}\label{ch27:fig:s42-19-6}
\end{minipage}\end{center}
}

\newcommand{\PrismDeviationDiagram}{%
\begin{center}\begin{minipage}{\linewidth}\centering
\begin{tikzpicture}[>=Latex,font=\small,every node/.style={fill=none}]
\draw[gray,thin,->] (0.0000000,0.0000000)--(6.1000000,0.0000000);
\draw[gray,thin,->] (0.0000000,0.0000000)--(0.0000000,3.3000000);
\draw[DeepBlue,thick] plot coordinates {(0.6601178,2.2650478) (0.6799715,1.8798761) (0.6998524,1.7231931) (0.7197610,1.6081283) (0.7396977,1.5145945) (0.7596629,1.4348059) (0.7796570,1.3647702) (0.7996804,1.3021169) (0.8197336,1.2453026) (0.8398170,1.1932560) (0.8599310,1.1451978) (0.8800762,1.1005403) (0.9002529,1.0588274) (0.9204616,1.0196967) (0.9407028,0.9828547) (0.9609769,0.9480600) (0.9812845,0.9151108) (1.0016261,0.8838370) (1.0220022,0.8540931) (1.0424132,0.8257539) (1.0628597,0.7987107) (1.0833423,0.7728683) (1.1038615,0.7481429) (1.1244178,0.7244604) (1.1450117,0.7017545) (1.1656440,0.6799662) (1.1863151,0.6590424) (1.2070256,0.6389350) (1.2277763,0.6196005) (1.2485675,0.6009996) (1.2694001,0.5830962) (1.2902746,0.5658573) (1.3111918,0.5492529) (1.3321521,0.5332550) (1.3531565,0.5178383) (1.3742054,0.5029790) (1.3952997,0.4886554) (1.4164401,0.4748471) (1.4376272,0.4615356) (1.4588620,0.4487032) (1.4801450,0.4363337) (1.5014772,0.4244120) (1.5228592,0.4129239) (1.5442920,0.4018562) (1.5657763,0.3911966) (1.5873131,0.3809334) (1.6089031,0.3710558) (1.6305473,0.3615536) (1.6522466,0.3524172) (1.6740020,0.3436378) (1.6958143,0.3352069) (1.7176845,0.3271166) (1.7396137,0.3193594) (1.7616029,0.3119284) (1.7836531,0.3048170) (1.8057654,0.2980190) (1.8279409,0.2915288) (1.8501807,0.2853407) (1.8724860,0.2794497) (1.8948580,0.2738511) (1.9172978,0.2685403) (1.9398067,0.2635132) (1.9623860,0.2587658) (1.9850371,0.2542945) (2.0077611,0.2500959) (2.0305596,0.2461668) (2.0534339,0.2425043) (2.0763854,0.2391057) (2.0994158,0.2359685) (2.1225265,0.2330903) (2.1457191,0.2304693) (2.1689952,0.2281034) (2.1923565,0.2259909) (2.2158047,0.2241305) (2.2393416,0.2225207) (2.2629690,0.2211605) (2.2866888,0.2200489) (2.3105029,0.2191851) (2.3344134,0.2185686) (2.3584222,0.2181988) (2.3825315,0.2180756) (2.4067435,0.2181988) (2.4310604,0.2185686) (2.4554847,0.2191851) (2.4800186,0.2200489) (2.5046648,0.2211605) (2.5294257,0.2225207) (2.5543040,0.2241305) (2.5793026,0.2259909) (2.6044243,0.2281034) (2.6296720,0.2304693) (2.6550488,0.2330903) (2.6805579,0.2359685) (2.7062026,0.2391057) (2.7319864,0.2425043) (2.7579128,0.2461668) (2.7839855,0.2500959) (2.8102084,0.2542945) (2.8365855,0.2587658) (2.8631209,0.2635132) (2.8898191,0.2685403) (2.9166846,0.2738511) (2.9437221,0.2794497) (2.9709365,0.2853407) (2.9983330,0.2915288) (3.0259171,0.2980190) (3.0536944,0.3048170) (3.0816707,0.3119284) (3.1098524,0.3193594) (3.1382460,0.3271166) (3.1668582,0.3352069) (3.1956962,0.3436378) (3.2247677,0.3524172) (3.2540806,0.3615536) (3.2836434,0.3710558) (3.3134647,0.3809334) (3.3435542,0.3911966) (3.3739215,0.4018562) (3.4045774,0.4129239) (3.4355329,0.4244120) (3.4667998,0.4363337) (3.4983907,0.4487032) (3.5303191,0.4615356) (3.5625992,0.4748471) (3.5952464,0.4886554) (3.6282771,0.5029790) (3.6617088,0.5178383) (3.6955604,0.5332550) (3.7298523,0.5492529) (3.7646065,0.5658573) (3.7998467,0.5830962) (3.8355988,0.6009996) (3.8718909,0.6196005) (3.9087535,0.6389350) (3.9462202,0.6590424) (3.9843279,0.6799662) (4.0231170,0.7017545) (4.0626325,0.7244604) (4.1029243,0.7481429) (4.1440479,0.7728683) (4.1860658,0.7987107) (4.2290484,0.8257539) (4.2730754,0.8540931) (4.3182380,0.8838370) (4.3646411,0.9151108) (4.4124064,0.9480600) (4.4616762,0.9828547) (4.5126190,1.0196967) (4.5654366,1.0588274) (4.6203741,1.1005403) (4.6777338,1.1451978) (4.7378963,1.1932560) (4.8013519,1.2453026) (4.8687504,1.3021169) (4.9409849,1.3647702) (5.0193421,1.4348059) (5.1057977,1.5145945) (5.2036792,1.6081283) (5.3194752,1.7231931) (5.4699253,1.8798761) (5.8107553,2.2650478)};
\fill (2.3825315,0.2180756) circle (1.2pt);
\node[below,fill=none,inner sep=2pt] at (2.3825315,0.2180756) {};
\draw[gray,dashed,thin] (2.3825315,0.2180756)--(2.4000000,1.1500000);
\node[above,fill=none,inner sep=2pt] at (2.4000000,1.2000000) {对称光路};
\node[below,fill=none,inner sep=2pt] at (0.8333333,-0.0800000) {$30^\circ$};
\node[below,fill=none,inner sep=2pt] at (2.5000000,-0.0800000) {$50^\circ$};
\node[below,fill=none,inner sep=2pt] at (4.1666667,-0.0800000) {$70^\circ$};
\node[below,fill=none,inner sep=2pt] at (5.8333333,-0.0800000) {$90^\circ$};
\node[left,fill=none,inner sep=2pt] at (-0.0500000,0.5000000) {$40^\circ$};
\node[left,fill=none,inner sep=2pt] at (-0.0500000,1.5000000) {$50^\circ$};
\node[left,fill=none,inner sep=2pt] at (-0.0500000,2.5000000) {$60^\circ$};
\node[right,fill=none,inner sep=2pt] at (6.1000000,0.0000000) {$i_1$};
\node[above,fill=none,inner sep=2pt] at (0.0000000,3.3000000) {$\delta$};
\node[above,fill=none,inner sep=2pt] at (4.0000000,2.8000000) {$A=60^\circ,\ n=1.5$};
\end{tikzpicture}
\figcap{偏向角随入射角的变化；最低点是两内部折射角相等的光路。\label{ch27:fig:deviation}}
\end{minipage}\end{center}
}

\newcommand{\FiberFermatPanel}{%
\begin{center}\begin{minipage}{\linewidth}\centering
\begin{tikzpicture}[>=Latex,font=\small,every node/.style={fill=none}]
\begin{scope}[shift={(-3.7000000,0.0000000)}]
\draw[thick] (-2.6000000,-0.6500000)--(2.6000000,-0.6500000);
\draw[thick] (-2.6000000,0.6500000)--(2.6000000,0.6500000);
\draw[gray,dashed,thin] (-2.6000000,-0.9000000)--(2.6000000,-0.9000000);
\draw[gray,dashed,thin] (-2.6000000,0.9000000)--(2.6000000,0.9000000);
\draw[DeepBlue,thick,->] (-2.6000000,0.0000000)--(-1.9500000,0.6500000);
\draw[DeepBlue,thick,->] (-1.9500000,0.6500000)--(-0.6500000,-0.6500000);
\draw[DeepBlue,thick,->] (-0.6500000,-0.6500000)--(0.6500000,0.6500000);
\draw[DeepBlue,thick,->] (0.6500000,0.6500000)--(1.9500000,-0.6500000);
\draw[DeepBlue,thick,->] (1.9500000,-0.6500000)--(2.6000000,0.0000000);
\node[above,fill=none,inner sep=2pt] at (0.0000000,1.0500000) {包层 $n_2<n_1$};
\node[below,fill=none,inner sep=2pt] at (0.0000000,-1.0500000) {芯层内，反射点均在边界};
\end{scope}
\begin{scope}[shift={(3.0000000,0.0000000)}]
\draw[thick] plot coordinates {(0.9844742,-1.0500000) (0.9383572,-1.0290000) (0.8944127,-1.0080000) (0.8524231,-0.9870000) (0.8122095,-0.9660000) (0.7736233,-0.9450000) (0.7365391,-0.9240000) (0.7008507,-0.9030000) (0.6664670,-0.8820000) (0.6333096,-0.8610000) (0.6013107,-0.8400000) (0.5704111,-0.8190000) (0.5405591,-0.7980000) (0.5117096,-0.7770000) (0.4838228,-0.7560000) (0.4568639,-0.7350000) (0.4308022,-0.7140000) (0.4056108,-0.6930000) (0.3812660,-0.6720000) (0.3577471,-0.6510000) (0.3350360,-0.6300000) (0.3131170,-0.6090000) (0.2919765,-0.5880000) (0.2716029,-0.5670000) (0.2519862,-0.5460000) (0.2331184,-0.5250000) (0.2149924,-0.5040000) (0.1976029,-0.4830000) (0.1809455,-0.4620000) (0.1650171,-0.4410000) (0.1498155,-0.4200000) (0.1353393,-0.3990000) (0.1215882,-0.3780000) (0.1085622,-0.3570000) (0.0962623,-0.3360000) (0.0846897,-0.3150000) (0.0738464,-0.2940000) (0.0637344,-0.2730000) (0.0543561,-0.2520000) (0.0457143,-0.2310000) (0.0378117,-0.2100000) (0.0306511,-0.1890000) (0.0242354,-0.1680000) (0.0185671,-0.1470000) (0.0136488,-0.1260000) (0.0094830,-0.1050000) (0.0060716,-0.0840000) (0.0034163,-0.0630000) (0.0015187,-0.0420000) (0.0003797,-0.0210000) (0.0000000,0.0000000) (0.0003797,0.0210000) (0.0015187,0.0420000) (0.0034163,0.0630000) (0.0060716,0.0840000) (0.0094830,0.1050000) (0.0136488,0.1260000) (0.0185671,0.1470000) (0.0242354,0.1680000) (0.0306511,0.1890000) (0.0378117,0.2100000) (0.0457143,0.2310000) (0.0543561,0.2520000) (0.0637344,0.2730000) (0.0738464,0.2940000) (0.0846897,0.3150000) (0.0962623,0.3360000) (0.1085622,0.3570000) (0.1215882,0.3780000) (0.1353393,0.3990000) (0.1498155,0.4200000) (0.1650171,0.4410000) (0.1809455,0.4620000) (0.1976029,0.4830000) (0.2149924,0.5040000) (0.2331184,0.5250000) (0.2519862,0.5460000) (0.2716029,0.5670000) (0.2919765,0.5880000) (0.3131170,0.6090000) (0.3350360,0.6300000) (0.3577471,0.6510000) (0.3812660,0.6720000) (0.4056108,0.6930000) (0.4308022,0.7140000) (0.4568639,0.7350000) (0.4838228,0.7560000) (0.5117096,0.7770000) (0.5405591,0.7980000) (0.5704111,0.8190000) (0.6013107,0.8400000) (0.6333096,0.8610000) (0.6664670,0.8820000) (0.7008507,0.9030000) (0.7365391,0.9240000) (0.7736233,0.9450000) (0.8122095,0.9660000) (0.8524231,0.9870000) (0.8944127,1.0080000) (0.9383572,1.0290000) (0.9844742,1.0500000)};
\draw[gray,thin,->] (-2.3000000,0.0000000)--(3.4000000,0.0000000);
\fill (-2.0000000,0.0000000) circle (1.2pt);
\node[below,fill=none,inner sep=2pt] at (-2.0000000,0.0000000) {$O$};
\fill (3.0000000,0.0000000) circle (1.2pt);
\node[below,fill=none,inner sep=2pt] at (3.0000000,0.0000000) {$I$};
\draw[DeepBlue,thick,->] (-2.0000000,0.0000000)--(0.6164078,0.8500000);
\draw[DeepBlue,thick,->] (0.6164078,0.8500000)--(3.0000000,0.0000000);
\fill (0.6164078,0.8500000) circle (1.2pt);
\node[above,fill=none,inner sep=2pt] at (0.6164078,0.8500000) {$P$};
\node[below,fill=none,inner sep=2pt] at (0.0000000,-1.3500000) {$OP+nPI=s_o+ns_i$};
\end{scope}
\end{tikzpicture}
\figcap{左为平直芯包界面导光；右为等光程非球面，离轴点 $P$ 也连接同一物、像点。}\label{ch27:fig:s42-19-10}
\end{minipage}\end{center}
}

\newcommand{\SphericalMirrorRotatingDiagram}{%
\begin{center}\begin{minipage}{\linewidth}\centering
\begin{tikzpicture}[>=Latex,font=\small,every node/.style={fill=none}]
\begin{scope}[shift={(-3.0000000,0.0000000)}]
\draw[thick] plot coordinates {(2.3640323,-1.8469844) (2.3883229,-1.8154652) (2.4121933,-1.7836265) (2.4356393,-1.7514739) (2.4586568,-1.7190133) (2.4812417,-1.6862501) (2.5033901,-1.6531903) (2.5250979,-1.6198396) (2.5463615,-1.5862039) (2.5671771,-1.5522892) (2.5875410,-1.5181013) (2.6074496,-1.4836463) (2.6268995,-1.4489303) (2.6458872,-1.4139593) (2.6644093,-1.3787396) (2.6824627,-1.3432773) (2.7000441,-1.3075786) (2.7171505,-1.2716499) (2.7337787,-1.2354974) (2.7499260,-1.1991276) (2.7655895,-1.1625468) (2.7807663,-1.1257614) (2.7954539,-1.0887780) (2.8096497,-1.0516030) (2.8233511,-1.0142429) (2.8365557,-0.9767045) (2.8492613,-0.9389941) (2.8614656,-0.9011186) (2.8731664,-0.8630845) (2.8843617,-0.8248986) (2.8950495,-0.7865675) (2.9052279,-0.7480981) (2.9148952,-0.7094970) (2.9240496,-0.6707711) (2.9326896,-0.6319271) (2.9408135,-0.5929720) (2.9484201,-0.5539126) (2.9555078,-0.5147557) (2.9620756,-0.4755082) (2.9681222,-0.4361771) (2.9736466,-0.3967692) (2.9786478,-0.3572915) (2.9831249,-0.3177509) (2.9870772,-0.2781545) (2.9905039,-0.2385091) (2.9934044,-0.1988217) (2.9957783,-0.1590994) (2.9976250,-0.1193490) (2.9989444,-0.0795777) (2.9997361,-0.0397923) (3.0000000,0.0000000) (2.9997361,0.0397923) (2.9989444,0.0795777) (2.9976250,0.1193490) (2.9957783,0.1590994) (2.9934044,0.1988217) (2.9905039,0.2385091) (2.9870772,0.2781545) (2.9831249,0.3177509) (2.9786478,0.3572915) (2.9736466,0.3967692) (2.9681222,0.4361771) (2.9620756,0.4755082) (2.9555078,0.5147557) (2.9484201,0.5539126) (2.9408135,0.5929720) (2.9326896,0.6319271) (2.9240496,0.6707711) (2.9148952,0.7094970) (2.9052279,0.7480981) (2.8950495,0.7865675) (2.8843617,0.8248986) (2.8731664,0.8630845) (2.8614656,0.9011186) (2.8492613,0.9389941) (2.8365557,0.9767045) (2.8233511,1.0142429) (2.8096497,1.0516030) (2.7954539,1.0887780) (2.7807663,1.1257614) (2.7655895,1.1625468) (2.7499260,1.1991276) (2.7337787,1.2354974) (2.7171505,1.2716499) (2.7000441,1.3075786) (2.6824627,1.3432773) (2.6644093,1.3787396) (2.6458872,1.4139593) (2.6268995,1.4489303) (2.6074496,1.4836463) (2.5875410,1.5181013) (2.5671771,1.5522892) (2.5463615,1.5862039) (2.5250979,1.6198396) (2.5033901,1.6531903) (2.4812417,1.6862501) (2.4586568,1.7190133) (2.4356393,1.7514739) (2.4121933,1.7836265) (2.3883229,1.8154652) (2.3640323,1.8469844)};
\draw[gray,thin,->] (-0.7000000,0.0000000)--(3.6000000,0.0000000);
\fill (0.0000000,0.0000000) circle (1.2pt);
\node[below,fill=none,inner sep=2pt] at (0.0000000,0.0000000) {$C$};
\fill (3.0000000,0.0000000) circle (1.2pt);
\node[below right,fill=none,inner sep=2pt] at (3.0000000,0.0000000) {$V$};
\fill (1.6550669,0.0000000) circle (1.2pt);
\node[below,fill=none,inner sep=2pt] at (1.6550669,0.0000000) {$M$};
\draw[DeepBlue,thick,->] (-1.0000000,1.2678548)--(2.7189234,1.2678548);
\draw[DeepBlue,thick,->] (2.7189234,1.2678548)--(1.6550669,0.0000000);
\draw[gray,dashed,thin] (0.0000000,0.0000000)--(2.7189234,1.2678548);
\fill (2.7189234,1.2678548) circle (1.2pt);
\node[below,fill=none,inner sep=2pt] at (2.7189234,1.2678548) {};
\draw[gray,dashed,thin] (2.7189234,1.2678548)--(2.1689234,1.7178548);
\node[above left,fill=none,inner sep=2pt] at (2.1189234,1.7678548) {$P$};
\end{scope}
\begin{scope}[shift={(4.0000000,0.0000000)}]
\draw[thick] (0.0000000,-0.3000000)--(0.0000000,2.6000000);
\draw[thick] (0.0000000,0.0000000)--(-0.9500000,2.3513325);
\fill (0.0000000,0.0000000) circle (1.2pt);
\node[below,fill=none,inner sep=2pt] at (0.0000000,0.0000000) {$O$};
\fill (0.0000000,1.0000000) circle (1.2pt);
\node[right,fill=none,inner sep=2pt] at (0.0000000,1.0000000) {$A$};
\fill (-0.4040262,1.0000000) circle (1.2pt);
\node[left,fill=none,inner sep=2pt] at (-0.4040262,1.0000000) {$C$};
\fill (0.0000000,1.3901636) circle (1.2pt);
\node[right,fill=none,inner sep=2pt] at (0.0000000,1.3901636) {$D$};
\draw[DeepBlue,thick,->] (0.0000000,1.0000000)--(-0.4040262,1.0000000);
\draw[DeepBlue,thick,->] (-0.4040262,1.0000000)--(0.0000000,1.3901636);
\draw[thick] plot coordinates {(0.0000000,0.4200000) (-0.0016127,0.4199969) (-0.0032253,0.4199876) (-0.0048379,0.4199721) (-0.0064505,0.4199505) (-0.0080629,0.4199226) (-0.0096752,0.4198885) (-0.0112874,0.4198483) (-0.0128994,0.4198019) (-0.0145113,0.4197492) (-0.0161229,0.4196904) (-0.0177343,0.4196254) (-0.0193454,0.4195542) (-0.0209562,0.4194769) (-0.0225667,0.4193933) (-0.0241769,0.4193036) (-0.0257867,0.4192076) (-0.0273962,0.4191055) (-0.0290052,0.4189973) (-0.0306138,0.4188828) (-0.0322220,0.4187622) (-0.0338297,0.4186353) (-0.0354369,0.4185024) (-0.0370435,0.4183632) (-0.0386497,0.4182179) (-0.0402552,0.4180664) (-0.0418602,0.4179088) (-0.0434645,0.4177449) (-0.0450682,0.4175750) (-0.0466713,0.4173988) (-0.0482736,0.4172166) (-0.0498752,0.4170281) (-0.0514761,0.4168335) (-0.0530763,0.4166328) (-0.0546756,0.4164260) (-0.0562742,0.4162129) (-0.0578719,0.4159938) (-0.0594688,0.4157685) (-0.0610648,0.4155371) (-0.0626599,0.4152996) (-0.0642541,0.4150559) (-0.0658473,0.4148061) (-0.0674395,0.4145503) (-0.0690308,0.4142882) (-0.0706210,0.4140201) (-0.0722102,0.4137459) (-0.0737984,0.4134656) (-0.0753854,0.4131792) (-0.0769713,0.4128867) (-0.0785561,0.4125881) (-0.0801398,0.4122834) (-0.0817222,0.4119727) (-0.0833035,0.4116558) (-0.0848835,0.4113329) (-0.0864623,0.4110040) (-0.0880398,0.4106690) (-0.0896160,0.4103279) (-0.0911909,0.4099808) (-0.0927644,0.4096276) (-0.0943366,0.4092684) (-0.0959074,0.4089031) (-0.0974767,0.4085319) (-0.0990447,0.4081546) (-0.1006111,0.4077713) (-0.1021761,0.4073819) (-0.1037396,0.4069866) (-0.1053015,0.4065853) (-0.1068619,0.4061780) (-0.1084207,0.4057646) (-0.1099780,0.4053453) (-0.1115336,0.4049201) (-0.1130875,0.4044888) (-0.1146398,0.4040516) (-0.1161904,0.4036085) (-0.1177393,0.4031593) (-0.1192864,0.4027043) (-0.1208318,0.4022433) (-0.1223754,0.4017764) (-0.1239172,0.4013035) (-0.1254572,0.4008248) (-0.1269954,0.4003401) (-0.1285316,0.3998495) (-0.1300660,0.3993530) (-0.1315984,0.3988507) (-0.1331289,0.3983424) (-0.1346575,0.3978283) (-0.1361840,0.3973083) (-0.1377086,0.3967825) (-0.1392311,0.3962508) (-0.1407515,0.3957133) (-0.1422699,0.3951699) (-0.1437862,0.3946207) (-0.1453004,0.3940657) (-0.1468124,0.3935049) (-0.1483223,0.3929383) (-0.1498300,0.3923659) (-0.1513354,0.3917877) (-0.1528387,0.3912037) (-0.1543396,0.3906140) (-0.1558384,0.3900185) (-0.1573348,0.3894172)};
\node[left,fill=none,inner sep=2pt] at (-0.8000000,0.3500000) {$\theta$};
\draw[gray,dashed,thin] (-0.7500000,0.4000000)--(-0.0800000,0.4200000);
\draw[gray,<->] (0.6000000,0.0000000)--(0.6000000,1.0000000);
\node[right,fill=none,inner sep=2pt] at (0.6000000,0.5000000) {$d$};
\end{scope}
\end{tikzpicture}
\figcap{左侧球面镜的 $P$ 在圆周上，$CP$ 为法线；右侧转镜的转轴 $O$ 位于墙上。}\label{ch27:fig:s42-19-11}
\end{minipage}\end{center}
}

\newcommand{\CylinderDiagram}{%
\begin{center}\begin{minipage}{\linewidth}\centering
\begin{tikzpicture}[>=Latex,font=\small,every node/.style={fill=none}]
\draw[thick] (0.0000000,0.0000000) circle (2);
\draw[thick] plot coordinates {(1.8131983,-0.9238707) (1.8218249,-0.9067409) (1.8302897,-0.8895306) (1.8385919,-0.8722413) (1.8467308,-0.8548745) (1.8547057,-0.8374317) (1.8625158,-0.8199146) (1.8701605,-0.8023246) (1.8776390,-0.7846634) (1.8849508,-0.7669325) (1.8920951,-0.7491335) (1.8990714,-0.7312679) (1.9058790,-0.7133373) (1.9125173,-0.6953434) (1.9189857,-0.6772878) (1.9252837,-0.6591719) (1.9314107,-0.6409976) (1.9373660,-0.6227663) (1.9431493,-0.6044796) (1.9487600,-0.5861393) (1.9541977,-0.5677469) (1.9594617,-0.5493041) (1.9645516,-0.5308125) (1.9694671,-0.5122737) (1.9742076,-0.4936894) (1.9787728,-0.4750613) (1.9831622,-0.4563910) (1.9873754,-0.4376801) (1.9914121,-0.4189304) (1.9952720,-0.4001435) (1.9989546,-0.3813210) (2.0024596,-0.3624646) (2.0057868,-0.3435761) (2.0089357,-0.3246570) (2.0119063,-0.3057091) (2.0146981,-0.2867340) (2.0173110,-0.2677335) (2.0197447,-0.2487092) (2.0219990,-0.2296627) (2.0240737,-0.2105959) (2.0259686,-0.1915104) (2.0276835,-0.1724079) (2.0292184,-0.1532900) (2.0305729,-0.1341586) (2.0317472,-0.1150152) (2.0327409,-0.0958616) (2.0335541,-0.0766995) (2.0341866,-0.0575306) (2.0346385,-0.0383566) (2.0349096,-0.0191791) (2.0350000,0.0000000) (2.0349096,0.0191791) (2.0346385,0.0383566) (2.0341866,0.0575306) (2.0335541,0.0766995) (2.0327409,0.0958616) (2.0317472,0.1150152) (2.0305729,0.1341586) (2.0292184,0.1532900) (2.0276835,0.1724079) (2.0259686,0.1915104) (2.0240737,0.2105959) (2.0219990,0.2296627) (2.0197447,0.2487092) (2.0173110,0.2677335) (2.0146981,0.2867340) (2.0119063,0.3057091) (2.0089357,0.3246570) (2.0057868,0.3435761) (2.0024596,0.3624646) (1.9989546,0.3813210) (1.9952720,0.4001435) (1.9914121,0.4189304) (1.9873754,0.4376801) (1.9831622,0.4563910) (1.9787728,0.4750613) (1.9742076,0.4936894) (1.9694671,0.5122737) (1.9645516,0.5308125) (1.9594617,0.5493041) (1.9541977,0.5677469) (1.9487600,0.5861393) (1.9431493,0.6044796) (1.9373660,0.6227663) (1.9314107,0.6409976) (1.9252837,0.6591719) (1.9189857,0.6772878) (1.9125173,0.6953434) (1.9058790,0.7133373) (1.8990714,0.7312679) (1.8920951,0.7491335) (1.8849508,0.7669325) (1.8776390,0.7846634) (1.8701605,0.8023246) (1.8625158,0.8199146) (1.8547057,0.8374317) (1.8467308,0.8548745) (1.8385919,0.8722413) (1.8302897,0.8895306) (1.8218249,0.9067409) (1.8131983,0.9238707)};
\draw[gray,thin,->] (-2.4000000,0.0000000)--(2.5000000,0.0000000);
\draw[DeepBlue,thick,->] (-3.2000000,1.0000000)--(-1.7320508,1.0000000);
\draw[DeepBlue,thick,->] (-1.7320508,1.0000000)--(2.0000000,0.0000000);
\draw[DeepBlue,thick,->] (2.0000000,0.0000000)--(-1.7320508,-1.0000000);
\draw[DeepBlue,thick,->] (-1.7320508,-1.0000000)--(-3.2000000,-1.0000000);
\draw[gray,dashed,thin] (-1.7320508,1.0000000)--(0.0000000,0.0000000);
\draw[gray,dashed,thin] (-1.7320508,-1.0000000)--(0.0000000,0.0000000);
\fill (-1.7320508,1.0000000) circle (1.2pt);
\node[above left,fill=none,inner sep=2pt] at (-1.7320508,1.0000000) {$A$};
\fill (-1.7320508,-1.0000000) circle (1.2pt);
\node[below left,fill=none,inner sep=2pt] at (-1.7320508,-1.0000000) {$B$};
\fill (2.0000000,0.0000000) circle (1.2pt);
\node[right,fill=none,inner sep=2pt] at (2.0000000,0.0000000) {$P$};
\fill (0.0000000,0.0000000) circle (1.2pt);
\node[below,fill=none,inner sep=2pt] at (0.0000000,0.0000000) {$C$};
\draw[gray,<->] (-3.5000000,-1.0000000)--(-3.5000000,1.0000000);
\node[left,fill=none,inner sep=2pt] at (-3.5000000,0.0000000) {$d$};
\node[below,fill=none,inner sep=2pt] at (1.5000000,-2.2000000) {右侧镀银；两外部光线反向平行};
\end{tikzpicture}
\figcap{$d$ 是入射与出射两平行线的间距；圆内两段只在反射面 $P$ 处折转。}\label{ch27:fig:s42-19-12}
\end{minipage}\end{center}
}

\newcommand{\AbsorbingShellDiagram}{%
% Independent panel: original physical paths and local coordinates retained.
\begin{center}\begin{minipage}{\linewidth}\centering
\begin{tikzpicture}[>=Latex,font=\small,every node/.style={fill=none}]
\draw[thick] (0.0000000,0.0000000) circle (2);
\draw[very thick] (0.0000000,0.0000000) circle (0.8);
\draw[DeepBlue,thick,->] (-3.0000000,1.2000000)--(-1.6000000,1.2000000);
\draw[DeepBlue,thick,->] (-1.6000000,1.2000000)--(0.1839273,0.7785697);
\draw[gray,dashed,thin] (-1.6000000,1.2000000)--(0.0000000,0.0000000);
\draw[gray,dashed,thin] (0.0000000,0.0000000)--(0.1839273,0.7785697);
\fill (0.1839273,0.7785697) circle (1.2pt);
\node[below,fill=none,inner sep=2pt] at (0.1839273,0.7785697) {};
\node[above right,fill=none,inner sep=2pt] at (0.4339273,1.0285697) {$T$};
\draw[gray,<->] (0.0000000,0.0000000)--(2.0000000,0.0000000);
\node[below,fill=none,inner sep=2pt] at (1.0000000,0.0000000) {$b$};
\node[above,fill=none,inner sep=2pt] at (-0.4000000,-0.2000000) {$a$};
\node[below,fill=none,inner sep=2pt] at (0.0000000,-2.2500000) {内球面为吸收层};
\end{tikzpicture}
\figcap{临界折射线在 $T$ 与半径为 $a$ 的吸收层相切，外球半径为 $b$。}\label{ch27:fig:s42-19-13}
\end{minipage}\end{center}
}

\newcommand{\GradientCircleDiagram}{%
% Independent panel: original physical paths and local coordinates retained.
\begin{center}\begin{minipage}{\linewidth}\centering
\begin{tikzpicture}[>=Latex,font=\small,every node/.style={fill=none}]
\draw[gray,thin,->] (-0.6000000,0.0000000)--(3.2000000,0.0000000);
\draw[gray,thin,->] (0.0000000,-0.3000000)--(0.0000000,3.2000000);
\node[right,fill=none,inner sep=2pt] at (3.2000000,0.0000000) {$x$};
\node[above,fill=none,inner sep=2pt] at (0.0000000,3.2000000) {$y$};
\draw[DeepBlue,thick,->] (-0.6000000,0.0000000)--(0.0000000,0.0000000);
\draw[DeepBlue,thick,->] plot coordinates {(0.0000000,0.0000000) (0.0347776,0.0002016) (0.0695505,0.0008063) (0.1043141,0.0018141) (0.1390637,0.0032249) (0.1737946,0.0050383) (0.2085021,0.0072543) (0.2431816,0.0098725) (0.2778284,0.0128925) (0.3124379,0.0163139) (0.3470054,0.0201364) (0.3815262,0.0243593) (0.4159958,0.0289821) (0.4504095,0.0340042) (0.4847626,0.0394249) (0.5190506,0.0452434) (0.5532689,0.0514591) (0.5874128,0.0580710) (0.6214777,0.0650783) (0.6554591,0.0724800) (0.6893525,0.0802752) (0.7231532,0.0884627) (0.7568567,0.0970415) (0.7904584,0.1060105) (0.8239540,0.1153683) (0.8573388,0.1251139) (0.8906084,0.1352458) (0.9237584,0.1457627) (0.9567841,0.1566632) (0.9896813,0.1679458) (1.0224455,0.1796090) (1.0550723,0.1916513) (1.0875573,0.2040710) (1.1198961,0.2168664) (1.1520844,0.2300358) (1.1841179,0.2435775) (1.2159923,0.2574897) (1.2477032,0.2717704) (1.2792465,0.2864178) (1.3106178,0.3014298) (1.3418130,0.3168045) (1.3728279,0.3325399) (1.4036583,0.3486337) (1.4343000,0.3650838) (1.4647490,0.3818880) (1.4950011,0.3990441) (1.5250523,0.4165497) (1.5548986,0.4344025) (1.5845359,0.4526001) (1.6139602,0.4711401) (1.6431677,0.4900199) (1.6721543,0.5092371) (1.7009162,0.5287889) (1.7294495,0.5486729) (1.7577504,0.5688863) (1.7858150,0.5894265) (1.8136397,0.6102906) (1.8412206,0.6314759) (1.8685541,0.6529794) (1.8956365,0.6747984) (1.9224641,0.6969299) (1.9490333,0.7193709) (1.9753406,0.7421184) (2.0013824,0.7651693) (2.0271553,0.7885206) (2.0526557,0.8121690) (2.0778803,0.8361115) (2.1028257,0.8603448) (2.1274884,0.8848657) (2.1518652,0.9096708) (2.1759528,0.9347568) (2.1997480,0.9601204) (2.2232476,0.9857581) (2.2464483,1.0116666) (2.2693472,1.0378422) (2.2919411,1.0642815) (2.3142270,1.0909810) (2.3362018,1.1179371) (2.3578627,1.1451460) (2.3792068,1.1726043) (2.4002310,1.2003081) (2.4209327,1.2282538) (2.4413091,1.2564376) (2.4613574,1.2848557) (2.4810748,1.3135043) (2.5004589,1.3423796) (2.5195069,1.3714776) (2.5382163,1.4007946) (2.5565846,1.4303264) (2.5746093,1.4600691) (2.5922880,1.4900189) (2.6096183,1.5201715) (2.6265979,1.5505231) (2.6432245,1.5810694) (2.6594959,1.6118064) (2.6754099,1.6427300) (2.6909643,1.6738360) (2.7061571,1.7051202) (2.7209862,1.7365784) (2.7354496,1.7682064) (2.7495454,1.8000000)};
\draw[gray,dashed,thin] (0.0000000,3.0000000)--(2.7495454,1.8000000);
\draw[gray,dashed,thin] (0.0000000,3.0000000)--(0.0000000,0.0000000);
\fill (0.0000000,3.0000000) circle (1.2pt);
\node[left,fill=none,inner sep=2pt] at (0.0000000,3.0000000) {$C$};
\node[below,fill=none,inner sep=2pt] at (0.8000000,2.1000000) {$R$};
\node[below,fill=none,inner sep=2pt] at (1.9000000,-0.3000000) {$x=R\sin\theta$};
\node[below,fill=none,inner sep=2pt] at (1.9000000,-0.7500000) {$y=R(1-\cos\theta)$};
\end{tikzpicture}
\figcap{光线从原点水平入射，沿圆心为 $C$ 的圆弧向上弯曲；圆弧切线确定传播方向。}\label{ch27:fig:s42-19-14}
\end{minipage}\end{center}
}

\newcommand{\FisherBayDiagram}{%
% Independent panel: original physical paths and local coordinates retained.
\begin{center}\begin{minipage}{\linewidth}\centering
\begin{tikzpicture}[>=Latex,font=\small,every node/.style={fill=none}]
\draw[thick] (-0.9576305,-3.5739256)--(0.0000000,0.0000000)--(0.9576305,-3.5739256);
\draw[gray,dashed,thin] (0.0000000,0.0000000)--(0.0000000,-3.6000000);
\draw[gray,dashed,thin] (0.0000000,0.0000000)--(-1.4974340,-2.4975371);
\draw[gray,dashed,thin] (-1.4974340,-2.4975371)--(-0.7246933,-2.7045923);
\fill (0.0000000,0.0000000) circle (1.2pt);
\node[above,fill=none,inner sep=2pt] at (0.0000000,0.0000000) {$D$};
\fill (-1.4974340,-2.4975371) circle (1.2pt);
\node[left,fill=none,inner sep=2pt] at (-1.4974340,-2.4975371) {$A$};
\fill (1.4974340,-2.4975371) circle (1.2pt);
\node[right,fill=none,inner sep=2pt] at (1.4974340,-2.4975371) {$B$};
\fill (-0.7246933,-2.7045923) circle (1.2pt);
\node[below right,fill=none,inner sep=2pt] at (-0.7246933,-2.7045923) {$C$};
\fill (-0.5994248,-2.2370840) circle (1.2pt);
\node[below,fill=none,inner sep=2pt] at (-0.5994248,-2.2370840) {};
\fill (0.5994248,-2.2370840) circle (1.2pt);
\node[below,fill=none,inner sep=2pt] at (0.5994248,-2.2370840) {};
\node[left,fill=none,inner sep=2pt] at (-1.6500000,-1.8000000) {$P$};
\draw[gray,dashed,thin] (-1.6000000,-1.8000000)--(-0.5994248,-2.2370840);
\node[above right,fill=none,inner sep=2pt] at (0.7794248,-1.9370840) {$Q$};
\draw[DeepBlue,thick,->] (-1.4974340,-2.4975371)--(-0.5994248,-2.2370840);
\draw[DeepBlue,thick,->] (-0.5994248,-2.2370840)--(0.5994248,-2.2370840);
\draw[DeepBlue,thick,->] (0.5994248,-2.2370840)--(1.4974340,-2.4975371);
\node[left,fill=none,inner sep=2pt] at (-1.6500000,-2.8500000) {$h$};
\draw[gray,dashed,thin] (-1.6000000,-2.8500000)--(-1.1110637,-2.6010647);
\node[left,fill=none,inner sep=2pt] at (-1.0000000,-1.2487685) {$l$};
\draw[thick] plot coordinates {(-0.1682324,-0.6278518) (-0.1649427,-0.6287240) (-0.1616484,-0.6295791) (-0.1583498,-0.6304168) (-0.1550467,-0.6312373) (-0.1517395,-0.6320404) (-0.1484281,-0.6328263) (-0.1451126,-0.6335948) (-0.1417931,-0.6343459) (-0.1384698,-0.6350796) (-0.1351426,-0.6357959) (-0.1318117,-0.6364948) (-0.1284773,-0.6371763) (-0.1251393,-0.6378402) (-0.1217979,-0.6384867) (-0.1184531,-0.6391157) (-0.1151051,-0.6397271) (-0.1117539,-0.6403211) (-0.1083997,-0.6408974) (-0.1050425,-0.6414562) (-0.1016824,-0.6419974) (-0.0983195,-0.6425210) (-0.0949540,-0.6430270) (-0.0915858,-0.6435154) (-0.0882151,-0.6439861) (-0.0848420,-0.6444392) (-0.0814666,-0.6448746) (-0.0780889,-0.6452923) (-0.0747091,-0.6456923) (-0.0713273,-0.6460746) (-0.0679435,-0.6464392) (-0.0645578,-0.6467861) (-0.0611704,-0.6471153) (-0.0577813,-0.6474267) (-0.0543906,-0.6477204) (-0.0509984,-0.6479963) (-0.0476048,-0.6482544) (-0.0442099,-0.6484948) (-0.0408138,-0.6487174) (-0.0374166,-0.6489222) (-0.0340184,-0.6491092) (-0.0306192,-0.6492784) (-0.0272192,-0.6494298) (-0.0238184,-0.6495635) (-0.0204170,-0.6496793) (-0.0170150,-0.6497773) (-0.0136126,-0.6498574) (-0.0102098,-0.6499198) (-0.0068067,-0.6499644) (-0.0034034,-0.6499911) (0.0000000,-0.6500000) (0.0034034,-0.6499911) (0.0068067,-0.6499644) (0.0102098,-0.6499198) (0.0136126,-0.6498574) (0.0170150,-0.6497773) (0.0204170,-0.6496793) (0.0238184,-0.6495635) (0.0272192,-0.6494298) (0.0306192,-0.6492784) (0.0340184,-0.6491092) (0.0374166,-0.6489222) (0.0408138,-0.6487174) (0.0442099,-0.6484948) (0.0476048,-0.6482544) (0.0509984,-0.6479963) (0.0543906,-0.6477204) (0.0577813,-0.6474267) (0.0611704,-0.6471153) (0.0645578,-0.6467861) (0.0679435,-0.6464392) (0.0713273,-0.6460746) (0.0747091,-0.6456923) (0.0780889,-0.6452923) (0.0814666,-0.6448746) (0.0848420,-0.6444392) (0.0882151,-0.6439861) (0.0915858,-0.6435154) (0.0949540,-0.6430270) (0.0983195,-0.6425210) (0.1016824,-0.6419974) (0.1050425,-0.6414562) (0.1083997,-0.6408974) (0.1117539,-0.6403211) (0.1151051,-0.6397271) (0.1184531,-0.6391157) (0.1217979,-0.6384867) (0.1251393,-0.6378402) (0.1284773,-0.6371763) (0.1318117,-0.6364948) (0.1351426,-0.6357959) (0.1384698,-0.6350796) (0.1417931,-0.6343459) (0.1451126,-0.6335948) (0.1484281,-0.6328263) (0.1517395,-0.6320404) (0.1550467,-0.6312373) (0.1583498,-0.6304168) (0.1616484,-0.6295791) (0.1649427,-0.6287240) (0.1682324,-0.6278518)};
\draw[gray,dashed,thin] (0.7000000,-0.8000000)--(0.0500000,-0.5500000);
\node[right,fill=none,inner sep=2pt] at (0.7500000,-0.8500000) {$\alpha$};
\node[below,fill=none,inner sep=2pt] at (0.0000000,-3.8000000) {湖湾内划船，岸外步行};
\end{tikzpicture}
\figcap{湖湾顶角为完整的 $\alpha$，$A,B$ 在陆地；$P,Q$ 为步行与划船的转换点。}\label{ch27:fig:s42-19-15}
\end{minipage}\end{center}
}

\newcommand{\ParallelPlateDiagram}{%
% Independent panel: original physical paths and local coordinates retained.
\begin{center}\begin{minipage}{\linewidth}\centering
\begin{tikzpicture}[>=Latex,font=\small,every node/.style={fill=none}]
\draw[thick] (-0.5000000,-0.6000000)--(-0.5000000,2.0000000);
\draw[thick] (1.0000000,-0.6000000)--(1.0000000,2.0000000);
\draw[gray,thin,->] (-2.7000000,0.0000000)--(3.0000000,0.0000000);
\fill (-2.5000000,0.0000000) circle (1.2pt);
\node[below,fill=none,inner sep=2pt] at (-2.5000000,0.0000000) {$O$};
\fill (-1.9809036,0.0000000) circle (1.2pt);
\node[below,fill=none,inner sep=2pt] at (-1.9809036,0.0000000) {$I$};
\draw[DeepBlue,thick,->] (-2.5000000,0.0000000)--(-0.5000000,0.5320431);
\draw[DeepBlue,thick,->] (-0.5000000,0.5320431)--(1.0000000,0.7929846);
\draw[DeepBlue,thick,->] (1.0000000,0.7929846)--(2.8000000,1.2718234);
\draw[gray,dashed,thin] (-1.9809036,0.0000000)--(1.0000000,0.7929846);
\node[above,fill=none,inner sep=2pt] at (2.8000000,1.2718234) {眼};
\node[above,fill=none,inner sep=2pt] at (0.2500000,1.6000000) {$n$};
\draw[gray,<->] (-2.5000000,-0.6500000)--(-0.5000000,-0.6500000);
\node[below,fill=none,inner sep=2pt] at (-1.5000000,-0.6500000) {$S$};
\draw[gray,<->] (-0.5000000,-0.6500000)--(1.0000000,-0.6500000);
\node[below,fill=none,inner sep=2pt] at (0.2500000,-0.6500000) {$D$};
\draw[gray,<->] (1.0000000,-0.6500000)--(2.8000000,-0.6500000);
\node[below,fill=none,inner sep=2pt] at (1.9000000,-0.6500000) {$L$};
\end{tikzpicture}
\figcap{玻璃板的入出射方向平行，视像 $I$ 在末段反向延长线上；$S,D,L$ 是三段轴向长度。}\label{ch27:fig:s42-19-16}
\end{minipage}\end{center}
}

\newcommand{\RainbowDiagram}{%
\begin{center}\begin{minipage}{\linewidth}\centering
\begin{tikzpicture}[>=Latex,font=\small,every node/.style={fill=none}]
\begin{scope}[shift={(-3.0000000,0.0000000)}]
\draw[thick] (0.0000000,0.0000000) circle (2);
\draw[DeepBlue,thick,->] (-2.1183502,1.7213259)--(-1.0183502,1.7213259);
\draw[gray,dashed,thin] (0.0000000,0.0000000)--(-1.0183502,1.7213259);
\fill (-1.0183502,1.7213259) circle (1.2pt);
\node[below,fill=none,inner sep=2pt] at (-1.0183502,1.7213259) {};
\node[above right,fill=none,inner sep=2pt] at (-0.4183502,2.3713259) {$A$};
\draw[gray,dashed,thin] (-1.0183502,1.7213259)--(-0.4683502,2.3213259);
\draw[gray,dashed,thin] (-1.0183502,1.7213259)--(-1.3493140,2.2807569);
\draw[thick] plot coordinates {(-1.1812862,1.9967381) (-1.1841322,1.9950344) (-1.1869604,1.9933012) (-1.1897705,1.9915389) (-1.1925622,1.9897475) (-1.1953352,1.9879273) (-1.1980892,1.9860784) (-1.2008238,1.9842011) (-1.2035388,1.9822956) (-1.2062340,1.9803620) (-1.2089089,1.9784005) (-1.2115634,1.9764115) (-1.2141971,1.9743950) (-1.2168098,1.9723514) (-1.2194011,1.9702808) (-1.2219709,1.9681834) (-1.2245187,1.9660595) (-1.2270444,1.9639093) (-1.2295477,1.9617330) (-1.2320283,1.9595310) (-1.2344860,1.9573033) (-1.2369204,1.9550502) (-1.2393313,1.9527721) (-1.2417185,1.9504691) (-1.2440817,1.9481414) (-1.2464206,1.9457894) (-1.2487351,1.9434133) (-1.2510247,1.9410133) (-1.2532894,1.9385897) (-1.2555288,1.9361428) (-1.2577428,1.9336728) (-1.2599310,1.9311799) (-1.2620933,1.9286645) (-1.2642293,1.9261269) (-1.2663390,1.9235672) (-1.2684220,1.9209858) (-1.2704782,1.9183830) (-1.2725072,1.9157589) (-1.2745090,1.9131140) (-1.2764832,1.9104485) (-1.2784297,1.9077627) (-1.2803482,1.9050568) (-1.2822386,1.9023312) (-1.2841006,1.8995861) (-1.2859341,1.8968219) (-1.2877388,1.8940388) (-1.2895146,1.8912372) (-1.2912613,1.8884173) (-1.2929786,1.8855794) (-1.2946664,1.8827240) (-1.2963245,1.8798511) (-1.2979528,1.8769613) (-1.2995510,1.8740547) (-1.3011190,1.8711317) (-1.3026567,1.8681926) (-1.3041637,1.8652377) (-1.3056401,1.8622674) (-1.3070856,1.8592819) (-1.3085001,1.8562816) (-1.3098834,1.8532668) (-1.3112353,1.8502378) (-1.3125558,1.8471950) (-1.3138447,1.8441386) (-1.3151019,1.8410691) (-1.3163271,1.8379866) (-1.3175204,1.8348917) (-1.3186815,1.8317845) (-1.3198103,1.8286655) (-1.3209067,1.8255350) (-1.3219707,1.8223932) (-1.3230020,1.8192406) (-1.3240005,1.8160775) (-1.3249663,1.8129041) (-1.3258990,1.8097210) (-1.3267988,1.8065283) (-1.3276654,1.8033265) (-1.3284987,1.8001159) (-1.3292988,1.7968968) (-1.3300654,1.7936696) (-1.3307985,1.7904346) (-1.3314981,1.7871922) (-1.3321640,1.7839428) (-1.3327962,1.7806866) (-1.3333946,1.7774240) (-1.3339592,1.7741554) (-1.3344898,1.7708811) (-1.3349865,1.7676015) (-1.3354492,1.7643169) (-1.3358777,1.7610277) (-1.3362722,1.7577342) (-1.3366325,1.7544368) (-1.3369586,1.7511359) (-1.3372505,1.7478317) (-1.3375081,1.7445247) (-1.3377315,1.7412153) (-1.3379205,1.7379036) (-1.3380751,1.7345902) (-1.3381954,1.7312754) (-1.3382814,1.7279595) (-1.3383330,1.7246429) (-1.3383502,1.7213259)};
\draw[thick] plot coordinates {(-0.7994049,1.3512409) (-0.7968135,1.3527862) (-0.7942330,1.3543498) (-0.7916636,1.3559314) (-0.7891053,1.3575309) (-0.7865583,1.3591484) (-0.7840227,1.3607838) (-0.7814987,1.3624368) (-0.7789863,1.3641076) (-0.7764857,1.3657959) (-0.7739970,1.3675018) (-0.7715204,1.3692250) (-0.7690559,1.3709656) (-0.7666036,1.3727235) (-0.7641638,1.3744985) (-0.7617365,1.3762905) (-0.7593218,1.3780996) (-0.7569199,1.3799256) (-0.7545308,1.3817684) (-0.7521547,1.3836279) (-0.7497918,1.3855040) (-0.7474420,1.3873966) (-0.7451056,1.3893057) (-0.7427826,1.3912312) (-0.7404733,1.3931729) (-0.7381776,1.3951307) (-0.7358956,1.3971046) (-0.7336276,1.3990945) (-0.7313736,1.4011003) (-0.7291338,1.4031218) (-0.7269082,1.4051589) (-0.7246969,1.4072117) (-0.7225001,1.4092799) (-0.7203178,1.4113634) (-0.7181503,1.4134623) (-0.7159975,1.4155763) (-0.7138596,1.4177053) (-0.7117367,1.4198493) (-0.7096288,1.4220081) (-0.7075362,1.4241817) (-0.7054589,1.4263699) (-0.7033970,1.4285726) (-0.7013506,1.4307897) (-0.6993198,1.4330212) (-0.6973047,1.4352668) (-0.6953054,1.4375265) (-0.6933220,1.4398002) (-0.6913547,1.4420878) (-0.6894034,1.4443891) (-0.6874683,1.4467040) (-0.6855495,1.4490324) (-0.6836471,1.4513743) (-0.6817612,1.4537295) (-0.6798918,1.4560978) (-0.6780392,1.4584792) (-0.6762032,1.4608735) (-0.6743841,1.4632806) (-0.6725820,1.4657005) (-0.6707968,1.4681329) (-0.6690288,1.4705778) (-0.6672780,1.4730351) (-0.6655445,1.4755046) (-0.6638283,1.4779861) (-0.6621296,1.4804797) (-0.6604484,1.4829851) (-0.6587849,1.4855023) (-0.6571390,1.4880310) (-0.6555110,1.4905713) (-0.6539008,1.4931229) (-0.6523085,1.4956857) (-0.6507343,1.4982597) (-0.6491781,1.5008446) (-0.6476402,1.5034404) (-0.6461205,1.5060469) (-0.6446191,1.5086640) (-0.6431361,1.5112916) (-0.6416716,1.5139296) (-0.6402256,1.5165777) (-0.6387983,1.5192359) (-0.6373896,1.5219041) (-0.6359997,1.5245821) (-0.6346286,1.5272698) (-0.6332765,1.5299670) (-0.6319432,1.5326736) (-0.6306290,1.5353896) (-0.6293339,1.5381147) (-0.6280580,1.5408488) (-0.6268012,1.5435918) (-0.6255638,1.5463435) (-0.6243456,1.5491039) (-0.6231469,1.5518727) (-0.6219676,1.5546499) (-0.6208079,1.5574353) (-0.6196677,1.5602288) (-0.6185471,1.5630302) (-0.6174463,1.5658394) (-0.6163651,1.5686562) (-0.6153038,1.5714806) (-0.6142623,1.5743123) (-0.6132407,1.5771513) (-0.6122390,1.5799974)};
\node[above,fill=none,inner sep=2pt] at (-1.6698526,2.0928693) {$i$};
\node[right,fill=none,inner sep=2pt] at (-0.3991862,1.2147334) {$r$};
\draw[DeepBlue,thick,->] (-1.0183502,1.7213259)--(1.8669753,0.7172191);
\draw[gray,dashed,thin] (0.0000000,0.0000000)--(1.8669753,0.7172191);
\draw[DeepBlue,thick,->] (1.8669753,0.7172191)--(0.3960251,-1.9603990);
\draw[gray,dashed,thin] (0.0000000,0.0000000)--(0.3960251,-1.9603990);
\draw[DeepBlue,thick,->] (0.3960251,-1.9603990)--(-0.6438926,-2.8977203);
\node[below,fill=none,inner sep=2pt] at (0.0000000,-3.3000000) {主虹：一次内反射};
\end{scope}
\begin{scope}[shift={(3.0000000,0.0000000)}]
\draw[thick] (0.0000000,0.0000000) circle (2);
\draw[DeepBlue,thick,->] (-1.7236096,1.9002924)--(-0.6236096,1.9002924);
\draw[gray,dashed,thin] (0.0000000,0.0000000)--(-0.6236096,1.9002924);
\fill (-0.6236096,1.9002924) circle (1.2pt);
\node[below,fill=none,inner sep=2pt] at (-0.6236096,1.9002924) {};
\node[above right,fill=none,inner sep=2pt] at (-0.0236096,2.5502924) {$A$};
\draw[gray,dashed,thin] (-0.6236096,1.9002924)--(-0.0736096,2.5002924);
\draw[gray,dashed,thin] (-0.6236096,1.9002924)--(-0.8262827,2.5178874);
\draw[thick] plot coordinates {(-0.7233871,2.2043392) (-0.7271910,2.2030644) (-0.7309786,2.2017420) (-0.7347494,2.2003723) (-0.7385027,2.1989554) (-0.7422379,2.1974915) (-0.7459545,2.1959809) (-0.7496518,2.1944239) (-0.7533293,2.1928206) (-0.7569865,2.1911714) (-0.7606227,2.1894764) (-0.7642373,2.1877360) (-0.7678299,2.1859504) (-0.7713998,2.1841199) (-0.7749464,2.1822448) (-0.7784693,2.1803254) (-0.7819678,2.1783619) (-0.7854414,2.1763548) (-0.7888896,2.1743042) (-0.7923118,2.1722106) (-0.7957076,2.1700743) (-0.7990762,2.1678955) (-0.8024173,2.1656747) (-0.8057303,2.1634122) (-0.8090146,2.1611083) (-0.8122698,2.1587635) (-0.8154954,2.1563780) (-0.8186908,2.1539522) (-0.8218555,2.1514866) (-0.8249891,2.1489815) (-0.8280910,2.1464374) (-0.8311608,2.1438545) (-0.8341979,2.1412333) (-0.8372020,2.1385743) (-0.8401725,2.1358778) (-0.8431089,2.1331443) (-0.8460109,2.1303742) (-0.8488778,2.1275680) (-0.8517094,2.1247260) (-0.8545052,2.1218487) (-0.8572646,2.1189366) (-0.8599873,2.1159902) (-0.8626729,2.1130098) (-0.8653209,2.1099961) (-0.8679309,2.1069493) (-0.8705025,2.1038701) (-0.8730353,2.1007589) (-0.8755289,2.0976161) (-0.8779829,2.0944424) (-0.8803969,2.0912381) (-0.8827705,2.0880038) (-0.8851034,2.0847401) (-0.8873953,2.0814473) (-0.8896456,2.0781261) (-0.8918542,2.0747769) (-0.8940205,2.0714002) (-0.8961444,2.0679967) (-0.8982255,2.0645669) (-0.9002634,2.0611112) (-0.9022577,2.0576302) (-0.9042083,2.0541245) (-0.9061148,2.0505946) (-0.9079769,2.0470411) (-0.9097943,2.0434646) (-0.9115668,2.0398655) (-0.9132939,2.0362445) (-0.9149755,2.0326021) (-0.9166114,2.0289389) (-0.9182012,2.0252555) (-0.9197446,2.0215525) (-0.9212416,2.0178304) (-0.9226917,2.0140899) (-0.9240949,2.0103314) (-0.9254508,2.0065557) (-0.9267592,2.0027632) (-0.9280201,1.9989547) (-0.9292330,1.9951306) (-0.9303980,1.9912916) (-0.9315147,1.9874384) (-0.9325830,1.9835714) (-0.9336028,1.9796913) (-0.9345738,1.9757988) (-0.9354960,1.9718944) (-0.9363691,1.9679787) (-0.9371931,1.9640524) (-0.9379678,1.9601161) (-0.9386931,1.9561704) (-0.9393689,1.9522159) (-0.9399950,1.9482532) (-0.9405714,1.9442830) (-0.9410980,1.9403059) (-0.9415747,1.9363225) (-0.9420014,1.9323334) (-0.9423781,1.9283393) (-0.9427046,1.9243408) (-0.9429811,1.9203385) (-0.9432073,1.9163331) (-0.9433833,1.9123251) (-0.9435090,1.9083152) (-0.9435844,1.9043041) (-0.9436096,1.9002924)};
\draw[thick] plot coordinates {(-0.4895335,1.4917295) (-0.4862970,1.4928059) (-0.4830691,1.4939078) (-0.4798501,1.4950354) (-0.4766401,1.4961885) (-0.4734394,1.4973669) (-0.4702481,1.4985708) (-0.4670664,1.4997999) (-0.4638946,1.5010542) (-0.4607329,1.5023336) (-0.4575814,1.5036381) (-0.4544403,1.5049675) (-0.4513099,1.5063218) (-0.4481904,1.5077008) (-0.4450818,1.5091046) (-0.4419846,1.5105330) (-0.4388987,1.5119859) (-0.4358244,1.5134633) (-0.4327620,1.5149650) (-0.4297116,1.5164909) (-0.4266734,1.5180410) (-0.4236475,1.5196151) (-0.4206343,1.5212132) (-0.4176338,1.5228351) (-0.4146463,1.5244808) (-0.4116719,1.5261501) (-0.4087109,1.5278430) (-0.4057633,1.5295593) (-0.4028295,1.5312989) (-0.3999096,1.5330618) (-0.3970038,1.5348477) (-0.3941122,1.5366566) (-0.3912350,1.5384885) (-0.3883725,1.5403430) (-0.3855248,1.5422203) (-0.3826920,1.5441200) (-0.3798744,1.5460422) (-0.3770722,1.5479867) (-0.3742854,1.5499533) (-0.3715144,1.5519419) (-0.3687592,1.5539525) (-0.3660200,1.5559849) (-0.3632971,1.5580389) (-0.3605905,1.5601145) (-0.3579005,1.5622115) (-0.3552272,1.5643297) (-0.3525708,1.5664691) (-0.3499314,1.5686295) (-0.3473092,1.5708108) (-0.3447045,1.5730128) (-0.3421173,1.5752353) (-0.3395477,1.5774784) (-0.3369961,1.5797417) (-0.3344625,1.5820252) (-0.3319471,1.5843288) (-0.3294500,1.5866522) (-0.3269715,1.5889953) (-0.3245116,1.5913580) (-0.3220706,1.5937402) (-0.3196485,1.5961417) (-0.3172455,1.5985623) (-0.3148618,1.6010019) (-0.3124975,1.6034603) (-0.3101528,1.6059374) (-0.3078279,1.6084330) (-0.3055228,1.6109469) (-0.3032377,1.6134791) (-0.3009727,1.6160293) (-0.2987281,1.6185974) (-0.2965039,1.6211833) (-0.2943003,1.6237867) (-0.2921174,1.6264074) (-0.2899554,1.6290455) (-0.2878143,1.6317006) (-0.2856944,1.6343725) (-0.2835958,1.6370613) (-0.2815185,1.6397665) (-0.2794628,1.6424882) (-0.2774287,1.6452261) (-0.2754164,1.6479800) (-0.2734260,1.6507499) (-0.2714576,1.6535354) (-0.2695114,1.6563364) (-0.2675875,1.6591528) (-0.2656860,1.6619844) (-0.2638070,1.6648309) (-0.2619506,1.6676923) (-0.2601170,1.6705683) (-0.2583063,1.6734588) (-0.2565185,1.6763635) (-0.2547538,1.6792823) (-0.2530124,1.6822150) (-0.2512942,1.6851615) (-0.2495995,1.6881215) (-0.2479284,1.6910948) (-0.2462808,1.6940813) (-0.2446570,1.6970808) (-0.2430570,1.7000930) (-0.2414810,1.7031179) (-0.2399291,1.7061551) (-0.2384012,1.7092046)};
\node[above,fill=none,inner sep=2pt] at (-1.2310181,2.3402411) {$i$};
\node[right,fill=none,inner sep=2pt] at (-0.0999029,1.2955351) {$r$};
\draw[DeepBlue,thick,->] (-0.6236096,1.9002924)--(1.8903165,0.6532255);
\draw[gray,dashed,thin] (0.0000000,0.0000000)--(1.8903165,0.6532255);
\draw[DeepBlue,thick,->] (1.8903165,0.6532255)--(0.6826820,-1.8798791);
\draw[gray,dashed,thin] (0.0000000,0.0000000)--(0.6826820,-1.8798791);
\draw[DeepBlue,thick,->] (0.6826820,-1.8798791)--(-1.8689827,-0.7119717);
\draw[gray,dashed,thin] (0.0000000,0.0000000)--(-1.8689827,-0.7119717);
\draw[DeepBlue,thick,->] (-1.8689827,-0.7119717)--(-2.7504472,0.3756957);
\node[below,fill=none,inner sep=2pt] at (0.0000000,-3.3000000) {次虹：两次内反射};
\end{scope}
\end{tikzpicture}
\figcap{主虹与次虹的球内路径；虚线为法线，$i,r$ 在入射点 $A$ 测量。离球后的箭头指向接收这束光的观察者。\label{ch27:fig:rainbow}}
\end{minipage}\end{center}
}


光学是一门古老而至今仍十分活跃的学科。现存最早的人工透镜可追溯到公元前七世纪；从以牛顿为代表的近代物理学开始，光学理论和光学仪器的制造一同进入快速发展时期。今天，光学不仅是基础物理的重要组成部分，也深刻影响着通信、精密测量、医学、天文、材料和信息技术。

贯穿光学发展史的核心问题是：光究竟是什么？牛顿（Isaac Newton）曾把光理解为高速运动的粒子流；同一时期，惠更斯（Christiaan Huygens）主张光是一种波。十九世纪初，杨（Thomas Young）的干涉实验和菲涅尔（Augustin-Jean Fresnel）的理论使波动说重新兴起；麦克斯韦（James Clerk Maxwell）随后指出光是横向电磁波，不同颜色对应不同频率或波长，而真空中的传播速度相同。二十世纪初，爱因斯坦（Albert Einstein）的光量子理论揭示了能量交换的离散性。波粒二象性并不表示单个光子会在“传播时是经典波、探测时是经典粒子”之间切换：光应由量子态或量子化电磁场统一描述，概率振幅可以叠加并产生干涉，探测器则记录一次次离散的能量交换事件。经典波和经典粒子只是不同实验条件下的近似图像；量子电动力学中，光子是电磁场的量子化激发。


光学仪器的发展构成了另一条同样重要的主线。平面镜的平整度、透镜和反射镜的面形精度、材料纯度以及装调误差都会改变仪器的性能。现实中不存在绝对完美的光学系统，设计者必须在功能、精度、口径、重量和成本之间权衡。历史人物与代表性器件提示我们：光的理论与实验仪器始终相互推动。

\historyfigure{optics-history-built-in-v1.png}{光学史人物与代表性实验器件：粒子说、波动说、棱镜色散、望远镜与双缝实验。}

当光学元件和传播尺度远大于光的波长时，波动光学可以取极短波长极限，得到更简单的\textbf{几何光学}。几何光学用“光线”表示能量传播方向，但光线不是有实体的细线，也不是经典粒子的真实轨迹；在各向同性介质中，它是波面法线在空间中的积分曲线。

因此几何光学有明确的适用边界：器件尺寸、曲率半径和传播中被分辨的细节都应远大于波长。接近小孔、尖锐边缘或亚波长结构时，单靠光线会漏掉衍射；两路光重新相遇时，单靠路径也不能判断明暗。后两讲将处理这些波动效应，本讲则先在短波长极限内建立可用于追踪方向的规律。

在任一时刻，把光场中相位相同的点连接起来，就得到等相位面（波面）。在各向同性透明介质中，光能沿波面的法线方向传播。点光源在均匀各向同性介质中的波面是一组同心球面，所以光线是从光源出发的直线。这给出几何光学的第一条基本规律：\textbf{光在均匀介质中沿直线传播}。

若介质不均匀，波速随位置改变，波面各部分前进速度不同，光线便会弯曲。图~\ref{ch27:fig:s42-19-1} 中，大气密度通常向地面方向增大，折射率也相应增大；接近地平线的太阳光逐层向法线方向偏折。观察者沿最后一段光线反向延长，会把太阳判断在真实位置之上，因此太阳已经略低于几何地平线时仍可能被看见。

\AtmosphereWavefrontPanel



\subsection{从波长、折射率到光程}
光的频率记为 $\nu$，角频率为 $\omega=2\pi\nu$。真空波长 $\lambda_0=c/\nu$ 与介质内波长不是同一个量。在静止界面上，入射波、反射波和折射波必须以相同频率振动，否则界面上的电场边界条件不能在每个时刻同时满足。若透明介质的相折射率为 $n$，则相速度 $v_p=c/n$，介质内波长为
\[
 \lambda=\frac{v_p}{\nu}=\frac{\lambda_0}{n}.
\]
所以“进入玻璃后波长变短”不表示光的频率降低。色散指 $n$ 随频率而变；比较不同颜色时，要分别使用各自的折射率。

一段长度为 $s$ 的均匀介质，使单色波积累的传播相位为
\[
 ks=\frac{2\pi}{\lambda}s=\frac{2\pi}{\lambda_0}ns.
\]
因此把 $ns$ 称为这段路的光程：它把“走过多长”和“介质使相位变化多快”合在一起。依次通过若干介质时，光程相加为 $\sum_jn_js_j$；折射率连续变化时才把求和改为积分。光程的单位仍是米，折射率没有量纲。它不是把折射后的折线路径拉直所得到的几何长度。

本章的普通玻璃、水和空气先按各向同性介质处理。在这种介质中，波面法线、波矢与能量传播方向一致；晶体中的非寻常光则不一定如此，第30章再区分波法线与能流方向。脉冲通过明显色散的介质时，能量包络的传播时间涉及群速度，也不能直接把相折射率代入 $t=ns/c$。


\section{反射定律与镜面光路}



光照到平整界面时，一部分能量返回原介质，称为反射。在入射点作界面的法线；入射光线与法线的夹角称为入射角 $i$，反射光线与法线的夹角称为反射角 $r$。入射光线、法线和反射光线共面，并且
\begin{equation}
  i=r.
\end{equation}
角度必须相对\textbf{法线}测量。对曲面，应先在入射点作切平面，再以切平面的垂线作为法线。

两块平面镜夹角为 $\alpha$ 时，连续两次反射后，出射方向相对入射方向的总转角为
\begin{equation}
  \Delta\varphi=2\alpha,
\end{equation}
其大小与第一束光的入射角无关。$\alpha=0$ 对应平行镜，可用于平移光路；$\alpha=\pi/2$ 时出射方向与入射方向相反，构成直角反射器。月球激光测距所用的角锥反射器正是这一思想的三维推广。

这个结论可以直接用方向角核验。设一条光线相对固定轴的有向角为 $\varphi$，镜面方向角为 $\beta$。镜面反射把光线关于镜面对称，故一次反射后
\[
 \varphi'=2\beta-\varphi.
\]
若两镜方向角依次为 $\beta_1,\beta_2$，两次反射后
\[
 \varphi''=2\beta_2-(2\beta_1-\varphi)
 =\varphi+2(\beta_2-\beta_1).
\]
因此方向转角确为 $2\alpha$，而不是 $\pi-2\alpha$。后一表达式有时表示两条无向直线间取补角后的夹角，不能当作光的有向传播方向变化。图~\ref{ch27:fig:s42-19-3} 中入射箭头、两次反射点和出射箭头正是按这一关系绘制的。

\historyfigure[0.96\linewidth]{reflection-applications-built-in-v1.png}{反射定律实验；潜望镜与月面激光反射器。\label{ch27:fig:s42-19-2}}
图~\ref{ch27:fig:s42-19-2} 将反射定律放到实验和装置中：潜望镜利用多次反射改变光路，月面反射器用于激光往返测距；逐个反射面都要以该面的法线量角。

\subsection{先定法线，再追踪反射方向}
入射角从法线量起，不从镜面量起。例如光线与镜面成 $20^\circ$，入射角就是 $70^\circ$。角度数值正确还不够：反射光必须留在入射侧，并沿离开反射点的方向传播。把一条直线的两个相反方向混为一谈，往往会得到“线平行但箭头错误”的光路。

设入射传播方向的单位矢量为 $\vect s$，镜面单位法向为 $\vect n$。分解成沿面与垂直面两部分：
\[
 \vect s=\bigl[\vect s-(\vect s\cdot\vect n)\vect n\bigr]
       +(\vect s\cdot\vect n)\vect n.
\]
反射后沿面分量不变，法向分量反号，因而
\begin{equation}\label{ch27:eq:depth-reflection-vector}
 \vect s_{\rm r}=\vect s-2(\vect s\cdot\vect n)\vect n.
\end{equation}
换用 $-\vect n$ 不改变结果。这一式既可检查平面镜，也可检查曲面镜，只需把 $\vect n$ 换成入射点处的局部法线。球面法线必须经过球心；它一般不是图上的竖直线或水平线。

连续两次反射时，应把第一次反射后的传播方向作为第二次的入射方向。用镜像法展开折线只是帮助计算：镜后虚像和虚线延长线并不表示光实际穿进镜后。求光斑运动时还要先确定最后一条反射线与屏的交点，再对交点位置求时间导数，不能把镜子的角速度直接当作光斑的线速度。


\ReflectionPanel


\SphericalMirrorRotatingDiagram
\Needspace{7\baselineskip}\begin{origexamplebox}{【球面镜】}\label{ch27:ex:s42-19-1}
有一个半径为 $R$ 的球形反射面，圆心位于 $C$，$CV$ 为其对称轴。一束与对称轴平行、距轴为 $d$ 的入射光经球面反射后与对称轴交于 $M$。求 $M$ 与球面顶点 $V$ 之间的距离 $s$。
\end{origexamplebox}\par\nopagebreak[3]
\begin{xesolution}
取子午截面。设入射点为 $P$，令 $u=\sqrt{R^2-d^2}$。从顶点 $V$ 向球心方向量，$P$ 的轴向坐标为 $R-u$。半径 $CP$ 是法线，记它与轴的夹角为 $\alpha$，则 $\sin\alpha=d/R$、$\cos\alpha=u/R$。入射光平行于轴，故反射光与轴的夹角为 $2\alpha$。由图~\ref{ch27:fig:s42-19-11} 左图中的直角三角形，
\[
 s=(R-u)+\frac{d}{\tan 2\alpha}.
\]
又 $\tan2\alpha=2du/(u^2-d^2)$，整理得
\[
 \boxed{s=R-\frac{R^2}{2\sqrt{R^2-d^2}}}.
\]
上式的 $s$ 是从 $V$ 向 $C$ 为正的有向坐标；题目所问几何距离为 $|s|$，参数域为 $|d|<R$。当 $|d|<\sqrt3R/2$ 时 $s>0$，超过此高度时轴交点移到顶点另一侧，$s<0$。当 $d\to0$ 时 $s\to R/2$；不同高度给出不同交点，说明球面镜存在球差。$d=0$ 的轴上光与轴重合，焦点位置须由近轴极限定义。
\end{xesolution}

\Needspace{7\baselineskip}\begin{origexamplebox}{【转动的镜子】}\label{ch27:ex:s42-19-2}
一竖直墙上 $A$ 点发出水平光束，平面镜 $OB$ 绕过 $O$ 的水平轴以角速度 $\omega$ 转动。光束经镜面反射后回到墙上形成光点 $D$。已知此刻镜面与墙的夹角为 $\theta$，$AO=d$，忽略光的传播时间，求光点在墙上的移动速度。
\end{origexamplebox}\par\nopagebreak[3]
\begin{xesolution}
%<optics-review:mirror-domain>
先限定能够在墙上看到反射光点的姿态。取图示镜面位于墙左侧、$0<\theta<\pi/2$，入射光从墙向左传播；以下把 $\omega=\dot\theta$ 作为带符号角速度。若题给的是角速度大小，须先由转向确定正负。反射方向相对向右的水平线转过 $2\theta$，其单位方向矢量为 $(\cos2\theta,\sin2\theta)$。因此光要从左侧镜面返回右侧墙面，必须有
\[
 \cos2\theta>0,\qquad 0<\theta<\frac\pi4.
\]
在 $\theta=\pi/4$ 时反射光平行于墙；超过这个角度，反射光向左离开，直线延长线与墙的交点不能再解释成实际光点。这里还默认入射点位于镜面的有效长度内。

%</optics-review:mirror-domain>

在垂直转轴的截面内令墙为 $x=0$，$O=(0,0)$，$A=(0,d)$。镜面直线可写为 $x=-y\tan\theta$，水平入射线 $y=d$ 与镜面交于 $C=(-d\tan\theta,d)$。镜面转过 $\theta$ 时反射方向转过 $2\theta$，故从 $C$ 到墙的竖直位移为 $d\tan\theta\tan2\theta$。于是
\[
 y_D=d+d\tan\theta\tan2\theta.
\]

%<optics-review:mirror-simplify>
这里 $d\tan\theta$ 是反射点到墙的水平距离，乘上反射方向的斜率 $\tan2\theta$ 才是竖直位移。利用
\[
 1+\tan\theta\tan2\theta
 =\frac{\cos\theta\cos2\theta+\sin\theta\sin2\theta}
 {\cos\theta\cos2\theta}
 =\frac1{\cos2\theta},
\]
也可把位置写成 $y_D=d\sec2\theta$，从而直接得到
\[
 v_D=2d\omega\sec2\theta\tan2\theta.
\]
这是下式的等价形式。在上述允许区间内，$\omega>0$ 时光点上移，$\omega<0$ 时下移；其量纲为长度乘角速度，即速度。

%</optics-review:mirror-simplify>
利用 $\dot\theta=\omega$ 求导，得到沿墙的瞬时速度
\[
 \boxed{v_D=d\omega\!\left(\sec^2\theta\tan2\theta+2\tan\theta\sec^22\theta\right)}.
\]
式子的正负表示移动方向；题目若只问速率，应取绝对值。$\theta$ 接近使反射光平行于墙的位置时，$\tan2\theta$ 发散，光点速度可很大，这并不表示物质或信息以该速度传播。
\end{xesolution}

\section{折射定律与多界面光路}



光从一种透明介质进入另一种透明介质时，传播方向一般会改变。折射光线仍位于入射平面内。若界面两侧折射率分别为 $n_1,n_2$，入射角和折射角分别为 $\theta_1,\theta_2$，则斯涅尔定律（Snell's law）为
\begin{equation}
  n_1\sin\theta_1=n_2\sin\theta_2.
\end{equation}
折射率定义为 $n=c/v$，其中 $c$ 为真空光速，$v$ 为介质中的相速度。在本讲讨论的普通透明介质中 $n>1$。同一介质的折射率通常随波长变化，所以复色光会发生色散。近轴条件 $|\theta|\ll 1$（角度用弧度表示）下，$\sin\theta\simeq\theta$，斯涅尔定律近似为
\begin{equation}
  n_1\theta_1\simeq n_2\theta_2.
\end{equation}



\subsection{多次折射中什么可以消去}
折射角总是相对于当前界面的法线定义。若各界面互相平行，法线方向相同，一段介质内的出射角就等于下一界面的入射角，于是
\[
 n_1\sin i_1=n_2\sin i_2=\cdots=n_k\sin i_k.
\]
中间折射率可以从“最终方向”的关系中消去，但中间层的厚度、侧向位移和光程不能消去。若两界面不平行，例如棱镜的两个侧面，必须先用几何关系换算第二次入射角，再用折射定律。

以厚度为 $t$ 的平行玻璃板为例，板外同为空气，空气侧角为 $i$，板内角为 $r$。第二次折射满足 $n\sin r=\sin i'$，因此 $i'=i$：出射光与入射光同向平行。板内光线的几何长度为 $t/\cos r$，相对原直线的垂直侧移量则为
\[
 b=\frac{t\sin(i-r)}{\cos r}.
\]
这里 $t\tan i-t\tan r$ 是在平行于板面的方向上比较得到的交点差，再乘 $\cos i$ 才得到两条平行光线的最短距离 $b$。当 $n\to1$ 或 $t\to0$ 时，$b\to0$；正入射时也没有侧移，但仍会产生光程变化。

处理球壳或圆柱时，先连“球心或圆心—入射点”得到法线。光线到中心的垂距、光线与半径的夹角和球半径之间满足直角三角形关系；先建立这一关系，才能把折射定律与“碰到内壁”“与内球相切”等条件连接起来。


\CylinderDiagram

\Needspace{7\baselineskip}\begin{origexamplebox}{【圆柱折射】}\label{ch27:ex:s42-19-3}
一个半径 $R=2\,\mathrm m$ 的透明圆柱体，其一侧是反射面。一束光从左侧空气入射，经折射和反射后，出射光与入射光平行；两平行光线的距离 $d=2\,\mathrm m$。求透明介质的折射率。
\end{origexamplebox}\par\nopagebreak[3]
\begin{xesolution}
在垂直圆柱轴的截面内分析。图~\ref{ch27:fig:s42-19-12} 中两外部光线关于过圆心 $C$ 的水平轴对称，入射点 $A$ 的高度为 $d/2=R/2$。取 $C=(0,0)$，便有 $A=(-\sqrt3R/2,R/2)$、反射点 $P=(R,0)$，因此相对向内法线 $AC$ 的入射角 $i=30^\circ$。内部线段 $AP$ 相对水平轴向下的倾角为 $\beta$，其中
\[
 \tan\beta=\frac{R/2}{R+\sqrt3R/2}=2-\sqrt3=\tan15^\circ.
\]
所以 $\beta=15^\circ$，折射角为法线与内部光线的夹角 $r=i-\beta=15^\circ$。在 $P$ 处按水平法线反射后，光线到达对称点 $B$，再折射为水平向左的出射光，方向也与题设一致。斯涅尔定律给出
\[
 1\cdot\sin30^\circ=n\sin15^\circ,\qquad
 \boxed{n=\frac{\sin30^\circ}{\sin15^\circ}=\frac{\sqrt6+\sqrt2}{2}\approx1.93}.
\]
图~\ref{ch27:fig:s42-19-12} 同时画出了两次折射点和右侧反射点；若只凭“平行”而忽略圆柱半径法线的位置关系，角度会误判。
\end{xesolution}

% Relocated AbsorbingShellDiagram to its own solution.

\Needspace{7\baselineskip}\begin{origexamplebox}{【大小球】}\label{ch27:ex:s42-19-4}
用折射率为 $n$ 的透明物质制成内、外半径分别为 $a,b$ 的空心球，内表面涂有能完全吸收光的物质。平行光射向此球时，球吸收的入射光束横截面积多大？
\end{origexamplebox}\par\nopagebreak[3]
\begin{xesolution}
\AbsorbingShellDiagram

先采用几何光学模型，忽略外表面的菲涅尔反射损耗，求能到达吸收层的入射光束面积，而不是包含反射损耗的功率吸收截面。取图~\ref{ch27:fig:s42-19-13} 中恰与内球在 $T$ 相切的临界光线。设它进入外球面后的折射角为 $r$，则直线到球心的距离为 $b\sin r=a$。设这条光线在入射平面内的入射高度（入射线到球心的垂直距离）为 $\rho$，空气侧有 $\rho=b\sin i$。由斯涅尔定律 $\sin i=n\sin r$，所以
\[
 \rho=b\sin i=nb\sin r=na.
\]
入射高度小于 $\rho$ 的光线会碰到吸收层，故
\[
 \boxed{S_{\rm abs}=\pi\min\{b^2,n^2a^2\}}.
\]
若 $na\ge b$，所有射入外球几何截面的光都能到达内层，面积饱和为 $\pi b^2$；这一截断条件不能省略。
\end{xesolution}

\Needspace{7\baselineskip}\begin{origexamplebox}{【平板玻璃砖】}\label{ch27:ex:s42-19-9}
观察者距折射率为 $n$、厚度为 $D$ 的平行玻璃板近表面为 $L$。发光点位于玻璃板另一侧，距远表面为 $S$。求观察者看到的发散光线反向延长后的视距离。
\end{origexamplebox}\par\nopagebreak[3]
\begin{xesolution}
\ParallelPlateDiagram

%<optics-review:plate-reference>
参看图~\ref{ch27:fig:s42-19-16}：$O$ 是光源，$I$ 在玻璃出射光的反向延长线上，眼睛在玻璃另一侧。$S,D,L$ 分别是光源到玻璃、玻璃内部、玻璃到眼睛的轴向长度。应从眼睛收到的末段方向反推位置，不是把各段的相位光程 $n\ell$ 相加。

%</optics-review:plate-reference>

取一条近轴光线。在空气中角度为 $\theta$，玻璃中角度为 $\theta_g$。由近轴斯涅尔定律 $\theta_g\simeq\theta/n$。光线在三个区段产生的横向位移为
\[
 y\simeq L\theta+D\theta_g+S\theta
  =\theta\left(L+\frac Dn+S\right).
\]
观察者按空气中的末段方向反向延长，虚像点到观察者的距离就是 $y/\theta$：
\[
 \boxed{L_{\rm view}=L+\frac Dn+S}.
\]
所以玻璃的视厚度为 $D/n$，比真实厚度小。该结果依赖近轴条件；大角度光线的反向延长线不再严格交于同一点。
\end{xesolution}

\section{棱镜偏向、色散与彩虹}

对顶角为 $A$、置于空气中的三棱镜，记两面外侧入射/出射角为 $i_1,i_2$，棱镜内的折射角为 $r_1,r_2$。由几何关系
\begin{equation}
 r_1+r_2=A,\qquad \delta=i_1+i_2-A.
\end{equation}
在对称光路 $i_1=i_2$、$r_1=r_2=A/2$ 处偏向角取极小值。普通玻璃中蓝光折射率较大，因而偏折通常比红光更强。牛顿的棱镜实验正是由此说明白光由多种颜色组成。

棱镜计算必须把两个界面依次写全。空气折射率取 $1$ 时，第一面与第二面分别满足
\[
 \sin i_1=n\sin r_1,\qquad n\sin r_2=\sin i_2.
\]
给定 $A,n$ 和 $i_1$ 后，先由第一式求 $r_1$，再用 $r_2=A-r_1$ 求 $i_2$，最后才代入 $\delta=i_1+i_2-A$。图示通常取两内部角均为正；若光线越过法线，应使用有向角。令 $r_c=\arcsin(1/n)$，从指定两面贯穿的必要角度条件是 $|r_1|,|r_2|\le r_c$，同时要确认中间光段没有先碰到底面。若第二面上 $|r_2|>r_c$，便发生全反射。等号表示掠射极限，不能把它当作具有有限法向能流的出射光。

把入、出方向交换，光路可逆，因此交换 $r_1,r_2$ 不改变偏向角；但这种对称性本身并不能证明极小。令 $f(r)=\arcsin(n\sin r)$，则 $\delta=f(r_1)+f(A-r_1)-A$。函数 $f'(r)=n\cos r/\sqrt{1-n^2\sin^2r}$ 为偶函数，并随 $|r|$ 增大而增大。由于 $A>0$，$r_1<A/2$ 时有 $|r_1|<|A-r_1|$，故 $\delta'<0$；$r_1>A/2$ 时相反。于是当 $A<2r_c$、对称光路可贯穿棱镜时，唯一极小值在 $r_1=r_2=A/2$，即
\[
 \delta_{\min}=2\arcsin\!\left(n\sin\frac A2\right)-A.
\]
实验中缓慢转动棱镜时，光斑先向一个方向移动、到达最小偏向位置后再反向；这比仅凭一张静态光路图更容易辨认极小值。

\RefractionPrismPanel

图\ref{ch27:fig:deviation}给出固定折射率下的偏向角曲线。入射角增大时偏向角先减小后增大，最低点对应内部对称光路；它与色散图回答的是两个不同问题。
\subsection{最小偏向与色散测量的条件}
棱镜公式中的 $A$ 是两个工作面的夹角，不一定等于图上最醒目的角。设内部两次折射角分别为 $r_1,r_2$，有 $r_1+r_2=A$。在可逆、对称的最小偏向光路中，$r_1=r_2=A/2$、$i_1=i_2=i$，故
\[
 \delta_{\min}=2i-A,\qquad
 n=\frac{\sin[(A+\delta_{\min})/2]}{\sin(A/2)}
\]
（外界为空气）。实验上固定棱镜并只改变波长时，不同颜色一般不同时满足最小偏向条件；不能把同一个入射角问题误做成“每种颜色都重新调整到最小偏向”。

薄棱镜近似还要求 $A,i,r$ 都足够小。把正弦近似成角度后，$i_1\simeq nr_1$、$i_2\simeq nr_2$，于是
\[
 \delta=i_1+i_2-A\simeq(n-1)A,
 \qquad \delta_{\rm violet}-\delta_{\rm red}
 \simeq(n_{\rm violet}-n_{\rm red})A.
\]
这说明色散不仅来自“有棱镜”，更来自不同频率的折射率不同。若需要较大的偏向角或高精度结果，必须回到两次折射的精确关系，不能继续用线性近似。


\PrismDeviationDiagram



\Needspace{7\baselineskip}\begin{origexamplebox}{【最小偏向角】}\label{ch27:ex:s42-19-11}
一三棱镜顶角为 $A$、折射率为 $n$，光由一个面入射并从另一面射出。求偏向角的可能范围。
\end{origexamplebox}\par\nopagebreak[3]
\begin{xesolution}
取空气中 $n>1$、$0<A<\pi$，讨论从指定两侧面贯穿且未被底面截断的光路。为覆盖小顶角情形，内部角 $r_1,r_2$ 采用有向角，满足 $r_1+r_2=A$；负值表示光线越过对应法线。外侧角为 $i_j=\arcsin(n\sin r_j)$，偏向角 $\delta=i_1+i_2-A$。
令 $r_c=\arcsin(1/n)$，两面透射要求 $|r_j|\le r_c$，于是
\[
 A-r_c\le r_1\le r_c.
\]
若 $A>2r_c$，区间为空，没有贯穿两面的透射光路；若 $A=2r_c$，只剩双掠射极限。以下取 $0<A<2r_c$。
写 $F(r)=\arcsin(n\sin r)$，其导数
\[
 F'(r)=\frac{n\cos r}{\sqrt{1-n^2\sin^2r}}
\]
为偶函数，随 $|r|$ 严格增加（$n>1$）。因此
$\delta'(r_1)=F'(r_1)-F'(A-r_1)$ 在 $r_1<A/2$ 时为负，在 $r_1>A/2$ 时为正。
最小值确实在 $r_1=r_2=A/2$，而不只由可逆性猜出：
\[
 \boxed{\delta_{\min}=2\arcsin\!\left(n\sin\frac A2\right)-A}.
\]
偏向角上界出现在某一外侧光线掠射界面。取 $r_1=r_c$，$r_2=A-r_c$，得到
\[
 \boxed{\delta_{\max}=\frac\pi2+\arcsin[n\sin(A-r_c)]-A}.
\]
这一式在 $A<r_c$ 时仍成立，此时第二个反正弦为负，不能擅自截为零。实际非掠射光路的范围为 $\delta_{\min}\le\delta<\delta_{\max}$；若允许极限光线，则右端也可写等号。特定入射位置、有限底面和光阑还会缩小可达范围，以上是可选择入射位置的两侧面几何范围。角度统一用弧度。
\end{xesolution}

\RainbowDiagram

\Needspace{7\baselineskip}\begin{origexamplebox}{【彩虹成因】}\label{ch27:ex:s42-19-12}
雨后出现彩虹。把空中的水滴看作折射率为 $n$ 的球，解释主虹的形成，并求其特征方向。
\end{origexamplebox}\par\nopagebreak[3]
\begin{xesolution}
主虹对应“入射折射—一次内反射—出射折射”的光路，如图\ref{ch27:fig:rainbow}左侧路径。入射点 $A$ 的半径是法线。设空气侧入射角为 $i$、水滴内折射角为 $r$，$\sin i=n\sin r$。光线相对入射方向的总偏折为
\[
 D_1=\pi+2i-4r.
\]

%<optics-review:rainbow-turns>
这个偏折角可逐段相加：两次折射各使方向转过 $i-r$，球内一次反射使方向转过 $\pi-2r$，所以 $D_1=2(i-r)+(\pi-2r)$。图中不同入射高度对应不同 $i$，并非每条光都偏折相同角度。

%</optics-review:rainbow-turns>
大量邻近光线在 $\dd D_1/\dd i=0$ 的方向聚集，形成明亮的虹带。由 $\cos i=n\cos r\,\dd r/\dd i$ 得
\[
 2-4\frac{\cos i}{n\cos r}=0
 \quad\Longrightarrow\quad
 \sin^2i=\frac{4-n^2}{3}.
\]

%<optics-review:rainbow-angles>
上面的代数步骤是：驻值条件给出 $2\cos i=n\cos r$，两边平方，再用 $n\sin r=\sin i$，便有
\[
 4(1-\sin^2i)=n^2-\sin^2i,
 \qquad 3\sin^2i=4-n^2.
\]
普通水滴满足 $1<n<2$，故驻值在非掠射的允许角域内。它对应最小偏折：在 $i=0$ 附近 $D_1'=2-4/n<0$，接近掠射时 $D_1'\to2>0$，而驻值方程在 $0<i<\pi/2$ 仅有一个解。

观察到的虹半径不是 $D_1$，而是视线与反太阳方向的夹角
\[
 \rho=\pi-D_1=4r-2i.
\]
取 $n=4/3$，依次求得 $i\approx59.39^\circ$、$r\approx40.20^\circ$，于是 $D_1\approx137.97^\circ$、$\rho\approx42.03^\circ$。驻值附近入射角改变较多而出射方向改变很少，光线因此集中；实际虹的亮度与宽度还受太阳有限角宽、水滴大小及波动效应影响，并非无限亮的一条线。

%</optics-review:rainbow-angles>
对水的可见光折射率代入，主虹的观察半径约为 $42^\circ$。折射率随波长变化，红光偏折较小，故主虹红色在外、紫色在内。单个水滴只把某一颜色送入某一方向；整条虹来自视线方向不同的大量水滴。
\end{xesolution}

\Needspace{7\baselineskip}\begin{origexamplebox}{【广义彩虹】}\label{ch27:ex:s42-19-13}
网页和自然界中还能看到多种类似彩虹的彩色弧带或光环。怎样判断并解释这些图样的成因？
\end{origexamplebox}\par\nopagebreak[3]
\begin{xesolution}
不能把所有彩色圆弧都归结为同一种几何光路，应先看光源方向、弧心位置、颜色次序和介质形态。
\begin{itemize}
  \item 次虹由水滴内两次反射形成，偏折 $D_2=2\pi+2i-6r$；它位于主虹外侧、亮度较弱，颜色次序与主虹相反，主虹与次虹之间形成较暗的 Alexander 带。
  \item 雾虹仍与水滴有关，但液滴很小，衍射展宽压过色散，因而弧带宽而近白色。
  \item 日晕和月晕主要由高空六角冰晶的折射形成，常见半径约为 $22^\circ$ 或 $46^\circ$；其圆心在太阳或月亮，而普通虹的圆心在反太阳点。
  \item “宝光”或 glory 出现在反太阳方向，涉及小液滴的后向散射、表面波和衍射，不能仅用一次或两次几何反射解释。
\end{itemize}
因此，几何光学能够给出许多弧带的主方向和色序，但小尺度粒子造成的展宽、干涉和衍射必须留到波动光学中处理。
\end{xesolution}

\Needspace{8\baselineskip}\begin{origexamplebox}{固定棱镜中的不同颜色}
\label{ch27:ex:depth-abbe}
光学玻璃的阿贝数定义为 $\nu=(n_D-1)/(n_F-n_C)$。顶角 $A=60^\circ$ 的棱镜有 $n_D=1.500$、$n_C=1.495$、$\nu=40$，置于空气中。白光以同一方向入射，调节棱镜，使 $D$ 光在棱镜内的传播方向垂直于顶角平分面。此后固定棱镜，求出射的 $C$ 光和 $F$ 光之间的夹角。
\end{origexamplebox}\par\nopagebreak[3]
\begin{xesolution}
先区分三种折射率。阿贝数给出的是两条谱线的折射率差，而不是折射率与波长的比值：
\[
 n_F-n_C=\frac{n_D-1}{\nu}=\frac{0.500}{40}=0.0125,
 \qquad n_F=1.5075.
\]
$D$ 光沿对称光路传播，因此在两个侧面处的内部折射角都是 $A/2=30^\circ$。参照图~\ref{ch27:fig:s42-19-6} 右图的两个工作面，分别相对该面法线计算折射角，第一面的共同入射角为
\[
 \sin i=n_D\sin(A/2)=0.750,
 \qquad i=48.5904^\circ.
\]
固定棱镜以后，$i$ 对三种颜色相同，但内部折射角并不相同。对任意折射率 $n$，逐次计算
\[
 r_1(n)=\arcsin\!\frac{\sin i}{n},\qquad
 r_2(n)=A-r_1(n),\qquad
 e(n)=\arcsin\!\bigl[n\sin r_2(n)\bigr].
\]
角 $e$ 从第二面外法线量起。两种颜色都从第二面透出，因为各自的 $n\sin r_2<1$；不能省略这一步存在性检查。代入得
\[
 e_C=48.1589^\circ,\qquad e_F=49.2438^\circ,
 \qquad\boxed{e_F-e_C=1.0849^\circ}.
\]
两束出射光共用同一个第二面法线，故出射角之差就是它们的夹角；总偏向角中的共同 $i-A$ 会抵消。$n_F>n_C$ 对应更大的偏向，与正常色散的方向一致。若令 $n_F=n_C$，两条计算链完全相同，夹角应为零。

本题 $A=60^\circ$，不能用薄棱镜式 $\delta\simeq(n-1)A$ 代替精确计算；也不能分别把 $C,F$ 两色都取为各自的最小偏向光路，因为实验中棱镜只调节了一次。
\end{xesolution}


\section{全反射、临界条件与光纤}



图~\ref{ch27:fig:s42-19-8} 的 A、B、C 依次为麦克斯韦、爱因斯坦和费曼：他们分别将光与电磁波、光量子和量子电动力学联系起来。D 为全反射实验情景，光从半圆形介质的圆弧面沿半径入射，再在平面边界反射；精确的角度条件以下面的临界角公式判定。
\historyfigure[0.88\linewidth]{optics-history-tir.png}{电磁理论与量子光学的代表人物，以及半圆形玻璃的全反射实验。\label{ch27:fig:s42-19-8}}

当光由高折射率介质 $n_1$ 射向低折射率介质 $n_2$，折射角随入射角增大得更快。折射角达到 $90^\circ$ 时的入射角称为临界角，满足
\begin{equation}
  \sin\theta_c=\frac{n_2}{n_1},\qquad n_1>n_2.
\end{equation}
当 $i>\theta_c$ 时不再存在传播的折射光，理想界面上发生全反射。

光纤利用芯层 $n_1$ 与包层 $n_2<n_1$ 的界面使光连续全反射，从而让传播方向跟随纤维弯曲。原理虽然简单，但长距离通信要求极低损耗。高锟指出早期光纤的主要损耗来自材料杂质，而非玻璃的不可避免吸收；高纯度材料由此把光纤变成实用通信媒介，他于 2009 年获得诺贝尔物理学奖。

全反射只限制芯层内光线在界面处的角度，并不意味着任意入射光都能被光纤接收。若外界折射率为 $n_0$，子午光线相对光纤轴的最大接收角 $\theta_a$ 满足
\[
 n_0\sin\theta_a=\sqrt{n_1^2-n_2^2}.
\]
右端称为数值孔径。弯曲过急会使外弯侧的界面入射角降到临界角以下，光便从包层泄漏；杂质吸收、散射和弯曲损耗则决定能够传输的实际距离。图~\ref{ch27:fig:s42-19-9} 展示光纤应用，图~\ref{ch27:fig:s42-19-10} 左侧则把芯层、包层和反射位置分开，说明“导光”不是光在单一玻璃棒中无条件反弹。

\historyfigure[0.92\linewidth]{fiber-pioneer-built-in-v1.png}{高锟、低损耗光纤与长距离光通信。\label{ch27:fig:s42-19-9}}

\subsection{从界面临界角推到光纤接收角}
设直光纤芯层折射率为 $n_1$、包层为 $n_2<n_1$，端面外介质为 $n_0$。只讨论穿过光纤轴的子午光线。光从端面进入后与轴成角 $\alpha$；由于端面法线沿轴，端面折射满足
\[
 n_0\sin\theta=n_1\sin\alpha.
\]
在圆柱侧壁处，法线垂直于轴，所以侧壁入射角为 $i=\pi/2-\alpha$。全反射要求
\[
 \sin i=\cos\alpha>\frac{n_2}{n_1}.
\]
将 $\sin^2\alpha=1-\cos^2\alpha$ 代回端面关系，得到
\[
 n_0^2\sin^2\theta<n_1^2-n_2^2.
\]
因此 $\sqrt{n_1^2-n_2^2}$ 称为相应子午光线模型的数值孔径。若其与 $n_0$ 的比值超过1，几何接收角的上限只能是 $90^\circ$，不能对大于1的数取反正弦。等号给临界极限，实际稳定导光还需留出余量；光纤弯曲、端面倾斜与三维斜光线会改变上述简单几何。

这也给出物理检验：$n_1\to n_2$ 时，接收范围缩为零；增大芯层和包层的折射率差，接收范围变大。全反射是界面上不存在传播型折射波，不是说低折射率一侧的电磁场严格为零；后面的隐失场正用来补充这一点。


\FiberFermatPanel



\Needspace{7\baselineskip}\begin{origexamplebox}{【全反射】}\label{ch27:ex:s42-19-10}
为什么从玻璃砖的一个侧面通常看不到由相邻侧面射入砖内的光线？玻璃折射率最大为多少时，才可能从侧面看见相邻侧面入射的光？
\end{origexamplebox}\par\nopagebreak[3]
\begin{xesolution}
设光从空气经一个面进入玻璃，玻璃内相对该面法线的折射角为 $r$。相邻面法线与原法线垂直，所以光到相邻面时入射角为 $90^\circ-r$。要从相邻面射出，必须满足
\[
 90^\circ-r\le\theta_c,\qquad \sin\theta_c=\frac1n.
\]
而空气侧入射角最大为 $90^\circ$，故玻璃内最大折射角满足 $\sin r_{\max}=1/n$，即 $r_{\max}=\theta_c$。存在可出射光线的条件是 $90^\circ-\theta_c\le\theta_c$，从而
\[
 \theta_c\ge45^\circ,\qquad \boxed{n\le\sqrt2}.
\]
这里的等号只是双掠射极限：$n=\sqrt2$ 要求入射和出射的空气侧光线都沿表面掠行，不能提供有限的贯穿通光量。因此实际非零通光量的条件为 $1<n<\sqrt2$，上确界为 $\sqrt2$。常见玻璃 $n\approx1.5>\sqrt2$，相邻面上发生全反射，不能沿该路径直接看到入射光束。
\end{xesolution}

\section{费马原理与非均匀介质中的光路}



逐个界面应用直线传播、反射和折射定律，原则上可以追踪任意光线；但复杂曲面会使法线和面形设计变得困难。费马原理提供了统一表述：单色光的真实光路使光程相对于邻近光路取驻值；在色散可忽略时，也可表述为传播时间取驻值。对本讲中的常规问题，它通常表现为极小值。

在色散可忽略的静止介质中，一段几何路程 $\dd s$ 的传播时间为 $\dd t=n\dd s/c$。定义光程
\begin{equation}
  \mathcal L=\int n\,\dd s,\qquad t=\frac{\mathcal L}{c}.
\end{equation}
真实光路使 $\delta\mathcal L=0$。这一表述可推导反射定律和折射定律，也解释为什么光线宁可在低折射率介质中走较长的几何距离。若同一点发出的不同光线经光学系统后严格会聚于同一点，则各条光路的总光程必须相同，否则会聚点处相位不能一致。

以平界面为例，界面两侧固定起点与终点到界面的垂直距离分别为 $a,b$，光线在界面上的交点横坐标为 $x$，两点的水平间隔为 $L$。光程
\[
 \mathcal L(x)=n_1\sqrt{a^2+x^2}
 +n_2\sqrt{b^2+(L-x)^2}.
\]
令 $\dd\mathcal L/\dd x=0$，得到
\[
 n_1\frac{x}{\sqrt{a^2+x^2}}
 =n_2\frac{L-x}{\sqrt{b^2+(L-x)^2}},
\]
两边的几何比值分别就是 $\sin i$ 与 $\sin r$，于是恢复斯涅尔定律。若两段都在同一介质而界面改为镜面，用镜像点把折线路径展开成直线，便得到 $i=r$。因此直线传播、反射、折射并不是彼此孤立的三条经验规则，而是同一个驻光程原理在不同约束下的结果。

例如，物点 $O=(-s_o,0)$、像点 $I=(s_i,0)$，界面顶点为原点，左侧折射率为 $1$、右侧为 $n$。界面点 $P=(x,y)$ 必须满足
\begin{equation}
 \sqrt{(x+s_o)^2+y^2}+n\sqrt{(x-s_i)^2+y^2}=s_o+n s_i.
\end{equation}
这是一类高次非球面。它说明即使“点到点严格成像”这一基本要求，也往往需要复杂且高精度的曲面；人的角膜与晶状体系统则是自然形成的精密光学系统。





\subsection{驻值方程如何成为折射定律}
将平界面取为横轴，起点在 $(0,a)$、终点在 $(L,-b)$，两侧折射率为 $n_1,n_2$，其中 $a,b>0$。经过界面点 $(x,0)$ 的候选折线路径，其光程为
\[
 \mathcal L(x)=n_1\sqrt{a^2+x^2}+n_2\sqrt{b^2+(L-x)^2}.
\]
求导时，第二段中的 $L-x$ 对 $x$ 求导为 $-1$，因此
\[
 \mathcal L'(x)=n_1\frac{x}{\sqrt{a^2+x^2}}
              -n_2\frac{L-x}{\sqrt{b^2+(L-x)^2}}.
\]
两个分式分别是两侧光线与法线夹角的正弦。令 $\mathcal L'=0$ 就得到斯涅尔定律，而不是一个独立的新折射公式。再求导，
\[
 \mathcal L''(x)=\frac{n_1a^2}{(a^2+x^2)^{3/2}}
                +\frac{n_2b^2}{[b^2+(L-x)^2]^{3/2}}>0,
\]
所以这一具体问题的驻值确实是唯一极小值。一般光学系统中的驻值不一定是最短路；不能把这里的正二阶导数结论推广到所有反射面、焦点和受约束光路。

分层介质中，每一薄层界面均平行，故相对层面法线的角 $i$ 满足 $n\sin i=\text{常量}$。如果改用相对层面的角 $\alpha$，同一规律写成 $n\cos\alpha=\text{常量}$。两式并不矛盾，差别只是角度的起算方向。建立微分方程前，应把图中的角与公式中的角核对一次。


% Relocated FisherBayDiagram to its own solution.

\Needspace{7\baselineskip}\begin{origexamplebox}{【运动学？】}\label{ch27:ex:s42-19-6}
湖湾由从顶点 $D$ 出发的两条直岸围成，夹角为 $\alpha$，$0<\alpha<\pi$。渔人的房子 $A$ 在一岸的陆地一侧，$C$ 是 $A$ 到该岸的垂足，$AC=h$，$AD=l$，且 $0<h<l$、$C$ 在岸线射线上。朋友的房子 $B$ 关于湖湾角平分线与 $A$ 对称。陆上步行速度为 $v$，划船速度为 $v/2$，可在岸上任意位置上下船。规划从 $A$ 到 $B$ 的最省时路线。
\end{origexamplebox}\par\nopagebreak[3]
\begin{xesolution}
\FisherBayDiagram

先分清陆地和水域。图~\ref{ch27:fig:s42-19-15} 中 $DC$ 是一侧岸线，$A,B$ 都不在水中；因此不能从房子直接划船。令
\[
 c=DC=\sqrt{l^2-h^2},\qquad \beta=\alpha/2.
\]
两岸上船、下船点记为 $P,Q$，$DP=t_1,DQ=t_2$。直线步行段均在各自陆地半平面内，船的直线段在湖湾内。其时间为
\[
 vT=\sqrt{h^2+(t_1-c)^2}+\sqrt{h^2+(t_2-c)^2}
 +2\sqrt{t_1^2+t_2^2-2t_1t_2\cos\alpha}.
\]
为什么可以取对称路线？前两项是严格凸函数，第三项是线性坐标差的欧氏范数，也是凸函数；交换 $t_1,t_2$ 不改变时间。将一条路线与它的对称路线取参数平均，凸性给出
$T((t_1+t_2)/2,(t_1+t_2)/2)\le T(t_1,t_2)$。
所以最优点必有 $t_1=t_2=t\ge0$。这不是仅凭图形对称猜测。
此时 $PQ=2t\sin\beta$，有
\[
 T(t)=\frac2v\left[\sqrt{h^2+(t-c)^2}+2t\sin\beta\right],
 \quad
 \frac v2T'(t)=\frac{t-c}{\sqrt{h^2+(t-c)^2}}+2\sin\beta.
\]
令 $k=2\sin\beta$，且
$T''(t)=2h^2/[v(h^2+(t-c)^2)^{3/2}]>0$。
由于 $T'(0)=2(k-c/l)/v$，答案完整分为
\[
 \boxed{\begin{cases}
 k\ge c/l:&t_*=0,\quad T_{\min}=2l/v;\\[2pt]
 0<k<c/l:&t_*=c-\dfrac{kh}{\sqrt{1-k^2}},\quad
 T_{\min}=\dfrac2v\left[h\sqrt{1-k^2}+kc\right].
 \end{cases}}
\]
第一种情形沿陆地 $A\to D\to B$ 行走，不划船；第二种情形沿 $A\to P$ 步行、$P\to Q$ 划船、$Q\to B$ 步行。
后一条件自动保证 $k<1$ 和 $t_*>0$。两式在 $k=c/l$ 处连续，时间单位均为长度除以速度；$h$ 不由 $l\sin\alpha$ 决定。
同岸的多次上下船不会更快，因为把这段船路改为沿同岸直线步行既不更长、速度又更大。于是上述一次横渡或完全绕岸的比较已覆盖最优路线。
它体现了费马最短时间思想：总距离最短不等于总时间最短。
\end{xesolution}

\Needspace{7\baselineskip}\begin{origexamplebox}{【严格聚焦】}\label{ch27:ex:s42-19-7}
球面镜不能把所有平行于轴的入射光严格会聚于一点。能够完成这一功能的反射面应具有何种形状？
\end{origexamplebox}\par\nopagebreak[3]
\begin{xesolution}
取顶点为原点、轴向右为 $x$ 正方向，焦点 $F=(f,0)$，$f>0$。平行光必须从右向左入射，照到向右开口的反射面。参考平面取 $x=a$，位于整个镜面右侧；反射点为 $P=(x,y)$。轴上光路长为 $a+f$，等光程条件为
\[
 (a-x)+\sqrt{(x-f)^2+y^2}=a+f.
\]
故 $PF=x+f$。平方并消去 $x^2+f^2$，得到 $y^2=4fx$；它也表明 $PF$ 等于 $P$ 到准线 $x=-f$ 的距离。因此子午截面是抛物线，旋转后为抛物面：
\[
 \boxed{y^2=4fx}
\]
也可直接检验反射定律。曲面法向可取 $\boldsymbol N=(-2f,y)$，入射单位方向为 $\boldsymbol k=(-1,0)$。反射方向
\[
 \boldsymbol k'=\boldsymbol k-2\frac{\boldsymbol k\cdot\boldsymbol N}{\boldsymbol N\cdot\boldsymbol N}\boldsymbol N
 =\frac{(4f^2-y^2,-4fy)}{4f^2+y^2}
 =\frac{F-P}{|F-P|}.
\]
所以反射光真正射向 $F$，而非仅延长线经过它。抛物面反射镜对轴上平行光严格无球差；有限口径带来衍射，制造误差和离轴入射仍会带来像差。
\end{xesolution}

\Needspace{7\baselineskip}\begin{origexamplebox}{【变折射率折射】}\label{ch27:ex:s42-19-5}
介质位于 $x>0,y>0$，折射率仅随 $y$ 连续变化。一束细光束沿 $x$ 方向垂直入射，并沿半径为 $R$ 的圆弧传播。求 $n(y)$。已知 $n(0)=n_0=1$，材料最大折射率为 $2.5$，求圆弧允许的最大圆心角。
\end{origexamplebox}\par\nopagebreak[3]
\begin{xesolution}
\GradientCircleDiagram

%<optics-review:gradient-geometry>
先把圆弧的角度与高度联系起来。以入射点为 $(0,0)$、圆心为 $(0,R)$，沿光的传播方向用参数 $\theta$ 写圆弧：
\[
 x=R\sin\theta,\qquad y=R(1-\cos\theta).
\]
微分得到 $(\dd x,\dd y)=R(\cos\theta,\sin\theta)\dd\theta$，所以 $\theta$ 同时就是切线与水平线的夹角。这便给出后面所用的 $\cos\theta=1-y/R$。水平薄层的法线竖直，斯涅尔定律中的入射角为 $\pi/2-\theta$，正弦因而变为 $\cos\theta$，不能在此误写成 $n\sin\theta$ 守恒。允许轨迹在 $0\le\theta<\pi/2$ 内；材料上限最终把高度限制到 $0\le y\le0.6R$，尚未到圆弧最高点。

%</optics-review:gradient-geometry>

把介质分成许多水平薄层。斯涅尔定律在各层界面上连续应用，得到守恒量 $n(y)\cos\theta=n_0$，其中 $\theta$ 是光线切线与 $x$ 轴的夹角。对图~\ref{ch27:fig:s42-19-14} 的圆弧，圆心在入射点正上方，有
\[
 \cos\theta=1-\frac{y}{R}.
\]
因此
\[
 \boxed{n(y)=\frac{n_0}{1-y/R}}.
\]
材料条件 $n\le2.5$ 给出 $\cos\theta\ge1/2.5=0.4$，故
\[
 \boxed{\theta_{\max}=\arccos0.4\approx66.4^\circ}.
\]
该角正是从入射半径转到当前半径的圆心角；若把它错当成相对法线的角，会得到互余的错误答案。
\end{xesolution}

\Needspace{7\baselineskip}\begin{origexamplebox}{【圆形光线】}\label{ch27:ex:s42-19-8}
一颗半径为 $R$ 的行星，其大气折射率随高度 $h$ 线性变化：$n=n_0-\alpha h$，其中 $n_0$ 是表面折射率。若某高度上沿行星表面切向发出的光恰能绕行星作圆周传播，求该高度 $h_0$。
\end{origexamplebox}\par\nopagebreak[3]
\begin{xesolution}
%<optics-review:circular-ray-proof>
关键是先解释圆轨道条件从何而来。设 $s$ 是沿光路的弧长，$\boldsymbol t=\dd\boldsymbol r/\dd s$ 为单位切向量。固定光路两端，只让中间各点有微小位移 $\delta\boldsymbol r$。光程 $\mathcal L=\int n\,\dd s$ 的一阶变化为
\[
 \delta\mathcal L=\int\!\left(\nabla n\cdot\delta\boldsymbol r
       +n\boldsymbol t\cdot\frac{\dd(\delta\boldsymbol r)}{\dd s}\right)\dd s.
\]
第二项分部积分，端点项因两端固定而为零。任意中间位移都要满足驻值条件，故
\[
 \frac{\dd(n\boldsymbol t)}{\dd s}=\nabla n.
\]
这说明折射率梯度怎样使光线弯曲。对半径为 $r$ 的圆，$n$ 沿圆周不变，单位切向量每走弧长 $\dd s$ 就向圆心转过 $\dd s/r$，即
\[
 \frac{\dd\boldsymbol t}{\dd s}=-\frac{\boldsymbol e_r}{r},
 \qquad \nabla n=\frac{\dd n}{\dd r}\boldsymbol e_r.
\]
因此 $-n/r=\dd n/\dd r$，整理即 $\dd(nr)/\dd r=0$。这个条件来自光线的局部弯曲规律；只说“沿圆周 $nr$ 为常量”还不足以推出它。

%</optics-review:circular-ray-proof>

球对称介质的射线守恒量为 $nr\sin\psi$，$\psi$ 是光线与径向的夹角。圆形轨道要求 $\psi=90^\circ$，并且 $n(r)r$ 在该半径处取驻值：
\[
 \frac{\dd}{\dd r}[n(r)r]=0,\qquad r=R+h.
\]
代入 $n=n_0-\alpha h$ 得
\[
 n_0-\alpha h-\alpha(R+h)=0,\qquad
 \boxed{h_0=\frac{n_0-\alpha R}{2\alpha}}.
\]
结果还必须满足 $h_0\ge0$ 且 $n(h_0)>0$；若材料参数不满足这些条件，物理解不存在。

%<optics-review:circular-ray-domain>
若 $n_0>0$，有限半径圆轨道要求折射率向外减小，即 $\alpha>0$；$\alpha=0$ 时光在均匀介质中沿直线传播，不能代入含 $1/\alpha$ 的结果。此时 $h_0\ge0$ 等价于 $n_0\ge\alpha R$，并且
\[
 r_0=\frac{n_0+\alpha R}{2\alpha},\qquad
 n(h_0)=\frac{n_0+\alpha R}{2}=\alpha r_0>0.
\]
还应确认算出的高度处于题设线性大气模型的适用范围；若要求严格在表面上方，则需 $n_0>\alpha R$。$n$ 无量纲，$\alpha$ 的量纲为长度的倒数，所得 $h_0$ 确为长度。本题求的是恰好切向发射的圆轨道，不意味着偏离该条件后仍能绕圆传播。

%</optics-review:circular-ray-domain>
\end{xesolution}

\section{界面上的电磁波与隐失场（选读）}

\subsection{反射、透射与偏振分量}


光线模型给出反射角和折射角，却不能回答反射光有多强、相位是否突变，以及为何反射率与偏振有关。
设折射率为 $n_1$ 的非磁性介质中的平面波，以入射角 $\theta_i$ 射向折射率为 $n_2$ 的介质，折射角为 $\theta_t$。
界面平移对称使频率和切向波矢连续，因而
\begin{equation}
 n_1\sin\theta_i=n_2\sin\theta_t.
\end{equation}
入射光、反射光和法线所在平面称入射面。电场垂直入射面的模式称 s 偏振，电场平行入射面的模式称 p 偏振。

\begin{center}\begin{minipage}{\linewidth}\centering
\begin{tikzpicture}[scale=.92,>=Latex]
  \fill[DeepBlue!6] (-4,-2.55) rectangle (4,0);
  \draw[very thick] (-4,0)--(4,0) node[right] {界面};
  \draw[dashed,guide] (0,-2.45)--(0,2.65) node[above] {法线};
  \coordinate (O) at (0,0);
  \draw[->,DeepBlue,very thick] (-3.0,2.2)--(O);
  \draw[->,curve,very thick] (O)--(3.0,2.2);
  \draw[->,DeepBlue,very thick] (O)--(1.55,-2.30);
  \node[above] at (-2.3,2.35) {入射};
  \node[above] at (2.4,2.35) {反射};
  \node[right] at (1.8,-2.3) {折射};
  \draw[->,guide] (0,.78) arc[start angle=90,end angle=144,radius=.78];
  \node at (-.48,.92) {$\theta_i$};
  \draw[->,guide] (0,-.80) arc[start angle=-90,end angle=-56,radius=.80];
  \node[left] at (-.65,-1.35) {$\theta_t$};
  \draw[guide,thin] (-.60,-1.32)--(.22,-.62);
  \node[left] at (-3.85,1.72) {$n_1$};
  \node[left] at (-3.85,-1.72) {$n_2$};
  \node[DeepBlue] at (-2.55,-.55) {$\odot\;s$：电场垂直纸面};
  \draw[->,DeepBlue,thick] (2.15,-1.05)--(2.78,-.62);
  \node[DeepBlue,right] at (2.12,-1.22) {$p$：电场在入射面内};
\end{tikzpicture}
\figcap{s、p 偏振的定义依赖入射面，而不是依赖纸张本身。图中折射光向法线偏折，对应 $n_2>n_1$。}
\end{minipage}\end{center}

规定界面法向 $\boldsymbol N$ 指向第二介质，$\boldsymbol s$ 垂直入射面，对入射、反射、透射各束分别取
$\boldsymbol p_j=\boldsymbol s\times\hat{\boldsymbol k}_j$。
这一右手构造使反射 p 基的切向方向与入射 p 基相反；因此振幅系数的正负必须连同基矢理解。
无自由面电荷、面电流且介质非磁性时，切向 $\boldsymbol E,\boldsymbol H$ 连续。用 $H=nE/Z_0$，s 分量满足
\[
 1+r_s=t_s,\qquad n_1\cos\theta_i(1-r_s)=n_2\cos\theta_t\,t_s.
\]
消去 $t_s$，可得下列第一行。p 分量则满足
\[
 \cos\theta_i(1-r_p)=\cos\theta_t\,t_p,\qquad
 n_1(1+r_p)=n_2t_p.
\]
第二式给 $t_p=n_1(1+r_p)/n_2$，代回第一式即得第二行。这就是菲涅尔公式（Fresnel equations）：
\begin{align}
 r_s&=\frac{n_1\cos\theta_i-n_2\cos\theta_t}
 {n_1\cos\theta_i+n_2\cos\theta_t},&
 t_s&=\frac{2n_1\cos\theta_i}
 {n_1\cos\theta_i+n_2\cos\theta_t},\\
 r_p&=\frac{n_2\cos\theta_i-n_1\cos\theta_t}
 {n_2\cos\theta_i+n_1\cos\theta_t},&
 t_p&=\frac{2n_1\cos\theta_i}
 {n_2\cos\theta_i+n_1\cos\theta_t}.
\end{align}
这里 $r,t$ 是按各自 s、p 正交基定义的电场复振幅比；p 基在反射前后的指向约定会影响 $r_p$ 的符号，但不影响可观测反射率。
对无吸收介质，能流反射率和透射率为
\begin{equation}
 R_{s,p}=|r_{s,p}|^2,
 \qquad
 T_{s,p}=\frac{n_2\cos\theta_t}{n_1\cos\theta_i}|t_{s,p}|^2,
 \qquad R_{s,p}+T_{s,p}=1.
\end{equation}
透射率中不能漏掉折射率和光束投影面积因子。正入射且 $n_2>n_1$ 时，反射电场相对入射电场发生 $\pi$ 相位突变，
这就是薄膜干涉中常说的“半波损失”；由高折射率侧向低折射率侧反射则没有这一突变。

对 p 偏振，当 $r_p=0$ 时反射光消失。结合折射定律可得
\begin{equation}
 \theta_i+\theta_t=\frac\pi2,
 \qquad
 \boxed{\tan\theta_B=\frac{n_2}{n_1}}.
\end{equation}
$\theta_B$ 称布儒斯特角（Brewster angle）。自然光在该角附近反射后，以 s 偏振为主；气体激光器中的布儒斯特窗正利用
这一性质降低 p 偏振的反射损耗。

\subsection{全反射中的隐失场}


当 $n_1>n_2$ 且 $\theta_i>\theta_c=\arcsin(n_2/n_1)$ 时，第二介质中不再有向法线方向传播的平面波，
但界面下方的电磁场并不为零。取界面为 $y=0$，第二介质位于 $y<0$，切向波数为
\[
 k_x=k_0n_1\sin\theta_i,
\]
其中 $k_0=\omega/c=2\pi/\lambda_0$，$\lambda_0$ 是真空波长，统一采用时间因子 $e^{-\ii\omega t}$。由 $k_x^2+k_y^2=n_2^2k_0^2$，第二介质中的法向波数成为纯虚数：
\begin{equation}
 k_y=-\ii\kappa,
 \qquad
 \kappa=k_0\sqrt{n_1^2\sin^2\theta_i-n_2^2}.
\end{equation}
于是透入第二介质的场具有
\begin{equation}
 E_t\propto e^{\ii k_xx}e^{\kappa y},
 \qquad
 \delta=\frac1\kappa
\end{equation}
的形式，$\delta$ 是振幅衰减到 $1/e$ 的穿透深度，强度相应衰减到 $e^{-2}$。选择 $k_y=-\ii\kappa$ 是为了使 $y\to-\infty$ 时场衰减，而不是无源增长。远离临界角且折射率差不小时，$\delta$ 常与波长同量级；当 $\theta_i\to\theta_c^+$ 时，$\kappa\to0$、$\delta\to\infty$，不可概括为始终只有一个波长。

\begin{center}\begin{minipage}{\linewidth}\centering
\begin{tikzpicture}[scale=.90,>=Latex]
  \fill[DeepBlue!6] (-4,-2.55) rectangle (4,0);
  \draw[very thick] (-4,0)--(4,0);
  \draw[dashed,guide] (0,-2.35)--(0,2.55) node[above] {法线};
  \draw[->,DeepBlue,very thick] (-3.0,2.05)--(0,0);
  \draw[->,curve,very thick] (0,0)--(3.0,2.05);
  \draw[DeepBlue,very thick] plot[smooth,domain=0:-2.25,samples=80,variable=\y]
    ({.9*exp(1.15*\y)},\y);
  \draw[->,DeepBlue,thick] (1.5,-1.45)--(3.4,-1.45);
  \node at (2.5,-1.8) {沿界面相位传播};
  \node[anchor=west] at (1.3,-.45) {$|E(y)|\propto e^{\kappa y}$};
  \node[anchor=west] at (1.3,-.90) {横向偏移表示振幅};
  \draw[<->,guide] (-1.15,0)--(-1.15,-.87) node[midway,left] {$\delta$};
  \node[left] at (-3.75,1.50) {$n_1$};
  \node[left] at (-3.75,-1.60) {$n_2<n_1$};
\end{tikzpicture}
\figcap{全反射与法向振幅包络：低折射率介质内的实线不是光线轨迹，其横向偏移表示振幅，法向不发生正弦振荡。\label{ch27:fig:evanescent}}
\end{minipage}\end{center}

图\ref{ch27:fig:evanescent}中 $y=-\delta$ 的振幅为界面处的 $1/e$。对于无损、无限厚第二介质，隐失场沿法向没有平均净能流，却可沿界面输运能量和相位。若在距离界面数个 $\delta$ 内放入另一块高折射率介质，
隐失场可耦合成传播波，出现受抑全反射；从量子图像看，它与隧穿具有相同的指数衰减结构。

\begin{starread}{有限光束在全反射中的微小位移}
平面波全反射系数的模为 $1$，但相位 $\phi_{s,p}$ 随入射角变化。真实窄光束是邻近切向波数 $k_x$ 的角谱叠加，
各分量反射后得到略有差别的相位；用驻相条件求反射光束中心，得到古斯--汉欣位移（Goos--H\"anchen shift）
\[
 \Delta_{\rm GH}=-\frac{\partial\phi_{s,p}}{\partial k_x}.
\]
它表现为反射点沿入射面方向移动，数量级通常接近波长。若进一步区分圆偏振分量，还会出现垂直入射面的
安贝尔--费多罗夫位移（Imbert--Fedorov shift）；它与光的自旋角动量和界面处的角动量守恒有关。
\end{starread}



\section{程函与光线方程（选读）}



在折射率缓慢变化的介质中，把单色标量场写成
\[
 \Psi(\vect r)=A(\vect r)e^{\ii k_0S(\vect r)}.
\]
这里 $S$ 的量纲是长度，$k_0S$ 才是无量纲相位。在无损、各向同性且 $n$ 缓变的标量近似中，亥姆霍兹方程为
$\nabla^2\Psi+k_0^2n^2\Psi=0$。代入上式并除去指数因子，逐项得到
\[
 \nabla^2A+2\ii k_0\nabla A\cdot\nabla S+\ii k_0A\nabla^2S
 +k_0^2A[n^2-|\nabla S|^2]=0.
\]
当振幅和折射率变化尺度远大于波长时，最高阶 $k_0^2$ 项先满足零，得到程函方程（eikonal equation）
\begin{equation}
 |\nabla S|=n(\vect r).
\end{equation}
等相位面法线给出光线方向。令弧长参数为 $s$、单位切向为
$\vect t=\nabla S/n$。沿光线求导为 $\dd/\dd s=\vect t\cdot\nabla$，再用梯度场无旋的恒等式
\[
 (\nabla S\cdot\nabla)\nabla S
 =\tfrac12\nabla|\nabla S|^2=\tfrac12\nabla n^2=n\nabla n,
\]
便有 $\dd(n\vect t)/\dd s=(\nabla S\cdot\nabla)\nabla S/n=\nabla n$，即光线方程
\begin{equation}
 \boxed{\frac{\dd}{\dd s}(n\vect t)=\nabla n}.
\end{equation}
它的切向分量是恒等式，法向分量说明光线向折射率增大的方向弯曲。若 $n=n(y)$，沿 $x$ 方向的平移对称性给出
\begin{equation}
 n(y)\sin\theta=\text{常量},
\end{equation}
其中 $\theta$ 是光线与层面法线的夹角；这就是分层介质中的连续折射定律。

同一光线方程也可由费马原理（Fermat principle）得到：固定端点之间的实际光路使光程取平稳值，
\begin{equation}
 \delta\int_A^B n(\vect r)\,\dd s=0.
\end{equation}
“平稳”不等同于“必为最短”，反射和某些多路径系统中可能取得极大值或鞍点；正确判据是一阶变分为零。

\Needspace{7\baselineskip}\begin{origexamplebox}{梯度折射率介质中的正弦光线}
平面分层介质的折射率在轴附近为
\[
 n(y)=n_0\left(1-\frac12g^2y^2\right),
 \qquad |gy|\ll1,
\]
其中 $g$ 的量纲为长度的倒数。一条近轴光线从 $x=0$ 处以 $y(0)=y_0$、$y'(0)=\alpha_0$ 入射。
求其轨迹，并说明这种介质为何能周期性会聚光束。

\begin{center}\begin{minipage}{\linewidth}\centering
\begin{tikzpicture}[x=.90cm,y=1.0cm,>=Latex]
  \fill[DeepBlue!6] (0,-2.0) rectangle (9,2.0);
  \draw[axis,->] (0,0)--(9.35,0) node[right] {$x$};
  \draw[axis,->] (0,-2.05)--(0,2.25) node[above] {$y$};
  \foreach \yy in {-1.5,-1,-.5,.5,1,1.5}{\draw[guide,dashed] (0,\yy)--(9,\yy);}
  \draw[curve,very thick,->,domain=0:8.8,samples=160]
    plot (\x,{1.22*cos(60*\x)+.20*sin(60*\x)});
  \draw[curve,thick,densely dashed,->,domain=0:8.8,samples=160]
    plot (\x,{.70*cos(60*\x)+.20*sin(60*\x)});
  \node[left] at (-.15,1.22) {$y_0$};
  \fill[DeepBlue] (1.5,.2) circle (1.5pt);
  \node[above] at (4.5,2.02) {$n$ 在轴上最大};
  \draw[<->,guide] (0,-2.35)--(6,-2.35) node[midway,below] {$2\pi/g$};
\end{tikzpicture}
\figcap{两条初始倾角相同、初始高度不同的光线在四分之一周期处相交；纵向尺度放大以辨认弯曲。空间周期为 $2\pi/g$。\label{ch27:fig:grin}}
\end{minipage}\end{center}
\end{origexamplebox}\par\nopagebreak[3]
\begin{xesolution}
近轴时 $s\simeq x$、$\vect t\simeq(1,y')$。光线方程的 $y$ 分量给出
\[
 \frac{\dd}{\dd x}\left(n\frac{\dd y}{\dd x}\right)=\frac{\dd n}{\dd y}.
\]
保留到 $y$ 和倾角的一阶量，左边可把 $n$ 取为 $n_0$，右边
$\dd n/\dd y=-n_0g^2y$，所以
\[
 y''+g^2y=0.
\]
结合初始条件，
\begin{equation}
 \boxed{y(x)=y_0\cos(gx)+\frac{\alpha_0}{g}\sin(gx)}.
\end{equation}
图\ref{ch27:fig:grin}中两条光线具有相同的 $\alpha_0$，在 $x=\pi/(2g)$ 处均有 $y=\alpha_0/g$，所以会聚点不必在轴上。光线以 $2\pi/g$ 为空间周期反复穿过轴附近。若入射的是平行于轴的光束，各条光线 $\alpha_0=0$ 而 $y_0$ 不同，则在
$x=(\pi/2+j\pi)/g$ 处同时到达 $y=0$，形成周期会聚；若来自同一轴上物点，则 $y_0=0$、$\alpha_0$ 不同，在 $x=j\pi/g$ 再会聚。一般任意一束光并不自动在同一点相交，必须说明其初始位置和倾角之间的关系。$g$ 的量纲为 $\mathrm{m^{-1}}$，因此 $gx$ 无量纲，轨迹也通过量纲检查。
若离轴过远使 $|gy|$ 不再远小于 $1$，二次折射率近似与正弦结论都会失效。
\end{xesolution}

\endgroup
\begingroup\tcbset{breakable/.default=false}
\renewcommand{\Needspace}[1]{\par\begingroup
\dimen0=#1\relax\vskip0pt plus\dimen0\penalty-100
\vskip0pt plus-\dimen0\vskip\dimen0\penalty9999\vskip-\dimen0
\vskip0pt\endgroup}
\chapter{成像与光学仪器}


\graphicspath{{/Users/dongxingxi/Documents/物理竞赛一轮资料/sample/xes-original/forty-third-section/}{/Users/dongxingxi/Documents/物理竞赛一轮资料/sample/xes-original/forty-third-section/assets/}}


\section{成像概念、光学器件与近轴条件}


\newcommand{\SphereRefractionPanel}{%
\begin{center}\begin{minipage}{\linewidth}\centering
\begin{tikzpicture}[>=Latex,font=\small,every node/.style={fill=none}]
\draw[thick] plot coordinates {(0.3616959,1.1471529) (0.3478034,1.1270521) (0.3341575,1.1067831) (0.3207603,1.0863489) (0.3076137,1.0657526) (0.2947197,1.0449971) (0.2820802,1.0240857) (0.2696972,1.0030215) (0.2575724,0.9818075) (0.2457077,0.9604470) (0.2341048,0.9389431) (0.2227655,0.9172991) (0.2116915,0.8955182) (0.2008844,0.8736036) (0.1903459,0.8515586) (0.1800775,0.8293865) (0.1700807,0.8070906) (0.1603570,0.7846742) (0.1509079,0.7621408) (0.1417349,0.7394935) (0.1328391,0.7167359) (0.1242221,0.6938713) (0.1158851,0.6709031) (0.1078293,0.6478348) (0.1000559,0.6246698) (0.0925661,0.6014116) (0.0853610,0.5780636) (0.0784417,0.5546293) (0.0718092,0.5311122) (0.0654645,0.5075159) (0.0594085,0.4838438) (0.0536423,0.4600995) (0.0481665,0.4362865) (0.0429820,0.4124084) (0.0380897,0.3884687) (0.0334902,0.3644711) (0.0291842,0.3404190) (0.0251724,0.3163161) (0.0214553,0.2921661) (0.0180336,0.2679724) (0.0149077,0.2437387) (0.0120781,0.2194686) (0.0095452,0.1951658) (0.0073094,0.1708338) (0.0053710,0.1464764) (0.0037304,0.1220971) (0.0023877,0.0976995) (0.0013432,0.0732874) (0.0005970,0.0488644) (0.0001493,0.0244340) (0.0000000,0.0000000) (0.0001493,-0.0244340) (0.0005970,-0.0488644) (0.0013432,-0.0732874) (0.0023877,-0.0976995) (0.0037304,-0.1220971) (0.0053710,-0.1464764) (0.0073094,-0.1708338) (0.0095452,-0.1951658) (0.0120781,-0.2194686) (0.0149077,-0.2437387) (0.0180336,-0.2679724) (0.0214553,-0.2921661) (0.0251724,-0.3163161) (0.0291842,-0.3404190) (0.0334902,-0.3644711) (0.0380897,-0.3884687) (0.0429820,-0.4124084) (0.0481665,-0.4362865) (0.0536423,-0.4600995) (0.0594085,-0.4838438) (0.0654645,-0.5075159) (0.0718092,-0.5311122) (0.0784417,-0.5546293) (0.0853610,-0.5780636) (0.0925661,-0.6014116) (0.1000559,-0.6246698) (0.1078293,-0.6478348) (0.1158851,-0.6709031) (0.1242221,-0.6938713) (0.1328391,-0.7167359) (0.1417349,-0.7394935) (0.1509079,-0.7621408) (0.1603570,-0.7846742) (0.1700807,-0.8070906) (0.1800775,-0.8293865) (0.1903459,-0.8515586) (0.2008844,-0.8736036) (0.2116915,-0.8955182) (0.2227655,-0.9172991) (0.2341048,-0.9389431) (0.2457077,-0.9604470) (0.2575724,-0.9818075) (0.2696972,-1.0030215) (0.2820802,-1.0240857) (0.2947197,-1.0449971) (0.3076137,-1.0657526) (0.3207603,-1.0863489) (0.3341575,-1.1067831) (0.3478034,-1.1270521) (0.3616959,-1.1471529)};
\draw[gray,thin,->] (-4.6000000,0.0000000)--(7.8995696,0.0000000);
\fill (-4.0000000,0.0000000) circle (1.2pt);
\node[below,fill=none,inner sep=2pt] at (-4.0000000,0.0000000) {$S$};
\fill (0.0000000,0.0000000) circle (1.2pt);
\node[below left,fill=none,inner sep=2pt] at (0.0000000,0.0000000) {$V$};
\fill (2.0000000,0.0000000) circle (1.2pt);
\node[below,fill=none,inner sep=2pt] at (2.0000000,0.0000000) {$C$};
\fill (7.4995696,0.0000000) circle (1.2pt);
\node[below,fill=none,inner sep=2pt] at (7.4995696,0.0000000) {$P$};
\fill (0.0482051,0.4364592) circle (1.2pt);
\node[above left,fill=none,inner sep=2pt] at (0.0482051,0.4364592) {$A$};
\draw[gray,dashed,thin] (2.0000000,0.0000000)--(0.0482051,0.4364592);
\draw[DeepBlue,thick,->] (-4.0000000,0.0000000)--(0.0482051,0.4364592);
\draw[DeepBlue,thick,->] (0.0482051,0.4364592)--(7.4995696,0.0000000);
\node[above,fill=none,inner sep=2pt] at (-2.0000000,1.0000000) {$n_1=1$};
\node[above,fill=none,inner sep=2pt] at (3.0000000,1.0000000) {$n_2=2$};
\end{tikzpicture}
\figcap{球面顶点 $V$、球心 $C$ 与球面点 $A$；图中 $P$ 为所示有限孔径光线的轴交点。}\label{ch28:fig:s43-20-1}
\end{minipage}\end{center}
}

\newcommand{\VirtualSignPanel}{%
\begin{center}\begin{minipage}{\linewidth}\centering
\begin{tikzpicture}[>=Latex,font=\small,every node/.style={fill=none}]
\begin{scope}[shift={(-3.2000000,0.0000000)}]
\draw[thick] (0.0000000,-1.4000000)--(0.0000000,1.4000000);
\draw[gray,thin,->] (-2.7000000,0.0000000)--(2.7000000,0.0000000);
\draw[DeepBlue,thick,->] (-2.0000000,0.0000000)--(0.0000000,-0.5500000);
\draw[DeepBlue,thick,->] (0.0000000,-0.5500000)--(2.4606300,-0.9919275);
\draw[gray,dashed,thin] (-3.0623724,0.0000000)--(0.0000000,-0.5500000);
\draw[DeepBlue,thick,->] (-2.0000000,0.0000000)--(0.0000000,0.5500000);
\draw[DeepBlue,thick,->] (0.0000000,0.5500000)--(2.4606300,0.9919275);
\draw[gray,dashed,thin] (-3.0623724,0.0000000)--(0.0000000,0.5500000);
\fill (-3.0623724,0.0000000) circle (1.2pt);
\node[below,fill=none,inner sep=2pt] at (-3.0623724,0.0000000) {};
\draw[gray,dashed,thin] (-3.0623724,0.0000000)--(-3.0623724,-0.6000000);
\node[below,fill=none,inner sep=2pt] at (-3.0623724,-0.6500000) {虚像};
\fill (-2.0000000,0.0000000) circle (1.2pt);
\node[below,fill=none,inner sep=2pt] at (-2.0000000,0.0000000) {};
\draw[gray,dashed,thin] (-2.0000000,0.0000000)--(-2.0000000,0.8500000);
\node[above,fill=none,inner sep=2pt] at (-2.0000000,0.9000000) {实物};
\node[below,fill=none,inner sep=2pt] at (0.0000000,-1.6000000) {空气到玻璃};
\end{scope}
\begin{scope}[shift={(3.2000000,0.0000000)}]
\draw[thick] (0.0000000,-1.4000000)--(0.0000000,1.4000000);
\draw[gray,thin,->] (-2.7000000,0.0000000)--(2.7000000,0.0000000);
\draw[DeepBlue,thick,->] (-2.4105135,-1.2128912)--(0.0000000,-0.5500000);
\draw[gray,dashed,thin] (0.0000000,-0.5500000)--(2.0000000,0.0000000);
\draw[DeepBlue,thick,->] (0.0000000,-0.5500000)--(1.2687483,0.0000000);
\draw[DeepBlue,thick,->] (-2.4105135,1.2128912)--(0.0000000,0.5500000);
\draw[gray,dashed,thin] (0.0000000,0.5500000)--(2.0000000,0.0000000);
\draw[DeepBlue,thick,->] (0.0000000,0.5500000)--(1.2687483,0.0000000);
\fill (1.2687483,0.0000000) circle (1.2pt);
\node[below,fill=none,inner sep=2pt] at (1.2687483,0.0000000) {};
\draw[gray,dashed,thin] (1.2687483,0.0000000)--(1.2687483,-0.6000000);
\node[below,fill=none,inner sep=2pt] at (1.2687483,-0.6500000) {实像};
\fill (2.0000000,0.0000000) circle (1.2pt);
\node[below,fill=none,inner sep=2pt] at (2.0000000,0.0000000) {};
\draw[gray,dashed,thin] (2.0000000,0.0000000)--(2.0000000,0.8500000);
\node[above,fill=none,inner sep=2pt] at (2.0000000,0.9000000) {虚物};
\node[below,fill=none,inner sep=2pt] at (0.0000000,-1.6000000) {玻璃到空气};
\end{scope}
\end{tikzpicture}
\figcap{左图出射光的反向延长线交于虚像；右图虚线为未遇界面时的会聚路径。有限孔径轴交点随角度略变。}\label{ch28:fig:s43-20-2}
\end{minipage}\end{center}
}

\newcommand{\EyePanel}{%
\begin{center}\begin{minipage}{\linewidth}\centering
\begin{tikzpicture}[>=Latex,font=\small,every node/.style={fill=none}]
\draw[thick] (0,0) ellipse (2.4 and 1.3);
\draw[thick] (-1.3,-.7) .. controls (-1.65,-.2) and (-1.65,.2) .. (-1.3,.7) .. controls (-.95,.2) and (-.95,-.2) .. cycle;
\draw[DeepBlue,thick,->] (-3.6000000,-0.5500000)--(-1.3000000,-0.5500000);
\draw[DeepBlue,thick,->] (-1.3000000,-0.5500000)--(2.4000000,0.0000000);
\draw[DeepBlue,thick,->] (-3.6000000,0.5500000)--(-1.3000000,0.5500000);
\draw[DeepBlue,thick,->] (-1.3000000,0.5500000)--(2.4000000,0.0000000);
\draw[gray,dashed,thin] (-1.3000000,-1.0000000)--(-1.3000000,1.0000000);
\node[above,fill=none,inner sep=2pt] at (-2.3000000,1.1500000) {角膜};
\node[above,fill=none,inner sep=2pt] at (-0.8000000,1.5000000) {晶状体};
\draw[gray,dashed,thin] (-0.8000000,1.4000000)--(-1.1000000,0.4500000);
\fill (2.4000000,0.0000000) circle (1.2pt);
\node[right,fill=none,inner sep=2pt] at (2.4000000,0.0000000) {视网膜};
\node[below,fill=none,inner sep=2pt] at (0.0000000,-1.5500000) {虚线平面表示眼的等效折射系统};
\end{tikzpicture}
\figcap{远物平行光在正常眼视网膜同一点会聚。光线按等效薄系统画，解剖轮廓用于辨认结构。}\label{ch28:fig:s43-20-3}
\end{minipage}\end{center}
}

\newcommand{\ApparentDepthDiagram}{%
\begin{center}\begin{minipage}{\linewidth}\centering
\begin{tikzpicture}[>=Latex,font=\small,every node/.style={fill=none}]
\begin{scope}[shift={(-3.0000000,0.0000000)}]
\draw[thick] (-2.0000000,0.0000000)--(2.0000000,0.0000000);
\fill (0.0000000,-1.8000000) circle (1.2pt);
\node[below,fill=none,inner sep=2pt] at (0.0000000,-1.8000000) {};
\draw[gray,dashed,thin] (0.0000000,-1.8000000)--(-0.7000000,-1.8000000);
\node[left,fill=none,inner sep=2pt] at (-0.7500000,-1.8000000) {$O$};
\draw[gray,dashed,thin] (0.0000000,-2.0000000)--(0.0000000,1.6000000);
\draw[DeepBlue,thick,->] (0.0000000,-1.8000000)--(-0.6000000,0.0000000);
\draw[DeepBlue,thick,->] (-0.6000000,0.0000000)--(-1.1902918,1.2694706);
\draw[gray,dashed,thin] (-0.6000000,0.0000000)--(0.0000000,-1.2903488);
\draw[DeepBlue,thick,->] (0.0000000,-1.8000000)--(0.6000000,0.0000000);
\draw[DeepBlue,thick,->] (0.6000000,0.0000000)--(1.1902918,1.2694706);
\draw[gray,dashed,thin] (0.6000000,0.0000000)--(0.0000000,-1.2903488);
\fill (0.0000000,-1.2903488) circle (1.2pt);
\node[below,fill=none,inner sep=2pt] at (0.0000000,-1.2903488) {};
\draw[gray,dashed,thin] (0.0000000,-1.2903488)--(0.7000000,-1.2903488);
\node[right,fill=none,inner sep=2pt] at (0.7500000,-1.2903488) {$I$};
\node[left,fill=none,inner sep=2pt] at (-1.3000000,-1.0000000) {水};
\node[left,fill=none,inner sep=2pt] at (-1.3000000,0.6000000) {空气中的眼};
\end{scope}
\begin{scope}[shift={(3.0000000,0.0000000)}]
\draw[thick] (-2.0000000,0.0000000)--(2.0000000,0.0000000);
\fill (0.0000000,-1.8000000) circle (1.2pt);
\node[below,fill=none,inner sep=2pt] at (0.0000000,-1.8000000) {};
\draw[gray,dashed,thin] (0.0000000,-1.8000000)--(-0.7000000,-1.8000000);
\node[left,fill=none,inner sep=2pt] at (-0.7500000,-1.8000000) {$O$};
\draw[gray,dashed,thin] (0.0000000,-2.0000000)--(0.0000000,1.6000000);
\draw[DeepBlue,thick,->] (0.0000000,-1.8000000)--(-0.6000000,0.0000000);
\draw[DeepBlue,thick,->] (-0.6000000,0.0000000)--(-0.9320392,1.3600551);
\draw[gray,dashed,thin] (-0.6000000,0.0000000)--(0.0000000,-2.4576411);
\draw[DeepBlue,thick,->] (0.0000000,-1.8000000)--(0.6000000,0.0000000);
\draw[DeepBlue,thick,->] (0.6000000,0.0000000)--(0.9320392,1.3600551);
\draw[gray,dashed,thin] (0.6000000,0.0000000)--(0.0000000,-2.4576411);
\fill (0.0000000,-2.4576411) circle (1.2pt);
\node[below,fill=none,inner sep=2pt] at (0.0000000,-2.4576411) {};
\draw[gray,dashed,thin] (0.0000000,-2.4576411)--(0.7000000,-2.4576411);
\node[right,fill=none,inner sep=2pt] at (0.7500000,-2.4576411) {$I$};
\node[left,fill=none,inner sep=2pt] at (-1.3000000,-1.0000000) {空气};
\node[left,fill=none,inner sep=2pt] at (-1.3000000,0.6000000) {水中的眼};
\end{scope}
\end{tikzpicture}
\figcap{从空气看水中（左）和从水中看空气（右）。虚像取所示两斜光线的反向交点；近轴极限为 $|v|=n_2u/n_1$。}\label{ch28:fig:s43-20-4}
\end{minipage}\end{center}
}

\newcommand{\ThinLensPanel}{%
\begin{center}\begin{minipage}{\linewidth}\centering
\begin{tikzpicture}[>=Latex,font=\small,every node/.style={fill=none}]
\begin{scope}[shift={(-5.0000000,0.0000000)}]
\draw[thick,fill=DeepBlue!5] (-.12,-1).. controls (-.65,-.3) and (-.65,.3)..(-.12,1)--(.12,1)..controls(.65,.3)and(.65,-.3)..(.12,-1)--cycle;
\node[below,fill=none,inner sep=2pt] at (0.0000000,-1.3000000) {双凸};
\end{scope}
\begin{scope}[shift={(-3.0000000,0.0000000)}]
\draw[thick,fill=DeepBlue!5] (0,-1)--(0,1)--(.12,1)..controls(.65,.3)and(.65,-.3)..(.12,-1)--cycle;
\node[below,fill=none,inner sep=2pt] at (0.0000000,-1.3000000) {平凸};
\end{scope}
\begin{scope}[shift={(-1.0000000,0.0000000)}]
\draw[thick,fill=DeepBlue!5] (-.1,-1)..controls(.3,-.3)and(.3,.3)..(-.1,1)--(.1,1)..controls(.9,.3)and(.9,-.3)..(.1,-1)--cycle;
\node[below,fill=none,inner sep=2pt] at (0.0000000,-1.3000000) {会聚弯月};
\end{scope}
\begin{scope}[shift={(1.0000000,0.0000000)}]
\draw[thick,fill=DeepBlue!5] (-.55,-1)..controls(-.05,-.3)and(-.05,.3)..(-.55,1)--(.55,1)..controls(.05,.3)and(.05,-.3)..(.55,-1)--cycle;
\node[below,fill=none,inner sep=2pt] at (0.0000000,-1.3000000) {双凹};
\end{scope}
\begin{scope}[shift={(3.0000000,0.0000000)}]
\draw[thick,fill=DeepBlue!5] (0,-1)--(0,1)--(.65,1)..controls(.1,.3)and(.1,-.3)..(.65,-1)--cycle;
\node[below,fill=none,inner sep=2pt] at (0.0000000,-1.3000000) {平凹};
\end{scope}
\begin{scope}[shift={(5.0000000,0.0000000)}]
\draw[thick,fill=DeepBlue!5] (-.45,-1)..controls(.45,-.3)and(.45,.3)..(-.45,1)--(.2,1)..controls(.6,.3)and(.6,-.3)..(.2,-1)--cycle;
\node[below,fill=none,inner sep=2pt] at (0.0000000,-1.3000000) {发散弯月};
\end{scope}
\begin{scope}[shift={(0.0000000,-4.0000000)}]
\draw[thick,fill=DeepBlue!5] (0.1669697,-0.8000000)--(0.1600696,-0.7840000)--(0.1533338,-0.7680000)--(0.1467607,-0.7520000)--(0.1403484,-0.7360000)--(0.1340954,-0.7200000)--(0.1280000,-0.7040000)--(0.1220607,-0.6880000)--(0.1162760,-0.6720000)--(0.1106446,-0.6560000)--(0.1051649,-0.6400000)--(0.0998358,-0.6240000)--(0.0946559,-0.6080000)--(0.0896241,-0.5920000)--(0.0847392,-0.5760000)--(0.0800000,-0.5600000)--(0.0754055,-0.5440000)--(0.0709546,-0.5280000)--(0.0666464,-0.5120000)--(0.0624799,-0.4960000)--(0.0584542,-0.4800000)--(0.0545684,-0.4640000)--(0.0508217,-0.4480000)--(0.0472133,-0.4320000)--(0.0437423,-0.4160000)--(0.0404082,-0.4000000)--(0.0372101,-0.3840000)--(0.0341475,-0.3680000)--(0.0312197,-0.3520000)--(0.0284260,-0.3360000)--(0.0257660,-0.3200000)--(0.0232390,-0.3040000)--(0.0208446,-0.2880000)--(0.0185823,-0.2720000)--(0.0164517,-0.2560000)--(0.0144522,-0.2400000)--(0.0125836,-0.2240000)--(0.0108454,-0.2080000)--(0.0092373,-0.1920000)--(0.0077591,-0.1760000)--(0.0064103,-0.1600000)--(0.0051907,-0.1440000)--(0.0041002,-0.1280000)--(0.0031385,-0.1120000)--(0.0023053,-0.0960000)--(0.0016006,-0.0800000)--(0.0010243,-0.0640000)--(0.0005761,-0.0480000)--(0.0002560,-0.0320000)--(0.0000640,-0.0160000)--(0.0000000,0.0000000)--(0.0000640,0.0160000)--(0.0002560,0.0320000)--(0.0005761,0.0480000)--(0.0010243,0.0640000)--(0.0016006,0.0800000)--(0.0023053,0.0960000)--(0.0031385,0.1120000)--(0.0041002,0.1280000)--(0.0051907,0.1440000)--(0.0064103,0.1600000)--(0.0077591,0.1760000)--(0.0092373,0.1920000)--(0.0108454,0.2080000)--(0.0125836,0.2240000)--(0.0144522,0.2400000)--(0.0164517,0.2560000)--(0.0185823,0.2720000)--(0.0208446,0.2880000)--(0.0232390,0.3040000)--(0.0257660,0.3200000)--(0.0284260,0.3360000)--(0.0312197,0.3520000)--(0.0341475,0.3680000)--(0.0372101,0.3840000)--(0.0404082,0.4000000)--(0.0437423,0.4160000)--(0.0472133,0.4320000)--(0.0508217,0.4480000)--(0.0545684,0.4640000)--(0.0584542,0.4800000)--(0.0624799,0.4960000)--(0.0666464,0.5120000)--(0.0709546,0.5280000)--(0.0754055,0.5440000)--(0.0800000,0.5600000)--(0.0847392,0.5760000)--(0.0896241,0.5920000)--(0.0946559,0.6080000)--(0.0998358,0.6240000)--(0.1051649,0.6400000)--(0.1106446,0.6560000)--(0.1162760,0.6720000)--(0.1220607,0.6880000)--(0.1280000,0.7040000)--(0.1340954,0.7200000)--(0.1403484,0.7360000)--(0.1467607,0.7520000)--(0.1533338,0.7680000)--(0.1600696,0.7840000)--(0.1669697,0.8000000)--(0.4330303,0.8000000)--(0.4399304,0.7840000)--(0.4466662,0.7680000)--(0.4532393,0.7520000)--(0.4596516,0.7360000)--(0.4659046,0.7200000)--(0.4720000,0.7040000)--(0.4779393,0.6880000)--(0.4837240,0.6720000)--(0.4893554,0.6560000)--(0.4948351,0.6400000)--(0.5001642,0.6240000)--(0.5053441,0.6080000)--(0.5103759,0.5920000)--(0.5152608,0.5760000)--(0.5200000,0.5600000)--(0.5245945,0.5440000)--(0.5290454,0.5280000)--(0.5333536,0.5120000)--(0.5375201,0.4960000)--(0.5415458,0.4800000)--(0.5454316,0.4640000)--(0.5491783,0.4480000)--(0.5527867,0.4320000)--(0.5562577,0.4160000)--(0.5595918,0.4000000)--(0.5627899,0.3840000)--(0.5658525,0.3680000)--(0.5687803,0.3520000)--(0.5715740,0.3360000)--(0.5742340,0.3200000)--(0.5767610,0.3040000)--(0.5791554,0.2880000)--(0.5814177,0.2720000)--(0.5835483,0.2560000)--(0.5855478,0.2400000)--(0.5874164,0.2240000)--(0.5891546,0.2080000)--(0.5907627,0.1920000)--(0.5922409,0.1760000)--(0.5935897,0.1600000)--(0.5948093,0.1440000)--(0.5958998,0.1280000)--(0.5968615,0.1120000)--(0.5976947,0.0960000)--(0.5983994,0.0800000)--(0.5989757,0.0640000)--(0.5994239,0.0480000)--(0.5997440,0.0320000)--(0.5999360,0.0160000)--(0.6000000,0.0000000)--(0.5999360,-0.0160000)--(0.5997440,-0.0320000)--(0.5994239,-0.0480000)--(0.5989757,-0.0640000)--(0.5983994,-0.0800000)--(0.5976947,-0.0960000)--(0.5968615,-0.1120000)--(0.5958998,-0.1280000)--(0.5948093,-0.1440000)--(0.5935897,-0.1600000)--(0.5922409,-0.1760000)--(0.5907627,-0.1920000)--(0.5891546,-0.2080000)--(0.5874164,-0.2240000)--(0.5855478,-0.2400000)--(0.5835483,-0.2560000)--(0.5814177,-0.2720000)--(0.5791554,-0.2880000)--(0.5767610,-0.3040000)--(0.5742340,-0.3200000)--(0.5715740,-0.3360000)--(0.5687803,-0.3520000)--(0.5658525,-0.3680000)--(0.5627899,-0.3840000)--(0.5595918,-0.4000000)--(0.5562577,-0.4160000)--(0.5527867,-0.4320000)--(0.5491783,-0.4480000)--(0.5454316,-0.4640000)--(0.5415458,-0.4800000)--(0.5375201,-0.4960000)--(0.5333536,-0.5120000)--(0.5290454,-0.5280000)--(0.5245945,-0.5440000)--(0.5200000,-0.5600000)--(0.5152608,-0.5760000)--(0.5103759,-0.5920000)--(0.5053441,-0.6080000)--(0.5001642,-0.6240000)--(0.4948351,-0.6400000)--(0.4893554,-0.6560000)--(0.4837240,-0.6720000)--(0.4779393,-0.6880000)--(0.4720000,-0.7040000)--(0.4659046,-0.7200000)--(0.4596516,-0.7360000)--(0.4532393,-0.7520000)--(0.4466662,-0.7680000)--(0.4399304,-0.7840000)--(0.4330303,-0.8000000)--(0.1669697,-0.8000000);
\draw[gray,thin,->] (-4.4000000,0.0000000)--(4.2666120,0.0000000);
\fill (-4.0000000,0.0000000) circle (1.2pt);
\node[above,fill=none,inner sep=2pt] at (-4.0000000,0.0000000) {$O$};
\fill (3.8166120,0.0000000) circle (1.2pt);
\node[above,fill=none,inner sep=2pt] at (3.8166120,0.0000000) {$I$};
\fill (0.0308631,0.3500000) circle (1.2pt);
\node[above left,fill=none,inner sep=2pt] at (0.0308631,0.3500000) {$P$};
\fill (0.5697641,0.3464526) circle (1.2pt);
\node[above right,fill=none,inner sep=2pt] at (0.5697641,0.3464526) {$Q$};
\draw[DeepBlue,thick,->] (-4.0000000,0.0000000)--(0.0308631,0.3500000);
\draw[DeepBlue,thick,->] (0.0308631,0.3500000)--(0.5697641,0.3464526);
\draw[DeepBlue,thick,->] (0.5697641,0.3464526)--(3.8166120,0.0000000);
\draw[gray,dashed,thin] (2.0000000,0.0000000)--(0.0308631,0.3500000);
\draw[gray,dashed,thin] (-1.4000000,0.0000000)--(0.5697641,0.3464526);
\fill (2.0000000,0.0000000) circle (1.2pt);
\node[below,fill=none,inner sep=2pt] at (2.0000000,0.0000000) {};
\fill (-1.4000000,0.0000000) circle (1.2pt);
\node[below,fill=none,inner sep=2pt] at (-1.4000000,0.0000000) {};
\draw[gray,dashed,thin] (2.0000000,0.0000000)--(2.0000000,-0.9000000);
\draw[gray,dashed,thin] (-1.4000000,0.0000000)--(-1.4000000,-0.9000000);
\node[below,fill=none,inner sep=2pt] at (2.0000000,-0.9500000) {$C_1$};
\node[below,fill=none,inner sep=2pt] at (-1.4000000,-0.9500000) {$C_2$};
\node[above,fill=none,inner sep=2pt] at (-1.2000000,1.0000000) {$C_1P=R_1$};
\node[above,fill=none,inner sep=2pt] at (2.2000000,1.0000000) {$C_2Q=|R_2|$};
\draw[gray,<->] (-4.0000000,-1.5000000)--(0.0000000,-1.5000000);
\node[below,fill=none,inner sep=2pt] at (-2.0000000,-1.5000000) {$u$};
\draw[gray,<->] (0.6000000,-1.5000000)--(3.8166120,-1.5000000);
\node[below,fill=none,inner sep=2pt] at (2.2083060,-1.5000000) {$v$};
\end{scope}
\end{tikzpicture}
\figcap{上为六种闭合透镜截面，置于空气中时左三种会聚、右三种发散；下为两球面分别折射的光路，$C_1P,C_2Q$ 是两面的半径和法线。}\label{ch28:fig:s43-20-5}
\end{minipage}\end{center}
}

\newcommand{\LensImagingDiagram}{%
\begin{center}\begin{minipage}{\linewidth}\centering
\begin{tikzpicture}[>=Latex,font=\small,every node/.style={fill=none}]
\draw[gray,thin,->] (-4.8000000,0.0000000)--(4.0000000,0.0000000);
\draw[DeepBlue,thick,<->] (0.0000000,-1.8500000)--(0.0000000,1.8500000);
\draw[black,thick,->] (-4.2000000,0.0000000)--(-4.2000000,1.5000000);
\node[above,fill=none,inner sep=2pt] at (-4.2000000,1.5000000) {物};
\draw[black,thick,->] (3.3123028,0.0000000)--(3.3123028,-1.1829653);
\node[below,fill=none,inner sep=2pt] at (3.3123028,-1.1829653) {像};
\fill (-3.7037037,0.0000000) circle (1.2pt);
\node[below,fill=none,inner sep=2pt] at (-3.7037037,0.0000000) {};
\node[below,fill=none,inner sep=2pt] at (-3.7037037,-0.5000000) {$2F$};
\fill (-1.8518519,0.0000000) circle (1.2pt);
\node[below,fill=none,inner sep=2pt] at (-1.8518519,0.0000000) {};
\node[below,fill=none,inner sep=2pt] at (-1.8518519,-0.5000000) {$F$};
\fill (1.8518519,0.0000000) circle (1.2pt);
\node[below,fill=none,inner sep=2pt] at (1.8518519,0.0000000) {};
\node[above,fill=none,inner sep=2pt] at (1.8518519,0.2000000) {$F'$};
\fill (3.7037037,0.0000000) circle (1.2pt);
\node[below,fill=none,inner sep=2pt] at (3.7037037,0.0000000) {};
\node[above,fill=none,inner sep=2pt] at (3.7037037,0.2000000) {$2F'$};
\draw[DeepBlue,thick,->] (-4.2000000,1.5000000)--(0.0000000,1.5000000);
\draw[DeepBlue,thick,->] (0.0000000,1.5000000)--(3.3123028,-1.1829653);
\draw[DeepBlue,thick,->] (-4.2000000,1.5000000)--(0.0000000,0.0000000);
\draw[DeepBlue,thick,->] (0.0000000,0.0000000)--(3.3123028,-1.1829653);
\draw[DeepBlue,thick,->] (-4.2000000,1.5000000)--(0.0000000,-1.1829653);
\draw[DeepBlue,thick,->] (0.0000000,-1.1829653)--(3.3123028,-1.1829653);
\end{tikzpicture}
\figcap{$u=70.0\,\mathrm{cm}$、$f\approx30.9\,\mathrm{cm}$、$v\approx55.2\,\mathrm{cm}$ 时的三条近轴特征光线；通过物方焦点的那条出射后平行光轴。}\label{ch28:fig:s43-20-6}
\end{minipage}\end{center}
}

\newcommand{\MagnifierDiagram}{%
\begin{center}\begin{minipage}{\linewidth}\centering
\begin{tikzpicture}[>=Latex,font=\small,every node/.style={fill=none}]
\begin{scope}[shift={(-2.6000000,0.0000000)}]
\draw[gray,thin,->] (-3.6000000,0.0000000)--(1.7000000,0.0000000);
\draw[DeepBlue,thick,<->] (0.0000000,-1.2000000)--(0.0000000,1.2000000);
\fill (-1.5000000,0.0000000) circle (1.2pt);
\node[below,fill=none,inner sep=2pt] at (-1.5000000,0.0000000) {$F$};
\draw[black,thick,->] (-1.0000000,0.0000000)--(-1.0000000,0.6500000);
\node[left,fill=none,inner sep=2pt] at (-1.0000000,0.6500000) {物};
\draw[black,thick,->] (-3.0000000,0.0000000)--(-3.0000000,1.9500000);
\node[above,fill=none,inner sep=2pt] at (-3.0000000,1.9500000) {虚像};
\draw[DeepBlue,thick,->] (-1.0000000,0.6500000)--(0.0000000,0.6500000);
\draw[DeepBlue,thick,->] (0.0000000,0.6500000)--(1.2000000,0.1300000);
\draw[gray,dashed,thin] (-3.0000000,1.9500000)--(0.0000000,0.6500000);
\draw[DeepBlue,thick,->] (-1.0000000,0.6500000)--(0.0000000,0.0000000);
\draw[DeepBlue,thick,->] (0.0000000,0.0000000)--(1.2000000,-0.7800000);
\draw[gray,dashed,thin] (-3.0000000,1.9500000)--(0.0000000,0.0000000);
\draw[Gold,very thick] (0.4000000,-0.3400000)--(0.4000000,0.5000000);
\node[right,fill=none,inner sep=2pt] at (0.5000000,1.2500000) {眼瞳位置};
\draw[gray,dashed,thin] (0.8000000,1.2000000)--(0.4000000,0.5000000);
\end{scope}
\begin{scope}[shift={(3.8000000,0.7000000)}]
\draw[black,thick,->] (-1.6000000,0.0000000)--(-1.6000000,0.6500000);
\draw[DeepBlue,thick,->] (-1.6000000,0.6500000)--(1.4000000,0.0000000);
\draw[gray,dashed,thin] (-1.6000000,0.0000000)--(1.4000000,0.0000000);
\node[right,fill=none,inner sep=2pt] at (1.4000000,0.0000000) {裸眼};
\draw[gray,<->] (-1.6000000,-0.4500000)--(1.4000000,-0.4500000);
\node[below,fill=none,inner sep=2pt] at (-0.1000000,-0.4500000) {$L_0$};
\node[above,fill=none,inner sep=2pt] at (0.0000000,1.2000000) {$\theta_0\simeq H/L_0$};
\end{scope}
\end{tikzpicture}
\figcap{左为放大镜的正立虚像及进入眼瞳的光；右为裸眼在明视距离观察同样大小物体的基准视角。}\label{ch28:fig:s43-20-7}
\end{minipage}\end{center}
}

\newcommand{\SphericalMirrorDiagram}{%
\begin{center}\begin{minipage}{\linewidth}\centering
\begin{tikzpicture}[>=Latex,font=\small,every node/.style={fill=none}]
\draw[thick] plot coordinates {(2.6488428,-1.4084147) (2.6624817,-1.3824585) (2.6758662,-1.3563702) (2.6889951,-1.3301523) (2.7018672,-1.3038074) (2.7144812,-1.2773379) (2.7268358,-1.2507464) (2.7389300,-1.2240354) (2.7507625,-1.1972075) (2.7623322,-1.1702652) (2.7736381,-1.1432111) (2.7846790,-1.1160479) (2.7954539,-1.0887780) (2.8059618,-1.0614041) (2.8162016,-1.0339288) (2.8261724,-1.0063547) (2.8358732,-0.9786845) (2.8453031,-0.9509208) (2.8544612,-0.9230663) (2.8633466,-0.8951236) (2.8719585,-0.8670954) (2.8802960,-0.8389843) (2.8883584,-0.8107932) (2.8961449,-0.7825245) (2.9036547,-0.7541811) (2.9108872,-0.7257657) (2.9178416,-0.6972809) (2.9245172,-0.6687295) (2.9309135,-0.6401143) (2.9370297,-0.6114379) (2.9428655,-0.5827031) (2.9484201,-0.5539126) (2.9536930,-0.5250692) (2.9586838,-0.4961756) (2.9633919,-0.4672347) (2.9678170,-0.4382491) (2.9719586,-0.4092216) (2.9758162,-0.3801551) (2.9793896,-0.3510522) (2.9826784,-0.3219158) (2.9856822,-0.2927487) (2.9884008,-0.2635536) (2.9908340,-0.2343333) (2.9929814,-0.2050906) (2.9948430,-0.1758284) (2.9964184,-0.1465493) (2.9977076,-0.1172563) (2.9987105,-0.0879520) (2.9994269,-0.0586393) (2.9998567,-0.0293211) (3.0000000,0.0000000) (2.9998567,0.0293211) (2.9994269,0.0586393) (2.9987105,0.0879520) (2.9977076,0.1172563) (2.9964184,0.1465493) (2.9948430,0.1758284) (2.9929814,0.2050906) (2.9908340,0.2343333) (2.9884008,0.2635536) (2.9856822,0.2927487) (2.9826784,0.3219158) (2.9793896,0.3510522) (2.9758162,0.3801551) (2.9719586,0.4092216) (2.9678170,0.4382491) (2.9633919,0.4672347) (2.9586838,0.4961756) (2.9536930,0.5250692) (2.9484201,0.5539126) (2.9428655,0.5827031) (2.9370297,0.6114379) (2.9309135,0.6401143) (2.9245172,0.6687295) (2.9178416,0.6972809) (2.9108872,0.7257657) (2.9036547,0.7541811) (2.8961449,0.7825245) (2.8883584,0.8107932) (2.8802960,0.8389843) (2.8719585,0.8670954) (2.8633466,0.8951236) (2.8544612,0.9230663) (2.8453031,0.9509208) (2.8358732,0.9786845) (2.8261724,1.0063547) (2.8162016,1.0339288) (2.8059618,1.0614041) (2.7954539,1.0887780) (2.7846790,1.1160479) (2.7736381,1.1432111) (2.7623322,1.1702652) (2.7507625,1.1972075) (2.7389300,1.2240354) (2.7268358,1.2507464) (2.7144812,1.2773379) (2.7018672,1.3038074) (2.6889951,1.3301523) (2.6758662,1.3563702) (2.6624817,1.3824585) (2.6488428,1.4084147)};
\draw[gray,thin,->] (-3.6000000,0.0000000)--(3.6000000,0.0000000);
\fill (-3.0000000,0.0000000) circle (1.2pt);
\node[below,fill=none,inner sep=2pt] at (-3.0000000,0.0000000) {$O$};
\fill (0.0000000,0.0000000) circle (1.2pt);
\node[below,fill=none,inner sep=2pt] at (0.0000000,0.0000000) {$C$};
\fill (3.0000000,0.0000000) circle (1.2pt);
\node[below right,fill=none,inner sep=2pt] at (3.0000000,0.0000000) {$V$};
\fill (2.9108872,0.7257657) circle (1.2pt);
\node[below,fill=none,inner sep=2pt] at (2.9108872,0.7257657) {};
\draw[gray,dashed,thin] (2.9108872,0.7257657)--(2.3608872,1.1757657);
\node[above left,fill=none,inner sep=2pt] at (2.3108872,1.2257657) {$P$};
\fill (1.0202029,0.0000000) circle (1.2pt);
\node[below,fill=none,inner sep=2pt] at (1.0202029,0.0000000) {$I$};
\draw[gray,dashed,thin] (0.0000000,0.0000000)--(2.9108872,0.7257657);
\draw[DeepBlue,thick,->] (-3.0000000,0.0000000)--(2.9108872,0.7257657);
\draw[DeepBlue,thick,->] (2.9108872,0.7257657)--(1.0202029,0.0000000);
\end{tikzpicture}
\figcap{球面镜真实反射光路，$CP$ 是法线。所示有限高度轴交点趋向近轴像点，但不与所有孔径光线严格重合。}\label{ch28:fig:s43-20-8}
\end{minipage}\end{center}
}

\newcommand{\AplanaticLensPanel}{%
% Independent panel: original physical paths and local coordinates retained.
\begin{center}\begin{minipage}{\linewidth}\centering
\begin{tikzpicture}[>=Latex,font=\small,every node/.style={fill=none}]
\draw[thick] (1.2122511,-0.9000000)--(1.2200945,-0.8775000)--(1.2277163,-0.8550000)--(1.2351183,-0.8325000)--(1.2423025,-0.8100000)--(1.2492707,-0.7875000)--(1.2560247,-0.7650000)--(1.2625661,-0.7425000)--(1.2688965,-0.7200000)--(1.2750174,-0.6975000)--(1.2809304,-0.6750000)--(1.2866368,-0.6525000)--(1.2921381,-0.6300000)--(1.2974353,-0.6075000)--(1.3025299,-0.5850000)--(1.3074230,-0.5625000)--(1.3121156,-0.5400000)--(1.3166089,-0.5175000)--(1.3209038,-0.4950000)--(1.3250014,-0.4725000)--(1.3289026,-0.4500000)--(1.3326081,-0.4275000)--(1.3361189,-0.4050000)--(1.3394356,-0.3825000)--(1.3425590,-0.3600000)--(1.3454899,-0.3375000)--(1.3482287,-0.3150000)--(1.3507762,-0.2925000)--(1.3531327,-0.2700000)--(1.3552990,-0.2475000)--(1.3572753,-0.2250000)--(1.3590622,-0.2025000)--(1.3606600,-0.1800000)--(1.3620690,-0.1575000)--(1.3632896,-0.1350000)--(1.3643219,-0.1125000)--(1.3651662,-0.0900000)--(1.3658228,-0.0675000)--(1.3662916,-0.0450000)--(1.3665729,-0.0225000)--(1.3666667,0.0000000)--(1.3665729,0.0225000)--(1.3662916,0.0450000)--(1.3658228,0.0675000)--(1.3651662,0.0900000)--(1.3643219,0.1125000)--(1.3632896,0.1350000)--(1.3620690,0.1575000)--(1.3606600,0.1800000)--(1.3590622,0.2025000)--(1.3572753,0.2250000)--(1.3552990,0.2475000)--(1.3531327,0.2700000)--(1.3507762,0.2925000)--(1.3482287,0.3150000)--(1.3454899,0.3375000)--(1.3425590,0.3600000)--(1.3394356,0.3825000)--(1.3361189,0.4050000)--(1.3326081,0.4275000)--(1.3289026,0.4500000)--(1.3250014,0.4725000)--(1.3209038,0.4950000)--(1.3166089,0.5175000)--(1.3121156,0.5400000)--(1.3074230,0.5625000)--(1.3025299,0.5850000)--(1.2974353,0.6075000)--(1.2921381,0.6300000)--(1.2866368,0.6525000)--(1.2809304,0.6750000)--(1.2750174,0.6975000)--(1.2688965,0.7200000)--(1.2625661,0.7425000)--(1.2560247,0.7650000)--(1.2492707,0.7875000)--(1.2423025,0.8100000)--(1.2351183,0.8325000)--(1.2277163,0.8550000)--(1.2200945,0.8775000)--(1.2122511,0.9000000)--(1.7860571,0.9000000)--(1.7972183,0.8775000)--(1.8080307,0.8550000)--(1.8185004,0.8325000)--(1.8286334,0.8100000)--(1.8384351,0.7875000)--(1.8479110,0.7650000)--(1.8570659,0.7425000)--(1.8659046,0.7200000)--(1.8744316,0.6975000)--(1.8826511,0.6750000)--(1.8905670,0.6525000)--(1.8981833,0.6300000)--(1.9055035,0.6075000)--(1.9125310,0.5850000)--(1.9192691,0.5625000)--(1.9257206,0.5400000)--(1.9318886,0.5175000)--(1.9377758,0.4950000)--(1.9433846,0.4725000)--(1.9487175,0.4500000)--(1.9537768,0.4275000)--(1.9585645,0.4050000)--(1.9630827,0.3825000)--(1.9673332,0.3600000)--(1.9713178,0.3375000)--(1.9750380,0.3150000)--(1.9784953,0.2925000)--(1.9816912,0.2700000)--(1.9846269,0.2475000)--(1.9873034,0.2250000)--(1.9897220,0.2025000)--(1.9918835,0.1800000)--(1.9937888,0.1575000)--(1.9954385,0.1350000)--(1.9968334,0.1125000)--(1.9979740,0.0900000)--(1.9988606,0.0675000)--(1.9994937,0.0450000)--(1.9998734,0.0225000)--(2.0000000,0.0000000)--(1.9998734,-0.0225000)--(1.9994937,-0.0450000)--(1.9988606,-0.0675000)--(1.9979740,-0.0900000)--(1.9968334,-0.1125000)--(1.9954385,-0.1350000)--(1.9937888,-0.1575000)--(1.9918835,-0.1800000)--(1.9897220,-0.2025000)--(1.9873034,-0.2250000)--(1.9846269,-0.2475000)--(1.9816912,-0.2700000)--(1.9784953,-0.2925000)--(1.9750380,-0.3150000)--(1.9713178,-0.3375000)--(1.9673332,-0.3600000)--(1.9630827,-0.3825000)--(1.9585645,-0.4050000)--(1.9537768,-0.4275000)--(1.9487175,-0.4500000)--(1.9433846,-0.4725000)--(1.9377758,-0.4950000)--(1.9318886,-0.5175000)--(1.9257206,-0.5400000)--(1.9192691,-0.5625000)--(1.9125310,-0.5850000)--(1.9055035,-0.6075000)--(1.8981833,-0.6300000)--(1.8905670,-0.6525000)--(1.8826511,-0.6750000)--(1.8744316,-0.6975000)--(1.8659046,-0.7200000)--(1.8570659,-0.7425000)--(1.8479110,-0.7650000)--(1.8384351,-0.7875000)--(1.8286334,-0.8100000)--(1.8185004,-0.8325000)--(1.8080307,-0.8550000)--(1.7972183,-0.8775000)--(1.7860571,-0.9000000)--(1.2122511,-0.9000000);
\draw[gray,thin,->] (-3.4000000,0.0000000)--(3.5000000,0.0000000);
\fill (-3.0000000,0.0000000) circle (1.2pt);
\node[below,fill=none,inner sep=2pt] at (-3.0000000,0.0000000) {$Q$};
\fill (-1.3333333,0.0000000) circle (1.2pt);
\draw[gray,thin](-1.3333333,.04)--(-1.3333333,.52);\node[above,fill=none,inner sep=2pt] at (-1.3333333,.56) {$C_1$};
\fill (0.0000000,0.0000000) circle (1.2pt);
\node[below,fill=none,inner sep=2pt] at (0.0000000,0.0000000) {};
\draw[gray,dashed,thin] (0.0000000,0.0000000)--(0.0000000,-0.6500000);
\node[below,fill=none,inner sep=2pt] at (0.0000000,-0.7000000) {$C_2$};
\draw[DeepBlue,thick,->] (-1.3333333,0.0000000)--(1.2948880,-0.6184277);
\draw[DeepBlue,thick,->] (1.2948880,-0.6184277)--(1.8540496,-0.7500000);
\draw[DeepBlue,thick,->] (1.8540496,-0.7500000)--(2.9411498,-0.9179680);
\draw[gray,dashed,thin] (-3.0000000,0.0000000)--(1.8540496,-0.7500000);
\draw[gray,dashed,thin] (0.0000000,0.0000000)--(1.8540496,-0.7500000);
\draw[DeepBlue,thick,->] (-1.3333333,0.0000000)--(1.2948880,0.6184277);
\draw[DeepBlue,thick,->] (1.2948880,0.6184277)--(1.8540496,0.7500000);
\draw[DeepBlue,thick,->] (1.8540496,0.7500000)--(2.9411498,0.9179680);
\draw[gray,dashed,thin] (-3.0000000,0.0000000)--(1.8540496,0.7500000);
\draw[gray,dashed,thin] (0.0000000,0.0000000)--(1.8540496,0.7500000);
\node[left,fill=none,inner sep=2pt] at (1.4000000,1.0500000) {$S_1$};
\node[right,fill=none,inner sep=2pt] at (2.1000000,1.0500000) {$S_2$};
\fill(1.8540496,.75)circle(1.2pt);\node[right]at(2.65,.4){$P$};\draw[thin](1.8540496,.75)--(2.55,.35);
\node[align=center]at(-.6,-1.6){$R_1$（左），$R_2$（右）\\玻璃折射率 $n$};
\end{tikzpicture}
\figcap{无球差透镜的 $Q$ 为虚像；虚线是出射光反向延长线，$C_1,C_2$ 是两球面曲率中心。}\label{ch28:fig:s43-20-9}
\end{minipage}\end{center}
}

\newcommand{\CatEyePanel}{%
% Independent panel: original physical paths and local coordinates retained.
\begin{center}\begin{minipage}{\linewidth}\centering
\begin{tikzpicture}[>=Latex,font=\small,every node/.style={fill=none}]
\draw[pattern=north east lines] (-1.8,.9) rectangle (1.8,1.6);\draw[pattern=north east lines] (-1.8,-1.6) rectangle (1.8,-.9);
\begin{scope}[shift={(-1.0000000,0.0000000)}]
\draw[thick,fill=DeepBlue!5] (-.55,-1)..controls(-.05,-.3)and(-.05,.3)..(-.55,1)--(.55,1)..controls(.05,.3)and(.05,-.3)..(.55,-1)--cycle;
\end{scope}
\begin{scope}[shift={(1.1000000,0.0000000)}]
\draw[thick,fill=DeepBlue!5] (-.12,-1).. controls (-.65,-.3) and (-.65,.3)..(-.12,1)--(.12,1)..controls(.65,.3)and(.65,-.3)..(.12,-1)--cycle;
\end{scope}
\node[below,fill=none,inner sep=2pt] at (-1.0000000,-2.5500000) {$f_1=-1\,\mathrm{cm}$};
\node[below,fill=none,inner sep=2pt] at (1.1000000,-2.5500000) {$f_2=3.5\,\mathrm{cm}$};
\node[left,fill=none,inner sep=2pt] at (-2.1000000,0.0000000) {室外};
\node[right,fill=none,inner sep=2pt] at (2.1000000,0.0000000) {室内};
\draw[gray,<->] (-1.0000000,-1.9500000)--(1.1000000,-1.9500000);
\node[below,fill=none,inner sep=2pt] at (0.0500000,-1.9500000) {$2.1\,\mathrm{cm}$};
\end{tikzpicture}
\figcap{猫眼的两透镜装在同一贯通门孔中；正向观察时光从室外凹透镜传向室内凸透镜。}\label{ch28:fig:s43-20-10}
\end{minipage}\end{center}
}

\newcommand{\SilveredLensPanel}{%
% Independent panel: original physical paths and local coordinates retained.
\begin{center}\begin{minipage}{\linewidth}\centering
\begin{tikzpicture}[>=Latex,font=\small,every node/.style={fill=none}]
\draw[gray,thin,->] (-3.5000000,0.0000000)--(0.5000000,0.0000000);
\draw[very thick,double] (0.0000000,-1.4000000)--(0.0000000,1.4000000);
\draw[black,thick,->] (-3.0000000,0.0000000)--(-3.0000000,0.9000000);
\draw[black,thick,->] (-1.0000000,0.0000000)--(-1.0000000,-0.3000000);
\draw[DeepBlue,thick,->] (-3.0000000,0.9000000)--(0.0000000,0.9000000);
\draw[DeepBlue,thick,->] (0.0000000,0.9000000)--(-1.0000000,-0.3000000);
\draw[DeepBlue,thick,->] (-3.0000000,0.9000000)--(0.0000000,0.0000000);
\draw[DeepBlue,thick,->] (0.0000000,0.0000000)--(-1.0000000,-0.3000000);
\node[above,fill=none,inner sep=2pt] at (-3.0000000,0.9000000) {$H$};
\node[below,fill=none,inner sep=2pt] at (-1.0000000,-0.5000000) {$-H/3$};
\node[above,fill=none,inner sep=2pt] at (0.0000000,1.5000000) {薄透镜与银面重合的等效平面};
\node[below,fill=none,inner sep=2pt] at (-2.0000000,-1.2000000) {$u=2f,\ v=2f/3$};
\end{tikzpicture}
\figcap{镀银薄透镜的等效折返光路，实像倒立，像高为物高的三分之一。}\label{ch28:fig:s43-20-11}
\end{minipage}\end{center}
}

\newcommand{\PrismTrainPanel}{%
% Independent panel: original physical paths and local coordinates retained.
\begin{center}\begin{minipage}{\linewidth}\centering
\begin{tikzpicture}[>=Latex,font=\small,every node/.style={fill=none}]
\draw[thick] (0.0000000,1.2000000)--(1.2000000,1.2000000)--(1.2000000,0.0000000)--(0.0000000,1.2000000);
\draw[DeepBlue,thick,->] (0.6000000,2.4000000)--(0.6000000,0.6000000);
\draw[DeepBlue,thick,->] (0.6000000,0.6000000)--(4.8000000,0.6000000);
\draw[black,thick,->] (0.6000000,2.4000000)--(0.9000000,2.4000000);
\node[above left,fill=none,inner sep=2pt] at (0.6000000,2.4000000) {$B$};
\node[above right,fill=none,inner sep=2pt] at (0.9000000,2.4000000) {$A$};
\begin{scope}[shift={(0.0000000,0.6000000)}]
\draw[DeepBlue,thick,<->] (2.4000000,-1.4000000)--(2.4000000,1.4000000);
\end{scope}
\draw[DeepBlue,thick] (4.2000000,-0.8000000)--(4.2000000,2.0000000);
\draw[DeepBlue,thick,->] (4.2000000,2.0000000)--(4.2000000,1.7000000);
\draw[DeepBlue,thick,->] (4.2000000,-0.8000000)--(4.2000000,-0.5000000);
\node[above,fill=none,inner sep=2pt] at (2.4000000,2.1500000) {$f_1=10$};
\node[above,fill=none,inner sep=2pt] at (4.2000000,2.1500000) {$f_2=-10$};
\draw[gray,<->] (-0.2500000,1.2000000)--(-0.2500000,2.4000000);
\node[left,fill=none,inner sep=2pt] at (-0.2500000,1.8000000) {$6$};
\draw[gray,<->] (-0.2500000,0.0000000)--(-0.2500000,1.2000000);
\node[left,fill=none,inner sep=2pt] at (-0.2500000,0.6000000) {$6$};
\draw[gray,<->] (1.2000000,-0.6000000)--(2.4000000,-0.6000000);
\node[below,fill=none,inner sep=2pt] at (1.8000000,-0.6000000) {$10$};
\draw[gray,<->] (2.4000000,-0.6000000)--(4.2000000,-0.6000000);
\node[below,fill=none,inner sep=2pt] at (3.3000000,-0.6000000) {$15$};
\node[below,fill=none,inner sep=2pt] at (2.4000000,-1.6000000) {长度单位：cm（横向比例压缩）};
\end{tikzpicture}
\figcap{累次成像的光轴在棱镜斜面上转过直角；物体 $AB$ 位于棱镜上方。}\label{ch28:fig:prism-train}
\end{minipage}\end{center}
}

\newcommand{\SequentialImagingDiagram}{%
\begin{center}\begin{minipage}{\linewidth}\centering
\begin{tikzpicture}[>=Latex,font=\small,every node/.style={fill=none}]
\draw[gray,thin,->] (-3.4000000,0.0000000)--(6.2000000,0.0000000);
\draw[DeepBlue,thick,<->] (0.0000000,-1.8000000)--(0.0000000,1.8000000);
\draw[DeepBlue,thick,<->] (4.0000000,-1.8000000)--(4.0000000,1.8000000);
\draw[black,thick,->] (-3.0000000,0.0000000)--(-3.0000000,0.6000000);
\node[above,fill=none,inner sep=2pt] at (-3.0000000,0.6000000) {物 $O$};
\draw[black,thick,->] (1.5000000,0.0000000)--(1.5000000,-0.3000000);
\draw[gray,dashed,thin] (1.5000000,-0.3000000)--(1.5000000,1.0000000);
\node[above,fill=none,inner sep=2pt] at (1.5000000,1.0500000) {中间像 $I_1$};
\draw[black,thick,->] (5.6666667,0.0000000)--(5.6666667,0.2000000);
\node[above,fill=none,inner sep=2pt] at (5.6666667,0.2000000) {最终像 $I_2$};
\draw[DeepBlue,thick,->] (-3.0000000,0.6000000)--(0.0000000,0.4000000);
\draw[DeepBlue,thick,->] (0.0000000,0.4000000)--(4.0000000,-1.4666667);
\draw[DeepBlue,thick,->] (4.0000000,-1.4666667)--(5.6666667,0.2000000);
\draw[DeepBlue,thick,->] (-3.0000000,0.6000000)--(0.0000000,0.0000000);
\draw[DeepBlue,thick,->] (0.0000000,0.0000000)--(4.0000000,-0.8000000);
\draw[DeepBlue,thick,->] (4.0000000,-0.8000000)--(5.6666667,0.2000000);
\node[above,fill=none,inner sep=2pt] at (0.0000000,1.9000000) {$f_1$};
\node[above,fill=none,inner sep=2pt] at (4.0000000,1.9000000) {$f_2$};
\draw[gray,<->] (0.0000000,-2.0000000)--(4.0000000,-2.0000000);
\node[below,fill=none,inner sep=2pt] at (2.0000000,-2.0000000) {$D$};
\end{tikzpicture}
\figcap{两透镜逐级成像：$I_1$ 同时是第二透镜的物，$u_2=D-v_1$；实际光线穿过 $I_1$ 时没有折转。\label{ch28:fig:sequential}}
\end{minipage}\end{center}
}

% Source revision 2026-09-23: physical optics; 12 cm print width, >=9 pt.
% Layout: two independent horizontal return paths, f=1.5, L=3.
% [S at front focus -> lens -> mirror] | [S -> lens -> image on mirror].
% Rays are blue; return rays gold; dashed lines are axes only. Transparent labels.
\newcommand{\SelfConjugateDiagram}{%
\begin{center}\begin{minipage}{\linewidth}\centering\begin{tikzpicture}[>=Latex,font=\small,scale=.88]
\begin{scope}
 \draw[gray,dashed](-1.8,0)--(3.3,0);
 \draw[thick,<->](0,-1.2)--(0,1.2);\draw[very thick](3,-1.2)--(3,1.2);
 \fill(-1.5,0)circle(1.4pt);\node[above left=4pt]at(-1.5,0){$S=F$};
 \foreach \h in {-.65,.65}{
  \draw[DeepBlue,thick,->](-1.5,0)--(0,\h)--(3,\h);
  \draw[Gold!80!black,thick,->](2.6,\h)--(.2,\h);
 }
 \node[above]at(0,1.3){透镜};\node[above]at(3,1.3){平面镜};
 \draw[<->](-1.5,-1.5)--(0,-1.5)node[midway,below]{$f$};
 \node at(.65,-2.2){(a) 平行光自准直};
\end{scope}
\begin{scope}[xshift=7cm]
 \draw[gray,dashed](-3.3,0)--(3.3,0);
 \draw[thick,<->](0,-1.2)--(0,1.2);\draw[very thick](3,-1.2)--(3,1.2);
 \fill(-3,0)circle(1.4pt)node[below]{$S$};\fill(3,0)circle(1.4pt)node[right]{$I_1$};
 \draw[DeepBlue,thick,->](-3,0)--(0,.65)--(3,0);
 \draw[Gold!80!black,thick,->](3,0)--(0,-.65)--(-3,0);
 \draw[<->](-3,-1.5)--(0,-1.5)node[midway,below]{$u$};
 \draw[<->](0,-1.5)--(3,-1.5)node[midway,below]{$L$};
 \node[above]at(0,1.3){透镜};\node[above]at(3,1.3){平面镜};
 \node at(0,-2.2){(b) 第一次像落在镜面};
\end{scope}
\end{tikzpicture}\figcap{透镜和平面镜的两种自共轭。右图用 $L=2f$ 示范，返程光线与另一条去程光线互换。\label{ch28:fig:self-conjugate}}\end{minipage}\end{center}}


几乎所有光学系统都有一个共同任务：收集从物点发出的部分光线，经过反射或折射后，使它们严格或近似会聚于另一个点；或者使出射光看起来像从某一点发出。从波动观点说，光学元件改变入射波面的形状，使各部分在像点具有相同相位。完美成像一般需要复杂非球面和极高加工精度；实际系统往往限制孔径和视场，并在像质、通光量、体积与成本之间折中。

这项任务包含两个不可混淆的要求：一是“同一个物点”的各条光线应在同一个像点重合，二是不同物点的像还应保持可识别的相对位置。前者若不满足会形成模糊斑，后者若不满足会产生畸变。增加孔径能收集更多光、提高衍射分辨率，却会让球面像差更显著；缩小孔径可改善近轴近似，却降低亮度。所谓光学系统的像质，正是在这些条件和加工误差共同作用下评定，而不是只看某一条特征光线是否穿过焦点。

球面容易加工的原因也很直接：两个相配曲面在研磨时，高出的区域优先被磨去，长期相对运动会趋向共同球面。非球面能够更好校正像差，但面形检测和装调更难。图\ref{ch28:fig:contexts}A 的研磨情景解释为什么实际器件偏爱球面；后面的图~\ref{ch28:fig:s43-20-1} 则说明使用球面所付出的代价：一般只能在一定孔径内近似成像。

实际设计还必须限定视场。轴上物点的近轴光也许能很好会聚，离轴物点却可能出现彗差、像散和像面弯曲；因此“一个焦距”只描述系统的一阶性质，并不等于整幅画面处处清晰。镜头组通过不同材料和不同曲率的元件互相补偿这些误差，代价是表面数量、重量、装调难度和成本增加。

成像结论必须同时注明孔径、视场和近轴条件。

\historyfigure{optical-instruments-built-in-v1.png}{成像器件与仪器：光具座、暗箱、透镜、显微镜、望远镜和人眼。}

图\ref{ch28:fig:contexts}把加工、折射现象、成像实验和仪器放在同一条学习线上。A 的玻璃坯与研磨盘以配合球面接触；B 的刚性鸭形模型跨越水面，水下部分的视位置改变；C 的光源、透镜和屏共用一条光具座；D 的反光镜将镜头来光折向上方取景器。下面分别用球面、平面界面和逐级成像解释这些现象，场景图不用于量取角度或焦距。
\historyfigure[0.94\linewidth]{imaging-contexts.png}{研磨球面（A）、水面折射演示（B）、透镜成像实验（C）和单反相机剖面（D）。\label{ch28:fig:contexts}}



\subsection{成像不是画一条“代表光线”}
一点发出的光通常占据许多方向。只有这些光线经过系统后在另一点附近会聚，或其反向延长线在另一点附近相交，才可以谈这个物点的像。只追踪一条光线不能确定像的位置；至少要取两条独立光线，并检查邻近光线是否也满足同一成像关系。

所谓近轴，是相对于每个元件自身主轴而言，入射高度和倾角都足够小，使 $\sin\theta\simeq\tan\theta\simeq\theta$。它不等于“整张图画得很小”。一条相对图纸水平线倾角很小的光线，在强倾斜的曲面上仍可能大角度入射。使用球面或薄透镜公式时，既要检查光线相对于轴的角度，也要检查孔径相对于曲率半径是否足够小。

本章的物距与像距采用实物、实像为正的约定。折射后出射光真实相交处为实像；出射光发散而其反向延长线相交处为虚像。前一元件的像若位于后一元件的下游，后一元件截获的是尚未会聚的光束，这个点对后一元件而言就是虚物。先判断会聚、发散和传播方向，再赋予符号，比根据“左负右正”机械代数更可靠。


\section{平面界面与视深}



平面界面相当于 $R\to\infty$ 的球面折射面，故
\begin{equation}
 \frac{n_1}{u}+\frac{n_2}{v}=0,\qquad v=-\frac{n_2}{n_1}u.
\end{equation}
实物在单一平面折射后成为虚像。观察者在低折射率介质中看高折射率介质里的物体，视深为真实深度乘 $n_2/n_1$，所以水底看起来变浅；反过来从水中看空气里的物体，视距离增大。结论只对近轴光线成立，大视角时不同光线的反向延长线不严格共点。

图~\ref{ch28:fig:s43-20-4} 可直接给出这一比例。物体在介质 $n_1$ 中的真实深度为 $u$，观察者在介质 $n_2$ 中。取一条靠近法线的光，横向位移 $y$ 同时满足
\[
 y\simeq u\theta_1\simeq |v|\theta_2,
 \qquad n_1\theta_1\simeq n_2\theta_2.
\]
消去角度得
\[
 |v|=\frac{n_2}{n_1}u.
\]
从空气看水中，$n_2/n_1\approx1/1.33$，所以视深约为真实深度的四分之三；从水中看空气，比例反过来。这个推导同时说明观察方向若偏离法线很大，$\tan\theta\simeq\theta$ 与 $\sin\theta\simeq\theta$ 不再成立，所谓唯一“视深”只是近似。

\subsection{为何视深公式只能近轴使用}
设物点在平界面下深度 $h$ 处，物方折射率为 $n_1$，观察侧折射率为 $n_2$。选一条到达界面的光线，与法线夹角分别为 $i,r$，其界面交点离轴距离为 $\rho=h\tan i$。把出射光向后延长到轴，它与轴的交点深度为
\[
 h'=\frac{\rho}{\tan r}=h\frac{\tan i}{\tan r}.
\]
结合 $n_1\sin i=n_2\sin r$，有
\[
 \frac{h'}h=\frac{n_2\cos r}{n_1\cos i}
 \simeq\frac{n_2}{n_1}\quad(i,r\ll1).
\]
在水下向空气看，$n_1>n_2$，所以近轴视深小于真实深度。大角度时余弦比不可略去，不同方向的光线延长后与轴的交点不同，因而不能说平界面在任意孔径下都形成一个严格的点像。

对多层平行透明介质，逐层相加的是横向位移。若第 $j$ 层厚度为 $t_j$、折射率为 $n_j$，最终观察介质的折射率为 $n_o$，近轴时 $n_j\theta_j=n_o\theta_o$，故
\[
 \rho=\sum_jt_j\theta_j
 =\theta_o\sum_j\frac{n_ot_j}{n_j},
 \qquad h_{\rm app}=\sum_j\frac{n_ot_j}{n_j}.
\]
观察者到最外界面的距离若也计入总视距离，须另加这一段；它不是某一层介质的视厚度。全部 $n_j=n_o$ 时，视距离应回到几何距离，这是检验层数、起点和终点有没有算错的直接办法。


\ApparentDepthDiagram



\Needspace{8\baselineskip}\begin{origexamplebox}{水中直杆的最大表观弯折}
\label{ch28:ex:depth-stick}
一根直杆与水平液面成角 $\alpha$，液体对空气的折射率为 $n>1$。从近于正上方观察，采用近轴视深模型，杆在液体内的像与杆的真实方向相差 $\beta$。求使 $\beta$ 最大的 $\alpha$，并求最大弯折角。
\end{origexamplebox}\par\nopagebreak[3]
\begin{xesolution}
这里求的是观察到的杆像方向，不是某一条折射光线转过的角度。图~\ref{ch28:fig:depth-stick} 中杆与液面交于 $O$，杆上一点 $P$ 的水平距离为 $x$、真实深度为 $h$。正上方的近轴观察使横向位置保持不变，而深度变为 $h'=h/n$，所以它的虚像在 $P'$。

% Figure design: geometric optics; original drawing, no reference-image trace.
% Print width <=11cm; 9-10pt labels, transparent. One vertical section.
% Nodes O=(0,0), Q=(2.2,0), P=(2.2,-3.11127), P'=(2.2,-1.55563).
% Example geometry n=2, tan(alpha)=sqrt(2); h/h'=2 exactly.
% Solid OP=real rod, dashed OP'=apparent rod; vertical dimensions on right.
% Angles: alpha at r=.65, alpha' at r=1.20, beta between rays at r=1.8.
% Minimal figure, no decorative panels. Whitespace above water for view arrow.
% ASCII: observer up | O----Q----water; O\P' (dashed); O\P (solid); h,h' right.
\begin{center}\begin{minipage}{\linewidth}\centering
\begin{tikzpicture}[>=Latex,every node/.style={fill=none,font=\small,inner sep=2pt}]
 \coordinate (O) at (0,0);\coordinate (Q) at (2.2,0);
 \coordinate (P) at (2.2,{-2.2*sqrt(2)});
 \coordinate (V) at (2.2,{-1.1*sqrt(2)});
 \draw[thick] (-.6,0)--(4.45,0);
 \node[left] at (-.6,0){液面};
 \draw[very thick,DeepBlue] (O)--(P);
 \draw[thick,dashed,DeepBlue] (O)--(V);
 \draw[guide] (Q)--(P);
 \fill (O) circle(1.3pt);\fill (P) circle(1.3pt);\fill (V) circle(1.3pt);
 \node[above left] at (O){$O$};\node[above left] at (Q){$Q$};
 \node[right=5pt] at (P){$P$};\node[right=5pt] at (V){$P'$};
 \draw[->] (2.2,.4)--(2.2,1.05);
 \node[above] at (2.2,1.12){近于正上方观察};
 \draw (.65,0) arc[start angle=0,end angle=-54.73561,radius=.65];
 \node at (-18:.91){$\alpha$};
 \draw (1.2,0) arc[start angle=0,end angle=-35.26439,radius=1.2];
 \node at (-13:1.56){$\alpha'$};
 \draw (-35.26439:1.8) arc[start angle=-35.26439,end angle=-54.73561,radius=1.8];
 \node at (-45:2.12){$\beta$};
 \draw[<->] (3.2,0)--(3.2,{-2.2*sqrt(2)});
 \node[right] at (3.2,-2.45){$h$};
 \draw[<->] (4.05,0)--(4.05,{-1.1*sqrt(2)});
 \node[right] at (4.05,-.78){$h'$};
 \node[below] at (1.1,-3.65){实线为杆，虚线为杆的像，不是光线};
\end{tikzpicture}
\figcap{虚像 $P'$ 比 $P$ 浅，水平坐标不变。图取 $n=2$，角度关系由同一组点确定。\label{ch28:fig:depth-stick}}
\end{minipage}\end{center}

记像杆与液面的角为 $\alpha'$，则
\[
 \tan\alpha=\frac hx,\qquad
 \tan\alpha'=\frac{h'}x=\frac{\tan\alpha}{n},\qquad
 \beta=\alpha-\arctan\frac{\tan\alpha}{n}.
\]
对 $0<\alpha<\pi/2$ 求导时使用弧度。令 $t=\tan\alpha$，有
\[
 \frac{\dd\beta}{\dd\alpha}
 =1-\frac{n(1+t^2)}{n^2+t^2}
 =\frac{(n-1)(n-t^2)}{n^2+t^2}.
\]
分母为正且 $n>1$，故 $t^2<n$ 时递增，$t^2>n$ 时递减。唯一最大值发生在
\[
 \boxed{\alpha=\arctan\sqrt n},\qquad
 \alpha'=\arctan\frac1{\sqrt n}=\frac\pi2-\alpha.
\]
由此
\[
 \boxed{\beta_{\max}=2\arctan\sqrt n-\frac\pi2
 =\arcsin\frac{n-1}{n+1}}.
\]
当杆趋于水平或竖直时，弯折角都趋于零；当 $n\to1$ 时也趋于零。以水的 $n=4/3$ 为例，$\alpha\approx49.11^\circ$、$\beta_{\max}\approx8.21^\circ$。结论建立在近于正上方、局部近轴的观察模型上，不能用来描述任意斜视时整根杆的严格像。
\end{xesolution}


\section{球面折射与无球差点}



透镜是最常见的成像元件，人眼本身也包含折射系统。严格点到点成像通常要求非球面，但球面易于研磨加工；在限制孔径、只取近轴光线后，球面可以近似成像。

图~\ref{ch28:fig:s43-20-1} 中，物点 $S$ 和像点 $P$ 位于球面顶点 $V$ 的两侧，顶点距离分别为 $u,v$；球面半径为 $R$，两侧折射率为 $n_1,n_2$。若入射点 $A$ 对应球心角 $\varphi$，两段几何长度为
\[
 \ell_1=\sqrt{(u+R)^2+R^2-2R(u+R)\cos\varphi},
\]
\[
 \ell_2=\sqrt{(v-R)^2+R^2+2R(v-R)\cos\varphi}.
\]
上述几何先针对 $u,v,R>0$ 的实物实像配置：令 $V=(0,0)$、$C=(R,0)$，则 $A=(R-R\cos\varphi,R\sin\varphi)$，$S=(-u,0)$、$P=(v,0)$。直接计算 $|A-S|$ 和 $|P-A|$ 就得到两式；特别地，$\varphi=0$ 时分别回到 $u,v$。
总光程 $\mathcal L=n_1\ell_1+n_2\ell_2$。对非轴上入射点求导，费马条件给出
\[
 \frac{\partial\mathcal L}{\partial\varphi}
 =R\sin\varphi\left[\frac{n_1(u+R)}{\ell_1}-\frac{n_2(v-R)}{\ell_2}\right]=0.
\]
括号内一般随 $\varphi$ 变化，故不能用一个固定的 $v$ 同时满足所有孔径光线。近轴条件 $|\varphi|\ll1$ 下，令 $\ell_1\to u,\ell_2\to v$，得到
$n_1(1+R/u)=n_2(1-R/v)$，即
\begin{equation}
 \frac{n_1}{u}+\frac{n_2}{v}=\frac{n_2-n_1}{R}.
\end{equation}
将近轴式推广到虚物和虚像时，采用“实正虚负”的物像距约定，而不是把左右坐标都按同一方向读数：发散入射光对应实物 $u>0$，会聚入射光对应虚物 $u<0$；出射光真实会聚对应实像 $v>0$，出射光发散而反向延长线相交对应虚像 $v<0$。$R$ 在球心位于出射侧时为正，位于入射侧时为负。虚共轭须使用有向光程或延长线，不能把上面的正几何长度直接当成负长度。只要在每个界面重新识别物方、像方和传播方向，这一约定可连续使用。

近轴式并不是把严格关系中的角度直接“删掉”，而是保留光程随孔径高度 $y$ 的最低非平凡项。球面在顶点附近的弧高为 $x\simeq y^2/(2R)$。物点到界面的距离与界面到像点的距离分别展开为
\[
 \ell_1\simeq u+x+\frac{y^2}{2u},\qquad
 \ell_2\simeq v-x+\frac{y^2}{2v}.
\]
代入 $\mathcal L=n_1\ell_1+n_2\ell_2$ 后，若所有近轴光线到达像点的光程相同，$y^2$ 项的系数必须为零：
\[
 \frac{n_1}{2u}+\frac{n_2}{2v}+\frac{n_1-n_2}{2R}=0,
\]
这就得到球面折射公式。被略去的 $y^4,y^6,\ldots$ 项正是大孔径时球差等像差的来源，所以“球面能够成像”必须连同近轴、窄孔径条件一起表述。

严格关系中，$\ell_1,\ell_2$ 本身随入射点角度 $\varphi$ 改变；因此即使对某一条光线找到了令光程驻定的入射点，也不能保证所有 $\varphi$ 都给出同一个 $P$。近轴展开的作用，是让最低阶光程差只由 $u,v,R,n_1,n_2$ 决定而不再显含孔径高度。若要让大角度光也严格会聚，就必须改变曲面形状，使从物点经界面到像点的总光程对整个有效口径都相等；这一般给出非球面。例~\ref{ch28:ex:s43-20-6} 的特殊共轭点则利用几何与折射率关系，在球面上实现了少见的无球差成像。

\SphereRefractionPanel
\VirtualSignPanel

虚像点没有真实能量会聚，不能在该点直接放胶片或 CCD 接收，只能让另一光学系统或眼睛再次处理这些发散光。虚物则不会由普通发光物体单独产生，它只会出现在多元件系统中：前一元件产生的会聚光尚未到达会聚点，就先遇到下一界面。

“虚”描述的是到达当前元件之前或离开当前元件之后的光束状态，而不是说这个点没有物理意义。实物发出的是发散束；虚物对应尚未完成会聚的入射束。实像处确有光能会聚，可以放屏接收；虚像处没有出射光通过，只能由反向延长线定位。判断顺序应当是先看箭头方向，再看光束是会聚还是发散，最后决定 $u,v$ 的正负。仅凭某一点画在元件左边还是右边，在反射系统或折返光路中很容易出错。

例如，把一个新的折射面插入原本会在其后方会聚的光束中，这个会聚点就是该折射面的虚物，$u<0$；折射后若光仍在后方真实相交，则 $v>0$。反过来，普通物体经放大镜后出射光发散，延长线在物方相交，所以 $u>0,v<0$。图~\ref{ch28:fig:s43-20-2} 左图为实物经平界面形成虚像，右图为虚物经平界面形成实像。虚线分别表示出射光的反向延长线和未遇界面时的入射光会聚路径，不是真实传播光线。

\Needspace{8\baselineskip}\begin{origexamplebox}{【球透镜】}\label{ch28:ex:s43-20-1}
空气中半径为 $R$、折射率为 $n>1$ 的玻璃球，在近轴、窄孔径条件下把球外的小明成一个球外实像。求小明到其像的最小距离。
\end{origexamplebox}\par\nopagebreak[3]
\begin{xesolution}
设物到第一球面的距离为 $x$。第一次折射（空气到玻璃，半径 $R$）给出
\[
 \frac1x+\frac n{v_1}=\frac{n-1}{R},\qquad
 v_1=\frac{Rn x}{(n-1)x-R}.
\]
第二球面的物距为 $u_2=2R-v_1$，其有向半径为 $-R$。再次折射后得到球外实像距第二面
\[
 v_2=-R\frac{(n-2)x-2R}{2(n-1)x+(n-2)R}.
\]

%<optics-review:ball-second-surface>
第二次折射的代入过程可以展开为
\[
 \frac n{u_2}+\frac1{v_2}=\frac{1-n}{-R}=\frac{n-1}{R},
 \qquad
 u_2=R\frac{(n-2)x-2R}{(n-1)x-R}.
\]
于是
\[
 \frac1{v_2}=\frac{n-1}{R}
 -\frac{n[(n-1)x-R]}{R[(n-2)x-2R]}
 =-\frac{2(n-1)x+(n-2)R}{R[(n-2)x-2R]},
\]
取倒数即得上述 $v_2$。如果 $v_1>2R$，第一面折射后的光尚未会聚就遇到第二面，故第二面的物为虚物、$u_2<0$；不能把 $u_2$ 改成绝对值。

%</optics-review:ball-second-surface>
小明到像的总距离化简为
\[
 L=x+2R+v_2
 =\frac{2(n-1)(R+x)^2}{(n-2)R+2(n-1)x}.
\]
记分母 $D_x=(n-2)R+2(n-1)x$。实物及球外实像要求 $x>0$、$v_2>0$：当 $1<n<2$ 时，$x>(2-n)R/[2(n-1)]$；当 $n=2$ 时任意 $x>0$；当 $n>2$ 时，$0<x<2R/(n-2)$。
这些区间内均有 $D_x>0$。逐项求导得
\[
 L'(x)=\frac{4(n-1)(R+x)\{D_x-(n-1)(R+x)\}}{D_x^2}
 =\frac{4(n-1)(R+x)[(n-1)x-R]}{D_x^2}.
\]
所以 $x<R/(n-1)$ 时递减，$x>R/(n-1)$ 时递增，而驻点在上述每个允许区间内。因此全局最小值为
\[
 \boxed{L_{\min}=\frac{2nR}{n-1}}.
\]
此时光路关于球心对称；两次折射的符号均已按各自传播方向重新定义。

%<optics-review:ball-parallel-limit>
还需检查一个容易漏掉的特殊情况：最优物距 $x_*=R/(n-1)$ 恰使 $v_1$ 表达式的分母为零。这不是物理过程失效，而是第一面出射光变为平行光，中间像在无穷远。直接令第一面公式中的 $n/v_1=0$；第二面也有 $n/u_2=0$，因而
\[
 \frac1{v_2}=\frac{n-1}{R},\qquad
 v_2=x_*=\frac R{n-1}.
\]
所以 $L=x_*+2R+v_2=2nR/(n-1)$，正好补上中间代数除法没有覆盖的这一点。$n\to1^+$ 时最小距离趋于无穷大，符合玻璃球与空气折射率相同时不再具有聚光能力；整个结果仍限于近轴、有效孔径较小的成像模型。

%</optics-review:ball-parallel-limit>
\end{xesolution}

% Relocated AplanaticLensPanel to its own solution.

\Needspace{8\baselineskip}\begin{origexamplebox}{【大角度成像】}\label{ch28:ex:s43-20-6}
显微镜物镜中有一枚两面均为球面的弯月形厚透镜。左表面 $S_1$ 球心为 $C_1$、半径 $R_1$；右表面 $S_2$ 球心为 $C_2$、半径 $R_2$。两球心都在相应表面的左侧，$C_1$ 又在 $C_2$ 左侧；玻璃对空气折射率为 $n>1$，且 $C_1C_2=R_2/n$。物位于 $C_1$。证明：（1）所有通过有效口径的出射光线的反向延长线严格交于一点 $Q$；（2）$QC_2=nR_2$，并判断像的虚实。
\end{origexamplebox}\par\nopagebreak[3]
\begin{xesolution}
\AplanaticLensPanel

从 $C_1$ 发出的任意光线到达 $S_1$ 时沿该球面的半径入射，所以第一面入射角为零，不发生偏折。于是玻璃内所有光线仍像从 $C_1$ 发出。

把图~\ref{ch28:fig:s43-20-9} 中第二球心取为原点：$C_2=(0,0)$、$C_1=(-R_2/n,0)$，在同侧选 $Q=(-nR_2,0)$。设第二面点 $P=R_2(\cos\varphi,\sin\varphi)$，并记 $a=|P-C_1|$、$b=|P-Q|$。直接计算
\[
 a^2=R_2^2(1+n^{-2}+2n^{-1}\cos\varphi),\qquad
 b^2=R_2^2(1+n^2+2n\cos\varphi)=n^2a^2.
\]
因此 $b=na$。外法向为 $\boldsymbol N=P/R_2$，入射方向 $(P-C_1)/a$，候选出射方向 $(P-Q)/b$。用二维叉积求它们与法线夹角的正弦：
\[
 \sin i=\frac{R_2|\sin\varphi|}{na},\qquad
 \sin r=\frac{nR_2|\sin\varphi|}{b}
 =\frac{R_2|\sin\varphi|}{a}=n\sin i.
\]
两方向的切向分量同号，符合玻璃到空气的折射定律。取通常朝右的有效口径 $|\varphi|<\pi/2$ 时，两方向与外法线的点积均为正，确为向外传播且不发生全反射；还需选取两球面间有正厚度、入射线先穿过 $S_1$ 的口径。
出射方向是从 $Q$ 指向 $P$ 再向右延伸，\textbf{不是从 $P$ 射向 $Q$}。故只有反向延长线在 $Q$ 相交，形成严格无球差虚像，且
\[
 \boxed{QC_2=nR_2}.
\]
证明没有使用小角近似，所以 $C_1,Q$ 是该球面的两个共轭无球差点；后续物镜可利用这束仿佛从 $Q$ 发出的光继续成像。图中到 $Q$ 的虚线表示延长线，在 $Q$ 处不能直接放屏接收此像。
\end{xesolution}

\section{球面镜、薄透镜与物像关系}



非球面反射镜造价较高，球面反射镜在近轴范围内同样可近似成像。图~\ref{ch28:fig:s43-20-8} 中 $CP$ 是法线，反射角等于入射角；对三角形作小角近似可得
\begin{equation}
 \frac1u+\frac1v=\frac2R=\frac1f,\qquad f=\frac R2.
\end{equation}
实物、虚物、实像和虚像的符号意义与透镜一致，只需注意反射后传播方向反转。横向放大率仍为 $m=-v/u$。凸球面镜可取 $R<0$、$f<0$，其作用类似发散透镜。

几何推导中，入射点 $P$、球心 $C$ 和像点必须位于同一子午平面。设物点、像点和球心在主轴上的距离分别为 $u,v,R$，而 $P$ 相对轴的高度很小。半径 $CP$ 是法线，反射定律给出法线两侧角度相等。把三角形中的正弦用小角本身代替，并用“物点方向角 $+$ 像点方向角 $=2\times$ 法线方向角”，就得到
\[
 \frac yu+\frac yv\simeq2\frac yR.
\]
约去非零的 $y$ 后即为球面镜公式。这个过程明确显示了因子 $2$ 来自法线把方向转过两倍，而不是来自光程“往返两次”。

平面镜是 $R\to\infty$ 的极限，因此
\begin{equation}
 v=-u.
\end{equation}
普通实物经平面镜成等大的虚像；一束原本将会聚于镜后虚物点的光，经平面镜也可以在镜前形成实像。单镜头反光相机的折转光路中就包含类似过程。

图\ref{ch28:fig:contexts}D 中反光镜处于镜头预定焦面之前，它把仍会聚的光束折向取景光路，并不是把镜头的会聚能力消除；拍摄时反光镜让开，光才能到达后方感光面。

球面镜公式同样只在近轴范围成立。对凹面镜，平行轴的边缘光与近轴光通常不在同一点会聚，这是球差；把孔径缩小能减轻球差，却同时降低通光量并增强衍射。凹面镜的 $f=R/2$ 只是顶点附近的傍轴焦距，不应把图~\ref{ch28:fig:s43-20-8} 上任意高度的反射点都强行连到这一点。平面镜则没有这一曲率近似：用镜像法可精确证明物、像关于镜面对称，但虚像一侧并无真实光线通过。

镜像法的精确性可由任意一点验证。把物点 $O$ 关于镜面对称到 $O'$，对镜面上任一反射点 $P$ 都有 $PO=PO'$；反射定律又使折线 $O\to P\to$ 观察者等价于从 $O'$ 发出的直线。所有反射光的反向延长线因此严格通过同一个 $O'$，且 $|v|=u$、横向倍率为 $1$。若入射光原本向镜后一点会聚，把该点视为虚物并作同样的对称，反射光会在镜前真实会聚，这说明平面镜也能“虚物成实像”。

\SphericalMirrorDiagram



传统透镜至少包含一个曲面。中心比边缘厚的透镜通常会聚，中心较薄的透镜通常发散。忽略厚度的薄透镜可看作两个球面折射面的组合。对置于空气中的折射率 $n$ 的薄透镜，两个有向曲率半径为 $R_1,R_2$，连续应用球面折射式并消去中间像，得到制镜者公式和成像公式
\begin{equation}
 \frac1f=(n-1)\left(\frac1{R_1}-\frac1{R_2}\right),
 \qquad
 \frac1u+\frac1v=\frac1f.
\end{equation}
会聚透镜 $f>0$，发散透镜 $f<0$。横向放大率（带方向）为 $m=-v/u$，其绝对值
\begin{equation}
 M_T=\left|\frac vu\right|=\left|\frac f{u-f}\right|.
\end{equation}
负号表示实像倒立。作图时可用三条近轴特征光线：过光心近似不偏折；入射平行轴的光经像方焦点；入射通过物方焦点的光出射平行轴。

把两个折射面的中间步骤写出，可以避免把 $R_2$ 的符号背错。第一面由空气进入玻璃：
\[
 \frac1u+\frac n{v_1}=\frac{n-1}{R_1}.
\]
薄透镜厚度忽略后，第一面的像是第二面的虚物，故第二面物距 $u_2=-v_1$。从玻璃射回空气：
\[
 \frac n{u_2}+\frac1v=\frac{1-n}{R_2}.
\]
两式相加时 $v_1$ 恰好消去，才得到制镜者公式。双凸透镜按从左向右传播的约定通常有 $R_1>0,R_2<0$，所以两项共同贡献正光焦度；若只把两个半径都当作无符号长度，第二项就会误消掉第一项。

透镜成像图还提供一套数值自检：$u>2f$ 时为缩小倒立实像；$u=2f$ 时等大；$f<u<2f$ 时为放大倒立实像；$0<u<f$ 时为正立放大虚像。$u\to\infty$ 应给出 $v\to f$，$u\to f^+$ 时 $v\to+\infty$。这些极限若与计算结果不符，应先检查物距、焦距和传播方向的符号，而不是修改作图去迎合错误答案。

“中心厚”与“会聚”是对普通空气中、普通正折射率玻璃的经验说法；决定光焦度的仍是折射率差和两个有向曲率。把同一玻璃透镜浸入折射率更高的液体，光焦度可能减小、为零甚至反号。薄透镜近似还要求厚度相对曲率半径和物像距离足够小，使两个主平面可合并为一个平面；厚透镜中必须保留两面的间隔，中间像位置也不能简单取 $u_2=-v_1$。

图~\ref{ch28:fig:s43-20-5} 下图把两个界面分开：$P$ 处以 $C_1P$ 为法线折射，玻璃内部的 $PQ$ 保持直线；到 $Q$ 才以 $C_2Q$ 为法线再次折射。上面的薄透镜公式，是把这两个依次发生的过程在厚度可忽略时合并，而不是让光线在玻璃内部任意拐弯。
\ThinLensPanel
\LensImagingDiagram
图~\ref{ch28:fig:s43-20-6} 的三条特征光线在同一像点会合；通过物方焦点的光线出射后平行于光轴，不能把它误画成继续通过像方焦点。

图\ref{ch28:fig:contexts}C 的实验中，移动屏直到光斑最清晰，所测的屏距应与物距一起满足成像公式；若像为虚像，不能靠移动出射侧屏找到它。

例如玻璃折射率 $n=1.54$，双凸透镜两面半径大小分别为 $100\,\mathrm{cm}$ 和 $20\,\mathrm{cm}$（按方向取 $R_1=100\,\mathrm{cm},R_2=-20\,\mathrm{cm}$），则 $f\approx30.9\,\mathrm{cm}$。物距 $u=70.0\,\mathrm{cm}$ 时，$v=fu/(u-f)\approx55.2\,\mathrm{cm}$，$|m|\approx0.788$，形成缩小倒立的实像。

横向放大率来自过光心近似不偏折的两相似三角形，故带符号形式为 $m=-v/u$。它只描述垂直主轴的微小线段，不保证任意平面图形经透镜后保持相似：图形上不同点若有不同轴向坐标，就具有不同物距和放大率。圆、抛物线以及运动轨迹经透镜后是否仍保持原曲线类型，需要把每个点分别按 $X=fx/(x-f)$、$Y=-fy/(x-f)$ 变换，不能把一个统一倍率套到整幅有深度的图形上。这正是例~\ref{ch28:ex:s43-20-3} 的出发点。



\Needspace{8\baselineskip}\begin{origexamplebox}{【透镜成像】}\label{ch28:ex:s43-20-3}
半径为 $R$ 的圆经焦距为 $f$ 的凸透镜成像。（1）圆心放在哪里时像仍为圆？（2）抛物线能否仍成抛物线？（3）透镜中看到的抛体运动能否仍像抛体运动？
\end{origexamplebox}\par\nopagebreak[3]
\begin{xesolution}
取光轴为 $x$ 轴，物点 $(x,y)$ 的近轴像为
\[
 X=\frac{fx}{x-f},\qquad Y=-\frac{fy}{x-f}.
\]
这里先讨论图形位于包含光轴的子午面、圆心在轴上的情形。$x>0$ 表示物点位于透镜左侧距离 $x$ 处，$X>0$ 表示像点在右侧；两侧横向坐标均向上为正。反解得 $x=fX/(X-f)$、$y=-fY/(X-f)$。圆 $(x-a)^2+y^2=R^2$ 变换后为
\[
 [(f-a)X+af]^2+f^2Y^2=R^2(X-f)^2.
\]
展开后 $X^2$ 的系数为 $(f-a)^2-R^2$，$Y^2$ 的系数为 $f^2$，两者相等且非零才仍为圆，得到
\[
 (a-f)^2=R^2+f^2.
\]
对透镜同侧的实物圆，取
\[
 \boxed{a=f+\sqrt{f^2+R^2}};
\]
另一代数根描述另一侧或虚物配置。

对轴向抛物线 $y^2=4p(x-a)$，同样代入得
\[
 f^2Y^2=4p[(f-a)X+af](X-f).
\]
当 $a=f$ 时右侧化为 $4pf^2(X-f)$，所以 $Y^2=4p(X-f)$：\textbf{存在特殊放置}使抛物线仍成抛物线。一般放置则会变成椭圆或双曲型圆锥曲线。这里讨论各曲线上满足近轴与有效口径的有限部分，不意味着无穷大的图形都能通过透镜。

抛体的\textbf{轨迹形状}与\textbf{运动规律}要分开。若子午面内 $x(t)=x_0+v_xt$，则
\[
 X(t)=f+\frac{f^2}{x_0+v_xt-f},\qquad
 X''(t)=\frac{2f^2v_x^2}{(x_0+v_xt-f)^3}.
\]
当 $v_x\ne0$ 时，纵向像加速度随时间变化，即使像轨迹恰为抛物线，也一般不是匀加速度抛体。若圆心偏离光轴为 $(a,b)$，代入后的圆锥曲线另含 $2bfXY$ 项；要在这个子午面仍为圆，还须 $b=0$。这补足了第一问中圆心的横向位置条件。

题目若不限定图形朝向，还存在另一类简单答案：将圆或抛体轨迹放在\textbf{垂直光轴的同一物平面}内，所有点物距均为 $u_0$，倍率 $m=-f/(u_0-f)$ 为常数。圆自然仍成圆，无须前述特殊轴向位置。若该平面中的抛体坐标为
\[
 y(t)=y_0+v_yt,\qquad z(t)=z_0+v_zt-\tfrac12gt^2,
\]
像坐标为 $Y=my,Z=mz$，故 $Y''=0,Z''=-mg$。忽略光传播时间差，在近轴有效视场内，像确实仍作非退化的匀加速度抛体运动。因此第三问“能否”的回答是\textbf{能}，上述横向平面就是明确例子；其等效加速度通常不等于当地重力加速度。
\end{xesolution}

\subsection{从平行光到近物：七种入射光束的像}
式 $1/u+1/v=1/f$ 不只给数值，也控制出射光束是会聚、平行还是发散。图~\ref{ch28:fig:lens-states-real}、图~\ref{ch28:fig:lens-states-virtual} 将同一会聚薄透镜的七种状态逐一画出。光沿左向右传播，$F,F'$ 分别是物方、像方焦点；图中物点和像点都在轴上。蓝色实线是实际光线，灰色虚线仅是出射光的反向延长线，不能把虚线看成有光沿它反向传播。

% Figure contract: seven axial bundles, deterministic paraxial lens, f=1.2.
% Purpose: distinguish real/virtual/infinite images, restore seven source states.
% Print width 14cm; small text (9-10pt), no opaque label backgrounds.
% Shared coordinates: lens x=0, incident slope y/u, outgoing slope y/u-y/f.
% Layout: four real-image rows, then three parallel/divergent rows, 2.35cm pitch.
% Diagram x=-4.8..4.8; annotations x=5.0; ticks and labels below beam region.
% Rays are straight inside each medium; only the lens changes direction.
% ASCII: object * --> | lens | --> * image ; row title right, focal ticks below.
% For u=0 sentinel use parallel incidence; u=f is calculated without 1/(u-f).
% Quality: common f and aperture; seven cases are not decorative variants.
\newcommand{\LensStateRow}[5]{%
\begin{scope}[yshift=#1cm]
 \draw[gray,thin](-4.8,0)--(4.7,0);
 \draw[thick,<->](0,-.55)--(0,.55);
 \foreach \x/\lab in {-2.4/2F,-1.2/F,1.2/F',2.4/2F'}{
  \draw[gray] (\x,-.04)--(\x,.04);
  \draw[gray,thin] (\x,-.6)--(\x,-.73);
  \node[below,font=\scriptsize]at(\x,-.73){$\lab$};
 }
 \foreach \rawheight in {-.32,-.16,.16,.32}{
  \pgfmathsetmacro{\h}{ifthenelse(and(#2>0,#2<.6),.6875*\rawheight,\rawheight)}
  \pgfmathsetmacro{\slopein}{ifthenelse(#2==0,0,\h/max(#2,.001))}
  \pgfmathsetmacro{\slopeout}{\slopein-\h/1.2}
  \pgfmathsetmacro{\leftx}{ifthenelse(#2==0,-4.6,-#2)}
  \draw[DeepBlue,thin,->](\leftx,{\h+\slopein*\leftx})--(0,\h);
  \draw[DeepBlue,thin,->](0,\h)--(#3,{\h+\slopeout*#3});
  \ifnum#4=1
   \pgfmathsetmacro{\virtualx}{-\h/\slopeout}
   \draw[gray,dashed,thin](\virtualx,0)--(0,\h);
  \fi
 }
 \node[right,align=left,text width=3.6cm]at(5,.03){#5};
\end{scope}}
\begin{center}\begin{minipage}{\linewidth}\centering
\begin{tikzpicture}[>=Latex,every node/.style={fill=none,inner sep=2pt,font=\small},scale=.9]
 \LensStateRow{0}{0}{2.1}{0}{(a) $u\to\infty$\\平行入射，$v=f$}
 \LensStateRow{-2.35}{3.6}{3.2}{0}{(b) $u>2f$\\$f<v<2f$}
 \LensStateRow{-4.7}{2.4}{3.6}{0}{(c) $u=2f$\\$v=2f$}
 \LensStateRow{-7.05}{1.8}{4.6}{0}{(d) $f<u<2f$\\$v>2f$}
\end{tikzpicture}
\figcap{四种形成实像的轴上光束。所有会聚点都由 $v=fu/(u-f)$ 决定，而不是凭图形外观选取。\label{ch28:fig:lens-states-real}}
\end{minipage}\end{center}

图~\ref{ch28:fig:lens-states-real}(b)—(d) 中，物点靠近焦点时，像点向右远离透镜。代入 $u=3f,2f,3f/2$，分别得 $v=3f/2,2f,3f$。当 $u\to f^+$，$v\to+\infty$；因此不能在有限屏幕上继续寻找一个焦点附近的像。

\begin{center}\begin{minipage}{\linewidth}\centering
\begin{tikzpicture}[>=Latex,every node/.style={fill=none,inner sep=2pt,font=\small},scale=.9]
 \LensStateRow{0}{1.2}{4.6}{0}{(e) $u=f$\\出射平行\\像在无穷远}
 \LensStateRow{-2.35}{.9}{2.5}{1}{(f) $f/2<u<f$\\$v<0,\ |v|>f$}
 \LensStateRow{-4.7}{.55}{1.8}{1}{(g) $0<u<f/2$\\$v<0,\ |v|<f$}
\end{tikzpicture}
\figcap{从焦点移向透镜，出射光由平行变为发散；后两图的虚像由反向延长线相交确定。\label{ch28:fig:lens-states-virtual}}
\end{minipage}\end{center}

图~\ref{ch28:fig:lens-states-virtual}(f)、(g) 分开了虚像在物方焦点之外和之内的两种情况。由
\[
 |v|=\frac{fu}{f-u},\qquad
 |v|>f\ \Longleftrightarrow\ u>\frac f2
\]
可见分界不是第二种透镜，而是同一透镜的物距 $u=f/2$；这时虚像恰在 $F$。例如 $u=3f/4$ 时 $v=-3f$，$u=f/3$ 时 $v=-f/2$。虚像仍位于物点的左侧，但光并未真的会聚在那里。实际有限高度物体的正倒与大小，需再用横向放大率 $m=-v/u$ 判断：(b)—(d) 的实像为倒立，(f)、(g) 的虚像为正立放大。

\section{组合系统与逐次成像}



多元件系统应沿传播方向逐级成像：前一元件的像就是后一元件的物。两薄透镜焦距为 $f_1,f_2$、间距为 $D$，第一透镜物距为 $u_1$ 时
\[
 v_1=\frac{f_1u_1}{u_1-f_1},\qquad
 u_2=D-v_1,\qquad
 v_2=\frac{f_2u_2}{u_2-f_2}.
\]
$v_1>D$ 时，第二透镜前的光正在会聚，故 $u_2<0$ 是虚物；$v_1<0$ 时第一透镜产生虚像，第二透镜看到的物距为 $D+|v_1|$。总横向放大率是各级带符号放大率之积。只背数值而不画传播方向，很容易把虚物和虚像混淆。

\SequentialImagingDiagram

图\ref{ch28:fig:sequential}的 $I_1$ 位于两镜之间，所以第二透镜的物距为正；若把第二透镜移到 $I_1$ 左侧，$I_1$ 就成为第二透镜的虚物。这个正负转换来自光束状态，而不是把中间像删掉。

逐级计算时建议在每个元件前写三行：光线此刻向哪边传播、入射束是发散还是会聚、若没有该元件它将从哪里来或到哪里去。第一问决定正方向，第二问决定实物或虚物，第三问给出物距的几何大小。中间像即使落在两个元件之外，也不能被删除；它仍通过入射束的会聚或发散状态影响下一次成像。最后还应检查总放大率 $m=m_1m_2\cdots$ 的符号与图中倒立次数是否一致。

显微镜和望远镜体现了两种典型组合方式。显微镜的正物镜把近处小物形成放大的中间实像，正目镜再把这个中间像当作焦内实物，给眼睛形成放大的虚像；望远镜的物镜把远处目标形成焦面附近的中间实像，目镜主要放大它对眼睛的视角。两者都需要中间像进入目镜的有效孔径，否则即使成像公式给出位置，观察者也收不到完整光束。伽利略望远镜用负目镜在物镜光束尚未会聚前截取虚物，系统短且成正像，却没有可放屏或十字丝的中间实像，视场也较受限制。

评价组合系统不能只看倍率。物镜口径决定收集光的能力和衍射分辨率，焦距组合决定角放大率，孔径光阑决定哪些倾斜光束能通过，而像差、镀膜和装调决定这些理论性能能否实现。倍率超过物镜能够提供的空间信息时，只是把衍射斑和像差一同放大。于是“高级望远镜”的简易判断应先看有效口径、像质和机械稳定性，再看可用倍率，而不能只比较广告中的最大倍数。

\Needspace{14\baselineskip}





\subsection{中间像、虚物和无穷远像的衔接}
两个薄透镜沿同一方向依次排列，间距为 $d$。第一镜求出的像距为 $v_1$ 时，第二镜的物距是 $u_2=d-v_1$，不是 $d+v_1$。当 $v_1>d$，第一镜单独形成的像位于第二镜之后，入射到第二镜的光尚在会聚，于是 $u_2<0$，恰好对应虚物。

这个关系来自同一个轴向坐标：第一镜在 $x=0$，第二镜在 $x=d$，中间像在 $x=v_1$。因此第二镜到它的有向物距为 $d-v_1$。每一次代入前都把这一位置关系写出，能避免“中间像是实像，所以对下一镜也必是实物”的错误。

若物恰在第一镜物方焦平面上，第一镜出射光平行，不能把 $v_1=\infty$ 当普通有限数计算。轴外物点高度为 $h$ 时，出射光束倾角为 $\theta=-h/f_1$。第二镜把这束平行光聚到自身后焦平面，高度
\[
 h'=f_2\theta=-\frac{f_2}{f_1}h.
\]
这里通过“位置—角度—位置”完成成像，虽然中间像在无穷远，横向倍率仍是有限的 $-f_2/f_1$。两镜间距影响有限孔径能否接收整个光束，但不改变这个理想焦平面到焦平面的近轴倍率。

反射折返系统还多一步：反射之后要沿新的传播方向重新确定下一个物距。推荐先沿实际光路给元件编号，再逐次记录“当前元件—入射束会聚点—出射束会聚点”，最后才相乘各级横向倍率。屏的位置、像的位置、透镜的位置是三种不同对象，不能用同一个字母代替。

\Needspace{8\baselineskip}\begin{origexamplebox}{【奇妙的成像】}\label{ch28:ex:s43-20-2}
用焦距为 $f$ 的凸透镜和平面镜组合成像。点光源经透镜、镜面反射并再次通过透镜后，像回到光源。光源位于主轴上，平面镜垂直于主轴；若平面镜到透镜的距离为 $L$，求光源到透镜的两个可能距离，并说明真实光源何时可实现。
\end{origexamplebox}\par\nopagebreak[3]
\begin{xesolution}
把平面镜后的返程光路关于镜面展开，问题等价于两枚相同薄透镜相距 $2L$，物和最终像分别位于系统两端、距各自透镜均为 $u$。用空气中光线列向量 $(y,\theta)^{\mathsf T}$，传播矩阵为 $T(d)=\left(\begin{smallmatrix}1&d\\0&1\end{smallmatrix}\right)$，薄透镜矩阵为 $F=\left(\begin{smallmatrix}1&0\\-1/f&1\end{smallmatrix}\right)$。依次相乘得
\[
 FT(2L)F=\begin{pmatrix}1-2L/f&2L\\2(L-f)/f^2&1-2L/f\end{pmatrix}.
\]
再在左右各乘 $T(u)$，总矩阵的右上元为
\[
 B=2u(1-2L/f)+2L+\frac{2(L-f)u^2}{f^2}.
\]

%<optics-review:matrix-meaning>
为什么检查右上元 $B$ 就够了？从轴上物点发出一条倾角为 $\theta_0$ 的光，其初始列向量为 $(0,\theta_0)^{\mathsf T}$。若总矩阵为 $\left(\begin{smallmatrix}A&B\\C&D\end{smallmatrix}\right)$，到候选像面时光线高度便是 $y=B\theta_0$。不同 $\theta_0$ 都必须给出 $y=0$，所以必须 $B=0$。例如将上面的两镜矩阵记为 $\left(\begin{smallmatrix}a&b\\c&a\end{smallmatrix}\right)$，左右各乘传播矩阵后，右上元就是 $b+2au+cu^2$；代入 $a=1-2L/f$、$b=2L$、$c=2(L-f)/f^2$，便得到刚才的二次式。这里的每一项都有长度量纲。

%</optics-review:matrix-meaning>
同一物点发出的不同倾角光必须在同一像点重合，所以 $B=0$。乘去分母并分解，回到原点的条件即
\[
 (u-f)\,[Lf-Lu+fu]=0.
\]
故
\[
 \boxed{u_1=f,\qquad u_2=\frac{Lf}{L-f}}.
\]
第一种情况下透镜把轴上点源的光变为平行于轴的光，垂直放置的平面镜令其原路返回。第二种可不用矩阵理解：第一次透镜成像恰在镜面上，即 $v_1=L$，于是 $1/u+1/L=1/f$。镜面反射后，这个镜面上的像点成为返程实物，再经透镜成像回原点。两条不同的光线在镜面相交后互换返程路径，不能说每一条都沿自身原路返回。
\SelfConjugateDiagram
图~\ref{ch28:fig:self-conjugate} 中 $S=F$ 表示物点恰在透镜的物方焦点；第一次出射为平行光束，折返后再由透镜会聚。
第二根对应有限共轭位置。$u_2$ 要成为镜前有限距离实物位置需 $L>f$；$L<f$ 时第二根对应虚物，不能靠移动真实点光源实现；$L=f$ 时第二根趋于无穷远，有限物距方程只剩 $u=f$。
\end{xesolution}

\Needspace{8\baselineskip}\begin{origexamplebox}{【猫眼探秘】}\label{ch28:ex:s43-20-7}
门上的猫眼由室外侧焦距 $-1.0\,\mathrm{cm}$ 的凹透镜和室内侧焦距 $3.5\,\mathrm{cm}$ 的凸透镜组成，间距 $2.1\,\mathrm{cm}$。说明为什么室外看不清室内而室内能看清室外；若试图从室外看室内应怎样做？
\end{origexamplebox}\par\nopagebreak[3]
\begin{xesolution}
\CatEyePanel

%<optics-review:cat-eye-reference>
图~\ref{ch28:fig:s43-20-10} 标出了同一门孔中的室外凹透镜与室内凸透镜。以下先以远处景物作近轴计算；正向光按“凹透镜—凸透镜”传播，反向则按“凸透镜—凹透镜”传播。每到一个元件都沿新的传播方向判定实物与虚物，不能只交换两个焦距而保留原来的物距符号。

%</optics-review:cat-eye-reference>

室外远物先经凹透镜，成为位于凹透镜室外侧 $1.0\,\mathrm{cm}$ 的虚像。对凸透镜，这个物的距离为 $u=2.1+1.0=3.1\,\mathrm{cm}<3.5\,\mathrm{cm}$，故
\[
 v=\frac{3.5\times3.1}{3.1-3.5}\approx-27.1\,\mathrm{cm}.
\]
凸透镜把前一级的小虚像再放大，最终虚像位于可调焦的约 $27.1\,\mathrm{cm}$ 处。这里的“放大”针对中间像，并不表示门外景物的总角放大率大于一；广角门镜通常压缩景物的角尺度，实际视场还受孔径和像差限制。

反向观察远处室内景物时，凸透镜先把平行光会聚于其室外侧 $3.5\,\mathrm{cm}$ 处；凹透镜位于途中，面对的是其外侧 $1.4\,\mathrm{cm}$ 处的虚物，即 $u=-1.4\,\mathrm{cm}$。由 $f=-1.0\,\mathrm{cm}$ 得 $v=-3.5\,\mathrm{cm}$，反向光束在凹透镜后发散，反向延长线交于室内侧距该凹透镜 $3.5\,\mathrm{cm}$ 处。能观察到多大范围，还须检查门孔、两片透镜和眼瞳对斜光束的截取，不能仅凭像距断言视场大小。

反向观察还有更直接的调焦困难：若眼睛几乎贴着室外透镜，虚像到眼睛只有约 $3.5\,\mathrm{cm}$，远小于通常约 $25\,\mathrm{cm}$ 的近点；仅仅贴近孔径、遮挡杂散光不能让正常眼清楚调焦。眼距为 $e$ 时虚像距眼约 $e+3.5\,\mathrm{cm}$，后退可能满足近点条件，却会严重损失视场和接收光量。
从光学原理上，可用适当会聚光学系统重新成像，使最终像处于观察眼可调焦的范围；例如忽略元件间距时，让焦距约 $3.5\,\mathrm{cm}$ 的会聚元件以该虚像为实物，可使出射光近似平行。这只解决调焦，不消除门镜原有的有限孔径、像差与视场限制。由此可见，\textbf{瞳孔接光}和\textbf{眼睛调焦}是两个必须分别满足的条件。
\end{xesolution}

% Relocated SilveredLensPanel to its own solution.

\Needspace{8\baselineskip}\begin{origexamplebox}{【镀银透镜】}\label{ch28:ex:s43-20-8}
一平凸透镜焦距为 $f$，其平面镀银。凸面一侧距透镜 $2f$ 处放置高为 $H$ 的物，物垂直主轴且下端在轴上。（1）作图判断像的虚实；（2）计算像的位置和大小。
\end{origexamplebox}\par\nopagebreak[3]
\begin{xesolution}
\SilveredLensPanel

光先通过透镜，遇平面镀银层反射，再反向通过同一透镜。薄元件近似下，透镜光焦度 $1/f$ 被使用两次，平面镜本身不增加光焦度，所以组合等效为焦距
\[
 F=\frac f2
\]
的凹面镜。图~\ref{ch28:fig:s43-20-11}用等效薄平面画两条代表光线，它们返回物方并真实相交，故为实像。逐级追踪也得到同一结论：第一次透镜给 $v_1=2f$；在光尚未会聚前由紧贴的平面镜反射，会聚点关于镜面对称到物方 $2f$ 处。这个点对返程透镜是虚物，$u_2=-2f$，于是 $-1/(2f)+1/v=1/f$，得 $v=2f/3$。第一次倍率 $-1$、平面镜横向倍率 $+1$、返程透镜倍率 $+1/3$，乘积为 $-1/3$。

等效镜的物距 $u=2f=4F$，成像公式给出
\[
 \frac1{2f}+\frac1v=\frac1F=\frac2f,\qquad
 \boxed{v=\frac{2f}{3}}.
\]
像位于凸面外物方 $2f/3$ 处。横向放大率 $m=-v/u=-1/3$，所以
\[
 \boxed{H'=\frac H3},
\]
方向倒立。把镀银面误当成“成一次平面镜虚像”而停止追踪，会漏掉返程的第二次折射。
\end{xesolution}

\Needspace{22\baselineskip}\begin{origexamplebox}{【累次成像】}\label{ch28:ex:s43-20-9}
折射率 $n=1.5$、直角边长 $6\,\mathrm{cm}$ 的全反射棱镜上方 $6\,\mathrm{cm}$ 放置物体 $AB$。棱镜右侧 $10\,\mathrm{cm}$ 处为焦距 $f_1=10\,\mathrm{cm}$ 的凸透镜，再向右 $15\,\mathrm{cm}$ 为焦距大小 $10\,\mathrm{cm}$ 的凹透镜。求最终像的位置和横向放大率。
\end{origexamplebox}\par\nopagebreak[3]
\begin{xesolution}
\PrismTrainPanel

先沿图~\ref{ch28:fig:prism-train}把棱镜中的折线光路展开。全反射只把光轴转过 $90^\circ$，不改变会聚性质；棱镜内从上表面经斜面到右表面的轴向路程为 $3+3=6\,\mathrm{cm}$。近轴传播时用约化倾角 $p=n\theta$，高度增量为 $\Delta y=dp/n$，所以该玻璃段的等效空气传播长度是 $6/n=4\,\mathrm{cm}$，不是相位光程 $6n$。故从棱镜右面看，物的等效距离为 $6+4=10\,\mathrm{cm}$；再加棱镜到凸透镜的 $10\,\mathrm{cm}$，有 $u_1=20\,\mathrm{cm}$。

凸透镜：
\[
 v_1=\frac{10\times20}{20-10}=20\,\mathrm{cm},\qquad m_1=-1.
\]
凹透镜位于凸透镜右侧 $15\,\mathrm{cm}$，所以到达它的光原本还要在其右侧 $5\,\mathrm{cm}$ 会聚，即 $u_2=-5\,\mathrm{cm}$。取 $f_2=-10\,\mathrm{cm}$：
\[
 v_2=\frac{(-10)(-5)}{-5-(-10)}=10\,\mathrm{cm},\qquad m_2=-\frac{v_2}{u_2}=2.
\]
因此最终形成在凹透镜右侧
\[
 \boxed{10\,\mathrm{cm}}
\]
处的实像，总横向放大率
\[
 \boxed{m=m_1m_2=-2,\qquad |M|=2}.
\]
负号表示相对棱镜展开后的等效物方向倒立；折回实际空间时还应连同棱镜造成的 $90^\circ$ 方位旋转理解。
\end{xesolution}


\subsection{近轴矩阵方法}


在近轴条件下，用列向量 $(y,p)^{\mathsf T}$ 表示一条光线，其中 $y$ 向上为正，$\theta$ 从向右光轴逆时针量为正，$p=n\theta$ 是约化倾角。这里不用 $u$，以免与物距混淆。
折射率为 $n$ 的均匀介质中传播距离 $d$，以及通过光焦度为 $\Phi$ 的薄透镜，分别对应
\begin{equation}
 \mathbf T(d,n)=
 \begin{pmatrix}1&d/n\\0&1\end{pmatrix},
 \qquad
 \mathbf L(\Phi)=
 \begin{pmatrix}1&0\\-\Phi&1\end{pmatrix}.
\end{equation}
若透镜两侧处于相同折射率 $n$ 的介质中，像方焦距为 $f$，则 $\Phi=n/f$；只有空气中才近似为 $1/f$。各矩阵左乘光线列向量，\textbf{先经过的元件写在最右侧}。例如先传播 $d_1$、过透镜、再传播 $d_2$，总矩阵是 $\mathbf T(d_2,n)\mathbf L(\Phi)\mathbf T(d_1,n)$。一般写为
\begin{equation}
 \begin{pmatrix}y'\\p'\end{pmatrix}
 =\begin{pmatrix}A&B\\C&D\end{pmatrix}
 \begin{pmatrix}y\\p\end{pmatrix},
 \qquad AD-BC=1.
\end{equation}
若固定物平面上的同一点发出不同倾角的光线，在像平面仍交于同一点，则输出高度不能依赖初始倾角，成像条件就是 $B=0$，
横向放大率为 $A$。若平行入射光经过系统后仍保持平行，系统为无焦系统（afocal system），相应条件为 $C=0$。
例如单薄透镜的物平面到像平面矩阵为 $\mathbf T(v,n)\mathbf L(n/f)\mathbf T(u,n)$。乘出右上元为
\[
 B=\frac{u+v-uv/f}{n}.
\]
所以 $B=0$ 恰好恢复 $1/u+1/v=1/f$；左上元 $A=1-v/f=-v/u$ 也恢复带符号倍率。$B$ 的量纲为长度，$C$ 为长度倒数，$A,D$ 无量纲，行列式为一。厚透镜和多透镜组的主平面、焦点与节点，都可从同一矩阵分解得到；这避免了每增加一个折射面就重新建立一套符号法则。



\section{眼睛、视角与放大镜}



把人眼粗略看成空气到折射率 $n=4/3$、曲率半径 $R=5.56\,\mathrm{mm}$ 的单球面系统。远处物体 $u\to\infty$ 时，平行光的像距
\begin{equation}
 f=\frac{nR}{n-1}\approx2.22\,\mathrm{cm},
\end{equation}
与角膜到视网膜的尺度接近。真实人眼还包含角膜、房水、晶状体和玻璃体；水晶体由睫状肌调节曲率，使不同物距的像仍落在固定的视网膜上。水下看不清物体，是因为水的折射率约 $1.33$，与角膜约 $1.376$ 过于接近，角膜折射能力显著下降。

动物眼大体可分为两种工作思路：人眼一类用折射系统在感光面上形成倒立实像；昆虫复眼则由许多小眼单元分别接收不同方向的光。对人眼而言，视网膜位置几乎固定，调焦不能靠前后移动“屏”，只能由睫状肌改变晶状体曲率和等效焦距。看远处时睫状肌较放松；物体移近时，为保持同一像距，系统光焦度必须增大。晶状体形变能力有限，于是出现近点。图~\ref{ch28:fig:s43-20-3} 中角膜、晶状体、视网膜以及光线最终会聚的位置，是理解这一调节链条所必需的四项信息。

\EyePanel

能清楚聚焦的最近点称为近点。青年人的典型明视距离取
\begin{equation}
 L_0=25\,\mathrm{cm}.
\end{equation}
它是目视仪器设计的重要尺度。近点随年龄和屈光状态变化；近视者的远点有限，远视或老视者的近点则较远。$20\,\mathrm{cm}$ 与 $30\,\mathrm{cm}$ 只能作直观示例；实际诊断还需结合远点和屈光度，不能只凭一个距离。

正常眼完全放松时，远处平行光应落在视网膜上。近视眼的系统光焦度相对过强或眼轴偏长，平行光在视网膜前会聚，需要负透镜先使光束稍发散；远视眼或部分老视状态则需正透镜提供额外会聚。显微镜、望远镜的最终像通常设计在无穷远附近，让眼睛放松观察；若使用者近视且不戴镜，调节目镜会把最终虚像移到其有限远点。后面的显微镜例题正是通过“最终虚像应落在哪里”判断镜筒移动方向。



人眼感受到的“大小”由视角决定。尺度 $H$、距离 $L$ 的小物体张角 $\theta\simeq H/L$。光学仪器的角放大率定义为使用仪器后的视角与裸眼基准视角之比。

放大镜是一枚会聚透镜。把物放在焦点以内，透镜形成正立放大的虚像，眼睛靠近透镜观察。若物接近焦点，出射光近似平行，视角 $\theta\simeq H/f$；裸眼把物放在明视距离 $L_0$ 时，$\theta_0\simeq H/L_0$，所以放松眼调焦的角放大率为
\begin{equation}
 M=\frac{\theta}{\theta_0}=\frac{L_0}{f}.
\end{equation}
放大镜的本质，是让眼睛能在小于明视距离的位置清楚观察物体。望远镜、显微镜等目视系统的目镜也起同样作用：它观察前级系统形成的中间像，并进一步放大视角。

若不要求眼睛放松，而把最终虚像调到近点 $v=-L_0$，薄透镜公式给出 $1/u=1/f+1/L_0$。眼睛紧贴透镜时 $\theta\simeq H/u$，相对裸眼近点视角 $H/L_0$ 的角放大率为
\[
 M_{\rm near}=\frac{L_0}{u}=1+\frac{L_0}{f}.
\]
因此教材中常见的 $L_0/f$ 与 $1+L_0/f$ 对应两种不同调焦约定，并不矛盾。前者观察舒适，后者放大略大却更易疲劳。图~\ref{ch28:fig:s43-20-7} 的虚线反向延长线必须交于物体同侧的虚像，不能把放大镜误画成在屏上形成实像。

\MagnifierDiagram



前面的放大镜成像关系还可以统一为目视仪器的基准：这里正常眼明视距离记为 $D_0=L_0$，以免与口径 $D$ 混淆。光学仪器的“放大”通常指视角增大，而不是在屏上形成更大的实像。放大镜把物体放在焦点附近，
使眼睛在较舒适状态下观察放大的虚像。最终像在无穷远时角放大率
\begin{equation}
 M_\infty=\frac {D_0}f;
\end{equation}
最终像在近点时
\begin{equation}
 M_{D_0}=1+\frac {D_0}f.
\end{equation}
两式对应不同调焦条件，不能混用。



\subsection{角放大率的比较对象}
放大镜产生正立放大的虚像，但“虚像高变大了多少”并不是眼睛感到放大多少的完整答案。眼睛接收的是方向范围，即物体或像对眼睛张开的视角。比较时应固定同一物体，并明确裸眼基准通常取明视距离 $D$，于是裸眼视角为 $\alpha_0\simeq h/D$。

眼睛紧靠薄透镜、最终像在无穷远时，物位于前焦面，出射束视角为 $\alpha\simeq h/f$，所以 $M=\alpha/\alpha_0=D/f$。眼睛放松与物在焦面是同一个光路条件的两种表述，不能同时把物距取为任意值。

若最终虚像在距离透镜约 $D$ 的近点，取 $v=-D$，成像公式给出
\[
 \frac1u=\frac1f+\frac1D.
\]
眼睛靠近透镜时 $\alpha\simeq h/u$，因此 $M=D/u=1+D/f$。两个常用公式不是互相矛盾的近似，而是对应不同调焦状态。眼睛与透镜间距离不可忽略时，近点应从眼睛量起，须先把这一距离加入几何关系后再求视角。

视力分辨、角放大率和视场也不是一回事。放大率可以增加，但像的细节若已被衍射或像差模糊，继续放大只会放大模糊斑。有限镜片口径还会限制能同时看到的物体范围；判断视场需要画边缘光线，不能只看过光心的一条线。


\section{显微镜与望远镜}

\subsection{显微镜的中间像与调焦}


复合显微镜由短焦物镜和目镜组成。物镜先形成放大的倒立实像，目镜再把中间像作为物体形成放大虚像。图\ref{ch28:fig:microscope}中，光学筒长 $L$ 定义为物镜后焦点 $F'_o$ 到目镜前焦面（中间像平面）的轴向距离，不是两个镜片之间的机械距离。物镜像距为 $v_o=f_o+L$，从薄透镜公式有 $m_o=1-v_o/f_o=-L/f_o$。目镜角放大率为 $D_0/f_e$，因此最终像在无穷远的近轴角放大率为
\begin{equation}
 M\simeq-\frac{L}{f_o}\frac{D_0}{f_e}.
\end{equation}
负号表示最终像相对原物倒立。提高几何放大率不能突破衍射分辨极限；有用放大应与物镜数值孔径相匹配。

\begin{center}\begin{minipage}{\linewidth}\centering
\begin{tikzpicture}[>=Latex,font=\small,every node/.style={fill=none}]
\draw[gray,thin,->] (-2.0000000,0.0000000)--(8.2000000,0.0000000);
\draw[DeepBlue,thick,<->] (0.0000000,-2.1000000)--(0.0000000,2.1000000);
\draw[DeepBlue,thick,<->] (6.2000000,-2.1000000)--(6.2000000,2.1000000);
\node[above,fill=none,inner sep=2pt] at (0.0000000,2.2000000) {物镜};
\node[above,fill=none,inner sep=2pt] at (6.2000000,2.2000000) {目镜};
\draw[black,thick,->] (-1.2500000,0.0000000)--(-1.2500000,0.3000000);
\fill (5.0000000,-1.2000000) circle (1.2pt);
\node[below,fill=none,inner sep=2pt] at (5.0000000,-1.2000000) {};
\draw[gray,dashed,thin] (5.0000000,-1.2000000)--(4.5000000,1.1500000);
\node[above,fill=none,inner sep=2pt] at (4.5000000,1.2000000) {中间实像};
\draw[gray,dashed,thin] (5.0000000,-1.5000000)--(5.0000000,1.4000000);
\draw[DeepBlue,thick,->] (-1.2500000,0.3000000)--(0.0000000,0.6500000);
\draw[DeepBlue,thick,->] (0.0000000,0.6500000)--(6.2000000,-1.6440000);
\draw[DeepBlue,thick,->] (6.2000000,-1.6440000)--(8.0000000,0.1560000);
\draw[DeepBlue,thick,->] (-1.2500000,0.3000000)--(0.0000000,-0.1500000);
\draw[DeepBlue,thick,->] (0.0000000,-0.1500000)--(6.2000000,-1.4520000);
\draw[DeepBlue,thick,->] (6.2000000,-1.4520000)--(8.0000000,0.3480000);
\draw[gray,<->] (1.0000000,-1.8500000)--(5.0000000,-1.8500000);
\node[below,fill=none,inner sep=2pt] at (3.0000000,-1.8500000) {$L$};
\fill (1.0000000,0.0000000) circle (1.2pt);
\node[below,fill=none,inner sep=2pt] at (1.0000000,0.0000000) {};
\draw[gray,dashed,thin] (1.0000000,0.0000000)--(1.0000000,-0.9500000);
\node[below,fill=none,inner sep=2pt] at (1.0000000,-1.0000000) {$F'_o$};
\draw[gray,<->] (5.0000000,-2.5000000)--(6.2000000,-2.5000000);
\node[below,fill=none,inner sep=2pt] at (5.6000000,-2.5000000) {$f_e$};
\end{tikzpicture}
\figcap{显微镜的中间像在目镜前焦面；两透镜之间不发生折转。$L$ 从物镜后焦点量到中间像平面。\label{ch28:fig:microscope}}
\end{minipage}\end{center}



\Needspace{8\baselineskip}\begin{origexamplebox}{【小明逞强】}\label{ch28:ex:s43-20-4}
近视的小明不戴眼镜，与视力正常的小红使用同一显微镜。小红调清楚后，小明应上调还是下调镜筒才能看清？
\end{origexamplebox}\par\nopagebreak[3]
\begin{xesolution}
近视眼不戴镜时远点较近，需要显微镜给出较近的最终虚像。目镜焦距为 $f_e$，中间像在目镜焦内、物距为 $u_e$，虚像距离大小 $|v_e|=f_eu_e/(f_e-u_e)$；要使 $|v_e|$ 变小，应使 $u_e$ 变小。

下调镜筒使物镜靠近标本。对物镜而言，物距更接近焦距，所成中间实像距物镜变大；目镜与物镜间距固定，因此中间像到目镜的距离 $u_e$ 变小，最终虚像变近。故应
\[
 \boxed{\text{下调镜筒（使物镜略靠近标本）}}.
\]
把调节方向写成导数也能自检。镜筒向标本移动 $\varepsilon>0$ 时，物镜物距变化 $\Delta u_o=-\varepsilon$，而
\[
 \frac{\dd v_o}{\dd u_o}=-\frac{f_o^2}{(u_o-f_o)^2}<0,
 \qquad u_e=\ell-v_o,
\]
其中 $\ell$ 为固定的两镜间距。因此 $v_o$ 增大而 $u_e$ 减小。若小红原来按放松眼调到无穷远，而小明远点距目镜为 $F_R$（忽略眼距），所需 $u_e=f_eF_R/(F_R+f_e)<f_e$，与上述方向一致。这里比较两人都采用舒适调焦的状态；若眼睛主动调节或佩戴矫正镜，应重新指定最终像的位置。
实际操作应缓慢并观察侧面距离，避免物镜碰到玻片。
\end{xesolution}
\subsection{望远镜的无焦系统与有效倍率}


开普勒望远镜（Keplerian telescope）用正焦距物镜和正焦距目镜。远物经物镜在焦平面形成实像，
调节为无焦系统时两透镜焦平面重合，角放大率
\begin{equation}
 M=-\frac{f_o}{f_e},
\end{equation}
镜筒长约 $f_o+f_e$。图\ref{ch28:fig:telescope}中，入射倾角 $\theta$ 对应中间像高 $f_o\theta$；此像在目镜前焦面，所以出射倾角为 $-f_o\theta/f_e$，两条出射光平行。伽利略望远镜（Galilean telescope）使用负焦距目镜，成正立虚像，镜筒更短，但视场受限。

\begin{center}\begin{minipage}{\linewidth}\centering
\begin{tikzpicture}[>=Latex,font=\small,every node/.style={fill=none}]
\draw[gray,thin,->] (-2.0000000,0.0000000)--(7.5000000,0.0000000);
\draw[DeepBlue,thick,<->] (0.0000000,-2.1000000)--(0.0000000,2.1000000);
\draw[DeepBlue,thick,<->] (5.5000000,-2.1000000)--(5.5000000,2.1000000);
\node[above,fill=none,inner sep=2pt] at (0.0000000,2.2000000) {物镜};
\node[above,fill=none,inner sep=2pt] at (5.5000000,2.2000000) {目镜};
\fill (4.0000000,-0.4800000) circle (1.2pt);
\node[below,fill=none,inner sep=2pt] at (4.0000000,-0.4800000) {};
\draw[gray,dashed,thin] (4.0000000,-0.4800000)--(3.5000000,1.1500000);
\node[above,fill=none,inner sep=2pt] at (3.5000000,1.2000000) {中间像点};
\draw[gray,dashed,thin] (4.0000000,-1.5000000)--(4.0000000,1.4000000);
\draw[DeepBlue,thick,->] (-1.6000000,0.8420000)--(0.0000000,0.6500000);
\draw[DeepBlue,thick,->] (0.0000000,0.6500000)--(5.5000000,-0.9037500);
\draw[DeepBlue,thick,->] (5.5000000,-0.9037500)--(7.3000000,-0.3277500);
\draw[DeepBlue,thick,->] (-1.6000000,0.0420000)--(0.0000000,-0.1500000);
\draw[DeepBlue,thick,->] (0.0000000,-0.1500000)--(5.5000000,-0.6037500);
\draw[DeepBlue,thick,->] (5.5000000,-0.6037500)--(7.3000000,-0.0277500);
\draw[gray,<->] (4.0000000,-2.5000000)--(5.5000000,-2.5000000);
\node[below,fill=none,inner sep=2pt] at (4.7500000,-2.5000000) {$f_e$};
\end{tikzpicture}
\figcap{无焦望远镜：公共焦面处光线相交但不折转；目镜后的两条出射线严格平行。\label{ch28:fig:telescope}}
\end{minipage}\end{center}

望远镜分辨本领取决于物镜口径而非焦距。圆孔的瑞利判据（Rayleigh criterion）为
\begin{equation}
 \theta_{\min}=1.22\frac\lambda D.
\end{equation}
增大口径既提高集光能力，也提高角分辨率；仅增大目镜倍率不会恢复物镜没有记录的细节。

\Needspace{8\baselineskip}\begin{origexamplebox}{望远镜口径的分辨要求}
两颗星的角距离为 $0.50''$，用波长 $550\,\mathrm{nm}$ 的光观察。按瑞利判据估算望远镜最小口径。
\end{origexamplebox}\par\nopagebreak[3]
\begin{xesolution}
先把角秒换成弧度：
\[
 \theta=0.50\frac{\pi}{180\times3600}
 \approx2.42\times10^{-6}\,\mathrm{rad}.
\]
所以
\[
 D_{\min}=1.22\frac\lambda\theta
 \approx1.22\frac{550\times10^{-9}}{2.42\times10^{-6}}
 \approx0.277\,\mathrm m.
\]
即 $\boxed{D_{\min}\approx28\,\mathrm{cm}}$。真实观测还会受大气视宁度、像差和信噪比限制。
\end{xesolution}



\Needspace{8\baselineskip}\begin{origexamplebox}{【天文望远镜】}\label{ch28:ex:s43-20-5}
天文台常见双凸式透射望远镜。（1）为什么不采用凹凸式（伽利略式）作主要天文望远镜？（2）怎样简易判断望远镜的性能高低？
\end{origexamplebox}\par\nopagebreak[3]
\begin{xesolution}
双凸的开普勒结构由正物镜先形成实的中间像，再由正目镜放大。它虽然给出倒像，但视场较大，容易在中间焦面放置十字丝、相机、光谱狭缝等。凹凸伽利略结构成正像且镜筒短，却没有可接收的中间实像，出瞳位于系统内部，视场较小，不适合主要天文观测附件。

简易比较不能只看标称“倍数”。在大气视宁度和加工质量相当时，物镜有效口径 $D$ 决定集光能力 $\propto D^2$，并决定衍射角分辨率 $\theta\simeq1.22\lambda/D$；因此先看\textbf{有效口径和像质}，再看焦距、消色差、镀膜、支架稳定性和目镜配置。超过口径所允许的倍率只会把模糊斑放大。
\end{xesolution}

\section{仪器分辨本领与光能探测}



棱镜或光栅把不同波长分到不同方向，准直镜把入射狭缝变成平行光，成像镜把各方向重新聚焦到探测器。
平面光栅正入射时 $d\sin\theta_m=m\lambda$。若 $N$ 条刻线均匀相干照明，忽略像差，在非零级次 $m$ 的主极大附近，用瑞利判据得到理论分辨本领
\begin{equation}
 \mathcal R=\frac\lambda{\Delta\lambda}=|m|N.
\end{equation}
这是因为两波长的主极大角间距为 $|m|\Delta\lambda/(d\cos\theta_m)$，单条谱线主极大到第一极小的角距为 $\lambda/(Nd\cos\theta_m)$；令二者相等即可消去共同几何因子。$m=0$ 没有通常意义的角分光能力，不能作为上述分辨公式的工作级次；接近掠射时也需谨慎使用局部小角展开。狭缝宽度、像差和探测器采样会降低实际分辨率；狭缝无限变窄又会显著损失光通量。



辐射度学（radiometry）按电磁辐射携带的功率计量；光度学（photometry）还要按明视觉光谱响应 $V(\lambda)$ 加权。
若光谱辐射通量为 $\Phi_{e,\lambda}$，光通量为
\begin{equation}
 \Phi_v=K_m\int_0^\infty V(\lambda)\Phi_{e,\lambda}\,\dd\lambda,
 \qquad K_m\simeq683\,\mathrm{lm/W}.
\end{equation}
国际单位制用频率 $540\,\mathrm{THz}$ 的单色辐射把光视效能固定为 $683\,\mathrm{lm/W}$；其他波长即使辐射功率相同，
视觉亮度也会因 $V(\lambda)$ 不同而改变。常用光度量为
\begin{equation}
 I_v=\frac{\dd\Phi_v}{\dd\Omega},\qquad
 E_v=\frac{\dd\Phi_v}{\dd A},\qquad
 L_v=\frac{\dd^2\Phi_v}{\dd A\cos\theta\,\dd\Omega},
\end{equation}
分别称发光强度、照度和亮度，单位依次为坎德拉、勒克斯和坎德拉每平方米。光电二极管、CCD/CMOS 等探测器有各自的
光谱响应曲线；把“人眼看起来一样亮”直接当成“探测器接收功率相同”通常是错误的。

\tcbset{breakable/.default=true}\begin{starread}{三刺激值与 CIE 1931 色度图}
人眼明视觉主要由三类视锥细胞提供长、中、短波段响应。CIE 1931 标准把光谱 $P(\lambda)$ 映射为三刺激值
\[
 X=\int\bar x(\lambda)P(\lambda)\,\dd\lambda,\quad
 Y=\int\bar y(\lambda)P(\lambda)\,\dd\lambda,\quad
 Z=\int\bar z(\lambda)P(\lambda)\,\dd\lambda,
\]
再用 $x=X/(X+Y+Z)$、$y=Y/(X+Y+Z)$ 表示与总亮度分离的色度。两束光相加时，三刺激值线性相加；
在图\ref{ch28:fig:chromaticity}的色度平面上，二者混合色位于两点连线上。显示器选定三原色后，只能覆盖三原色三角形内部的颜色，并不能覆盖整个可见色域。

\begin{center}\begin{minipage}{\linewidth}\centering
\begin{tikzpicture}
\begin{axis}[width=.70\linewidth,height=6.0cm,axis lines=left,
  xlabel={$x$},ylabel={$y$},xmin=0,xmax=.80,ymin=0,ymax=.88,
  xtick={0,.2,.4,.6,.8},ytick={0,.2,.4,.6,.8},
  every axis label/.style={fill=none}]
  \addplot[curve,very thick] coordinates{
    (.174,.005)(.157,.018)(.124,.058)(.091,.133)(.045,.295)(.008,.538)
    (.014,.750)(.074,.834)(.155,.806)(.230,.754)(.302,.692)(.373,.625)
    (.444,.555)(.513,.487)(.575,.424)(.627,.373)(.666,.334)(.692,.308)
    (.708,.292)(.719,.281)(.726,.274)(.735,.265)(.174,.005)};
  \addplot[DeepBlue,thick] coordinates{(.64,.33)(.30,.60)(.15,.06)(.64,.33)};
  \addplot[only marks,mark=*,mark size=1.7pt] coordinates{(.333,.333)};
  \node[anchor=west] at (axis cs:.35,.35) {$E$};
  \node[DeepBlue] at (axis cs:.52,.10) {RGB 色域};
  \draw[->,thin] (axis cs:.49,.13)--(axis cs:.43,.30);
  \node at (axis cs:.20,.66) {可见色域};
\end{axis}
\end{tikzpicture}
\figcap{CIE 1931 色度示意：弯曲边界为近似光谱轨迹，两端的直线为紫线（非单色光）。E 为等能光谱白点 $(1/3,1/3)$，不等同于任意 RGB 标准的白点。\label{ch28:fig:chromaticity}}
\end{minipage}\end{center}
\end{starread}\tcbset{breakable/.default=false}

\tcbset{breakable/.default=true}\begin{starread}{光阑、光瞳、像差与展度守恒}
孔径光阑限制轴上光束锥角，视场光阑限制可见物方范围。孔径光阑经其前方光具组成的像称入射光瞳，经其后方光具组成的像称出射光瞳；
望远镜目镜应让出射光瞳与人眼瞳孔匹配。球差、彗差、像散、场曲、畸变和色差具有不同的孔径与视场依赖，
不能只用“焦点不准”概括。

无损系统中的微分展度
\[
 \dd G=n^2\cos\theta\,\dd A\,\dd\Omega
\]
沿光路守恒。这意味着被动光学系统不能同时把光束横截面积和发散立体角任意压缩；提高照度、扩大视场、增大口径和保持像质之间
必然存在约束，而不是只要增加透镜片数就能全部兼得。
\end{starread}\tcbset{breakable/.default=false}

\endgroup
\subsection{将放大、分辨和亮度分开判断}
用望远镜观测两个角距为 $\Delta\theta$ 的远处点源时，物镜首先把它们形成相邻的衍射斑。口径为 $D_o$ 的圆形物镜，其瑞利判据为 $\Delta\theta\gtrsim1.22\lambda/D_o$。目镜只是把已经形成的像以更大的视角交给眼睛；若两斑在物镜焦面上不能分开，仅更换更短焦距的目镜不能恢复丢失的细节。

相机还需比较衍射斑和像素。焦距为 $f$ 时，艾里斑第一暗环半径约为 $1.22\lambda f/D_o$。减小孔径可以改善部分几何像差和增加景深，却会增大这个衍射尺度。由此可见“缩光圈越小越清晰”并非普遍规律，清晰度取决于像差、衍射、采样和对焦误差中哪一项占主导。

最后检查能量量纲：光功率 $\Phi_e$ 的单位为瓦，照度型的辐照度 $E_e=\dd\Phi_e/\dd A$ 的单位为瓦每平方米，曝光量 $H_e=\int E_e\dd t$ 的单位为焦耳每平方米。像变大并不创造能量；在接收总功率近似不变时，同样的功率分布到更大像面积上，单位面积得到的能量反而减少。光度量还经过人眼光谱响应加权，不能把流明与瓦不加条件地互换。


% EMOpticsQuality20261023
% QualityRevision20261023
\begingroup\tcbset{breakable/.default=false}
\renewcommand{\Needspace}[1]{\par\begingroup
\dimen0=#1\relax\vskip0pt plus\dimen0\penalty-100
\vskip0pt plus-\dimen0\vskip\dimen0\penalty9999\vskip-\dimen0
\vskip0pt\endgroup}
\chapter{光的干涉}


\graphicspath{{../../../sample/xes-original/forty-fourth-section/}{../../../sample/xes-original/forty-fourth-section/assets/}}


\section{干涉现象与相干条件}


% 图契约：光路由界面、成像和反射定律确定；标签不覆盖图线。
% 光学纯几何，印刷宽度<=14.5cm，标签>=8pt；透明标签，克制配色。
% 草图：光源 -> 分束元件 -> 两条相干光路 -> 共同观察区。
% 实线是真实光路，虚线是延长线/法线；尺寸在光路外侧。
\newcommand{\InterferencePhenomenaPanel}{%
\begin{center}\begin{minipage}{\linewidth}\centering\begin{tikzpicture}
 \shade[left color=violet,right color=blue] (0,0) rectangle (2,.55);
 \shade[left color=blue,right color=green!65!black] (2,0) rectangle (4,.55);
 \shade[left color=green!65!black,right color=yellow] (4,0) rectangle (6,.55);
 \shade[left color=yellow,right color=red] (6,0) rectangle (9,.55);
 \draw[->] (0,-.35)--(9.5,-.35) node[right]{$\lambda_0$};
 \foreach \x/\v in {0/380,2/470,4/560,6/650,9/780}{\draw(\x,-.29)--(\x,-.42)node[below]{$\v$};}
 \node[above] at (4.5,.65){可见光连续谱（单位：$\mathrm{nm}$）};
\end{tikzpicture}\figcap{薄膜颜色取决于反射增强的波段；可见光谱连续，并非只有七条离散谱线。\label{ch29:fig:s44-21-2}}\end{minipage}\end{center}}

\newcommand{\YoungPanel}{%
\begin{center}\begin{minipage}{\linewidth}\centering\begin{tikzpicture}[scale=.85,>=Latex]
 \coordinate(S)at(-3,0);\coordinate(A)at(0,.65);\coordinate(B)at(0,-.65);\coordinate(P)at(5,1.45);
 \fill(S)circle(1.5pt)node[left]{$S$};
 \draw[very thick](0,-1.8)--(0,-.76) (0,-.54)--(0,.54) (0,.76)--(0,1.8);
 \draw[very thick](5,-1.8)--(5,2.1);\fill(P)circle(1.5pt)node[above left]{$P(L,y)$};
 \draw[DeepBlue,thick](S)--(A)--(P) (S)--(B)--(P);
 \node[above left]at(A){$S_1$};\node[below left]at(B){$S_2$};
 \draw[guide](-3.2,0)--(5.4,0);\node[below right]at(5,0){$O$};
 \draw[<->](-.45,-.65)--(-.45,.65)node[midway,left]{$d$};
 \draw[<->](0,-2.1)--(5,-2.1)node[midway,below]{$L$};
 \draw[<->](5.65,0)--(5.65,1.45)node[midway,right]{$y$};
 \node at(2.8,1.45){$r_1$};\node at(2.8,.22){$r_2$};
 \draw[guide](0,0)--(P);
 \draw[TextGray](1,0)arc(0:16.173:1);\node at(.65,.52){$\theta$};
 \draw[TextGray,thin](.85,.48)--(1.08,.15);
 \begin{scope}[shift={(0,-4.3)}]
  \draw[very thick](0,-.9)--(0,-.5) (0,-.3)--(0,.3) (0,.5)--(0,.9);
  \foreach \r in {.45,.85,1.25,1.65}{
   \draw[DeepBlue!75] (0,.4) ++(-65:\r) arc[start angle=-65,end angle=65,radius=\r];
   \draw[Blue!70] (0,-.4) ++(-65:\r) arc[start angle=-65,end angle=65,radius=\r];}
  \node[anchor=west,align=left]at(2.3,0){两个次级波面叠加\\实线弧表示同相位波面};
 \end{scope}
\end{tikzpicture}\figcap{双缝的相干照明、屏上距离与次级波面。$S$ 是初级狭缝；下图显示双缝后相互重叠的波面。\label{ch29:fig:s44-21-3}}\end{minipage}\end{center}}

\newcommand{\YoungPatternPanel}{%
\begin{center}\begin{minipage}{\linewidth}\centering\begin{tikzpicture}[x=.9cm,y=1cm]
 \foreach \x in {-5,-4.5,...,-1}{\fill[DeepBlue!80](\x,0)rectangle+(0.17,.8);}
 \node[below,align=center]at(-3,-.25){单色光\\等间距明暗条纹};
 \fill[black!75](.3,0)rectangle(6.1,.8);
 \fill[white](3.08,0)rectangle(3.32,.8);
 \foreach \s in {-1,1}{
  \fill[blue!80!white]({3.2+\s*.65-.08},0)rectangle+(.16,.8);
  \fill[green!80!white]({3.2+\s*.95-.08},0)rectangle+(.16,.8);
  \fill[red!85!white]({3.2+\s*1.25-.08},0)rectangle+(.16,.8);}
 \node[below,align=center]at(3.2,-.25){白光零级白纹\\同一级：紫蓝在内，红在外};
\end{tikzpicture}\figcap{单色光与白光双缝图样。右图只展开零级和两侧第一级的颜色次序，宽谱高阶条纹很快相互重叠。\label{ch29:fig:s44-21-4}}\end{minipage}\end{center}}

\newcommand{\FilmPhenomenaPanel}{%
\historyfigure{thin-film-scenes-20260923.png}{左侧泡膜和右侧油膜的彩纹。膜厚与观察方向改变上下表面反射光的相位差。\label{ch29:fig:s44-21-5}}}

\newcommand{\ThinFilmDiagram}{%
% A=(-1,.6), D=1.5, r=20deg, n=1.5；B,C,E由同一组参数确定。
\begin{center}\begin{minipage}{\linewidth}\centering\begin{tikzpicture}[>=Latex]
 \pgfmathsetmacro{\rr}{20}\pgfmathsetmacro{\ia}{asin(1.5*sin(\rr))}
 \coordinate(A)at(-1,.6);\coordinate(B)at({-1+1.5*tan(\rr)},-.9);
 \coordinate(C)at({-1+3*tan(\rr)},.6);
 \pgfmathsetmacro{\ae}{3*tan(\rr)*sin(\ia)}
 \coordinate(E)at({-1+\ae*sin(\ia)},{.6+\ae*cos(\ia)});
 \fill[Blue!10](-2,-.9)rectangle(3,.6);\draw[very thick](-2,.6)--(3,.6) (-2,-.9)--(3,-.9);
 \draw[DeepBlue,thick]({-1-2*sin(\ia)},{.6+2*cos(\ia)})--(A)--(B)--(C);
 \draw[DeepBlue,thick,->](A)--({-1+2.4*sin(\ia)},{.6+2.4*cos(\ia)});
 \draw[Blue,thick,->](C)--({-1+3*tan(\rr)+2.4*sin(\ia)},{.6+2.4*cos(\ia)});
 \draw[guide](A)--(-1,2.35) (A)--(-1,-1.15);\draw[guide](E)--(C);
 \draw[TextGray](-1,1.23)arc[start angle=90,end angle={90+\ia},radius=.63];
 \node at(-1.26,1.7){$i$};
 \draw[TextGray](-1,.05)arc[start angle=-90,end angle={-90+\rr},radius=.55];
 \node at(-.89,-.05){$r$};
 \node[above left]at(A){$A$};\node[below]at(B){$B$};\node[below right]at(C){$C$};
 \node[above left]at(E){$E$};\node at(1.9,-.16){$n$};
 \draw[<->](3.35,-.9)--(3.35,.6)node[midway,right]{$D$};
 \node[anchor=west,align=left]at(2.4,2.1){$CE\perp AE$\\共同出射波面};
\end{tikzpicture}\figcap{在共同波面 $CE$ 比较相位。两路分别走 $AE$ 与 $ABC$，传播光程差为 $n(AB+BC)-AE$；两束出射光平行。\label{ch29:fig:s44-21-6}}\end{minipage}\end{center}}

\newcommand{\MichelsonDiagram}{%
% 分光面正斜率：+x反射到+y，-x反射到-y；镜面就是折返点。
\begin{center}\begin{minipage}{\linewidth}\centering\begin{tikzpicture}[scale=.85,>=Latex]
 \draw[very thick](-.65,-.65)--(.65,.65);\node[below right]at(.55,-.4){分光镜 $M$};
 \draw[very thick](4,-.75)--(4,.75);\node[right]at(4,.8){$M_1$};
 \draw[very thick](-.75,3.4)--(.75,3.4);\node[right]at(.8,3.4){$M_2$};
 \draw[thick](-.4,1.15)--(.4,1.95) (-.47,1.22)--(.33,2.02);
 \node[anchor=west]at(.65,1.75){补偿板 $P$};
 \draw[DeepBlue,very thick](-3,0)--(4,0) (0,-2)--(0,3.4);
 \draw[->,DeepBlue](-2.4,0)--(-1.4,0);
 \draw[->,DeepBlue](1,0)--(2.2,0);\draw[->,DeepBlue](3.2,.13)--(2.3,.13);
 \draw[->,DeepBlue](0,.65)--(0,1.1);\draw[->,DeepBlue](.13,2.95)--(.13,2.25);
 \draw[->,DeepBlue](0,-1)--(0,-1.7);
 \node[left]at(-3,0){光源};\node[below]at(0,-2.2){探测器};
 \node[anchor=west,align=left]at(1.6,2.6){端镜移动 $x$\\往返光程变化 $2x$};
\end{tikzpicture}\figcap{迈克耳孙干涉仪。往返光路共线，短箭头表示往返方向；薄板近似下未画玻璃内的微小侧移。\label{ch29:fig:s44-21-7}}\end{minipage}\end{center}}

\newcommand{\MovingSourceDiagram}{%
% Independent panel: original physical paths and local coordinates retained.
\begin{center}\begin{minipage}{\linewidth}\centering
\begin{tikzpicture}[scale=.83,>=Latex]
  \coordinate(S)at(-2,0);\coordinate(T)at(-2,.8);\coordinate(P)at(3,-1.2);
  \draw[very thick](0,-1.6)--(0,-.5) (0,-.3)--(0,.3) (0,.5)--(0,1.6);
  \draw[very thick](3,-1.6)--(3,1.2);\draw[guide](-2.3,0)--(3.3,0);
  \draw[->,DeepBlue](S)--(T);\fill(T)circle(1.5pt)node[above]{$S_{\rm new}$};
  \draw[Blue,thick](T)--(0,.4)--(P) (T)--(0,-.4)--(P);
  \fill(P)circle(1.5pt)node[right]{$P_0$};\node[left]at(S){$S$};
  \node[below,align=center]at(.5,-1.95){源上移，零光程差点下移};
 \end{tikzpicture}
\figcap{光源上移使零光程差点下移，图中 $P_0$ 表示移动后的中央明纹位置。\label{ch29:fig:moving-source}}
\end{minipage}\end{center}
}

\newcommand{\FiniteLloydDiagram}{%
% Independent panel: original physical paths and local coordinates retained.
\begin{center}\begin{minipage}{\linewidth}\centering
\begin{tikzpicture}[scale=.83,>=Latex]
  \coordinate(S)at(0,.45);\coordinate(V)at(0,-.45);
  \draw[guide](-.2,0)--(6,0);\draw[line width=1.8pt](1.2,0)--(3.3,0);
  \draw[very thick](6,-.65)--(6,2.25);\fill(S)circle(1.5pt)node[above left]{$S$};
  \fill(V)circle(1.3pt)node[below left]{$S'$};
  \foreach \x in {1.2,3.3}{\pgfmathsetmacro{\yy}{6*.45/\x-.45}
   \draw[DeepBlue,thick](S)--(\x,0)--(6,\yy);\draw[guide](V)--(\x,0);}
  \draw[Blue,thick](S)--(6,1.05);\draw[<->](-.5,0)--(-.5,.45)node[midway,left]{$d$};
  \draw[<->](0,-.9)--(1.2,-.9)node[midway,below]{$b$};
  \draw[<->](1.2,-.9)--(3.3,-.9)node[midway,below]{$l$};
  \draw[<->](0,-1.6)--(6,-1.6)node[midway,below]{$L$};
  \draw[DeepBlue,line width=2pt](6,{6*.45/3.3-.45})--(6,{6*.45/1.2-.45});
  \node[anchor=west]at(6.25,1.35){重叠区};
 \end{tikzpicture}
\figcap{有限长劳埃德镜：只有粗线段是真实镜面；两端点反射的边界光线限定屏上的重叠区。\label{ch29:fig:s44-21-8}}
\end{minipage}\end{center}
}

\newcommand{\WedgeInterferenceDiagram}{%
\begin{center}\begin{minipage}{\linewidth}\centering\begin{tikzpicture}[scale=.9,>=Latex]
 \fill[Blue!10](0,0)--(0,2.8)--(1.4,2.8)--cycle;\draw[thick](0,0)--(0,2.8)--(1.4,2.8)--cycle;
 \draw[axis](-.3,0)--(2,0)node[below right]{$x$};\draw[axis](0,-.3)--(0,3.4)node[above]{$y$};
 \draw[very thick](2,-.1)--(2,-.06) (2,.06)--(2,.28) (2,.42)--(2,1.23) (2,1.37)--(2,2.43) (2,2.57)--(2,3);
 \draw[<->](3.3,-.1)--(3.3,3.0);\fill(5.6,1.4)circle(1.5pt)node[right]{$F$};
 \foreach \yy in {0,.35,1.3,2.5}{
  \draw[DeepBlue,thick](-1.6,\yy)--(3.3,\yy)--(5.6,1.4);}
 \draw[<->](0,3.08)--(1.4,3.08)node[midway,above]{$h$};
 \draw[<->](0,-.55)--(.65,-.55)node[midway,below]{$x_y$};
 \draw[guide](.65,1.3)--(.65,-.6);\node[left]at(-1.55,1.9){平行光};
 \node[below]at(2,-.85){开缝挡板};\node[below]at(3.6,-.45){理想透镜};
 \draw[TextGray](0,.7)arc[start angle=90,end angle=63.435,radius=.7];\node at(.29,.93){$\theta$};
 \begin{scope}[shift={(8,1)}]
  \draw[guide](.35,.3)--(.35,1.8)--(1.25,1.8);
  \draw[guide](.35,.3)--(.65,.3)--(.65,.8)--(.95,.8)--(.95,1.3)--(1.25,1.3)--(1.25,1.8);
  \foreach \x/\y in {0/0,0/1.5,.3/0,.3/.5,.6/.5,.6/1,.9/1,.9/1.5}{
   \draw[TextGray!60](\x,\y)--({\x+.35},{\y+.3});}
  \draw[thick](0,0)--(0,1.5)--(.9,1.5);
  \draw[thick](0,0)--(.3,0)--(.3,.5)--(.6,.5)--(.6,1)--(.9,1)--(.9,1.5);
  \draw[DeepBlue,->](-.7,.25)--(.3,.25)--(1.4,.25);
  \draw[DeepBlue,->](-.7,.75)--(.6,.75)--(1.4,.75);
  \node[below,align=center]at(.4,-.35){台阶沿 $z$ 延伸\\小面垂直于光线};
 \end{scope}
\end{tikzpicture}\figcap{劈尖与台阶模型。斜线给出厚度包络 $x_y=y\tan\theta$；右图说明微小台阶如何保持直线传播。$h$ 沿 $x$ 方向量取。\label{ch29:fig:wedge}}\end{minipage}\end{center}}

\newcommand{\PlaneWaveDiagram}{%
\begin{center}\begin{minipage}{\linewidth}\centering\begin{tikzpicture}[scale=.85,>=Latex]
 \draw[very thick](2,-1.8)--(2,1.8);\coordinate(P)at(2,0);
 \draw[guide](-2.6,0)--(2.5,0);\node[below]at(-1.7,-.1){屏法线};
 \draw[DeepBlue,thick,->](-2,{-4*tan(25)})--(P);
 \draw[Blue,thick,->](-2,{4*tan(20)})--(P);
 \draw(1.2,0)arc[start angle=180,end angle=205,radius=.8];
 \draw(1,0)arc[start angle=180,end angle=160,radius=1];
 \node at(.76,-.43){$\alpha$};\node at(.65,.4){$\beta$};
 \node[right]at(P){$P$};\node[above]at(2,1.9){$+y$};
 \node[anchor=west,align=left]at(3,.65){$k_{1y}=k\sin\alpha$\\$k_{2y}=-k\sin\beta$};
 \node[anchor=west]at(3,-.8){$\Delta y=\lambda/(\sin\alpha+\sin\beta)$};
\end{tikzpicture}\figcap{两波矢在屏面方向的分量符号相反。条纹间距由相位梯度决定，不是两条示意光线之间的距离。\label{ch29:fig:s44-21-9}}\end{minipage}\end{center}}

\newcommand{\BiprismDiagram}{%
\begin{center}\begin{minipage}{\linewidth}\centering\begin{tikzpicture}[scale=.9,>=Latex]
 \draw[thick,fill=Blue!8](0,0)--(0,1.5)--(.45,0)--(0,-1.5)--cycle;
 \draw[guide](-3.4,0)--(6.5,0);\fill(-3,0)circle(1.5pt)node[left]{$S$};
 \draw[very thick](6,-1.6)--(6,1.6);
 \foreach \s in {-1,1}{
  \draw[DeepBlue,thick](-3,0)--(0,{\s*.3})--(.352960321,{\s*.323465596})--(6,{\s*.041448685});
  \draw[guide](-3,{\s*.490914661})--(.352960321,{\s*.323465596});
  \fill(-3,{\s*.490914661})circle(1.2pt);}
 \node[left]at(-3,.62){$S_1'$};\node[left]at(-3,-.62){$S_2'$};
 \draw[<->](-3,-1.9)--(0,-1.9)node[midway,below]{$L_1$};
 \draw[<->](0,-1.9)--(6,-1.9)node[midway,below]{$L_2$};
 \draw[thin](0,0)--(.45,0);
 \draw[thin](0,1.05)arc(-90:-73.3:.45);\node[left]at(0,1.12){$\alpha$};
 \draw[thin](.16,-.75)--(1.1,-1.05)node[right]{$n$};
 \draw[fill=Blue!8](1.4,.02)rectangle(1.56,.65);
 \draw[thin](1.48,.65)--(2.3,1.1)node[right]{膜 $n',t$};
 \node[above,align=center]at(2.5,1.85){上、下棱镜使光偏向公共底面\\出射线反向延长确定虚源};
\end{tikzpicture}\figcap{菲涅耳双棱镜。上、下两半在公共底面相接；光线在界面折转，虚线与出射线共线。膜仅遮住上半束，忽略其横向位移而计入相位延迟；虚源公式采用薄棱镜近轴近似。\label{ch29:fig:s44-21-10}}\end{minipage}\end{center}}

\newcommand{\DoubleMirrorDiagram}{%
% alpha=12deg,R=2,L=5；两镜方向-51,-39deg；几何角度夸大。
\begin{center}\begin{minipage}{\linewidth}\centering\begin{tikzpicture}[scale=.85,>=Latex]
 \coordinate(O)at(0,0);\coordinate(S)at(0,2);
 \coordinate(Vone)at({-2*cos(12)},{2*sin(12)});\coordinate(Vtwo)at({-2*cos(12)},{-2*sin(12)});
 \draw[very thick](-1.4,{1.4*tan(51)})--(O)--(1.65,{-1.65*tan(39)});
 \draw[guide](-1.96,0)--(5.4,0);\draw[very thick](5,-1.7)--(5,2.2);
 \draw[DeepBlue,line width=2pt](5,{-5*tan(12)})--(5,{5*tan(12)});
 \foreach \V in {Vone,Vtwo}{\draw[guide](\V)--(O);}
 \draw[guide](O)--(5,{-5*tan(12)}) (O)--(5,{5*tan(12)});
 \coordinate(A)at(-.254342226,.314086491);\coordinate(B)at(.343717036,-.278336567);
 \draw[DeepBlue,thick](S)--(A)--(5,0) (S)--(B)--(5,0);
 \draw[guide](Vone)--(A) (Vtwo)--(B);
 \fill(S)circle(1.5pt)node[above]{$S$};\fill(O)circle(1.2pt)node[below left]{$O$};
 \fill(Vone)circle(1.2pt)node[left]{$S_1'$};\fill(Vtwo)circle(1.2pt)node[left]{$S_2'$};
 \node[left]at(-1.75,1.5){$M_1$};\node[below]at(1.45,-1.3){$M_2$};
 \draw[<->](5.45,{-5*tan(12)})--(5.45,{5*tan(12)})node[midway,right]{$W$};
 \draw[<->](0,-2)--(5,-2)node[midway,below]{$L$};
 \draw[<->]({-2*cos(12)},-2.7)--(5,-2.7)node[midway,below]{$L+R\cos\alpha$};
 \node[anchor=west]at(.65,2.75){$OS=R$，小夹角 $\alpha$ 已夸大};
\end{tikzpicture}\figcap{双面镜的两条反射光路在屏上重合。两虚源都在镜后，屏距从虚源平面量起；经过 $O$ 的边界射线限定 $W$。\label{ch29:fig:s44-21-11}}\end{minipage}\end{center}}

\newcommand{\BilletSourceDiagram}{%
\begin{center}\begin{minipage}{\linewidth}\centering\begin{tikzpicture}[scale=.7,>=Latex]
 \draw[guide](-3.2,0)--(8.3,0);\fill(-3,0)circle(1.5pt)node[left]{$S$};
 \foreach \s in {-1,1}{
  \draw[very thick](0,{\s*.3})--(0,{\s*1.8});\draw[thin](-.12,{\s*.3})--(.12,{\s*.3});
  \foreach \h in {.3,1.8}{\draw[DeepBlue,thick](-3,0)--(0,{\s*\h})--(4,{\s*.7})--(8,{\s*(1.4-\h)});}
  \fill(4,{\s*.7})circle(1.5pt);}
 \draw[very thick](8,-1.9)--(8,1.9);\draw[Blue,line width=2pt](8,-.4)--(8,.4);
 \node[above]at(4,.7){$S_1'$};\node[below]at(4,-.7){$S_2'$};
 \draw[<->](-.5,-.3)--(-.5,.3)node[midway,left]{$\delta$};
 \draw[<->](4.8,-.7)--(4.8,.7)node[midway,right]{$d$};
 \draw[<->](8.6,-.4)--(8.6,.4)node[midway,right]{$D$};
 \draw[<->](-3,-2.1)--(0,-2.1)node[midway,below]{$l_1$};
 \draw[<->](0,-2.1)--(4,-2.1)node[midway,below]{$l_2$};
 \draw[<->](4,-2.1)--(8,-2.1)node[midway,below]{$L-l_2$};
\end{tikzpicture}\figcap{解图：比耶对切透镜、两实像源与公共照明区。竖线是薄透镜所在平面，切口之间留有间隙。\label{ch29:fig:s44-ex-21-8}\label{ch29:fig:billet}}\end{minipage}\end{center}}
\newcommand{\MeslinSourceDiagram}{%
\begin{center}\begin{minipage}{\linewidth}\centering\begin{tikzpicture}[scale=.9,>=Latex]
 \draw[guide](-3.2,0)--(6,0);\fill(-3,0)circle(1.5pt)node[left]{$S$};
 \draw[very thick](0,0)--(0,1.5) (2,-1.5)--(2,0);
 \draw[DeepBlue,thick](-3,0)--(0,1.2)--(3,0)--(5,-.8);
 \draw[Blue,thick](-3,0)--(2,-1.2)--({29/7},0)--(5,.48);
 \draw[very thick]({25/7},-1.2)--({25/7},1.2);\fill(3,0)circle(1.5pt);\fill({29/7},0)circle(1.5pt);
 \node[above left]at(3,0){$S_1'$};\node[above right]at({29/7},0){$S_2'$};
 \node[above]at({25/7},1.3){中点屏};\node[below]at(1,-1.7){两半沿轴错开};
 \draw[<->](3,-1.5)--({29/7},-1.5)node[midway,below]{$a$};
 \begin{scope}[shift={(7.4,0)}]
  \foreach \r in {.4,.7,1.0}{\draw[DeepBlue](180:\r)arc[start angle=180,end angle=360,radius=\r];}
  \draw[guide](-1.1,0)--(1.1,0);\node[above]at(0,.2){屏的正视图};
 \end{scope}
\end{tikzpicture}\figcap{解图：梅斯林装置的两像都在光轴上，屏置于两像之间。重叠光只占一侧半屏，条纹为同心半圆弧；轴向间隔已夸大。\label{ch29:fig:s44-ex-21-9}\label{ch29:fig:meslin}}\end{minipage}\end{center}}

\newcommand{\NewtonRingsDiagram}{%
\begin{center}\begin{minipage}{\linewidth}\centering\begin{tikzpicture}[scale=.92,>=Latex]
 \begin{scope}
  \draw[very thick](-2.5,0)--(2.5,0);
  \draw[very thick,domain=-2.2:2.2,samples=101]plot(\x,{4-sqrt(16-\x*\x)});
  \draw[thick](-2.2,{4-sqrt(16-2.2*2.2)})--(-2.2,1.1)--(2.2,1.1)--(2.2,{4-sqrt(16-2.2*2.2)});
  \fill(0,0)circle(1.3pt)node[below left]{$O$};
  \draw[<->](1.8,0)--(1.8,{4-sqrt(16-1.8*1.8)})node[midway,right=4pt]{$t(r)$};
  \draw[<->](0,-.65)--(1.8,-.65)node[midway,below]{$r$};
  \node[above]at(0,1.15){平凸透镜};\node[below]at(-1.1,-1.15){底板（透明玻璃特例）};
 \end{scope}
 \begin{scope}[shift={(6,0)}]
  \foreach \m in {1,...,7}{\draw[DeepBlue](0,0)circle({.55*sqrt(\m)});}
  \fill(0,0)circle(.06);\node[below]at(0,-1.65){暗环满足 $r_m^2=m\lambda R$};
 \end{scope}
\end{tikzpicture}\figcap{牛顿环的接触几何与暗环半径。球面在 $O$ 与平板相切，$t(r)$ 两端分别落在平板和球面；环半径按平方级次生成。\label{ch29:fig:newton}}\end{minipage}\end{center}}


波动光学研究几何光学不能解释的光学性质，其中最典型的是干涉、衍射和偏振。光是电磁波，两束光在同一点相遇时，先叠加的是电场矢量，而探测到的能流却与电场平方的时间平均成正比。正如 $(a+b)^2$ 中除 $a^2+b^2$ 外还有交叉项，多个光场叠加后的总强度一般不等于各自强度的简单相加；能量在空间重新分布，就是干涉。

日常所见的肥皂泡、油膜和镜头镀膜颜色，都要求同一列光先被分开，沿不同光程传播后再在眼睛或探测器处重合。几何光学能确定每束光的路线，却不能判断重合后的明暗和颜色；这些必须由相位关系决定。

按“怎样把同一列光分开”可把常见装置分为两类。杨氏双缝、菲涅耳双棱镜和双面镜从同一波面的不同区域产生次级波，称为分波阵面干涉；薄膜和迈克耳孙干涉仪利用部分反射、部分透射把一束光的振幅分开，称为分振幅干涉。两类装置的几何结构不同，但计算核心完全相同：列出到观察点的全部传播光程和界面相变，求相位差，再由振幅叠加得到强度。图~\ref{ch29:fig:s44-21-1} 的同心环纹正是几何光路相近、相位结果却随膜厚和波长改变的典型证据。

\historyfigure{newton-rings-20260923.png}{牛顿环的俯视情景：同心环向外逐渐变密。定量几何见本章牛顿环解图。\label{ch29:fig:s44-21-1}}
\InterferencePhenomenaPanel
图~\ref{ch29:fig:s44-21-2} 把薄膜颜色与连续光谱并列：某处显色来自不同波长的增强或削弱，而不是薄膜只反射一种固定颜色。



\section{电场叠加、光程差与干涉强度}



设同频、同偏振的两束单色光在一点产生
\[
 E_1(t)=E_{01}\cos(\omega t+\varphi_1),\qquad
 E_2(t)=E_{02}\cos(\omega t+\varphi_2).
\]
电磁场满足线性叠加，$E=E_1+E_2$。对远长于光周期的探测时间取平均，得到
\begin{equation}
 I=I_1+I_2+2\sqrt{I_1I_2}\cos\Delta\varphi,\qquad
 \Delta\varphi=\varphi_1-\varphi_2.
\end{equation}
最后一项称为干涉项。若 $I_1=I_2=I_0$，则 $I_{\max}=4I_0$、$I_{\min}=0$；若振幅不同，即使相消位置也不能完全变黑。干涉可见度还要求频率相同或带宽足够窄、偏振方向不正交，且相位差在探测时间内稳定。两个互不相关的普通光源相位快速漂移，干涉项平均为零。

上述平均式可从瞬时场逐项看清。展开 $(E_1+E_2)^2$ 后，$\cos^2(\omega t+\varphi)$ 在许多光周期上的平均为 $1/2$，交叉项的平均为 $\tfrac12\cos(\varphi_1-\varphi_2)$；探测器来不及跟随约 $10^{15}\,\mathrm{Hz}$ 的电场振荡，只记录这个时间平均。若定义条纹可见度
\[
 \mathcal V=\frac{I_{\max}-I_{\min}}{I_{\max}+I_{\min}},
\]
则理想同偏振两束光有
\[
 \mathcal V=\frac{2\sqrt{I_1I_2}}{I_1+I_2}.
\]
所以两路亮度越接近，明暗反差越大；“完全相消”不仅要求相位差为 $\pi$，也要求两振幅相等。偏振正交或相位随机时，交叉项消失，即使两束光在空间重合也看不到稳定条纹。

能量守恒并未被干涉破坏。亮处可能达到单束光强的四倍，是因为能流被重新分配。对无吸收的完整干涉装置，把所有出射端口的功率相加，仍等于总入射功率；例如分光镜的一个端口变暗时，另一个端口可以变亮。不能任取一小块屏面就断言其干涉项积分为零。对于包含许多完整周期、包络变化缓慢的双缝图样，亮暗区的交叉项近似相互抵消。干涉改变能流分布，并不创造能量。

光程差 $\Delta\mathcal L$ 与相位差的关系为
\begin{equation}
 \Delta\varphi=\frac{2\pi}{\lambda_0}\Delta\mathcal L,\qquad
 \mathcal L=\int n\,\dd s,
\end{equation}
其中 $\lambda_0$ 是真空波长。除传播光程外，还要计入反射相位。对透明、无吸收介质的正入射或近正入射，按同一切向电场方向比较，从低折射率侧向高折射率侧反射有 $\pi$ 相变，反向反射无此相变。这就是常用的“半波损失”规则；全反射、金属反射及大角度不同偏振的情形应使用相应的菲涅尔反射系数，不能一概套用。不同路线中的传播相位和界面相位必须作为整体比较。

可见光只是电磁谱中的一小段，真空波长大致从紫光的 $400\,\mathrm{nm}$ 延伸到红光的 $700\,\mathrm{nm}$。同一频率的光进入折射率 $n$ 的介质后，频率不变、相速变为 $c/n$，介质波长变为 $\lambda_0/n$；所以几何长度 $l$ 中累积的相位等于真空中长度 $nl$ 的相位，这就是光程 $nl$ 的由来。比较两路相位时必须统一使用真空波长与光程，或分别使用各介质波长与几何程，不能把两套写法混在同一式中。



\subsection{为什么相加的是电场而不是亮度}
设两束同频、同偏振光到达同一点的电场为
\[
 E_1=A_1\cos\omega t,\qquad E_2=A_2\cos(\omega t-\delta).
\]
探测器响应于总电场平方的时间平均，而不是响应于某一时刻的电场。展开平方并分别平均：
\[
 \langle(E_1+E_2)^2\rangle
 =\frac{A_1^2}2+\frac{A_2^2}2+A_1A_2\cos\delta.
\]
因此强度除了 $I_1+I_2$，还多出含相位差的交叉项。相位稳定时，交叉项可以为正，也可以为负；相位在探测时间内均匀漂移时，$\langle\cos\delta\rangle=0$，才退化为直接相加强度。所谓“频率相同”只是必要条件之一，不能保证两个独立普通光源长期保持固定相位差。

完全相干且同偏振时，
\[
 I_{\max}=(\sqrt{I_1}+\sqrt{I_2})^2,\qquad
 I_{\min}=(\sqrt{I_1}-\sqrt{I_2})^2.
\]
只有 $I_1=I_2$ 时暗纹才降为零。改变两路强度比例会改变条纹反差，却不一定改变条纹位置。若两电场偏振方向成角 $\psi$，交叉项还要乘 $\cos\psi$；互相正交的偏振在不加检偏器时不产生这一标量干涉项。

\subsection{统一光程差和附加相位的记账方式}
指定第二路减第一路，记几何传播光程差为 $\Delta=\mathcal L_2-\mathcal L_1$。对给定时间因子约定，把两束在观察点的相位差写为
\[
 \delta=\frac{2\pi\Delta}{\lambda_0}+\delta_{\rm extra}.
\]
其中 $\delta_{\rm extra}$ 收集分束、反射以及源端的相对相位。必须整套使用同一约定；若调换两路编号，光程差和附加相位都应反号，最终强度不变。明纹条件是 $\delta=2m\pi$，不是在所有场合都直接写 $\Delta=m\lambda_0$。薄膜和劳埃德镜中的明暗翻转，正来自不能遗漏的附加相位。


\PlaneWaveDiagram
图~\ref{ch29:fig:s44-21-9} 中两波矢的横向分量相反，因此沿屏移动时两波相位之差持续变化；条纹周期应由这个相位差的变化率求出。

\Needspace{8\baselineskip}\begin{origexamplebox}{【平行光干涉】}\label{ch29:ex:s44-21-5}
两束波长均为 $\lambda$ 的平行光从屏法线两侧入射，与法线分别成 $\alpha,\beta$。干涉条纹间距为（
A. $\lambda/(\sin\alpha+\sin\beta)$；
B. $\lambda/(\cos\alpha+\cos\beta)$；
C. $\lambda/(\sin\alpha-\sin\beta)$；
D. $\lambda/|\cos\alpha-\cos\beta|$）。

\begin{center}\begin{minipage}{.96\linewidth}\centering
\begin{tikzpicture}[>=Latex,font=\small,every node/.style={fill=none},scale=.88]

\draw[thick] (3,-1.6)--(3,1.6) node[above]{屏};\draw[dashed,->] (0,0)--(4,0) node[right]{法线 $x$};
\draw[DeepBlue,thick,->] (0,-1.2)--(3,0);\draw[DeepBlue,thick,->] (0,1)--(3,0);
\node at (1.8,-.25){$\alpha$};\node at (1.7,.25){$\beta$};\draw[->] (3.6,.2)--(3.6,1.2) node[right]{$y$};

\end{tikzpicture}
\figcap{题图：两束等波长平面波位于屏法线的两侧，沿屏的正方向取 $+y$。\label{ch29:quality:q1}}
\end{minipage}\end{center}
\end{origexamplebox}\par\nopagebreak[3]
\begin{xesolution}
\begin{center}\begin{minipage}{.96\linewidth}\centering
\begin{tikzpicture}[>=Latex,font=\small,every node/.style={fill=none},scale=.88]

\draw[thick] (3,-1.6)--(3,1.6) node[above]{屏};\draw[dashed,->] (0,0)--(4,0) node[right]{法线 $x$};
\draw[DeepBlue,thick,->] (0,-1.2)--(3,0);\draw[DeepBlue,thick,->] (0,1)--(3,0);
\node at (1.8,-.25){$\alpha$};\node at (1.7,.25){$\beta$};\draw[->] (3.6,.2)--(3.6,1.2) node[right]{$y$};

\end{tikzpicture}
\figcap{解图：条纹间距取决于波矢差沿屏面的投影。\label{ch29:quality:a1}}
\end{minipage}\end{center}

沿屏面移动 $y$，两束平面波的相位分别改变 $k y\sin\alpha$ 与 $-k y\sin\beta$，故相位差梯度为
\[
 \frac{\dd\Delta\varphi}{\dd y}=k(\sin\alpha+\sin\beta).
\]
相位差增加 $2\pi$ 对应一条纹间距：
\[
 \boxed{\Delta y=\frac{\lambda}{\sin\alpha+\sin\beta}}.
\]
对应选项 A。若两光位于法线同一侧，正负号应按波矢在屏面上的投影重新取值，而不能机械使用相加式。

图~\ref{ch29:quality:a1} 中，下方入射波的波矢沿屏投影为 $+k\sin\alpha$，上方入射波的投影为 $-k\sin\beta$；相位差的屏面梯度因此是两者之差。
\end{xesolution}

\section{杨氏双缝与条纹移动}



杨氏实验用初级狭缝 $S$ 同相照亮 $S_1,S_2$，从同一波面上取出两个相干次级波源。两缝间距为 $d$，屏距为 $L$，屏上点 $P$ 与主轴夹角为 $\theta$。两路几何程差
\[
 \Delta r=r_2-r_1,\qquad
 \Delta\varphi=\frac{2\pi}{\lambda}\Delta r.
\]
远场条件 $L\gg d$ 且观察角不大时，
\begin{equation}
 \Delta r\simeq d\sin\theta\simeq d\frac yL.
\end{equation}
明纹满足 $\Delta r=m\lambda$，暗纹满足 $\Delta r=(m+1/2)\lambda$，相邻明纹间距
\begin{equation}
 \Delta y=\frac{\lambda L}{d}.
\end{equation}
改变两缝的相对初相位只会平移图样，不改变斜率 $\dd\Delta r/\dd y$，所以条纹间距不变。杨氏实验不仅确立了光的波动性，其“路径不可区分时振幅相加”的结构也成为量子物理的重要原型。

远场近似可以从精确距离展开得到。若两缝坐标为 $\pm d/2$，屏点坐标为 $(L,y)$，则
\[
 r_{1,2}=\sqrt{L^2+(y\mp d/2)^2}.
\]
在 $L$ 远大于 $|y|,d$ 时，用 $\sqrt{1+u}=1+u/2+O(u^2)$，有
\[
 r_1\simeq L+\frac{(y-d/2)^2}{2L},\qquad
 r_2\simeq L+\frac{(y+d/2)^2}{2L}.
\]
两式相减得 $r_2-r_1\simeq yd/L$，正负号与前面的定义一致：$y>0$ 时，下缝到观察点更远。再用 $y/L\simeq\sin\theta$ 得到上式。被略去的高阶项会改变远离主轴的条纹间距，因此 $\Delta y=\lambda L/d$ 不是任意近屏、大角度几何下的精确公式。

图~\ref{ch29:fig:s44-21-3} 中初级狭缝 $S$ 不是装饰。若直接用两个彼此独立的灯照亮 $S_1,S_2$，两灯相位各自漂移，屏上的交叉项会在观察时间内平均掉；让同一束光先通过 $S$，再同时照亮两缝，才能获得稳定的相位关系。若 $S$ 到两缝距离不同，只会在 $S_1,S_2$ 间加入固定初相位差 $\varphi_0$，屏上明纹位置变为
\[
 d\frac yL+\frac{\lambda\varphi_0}{2\pi}=m\lambda,
\]
条纹整体平移而间距仍为 $\lambda L/d$。

激光相干长度长、方向性强，图~\ref{ch29:fig:s44-21-4} 中条纹通常清晰且可延伸较远；普通宽带光源的不同波长分别形成不同间距的条纹，只有零光程差附近能够同时保持高可见度。源缝变宽还会把来自不同源点、位置略有错开的条纹叠加，从而降低反差。这些实验条件解释了为什么“有两条光路”只是必要条件，不是出现清晰干涉的充分条件。

\subsection{先写两段路，再使用小角近似}
双缝间距为 $d$，缝位于 $y=\pm d/2$，屏距为 $L$。屏上坐标为 $y$ 的点到下、上两缝的距离分别为
\[
 r_-=\sqrt{L^2+(y+d/2)^2},\qquad
 r_+=\sqrt{L^2+(y-d/2)^2}.
\]
有理化给出精确关系
\[
 r_--r_+=\frac{2dy}{r_-+r_+}.
\]
当 $|y|,d\ll L$ 时，分母可取 $2L$，于是 $\Delta\simeq dy/L$。相邻明纹光程差相差一个波长，因此 $\Delta y=\lambda L/d$。这里“等间距”是近轴屏坐标中的近似；大角度观察时应先求满足 $d\sin\theta=m\lambda$ 的角度，再用 $y=L\tan\theta$，不能把 $\tan\theta$ 与 $\sin\theta$ 混用。

源缝不在中垂线上时，光从源到双缝之前已经累积了路差，所以完整差值为
\[
 \Delta=(\mathcal L_{S\to -}-\mathcal L_{S\to +})+(r_--r_+).
\]
插片、移动源或改变一路介质，本质上都是改变第一项或增加一个常量程差。求条纹向哪侧移动时，先判断哪一路变长，再判断屏上往哪侧移动可以使该路变短以补偿，而不要只记“向薄片侧移动”的口令；若装置换了成像方向，口令可能失效。


\YoungPanel
\YoungPatternPanel
% Contract: two normalized grayscale interference strips; print width 12cm.
% Physical model: equal local beam intensities, I/(I1+I2)=1+V(x)cos(2pi*x).
% Top V=.98; bottom V=.98 exp[-(x/1.4)^2]; identical fringe spatial period.
% Smooth narrow-band illustrative coherence envelope, not a measured source spectrum.
% Layout: two 8cm strips stacked 1.6cm apart, descriptive text to their right.
% No labels over pixels; no white masks; no inferred color from a grayscale scan.
% ASCII: [high contrast many fringes] laser; [central clear, outer washed] ordinary.
\begin{center}\begin{minipage}{\linewidth}\centering
\begin{tikzpicture}[every node/.style={fill=none,font=\small},x=1cm,y=1cm]
 \foreach \row in {0,1}{
  \foreach \j in {0,...,319}{
   \pgfmathsetmacro{\xx}{-4+(\j+.5)/40}
   \pgfmathsetmacro{\vis}{ifthenelse(\row==0,.98,.98*exp(-(\xx/1.4)^2))}
   \pgfmathsetmacro{\shadevalue}{50*(1-\vis*cos(360*\xx))}
   \fill[black!\shadevalue]({-4+\j/40},{-1.6*\row})rectangle({-4+(\j+1)/40+.001},{.65-1.6*\row});
  }
 }
 \node[right,align=left,text width=4cm]at(4.3,.32){窄谱、良好相干照明\\较大范围仍有清晰条纹};
 \node[right,align=left,text width=4cm]at(4.3,-1.28){普通光源经狭缝处理\\中心附近条纹较清晰};
\end{tikzpicture}
\figcap{相同条纹间距下的反差比较。灰度只表示归一化强度，不用于比较两种光源的绝对亮度。\label{ch29:fig:coherence-contrast}}
\end{minipage}\end{center}

图~\ref{ch29:fig:coherence-contrast} 保留了另一层信息：普通光源经狭缝获得部分空间相干性后，也能产生干涉，并非“只有激光能干涉”。但光谱宽度及有效光源尺寸会限制反差；当两路光程差增大，不同频率产生的条纹逐渐错开，时间平均后外侧条纹变淡。图中用随离轴距离增大而减小的可见度示意这一过程。实际反差包络由光谱、光源尺寸和仪器决定，不必恰好是高斯形。白光的颜色次序则应另看图~\ref{ch29:fig:s44-21-4}，不能用颜色次序代替反差的比较。


\Needspace{8\baselineskip}\begin{origexamplebox}{【移动光源】}\label{ch29:ex:s44-21-1}
杨氏双缝实验中，单色点光源沿垂直于主轴方向向上移动一小段，屏上条纹如何变化？A. 上移、间距不变；B. 上移、间距变大；C. 下移、间距不变；D. 下移、间距变大。

\begin{center}\begin{minipage}{.96\linewidth}\centering
\begin{tikzpicture}[>=Latex,font=\small,every node/.style={fill=none},scale=.88]

\fill (-2,0) circle (1.7pt) node[left]{$S$};\draw[thick] (0,-1.3)--(0,-.45) (0,-.3)--(0,.3) (0,.45)--(0,1.3);
\fill (0,.375) circle (1.5pt) node[right]{$S_1$};\fill (0,-.375) circle (1.5pt) node[right]{$S_2$};
\draw[DeepBlue,->] (-2,0)--(0,.375);\draw[DeepBlue,->] (-2,0)--(0,-.375);
\draw[thick] (4,-1.3)--(4,1.3) node[above]{屏};\draw[dashed] (-2,0)--(4,0);
\draw[<->] (0,-1.6)--(4,-1.6) node[midway,below]{$L$};\draw[<->] (-.4,-.375)--(-.4,.375) node[midway,left]{$d$};
\draw[->] (-2,.15)--(-2,.9) node[above]{$\delta s$};
\end{tikzpicture}
\figcap{题图：点光源沿双缝连线方向上移，两缝位置与屏距保持不变。\label{ch29:quality:q2}}
\end{minipage}\end{center}
\end{origexamplebox}\par\nopagebreak[3]
\begin{xesolution}
\MovingSourceDiagram

光源上移后，到上缝的距离变短、到下缝的距离变长，两个次级光源获得一个固定初相位差。为了重新满足零光程差，屏上观察点必须向下移动，以用“下缝到屏更短”的几何程差补偿源端程差。由于光源移动不改变两缝间距 $d$、屏距 $L$ 和波长，$\Delta y=\lambda L/d$ 不变。
\[
 \boxed{\text{C：条纹向下移动，间距不变。}}
\]
进一步设光源到双缝面的距离为 $a$，上移量为 $s>0$，屏坐标 $y$ 向上为正。上、下缝分别在 $d/2,-d/2$，在 $s,d\ll a$、$y,d\ll L$ 的近轴条件下，两条完整路径之差为
\[
 \Delta\mathcal L=
 \left[\sqrt{a^2+(s+d/2)^2}-\sqrt{a^2+(s-d/2)^2}\right]
 +\left[\sqrt{L^2+(y+d/2)^2}-\sqrt{L^2+(y-d/2)^2}\right]
 \simeq\frac{ds}{a}+\frac{dy}{L}.
\]
中央明纹满足 $\Delta\mathcal L=0$，故 $y_0=-Ls/a<0$；相邻明纹满足程差相差 $\lambda$，仍有 $\Delta y=\lambda L/d$。图~\ref{ch29:fig:moving-source} 画出这种向下补偿的方向；该结论使用近轴近似。
\end{xesolution}

\Needspace{8\baselineskip}\begin{origexamplebox}{【白光干涉】}\label{ch29:ex:s44-21-11}
白光通过杨氏双缝后为什么出现彩色条纹且级数很少？为什么这种现象难观察？

\begin{center}\begin{minipage}{.96\linewidth}\centering
\begin{tikzpicture}[>=Latex,font=\small,every node/.style={fill=none},scale=.88]

\fill (-2,0) circle (1.7pt) node[left]{白光源 $S$};\draw[thick] (0,-1.3)--(0,-.45) (0,-.3)--(0,.3) (0,.45)--(0,1.3);
\fill (0,.375) circle (1.5pt) node[right]{$S_1$};\fill (0,-.375) circle (1.5pt) node[right]{$S_2$};
\draw[DeepBlue,->] (-2,0)--(0,.375);\draw[DeepBlue,->] (-2,0)--(0,-.375);
\draw[thick] (4,-1.3)--(4,1.3) node[above]{屏};\draw[dashed] (-2,0)--(4,0);
\draw[<->] (0,-1.6)--(4,-1.6) node[midway,below]{$L$};\draw[<->] (-.4,-.375)--(-.4,.375) node[midway,left]{$d$};

\end{tikzpicture}
\figcap{题图：宽光谱通过同一对双缝，各波长具有相同的几何光程差。\label{ch29:quality:q3}}
\end{minipage}\end{center}
\end{origexamplebox}\par\nopagebreak[3]
\begin{xesolution}
\begin{center}\begin{minipage}{.96\linewidth}\centering
\begin{tikzpicture}[>=Latex,font=\small,every node/.style={fill=none},scale=.88]

\fill (-2,0) circle (1.7pt) node[left]{白光源 $S$};\draw[thick] (0,-1.3)--(0,-.45) (0,-.3)--(0,.3) (0,.45)--(0,1.3);
\fill (0,.375) circle (1.5pt) node[right]{$S_1$};\fill (0,-.375) circle (1.5pt) node[right]{$S_2$};
\draw[DeepBlue,->] (-2,0)--(0,.375);\draw[DeepBlue,->] (-2,0)--(0,-.375);
\draw[thick] (4,-1.3)--(4,1.3) node[above]{屏};\draw[dashed] (-2,0)--(4,0);
\draw[<->] (0,-1.6)--(4,-1.6) node[midway,below]{$L$};\draw[<->] (-.4,-.375)--(-.4,.375) node[midway,left]{$d$};

\end{tikzpicture}
\figcap{解图：用同一双缝几何比较各波长的相位差及条纹尺度。\label{ch29:quality:a3}}
\end{minipage}\end{center}

中央零级处所有波长程差均为零，叠加成白纹。离开中心后，明纹位置 $y_m=m\lambda L/d$ 随波长不同：紫蓝光靠内、红光靠外，形成彩边。级次增大时，不同波长的不同级条纹互相覆盖，综合色趋于均匀，很快失去可见度。

普通白光的相干长度很短，只有两路光程差接近零时相位关系才稳定；双缝又很窄，亮度低，且源缝、双缝和平衡光路必须精确调整。所以白光条纹只在零光程差附近少量出现，远比激光双缝难观察。

图~\ref{ch29:quality:a3} 中，$d,L$ 对各色相同，而条纹间距正比于波长；零光程差处各色明纹重合，离中心后不同波长的级次逐渐错开并互相冲淡。
\end{xesolution}

\section{虚光源干涉与公共照明区}




\subsection{等效光源解决路差，边缘光线决定能否看见}
双棱镜、双面镜和对切透镜看似不同，都可以分成三步：先用反射或成像规律找到一对相干的等效源；再从它们到屏的距离计算相位差；最后追踪器件边缘所允许的真实光线，确定屏上两束共同到达的区域。

第二步给出的条纹是一个数学上的延拓图样。只有落在第三步公共照明区内的部分才真正可见。把虚源到屏的连线反向延长时，还必须检查它是否穿过对应镜面或透镜的有效口径；否则虽然能算出路差，实际却没有那一路光。这是题目中的有限镜长、半透镜边缘和挡光片不能略去的原因。

若近轴明纹中心为 $y_m=y_0+m\Delta y$，公共照明区为闭区间 $[a,b]$，那么允许级次满足
\[
 \frac{a-y_0}{\Delta y}\le m\le\frac{b-y_0}{\Delta y}.
\]
用上取整与下取整分别确定两端，得到
\begin{equation}\label{ch29:eq:depth-fringe-count}
 N=\max\!\left\{0,
 \left\lfloor\frac{b-y_0}{\Delta y}\right\rfloor
 -\left\lceil\frac{a-y_0}{\Delta y}\right\rceil+1\right\}.
\end{equation}
所以仅给宽度 $b-a$，而不给条纹相对窗口的位置，通常不能确定条数。若题目明确允许调节相位，才可另求平移图样时的最大数。边缘不计入、或要求亮带完整落入窗口时，端点不等式也要相应改变。

劳埃德镜只有一路经过反射，而双面镜的两路通常各反射一次；前者常需计入相对半波损失，后者在同样反射条件下这部分相位会抵消。等效源法只替代了路径几何，并不会自动替你处理反射相位。

\Needspace{8\baselineskip}\begin{origexamplebox}{【劳埃德镜（Lloyd's mirror）】}\label{ch29:ex:s44-21-3}
劳埃德镜长 $l=7.5\,\mathrm{cm}$；点光源 $S$ 离镜面 $d=0.15\,\mathrm{mm}$，距镜左端 $b=4.5\,\mathrm{cm}$；屏与镜垂直，距 $S$ 的水平距离 $L=1.2\,\mathrm m$。$\lambda=6.0\times10^{-7}\,\mathrm m$。求：（1）干涉区；（2）明纹间距；（3）该区第一条暗纹位置。

\begin{center}\begin{minipage}{.96\linewidth}\centering
\begin{tikzpicture}[>=Latex,font=\small,every node/.style={fill=none},scale=.88]

\draw[very thick] (0,0)--(2,0);\node[below] at (1,0){有限镜长 $l$};
\fill (-1,.55) circle (1.7pt) node[above]{$S$};\draw[thick] (5,-.1)--(5,2) node[above]{屏};
\draw[<->] (-1,.15)--(0,.15) node[midway,above]{$b$};\draw[<->] (-1,-.7)--(5,-.7) node[midway,below]{$L$};
\draw[<->] (-1.5,0)--(-1.5,.55) node[midway,left]{$d$};

\end{tikzpicture}
\figcap{题图：光源高于镜面 $d$，镜面只有从距光源水平位置 $b$ 到 $b+l$ 的有限长度。\label{ch29:quality:q4}}
\end{minipage}\end{center}
\end{origexamplebox}\par\nopagebreak[3]
\begin{xesolution}
\FiniteLloydDiagram
图~\ref{ch29:fig:s44-21-8} 中的虚源 $S'$、有限镜段与屏上边界确定下面的不等式。

反射光等效来自镜面对称虚源 $S'$，但反射带来 $\pi$ 相变。连接 $S'$ 与屏上高度 $y$ 的点，其与镜面的交点横坐标为 $x_0=Ld/(y+d)$。要落在有限镜段，须 $b\le x_0\le b+l$，故
\[
 \frac{Ld}{b+l}-d\le y\le\frac{Ld}{b}-d.
\]
代入数值得
\[
 \boxed{1.35\,\mathrm{mm}\le y\le3.85\,\mathrm{mm}}.
\]
两等效源间距为 $2d$，所以
\[
 \boxed{\Delta y=\frac{\lambda L}{2d}=2.40\,\mathrm{mm}}.
\]
由于反射半波损失，几何程差 $2dy/L=m\lambda$ 的位置是暗纹，即 $y_m=m\Delta y$。干涉区内最小的正整数 $m$ 为 $1$，故第一条暗纹在
\[
 \boxed{y=2.40\,\mathrm{mm}}
\]
处。若漏掉反射相变，会把明暗次序完全颠倒。
\end{xesolution}

% Biprism panel relocated to its solution.

\Needspace{22\baselineskip}\begin{origexamplebox}{【菲涅耳双棱镜】}\label{ch29:ex:s44-21-6}
双棱镜折射率为 $n$，小棱角为 $\alpha$；缝光源到棱镜距离为 $L_1$，屏到棱镜为 $L_2$。（1）求条纹间距；（2）用折射率 $n'$ 的肥皂膜遮住一半，零级条纹移动 $a$，求膜厚。

\begin{center}\begin{minipage}{.96\linewidth}\centering
\begin{tikzpicture}[>=Latex,font=\small,every node/.style={fill=none},scale=.88]

\fill (-2,0) circle (1.7pt) node[left]{$S$};\draw[thick] (0,-1.3)--(.2,0)--(0,1.3)--(-.2,0)--cycle;
\node[above] at (0,1.3){双棱镜 $n,\alpha$};\draw[thick] (4,-1.3)--(4,1.3) node[above]{屏};
\draw[dashed] (-2,0)--(4,0);\draw[<->] (-2,-1.7)--(0,-1.7) node[midway,below]{$L_1$};\draw[<->] (0,-1.7)--(4,-1.7) node[midway,below]{$L_2$};
\draw[gray] (.7,.2) rectangle (.85,1);\node[right] at (.85,.8){可加肥皂膜 $n'$};

\end{tikzpicture}
\figcap{题图：双棱镜的上下两半由同一缝光源照明，一半光路可插入薄膜。\label{ch29:quality:q5}}
\end{minipage}\end{center}
\end{origexamplebox}\par\nopagebreak[3]
\begin{xesolution}
\BiprismDiagram
图~\ref{ch29:fig:s44-21-10} 的两条反向延长线确定虚源；肥皂膜只覆盖其中一束，因而只给这一束增加光程。

每半棱镜使近轴光偏转 $\delta\simeq(n-1)\alpha$，两虚源间距
\[
 d=2L_1(n-1)\alpha.
\]
虚源到屏的距离近似 $L_1+L_2$，故
\[
 \boxed{\Delta y=\frac{\lambda(L_1+L_2)}{2L_1(n-1)\alpha}}.
\]
肥皂膜给一路增加光程 $(n'-1)t$。条纹移动 $a$ 对应的原几何程差为 $da/(L_1+L_2)$，补偿条件给出
\[
 \boxed{t=\frac{ad}{(L_1+L_2)(n'-1)}
 =\frac{2aL_1(n-1)\alpha}{(L_1+L_2)(n'-1)}}.
\]
移动方向由哪一半被遮决定；厚度取正值。
\end{xesolution}

\DoubleMirrorDiagram

\Needspace{8\baselineskip}\begin{origexamplebox}{【菲涅耳双面镜】}\label{ch29:ex:s44-21-7}
两镜相对共面位置的小夹角 $\alpha=20'$，缝光源到镜交线距离 $R=10\,\mathrm{cm}$，$\lambda=600\,\mathrm{nm}$，屏距镜交线 $L=210\,\mathrm{cm}$。采用图示对称配置：截面上 $O=(0,0)$、光源 $S=(0,R)$、屏为 $x=L$，两镜方向为 $-45^\circ\mp\alpha/2$，反射面朝向光源，尺寸足以接收公共照明区的光线；忽略边缘衍射，两次反射具有相同相位变化。（1）求间距与最多可见条纹数；（2）$R$ 加倍时怎样变化；（3）光源沿 $+x$ 横移 $\delta s$ 时怎样变化。

\begin{center}\begin{minipage}{.96\linewidth}\centering
\begin{tikzpicture}[>=Latex,font=\small,every node/.style={fill=none},scale=.88]

\draw[thick] (0,0)--(2,-2.2) (0,0)--(2.2,-2);
\fill (0,0) circle (1.5pt) node[left]{$O$};\fill (0,1.2) circle (1.7pt) node[above]{$S$};
\draw[dashed] (0,0)--(0,1.2) node[midway,left]{$R$};\draw[thick] (4,-1.6)--(4,1.6) node[above]{$x=L$};
\draw[dashed,->] (-.5,0)--(4.6,0) node[right]{$x$};\node[right] at (.65,-.65){小夹角 $\alpha$};

\end{tikzpicture}
\figcap{题图：双镜方向接近 $-45^\circ$，分别偏转 $\mp\alpha/2$；图中小夹角作了夸大。\label{ch29:quality:q6}}
\end{minipage}\end{center}
\end{origexamplebox}\par\nopagebreak[3]
\begin{xesolution}
\begin{center}\begin{minipage}{.96\linewidth}\centering
\begin{tikzpicture}[scale=.85,>=Latex]
 \coordinate(O)at(0,0);\coordinate(S)at(0,2);
 \coordinate(Vone)at({-2*cos(12)},{2*sin(12)});\coordinate(Vtwo)at({-2*cos(12)},{-2*sin(12)});
 \draw[very thick](-1.4,{1.4*tan(51)})--(O)--(1.65,{-1.65*tan(39)});
 \draw[guide](-1.96,0)--(5.4,0);\draw[very thick](5,-1.7)--(5,2.2);
 \draw[DeepBlue,line width=2pt](5,{-5*tan(12)})--(5,{5*tan(12)});
 \foreach \V in {Vone,Vtwo}{\draw[guide](\V)--(O);}
 \draw[guide](O)--(5,{-5*tan(12)}) (O)--(5,{5*tan(12)});
 \coordinate(A)at(-.254342226,.314086491);\coordinate(B)at(.343717036,-.278336567);
 \draw[DeepBlue,thick](S)--(A)--(5,0) (S)--(B)--(5,0);
 \draw[guide](Vone)--(A) (Vtwo)--(B);
 \fill(S)circle(1.5pt)node[above]{$S$};\fill(O)circle(1.2pt)node[below left]{$O$};
 \fill(Vone)circle(1.2pt)node[left]{$S_1'$};\fill(Vtwo)circle(1.2pt)node[left]{$S_2'$};
 \node[left]at(-1.75,1.5){$M_1$};\node[below]at(1.45,-1.3){$M_2$};
 \draw[<->](5.45,{-5*tan(12)})--(5.45,{5*tan(12)})node[midway,right]{$W$};
 \draw[<->](0,-2)--(5,-2)node[midway,below]{$L$};
 \draw[<->]({-2*cos(12)},-2.7)--(5,-2.7)node[midway,below]{$L+R\cos\alpha$};
 \node[anchor=west]at(.65,2.75){$OS=R$，小夹角 $\alpha$ 已夸大};
\end{tikzpicture}
\figcap{解图：用两个虚光源及有限公共照明区确定间距和条纹数。\label{ch29:quality:a6}}
\end{minipage}\end{center}

先统一几何距离。图~\ref{ch29:fig:s44-21-11} 中镜交线的截点为 $O$，屏位于 $x=L$；取光源 $S=(0,R)$，两镜所在直线的方向角分别为 $-45^\circ-\alpha/2$ 与 $-45^\circ+\alpha/2$。这里 $\alpha$ 是两镜相对于共面位置的小夹角。光源关于两镜的像分别为
\[
 S_1'=(-R\cos\alpha,R\sin\alpha),\qquad
 S_2'=(-R\cos\alpha,-R\sin\alpha).
\]
两镜分别形成相干虚源，间距
\[
 d=2R\sin\alpha\simeq2R\alpha.
\]
虚源平面位于 $O$ 后方 $R\cos\alpha$，所以虚源到屏的距离是 $L+R\cos\alpha$，不是 $L$。因此
\[
 \boxed{\Delta y=\frac{\lambda(L+R\cos\alpha)}{2R\sin\alpha}
 \approx1.134\,\mathrm{mm}}.
\]
两束反射光只在共同照明区内形成条纹。取镜面足够宽、外边缘不另行截光的理想情况，图~\ref{ch29:fig:s44-21-11} 中从 $S_1',S_2'$ 经过 $O$ 的边界射线在屏上交于 $y=\mp L\tan\alpha$，故重叠宽度 $W=2L\tan\alpha\approx24.435\,\mathrm{mm}$。两次反射的相位相同，中心为明纹。固定对称图样的亮纹中心满足 $y_m=m\Delta y$，于是
\[
 N_{\rm bright}=2\left\lfloor\frac W{2\Delta y}\right\rfloor+1
 =\boxed{21}.
\]
对应 $m=-10,\ldots,10$；暗纹中心有22条。若题目的“最多”还允许调节两路恒定相位、把条纹在同宽窗口内平移，则亮纹中心的最大数为 $\lfloor W/\Delta y\rfloor+1=22$。这不是当前中心为亮的对称图样的条数。镜面外缘、挡光片或不计边界纹会进一步限制可见数。

$R$ 加倍后，虚源间距加倍，但虚源到屏的距离也改变：
\[
 \Delta y'=\frac{\lambda(L+2R\cos\alpha)}{4R\sin\alpha}
 \approx0.593\,\mathrm{mm},\qquad
 \frac{\Delta y'}{\Delta y}=\frac{L+2R\cos\alpha}{2(L+R\cos\alpha)}.
\]
比例约为0.523，不是严格的一半。理想重叠宽度不变，固定对称图样有41条亮纹中心，允许调相位时最大42条。

最后必须说明“横移”的方向。取光源沿图中 $+x$ 方向移动 $\delta s$，即 $S=(\delta s,R)$。两虚像坐标变为
\[
 \begin{split}
 S_1'&=(-R\cos\alpha-\delta s\sin\alpha,\ R\sin\alpha-\delta s\cos\alpha),\\
 S_2'&=(-R\cos\alpha+\delta s\sin\alpha,-R\sin\alpha-\delta s\cos\alpha).
 \end{split}
\]
到两虚像等距的点满足中垂线方程 $Ry=\delta s\,x$，故中央亮纹在屏上的位移为
\[
 \boxed{\delta y_0=\frac LR\,\delta s}.
\]
在这一坐标约定下是向 $+y$ 移动；若将源的正方向反取，结果的符号也反取。小位移的一阶近似下条纹间距不变，但公共照明区边界也应随新光路检查。量纲上 $\lambda(L+R)/R$ 是长度；$L\gg R$ 时才恢复常见的 $\lambda L/(2R\alpha)$ 近似。

图~\ref{ch29:quality:a6} 中，两面镜反射所得的虚光源到交线距离均为 $R$，公共照明区还受两实际反射半平面的限制；条纹数须在该重叠区内计数。
\end{xesolution}


\Needspace{8\baselineskip}\begin{origexamplebox}{【比耶对切透镜】}\label{ch29:ex:s44-21-8}
焦距为 $f$ 的透镜前 $l_1>f$ 处为波长 $\lambda$ 的缝光源，透镜后 $L$ 处为屏，且 $L>l_2=fl_1/(l_1-f)$。透镜沿直径切为两半并沿切口垂直方向对称移开总距离 $\delta>0$。已知屏上的公共重叠区为 $-D/2\le y\le D/2$，关于中央明纹对称。求条纹间距、亮纹中心的条数，并说明只计两侧暗纹之间完整亮带时如何计数。

\begin{center}\begin{minipage}{.96\linewidth}\centering
\begin{tikzpicture}[>=Latex,font=\small,every node/.style={fill=none},scale=.88]

\fill (-2,0) circle (1.7pt) node[left]{$S$};\draw[<->,thick] (0,.35)--(0,1.5);\draw[<->,thick] (0,-1.5)--(0,-.35);
\draw[<->] (.4,-.35)--(.4,.35) node[midway,right]{$\delta$};\draw[thick] (4,-1.4)--(4,1.4) node[above]{屏};
\draw[dashed] (-2,0)--(4,0);\node[above] at (-1,1.1){$l_1>f$};\node[below] at (2,-1.6){屏距 $L>l_2$，公共区全宽 $D$};

\end{tikzpicture}
\figcap{题图：透镜沿直径分开并横向对称移开，两半的焦距相同。\label{ch29:quality:q7}}
\end{minipage}\end{center}
\end{origexamplebox}\par\nopagebreak[3]
\begin{xesolution}
\BilletSourceDiagram
未切透镜的像距
\[
 l_2=\frac{fl_1}{l_1-f}.
\]
设某半透镜的光轴平移到 $y=c$。相对于它的新光轴，源的物高为 $-c$，像高为 $(-l_2/l_1)(-c)$；换回原光轴，像的坐标为 $c(1+l_2/l_1)=cl_2/f$。取 $c=\pm\delta/2$，两实像源间距
\[
 d=\delta\frac{l_2}{f}.
\]
屏在像面之后时，相干源到屏的距离为 $L-l_2$，所以
\[
 \boxed{\Delta y=\frac{\lambda(L-l_2)}d
 =\frac{\lambda f(L-l_2)}{\delta l_2}}.
\]
图~\ref{ch29:fig:billet} 中 $d$ 是两实像源间距，$D$ 是屏上的公共照明宽度，二者不可混用。由中心亮纹位于 $y=0$，亮纹中心 $y_m=m\Delta y$ 必须满足 $|m|\Delta y\le D/2$，所以包括端点的亮纹中心条数是
\[
 \boxed{N_{\rm center}=2\left\lfloor\frac{D}{2\Delta y}\right\rfloor+1}.
\]
若端点不计入，且 $D/(2\Delta y)$ 恰为整数，应去掉这两条端点纹。若“完整亮纹”指相邻两个暗纹之间的整条亮带，则第 $m$ 条的两侧暗纹位于 $(m\pm\tfrac12)\Delta y$，要求 $(|m|+\tfrac12)\Delta y\le D/2$，因此
\[
 \boxed{N_{\rm full}=\max\!\left\{0,\,
 2\left\lfloor\frac{D}{2\Delta y}-\frac12\right\rfloor+1\right\}.}
\]
如何由实际口径确定题给的 $D$？设每半透镜从切口到外缘的有效高度为 $a_p$，取上半片轴高 $c=\delta/2$、入射高度 $h\in[c,c+a_p]$。屏上光线高度为
\[
 Y=\frac{L}{l_2}\frac{cl_2}{f}+\left(1-\frac{L}{l_2}\right)h.
\]
于是上半束覆盖 $[q-sa_p,q]$，下半束覆盖 $[-q,-q+sa_p]$，其中
$q=\frac\delta2(1+L/l_1)$、$s=L/l_2-1>0$。两者交集给出
\[
 \boxed{D=2\max\{0,\min(q,sa_p-q)\}.}
\]
外缘足够宽时才可化为 $D=2q=\delta(1+L/l_1)$；若 $sa_p\le q$，没有非零宽度的共同照明区，不能直接套条纹计数式。这里将恰在边界上的暗纹算作完整亮带的边界。若孔径使边缘光强趋零，实际能分辨的条数还可能更少。只有在允许条纹任意平移并求最大中心数时，才使用 $\lfloor D/\Delta y\rfloor+1$。
\end{xesolution}

\Needspace{20\baselineskip}\begin{origexamplebox}{【梅斯林干涉（Meslin interference）】}\label{ch29:ex:s44-21-9}
凸透镜沿直径切开，两半沿光轴错开。轴上单色点源经两半分别成像于 $S_1',S_2'$，观察屏垂直于轴。（1）求条纹形状；（2）$f=30\,\mathrm{cm}$，物到近半透镜 $60\,\mathrm{cm}$，两半错开 $8\,\mathrm{cm}$，$\lambda=500\,\mathrm{nm}$，屏置于两像中点，求近轴条纹尺度。

\begin{center}\begin{minipage}{.96\linewidth}\centering
\begin{tikzpicture}[>=Latex,font=\small,every node/.style={fill=none},scale=.88]

\fill (-2,0) circle (1.7pt) node[left]{$S$};\draw[<->,thick] (0,.15)--(0,1.3);\draw[<->,thick] (1,-1.3)--(1,-.15);
\draw[dashed] (-2,0)--(5,0);\draw[thick] (4,-1.4)--(4,1.4) node[above]{观察屏};
\draw[<->] (0,-1.7)--(1,-1.7) node[midway,below]{轴向错开 $8\,\mathrm{cm}$};\node[above] at (1,1.6){两半焦距相同};

\end{tikzpicture}
\figcap{题图：两半透镜沿光轴错开，两个共轭像点的轴向位置不同；屏位于题设两像的中点。\label{ch29:quality:q8}}
\end{minipage}\end{center}
\end{origexamplebox}\par\nopagebreak[3]
\begin{xesolution}
\MeslinSourceDiagram
图~\ref{ch29:fig:meslin} 中两相干像源在同一轴上，屏在两像之间。若只考虑相位条件，等相位差轨迹是同心圆；但上、下半透镜各只透过半个波面。先成像的一束过轴后进入另一半空间，与尚未成像的一束只在该半屏重叠，所以实际可见的是\boxed{\text{同心半圆弧}}，而不是整圈。

近半透镜像距 $v_1=30\times60/(60-30)=60\,\mathrm{cm}$。远半透镜到物的距离 $68\,\mathrm{cm}$，像距 $v_2=30\times68/(68-30)=53.684\,\mathrm{cm}$，按同一原点计其像坐标为 $8+53.684=61.684\,\mathrm{cm}$。两轴上像源间距
\[
 a=1.684\,\mathrm{cm}.
\]
屏位于中点，一束波在过焦点后发散，另一束仍在向焦点会聚；两波面曲率半径分别为 $+a/2$ 与 $-a/2$。近轴径向相位差增量
\[
 \Delta\varphi(r)-\Delta\varphi(0)=\frac{2kr^2}{a}.
\]
相邻圆纹满足
\[
 \boxed{r_{m+1}^2-r_m^2=\frac{\lambda a}{2}
 =4.21\times10^{-9}\,\mathrm{m^2}=0.00421\,\mathrm{mm^2}}.
\]
上式可直接由球面波相位展开解释：屏到两像的轴向距离均为 $a/2$，离轴点到像的距离增加量约为 $r^2/a$。一束已发散，另一束仍会聚，曲率符号相反，故相位差增加量为 $2kr^2/a$。实际相邻圆弧的径向间距为
\[
 \boxed{r_{m+1}-r_m=\frac{\lambda a}{2(r_{m+1}+r_m)}}.
\]
它随半径增大而减小；$\sqrt{\lambda a/2}\approx64.9\,\mu\mathrm m$ 是平方半径级差对应的特征尺度，不能不考虑轴上相位就把它当作“第一亮环半径”。过焦点相位和两路恒定相位决定中心明暗，不影响上述相邻级次的平方半径差。
\end{xesolution}

\Needspace{8\baselineskip}\begin{origexamplebox}{光源宽度怎样降低条纹反差}
\label{ch29:ex:depth-source-width}
沿用例~\ref{ch29:ex:s44-21-7} 的双面镜对称配置，取 $R=0.10\,\mathrm m$、$L=2.10\,\mathrm m$、$\alpha=20'$、$\lambda=600\,\mathrm{nm}$。把点光源改为沿 $x$ 方向、以 $(0,R)$ 为中心、全宽为 $w$ 的均匀线光源，不同源点彼此不相干。忽略源宽引起的振幅差和渐晕，仅考察所有源点均能照亮的近轴区域。（1）求条纹可见度随 $w$ 的关系；（2）求条纹第一次完全消失的源宽，并说明“能看见条纹”能否定义一个唯一最大源宽。

\begin{center}\begin{minipage}{.96\linewidth}\centering
\begin{tikzpicture}[>=Latex,font=\small,every node/.style={fill=none},scale=.88]

\draw[thick] (0,0)--(2,-2.2) (0,0)--(2.2,-2);
\fill (0,0) circle (1.5pt) node[left]{$O$};\draw[very thick] (-.45,1.2)--(.45,1.2);\node[above] at (0,1.2){均匀非相干线源，全宽 $w$};
\draw[dashed] (0,0)--(0,1.2) node[midway,left]{$R$};\draw[thick] (4,-1.6)--(4,1.6) node[above]{$x=L$};
\draw[dashed,->] (-.5,0)--(4.6,0) node[right]{$x$};\node[right] at (.65,-.65){小夹角 $\alpha$};

\end{tikzpicture}
\figcap{题图：线源沿 $x$ 方向展开，每个源点分别产生一套条纹，不同源点之间不相干。\label{ch29:quality:q9}}
\end{minipage}\end{center}
\end{origexamplebox}\par\nopagebreak[3]
\begin{xesolution}
一个源点经两镜形成的两束光彼此相干，但两个不同源点的场不能保持稳定相位，所以应先求每个源点的干涉强度，再把强度相加。例~\ref{ch29:ex:s44-21-7} 已得横移 $s$ 时中央亮纹移动 $Ms$，其中
\[
 M=\frac LR=21,\qquad
 \Delta y=\frac{\lambda(L+R\cos\alpha)}{2R\sin\alpha}
 \approx1.134\,\mathrm{mm}.
\]
这说明源上很小的横向距离可以在屏上变为较大的条纹错位。记 $k_y=2\pi/\Delta y$，来自源元 $\dd s$ 的强度可写为
\[
 \dd I=C[1+\cos k_y(y-Ms)]\dd s,
 \qquad -w/2\le s\le w/2,
\]
其中 $C$ 为单位源宽贡献的平均强度。积分中的正弦项因关于 $s=0$ 对称而抵消，余下
\[
 I(y)=Cw\left[1+
 \frac{\sin(k_yMw/2)}{k_yMw/2}\cos(k_yy)\right].
\]
因此
\begin{equation}\label{ch29:eq:depth-source-visibility}
 \boxed{\mathcal V(w)=\left|\frac{\sin(\pi Mw/\Delta y)}{\pi Mw/\Delta y}\right|}.
\end{equation}
当 $w\to0$ 时，分式趋于1，回到理想点光源；第一次为零时 $Mw=\Delta y$，故
\[
 \boxed{w_0=\frac{\Delta y}{M}=54.0\,\mu\mathrm m}.
\]

% Figure design: optical coherence, analytic plot rather than invented trend.
% Print width 9cm, minimum label 9pt; x=w/w0 in [0,2.5], V in [0,1].
% Main blue curve abs(sin(pi*x)/(pi*x)), value at 0 defined by limit.
% Dashed guides only at first zero x=1; zero at x=2 also labelled.
% ASCII: V |\_ /\_ / ; horizontal axis w/w0. No boxes or opaque tags.
\begin{center}\begin{minipage}{\linewidth}\centering
\begin{tikzpicture}[x=2.6cm,y=2.5cm,>=Latex,every node/.style={fill=none,font=\small,inner sep=2pt}]
 \draw[->] (0,0)--(2.65,0) node[right]{$w/w_0$};
 \draw[->] (0,0)--(0,1.17) node[above]{$\mathcal V$};
 \draw[DeepBlue,thick] (0,1)--(.004,{abs(sin(180*.004)/(pi*.004))});
 \draw[DeepBlue,thick,domain=.004:2.5,samples=301]
 plot(\x,{abs(sin(180*\x)/(pi*\x))});
 \draw[guide] (1,0)--(1,1);
 \node[left] at (0,1){$1$};\node[below left] at (0,0){$0$};
 \node[below] at (1,0){$1$};\node[below] at (2,0){$2$};
 \node[above,align=center] at (1.72,.61){第一次零点后\\还有低反差条纹};
\end{tikzpicture}
\figcap{均匀非相干线光源的宽度效应。横轴按第一次零点的源宽归一化。\label{ch29:fig:depth-source-visibility}}
\end{minipage}\end{center}

图~\ref{ch29:fig:depth-source-visibility} 显示 $w>w_0$ 后可见度不总为零，还存在较弱的旁瓣。实际“可见”要指定反差阈值、噪声和渐晕，因此题目若只说“具有一定反差”而不给标准，就没有唯一的最大允许宽度。这里的 $w_0$ 是清楚定义的第一次零点，不应偷换成普遍的可见性界限；其他源强分布也会给出不同的可见度函数。
\end{xesolution}


\section{薄膜干涉与增透膜}



分波阵面装置把一个波面的不同区域变成次级波源；分振幅装置则由部分反射/透射把同一束光的振幅分开，再使各束重合。薄膜反射是最常见的分振幅干涉。

图~\ref{ch29:fig:s44-21-6} 中，折射率为 $n$、厚度为 $D$ 的膜置于空气中，入射角为 $i$、膜内折射角为 $r$。上下表面两束反射光的传播光程差为
\[
 \Delta\mathcal L_{\rm geom}=2nD\cos r.
\]
若两次反射中恰有一次发生半波损失，总相位差可写成
\begin{equation}
 \Delta\varphi=\frac{2\pi}{\lambda_0}(2nD\cos r)+\pi.
\end{equation}
因此在\textbf{反射光}中，带一次半波损失时，暗纹满足 $2nD\cos r=m\lambda_0$，明纹满足 $2nD\cos r=(m+1/2)\lambda_0$。若零次或两次相变，明暗条件互换。正入射时只需令 $\cos r=1$。

图~\ref{ch29:fig:s44-21-6} 中传播光程差可逐段核算。进入薄膜的一束在膜内往返的光程为 $2nD/\cos r$；与此同时，直接在上表面反射的另一束在膜外多走了一段，其在出射方向上的投影光程为 $2D\tan r\sin i$。利用 $\sin i=n\sin r$，两者相减：
\[
 \frac{2nD}{\cos r}-2D\tan r\sin i
 =2nD\cos r.
\]
因此 $2nD/\cos r$ 不是最终程差，漏减外部补偿段会把角度依赖写反。相位突变应在传播光程算完后另行加入：只有上下两次反射的相变次数相差奇数，才等效多出 $\lambda_0/2$。

对楔形膜或缓慢变厚的油膜，$D$ 随位置变化，等光程差的位置形成条纹；对白光，不同 $\lambda_0$ 的加强位置不同，于是出现颜色。若膜厚超过光源相干长度，来自不同往返路径的相位关系不再稳定，彩纹会被平均掉。薄膜现象图和薄膜光路图中膜的两个界面、两束最终平行的出射光以及一次相变的位置，必须同时参与判断。

\FilmPhenomenaPanel
图~\ref{ch29:fig:s44-21-5} 中的薄膜现象都需要比较上下表面反射的两束光；以下计算既计膜内传播，也计反射相变。
\ThinFilmDiagram

太阳能电池、眼镜和相机镜头的增透膜利用反射光相消。取空气、透明膜、透明基底三个介质，折射率满足 $n_0<n_{\rm f}<n_{\rm s}$；正入射时上下界面的反射均有 $\pi$ 相变，二者的相对界面相位为零。要使传播相位相差 $\pi$，最薄的膜须满足 $2n_{\rm f}D=\lambda_0/2$，即 $D=\lambda_0/(4n_{\rm f})$。

\begin{center}\begin{minipage}{\linewidth}\centering\begin{tikzpicture}[>=Latex,font=\small]
\fill[Blue!10](-3,-.6)rectangle(3,.6);
\fill[SoftGray](-3,-1.5)rectangle(3,-.6);
\draw[thick](-3,.6)--(3,.6) (-3,-.6)--(3,-.6);
\node[anchor=west]at(3.3,1.1){空气 $n_0=1$};
\node[anchor=west]at(3.3,0){膜 $n_{\rm f}=1.45$};
\node[anchor=west]at(3.3,-1.1){基底 $n_{\rm s}=3.5$};
\draw[DeepBlue,->](-1.4,1.8)--(-1.4,-.6);
\draw[Blue,->](-.8,.6)--(-.8,1.8);
\draw[Gold,->](.3,-.6)--(.3,1.8);
\node[above]at(-1.4,1.8){入射};\node[above]at(-.5,2.3){两界面反射};
\draw[<->](2.6,-.6)--(2.6,.6)node[midway,left]{$D$};
\end{tikzpicture}\figcap{增透膜的三介质模型。正入射光线为便于辨认而横向分列，实际往返光路重合。\label{ch29:fig:coating}}\end{minipage}\end{center}

例如 $\lambda_0=550\,\mathrm{nm}$、$n_{\rm f}=1.45$ 时，$D\approx94.8\,\mathrm{nm}$。图~\ref{ch29:fig:coating} 的无吸收模型中，未镀膜反射率为 $[(n_{\rm s}-n_0)/(n_{\rm s}+n_0)]^2\approx30.9\%$；计入膜内全部往返反射，四分之一波膜的反射率为
\[
 R=\left(\frac{n_0n_{\rm s}-n_{\rm f}^2}{n_0n_{\rm s}+n_{\rm f}^2}\right)^2\approx6.22\%.
\]
只有 $n_{\rm f}^2=n_0n_{\rm s}$ 时理想反射率才恰为零。这里的基底折射率仅为计算模型；真实硅等材料还具有吸收与色散，不能直接拿此结果替代实测。真实多层膜还要匹配入射角和宽谱响应。镜头表面呈紫红或绿色，是其余波段反射光的综合色。

四分之一波厚度只保证两束主要反射光相差 $\pi$，要做到完全相消还需要两束反射振幅相等。单层膜在正入射下的理想振幅匹配要求膜折射率约为两侧折射率的几何平均；现实材料未必恰好满足，因此反射率只能被显著压低而非严格为零。入射角或波长改变时，$2nD\cos r$ 也随之改变，所以为某一中心波长设计的膜在反射光中会显出其余颜色，多层膜则用多组厚度扩大有效波段。

反射光与透射光的明暗条件互补：忽略吸收时，某波长在反射方向被相消，能量便主要进入透射方向。增透膜并不是把这部分光“吸收掉”，而是利用相位把能流重新分配；高反膜则反过来设计多层界面，使目标波段的反射振幅同相叠加。




\subsection{增透为何既要求膜厚，也要求折射率匹配}
膜厚决定相位，界面折射率决定反射振幅。只有相位相反而振幅不等，反射仍不能完全消失。以下取正入射、各介质无吸收且非磁性，入射介质、单层膜和半无限基底的折射率依次为 $n_0,n,n_s$。

用电场振幅的比值定义反射系数，两个界面的系数为
\[
 r_{01}=\frac{n_0-n}{n_0+n},\qquad
 r_{12}=\frac{n-n_s}{n+n_s}.
\]
负号已经包含反射的 $\pi$ 相变，使用这组带符号系数时不能再额外加一次“半波损失”。从膜内向入射侧反射有 $r_{10}=-r_{01}$；两次透射系数满足
\[
 t_{01}t_{10}=\frac{2n_0}{n_0+n}\frac{2n}{n_0+n}
 =1-r_{01}^2.
\]
设膜厚为 $t$，单程相位为 $\phi=2\pi nt/\lambda_0$，往返相位因子为 $q=\exp(2\ii\phi)$。第一束反射振幅为 $r_{01}$，进入膜后从底面反射并透回来的第一束为 $t_{01}t_{10}r_{12}q$，以后每多往返一次就乘 $r_{10}r_{12}q$。把等比级数相加，得到总反射系数
\[
 r=r_{01}+\frac{(1-r_{01}^2)r_{12}q}{1+r_{01}r_{12}q}
 =\frac{r_{01}+r_{12}q}{1+r_{01}r_{12}q}.
\]
最薄的四分之一波膜满足 $nt=\lambda_0/4$，因此 $q=-1$。要使分子为零，还必须有 $r_{01}=r_{12}$。代入并交叉相乘：
\[
 (n_0-n)(n+n_s)=(n-n_s)(n_0+n)
 \quad\Longrightarrow\quad n^2=n_0n_s.
\]
所以理想单层增透膜的两个条件为 $n=\sqrt{n_0n_s}$ 和 $t=\lambda_0/(4n)$。仅满足厚度条件时通常只能减弱反射；改变波长或入射角后，相位和振幅匹配也会改变，不能认为这一层膜对所有颜色和方向都完全增透。

\Needspace{8\baselineskip}\begin{origexamplebox}{【油膜干涉】}\label{ch29:ex:s44-21-12}
地面油膜为什么呈五彩斑斓？这说明了什么？

\begin{center}\begin{minipage}{.96\linewidth}\centering
\begin{tikzpicture}[>=Latex,font=\small,every node/.style={fill=none},scale=.88]

\draw[thick] (-1,0)--(4,0);\draw[thick] (-1,-.35)--(4,-.65);\node[right] at (4,-.35){油膜 $n$};
\draw[DeepBlue,->] (-.8,1.4)--(.5,0);\draw[DeepBlue,->] (.5,0)--(1.8,1.4);
\draw[DeepBlue,->] (.5,0)--(.85,-.46)--(1.3,0)--(2.7,1.4);
\draw[<->] (3,0)--(3,-.59) node[midway,left]{$t(x)$};\node[above] at (2.5,1.6){两界面反射光};

\end{tikzpicture}
\figcap{题图：油膜两个界面的反射光，膜厚随位置变化。\label{ch29:quality:q10}}
\end{minipage}\end{center}
\end{origexamplebox}\par\nopagebreak[3]
\begin{xesolution}
\begin{center}\begin{minipage}{.96\linewidth}\centering
\begin{tikzpicture}[>=Latex,font=\small,every node/.style={fill=none},scale=.88]

\draw[thick] (-1,0)--(4,0);\draw[thick] (-1,-.35)--(4,-.65);\node[right] at (4,-.35){油膜 $n$};
\draw[DeepBlue,->] (-.8,1.4)--(.5,0);\draw[DeepBlue,->] (.5,0)--(1.8,1.4);
\draw[DeepBlue,->] (.5,0)--(.85,-.46)--(1.3,0)--(2.7,1.4);
\draw[<->] (3,0)--(3,-.59) node[midway,left]{$t(x)$};\node[above] at (2.5,1.6){两界面反射光};

\end{tikzpicture}
\figcap{解图：在同一出射方向比较上、下两界面的反射光。\label{ch29:quality:a10}}
\end{minipage}\end{center}

白光照到油膜，上、下表面反射光重新重合。对给定位置和观察角，反射相位差约由 $2nD\cos r$ 加界面半波损失决定；不同波长的加强条件不同，所以反射光只突出某些颜色。油膜厚度和表面坡度随位置变化，颜色也随位置改变。

这说明油层薄到与可见光波长同量级，并且在至少一个相干面积内厚度变化缓慢；同时也直接显示光的波动叠加性质。若油层远厚于白光相干长度，众多反射次序平均后便不会形成清晰彩纹。

图~\ref{ch29:quality:a10} 中，两束反射光的几何光程差含 $2nt\cos r$，还要加入两界面的反射相位差；不同位置的 $t$ 选择不同增强波长，故出现彩色。
\end{xesolution}

\Needspace{8\baselineskip}\begin{origexamplebox}{【薄膜干涉】}\label{ch29:ex:s44-21-13}
竖直的未吹泡泡液膜从上到下呈不同颜色，颜色缓慢变化并最终变黑，为什么？

\begin{center}\begin{minipage}{.96\linewidth}\centering
\begin{tikzpicture}[>=Latex,font=\small,every node/.style={fill=none},scale=.88]

\draw[thick] (0,2.2)--(0,-.4);\draw[thick] (.1,2.2)--(.65,-.4);\node[right] at (.7,1.1){竖直液膜};
\draw[->] (2,1.9)--(2,.7) node[midway,right]{$\boldsymbol g$};
\draw[DeepBlue,->] (-1,1.8)--(0,1.4);\draw[DeepBlue,->] (-1,.3)--(0,0);
\node[left] at (-.2,2.35){顶端逐渐变薄};\node[below] at (.3,-.6){空气/液体/空气};

\end{tikzpicture}
\figcap{题图：液体在重力下排向下方，膜顶更薄；图中为两侧空气的自由液膜。\label{ch29:quality:q11}}
\end{minipage}\end{center}
\end{origexamplebox}\par\nopagebreak[3]
\begin{xesolution}
\begin{center}\begin{minipage}{.96\linewidth}\centering
\begin{tikzpicture}[>=Latex,font=\small,every node/.style={fill=none},scale=.88]

\draw[thick] (0,2.2)--(0,-.4);\draw[thick] (.1,2.2)--(.65,-.4);\node[right] at (.7,1.1){竖直液膜};
\draw[->] (2,1.9)--(2,.7) node[midway,right]{$\boldsymbol g$};
\draw[DeepBlue,->] (-1,1.8)--(0,1.4);\draw[DeepBlue,->] (-1,.3)--(0,0);
\node[left] at (-.2,2.35){顶端逐渐变薄};\node[below] at (.3,-.6){空气/液体/空气};

\end{tikzpicture}
\figcap{解图：厚度随高度和时间改变，局部两界面反射决定颜色。\label{ch29:quality:a11}}
\end{minipage}\end{center}

重力使液体向下排流，膜通常上薄下厚；每一高度对应不同 $D$，反射增强的波长不同，因而形成近似水平的等厚彩带。排流继续时，同一高度的厚度随时间减小，颜色便缓慢迁移。

当顶部膜厚降到几十纳米量级，$2nD$ 对全部可见波长都很小；上下表面反射中恰有一次 $\pi$ 相变，两束反射光近乎相消，形成“黑膜”。黑色不是液体变黑，而是可见反射光被压低；膜已经很薄，但黑膜可因表面力而维持一段时间，变黑本身不等于已经破裂或必将立即破裂。

图~\ref{ch29:quality:a11} 中，同一高度的膜厚随排液而减小，使增强波长依次变化；极薄时两反射光之间的半波损失主导，低级反射趋于相消而发黑。
\end{xesolution}

\section{劈尖与牛顿环}




\subsection{把条纹位置转化成膜厚位置}
在单色光近于垂直入射、恰有一次相对半波损失时，空气膜的反射暗纹满足 $2t=m\lambda_0$。劈尖小角为 $\alpha$，距接触线的坐标为 $x$，则 $t=x\tan\alpha\simeq\alpha x$，从而
\[
 x_m=\frac{m\lambda_0}{2\alpha},\qquad
 x_{m+1}-x_m=\frac{\lambda_0}{2\alpha}.
\]
因此直而等间距的条纹意味着所测截面上的厚度近似线性增加。若有局部凸起，条纹要沿保持同一厚度的位置弯曲；判断凸凹时必须说明观察的是哪一界面以及空气膜变厚还是变薄。

牛顿环中，球面曲率半径为 $R$、离轴半径为 $\rho$。几何关系不是 $t=\rho^2/R$，而是
\[
 \rho^2+(R-t)^2=R^2,
 \quad \rho^2=2Rt-t^2,
 \quad t\simeq\frac{\rho^2}{2R}\ (t\ll R).
\]
于是反射暗环半径满足 $\rho_m^2=m\lambda_0R$。相邻暗环的半径差并不恒定，恒定的是半径平方差。测量时使用
\[
 D_{m+p}^2-D_m^2=4p\lambda_0R
\]
（$D_m=2\rho_m$），能避开把中央某个环误认作第零级的影响。

若透镜并未真正接触平板，中心存在间隙 $t_0$，应写 $t=t_0+\rho^2/(2R)$。此时中心未必为暗，环的绝对级次也改变，但相隔 $p$ 级的直径平方差公式仍成立。把“中心必须暗”当作不带条件的规律，会把灰尘、接触不良和膜厚偏置误解释成新现象。

\Needspace{8\baselineskip}\begin{origexamplebox}{【劈尖的干涉】}\label{ch29:ex:s44-21-4}
透明劈尖的尖角为 $\theta$，沿传播方向 $x$ 的最大厚度为 $h$，折射率为 $n(x)=1+bx$（$b>0$）。将斜面理解为许多微小台阶，各台阶的入射、出射面均垂直于 $x$ 轴，忽略台阶边缘的衍射及反射损失。波长 $\lambda$ 的平行光沿 $+x$ 入射，各光线在介质中直线传播。其后紧贴不透明挡板，开有平行于 $z$ 轴的狭缝；缝在 $y$ 方向排列，第一缝位于尖端 $y=0$，该处厚度为零。理想会聚透镜把各缝光汇聚到后焦点，各缝至焦点的其余光程相同。求其余缝的允许坐标，并判断是否可以选择等距的缝，使焦点处各束光同相加强。

\begin{center}\begin{minipage}{.96\linewidth}\centering
\begin{tikzpicture}[>=Latex,font=\small,every node/.style={fill=none},scale=.88]

\fill[DeepBlue!7] (0,0)--(0,2.3)--(1.5,2.3)--cycle;\draw[thick] (0,0)--(0,2.3)--(1.5,2.3)--cycle;
\draw[->] (-1,0)--(3.2,0) node[right]{$x$};\draw[->] (-.5,-.2)--(-.5,2.7) node[above]{$y$};
\foreach \y in {.5,1,1.5,2}{\pgfmathsetmacro{\xx}{1.5*\y/2.3}\draw[DeepBlue,->] (-1,\y)--(\xx,\y);\draw[thick] (\xx,\y-.08)--(\xx,\y+.08);}
\draw[<->] (0,2.55)--(1.5,2.55) node[midway,above]{$h$};\node[right] at (1.6,1.5){挡板各狭缝};
\node[right] at (1.7,.6){$x_y=y\tan\theta$};

\end{tikzpicture}
\figcap{题图：各高度光线经历的厚度 $x_y$ 不同，挡板缝的坐标沿 $y$ 量取；微台阶界面按题设处理。\label{ch29:quality:q12}}
\end{minipage}\end{center}
\end{origexamplebox}\par\nopagebreak[3]
\begin{xesolution}
图~\ref{ch29:fig:wedge} 中，$y$ 是缝沿挡板的坐标，$x_y$ 才是该光线穿过介质的厚度。由尖角几何，$x_y=y\tan\theta$。
\WedgeInterferenceDiagram
相对于同样几何长度的空气，该光线增加的光程为
\[
 \Delta\mathcal L(y)=\int_0^{x_y}[n(x)-1]\,\dd x
 =\int_0^{x_y}bx\,\dd x=\frac b2y^2\tan^2\theta.
\]
物到像的其余光程相等，故各缝要在透镜焦点同相加强，须 $\Delta\mathcal L(y_m)=m\lambda$：
\[
 \boxed{y_m=\cot\theta\sqrt{\frac{2m\lambda}{b}},\qquad m=0,1,2,\ldots}
\]
并受 $x_y\le h$ 限制，即 $0\le m\le bh^2/(2\lambda)$。注意 $[b]=\mathrm m^{-1}$，所以 $\sqrt{\lambda/b}$ 的量纲为长度。

若把每一个相邻整数级次都开成一条缝，则
\[
 y_{m+1}-y_m=\cot\theta\sqrt{\frac{2\lambda}{b}}
 (\sqrt{m+1}-\sqrt m),
\]
间距确实逐渐减小。但题目并未要求所有允许位置都开缝：\textbf{选取部分级次可以得到等距缝列。}
取任意正整数 $q$，只选 $m=qj^2$，就有
\[
 \boxed{y_j=j\cot\theta\sqrt{\frac{2q\lambda}{b}},\quad
 j=0,1,\ldots,\left\lfloor h\sqrt{\frac{b}{2q\lambda}}\right\rfloor.}
\]
此时第 $j$ 条缝的附加光程为 $qj^2\lambda$，与第一缝的相位差为 $2\pi qj^2$，每束仍同相。最小等距间隔对应 $q=1$；至少能开两条缝还要求 $bh^2\ge2\lambda$。这区分了“连续级次不等距”和“选取级次可以等距”两件事。
\end{xesolution}

\NewtonRingsDiagram
图~\ref{ch29:fig:newton} 中透镜与平板之间的空气层厚度随离接触点的距离增大；等厚位置成圆，因而干涉条纹呈同心环。

\Needspace{8\baselineskip}\begin{origexamplebox}{【牛顿环】}\label{ch29:ex:s44-21-14}
把凸面半径为 $R$ 的平凸透镜以凸面向下放在平面镜上，中心恰好接触，从反射方向观察空气膜的明暗同心环。（1）求第 $m$ 级条纹半径；（2）判断中心明暗；（3）说明这一现象的早期观察者。题目只写“平面镜”，未指定其反射镀层；请先说明所需的反射相位条件，再给出常用透明玻璃底板模型的结果。

\begin{center}\begin{minipage}{.96\linewidth}\centering
\begin{tikzpicture}[>=Latex,font=\small,every node/.style={fill=none},scale=.88]

\draw[thick] (-2,0)--(2,0);\draw[thick,domain=-1.8:1.8,smooth] plot (\x,{.18*\x*\x});\draw[thick] (-1.8,.9)--(1.8,.9);
\fill (0,0) circle (1.5pt) node[below]{$O$ 接触点};\node[right] at (1.5,.3){空气膜};
\draw[DeepBlue,->] (-1.4,1.6)--(-.8,.3);\draw[<->] (1.1,0)--(1.1,.218) node[right]{$t(r)$};
\draw[<->] (0,-.6)--(1.1,-.6) node[midway,below]{$r$};

\end{tikzpicture}
\figcap{题图：凸面与底板之间的空气膜及接触中心；底板的反射相位仍需按题设材料确定。\label{ch29:quality:q13}}
\end{minipage}\end{center}
\end{origexamplebox}\par\nopagebreak[3]
\begin{xesolution}
\begin{center}\begin{minipage}{.96\linewidth}\centering
\begin{tikzpicture}[scale=.92,>=Latex]
 \begin{scope}
  \draw[very thick](-2.5,0)--(2.5,0);
  \draw[very thick,domain=-2.2:2.2,samples=101]plot(\x,{4-sqrt(16-\x*\x)});
  \draw[thick](-2.2,{4-sqrt(16-2.2*2.2)})--(-2.2,1.1)--(2.2,1.1)--(2.2,{4-sqrt(16-2.2*2.2)});
  \fill(0,0)circle(1.3pt)node[below left]{$O$};
  \draw[<->](1.8,0)--(1.8,{4-sqrt(16-1.8*1.8)})node[midway,right=4pt]{$t(r)$};
  \draw[<->](0,-.65)--(1.8,-.65)node[midway,below]{$r$};
  \node[above]at(0,1.15){平凸透镜};\node[below]at(-1.1,-1.15){底板（透明玻璃特例）};
 \end{scope}
 \begin{scope}[shift={(6,0)}]
  \foreach \m in {1,...,7}{\draw[DeepBlue](0,0)circle({.55*sqrt(\m)});}
  \fill(0,0)circle(.06);\node[below]at(0,-1.65){暗环满足 $r_m^2=m\lambda R$};
 \end{scope}
\end{tikzpicture}
\figcap{解图：由接触几何确定膜厚，再连同反射相位判断明暗级次。\label{ch29:quality:a13}}
\end{minipage}\end{center}

“平面镜”只确定几何形状，未确定反射相位。如果是金属镀层，不能从题给的 $R,\lambda$ 唯一确定中心明暗。下面保留一般相位 $\delta_r$，并把透明玻璃底板作为条件明确的特例。
离中心半径 $r$ 处的空气膜厚度由球面几何给出
\[
 t=R-\sqrt{R^2-r^2}\simeq\frac{r^2}{2R}.
\]
图中的 $O$ 是真实接触点，厚度从 $O$ 向外增大。透明玻璃模型下，玻璃到空气的反射无半波损失，空气到玻璃的反射有半波损失，故反射暗环满足 $2t=m\lambda$：
\[
 \boxed{r_m^{(\rm dark)}=\sqrt{m\lambda R},\qquad m=0,1,2,\ldots}
\]
反射明环满足 $2t=(m+1/2)\lambda$：
\[
 \boxed{r_m^{(\rm bright)}=\sqrt{\left(m+\tfrac12\right)\lambda R}}.
\]
$r=0$ 处几何程差为零但有半波损失，因此\boxed{\text{中心为暗纹}}。若透镜与平面不真正接触，中心会因残余膜厚而改变。

如果底板改为金属镜，反射系数一般为复数。设两束反射引入的相对相位为 $\delta_r$，则总相位差为 $4\pi t/\lambda+\delta_r$，相消位置满足
\[
 \frac{4\pi t}{\lambda}+\delta_r=(2m+1)\pi.
\]
在 $r\ll R$ 下，极小与极大的平方半径分别为
\[
 r_{m,\min}^2=\lambda R\left(m+\frac12-\frac{\delta_r}{2\pi}\right),\qquad
 r_{m,\max}^2=\lambda R\left(m-\frac{\delta_r}{2\pi}\right),
\]
只保留右端非负的整数 $m$。中心强度为 $I(0)=I_1+I_2+2\sqrt{I_1I_2}\cos\delta_r$。所以中心是否为暗纹要由 $\delta_r$ 决定，不能直接沿用透明玻璃的半波损失判断；两束振幅不等时，相消处也只是极小而不是完全无光。几何关系 $t\simeq r^2/(2R)$ 仍成立。

这类颜色和环纹在牛顿之前已由罗伯特·胡克（Robert Hooke）观察并记载；牛顿后来作了系统、定量研究，“牛顿环”由此得名。

图~\ref{ch29:quality:a13} 中，空气膜厚度近轴为 $r^2/(2R)$，所以等光程差线为同心圆；透明玻璃底板产生一次半波损失，中心为暗，金属镀层须另行指定相位。
\end{xesolution}

\section{迈克耳孙干涉仪与测量}



迈克耳孙干涉仪用半透半反分光镜把光分入两条近乎正交的臂，经 $M_1,M_2$ 反射后再次合束。补偿板 $P$ 使两束光穿过相同厚度的玻璃。把一面反射镜关于分光镜成像，两臂等效为两面几乎平行或略有夹角的镜；微小位移会让干涉条纹移动。

若某臂反射镜移动 $x$，往返光程改变 $2x$，每改变一个波长就有一条条纹扫过固定点。这使迈克耳孙装置成为高灵敏位移和波长测量工具。迈克耳孙的工作获 1907 年诺贝尔物理学奖；现代大型激光干涉仪沿用相同相位测量思想，通过长臂、多次反射、真空和隔振检测极小长度变化，引力波探测是其著名应用。

大型引力波装置可以看作把这一原理推到极端：两条互相垂直的真空臂长达千米，光在臂腔中多次往返以增加等效光程，隔振与双站符合则用来排除局部扰动。引力波通过时，两臂长度以相反方式产生极微小相对变化，合束处的相位和光强随之改变。它测量的不是镜子“被光推着移动”，而是两条光路长度差的时间变化；这一区别也说明迈克耳孙图中两臂、分光镜、反射镜和探测端的连接关系都不能省略。

分光镜使一路先反射后透射，另一路先透射后反射；若没有补偿板，两路穿过玻璃的次数和厚度不同，宽带光中会出现额外色散，零光程差难以同时满足。把反射镜 $M_2$ 关于分光镜作虚像后，仪器可等效为两面“虚镜”形成的空气薄膜：虚镜平行时得到等倾圆纹，略有夹角时得到近似直条纹。镜移动 $\lambda/2$ 会使往返程差变化一个波长，因此固定点恰有一条纹通过，这是后面例~\ref{ch29:ex:s44-21-10} 速度测量的直接依据。

\MichelsonDiagram
图~\ref{ch29:fig:s44-21-7} 的两臂都往返传播；反射镜沿臂移动距离 $x$，该臂的几何路程改变 $2x$，而不是 $x$。

激光干涉引力波天文台（LIGO）在美国汉福德和利文斯顿各设一台干涉仪，每台的两条真空臂几何长度均为 $4\,\mathrm{km}$。图~\ref{ch29:fig:ligo-sites} 中左侧为干旱地区的汉福德站，右侧为森林地区的利文斯顿站；每站两臂互相垂直。不同站点的环境扰动通常不相关，比较两站信号到达时间和波形，是判断天体来源的重要步骤。

\historyfigure{ligo-two-sites-20260923.png}{引力波双站的环境与两条正交长臂；情景图不用于量取尺寸。\label{ch29:fig:ligo-sites}}

\begin{center}\begin{minipage}{\linewidth}\centering\begin{tikzpicture}[x=.19cm,y=.13cm,font=\small]
\draw[RuleGray](-125,25)rectangle(-65,50);
\foreach \x in {-120,-100,-80}{\draw[RuleGray](\x,25)--(\x,50);\node[below]at(\x,25){$\x^\circ$};}
\foreach \y in {30,40,50}{\draw[RuleGray](-125,\y)--(-65,\y);\node[left]at(-125,\y){$\y^\circ$};}
\coordinate(H)at(-119.407,46.455);\coordinate(L)at(-90.774,30.563);
\fill[DeepBlue](H)circle(2pt);\fill[DeepBlue](L)circle(2pt);
\draw[Blue,dashed](H)--(L);
\node[above right]at(H){汉福德（美国西北部）};
\node[above right]at(L){利文斯顿（美国南部）};
\end{tikzpicture}\figcap{双站近似经纬位置；横坐标是东经有向角，负值表示西经，纵坐标为北纬。图上直线不用于量取球面距离。\label{ch29:fig:ligomap}}\end{minipage}\end{center}

图~\ref{ch29:fig:ligomap} 中汉福德位于利文斯顿的西北方向。两站具有不同的局部环境，远距离布置使共同信号与单站噪声的区分更有依据；信号到达的先后还取决于波源方向。

首次直接探测到的引力波信号于2015年9月14日到达两站。若应变的量级为 $h\sim10^{-21}$，则单臂长度变化的量级是 $hL\sim4\times10^{-18}\,\mathrm m$，而不是毫米或纳米。臂内的法布里--珀罗腔使光反复往返，低频相位响应约增强300倍；这是光学响应的增强，不能把实体臂长写成几百倍的 $4\,\mathrm{km}$，也不能把增强系数不加频率条件地视为常数。

\subsection{移动镜面、计数与折射率测量}
镜面沿其法线移动 $\Delta l$ 时，一路光往返的几何长度改变 $2\Delta l$。在空气近似为真空且近轴时，若观察点有 $N$ 个完整条纹周期通过，便有
\[
 2|\Delta l|=N\lambda_0.
\]
因子2来自同一路的去程和回程，不来自仪器有两条臂。计数给出的是变化量，通常不能确定初始绝对级次；镜面向哪侧移动，则必须另外通过机械位移或条纹运动方向判断。

若在一路放入长度为 $l$ 的气室，折射率从 $n_a$ 变为 $n_b$，光在气室中往返，其他条件不变，则光程变化为 $2l(n_b-n_a)$，于是
\[
 |n_b-n_a|=\frac{N\lambda_0}{2l}.
\]
若装置只单程通过待测片，分母便没有2。做题时先把实际路径逐段列出来，再决定“一次”还是“两次”，比记忆仪器名称更可靠。

把一块斜放玻璃板插入一路时，除了改变相位，还可能引起横向位移和色散。白光干涉要求两臂不仅几何长度接近，所经过玻璃的厚度也应匹配；补偿板的作用就是让不同频率的相位能够同时接近匹配，不能仅凭某一单色光的零程差认定白光条件已满足。


% Moving-source panel relocated to its solution.






\Needspace{8\baselineskip}\begin{origexamplebox}{【变化臂长】}\label{ch29:ex:s44-21-10}
迈克耳孙干涉仪一臂反射镜以速度 $v$ 匀速平移。（1）$\lambda=600\,\mathrm{nm}$，固定取样窗中的信号频率 $50\,\mathrm{Hz}$，求镜速；（2）用平均波长 $589.3\,\mathrm{nm}$ 的钠黄光，镜速同前，拍频 $5.2\times10^{-2}\,\mathrm{Hz}$，求两谱线波长差。

\begin{center}\begin{minipage}{.96\linewidth}\centering
\begin{tikzpicture}[>=Latex,font=\small,every node/.style={fill=none},scale=.88]

\draw[thick] (-.35,-.35)--(.35,.35);\draw[very thick] (3,-.7)--(3,.7);\draw[very thick] (-.7,2)--(.7,2);
\draw[DeepBlue,->] (-2,0)--(0,0);\draw[DeepBlue,->] (0,0)--(2.8,0);\draw[DeepBlue,->] (0,0)--(0,1.8);\draw[->] (0,0)--(0,-1.1);
\node[below] at (-1.5,0){光源};\node[right] at (0,-1){探测窗};\draw[->] (3.2,1)--(4.2,1) node[right]{$v$};
\node[above] at (3,.7){动镜};\node[above] at (0,2){定镜};

\end{tikzpicture}
\figcap{题图：迈克耳孙两臂，动镜沿一臂传播方向匀速移动。\label{ch29:quality:q14}}
\end{minipage}\end{center}
\end{origexamplebox}\par\nopagebreak[3]
\begin{xesolution}
\begin{center}\begin{minipage}{.96\linewidth}\centering
\begin{tikzpicture}[>=Latex,font=\small,every node/.style={fill=none},scale=.88]

\draw[thick] (-.35,-.35)--(.35,.35);\draw[very thick] (3,-.7)--(3,.7);\draw[very thick] (-.7,2)--(.7,2);
\draw[DeepBlue,->] (-2,0)--(0,0);\draw[DeepBlue,->] (0,0)--(2.8,0);\draw[DeepBlue,->] (0,0)--(0,1.8);\draw[->] (0,0)--(0,-1.1);
\node[below] at (-1.5,0){光源};\node[right] at (0,-1){探测窗};\draw[->] (3.2,1)--(4.2,1) node[right]{$v$};
\node[above] at (3,.7){动镜};\node[above] at (0,2){定镜};

\end{tikzpicture}
\figcap{解图：动镜所在光臂的往返光程变化包含两次镜位移。\label{ch29:quality:a14}}
\end{minipage}\end{center}

镜移动 $v\dd t$ 使往返光程改变 $2v\dd t$，故条纹频率 $\nu=2v/\lambda$：
\[
 \boxed{v=\frac{\nu\lambda}{2}=1.50\times10^{-5}\,\mathrm{m/s}=15.0\,\mu\mathrm{m/s}}.
\]
两近邻波长产生的时间频率差
\[
 \Delta\nu=2v\left|\frac1{\lambda_1}-\frac1{\lambda_2}\right|
 \simeq\frac{2v\Delta\lambda}{\bar\lambda^2}.
\]
所以
\[
 \boxed{\Delta\lambda=\frac{\Delta\nu\,\bar\lambda^2}{2v}
 \approx6.02\times10^{-10}\,\mathrm m=0.602\,\mathrm{nm}}.
\]
这里两色的光学频率相差很大，探测器并不是记录光场在光学频率上的拍频，而是先分别得到两列随镜位移变化的低频强度信号，再叠加：
\[
 \cos(2\pi\nu_1t)+\cos(2\pi\nu_2t)
 =2\cos[\pi(\nu_1-\nu_2)t]\cos[\pi(\nu_1+\nu_2)t].
\]
包络的绝对值每隔 $1/|\nu_1-\nu_2|$ 再现一次，故题给拍频取两条纹频率之差，不能再多乘或少乘一个 $2$。两谱线强度不等会降低包络谷的深度，但不会改变其再现周期。
这里 $50\,\mathrm{Hz}$ 与 $5.2\times10^{-2}\,\mathrm{Hz}$ 的量级相差三阶，代入时必须保留指数。

图~\ref{ch29:quality:a14} 中，动镜单程位移 $\Delta x$ 在往返光路中改变 $2\Delta x$，因此单色条纹频率为 $2v/\lambda$；两相近波长的拍频是各自条纹频率之差。
\end{xesolution}

\section{复振幅与相干性的定量描述（选读）}



标量波模型满足
\begin{equation}
 \left(\nabla^2-\frac1{v^2}\frac{\partial^2}{\partial t^2}\right)\psi(\vect r,t)=0.
\end{equation}
对角频率为 $\omega$ 的单色波，取时间因子 $e^{-\ii\omega t}$，写成
\begin{equation}
 \psi(\vect r,t)=\Re\!\left[U(\vect r)e^{-\ii\omega t}\right],
 \qquad (\nabla^2+k^2)U=0,
\end{equation}
其中 $k=\omega/v$，第二式是亥姆霍兹方程（Helmholtz equation）。改变时间因子约定会同时改变传播相因子的符号，但可测的强度和相位差不变。

下文把复振幅归一化为 $I_j=\langle|U_j|^2\rangle$；若直接使用电场复振幅，在同一无吸收介质中还需乘共同系数 $n\varepsilon_0c/2$。对同偏振光，探测器在远长于光学周期的时间内平均，得到
\begin{equation}
 I=I_1+I_2+2\Re\Gamma_{12},
 \qquad
 \Gamma_{12}=\left\langle U_1^*U_2\right\rangle.
\end{equation}
当 $I_1,I_2$ 均非零时，归一化复相干度定义为
\begin{equation}
 \gamma_{12}=\frac{\Gamma_{12}}{\sqrt{I_1I_2}},
 \qquad 0\le |\gamma_{12}|\le1.
\end{equation}
改变两路相位差并记录条纹最大、最小强度，条纹可见度为
\begin{equation}
 \mathcal V=\frac{I_{\max}-I_{\min}}{I_{\max}+I_{\min}}
 =\frac{2\sqrt{I_1I_2}}{I_1+I_2}|\gamma_{12}|.
\end{equation}
因此条纹变浅既可能是相干度下降，也可能只是两束强度相差过大；仅凭“看不见条纹”不能断言两束完全不相干。

时间相干性比较同一点不同时刻的场，谱宽为 $\Delta\nu$ 的准单色光，其相干时间和真空相干长度量级为
\begin{equation}
 \tau_c\sim\frac1{\Delta\nu},
 \qquad \ell_c\sim c\tau_c.
\end{equation}
空间相干性比较同一波面不同位置的场，扩展光源越大，不同源点形成的条纹错位越明显，相干面积通常越小。这分别解释了宽带光对光程差敏感、宽光源对双缝间距敏感的原因。

\tcbset{breakable/.default=true}\begin{starread}{多光束干涉与法布里--珀罗腔（Fabry--P\'erot cavity）}
两块平行的部分反射镜使光在腔内多次往返。相邻透射束的往返相位差为
\[
 \delta=\frac{4\pi nd\cos r}{\lambda_0}.
\]
若两镜强度反射率相同且为 $R$，忽略吸收，每通过两镜一次，振幅因子为 $1-R$；
每多往返一周，振幅再乘 $R\e^{\ii\delta}$。因此全部透射束的复振幅和为
\[
 \frac{E_t}{E_0}
 =(1-R)\e^{\ii\delta/2}
 \sum_{j=0}^{\infty}(R\e^{\ii\delta})^j
 =\frac{(1-R)\e^{\ii\delta/2}}{1-R\e^{\ii\delta}}.
\]
取模平方，并利用
$|1-R\e^{\ii\delta}|^2=(1-R)^2+4R\sin^2(\delta/2)$，得到艾里公式（Airy formula）：
\[
 \frac{I_t}{I_0}=\frac1{1+\mathcal F\sin^2(\delta/2)},
 \qquad \mathcal F=\frac{4R}{(1-R)^2}.
\]
当 $\delta=2m\pi$ 时各束同相，透射达到极大；$R$ 越接近1，峰越窄，仪器分辨相近谱线的能力越强。

\begin{center}\begin{minipage}{\linewidth}\centering
\begin{tikzpicture}[scale=.86,>=Latex]
  \draw[very thick] (0,-1.55)--(0,2.05);
  \draw[very thick] (5,-1.55)--(5,2.05);
  \node[below] at (0,-1.65){部分反射镜};
  \node[below] at (5,-1.65){部分反射镜};
  \draw[<->,guide] (.1,2.18)--(4.9,2.18) node[midway,above]{$d$};
  \draw[thick,DeepBlue] (-1.4,-1.254)--(0,-1.10)--(5,-.55)--(0,0)--(5,.55)--(0,1.10)--(5,1.65);
  \draw[-{Latex[length=2mm]},thick,DeepBlue] (5,-.55)--(6.4,-.396);
  \draw[-{Latex[length=2mm]},thick,DeepBlue] (5,.55)--(6.4,.704);
  \draw[-{Latex[length=2mm]},thick,DeepBlue] (5,1.65)--(6.4,1.804);
  \node[right] at (6.4,.704){逐次透射束};
\end{tikzpicture}
\figcap{法布里--珀罗腔中的多次反射与透射。右侧各束平行；用会聚透镜在后焦面叠加同方向的光束，或在远场观察其共同角谱。}
\end{minipage}\end{center}

相邻纵模满足 $2nd\cos r=m\lambda_0$。近轴且介质色散可忽略时，相邻共振频率之差（自由光谱范围）为
\[
 \Delta\nu_{\rm FSR}=\frac{c}{2nd}.
\]
在高反射率近似下，令 $\delta=2m\pi+\varepsilon$ 且
$\sin(\varepsilon/2)\simeq\varepsilon/2$。透射强度降到峰值一半时
$\mathcal F\varepsilon^2/4=1$，故单个透射峰的频率半高全宽为
\[
 \delta\nu_{1/2}=\frac{\Delta\nu_{\rm FSR}}{\mathcal F_{\rm in}},
 \qquad
 \mathcal F_{\rm in}=\frac{\pi\sqrt R}{1-R},
\]
其中 $\mathcal F_{\rm in}$ 称为精细度（finesse）。注意艾里公式中的系数 $\mathcal F$ 与精细度
$\mathcal F_{\rm in}$ 不是同一个量；前者控制峰形，后者是“相邻峰间距与单峰宽度之比”。

例如 $d=1.00\,\mathrm{cm}$、$n=1$、$R=0.95$ 时，
$\Delta\nu_{\rm FSR}=15.0\,\mathrm{GHz}$，$\mathcal F_{\rm in}\approx61.2$，
因而 $\delta\nu_{1/2}\approx245\,\mathrm{MHz}$。若两条谱线的频差远小于这个峰宽，它们在该腔的透射谱上不能清楚分开；
增大 $R$ 会压窄峰，但也会增加建立稳定腔场所需的往返次数，并提高对镜面平行度和机械稳定性的要求。

这个频域结论还有直接的时域解释。关掉入射光后，腔内光每经过一个往返时间
$t_{\rm rt}=2nd/c$，能量乘以 $R^2$，所以
\[
 U(t)=U(0)\exp\!\left(-\frac{t}{\tau_E}\right),
 \qquad
 \tau_E=-\frac{t_{\rm rt}}{2\ln R}
 \simeq\frac{nd}{c(1-R)}.
\]
高反射率近似下，谱峰半高全宽满足
\[
 \boxed{\delta\nu_{1/2}\simeq\frac1{2\pi\tau_E}}.
\]
因此“峰越窄”与“腔内能量衰减越慢”是同一物理事实的频域和时域表述；测量关断后的环降曲线，
也能反推出腔损耗和精细度。
\end{starread}\tcbset{breakable/.default=false}


\section{从相位平均理解条纹能否观察}
写下光程差只能确定理想相干分量的明暗位置。实际探测器还会对曝光时间、光谱宽度和光源的空间范围平均；只有这些平均之后的振荡项没有被抵消，条纹才具有可见反差。
沿用前文 $\Gamma_{12}=\langle U_1^*U_2\rangle$ 的定义，它已经包含两路的全部相位。若把可控的传播相位差单独记为 $\Delta\varphi$，就须另外定义去掉这一相位后的复相干度 $\gamma_{12}^{(0)}$，使
\[
 \Gamma_{12}=\sqrt{I_1I_2}\,\gamma_{12}^{(0)}e^{\mathrm i\Delta\varphi},\qquad
 \gamma_{12}=\gamma_{12}^{(0)}e^{\mathrm i\Delta\varphi}.
\]
因此平均强度为
\[
 I=I_1+I_2+2\sqrt{I_1I_2}\,|\gamma_{12}^{(0)}|
 \cos\!\left(\Delta\varphi+\arg\gamma_{12}^{(0)}\right).
\]
两种复相干度的模相同；不能把已包含传播相位的 $\arg\gamma_{12}$ 再与同一 $\Delta\varphi$ 相加。
于是可见度是
\[
 \mathcal V=\frac{I_{\max}-I_{\min}}{I_{\max}+I_{\min}}
 =\frac{2\sqrt{I_1I_2}}{I_1+I_2}|\gamma_{12}|.
\]
即使完全相干，严重不等强的两束光也只有很低的可见度。反过来，平均强度相等也不能保证相干：独立源点之间若没有固定相位，应叠加强度，各自条纹的相位平移会冲淡总图样。

有限谱宽造成的是时间相干性的限制。若两条光路相差 $\Delta\mathcal L$，相应时间延迟为 $\Delta t=\Delta\mathcal L/c$；不同频率分量的相位差随频率改变。当 $\Delta\nu\,|\Delta t|$ 达到一的量级，各频率的条纹不再对齐，反差通常明显下降。因此常用的相干长度量级为 $\ell_c\sim c/\Delta\nu$，但精确数值和衰减曲线取决于谱线形状与采用的宽度定义。

这些限制解释了白光条纹为什么主要出现在近零光程差区域，也解释了为什么窄光谱并不能补救任意宽的空间非相干光源。做题时应先确定几何光程和反射相位，再按题给光源模型进行空间或频率平均；把所有实验困难统称为“相位不同”会遗漏具体可计算的约束。

\endgroup
\begingroup\tcbset{breakable/.default=false}
\renewcommand{\Needspace}[1]{\par\begingroup
\dimen0=#1\relax\vskip0pt plus\dimen0\penalty-100
\vskip0pt plus-\dimen0\vskip\dimen0\penalty9999\vskip-\dimen0
\vskip0pt\endgroup}
\chapter{衍射、偏振与晶体光学}


\graphicspath{{/Users/dongxingxi/Documents/物理竞赛一轮资料/sample/xes-original/forty-fifth-section/}{/Users/dongxingxi/Documents/物理竞赛一轮资料/sample/xes-original/forty-fifth-section/assets/}}


\section{衍射现象与波动模型}


% Geometry: shared coordinates; transparent labels; numerical profiles where required.
\newcommand{\BasicDiffractionPanel}{%
\begin{center}\begin{minipage}{\linewidth}\centering\begin{tikzpicture}[scale=.8,>=Latex,font=\small]
\begin{scope}[xshift=-4cm]
 \fill[SoftGray](-.12,-.8)rectangle(.12,.8);
 \draw[very thick](0,-.8)--(0,.8);
 \fill(-2.5,0)circle(.04)node[left]{$S$};
 \foreach \y in {-.8,.8}{\draw[DeepBlue,->](-2.5,0)--(0,\y);
 \draw[dashed,TextGray](0,\y)--(3,{2.2*\y});
 \foreach \r in {.4,.8,1.2}{\draw[Blue!65]([shift={(-50:\r)}]0,\y)arc(-50:50:\r);}}
 \draw[very thick](3,-2)--(3,2)node[above]{屏};
 \node[below]at(.3,-2.3){次级波来自开放区域及边缘};
\end{scope}
\begin{scope}[xshift=3.7cm]
 \draw[very thick](0,-1.8)--(0,-.3) (0,.3)--(0,1.8);
 \foreach \y in {-.2,0,.2}{\draw[DeepBlue,->](-2,\y)--(0,\y);}
 \foreach \r in {.55,1,1.5,2}{\draw[Blue!70](40:\r)arc(40:-40:\r);}
 \node[above]at(0,1.9){孔径};
 \node[below,align=center]at(.2,-2.3){孔内各微元波相干叠加\\不能当作几条折射光线};
\end{scope}
\end{tikzpicture}\figcap{不透明障碍与孔径改变开放波面的范围；虚线只标几何阴影边界。}\label{ch30:fig:s45-22-1}\end{minipage}\end{center}
}

% Geometry: shared coordinates; transparent labels; numerical profiles where required.
\newcommand{\GratingDiagram}{%
\begin{center}\begin{minipage}{\linewidth}\centering\begin{tikzpicture}[scale=.9,>=Latex,font=\small]
 \def\ang{35}\coordinate(A)at(0,0);\coordinate(B)at(2,0);
 \coordinate(H)at({2*sin(\ang)^2},{2*sin(\ang)*cos(\ang)});
 \foreach \x in {-2,0,2,4}{\draw[very thick](\x-.35,0)--(\x-.09,0) (\x+.09,0)--(\x+.35,0);
 \draw[DeepBlue,->](\x,-1.25)--(\x,-.1);
 \draw[Blue,->](\x,0)--({\x+2.5*sin(\ang)},{2.5*cos(\ang)});}
 \draw[dashed](0,0)--(0,2.7);
 \draw[dashed,TextGray](B)--(H);
 \draw[Gold,very thick](A)--(H);
 \draw(0,.8)arc(90:55:.8);\node at(.18,.52){$\theta$};
 \node[above left]at(A){$A$};\node[above right]at(B){$B$};
 \node[above left]at(H){$H$};
 \draw[<->](0,-1.65)--(2,-1.65)node[midway,below]{$AB=d$};
 \node[anchor=west,align=left]at(3,2.75){$BH\perp AH$\\$AH=d\sin\theta$};
\end{tikzpicture}\figcap{在共同出射波前 $BH$ 上比较相位，程差的大小等于投影 $AH$。}\label{ch30:fig:s45-22-2}\end{minipage}\end{center}
}

\newcommand{\SpectrumDiagram}{%
\begin{center}\begin{minipage}{\linewidth}\centering\begin{tikzpicture}[scale=.82,>=Latex]
 \draw[very thick](0,-1.5)--(0,1.5);\draw[DeepBlue,very thick,->](-3,0)--(0,0);
 \draw[TextGray,thick,->](0,0)--(3.8,0)node[right]{零级};
 \foreach \a/\c in {-35/red!70!black,-25/green!50!black,-18/blue!70!black,18/blue!70!black,25/green!50!black,35/red!70!black}{\draw[very thick,\c,->](0,0)--({3*cos(\a)},{3*sin(\a)});}
 \node[red!70!black,anchor=west] at (2.65,1.95) {红光};
 \node[green!50!black,anchor=west] at (3.05,1.28) {绿光};
 \node[blue!70!black,anchor=west] at (3.25,.63) {紫蓝光};
 \node[blue!70!black,anchor=west] at (3.25,-.63) {紫蓝光};
 \node[green!50!black,anchor=west] at (3.05,-1.28) {绿光};
 \node[red!70!black,anchor=west] at (2.65,-1.95) {红光};
 \node[below]at(0,-2.65){同一级中长波长偏转角更大};
\end{tikzpicture}\figcap{白光通过光栅后按波长分离。}\label{ch30:fig:s45-22-3}\end{minipage}\end{center}}

% Geometry: shared coordinates; transparent labels; numerical profiles where required.
\newcommand{\GratingCurves}{%
\begin{center}\begin{minipage}{\linewidth}\centering\begin{tikzpicture}[x=1cm,y=1.1cm,>=Latex,font=\small]
 \foreach \shift/\N in {0/5,4.5/10,9/50}{
 \begin{scope}[xshift=\shift cm]
 \draw[->](-1.5,0)--(1.6,0)node[right]{$u$};
 \draw[->](-1.5,0)--(-1.5,1.8);
 \draw[DeepBlue,thick]plot[samples=2001,domain=-1.4:1.4,variable=\t]
 (\t,{1.5*(abs(sin(180*\t))<0.000001 ? 1 : (sin(180*\N*\t)/(\N*sin(180*\t)))^2)});
 \foreach \u in {-1,0,1}{\draw(\u,0)--(\u,-.05)node[below]{$\u$};}
 \node at(0,2.05){$N=\N$};
 \end{scope}}
\end{tikzpicture}\figcap{$u=\delta/(2\pi)$，纵坐标为 $I/(N^2I_1)$；相邻整数位置为主极大，其间有 $N-1$ 个零点与 $N-2$ 个次极大。}\label{ch30:fig:s45-22-4}\end{minipage}\end{center}
}

% Geometry: shared coordinates; transparent labels; numerical profiles where required.
\newcommand{\BraggDiagram}{%
\begin{center}\begin{minipage}{\linewidth}\centering\begin{tikzpicture}[scale=.9,>=Latex,font=\small]
 \coordinate(A)at(0,0);\coordinate(B)at(0,-1.2);
 \coordinate(H)at(-.6,-.6);\coordinate(K)at(.6,-.6);
 \foreach \y in {0,-1.2,-2.4}{\draw[RuleGray](-3.5,\y)--(3.5,\y);
 \foreach \x in {-3,-2.4,...,3}{\fill[DeepBlue](\x,\y)circle(.045);}}
 \draw[DeepBlue,very thick,->](-2.3,2.3)--(A)--(2.3,2.3);
 \draw[Blue,very thick,->](-2.3,1.1)--(B)--(2.3,1.1);
 \draw[dashed](A)--(H) (A)--(K);
 \draw[Gold,very thick](H)--(B)--(K);
 \draw[<->](-3.1,0)--(-3.1,-1.2)node[midway,left]{$d$};
 \draw(-.7,0)arc(180:135:.7);\node at(-.85,.4){$\theta$};
 \node[above]at(A){$A$};\node[below]at(B){$B$};
 \node[left]at(H){$H$};\node[right]at(K){$K$};
 \node[align=left,anchor=west]at(3,1.4){$AH\perp HB$\\$AK\perp BK$\\$HB+BK=2d\sin\theta$};
\end{tikzpicture}\figcap{布拉格几何中，$\theta$ 从晶面量起；$HB$ 与 $BK$ 是相对共同入射、出射波前的两段额外光程。}\label{ch30:fig:s45-22-7}\end{minipage}\end{center}
}

\newcommand{\SingleSlitCurve}{%
\begin{center}\begin{minipage}{\linewidth}\centering\begin{tikzpicture}[x=.8cm,y=1.5cm]
 \draw[->](-6,0)--(6,0)node[right]{$D\sin\theta/\lambda$};\draw[->](0,0)--(0,1.25);
 \node at(-.7,1.1){$I/I_0$};
 \draw[DeepBlue,very thick]plot[samples=401,domain=-5.8:5.8,variable=\t]
   (\t,{abs(\t)<0.001 ? 1 : (sin(180*\t)/(3.14159*\t))^2});
 \foreach \x in {-3,-2,-1,1,2,3}{\draw(\x,.04)--(\x,-.04)node[below]{$\x$};}
 \foreach \j in {-115,...,115}{
 \pgfmathsetmacro{\u}{\j/20}
 \ifnum\j=0\def\bright{100}\else
 \pgfmathsetmacro{\bright}{100*(sin(180*\u)/(pi*\u))^2}\fi
 \fill[white!\bright!black](\u,1.4)rectangle({\u+.051},1.56);}
 \node[below]at(0,-.52){中央主极大两侧第一零点为 $-1,+1$};
\end{tikzpicture}\figcap{矩形单缝的 $\operatorname{sinc}^2$ 强度分布。}\label{ch30:fig:s45-22-8}\end{minipage}\end{center}}

\newcommand{\ResolutionPanel}{%
\begin{center}\begin{minipage}{\linewidth}\centering\begin{tikzpicture}[x=.65cm,y=1.1cm,>=Latex,font=\small]
\begin{scope}[xshift=0cm]
\fill[DeepBlue!7.16](-2.0000,1.5500)rectangle(-1.9050,1.6085);
\fill[DeepBlue!8.14](-2.0000,1.6065)rectangle(-1.9050,1.6650);
\fill[DeepBlue!8.88](-2.0000,1.6630)rectangle(-1.9050,1.7216);
\fill[DeepBlue!9.36](-2.0000,1.7196)rectangle(-1.9050,1.7781);
\fill[DeepBlue!9.59](-2.0000,1.7761)rectangle(-1.9050,1.8346);
\fill[DeepBlue!9.59](-2.0000,1.8326)rectangle(-1.9050,1.8911);
\fill[DeepBlue!9.41](-2.0000,1.8891)rectangle(-1.9050,1.9477);
\fill[DeepBlue!9.12](-2.0000,1.9457)rectangle(-1.9050,2.0042);
\fill[DeepBlue!8.78](-2.0000,2.0022)rectangle(-1.9050,2.0607);
\fill[DeepBlue!8.46](-2.0000,2.0587)rectangle(-1.9050,2.1172);
\fill[DeepBlue!8.21](-2.0000,2.1152)rectangle(-1.9050,2.1737);
\fill[DeepBlue!8.07](-2.0000,2.1717)rectangle(-1.9050,2.2303);
\fill[DeepBlue!8.07](-2.0000,2.2283)rectangle(-1.9050,2.2868);
\fill[DeepBlue!8.21](-2.0000,2.2848)rectangle(-1.9050,2.3433);
\fill[DeepBlue!8.46](-2.0000,2.3413)rectangle(-1.9050,2.3998);
\fill[DeepBlue!8.78](-2.0000,2.3978)rectangle(-1.9050,2.4563);
\fill[DeepBlue!9.12](-2.0000,2.4543)rectangle(-1.9050,2.5129);
\fill[DeepBlue!9.41](-2.0000,2.5109)rectangle(-1.9050,2.5694);
\fill[DeepBlue!9.59](-2.0000,2.5674)rectangle(-1.9050,2.6259);
\fill[DeepBlue!9.59](-2.0000,2.6239)rectangle(-1.9050,2.6824);
\fill[DeepBlue!9.36](-2.0000,2.6804)rectangle(-1.9050,2.7390);
\fill[DeepBlue!8.88](-2.0000,2.7370)rectangle(-1.9050,2.7955);
\fill[DeepBlue!8.14](-2.0000,2.7935)rectangle(-1.9050,2.8520);
\fill[DeepBlue!8.32](-1.9070,1.5500)rectangle(-1.8120,1.6085);
\fill[DeepBlue!9.05](-1.9070,1.6065)rectangle(-1.8120,1.6650);
\fill[DeepBlue!9.45](-1.9070,1.6630)rectangle(-1.8120,1.7216);
\fill[DeepBlue!9.51](-1.9070,1.7196)rectangle(-1.8120,1.7781);
\fill[DeepBlue!9.26](-1.9070,1.7761)rectangle(-1.8120,1.8346);
\fill[DeepBlue!8.76](-1.9070,1.8326)rectangle(-1.8120,1.8911);
\fill[DeepBlue!8.09](-1.9070,1.8891)rectangle(-1.8120,1.9477);
\fill[DeepBlue!7.32](-1.9070,1.9457)rectangle(-1.8120,2.0042);
\fill[DeepBlue!6.57](-1.9070,2.0022)rectangle(-1.8120,2.0607);
\fill[DeepBlue!5.91](-1.9070,2.0587)rectangle(-1.8120,2.1172);
\fill[DeepBlue!5.43](-1.9070,2.1152)rectangle(-1.8120,2.1737);
\fill[DeepBlue!5.17](-1.9070,2.1717)rectangle(-1.8120,2.2303);
\fill[DeepBlue!5.17](-1.9070,2.2283)rectangle(-1.8120,2.2868);
\fill[DeepBlue!5.43](-1.9070,2.2848)rectangle(-1.8120,2.3433);
\fill[DeepBlue!5.91](-1.9070,2.3413)rectangle(-1.8120,2.3998);
\fill[DeepBlue!6.57](-1.9070,2.3978)rectangle(-1.8120,2.4563);
\fill[DeepBlue!7.32](-1.9070,2.4543)rectangle(-1.8120,2.5129);
\fill[DeepBlue!8.09](-1.9070,2.5109)rectangle(-1.8120,2.5694);
\fill[DeepBlue!8.76](-1.9070,2.5674)rectangle(-1.8120,2.6259);
\fill[DeepBlue!9.26](-1.9070,2.6239)rectangle(-1.8120,2.6824);
\fill[DeepBlue!9.51](-1.9070,2.6804)rectangle(-1.8120,2.7390);
\fill[DeepBlue!9.45](-1.9070,2.7370)rectangle(-1.8120,2.7955);
\fill[DeepBlue!9.05](-1.9070,2.7935)rectangle(-1.8120,2.8520);
\fill[DeepBlue!9.06](-1.8140,1.5500)rectangle(-1.7189,1.6085);
\fill[DeepBlue!9.42](-1.8140,1.6065)rectangle(-1.7189,1.6650);
\fill[DeepBlue!9.34](-1.8140,1.6630)rectangle(-1.7189,1.7216);
\fill[DeepBlue!8.86](-1.8140,1.7196)rectangle(-1.7189,1.7781);
\fill[DeepBlue!8.01](-1.8140,1.7761)rectangle(-1.7189,1.8346);
\fill[DeepBlue!6.89](-1.8140,1.8326)rectangle(-1.7189,1.8911);
\fill[DeepBlue!5.61](-1.8140,1.8891)rectangle(-1.7189,1.9477);
\fill[DeepBlue!4.29](-1.8140,1.9457)rectangle(-1.7189,2.0042);
\fill[DeepBlue!3.06](-1.8140,2.0022)rectangle(-1.7189,2.0607);
\fill[DeepBlue!2.04](-1.8140,2.0587)rectangle(-1.7189,2.1172);
\fill[DeepBlue!1.33](-1.8140,2.1152)rectangle(-1.7189,2.1737);
\fill[DeepBlue!1.00](-1.8140,2.1717)rectangle(-1.7189,2.2303);
\fill[DeepBlue!1.00](-1.8140,2.2283)rectangle(-1.7189,2.2868);
\fill[DeepBlue!1.33](-1.8140,2.2848)rectangle(-1.7189,2.3433);
\fill[DeepBlue!2.04](-1.8140,2.3413)rectangle(-1.7189,2.3998);
\fill[DeepBlue!3.06](-1.8140,2.3978)rectangle(-1.7189,2.4563);
\fill[DeepBlue!4.29](-1.8140,2.4543)rectangle(-1.7189,2.5129);
\fill[DeepBlue!5.61](-1.8140,2.5109)rectangle(-1.7189,2.5694);
\fill[DeepBlue!6.89](-1.8140,2.5674)rectangle(-1.7189,2.6259);
\fill[DeepBlue!8.01](-1.8140,2.6239)rectangle(-1.7189,2.6824);
\fill[DeepBlue!8.86](-1.8140,2.6804)rectangle(-1.7189,2.7390);
\fill[DeepBlue!9.34](-1.8140,2.7370)rectangle(-1.7189,2.7955);
\fill[DeepBlue!9.42](-1.8140,2.7935)rectangle(-1.7189,2.8520);
\fill[DeepBlue!9.36](-1.7209,1.5500)rectangle(-1.6259,1.6085);
\fill[DeepBlue!9.23](-1.7209,1.6065)rectangle(-1.6259,1.6650);
\fill[DeepBlue!8.57](-1.7209,1.6630)rectangle(-1.6259,1.7216);
\fill[DeepBlue!7.44](-1.7209,1.7196)rectangle(-1.6259,1.7781);
\fill[DeepBlue!5.91](-1.7209,1.7761)rectangle(-1.6259,1.8346);
\fill[DeepBlue!4.14](-1.7209,1.8326)rectangle(-1.6259,1.8911);
\fill[DeepBlue!2.37](-1.7209,1.8891)rectangle(-1.6259,1.9477);
\fill[DeepBlue!1.43](-1.7209,1.9457)rectangle(-1.6259,2.0042);
\fill[DeepBlue!2.40](-1.7209,2.0022)rectangle(-1.6259,2.0607);
\fill[DeepBlue!3.72](-1.7209,2.0587)rectangle(-1.6259,2.1172);
\fill[DeepBlue!4.74](-1.7209,2.1152)rectangle(-1.6259,2.1737);
\fill[DeepBlue!5.28](-1.7209,2.1717)rectangle(-1.6259,2.2303);
\fill[DeepBlue!5.28](-1.7209,2.2283)rectangle(-1.6259,2.2868);
\fill[DeepBlue!4.74](-1.7209,2.2848)rectangle(-1.6259,2.3433);
\fill[DeepBlue!3.72](-1.7209,2.3413)rectangle(-1.6259,2.3998);
\fill[DeepBlue!2.40](-1.7209,2.3978)rectangle(-1.6259,2.4563);
\fill[DeepBlue!1.43](-1.7209,2.4543)rectangle(-1.6259,2.5129);
\fill[DeepBlue!2.37](-1.7209,2.5109)rectangle(-1.6259,2.5694);
\fill[DeepBlue!4.14](-1.7209,2.5674)rectangle(-1.6259,2.6259);
\fill[DeepBlue!5.91](-1.7209,2.6239)rectangle(-1.6259,2.6824);
\fill[DeepBlue!7.44](-1.7209,2.6804)rectangle(-1.6259,2.7390);
\fill[DeepBlue!8.57](-1.7209,2.7370)rectangle(-1.6259,2.7955);
\fill[DeepBlue!9.23](-1.7209,2.7935)rectangle(-1.6259,2.8520);
\fill[DeepBlue!9.24](-1.6279,1.5500)rectangle(-1.5329,1.6085);
\fill[DeepBlue!8.54](-1.6279,1.6065)rectangle(-1.5329,1.6650);
\fill[DeepBlue!7.24](-1.6279,1.6630)rectangle(-1.5329,1.7216);
\fill[DeepBlue!5.45](-1.6279,1.7196)rectangle(-1.5329,1.7781);
\fill[DeepBlue!3.42](-1.6279,1.7761)rectangle(-1.5329,1.8346);
\fill[DeepBlue!2.23](-1.6279,1.8326)rectangle(-1.5329,1.8911);
\fill[DeepBlue!3.63](-1.6279,1.8891)rectangle(-1.5329,1.9477);
\fill[DeepBlue!6.01](-1.6279,1.9457)rectangle(-1.5329,2.0042);
\fill[DeepBlue!8.34](-1.6279,2.0022)rectangle(-1.5329,2.0607);
\fill[DeepBlue!10.30](-1.6279,2.0587)rectangle(-1.5329,2.1172);
\fill[DeepBlue!11.70](-1.6279,2.1152)rectangle(-1.5329,2.1737);
\fill[DeepBlue!12.43](-1.6279,2.1717)rectangle(-1.5329,2.2303);
\fill[DeepBlue!12.43](-1.6279,2.2283)rectangle(-1.5329,2.2868);
\fill[DeepBlue!11.70](-1.6279,2.2848)rectangle(-1.5329,2.3433);
\fill[DeepBlue!10.30](-1.6279,2.3413)rectangle(-1.5329,2.3998);
\fill[DeepBlue!8.34](-1.6279,2.3978)rectangle(-1.5329,2.4563);
\fill[DeepBlue!6.01](-1.6279,2.4543)rectangle(-1.5329,2.5129);
\fill[DeepBlue!3.63](-1.6279,2.5109)rectangle(-1.5329,2.5694);
\fill[DeepBlue!2.23](-1.6279,2.5674)rectangle(-1.5329,2.6259);
\fill[DeepBlue!3.42](-1.6279,2.6239)rectangle(-1.5329,2.6824);
\fill[DeepBlue!5.45](-1.6279,2.6804)rectangle(-1.5329,2.7390);
\fill[DeepBlue!7.24](-1.6279,2.7370)rectangle(-1.5329,2.7955);
\fill[DeepBlue!8.54](-1.6279,2.7935)rectangle(-1.5329,2.8520);
\fill[DeepBlue!8.77](-1.5349,1.5500)rectangle(-1.4399,1.6085);
\fill[DeepBlue!7.48](-1.5349,1.6065)rectangle(-1.4399,1.6650);
\fill[DeepBlue!5.60](-1.5349,1.6630)rectangle(-1.4399,1.7216);
\fill[DeepBlue!3.58](-1.5349,1.7196)rectangle(-1.4399,1.7781);
\fill[DeepBlue!3.25](-1.5349,1.7761)rectangle(-1.4399,1.8346);
\fill[DeepBlue!5.67](-1.5349,1.8326)rectangle(-1.4399,1.8911);
\fill[DeepBlue!8.97](-1.5349,1.8891)rectangle(-1.4399,1.9477);
\fill[DeepBlue!12.32](-1.5349,1.9457)rectangle(-1.4399,2.0042);
\fill[DeepBlue!15.37](-1.5349,2.0022)rectangle(-1.4399,2.0607);
\fill[DeepBlue!17.88](-1.5349,2.0587)rectangle(-1.4399,2.1172);
\fill[DeepBlue!19.66](-1.5349,2.1152)rectangle(-1.4399,2.1737);
\fill[DeepBlue!20.58](-1.5349,2.1717)rectangle(-1.4399,2.2303);
\fill[DeepBlue!20.58](-1.5349,2.2283)rectangle(-1.4399,2.2868);
\fill[DeepBlue!19.66](-1.5349,2.2848)rectangle(-1.4399,2.3433);
\fill[DeepBlue!17.88](-1.5349,2.3413)rectangle(-1.4399,2.3998);
\fill[DeepBlue!15.37](-1.5349,2.3978)rectangle(-1.4399,2.4563);
\fill[DeepBlue!12.32](-1.5349,2.4543)rectangle(-1.4399,2.5129);
\fill[DeepBlue!8.97](-1.5349,2.5109)rectangle(-1.4399,2.5694);
\fill[DeepBlue!5.67](-1.5349,2.5674)rectangle(-1.4399,2.6259);
\fill[DeepBlue!3.25](-1.5349,2.6239)rectangle(-1.4399,2.6824);
\fill[DeepBlue!3.58](-1.5349,2.6804)rectangle(-1.4399,2.7390);
\fill[DeepBlue!5.60](-1.5349,2.7370)rectangle(-1.4399,2.7955);
\fill[DeepBlue!7.48](-1.5349,2.7935)rectangle(-1.4399,2.8520);
\fill[DeepBlue!8.08](-1.4419,1.5500)rectangle(-1.3468,1.6085);
\fill[DeepBlue!6.27](-1.4419,1.6065)rectangle(-1.3468,1.6650);
\fill[DeepBlue!4.24](-1.4419,1.6630)rectangle(-1.3468,1.7216);
\fill[DeepBlue!3.95](-1.4419,1.7196)rectangle(-1.3468,1.7781);
\fill[DeepBlue!6.73](-1.4419,1.7761)rectangle(-1.3468,1.8346);
\fill[DeepBlue!10.76](-1.4419,1.8326)rectangle(-1.3468,1.8911);
\fill[DeepBlue!15.08](-1.4419,1.8891)rectangle(-1.3468,1.9477);
\fill[DeepBlue!19.27](-1.4419,1.9457)rectangle(-1.3468,2.0042);
\fill[DeepBlue!23.02](-1.4419,2.0022)rectangle(-1.3468,2.0607);
\fill[DeepBlue!26.08](-1.4419,2.0587)rectangle(-1.3468,2.1172);
\fill[DeepBlue!28.23](-1.4419,2.1152)rectangle(-1.3468,2.1737);
\fill[DeepBlue!29.35](-1.4419,2.1717)rectangle(-1.3468,2.2303);
\fill[DeepBlue!29.35](-1.4419,2.2283)rectangle(-1.3468,2.2868);
\fill[DeepBlue!28.23](-1.4419,2.2848)rectangle(-1.3468,2.3433);
\fill[DeepBlue!26.08](-1.4419,2.3413)rectangle(-1.3468,2.3998);
\fill[DeepBlue!23.02](-1.4419,2.3978)rectangle(-1.3468,2.4563);
\fill[DeepBlue!19.27](-1.4419,2.4543)rectangle(-1.3468,2.5129);
\fill[DeepBlue!15.08](-1.4419,2.5109)rectangle(-1.3468,2.5694);
\fill[DeepBlue!10.76](-1.4419,2.5674)rectangle(-1.3468,2.6259);
\fill[DeepBlue!6.73](-1.4419,2.6239)rectangle(-1.3468,2.6824);
\fill[DeepBlue!3.95](-1.4419,2.6804)rectangle(-1.3468,2.7390);
\fill[DeepBlue!4.24](-1.4419,2.7370)rectangle(-1.3468,2.7955);
\fill[DeepBlue!6.27](-1.4419,2.7935)rectangle(-1.3468,2.8520);
\fill[DeepBlue!7.27](-1.3488,1.5500)rectangle(-1.2538,1.6085);
\fill[DeepBlue!5.20](-1.3488,1.6065)rectangle(-1.2538,1.6650);
\fill[DeepBlue!4.24](-1.3488,1.6630)rectangle(-1.2538,1.7216);
\fill[DeepBlue!6.68](-1.3488,1.7196)rectangle(-1.2538,1.7781);
\fill[DeepBlue!11.10](-1.3488,1.7761)rectangle(-1.2538,1.8346);
\fill[DeepBlue!16.20](-1.3488,1.8326)rectangle(-1.2538,1.8911);
\fill[DeepBlue!21.44](-1.3488,1.8891)rectangle(-1.2538,1.9477);
\fill[DeepBlue!26.44](-1.3488,1.9457)rectangle(-1.2538,2.0042);
\fill[DeepBlue!30.87](-1.3488,2.0022)rectangle(-1.2538,2.0607);
\fill[DeepBlue!34.46](-1.3488,2.0587)rectangle(-1.2538,2.1172);
\fill[DeepBlue!36.99](-1.3488,2.1152)rectangle(-1.2538,2.1737);
\fill[DeepBlue!38.29](-1.3488,2.1717)rectangle(-1.2538,2.2303);
\fill[DeepBlue!38.29](-1.3488,2.2283)rectangle(-1.2538,2.2868);
\fill[DeepBlue!36.99](-1.3488,2.2848)rectangle(-1.2538,2.3433);
\fill[DeepBlue!34.46](-1.3488,2.3413)rectangle(-1.2538,2.3998);
\fill[DeepBlue!30.87](-1.3488,2.3978)rectangle(-1.2538,2.4563);
\fill[DeepBlue!26.44](-1.3488,2.4543)rectangle(-1.2538,2.5129);
\fill[DeepBlue!21.44](-1.3488,2.5109)rectangle(-1.2538,2.5694);
\fill[DeepBlue!16.20](-1.3488,2.5674)rectangle(-1.2538,2.6259);
\fill[DeepBlue!11.10](-1.3488,2.6239)rectangle(-1.2538,2.6824);
\fill[DeepBlue!6.68](-1.3488,2.6804)rectangle(-1.2538,2.7390);
\fill[DeepBlue!4.24](-1.3488,2.7370)rectangle(-1.2538,2.7955);
\fill[DeepBlue!5.20](-1.3488,2.7935)rectangle(-1.2538,2.8520);
\fill[DeepBlue!6.43](-1.2558,1.5500)rectangle(-1.1608,1.6085);
\fill[DeepBlue!4.61](-1.2558,1.6065)rectangle(-1.1608,1.6650);
\fill[DeepBlue!5.73](-1.2558,1.6630)rectangle(-1.1608,1.7216);
\fill[DeepBlue!10.00](-1.2558,1.7196)rectangle(-1.1608,1.7781);
\fill[DeepBlue!15.54](-1.2558,1.7761)rectangle(-1.1608,1.8346);
\fill[DeepBlue!21.58](-1.2558,1.8326)rectangle(-1.1608,1.8911);
\fill[DeepBlue!27.66](-1.2558,1.8891)rectangle(-1.1608,1.9477);
\fill[DeepBlue!33.41](-1.2558,1.9457)rectangle(-1.1608,2.0042);
\fill[DeepBlue!38.48](-1.2558,2.0022)rectangle(-1.1608,2.0607);
\fill[DeepBlue!42.58](-1.2558,2.0587)rectangle(-1.1608,2.1172);
\fill[DeepBlue!45.46](-1.2558,2.1152)rectangle(-1.1608,2.1737);
\fill[DeepBlue!46.94](-1.2558,2.1717)rectangle(-1.1608,2.2303);
\fill[DeepBlue!46.94](-1.2558,2.2283)rectangle(-1.1608,2.2868);
\fill[DeepBlue!45.46](-1.2558,2.2848)rectangle(-1.1608,2.3433);
\fill[DeepBlue!42.58](-1.2558,2.3413)rectangle(-1.1608,2.3998);
\fill[DeepBlue!38.48](-1.2558,2.3978)rectangle(-1.1608,2.4563);
\fill[DeepBlue!33.41](-1.2558,2.4543)rectangle(-1.1608,2.5129);
\fill[DeepBlue!27.66](-1.2558,2.5109)rectangle(-1.1608,2.5694);
\fill[DeepBlue!21.58](-1.2558,2.5674)rectangle(-1.1608,2.6259);
\fill[DeepBlue!15.54](-1.2558,2.6239)rectangle(-1.1608,2.6824);
\fill[DeepBlue!10.00](-1.2558,2.6804)rectangle(-1.1608,2.7390);
\fill[DeepBlue!5.73](-1.2558,2.7370)rectangle(-1.1608,2.7955);
\fill[DeepBlue!4.61](-1.2558,2.7935)rectangle(-1.1608,2.8520);
\fill[DeepBlue!5.62](-1.1628,1.5500)rectangle(-1.0678,1.6085);
\fill[DeepBlue!4.64](-1.1628,1.6065)rectangle(-1.0678,1.6650);
\fill[DeepBlue!7.73](-1.1628,1.6630)rectangle(-1.0678,1.7216);
\fill[DeepBlue!13.21](-1.1628,1.7196)rectangle(-1.0678,1.7781);
\fill[DeepBlue!19.67](-1.1628,1.7761)rectangle(-1.0678,1.8346);
\fill[DeepBlue!26.54](-1.1628,1.8326)rectangle(-1.0678,1.8911);
\fill[DeepBlue!33.37](-1.1628,1.8891)rectangle(-1.0678,1.9477);
\fill[DeepBlue!39.79](-1.1628,1.9457)rectangle(-1.0678,2.0042);
\fill[DeepBlue!45.44](-1.1628,2.0022)rectangle(-1.0678,2.0607);
\fill[DeepBlue!49.98](-1.1628,2.0587)rectangle(-1.0678,2.1172);
\fill[DeepBlue!53.17](-1.1628,2.1152)rectangle(-1.0678,2.1737);
\fill[DeepBlue!54.81](-1.1628,2.1717)rectangle(-1.0678,2.2303);
\fill[DeepBlue!54.81](-1.1628,2.2283)rectangle(-1.0678,2.2868);
\fill[DeepBlue!53.17](-1.1628,2.2848)rectangle(-1.0678,2.3433);
\fill[DeepBlue!49.98](-1.1628,2.3413)rectangle(-1.0678,2.3998);
\fill[DeepBlue!45.44](-1.1628,2.3978)rectangle(-1.0678,2.4563);
\fill[DeepBlue!39.79](-1.1628,2.4543)rectangle(-1.0678,2.5129);
\fill[DeepBlue!33.37](-1.1628,2.5109)rectangle(-1.0678,2.5694);
\fill[DeepBlue!26.54](-1.1628,2.5674)rectangle(-1.0678,2.6259);
\fill[DeepBlue!19.67](-1.1628,2.6239)rectangle(-1.0678,2.6824);
\fill[DeepBlue!13.21](-1.1628,2.6804)rectangle(-1.0678,2.7390);
\fill[DeepBlue!7.73](-1.1628,2.7370)rectangle(-1.0678,2.7955);
\fill[DeepBlue!4.64](-1.1628,2.7935)rectangle(-1.0678,2.8520);
\fill[DeepBlue!4.77](-1.0698,1.5500)rectangle(-0.9747,1.6085);
\fill[DeepBlue!5.03](-1.0698,1.6065)rectangle(-0.9747,1.6650);
\fill[DeepBlue!9.61](-1.0698,1.6630)rectangle(-0.9747,1.7216);
\fill[DeepBlue!16.00](-1.0698,1.7196)rectangle(-0.9747,1.7781);
\fill[DeepBlue!23.21](-1.0698,1.7761)rectangle(-0.9747,1.8346);
\fill[DeepBlue!30.76](-1.0698,1.8326)rectangle(-0.9747,1.8911);
\fill[DeepBlue!38.23](-1.0698,1.8891)rectangle(-0.9747,1.9477);
\fill[DeepBlue!45.21](-1.0698,1.9457)rectangle(-0.9747,2.0042);
\fill[DeepBlue!51.33](-1.0698,2.0022)rectangle(-0.9747,2.0607);
\fill[DeepBlue!56.25](-1.0698,2.0587)rectangle(-0.9747,2.1172);
\fill[DeepBlue!59.69](-1.0698,2.1152)rectangle(-0.9747,2.1737);
\fill[DeepBlue!61.46](-1.0698,2.1717)rectangle(-0.9747,2.2303);
\fill[DeepBlue!61.46](-1.0698,2.2283)rectangle(-0.9747,2.2868);
\fill[DeepBlue!59.69](-1.0698,2.2848)rectangle(-0.9747,2.3433);
\fill[DeepBlue!56.25](-1.0698,2.3413)rectangle(-0.9747,2.3998);
\fill[DeepBlue!51.33](-1.0698,2.3978)rectangle(-0.9747,2.4563);
\fill[DeepBlue!45.21](-1.0698,2.4543)rectangle(-0.9747,2.5129);
\fill[DeepBlue!38.23](-1.0698,2.5109)rectangle(-0.9747,2.5694);
\fill[DeepBlue!30.76](-1.0698,2.5674)rectangle(-0.9747,2.6259);
\fill[DeepBlue!23.21](-1.0698,2.6239)rectangle(-0.9747,2.6824);
\fill[DeepBlue!16.00](-1.0698,2.6804)rectangle(-0.9747,2.7390);
\fill[DeepBlue!9.61](-1.0698,2.7370)rectangle(-0.9747,2.7955);
\fill[DeepBlue!5.03](-1.0698,2.7935)rectangle(-0.9747,2.8520);
\fill[DeepBlue!3.79](-0.9767,1.5500)rectangle(-0.8817,1.6085);
\fill[DeepBlue!5.42](-0.9767,1.6065)rectangle(-0.8817,1.6650);
\fill[DeepBlue!11.07](-0.9767,1.6630)rectangle(-0.8817,1.7216);
\fill[DeepBlue!18.13](-0.9767,1.7196)rectangle(-0.8817,1.7781);
\fill[DeepBlue!25.92](-0.9767,1.7761)rectangle(-0.8817,1.8346);
\fill[DeepBlue!33.99](-0.9767,1.8326)rectangle(-0.8817,1.8911);
\fill[DeepBlue!41.93](-0.9767,1.8891)rectangle(-0.8817,1.9477);
\fill[DeepBlue!49.34](-0.9767,1.9457)rectangle(-0.8817,2.0042);
\fill[DeepBlue!55.81](-0.9767,2.0022)rectangle(-0.8817,2.0607);
\fill[DeepBlue!61.01](-0.9767,2.0587)rectangle(-0.8817,2.1172);
\fill[DeepBlue!64.64](-0.9767,2.1152)rectangle(-0.8817,2.1737);
\fill[DeepBlue!66.51](-0.9767,2.1717)rectangle(-0.8817,2.2303);
\fill[DeepBlue!66.51](-0.9767,2.2283)rectangle(-0.8817,2.2868);
\fill[DeepBlue!64.64](-0.9767,2.2848)rectangle(-0.8817,2.3433);
\fill[DeepBlue!61.01](-0.9767,2.3413)rectangle(-0.8817,2.3998);
\fill[DeepBlue!55.81](-0.9767,2.3978)rectangle(-0.8817,2.4563);
\fill[DeepBlue!49.34](-0.9767,2.4543)rectangle(-0.8817,2.5129);
\fill[DeepBlue!41.93](-0.9767,2.5109)rectangle(-0.8817,2.5694);
\fill[DeepBlue!33.99](-0.9767,2.5674)rectangle(-0.8817,2.6259);
\fill[DeepBlue!25.92](-0.9767,2.6239)rectangle(-0.8817,2.6824);
\fill[DeepBlue!18.13](-0.9767,2.6804)rectangle(-0.8817,2.7390);
\fill[DeepBlue!11.07](-0.9767,2.7370)rectangle(-0.8817,2.7955);
\fill[DeepBlue!5.42](-0.9767,2.7935)rectangle(-0.8817,2.8520);
\fill[DeepBlue!2.57](-0.8837,1.5500)rectangle(-0.7887,1.6085);
\fill[DeepBlue!5.60](-0.8837,1.6065)rectangle(-0.7887,1.6650);
\fill[DeepBlue!11.99](-0.8837,1.6630)rectangle(-0.7887,1.7216);
\fill[DeepBlue!19.49](-0.8837,1.7196)rectangle(-0.7887,1.7781);
\fill[DeepBlue!27.64](-0.8837,1.7761)rectangle(-0.7887,1.8346);
\fill[DeepBlue!36.03](-0.8837,1.8326)rectangle(-0.7887,1.8911);
\fill[DeepBlue!44.27](-0.8837,1.8891)rectangle(-0.7887,1.9477);
\fill[DeepBlue!51.93](-0.8837,1.9457)rectangle(-0.7887,2.0042);
\fill[DeepBlue!58.62](-0.8837,2.0022)rectangle(-0.7887,2.0607);
\fill[DeepBlue!63.98](-0.8837,2.0587)rectangle(-0.7887,2.1172);
\fill[DeepBlue!67.73](-0.8837,2.1152)rectangle(-0.7887,2.1737);
\fill[DeepBlue!69.66](-0.8837,2.1717)rectangle(-0.7887,2.2303);
\fill[DeepBlue!69.66](-0.8837,2.2283)rectangle(-0.7887,2.2868);
\fill[DeepBlue!67.73](-0.8837,2.2848)rectangle(-0.7887,2.3433);
\fill[DeepBlue!63.98](-0.8837,2.3413)rectangle(-0.7887,2.3998);
\fill[DeepBlue!58.62](-0.8837,2.3978)rectangle(-0.7887,2.4563);
\fill[DeepBlue!51.93](-0.8837,2.4543)rectangle(-0.7887,2.5129);
\fill[DeepBlue!44.27](-0.8837,2.5109)rectangle(-0.7887,2.5694);
\fill[DeepBlue!36.03](-0.8837,2.5674)rectangle(-0.7887,2.6259);
\fill[DeepBlue!27.64](-0.8837,2.6239)rectangle(-0.7887,2.6824);
\fill[DeepBlue!19.49](-0.8837,2.6804)rectangle(-0.7887,2.7390);
\fill[DeepBlue!11.99](-0.8837,2.7370)rectangle(-0.7887,2.7955);
\fill[DeepBlue!5.60](-0.8837,2.7935)rectangle(-0.7887,2.8520);
\fill[DeepBlue!1.06](-0.7907,1.5500)rectangle(-0.6957,1.6085);
\fill[DeepBlue!5.58](-0.7907,1.6065)rectangle(-0.6957,1.6650);
\fill[DeepBlue!12.33](-0.7907,1.6630)rectangle(-0.6957,1.7216);
\fill[DeepBlue!20.00](-0.7907,1.7196)rectangle(-0.6957,1.7781);
\fill[DeepBlue!28.28](-0.7907,1.7761)rectangle(-0.6957,1.8346);
\fill[DeepBlue!36.77](-0.7907,1.8326)rectangle(-0.6957,1.8911);
\fill[DeepBlue!45.09](-0.7907,1.8891)rectangle(-0.6957,1.9477);
\fill[DeepBlue!52.83](-0.7907,1.9457)rectangle(-0.6957,2.0042);
\fill[DeepBlue!59.58](-0.7907,2.0022)rectangle(-0.6957,2.0607);
\fill[DeepBlue!64.98](-0.7907,2.0587)rectangle(-0.6957,2.1172);
\fill[DeepBlue!68.76](-0.7907,2.1152)rectangle(-0.6957,2.1737);
\fill[DeepBlue!70.71](-0.7907,2.1717)rectangle(-0.6957,2.2303);
\fill[DeepBlue!70.71](-0.7907,2.2283)rectangle(-0.6957,2.2868);
\fill[DeepBlue!68.76](-0.7907,2.2848)rectangle(-0.6957,2.3433);
\fill[DeepBlue!64.98](-0.7907,2.3413)rectangle(-0.6957,2.3998);
\fill[DeepBlue!59.58](-0.7907,2.3978)rectangle(-0.6957,2.4563);
\fill[DeepBlue!52.83](-0.7907,2.4543)rectangle(-0.6957,2.5129);
\fill[DeepBlue!45.09](-0.7907,2.5109)rectangle(-0.6957,2.5694);
\fill[DeepBlue!36.77](-0.7907,2.5674)rectangle(-0.6957,2.6259);
\fill[DeepBlue!28.28](-0.7907,2.6239)rectangle(-0.6957,2.6824);
\fill[DeepBlue!20.00](-0.7907,2.6804)rectangle(-0.6957,2.7390);
\fill[DeepBlue!12.33](-0.7907,2.7370)rectangle(-0.6957,2.7955);
\fill[DeepBlue!5.58](-0.7907,2.7935)rectangle(-0.6957,2.8520);
\fill[DeepBlue!0.77](-0.6977,1.5500)rectangle(-0.6027,1.6085);
\fill[DeepBlue!5.50](-0.6977,1.6065)rectangle(-0.6027,1.6650);
\fill[DeepBlue!12.15](-0.6977,1.6630)rectangle(-0.6027,1.7216);
\fill[DeepBlue!19.69](-0.6977,1.7196)rectangle(-0.6027,1.7781);
\fill[DeepBlue!27.82](-0.6977,1.7761)rectangle(-0.6027,1.8346);
\fill[DeepBlue!36.17](-0.6977,1.8326)rectangle(-0.6027,1.8911);
\fill[DeepBlue!44.35](-0.6977,1.8891)rectangle(-0.6027,1.9477);
\fill[DeepBlue!51.96](-0.6977,1.9457)rectangle(-0.6027,2.0042);
\fill[DeepBlue!58.61](-0.6977,2.0022)rectangle(-0.6027,2.0607);
\fill[DeepBlue!63.93](-0.6977,2.0587)rectangle(-0.6027,2.1172);
\fill[DeepBlue!67.65](-0.6977,2.1152)rectangle(-0.6027,2.1737);
\fill[DeepBlue!69.56](-0.6977,2.1717)rectangle(-0.6027,2.2303);
\fill[DeepBlue!69.56](-0.6977,2.2283)rectangle(-0.6027,2.2868);
\fill[DeepBlue!67.65](-0.6977,2.2848)rectangle(-0.6027,2.3433);
\fill[DeepBlue!63.93](-0.6977,2.3413)rectangle(-0.6027,2.3998);
\fill[DeepBlue!58.61](-0.6977,2.3978)rectangle(-0.6027,2.4563);
\fill[DeepBlue!51.96](-0.6977,2.4543)rectangle(-0.6027,2.5129);
\fill[DeepBlue!44.35](-0.6977,2.5109)rectangle(-0.6027,2.5694);
\fill[DeepBlue!36.17](-0.6977,2.5674)rectangle(-0.6027,2.6259);
\fill[DeepBlue!27.82](-0.6977,2.6239)rectangle(-0.6027,2.6824);
\fill[DeepBlue!19.69](-0.6977,2.6804)rectangle(-0.6027,2.7390);
\fill[DeepBlue!12.15](-0.6977,2.7370)rectangle(-0.6027,2.7955);
\fill[DeepBlue!5.50](-0.6977,2.7935)rectangle(-0.6027,2.8520);
\fill[DeepBlue!2.83](-0.6047,1.5500)rectangle(-0.5096,1.6085);
\fill[DeepBlue!5.68](-0.6047,1.6065)rectangle(-0.5096,1.6650);
\fill[DeepBlue!11.59](-0.6047,1.6630)rectangle(-0.5096,1.7216);
\fill[DeepBlue!18.63](-0.6047,1.7196)rectangle(-0.5096,1.7781);
\fill[DeepBlue!26.31](-0.6047,1.7761)rectangle(-0.5096,1.8346);
\fill[DeepBlue!34.26](-0.6047,1.8326)rectangle(-0.5096,1.8911);
\fill[DeepBlue!42.08](-0.6047,1.8891)rectangle(-0.5096,1.9477);
\fill[DeepBlue!49.36](-0.6047,1.9457)rectangle(-0.5096,2.0042);
\fill[DeepBlue!55.74](-0.6047,2.0022)rectangle(-0.5096,2.0607);
\fill[DeepBlue!60.85](-0.6047,2.0587)rectangle(-0.5096,2.1172);
\fill[DeepBlue!64.42](-0.6047,2.1152)rectangle(-0.5096,2.1737);
\fill[DeepBlue!66.26](-0.6047,2.1717)rectangle(-0.5096,2.2303);
\fill[DeepBlue!66.26](-0.6047,2.2283)rectangle(-0.5096,2.2868);
\fill[DeepBlue!64.42](-0.6047,2.2848)rectangle(-0.5096,2.3433);
\fill[DeepBlue!60.85](-0.6047,2.3413)rectangle(-0.5096,2.3998);
\fill[DeepBlue!55.74](-0.6047,2.3978)rectangle(-0.5096,2.4563);
\fill[DeepBlue!49.36](-0.6047,2.4543)rectangle(-0.5096,2.5129);
\fill[DeepBlue!42.08](-0.6047,2.5109)rectangle(-0.5096,2.5694);
\fill[DeepBlue!34.26](-0.6047,2.5674)rectangle(-0.5096,2.6259);
\fill[DeepBlue!26.31](-0.6047,2.6239)rectangle(-0.5096,2.6824);
\fill[DeepBlue!18.63](-0.6047,2.6804)rectangle(-0.5096,2.7390);
\fill[DeepBlue!11.59](-0.6047,2.7370)rectangle(-0.5096,2.7955);
\fill[DeepBlue!5.68](-0.6047,2.7935)rectangle(-0.5096,2.8520);
\fill[DeepBlue!5.04](-0.5116,1.5500)rectangle(-0.4166,1.6085);
\fill[DeepBlue!6.33](-0.5116,1.6065)rectangle(-0.4166,1.6650);
\fill[DeepBlue!10.84](-0.5116,1.6630)rectangle(-0.4166,1.7216);
\fill[DeepBlue!16.94](-0.5116,1.7196)rectangle(-0.4166,1.7781);
\fill[DeepBlue!23.86](-0.5116,1.7761)rectangle(-0.4166,1.8346);
\fill[DeepBlue!31.14](-0.5116,1.8326)rectangle(-0.4166,1.8911);
\fill[DeepBlue!38.37](-0.5116,1.8891)rectangle(-0.4166,1.9477);
\fill[DeepBlue!45.14](-0.5116,1.9457)rectangle(-0.4166,2.0042);
\fill[DeepBlue!51.09](-0.5116,2.0022)rectangle(-0.4166,2.0607);
\fill[DeepBlue!55.87](-0.5116,2.0587)rectangle(-0.4166,2.1172);
\fill[DeepBlue!59.22](-0.5116,2.1152)rectangle(-0.4166,2.1737);
\fill[DeepBlue!60.95](-0.5116,2.1717)rectangle(-0.4166,2.2303);
\fill[DeepBlue!60.95](-0.5116,2.2283)rectangle(-0.4166,2.2868);
\fill[DeepBlue!59.22](-0.5116,2.2848)rectangle(-0.4166,2.3433);
\fill[DeepBlue!55.87](-0.5116,2.3413)rectangle(-0.4166,2.3998);
\fill[DeepBlue!51.09](-0.5116,2.3978)rectangle(-0.4166,2.4563);
\fill[DeepBlue!45.14](-0.5116,2.4543)rectangle(-0.4166,2.5129);
\fill[DeepBlue!38.37](-0.5116,2.5109)rectangle(-0.4166,2.5694);
\fill[DeepBlue!31.14](-0.5116,2.5674)rectangle(-0.4166,2.6259);
\fill[DeepBlue!23.86](-0.5116,2.6239)rectangle(-0.4166,2.6824);
\fill[DeepBlue!16.94](-0.5116,2.6804)rectangle(-0.4166,2.7390);
\fill[DeepBlue!10.84](-0.5116,2.7370)rectangle(-0.4166,2.7955);
\fill[DeepBlue!6.33](-0.5116,2.7935)rectangle(-0.4166,2.8520);
\fill[DeepBlue!7.23](-0.4186,1.5500)rectangle(-0.3236,1.6085);
\fill[DeepBlue!7.44](-0.4186,1.6065)rectangle(-0.3236,1.6650);
\fill[DeepBlue!10.12](-0.4186,1.6630)rectangle(-0.3236,1.7216);
\fill[DeepBlue!14.80](-0.4186,1.7196)rectangle(-0.3236,1.7781);
\fill[DeepBlue!20.61](-0.4186,1.7761)rectangle(-0.3236,1.8346);
\fill[DeepBlue!26.97](-0.4186,1.8326)rectangle(-0.3236,1.8911);
\fill[DeepBlue!33.39](-0.4186,1.8891)rectangle(-0.3236,1.9477);
\fill[DeepBlue!39.49](-0.4186,1.9457)rectangle(-0.3236,2.0042);
\fill[DeepBlue!44.88](-0.4186,2.0022)rectangle(-0.3236,2.0607);
\fill[DeepBlue!49.24](-0.4186,2.0587)rectangle(-0.3236,2.1172);
\fill[DeepBlue!52.30](-0.4186,2.1152)rectangle(-0.3236,2.1737);
\fill[DeepBlue!53.87](-0.4186,2.1717)rectangle(-0.3236,2.2303);
\fill[DeepBlue!53.87](-0.4186,2.2283)rectangle(-0.3236,2.2868);
\fill[DeepBlue!52.30](-0.4186,2.2848)rectangle(-0.3236,2.3433);
\fill[DeepBlue!49.24](-0.4186,2.3413)rectangle(-0.3236,2.3998);
\fill[DeepBlue!44.88](-0.4186,2.3978)rectangle(-0.3236,2.4563);
\fill[DeepBlue!39.49](-0.4186,2.4543)rectangle(-0.3236,2.5129);
\fill[DeepBlue!33.39](-0.4186,2.5109)rectangle(-0.3236,2.5694);
\fill[DeepBlue!26.97](-0.4186,2.5674)rectangle(-0.3236,2.6259);
\fill[DeepBlue!20.61](-0.4186,2.6239)rectangle(-0.3236,2.6824);
\fill[DeepBlue!14.80](-0.4186,2.6804)rectangle(-0.3236,2.7390);
\fill[DeepBlue!10.12](-0.4186,2.7370)rectangle(-0.3236,2.7955);
\fill[DeepBlue!7.44](-0.4186,2.7935)rectangle(-0.3236,2.8520);
\fill[DeepBlue!9.25](-0.3256,1.5500)rectangle(-0.2306,1.6085);
\fill[DeepBlue!8.75](-0.3256,1.6065)rectangle(-0.2306,1.6650);
\fill[DeepBlue!9.59](-0.3256,1.6630)rectangle(-0.2306,1.7216);
\fill[DeepBlue!12.40](-0.3256,1.7196)rectangle(-0.2306,1.7781);
\fill[DeepBlue!16.73](-0.3256,1.7761)rectangle(-0.2306,1.8346);
\fill[DeepBlue!21.91](-0.3256,1.8326)rectangle(-0.2306,1.8911);
\fill[DeepBlue!27.38](-0.3256,1.8891)rectangle(-0.2306,1.9477);
\fill[DeepBlue!32.68](-0.3256,1.9457)rectangle(-0.2306,2.0042);
\fill[DeepBlue!37.43](-0.3256,2.0022)rectangle(-0.2306,2.0607);
\fill[DeepBlue!41.30](-0.3256,2.0587)rectangle(-0.2306,2.1172);
\fill[DeepBlue!44.03](-0.3256,2.1152)rectangle(-0.2306,2.1737);
\fill[DeepBlue!45.45](-0.3256,2.1717)rectangle(-0.2306,2.2303);
\fill[DeepBlue!45.45](-0.3256,2.2283)rectangle(-0.2306,2.2868);
\fill[DeepBlue!44.03](-0.3256,2.2848)rectangle(-0.2306,2.3433);
\fill[DeepBlue!41.30](-0.3256,2.3413)rectangle(-0.2306,2.3998);
\fill[DeepBlue!37.43](-0.3256,2.3978)rectangle(-0.2306,2.4563);
\fill[DeepBlue!32.68](-0.3256,2.4543)rectangle(-0.2306,2.5129);
\fill[DeepBlue!27.38](-0.3256,2.5109)rectangle(-0.2306,2.5694);
\fill[DeepBlue!21.91](-0.3256,2.5674)rectangle(-0.2306,2.6259);
\fill[DeepBlue!16.73](-0.3256,2.6239)rectangle(-0.2306,2.6824);
\fill[DeepBlue!12.40](-0.3256,2.6804)rectangle(-0.2306,2.7390);
\fill[DeepBlue!9.59](-0.3256,2.7370)rectangle(-0.2306,2.7955);
\fill[DeepBlue!8.75](-0.3256,2.7935)rectangle(-0.2306,2.8520);
\fill[DeepBlue!10.94](-0.2326,1.5500)rectangle(-0.1375,1.6085);
\fill[DeepBlue!10.00](-0.2326,1.6065)rectangle(-0.1375,1.6650);
\fill[DeepBlue!9.33](-0.2326,1.6630)rectangle(-0.1375,1.7216);
\fill[DeepBlue!9.96](-0.2326,1.7196)rectangle(-0.1375,1.7781);
\fill[DeepBlue!12.39](-0.2326,1.7761)rectangle(-0.1375,1.8346);
\fill[DeepBlue!16.18](-0.2326,1.8326)rectangle(-0.1375,1.8911);
\fill[DeepBlue!20.60](-0.2326,1.8891)rectangle(-0.1375,1.9477);
\fill[DeepBlue!25.09](-0.2326,1.9457)rectangle(-0.1375,2.0042);
\fill[DeepBlue!29.20](-0.2326,2.0022)rectangle(-0.1375,2.0607);
\fill[DeepBlue!32.60](-0.2326,2.0587)rectangle(-0.1375,2.1172);
\fill[DeepBlue!35.02](-0.2326,2.1152)rectangle(-0.1375,2.1737);
\fill[DeepBlue!36.28](-0.2326,2.1717)rectangle(-0.1375,2.2303);
\fill[DeepBlue!36.28](-0.2326,2.2283)rectangle(-0.1375,2.2868);
\fill[DeepBlue!35.02](-0.2326,2.2848)rectangle(-0.1375,2.3433);
\fill[DeepBlue!32.60](-0.2326,2.3413)rectangle(-0.1375,2.3998);
\fill[DeepBlue!29.20](-0.2326,2.3978)rectangle(-0.1375,2.4563);
\fill[DeepBlue!25.09](-0.2326,2.4543)rectangle(-0.1375,2.5129);
\fill[DeepBlue!20.60](-0.2326,2.5109)rectangle(-0.1375,2.5694);
\fill[DeepBlue!16.18](-0.2326,2.5674)rectangle(-0.1375,2.6259);
\fill[DeepBlue!12.39](-0.2326,2.6239)rectangle(-0.1375,2.6824);
\fill[DeepBlue!9.96](-0.2326,2.6804)rectangle(-0.1375,2.7390);
\fill[DeepBlue!9.33](-0.2326,2.7370)rectangle(-0.1375,2.7955);
\fill[DeepBlue!10.00](-0.2326,2.7935)rectangle(-0.1375,2.8520);
\fill[DeepBlue!12.15](-0.1395,1.5500)rectangle(-0.0445,1.6085);
\fill[DeepBlue!10.96](-0.1395,1.6065)rectangle(-0.0445,1.6650);
\fill[DeepBlue!9.27](-0.1395,1.6630)rectangle(-0.0445,1.7216);
\fill[DeepBlue!7.79](-0.1395,1.7196)rectangle(-0.0445,1.7781);
\fill[DeepBlue!7.81](-0.1395,1.7761)rectangle(-0.0445,1.8346);
\fill[DeepBlue!9.99](-0.1395,1.8326)rectangle(-0.0445,1.8911);
\fill[DeepBlue!13.49](-0.1395,1.8891)rectangle(-0.0445,1.9477);
\fill[DeepBlue!17.36](-0.1395,1.9457)rectangle(-0.0445,2.0042);
\fill[DeepBlue!21.03](-0.1395,2.0022)rectangle(-0.0445,2.0607);
\fill[DeepBlue!24.11](-0.1395,2.0587)rectangle(-0.0445,2.1172);
\fill[DeepBlue!26.31](-0.1395,2.1152)rectangle(-0.0445,2.1737);
\fill[DeepBlue!27.46](-0.1395,2.1717)rectangle(-0.0445,2.2303);
\fill[DeepBlue!27.46](-0.1395,2.2283)rectangle(-0.0445,2.2868);
\fill[DeepBlue!26.31](-0.1395,2.2848)rectangle(-0.0445,2.3433);
\fill[DeepBlue!24.11](-0.1395,2.3413)rectangle(-0.0445,2.3998);
\fill[DeepBlue!21.03](-0.1395,2.3978)rectangle(-0.0445,2.4563);
\fill[DeepBlue!17.36](-0.1395,2.4543)rectangle(-0.0445,2.5129);
\fill[DeepBlue!13.49](-0.1395,2.5109)rectangle(-0.0445,2.5694);
\fill[DeepBlue!9.99](-0.1395,2.5674)rectangle(-0.0445,2.6259);
\fill[DeepBlue!7.81](-0.1395,2.6239)rectangle(-0.0445,2.6824);
\fill[DeepBlue!7.79](-0.1395,2.6804)rectangle(-0.0445,2.7390);
\fill[DeepBlue!9.27](-0.1395,2.7370)rectangle(-0.0445,2.7955);
\fill[DeepBlue!10.96](-0.1395,2.7935)rectangle(-0.0445,2.8520);
\fill[DeepBlue!12.79](-0.0465,1.5500)rectangle(0.0485,1.6085);
\fill[DeepBlue!11.48](-0.0465,1.6065)rectangle(0.0485,1.6650);
\fill[DeepBlue!9.28](-0.0465,1.6630)rectangle(0.0485,1.7216);
\fill[DeepBlue!6.40](-0.0465,1.7196)rectangle(0.0485,1.7781);
\fill[DeepBlue!3.50](-0.0465,1.7761)rectangle(0.0485,1.8346);
\fill[DeepBlue!3.72](-0.0465,1.8326)rectangle(0.0485,1.8911);
\fill[DeepBlue!7.26](-0.0465,1.8891)rectangle(0.0485,1.9477);
\fill[DeepBlue!11.25](-0.0465,1.9457)rectangle(0.0485,2.0042);
\fill[DeepBlue!14.94](-0.0465,2.0022)rectangle(0.0485,2.0607);
\fill[DeepBlue!17.99](-0.0465,2.0587)rectangle(0.0485,2.1172);
\fill[DeepBlue!20.16](-0.0465,2.1152)rectangle(0.0485,2.1737);
\fill[DeepBlue!21.29](-0.0465,2.1717)rectangle(0.0485,2.2303);
\fill[DeepBlue!21.29](-0.0465,2.2283)rectangle(0.0485,2.2868);
\fill[DeepBlue!20.16](-0.0465,2.2848)rectangle(0.0485,2.3433);
\fill[DeepBlue!17.99](-0.0465,2.3413)rectangle(0.0485,2.3998);
\fill[DeepBlue!14.94](-0.0465,2.3978)rectangle(0.0485,2.4563);
\fill[DeepBlue!11.25](-0.0465,2.4543)rectangle(0.0485,2.5129);
\fill[DeepBlue!7.26](-0.0465,2.5109)rectangle(0.0485,2.5694);
\fill[DeepBlue!3.72](-0.0465,2.5674)rectangle(0.0485,2.6259);
\fill[DeepBlue!3.50](-0.0465,2.6239)rectangle(0.0485,2.6824);
\fill[DeepBlue!6.40](-0.0465,2.6804)rectangle(0.0485,2.7390);
\fill[DeepBlue!9.28](-0.0465,2.7370)rectangle(0.0485,2.7955);
\fill[DeepBlue!11.48](-0.0465,2.7935)rectangle(0.0485,2.8520);
\fill[DeepBlue!12.79](0.0465,1.5500)rectangle(0.1415,1.6085);
\fill[DeepBlue!11.48](0.0465,1.6065)rectangle(0.1415,1.6650);
\fill[DeepBlue!9.28](0.0465,1.6630)rectangle(0.1415,1.7216);
\fill[DeepBlue!6.40](0.0465,1.7196)rectangle(0.1415,1.7781);
\fill[DeepBlue!3.50](0.0465,1.7761)rectangle(0.1415,1.8346);
\fill[DeepBlue!3.72](0.0465,1.8326)rectangle(0.1415,1.8911);
\fill[DeepBlue!7.26](0.0465,1.8891)rectangle(0.1415,1.9477);
\fill[DeepBlue!11.25](0.0465,1.9457)rectangle(0.1415,2.0042);
\fill[DeepBlue!14.94](0.0465,2.0022)rectangle(0.1415,2.0607);
\fill[DeepBlue!17.99](0.0465,2.0587)rectangle(0.1415,2.1172);
\fill[DeepBlue!20.16](0.0465,2.1152)rectangle(0.1415,2.1737);
\fill[DeepBlue!21.29](0.0465,2.1717)rectangle(0.1415,2.2303);
\fill[DeepBlue!21.29](0.0465,2.2283)rectangle(0.1415,2.2868);
\fill[DeepBlue!20.16](0.0465,2.2848)rectangle(0.1415,2.3433);
\fill[DeepBlue!17.99](0.0465,2.3413)rectangle(0.1415,2.3998);
\fill[DeepBlue!14.94](0.0465,2.3978)rectangle(0.1415,2.4563);
\fill[DeepBlue!11.25](0.0465,2.4543)rectangle(0.1415,2.5129);
\fill[DeepBlue!7.26](0.0465,2.5109)rectangle(0.1415,2.5694);
\fill[DeepBlue!3.72](0.0465,2.5674)rectangle(0.1415,2.6259);
\fill[DeepBlue!3.50](0.0465,2.6239)rectangle(0.1415,2.6824);
\fill[DeepBlue!6.40](0.0465,2.6804)rectangle(0.1415,2.7390);
\fill[DeepBlue!9.28](0.0465,2.7370)rectangle(0.1415,2.7955);
\fill[DeepBlue!11.48](0.0465,2.7935)rectangle(0.1415,2.8520);
\fill[DeepBlue!12.15](0.1395,1.5500)rectangle(0.2346,1.6085);
\fill[DeepBlue!10.96](0.1395,1.6065)rectangle(0.2346,1.6650);
\fill[DeepBlue!9.27](0.1395,1.6630)rectangle(0.2346,1.7216);
\fill[DeepBlue!7.79](0.1395,1.7196)rectangle(0.2346,1.7781);
\fill[DeepBlue!7.81](0.1395,1.7761)rectangle(0.2346,1.8346);
\fill[DeepBlue!9.99](0.1395,1.8326)rectangle(0.2346,1.8911);
\fill[DeepBlue!13.49](0.1395,1.8891)rectangle(0.2346,1.9477);
\fill[DeepBlue!17.36](0.1395,1.9457)rectangle(0.2346,2.0042);
\fill[DeepBlue!21.03](0.1395,2.0022)rectangle(0.2346,2.0607);
\fill[DeepBlue!24.11](0.1395,2.0587)rectangle(0.2346,2.1172);
\fill[DeepBlue!26.31](0.1395,2.1152)rectangle(0.2346,2.1737);
\fill[DeepBlue!27.46](0.1395,2.1717)rectangle(0.2346,2.2303);
\fill[DeepBlue!27.46](0.1395,2.2283)rectangle(0.2346,2.2868);
\fill[DeepBlue!26.31](0.1395,2.2848)rectangle(0.2346,2.3433);
\fill[DeepBlue!24.11](0.1395,2.3413)rectangle(0.2346,2.3998);
\fill[DeepBlue!21.03](0.1395,2.3978)rectangle(0.2346,2.4563);
\fill[DeepBlue!17.36](0.1395,2.4543)rectangle(0.2346,2.5129);
\fill[DeepBlue!13.49](0.1395,2.5109)rectangle(0.2346,2.5694);
\fill[DeepBlue!9.99](0.1395,2.5674)rectangle(0.2346,2.6259);
\fill[DeepBlue!7.81](0.1395,2.6239)rectangle(0.2346,2.6824);
\fill[DeepBlue!7.79](0.1395,2.6804)rectangle(0.2346,2.7390);
\fill[DeepBlue!9.27](0.1395,2.7370)rectangle(0.2346,2.7955);
\fill[DeepBlue!10.96](0.1395,2.7935)rectangle(0.2346,2.8520);
\fill[DeepBlue!10.94](0.2326,1.5500)rectangle(0.3276,1.6085);
\fill[DeepBlue!10.00](0.2326,1.6065)rectangle(0.3276,1.6650);
\fill[DeepBlue!9.33](0.2326,1.6630)rectangle(0.3276,1.7216);
\fill[DeepBlue!9.96](0.2326,1.7196)rectangle(0.3276,1.7781);
\fill[DeepBlue!12.39](0.2326,1.7761)rectangle(0.3276,1.8346);
\fill[DeepBlue!16.18](0.2326,1.8326)rectangle(0.3276,1.8911);
\fill[DeepBlue!20.60](0.2326,1.8891)rectangle(0.3276,1.9477);
\fill[DeepBlue!25.09](0.2326,1.9457)rectangle(0.3276,2.0042);
\fill[DeepBlue!29.20](0.2326,2.0022)rectangle(0.3276,2.0607);
\fill[DeepBlue!32.60](0.2326,2.0587)rectangle(0.3276,2.1172);
\fill[DeepBlue!35.02](0.2326,2.1152)rectangle(0.3276,2.1737);
\fill[DeepBlue!36.28](0.2326,2.1717)rectangle(0.3276,2.2303);
\fill[DeepBlue!36.28](0.2326,2.2283)rectangle(0.3276,2.2868);
\fill[DeepBlue!35.02](0.2326,2.2848)rectangle(0.3276,2.3433);
\fill[DeepBlue!32.60](0.2326,2.3413)rectangle(0.3276,2.3998);
\fill[DeepBlue!29.20](0.2326,2.3978)rectangle(0.3276,2.4563);
\fill[DeepBlue!25.09](0.2326,2.4543)rectangle(0.3276,2.5129);
\fill[DeepBlue!20.60](0.2326,2.5109)rectangle(0.3276,2.5694);
\fill[DeepBlue!16.18](0.2326,2.5674)rectangle(0.3276,2.6259);
\fill[DeepBlue!12.39](0.2326,2.6239)rectangle(0.3276,2.6824);
\fill[DeepBlue!9.96](0.2326,2.6804)rectangle(0.3276,2.7390);
\fill[DeepBlue!9.33](0.2326,2.7370)rectangle(0.3276,2.7955);
\fill[DeepBlue!10.00](0.2326,2.7935)rectangle(0.3276,2.8520);
\fill[DeepBlue!9.25](0.3256,1.5500)rectangle(0.4206,1.6085);
\fill[DeepBlue!8.75](0.3256,1.6065)rectangle(0.4206,1.6650);
\fill[DeepBlue!9.59](0.3256,1.6630)rectangle(0.4206,1.7216);
\fill[DeepBlue!12.40](0.3256,1.7196)rectangle(0.4206,1.7781);
\fill[DeepBlue!16.73](0.3256,1.7761)rectangle(0.4206,1.8346);
\fill[DeepBlue!21.91](0.3256,1.8326)rectangle(0.4206,1.8911);
\fill[DeepBlue!27.38](0.3256,1.8891)rectangle(0.4206,1.9477);
\fill[DeepBlue!32.68](0.3256,1.9457)rectangle(0.4206,2.0042);
\fill[DeepBlue!37.43](0.3256,2.0022)rectangle(0.4206,2.0607);
\fill[DeepBlue!41.30](0.3256,2.0587)rectangle(0.4206,2.1172);
\fill[DeepBlue!44.03](0.3256,2.1152)rectangle(0.4206,2.1737);
\fill[DeepBlue!45.45](0.3256,2.1717)rectangle(0.4206,2.2303);
\fill[DeepBlue!45.45](0.3256,2.2283)rectangle(0.4206,2.2868);
\fill[DeepBlue!44.03](0.3256,2.2848)rectangle(0.4206,2.3433);
\fill[DeepBlue!41.30](0.3256,2.3413)rectangle(0.4206,2.3998);
\fill[DeepBlue!37.43](0.3256,2.3978)rectangle(0.4206,2.4563);
\fill[DeepBlue!32.68](0.3256,2.4543)rectangle(0.4206,2.5129);
\fill[DeepBlue!27.38](0.3256,2.5109)rectangle(0.4206,2.5694);
\fill[DeepBlue!21.91](0.3256,2.5674)rectangle(0.4206,2.6259);
\fill[DeepBlue!16.73](0.3256,2.6239)rectangle(0.4206,2.6824);
\fill[DeepBlue!12.40](0.3256,2.6804)rectangle(0.4206,2.7390);
\fill[DeepBlue!9.59](0.3256,2.7370)rectangle(0.4206,2.7955);
\fill[DeepBlue!8.75](0.3256,2.7935)rectangle(0.4206,2.8520);
\fill[DeepBlue!7.23](0.4186,1.5500)rectangle(0.5136,1.6085);
\fill[DeepBlue!7.44](0.4186,1.6065)rectangle(0.5136,1.6650);
\fill[DeepBlue!10.12](0.4186,1.6630)rectangle(0.5136,1.7216);
\fill[DeepBlue!14.80](0.4186,1.7196)rectangle(0.5136,1.7781);
\fill[DeepBlue!20.61](0.4186,1.7761)rectangle(0.5136,1.8346);
\fill[DeepBlue!26.97](0.4186,1.8326)rectangle(0.5136,1.8911);
\fill[DeepBlue!33.39](0.4186,1.8891)rectangle(0.5136,1.9477);
\fill[DeepBlue!39.49](0.4186,1.9457)rectangle(0.5136,2.0042);
\fill[DeepBlue!44.88](0.4186,2.0022)rectangle(0.5136,2.0607);
\fill[DeepBlue!49.24](0.4186,2.0587)rectangle(0.5136,2.1172);
\fill[DeepBlue!52.30](0.4186,2.1152)rectangle(0.5136,2.1737);
\fill[DeepBlue!53.87](0.4186,2.1717)rectangle(0.5136,2.2303);
\fill[DeepBlue!53.87](0.4186,2.2283)rectangle(0.5136,2.2868);
\fill[DeepBlue!52.30](0.4186,2.2848)rectangle(0.5136,2.3433);
\fill[DeepBlue!49.24](0.4186,2.3413)rectangle(0.5136,2.3998);
\fill[DeepBlue!44.88](0.4186,2.3978)rectangle(0.5136,2.4563);
\fill[DeepBlue!39.49](0.4186,2.4543)rectangle(0.5136,2.5129);
\fill[DeepBlue!33.39](0.4186,2.5109)rectangle(0.5136,2.5694);
\fill[DeepBlue!26.97](0.4186,2.5674)rectangle(0.5136,2.6259);
\fill[DeepBlue!20.61](0.4186,2.6239)rectangle(0.5136,2.6824);
\fill[DeepBlue!14.80](0.4186,2.6804)rectangle(0.5136,2.7390);
\fill[DeepBlue!10.12](0.4186,2.7370)rectangle(0.5136,2.7955);
\fill[DeepBlue!7.44](0.4186,2.7935)rectangle(0.5136,2.8520);
\fill[DeepBlue!5.04](0.5116,1.5500)rectangle(0.6067,1.6085);
\fill[DeepBlue!6.33](0.5116,1.6065)rectangle(0.6067,1.6650);
\fill[DeepBlue!10.84](0.5116,1.6630)rectangle(0.6067,1.7216);
\fill[DeepBlue!16.94](0.5116,1.7196)rectangle(0.6067,1.7781);
\fill[DeepBlue!23.86](0.5116,1.7761)rectangle(0.6067,1.8346);
\fill[DeepBlue!31.14](0.5116,1.8326)rectangle(0.6067,1.8911);
\fill[DeepBlue!38.37](0.5116,1.8891)rectangle(0.6067,1.9477);
\fill[DeepBlue!45.14](0.5116,1.9457)rectangle(0.6067,2.0042);
\fill[DeepBlue!51.09](0.5116,2.0022)rectangle(0.6067,2.0607);
\fill[DeepBlue!55.87](0.5116,2.0587)rectangle(0.6067,2.1172);
\fill[DeepBlue!59.22](0.5116,2.1152)rectangle(0.6067,2.1737);
\fill[DeepBlue!60.95](0.5116,2.1717)rectangle(0.6067,2.2303);
\fill[DeepBlue!60.95](0.5116,2.2283)rectangle(0.6067,2.2868);
\fill[DeepBlue!59.22](0.5116,2.2848)rectangle(0.6067,2.3433);
\fill[DeepBlue!55.87](0.5116,2.3413)rectangle(0.6067,2.3998);
\fill[DeepBlue!51.09](0.5116,2.3978)rectangle(0.6067,2.4563);
\fill[DeepBlue!45.14](0.5116,2.4543)rectangle(0.6067,2.5129);
\fill[DeepBlue!38.37](0.5116,2.5109)rectangle(0.6067,2.5694);
\fill[DeepBlue!31.14](0.5116,2.5674)rectangle(0.6067,2.6259);
\fill[DeepBlue!23.86](0.5116,2.6239)rectangle(0.6067,2.6824);
\fill[DeepBlue!16.94](0.5116,2.6804)rectangle(0.6067,2.7390);
\fill[DeepBlue!10.84](0.5116,2.7370)rectangle(0.6067,2.7955);
\fill[DeepBlue!6.33](0.5116,2.7935)rectangle(0.6067,2.8520);
\fill[DeepBlue!2.83](0.6047,1.5500)rectangle(0.6997,1.6085);
\fill[DeepBlue!5.68](0.6047,1.6065)rectangle(0.6997,1.6650);
\fill[DeepBlue!11.59](0.6047,1.6630)rectangle(0.6997,1.7216);
\fill[DeepBlue!18.63](0.6047,1.7196)rectangle(0.6997,1.7781);
\fill[DeepBlue!26.31](0.6047,1.7761)rectangle(0.6997,1.8346);
\fill[DeepBlue!34.26](0.6047,1.8326)rectangle(0.6997,1.8911);
\fill[DeepBlue!42.08](0.6047,1.8891)rectangle(0.6997,1.9477);
\fill[DeepBlue!49.36](0.6047,1.9457)rectangle(0.6997,2.0042);
\fill[DeepBlue!55.74](0.6047,2.0022)rectangle(0.6997,2.0607);
\fill[DeepBlue!60.85](0.6047,2.0587)rectangle(0.6997,2.1172);
\fill[DeepBlue!64.42](0.6047,2.1152)rectangle(0.6997,2.1737);
\fill[DeepBlue!66.26](0.6047,2.1717)rectangle(0.6997,2.2303);
\fill[DeepBlue!66.26](0.6047,2.2283)rectangle(0.6997,2.2868);
\fill[DeepBlue!64.42](0.6047,2.2848)rectangle(0.6997,2.3433);
\fill[DeepBlue!60.85](0.6047,2.3413)rectangle(0.6997,2.3998);
\fill[DeepBlue!55.74](0.6047,2.3978)rectangle(0.6997,2.4563);
\fill[DeepBlue!49.36](0.6047,2.4543)rectangle(0.6997,2.5129);
\fill[DeepBlue!42.08](0.6047,2.5109)rectangle(0.6997,2.5694);
\fill[DeepBlue!34.26](0.6047,2.5674)rectangle(0.6997,2.6259);
\fill[DeepBlue!26.31](0.6047,2.6239)rectangle(0.6997,2.6824);
\fill[DeepBlue!18.63](0.6047,2.6804)rectangle(0.6997,2.7390);
\fill[DeepBlue!11.59](0.6047,2.7370)rectangle(0.6997,2.7955);
\fill[DeepBlue!5.68](0.6047,2.7935)rectangle(0.6997,2.8520);
\fill[DeepBlue!0.77](0.6977,1.5500)rectangle(0.7927,1.6085);
\fill[DeepBlue!5.50](0.6977,1.6065)rectangle(0.7927,1.6650);
\fill[DeepBlue!12.15](0.6977,1.6630)rectangle(0.7927,1.7216);
\fill[DeepBlue!19.69](0.6977,1.7196)rectangle(0.7927,1.7781);
\fill[DeepBlue!27.82](0.6977,1.7761)rectangle(0.7927,1.8346);
\fill[DeepBlue!36.17](0.6977,1.8326)rectangle(0.7927,1.8911);
\fill[DeepBlue!44.35](0.6977,1.8891)rectangle(0.7927,1.9477);
\fill[DeepBlue!51.96](0.6977,1.9457)rectangle(0.7927,2.0042);
\fill[DeepBlue!58.61](0.6977,2.0022)rectangle(0.7927,2.0607);
\fill[DeepBlue!63.93](0.6977,2.0587)rectangle(0.7927,2.1172);
\fill[DeepBlue!67.65](0.6977,2.1152)rectangle(0.7927,2.1737);
\fill[DeepBlue!69.56](0.6977,2.1717)rectangle(0.7927,2.2303);
\fill[DeepBlue!69.56](0.6977,2.2283)rectangle(0.7927,2.2868);
\fill[DeepBlue!67.65](0.6977,2.2848)rectangle(0.7927,2.3433);
\fill[DeepBlue!63.93](0.6977,2.3413)rectangle(0.7927,2.3998);
\fill[DeepBlue!58.61](0.6977,2.3978)rectangle(0.7927,2.4563);
\fill[DeepBlue!51.96](0.6977,2.4543)rectangle(0.7927,2.5129);
\fill[DeepBlue!44.35](0.6977,2.5109)rectangle(0.7927,2.5694);
\fill[DeepBlue!36.17](0.6977,2.5674)rectangle(0.7927,2.6259);
\fill[DeepBlue!27.82](0.6977,2.6239)rectangle(0.7927,2.6824);
\fill[DeepBlue!19.69](0.6977,2.6804)rectangle(0.7927,2.7390);
\fill[DeepBlue!12.15](0.6977,2.7370)rectangle(0.7927,2.7955);
\fill[DeepBlue!5.50](0.6977,2.7935)rectangle(0.7927,2.8520);
\fill[DeepBlue!1.06](0.7907,1.5500)rectangle(0.8857,1.6085);
\fill[DeepBlue!5.58](0.7907,1.6065)rectangle(0.8857,1.6650);
\fill[DeepBlue!12.33](0.7907,1.6630)rectangle(0.8857,1.7216);
\fill[DeepBlue!20.00](0.7907,1.7196)rectangle(0.8857,1.7781);
\fill[DeepBlue!28.28](0.7907,1.7761)rectangle(0.8857,1.8346);
\fill[DeepBlue!36.77](0.7907,1.8326)rectangle(0.8857,1.8911);
\fill[DeepBlue!45.09](0.7907,1.8891)rectangle(0.8857,1.9477);
\fill[DeepBlue!52.83](0.7907,1.9457)rectangle(0.8857,2.0042);
\fill[DeepBlue!59.58](0.7907,2.0022)rectangle(0.8857,2.0607);
\fill[DeepBlue!64.98](0.7907,2.0587)rectangle(0.8857,2.1172);
\fill[DeepBlue!68.76](0.7907,2.1152)rectangle(0.8857,2.1737);
\fill[DeepBlue!70.71](0.7907,2.1717)rectangle(0.8857,2.2303);
\fill[DeepBlue!70.71](0.7907,2.2283)rectangle(0.8857,2.2868);
\fill[DeepBlue!68.76](0.7907,2.2848)rectangle(0.8857,2.3433);
\fill[DeepBlue!64.98](0.7907,2.3413)rectangle(0.8857,2.3998);
\fill[DeepBlue!59.58](0.7907,2.3978)rectangle(0.8857,2.4563);
\fill[DeepBlue!52.83](0.7907,2.4543)rectangle(0.8857,2.5129);
\fill[DeepBlue!45.09](0.7907,2.5109)rectangle(0.8857,2.5694);
\fill[DeepBlue!36.77](0.7907,2.5674)rectangle(0.8857,2.6259);
\fill[DeepBlue!28.28](0.7907,2.6239)rectangle(0.8857,2.6824);
\fill[DeepBlue!20.00](0.7907,2.6804)rectangle(0.8857,2.7390);
\fill[DeepBlue!12.33](0.7907,2.7370)rectangle(0.8857,2.7955);
\fill[DeepBlue!5.58](0.7907,2.7935)rectangle(0.8857,2.8520);
\fill[DeepBlue!2.57](0.8837,1.5500)rectangle(0.9787,1.6085);
\fill[DeepBlue!5.60](0.8837,1.6065)rectangle(0.9787,1.6650);
\fill[DeepBlue!11.99](0.8837,1.6630)rectangle(0.9787,1.7216);
\fill[DeepBlue!19.49](0.8837,1.7196)rectangle(0.9787,1.7781);
\fill[DeepBlue!27.64](0.8837,1.7761)rectangle(0.9787,1.8346);
\fill[DeepBlue!36.03](0.8837,1.8326)rectangle(0.9787,1.8911);
\fill[DeepBlue!44.27](0.8837,1.8891)rectangle(0.9787,1.9477);
\fill[DeepBlue!51.93](0.8837,1.9457)rectangle(0.9787,2.0042);
\fill[DeepBlue!58.62](0.8837,2.0022)rectangle(0.9787,2.0607);
\fill[DeepBlue!63.98](0.8837,2.0587)rectangle(0.9787,2.1172);
\fill[DeepBlue!67.73](0.8837,2.1152)rectangle(0.9787,2.1737);
\fill[DeepBlue!69.66](0.8837,2.1717)rectangle(0.9787,2.2303);
\fill[DeepBlue!69.66](0.8837,2.2283)rectangle(0.9787,2.2868);
\fill[DeepBlue!67.73](0.8837,2.2848)rectangle(0.9787,2.3433);
\fill[DeepBlue!63.98](0.8837,2.3413)rectangle(0.9787,2.3998);
\fill[DeepBlue!58.62](0.8837,2.3978)rectangle(0.9787,2.4563);
\fill[DeepBlue!51.93](0.8837,2.4543)rectangle(0.9787,2.5129);
\fill[DeepBlue!44.27](0.8837,2.5109)rectangle(0.9787,2.5694);
\fill[DeepBlue!36.03](0.8837,2.5674)rectangle(0.9787,2.6259);
\fill[DeepBlue!27.64](0.8837,2.6239)rectangle(0.9787,2.6824);
\fill[DeepBlue!19.49](0.8837,2.6804)rectangle(0.9787,2.7390);
\fill[DeepBlue!11.99](0.8837,2.7370)rectangle(0.9787,2.7955);
\fill[DeepBlue!5.60](0.8837,2.7935)rectangle(0.9787,2.8520);
\fill[DeepBlue!3.79](0.9767,1.5500)rectangle(1.0718,1.6085);
\fill[DeepBlue!5.42](0.9767,1.6065)rectangle(1.0718,1.6650);
\fill[DeepBlue!11.07](0.9767,1.6630)rectangle(1.0718,1.7216);
\fill[DeepBlue!18.13](0.9767,1.7196)rectangle(1.0718,1.7781);
\fill[DeepBlue!25.92](0.9767,1.7761)rectangle(1.0718,1.8346);
\fill[DeepBlue!33.99](0.9767,1.8326)rectangle(1.0718,1.8911);
\fill[DeepBlue!41.93](0.9767,1.8891)rectangle(1.0718,1.9477);
\fill[DeepBlue!49.34](0.9767,1.9457)rectangle(1.0718,2.0042);
\fill[DeepBlue!55.81](0.9767,2.0022)rectangle(1.0718,2.0607);
\fill[DeepBlue!61.01](0.9767,2.0587)rectangle(1.0718,2.1172);
\fill[DeepBlue!64.64](0.9767,2.1152)rectangle(1.0718,2.1737);
\fill[DeepBlue!66.51](0.9767,2.1717)rectangle(1.0718,2.2303);
\fill[DeepBlue!66.51](0.9767,2.2283)rectangle(1.0718,2.2868);
\fill[DeepBlue!64.64](0.9767,2.2848)rectangle(1.0718,2.3433);
\fill[DeepBlue!61.01](0.9767,2.3413)rectangle(1.0718,2.3998);
\fill[DeepBlue!55.81](0.9767,2.3978)rectangle(1.0718,2.4563);
\fill[DeepBlue!49.34](0.9767,2.4543)rectangle(1.0718,2.5129);
\fill[DeepBlue!41.93](0.9767,2.5109)rectangle(1.0718,2.5694);
\fill[DeepBlue!33.99](0.9767,2.5674)rectangle(1.0718,2.6259);
\fill[DeepBlue!25.92](0.9767,2.6239)rectangle(1.0718,2.6824);
\fill[DeepBlue!18.13](0.9767,2.6804)rectangle(1.0718,2.7390);
\fill[DeepBlue!11.07](0.9767,2.7370)rectangle(1.0718,2.7955);
\fill[DeepBlue!5.42](0.9767,2.7935)rectangle(1.0718,2.8520);
\fill[DeepBlue!4.77](1.0698,1.5500)rectangle(1.1648,1.6085);
\fill[DeepBlue!5.03](1.0698,1.6065)rectangle(1.1648,1.6650);
\fill[DeepBlue!9.61](1.0698,1.6630)rectangle(1.1648,1.7216);
\fill[DeepBlue!16.00](1.0698,1.7196)rectangle(1.1648,1.7781);
\fill[DeepBlue!23.21](1.0698,1.7761)rectangle(1.1648,1.8346);
\fill[DeepBlue!30.76](1.0698,1.8326)rectangle(1.1648,1.8911);
\fill[DeepBlue!38.23](1.0698,1.8891)rectangle(1.1648,1.9477);
\fill[DeepBlue!45.21](1.0698,1.9457)rectangle(1.1648,2.0042);
\fill[DeepBlue!51.33](1.0698,2.0022)rectangle(1.1648,2.0607);
\fill[DeepBlue!56.25](1.0698,2.0587)rectangle(1.1648,2.1172);
\fill[DeepBlue!59.69](1.0698,2.1152)rectangle(1.1648,2.1737);
\fill[DeepBlue!61.46](1.0698,2.1717)rectangle(1.1648,2.2303);
\fill[DeepBlue!61.46](1.0698,2.2283)rectangle(1.1648,2.2868);
\fill[DeepBlue!59.69](1.0698,2.2848)rectangle(1.1648,2.3433);
\fill[DeepBlue!56.25](1.0698,2.3413)rectangle(1.1648,2.3998);
\fill[DeepBlue!51.33](1.0698,2.3978)rectangle(1.1648,2.4563);
\fill[DeepBlue!45.21](1.0698,2.4543)rectangle(1.1648,2.5129);
\fill[DeepBlue!38.23](1.0698,2.5109)rectangle(1.1648,2.5694);
\fill[DeepBlue!30.76](1.0698,2.5674)rectangle(1.1648,2.6259);
\fill[DeepBlue!23.21](1.0698,2.6239)rectangle(1.1648,2.6824);
\fill[DeepBlue!16.00](1.0698,2.6804)rectangle(1.1648,2.7390);
\fill[DeepBlue!9.61](1.0698,2.7370)rectangle(1.1648,2.7955);
\fill[DeepBlue!5.03](1.0698,2.7935)rectangle(1.1648,2.8520);
\fill[DeepBlue!5.62](1.1628,1.5500)rectangle(1.2578,1.6085);
\fill[DeepBlue!4.64](1.1628,1.6065)rectangle(1.2578,1.6650);
\fill[DeepBlue!7.73](1.1628,1.6630)rectangle(1.2578,1.7216);
\fill[DeepBlue!13.21](1.1628,1.7196)rectangle(1.2578,1.7781);
\fill[DeepBlue!19.67](1.1628,1.7761)rectangle(1.2578,1.8346);
\fill[DeepBlue!26.54](1.1628,1.8326)rectangle(1.2578,1.8911);
\fill[DeepBlue!33.37](1.1628,1.8891)rectangle(1.2578,1.9477);
\fill[DeepBlue!39.79](1.1628,1.9457)rectangle(1.2578,2.0042);
\fill[DeepBlue!45.44](1.1628,2.0022)rectangle(1.2578,2.0607);
\fill[DeepBlue!49.98](1.1628,2.0587)rectangle(1.2578,2.1172);
\fill[DeepBlue!53.17](1.1628,2.1152)rectangle(1.2578,2.1737);
\fill[DeepBlue!54.81](1.1628,2.1717)rectangle(1.2578,2.2303);
\fill[DeepBlue!54.81](1.1628,2.2283)rectangle(1.2578,2.2868);
\fill[DeepBlue!53.17](1.1628,2.2848)rectangle(1.2578,2.3433);
\fill[DeepBlue!49.98](1.1628,2.3413)rectangle(1.2578,2.3998);
\fill[DeepBlue!45.44](1.1628,2.3978)rectangle(1.2578,2.4563);
\fill[DeepBlue!39.79](1.1628,2.4543)rectangle(1.2578,2.5129);
\fill[DeepBlue!33.37](1.1628,2.5109)rectangle(1.2578,2.5694);
\fill[DeepBlue!26.54](1.1628,2.5674)rectangle(1.2578,2.6259);
\fill[DeepBlue!19.67](1.1628,2.6239)rectangle(1.2578,2.6824);
\fill[DeepBlue!13.21](1.1628,2.6804)rectangle(1.2578,2.7390);
\fill[DeepBlue!7.73](1.1628,2.7370)rectangle(1.2578,2.7955);
\fill[DeepBlue!4.64](1.1628,2.7935)rectangle(1.2578,2.8520);
\fill[DeepBlue!6.43](1.2558,1.5500)rectangle(1.3508,1.6085);
\fill[DeepBlue!4.61](1.2558,1.6065)rectangle(1.3508,1.6650);
\fill[DeepBlue!5.73](1.2558,1.6630)rectangle(1.3508,1.7216);
\fill[DeepBlue!10.00](1.2558,1.7196)rectangle(1.3508,1.7781);
\fill[DeepBlue!15.54](1.2558,1.7761)rectangle(1.3508,1.8346);
\fill[DeepBlue!21.58](1.2558,1.8326)rectangle(1.3508,1.8911);
\fill[DeepBlue!27.66](1.2558,1.8891)rectangle(1.3508,1.9477);
\fill[DeepBlue!33.41](1.2558,1.9457)rectangle(1.3508,2.0042);
\fill[DeepBlue!38.48](1.2558,2.0022)rectangle(1.3508,2.0607);
\fill[DeepBlue!42.58](1.2558,2.0587)rectangle(1.3508,2.1172);
\fill[DeepBlue!45.46](1.2558,2.1152)rectangle(1.3508,2.1737);
\fill[DeepBlue!46.94](1.2558,2.1717)rectangle(1.3508,2.2303);
\fill[DeepBlue!46.94](1.2558,2.2283)rectangle(1.3508,2.2868);
\fill[DeepBlue!45.46](1.2558,2.2848)rectangle(1.3508,2.3433);
\fill[DeepBlue!42.58](1.2558,2.3413)rectangle(1.3508,2.3998);
\fill[DeepBlue!38.48](1.2558,2.3978)rectangle(1.3508,2.4563);
\fill[DeepBlue!33.41](1.2558,2.4543)rectangle(1.3508,2.5129);
\fill[DeepBlue!27.66](1.2558,2.5109)rectangle(1.3508,2.5694);
\fill[DeepBlue!21.58](1.2558,2.5674)rectangle(1.3508,2.6259);
\fill[DeepBlue!15.54](1.2558,2.6239)rectangle(1.3508,2.6824);
\fill[DeepBlue!10.00](1.2558,2.6804)rectangle(1.3508,2.7390);
\fill[DeepBlue!5.73](1.2558,2.7370)rectangle(1.3508,2.7955);
\fill[DeepBlue!4.61](1.2558,2.7935)rectangle(1.3508,2.8520);
\fill[DeepBlue!7.27](1.3488,1.5500)rectangle(1.4439,1.6085);
\fill[DeepBlue!5.20](1.3488,1.6065)rectangle(1.4439,1.6650);
\fill[DeepBlue!4.24](1.3488,1.6630)rectangle(1.4439,1.7216);
\fill[DeepBlue!6.68](1.3488,1.7196)rectangle(1.4439,1.7781);
\fill[DeepBlue!11.10](1.3488,1.7761)rectangle(1.4439,1.8346);
\fill[DeepBlue!16.20](1.3488,1.8326)rectangle(1.4439,1.8911);
\fill[DeepBlue!21.44](1.3488,1.8891)rectangle(1.4439,1.9477);
\fill[DeepBlue!26.44](1.3488,1.9457)rectangle(1.4439,2.0042);
\fill[DeepBlue!30.87](1.3488,2.0022)rectangle(1.4439,2.0607);
\fill[DeepBlue!34.46](1.3488,2.0587)rectangle(1.4439,2.1172);
\fill[DeepBlue!36.99](1.3488,2.1152)rectangle(1.4439,2.1737);
\fill[DeepBlue!38.29](1.3488,2.1717)rectangle(1.4439,2.2303);
\fill[DeepBlue!38.29](1.3488,2.2283)rectangle(1.4439,2.2868);
\fill[DeepBlue!36.99](1.3488,2.2848)rectangle(1.4439,2.3433);
\fill[DeepBlue!34.46](1.3488,2.3413)rectangle(1.4439,2.3998);
\fill[DeepBlue!30.87](1.3488,2.3978)rectangle(1.4439,2.4563);
\fill[DeepBlue!26.44](1.3488,2.4543)rectangle(1.4439,2.5129);
\fill[DeepBlue!21.44](1.3488,2.5109)rectangle(1.4439,2.5694);
\fill[DeepBlue!16.20](1.3488,2.5674)rectangle(1.4439,2.6259);
\fill[DeepBlue!11.10](1.3488,2.6239)rectangle(1.4439,2.6824);
\fill[DeepBlue!6.68](1.3488,2.6804)rectangle(1.4439,2.7390);
\fill[DeepBlue!4.24](1.3488,2.7370)rectangle(1.4439,2.7955);
\fill[DeepBlue!5.20](1.3488,2.7935)rectangle(1.4439,2.8520);
\fill[DeepBlue!8.08](1.4419,1.5500)rectangle(1.5369,1.6085);
\fill[DeepBlue!6.27](1.4419,1.6065)rectangle(1.5369,1.6650);
\fill[DeepBlue!4.24](1.4419,1.6630)rectangle(1.5369,1.7216);
\fill[DeepBlue!3.95](1.4419,1.7196)rectangle(1.5369,1.7781);
\fill[DeepBlue!6.73](1.4419,1.7761)rectangle(1.5369,1.8346);
\fill[DeepBlue!10.76](1.4419,1.8326)rectangle(1.5369,1.8911);
\fill[DeepBlue!15.08](1.4419,1.8891)rectangle(1.5369,1.9477);
\fill[DeepBlue!19.27](1.4419,1.9457)rectangle(1.5369,2.0042);
\fill[DeepBlue!23.02](1.4419,2.0022)rectangle(1.5369,2.0607);
\fill[DeepBlue!26.08](1.4419,2.0587)rectangle(1.5369,2.1172);
\fill[DeepBlue!28.23](1.4419,2.1152)rectangle(1.5369,2.1737);
\fill[DeepBlue!29.35](1.4419,2.1717)rectangle(1.5369,2.2303);
\fill[DeepBlue!29.35](1.4419,2.2283)rectangle(1.5369,2.2868);
\fill[DeepBlue!28.23](1.4419,2.2848)rectangle(1.5369,2.3433);
\fill[DeepBlue!26.08](1.4419,2.3413)rectangle(1.5369,2.3998);
\fill[DeepBlue!23.02](1.4419,2.3978)rectangle(1.5369,2.4563);
\fill[DeepBlue!19.27](1.4419,2.4543)rectangle(1.5369,2.5129);
\fill[DeepBlue!15.08](1.4419,2.5109)rectangle(1.5369,2.5694);
\fill[DeepBlue!10.76](1.4419,2.5674)rectangle(1.5369,2.6259);
\fill[DeepBlue!6.73](1.4419,2.6239)rectangle(1.5369,2.6824);
\fill[DeepBlue!3.95](1.4419,2.6804)rectangle(1.5369,2.7390);
\fill[DeepBlue!4.24](1.4419,2.7370)rectangle(1.5369,2.7955);
\fill[DeepBlue!6.27](1.4419,2.7935)rectangle(1.5369,2.8520);
\fill[DeepBlue!8.77](1.5349,1.5500)rectangle(1.6299,1.6085);
\fill[DeepBlue!7.48](1.5349,1.6065)rectangle(1.6299,1.6650);
\fill[DeepBlue!5.60](1.5349,1.6630)rectangle(1.6299,1.7216);
\fill[DeepBlue!3.58](1.5349,1.7196)rectangle(1.6299,1.7781);
\fill[DeepBlue!3.25](1.5349,1.7761)rectangle(1.6299,1.8346);
\fill[DeepBlue!5.67](1.5349,1.8326)rectangle(1.6299,1.8911);
\fill[DeepBlue!8.97](1.5349,1.8891)rectangle(1.6299,1.9477);
\fill[DeepBlue!12.32](1.5349,1.9457)rectangle(1.6299,2.0042);
\fill[DeepBlue!15.37](1.5349,2.0022)rectangle(1.6299,2.0607);
\fill[DeepBlue!17.88](1.5349,2.0587)rectangle(1.6299,2.1172);
\fill[DeepBlue!19.66](1.5349,2.1152)rectangle(1.6299,2.1737);
\fill[DeepBlue!20.58](1.5349,2.1717)rectangle(1.6299,2.2303);
\fill[DeepBlue!20.58](1.5349,2.2283)rectangle(1.6299,2.2868);
\fill[DeepBlue!19.66](1.5349,2.2848)rectangle(1.6299,2.3433);
\fill[DeepBlue!17.88](1.5349,2.3413)rectangle(1.6299,2.3998);
\fill[DeepBlue!15.37](1.5349,2.3978)rectangle(1.6299,2.4563);
\fill[DeepBlue!12.32](1.5349,2.4543)rectangle(1.6299,2.5129);
\fill[DeepBlue!8.97](1.5349,2.5109)rectangle(1.6299,2.5694);
\fill[DeepBlue!5.67](1.5349,2.5674)rectangle(1.6299,2.6259);
\fill[DeepBlue!3.25](1.5349,2.6239)rectangle(1.6299,2.6824);
\fill[DeepBlue!3.58](1.5349,2.6804)rectangle(1.6299,2.7390);
\fill[DeepBlue!5.60](1.5349,2.7370)rectangle(1.6299,2.7955);
\fill[DeepBlue!7.48](1.5349,2.7935)rectangle(1.6299,2.8520);
\fill[DeepBlue!9.24](1.6279,1.5500)rectangle(1.7229,1.6085);
\fill[DeepBlue!8.54](1.6279,1.6065)rectangle(1.7229,1.6650);
\fill[DeepBlue!7.24](1.6279,1.6630)rectangle(1.7229,1.7216);
\fill[DeepBlue!5.45](1.6279,1.7196)rectangle(1.7229,1.7781);
\fill[DeepBlue!3.42](1.6279,1.7761)rectangle(1.7229,1.8346);
\fill[DeepBlue!2.23](1.6279,1.8326)rectangle(1.7229,1.8911);
\fill[DeepBlue!3.63](1.6279,1.8891)rectangle(1.7229,1.9477);
\fill[DeepBlue!6.01](1.6279,1.9457)rectangle(1.7229,2.0042);
\fill[DeepBlue!8.34](1.6279,2.0022)rectangle(1.7229,2.0607);
\fill[DeepBlue!10.30](1.6279,2.0587)rectangle(1.7229,2.1172);
\fill[DeepBlue!11.70](1.6279,2.1152)rectangle(1.7229,2.1737);
\fill[DeepBlue!12.43](1.6279,2.1717)rectangle(1.7229,2.2303);
\fill[DeepBlue!12.43](1.6279,2.2283)rectangle(1.7229,2.2868);
\fill[DeepBlue!11.70](1.6279,2.2848)rectangle(1.7229,2.3433);
\fill[DeepBlue!10.30](1.6279,2.3413)rectangle(1.7229,2.3998);
\fill[DeepBlue!8.34](1.6279,2.3978)rectangle(1.7229,2.4563);
\fill[DeepBlue!6.01](1.6279,2.4543)rectangle(1.7229,2.5129);
\fill[DeepBlue!3.63](1.6279,2.5109)rectangle(1.7229,2.5694);
\fill[DeepBlue!2.23](1.6279,2.5674)rectangle(1.7229,2.6259);
\fill[DeepBlue!3.42](1.6279,2.6239)rectangle(1.7229,2.6824);
\fill[DeepBlue!5.45](1.6279,2.6804)rectangle(1.7229,2.7390);
\fill[DeepBlue!7.24](1.6279,2.7370)rectangle(1.7229,2.7955);
\fill[DeepBlue!8.54](1.6279,2.7935)rectangle(1.7229,2.8520);
\fill[DeepBlue!9.36](1.7209,1.5500)rectangle(1.8160,1.6085);
\fill[DeepBlue!9.23](1.7209,1.6065)rectangle(1.8160,1.6650);
\fill[DeepBlue!8.57](1.7209,1.6630)rectangle(1.8160,1.7216);
\fill[DeepBlue!7.44](1.7209,1.7196)rectangle(1.8160,1.7781);
\fill[DeepBlue!5.91](1.7209,1.7761)rectangle(1.8160,1.8346);
\fill[DeepBlue!4.14](1.7209,1.8326)rectangle(1.8160,1.8911);
\fill[DeepBlue!2.37](1.7209,1.8891)rectangle(1.8160,1.9477);
\fill[DeepBlue!1.43](1.7209,1.9457)rectangle(1.8160,2.0042);
\fill[DeepBlue!2.40](1.7209,2.0022)rectangle(1.8160,2.0607);
\fill[DeepBlue!3.72](1.7209,2.0587)rectangle(1.8160,2.1172);
\fill[DeepBlue!4.74](1.7209,2.1152)rectangle(1.8160,2.1737);
\fill[DeepBlue!5.28](1.7209,2.1717)rectangle(1.8160,2.2303);
\fill[DeepBlue!5.28](1.7209,2.2283)rectangle(1.8160,2.2868);
\fill[DeepBlue!4.74](1.7209,2.2848)rectangle(1.8160,2.3433);
\fill[DeepBlue!3.72](1.7209,2.3413)rectangle(1.8160,2.3998);
\fill[DeepBlue!2.40](1.7209,2.3978)rectangle(1.8160,2.4563);
\fill[DeepBlue!1.43](1.7209,2.4543)rectangle(1.8160,2.5129);
\fill[DeepBlue!2.37](1.7209,2.5109)rectangle(1.8160,2.5694);
\fill[DeepBlue!4.14](1.7209,2.5674)rectangle(1.8160,2.6259);
\fill[DeepBlue!5.91](1.7209,2.6239)rectangle(1.8160,2.6824);
\fill[DeepBlue!7.44](1.7209,2.6804)rectangle(1.8160,2.7390);
\fill[DeepBlue!8.57](1.7209,2.7370)rectangle(1.8160,2.7955);
\fill[DeepBlue!9.23](1.7209,2.7935)rectangle(1.8160,2.8520);
\fill[DeepBlue!9.06](1.8140,1.5500)rectangle(1.9090,1.6085);
\fill[DeepBlue!9.42](1.8140,1.6065)rectangle(1.9090,1.6650);
\fill[DeepBlue!9.34](1.8140,1.6630)rectangle(1.9090,1.7216);
\fill[DeepBlue!8.86](1.8140,1.7196)rectangle(1.9090,1.7781);
\fill[DeepBlue!8.01](1.8140,1.7761)rectangle(1.9090,1.8346);
\fill[DeepBlue!6.89](1.8140,1.8326)rectangle(1.9090,1.8911);
\fill[DeepBlue!5.61](1.8140,1.8891)rectangle(1.9090,1.9477);
\fill[DeepBlue!4.29](1.8140,1.9457)rectangle(1.9090,2.0042);
\fill[DeepBlue!3.06](1.8140,2.0022)rectangle(1.9090,2.0607);
\fill[DeepBlue!2.04](1.8140,2.0587)rectangle(1.9090,2.1172);
\fill[DeepBlue!1.33](1.8140,2.1152)rectangle(1.9090,2.1737);
\fill[DeepBlue!1.00](1.8140,2.1717)rectangle(1.9090,2.2303);
\fill[DeepBlue!1.00](1.8140,2.2283)rectangle(1.9090,2.2868);
\fill[DeepBlue!1.33](1.8140,2.2848)rectangle(1.9090,2.3433);
\fill[DeepBlue!2.04](1.8140,2.3413)rectangle(1.9090,2.3998);
\fill[DeepBlue!3.06](1.8140,2.3978)rectangle(1.9090,2.4563);
\fill[DeepBlue!4.29](1.8140,2.4543)rectangle(1.9090,2.5129);
\fill[DeepBlue!5.61](1.8140,2.5109)rectangle(1.9090,2.5694);
\fill[DeepBlue!6.89](1.8140,2.5674)rectangle(1.9090,2.6259);
\fill[DeepBlue!8.01](1.8140,2.6239)rectangle(1.9090,2.6824);
\fill[DeepBlue!8.86](1.8140,2.6804)rectangle(1.9090,2.7390);
\fill[DeepBlue!9.34](1.8140,2.7370)rectangle(1.9090,2.7955);
\fill[DeepBlue!9.42](1.8140,2.7935)rectangle(1.9090,2.8520);
\fill[DeepBlue!8.32](1.9070,1.5500)rectangle(2.0020,1.6085);
\fill[DeepBlue!9.05](1.9070,1.6065)rectangle(2.0020,1.6650);
\fill[DeepBlue!9.45](1.9070,1.6630)rectangle(2.0020,1.7216);
\fill[DeepBlue!9.51](1.9070,1.7196)rectangle(2.0020,1.7781);
\fill[DeepBlue!9.26](1.9070,1.7761)rectangle(2.0020,1.8346);
\fill[DeepBlue!8.76](1.9070,1.8326)rectangle(2.0020,1.8911);
\fill[DeepBlue!8.09](1.9070,1.8891)rectangle(2.0020,1.9477);
\fill[DeepBlue!7.32](1.9070,1.9457)rectangle(2.0020,2.0042);
\fill[DeepBlue!6.57](1.9070,2.0022)rectangle(2.0020,2.0607);
\fill[DeepBlue!5.91](1.9070,2.0587)rectangle(2.0020,2.1172);
\fill[DeepBlue!5.43](1.9070,2.1152)rectangle(2.0020,2.1737);
\fill[DeepBlue!5.17](1.9070,2.1717)rectangle(2.0020,2.2303);
\fill[DeepBlue!5.17](1.9070,2.2283)rectangle(2.0020,2.2868);
\fill[DeepBlue!5.43](1.9070,2.2848)rectangle(2.0020,2.3433);
\fill[DeepBlue!5.91](1.9070,2.3413)rectangle(2.0020,2.3998);
\fill[DeepBlue!6.57](1.9070,2.3978)rectangle(2.0020,2.4563);
\fill[DeepBlue!7.32](1.9070,2.4543)rectangle(2.0020,2.5129);
\fill[DeepBlue!8.09](1.9070,2.5109)rectangle(2.0020,2.5694);
\fill[DeepBlue!8.76](1.9070,2.5674)rectangle(2.0020,2.6259);
\fill[DeepBlue!9.26](1.9070,2.6239)rectangle(2.0020,2.6824);
\fill[DeepBlue!9.51](1.9070,2.6804)rectangle(2.0020,2.7390);
\fill[DeepBlue!9.45](1.9070,2.7370)rectangle(2.0020,2.7955);
\fill[DeepBlue!9.05](1.9070,2.7935)rectangle(2.0020,2.8520);
\draw[->](-2.2,0)--(2.3,0)node[right]{$u$};
\draw[TextGray,dashed,thin]plot coordinates{(-2.200,0.001576)(-2.180,0.001535)(-2.160,0.001477)(-2.140,0.001402)(-2.120,0.001311)(-2.100,0.001206)(-2.080,0.001090)(-2.060,0.000965)(-2.040,0.000835)(-2.020,0.000702)(-2.000,0.000571)(-1.980,0.000446)(-1.960,0.000329)(-1.940,0.000224)(-1.920,0.000136)(-1.900,0.000067)(-1.880,0.000021)(-1.860,0.000001)(-1.840,0.000008)(-1.820,0.000045)(-1.800,0.000113)(-1.780,0.000211)(-1.760,0.000341)(-1.740,0.000502)(-1.720,0.000690)(-1.700,0.000906)(-1.680,0.001144)(-1.660,0.001402)(-1.640,0.001676)(-1.620,0.001959)(-1.600,0.002249)(-1.580,0.002537)(-1.560,0.002819)(-1.540,0.003089)(-1.520,0.003340)(-1.500,0.003567)(-1.480,0.003763)(-1.460,0.003925)(-1.440,0.004047)(-1.420,0.004125)(-1.400,0.004157)(-1.380,0.004141)(-1.360,0.004075)(-1.340,0.003959)(-1.320,0.003796)(-1.300,0.003588)(-1.280,0.003338)(-1.260,0.003053)(-1.240,0.002737)(-1.220,0.002400)(-1.200,0.002049)(-1.180,0.001694)(-1.160,0.001346)(-1.140,0.001014)(-1.120,0.000712)(-1.100,0.000448)(-1.080,0.000236)(-1.060,0.000086)(-1.040,0.000009)(-1.020,0.000013)(-1.000,0.000108)(-0.980,0.000300)(-0.960,0.000594)(-0.940,0.000995)(-0.920,0.001503)(-0.900,0.002118)(-0.880,0.002837)(-0.860,0.003653)(-0.840,0.004560)(-0.820,0.005547)(-0.800,0.006600)(-0.780,0.007705)(-0.760,0.008843)(-0.740,0.009997)(-0.720,0.011145)(-0.700,0.012265)(-0.680,0.013334)(-0.660,0.014328)(-0.640,0.015225)(-0.620,0.016001)(-0.600,0.016635)(-0.580,0.017107)(-0.560,0.017400)(-0.540,0.017498)(-0.520,0.017391)(-0.500,0.017071)(-0.480,0.016537)(-0.460,0.015792)(-0.440,0.014844)(-0.420,0.013709)(-0.400,0.012407)(-0.380,0.010968)(-0.360,0.009427)(-0.340,0.007826)(-0.320,0.006215)(-0.300,0.004652)(-0.280,0.003199)(-0.260,0.001929)(-0.240,0.000916)(-0.220,0.000244)(-0.200,0.000000)(-0.180,0.000275)(-0.160,0.001165)(-0.140,0.002766)(-0.120,0.005177)(-0.100,0.008497)(-0.080,0.012824)(-0.060,0.018254)(-0.040,0.024879)(-0.020,0.032785)(0.000,0.042055)(0.020,0.052761)(0.040,0.064969)(0.060,0.078734)(0.080,0.094100)(0.100,0.111099)(0.120,0.129748)(0.140,0.150053)(0.160,0.172003)(0.180,0.195571)(0.200,0.220718)(0.220,0.247383)(0.240,0.275494)(0.260,0.304960)(0.280,0.335675)(0.300,0.367516)(0.320,0.400347)(0.340,0.434018)(0.360,0.468364)(0.380,0.503209)(0.400,0.538368)(0.420,0.573645)(0.440,0.608836)(0.460,0.643732)(0.480,0.678122)(0.500,0.711788)(0.520,0.744518)(0.540,0.776096)(0.560,0.806314)(0.580,0.834967)(0.600,0.861861)(0.620,0.886808)(0.640,0.909635)(0.660,0.930179)(0.680,0.948295)(0.700,0.963852)(0.720,0.976738)(0.740,0.986859)(0.760,0.994142)(0.780,0.998533)(0.800,1.000000)(0.820,0.998533)(0.840,0.994142)(0.860,0.986859)(0.880,0.976738)(0.900,0.963852)(0.920,0.948295)(0.940,0.930179)(0.960,0.909635)(0.980,0.886808)(1.000,0.861861)(1.020,0.834967)(1.040,0.806314)(1.060,0.776096)(1.080,0.744518)(1.100,0.711788)(1.120,0.678122)(1.140,0.643732)(1.160,0.608836)(1.180,0.573645)(1.200,0.538368)(1.220,0.503209)(1.240,0.468364)(1.260,0.434018)(1.280,0.400347)(1.300,0.367516)(1.320,0.335675)(1.340,0.304960)(1.360,0.275494)(1.380,0.247383)(1.400,0.220718)(1.420,0.195571)(1.440,0.172003)(1.460,0.150053)(1.480,0.129748)(1.500,0.111099)(1.520,0.094100)(1.540,0.078734)(1.560,0.064969)(1.580,0.052761)(1.600,0.042055)(1.620,0.032785)(1.640,0.024879)(1.660,0.018254)(1.680,0.012824)(1.700,0.008497)(1.720,0.005177)(1.740,0.002766)(1.760,0.001165)(1.780,0.000275)(1.800,0.000000)(1.820,0.000244)(1.840,0.000916)(1.860,0.001929)(1.880,0.003199)(1.900,0.004652)(1.920,0.006215)(1.940,0.007826)(1.960,0.009427)(1.980,0.010968)(2.000,0.012407)(2.020,0.013709)(2.040,0.014844)(2.060,0.015792)(2.080,0.016537)(2.100,0.017071)(2.120,0.017391)(2.140,0.017498)(2.160,0.017400)(2.180,0.017107)(2.200,0.016635)};
\draw[TextGray,dashed,thin]plot coordinates{(-2.200,0.016635)(-2.180,0.017107)(-2.160,0.017400)(-2.140,0.017498)(-2.120,0.017391)(-2.100,0.017071)(-2.080,0.016537)(-2.060,0.015792)(-2.040,0.014844)(-2.020,0.013709)(-2.000,0.012407)(-1.980,0.010968)(-1.960,0.009427)(-1.940,0.007826)(-1.920,0.006215)(-1.900,0.004652)(-1.880,0.003199)(-1.860,0.001929)(-1.840,0.000916)(-1.820,0.000244)(-1.800,0.000000)(-1.780,0.000275)(-1.760,0.001165)(-1.740,0.002766)(-1.720,0.005177)(-1.700,0.008497)(-1.680,0.012824)(-1.660,0.018254)(-1.640,0.024879)(-1.620,0.032785)(-1.600,0.042055)(-1.580,0.052761)(-1.560,0.064969)(-1.540,0.078734)(-1.520,0.094100)(-1.500,0.111099)(-1.480,0.129748)(-1.460,0.150053)(-1.440,0.172003)(-1.420,0.195571)(-1.400,0.220718)(-1.380,0.247383)(-1.360,0.275494)(-1.340,0.304960)(-1.320,0.335675)(-1.300,0.367516)(-1.280,0.400347)(-1.260,0.434018)(-1.240,0.468364)(-1.220,0.503209)(-1.200,0.538368)(-1.180,0.573645)(-1.160,0.608836)(-1.140,0.643732)(-1.120,0.678122)(-1.100,0.711788)(-1.080,0.744518)(-1.060,0.776096)(-1.040,0.806314)(-1.020,0.834967)(-1.000,0.861861)(-0.980,0.886808)(-0.960,0.909635)(-0.940,0.930179)(-0.920,0.948295)(-0.900,0.963852)(-0.880,0.976738)(-0.860,0.986859)(-0.840,0.994142)(-0.820,0.998533)(-0.800,1.000000)(-0.780,0.998533)(-0.760,0.994142)(-0.740,0.986859)(-0.720,0.976738)(-0.700,0.963852)(-0.680,0.948295)(-0.660,0.930179)(-0.640,0.909635)(-0.620,0.886808)(-0.600,0.861861)(-0.580,0.834967)(-0.560,0.806314)(-0.540,0.776096)(-0.520,0.744518)(-0.500,0.711788)(-0.480,0.678122)(-0.460,0.643732)(-0.440,0.608836)(-0.420,0.573645)(-0.400,0.538368)(-0.380,0.503209)(-0.360,0.468364)(-0.340,0.434018)(-0.320,0.400347)(-0.300,0.367516)(-0.280,0.335675)(-0.260,0.304960)(-0.240,0.275494)(-0.220,0.247383)(-0.200,0.220718)(-0.180,0.195571)(-0.160,0.172003)(-0.140,0.150053)(-0.120,0.129748)(-0.100,0.111099)(-0.080,0.094100)(-0.060,0.078734)(-0.040,0.064969)(-0.020,0.052761)(0.000,0.042055)(0.020,0.032785)(0.040,0.024879)(0.060,0.018254)(0.080,0.012824)(0.100,0.008497)(0.120,0.005177)(0.140,0.002766)(0.160,0.001165)(0.180,0.000275)(0.200,0.000000)(0.220,0.000244)(0.240,0.000916)(0.260,0.001929)(0.280,0.003199)(0.300,0.004652)(0.320,0.006215)(0.340,0.007826)(0.360,0.009427)(0.380,0.010968)(0.400,0.012407)(0.420,0.013709)(0.440,0.014844)(0.460,0.015792)(0.480,0.016537)(0.500,0.017071)(0.520,0.017391)(0.540,0.017498)(0.560,0.017400)(0.580,0.017107)(0.600,0.016635)(0.620,0.016001)(0.640,0.015225)(0.660,0.014328)(0.680,0.013334)(0.700,0.012265)(0.720,0.011145)(0.740,0.009997)(0.760,0.008843)(0.780,0.007705)(0.800,0.006600)(0.820,0.005547)(0.840,0.004560)(0.860,0.003653)(0.880,0.002837)(0.900,0.002118)(0.920,0.001503)(0.940,0.000995)(0.960,0.000594)(0.980,0.000300)(1.000,0.000108)(1.020,0.000013)(1.040,0.000009)(1.060,0.000086)(1.080,0.000236)(1.100,0.000448)(1.120,0.000712)(1.140,0.001014)(1.160,0.001346)(1.180,0.001694)(1.200,0.002049)(1.220,0.002400)(1.240,0.002737)(1.260,0.003053)(1.280,0.003338)(1.300,0.003588)(1.320,0.003796)(1.340,0.003959)(1.360,0.004075)(1.380,0.004141)(1.400,0.004157)(1.420,0.004125)(1.440,0.004047)(1.460,0.003925)(1.480,0.003763)(1.500,0.003567)(1.520,0.003340)(1.540,0.003089)(1.560,0.002819)(1.580,0.002537)(1.600,0.002249)(1.620,0.001959)(1.640,0.001676)(1.660,0.001402)(1.680,0.001144)(1.700,0.000906)(1.720,0.000690)(1.740,0.000502)(1.760,0.000341)(1.780,0.000211)(1.800,0.000113)(1.820,0.000045)(1.840,0.000008)(1.860,0.000001)(1.880,0.000021)(1.900,0.000067)(1.920,0.000136)(1.940,0.000224)(1.960,0.000329)(1.980,0.000446)(2.000,0.000571)(2.020,0.000702)(2.040,0.000835)(2.060,0.000965)(2.080,0.001090)(2.100,0.001206)(2.120,0.001311)(2.140,0.001402)(2.160,0.001477)(2.180,0.001535)(2.200,0.001576)};
\draw[DeepBlue,thick]plot coordinates{(-2.200,0.018211)(-2.180,0.018643)(-2.160,0.018877)(-2.140,0.018900)(-2.120,0.018701)(-2.100,0.018277)(-2.080,0.017627)(-2.060,0.016757)(-2.040,0.015679)(-2.020,0.014411)(-2.000,0.012979)(-1.980,0.011414)(-1.960,0.009755)(-1.940,0.008050)(-1.920,0.006351)(-1.900,0.004719)(-1.880,0.003221)(-1.860,0.001929)(-1.840,0.000924)(-1.820,0.000289)(-1.800,0.000113)(-1.780,0.000487)(-1.760,0.001506)(-1.740,0.003267)(-1.720,0.005867)(-1.700,0.009403)(-1.680,0.013969)(-1.660,0.019656)(-1.640,0.026554)(-1.620,0.034744)(-1.600,0.044303)(-1.580,0.055298)(-1.560,0.067788)(-1.540,0.081823)(-1.520,0.097440)(-1.500,0.114665)(-1.480,0.133511)(-1.460,0.153978)(-1.440,0.176049)(-1.420,0.199697)(-1.400,0.224875)(-1.380,0.251524)(-1.360,0.279569)(-1.340,0.308919)(-1.320,0.339471)(-1.300,0.371104)(-1.280,0.403686)(-1.260,0.437071)(-1.240,0.471101)(-1.220,0.505610)(-1.200,0.540418)(-1.180,0.575339)(-1.160,0.610182)(-1.140,0.644747)(-1.120,0.678833)(-1.100,0.712237)(-1.080,0.744754)(-1.060,0.776182)(-1.040,0.806322)(-1.020,0.834980)(-1.000,0.861969)(-0.980,0.887108)(-0.960,0.910229)(-0.940,0.931174)(-0.920,0.949798)(-0.900,0.965970)(-0.880,0.979574)(-0.860,0.990512)(-0.840,0.998702)(-0.820,1.004079)(-0.800,1.006600)(-0.780,1.006237)(-0.760,1.002985)(-0.740,0.996856)(-0.720,0.987883)(-0.700,0.976117)(-0.680,0.961629)(-0.660,0.944507)(-0.640,0.924860)(-0.620,0.902810)(-0.600,0.878496)(-0.580,0.852075)(-0.560,0.823714)(-0.540,0.793594)(-0.520,0.761908)(-0.500,0.728860)(-0.480,0.694659)(-0.460,0.659524)(-0.440,0.623680)(-0.420,0.587353)(-0.400,0.550776)(-0.380,0.514178)(-0.360,0.477791)(-0.340,0.441844)(-0.320,0.406562)(-0.300,0.372168)(-0.280,0.338874)(-0.260,0.306889)(-0.240,0.276410)(-0.220,0.247627)(-0.200,0.220718)(-0.180,0.195847)(-0.160,0.173167)(-0.140,0.152819)(-0.120,0.134925)(-0.100,0.119596)(-0.080,0.106925)(-0.060,0.096989)(-0.040,0.089848)(-0.020,0.085546)(0.000,0.084109)(0.020,0.085546)(0.040,0.089848)(0.060,0.096989)(0.080,0.106925)(0.100,0.119596)(0.120,0.134925)(0.140,0.152819)(0.160,0.173167)(0.180,0.195847)(0.200,0.220718)(0.220,0.247627)(0.240,0.276410)(0.260,0.306889)(0.280,0.338874)(0.300,0.372168)(0.320,0.406562)(0.340,0.441844)(0.360,0.477791)(0.380,0.514178)(0.400,0.550776)(0.420,0.587353)(0.440,0.623680)(0.460,0.659524)(0.480,0.694659)(0.500,0.728860)(0.520,0.761908)(0.540,0.793594)(0.560,0.823714)(0.580,0.852075)(0.600,0.878496)(0.620,0.902810)(0.640,0.924860)(0.660,0.944507)(0.680,0.961629)(0.700,0.976117)(0.720,0.987883)(0.740,0.996856)(0.760,1.002985)(0.780,1.006237)(0.800,1.006600)(0.820,1.004079)(0.840,0.998702)(0.860,0.990512)(0.880,0.979574)(0.900,0.965970)(0.920,0.949798)(0.940,0.931174)(0.960,0.910229)(0.980,0.887108)(1.000,0.861969)(1.020,0.834980)(1.040,0.806322)(1.060,0.776182)(1.080,0.744754)(1.100,0.712237)(1.120,0.678833)(1.140,0.644747)(1.160,0.610182)(1.180,0.575339)(1.200,0.540418)(1.220,0.505610)(1.240,0.471101)(1.260,0.437071)(1.280,0.403686)(1.300,0.371104)(1.320,0.339471)(1.340,0.308919)(1.360,0.279569)(1.380,0.251524)(1.400,0.224875)(1.420,0.199697)(1.440,0.176049)(1.460,0.153978)(1.480,0.133511)(1.500,0.114665)(1.520,0.097440)(1.540,0.081823)(1.560,0.067788)(1.580,0.055298)(1.600,0.044303)(1.620,0.034744)(1.640,0.026554)(1.660,0.019656)(1.680,0.013969)(1.700,0.009403)(1.720,0.005867)(1.740,0.003267)(1.760,0.001506)(1.780,0.000487)(1.800,0.000113)(1.820,0.000289)(1.840,0.000924)(1.860,0.001929)(1.880,0.003221)(1.900,0.004719)(1.920,0.006351)(1.940,0.008050)(1.960,0.009755)(1.980,0.011414)(2.000,0.012979)(2.020,0.014411)(2.040,0.015679)(2.060,0.016757)(2.080,0.017627)(2.100,0.018277)(2.120,0.018701)(2.140,0.018900)(2.160,0.018877)(2.180,0.018643)(2.200,0.018211)};
\node[below,align=center]at(0,-.2){$\Delta\theta/\theta_{\rm A}=1.6$\\可分辨};
\end{scope}
\begin{scope}[xshift=4.6cm]
\fill[DeepBlue!0.83](-2.0000,1.5500)rectangle(-1.9050,1.6085);
\fill[DeepBlue!1.83](-2.0000,1.6065)rectangle(-1.9050,1.6650);
\fill[DeepBlue!2.95](-2.0000,1.6630)rectangle(-1.9050,1.7216);
\fill[DeepBlue!4.02](-2.0000,1.7196)rectangle(-1.9050,1.7781);
\fill[DeepBlue!4.98](-2.0000,1.7761)rectangle(-1.9050,1.8346);
\fill[DeepBlue!5.81](-2.0000,1.8326)rectangle(-1.9050,1.8911);
\fill[DeepBlue!6.51](-2.0000,1.8891)rectangle(-1.9050,1.9477);
\fill[DeepBlue!7.07](-2.0000,1.9457)rectangle(-1.9050,2.0042);
\fill[DeepBlue!7.50](-2.0000,2.0022)rectangle(-1.9050,2.0607);
\fill[DeepBlue!7.81](-2.0000,2.0587)rectangle(-1.9050,2.1172);
\fill[DeepBlue!8.01](-2.0000,2.1152)rectangle(-1.9050,2.1737);
\fill[DeepBlue!8.10](-2.0000,2.1717)rectangle(-1.9050,2.2303);
\fill[DeepBlue!8.10](-2.0000,2.2283)rectangle(-1.9050,2.2868);
\fill[DeepBlue!8.01](-2.0000,2.2848)rectangle(-1.9050,2.3433);
\fill[DeepBlue!7.81](-2.0000,2.3413)rectangle(-1.9050,2.3998);
\fill[DeepBlue!7.50](-2.0000,2.3978)rectangle(-1.9050,2.4563);
\fill[DeepBlue!7.07](-2.0000,2.4543)rectangle(-1.9050,2.5129);
\fill[DeepBlue!6.51](-2.0000,2.5109)rectangle(-1.9050,2.5694);
\fill[DeepBlue!5.81](-2.0000,2.5674)rectangle(-1.9050,2.6259);
\fill[DeepBlue!4.98](-2.0000,2.6239)rectangle(-1.9050,2.6824);
\fill[DeepBlue!4.02](-2.0000,2.6804)rectangle(-1.9050,2.7390);
\fill[DeepBlue!2.95](-2.0000,2.7370)rectangle(-1.9050,2.7955);
\fill[DeepBlue!1.83](-2.0000,2.7935)rectangle(-1.9050,2.8520);
\fill[DeepBlue!2.67](-1.9070,1.5500)rectangle(-1.8120,1.6085);
\fill[DeepBlue!3.99](-1.9070,1.6065)rectangle(-1.8120,1.6650);
\fill[DeepBlue!5.21](-1.9070,1.6630)rectangle(-1.8120,1.7216);
\fill[DeepBlue!6.28](-1.9070,1.7196)rectangle(-1.8120,1.7781);
\fill[DeepBlue!7.19](-1.9070,1.7761)rectangle(-1.8120,1.8346);
\fill[DeepBlue!7.94](-1.9070,1.8326)rectangle(-1.8120,1.8911);
\fill[DeepBlue!8.51](-1.9070,1.8891)rectangle(-1.8120,1.9477);
\fill[DeepBlue!8.95](-1.9070,1.9457)rectangle(-1.8120,2.0042);
\fill[DeepBlue!9.25](-1.9070,2.0022)rectangle(-1.8120,2.0607);
\fill[DeepBlue!9.45](-1.9070,2.0587)rectangle(-1.8120,2.1172);
\fill[DeepBlue!9.58](-1.9070,2.1152)rectangle(-1.8120,2.1737);
\fill[DeepBlue!9.63](-1.9070,2.1717)rectangle(-1.8120,2.2303);
\fill[DeepBlue!9.63](-1.9070,2.2283)rectangle(-1.8120,2.2868);
\fill[DeepBlue!9.58](-1.9070,2.2848)rectangle(-1.8120,2.3433);
\fill[DeepBlue!9.45](-1.9070,2.3413)rectangle(-1.8120,2.3998);
\fill[DeepBlue!9.25](-1.9070,2.3978)rectangle(-1.8120,2.4563);
\fill[DeepBlue!8.95](-1.9070,2.4543)rectangle(-1.8120,2.5129);
\fill[DeepBlue!8.51](-1.9070,2.5109)rectangle(-1.8120,2.5694);
\fill[DeepBlue!7.94](-1.9070,2.5674)rectangle(-1.8120,2.6259);
\fill[DeepBlue!7.19](-1.9070,2.6239)rectangle(-1.8120,2.6824);
\fill[DeepBlue!6.28](-1.9070,2.6804)rectangle(-1.8120,2.7390);
\fill[DeepBlue!5.21](-1.9070,2.7370)rectangle(-1.8120,2.7955);
\fill[DeepBlue!3.99](-1.9070,2.7935)rectangle(-1.8120,2.8520);
\fill[DeepBlue!4.86](-1.8140,1.5500)rectangle(-1.7189,1.6085);
\fill[DeepBlue!6.18](-1.8140,1.6065)rectangle(-1.7189,1.6650);
\fill[DeepBlue!7.33](-1.8140,1.6630)rectangle(-1.7189,1.7216);
\fill[DeepBlue!8.27](-1.8140,1.7196)rectangle(-1.7189,1.7781);
\fill[DeepBlue!9.00](-1.8140,1.7761)rectangle(-1.7189,1.8346);
\fill[DeepBlue!9.52](-1.8140,1.8326)rectangle(-1.7189,1.8911);
\fill[DeepBlue!9.86](-1.8140,1.8891)rectangle(-1.7189,1.9477);
\fill[DeepBlue!10.07](-1.8140,1.9457)rectangle(-1.7189,2.0042);
\fill[DeepBlue!10.16](-1.8140,2.0022)rectangle(-1.7189,2.0607);
\fill[DeepBlue!10.19](-1.8140,2.0587)rectangle(-1.7189,2.1172);
\fill[DeepBlue!10.19](-1.8140,2.1152)rectangle(-1.7189,2.1737);
\fill[DeepBlue!10.18](-1.8140,2.1717)rectangle(-1.7189,2.2303);
\fill[DeepBlue!10.18](-1.8140,2.2283)rectangle(-1.7189,2.2868);
\fill[DeepBlue!10.19](-1.8140,2.2848)rectangle(-1.7189,2.3433);
\fill[DeepBlue!10.19](-1.8140,2.3413)rectangle(-1.7189,2.3998);
\fill[DeepBlue!10.16](-1.8140,2.3978)rectangle(-1.7189,2.4563);
\fill[DeepBlue!10.07](-1.8140,2.4543)rectangle(-1.7189,2.5129);
\fill[DeepBlue!9.86](-1.8140,2.5109)rectangle(-1.7189,2.5694);
\fill[DeepBlue!9.52](-1.8140,2.5674)rectangle(-1.7189,2.6259);
\fill[DeepBlue!9.00](-1.8140,2.6239)rectangle(-1.7189,2.6824);
\fill[DeepBlue!8.27](-1.8140,2.6804)rectangle(-1.7189,2.7390);
\fill[DeepBlue!7.33](-1.8140,2.7370)rectangle(-1.7189,2.7955);
\fill[DeepBlue!6.18](-1.8140,2.7935)rectangle(-1.7189,2.8520);
\fill[DeepBlue!6.91](-1.7209,1.5500)rectangle(-1.6259,1.6085);
\fill[DeepBlue!8.10](-1.7209,1.6065)rectangle(-1.6259,1.6650);
\fill[DeepBlue!9.03](-1.7209,1.6630)rectangle(-1.6259,1.7216);
\fill[DeepBlue!9.70](-1.7209,1.7196)rectangle(-1.6259,1.7781);
\fill[DeepBlue!10.11](-1.7209,1.7761)rectangle(-1.6259,1.8346);
\fill[DeepBlue!10.28](-1.7209,1.8326)rectangle(-1.6259,1.8911);
\fill[DeepBlue!10.27](-1.7209,1.8891)rectangle(-1.6259,1.9477);
\fill[DeepBlue!10.14](-1.7209,1.9457)rectangle(-1.6259,2.0042);
\fill[DeepBlue!9.94](-1.7209,2.0022)rectangle(-1.6259,2.0607);
\fill[DeepBlue!9.73](-1.7209,2.0587)rectangle(-1.6259,2.1172);
\fill[DeepBlue!9.56](-1.7209,2.1152)rectangle(-1.6259,2.1737);
\fill[DeepBlue!9.47](-1.7209,2.1717)rectangle(-1.6259,2.2303);
\fill[DeepBlue!9.47](-1.7209,2.2283)rectangle(-1.6259,2.2868);
\fill[DeepBlue!9.56](-1.7209,2.2848)rectangle(-1.6259,2.3433);
\fill[DeepBlue!9.73](-1.7209,2.3413)rectangle(-1.6259,2.3998);
\fill[DeepBlue!9.94](-1.7209,2.3978)rectangle(-1.6259,2.4563);
\fill[DeepBlue!10.14](-1.7209,2.4543)rectangle(-1.6259,2.5129);
\fill[DeepBlue!10.27](-1.7209,2.5109)rectangle(-1.6259,2.5694);
\fill[DeepBlue!10.28](-1.7209,2.5674)rectangle(-1.6259,2.6259);
\fill[DeepBlue!10.11](-1.7209,2.6239)rectangle(-1.6259,2.6824);
\fill[DeepBlue!9.70](-1.7209,2.6804)rectangle(-1.6259,2.7390);
\fill[DeepBlue!9.03](-1.7209,2.7370)rectangle(-1.6259,2.7955);
\fill[DeepBlue!8.10](-1.7209,2.7935)rectangle(-1.6259,2.8520);
\fill[DeepBlue!8.60](-1.6279,1.5500)rectangle(-1.5329,1.6085);
\fill[DeepBlue!9.52](-1.6279,1.6065)rectangle(-1.5329,1.6650);
\fill[DeepBlue!10.11](-1.6279,1.6630)rectangle(-1.5329,1.7216);
\fill[DeepBlue!10.36](-1.6279,1.7196)rectangle(-1.5329,1.7781);
\fill[DeepBlue!10.30](-1.6279,1.7761)rectangle(-1.5329,1.8346);
\fill[DeepBlue!10.00](-1.6279,1.8326)rectangle(-1.5329,1.8911);
\fill[DeepBlue!9.53](-1.6279,1.8891)rectangle(-1.5329,1.9477);
\fill[DeepBlue!8.96](-1.6279,1.9457)rectangle(-1.5329,2.0042);
\fill[DeepBlue!8.39](-1.6279,2.0022)rectangle(-1.5329,2.0607);
\fill[DeepBlue!7.90](-1.6279,2.0587)rectangle(-1.5329,2.1172);
\fill[DeepBlue!7.53](-1.6279,2.1152)rectangle(-1.5329,2.1737);
\fill[DeepBlue!7.34](-1.6279,2.1717)rectangle(-1.5329,2.2303);
\fill[DeepBlue!7.34](-1.6279,2.2283)rectangle(-1.5329,2.2868);
\fill[DeepBlue!7.53](-1.6279,2.2848)rectangle(-1.5329,2.3433);
\fill[DeepBlue!7.90](-1.6279,2.3413)rectangle(-1.5329,2.3998);
\fill[DeepBlue!8.39](-1.6279,2.3978)rectangle(-1.5329,2.4563);
\fill[DeepBlue!8.96](-1.6279,2.4543)rectangle(-1.5329,2.5129);
\fill[DeepBlue!9.53](-1.6279,2.5109)rectangle(-1.5329,2.5694);
\fill[DeepBlue!10.00](-1.6279,2.5674)rectangle(-1.5329,2.6259);
\fill[DeepBlue!10.30](-1.6279,2.6239)rectangle(-1.5329,2.6824);
\fill[DeepBlue!10.36](-1.6279,2.6804)rectangle(-1.5329,2.7390);
\fill[DeepBlue!10.11](-1.6279,2.7370)rectangle(-1.5329,2.7955);
\fill[DeepBlue!9.52](-1.6279,2.7935)rectangle(-1.5329,2.8520);
\fill[DeepBlue!9.78](-1.5349,1.5500)rectangle(-1.4399,1.6085);
\fill[DeepBlue!10.29](-1.5349,1.6065)rectangle(-1.4399,1.6650);
\fill[DeepBlue!10.39](-1.5349,1.6630)rectangle(-1.4399,1.7216);
\fill[DeepBlue!10.09](-1.5349,1.7196)rectangle(-1.4399,1.7781);
\fill[DeepBlue!9.44](-1.5349,1.7761)rectangle(-1.4399,1.8346);
\fill[DeepBlue!8.55](-1.5349,1.8326)rectangle(-1.4399,1.8911);
\fill[DeepBlue!7.52](-1.5349,1.8891)rectangle(-1.4399,1.9477);
\fill[DeepBlue!6.47](-1.5349,1.9457)rectangle(-1.4399,2.0042);
\fill[DeepBlue!5.53](-1.5349,2.0022)rectangle(-1.4399,2.0607);
\fill[DeepBlue!4.81](-1.5349,2.0587)rectangle(-1.4399,2.1172);
\fill[DeepBlue!4.34](-1.5349,2.1152)rectangle(-1.4399,2.1737);
\fill[DeepBlue!4.13](-1.5349,2.1717)rectangle(-1.4399,2.2303);
\fill[DeepBlue!4.13](-1.5349,2.2283)rectangle(-1.4399,2.2868);
\fill[DeepBlue!4.34](-1.5349,2.2848)rectangle(-1.4399,2.3433);
\fill[DeepBlue!4.81](-1.5349,2.3413)rectangle(-1.4399,2.3998);
\fill[DeepBlue!5.53](-1.5349,2.3978)rectangle(-1.4399,2.4563);
\fill[DeepBlue!6.47](-1.5349,2.4543)rectangle(-1.4399,2.5129);
\fill[DeepBlue!7.52](-1.5349,2.5109)rectangle(-1.4399,2.5694);
\fill[DeepBlue!8.55](-1.5349,2.5674)rectangle(-1.4399,2.6259);
\fill[DeepBlue!9.44](-1.5349,2.6239)rectangle(-1.4399,2.6824);
\fill[DeepBlue!10.09](-1.5349,2.6804)rectangle(-1.4399,2.7390);
\fill[DeepBlue!10.39](-1.5349,2.7370)rectangle(-1.4399,2.7955);
\fill[DeepBlue!10.29](-1.5349,2.7935)rectangle(-1.4399,2.8520);
\fill[DeepBlue!10.34](-1.4419,1.5500)rectangle(-1.3468,1.6085);
\fill[DeepBlue!10.32](-1.4419,1.6065)rectangle(-1.3468,1.6650);
\fill[DeepBlue!9.79](-1.4419,1.6630)rectangle(-1.3468,1.7216);
\fill[DeepBlue!8.81](-1.4419,1.7196)rectangle(-1.3468,1.7781);
\fill[DeepBlue!7.47](-1.4419,1.7761)rectangle(-1.3468,1.8346);
\fill[DeepBlue!5.91](-1.4419,1.8326)rectangle(-1.3468,1.8911);
\fill[DeepBlue!4.34](-1.4419,1.8891)rectangle(-1.3468,1.9477);
\fill[DeepBlue!3.12](-1.4419,1.9457)rectangle(-1.3468,2.0042);
\fill[DeepBlue!2.75](-1.4419,2.0022)rectangle(-1.3468,2.0607);
\fill[DeepBlue!3.19](-1.4419,2.0587)rectangle(-1.3468,2.1172);
\fill[DeepBlue!3.82](-1.4419,2.1152)rectangle(-1.3468,2.1737);
\fill[DeepBlue!4.21](-1.4419,2.1717)rectangle(-1.3468,2.2303);
\fill[DeepBlue!4.21](-1.4419,2.2283)rectangle(-1.3468,2.2868);
\fill[DeepBlue!3.82](-1.4419,2.2848)rectangle(-1.3468,2.3433);
\fill[DeepBlue!3.19](-1.4419,2.3413)rectangle(-1.3468,2.3998);
\fill[DeepBlue!2.75](-1.4419,2.3978)rectangle(-1.3468,2.4563);
\fill[DeepBlue!3.12](-1.4419,2.4543)rectangle(-1.3468,2.5129);
\fill[DeepBlue!4.34](-1.4419,2.5109)rectangle(-1.3468,2.5694);
\fill[DeepBlue!5.91](-1.4419,2.5674)rectangle(-1.3468,2.6259);
\fill[DeepBlue!7.47](-1.4419,2.6239)rectangle(-1.3468,2.6824);
\fill[DeepBlue!8.81](-1.4419,2.6804)rectangle(-1.3468,2.7390);
\fill[DeepBlue!9.79](-1.4419,2.7370)rectangle(-1.3468,2.7955);
\fill[DeepBlue!10.32](-1.4419,2.7935)rectangle(-1.3468,2.8520);
\fill[DeepBlue!10.22](-1.3488,1.5500)rectangle(-1.2538,1.6085);
\fill[DeepBlue!9.56](-1.3488,1.6065)rectangle(-1.2538,1.6650);
\fill[DeepBlue!8.30](-1.3488,1.6630)rectangle(-1.2538,1.7216);
\fill[DeepBlue!6.54](-1.3488,1.7196)rectangle(-1.2538,1.7781);
\fill[DeepBlue!4.43](-1.3488,1.7761)rectangle(-1.2538,1.8346);
\fill[DeepBlue!2.40](-1.3488,1.8326)rectangle(-1.2538,1.8911);
\fill[DeepBlue!2.27](-1.3488,1.8891)rectangle(-1.2538,1.9477);
\fill[DeepBlue!4.28](-1.3488,1.9457)rectangle(-1.2538,2.0042);
\fill[DeepBlue!6.46](-1.3488,2.0022)rectangle(-1.2538,2.0607);
\fill[DeepBlue!8.32](-1.3488,2.0587)rectangle(-1.2538,2.1172);
\fill[DeepBlue!9.66](-1.3488,2.1152)rectangle(-1.2538,2.1737);
\fill[DeepBlue!10.36](-1.3488,2.1717)rectangle(-1.2538,2.2303);
\fill[DeepBlue!10.36](-1.3488,2.2283)rectangle(-1.2538,2.2868);
\fill[DeepBlue!9.66](-1.3488,2.2848)rectangle(-1.2538,2.3433);
\fill[DeepBlue!8.32](-1.3488,2.3413)rectangle(-1.2538,2.3998);
\fill[DeepBlue!6.46](-1.3488,2.3978)rectangle(-1.2538,2.4563);
\fill[DeepBlue!4.28](-1.3488,2.4543)rectangle(-1.2538,2.5129);
\fill[DeepBlue!2.27](-1.3488,2.5109)rectangle(-1.2538,2.5694);
\fill[DeepBlue!2.40](-1.3488,2.5674)rectangle(-1.2538,2.6259);
\fill[DeepBlue!4.43](-1.3488,2.6239)rectangle(-1.2538,2.6824);
\fill[DeepBlue!6.54](-1.3488,2.6804)rectangle(-1.2538,2.7390);
\fill[DeepBlue!8.30](-1.3488,2.7370)rectangle(-1.2538,2.7955);
\fill[DeepBlue!9.56](-1.3488,2.7935)rectangle(-1.2538,2.8520);
\fill[DeepBlue!9.46](-1.2558,1.5500)rectangle(-1.1608,1.6085);
\fill[DeepBlue!8.07](-1.2558,1.6065)rectangle(-1.1608,1.6650);
\fill[DeepBlue!6.00](-1.2558,1.6630)rectangle(-1.1608,1.7216);
\fill[DeepBlue!3.37](-1.2558,1.7196)rectangle(-1.1608,1.7781);
\fill[DeepBlue!0.94](-1.2558,1.7761)rectangle(-1.1608,1.8346);
\fill[DeepBlue!3.48](-1.2558,1.8326)rectangle(-1.1608,1.8911);
\fill[DeepBlue!6.93](-1.2558,1.8891)rectangle(-1.1608,1.9477);
\fill[DeepBlue!10.29](-1.2558,1.9457)rectangle(-1.1608,2.0042);
\fill[DeepBlue!13.29](-1.2558,2.0022)rectangle(-1.1608,2.0607);
\fill[DeepBlue!15.75](-1.2558,2.0587)rectangle(-1.1608,2.1172);
\fill[DeepBlue!17.48](-1.2558,2.1152)rectangle(-1.1608,2.1737);
\fill[DeepBlue!18.38](-1.2558,2.1717)rectangle(-1.1608,2.2303);
\fill[DeepBlue!18.38](-1.2558,2.2283)rectangle(-1.1608,2.2868);
\fill[DeepBlue!17.48](-1.2558,2.2848)rectangle(-1.1608,2.3433);
\fill[DeepBlue!15.75](-1.2558,2.3413)rectangle(-1.1608,2.3998);
\fill[DeepBlue!13.29](-1.2558,2.3978)rectangle(-1.1608,2.4563);
\fill[DeepBlue!10.29](-1.2558,2.4543)rectangle(-1.1608,2.5129);
\fill[DeepBlue!6.93](-1.2558,2.5109)rectangle(-1.1608,2.5694);
\fill[DeepBlue!3.48](-1.2558,2.5674)rectangle(-1.1608,2.6259);
\fill[DeepBlue!0.94](-1.2558,2.6239)rectangle(-1.1608,2.6824);
\fill[DeepBlue!3.37](-1.2558,2.6804)rectangle(-1.1608,2.7390);
\fill[DeepBlue!6.00](-1.2558,2.7370)rectangle(-1.1608,2.7955);
\fill[DeepBlue!8.07](-1.2558,2.7935)rectangle(-1.1608,2.8520);
\fill[DeepBlue!8.19](-1.1628,1.5500)rectangle(-1.0678,1.6085);
\fill[DeepBlue!6.04](-1.1628,1.6065)rectangle(-1.0678,1.6650);
\fill[DeepBlue!3.12](-1.1628,1.6630)rectangle(-1.0678,1.7216);
\fill[DeepBlue!0.49](-1.1628,1.7196)rectangle(-1.0678,1.7781);
\fill[DeepBlue!4.53](-1.1628,1.7761)rectangle(-1.0678,1.8346);
\fill[DeepBlue!8.88](-1.1628,1.8326)rectangle(-1.0678,1.8911);
\fill[DeepBlue!13.26](-1.1628,1.8891)rectangle(-1.0678,1.9477);
\fill[DeepBlue!17.42](-1.1628,1.9457)rectangle(-1.0678,2.0042);
\fill[DeepBlue!21.12](-1.1628,2.0022)rectangle(-1.0678,2.0607);
\fill[DeepBlue!24.11](-1.1628,2.0587)rectangle(-1.0678,2.1172);
\fill[DeepBlue!26.22](-1.1628,2.1152)rectangle(-1.0678,2.1737);
\fill[DeepBlue!27.31](-1.1628,2.1717)rectangle(-1.0678,2.2303);
\fill[DeepBlue!27.31](-1.1628,2.2283)rectangle(-1.0678,2.2868);
\fill[DeepBlue!26.22](-1.1628,2.2848)rectangle(-1.0678,2.3433);
\fill[DeepBlue!24.11](-1.1628,2.3413)rectangle(-1.0678,2.3998);
\fill[DeepBlue!21.12](-1.1628,2.3978)rectangle(-1.0678,2.4563);
\fill[DeepBlue!17.42](-1.1628,2.4543)rectangle(-1.0678,2.5129);
\fill[DeepBlue!13.26](-1.1628,2.5109)rectangle(-1.0678,2.5694);
\fill[DeepBlue!8.88](-1.1628,2.5674)rectangle(-1.0678,2.6259);
\fill[DeepBlue!4.53](-1.1628,2.6239)rectangle(-1.0678,2.6824);
\fill[DeepBlue!0.49](-1.1628,2.6804)rectangle(-1.0678,2.7390);
\fill[DeepBlue!3.12](-1.1628,2.7370)rectangle(-1.0678,2.7955);
\fill[DeepBlue!6.04](-1.1628,2.7935)rectangle(-1.0678,2.8520);
\fill[DeepBlue!6.63](-1.0698,1.5500)rectangle(-0.9747,1.6085);
\fill[DeepBlue!3.85](-1.0698,1.6065)rectangle(-0.9747,1.6650);
\fill[DeepBlue!1.39](-1.0698,1.6630)rectangle(-0.9747,1.7216);
\fill[DeepBlue!4.79](-1.0698,1.7196)rectangle(-0.9747,1.7781);
\fill[DeepBlue!9.65](-1.0698,1.7761)rectangle(-0.9747,1.8346);
\fill[DeepBlue!14.85](-1.0698,1.8326)rectangle(-0.9747,1.8911);
\fill[DeepBlue!20.06](-1.0698,1.8891)rectangle(-0.9747,1.9477);
\fill[DeepBlue!24.98](-1.0698,1.9457)rectangle(-0.9747,2.0042);
\fill[DeepBlue!29.32](-1.0698,2.0022)rectangle(-0.9747,2.0607);
\fill[DeepBlue!32.83](-1.0698,2.0587)rectangle(-0.9747,2.1172);
\fill[DeepBlue!35.30](-1.0698,2.1152)rectangle(-0.9747,2.1737);
\fill[DeepBlue!36.57](-1.0698,2.1717)rectangle(-0.9747,2.2303);
\fill[DeepBlue!36.57](-1.0698,2.2283)rectangle(-0.9747,2.2868);
\fill[DeepBlue!35.30](-1.0698,2.2848)rectangle(-0.9747,2.3433);
\fill[DeepBlue!32.83](-1.0698,2.3413)rectangle(-0.9747,2.3998);
\fill[DeepBlue!29.32](-1.0698,2.3978)rectangle(-0.9747,2.4563);
\fill[DeepBlue!24.98](-1.0698,2.4543)rectangle(-0.9747,2.5129);
\fill[DeepBlue!20.06](-1.0698,2.5109)rectangle(-0.9747,2.5694);
\fill[DeepBlue!14.85](-1.0698,2.5674)rectangle(-0.9747,2.6259);
\fill[DeepBlue!9.65](-1.0698,2.6239)rectangle(-0.9747,2.6824);
\fill[DeepBlue!4.79](-1.0698,2.6804)rectangle(-0.9747,2.7390);
\fill[DeepBlue!1.39](-1.0698,2.7370)rectangle(-0.9747,2.7955);
\fill[DeepBlue!3.85](-1.0698,2.7935)rectangle(-0.9747,2.8520);
\fill[DeepBlue!5.25](-0.9767,1.5500)rectangle(-0.8817,1.6085);
\fill[DeepBlue!2.84](-0.9767,1.6065)rectangle(-0.8817,1.6650);
\fill[DeepBlue!4.46](-0.9767,1.6630)rectangle(-0.8817,1.7216);
\fill[DeepBlue!9.21](-0.9767,1.7196)rectangle(-0.8817,1.7781);
\fill[DeepBlue!14.85](-0.9767,1.7761)rectangle(-0.8817,1.8346);
\fill[DeepBlue!20.83](-0.9767,1.8326)rectangle(-0.8817,1.8911);
\fill[DeepBlue!26.80](-0.9767,1.8891)rectangle(-0.8817,1.9477);
\fill[DeepBlue!32.42](-0.9767,1.9457)rectangle(-0.8817,2.0042);
\fill[DeepBlue!37.37](-0.9767,2.0022)rectangle(-0.8817,2.0607);
\fill[DeepBlue!41.36](-0.9767,2.0587)rectangle(-0.8817,2.1172);
\fill[DeepBlue!44.16](-0.9767,2.1152)rectangle(-0.8817,2.1737);
\fill[DeepBlue!45.60](-0.9767,2.1717)rectangle(-0.8817,2.2303);
\fill[DeepBlue!45.60](-0.9767,2.2283)rectangle(-0.8817,2.2868);
\fill[DeepBlue!44.16](-0.9767,2.2848)rectangle(-0.8817,2.3433);
\fill[DeepBlue!41.36](-0.9767,2.3413)rectangle(-0.8817,2.3998);
\fill[DeepBlue!37.37](-0.9767,2.3978)rectangle(-0.8817,2.4563);
\fill[DeepBlue!32.42](-0.9767,2.4543)rectangle(-0.8817,2.5129);
\fill[DeepBlue!26.80](-0.9767,2.5109)rectangle(-0.8817,2.5694);
\fill[DeepBlue!20.83](-0.9767,2.5674)rectangle(-0.8817,2.6259);
\fill[DeepBlue!14.85](-0.9767,2.6239)rectangle(-0.8817,2.6824);
\fill[DeepBlue!9.21](-0.9767,2.6804)rectangle(-0.8817,2.7390);
\fill[DeepBlue!4.46](-0.9767,2.7370)rectangle(-0.8817,2.7955);
\fill[DeepBlue!2.84](-0.9767,2.7935)rectangle(-0.8817,2.8520);
\fill[DeepBlue!4.74](-0.8837,1.5500)rectangle(-0.7887,1.6085);
\fill[DeepBlue!4.36](-0.8837,1.6065)rectangle(-0.7887,1.6650);
\fill[DeepBlue!7.92](-0.8837,1.6630)rectangle(-0.7887,1.7216);
\fill[DeepBlue!13.40](-0.8837,1.7196)rectangle(-0.7887,1.7781);
\fill[DeepBlue!19.71](-0.8837,1.7761)rectangle(-0.7887,1.8346);
\fill[DeepBlue!26.39](-0.8837,1.8326)rectangle(-0.7887,1.8911);
\fill[DeepBlue!33.02](-0.8837,1.8891)rectangle(-0.7887,1.9477);
\fill[DeepBlue!39.26](-0.8837,1.9457)rectangle(-0.7887,2.0042);
\fill[DeepBlue!44.74](-0.8837,2.0022)rectangle(-0.7887,2.0607);
\fill[DeepBlue!49.15](-0.8837,2.0587)rectangle(-0.7887,2.1172);
\fill[DeepBlue!52.25](-0.8837,2.1152)rectangle(-0.7887,2.1737);
\fill[DeepBlue!53.84](-0.8837,2.1717)rectangle(-0.7887,2.2303);
\fill[DeepBlue!53.84](-0.8837,2.2283)rectangle(-0.7887,2.2868);
\fill[DeepBlue!52.25](-0.8837,2.2848)rectangle(-0.7887,2.3433);
\fill[DeepBlue!49.15](-0.8837,2.3413)rectangle(-0.7887,2.3998);
\fill[DeepBlue!44.74](-0.8837,2.3978)rectangle(-0.7887,2.4563);
\fill[DeepBlue!39.26](-0.8837,2.4543)rectangle(-0.7887,2.5129);
\fill[DeepBlue!33.02](-0.8837,2.5109)rectangle(-0.7887,2.5694);
\fill[DeepBlue!26.39](-0.8837,2.5674)rectangle(-0.7887,2.6259);
\fill[DeepBlue!19.71](-0.8837,2.6239)rectangle(-0.7887,2.6824);
\fill[DeepBlue!13.40](-0.8837,2.6804)rectangle(-0.7887,2.7390);
\fill[DeepBlue!7.92](-0.8837,2.7370)rectangle(-0.7887,2.7955);
\fill[DeepBlue!4.36](-0.8837,2.7935)rectangle(-0.7887,2.8520);
\fill[DeepBlue!5.42](-0.7907,1.5500)rectangle(-0.6957,1.6085);
\fill[DeepBlue!6.68](-0.7907,1.6065)rectangle(-0.6957,1.6650);
\fill[DeepBlue!11.03](-0.7907,1.6630)rectangle(-0.6957,1.7216);
\fill[DeepBlue!17.02](-0.7907,1.7196)rectangle(-0.6957,1.7781);
\fill[DeepBlue!23.87](-0.7907,1.7761)rectangle(-0.6957,1.8346);
\fill[DeepBlue!31.10](-0.7907,1.8326)rectangle(-0.6957,1.8911);
\fill[DeepBlue!38.28](-0.7907,1.8891)rectangle(-0.6957,1.9477);
\fill[DeepBlue!45.03](-0.7907,1.9457)rectangle(-0.6957,2.0042);
\fill[DeepBlue!50.95](-0.7907,2.0022)rectangle(-0.6957,2.0607);
\fill[DeepBlue!55.71](-0.7907,2.0587)rectangle(-0.6957,2.1172);
\fill[DeepBlue!59.05](-0.7907,2.1152)rectangle(-0.6957,2.1737);
\fill[DeepBlue!60.77](-0.7907,2.1717)rectangle(-0.6957,2.2303);
\fill[DeepBlue!60.77](-0.7907,2.2283)rectangle(-0.6957,2.2868);
\fill[DeepBlue!59.05](-0.7907,2.2848)rectangle(-0.6957,2.3433);
\fill[DeepBlue!55.71](-0.7907,2.3413)rectangle(-0.6957,2.3998);
\fill[DeepBlue!50.95](-0.7907,2.3978)rectangle(-0.6957,2.4563);
\fill[DeepBlue!45.03](-0.7907,2.4543)rectangle(-0.6957,2.5129);
\fill[DeepBlue!38.28](-0.7907,2.5109)rectangle(-0.6957,2.5694);
\fill[DeepBlue!31.10](-0.7907,2.5674)rectangle(-0.6957,2.6259);
\fill[DeepBlue!23.87](-0.7907,2.6239)rectangle(-0.6957,2.6824);
\fill[DeepBlue!17.02](-0.7907,2.6804)rectangle(-0.6957,2.7390);
\fill[DeepBlue!11.03](-0.7907,2.7370)rectangle(-0.6957,2.7955);
\fill[DeepBlue!6.68](-0.7907,2.7935)rectangle(-0.6957,2.8520);
\fill[DeepBlue!6.73](-0.6977,1.5500)rectangle(-0.6027,1.6085);
\fill[DeepBlue!8.75](-0.6977,1.6065)rectangle(-0.6027,1.6650);
\fill[DeepBlue!13.45](-0.6977,1.6630)rectangle(-0.6027,1.7216);
\fill[DeepBlue!19.76](-0.6977,1.7196)rectangle(-0.6027,1.7781);
\fill[DeepBlue!26.99](-0.6977,1.7761)rectangle(-0.6027,1.8346);
\fill[DeepBlue!34.62](-0.6977,1.8326)rectangle(-0.6027,1.8911);
\fill[DeepBlue!42.21](-0.6977,1.8891)rectangle(-0.6027,1.9477);
\fill[DeepBlue!49.34](-0.6977,1.9457)rectangle(-0.6027,2.0042);
\fill[DeepBlue!55.60](-0.6977,2.0022)rectangle(-0.6027,2.0607);
\fill[DeepBlue!60.64](-0.6977,2.0587)rectangle(-0.6027,2.1172);
\fill[DeepBlue!64.18](-0.6977,2.1152)rectangle(-0.6027,2.1737);
\fill[DeepBlue!65.99](-0.6977,2.1717)rectangle(-0.6027,2.2303);
\fill[DeepBlue!65.99](-0.6977,2.2283)rectangle(-0.6027,2.2868);
\fill[DeepBlue!64.18](-0.6977,2.2848)rectangle(-0.6027,2.3433);
\fill[DeepBlue!60.64](-0.6977,2.3413)rectangle(-0.6027,2.3998);
\fill[DeepBlue!55.60](-0.6977,2.3978)rectangle(-0.6027,2.4563);
\fill[DeepBlue!49.34](-0.6977,2.4543)rectangle(-0.6027,2.5129);
\fill[DeepBlue!42.21](-0.6977,2.5109)rectangle(-0.6027,2.5694);
\fill[DeepBlue!34.62](-0.6977,2.5674)rectangle(-0.6027,2.6259);
\fill[DeepBlue!26.99](-0.6977,2.6239)rectangle(-0.6027,2.6824);
\fill[DeepBlue!19.76](-0.6977,2.6804)rectangle(-0.6027,2.7390);
\fill[DeepBlue!13.45](-0.6977,2.7370)rectangle(-0.6027,2.7955);
\fill[DeepBlue!8.75](-0.6977,2.7935)rectangle(-0.6027,2.8520);
\fill[DeepBlue!8.01](-0.6047,1.5500)rectangle(-0.5096,1.6085);
\fill[DeepBlue!10.20](-0.6047,1.6065)rectangle(-0.5096,1.6650);
\fill[DeepBlue!14.96](-0.6047,1.6630)rectangle(-0.5096,1.7216);
\fill[DeepBlue!21.40](-0.6047,1.7196)rectangle(-0.5096,1.7781);
\fill[DeepBlue!28.82](-0.6047,1.7761)rectangle(-0.5096,1.8346);
\fill[DeepBlue!36.70](-0.6047,1.8326)rectangle(-0.5096,1.8911);
\fill[DeepBlue!44.56](-0.6047,1.8891)rectangle(-0.5096,1.9477);
\fill[DeepBlue!51.95](-0.6047,1.9457)rectangle(-0.5096,2.0042);
\fill[DeepBlue!58.45](-0.6047,2.0022)rectangle(-0.5096,2.0607);
\fill[DeepBlue!63.68](-0.6047,2.0587)rectangle(-0.5096,2.1172);
\fill[DeepBlue!67.35](-0.6047,2.1152)rectangle(-0.5096,2.1737);
\fill[DeepBlue!69.24](-0.6047,2.1717)rectangle(-0.5096,2.2303);
\fill[DeepBlue!69.24](-0.6047,2.2283)rectangle(-0.5096,2.2868);
\fill[DeepBlue!67.35](-0.6047,2.2848)rectangle(-0.5096,2.3433);
\fill[DeepBlue!63.68](-0.6047,2.3413)rectangle(-0.5096,2.3998);
\fill[DeepBlue!58.45](-0.6047,2.3978)rectangle(-0.5096,2.4563);
\fill[DeepBlue!51.95](-0.6047,2.4543)rectangle(-0.5096,2.5129);
\fill[DeepBlue!44.56](-0.6047,2.5109)rectangle(-0.5096,2.5694);
\fill[DeepBlue!36.70](-0.6047,2.5674)rectangle(-0.5096,2.6259);
\fill[DeepBlue!28.82](-0.6047,2.6239)rectangle(-0.5096,2.6824);
\fill[DeepBlue!21.40](-0.6047,2.6804)rectangle(-0.5096,2.7390);
\fill[DeepBlue!14.96](-0.6047,2.7370)rectangle(-0.5096,2.7955);
\fill[DeepBlue!10.20](-0.6047,2.7935)rectangle(-0.5096,2.8520);
\fill[DeepBlue!8.92](-0.5116,1.5500)rectangle(-0.4166,1.6085);
\fill[DeepBlue!10.86](-0.5116,1.6065)rectangle(-0.4166,1.6650);
\fill[DeepBlue!15.39](-0.5116,1.6630)rectangle(-0.4166,1.7216);
\fill[DeepBlue!21.77](-0.5116,1.7196)rectangle(-0.4166,1.7781);
\fill[DeepBlue!29.24](-0.5116,1.7761)rectangle(-0.4166,1.8346);
\fill[DeepBlue!37.23](-0.5116,1.8326)rectangle(-0.4166,1.8911);
\fill[DeepBlue!45.23](-0.5116,1.8891)rectangle(-0.4166,1.9477);
\fill[DeepBlue!52.77](-0.5116,1.9457)rectangle(-0.4166,2.0042);
\fill[DeepBlue!59.41](-0.5116,2.0022)rectangle(-0.4166,2.0607);
\fill[DeepBlue!64.76](-0.5116,2.0587)rectangle(-0.4166,2.1172);
\fill[DeepBlue!68.52](-0.5116,2.1152)rectangle(-0.4166,2.1737);
\fill[DeepBlue!70.45](-0.5116,2.1717)rectangle(-0.4166,2.2303);
\fill[DeepBlue!70.45](-0.5116,2.2283)rectangle(-0.4166,2.2868);
\fill[DeepBlue!68.52](-0.5116,2.2848)rectangle(-0.4166,2.3433);
\fill[DeepBlue!64.76](-0.5116,2.3413)rectangle(-0.4166,2.3998);
\fill[DeepBlue!59.41](-0.5116,2.3978)rectangle(-0.4166,2.4563);
\fill[DeepBlue!52.77](-0.5116,2.4543)rectangle(-0.4166,2.5129);
\fill[DeepBlue!45.23](-0.5116,2.5109)rectangle(-0.4166,2.5694);
\fill[DeepBlue!37.23](-0.5116,2.5674)rectangle(-0.4166,2.6259);
\fill[DeepBlue!29.24](-0.5116,2.6239)rectangle(-0.4166,2.6824);
\fill[DeepBlue!21.77](-0.5116,2.6804)rectangle(-0.4166,2.7390);
\fill[DeepBlue!15.39](-0.5116,2.7370)rectangle(-0.4166,2.7955);
\fill[DeepBlue!10.86](-0.5116,2.7935)rectangle(-0.4166,2.8520);
\fill[DeepBlue!9.33](-0.4186,1.5500)rectangle(-0.3236,1.6085);
\fill[DeepBlue!10.64](-0.4186,1.6065)rectangle(-0.3236,1.6650);
\fill[DeepBlue!14.70](-0.4186,1.6630)rectangle(-0.3236,1.7216);
\fill[DeepBlue!20.86](-0.4186,1.7196)rectangle(-0.3236,1.7781);
\fill[DeepBlue!28.26](-0.4186,1.7761)rectangle(-0.3236,1.8346);
\fill[DeepBlue!36.26](-0.4186,1.8326)rectangle(-0.3236,1.8911);
\fill[DeepBlue!44.30](-0.4186,1.8891)rectangle(-0.3236,1.9477);
\fill[DeepBlue!51.91](-0.4186,1.9457)rectangle(-0.3236,2.0042);
\fill[DeepBlue!58.62](-0.4186,2.0022)rectangle(-0.3236,2.0607);
\fill[DeepBlue!64.04](-0.4186,2.0587)rectangle(-0.3236,2.1172);
\fill[DeepBlue!67.84](-0.4186,2.1152)rectangle(-0.3236,2.1737);
\fill[DeepBlue!69.80](-0.4186,2.1717)rectangle(-0.3236,2.2303);
\fill[DeepBlue!69.80](-0.4186,2.2283)rectangle(-0.3236,2.2868);
\fill[DeepBlue!67.84](-0.4186,2.2848)rectangle(-0.3236,2.3433);
\fill[DeepBlue!64.04](-0.4186,2.3413)rectangle(-0.3236,2.3998);
\fill[DeepBlue!58.62](-0.4186,2.3978)rectangle(-0.3236,2.4563);
\fill[DeepBlue!51.91](-0.4186,2.4543)rectangle(-0.3236,2.5129);
\fill[DeepBlue!44.30](-0.4186,2.5109)rectangle(-0.3236,2.5694);
\fill[DeepBlue!36.26](-0.4186,2.5674)rectangle(-0.3236,2.6259);
\fill[DeepBlue!28.26](-0.4186,2.6239)rectangle(-0.3236,2.6824);
\fill[DeepBlue!20.86](-0.4186,2.6804)rectangle(-0.3236,2.7390);
\fill[DeepBlue!14.70](-0.4186,2.7370)rectangle(-0.3236,2.7955);
\fill[DeepBlue!10.64](-0.4186,2.7935)rectangle(-0.3236,2.8520);
\fill[DeepBlue!9.26](-0.3256,1.5500)rectangle(-0.2306,1.6085);
\fill[DeepBlue!9.56](-0.3256,1.6065)rectangle(-0.2306,1.6650);
\fill[DeepBlue!12.92](-0.3256,1.6630)rectangle(-0.2306,1.7216);
\fill[DeepBlue!18.79](-0.3256,1.7196)rectangle(-0.2306,1.7781);
\fill[DeepBlue!26.08](-0.3256,1.7761)rectangle(-0.2306,1.8346);
\fill[DeepBlue!34.05](-0.3256,1.8326)rectangle(-0.2306,1.8911);
\fill[DeepBlue!42.10](-0.3256,1.8891)rectangle(-0.2306,1.9477);
\fill[DeepBlue!49.73](-0.3256,1.9457)rectangle(-0.2306,2.0042);
\fill[DeepBlue!56.47](-0.3256,2.0022)rectangle(-0.2306,2.0607);
\fill[DeepBlue!61.92](-0.3256,2.0587)rectangle(-0.2306,2.1172);
\fill[DeepBlue!65.74](-0.3256,2.1152)rectangle(-0.2306,2.1737);
\fill[DeepBlue!67.72](-0.3256,2.1717)rectangle(-0.2306,2.2303);
\fill[DeepBlue!67.72](-0.3256,2.2283)rectangle(-0.2306,2.2868);
\fill[DeepBlue!65.74](-0.3256,2.2848)rectangle(-0.2306,2.3433);
\fill[DeepBlue!61.92](-0.3256,2.3413)rectangle(-0.2306,2.3998);
\fill[DeepBlue!56.47](-0.3256,2.3978)rectangle(-0.2306,2.4563);
\fill[DeepBlue!49.73](-0.3256,2.4543)rectangle(-0.2306,2.5129);
\fill[DeepBlue!42.10](-0.3256,2.5109)rectangle(-0.2306,2.5694);
\fill[DeepBlue!34.05](-0.3256,2.5674)rectangle(-0.2306,2.6259);
\fill[DeepBlue!26.08](-0.3256,2.6239)rectangle(-0.2306,2.6824);
\fill[DeepBlue!18.79](-0.3256,2.6804)rectangle(-0.2306,2.7390);
\fill[DeepBlue!12.92](-0.3256,2.7370)rectangle(-0.2306,2.7955);
\fill[DeepBlue!9.56](-0.3256,2.7935)rectangle(-0.2306,2.8520);
\fill[DeepBlue!8.82](-0.2326,1.5500)rectangle(-0.1375,1.6085);
\fill[DeepBlue!7.75](-0.2326,1.6065)rectangle(-0.1375,1.6650);
\fill[DeepBlue!10.20](-0.2326,1.6630)rectangle(-0.1375,1.7216);
\fill[DeepBlue!15.86](-0.2326,1.7196)rectangle(-0.1375,1.7781);
\fill[DeepBlue!23.14](-0.2326,1.7761)rectangle(-0.1375,1.8346);
\fill[DeepBlue!31.12](-0.2326,1.8326)rectangle(-0.1375,1.8911);
\fill[DeepBlue!39.19](-0.2326,1.8891)rectangle(-0.1375,1.9477);
\fill[DeepBlue!46.84](-0.2326,1.9457)rectangle(-0.1375,2.0042);
\fill[DeepBlue!53.61](-0.2326,2.0022)rectangle(-0.1375,2.0607);
\fill[DeepBlue!59.08](-0.2326,2.0587)rectangle(-0.1375,2.1172);
\fill[DeepBlue!62.91](-0.2326,2.1152)rectangle(-0.1375,2.1737);
\fill[DeepBlue!64.89](-0.2326,2.1717)rectangle(-0.1375,2.2303);
\fill[DeepBlue!64.89](-0.2326,2.2283)rectangle(-0.1375,2.2868);
\fill[DeepBlue!62.91](-0.2326,2.2848)rectangle(-0.1375,2.3433);
\fill[DeepBlue!59.08](-0.2326,2.3413)rectangle(-0.1375,2.3998);
\fill[DeepBlue!53.61](-0.2326,2.3978)rectangle(-0.1375,2.4563);
\fill[DeepBlue!46.84](-0.2326,2.4543)rectangle(-0.1375,2.5129);
\fill[DeepBlue!39.19](-0.2326,2.5109)rectangle(-0.1375,2.5694);
\fill[DeepBlue!31.12](-0.2326,2.5674)rectangle(-0.1375,2.6259);
\fill[DeepBlue!23.14](-0.2326,2.6239)rectangle(-0.1375,2.6824);
\fill[DeepBlue!15.86](-0.2326,2.6804)rectangle(-0.1375,2.7390);
\fill[DeepBlue!10.20](-0.2326,2.7370)rectangle(-0.1375,2.7955);
\fill[DeepBlue!7.75](-0.2326,2.7935)rectangle(-0.1375,2.8520);
\fill[DeepBlue!8.24](-0.1395,1.5500)rectangle(-0.0445,1.6085);
\fill[DeepBlue!5.49](-0.1395,1.6065)rectangle(-0.0445,1.6650);
\fill[DeepBlue!6.85](-0.1395,1.6630)rectangle(-0.0445,1.7216);
\fill[DeepBlue!12.68](-0.1395,1.7196)rectangle(-0.0445,1.7781);
\fill[DeepBlue!20.15](-0.1395,1.7761)rectangle(-0.0445,1.8346);
\fill[DeepBlue!28.25](-0.1395,1.8326)rectangle(-0.0445,1.8911);
\fill[DeepBlue!36.40](-0.1395,1.8891)rectangle(-0.0445,1.9477);
\fill[DeepBlue!44.10](-0.1395,1.9457)rectangle(-0.0445,2.0042);
\fill[DeepBlue!50.89](-0.1395,2.0022)rectangle(-0.0445,2.0607);
\fill[DeepBlue!56.38](-0.1395,2.0587)rectangle(-0.0445,2.1172);
\fill[DeepBlue!60.23](-0.1395,2.1152)rectangle(-0.0445,2.1737);
\fill[DeepBlue!62.22](-0.1395,2.1717)rectangle(-0.0445,2.2303);
\fill[DeepBlue!62.22](-0.1395,2.2283)rectangle(-0.0445,2.2868);
\fill[DeepBlue!60.23](-0.1395,2.2848)rectangle(-0.0445,2.3433);
\fill[DeepBlue!56.38](-0.1395,2.3413)rectangle(-0.0445,2.3998);
\fill[DeepBlue!50.89](-0.1395,2.3978)rectangle(-0.0445,2.4563);
\fill[DeepBlue!44.10](-0.1395,2.4543)rectangle(-0.0445,2.5129);
\fill[DeepBlue!36.40](-0.1395,2.5109)rectangle(-0.0445,2.5694);
\fill[DeepBlue!28.25](-0.1395,2.5674)rectangle(-0.0445,2.6259);
\fill[DeepBlue!20.15](-0.1395,2.6239)rectangle(-0.0445,2.6824);
\fill[DeepBlue!12.68](-0.1395,2.6804)rectangle(-0.0445,2.7390);
\fill[DeepBlue!6.85](-0.1395,2.7370)rectangle(-0.0445,2.7955);
\fill[DeepBlue!5.49](-0.1395,2.7935)rectangle(-0.0445,2.8520);
\fill[DeepBlue!7.84](-0.0465,1.5500)rectangle(0.0485,1.6085);
\fill[DeepBlue!3.47](-0.0465,1.6065)rectangle(0.0485,1.6650);
\fill[DeepBlue!3.60](-0.0465,1.6630)rectangle(0.0485,1.7216);
\fill[DeepBlue!10.36](-0.0465,1.7196)rectangle(0.0485,1.7781);
\fill[DeepBlue!18.17](-0.0465,1.7761)rectangle(0.0485,1.8346);
\fill[DeepBlue!26.42](-0.0465,1.8326)rectangle(0.0485,1.8911);
\fill[DeepBlue!34.65](-0.0465,1.8891)rectangle(0.0485,1.9477);
\fill[DeepBlue!42.40](-0.0465,1.9457)rectangle(0.0485,2.0042);
\fill[DeepBlue!49.23](-0.0465,2.0022)rectangle(0.0485,2.0607);
\fill[DeepBlue!54.73](-0.0465,2.0587)rectangle(0.0485,2.1172);
\fill[DeepBlue!58.60](-0.0465,2.1152)rectangle(0.0485,2.1737);
\fill[DeepBlue!60.59](-0.0465,2.1717)rectangle(0.0485,2.2303);
\fill[DeepBlue!60.59](-0.0465,2.2283)rectangle(0.0485,2.2868);
\fill[DeepBlue!58.60](-0.0465,2.2848)rectangle(0.0485,2.3433);
\fill[DeepBlue!54.73](-0.0465,2.3413)rectangle(0.0485,2.3998);
\fill[DeepBlue!49.23](-0.0465,2.3978)rectangle(0.0485,2.4563);
\fill[DeepBlue!42.40](-0.0465,2.4543)rectangle(0.0485,2.5129);
\fill[DeepBlue!34.65](-0.0465,2.5109)rectangle(0.0485,2.5694);
\fill[DeepBlue!26.42](-0.0465,2.5674)rectangle(0.0485,2.6259);
\fill[DeepBlue!18.17](-0.0465,2.6239)rectangle(0.0485,2.6824);
\fill[DeepBlue!10.36](-0.0465,2.6804)rectangle(0.0485,2.7390);
\fill[DeepBlue!3.60](-0.0465,2.7370)rectangle(0.0485,2.7955);
\fill[DeepBlue!3.47](-0.0465,2.7935)rectangle(0.0485,2.8520);
\fill[DeepBlue!7.84](0.0465,1.5500)rectangle(0.1415,1.6085);
\fill[DeepBlue!3.47](0.0465,1.6065)rectangle(0.1415,1.6650);
\fill[DeepBlue!3.60](0.0465,1.6630)rectangle(0.1415,1.7216);
\fill[DeepBlue!10.36](0.0465,1.7196)rectangle(0.1415,1.7781);
\fill[DeepBlue!18.17](0.0465,1.7761)rectangle(0.1415,1.8346);
\fill[DeepBlue!26.42](0.0465,1.8326)rectangle(0.1415,1.8911);
\fill[DeepBlue!34.65](0.0465,1.8891)rectangle(0.1415,1.9477);
\fill[DeepBlue!42.40](0.0465,1.9457)rectangle(0.1415,2.0042);
\fill[DeepBlue!49.23](0.0465,2.0022)rectangle(0.1415,2.0607);
\fill[DeepBlue!54.73](0.0465,2.0587)rectangle(0.1415,2.1172);
\fill[DeepBlue!58.60](0.0465,2.1152)rectangle(0.1415,2.1737);
\fill[DeepBlue!60.59](0.0465,2.1717)rectangle(0.1415,2.2303);
\fill[DeepBlue!60.59](0.0465,2.2283)rectangle(0.1415,2.2868);
\fill[DeepBlue!58.60](0.0465,2.2848)rectangle(0.1415,2.3433);
\fill[DeepBlue!54.73](0.0465,2.3413)rectangle(0.1415,2.3998);
\fill[DeepBlue!49.23](0.0465,2.3978)rectangle(0.1415,2.4563);
\fill[DeepBlue!42.40](0.0465,2.4543)rectangle(0.1415,2.5129);
\fill[DeepBlue!34.65](0.0465,2.5109)rectangle(0.1415,2.5694);
\fill[DeepBlue!26.42](0.0465,2.5674)rectangle(0.1415,2.6259);
\fill[DeepBlue!18.17](0.0465,2.6239)rectangle(0.1415,2.6824);
\fill[DeepBlue!10.36](0.0465,2.6804)rectangle(0.1415,2.7390);
\fill[DeepBlue!3.60](0.0465,2.7370)rectangle(0.1415,2.7955);
\fill[DeepBlue!3.47](0.0465,2.7935)rectangle(0.1415,2.8520);
\fill[DeepBlue!8.24](0.1395,1.5500)rectangle(0.2346,1.6085);
\fill[DeepBlue!5.49](0.1395,1.6065)rectangle(0.2346,1.6650);
\fill[DeepBlue!6.85](0.1395,1.6630)rectangle(0.2346,1.7216);
\fill[DeepBlue!12.68](0.1395,1.7196)rectangle(0.2346,1.7781);
\fill[DeepBlue!20.15](0.1395,1.7761)rectangle(0.2346,1.8346);
\fill[DeepBlue!28.25](0.1395,1.8326)rectangle(0.2346,1.8911);
\fill[DeepBlue!36.40](0.1395,1.8891)rectangle(0.2346,1.9477);
\fill[DeepBlue!44.10](0.1395,1.9457)rectangle(0.2346,2.0042);
\fill[DeepBlue!50.89](0.1395,2.0022)rectangle(0.2346,2.0607);
\fill[DeepBlue!56.38](0.1395,2.0587)rectangle(0.2346,2.1172);
\fill[DeepBlue!60.23](0.1395,2.1152)rectangle(0.2346,2.1737);
\fill[DeepBlue!62.22](0.1395,2.1717)rectangle(0.2346,2.2303);
\fill[DeepBlue!62.22](0.1395,2.2283)rectangle(0.2346,2.2868);
\fill[DeepBlue!60.23](0.1395,2.2848)rectangle(0.2346,2.3433);
\fill[DeepBlue!56.38](0.1395,2.3413)rectangle(0.2346,2.3998);
\fill[DeepBlue!50.89](0.1395,2.3978)rectangle(0.2346,2.4563);
\fill[DeepBlue!44.10](0.1395,2.4543)rectangle(0.2346,2.5129);
\fill[DeepBlue!36.40](0.1395,2.5109)rectangle(0.2346,2.5694);
\fill[DeepBlue!28.25](0.1395,2.5674)rectangle(0.2346,2.6259);
\fill[DeepBlue!20.15](0.1395,2.6239)rectangle(0.2346,2.6824);
\fill[DeepBlue!12.68](0.1395,2.6804)rectangle(0.2346,2.7390);
\fill[DeepBlue!6.85](0.1395,2.7370)rectangle(0.2346,2.7955);
\fill[DeepBlue!5.49](0.1395,2.7935)rectangle(0.2346,2.8520);
\fill[DeepBlue!8.82](0.2326,1.5500)rectangle(0.3276,1.6085);
\fill[DeepBlue!7.75](0.2326,1.6065)rectangle(0.3276,1.6650);
\fill[DeepBlue!10.20](0.2326,1.6630)rectangle(0.3276,1.7216);
\fill[DeepBlue!15.86](0.2326,1.7196)rectangle(0.3276,1.7781);
\fill[DeepBlue!23.14](0.2326,1.7761)rectangle(0.3276,1.8346);
\fill[DeepBlue!31.12](0.2326,1.8326)rectangle(0.3276,1.8911);
\fill[DeepBlue!39.19](0.2326,1.8891)rectangle(0.3276,1.9477);
\fill[DeepBlue!46.84](0.2326,1.9457)rectangle(0.3276,2.0042);
\fill[DeepBlue!53.61](0.2326,2.0022)rectangle(0.3276,2.0607);
\fill[DeepBlue!59.08](0.2326,2.0587)rectangle(0.3276,2.1172);
\fill[DeepBlue!62.91](0.2326,2.1152)rectangle(0.3276,2.1737);
\fill[DeepBlue!64.89](0.2326,2.1717)rectangle(0.3276,2.2303);
\fill[DeepBlue!64.89](0.2326,2.2283)rectangle(0.3276,2.2868);
\fill[DeepBlue!62.91](0.2326,2.2848)rectangle(0.3276,2.3433);
\fill[DeepBlue!59.08](0.2326,2.3413)rectangle(0.3276,2.3998);
\fill[DeepBlue!53.61](0.2326,2.3978)rectangle(0.3276,2.4563);
\fill[DeepBlue!46.84](0.2326,2.4543)rectangle(0.3276,2.5129);
\fill[DeepBlue!39.19](0.2326,2.5109)rectangle(0.3276,2.5694);
\fill[DeepBlue!31.12](0.2326,2.5674)rectangle(0.3276,2.6259);
\fill[DeepBlue!23.14](0.2326,2.6239)rectangle(0.3276,2.6824);
\fill[DeepBlue!15.86](0.2326,2.6804)rectangle(0.3276,2.7390);
\fill[DeepBlue!10.20](0.2326,2.7370)rectangle(0.3276,2.7955);
\fill[DeepBlue!7.75](0.2326,2.7935)rectangle(0.3276,2.8520);
\fill[DeepBlue!9.26](0.3256,1.5500)rectangle(0.4206,1.6085);
\fill[DeepBlue!9.56](0.3256,1.6065)rectangle(0.4206,1.6650);
\fill[DeepBlue!12.92](0.3256,1.6630)rectangle(0.4206,1.7216);
\fill[DeepBlue!18.79](0.3256,1.7196)rectangle(0.4206,1.7781);
\fill[DeepBlue!26.08](0.3256,1.7761)rectangle(0.4206,1.8346);
\fill[DeepBlue!34.05](0.3256,1.8326)rectangle(0.4206,1.8911);
\fill[DeepBlue!42.10](0.3256,1.8891)rectangle(0.4206,1.9477);
\fill[DeepBlue!49.73](0.3256,1.9457)rectangle(0.4206,2.0042);
\fill[DeepBlue!56.47](0.3256,2.0022)rectangle(0.4206,2.0607);
\fill[DeepBlue!61.92](0.3256,2.0587)rectangle(0.4206,2.1172);
\fill[DeepBlue!65.74](0.3256,2.1152)rectangle(0.4206,2.1737);
\fill[DeepBlue!67.72](0.3256,2.1717)rectangle(0.4206,2.2303);
\fill[DeepBlue!67.72](0.3256,2.2283)rectangle(0.4206,2.2868);
\fill[DeepBlue!65.74](0.3256,2.2848)rectangle(0.4206,2.3433);
\fill[DeepBlue!61.92](0.3256,2.3413)rectangle(0.4206,2.3998);
\fill[DeepBlue!56.47](0.3256,2.3978)rectangle(0.4206,2.4563);
\fill[DeepBlue!49.73](0.3256,2.4543)rectangle(0.4206,2.5129);
\fill[DeepBlue!42.10](0.3256,2.5109)rectangle(0.4206,2.5694);
\fill[DeepBlue!34.05](0.3256,2.5674)rectangle(0.4206,2.6259);
\fill[DeepBlue!26.08](0.3256,2.6239)rectangle(0.4206,2.6824);
\fill[DeepBlue!18.79](0.3256,2.6804)rectangle(0.4206,2.7390);
\fill[DeepBlue!12.92](0.3256,2.7370)rectangle(0.4206,2.7955);
\fill[DeepBlue!9.56](0.3256,2.7935)rectangle(0.4206,2.8520);
\fill[DeepBlue!9.33](0.4186,1.5500)rectangle(0.5136,1.6085);
\fill[DeepBlue!10.64](0.4186,1.6065)rectangle(0.5136,1.6650);
\fill[DeepBlue!14.70](0.4186,1.6630)rectangle(0.5136,1.7216);
\fill[DeepBlue!20.86](0.4186,1.7196)rectangle(0.5136,1.7781);
\fill[DeepBlue!28.26](0.4186,1.7761)rectangle(0.5136,1.8346);
\fill[DeepBlue!36.26](0.4186,1.8326)rectangle(0.5136,1.8911);
\fill[DeepBlue!44.30](0.4186,1.8891)rectangle(0.5136,1.9477);
\fill[DeepBlue!51.91](0.4186,1.9457)rectangle(0.5136,2.0042);
\fill[DeepBlue!58.62](0.4186,2.0022)rectangle(0.5136,2.0607);
\fill[DeepBlue!64.04](0.4186,2.0587)rectangle(0.5136,2.1172);
\fill[DeepBlue!67.84](0.4186,2.1152)rectangle(0.5136,2.1737);
\fill[DeepBlue!69.80](0.4186,2.1717)rectangle(0.5136,2.2303);
\fill[DeepBlue!69.80](0.4186,2.2283)rectangle(0.5136,2.2868);
\fill[DeepBlue!67.84](0.4186,2.2848)rectangle(0.5136,2.3433);
\fill[DeepBlue!64.04](0.4186,2.3413)rectangle(0.5136,2.3998);
\fill[DeepBlue!58.62](0.4186,2.3978)rectangle(0.5136,2.4563);
\fill[DeepBlue!51.91](0.4186,2.4543)rectangle(0.5136,2.5129);
\fill[DeepBlue!44.30](0.4186,2.5109)rectangle(0.5136,2.5694);
\fill[DeepBlue!36.26](0.4186,2.5674)rectangle(0.5136,2.6259);
\fill[DeepBlue!28.26](0.4186,2.6239)rectangle(0.5136,2.6824);
\fill[DeepBlue!20.86](0.4186,2.6804)rectangle(0.5136,2.7390);
\fill[DeepBlue!14.70](0.4186,2.7370)rectangle(0.5136,2.7955);
\fill[DeepBlue!10.64](0.4186,2.7935)rectangle(0.5136,2.8520);
\fill[DeepBlue!8.92](0.5116,1.5500)rectangle(0.6067,1.6085);
\fill[DeepBlue!10.86](0.5116,1.6065)rectangle(0.6067,1.6650);
\fill[DeepBlue!15.39](0.5116,1.6630)rectangle(0.6067,1.7216);
\fill[DeepBlue!21.77](0.5116,1.7196)rectangle(0.6067,1.7781);
\fill[DeepBlue!29.24](0.5116,1.7761)rectangle(0.6067,1.8346);
\fill[DeepBlue!37.23](0.5116,1.8326)rectangle(0.6067,1.8911);
\fill[DeepBlue!45.23](0.5116,1.8891)rectangle(0.6067,1.9477);
\fill[DeepBlue!52.77](0.5116,1.9457)rectangle(0.6067,2.0042);
\fill[DeepBlue!59.41](0.5116,2.0022)rectangle(0.6067,2.0607);
\fill[DeepBlue!64.76](0.5116,2.0587)rectangle(0.6067,2.1172);
\fill[DeepBlue!68.52](0.5116,2.1152)rectangle(0.6067,2.1737);
\fill[DeepBlue!70.45](0.5116,2.1717)rectangle(0.6067,2.2303);
\fill[DeepBlue!70.45](0.5116,2.2283)rectangle(0.6067,2.2868);
\fill[DeepBlue!68.52](0.5116,2.2848)rectangle(0.6067,2.3433);
\fill[DeepBlue!64.76](0.5116,2.3413)rectangle(0.6067,2.3998);
\fill[DeepBlue!59.41](0.5116,2.3978)rectangle(0.6067,2.4563);
\fill[DeepBlue!52.77](0.5116,2.4543)rectangle(0.6067,2.5129);
\fill[DeepBlue!45.23](0.5116,2.5109)rectangle(0.6067,2.5694);
\fill[DeepBlue!37.23](0.5116,2.5674)rectangle(0.6067,2.6259);
\fill[DeepBlue!29.24](0.5116,2.6239)rectangle(0.6067,2.6824);
\fill[DeepBlue!21.77](0.5116,2.6804)rectangle(0.6067,2.7390);
\fill[DeepBlue!15.39](0.5116,2.7370)rectangle(0.6067,2.7955);
\fill[DeepBlue!10.86](0.5116,2.7935)rectangle(0.6067,2.8520);
\fill[DeepBlue!8.01](0.6047,1.5500)rectangle(0.6997,1.6085);
\fill[DeepBlue!10.20](0.6047,1.6065)rectangle(0.6997,1.6650);
\fill[DeepBlue!14.96](0.6047,1.6630)rectangle(0.6997,1.7216);
\fill[DeepBlue!21.40](0.6047,1.7196)rectangle(0.6997,1.7781);
\fill[DeepBlue!28.82](0.6047,1.7761)rectangle(0.6997,1.8346);
\fill[DeepBlue!36.70](0.6047,1.8326)rectangle(0.6997,1.8911);
\fill[DeepBlue!44.56](0.6047,1.8891)rectangle(0.6997,1.9477);
\fill[DeepBlue!51.95](0.6047,1.9457)rectangle(0.6997,2.0042);
\fill[DeepBlue!58.45](0.6047,2.0022)rectangle(0.6997,2.0607);
\fill[DeepBlue!63.68](0.6047,2.0587)rectangle(0.6997,2.1172);
\fill[DeepBlue!67.35](0.6047,2.1152)rectangle(0.6997,2.1737);
\fill[DeepBlue!69.24](0.6047,2.1717)rectangle(0.6997,2.2303);
\fill[DeepBlue!69.24](0.6047,2.2283)rectangle(0.6997,2.2868);
\fill[DeepBlue!67.35](0.6047,2.2848)rectangle(0.6997,2.3433);
\fill[DeepBlue!63.68](0.6047,2.3413)rectangle(0.6997,2.3998);
\fill[DeepBlue!58.45](0.6047,2.3978)rectangle(0.6997,2.4563);
\fill[DeepBlue!51.95](0.6047,2.4543)rectangle(0.6997,2.5129);
\fill[DeepBlue!44.56](0.6047,2.5109)rectangle(0.6997,2.5694);
\fill[DeepBlue!36.70](0.6047,2.5674)rectangle(0.6997,2.6259);
\fill[DeepBlue!28.82](0.6047,2.6239)rectangle(0.6997,2.6824);
\fill[DeepBlue!21.40](0.6047,2.6804)rectangle(0.6997,2.7390);
\fill[DeepBlue!14.96](0.6047,2.7370)rectangle(0.6997,2.7955);
\fill[DeepBlue!10.20](0.6047,2.7935)rectangle(0.6997,2.8520);
\fill[DeepBlue!6.73](0.6977,1.5500)rectangle(0.7927,1.6085);
\fill[DeepBlue!8.75](0.6977,1.6065)rectangle(0.7927,1.6650);
\fill[DeepBlue!13.45](0.6977,1.6630)rectangle(0.7927,1.7216);
\fill[DeepBlue!19.76](0.6977,1.7196)rectangle(0.7927,1.7781);
\fill[DeepBlue!26.99](0.6977,1.7761)rectangle(0.7927,1.8346);
\fill[DeepBlue!34.62](0.6977,1.8326)rectangle(0.7927,1.8911);
\fill[DeepBlue!42.21](0.6977,1.8891)rectangle(0.7927,1.9477);
\fill[DeepBlue!49.34](0.6977,1.9457)rectangle(0.7927,2.0042);
\fill[DeepBlue!55.60](0.6977,2.0022)rectangle(0.7927,2.0607);
\fill[DeepBlue!60.64](0.6977,2.0587)rectangle(0.7927,2.1172);
\fill[DeepBlue!64.18](0.6977,2.1152)rectangle(0.7927,2.1737);
\fill[DeepBlue!65.99](0.6977,2.1717)rectangle(0.7927,2.2303);
\fill[DeepBlue!65.99](0.6977,2.2283)rectangle(0.7927,2.2868);
\fill[DeepBlue!64.18](0.6977,2.2848)rectangle(0.7927,2.3433);
\fill[DeepBlue!60.64](0.6977,2.3413)rectangle(0.7927,2.3998);
\fill[DeepBlue!55.60](0.6977,2.3978)rectangle(0.7927,2.4563);
\fill[DeepBlue!49.34](0.6977,2.4543)rectangle(0.7927,2.5129);
\fill[DeepBlue!42.21](0.6977,2.5109)rectangle(0.7927,2.5694);
\fill[DeepBlue!34.62](0.6977,2.5674)rectangle(0.7927,2.6259);
\fill[DeepBlue!26.99](0.6977,2.6239)rectangle(0.7927,2.6824);
\fill[DeepBlue!19.76](0.6977,2.6804)rectangle(0.7927,2.7390);
\fill[DeepBlue!13.45](0.6977,2.7370)rectangle(0.7927,2.7955);
\fill[DeepBlue!8.75](0.6977,2.7935)rectangle(0.7927,2.8520);
\fill[DeepBlue!5.42](0.7907,1.5500)rectangle(0.8857,1.6085);
\fill[DeepBlue!6.68](0.7907,1.6065)rectangle(0.8857,1.6650);
\fill[DeepBlue!11.03](0.7907,1.6630)rectangle(0.8857,1.7216);
\fill[DeepBlue!17.02](0.7907,1.7196)rectangle(0.8857,1.7781);
\fill[DeepBlue!23.87](0.7907,1.7761)rectangle(0.8857,1.8346);
\fill[DeepBlue!31.10](0.7907,1.8326)rectangle(0.8857,1.8911);
\fill[DeepBlue!38.28](0.7907,1.8891)rectangle(0.8857,1.9477);
\fill[DeepBlue!45.03](0.7907,1.9457)rectangle(0.8857,2.0042);
\fill[DeepBlue!50.95](0.7907,2.0022)rectangle(0.8857,2.0607);
\fill[DeepBlue!55.71](0.7907,2.0587)rectangle(0.8857,2.1172);
\fill[DeepBlue!59.05](0.7907,2.1152)rectangle(0.8857,2.1737);
\fill[DeepBlue!60.77](0.7907,2.1717)rectangle(0.8857,2.2303);
\fill[DeepBlue!60.77](0.7907,2.2283)rectangle(0.8857,2.2868);
\fill[DeepBlue!59.05](0.7907,2.2848)rectangle(0.8857,2.3433);
\fill[DeepBlue!55.71](0.7907,2.3413)rectangle(0.8857,2.3998);
\fill[DeepBlue!50.95](0.7907,2.3978)rectangle(0.8857,2.4563);
\fill[DeepBlue!45.03](0.7907,2.4543)rectangle(0.8857,2.5129);
\fill[DeepBlue!38.28](0.7907,2.5109)rectangle(0.8857,2.5694);
\fill[DeepBlue!31.10](0.7907,2.5674)rectangle(0.8857,2.6259);
\fill[DeepBlue!23.87](0.7907,2.6239)rectangle(0.8857,2.6824);
\fill[DeepBlue!17.02](0.7907,2.6804)rectangle(0.8857,2.7390);
\fill[DeepBlue!11.03](0.7907,2.7370)rectangle(0.8857,2.7955);
\fill[DeepBlue!6.68](0.7907,2.7935)rectangle(0.8857,2.8520);
\fill[DeepBlue!4.74](0.8837,1.5500)rectangle(0.9787,1.6085);
\fill[DeepBlue!4.36](0.8837,1.6065)rectangle(0.9787,1.6650);
\fill[DeepBlue!7.92](0.8837,1.6630)rectangle(0.9787,1.7216);
\fill[DeepBlue!13.40](0.8837,1.7196)rectangle(0.9787,1.7781);
\fill[DeepBlue!19.71](0.8837,1.7761)rectangle(0.9787,1.8346);
\fill[DeepBlue!26.39](0.8837,1.8326)rectangle(0.9787,1.8911);
\fill[DeepBlue!33.02](0.8837,1.8891)rectangle(0.9787,1.9477);
\fill[DeepBlue!39.26](0.8837,1.9457)rectangle(0.9787,2.0042);
\fill[DeepBlue!44.74](0.8837,2.0022)rectangle(0.9787,2.0607);
\fill[DeepBlue!49.15](0.8837,2.0587)rectangle(0.9787,2.1172);
\fill[DeepBlue!52.25](0.8837,2.1152)rectangle(0.9787,2.1737);
\fill[DeepBlue!53.84](0.8837,2.1717)rectangle(0.9787,2.2303);
\fill[DeepBlue!53.84](0.8837,2.2283)rectangle(0.9787,2.2868);
\fill[DeepBlue!52.25](0.8837,2.2848)rectangle(0.9787,2.3433);
\fill[DeepBlue!49.15](0.8837,2.3413)rectangle(0.9787,2.3998);
\fill[DeepBlue!44.74](0.8837,2.3978)rectangle(0.9787,2.4563);
\fill[DeepBlue!39.26](0.8837,2.4543)rectangle(0.9787,2.5129);
\fill[DeepBlue!33.02](0.8837,2.5109)rectangle(0.9787,2.5694);
\fill[DeepBlue!26.39](0.8837,2.5674)rectangle(0.9787,2.6259);
\fill[DeepBlue!19.71](0.8837,2.6239)rectangle(0.9787,2.6824);
\fill[DeepBlue!13.40](0.8837,2.6804)rectangle(0.9787,2.7390);
\fill[DeepBlue!7.92](0.8837,2.7370)rectangle(0.9787,2.7955);
\fill[DeepBlue!4.36](0.8837,2.7935)rectangle(0.9787,2.8520);
\fill[DeepBlue!5.25](0.9767,1.5500)rectangle(1.0718,1.6085);
\fill[DeepBlue!2.84](0.9767,1.6065)rectangle(1.0718,1.6650);
\fill[DeepBlue!4.46](0.9767,1.6630)rectangle(1.0718,1.7216);
\fill[DeepBlue!9.21](0.9767,1.7196)rectangle(1.0718,1.7781);
\fill[DeepBlue!14.85](0.9767,1.7761)rectangle(1.0718,1.8346);
\fill[DeepBlue!20.83](0.9767,1.8326)rectangle(1.0718,1.8911);
\fill[DeepBlue!26.80](0.9767,1.8891)rectangle(1.0718,1.9477);
\fill[DeepBlue!32.42](0.9767,1.9457)rectangle(1.0718,2.0042);
\fill[DeepBlue!37.37](0.9767,2.0022)rectangle(1.0718,2.0607);
\fill[DeepBlue!41.36](0.9767,2.0587)rectangle(1.0718,2.1172);
\fill[DeepBlue!44.16](0.9767,2.1152)rectangle(1.0718,2.1737);
\fill[DeepBlue!45.60](0.9767,2.1717)rectangle(1.0718,2.2303);
\fill[DeepBlue!45.60](0.9767,2.2283)rectangle(1.0718,2.2868);
\fill[DeepBlue!44.16](0.9767,2.2848)rectangle(1.0718,2.3433);
\fill[DeepBlue!41.36](0.9767,2.3413)rectangle(1.0718,2.3998);
\fill[DeepBlue!37.37](0.9767,2.3978)rectangle(1.0718,2.4563);
\fill[DeepBlue!32.42](0.9767,2.4543)rectangle(1.0718,2.5129);
\fill[DeepBlue!26.80](0.9767,2.5109)rectangle(1.0718,2.5694);
\fill[DeepBlue!20.83](0.9767,2.5674)rectangle(1.0718,2.6259);
\fill[DeepBlue!14.85](0.9767,2.6239)rectangle(1.0718,2.6824);
\fill[DeepBlue!9.21](0.9767,2.6804)rectangle(1.0718,2.7390);
\fill[DeepBlue!4.46](0.9767,2.7370)rectangle(1.0718,2.7955);
\fill[DeepBlue!2.84](0.9767,2.7935)rectangle(1.0718,2.8520);
\fill[DeepBlue!6.63](1.0698,1.5500)rectangle(1.1648,1.6085);
\fill[DeepBlue!3.85](1.0698,1.6065)rectangle(1.1648,1.6650);
\fill[DeepBlue!1.39](1.0698,1.6630)rectangle(1.1648,1.7216);
\fill[DeepBlue!4.79](1.0698,1.7196)rectangle(1.1648,1.7781);
\fill[DeepBlue!9.65](1.0698,1.7761)rectangle(1.1648,1.8346);
\fill[DeepBlue!14.85](1.0698,1.8326)rectangle(1.1648,1.8911);
\fill[DeepBlue!20.06](1.0698,1.8891)rectangle(1.1648,1.9477);
\fill[DeepBlue!24.98](1.0698,1.9457)rectangle(1.1648,2.0042);
\fill[DeepBlue!29.32](1.0698,2.0022)rectangle(1.1648,2.0607);
\fill[DeepBlue!32.83](1.0698,2.0587)rectangle(1.1648,2.1172);
\fill[DeepBlue!35.30](1.0698,2.1152)rectangle(1.1648,2.1737);
\fill[DeepBlue!36.57](1.0698,2.1717)rectangle(1.1648,2.2303);
\fill[DeepBlue!36.57](1.0698,2.2283)rectangle(1.1648,2.2868);
\fill[DeepBlue!35.30](1.0698,2.2848)rectangle(1.1648,2.3433);
\fill[DeepBlue!32.83](1.0698,2.3413)rectangle(1.1648,2.3998);
\fill[DeepBlue!29.32](1.0698,2.3978)rectangle(1.1648,2.4563);
\fill[DeepBlue!24.98](1.0698,2.4543)rectangle(1.1648,2.5129);
\fill[DeepBlue!20.06](1.0698,2.5109)rectangle(1.1648,2.5694);
\fill[DeepBlue!14.85](1.0698,2.5674)rectangle(1.1648,2.6259);
\fill[DeepBlue!9.65](1.0698,2.6239)rectangle(1.1648,2.6824);
\fill[DeepBlue!4.79](1.0698,2.6804)rectangle(1.1648,2.7390);
\fill[DeepBlue!1.39](1.0698,2.7370)rectangle(1.1648,2.7955);
\fill[DeepBlue!3.85](1.0698,2.7935)rectangle(1.1648,2.8520);
\fill[DeepBlue!8.19](1.1628,1.5500)rectangle(1.2578,1.6085);
\fill[DeepBlue!6.04](1.1628,1.6065)rectangle(1.2578,1.6650);
\fill[DeepBlue!3.12](1.1628,1.6630)rectangle(1.2578,1.7216);
\fill[DeepBlue!0.49](1.1628,1.7196)rectangle(1.2578,1.7781);
\fill[DeepBlue!4.53](1.1628,1.7761)rectangle(1.2578,1.8346);
\fill[DeepBlue!8.88](1.1628,1.8326)rectangle(1.2578,1.8911);
\fill[DeepBlue!13.26](1.1628,1.8891)rectangle(1.2578,1.9477);
\fill[DeepBlue!17.42](1.1628,1.9457)rectangle(1.2578,2.0042);
\fill[DeepBlue!21.12](1.1628,2.0022)rectangle(1.2578,2.0607);
\fill[DeepBlue!24.11](1.1628,2.0587)rectangle(1.2578,2.1172);
\fill[DeepBlue!26.22](1.1628,2.1152)rectangle(1.2578,2.1737);
\fill[DeepBlue!27.31](1.1628,2.1717)rectangle(1.2578,2.2303);
\fill[DeepBlue!27.31](1.1628,2.2283)rectangle(1.2578,2.2868);
\fill[DeepBlue!26.22](1.1628,2.2848)rectangle(1.2578,2.3433);
\fill[DeepBlue!24.11](1.1628,2.3413)rectangle(1.2578,2.3998);
\fill[DeepBlue!21.12](1.1628,2.3978)rectangle(1.2578,2.4563);
\fill[DeepBlue!17.42](1.1628,2.4543)rectangle(1.2578,2.5129);
\fill[DeepBlue!13.26](1.1628,2.5109)rectangle(1.2578,2.5694);
\fill[DeepBlue!8.88](1.1628,2.5674)rectangle(1.2578,2.6259);
\fill[DeepBlue!4.53](1.1628,2.6239)rectangle(1.2578,2.6824);
\fill[DeepBlue!0.49](1.1628,2.6804)rectangle(1.2578,2.7390);
\fill[DeepBlue!3.12](1.1628,2.7370)rectangle(1.2578,2.7955);
\fill[DeepBlue!6.04](1.1628,2.7935)rectangle(1.2578,2.8520);
\fill[DeepBlue!9.46](1.2558,1.5500)rectangle(1.3508,1.6085);
\fill[DeepBlue!8.07](1.2558,1.6065)rectangle(1.3508,1.6650);
\fill[DeepBlue!6.00](1.2558,1.6630)rectangle(1.3508,1.7216);
\fill[DeepBlue!3.37](1.2558,1.7196)rectangle(1.3508,1.7781);
\fill[DeepBlue!0.94](1.2558,1.7761)rectangle(1.3508,1.8346);
\fill[DeepBlue!3.48](1.2558,1.8326)rectangle(1.3508,1.8911);
\fill[DeepBlue!6.93](1.2558,1.8891)rectangle(1.3508,1.9477);
\fill[DeepBlue!10.29](1.2558,1.9457)rectangle(1.3508,2.0042);
\fill[DeepBlue!13.29](1.2558,2.0022)rectangle(1.3508,2.0607);
\fill[DeepBlue!15.75](1.2558,2.0587)rectangle(1.3508,2.1172);
\fill[DeepBlue!17.48](1.2558,2.1152)rectangle(1.3508,2.1737);
\fill[DeepBlue!18.38](1.2558,2.1717)rectangle(1.3508,2.2303);
\fill[DeepBlue!18.38](1.2558,2.2283)rectangle(1.3508,2.2868);
\fill[DeepBlue!17.48](1.2558,2.2848)rectangle(1.3508,2.3433);
\fill[DeepBlue!15.75](1.2558,2.3413)rectangle(1.3508,2.3998);
\fill[DeepBlue!13.29](1.2558,2.3978)rectangle(1.3508,2.4563);
\fill[DeepBlue!10.29](1.2558,2.4543)rectangle(1.3508,2.5129);
\fill[DeepBlue!6.93](1.2558,2.5109)rectangle(1.3508,2.5694);
\fill[DeepBlue!3.48](1.2558,2.5674)rectangle(1.3508,2.6259);
\fill[DeepBlue!0.94](1.2558,2.6239)rectangle(1.3508,2.6824);
\fill[DeepBlue!3.37](1.2558,2.6804)rectangle(1.3508,2.7390);
\fill[DeepBlue!6.00](1.2558,2.7370)rectangle(1.3508,2.7955);
\fill[DeepBlue!8.07](1.2558,2.7935)rectangle(1.3508,2.8520);
\fill[DeepBlue!10.22](1.3488,1.5500)rectangle(1.4439,1.6085);
\fill[DeepBlue!9.56](1.3488,1.6065)rectangle(1.4439,1.6650);
\fill[DeepBlue!8.30](1.3488,1.6630)rectangle(1.4439,1.7216);
\fill[DeepBlue!6.54](1.3488,1.7196)rectangle(1.4439,1.7781);
\fill[DeepBlue!4.43](1.3488,1.7761)rectangle(1.4439,1.8346);
\fill[DeepBlue!2.40](1.3488,1.8326)rectangle(1.4439,1.8911);
\fill[DeepBlue!2.27](1.3488,1.8891)rectangle(1.4439,1.9477);
\fill[DeepBlue!4.28](1.3488,1.9457)rectangle(1.4439,2.0042);
\fill[DeepBlue!6.46](1.3488,2.0022)rectangle(1.4439,2.0607);
\fill[DeepBlue!8.32](1.3488,2.0587)rectangle(1.4439,2.1172);
\fill[DeepBlue!9.66](1.3488,2.1152)rectangle(1.4439,2.1737);
\fill[DeepBlue!10.36](1.3488,2.1717)rectangle(1.4439,2.2303);
\fill[DeepBlue!10.36](1.3488,2.2283)rectangle(1.4439,2.2868);
\fill[DeepBlue!9.66](1.3488,2.2848)rectangle(1.4439,2.3433);
\fill[DeepBlue!8.32](1.3488,2.3413)rectangle(1.4439,2.3998);
\fill[DeepBlue!6.46](1.3488,2.3978)rectangle(1.4439,2.4563);
\fill[DeepBlue!4.28](1.3488,2.4543)rectangle(1.4439,2.5129);
\fill[DeepBlue!2.27](1.3488,2.5109)rectangle(1.4439,2.5694);
\fill[DeepBlue!2.40](1.3488,2.5674)rectangle(1.4439,2.6259);
\fill[DeepBlue!4.43](1.3488,2.6239)rectangle(1.4439,2.6824);
\fill[DeepBlue!6.54](1.3488,2.6804)rectangle(1.4439,2.7390);
\fill[DeepBlue!8.30](1.3488,2.7370)rectangle(1.4439,2.7955);
\fill[DeepBlue!9.56](1.3488,2.7935)rectangle(1.4439,2.8520);
\fill[DeepBlue!10.34](1.4419,1.5500)rectangle(1.5369,1.6085);
\fill[DeepBlue!10.32](1.4419,1.6065)rectangle(1.5369,1.6650);
\fill[DeepBlue!9.79](1.4419,1.6630)rectangle(1.5369,1.7216);
\fill[DeepBlue!8.81](1.4419,1.7196)rectangle(1.5369,1.7781);
\fill[DeepBlue!7.47](1.4419,1.7761)rectangle(1.5369,1.8346);
\fill[DeepBlue!5.91](1.4419,1.8326)rectangle(1.5369,1.8911);
\fill[DeepBlue!4.34](1.4419,1.8891)rectangle(1.5369,1.9477);
\fill[DeepBlue!3.12](1.4419,1.9457)rectangle(1.5369,2.0042);
\fill[DeepBlue!2.75](1.4419,2.0022)rectangle(1.5369,2.0607);
\fill[DeepBlue!3.19](1.4419,2.0587)rectangle(1.5369,2.1172);
\fill[DeepBlue!3.82](1.4419,2.1152)rectangle(1.5369,2.1737);
\fill[DeepBlue!4.21](1.4419,2.1717)rectangle(1.5369,2.2303);
\fill[DeepBlue!4.21](1.4419,2.2283)rectangle(1.5369,2.2868);
\fill[DeepBlue!3.82](1.4419,2.2848)rectangle(1.5369,2.3433);
\fill[DeepBlue!3.19](1.4419,2.3413)rectangle(1.5369,2.3998);
\fill[DeepBlue!2.75](1.4419,2.3978)rectangle(1.5369,2.4563);
\fill[DeepBlue!3.12](1.4419,2.4543)rectangle(1.5369,2.5129);
\fill[DeepBlue!4.34](1.4419,2.5109)rectangle(1.5369,2.5694);
\fill[DeepBlue!5.91](1.4419,2.5674)rectangle(1.5369,2.6259);
\fill[DeepBlue!7.47](1.4419,2.6239)rectangle(1.5369,2.6824);
\fill[DeepBlue!8.81](1.4419,2.6804)rectangle(1.5369,2.7390);
\fill[DeepBlue!9.79](1.4419,2.7370)rectangle(1.5369,2.7955);
\fill[DeepBlue!10.32](1.4419,2.7935)rectangle(1.5369,2.8520);
\fill[DeepBlue!9.78](1.5349,1.5500)rectangle(1.6299,1.6085);
\fill[DeepBlue!10.29](1.5349,1.6065)rectangle(1.6299,1.6650);
\fill[DeepBlue!10.39](1.5349,1.6630)rectangle(1.6299,1.7216);
\fill[DeepBlue!10.09](1.5349,1.7196)rectangle(1.6299,1.7781);
\fill[DeepBlue!9.44](1.5349,1.7761)rectangle(1.6299,1.8346);
\fill[DeepBlue!8.55](1.5349,1.8326)rectangle(1.6299,1.8911);
\fill[DeepBlue!7.52](1.5349,1.8891)rectangle(1.6299,1.9477);
\fill[DeepBlue!6.47](1.5349,1.9457)rectangle(1.6299,2.0042);
\fill[DeepBlue!5.53](1.5349,2.0022)rectangle(1.6299,2.0607);
\fill[DeepBlue!4.81](1.5349,2.0587)rectangle(1.6299,2.1172);
\fill[DeepBlue!4.34](1.5349,2.1152)rectangle(1.6299,2.1737);
\fill[DeepBlue!4.13](1.5349,2.1717)rectangle(1.6299,2.2303);
\fill[DeepBlue!4.13](1.5349,2.2283)rectangle(1.6299,2.2868);
\fill[DeepBlue!4.34](1.5349,2.2848)rectangle(1.6299,2.3433);
\fill[DeepBlue!4.81](1.5349,2.3413)rectangle(1.6299,2.3998);
\fill[DeepBlue!5.53](1.5349,2.3978)rectangle(1.6299,2.4563);
\fill[DeepBlue!6.47](1.5349,2.4543)rectangle(1.6299,2.5129);
\fill[DeepBlue!7.52](1.5349,2.5109)rectangle(1.6299,2.5694);
\fill[DeepBlue!8.55](1.5349,2.5674)rectangle(1.6299,2.6259);
\fill[DeepBlue!9.44](1.5349,2.6239)rectangle(1.6299,2.6824);
\fill[DeepBlue!10.09](1.5349,2.6804)rectangle(1.6299,2.7390);
\fill[DeepBlue!10.39](1.5349,2.7370)rectangle(1.6299,2.7955);
\fill[DeepBlue!10.29](1.5349,2.7935)rectangle(1.6299,2.8520);
\fill[DeepBlue!8.60](1.6279,1.5500)rectangle(1.7229,1.6085);
\fill[DeepBlue!9.52](1.6279,1.6065)rectangle(1.7229,1.6650);
\fill[DeepBlue!10.11](1.6279,1.6630)rectangle(1.7229,1.7216);
\fill[DeepBlue!10.36](1.6279,1.7196)rectangle(1.7229,1.7781);
\fill[DeepBlue!10.30](1.6279,1.7761)rectangle(1.7229,1.8346);
\fill[DeepBlue!10.00](1.6279,1.8326)rectangle(1.7229,1.8911);
\fill[DeepBlue!9.53](1.6279,1.8891)rectangle(1.7229,1.9477);
\fill[DeepBlue!8.96](1.6279,1.9457)rectangle(1.7229,2.0042);
\fill[DeepBlue!8.39](1.6279,2.0022)rectangle(1.7229,2.0607);
\fill[DeepBlue!7.90](1.6279,2.0587)rectangle(1.7229,2.1172);
\fill[DeepBlue!7.53](1.6279,2.1152)rectangle(1.7229,2.1737);
\fill[DeepBlue!7.34](1.6279,2.1717)rectangle(1.7229,2.2303);
\fill[DeepBlue!7.34](1.6279,2.2283)rectangle(1.7229,2.2868);
\fill[DeepBlue!7.53](1.6279,2.2848)rectangle(1.7229,2.3433);
\fill[DeepBlue!7.90](1.6279,2.3413)rectangle(1.7229,2.3998);
\fill[DeepBlue!8.39](1.6279,2.3978)rectangle(1.7229,2.4563);
\fill[DeepBlue!8.96](1.6279,2.4543)rectangle(1.7229,2.5129);
\fill[DeepBlue!9.53](1.6279,2.5109)rectangle(1.7229,2.5694);
\fill[DeepBlue!10.00](1.6279,2.5674)rectangle(1.7229,2.6259);
\fill[DeepBlue!10.30](1.6279,2.6239)rectangle(1.7229,2.6824);
\fill[DeepBlue!10.36](1.6279,2.6804)rectangle(1.7229,2.7390);
\fill[DeepBlue!10.11](1.6279,2.7370)rectangle(1.7229,2.7955);
\fill[DeepBlue!9.52](1.6279,2.7935)rectangle(1.7229,2.8520);
\fill[DeepBlue!6.91](1.7209,1.5500)rectangle(1.8160,1.6085);
\fill[DeepBlue!8.10](1.7209,1.6065)rectangle(1.8160,1.6650);
\fill[DeepBlue!9.03](1.7209,1.6630)rectangle(1.8160,1.7216);
\fill[DeepBlue!9.70](1.7209,1.7196)rectangle(1.8160,1.7781);
\fill[DeepBlue!10.11](1.7209,1.7761)rectangle(1.8160,1.8346);
\fill[DeepBlue!10.28](1.7209,1.8326)rectangle(1.8160,1.8911);
\fill[DeepBlue!10.27](1.7209,1.8891)rectangle(1.8160,1.9477);
\fill[DeepBlue!10.14](1.7209,1.9457)rectangle(1.8160,2.0042);
\fill[DeepBlue!9.94](1.7209,2.0022)rectangle(1.8160,2.0607);
\fill[DeepBlue!9.73](1.7209,2.0587)rectangle(1.8160,2.1172);
\fill[DeepBlue!9.56](1.7209,2.1152)rectangle(1.8160,2.1737);
\fill[DeepBlue!9.47](1.7209,2.1717)rectangle(1.8160,2.2303);
\fill[DeepBlue!9.47](1.7209,2.2283)rectangle(1.8160,2.2868);
\fill[DeepBlue!9.56](1.7209,2.2848)rectangle(1.8160,2.3433);
\fill[DeepBlue!9.73](1.7209,2.3413)rectangle(1.8160,2.3998);
\fill[DeepBlue!9.94](1.7209,2.3978)rectangle(1.8160,2.4563);
\fill[DeepBlue!10.14](1.7209,2.4543)rectangle(1.8160,2.5129);
\fill[DeepBlue!10.27](1.7209,2.5109)rectangle(1.8160,2.5694);
\fill[DeepBlue!10.28](1.7209,2.5674)rectangle(1.8160,2.6259);
\fill[DeepBlue!10.11](1.7209,2.6239)rectangle(1.8160,2.6824);
\fill[DeepBlue!9.70](1.7209,2.6804)rectangle(1.8160,2.7390);
\fill[DeepBlue!9.03](1.7209,2.7370)rectangle(1.8160,2.7955);
\fill[DeepBlue!8.10](1.7209,2.7935)rectangle(1.8160,2.8520);
\fill[DeepBlue!4.86](1.8140,1.5500)rectangle(1.9090,1.6085);
\fill[DeepBlue!6.18](1.8140,1.6065)rectangle(1.9090,1.6650);
\fill[DeepBlue!7.33](1.8140,1.6630)rectangle(1.9090,1.7216);
\fill[DeepBlue!8.27](1.8140,1.7196)rectangle(1.9090,1.7781);
\fill[DeepBlue!9.00](1.8140,1.7761)rectangle(1.9090,1.8346);
\fill[DeepBlue!9.52](1.8140,1.8326)rectangle(1.9090,1.8911);
\fill[DeepBlue!9.86](1.8140,1.8891)rectangle(1.9090,1.9477);
\fill[DeepBlue!10.07](1.8140,1.9457)rectangle(1.9090,2.0042);
\fill[DeepBlue!10.16](1.8140,2.0022)rectangle(1.9090,2.0607);
\fill[DeepBlue!10.19](1.8140,2.0587)rectangle(1.9090,2.1172);
\fill[DeepBlue!10.19](1.8140,2.1152)rectangle(1.9090,2.1737);
\fill[DeepBlue!10.18](1.8140,2.1717)rectangle(1.9090,2.2303);
\fill[DeepBlue!10.18](1.8140,2.2283)rectangle(1.9090,2.2868);
\fill[DeepBlue!10.19](1.8140,2.2848)rectangle(1.9090,2.3433);
\fill[DeepBlue!10.19](1.8140,2.3413)rectangle(1.9090,2.3998);
\fill[DeepBlue!10.16](1.8140,2.3978)rectangle(1.9090,2.4563);
\fill[DeepBlue!10.07](1.8140,2.4543)rectangle(1.9090,2.5129);
\fill[DeepBlue!9.86](1.8140,2.5109)rectangle(1.9090,2.5694);
\fill[DeepBlue!9.52](1.8140,2.5674)rectangle(1.9090,2.6259);
\fill[DeepBlue!9.00](1.8140,2.6239)rectangle(1.9090,2.6824);
\fill[DeepBlue!8.27](1.8140,2.6804)rectangle(1.9090,2.7390);
\fill[DeepBlue!7.33](1.8140,2.7370)rectangle(1.9090,2.7955);
\fill[DeepBlue!6.18](1.8140,2.7935)rectangle(1.9090,2.8520);
\fill[DeepBlue!2.67](1.9070,1.5500)rectangle(2.0020,1.6085);
\fill[DeepBlue!3.99](1.9070,1.6065)rectangle(2.0020,1.6650);
\fill[DeepBlue!5.21](1.9070,1.6630)rectangle(2.0020,1.7216);
\fill[DeepBlue!6.28](1.9070,1.7196)rectangle(2.0020,1.7781);
\fill[DeepBlue!7.19](1.9070,1.7761)rectangle(2.0020,1.8346);
\fill[DeepBlue!7.94](1.9070,1.8326)rectangle(2.0020,1.8911);
\fill[DeepBlue!8.51](1.9070,1.8891)rectangle(2.0020,1.9477);
\fill[DeepBlue!8.95](1.9070,1.9457)rectangle(2.0020,2.0042);
\fill[DeepBlue!9.25](1.9070,2.0022)rectangle(2.0020,2.0607);
\fill[DeepBlue!9.45](1.9070,2.0587)rectangle(2.0020,2.1172);
\fill[DeepBlue!9.58](1.9070,2.1152)rectangle(2.0020,2.1737);
\fill[DeepBlue!9.63](1.9070,2.1717)rectangle(2.0020,2.2303);
\fill[DeepBlue!9.63](1.9070,2.2283)rectangle(2.0020,2.2868);
\fill[DeepBlue!9.58](1.9070,2.2848)rectangle(2.0020,2.3433);
\fill[DeepBlue!9.45](1.9070,2.3413)rectangle(2.0020,2.3998);
\fill[DeepBlue!9.25](1.9070,2.3978)rectangle(2.0020,2.4563);
\fill[DeepBlue!8.95](1.9070,2.4543)rectangle(2.0020,2.5129);
\fill[DeepBlue!8.51](1.9070,2.5109)rectangle(2.0020,2.5694);
\fill[DeepBlue!7.94](1.9070,2.5674)rectangle(2.0020,2.6259);
\fill[DeepBlue!7.19](1.9070,2.6239)rectangle(2.0020,2.6824);
\fill[DeepBlue!6.28](1.9070,2.6804)rectangle(2.0020,2.7390);
\fill[DeepBlue!5.21](1.9070,2.7370)rectangle(2.0020,2.7955);
\fill[DeepBlue!3.99](1.9070,2.7935)rectangle(2.0020,2.8520);
\draw[->](-2.2,0)--(2.3,0)node[right]{$u$};
\draw[TextGray,dashed,thin]plot coordinates{(-2.200,0.000067)(-2.180,0.000021)(-2.160,0.000001)(-2.140,0.000008)(-2.120,0.000045)(-2.100,0.000113)(-2.080,0.000211)(-2.060,0.000341)(-2.040,0.000502)(-2.020,0.000690)(-2.000,0.000906)(-1.980,0.001144)(-1.960,0.001402)(-1.940,0.001676)(-1.920,0.001959)(-1.900,0.002249)(-1.880,0.002537)(-1.860,0.002819)(-1.840,0.003089)(-1.820,0.003340)(-1.800,0.003567)(-1.780,0.003763)(-1.760,0.003925)(-1.740,0.004047)(-1.720,0.004125)(-1.700,0.004157)(-1.680,0.004141)(-1.660,0.004075)(-1.640,0.003959)(-1.620,0.003796)(-1.600,0.003588)(-1.580,0.003338)(-1.560,0.003053)(-1.540,0.002737)(-1.520,0.002400)(-1.500,0.002049)(-1.480,0.001694)(-1.460,0.001346)(-1.440,0.001014)(-1.420,0.000712)(-1.400,0.000448)(-1.380,0.000236)(-1.360,0.000086)(-1.340,0.000009)(-1.320,0.000013)(-1.300,0.000108)(-1.280,0.000300)(-1.260,0.000594)(-1.240,0.000995)(-1.220,0.001503)(-1.200,0.002118)(-1.180,0.002837)(-1.160,0.003653)(-1.140,0.004560)(-1.120,0.005547)(-1.100,0.006600)(-1.080,0.007705)(-1.060,0.008843)(-1.040,0.009997)(-1.020,0.011145)(-1.000,0.012265)(-0.980,0.013334)(-0.960,0.014328)(-0.940,0.015225)(-0.920,0.016001)(-0.900,0.016635)(-0.880,0.017107)(-0.860,0.017400)(-0.840,0.017498)(-0.820,0.017391)(-0.800,0.017071)(-0.780,0.016537)(-0.760,0.015792)(-0.740,0.014844)(-0.720,0.013709)(-0.700,0.012407)(-0.680,0.010968)(-0.660,0.009427)(-0.640,0.007826)(-0.620,0.006215)(-0.600,0.004652)(-0.580,0.003199)(-0.560,0.001929)(-0.540,0.000916)(-0.520,0.000244)(-0.500,0.000000)(-0.480,0.000275)(-0.460,0.001165)(-0.440,0.002766)(-0.420,0.005177)(-0.400,0.008497)(-0.380,0.012824)(-0.360,0.018254)(-0.340,0.024879)(-0.320,0.032785)(-0.300,0.042055)(-0.280,0.052761)(-0.260,0.064969)(-0.240,0.078734)(-0.220,0.094100)(-0.200,0.111099)(-0.180,0.129748)(-0.160,0.150053)(-0.140,0.172003)(-0.120,0.195571)(-0.100,0.220718)(-0.080,0.247383)(-0.060,0.275494)(-0.040,0.304960)(-0.020,0.335675)(0.000,0.367516)(0.020,0.400347)(0.040,0.434018)(0.060,0.468364)(0.080,0.503209)(0.100,0.538368)(0.120,0.573645)(0.140,0.608836)(0.160,0.643732)(0.180,0.678122)(0.200,0.711788)(0.220,0.744518)(0.240,0.776096)(0.260,0.806314)(0.280,0.834967)(0.300,0.861861)(0.320,0.886808)(0.340,0.909635)(0.360,0.930179)(0.380,0.948295)(0.400,0.963852)(0.420,0.976738)(0.440,0.986859)(0.460,0.994142)(0.480,0.998533)(0.500,1.000000)(0.520,0.998533)(0.540,0.994142)(0.560,0.986859)(0.580,0.976738)(0.600,0.963852)(0.620,0.948295)(0.640,0.930179)(0.660,0.909635)(0.680,0.886808)(0.700,0.861861)(0.720,0.834967)(0.740,0.806314)(0.760,0.776096)(0.780,0.744518)(0.800,0.711788)(0.820,0.678122)(0.840,0.643732)(0.860,0.608836)(0.880,0.573645)(0.900,0.538368)(0.920,0.503209)(0.940,0.468364)(0.960,0.434018)(0.980,0.400347)(1.000,0.367516)(1.020,0.335675)(1.040,0.304960)(1.060,0.275494)(1.080,0.247383)(1.100,0.220718)(1.120,0.195571)(1.140,0.172003)(1.160,0.150053)(1.180,0.129748)(1.200,0.111099)(1.220,0.094100)(1.240,0.078734)(1.260,0.064969)(1.280,0.052761)(1.300,0.042055)(1.320,0.032785)(1.340,0.024879)(1.360,0.018254)(1.380,0.012824)(1.400,0.008497)(1.420,0.005177)(1.440,0.002766)(1.460,0.001165)(1.480,0.000275)(1.500,0.000000)(1.520,0.000244)(1.540,0.000916)(1.560,0.001929)(1.580,0.003199)(1.600,0.004652)(1.620,0.006215)(1.640,0.007826)(1.660,0.009427)(1.680,0.010968)(1.700,0.012407)(1.720,0.013709)(1.740,0.014844)(1.760,0.015792)(1.780,0.016537)(1.800,0.017071)(1.820,0.017391)(1.840,0.017498)(1.860,0.017400)(1.880,0.017107)(1.900,0.016635)(1.920,0.016001)(1.940,0.015225)(1.960,0.014328)(1.980,0.013334)(2.000,0.012265)(2.020,0.011145)(2.040,0.009997)(2.060,0.008843)(2.080,0.007705)(2.100,0.006600)(2.120,0.005547)(2.140,0.004560)(2.160,0.003653)(2.180,0.002837)(2.200,0.002118)};
\draw[TextGray,dashed,thin]plot coordinates{(-2.200,0.002118)(-2.180,0.002837)(-2.160,0.003653)(-2.140,0.004560)(-2.120,0.005547)(-2.100,0.006600)(-2.080,0.007705)(-2.060,0.008843)(-2.040,0.009997)(-2.020,0.011145)(-2.000,0.012265)(-1.980,0.013334)(-1.960,0.014328)(-1.940,0.015225)(-1.920,0.016001)(-1.900,0.016635)(-1.880,0.017107)(-1.860,0.017400)(-1.840,0.017498)(-1.820,0.017391)(-1.800,0.017071)(-1.780,0.016537)(-1.760,0.015792)(-1.740,0.014844)(-1.720,0.013709)(-1.700,0.012407)(-1.680,0.010968)(-1.660,0.009427)(-1.640,0.007826)(-1.620,0.006215)(-1.600,0.004652)(-1.580,0.003199)(-1.560,0.001929)(-1.540,0.000916)(-1.520,0.000244)(-1.500,0.000000)(-1.480,0.000275)(-1.460,0.001165)(-1.440,0.002766)(-1.420,0.005177)(-1.400,0.008497)(-1.380,0.012824)(-1.360,0.018254)(-1.340,0.024879)(-1.320,0.032785)(-1.300,0.042055)(-1.280,0.052761)(-1.260,0.064969)(-1.240,0.078734)(-1.220,0.094100)(-1.200,0.111099)(-1.180,0.129748)(-1.160,0.150053)(-1.140,0.172003)(-1.120,0.195571)(-1.100,0.220718)(-1.080,0.247383)(-1.060,0.275494)(-1.040,0.304960)(-1.020,0.335675)(-1.000,0.367516)(-0.980,0.400347)(-0.960,0.434018)(-0.940,0.468364)(-0.920,0.503209)(-0.900,0.538368)(-0.880,0.573645)(-0.860,0.608836)(-0.840,0.643732)(-0.820,0.678122)(-0.800,0.711788)(-0.780,0.744518)(-0.760,0.776096)(-0.740,0.806314)(-0.720,0.834967)(-0.700,0.861861)(-0.680,0.886808)(-0.660,0.909635)(-0.640,0.930179)(-0.620,0.948295)(-0.600,0.963852)(-0.580,0.976738)(-0.560,0.986859)(-0.540,0.994142)(-0.520,0.998533)(-0.500,1.000000)(-0.480,0.998533)(-0.460,0.994142)(-0.440,0.986859)(-0.420,0.976738)(-0.400,0.963852)(-0.380,0.948295)(-0.360,0.930179)(-0.340,0.909635)(-0.320,0.886808)(-0.300,0.861861)(-0.280,0.834967)(-0.260,0.806314)(-0.240,0.776096)(-0.220,0.744518)(-0.200,0.711788)(-0.180,0.678122)(-0.160,0.643732)(-0.140,0.608836)(-0.120,0.573645)(-0.100,0.538368)(-0.080,0.503209)(-0.060,0.468364)(-0.040,0.434018)(-0.020,0.400347)(0.000,0.367516)(0.020,0.335675)(0.040,0.304960)(0.060,0.275494)(0.080,0.247383)(0.100,0.220718)(0.120,0.195571)(0.140,0.172003)(0.160,0.150053)(0.180,0.129748)(0.200,0.111099)(0.220,0.094100)(0.240,0.078734)(0.260,0.064969)(0.280,0.052761)(0.300,0.042055)(0.320,0.032785)(0.340,0.024879)(0.360,0.018254)(0.380,0.012824)(0.400,0.008497)(0.420,0.005177)(0.440,0.002766)(0.460,0.001165)(0.480,0.000275)(0.500,0.000000)(0.520,0.000244)(0.540,0.000916)(0.560,0.001929)(0.580,0.003199)(0.600,0.004652)(0.620,0.006215)(0.640,0.007826)(0.660,0.009427)(0.680,0.010968)(0.700,0.012407)(0.720,0.013709)(0.740,0.014844)(0.760,0.015792)(0.780,0.016537)(0.800,0.017071)(0.820,0.017391)(0.840,0.017498)(0.860,0.017400)(0.880,0.017107)(0.900,0.016635)(0.920,0.016001)(0.940,0.015225)(0.960,0.014328)(0.980,0.013334)(1.000,0.012265)(1.020,0.011145)(1.040,0.009997)(1.060,0.008843)(1.080,0.007705)(1.100,0.006600)(1.120,0.005547)(1.140,0.004560)(1.160,0.003653)(1.180,0.002837)(1.200,0.002118)(1.220,0.001503)(1.240,0.000995)(1.260,0.000594)(1.280,0.000300)(1.300,0.000108)(1.320,0.000013)(1.340,0.000009)(1.360,0.000086)(1.380,0.000236)(1.400,0.000448)(1.420,0.000712)(1.440,0.001014)(1.460,0.001346)(1.480,0.001694)(1.500,0.002049)(1.520,0.002400)(1.540,0.002737)(1.560,0.003053)(1.580,0.003338)(1.600,0.003588)(1.620,0.003796)(1.640,0.003959)(1.660,0.004075)(1.680,0.004141)(1.700,0.004157)(1.720,0.004125)(1.740,0.004047)(1.760,0.003925)(1.780,0.003763)(1.800,0.003567)(1.820,0.003340)(1.840,0.003089)(1.860,0.002819)(1.880,0.002537)(1.900,0.002249)(1.920,0.001959)(1.940,0.001676)(1.960,0.001402)(1.980,0.001144)(2.000,0.000906)(2.020,0.000690)(2.040,0.000502)(2.060,0.000341)(2.080,0.000211)(2.100,0.000113)(2.120,0.000045)(2.140,0.000008)(2.160,0.000001)(2.180,0.000021)(2.200,0.000067)};
\draw[DeepBlue,thick]plot coordinates{(-2.200,0.002185)(-2.180,0.002858)(-2.160,0.003654)(-2.140,0.004568)(-2.120,0.005592)(-2.100,0.006712)(-2.080,0.007916)(-2.060,0.009185)(-2.040,0.010499)(-2.020,0.011836)(-2.000,0.013171)(-1.980,0.014478)(-1.960,0.015731)(-1.940,0.016901)(-1.920,0.017961)(-1.900,0.018884)(-1.880,0.019644)(-1.860,0.020219)(-1.840,0.020587)(-1.820,0.020730)(-1.800,0.020638)(-1.780,0.020301)(-1.760,0.019717)(-1.740,0.018891)(-1.720,0.017834)(-1.700,0.016565)(-1.680,0.015109)(-1.660,0.013501)(-1.640,0.011785)(-1.620,0.010011)(-1.600,0.008240)(-1.580,0.006538)(-1.560,0.004981)(-1.540,0.003654)(-1.520,0.002644)(-1.500,0.002049)(-1.480,0.001970)(-1.460,0.002511)(-1.440,0.003780)(-1.420,0.005889)(-1.400,0.008946)(-1.380,0.013061)(-1.360,0.018340)(-1.340,0.024887)(-1.320,0.032798)(-1.300,0.042162)(-1.280,0.053061)(-1.260,0.065563)(-1.240,0.079729)(-1.220,0.095603)(-1.200,0.113217)(-1.180,0.132585)(-1.160,0.153706)(-1.140,0.176563)(-1.120,0.201118)(-1.100,0.227317)(-1.080,0.255088)(-1.060,0.284338)(-1.040,0.314957)(-1.020,0.346820)(-1.000,0.379781)(-0.980,0.413681)(-0.960,0.448346)(-0.940,0.483589)(-0.920,0.519211)(-0.900,0.555004)(-0.880,0.590752)(-0.860,0.626235)(-0.840,0.661230)(-0.820,0.695512)(-0.800,0.728860)(-0.780,0.761055)(-0.760,0.791888)(-0.740,0.821158)(-0.720,0.848676)(-0.700,0.874268)(-0.680,0.897776)(-0.660,0.919061)(-0.640,0.938005)(-0.620,0.954510)(-0.600,0.968503)(-0.580,0.979937)(-0.560,0.988787)(-0.540,0.995058)(-0.520,0.998777)(-0.500,1.000000)(-0.480,0.998808)(-0.460,0.995306)(-0.440,0.989625)(-0.420,0.981915)(-0.400,0.972349)(-0.380,0.961119)(-0.360,0.948433)(-0.340,0.934513)(-0.320,0.919593)(-0.300,0.903915)(-0.280,0.887728)(-0.260,0.871283)(-0.240,0.854830)(-0.220,0.838618)(-0.200,0.822887)(-0.180,0.807870)(-0.160,0.793785)(-0.140,0.780838)(-0.120,0.769216)(-0.100,0.759086)(-0.080,0.750593)(-0.060,0.743858)(-0.040,0.738978)(-0.020,0.736022)(0.000,0.735032)(0.020,0.736022)(0.040,0.738978)(0.060,0.743858)(0.080,0.750593)(0.100,0.759086)(0.120,0.769216)(0.140,0.780838)(0.160,0.793785)(0.180,0.807870)(0.200,0.822887)(0.220,0.838618)(0.240,0.854830)(0.260,0.871283)(0.280,0.887728)(0.300,0.903915)(0.320,0.919593)(0.340,0.934513)(0.360,0.948433)(0.380,0.961119)(0.400,0.972349)(0.420,0.981915)(0.440,0.989625)(0.460,0.995306)(0.480,0.998808)(0.500,1.000000)(0.520,0.998777)(0.540,0.995058)(0.560,0.988787)(0.580,0.979937)(0.600,0.968503)(0.620,0.954510)(0.640,0.938005)(0.660,0.919061)(0.680,0.897776)(0.700,0.874268)(0.720,0.848676)(0.740,0.821158)(0.760,0.791888)(0.780,0.761055)(0.800,0.728860)(0.820,0.695512)(0.840,0.661230)(0.860,0.626235)(0.880,0.590752)(0.900,0.555004)(0.920,0.519211)(0.940,0.483589)(0.960,0.448346)(0.980,0.413681)(1.000,0.379781)(1.020,0.346820)(1.040,0.314957)(1.060,0.284338)(1.080,0.255088)(1.100,0.227317)(1.120,0.201118)(1.140,0.176563)(1.160,0.153706)(1.180,0.132585)(1.200,0.113217)(1.220,0.095603)(1.240,0.079729)(1.260,0.065563)(1.280,0.053061)(1.300,0.042162)(1.320,0.032798)(1.340,0.024887)(1.360,0.018340)(1.380,0.013061)(1.400,0.008946)(1.420,0.005889)(1.440,0.003780)(1.460,0.002511)(1.480,0.001970)(1.500,0.002049)(1.520,0.002644)(1.540,0.003654)(1.560,0.004981)(1.580,0.006538)(1.600,0.008240)(1.620,0.010011)(1.640,0.011785)(1.660,0.013501)(1.680,0.015109)(1.700,0.016565)(1.720,0.017834)(1.740,0.018891)(1.760,0.019717)(1.780,0.020301)(1.800,0.020638)(1.820,0.020730)(1.840,0.020587)(1.860,0.020219)(1.880,0.019644)(1.900,0.018884)(1.920,0.017961)(1.940,0.016901)(1.960,0.015731)(1.980,0.014478)(2.000,0.013171)(2.020,0.011836)(2.040,0.010499)(2.060,0.009185)(2.080,0.007916)(2.100,0.006712)(2.120,0.005592)(2.140,0.004568)(2.160,0.003654)(2.180,0.002858)(2.200,0.002185)};
\node[below,align=center]at(0,-.2){$\Delta\theta/\theta_{\rm A}=1$\\瑞利临界};
\end{scope}
\begin{scope}[xshift=9.2cm]
\fill[DeepBlue!4.30](-2.0000,1.5500)rectangle(-1.9050,1.6085);
\fill[DeepBlue!4.14](-2.0000,1.6065)rectangle(-1.9050,1.6650);
\fill[DeepBlue!4.02](-2.0000,1.6630)rectangle(-1.9050,1.7216);
\fill[DeepBlue!3.97](-2.0000,1.7196)rectangle(-1.9050,1.7781);
\fill[DeepBlue!3.99](-2.0000,1.7761)rectangle(-1.9050,1.8346);
\fill[DeepBlue!4.09](-2.0000,1.8326)rectangle(-1.9050,1.8911);
\fill[DeepBlue!4.24](-2.0000,1.8891)rectangle(-1.9050,1.9477);
\fill[DeepBlue!4.41](-2.0000,1.9457)rectangle(-1.9050,2.0042);
\fill[DeepBlue!4.58](-2.0000,2.0022)rectangle(-1.9050,2.0607);
\fill[DeepBlue!4.72](-2.0000,2.0587)rectangle(-1.9050,2.1172);
\fill[DeepBlue!4.82](-2.0000,2.1152)rectangle(-1.9050,2.1737);
\fill[DeepBlue!4.88](-2.0000,2.1717)rectangle(-1.9050,2.2303);
\fill[DeepBlue!4.88](-2.0000,2.2283)rectangle(-1.9050,2.2868);
\fill[DeepBlue!4.82](-2.0000,2.2848)rectangle(-1.9050,2.3433);
\fill[DeepBlue!4.72](-2.0000,2.3413)rectangle(-1.9050,2.3998);
\fill[DeepBlue!4.58](-2.0000,2.3978)rectangle(-1.9050,2.4563);
\fill[DeepBlue!4.41](-2.0000,2.4543)rectangle(-1.9050,2.5129);
\fill[DeepBlue!4.24](-2.0000,2.5109)rectangle(-1.9050,2.5694);
\fill[DeepBlue!4.09](-2.0000,2.5674)rectangle(-1.9050,2.6259);
\fill[DeepBlue!3.99](-2.0000,2.6239)rectangle(-1.9050,2.6824);
\fill[DeepBlue!3.97](-2.0000,2.6804)rectangle(-1.9050,2.7390);
\fill[DeepBlue!4.02](-2.0000,2.7370)rectangle(-1.9050,2.7955);
\fill[DeepBlue!4.14](-2.0000,2.7935)rectangle(-1.9050,2.8520);
\fill[DeepBlue!4.19](-1.9070,1.5500)rectangle(-1.8120,1.6085);
\fill[DeepBlue!4.13](-1.9070,1.6065)rectangle(-1.8120,1.6650);
\fill[DeepBlue!4.19](-1.9070,1.6630)rectangle(-1.8120,1.7216);
\fill[DeepBlue!4.35](-1.9070,1.7196)rectangle(-1.8120,1.7781);
\fill[DeepBlue!4.62](-1.9070,1.7761)rectangle(-1.8120,1.8346);
\fill[DeepBlue!4.93](-1.9070,1.8326)rectangle(-1.8120,1.8911);
\fill[DeepBlue!5.26](-1.9070,1.8891)rectangle(-1.8120,1.9477);
\fill[DeepBlue!5.56](-1.9070,1.9457)rectangle(-1.8120,2.0042);
\fill[DeepBlue!5.83](-1.9070,2.0022)rectangle(-1.8120,2.0607);
\fill[DeepBlue!6.04](-1.9070,2.0587)rectangle(-1.8120,2.1172);
\fill[DeepBlue!6.18](-1.9070,2.1152)rectangle(-1.8120,2.1737);
\fill[DeepBlue!6.25](-1.9070,2.1717)rectangle(-1.8120,2.2303);
\fill[DeepBlue!6.25](-1.9070,2.2283)rectangle(-1.8120,2.2868);
\fill[DeepBlue!6.18](-1.9070,2.2848)rectangle(-1.8120,2.3433);
\fill[DeepBlue!6.04](-1.9070,2.3413)rectangle(-1.8120,2.3998);
\fill[DeepBlue!5.83](-1.9070,2.3978)rectangle(-1.8120,2.4563);
\fill[DeepBlue!5.56](-1.9070,2.4543)rectangle(-1.8120,2.5129);
\fill[DeepBlue!5.26](-1.9070,2.5109)rectangle(-1.8120,2.5694);
\fill[DeepBlue!4.93](-1.9070,2.5674)rectangle(-1.8120,2.6259);
\fill[DeepBlue!4.62](-1.9070,2.6239)rectangle(-1.8120,2.6824);
\fill[DeepBlue!4.35](-1.9070,2.6804)rectangle(-1.8120,2.7390);
\fill[DeepBlue!4.19](-1.9070,2.7370)rectangle(-1.8120,2.7955);
\fill[DeepBlue!4.13](-1.9070,2.7935)rectangle(-1.8120,2.8520);
\fill[DeepBlue!4.30](-1.8140,1.5500)rectangle(-1.7189,1.6085);
\fill[DeepBlue!4.47](-1.8140,1.6065)rectangle(-1.7189,1.6650);
\fill[DeepBlue!4.78](-1.8140,1.6630)rectangle(-1.7189,1.7216);
\fill[DeepBlue!5.19](-1.8140,1.7196)rectangle(-1.7189,1.7781);
\fill[DeepBlue!5.66](-1.8140,1.7761)rectangle(-1.7189,1.8346);
\fill[DeepBlue!6.12](-1.8140,1.8326)rectangle(-1.7189,1.8911);
\fill[DeepBlue!6.54](-1.8140,1.8891)rectangle(-1.7189,1.9477);
\fill[DeepBlue!6.91](-1.8140,1.9457)rectangle(-1.7189,2.0042);
\fill[DeepBlue!7.21](-1.8140,2.0022)rectangle(-1.7189,2.0607);
\fill[DeepBlue!7.43](-1.8140,2.0587)rectangle(-1.7189,2.1172);
\fill[DeepBlue!7.58](-1.8140,2.1152)rectangle(-1.7189,2.1737);
\fill[DeepBlue!7.65](-1.8140,2.1717)rectangle(-1.7189,2.2303);
\fill[DeepBlue!7.65](-1.8140,2.2283)rectangle(-1.7189,2.2868);
\fill[DeepBlue!7.58](-1.8140,2.2848)rectangle(-1.7189,2.3433);
\fill[DeepBlue!7.43](-1.8140,2.3413)rectangle(-1.7189,2.3998);
\fill[DeepBlue!7.21](-1.8140,2.3978)rectangle(-1.7189,2.4563);
\fill[DeepBlue!6.91](-1.8140,2.4543)rectangle(-1.7189,2.5129);
\fill[DeepBlue!6.54](-1.8140,2.5109)rectangle(-1.7189,2.5694);
\fill[DeepBlue!6.12](-1.8140,2.5674)rectangle(-1.7189,2.6259);
\fill[DeepBlue!5.66](-1.8140,2.6239)rectangle(-1.7189,2.6824);
\fill[DeepBlue!5.19](-1.8140,2.6804)rectangle(-1.7189,2.7390);
\fill[DeepBlue!4.78](-1.8140,2.7370)rectangle(-1.7189,2.7955);
\fill[DeepBlue!4.47](-1.8140,2.7935)rectangle(-1.7189,2.8520);
\fill[DeepBlue!4.72](-1.7209,1.5500)rectangle(-1.6259,1.6085);
\fill[DeepBlue!5.16](-1.7209,1.6065)rectangle(-1.6259,1.6650);
\fill[DeepBlue!5.71](-1.7209,1.6630)rectangle(-1.6259,1.7216);
\fill[DeepBlue!6.29](-1.7209,1.7196)rectangle(-1.6259,1.7781);
\fill[DeepBlue!6.86](-1.7209,1.7761)rectangle(-1.6259,1.8346);
\fill[DeepBlue!7.36](-1.7209,1.8326)rectangle(-1.6259,1.8911);
\fill[DeepBlue!7.79](-1.7209,1.8891)rectangle(-1.6259,1.9477);
\fill[DeepBlue!8.14](-1.7209,1.9457)rectangle(-1.6259,2.0042);
\fill[DeepBlue!8.40](-1.7209,2.0022)rectangle(-1.6259,2.0607);
\fill[DeepBlue!8.58](-1.7209,2.0587)rectangle(-1.6259,2.1172);
\fill[DeepBlue!8.69](-1.7209,2.1152)rectangle(-1.6259,2.1737);
\fill[DeepBlue!8.75](-1.7209,2.1717)rectangle(-1.6259,2.2303);
\fill[DeepBlue!8.75](-1.7209,2.2283)rectangle(-1.6259,2.2868);
\fill[DeepBlue!8.69](-1.7209,2.2848)rectangle(-1.6259,2.3433);
\fill[DeepBlue!8.58](-1.7209,2.3413)rectangle(-1.6259,2.3998);
\fill[DeepBlue!8.40](-1.7209,2.3978)rectangle(-1.6259,2.4563);
\fill[DeepBlue!8.14](-1.7209,2.4543)rectangle(-1.6259,2.5129);
\fill[DeepBlue!7.79](-1.7209,2.5109)rectangle(-1.6259,2.5694);
\fill[DeepBlue!7.36](-1.7209,2.5674)rectangle(-1.6259,2.6259);
\fill[DeepBlue!6.86](-1.7209,2.6239)rectangle(-1.6259,2.6824);
\fill[DeepBlue!6.29](-1.7209,2.6804)rectangle(-1.6259,2.7390);
\fill[DeepBlue!5.71](-1.7209,2.7370)rectangle(-1.6259,2.7955);
\fill[DeepBlue!5.16](-1.7209,2.7935)rectangle(-1.6259,2.8520);
\fill[DeepBlue!5.43](-1.6279,1.5500)rectangle(-1.5329,1.6085);
\fill[DeepBlue!6.09](-1.6279,1.6065)rectangle(-1.5329,1.6650);
\fill[DeepBlue!6.77](-1.6279,1.6630)rectangle(-1.5329,1.7216);
\fill[DeepBlue!7.42](-1.6279,1.7196)rectangle(-1.5329,1.7781);
\fill[DeepBlue!7.98](-1.6279,1.7761)rectangle(-1.5329,1.8346);
\fill[DeepBlue!8.43](-1.6279,1.8326)rectangle(-1.5329,1.8911);
\fill[DeepBlue!8.77](-1.6279,1.8891)rectangle(-1.5329,1.9477);
\fill[DeepBlue!9.01](-1.6279,1.9457)rectangle(-1.5329,2.0042);
\fill[DeepBlue!9.16](-1.6279,2.0022)rectangle(-1.5329,2.0607);
\fill[DeepBlue!9.25](-1.6279,2.0587)rectangle(-1.5329,2.1172);
\fill[DeepBlue!9.30](-1.6279,2.1152)rectangle(-1.5329,2.1737);
\fill[DeepBlue!9.32](-1.6279,2.1717)rectangle(-1.5329,2.2303);
\fill[DeepBlue!9.32](-1.6279,2.2283)rectangle(-1.5329,2.2868);
\fill[DeepBlue!9.30](-1.6279,2.2848)rectangle(-1.5329,2.3433);
\fill[DeepBlue!9.25](-1.6279,2.3413)rectangle(-1.5329,2.3998);
\fill[DeepBlue!9.16](-1.6279,2.3978)rectangle(-1.5329,2.4563);
\fill[DeepBlue!9.01](-1.6279,2.4543)rectangle(-1.5329,2.5129);
\fill[DeepBlue!8.77](-1.6279,2.5109)rectangle(-1.5329,2.5694);
\fill[DeepBlue!8.43](-1.6279,2.5674)rectangle(-1.5329,2.6259);
\fill[DeepBlue!7.98](-1.6279,2.6239)rectangle(-1.5329,2.6824);
\fill[DeepBlue!7.42](-1.6279,2.6804)rectangle(-1.5329,2.7390);
\fill[DeepBlue!6.77](-1.6279,2.7370)rectangle(-1.5329,2.7955);
\fill[DeepBlue!6.09](-1.6279,2.7935)rectangle(-1.5329,2.8520);
\fill[DeepBlue!6.31](-1.5349,1.5500)rectangle(-1.4399,1.6085);
\fill[DeepBlue!7.09](-1.5349,1.6065)rectangle(-1.4399,1.6650);
\fill[DeepBlue!7.80](-1.5349,1.6630)rectangle(-1.4399,1.7216);
\fill[DeepBlue!8.40](-1.5349,1.7196)rectangle(-1.4399,1.7781);
\fill[DeepBlue!8.84](-1.5349,1.7761)rectangle(-1.4399,1.8346);
\fill[DeepBlue!9.14](-1.5349,1.8326)rectangle(-1.4399,1.8911);
\fill[DeepBlue!9.30](-1.5349,1.8891)rectangle(-1.4399,1.9477);
\fill[DeepBlue!9.35](-1.5349,1.9457)rectangle(-1.4399,2.0042);
\fill[DeepBlue!9.34](-1.5349,2.0022)rectangle(-1.4399,2.0607);
\fill[DeepBlue!9.29](-1.5349,2.0587)rectangle(-1.4399,2.1172);
\fill[DeepBlue!9.24](-1.5349,2.1152)rectangle(-1.4399,2.1737);
\fill[DeepBlue!9.21](-1.5349,2.1717)rectangle(-1.4399,2.2303);
\fill[DeepBlue!9.21](-1.5349,2.2283)rectangle(-1.4399,2.2868);
\fill[DeepBlue!9.24](-1.5349,2.2848)rectangle(-1.4399,2.3433);
\fill[DeepBlue!9.29](-1.5349,2.3413)rectangle(-1.4399,2.3998);
\fill[DeepBlue!9.34](-1.5349,2.3978)rectangle(-1.4399,2.4563);
\fill[DeepBlue!9.35](-1.5349,2.4543)rectangle(-1.4399,2.5129);
\fill[DeepBlue!9.30](-1.5349,2.5109)rectangle(-1.4399,2.5694);
\fill[DeepBlue!9.14](-1.5349,2.5674)rectangle(-1.4399,2.6259);
\fill[DeepBlue!8.84](-1.5349,2.6239)rectangle(-1.4399,2.6824);
\fill[DeepBlue!8.40](-1.5349,2.6804)rectangle(-1.4399,2.7390);
\fill[DeepBlue!7.80](-1.5349,2.7370)rectangle(-1.4399,2.7955);
\fill[DeepBlue!7.09](-1.5349,2.7935)rectangle(-1.4399,2.8520);
\fill[DeepBlue!7.25](-1.4419,1.5500)rectangle(-1.3468,1.6085);
\fill[DeepBlue!8.04](-1.4419,1.6065)rectangle(-1.3468,1.6650);
\fill[DeepBlue!8.67](-1.4419,1.6630)rectangle(-1.3468,1.7216);
\fill[DeepBlue!9.11](-1.4419,1.7196)rectangle(-1.3468,1.7781);
\fill[DeepBlue!9.34](-1.4419,1.7761)rectangle(-1.3468,1.8346);
\fill[DeepBlue!9.39](-1.4419,1.8326)rectangle(-1.3468,1.8911);
\fill[DeepBlue!9.30](-1.4419,1.8891)rectangle(-1.3468,1.9477);
\fill[DeepBlue!9.12](-1.4419,1.9457)rectangle(-1.3468,2.0042);
\fill[DeepBlue!8.90](-1.4419,2.0022)rectangle(-1.3468,2.0607);
\fill[DeepBlue!8.69](-1.4419,2.0587)rectangle(-1.3468,2.1172);
\fill[DeepBlue!8.53](-1.4419,2.1152)rectangle(-1.3468,2.1737);
\fill[DeepBlue!8.45](-1.4419,2.1717)rectangle(-1.3468,2.2303);
\fill[DeepBlue!8.45](-1.4419,2.2283)rectangle(-1.3468,2.2868);
\fill[DeepBlue!8.53](-1.4419,2.2848)rectangle(-1.3468,2.3433);
\fill[DeepBlue!8.69](-1.4419,2.3413)rectangle(-1.3468,2.3998);
\fill[DeepBlue!8.90](-1.4419,2.3978)rectangle(-1.3468,2.4563);
\fill[DeepBlue!9.12](-1.4419,2.4543)rectangle(-1.3468,2.5129);
\fill[DeepBlue!9.30](-1.4419,2.5109)rectangle(-1.3468,2.5694);
\fill[DeepBlue!9.39](-1.4419,2.5674)rectangle(-1.3468,2.6259);
\fill[DeepBlue!9.34](-1.4419,2.6239)rectangle(-1.3468,2.6824);
\fill[DeepBlue!9.11](-1.4419,2.6804)rectangle(-1.3468,2.7390);
\fill[DeepBlue!8.67](-1.4419,2.7370)rectangle(-1.3468,2.7955);
\fill[DeepBlue!8.04](-1.4419,2.7935)rectangle(-1.3468,2.8520);
\fill[DeepBlue!8.16](-1.3488,1.5500)rectangle(-1.2538,1.6085);
\fill[DeepBlue!8.86](-1.3488,1.6065)rectangle(-1.2538,1.6650);
\fill[DeepBlue!9.30](-1.3488,1.6630)rectangle(-1.2538,1.7216);
\fill[DeepBlue!9.49](-1.3488,1.7196)rectangle(-1.2538,1.7781);
\fill[DeepBlue!9.43](-1.3488,1.7761)rectangle(-1.2538,1.8346);
\fill[DeepBlue!9.19](-1.3488,1.8326)rectangle(-1.2538,1.8911);
\fill[DeepBlue!8.83](-1.3488,1.8891)rectangle(-1.2538,1.9477);
\fill[DeepBlue!8.43](-1.3488,1.9457)rectangle(-1.2538,2.0042);
\fill[DeepBlue!8.06](-1.3488,2.0022)rectangle(-1.2538,2.0607);
\fill[DeepBlue!7.78](-1.3488,2.0587)rectangle(-1.2538,2.1172);
\fill[DeepBlue!7.59](-1.3488,2.1152)rectangle(-1.2538,2.1737);
\fill[DeepBlue!7.50](-1.3488,2.1717)rectangle(-1.2538,2.2303);
\fill[DeepBlue!7.50](-1.3488,2.2283)rectangle(-1.2538,2.2868);
\fill[DeepBlue!7.59](-1.3488,2.2848)rectangle(-1.2538,2.3433);
\fill[DeepBlue!7.78](-1.3488,2.3413)rectangle(-1.2538,2.3998);
\fill[DeepBlue!8.06](-1.3488,2.3978)rectangle(-1.2538,2.4563);
\fill[DeepBlue!8.43](-1.3488,2.4543)rectangle(-1.2538,2.5129);
\fill[DeepBlue!8.83](-1.3488,2.5109)rectangle(-1.2538,2.5694);
\fill[DeepBlue!9.19](-1.3488,2.5674)rectangle(-1.2538,2.6259);
\fill[DeepBlue!9.43](-1.3488,2.6239)rectangle(-1.2538,2.6824);
\fill[DeepBlue!9.49](-1.3488,2.6804)rectangle(-1.2538,2.7390);
\fill[DeepBlue!9.30](-1.3488,2.7370)rectangle(-1.2538,2.7955);
\fill[DeepBlue!8.86](-1.3488,2.7935)rectangle(-1.2538,2.8520);
\fill[DeepBlue!8.98](-1.2558,1.5500)rectangle(-1.1608,1.6085);
\fill[DeepBlue!9.48](-1.2558,1.6065)rectangle(-1.1608,1.6650);
\fill[DeepBlue!9.66](-1.2558,1.6630)rectangle(-1.1608,1.7216);
\fill[DeepBlue!9.54](-1.2558,1.7196)rectangle(-1.1608,1.7781);
\fill[DeepBlue!9.17](-1.2558,1.7761)rectangle(-1.1608,1.8346);
\fill[DeepBlue!8.68](-1.2558,1.8326)rectangle(-1.1608,1.8911);
\fill[DeepBlue!8.19](-1.2558,1.8891)rectangle(-1.1608,1.9477);
\fill[DeepBlue!7.82](-1.2558,1.9457)rectangle(-1.1608,2.0042);
\fill[DeepBlue!7.63](-1.2558,2.0022)rectangle(-1.1608,2.0607);
\fill[DeepBlue!7.60](-1.2558,2.0587)rectangle(-1.1608,2.1172);
\fill[DeepBlue!7.67](-1.2558,2.1152)rectangle(-1.1608,2.1737);
\fill[DeepBlue!7.73](-1.2558,2.1717)rectangle(-1.1608,2.2303);
\fill[DeepBlue!7.73](-1.2558,2.2283)rectangle(-1.1608,2.2868);
\fill[DeepBlue!7.67](-1.2558,2.2848)rectangle(-1.1608,2.3433);
\fill[DeepBlue!7.60](-1.2558,2.3413)rectangle(-1.1608,2.3998);
\fill[DeepBlue!7.63](-1.2558,2.3978)rectangle(-1.1608,2.4563);
\fill[DeepBlue!7.82](-1.2558,2.4543)rectangle(-1.1608,2.5129);
\fill[DeepBlue!8.19](-1.2558,2.5109)rectangle(-1.1608,2.5694);
\fill[DeepBlue!8.68](-1.2558,2.5674)rectangle(-1.1608,2.6259);
\fill[DeepBlue!9.17](-1.2558,2.6239)rectangle(-1.1608,2.6824);
\fill[DeepBlue!9.54](-1.2558,2.6804)rectangle(-1.1608,2.7390);
\fill[DeepBlue!9.66](-1.2558,2.7370)rectangle(-1.1608,2.7955);
\fill[DeepBlue!9.48](-1.2558,2.7935)rectangle(-1.1608,2.8520);
\fill[DeepBlue!9.66](-1.1628,1.5500)rectangle(-1.0678,1.6085);
\fill[DeepBlue!9.90](-1.1628,1.6065)rectangle(-1.0678,1.6650);
\fill[DeepBlue!9.77](-1.1628,1.6630)rectangle(-1.0678,1.7216);
\fill[DeepBlue!9.34](-1.1628,1.7196)rectangle(-1.0678,1.7781);
\fill[DeepBlue!8.78](-1.1628,1.7761)rectangle(-1.0678,1.8346);
\fill[DeepBlue!8.31](-1.1628,1.8326)rectangle(-1.0678,1.8911);
\fill[DeepBlue!8.15](-1.1628,1.8891)rectangle(-1.0678,1.9477);
\fill[DeepBlue!8.39](-1.1628,1.9457)rectangle(-1.0678,2.0042);
\fill[DeepBlue!8.95](-1.1628,2.0022)rectangle(-1.0678,2.0607);
\fill[DeepBlue!9.64](-1.1628,2.0587)rectangle(-1.0678,2.1172);
\fill[DeepBlue!10.23](-1.1628,2.1152)rectangle(-1.0678,2.1737);
\fill[DeepBlue!10.57](-1.1628,2.1717)rectangle(-1.0678,2.2303);
\fill[DeepBlue!10.57](-1.1628,2.2283)rectangle(-1.0678,2.2868);
\fill[DeepBlue!10.23](-1.1628,2.2848)rectangle(-1.0678,2.3433);
\fill[DeepBlue!9.64](-1.1628,2.3413)rectangle(-1.0678,2.3998);
\fill[DeepBlue!8.95](-1.1628,2.3978)rectangle(-1.0678,2.4563);
\fill[DeepBlue!8.39](-1.1628,2.4543)rectangle(-1.0678,2.5129);
\fill[DeepBlue!8.15](-1.1628,2.5109)rectangle(-1.0678,2.5694);
\fill[DeepBlue!8.31](-1.1628,2.5674)rectangle(-1.0678,2.6259);
\fill[DeepBlue!8.78](-1.1628,2.6239)rectangle(-1.0678,2.6824);
\fill[DeepBlue!9.34](-1.1628,2.6804)rectangle(-1.0678,2.7390);
\fill[DeepBlue!9.77](-1.1628,2.7370)rectangle(-1.0678,2.7955);
\fill[DeepBlue!9.90](-1.1628,2.7935)rectangle(-1.0678,2.8520);
\fill[DeepBlue!10.18](-1.0698,1.5500)rectangle(-0.9747,1.6085);
\fill[DeepBlue!10.12](-1.0698,1.6065)rectangle(-0.9747,1.6650);
\fill[DeepBlue!9.68](-1.0698,1.6630)rectangle(-0.9747,1.7216);
\fill[DeepBlue!9.08](-1.0698,1.7196)rectangle(-0.9747,1.7781);
\fill[DeepBlue!8.66](-1.0698,1.7761)rectangle(-0.9747,1.8346);
\fill[DeepBlue!8.78](-1.0698,1.8326)rectangle(-0.9747,1.8911);
\fill[DeepBlue!9.61](-1.0698,1.8891)rectangle(-0.9747,1.9477);
\fill[DeepBlue!10.99](-1.0698,1.9457)rectangle(-0.9747,2.0042);
\fill[DeepBlue!12.58](-1.0698,2.0022)rectangle(-0.9747,2.0607);
\fill[DeepBlue!14.06](-1.0698,2.0587)rectangle(-0.9747,2.1172);
\fill[DeepBlue!15.18](-1.0698,2.1152)rectangle(-0.9747,2.1737);
\fill[DeepBlue!15.78](-1.0698,2.1717)rectangle(-0.9747,2.2303);
\fill[DeepBlue!15.78](-1.0698,2.2283)rectangle(-0.9747,2.2868);
\fill[DeepBlue!15.18](-1.0698,2.2848)rectangle(-0.9747,2.3433);
\fill[DeepBlue!14.06](-1.0698,2.3413)rectangle(-0.9747,2.3998);
\fill[DeepBlue!12.58](-1.0698,2.3978)rectangle(-0.9747,2.4563);
\fill[DeepBlue!10.99](-1.0698,2.4543)rectangle(-0.9747,2.5129);
\fill[DeepBlue!9.61](-1.0698,2.5109)rectangle(-0.9747,2.5694);
\fill[DeepBlue!8.78](-1.0698,2.5674)rectangle(-0.9747,2.6259);
\fill[DeepBlue!8.66](-1.0698,2.6239)rectangle(-0.9747,2.6824);
\fill[DeepBlue!9.08](-1.0698,2.6804)rectangle(-0.9747,2.7390);
\fill[DeepBlue!9.68](-1.0698,2.7370)rectangle(-0.9747,2.7955);
\fill[DeepBlue!10.12](-1.0698,2.7935)rectangle(-0.9747,2.8520);
\fill[DeepBlue!10.53](-0.9767,1.5500)rectangle(-0.8817,1.6085);
\fill[DeepBlue!10.15](-0.9767,1.6065)rectangle(-0.8817,1.6650);
\fill[DeepBlue!9.50](-0.9767,1.6630)rectangle(-0.8817,1.7216);
\fill[DeepBlue!9.03](-0.9767,1.7196)rectangle(-0.8817,1.7781);
\fill[DeepBlue!9.30](-0.9767,1.7761)rectangle(-0.8817,1.8346);
\fill[DeepBlue!10.62](-0.9767,1.8326)rectangle(-0.8817,1.8911);
\fill[DeepBlue!12.80](-0.9767,1.8891)rectangle(-0.8817,1.9477);
\fill[DeepBlue!15.38](-0.9767,1.9457)rectangle(-0.8817,2.0042);
\fill[DeepBlue!17.93](-0.9767,2.0022)rectangle(-0.8817,2.0607);
\fill[DeepBlue!20.13](-0.9767,2.0587)rectangle(-0.8817,2.1172);
\fill[DeepBlue!21.73](-0.9767,2.1152)rectangle(-0.8817,2.1737);
\fill[DeepBlue!22.57](-0.9767,2.1717)rectangle(-0.8817,2.2303);
\fill[DeepBlue!22.57](-0.9767,2.2283)rectangle(-0.8817,2.2868);
\fill[DeepBlue!21.73](-0.9767,2.2848)rectangle(-0.8817,2.3433);
\fill[DeepBlue!20.13](-0.9767,2.3413)rectangle(-0.8817,2.3998);
\fill[DeepBlue!17.93](-0.9767,2.3978)rectangle(-0.8817,2.4563);
\fill[DeepBlue!15.38](-0.9767,2.4543)rectangle(-0.8817,2.5129);
\fill[DeepBlue!12.80](-0.9767,2.5109)rectangle(-0.8817,2.5694);
\fill[DeepBlue!10.62](-0.9767,2.5674)rectangle(-0.8817,2.6259);
\fill[DeepBlue!9.30](-0.9767,2.6239)rectangle(-0.8817,2.6824);
\fill[DeepBlue!9.03](-0.9767,2.6804)rectangle(-0.8817,2.7390);
\fill[DeepBlue!9.50](-0.9767,2.7370)rectangle(-0.8817,2.7955);
\fill[DeepBlue!10.15](-0.9767,2.7935)rectangle(-0.8817,2.8520);
\fill[DeepBlue!10.69](-0.8837,1.5500)rectangle(-0.7887,1.6085);
\fill[DeepBlue!10.02](-0.8837,1.6065)rectangle(-0.7887,1.6650);
\fill[DeepBlue!9.36](-0.8837,1.6630)rectangle(-0.7887,1.7216);
\fill[DeepBlue!9.47](-0.8837,1.7196)rectangle(-0.7887,1.7781);
\fill[DeepBlue!10.97](-0.8837,1.7761)rectangle(-0.7887,1.8346);
\fill[DeepBlue!13.74](-0.8837,1.8326)rectangle(-0.7887,1.8911);
\fill[DeepBlue!17.24](-0.8837,1.8891)rectangle(-0.7887,1.9477);
\fill[DeepBlue!20.93](-0.8837,1.9457)rectangle(-0.7887,2.0042);
\fill[DeepBlue!24.39](-0.8837,2.0022)rectangle(-0.7887,2.0607);
\fill[DeepBlue!27.27](-0.8837,2.0587)rectangle(-0.7887,2.1172);
\fill[DeepBlue!29.33](-0.8837,2.1152)rectangle(-0.7887,2.1737);
\fill[DeepBlue!30.40](-0.8837,2.1717)rectangle(-0.7887,2.2303);
\fill[DeepBlue!30.40](-0.8837,2.2283)rectangle(-0.7887,2.2868);
\fill[DeepBlue!29.33](-0.8837,2.2848)rectangle(-0.7887,2.3433);
\fill[DeepBlue!27.27](-0.8837,2.3413)rectangle(-0.7887,2.3998);
\fill[DeepBlue!24.39](-0.8837,2.3978)rectangle(-0.7887,2.4563);
\fill[DeepBlue!20.93](-0.8837,2.4543)rectangle(-0.7887,2.5129);
\fill[DeepBlue!17.24](-0.8837,2.5109)rectangle(-0.7887,2.5694);
\fill[DeepBlue!13.74](-0.8837,2.5674)rectangle(-0.7887,2.6259);
\fill[DeepBlue!10.97](-0.8837,2.6239)rectangle(-0.7887,2.6824);
\fill[DeepBlue!9.47](-0.8837,2.6804)rectangle(-0.7887,2.7390);
\fill[DeepBlue!9.36](-0.8837,2.7370)rectangle(-0.7887,2.7955);
\fill[DeepBlue!10.02](-0.8837,2.7935)rectangle(-0.7887,2.8520);
\fill[DeepBlue!10.64](-0.7907,1.5500)rectangle(-0.6957,1.6085);
\fill[DeepBlue!9.75](-0.7907,1.6065)rectangle(-0.6957,1.6650);
\fill[DeepBlue!9.37](-0.7907,1.6630)rectangle(-0.6957,1.7216);
\fill[DeepBlue!10.54](-0.7907,1.7196)rectangle(-0.6957,1.7781);
\fill[DeepBlue!13.52](-0.7907,1.7761)rectangle(-0.6957,1.8346);
\fill[DeepBlue!17.72](-0.7907,1.8326)rectangle(-0.6957,1.8911);
\fill[DeepBlue!22.46](-0.7907,1.8891)rectangle(-0.6957,1.9477);
\fill[DeepBlue!27.19](-0.7907,1.9457)rectangle(-0.6957,2.0042);
\fill[DeepBlue!31.50](-0.7907,2.0022)rectangle(-0.6957,2.0607);
\fill[DeepBlue!35.05](-0.7907,2.0587)rectangle(-0.6957,2.1172);
\fill[DeepBlue!37.57](-0.7907,2.1152)rectangle(-0.6957,2.1737);
\fill[DeepBlue!38.88](-0.7907,2.1717)rectangle(-0.6957,2.2303);
\fill[DeepBlue!38.88](-0.7907,2.2283)rectangle(-0.6957,2.2868);
\fill[DeepBlue!37.57](-0.7907,2.2848)rectangle(-0.6957,2.3433);
\fill[DeepBlue!35.05](-0.7907,2.3413)rectangle(-0.6957,2.3998);
\fill[DeepBlue!31.50](-0.7907,2.3978)rectangle(-0.6957,2.4563);
\fill[DeepBlue!27.19](-0.7907,2.4543)rectangle(-0.6957,2.5129);
\fill[DeepBlue!22.46](-0.7907,2.5109)rectangle(-0.6957,2.5694);
\fill[DeepBlue!17.72](-0.7907,2.5674)rectangle(-0.6957,2.6259);
\fill[DeepBlue!13.52](-0.7907,2.6239)rectangle(-0.6957,2.6824);
\fill[DeepBlue!10.54](-0.7907,2.6804)rectangle(-0.6957,2.7390);
\fill[DeepBlue!9.37](-0.7907,2.7370)rectangle(-0.6957,2.7955);
\fill[DeepBlue!9.75](-0.7907,2.7935)rectangle(-0.6957,2.8520);
\fill[DeepBlue!10.36](-0.6977,1.5500)rectangle(-0.6027,1.6085);
\fill[DeepBlue!9.36](-0.6977,1.6065)rectangle(-0.6027,1.6650);
\fill[DeepBlue!9.61](-0.6977,1.6630)rectangle(-0.6027,1.7216);
\fill[DeepBlue!12.16](-0.6977,1.7196)rectangle(-0.6027,1.7781);
\fill[DeepBlue!16.63](-0.6977,1.7761)rectangle(-0.6027,1.8346);
\fill[DeepBlue!22.16](-0.6977,1.8326)rectangle(-0.6027,1.8911);
\fill[DeepBlue!28.05](-0.6977,1.8891)rectangle(-0.6027,1.9477);
\fill[DeepBlue!33.78](-0.6977,1.9457)rectangle(-0.6027,2.0042);
\fill[DeepBlue!38.92](-0.6977,2.0022)rectangle(-0.6027,2.0607);
\fill[DeepBlue!43.12](-0.6977,2.0587)rectangle(-0.6027,2.1172);
\fill[DeepBlue!46.08](-0.6977,2.1152)rectangle(-0.6027,2.1737);
\fill[DeepBlue!47.62](-0.6977,2.1717)rectangle(-0.6027,2.2303);
\fill[DeepBlue!47.62](-0.6977,2.2283)rectangle(-0.6027,2.2868);
\fill[DeepBlue!46.08](-0.6977,2.2848)rectangle(-0.6027,2.3433);
\fill[DeepBlue!43.12](-0.6977,2.3413)rectangle(-0.6027,2.3998);
\fill[DeepBlue!38.92](-0.6977,2.3978)rectangle(-0.6027,2.4563);
\fill[DeepBlue!33.78](-0.6977,2.4543)rectangle(-0.6027,2.5129);
\fill[DeepBlue!28.05](-0.6977,2.5109)rectangle(-0.6027,2.5694);
\fill[DeepBlue!22.16](-0.6977,2.5674)rectangle(-0.6027,2.6259);
\fill[DeepBlue!16.63](-0.6977,2.6239)rectangle(-0.6027,2.6824);
\fill[DeepBlue!12.16](-0.6977,2.6804)rectangle(-0.6027,2.7390);
\fill[DeepBlue!9.61](-0.6977,2.7370)rectangle(-0.6027,2.7955);
\fill[DeepBlue!9.36](-0.6977,2.7935)rectangle(-0.6027,2.8520);
\fill[DeepBlue!9.83](-0.6047,1.5500)rectangle(-0.5096,1.6085);
\fill[DeepBlue!8.85](-0.6047,1.6065)rectangle(-0.5096,1.6650);
\fill[DeepBlue!10.07](-0.6047,1.6630)rectangle(-0.5096,1.7216);
\fill[DeepBlue!14.12](-0.6047,1.7196)rectangle(-0.5096,1.7781);
\fill[DeepBlue!19.99](-0.6047,1.7761)rectangle(-0.5096,1.8346);
\fill[DeepBlue!26.73](-0.6047,1.8326)rectangle(-0.5096,1.8911);
\fill[DeepBlue!33.69](-0.6047,1.8891)rectangle(-0.5096,1.9477);
\fill[DeepBlue!40.36](-0.6047,1.9457)rectangle(-0.5096,2.0042);
\fill[DeepBlue!46.30](-0.6047,2.0022)rectangle(-0.5096,2.0607);
\fill[DeepBlue!51.12](-0.6047,2.0587)rectangle(-0.5096,2.1172);
\fill[DeepBlue!54.51](-0.6047,2.1152)rectangle(-0.5096,2.1737);
\fill[DeepBlue!56.26](-0.6047,2.1717)rectangle(-0.5096,2.2303);
\fill[DeepBlue!56.26](-0.6047,2.2283)rectangle(-0.5096,2.2868);
\fill[DeepBlue!54.51](-0.6047,2.2848)rectangle(-0.5096,2.3433);
\fill[DeepBlue!51.12](-0.6047,2.3413)rectangle(-0.5096,2.3998);
\fill[DeepBlue!46.30](-0.6047,2.3978)rectangle(-0.5096,2.4563);
\fill[DeepBlue!40.36](-0.6047,2.4543)rectangle(-0.5096,2.5129);
\fill[DeepBlue!33.69](-0.6047,2.5109)rectangle(-0.5096,2.5694);
\fill[DeepBlue!26.73](-0.6047,2.5674)rectangle(-0.5096,2.6259);
\fill[DeepBlue!19.99](-0.6047,2.6239)rectangle(-0.5096,2.6824);
\fill[DeepBlue!14.12](-0.6047,2.6804)rectangle(-0.5096,2.7390);
\fill[DeepBlue!10.07](-0.6047,2.7370)rectangle(-0.5096,2.7955);
\fill[DeepBlue!8.85](-0.6047,2.7935)rectangle(-0.5096,2.8520);
\fill[DeepBlue!9.03](-0.5116,1.5500)rectangle(-0.4166,1.6085);
\fill[DeepBlue!8.22](-0.5116,1.6065)rectangle(-0.4166,1.6650);
\fill[DeepBlue!10.69](-0.5116,1.6630)rectangle(-0.4166,1.7216);
\fill[DeepBlue!16.20](-0.5116,1.7196)rectangle(-0.4166,1.7781);
\fill[DeepBlue!23.31](-0.5116,1.7761)rectangle(-0.4166,1.8346);
\fill[DeepBlue!31.16](-0.5116,1.8326)rectangle(-0.4166,1.8911);
\fill[DeepBlue!39.11](-0.5116,1.8891)rectangle(-0.4166,1.9477);
\fill[DeepBlue!46.65](-0.5116,1.9457)rectangle(-0.4166,2.0042);
\fill[DeepBlue!53.32](-0.5116,2.0022)rectangle(-0.4166,2.0607);
\fill[DeepBlue!58.72](-0.5116,2.0587)rectangle(-0.4166,2.1172);
\fill[DeepBlue!62.51](-0.5116,2.1152)rectangle(-0.4166,2.1737);
\fill[DeepBlue!64.46](-0.5116,2.1717)rectangle(-0.4166,2.2303);
\fill[DeepBlue!64.46](-0.5116,2.2283)rectangle(-0.4166,2.2868);
\fill[DeepBlue!62.51](-0.5116,2.2848)rectangle(-0.4166,2.3433);
\fill[DeepBlue!58.72](-0.5116,2.3413)rectangle(-0.4166,2.3998);
\fill[DeepBlue!53.32](-0.5116,2.3978)rectangle(-0.4166,2.4563);
\fill[DeepBlue!46.65](-0.5116,2.4543)rectangle(-0.4166,2.5129);
\fill[DeepBlue!39.11](-0.5116,2.5109)rectangle(-0.4166,2.5694);
\fill[DeepBlue!31.16](-0.5116,2.5674)rectangle(-0.4166,2.6259);
\fill[DeepBlue!23.31](-0.5116,2.6239)rectangle(-0.4166,2.6824);
\fill[DeepBlue!16.20](-0.5116,2.6804)rectangle(-0.4166,2.7390);
\fill[DeepBlue!10.69](-0.5116,2.7370)rectangle(-0.4166,2.7955);
\fill[DeepBlue!8.22](-0.5116,2.7935)rectangle(-0.4166,2.8520);
\fill[DeepBlue!7.98](-0.4186,1.5500)rectangle(-0.3236,1.6085);
\fill[DeepBlue!7.48](-0.4186,1.6065)rectangle(-0.3236,1.6650);
\fill[DeepBlue!11.37](-0.4186,1.6630)rectangle(-0.3236,1.7216);
\fill[DeepBlue!18.19](-0.4186,1.7196)rectangle(-0.3236,1.7781);
\fill[DeepBlue!26.41](-0.4186,1.7761)rectangle(-0.3236,1.8346);
\fill[DeepBlue!35.22](-0.4186,1.8326)rectangle(-0.3236,1.8911);
\fill[DeepBlue!44.04](-0.4186,1.8891)rectangle(-0.3236,1.9477);
\fill[DeepBlue!52.37](-0.4186,1.9457)rectangle(-0.3236,2.0042);
\fill[DeepBlue!59.69](-0.4186,2.0022)rectangle(-0.3236,2.0607);
\fill[DeepBlue!65.60](-0.4186,2.0587)rectangle(-0.3236,2.1172);
\fill[DeepBlue!69.75](-0.4186,2.1152)rectangle(-0.3236,2.1737);
\fill[DeepBlue!71.88](-0.4186,2.1717)rectangle(-0.3236,2.2303);
\fill[DeepBlue!71.88](-0.4186,2.2283)rectangle(-0.3236,2.2868);
\fill[DeepBlue!69.75](-0.4186,2.2848)rectangle(-0.3236,2.3433);
\fill[DeepBlue!65.60](-0.4186,2.3413)rectangle(-0.3236,2.3998);
\fill[DeepBlue!59.69](-0.4186,2.3978)rectangle(-0.3236,2.4563);
\fill[DeepBlue!52.37](-0.4186,2.4543)rectangle(-0.3236,2.5129);
\fill[DeepBlue!44.04](-0.4186,2.5109)rectangle(-0.3236,2.5694);
\fill[DeepBlue!35.22](-0.4186,2.5674)rectangle(-0.3236,2.6259);
\fill[DeepBlue!26.41](-0.4186,2.6239)rectangle(-0.3236,2.6824);
\fill[DeepBlue!18.19](-0.4186,2.6804)rectangle(-0.3236,2.7390);
\fill[DeepBlue!11.37](-0.4186,2.7370)rectangle(-0.3236,2.7955);
\fill[DeepBlue!7.48](-0.4186,2.7935)rectangle(-0.3236,2.8520);
\fill[DeepBlue!6.69](-0.3256,1.5500)rectangle(-0.2306,1.6085);
\fill[DeepBlue!6.67](-0.3256,1.6065)rectangle(-0.2306,1.6650);
\fill[DeepBlue!12.02](-0.3256,1.6630)rectangle(-0.2306,1.7216);
\fill[DeepBlue!19.96](-0.3256,1.7196)rectangle(-0.2306,1.7781);
\fill[DeepBlue!29.08](-0.3256,1.7761)rectangle(-0.2306,1.8346);
\fill[DeepBlue!38.71](-0.3256,1.8326)rectangle(-0.2306,1.8911);
\fill[DeepBlue!48.28](-0.3256,1.8891)rectangle(-0.2306,1.9477);
\fill[DeepBlue!57.26](-0.3256,1.9457)rectangle(-0.2306,2.0042);
\fill[DeepBlue!65.14](-0.3256,2.0022)rectangle(-0.2306,2.0607);
\fill[DeepBlue!71.48](-0.3256,2.0587)rectangle(-0.2306,2.1172);
\fill[DeepBlue!75.93](-0.3256,2.1152)rectangle(-0.2306,2.1737);
\fill[DeepBlue!78.22](-0.3256,2.1717)rectangle(-0.2306,2.2303);
\fill[DeepBlue!78.22](-0.3256,2.2283)rectangle(-0.2306,2.2868);
\fill[DeepBlue!75.93](-0.3256,2.2848)rectangle(-0.2306,2.3433);
\fill[DeepBlue!71.48](-0.3256,2.3413)rectangle(-0.2306,2.3998);
\fill[DeepBlue!65.14](-0.3256,2.3978)rectangle(-0.2306,2.4563);
\fill[DeepBlue!57.26](-0.3256,2.4543)rectangle(-0.2306,2.5129);
\fill[DeepBlue!48.28](-0.3256,2.5109)rectangle(-0.2306,2.5694);
\fill[DeepBlue!38.71](-0.3256,2.5674)rectangle(-0.2306,2.6259);
\fill[DeepBlue!29.08](-0.3256,2.6239)rectangle(-0.2306,2.6824);
\fill[DeepBlue!19.96](-0.3256,2.6804)rectangle(-0.2306,2.7390);
\fill[DeepBlue!12.02](-0.3256,2.7370)rectangle(-0.2306,2.7955);
\fill[DeepBlue!6.67](-0.3256,2.7935)rectangle(-0.2306,2.8520);
\fill[DeepBlue!5.22](-0.2326,1.5500)rectangle(-0.1375,1.6085);
\fill[DeepBlue!5.84](-0.2326,1.6065)rectangle(-0.1375,1.6650);
\fill[DeepBlue!12.56](-0.2326,1.6630)rectangle(-0.1375,1.7216);
\fill[DeepBlue!21.38](-0.2326,1.7196)rectangle(-0.1375,1.7781);
\fill[DeepBlue!31.21](-0.2326,1.7761)rectangle(-0.1375,1.8346);
\fill[DeepBlue!41.48](-0.2326,1.8326)rectangle(-0.1375,1.8911);
\fill[DeepBlue!51.63](-0.2326,1.8891)rectangle(-0.1375,1.9477);
\fill[DeepBlue!61.12](-0.2326,1.9457)rectangle(-0.1375,2.0042);
\fill[DeepBlue!69.43](-0.2326,2.0022)rectangle(-0.1375,2.0607);
\fill[DeepBlue!76.12](-0.2326,2.0587)rectangle(-0.1375,2.1172);
\fill[DeepBlue!80.80](-0.2326,2.1152)rectangle(-0.1375,2.1737);
\fill[DeepBlue!83.21](-0.2326,2.1717)rectangle(-0.1375,2.2303);
\fill[DeepBlue!83.21](-0.2326,2.2283)rectangle(-0.1375,2.2868);
\fill[DeepBlue!80.80](-0.2326,2.2848)rectangle(-0.1375,2.3433);
\fill[DeepBlue!76.12](-0.2326,2.3413)rectangle(-0.1375,2.3998);
\fill[DeepBlue!69.43](-0.2326,2.3978)rectangle(-0.1375,2.4563);
\fill[DeepBlue!61.12](-0.2326,2.4543)rectangle(-0.1375,2.5129);
\fill[DeepBlue!51.63](-0.2326,2.5109)rectangle(-0.1375,2.5694);
\fill[DeepBlue!41.48](-0.2326,2.5674)rectangle(-0.1375,2.6259);
\fill[DeepBlue!31.21](-0.2326,2.6239)rectangle(-0.1375,2.6824);
\fill[DeepBlue!21.38](-0.2326,2.6804)rectangle(-0.1375,2.7390);
\fill[DeepBlue!12.56](-0.2326,2.7370)rectangle(-0.1375,2.7955);
\fill[DeepBlue!5.84](-0.2326,2.7935)rectangle(-0.1375,2.8520);
\fill[DeepBlue!3.72](-0.1395,1.5500)rectangle(-0.0445,1.6085);
\fill[DeepBlue!5.13](-0.1395,1.6065)rectangle(-0.0445,1.6650);
\fill[DeepBlue!12.94](-0.1395,1.6630)rectangle(-0.0445,1.7216);
\fill[DeepBlue!22.37](-0.1395,1.7196)rectangle(-0.0445,1.7781);
\fill[DeepBlue!32.69](-0.1395,1.7761)rectangle(-0.0445,1.8346);
\fill[DeepBlue!43.39](-0.1395,1.8326)rectangle(-0.0445,1.8911);
\fill[DeepBlue!53.94](-0.1395,1.8891)rectangle(-0.0445,1.9477);
\fill[DeepBlue!63.78](-0.1395,1.9457)rectangle(-0.0445,2.0042);
\fill[DeepBlue!72.40](-0.1395,2.0022)rectangle(-0.0445,2.0607);
\fill[DeepBlue!79.32](-0.1395,2.0587)rectangle(-0.0445,2.1172);
\fill[DeepBlue!84.17](-0.1395,2.1152)rectangle(-0.0445,2.1737);
\fill[DeepBlue!86.66](-0.1395,2.1717)rectangle(-0.0445,2.2303);
\fill[DeepBlue!86.66](-0.1395,2.2283)rectangle(-0.0445,2.2868);
\fill[DeepBlue!84.17](-0.1395,2.2848)rectangle(-0.0445,2.3433);
\fill[DeepBlue!79.32](-0.1395,2.3413)rectangle(-0.0445,2.3998);
\fill[DeepBlue!72.40](-0.1395,2.3978)rectangle(-0.0445,2.4563);
\fill[DeepBlue!63.78](-0.1395,2.4543)rectangle(-0.0445,2.5129);
\fill[DeepBlue!53.94](-0.1395,2.5109)rectangle(-0.0445,2.5694);
\fill[DeepBlue!43.39](-0.1395,2.5674)rectangle(-0.0445,2.6259);
\fill[DeepBlue!32.69](-0.1395,2.6239)rectangle(-0.0445,2.6824);
\fill[DeepBlue!22.37](-0.1395,2.6804)rectangle(-0.0445,2.7390);
\fill[DeepBlue!12.94](-0.1395,2.7370)rectangle(-0.0445,2.7955);
\fill[DeepBlue!5.13](-0.1395,2.7935)rectangle(-0.0445,2.8520);
\fill[DeepBlue!2.56](-0.0465,1.5500)rectangle(0.0485,1.6085);
\fill[DeepBlue!4.69](-0.0465,1.6065)rectangle(0.0485,1.6650);
\fill[DeepBlue!13.14](-0.0465,1.6630)rectangle(0.0485,1.7216);
\fill[DeepBlue!22.87](-0.0465,1.7196)rectangle(0.0485,1.7781);
\fill[DeepBlue!33.44](-0.0465,1.7761)rectangle(0.0485,1.8346);
\fill[DeepBlue!44.37](-0.0465,1.8326)rectangle(0.0485,1.8911);
\fill[DeepBlue!55.12](-0.0465,1.8891)rectangle(0.0485,1.9477);
\fill[DeepBlue!65.15](-0.0465,1.9457)rectangle(0.0485,2.0042);
\fill[DeepBlue!73.92](-0.0465,2.0022)rectangle(0.0485,2.0607);
\fill[DeepBlue!80.96](-0.0465,2.0587)rectangle(0.0485,2.1172);
\fill[DeepBlue!85.88](-0.0465,2.1152)rectangle(0.0485,2.1737);
\fill[DeepBlue!88.42](-0.0465,2.1717)rectangle(0.0485,2.2303);
\fill[DeepBlue!88.42](-0.0465,2.2283)rectangle(0.0485,2.2868);
\fill[DeepBlue!85.88](-0.0465,2.2848)rectangle(0.0485,2.3433);
\fill[DeepBlue!80.96](-0.0465,2.3413)rectangle(0.0485,2.3998);
\fill[DeepBlue!73.92](-0.0465,2.3978)rectangle(0.0485,2.4563);
\fill[DeepBlue!65.15](-0.0465,2.4543)rectangle(0.0485,2.5129);
\fill[DeepBlue!55.12](-0.0465,2.5109)rectangle(0.0485,2.5694);
\fill[DeepBlue!44.37](-0.0465,2.5674)rectangle(0.0485,2.6259);
\fill[DeepBlue!33.44](-0.0465,2.6239)rectangle(0.0485,2.6824);
\fill[DeepBlue!22.87](-0.0465,2.6804)rectangle(0.0485,2.7390);
\fill[DeepBlue!13.14](-0.0465,2.7370)rectangle(0.0485,2.7955);
\fill[DeepBlue!4.69](-0.0465,2.7935)rectangle(0.0485,2.8520);
\fill[DeepBlue!2.56](0.0465,1.5500)rectangle(0.1415,1.6085);
\fill[DeepBlue!4.69](0.0465,1.6065)rectangle(0.1415,1.6650);
\fill[DeepBlue!13.14](0.0465,1.6630)rectangle(0.1415,1.7216);
\fill[DeepBlue!22.87](0.0465,1.7196)rectangle(0.1415,1.7781);
\fill[DeepBlue!33.44](0.0465,1.7761)rectangle(0.1415,1.8346);
\fill[DeepBlue!44.37](0.0465,1.8326)rectangle(0.1415,1.8911);
\fill[DeepBlue!55.12](0.0465,1.8891)rectangle(0.1415,1.9477);
\fill[DeepBlue!65.15](0.0465,1.9457)rectangle(0.1415,2.0042);
\fill[DeepBlue!73.92](0.0465,2.0022)rectangle(0.1415,2.0607);
\fill[DeepBlue!80.96](0.0465,2.0587)rectangle(0.1415,2.1172);
\fill[DeepBlue!85.88](0.0465,2.1152)rectangle(0.1415,2.1737);
\fill[DeepBlue!88.42](0.0465,2.1717)rectangle(0.1415,2.2303);
\fill[DeepBlue!88.42](0.0465,2.2283)rectangle(0.1415,2.2868);
\fill[DeepBlue!85.88](0.0465,2.2848)rectangle(0.1415,2.3433);
\fill[DeepBlue!80.96](0.0465,2.3413)rectangle(0.1415,2.3998);
\fill[DeepBlue!73.92](0.0465,2.3978)rectangle(0.1415,2.4563);
\fill[DeepBlue!65.15](0.0465,2.4543)rectangle(0.1415,2.5129);
\fill[DeepBlue!55.12](0.0465,2.5109)rectangle(0.1415,2.5694);
\fill[DeepBlue!44.37](0.0465,2.5674)rectangle(0.1415,2.6259);
\fill[DeepBlue!33.44](0.0465,2.6239)rectangle(0.1415,2.6824);
\fill[DeepBlue!22.87](0.0465,2.6804)rectangle(0.1415,2.7390);
\fill[DeepBlue!13.14](0.0465,2.7370)rectangle(0.1415,2.7955);
\fill[DeepBlue!4.69](0.0465,2.7935)rectangle(0.1415,2.8520);
\fill[DeepBlue!3.72](0.1395,1.5500)rectangle(0.2346,1.6085);
\fill[DeepBlue!5.13](0.1395,1.6065)rectangle(0.2346,1.6650);
\fill[DeepBlue!12.94](0.1395,1.6630)rectangle(0.2346,1.7216);
\fill[DeepBlue!22.37](0.1395,1.7196)rectangle(0.2346,1.7781);
\fill[DeepBlue!32.69](0.1395,1.7761)rectangle(0.2346,1.8346);
\fill[DeepBlue!43.39](0.1395,1.8326)rectangle(0.2346,1.8911);
\fill[DeepBlue!53.94](0.1395,1.8891)rectangle(0.2346,1.9477);
\fill[DeepBlue!63.78](0.1395,1.9457)rectangle(0.2346,2.0042);
\fill[DeepBlue!72.40](0.1395,2.0022)rectangle(0.2346,2.0607);
\fill[DeepBlue!79.32](0.1395,2.0587)rectangle(0.2346,2.1172);
\fill[DeepBlue!84.17](0.1395,2.1152)rectangle(0.2346,2.1737);
\fill[DeepBlue!86.66](0.1395,2.1717)rectangle(0.2346,2.2303);
\fill[DeepBlue!86.66](0.1395,2.2283)rectangle(0.2346,2.2868);
\fill[DeepBlue!84.17](0.1395,2.2848)rectangle(0.2346,2.3433);
\fill[DeepBlue!79.32](0.1395,2.3413)rectangle(0.2346,2.3998);
\fill[DeepBlue!72.40](0.1395,2.3978)rectangle(0.2346,2.4563);
\fill[DeepBlue!63.78](0.1395,2.4543)rectangle(0.2346,2.5129);
\fill[DeepBlue!53.94](0.1395,2.5109)rectangle(0.2346,2.5694);
\fill[DeepBlue!43.39](0.1395,2.5674)rectangle(0.2346,2.6259);
\fill[DeepBlue!32.69](0.1395,2.6239)rectangle(0.2346,2.6824);
\fill[DeepBlue!22.37](0.1395,2.6804)rectangle(0.2346,2.7390);
\fill[DeepBlue!12.94](0.1395,2.7370)rectangle(0.2346,2.7955);
\fill[DeepBlue!5.13](0.1395,2.7935)rectangle(0.2346,2.8520);
\fill[DeepBlue!5.22](0.2326,1.5500)rectangle(0.3276,1.6085);
\fill[DeepBlue!5.84](0.2326,1.6065)rectangle(0.3276,1.6650);
\fill[DeepBlue!12.56](0.2326,1.6630)rectangle(0.3276,1.7216);
\fill[DeepBlue!21.38](0.2326,1.7196)rectangle(0.3276,1.7781);
\fill[DeepBlue!31.21](0.2326,1.7761)rectangle(0.3276,1.8346);
\fill[DeepBlue!41.48](0.2326,1.8326)rectangle(0.3276,1.8911);
\fill[DeepBlue!51.63](0.2326,1.8891)rectangle(0.3276,1.9477);
\fill[DeepBlue!61.12](0.2326,1.9457)rectangle(0.3276,2.0042);
\fill[DeepBlue!69.43](0.2326,2.0022)rectangle(0.3276,2.0607);
\fill[DeepBlue!76.12](0.2326,2.0587)rectangle(0.3276,2.1172);
\fill[DeepBlue!80.80](0.2326,2.1152)rectangle(0.3276,2.1737);
\fill[DeepBlue!83.21](0.2326,2.1717)rectangle(0.3276,2.2303);
\fill[DeepBlue!83.21](0.2326,2.2283)rectangle(0.3276,2.2868);
\fill[DeepBlue!80.80](0.2326,2.2848)rectangle(0.3276,2.3433);
\fill[DeepBlue!76.12](0.2326,2.3413)rectangle(0.3276,2.3998);
\fill[DeepBlue!69.43](0.2326,2.3978)rectangle(0.3276,2.4563);
\fill[DeepBlue!61.12](0.2326,2.4543)rectangle(0.3276,2.5129);
\fill[DeepBlue!51.63](0.2326,2.5109)rectangle(0.3276,2.5694);
\fill[DeepBlue!41.48](0.2326,2.5674)rectangle(0.3276,2.6259);
\fill[DeepBlue!31.21](0.2326,2.6239)rectangle(0.3276,2.6824);
\fill[DeepBlue!21.38](0.2326,2.6804)rectangle(0.3276,2.7390);
\fill[DeepBlue!12.56](0.2326,2.7370)rectangle(0.3276,2.7955);
\fill[DeepBlue!5.84](0.2326,2.7935)rectangle(0.3276,2.8520);
\fill[DeepBlue!6.69](0.3256,1.5500)rectangle(0.4206,1.6085);
\fill[DeepBlue!6.67](0.3256,1.6065)rectangle(0.4206,1.6650);
\fill[DeepBlue!12.02](0.3256,1.6630)rectangle(0.4206,1.7216);
\fill[DeepBlue!19.96](0.3256,1.7196)rectangle(0.4206,1.7781);
\fill[DeepBlue!29.08](0.3256,1.7761)rectangle(0.4206,1.8346);
\fill[DeepBlue!38.71](0.3256,1.8326)rectangle(0.4206,1.8911);
\fill[DeepBlue!48.28](0.3256,1.8891)rectangle(0.4206,1.9477);
\fill[DeepBlue!57.26](0.3256,1.9457)rectangle(0.4206,2.0042);
\fill[DeepBlue!65.14](0.3256,2.0022)rectangle(0.4206,2.0607);
\fill[DeepBlue!71.48](0.3256,2.0587)rectangle(0.4206,2.1172);
\fill[DeepBlue!75.93](0.3256,2.1152)rectangle(0.4206,2.1737);
\fill[DeepBlue!78.22](0.3256,2.1717)rectangle(0.4206,2.2303);
\fill[DeepBlue!78.22](0.3256,2.2283)rectangle(0.4206,2.2868);
\fill[DeepBlue!75.93](0.3256,2.2848)rectangle(0.4206,2.3433);
\fill[DeepBlue!71.48](0.3256,2.3413)rectangle(0.4206,2.3998);
\fill[DeepBlue!65.14](0.3256,2.3978)rectangle(0.4206,2.4563);
\fill[DeepBlue!57.26](0.3256,2.4543)rectangle(0.4206,2.5129);
\fill[DeepBlue!48.28](0.3256,2.5109)rectangle(0.4206,2.5694);
\fill[DeepBlue!38.71](0.3256,2.5674)rectangle(0.4206,2.6259);
\fill[DeepBlue!29.08](0.3256,2.6239)rectangle(0.4206,2.6824);
\fill[DeepBlue!19.96](0.3256,2.6804)rectangle(0.4206,2.7390);
\fill[DeepBlue!12.02](0.3256,2.7370)rectangle(0.4206,2.7955);
\fill[DeepBlue!6.67](0.3256,2.7935)rectangle(0.4206,2.8520);
\fill[DeepBlue!7.98](0.4186,1.5500)rectangle(0.5136,1.6085);
\fill[DeepBlue!7.48](0.4186,1.6065)rectangle(0.5136,1.6650);
\fill[DeepBlue!11.37](0.4186,1.6630)rectangle(0.5136,1.7216);
\fill[DeepBlue!18.19](0.4186,1.7196)rectangle(0.5136,1.7781);
\fill[DeepBlue!26.41](0.4186,1.7761)rectangle(0.5136,1.8346);
\fill[DeepBlue!35.22](0.4186,1.8326)rectangle(0.5136,1.8911);
\fill[DeepBlue!44.04](0.4186,1.8891)rectangle(0.5136,1.9477);
\fill[DeepBlue!52.37](0.4186,1.9457)rectangle(0.5136,2.0042);
\fill[DeepBlue!59.69](0.4186,2.0022)rectangle(0.5136,2.0607);
\fill[DeepBlue!65.60](0.4186,2.0587)rectangle(0.5136,2.1172);
\fill[DeepBlue!69.75](0.4186,2.1152)rectangle(0.5136,2.1737);
\fill[DeepBlue!71.88](0.4186,2.1717)rectangle(0.5136,2.2303);
\fill[DeepBlue!71.88](0.4186,2.2283)rectangle(0.5136,2.2868);
\fill[DeepBlue!69.75](0.4186,2.2848)rectangle(0.5136,2.3433);
\fill[DeepBlue!65.60](0.4186,2.3413)rectangle(0.5136,2.3998);
\fill[DeepBlue!59.69](0.4186,2.3978)rectangle(0.5136,2.4563);
\fill[DeepBlue!52.37](0.4186,2.4543)rectangle(0.5136,2.5129);
\fill[DeepBlue!44.04](0.4186,2.5109)rectangle(0.5136,2.5694);
\fill[DeepBlue!35.22](0.4186,2.5674)rectangle(0.5136,2.6259);
\fill[DeepBlue!26.41](0.4186,2.6239)rectangle(0.5136,2.6824);
\fill[DeepBlue!18.19](0.4186,2.6804)rectangle(0.5136,2.7390);
\fill[DeepBlue!11.37](0.4186,2.7370)rectangle(0.5136,2.7955);
\fill[DeepBlue!7.48](0.4186,2.7935)rectangle(0.5136,2.8520);
\fill[DeepBlue!9.03](0.5116,1.5500)rectangle(0.6067,1.6085);
\fill[DeepBlue!8.22](0.5116,1.6065)rectangle(0.6067,1.6650);
\fill[DeepBlue!10.69](0.5116,1.6630)rectangle(0.6067,1.7216);
\fill[DeepBlue!16.20](0.5116,1.7196)rectangle(0.6067,1.7781);
\fill[DeepBlue!23.31](0.5116,1.7761)rectangle(0.6067,1.8346);
\fill[DeepBlue!31.16](0.5116,1.8326)rectangle(0.6067,1.8911);
\fill[DeepBlue!39.11](0.5116,1.8891)rectangle(0.6067,1.9477);
\fill[DeepBlue!46.65](0.5116,1.9457)rectangle(0.6067,2.0042);
\fill[DeepBlue!53.32](0.5116,2.0022)rectangle(0.6067,2.0607);
\fill[DeepBlue!58.72](0.5116,2.0587)rectangle(0.6067,2.1172);
\fill[DeepBlue!62.51](0.5116,2.1152)rectangle(0.6067,2.1737);
\fill[DeepBlue!64.46](0.5116,2.1717)rectangle(0.6067,2.2303);
\fill[DeepBlue!64.46](0.5116,2.2283)rectangle(0.6067,2.2868);
\fill[DeepBlue!62.51](0.5116,2.2848)rectangle(0.6067,2.3433);
\fill[DeepBlue!58.72](0.5116,2.3413)rectangle(0.6067,2.3998);
\fill[DeepBlue!53.32](0.5116,2.3978)rectangle(0.6067,2.4563);
\fill[DeepBlue!46.65](0.5116,2.4543)rectangle(0.6067,2.5129);
\fill[DeepBlue!39.11](0.5116,2.5109)rectangle(0.6067,2.5694);
\fill[DeepBlue!31.16](0.5116,2.5674)rectangle(0.6067,2.6259);
\fill[DeepBlue!23.31](0.5116,2.6239)rectangle(0.6067,2.6824);
\fill[DeepBlue!16.20](0.5116,2.6804)rectangle(0.6067,2.7390);
\fill[DeepBlue!10.69](0.5116,2.7370)rectangle(0.6067,2.7955);
\fill[DeepBlue!8.22](0.5116,2.7935)rectangle(0.6067,2.8520);
\fill[DeepBlue!9.83](0.6047,1.5500)rectangle(0.6997,1.6085);
\fill[DeepBlue!8.85](0.6047,1.6065)rectangle(0.6997,1.6650);
\fill[DeepBlue!10.07](0.6047,1.6630)rectangle(0.6997,1.7216);
\fill[DeepBlue!14.12](0.6047,1.7196)rectangle(0.6997,1.7781);
\fill[DeepBlue!19.99](0.6047,1.7761)rectangle(0.6997,1.8346);
\fill[DeepBlue!26.73](0.6047,1.8326)rectangle(0.6997,1.8911);
\fill[DeepBlue!33.69](0.6047,1.8891)rectangle(0.6997,1.9477);
\fill[DeepBlue!40.36](0.6047,1.9457)rectangle(0.6997,2.0042);
\fill[DeepBlue!46.30](0.6047,2.0022)rectangle(0.6997,2.0607);
\fill[DeepBlue!51.12](0.6047,2.0587)rectangle(0.6997,2.1172);
\fill[DeepBlue!54.51](0.6047,2.1152)rectangle(0.6997,2.1737);
\fill[DeepBlue!56.26](0.6047,2.1717)rectangle(0.6997,2.2303);
\fill[DeepBlue!56.26](0.6047,2.2283)rectangle(0.6997,2.2868);
\fill[DeepBlue!54.51](0.6047,2.2848)rectangle(0.6997,2.3433);
\fill[DeepBlue!51.12](0.6047,2.3413)rectangle(0.6997,2.3998);
\fill[DeepBlue!46.30](0.6047,2.3978)rectangle(0.6997,2.4563);
\fill[DeepBlue!40.36](0.6047,2.4543)rectangle(0.6997,2.5129);
\fill[DeepBlue!33.69](0.6047,2.5109)rectangle(0.6997,2.5694);
\fill[DeepBlue!26.73](0.6047,2.5674)rectangle(0.6997,2.6259);
\fill[DeepBlue!19.99](0.6047,2.6239)rectangle(0.6997,2.6824);
\fill[DeepBlue!14.12](0.6047,2.6804)rectangle(0.6997,2.7390);
\fill[DeepBlue!10.07](0.6047,2.7370)rectangle(0.6997,2.7955);
\fill[DeepBlue!8.85](0.6047,2.7935)rectangle(0.6997,2.8520);
\fill[DeepBlue!10.36](0.6977,1.5500)rectangle(0.7927,1.6085);
\fill[DeepBlue!9.36](0.6977,1.6065)rectangle(0.7927,1.6650);
\fill[DeepBlue!9.61](0.6977,1.6630)rectangle(0.7927,1.7216);
\fill[DeepBlue!12.16](0.6977,1.7196)rectangle(0.7927,1.7781);
\fill[DeepBlue!16.63](0.6977,1.7761)rectangle(0.7927,1.8346);
\fill[DeepBlue!22.16](0.6977,1.8326)rectangle(0.7927,1.8911);
\fill[DeepBlue!28.05](0.6977,1.8891)rectangle(0.7927,1.9477);
\fill[DeepBlue!33.78](0.6977,1.9457)rectangle(0.7927,2.0042);
\fill[DeepBlue!38.92](0.6977,2.0022)rectangle(0.7927,2.0607);
\fill[DeepBlue!43.12](0.6977,2.0587)rectangle(0.7927,2.1172);
\fill[DeepBlue!46.08](0.6977,2.1152)rectangle(0.7927,2.1737);
\fill[DeepBlue!47.62](0.6977,2.1717)rectangle(0.7927,2.2303);
\fill[DeepBlue!47.62](0.6977,2.2283)rectangle(0.7927,2.2868);
\fill[DeepBlue!46.08](0.6977,2.2848)rectangle(0.7927,2.3433);
\fill[DeepBlue!43.12](0.6977,2.3413)rectangle(0.7927,2.3998);
\fill[DeepBlue!38.92](0.6977,2.3978)rectangle(0.7927,2.4563);
\fill[DeepBlue!33.78](0.6977,2.4543)rectangle(0.7927,2.5129);
\fill[DeepBlue!28.05](0.6977,2.5109)rectangle(0.7927,2.5694);
\fill[DeepBlue!22.16](0.6977,2.5674)rectangle(0.7927,2.6259);
\fill[DeepBlue!16.63](0.6977,2.6239)rectangle(0.7927,2.6824);
\fill[DeepBlue!12.16](0.6977,2.6804)rectangle(0.7927,2.7390);
\fill[DeepBlue!9.61](0.6977,2.7370)rectangle(0.7927,2.7955);
\fill[DeepBlue!9.36](0.6977,2.7935)rectangle(0.7927,2.8520);
\fill[DeepBlue!10.64](0.7907,1.5500)rectangle(0.8857,1.6085);
\fill[DeepBlue!9.75](0.7907,1.6065)rectangle(0.8857,1.6650);
\fill[DeepBlue!9.37](0.7907,1.6630)rectangle(0.8857,1.7216);
\fill[DeepBlue!10.54](0.7907,1.7196)rectangle(0.8857,1.7781);
\fill[DeepBlue!13.52](0.7907,1.7761)rectangle(0.8857,1.8346);
\fill[DeepBlue!17.72](0.7907,1.8326)rectangle(0.8857,1.8911);
\fill[DeepBlue!22.46](0.7907,1.8891)rectangle(0.8857,1.9477);
\fill[DeepBlue!27.19](0.7907,1.9457)rectangle(0.8857,2.0042);
\fill[DeepBlue!31.50](0.7907,2.0022)rectangle(0.8857,2.0607);
\fill[DeepBlue!35.05](0.7907,2.0587)rectangle(0.8857,2.1172);
\fill[DeepBlue!37.57](0.7907,2.1152)rectangle(0.8857,2.1737);
\fill[DeepBlue!38.88](0.7907,2.1717)rectangle(0.8857,2.2303);
\fill[DeepBlue!38.88](0.7907,2.2283)rectangle(0.8857,2.2868);
\fill[DeepBlue!37.57](0.7907,2.2848)rectangle(0.8857,2.3433);
\fill[DeepBlue!35.05](0.7907,2.3413)rectangle(0.8857,2.3998);
\fill[DeepBlue!31.50](0.7907,2.3978)rectangle(0.8857,2.4563);
\fill[DeepBlue!27.19](0.7907,2.4543)rectangle(0.8857,2.5129);
\fill[DeepBlue!22.46](0.7907,2.5109)rectangle(0.8857,2.5694);
\fill[DeepBlue!17.72](0.7907,2.5674)rectangle(0.8857,2.6259);
\fill[DeepBlue!13.52](0.7907,2.6239)rectangle(0.8857,2.6824);
\fill[DeepBlue!10.54](0.7907,2.6804)rectangle(0.8857,2.7390);
\fill[DeepBlue!9.37](0.7907,2.7370)rectangle(0.8857,2.7955);
\fill[DeepBlue!9.75](0.7907,2.7935)rectangle(0.8857,2.8520);
\fill[DeepBlue!10.69](0.8837,1.5500)rectangle(0.9787,1.6085);
\fill[DeepBlue!10.02](0.8837,1.6065)rectangle(0.9787,1.6650);
\fill[DeepBlue!9.36](0.8837,1.6630)rectangle(0.9787,1.7216);
\fill[DeepBlue!9.47](0.8837,1.7196)rectangle(0.9787,1.7781);
\fill[DeepBlue!10.97](0.8837,1.7761)rectangle(0.9787,1.8346);
\fill[DeepBlue!13.74](0.8837,1.8326)rectangle(0.9787,1.8911);
\fill[DeepBlue!17.24](0.8837,1.8891)rectangle(0.9787,1.9477);
\fill[DeepBlue!20.93](0.8837,1.9457)rectangle(0.9787,2.0042);
\fill[DeepBlue!24.39](0.8837,2.0022)rectangle(0.9787,2.0607);
\fill[DeepBlue!27.27](0.8837,2.0587)rectangle(0.9787,2.1172);
\fill[DeepBlue!29.33](0.8837,2.1152)rectangle(0.9787,2.1737);
\fill[DeepBlue!30.40](0.8837,2.1717)rectangle(0.9787,2.2303);
\fill[DeepBlue!30.40](0.8837,2.2283)rectangle(0.9787,2.2868);
\fill[DeepBlue!29.33](0.8837,2.2848)rectangle(0.9787,2.3433);
\fill[DeepBlue!27.27](0.8837,2.3413)rectangle(0.9787,2.3998);
\fill[DeepBlue!24.39](0.8837,2.3978)rectangle(0.9787,2.4563);
\fill[DeepBlue!20.93](0.8837,2.4543)rectangle(0.9787,2.5129);
\fill[DeepBlue!17.24](0.8837,2.5109)rectangle(0.9787,2.5694);
\fill[DeepBlue!13.74](0.8837,2.5674)rectangle(0.9787,2.6259);
\fill[DeepBlue!10.97](0.8837,2.6239)rectangle(0.9787,2.6824);
\fill[DeepBlue!9.47](0.8837,2.6804)rectangle(0.9787,2.7390);
\fill[DeepBlue!9.36](0.8837,2.7370)rectangle(0.9787,2.7955);
\fill[DeepBlue!10.02](0.8837,2.7935)rectangle(0.9787,2.8520);
\fill[DeepBlue!10.53](0.9767,1.5500)rectangle(1.0718,1.6085);
\fill[DeepBlue!10.15](0.9767,1.6065)rectangle(1.0718,1.6650);
\fill[DeepBlue!9.50](0.9767,1.6630)rectangle(1.0718,1.7216);
\fill[DeepBlue!9.03](0.9767,1.7196)rectangle(1.0718,1.7781);
\fill[DeepBlue!9.30](0.9767,1.7761)rectangle(1.0718,1.8346);
\fill[DeepBlue!10.62](0.9767,1.8326)rectangle(1.0718,1.8911);
\fill[DeepBlue!12.80](0.9767,1.8891)rectangle(1.0718,1.9477);
\fill[DeepBlue!15.38](0.9767,1.9457)rectangle(1.0718,2.0042);
\fill[DeepBlue!17.93](0.9767,2.0022)rectangle(1.0718,2.0607);
\fill[DeepBlue!20.13](0.9767,2.0587)rectangle(1.0718,2.1172);
\fill[DeepBlue!21.73](0.9767,2.1152)rectangle(1.0718,2.1737);
\fill[DeepBlue!22.57](0.9767,2.1717)rectangle(1.0718,2.2303);
\fill[DeepBlue!22.57](0.9767,2.2283)rectangle(1.0718,2.2868);
\fill[DeepBlue!21.73](0.9767,2.2848)rectangle(1.0718,2.3433);
\fill[DeepBlue!20.13](0.9767,2.3413)rectangle(1.0718,2.3998);
\fill[DeepBlue!17.93](0.9767,2.3978)rectangle(1.0718,2.4563);
\fill[DeepBlue!15.38](0.9767,2.4543)rectangle(1.0718,2.5129);
\fill[DeepBlue!12.80](0.9767,2.5109)rectangle(1.0718,2.5694);
\fill[DeepBlue!10.62](0.9767,2.5674)rectangle(1.0718,2.6259);
\fill[DeepBlue!9.30](0.9767,2.6239)rectangle(1.0718,2.6824);
\fill[DeepBlue!9.03](0.9767,2.6804)rectangle(1.0718,2.7390);
\fill[DeepBlue!9.50](0.9767,2.7370)rectangle(1.0718,2.7955);
\fill[DeepBlue!10.15](0.9767,2.7935)rectangle(1.0718,2.8520);
\fill[DeepBlue!10.18](1.0698,1.5500)rectangle(1.1648,1.6085);
\fill[DeepBlue!10.12](1.0698,1.6065)rectangle(1.1648,1.6650);
\fill[DeepBlue!9.68](1.0698,1.6630)rectangle(1.1648,1.7216);
\fill[DeepBlue!9.08](1.0698,1.7196)rectangle(1.1648,1.7781);
\fill[DeepBlue!8.66](1.0698,1.7761)rectangle(1.1648,1.8346);
\fill[DeepBlue!8.78](1.0698,1.8326)rectangle(1.1648,1.8911);
\fill[DeepBlue!9.61](1.0698,1.8891)rectangle(1.1648,1.9477);
\fill[DeepBlue!10.99](1.0698,1.9457)rectangle(1.1648,2.0042);
\fill[DeepBlue!12.58](1.0698,2.0022)rectangle(1.1648,2.0607);
\fill[DeepBlue!14.06](1.0698,2.0587)rectangle(1.1648,2.1172);
\fill[DeepBlue!15.18](1.0698,2.1152)rectangle(1.1648,2.1737);
\fill[DeepBlue!15.78](1.0698,2.1717)rectangle(1.1648,2.2303);
\fill[DeepBlue!15.78](1.0698,2.2283)rectangle(1.1648,2.2868);
\fill[DeepBlue!15.18](1.0698,2.2848)rectangle(1.1648,2.3433);
\fill[DeepBlue!14.06](1.0698,2.3413)rectangle(1.1648,2.3998);
\fill[DeepBlue!12.58](1.0698,2.3978)rectangle(1.1648,2.4563);
\fill[DeepBlue!10.99](1.0698,2.4543)rectangle(1.1648,2.5129);
\fill[DeepBlue!9.61](1.0698,2.5109)rectangle(1.1648,2.5694);
\fill[DeepBlue!8.78](1.0698,2.5674)rectangle(1.1648,2.6259);
\fill[DeepBlue!8.66](1.0698,2.6239)rectangle(1.1648,2.6824);
\fill[DeepBlue!9.08](1.0698,2.6804)rectangle(1.1648,2.7390);
\fill[DeepBlue!9.68](1.0698,2.7370)rectangle(1.1648,2.7955);
\fill[DeepBlue!10.12](1.0698,2.7935)rectangle(1.1648,2.8520);
\fill[DeepBlue!9.66](1.1628,1.5500)rectangle(1.2578,1.6085);
\fill[DeepBlue!9.90](1.1628,1.6065)rectangle(1.2578,1.6650);
\fill[DeepBlue!9.77](1.1628,1.6630)rectangle(1.2578,1.7216);
\fill[DeepBlue!9.34](1.1628,1.7196)rectangle(1.2578,1.7781);
\fill[DeepBlue!8.78](1.1628,1.7761)rectangle(1.2578,1.8346);
\fill[DeepBlue!8.31](1.1628,1.8326)rectangle(1.2578,1.8911);
\fill[DeepBlue!8.15](1.1628,1.8891)rectangle(1.2578,1.9477);
\fill[DeepBlue!8.39](1.1628,1.9457)rectangle(1.2578,2.0042);
\fill[DeepBlue!8.95](1.1628,2.0022)rectangle(1.2578,2.0607);
\fill[DeepBlue!9.64](1.1628,2.0587)rectangle(1.2578,2.1172);
\fill[DeepBlue!10.23](1.1628,2.1152)rectangle(1.2578,2.1737);
\fill[DeepBlue!10.57](1.1628,2.1717)rectangle(1.2578,2.2303);
\fill[DeepBlue!10.57](1.1628,2.2283)rectangle(1.2578,2.2868);
\fill[DeepBlue!10.23](1.1628,2.2848)rectangle(1.2578,2.3433);
\fill[DeepBlue!9.64](1.1628,2.3413)rectangle(1.2578,2.3998);
\fill[DeepBlue!8.95](1.1628,2.3978)rectangle(1.2578,2.4563);
\fill[DeepBlue!8.39](1.1628,2.4543)rectangle(1.2578,2.5129);
\fill[DeepBlue!8.15](1.1628,2.5109)rectangle(1.2578,2.5694);
\fill[DeepBlue!8.31](1.1628,2.5674)rectangle(1.2578,2.6259);
\fill[DeepBlue!8.78](1.1628,2.6239)rectangle(1.2578,2.6824);
\fill[DeepBlue!9.34](1.1628,2.6804)rectangle(1.2578,2.7390);
\fill[DeepBlue!9.77](1.1628,2.7370)rectangle(1.2578,2.7955);
\fill[DeepBlue!9.90](1.1628,2.7935)rectangle(1.2578,2.8520);
\fill[DeepBlue!8.98](1.2558,1.5500)rectangle(1.3508,1.6085);
\fill[DeepBlue!9.48](1.2558,1.6065)rectangle(1.3508,1.6650);
\fill[DeepBlue!9.66](1.2558,1.6630)rectangle(1.3508,1.7216);
\fill[DeepBlue!9.54](1.2558,1.7196)rectangle(1.3508,1.7781);
\fill[DeepBlue!9.17](1.2558,1.7761)rectangle(1.3508,1.8346);
\fill[DeepBlue!8.68](1.2558,1.8326)rectangle(1.3508,1.8911);
\fill[DeepBlue!8.19](1.2558,1.8891)rectangle(1.3508,1.9477);
\fill[DeepBlue!7.82](1.2558,1.9457)rectangle(1.3508,2.0042);
\fill[DeepBlue!7.63](1.2558,2.0022)rectangle(1.3508,2.0607);
\fill[DeepBlue!7.60](1.2558,2.0587)rectangle(1.3508,2.1172);
\fill[DeepBlue!7.67](1.2558,2.1152)rectangle(1.3508,2.1737);
\fill[DeepBlue!7.73](1.2558,2.1717)rectangle(1.3508,2.2303);
\fill[DeepBlue!7.73](1.2558,2.2283)rectangle(1.3508,2.2868);
\fill[DeepBlue!7.67](1.2558,2.2848)rectangle(1.3508,2.3433);
\fill[DeepBlue!7.60](1.2558,2.3413)rectangle(1.3508,2.3998);
\fill[DeepBlue!7.63](1.2558,2.3978)rectangle(1.3508,2.4563);
\fill[DeepBlue!7.82](1.2558,2.4543)rectangle(1.3508,2.5129);
\fill[DeepBlue!8.19](1.2558,2.5109)rectangle(1.3508,2.5694);
\fill[DeepBlue!8.68](1.2558,2.5674)rectangle(1.3508,2.6259);
\fill[DeepBlue!9.17](1.2558,2.6239)rectangle(1.3508,2.6824);
\fill[DeepBlue!9.54](1.2558,2.6804)rectangle(1.3508,2.7390);
\fill[DeepBlue!9.66](1.2558,2.7370)rectangle(1.3508,2.7955);
\fill[DeepBlue!9.48](1.2558,2.7935)rectangle(1.3508,2.8520);
\fill[DeepBlue!8.16](1.3488,1.5500)rectangle(1.4439,1.6085);
\fill[DeepBlue!8.86](1.3488,1.6065)rectangle(1.4439,1.6650);
\fill[DeepBlue!9.30](1.3488,1.6630)rectangle(1.4439,1.7216);
\fill[DeepBlue!9.49](1.3488,1.7196)rectangle(1.4439,1.7781);
\fill[DeepBlue!9.43](1.3488,1.7761)rectangle(1.4439,1.8346);
\fill[DeepBlue!9.19](1.3488,1.8326)rectangle(1.4439,1.8911);
\fill[DeepBlue!8.83](1.3488,1.8891)rectangle(1.4439,1.9477);
\fill[DeepBlue!8.43](1.3488,1.9457)rectangle(1.4439,2.0042);
\fill[DeepBlue!8.06](1.3488,2.0022)rectangle(1.4439,2.0607);
\fill[DeepBlue!7.78](1.3488,2.0587)rectangle(1.4439,2.1172);
\fill[DeepBlue!7.59](1.3488,2.1152)rectangle(1.4439,2.1737);
\fill[DeepBlue!7.50](1.3488,2.1717)rectangle(1.4439,2.2303);
\fill[DeepBlue!7.50](1.3488,2.2283)rectangle(1.4439,2.2868);
\fill[DeepBlue!7.59](1.3488,2.2848)rectangle(1.4439,2.3433);
\fill[DeepBlue!7.78](1.3488,2.3413)rectangle(1.4439,2.3998);
\fill[DeepBlue!8.06](1.3488,2.3978)rectangle(1.4439,2.4563);
\fill[DeepBlue!8.43](1.3488,2.4543)rectangle(1.4439,2.5129);
\fill[DeepBlue!8.83](1.3488,2.5109)rectangle(1.4439,2.5694);
\fill[DeepBlue!9.19](1.3488,2.5674)rectangle(1.4439,2.6259);
\fill[DeepBlue!9.43](1.3488,2.6239)rectangle(1.4439,2.6824);
\fill[DeepBlue!9.49](1.3488,2.6804)rectangle(1.4439,2.7390);
\fill[DeepBlue!9.30](1.3488,2.7370)rectangle(1.4439,2.7955);
\fill[DeepBlue!8.86](1.3488,2.7935)rectangle(1.4439,2.8520);
\fill[DeepBlue!7.25](1.4419,1.5500)rectangle(1.5369,1.6085);
\fill[DeepBlue!8.04](1.4419,1.6065)rectangle(1.5369,1.6650);
\fill[DeepBlue!8.67](1.4419,1.6630)rectangle(1.5369,1.7216);
\fill[DeepBlue!9.11](1.4419,1.7196)rectangle(1.5369,1.7781);
\fill[DeepBlue!9.34](1.4419,1.7761)rectangle(1.5369,1.8346);
\fill[DeepBlue!9.39](1.4419,1.8326)rectangle(1.5369,1.8911);
\fill[DeepBlue!9.30](1.4419,1.8891)rectangle(1.5369,1.9477);
\fill[DeepBlue!9.12](1.4419,1.9457)rectangle(1.5369,2.0042);
\fill[DeepBlue!8.90](1.4419,2.0022)rectangle(1.5369,2.0607);
\fill[DeepBlue!8.69](1.4419,2.0587)rectangle(1.5369,2.1172);
\fill[DeepBlue!8.53](1.4419,2.1152)rectangle(1.5369,2.1737);
\fill[DeepBlue!8.45](1.4419,2.1717)rectangle(1.5369,2.2303);
\fill[DeepBlue!8.45](1.4419,2.2283)rectangle(1.5369,2.2868);
\fill[DeepBlue!8.53](1.4419,2.2848)rectangle(1.5369,2.3433);
\fill[DeepBlue!8.69](1.4419,2.3413)rectangle(1.5369,2.3998);
\fill[DeepBlue!8.90](1.4419,2.3978)rectangle(1.5369,2.4563);
\fill[DeepBlue!9.12](1.4419,2.4543)rectangle(1.5369,2.5129);
\fill[DeepBlue!9.30](1.4419,2.5109)rectangle(1.5369,2.5694);
\fill[DeepBlue!9.39](1.4419,2.5674)rectangle(1.5369,2.6259);
\fill[DeepBlue!9.34](1.4419,2.6239)rectangle(1.5369,2.6824);
\fill[DeepBlue!9.11](1.4419,2.6804)rectangle(1.5369,2.7390);
\fill[DeepBlue!8.67](1.4419,2.7370)rectangle(1.5369,2.7955);
\fill[DeepBlue!8.04](1.4419,2.7935)rectangle(1.5369,2.8520);
\fill[DeepBlue!6.31](1.5349,1.5500)rectangle(1.6299,1.6085);
\fill[DeepBlue!7.09](1.5349,1.6065)rectangle(1.6299,1.6650);
\fill[DeepBlue!7.80](1.5349,1.6630)rectangle(1.6299,1.7216);
\fill[DeepBlue!8.40](1.5349,1.7196)rectangle(1.6299,1.7781);
\fill[DeepBlue!8.84](1.5349,1.7761)rectangle(1.6299,1.8346);
\fill[DeepBlue!9.14](1.5349,1.8326)rectangle(1.6299,1.8911);
\fill[DeepBlue!9.30](1.5349,1.8891)rectangle(1.6299,1.9477);
\fill[DeepBlue!9.35](1.5349,1.9457)rectangle(1.6299,2.0042);
\fill[DeepBlue!9.34](1.5349,2.0022)rectangle(1.6299,2.0607);
\fill[DeepBlue!9.29](1.5349,2.0587)rectangle(1.6299,2.1172);
\fill[DeepBlue!9.24](1.5349,2.1152)rectangle(1.6299,2.1737);
\fill[DeepBlue!9.21](1.5349,2.1717)rectangle(1.6299,2.2303);
\fill[DeepBlue!9.21](1.5349,2.2283)rectangle(1.6299,2.2868);
\fill[DeepBlue!9.24](1.5349,2.2848)rectangle(1.6299,2.3433);
\fill[DeepBlue!9.29](1.5349,2.3413)rectangle(1.6299,2.3998);
\fill[DeepBlue!9.34](1.5349,2.3978)rectangle(1.6299,2.4563);
\fill[DeepBlue!9.35](1.5349,2.4543)rectangle(1.6299,2.5129);
\fill[DeepBlue!9.30](1.5349,2.5109)rectangle(1.6299,2.5694);
\fill[DeepBlue!9.14](1.5349,2.5674)rectangle(1.6299,2.6259);
\fill[DeepBlue!8.84](1.5349,2.6239)rectangle(1.6299,2.6824);
\fill[DeepBlue!8.40](1.5349,2.6804)rectangle(1.6299,2.7390);
\fill[DeepBlue!7.80](1.5349,2.7370)rectangle(1.6299,2.7955);
\fill[DeepBlue!7.09](1.5349,2.7935)rectangle(1.6299,2.8520);
\fill[DeepBlue!5.43](1.6279,1.5500)rectangle(1.7229,1.6085);
\fill[DeepBlue!6.09](1.6279,1.6065)rectangle(1.7229,1.6650);
\fill[DeepBlue!6.77](1.6279,1.6630)rectangle(1.7229,1.7216);
\fill[DeepBlue!7.42](1.6279,1.7196)rectangle(1.7229,1.7781);
\fill[DeepBlue!7.98](1.6279,1.7761)rectangle(1.7229,1.8346);
\fill[DeepBlue!8.43](1.6279,1.8326)rectangle(1.7229,1.8911);
\fill[DeepBlue!8.77](1.6279,1.8891)rectangle(1.7229,1.9477);
\fill[DeepBlue!9.01](1.6279,1.9457)rectangle(1.7229,2.0042);
\fill[DeepBlue!9.16](1.6279,2.0022)rectangle(1.7229,2.0607);
\fill[DeepBlue!9.25](1.6279,2.0587)rectangle(1.7229,2.1172);
\fill[DeepBlue!9.30](1.6279,2.1152)rectangle(1.7229,2.1737);
\fill[DeepBlue!9.32](1.6279,2.1717)rectangle(1.7229,2.2303);
\fill[DeepBlue!9.32](1.6279,2.2283)rectangle(1.7229,2.2868);
\fill[DeepBlue!9.30](1.6279,2.2848)rectangle(1.7229,2.3433);
\fill[DeepBlue!9.25](1.6279,2.3413)rectangle(1.7229,2.3998);
\fill[DeepBlue!9.16](1.6279,2.3978)rectangle(1.7229,2.4563);
\fill[DeepBlue!9.01](1.6279,2.4543)rectangle(1.7229,2.5129);
\fill[DeepBlue!8.77](1.6279,2.5109)rectangle(1.7229,2.5694);
\fill[DeepBlue!8.43](1.6279,2.5674)rectangle(1.7229,2.6259);
\fill[DeepBlue!7.98](1.6279,2.6239)rectangle(1.7229,2.6824);
\fill[DeepBlue!7.42](1.6279,2.6804)rectangle(1.7229,2.7390);
\fill[DeepBlue!6.77](1.6279,2.7370)rectangle(1.7229,2.7955);
\fill[DeepBlue!6.09](1.6279,2.7935)rectangle(1.7229,2.8520);
\fill[DeepBlue!4.72](1.7209,1.5500)rectangle(1.8160,1.6085);
\fill[DeepBlue!5.16](1.7209,1.6065)rectangle(1.8160,1.6650);
\fill[DeepBlue!5.71](1.7209,1.6630)rectangle(1.8160,1.7216);
\fill[DeepBlue!6.29](1.7209,1.7196)rectangle(1.8160,1.7781);
\fill[DeepBlue!6.86](1.7209,1.7761)rectangle(1.8160,1.8346);
\fill[DeepBlue!7.36](1.7209,1.8326)rectangle(1.8160,1.8911);
\fill[DeepBlue!7.79](1.7209,1.8891)rectangle(1.8160,1.9477);
\fill[DeepBlue!8.14](1.7209,1.9457)rectangle(1.8160,2.0042);
\fill[DeepBlue!8.40](1.7209,2.0022)rectangle(1.8160,2.0607);
\fill[DeepBlue!8.58](1.7209,2.0587)rectangle(1.8160,2.1172);
\fill[DeepBlue!8.69](1.7209,2.1152)rectangle(1.8160,2.1737);
\fill[DeepBlue!8.75](1.7209,2.1717)rectangle(1.8160,2.2303);
\fill[DeepBlue!8.75](1.7209,2.2283)rectangle(1.8160,2.2868);
\fill[DeepBlue!8.69](1.7209,2.2848)rectangle(1.8160,2.3433);
\fill[DeepBlue!8.58](1.7209,2.3413)rectangle(1.8160,2.3998);
\fill[DeepBlue!8.40](1.7209,2.3978)rectangle(1.8160,2.4563);
\fill[DeepBlue!8.14](1.7209,2.4543)rectangle(1.8160,2.5129);
\fill[DeepBlue!7.79](1.7209,2.5109)rectangle(1.8160,2.5694);
\fill[DeepBlue!7.36](1.7209,2.5674)rectangle(1.8160,2.6259);
\fill[DeepBlue!6.86](1.7209,2.6239)rectangle(1.8160,2.6824);
\fill[DeepBlue!6.29](1.7209,2.6804)rectangle(1.8160,2.7390);
\fill[DeepBlue!5.71](1.7209,2.7370)rectangle(1.8160,2.7955);
\fill[DeepBlue!5.16](1.7209,2.7935)rectangle(1.8160,2.8520);
\fill[DeepBlue!4.30](1.8140,1.5500)rectangle(1.9090,1.6085);
\fill[DeepBlue!4.47](1.8140,1.6065)rectangle(1.9090,1.6650);
\fill[DeepBlue!4.78](1.8140,1.6630)rectangle(1.9090,1.7216);
\fill[DeepBlue!5.19](1.8140,1.7196)rectangle(1.9090,1.7781);
\fill[DeepBlue!5.66](1.8140,1.7761)rectangle(1.9090,1.8346);
\fill[DeepBlue!6.12](1.8140,1.8326)rectangle(1.9090,1.8911);
\fill[DeepBlue!6.54](1.8140,1.8891)rectangle(1.9090,1.9477);
\fill[DeepBlue!6.91](1.8140,1.9457)rectangle(1.9090,2.0042);
\fill[DeepBlue!7.21](1.8140,2.0022)rectangle(1.9090,2.0607);
\fill[DeepBlue!7.43](1.8140,2.0587)rectangle(1.9090,2.1172);
\fill[DeepBlue!7.58](1.8140,2.1152)rectangle(1.9090,2.1737);
\fill[DeepBlue!7.65](1.8140,2.1717)rectangle(1.9090,2.2303);
\fill[DeepBlue!7.65](1.8140,2.2283)rectangle(1.9090,2.2868);
\fill[DeepBlue!7.58](1.8140,2.2848)rectangle(1.9090,2.3433);
\fill[DeepBlue!7.43](1.8140,2.3413)rectangle(1.9090,2.3998);
\fill[DeepBlue!7.21](1.8140,2.3978)rectangle(1.9090,2.4563);
\fill[DeepBlue!6.91](1.8140,2.4543)rectangle(1.9090,2.5129);
\fill[DeepBlue!6.54](1.8140,2.5109)rectangle(1.9090,2.5694);
\fill[DeepBlue!6.12](1.8140,2.5674)rectangle(1.9090,2.6259);
\fill[DeepBlue!5.66](1.8140,2.6239)rectangle(1.9090,2.6824);
\fill[DeepBlue!5.19](1.8140,2.6804)rectangle(1.9090,2.7390);
\fill[DeepBlue!4.78](1.8140,2.7370)rectangle(1.9090,2.7955);
\fill[DeepBlue!4.47](1.8140,2.7935)rectangle(1.9090,2.8520);
\fill[DeepBlue!4.19](1.9070,1.5500)rectangle(2.0020,1.6085);
\fill[DeepBlue!4.13](1.9070,1.6065)rectangle(2.0020,1.6650);
\fill[DeepBlue!4.19](1.9070,1.6630)rectangle(2.0020,1.7216);
\fill[DeepBlue!4.35](1.9070,1.7196)rectangle(2.0020,1.7781);
\fill[DeepBlue!4.62](1.9070,1.7761)rectangle(2.0020,1.8346);
\fill[DeepBlue!4.93](1.9070,1.8326)rectangle(2.0020,1.8911);
\fill[DeepBlue!5.26](1.9070,1.8891)rectangle(2.0020,1.9477);
\fill[DeepBlue!5.56](1.9070,1.9457)rectangle(2.0020,2.0042);
\fill[DeepBlue!5.83](1.9070,2.0022)rectangle(2.0020,2.0607);
\fill[DeepBlue!6.04](1.9070,2.0587)rectangle(2.0020,2.1172);
\fill[DeepBlue!6.18](1.9070,2.1152)rectangle(2.0020,2.1737);
\fill[DeepBlue!6.25](1.9070,2.1717)rectangle(2.0020,2.2303);
\fill[DeepBlue!6.25](1.9070,2.2283)rectangle(2.0020,2.2868);
\fill[DeepBlue!6.18](1.9070,2.2848)rectangle(2.0020,2.3433);
\fill[DeepBlue!6.04](1.9070,2.3413)rectangle(2.0020,2.3998);
\fill[DeepBlue!5.83](1.9070,2.3978)rectangle(2.0020,2.4563);
\fill[DeepBlue!5.56](1.9070,2.4543)rectangle(2.0020,2.5129);
\fill[DeepBlue!5.26](1.9070,2.5109)rectangle(2.0020,2.5694);
\fill[DeepBlue!4.93](1.9070,2.5674)rectangle(2.0020,2.6259);
\fill[DeepBlue!4.62](1.9070,2.6239)rectangle(2.0020,2.6824);
\fill[DeepBlue!4.35](1.9070,2.6804)rectangle(2.0020,2.7390);
\fill[DeepBlue!4.19](1.9070,2.7370)rectangle(2.0020,2.7955);
\fill[DeepBlue!4.13](1.9070,2.7935)rectangle(2.0020,2.8520);
\draw[->](-2.2,0)--(2.3,0)node[right]{$u$};
\draw[TextGray,dashed,thin]plot coordinates{(-2.200,0.001537)(-2.180,0.001817)(-2.160,0.002104)(-2.140,0.002393)(-2.120,0.002679)(-2.100,0.002956)(-2.080,0.003217)(-2.060,0.003457)(-2.040,0.003669)(-2.020,0.003849)(-2.000,0.003991)(-1.980,0.004092)(-1.960,0.004147)(-1.940,0.004155)(-1.920,0.004114)(-1.900,0.004023)(-1.880,0.003884)(-1.860,0.003698)(-1.840,0.003468)(-1.820,0.003200)(-1.800,0.002898)(-1.780,0.002571)(-1.760,0.002226)(-1.740,0.001872)(-1.720,0.001519)(-1.700,0.001177)(-1.680,0.000859)(-1.660,0.000574)(-1.640,0.000335)(-1.620,0.000153)(-1.600,0.000038)(-1.580,0.000000)(-1.560,0.000049)(-1.540,0.000191)(-1.520,0.000434)(-1.500,0.000781)(-1.480,0.001235)(-1.460,0.001797)(-1.440,0.002465)(-1.420,0.003233)(-1.400,0.004096)(-1.380,0.005044)(-1.360,0.006066)(-1.340,0.007147)(-1.320,0.008271)(-1.300,0.009420)(-1.280,0.010573)(-1.260,0.011710)(-1.240,0.012807)(-1.220,0.013842)(-1.200,0.014791)(-1.180,0.015630)(-1.160,0.016337)(-1.140,0.016893)(-1.120,0.017277)(-1.100,0.017474)(-1.080,0.017470)(-1.060,0.017258)(-1.040,0.016831)(-1.020,0.016191)(-1.000,0.015343)(-0.980,0.014299)(-0.960,0.013077)(-0.940,0.011703)(-0.920,0.010208)(-0.900,0.008631)(-0.880,0.007018)(-0.860,0.005424)(-0.840,0.003907)(-0.820,0.002537)(-0.800,0.001385)(-0.780,0.000532)(-0.760,0.000063)(-0.740,0.000067)(-0.720,0.000637)(-0.700,0.001870)(-0.680,0.003864)(-0.660,0.006717)(-0.640,0.010529)(-0.620,0.015396)(-0.600,0.021411)(-0.580,0.028666)(-0.560,0.037245)(-0.540,0.047224)(-0.520,0.058674)(-0.500,0.071654)(-0.480,0.086215)(-0.460,0.102394)(-0.440,0.120216)(-0.420,0.139694)(-0.400,0.160823)(-0.380,0.183587)(-0.360,0.207950)(-0.340,0.233865)(-0.320,0.261263)(-0.300,0.290064)(-0.280,0.320168)(-0.260,0.351463)(-0.240,0.383817)(-0.220,0.417088)(-0.200,0.451117)(-0.180,0.485736)(-0.160,0.520762)(-0.140,0.556004)(-0.120,0.591264)(-0.100,0.626334)(-0.080,0.661004)(-0.060,0.695059)(-0.040,0.728284)(-0.020,0.760464)(0.000,0.791388)(0.020,0.820849)(0.040,0.848646)(0.060,0.874589)(0.080,0.898497)(0.100,0.920202)(0.120,0.939549)(0.140,0.956400)(0.160,0.970635)(0.180,0.982149)(0.200,0.990859)(0.220,0.996701)(0.240,0.999633)(0.260,0.999633)(0.280,0.996701)(0.300,0.990859)(0.320,0.982149)(0.340,0.970635)(0.360,0.956400)(0.380,0.939549)(0.400,0.920202)(0.420,0.898497)(0.440,0.874589)(0.460,0.848646)(0.480,0.820849)(0.500,0.791388)(0.520,0.760464)(0.540,0.728284)(0.560,0.695059)(0.580,0.661004)(0.600,0.626334)(0.620,0.591264)(0.640,0.556004)(0.660,0.520762)(0.680,0.485736)(0.700,0.451117)(0.720,0.417088)(0.740,0.383817)(0.760,0.351463)(0.780,0.320168)(0.800,0.290064)(0.820,0.261263)(0.840,0.233865)(0.860,0.207950)(0.880,0.183587)(0.900,0.160823)(0.920,0.139694)(0.940,0.120216)(0.960,0.102394)(0.980,0.086215)(1.000,0.071654)(1.020,0.058674)(1.040,0.047224)(1.060,0.037245)(1.080,0.028666)(1.100,0.021411)(1.120,0.015396)(1.140,0.010529)(1.160,0.006717)(1.180,0.003864)(1.200,0.001870)(1.220,0.000637)(1.240,0.000067)(1.260,0.000063)(1.280,0.000532)(1.300,0.001385)(1.320,0.002537)(1.340,0.003907)(1.360,0.005424)(1.380,0.007018)(1.400,0.008631)(1.420,0.010208)(1.440,0.011703)(1.460,0.013077)(1.480,0.014299)(1.500,0.015343)(1.520,0.016191)(1.540,0.016831)(1.560,0.017258)(1.580,0.017470)(1.600,0.017474)(1.620,0.017277)(1.640,0.016893)(1.660,0.016337)(1.680,0.015630)(1.700,0.014791)(1.720,0.013842)(1.740,0.012807)(1.760,0.011710)(1.780,0.010573)(1.800,0.009420)(1.820,0.008271)(1.840,0.007147)(1.860,0.006066)(1.880,0.005044)(1.900,0.004096)(1.920,0.003233)(1.940,0.002465)(1.960,0.001797)(1.980,0.001235)(2.000,0.000781)(2.020,0.000434)(2.040,0.000191)(2.060,0.000049)(2.080,0.000000)(2.100,0.000038)(2.120,0.000153)(2.140,0.000335)(2.160,0.000574)(2.180,0.000859)(2.200,0.001177)};
\draw[TextGray,dashed,thin]plot coordinates{(-2.200,0.001177)(-2.180,0.000859)(-2.160,0.000574)(-2.140,0.000335)(-2.120,0.000153)(-2.100,0.000038)(-2.080,0.000000)(-2.060,0.000049)(-2.040,0.000191)(-2.020,0.000434)(-2.000,0.000781)(-1.980,0.001235)(-1.960,0.001797)(-1.940,0.002465)(-1.920,0.003233)(-1.900,0.004096)(-1.880,0.005044)(-1.860,0.006066)(-1.840,0.007147)(-1.820,0.008271)(-1.800,0.009420)(-1.780,0.010573)(-1.760,0.011710)(-1.740,0.012807)(-1.720,0.013842)(-1.700,0.014791)(-1.680,0.015630)(-1.660,0.016337)(-1.640,0.016893)(-1.620,0.017277)(-1.600,0.017474)(-1.580,0.017470)(-1.560,0.017258)(-1.540,0.016831)(-1.520,0.016191)(-1.500,0.015343)(-1.480,0.014299)(-1.460,0.013077)(-1.440,0.011703)(-1.420,0.010208)(-1.400,0.008631)(-1.380,0.007018)(-1.360,0.005424)(-1.340,0.003907)(-1.320,0.002537)(-1.300,0.001385)(-1.280,0.000532)(-1.260,0.000063)(-1.240,0.000067)(-1.220,0.000637)(-1.200,0.001870)(-1.180,0.003864)(-1.160,0.006717)(-1.140,0.010529)(-1.120,0.015396)(-1.100,0.021411)(-1.080,0.028666)(-1.060,0.037245)(-1.040,0.047224)(-1.020,0.058674)(-1.000,0.071654)(-0.980,0.086215)(-0.960,0.102394)(-0.940,0.120216)(-0.920,0.139694)(-0.900,0.160823)(-0.880,0.183587)(-0.860,0.207950)(-0.840,0.233865)(-0.820,0.261263)(-0.800,0.290064)(-0.780,0.320168)(-0.760,0.351463)(-0.740,0.383817)(-0.720,0.417088)(-0.700,0.451117)(-0.680,0.485736)(-0.660,0.520762)(-0.640,0.556004)(-0.620,0.591264)(-0.600,0.626334)(-0.580,0.661004)(-0.560,0.695059)(-0.540,0.728284)(-0.520,0.760464)(-0.500,0.791388)(-0.480,0.820849)(-0.460,0.848646)(-0.440,0.874589)(-0.420,0.898497)(-0.400,0.920202)(-0.380,0.939549)(-0.360,0.956400)(-0.340,0.970635)(-0.320,0.982149)(-0.300,0.990859)(-0.280,0.996701)(-0.260,0.999633)(-0.240,0.999633)(-0.220,0.996701)(-0.200,0.990859)(-0.180,0.982149)(-0.160,0.970635)(-0.140,0.956400)(-0.120,0.939549)(-0.100,0.920202)(-0.080,0.898497)(-0.060,0.874589)(-0.040,0.848646)(-0.020,0.820849)(0.000,0.791388)(0.020,0.760464)(0.040,0.728284)(0.060,0.695059)(0.080,0.661004)(0.100,0.626334)(0.120,0.591264)(0.140,0.556004)(0.160,0.520762)(0.180,0.485736)(0.200,0.451117)(0.220,0.417088)(0.240,0.383817)(0.260,0.351463)(0.280,0.320168)(0.300,0.290064)(0.320,0.261263)(0.340,0.233865)(0.360,0.207950)(0.380,0.183587)(0.400,0.160823)(0.420,0.139694)(0.440,0.120216)(0.460,0.102394)(0.480,0.086215)(0.500,0.071654)(0.520,0.058674)(0.540,0.047224)(0.560,0.037245)(0.580,0.028666)(0.600,0.021411)(0.620,0.015396)(0.640,0.010529)(0.660,0.006717)(0.680,0.003864)(0.700,0.001870)(0.720,0.000637)(0.740,0.000067)(0.760,0.000063)(0.780,0.000532)(0.800,0.001385)(0.820,0.002537)(0.840,0.003907)(0.860,0.005424)(0.880,0.007018)(0.900,0.008631)(0.920,0.010208)(0.940,0.011703)(0.960,0.013077)(0.980,0.014299)(1.000,0.015343)(1.020,0.016191)(1.040,0.016831)(1.060,0.017258)(1.080,0.017470)(1.100,0.017474)(1.120,0.017277)(1.140,0.016893)(1.160,0.016337)(1.180,0.015630)(1.200,0.014791)(1.220,0.013842)(1.240,0.012807)(1.260,0.011710)(1.280,0.010573)(1.300,0.009420)(1.320,0.008271)(1.340,0.007147)(1.360,0.006066)(1.380,0.005044)(1.400,0.004096)(1.420,0.003233)(1.440,0.002465)(1.460,0.001797)(1.480,0.001235)(1.500,0.000781)(1.520,0.000434)(1.540,0.000191)(1.560,0.000049)(1.580,0.000000)(1.600,0.000038)(1.620,0.000153)(1.640,0.000335)(1.660,0.000574)(1.680,0.000859)(1.700,0.001177)(1.720,0.001519)(1.740,0.001872)(1.760,0.002226)(1.780,0.002571)(1.800,0.002898)(1.820,0.003200)(1.840,0.003468)(1.860,0.003698)(1.880,0.003884)(1.900,0.004023)(1.920,0.004114)(1.940,0.004155)(1.960,0.004147)(1.980,0.004092)(2.000,0.003991)(2.020,0.003849)(2.040,0.003669)(2.060,0.003457)(2.080,0.003217)(2.100,0.002956)(2.120,0.002679)(2.140,0.002393)(2.160,0.002104)(2.180,0.001817)(2.200,0.001537)};
\draw[DeepBlue,thick]plot coordinates{(-2.200,0.002715)(-2.180,0.002675)(-2.160,0.002678)(-2.140,0.002728)(-2.120,0.002832)(-2.100,0.002994)(-2.080,0.003217)(-2.060,0.003505)(-2.040,0.003860)(-2.020,0.004283)(-2.000,0.004772)(-1.980,0.005327)(-1.960,0.005945)(-1.940,0.006620)(-1.920,0.007347)(-1.900,0.008119)(-1.880,0.008928)(-1.860,0.009763)(-1.840,0.010615)(-1.820,0.011471)(-1.800,0.012318)(-1.780,0.013144)(-1.760,0.013936)(-1.740,0.014679)(-1.720,0.015361)(-1.700,0.015968)(-1.680,0.016489)(-1.660,0.016912)(-1.640,0.017228)(-1.620,0.017430)(-1.600,0.017512)(-1.580,0.017471)(-1.560,0.017306)(-1.540,0.017022)(-1.520,0.016624)(-1.500,0.016124)(-1.480,0.015534)(-1.460,0.014874)(-1.440,0.014167)(-1.420,0.013441)(-1.400,0.012727)(-1.380,0.012062)(-1.360,0.011489)(-1.340,0.011054)(-1.320,0.010807)(-1.300,0.010805)(-1.280,0.011105)(-1.260,0.011773)(-1.240,0.012874)(-1.220,0.014479)(-1.200,0.016661)(-1.180,0.019494)(-1.160,0.023055)(-1.140,0.027422)(-1.120,0.032673)(-1.100,0.038885)(-1.080,0.046137)(-1.060,0.054502)(-1.040,0.064055)(-1.020,0.074864)(-1.000,0.086997)(-0.980,0.100514)(-0.960,0.115471)(-0.940,0.131919)(-0.920,0.149901)(-0.900,0.169454)(-0.880,0.190605)(-0.860,0.213374)(-0.840,0.237772)(-0.820,0.263800)(-0.800,0.291449)(-0.780,0.320701)(-0.760,0.351525)(-0.740,0.383884)(-0.720,0.417725)(-0.700,0.452987)(-0.680,0.489600)(-0.660,0.527479)(-0.640,0.566533)(-0.620,0.606660)(-0.600,0.647746)(-0.580,0.689670)(-0.560,0.732303)(-0.540,0.775507)(-0.520,0.819138)(-0.500,0.863042)(-0.480,0.907063)(-0.460,0.951040)(-0.440,0.994805)(-0.420,1.038191)(-0.400,1.081025)(-0.380,1.123136)(-0.360,1.164351)(-0.340,1.204499)(-0.320,1.243412)(-0.300,1.280923)(-0.280,1.316870)(-0.260,1.351096)(-0.240,1.383450)(-0.220,1.413789)(-0.200,1.441976)(-0.180,1.467884)(-0.160,1.491397)(-0.140,1.512405)(-0.120,1.530813)(-0.100,1.546536)(-0.080,1.559501)(-0.060,1.569648)(-0.040,1.576930)(-0.020,1.581312)(0.000,1.582776)(0.020,1.581312)(0.040,1.576930)(0.060,1.569648)(0.080,1.559501)(0.100,1.546536)(0.120,1.530813)(0.140,1.512405)(0.160,1.491397)(0.180,1.467884)(0.200,1.441976)(0.220,1.413789)(0.240,1.383450)(0.260,1.351096)(0.280,1.316870)(0.300,1.280923)(0.320,1.243412)(0.340,1.204499)(0.360,1.164351)(0.380,1.123136)(0.400,1.081025)(0.420,1.038191)(0.440,0.994805)(0.460,0.951040)(0.480,0.907063)(0.500,0.863042)(0.520,0.819138)(0.540,0.775507)(0.560,0.732303)(0.580,0.689670)(0.600,0.647746)(0.620,0.606660)(0.640,0.566533)(0.660,0.527479)(0.680,0.489600)(0.700,0.452987)(0.720,0.417725)(0.740,0.383884)(0.760,0.351525)(0.780,0.320701)(0.800,0.291449)(0.820,0.263800)(0.840,0.237772)(0.860,0.213374)(0.880,0.190605)(0.900,0.169454)(0.920,0.149901)(0.940,0.131919)(0.960,0.115471)(0.980,0.100514)(1.000,0.086997)(1.020,0.074864)(1.040,0.064055)(1.060,0.054502)(1.080,0.046137)(1.100,0.038885)(1.120,0.032673)(1.140,0.027422)(1.160,0.023055)(1.180,0.019494)(1.200,0.016661)(1.220,0.014479)(1.240,0.012874)(1.260,0.011773)(1.280,0.011105)(1.300,0.010805)(1.320,0.010807)(1.340,0.011054)(1.360,0.011489)(1.380,0.012062)(1.400,0.012727)(1.420,0.013441)(1.440,0.014167)(1.460,0.014874)(1.480,0.015534)(1.500,0.016124)(1.520,0.016624)(1.540,0.017022)(1.560,0.017306)(1.580,0.017471)(1.600,0.017512)(1.620,0.017430)(1.640,0.017228)(1.660,0.016912)(1.680,0.016489)(1.700,0.015968)(1.720,0.015361)(1.740,0.014679)(1.760,0.013936)(1.780,0.013144)(1.800,0.012318)(1.820,0.011471)(1.840,0.010615)(1.860,0.009763)(1.880,0.008928)(1.900,0.008119)(1.920,0.007347)(1.940,0.006620)(1.960,0.005945)(1.980,0.005327)(2.000,0.004772)(2.020,0.004283)(2.040,0.003860)(2.060,0.003505)(2.080,0.003217)(2.100,0.002994)(2.120,0.002832)(2.140,0.002728)(2.160,0.002678)(2.180,0.002675)(2.200,0.002715)};
\node[below,align=center]at(0,-.2){$\Delta\theta/\theta_{\rm A}=0.5$\\低于瑞利判据};
\end{scope}
\end{tikzpicture}\figcap{等强、不相干点源的艾里强度相加。上为二维亮度，下为中心剖面；虚线是单点强度，实线是总强度，$u=\theta/\theta_{\rm A}$、$\theta_{\rm A}=1.22\lambda/D$。}\label{ch30:fig:s45-22-9}\end{minipage}\end{center}
}

\newcommand{\ArrayDiagram}{%
\begin{center}\begin{minipage}{\linewidth}\centering\begin{tikzpicture}[x=.064cm,y=.035cm,font=\small]
\draw[RuleGray](-180,-90)rectangle(20,60);
\foreach \x/\txt in {-180/180^\circ,-120/120^\circ W,-60/60^\circ W,0/0^\circ}{
\draw[RuleGray](\x,-90)--(\x,60);\node[above]at(\x,60){$\txt$};}
\foreach \y in {-60,-30,0,30}{\draw[RuleGray](-180,\y)--(20,\y);\node[left]at(-180,\y){$\y^\circ$};}
\coordinate(H)at(-155,20);\coordinate(A)at(-110,33);\coordinate(M)at(-97,19);
\coordinate(C)at(-68,-23);\coordinate(S)at(-3,37);\coordinate(P)at(0,-90);
\draw[Blue!50](H)--(S)--(C)--(H) (A)--(C) (M)--(S) (C)--(P) (H)--(P);
\foreach \n in {H,A,M,C,S,P}{\fill[DeepBlue](\n)circle(2pt);}
\node[align=center]at(-130,-113){夏威夷\\SMA / JCMT};
\draw[TextGray](H)--(-173,0)--(-173,-105)--(-160,-113);
\node[above left]at(A){亚利桑那 SMT};
\node at(-83,48){墨西哥 LMT};
\draw[TextGray](M)--(-88,37)--(-83,39);
\node[anchor=west,align=left]at(-160,-62){智利\\ALMA / APEX};
\draw[TextGray](C)--(-115,-45)--(-143,-51);
\node[above left]at(S){西班牙 IRAM};
\node[below,align=center]at(-30,-100){南极 SPT（极点经度任意）};
\draw[TextGray](P)--(-15,-99);
\end{tikzpicture}\figcap{2017年事件视界望远镜的8个台站分布在6个地理位置。点位按近似经纬度标出，线仅连接示例基线，不代表地面大圆路径或连续填满的口径。}\label{ch30:fig:s45-22-10}\end{minipage}\end{center}}
\newcommand{\LaserMoonDiagram}{%
\begin{center}\begin{minipage}{\linewidth}\centering\begin{tikzpicture}[scale=.8,>=Latex,font=\small]
 \begin{scope}[xshift=-3.5cm]
 \draw[very thick](-3,-.2)rectangle(-2.5,.2);
 \draw[DeepBlue](-2.5,0)--(2,1.7) (-2.5,0)--(2,-1.7);
 \draw[dashed](-2.5,0)--(2.5,0);
 \draw[thick](2,-1.8)--(2,1.8);
 \draw[<->](-2.5,-2.1)--(2,-2.1)node[midway,below]{$L$};
 \node[above]at(-2.5,.4){孔径 $D$};
 \node[above,align=center]at(2,1.9){月球距离处\\垂直截面};
 \end{scope}
 \begin{scope}[xshift=4cm]
 \draw[DeepBlue,thick](0,0)circle(1.7);
 \filldraw[fill=Gold!20,draw=TextGray](0,0)circle(.996);
 \draw[->](0,0)--(0,1.7)node[midway,right]{$r_1$};
 \node at(0,-.2){月球投影};
 \node[below,align=center]at(0,-1.95){$r_1=2.965\,\mathrm{Mm}>R_{\rm M}$\\$R_{\rm M}=1.737\,\mathrm{Mm}$};
 \end{scope}
\end{tikzpicture}\figcap{右图按半径比例画出中央斑第一暗环与月球投影；左图传播距离不按比例。}\label{ch30:fig:s45-22-11}\end{minipage}\end{center}
}

% Geometry: shared coordinates; transparent labels; numerical profiles where required.
\newcommand{\MolecularLimitDiagram}{%
\historyfigure{ribosome-20260923.png}{核糖体的大、小亚基及核糖核酸链。整体尺度约为几十纳米，远小于普通可见光显微镜约 $200\,\mathrm{nm}$ 的分辨尺度；可探测散射信号不等于可分辨内部结构。\label{ch30:fig:s45-22-12}}
}

% Geometry: shared coordinates; transparent labels; numerical profiles where required.
\newcommand{\GridPatternDiagram}{%
\begin{center}\begin{minipage}{\linewidth}\centering\begin{tikzpicture}[scale=.8,font=\small]
 \begin{scope}[xshift=-3.2cm]
 \foreach \x in {-1.5,-1.2,...,1.5}{\draw[TextGray](\x,-1.5)--(\x,1.5);}
 \foreach \y in {-1.5,-1.2,...,1.5}{\draw[TextGray](-1.5,\y)--(1.5,\y);}
 \node[below]at(0,-1.8){经纬正交周期孔阵};
 \end{scope}
 \begin{scope}[xshift=3.2cm]
 \foreach \m in {-2,-1,0,1,2}{\foreach \n in {-2,-1,0,1,2}{
 \ifnum\m=0\def\opa{1}\else\pgfmathsetmacro{\opa}{(sin(72*\m)/(.4*pi*\m))^2}\fi
 \ifnum\n=0\def\opb{1}\else\pgfmathsetmacro{\opb}{(sin(72*\n)/(.4*pi*\n))^2}\fi
 \pgfmathsetmacro{\op}{\opa*\opb}
 \fill[DeepBlue,opacity=\op](.75*\m,.75*\n)circle(.11);}}
 \node[below]at(0,-1.8){亮度受单个孔的衍射包络调制};
 \end{scope}
\end{tikzpicture}\figcap{二维正方形孔阵的衍射级次；点距由周期决定，强度并不相同。}\label{ch30:fig:s45-22-13}\end{minipage}\end{center}
}

% Geometry: shared coordinates; transparent labels; numerical profiles where required.
\newcommand{\ObliqueGratingDiagram}{%
\begin{center}\begin{minipage}{\linewidth}\centering\begin{tikzpicture}[scale=.9,>=Latex,font=\small]
 \coordinate(A)at(0,0);\coordinate(B)at(0,1.5);
 \draw[very thick](0,-1.15)--(0,-.08) (0,.08)--(0,1.42) (0,1.58)--(0,2.65);
 \foreach \y in {0,1.5}{
 \draw[DeepBlue,->](-2.7,{\y-2.7*tan(20)})--(0,\y);
 \draw[Blue,->](0,\y)--(2.7,{\y+2.7*tan(35)});
 \draw[dashed](-.3,\y)--(3,\y);}
 \draw[DeepBlue,densely dotted](0,0)--(1.5,{1.5*tan(20)});
 \draw(.9,0)arc(0:20:.9);\node at(1.2,.12){$\theta_i$};
 \draw(.65,1.5)arc(0:35:.65);\node at(.9,1.78){$\theta_m$};
 \node[above left=4pt]at(A){$A$};\node[above left=4pt]at(B){$B$};
 \draw[<->](-3.2,0)--(-3.2,1.5)node[midway,left]{$a$};
 \node[below,align=center]at(0,-1.5){有向角均从 $+x$ 法线向 $+y$ 量取\\入射贡献 $ka\sin\theta_i$，远场传播贡献 $-ka\sin\theta_m$};
\end{tikzpicture}\figcap{解图：比较刻线 $A,B$，入射相位差与传播相位差必须一起计入。}\label{ch30:fig:s45-ex-22-9}\end{minipage}\end{center}
}


衍射是光的波动性带来的另一类基本现象。几何光学预言障碍物应产生边界锐利的影子、光通过小孔后只留下孔的几何投影；实际观察却会在边缘看到明暗细纹，在圆孔后看到中央亮斑和环纹。凡是传播方向或空间强度分布偏离几何光学预言、并由有限波长引起的现象，都可归入衍射。

可见光波长很短，所以日常大尺度物体附近衍射角通常很小；当孔径、障碍尺度接近波长，或传播距离足够长时，衍射便不可忽略。完整衍射理论需要积分波动方程，本讲从周期结构、单缝和圆孔几个可计算模型建立物理图像。

\BasicDiffractionPanel
图~\ref{ch30:fig:s45-22-1} 中虚线只表示几何光学的阴影边界；边界附近的明暗变化需要开放波面上各部分的波相干叠加才能解释。



\section{单缝衍射与圆孔图样}



宽度为 $D$ 的矩形孔可视为连续分布的次级波源，是 $N\to\infty$、$d\to0$ 且 $Nd=D$ 的光栅极限。令
\[
 \beta=\frac{\pi D}{\lambda}\sin\theta,
\]
夫琅禾费远场（Fraunhofer far field）强度为
\begin{equation}
 I(\theta)=I_0\left(\frac{\sin\beta}{\beta}\right)^2.
\end{equation}
第一暗纹满足 $D\sin\theta=\lambda$。小角时，中央光束的半角量级
\begin{equation}
 \theta\sim\frac{\lambda}{D}.
\end{equation}
孔越小，衍射角越大；这与“缩小孔能无限提高几何清晰度”的直觉相反。圆孔产生艾里图样（Airy pattern），第一暗环角半径为
\begin{equation}
 \theta_A=1.22\frac{\lambda}{D}.
\end{equation}

连续极限也可直接积分得到。孔内横坐标为 $x$ 的窄条相对中央窄条有相位 $kx\sin\theta$，所以
\[
 \widetilde A(\theta)\propto
 \int_{-D/2}^{D/2}\e^{\mathrm ikx\sin\theta}\,\dd x
 =D\,\frac{\sin\beta}{\beta},
 \qquad \beta=\frac{\pi D\sin\theta}{\lambda}.
\]
强度取振幅模平方，才得到 $\operatorname{sinc}^2$ 曲线。中央 $\beta=0$ 处同样应用极限 $\sin\beta/\beta\to1$；图~\ref{ch30:fig:s45-22-8} 的中央峰必须最高，若曲线在原点落到零，就是数值绘图把可去奇点处理错了。

中央亮斑位于两侧第一暗纹之间，角宽约 $2\lambda/D$，而常说的“衍射张角 $\lambda/D$”是半角。题目若问屏上直径，应再乘传播距离并乘 $2$；若问圆孔第一暗环，则还要使用 $1.22$。这三个量只差常数因子，却对应不同几何定义，必须结合图中的半径或直径标记判断。

\subsection{从整条缝的相干叠加得到暗纹}
设缝宽为 $a$，横向坐标为 $x\in[-a/2,a/2]$，入射光在缝内振幅均匀、相位相同。远场观察方向与轴成角 $\theta$，缝内相隔 $x$ 的两点向该方向传播时相差光程 $x\sin\theta$。把振幅密度记为常量 $C$，则
\[
 U(\theta)=C\int_{-a/2}^{a/2}
 \exp\!\left(\frac{2\pi\ii x\sin\theta}{\lambda}\right)\dd x.
\]
积分后，以 $U(0)=Ca$ 归一化，得到
\[
 \frac{U(\theta)}{U(0)}=\frac{\sin\beta}{\beta},
 \qquad \beta=\frac{\pi a\sin\theta}{\lambda}.
\]
因此强度为 $I(0)(\sin\beta/\beta)^2$。在 $\beta=0$ 处应取极限1，不是把“分母为零”解释为强度无穷大。暗纹来自 $\beta=m\pi$，但必须排除 $m=0$，即 $a\sin\theta=m\lambda$、$m=\pm1,\pm2,\ldots$。

第一对暗纹围成中央主极大。近轴屏距为 $L$ 时，它的全宽约为 $2\lambda L/a$；在焦距为 $f$ 的透镜后焦面上，改为 $2\lambda f/a$。圆孔的图样来自对圆形面积积分，不能把缝宽直接改叫孔径而继续用单缝系数；第一暗环对应 $1.22\lambda/D$ 的角半径。


\SingleSlitCurve


\Needspace{7\baselineskip}\begin{origexamplebox}{【不同波长的中央亮斑】}\label{ch30:ex:s45-22-0}
单色平行光照射宽度可调的单缝，缝后远处放置接收屏。红光波长 $655\,\mathrm{nm}$，绿光波长 $532\,\mathrm{nm}$。判断下列说法：

A. 红光通过 $0.8\,\mathrm{mm}$ 缝、绿光通过 $0.7\,\mathrm{mm}$ 缝时，绿光中央亮斑更宽；

B. 上述配置中红光中央亮斑更宽；

C. 红光通过 $0.7\,\mathrm{mm}$ 缝、绿光通过 $0.8\,\mathrm{mm}$ 缝时，绿光中央亮斑更宽；

D. 上述配置中红光中央亮斑更宽。
\end{origexamplebox}\par\nopagebreak[3]
\begin{xesolution}
单缝中央亮斑宽度 $w\simeq2L\lambda/a$，只需比较 $\lambda/a$。第一组
\[
 \frac{655}{0.8}=818.8>\frac{532}{0.7}=760,
\]
故红光更宽，对应B。第二组
\[
 \frac{655}{0.7}=935.7>\frac{532}{0.8}=665,
\]
仍是红光更宽，对应D。因此
\[
 \boxed{\text{B、D正确。}}
\]
条纹尺度由波长与孔径宽度的比值共同决定，不能只比较波长或只比较缝宽。
\end{xesolution}

\section{孔径、分辨率与成像限制}

任何点物经有限圆孔成像后都成为艾里斑。两个点的斑重叠过多时便无法分辨。瑞利判据（Rayleigh criterion）把一个艾里斑的主极大落在另一个的第一暗环处作为临界，因此口径 $D$ 的理想系统角分辨率约为 $1.22\lambda/D$。实际仪器还受像差、传感器采样、噪声和大气扰动限制，但有限孔径的衍射不能仅靠改善表面打磨消除；这里采用瑞利判据，并不把它当成所有探测和参数估计方法的绝对信息极限。

\ResolutionPanel
图~\ref{ch30:fig:s45-22-9} 的实线是两个不相干点源的强度之和；比较两峰与中间凹陷，就能看出点源间距改变时分辨程度的变化。

单台望远镜口径受制造、质量和支架限制。多台望远镜通过相位同步和相关处理可合成长基线，提高特定方向的角分辨率。长基线光学/射电阵列以及事件视界望远镜的全球基线，都体现“有效口径由最大相干基线决定”的思想；阵列的稀疏采样和成像重建仍需单独处理，不能简单等同于一面实心巨镜。

分辨率与像素采样是两层限制。衍射先决定光学系统能传递到像面的最细结构，传感器再对它采样；像素太粗会丢失本可分辨的细节，像素远小于衍射斑则只会把同一模糊斑分成更多格。缩小光圈可减轻部分几何像差并增加景深，但同时增大 $\lambda/D$，所以相机往往在中等光圈获得最佳实际清晰度。后面的“视网膜屏”和“数码相机”两题分别从人眼采样与传感器采样检验这一关系。

\historyfigure{astronomy-contexts-20260923.png}{射电望远镜阵列与月球激光实验的情景。左图强调多台天线协作，右图强调地面发射、接收端；图不按距离与束宽比例绘制。\label{ch30:fig:astronomy-context}}
\ArrayDiagram
图~\ref{ch30:fig:s45-22-10} 中连接不同台站的基线给出相位比较的空间尺度；这些连接线不表示地面上存在连续的反射面。

实验情景图~\ref{ch30:fig:astronomy-context} 左侧强调各天线分别接收信号，再按时延相关合成；右侧则是激光往返传播的地面实验端。全球站点图把台站分离距离与单台口径区分开：分辨角的量级由基线对波前的投影 $B_\perp$ 决定，即 $\lambda/B_\perp$；稀疏阵列并不等于具有同样直径、连续填满的大镜面。


\subsection{“看见”“分开”和“成像清晰”是不同问题}
一个点源即使远小于分辨尺度，只要送到探测器的光足够多，仍可能被检测到；分辨两个相邻点，则要求它们各自的像斑没有重叠到难以区分。讨论分辨本领时，必须先说明被比较的是两点的间距、两点的角距，还是一个物体本身的大小。

对相距 $s$、距离约为 $L$ 的两个远处物点，小角近似给出 $\Delta\theta\simeq s/L$。圆孔瑞利判据于是写成
\[
 \frac{s}{L}\gtrsim1.22\frac{\lambda}{D},
 \qquad D\gtrsim1.22\frac{\lambda L}{s}.
\]
这里 $s,L,D,\lambda$ 都是长度，两边角度均为无量纲。若结果要求几米口径，却画成手持放大镜大小，应先检查纳米到米、角秒到弧度的换算。

瑞利判据主要针对两个非相干点源、相近亮度与理想圆孔。它是常用的可分辨标准，不是任何图像处理方法都不能突破的绝对信息边界。实际还要考虑大气扰动、像差、噪声和采样；在竞赛题明确要求“仅计衍射限制”时，才忽略这些因素。

\Needspace{7\baselineskip}\begin{origexamplebox}{【视网膜屏】}\label{ch30:ex:s45-22-1}
手机号称“视网膜屏”。（1）若近视者的明视距离仅为 $10\,\mathrm{cm}$，像素密度需达到多少才使单个像素不易被眼分辨？（2）继续增加像素是否有用？
\end{origexamplebox}\par\nopagebreak[3]
\begin{xesolution}
正常眼的典型角分辨率约为 $1'$，即
\[
 \theta_{\min}\approx\frac{\pi}{180\times60}=2.91\times10^{-4}\,\mathrm{rad}.
\]
在 $L=0.10\,\mathrm m$ 处，对应线尺度
\[
 p=L\theta_{\min}\approx2.91\times10^{-5}\,\mathrm m=29\,\mu\mathrm m.
\]
换算为每英寸像素数
\[
 \boxed{\mathrm{ppi}\approx\frac{25.4\,\mathrm{mm}}{0.029\,\mathrm{mm}}\approx8.7\times10^2},
\]
即约 $900\,\mathrm{ppi}$。这是量级估计，个体视力、对比度、子像素排列和观看距离都会改变结果。

当像素角尺度已明显小于眼的分辨角，再增加像素不会在该距离增加可见空间细节，只会提高数据量和功耗；它仍可能有利于更近观看、图像缩放、抗锯齿、彩色子像素采样或后期裁剪，所以“无用”只限于固定距离下的直接视觉分辨。
\end{xesolution}

\Needspace{7\baselineskip}\begin{origexamplebox}{【六角星芒】}\label{ch30:ex:s45-22-2}
拍摄点光源时出现六角星芒。（1）解释成因并判断可推知什么；（2）动画中角色飞走后留下星芒是否合理？
\end{origexamplebox}\par\nopagebreak[3]
\begin{xesolution}
相机光圈不是理想圆孔时，直边会产生方向性衍射：垂直于一条边的两个相反方向形成一对亮芒，平行边产生的方向会重合。六边形光阑的三组平行对边因此可形成六条星芒；但三角形孔径的三条边也可给出六个方向，望远镜的支撑结构也能产生相似图样。因此\boxed{\text{六条星芒不能唯一确定叶片数}}。只能结合孔径结构、星芒方向及曝光条件判断；六边形光圈是一个可能模型，而不是唯一结论。圆整孔径减弱这种定向结构，光学表面划痕和传感器溢出也可能叠加其他伪影。

角色本身飞走不会在空气中凭空留下固定星芒。只有画面经过具有多边形孔径的“摄影机”、且角色或背景含足够亮的近点光源时，星芒作为成像系统的衍射效果才合理；动画中的星芒通常是表现速度或闪光的艺术符号。
\end{xesolution}

\Needspace{7\baselineskip}\begin{origexamplebox}{【提升视力】}\label{ch30:ex:s45-22-3}
看不清远处公交车号时，用拇指和食指做一个很小的缝，透过缝反而看清，为什么？
\end{origexamplebox}\par\nopagebreak[3]
\begin{xesolution}
屈光不正或调焦误差使大瞳孔内不同位置进入的光不能在同一视网膜点会聚，形成较大的弥散圆。小缝挡掉离轴边缘光，只保留方向范围窄的近轴光，弥散圆缩小、景深增大，所以暂时变清楚；这就是针孔效应。

但孔继续缩小时，衍射角 $\sim\lambda/D$ 增大，艾里斑又会变宽，同时亮度下降。因此存在几何模糊与衍射模糊的折中，小孔不能真正治疗近视，也不能无限提升分辨率。
\end{xesolution}

\Needspace{25\baselineskip}
\Needspace{7\baselineskip}\begin{origexamplebox}{【激光照月】}\label{ch30:ex:s45-22-4}
口径 $D=100\,\mu\mathrm m$ 的 He--Ne 激光器发射 $\lambda=632.8\,\mathrm{nm}$ 单色光，对准月球。地月距离 $L=3.84\times10^8\,\mathrm m$，估算到达月球时覆盖面积。
\end{origexamplebox}\par\nopagebreak[3]
\begin{xesolution}
先区分垂直于传播方向的束斑截面与弯曲的月球表面。图~\ref{ch30:fig:s45-22-11} 右图中，月球投影和第一暗环使用相同的长度比例。
\LaserMoonDiagram
把出光口视为圆孔，中央斑角半径约 $1.22\lambda/D$：
\[
 \theta_A=1.22\frac{632.8\times10^{-9}}{100\times10^{-6}}
 \approx7.72\times10^{-3}\,\mathrm{rad}.
\]
在距离 $L$、垂直于光轴的假想平面上，中央斑第一暗环半径为
\[
 r=L\theta_A\approx2.96\times10^6\,\mathrm m,
\]
相应的平面光斑面积为
\[
 \boxed{A\approx\pi r^2\approx2.75\times10^{13}\,\mathrm{m^2}}.
\]
若只作量级估计、去掉圆孔系数 $1.22$，得到约 $1.9\times10^{13}\,\mathrm{m^2}$。但这不是月球表面的受照面积：月球半径约 $R_M=1.737\times10^6\,\mathrm m$，小于上述光斑半径，因此瞄准月心时，整个月球的迎光投影圆盘都落在中央斑以内。投影面积只有 $\pi R_M^2\approx9.48\times10^{12}\,\mathrm{m^2}$；近似平行照明下迎光半球的曲面面积约为 $2\pi R_M^2$，又是另一个量，且各处照度并不相同。题设光束到月球已扩展至数千公里尺度，不能当作几何平行细束；实际月球激光测距须用大得多的发射口径。
\end{xesolution}

\Needspace{7\baselineskip}\begin{origexamplebox}{【数码相机】}\label{ch30:ex:s45-22-5}
边长 $L=35\,\mathrm{mm}$、像素数 $N_p=5\,\mathrm{Mpix}$ 的方形 CCD 配焦距 $f=38\,\mathrm{mm}$ 镜头。光圈数 $F\#=f/D$ 可取 $2,2.8,4,5.6,8,11,16,22$。（1）估计像面最佳空间分辨率 $\Delta x_{\min}$；（2）估计超过何种像素总数后只会造成采样浪费。
\end{origexamplebox}\par\nopagebreak[3]
\begin{xesolution}
衍射斑角半径 $1.22\lambda/D$，落在焦平面的线半径
\[
 \Delta x=1.22\lambda\frac fD=1.22\lambda(F\#).
\]
最佳分辨率取最大光圈 $F\#=2$。以 $\lambda=550\,\mathrm{nm}$ 估算
\[
 \boxed{\Delta x_{\min}\approx1.34\,\mu\mathrm m}.
\]
以每个独立分辨单元占一格作本题量级估算，方形感光面最多包含
\[
 N_{\max}\sim\left(\frac{35\,\mathrm{mm}}{1.34\,\mu\mathrm m}\right)^2
 \approx\boxed{6.8\times10^8\ \text{像素，约 }680\,\mathrm{Mpix}}.
\]
这个 $680\,\mathrm{Mpix}$ 只是按“一格对应一个瑞利间隔”作出的数量级估算，\textbf{不是超过它像素就严格无用的阈值}。以非相干成像的衍射截止频率 $f_c=1/(\lambda F\#)$ 为准，奈奎斯特采样需要 $p\le1/(2f_c)=\lambda F\#/2$；在本题理想参数下为 $0.55\,\mu\mathrm m$，对应约 $4.0\times10^9$ 个方形像素。两种数字不同，是因为“分辨两个点”和“采样全部可传递空间频率”不是同一判据。实际系统还受彩色滤镜、噪声、像差、对比度和数据处理影响，不能给出与这些条件无关的统一“浪费界线”。现有 $5\,\mathrm{Mpix}$ 的像素间距约 $15.7\,\mu\mathrm m$，远大于上述理想衍射尺度。
\end{xesolution}

\MolecularLimitDiagram
图~\ref{ch30:fig:s45-22-12} 对比分子过程的尺度与普通可见光显微镜的分辨尺度；下面讨论的是能否分开局部结构，不是能否检测到分子发出的光。

\Needspace{7\baselineskip}\begin{origexamplebox}{【光学极限】}\label{ch30:ex:s45-22-6}
解释为什么无论怎样改造普通基于可见光的光学显微镜，都无法直接观察 RNA 合成蛋白质的分子过程。
\end{origexamplebox}\par\nopagebreak[3]
\begin{xesolution}
普通远场显微镜的横向衍射极限约为
\[
 \Delta x\simeq0.61\frac{\lambda}{\mathrm{NA}}.
\]
可见光 $\lambda\sim400$--$700\,\mathrm{nm}$，即使高数值孔径 $\mathrm{NA}\sim1.4$，分辨尺度仍约 $200\,\mathrm{nm}$。RNA、核糖体上的局部结构和氨基酸连接过程尺度是数纳米甚至更小，相差两至三数量级，两个分子态的普通光学点扩散斑完全重叠。

增大倍率只会放大同一个模糊斑，不能恢复被孔径截止的空间频率。电子显微镜、X 射线、单分子荧光和超分辨技术可间接研究相关过程，但它们不再是本题所指的普通可见光衍射极限显微镜，且实时、活体和化学标记条件各有约束。
\end{xesolution}

\Needspace{7\baselineskip}\begin{origexamplebox}{【充分遐想】}\label{ch30:ex:s45-22-7}
非鸟类恐龙约在 6600 万年前灭绝。距地球一亿光年以上的文明若想用可见光望远镜分辨地球上的恐龙，估算所需口径。
\end{origexamplebox}\par\nopagebreak[3]
\begin{xesolution}
取恐龙特征尺度 $s\sim10\,\mathrm m$，距离 $L=10^8\,\mathrm{ly}\approx9.46\times10^{23}\,\mathrm m$，角尺度
\[
 \theta\sim\frac sL\approx1.06\times10^{-23}\,\mathrm{rad}.
\]
以 $\lambda=550\,\mathrm{nm}$ 和瑞利判据估算
\[
 D\gtrsim1.22\frac\lambda\theta
 \approx6.3\times10^{16}\,\mathrm m
 \approx\boxed{6.7\ \text{光年}}.
\]
这只是分辨一个十米亮暗结构的理论口径；真实恐龙不自行发光，反射光子极少，恒星、行星大气和背景噪声使要求更加夸张。图~\ref{ch30:fig:distant-past} 的两个场景用于区分距离与所见时代：在忽略宇宙膨胀的估算中，距地球一亿光年的观察者现在接收到约一亿年前离开地球的光，所见仍处于恐龙时代；若要恰看灭绝前后，则距离约为6600万光年。题设“一亿光年以上”是下界，不能据此断言任何更远距离都恰好对应恐龙时代。
\historyfigure{galaxy-dinosaurs-20260923.png}{遥远星系与远古地球的情景示意。所见时代由光行时决定，能否看清个体则取决于角分辨率和光子数。\label{ch30:fig:distant-past}}
\end{xesolution}

\section{多缝干涉与光栅衍射}



设直线上有 $N$ 个等间距散射中心，间距为 $d$，单色平面波正入射。所有中心同相受激，向远场方向 $\theta$ 发出次级波。相邻中心到探测器的程差为 $d\sin\theta$，相位差
\begin{equation}
 \delta=\frac{2\pi}{\lambda}d\sin\theta.
\end{equation}
若每个中心在该方向贡献振幅 $A_0$，总复振幅是等比级数
\[
 \widetilde A=A_0\sum_{j=0}^{N-1}\e^{\mathrm i j\delta}
 =A_0\e^{\mathrm i(N-1)\delta/2}\frac{\sin(N\delta/2)}{\sin(\delta/2)}.
\]
所以
\begin{equation}
 I(\theta)=I_1\left[\frac{\sin(N\delta/2)}{\sin(\delta/2)}\right]^2.
\end{equation}
当 $\delta=2m\pi$ 时，每条刻线的振幅同相相加，取主极大
\begin{equation}
 d\sin\theta_m=m\lambda,\qquad m=0,\pm1,\pm2,\ldots
\end{equation}
且峰值为单刻线强度的 $N^2$ 倍。两个相邻主极大之间有 $N-1$ 个零点和 $N-2$ 个次极大；$N$ 越大，主极大越窄、方向选择越精确。图~\ref{ch30:fig:s45-22-4} 中曲线作了峰值归一化，用来比较宽度而非绝对能量。

等比求和中的整体相因子 $\e^{\mathrm i(N-1)\delta/2}$ 只改变合振幅的相位，不改变强度；真正决定角分布的是正弦比。主极大处分子、分母同时为零，不能把它直接代成 $0/0$ 后说强度为零。令 $\delta=2m\pi+\varepsilon$ 并取 $\varepsilon\to0$，有
\[
 \frac{\sin(N\delta/2)}{\sin(\delta/2)}\longrightarrow(-1)^{m(N-1)}N,
\]
这是因为分子的线性项为 $(-1)^{Nm}N\varepsilon/2$，分母为 $(-1)^m\varepsilon/2$。取模平方后符号消失，所以峰值仍为 $N^2I_1$；与前面的整体相因子相乘后，主极大的合振幅恰为 $NA_0$。相邻主峰之间，分子在 $\delta=2m\pi+2q\pi/N$（$q=1,\ldots,N-1$）处为零，给出 $N-1$ 个暗方向，这也是图~\ref{ch30:fig:s45-22-4} 中主峰随 $N$ 增大而变窄的原因。

若入射光并非正入射，相邻刻线在被照亮时已经带有初始相位差；出射路径又贡献另一部分相位差。两者必须带符号相减，得到本节例~\ref{ch30:ex:s45-22-9} 的斜入射公式。正入射只是在入射项为零时的特殊情况。

\GratingDiagram
图~\ref{ch30:fig:s45-22-2} 在共同出射波前 $BH$ 上比较两束光：投影段 $AH$ 给出相邻缝的程差，因而必须先确定角度的基准线。
\SpectrumDiagram
\GratingCurves

由于 $|\sin\theta_m|\le1$，允许级次满足 $|m|\le d/\lambda$。若 $d<\lambda$，除 $m=0$ 外没有传播级；若 $d\gg\lambda$，会出现许多级；宽波段观察时更需防止不同级次的光谱重叠。这与同一级内的分辨本领不是一回事。光栅用于光谱仪时，通常让间距与目标波长处于合适量级。不同颜色在同一级中的角度不同，$\lambda$ 越长，$|\theta_m|$ 越大。

这里的“没有衍射解”专指没有 $m\ne0$ 的传播级，零级直射或镜面方向仍然存在。对同一正级次，红光波长长，偏离零级更远；紫蓝光更靠近中央。不同级次还可能重叠，例如某波长的二级可与另一波长的三级落在相近角度，所以实际光谱仪还需限制工作级次和波段。图~\ref{ch30:fig:s45-22-3} 中红光位于同一级光谱的外侧，紫蓝光位于内侧。



\subsection{多缝峰的位置、宽度和缺级分别由什么决定}
设共有 $N$ 条相同狭缝，缝宽为 $a$、中心间距为 $d$，正入射单色光。相邻缝向方向 $\theta$ 发出的光相位差为 $\varphi=2\pi d\sin\theta/\lambda$。若单缝的振幅为 $U_1(\theta)$，总振幅是等比和
\[
 U_N=U_1\sum_{j=0}^{N-1}e^{\ii j\varphi}
 =U_1e^{\ii(N-1)\varphi/2}
 \frac{\sin(N\varphi/2)}{\sin(\varphi/2)}.
\]
取模平方可看出两个不同因素：单缝衍射包络决定大范围内的强弱，多缝干涉因子在包络内形成狭窄主峰。主峰方向满足 $d\sin\theta=m\lambda$，只与间距有关；主峰角宽随照亮的缝数增加而减小。

在第 $m$ 级主峰附近，令 $\varphi=2m\pi+\varepsilon$。最近暗纹满足 $N\varepsilon/2=\pm\pi$，故主峰到一侧暗纹的角距近似为
\[
 \delta\theta\simeq\frac{\lambda}{Nd\cos\theta_m}.
\]
如果某个主峰恰好遇到单缝包络的零点，则同时有
\[
 d\sin\theta=m\lambda,\qquad a\sin\theta=p\lambda\quad(p\ne0),
\]
所以 $ma/d=p$ 时该级缺失。不能先按 $|m\lambda/d|\le1$ 统计可传播级次，再把它们全当成可见亮峰。

对两条相近谱线，主峰位置之差约为 $\Delta\theta=m\Delta\lambda/(d\cos\theta_m)$。令这一差值达到单个主峰的半宽尺度，得到光栅的瑞利分辨本领
\[
 \frac{\lambda}{\Delta\lambda}\simeq |m|N.
\]
它说明照亮的总缝数与级次共同决定分辨能力。增大 $d$ 并不会自动提高分辨本领；还需考虑在固定光栅宽度下 $N$ 会如何变化、该级是否缺级以及不同级光谱是否重叠。


\GridPatternDiagram

\Needspace{7\baselineskip}\begin{origexamplebox}{【周期结构】}\label{ch30:ex:s45-22-8}
通过经纬正交织物观看 $20\,\mathrm m$ 外、波长 $550\,\mathrm{nm}$ 的点光源，周围出现正方形亮斑阵列，最近邻表观距离 $12\,\mathrm{cm}$。求织物线间距。
\end{origexamplebox}\par\nopagebreak[3]
\begin{xesolution}
经、纬线分别构成一维光栅，两组衍射的乘积形成二维正方点阵。最近邻一级斑的角度很小：
\[
 \theta\simeq\frac{0.12}{20}=6.0\times10^{-3}\,\mathrm{rad}.
\]
一级光栅条件 $d\sin\theta=\lambda$，故
\[
 \boxed{d\simeq\frac{550\times10^{-9}}{6.0\times10^{-3}}
 =9.17\times10^{-5}\,\mathrm m=0.092\,\mathrm{mm}}.
\]
两正交方向间距若不同，亮斑阵列会成为矩形而非正方形；这就是图~\ref{ch30:fig:s45-22-13} 中位置关系所表达的信息。
\end{xesolution}

\ObliqueGratingDiagram
图~\ref{ch30:fig:s45-ex-22-9} 中 $A,B$ 是相邻透光缝；入射波到达两缝时已有相位差，不能只比较两缝到屏的传播路程。

\Needspace{7\baselineskip}\begin{origexamplebox}{【光栅方程】}\label{ch30:ex:s45-22-9}
证明斜入射透射光栅公式
\[
 a(\sin\theta_m-\sin\theta_i)=m\lambda,
\]
其中 $a$ 为刻线间距，$\theta_i,\theta_m$ 分别为入射光与第 $m$ 级衍射光相对光栅法线的有向角。
\end{origexamplebox}\par\nopagebreak[3]
\begin{xesolution}
取相邻两缝 $A,B$，沿刻线间距方向的位移为 $a$。斜入射波到达 $A,B$ 时已经有初始光程差 $a\sin\theta_i$；从两缝向观察方向 $\theta_m$ 传播时，又产生出射几何程差 $a\sin\theta_m$。按图中的有向角约定，总有效程差为
\[
 \Delta\mathcal L=a\sin\theta_m-a\sin\theta_i.
\]
所有刻线同相加强要求 $\Delta\mathcal L=m\lambda$，于是
\[
 \boxed{a(\sin\theta_m-\sin\theta_i)=m\lambda}.
\]
正入射 $\theta_i=0$ 时退化为 $a\sin\theta_m=m\lambda$。若把入射光画在法线另一侧，$\theta_i$ 本身为负，公式仍自动给出正确的“正弦相加”形式；因此必须使用有向角，而不是事后凭图改正负号。
\end{xesolution}




\section{晶体衍射与布拉格定律}



X 射线是极短波长电磁波，其波长与原子间距相当。晶体中的原子周期排列，每个原子在入射场驱动下成为次级散射源，因而晶体是天然三维光栅。晶体 X 射线衍射既证明了 X 射线的波动性，也提供了研究原子排列和分子结构的工具；冯·劳厄因晶体 X 射线衍射工作获 1914 年诺贝尔物理学奖。

在晶体衍射被发现之前，X 射线虽已用于成像，但其波动本性和波长尺度尚缺少直接证据。晶体一方面提供原子尺度的周期，另一方面又会把散射波集中到少数方向，因此衍射斑的位置可反推晶面间距和取向。反过来，若晶格常数已知，也能由斑点位置测出 X 射线波长。这种“双向标尺”是 X 射线晶体学能够研究微观结构的基础。

把晶体看作间距 $d$ 的平行晶面。射线与晶面夹角为 $\theta$ 时，相邻晶面反射波程差为 $2d\sin\theta$。同相加强条件
\begin{equation}
 2d\sin\theta=n\lambda,\qquad n=1,2,3,\ldots
\end{equation}
称为布拉格定律。这里 $\theta$ 是相对\textbf{晶面}的掠射角，不是相对晶面法线的角；两者互余，混用会导致正弦和余弦互换。

\historyfigure{von-laue-diffraction-v1.png}{马克斯·冯·劳厄（Max von Laue，1879—1960）与晶体 X 射线衍射斑。他因发现晶体中的 X 射线衍射获得1914年诺贝尔物理学奖。\label{ch30:fig:s45-22-5}}
图~\ref{ch30:fig:s45-22-5} 的衍射斑点对应满足相位匹配的方向；劳厄实验利用这些离散方向揭示晶体内部的周期结构。

\begin{center}\begin{minipage}{\linewidth}\centering\begin{tikzpicture}[font=\small]
\begin{scope}[xshift=-3cm]
\fill[black!90](-1.6,-1.6)rectangle(1.6,1.6);
\foreach \r in {.28,.52,.77,1.04,1.35}{\draw[white!70!black,line width=.7pt](0,0)circle(\r);}
\node[below]at(0,-1.8){多取向微晶的环状衍射};
\end{scope}
\begin{scope}[xshift=3cm]
\fill[black!90](-1.6,-1.6)rectangle(1.6,1.6);
\foreach \x in {-1.2,-.6,0,.6,1.2}{\foreach \y in {-1.2,-.6,0,.6,1.2}{
\fill[white!80!black](\x,\y)circle(.035);}}
\node[below]at(0,-1.8){固定取向晶体的离散点斑};
\end{scope}
\end{tikzpicture}\figcap{环纹和点斑是两类不同的衍射结构；这里展示形态，不代表某一种样品的测量数据。}\label{ch30:fig:xray-patterns}\end{minipage}\end{center}

图~\ref{ch30:fig:xray-patterns} 中，大量随机取向微晶把相同散射角的信号分布到圆锥面上，平面探测器截得环纹；固定取向单晶则只在满足相位匹配的方向产生离散斑点。不能把所有晶体的衍射照片都画成相同方阵。

\begin{center}\begin{minipage}{\linewidth}\centering\begin{tikzpicture}[x={(1cm,0cm)},y={(.38cm,.28cm)},z={(0cm,1cm)},font=\small]
\foreach \sh/\plane in {0/1,5/2}{
\begin{scope}[shift={(\sh,0,0)}]
\foreach \a in {0,2}{\foreach \b in {0,2}{
\draw[RuleGray](0,\a,\b)--(2,\a,\b);
\draw[RuleGray](\a,0,\b)--(\a,2,\b);
\draw[RuleGray](\a,\b,0)--(\a,\b,2);}}
\ifnum\plane=1
\filldraw[fill=Blue!16,draw=Blue!75!black,thin](0,1,0)--(2,1,0)--(2,1,2)--(0,1,2)--cycle;
\else
\filldraw[fill=Gold!22,draw=Gold!70!black,thin](0,2,0)--(2,0,0)--(2,0,2)--(0,2,2)--cycle;
\fi
\foreach \y in {2,1,0}{\foreach \z in {0,1,2}{\foreach \x in {0,1,2}{
\pgfmathtruncatemacro{\parity}{mod(\x+\y+\z,2)}
\ifnum\parity=0\filldraw[fill=Blue!45,draw=DeepBlue](\x,\y,\z)circle(2.8pt);
\else\filldraw[fill=Gold!55,draw=Gold!80!black](\x,\y,\z)circle(1.8pt);\fi}}}
\end{scope}}
\begin{scope}[shift={(10,0,0)},x={(.5cm,0cm)},y={(0cm,.5cm)}]
\begin{scope}
\clip(-.2,-.2)rectangle(4.2,4.2);
\foreach \j in {0,1,2,3,4}{\draw[Blue!60](0,\j)--(4,\j);}
\foreach \j in {-4,-2,0,2,4}{\draw[TextGray,dashed](0,\j)--(4,{\j+4});}
\foreach \j in {-2,0,2,4}{\draw[DeepBlue,dotted](0,\j)--(4,{\j+2});}
\end{scope}
\foreach \x in {0,...,4}{\foreach \y in {0,...,4}{
\pgfmathtruncatemacro{\parity}{mod(\x+\y,2)}
\ifnum\parity=0\fill[DeepBlue](\x,\y)circle(1.7pt);
\else\fill[TextGray!60](\x,\y)circle(1.2pt);\fi}}
\node[below,align=center]at(2,-.9){二维截面中的\\不同平行行族};
\end{scope}
\node at(1,0,-.55){$\{100\}$ 类晶面};
\node at(6,0,-.55){$\{110\}$ 类晶面};
\node[anchor=west]at(0,0,3.15){深色大球：$\mathrm{Cl}^-$；浅色小球：$\mathrm{Na}^+$；球大小仅用于区分离子};
\end{tikzpicture}\figcap{氯化钠晶格的同一常规晶胞、两类晶面及二维截面中的平行行族。沿立方轴相邻离子间距为晶格常数的一半，不把立方格线当作共价键。}\label{ch30:fig:nacl}\end{minipage}\end{center}

图~\ref{ch30:fig:nacl} 把钠、氯离子交替排列在立方细网格上；同种离子各自组成面心立方子晶格。不同方向的平行晶面具有不同间距，立方晶格满足 $d_{hkl}=a_{\rm c}/\sqrt{h^2+k^2+l^2}$，其中 $a_{\rm c}$ 是常规晶胞边长。例如 $d_{100}=a_{\rm c}$，$d_{110}=a_{\rm c}/\sqrt2$；具体哪一级反射出现还取决于基元的散射相位，不能只凭几何间距断言每一峰都存在。
\BraggDiagram
图~\ref{ch30:fig:s45-22-7} 的 $HB$ 与 $BK$ 是下层光路多出的两段；各为 $d\sin\theta$，相加才得到布拉格公式中的因子 $2$。



\subsection{布拉格角从晶面量起}
晶面间距为 $d$，光线与晶面的夹角为 $\theta_B$。来自相邻晶面、沿同一方向出射的两束波，下层一路多走的入射段和出射段各为 $d\sin\theta_B$，因此程差是 $2d\sin\theta_B$。相长条件为 $2d\sin\theta_B=m\lambda$。

这里的 $\theta_B$ 不是光线与晶面法线的夹角。若题目给的是相对法线角 $i$，同一公式须写成 $2d\cos i=m\lambda$；若给的是入射束与衍射束之间的散射角，则通常为 $2\theta_B$。三种角之间的换算必须先完成，再代入数值。

晶面反射只是描述晶体周期散射的有效几何模型，并不是每一晶面都像涂银镜面那样独立反射。布拉格条件给出可能相长的方向；晶胞中原子的分布还可能使某些满足该条件的峰消失。仅要求几何衍射角时可以不计算结构因子，但不能据此声称所有级次强度都相同。


\section{偏振态、偏振片与强度}



沿 $z$ 传播的单色平面波可写成两个正交分量
\begin{align}
 E_x&=A_x\cos(kz-\omega t),\\
 E_y&=A_y\cos(kz-\omega t+\delta).
\end{align}
$\delta=0$ 或 $\pi$ 时电场端点沿固定直线运动，是线偏振；$A_x=A_y$ 且
$\delta=\pm\pi/2$ 时为圆偏振；一般情况为椭圆偏振。旋向必须同时说明观察方向或传播方向约定，
只写“左旋、右旋”而不说明视向可能产生歧义。

\begin{center}\begin{minipage}{\linewidth}\centering
\begin{tikzpicture}[scale=.82,>=Latex]
  \begin{scope}[xshift=-4.2cm]
    \draw[axis] (-1.6,0)--(1.7,0) node[right]{$E_x$};
    \draw[axis] (0,-1.4)--(0,1.5) node[above]{$E_y$};
    \draw[curve] (-1.25,-.75)--(1.25,.75);
    \node[below] at (0,-1.65){线偏振};
  \end{scope}
  \begin{scope}
    \draw[axis] (-1.6,0)--(1.7,0) node[right]{$E_x$};
    \draw[axis] (0,-1.4)--(0,1.5) node[above]{$E_y$};
    \draw[curve] (0,0) circle (1.05);
    \draw[->,DeepBlue,thick] (1.05,0) arc[start angle=0,end angle=45,radius=1.05];
    \node[below] at (0,-1.65){圆偏振};
  \end{scope}
  \begin{scope}[xshift=4.2cm]
    \draw[axis] (-1.6,0)--(1.7,0) node[right]{$E_x$};
    \draw[axis] (0,-1.4)--(0,1.5) node[above]{$E_y$};
    \draw[curve,rotate=25] (0,0) ellipse (1.3 and .65);
    \node[below] at (0,-1.65){椭圆偏振};
  \end{scope}
\end{tikzpicture}
\figcap{固定空间点处电场矢量端点的三种轨迹。圆图取 $\delta=+\pi/2$，$+x$ 向右、$+y$ 向上；从 $+z$ 侧迎着传播方向看，随时间增加端点从 $+x$ 转向 $+y$。坐标轴是电场分量，不是光路。}
\end{minipage}\end{center}



理想线偏振片只透过沿透光轴的电场分量。振幅投影为 $E=E_0\cos\theta$，所以强度满足
\begin{equation}
 I=I_0\cos^2\theta.
\end{equation}
自然光通过第一片理想偏振片后强度减半，之后再用马吕斯定律（Malus' law）。两片互相垂直的偏振片之间插入第三片时，
反而会出现透射光，因为中间偏振片先把状态投影到新的方向。


\subsection{先区分自然光、线偏振光和部分偏振光}
自然光没有长期固定的横向振动方向。通过理想线偏振片后，平均透射强度为原来的一半，这是对全部方向取平均的结果。已经线偏振的光通过与其方向成角 $\alpha$ 的检偏器时，振幅先投影为 $E\cos\alpha$，再平方得到 $I=I_0\cos^2\alpha$。这两个公式适用的入射态不同，不能在每经过一片偏振片时都机械地再乘 $1/2$。

两片正交偏振片之间加第三片时，第一片输出已经是线偏振。中间片把它投影到新的透振方向，最后一片再作一次投影，所以透射强度可以从零变成非零。这不表示新偏振片产生了额外光能：每一步的透射强度仍不超过上一步，改变的是送到最后一片的偏振方向。

圆偏振光与自然光通过旋转线偏振片时，透射强度都可以不随角度变化。仅凭这一实验不能区分二者，因为它只比较两个正交方向上的平均强度，丢掉了相位关系。加四分之一波片把相位信息转成振幅方向信息，才可以进一步判别；后面波片一节将给出具体计算。

\Needspace{7\baselineskip}\begin{origexamplebox}{三片偏振片}
自然光强度为 $I_0$，依次通过透光轴为 $0^\circ,45^\circ,90^\circ$ 的三片理想偏振片。求最终强度，并与只放首尾两片比较。
\end{origexamplebox}\par\nopagebreak[3]
\begin{xesolution}
第一片后 $I_1=I_0/2$。第二片与第一片夹角 $45^\circ$：
\[
 I_2=I_1\cos^245^\circ=\frac{I_0}{4}.
\]
第三片与第二片仍夹角 $45^\circ$，所以
\[
 \boxed{I_3=I_2\cos^245^\circ=\frac{I_0}{8}}.
\]
若只放 $0^\circ$ 与 $90^\circ$ 两片，第二片前的线偏振方向与其轴垂直，理想透射强度为零。
\end{xesolution}

界面反射也能选择偏振：第27章已经由菲涅尔公式推出布儒斯特角，并说明该角度处 p 偏振反射消失。
因此水面、玻璃等自然反光常以 s 偏振为主，摄影偏振镜可通过转动压低相应方向的眩光。



\section{波片与相位延迟}



双折射材料对两条本征偏振方向具有不同折射率 $n_f,n_s$。厚度为 $d$ 的波片引入相位差
\begin{equation}
 \Delta\phi=\frac{2\pi}{\lambda_0}(n_s-n_f)d.
\end{equation}
$\Delta\phi=\pi/2$ 为四分之一波片，可在线偏振与椭圆/圆偏振之间转换；$\Delta\phi=\pi$ 为半波片，可把线偏振方向旋转两倍于入射偏振与快轴的夹角。


\subsection{振幅分解与相位延迟不能混成一步}
以波片两条本征偏振方向为坐标轴。入射线偏振与第一轴成角 $\alpha$，首先作振幅分解：$E_x=E_0\cos\alpha$、$E_y=E_0\sin\alpha$。无吸收理想波片不会把这两个分量的振幅任意放大，而是在厚度 $t$ 上让它们积累不同相位：
\[
 \delta=\frac{2\pi}{\lambda_0}(n_y-n_x)t.
\]
输出是否线偏振、椭圆偏振或圆偏振，取决于振幅比与 $\delta$ 两项。四分之一波片只保证相差为 $\pi/2$ 的奇数倍；若入射方向不是与两轴成 $45^\circ$，两个振幅不等，输出一般是椭圆而非圆。

为避免旋向歧义，写电场参数式时还应规定观察方向与时间因子。例如沿传播方向看、取
$E_x=A\cos\omega t$、$E_y=B\cos(\omega t-\delta)$，就可以代入 $t=0$ 及一个微小正时间，直接判断电矢量怎样转动。只写“左旋”而不说沿光看还是迎光看，可能把同一物理态叫成相反的名称。

\Needspace{7\baselineskip}\begin{origexamplebox}{四分之一波片产生圆偏振}
线偏振光入射四分之一波片，偏振方向与快轴夹角为 $45^\circ$。证明出射光为圆偏振，并说明条件改变后的结果。
\end{origexamplebox}\par\nopagebreak[3]
\begin{xesolution}
令快轴为 $x$、慢轴为 $y$，入射偏振的有向角为 $\alpha$。按前面的 $kz-\omega t$ 约定，去掉两分量的共同相位后，可写
\[
 E_x=E_0\cos\alpha\cos\psi,\qquad
 E_y=E_0\sin\alpha\cos(\psi+\pi/2)=-E_0\sin\alpha\sin\psi.
\]
当 $\sin\alpha\cos\alpha\ne0$ 时，消去 $\psi$ 得
\[
 \frac{E_x^2}{E_0^2\cos^2\alpha}
 +\frac{E_y^2}{E_0^2\sin^2\alpha}=1.
\]
在 $\alpha=45^\circ$ 时，两半轴均为 $E_0/\sqrt2$，因此轨迹为圆。完整的分类是：
$\alpha=45^\circ+90^\circ m$ 时为圆偏振；$\alpha=90^\circ m$ 时只激发一个本征分量，仍是线偏振；其余角度为非退化椭圆偏振。这里 $m$ 为整数。$135^\circ$ 同样能产生圆偏振，但两分量的相对符号改变，旋向与 $45^\circ$ 时相反。

旋向可用固定点的时间变化判断，而不依赖含混的左右旋名称：在上述约定和 $\alpha=45^\circ$ 下，取 $\psi=0$ 的时刻，$E_x>0,E_y=0$；随时间增加，$\psi$ 减小，$E_y$ 转为正值，故端点从 $+x$ 向 $+y$ 转动。若波长改变，相位延迟一般不再恰为 $\pi/2$，还须同时检验等幅与正交相位两项条件。
\end{xesolution}



同一真空波长 $\lambda_0$ 下，相位积累为 $2\pi nd/\lambda_0$。折射率较小的本征方向相速度较大，称快轴；
折射率较大的方向称慢轴。若 $\Delta n=n_s-n_f>0$，厚度 $d$ 引入的慢轴相对快轴延迟为
\begin{equation}
 \delta=\frac{2\pi}{\lambda_0}\Delta n\,d.
\end{equation}
实际波片允许延迟多加 $2\pi m$，所以四分之一波片和半波片可以是零级或多级。多级波片对波长和温度更敏感。

\begin{center}\begin{minipage}{\linewidth}\centering
\begin{tikzpicture}[scale=.9,>=Latex]
  \draw[axis] (-.2,0)--(5.5,0) node[right]{快轴 $f$};
  \draw[axis] (0,-.2)--(0,3.5) node[above]{慢轴 $s$};
  \draw[vector] (0,0)--(3.5,2.4) node[above right]{$\vect E_{\rm in}$};
  \draw[component] (0,0)--(3.5,0) node[midway,below]{$E_f$};
  \draw[component] (0,0)--(0,2.4) node[midway,left]{$E_s$};
  \draw[guide] (3.5,0)--(3.5,2.4)--(0,2.4);
  \draw[->,DeepBlue] (1.1,0) arc[start angle=0,end angle=34.5,radius=1.1];
  \node[clearlabel] at (1.3,.36) {$\alpha$};
\end{tikzpicture}
\figcap{入射线偏振在波片快、慢轴上的分解。波片改变两分量相位，不改变理想情况下的分量振幅。}\label{ch30:fig:waveplate-components}
\end{minipage}\end{center}
图~\ref{ch30:fig:waveplate-components} 的两个投影给出快、慢轴上的分量振幅。仅当相对相位为 $\pm\pi/2$ 时，这两个方向才是偏振椭圆的主轴；一般相位差下椭圆会倾斜，不能把任意坐标分量直接称为椭圆半轴。

\Needspace{7\baselineskip}\begin{origexamplebox}{零级四分之一波片的厚度}
材料在波长 $\lambda_0=600\,\mathrm{nm}$ 时有 $n_s-n_f=0.0090$。求最薄四分之一波片的厚度；若制成厚度为该值的5倍，判断其等效延迟。
\end{origexamplebox}\par\nopagebreak[3]
\begin{xesolution}
最薄四分之一波片要求 $\delta=\pi/2$：
\[
 d_0=\frac{\lambda_0}{4\Delta n}
 =\frac{600\,\mathrm{nm}}{4\times0.0090}
 =\boxed{16.7\,\mu\mathrm m}.
\]
厚度 $5d_0$ 的延迟为 $5\pi/2=2\pi+\pi/2$，对单色完全偏振光等效于四分之一波片；但它是更高阶波片，波长改变时额外的 $2\pi$ 也会变化，因此色散敏感性不同。
量纲上波长除以无量纲折射率差仍为长度。
\end{xesolution}

\begin{center}\begin{minipage}{\linewidth}\centering
\begin{tikzpicture}[scale=.9,>=Latex]
  % 纯几何：两方向关于水平快轴镜像，角度严格为正负20度。
  % 横向8cm，字号沿用正文图样；两角标位于弧外，直线不穿标签。
  \draw[axis] (-3,0)--(3.4,0) node[right]{快轴 $x$};
  \draw[axis] (0,-1.65)--(0,1.8) node[above]{慢轴 $y$};
  \draw[DeepBlue,very thick] (-2.6,{-2.6*tan(20)})--(2.6,{2.6*tan(20)});
  \draw[TextGray,very thick,dashed] (-2.6,{2.6*tan(20)})--(2.6,{-2.6*tan(20)});
  \draw[DeepBlue] (1.0,0) arc[start angle=0,end angle=20,radius=1.0];
  \draw[TextGray] (1.5,0) arc[start angle=0,end angle=-20,radius=1.5];
  \node[anchor=west] at (1.2,1.05) {$+\alpha$};
  \node[anchor=west] at (1.6,-1.2) {$-\alpha$};
  \node at (-2.1,1.55) {出射偏振};
  \node at (-2.1,-1.55) {入射偏振};
\end{tikzpicture}
\figcap{半波片把入射线偏振关于快轴镜像，因此偏振方位转过 $2\alpha$。}\label{ch30:fig:half-wave}
\end{minipage}\end{center}
图~\ref{ch30:fig:half-wave} 中实线与虚线关于快轴对称，故两者方位差为 $2\alpha$。

\Needspace{7\baselineskip}\begin{origexamplebox}{半波片旋转偏振方向}
线偏振方向与半波片快轴夹角为 $\alpha=20^\circ$。半波片后接一检偏器；为使透射强度最大，检偏器相对原偏振方向应转多少度？若检偏器保持原方向，透射强度是入射到检偏器前强度的多少？
\end{origexamplebox}\par\nopagebreak[3]
\begin{xesolution}
把入射场分解到快、慢轴，半波片使两分量相差 $\pi$，等价于偏振方向关于快轴反射。因此出射偏振相对原方向转过
\[
 2\alpha=\boxed{40^\circ}.
\]
检偏器转到新方向时透射最大。若仍沿原方向，根据马吕斯定律
\[
 \boxed{\frac I{I_{\rm before}}=\cos^2(40^\circ)\approx0.587}.
\]
$\alpha\to0$ 时出射方向不变；快轴转动 $\Delta\alpha$ 会使偏振方向转动 $2\Delta\alpha$，与半波片常用于偏振旋转相符。
\end{xesolution}

\tcbset{breakable/.default=true}\begin{starread}{晶体光学与琼斯矩阵}
各向异性晶体中介电响应是张量，波矢方向与能流方向一般不重合。单轴晶体具有寻常光和非寻常光，后者折射率随传播方向变化，
因此一个入射光束可分裂成两束。完全偏振的相干变换可用琼斯矢量和矩阵表示：
\[
 \vect E=\begin{pmatrix}E_x\\E_y\end{pmatrix},\qquad
 \mathbf J_{\rm pol}(\theta)=\vect e_\theta\vect e_\theta^{\mathsf T}.
\]
这里的琼斯（Jones）矢量只适用于完全偏振光；部分偏振光不能由单个琼斯矢量完整描述，需要斯托克斯参数或相干矩阵。吸收、散射和退偏器件还需更一般的穆勒（Mueller）矩阵。
\end{starread}\tcbset{breakable/.default=false}


\Needspace{8\baselineskip}\begin{origexamplebox}{用波片测量部分偏振光}
\label{ch30:ex:depth-cai67}
部分椭圆偏振光沿 $z$ 方向传播，通过理想线偏振检偏器。透振轴沿 $x$ 轴时透射强度最大，为 $1.5I_0$；沿 $y$ 轴时最小，为 $I_0$。
（1）透振轴与 $x$ 轴成角 $\theta$ 时，透射强度怎样变化？是否依赖光束中自然光所占比例？
（2）在检偏器前加入两本征轴沿 $x,y$ 的理想四分之一波片，发现检偏器透振轴与 $x$ 轴成 $30^\circ$ 时透射强度最大。求这个最大值及原入射光中非偏振成分所占比例。把部分偏振光看成一束完全偏振光与一束互不相干的自然光的叠加。
\end{origexamplebox}\par\nopagebreak[3]
\begin{xesolution}
第一问不能直接把 $1.5I_0$ 代入马吕斯定律，因为光束不是纯线偏振。用 $X,Y$ 表示完全偏振部分沿椭圆两主轴的分量强度，$U$ 表示自然光部分的总强度。自然光通过任意方向的理想检偏器都贡献 $U/2$。由两个极值方向得
\[
 X+\frac U2=1.5I_0,\qquad
 Y+\frac U2=I_0,
 \qquad X-Y=0.5I_0.
\]
椭圆主轴分量相差 $\pi/2$，投影后的交叉项时间平均为零。因此
\[
 I(\theta)=X\cos^2\theta+Y\sin^2\theta+\frac U2
 =\boxed{I_0+0.5I_0\cos^2\theta}.
\]
也可写为 $1.25I_0+0.25I_0\cos2\theta$。给定测得的两个极值后，曲线不再显含 $U$；但是这不表示 $U=0$。仅旋转线偏振片，还不足以确定有多少自然光。

四分之一波片使两个主轴分量的相对相位再改变 $\pm\pi/2$，于是完全偏振部分变成线偏振；自然光部分仍为自然光，总强度也不变。图~\ref{ch30:fig:depth-cai67} 将两个阶段画在同一横向坐标约定下。最大透射方向为 $+30^\circ$，说明波片后线偏振方向为该方向（模 $\pi$），故
\[
 \tan30^\circ=\sqrt{Y/X},\qquad Y=X/3.
\]
和 $X-Y=0.5I_0$ 联立，得到
\[
 X=0.75I_0,\qquad Y=0.25I_0,\qquad U=1.5I_0.
\]

% Figure design: transverse polarization states, not a longitudinal ray diagram.
% Width <=12cm; two equal 4cm panels, separated by 1.6cm. Labels >=9pt.
% Left ellipse semi-axes sqrt(.75)*1.7 and sqrt(.25)*1.7, ratio sqrt(3).
% Right line at 30 degrees, same energy radius sqrt(X+Y)*1.7.
% No rotation arrow: incident helicity is not inferred by this measurement.
% Solid blue field loci, thin gray axes, transparent labels, no connectors.
% ASCII: [Ex horizontal, Ey vertical: ellipse]  [same axes: / line at 30deg].
\begin{center}\begin{minipage}{\linewidth}\centering
\begin{tikzpicture}[>=Latex,every node/.style={fill=none,font=\small,inner sep=2pt}]
 \begin{scope}
  \draw[->,gray] (-1.85,0)--(1.95,0) node[right]{$E_x$};
  \draw[->,gray] (0,-1.18)--(0,1.48) node[above]{$E_y$};
  \draw[DeepBlue,thick] (0,0) ellipse[x radius={1.7*sqrt(.75)},y radius=.85];
  \node[below,align=center] at (0,-1.5){波片前的完全偏振部分\\振幅比 $\sqrt{X/Y}=\sqrt3$};
 \end{scope}
 \begin{scope}[xshift=6.4cm]
  \draw[->,gray] (-1.85,0)--(1.95,0) node[right]{$E_x$};
  \draw[->,gray] (0,-1.18)--(0,1.48) node[above]{$E_y$};
  \draw[DeepBlue,thick] (210:1.7)--(30:1.7);
  \draw (.75,0) arc[start angle=0,end angle=30,radius=.75];
  \node at (15:1.18){$30^\circ$};
  \node[below,align=center] at (0,-1.5){波片后的完全偏振部分\\两个分量同相，成为线偏振};
 \end{scope}
\end{tikzpicture}
\figcap{蓝线为电矢量端点轨迹，不是光路；未画出的自然光经波片后仍非偏振。\label{ch30:fig:depth-cai67}}
\end{minipage}\end{center}

原入射总强度为 $X+Y+U=2.5I_0$，所以非偏振成分比例为
\[
 \boxed{\frac{U}{X+Y+U}=\frac{1.5}{2.5}=60\%}.
\]
检偏器沿新线偏振方向时，完全偏振部分 $X+Y=I_0$ 全部透过，自然光透过一半，于是
\[
 \boxed{I_{\max}'=X+Y+\frac U2=1.75I_0}.
\]
与其垂直时只有自然光的一半通过，$I_{\min}'=0.75I_0$；两者之和仍为 $2.5I_0$，与总入射强度一致。波片提高了旋转检偏器时的最大透射值，但没有创造能量，只把原先藏在相位关系中的信息转成可选择透过的线偏振方向。
\end{xesolution}


\section{单轴晶体与双折射}



单轴晶体有一条光轴。这里“光轴”是沿该方向传播不发生双折射的方向，不是一根可见的几何轴，也不等同于器件的中心线。
主截面是光轴与波矢共同张成的平面。寻常光（ordinary wave，o光）的折射率 $n_o$ 与传播方向无关；非寻常光（extraordinary wave，e光）的有效折射率随波矢与光轴夹角 $\theta$ 改变：
\begin{equation}
 \frac1{n_e^2(\theta)}
 =\frac{\cos^2\theta}{n_o^2}+\frac{\sin^2\theta}{n_e^2}.
\end{equation}
式中右端的 $n_e$ 是波矢垂直光轴时的主折射率。$\theta=0$ 时 $n_e(0)=n_o$，所以沿光轴传播两束简并。
$n_e>n_o$ 称正单轴晶体，$n_e<n_o$ 称负单轴晶体。

寻常光的电场垂直主截面，非寻常光的电位移 $\vect D$ 位于主截面内。在各向异性介质中，坡印廷矢量
$\vect S$ 一般不平行于波矢 $\vect k$，所以e光的能流方向会相对波法线发生走离；画双折射光路时必须区分“波矢折射”和“光束能流方向”。

\begin{center}\begin{minipage}{\linewidth}\centering
\begin{tikzpicture}[scale=.9,>=Latex]
  \draw[axis] (-3.1,0)--(3.2,0) node[right]{$k_x$};
  \draw[axis] (0,-2.2)--(0,2.4) node[above]{$k_z$（光轴）};
  \draw[DeepBlue,very thick] (0,0) circle (1.65);
  \draw[TextGray,very thick] (0,0) ellipse (2.25 and 1.65);
  \draw[vector] (0,0)--(1.55,1.35) node[above right]{$\vect k$};
  \draw[->,DeepBlue] (0,1.0) arc[start angle=90,end angle=41,radius=1.0];
  \node[clearlabel] at (.42,1.02) {$\theta$};
  \node[clearlabel,anchor=east] at (-2.0,1.65) {o光};
  \draw[thin,TextGray] (-1.94,1.54)--(-1.01,1.30);
  \node[clearlabel,anchor=west] at (2.45,-1.0) {e光};
  \draw[thin,TextGray] (2.39,-.96)--(1.94,-.84);
\end{tikzpicture}
\figcap{正单轴晶体的等频波矢面截面。圆对应o光，横向较长的椭圆对应e光；两曲线在光轴方向相切。}\label{ch30:fig:index-ellipse}
\end{minipage}\end{center}
图~\ref{ch30:fig:index-ellipse} 中沿光轴的两个截距相同，直观对应 $n_e(0)=n_o$。

下题改用负单轴晶体，此时e光椭圆的横向半轴小于o光圆半径，与图中的正单轴情形相反；两者沿光轴仍简并。

\Needspace{7\baselineskip}\begin{origexamplebox}{非寻常光的方向折射率}
某负单轴晶体 $n_o=1.66$、$n_e=1.49$。波矢与光轴夹角为 $45^\circ$。求e光的有效折射率，并检验它应落在哪个范围。
\end{origexamplebox}\par\nopagebreak[3]
\begin{xesolution}
代入方向折射率公式：
\[
 \frac1{n_e^2(45^\circ)}
 =\frac12\left(\frac1{1.66^2}+\frac1{1.49^2}\right).
\]
因此
\[
 \boxed{n_e(45^\circ)\approx1.57}.
\]
结果位于 $1.49$ 与 $1.66$ 之间；若算出超出两主折射率范围，通常是把折射率椭圆的半轴或角度定义颠倒了。
$\theta\to0$ 时公式给出 $n_o$，$\theta\to90^\circ$ 时给出主值 $n_e$。
\end{xesolution}

\begin{center}\begin{minipage}{\linewidth}\centering
\begin{tikzpicture}[scale=.92,>=Latex]
  % 纯几何图：两束从同一入射点分开，出射面横移由同一组 d,rho 计算。
  % 布局：入射 -> 晶片(0<=z<=5) -> 平行出射；标注在外侧，无白底。
  \fill[DeepBlue!5] (0,-1.4) rectangle (5,2.1);
  \draw[very thick] (0,-1.4)--(0,2.1) (5,-1.4)--(5,2.1);
  \coordinate (A) at (0,0); \coordinate (B) at (5,0);
  \coordinate (C) at (5,{5*tan(15)});
  \draw[vector] (-2,0)--(A) node[midway,above=5pt]{入射光};
  \draw[DeepBlue,thick] (A)--(C);
  \draw[vector] (C)--(7,{5*tan(15)}) node[right]{e光能流};
  \draw[component] (A)--(B);
  \draw[vector] (B)--(7,0) node[right]{o光};
  \draw[guide] (0,-.4)--(0,2.4);
  \draw[TextGray] (1.1,0) arc[start angle=0,end angle=15,radius=1.1];
  \node at (1.5,.50) {$\rho$};
  \draw[<->] (5.65,0)--(5.65,{5*tan(15)}) node[midway,right=4pt]{$\Delta x$};
  \draw[<->] (0,-1.7)--(5,-1.7) node[midway,below]{$d$};
  \draw[TextGray,thick] (2.5,-1.05)--(3.15,-.40);
  \node[below] at (2.8,-1.05) {光轴方向};
  \node[above left] at (A){$A$};\node[below left] at (B){$B$};
  \node[above left] at (C){$C$};
\end{tikzpicture}
\figcap{正入射单轴晶片的一种双折射几何。o光不偏折，e光能流可因走离而产生横向分离。}\label{ch30:fig:double-refraction}
\end{minipage}\end{center}

\Needspace{7\baselineskip}\begin{origexamplebox}{晶片后的两束横向分离}
厚度为 $d$ 的平行晶片中，o光沿法线传播，e光的能流与法线夹角为走离角 $\rho$。忽略出射面的额外横移，求两束在出射面处的间距，并给出小角近似。
\end{origexamplebox}\par\nopagebreak[3]
\begin{xesolution}
图~\ref{ch30:fig:double-refraction} 中 $A$ 是共同入射点，$AB=d$ 沿晶片法线，$BC=\Delta x$ 为出射面内的横移。由直角三角形 $ABC$ 得
\[
 \boxed{\Delta x=d\tan\rho}.
\]
当 $|\rho|\ll1$（弧度）时
\[
 \boxed{\Delta x\simeq d\rho}.
\]
间距符号由走离方向决定，通常报告其绝对值。$d$ 乘无量纲三角函数仍是长度；$\rho\to0$ 时两束重合，符合沿光轴或无各向异性时的极限。
\end{xesolution}



\section{衍射积分与部分相干光（选读）}

\subsection{衍射积分与傅里叶光学}


惠更斯--菲涅耳原理（Huygens--Fresnel principle）把孔径面上的每一点看成次级波源。省略只影响整体相位和归一化的常数后，观察点 $P$ 的场可写成
\begin{equation}
 U(P)\simeq\frac1{\ii\lambda}\iint_{\mathcal A}
 U(Q)\,K(\theta,\theta')\frac{e^{\ii k r}}r\dd S,
\end{equation}
其中 $r=|PQ|$，$K$ 是倾斜因子。这个公式不是“光线绕过障碍物”的几何规则，而是孔径面各微元复振幅的相干积分；遮挡改变积分区域，进而改变整个屏上的场。

在夫琅禾费远场（Fraunhofer region），或把观察屏放在理想透镜后焦面时，远场振幅与孔径面出射场的二维傅里叶变换成正比。若入射场为 $U_{\rm in}(x,y)$、孔径透射函数为 $t(x,y)$，则
\begin{equation}
 U(\theta_x,\theta_y)\propto
 \iint U_{\rm in}(x,y)t(x,y)e^{-\ii k(x\sin\theta_x+y\sin\theta_y)}\dd x\dd y.
\end{equation}
只有均匀平面波正入射时，$U_{\rm in}$ 才可提出积分号外。此时宽为 $a$、高为 $b$ 的矩形孔产生
\begin{equation}
 \frac I{I_0}=
 \left(\frac{\sin\alpha}{\alpha}\right)^2
 \left(\frac{\sin\beta}{\beta}\right)^2,
 \quad
 \alpha=\frac{\pi a\sin\theta_x}{\lambda},
 \quad
 \beta=\frac{\pi b\sin\theta_y}{\lambda}.
\end{equation}
孔径在某一方向越窄，远场在该方向展开得越宽；这是一对傅里叶共轭尺度，而不是成像系统失焦。

\tcbset{breakable/.default=true}\begin{starread}{菲涅耳波带与波带片}
在近场计算轴上点 $P$ 的光强时，可把波面按到 $P$ 的光程每增加 $\lambda/2$ 分成一个半波带。相邻带的贡献相位近似相反，开放或遮住某些带就能控制轴上叠加。

对正入射平面波和焦距 $f$，平面波带片第 $n$ 条边界的半径近似满足
\[
 \sqrt{f^2+r_n^2}-f=\frac{n\lambda}{2},
 \qquad r_n^2=n\lambda f+\frac{n^2\lambda^2}{4}
 \simeq n\lambda f.
\]
交替透光的振幅型波带片在 $f$ 处形成主焦点，并伴随较弱的高阶实、虚焦点；它依靠衍射聚焦，不能简单套用玻璃薄透镜的折射率公式。

\begin{center}\begin{minipage}{\linewidth}\centering
\begin{tikzpicture}[scale=.86,>=Latex]
  \begin{scope}[xshift=-3.8cm]
    \foreach \n in {2,4,6}{
      \fill[DeepBlue!22,even odd rule] (0,0) circle ({.65*sqrt(\n)})
        (0,0) circle ({.65*sqrt(\n-1)});}
    \foreach \n in {1,...,6}{\draw[RuleGray](0,0)circle({.65*sqrt(\n)});}
    \node[below] at (0,-1.85){正视图：交替透光区};
  \end{scope}
  \draw[very thick] (0,-1.65)--(0,1.65) node[above]{波带片};
  \coordinate (F) at (5,0);
  \fill[Ink] (F) circle (1.6pt) node[right]{$F$};
  \foreach \y in {-1.379,-1.028,0,1.028,1.379}{
    \draw[-{Latex[length=2mm]},thick,DeepBlue] (-1.3,\y)--(0,\y)--(F);
  }
  \draw[<->,guide] (.12,-1.48)--(4.88,-1.48) node[midway,below]{$f$};
\end{tikzpicture}
\figcap{振幅型波带片的区域结构与主焦点。侧视图只画开放区贡献的代表性光线，实际计算是复振幅积分。}
\end{minipage}\end{center}
\end{starread}\tcbset{breakable/.default=false}

\subsection{相干矩阵与部分偏振}


完全偏振光可由一个确定的复振幅矢量描述；自然光或部分偏振光还含随机涨落，需要使用相干矩阵
\begin{equation}
 \mathbf J=
 \left\langle
 \begin{pmatrix}E_x\\E_y\end{pmatrix}
 \begin{pmatrix}E_x^*&E_y^*\end{pmatrix}
 \right\rangle
 =\begin{pmatrix}J_{xx}&J_{xy}\\J_{yx}&J_{yy}\end{pmatrix}.
\end{equation}
对角元是两个正交分量的强度，$\operatorname{tr}\mathbf J=I$ 为总强度，非对角元记录两分量的相干性。取一种常用旋向约定，斯托克斯参数（Stokes parameters）为
\begin{equation}
 S_0=J_{xx}+J_{yy},\quad S_1=J_{xx}-J_{yy},\quad
 S_2=2\Re J_{xy},\quad S_3=2\Im J_{xy}.
\end{equation}
偏振度
\begin{equation}
 \mathcal P=\frac{\sqrt{S_1^2+S_2^2+S_3^2}}{S_0}
\end{equation}
在 $0$ 与 $1$ 之间。$\mathcal P=1$ 等价于 $\det\mathbf J=0$，表示完全偏振；理想自然光满足 $\mathbf J=(I/2)\mathbf 1$，故 $\mathcal P=0$。这也说明“每一瞬间电场都有方向”并不意味着长时间平均后的光必为完全偏振。


\endgroup
% EMOpticsQuality20261023
% QualityRevision20261023
% ModernLayout
\begingroup
\hypersetup{bookmarksdepth=3}
\tcbset{breakable/.default=false,after={\par\nobreak\smallskip}}
\newcommand{\ModernAnswerLabel}{{\bfseries\color{DeepBlue}解答}}
\RenewDocumentEnvironment{xesolution}{+b}{%
  \par\nobreak\smallskip\nobreak\begingroup\small
  \clubpenalties 2 10000 10000\displaywidowpenalties 2 10000 10000\noindent
  #1\par\endgroup\smallskip
}{}
\chapter{狭义相对论}

\section{相对性原理与光速不变}

狭义相对论首先用惯性参考系建立时空关系，也可在惯性坐标中描述加速粒子；它不把引力造成的时空弯曲包含在内。它建立在两条实验支持的原则上：物理定律在所有惯性系中形式相同；真空光速
$c$ 与光源和惯性观察者的运动无关。它并不是在牛顿力学（Newtonian mechanics）上额外增加一个“高速修正力”，而是改变不同惯性系之间
时空坐标的变换关系。

十九世纪末的干涉实验试图比较相互垂直方向的往返光程。图~\ref{ch31:fig:mm-lab}展示实验操作的历史情景；
分光镜把入射光分入两条近乎正交的臂，光在端镜反射后返回并合束。两臂的往返时间差会改变干涉相位；转动装置可比较不同方向的结果。
精确光路见第29章“迈克耳孙干涉仪”一节，不能从情景图中量取几何数据。
\historyfigure[0.82\linewidth]{assets/generated/michelson-morley-lab.png}{
迈克耳孙—莫雷干涉实验的复原情景示意（非史料照片）。厚重平台用于隔振，转动支撑用于比较不同方位的光程。\label{ch31:fig:mm-lab}}

设 $S'$ 相对 $S$ 沿 $x$ 轴以速度 $v$ 运动，且两原点在 $t=t'=0$ 重合。洛伦兹变换（Lorentz transformation）为
\begin{equation}
 x'=\gamma(x-vt),\qquad
 t'=\gamma\left(t-\frac{vx}{c^2}\right),
 \qquad y'=y,\quad z'=z,
\end{equation}
这里的 $t$ 不是某一处观察者“看到光到达眼睛”的时刻，而是 $S$ 系同步钟网在事件发生地点记录的坐标时间。
各系可用光信号同步：从钟A发光、在钟B立即反射、再回到A，令B反射时的读数等于A发射与接收读数的平均值。
两系采用同样的尺和钟，坐标轴平行，且 $|v|<c$。其中
\begin{equation}
 \gamma=\frac1{\sqrt{1-v^2/c^2}}.
\end{equation}
为什么时间变换含有位置 $x$？时空均匀性要求惯性坐标之间为线性变换。
$S'$ 原点满足 $x=vt$，所以可先写 $x'=a(x-vt)$、$t'=bt+dx$。
从共同原点分别向左右发光，光速不变要求
\[
 a(c-v)=c(b+dc),\qquad a(-c-v)=-c(b-dc).
\]
两式相加、相减，得到 $b=a$、$d=-av/c^2$，于是
$t'=a(t-vx/c^2)$。由相对性原理和空间反向对称性，逆变换的同一系数仍为 $a$，但 $v$ 换成 $-v$：
\[
 x=a(x'+vt')=a^2(1-v^2/c^2)x.
\]
取与 $v=0$ 连续相接的正根，便有 $a=\gamma$。这也说明不能在保留光速不变的同时继续要求 $t'=t$。

将变换式代入并展开，交叉项相消，有
\[
 c^2t'^2-x'^2=\gamma^2(1-v^2/c^2)(c^2t^2-x^2)=c^2t^2-x^2.
\]
因此它保持任意两个事件的时空间隔
\begin{equation}
 c^2\Delta t^2-\Delta x^2-\Delta y^2-\Delta z^2
\end{equation}
不变。间隔为正称类时、为零称类光、为负称类空。类时事件可由低于光速的物体连接；类光事件可由光信号连接；类空事件不能相互传递不超光速的因果信号。
低速极限 $v/c\to0$ 时，$\gamma\to1$，恢复伽利略变换（Galilean transformation）。


\subsection{先选事件，再选参考系}
洛伦兹变换变换的是同一事件的坐标。解题时先为事件命名，再列出哪一系中哪个坐标差为零。常用的三个约束不能相互替代：匀速钟的两次读数在该钟的静止惯性系满足 $\Delta x'=0$；运动杆在实验室测长要求 $\Delta t=0$；杆的固有长度要求在杆系同时测两端。
若 $\Delta x=600\,\mathrm m$、$\Delta t=1.00\,\mu\mathrm s$，取 $c\simeq3.00\times10^8\,\mathrm{m/s}$，则 $c\Delta t=300\,\mathrm m<\Delta x$，两事件类空。
使两事件同时的换系速度由 $\Delta t'=0$ 得到
\[
 v=\frac{c^2\Delta t}{\Delta x}=0.50c.
\]
再改用 $v=0.80c$，有
\[
 \Delta t'=\frac53(1.00-1.60)\,\mu\mathrm s=-1.00\,\mu\mathrm s.
\]
事件先后可以反转，但两事件没有亚光速因果联系，因此不构成因果矛盾。反过来，若两事件类时，使它们同地的速度为 $v=\Delta x/\Delta t$，其绝对值小于 $c$；此时才存在连接两事件的钟的静止系。

\subsection{光钟与时间延缓}

一只光钟由相距 $L$ 的两面平行镜组成，光脉冲在镜间往返一次记一个周期。钟的静止系中半周期为
$\Delta\tau/2=L/c$。这里固有周期 $\Delta\tau$ 对应下镜处发射A、返回C这两个事件，而不是空间分离的发射A、反射B之间的间隔。
若钟沿平行镜面、垂直于静止系光路的方向以速度 $v$ 运动，实验室看到光走斜边。横向镜间距不变，往返两段对称，所以
\begin{equation}
 \left(c\frac{\Delta t}{2}\right)^2
 =L^2+\left(v\frac{\Delta t}{2}\right)^2,
\end{equation}
从而 $\Delta t=\gamma\Delta\tau$。这不是光在实验室“走得更快”，而是光速仍为 $c$、路径变长。

\begin{center}\begin{minipage}{\linewidth}\centering
% 核心结论：同一发射反射事件在运动系形成光程直角三角形。
% 原型：纯几何对比；保留现行蓝灰风格，宽约12cm，字号10pt。
% 草图：静止 A↑B | 运动 A↗B；下方尺寸与标题分两行，L标在右侧。
\begin{tikzpicture}[font=\small,>=Latex]
 \begin{scope}
  \coordinate (A) at (1,0); \coordinate (B) at (1,2.4);
  \draw[very thick] (.2,0)--(1.8,0) (.2,2.4)--(1.8,2.4);
  \draw[vector] (A)--(B);
  \node[left=8pt] at (1,1.2) {$L=c\Delta\tau/2$};
  \node[below right=5pt] at (A) {A};\node[above right=5pt] at (B) {B};
  \node at (1,-1.45) {钟的静止系};
 \end{scope}
 \begin{scope}[xshift=5.7cm]
  \coordinate (A) at (0,0);\coordinate (B) at (1.8,2.4);\coordinate (D) at (1.8,0);
  \draw[very thick] (-.65,0)--(.65,0) (1.15,2.4)--(2.45,2.4);
  \draw[vector] (A)--(B);
  \node[above left=5pt] at (.9,1.2) {$c\Delta t/2$};
  \draw[guide] (B)--(D) (A)--(D);
  \draw[TextGray,<->] (2.85,0)--(2.85,2.4) node[midway,right=4pt] {$L$};
  \draw[TextGray,<->] (0,-.55)--(1.8,-.55) node[midway,below=3pt] {$v\Delta t/2$};
  \node[below left=4pt] at (A) {A};\node[above right=4pt] at (B) {B};
  \node at (1,-1.45) {实验室系};
 \end{scope}
\end{tikzpicture}

\figcap{横向光钟的半个周期：A发射、B反射。右图两段粗镜面分别表示对应事件时刻的位置，不是同一时刻的完整钟。几何示意未按下题数值绘制。\label{ch31:fig:light-clock}}
\end{minipage}\end{center}
图~\ref{ch31:fig:light-clock} 右侧斜边、竖边和横边组成直角三角形，直接给出时间延缓因子。

\begin{origexamplebox}{光钟的实验室周期}
光钟两镜间距 $L=0.50\,\mathrm m$，钟以 $0.80c$ 相对实验室横向运动。求钟的固有周期与实验室所见周期，并求半周期内光在实验室走过的路程。

\begin{center}\begin{minipage}{.96\linewidth}\centering
\begin{tikzpicture}[>=Latex,font=\small,every node/.style={fill=none},scale=.88]

\draw[very thick] (-.8,0)--(.8,0) (-.8,1.8)--(.8,1.8);
\draw[DeepBlue,->] (0,0)--(0,1.7);\draw[DeepBlue,->] (.2,1.8)--(.2,.1);
\draw[<->] (-1.1,0)--(-1.1,1.8) node[midway,left]{$L$};
\draw[->] (1.3,.9)--(3,.9) node[midway,above]{$v=0.80c$};\node[below] at (0,-.35){光钟自身的两面镜};

\end{tikzpicture}
\figcap{题图：两镜间距沿垂直运动方向量取，整只光钟相对实验室横向运动。\label{ch31:quality:q1}}
\end{minipage}\end{center}
\end{origexamplebox}
\begin{xesolution}
\begin{center}\begin{minipage}{.96\linewidth}\centering
\begin{tikzpicture}[>=Latex,font=\small,scale=.88]
\draw[DeepBlue,thick,->](0,0)--(2,1.5);
\draw[DeepBlue,thick,->](2,1.5)--(4,0);
\draw[dashed](0,0)--(4,0)(2,0)--(2,1.5);
\fill(0,0)circle(1.5pt)node[below]{$A$};
\fill(2,1.5)circle(1.5pt)node[above]{$B$};
\fill(4,0)circle(1.5pt)node[below]{$C$};
\draw[<->](0,-.6)--(2,-.6)node[midway,below]{$v\Delta t/2$};
\draw[dashed](2,1.5)--(4.7,1.5)(4,0)--(4.7,0);
\draw[<->](4.7,0)--(4.7,1.5)node[midway,right]{$L$};
\node[above left]at(1,.75){$c\Delta t/2$};
\end{tikzpicture}
\figcap{解图：实验室中光的往返路径。$A,B,C$ 属于不同时间的三个事件。\label{ch31:quality:a1}}
\end{minipage}\end{center}

\ModernAnswerLabel\quad 固有周期
\[
 \Delta\tau=\frac{2L}{c}=\frac{1.0\,\mathrm m}{c}\approx3.34\,\mathrm{ns}.
\]
$v=0.80c$ 时 $\gamma=1/\sqrt{1-0.8^2}=5/3$，故
\[
 \boxed{\Delta t=\gamma\Delta\tau\approx5.56\,\mathrm{ns}}.
\]
半周期内光程
\[
 \boxed{s=c\Delta t/2=\gamma L\approx0.833\,\mathrm m}.
\]
用直角三角形检查：横移 $v\Delta t/2\approx0.667\,\mathrm m$，与竖边 $0.50\,\mathrm m$ 合成的斜边确为 $0.833\,\mathrm m$。

图~\ref{ch31:quality:a1} 中，实验室半周期 $\Delta t/2$ 内，光程为三角形斜边 $c\Delta t/2$，镜面横移为 $v\Delta t/2$；两镜的横向间距 $L$ 不受纵向尺缩影响。
\end{xesolution}

\begin{center}\begin{minipage}{\linewidth}\centering
% 核心结论：时间轴为运动原点世界线，空间轴为同时线。
% 纯几何，约11x6cm，字号10pt；ct↑，x→，双向光线；P投影到直角轴。
% beta=0.4；事件P=(4,3)远离斜轴；所有标签透明。
\begin{tikzpicture}[font=\small,scale=.93,>=Latex]
 \draw[axis] (-3.5,0)--(5.6,0) node[right] {$x$};
 \draw[axis] (0,-.4)--(0,4.9) node[above] {$ct$};
 \draw[DeepBlue,dashdotted] (0,0)--(3.8,3.8);
 \draw[DeepBlue,dashdotted] (0,0)--(-3.2,3.2);
 \node[above left] at (-3.2,3.2) {$x=-ct$};
 \node[above right] at (3.8,3.8) {$x=ct$};
 \draw[Blue,thick] (0,0)--(1.8,4.5) node[above=4pt] {$ct'$};
 \draw[Blue,thick] (0,0)--(5.1,2.04) node[right=4pt] {$x'$};
 \coordinate (P) at (4,3);
 \draw[guide] (0,3)--(P)--(4,0);
 \fill (P) circle (1.6pt);\node[right=5pt] at (P) {P};
 \node[left=4pt] at (0,3) {$ct_P$};\node[below=4pt] at (4,0) {$x_P$};
 \node[below left] at (0,0) {O};
\end{tikzpicture}

\figcap{两个惯性系的时空轴。虚线将事件P投影到地面系的x、ct轴；两条点划线为向左右传播的光信号。\label{ch31:fig:spacetime-axes}}
\end{minipage}\end{center}
图~\ref{ch31:fig:spacetime-axes} 中，$ct'$ 轴是 $x'=0$ 的世界线，即 $x=vt$；$x'$ 轴是 $t'=0$ 的同时线，即 $ct=(v/c)x$。
采用相同的 $x,ct$ 长度比例尺时，世界线的横向变化率为 $\dd x/\dd(ct)=v/c$；图上通常所说的纵横斜率则是 $\dd(ct)/\dd x=c/v$（$v\ne0$）。两条光线的斜率均为 $\pm1$。
图上的欧氏距离却不是时空间隔，也不能用向斜轴作欧氏垂线的方法读取 $x',ct'$；斜坐标须按洛伦兹变换标定。

\section{固有时、时间延缓与长度收缩}

沿同一只钟世界线的两个读数事件，钟本身记录的时间称固有时 $\Delta\tau$。若该钟在 $S$ 中匀速运动，两事件满足
$\Delta x=v\Delta t$，在钟的静止系中则 $\Delta x'=0$。由间隔不变，
\[
 c^2\Delta\tau^2=c^2\Delta t^2-v^2\Delta t^2,
 \qquad \Delta\tau=\Delta t\sqrt{1-v^2/c^2},
\]
故
\begin{equation}
 \Delta t=\gamma\Delta\tau.
\end{equation}
运动钟在该惯性系坐标时间中的读数间隔更小。这里比较的是同一对事件，不能把不同位置的两只未同步时钟混入。

物体在自身静止系 $S'$ 测得的长度为固有长度 $L_0$。在地面系 $S$ 同时记录两端位置，必须取 $\Delta t=0$，令其位置差为 $L$。
对这两个测量事件，静止杆两端的坐标差仍为 $L_0$，因此
\[
 L_0=\Delta x'=\gamma(\Delta x-v\Delta t)=\gamma L,
\]
得到
\begin{equation}
 L=\frac{L_0}{\gamma}.
\end{equation}
长度收缩只发生在运动方向。它与同时性的相对性不可分：不同惯性系“同时测两端”对应不同的两件事件。

\begin{origexamplebox}{高速不稳定粒子的寿命与飞行距离}
某不稳定粒子在自身静止系平均寿命为 $\tau_0=2.2\,\mu\mathrm s$，在实验室以 $0.98c$ 运动。求实验室测得平均寿命和平均飞行距离。

\begin{center}\begin{minipage}{.96\linewidth}\centering
\begin{tikzpicture}[>=Latex,font=\small,every node/.style={fill=none},scale=.88]

\draw[->] (-.5,0)--(5,0) node[right]{$x$};\fill (0,0) circle (1.8pt) node[below]{产生事件};
\draw[DeepBlue,->] (0,.2)--(3.5,.2) node[midway,above]{$v=0.98c$};\draw[dashed] (4,-.5)--(4,.8) node[above]{衰变事件（平均位置待求）};
\node[below] at (2,-.7){粒子自身静止系平均寿命 $\tau_0$};

\end{tikzpicture}
\figcap{题图：同一粒子的产生和衰变事件，固有寿命在粒子静止系测量。\label{ch31:quality:q2}}
\end{minipage}\end{center}
\end{origexamplebox}
\begin{xesolution}
\ModernAnswerLabel\quad \[
 \gamma=\frac1{\sqrt{1-0.98^2}}\approx5.03.
\]
实验室寿命
\[
 \boxed{\tau=\gamma\tau_0\approx11.1\,\mu\mathrm s}.
\]
平均飞行距离
\[
 \boxed{\ell=v\gamma\tau_0
 =\frac{0.98c\times2.2\,\mu\mathrm s}{\sqrt{1-0.98^2}}
 \approx3.25\,\mathrm{km}}.
\]
在粒子静止系中也可得到相同相遇事件：实验室到探测位置的距离收缩为原来的 $1/\gamma$。
具体地，把产生事件记为A、衰变事件记为B；平均寿命是大量粒子的统计平均，
并非每个粒子都恰好在 $11.1\,\mu\mathrm s$ 后衰变。忽略减速时，每次飞行都有
$\Delta x=v\Delta t$，取平均后仍有 $\langle\Delta x\rangle=v\langle\Delta t\rangle$。
\begin{center}\begin{minipage}{\linewidth}\centering
% Contract: two simultaneous spatial snapshots, not two world-line diagrams.
% Same events A (creation) and B (decay); lengths refer to ground markers P,Q.
% Transparent labels, 12cm print width; no decorative backgrounds.
\begin{tikzpicture}[>=Latex,font=\small]
 \draw[->] (0,0)--(5.1,0) node[right]{$x$};
 \fill (0,0) circle(2pt) (4.4,0) circle(2pt);
 \node[below=5pt] at (0,0){P};\node[below=5pt] at (4.4,0){Q};
 \draw[DeepBlue,thick,->] (.3,.65)--(1.8,.65) node[midway,above=3pt]{$v$};
 \node[above=5pt] at (3.2,.65){粒子运动};
 \draw[<->] (0,-.65)--(4.4,-.65) node[midway,below=3pt]{$\ell=v\tau$};
 \node at (2.2,1.45){实验室系：P、Q静止};
 \begin{scope}[yshift=-3.1cm]
  \draw[->] (0,0)--(5.1,0) node[right]{$x'$};
  \fill (0,0) circle(2pt) (2.4,0) circle(2pt);
  \node[below=5pt] at (0,0){P};\node[below=5pt] at (2.4,0){Q};
  \draw[DeepBlue,thick,->] (2.4,.7)--(1.1,.7) node[midway,above=3pt]{$-v$};
  \node[above=4pt] at (4,.5){地面标记左移};
  \draw[<->] (0,-.65)--(2.4,-.65) node[midway,below=3pt]{$\ell'=\ell/\gamma=v\tau_0$};
  \node at (2.2,1.45){粒子系：粒子留在 $x'=0$};
 \end{scope}
\end{tikzpicture}

\figcap{产生时刻两系中的地面标记。P是产生位置，Q取在平均飞行距离处；长度仅作示意。\label{ch31:fig:particle-lifetime}}
\end{minipage}\end{center}
图~\ref{ch31:fig:particle-lifetime} 中，地面标记P、Q的固有间距为 $\ell$。
粒子系中它们整体向左运动，同时测得间距
$\ell'=\ell/\gamma\approx0.646\,\mathrm{km}$。
Q到达静止粒子所需时间为
\[
 t'=\frac{\ell'}v=\frac{\ell}{\gamma v}=\tau_0=2.2\,\mu\mathrm s.
\]
这样两系对同一相遇事件给出一致结果：实验室使用延长的寿命，粒子系使用收缩的地面距离。
\end{xesolution}


\subsection{加速运动中的固有时累积}
匀速公式用于整段旅程之前，要检查钟是否全程相对所选惯性系以同一速率运动，即 $\gamma$ 是否恒定。速度随时间变化时，每个足够短的时间段满足
\[
 \dd\tau=\sqrt{1-v(t)^2/c^2}\,\dd t,\qquad
 \tau=\int_{t_A}^{t_B}\sqrt{1-v(t)^2/c^2}\,\dd t.
\]
这给出理想钟沿既定世界线的累计读数；不是把全程平均速度放入 $\gamma$。只要各段之间钟的读数连续，便可逐段相加。
例如地面与飞船在A相遇，飞船以 $0.60c$ 匀速飞出3.00年（地面计时），在远端短时掉头，再以同速返回并在B重逢。忽略掉头所占时间，地面钟累计6.00年，飞船钟累计
\[
 \tau=2(3.00\,\mathrm{yr})\sqrt{1-0.60^2}=4.80\,\mathrm{yr}.
\]
A、B是同一对重逢事件，两只钟的累计读数可以直接比较。飞船去、回两段对应不同惯性系，不能把“运动钟慢”的互易说法在中途无条件拼接；差异应由各自世界线上的固有时积分计算。实际有限时长掉头还需把掉头段的积分加入。

\section{同时性的相对性与速度变换}

在 $S$ 中同时发生、相隔 $\Delta x$ 的两事件，在 $S'$ 中的时间差为
\begin{equation}
 \Delta t'=-\gamma\frac{v\Delta x}{c^2}.
\end{equation}
所以空间分离事件的同时性不是绝对的；但因果相关的类时或光速间隔事件次序不会因惯性系改变。

若物体在 $S$ 中速度为 $u_x,u_y$，先对同一段轨迹作微分：
\[
 \dd x'=\gamma(\dd x-v\dd t),\quad
 \dd y'=\dd y,\quad
 \dd t'=\gamma\left(\dd t-\frac{v\dd x}{c^2}\right)
 =\gamma(1-vu_x/c^2)\dd t.
\]
分别用位移微分除以同一个 $\dd t'$，得到
\begin{equation}
 u_x'=\frac{u_x-v}{1-u_xv/c^2},\qquad
 u_y'=\frac{u_y}{\gamma(1-u_xv/c^2)}.
\end{equation}
当 $u_x=c$ 时仍有 $u_x'=c$；低速极限恢复 $u_x'=u_x-v$。

\begin{origexamplebox}{车厢两端闪光的同时性}
在地面系 $S$ 中，一列沿 $+x$ 方向以 $v$ 运动的车厢前、后端在同一地面时刻发出闪光，两事件相距车厢在地面系的长度 $L$。求车厢系 $S'$ 中两次发光的时间差，并判断哪一端先发光。
\begin{center}\begin{minipage}{\linewidth}\centering
% 核心结论：题图给地面同时发光、地面长度L和车速，不给同时线或光相遇点。
% 纯几何，蓝灰，10pt，目标宽10cm；草图 A|------L------|B，v向右，t=0。
\begin{tikzpicture}[font=\small,>=Latex]
 \draw[thick] (0,0) rectangle (6,1.1);
 \fill[DeepBlue] (0,.55) circle(2pt);\fill[DeepBlue] (6,.55) circle(2pt);
 \node[left=5pt,align=center] at (0,.55) {后端A\\发光};
 \node[right=5pt,align=center] at (6,.55) {前端B\\发光};
 \draw[vector] (2,1.65)--(4,1.65) node[midway,above] {$v$};
 \draw[TextGray,<->] (0,-.55)--(6,-.55) node[midway,below] {$L$（地面同时测量）};
 \node at (3,.55) {地面系 $t=0$};
\end{tikzpicture}

\figcap{地面同一时刻的车厢位置及两端发光事件；箭头只给车速方向。\label{ch31:fig:train-question}}
\end{minipage}\end{center}
\end{origexamplebox}
\begin{xesolution}
\begin{center}\begin{minipage}{\linewidth}\centering\raggedright\ModernAnswerLabel\par\smallskip\centering
% 核心结论：地面同时的端点发光在车厢系不同时；光相遇点并非车中点。
% 纯几何时空图，约10x6cm，10pt；L=4,beta=1/3,ctmax=4.2。
% 后端x=ct/3，前端x=4+ct/3；光线A->P<-B，P=(2,2)。
\begin{tikzpicture}[font=\small,>=Latex]
 \def\trainL{4}\def\trainBeta{.3333333333}\def\trainT{4.2}
 \coordinate (A) at (0,0);\coordinate (B) at (\trainL,0);
 \coordinate (P) at ({\trainL/2},{\trainL/2});
 \draw[axis] (-.35,0)--(6.6,0) node[right] {$x$};
 \draw[axis] (0,0)--(0,4.9) node[above] {$ct$};
 \draw[TextGray,thick] (A)--({\trainBeta*\trainT},\trainT) node[above=6pt] {后端};
 \draw[TextGray,thick] (B)--({\trainL+\trainBeta*\trainT},\trainT) node[above=6pt] {前端};
 \draw[DeepBlue,very thick] (A)--(P)--(B);
 \draw[vector] (.65,.65)--(1.3,1.3);
 \draw[vector] (3.35,.65)--(2.7,1.3);
 \draw[Blue,dashed] (A)--(5.8,{5.8*\trainBeta});
 \node[right=5pt] at (5.8,{5.8*\trainBeta}) {$t'=0$};
 \fill (P) circle (1.6pt);\node[above=5pt] at (P) {P};
 \fill (A) circle (1.4pt);\fill (B) circle (1.4pt);
 \node[below left=3pt] at (A) {A};\node[below right=3pt] at (B) {B};
 \draw[TextGray,<->] (0,-.65)--(\trainL,-.65) node[midway,below=3pt] {$L$};
\end{tikzpicture}

\figcap{车厢两端在地面系同时发光。灰线为同向平行的车端世界线；蓝线在P相遇，但P不是运动中的车厢中点。\label{ch31:fig:simultaneity}}
\end{minipage}\end{center}


取前端事件减后端事件，则在地面系 $\Delta t=0$、$\Delta x=L>0$。洛伦兹变换给出
\[
 \Delta t'=\gamma\left(\Delta t-\frac{v\Delta x}{c^2}\right)
 =\boxed{-\gamma\frac{vL}{c^2}}.
\]
负号表示前端事件在车厢系更早。图~\ref{ch31:fig:simultaneity} 取后端发光事件A为原点、前端事件B为 $(L,0)$。
过A的车厢系同时线为 $ct=\beta x$，B在这条线下方，所以 $t'_B<t'_A=0$。
地面系中两光线为 $x=ct$ 与 $x=L-ct$，故相遇点P满足 $ct=L/2,x=L/2$；该时刻车厢中点却位于 $x=L/2+vL/(2c)$。
不能把光线相遇点、车厢中点及某一端的接收事件混作同一个事件。
$v\to0$ 或 $L\to0$ 时差异消失；$vL/c^2$ 具有时间量纲。
\end{xesolution}

\begin{center}\begin{minipage}{\linewidth}\centering
\begin{tikzpicture}[scale=.95,>=Latex]
  \draw[axis] (-.2,0)--(5.6,0) node[right]{$x$};
  \draw[axis] (0,-.2)--(0,4.3) node[above]{$y$};
  \draw[vector] (0,0)--(4.3,3.2) node[above right]{$\vect u$};
  \draw[component] (0,0)--(4.3,0) node[midway,below]{$u_x$};
  \draw[component] (4.3,0)--(4.3,3.2) node[midway,right]{$u_y$};
  \draw[vector] (.5,-.65)--(2.5,-.65) node[midway,below]{$\vect v$（换系）};
\end{tikzpicture}
\figcap{二维速度变换的分量约定示意，不按下题数值绘制。纵向、横向分量的变换分母相同，但横向还含 $1/\gamma$。\label{ch31:fig:velocity-components}}
\end{minipage}\end{center}
图~\ref{ch31:fig:velocity-components} 中换系速度只沿 $x$，所以 $u_x,u_y$ 必须分别代入不同的变换式。

\begin{origexamplebox}{横向速度分量的变换}
在实验室系中，粒子速度为 $u_x=0.60c$、$u_y=0.30c$。求相对实验室沿 $+x$ 方向以 $v=0.50c$ 运动的惯性系所测速度分量和速率。

\begin{center}\begin{minipage}{.96\linewidth}\centering
\begin{tikzpicture}[>=Latex,font=\small,every node/.style={fill=none},scale=.88]

\draw[->] (0,0)--(3,0) node[right]{$x$};\draw[->] (0,0)--(0,2) node[above]{$y$};
\draw[DeepBlue,thick,->] (0,0)--(2,1) node[right]{$\boldsymbol u=(.60c,.30c)$};
\draw[->] (.6,-.7)--(2.5,-.7) node[midway,below]{$S'\text{ 沿 }+x\text{ 以 }.50c\text{ 运动}$};

\end{tikzpicture}
\figcap{题图：实验室中的二维速度与沿 $+x$ 的参考系变换。\label{ch31:quality:q4}}
\end{minipage}\end{center}
\end{origexamplebox}
\begin{xesolution}
\ModernAnswerLabel\quad $\gamma_v=1/\sqrt{1-0.5^2}=2/\sqrt3$，共同分母
$1-u_xv/c^2=1-0.30=0.70$。故
\[
 u_x'=\frac{0.60c-0.50c}{0.70}=0.1429c,
\]
\[
 u_y'=\frac{0.30c}{(2/\sqrt3)(0.70)}=0.3712c.
\]
所以
\[
 \boxed{u'=\sqrt{u_x'^2+u_y'^2}\approx0.398c}.
\]
结果小于 $c$。若用伽利略变换而错误地令 $u_y'=u_y$，会漏掉横向分量的时间与同时性变换因子。
图~\ref{ch31:fig:velocity-components} 中的 $u_y$ 来自粒子相邻两位置的横向坐标差，
虽然 $\dd y'=\dd y$，计时却变成
\[
 \dd t'=\frac2{\sqrt3}\,(0.70)\dd t\approx0.8083\dd t.
\]
同样的横向位移除以更短的坐标时间，就有 $u_y'>u_y$。
这不是横向又施加了一个力，而是对同一段运动更换测量参考系。
\begin{center}\begin{minipage}{\linewidth}\centering
% Contract: same scale 6cm per c, components (.6,.3) and (.142857,.371154).
% Axes form mathematical 2D x/y orientation; dashed legs are components.
\begin{tikzpicture}[>=Latex,font=\small]
 \foreach \shift/\dx/\dy/\name in {0/3.6/1.8/S,6/.857143/2.226922/{S'}}{
  \begin{scope}[xshift=\shift cm]
   \draw[->] (-.2,0)--(4,0);\draw[->] (0,-.2)--(0,2.65);
   \draw[DeepBlue,very thick,->] (0,0)--(\dx,\dy);
   \draw[dashed,gray] (\dx,\dy)--(\dx,0);
   \node at (1.8,-.85){$\name$};
  \end{scope}
 }
 \node[above right=4pt] at (3.6,1.8){$\vect u$};
 \node[below=4pt] at (1.8,0){$0.60c$};
 \node[right=5pt] at (3.6,.9){$0.30c$};
 \node[above right=4pt] at (6.857143,2.226922){$\vect u'$};
 \node[below=5pt] at (6.45,0){$0.1429c$};
 \node[right=5pt] at (6.857143,1.15){$0.3712c$};
\end{tikzpicture}

\figcap{本题同一粒子在两系中的速度，左右使用同一比例尺。水平、竖直方向分别是各系的 $x,y$ 方向。\label{ch31:fig:velocity-two-frames}}
\end{minipage}\end{center}
图~\ref{ch31:fig:velocity-two-frames} 显示 $u_x'$ 减小而 $u_y'$ 增大，
同时总速率从 $\sqrt{0.60^2+0.30^2}c\approx0.671c$ 变为 $0.398c$。
\end{xesolution}

\begin{origexamplebox}{相向飞船看到的相对速度}
两飞船在地面系中分别以 $0.80c$ 向右和 $0.60c$ 向左运动。求右行飞船测得左行飞船的速度。

\begin{center}\begin{minipage}{.96\linewidth}\centering
\begin{tikzpicture}[>=Latex,font=\small,every node/.style={fill=none},scale=.88]

\draw[->] (-1,0)--(5,0) node[right]{地面系 $x$};
\fill (0,0) circle (2pt) node[below]{$A$};\fill (4,0) circle (2pt) node[below]{$B$};
\draw[DeepBlue,->] (0,.4)--(1.7,.4) node[midway,above]{$+0.80c$};\draw[DeepBlue,->] (4,.4)--(2.6,.4) node[midway,above]{$-0.60c$};

\end{tikzpicture}
\figcap{题图：两船速度先都相对地面取值，右行船 $A$ 是所求新参考系。\label{ch31:quality:q5}}
\end{minipage}\end{center}
\end{origexamplebox}
\begin{xesolution}
\begin{center}\begin{minipage}{.96\linewidth}\centering
\begin{tikzpicture}[>=Latex,font=\small,every node/.style={fill=none},scale=.88]

\draw[->] (-1,0)--(5,0) node[right]{地面系 $x$};
\fill (0,0) circle (2pt) node[below]{$A$};\fill (4,0) circle (2pt) node[below]{$B$};
\draw[DeepBlue,->] (0,.4)--(1.7,.4) node[midway,above]{$+0.80c$};\draw[DeepBlue,->] (4,.4)--(2.6,.4) node[midway,above]{$-0.60c$};

\end{tikzpicture}
\figcap{解图：先固定谁是观察者，再把另一船的有号地面速度变换过去。\label{ch31:quality:a5}}
\end{minipage}\end{center}

\ModernAnswerLabel\quad 取右行为正。地面系中左行飞船 $u_x=-0.60c$，右行飞船参考系速度参数 $v=0.80c$。所以
\[
 u_x'=\frac{-0.60c-0.80c}{1-(-0.60)(0.80)}
 =-\frac{1.40}{1.48}c\approx-0.946c.
\]
故相对速率为 $\boxed{0.946c}$，小于 $c$。直接相加得到 $1.40c$ 是伽利略速度变换在高速下的错误使用。

图~\ref{ch31:quality:a5} 中，$A$ 系由沿 $+x$ 以 $0.80c$ 的洛伦兹变换得到，故把 $B$ 的地面速度 $-0.60c$ 代入同一速度合成式；所得负号表示在 $A$ 系仍向左。
\end{xesolution}

\section{光行差与相对论多普勒效应}

平面波相位
\begin{equation}
 \Phi=\omega t-\vect k\cdot\vect r
\end{equation}
必须在同一事件给出相同相位。将逆变换 $t=\gamma(t'+vx'/c^2)$、$x=\gamma(x'+vt')$ 代入，
\[
 \Phi=\gamma(\omega-vk_x)t'
 -\gamma\left(k_x-\frac{v\omega}{c^2}\right)x'-k_yy'-k_zz'.
\]
与 $\Phi=\omega't'-\vect k'\cdot\vect r'$ 比较系数，便得到
\begin{equation}
 \omega'=\gamma(\omega-vk_x),\qquad
 k_x'=\gamma\left(k_x-\frac{v\omega}{c^2}\right),\qquad
 k_y'=k_y.
\end{equation}
对真空光波有 $|\vect k|=\omega/c$。若 $\theta$ 是 $S$ 中光线与 $+x$ 轴的夹角，则
\begin{equation}
 \cos\theta'=\frac{\cos\theta-\beta}{1-\beta\cos\theta},\qquad
 \beta=\frac vc,
\end{equation}
其中 $\theta'$ 是 $S'$ 中光线与 $+x'$ 轴的夹角；此式来自 $\cos\theta'=ck'_x/\omega'$。
频率同时满足 $f'=\gamma f(1-\beta\cos\theta)$，这两式必须使用同一套角度约定。
将 $(\omega/c,\vect k)$ 合写称为波四矢量（wave four-vector），四矢量的一般定义见本章后面的选读。

沿视线方向，令 $f_0$ 为光源静止系发射频率、$\beta>0$ 为源与观察者的相对速率。在光源系取光沿 $+x$ 传播。
观察者沿 $+x$ 远离时，代入 $\cos\theta=1$ 有
$f_{\rm rec}=\gamma f_0(1-\beta)$；观察者迎光运动则把换系速度改成 $-v$。利用
$\gamma(1\mp\beta)=\sqrt{(1\mp\beta)/(1\pm\beta)}$，得到
\begin{equation}
 f_{\rm rec}=f_0\sqrt{\frac{1-\beta}{1+\beta}},
 \qquad
 f_{\rm app}=f_0\sqrt{\frac{1+\beta}{1-\beta}},
\end{equation}
下标分别表示远离和接近。例如相对速率为 $0.60c$，远离时接收频率为 $f_0/2$，接近时为 $2f_0$。
$\beta\to0$ 时二者都回到 $f_0$。它同时含有经典传播距离效应和时间延缓，不能只用“波峰追赶”得到完整结果。




\subsection{从谱线频移反求速度}
已知光源静止系的谱线波长 $\lambda_0$，观察者测得沿视线远离的光源谱线为 $\lambda_{\rm obs}$。先取比值 $r=\lambda_{\rm obs}/\lambda_0=f_0/f_{\rm rec}$，再反解：
\[
 r^2=\frac{1+\beta}{1-\beta},\qquad
 \beta=\frac{r^2-1}{r^2+1}.
\]
例如 $\lambda_{\rm obs}=1.50\lambda_0$ 时，$\beta=5/13\approx0.385$，而把红移 $z=r-1=0.50$ 直接当作 $v/c$ 会高估速度。$z\ll1$ 时才有 $\beta\simeq z$。
这里把频移全部归因于狭义相对论中的相对运动；不含引力红移和宇宙膨胀的传播效应。


\paragraph{到达间隔包含钟慢与传播路程两部分}
在观察者惯性系中，光源沿视线以 $v>0$ 远离。两次发射在光源自身钟上相隔 $\Delta\tau=1/f_0$，在观察者系中相隔 $\Delta t=\gamma\Delta\tau$。第二次发射时光源又远离了 $v\Delta t$，第二束光还要多走这段路，所以
\[
 \Delta t_{\rm rec}=\Delta t+\frac{v\Delta t}{c}
 =\gamma(1+\beta)\Delta\tau,\qquad
 r=\frac{f_0}{f_{\rm rec}}=\gamma(1+\beta).
\]
平方并用 $\gamma^2=(1-\beta^2)^{-1}$，即可得到前面的 $r^2=(1+\beta)/(1-\beta)$。仅使用时间延缓会漏掉传播路程变化，也就无法分辨靠近与远离。
对沿视线靠近的源，可在同一约定下取 $\beta<0$，此时 $r<1$，反解式自然给出负速度；$r=1$ 给出零径向速度，$r\to\infty$ 时 $\beta\to1^-$，不会反推出超光速运动。非沿视线运动还须使用含观测角的完整多普勒关系。


\subsection{“横向”必须注明在哪个参考系}
取 $S$ 为光源静止系，$S'$ 观察者沿 $+x$ 以 $v$ 运动。两系的光线角度分别记为 $\theta,\theta'$。若光源系中横向发射，即 $\theta=90^\circ$，则
\[
 f'=\gamma f_0,\qquad \cos\theta'=-\beta.
\]
若观察者系中光线横向，即 $\theta'=90^\circ$，光行差式却要求 $\cos\theta=\beta$，从而
\[
 f'=\gamma f_0(1-\beta^2)=\frac{f_0}{\gamma}.
\]
两种频移相反，是因为所选的光线方向不同。在 $\beta=0.60$ 时，前者为 $1.25f_0$、后者为 $0.80f_0$；代入同一套角度变换即可互相核对。沿视线的接近、远离公式不能直接用于这种斜向问题。

\section{相对论动量与能量}

为了使动量守恒在所有惯性系中保持相同形式，质量为 $m$ 的粒子具有
\begin{equation}
 \vect p=\gamma m\vect v,
 \qquad E=\gamma mc^2.
\end{equation}
静止能量 $E_0=mc^2$，动能
\begin{equation}
 K=(\gamma-1)mc^2.
\end{equation}
能量动量关系为
\begin{equation}
 E^2=p^2c^2+m^2c^4.
\end{equation}
无静质量粒子满足 $E=pc$ 并以 $c$ 传播。低速展开
$K\simeq mv^2/2+3mv^4/(8c^2)+\cdots$，首项回到经典动能。

三维力仍定义为 $\vect F=\dd\vect p/\dd t$。对恒定静质量粒子的能量动量关系求导，
\[
 2E\frac{\dd E}{\dd t}=2c^2\vect p\cdot\frac{\dd\vect p}{\dd t},
 \qquad \frac{c^2\vect p}{E}=\vect v,
\]
两式相结合得到
\begin{equation}
 \frac{\dd E}{\dd t}=\vect F\cdot\vect v.
\end{equation}
因此功率关系保持熟悉的形式，但一般不能再写成 $\vect F=m\vect a$：平行和垂直于速度的加速度响应不同。守恒计算应优先使用四动量，而不是人为引入“随速度变化的质量”。
具体地，由 $\gamma=(1-v^2/c^2)^{-1/2}$ 求导：
\[
 \frac{\dd\gamma}{\dd t}=\gamma^3\frac{\vect v\cdot\vect a}{c^2},
 \quad \vect F=m\frac{\dd(\gamma\vect v)}{\dd t}
 =m\gamma\vect a+m\gamma^3\frac{\vect v\cdot\vect a}{c^2}\vect v.
\]
垂直于速度的分量中第二项为零；平行分量中用 $1+\gamma^2v^2/c^2=\gamma^2$，得到
\begin{equation}
 F_{\parallel}=\gamma^3ma_{\parallel},
 \qquad
 F_{\perp}=\gamma ma_{\perp}.
\end{equation}
低速极限两式都回到 $F=ma$；高速时继续增加平行速度尤其困难，而垂直力主要改变动量方向。

\begin{origexamplebox}{相对论动能反求速度}
粒子的动能等于其静止能量，即 $K=mc^2$。求粒子速度。

\begin{center}\begin{minipage}{.96\linewidth}\centering
\begin{tikzpicture}[>=Latex,font=\small,every node/.style={fill=none},scale=.88]

\draw[thick] (0,0) rectangle (1.5,.65);\node at (.75,.325){$mc^2$};
\draw[DeepBlue,thick] (1.5,0) rectangle (3,.65);\node at (2.25,.325){$K=mc^2$};
\node[above] at (1.5,.9){粒子总能量的两部分};\node[below] at (1.5,-.35){静质量 $m$，速度 $v$ 待求};

\end{tikzpicture}
\figcap{题图：题设给出动能与静止能相等，两部分组成粒子的总能量。\label{ch31:quality:q6}}
\end{minipage}\end{center}
\end{origexamplebox}
\begin{xesolution}
\begin{center}\begin{minipage}{.96\linewidth}\centering
\begin{tikzpicture}[>=Latex,font=\small,every node/.style={fill=none},scale=.88]

\draw[thick] (0,0) rectangle (1.5,.65);\node at (.75,.325){$mc^2$};
\draw[DeepBlue,thick] (1.5,0) rectangle (3,.65);\node at (2.25,.325){$K=mc^2$};
\node[above] at (1.5,.9){粒子总能量的两部分};\node[below] at (1.5,-.35){静质量 $m$，速度 $v$ 待求};

\end{tikzpicture}
\figcap{解图：区分静止能、动能和总能量，避免把 $K$ 直接等同于 $\gamma mc^2$。\label{ch31:quality:a6}}
\end{minipage}\end{center}

\ModernAnswerLabel\quad 由 $K=(\gamma-1)mc^2$ 得 $\gamma-1=1$，所以 $\gamma=2$。于是
\[
 \frac1{\sqrt{1-v^2/c^2}}=2,
 \qquad
 \boxed{v=\frac{\sqrt3}{2}c\approx0.866c}.
\]
若误用经典动能 $mv^2/2=mc^2$ 会得到 $v=\sqrt2c$，直接暴露出经典公式已超出适用范围。

图~\ref{ch31:quality:a6} 中，左块为静止能，右块为题给动能，总能量才等于 $\gamma mc^2$；因此先得到 $\gamma=2$，再由 $1-v^2/c^2=1/\gamma^2$ 求速率。
\end{xesolution}


\subsection{电场做功与磁场偏转的不同作用}
对于电荷量为 $q$、静质量为 $m$ 的粒子，静电场中的能量增量为 $q(\varphi_i-\varphi_f)$。电子从静止经加速电压大小 $U$ 获得 $K=eU$，所以
\[
 \gamma=1+\frac{eU}{mc^2},\quad
 v=c\sqrt{1-\left(1+\frac{eU}{mc^2}\right)^{-2}},\quad
 pc=\sqrt{eU(eU+2mc^2)}.
\]
磁场力与速度垂直，故不改变总能量。若速度垂直于匀强磁场，动量方向以角速度 $v/r$ 转动，$|\dd\vect p/\dd t|=pv/r$；与 $|q|vB$ 比较得
\[
 r=\frac p{|q|B},\qquad \omega=\frac{|q|B}{\gamma m}.
\]
高速时仍可用“动量除以电荷磁场”求半径，但不能再把动量直接写成 $mv$，也不能把回旋频率当作与能量无关的常数。低能极限两式恢复经典结果。

\section{系统不变质量、碰撞与反应阈值}
\begingroup\tcbset{breakable/.default=true,after={\par\smallskip\penalty0}}\begin{starread}{四维矢量、快度与动量中心系}
定义四维位置 $x^\mu=(ct,\vect x)$、四维动量 $p^\mu=(E/c,\vect p)$。采用度规符号
$(+,-,-,-)$ 时
\[
 p^\mu p_\mu=\frac{E^2}{c^2}-p^2=m^2c^2.
\]
能量守恒和三维动量守恒可统一写成四动量守恒。四速度 $U^\mu=\dd x^\mu/\dd\tau$ 的模恒为 $c$；
对于恒定静质量，四力 $K^\mu=\dd p^\mu/\dd\tau=m\,\dd U^\mu/\dd\tau$ 与四速度正交：对 $U^\mu U_\mu=c^2$ 求导，得到 $U_\mu\dd U^\mu/\dd\tau=0$。
四维语言的价值在于把参考系变换协变地组织起来，并不意味着时间和空间在测量意义上完全相同。
四加速度、本征力与加速杆的进一步讨论见高阶选读A6“常见四矢量与相对论动力学”。

沿同一直线的速度可用快度（rapidity）$\eta$ 表示：$\beta=\tanh\eta$、$\gamma=\cosh\eta$、$\gamma\beta=\sinh\eta$。
连续两次同向洛伦兹推促对应快度相加，所以自动给出相对论速度合成律。

% EMSpacetimeBoxSplit
\end{starread}\par\penalty0
\begin{starread}{系统不变质量与四电流}
对一个封闭粒子系，令总四动量为
\[
 P^\mu=\left(\frac{E}{c},\vect P\right).
\]
若 $P^\mu$ 为类时矢量，总动量为零的参考系称动量中心系（center-of-momentum frame），其相对实验室的速度为
\[
 \vect V_{\rm cm}=\frac{c^2\vect P}{E}.
\]
系统的不变质量 $M$ 由
\[
 M^2c^4=E^2-P^2c^2
\]
定义。在动量中心系中 $E_{\rm cm}=Mc^2$。碰撞或粒子反应的阈值应由总四动量守恒和这个不变量求出；把各粒子的静质量简单相加，通常会漏掉系统内部的相对动能与结合能。

% CompletenessEM:ch31-current-force
\noindent\textbf{四电流为什么按四矢量变换？}
先看只有一种载流子的连续介质。设其自身静止系中的有符号电荷密度为
$\rho_*$，在当前参考系中的速度为 $\vect u$。电荷不变而沿运动方向的体积缩短，
故 $\rho=\gamma_u\rho_*$、$\vect j=\rho\vect u$，从而
\[
 J^\mu=(c\rho,\vect j)=\rho_*U^\mu,
 \qquad U^\mu=\gamma_u(c,\vect u).
\]
右边是标量乘四速度，所以是四矢量；多种载流子的贡献相加，结论不变。
沿 $x$ 轴以速度 $v$ 换系时，直接使用与 $(ct,x)$ 相同的变换矩阵：
\[
 \begin{pmatrix}c\rho'\\j_x'\end{pmatrix}
 =\gamma_v\begin{pmatrix}1&-v/c\\-v/c&1\end{pmatrix}
 \begin{pmatrix}c\rho\\j_x\end{pmatrix}.
\]
因此
\[
 \boxed{\rho'=\gamma_v\left(\rho-\frac{vj_x}{c^2}\right),\qquad
 j_x'=\gamma_v(j_x-v\rho)},\qquad j_y'=j_y,\quad j_z'=j_z.
\]
连续性方程可统一写成 $\partial_\mu J^\mu=0$，其意义是同一个电荷守恒律在各系中成立，
并不要求各系在各自同时面上测得相同的局部电荷密度。

\begin{qualityfigure}
\begin{tikzpicture}[>=Latex,font=\small,scale=.85]
 \draw[thick](-.4,.3)rectangle(8,1.85);
 \draw[thick](-.4,-2.1)rectangle(8,-.55);
 \node[left]at(-.5,1.05){$S$};\node[left]at(-.5,-1.3){$S'$};
 \node[anchor=west]at(-.15,1.5){离子：$+\rho_0$，静止};
 \node[anchor=west]at(-.15,.65){电子：$-\rho_0$};
 \draw[->,DeepBlue,thick](7.5,.7)--(6.3,.7)node[midway,above]{$u$};
 \node[anchor=west]at(-.15,-.85){离子：$\gamma_v\rho_0$};
 \draw[->,DeepBlue](7.5,-.9)--(6.3,-.9)node[midway,above]{$v$};
 \node[anchor=west]at(-.15,-1.65){电子：$-\gamma_v\rho_0(1+vu/c^2)$};
 \draw[->,DeepBlue](7.5,-1.75)--(6.3,-1.75)node[midway,above]{$|u_e'|$};
 \draw[->](-.2,2.45)--(2.1,2.45)node[right]{$x$，$S'$相对$S$的速度$v$};
\end{tikzpicture}
\figcap{中性导线的两种载流子分别换系；框内标的是各系测得的体密度，不是相同同时面上的粒子计数。\label{complete:em:current-boost}}
\end{qualityfigure}

\noindent\textbf{算例：中性导线换系后为何带电？}
在图~\ref{complete:em:current-boost} 的 $S$ 系中，离子静止，密度 $+\rho_0$；
电子以 $-u\vect e_x$ 运动，密度 $-\rho_0$。所以总密度为零，但
$j_x=\rho_0u>0$。分别变换而非把两种物质当成同一个静止介质，得到
\[
 \begin{aligned}
 \rho_+'&=\gamma_v\rho_0,&j_+'&=-\gamma_vv\rho_0,\\
 \rho_e'&=-\gamma_v\rho_0(1+vu/c^2),&
 j_e'&=\gamma_v\rho_0(u+v).
 \end{aligned}
\]
电子速度也满足 $u_e'=-(u+v)/(1+vu/c^2)$，与 $j_e'=\rho_e'u_e'$ 一致。
两项相加，
\[
 \rho'=-\frac{\gamma_vv\rho_0u}{c^2},\qquad
 j_x'=\gamma_v\rho_0u.
\]
例如 $u=0.10c$、$v=0.60c$ 时 $\gamma_v=1.25$，
离子密度为 $1.25\rho_0$，电子密度为 $-1.325\rho_0$，总密度为
$-0.075\rho_0$，而 $j_x'=1.25j_x$。这里 $\rho_0$ 是两种载流子在 $S$ 系中的密度大小，
不是同时声称它们在各自静止系中也有相同固有密度。

\noindent\textbf{四力密度及其检验。}
沿用度规 $(+---)$，取电磁场张量分量
$F^{0i}=-E_i/c$、$F^{ij}=-\varepsilon_{ijk}B_k$，且
$J_\nu=(c\rho,-\vect j)$。缩并得到
\[
 f^\mu=F^{\mu\nu}J_\nu
 =\left(\frac{\vect j\cdot\vect E}{c},\,
 \rho\vect E+\vect j\times\vect B\right)
 =\left(\frac wc,\vect f\right).
\]
因张量与四矢量缩并仍是四矢量，同一矩阵作用于 $(w/c,f_x)$，所以
\[
 \boxed{w'=\gamma_v(w-vf_x),\qquad
 f_x'=\gamma_v\left(f_x-\frac{vw}{c^2}\right)},\qquad
 f_y'=f_y,\quad f_z'=f_z.
\]
用前面的中性导线做直接检验：只计外加纵向电场 $\vect E=E_x\vect e_x$、外加 $\vect B=0$
造成的贡献，则 $f_x=0$、$w=j_xE_x$。
沿电场换系后 $E_x'=E_x$、$\vect B'=0$，故由电荷密度直接算出
\[
 f_x'=\rho'E_x=-\frac{\gamma_vv j_xE_x}{c^2},\qquad
 w'=j_x'E_x=\gamma_vj_xE_x,
\]
与四矢量变换逐项相符。这个检验只隔离外场贡献，导线自身场另行叠加。
此处 $\vect f$ 是单位体积的力；它与单粒子三维力 $\vect F$ 的换系规律不同，不能混用。


\end{starread}\endgroup
\begin{origexamplebox}{对称完全非弹性碰撞后的静质量}
两个静质量均为 $m$ 的粒子在动量中心系中以相同速率 $v$ 相向运动，碰撞后结合成一个静止复合体，能量不向外辐射。求复合体静质量 $M$，并解释为何一般大于 $2m$。
\begin{center}\begin{minipage}{\linewidth}\centering
% 核心结论：已知两等质量粒子同速相向结合，待求末态静质量。
% 纯几何两行，蓝灰10pt，目标宽9cm；m -> <- m / M=?，无结果或辅助四动量。
\begin{tikzpicture}[font=\small,>=Latex]
 \node[draw,circle,minimum size=9mm] (a) at (-2.6,1) {$m$};
 \node[draw,circle,minimum size=9mm] (b) at (2.6,1) {$m$};
 \draw[vector] (a.east)--(-.5,1) node[midway,above] {$v$};
 \draw[vector] (b.west)--(.5,1) node[midway,above] {$v$};
 \node[left=8pt] at (a.west) {初态};
 \node[draw,minimum width=2cm,minimum height=.8cm] (f) at (0,-.65) {$M\;?$};
 \node[right=8pt] at (f.east) {结合后静止};
 \node[left=8pt] at (-3.05,-.65) {末态};
\end{tikzpicture}

\figcap{两粒子的初态及未知末态质量；动量中心系中初态总动量为零。\label{ch31:fig:collision-question}}
\end{minipage}\end{center}
\end{origexamplebox}
\begin{xesolution}
\begin{center}\begin{minipage}{\linewidth}\centering\raggedright\ModernAnswerLabel\par\smallskip\centering
% 核心结论：相向动量抵消，但初动能进入末态不变质量。
% 纯对比，上排初态m -> <- m，下排末态M；约10x4cm，10pt。
% 反应连接箭头与动量箭头分层，文字不骑线。
\begin{tikzpicture}[font=\small,>=Latex]
 \node[draw,circle,minimum size=10mm,inner sep=3pt] (a) at (-2.7,1.2) {$m$};
 \node[draw,circle,minimum size=10mm,inner sep=3pt] (b) at (2.7,1.2) {$m$};
 \draw[vector] (a.east)--(-.8,1.2) node[midway,above=5pt] {$+p$};
 \draw[vector] (b.west)--(.8,1.2) node[midway,above=5pt] {$-p$};
 \node[left=8pt] at (a.west) {初态};
 \node[draw,minimum width=3cm,minimum height=1cm,inner sep=6pt] (f) at (0,-1.2) {$M=2\gamma m$};
 \node[left=8pt] at (-3.2,-1.2) {末态};
 \draw[TextGray,->] (0,.65)--(f.north);
 \node[right=8pt] at (0,.25) {结合且无能量逸出};
 \node[below=6pt] at (f.south) {$\vect P=0$};
\end{tikzpicture}

\figcap{动量中心系中两个等质量粒子对撞并结合。总三动量为零，但系统静质量包含初始动能。\label{ch31:fig:inelastic-collision}}
\end{minipage}\end{center}
图~\ref{ch31:fig:inelastic-collision} 中两支初动量相消，末态复合体静止，因此能量守恒可直接求其不变质量。


初态总三动量为零，总能量为
\[
 E_i=2\gamma mc^2.
\]
末态复合体静止，能量为 $Mc^2$。能量守恒给出
\[
 \boxed{M=2\gamma m}.
\]
当 $v>0$ 时 $\gamma>1$，所以 $M>2m$；增加的静质量来自被困在复合体内部的初始动能（内能）。
对应的能量等式为 $(M-2m)c^2=2(\gamma-1)mc^2$：左侧是静能增加，右侧是初始总动能。
$v\to0$ 时 $M\to2m$。若碰撞中有光子或其他粒子带走能量，必须把其四动量计入，不能仍用此式。
\end{xesolution}





\begin{origexamplebox}{固定靶粒子反应的阈能}
静质量为 $m_a>0$ 的入射粒子撞击实验室中静止、静质量为 $m_b>0$ 的靶粒子。反应末态至少有两个有质量粒子，其静质量之和为 $M_f>m_a+m_b$，即反应为吸能反应。
忽略外力，只讨论由能量和动量守恒决定的运动学阈值，不考虑库仑势垒及反应概率。求最小入射动能，并说明放能反应为什么不能直接使用所得负值作为动能下限。

\begin{center}\begin{minipage}{.96\linewidth}\centering
\begin{tikzpicture}[>=Latex,font=\small,every node/.style={fill=none},scale=.88]

\fill (0,0) circle (2pt) node[below]{$a,m_a$};\draw[DeepBlue,->] (0,.3)--(1.7,.3) node[midway,above]{$\boldsymbol p_a$};
\fill (2.5,0) circle (3pt) node[below]{$b,m_b,\ v_b=0$};
\draw[->] (3.1,0)--(4.3,0);\draw[dashed] (5,0) circle (.5);\node[above] at (5,.7){末态总静质量至少 $M_f$};

\end{tikzpicture}
\figcap{题图：实验室固定靶反应的初态；入射总动量非零，末态在实验室不能全部静止。\label{ch31:quality:q8}}
\end{minipage}\end{center}
\end{origexamplebox}
\begin{xesolution}
\ModernAnswerLabel\quad 先在动量中心系中寻找生成末态所需的最少能量。在这个系中，总三动量为零，而每个有质量粒子的能量满足
\[
 E_i^*=\sqrt{m_i^2c^4+p_i^{*2}c^2}\ge m_ic^2.
\]
相加得 $E_{\rm final}^*\ge\sum_i m_ic^2=M_fc^2$。
只有所有 $p_i^*=0$ 时取等号，这就是阈值末态“没有相对动能”的由来。
注意它们只在质心系共同静止：实验室初态总动量不为零，末态必须整体运动来承接这部分动量，不能在实验室把末态全画成静止。
\begin{center}\begin{minipage}{\linewidth}\centering
% Contract: generic two-mass final state at threshold. All lab momenta forward;
% CM final particles at rest. Directions only; lengths not magnitudes.
\begin{tikzpicture}[>=Latex,font=\small,y=.75cm]
 \node[anchor=east] at (0,.8){实验室初态};
 \fill (.7,.8) circle(2.5pt) (3.7,.8) circle(2.5pt);
 \node[below=5pt] at (.7,.8){$a$};\node[below=5pt] at (3.7,.8){$b$ 静止};
 \draw[DeepBlue,thick,->] (.9,.8)--(2.5,.8) node[midway,above=4pt]{$p_a$};
 \node[anchor=east] at (0,-.9){实验室阈值末态};
 \foreach \x in {.7,3.7}{\fill (\x,-.9) circle(2.5pt);\draw[DeepBlue,thick,->] (\x,-.9)--++(.9,0);}
 \node[above=5pt] at (.7,-.9){$1$};\node[above=5pt] at (3.7,-.9){$2$};
 \node at (2.2,-1.6){同一速度 $V_c$，整体运动};
 \node[anchor=east] at (0,-2.8){质心系阈值末态};
 \fill (.7,-2.8) circle(2.5pt) (3.7,-2.8) circle(2.5pt);
 \node[below=5pt] at (.7,-2.8){$p_1^*=0$};\node[below=5pt] at (3.7,-2.8){$p_2^*=0$};
\end{tikzpicture}

\figcap{以两个末态粒子为例的阈值几何，箭头只表示动量方向。\label{ch31:fig:threshold-frames}}
\end{minipage}\end{center}
图~\ref{ch31:fig:threshold-frames} 第二行各末态粒子具有同一实验室速度 $V_c$，
但质量不同则动量不同；第三行才是所有动量均为零的质心系。
因此阈值末态总四动量的不变质量为 $M_f$。设入射粒子的实验室总能量、动量为 $E_a,p_a$。初态总能量和总动量为
\[
 E_i=E_a+m_bc^2,
 \qquad P_i=p_a.
\]
利用 $E_a^2-p_a^2c^2=m_a^2c^4$，初态不变量为
\begin{align}
 E_i^2-P_i^2c^2
 &=m_a^2c^4+m_b^2c^4+2m_bc^2E_a.
\end{align}
在阈值处令它等于 $M_f^2c^4$，得到
\[
 E_{a,\rm th}=\frac{M_f^2-m_a^2-m_b^2}{2m_b}c^2.
\]
入射动能是总能量减静止能量，所以
\begin{equation}
 \boxed{K_{\rm th}=\frac{M_f^2-(m_a+m_b)^2}{2m_b}c^2}.
\end{equation}
若末态只是原粒子的弹性散射，$M_f=m_a+m_b$，阈能为零；若产生了新的静质量，阈能不仅要提供静能差，还要保证实验室系总动量守恒，因此通常大于简单的静能差。量纲上括号内是质量，乘 $c^2$ 后为能量。
若改成具有连续末态相空间的放能反应，$M_f<m_a+m_b$，入射动能趋于零时初态就已有足够的不变能量，释放的静能差成为末态相对动能，不能再令末态相对动能为零。
因此这类反应仅从运动学考虑的下限为0。把两种情形合写为
\[
 K_{\min}=\max\left\{0,\frac{M_f^2-(m_a+m_b)^2}{2m_b}c^2\right\}.
\]
例如 $m_a=m_b=m$、$M_f=1.5m$ 时，代数表达式为 $-0.875mc^2$，只说明反应放能，不表示可以有负入射动能。
单一粒子末态属于另外的能量匹配问题，不能把上面的连续末态阈值讨论直接套用。
\end{xesolution}

\subsection{用光子系统检验不变质量}
单个真空光子的静质量为零，并不意味着多个光子组成的系统不变质量也为零。两光子的能量分别为 $E_1,E_2$，传播方向夹角为 $\theta$，总动量模平方为
\[
 P^2c^2=E_1^2+E_2^2+2E_1E_2\cos\theta.
\]
由总能量 $E=E_1+E_2$，得到
\[
 M^2c^4=E^2-P^2c^2=2E_1E_2(1-\cos\theta).
\]
同向时 $M=0$，不存在总动量为零的惯性系；反向且等能 $E_1=E_2=\varepsilon$ 时，总动量为零，$M=2\varepsilon/c^2$。这里的质量属于系统，不是给每个光子分配了静质量。计算碰撞阈值时使用总四动量，正是为了保留这类内部相对运动的能量。


\hypersetup{bookmarksdepth=2}
\endgroup

% ModernLayout
\begingroup
\hypersetup{bookmarksdepth=3}
\tcbset{breakable/.default=false,after={\par\nobreak\smallskip}}
\newcommand{\ModernAnswerLabel}{{\bfseries\color{DeepBlue}解答}}
\RenewDocumentEnvironment{xesolution}{+b}{%
  \par\nobreak\smallskip\nobreak\begingroup\small
  \clubpenalties 2 10000 10000\displaywidowpenalties 2 10000 10000\noindent
  #1\par\endgroup\smallskip
}{}
\chapter{量子现象、原子与原子核物理}

\section{热辐射与能量量子}

黑体是在给定温度下对各波长辐射具有最大吸收率的理想系统。经典能量均分与连续电磁模式结合会在高频给出发散，
普朗克（Max Planck）假定频率为 $\nu$ 的模式与物质交换能量时以 $h\nu$ 为单位，得到正确的黑体谱。单位体积、单位频率间隔的谱能量密度为
\begin{equation}
 u_\nu(T)=\frac{8\pi h\nu^3}{c^3}\frac1{e^{h\nu/(k_BT)}-1}.
\end{equation}
低频极限 $h\nu\ll k_BT$ 时，令 $x=h\nu/(k_BT)$，利用 $e^x-1\simeq x$，得到
\[
 u_\nu\simeq\frac{8\pi\nu^2k_BT}{c^3}.
\]
这就是瑞利--金斯形式（Rayleigh--Jeans law）。高频端的指数因子压过三次幂，因而不存在经典理论的紫外发散。
对频率积分得到的是单位体积内的总辐射能量 $u=\int_0^\infty u_\nu\,\dd\nu$，单位为 $\mathrm{J/m^3}$，并不是功率。
各向同性辐射中，每单位立体角分配的能量密度为 $u/(4\pi)$。穿过单位面积的能流还要乘传播速度 $c$ 和法向投影 $\cos\theta$，因此出射半球给出
\[
 j=\int_0^{2\pi}\!\int_0^{\pi/2}\frac{cu}{4\pi}\cos\theta\sin\theta\,\dd\theta\,\dd\varphi
 =\frac{cu}{4}.
\]
同一角积分还说明单位面积向立体角元 $\dd\Omega$ 发射的功率满足 $\dd j/\dd\Omega=(j/\pi)\cos\theta$，即朗伯余弦定律；$\theta$ 从表面外法线量起。
所以黑体的总辐出度（单位面积的出射功率）为
\begin{equation}
 j=\sigma T^4,
\end{equation}
面积为 $S$ 的均匀黑体表面的出射功率为 $P=Sj$。这里 $j$ 的单位为 $\mathrm{W/m^2}$；若计算与温度为 $T_0$ 的黑体环境之间的净交换功率，则为 $\sigma S(T^4-T_0^4)$。
单位波长间隔的能量密度需作变量变换，而不能只把 $\nu$ 换成 $c/\lambda$：
\[
 u_\lambda=u_\nu(c/\lambda)\left|\frac{\dd\nu}{\dd\lambda}\right|
 =\frac{8\pi hc}{\lambda^5}\frac1{e^{hc/(\lambda k_BT)}-1},\qquad M_\lambda=\frac c4u_\lambda.
\]
令 $z=hc/(\lambda k_BT)$。固定温度时，$M_\lambda$ 的极大值对应 $z^5/(e^z-1)$ 的极大值，其导数为零要求
\[
 5(e^z-1)-ze^z=0,\qquad 5(1-e^{-z})=z.
\]
非零解 $z\simeq4.9651$ 给出维恩位移律（Wien's displacement law）$\lambda_{\max}T=b=hc/(4.9651k_B)$。
按频率和按波长表示的谱密度，其峰值不是简单的倒数关系。谱峰波长随温度升高向短波移动，但不能把
$hc/\lambda_{\max}$ 直接当作所有光子的平均能量。

热平衡时，同一频率和方向上的发射率等于吸收率，这就是基尔霍夫热辐射定律（Kirchhoff's law）。黑体的吸收率为1；实际物体的发射谱还必须乘上不超过1且依赖材料、波长和方向的发射率，不能只凭温度把所有表面都当成理想黑体。


\subsection{从谱积分到温度四次方}
在普朗克谱积分中令 $x=h\nu/(k_BT)$，则 $\dd\nu=(k_BT/h)\dd x$，所有温度因子可以提出积分号：
\[
 u=\frac{8\pi k_B^4T^4}{h^3c^3}
 \int_0^\infty\frac{x^3}{e^x-1}\,\dd x
 =\frac{8\pi^5k_B^4}{15h^3c^3}T^4.
\]
这里使用无量纲积分 $\int_0^\infty x^3/(e^x-1)\,\dd x=\pi^4/15$。结合 $j=cu/4$，得
$\sigma=2\pi^5k_B^4/(15h^3c^2)$。温度四次方既来自谱中的三次幂，也来自频率积分的一个尺度因子。

若球形天体半径为 $R$、距离为 $d\gg R$，忽略吸收且均匀辐射，则总光度 $L_*=4\pi R^2\sigma T^4$，探测器处辐照度
\[
 F=\frac{L_*}{4\pi d^2}=\sigma T^4\left(\frac Rd\right)^2.
\]
表面的辐出度与远处接收的辐照度同为 $\mathrm{W/m^2}$，却对应不同面积，不能因单位相同而直接相等。实际恒星大气、选择性吸收会改变谱形，本节只使用黑体近似。

\subsection{黑体谱的温标}

\begin{origexamplebox}{恒星表面温度与辐出度}
把一颗恒星表面近似为黑体，其单位波长间隔的辐射谱峰值在 $500\,\mathrm{nm}$。用
$b=2.898\times10^{-3}\,\mathrm{m\,K}$、$\sigma=5.67\times10^{-8}\,\mathrm{W\,m^{-2}K^{-4}}$，求表面温度和单位面积总辐射功率。
\end{origexamplebox}
\begin{xesolution}
\begin{center}\begin{minipage}{\linewidth}\centering\raggedright\ModernAnswerLabel\par\smallskip\centering
% Contract: common absolute 普朗克 scale; T=4000,6000 K; peaks b/T.
\begin{tikzpicture}
\begin{axis}[width=.83\linewidth,height=6cm,axis lines=left,
 xlabel={$\lambda\;(\mu\mathrm m)$},ylabel={$M_\lambda/(10^{13}\,\mathrm{W\,m^{-3}})$},
 xmin=.15,xmax=2.5,ymin=0,ymax=11.8,xtick={.5,1,1.5,2},ytick={0,2,4,6,8,10},
 legend style={draw=none,fill=none,at={(.97,.95)},anchor=north east},
 every axis label/.style={fill=none}]
 \addplot[DeepBlue,very thick,domain=.15:2.5,samples=300]{37.4177/(x^5*(exp(14387.77/(6000*x))-1))};
 \addlegendentry{$6000\,\mathrm K$}
 \addplot[TextGray,very thick,dashed,domain=.15:2.5,samples=300]{37.4177/(x^5*(exp(14387.77/(4000*x))-1))};
 \addlegendentry{$4000\,\mathrm K$}
 \draw[guide] (axis cs:.48296203365,0)--(axis cs:.48296203365,10.0053244776);
 \draw[guide] (axis cs:.72444305047,0)--(axis cs:.72444305047,1.31757359375);
\end{axis}
\end{tikzpicture}

\figcap{黑体单色辐出度随温度的变化：升温后总面积增大，峰值向短波移动。}\label{ch32:fig:blackbody}
\end{minipage}\end{center}


维恩位移律给出
\[
 T=\frac b{\lambda_{\max}}
 =\frac{2.898\times10^{-3}}{500\times10^{-9}}
 \approx\boxed{5.80\times10^3\,\mathrm K}.
\]
斯蒂芬--玻尔兹曼定律给出
\[
 j=\sigma T^4\approx5.67\times10^{-8}(5.80\times10^3)^4
 \approx\boxed{6.4\times10^7\,\mathrm{W/m^2}}.
\]
峰值波长乘温度的单位为 $\mathrm{m\,K}$；温度升高时 $\lambda_{\max}$ 变短且总功率按 $T^4$ 增大，与图~\ref{ch32:fig:blackbody} 一致。
\end{xesolution}



\section{光子、光电效应与康普顿散射}

光的量子具有
\begin{equation}
 E=h\nu=\frac{hc}{\lambda},\qquad p=\frac E c=\frac h\lambda.
\end{equation}
光电效应中，一个电子吸收一个光子的最简能量关系为
\begin{equation}
 K_{\max}=h\nu-W,
\end{equation}
其中 $W$ 为逸出功，只有 $h\nu\ge W$ 才能发生这一单光子光电过程。截止电压 $U_s$ 指使最快光电子也不能到达收集极的反向电压大小。
设发射极电势为零，收集极电势为 $-U_s$；电子电荷为 $-e$，到达收集极需要增加的电势能为 $(-e)(-U_s)=eU_s$。临界时末动能为零，所以 $eU_s=K_{\max}$。
频率决定单个光子的能量和截止电压，光强主要改变单位时间到达的光子数；在多光子或强场条件下还会出现更复杂过程。

\begin{origexamplebox}{光电效应截止电压}
金属逸出功 $W=2.2\,\mathrm{eV}$，用波长 $400\,\mathrm{nm}$ 的单色光照射。求光电子最大动能和截止电压。
\end{origexamplebox}
\begin{xesolution}
\ModernAnswerLabel\quad 使用 $hc\simeq1240\,\mathrm{eV\,nm}$：
\[
 E_\gamma=\frac{1240}{400}\,\mathrm{eV}=3.10\,\mathrm{eV}.
\]
所以
\[
 K_{\max}=3.10-2.20=0.90\,\mathrm{eV},
 \qquad \boxed{U_s=0.90\,\mathrm V}.
\]
这里第二个数值不是把能量直接改写成电压。电子伏特的定义是
$1\,\mathrm{eV}=e\times1\,\mathrm V$，因此
\[
 U_s=\frac{K_{\max}}e
 =\frac{0.90\times1.602\times10^{-19}\,\mathrm J}
 {1.602\times10^{-19}\,\mathrm C}=0.90\,\mathrm V.
\]
\begin{center}\begin{minipage}{\linewidth}\centering
% Contract: collector has lower potential; electron travels right and slows.
% Plates are electrodes, rays enter cathode from upper left; not a circuit schematic.
\begin{tikzpicture}[>=Latex,font=\small]
 \draw[very thick] (0,-.6)--(0,1.3) (5,-.6)--(5,1.3);
 \node[below=6pt,align=center] at (0,-.6){发射极K\\$\varphi_K=0$};
 \node[below=6pt,align=center] at (5,-.6){收集极A\\$\varphi_A=-U_s$};
 \draw[DeepBlue,thick,->] (-1.2,2)--(0,.7) node[pos=.22,left=3pt]{$h\nu$};
 \draw[DeepBlue,thick,->] (.45,.55)--(2.4,.55) node[midway,above=5pt]{$e^-$ 向右运动};
 \draw[->] (3,-.05)--(4.3,-.05) node[midway,below=4pt]{$\vect E$};
 \draw[->] (2.4,1.55)--(1.1,1.55) node[midway,above=4pt]{$\vect F=-e\vect E$};
\end{tikzpicture}

\figcap{截止电压的电势约定。K接较高电势，A接较低电势；图示电场和受力方向相反。}\label{ch32:fig:photoelectric-apparatus}
\end{minipage}\end{center}
如图~\ref{ch32:fig:photoelectric-apparatus}，电子向右时电场力向左，故减速。
从K到A电势降低 $U_s$，电子电势能却增加 $eU_s$；
最快电子恰好耗尽初动能时到达临界截止，收集极电势为 $-0.90\,\mathrm V$。
若增加同频光强，饱和光电流可增大，但理想模型中的截止电压不变。
\end{xesolution}

\subsection{光电效应的直线判据}

\begin{origexamplebox}{由截止电压数据求普朗克常数}
同一种金属在频率 $f_1=6.0\times10^{14}\,\mathrm{Hz}$ 时截止电压为 $0.50\,\mathrm V$，
在 $f_2=8.0\times10^{14}\,\mathrm{Hz}$ 时为 $1.33\,\mathrm V$。估算 $h$、阈频和逸出功。
\end{origexamplebox}
\begin{xesolution}
\begin{center}\begin{minipage}{\linewidth}\centering\raggedright\ModernAnswerLabel\par\smallskip\centering
\begin{tikzpicture}
\begin{axis}[width=.72\linewidth,height=4.7cm,axis lines=left,
  xlabel={$f\;(10^{14}\,\mathrm{Hz})$},ylabel={$U_s\;(\mathrm V)$},xmin=3.5,xmax=9,ymin=0,ymax=1.8,
  xtick={4.8,6,8},xticklabels={$f_0$,$f_1$,$f_2$},ytick={.5,1.33},
  every axis label/.style={fill=none}]
  \addplot[DeepBlue,very thick,domain=4.795180723:8.8]{.415*x-1.99};
  \addplot[only marks,mark=*,DeepBlue] coordinates{(6,.5)(8,1.33)};
  \draw[guide] (axis cs:6,0)--(axis cs:6,.5);
  \draw[guide] (axis cs:8,0)--(axis cs:8,1.33);
\end{axis}
\end{tikzpicture}
\figcap{截止电压对频率的直线关系 $U_s=(h/e)f-W/e$。斜率决定普朗克常数，横截距决定阈频。}\label{ch32:fig:photoelectric-line}
\end{minipage}\end{center}
图~\ref{ch32:fig:photoelectric-line} 标出的两测量点用于本题的斜率计算，不能用连线在纵轴上的截距误认逸出功本身。


由 $eU_s=hf-W$，直线斜率为
（把两个频率对应的等式相减，逸出功 $W$ 消去）
\[
 \frac he=\frac{1.33-0.50}{(8.0-6.0)\times10^{14}}
 =4.15\times10^{-15}\,\mathrm{V\,s}.
\]
取 $e=1.602\times10^{-19}\,\mathrm C$，
\[
 \boxed{h\approx6.65\times10^{-34}\,\mathrm{J\,s}}.
\]
阈频
\[
 f_0=f_1-\frac{U_{s1}}{h/e}
 \approx\boxed{4.80\times10^{14}\,\mathrm{Hz}}.
\]
以电子伏特表示，$W/e=(h/e)f_0\approx1.99\,\mathrm V$，故
$\boxed{W\approx1.99\,\mathrm{eV}}$。两组数据的斜率单位为 $\mathrm{V\,s}$，乘电荷后为 $\mathrm{J\,s}$。
\end{xesolution}

\subsection{散射中的能量与动量守恒}
康普顿散射（Compton scattering）中，取初始自由电子静止，入射光子沿 $x$ 轴，散射光子与它成角 $\theta$。
设两光子能量分别为 $E,E'$，末态电子总能量为 $E_e$。能量和两个方向的动量守恒分别给出
\[
 E+m_ec^2=E'+E_e,\qquad p_{ex}=\frac{E-E'\cos\theta}{c},\qquad
 p_{ey}=-\frac{E'\sin\theta}{c}.
\]
因此 $p_e^2c^2=E^2+E'^2-2EE'\cos\theta$。再用 $E_e^2=m_e^2c^4+p_e^2c^2$，代入 $E_e=m_ec^2+E-E'$，消去两边共同项：
\[
 2m_ec^2(E-E')=2EE'(1-\cos\theta),\qquad
 \frac1{E'}-\frac1E=\frac{1-\cos\theta}{m_ec^2}.
\]
最后代入 $E=hc/\lambda$，得到波长改变
\begin{equation}
 \lambda'-\lambda=\frac{h}{m_ec}(1-\cos\theta).
\end{equation}
它说明高能电磁辐射交换能量和动量时具有不可忽略的量子性。

\subsection{康普顿散射几何}

\begin{origexamplebox}{九十度康普顿散射}
波长 $\lambda=0.0710\,\mathrm{nm}$ 的X射线被近似静止的自由电子散射，光子偏转 $90^\circ$。
取电子康普顿波长 $h/(m_ec)=0.002426\,\mathrm{nm}$，求散射光波长及散射前后光子能量。
\end{origexamplebox}
\begin{xesolution}
\begin{center}\begin{minipage}{\linewidth}\centering\raggedright\ModernAnswerLabel\par\smallskip\centering
% Contract: 90-degree photon scattering; recoil direction from p_e=p-p'.
\begin{tikzpicture}[font=\small,>=Latex]
 \coordinate (O) at (0,0);
 \draw[vector] (-3.6,0)--(O) node[midway,above=6pt] {入射光子 $\lambda$};
 \draw[guide] (O)--(3.1,0);
 \draw[vector] (O)--(0,2.7) node[above] {散射光子 $\lambda'$};
 \draw[vector] (O)--(2.7,-2.6108) node[right=5pt] {$\vect p_e$};
 \fill (O) circle (1.5pt);
 \node[below left=8pt] at (O) {$O$};
 \draw[DeepBlue,->] (.7,0) arc[start angle=0,end angle=90,radius=.7];
 \node at (1.2,1.05) {$\theta=90^\circ$};
 \node[align=center] at (-2,-1.1) {电子初始静止\\散射点为 $O$};
\end{tikzpicture}

\figcap{康普顿散射中的动量矢量方向，箭长不表示动量大小。散射角 $\theta$ 是入射与散射光子方向的夹角。}\label{ch32:fig:compton}
\end{minipage}\end{center}
图~\ref{ch32:fig:compton} 对应本题的 $90^\circ$ 散射：散射光子向上带走动量，反冲电子便必须具有向下的动量分量，同时保有向右的动量分量。电子箭头的方向由 $\vect p_e=\vect p-\vect p'$ 决定，而不是任意画出。


$\theta=90^\circ$ 时 $1-\cos\theta=1$，故
\[
 \boxed{\lambda'=0.0710+0.002426=0.07343\,\mathrm{nm}}.
\]
用 $hc=1.23984\,\mathrm{keV\,nm}$：
\[
 E=\frac{1.23984}{0.0710}\approx17.46\,\mathrm{keV},
 \qquad
 E'=\frac{1.23984}{0.07343}\approx16.89\,\mathrm{keV}.
\]
保留未舍入数值作差，光子损失约 $0.577\,\mathrm{keV}$，转化为电子动能。$\theta=0$ 时波长不变；$\theta=\pi$ 时改变量达到两倍康普顿波长。
图~\ref{ch32:fig:compton} 的出射光子向上，故按前面的动量守恒式有
\[
 p_{ex}=E/c>0,\qquad p_{ey}=-E'/c<0,
 \qquad \tan\phi=\frac{|p_{ey}|}{p_{ex}}=\frac{E'}E\approx0.967.
\]
电子在入射方向下方约 $44.0^\circ$ 反冲，与图中的箭头一致。
还可作独立能量检查：$p_e^2c^2=E^2+E'^2$，所以
\[
 K_e=\sqrt{(m_ec^2)^2+E^2+E'^2}-m_ec^2\approx0.577\,\mathrm{keV}.
\]
它与 $E-E'$ 一致；若把两光子的能量差全部当作电子的总能量，就会漏掉电子原有的静止能量。
\end{xesolution}



\section{物质波与量子态}

德布罗意关系（de Broglie relation）把粒子动量与波长联系起来：
\begin{equation}
 \lambda=\frac hp.
\end{equation}
电子、中子、原子乃至更大复合粒子都能在适当相干条件下产生衍射。这里“波”指量子态的概率振幅具有相位并可叠加，
“粒子”指探测时能量、电荷等以局域事件交换。量子对象不是一颗始终沿唯一轨迹运动的经典小球，也不是在三维空间中可直接测量的经典物质波；
实验装置和测量问题决定哪些概率分布可被观察。

记约化普朗克常数为 $\hbar=h/(2\pi)$。位置与动量的不确定关系为
\begin{equation}
 \Delta x\,\Delta p_x\ge\frac\hbar2.
\end{equation}
这里 $\Delta x=\sqrt{\langle x^2\rangle-\langle x\rangle^2}$，$\Delta p_x$ 同样是标准差，而不是两次测量值之差。它是量子态允许的概率分布宽度关系，不等同于仪器粗糙造成的普通测量误差。把波包限制得更窄需要更宽的动量成分。

\begin{origexamplebox}{电子经电势差加速后的波长}
电子从静止经电势差 $U$ 加速，且 $eU\ll m_ec^2$。求德布罗意波长，并说明何时需相对论修正。
\end{origexamplebox}
\begin{xesolution}
\ModernAnswerLabel\quad 非相对论能量关系 $eU=p^2/(2m_e)$ 给出
\[
 \boxed{\lambda=\frac h{\sqrt{2m_eeU}}}.
\]
若 $eU$ 与 $m_ec^2=511\,\mathrm{keV}$ 不再相比很小，应使用
\[
 (pc)^2=(m_ec^2+eU)^2-m_e^2c^4
\]
求动量。展开右侧并取正根，得到
\[
 pc=\sqrt{2m_ec^2eU+(eU)^2},\qquad
 \lambda=\frac{h}{\sqrt{2m_eeU}\sqrt{1+eU/(2m_ec^2)}}.
\]
当 $eU/(m_ec^2)\to0$ 时第二个根式趋于1，回到前面的低压公式。低压近似的适用条件来自能量尺度比较，而不是只看电压数值是否“很大”。
\end{xesolution}

\begingroup\tcbset{breakable/.default=true}\begin{starread}{薛定谔方程与自旋}
非相对论单粒子量子态由波函数 $\psi(\vect r,t)$ 表示，满足
\[
 \ii\hbar\frac{\partial\psi}{\partial t}
 =\left[-\frac{\hbar^2}{2m}\nabla^2+V(\vect r,t)\right]\psi.
\]
$|\psi|^2$ 是位置概率密度，满足 $\int |\psi|^2\,\dd^3r=1$。定态问题化为哈密顿（Hamilton）算符的本征值问题；边界条件使束缚态能量离散。
自旋是粒子的内禀角动量，不等同于带电小球绕自身轴作经典旋转。电子自旋测量沿任意轴只得到
$\pm\hbar/2$，原子中的离散单电子能级与泡利不相容原理决定基本的壳层填充：每个空间轨道至多容纳两个自旋态不同的电子。自旋磁矩则参与磁场中的能级分裂等磁相互作用，不能把磁矩本身说成壳层填充的原因。
\end{starread}\endgroup
\subsection{由边界条件认识能级离散}
以一维无限深势阱 $0<x<a$ 为模型：阱内 $V=0$，阱外概率为零，波函数在两壁满足 $\psi(0)=\psi(a)=0$。定态方程化为
\[
 \psi''+k^2\psi=0,\qquad E=\frac{\hbar^2k^2}{2m}.
\]
第一边界排除余弦项，第二边界要求 $\sin(ka)=0$。非零波函数对应 $k_n=n\pi/a$，$n=1,2,\ldots$，归一化后
\[
 \psi_n(x)=\sqrt{\frac2a}\sin\frac{n\pi x}{a},\qquad
 E_n=\frac{n^2\pi^2\hbar^2}{2ma^2}=\frac{n^2h^2}{8ma^2}.
\]
$n=0$ 只给处处为零的函数，不能归一化，所以基态能量不为零。这个模型说明离散值来自波动方程与边界条件，而不是任意把经典能量挑几个数值留下。
对电子和 $a=1.00\,\mathrm{nm}$，$E_1\approx0.376\,\mathrm{eV}$，$E_2-E_1=3E_1\approx1.13\,\mathrm{eV}$。阱宽减半时所有能级增为四倍。
无限势垒是理想边界；有限势垒会出现穿透尾部，不能继续强令阱外波函数处处为零。

衍射计算则先由加速能量求 $p$ 和 $\lambda=h/p$，再使用晶面的相位条件。加速电压、波长、散射角属于同一计算链；晶面夹角与法线夹角互为余角，使用布拉格式前应明确测角方式。

\subsection{电子衍射}

\begin{origexamplebox}{低能电子的德布罗意波长与布拉格角}
电子从静止经 $54.0\,\mathrm V$ 加速后入射晶体，晶面间距为 $0.200\,\mathrm{nm}$。
电子波长可用下式计算（$\lambda$ 以纳米计，$U$ 以伏特计）：
\[
 \lambda(\mathrm{nm})\approx\frac{1.227}{\sqrt{U(\mathrm V)}}.
\]
求电子波长和一级布拉格角。
\end{origexamplebox}
\begin{xesolution}
\begin{center}\begin{minipage}{\linewidth}\centering\raggedright\ModernAnswerLabel\par\smallskip\centering
% Contract: planes y=0,1.4; grazing angle 30deg schematic; AC,AD normal to rays.
\begin{tikzpicture}[font=\small,>=Latex]
 \coordinate (A) at (0,1.4); \coordinate (B) at (0,0);
 \coordinate (C) at (-.606218,.35); \coordinate (D) at (.606218,.35);
 \foreach \y in {0,1.4}{\draw[TextGray] (-3.2,\y)--(3.4,\y);}
 \draw[vector] (-2.598076,2.9)--(A);
 \draw[vector] (A)--(2.598076,2.9);
 \draw[vector] (-2.598076,1.5)--(B);
 \draw[vector] (B)--(2.598076,1.5);
 \draw[guide] (A)--(C); \draw[guide] (A)--(D);
 \foreach \p in {A,B,C,D}{\fill (\p) circle (1.3pt);}
 \node[above=6pt] at (A) {$A$}; \node[below=7pt] at (B) {$B$};
 \node[left=7pt] at (C) {$C$}; \node[right=7pt] at (D) {$D$};
 \draw[<->] (3.8,0)--(3.8,1.4) node[midway,right=5pt] {$d$};
 \draw (1,1.4) arc[start angle=0,end angle=30,radius=1];
 \node at (1.3,1.72) {$\theta$};
 \node[above] at (-2.4,2.9) {入射方向};
 \node[above] at (2.35,2.9) {出射方向};
\end{tikzpicture}

\figcap{晶面间距为 $d$ 的布拉格衍射几何。相邻晶面反射波程差为 $2d\sin\theta$。}\label{ch32:fig:electron-diffraction}
\end{minipage}\end{center}
图~\ref{ch32:fig:electron-diffraction} 中 $\theta$ 是电子束与晶面的掠射角，而不是与晶面法线的夹角。两束平行波分别在上层的 $A$ 和下层的 $B$ 处散射。
从 $A$ 向下层入射线、出射线作垂线，垂足分别为 $C,D$。$AC$ 是入射等相位面的截线，$AD$ 用来比较出射相位；下层波多走的路程正是 $CB+BD$。
因为 $AB=d$，两个直角三角形分别给出 $CB=d\sin\theta$、$BD=d\sin\theta$。
两段相加得到波程差 $2d\sin\theta$；同相叠加要求此值等于整数倍波长。图中角度为几何示意，本题另由数据计算。


非相对论德布罗意波长
\[
 \lambda=\frac{1.227}{\sqrt{54.0}}\,\mathrm{nm}
 \approx\boxed{0.167\,\mathrm{nm}}.
\]
一级布拉格条件 $2d\sin\theta=\lambda$ 给出
\[
 \sin\theta=\frac{1.227/\sqrt{54.0}}{0.400}\approx0.4174,
 \qquad \boxed{\theta\approx24.7^\circ}.
\]
$\sin\theta<1$，一级衍射允许。加速电压远小于 $511\,\mathrm{kV}$，非相对论近似自洽；若电压接近该尺度须修正动量关系。
\end{xesolution}



\section{原子光谱与玻尔模型}

氢原子光谱满足里德伯关系（Rydberg formula）
\begin{equation}
 \frac1\lambda=R_H\left(\frac1{n_f^2}-\frac1{n_i^2}\right),
 \qquad n_i>n_f.
\end{equation}
玻尔模型（Bohr model）先近似把质子视为静止，记 $k_e=1/(4\pi\varepsilon_0)$。库仑力提供向心力，角动量取离散值：
\[
 \frac{k_ee^2}{r_n^2}=\frac{m_ev_n^2}{r_n},\qquad m_ev_nr_n=n\hbar.
\]
把 $v_n=n\hbar/(m_er_n)$ 代入第一式，约去公共因子，得
\[
 r_n=\frac{n^2\hbar^2}{m_ek_ee^2}=a_0n^2,\qquad v_n=\frac{k_ee^2}{n\hbar}.
\]
以无限远处电子与质子静止时的能量为零，势能 $V=-k_ee^2/r_n$，动能 $K=m_ev_n^2/2=k_ee^2/(2r_n)$，故总能量
\[
 E_n=K+V=-\frac{k_ee^2}{2r_n}=-\frac{m_ek_e^2e^4}{2\hbar^2n^2}.
\]
代入常数后得到
\begin{equation}
 r_n=a_0n^2,\qquad E_n=-\frac{13.6\,\mathrm{eV}}{n^2}.
\end{equation}
跃迁光子满足 $h\nu=E_i-E_f$。玻尔模型正确解释氢样离子能级的主要结构，却不能完整描述多电子原子、谱线强度、精细结构和自旋；
这些需要量子力学。

原子尺度上电磁相互作用的无量纲强度由精细结构常数
\begin{equation}
 \alpha=\frac{e^2}{4\pi\varepsilon_0\hbar c}\approx\frac1{137}
\end{equation}
刻画。玻尔半径可写成 $a_0=\hbar/(m_ec\alpha)$，轨道速率为 $v_n=\alpha c/n$。
约化康普顿波长 $\bar\lambda_C=\hbar/(m_ec)$ 与经典电子半径 $r_e=k_ee^2/(m_ec^2)$ 因而满足
\[
 \bar\lambda_C=\alpha a_0,\qquad r_e=\alpha\bar\lambda_C=\alpha^2a_0.
\]
这里 $r_e$ 是由能量尺度定义的长度，不表示电子是半径为 $r_e$ 的刚性球。精细结构能量相对主能级通常带有 $\alpha^2$ 量级的小因子。这种无量纲参数比单独比较带单位常数更能显示修正项的相对大小。

\begin{origexamplebox}{氢原子的电离与谱线}
处于 $n=3$ 能级的氢原子跃迁到 $n=2$。（1）求光子能量；（2）从 $n=3$ 状态电离至少需要多少能量？
\end{origexamplebox}
\begin{xesolution}
\ModernAnswerLabel\quad \[
 E_3=-\frac{13.6}{9}\,\mathrm{eV},\qquad E_2=-\frac{13.6}{4}\,\mathrm{eV}.
\]
发射光子能量
\[
 \boxed{E_\gamma=E_3-E_2=13.6\left(\frac14-\frac19\right)\mathrm{eV}
 \approx1.89\,\mathrm{eV}}.
\]
从 $n=3$ 到连续谱零能级所需
\[
 \boxed{E_{\rm ion}=0-E_3=\frac{13.6}{9}\,\mathrm{eV}\approx1.51\,\mathrm{eV}}.
\]
两问对应不同终态，不能把跃迁光子能量当成该状态的电离能。
\begin{center}\begin{minipage}{\linewidth}\centering
% Contract: E2=-3.4, E3=-1.511, continuum=0; arrow directions encode processes.
\begin{tikzpicture}[>=Latex,font=\small,y=.85cm]
 \draw[->] (-.55,-3.9)--(-.55,.85) node[above]{$E/\mathrm{eV}$};
 \foreach \y/\lab in {-3.4/{n=2},-1.5111/{n=3},0/{n\to\infty}}{
  \draw[thick] (0,\y)--(6,\y);
  \node[right=5pt] at (6,\y){$\lab$};
 }
 \node[left=4pt] at (-.55,-3.4){$-3.40$};
 \node[left=4pt] at (-.55,-1.5111){$-1.51$};
 \node[left=4pt] at (-.55,0){$0$};
 \draw[DeepBlue,very thick,->] (1.1,-1.5111)--(1.1,-3.4);
 \node[right=6pt,align=left] at (1.1,-2.45){发光\\$1.89\,\mathrm{eV}$};
 \draw[DeepBlue,very thick,->] (4,-1.5111)--(4,0);
 \node[left=6pt,align=right] at (4,-.7){吸能电离\\$1.51\,\mathrm{eV}$};
\end{tikzpicture}

\figcap{从同一 $n=3$ 能级出发的两个不同过程，能级高度按能量绘制。}\label{ch32:fig:hydrogen-two-processes}
\end{minipage}\end{center}
图~\ref{ch32:fig:hydrogen-two-processes} 左箭头向下，原子的能量减少，
光子获得正能量 $E_3-E_2$；右箭头向上，原子需要吸收正能量 $0-E_3$。
零能量线表示电子与质子在无限远处、相对动能为零的电离阈值，不是又一个有限半径的束缚轨道。
\end{xesolution}


\subsection{氢样离子、约化质量与能级极限}
只有一个电子、核电荷为 $+Ze$ 的氢样离子，可以把库仑常数因子中的 $e^2$ 换成 $Ze^2$。若核质量为 $M$，电子和核均绕质心运动，相对坐标的质量为约化质量
\[
 \mu=\frac{m_eM}{m_e+M}.
\]
相应的轨道半径（两粒子的间距）和能级为
\[
 r_n=\frac{m_e}{\mu}\frac{a_0n^2}{Z},\qquad
 E_n=-\frac{\mu}{m_e}\frac{13.6Z^2}{n^2}\,\mathrm{eV}.
\]
因此氢样氦离子 $\mathrm{He}^+$ 的基态电离能约为 $54.4\,\mathrm{eV}$；中性氦有两个电子，不能直接套这条单电子公式。对同一 $Z$ 的同位素，较重的核使 $\mu$ 更接近 $m_e$，同一跃迁的光子能量略增、波长略短。
这里 $13.6\,\mathrm{eV}$ 已取近似值，精密同位素位移还需一致使用常数精度及其他修正。玻尔近似的速度尺度为 $Z\alpha c/n$；当它不再远小于 $c$ 时，非相对论轨道模型不适用。

\subsection{氢光谱的级次}

\begin{origexamplebox}{巴耳末系最长波长与系限}
取里德伯常数 $R_H=1.097\times10^7\,\mathrm{m^{-1}}$。求氢原子巴耳末系（终态 $n_f=2$）的最长波长和短波极限。
\end{origexamplebox}
\begin{xesolution}
\begin{center}\begin{minipage}{\linewidth}\centering\raggedright\ModernAnswerLabel\par\smallskip\centering
% Contract: upper levels proportional in eV; explicitly broken axis below n=2.
\begin{tikzpicture}[font=\small,>=Latex]
 \draw[axis] (-2.6,-.35)--(-2.6,.7);
 \draw[axis] (-2.6,1.25)--(-2.6,6.15) node[above] {$E\;(\mathrm{eV})$};
 \draw (-2.72,.86)--(-2.48,1.00); \draw (-2.72,1.06)--(-2.48,1.20);
 \foreach \n/\y/\en in {1/0/-13.6,2/2.2/-3.40,3/4.089/-1.51,4/4.75/-0.85}{
  \draw[DeepBlue,thick] (-.1,\y)--(4.8,\y);
  \node[left=6pt] at (-.1,\y) {$n=\n$};
  \node[left=5pt] at (-2.6,\y) {$\en$};
 }
 \draw[TextGray,dashed] (-.1,5.6)--(4.8,5.6);
 \node[left=6pt] at (-.1,5.6) {$n=\infty$};
 \node[left=5pt] at (-2.6,5.6) {$0$};
 \draw[vector] (1.4,5.6)--(1.4,2.2);
 \draw[vector] (3.8,4.089)--(3.8,2.2);
 \node[below=8pt,align=center] at (1.4,2.2) {巴耳末系限};
 \node[below=8pt] at (3.8,2.2) {$3\to2$};
 \node at (2.4,.9) {纵轴断开};
\end{tikzpicture}

\figcap{氢原子能级与巴耳末系。纵轴在基态与第一激发态之间断开；上部各能级按能量比例绘制，越靠近电离极限，间隔越小。}\label{ch32:fig:hydrogen-levels}
\end{minipage}\end{center}


最长波长对应最小能级差，即 $n_i=3$：
\[
 \frac1{\lambda_{\max}}=R_H\left(\frac14-\frac19\right)
 =R_H\frac5{36},
\]
所以 $\boxed{\lambda_{\max}\approx656\,\mathrm{nm}}$。短波极限对应 $n_i\to\infty$：
\[
 \frac1{\lambda_{\min}}=\frac{R_H}{4},
 \qquad \boxed{\lambda_{\min}\approx364.6\,\mathrm{nm}}.
\]
图~\ref{ch32:fig:hydrogen-levels} 显示高能级向电离极限聚集，故跃迁光子能量趋于有限上限而不是无限增大。
\end{xesolution}



\section{原子核、结合能与核反应}

原子的典型尺度约为 $10^{-10}\,\mathrm m$；原子核半径通常在 $10^{-15}$--$10^{-14}\,\mathrm m$ 范围，
经验关系为
\begin{equation}
 R\simeq r_0A^{1/3},\qquad r_0\approx1.2\,\mathrm{fm}.
\end{equation}
核质量小于自由核子质量之和，质量亏损对应结合能
\begin{equation}
 B=\Delta mc^2.
\end{equation}
平均每核子结合能反映核的相对稳定性。轻核聚变和重核裂变都可能向更高平均结合能方向释放能量。


\subsection{质量表、反应Q值与两体反冲}
记核电荷数为 $Z$、中子数为 $N$、质量数为 $A=Z+N$。核反应前后先核对总电荷与总核子数，再计算质量差。若用裸核质量 $M_N$，结合能为
\[
 B=[Zm_p+(A-Z)m_n-M_N]c^2.
\]
若给的是中性原子质量 $M_{\rm at}$，并忽略电子结合能差，宜改用氢原子质量 $m_H$：
\[
 B\simeq[Zm_H+(A-Z)m_n-M_{\rm at}]c^2.
\]
两式中的质量对象必须一致；把原子质量和裸质子质量直接混在一起会多减去 $Z$ 个电子静质量。
一般核反应的 $Q=(\sum M_i-\sum M_f)c^2$，并满足 $K_f-K_i=Q$。$Q<0$ 时固定靶入射阈能通常大于 $|Q|$，因为末态还需携带总动量；精确计算用第31章的系统不变量。

当静止母核发生两体衰变，产物动量相反且大小相等。若两产物均非相对论，有 $K_1=p^2/(2m_1)$、$K_2=p^2/(2m_2)$，从而
\[
 K_1=Q\frac{m_2}{m_1+m_2},\qquad
 K_2=Q\frac{m_1}{m_1+m_2}.
\]
较轻产物带走较多动能。若一产物为光子，须用 $p_\gamma=E_\gamma/c$，不能把光子代入 $p^2/(2m)$。母体静质量 $M_i$、子体质量 $M_f$ 的两体辐射跃迁精确给出
\[
 E_\gamma=\frac{M_i^2-M_f^2}{2M_i}c^2<Q=(M_i-M_f)c^2,
\]
差值用于子体反冲。忽略反冲是建立在 $Q\ll M_ic^2$ 上的近似。

\subsection{结合能曲线与裂变能}

\begin{origexamplebox}{由平均结合能估算裂变能}
粗略取 $^{235}\mathrm U$ 的平均每核子结合能为 $7.6\,\mathrm{MeV}$，裂变产物平均为
$8.5\,\mathrm{MeV}$。忽略裂变中子与产物间的细小差别，估算一次裂变释放的能量，并换算为焦耳。
\end{origexamplebox}
\begin{xesolution}
\begin{center}\begin{minipage}{\linewidth}\centering\raggedright\ModernAnswerLabel\par\smallskip\centering
% Contract: schematic binding trend, arrows separate from approximate data.
\begin{tikzpicture}
\begin{axis}[width=.83\linewidth,height=5.8cm,axis lines=left,
 xlabel={$A$},ylabel={$B/A\;(\mathrm{MeV})$},xmin=0,xmax=250,ymin=0,ymax=11,
 xtick={56,120,235},ytick={0,4,8},every axis label/.style={fill=none}]
 \addplot[DeepBlue,very thick] coordinates{(1,0)(2,1.1)(4,7.1)(12,7.7)(20,8.0)(40,8.5)(56,8.8)(90,8.7)(120,8.5)(160,8.2)(200,7.9)(235,7.6)};
 \node[above=6pt] at (axis cs:56,8.8) {铁、镍附近};
 \draw[vector] (axis cs:220,5.7)--(axis cs:130,6.7);
 \node[below=6pt] at (axis cs:175,5.7) {重核裂变};
 \draw[vector] (axis cs:15,4.0)--(axis cs:50,6.1);
 \node[right=5pt] at (axis cs:42,4.5) {轻核聚变};
\end{axis}
\end{tikzpicture}

\figcap{平均每核子结合能的典型趋势。轻核聚变和重核裂变都可向更高结合能方向释放能量。}\label{ch32:fig:binding-curve}
\end{minipage}\end{center}
图~\ref{ch32:fig:binding-curve} 中箭头只指示反应趋向，不表示精确核素之间的能量差；折线为典型趋势，不能代替核质量表。本题用题给的反应前后平均结合能之差，乘参与反应的核子数估算释放能量。


总结合能增加约
\[
 Q\approx235(8.5-7.6)\,\mathrm{MeV}
 =\boxed{2.12\times10^2\,\mathrm{MeV}}.
\]
用 $1\,\mathrm{MeV}=1.602\times10^{-13}\,\mathrm J$，
\[
 \boxed{Q\approx3.39\times10^{-11}\,\mathrm J}.
\]
这是数量级估算，精确值应使用各核素静质量和实际产物分支。符号为正，因为末态总结合能更大、总静质量更小。
\end{xesolution}

\subsection{核反应的Q值}

\begin{origexamplebox}{氘氚聚变的Q值}
反应 $^2\mathrm H+{}^3\mathrm H\to{}^4\mathrm{He}+n$ 中，氘、氚、氦的中性原子质量以及中子的质量依次为
$2.014102\,\mathrm u$、$3.016049\,\mathrm u$、$4.002603\,\mathrm u$、$1.008665\,\mathrm u$。
取 $1\,\mathrm u\,c^2=931.5\,\mathrm{MeV}$，求Q值。
\end{origexamplebox}
\begin{xesolution}
\begin{center}\begin{minipage}{\linewidth}\centering\raggedright\ModernAnswerLabel\par\smallskip\centering
% Contract: reaction bookkeeping, not spatial particle trajectories.
\begin{tikzpicture}[font=\small,>=Latex]
 \node[draw,circle,minimum size=11mm] (d) at (-4,0) {$^2\mathrm H$};
 \node at (-2.95,0) {$+$};
 \node[draw,circle,minimum size=11mm] (t) at (-1.9,0) {$^3\mathrm H$};
 \draw[vector] (-.95,0)--(.95,0);
 \node[draw,circle,minimum size=13mm] (he) at (2,0) {$^4\mathrm{He}$};
 \node at (3.15,0) {$+$};
 \node[draw,circle,minimum size=8mm] (n) at (4.15,0) {$n$};
 \node at (-2.95,1.1) {反应物}; \node at (3.05,1.1) {产物};
 \node[align=center] at (0,-1.3) {$Q=(m_D+m_T-m_{\rm He}-m_n)c^2>0$\\总动能增加，静质量减少};
\end{tikzpicture}

\figcap{氘氚聚变的质量--能量收支。核子数和电荷数在反应前后分别守恒。}\label{ch32:fig:fusion-q}
\end{minipage}\end{center}
图~\ref{ch32:fig:fusion-q} 把反应物和产物分列两侧，下面计算先用质量差判断 $Q$ 值符号，再核验两个守恒数。


两侧电子数相同，电子静质量在相减时抵消；电子结合能仍有很小差异，
但在本题约 $17.6\,\mathrm{MeV}$ 的精度下可以忽略。因此用所给中性原子质量计算。质量差
\[
 \Delta m=(2.014102+3.016049-4.002603-1.008665)\,\mathrm u
 =0.018883\,\mathrm u.
\]
故
\[
 \boxed{Q=\Delta mc^2\approx0.018883\times931.5
 \approx17.6\,\mathrm{MeV}}.
\]
$Q>0$，静质量减少量转化为产物动能。核子数 $2+3=4+1$、电荷数 $1+1=2+0$ 均守恒，是写反应式时的第一层检查。
\end{xesolution}



\section{放射性衰变与粒子守恒}

放射性核数满足
\begin{equation}
 N(t)=N_0e^{-\lambda t},\qquad
 T_{1/2}=\frac{\ln2}{\lambda},\qquad
 \mathcal A(t)=\lambda N(t).
\end{equation}
活度记为 $\mathcal A$，单位为贝可勒尔（$\mathrm{Bq}=\mathrm{s^{-1}}$），以免与质量数 $A$ 混淆。指数规律来自 $\dd N/\dd t=-\lambda N$：分离变量并积分得 $\ln[N(t)/N_0]=-\lambda t$；令 $N/N_0=1/2$ 即得半衰期。半衰期是统计规律，不表示某个原子核在半衰期时刻必然衰变。


\begin{origexamplebox}{经过三个半衰期的剩余量}
某样品半衰期为 $T_{1/2}$。经过 $3T_{1/2}$ 后，未衰变核数和活度分别为初值的多少？
\end{origexamplebox}
\begin{xesolution}
\ModernAnswerLabel\quad 每经过一个半衰期剩余一半，故
\[
 \boxed{\frac N{N_0}=2^{-3}=\frac18}.
\]
同一核素的 $\lambda$ 不变，活度 $\mathcal A=\lambda N$，所以
\[
 \boxed{\frac{\mathcal A}{\mathcal A_0}=\frac18}.
\]
若衰变产物也有放射性，总测得活度还要加上子核贡献，不能只用单指数。
\end{xesolution}


\subsection{衰变方式与计数模型}
常见衰变方式可用守恒量区分：
\begin{align*}
 {}^A_ZX&\to{}^{A-4}_{Z-2}Y+{}^4_2\mathrm{He} &&(\alpha),\\
 {}^A_ZX&\to{}^A_{Z+1}Y+e^-+\bar\nu_e &&(\beta^-),\\
 {}^A_ZX&\to{}^A_{Z-1}Y+e^++\nu_e &&(\beta^+),\\
 {}^A_ZX^*&\to{}^A_ZX+\gamma &&(\gamma).
\end{align*}
$\beta$ 衰变中的电子或正电子是在相互作用中产生的，不是从核内预先存放的经典电子中取出。三体末态中能量在带电轻子、中微子与反冲核之间分配，所以电子能谱连续；这不表示每次衰变的总能量不守恒。$\gamma$ 跃迁不改变核的 $A,Z$。
若用中性原子质量并忽略电子结合能差，$\beta^-$ 衰变约有 $Q=[M_X-M_Y]c^2$，而 $\beta^+$ 衰变约有 $Q=[M_X-M_Y-2m_e]c^2$；后者额外需要两个电子静能的质量差门槛。不能把两式混用。

探测器的净计数率还与效率 $\varepsilon$ 有关：$R(t)=\varepsilon\mathcal A(t)$。若效率稳定、只计单一核素且背景已扣除，间隔 $\Delta t$ 的计数率比满足
\[
 \lambda=\frac1{\Delta t}\ln\frac{R(t)}{R(t+\Delta t)}.
\]
从 $t=0$ 到 $\Delta t$ 的期望净计数应积分：$C=\varepsilon N_0(1-e^{-\lambda\Delta t})$，只有 $\lambda\Delta t\ll1$ 时才可近似为初始计数率乘时间。任意区间 $[t,t+\Delta t]$ 应将上式的 $N_0$ 换为 $N(t)$。若有恒定背景，应先减去背景再取对数；若子核也能被计数，应先建立衰变链模型。

\subsection{两级放射性衰变链}

\begin{origexamplebox}{母核--子核衰变链的峰值}
母核1以衰变常数 $\lambda_1$ 变成放射性子核2，子核2再以 $\lambda_2=2\lambda_1$ 衰变为稳定核。
初始有 $N_{10}$ 个母核且无子核。求 $N_2(t)$、子核数达到最大值的时刻和最大值。
\end{origexamplebox}
\begin{xesolution}
\begin{center}\begin{minipage}{\linewidth}\centering\raggedright\ModernAnswerLabel\par\smallskip\centering
% Contract: lambda2=2lambda1, N2 peak (ln2,1/4), legend outside curves.
\begin{tikzpicture}
\begin{axis}[width=.8\linewidth,height=5.4cm,axis lines=left,
 xlabel={$\lambda_1t$},ylabel={$N/N_{10}$},xmin=0,xmax=5,ymin=0,ymax=1.08,
 xtick={.693147,2,4},xticklabels={$\ln2$,$2$,$4$},ytick={.25,.5,1},
 legend style={draw=none,fill=none,at={(.96,.94)},anchor=north east},
 every axis label/.style={fill=none}]
 \addplot[DeepBlue,very thick,domain=0:5,samples=200]{exp(-x)}; \addlegendentry{母核 $N_1$}
 \addplot[TextGray,very thick,dashed,domain=0:5,samples=200]{exp(-x)-exp(-2*x)}; \addlegendentry{子核 $N_2$}
 \draw[guide] (axis cs:.693147,0)--(axis cs:.693147,.25)--(axis cs:0,.25);
 \addplot[only marks,mark=*,DeepBlue,forget plot] coordinates {(.693147,.25)};
\end{axis}
\end{tikzpicture}

\figcap{$\lambda_2=2\lambda_1$ 时母核与子核数目的变化。子核先积累，达到峰值后再衰减。}\label{ch32:fig:decay-chain}
\end{minipage}\end{center}


母核数 $N_1=N_{10}e^{-\lambda_1t}$。子核满足
\[
 \dot N_2=\lambda_1N_1-\lambda_2N_2,
 \qquad N_2(0)=0.
\]
将方程写成 $\dot N_2+\lambda_2N_2=\lambda_1N_{10}e^{-\lambda_1t}$。两侧乘积分因子 $e^{\lambda_2t}$，左侧恰好成为乘积的导数：
\[
 \frac{\dd}{\dd t}(N_2e^{\lambda_2t})=\lambda_1N_{10}e^{(\lambda_2-\lambda_1)t}.
\]
从0积分到 $t$，利用 $N_2(0)=0$，当 $\lambda_2\ne\lambda_1$ 时有
\[
 N_2e^{\lambda_2t}=\lambda_1N_{10}\frac{e^{(\lambda_2-\lambda_1)t}-1}{\lambda_2-\lambda_1}.
\]
两侧乘 $e^{-\lambda_2t}$，得到一般的巴特曼形式
\[
 N_2=N_{10}\frac{\lambda_1}{\lambda_2-\lambda_1}
 \left(e^{-\lambda_1t}-e^{-\lambda_2t}\right).
\]
代入 $\lambda_2=2\lambda_1$：
\[
 \boxed{N_2=N_{10}\left(e^{-\lambda_1t}-e^{-2\lambda_1t}\right)}.
\]
求导得 $\dot N_2=\lambda_1N_{10}e^{-\lambda_1t}[-1+2e^{-\lambda_1t}]$。
前面的因子始终为正，方括号由正变负，故唯一零点就是最大值点。令 $e^{-\lambda_1t}=1/2$，所以
\[
 \boxed{t_{\max}=\frac{\ln2}{\lambda_1}},
 \qquad \boxed{N_{2,\max}=\frac{N_{10}}4}.
\]
$t=0$ 与 $t\to\infty$ 时 $N_2$ 都为零，峰值位于中间，和图~\ref{ch32:fig:decay-chain} 一致。
\end{xesolution}

\subsection{衰变链的相等常数极限}
一般两级模型中令 $\lambda_2\to\lambda_1=\lambda$，巴特曼式表面上为 $0/0$，应先取极限而非宣布发散。用导数或直接积分得
\[
 N_2(t)=\lambda N_{10}t e^{-\lambda t},\qquad
 t_{\max}=\frac1\lambda,\qquad N_{2,\max}=\frac{N_{10}}e.
\]
无论两常数是否相等，子核峰值时都满足生产率等于损失率：$\lambda_1N_1=\lambda_2N_2$。相等常数极限、初始导数 $\dot N_2(0)=\lambda_1N_{10}$、长时间趋零，是检查衰变链公式的三个独立办法。若母核只有比例 $b$ 的分支通向所计子核，源项应改为 $b\lambda_1N_1$；若初始已有子核，还需加 $N_{20}e^{-\lambda_2t}$。

\subsection{中微子与粒子物理概览}

连续的 $\beta$ 衰变能谱促使泡利（Wolfgang Pauli）提出中微子，以保持逐次事件中的能量、动量和角动量守恒。中微子只参与弱相互作用和引力，
探测率极低。2002年诺贝尔物理学奖的一半授予戴维斯和小柴昌俊，表彰宇宙中微子探测方面的开创性工作；另一半授予贾科尼，表彰宇宙X射线源研究。2015年奖项表彰中微子振荡的发现，
这一现象表明中微子具有非零质量。两次奖项对应不同成果，不能合并成一次“发现中微子质量”的奖项。

标准模型把物质基本粒子分为夸克和轻子，以规范玻色子传递强、电磁和弱相互作用。质子、中子由夸克和胶子组成；
夸克具有分数电荷但在通常条件下不能孤立观测。标准模型不包含完整的量子引力理论，也不能单独解释暗物质等所有现象。




\hypersetup{bookmarksdepth=2}
\endgroup


\appendix
\setcounter{chapter}{4}
% AppendixQuality20261023
\begingroup\tcbset{breakable/.default=false}
\chapter{电磁学补充专题}
% Editable, problem-specific diagrams for the electromagnetic appendix.
\tikzset{resistor/.style={thick,decorate,decoration={zigzag,segment length=5pt,amplitude=2pt}}}
\newcommand{\QAENamed}[3]{\begin{qualityfigure}\begingroup
  \renewcommand{\figcap}[1]{}\renewcommand{\label}[1]{}#1\endgroup
  \figcap{#2}\label{#3}\end{qualityfigure}}
\newcommand{\QAESources}{\begin{tikzpicture}[>=Latex,scale=.74,font=\small]
 \begin{scope}[xshift=-6cm]
  \draw[very thick] (0,-1.3)--(0,1.3);\node[above] at (0,1.3) {$\lambda$};
  \draw[guide] (0,0)--(1.7,0);\fill (1.7,0) circle(2pt) node[right] {$P$};
  \node[below] at (.85,0) {$a$};\node at (.8,-1.9) {无限直线};
 \end{scope}
 \begin{scope}[xshift=-2.6cm]
  \draw[very thick] (0,-.8)--(0,1.4);\coordinate(P) at(1.7,0);\coordinate(O) at(0,0);
  \draw[guide] (P)--(0,1.4) (P)--(0,-.8) (P)--(O);
  \node[left] at (0,1.4) {$T_1$};\node[left] at (0,-.8) {$T_2$};
  \node at(1.1,.3){$\theta_1$};\node at(1.05,-.25){$\theta_2$};
  \fill(P) circle(2pt) node[right] {$P$};\node[below] at(.8,0) {$a$};
  \node at (.7,-1.9) {有限直线：两端张角};
 \end{scope}
 \begin{scope}[xshift=1.4cm]
  \fill[PaleBlue] (-.8,-.4)--(.7,-.8)--(1.6,.2)--(.1,.6)--cycle;
  \draw[guide] (.4,-.1)--(.4,1.4);\node[left]at(.4,.75){$a$};\fill(.4,1.4) circle(2pt) node[right] {$P$};
  \node at(.4,-1.2) {$\sigma$};\node at(.4,-1.9) {无限平面／有限宽条带};
 \end{scope}
 \begin{scope}[xshift=6.0cm]
  \draw[fill=PaleBlue] (0,0) ellipse(.5 and 1.15);\draw[axis](0,0)--(2.1,0);
  \draw[guide](0,1.15)--(1.65,0);\fill(1.65,0) circle(2pt) node[above] {$P$};
  \node[below] at(.9,0) {$a$};\node[left] at(0,.8) {$R$};\node at(.7,-1.9) {圆盘轴线};
 \end{scope}
\end{tikzpicture}}
\newcommand{\QAEStrip}{\begin{tikzpicture}[>=Latex,scale=.85,font=\small]
 \begin{scope}[xshift=-3.8cm]
  \draw[axis](-2,0)--(2.1,0) node[right] {$x$};\draw[very thick](-1.3,0)--(1.8,0);
  \fill[PaleBlue] (.8,-.13) rectangle(1,.13);\node[below]at(.9,-.15){$\dd x$};
  \coordinate(P)at(0,1.8);\draw[guide](-1.3,0)--(P)--(1.8,0) (0,0)--(P);
  \fill(P) circle(2pt) node[above] {$P$};\node[left]at(0,.55){$a$};
  \node at(.48,1.06){$\theta_1$};\node at(-.4,1.1){$\theta_2$};
  \draw[vector](P)--(-.4,2.55) node[left]{$\dd\vect E$};
  \node[below,align=center]at(0,-.65){条带沿纸外方向无限延伸\\每条线元 $\dd\lambda=\sigma\dd x$};
 \end{scope}
 \begin{scope}[xshift=3.7cm]
  \fill[PaleBlue](-1.6,-.15)rectangle(1.6,.15);\draw[guide](-.9,-1)--(.9,-1)--(.9,1)--(-.9,1)--cycle;
  \draw[vector](0,.15)--(0,1.65)node[right]{$\vect E_+$};
  \draw[vector](0,-.15)--(0,-1.65)node[right]{$\vect E_-$};
  \node at(0,-2.15){跨平面的柱形高斯面：两个端面通量};
 \end{scope}
\end{tikzpicture}}
\newcommand{\QAEZeroField}{\begin{tikzpicture}[>=Latex,scale=.8,font=\small]
 \draw[axis](-3.8,0)--(5.3,0)node[right]{$x$};
 \fill[DeepBlue](-2,0)circle(3pt)node[below]{$+Q$};\fill[Gold](0,0)circle(3pt)node[below]{$-q$};
 \fill(3,0)circle(1.8pt)node[below right]{$x_0$};
 \draw[vector](3.05,0)--(4.5,0)node[above]{$\delta E_x>0$};\draw[vector](2.95,0)--(1.6,0);
 \draw[vector](3,1.4)--(3,.25);\draw[vector](3,-1.4)--(3,-.25);
 \draw[guide](-2,1.9)--(3,1.9) (0,1.35)--(3,1.35);
 \node[above]at(.4,1.9){$x_0+a$};\node[above]at(1.4,1.35){$x_0-a$};
 \node[below,align=center]at(0,-2){轴向扰动离开零场点，横向扰动趋向轴线\\箭头表示零场点邻域的场，非两电荷之间的全局场线};
\end{tikzpicture}}
\newcommand{\QAECavityVectors}{\begin{tikzpicture}[>=Latex,scale=.9]
 \draw[fill=PaleBlue](0,0)circle(1.8);\draw[fill=white](.65,.2)circle(.8);
 \coordinate(O)at(0,0);\coordinate(C)at(.65,.2);\coordinate(P)at(.9,.8);
 \draw[vector](O)--(C)node[midway,below]{$\vect d$};\draw[vector](C)--(P)node[right]{$\vect r'$};
 \draw[vector](O)--(P)node[midway,left]{$\vect r$};
 \foreach\p in {O,C,P}{\fill(\p)circle(1.5pt)node[below right]{$\p$};}
 \node[align=left,anchor=west]at(2.5,.4){完整大球：$+\rho$\\补偿小球：$-\rho$\\$\vect r-\vect r'=\vect d$};
\end{tikzpicture}}
\newcommand{\QAEDipoleComponents}{\begin{tikzpicture}[>=Latex,scale=.9]
 \draw[axis](-1,0)--(3.8,0)node[right]{$\vect p$};\coordinate(O)at(0,0);\coordinate(P)at(2.2,1.5);
 \draw[guide](O)--(P)node[midway,above left]{$r$};\node[below]at(O){$O$};\fill(P)circle(2pt)node[above]{$P$};
 \draw[vector](P)--++(1.1,.75)node[right]{$E_r$};\draw[vector](P)--++(-.6,.88)node[above]{$E_\theta$};
 \draw (0.65,0) arc(0:34.3:.65)node[midway,right]{$\theta$};
 \node[below,align=center]at(1,-1){极角由 $\vect p$ 向场点量起\\$E_r,E_\theta$ 的正基矢按极坐标定义};
\end{tikzpicture}}
\newcommand{\QAEPotentialPlot}{\begin{tikzpicture}[>=Latex,scale=.75,font=\small]
 \begin{scope}[xshift=-4cm]
  \draw[axis](0,-1.4)--(0,1.5)node[above]{$\varphi$};\draw[axis](0,0)--(3.2,0)node[right]{$r/r_0$};
  \draw[DeepBlue,thick,domain=.3:3,samples=60]plot(\x,{-ln(\x)});\draw[guide](1,0)--(1,-.2)node[below]{$1$};
  \node[below]at(1.6,-1.9){直线电荷：$\varphi\propto-\ln(r/r_0)$};
 \end{scope}
 \begin{scope}[xshift=3cm]
  \draw[axis](-2,0)--(2.3,0)node[right]{$z$};\draw[axis](0,-1.7)--(0,1.4)node[above]{$\varphi$};
  \draw[DeepBlue,thick](-1.6,-1.3)--(0,0)--(1.6,-1.3);
  \node[below]at(0,-1.9){面电荷：$\varphi\propto-|z|$};
 \end{scope}
\end{tikzpicture}}
\newcommand{\QAECapacitors}[1]{\begin{tikzpicture}[>=Latex,scale=.76,font=\small]
 \begin{scope}[xshift=-5cm]
  \draw[very thick](-1.1,.65)--(1.1,.65) (-1.1,-.65)--(1.1,-.65);
  \node[above]at(0,.65){$+Q,\ S=l^2$};\node[below]at(0,-.65){$-Q$};\draw[<->](1.4,-.65)--(1.4,.65)node[midway,right]{$d$};
  \ifnum#1=1 \draw[vector](0,.45)--(0,-.45)node[right]{$E$};\fi
 \end{scope}
 \begin{scope}
  \draw(0,0)circle(.7);\draw(0,0)circle(1.3);\draw[guide](0,0)--(1.3,0);
  \node[above]at(.45,0){$R_1$};\node[below]at(1,0){$R_2$};\node[left]at(-.7,0){$+Q$};\node[above]at(0,1.3){$-Q$};
  \ifnum#1=1\draw[vector](.78,-.6)--(1.03,-.8)node[right]{$E(r)$};\fi
 \end{scope}
 \begin{scope}[xshift=5cm]
  \draw(0,0)circle(.55);\draw(0,0)circle(1.25);\draw[guide](0,0)--(1.25,0);
  \node[above]at(.3,0){$a$};\node[below]at(.95,0){$b$};\node[above,align=center]at(0,1.3){同轴柱横截面\\轴向长度 $l$};
  \ifnum#1=1\draw[vector](0,-.65)--(0,-1.15)node[right]{$E(r)$};\fi
 \end{scope}
\end{tikzpicture}}
\newcommand{\QAESplitQuestion}{\begin{tikzpicture}[>=Latex,scale=.8]
 \foreach\h in {-3.5,3.5}{\draw[very thick](\h-1.6,1)--(\h+1.6,1) (\h-1.6,-1)--(\h+1.6,-1);}
 \fill[PaleBlue](-5.1,-.95)rectangle(-3.5,.95);\draw[guide](-3.5,-1)--(-3.5,1);
 \node at(-4.3,0){$\varepsilon_r$};\node at(-2.7,0){真空};\node at(-3.5,-1.5){各占面积 $S/2$};
 \fill[PaleBlue](1.9,0)rectangle(5.1,.95);\node at(3.5,.5){$\varepsilon_r$};\node at(3.5,-.5){真空};
 \draw[<->](5.4,0)--(5.4,1)node[midway,right]{$d/2$};\node at(3.5,-1.5){各层面积 $S$};
\end{tikzpicture}}
\newcommand{\QAEInsertQuestion}{\begin{tikzpicture}[>=Latex,scale=.82]
 \draw[very thick](-3.3,.7)--(3.3,.7) (-3.3,-.7)--(3.3,-.7);
 \fill[PaleBlue](-3.3,-.65)rectangle(-.3,.65);\node at(-1.8,0){$\varepsilon_r$};\node at(1.4,0){真空};
 \draw[<->](-3.3,-1.1)--(-.3,-1.1)node[midway,below]{$x$};\draw[<->](-.3,-1.1)--(3.3,-1.1)node[midway,below]{$a-x$};
 \draw[<->](3.6,-.7)--(3.6,.7)node[midway,right]{$d$};\node[above]at(0,1.1){垂直纸面宽度 $b$};
\end{tikzpicture}}
\newcommand{\QAERCQuestion}{\begin{tikzpicture}[>=Latex,scale=.8]
 \draw(-3,0)--(-2.3,0);\draw[very thick](-2.3,-.5)--(-2.3,.5);\draw(-1.95,-.3)--(-1.95,.3);
 \draw(-1.95,0)--(-.7,0);\draw[resistor](-.7,0)--(1.1,0)node[midway,above]{$R$};
 \draw(1.1,0)--(2,0);\draw[very thick](2,-.5)--(2,.5) (2.4,-.5)--(2.4,.5);
 \draw(2.4,0)--(3,0)--(3,-1.4)--(-3,-1.4)--(-3,0);
 \node[below]at(-2.15,-.55){$\mathcal E$};\node[above]at(2.2,.6){$C,\ Q(0)=0$};
\end{tikzpicture}}
\newcommand{\QAEEarthQuestion}{\begin{tikzpicture}[>=Latex,scale=.8]
 \fill[SoftGray](0,0)circle(1.4);\draw(0,0)circle(1.4);\draw[DeepBlue](0,0)circle(2);
 \draw[->](0,0)--(1.4,0)node[midway,below]{$R_\oplus$};\draw[<->](1.4,0)--(2,0)node[midway,above]{$h$};
 \node at(0,.55){地球 $-Q$};\node[above]at(0,2){等效外壳 $+Q$};
 \node[anchor=west,align=left]at(2.8,.3){薄层 $h\ll R_\oplus$\\空气电阻率 $\rho$ 均匀\\$h$ 在图中放大};
\end{tikzpicture}}
\newcommand{\QAEVacuumQuestion}{\begin{tikzpicture}[>=Latex,scale=.85]
 \draw[very thick](-3,-1)--(-3,1) (3,-1)--(3,1);
 \node[above]at(-3,1){阴极 $0$};\node[above]at(3,1){阳极 $U>0$};
 \draw[axis](-3,-1.5)--(3.8,-1.5)node[right]{$x$};\node[below]at(-3,-1.5){$0$};\node[below]at(3,-1.5){$L$};
 \draw[vector](-1,0)--(1.3,0)node[midway,above]{$e^-,\ v$};\node at(0,-2.2){电极面积 $S$；忽略初始速率与边缘效应};
\end{tikzpicture}}
\newcommand{\QAEIonQuestion}{\begin{tikzpicture}[>=Latex,scale=.85]
 \draw[very thick](-3,-1)--(-3,1) (3,-1)--(3,1);
 \node[above]at(-3,1){阳极：读出};\node[above]at(3,1){阴极：$-U_0$};
 \draw[guide](-1,-1)--(-1,1);\node[above]at(-1,1){$N$ 对};
 \draw[vector](-1,.35)--(-2.4,.35)node[midway,above]{$v^-$};\draw[vector](-1,-.35)--(1,-.35)node[midway,below]{$v^+$};
 \draw[<->](-3,-1.5)--(-1,-1.5)node[midway,below]{$x_0$};\draw[<->](-1,-1.5)--(3,-1.5)node[midway,below]{$D-x_0$};
 \draw[axis](-3,1.7)--(3.6,1.7)node[right]{$x$};
\end{tikzpicture}}
\newcommand{\QAERingQuestion}{\begin{tikzpicture}[>=Latex,scale=.85]
 \draw[thick](0,0)circle(1.3);\fill(-1.3,0)circle(2pt)node[left]{$a$};\fill(.65,1.126)circle(2pt)node[above right]{$b$};
 \draw[->](.92,.92)arc(45:20:1.3);\node[anchor=west,align=left]at(2,.3){均匀电阻丝：总电阻 $R$\\固定 $a$，沿环移动 $b$\\电源仅接于 $a,b$};
\end{tikzpicture}}
\newcommand{\QAECube}{\begin{tikzpicture}[scale=.85,font=\small]
 \coordinate(A)at(-2,-1.2);\coordinate(B)at(.8,-1.2);\coordinate(C)at(.8,1.2);\coordinate(D)at(-2,1.2);
 \coordinate(E)at(-.8,-.3);\coordinate(F)at(2,-.3);\coordinate(G)at(2,2.1);\coordinate(H)at(-.8,2.1);
 \foreach\u/\v in {A/B,B/C,C/D,D/A,E/F,F/G,G/H,H/E,A/E,B/F,C/G,D/H}{\draw[resistor](\u)--(\v);}
 \foreach\u/\pos in {A/below left,B/below,C/above left,D/left,E/below right,F/right,G/above right,H/above}{\fill(\u)circle(1.6pt)node[\pos]{$\u$};}
 \node[anchor=west,align=left]at(3,.4){每条棱均为 $R$\\三组端口：$AB,AC,AG$\\交叉处无黑点则不相连};
\end{tikzpicture}}
\newcommand{\QAELadder}[1]{\begin{tikzpicture}[scale=.85,>=Latex]
 \node[left]at(0,1){$a$};\node[left]at(0,-1){$b$};
 \foreach\i in {0,1,2}{\foreach\y in {-1,1}{\draw[resistor](2*\i,\y)--(2*\i+1.4,\y);\draw(2*\i+1.4,\y)--(2*\i+2,\y);}\draw[resistor](2*\i+1.4,-1)--(2*\i+1.4,1);}
 \draw[dotted,thick](6,1)--(6.8,1) (6,-1)--(6.8,-1);
 \ifnum#1=1\draw[dashed,DeepBlue](1.9,-1.5)rectangle(7,1.5);\node[above,DeepBlue]at(4.4,1.5){尾部仍为 $X$};\fi
 \node[below]at(3.3,-1.8){每只电阻均为 $R$；端口前没有横接电阻};
\end{tikzpicture}}
\newcommand{\QAEThreeRail}[1]{\begin{tikzpicture}[scale=.85,font=\small]
 \foreach\y/\n in {1.5/A,0/B,-1.5/C}{\node[left]at(0,\y){$\n$};}
 \foreach\i in {0,1,2}{\foreach\y in {-1.5,0,1.5}{\draw[resistor](2*\i,\y)--(2*\i+1.4,\y);\draw(2*\i+1.4,\y)--(2*\i+2,\y);}\draw[resistor](2*\i+1.4,-1.5)--(2*\i+1.4,0);\draw[resistor](2*\i+1.4,0)--(2*\i+1.4,1.5);}
 \foreach\y in {-1.5,0,1.5}{\draw[dotted,thick](6,\y)--(6.8,\y);}
 \ifnum#1=1\draw[dashed,DeepBlue](1.9,-2)rectangle(7,2);\node[above,DeepBlue]at(4.3,2){三端尾部：$p,p,q$};\fi
 \node[below]at(3,-2.3){一节含三只水平电阻和两只竖直电阻，均为 $R$};
\end{tikzpicture}}
\newcommand{\QAEThreeReduction}{\begin{tikzpicture}[>=Latex,scale=.85]
 \coordinate(A)at(-1,1.4);\coordinate(B)at(1.5,0);\coordinate(C)at(-1,-1.4);
 \draw[resistor](A)--(B)node[midway,above right]{$p/R$};\draw[resistor](B)--(C)node[midway,below right]{$p/R$};\draw[resistor](C)--(A)node[midway,left]{$q/R$};
 \foreach\n in {A,B,C}{\fill(\n)circle(2pt)node[above left]{$\n$};}
 \node[anchor=west,align=left]at(3,.4){标注量是电导\\$AB,BC$：$p/R$\\$AC$：$q/R$\\第三端口悬空，仍参与分流};
\end{tikzpicture}}
\newcommand{\QAESierpinskiPart}[1]{%
 \ifnum#1=0
  \filldraw[fill=PaleBlue,draw=DeepBlue] (0,0)--(2,0)--(1,1.7320508)--cycle;
 \else
  \pgfmathtruncatemacro{\QAENextDepth}{#1-1}
  \begin{scope}[scale=.5]\QAESierpinskiPart{\QAENextDepth}\end{scope}
  \begin{scope}[shift={(1,0)},scale=.5]\QAESierpinskiPart{\QAENextDepth}\end{scope}
  \begin{scope}[shift={(.5,.8660254)},scale=.5]\QAESierpinskiPart{\QAENextDepth}\end{scope}
 \fi}
\newcommand{\QAESierpinski}{\begin{tikzpicture}[scale=1.15,font=\small]
 \foreach\n in {0,1,2}{\begin{scope}[xshift=3.2*\n cm]
  \QAESierpinskiPart{\n}
  \node[below]at(1,-.3){第 $\n$ 阶：$3^{\n}$ 个小三角形};
 \end{scope}}
\end{tikzpicture}}



本附录补充第19—26章的电磁学专题与例题，重点讨论连续电荷分布、电势、电容器、恒定电流和复杂电阻网络。各题配有完整解答与图示，可结合正文选读。

\Needspace{18\baselineskip}

\section{静电场：连续分布与场线}

\begingroup

\graphicspath{{../../../sample/xes-original/eighteenth-section/}{../../../sample/xes-original/eighteenth-section/assets/}}

% Restrained TikZ redraws for Lesson 29.

\newcommand{\matterStructureDiagram}{%
\begin{center}
\begin{tikzpicture}[scale=0.78,>=Latex]
  % Atom
  \begin{scope}[xshift=-6.0cm]
    \fill[DeepBlue] (0,0) circle (3pt);
    \draw[TextGray,rotate=18] (0,0) ellipse (1.25 and 0.48);
    \draw[TextGray,rotate=-42] (0,0) ellipse (1.25 and 0.48);
    \draw[TextGray,rotate=78] (0,0) ellipse (1.25 and 0.48);
    \fill[Ink] (1.05,0.28) circle (2.2pt);
    \node[TextGray] at (0,-1.55) {原子 $\sim10^{-10}\,\mathrm m$};
  \end{scope}
  % Nucleus
  \begin{scope}[xshift=-2.0cm]
    \foreach \x/\y in {-0.50/0.20,-0.15/0.52,0.25/0.42,0.55/0.05,0.22/-0.35,-0.28/-0.32}{
      \fill[DeepBlue!70] (\x,\y) circle (0.28);
      \draw[Ink] (\x,\y) circle (0.28);
    }
    \node[TextGray] at (0,-1.55) {原子核 $\sim10^{-14}\,\mathrm m$};
  \end{scope}
  % Nucleon
  \begin{scope}[xshift=2.0cm]
    \draw[curve,fill=PaleBlue] (0,0) circle (0.95);
    \foreach \a in {20,80,...,320}{
      \draw[TextGray] (0,0)--({0.68*cos(\a)},{0.68*sin(\a)});
    }
    \node[TextGray] at (0,-1.55) {核子 $\sim10^{-15}\,\mathrm m$};
  \end{scope}
  % Quark content
  \begin{scope}[xshift=6.0cm]
    \draw[curve,fill=PaleBlue] (0,0) circle (0.95);
    \fill[DeepBlue] (-0.38,-0.18) circle (0.22) node[below left] {$u$};
    \fill[Gold] (0.38,-0.18) circle (0.22) node[below right] {$u$};
    \fill[Ink] (0,0.43) circle (0.22) node[above] {$d$};
    \draw[TextGray] (-0.38,-0.18)--(0.38,-0.18)--(0,0.43)--cycle;
    \node[TextGray] at (0,-1.55) {夸克尺度 $<10^{-18}\,\mathrm m$};
  \end{scope}
  \foreach \x in {-4.0,0,4.0}{\draw[vector] (\x-0.55,0)--(\x+0.55,0);}
\end{tikzpicture}

\figcap{从原子到核子内部的物质层级；长度标度为数量级示意}\label{appE:fig:s18-legacy-01}
\end{center}}

\newcommand{\coulombDiagram}{%
\begin{center}
\begin{tikzpicture}[scale=0.9,>=Latex]
  \coordinate (A) at (-2.7,0); \coordinate (B) at (2.7,0);
  \fill[DeepBlue] (A) circle (3pt) node[below=4pt] {$q_1$};
  \fill[DeepBlue] (B) circle (3pt) node[below=4pt] {$q_2$};
  \draw[guide] (A)--(B);
  \draw[<->] (-2.55,-0.75)--node[fill=white,inner sep=1pt] {$r$}(2.55,-0.75);
  \draw[vector] (A)--++(-1.6,0) node[above] {$\vect F_{21}$};
  \draw[vector] (B)--++(1.6,0) node[above] {$\vect F_{12}$};
  \node[TextGray] at (0,1.0) {图示为同号电荷；异号时两力同时反向};
\end{tikzpicture}

\figcap{点电荷间的库仑力沿两电荷连线，且满足作用力与反作用力关系}\label{appE:fig:s18-legacy-02}
\end{center}}

\newcommand{\coulombHistoryContextDiagram}{%
\begin{center}
\begin{tikzpicture}[scale=0.82,>=Latex]
  \node[inner sep=0] at (-4.4,0)
    {\includegraphics[height=3.8cm]{coulomb-portrait-redraw.png}};
  \node[TextGray] at (-4.4,-2.25) {Charles-Augustin de 库仑};
  \begin{scope}[xshift=1.2cm]
    \node[TextGray,anchor=west] at (-2.25,1.75) {俯视};
    \draw[thick] (-2.0,2.0)--(2.0,2.0);
    \draw[thick] (0,2.0)--(0,0.35);
    \draw[very thick] (-2.0,0.25)--(2.0,0.25);
    \fill[DeepBlue] (-2.0,0.25) circle (5pt) node[below] {$q_1$};
    \fill[DeepBlue] (2.0,0.25) circle (5pt) node[below] {$q_2$};
    \fill[Ink] (-2.0,1.35) circle (5pt) node[left] {$Q$};
    \draw[guide] (-2.0,1.35)--(-2.0,0.25);
    % Q is placed tangentially to the suspended arm, so the 库仑 force
    % produces a nonzero torque about the suspension axis.
    \draw[vector] (-2.0,0.25)--(-2.0,-0.62) node[left] {$\vect F$};
    \draw[<->] (0,-0.82)--node[fill=white,inner sep=1pt] {$d$}(-2.0,-0.82);
    \draw[DeepBlue,thick,->] (0,0.75) arc[start angle=90,end angle=25,radius=0.65];
    \node[DeepBlue] at (0.65,0.9) {$\theta$};
    \node[TextGray,align=left,anchor=west] at (3.0,0.55)
      {悬丝扭转\\$Fd=\kappa\theta$\\$\theta\propto r^{-2}$};
  \end{scope}
\end{tikzpicture}

\figcap{库仑与扭秤测力原理：静电力矩由悬丝的回复力矩平衡}\label{appE:fig:s18-legacy-03}
\end{center}}

\newcommand{\fieldLineDiagram}{%
\begin{center}
\begin{tikzpicture}[scale=0.72,>=Latex]
  \begin{scope}[xshift=-6.0cm]
    \fill[DeepBlue] (-0.75,0) circle (4pt) node[below] {$+Q$};
    \fill[DeepBlue] (0.75,0) circle (4pt) node[below] {$+Q$};
    \draw[->,DeepBlue,thick] (-0.98,0)--(-2.35,0);
    \draw[->,DeepBlue,thick] (0.98,0)--(2.35,0);
    \foreach \s in {-1,1}{
      \draw[->,DeepBlue,thick] ({0.75*\s},{0.18}) .. controls ({1.25*\s},0.65) and ({1.45*\s},1.45) .. ({1.55*\s},2.1);
      \draw[->,DeepBlue,thick] ({0.75*\s},{-0.18}) .. controls ({1.25*\s},-0.65) and ({1.45*\s},-1.45) .. ({1.55*\s},-2.1);
      \draw[->,DeepBlue,thick] ({0.62*\s},0.16) .. controls ({0.35*\s},0.85) and ({0.25*\s},1.55) .. ({0.18*\s},2.15);
      \draw[->,DeepBlue,thick] ({0.62*\s},-0.16) .. controls ({0.35*\s},-0.85) and ({0.25*\s},-1.55) .. ({0.18*\s},-2.15);
    }
    \node[TextGray] at (0,-2.55) {等量同号电荷};
  \end{scope}
  \begin{scope}
    \fill[DeepBlue] (-1.4,0) circle (4pt) node[below] {$+Q$};
    \fill[Ink] (1.4,0) circle (4pt) node[below] {$-Q$};
    \foreach \h in {0.28,0.62,1.0,1.45}{
      \draw[->,DeepBlue,thick]
        (-1.16,{0.55*\h}) .. controls (-0.35,{1.45*\h}) and (0.35,{1.45*\h}) .. (1.16,{0.55*\h});
      \draw[->,DeepBlue,thick]
        (-1.16,{-0.55*\h}) .. controls (-0.35,{-1.45*\h}) and (0.35,{-1.45*\h}) .. (1.16,{-0.55*\h});
    }
    \draw[->,DeepBlue,thick] (-1.05,0)--(1.05,0);
    \node[TextGray] at (0,-2.55) {等量异号电荷};
  \end{scope}
  \begin{scope}[xshift=6.0cm]
    \fill[DeepBlue] (0,0) circle (4pt) node[below right] {$q(t)$};
    \draw[guide,dashed] (0,-2.1)--(0,2.1) node[above] {振动轴};
    \draw[->,DeepBlue,thick] (0.15,0.05) .. controls (0.55,0.35) and (0.25,0.95) .. (0.95,1.30)
      .. controls (1.55,1.60) and (1.25,2.05) .. (2.25,2.25);
    \draw[->,DeepBlue,thick] (-0.15,0.05) .. controls (-0.55,0.35) and (-0.25,0.95) .. (-0.95,1.30)
      .. controls (-1.55,1.60) and (-1.25,2.05) .. (-2.25,2.25);
    \draw[->,DeepBlue,thick] (0.15,-0.05) .. controls (0.55,-0.35) and (0.25,-0.95) .. (0.95,-1.30)
      .. controls (1.55,-1.60) and (1.25,-2.05) .. (2.25,-2.25);
    \draw[->,DeepBlue,thick] (-0.15,-0.05) .. controls (-0.55,-0.35) and (-0.25,-0.95) .. (-0.95,-1.30)
      .. controls (-1.55,-1.60) and (-1.25,-2.05) .. (-2.25,-2.25);
    \node[TextGray] at (0,-2.55) {振荡电荷：场线随时间演化};
  \end{scope}
\end{tikzpicture}

\figcap{同号、异号与振荡电荷的典型电场线；第三幅不能按静电场理解}\label{appE:fig:s18-legacy-04}
\end{center}}

\newcommand{\fluxDiagram}{%
\begin{center}
\begin{tikzpicture}[scale=0.9,>=Latex]
  \coordinate (O) at (0,0);
  \coordinate (Etip) at (2.45,1.25);
  \coordinate (Stip) at (0.7,2.15);
  \draw[fill=PaleBlue,draw=DeepBlue,thick] (-1.8,-0.85)--(1.8,-0.30)--(1.8,1.05)--(-1.8,0.50)--cycle;
  \draw[vector] (O)--++(0.7,2.15) node[above] {$\vect S$};
  \draw[vector] (O)--++(2.45,1.25) node[right] {$\vect E$};
  \draw pic[draw,"$\theta$",angle radius=8mm,angle eccentricity=1.3]
    {angle=Etip--O--Stip};
  \node[TextGray] at (0,-1.35) {$\Phi_E=\vect E\cdot\vect S=ES\cos\theta$};
\end{tikzpicture}
\end{center}}

\newcommand{\gaussDiagram}{%
\begin{center}
\begin{tikzpicture}[scale=0.88,>=Latex]
  \draw[DeepBlue,thick] (0,0) ellipse (2.7 and 1.75);
  \fill[DeepBlue] (-0.45,0.15) circle (4pt) node[below right] {$q_{\rm in}$};
  \fill[Ink] (3.9,0.65) circle (4pt) node[right] {$q_{\rm out}$};
  \foreach \a in {10,55,...,325}{
    \draw[->,TextGray,thin] (-0.45,0.15)--++({1.9*cos(\a)},{1.35*sin(\a)});
  }
  \draw[vector] (2.35,0.88)--++(0.85,0.35) node[right] {$\dd\vect S$};
  \node[TextGray,align=left,anchor=west] at (5.0,0)
    {闭合面的净通量只由面内净电荷决定；\\面外电荷会改变局部电场，但净通量贡献为零。};
\end{tikzpicture}

\figcap{高斯定律中的闭合曲面}\label{appE:fig:s18-legacy-05}
\end{center}}

\newcommand{\gaussContextDiagram}{%
\begin{center}
\begin{tikzpicture}[scale=0.78,>=Latex]
  \node[inner sep=0] at (-6.0,0.2)
    {\includegraphics[width=2.2cm]{gauss-portrait-redraw.png}};
  \node[TextGray] at (-6.0,-1.65) {高斯};
  \begin{scope}[xshift=-2.8cm]
    \fill[PaleBlue] (-1.15,-0.6)--(1.15,-0.25)--(0.9,0.8)--(-1.3,0.5)--cycle;
    \draw[thick] (-1.15,-0.6)--(1.15,-0.25)--(0.9,0.8)--(-1.3,0.5)--cycle;
    \draw[vector] (0,0)--(0.55,1.55) node[above] {$\dd\vect S$};
    \draw[vector] (0,0)--(1.65,0.65) node[right] {$\vect E$};
    \coordinate (Ee) at (1.65,0.65);
    \coordinate (Oo) at (0,0);
    \coordinate (Ss) at (0.55,1.55);
    \draw pic[draw,"$\theta$",angle radius=6mm,angle eccentricity=1.3]
      {angle=Ee--Oo--Ss};
    \node[DeepBlue] at (0,-1.55) {$\dd\Phi_E=E\dd S\cos\theta$};
  \end{scope}
  \begin{scope}[xshift=1.2cm]
    \draw[DeepBlue,thick,rounded corners=12pt] plot[smooth cycle] coordinates
      {(-1.45,0.1) (-0.8,1.25) (0.7,1.1) (1.45,0.2) (0.8,-1.2) (-0.9,-1.0)};
    \fill[Ink] (0,0) circle (3pt) node[below right] {$q$};
    \foreach \a in {15,60,...,330}{
      \draw[->,TextGray] ({0.25*cos(\a)},{0.25*sin(\a)})--({1.05*cos(\a)},{0.85*sin(\a)});
    }
    \node[DeepBlue] at (0,-1.55) {$\Phi_E=q/\varepsilon_0$};
  \end{scope}
  \begin{scope}[xshift=5.35cm]
    \draw[DeepBlue,thick,rounded corners=12pt] plot[smooth cycle] coordinates
      {(-1.35,0.1) (-0.75,1.15) (0.75,1.0) (1.4,0.15) (0.75,-1.15) (-0.8,-1.0)};
    \fill[Ink] (-2.25,0) circle (3pt) node[below] {$q$};
    \foreach \y in {-0.65,-0.2,0.25,0.7}{
      \draw[->,TextGray] (-2.0,\y)--(1.0,{0.75*\y});
    }
    \node[DeepBlue] at (0,-1.55) {$\Phi_E=0$};
  \end{scope}
\end{tikzpicture}

\figcap{面积元通量、面内电荷与面外电荷三种几何情形}\label{appE:fig:s18-legacy-06}
\end{center}}

\newcommand{\gaussDifferentialCubeDiagram}{%
\begin{center}
\begin{tikzpicture}[scale=0.88,>=Latex]
  \coordinate (A) at (0,0); \coordinate (B) at (3.4,0.35);
  \coordinate (C) at (4.65,1.35); \coordinate (D) at (1.25,1.0);
  \coordinate (E) at (0,2.55); \coordinate (F) at (3.4,2.90);
  \coordinate (G) at (4.65,3.90); \coordinate (H) at (1.25,3.55);
  \draw[curve] (A)--(B)--(C)--(D)--cycle;
  \draw[curve] (E)--(F)--(G)--(H)--cycle;
  \draw[curve] (A)--(E) (B)--(F) (C)--(G) (D)--(H);
  \draw[axis] (A)--++(5.5,0.55);
  \node[below] at (5.70,0.10) {$x$};
  \draw[axis] (A)--++(1.65,1.33) node[above right] {$y$};
  \draw[axis] (A)--++(0,4.45) node[above] {$z$};
  \draw[vector] (B)--++(1.25,0.13);
  \node[anchor=west,fill=white,inner sep=0.8pt] at (5.10,1.02)
    {$E_x(x+\Delta x)$};
  \draw[TextGray,thin] (5.10,0.91)--(4.68,0.51);
  \draw[vector] (A)--++(-1.1,-0.11) node[below] {$E_x(x)$};
  \draw[vector] (H)--++(0,1.15) node[left] {$E_z(z+\Delta z)$};
  \node[TextGray,align=left,anchor=west] at (6.45,2.55)
    {六个面的净通量相加\\保留一阶差分\\$\displaystyle \nabla\!\cdot\!\vect E=\rho/\varepsilon_0$};
  \node[DeepBlue] at (2.0,-0.55) {$\Delta x$};
  \node[DeepBlue,fill=white,inner sep=0.8pt] at (4.85,1.72) {$\Delta y$};
  \node[DeepBlue] at (-0.35,1.3) {$\Delta z$};
\end{tikzpicture}

\figcap{无穷小长方体上的通量收支给出高斯定律的微分形式}\label{appE:fig:s18-legacy-07}
\end{center}}

\newcommand{\lineFieldDiagram}{%
\begin{center}
\begin{tikzpicture}[scale=0.84,>=Latex]
  \draw[very thick] (0,-2.2)--(0,2.2);
  \node[left] at (0,2.0) {$\lambda$};
  \coordinate (P) at (3.0,0);
  \fill[DeepBlue] (P) circle (2.4pt) node[right] {$P$};
  \draw[<->] (0,-0.45)--node[fill=white,inner sep=1pt] {$a$}(3.0,-0.45);
  \draw[vector] (P)--++(1.5,0) node[above] {$\vect E$};
\end{tikzpicture}
\end{center}}

\newcommand{\finiteLineDiagram}{%
\begin{center}
\begin{tikzpicture}[scale=0.9,>=Latex]
  \coordinate (O) at (0,0); \coordinate (P) at (3.2,0);
  \coordinate (T) at (0,2.3); \coordinate (D) at (0,-2.3);
  \draw[very thick] (0,-2.3)--(0,2.3);
  \fill[DeepBlue] (P) circle (2.4pt) node[right] {$P$};
  \draw[<->] (0,-0.5)--node[fill=white,inner sep=1pt] {$a$}(3.2,-0.5);
  \draw[guide] (P)--(0,2.3) (P)--(0,-2.3);
  \draw pic[draw,"$\theta_1$",angle radius=8mm,angle eccentricity=1.35]
    {angle=T--P--O};
  \draw pic[draw,"$\theta_2$",angle radius=11mm,angle eccentricity=1.25]
    {angle=O--P--D};
  \draw[vector] (P)--++(1.35,0) node[above] {$\vect E$};
\end{tikzpicture}

\figcap{有限长带电直线中垂线上的场点}\label{appE:fig:s18-legacy-08}
\end{center}}

\newcommand{\diskFieldDiagram}{%
\begin{center}
\begin{tikzpicture}[scale=0.88,>=Latex]
  \coordinate (O) at (0,0); \coordinate (P) at (4.0,0);
  \coordinate (T) at (0,2.1);
  \draw[fill=PaleBlue,draw=DeepBlue,thick] (O) ellipse (0.65 and 2.1);
  \draw[DeepBlue] (O) ellipse (0.38 and 1.2);
  \fill (O) circle (1.6pt) node[below left] {$O$};
  \fill[DeepBlue] (P) circle (2.4pt) node[right] {$P$};
  \draw[axis] (-0.8,0)--(5.3,0) node[right] {$x$};
  \draw[<->] (0,-2.5)--node[below] {$a$}(4.0,-2.5);
  \draw[guide] (P)--(0,2.1);
  \draw pic[draw,"$\theta$",angle radius=9mm,angle eccentricity=1.35]
    {angle=T--P--O};
  \node[left] at (-0.65,1.7) {$\sigma$};
\end{tikzpicture}

\figcap{均匀带电圆盘轴线上的场点；图中 $a$ 表示场点到圆盘的距离}\label{appE:fig:s18-legacy-09}
\end{center}}

\newcommand{\unequalChargeDiagram}{%
\begin{center}
\begin{tikzpicture}[scale=0.9,>=Latex]
  \coordinate (A) at (-2.1,0); \coordinate (B) at (2.1,0); \coordinate (N) at (5.4,0);
  \fill[DeepBlue] (A) circle (4pt) node[below] {$+Q$};
  \fill[Ink] (B) circle (4pt) node[below] {$-q$};
  \fill[DeepBlue] (N) circle (2.4pt) node[below] {$N$};
  \draw[axis] (-3.2,0)--(6.6,0) node[right] {$x$};
  \draw[<->] (-2.1,-0.8)--node[fill=white,inner sep=1pt] {$2a$}(2.1,-0.8);
  \draw[<->] (2.1,0.75)--node[fill=white,inner sep=1pt] {$x_0-a$}(5.4,0.75);
  \node[TextGray] at (1.65,1.45) {$Q>q>0$，零场点位于小电荷外侧};
\end{tikzpicture}

\figcap{不等量异号点电荷的零场点}\label{appE:fig:s18-legacy-10}
\end{center}}

\newcommand{\cavityDiagram}{%
\begin{center}
\begin{tikzpicture}[scale=0.9,>=Latex]
  \coordinate (O) at (0,0); \coordinate (C) at (1.15,0.40); \coordinate (P) at (1.82,1.12);
  \draw[fill=PaleBlue,draw=DeepBlue,thick] (O) circle (2.40);
  \draw[fill=white,draw=Ink,thick] (C) circle (1.02);
  \fill (O) circle (1.5pt) node[below left=2pt,fill=white,inner sep=0.5pt] {$O$};
  \fill (C) circle (1.5pt) node[below=3pt,fill=white,inner sep=0.5pt] {$C$};
  \fill[DeepBlue] (P) circle (2.2pt) node[above right=2pt,fill=white,inner sep=0.5pt] {$P$};
  \draw[vector] (O)--(C) node[pos=0.52,below=3pt,fill=white,inner sep=0.7pt] {$\vect d$};
  \draw[vector] (C)--(P) node[pos=0.55,right=4pt,fill=white,inner sep=0.7pt] {$\vect r'$};
  \draw[vector] (O)--(P) node[pos=0.48,above left=3pt,fill=white,inner sep=0.7pt] {$\vect r$};
  \node[TextGray,align=left,anchor=west] at (3.55,0.25)
    {均匀带正电球体\\叠加同位置、同大小、带负电的球体\\以挖去偏心球形空腔};
\end{tikzpicture}

\figcap{球形空腔的叠加模型，$\vect r=\vect d+\vect r'$}\label{appE:fig:s18-legacy-11}
\end{center}}

\newcommand{\multipoleDiagram}{%
\begin{center}
\begin{tikzpicture}[scale=0.9,>=Latex]
  \begin{scope}[xshift=-3.8cm]
    \coordinate (O1) at (0,0); \coordinate (X1) at (2.7,0); \coordinate (P1) at (2.7,2.3);
    \fill[Ink] (-0.75,0) circle (3pt) node[below] {$-q$};
    \fill[DeepBlue] (0.75,0) circle (3pt) node[below] {$+q$};
    \draw[<->] (-0.7,-0.75)--node[fill=white,inner sep=1pt] {$l$}(0.7,-0.75);
    \draw[axis] (0,0)--(2.7,2.3) node[right] {$P(r,\theta)$};
    \draw pic[draw,"$\theta$",angle radius=8mm,angle eccentricity=1.35]
      {angle=X1--O1--P1};
    \node[TextGray] at (0,-1.45) {电偶极子};
  \end{scope}
  \begin{scope}[xshift=4.0cm]
    \coordinate (O2) at (0,0); \coordinate (X2) at (2.7,0); \coordinate (P2) at (2.6,2.3);
    \fill[Ink] (-1.0,0) circle (3pt) node[below] {$-q$};
    \fill[DeepBlue] (0,0) circle (3pt) node[below] {$+2q$};
    \fill[Ink] (1.0,0) circle (3pt) node[below] {$-q$};
    \draw[<->] (-0.95,-0.75)--node[fill=white,inner sep=1pt] {$l$}(-0.05,-0.75);
    \draw[<->] (0.05,-0.75)--node[fill=white,inner sep=1pt] {$l$}(0.95,-0.75);
    \draw[axis] (0,0)--(2.6,2.3) node[right] {$P(r,\theta)$};
    \draw pic[draw,"$\theta$",angle radius=8mm,angle eccentricity=1.35]
      {angle=X2--O2--P2};
    \node[TextGray] at (0,-1.45) {直线电四极子};
  \end{scope}
\end{tikzpicture}

\figcap{例题 29.4 与例题 29.5 的几何关系}\label{appE:fig:s18-legacy-12}
\end{center}}

\newcommand{\dipoleDiagram}{%
\begin{center}
\begin{tikzpicture}[scale=0.95,>=Latex]
  \coordinate (O) at (0,0); \coordinate (X) at (3.3,0); \coordinate (P) at (3.3,2.5);
  \fill[Ink] (-0.85,0) circle (3pt) node[below] {$-q$};
  \fill[DeepBlue] (0.85,0) circle (3pt) node[below] {$+q$};
  \draw[<->] (-0.8,-0.75)--node[fill=white,inner sep=1pt] {$l$}(0.8,-0.75);
  \draw[axis] (O)--(P) node[right] {$P(r,\theta)$};
  \draw[guide] (O)--(X);
  \draw pic[draw,"$\theta$",angle radius=9mm,angle eccentricity=1.35]
    {angle=X--O--P};
  \draw[vector] (-0.75,0.55)--(0.75,0.55) node[midway,above] {$\vect p$};
\end{tikzpicture}

\figcap{电偶极矩 $\vect p$ 从负电荷指向正电荷}\label{appE:fig:s18-legacy-13}
\end{center}}

\newcommand{\quadrupoleDiagram}{%
\begin{center}
\begin{tikzpicture}[scale=0.95,>=Latex]
  \coordinate (O) at (0,0); \coordinate (X) at (3.3,0); \coordinate (P) at (3.1,2.5);
  \fill[Ink] (-1.1,0) circle (3pt) node[below] {$-q$};
  \fill[DeepBlue] (0,0) circle (3pt) node[below] {$+2q$};
  \fill[Ink] (1.1,0) circle (3pt) node[below] {$-q$};
  \draw[<->] (-1.05,-0.75)--node[fill=white,inner sep=1pt] {$l$}(-0.05,-0.75);
  \draw[<->] (0.05,-0.75)--node[fill=white,inner sep=1pt] {$l$}(1.05,-0.75);
  \draw[axis] (O)--(P) node[right] {$P(r,\theta)$};
  \draw[guide] (O)--(X);
  \draw pic[draw,"$\theta$",angle radius=9mm,angle eccentricity=1.35]
    {angle=X--O--P};
\end{tikzpicture}

\figcap{直线电四极子的几何关系}\label{appE:fig:s18-legacy-14}
\end{center}}

\newcommand{\slabDiagram}{%
\begin{center}
\begin{tikzpicture}[scale=0.9,>=Latex]
  \coordinate (A) at (2.65,0.52); \coordinate (B) at (3.85,0.52); \coordinate (C) at (3.85,0.35);
  \fill[PaleBlue] (-2.1,-1.15) rectangle (2.1,1.15);
  \draw[DeepBlue,thick] (-2.1,-1.15) rectangle (2.1,1.15);
  \foreach \x in {-1.7,-1.0,...,1.8}
    \foreach \y in {-0.75,0,0.75}
      \fill[DeepBlue] (\x,\y) circle (1.2pt);
  \draw[axis] (-4.0,0)--(4.0,0) node[right] {$x$};
  \node[below] at (-2.1,-1.2) {$-d$}; \node[below] at (2.1,-1.2) {$d$};
  \fill[Ink] (0,0) circle (2.8pt) node[above] {$m,q$};
  \draw[vector] (0,0)--++(1.35,0) node[above] {$\vect v_0$};
  \node[TextGray] at (0,1.55) {体电荷密度 $\rho>0$};
\end{tikzpicture}

\figcap{仅在 $-d<x<d$ 内均匀带电的无穷大厚板}\label{appE:fig:s18-legacy-15}
\end{center}}

\newcommand{\electronLensDiagram}{%
\begin{center}
\begin{tikzpicture}[scale=0.9,>=Latex]
  \draw[very thick] (0,-2.0)--(0,2.0) node[above] {正线电荷 $\lambda$};
  \draw[axis] (-4.4,0)--(4.5,0) node[right] {$x$};
  \draw[guide] (-3.8,1.05)--(3.8,1.05);
  \draw[DeepBlue,very thick,-{Latex[length=2.5mm]}]
    (-3.8,1.05) .. controls (-1.2,1.05) and (1.2,0.66) .. (3.8,0.35);
  \draw[<->] (-3.1,0)--node[left] {$b$}(-3.1,1.05);
  \draw[vector] (-3.7,1.05)--++(1.35,0) node[above] {$\vect v$};
  \draw[guide] (2.65,0.52)--++(1.2,0);
  \draw pic[draw,"$\alpha$",angle radius=8mm,angle eccentricity=1.3]
    {angle=B--A--C};
\end{tikzpicture}

\figcap{电子掠过正线电荷时受到指向金属线的微弱偏转}\label{appE:fig:s18-legacy-16}
\end{center}}


\begin{origexamplebox}{【统统算出来】}\label{appE:ex:s18-29-1}
利用积分或高斯定律计算下列均匀带电系统产生的电场：
\begin{enumerate}[label=(\arabic*),leftmargin=2.2em]
  \item 无限长直导线外距导线 $a$ 处，线电荷密度为 $\lambda$；
  \item 有限长直导线所在直线外、距垂足 $a$ 处（垂足在线段内），两端相对过垂足法线的张角分别为 $\theta_1,\theta_2$；
  \item 无限大均匀带电平面外，面电荷密度为 $\sigma$；
  \item 仅一个方向无限延伸的均匀带电条带所在平面外距平面 $a$ 处（垂足在条带内），两侧边缘所张角为
  $\theta_1,\theta_2$；
  \item 半径为 $R$、面电荷密度为 $\sigma$ 的圆盘轴线上，距圆盘中心 $a$ 处。
\end{enumerate}
\QAENamed{\QAESources}{直线、平面、条带与圆盘的题设几何；场点与电荷均不重合，距离为正。}{quality:appE:q1}
\end{origexamplebox}

\begin{xesolution}
记 $k=1/(4\pi\varepsilon_0)$。

\textbf{(1) 无限长直导线。}取与导线同轴、半径 $a$、长度 $L$ 的圆柱面。电场沿径向，
且仅圆柱侧面有通量。由高斯定律
\[
  E(2\pi aL)=\frac{\lambda L}{\varepsilon_0},
\]
故
\[
  \boxed{E=\frac{\lambda}{2\pi\varepsilon_0a}}.
\]
当 $\lambda>0$ 时方向背离导线，$\lambda<0$ 时方向指向导线。
\QAENamed{\lineFieldDiagram}{同轴高斯面只在侧面产生通量。}{quality:appE:a1-line}

\textbf{(2) 有限长直导线。}以过垂足指向场点的法线为法向正方向，沿导线取坐标 $x$，电荷元
$\dd q=\lambda\dd x$ 到场点的距离为 $r=\sqrt{a^2+x^2}$。它产生的场元为
\[
  \dd E=\frac{k\lambda\dd x}{a^2+x^2}.
\]
将场元分解到法向和导线方向，并令 $x=a\tan\theta$，得到
\[
  \boxed{E_\perp=\frac{k\lambda}{a}
  (\sin\theta_1+\sin\theta_2)},
\]
\[
  \boxed{E_\parallel=\frac{k\lambda}{a}
  (\cos\theta_1-\cos\theta_2)}.
\]
$E_\parallel$ 的正方向取向张角为 $\theta_1$ 的一端；若线段关于过垂足的法线对称，
$\theta_1=\theta_2=\theta$，则切向分量抵消，
\[
  \boxed{E=\frac{2k\lambda}{a}\sin\theta}.
\]
令 $\theta\to\pi/2$，即回到无限长直导线的结果。

\textbf{(3) 无限大平面。}取跨过平面的薄圆柱高斯面。由平面对称性，两底面的电场大小
相同且垂直平面，侧面无通量，因而
\[
  2ES=\frac{\sigma S}{\varepsilon_0},
  \qquad
  \boxed{E=\frac{\sigma}{2\varepsilon_0}}.
\]

\textbf{(4) 无限长条带。}把条带分成许多与无限方向平行的细线，每条细线的线电荷密度
$\dd\lambda=\sigma\dd x$。设场点到条带所在平面的垂足为坐标原点，条带两边分别位于
$x=-a\tan\theta_2$ 和 $x=a\tan\theta_1$。叠加无限长直线的电场，有
\[
  \boxed{E_\perp=\frac{\sigma}{2\pi\varepsilon_0}
  (\theta_1+\theta_2)},
\]
\[
  \boxed{E_\parallel=
  \frac{\sigma}{2\pi\varepsilon_0}
  \ln\frac{\cos\theta_1}{\cos\theta_2}}.
\]
这里 $E_\parallel$ 的正方向取向张角为 $\theta_1$ 的一侧。
若场点位于条带中垂面，$\theta_1=\theta_2=\theta$，则
\[
  \boxed{E=\frac{\sigma\theta}{\pi\varepsilon_0}}.
\]
当条带宽度趋于无穷，$\theta\to\pi/2$，再次得到 $E=\sigma/(2\varepsilon_0)$。

\textbf{(5) 均匀带电圆盘。}把圆盘分成半径 $s$、宽度 $\dd s$ 的同心圆环，
$\dd q=2\pi\sigma s\dd s$。圆环在轴线场点处的横向分量相消，轴向场元为
\[
  \dd E_x=k\frac{a\,\dd q}{(a^2+s^2)^{3/2}}.
\]
从 $s=0$ 积分到 $R$，得到
\[
  \boxed{E_x=\frac{\sigma}{2\varepsilon_0}
  \left(1-\frac{a}{\sqrt{a^2+R^2}}\right)}.
\]
若圆盘边缘对场点的半张角为 $\theta$，则
\[
  \boxed{E_x=\frac{\sigma}{2\varepsilon_0}(1-\cos\theta)}.
\]
\QAENamed{\diskFieldDiagram}{圆盘分割为同心圆环，各环横向分量相消。}{quality:appE:a1-disk}

\QAENamed{\QAEStrip}{条带的线元分割与无限平面的双端面高斯面。}{quality:appE:a1}
图~\ref{quality:appE:a1}左图用 $x=a\tan\theta$ 将条带线元换成角度；每条线元的水平场沿场点到线元方向的反向，因此积分后水平分量含对数而非角度和。右图两个端面都要计入通量。图~\ref{quality:appE:a1-line}与图~\ref{quality:appE:a1-disk}分别对应同轴高斯法和同心环积分法。
\end{xesolution}

\begin{origexamplebox}{【奇妙的电场线】}\label{appE:ex:s18-29-2}
点电荷 $+Q$、$-q$ 分别位于 $(-a,0,0)$、$(a,0,0)$，且 $Q>q>0$。系统关于两电荷连线
旋转对称。回答：
\begin{enumerate}[label=(\arabic*),leftmargin=2.2em]
  \item $+Q$ 发出的电场线能否到达无穷远？$-q$ 接收的电场线是否来自无穷远？
  \item 空间中是否存在电场为零的点？该点附近的电场线有何特征？
  \item 无穷远处的电场有什么特征？
  \item 写出电场线方程的一种隐式表示。
\end{enumerate}
\QAENamed{\unequalChargeDiagram}{不等量异号电荷位于同一轴线上，$Q>q>0$；场线数仅表示通量比例。}{quality:appE:q2}
\end{origexamplebox}

\begin{xesolution}
\textbf{(1)} 电场线条数只能用来表示通量比例，不是绝对数量。由于净电荷为 $Q-q>0$，
$+Q$ 发出的电通量中有一部分终止于 $-q$，其余必然延伸到无穷远。$-q$ 所接收的全部电场线
都可来自 $+Q$，不需要由无穷远进入。

\textbf{(2)} 零场点只能在两电荷连线延长线上。两电荷之间的电场同向，不可能抵消；
在 $+Q$ 左侧，大电荷又更近，也不能抵消。因此零场点在 $-q$ 的右侧。设其坐标为 $x_0>a$，
则
\[
  \frac{kQ}{(x_0+a)^2}=\frac{kq}{(x_0-a)^2},
\]
解得
\[
  \boxed{x_0=a\frac{\sqrt Q+\sqrt q}{\sqrt Q-\sqrt q}}.
\]
零场点附近的场具有鞍点结构：沿两电荷连线方向，电场线从零场点两侧离开；沿垂直方向，
电场线向该轴线汇聚。电场线本身不穿过电场为零的点，因为该点没有确定的切向方向。

\textbf{(3)} 当 $r\gg a$ 时，观察者分辨不出两电荷的间距，最低阶电场等效于位于原点的
净电荷 $Q-q$：
\[
  \boxed{\vect E\simeq k\frac{Q-q}{r^2}\uvec r}.
\]
因此无穷远处的电场线近似沿径向向外，偶极修正比主项多衰减一个 $1/r$。

\textbf{(4)} 在包含两电荷的任一子午面内，设场点相对于 $+Q$、$-q$ 的极角分别为
$\theta_+$、$\theta_-$，二者的极轴都取沿 $+x$ 方向。由轴对称电场的通量函数可得，
每一条电场线满足
\[
  \boxed{Q(1-\cos\theta_+)-q(1-\cos\theta_-)=C},
\]
其中 $C$ 由所选电场线决定。等价地，也可写成
\[
  Q\cos\theta_+-q\cos\theta_-=C'.
\]
这是一种隐式方程；绕 $x$ 轴旋转即可得到三维电场线族。

\QAENamed{\QAEZeroField}{零场点两侧的轴向场与垂直扰动方向；只画零场点邻域的线性结构。}{quality:appE:a2}

图~\ref{quality:appE:a2}中较小负电荷距 $x_0$ 更近，才能与较大的正电荷抵消。令 $A=\partial E_x/\partial x|_{x_0}>0$，无源邻域的轴对称性和 $\nabla\cdot\vect E=0$ 给出横向线性系数 $-A/2$，从而得到图示鞍形方向。电荷之间并不存在这样的零场点。
\end{xesolution}

\begin{origexamplebox}{【球壳空洞】}\label{appE:ex:s18-29-3}
一个球体内部挖去偏心球形空腔，未挖去部分具有均匀体电荷密度 $\rho$。
\begin{enumerate}[label=(\arabic*),leftmargin=2.2em]
  \item 求空腔内的电场；
  \item 分析它与万有引力中同类模型的联系，并讨论题目若把球体称为“金属球体”会有什么问题。
\end{enumerate}
\QAENamed{\cavityDiagram}{均匀带电绝缘球中的偏心球形空腔；空腔完全处于大球内。}{quality:appE:q3}
\end{origexamplebox}

\begin{xesolution}
设完整大球球心为 $O$，空腔球心为 $C$，$\vect d=\overrightarrow{OC}$。把实际电荷分布分解为：
以 $O$ 为球心、密度为 $\rho$ 的完整均匀带电球，加上以 $C$ 为球心、密度为 $-\rho$ 的小球。
均匀带电实心球内部距球心位置矢量为 $\vect r$ 的点，电场为
\[
  \vect E=\frac{\rho}{3\varepsilon_0}\vect r.
\]
对空腔内任一点 $P$，记 $\vect r=\overrightarrow{OP}$、$\vect r'=\overrightarrow{CP}$，
则 $\vect r=\vect d+\vect r'$。叠加两球的电场：
\[
  \vect E=\frac{\rho}{3\varepsilon_0}\vect r
  -\frac{\rho}{3\varepsilon_0}\vect r'
  =\boxed{\frac{\rho}{3\varepsilon_0}\vect d}.
\]
所以空腔内电场是匀强电场，当 $\rho>0$ 时方向由大球球心指向空腔球心。

万有引力也服从平方反比定律，均匀球形空腔可用完全相同的叠加方法计算，只需把电学常量与
质量密度替换为引力常量与质量密度；引力只有吸引，最终场方向相应改变。

若把球体称为“金属球体”并同时规定内部均匀体电荷密度，这在静电平衡下不成立。导体内
自由电荷会重新分布，导体材料内部电场必须为零。因此上述答案适用于均匀体电荷被冻结的
绝缘球体；若确为处于静电平衡的导体，必须按导体边界条件重新求解。

\QAENamed{\QAECavityVectors}{带正、负密度的两个完整球叠加；$\vect r=\vect d+\vect r'$。}{quality:appE:a3}

图~\ref{quality:appE:a3}的两条位置矢量都终止于同一 $P$，相减后只剩与 $P$ 无关的 $\vect d$，这正是空腔内场均匀的原因。若 $\rho<0$，方向随密度符号反转；万有引力对应 $\vect g=-4\pi G\rho_m\vect d/3$。
\end{xesolution}

\begin{origexamplebox}{【电偶极子】}\label{appE:ex:s18-29-4}
一正一负两个点电荷的电荷量分别为 $+q$、$-q$，间距为 $l$。求 $r\gg l$ 处电场强度的
大小和方向，并求在该近似下的电场线方程。
\QAENamed{\dipoleDiagram}{两个等量异号点电荷与远场位置；偶极矩方向由负电荷指向正电荷。}{quality:appE:q4}
\end{origexamplebox}

\begin{xesolution}
定义电偶极矩
\[
  \vect p=q\vect l,
\]
方向由负电荷指向正电荷。取偶极子中心为极点，以 $\vect p$ 方向为极轴。当 $r\gg l$ 时，
两点电荷势的最低阶非零项为
\[
  \varphi(r,\theta)=k\frac{p\cos\theta}{r^2}.
\]
由 $\vect E=-\nabla\varphi$ 得
\[
  \boxed{E_r=\frac{2kp\cos\theta}{r^3}},
  \qquad
  \boxed{E_\theta=\frac{kp\sin\theta}{r^3}}.
\]
因此
\[
  \boxed{E=\frac{kp}{r^3}\sqrt{1+3\cos^2\theta}}.
\]
方向由 $E_r,E_\theta$ 共同确定；轴线上电场沿 $\vect p$，中垂面上电场与 $\vect p$ 反向。

电场线满足
\[
  \frac{\dd r}{r\dd\theta}=\frac{E_r}{E_\theta}=2\cot\theta.
\]
积分后得到
\[
  \boxed{r=C\sin^2\theta},
\]
其中 $C>0$ 标记不同的电场线。

\QAENamed{\QAEDipoleComponents}{远场中的径向、极角方向及电场分量；图示第一象限。}{quality:appE:a4}

图~\ref{quality:appE:a4}中 $\theta$ 沿从正极轴转向 $P$ 的方向增大，故 $E_\theta$ 的正方向是逆时针的局部切向。把 $E_r,E_\theta$ 按这两条正交基矢合成，才能得到空间方向。$r=C\sin^2\theta$ 只在理想偶极或远场有效，曲线向 $r\sim l$ 延伸的部分不再代表实际两电荷的近场。
\end{xesolution}

\endgroup

\Needspace{18\baselineskip}

\section{电势与相互作用能}

\begingroup

\graphicspath{{../../../sample/xes-original/nineteenth-section/}{../../../sample/xes-original/nineteenth-section/assets/}}

% Restrained TikZ redraws for Lesson 30.

\newcommand{\potentialGradientDiagram}{%
\begin{center}
\begin{tikzpicture}[scale=0.72,>=Latex]
  \begin{scope}[xshift=-6.0cm]
    \foreach \x in {-1.8,-0.9,0,0.9,1.8}{\draw[DeepBlue,thick] (\x,-1.65)--(\x,1.65);}
    \foreach \y in {-1.15,-0.38,0.38,1.15}{\draw[->,TextGray] (-2.15,\y)--(2.15,\y);}
    \node[TextGray] at (0,-2.1) {匀强电场};
  \end{scope}
  \begin{scope}
    \fill[DeepBlue] (0,0) circle (4pt) node[below right] {$+Q$};
    \foreach \r in {0.55,1.05,1.55}{\draw[DeepBlue,thick] (0,0) circle (\r);}
    \foreach \a in {0,45,...,315}{\draw[->,TextGray] ({0.35*cos(\a)},{0.35*sin(\a)})--({2.05*cos(\a)},{2.05*sin(\a)});}
    \node[TextGray] at (0,-2.1) {孤立正点电荷};
  \end{scope}
  \begin{scope}[xshift=6.0cm]
    \fill[DeepBlue] (-1.15,0) circle (4pt) node[below] {$+q$};
    \fill[Ink] (1.15,0) circle (4pt) node[below] {$-q$};
    \foreach \h in {0.35,0.75,1.15}{
      \draw[->,TextGray] (-0.95,{0.55*\h}) .. controls (-0.25,{1.45*\h}) and (0.25,{1.45*\h}) .. (0.95,{0.55*\h});
      \draw[->,TextGray] (-0.95,{-0.55*\h}) .. controls (-0.25,{-1.45*\h}) and (0.25,{-1.45*\h}) .. (0.95,{-0.55*\h});
    }
    \draw[DeepBlue,thick,dashed] (0,-1.85)--(0,1.85) node[above] {$\varphi=0$};
    \node[TextGray] at (0,-2.1) {等量异号点电荷};
  \end{scope}
\end{tikzpicture}

\figcap{三类典型电场的等势线（蓝）与电场线（灰）；两者处处正交}\label{appE:fig:s19-legacy-01}
\end{center}}

\newcommand{\dipoleEquipotentialDiagram}{%
\begin{center}
\begin{tikzpicture}[scale=0.82,>=Latex]
  \fill[DeepBlue] (-1.35,0) circle (4pt) node[below] {$+q$};
  \fill[Ink] (1.35,0) circle (4pt) node[below] {$-q$};
  \foreach \h in {0.35,0.7,1.05,1.4}{
    \draw[DeepBlue,thick] (-1.12,{0.5*\h}) .. controls (-0.35,{1.45*\h}) and (0.35,{1.45*\h}) .. (1.12,{0.5*\h});
    \draw[DeepBlue,thick] (-1.12,{-0.5*\h}) .. controls (-0.35,{-1.45*\h}) and (0.35,{-1.45*\h}) .. (1.12,{-0.5*\h});
  }
  \draw[TextGray,thick,dashed] (0,-2.4)--(0,2.4) node[above] {$\varphi=0$};
  \draw[TextGray] (-1.35,0) circle (0.55);
  \draw[TextGray] (1.35,0) circle (0.55);
\end{tikzpicture}

\figcap{等量异号点电荷的电场线与零等势面}\label{appE:fig:s19-legacy-02}
\end{center}}

\newcommand{\ionizationEnergyDiagram}{%
\begin{center}
\begin{tikzpicture}[x=0.54cm,y=0.12cm,>=Latex]
  \draw[axis] (0,0)--(19.4,0) node[right] {原子序数 $Z$};
  \draw[axis] (0,0)--(0,28) node[above] {第一电离能/eV};
  \foreach \y in {5,10,15,20,25}{
    \draw[TextGray,thin] (0,\y)--(18.6,\y);
    \node[TextGray,left] at (0,\y) {\y};
  }
  \draw[DeepBlue,thick] plot coordinates {
    (1,13.60) (2,24.59) (3,5.39) (4,9.32) (5,8.30) (6,11.26)
    (7,14.53) (8,13.62) (9,17.42) (10,21.56) (11,5.14) (12,7.65)
    (13,5.99) (14,8.15) (15,10.49) (16,10.36) (17,12.97) (18,15.76)};
  \foreach \x/\y in {1/13.60,2/24.59,3/5.39,4/9.32,5/8.30,6/11.26,7/14.53,8/13.62,9/17.42,10/21.56,11/5.14,12/7.65,13/5.99,14/8.15,15/10.49,16/10.36,17/12.97,18/15.76}{
    \fill[DeepBlue] (\x,\y) circle (2pt);
  }
  \foreach \x/\lab in {1/H,2/He,3/Li,6/C,8/O,10/Ne,11/Na,14/Si,17/Cl,18/Ar}{
    \node[TextGray,below,rotate=55,anchor=east] at (\x,-0.3) {\lab};
  }
  \node[TextGray,align=left,anchor=west] at (19.9,16)
    {稀有气体形成峰值\\碱金属形成低谷\\多数原子处于 $5$--$25\,\mathrm{eV}$};
\end{tikzpicture}

\figcap{前18号元素的第一电离能；折线连接仅用于显示周期性趋势}\label{appE:fig:s19-legacy-03}
\end{center}}

\newcommand{\sphereEnergyDiagram}{%
\begin{center}
\begin{tikzpicture}[scale=0.88,>=Latex]
  \coordinate (O) at (0,0);
  \draw[curve] (O) circle (1.55);
  \foreach \a in {0,30,...,330}{
    \fill[DeepBlue] ({1.55*cos(\a)},{1.55*sin(\a)}) circle (1.2pt);
  }
  \draw[vector] (O)--(1.1,1.1) node[midway,above left] {$R$};
  \draw[vector] (1.55,0)--++(1.4,0) node[above] {$E(R)$};
  \draw[vector] (1.05,1.14)--++(0.95,1.02) node[right] {$p$};
  \node[TextGray,align=left,anchor=west] at (4.1,0.35)
    {$Q$ 均匀分布在球面\\表面电场 $E(R)=Q/(4\pi\varepsilon_0R^2)$\\静电压强 $p=\varepsilon_0E^2/2$};
\end{tikzpicture}

\figcap{均匀带电球壳的电场能量与静电压强}\label{appE:fig:s19-legacy-04}
\end{center}}

\newcommand{\zeroPotentialSphereDiagram}{%
\begin{center}
\begin{tikzpicture}[scale=0.85,>=Latex]
  \draw[axis] (-4.6,0)--(5.25,0) node[right] {$x$};
  \fill[Ink] (0,0) circle (4pt) node[below] {$-q$};
  \fill[DeepBlue] (4.2,0) circle (4pt) node[below] {$+Q$};
  % Draw the exact Q=2q member of the Apollonius family: centre=-a/3,
  % radius=2a/3.  It therefore encloses the smaller charge at the origin.
  \draw[DeepBlue,thick] (-1.4,0) circle (2.8);
  \node[TextGray] at (-1.4,3.25) {$\varphi=0$（示意取 $Q=2q$）};
  \draw[<->] (0,-0.8)--node[fill=white,inner sep=1pt] {$a$}(4.2,-0.8);
  \node[TextGray] at (-1.4,-3.25)
    {当 $Q>q$ 时，零等势面包围较小的负电荷};
\end{tikzpicture}

\figcap{不等量异号点电荷的零等势面是 Apollonius 球面}\label{appE:fig:s19-legacy-05}
\end{center}}

\newcommand{\twoLineDiagram}{%
\begin{center}
\begin{tikzpicture}[scale=0.9,>=Latex]
  \draw[axis] (-4.6,0)--(4.6,0) node[right] {$x$};
  \draw[axis] (0,-2.7)--(0,2.7) node[above] {$y$};
  \fill[DeepBlue] (1.8,0) circle (4pt) node[below] {$+\lambda$};
  \fill[Ink] (-1.8,0) circle (4pt) node[below] {$-\lambda$};
  \draw[TextGray,dashed] (0,-2.4)--(0,2.4) node[above right] {$\varphi=0$};
  % Equipotentials: Apollonius circles surrounding one line charge.
  % K=3 and K=5 members.  Every circle surrounds the corresponding
  % line charge at x=+/-a rather than sitting between the charges.
  \draw[DeepBlue,thick] (2.25,0) circle (1.35);
  \draw[DeepBlue,thick] (1.95,0) circle (0.75);
  \draw[DeepBlue,thick] (-2.25,0) circle (1.35);
  \draw[DeepBlue,thick] (-1.95,0) circle (0.75);
  % Field lines: circles through both line charges (only representative arcs).
  \draw[TextGray,thick,dashed] (0,1.65) circle (2.44);
  \draw[TextGray,thick,dashed] (0,-1.65) circle (2.44);
  \node[TextGray] at (0,-3.2) {实线是等势线；虚线是与它们正交、并通过两导线的电场线};
\end{tikzpicture}

\figcap{两根异号无限长带电直线的横截面}\label{appE:fig:s19-legacy-06}
\end{center}}

\newcommand{\finiteLineDiagramThirty}{%
\begin{center}
\begin{tikzpicture}[scale=0.92,>=Latex]
  \draw[axis] (-4.2,0)--(4.2,0) node[right] {$x$};
  \draw[axis] (0,-2.4)--(0,2.4) node[above] {$\rho$};
  \draw[DeepBlue,very thick] (-1.7,0)--(1.7,0);
  \fill (0,0) circle (1.5pt) node[below] {$O$};
  \node[below] at (-1.7,0) {$-a$}; \node[below] at (1.7,0) {$a$};
  \coordinate (P) at (2.7,1.6);
  \fill[DeepBlue] (P) circle (2.5pt) node[right] {$P(x,\rho)$};
  \draw[guide] (P)--(-1.7,0) (P)--(1.7,0);
  \node[TextGray] at (0,-2.95) {$r_++r_-=\text{常量}$ 为旋转椭球；$r_--r_+=\text{常量}$ 为旋转双曲面};
\end{tikzpicture}

\figcap{有限长均匀带电直线的轴对称几何}\label{appE:fig:s19-legacy-07}
\end{center}}

\newcommand{\repulsionDiagram}{%
\begin{center}
\begin{tikzpicture}[scale=0.95,>=Latex]
  \draw[thick] (-1.8,3.5)--(1.8,3.5);
  \foreach \x in {-1.6,-1.2,...,1.6}{\draw[TextGray] (\x,3.5)--++(-0.18,0.18);}
  \coordinate (O) at (0,3.5); \coordinate (A) at (-1.55,1.05); \coordinate (B) at (0,0.42);
  \draw[curve] (O)--(A);
  \fill[DeepBlue] (A) circle (4pt) node[below left=3pt,fill=white,inner sep=0.7pt] {$A(+q)$};
  \fill[DeepBlue] (B) circle (4pt) node[right=4pt,fill=white,inner sep=0.7pt] {$B(+q)$};
  \draw[guide] (O)--(B);
  \draw[<->] (-2.25,0.42)--node[left] {$h$}(-2.25,1.05);
  \draw[vector] (A)--++(0,-1.15) node[below right=2pt,fill=white,inner sep=0.7pt] {$m\vect g$};
  \draw[vector] (A)--++(0.75,1.18) node[above right=2pt,fill=white,inner sep=0.7pt] {$\vect T$};
  \draw[vector] (A)--++(-1.25,0.51) node[above left=2pt,fill=white,inner sep=0.7pt] {$\vect F_e$};
  \node[TextGray,align=left,anchor=west] at (3.0,1.55)
    {$B$ 最终位于 $A$ 原来的最低点\\慢速移动时可用能量关系};
\end{tikzpicture}

\figcap{同号小球排斥过程的始末几何}\label{appE:fig:s19-legacy-08}
\end{center}}

\newcommand{\bohrDiagram}{%
\begin{center}
  \begin{tikzpicture}[scale=0.9,>=Latex]
  \fill[DeepBlue] (0,0) circle (5pt) node[below left=4pt,fill=white,inner sep=0.7pt] {$p(+e)$};
  \draw[DeepBlue,thick] (0,0) circle (2.05);
  \fill[Ink] (2.05,0) circle (4pt) node[below right=4pt,fill=white,inner sep=0.7pt] {$e^-(-e)$};
  \draw[vector] (2.05,0)--++(0,1.45) node[above right=2pt,fill=white,inner sep=0.7pt] {$\vect v$};
  \draw[vector] (2.05,0)--++(-1.25,0) node[pos=0.55,above=4pt,fill=white,inner sep=0.7pt] {$\vect F_e$};
  \draw[<->] (0,-1.05)--node[fill=white,inner sep=1pt] {$r$}(2.05,-1.05);
\end{tikzpicture}

\figcap{玻尔基态模型中的电子圆周运动}\label{appE:fig:s19-legacy-09}
\end{center}}

\newcommand{\rutherfordDiagram}{%
\begin{center}
\begin{tikzpicture}[scale=0.9,>=Latex]
  \fill[DeepBlue] (0,0) circle (5pt) node[below right] {$Ze$};
  \draw[axis] (-4.8,0)--(4.8,0) node[right] {$x$};
  \draw[guide] (-4.2,1.25)--(0,1.25);
  \draw[DeepBlue,very thick,-{Latex[length=2.7mm]}]
    (-4.2,1.25) .. controls (-1.5,1.25) and (-0.8,2.2) .. (3.4,3.0);
  \draw[<->] (-3.4,0)--node[left] {$b$}(-3.4,1.25);
  \draw[vector] (-4.0,1.25)--++(1.3,0) node[above] {$\vect v_\infty$};
  \draw[guide] (1.9,2.55)--++(1.7,0);
  \coordinate (A) at (1.9,2.55); \coordinate (B) at (3.6,2.55); \coordinate (C) at (3.4,3.0);
  \draw pic[draw,"$\theta$",angle radius=9mm,angle eccentricity=1.25]
    {angle=B--A--C};
\end{tikzpicture}

\figcap{库仑斥力下的双曲线散射轨迹}\label{appE:fig:s19-legacy-10}
\end{center}}


\begin{origexamplebox}{【电势算电场】}\label{appE:ex:s19-30-1}
利用“先求电势，再求电场”的方法，重新计算上一讲中的典型带电系统：点电荷系、有限带电体
和无限长或无限大的电荷分布。分析哪些情况下不能取无穷远为零电势，并给出处理方法。

\QAENamed{\QAESources}{有限带电体与无限分布的比较；先确定几何和有限参考点，再求电势差。}{quality:appE:q5}
\end{origexamplebox}

\begin{xesolution}
方法的核心是
\[
  \varphi(\vect r)=k\int\frac{\dd q}{|\vect r-\vect r'|},
  \qquad
  \vect E=-\nabla\varphi,
  \qquad k=\frac1{4\pi\varepsilon_0}.
\]

\textbf{点电荷系。}直接标量叠加
\[
  \varphi=k\sum_i\frac{Q_i}{r_i},
\]
然后求梯度即得
\[
  \vect E=k\sum_i\frac{Q_i}{r_i^3}\vect r_i.
\]
这种方法对偶极子、多极子尤其方便，因为远场展开先在标量电势上进行，再求导即可。

\textbf{均匀带电圆盘。}半径 $R$、面密度 $\sigma$ 的圆盘，在轴线上距圆心 $z>0$ 处的电势为
\[
  \varphi(z)=k\int_0^R\frac{2\pi\sigma s\dd s}{\sqrt{z^2+s^2}}
  =\frac{\sigma}{2\varepsilon_0}
  \bigl(\sqrt{z^2+R^2}-z\bigr).
\]
因此
\[
  \boxed{E_z=-\frac{\dd\varphi}{\dd z}
  =\frac{\sigma}{2\varepsilon_0}
  \left(1-\frac{z}{\sqrt{z^2+R^2}}\right)}.
\]

\textbf{无限长直线。}若线电荷密度为 $\lambda$，$E_r=\lambda/(2\pi\varepsilon_0r)$。
从有限参考半径 $r_0$ 积分得
\[
  \boxed{\varphi(r)-\varphi(r_0)
  =-\frac{\lambda}{2\pi\varepsilon_0}\ln\frac r{r_0}}.
\]
因为 $r\to\infty$ 时对数发散，所以不能令无穷远处电势为零。

\textbf{无限大平面。}若面电荷密度为 $\sigma$，取平面为 $z=0$，则
\[
  E_z=\frac{\sigma}{2\varepsilon_0}\operatorname{sgn}(z),
\]
取 $\varphi(0)=0$ 后
\[
  \boxed{\varphi(z)=-\frac{\sigma}{2\varepsilon_0}|z|}.
\]
它同样不能用无穷远作为零电势。对非局域电荷分布，正确做法是选定有限参考点，始终只计算
电势差；若多个无限分布组合后总场在远处衰减足够快，也可先把组合视为整体再选零点。

\QAENamed{\QAEPotentialPlot}{正电荷密度时的电势形状；两图的纵向正尺度分别为 $\lambda/(2\pi\varepsilon_0)$、$\sigma/(2\varepsilon_0)$。}{quality:appE:a5}

图~\ref{quality:appE:a5}中对数和绝对值曲线在远处均不趋于有限常数，所以不能给它们硬加“无穷远零电势”。求导会消去有限参考点带来的常数，电场不随零点选择而变。平面 $z=0$ 处电势连续、法向导数跳变，与面电荷边界相符。
\end{xesolution}

\endgroup

\Needspace{18\baselineskip}

\section{电介质与电容器}

\begingroup

\graphicspath{{../../../sample/xes-original/twenty-first-section/}{../../../sample/xes-original/twenty-first-section/assets/}}

% Restrained TikZ redraws for Lesson 32.

\newcommand{\arbitraryCapacitorDiagram}{%
\begin{center}
\begin{tikzpicture}[x=1cm,y=1cm,>=Latex]
  \fill[PaleBlue] plot[smooth cycle,tension=0.78] coordinates
    {(-3.7,0.1) (-3.3,1.35) (-2.05,1.55) (-1.45,0.55) (-1.8,-0.75) (-3.05,-1.25) (-3.95,-0.55)};
  \draw[DeepBlue,very thick] plot[smooth cycle,tension=0.78] coordinates
    {(-3.7,0.1) (-3.3,1.35) (-2.05,1.55) (-1.45,0.55) (-1.8,-0.75) (-3.05,-1.25) (-3.95,-0.55)};
  \fill[SoftGray] plot[smooth cycle,tension=0.78] coordinates
    {(3.7,0.15) (3.15,1.45) (1.85,1.25) (1.4,0.15) (1.95,-1.05) (3.35,-1.25) (4.0,-0.45)};
  \draw[Ink,very thick] plot[smooth cycle,tension=0.78] coordinates
    {(3.7,0.15) (3.15,1.45) (1.85,1.25) (1.4,0.15) (1.95,-1.05) (3.35,-1.25) (4.0,-0.45)};
  \foreach \p in {(-3.35,0.65),(-2.65,1.05),(-2.15,0.15),(-3.1,-0.55)}
    \node[DeepBlue] at \p {$+$};
  \foreach \p in {(3.3,0.7),(2.45,0.95),(1.95,0.05),(2.85,-0.55)}
    \node[Ink] at \p {$-$};
  \draw[component] (-1.65,0.82) .. controls (-0.5,1.5) and (0.5,1.5) .. (1.65,0.82);
  \draw[component] (-1.45,0.15) .. controls (-0.45,0.65) and (0.45,0.65) .. (1.45,0.15);
  \draw[component] (-1.62,-0.55) .. controls (-0.5,-1.15) and (0.5,-1.15) .. (1.72,-0.55);
  \node[DeepBlue] at (-2.7,-1.65) {$+Q,\ \varphi_+$};
  \node[Ink] at (2.7,-1.65) {$-Q,\ \varphi_-$};
  \draw[<->,TextGray] (-1.05,-1.9)--node[below] {$U=\varphi_+-\varphi_-$}(1.05,-1.9);
\end{tikzpicture}

\figcap{任意形状的两块导体组成双极电容器}\label{appE:fig:s21-legacy-01}
\end{center}}

\newcommand{\capacitorGeometryDiagram}{%
\begin{center}
\begin{tikzpicture}[scale=0.82,>=Latex]
  % parallel plates
  \draw[DeepBlue,very thick] (-5.2,0.8)--(-2.6,0.8);
  \draw[DeepBlue,very thick] (-5.2,-0.8)--(-2.6,-0.8);
  \foreach \x in {-4.9,-4.4,...,-2.9}{
    \node[DeepBlue] at (\x,0.48) {$+$}; \node[Ink] at (\x,-0.48) {$-$};
    \draw[component] (\x,0.32)--(\x,-0.32);
  }
  \draw[<->] (-5.55,-0.8)--node[left] {$d$}(-5.55,0.8);
  \node[TextGray] at (-3.9,-1.45) {平行板};
  % spherical
  \draw[curve] (0,0) circle (0.73); \draw[curve] (0,0) circle (1.28);
  \draw[<->] (0,0)--node[above] {$R_1$}(0.73,0);
  \draw[<->] (0,-0.28)--node[below] {$R_2$}(1.28,-0.28);
  \node[TextGray] at (0,-1.45) {同心球};
  % coaxial
  \draw[curve] (3.6,0) circle (0.62); \draw[curve] (3.6,0) circle (1.20);
  \foreach \a in {0,45,...,315}{\draw[component] ({3.6+0.66*cos(\a)},{0.66*sin(\a)})--({3.6+1.12*cos(\a)},{1.12*sin(\a)});}
  \node at (3.6,0) {$a$}; \node[above right] at (4.35,0.75) {$b$};
  \node[TextGray] at (3.6,-1.45) {同轴圆柱截面};
\end{tikzpicture}

\figcap{平行板、同心球与同轴圆柱电容器}\label{appE:fig:s21-legacy-02}
\end{center}}

\newcommand{\seriesParallelDiagram}{%
\begin{center}
\begin{tikzpicture}[scale=0.86,>=Latex]
  % parallel
  \draw[thick] (-5.0,1.25)--(-4.2,1.25) (-5.0,-1.25)--(-4.2,-1.25);
  \draw[thick] (-4.2,1.25)--(-4.2,-1.25) (-1.5,1.25)--(-1.5,-1.25);
  \draw[thick] (-4.2,0.72)--(-3.3,0.72) (-2.4,0.72)--(-1.5,0.72);
  \draw[very thick,DeepBlue] (-3.3,0.38)--(-3.3,1.06) (-2.4,0.38)--(-2.4,1.06);
  \draw[thick] (-4.2,-0.72)--(-3.3,-0.72) (-2.4,-0.72)--(-1.5,-0.72);
  \draw[very thick,DeepBlue] (-3.3,-1.06)--(-3.3,-0.38) (-2.4,-1.06)--(-2.4,-0.38);
  \node at (-2.85,1.28) {$C_1$}; \node at (-2.85,-1.28) {$C_2$};
  \node[TextGray] at (-3.25,-1.85) {并联：电压相同};
  % series
  \draw[thick] (0.8,0)--(1.8,0) (2.55,0)--(3.45,0) (4.2,0)--(5.2,0);
  \draw[very thick,DeepBlue] (1.8,-0.55)--(1.8,0.55) (2.55,-0.55)--(2.55,0.55);
  \draw[very thick,DeepBlue] (3.45,-0.55)--(3.45,0.55) (4.2,-0.55)--(4.2,0.55);
  \node at (2.18,0.95) {$C_1$}; \node at (3.82,0.95) {$C_2$};
  \node[TextGray] at (3.0,-1.85) {串联：中间孤立导体净电荷不变};
\end{tikzpicture}

\figcap{电容器的并联与串联}\label{appE:fig:s21-legacy-03}
\end{center}}

\newcommand{\polarizationDiagram}{%
\begin{center}
\begin{tikzpicture}[scale=0.87,>=Latex]
  \fill[PaleBlue] (-3.3,-1.5) rectangle (3.3,1.5);
  \draw[TextGray,thick] (-3.3,-1.5) rectangle (3.3,1.5);
  \foreach \x in {-2.6,-1.3,0,1.3,2.6}{
    \foreach \y in {-0.85,0,0.85}{
      \node[Ink] at (\x-0.16,\y) {$-$}; \node[DeepBlue] at (\x+0.16,\y) {$+$};
    }
  }
  \foreach \y in {-1.05,-0.35,0.35,1.05}{
    \node[Ink] at (-3.55,\y) {$-$}; \node[DeepBlue] at (3.55,\y) {$+$};
  }
  \draw[vector] (-1.0,2.05)--(1.0,2.05) node[midway,above] {$\vect E$};
  \draw[vector] (-1.0,-2.05)--(1.0,-2.05) node[midway,below] {$\vect P$};
  \node[TextGray,align=left,anchor=west] at (4.25,0.45)
    {分子偶极矩沿外场排列\\两端出现束缚面电荷};
\end{tikzpicture}

\figcap{外场作用下电介质的极化与束缚面电荷}\label{appE:fig:s21-legacy-04}
\end{center}}

\newcommand{\reciprocitySphereDiagram}{%
\begin{center}
\begin{tikzpicture}[x=1cm,y=1cm,>=Latex]
  \draw[curve] (-2.5,0) circle (1.18);
  \fill (-2.5,0) circle (1.8pt) node[below=4pt] {$O$};
  \fill[DeepBlue] (2.5,0) circle (2.6pt) node[above=5pt] {$P$};
  \draw[<->] (-2.5,0.25)--node[above] {$r$}(2.5,0.25);
  \node[align=center] at (-2.5,-1.65) {状态 I：球带单位电荷\\状态 II：$P$ 点放电荷 $q$};
  \node[align=left,anchor=west,TextGray] at (4.0,0.45)
    {$\displaystyle \sum_i Q_i\varphi_i'=\sum_iQ_i'\varphi_i$\\[4pt]
     交换“源”与“测量位置”，\\倒易关系保持不变};
\end{tikzpicture}

\figcap{格林倒易定理：导体球与球外一点的两种静电状态}\label{appE:fig:s21-legacy-05}
\end{center}}

\newcommand{\reciprocityConductorsDiagram}{%
\begin{center}
\begin{tikzpicture}[x=1cm,y=1cm,>=Latex]
  \fill[PaleBlue] plot[smooth cycle,tension=0.8] coordinates
    {(-4.1,0.2) (-3.5,1.25) (-2.15,1.05) (-1.65,0) (-2.3,-1.1) (-3.7,-0.95)};
  \draw[DeepBlue,thick] plot[smooth cycle,tension=0.8] coordinates
    {(-4.1,0.2) (-3.5,1.25) (-2.15,1.05) (-1.65,0) (-2.3,-1.1) (-3.7,-0.95)};
  \fill[SoftGray] plot[smooth cycle,tension=0.8] coordinates
    {(4.05,0.15) (3.35,1.2) (2.0,0.9) (1.55,-0.15) (2.35,-1.2) (3.75,-0.85)};
  \draw[Ink,thick] plot[smooth cycle,tension=0.8] coordinates
    {(4.05,0.15) (3.35,1.2) (2.0,0.9) (1.55,-0.15) (2.35,-1.2) (3.75,-0.85)};
  \node[align=center] at (-2.85,0) {$Q_1$\\$\varphi_1$};
  \node[align=center] at (2.85,0) {$Q_2$\\$\varphi_2$};
  \draw[<->,TextGray] (-1.15,0)--node[above] {静电感应}(1.15,0);
  \node[align=center,TextGray] at (0,-1.7)
    {第二状态用 $Q_1',Q_2',\varphi_1',\varphi_2'$ 表示；导体几何保持不变};
\end{tikzpicture}

\figcap{格林倒易定理：两个任意导体的两种静电平衡状态}\label{appE:fig:s21-legacy-06}
\end{center}}

\newcommand{\dielectricExperimentDiagram}{%
\begin{center}
\begin{tikzpicture}[x=0.86cm,y=0.86cm,>=Latex]
  \foreach \shift/\lab/\needle in {-4.3/{未插入介质：$C_0,U_0$}/32,4.3/{插入介质：$C=\varepsilon_rC_0$}/-8}{
    \begin{scope}[shift={(\shift,0)}]
      \draw[Ink,very thick] (-1.15,1.0)--(-1.15,3.1);
      \draw[Ink,very thick] (1.15,1.0)--(1.15,3.1);
      \node[DeepBlue] at (-1.45,2.1) {$+$}; \node[Ink] at (1.45,2.1) {$-$};
      \ifdim\shift pt>0pt
        \fill[DeepBlue!18] (-0.95,1.15) rectangle (0.95,2.95);
        \draw[DeepBlue,thick] (-0.95,1.15) rectangle (0.95,2.95);
        \node[DeepBlue,rotate=90] at (0,2.05) {电介质};
      \fi
      \draw[Ink,thick] (-1.15,1.0)--(-1.15,0.45)--(-0.75,0.45);
      \draw[Ink,thick] (1.15,1.0)--(1.15,0.45)--(0.75,0.45);
      \draw[rounded corners=3pt,fill=SoftGray,draw=Ink,thick] (-0.75,-1.5) rectangle (0.75,0.55);
      \draw[Ink] (0,-0.25) circle (0.48);
      \draw[DeepBlue,very thick] (0,-0.25)--({0.38*cos(\needle)},{-0.25+0.38*sin(\needle)});
      \foreach \a in {20,50,...,160}{\draw[TextGray] ({0.38*cos(\a)},{-0.25+0.38*sin(\a)})--({0.45*cos(\a)},{-0.25+0.45*sin(\a)});}
      \node[align=center] at (0,-2.05) {\lab};
    \end{scope}
  }
  \draw[vector] (-1.15,0.2)--(1.15,0.2) node[midway,above] {插入介质};
\end{tikzpicture}

\figcap{保持电荷不变时，插入电介质使电势差减小、电容增大}\label{appE:fig:s21-legacy-07}
\end{center}}

\newcommand{\touchingSpheresDiagram}{%
\begin{center}
\begin{tikzpicture}[scale=0.92,>=Latex]
  \draw[curve] (-1.3,0) circle (1.3); \draw[curve] (1.3,0) circle (1.3);
  \fill (0,0) circle (1.4pt) node[below] {接触点};
  \draw[<->] (1.3,0)--node[above] {$R$}(2.6,0);
  \foreach \x/\r in {-1.75/2.3,-0.95/1.8,-0.48/1.25,0.48/1.25,0.95/1.8,1.75/2.3}{
    \fill[DeepBlue] (\x,0) circle (\r pt);
  }
  \node[TextGray] at (0,-2.05) {接触极限可由双球坐标或无穷像电荷序列求得};
\end{tikzpicture}

\figcap{两个等半径相切球连成一个等势导体}\label{appE:fig:s21-legacy-08}
\end{center}}

\newcommand{\twoCylinderDiagram}{%
\begin{center}
\begin{tikzpicture}[scale=0.9,>=Latex]
  \draw[curve] (-2.2,0) circle (0.85); \draw[curve] (2.2,0) circle (0.85);
  \node at (-2.2,0) {$+\lambda$}; \node at (2.2,0) {$-\lambda$};
  \fill ( -2.2,0) circle (1.3pt); \fill (2.2,0) circle (1.3pt);
  \draw[<->] (-2.2,-1.28)--node[fill=white,inner sep=1pt] {$L$}(2.2,-1.28);
  \draw[<->] (-2.2,0)--node[above] {$R$}(-1.35,0);
  \draw[DeepBlue,thick] (-1.55,0.8) .. controls (0,2.0) .. (1.55,0.8);
  \draw[DeepBlue,thick] (-1.55,-0.8) .. controls (0,-2.0) .. (1.55,-0.8);
  \draw[DeepBlue,thick] (-1.35,0.25) .. controls (0,0.78) .. (1.35,0.25);
  \draw[DeepBlue,thick] (-1.35,-0.25) .. controls (0,-0.78) .. (1.35,-0.25);
  \node[TextGray] at (0,-2.35) {$\cosh\tau_0=L/(2R)$，两圆柱面对应 $\tau=\pm\tau_0$};
\end{tikzpicture}

\figcap{两根等半径平行圆柱导体的横截面}\label{appE:fig:s21-legacy-09}
\end{center}}

\newcommand{\splitDielectricDiagram}{%
\begin{center}
\begin{tikzpicture}[scale=0.83,>=Latex]
  % side-by-side
  \draw[DeepBlue,very thick] (-5.1,1.0)--(-1.9,1.0) (-5.1,-1.0)--(-1.9,-1.0);
  \fill[PaleBlue] (-5.1,-0.95) rectangle (-3.5,0.95);
  \draw[guide] (-3.5,-1.0)--(-3.5,1.0);
  \node at (-4.3,0) {$\varepsilon$}; \node at (-2.7,0) {$\varepsilon_0$};
  \node[TextGray] at (-3.5,-1.55) {按面积分割：并联};
  % layered
  \draw[DeepBlue,very thick] (1.2,1.0)--(4.4,1.0) (1.2,-1.0)--(4.4,-1.0);
  \fill[PaleBlue] (1.2,0) rectangle (4.4,0.95);
  \draw[guide] (1.2,0)--(4.4,0);
  \node at (2.8,0.5) {$\varepsilon$}; \node at (2.8,-0.5) {$\varepsilon_0$};
  \node[TextGray] at (2.8,-1.55) {按厚度分割：串联};
\end{tikzpicture}

\figcap{部分填充介质的两种等效方式}\label{appE:fig:s21-legacy-10}
\end{center}}

\newcommand{\dielectricPullDiagram}{%
\begin{center}
\begin{tikzpicture}[scale=0.88,>=Latex]
  \draw[DeepBlue,very thick] (-3.8,0.8)--(3.8,0.8) (-3.8,-0.8)--(3.8,-0.8);
  \fill[PaleBlue] (-3.8,-0.72) rectangle (0.6,0.72);
  \draw[TextGray,thick] (-3.8,-0.72) rectangle (0.6,0.72);
  \draw[<->] (-3.8,-1.22)--node[fill=white,inner sep=1pt] {$x$}(0.6,-1.22);
  \draw[<->] (0.6,-1.22)--node[fill=white,inner sep=1pt] {$a-x$}(3.8,-1.22);
  \draw[<->] (4.15,-0.8)--node[right] {$d$}(4.15,0.8);
  \draw[vector] (0.55,0)--(2.2,0) node[midway,above] {$\vect F$};
  \node[TextGray] at (0,-1.95) {介质进入越深，等效电容越大，电场把介质继续拉入};
\end{tikzpicture}

\figcap{部分插入平行板电容器的介质板}\label{appE:fig:s21-legacy-11}
\end{center}}

\newcommand{\dielectricSphereDiagram}{%
\begin{center}
\begin{tikzpicture}[scale=0.88,>=Latex]
  \fill[PaleBlue] (0,0) circle (1.35); \draw[curve] (0,0) circle (1.35);
  \foreach \a in {-60,-30,0,30,60}{
    \node[Ink] at ({-1.08*cos(\a)},{1.08*sin(\a)}) {$-$};
    \node[DeepBlue] at ({1.08*cos(\a)},{1.08*sin(\a)}) {$+$};
  }
  \draw[vector] (-4.3,1.95)--(-2.8,1.95) node[midway,above] {$\vect E_0$};
  \draw[DeepBlue,thick,-{Latex[length=2mm]}] (-4.2,1.72)--(4.2,1.72);
  \draw[DeepBlue,thick,-{Latex[length=2mm]}] (-4.2,-1.72)--(4.2,-1.72);
  \draw[DeepBlue,thick] (-4.2,0.92) .. controls (-2.6,0.92) and (-1.8,0.62) .. (-1.26,0.48);
  \draw[DeepBlue,thick,-{Latex[length=2mm]}] (-1.26,0.48)--(1.26,0.48);
  \draw[DeepBlue,thick,-{Latex[length=2mm]}] (1.26,0.48) .. controls (1.8,0.62) and (2.6,0.92) .. (4.2,0.92);
  \draw[DeepBlue,thick] (-4.2,-0.92) .. controls (-2.6,-0.92) and (-1.8,-0.62) .. (-1.26,-0.48);
  \draw[DeepBlue,thick,-{Latex[length=2mm]}] (-1.26,-0.48)--(1.26,-0.48);
  \draw[DeepBlue,thick,-{Latex[length=2mm]}] (1.26,-0.48) .. controls (1.8,-0.62) and (2.6,-0.92) .. (4.2,-0.92);
  \draw[DeepBlue,thick] (-4.2,0)--(-1.35,0);
  \draw[DeepBlue,thick,-{Latex[length=2mm]}] (-1.35,0)--(1.35,0);
  \draw[DeepBlue,thick,-{Latex[length=2mm]}] (1.35,0)--(4.2,0);
  \node at (0,-0.82) {$\varepsilon$};
  \node[TextGray] at (0,-2.35) {$\vect E_{\rm in}=\dfrac{3\varepsilon_0}{\varepsilon+2\varepsilon_0}\vect E_0$};
\end{tikzpicture}

\figcap{匀强电场中的线性电介质球}\label{appE:fig:s21-legacy-12}
\end{center}}

\newcommand{\chargingDiagram}{%
\begin{center}
\begin{tikzpicture}[scale=0.9,>=Latex]
  \draw[thick] (-3.8,0)--(-2.9,0);
  \draw[thick] (-2.9,-0.55)--(-2.9,0.55); \draw[very thick,DeepBlue] (-2.55,-0.35)--(-2.55,0.35);
  \node[below] at (-2.72,-0.62) {$\mathcal E$};
  \draw[thick] (-2.55,0)--(-1.7,0);
  \draw[thick,decorate,decoration={zigzag,segment length=6pt,amplitude=2.5pt}] (-1.7,0)--(0.2,0);
  \node[above] at (-0.75,0.35) {$R$};
  \draw[thick] (0.2,0)--(1.25,0);
  \draw[very thick,DeepBlue] (1.25,-0.55)--(1.25,0.55) (1.85,-0.55)--(1.85,0.55);
  \node[above] at (1.55,0.75) {$C$};
  \draw[thick] (1.85,0)--(2.8,0)--(2.8,-1.35)--(-3.8,-1.35)--(-3.8,0);
  \node[TextGray,align=left,anchor=west] at (4.0,-0.1)
    {电源功 $C\mathcal E^2$\\电场能 $C\mathcal E^2/2$\\焦耳热 $C\mathcal E^2/2$};
\end{tikzpicture}

\figcap{电阻参与的电容器充电过程}\label{appE:fig:s21-legacy-13}
\end{center}}

\newcommand{\starCapacitorDiagram}{%
\begin{center}
\begin{tikzpicture}[scale=0.94,>=Latex]
  \fill[DeepBlue] (0,0) circle (2pt) node[below right] {$\varphi_0$};
  \foreach \a/\c/\p in {90/C_1/\varphi_1,210/C_2/\varphi_2,330/C_3/\varphi_3}{
    \draw[thick] (0,0)--({0.95*cos(\a)},{0.95*sin(\a)});
    \draw[very thick,DeepBlue] ({0.78*cos(\a)-0.35*sin(\a)},{0.78*sin(\a)+0.35*cos(\a)})--({0.78*cos(\a)+0.35*sin(\a)},{0.78*sin(\a)-0.35*cos(\a)});
    \draw[very thick,DeepBlue] ({1.12*cos(\a)-0.35*sin(\a)},{1.12*sin(\a)+0.35*cos(\a)})--({1.12*cos(\a)+0.35*sin(\a)},{1.12*sin(\a)-0.35*cos(\a)});
    \draw[thick] ({1.12*cos(\a)},{1.12*sin(\a)})--({2.2*cos(\a)},{2.2*sin(\a)});
    \node at ({1.48*cos(\a)+0.45*sin(\a)},{1.48*sin(\a)-0.45*cos(\a)}) {$\c$};
    \node at ({2.55*cos(\a)},{2.55*sin(\a)}) {$\p$};
  }
  \node[TextGray,align=left,anchor=west] at (4.0,0.45)
    {中心导体孤立且初始中性\\$\sum C_i(\varphi_0-\varphi_i)=0$};
\end{tikzpicture}

\figcap{三电容星形连接的浮置中心节点}\label{appE:fig:s21-legacy-14}
\end{center}}

\newcommand{\atmosphereCapacitorDiagram}{%
\begin{center}
\begin{tikzpicture}[scale=0.83,>=Latex]
  \fill[PaleBlue] (0,0) circle (2.18); \fill[white] (0,0) circle (1.86);
  \fill[SoftGray] (0,0) circle (1.62);
  \draw[curve] (0,0) circle (1.62); \draw[DeepBlue,thick] (0,0) circle (2.18);
  \node at (0,0) {地球};
  \foreach \a in {0,35,...,325}{
    \node[Ink] at ({1.72*cos(\a)},{1.72*sin(\a)}) {$-$};
    \node[DeepBlue] at ({2.05*cos(\a)},{2.05*sin(\a)}) {$+$};
  }
  \draw[<->] (2.35,0)--node[right] {$h$}(2.35,0.56);
  \node[TextGray,align=left,anchor=west] at (4.1,0.7)
    {内半径 $R_\oplus$\\外半径 $R_\oplus+h$\\大气同时提供电容与漏电电阻};
\end{tikzpicture}

\figcap{理想化的大气球形电容器模型}\label{appE:fig:s21-legacy-15}
\end{center}}


\begin{origexamplebox}{【统统算出来】}\label{appE:ex:s21-32-1}
求下列三种电容器的电容：
\begin{enumerate}[label=(\arabic*),leftmargin=2.2em]
  \item 面积 $S=l^2$、间距 $d$ 的平行板电容器，$l\gg d$；
  \item 内、外半径为 $R_1,R_2$ 的同心球形电容器，$R_1<R_2$；
  \item 内、外半径为 $a,b$、长度为 $l$ 的同轴圆柱形电容器，$l\gg b$。
\end{enumerate}

\QAENamed{\QAECapacitors{0}}{平行板、同心球与同轴圆柱三种电容器的尺寸；后两图均为中心横截面。}{quality:appE:q6}
\end{origexamplebox}

\begin{xesolution}
三题都让内、外电极分别带 $+Q$、$-Q$，先由高斯定律求两极间电场，再积分求电压。

\textbf{(1) 平行板。}忽略边缘效应，板间场为
\[
  E=\frac{Q}{\varepsilon_0S},\qquad
  U=Ed=\frac{Qd}{\varepsilon_0S}.
\]
因此
\begin{equation}
  \boxed{C_{\rm pp}=\frac QU=\frac{\varepsilon_0S}{d}}.
\end{equation}

\textbf{(2) 同心球。}在 $R_1<r<R_2$ 内，
\[
  E(r)4\pi r^2=\frac Q{\varepsilon_0},
  \qquad E(r)=\frac{Q}{4\pi\varepsilon_0r^2}.
\]
故
\[
  U=\int_{R_1}^{R_2}E(r)\dd r
  =\frac{Q}{4\pi\varepsilon_0}\left(\frac1{R_1}-\frac1{R_2}\right),
\]
从而
\begin{equation}
  \boxed{C_{\rm sph}=4\pi\varepsilon_0
  \frac{R_1R_2}{R_2-R_1}}.
\end{equation}
令 $R_2\to\infty$，得到孤立球的 $C=4\pi\varepsilon_0R_1$。

\textbf{(3) 同轴圆柱。}柱面高斯面给出
\[
  E(r)2\pi rl=\frac Q{\varepsilon_0},
  \qquad E(r)=\frac{Q}{2\pi\varepsilon_0lr}.
\]
所以
\[
  U=\int_a^bE(r)\dd r=\frac{Q}{2\pi\varepsilon_0l}\ln\frac ba,
\]
以及
\begin{equation}
  \boxed{C_{\rm cyl}=\frac{2\pi\varepsilon_0l}{\ln(b/a)}}.
\end{equation}
三种推导中的近似只用于忽略开口或极板边缘的非均匀场；在理想同心球中不存在边缘近似。

\QAENamed{\QAECapacitors{1}}{积分路径沿两极间的场线；同心球与同轴柱使用不同的高斯面积。}{quality:appE:a6}

图~\ref{quality:appE:a6}内外电极之间分别使用平面、球面和柱面对称性。球面面积随 $r^2$ 增长，柱面面积随 $r$ 增长，因而径向积分一个给倒数差、一个给对数；不能仅凭外形把两式套换。
\end{xesolution}

\begin{origexamplebox}{【改造电容器】}\label{appE:ex:s21-32-4}
平行板电容器的极板面积为 $S$、间距为 $d$。
\begin{enumerate}[label=(\arabic*),leftmargin=2.2em]
  \item 沿面积方向，在左半边完全填入相对介电常数为 $\varepsilon_r>1$ 的介质；
  \item 沿厚度方向，在上半层填入相对介电常数为 $\varepsilon_r$ 的介质。
\end{enumerate}
分别求改造后的电容，并判断增减。

忽略边缘效应对总电容的修正，介质为均匀线性电介质。

\QAENamed{\QAESplitQuestion}{半面积填充与半厚度填充；蓝色表示介质，白色表示真空。}{quality:appE:q7}
\end{origexamplebox}

\begin{xesolution}
\textbf{(1) 按面积分成左右两区。}两区具有相同板间距与电压，相当于并联：
\[
  C_1=\frac{\varepsilon_r\varepsilon_0(S/2)}d,
  \qquad C_2=\frac{\varepsilon_0(S/2)}d.
\]
故
\begin{equation}
  \boxed{C_{\rm left}=C_1+C_2
  =\frac{\varepsilon_0S}{2d}(\varepsilon_r+1)}.
\end{equation}

\textbf{(2) 按厚度分成上下两层。}两层面积相同，自由面电荷密度相同，电位移的法向分量连续，
相当于串联：
\[
  C_1=\frac{\varepsilon_r\varepsilon_0S}{d/2},
  \qquad C_2=\frac{\varepsilon_0S}{d/2}.
\]
于是
\begin{equation}
  \boxed{C_{\rm upper}=\left(\frac1{C_1}+\frac1{C_2}\right)^{-1}
  =\frac{2\varepsilon_r}{\varepsilon_r+1}\frac{\varepsilon_0S}{d}}.
\end{equation}
因 $\varepsilon_r>1$，两种改造都使电容变大；但一般有
\[
  C_{\rm left}-C_{\rm upper}
  =\frac{\varepsilon_0S}{d}\frac{(\varepsilon_r-1)^2}{2(\varepsilon_r+1)}\ge0,
\]
所以相同介质量按面积并排放置，增容更明显。

\QAENamed{\splitDielectricDiagram}{两区域等电压形成并联，两层等法向电位移形成串联；$\varepsilon=\varepsilon_r\varepsilon_0$。}{quality:appE:a7}

图~\ref{quality:appE:a7}左图沿面积切开，两块介质都连接同一对电极；右图沿厚度切开，中间界面没有自由面电荷，故两层法向 $D$ 相同。这些边界条件分别决定并联和串联，而不是按介质块数判断。
\end{xesolution}

\begin{origexamplebox}{【奇妙的吸引力】}\label{appE:ex:s21-32-5}
小明又需要你的帮助了。他把一块相对介电常数为 $\varepsilon_r>1$ 的电介质插入平行板电容器中。
极板面积为 $S=ab$、间距为 $d$，介质厚度等于板间距，并沿边长 $a$ 的方向插入深度 $x$。
电介质会受到继续向电容器内部运动的吸引力；这个力从微观上怎样产生并不容易直接看清，请利用虚功法
把它算出来。
\begin{enumerate}[label=(\arabic*),leftmargin=2.2em]
  \item 电容器电荷 $Q$ 保持不变，用介质内电场 $E$ 表示吸入力；
  \item 电容器接在恒压电源上，用介质内电场 $E$ 表示吸入力。
\end{enumerate}

忽略边缘效应对总电容的修正，介质为均匀线性电介质。

\QAENamed{\QAEInsertQuestion}{介质沿 $+x$ 方向进入，插入长度 $0<x<a$；图中尚未给出受力。}{quality:appE:q8}
\end{origexamplebox}

\begin{xesolution}
在薄间隙、介质前缘远离极板两端的内部插入区间，忽略边缘效应对总电容的高阶修正，已填充区域与真空区域相当于并联：
\begin{equation}
  C(x)=\frac{\varepsilon_0b}{d}\,[a+(\varepsilon_r-1)x],
  \qquad
  \frac{\dd C}{\dd x}=\frac{\varepsilon_0b}{d}(\varepsilon_r-1).
\end{equation}

\textbf{(1) 固定电荷。}系统静电能 $W=Q^2/[2C(x)]$，介质所受力沿 $x$ 增大方向：
\[
  F=-\left(\frac{\partial W}{\partial x}\right)_Q
  =\frac{Q^2}{2C^2}\frac{\dd C}{\dd x}
  =\frac{U^2}{2}\frac{\dd C}{\dd x}.
\]
介质内场 $E=U/d$，故
\begin{equation}
  \boxed{F_Q=\frac12\varepsilon_0(\varepsilon_r-1)bdE^2}.
\end{equation}

\textbf{(2) 固定电压。}电容器能量 $CU^2/2$ 随 $C$ 增大而增大，但电源还会输送电荷并做功
$U\dd Q=U^2\dd C$。把电源计入后，机械势能的增量是
$\dd(CU^2/2)-U\dd Q=-U^2\dd C/2$，所以
\[
  F=\frac{U^2}{2}\frac{\dd C}{\dd x}.
\]
仍用 $E=U/d$ 表示，得到
\begin{equation}
  \boxed{F_U=\frac12\varepsilon_0(\varepsilon_r-1)bdE^2}.
\end{equation}
两式形式相同，是因为题目要求用各自状态下的 $E$ 表示；固定 $Q$ 时 $E$ 随 $x$ 改变，固定
$U$ 时 $E$ 不变。力都把介质拉入电容器，因为介质进入会增大电容。

\QAENamed{\dielectricPullDiagram}{吸引力指向插入深度增加方向；固定电压时须连同电源功计算。}{quality:appE:a8}

图~\ref{quality:appE:a8}的箭头沿介质右端向电容器内部，即 $x$ 增大方向。力来自插入前缘附近的非均匀边缘场；忽略边缘场对总电容的高阶修正，并不意味着可以在微观受力分析中删去边缘场。虚功法用 $C'(x)$ 保留了该作用的整体效果。
\end{xesolution}

\begin{origexamplebox}{【充电过程】}\label{appE:ex:s21-32-8}
用电动势为 $\mathcal E$ 的理想电源给初始未充电的电容 $C$ 充电。
\begin{enumerate}[label=(\arabic*),leftmargin=2.2em]
  \item 电路中串联电阻 $R$，充电完成前电源做功和电阻发热各为多少；
  \item 令 $R$ 趋于零，两个结果怎样变化；
  \item 若理想电路中完全没有电阻，应如何理解？
\end{enumerate}

\QAENamed{\QAERCQuestion}{初始未充电的电容器与串联电阻、电源；讨论闭合后的过程。}{quality:appE:q9}
\end{origexamplebox}

\begin{xesolution}
\textbf{(1)} 充电完成后 $Q=C\mathcal E$。电源搬运总电荷 $Q$，故
\begin{equation}
  \boxed{W_{\rm source}=\mathcal EQ=C\mathcal E^2}.
\end{equation}
电容器最终储能
\[
  W_C=\frac12C\mathcal E^2.
\]
其余能量变成焦耳热：
\begin{equation}
  \boxed{Q_R=W_{\rm source}-W_C=\frac12C\mathcal E^2}.
\end{equation}
也可由 $i(t)=\mathcal E e^{-t/(RC)}/R$ 积分验证：
$\int_0^\infty i^2R\dd t=C\mathcal E^2/2$。

\textbf{(2)} 当 $R\to0^+$，充电时间常数 $RC\to0$、初始电流峰值趋于无穷，但热量积分仍为
$C\mathcal E^2/2$，与 $R$ 的数值无关。这是先对有限 $R$ 完成积分、再取极限的结果。

\textbf{(3)} 完全没有电阻、电感和辐射的理想回路不能由一阶 $RC$ 方程描述；它要求电容电压在
零时间内突变，产生无穷大电流，数学模型自相矛盾。真实电路总有寄生电感、导线电阻和电磁辐射。
若保留电感而不含耗散，能量会在电容电场与电感磁场之间振荡，电容电压并不单调趋于
$\mathcal E$；微小电阻或辐射最终耗散那一半能量。因而“零电阻”只能作为 $R\to0^+$ 的极限，
不能把所有寄生参数同时删掉后仍期待同样的稳态充电过程。

\QAENamed{\chargingDiagram}{有限电阻充电的三项能量账；先积分，再取电阻趋零极限。}{quality:appE:a9}

图~\ref{quality:appE:a9}将电源功拆成电场储能与焦耳热。初始未充电是两项各占一半的前提；若电容预先有电压，必须重新计算搬运的电荷和初末储能，不能仍套用这个数值。
\end{xesolution}

\begin{origexamplebox}{【超大电容】}\label{appE:ex:s21-32-11}
考虑地球与大气的理想化电容器模型：地球表面是其中一个极板，带负电荷；
高空的大气等效为另一个极板，带正电荷。把地球表面和高空 $h=\SI{5}{km}$ 处的大气层近似为
同心球形电容器。取地球半径
$R_\oplus=\SI{6.4e6}{m}$，地球带电量绝对值 $Q=\SI{5e5}{C}$，空气电阻率
$\rho=\SI{3e13}{\ohm.m}$。
\begin{enumerate}[label=(\arabic*),leftmargin=2.2em]
  \item 求球形电容器的电容；
  \item 求地表电场能量密度和电容器总能量；
  \item 求大气层的径向漏电电阻；
  \item 求 $RC$ 放电时间常数；
  \item 若每次雷电平均向地面补充 $\SI{25}{C}$ 电荷，维持该电荷量每天约需多少次雷电？
\end{enumerate}

\QAENamed{\QAEEarthQuestion}{地表与等效大气外壳的同心球模型；图中的大气层厚度已放大。}{quality:appE:q10}
\end{origexamplebox}

\begin{xesolution}
令内、外半径分别为 $R=R_\oplus$、$R+h$。

\textbf{(1) 电容。}
\begin{equation}
  C=4\pi\varepsilon_0\frac{R(R+h)}h
  \approx\boxed{\SI{0.912}{F}}.
\end{equation}
因 $h\ll R$，也可用薄层近似 $C\approx4\pi\varepsilon_0R^2/h$。

\textbf{(2) 地表场与能量。}地表外侧
\[
  E(R)=\frac{Q}{4\pi\varepsilon_0R^2}\approx\SI{1.10e2}{V.m^{-1}}.
\]
故
\begin{equation}
  \boxed{u(R)=\frac12\varepsilon_0E^2\approx\SI{5.33e-8}{J.m^{-3}}}.
\end{equation}
总能量
\begin{equation}
  \boxed{W=\frac{Q^2}{2C}\approx\SI{1.37e11}{J}}.
\end{equation}

\textbf{(3) 径向电阻。}半径 $r$ 处厚度 $\dd r$ 的球壳电阻
$\dd R_{\rm el}=\rho\dd r/(4\pi r^2)$，所以
\begin{equation}
  R_{\rm el}=\frac\rho{4\pi}\left(\frac1R-\frac1{R+h}\right)
  \approx\boxed{\SI{2.91e2}{\ohm}}.
\end{equation}

\textbf{(4) 时间常数。}相乘时几何因子恰好消去：
\begin{equation}
  \boxed{\tau=R_{\rm el}C=\rho\varepsilon_0
  \approx\SI{2.66e2}{s}\approx\SI{4.43}{min}}.
\end{equation}

\textbf{(5) 雷电补充率。}漏电电流的平均量级为
\[
  I=\frac Q\tau\approx\SI{1.88e3}{A}.
\]
每天漏失电荷 $I\times\SI{86400}{s}$，所需雷电次数
\begin{equation}
  \boxed{N=\frac{Q}{\tau}\frac{\SI{86400}{s}}{\SI{25}{C}}
  \approx6.5\times10^6\ \text{次/天}}.
\end{equation}
这是题设参数下的数量级模型；真实全球电路的电荷分布、天气电导率和雷暴补充机制更复杂。

\QAENamed{\QAEEarthQuestion}{同一球壳面积 $4\pi r^2$ 分别进入场、电容和漏电电阻，形成几何因子的抵消。}{quality:appE:a10}

图~\ref{quality:appE:a10}中的径向薄壳具有面积 $4\pi r^2$。电阻积分类似串联薄层，电势积分也沿相同半径进行，故 $R_{\rm el}C=\rho\varepsilon_0$ 与球壳的两端半径无关；这要求电阻率在整个壳层内均匀。
\end{xesolution}

\endgroup

\Needspace{18\baselineskip}

\section{恒定电流}

\begingroup

\graphicspath{{../../../sample/xes-original/twenty-second-section/}{../../../sample/xes-original/twenty-second-section/assets/}}

% Restrained TikZ redraws for Lesson 33.

\newcommand{\currentDensityDiagram}{%
\begin{center}
\begin{tikzpicture}[scale=0.88,>=Latex]
  \fill[PaleBlue] (-3.6,-1.15) rectangle (3.6,1.15);
  \draw[curve] (-3.6,-1.15) rectangle (3.6,1.15);
  \draw[DeepBlue,very thick] (0,-1.15)--(0,1.15);
  \foreach \y in {-0.75,-0.25,0.25,0.75}{
    \draw[vector] (-2.8,\y)--(-0.3,\y);
    \draw[vector] (0.3,\y)--(2.8,\y);
  }
  \node[above] at (0,1.2) {截面 $S$};
  \node[TextGray] at (0,-1.7) {$I=\displaystyle\int_S\vect j\cdot\dd\vect S$，箭头表示约定的电流方向};
\end{tikzpicture}

\figcap{补图 33-A\quad 电流是电流密度穿过截面的通量}\label{appE:fig:s22-legacy-01}
\end{center}}

\newcommand{\ivCharacteristicDiagram}{%
\begin{center}
\begin{tikzpicture}[scale=0.88,>=Latex]
  \draw[axis] (-3.5,0)--(3.7,0) node[right] {$U$};
  \draw[axis] (0,-2.6)--(0,2.8) node[above] {$I$};
  \draw[DeepBlue,very thick] (-2.7,-2.1)--(2.7,2.1) node[above right] {理想电阻};
  \draw[Ink,very thick] (-2.9,-2.15) .. controls (-1.8,-0.45) and (-0.7,-0.25) .. (0,0)
    .. controls (0.7,0.25) and (1.8,0.45) .. (2.9,2.15);
  \node[Ink,anchor=west] at (2.35,1.2) {非欧姆元件};
  \draw[guide] (2.0,0)--(2.0,1.56);
  \node[TextGray] at (0,-3.15) {选定端口电压、电流正方向后，曲线完整描述二端元件的外特性};
\end{tikzpicture}

\figcap{二端元件及其伏安特性曲线}\label{appE:fig:s22-legacy-02}
\end{center}}

\newcommand{\ohmHistoryContextDiagram}{%
\begin{center}
\begin{tikzpicture}[scale=0.82,>=Latex]
  \node[inner sep=0] at (-3.7,0)
    {\includegraphics[height=3.4cm]{ohm-portrait-redraw.png}};
  \node[TextGray] at (-3.7,-2.05) {格奥尔格·西蒙·欧姆};
  \begin{scope}[xshift=2.0cm]
    \draw[thick,rounded corners=4pt] (-2.3,-1.45) rectangle (3.0,1.45);
    \foreach \y in {-1.2,-0.8,...,1.2}{
      \draw[TextGray] (-2.42,\y)--(-2.3,\y);
      \draw[TextGray] (3.0,\y)--(3.12,\y);
    }
    \draw[thick] (-1.6,0)--(-0.85,0);
    \draw[fill=SoftGray,thick,rounded corners=2pt] (-0.85,-0.38) rectangle (1.5,0.38);
    \draw[thick] (1.5,0)--(2.25,0);
    \foreach \x in {-0.45,0.05,0.55,1.05}{\draw[DeepBlue,very thick] (\x,-0.38)--(\x,0.38);}
    \node[DeepBlue,font=\large] at (0.3,0.95) {$U=IR$};
    \node[TextGray,font=\small] at (0.3,-0.95) {色环电阻与欧姆定律};
  \end{scope}
\end{tikzpicture}
\figcap{欧姆肖像与纪念欧姆定律的邮票图案}\label{appE:fig:s22-legacy-03}
\end{center}}

\newcommand{\uniformResistorDiagram}{%
\begin{center}
\begin{tikzpicture}[x=1cm,y=1cm,>=Latex]
  \fill[SoftGray] (-3.2,-0.72) rectangle (3.2,0.72);
  \draw[Ink,thick] (-3.2,-0.72) rectangle (3.2,0.72);
  \draw[Ink,thick] (-3.2,0) ellipse (0.24 and 0.72);
  \draw[Ink,thick] (3.2,0) ellipse (0.24 and 0.72);
  \draw[vector] (-2.5,1.25)--(2.5,1.25) node[midway,above] {$I$};
  \draw[<->] (-3.2,-1.2)--node[below] {$l$}(3.2,-1.2);
  \node at (-3.85,0.25) {$A$};
  \node[TextGray,align=left,anchor=west] at (4.3,0.35)
    {均匀材料：电阻率 $\rho$\\截面积 $A$，长度 $l$\\$\displaystyle R=\rho\frac lA$};
\end{tikzpicture}

\figcap{均匀圆柱形电阻的几何参数}\label{appE:fig:s22-legacy-04}
\end{center}}

\newcommand{\gasDischargeContextDiagram}{%
\begin{center}
\begin{tikzpicture}[x=1.03cm,y=0.92cm,>=Latex]
  \draw[axis] (0,0)--(8.1,0) node[right] {$I$};
  \draw[axis] (0,0)--(0,4.15) node[above] {$U$};
  \foreach \x/\lab in {0.8/$10^{-15}$,2.5/$10^{-10}$,4.5/$10^{-5}$,6.4/$10^0$}{
    \draw (\x,0.08)--(\x,-0.08) node[below] {\lab};
  }
  \foreach \y/\lab in {1.0/300,2.0/600,3.35/1000}{
    \draw (0.08,\y)--(-0.08,\y) node[left] {\lab};
  }
  \draw[guide,dashed] (4.25,0)--(4.25,3.8);
  \draw[guide,dashed] (6.25,0)--(6.25,3.8);
  \node[TextGray] at (2.1,3.75) {暗放电};
  \node[TextGray] at (5.2,3.75) {辉光放电};
  \node[TextGray] at (7.1,3.75) {弧光放电};
  \draw[DeepBlue,very thick]
    (0.45,0.25) .. controls (0.7,0.28) and (0.82,0.55) .. (0.92,1.75)
    .. controls (1.02,2.05) and (1.18,2.12) .. (1.35,2.15)
    -- (3.85,2.3)
    .. controls (4.35,2.28) and (4.5,2.05) .. (4.8,1.75)
    .. controls (5.1,1.35) and (5.4,0.9) .. (5.75,0.95)
    .. controls (6.05,1.0) and (6.1,1.75) .. (6.22,2.25)
    .. controls (6.3,2.75) and (6.36,3.45) .. (6.45,3.5)
    .. controls (6.55,3.3) and (6.42,0.6) .. (6.75,0.42)
    .. controls (7.1,0.28) and (7.55,0.48) .. (7.85,0.58);
  \fill[DeepBlue] (3.85,2.3) circle (2.4pt);
  \fill[DeepBlue] (5.75,0.95) circle (2.4pt);
  \fill[DeepBlue] (6.45,3.5) circle (2.4pt);
  \fill[DeepBlue] (6.75,0.42) circle (2.4pt);
  \draw[<-,Ink,thick] (5.25,1.05)--(4.55,0.45)
    node[below left,align=center] {$\dd U/\dd I<0$\\负微分电阻};
\end{tikzpicture}

\figcap{典型气体放电的电压—电流曲线；横轴按对数刻度标示}\label{appE:fig:s22-legacy-05}
\end{center}}

\newcommand{\drudeDiagram}{%
\begin{center}
\begin{tikzpicture}[scale=0.85,>=Latex]
  \fill[SoftGray] (-4.0,-1.35) rectangle (4.0,1.35);
  \foreach \x in {-3.4,-2.3,...,3.3}{\foreach \y in {-0.85,0,0.85}{\fill[TextGray] (\x,\y) circle (2.4pt);}}
  \draw[DeepBlue,thick,-{Latex[length=2mm]}] (-3.5,0.55) .. controls (-2.7,1.05) and (-2.1,-0.8) .. (-1.2,-0.25)
    .. controls (-0.4,0.45) and (0.5,-0.8) .. (1.25,0.1)
    .. controls (1.9,0.85) and (2.65,-0.45) .. (3.45,0.35);
  \draw[vector] (-1.1,2.0)--(1.1,2.0) node[midway,above] {$\vect E$};
  \draw[vector] (1.0,-1.95)--(-1.0,-1.95) node[midway,below] {电子平均漂移};
  \node[TextGray] at (0,-2.55) {高速无规则运动叠加很小的定向漂移；碰撞使动量不断弛豫};
\end{tikzpicture}

\figcap{导体内电子的高速无规则运动与微小定向漂移}\label{appE:fig:s22-legacy-06}
\end{center}}

\newcommand{\resistorConnectionDiagram}{%
\begin{center}
\begin{tikzpicture}[scale=0.85,>=Latex]
  \draw[thick] (-5.0,0)--(-4.2,0);
  \draw[thick,decorate,decoration={zigzag,segment length=6pt,amplitude=2.5pt}] (-4.2,0)--(-2.8,0);
  \draw[thick] (-2.8,0)--(-2.2,0);
  \draw[thick,decorate,decoration={zigzag,segment length=6pt,amplitude=2.5pt}] (-2.2,0)--(-0.8,0);
  \draw[thick] (-0.8,0)--(-0.1,0);
  \node[above] at (-3.5,0.45) {$R_1$}; \node[above] at (-1.5,0.45) {$R_2$};
  \node[TextGray] at (-2.55,-1.0) {串联：电流相同};
  \draw[thick] (1.1,1.05)--(1.8,1.05) (1.1,-1.05)--(1.8,-1.05);
  \draw[thick] (1.8,1.05)--(1.8,-1.05) (5.2,1.05)--(5.2,-1.05);
  \draw[thick] (1.8,0.62)--(2.45,0.62);
  \draw[thick,decorate,decoration={zigzag,segment length=6pt,amplitude=2.5pt}] (2.45,0.62)--(4.5,0.62);
  \draw[thick] (4.5,0.62)--(5.2,0.62);
  \draw[thick] (1.8,-0.62)--(2.45,-0.62);
  \draw[thick,decorate,decoration={zigzag,segment length=6pt,amplitude=2.5pt}] (2.45,-0.62)--(4.5,-0.62);
  \draw[thick] (4.5,-0.62)--(5.2,-0.62);
  \node[above] at (3.48,0.88) {$R_1$}; \node[below] at (3.48,-0.88) {$R_2$};
  \node[TextGray] at (3.15,-1.65) {并联：电压相同};
\end{tikzpicture}

\figcap{电阻的串联与并联}\label{appE:fig:s22-legacy-07}
\end{center}}

\newcommand{\circuitTerminologyDiagram}{%
\begin{center}
\begin{tikzpicture}[x=1cm,y=1cm,>=Latex]
  \coordinate (A) at (-3.6,1.35); \coordinate (B) at (0,1.35);
  \coordinate (C) at (3.6,1.35); \coordinate (D) at (-3.6,-1.35);
  \coordinate (E) at (0,-1.35); \coordinate (F) at (3.6,-1.35);
  \foreach \p/\n in {A/$a$,B/$b$,C/$c$,D/$d$,E/$e$,F/$f$}{
    \fill[DeepBlue] (\p) circle (2.7pt) node[above=4pt] {\n};
  }
  \draw[thick] (A)--node[above] {元件 1}(B)--node[above] {元件 2}(C);
  \draw[thick] (D)--node[below] {元件 4}(E)--node[below] {元件 5}(F);
  \draw[thick] (A)--node[left] {支路}(D);
  \draw[thick] (B)--node[right] {元件 3}(E);
  \draw[thick] (C)--(F);
  \draw[DeepBlue,very thick,rounded corners=3pt,-{Latex[length=2mm]}]
    (-3.95,1.7)--(0.35,1.7)--(0.35,-1.0)--(-3.95,-1.0)--cycle;
  \node[DeepBlue] at (-1.8,0.25) {一个回路};
  \node[TextGray,align=left,anchor=west] at (4.45,0.55)
    {节点：支路的连接点\\支路：相邻节点间的通路\\回路：从节点出发并返回的闭合路径};
\end{tikzpicture}

\figcap{电路网络中的元件、节点、支路与回路}\label{appE:fig:s22-legacy-08}
\end{center}}

\newcommand{\kirchhoffDiagram}{%
\begin{center}
\begin{tikzpicture}[scale=0.9,>=Latex]
  \fill[DeepBlue] (0,0) circle (3pt) node[below right] {节点};
  \draw[vector] (-3.2,1.4)--(-0.2,0.1) node[midway,above] {$I_1$};
  \draw[vector] (-3.2,-1.4)--(-0.2,-0.1) node[midway,below] {$I_2$};
  \draw[vector] (0.2,0.12)--(3.2,1.35) node[midway,above] {$I_3$};
  \draw[vector] (0.2,-0.12)--(3.2,-1.35) node[midway,below] {$I_4$};
  \node[TextGray] at (0,-2.15) {KCL：$I_1+I_2=I_3+I_4$};
  \draw[DeepBlue,thick,rounded corners] (5.1,-1.25) rectangle (9.2,1.25);
  \draw[vector] (5.45,-0.85)--(7.0,-0.85);
  \node[TextGray] at (7.15,-2.15) {KVL：沿闭合回路电压代数和为零};
\end{tikzpicture}

\figcap{基尔霍夫电流定律与电压定律}\label{appE:fig:s22-legacy-09}
\end{center}}

\newcommand{\nodeMethodDiagram}{%
\begin{center}
\begin{tikzpicture}[x=1cm,y=1cm,>=Latex]
  \coordinate (a) at (-3.4,1.35); \coordinate (b) at (0,1.35); \coordinate (c) at (3.4,1.35);
  \coordinate (g1) at (-3.4,-1.35); \coordinate (g2) at (0,-1.35); \coordinate (g3) at (3.4,-1.35);
  \draw[thick] (a)--(b)--(c); \draw[thick] (g1)--(g2)--(g3);
  \draw[thick] (a)--(-3.4,0.55);
  \draw[very thick,DeepBlue] (-3.75,0.38)--(-3.05,0.38) (-3.62,0.03)--(-3.18,0.03);
  \draw[thick] (-3.4,0.03)--(g1); \node[left] at (-3.85,0.2) {$\mathcal E_0$};
  \draw[thick,decorate,decoration={zigzag,segment length=5pt,amplitude=2.2pt}] (b)--(g2);
  \draw[thick,decorate,decoration={zigzag,segment length=5pt,amplitude=2.2pt}] (c)--(g3);
  \draw[thick,decorate,decoration={zigzag,segment length=5pt,amplitude=2.2pt}] (a)--(b);
  \node[above] at (-1.7,1.65) {$R_1$}; \node[right] at (0.35,0) {$R_2$}; \node[right] at (3.75,0) {$R_3$};
  \fill[DeepBlue] (a) circle (2.7pt) node[above left] {$a:\ \varphi_a=\mathcal E_0$};
  \fill[DeepBlue] (b) circle (2.7pt) node[above right] {$b:\ \varphi_b=\varphi$};
  \fill[Ink] (g2) circle (2.7pt) node[below] {参考节点 $\varphi=0$};
  \node[TextGray,align=left,anchor=west] at (5.0,0.45)
    {选参考节点后，各支路电压\\自动等于节点电势之差；\\只需对未知节点列 KCL};
\end{tikzpicture}

\figcap{节点电压法：先求节点电势，再由伏安特性求支路电流}\label{appE:fig:s22-legacy-10}
\end{center}}

\newcommand{\tolmanDiagram}{%
\begin{center}
\begin{tikzpicture}[scale=0.9,>=Latex]
  \draw[curve] (0,0) ellipse (2.25 and 1.35);
  \draw[vector] (0,0)--(0.95,0.8) node[midway,above left] {$r$};
  \draw[DeepBlue,very thick,-{Latex[length=2.5mm]}] (1.65,0.92) arc[start angle=35,end angle=135,x radius=2.05,y radius=1.28];
  \node[DeepBlue] at (0,1.95) {$\omega$};
  \draw[vector] (2.25,0)--(2.25,-1.05) node[right] {$\vect v_e(0)$};
  \draw[very thick] (-0.15,1.35)--(0.15,1.35);
  \node[TextGray,align=left,anchor=west] at (4.0,0.55)
    {晶格突然停止\\电子因惯性继续运动\\随后在碰撞中指数衰减};
\end{tikzpicture}

\figcap{托尔曼--斯图尔特转动铜环实验}\label{appE:fig:s22-legacy-11}
\end{center}}

\newcommand{\transientCurrentDiagram}{%
\begin{center}
\begin{tikzpicture}[scale=0.9,>=Latex]
  \draw[axis] (-0.4,0)--(6.0,0) node[right] {$t$};
  \draw[axis] (0,-0.2)--(0,3.2) node[above] {$I$};
  \draw[guide] (0,2.55)--(5.7,2.55) node[right] {$U/R$};
  \draw[DeepBlue,very thick] (0,0) .. controls (0.75,1.75) and (1.7,2.35) .. (5.5,2.54);
  \draw[<->] (0,-0.55)--node[below] {$\tau$}(1.55,-0.55);
  \node[TextGray] at (2.8,-1.15) {$I(t)=\dfrac UR(1-\e^{-t/\tau})$，载流子惯性使电流不能瞬间建立};
\end{tikzpicture}

\figcap{突加电压后的微观暂态电流}\label{appE:fig:s22-legacy-12}
\end{center}}

\newcommand{\rotatingRingDiagram}{%
\begin{samepage}
\begin{center}
\begin{tikzpicture}[scale=0.70,>=Latex]
  \draw[curve] (0,0) circle (1.65);
  \foreach \a in {0,30,...,330}{\node[DeepBlue] at ({1.65*cos(\a)},{1.65*sin(\a)}) {$+$};}
  \draw[DeepBlue,very thick,-{Latex[length=2.5mm]}] (1.35,0.95) arc[start angle=35,end angle=135,radius=1.65];
  \draw[vector] (0,0)--(1.1,0) node[midway,above] {$R$};
  \node[TextGray,align=left,anchor=west] at (3.4,0.6)
    {每转一周，任一截面通过总电荷 $Q$\\周期 $T=2\pi/\omega$};
\end{tikzpicture}

\figcap{转动带电环与等效环电流}\label{appE:fig:s22-legacy-13}
\end{center}
\end{samepage}}

\newcommand{\vacuumDiodeDiagram}{%
\begin{center}
\begin{tikzpicture}[scale=0.86,>=Latex]
  \draw[DeepBlue,very thick] (-3.5,-1.2)--(-3.5,1.2) node[above] {阴极 $\varphi=0$};
  \draw[DeepBlue,very thick] (3.5,-1.2)--(3.5,1.2) node[above] {阳极 $\varphi=U$};
  \foreach \x in {-2.7,-1.9,...,2.9}{\node[Ink] at (\x,0) {$-$};}
  \draw[vector] (-2.8,-0.75)--(2.7,-0.75) node[midway,below] {电子运动};
  \draw[axis] (-3.5,-1.75)--(4.1,-1.75) node[right] {$x$};
  \draw[<->] (-3.5,-2.25)--node[fill=white,inner sep=1pt] {$L$}(3.5,-2.25);
  \node[TextGray] at (0,-2.95) {$\varphi(x)=U(x/L)^{4/3}$，空间电荷使电势不是线性分布};
\end{tikzpicture}

\figcap{平板真空二极管中的空间电荷限制电流}\label{appE:fig:s22-legacy-14}
\end{center}}

\newcommand{\nonlinearShellDiagram}{%
\begin{samepage}
\begin{center}
\begin{tikzpicture}[scale=0.70,>=Latex]
  \fill[PaleBlue] (0,0) circle (1.85); \fill[white] (0,0) circle (0.82);
  \draw[curve] (0,0) circle (0.82); \draw[curve] (0,0) circle (1.85);
  \foreach \a in {0,45,...,315}{\draw[component] ({0.92*cos(\a)},{0.92*sin(\a)})--({1.7*cos(\a)},{1.7*sin(\a)});}
  \draw[<->] (0,0)--node[above] {$a$}(0.82,0);
  \draw[<->] (0,-0.35)--node[below] {$b$}(1.85,-0.35);
  \node[TextGray,align=left,anchor=west] at (3.5,0.6)
    {$j(r)=I/(4\pi r^2)$\\$j=kE^2$\\故 $E\propto1/r$};
\end{tikzpicture}

\figcap{电导率满足 $\sigma=kE$ 的厚球壳}\label{appE:fig:s22-legacy-15}
\end{center}
\end{samepage}}

\newcommand{\ionChamberDiagram}{%
\begin{center}
\begin{tikzpicture}[scale=0.83,>=Latex]
  \draw[DeepBlue,very thick] (-4.1,1.35)--(4.1,1.35) node[right] {阴极 $-U_0$};
  \draw[DeepBlue,very thick] (-4.1,-1.35)--(4.1,-1.35) node[right] {阳极};
  \draw[guide] (-3.2,-0.35)--(3.1,-0.35) node[right] {粒子径迹};
  \foreach \x in {-2.7,-1.9,...,2.9}{\node[DeepBlue] at (\x,-0.15) {$+$}; \node[Ink] at (\x,-0.55) {$-$};}
  \draw[vector] (-1.4,-0.05)--(-1.4,0.95) node[right] {$v^+$};
  \draw[vector] (0.8,-0.55)--(0.8,-1.12) node[right] {$v^-$};
  \draw[<->] (-4.55,-1.35)--node[left] {$D$}(-4.55,1.35);
  \draw[<->] (3.55,-1.35)--node[right] {$x_0$}(3.55,-0.35);
  \node[TextGray] at (0,-2.15) {电子先到达阳极形成快分量，正离子随后形成慢分量};
\end{tikzpicture}

\figcap{平行板电离室中的电子--离子漂移}\label{appE:fig:s22-legacy-16}
\end{center}}

\newcommand{\ringResistanceDiagram}{%
\begin{center}
\begin{tikzpicture}[scale=0.9,>=Latex]
  \coordinate (RingA) at (-2.7,0);
  \coordinate (RingO) at (0,0);
  \coordinate (RingB) at (2.13,-1.49);
  \draw[curve] (0,0) circle (1.75);
  \fill[DeepBlue] (-1.75,0) circle (3pt) node[left] {$a$};
  \fill[DeepBlue] ({1.75*cos(-35)},{1.75*sin(-35)}) circle (3pt) node[right] {$b$};
  \draw[axis] (0,0)--(-2.7,0);
  \draw[axis] (0,0)--({2.6*cos(-35)},{2.6*sin(-35)});
  \draw pic[draw,"$\theta$",angle radius=11mm,angle eccentricity=1.35]
    {angle=RingA--RingO--RingB};
  \node[TextGray,align=left,anchor=west] at (4.1,0.5)
    {两段电阻 $R_1+R_2=R$\\在端点 $a,b$ 之间并联};
\end{tikzpicture}

\figcap{圆环电阻被接点分成两条并联支路}\label{appE:fig:s22-legacy-17}
\end{center}}

\newcommand{\dualSourceDiagram}{%
\begin{center}
\begin{tikzpicture}[scale=0.88,>=Latex]
  \draw[thick] (-4.0,1.5)--(4.0,1.5) (-4.0,-1.5)--(4.0,-1.5);
  \draw[thick] (-3.2,1.5)--(-3.2,0.95);
  \draw[very thick,DeepBlue] (-3.55,0.75)--(-2.85,0.75) (-3.43,0.38)--(-2.97,0.38);
  \draw[thick] (-3.2,0.38)--(-3.2,-0.1);
  \draw[thick,decorate,decoration={zigzag,segment length=5pt,amplitude=2.2pt}] (-3.2,-0.1)--(-3.2,-1.05);
  \draw[thick] (-3.2,-1.05)--(-3.2,-1.5);
  \node[left] at (-3.65,0.55) {$\mathcal E_1$}; \node[left] at (-3.65,-0.55) {$r_1$};
  \draw[thick] (0,1.5)--(0,0.95);
  \draw[very thick,DeepBlue] (-0.35,0.75)--(0.35,0.75) (-0.23,0.38)--(0.23,0.38);
  \draw[thick] (0,0.38)--(0,-0.1);
  \draw[thick,decorate,decoration={zigzag,segment length=5pt,amplitude=2.2pt}] (0,-0.1)--(0,-1.05);
  \draw[thick] (0,-1.05)--(0,-1.5);
  \node[right] at (0.45,0.55) {$\mathcal E_2$}; \node[right] at (0.45,-0.55) {$r_2$};
  \draw[thick] (3.2,1.5)--(3.2,0.95);
  \draw[thick,decorate,decoration={zigzag,segment length=5pt,amplitude=2.2pt}] (3.2,0.95)--(3.2,-0.95);
  \draw[thick] (3.2,-0.95)--(3.2,-1.5);
  \node[right] at (3.55,0) {$R$};
  \draw[vector] (2.75,1.85)--(3.65,1.85) node[midway,above] {$I$};
\end{tikzpicture}

\figcap{两个有内阻电源并联向负载供电}\label{appE:fig:s22-legacy-18}
\end{center}}


\begin{origexamplebox}{【真空二极管】}\label{appE:ex:s22-33-4}
平板真空二极管的阴极、阳极面积均为 $S$，间距为 $L$。阴极位于 $x=0$ 且电势为零，阳极位于
$x=L$ 且电势为 $U>0$。电子从阴极以可忽略的初速度发射，稳恒电流为 $I$。
\begin{enumerate}[label=(\arabic*),leftmargin=2.2em]
  \item 写出板间电势满足的方程；
  \item 用 $\varphi(x)$ 表示电子速度；
  \item 用 $I$、速度表示电子数密度；
  \item 消去速度和数密度，写出电势的微分方程与空间电荷限制状态的边界条件；
  \item 求伏安关系并证明查尔德--朗缪尔定律 $I=KU^{3/2}$。
\end{enumerate}

\QAENamed{\QAEVacuumQuestion}{阴极、阳极、电势边界与电子运动方向；约定电流大小为正。}{quality:appE:q11}
\end{origexamplebox}

\begin{xesolution}
\textbf{(1)} 电子数密度为 $n(x)$，体电荷密度 $\rho_q=-en$。泊松方程给出
\begin{equation}
  \frac{\dd^2\varphi}{\dd x^2}=-\frac{\rho_q}{\varepsilon_0}
  =\frac{en}{\varepsilon_0}.
\end{equation}

\textbf{(2)} 电子从阴极由静止出发。能量守恒
\[
  \frac12mv^2=e\varphi,
\]
所以
\begin{equation}
  \boxed{v(x)=\sqrt{\frac{2e\varphi(x)}m}}.
\end{equation}

\textbf{(3)} 稳恒电流在各截面相同，取电流大小有
\begin{equation}
  I=en(x)v(x)S,
  \qquad
  \boxed{n(x)=\frac{I}{eSv(x)}}.
\end{equation}

\textbf{(4)} 消去 $n,v$：
\begin{equation}
  \frac{\dd^2\varphi}{\dd x^2}
  =\frac{I}{\varepsilon_0S}\sqrt{\frac{m}{2e}}\,\varphi^{-1/2}.
\end{equation}
空间电荷限制的最大稳定电流状态下，阴极表面电场恰为零；若阴极场反向，发出的电子会被推回。
边界条件为
\begin{equation}
  \varphi(0)=0,
  \qquad \varphi'(0)=0,
  \qquad \varphi(L)=U.
\end{equation}

\textbf{(5)} 记
\[
  A_0=\frac{I}{\varepsilon_0S}\sqrt{\frac{m}{2e}}.
\]
用 $\varphi'$ 乘微分方程并从阴极积分：
\[
  \frac12(\varphi')^2=2A_0\sqrt\varphi,
  \qquad
  \varphi'=2\sqrt{A_0}\,\varphi^{1/4}.
\]
再次积分得
\[
  \varphi(x)=\left(\frac32\sqrt{A_0}\,x\right)^{4/3}.
\]
代入 $\varphi(L)=U$，得到电势分布和电流：
\begin{equation}
  \boxed{\varphi(x)=U\left(\frac xL\right)^{4/3}},
\end{equation}
\begin{equation}
  \boxed{I=\frac49\varepsilon_0S\sqrt{\frac{2e}{m}}
  \frac{U^{3/2}}{L^2}}.
\end{equation}
因此
\begin{equation}
  \boxed{K=\frac49\frac{\varepsilon_0S}{L^2}\sqrt{\frac{2e}{m}}}.
\end{equation}
这一定律建立在平板近似、零初速度、充分电子发射和稳恒一维空间电荷等条件上。
图~\ref{quality:appE:a11} 中电子从零电势阴极向高电势阳极运动；$\varphi\propto x^{4/3}$ 也表明板间电场
$E_x=-\varphi'$、$|\vect E|=\varphi'$ 并非均匀场。

\QAENamed{\vacuumDiodeDiagram}{电子向高电势端加速；电势为非线性分布，电场指向阴极。}{quality:appE:a11}

图~\ref{quality:appE:a11}中电子速率沿 $+x$ 增大，电势也随 $x$ 增大，但 $E_x=-\varphi'<0$，负电荷所受电力仍沿 $+x$。必须区分电场方向、电子速度方向和电流的大小。
\end{xesolution}

\begin{origexamplebox}{【电离计数器】}\label{appE:ex:s22-33-6}
电离计数器测量带电粒子穿过气体时产生的电离电荷。电场把电离产物中的电子和正离子分别引向
阳极和阴极；若入射粒子在气体中被完全吸收，还可由总电离量估计粒子能量。最简单的电离室由
一对平行板电极构成，密封容器中充有工作气体。板间距为 $D$、面积为 $S$，阴极电势为 $-U_0$，
阳极接读出电路。带电粒子在距阳极 $x_0$ 处沿平行于极板的径迹产生 $N$ 对电子--离子；若粒子
沉积能量为 $E_{\rm dep}$、气体中产生一对电子--离子所需平均能量为 $W$，则
$N=E_{\rm dep}/W$。正离子与电子的漂移速率分别为 $v^+,v^-$，假设电荷全部被收集，忽略复合、
电子俘获和次级电离。
\begin{enumerate}[label=(\arabic*),leftmargin=2.2em]
  \item 求均匀电场；
  \item 求极板电容；
  \item 当负载电阻很大时，求输出电压变化 $\Delta U(t)$；
  \item 在 $v^-\gg v^+$ 时描述波形；
  \item 若负载电阻 $R$ 不可忽略，写出输出电压的时间关系。
\end{enumerate}

采用小信号近似，阳极初始电势为零，脉冲幅值 $\Delta U=-\delta V_{\rm anode}$；要求 $|\delta V_{\rm anode}|\ll U_0$，从而可将漂移速率视为常量。

\QAENamed{\QAEIonQuestion}{平行板电离室中电子与离子的两段漂移距离及读出电极。}{quality:appE:q12}
\end{origexamplebox}

\begin{xesolution}
取 $x$ 从阳极指向阴极，阳极 $x=0$、阴极 $x=D$。

\textbf{(1)} 忽略边缘效应，
\begin{equation}
  \boxed{E=\frac{U_0}{D}}.
\end{equation}
场方向由阳极指向阴极；电子向阳极漂移，正离子向阴极漂移。

\textbf{(2)} 气体介电常数近似为 $\varepsilon_0$，
\begin{equation}
  \boxed{C=\frac{\varepsilon_0S}{D}}.
\end{equation}

\textbf{(3)} 平行板的读出权重场为 $1/D$。每个仍在漂移的电荷产生大小为 $ev/D$ 的感应电流。
令
\[
  t_- =\frac{x_0}{v^-},
  \qquad
  t_+ =\frac{D-x_0}{v^+}
\]
分别为电子到阳极、离子到阴极的时间。感应电流大小为
\begin{equation}
  i_{\rm ind}(t)=\frac{Ne}{D}
  \left[v^-H(t_--t)+v^+H(t_+-t)\right],
\end{equation}
其中 $H$ 为阶跃函数。负载近似开路时，感应电荷全部给电容充电：
\begin{equation}
  \boxed{\Delta U(t)=\frac{Ne}{CD}\left[
  v^-\min(t,t_-)+v^+\min(t,t_+)\right]}.
\end{equation}
当两类载流子全部到达电极，$\Delta U(\infty)=Ne/C$，与电离位置 $x_0$ 无关。

\textbf{(4)} 当电离位置不紧靠阴极且 $t_-\ll t_+$ 时，电子在短时间 $t_-$ 内形成陡峭快分量
$Ne x_0/(CD)$；随后正离子在较长时间 $t_+$ 内形成缓慢上升分量
$Ne(D-x_0)/(CD)$。两部分相加仍为 $Ne/C$。

\textbf{(5)} 负载电阻使已形成的电压以时间常数 $\tau_c=RC$ 泄放。电路方程为
\begin{equation}
  \boxed{C\frac{\dd\Delta U}{\dd t}+\frac{\Delta U}{R}=i_{\rm ind}(t)}.
\end{equation}
取 $\Delta U(0)=0$，统一写成卷积形式
\begin{equation}
  \boxed{\Delta U(t)=\frac1C\int_0^t
  i_{\rm ind}(t')\e^{-(t-t')/(RC)}\dd t'}.
\end{equation}
它表明每一小段新产生的感应电压都从产生时刻起指数衰减。若粒子在气体中完全损失能量
$E_p$，且产生一对电子--离子的平均能量为 $W$，还可用 $N=E_p/W$ 反推粒子能量。
图~\ref{quality:appE:a12} 把两段漂移路程直接标成 $x_0$ 与 $D-x_0$：它们分别决定快分量和慢分量的持续时间
及最终幅度。

\QAENamed{\ionChamberDiagram}{电子与正离子的漂移分别形成快、慢脉冲分量；输出取负电压脉冲的正幅值。}{quality:appE:a12}

图~\ref{quality:appE:a12}中电子向阳极、正离子向阴极的漂移对负电压脉冲幅值均作正贡献。只有两类载流子全部到达后，路径长度之和才为 $D$，于是开路总幅值与 $x_0$ 无关；有限负载下较早形成的分量已泄放，峰值通常重新依赖位置。
\end{xesolution}

\begin{origexamplebox}{【圆环电阻】}\label{appE:ex:s22-33-7}
总电阻为 $R$ 的均匀电阻丝弯成圆环，在环上两点 $a,b$ 之间接入电源。固定 $a$，移动 $b$。
\begin{enumerate}[label=(\arabic*),leftmargin=2.2em]
  \item 若电源保持电压 $U$，何时总发热功率最小，最小值是多少；
  \item 若改用保持电流 $I$ 的理想恒流源，结果怎样？
\end{enumerate}

\QAENamed{\QAERingQuestion}{圆环的两个电源接点；接点间不另接短路导线。}{quality:appE:q13}
\end{origexamplebox}

\begin{xesolution}
两接点把圆环分成电阻 $R_1,R_2$ 的两段，且 $R_1+R_2=R$。两段在 $a,b$ 之间并联：
\begin{equation}
  R_{\rm eq}=\frac{R_1R_2}{R_1+R_2}=\frac{R_1R_2}{R}.
\end{equation}

\textbf{(1) 固定电压。}
\[
  P_U=\frac{U^2}{R_{\rm eq}}=\frac{U^2R}{R_1R_2}.
\]
在 $R_1+R_2=R$ 下，乘积 $R_1R_2$ 在 $R_1=R_2=R/2$ 时最大，故功率最小。此时 $b$ 与 $a$
位于直径两端，
\begin{equation}
  \boxed{R_{\rm eq}=\frac R4},
  \qquad
  \boxed{P_{U,\min}=\frac{4U^2}{R}}.
\end{equation}

\textbf{(2) 固定电流。}
\[
  P_I=I^2R_{\rm eq}=I^2\frac{R_1R_2}{R}.
\]
此时要使乘积 $R_1R_2$ 尽量小，因此 $b$ 应尽量靠近 $a$。在允许两接点重合的理想模型中
\begin{equation}
  \boxed{P_{I,\min}=0\quad(b=a)}.
\end{equation}
若规定 $a,b$ 必须不同，则没有被实际取到的正最小值，只有当 $b\to a$ 时趋于零的下确界。固定
电压与固定电流给出相反的优化方向，是因为一个最小化 $1/R_{\rm eq}$，另一个最小化
$R_{\rm eq}$。
图~\ref{quality:appE:a13} 中两条圆弧正是并联的 $R_1,R_2$；把 $b$ 移到 $a$ 的对径点，两条圆弧等长，
对应 $R_1=R_2$。

\QAENamed{\ringResistanceDiagram}{两段圆弧是同端口并联支路；优化变量满足 $R_1+R_2=R$。}{quality:appE:a13}

图~\ref{quality:appE:a13}中的两个支路共用 $a,b$，即使两支路沿不同圆弧也仍是并联。恒流和恒压分别约束电源的不同量，所以同一个 $R_{\rm eq}$ 的极大值在两种优化中具有相反的作用。
\end{xesolution}

\endgroup

\Needspace{18\baselineskip}

\section{复杂电阻网络}

\begingroup

\graphicspath{{../../../sample/xes-original/twenty-third-section/}{../../../sample/xes-original/twenty-third-section/assets/}}

\tikzset{
  resistor/.style={thick,decorate,decoration={zigzag,segment length=4.2pt,amplitude=1.8pt}},
  circuitwire/.style={thick,Ink},
  samenode/.style={circle,draw=DeepBlue,fill=PaleBlue,inner sep=1.6pt},
  terminal/.style={circle,fill=DeepBlue,inner sep=2.2pt}
}

\newcommand{\sourceEquivalentDiagram}{%
\begin{center}
\begin{tikzpicture}[scale=0.88,>=Latex]
  \begin{scope}[xshift=-3.4cm]
    \draw[circuitwire] (-1.5,1)--(-0.45,1);
    \draw[very thick,DeepBlue] (-0.45,1.35)--(-0.45,0.65) (-0.08,1.24)--(-0.08,0.76);
    \draw[resistor] (-0.08,1)--(1.25,1);
    \draw[circuitwire] (1.25,1)--(1.55,1) (-1.5,-1)--(1.55,-1);
    \node at (-0.25,1.65) {$\mathcal E$}; \node at (0.65,1.45) {$r$};
    \node[terminal] at (-1.5,1){}; \node[terminal] at (-1.5,-1){};
  \end{scope}
  \draw[<->,DeepBlue,very thick] (-1.0,0)--(1.0,0)
    node[midway,above] {$\mathcal E=I_{\rm N}r$};
  \begin{scope}[xshift=3.4cm]
    \draw[circuitwire] (-1.5,1)--(1.45,1) (-1.5,-1)--(1.45,-1);
    \draw[circuitwire] (0.15,1)--(0.15,0.45) (0.15,-0.45)--(0.15,-1);
    \draw[resistor] (0.15,0.45)--(0.15,-0.45);
    \draw[circuitwire] (1.05,1)--(1.05,0.48) (1.05,-0.48)--(1.05,-1);
    \draw[DeepBlue,very thick,->] (1.05,-0.48)--(1.05,0.48);
    \node[left] at (-0.05,0) {$r$}; \node[right] at (1.25,0) {$I_{\rm N}$};
    \node[terminal] at (-1.5,1){}; \node[terminal] at (-1.5,-1){};
  \end{scope}
\end{tikzpicture}

\figcap{实际电压源与实际电流源的端口等效}\label{appE:fig:s23-legacy-01}
\end{center}}

\newcommand{\theveninDiagram}{%
\begin{center}
\begin{tikzpicture}[scale=0.88,>=Latex]
  \draw[rounded corners=2mm,fill=PaleBlue,draw=RuleGray] (-4,-1.35) rectangle (-1.25,1.35);
  \node[align=center] at (-2.62,0) {任意线性\\有源二端网络};
  \draw[circuitwire] (-1.25,0.75)--(-0.65,0.75) (-1.25,-0.75)--(-0.65,-0.75);
  \node[terminal] at (-0.65,0.75){}; \node[terminal] at (-0.65,-0.75){};
  \draw[->,DeepBlue,very thick] (0,0)--(1.1,0);
  \draw[circuitwire] (1.6,0.75)--(2.1,0.75);
  \draw[very thick,DeepBlue] (2.1,1.1)--(2.1,0.4) (2.45,0.98)--(2.45,0.52);
  \draw[resistor] (2.45,0.75)--(3.75,0.75);
  \draw[circuitwire] (3.75,0.75)--(4.1,0.75) (1.6,-0.75)--(4.1,-0.75);
  \node at (2.3,1.35) {$U_{\rm oc}$}; \node at (3.15,1.2) {$R_{\rm th}$};
  \node[terminal] at (4.1,0.75){}; \node[terminal] at (4.1,-0.75){};
\end{tikzpicture}

\figcap{戴维南等效只保留开路电压与端口内阻}\label{appE:fig:s23-legacy-02}
\end{center}}

\newcommand{\bridgeDiagram}{%
\begin{center}
\begin{tikzpicture}[scale=0.9,>=Latex]
  \coordinate (L) at (-3,0); \coordinate (T) at (0,1.65);
  \coordinate (B) at (0,-1.65); \coordinate (R) at (3,0);
  \draw[resistor] (L)--node[above left] {$R_1$}(T);
  \draw[resistor] (T)--node[above right] {$R_2$}(R);
  \draw[resistor] (L)--node[below left] {$R_3$}(B);
  \draw[resistor] (B)--node[below right] {$R_4$}(R);
  \draw[resistor,DeepBlue] (T)--node[right] {$R_5$}(B);
  \fill[DeepBlue] (L) circle (2.5pt) node[left] {$A$};
  \fill[DeepBlue] (R) circle (2.5pt) node[right] {$B$};
  \node[TextGray,align=left,anchor=west] at (4.0,0.35)
    {$R_1/R_2=R_3/R_4$\\$\Rightarrow\ U_T=U_B$\\$\Rightarrow\ I_5=0$};
\end{tikzpicture}

\figcap{平衡电桥的中间支路可断开，也可将两端等势点合并}\label{appE:fig:s23-legacy-03}
\end{center}}

\newcommand{\starDeltaDiagram}{%
\begin{center}
\begin{tikzpicture}[scale=0.86,>=Latex]
  \begin{scope}[xshift=-3.1cm]
    \coordinate (A) at (0,1.65); \coordinate (B) at (-1.7,-1.1); \coordinate (C) at (1.7,-1.1);
    \draw[resistor] (A)--node[left] {$R_{AB}$}(B);
    \draw[resistor] (B)--node[below] {$R_{BC}$}(C);
    \draw[resistor] (C)--node[right] {$R_{CA}$}(A);
    \node[above] at (A) {$A$}; \node[left] at (B) {$B$}; \node[right] at (C) {$C$};
  \end{scope}
  \draw[<->,DeepBlue,very thick] (-0.55,0)--(0.55,0);
  \begin{scope}[xshift=3.1cm]
    \coordinate (A) at (0,1.65); \coordinate (B) at (-1.7,-1.1); \coordinate (C) at (1.7,-1.1); \coordinate (O) at (0,0);
    \draw[resistor] (A)--node[right] {$R_A$}(O);
    \draw[resistor] (B)--node[above left] {$R_B$}(O);
    \draw[resistor] (C)--node[above right] {$R_C$}(O);
    \node[above] at (A) {$A$}; \node[left] at (B) {$B$}; \node[right] at (C) {$C$};
    \fill[DeepBlue] (O) circle (2.2pt);
  \end{scope}
\end{tikzpicture}

\figcap{三角形连接与星形连接的端口等效}\label{appE:fig:s23-legacy-04}
\end{center}}

\newcommand{\gridDiagram}{%
\begin{center}
\begin{tikzpicture}[scale=0.8]
  \foreach \x in {0,1,2}{\draw[curve] (\x,0)--(\x,2);}
  \foreach \y in {0,1,2}{\draw[curve] (0,\y)--(2,\y);}
  \foreach \x in {0,1,2}{\foreach \y in {0,1,2}{\fill[Ink] (\x,\y) circle (1.5pt);}}
  \node[left] at (0,0) {$A$}; \node[right] at (2,0) {$B$}; \node[right] at (2,2) {$C$};
  \draw[->,DeepBlue] (3.0,1)--(4.25,1);
  \node[TextGray,align=left,anchor=west] at (4.6,1)
    {先按端点对称性分组\\再合并等势节点};
\end{tikzpicture}

\figcap{田字格网络的节点与端点标记}\label{appE:fig:s23-legacy-05}
\end{center}}

\newcommand{\tetrahedronDiagram}{%
\begin{center}
\begin{tikzpicture}[scale=0.88]
  \coordinate (A) at (0,2); \coordinate (B) at (-2,-1.25);
  \coordinate (C) at (2,-1.25); \coordinate (D) at (0,-0.35);
  \foreach \u/\v in {A/B,A/C,A/D,B/C,B/D,C/D}{\draw[curve] (\u)--(\v);}
  \foreach \u/\pos in {A/above,B/left,C/right,D/below}{\fill[DeepBlue] (\u) circle (2.3pt) node[\pos] {$\u$};}
  \node[TextGray,anchor=west,align=left] at (3.0,0.35)
    {从 $A$ 流入、从 $B$ 流出时\\$C,D$ 关于交换对称，故等势};
\end{tikzpicture}

\figcap{正四面体六条棱均接电阻 $R$}\label{appE:fig:s23-legacy-06}
\end{center}}

\newcommand{\completeGraphDiagram}{%
\begin{center}
\begin{tikzpicture}[scale=0.82]
  \foreach \i in {0,...,4}{\coordinate (P\i) at ({90+72*\i}:1.8);}
  \foreach \i in {0,...,4}{\foreach \j in {0,...,4}{\ifnum\i<\j\draw[curve] (P\i)--(P\j);\fi}}
  \foreach \i in {0,...,4}{\fill[DeepBlue] (P\i) circle (2.2pt);}
  \node[TextGray,anchor=west,align=left] at (3.0,0.35)
    {$K_N$ 中任意两点都有一只 $R$\\其余 $N-2$ 个节点彼此等势};
\end{tikzpicture}

\figcap{$N=5$ 时的完全图电阻网络}\label{appE:fig:s23-legacy-07}
\end{center}}

\newcommand{\octagonDiagram}{%
\begin{center}
\begin{tikzpicture}[scale=0.75]
  \foreach \i in {0,...,7}{\coordinate (O\i) at ({22.5+45*\i}:2.25);}
  \foreach \i in {0,...,7}{\foreach \j in {0,...,7}{
    \pgfmathtruncatemacro{\d}{mod(\j-\i+8,8)}
    \ifnum\i<\j\ifnum\d=4\else\draw[thin,TextGray] (O\i)--(O\j);\fi\fi}}
  \foreach \i in {0,...,7}{\fill[Ink] (O\i) circle (1.7pt);}
  \fill[DeepBlue] (O0) circle (3pt) node[right] {$A$};
  \fill[DeepBlue] (O4) circle (3pt) node[left] {$B$};
  \node[TextGray,anchor=west,align=left] at (3.7,0.3)
    {删去四条最长对角线\\求对径端点 $A,B$ 间电阻};
\end{tikzpicture}

\figcap{变种八边形：除对径点外任意两点均以 $R$ 相连}\label{appE:fig:s23-legacy-08}
\end{center}}

\newcommand{\cubeDiagram}{%
\begin{center}
\begin{tikzpicture}[scale=0.82]
  \coordinate (A) at (-2,-1.3); \coordinate (B) at (1,-1.3);
  \coordinate (C) at (1,1.5); \coordinate (D) at (-2,1.5);
  \coordinate (E) at (-0.8,-0.35); \coordinate (F) at (2.2,-0.35);
  \coordinate (G) at (2.2,2.45); \coordinate (H) at (-0.8,2.45);
  \foreach \u/\v in {A/B,B/C,C/D,D/A,E/F,F/G,G/H,H/E,A/E,B/F,C/G,D/H}{\draw[curve] (\u)--(\v);}
  \draw[guide] (E)--(F) (E)--(H) (A)--(E);
  \foreach \u in {A,B,C,D,E,F,G,H}{\fill[Ink] (\u) circle (1.7pt);}
  \node[left] at (A) {$A$}; \node[below] at (B) {$B$}; \node[right] at (C) {$C$}; \node[above] at (G) {$C_1$};
  \node[TextGray,anchor=west,align=left] at (3.25,0.55)
    {棱相邻：$A,B$\\面对角：$A,C$\\体对角：$A,G$};
\end{tikzpicture}

\figcap{正方体十二条棱均接电阻 $R$}\label{appE:fig:s23-legacy-09}
\end{center}}

\newcommand{\ladderDiagram}{%
\begin{center}
\begin{tikzpicture}[scale=0.78]
  \node[left] at (0,1.2) {$a$}; \node[left] at (0,-1.2) {$b$};
  \foreach \i in {0,1,2,3}{
    \draw[resistor] (1.45*\i,1.2)--(1.45*\i+1.05,1.2);
    \draw[resistor] (1.45*\i,-1.2)--(1.45*\i+1.05,-1.2);
    \draw[resistor] (1.45*\i+1.05,1.2)--(1.45*\i+1.05,-1.2);
    \fill[Ink] (1.45*\i+1.05,1.2) circle (1.4pt);
    \fill[Ink] (1.45*\i+1.05,-1.2) circle (1.4pt);
  }
  \draw[dotted,thick] (5.7,1.2)--(6.45,1.2) (5.7,-1.2)--(6.45,-1.2);
\end{tikzpicture}

\figcap{每一节含上下两个串联电阻和一只横接电阻}\label{appE:fig:s23-legacy-10}
\end{center}}

\newcommand{\threeRailDiagram}{%
\begin{center}
\begin{tikzpicture}[scale=0.72]
  \foreach \row/\lab in {1.6/A,0/B,-1.6/C}{\node[left] at (0,\row) {$\lab$};}
  \foreach \i in {0,1,2,3}{
    \foreach \row in {1.6,0,-1.6}{\draw[resistor] (1.5*\i,\row)--(1.5*\i+1.05,\row);}
    \draw[resistor] (1.5*\i+1.05,1.6)--(1.5*\i+1.05,0);
    \draw[resistor] (1.5*\i+1.05,0)--(1.5*\i+1.05,-1.6);
  }
  \foreach \row in {1.6,0,-1.6}{\draw[dotted,thick] (5.6,\row)--(6.35,\row);}
\end{tikzpicture}

\figcap{三端无限方格网络及端点 $A,B,C$}\label{appE:fig:s23-legacy-11}
\end{center}}

\newcommand{\nestedTriangleDiagram}{%
\begin{center}
\begin{tikzpicture}[scale=0.85]
  \coordinate (A) at (-3,-1.55); \coordinate (B) at (3,-1.55); \coordinate (C) at (0,3.65);
  \draw[curve] (A)--(B)--(C)--cycle;
  \coordinate (L1) at (-1.5,1.05); \coordinate (R1) at (1.5,1.05); \coordinate (M1) at (0,-1.55);
  \draw[curve] (L1)--(R1)--(M1)--cycle;
  \coordinate (L2) at (-0.75,-0.25); \coordinate (R2) at (0.75,-0.25); \coordinate (M2) at (0,1.05);
  \draw[curve] (L2)--(R2)--(M2)--cycle;
  \coordinate (L3) at (-0.375,0.40); \coordinate (R3) at (0.375,0.40); \coordinate (M3) at (0,-0.25);
  \draw[curve] (L3)--(R3)--(M3)--cycle;
  \node[left] at (A) {$A$}; \node[right] at (B) {$B$}; \node[above] at (C) {$C$};
  \node[TextGray,anchor=west,align=left] at (4.0,0.65)
    {边长依次为\\$a,a/2,a/4,\ldots$};
\end{tikzpicture}

\figcap{相邻内接等边三角形在三个中点处与外层相接}\label{appE:fig:s23-legacy-12}
\end{center}}

\newcommand{\sierpinskiDiagram}{%
\begin{center}
\begin{tikzpicture}[scale=0.78]
  \begin{scope}[xshift=-3.4cm]
    \fill[PaleBlue] (-1.7,-1.1)--(1.7,-1.1)--(0,1.84)--cycle;
    \node[below] at (0,-1.45) {第 0 阶};
  \end{scope}
  \begin{scope}
    \fill[PaleBlue] (-1.7,-1.1)--(0,-1.1)--(-0.85,0.37)--cycle;
    \fill[PaleBlue] (0,-1.1)--(1.7,-1.1)--(0.85,0.37)--cycle;
    \fill[PaleBlue] (-0.85,0.37)--(0.85,0.37)--(0,1.84)--cycle;
    \draw[curve] (-1.7,-1.1)--(1.7,-1.1)--(0,1.84)--cycle;
    \node[below] at (0,-1.45) {第 1 阶};
  \end{scope}
  \begin{scope}[xshift=3.4cm]
    \foreach \sx/\sy in {-1.275/-0.735,0/-0.735,-0.6375/0.3675}{
      \fill[PaleBlue] (\sx-0.425,\sy-0.3675)--(\sx+0.425,\sy-0.3675)--(\sx,\sy+0.3675)--cycle;
    }
    \foreach \sx/\sy in {-0.425/-0.735,0.85/-0.735,0.425/0.3675}{
      \fill[PaleBlue] (\sx-0.425,\sy-0.3675)--(\sx+0.425,\sy-0.3675)--(\sx,\sy+0.3675)--cycle;
    }
    \fill[PaleBlue] (-0.425,0.735)--(0.425,0.735)--(0,1.47)--cycle;
    \draw[curve] (-1.7,-1.1)--(1.7,-1.1)--(0,1.84)--cycle;
    \node[below] at (0,-1.45) {第 2 阶};
  \end{scope}
\end{tikzpicture}

\figcap{谢尔宾斯基镂垫的前几阶构造}\label{appE:fig:s23-legacy-13}
\end{center}}


\begin{origexamplebox}{【正方体电阻】}\label{appE:ex:s23-34-6}
正方体框架的十二条棱均接电阻 $R$。求两顶点分别为下列位置关系时的等效电阻：
\begin{enumerate}[label=(\arabic*),leftmargin=2.2em]
  \item 棱相邻；
  \item 面对角线两端；
  \item 体对角线两端。
\end{enumerate}

\QAENamed{\QAECube}{正方体的十二条电阻棱与八个节点；三种待求端口共用该编号。}{quality:appE:q14}
\end{origexamplebox}

\begin{xesolution}
这里以 $1\,\mathrm A$ 电流法统一处理。将流出端取零电势，按保持两个端点不动的对称变换把其余
顶点分组，列每组的 KCL；流入端电势数值就是等效电阻。

\textbf{(1) 棱相邻。}
除端点外的六个顶点按对称性合并后，只需列少量节点方程。解得流入端电势为
$7R/12$，故
\begin{equation}
  \boxed{R_{\rm edge}=\frac{7}{12}R}.
\end{equation}

\textbf{(2) 面对角。}
交换端点 $A,C$ 的垂直平分镜面使电势变为 $V_A-V$；$B,D,F,H$ 均位于该镜面上，故都为 $V_A/2$，而 $E,G$ 分别保留。节点方程给出
\begin{equation}
  \boxed{R_{\rm face}=\frac34R}.
\end{equation}

\textbf{(3) 体对角。}
与流入端相邻的三个顶点等势，与流出端相邻的三个顶点也等势。网络直接化成
\[
  \frac R3+\frac R6+\frac R3,
\]
其中中间两组之间有六只 $R$ 并联。因此
\begin{equation}
  \boxed{R_{\rm body}=\frac56R}.
\end{equation}
三个结果满足“端点在图上的最短距离越远，等效电阻越大”。图~\ref{quality:appE:a14} 用实线画出全部十二条棱；
透视图中的交叉不增加节点。

\QAENamed{\QAECube}{按图中节点写基尔霍夫电流定律；表中的电势依次对应三组端口。}{quality:appE:a14}

图~\ref{quality:appE:a14}的八节点在每种端口选择下均满足相同邻接关系。令 $V_{\rm out}=0$、注入电流为 $I_0$，以 $I_0R$ 归一化电势 $u_i=V_i/(I_0R)$；流入端满足 $3u_A-\sum_{j\sim A}u_j=1$，其余非流出节点满足 $3u_i-\sum_{j\sim i}u_j=0$。例如 $AB$ 端口有 $3u_A-u_B-u_D-u_E=1$、$3u_C-u_B-u_D-u_G=0$。解该七元线性方程组得到：\par\begin{center}\begin{tabular}{c|cccccccc}端口 & $u_A$ & $u_B$ & $u_C$ & $u_D$ & $u_E$ & $u_F$ & $u_G$ & $u_H$\\\hline $AB$ & $7/12$&$0$&$5/24$&$3/8$&$3/8$&$5/24$&$1/4$&$1/3$\\ $AC$ & $3/4$&$3/8$&$0$&$3/8$&$1/2$&$3/8$&$1/4$&$3/8$\\ $AG$ & $5/6$&$1/2$&$1/3$&$1/2$&$1/2$&$1/3$&$0$&$1/3$\end{tabular}\end{center}注入端的 $u_A$ 就是 $R_{\rm eq}/R$；逐行代回节点式，可独立复核三个答案。
\end{xesolution}

\begin{origexamplebox}{【无穷网络】}\label{appE:ex:s23-34-7}
无限梯形网络如图~\ref{quality:appE:q15}，每只电阻均为 $R$。每一节的上、下支路各串联一只 $R$，随后用一只
$R$ 横接两支路。求最左端 $a,b$ 间的等效电阻。

\QAENamed{\QAELadder{0}}{连接完整的二端梯形网络；横接电阻位于每一节两只串联电阻之后。}{quality:appE:q15}
\end{origexamplebox}

\begin{xesolution}
记整个无限网络的等效电阻为 $X$。越过最左端上下两只串联电阻后，后面的网络与原网络完全相同，
等效电阻仍为 $X$；它与该处横接的电阻 $R$ 并联。因此
\begin{equation}
  X=2R+(R\parallel X)=2R+\frac{RX}{R+X}.
\end{equation}
移项得
\[
  X^2-2RX-2R^2=0.
\]
电阻必须为正，取正根
\begin{equation}
  \boxed{X=(1+\sqrt3)R}.
\end{equation}
图~\ref{quality:appE:a15} 中切去第一节后留下的虚线部分，与切去前的整张图在拓扑和每只电阻值上完全相同；
这正是方程中再次出现 $X$ 的原因。

\QAENamed{\QAELadder{1}}{虚线尾部与原网络相同；第一节之后的横接 $R$ 与尾部 $X$ 并联。}{quality:appE:a15}

图~\ref{quality:appE:a15}的虚线框从第一只横接电阻之后开始，故它与该横接电阻是并联关系，端口最前面的两只 $R$ 才是串联。有限节开路截断从 $X_1=3R$ 出发，以 $X_{n+1}=2R+RX_n/(R+X_n)$ 递推，单调下降且有下界 $2R$，所以极限存在并选择正文的正根。
\end{xesolution}

\begin{origexamplebox}{【三端无穷网络】}\label{appE:ex:s23-34-8}
图~\ref{quality:appE:q16} 的三端无限方格网络中，每只电阻均为 $R$。求 $A,C$ 间和 $A,B$ 间的等效电阻。

\QAENamed{\QAEThreeRail{0}}{连接完整的三端无限条带；左端三个端口之间没有直接的竖直电阻。}{quality:appE:q16}
\end{origexamplebox}

\begin{xesolution}
先把三端无限尾部等效成一个三角形电导网络。以 $1/R$ 为电导单位，由上下反射对称性，设尾部
$AB,BC$ 两边的等效电导均为 $p$，$AC$ 边电导为 $q$。在它左边再添一列三只水平电阻和两只
竖直电阻后，网络应与原网络相同。

对新增的三个内部节点作消元，端口等效电导变换为
\begin{equation}
  p'=\frac{p+1}{3p+4},
\end{equation}
\begin{equation}
  q'=\frac{p^2+2pq+2p+3q+1}{(3p+4)(p+2q+2)}.
\end{equation}
自相似要求 $p'=p, q'=q$。取正电导解，可把最终所需电阻化简为
\begin{equation}
  \boxed{R_{AC}=(1+\sqrt5)R},
\end{equation}
\begin{equation}
  \boxed{R_{AB}=\left(1+\frac{\sqrt5+\sqrt{21}}4\right)R}.
\end{equation}
数值上 $R_{AC}\approx3.236R$，$R_{AB}\approx2.705R$。相邻端点的等效电阻较小，符合电流可在
较短路径间分流的直觉。图~\ref{quality:appE:a16} 把尾部的三个端口都保留下来，所以不能像二端梯形网络那样只设
一个未知电阻；这里必须用两个独立参数 $p,q$ 描述三端口。

\QAENamed{\QAEThreeRail{1}\par\QAEThreeReduction}{三端尾部及其等效三角形；$p/R,q/R$ 均为电导，不能误当电阻。}{quality:appE:a16}

图~\ref{quality:appE:a16}保留尾部的三个端口。还可把三端电导矩阵按上下反射分解：反对称电势型 $(1,0,-1)$ 的特征电导为 $\lambda_1=p+2q$，对称型 $(1,-2,1)$ 的为 $\lambda_2=3p$。一节的两根竖直电阻分别给两模增加 $1,3$，三根水平电阻再各串联单位电阻，故 $\lambda_1=(\lambda_1+1)/(\lambda_1+2)$、$\lambda_2=(\lambda_2+3)/(\lambda_2+4)$。取正根可得 $\lambda_1=(\sqrt5-1)/2$、$\lambda_2=(\sqrt{21}-3)/2$，以及 $p=(\sqrt{21}-3)/6$、$q=(3\sqrt5-\sqrt{21})/12$。端口电流分解到两模后，$R_{AC}/R=2/\lambda_1$，$R_{AB}/R=1/(2\lambda_1)+3/(2\lambda_2)$，从而得到上述两个根式。
\end{xesolution}

\begin{origexamplebox}{【分形电阻】}\label{appE:ex:s23-34-10}
边长为 $a$ 的正三角形 $ABC$，每一步都将每个保留三角形的三边二等分，并挖去中央倒立三角形，
无限进行得到谢尔宾斯基镂垫。
\begin{enumerate}[label=(\arabic*),leftmargin=2.2em]
  \item 求最终剩余面积；
  \item 求所有孔洞的总边长，不计最外边界；
  \item 若保留下来的所有边均由单位长度电阻为 $\rho_\ell$ 的金属丝构成，求 $A,B$ 间等效电阻；
  \item 讨论这些结果所反映的分形性质。
\end{enumerate}

\QAENamed{\QAESierpinski}{谢尔宾斯基网络的有限阶构造；各阶均采用相同的单位长度电阻。}{quality:appE:q17}
\end{origexamplebox}

\begin{xesolution}
设原三角形面积为 $S_0=\sqrt3a^2/4$。

\textbf{(1) 面积。}
第 $n$ 步保留 $3^n$ 个边长 $a/2^n$ 的小三角形，故
\begin{equation}
  S_n=S_0\left(\frac34\right)^n
  \longrightarrow\boxed{0}.
\end{equation}

\textbf{(2) 孔洞总边长。}
第 $n\ge1$ 步新挖出 $3^{n-1}$ 个边长 $a/2^n$ 的孔，新增周长
\[
  L_n=3^{n-1}\cdot3\frac{a}{2^n}=a\left(\frac32\right)^n.
\]
因此
\begin{equation}
  \sum_{n=1}^{\infty}L_n
  =a\sum_{n=1}^{\infty}\left(\frac32\right)^n
  =\boxed{+\infty}.
\end{equation}

\textbf{(3) 等效电阻。}
记第 $n$ 阶有限网络两顶点间电阻为 $X_n$。第 $n+1$ 阶由三个第 $n$ 阶副本组成，每个副本的
线性尺度减半，所以若单位长度电阻不变，其全部电阻同时乘 $1/2$。由于三端对称，一个副本若两端
等效电阻为 $X_n$，其等效三角形每边电阻为 $3X_n/2$；缩小后每边为 $3X_n/4$。把三个这样的
三角形按图~\ref{quality:appE:a17} 连接，对中点节点列 KCL，得到新网络的两端电阻为单个三角形边电阻的 $10/9$，
于是
\begin{equation}
  X_{n+1}=\frac{10}{9}\cdot\frac{3X_n}{4}=\frac56X_n.
\end{equation}
第 0 阶只是三边电阻均为 $\rho_\ell a$ 的三角形，
\[
  X_0=\rho_\ell a\parallel2\rho_\ell a=\frac23\rho_\ell a.
\]
所以
\begin{equation}
  X_n=\frac23\rho_\ell a\left(\frac56\right)^n
  \longrightarrow\boxed{0}.
\end{equation}

\textbf{(4) 讨论。}
面积趋零而内部边界总长发散，表明“占据面积”与“边界复杂度”可以朝相反方向演化；网络电阻的递推则由拓扑因子与长度缩放共同决定，本题每步乘 $5/6$，最终等效电阻趋零。图~\ref{quality:appE:a17} 的前三阶
同时显示：每一步都复制三个半尺度副本，这正是三组几何级数与电阻递推的共同来源。

\QAENamed{\QAESierpinski}{三个半尺度副本组成下一阶；相同单位长度电阻意味着每个副本电阻同时缩小一半。}{quality:appE:a17}

图~\ref{quality:appE:a17}每一步都把三个半尺度副本连接成新网络。若小段电阻不随尺度缩短，同样拓扑的递推因子会是 $5/3$；本题固定单位长度电阻，另有长度因子 $1/2$，乘积才为 $5/6$。结论描述这一特定导线模型，不能推广到不同电阻缩放规则的分形材料。
\end{xesolution}

\endgroup

\Needspace{18\baselineskip}
\section{同题异解中的进一步结果}

下列两项与主章题目采用同一模型，但旧版推导给出了主线解答之外仍有价值的推广，
因此单独保留；主章中的基本结论不在这里重复计算。

\subsection{中性导体球的极化率与非均匀场受力}

第20章已经得到：半径为 $R$ 的中性导体球置于局部匀强电场
$\vect E_0$ 中时，球外附加电势等效于球心偶极子，
\[
 \vect p=4\pi\varepsilon_0R^3\vect E_0.
\]
所以导体球的静电极化率为
\[
 \boxed{\alpha=4\pi\varepsilon_0R^3},
 \qquad \vect p=\alpha\vect E_0.
\]
由球面感应电荷密度
$\sigma(\theta)=3\varepsilon_0E_0\cos\theta$ 还可直接检验中性条件：
\[
 Q_+=\int_0^{2\pi}\!\int_0^{\pi/2}
 3\varepsilon_0E_0\cos\theta\,R^2\sin\theta
 \,\dd\theta\dd\phi
 =3\pi\varepsilon_0E_0R^2,
 \qquad Q_-=-Q_+.
\]
净电荷虽然为零，正、负电荷中心的分离仍产生非零偶极矩。

若外场在小球尺度内缓慢变化，把小球从无场处准静态移到场强 $E$ 处，
建立感应偶极矩所需的相互作用能为
\[
 U_{\rm ind}=-\int_0^E\alpha E'\,\dd E'
 =-\frac12\alpha E^2.
\]
因此低阶受力为
\[
 \boxed{\vect F=-\nabla U_{\rm ind}
 =\frac12\alpha\nabla(E^2)}.
\]
它指向场强较大的区域。该近似要求 $R$ 远小于外场的空间变化尺度；
若外场严格均匀，合力为零，只剩内部应力。

\subsection{两等势导体球的有限间距推广}

% CompletenessThermalFields20261024: finite-distance Kelvin images.
设两个等半径导体球的半径为 $R$、球心距为 $L>2R$，二者导通并保持共同电势 $V$；
无穷远取零电势。以 $O_A,O_B$ 表示球心，沿 $O_AO_B$ 的轴线放置像电荷。
定义 $\alpha>0$ 使 $\cosh\alpha=L/(2R)$。图~\ref{appE:complete:two-sphere-images}
标出第一条反演链：$q_0$ 从 $A$ 球心开始，$q_1$ 位于 $B$ 内，再把 $q_1$ 对 $A$ 反演得到 $q_2$；
另有一条从 $B$ 球心开始的镜面对称链。

\begin{qualityfigure}
\centering
\begin{tikzpicture}[>=Latex,font=\small,scale=1.0]
 \draw[thick](-2,0)circle(1.15);\draw[thick](2,0)circle(1.15);
 \draw[dashed](-3.5,0)--(3.5,0);
 \fill[DeepBlue](-2,0)circle(2.2pt);\node[below]at(-2,-.08){$q_0$};
 \fill[TextGray](1.6694,0)circle(2.2pt);\node[above]at(1.60,.12){$q_1$};
 \fill[DeepBlue](-1.6396,0)circle(2.2pt);\node[above]at(-1.60,.12){$q_2$};
 \node[left]at(-3.15,.5){球 $A$};\node[right]at(3.15,.5){球 $B$};
 \node[below]at(-2,-1.19){$O_A$};\node[below]at(2,-1.19){$O_B$};
 \draw[->,DeepBlue](-2,.08)to[bend left=37]node[above]{$q_1=-Rq_0/L$}(1.6694,.08);
 \draw[->,TextGray](1.65,-.08)to[bend left=28]node[below]{$q_2=-Rq_1/(L-d_1)$}(-1.6396,-.08);
 \draw[<->](-2,-1.65)--node[below]{$L$}(2,-1.65);
 \draw[<->](2,.7)--node[above]{$d_1$}(1.6694,.7);
 \draw[->](-2,0)--(-2.8132,-.8132)node[midway,below left]{$R$};
\end{tikzpicture}
\figcap{有限间距双球的一条交错像电荷链。$d_n$ 从该像所在球的球心向另一球量取；
图按 $R/L=1.15/4$ 绘出 $d_1=R^2/L$ 与 $d_2=R^2/(L-d_1)$。
镜面对称的第二条链未重复画出。\label{appE:complete:two-sphere-images}}
\end{qualityfigure}

\textbf{球面反演为什么给出递推。}
一个距球心 $s>R$ 的外部点电荷 $q$，在该球内的 Kelvin 像为
$q'=-Rq/s$，像点距球心 $d'=R^2/s$；在球面上，两者的电势逐点相消。
从球 $A$ 中心的 $q_0=4\pi\varepsilon_0RV$ 开始，它在 $A$ 球面单独给出常数 $V$。
对 $B$ 反演产生 $q_1$，消除 $q_0$ 在 $B$ 球面的电势；再对 $A$ 反演产生 $q_2$，消除 $q_1$ 在 $A$ 球面的电势。
以后继续交替。若 $d_n$ 是第 $n$ 个像距所在球心、向另一球的距离，则
\begin{equation}
 d_0=0,\qquad
 d_{n+1}=\frac{R^2}{L-d_n},\qquad
 q_{n+1}=-\frac{R}{L-d_n}q_n.
 \label{appE:eq:complete-image-recurrence}
\end{equation}
图中 $q_0$ 与 $q_1$ 在球 $B$ 上相消，$q_1$ 与 $q_2$ 在球 $A$ 上相消，清楚显示了配对方式。
单独这条链给 $A$ 球电势 $V$、$B$ 球电势零；再叠加从 $B$ 球心开始的对称链，才得到两球同时为 $V$ 的边界条件。

\textbf{解出位置和电量。}
利用 $2\cosh\alpha\,\sinh[(n+1)\alpha]-\sinh(n\alpha)=\sinh[(n+2)\alpha]$，
从 $n=0$ 归纳可得
\[
 \frac{d_n}{R}=\frac{\sinh(n\alpha)}{\sinh[(n+1)\alpha]},\qquad
 \frac{q_n}{q_0}=(-1)^n\frac{\sinh\alpha}{\sinh[(n+1)\alpha]}.
\]
例如把第一个位置式代入 $L-d_n$，恰得
$R\sinh[(n+2)\alpha]/\sinh[(n+1)\alpha]$，故同时恢复式~\eqref{appE:eq:complete-image-recurrence}的两个递推。
$L>2R$ 时 $\alpha>0$，$|q_n|$ 按指数衰减；像点留在各球内部，并向
$d_\infty=Re^{-\alpha}<R$ 聚集，所以球外电势级数收敛。
配对相消验证了真实球面的边界条件，无穷远电势也为零；由外部 Dirichlet 问题的唯一性，
该像系统在球外等效于真实导体。包住两球的远处高斯面又说明真实总电荷等于全部像电荷之和。

两条链的总和为 $Q_{\rm tot}=2\sum_{n=0}^\infty q_n$，因此
\[
 \boxed{C(L)=\frac{Q_{\rm tot}}V
 =8\pi\varepsilon_0R\sinh\alpha
 \sum_{n=1}^\infty\frac{(-1)^{n+1}}{\sinh(n\alpha)}}.
\]
电容的量纲为 $\varepsilon_0R$。当 $L\to\infty$，$\alpha\to\infty$，除首项外的比值均趋零，
故 $C(L)\to8\pi\varepsilon_0R$，等于两个远离的孤立球电容之和。

相切极限 $L\to2R^+$ 对应 $\alpha\to0^+$。对每个固定 $n$，
$\sinh\alpha/\sinh(n\alpha)\to1/n$。
这里还要允许极限与级数交换：这些正系数随 $n$ 单调下降，且
$\sinh(n\alpha)\ge n\sinh\alpha$，所以第 $N$ 项之后的交错余项不超过 $1/(N+1)$，
这个界与 $\alpha$ 无关。先控制余项、再对有限项取极限，便有
\[
 C(2R)=8\pi\varepsilon_0R\sum_{n=1}^\infty\frac{(-1)^{n+1}}n
 =\boxed{8\pi\varepsilon_0R\ln2}.
\]
这既恢复第20章的相切结果，也说明有限间距的双曲函数递推怎样连续过渡为 $q_n\propto(-1)^n/(n+1)$。


\endgroup


\cleardoublepage
\phantomsection
\addcontentsline{toc}{part}{第7卷　高阶选读}
\part*{第7卷　高阶选读}
% AdvancedMechanicsAssumptions20261023-v1
% AdvancedMechanicsQuality20261023-v1
\begingroup
\newif\ifAdvHideHeading
\RenewDocumentEnvironment{advexamplebox}{m}{%
  \refstepcounter{advexample}%
  \begin{tcolorbox}[
    enhanced,unbreakable,colback=white,colframe=RuleGray,boxrule=0.5pt,arc=0mm,
    title={例~\theadvexample\quad #1},fonttitle=\bfseries,
    coltitle=Ink,colbacktitle=PaleBlue,borderline west={1.4pt}{0pt}{DeepBlue}]
}{\end{tcolorbox}}

\RenewDocumentEnvironment{advsolution}{+b}{%
  \par\smallskip\begingroup\small\noindent
  \ifAdvHideHeading\global\AdvHideHeadingfalse\else\LectureSolutionHeading\fi #1\par\endgroup\smallskip
}{}

\RenewDocumentEnvironment{advdefinition}{m +b}{%
  \refstepcounter{advdefinition}%
  \begin{tcolorbox}[
    enhanced,unbreakable,colback=white,colframe=RuleGray,boxrule=0.5pt,arc=0mm,
    title={定义~\theadvdefinition\quad #1},fonttitle=\bfseries,
    coltitle=Ink,colbacktitle=PaleGold,borderline west={1.2pt}{0pt}{Gold}]
  #2
  \end{tcolorbox}
}{}

\RenewDocumentEnvironment{advtheorem}{m +b}{%
  \refstepcounter{advtheorem}%
  \begin{tcolorbox}[
    enhanced,unbreakable,colback=white,colframe=RuleGray,boxrule=0.5pt,arc=0mm,
    title={定理~\theadvtheorem\quad #1},fonttitle=\bfseries,
    coltitle=Ink,colbacktitle=PaleBlue,borderline west={1.2pt}{0pt}{DeepBlue}]
  #2
  \end{tcolorbox}
}{}


\AdvancedUnit{1}{力学}

\section{分析力学基础}

\subsection{广义坐标、约束与拉格朗日方程}

牛顿方程
\[
  m_i\ddot{\vect r}_i=\vect F_i,\qquad \dot{\vect p}_i=\vect F_i
\]
对直角坐标很直接，但在多体系统、曲线坐标和带约束系统中，它会引入大量未知约束力。
分析力学的做法是：先用适合体系几何结构的坐标消去约束，再用标量函数统一编码动力学。

设系统含有 $n$ 个质点，完整约束写成
\begin{equation}\label{advA1:eq:constraints}
 f^A(\vect r_1,\ldots,\vect r_n,t)=0,
 \qquad A=1,\ldots,N .
\end{equation}
若它们相互独立，系统自由度为 $f=3n-N$。对应的 $f$维可运动空间称为位形空间。
广义坐标 $q_1,\ldots,q_f$ 的选取应使映射
\begin{equation}\label{advA1:eq:point-map}
 \vect r_i=\vect r_i(q_1,\ldots,q_f,t)
\end{equation}
自动满足式\eqref{advA1:eq:constraints}。约束于是不再以附加方程出现，而是被吸收到函数\eqref{advA1:eq:point-map}的结构中。不显含时间的约束称为稳定约束；显含时间的称为非稳定约束。本节不处理无法积分为位形约束的非完整约束。

\begin{advexamplebox}{刚性椭球的十二个约束}\label{advA1:ex:ellipsoid}
六个质点固定在椭球的六个顶点 $A,A',B,B',C,C'$ 上，
$AA'=2a,BB'=2b,CC'=2c$。用十八个直角坐标写出一组相互独立的十二个约束。

\begin{qualityfigure}
\begin{tikzpicture}[>=Latex,font=\small]
\draw[DeepBlue](0,0)ellipse(2.5 and 1.5);\draw[dashed](-2.5,0)arc(180:360:2.5 and .45);\draw(-2.5,0)arc(180:0:2.5 and .45);
\foreach\x/\y/\t in {2.5/0/A,-2.5/0/A',.9/.45/B,-.9/-.45/B',0/1.5/C,0/-1.5/C'}{\fill(\x,\y)circle(2pt)node[above right]{$\t$};}
\draw[dashed](-2.5,0)--(2.5,0)(-.9,-.45)--(.9,.45)(0,-1.5)--(0,1.5);\draw[<->](-2.5,-2.1)--(2.5,-2.1)node[midway,below]{$AA'=2a$};\draw[<->](-3,-1.5)--(-3,1.5)node[midway,left]{$CC'=2c$};\node at(1.7,.85){$BB'=2b$};
\end{tikzpicture}
\figcap{六个对径顶点和三条主轴（题设几何）}\label{advA1:fig:quality-ellipsoid-q}
\end{qualityfigure}
\end{advexamplebox}
\begin{advsolution}
\begin{qualityfigure}
\begin{tikzpicture}[>=Latex,font=\small]
\draw[DeepBlue](0,0)ellipse(2.5 and 1.5);\draw[dashed](-2.5,0)arc(180:360:2.5 and .45);\draw(-2.5,0)arc(180:0:2.5 and .45);
\foreach\x/\y/\t in {2.5/0/A,-2.5/0/A',.9/.45/B,-.9/-.45/B',0/1.5/C,0/-1.5/C'}{\fill(\x,\y)circle(2pt)node[above right]{$\t$};}
\draw[dashed](-2.5,0)--(2.5,0)(-.9,-.45)--(.9,.45)(0,-1.5)--(0,1.5);\draw[<->](-2.5,-2.1)--(2.5,-2.1)node[midway,below]{$AA'=2a$};\draw[<->](-3,-1.5)--(-3,1.5)node[midway,left]{$CC'=2c$};\node at(1.7,.85){$BB'=2b$};\node[align=left]at(6.5,0){三组对径端点共用同一中点\\$|\vect u|=2a,\ |\vect v|=2b,\ |\vect w|=2c$\\三条对径矢量两两正交};
\end{tikzpicture}
\figcap{六个对径顶点和三条主轴（解答辅助图）}\label{advA1:fig:quality-ellipsoid-a}
\end{qualityfigure}
图~\ref{advA1:fig:quality-ellipsoid-a}中，$A,A'$、$B,B'$、$C,C'$ 各自的中点必须相同；长度条件与正交条件是另外六个标量约束，不能把六个端点各自固定在实验室坐标中。

引入三条对径矢量
\[
 \vect u=\vect r_A-\vect r_{A'},\quad
 \vect v=\vect r_B-\vect r_{B'},\quad
 \vect w=\vect r_C-\vect r_{C'}.
\]
先要求三条对径共中点：
\[
 \vect r_A+\vect r_{A'}-\vect r_B-\vect r_{B'}=\vect0,
 \quad
 \vect r_A+\vect r_{A'}-\vect r_C-\vect r_{C'}=\vect0 .
\]
这是 $6$ 个标量方程。再规定对径长度和两两正交：
\[
 \vect u^2=4a^2,\qquad\vect v^2=4b^2,\qquad\vect w^2=4c^2,
 \qquad
 \vect u\!\cdot\!\vect v=\vect v\!\cdot\!\vect w
 =\vect w\!\cdot\!\vect u=0 .
\]
共有 $6+3+3=12$ 个独立标量约束，因而自由度为 $18-12=6$，与一般刚体的三个平动和三个转动自由度一致。
\end{advsolution}


位形空间中的光滑轨道 $q_\alpha=q_\alpha(t)$ 确定了系统的运动。一般动力学量是
$A(q_\alpha,\dot q_\alpha,t)$，沿轨道的全导数为
\begin{equation}\label{advA1:eq:total-derivative}
 \dot A=\sum_\alpha\frac{\partial A}{\partial q_\alpha}\dot q_\alpha
 +\sum_\alpha\frac{\partial A}{\partial\dot q_\alpha}\ddot q_\alpha
 +\frac{\partial A}{\partial t}.
\end{equation}
由式\eqref{advA1:eq:point-map}，
\[
 \dot{\vect r}_i=\sum_\alpha
 \frac{\partial\vect r_i}{\partial q_\alpha}\dot q_\alpha
 +\frac{\partial\vect r_i}{\partial t}.
\]
因此动能对广义速度总是二次多项式：
\begin{align}
 T&=T_2+T_1+T_0,\label{advA1:eq:T-decomp}\\
 T_2&=\frac12\sum_{\alpha\beta}m_{\alpha\beta}\dot q_\alpha\dot q_\beta,
 &m_{\alpha\beta}&=\sum_i m_i
 \frac{\partial\vect r_i}{\partial q_\alpha}\!\cdot\!
 \frac{\partial\vect r_i}{\partial q_\beta},\notag\\
 T_1&=\sum_\alpha\pi_\alpha\dot q_\alpha,
 &\pi_\alpha&=\sum_i m_i
 \frac{\partial\vect r_i}{\partial q_\alpha}\!\cdot\!
 \frac{\partial\vect r_i}{\partial t},\notag\\
 T_0&=\frac12\tau,
 &\tau&=\sum_i m_i\left(\frac{\partial\vect r_i}{\partial t}\right)^2.\notag
\end{align}
稳定约束且坐标变换不含时时，$T_1=T_0=0$。

\begin{advdefinition}{广义机械动量}\label{advA1:def:momentum}
\[
 p_\alpha^{(\mathrm{mech})}=\frac{\partial T}{\partial\dot q_\alpha}
 =\sum_\beta m_{\alpha\beta}\dot q_\beta+\pi_\alpha .
\]
\end{advdefinition}

\begin{advdefinition}{广义力}\label{advA1:def:force}
对虚位移 $\dd\vect r_i=\sum_\alpha
(\partial\vect r_i/\partial q_\alpha)\dd q_\alpha$，定义
\[
 Q_\alpha=\sum_i\vect F_i\!\cdot\!
 \frac{\partial\vect r_i}{\partial q_\alpha}.
\]
若坐标变换显含时间，实际元功还包含
\[
 \dd W=\sum_\alpha Q_\alpha\dd q_\alpha+h\dd t,
 \qquad h=\sum_i\vect F_i\!\cdot\!\frac{\partial\vect r_i}{\partial t}.
\]
\end{advdefinition}

\begin{advexamplebox}{在膨胀球面上运动的质点}\label{advA1:ex:expanding-sphere}
设 $v>0,t>0$。质量为 $m$ 的质点被约束在 $x^2+y^2+z^2-v^2t^2=0$ 上。以球坐标角
$\theta,\varphi$ 为广义坐标，求位矢、动能、广义动量及其全导数、广义力和元功。

\begin{qualityfigure}
\begin{tikzpicture}[>=Latex,font=\small]
\draw[->](0,0)--(3.5,0)node[right]{$x$};\draw[->](0,0)--(0,3)node[above]{$z$};\draw[->](0,0)--(-1.3,-1)node[left]{$y$};\node[below left]at(0,0){$O$};\draw[DeepBlue](0,0)circle(2);\coordinate(P)at(1.3,1.3);\fill(P)circle(2pt)node[right]{$P$};\draw[thick](0,0)--(P)node[pos=.65,below right=3pt]{$r$};
\draw[dashed](P)--(1.3,-.45)--(0,0);\draw(0,.7)arc(90:45:.7)node[midway,above]{$\theta$};\draw(.7,0)arc(0:-19.1:.7)node[midway,below]{$\varphi$};\node at(3.5,2){$r(t)=vt$};
\end{tikzpicture}
\figcap{膨胀球面的瞬时半径和两角坐标（题设几何）}\label{advA1:fig:quality-expanding-sphere-q}
\end{qualityfigure}
\end{advexamplebox}
\begin{advsolution}
\begin{qualityfigure}
\begin{tikzpicture}[>=Latex,font=\small]
\draw[->](0,0)--(3.5,0)node[right]{$x$};\draw[->](0,0)--(0,3)node[above]{$z$};\draw[->](0,0)--(-1.3,-1)node[left]{$y$};\node[below left]at(0,0){$O$};\draw[DeepBlue](0,0)circle(2);\coordinate(P)at(1.3,1.3);\fill(P)circle(2pt)node[right]{$P$};\draw[thick](0,0)--(P)node[pos=.65,below right=3pt]{$r$};
\draw[dashed](P)--(1.3,-.45)--(0,0);\draw(0,.7)arc(90:45:.7)node[midway,above]{$\theta$};\draw(.7,0)arc(0:-19.1:.7)node[midway,below]{$\varphi$};\node at(3.5,2){$r(t)=vt$};\draw[->,Gold](P)--(2.1,2.1)node[right]{$\vect e_r$};\draw[->,Gold](P)--(2.1,.5)node[right]{$\vect e_\theta$};\node[align=left]at(4,-1.4){$\dd s_r=\dd r$\\$\dd s_\theta=r\dd\theta$\\$\dd s_\varphi=r\sin\theta\dd\varphi$};
\end{tikzpicture}
\figcap{膨胀球面的瞬时半径和两角坐标（解答辅助图）}\label{advA1:fig:quality-expanding-sphere-a}
\end{qualityfigure}
图~\ref{advA1:fig:quality-expanding-sphere-a}中，$P$ 的径向距离已由非稳定约束固定为 $vt$；仍可独立变化的是 $\theta,\varphi$，而实际微位移还含球面膨胀的径向部分。

\[
 \vect r=vt(\sin\theta\cos\varphi,\sin\theta\sin\varphi,\cos\theta),
\]
\[
 T=\frac m2\left[v^2+v^2t^2
 (\dot\theta^2+\sin^2\theta\,\dot\varphi^2)\right].
\]
于是
\[
 p_\theta=mv^2t^2\dot\theta,
 \qquad p_\varphi=mv^2t^2\sin^2\theta\,\dot\varphi,
\]
\[
 \dot p_\theta=mv^2(2t\dot\theta+t^2\ddot\theta),
\]
\[
 \dot p_\varphi=mv^2\!\left[
 2t\sin^2\theta\dot\varphi
 +2t^2\sin\theta\cos\theta\dot\theta\dot\varphi
 +t^2\sin^2\theta\ddot\varphi\right].
\]
若合力的球坐标分量为 $(F_r,F_\theta,F_\varphi)$，则
\[
 Q_\theta=vtF_\theta,\qquad Q_\varphi=vt\sin\theta F_\varphi,
 \qquad h=vF_r,
\]
\[
 \dd W=Q_\theta\dd\theta+Q_\varphi\dd\varphi+vF_r\dd t.
\]
最后一项是球面本身膨胀方向上的做功，不能误并入两个角向广义力。
\end{advsolution}

\begin{advtheorem}{拉格朗日方程的引理}\label{advA1:thm:lagrange-lemma}
沿任一满足牛顿方程的运动，
\begin{equation}\label{advA1:eq:lagrange-lemma}
 \frac{\dd}{\dd t}\frac{\partial T}{\partial\dot q_\alpha}
 -\frac{\partial T}{\partial q_\alpha}=Q_\alpha .
\end{equation}
\end{advtheorem}
证明不再省略。由链式法则，
\[
 \frac{\partial T}{\partial\dot q_\alpha}
 =\sum_i m_i\dot{\vect r}_i\!\cdot\!
 \frac{\partial\vect r_i}{\partial q_\alpha}.
\]
对时间求全导数，并注意
$\partial\dot{\vect r}_i/\partial q_\alpha
=\dd(\partial\vect r_i/\partial q_\alpha)/\dd t$，得
\[
 \frac{\dd}{\dd t}\frac{\partial T}{\partial\dot q_\alpha}
 =\sum_i m_i\ddot{\vect r}_i\!\cdot\!
 \frac{\partial\vect r_i}{\partial q_\alpha}
 +\sum_i m_i\dot{\vect r}_i\!\cdot\!
 \frac{\partial\dot{\vect r}_i}{\partial q_\alpha}.
\]
第二项正是 $\partial T/\partial q_\alpha$，两式相减，再用
$m_i\ddot{\vect r}_i=\vect F_i$ 便得式\eqref{advA1:eq:lagrange-lemma}。

\begin{advexamplebox}{二维极坐标中的动力学方程}\label{advA1:ex:polar}
由拉格朗日引理推导平面极坐标中的质点动力学方程。

\begin{qualityfigure}
\begin{tikzpicture}[>=Latex,font=\small]
\draw[->](0,0)--(4.6,0)node[right]{$x$};\draw[->](0,0)--(0,3.4)node[above]{$y$};\node[below left]at(0,0){$O$};\coordinate(P)at(2.5,1.8);\fill(P)circle(2pt)node[left]{$P$};\draw[thick](0,0)--(P)node[midway,below]{$r$};\draw(.9,0)arc(0:35.75:.9)node[midway,right]{$\theta$};
\end{tikzpicture}
\figcap{极坐标中随位置旋转的正交基矢（题设几何）}\label{advA1:fig:quality-polar-q}
\end{qualityfigure}
\end{advexamplebox}
\begin{advsolution}
\begin{qualityfigure}
\begin{tikzpicture}[>=Latex,font=\small]
\draw[->](0,0)--(4.6,0)node[right]{$x$};\draw[->](0,0)--(0,3.4)node[above]{$y$};\node[below left]at(0,0){$O$};\coordinate(P)at(2.5,1.8);\fill(P)circle(2pt)node[left]{$P$};\draw[thick](0,0)--(P)node[midway,below]{$r$};\draw(.9,0)arc(0:35.75:.9)node[midway,right]{$\theta$};\draw[->,DeepBlue](P)--(3.5,2.52)node[right]{$\vect e_r$};\draw[->,DeepBlue](P)--(1.78,2.8)node[left]{$\vect e_\theta$};
\end{tikzpicture}
\figcap{极坐标中随位置旋转的正交基矢（解答辅助图）}\label{advA1:fig:quality-polar-a}
\end{qualityfigure}
图~\ref{advA1:fig:quality-polar-a}中，径向与横向单位矢量随 $\theta$ 改变。先把速度写成径向、横向两项，再求导时须同时对基矢求导。

$T=m(\dot r^2+r^2\dot\theta^2)/2$，而 $Q_r=F_r, Q_\theta=rF_\theta$。代入式\eqref{advA1:eq:lagrange-lemma}：
\[
 m(\ddot r-r\dot\theta^2)=F_r,
 \qquad
 m(r\ddot\theta+2\dot r\dot\theta)=F_\theta.
\]
第二式括号内每项均为加速度；乘以质量 $m$ 后与右端力同量纲。若漏掉 $r\ddot\theta$ 中的 $r$，量纲会立即失配。
\end{advsolution}

\begin{advexamplebox}{三维球坐标中的动力学方程}\label{advA1:ex:spherical}
推导球坐标 $(r,\theta,\varphi)$ 中的三个动力学方程，并与加速度的矢量分解比较。

\begin{qualityfigure}
\begin{tikzpicture}[>=Latex,font=\small]
\draw[->](0,0)--(3.5,0)node[right]{$x$};\draw[->](0,0)--(0,3)node[above]{$z$};\draw[->](0,0)--(-1.3,-1)node[left]{$y$};\node[below left]at(0,0){$O$};\draw[DeepBlue](0,0)circle(2);\coordinate(P)at(1.3,1.3);\fill(P)circle(2pt)node[right]{$P$};\draw[thick](0,0)--(P)node[pos=.65,below right=3pt]{$r$};
\draw[dashed](P)--(1.3,-.45)--(0,0);\draw(0,.7)arc(90:45:.7)node[midway,above]{$\theta$};\draw(.7,0)arc(0:-19.1:.7)node[midway,below]{$\varphi$};
\end{tikzpicture}
\figcap{球坐标的径向、极角和方位角位移（题设几何）}\label{advA1:fig:quality-spherical-q}
\end{qualityfigure}
\end{advexamplebox}
\begin{advsolution}
\begin{qualityfigure}
\begin{tikzpicture}[>=Latex,font=\small]
\draw[->](0,0)--(3.5,0)node[right]{$x$};\draw[->](0,0)--(0,3)node[above]{$z$};\draw[->](0,0)--(-1.3,-1)node[left]{$y$};\node[below left]at(0,0){$O$};\draw[DeepBlue](0,0)circle(2);\coordinate(P)at(1.3,1.3);\fill(P)circle(2pt)node[right]{$P$};\draw[thick](0,0)--(P)node[pos=.65,below right=3pt]{$r$};
\draw[dashed](P)--(1.3,-.45)--(0,0);\draw(0,.7)arc(90:45:.7)node[midway,above]{$\theta$};\draw(.7,0)arc(0:-19.1:.7)node[midway,below]{$\varphi$};\draw[->,Gold](P)--(2.1,2.1)node[right]{$\vect e_r$};\draw[->,Gold](P)--(2.1,.5)node[right]{$\vect e_\theta$};\node[align=left]at(4,-1.4){$\dd s_r=\dd r$\\$\dd s_\theta=r\dd\theta$\\$\dd s_\varphi=r\sin\theta\dd\varphi$};
\end{tikzpicture}
\figcap{球坐标的径向、极角和方位角位移（解答辅助图）}\label{advA1:fig:quality-spherical-a}
\end{qualityfigure}
图~\ref{advA1:fig:quality-spherical-a}中，三个线元分别为 $\dd r,r\dd\theta,r\sin\theta\dd\varphi$；第三项的纬圈半径是 $r\sin\theta$，不能漏掉这个度量因子。

\[
 T=\frac m2[\dot r^2+r^2\dot\theta^2+r^2\sin^2\theta\dot\varphi^2],
\quad
 (Q_r,Q_\theta,Q_\varphi)=(F_r,rF_\theta,r\sin\theta F_\varphi).
\]
逐个代入引理得
\begin{align*}
 m[\ddot r-r(\dot\theta^2+\sin^2\theta\dot\varphi^2)]&=F_r,\\
 m[r\ddot\theta+2\dot r\dot\theta-r\sin\theta\cos\theta\dot\varphi^2]&=F_\theta,\\
 m[r\sin\theta\ddot\varphi+2\dot r\sin\theta\dot\varphi
 +2r\cos\theta\dot\theta\dot\varphi]&=F_\varphi.
\end{align*}
直接对 $\vect r=r\vect e_r$ 求两次导数，再投影到
$(\vect e_r,\vect e_\theta,\vect e_\varphi)$，会得到完全相同的三个分量，这正是坐标协变性的具体检验。
\end{advsolution}

\begin{advexamplebox}{平移且转动的参考系}\label{advA1:ex:moving-frame}
建立坐标变换
\[
 \binom{\bar x}{\bar y}=
 \begin{pmatrix}\cos\theta&\sin\theta\\-\sin\theta&\cos\theta\end{pmatrix}
 \left[\binom{x}{y}-\binom{X}{Y}\right],
\]
其中 $X,Y,\theta$ 均为时间函数。解释其几何意义，写出元功并推导新系中的运动方程。

\begin{qualityfigure}
\begin{tikzpicture}[>=Latex,font=\small]
\draw[->](0,0)--(4.2,0)node[right]{$x$};\draw[->](0,0)--(0,3.2)node[above]{$y$};\node[below left]at(0,0){$O$};\coordinate(C)at(1.8,1);\fill(C)circle(2pt)node[below right]{$O'$};\draw[->,Gold](0,0)--(C)node[midway,below]{$\vect R$};\draw[->,DeepBlue](C)--(3.8,2.2)node[right]{$\bar x$};\draw[->,DeepBlue](C)--(.8,2.67)node[above]{$\bar y$};\draw(2.4,1)arc(0:31:.6)node[midway,right]{$\theta$};
\end{tikzpicture}
\figcap{原点平移与坐标基矢转动（题设几何）}\label{advA1:fig:quality-moving-frame-q}
\end{qualityfigure}
\end{advexamplebox}
\begin{advsolution}
\begin{qualityfigure}
\begin{tikzpicture}[>=Latex,font=\small]
\draw[->](0,0)--(4.2,0)node[right]{$x$};\draw[->](0,0)--(0,3.2)node[above]{$y$};\node[below left]at(0,0){$O$};\coordinate(C)at(1.8,1);\fill(C)circle(2pt)node[below right]{$O'$};\draw[->,Gold](0,0)--(C)node[midway,below]{$\vect R$};\draw[->,DeepBlue](C)--(3.8,2.2)node[right]{$\bar x$};\draw[->,DeepBlue](C)--(.8,2.67)node[above]{$\bar y$};\draw(2.4,1)arc(0:31:.6)node[midway,right]{$\theta$};\node[align=left]at(4,-1.2){平移原点后再旋转分量\\相同物理矢量使用不同基矢表示};
\end{tikzpicture}
\figcap{原点平移与坐标基矢转动（解答辅助图）}\label{advA1:fig:quality-moving-frame-a}
\end{qualityfigure}
图~\ref{advA1:fig:quality-moving-frame-a}中，先减去原点位移 $\vect R=(X,Y)$，再将同一个相对位置投影到转过 $\theta$ 的新轴上。对坐标变换求导时，平移和基矢转动均随时间变化。

这是原点位于 $\vect R=(X,Y)$、坐标轴相对惯性系转过 $\theta$ 的非惯性系。令
$\vect\rho=(\bar x,\bar y)$，$\vect\Omega=\dot\theta\vect e_z$，则
\[
 \vect r=\vect R+\mathsf R(\theta)\vect\rho,
\]
\[
 \dd W=\vect F'\!\cdot\!\dd\vect\rho
 +\vect F\!\cdot\![\dot{\vect R}+\vect\Omega\times
 \mathsf R(\theta)\vect\rho]\dd t,
\]
其中 $\vect F'=\mathsf R^{\mathsf T}\vect F$。两次求导得
\[
 m\ddot{\vect\rho}=\vect F'-m\vect a_0'
 -2m\vect\Omega\times\dot{\vect\rho}
 -m\dot{\vect\Omega}\times\vect\rho
 -m\vect\Omega\times(\vect\Omega\times\vect\rho).
\]
右边依次是真实力、平移惯性力、科里奥利力、欧拉力和离心力。所有项都是力，量纲为 $MLT^{-2}$。
\end{advsolution}

\begin{advdefinition}{理想约束}\label{advA1:def:ideal-constraint}
若约束 $f^A=0$ 对第 $i$ 个质点产生的约束力为
\[
 \vect P_i^{,A}=\lambda_A\nabla_i f^A,
\]
则称其为理想约束。由于
$\sum_i\nabla_i f^A\cdot\partial\vect r_i/\partial q_\alpha=0$，它对任一允许的虚位移都不做虚功，所以其广义力恒为零。
\end{advdefinition}

\begin{advdefinition}{广义势}\label{advA1:def:generalized-potential}
若存在 $V(q_\alpha,\dot q_\alpha,t)$ 使主动力的广义力满足
\[
 Q_\alpha=\frac{\dd}{\dd t}\frac{\partial V}{\partial\dot q_\alpha}
 -\frac{\partial V}{\partial q_\alpha},
\]
则称相互作用可由广义势描述。若 $V$ 不含速度，就回到 $Q_\alpha=-\partial V/\partial q_\alpha$。
\end{advdefinition}

\begin{advdefinition}{拉格朗日量}\label{advA1:def:lagrangian}
\[
 L(q_\alpha,\dot q_\alpha,t)=T(q_\alpha,\dot q_\alpha,t)
 -V(q_\alpha,\dot q_\alpha,t),
 \qquad p_\alpha=\frac{\partial L}{\partial\dot q_\alpha}.
\]
含速势时，$p_\alpha$ 是正则动量，一般不等于机械动量。
\end{advdefinition}

\begin{advtheorem}{拉格朗日方程}\label{advA1:thm:lagrange}
对可由广义势描述的理想约束系统，
\begin{equation}\label{advA1:eq:lagrange}
 \frac{\dd}{\dd t}\frac{\partial L}{\partial\dot q_\alpha}
 -\frac{\partial L}{\partial q_\alpha}=0,
 \qquad \alpha=1,\ldots,f .
\end{equation}
\end{advtheorem}
证明只需把定义\ref{advA1:def:generalized-potential}代入引理\ref{advA1:thm:lagrange-lemma}，再将动能项和势函数项合并。这一步也说明，方程的普遍性不是额外假设，而是由牛顿方程和点变换的链式法则得到的。

\begin{advexamplebox}{收短摆绳的圆锥摆}\label{advA1:ex:conical-pendulum}
均匀重力场中，质量 $m$ 的摆球以竖直向下方向为极轴，摆角为 $\theta$、方位角为 $\varphi$。摆长以常速 $v$ 收短，$l(t)=l_0-vt$。写出 $\theta,\varphi$ 的运动方程。

\begin{qualityfigure}
\begin{tikzpicture}[>=Latex,font=\small]
\fill(0,2.5)circle(2pt)node[above]{$O$};\draw[dashed](0,2.5)--(0,-1);\coordinate(P)at(1.5,.1);\draw[thick](0,2.5)--(P);\fill(P)circle(3pt)node[right]{$m$};\draw(0,1.6)arc(-90:-58:.9)node[midway,below right]{$\theta$};\draw[dashed](0,.1)ellipse(1.5 and .4);\node[above right]at(.75,1.3){$l(t)=l_0-vt$};
\end{tikzpicture}
\figcap{收短摆绳的摆角和瞬时轨道（题设几何）}\label{advA1:fig:quality-conical-pendulum-q}
\end{qualityfigure}
\end{advexamplebox}
\AdvHideHeadingtrue
\begin{advsolution}题图~\ref{advA1:fig:quality-conical-pendulum-q}中，摆长沿绳方向改变，方位角对应的水平轨道半径为 $l\sin\theta$；因而 $\varphi$ 的广义动量含 $l^2\sin^2\theta$，收绳时不能视为常系数。

\noindent\begin{minipage}{\linewidth}{\bfseries\color{DeepBlue}解答}\par
\begin{center}\begin{minipage}{\linewidth}\centering
\begin{tikzpicture}[scale=.95]
 \coordinate (O) at (0,1.5);\coordinate (P) at (1.3,-.5);
 \draw[axis] (0,1.8)--(0,-1.5) node[below]{$z$};
 \draw[guide] (0,-.5) ellipse (1.3 and .38);
 \draw[curve] (O)--(P) node[midway,right=5pt]{$l(t)$};
 \fill (O) circle(1.5pt) node[left]{$O$};\fill (P) circle(2.2pt) node[right=5pt]{$m$};
 \draw[DeepBlue,->] (0,.85) arc[start angle=-90,end angle=-56.98,radius=.65];
 \node at (.32,.70) {$\theta$};
 \draw[TextGray,->] ({1.3*cos(220)},{-.5+.38*sin(220)}) arc[start angle=220,end angle=300,x radius=1.3,y radius=.38];
 \node at (-.85,-1.08) {$\varphi$ 正向};
\end{tikzpicture}
\bookcaption{收短摆绳时的两个广义坐标}\label{advA1:fig:conical}
\end{minipage}\end{center}\end{minipage}

如图\ref{advA1:fig:conical}，
\[
 L=\frac m2[\dot l^2+l^2\dot\theta^2+l^2\sin^2\theta\dot\varphi^2]
 +mgl\cos\theta .
\]
$l$ 是给定函数而非待求自由度。对 $\theta,\varphi$ 列方程：
\[
 l^2\ddot\theta+2l\dot l\dot\theta
 -l^2\sin\theta\cos\theta\dot\varphi^2+gl\sin\theta=0,
\]
\[
 \frac{\dd}{\dd t}(ml^2\sin^2\theta\dot\varphi)=0,
 \qquad \dot l=-v.
\]
第二式表明绕竖直轴的角动量守恒；收绳装置可以交换能量，所以摆球的机械能一般不守恒。
\end{advsolution}

\begin{advexamplebox}{缓慢变化引力常量的开普勒问题}\label{advA1:ex:kepler-time}
平面有心力 $F_r=-mk(t)/r^2$。忽略使 $k$ 改变的物质流、辐射等其他作用，写出运动方程并判断哪些守恒量仍然存在。

\begin{qualityfigure}
\begin{tikzpicture}[>=Latex,font=\small]
\draw[->](0,0)--(4.6,0)node[right]{$x$};\draw[->](0,0)--(0,3.4)node[above]{$y$};\node[below left]at(0,0){$O$};\coordinate(P)at(2.5,1.8);\fill(P)circle(2pt)node[left]{$P$};\draw[thick](0,0)--(P)node[midway,below]{$r$};\draw(.9,0)arc(0:35.75:.9)node[midway,right]{$\theta$};\draw[->,Gold,very thick](P)--(1.7,1.224)node[above left]{$F_r$};\node at(3,-1){$k=k(t)$};
\end{tikzpicture}
\figcap{随时间变化的中心引力和极径（题设几何）}\label{advA1:fig:quality-kepler-time-q}
\end{qualityfigure}
\end{advexamplebox}
\begin{advsolution}
\begin{qualityfigure}
\begin{tikzpicture}[>=Latex,font=\small]
\draw[->](0,0)--(4.6,0)node[right]{$x$};\draw[->](0,0)--(0,3.4)node[above]{$y$};\node[below left]at(0,0){$O$};\coordinate(P)at(2.5,1.8);\fill(P)circle(2pt)node[left]{$P$};\draw[thick](0,0)--(P)node[midway,below]{$r$};\draw(.9,0)arc(0:35.75:.9)node[midway,right]{$\theta$};\draw[->,Gold,very thick](P)--(1.7,1.224)node[above left]{$F_r$};\node at(3,-1){$k=k(t)$};\draw[->,DeepBlue](P)--(3.5,2.52)node[right]{$\vect e_r$};\draw[->,DeepBlue](P)--(1.78,2.8)node[left]{$\vect e_\theta$};
\end{tikzpicture}
\figcap{随时间变化的中心引力和极径（解答辅助图）}\label{advA1:fig:quality-kepler-time-a}
\end{qualityfigure}
图~\ref{advA1:fig:quality-kepler-time-a}中，力始终沿 $PO$，所以对 $O$ 的力矩为零；但力的强度显含时间，角动量守恒与能量是否守恒须分别判断。

\[
 L=\frac m2(\dot r^2+r^2\dot\varphi^2)+\frac{mk(t)}r .
\]
拉格朗日方程给出
\[
 \ddot r-r\dot\varphi^2=-\frac{k(t)}{r^2},
 \qquad \frac{\dd}{\dd t}(mr^2\dot\varphi)=0.
\]
因为 $\varphi$ 仍是循环坐标，角动量严格守恒。哈密顿量
\[
 H=\frac m2(\dot r^2+r^2\dot\varphi^2)-\frac{mk(t)}r
\]
却满足
\[
 \dot H=-\frac{\partial L}{\partial t}=-\frac{m\dot k(t)}r,
\]
所以 $\dot k\ne0$ 时能量不守恒。量纲检查：$k$ 的量纲是 $L^3T^{-2}$，$mk/r$ 的量纲正是能量。
\end{advsolution}

\Needspace{12\baselineskip}
\subsection{守恒量、对称性与罗斯约化}

\begin{advdefinition}{哈密顿量}\label{advA1:def:hamiltonian}
\begin{equation}\label{advA1:eq:H}
 H(q,\dot q,t)=\sum_\alpha p_\alpha\dot q_\alpha-L,
 \qquad p_\alpha=\frac{\partial L}{\partial\dot q_\alpha}.
\end{equation}
在稳定约束、不含时坐标变换且 $V=V(q)$ 时，$H=T+V$。一般情形下不应把哈密顿量无条件地称为机械能。
\end{advdefinition}

\begin{advtheorem}{能量积分}\label{advA1:thm:energy}
沿真实运动，
\begin{equation}\label{advA1:eq:Hdot}
 \dot H=-\frac{\partial L}{\partial t}.
\end{equation}
因而当且仅当轨道上 $\partial L/\partial t=0$ 时，$H$ 守恒。
\end{advtheorem}
证明：对式\eqref{advA1:eq:H}求导，
\[
 \dot H=\sum_\alpha(\dot p_\alpha\dot q_\alpha+p_\alpha\ddot q_\alpha)
 -\sum_\alpha\left(\frac{\partial L}{\partial q_\alpha}\dot q_\alpha
 +\frac{\partial L}{\partial\dot q_\alpha}\ddot q_\alpha\right)
 -\frac{\partial L}{\partial t}.
\]
用 $\dot p_\alpha=\partial L/\partial q_\alpha$ 和
$p_\alpha=\partial L/\partial\dot q_\alpha$，前两组项逐项消去，即得式\eqref{advA1:eq:Hdot}。

\begin{advexamplebox}{系数随时间同比增长的谐振子}\label{advA1:ex:scaled-oscillator}
设 $m,k>0$；涉及缓变极限时取 $0<\gamma<\sqrt{k/m}$，并固定不随 $\gamma$ 改变的有限初始位置与速度。对 $L=(m\dot q^2-kq^2)/2$：
(1) 作坐标标度变换使动能系数为 $1$；
(2) 再考虑
\[
 L=\frac12me^{2\gamma t}\dot q^2-\frac12ke^{2\gamma t}q^2,
 \qquad \omega^2=k/m,
\]
求方程、哈密顿量及其全导数；
(3) 用 $\bar q=e^{\gamma t}q$ 找出守恒量；
(4) 证明在 $t=1/\gamma$ 的时间尺度上，$\gamma\to0$ 时哈密顿量的变化趋于零。

\begin{qualityfigure}
\begin{tikzpicture}[>=Latex,font=\small]
\draw[very thick](0,-.5)--(0,.6);\draw[decorate,decoration={coil,segment length=5pt,amplitude=3pt}](0,0)--(3,0);\draw[fill=PaleBlue](3,-.4)rectangle(4,.4);\draw[->](3.5,-.9)--(5,-.9)node[right]{$q$};\node at(2,1.2){拉格朗日量中的系数同比增长};
\end{tikzpicture}
\figcap{谐振坐标与指数标度后的振动（题设几何）}\label{advA1:fig:quality-scaled-oscillator-q}
\end{qualityfigure}
\end{advexamplebox}
\begin{advsolution}
\begin{qualityfigure}
\begin{tikzpicture}[>=Latex,font=\small]
\draw[very thick](0,-.5)--(0,.6);\draw[decorate,decoration={coil,segment length=5pt,amplitude=3pt}](0,0)--(3,0);\draw[fill=PaleBlue](3,-.4)rectangle(4,.4);\draw[->](3.5,-.9)--(5,-.9)node[right]{$q$};\node at(2,1.2){拉格朗日量中的系数同比增长};\draw[->](0,-2.4)--(5,-2.4)node[right]{$t$};\draw[DeepBlue,domain=0:5,samples=90]plot(\x,{-2.4+.6*exp(-.3*\x)*cos(240*\x)});\node at(2,-3.3){$q=e^{-\gamma t}\bar q$};
\end{tikzpicture}
\figcap{谐振坐标与指数标度后的振动（解答辅助图）}\label{advA1:fig:quality-scaled-oscillator-a}
\end{qualityfigure}
图~\ref{advA1:fig:quality-scaled-oscillator-a}中，原坐标 $q$ 具有指数包络，标度坐标 $\bar q=e^{\gamma t}q$ 消去了包络；这项坐标变化不能直接把原哈密顿量当成新系统的守恒能量。

(1) 取 $Q=\sqrt m\,q$，则
$L=(\dot Q^2-\omega^2Q^2)/2$。

(2) 方程与哈密顿量为
\[
 \ddot q+2\gamma\dot q+\omega^2q=0,
 \qquad H=\frac m2e^{2\gamma t}(\dot q^2+\omega^2q^2),
\]
\[
 \dot H=-\frac{\partial L}{\partial t}
 =\gamma me^{2\gamma t}(\omega^2q^2-\dot q^2).
\]

(3) 代入 $q=e^{-\gamma t}\bar q$ 后
\[
 L=\frac m2[(\dot{\bar q}-\gamma\bar q)^2-\omega^2\bar q^2]
\]
不显含时间，因而
\[
 \mathcal E=\frac m2[\dot{\bar q}^{\,2}+(\omega^2-\gamma^2)\bar q^2]
 =H+m\gamma e^{2\gamma t}q\dot q
\]
守恒。

(4) 当 $\gamma<\omega$ 时，
\[
 q=e^{-\gamma t}C\cos(\Omega t+\phi),
 \qquad \Omega=\sqrt{\omega^2-\gamma^2}.
\]
代入 $H$ 可见其非常数部分的振幅为 $\mathcal O(\gamma)$，故即使取
$t=1/\gamma$，仍有
$|H(1/\gamma)-H(0)|\le C_1\gamma+C_2\gamma^2\to0$。这是缓变极限，不是对任意有限 $\gamma$ 的精确能量守恒。
还可给出对全部时间成立的界，避免把短时估计直接外推到 $1/\gamma$。
由守恒量 $\mathcal E=m(\dot{\bar q}^{\,2}+\Omega^2\bar q^2)/2$，有
\[
 H-\mathcal E=-m\gamma\bar q\dot{\bar q}+m\gamma^2\bar q^2,
 \quad |H-\mathcal E|\le\frac{\gamma\mathcal E}{\Omega}
 +\frac{2\gamma^2\mathcal E}{\Omega^2}.
\]
因而 $|H(t)-H(0)|\le2\gamma\mathcal E/\Omega+4\gamma^2\mathcal E/\Omega^2$。
在题设固定初值且 $\gamma\to0^+$ 时，$\mathcal E$ 有界、$\Omega\to\omega>0$，
故这个一致界在 $t=1/\gamma$ 时仍趋于零。

\end{advsolution}

\begin{advtheorem}{循环坐标与动量积分}\label{advA1:thm:cyclic}
若 $\partial L/\partial q_f=0$，则 $q_f$ 是循环坐标，且
\[
 p_f=\frac{\partial L}{\partial\dot q_f}=\text{常量}.
\]
反之，若 $p_f$ 对所有运动守恒，则拉格朗日方程给出 $\partial L/\partial q_f=0$。
\end{advtheorem}

\begin{advtheorem}{连续对称性的诺特守恒量}\label{advA1:thm:noether}
设单参数点变换
$\bar q_\alpha=\Lambda_\alpha(q,t,\lambda)$满足
$\Lambda_\alpha(q,t,0)=q_\alpha$，并使拉格朗日量不变，则
\begin{equation}\label{advA1:eq:noether-charge}
 \mathcal Q=\sum_\alpha p_\alpha
 \left.\frac{\partial\Lambda_\alpha}{\partial\lambda}\right|_{\lambda=0}
\end{equation}
是守恒量。
\end{advtheorem}
证明：记
$\delta q_\alpha=(\partial\Lambda_\alpha/\partial\lambda)_0\delta\lambda$。不变性给出
\[
 0=\delta L=\sum_\alpha\left(
 \frac{\partial L}{\partial q_\alpha}\delta q_\alpha
 +p_\alpha\frac{\dd}{\dd t}\delta q_\alpha\right)
 =\frac{\dd}{\dd t}\sum_\alpha p_\alpha\delta q_\alpha,
\]
其中用了 $\dot p_\alpha=\partial L/\partial q_\alpha$。除以常数 $\delta\lambda$ 即得式\eqref{advA1:eq:noether-charge}。
空间平移对应线动量，旋转对应角动量，时间平移对应能量。

\begin{advexamplebox}{对称解绕的旋转圆筒}\label{advA1:ex:unwinding-cylinder}
质量 $M$、半径 $R$ 的薄壁圆筒可绕定轴自由转动。两条相同轻绳对称缠绕，自由端各系质量 $m$ 的小球。初始直绳段长度为零，整体角速度为 $\omega_0$。释放后绳不打滑。（1）求每侧打开长度 $s(t)$；（2）若两侧可打开绳长的总和为 $L$，讨论原题所问“打开后”绳相对固定在筒上的直径的夹角随时间的关系。须说明所选绳端连接模型与公式适用阶段。

\begin{qualityfigure}
\begin{tikzpicture}[>=Latex,font=\small]
\draw[thick](0,0)circle(1.3);\draw[dashed](-1.7,0)--(1.7,0);\draw[very thick](1.3,0)--(1.3,2);\draw[very thick](-1.3,0)--(-1.3,-2);\fill(1.3,2)circle(3pt)node[right]{$m$};\fill(-1.3,-2)circle(3pt)node[left]{$m$};\draw(0,0)--(1.3,0)node[midway,below]{$R$};\node at(0,.4){$M$};\draw[->,Gold](.7,0)arc(0:100:.7)node[midway,above]{$\omega_0$};\node at(3.5,.7){对称解绕；俯视图};
\end{tikzpicture}
\figcap{两侧对称解绕的切线绳段（题设几何）}\label{advA1:fig:quality-unwinding-cylinder-q}
\end{qualityfigure}
\end{advexamplebox}
\AdvHideHeadingtrue
\begin{advsolution}题图~\ref{advA1:fig:quality-unwinding-cylinder-q}中，圆筒两侧的直绳段与筒面相切，另一侧由中心对称得到。绳不打滑联系的是筒转角、切点方位和打开长度，公式只在仍有筒面接触的阶段有效。

\noindent\begin{minipage}{\linewidth}{\bfseries\color{DeepBlue}解答}\par
\begin{center}\begin{minipage}{\linewidth}\centering
\begin{tikzpicture}[scale=.88]
 \draw[very thick] (0,0) circle(1.05);
 \draw[guide] (-1.45,0)--(1.45,0) node[right]{初始固定直径};
 \coordinate (T) at ({1.05*cos(35)},{1.05*sin(35)});
 \draw[curve] (T)--++({2.0*cos(125)},{2.0*sin(125)}) coordinate (P);
 \fill (P) circle(2.5pt) node[above left]{$m$};
 \draw[->,DeepBlue] (.45,0) arc[start angle=0,end angle=35,radius=.45];
 \node at (.68,.19) {$\phi$};
 \draw[<->,TextGray] ($(T)+(.08,.05)$)--($(P)+(.08,.05)$) node[midway,right]{$s$};
 \fill (T) circle(1.8pt) node[right=4pt]{$T$};
 \node at (0,-1.38) {另一侧与图示关于圆心对称};
\end{tikzpicture}
\bookcaption{圆筒上的切点、直绳段和固定直径}\label{advA1:fig:cylinder}
\end{minipage}\end{center}\end{minipage}

如图\ref{advA1:fig:cylinder}，以圆筒转角 $\theta$ 和每侧直绳长 $s$ 为坐标。图中 $\phi$ 是切点半径相对初始固定直径的方位角，直绳切线的方位角为 $\phi+\pi/2$。绳不打滑意味着
$\phi=\theta-s/R$。两球与圆筒的总动能为
\[
 T=\frac12(M+2m)R^2\dot\theta^2
 +ms^2\left(\dot\theta-\frac{\dot s}{R}\right)^2.
\]
$\theta$ 是循环坐标。记
\[
 A=(M+2m)R^2+2ms^2,\quad B=-\frac{2ms^2}{R},
 \quad C=\frac{2ms^2}{R^2},
\]
则 $T=(A\dot\theta^2+2B\dot\theta\dot s+C\dot s^2)/2$。角动量和能量分别为
\[
 J=A\dot\theta+B\dot s=(M+2m)R^2\omega_0,
 \quad E=\frac12(M+2m)R^2\omega_0^2.
\]
消去 $\dot\theta$ 后
\[
 E=\frac12\left(C-\frac{B^2}{A}\right)\dot s^2+\frac{J^2}{2A}.
\]
代入 $A,B,C,J,E$ 化简，得 $\dot s=R\omega_0$，约定 $\omega_0>0$ 并取打开分支，
\[
 \boxed{s(t)=R\omega_0t},\qquad \dot\theta=\omega_0,
 \qquad \phi=\theta-s/R=0.
\]
所以直绳段在惯性系中方向不变。圆筒相对切点半径转过 $\theta-\phi=\omega_0t=s/R$，这是相对初态的转角；直绳与筒上随转直径的有向几何夹角则为
$\beta=\phi+\pi/2-\theta=\pi/2-\omega_0t$（直线夹角按模 $\pi$ 理解）。若每侧可打开的长度为 $L/2$，则打开时间
$t_*=L/(2R\omega_0)$，在这一阶段终点，累计相对转角为 $L/(2R)$，而绳与随转直径的夹角为 $\beta_*=\pi/2-L/(2R)$（模 $\pi$）。

原文没有明确绳端是否固定圆筒，也未交代全部解绕后两球与绳的连接、脱离或绷紧约束。因此以上 $s(t)$ 与相对夹角公式只用于仍在解绕的阶段 $0<t<t_*$；不能不加新条件延伸到 $t>t_*$。若“打开后”指已完全离开筒面，必须先补全绳端约束，再用新的几何约束与守恒量求运动，原给定信息不足以唯一确定该阶段的夹角。
\end{advsolution}

若 $q_f$ 是循环坐标且 $p_f=c_f$，对 $\dot q_f$ 作勒让德变换得罗斯函数
\begin{equation}\label{advA1:eq:routhian}
 \mathcal R=L-p_f\dot q_f.
\end{equation}
用 $p_f=c_f$ 解出 $\dot q_f$并代回后，其余 $f-1$ 个坐标仍满足拉格朗日形式的罗斯方程。

\begin{advexamplebox}{自由转动圆环上小珠的小振动}\label{advA1:ex:routh-hoop}
\emph{题源说明：原资料仅给主题或注明课堂描述；以下条件为依据主题所作的教学补全。}\par
半径 $a$ 的竖直圆环可绕竖直直径无摩擦转动，圆环对该轴的转动惯量为 $I$。质量 $m$ 的小珠在环上滑动。以珠从最低点向上的角位移 $\theta$ 和环的方位角 $\varphi$ 为坐标。在绕竖直轴角动量为 $J$ 时，求最低点附近的振动频率及稳定条件。

\begin{qualityfigure}
\begin{tikzpicture}[>=Latex,font=\small]
\draw[thick](0,0)ellipse(1.5 and 2);\draw[->](0,-2.4)--(0,2.5)node[above]{竖直转轴};\fill(0,0)circle(2pt)node[left]{$O$};\coordinate(P)at(.7949,-1.6961);\draw[dashed](0,0)--(P);\fill(P)circle(3pt)node[right]{$m$};\draw[DeepBlue](0,-.7)arc(-90:-58:.525 and .7);\node at(.34,-.8){$\theta$};\draw[->,Gold](-.7,.7)arc(180:0:.7 and .2)node[right]{$\varphi$};\node[align=left]at(3.3,0){环半径 $a$\\轴惯量 $I$};
\end{tikzpicture}
\figcap{可绕竖直直径转动的圆环与小珠（题设几何）}\label{advA1:fig:quality-routh-hoop-q}
\end{qualityfigure}
\end{advexamplebox}
\AdvHideHeadingtrue
\begin{advsolution}题图~\ref{advA1:fig:quality-routh-hoop-q}中，小珠到竖直轴的距离是 $a\sin\theta$；它使绕轴惯量成为 $I+ma^2\sin^2\theta$，消去循环坐标时应固定角动量 $J$。

\noindent\begin{minipage}{\linewidth}{\bfseries\color{DeepBlue}解答}\par
\begin{center}\begin{minipage}{\linewidth}\centering
\begin{tikzpicture}[scale=.9]
 \draw[axis] (0,2)--(0,-2.1) node[below]{转轴};
 \draw[very thick] (0,0) ellipse (1.35 and 1.8);
 \coordinate (P) at ({1.35*sin(32)},{-1.8*cos(32)});
 \draw[guide] (0,0)--(P);\fill (0,0) circle(1.5pt) node[left]{$O$};
 \fill (P) circle(2.4pt) node[right=5pt]{$m$};
 \draw[DeepBlue,->] (0,-.68) arc[start angle=-90,end angle=-58,x radius=.51,y radius=.68];
 \node at (.36,-.82) {$\theta$};
 \draw[TextGray,->] ({.75*cos(160)},{.65+.20*sin(160)}) arc[start angle=160,end angle=20,x radius=.75,y radius=.20];
 \node at (.90,.77) {$\varphi$};
\end{tikzpicture}
\bookcaption{可绕竖直直径转动的圆环与环上小珠}\label{advA1:fig:hoop}
\end{minipage}\end{center}\end{minipage}

如图\ref{advA1:fig:hoop}，
\[
 L=\frac12ma^2\dot\theta^2
 +\frac12(I+ma^2\sin^2\theta)\dot\varphi^2
 +mga\cos\theta .
\]
$\varphi$ 循环，
$J=(I+ma^2\sin^2\theta)\dot\varphi$。作罗斯约化：
\[
 \mathcal R=\frac12ma^2\dot\theta^2+mga\cos\theta
 -\frac{J^2}{2(I+ma^2\sin^2\theta)}.
\]
在 $\theta=0$ 附近保留二阶项，
\[
 \mathcal R=\text{常量}+\frac12ma^2\dot\theta^2
 -\frac12ma^2\left(\frac ga-\frac{J^2}{I^2}\right)\theta^2+\order{\theta^4}.
\]
因而
\[
 \boxed{\omega^2=\frac ga-\left(\frac JI\right)^2},
 \qquad \frac{|J|}I<\sqrt{\frac ga}.
\]
当 $|J|/I=\sqrt{g/a}$ 时二阶恢复力消失，但有效势的首个非零增量为
\[
 U_{\rm eff}(\theta)-U_{\rm eff}(0)
 =mga\left(\frac18+\frac{ma^2}{2I}\right)\theta^4+\order{\theta^6}>0
\]
（充分小的非零 $\theta$），所以最低点仍为稳定平衡，只是线性小振动频率为零，有限小振幅的周期须由非线性项决定。因此稳定范围含等号；上式严格不等号对应非零的线性回复频率。超过阈值时最低点变为不稳定点。
\end{advsolution}

\subsection{小振动、傅里叶变换与周期系统}

在稳定平衡点 $q_\alpha=0$ 附近，将动能和势能展开到二阶：
\begin{equation}\label{advA1:eq:small-L}
 L=\frac12\dot{\vect q}^{\mathsf T}\mathsf T\dot{\vect q}
 -\frac12\vect q^{\mathsf T}\mathsf V\vect q .
\end{equation}
$\mathsf T$ 是正定的惯性矩阵，$\mathsf V$ 是对称的弹性矩阵。被忽略的项至少含有三个小量，可用微扰迭代逐阶修正。

\begin{advdefinition}{线性变换}\label{advA1:def:linear-map}
向量空间上满足
$F(\lambda\vect x+\mu\vect y)=\lambda F(\vect x)+\mu F(\vect y)$
的映射称为线性变换。
\end{advdefinition}

\begin{advtheorem}{线性变换的矩阵表示}\label{advA1:thm:linear-matrix}
在选定基底后，每个有限维线性变换唯一对应一个矩阵 $\mathsf F$，使
$\vect y=F(\vect x)$ 等价于 $[\vect y]=\mathsf F[\vect x]$。
\end{advtheorem}
证明：只需把每个基矢的像作为矩阵的一列；线性性保证它对任意向量成立，基底展开的唯一性保证矩阵唯一。

\begin{advdefinition}{对称二次型}\label{advA1:def:quadratic-form}
若 $S(\vect x,\vect y)$ 对两个变量均线性，且
$S(\vect x,\vect y)=S(\vect y,\vect x)$，则称其为对称二次型对应的双线性型。
\end{advdefinition}

\begin{advtheorem}{二次型的对称矩阵表示}\label{advA1:thm:quadratic-matrix}
在选定基底后，存在唯一对称矩阵 $\mathsf S$，使
\[
 S(\vect x,\vect y)=\vect x^{\mathsf T}\mathsf S\vect y.
\]
\end{advtheorem}
证明：令 $S_{ij}=S(\vect e_i,\vect e_j)$，用双线性展开即得矩阵表示；对称性给出 $S_{ij}=S_{ji}$。

\begin{advdefinition}{合同变换}\label{advA1:def:congruence}
对非奇异矩阵 $\mathsf P$，
\[
 \mathsf S\longmapsto \mathsf P^{\mathsf T}\mathsf S\mathsf P
\]
称为对称矩阵的合同变换。它是二次型在坐标变换 $\vect q=\mathsf P\vect Q$ 下的自然变换。
\end{advdefinition}

\begin{advtheorem}{惯性矩阵与弹性矩阵的同时对角化}\label{advA1:thm:normal-modes}
若 $\mathsf T$ 对称正定、$\mathsf V$ 对称，则存在非奇异矩阵 $\mathsf A$，使
\begin{equation}\label{advA1:eq:sim-diag}
 \mathsf A^{\mathsf T}\mathsf T\mathsf A=\mathsf I,
 \qquad
 \mathsf A^{\mathsf T}\mathsf V\mathsf A
 =\operatorname{Diag}(\omega_1^2,\ldots,\omega_n^2).
\end{equation}
\end{advtheorem}
证明：因 $\mathsf T$ 正定，存在可逆平方根 $\mathsf T^{1/2}$。矩阵
$\mathsf K=\mathsf T^{-1/2}\mathsf V\mathsf T^{-1/2}$ 对称，因此可被正交矩阵 $\mathsf O$ 对角化。取
$\mathsf A=\mathsf T^{-1/2}\mathsf O$，立即得式\eqref{advA1:eq:sim-diag}。令 $\vect q=\mathsf A\vect Q$，式\eqref{advA1:eq:small-L}化为
\[
 L=\frac12\sum_j(\dot Q_j^2-\omega_j^2Q_j^2).
\]
$\omega_j^2>0$ 表示稳定振动，$\omega_j^2=0$ 表示零模，$\omega_j^2<0$ 表示指数不稳定。

\begin{advexamplebox}{刚体复摆的小振动}\label{advA1:ex:physical-pendulum}
\emph{题源说明：原资料仅给主题或注明课堂描述；以下条件为依据主题所作的教学补全。}\par
刚体质量为 $M$，质心到水平固定轴的距离为 $a$，对该轴的转动惯量为 $I$。求竖直向下稳定位置附近的周期。

\begin{qualityfigure}
\begin{tikzpicture}[>=Latex,font=\small]
\fill(0,2.5)circle(2pt)node[above]{$O$};\draw[dashed](0,2.5)--(0,-1);\coordinate(P)at(1.5,.1);\draw[thick](0,2.5)--(P);\fill(P)circle(3pt)node[right]{$G$};\draw(0,1.6)arc(-90:-58:.9)node[midway,below right]{$\theta$};\node at(3,1.2){$OG=a$};\node at(3,.4){轴惯量 $I$，质量 $M$};
\end{tikzpicture}
\figcap{固定轴与偏离铅垂线的刚体质心（题设几何）}\label{advA1:fig:quality-physical-pendulum-q}
\end{qualityfigure}
\end{advexamplebox}
\begin{advsolution}
\begin{qualityfigure}
\begin{tikzpicture}[>=Latex,font=\small]
\fill(0,2.5)circle(2pt)node[above]{$O$};\draw[dashed](0,2.5)--(0,-1);\coordinate(P)at(1.5,.1);\draw[thick](0,2.5)--(P);\fill(P)circle(3pt)node[right]{$G$};\draw(0,1.6)arc(-90:-58:.9)node[midway,below right]{$\theta$};\node at(3,1.2){$OG=a$};\node at(3,.4){轴惯量 $I$，质量 $M$};\draw[->,Gold](P)--++(0,-1)node[below]{$Mg$};
\end{tikzpicture}
\figcap{固定轴与偏离铅垂线的刚体质心（解答辅助图）}\label{advA1:fig:quality-physical-pendulum-a}
\end{qualityfigure}
图~\ref{advA1:fig:quality-physical-pendulum-a}中，质心 $G$ 距支点轴为 $a$，重力矩的力臂是 $a\sin\theta$；小角展开后用刚体对该固定轴的惯量 $I$，不能用质点式 $Ma^2$ 代替。

以偏角 $\theta$ 为坐标，
\[
 L=\frac12I\dot\theta^2+Mga\cos\theta.
\]
精确方程是 $I\ddot\theta+Mga\sin\theta=0$。小角度下
\[
 I\ddot\theta+Mga\theta=0,
 \qquad \boxed{\omega=\sqrt{\frac{Mga}{I}}},
 \qquad \boxed{T_0=2\pi\sqrt{\frac I{Mga}}}.
\]
若 $I=Ma^2$，回到长度 $a$ 的单摆；若 $a\to0$，重力恢复力矩消失，周期趋于无穷，两个极限均正确。
\end{advsolution}

\subsubsection{陀螺型速度耦合与傅科摆}

考虑二维谐振子
\[
 T=\frac m2(\dot x^2+\dot y^2),\qquad
 V_0=\frac k2(x^2+y^2),
\]
并受到速度型力
\[
 F_x=-m\Omega\dot y,\qquad F_y=m\Omega\dot x.
\]
它可由含速势
\[
 V_g=\frac{m\Omega}{2}(x\dot y-y\dot x)
\]
描述。令 $z=x+\ii y$，方程化为
\begin{equation}\label{advA1:eq:foucault-complex}
 \ddot z-\ii\Omega\dot z+\omega^2z=0,
 \qquad \omega^2=k/m.
\end{equation}
作变换 $z=e^{\ii\Omega t/2}\zeta$，得
\[
 \ddot\zeta+\omega_0^2\zeta=0,
 \qquad \omega_0^2=\omega^2+\frac{\Omega^2}{4}.
\]
两个圆偏振模的频率为 $\omega_0\pm\Omega/2$，线偏振叠加后的振动平面以 $\Omega/2$ 进动。对地球上的傅科摆，取 $x$ 向东、$y$ 向北、竖直向上为正，地球自转的竖直分量为 $\Omega_E\sin\lambda$。科里奥利力的水平分量为 $F_x=2m\Omega_E\sin\lambda\,\dot y$、$F_y=-2m\Omega_E\sin\lambda\,\dot x$，故按本节的力和复坐标约定，应代入
$\Omega=-2\Omega_E\sin\lambda$，振动平面的有符号进动角速度为 $-\Omega_E\sin\lambda$；北半球从上方看为顺时针，进动速率的大小为 $|\Omega_E\sin\lambda|$。

\begin{advtheorem}{周期函数的傅里叶级数}\label{advA1:thm:fourier-series}
对周期 $l$ 的光滑函数，
\[
 f(x)=\sum_{n=-\infty}^{\infty}c_n e^{\ii2\pi nx/l},
 \qquad c_n=\frac1l\int_0^l f(x)e^{-\ii2\pi nx/l}\dd x.
\]
\end{advtheorem}
证明要点是函数 $e^{\ii2\pi nx/l}$ 在 $[0,l]$ 上正交：
\[
 \frac1l\int_0^l e^{\ii2\pi(n-m)x/l}\dd x=\delta_{nm}.
\]
将展开式两边与第 $m$ 个基函数的共轭相乘并积分，即得 $c_m$。光滑性保证系数足够快衰减并使级数按所需意义收敛。

\begin{advtheorem}{平方可积函数的傅里叶变换}\label{advA1:thm:fourier-transform}
采用对称归一化，
\[
 g(k)=\frac1{\sqrt{2\pi}}\int_{-\infty}^{\infty}f(x)e^{-\ii kx}\dd x,
 \qquad
 f(x)=\frac1{\sqrt{2\pi}}\int_{-\infty}^{\infty}g(k)e^{\ii kx}\dd k.
\]
对 $f\in L^2(\mathbb R)$，等式按 $L^2$ 意义成立。
\end{advtheorem}
可先在长区间 $[-L/2,L/2]$ 上作傅里叶级数，再令
$L\to\infty$。此时离散波数间距 $\Delta k=2\pi/L\to0$，级数的黎曼和趋于积分，得到上式。更严格的表述依赖平方可积函数空间的完备性。

连续/离散、周期/非周期的四种组合分别对应傅里叶变换、傅里叶级数、离散时间傅里叶变换和离散傅里叶变换。它们的共同核心都是用平移对称性的本征函数实现解耦。

\subsubsection{一维单原子链}

在平衡位置 $x_n^{(0)}=nl$ 上放置质量 $m$ 的质点，纵向位移为 $u_n$，最近邻弹簧常量为 $\kappa$。
\[
 L=\sum_n\left[\frac12m\dot u_n^2-\frac12\kappa(u_{n+1}-u_n)^2\right].
\]
对离散坐标作变换
\[
 \phi(k,t)=\sum_{n=-\infty}^{\infty}u_n(t)e^{-\ii knl},
 \quad
 u_n(t)=\frac l{2\pi}\int_{-\pi/l}^{\pi/l}
 \phi(k,t)e^{\ii knl}\dd k.
\]
由于 $\phi(k+2\pi/l)=\phi(k)$，只需保留第一布里渊区 $[-\pi/l,\pi/l]$。代入并利用正交性，各个 $k$ 模彼此解耦，色散关系为
\begin{equation}\label{advA1:eq:lattice-dispersion}
 \boxed{\omega(k)=2\sqrt{\frac\kappa m}
 \left|\sin\frac{kl}{2}\right|}.
\end{equation}
长波极限 $|kl|\ll1$ 下，$\omega\simeq l\sqrt{\kappa/m}|k|$，群速度趋于
$l\sqrt{\kappa/m}$；在布里渊区边界 $k=\pm\pi/l$ 处群速度为零。

\begin{advexamplebox}{用极坐标求傅科摆进动}\label{advA1:ex:foucault-polar}
对 $L=m(\dot x^2+\dot y^2)/2-m\omega^2(x^2+y^2)/2-m\Omega(x\dot y-y\dot x)/2$，改用极坐标 $(r,\theta)$，求共转线偏振解的进动速率，并说明一般初态及静止释放时的区别。

\begin{qualityfigure}
\begin{tikzpicture}[>=Latex,font=\small]
\draw[->](-2.5,0)--(3,0)node[right]{$x$};\draw[->](0,-2)--(0,2.4)node[above]{$y$};\draw[DeepBlue,rotate=25](0,0)ellipse(2 and .5);\draw[dashed,rotate=25](-2.3,0)--(2.3,0)node[right]{共转长轴};\draw(.8,0)arc(0:25:.8)node[midway,right]{$\theta$};
\end{tikzpicture}
\figcap{共转坐标中椭圆的主轴与极角（题设几何）}\label{advA1:fig:quality-foucault-polar-q}
\end{qualityfigure}
\end{advexamplebox}
\begin{advsolution}
\begin{qualityfigure}
\begin{tikzpicture}[>=Latex,font=\small]
\draw[->](-2.5,0)--(3,0)node[right]{$x$};\draw[->](0,-2)--(0,2.4)node[above]{$y$};\draw[DeepBlue,rotate=25](0,0)ellipse(2 and .5);\draw[dashed,rotate=25](-2.3,0)--(2.3,0)node[right]{共转长轴};\draw(.8,0)arc(0:25:.8)node[midway,right]{$\theta$};\node[align=center]at(0,-2.6){共转椭圆的主轴以 $\Omega/2$ 转动\\一般初态不必是直线振动};
\end{tikzpicture}
\figcap{共转坐标中椭圆的主轴与极角（解答辅助图）}\label{advA1:fig:quality-foucault-polar-a}
\end{qualityfigure}
图~\ref{advA1:fig:quality-foucault-polar-a}中，椭圆的轴角与质点的瞬时极角不同；只有共转系中的特定线偏振解才有 $J=0$，一般初态应保留椭圆的横向运动。

因 $x\dot y-y\dot x=r^2\dot\theta$，
\[
 L=\frac m2(\dot r^2+r^2\dot\theta^2)-\frac k2r^2
 -\frac{m\Omega}{2}r^2\dot\theta,
 \qquad J=mr^2\left(\dot\theta-\frac\Omega2\right).
\]
令共转角 $\psi=\theta-\Omega t/2$，则 $J=mr^2\dot\psi$。
在共转坐标中选取线偏振振动才有 $J=0$；除穿过原点时极角的 $\pi$ 跳变外，
振动直线满足 $\dot\theta=\Omega/2$。
此时用可正可负的直线坐标 $q$ 代替非负半径，得到
\[
 \ddot q+\omega_0^2q=0,\qquad \omega_0^2=\omega^2+\Omega^2/4.
\]
一般初态在共转系中形成固定椭圆，该椭圆的轴随坐标系以 $\Omega/2$ 转动。
例如原坐标中从 $z(0)=a,\dot z(0)=0$ 静止释放，有
\[
 J=-\frac12m\Omega a^2\ne0,\qquad
 z(t)=a\e^{\ii\Omega t/2}
 \left[\cos\omega_0t-\frac{\ii\Omega}{2\omega_0}\sin\omega_0t\right].
\]
在 $|\Omega|\ll\omega_0$ 时椭圆很窄，才可近似称为单一振动平面进动。
因此 $J=0$ 是特定线偏振解的条件，不能由任意“静止释放”推出。

\end{advsolution}

\begin{advexamplebox}{无限一维电阻网络}\label{advA1:ex:resistor-chain}
\emph{题源说明：原资料仅给主题或注明课堂描述；以下条件为依据主题所作的教学补全。}\par
无限多个节点按整数 $n$ 排列，相邻节点之间接电阻 $R$，每个节点再通过电阻 $r$ 接公共地。求 $n=0$ 节点对地的等效电阻。

\begin{qualityfigure}
\begin{tikzpicture}[>=Latex,font=\small]
\foreach\i in{-2,-1,0,1,2}{\fill(1.5*\i,0)circle(2pt)node[above]{$\i$};\draw(1.5*\i,0)--(1.5*\i,-.3)--++(.12,-.12)--++(-.24,-.2)--++(.24,-.2)--++(-.24,-.2)--++(.12,-.12)--(1.5*\i,-1.5);\node[right]at(1.5*\i,-.8){$r$};}\draw(-3,-1.5)--(3,-1.5)node[right]{公共地};\foreach\i in{-2,-1,0,1}{\draw(1.5*\i,0)--++(.4,0)--++(.1,.15)--++(.2,-.3)--++(.2,.3)--++(.2,-.3)--++(.1,.15)--++(.3,0);\node at(1.5*\i+.75,.5){$R$};}\node at(-3.7,0){$\cdots$};\node at(3.7,0){$\cdots$};
\end{tikzpicture}
\figcap{无限电阻链的节点和对地支路（题设几何）}\label{advA1:fig:quality-resistor-chain-q}
\end{qualityfigure}
\end{advexamplebox}
\begin{advsolution}
\begin{qualityfigure}
\begin{tikzpicture}[>=Latex,font=\small]
\foreach\i in{-2,-1,0,1,2}{\fill(1.5*\i,0)circle(2pt)node[above]{$\i$};\draw(1.5*\i,0)--(1.5*\i,-.3)--++(.12,-.12)--++(-.24,-.2)--++(.24,-.2)--++(-.24,-.2)--++(.12,-.12)--(1.5*\i,-1.5);\node[right]at(1.5*\i,-.8){$r$};}\draw(-3,-1.5)--(3,-1.5)node[right]{公共地};\foreach\i in{-2,-1,0,1}{\draw(1.5*\i,0)--++(.4,0)--++(.1,.15)--++(.2,-.3)--++(.2,.3)--++(.2,-.3)--++(.1,.15)--++(.3,0);\node at(1.5*\i+.75,.5){$R$};}\node at(-3.7,0){$\cdots$};\node at(3.7,0){$\cdots$};\draw[Gold,thick](-.3,-1.7)rectangle(.3,.7);\node at(0,-2.3){在 $n=0$ 注入电流，两侧衰减解须在无穷远有界};
\end{tikzpicture}
\figcap{无限电阻链的节点和对地支路（解答辅助图）}\label{advA1:fig:quality-resistor-chain-a}
\end{qualityfigure}
图~\ref{advA1:fig:quality-resistor-chain-a}中，每个整数节点均通过 $r$ 接地，相邻节点通过 $R$ 相连；在零号节点注入电流后，两侧节点电势必须满足同一差分方程及无穷远有界条件。

向 $0$ 节点注入电流 $I$，节点电势为 $V_n$。对所有 $n$，
\[
 \frac{V_n}{r}+\frac{2V_n-V_{n-1}-V_{n+1}}R=I\delta_{n0}.
\]
作离散傅里叶变换，
\[
 V(k)=\frac I{r^{-1}+2R^{-1}(1-\cos k)}.
\]
因此
\[
 \frac{V_0}{I}=\frac1{2\pi}\int_{-\pi}^{\pi}
 \frac{\dd k}{r^{-1}+2R^{-1}(1-\cos k)}.
\]
% A1CompletenessInput20261024: standard-integral
% A1Completeness20261024: resistor-integral
下面把最后的积分算清，避免把电路答案藏在“标准积分”中。对$a>|b|$，设
\[
 J(a,b)=\frac1{2\pi}\int_{-\pi}^{\pi}\frac{\dd k}{a-b\cos k}.
\]
令$t=\tan(k/2)$，则$\cos k=(1-t^2)/(1+t^2)$、$\dd k=2\dd t/(1+t^2)$。
$k$从$-\pi$走到$\pi$时，$t$从$-\infty$走到$+\infty$，故
\[
 J=\frac1\pi\int_{-\infty}^{\infty}
 \frac{\dd t}{(a-b)+(a+b)t^2}.
\]
因$a-b,a+b$都为正，再令$u=t\sqrt{(a+b)/(a-b)}$：
\[
 J=\frac1{\pi\sqrt{a^2-b^2}}
 \int_{-\infty}^{\infty}\frac{\dd u}{1+u^2}
 =\frac{\arctan u|_{-\infty}^{\infty}}{\pi\sqrt{a^2-b^2}}
 =\frac1{\sqrt{a^2-b^2}}.
\]
\begin{qualityfigure}
\begin{tikzpicture}[>=Latex,font=\small]
 \begin{scope}
  \draw[->](-2.2,1)--(2.3,1)node[right]{$k$};
  \foreach\x/\lab in{-1.8/-\pi,0/0,1.8/\pi}{\draw(\x,.92)--(\x,1.08);\node[above]at(\x,1.1){$\lab$};}
  \draw[->](-2.2,-.7)--(2.3,-.7)node[right]{$t$};
  \node[below]at(-1.8,-.8){$-\infty$};\node[below]at(0,-.8){$0$};\node[below]at(1.8,-.8){$+\infty$};
  \foreach\x in{-1.8,0,1.8}{\draw[->,DeepBlue](\x,.65)--(\x,-.35);}
  \node at(0,-1.7){$t=\tan(k/2)$};
 \end{scope}
 \begin{scope}[shift={(6.2,-.4)}]
  \draw[->](-2.1,0)--(2.3,0)node[right]{$u$};
  \draw[->](0,0)--(0,1.65)node[above]{$(1+u^2)^{-1}$};
  \fill[PaleBlue,domain=-2:2,samples=80](-2,0)--plot(\x,{1/(1+\x*\x)})--(2,0)--cycle;
  \draw[DeepBlue,thick,domain=-2:2,samples=80]plot(\x,{1/(1+\x*\x)});
  \node at(0,-.65){全实轴曲线下的面积为$\pi$};
 \end{scope}
\end{tikzpicture}
\figcap{半角代换将一个周期的积分化为全实轴上的反正切积分；右图的两端尾部延伸到无穷远。}\label{advA1:fig:completeness-resistor-integral}
\end{qualityfigure}
图~\ref{advA1:fig:completeness-resistor-integral}左侧的三个对应点固定了积分区间，右侧面积由反正切的两个端点极限得到。本题取$a=r^{-1}+2R^{-1}$、$b=2R^{-1}$，电阻为正保证$a-b=r^{-1}>0$，代入得


\[
 \boxed{R_{\rm eq}=\frac1{\sqrt{r^{-2}+4/(rR)}}
 =\frac r{\sqrt{1+4r/R}}}.
\]
$R\to\infty$ 时各节点断开，$R_{\rm eq}\to r$；$r\to\infty$ 时整个网络没有到地通路，$R_{\rm eq}\to\infty$，极限检查通过。
\end{advsolution}

\begin{advexamplebox}{六原子环的简正模}\label{advA1:ex:six-atom-ring}
\emph{题源说明：原资料仅给主题或注明课堂描述；以下条件为依据主题所作的教学补全。}\par
六个质量均为 $m$ 的原子排成闭环，只考虑沿环方向的小位移 $u_j$，相邻原子之间的有效弹簧常量为 $\kappa$，$u_{j+6}=u_j$。求全部简正频率及简并度。

\begin{qualityfigure}
\begin{tikzpicture}[>=Latex,font=\small]
\foreach\i in{0,...,5}{\coordinate(P\i)at({60*\i}:1.8);\fill(P\i)circle(2.5pt)node[below right]{$j=\i$};\draw[->,Gold](P\i)--++({60*\i+90}:.55)node[above]{$u_\i$};}\draw[DeepBlue,thick](P0)--(P1)--(P2)--(P3)--(P4)--(P5)--cycle;\node[align=center]at(0,0){$m_j=m$\\$\kappa$};\node at(0,-2.7){沿环方向的小位移，$u_{j+6}=u_j$};
\end{tikzpicture}
\figcap{六原子的周期边界与切向小位移（题设几何）}\label{advA1:fig:quality-six-atom-ring-q}
\end{qualityfigure}
\end{advexamplebox}
\AdvHideHeadingtrue
\begin{advsolution}题图~\ref{advA1:fig:quality-six-atom-ring-q}中，相邻原子之间的力取决于位移差，首尾相接给出 $u_{j+6}=u_j$；允许波数由这个周期边界量子化，不能使用自由端链的边界条件。

\noindent\begin{minipage}{\linewidth}{\bfseries\color{DeepBlue}解答}\par
\begin{center}\begin{minipage}{\linewidth}\centering
\begin{tikzpicture}[scale=.92]
 \foreach \j in {0,...,5}{
   \coordinate (P\j) at ({90-60*\j}:1.42);
 }
 \foreach \j [evaluate=\j as \k using {int(mod(\j+1,6))}] in {0,...,5}{
   \draw[spring] (P\j)--(P\k);
 }
 \foreach \j in {0,...,5}{
   \fill (P\j) circle(3pt);
   \node at ($1.28*(P\j)$) {$u_{\j}$};
 }
 \draw[->,DeepBlue] (0,2.15) arc[start angle=90,end angle=25,radius=2.15];
 \node[above right] at (2.15,1.85) {正向};
\end{tikzpicture}
\bookcaption{六原子环中的周期边界与最近邻耦合}\label{advA1:fig:ring}
\end{minipage}\end{center}\end{minipage}

如图\ref{advA1:fig:ring}，方程为
\[
 m\ddot u_j=\kappa(u_{j+1}+u_{j-1}-2u_j).
\]
取 $u_j=Ae^{\ii(kj-\omega t)}$，周期边界给出
$k_n=2\pi n/6$，$n=0,\ldots,5$，而
\[
 \omega_n=2\sqrt{\frac\kappa m}\left|\sin\frac{k_n}{2}\right|.
\]
所以频谱为
\[
 \boxed{0\ (1),\quad \sqrt{\kappa/m}\ (2),\quad
 \sqrt{3\kappa/m}\ (2),\quad 2\sqrt{\kappa/m}\ (1)}.
\]
括号中是简并度。零模对应六个原子沿环同向等量移动，弹簧不变形；$n$ 与 $6-n$ 构成反向传播的一对简并模。
\end{advsolution}

\section{刚体}

\subsection{质量几何与矩生成函数}

对密度分布 $\rho(x,y,z)$，定义笛卡尔矩
\begin{equation}\label{advA1:eq:cartesian-moment}
 M_{pqr}=\iiint \rho(x,y,z)x^py^qz^r\dd V .
\end{equation}
零阶矩是总质量，一阶矩确定质心，二阶矩决定转动惯量张量。更高阶矩则继续编码质量分布的细节。

\begin{advexamplebox}{矩的平移定理}\label{advA1:ex:moment-translation}
设相对原点 $O$ 的矩为 $M_{pqr}$。把物体整体平移 $\vect a=(a,b,c)$ 后，证明新矩为
\[
 M'_{pqr}=\sum_{i=0}^{p}\sum_{j=0}^{q}\sum_{k=0}^{r}
 \binom pi\binom qj\binom rk
 a^{p-i}b^{q-j}c^{r-k}M_{ijk}.
\]
\end{advexamplebox}
\begin{advsolution}
平移后同一质量元的坐标由 $(x,y,z)$ 变为 $(x+a,y+b,z+c)$，而 $\rho\dd V$ 不变。代入定义\eqref{advA1:eq:cartesian-moment}，对三个因子分别用二项式定理：
\[
 (x+a)^p(y+b)^q(z+c)^r
 =\sum_{ijk}\binom pi\binom qj\binom rk
 x^iy^jz^k a^{p-i}b^{q-j}c^{r-k}.
\]
逐项积分即得所证公式。对一阶矩，它给出
$\vect M'_1=\vect M_1+M\vect a$；对二阶矩，它给出熟悉的平行轴项。
\end{advsolution}

引入矩生成函数
\begin{equation}\label{advA1:eq:moment-generator}
 P(s_1,s_2,s_3)=\iiint\rho(x,y,z)
 e^{s_1x+s_2y+s_3z}\dd V .
\end{equation}

\begin{advtheorem}{笛卡尔矩的生成}\label{advA1:thm:moment-generator}
\[
 M_{pqr}=\left.
 \frac{\partial^{p+q+r}P}
 {\partial s_1^p\partial s_2^q\partial s_3^r}
 \right|_{(0,0,0)}.
\]
\end{advtheorem}
证明：在积分号内对 $s_1,s_2,s_3$ 求导，每次分别带出一个 $x,y,z$；最后令 $s_1=s_2=s_3=0$，指数因子变为 $1$，正好得到式\eqref{advA1:eq:cartesian-moment}。因而在收敛条件满足时，全部矩就是 $P$ 在原点处的多元泰勒系数。

\begin{advexamplebox}{均匀正方形薄板的前四级矩}\label{advA1:ex:square-moments}
在 $0\le x\le l, 0\le y\le l$ 上面密度 $\sigma$ 均匀分布。求总阶数 $p+q=0,1,2,3$ 的全部平面笛卡尔矩 $M_{pq}$。
\end{advexamplebox}
\begin{advsolution}
\[
 M_{pq}=\sigma\int_0^l x^p\dd x\int_0^l y^q\dd y
 =\frac{\sigma l^{p+q+2}}{(p+1)(q+1)}.
\]
因此
\begin{align*}
 &M_{00}=\sigma l^2,\\
 &M_{10}=M_{01}=\frac12\sigma l^3,\\
 &M_{20}=M_{02}=\frac13\sigma l^4,
 \qquad M_{11}=\frac14\sigma l^4,\\
 &M_{30}=M_{03}=\frac14\sigma l^5,
 \qquad M_{21}=M_{12}=\frac16\sigma l^5.
\end{align*}
单位为质量乘长度的 $p+q$ 次方，与定义相符。
\end{advsolution}

\begin{advexamplebox}{卷积与矩生成函数}\label{advA1:ex:convolution}
定义三维卷积
\[
 \gamma(\vect r)=(\alpha*\beta)(\vect r)
 =\int\alpha(\vect r')\beta(\vect r-\vect r\,')\dd^3r'.
\]
若三个分布的矩生成函数分别为 $P,Q,R$，矩分别为 $L,M,N$，证明 $R=PQ$，并写出 $N_{pqr}$。
\end{advexamplebox}
\begin{advsolution}
对卷积作变量代换 $\vect u=\vect r-\vect r\,'$，
\begin{align*}
 R(\vect s)
 &=\int\dd^3r\int\dd^3r'\,\alpha(\vect r')
 \beta(\vect r-\vect r\,')e^{\vect s\cdot\vect r}\\
 &=\left[\int\alpha(\vect r')e^{\vect s\cdot\vect r'}\dd^3r'\right]
 \left[\int\beta(\vect u)e^{\vect s\cdot\vect u}\dd^3u\right]
 =P(\vect s)Q(\vect s).
\end{align*}
对乘积求 $p+q+r$ 阶偏导并用多重莱布尼茨公式：
\[
 N_{pqr}=\sum_{i=0}^{p}\sum_{j=0}^{q}\sum_{k=0}^{r}
 \binom pi\binom qj\binom rk
 L_{ijk}M_{p-i,q-j,r-k}.
\]
这一关系与“独立随机变量之和的矩生成函数相乘”是同一个结构。
\end{advsolution}

\subsection{欧拉运动学方程}

对定点转动，取初始时刻的固连正交基
$(\vect e_1,\vect e_2,\vect e_3)$。设形变映射为 $\vect R(\vect r,t)$，并定义
$\vect\varepsilon_i=\partial\vect R/\partial x_i$。

\begin{advdefinition}{刚体}\label{advA1:def:rigid-body}
刚体的形变映射在每一时刻都保长且保向：
\[
 \vect\varepsilon_i\cdot\vect\varepsilon_j=\delta_{ij},
 \qquad (\vect\varepsilon_1\times\vect\varepsilon_2)
 \cdot\vect\varepsilon_3=+1.
\]
\end{advdefinition}
固定原点后，这种映射由特殊正交矩阵 $\mathsf A(t)\in SO(3)$ 描述：
\[
 \vect R=\mathsf A(t)\vect r,
 \qquad \mathsf A^{\mathsf T}\mathsf A=\mathsf I,
 \qquad \det\mathsf A=1.
\]
“保长”给出正交性，“保向”排除行列式为 $-1$ 的镜像变换。连续的刚体运动不能穿过 $\det\mathsf A=0$，因而也不能在运动中突然改变左右手性。

\begin{advexamplebox}{三条坐标轴的转动矩阵}\label{advA1:ex:axis-rotations}
写出绕 $x,y,z$ 轴主动旋转 $\theta$ 的矩阵，检验不同轴转动是否对易，并讨论能否用一组实坐标同时将所有转动对角化。
\end{advexamplebox}
\begin{advsolution}
\[
 R_x=\begin{pmatrix}1&0&0\\0&c&-s\\0&s&c\end{pmatrix},
 \quad
 R_y=\begin{pmatrix}c&0&s\\0&1&0\\-s&0&c\end{pmatrix},
 \quad
 R_z=\begin{pmatrix}c&-s&0\\s&c&0\\0&0&1\end{pmatrix},
\]
其中 $c=\cos\theta, s=\sin\theta$。例如对矢量 $\vect e_x$，
$R_xR_y\vect e_x\ne R_yR_x\vect e_x$，所以一般 $R_xR_y\ne R_yR_x$。无穷小角时对易子是二阶小量，这也解释了小转角下“转角矢量近似可相加”只是一阶结论。

一般实转动的本征值为 $1,e^{\pm\ii\theta}$，除 $\theta=0,\pi$ 等特例外不能在实数域对角化。更不可能用同一组基同时对角化全部 $SO(3)$ 元素，因为可同时对角化的矩阵必彼此对易。
\end{advsolution}

定义刚体坐标中的角速度矩阵
\begin{equation}\label{advA1:eq:omega-matrix}
 \bm{\Omega}=\mathsf{A}^{\mathsf T}\dot{\mathsf{A}}.
\end{equation}

\begin{advtheorem}{$SO(3)$ 的生成元是反对称矩阵}\label{advA1:thm:so3-generator}
\[
 \bm{\Omega}^{\mathsf T}=-\bm{\Omega}.
\]
\end{advtheorem}
对 $\mathsf A^{\mathsf T}\mathsf A=\mathsf I$ 求导：
$\dot{\mathsf A}^{\mathsf T}\mathsf A+\mathsf A^{\mathsf T}\dot{\mathsf A}=0$。第一项正是第二项的转置，因而得证。三维反对称矩阵只有三个独立分量，可与伪矢量 $\vect\omega=(\omega_1,\omega_2,\omega_3)$ 对应：
\[
 [\vect\omega]_\times=
 \begin{pmatrix}
 0&-\omega_3&\omega_2\\
 \omega_3&0&-\omega_1\\
 -\omega_2&\omega_1&0
 \end{pmatrix},
 \qquad [\vect\omega]_\times\vect r=\vect\omega\times\vect r.
\]
所以刚体上的点在刚体坐标中满足 $\vect v=\vect\omega\times\vect r$。

\begin{advexamplebox}{矩阵指数与有限转动}\label{advA1:ex:matrix-exponential}
记 $[\vect n]_\times$ 为单位矢量 $\vect n$ 的叉乘矩阵。证明
\[
 R_{\vect n}(\theta)=e^{\theta[\vect n]_\times}
 =\lim_{N\to\infty}\left(\mathsf I+
 \frac\theta N[\vect n]_\times\right)^N,
\]
并得到罗德里格公式。
\end{advexamplebox}
\begin{advsolution}
因 $[\vect n]_\times^2=\vect n\vect n^{\mathsf T}-\mathsf I$，且
$[\vect n]_\times^3=-[\vect n]_\times$，指数级数的奇、偶次项可分别求和：
\[
 e^{\theta[\vect n]_\times}
 =\mathsf I+\sin\theta[\vect n]_\times
 +(1-\cos\theta)[\vect n]_\times^2.
\]
作用于任意 $\vect r$ 便得
\[
 R_{\vect n}(\theta)\vect r
 =\vect r\cos\theta+(\vect n\times\vect r)\sin\theta
 +\vect n(\vect n\cdot\vect r)(1-\cos\theta),
\]
即罗德里格公式。极限形式是矩阵指数的标准定义，表明有限转动可看成无穷多个无穷小转动的有序乘积。
\end{advsolution}

\begin{advexamplebox}{角速度矩阵与伪矢量的同变换}\label{advA1:ex:omega-transform}
对 $\mathsf O\in SO(3)$，证明
\[
 [\vect\omega']_\times=\mathsf O[\vect\omega]_\times\mathsf O^{-1}
 \quad\Longleftrightarrow\quad
 \vect\omega'=\mathsf O\vect\omega .
\]
\end{advexamplebox}
\begin{advsolution}
对任意矢量 $\vect x$，
\[
 \mathsf O[\vect\omega]_\times\mathsf O^{-1}\vect x
 =\mathsf O(\vect\omega\times\mathsf O^{-1}\vect x)
 =(\mathsf O\vect\omega)\times\vect x.
\]
第二个等号用了 $\det\mathsf O=+1$ 的保向性。因为等式对所有 $\vect x$ 成立，所以两个叉乘矩阵相等。反向推理同理。若把 $SO(3)$ 放宽为包含镜反射的 $O(3)$，就会多出 $\det\mathsf O$，这正是角速度称为伪矢量的原因。
\end{advsolution}

\subsubsection{欧拉角与角速度}

采用 $Z-X_1-Z_2$ 约定：先绕定系 $Z$ 轴转进动角 $\varphi$，再绕节线 $X_1$ 转章动角 $\theta$，最后绕刚体轴 $Z_2=z$ 转自转角 $\psi$。
\begin{center}\begin{minipage}{\linewidth}\centering
\begin{tikzpicture}[scale=.87]
 \begin{scope}[shift={(0,0)}]
  \draw[axis] (0,0)--(1.35,0) node[right]{$X$};
  \draw[curve,->] (0,0)--(1.10,.64) node[right]{$X_1$};
  \draw[curve,->] (0,0)--(0,1.35) node[above]{$Z_1=Z$};
  \draw[->,TextGray] (.48,.05) arc[start angle=6,end angle=30,radius=.50];
  \node at (.25,-1.05) {\strut(1)};
 \end{scope}
 \begin{scope}[shift={(5,0)}]
  \draw[curve,->] (0,0)--(1.35,0) node[below]{$X_2=X_1$};
  \draw[axis] (0,0)--(0,1.35) node[above]{$Z_1$};
  \draw[curve,->] (0,0)--(.67,1.16) node[above right]{$Z_2$};
  \draw[->,TextGray] (.05,.50) arc[start angle=84,end angle=60,radius=.50];
  \node at (.25,-1.05) {\strut(2)};
 \end{scope}
 \begin{scope}[shift={(10,0)}]
  \draw[axis] (0,0)--(1.35,0) node[right]{$X_2$};
  \draw[curve,->] (0,0)--(1.10,.64) node[right]{$x$};
  \draw[curve,->] (0,0)--(.30,1.32) node[above]{$z=Z_2$};
  \draw[->,TextGray] (.48,.05) arc[start angle=6,end angle=30,radius=.50];
  \node at (.25,-1.05) {\strut(3)};
 \end{scope}
 \draw[-{Latex[length=2.2mm]},TextGray,thick] (2.05,.20)--(3.60,.20);
 \draw[-{Latex[length=2.2mm]},TextGray,thick] (7.05,.20)--(8.60,.20);
\end{tikzpicture}
\bookcaption{欧拉角的三步有序转动；次序改变会改变最终姿态}\label{advA1:fig:euler-steps}
\end{minipage}\end{center}
图\ref{advA1:fig:euler-steps}将三次转动分开画出，避免在一张三维图中混淆中间轴。由矩阵乘积的莱布尼茨求导法则，总角速度是三步瞬时角速度的矢量和。在刚体主轴系中，
\begin{equation}\label{advA1:eq:euler-kinematic}
 \begin{aligned}
 \omega_1&=\dot\varphi\sin\theta\sin\psi+\dot\theta\cos\psi,\\
 \omega_2&=\dot\varphi\sin\theta\cos\psi-\dot\theta\sin\psi,\\
 \omega_3&=\dot\varphi\cos\theta+\dot\psi.
 \end{aligned}
\end{equation}
这是欧拉运动学方程。在 $\theta=0$ 处 $\varphi$ 与 $\psi$ 的作用重合，是坐标奇点，不是刚体姿态的物理奇点。

\subsection{惯量张量与欧拉动力学方程}

对绕定点以角速度 $\vect\omega$ 转动的刚体，
\[
 T=\frac12\int\rho(\vect\omega\times\vect r)^2\dd V,
 \qquad
 \vect L=\int\rho\vect r\times(\vect\omega\times\vect r)\dd V.
\]
用三重矢积公式，
\begin{equation}\label{advA1:eq:inertia-tensor}
 \vect L=\mathsf I\vect\omega,
 \quad
 \mathsf I=\int\rho(r^2\mathsf 1-\vect r\vect r^{\mathsf T})\dd V,
 \quad
 T=\frac12\vect\omega\cdot\mathsf I\vect\omega.
\end{equation}
若二阶质量矩矩阵为 $\mathsf M=\int\rho\vect r\vect r^{\mathsf T}\dd V$，则
$\mathsf I=\operatorname{tr}(\mathsf M)\mathsf1-\mathsf M$。

\begin{advexamplebox}{绕任意轴的转动惯量}\label{advA1:ex:axis-inertia}
刚体绕经定点的单位方向 $\vect n$ 以角速度 $\omega$ 转动。证明该轴的转动惯量为
\[
 I_{\vect n}=\vect n^{\mathsf T}\mathsf I\vect n.
\]
\end{advexamplebox}
\begin{advsolution}
令 $\vect\omega=\omega\vect n$，将式\eqref{advA1:eq:inertia-tensor}中的动能与定义 $T=I_{\vect n}\omega^2/2$ 比较，
\[
 \frac12I_{\vect n}\omega^2
 =\frac12\omega^2\vect n^{\mathsf T}\mathsf I\vect n.
\]
因而得证。同时
$\vect L\cdot\vect n=\omega\vect n^{\mathsf T}\mathsf I\vect n=I_{\vect n}\omega$，但若 $\vect n$ 不是主轴，$\vect L$ 本身不与 $\vect\omega$ 平行。
\end{advsolution}

\begin{advtheorem}{主轴定理}\label{advA1:thm:principal-axes}
实对称惯量张量存在一组正交主轴，在该基底下
\[
 \mathsf I=\operatorname{Diag}(I_1,I_2,I_3).
\]
\end{advtheorem}
这是实对称矩阵的谱定理。解
$\det(\mathsf I-\lambda\mathsf1)=0$得三个实本征值；不同本征值的本征矢量自动正交，简并子空间内再作格拉姆--施密特正交化。以三个归一化本征矢量为列构造正交矩阵，即得对角化。在主轴系中
\[
 L_i=I_i\omega_i,
 \qquad T=\frac12(I_1\omega_1^2+I_2\omega_2^2+I_3\omega_3^2).
\]

\begin{advtheorem}{欧拉动力学方程}\label{advA1:thm:euler-dynamics}
在刚体主轴系中，外力矩分量为 $M_i$ 时，
\begin{equation}\label{advA1:eq:euler-dynamics}
 \begin{aligned}
 I_1\dot\omega_1+(I_3-I_2)\omega_2\omega_3&=M_1,\\
 I_2\dot\omega_2+(I_1-I_3)\omega_3\omega_1&=M_2,\\
 I_3\dot\omega_3+(I_2-I_1)\omega_1\omega_2&=M_3.
 \end{aligned}
\end{equation}
\end{advtheorem}
证明：惯性系中 $\dd\vect L/\dd t=\vect M$。任意矢量的两系导数关系为
\[
 \left(\frac{\dd\vect L}{\dd t}\right)_{\rm inertial}
 =\left(\frac{\dd\vect L}{\dd t}\right)_{\rm body}
 +\vect\omega\times\vect L.
\]
在主轴系中 $\vect L=(I_1\omega_1,I_2\omega_2,I_3\omega_3)$，逐分量展开即得式\eqref{advA1:eq:euler-dynamics}。

把式\eqref{advA1:eq:euler-kinematic}代入动能，就可用欧拉角写拉格朗日方程。外力矩的功率为
\[
 \vect M\cdot\vect\omega=M_\varphi\dot\varphi
 +M_\theta\dot\theta+M_\psi\dot\psi,
\]
其中
\[
 M_\varphi=M_1\sin\theta\sin\psi+M_2\sin\theta\cos\psi+M_3\cos\theta,
\]
\[
 M_\theta=M_1\cos\psi-M_2\sin\psi,
 \qquad M_\psi=M_3.
\]
它们就是三个欧拉角对应的广义力。

\subsubsection{自由非对称刚体}

设 $I_1<I_2<I_3$。无外力矩时，角动量大小与动能守恒。记刚体系中的角动量分量为
$\Pi_i=I_i\omega_i$，则轨道是
\begin{equation}\label{advA1:eq:pi-intersection}
 \Pi_1^2+\Pi_2^2+\Pi_3^2=L^2,
 \qquad
 \frac{\Pi_1^2}{2I_1}+\frac{\Pi_2^2}{2I_2}
 +\frac{\Pi_3^2}{2I_3}=T.
\end{equation}
即角动量球与能量椭球的交线。
\begin{center}\begin{minipage}{\linewidth}\centering
\begin{tikzpicture}[scale=.88]
 \begin{scope}
  \draw[axis] (-2.3,0)--(2.4,0) node[right]{$\Pi_1$};
  \draw[axis] (0,-2.0)--(0,2.1) node[above]{$\Pi_3$};
  \foreach \sgn in {-1,1}{
   \draw[curve,domain=0:360,samples=121,variable=\t] plot ({.83666*cos(\t)+.460163*sin(\t)},{\sgn*1.8*sqrt(.3+.525*cos(\t)^2)+.167332*sin(\t)});
   \fill (0,{\sgn*1.8}) circle(1.8pt);
   \fill ({\sgn*2},0) circle(1.8pt);
  }
  \node at (0,-2.42){(a) 交线投影；点为稳定主轴方向};
 \end{scope}
 \begin{scope}[shift={(6,0)}]
  \draw[very thick] (0,0) ellipse (1.55 and 1.02);
  \draw[ground] (-2,-1.02)--(2,-1.02);
  \draw[vector] (0,0)--(0,-1.75) node[right]{$\vect L$};
  \fill (0,-1.02) circle(2pt);
  \draw[thin] (.10,-1.02)--(.65,-.70) node[right]{接触点};
  \fill (0,0) circle(1.4pt) node[above]{$O$};
  \node at (0,1.45){能量椭球};
  \node at (0,-2.42){(b) 定平面垂直于 $\vect L$};
 \end{scope}
\end{tikzpicture}
\bookcaption{自由刚体的角动量交线与能量椭球纯滚动图像}\label{advA1:fig:poinsot}
\end{minipage}\end{center}
图\ref{advA1:fig:poinsot}(a) 是三维交线的斜投影，未画出 $\Pi_2$ 轴；投影曲线的交叉不表示三维轨道相交。绕最小和最大主轴的转动稳定，绕中间主轴的转动不稳定，后者会表现为反复翻转，即所谓中间轴定理或扎尼别科夫效应。在角速度空间，能量椭球在惯性系中沿一个与守恒 $\vect L$ 垂直的定平面纯滚动，如图\ref{advA1:fig:poinsot}(b)，这就是庞锁构造。

\subsubsection{对称陀螺与章动}

对 $I_1=I_2=I, I_3=I_0$ 的对称陀螺，
\begin{equation}\label{advA1:eq:symmetric-top-T}
 T=\frac I2(\dot\theta^2+\dot\varphi^2\sin^2\theta)
 +\frac{I_0}{2}(\dot\psi+\dot\varphi\cos\theta)^2.
\end{equation}
若势能仅为 $V(\theta)$，$\varphi,\psi$ 都是循环坐标：
\[
 J_\varphi=I\dot\varphi\sin^2\theta+J_\psi\cos\theta,
 \qquad
 J_\psi=I_0(\dot\psi+\dot\varphi\cos\theta).
\]
同时罗斯约化两个循环坐标，
\begin{equation}\label{advA1:eq:top-effective}
 \mathcal R=\frac I2\dot\theta^2-V_{\rm eff}(\theta)-\frac{J_\psi^2}{2I_0},
 \quad
 V_{\rm eff}=\frac{(J_\varphi-J_\psi\cos\theta)^2}
 {2I\sin^2\theta}+V(\theta).
\end{equation}
定义约化能 $E_\theta=E-J_\psi^2/(2I_0)$，转向点由 $E_\theta=V_{\rm eff}$ 决定，再代回
\[
 \dot\varphi=\frac{J_\varphi-J_\psi\cos\theta}{I\sin^2\theta},
 \qquad
 \dot\psi=\frac{J_\psi}{I_0}-\dot\varphi\cos\theta
\]
即可求进动与自转。
\begin{center}\begin{minipage}{\linewidth}\centering
\begin{tikzpicture}[x=1.9cm,y=1.45cm]
 \draw[axis] (0,0)--(3.35,0) node[right]{$\theta$};
 \draw[axis] (0,0)--(0,2.35) node[above]{$V_{\rm eff}$};
 \draw[curve,domain=.50:2.64159,samples=101] plot (\x,{.5/(sin(\x r)^2)});
 \draw[TextGray,thick] (.4,1)--(2.75,1) node[right]{$E_\theta$};
 \foreach \x/\lab in {.785398/\theta_-,2.356194/\theta_+}{\draw[guide] (\x,0)--(\x,1);\fill (\x,1) circle(1.6pt);\node[below] at (\x,0){$\lab$};}
 \fill (1.570796,.5) circle(1.7pt) node[below=4pt]{$\theta_0$};
 \node[below] at (3.14159,0){$\pi$};
 \node[font=\small] at (1.6,2.7){示例：$V_{\rm eff}=1/(2\sin^2\theta)$};
\end{tikzpicture}
\bookcaption{对称陀螺的有效势、稳定无章动点与章动转向点}\label{advA1:fig:top-effective}
\end{minipage}\end{center}
图\ref{advA1:fig:top-effective}中 $\theta_0$ 是有效势极小点；能量稍高时，$\theta$ 在 $\theta_-$ 与 $\theta_+$ 之间章动。

\begin{advexamplebox}{无章动条件与小章动频率}\label{advA1:ex:nutation}
对式\eqref{advA1:eq:top-effective}的一般对称陀螺，求恒定 $\theta=\theta_0$ 的条件，并求其附近小章动的频率。
\end{advexamplebox}
\begin{advsolution}
无章动要求
\[
 V_{\rm eff}'(\theta_0)=0,
 \qquad E=V_{\rm eff}(\theta_0)+\frac{J_\psi^2}{2I_0}.
\]
令 $A(\theta)=J_\varphi-J_\psi\cos\theta$，第一个条件显式写成
\[
 \frac1I\left[
 \frac{A J_\psi}{\sin\theta}
 -\frac{A^2\cos\theta}{\sin^3\theta}\right]_{\theta_0}
 +V'(\theta_0)=0.
\]
令 $\theta=\theta_0+\eta$，将有效势展开到二阶：
\[
 I\ddot\eta+V_{\rm eff}''(\theta_0)\eta=0.
\]
因而
\[
 \boxed{\omega_{\rm nut}^2=\frac{V_{\rm eff}''(\theta_0)}I}.
\]
稳定无章动还要求 $V_{\rm eff}''(\theta_0)>0$；若小于零，则微小扰动指数增长。
\end{advsolution}

\begin{advexamplebox}{始终保持底边水平的等边三角板}\label{advA1:ex:triangle-top}
均匀等边三角形薄板 $OAB$ 的质量为 $m$、边长为 $l$，绕过顶点 $O$ 的竖直轴定点转动，且 $AB$ 始终水平。
(1) 当板以进动角速度 $\omega$ 作无章动运动时，求板中线 $OG$ 与竖直向下方向的夹角 $\phi$；
(2) 在不施加额外竖直轴力矩的情况下，求 $\phi$ 发生微小章动的周期。
\end{advexamplebox}
\AdvHideHeadingtrue
\begin{advsolution}
\noindent\begin{minipage}{\linewidth}{\bfseries\color{DeepBlue}解答}\par
\begin{center}\begin{minipage}{\linewidth}\centering
\begin{tikzpicture}[scale=.92]
 \draw[axis] (0,1.9)--(0,-2.0) node[below]{竖直轴};
 \coordinate (O) at (0,.85);
 \coordinate (G) at (.7,-.083333);
 \coordinate (A) at (.15,-.55);
 \coordinate (B) at (1.95,-.55);
 \fill[DeepBlue!7] (O)--(A)--(B)--cycle;
 \draw[very thick] (O)--(A)--(B)--cycle;
 \draw[guide] (O)--(0,-1.35);
 \draw[curve] (O)--(G) node[midway,above right]{$OG$};
 \fill (G) circle(2.3pt) node[right=4pt]{$G$};
 \draw[->,DeepBlue] ($(O)+(0,-.48)$) arc[start angle=-90,end angle=-53,radius=.48];
 \node at (.40,.20) {$\phi$};
 \draw[very thick,TextGray] (A)--(B) node[midway,below=4pt]{$AB$水平};
 \draw[->,TextGray] (-.55,.85) arc[start angle=180,end angle=20,x radius=.85,y radius=.25];
 \node at (-.82,1.15) {$\alpha$};
\end{tikzpicture}
\bookcaption{等边三角板的中线、水平底边和进动角}\label{advA1:fig:triangle}
\end{minipage}\end{center}\end{minipage}

如图\ref{advA1:fig:triangle}，记板的方位角为 $\alpha$。质心到 $O$ 的距离
\[
 a=OG=\frac{l}{\sqrt3}.
\]
取过 $O$ 且平行 $AB$ 的轴、板法线和 $OG$ 为三条主轴，其主惯量为
\[
 I_1=\frac{3}{8}ml^2,
 \qquad I_2=\frac{5}{12}ml^2,
 \qquad I_3=\frac1{24}ml^2,
 \qquad I_2-I_3=I_1.
\]
动能和势能为
\[
 T=\frac12I_1\dot\phi^2+
 \frac12A(\phi)\dot\alpha^2,
 \quad A(\phi)=I_2\sin^2\phi+I_3\cos^2\phi,
 \qquad V=-mga\cos\phi.
\]

(1) 无章动时 $\dot\alpha=\omega$，$\phi$ 方程给出
\[
 (I_2-I_3)\omega^2\sin\phi\cos\phi-mga\sin\phi=0.
\]
非竖直解为
\[
 \boxed{\cos\phi=\frac{mga}{(I_2-I_3)\omega^2}
 =\frac{8g}{3\sqrt3\,l\omega^2}}.
\]
它要求 $\omega^2\ge8g/(3\sqrt3\,l)$。

(2) 章动时 $\alpha$ 不再保持严格匀速，守恒的是
$p_\alpha=A(\phi)\dot\alpha=A(\phi_0)\omega$。罗斯函数为
\[
 \mathcal R=\frac12I_1\dot\phi^2+mga\cos\phi
 -\frac{p_\alpha^2}{2A(\phi)}.
\]
在平衡角 $\phi_0=\phi$ 附近线性化，并用第(1)问的平衡条件化简，得
\[
 \Omega_{\rm nut}^2
 =\omega^2\sin^2\phi
 \left[1+\frac{4I_1\cos^2\phi}{A(\phi)}\right]
 =\omega^2\sin^2\phi\,
 \frac{10+27\cos^2\phi}{10-9\cos^2\phi}.
\]
因而
\[
 \boxed{T_{\rm nut}=\frac{2\pi}{\omega\sin\phi}
 \sqrt{\frac{10-9\cos^2\phi}{10+27\cos^2\phi}}}.
\]
若误把章动过程中的 $\dot\alpha$ 仍当作外部强制的常量，会遗漏根号中的惯性耦合因子。
\end{advsolution}

\section{连续体}

\subsection{分析力学的变分表述}

有限维函数的极值由导数为零刻画；当未知对象本身是函数时，研究对象就变成泛函。设区域
$D\subset\mathbb R^n$ 上有若干函数 $f_\mu(\vect x)$，局域且只含一阶导数的密度泛函写成
\[
 \mathcal L=\mathcal L(f_\mu,\partial_i f_\mu,x_i),
 \qquad
 S[f_\mu]=\int_D\mathcal L\,\dd^n x .
\]

\begin{advdefinition}{泛函导数}\label{advA1:def:functional-derivative}
若对满足指定边界条件的任意测试函数 $g_\mu$，
\[
 \left.\frac{\dd}{\dd\epsilon}S[f_\mu+\epsilon g_\mu]\right|_{\epsilon=0}
 =\int_D\sum_\mu\frac{\delta S}{\delta f_\mu}g_\mu\,\dd^n x ,
\]
则称 $\delta S/\delta f_\mu$ 为泛函在 $f_\mu$ 处的泛函导数。
\end{advdefinition}

\begin{advtheorem}{欧拉--拉格朗日方程}\label{advA1:thm:field-el}
当边界值 $f_\mu|_{\partial D}$ 固定时，$S$ 的一阶变分为零等价于
\begin{equation}\label{advA1:eq:field-el}
 \frac{\delta S}{\delta f_\mu}
 =\frac{\partial\mathcal L}{\partial f_\mu}
 -\sum_i\partial_i\left(
 \frac{\partial\mathcal L}{\partial(\partial_i f_\mu)}\right)=0 .
\end{equation}
\end{advtheorem}
证明：令 $\delta f_\mu=\epsilon g_\mu$，展开到一阶：
\[
 \delta S=\int_D\sum_\mu\left[
 \frac{\partial\mathcal L}{\partial f_\mu}g_\mu+
 \sum_i\frac{\partial\mathcal L}{\partial(\partial_i f_\mu)}
 \partial_i g_\mu\right]\dd^n x .
\]
对第二项分部积分。边界上 $g_\mu=0$，边界项消失；区域内 $g_\mu$ 任意，故其系数必须处处为零，即得式\eqref{advA1:eq:field-el}。

这个推导同时给出三个常用变式。若边界值自由，分部积分留下自然边界条件
\[
 \left.\vect n\cdot
 \frac{\partial\mathcal L}{\partial(\nabla f_\mu)}\right|_{\partial D}=0.
\]
若密度含到 $n$ 阶导数，则反复分部积分得到欧拉--泊松方程
\[
 \sum_{k=0}^{n}(-1)^k\partial_{i_1}\cdots\partial_{i_k}
 \left[\frac{\partial\mathcal L}
 {\partial(\partial_{i_1}\cdots\partial_{i_k}f_\mu)}\right]=0.
\]
若还带有积分约束 $C_a[f]=0$，则对
$\mathcal L-\sum_a\lambda_a\mathcal C_a$ 使用同一方程，并与全部约束联立求 $\lambda_a$。

\begin{advexamplebox}{固定均值与能量时的最大熵分布}\label{advA1:ex:max-entropy}
求一维速度分布 $f(v)$，使
\[
 -\int_{\mathbb R}f\ln f\,\dd v
\]
在
\[
 \int f\,\dd v=1,\qquad \int vf\,\dd v=0,\qquad
 \int\frac12mv^2f\,\dd v=\frac12kT
\]
条件下最大。
\end{advexamplebox}
\begin{advsolution}
引入三个拉格朗日乘数，令
\[
 \mathcal F[f]=-\int f\ln f\,\dd v-\lambda_0\!\left(\int f\,\dd v-1\right)
 -\lambda_1\int vf\,\dd v
 -\lambda_2\!\left(\int\frac12mv^2f\,\dd v-\frac12kT\right).
\]
泛函导数为
\[
 \frac{\delta\mathcal F}{\delta f}
 =-(\ln f+1)-\lambda_0-\lambda_1v-\frac12\lambda_2mv^2=0.
\]
因此 $f=Ae^{-av^2-bv}$。零均值给出 $b=0$，归一化与能量约束给出
\[
 \boxed{f(v)=\sqrt{\frac{m}{2\pi kT}}
 \exp\!\left(-\frac{mv^2}{2kT}\right)}.
\]
二阶变分
$\delta^2(-\int f\ln f)=-\int(\delta f)^2/f\,\dd v<0$，故驻点确为最大值。其微分熵为
\[
 S_{\max}=\frac12\left[1+\ln\!\left(\frac{2\pi kT}{m}\right)\right].
\]
\end{advsolution}

\newcommand{\AOneElasticSecond}{%
\subsubsection{应力、各向同性本构关系与弹性波}

过一点取有向面元 $\dd\vect S=\vect n\,\dd S$。相邻两侧材料通过该面元传递的力在局域极限下与面元线性：
\[
 \dd\vect F=\bm\sigma\,\dd\vect S,
\]
其中 $\bm\sigma$ 是应力张量。对角元是正应力，非对角元是剪应力。

\begin{advtheorem}{经典连续体的应力张量对称}\label{advA1:thm:stress-symmetric}
若质元内部不存在独立的偶应力，局域角动量守恒要求
\[
 \boxed{\sigma_{ij}=\sigma_{ji}}.
\]
\end{advtheorem}
证明：取边长为 $\ell$ 的微小立方体。相对两个面上的应力合力对质心的力矩中，体力矩是
$\order{\ell^4}$，而剪应力对产生的力偶是 $\order{\ell^3}$。例如绕 $z$ 轴的领先项为
$(\sigma_{21}-\sigma_{12})\ell^3$。令 $\ell\to0$ 时角动量平衡必须逐阶成立，故
$\sigma_{12}=\sigma_{21}$；循环指标即得其余分量。
\begin{center}\begin{minipage}{\linewidth}\centering
\begin{tikzpicture}[scale=.92]
 \draw[very thick] (0,0) rectangle (2.2,1.6);
 \draw[vector] (0,1.25)--(0,.45) node[midway,left]{$\sigma_{21}$};
 \draw[vector] (2.2,.35)--(2.2,1.15) node[midway,right]{$\sigma_{21}$};
 \draw[vector] (1.25,0)--(.45,0) node[midway,below]{$\sigma_{12}$};
 \draw[vector] (.95,1.6)--(1.75,1.6) node[midway,above]{$\sigma_{12}$};
 \node at (1.1,.8) {微小质元};
\end{tikzpicture}
\bookcaption{剪应力不成对相等时会留下随体积缩放的净力矩}\label{advA1:fig:stress-symmetry}
\end{minipage}\end{center}
图\ref{advA1:fig:stress-symmetry}中的箭头作用在不同面上；相等的是互易剪应力，而不是把同一面上的两个方向混为一谈。

\begin{advtheorem}{圣维南--基尔霍夫各向同性线弹性关系}\label{advA1:thm:svk}
各向同性、线性小应变材料的弹性能密度和应力为
\begin{equation}\label{advA1:eq:svk}
 W=\frac{\lambda}{2}(\operatorname{tr}\bm\varepsilon)^2
 +\mu\,\operatorname{tr}(\bm\varepsilon^2),
 \qquad
 \bm\sigma=\lambda(\operatorname{tr}\bm\varepsilon)\mathsf1
 +2\mu\bm\varepsilon ,
\end{equation}
其中 $\lambda,\mu$ 是拉梅参数。
\end{advtheorem}
证明：二阶对称张量在 $SO(3)$ 下分为迹和对称无迹两个互不混合的不可约部分。各向同性二次型只能分别给这两个部分赋予一个标量系数，不能出现依赖方向的交叉项，故
$W=a(\operatorname{tr}\bm\varepsilon)^2+b\,\operatorname{tr}(\bm\varepsilon^2)$。把常数重命名为
$a=\lambda/2,b=\mu$，再用
$\sigma_{ij}=\partial W/\partial\varepsilon_{ij}$ 即得式\eqref{advA1:eq:svk}。

线弹性场的拉格朗日密度是
\[
 \mathcal L=\frac12\rho_m\dot{\vect\xi}^{\,2}-W,
\]
其中 $\rho_m$ 是质量密度。场方程为纳维--柯西方程
\begin{equation}\label{advA1:eq:navier-cauchy}
 \rho_m\ddot{\vect\xi}
 =(\lambda+\mu)\nabla(\nabla\cdot\vect\xi)
 +\mu\nabla^2\vect\xi+\vect f .
\end{equation}

\begin{advexamplebox}{由应力散度推出纳维--柯西方程}\label{advA1:ex:stress-divergence}
从应力张量的局域动量平衡
$\rho_m\ddot{\vect\xi}=\nabla\cdot\bm\sigma+\vect f$
出发，推导式\eqref{advA1:eq:navier-cauchy}。
\end{advexamplebox}
\begin{advsolution}
由式\eqref{advA1:eq:svk}及
$2\bm\varepsilon=\nabla\vect\xi+(\nabla\vect\xi)^{\mathsf T}$，
\[
 \bm\sigma=\lambda(\nabla\cdot\vect\xi)\mathsf1
 +\mu[\nabla\vect\xi+(\nabla\vect\xi)^{\mathsf T}].
\]
逐分量取散度：
\[
 (\nabla\cdot\bm\sigma)_i
 =\lambda\partial_i\partial_j\xi_j
 +\mu\partial_j\partial_j\xi_i
 +\mu\partial_j\partial_i\xi_j.
\]
交换光滑场的偏导次序，得
\[
 \nabla\cdot\bm\sigma
 =(\lambda+\mu)\nabla(\nabla\cdot\vect\xi)
 +\mu\nabla^2\vect\xi,
\]
代回动量平衡即得所求。边界牵引条件是
$\bm\sigma\cdot\vect n=\dd\vect F/\dd S$。
\end{advsolution}

令 $\vect\xi=\vect A e^{\ii(\vect k\cdot\vect r-\omega t)}$，把振幅分解为纵向
$\vect A_\parallel$ 与横向 $\vect A_\perp$。式\eqref{advA1:eq:navier-cauchy}给出
\[
 c_{\rm L}=\sqrt{\frac{\lambda+2\mu}{\rho_m}},
 \qquad
 c_{\rm T}=\sqrt{\frac{\mu}{\rho_m}}.
\]
纵波只有一个偏振，横波有两个彼此正交的偏振；因此固体中的首波（P 波）比次波（S 波）快。
\begin{center}\begin{minipage}{\linewidth}\centering
\begin{tikzpicture}[scale=.9]
 \begin{scope}
  \draw[axis] (-.2,0)--(4.3,0) node[right]{$\vect k$};
  \foreach \x in {0,.46,1.02,1.72,2.48,3.12,3.60}{\fill (\x,0) circle(2pt);}
  \draw[<->,DeepBlue,thick] (1.2,-.38)--(1.75,-.38);
  \node[DeepBlue,left] at (1.0,-.38) {$\vect A_\parallel$};
  \node at (2,-.9) {纵波};
 \end{scope}
 \begin{scope}[shift={(6.0,0)}]
  \draw[axis] (-.2,0)--(4.3,0) node[right]{$\vect k$};
  \foreach \x in {0,.6,...,3.6}{\fill (\x,{.45*sin(120*\x)}) circle(2pt);}
  \draw[<->,DeepBlue,thick] (1.2,-.35)--(1.2,.45) node[right]{$\vect A_\perp$};
  \node at (2,-.9) {横波};
 \end{scope}
\end{tikzpicture}
\bookcaption{弹性体纵波的位移平行于波矢，横波的位移垂直于波矢}\label{advA1:fig:elastic-waves}
\end{minipage}\end{center}
图\ref{advA1:fig:elastic-waves}把传播方向与位移方向同时标出，二者不能仅凭波形外观判断。

单轴拉伸定义杨氏模量 $E_Y$ 与泊松比 $\nu$。它们和拉梅参数的关系为
\[
 \nu=\frac{\lambda}{2(\lambda+\mu)},\qquad
 E_Y=\frac{\mu(3\lambda+2\mu)}{\lambda+\mu},
\]
\[
 \lambda=\frac{\nu E_Y}{(1+\nu)(1-2\nu)},\qquad
 \mu=\frac{E_Y}{2(1+\nu)}.
\]
体积模量与剪切模量为
\[
 K=\lambda+\frac{2\mu}{3},\qquad G=\mu.
\]
其中
$\operatorname{tr}\bm\sigma=3K\operatorname{tr}\bm\varepsilon$，而偏应力与偏应变满足
$\bm\Sigma=2G\bm E$。

\begin{advexamplebox}{梁的均匀弯曲}\label{advA1:ex:beam-bending}
长为 $L$、宽为 $b$、厚为 $h$ 的均匀矩形梁弯成总转角 $\Theta$ 的圆弧。质量密度为
$\rho_m$。求应变分布、弹性能和维持弯曲所需端力矩，并用拉梅参数表示结果。
\end{advexamplebox}
\begin{advsolution}
欧拉--伯努利假设下，中性面到距离 $y$ 处的纵向应变为
\[
 \varepsilon_{xx}=-\kappa y,\qquad \kappa=\frac{\Theta}{L}.
\]
侧面自由时应力是单轴的，
$\sigma_{xx}=E_Y\varepsilon_{xx}$。截面二次矩
$I_b=bh^3/12$，所以
\[
 M=\int_A(-y)\sigma_{xx}\dd A
 =E_YI_b\kappa
 =\boxed{\frac{E_Ybh^3}{12L}\Theta}.
\]
弹性能
\[
 U=\frac12\int_0^L E_YI_b\kappa^2\dd x
 =\boxed{\frac{E_Ybh^3}{24L}\Theta^2}.
\]
若用质量 $m=\rho_mbhL$ 消去 $b$，
\[
 U=\frac{E_Ymh^2}{24\rho_mL^2}\Theta^2,
 \qquad
 E_Y=\frac{\mu(3\lambda+2\mu)}{\lambda+\mu}.
\]
厚度以三次方进入弯曲刚度，这是最容易在尺度估算中遗漏的关系。
\begin{center}\begin{minipage}{\linewidth}\centering
\begin{tikzpicture}[scale=.9]
 \coordinate (C) at (0,0);
 \fill[DeepBlue!8] (35:2.7) arc(35:145:2.7)--(145:3.3) arc(145:35:3.3)--cycle;
 \draw[very thick] (35:2.7) arc(35:145:2.7)--(145:3.3) arc(145:35:3.3)--cycle;
 \draw[guide] (35:3) arc(35:145:3);
 \draw[guide] (C)--(35:3.6) (C)--(145:3.6);
 \draw[DeepBlue,->] (35:.85) arc(35:145:.85);
 \node at (0,1.08){$\Theta$};\fill(C) circle(1.5pt) node[below]{$C$};
 \draw[<->] (C)--(90:3) node[midway,right]{$R$};
 \draw[<->] (35:2.7)--(35:3.3) node[midway,right=5pt]{$h$};
 \node at (0,3.75){同心弯梁截面};
\end{tikzpicture}
\bookcaption{梁弯曲时中性层不伸长，上下表面的纵向应变符号相反}\label{advA1:fig:beam}
\end{minipage}\end{center}
\end{advsolution}

\begin{advexamplebox}{圆柱的扭转}\label{advA1:ex:cylinder-torsion}
半径 $R$、长度 $L$ 的均匀实心圆柱，底面固定，顶面刚性转过小角 $\Theta$。求位移场、剪应力、维持转动所需力矩和弹性能。
\end{advexamplebox}
\begin{advsolution}
取柱坐标，满足两端边界的位移为
\[
 \xi_\varphi(r,z)=\frac{\Theta z}{L}r,\qquad
 \xi_r=\xi_z=0.
\]
工程剪应变与剪应力为
\[
 \gamma_{z\varphi}=\frac{\partial\xi_\varphi}{\partial z}
 =\frac{\Theta r}{L},\qquad
 \tau_{z\varphi}=\mu\frac{\Theta r}{L}.
\]
顶面力矩
\[
 M=\int_A\tau_{z\varphi}r\,\dd A
 =\frac{\mu\Theta}{L}\int_0^Rr^2(2\pi r\dd r)
 =\boxed{\frac{\pi\mu R^4}{2L}\Theta}.
\]
因此
\[
 \boxed{U=\frac12M\Theta
 =\frac{\pi\mu R^4}{4L}\Theta^2}.
\]
空心圆柱只需把 $R^4$ 换为 $R_{\rm o}^4-R_{\rm i}^4$。力矩量纲
$[\mu R^4/L]=FL$，检查正确。
\begin{center}\begin{minipage}{\linewidth}\centering
\begin{tikzpicture}[scale=.86]
 \draw[very thick] (-1.2,-1.5)--(-1.2,1.35);
 \draw[very thick] (1.2,-1.5)--(1.2,1.35);
 \draw[very thick] (0,1.35) ellipse (1.2 and .35);
 \draw[very thick] (0,-1.5) ellipse (1.2 and .35);
 \draw[->,DeepBlue,thick] (.20,1.85) arc[start angle=105,end angle=-25,x radius=1.15,y radius=.42];
 \node at (1.45,1.72) {$\Theta$};
 \draw[<->,TextGray] (1.55,-1.5)--(1.55,1.35) node[midway,right]{$L$};
 \draw[ground] (-1.55,-1.86)--(1.55,-1.86);
\end{tikzpicture}
\bookcaption{圆柱顶面相对底面的刚性扭转与轴向长度}\label{advA1:fig:torsion}
\end{minipage}\end{center}
\end{advsolution}

}

\newcommand{\AOneFluidSection}{%
\subsection{流体}

流体没有由静态形变直接决定的恢复应力，其应力取决于形变速率，并且常与热学自由度耦合。因此用欧拉描述：速度、密度、压强和温度是空间与时间的场
$\vect v(\vect r,t),\rho(\vect r,t),p(\vect r,t),T(\vect r,t)$。质量守恒为
\begin{equation}\label{advA1:eq:continuity}
 \frac{\partial\rho}{\partial t}+\nabla\cdot(\rho\vect v)=0.
\end{equation}

\begin{advdefinition}{随体全导数}\label{advA1:def:material-derivative}
\[
 \frac{D}{Dt}=\frac{\partial}{\partial t}+\vect v\cdot\nabla.
\]
它是在跟随同一流体质元运动时测得的变化率。
\end{advdefinition}
由式\eqref{advA1:eq:continuity}，
\[
 \frac{D\rho}{Dt}=-\rho\nabla\cdot\vect v.
\]
速度梯度分解为
\[
 \nabla\vect v=\bm U+\bm W,\qquad
 \bm U=\frac12[\nabla\vect v+(\nabla\vect v)^{\mathsf T}],
 \quad
 \bm W=\frac12[\nabla\vect v-(\nabla\vect v)^{\mathsf T}],
\]
其中 $\bm U$ 是应变率张量，$\bm W$ 是涡张量；涡量矢量为
$\vect\omega=\nabla\times\vect v$。再去掉 $\bm U$ 的迹得到偏应变率
$\bm U'=\bm U-(\nabla\cdot\vect v)\mathsf1/3$。

\begin{advdefinition}{线性本构前提}\label{advA1:def:linear-constitutive}
本构方程中的材料参数视为与流体自由度无关的常数；应力对形变率只保留线性项。
\end{advdefinition}

\begin{advdefinition}{牛顿流体}\label{advA1:def:newtonian-fluid}
采用“压应力为负”的张量符号时，
\[
 \bm\sigma=-p\mathsf1+2\eta\bm U'
 +\zeta(\nabla\cdot\vect v)\mathsf1,
\]
其中 $\eta$ 是动力黏度，$\zeta$ 是体黏度（第二黏度）。
\end{advdefinition}

\begin{advdefinition}{正压流体}\label{advA1:def:barotropic}
若物态方程可写为 $\rho=\rho(p)$，或在指定熵、温度条件下把密度视为压强的单值函数，则称为正压流体。此时
$\nabla\rho\times\nabla p=0$。
\end{advdefinition}

\begin{advdefinition}{理想流体}\label{advA1:def:ideal-fluid}
忽略全部黏滞偏应力，即
\[
 \bm\sigma=-p\mathsf1.
\]
\end{advdefinition}

\begin{advdefinition}{不可压缩流体}\label{advA1:def:incompressible}
沿质元 $\rho$ 为常数；对均匀密度情形，
\[
 \rho=\text{常量},\qquad \nabla\cdot\vect v=0.
\]
\end{advdefinition}

\begin{advtheorem}{理想流体的欧拉方程}\label{advA1:thm:euler-fluid}
无体力时，动量守恒的守恒形式和随体形式分别是
\[
 \frac{\partial(\rho\vect v)}{\partial t}
 +\nabla\cdot(\rho\vect v\vect v+p\mathsf1)=0,
\]
\[
 \boxed{\rho\frac{D\vect v}{Dt}=-\nabla p}.
\]
\end{advtheorem}
证明：对任意固定控制体应用动量守恒。对流动量通量是
$\rho\vect v\vect v$，表面压强力用散度定理化为 $-\nabla p$。再用质量连续性方程消去
$\vect v[\partial_t\rho+\nabla\cdot(\rho\vect v)]$，便得到随体形式。二次通量使欧拉方程本质上是非线性的。

\begin{advtheorem}{开尔文环流定理}\label{advA1:thm:kelvin}
对无黏、正压且体力为保守力的流动，任意随体闭曲线 $C(t)$ 上的环流
\[
 \Gamma=\oint_{C(t)}\vect v\cdot\dd\vect l
\]
满足 $D\Gamma/Dt=0$。
\end{advtheorem}
证明：随体线元满足
$D(\dd\vect l)/Dt=(\dd\vect l\cdot\nabla)\vect v$，故
\[
 \frac{D\Gamma}{Dt}=\oint\left[
 \frac{D\vect v}{Dt}\cdot\dd\vect l+
 \vect v\cdot(\dd\vect l\cdot\nabla)\vect v\right].
\]
第二项是 $\oint\dd(v^2/2)=0$。正压流体中
$\nabla p/\rho=\nabla h(p)$，保守体力也是梯度，闭路积分均为零。

\begin{advexamplebox}{涡度方程}\label{advA1:ex:vorticity}
在无黏正压流体中证明
\[
 \frac{D\vect\omega}{Dt}
 =(\vect\omega\cdot\nabla)\vect v
 -(\nabla\cdot\vect v)\vect\omega .
\]
\end{advexamplebox}
\begin{advsolution}
对欧拉方程取旋度。正压条件使
$\nabla\times(\nabla p/\rho)=
\nabla(1/\rho)\times\nabla p=0$。利用
\[
 (\vect v\cdot\nabla)\vect v
 =\nabla\frac{v^2}{2}-\vect v\times\vect\omega
\]
得
\[
 \partial_t\vect\omega=\nabla\times(\vect v\times\vect\omega).
\]
展开叉积旋度，并用 $\nabla\cdot\vect\omega=0$：
\[
 \partial_t\vect\omega
 +(\vect v\cdot\nabla)\vect\omega
 =(\vect\omega\cdot\nabla)\vect v
 -(\nabla\cdot\vect v)\vect\omega.
\]
左边即随体导数。不可压缩流体中第二项消失；二维无伸展流动还会使涡量随质元保持不变。
\end{advsolution}

\begin{advexamplebox}{应变率张量的随体演化}\label{advA1:ex:strain-rate}
设理想流体满足 $\rho=\rho(p)$。证明
\[
 \frac{D\bm U}{Dt}
 =-\frac{\nabla\nabla p}{\rho}
 +\frac{\dd\rho}{\dd p}\frac{\nabla p\otimes\nabla p}{\rho^2}
 -\bm U^2-\bm W^2.
\]
\end{advexamplebox}
\begin{advsolution}
记 $\mathsf A=\nabla\vect v=\bm U+\bm W$。对速度梯度取随体导数时，随体导数与空间梯度不对易：
\[
 \frac{D\mathsf A}{Dt}
 =\nabla\left(\frac{D\vect v}{Dt}\right)-\mathsf A^2.
\]
欧拉方程给出
\[
 -\nabla\left(\frac{\nabla p}{\rho}\right)
 =-\frac{\nabla\nabla p}{\rho}
 +\frac{\dd\rho}{\dd p}\frac{\nabla p\otimes\nabla p}{\rho^2}.
\]
取矩阵的对称部分。因
$\operatorname{sym}(\mathsf A^2)=\bm U^2+\bm W^2$，交叉项
$\bm U\bm W+\bm W\bm U$ 是反对称的，故得到所证公式。
\end{advsolution}

若初始无旋且满足开尔文定理的条件，涡度方程没有产生涡量的源项，流动保持无旋，可写
$\vect v=\nabla\Phi$。利用
\[
 (\vect v\cdot\nabla)\vect v
 =\nabla(v^2/2)-\vect v\times(\nabla\times\vect v),
\]
正压流体的欧拉方程可积分为
\[
 \partial_t\Phi+\frac{v^2}{2}+\int^p\frac{\dd p'}{\rho(p')}=C(t).
\]

\begin{advexamplebox}{半圆柱形机库的升力}\label{advA1:ex:hangar-lift}
半圆柱形机库截面是上半圆 $r=R$，沿 $z$ 方向无限长。外部不可压缩无旋空气的速度势为
\[
 \Phi=-v_0\left(r+\frac{R^2}{r}\right)\cos\theta.
\]
若机库内部压强为远处静压 $p_\infty$，求单位长度机库受到的竖直升力。
\end{advexamplebox}
\begin{advsolution}
速度分量为
\[
 v_r=\partial_r\Phi=-v_0\left(1-\frac{R^2}{r^2}\right)\cos\theta,
\]
\[
 v_\theta=\frac1r\partial_\theta\Phi
 =v_0\left(1+\frac{R^2}{r^2}\right)\sin\theta.
\]
在表面 $r=R$，$v_r=0$、$v_\theta=2v_0\sin\theta$。伯努利方程给出外表面压强
\[
 p_{\rm out}=p_\infty+\frac12\rho
 (v_0^2-4v_0^2\sin^2\theta).
\]
内外压强差沿外法线产生的竖直分力为
\[
 \frac{F_y}{H}
 =R\int_0^\pi(p_\infty-p_{\rm out})\sin\theta\,\dd\theta
 =\frac12\rho v_0^2R\left(
 4\int_0^\pi\sin^3\theta\,\dd\theta-\int_0^\pi\sin\theta\,\dd\theta\right).
\]
用 $\int_0^\pi\sin^3\theta\,\dd\theta=4/3$，得
\[
 \boxed{\frac{F_y}{H}=\frac53\rho v_0^2R}.
\]
\begin{center}\begin{minipage}{\linewidth}\centering
\begin{tikzpicture}[scale=.92]
 \draw[ground] (-2.2,0)--(2.2,0);
 \draw[very thick] (-1.65,0) arc[start angle=180,end angle=0,radius=1.65];
 \foreach \a in {20,50,90,130,160}{
  \draw[vector]
    ({1.65*cos(\a)-.28*sin(\a)},{1.65*sin(\a)+.28*cos(\a)})
    --({1.65*cos(\a)+.28*sin(\a)},{1.65*sin(\a)-.28*cos(\a)});
 }
 \draw[vector] (-2.5,1.45)--(-1.6,1.45) node[midway,above]{$v_0$};
 \draw[axis] (0,0)--(0,2.15) node[above]{$y$};
 \node[fill=white] at (0,.55) {机库内部 $p_\infty$};
\end{tikzpicture}
\bookcaption{半圆柱表面速度增大导致外压降低，压力差的竖直分量形成升力}\label{advA1:fig:hangar}
\end{minipage}\end{center}
\end{advsolution}

热力学基本关系对随体质元写成
\[
 \frac{Du_s}{Dt}=-p\frac{Dv_s}{Dt}+T\frac{Ds_s}{Dt},
 \qquad v_s=1/\rho.
\]
准静态绝热、无耗散时 $Ds_s/Dt=0$，内能变化与速度散度耦合。

\begin{advtheorem}{理想流体的能量方程}\label{advA1:thm:bernoulli-energy}
令总能量密度
\[
 w=\rho u_s+\frac12\rho v^2.
\]
无体力、无黏且准静态绝热时，
\[
 \boxed{\frac{\partial w}{\partial t}
 +\nabla\cdot[(w+p)\vect v]=0}.
\]
\end{advtheorem}
证明：用欧拉方程点乘 $\vect v$ 得到宏观动能方程；用
$Du_s/Dt=-(p/\rho)\nabla\cdot\vect v$ 与质量连续性方程得到内能方程。两式相加，压强功整理为
$\nabla\cdot(p\vect v)$，即得守恒形式。

\begin{advexamplebox}{伯努利原理的两种形式}\label{advA1:ex:bernoulli-principle}
(1) 理想流体受单位质量势能 $\phi(\vect r)$，证明准静态绝热定常流中沿流线
\[
 u_s+\frac p\rho+\frac{v^2}{2}+\phi=C.
\]
(2) 对不可压缩理想流体，证明即使不使用准静态绝热条件，也有
$p/\rho+v^2/2+\phi$ 沿流线不变；若流动无旋，则常数在整个连通流区相同。
\end{advexamplebox}
\begin{advsolution}
(1) 把定常能量方程写成随体形式，并用热力学恒等式
$\dd h=T\dd s+\dd p/\rho$，其中 $h=u_s+p/\rho$。沿绝热流线
$\dd s=0$，故
\[
 \vect v\cdot\nabla\left(
 h+\frac{v^2}{2}+\phi\right)=0.
\]

(2) 不可压缩欧拉方程是
\[
 (\vect v\cdot\nabla)\vect v
 =-\nabla\left(\frac p\rho+\phi\right).
\]
与 $\vect v$ 点乘：
\[
 \vect v\cdot\nabla\left(
 \frac{v^2}{2}+\frac p\rho+\phi\right)=0,
\]
所以沿流线守恒。若 $\nabla\times\vect v=0$，则
$(\vect v\cdot\nabla)\vect v=\nabla(v^2/2)$，整个括号的梯度为零，故在连通流区全局为同一常数。
\end{advsolution}

\begin{advtheorem}{达朗贝尔佯谬}\label{advA1:thm:dalembert}
在无环流、不可压缩、无黏、定常势流模型中，匀速来流对封闭物体的合阻力为零。
\end{advtheorem}
证明可用包围物体的远处控制面，不要求物体前后对称。
对三维局域封闭物体，假定流体在无穷远趋于均匀来流且无源汇、无尾迹。
势流的远场扰动从偶极项开始，故
$\vect v=\vect U+\vect u$ 中 $\vect u=O(R^{-3})$。
Bernoulli 关系给出 $p-p_\infty=-\rho\vect U\cdot\vect u+O(u^2)$。
稳恒动量通量用 $p\mathsf1+\rho\vect v\vect v$ 表示；
均匀部分在闭合控制面上积分为零，扰动部分至多为
$O(R^2R^{-3})\to0$。物面不可穿透，因而总压强力也为零。
二维无环流的局域势流可作相应远场论证；若有环流，仍需区分升力与阻力。
前后几何对称只允许对特殊形状配对压强，不能作为一般证明。
现实阻力来自黏性边界层、分离和尾迹，恰是此理想模型删除的机制。

\subsubsection{不可压缩牛顿流体}

简单平行剪切中
\[
 \tau=\eta\frac{\dd v}{\dd y}.
\]
动力黏度单位为 $\si{Pa.s}$，运动黏度 $\nu=\eta/\rho$ 的单位为
$\si{m^2.s^{-1}}$。
\begin{center}\begin{minipage}{\linewidth}\centering
\begin{tikzpicture}[scale=.9]
 \draw[ground] (-.2,0)--(4.8,0);
 \draw[ground] (-.2,2.2)--(4.8,2.2);
 \foreach \y/\len in {.4/.7,.8/1.25,1.2/1.8,1.6/2.35,2.0/2.9}{
  \draw[vector] (.45,\y)--++(\len,0);
 }
 \draw[<->,TextGray] (4.2,.8)--(4.2,1.2) node[midway,right]{$\dd y$};
 \draw[guide] (1.7,.8)--(4.3,.8) (2.25,1.2)--(4.3,1.2);
\node at (2.2,2.52) {上板带动流体};
\end{tikzpicture}
\bookcaption{牛顿黏滞定律：相邻流层的速度差产生剪应力}\label{advA1:fig:newton-viscosity}
\end{minipage}\end{center}
图\ref{advA1:fig:newton-viscosity}中箭头长度表示速度，标签放在流层外侧，没有遮住速度线。

不可压缩牛顿流体满足纳维--斯托克斯方程
\[
 \rho\left(\frac{\partial\vect v}{\partial t}
 +\vect v\cdot\nabla\vect v\right)
 =-\nabla p+\eta\nabla^2\vect v+\vect f.
\]
可压缩情形还会出现体黏度与 $\nabla(\nabla\cdot\vect v)$ 项。圆管泊肃叶定律
$\Delta p=8\eta QL/(\pi R^4)$ 和低雷诺数球体斯托克斯阻力
$F=6\pi\eta Rv$ 都是其经典特解或渐近结果。

\begin{advexamplebox}{同轴圆柱间的旋转 Couette 流}\label{advA1:ex:cylindrical-couette}
高为 $H$ 的同轴圆柱半径为 $R_1<R_2$，角速度分别为
$\omega_1,\omega_2$。间隙内充满不可压缩牛顿流体，黏度为 $\eta$。求定常速度场及维持两圆柱转动所需外力矩。
\end{advexamplebox}
\begin{advsolution}
轴对称定常解只有方位分量，纳维--斯托克斯方程化为
\[
 \frac{\dd^2v_\theta}{\dd r^2}
 +\frac1r\frac{\dd v_\theta}{\dd r}
 -\frac{v_\theta}{r^2}=0,
\]
故
\[
 v_\theta=Ar+\frac Br,
\]
\[
 A=\frac{\omega_2R_2^2-\omega_1R_1^2}{R_2^2-R_1^2},
 \qquad
 B=\frac{R_1^2R_2^2(\omega_1-\omega_2)}
 {R_2^2-R_1^2}.
\]
剪应力为
\[
 \tau_{r\theta}=\eta\left(
 \frac{\dd v_\theta}{\dd r}-\frac{v_\theta}{r}\right)
 =-\frac{2\eta B}{r^2}.
\]
任一圆柱面上传递的力矩大小与半径无关：
\[
 |\mathcal M|=2\pi rH\cdot r|\tau_{r\theta}|
 =\boxed{4\pi\eta H|B|}.
\]
若 $\omega_1>\omega_2$，维持内筒需沿正方向施加
$\mathcal M_1=4\pi\eta HB$，维持外筒需施加
$\mathcal M_2=-4\pi\eta HB$；两者之和为零，功率差转化为黏滞热。
\begin{center}\begin{minipage}{\linewidth}\centering
\begin{tikzpicture}[scale=.88]
 \draw[very thick] (0,0) circle(.85);
 \draw[very thick] (0,0) circle(1.65);
 \draw[->,DeepBlue,thick] (.85,0) arc[start angle=0,end angle=105,radius=.85];
 \draw[->,TextGray,thick] (1.65,0) arc[start angle=0,end angle=55,radius=1.65];
 \draw[<->,TextGray] (0,0)--(.85,0) node[midway,below]{$R_1$};
 \draw[<->,TextGray] (0,0)--({1.65*cos(-35)},{1.65*sin(-35)}) node[midway,below]{$R_2$};
 \node at (.30,.56) {$\omega_1$};
 \node at (1.45,1.25) {$\omega_2$};
\end{tikzpicture}
\bookcaption{同轴圆柱之间的方位 Couette 流与两侧无滑移边界}\label{advA1:fig:couette}
\end{minipage}\end{center}
\end{advsolution}

\begin{advexamplebox}{环隙中稳定下落圆柱测黏度}\label{advA1:ex:falling-cylinder-viscosity}
半径 $R$ 的圆柱杯内盛密度 $\rho_0$、黏度未知的液体。半径 $r<R$、高 $a$、密度
$\rho>\rho_0$ 的小圆柱与杯同轴，以稳定速度 $v$ 下落。圆柱完全浸没，杯底封闭，液体不可压缩；取远离杯底的长环隙、低雷诺数极限，忽略端部局部损失，但保留连续性所要求的压差与回流。求黏度。
\end{advexamplebox}
\begin{advsolution}
取向下为 $z$ 正向，设圆柱完全浸没，液体不可压缩且杯底封闭。
记 $H=\ln(R/r)$、$D=R^2-r^2$，把静水压梯度扣除后令
$G=\dd p/\dd z-\rho_0g$。充分发展方程与无滑移边界为
\[
 \eta\frac1s\frac{\dd}{\dd s}\left(s\frac{\dd u}{\dd s}\right)=G,
 \qquad u(r)=v,\quad u(R)=0.
\]
积分得到剪切驱动与压差驱动的叠加：
\[
 u(s)=v\frac{\ln(R/s)}H+
 \frac G{4\eta}\left[s^2-R^2+\frac D H\ln\frac R s\right].
\]
圆柱每秒排开 $\pi r^2v$ 的液体，故环隙必须向上回流：
\[
 2\pi\int_r^R u(s)s\,\dd s=-\pi r^2v.
\]
代入速度分布，令 $K=H(R^2+r^2)-D>0$，得
\[
 G=\frac{4\eta v}K.
\]
除侧面剪力外，端面还受到相对于静水压的压差，二者都是向上阻力：
\[
 F_{\rm drag}=-2\pi ra\eta u'(r)+G\pi r^2a
 =\frac{2\pi\eta av(R^2+r^2)}K.
\]
把静水浮力另计，终端速度条件为
$(\rho-\rho_0)\pi r^2ag=F_{\rm drag}$，从而
\[
 \boxed{\eta=\frac{(\rho-\rho_0)r^2g}{2v}
 \left[\ln\frac Rr-\frac{R^2-r^2}{R^2+r^2}\right]}.
\]
方括号严格为正：令 $R/r=\e^H$，它就是 $H-\tanh H>0$。
窄隙 $h=R-r\ll r$ 时约为 $h^3/(3r^3)$，于是
$\eta\simeq(\rho-\rho_0)gh^3/(6rv)$。
原来的纯对数速度剖面仅有剪切流，会给出向下净流量，不满足封底杯的连续性条件。
本式保留压差但忽略入口、出口的局部损失，适用于长圆柱、低雷诺数和远离杯底的充分发展段；
“忽略端部绕流细节”不等于可以删去端面压差。
\begin{center}\begin{minipage}{\linewidth}\centering
\begin{tikzpicture}[scale=.88]
 \draw[very thick] (-2,-2)--(-2,2);
 \draw[very thick] (2,-2)--(2,2);
 \draw[ground] (-2,-2)--(2,-2);
 \fill[DeepBlue!6] (-2,-1.95) rectangle (2,1.15);
 \draw[very thick,fill=white] (-.75,-.65) rectangle (.75,1.05);
 \draw[vector] (0,.25)--(0,-.60) node[midway,right]{$v$};
 \draw[<->,TextGray] (-.75,1.38)--(.75,1.38) node[midway,above]{$2r$};
 \draw[<->,TextGray] (-2,-1.55)--(2,-1.55) node[midway,below]{$2R$};
 \draw[<->,TextGray] (1.05,-.65)--(1.05,1.05) node[midway,right]{$a$};
\end{tikzpicture}
\bookcaption{同轴下落圆柱与杯壁之间的轴向环隙层流}\label{advA1:fig:falling-cylinder}
\end{minipage}\end{center}
图\ref{advA1:fig:falling-cylinder}明确画出 $r<R$、共轴和侧面环隙；液体在环隙内向上净回流，端部局部损失按长环隙近似忽略。
\end{advsolution}

}

\begin{advexamplebox}{等周不等式}\label{advA1:ex:isoperimetric}
设光滑闭曲线长度为 $C$，证明其围成的有向面积满足
\[
 |A|=\left|\frac12\oint(x\,\dd y-y\,\dd x)\right|
 \le\frac{C^2}{4\pi},
\]
并指出等号条件。
\end{advexamplebox}
\begin{advsolution}
以弧长 $s\in[0,C]$ 参数化，并平移原点使
$\int_0^C x\,\dd s=\int_0^C y\,\dd s=0$。周期函数的维廷格不等式给出
\[
 \int_0^C(x^2+y^2)\dd s\le
 \left(\frac{C}{2\pi}\right)^2
 \int_0^C(x'^2+y'^2)\dd s.
\]
弧长参数满足 $x'^2+y'^2=1$。再用柯西不等式：
\[
 2|A|=\left|\int_0^C(xy'-yx')\dd s\right|
 \le\sqrt{\int_0^C(x^2+y^2)\dd s}
 \sqrt{\int_0^C(x'^2+y'^2)\dd s}
 \le\frac{C^2}{2\pi}.
\]
因此 $|A|\le C^2/(4\pi)$。两次不等式同时取等号要求
$x,y$ 只含第一傅里叶谐波且切矢量处处与位矢垂直，故曲线必须是圆。
\end{advsolution}

\begin{advexamplebox}{最速降线}\label{advA1:ex:brachistochrone}
重力沿 $y$ 轴正向，质点从 $(0,0)$ 静止释放，沿光滑轨道到达 $(l,h)$。忽略摩擦，求使下降时间
\[
 \mathcal T[y]=\int_0^l\sqrt{\frac{1+y'^2}{2gy}}\,\dd x
\]
最小的曲线。
\end{advexamplebox}
\begin{advsolution}
密度函数不显含 $x$，可用贝尔特拉米恒等式
\[
 \mathcal L-y'\frac{\partial\mathcal L}{\partial y'}
 =\frac1{\sqrt{2gy(1+y'^2)}}=C_0 .
\]
令 $y=a(1-\cos\theta)$，则由上式得到
\[
 x=a(\theta-\sin\theta)+C_1.
\]
取起点为 $\theta=0$：
\[
 \boxed{x=a(\theta-\sin\theta),\qquad
 y=a(1-\cos\theta)}.
\]
这是半径为 $a$ 的圆滚动时圆周上一点描出的旋轮线。参数 $a$ 与终点参数 $\theta_1$ 由
$l=a(\theta_1-\sin\theta_1)$、$h=a(1-\cos\theta_1)$ 唯一确定适用分支。
\begin{center}\begin{minipage}{\linewidth}\centering
\begin{tikzpicture}[x=1.05cm,y=.85cm]
 \draw[axis] (0,0)--(4.2,0) node[right]{$x$};
 \draw[axis] (0,0)--(0,-2.65) node[below]{$y$（向下）};
 \draw[curve,domain=0.03:2.8,samples=120,smooth,variable=\t]
   plot ({\t-sin(\t r)},{-(1-cos(\t r))});
 \fill (0,0) circle(2pt) node[above left]{$O$};
 \fill ({2.8-sin(2.8 r)},{-(1-cos(2.8 r))}) circle(2pt) node[right]{$P$};
 \draw[guide] (0,{-(1-cos(2.8 r))})--({2.8-sin(2.8 r)},{-(1-cos(2.8 r))});
\end{tikzpicture}
\bookcaption{最速降线是旋轮线，而不是直线或圆弧}\label{advA1:fig:brachistochrone}
\end{minipage}\end{center}
图\ref{advA1:fig:brachistochrone}中轨道在起点具有竖直切线，使质点先迅速获得速度，再用较高速度完成横向运动。
\end{advsolution}

\begin{advexamplebox}{最小耗散原理}\label{advA1:ex:least-dissipation}
区域 $D$ 内电导率为 $\sigma(\vect r)>0$，给定非静电力场 $\vect K$。在固定非静电力功率
\[
 -\int_D \vect K\cdot\sigma\nabla\phi\,\dd V=P
\]
的条件下，求使给定二次泛函
$Q[\phi]=\int_D\sigma|\nabla\phi|^2\dd V$ 取极值的函数，并解释它与稳恒导电方程的形式关系，辨明何时可解释为物理耗散。
\end{advexamplebox}
\begin{advsolution}
对
\[
 \mathcal F[\phi]=\int_D\sigma|\nabla\phi|^2\dd V
 +\lambda\left(\int_D\vect K\cdot\sigma\nabla\phi\,\dd V+P\right)
\]
变分。若边界电势固定，则
\[
 \delta\mathcal F
 =-\int_D\nabla\cdot\{\sigma(2\nabla\phi+\lambda\vect K)\}
 \,\delta\phi\,\dd V .
\]
所以
\[
 \nabla\cdot\{\sigma(2\nabla\phi+\lambda\vect K)\}=0.
\]
把电势作常数重标度 $V=2\phi/(-\lambda)$，便得到
\[
 \nabla\cdot\vect j=0,\qquad
 \vect j=\sigma(\vect K-\nabla V),
\]
即欧姆定律与稳恒电荷守恒的形式；重标度也改变边界电势，且只在 $\lambda\ne0$ 时有效。二阶变分
$\delta^2Q=2\int_D\sigma|\nabla\delta\phi|^2\dd V>0$，所以在固定边界、线性约束的可行集合内这是最小值。$\lambda$ 仍须由给定的 $P$ 求定；一般没有 $Q_{\min}=P$。例如零边值的一维区间 $[0,1]$，取无量纲 $\sigma=1,K=x$，则约束为 $\int_0^1\phi\,\dd x=P$，极小解为 $\phi=6Px(1-x)$，且 $Q_{\min}=12P^2$。

对于重标度后的物理电流，真正的焦耳功率是
\[
 Q_J=\int_D\frac{j^2}{\sigma}\dd V
 =\int_D\sigma|\vect K-\nabla V|^2\dd V,
\]
并非原泛函 $Q[\phi]$。由连续性可得
\[
 Q_J=\int_D\vect K\cdot\vect j\,\dd V
 -\oint_{\partial D}V\vect j\cdot\dd\vect S.
\]
只有边界电功率项为零时，非静电力的输入功率才等于焦耳热。
\end{advsolution}

\begin{advexamplebox}{中心力场中的最速降线}\label{advA1:ex:central-brachistochrone}
质点从半径 $r_0$ 处静止释放。分别对
(1) 平方反比引力场；(2) 均匀球形天体内部的线性引力场，写出最速降线的确定方程。
\end{advexamplebox}
\begin{advsolution}
若速率只依赖 $r$，下降时间是光学型泛函
\[
 \mathcal T=\int n(r)\,\dd s,\qquad n(r)=\frac1{v(r)}.
\]
旋转对称性给出守恒量
\[
 b=n(r)r\sin\alpha,
\]
其中 $\alpha$ 是轨道切线与径向的夹角。因此所有中心力最速降线均由
\begin{equation}\label{advA1:eq:central-brach-quad}
 \frac{\dd\varphi}{\dd r}
 =\frac{b}{r\sqrt{n^2(r)r^2-b^2}}
\end{equation}
以及端点条件确定。参数 $b$ 的物理意义与分层介质中的斯涅耳常量相同。

(1) 对 $U=-GMm/r$，
\[
 v(r)=\sqrt{2GM\left(\frac1r-\frac1{r_0}\right)},\qquad
 n(r)=\left[2GM\left(\frac1r-\frac1{r_0}\right)\right]^{-1/2}.
\]
代入式\eqref{advA1:eq:central-brach-quad}得到径向单调段；对终点 $r_1<r_0$ 且无径向转折的候选轨道，用角位移积分
\[
 \Delta\varphi=\int_{r_1}^{r_0}
 \frac{b\,\dd r}{r\sqrt{r^2/[2GM(1/r-1/r_0)]-b^2}}
\]
求 $b$，再由 $\mathcal T=\int n\,\dd s$ 求时间。

(2) 均匀球内 $g(r)=\Omega_g^2r$，$\Omega_g^2=GM/R^3$。从 $r_0$ 静止释放时
\[
 v(r)=\Omega_g\sqrt{r_0^2-r^2},\qquad
 n(r)=\frac1{\Omega_g\sqrt{r_0^2-r^2}}.
\]
令 $q=\Omega_g b$，轨道满足
\[
 \frac{\dd\varphi}{\dd r}
 =\frac{q\sqrt{r_0^2-r^2}}
 {r\sqrt{(1+q^2)r^2-q^2r_0^2}}.
\]
上述导数的正负号须随径向运动方向选择。根号要求轨道不越过由
$n(r)r=b$ 决定的最近心点 $r_*$。若先向内后向外，角位移须改为
\[
 |\Delta\varphi|=\int_{r_*}^{r_0}F(r;b)\,\dd r
 +\int_{r_*}^{r_1}F(r;b)\,\dd r,
 \quad F(r;b)=\frac{|b|}{r\sqrt{n^2(r)r^2-b^2}}.
\]
同半径异角端点取 $r_1=r_0$ 时，两段贡献相加而不相消。角条件可能给出多个驻时轨道，还须比较正的总时间并检查轨道是否位于所用引力模型的区域；不能把欧拉方程的一个积分分支直接宣称为唯一最速轨道。
\end{advsolution}

对力学轨道 $q_\alpha(t)$，作用量为
\[
 S[q]=\int_{t_1}^{t_2}L(q,\dot q,t)\dd t.
\]

\begin{advtheorem}{哈密顿原理}\label{advA1:thm:hamilton-principle}
在固定起止时刻与端点的条件下，真实轨道使作用量的一阶变分为零：
\[
 \delta S=0.
\]
它给出的欧拉--拉格朗日方程正是分析力学的拉格朗日方程。
\end{advtheorem}
“一阶变分为零”比“全局最小”更准确，因为真实轨道也可能对应局部最大或鞍点。

\begin{advexamplebox}{曲面上的测地线}\label{advA1:ex:geodesic}
曲面 $\vect r=\vect r(u^1,u^2)$ 的线元为
\[
 \dd s^2=g_{ij}(u)\dd u^i\dd u^j.
\]
求参数曲线 $u^i=u^i(t)$ 为测地线的条件。
\end{advexamplebox}
\begin{advsolution}
若 $t$ 是仿射参数，可对能量泛函
\[
 S=\frac12\int g_{ij}\dot u^i\dot u^j\dd t
\]
变分。欧拉--拉格朗日方程为
\[
 \frac{\dd}{\dd t}(g_{kj}\dot u^j)
 -\frac12\partial_k g_{ij}\dot u^i\dot u^j=0.
\]
左乘逆度规 $g^{\ell k}$，得
\[
 \boxed{\ddot u^\ell+\Gamma^\ell{}_{ij}\dot u^i\dot u^j=0},
\qquad
 \Gamma^\ell{}_{ij}=\frac12g^{\ell k}
 (\partial_i g_{kj}+\partial_j g_{ki}-\partial_k g_{ij}).
\]
若用任意参数直接变分弧长，则右边可出现
$f(t)\dot u^\ell$；重新参数化为弧长或其他仿射参数后该项消失。
\end{advsolution}

\begin{advexamplebox}{极小曲面的条件}\label{advA1:ex:minimal-surface}
曲面 $\vect r=\vect r(u^1,u^2)$ 是某空间区域的边界。求固定边界曲线时它使面积取驻值的条件，并写出图形曲面 $z=f(x,y)$ 的方程。
\end{advexamplebox}
\begin{advsolution}
面积泛函为
\[
 A[\vect r]=\int\sqrt{\det g}\,\dd u^1\dd u^2.
\]
沿法向作变分 $\delta\vect r=\eta\vect n$，一阶面积变分为
\[
 \delta A=-2\int H\eta\,\dd A+\oint_{\partial S}(\cdots),
\]
固定边界使线积分消失。由于 $\eta$ 任意，极小曲面的必要条件为平均曲率
\[
 \boxed{H=0}.
\]
对 $z=f(x,y)$，
\[
 A[f]=\iint\sqrt{1+f_x^2+f_y^2}\,\dd x\dd y.
\]
代入欧拉--拉格朗日方程，得到
\[
 \boxed{(1+f_y^2)f_{xx}-2f_xf_yf_{xy}
 +(1+f_x^2)f_{yy}=0}.
\]
平面显然满足该方程；悬链面等非平凡曲面也满足，但球面因 $H\ne0$ 不是无约束极小曲面。
\end{advsolution}

\subsection{弹性体与场论守恒律}

先看一维弹性棒。位移场为 $\xi(x,t)$，质量线密度与等效弹性系数记为 $\rho,E$：
\[
 \mathcal L=\frac12\rho(\partial_t\xi)^2-\frac12E(\partial_x\xi)^2.
\]
场方程、波速与能量连续性方程分别为
\[
 \rho\,\partial_t^2\xi-E\,\partial_x^2\xi=0,\qquad
 c=\sqrt{E/\rho},
\]
\[
 w=\frac12\rho\xi_t^2+\frac12E\xi_x^2,\qquad
 \Pi=-E\xi_t\xi_x,\qquad
 \partial_t w+\partial_x\Pi=0.
\]
相应的机械动量密度与动量流为
\[
 g=\rho\xi_t,\qquad \mathcal T=-E\xi_x,\qquad
 \partial_tg+\partial_x\mathcal T=0.
\]

\begin{advtheorem}{场论的诺特定理}\label{advA1:thm:field-noether}
若标量场 $\phi_\nu$ 的连续变换
$\bar\phi_\nu=\Lambda_\nu(\phi,\lambda)$ 在 $\lambda=0$ 的生成元为
$\Psi_\nu=(\partial\Lambda_\nu/\partial\lambda)_0$，且拉格朗日密度的变化是散度
\[
 \left.\frac{\partial\bar{\mathcal L}}{\partial\lambda}\right|_0
 =\partial_\mu B^\mu,
\]
则
\[
 j^\mu=B^\mu-\sum_\nu
 \frac{\partial\mathcal L}{\partial(\partial_\mu\phi_\nu)}\Psi_\nu,
 \qquad \partial_\mu j^\mu=0.
\]
\end{advtheorem}
证明：把 $\partial_\lambda\mathcal L$ 用链式法则展开，对含
$\partial_\mu\Psi_\nu$ 的项作乘积求导，再用场方程
$\partial_\mu[\partial\mathcal L/\partial(\partial_\mu\phi_\nu)]
=\partial\mathcal L/\partial\phi_\nu$，余项正好整理为 $\partial_\mu j^\mu=0$。

若 $\mathcal L$ 不显含时空坐标，四个时空平移参数给出能量--动量张量
\[
 T^\mu{}_\nu=
 \sum_a\frac{\partial\mathcal L}{\partial(\partial_\mu\phi_a)}
 \partial_\nu\phi_a-\delta^\mu{}_\nu\mathcal L,
 \qquad \partial_\mu T^\mu{}_\nu=0.
\]
其时间分量是能量密度与能流，空间分量是动量密度与动量流。

\begin{advexamplebox}{弹性棒的平移对称性与机械动量}\label{advA1:ex:rod-noether}
证明一维弹性棒的拉格朗日密度不显含 $\xi$ 所对应的连续对称性产生机械动量守恒。
\end{advexamplebox}
\begin{advsolution}
变换 $\xi\mapsto\xi+\epsilon$ 不改变 $\xi_t,\xi_x$，故 $\delta\mathcal L=0$，生成元为 $\Psi=1$。
诺特流可整体乘以 $-1$ 重新约定为
\[
 j^0=\frac{\partial\mathcal L}{\partial\xi_t}=\rho\xi_t,
 \qquad
 j^1=\frac{\partial\mathcal L}{\partial\xi_x}=-E\xi_x.
\]
连续性方程
\[
 \partial_t(\rho\xi_t)+\partial_x(-E\xi_x)=0
\]
正是波动方程，也就是机械动量密度 $g=\rho\xi_t$ 与动量流
$\mathcal T=-E\xi_x$ 的守恒式。
\end{advsolution}

三维弹性体用拉格朗日描述
$\vect R(\vect r,t)=\vect r+\vect\xi(\vect r,t)$。变形后的基矢
$\vect\varepsilon_i=\partial\vect R/\partial x_i$ 构成度规
$g_{ij}=\vect\varepsilon_i\cdot\vect\varepsilon_j$。

\begin{advdefinition}{应变张量}\label{advA1:def:strain}
\[
 \varepsilon_{ij}=\frac12(g_{ij}-\delta_{ij}).
\]
在线性小形变下，对角元是正应变，非对角元是工程剪应变的一半。
\end{advdefinition}

\begin{advexamplebox}{纯剪切把单位圆变成椭圆}\label{advA1:ex:shear-circle}
二维小应变张量为
\[
 \bm\varepsilon=\begin{pmatrix}0&\theta/2\\\theta/2&0\end{pmatrix}.
\]
求单位圆 $x^2+y^2=1$ 的变形后曲线，并对角化应变张量。
\end{advexamplebox}
\begin{advsolution}
取不含刚性转动的代表变形梯度
$\mathsf F=\mathsf1+\bm\varepsilon$：
\[
 x'=x+\frac\theta2y,\qquad
 y'=y+\frac\theta2x.
\]
保留到 $\theta$ 一阶，逆变换为
$x=x'-\theta y'/2,\ y=y'-\theta x'/2$，所以
\[
 \boxed{x'^2+y'^2-2\theta x'y'=1+\order{\theta^2}}.
\]
取旋转坐标
$u=(x+y)/\sqrt2,\ v=(x-y)/\sqrt2$，有
\[
 \mathsf O\bm\varepsilon\mathsf O^{\mathsf T}
 =\operatorname{Diag}(\theta/2,-\theta/2).
\]
故精确到线性阶，主轴方向的伸长率分别为 $1+\theta/2$ 与
$1-\theta/2$：
\[
 \frac{u'^2}{(1+\theta/2)^2}
 +\frac{v'^2}{(1-\theta/2)^2}=1.
\]
\begin{center}\begin{minipage}{\linewidth}\centering
\begin{tikzpicture}[scale=1.0]
 \draw[axis] (-2.0,0)--(2.0,0) node[right]{$x'$};
 \draw[axis] (0,-1.55)--(0,1.55) node[above]{$y'$};
 \draw[guide] (0,0) circle(1.0);
 \begin{scope}[rotate=45]
  \draw[curve] (0,0) ellipse (1.36 and .74);
  \draw[guide] (-1.55,0)--(1.55,0) node[right]{$u$};
  \draw[guide] (0,-1.05)--(0,1.05) node[above]{$v$};
 \end{scope}
\end{tikzpicture}
\bookcaption{纯剪应变的主轴沿 $x=\pm y$，单位圆变为转过 $45^\circ$ 的椭圆}\label{advA1:fig:shear-ellipse}
\end{minipage}\end{center}
图\ref{advA1:fig:shear-ellipse}中虚线圆仅作变形前参照，标签没有用白底遮挡曲线。
\end{advsolution}

\begin{advtheorem}{应变张量的位移场表示}\label{advA1:thm:strain-displacement}
在线性小形变近似下，
\[
 \boxed{\bm\varepsilon=\frac12[
 \nabla\vect\xi+(\nabla\vect\xi)^{\mathsf T}]}.
\]
\end{advtheorem}
证明：$\mathsf F=\mathsf1+\nabla\vect\xi$，故
\[
 \mathsf g=\mathsf F^{\mathsf T}\mathsf F
 =\mathsf1+\nabla\vect\xi+(\nabla\vect\xi)^{\mathsf T}
 +(\nabla\vect\xi)^{\mathsf T}\nabla\vect\xi.
\]
忽略位移梯度的二阶项，再用
$\bm\varepsilon=(\mathsf g-\mathsf1)/2$ 即得。反对称部分
\[
 \bm\varphi=\frac12[
 \nabla\vect\xi-(\nabla\vect\xi)^{\mathsf T}]
\]
只描述质元的刚性小转动，不改变线元长度。

\begin{advtheorem}{$SO(3)$ 张量积的不可约分解}\label{advA1:thm:so3-decomposition}
三维旋转群的整数角动量不可约实表示维数为 $2\ell+1$。对
$m=2\ell_1+1\ge n=2\ell_2+1$，
\[
 m\otimes n=(m+n-1)\oplus(m+n-3)\oplus\cdots\oplus(m-n+1).
\]
特别地，
\[
 3\otimes3=5\oplus3\oplus1.
\]
\end{advtheorem}
对二阶张量 $\mathsf P$，这三个部分分别是
\[
 \mathsf P_{\rm ST}
 =\frac{\mathsf P+\mathsf P^{\mathsf T}}2
 -\frac{\operatorname{tr}\mathsf P}{3}\mathsf1,\quad
 \mathsf P_{\rm A}=\frac{\mathsf P-\mathsf P^{\mathsf T}}2,\quad
 \mathsf P_{\rm tr}=\frac{\operatorname{tr}\mathsf P}{3}\mathsf1.
\]
它们的维数依次为 $5,3,1$。对称无迹应变描述定体积形变，迹
$\operatorname{tr}\bm\varepsilon$ 是体积相对变化，反对称部分是刚性转动。

\begin{advexamplebox}{三阶张量的不可约分解计数}\label{advA1:ex:rank3-decomp}
证明维数分解
\[
 3\otimes3\otimes3=7\oplus2\cdot5\oplus3\cdot3\oplus1.
\]
\end{advexamplebox}
\begin{advsolution}
先耦合前两个矢量表示：
\[
 3\otimes3=5\oplus3\oplus1
 \quad(\ell=2,1,0).
\]
再分别与第三个 $\ell=1$ 的矢量耦合：
\[
 5\otimes3=7\oplus5\oplus3,\qquad
 3\otimes3=5\oplus3\oplus1,\qquad
 1\otimes3=3.
\]
合并同类项即得
$7\oplus2\cdot5\oplus3\cdot3\oplus1$。维数检查
$7+10+9+1=27=3^3$，没有遗漏分量。
\end{advsolution}

\AOneElasticSecond
\AOneFluidSection


\clearpage\endgroup

\begingroup
\newif\ifAdvHideHeading
\RenewDocumentEnvironment{advexamplebox}{m}{%
  \refstepcounter{advexample}%
  \begin{tcolorbox}[
    enhanced,unbreakable,colback=white,colframe=RuleGray,boxrule=0.5pt,arc=0mm,
    title={例~\theadvexample\quad #1},fonttitle=\bfseries,
    coltitle=Ink,colbacktitle=PaleBlue,borderline west={1.4pt}{0pt}{DeepBlue}]
}{\end{tcolorbox}}

\RenewDocumentEnvironment{advsolution}{+b}{%
  \par\smallskip\begingroup\small\noindent
  \ifAdvHideHeading\global\AdvHideHeadingfalse\else\LectureSolutionHeading\fi #1\par\endgroup\smallskip
}{}

\RenewDocumentEnvironment{advdefinition}{m +b}{%
  \refstepcounter{advdefinition}%
  \begin{tcolorbox}[
    enhanced,unbreakable,colback=white,colframe=RuleGray,boxrule=0.5pt,arc=0mm,
    title={定义~\theadvdefinition\quad #1},fonttitle=\bfseries,
    coltitle=Ink,colbacktitle=PaleGold,borderline west={1.2pt}{0pt}{Gold}]
  #2
  \end{tcolorbox}
}{}

\RenewDocumentEnvironment{advtheorem}{m +b}{%
  \refstepcounter{advtheorem}%
  \begin{tcolorbox}[
    enhanced,unbreakable,colback=white,colframe=RuleGray,boxrule=0.5pt,arc=0mm,
    title={定理~\theadvtheorem\quad #1},fonttitle=\bfseries,
    coltitle=Ink,colbacktitle=PaleBlue,borderline west={1.2pt}{0pt}{DeepBlue}]
  #2
  \end{tcolorbox}
}{}


\AdvancedUnit{2}{热学}

\section{热力学}

\subsection{热力学定律与基本关系}

热力学（thermodynamics）研究宏观体系的平衡态与状态变化。研究对象可以是易拉罐中的饮料、大气中的气团，也可以是外磁场中的超导体；选定对象称为\emph{体系}，其余部分称为外界。第一定律排除第一类永动机，第二定律排除把单一热源吸收的热量全部变成功的循环装置，第三定律则说明绝对零度不能由有限步骤达到。不同表述的逻辑起点虽有差别，但最终必须给出相同的可观测预言。

体系处于平衡态时，内部各子体系之间以及体系与外界之间均无宏观净流。体系具有内能 $U$；第 $i$ 种物质用粒子数 $N_i$ 表示，辐射等粒子数不守恒的成分则不必把粒子数列为独立变量。

\begin{advtheorem}{宏观态的完备参数组}\label{advA2:thm:macrostate}
一个平衡宏观态可由内能、宏观广义坐标和各组分粒子数完整描述：
\[
 A=A(U,\{X_\alpha\},\{N_i\}).
\]
若只保留一个体积坐标，则写成 $A=A(U,V,\{N_i\})$。
\end{advtheorem}
这里“完整”意味着参数空间与宏观状态空间之间是一一对应：不存在遗漏状态，同一组参数也不代表两个不同的宏观态。实际过程若不能视为连续平衡态序列，称为非准静态过程；全过程都经过平衡态的理想极限称为准静态过程。准静态还不自动等于可逆，只有不存在摩擦、有限温差传热等耗散时才可逆。

均匀体系中的量分为广延量和强度量。广延量随体系复制而同比放大，
\[
 A(\lambda U,\lambda V,\lambda N)=\lambda A(U,V,N),
\]
强度量则满足 $a(\lambda U,\lambda V,\lambda N)=a(U,V,N)$。内能、体积和粒子数既是广延量，又对互不相互作用的子体系具有加性；大量相同体系构成系综。孤立体系满足 $U,V,N$ 均不与外界交换。

\begin{advtheorem}{熵原理}\label{advA2:thm:entropy-principle}
平衡热力学可用一个加性态函数
\[
 S=S(U,V,N)
\]
组织。在本节所讨论的正温度体系中，熵随内能单调增加；孤立体系的可逆过程总熵不变，不可逆过程总熵增加。
\end{advtheorem}

\begin{advdefinition}{温度、压强与化学势}\label{advA2:def:intensives}
对单组分简单体系定义
\[
 \left(\frac{\partial S}{\partial U}\right)_{V,N}=\frac1T,
 \quad
 \left(\frac{\partial S}{\partial V}\right)_{U,N}=\frac pT,
 \quad
 \left(\frac{\partial S}{\partial N}\right)_{U,V}=-\frac\mu T.
\]
因此
\begin{equation}\label{advA2:eq:entropy-fundamental}
 \dd S=\frac{\dd U}{T}+\frac pT\dd V-\frac\mu T\dd N,
 \qquad
 \dd U=T\dd S-p\dd V+\mu\dd N.
\end{equation}
\end{advdefinition}
式\eqref{advA2:eq:entropy-fundamental}立即解释三个自发交换方向：热量从高温流向低温，机械平衡要求两侧压强相等，可交换粒子时平衡要求化学势相等。

\begin{advexamplebox}{小压强一侧为何仍可能膨胀}\label{advA2:ex:pressure-temperature}
两个子体系满足 $p_1<p_2$ 且 $T_1<T_2$。在固定各子体系内能与粒子数的熵虚变分中，解释为什么体积从体系2转给体系1仍可能使总熵增加；说明这是否就是自由绝热活塞的实际运动判据。
\end{advexamplebox}
\begin{advsolution}
令 $\dd V_1=-\dd V_2>0$，并按题设在状态空间中固定 $U_1,U_2,N_1,N_2$，则
\[
 \dd S_{\rm tot}
 =\left(\frac{p_1}{T_1}-\frac{p_2}{T_2}\right)\dd V_1.
\]
决定方向的是 $p/T$，不是单独的 $p$。例如取
$p_1=1,T_1=1,p_2=2,T_2=3$，虽然 $p_1<p_2$，却有
$p_1/T_1>p_2/T_2$，所以 $\dd V_1>0$ 对应熵增。单位检查：$p\dd V/T$ 的单位为能量除以温度，正是熵。这里比较的是受约束的邻近平衡态。实际绝热活塞移动时两侧一般还发生压强功引起的内能变化；总容器绝热并不意味着各侧 $\dd U_i=0$。因此不能据此断言自由活塞会自发向高压侧移动。
\end{advsolution}

\newcommand{\ASecondThermodynamicsLater}{%
\subsection{热力学函数}

\begin{advdefinition}{五个常用热力学势}\label{advA2:def:potentials}
对单组分体系，常用势及其自然变量为
\[
\begin{array}{c|c|c}
\text{势}&\text{定义}&\text{全微分}\\ \hline
U(S,V,N)&U&\dd U=T\dd S-p\dd V+\mu\dd N\\
H(S,p,N)&U+pV&\dd H=T\dd S+V\dd p+\mu\dd N\\
F(T,V,N)&U-TS&\dd F=-S\dd T-p\dd V+\mu\dd N\\
G(T,p,N)&U-TS+pV&\dd G=-S\dd T+V\dd p+\mu\dd N\\
\Omega(T,V,\mu)&F-\mu N&\dd\Omega=-S\dd T-p\dd V-N\dd\mu
\end{array}
\]
均匀宏观体系有 $\Omega=-pV$。源资料把 $-\Omega=pV$ 记作巨势；本卷采用统计物理中更常见的符号约定，并在后文保持一致。
\end{advdefinition}

\begin{advtheorem}{麦克斯韦关系}\label{advA2:thm:maxwell}
由各热力学势的混合偏导可交换，得到
\[
 \left(\frac{\partial T}{\partial V}\right)_S
 =-\left(\frac{\partial p}{\partial S}\right)_V,
 \quad
 \left(\frac{\partial T}{\partial p}\right)_S
 =\left(\frac{\partial V}{\partial S}\right)_p,
\]
\[
 \left(\frac{\partial S}{\partial V}\right)_T
 =\left(\frac{\partial p}{\partial T}\right)_V,
 \quad
 \left(\frac{\partial S}{\partial p}\right)_T
 =-\left(\frac{\partial V}{\partial T}\right)_p.
\]
\end{advtheorem}
例如最后两式分别来自 $F_{TV}=F_{VT}$ 与 $G_{Tp}=G_{pT}$。这一方法同时说明每个偏导中必须固定哪个自然变量，不能只凭符号记忆。

\begin{advexamplebox}{吉布斯--亥姆霍兹方程}\label{advA2:ex:gibbs-helmholtz}
证明
\[
 \boxed{\left[\frac{\partial(G/T)}{\partial T}\right]_p=-\frac{H}{T^2}}.
\]
\end{advexamplebox}
\begin{advsolution}
在定压下 $G_T=-S$，而 $G=H-TS$。因此
\[
 \left[\frac{\partial(G/T)}{\partial T}\right]_p
 =\frac{TG_T-G}{T^2}
 =\frac{-TS-(H-TS)}{T^2}=-\frac H{T^2}.
\]
两边单位均为熵除以温度。
\end{advsolution}

\begin{advexamplebox}{液膜表面张力与表面热力学}\label{advA2:ex:surface-thermodynamics}
面积为 $A$ 的液膜，其表面自由能为 $F_s=\gamma(T)A$。
\begin{enumerate}
 \item 求表面熵与表面内能。
 \item 若 $\gamma=C(T_c-T)$，等温增加面积 $\Delta A$，求外力功与吸热。
 \item 若 $\gamma=CT_c(1-T/T_c)^a$ 且 $1<a<2$，保持面积不变加热时表面层吸热还是放热？这是否意味着整个液体具有负热容？
\end{enumerate}
\end{advexamplebox}
\begin{advsolution}
表面功项为 $\gamma\dd A$，故
\[
 \dd F_s=-S_s\dd T+\gamma\dd A,
 \quad
 \boxed{S_s=-A\gamma'(T)},
 \quad
 \boxed{U_s=F_s+TS_s=A[\gamma-T\gamma']}.
\]
等温拉伸时
\[
 \boxed{W_{\rm ext}=\gamma\Delta A},
 \qquad
 \boxed{Q=T\Delta S_s=-T\gamma'\Delta A}.
\]
对线性定律，$W_{\rm ext}=C(T_c-T)\Delta A$，$Q=CT\Delta A>0$。

定面积时表面热容
\[
 C_A=T\left(\frac{\partial S_s}{\partial T}\right)_A
 =-TA\gamma''(T).
\]
当 $1<a<2$ 时 $\gamma''>0$，故表面\emph{贡献}为负，即升温时表面层向体相释放能量。这不等于完整孤立液体的总热容为负：表面层不是可独立隔离的宏观体系，体相热容与表面--体相能量交换必须同时计入。临界点附近还要考虑涨落，经验幂律也不能任意外推。
\end{advsolution}

\begin{advexamplebox}{重力场中理想气体与悬挂球争议}\label{advA2:ex:gravity-heat-capacity}
质量为 $m$ 的经典分子在 $0\le z\le H$ 内运动，外场势能为 $mgz$。除竖直势能外，每个分子的二次型自由度给出热容 $fk/2$。
\begin{enumerate}
 \item 求 $N$ 个分子的 $U(g,T)$ 与定 $g$ 热容 $C_g$，并比较 $g=0$ 与 $g\to\infty$。
 \item 对一般处在外场 $g$ 中的物体，令总质量为 $M$、质心高度为 $z(T,g)$，定义
 $\alpha_z=z^{-1}(\partial z/\partial T)_g$。推导 $(\partial C_g/\partial g)_T$。
 \item 地面上承托的球与天花板下悬挂的球吸收相同热量时，能否仅凭“膨胀时重力做功的正负”断言谁升温更多？
\end{enumerate}
\end{advexamplebox}
\begin{advsolution}
单粒子的竖直构型配分因子为
\[
 Z_z=\int_0^H e^{-mgz/(kT)}\dd z
 =\frac{1-e^{-x}}{\beta mg},
 \qquad x=\frac{mgH}{kT}.
\]
因此平均重力势能为
\[
 \langle mgz\rangle=-\frac{\partial\ln Z_z}{\partial\beta}
 =kT\left(1-\frac{x}{e^x-1}\right).
\]
于是
\[
 \boxed{U=N\left[\frac f2kT+kT\left(1-\frac{x}{e^x-1}\right)\right]},
\]
\[
 \boxed{C_g=Nk\left[\frac f2+1-\frac{x^2e^x}{(e^x-1)^2}\right]}.
\]
$x\to0$ 时括号中的重力贡献趋于零；$x\to\infty$ 时趋于1，所以
$C_g\to Nfk/2$ 与 $Nk(f/2+1)$。极限连续且热容为正。

把重力势能包括在体系能量中，亥姆霍兹自由能满足
\[
 \dd F=-S\dd T+Mz\dd g.
\]
混合偏导给出
\[
 \left(\frac{\partial S}{\partial g}\right)_T
 =-M\left(\frac{\partial z}{\partial T}\right)_g=-Mz\alpha_z.
\]
再对 $C_g=T(\partial S/\partial T)_g$ 求导：
\[
 \boxed{\left(\frac{\partial C_g}{\partial g}\right)_T
 =-MT\frac{\partial(z\alpha_z)}{\partial T}
 =-MT\left(\frac{\partial^2z}{\partial T^2}\right)_g}.
\]
直接把上面的气体质心 $z=(kT/mg)-H/(e^x-1)$ 代入，可复现 $C_g$ 对 $g$ 的变化。

对两个球，支承方式改变了应力分布、质心高度函数和热容。升温量应由完整的
$\Delta T=Q/C_g$ 决定，不能把“内能增加”与“温度增加”直接等同。若假设仅靠支承方向便能无条件得到固定温差，把两球交替接触热库并互换支承即可构造从单一热源循环取功的装置，违背第二定律。因此争议推理中不严谨的是忽略状态方程、弹性能和热容，仅按重力功符号比较温升。

\begin{center}\begin{minipage}{\linewidth}\centering
\begin{tikzpicture}[scale=.9]
 \begin{scope}[xshift=-2.7cm]
  \draw[ground] (-1.2,0)--(1.2,0);
  \draw[very thick,fill=DeepBlue!6] (0,.72) circle(.72);
  \draw[vector] (0,.72)--(0,.15) node[midway,right]{$Mg$};
  \node at (0,-.45) {地面承托};
 \end{scope}
 \begin{scope}[xshift=2.7cm]
  \draw[ground] (-1.2,2.25)--(1.2,2.25);
  \draw[very thick] (0,2.25)--(0,1.45);
  \draw[very thick,fill=DeepBlue!6] (0,.73) circle(.72);
  \draw[vector] (0,.73)--(0,.16) node[midway,right]{$Mg$};
  \node at (0,-.45) {绳悬挂};
 \end{scope}
\end{tikzpicture}
\bookcaption{支承方式改变约束和应力，不能只按重力功的正负判断温升}\label{advA2:fig:supported-balls}
\end{minipage}\end{center}
图\ref{advA2:fig:supported-balls}只表达约束差别；比较温升时必须为两种约束分别求 $z(T,g)$ 与 $C_g$。
\end{advsolution}

\subsection{相平衡与相变}

同一温度、压强下，相共存要求两相的化学势相等。对单组分物质，$\dd\mu=-s\dd T+v\dd p$，故沿共存线
\[
 \frac{\dd p}{\dd T}=\frac{s_2-s_1}{v_2-v_1}
 =\frac{\ell}{T(v_2-v_1)}.
\]
这里 $s,v,\ell$ 必须统一为每摩尔量或每单位质量量。分母很小时不能直接使用“液体体积可忽略”的蒸气近似。

\paragraph{由自由能得到等面积条件。}
等温线上若状态方程出现回折，机械稳定要求 $(\partial p/\partial v)_T<0$；斜率为正的部分不是可稳定保持的均匀相。
记摩尔 Helmholtz 自由能为 $f$，则 $\dd f=-p\dd v$。两端相同 $p_*$ 且 $\mu=f+p_*v$ 相等，给出
\[
 \int_{v_1}^{v_2}[p(v,T)-p_*]\dd v=0.
\]
这就是 Maxwell 等面积条件，确定的是两相共存端点；$(\partial p/\partial v)_T=0$ 给出的旋节边界只标记均匀相局部稳定性丧失。临界点处两端合并，常见光滑状态方程满足一、二阶体积导数同时为零。解题时应先问所求是共存、局部稳定，还是临界条件，三者不能互换。


\begin{advtheorem}{范德瓦耳斯流体的临界点与对比状态方程}\label{advA2:thm:vdw}
对
\[
 \left(p+\frac{N^2a}{V^2}\right)(V-Nb)=NkT,
\]
临界常数和对比状态方程为
\[
 \boxed{V_c=3Nb,\quad p_c=\frac{a}{27b^2},\quad kT_c=\frac{8a}{27b}},
\]
\[
 \boxed{\left(\pi+\frac3{v^2}\right)(3v-1)=8t},
 \quad v=\frac V{V_c},\ \pi=\frac p{p_c},\ t=\frac T{T_c}.
\]
\end{advtheorem}
证明：把压强写成
$p=NkT/(V-Nb)-N^2a/V^2$。临界等温线的水平拐点满足
$(\partial p/\partial V)_T=(\partial^2p/\partial V^2)_T=0$；联立即得三项临界常数。代入无量纲变量后，材料参数 $a,b$ 全部消失，得到对比状态方程。

两相平衡要求 $T_1=T_2,p_1=p_2,\mu_1=\mu_2$。一级相变中 $\mu$ 的一阶导数，即摩尔熵和摩尔体积，可以跃变；连续相变的一阶导数连续而二阶导数可跃变。

\begin{advtheorem}{克拉珀龙与埃伦费斯关系}\label{advA2:thm:clapeyron}
一级相变共存曲线满足
\[
 \boxed{\frac{\dd p}{\dd T}=\frac{\Delta s}{\Delta v}=\frac{L}{T\Delta v}}.
\]
连续相变在相界上若 $s,v$ 连续，则
\[
 \boxed{\frac{\dd p}{\dd T}=\frac{\Delta c_p}{Tv\Delta\alpha}
 =\frac{\Delta\alpha}{\Delta\kappa_T}}.
\]
\end{advtheorem}
证明：两相化学势相等，沿共存线微分
$-s_1\dd T+v_1\dd p=-s_2\dd T+v_2\dd p$，即得克拉珀龙式。连续相变时再分别对 $\Delta s=0$、$\Delta v=0$ 沿相界求全微分，利用
$s_T=c_p/T,s_p=-v\alpha,v_T=v\alpha,v_p=-v\kappa_T$，便得两条埃伦费斯关系。

\begin{advexamplebox}{超导临界场、潜热与热容突变}\label{advA2:ex:superconducting-transition}
正常--超导相变的热力学临界场满足
\[
 H_c(T)=H_0\left(1-\frac{T^2}{T_c^2}\right).
\]
求相变潜热，并求零场二级相变点的单位体积热容突变。
\end{advexamplebox}
\begin{advsolution}
采用SI制，凝聚能密度满足
\[
 g_n-g_s=\frac{\mu_0H_c^2}{2}.
\]
因此
\[
 s_n-s_s=-\mu_0H_c\frac{\dd H_c}{\dd T},
 \qquad
 L=T(s_n-s_s)=-\mu_0TH_cH_c'.
\]
代入给定函数，
\[
 \boxed{L(T)=2\mu_0H_0^2\frac{T^2}{T_c^2}
 \left(1-\frac{T^2}{T_c^2}\right)}.
\]
在零场端点 $T=T_c$，$H_c=0$，所以潜热为零；但
\[
 c_s-c_n=T\frac{\dd(s_s-s_n)}{\dd T}
 =\mu_0T[(H_c')^2+H_cH_c'']
\]
给出
\[
 \boxed{\Delta c\big|_{T_c}=\frac{4\mu_0H_0^2}{T_c}}.
\]
若采用高斯制，只需把 $\mu_0H_c^2/2$ 换成 $H_c^2/(8\pi)$；物理结论不变。
\end{advsolution}

平均场海森堡模型把每个自旋受到的近邻交换作用替换成与平均磁化强度成正比的分子场。对自旋 $1/2$，令单自旋磁矩 $\mu_s=\hbar\gamma/2$，饱和磁矩 $M_0=N\mu_s$。若每个自旋有 $n$ 个近邻，交换能标为 $\hbar J$，零外场自洽方程可写成
\begin{equation}\label{advA2:eq:weiss-self-consistent}
 m=\tanh\!\left(\frac{T_c}{T}m\right),
 \qquad m=\frac M{M_0},\qquad
 kT_c=\frac{n\hbar J}{4}.
\end{equation}

\begin{advexamplebox}{外斯模型的临界温度}\label{advA2:ex:weiss-critical}
由自洽方程在正温度下，证明自发磁化的非零解只在 $T<T_c=n\hbar J/(4k)$ 时出现，并检查 $T=T_c$。
\end{advexamplebox}
\begin{advsolution}
$m=0$ 总是解。临界点附近展开 $\tanh x=x-x^3/3+\cdots$：
\[
 m=\frac{T_c}{T}m-\frac13\left(\frac{T_c}{T}m\right)^3+\cdots.
\]
非零解要求线性项斜率 $T_c/T>1$，即 $T<T_c$。在 $T=T_c$，由于 $|\tanh m|<|m|$ 对所有 $m\ne0$ 成立，仍只有 $m=0$。$T>T_c$ 时右端曲线在原点斜率小于1，且始终被 $m$ 直线压住，故没有别的交点。
\end{advsolution}

\begin{advexamplebox}{铁磁相变曲线与临界指数}\label{advA2:ex:weiss-curve}
证明零场相变曲线可写成
\[
 \frac{T}{T_c}=\frac{2m}{\ln[(1+m)/(1-m)]},
\]
求 $M_0$ 并确定 $T\to T_c^-$ 时磁化强度的临界指数。
\end{advexamplebox}
\begin{advsolution}
对式\eqref{advA2:eq:weiss-self-consistent}取反双曲正切：
\[
 \frac{T_c}{T}m=\operatorname{artanh}m
 =\frac12\ln\frac{1+m}{1-m},
\]
即得题式。自旋全同向时
$\boxed{M_0=N\hbar\gamma/2}$。临界附近由上例展开式，
\[
 m^2\simeq3\left(1-\frac T{T_c}\right),
 \quad
 \boxed{M\propto(T_c-T)^{1/2}}.
\]
平均场磁化临界指数为 $1/2$。
\begin{center}\begin{minipage}{\linewidth}\centering
\begin{tikzpicture}[scale=.9]
 \draw[axis] (0,0)--(4.5,0) node[right]{$T/T_c$};
 \draw[axis] (0,0)--(0,2.8) node[above]{$M/M_0$};
 \draw[curve] plot[smooth] coordinates {
   (0,2.4) (.72,2.40) (1.44,2.37) (2.16,2.18)
   (2.88,1.70) (3.24,1.26) (3.42,.91) (3.60,0)};
 \draw[guide] (3.6,0)--(3.6,2.55);
 \node[below] at (3.6,0) {$1$};
 \node[left] at (0,2.4) {$1$};
\end{tikzpicture}
\bookcaption{平均场磁化强度在临界点连续降为零，近临界区呈平方根律}\label{advA2:fig:weiss-curve}
\end{minipage}\end{center}
\end{advsolution}

\begin{advexamplebox}{居里--外斯定律}\label{advA2:ex:curie-weiss}
在线性弱场区求磁化率，并证明 $1/\chi$ 与 $T-T_c$ 成正比。
\end{advexamplebox}
\begin{advsolution}
把分子场写成 $B_{\rm eff}=B+\lambda M$。弱场高温下
\[
 M=N\mu_s\tanh\frac{\mu_s(B+\lambda M)}{kT}
 \simeq\frac{N\mu_s^2}{kT}(B+\lambda M).
\]
定义 $T_c=N\mu_s^2\lambda/k$，解得
\[
 \boxed{\chi=\frac{\partial M}{\partial B}
 =\frac{N\mu_s^2/k}{T-T_c}},
 \qquad
 \boxed{\frac1\chi=\frac{k}{N\mu_s^2}(T-T_c)}.
\]
当 $\lambda\to0$ 时退化为居里定律 $\chi\propto1/T$。
\end{advsolution}

朗道理论用序参量 $\phi$ 的对称性组织相变。若体系在 $\phi\to-\phi$ 下不变，最低阶展开可写成
\[
 F=F_0+a(T-T_c)\phi^2+b\phi^4+\cdots,
 \qquad b>0.
\]
$T<T_c$ 时平衡序参量为
$\phi_\pm=\pm\sqrt{-a(T-T_c)/(2b)}$。熵与热容分别为
\[
 S=-\left(\frac{\partial F}{\partial T}\right)_\phi,
 \qquad C=T\left(\frac{\partial S}{\partial T}\right),
\]
从而在平均场近似中热容出现有限跃变
$\Delta C=a^2T_c/(2b)$。

\begin{advexamplebox}{六次朗道势与一级相变}\label{advA2:ex:landau-first-order}
证明当
\[
 F(\phi)=F_0+a\phi^2-b\phi^4+c\phi^6,
 \qquad a,b,c>0,
\]
随控制参量改变时可以发生一级相变，并求两相等自由能时的非零序参量与 $a$。
\end{advexamplebox}
\begin{advsolution}
令 $x=\phi^2\ge0$。非零驻点满足
\[
 a-2bx+3cx^2=0.
\]
一级相变共存时，$\phi=0$ 与非零极小值具有相同自由能，故还需
$ax-bx^2+cx^3=0$。两式联立得
\[
 \boxed{x=\phi^2=\frac b{2c}},
 \qquad
 \boxed{a=\frac{b^2}{4c}}.
\]
在共存点序参量从0跳到 $\sqrt{b/(2c)}$，故是一级相变；$c>0$ 保证大 $|\phi|$ 时自由能有下界。若加入空间非均匀与电磁耦合，得到朗道--金兹堡泛函
\[
 F=F_n+\alpha|\psi|^2+\frac\beta2|\psi|^4
 +\frac1{2m}\left|(-i\hbar\nabla-2e\vect A)\psi\right|^2
 +\frac{|\vect B|^2}{2\mu_0}.
\]
\end{advsolution}

}

\begin{advexamplebox}{欧拉齐次函数定理}\label{advA2:ex:euler-homogeneous}
若 $A(\lambda x_1,\ldots,\lambda x_n)=\lambda^kA(x_1,\ldots,x_n)$，证明
\[
 \boxed{kA=\sum_i x_i\frac{\partial A}{\partial x_i}}.
\]
\end{advexamplebox}
\begin{advsolution}
固定 $x_i$，对缩放参数求导：
\[
 \frac{\dd}{\dd\lambda}A(\lambda x_i)
 =\sum_i x_i\frac{\partial A}{\partial(\lambda x_i)}
 =k\lambda^{k-1}A(x_i).
\]
令 $\lambda=1$ 即得结论。对广延热力学量 $k=1$；对强度量 $k=0$，相应加权和为零。
\end{advsolution}

\begin{advexamplebox}{萨克--泰特洛德方程给出的理想气体热力学}\label{advA2:ex:sackur-tetrode}
单原子经典理想气体满足
\[
 \frac{S}{Nk}=\ln\!\left[
 \frac VN\left(\frac{4\pi m}{3h^2}\frac UN\right)^{3/2}
 \right]+\frac52.
\]
以 $N,V,T$ 为自变量求 $U,S,p,\mu$。
\end{advexamplebox}
\begin{advsolution}
由熵基本关系，
\[
 \frac1T=\left(\frac{\partial S}{\partial U}\right)_{V,N}
 =\frac{3Nk}{2U},
\]
故
\[
 \boxed{U=\frac32NkT},\qquad
 \boxed{p=T\left(\frac{\partial S}{\partial V}\right)_{U,N}
 =\frac{NkT}{V}}.
\]
代回熵式，
\[
 \boxed{S=Nk\left\{\ln\!\left[
 \frac VN\left(\frac{2\pi mkT}{h^2}\right)^{3/2}\right]+\frac52\right\}}.
\]
再由 $\mu=-T(\partial S/\partial N)_{U,V}$，或用下述欧拉关系，得到
\[
 \boxed{\mu=-kT\ln\!\left[
 \frac VN\left(\frac{2\pi mkT}{h^2}\right)^{3/2}\right]}.
\]
$U$ 与 $pV$ 均与 $NkT$ 同量纲；稀薄极限 $V/N\to\infty$ 时 $\mu\to-\infty$，符合经典气体极限。
\end{advsolution}

熵式和能量式分别称为熵表象与能量表象的基本关系。绝对零点附近还需第三定律约束积分常数。

\begin{advtheorem}{热力学第三定律的能量表述}\label{advA2:thm:third-law}
在基态不具有宏观简并的通常情形，$T\to0$ 时熵趋于与外参量无关的常数，可取为零，且
\[
 \lim_{S\to0}\left(\frac{\partial U}{\partial S}\right)_{V,N}=0,
 \qquad
 U(S,V,N)=U_0(V,N)+\int_0^S T\,\dd S.
\]
\end{advtheorem}
广延性还给出 $U(S,V,N)=Nu(S/N,V/N)$，以后用小写字母表示每粒子量。

\begin{advexamplebox}{化学势关系与吉布斯--杜亥姆关系}\label{advA2:ex:gibbs-duhem}
证明
\[
 \mu=u-Ts+pv,
 \qquad U=TS-pV+\mu N,
\]
并写出吉布斯--杜亥姆关系。
\end{advexamplebox}
\begin{advsolution}
$U(S,V,N)$ 是一次齐次函数。对它应用例\ref{advA2:ex:euler-homogeneous}，再用
$U_S=T,U_V=-p,U_N=\mu$，得
\[
 \boxed{U=TS-pV+\mu N}.
\]
除以 $N$ 即有 $\mu=u-Ts+pv$。对欧拉式求微分，并与
$\dd U=T\dd S-p\dd V+\mu\dd N$ 比较，消去同类项：
\[
 \boxed{S\dd T-V\dd p+N\dd\mu=0}.
\]
这说明三个强度量不能任意独立改变。
\end{advsolution}

对定粒子数体系，焓 $H=U+pV$ 满足 $\dd H=T\dd S+V\dd p$。定义
\[
 C_p=T\left(\frac{\partial S}{\partial T}\right)_p
 =\left(\frac{\partial H}{\partial T}\right)_p,
 \quad
 C_V=T\left(\frac{\partial S}{\partial T}\right)_V
 =\left(\frac{\partial U}{\partial T}\right)_V,
\]
\[
 \alpha=\frac1V\left(\frac{\partial V}{\partial T}\right)_p,
 \quad
 \beta=\frac1p\left(\frac{\partial p}{\partial T}\right)_V,
 \quad
 \kappa_T=-\frac1V\left(\frac{\partial V}{\partial p}\right)_T,
 \quad
 \kappa_S=-\frac1V\left(\frac{\partial V}{\partial p}\right)_S.
\]

\begin{advtheorem}{能态定理}\label{advA2:thm:energy-state}
\[
 \boxed{\left(\frac{\partial U}{\partial V}\right)_T+p
 =T\left(\frac{\partial p}{\partial T}\right)_V}.
\]
\end{advtheorem}
证明：令 $F=U-TS$，则
$\dd F=-S\dd T-p\dd V$，因而
$(\partial S/\partial V)_T=(\partial p/\partial T)_V$。又
\[
 \left(\frac{\partial U}{\partial V}\right)_T
 =\left(\frac{\partial F}{\partial V}\right)_T
 +T\left(\frac{\partial S}{\partial V}\right)_T
 =-p+T\left(\frac{\partial p}{\partial T}\right)_V.
\]

\begin{advexamplebox}{膨胀系数、压强系数与等温压缩系数}\label{advA2:ex:alpha-beta-kappa}
证明 $\alpha=\beta\kappa_Tp$。
\end{advexamplebox}
\begin{advsolution}
把状态方程写成 $V=V(T,p)$。在定体积线上
\[
 0=\left(\frac{\partial V}{\partial T}\right)_p\dd T
 +\left(\frac{\partial V}{\partial p}\right)_T\dd p,
\]
所以
\[
 \left(\frac{\partial p}{\partial T}\right)_V
 =-\frac{V\alpha}{-V\kappa_T}=\frac{\alpha}{\kappa_T}.
\]
左侧又等于 $p\beta$，故 $\boxed{\alpha=p\beta\kappa_T}$。理想气体给出
$\alpha=1/T,\beta=1/T,\kappa_T=1/p$，校验成立。
\end{advsolution}

\begin{advexamplebox}{热容比与压缩系数比}\label{advA2:ex:cp-cv-kappa}
证明
\[
 \boxed{\frac{C_p}{C_V}=\frac{\kappa_T}{\kappa_S}}.
\]
\end{advexamplebox}
\begin{advsolution}
把熵写成 $S(T,V)$。由麦克斯韦关系
$(\partial S/\partial V)_T=(\partial p/\partial T)_V$，
\[
 \dd S=\frac{C_V}{T}\dd T
 +\left(\frac{\partial p}{\partial T}\right)_V\dd V.
\]
沿绝热线令 $\dd S=0$，再与状态方程微分联立，可得
\[
 \kappa_S=\kappa_T-\frac{TV\alpha^2}{C_p}.
\]
另一方面
$C_p-C_V=TV\alpha^2/\kappa_T$。消元便得
$\kappa_S=\kappa_TC_V/C_p$，即所求。稳定体系中
$C_p\ge C_V>0$，所以 $\kappa_T\ge\kappa_S>0$，方向也正确。
\end{advsolution}

\begin{advexamplebox}{两种热容之差}\label{advA2:ex:cp-minus-cv}
证明
\[
 \boxed{C_p-C_V=\alpha\beta pVT}.
\]
\end{advexamplebox}
\begin{advsolution}
由
\[
 C_p-C_V=T\left[\left(\frac{\partial S}{\partial T}\right)_p
 -\left(\frac{\partial S}{\partial T}\right)_V\right]
\]
及链式法则，
\[
 C_p-C_V
 =T\left(\frac{\partial S}{\partial V}\right)_T
 \left(\frac{\partial V}{\partial T}\right)_p
 =T(p\beta)(V\alpha).
\]
因此得到题式。用例\ref{advA2:ex:alpha-beta-kappa}也可写成
$C_p-C_V=TV\alpha^2/\kappa_T\ge0$；理想气体代入后为 $Nk$，与
$C_p=5Nk/2,C_V=3Nk/2$ 一致。
\end{advsolution}

\ASecondThermodynamicsLater
% AdvancedThermalQuality20261023-v1
% Editable statistical mechanics figures and the quantitative Bose heat curve.
\providecommand{\ThermalFigure}[3]{\begin{qualityfigure}\begin{tikzpicture}[>=Latex,font=\small,x=.8cm,y=.8cm]#1\end{tikzpicture}\figcap{#2}\label{#3}\end{qualityfigure}}
\newcommand{\AQMicrocanonical}{\ThermalFigure{
 \draw[->](-2.5,0)--(2.5,0)node[right]{$q$};\draw[->](0,-1.7)--(0,1.7)node[above]{$p$};
 \draw[DeepBlue,thick](0,0)ellipse(2 and 1.2);\draw[DeepBlue,thick](0,0)ellipse(1.7 and 1);\node[right]at(3,1){能壳上的相空间体积 $\omega$};
}{只在给定能壳上归一化概率密度；壳外不属于本次熵极值问题。}{advA2:quality:micro-q}}
\newcommand{\AAMicrocanonical}{\ThermalFigure{
 \foreach\x in {0,1,2,3}{\draw[fill=DeepBlue!12](\x,0)rectangle(\x+.8,1.2);}
 \foreach\x/\h in {6/.4,7/2,8/1.6,9/.8}{\draw[fill=TextGray!15](\x,0)rectangle(\x+.8,\h);}
 \node[below]at(1.9,-.4){等体积小胞：均匀密度};\node[below]at(7.9,-.4){同一归一化：非均匀密度};
}{把同一能壳分成等体积小胞，两组柱面积相同，仅分布不同。}{advA2:quality:micro-a}
图~\ref{advA2:quality:micro-a}比较的是同一能壳和同一测度下的密度；若改变能壳体积，就不是同一个极值问题。相对熵非负可由 $x\ln x-x+1\ge0$ 对均匀参考测度积分得到，等号要求密度比几乎处处等于1，因而也给出最大值的唯一性。
}
\newcommand{\AQGrand}{\ThermalFigure{
 \draw[thick](0,0)rectangle(4,2);\node at(2,1){系统，$V$ 固定};\node[draw](R)at(8,1){巨大热、粒子库 $T,\mu$};
 \draw[<->,DeepBlue,thick](4,1.5)--(6,1.5)node[midway,above]{能量};\draw[<->,thick](4,.5)--(6,.5)node[midway,below]{粒子};
}{巨正则体系允许能量和粒子数涨落，但体积作为外控参数固定。}{advA2:quality:grand-q}}
\newcommand{\AAGrand}{\ThermalFigure{
 \node[draw](X)at(4,1){$\ln\Xi(\mu,V,T)$};\node[draw](N)at(0,0){$\partial_\mu\ln\Xi$};\node[draw](P)at(4,-1){$\partial_V\ln\Xi$};\node[draw](S)at(8,0){$\partial_T\ln\Xi$};
 \draw[->](X)--(N);\draw[->](X)--(P);\draw[->](X)--(S);\node[above]at(4,2){每次偏导都须固定其余自然变量};
}{粒子数、压强、熵来自同一巨配分函数的不同参数响应。}{advA2:quality:grand-a}
图~\ref{advA2:quality:grand-a}中压强的一般公式是 $p=kT(\partial\ln\Xi/\partial V)_{T,\mu}$。只有均匀、广延的体相体系，忽略边界修正并满足 $\Omega=-pV$ 时，才可进一步写成 $p=(kT/V)\ln\Xi$。固定逸度求温度导数与固定化学势求导不同，因为 $z=e^{\beta\mu}$ 本身含温度。
}
\newcommand{\AQDensityStates}{\ThermalFigure{
 \draw[->](-2,0)--(2,0)node[right]{$p$};\draw[very thick](-1.2,0)--(1.2,0);\node[below]at(0,-.5){一维动量区间};
 \begin{scope}[shift={(5,0)}]\draw[thick](0,0)circle(1.3);\draw[DeepBlue](0,0)circle(1.1);\node[below]at(0,-1.7){二维动量圆壳};\end{scope}
 \begin{scope}[shift={(10,0)}]\draw[thick](0,0)circle(1.3);\draw[dashed](0,0)ellipse(1.3 and .4);\draw[->](0,0)--(45:1.3)node[right]{$p$};\node[below]at(0,-1.7){三维动量球壳};\end{scope}
}{先数能量不超过给定值的动量态，再由壳体积求能量态密度。}{advA2:quality:dos-q}}
\newcommand{\AADensityStates}{\ThermalFigure{
 \draw[->](0,0)--(4,0)node[right]{$\varepsilon$};\draw[->](0,0)--(0,3)node[above]{$g$};
 \draw[DeepBlue,domain=.25:3.5,samples=60]plot(\x,{1/sqrt(\x)});\draw[thick](.25,1)--(3.5,1);\draw[TextGray,domain=.25:3.5,samples=60]plot(\x,{sqrt(\x)});
 \node[right,DeepBlue]at(3.5,.53){$d=1$};\node[right]at(3.5,1){$d=2$};\node[right,TextGray]at(3.5,1.87){$d=3$};
 \node[above]at(2,3){非相对论：指数 $d/2-1$};
 \begin{scope}[shift={(7,0)}]\draw[->](0,0)--(4,0)node[right]{$\varepsilon$};\draw[->](0,0)--(0,3)node[above]{$g$};\draw[thick](.1,.5)--(3,.5);\draw[DeepBlue](0,0)--(3,1.5);\draw[TextGray,domain=0:3,samples=60]plot(\x,{.25*\x*\x});\node[above]at(2,3){无质量：指数 $d-1$};\node[right]at(3,.5){$d=1$};\node[right,DeepBlue]at(3,1.5){$d=2$};\node[right,TextGray]at(3,2.25){$d=3$};\end{scope}
}{两种色散均按 $d=1,2,3$ 比较幂律，曲线归一化系数仅为示意。}{advA2:quality:dos-a}
图~\ref{advA2:quality:dos-a}中的幂指数由两个因子合成：动量球壳给 $p^{d-1}\dd p$，从动量改成能量还要乘 $\dd p/\dd\varepsilon$。所以非相对论得到 $\varepsilon^{d/2-1}$，无质量粒子得到 $\varepsilon^{d-1}$；直接把球壳半径改名为能量会漏掉雅可比。
}
\newcommand{\AQFreePartition}{\ThermalFigure{
 \draw[->](0,0)--(4,0)node[right]{$p$};\draw[->](0,0)--(0,3)node[above]{$\varepsilon$};\draw[DeepBlue,thick,domain=0:3,samples=50]plot(\x,{.25*\x*\x});\node[right]at(4,1.5){$p^2/(2m)$};
 \begin{scope}[shift={(7,0)}]\draw[->](0,0)--(4,0)node[right]{$p$};\draw[->](0,0)--(0,3)node[above]{$\varepsilon$};\draw[DeepBlue,thick](0,0)--(3,2.4);\node[right]at(3,2.4){$pc$};\end{scope}
}{两种单粒子色散关系；三维空间和统计权重保持相同。}{advA2:quality:free-q}}
\newcommand{\AAFreePartition}{\ThermalFigure{
 \draw[->](0,0)--(6,0)node[right]{$\beta\varepsilon$};\draw[->](0,0)--(0,3)node[above]{积分权重};
 \draw[DeepBlue,thick,domain=.001:5.5,samples=90]plot(\x,{3*sqrt(\x)*exp(-\x)});
 \draw[thick,domain=0:5.5,samples=90]plot(\x,{2*\x*\x*exp(-\x)});
 \node[right,align=left]at(6.5,1.6){非相对论：$\varepsilon^{1/2}e^{-\beta\varepsilon}$\\无质量：$\varepsilon^2e^{-\beta\varepsilon}$};
}{配分积分把态密度与玻尔兹曼权重相乘，能量平均再多乘一因子 $\varepsilon$。}{advA2:quality:free-a}
图~\ref{advA2:quality:free-a}中两权重的幂次不同。令 $u=\beta\varepsilon$ 后，积分的温度依赖全部提出为 $T^{3/2}$ 或 $T^3$，而无量纲积分为常数；对 $\ln z_1$ 求温度或 $\beta$ 导数便得到对应均能。对这里的固定粒子数模型，两者均有 $pV=NkT$。
}
\newcommand{\AQClassicalEntropy}{\ThermalFigure{
 \draw[thick](0,0)rectangle(5,2);\foreach\x/\y in {1/.6,2.7/1.4,4/.5}{\fill[DeepBlue](\x,\y)circle(3pt);}
 \draw[<->,thick](1,.9)..controls(1.8,2.4)and(2.8,2.4)..(2.7,1.7);\node[right]at(6,1){交换全同粒子的标号};
}{同一物理构形可被不同粒子标号重复记录，经典全同粒子配分函数须处理这一重复。}{advA2:quality:classical-s-q}}
\newcommand{\AAClassicalEntropy}{\ThermalFigure{
 \node[draw,align=center](L)at(2,1){三粒子标号的六种排列\\$123,132,213,231,312,321$};\node[draw](R)at(9,1){一个无标号构形};\draw[->,thick](L)--node[above]{除以 $3!$}(R);
}{$N!$ 修正将标号重复还原为物理态计数，而不是删除真实可区分的构型。}{advA2:quality:classical-s-a}
图~\ref{advA2:quality:classical-s-a}解释 $1/N!$ 的来源。用斯特林公式 $\ln N!\simeq N\ln N-N$ 后，$S$ 才具有热力学极限下的广延性；有限 $N$ 仍有次广延修正。展开成 $\ln T$、$\ln(V/N)$ 的写法只是选定单位后的缩写，严谨的量纲检查应使用完整的无量纲组合 $V/(N\lambda_T^3)$。
}
\newcommand{\AQOscillators}{\ThermalFigure{
 \foreach\s/\i in {0/1,4/2,8/N}{\begin{scope}[shift={(\s,0)}]\draw[thick](0,0)--(0,1.2);\draw[decorate,decoration={coil,segment length=2mm,amplitude=1mm},thick](0,.6)--(1.7,.6);\draw[thick](1.7,.1)rectangle(2.5,1.1);\node at(2.1,.6){$\i$};\end{scope}}
 \node[above]at(5,2){固定在不同位置、彼此独立的振子，频率均为 $\omega$};
}{振子由不同位置标识，统计计数不能照搬自由全同粒子的 $1/N!$。}{advA2:quality:osc-q}}
\newcommand{\AAOscillators}{\ThermalFigure{
 \draw[->](-2.5,0)--(2.5,0)node[right]{$q$};\draw[->](0,-1.8)--(0,1.8)node[above]{$p$};\draw[DeepBlue,thick](0,0)ellipse(2 and 1.2);
 \node[right,align=left]at(3,1){$H=p^2/(2m)+m\omega^2q^2/2$\\等能线为椭圆\\两个独立二次项各贡献 $kT/2$};
}{相空间高斯积分在 $q,p$ 两方向因子化，再对独立振子取乘积。}{advA2:quality:osc-a}
图~\ref{advA2:quality:osc-a}中的椭圆不是位置空间轨道，而是单振子的相空间等能线。两维高斯积分分别产生 $T^{1/2}$，故 $z\propto T$；$N$ 个有固定标识的位置对应 $Z=z^N$，而不是再除以 $N!$。
}
\newcommand{\AQRotors}{\ThermalFigure{
 \draw[fill=DeepBlue!10,thick](0,0)ellipse(1.8 and 1);\foreach\x/\y/\lab in {3/0/1,0/2/2,-1.8/-1.5/3}{\draw[->,thick](0,0)--(\x,\y)node[right]{$I_{\lab}$ 主轴};}
 \node[right]at(4.5,1){非线性刚体，三个主惯量均非零};
}{主轴是固定于刚体的坐标方向，取向还需用整个旋转群计数。}{advA2:quality:rotor-q}}
\newcommand{\AARotors}{\ThermalFigure{
 \draw[thick](0,0)ellipse(2 and 1);\draw[dashed](0,0)ellipse(2 and .35);\draw[->](0,0)--(3,0)node[right]{$L_1$};\draw[->](0,0)--(0,2)node[above]{$L_2$};\draw[->](0,0)--(-1.5,-1.6)node[left]{$L_3$};
 \node[right,align=left]at(4,1){$H=\sum_iL_i^2/(2I_i)$\\角动量积分给 $T^{3/2}\sqrt{I_1I_2I_3}$\\取向体积为 $8\pi^2/\sigma$};
}{角动量空间等能面为椭球，取向对称数另外消除同一取向的重复。}{advA2:quality:rotor-a}
图~\ref{advA2:quality:rotor-a}中三个主轴二次项各独立贡献一个高斯积分。取向对称数 $\sigma$ 处理分子自身的旋转等价，$N!$ 则处理不同全同分子交换，两者不是同一种重复，不能相互替代。线性分子的轴向惯量为零，必须改用两个转动自由度的模型。
}
\newcommand{\AQDipoleStatistics}{\ThermalFigure{
 \draw[->,thick](0,0)--(0,3)node[above]{$\vect E$};\draw[->,DeepBlue,thick](0,0)--(2,1.7)node[right]{$\vect d$};\draw(0,1)arc(90:40.4:1)node[right]{$\theta$};\node[right]at(4,1.5){只计取向能 $-dE\cos\theta$};
}{偶极取向相对外场的极角为 $\theta$，方位角仍需积分。}{advA2:quality:dip-stat-q}}
\newcommand{\AADipoleStatistics}{\ThermalFigure{
 \draw[thick](0,0)circle(1.5);\draw[DeepBlue,thick](-1.2,.9)--(1.2,.9);\draw[->](0,0)--(1.2,.9);\draw[->](0,0)--(0,2)node[above]{$\vect E$};
 \node[right,align=left]at(3,1){取向带测度 $\dd\Omega=2\pi\sin\theta\dd\theta$\\令 $u=\cos\theta$，积分上下限为 $-1,1$};
}{取向球上的纬带面积与玻尔兹曼权重相乘；角度并不是均匀分布在 $[0,\pi]$。}{advA2:quality:dip-stat-a}
图~\ref{advA2:quality:dip-stat-a}中的 $\sin\theta$ 来自球面面积元，不能漏掉。弱场展开时球面对称使 $\langle\cos\theta\rangle_0=0$、$\langle\cos^2\theta\rangle_0=1/3$，直接给出线性响应中的 $1/3$，也是朗之万函数展开的独立核验。
}
\newcommand{\AQQuantumDistributions}{\ThermalFigure{
 \foreach\s/\lab in {0/费米态,6/玻色态}{\begin{scope}[shift={(\s,0)}]\draw[thick](-1,0)--(2,0);\node[above]at(.5,.7){能级 $\varepsilon_i$，\lab};\end{scope}}
 \node[below]at(.5,-.6){占据数只能为 $0,1$};\node[below]at(6.5,-.6){占据数可为 $0,1,2,\ldots$};
}{两种统计的区别先体现在同一个单粒子态允许多少粒子占据。}{advA2:quality:quantum-q}}
\newcommand{\AAQuantumDistributions}{\ThermalFigure{
 \draw[->](0,0)--(5,0)node[right]{$x=(\varepsilon-\mu)/(kT)$};\draw[->](0,0)--(0,4)node[above]{$\langle n\rangle$};
 \draw[DeepBlue,thick,domain=0:4.5,samples=70]plot(\x,{1/(exp(\x)+1)});
 \draw[thick,domain=.25:4.5,samples=100]plot(\x,{1/(exp(\x)-1)});
 \draw[dashed,domain=0:4.5,samples=70]plot(\x,{exp(-\x)});\node[right,align=left]at(6,2){上：玻色\\中：经典极限 $e^{-x}$\\下：费米（图中 $x>0$）};
}{同一能级的平均占据；两种量子分布在 $x\gg1$ 时趋于同一经典尾部。}{advA2:quality:quantum-a}
图~\ref{advA2:quality:quantum-a}中玻色支仅在 $x>0$ 有有限正占据。对有限系统的单基态几何级数，严格收敛要求 $\mu<\varepsilon_0$；热力学极限中的凝聚态需单独处理，才可说 $\mu$ 从下方趋于基态能。费米支则自动限制每个态的占据不超过1。
}
\newcommand{\AQFermiCapacity}{\ThermalFigure{
 \draw[->](0,0)--(6,0)node[right]{$\varepsilon$};\draw[->](0,0)--(0,2.5)node[above]{$f$};\draw[thick](0,2)--(3,2)--(3,0)--(5.5,0);\node[below]at(3,0){$\varepsilon_F$};\node[above]at(3,2.8){$T=0$ 的填充基态};
}{定粒子数费米气体的零温参考态；满占据以下粒子不能都独立获得经典热激发。}{advA2:quality:fermi-c-q}}
\newcommand{\AAFermiCapacity}{\ThermalFigure{
 \draw[->](0,0)--(6,0)node[right]{$\varepsilon$};\draw[->](0,0)--(0,2.5)node[above]{$f$};\draw[dashed](0,2)--(3,2)--(3,0)--(5.5,0);
 \draw[DeepBlue,thick,domain=0:5.7,samples=100]plot(\x,{2/(exp(3*(\x-3))+1)});\draw[<->](2.3,2.7)--node[above]{$O(kT)$}(3.7,2.7);\node[below]at(3,0){$\mu\simeq\varepsilon_F$};
}{只有费米面附近的窄能窗参与低温热激发，窗外分布几乎不变。}{advA2:quality:fermi-c-a}
图~\ref{advA2:quality:fermi-c-a}中的平滑阶跃还须保持总粒子数不变。索末菲展开分别给出
\[
 0=g(\varepsilon_F)\delta\mu+\frac{\pi^2}{6}(kT)^2g'(\varepsilon_F)+O(T^4),
\]
再对能量积分展开并消去 $\delta\mu$，只剩 $\delta U=(\pi^2/6)g(\varepsilon_F)(kT)^2$。这样明确了定 $N$ 条件的作用；若固定化学势，不能不加检查地使用同一展开式。
}
\newcommand{\AQBoseCapacity}{\ThermalFigure{
 \draw[thick](0,0)--(5,0);\draw[thick](0,1)--(5,1) (0,1.8)--(5,1.8) (0,2.4)--(5,2.4);
 \foreach\x in {.5,1,1.5,2,2.5}{\fill[DeepBlue](\x,0)circle(2pt);}\foreach\x/\y in {3/1,4/1.8,2.8/2.4}{\fill(\x,\y)circle(2pt);}
 \node[right]at(6,0){基态 $N_0$};\node[right]at(6,1.8){激发态 $N-N_0$};
}{三维均匀理想玻色气体，将基态占据与激发态连续谱分开。}{advA2:quality:bose-c-q}}
\newcommand{\AABoseCapacity}{\ThermalFigure{
 \draw[->](0,0)--(5.2,0)node[right]{$T/T_c$};\draw[->](0,0)--(0,2.5)node[above]{$C_V/(Nk)$};
 \draw[DeepBlue,thick,domain=.001:1,samples=80]plot({1.5*\x},{1.92567168*\x^(1.5)});
 \draw[DeepBlue,thick]plot coordinates{(1.500000,1.92567168)(1.515000,1.91801858)(1.537500,1.90697113)(1.575000,1.88963169)(1.650000,1.85850730)(1.800000,1.80760611)(2.025000,1.75118594)(2.250000,1.71034420)(2.625000,1.66310220)(3.000000,1.63138079)(3.750000,1.59211964)(4.500000,1.56923444)};
 \draw[dashed](1.5,0)--(1.5,2.1);\node[below]at(1.5,0){$1$};\node[below]at(3,0){$2$};\node[below]at(4.5,0){$3$};\draw[dashed](0,1.5)--(5,1.5)node[right]{$3/2$};
}{Bose积分给出的定粒子数热容；上支数值解粒子数方程，未用经验指数替代。}{advA2:quality:bose-c-a}
图~\ref{advA2:quality:bose-c-a}中的临界以上曲线须同时考虑逸度随温度改变。定义 $g_s(z)=\operatorname{Li}_s(z)$，在 $T>T_c$ 有
\[
 \frac NV=\lambda_T^{-3}g_{3/2}(z),\qquad
 \frac{C_V}{Nk}=\frac{15}{4}\frac{g_{5/2}(z)}{g_{3/2}(z)}
 -\frac94\frac{g_{3/2}(z)}{g_{1/2}(z)}.
\]
定 $N,V$ 求导时，利用 $z\,g_s'(z)=g_{s-1}(z)$ 消去 $\dd z/\dd T$，便得到第二项。$z\to1^-$ 时 $g_{1/2}$ 发散，第二项趋零，与下支连续相接；$z\to0$ 时各 $g_s\simeq z$，回到经典极限 $3/2$。
}
\newcommand{\AQQuantumVirial}{\ThermalFigure{
 \foreach\x in {0,3,6,9}{\fill[DeepBlue](\x,0)circle(2pt);\draw[dashed](\x,0)circle(.7);}
 \draw[<->](0,-1.3)--node[below]{平均间距 $n^{-1/3}$}(3,-1.3);\draw[<->](-.7,1.2)--node[above]{$\lambda_T$}(.7,1.2);
}{热波长与平均间距决定占据稀疏程度，内禀简并度 $g$ 还会分摊每态粒子数。}{advA2:quality:virial-q}}
\newcommand{\AAQuantumVirial}{\ThermalFigure{
 \draw[->](0,0)--(5,0)node[right]{$x=n\lambda_T^3/g$};\draw[->](0,0)--(0,3.5)node[above]{$p/(nkT)$};\node[left]at(0,1.8){$1$};
 \draw[dashed](0,1.8)--(4.5,1.8);\draw[DeepBlue,thick](0,1.8)--(4.5,2.8)node[right]{费米};\draw[thick](0,1.8)--(4.5,.8)node[right]{玻色};
 \node[above]at(3,3.6){只示意 $x\ll1$ 的一阶斜率，纵向修正已放大};
}{交换统计使弱简并压强分别高于或低于经典值，但不是粒子间真实势力。}{advA2:quality:virial-a}
图~\ref{advA2:quality:virial-a}中的两直线只在原点附近表示第二维里修正。将逸度展开反解成密度展开后再比较 $p$，才能确定修正系数；若直接把 $z$ 当作 $n\lambda_T^3/g$，会漏掉同阶的反解项，并把正确的 $2^{-5/2}$ 系数算错。
}

\section{统计力学}

\subsection{相空间、系综与配分函数}

考虑不显含时间的拉格朗日系统 $L(q_\alpha,\dot q_\alpha)$。定义
\[
 p_\alpha=\frac{\partial L}{\partial\dot q_\alpha},
 \qquad
 H(q,p)=\sum_\alpha p_\alpha\dot q_\alpha-L,
\]
哈密顿方程为
\[
 \dot q_\alpha=\frac{\partial H}{\partial p_\alpha},
 \qquad
 \dot p_\alpha=-\frac{\partial H}{\partial q_\alpha}.
\]
把 $f$ 个广义坐标和 $f$ 个广义动量合成相坐标
$\vect\xi=(q_1,\ldots,q_f,p_1,\ldots,p_f)$，则
\[
 \dot{\vect\xi}=\mathsf J\nabla_\xi H,
 \qquad
 \mathsf J=\begin{pmatrix}0&\mathsf1\\-\mathsf1&0\end{pmatrix}.
\]
与量子态计数协调的无量纲相体积元取为
$\dd\Omega=\dd^f q\,\dd^f p/h^f$。系综是在相空间内沿哈密顿轨道运动的大量代表点，其归一化密度记为 $\rho(\vect\xi,t)$。

\begin{advtheorem}{刘维尔定理}\label{advA2:thm:liouville}
哈密顿相流不可压缩，沿任一相轨道
\[
 \boxed{\frac{D\rho}{Dt}=0}.
\]
等价地，相空间任意随流体积保持不变。
\end{advtheorem}
证明：相空间概率守恒给出连续性方程
\[
 \frac{\partial\rho}{\partial t}
 +\nabla_\xi\cdot(\rho\dot{\vect\xi})=0.
\]
哈密顿方程使相流散度为
\[
 \nabla_\xi\cdot\dot{\vect\xi}
 =\sum_\alpha\left(
 \frac{\partial^2H}{\partial q_\alpha\partial p_\alpha}
 -\frac{\partial^2H}{\partial p_\alpha\partial q_\alpha}
 \right)=0.
\]
因此连续性方程化为
$\partial_t\rho+\dot{\vect\xi}\cdot\nabla_\xi\rho=D\rho/Dt=0$。这也补足了源资料中省略的证明。

玻尔兹曼在分子混沌等动力学假设下得到粗粒化的 $H$ 定理，
$H_{\rm B}=\int\rho\ln\rho\,\dd\Omega$ 随时间不增。要注意：刘维尔定理保证孤立哈密顿体系的\emph{细粒化}吉布斯熵
$-k\int\rho\ln\rho\,\dd\Omega$ 不变；不可逆熵增来自粗粒化或约化描述，二者并不矛盾。

微正则系综固定 $N,E,V$。令
\[
 \omega(N,E,V)=\int_{E\le H<E+\Delta E}\dd\Omega,
 \qquad
 \rho=\begin{cases}1/\omega,&E\le H<E+\Delta E,\\0,&\text{其余},\end{cases}
\]
并记 $\Omega_E=\int_{H\le E}\dd\Omega$，则窄能壳中
$\omega\simeq(\partial\Omega_E/\partial E)_{N,V}\Delta E$。系综平均为
$\langle A\rangle=\int A\rho\,\dd\Omega$，熵为
$S=-k\int\rho\ln\rho\,\dd\Omega=k\ln\omega$。

\begin{advexamplebox}{等概率分布使微正则熵最大}\label{advA2:ex:microcanonical-maximum}
证明：在给定能壳和归一化条件下，均匀分布 $\rho=1/\omega$ 使
$S=-k\int\rho\ln\rho\,\dd\Omega$ 最大。

\AQMicrocanonical
\end{advexamplebox}
\begin{advsolution}
\AAMicrocanonical

对任意归一化分布 $\rho$，与均匀分布 $\rho_0=1/\omega$ 的相对熵满足
\[
 D(\rho\Vert\rho_0)=\int\rho\ln\frac{\rho}{\rho_0}\dd\Omega\ge0.
\]
于是
\[
 -\int\rho\ln\rho\,\dd\Omega
 \le-\int\rho\ln\rho_0\,\dd\Omega=\ln\omega.
\]
等号只在 $\rho=\rho_0$ 几乎处处成立时取得。也可用拉格朗日乘子变分，二阶变分
$-k\int(\delta\rho)^2/(2\rho)\dd\Omega<0$ 保证是最大值。
\end{advsolution}

两个弱耦合体系交换能量时，总态数
$\omega_1(E_1)\omega_2(E-E_1)$ 在平衡处最大，给出
\[
 \frac1{T_i}=\left(\frac{\partial S_i}{\partial E_i}\right)_{V_i,N_i},
 \qquad T_1=T_2.
\]
同理，允许体积和粒子交换时分别得到
$(\partial S/\partial V)_{E,N}=p/T$、
$(\partial S/\partial N)_{E,V}=-\mu/T$，以及
$p_1=p_2,\mu_1=\mu_2$，从微观态计数恢复
$\dd E=T\dd S-p\dd V+\mu\dd N$。

小体系与大热库接触时，把热库熵在总能量附近展开：
\[
 \ln\omega_B(E_T-E)=\text{常数}-\frac{E}{kT}+\cdots,
\]
故正则系综权重为 $e^{-E/(kT)}$。若还与粒子库接触，展开中出现化学势项，巨正则权重为
$e^{-(E-\mu N)/(kT)}$。

\begin{advdefinition}{配分函数与巨配分函数}\label{advA2:def:partition-functions}
\[
 Z(N,V,T)=\int e^{-\beta H_N}\dd\Omega_N,
 \qquad
 \Xi(\mu,V,T)=\sum_{N=0}^\infty e^{\beta\mu N}Z(N,V,T),
 \quad \beta=\frac1{kT}.
\]
热力学势为
\[
 F=-kT\ln Z,
 \qquad
 \Omega=-kT\ln\Xi=-pV.
\]
\end{advdefinition}
由此
\[
 U=-\frac{\partial\ln Z}{\partial\beta},
 \quad
 \mu=F_N,\quad p=-F_V,\quad S=-F_T.
\]

\begin{advexamplebox}{由巨配分函数直接求热力学量}\label{advA2:ex:grand-observables}
用 $\Xi(\mu,V,T)$ 写出平均粒子数、压强、熵和内能。

\AQGrand
\end{advexamplebox}
\begin{advsolution}
\AAGrand

对 $z=e^{\beta\mu}$ 求导，
\[
 \boxed{\langle N\rangle
 =z\frac{\partial\ln\Xi}{\partial z}
 =kT\left(\frac{\partial\ln\Xi}{\partial\mu}\right)_{T,V}}.
\]
由 $\Omega=-kT\ln\Xi$ 及 $p=-(\partial\Omega/\partial V)_{T,\mu}$，
\[
 \boxed{p=kT\left(\frac{\partial\ln\Xi}{\partial V}\right)_{T,\mu}},
 \qquad
 \boxed{S=k\ln\Xi+kT
 \left(\frac{\partial\ln\Xi}{\partial T}\right)_{\mu,V}}.
\]
固定 $z,V$ 对 $\beta$ 求导得到
\[
 \boxed{U=-\left(\frac{\partial\ln\Xi}{\partial\beta}\right)_{z,V}},
\]
等价形式是在固定 $\mu$ 时
$U-\mu N=-(\partial\ln\Xi/\partial\beta)_{\mu,V}$。最后
$U=TS-pV+\mu N$ 可作一致性校验。
\end{advsolution}

\subsection{经典统计结论}

无内禀简并的自由粒子单粒子配分函数写成
\[
 z_1=\int_0^\infty g(\varepsilon)e^{-\beta\varepsilon}\dd\varepsilon.
\]

\begin{advexamplebox}{不同维数自由粒子的态密度}\label{advA2:ex:density-of-states}
求一、二、三维非相对论自由粒子和无质量相对论粒子的态密度 $g(\varepsilon)$。

\AQDensityStates
\end{advexamplebox}
\begin{advsolution}
\AADensityStates

每个平动量子态占据相空间体积 $(2\pi\hbar)^d$。先数动量球内态数，再对能量求导。
对 $\varepsilon=p^2/(2m)$，
\[
 \boxed{g_{3}(\varepsilon)=\frac{V(2m)^{3/2}}{4\pi^2\hbar^3}\varepsilon^{1/2}},
\]
\[
 \boxed{g_{2}(\varepsilon)=\frac{Am}{2\pi\hbar^2}},
 \qquad
 \boxed{g_{1}(\varepsilon)=\frac{L}{\pi\hbar}\sqrt{\frac{m}{2\varepsilon}}}.
\]
对 $\varepsilon=pc$，
\[
 \boxed{g_{3}^{\rm rel}=\frac{V\varepsilon^2}{2\pi^2\hbar^3c^3}},
 \quad
 \boxed{g_{2}^{\rm rel}=\frac{A\varepsilon}{2\pi\hbar^2c^2}},
 \quad
 \boxed{g_{1}^{\rm rel}=\frac{L}{\pi\hbar c}}.
\]
每一式的 $g\dd\varepsilon$ 都无量纲；若存在自旋简并度，还应整体乘以简并因子。
\end{advsolution}

\begin{advexamplebox}{自由粒子配分函数与物态方程}\label{advA2:ex:free-particle-partition}
计算三维非相对论和无质量相对论自由粒子的单粒子配分函数，并证明平均能量分别为
$3kT/2$、$3kT$，且分别满足 $pV=2E/3$、$pV=E/3$。

\AQFreePartition
\end{advexamplebox}
\begin{advsolution}
\AAFreePartition

令热德布罗意波长
\[
 \lambda_T=\sqrt{\frac{2\pi\hbar^2}{mkT}}.
\]
高斯积分给出非相对论单粒子
$z_1=V/\lambda_T^3\propto VT^{3/2}$；无质量粒子由
$\int_0^\infty\varepsilon^2e^{-\beta\varepsilon}\dd\varepsilon=2/\beta^3$ 得
\[
 z_1^{\rm rel}=\frac{V(kT)^3}{\pi^2\hbar^3c^3}\propto VT^3.
\]
于是 $E_1=-\partial_\beta\ln z_1$ 分别为 $3kT/2$ 和 $3kT$。两者都因 $\ln z_1$ 含 $\ln V$ 而给出
$p=kT(\partial\ln z_1/\partial V)_T=kT/V$，所以
\[
 \boxed{pV=\frac23E}\quad\text{和}\quad
 \boxed{pV=\frac13E}.
\]
这里 $E$ 可指单粒子平均能量；乘以粒子数后关系不变。
\end{advsolution}

经典全同粒子必须除去置换重复计数：$Z_N=z_1^N/N!$。斯特林公式
$\ln N!\simeq N\ln N-N$ 使自由能和熵恢复广延性。

\begin{advexamplebox}{单原子经典理想气体熵}\label{advA2:ex:classical-entropy}
由 $Z_N=(V/\lambda_T^3)^N/N!$ 证明
\[
 S=\frac32Nk\ln T+Nk\ln\frac VN+Nks_0,
 \quad
 s_0=\frac32\ln\frac{mk}{2\pi\hbar^2}+\frac52.
\]

\AQClassicalEntropy
\end{advexamplebox}
\begin{advsolution}
\AAClassicalEntropy

用 $F=-kT\ln Z_N$、$S=-(\partial F/\partial T)_{V,N}$：
\[
 S=Nk\left[\ln\frac{V}{N\lambda_T^3}+\frac52\right].
\]
展开 $\lambda_T$ 便得到题式。它与前面的萨克--泰特洛德结果完全一致；$1/N!$ 若遗漏，会多出非广延的 $Nk\ln N$，这就是吉布斯佯谬的经典表现。
\end{advsolution}

\begin{advexamplebox}{$N$ 个一维经典谐振子}\label{advA2:ex:harmonic-partition}
计算 $N$ 个频率为 $\omega$、彼此独立且固定在不同位置的一维经典谐振子的配分函数、内能与热容。

\AQOscillators
\end{advexamplebox}
\begin{advsolution}
\AAOscillators

单振子的哈密顿量 $H=p^2/(2m)+m\omega^2q^2/2$，故
\[
 z=\frac1h\int_{-\infty}^{\infty}\!\dd p\dd q\,e^{-\beta H}
 =\frac{kT}{\hbar\omega}.
\]
振子依附在不同位置，彼此可分辨，不除以 $N!$：
\[
 \boxed{Z=z^N=\left(\frac{kT}{\hbar\omega}\right)^N},
 \quad
 \boxed{U=NkT},\quad \boxed{C_V=Nk}.
\]
两个二次项各贡献 $kT/2$，与能量均分定理一致。低温时经典结果失效，量子配分函数应取
$e^{-\beta\hbar\omega/2}/(1-e^{-\beta\hbar\omega})$。
\end{advsolution}

\begin{advexamplebox}{$N$ 个三自由度刚体转子}\label{advA2:ex:rigid-rotor-partition}
主惯量为 $I_1,I_2,I_3$ 的非线性刚体分子只有转动自由度。求其高温经典配分函数，并给出转动内能与热容。

\AQRotors
\end{advexamplebox}
\begin{advsolution}
\AARotors

对欧拉角和三个共轭角动量作高斯积分，单分子转动配分函数为
\[
 z_{\rm rot}=\frac{\sqrt\pi}{\sigma}
 \left(\frac{8\pi^2kT}{h^2}\right)^{3/2}
 \sqrt{I_1I_2I_3},
\]
其中 $\sigma$ 是分子取向的对称数。若这些转子位于固定晶格点，$Z=z_{\rm rot}^N$；若作为全同自由分子的内部因子，则它进入
$Z=(z_{\rm trans}z_{\rm rot})^N/N!$。因 $z_{\rm rot}\propto T^{3/2}$，
\[
 \boxed{U_{\rm rot}=\frac32NkT},
 \qquad \boxed{C_{V,\rm rot}=\frac32Nk}.
\]
三项角动能各贡献 $kT/2$。
\end{advsolution}

\begin{advexamplebox}{匀强电场中的刚性偶极子}\label{advA2:ex:dipole-partition}
$N$ 个大小为 $d$ 的刚性电偶极子置于匀强电场 $E$ 中，忽略动能。求配分函数、平均偶极矩与弱场极化率。

\AQDipoleStatistics
\end{advexamplebox}
\begin{advsolution}
\AADipoleStatistics

取 $x=dE/(kT)$。单偶极子取向积分为
\[
 z_1=\int e^{x\cos\theta}\dd\Omega
 =4\pi\frac{\sinh x}{x},
 \qquad Z=z_1^N.
\]
沿场方向的平均偶极矩为
\[
 \langle d_\parallel\rangle
 =kT\frac{\partial\ln z_1}{\partial E}
 =d\mathcal L(x),
 \quad \mathcal L(x)=\coth x-\frac1x.
\]
故极化强度
$P=Nd\mathcal L(x)/V$。弱场 $\mathcal L(x)=x/3+O(x^3)$，
\[
 \boxed{\chi_e=\frac1{\varepsilon_0}\frac{\partial P}{\partial E}
 =\frac{Nd^2}{3\varepsilon_0VkT}}.
\]
强场 $x\to\infty$ 时 $P\to Nd/V$，正确再现饱和极化。
\end{advsolution}

\subsection{费米统计与玻色统计}

把单粒子态记为 $\varepsilon_i$、占据数记为 $n_i$，则
$E=\sum_i n_i\varepsilon_i,N=\sum_i n_i$。每个态的巨配分函数为
\[
 \Xi_i=\sum_{n_i}e^{-\beta(\varepsilon_i-\mu)n_i}.
\]
费米子有 $n_i=0,1$；玻色子有 $n_i=0,1,2,\ldots$。

\begin{advexamplebox}{费米--狄拉克与玻色--爱因斯坦分布}\label{advA2:ex:quantum-distributions}
求两种统计的平均占据数和巨势，并写出压强公式。

\AQQuantumDistributions
\end{advexamplebox}
\begin{advsolution}
\AAQuantumDistributions

令 $z_i=e^{-\beta(\varepsilon_i-\mu)}$。费米子
$\Xi_i=1+z_i$，玻色子 $\Xi_i=(1-z_i)^{-1}$。用
$\langle n_i\rangle=z_i\partial_{z_i}\ln\Xi_i$，
\[
 \boxed{\langle n_i\rangle=\frac1{e^{(\varepsilon_i-\mu)/(kT)}\pm1}},
\]
上号对应费米、下号对应玻色。又
\[
 \boxed{\Omega=\mp kT\sum_i\ln\!\left(1\pm e^{-\beta(\varepsilon_i-\mu)}\right)},
\]
因此
\[
 \boxed{pV=\pm kT\sum_i\ln\!\left(1\pm e^{-\beta(\varepsilon_i-\mu)}\right)}.
\]
有限系统的玻色巨配分级数严格要求 $\mu$ 低于基态能量（凝聚热力学极限另行处理）；费米式自动体现泡利不相容原理。
\end{advsolution}

三维理想费米气在零温下把单粒子态填充到费米能
\[
 \varepsilon_F=\frac{\hbar^2}{2m}
 \left(\frac{6\pi^2N}{gV}\right)^{2/3},
 \qquad p_F=\frac{2N\varepsilon_F}{5V},
\]
其中 $g$ 是内禀简并度。

\begin{advexamplebox}{金属电子的低温热容}\label{advA2:ex:fermi-heat-capacity}
证明 $T\ll T_F$ 时理想费米气的电子热容为
\[
 \boxed{C_V=\frac{\pi^2}{2}Nk\frac{kT}{\varepsilon_F}}.
\]

\AQFermiCapacity
\end{advexamplebox}
\begin{advsolution}
% CompletenessThermalFields20261024: fixed-N Sommerfeld derivation.
先说明低温展开怎样产生，再实施定粒子数约束。定义 $T_F=\varepsilon_F/k$，以下固定 $N,V$，
并令 $g(\varepsilon)$ 表示包含全部自旋简并的\emph{总态密度}，单位为能量的倒数。
费米分布为 $f_T(\varepsilon)=[\exp((\varepsilon-\mu)/(kT))+1]^{-1}$。
图~\ref{advA2:complete:sommerfeld}把它与同一化学势处的零温阶跃相减；这一步把整个费米海的积分
变成只在 $|\varepsilon-\mu|\sim kT$ 内有贡献的修正。

\begin{qualityfigure}
\centering
\begin{tikzpicture}[>=Latex,font=\small,x=1.0cm,y=2.8cm]
 \draw[->](-4.6,0)--(4.6,0)node[right]{$x=(\varepsilon-\mu)/(kT)$};
 \draw[->](0,-.62)--(0,.69)node[above]{$D(x)$};
 \fill[DeepBlue!13,domain=0:4,samples=70] (0,0)--plot(\x,{1/(exp(\x)+1)})--(4,0)--cycle;
 \fill[DeepBlue!13,domain=-4:0,samples=70] (-4,0)--plot(\x,{-1/(exp(-\x)+1)})--(0,0)--cycle;
 \draw[DeepBlue,thick,domain=0:4,samples=70]plot(\x,{1/(exp(\x)+1)});
 \draw[DeepBlue,thick,domain=-4:0,samples=70]plot(\x,{-1/(exp(-\x)+1)});
 \draw[dashed](0,-.5)--(0,.5);
 \node[above,align=center]at(2,.25){增占据\\$x>0$};
 \node[below,align=center]at(-2,-.25){减占据\\$x<0$};
 \node[below]at(3.7,-.47){$D(-x)=-D(x)$};
 \node[above]at(-2.6,.48){$\int D(x)\,\dd x=0$};
\end{tikzpicture}
\figcap{费米分布相对零温阶跃的差 $D(x)=(e^x+1)^{-1}-\Theta(-x)$。
两瓣的有向面积相消，但乘上 $x$ 后两侧贡献同号；这决定首项为 $T^2$ 修正。\label{advA2:complete:sommerfeld}}
\end{qualityfigure}

对在 $\mu$ 附近光滑、无带边或奇点的函数 $\Phi$，考虑
\[
 I_\Phi=\int_0^\infty\Phi(\varepsilon)f_T(\varepsilon)\,\dd\varepsilon.
\]
设 $\mu\gg kT$。作 $x=(\varepsilon-\mu)/(kT)$ 代换，并在费米面附近展开 $\Phi$，得到
\[
 I_\Phi-\int_0^\mu\Phi(\varepsilon)\,\dd\varepsilon
 =kT\int_{-\mu/(kT)}^\infty\Phi(\mu+kTx)D(x)\,\dd x.
\]
由于 $D$ 的尾部指数衰减，各项的矩积分可把下限延到 $-\infty$，所添误差指数小。
图中 $D$ 为奇函数，因此常数项和二次项不贡献；一次项的系数为
\[
 \int_{-\infty}^{\infty}xD(x)\,\dd x
 =2\int_0^\infty\frac{x\,\dd x}{e^x+1}
 =2\sum_{r=1}^\infty\frac{(-1)^{r+1}}{r^2}
 =\frac{\pi^2}{6}.
\]
这里先展开 $(e^x+1)^{-1}=\sum_{r\ge1}(-1)^{r+1}e^{-rx}$，逐项积分用
$\int_0^\infty xe^{-rx}\dd x=r^{-2}$；交错平方和是普通平方和 $\pi^2/6$ 的一半。
于是索末菲展开的前两项为
\begin{equation}
 I_\Phi=\int_0^\mu\Phi(\varepsilon)\,\dd\varepsilon
 +\frac{\pi^2}{6}(kT)^2\Phi'(\mu)+O(T^4).
 \label{advA2:eq:complete-sommerfeld}
\end{equation}
余项表示在费米面光滑、导数有界的低温渐近展开；不能用在 $kT$ 已与 $\varepsilon_F$ 同量级的区域。

\textbf{先固定粒子数。} 取 $\Phi=g$，写 $\mu=\varepsilon_F+\delta\mu$。由于
$N=\int_0^{\varepsilon_F}g\,\dd\varepsilon$，式~\eqref{advA2:eq:complete-sommerfeld}给出
\[
 0=g_F\delta\mu+\frac{\pi^2}{6}(kT)^2g'_F+O(T^4),
 \qquad
 \delta\mu=-\frac{\pi^2}{6}(kT)^2\frac{g'_F}{g_F}+O(T^4),
\]
下标 $F$ 表示在 $\varepsilon_F$ 处取值。三维自由粒子有 $g=A\sqrt\varepsilon$，因此
\[
 g_F=\frac{3N}{2\varepsilon_F},\qquad
 \mu=\varepsilon_F-\frac{\pi^2}{12}\frac{(kT)^2}{\varepsilon_F}+O(T^4).
\]
化学势的微小下移正是补偿图中态密度不恒定时两侧粒子数不完全相等的修正。

\textbf{再计算能量。} 取 $\Phi=\varepsilon g(\varepsilon)$，保留同阶的两项：
\begin{align*}
 U(T)-U(0)
 &=\varepsilon_Fg_F\delta\mu
   +\frac{\pi^2}{6}(kT)^2(g_F+\varepsilon_Fg'_F)+O(T^4)\\
 &=-\frac{\pi^2}{6}(kT)^2\varepsilon_Fg'_F
   +\frac{\pi^2}{6}(kT)^2(g_F+\varepsilon_Fg'_F)+O(T^4)\\
 &=\frac{\pi^2}{6}g_F(kT)^2+O(T^4)
  =\frac{\pi^2}{4}N\frac{(kT)^2}{\varepsilon_F}+O(T^4).
\end{align*}
第一项来自积分上限的移动；若误把 $\mu$ 固定，就会漏掉这一项并改变系数。
最后在固定 $N,V$ 下求导，得到
\[
 \boxed{C_V=\frac{\pi^2}{2}Nk\frac{kT}{\varepsilon_F}+O(T^3)}.
\]
$C_V$ 的单位为能量除以温度，$T\to0$ 时线性趋零；只有费米面附近的窄能窗能被热激发，
所以它比经典值 $3Nk/2$ 小一个 $T/T_F$ 因子。

\end{advsolution}

理想玻色气的激发态粒子数为
\[
 N_e=\int_0^\infty\frac{g(\varepsilon)\dd\varepsilon}
 {e^{(\varepsilon-\mu)/(kT)}-1}.
\]
三维均匀气体在临界温度以下有 $\mu=0$、
\[
 N_0=N\left[1-\left(\frac T{T_c}\right)^{3/2}\right].
\]

\begin{advexamplebox}{玻色--爱因斯坦凝聚的热容曲线}\label{advA2:ex:bec-heat-capacity}
定性并定量说明三维理想玻色气定容热容随温度的变化，并画出临界点附近曲线。

\AQBoseCapacity
\end{advexamplebox}
\begin{advsolution}
\AABoseCapacity

$T<T_c$ 时
\[
 U=\frac32kT\frac{V}{\lambda_T^3}\zeta(5/2),
\]
故
\[
 \boxed{C_V=\frac{15}{4}Nk
 \frac{\zeta(5/2)}{\zeta(3/2)}
 \left(\frac T{T_c}\right)^{3/2}}.
\]
它从零按 $T^{3/2}$ 增加，在 $T_c^-$ 达到约 $1.93Nk$。$T>T_c$ 时化学势转为负值，热容从同一数值连续接出，随后趋向经典极限 $3Nk/2$；一阶导数在 $T_c$ 处不连续，形成尖点而非潜热跃变。


图~\ref{advA2:quality:bose-c-a}的两支在临界点相接；基态在本模型中不贡献激发能，低温热容来自激发态占据随温度改变。
\end{advsolution}

\begin{advexamplebox}{稀薄量子气体的第二维里修正}\label{advA2:ex:quantum-virial}
对三维非相对论理想费米气和玻色气：
\begin{enumerate}
 \item 写出经典极限的化学势、内能和压强。
 \item 给出弱简并条件，并求到 $n\lambda_T^3/g$ 一阶的 $\mu,U,p$。
 \item 把物态方程与无排斥体积的范德瓦耳斯形式
 $(p+aN^2/V^2)V=NkT$ 比较，说明这种等价的物理边界。
\end{enumerate}

\AQQuantumVirial
\end{advexamplebox}
\begin{advsolution}
\AAQuantumVirial

令 $\eta=+1$ 表示费米统计、$\eta=-1$ 表示玻色统计，逸度
$z=e^{\mu/(kT)}$，并记 $x=n\lambda_T^3/g$。经典麦克斯韦--玻尔兹曼极限为
\[
 \boxed{\mu_{\rm MB}=kT\ln x},\qquad
 \boxed{U_{\rm MB}=\frac32NkT},\qquad
 \boxed{p_{\rm MB}=nkT}.
\]
弱简并条件是
\[
 \boxed{x=\frac{n\lambda_T^3}{g}\ll1},
\]
即平均每个热波包和每个内禀态中的粒子数远小于1。展开分布函数，
\[
 x=z-\eta\frac{z^2}{2^{3/2}}+O(z^3),
 \qquad
 \frac{p\lambda_T^3}{gkT}
 =z-\eta\frac{z^2}{2^{5/2}}+O(z^3).
\]
反解并保留一阶修正：
\[
 \begin{aligned}
 \mu&=kT\left[\ln x+\eta\frac{x}{2^{3/2}}\right],\\
 p&=nkT\left[1+\eta\frac{x}{2^{5/2}}\right],\qquad U=\frac32pV.
 \end{aligned}
\]
与 $p=nkT-an^2$ 比较，得到 $a(T)=-\eta kT\lambda_T^3/(2^{5/2}g)$。
费米统计给出正压修正，可看成“统计排斥”；玻色统计给出负压修正，可看成“统计聚集”。它只是第二维里系数的数学等价，并不意味着无相互作用量子粒子之间真的出现了范德瓦耳斯力；密度升高后高阶项不可忽略，等价立即失效。
\end{advsolution}


\clearpage\endgroup

\begingroup
\newif\ifAdvHideHeading
\RenewDocumentEnvironment{advexamplebox}{m}{%
  \refstepcounter{advexample}%
  \begin{tcolorbox}[
    enhanced,unbreakable,colback=white,colframe=RuleGray,boxrule=0.5pt,arc=0mm,
    title={例~\theadvexample\quad #1},fonttitle=\bfseries,
    coltitle=Ink,colbacktitle=PaleBlue,borderline west={1.4pt}{0pt}{DeepBlue}]
}{\end{tcolorbox}}

\RenewDocumentEnvironment{advsolution}{+b}{%
  \par\smallskip\begingroup\small\noindent
  \ifAdvHideHeading\global\AdvHideHeadingfalse\else\LectureSolutionHeading\fi #1\par\endgroup\smallskip
}{}

\RenewDocumentEnvironment{advdefinition}{m +b}{%
  \refstepcounter{advdefinition}%
  \begin{tcolorbox}[
    enhanced,unbreakable,colback=white,colframe=RuleGray,boxrule=0.5pt,arc=0mm,
    title={定义~\theadvdefinition\quad #1},fonttitle=\bfseries,
    coltitle=Ink,colbacktitle=PaleGold,borderline west={1.2pt}{0pt}{Gold}]
  #2
  \end{tcolorbox}
}{}

\RenewDocumentEnvironment{advtheorem}{m +b}{%
  \refstepcounter{advtheorem}%
  \begin{tcolorbox}[
    enhanced,unbreakable,colback=white,colframe=RuleGray,boxrule=0.5pt,arc=0mm,
    title={定理~\theadvtheorem\quad #1},fonttitle=\bfseries,
    coltitle=Ink,colbacktitle=PaleBlue,borderline west={1.2pt}{0pt}{DeepBlue}]
  #2
  \end{tcolorbox}
}{}


\AdvancedUnit{3}{数学工具}
% AdvancedThermalQuality20261023-v1
% Editable mathematical diagrams: index structure, orthogonal directions and boundaries.
\providecommand{\ThermalFigure}[3]{\begin{qualityfigure}\begin{tikzpicture}[>=Latex,font=\small,x=.8cm,y=.8cm]#1\end{tikzpicture}\figcap{#2}\label{#3}\end{qualityfigure}}
\newcommand{\AQEpsilon}{\ThermalFigure{
 \draw[thick](0,0)grid[step=.7](2.1,2.1);\foreach\x/\y in {.35/1.75,1.05/1.05,1.75/.35}{\node at(\x,\y){$1$};}
 \node[left]at(0,1.75){$[1]_r$};\node[left]at(0,1.05){$[2]_r$};\node[left]at(0,.35){$[3]_r$};
 \node[right,align=left]at(3,1.2){从单位矩阵依次选行 $i,j,k$\\未标格的元素为零};
}{每行是一个标准基矢的分量，指标重复与行重复相对应。}{advA3:quality:epsilon-q}}
\newcommand{\AAEpsilon}{\ThermalFigure{
 \node[draw,align=center](A)at(1,1){行序 $1,2,3$\\正定向};\node[draw,align=center](B)at(6,1){行序 $2,1,3$\\交换两行};
 \draw[<->,thick](A)--node[above]{一次换位}(B);\node[below]at(3.5,-.5){行列式 $+1\longleftrightarrow-1$；两行重复时为零};
}{定向改变对应全反对称符号；缩并后的 Kronecker 符号表达相同的指标配对。}{advA3:quality:epsilon-a}
图~\ref{advA3:quality:epsilon-a}中一次换位翻转定向，循环置换包含两次换位而保号。缩并 $\varepsilon_{ijk}\varepsilon_{imn}$ 时，若 $j=k$ 两边均为零；若 $j\ne k$，第一个因子只允许互补的一个 $i$，第二个因子再检验 $(m,n)$ 是同序还是反序，从而得到两个 $\delta$ 乘积之差。
}
\newcommand{\AQJacobi}{\ThermalFigure{
 \foreach\x/\y/\lab in {3/0/u,1/2.5/v,-1.5/1.2/w}{\draw[->,thick](0,0)--(\x,\y)node[above]{$\vect\lab$};}\fill(0,0)circle(2pt)node[below]{$O$};
 \node[right]at(4,1.5){同起点矢量的投影示意，不预设三矢量共面};
}{三个矢量在双重叉积中循环交换，括号的次序须保留。}{advA3:quality:jacobi-q}}
\newcommand{\AAJacobi}{\ThermalFigure{
 \node[draw,align=left]at(4,2){沿 $\vect u$：$-(\vect v\cdot\vect w)\vect u+(\vect w\cdot\vect v)\vect u$};
 \node[draw,align=left]at(4,.7){沿 $\vect v$：$(\vect u\cdot\vect w)\vect v-(\vect w\cdot\vect u)\vect v$};
 \node[draw,align=left]at(4,-.6){沿 $\vect w$：$-(\vect u\cdot\vect v)\vect w+(\vect v\cdot\vect u)\vect w$};
}{展开后三个方向的系数分别成对抵消，不需要三个矢量线性无关。}{advA3:quality:jacobi-a}
图~\ref{advA3:quality:jacobi-a}把六项按同一输出矢量分组。每组的两个标量内积相等、符号相反，故逐组为零；证明不依赖三矢量能否组成一组基。这也说明应使用双重叉积恒等式，不能将叉积误当作满足结合律的乘法。
}
\newcommand{\AQDerivations}{\ThermalFigure{
 \node[draw](W)at(0,1){$\vect w$};\node[draw](V)at(4,2){$D_{\vect v}\vect w$};\node[draw](U)at(4,0){$D_{\vect u}\vect w$};
 \draw[->](W)--node[above]{$D_{\vect v}$}(V);\draw[->](W)--node[below]{$D_{\vect u}$}(U);
 \node[right]at(7,1){继续复合并比较两种次序};
}{算子复合中右侧算子先作用；两条路径的差是对易子。}{advA3:quality:derivation-q}}
\newcommand{\AADerivations}{\ThermalFigure{
 \node[draw,align=center](A)at(1,1){$D_{\vect u}D_{\vect v}\vect w$\\$-D_{\vect v}D_{\vect u}\vect w$};
 \node[draw](B)at(8,1){$(\vect u\times\vect v)\times\vect w$};\draw[->,thick](A)--node[above]{雅可比重排}(B);
 \node[below]at(4.5,-.4){同一等式对每个输入 $\vect w$ 成立，才得到算子恒等式};
}{把对易子先落实到任意测试矢量，避免将形式符号当作交换乘法。}{advA3:quality:derivation-a}
图~\ref{advA3:quality:derivation-a}中的右端是一个新的叉积算子对 $\vect w$ 的作用。由于 $\vect w$ 任意，才可去掉测试输入得到算子等式。导子性质则将 $D_{\vect u}$ 分别作用于楔积的两个因子，输出仍为同一矢量空间中的矢量。
}
\newcommand{\AQSensors}{\ThermalFigure{
 \draw[thick](0,0)grid[step=.8](2.4,2.4);\fill[DeepBlue!12](0,0)--(2.4,2.4)--(2.4,0)--cycle;
 \node[right,align=left]at(3.2,1.6){二阶例：$T_{ij}=T_{ji}$\\一般 $n$ 阶：所有指标置换均对称};
}{完全对称性把指标顺序不同、出现次数相同的分量识别为同一个自由参数。}{advA3:quality:symmetric-q}}
\newcommand{\AASensors}{\ThermalFigure{
 \foreach\x in {0,.7,2.2,3.7,4.4}{\fill[DeepBlue](\x,0)circle(3pt);}\draw[very thick](1.4,-.5)--(1.4,.5) (2.9,-.5)--(2.9,.5);
 \node[below]at(.35,-.7){$r_1=2$};\node[below]at(2.2,-.7){$r_2=1$};\node[below]at(4.05,-.7){$r_3=2$};
 \node[above]at(2.2,1.1){五颗星、两块隔板：五阶三维的一种分量类型};
}{允许相邻隔板，表示某类指标可以出现零次；总位置数为 $n+m-1$。}{advA3:quality:symmetric-a}
图~\ref{advA3:quality:symmetric-a}中的点表示指标出现次数而不是不同粒子。完全对称性使每一组计数 $(r_1,\ldots,r_m)$ 只对应一个独立分量，因此从 $n+m-1$ 个位置选 $m-1$ 个放隔板即可。若张量只对部分指标对称，这一计数不再适用。
}
\newcommand{\AQExterior}{\ThermalFigure{
 \foreach\x/\y/\lab in {3/0/1,0/2.5/2,-1.8/-1.2/3}{\draw[->,thick](0,0)--(\x,\y)node[right]{$\vect e_{\lab}$};}
 \node[right]at(4,1.5){各阶基由不同基矢的楔积构成};
}{外代数同时含零阶标量、向量及更高阶有向面积与体积元。}{advA3:quality:exterior-q}}
\newcommand{\AAExterior}{\ThermalFigure{
 \foreach\x/\w/\lab in {0/1/{1},1/3/{e_1,e_2,e_3},4/3/{e_1\wedge e_2,e_1\wedge e_3,e_2\wedge e_3},7/1/{e_1\wedge e_2\wedge e_3}}{\draw[thick](\x,0)rectangle(\x+\w,1);}
 \node at(.5,.5){$1$};\node at(2.5,.5){$\Lambda^1$：3};\node at(5.5,.5){$\Lambda^2$：3};\node at(7.5,.5){1};
 \node[below]at(4,-.6){三维：$1+3+3+1$；每个基矢可选或不选};
}{每个递增指标子集唯一对应一个楔积基元，重复指标则为零。}{advA3:quality:exterior-a}
图~\ref{advA3:quality:exterior-a}按阶数分组计数，最左的一份必须包括空子集对应的标量单位。每个基矢独立有“选、不选”两种选择，给出 $2^n$ 个子集；固定选 $l$ 个时才得到 $\binom nl$，两种计数相互核验。
}
\newcommand{\AQHodge}{\ThermalFigure{
 \fill[DeepBlue!10](-2,-.5)--(1,-.5)--(2,.5)--(-1,.5)--cycle;\draw[thick](-2,-.5)--(1,-.5)--(2,.5)--(-1,.5)--cycle;
 \draw[->,thick](0,0)--(0,2.6)node[right]{$\vect v$};\draw[->,DeepBlue](0,0)--(2,2)node[right]{$\vect u$};\node[below]at(0,-1){$\star\vect v$ 所在有向平面};
}{三维欧氏空间中的矢量与有向面积对偶，采用固定正定向。}{advA3:quality:hodge-q}}
\newcommand{\AAHodge}{\ThermalFigure{
 \draw[thick](-2,0)--(3,0);\node[below]at(-1,-.2){$\star\vect v$ 的平面};\draw[->](0,0)--(0,2.5)node[left]{$\vect v$};
 \draw[->,DeepBlue,thick](0,0)--(2,2)node[right]{$\vect u$};\draw[dashed](2,2)--(0,2);\draw[->,thick](0,0)--(0,2)node[left]{$u_\parallel$};
}{楔积产生有向体积，只保留 $\vect u$ 沿平面法向的分量。}{advA3:quality:hodge-a}
图~\ref{advA3:quality:hodge-a}中切向分量与面积元线性相关，其楔积为零，法向分量则给出 $|u_\parallel||v|$ 的有向体积。最后的霍奇对偶将三阶体积元变成标量，结果正是内积。若改变空间定向，两次对偶的符号须一起改变，不能只更改一次。
}
\newcommand{\AQRadialHessian}{\ThermalFigure{
 \foreach\r in {.7,1.3,2}{\draw[TextGray](0,0)circle(\r);}\fill(0,0)circle(2pt)node[below]{$O$};\fill[DeepBlue](1.414,1.414)circle(2pt)node[right]{$P$};\draw[->,thick](0,0)--(1.414,1.414)node[midway,above left]{$\vect r$};
 \node[right]at(3.5,1){$f(r)$ 的等值面为同心球面};
}{径向函数只通过 $r$ 依赖各坐标，导数仍须按坐标分量逐一求取。}{advA3:quality:radial-hessian-q}}
\newcommand{\AARadialHessian}{\ThermalFigure{
 \draw[thick](0,0)circle(2);\fill(2,0)circle(2pt)node[below right]{$P$};\draw[->,DeepBlue,thick](2,0)--(3.3,0)node[right]{径向};\draw[->,thick](2,0)--(2,1.4)node[right]{切向};\draw[dashed](0,0)--(2,0);
 \node[right,align=left]at(4.8,1.5){径向本征值 $n(n-1)r^{n-2}$\\两个切向本征值 $nr^{n-2}$};
}{Hessian 在径向与两切向分量上分别作用，其迹由三个本征值相加。}{advA3:quality:radial-hessian-a}
图~\ref{advA3:quality:radial-hessian-a}将局部增量分为沿 $\vect r$ 与垂直 $\vect r$ 的方向。由 $\partial_i r=x_i/r$ 再求一次导数可得 $\partial_i\partial_jr^n=nr^{n-2}\delta_{ij}+n(n-2)r^{n-4}x_ix_j$；切向增量与 $\vect r$ 点乘为零，直接读出两个本征值。所有这些点态表达式都排除原点的分布项。
}
\newcommand{\AQGreen}{\ThermalFigure{
 \draw[thick,fill=DeepBlue!7](0,0)ellipse(2.5 and 1.5);\node at(0,0){$V:\ \phi,\psi$};\fill(2.5,0)circle(2pt);\draw[->,thick](2.5,0)--(3.7,0)node[right]{$\vect n$};\node[above]at(0,2){分片光滑边界 $\partial V$};
}{边界法向取向外，曲面元为 $\dd\vect S=\vect n\dd S$。}{advA3:quality:green-q}}
\newcommand{\AAGreen}{\ThermalFigure{
 \draw[fill=DeepBlue!7,thick](-3,-1.3)rectangle(3,1.3);\draw[thick](0,-1.3)--(0,1.3);\node at(-1.5,.6){$V_1$};\node at(1.5,.6){$V_2$};
 \draw[->,DeepBlue,thick](0,0)--(1,0)node[below]{$\vect n_1$};\draw[->,thick](0,0)--(-1,0)node[below]{$\vect n_2$};
 \draw[->](3,.7)--(4.2,.7)node[right]{外边界法向};
}{相邻子体积的内部通量相消，只有整个区域的外边界保留下来。}{advA3:quality:green-a}
图~\ref{advA3:quality:green-a}中内部面的两法向相反，同一光滑通量 $\phi\nabla\psi$ 的贡献因而抵消。第一恒等式来自这一通量的散度；交换 $\phi,\psi$ 后相减，两个相同的梯度内积也抵消，从而得到第二恒等式。场若含奇点，应先挖去奇点小邻域，并另外保留其内边界通量。
}

\section{矢量分析}

\subsection{指标语言与三维矢量}

指标记号把成组分量、求和和缩并统一到一套语法中。例如
\[
 a_i=(1,2,3),\qquad b_i=(4,5,6),\qquad
 c_{ij}=a_ib_j,\qquad a_ib_i=32.
\]
只出现一次的 $i,j,k$ 是\emph{自由指标}，等式两边必须逐一对应；同一项中恰好出现两次的指标是\emph{哑指标}，按爱因斯坦求和约定对其全部取值求和。不同自由指标独立取值，三维中含两个自由指标的量有 $3^2$ 个分量。哑指标可改名，却不能与同项已有自由指标冲突。

从四个最小表达式可以引入这套语言：
\[
 a_i=(1,2,3),\qquad b_i=(4,5,6),
\]
\[
 c_{ij}=a_ib_j=
 \begin{pmatrix}4&5&6\\8&10&12\\12&15&18\end{pmatrix},
 \qquad d=a_ib_i=1\cdot4+2\cdot5+3\cdot6=32.
\]
前两式各含一个自由指标，实际各代表三个分量；第三式的 $i,j$ 独立变化，代表九个分量；第四式的 $i$ 在同一单项式中出现两次，因而是哑指标，只代表一次求和。由此得到三条必须始终遵守的语法规则：

\begin{enumerate}
 \item 同一单项式中，一个指标只能出现一次（自由）或两次（缩并），不能出现三次，例如 $a_ib_ic_i$ 没有定义；
 \item 等式两边的自由指标必须相同，例如 $A_{ij}=B_{ik}$ 不成立，因为右侧还含未缩并的 $k$；
 \item 不同单项式中的哑指标可以独立改名，例如
 $a_ib_i+c_jd_j=a_kb_k+c_ld_l$，但不能把它改成与同项自由指标冲突的名字。
\end{enumerate}

这些规则并非排版习惯，而是保证一个短式能够无歧义地还原成分量方程。比如
$A_{ij}x_j=b_i$ 代表三条方程，每条方程内部对 $j$ 求和；而
$A_{ij}x_i=b_j$ 则是另一个转置后的线性映射。指标的上下位置在一般坐标系中还区分逆变与协变分量；本节限于欧氏正交归一基，可用度量 $\delta_{ij}$ 升降指标，所以数值上 $x_i=x^i$，但概念上仍保留位置区别。

\begin{advdefinition}{克罗内克符号}\label{advA3:def:kronecker}
\[
 \delta_{ij}=\begin{cases}1,&i=j,\\0,&i\ne j.
 \end{cases}
\]
它是欧氏正交归一基中的单位张量分量。
\end{advdefinition}

\begin{advdefinition}{列维--奇维塔符号}\label{advA3:def:levi-civita}
三维中
\[
 \varepsilon_{ijk}=\begin{cases}
 +1,&(ijk)\text{ 是 }(123)\text{ 的偶置换},\\
 -1,&(ijk)\text{ 是 }(123)\text{ 的奇置换},\\
 0,&i,j,k\text{ 中有重复}.
 \end{cases}
\]
\end{advdefinition}
$\delta_{ij}$ 缩并时起换指标作用，例如
$a_ib_j\delta_{ij}=a_ib_i$、
$\delta_{ij}\delta_{kl}A_{ikm}=A_{jlm}$。

把换指标过程逐项写开可看得更清楚。对固定自由指标 $i$，
\[
 \delta_{ij}v_j=\delta_{i1}v_1+\delta_{i2}v_2+\delta_{i3}v_3=v_i.
\]
因此 $\delta_{ij}$ 是单位映射的分量。它的迹为
$\delta_{ii}=3$；而 $\delta_{ij}\delta_{jk}=\delta_{ik}$ 表明单位映射复合后仍是自身。列维--奇维塔符号则记录三维空间的取向：只有三个指标互异时才非零，循环置换 $(123)\to(231)\to(312)$ 保持正号，任意一次对换改变符号。

\begin{advtheorem}{列维--奇维塔缩并恒等式}\label{advA3:thm:epsilon-identities}
\[
 \varepsilon_{ijk}\varepsilon_{ijk}=6,
 \qquad
 \varepsilon_{ijk}\varepsilon_{ijn}=2\delta_{kn},
\]
\[
 \boxed{\varepsilon_{ijk}\varepsilon_{imn}
 =\delta_{jm}\delta_{kn}-\delta_{jn}\delta_{km}},
\]
\[
 \varepsilon_{ijk}\varepsilon_{lmn}
 =\det\!\begin{pmatrix}
 \delta_{il}&\delta_{im}&\delta_{in}\\
 \delta_{jl}&\delta_{jm}&\delta_{jn}\\
 \delta_{kl}&\delta_{km}&\delta_{kn}
 \end{pmatrix}.
\]
\end{advtheorem}
最后一式是两个有向体积形式的内积展开；令 $l=i$ 并缩并，依次得到前三式。

还可以不用预设结论来检查最常用的二指标缩并。固定 $j,k,m,n$，对 $i$ 求和。若 $j=k$，左边因第一个 $\varepsilon$ 为零，右边也成为
$\delta_{jm}\delta_{jn}-\delta_{jn}\delta_{jm}=0$。若 $j\ne k$，只有唯一一个与 $j,k$ 互补的 $i$ 可能贡献；第二个 $\varepsilon_{imn}$ 也只有当 $(m,n)$ 等于 $(j,k)$ 或 $(k,j)$ 时非零，正好分别对应右边的 $+1$ 与 $-1$。这给出独立的分量核验。

\begin{advexamplebox}{用行列式构造列维--奇维塔符号}\label{advA3:ex:epsilon-determinant}
令 $[i]_r$、$[i]_c$ 分别为三阶单位矩阵的第 $i$ 行和第 $i$ 列。
\begin{enumerate}
 \item 证明
 $\varepsilon_{ijk}=\det([i]_r,[j]_r,[k]_r)^{\mathsf T}$；
 \item 由此证明 $\varepsilon_{ijk}$ 全反对称；
 \item 证明定理\ref{advA3:thm:epsilon-identities}。
\end{enumerate}

\AQEpsilon
\end{advexamplebox}
\begin{advsolution}
\AAEpsilon

若 $i,j,k$ 有重复，行列式有两行相同而为零；若互异，这三行恰是单位阵三行的一个置换，行列式等于置换符号，正是 $\varepsilon_{ijk}$。交换任意两个指标等于交换两行，故符号反号。

对两个行列式用柯西--比内公式，
\[
 \varepsilon_{ijk}\varepsilon_{lmn}
 =\det([i]_r,[j]_r,[k]_r)^{\mathsf T}
  \det([l]_c,[m]_c,[n]_c),
\]
矩阵乘积的九个元素就是相应的 $\delta$，得到定理最后一式。再令
$l=i$ 求和；只有与固定 $(j,k)$ 互补的一个 $i$ 有贡献，展开行列式即得
$\delta_{jm}\delta_{kn}-\delta_{jn}\delta_{km}$。继续令 $m=j$、$n=k$ 缩并，分别得到 $2\delta_{kn}$ 和6。
\end{advsolution}

\begin{advdefinition}{三维欧氏矢量空间}\label{advA3:def:vector-space}
取线性无关基 $\vect e_1,\vect e_2,\vect e_3$，三维实矢量空间为
\[
 \mathbb R^3=\{x^i\vect e_i\mid x^i\in\mathbb R\}.
\]
元素 $\vect x=x^i\vect e_i$ 称为矢量；$x^i$ 只是该矢量在选定基下的分量。抽象指标写法 $x^a$ 强调“这是一个矢量”而不指定基。
\end{advdefinition}

这里要严格区分三个对象：$\vect e_i$ 是基矢，$x^i$ 是三个实数分量，$x^a$ 是与选基无关的抽象矢量。若换基
$\vect e'_i=\vect e_jM^j{}_i$，其中矩阵 $M$ 的第 $i$ 列为新基矢 $\vect e'_i$ 在旧基中的分量。同一个几何矢量的坐标列必须满足 $x=Mx'$，因而反向变化，
\[
 x'^i=(M^{-1})^i{}_jx^j,
 \qquad x'^i\vect e'_i=x^j\vect e_j.
\]
所以“把三个数排成一列”只有连同它们所依赖的基才表示矢量。基在定义线性空间时可以先视为“忘却结构的符号”：暂不赋予长度、角度或方向，随后再由点积与取向结构赋予欧氏几何含义。

抽象指标 $a,b,c$ 不取 $1,2,3$ 中的数值，它只标明对象所在的张量空间。例如 $u^av^b$ 是二阶张量，$u^av_a$ 是缩并得到的标量；把抽象指标换成具体指标并补上基，才得到可计算的分量式。后续若不讨论换基，通常仍用粗体矢量记号提高可读性，但每一个矢量恒等式都可以退回指标式审查。

\begin{advdefinition}{点积}\label{advA3:def:dot}
点积是双线性、对称、正定映射 $\mathbb R^3\times\mathbb R^3\to\mathbb R$，在正交归一基中
$\vect e_i\cdot\vect e_j=\delta_{ij}$，所以
$\vect x\cdot\vect y=x^iy^j\delta_{ij}=x^iy^i$。
\end{advdefinition}

\begin{advdefinition}{叉积}\label{advA3:def:cross}
在定向三维欧氏空间中，叉积是双线性映射
$\mathbb R^3\times\mathbb R^3\to\mathbb R^3$，由
\[
 \vect e_i\times\vect e_j=\varepsilon_{ijk}\vect e_k
\]
确定方向与大小。
\end{advdefinition}

\begin{advtheorem}{点积交换、叉积反交换}\label{advA3:thm:dot-cross-symmetry}
\[
 \vect u\cdot\vect v=\vect v\cdot\vect u,
 \qquad
 \vect u\times\vect v=-\vect v\times\vect u.
\]
\end{advtheorem}
证明只需分别使用 $\delta_{ij}=\delta_{ji}$ 与
$\varepsilon_{ijk}=-\varepsilon_{jik}$。

分量计算同时还原了熟悉的坐标公式：
\[
 \vect x\cdot\vect y=x_1y_1+x_2y_2+x_3y_3,
\]
\[
 \vect x\times\vect y=
 (x_2y_3-x_3y_2)\vect e_1
 +(x_3y_1-x_1y_3)\vect e_2
 +(x_1y_2-x_2y_1)\vect e_3.
\]
由 Cauchy--Schwarz 不等式可定义夹角
$\cos\theta=(\vect x\cdot\vect y)/(|\vect x||\vect y|)$；而叉积大小满足
\[
 |\vect x\times\vect y|^2
 =x_iy_jx_my_n\varepsilon_{ijk}\varepsilon_{mnk}
 =|\vect x|^2|\vect y|^2-(\vect x\cdot\vect y)^2,
\]
所以 $|\vect x\times\vect y|=|\vect x||\vect y|\sin\theta$，方向由空间取向和右手规则确定。

\begin{advdefinition}{三重标积}\label{advA3:def:triple-scalar}
\[
 \mathcal V(\vect u,\vect v,\vect w)
 =(\vect u\times\vect v)\cdot\vect w
 =u_iv_jw_k\varepsilon_{ijk}.
\]
它是三矢量张成平行六面体的有向体积。
\end{advdefinition}

\begin{advtheorem}{三重标积的刻画}\label{advA3:thm:triple-properties}
$\mathcal V$ 对三个变量分别线性，任意两个输入相等时为零，循环置换不变，并满足
$\mathcal V(\vect e_1,\vect e_2,\vect e_3)=1$。反过来，这四条性质唯一确定三重标积。
\end{advtheorem}
展开到基矢后，交错性把所有含重复指标的项消去，正规性和循环性把其余六项化成置换符号，因此必为
$u_iv_jw_k\varepsilon_{ijk}$。

用分量矩阵写就是
\[
 \mathcal V(\vect u,\vect v,\vect w)=
 \det\!\begin{pmatrix}
 u_1&u_2&u_3\\v_1&v_2&v_3\\w_1&w_2&w_3
 \end{pmatrix}.
\]
交换两行改变取向，所以反号；若三个矢量共面，则三行线性相关，体积为零。反过来，只要这一行列式非零，三个矢量就构成一组基。三重标积因而同时承担几何体积、线性无关判据和坐标变换 Jacobian 三种角色。

\begin{advtheorem}{三重矢积公式}\label{advA3:thm:triple-vector}
\[
 (\vect u\times\vect v)\times\vect w
 =(\vect u\cdot\vect w)\vect v-(\vect v\cdot\vect w)\vect u,
\]
\[
 \vect u\times(\vect v\times\vect w)
 =(\vect u\cdot\vect w)\vect v-(\vect u\cdot\vect v)\vect w.
\]
\end{advtheorem}
例如第一式的第 $l$ 分量为
$u_iv_jw_k\varepsilon_{ijm}\varepsilon_{mkl}$；用定理\ref{advA3:thm:epsilon-identities}化为
$u_iw_iv_l-v_iw_iu_l$，即得。

\begin{advexamplebox}{叉积的雅可比恒等式}\label{advA3:ex:jacobi-cross}
证明
\[
 (\vect u\times\vect v)\times\vect w
 +(\vect w\times\vect u)\times\vect v
 +(\vect v\times\vect w)\times\vect u=0.
\]

\AQJacobi
\end{advexamplebox}
\begin{advsolution}
\AAJacobi

逐项使用定理\ref{advA3:thm:triple-vector}：
\[
 (u\cdot w)v-(v\cdot w)u
 +(w\cdot v)u-(u\cdot v)w
 +(v\cdot u)w-(w\cdot u)v=0.
\]
每一方向的两项因点积交换律相消。
\end{advsolution}

\begin{advexamplebox}{叉积导子与对易子}\label{advA3:ex:cross-derivation}
定义 $D_{\vect u}(\vect v)=\vect u\times\vect v$，并令
$[L_1,L_2]=L_1L_2-L_2L_1$。证明
\[
 [D_{\vect u},D_{\vect v}]=D_{\vect u\times\vect v},
 \qquad
 D_{\vect u}(\vect v\times\vect w)
 =(D_{\vect u}\vect v)\times\vect w
 +\vect v\times(D_{\vect u}\vect w).
\]

\AQDerivations
\end{advexamplebox}
\begin{advsolution}
\AADerivations

对任意 $\vect w$，雅可比恒等式重排为
\[
 \vect u\times(\vect v\times\vect w)
 -\vect v\times(\vect u\times\vect w)
 =(\vect u\times\vect v)\times\vect w,
\]
即第一式。第二式直接展开右边并用三重矢积：
\[
 (u\times v)\times w+v\times(u\times w)
 =(u\cdot w)v-(v\cdot w)u+(v\cdot w)u-(u\cdot v)w,
\]
等于 $u\times(v\times w)$。
\end{advsolution}

\subsection{张量与外代数}

\begin{advdefinition}{张量与张量积}\label{advA3:def:tensor-product}
二阶张量空间是
\[
 \mathbb R^3\otimes\mathbb R^3
 =\{T^{ij}\vect e_i\vect e_j\},
\]
张量积由双线性和
$\vect e_i\otimes\vect e_j=\vect e_i\vect e_j$ 定义。高阶张量用更多有序指标表示；零阶和一阶张量分别是标量和矢量。
\end{advdefinition}
张量积一般不交换：$u^av^b$ 与 $v^au^b$ 的分量顺序不同。各阶张量组成分级代数
$T(\mathbb R^3)=\bigoplus_{n\ge0}T^n(\mathbb R^3)$。

二阶张量在给定基下有九个分量，
\[
 \mathsf T=T^{ij}\vect e_i\otimes\vect e_j.
\]
它可以看成双线性型，也可以在用度量识别矢量和对偶矢量后看成线性映射。二者分量写法相近，但指标作用不同，因而不应仅凭“像矩阵”就随意交换次序。三个矢量的张量积
$u^av^bw^c$ 属于三阶空间；一般有
\[
 T(\mathbb R^3)=T^0\oplus T^1\oplus T^2\oplus\cdots,
 \qquad T^l\otimes T^m\subset T^{l+m}.
\]
标量是零阶张量，矢量是一阶张量。这个分级解释了为什么张量积会增加指标，而缩并会减少两个指标。

\begin{advdefinition}{张量转置}\label{advA3:def:tensor-transpose}
$t_{pq}$ 交换张量第 $p,q$ 个指标。例如
$(t_{12}T)^{ab}=T^{ba}$；矩阵转置只是二阶张量的这一特例。
\end{advdefinition}

\begin{advdefinition}{张量缩并}\label{advA3:def:contraction}
缩并 $c_{pq}$ 把第 $p,q$ 个一上一下指标相等并求和，从而使阶数降低2。例如
对未缩并的二阶张量 $T^a{}_b=u^av_b$，有 $c_{12}T=T^a{}_a=u^av_a=\vect u\cdot\vect v$。二阶混合张量的缩并就是迹，不能把原来两个不同矢量误写成同一矢量的模平方。
\end{advdefinition}
抽象指标能无歧义地区分矩阵乘法的不同转置组合：
$A^{ab}B_{bc}$、$A^{ab}B_{cb}$、$A^{ba}B_{bc}$ 分别代表不同缩并。双点积也是两个指标同时缩并。

更具体地，把 $A,B$ 的分量看作矩阵，四种常见乘积分别为
\[
 A^{ab}B_{bc}\leftrightarrow AB,\qquad
 A^{ab}B_{cb}\leftrightarrow AB^{\mathsf T},
\]
\[
 A^{ba}B_{bc}\leftrightarrow A^{\mathsf T}B,\qquad
 A^{ba}B_{cb}\leftrightarrow A^{\mathsf T}B^{\mathsf T}.
\]
而
$A^{ab}B_{ba}=\operatorname{tr}(AB)$ 与
$A^{ab}B_{ab}=\operatorname{tr}(A^{\mathsf T}B)$
是两种不同的双缩并。抽象指标的价值正在于：每一次转置和缩并都直接写在指标位置上，不依赖另造符号。

\begin{advdefinition}{对称化与反对称化}\label{advA3:def:symmetrization}
\[
 T^{(a_1\cdots a_n)}=\frac1{n!}\sum_{\sigma\in S_n}
 T^{a_{\sigma(1)}\cdots a_{\sigma(n)}},
\]
\[
 T^{[a_1\cdots a_n]}=\frac1{n!}\sum_{\sigma\in S_n}
 \operatorname{sgn}(\sigma)T^{a_{\sigma(1)}\cdots a_{\sigma(n)}}.
\]
两者都是投影算子：$\mathcal S^2=\mathcal S$、$\mathcal A^2=\mathcal A$。
\end{advdefinition}

二阶情形尤其直观：
\[
 T^{(ab)}=\frac12(T^{ab}+T^{ba}),\qquad
 T^{[ab]}=\frac12(T^{ab}-T^{ba}),
\]
\[
 T^{ab}=T^{(ab)}+T^{[ab]}.
\]
对称部分在转置下不变，反对称部分在转置下反号。除以 $n!$ 是必要的：若 $T$ 已完全对称，再做一次对称化必须仍等于 $T$，而不是变成 $n!T$。完全反对称张量只保留指标互异的分量，其余含重复指标的分量全部为零。

\begin{advexamplebox}{对称张量空间的维数}\label{advA3:ex:symmetric-dimension}
证明
\[
 \dim S^n(\mathbb R^3)=\frac{(n+1)(n+2)}2,
\]
并推广到 $S^n(\mathbb R^m)$。

\AQSensors
\end{advexamplebox}
\begin{advsolution}
\AASensors

完全对称分量只由指标1、2、3各出现多少次决定。令出现次数为
$r_1,r_2,r_3\ge0$ 且 $r_1+r_2+r_3=n$，不同非负整数解的个数由插板法为
\[
 \binom{n+2}{2}=\frac{(n+1)(n+2)}2.
\]
$m$ 维时同理数 $r_1+\cdots+r_m=n$ 的非负整数解：
\[
 \boxed{\dim S^n(\mathbb R^m)=\binom{n+m-1}{m-1}}.
\]
\end{advsolution}

三维中 $\dim\Lambda^0=1,\dim\Lambda^1=3,\dim\Lambda^2=3,\dim\Lambda^3=1$。二阶交错基可取
$\vect e_1\wedge\vect e_2$、$\vect e_2\wedge\vect e_3$、$\vect e_3\wedge\vect e_1$；三阶只有体积元
$\vect e_1\wedge\vect e_2\wedge\vect e_3$。

按照本节楔积的归一化，
\[
 \vect e_i\wedge\vect e_j
 =\vect e_i\otimes\vect e_j-\vect e_j\otimes\vect e_i,
\]
而三阶体积元是六个排列项按置换符号求和：
\[
 \begin{aligned}
 \vect e_1\wedge\vect e_2\wedge\vect e_3
 ={}&\vect e_1\vect e_2\vect e_3
 +\vect e_2\vect e_3\vect e_1
 +\vect e_3\vect e_1\vect e_2\\
 &-\vect e_3\vect e_2\vect e_1
 -\vect e_1\vect e_3\vect e_2
 -\vect e_2\vect e_1\vect e_3.
 \end{aligned}
\]
由这些显式展开即可自然过渡到一般交错代数。

\begin{advtheorem}{三维四阶以上交错张量为零}\label{advA3:thm:alternating-zero}
\[
 n\ge4\quad\Longrightarrow\quad \Lambda^n(\mathbb R^3)=\{0\}.
\]
\end{advtheorem}
证明：任意基张量含 $n\ge4$ 个取值只可能为1、2、3的指标，抽屉原理保证至少两个相同。交换这两个位置既不改变基张量，又因全反对称性使其反号，故该项为零。

\begin{advdefinition}{楔积}\label{advA3:def:wedge}
若 $\alpha\in\Lambda^l$、$\beta\in\Lambda^m$，定义
\[
 \alpha\wedge\beta=\frac{(l+m)!}{l!m!}\,
 \mathcal A(\alpha\otimes\beta).
\]
它满足结合律及分次交换律
$\alpha\wedge\beta=(-1)^{lm}\beta\wedge\alpha$。
\end{advdefinition}

分次交换律中的 $(-1)^{lm}$ 来自把 $\alpha$ 的 $l$ 个指标逐个越过 $\beta$ 的 $m$ 个指标，一共需要 $lm$ 次对换。因而两个1-形式反交换，一个1-形式与一个2-形式交换，两个2-形式也交换。特别地，对任意1-形式 $u$，
$u\wedge u=0$。结合律保证可以无括号地写
$u\wedge v\wedge w$，它等于三者张量积的完全反对称部分乘以 $3!$。

\begin{advdefinition}{三维霍奇对偶}\label{advA3:def:hodge}
霍奇算子 $\star:\Lambda^l\to\Lambda^{3-l}$ 线性，并由
\[
 \star1=\vect e_1\wedge\vect e_2\wedge\vect e_3,
 \quad
 \star\vect e_1=\vect e_2\wedge\vect e_3
\]
及循环关系确定；反向又有
$\star(\vect e_2\wedge\vect e_3)=\vect e_1$。三维欧氏正定情形满足 $\star\star=1$。
\end{advdefinition}
由定义
\[
 \star(\vect x\wedge\vect y)=\vect x\times\vect y,
 \qquad
 \star(\vect x\wedge\vect y\wedge\vect z)
 =\mathcal V(\vect x,\vect y,\vect z).
\]
这说明叉积是二矢量交错张量积的霍奇对偶，三重标积是三矢量楔积的霍奇对偶。

逐个基矢核验第一式：
\[
 \star(\vect e_i\wedge\vect e_j)=\varepsilon_{ijk}\vect e_k,
\]
恰与叉积定义相同；由双线性即可推广到任意 $\vect x,\vect y$。再楔入 $\vect z$ 并取对偶，
\[
 \star(\vect x\wedge\vect y\wedge\vect z)
 =x_iy_jz_k\varepsilon_{ijk},
\]
就是三重标积。于是三维叉积并不是一般维数都自然存在的基本乘法，而是“三维二阶交错张量与一阶矢量恰好同维”带来的特殊对偶。

若记体积元
$\Omega=\vect e_1\wedge\vect e_2\wedge\vect e_3$，则内插乘积满足
\[
 \iota_{\vect u}\Omega=\star\vect u,
 \qquad \iota_{\vect v}\iota_{\vect u}\Omega
 =\star(\vect u\wedge\vect v).
\]
也可用记号 $\vect e_i\cdot\Omega=\star\vect e_i$ 表达同一关系，并由此把三重矢积重写为
\[
 (\vect x\times\vect y)\times\vect z
 =\iota_{\vect z}(\vect x\wedge\vect y)
 =(\vect x\cdot\vect z)\vect y-(\vect y\cdot\vect z)\vect x.
\]
这里右式也再次说明三重矢积不满足结合律。

\begin{advexamplebox}{外代数总维数}\label{advA3:ex:exterior-dimension}
证明 $n$ 维空间的外代数满足
\[
 \dim\Lambda(\mathbb R^n)=2^n.
\]

\AQExterior
\end{advexamplebox}
\begin{advsolution}
\AAExterior

$l$ 阶基由 $n$ 个基矢中任选 $l$ 个并按递增次序楔积得到，故
$\dim\Lambda^l=\binom nl$。于是
\[
 \dim\Lambda=\sum_{l=0}^n\binom nl=(1+1)^n=2^n.
\]
这也给出三维的 $1+3+3+1=8$。
\end{advsolution}

\begin{advexamplebox}{霍奇对偶恢复点积}\label{advA3:ex:hodge-dot}
证明对三维矢量 $\vect u,\vect v$，
\[
 \boxed{\star(\vect u\wedge\star\vect v)=\vect u\cdot\vect v}.
\]

\AQHodge
\end{advexamplebox}
\begin{advsolution}
\AAHodge

写 $\star\vect v=\tfrac12v_i\varepsilon_{ijk}\vect e_j\wedge\vect e_k$，则
\[
 \star(\vect u\wedge\star\vect v)
 =\frac12u_lv_i\varepsilon_{ijk}\varepsilon_{ljk}
 =u_lv_i\delta_{li}=u_iv_i.
\]
第二步用了 $\varepsilon_{ijk}\varepsilon_{ljk}=2\delta_{il}$。结果是标量，阶数也由
$1+(3-1)=3$ 再对偶到0阶，检查一致。
\end{advsolution}

\subsection{三维矢量微分}

在每点赋一个标量、矢量或张量分别得到标量场 $f(x)$、矢量场
$\vect F=F_i(x)\vect e_i$ 和张量场。直角坐标固定基中
\[
 \nabla f=(\partial_i f)\vect e_i,
 \quad
 \nabla\vect F=(\partial_iF_j)\vect e_i\vect e_j,
 \quad
 \nabla\cdot\vect F=\partial_iF_i,
 \quad
 \nabla\times\vect F=\varepsilon_{ijk}\partial_iF_j\vect e_k.
\]

从全微分出发，
\[
 \dd f=(\partial_if)\dd x^i=\nabla f\cdot\dd\vect x,
\]
所以梯度是“与任意位移点乘后给出一阶变化”的唯一矢量。继续微分得到 Hessian：
\[
 \dd^2f=(\partial_i\partial_jf)\dd x^i\dd x^j
 =(\dd\vect x\,\dd\vect x):\nabla\nabla f.
\]
矢量场的微分为
\[
 \dd\vect F=(\partial_iF_j)\dd x^i\vect e_j
 =\dd\vect x\cdot\nabla\vect F.
\]
这三式解释了为什么 $\nabla f$ 是矢量、$\nabla\vect F$ 是二阶张量，也解释了张量指标的次序：第一个指标来自求导方向，第二个指标来自场的分量方向。

\begin{advtheorem}{梯度的乘法与链式规则}\label{advA3:thm:gradient-rules}
\[
 \nabla(\lambda f+\mu g)=\lambda\nabla f+\mu\nabla g,
 \quad \nabla(fg)=f\nabla g+g\nabla f,
\]
\[
 \nabla F(g)=F'(g)\nabla g,
 \quad
 \nabla(f\vect F)=f\nabla\vect F+(\nabla f)\vect F,
\]
\[
 \nabla(\vect F\cdot\vect G)
 =(\nabla\vect F)\cdot\vect G+(\nabla\vect G)\cdot\vect F.
\]
\end{advtheorem}
证明逐分量应用普通乘法和链式法则即可。最后两式要保留张量积的指标顺序，不能把 $(\nabla f)\vect F$ 随意换成 $\vect F\nabla f$。

例如最后两式的分量推导为
\[
 \partial_i(fF_j)=f\partial_iF_j+(\partial_if)F_j,
\]
\[
 \partial_i(F_jG_j)=(\partial_iF_j)G_j+(\partial_iG_j)F_j.
\]
第一式右端第二项确实是 $(\nabla f)\vect F$，而非 $\vect F(\nabla f)$；后者的分量是 $F_i\partial_jf$，除非张量恰好对称，否则两者不同。许多矢量恒等式错误都来自把张量积误当成普通数乘，因此这里必须保留指标顺序。

\begin{advexamplebox}{径向函数的梯度与Hessian}\label{advA3:ex:radial-gradients}
令 $r=(x_ix_i)^{1/2}$、$\vect r=x_i\vect e_i$。计算
$\nabla x_i,\nabla r^2,\nabla r,\nabla(1/r),\nabla(\vect r/r^3)$ 以及 $\nabla\nabla r^n$（$r\ne0$）。

\AQRadialHessian
\end{advexamplebox}
\begin{advsolution}
\AARadialHessian

\[
 \boxed{\nabla x_i=\vect e_i},\qquad
 \boxed{\nabla r^2=2\vect r},\qquad
 \boxed{\nabla r=\frac{\vect r}{r}},
\]
\[
 \boxed{\nabla\frac1r=-\frac{\vect r}{r^3}},
 \qquad
 \boxed{\nabla\cdot\frac{\vect r}{r^3}=0\quad(r\ne0)}.
\]
最后一式若把“$\nabla$ 作用于矢量”理解为完整梯度，则
\[
 \nabla\left(\frac{\vect r}{r^3}\right)
 =\frac{\mathsf1}{r^3}-3\frac{\vect r\vect r}{r^5}.
\]
一般地
\[
 \boxed{\nabla\nabla r^n
 =nr^{n-2}\mathsf1+n(n-2)r^{n-4}\vect r\vect r}.
\]
其迹给出 $\nabla^2r^n=n(n+1)r^{n-2}$。$r=0$ 需按分布处理，例如
$\nabla\cdot(\vect r/r^3)=4\pi\delta^{(3)}(\vect r)$。
\end{advsolution}

外导数可写成 $\nabla\wedge(\cdot)=\vect e_i\wedge\partial_i(\cdot)$。
对标量，它就是梯度；对矢量取霍奇对偶得到旋度；对矢量的霍奇对偶再外导并对偶，得到散度。

\begin{advtheorem}{庞加莱引理}\label{advA3:thm:poincare}
外导数恒满足 $\dd^2=0$。在星形区域上，每个次数 $k\ge1$ 的闭微分形式都是恰当形式；对一般区域，这一结论局部成立。零次闭形式只是局部常数，不属于上述正次数结论。特别地
\[
 \nabla\times\nabla f=0,
 \qquad
 \nabla\cdot(\nabla\times\vect F)=0.
\]
\end{advtheorem}
$\dd^2=0$ 直接来自偏导可交换与楔积反对称：
$\partial_i\partial_j=\partial_j\partial_i$，而
$\dd x^i\wedge\dd x^j=-\dd x^j\wedge\dd x^i$。闭形式局部恰当的构造可沿从星形中心到 $x$ 的径向线积分给出。

\begin{advtheorem}{广义斯托克斯定理}\label{advA3:thm:general-stokes}
对有向光滑流形 $M$ 及紧支撑微分形式 $\omega$，
\[
 \boxed{\int_M\dd\omega=\int_{\partial M}\omega}.
\]
\end{advtheorem}
它在坐标小块上由一维微积分基本定理反复使用得到，再用单位分解拼接；内部相邻小块的公共边界因取向相反而抵消，只留下 $\partial M$。

\begin{advtheorem}{散度定理}\label{advA3:thm:divergence-theorem}
\[
 \boxed{\oint_{\partial V}\vect F\cdot\dd\vect S
 =\int_V\nabla\cdot\vect F\,\dd V}.
\]
\end{advtheorem}

\begin{advtheorem}{开尔文--斯托克斯定理}\label{advA3:thm:curl-stokes}
\[
 \boxed{\oint_{\partial S}\vect F\cdot\dd\vect l
 =\int_S(\nabla\times\vect F)\cdot\dd\vect S}.
\]
\end{advtheorem}
两式分别把广义斯托克斯定理应用于 $\star\vect F$ 与矢量场对应的1-形式；边界方向与面法向必须满足右手定向。

\begin{advtheorem}{散度运算公式}\label{advA3:thm:divergence-rules}
\[
 \nabla\cdot(f\vect F)=\nabla f\cdot\vect F+f\nabla\cdot\vect F,
\]
若并矢 $(\vect F\vect G)_{ij}=F_iG_j$ 且散度缩并第一个指标，则
\[
 \nabla\cdot(\vect F\vect G)
 =(\nabla\cdot\vect F)\vect G+(\vect F\cdot\nabla)\vect G.
\]
\[
 \nabla\cdot(\vect F\times\vect G)
 =\vect G\cdot\nabla\times\vect F-\vect F\cdot\nabla\times\vect G,
\]
并且对二阶张量 $\mathsf T$，$(\nabla\cdot\mathsf T)_j=\partial_iT_{ij}$。线性规则逐项成立。
\end{advtheorem}

这些公式都可直接审查指标。例如
\[
 [\nabla\cdot(\vect F\vect G)]_j
 =\partial_i(F_iG_j)=(\partial_iF_i)G_j+F_i\partial_iG_j,
\]
而
\[
 \nabla\cdot(\vect F\times\vect G)
 =\partial_i(\varepsilon_{jki}F_jG_k)
 =G_k\varepsilon_{kij}\partial_iF_j
 -F_j\varepsilon_{jik}\partial_iG_k.
\]
这正是 $\vect G\cdot\nabla\times\vect F-\vect F\cdot\nabla\times\vect G$。指标推导还明确了散度对并矢缩并的是哪个位置，避免把 $(\vect F\cdot\nabla)\vect G$ 错写成 $(\nabla\vect G)\cdot\vect F$ 后再随意转置。

\begin{advtheorem}{旋度与拉普拉斯公式}\label{advA3:thm:curl-rules}
\[
 \nabla\times(f\vect F)=\nabla f\times\vect F+f\nabla\times\vect F,
\]
\[
 \boxed{\nabla\times(\nabla\times\vect F)
 =\nabla(\nabla\cdot\vect F)-\nabla^2\vect F},
\]
\[
 \nabla(\vect F\cdot\vect G)
 =(\vect F\cdot\nabla)\vect G+(\vect G\cdot\nabla)\vect F
 +\vect F\times(\nabla\times\vect G)
 +\vect G\times(\nabla\times\vect F),
\]
\[
 \nabla\times(\vect F\times\vect G)
 =(\vect G\cdot\nabla)\vect F+(\nabla\cdot\vect G)\vect F
 -(\vect F\cdot\nabla)\vect G-(\nabla\cdot\vect F)\vect G.
\]
\end{advtheorem}
证明都可把叉积写成 $\varepsilon_{ijk}$ 并使用
$\varepsilon_{ijk}\varepsilon_{imn}=\delta_{jm}\delta_{kn}-\delta_{jn}\delta_{km}$；这同时避免凭图形直觉误换指标顺序。

以最关键的双旋度为例，第 $i$ 分量为
\[
 [\nabla\times(\nabla\times\vect F)]_i
 =\varepsilon_{jki}\partial_j
 (\varepsilon_{lmk}\partial_lF_m).
\]
缩并两个 $\varepsilon$ 后得到
\[
 (\delta_{il}\delta_{jm}-\delta_{im}\delta_{jl})
 \partial_j\partial_lF_m
 =\partial_i(\partial_mF_m)-\partial_j\partial_jF_i.
\]
因此
$\nabla\times\nabla\times\vect F=\nabla(\nabla\cdot\vect F)-\nabla^2\vect F$。
其余两个含叉积的公式以同一恒等式展开；每项都只使用乘法法则与偏导交换，故在直角坐标固定基中成立。

\begin{advtheorem}{标量、矢量与张量的边界积分形式}\label{advA3:thm:boundary-integrals}
在分段光滑区域上，
\[
 \oint_{\partial V}f\,\dd\vect S=\int_V\nabla f\,\dd V,
 \qquad
 \oint_{\partial V}\dd\vect S\times\vect F
 =\int_V\nabla\times\vect F\,\dd V,
\]
\[
 \oint_{\partial V}\dd\vect S\,\vect F
 =\int_V\nabla\vect F\,\dd V,
 \qquad
 \oint_{\partial V}\dd\vect S\cdot\vect F
 =\int_V\nabla\cdot\vect F\,\dd V,
\]
以及
\[
 \oint_{\partial S}f\,\dd\vect l
 =\int_S\dd\vect S\times\nabla f,
 \qquad
 \oint_{\partial S}\dd\vect l\times\vect F
 =\int_S\bigl[(\nabla\vect F)\cdot\dd\vect S
 -(\nabla\cdot\vect F)\dd\vect S\bigr],
 \]
\[
 \oint_{\partial S}\vect F\cdot\dd\vect l
 =\int_S(\nabla\times\vect F)\cdot\dd\vect S.
\]
\end{advtheorem}
证明：前四式分别把散度定理应用于 $f\vect a$、
$\vect a\cdot(\vect F\times\cdot)$、$T_{ij}$ 的每个分量和 $\vect F$，其中 $\vect a$ 为任意常矢量；任意性保证矢量或张量等式。后两式同理由开尔文--斯托克斯定理逐分量得到。

中间的线积分叉积公式可用任意常矢量 $\vect a$ 核验：
\[
 \vect a\cdot\oint\dd\vect l\times\vect F
 =\oint(\vect F\times\vect a)\cdot\dd\vect l
 =\int_S\nabla\times(\vect F\times\vect a)\cdot\dd\vect S.
\]
由于 $\vect a$ 为常量，
$\nabla\times(\vect F\times\vect a)
=(\vect a\cdot\nabla)\vect F-\vect a(\nabla\cdot\vect F)$，把 $\vect a$ 提出即得原矢量等式。这个证明也固定了两个项的正负号和张量乘法次序。

\begin{advexamplebox}{格林第一、第二恒等式}\label{advA3:ex:green-identities}
证明
\[
 \int_V(\phi\nabla^2\psi+\nabla\phi\cdot\nabla\psi)\dd V
 =\oint_{\partial V}\phi\nabla\psi\cdot\dd\vect S,
\]
\[
 \int_V(\phi\nabla^2\psi-\psi\nabla^2\phi)\dd V
 =\oint_{\partial V}(\phi\nabla\psi-\psi\nabla\phi)\cdot\dd\vect S.
\]

\AQGreen
\end{advexamplebox}
\begin{advsolution}
\AAGreen

乘法法则给出
\[
 \nabla\cdot(\phi\nabla\psi)
 =\nabla\phi\cdot\nabla\psi+\phi\nabla^2\psi.
\]
在 $V$ 上积分并用散度定理即得第一式。再交换 $\phi,\psi$ 得到另一份第一式，两式相减，交叉梯度项相消，得到第二式。若边界是分片光滑的，按每片外法向积分即可；内部人为分界面仍因法向相反而抵消。
\end{advsolution}

\section{偏微分方程}

偏微分方程把场在空间各点之间的联系与边界、初始条件同时编码。处理这类问题的主线有三条：先用自伴本征值问题建立正交函数系，再用分离变量把偏微分方程化成常微分方程，最后用正交展开或傅里叶变换确定各阶模态的系数。本节依次讨论 Sturm--Liouville 理论、三族正交多项式、波动方程与热方程，并给出所用定理的证明与例题的完整解答。

\subsection{Sturm--Liouville 问题与正交函数系}

先看固定在 $0<x<l$ 上的弦。分离变量后空间函数满足
\[
 y''+\lambda y=0,\qquad y(0)=y(l)=0.
\]
非零解只在 $\lambda_n=(n\pi/l)^2$ 时存在。取
\[
 y_n(x)=\sqrt{\frac2l}\sin\frac{n\pi x}{l},\qquad n=1,2,\ldots,
\]
便有
\[
 \int_0^l y_ny_m\,\dd x=\delta_{nm}.
\]
若 $y=\sum_n c_ny_n$，则 $c_n=\int_0^l yy_n\,\dd x$；在完备性成立时，Parseval 等式给出
$\int_0^l|y|^2\dd x=\sum_n|c_n|^2$。这正是有限维正交投影在函数空间中的延伸。

\begin{advdefinition}{带权内积、范数与完备正交系}\label{advA3:def:weighted-inner-product}
在区间 $(a,b)$ 上给定连续正权函数 $\rho(x)$，定义
\[
 (f,g)_\rho=\int_a^b f^*(x)g(x)\rho(x)\,\dd x,
 \qquad \|f\|_\rho=\sqrt{(f,f)_\rho}.
\]
若 $(z_i,z_j)_\rho=\delta_{ij}$，则 $\{z_i\}$ 是正交归一函数系；若每个属于所讨论函数空间的 $f$ 都能按范数收敛地写成 $f=\sum_i(z_i,f)_\rho z_i$，则称该函数系完备。
\end{advdefinition}

内积具有共轭对称性、线性和正定性。实函数情形可以去掉共轭号。完备性不是“每两个函数正交”自动带来的性质，它还要求这一组函数没有漏掉空间中的任何独立方向。

\begin{advexamplebox}{Cauchy--Schwarz 不等式与 Gram--Schmidt 正交化}\label{advA3:ex:cauchy-schmidt}
证明
\[
 |(f,g)_\rho|\le \|f\|_\rho\|g\|_\rho,
\]
并说明如何把线性无关函数 $z_1,\ldots,z_n$ 化为正交归一组。
\end{advexamplebox}
\begin{advsolution}
若 $g=0$ 结论显然。若 $g\ne0$，对任意复数 $c$ 有
\[
 0\le\|f-cg\|^2
 =\|f\|^2-c(f,g)-c^*(g,f)+|c|^2\|g\|^2.
\]
取 $c=(g,f)/\|g\|^2$，得到
\[
 0\le\|f\|^2-\frac{|(f,g)|^2}{\|g\|^2},
\]
即所证不等式；等号当且仅当 $f,cg$ 线性相关。

正交化递推为
\[
 u_1=z_1,\qquad y_1=\frac{u_1}{\|u_1\|},
\]
\[
 u_k=z_k-\sum_{j=1}^{k-1}(y_j,z_k)y_j,\qquad
 y_k=\frac{u_k}{\|u_k\|}.
\]
对 $i<k$，$(y_i,u_k)=(y_i,z_k)-(y_i,z_k)=0$。线性无关保证 $u_k\ne0$，因此归一化合法。每一步只是从 $z_k$ 中扣掉已知方向的投影，所以张成空间不变，最终 $\{y_k\}$ 是原空间的一组正交归一基。
\end{advsolution}

把 $1,x,x^2,\ldots$ 依次正交化可得到正交多项式。若 $Q_n$ 是 $n$ 次正交多项式，则它在 $(a,b)$ 内恰有 $n$ 个互异零点：若变号零点只有 $m<n$ 个，令 $R_m$ 为这些零点对应的一次因子之积，则 $Q_nR_m$ 在区间内保持同号，故 $(Q_n,R_m)_\rho\ne0$，与 $Q_n$ 和全部低于 $n$ 次多项式正交矛盾。常用三族为
\[
\begin{array}{c|c|c}
\text{多项式}&\text{区间}&\rho(x)\\\hline
\text{Legendre}&(-1,1)&1\\
\text{Hermite}&(-\infty,\infty)&\e^{-x^2}\\
\text{Laguerre}&(0,\infty)&x^\mu\e^{-x},\ \mu>-1
\end{array}
\]

\begin{advdefinition}{Sturm--Liouville 问题}\label{advA3:def:sl-problem}
给定区间 $[a,b]$ 上 $p(x)>0$、$\rho(x)>0$ 和实函数 $q(x)$，求
\[
 \frac{\dd}{\dd x}\!\left(p\frac{\dd y}{\dd x}\right)
 +q(x)y+\lambda\rho(x)y=0                                      \tag{A3.SL}
\]
在分离边界条件
\[
 c_1y(a)+c'_1y'(a)=0,\qquad c_2y(b)+c'_2y'(b)=0
\]
或周期边界条件 $y(a)=y(b),\ p(a)y'(a)=p(b)y'(b)$ 下的非零解。使非零解存在的 $\lambda$ 称本征值，相应的 $y$ 称本征函数。无穷区间还要求 $\int |y|^2\rho\dd x<\infty$。
\end{advdefinition}

定义算符
\[
 L=-\frac1\rho\left[\frac{\dd}{\dd x}
 \left(p\frac{\dd}{\dd x}\right)+q\right],
\]
则方程写成 $Ly=\lambda y$。固定弦、Legendre、Hermite 和 Laguerre 方程分别对应
\[
\begin{array}{c|c}
\text{问题}&\text{自伴形式}\\\hline
\text{固定弦}&y''+\lambda y=0\\
\text{Legendre}&[(1-x^2)y']'+\lambda y=0\\
\text{Hermite}&(\e^{-x^2}y')'+\lambda\e^{-x^2}y=0\\
\text{Laguerre}&(x^{\mu+1}\e^{-x}y')'+\lambda x^\mu\e^{-x}y=0
\end{array}
\]

\begin{advtheorem}{Sturm--Liouville 算符的自伴性}\label{advA3:thm:sl-self-adjoint}
在上述边界条件及允许的函数空间内，
\[
 (f,Lg)_\rho=(Lf,g)_\rho.
\]
\end{advtheorem}
\noindent\textbf{证明}\quad 两次分部积分等价于 Lagrange 恒等式
\[
 (f,Lg)_\rho-(Lf,g)_\rho
 =\bigl[p(f'^*g-f^*g')\bigr]_a^b.                         \tag{A3.1}
\]
分离边界条件在每个端点都把 $f,f'$ 与 $g,g'$ 限制为同一线性关系，因此括号为零；周期边界条件下两端贡献相消。奇异端点（例如 Legendre 的 $x=\pm1$）由 $p=0$ 消去边界项；无穷端点则由平方可积性和允许的衰减条件消去边界项，故式(A3.1)右端为零。\hfill$\square$

\begin{advtheorem}{正则 Sturm--Liouville 谱的存在性与完备性}\label{advA3:thm:sl-spectrum}
正则分离边界问题的本征值为离散实数列
\[
 \lambda_1<\lambda_2<\cdots<\lambda_n<\cdots,\qquad \lambda_n\to+\infty.
\]
每个本征值的本征空间为一维；可取实归一化本征函数 $\phi_n$，它在 $(a,b)$ 内恰有 $n-1$ 个零点。$\{\phi_n\}$ 在相应带权 $L^2$ 空间中完备。
\end{advtheorem}
这个谱定理保证后面的模态展开不是形式猜测。正则问题的一维本征空间来自二阶方程的 Wronskian 与两端边界条件；离散性和完备性可由自伴算符的 Green 函数是紧算符来证明。

\begin{advtheorem}{不同本征值的本征函数正交}\label{advA3:thm:sl-orthogonality}
归一化后
\[
 (\phi_n,\phi_m)_\rho=\delta_{nm}.
\]
\end{advtheorem}
\noindent\textbf{证明}\quad 由自伴性和 $L\phi_j=\lambda_j\phi_j$，
\[
 \lambda_m(\phi_n,\phi_m)=(\phi_n,L\phi_m)
 =(L\phi_n,\phi_m)=\lambda_n(\phi_n,\phi_m).
\]
若 $n\ne m$，则 $\lambda_n\ne\lambda_m$，内积只能为零；再逐个归一化即得。\hfill$\square$

\begin{advtheorem}{正交多项式的特征方程}\label{advA3:thm:orthogonal-polynomial-ode}
若权函数满足 Pearson 方程
\[
 \frac{\rho'}\rho=\frac{\alpha(x)}{\beta(x)}
 =\frac{\alpha_0+\alpha_1x}{\beta_0+\beta_1x+\beta_2x^2},
\]
且 $p=\rho\beta$ 在两端消失，则 $n$ 次正交多项式 $Q_n$ 满足
\[
 \frac{\dd}{\dd x}\left(pQ_n'\right)+\lambda_n\rho Q_n=0,
 \qquad \boxed{\lambda_n=-n(n+1)\beta_2-n\alpha_1}.
\]
\end{advtheorem}
\noindent\textbf{证明}\quad 对 $k<n$ 两次分部积分，边界项由 $p(a)=p(b)=0$ 消失：
\[
 \int_a^b (pQ_n')'x^k\dd x
 =k\int_a^bQ_n\{(\beta'+\alpha)x^{k-1}+(k-1)\beta x^{k-2}\}\rho\dd x=0.
\]
括号内至多为 $k$ 次多项式。因此 $(pQ_n')'/\rho$ 与所有低于 $n$ 次的多项式正交，它本身又至多是 $n$ 次，只能等于 $-\lambda_nQ_n$。比较 $x^n$ 的最高次项，得到
$\beta_2n(n-1)+(2\beta_2+\alpha_1)n+\lambda_n=0$，即所列本征值。\hfill$\square$

\begin{advexamplebox}{把一般二阶常微分方程化为自伴形式}\label{advA3:ex:ode-to-sl}
设
\[
 P(x)y''+Q(x)y'+[R_0(x)+\lambda R_1(x)]y=0,
 \]
其中 $P\ne0$。求等价的 Sturm--Liouville 形式。
\end{advexamplebox}
\begin{advsolution}
先除以 $P$，再乘积分因子 $\mu$。要使 $y''$ 与 $y'$ 合成 $(py')'$，须有
\[
 p=\mu,\qquad p'=\mu\frac QP,\qquad
 \frac{\mu'}\mu=\frac QP.
\]
故
\[
 \mu(x)=C\exp\!\left(\int^x\frac{Q(s)}{P(s)}\dd s\right).
\]
于是
\[
 (\mu y')'+\mu\frac{R_0}{P}y+\lambda\mu\frac{R_1}{P}y=0,
\]
即
\[
 \boxed{p=\mu,\quad q=\mu R_0/P,\quad \rho=\mu R_1/P}.
\]
若要成为正则 SL 问题，还需检查所选区间上 $p,\rho>0$ 以及边界条件确实消去边界型。各项量纲都等于 $\mu y''$ 的量纲。
\end{advsolution}

\begin{advexamplebox}{Bessel 方程的自伴形式}\label{advA3:ex:bessel-sl}
将
\[
 x^2y''+xy'+(x^2-\alpha^2)y=0
\]
化为自伴形式，并判断它是否存在一族有限次正交多项式解。
\end{advexamplebox}
\begin{advsolution}
除以 $x$ 即得
\[
 \boxed{(xy')'+\left(x-\frac{\alpha^2}{x}\right)y=0}.
\]
若把 $\alpha^2$ 当本征参数，可写成
$(xy')'+xy-\alpha^2x^{-1}y=0$；若把固定区间边界条件决定的径向波数 $k$ 当本征参数，则标准形式是
\[
 (xy')'+\left(k^2x-\frac{\alpha^2}{x}\right)y=0,
\]
权函数为 $\rho=x$。它不满足定理\ref{advA3:thm:orthogonal-polynomial-ode}所需的 $q=0$（或可并入常数本征值）条件，幂级数也不会在一般阶次截断，故解为 Bessel 函数而不是有限次正交多项式。只有特殊退化参数可能偶然化简，不能构成完整多项式族。
\end{advsolution}

\begin{advexamplebox}{Chebyshev 多项式}\label{advA3:ex:chebyshev}
讨论
\[
 (1-x^2)y''-xy'+\lambda y=0,\qquad -1<x<1
\]
在自然边界条件 $\lim_{x\to\pm1}\sqrt{1-x^2}\,y'=0$ 下的本征值，并写出零至四次的归一化多项式解。
\end{advexamplebox}
\begin{advsolution}
令 $x=\cos\theta$，则
\[
 (1-x^2)y''-xy'=\frac{\dd^2y}{\dd\theta^2}.
\]
边界条件转化为 $y_\theta(0)=y_\theta(\pi)=0$。对 $y_{\theta\theta}+\lambda y=0$ 求解可得 $\lambda\ge0$，且 $\sqrt\lambda$ 必须为整数，故
$y=\cos n\theta$、$\lambda_n=n^2$。因此第一类 Chebyshev 多项式
\[
 T_n(x)=\cos(n\arccos x)
\]
的前五个为
\[
 T_0=1, T_1=x, T_2=2x^2-1, T_3=4x^3-3x,
\ T_4=8x^4-8x^2+1.
\]
仅要求端点有限不足以确定这些本征值，例如 $\cos(\alpha\arccos x)$ 对非整数 $\alpha$ 也有限，但不满足上述两端条件。其权函数为 $\rho=(1-x^2)^{-1/2}$，且
\[
 \int_{-1}^1\frac{T_nT_m}{\sqrt{1-x^2}}\dd x
 =\begin{cases}\pi,&n=m=0,\\\pi/2,&n=m\ge1,\\0,&n\ne m.
 \end{cases}
\]
故归一化解是 $\widehat T_0=T_0/\sqrt\pi$、
$\widehat T_n=\sqrt{2/\pi}\,T_n$（$n\ge1$）。代回原方程可逐项检查；$\lambda_n$ 无量纲，与角变量二阶导数一致。
\end{advsolution}

\begin{advtheorem}{Rodrigues 公式}\label{advA3:thm:rodrigues}
在定理\ref{advA3:thm:orthogonal-polynomial-ode}的条件下，一组与 $Q_n$ 成比例的正交多项式为
\[
 \boxed{Q_n(x)=\frac1{\rho(x)}\frac{\dd^n}{\dd x^n}
 \bigl[\rho(x)\beta^n(x)\bigr]}.
\]
\end{advtheorem}
\noindent\textbf{证明}\quad 令
\[
 Q_n^{(k)}=\frac1{\rho\beta^{n-k}}
 \frac{\dd^k}{\dd x^k}(\rho\beta^n),\qquad k=0,1,\ldots,n.
\]
由 $\rho'\beta=\rho\alpha$，
\[
 Q_n^{(k+1)}=[\alpha+(n-k)\beta']Q_n^{(k)}
 +\beta(Q_n^{(k)})',
\]
所以其次数逐步增加，$Q_n^{(n)}$ 是 $n$ 次多项式。对 $m<n$ 分部积分 $n$ 次，端点项由 $\rho\beta^j$ 的衰减消失，
\[
 \int_a^b x^mQ_n\rho\dd x
 =(-1)^n\int_a^b\frac{\dd^n x^m}{\dd x^n}\rho\beta^n\dd x=0.
\]
故它与所有低次多项式正交。\hfill$\square$

\begin{advtheorem}{正交多项式的母函数}\label{advA3:thm:generating-function}
令 $z_1=z_1(z,\varepsilon)$ 是方程
\[
 \zeta=z+\varepsilon\beta(\zeta)
\]
中当 $\varepsilon\to0$ 时趋近 $z$ 的根，则 Rodrigues 多项式的指数型母函数为
\[
 \boxed{\Psi(z,\varepsilon)=
 \frac{\rho(z_1)}{\rho(z)}\frac1{1-\varepsilon\beta'(z_1)}
 =\sum_{n=0}^\infty\frac{Q_n(z)}{n!}\varepsilon^n}.
\]
\end{advtheorem}
\noindent\textbf{证明}\quad 把 Rodrigues 公式代入级数并用 Cauchy 积分公式，
\[
 \Psi=\frac1{2\pi\ii\rho(z)}\oint
 \frac{\rho(\zeta)}{\zeta-z-\varepsilon\beta(\zeta)}\dd\zeta.
\]
取只包围邻近根 $z_1$ 的小围道。分母对 $\zeta$ 的导数为
$1-\varepsilon\beta'(z_1)$，由留数定理立刻得到结论。这个证明同时说明级数的收敛范围由另一个根或 $\rho$ 的奇点决定。\hfill$\square$

\subsection{Legendre、Hermite 与 Laguerre 多项式}

对 Legendre 方程，$\rho=1$、$\beta=1-x^2$，故
\[
 \lambda_n=n(n+1),\qquad
 P_n(x)=\frac1{2^nn!}\frac{\dd^n}{\dd x^n}(x^2-1)^n,
\]
\[
 \frac1{\sqrt{1-2xt+t^2}}=\sum_{n=0}^\infty P_n(x)t^n.
\]
对母函数求导并比较 $t^n$ 系数，得到 Bonnet 递推式
\[
 (2n+1)xP_n=(n+1)P_{n+1}+nP_{n-1}.
\]
又由 $P_n(1)=1$、$P_n(-1)=(-1)^n$ 和正交性，
\[
 \int_{-1}^1P_nP_m\dd x=\frac2{2n+1}\delta_{nm}.
\]
前六个多项式为
\[
 P_0=1,\quad P_1=x,\quad P_2=\frac{3x^2-1}{2},\quad
 P_3=\frac{5x^3-3x}{2},
\]
\[
 P_4=\frac{35x^4-30x^2+3}{8},\qquad
 P_5=\frac{63x^5-70x^3+15x}{8}.
\]

\begin{center}\begin{minipage}{\linewidth}\centering
\centering
\begin{tikzpicture}
\begin{axis}[width=.82\textwidth,height=6.2cm,domain=-1:1,xmin=-1,xmax=1,ymin=-1.08,ymax=1.08,
 axis lines=middle,grid=major,xlabel={$x$},ylabel={$P_n(x)$},
 legend style={font=\scriptsize,at={(0.5,-0.18)},anchor=north,legend columns=3},
 every axis plot/.append style={line width=.8pt},samples=181]
 \addplot[DeepBlue] {1}; \addlegendentry{$P_0$}
 \addplot[orange!80!black] {x}; \addlegendentry{$P_1$}
 \addplot[green!55!black] {(3*x^2-1)/2}; \addlegendentry{$P_2$}
 \addplot[red!75!black] {(5*x^3-3*x)/2}; \addlegendentry{$P_3$}
 \addplot[violet] {(35*x^4-30*x^2+3)/8}; \addlegendentry{$P_4$}
 \addplot[brown!80!black] {(63*x^5-70*x^3+15*x)/8}; \addlegendentry{$P_5$}
\end{axis}
\end{tikzpicture}
\bookcaption{前六个 Legendre 多项式。每个 $P_n$ 在 $(-1,1)$ 内恰有 $n$ 个零点，端点值固定为 $\pm1$。}\label{advA3:fig:legendre}
\end{minipage}\end{center}

Hermite 方程给出
\[
 \lambda_n=2n,\qquad H_n(x)=(-1)^n\e^{x^2}\frac{\dd^n}{\dd x^n}\e^{-x^2},
\]
\[
 \e^{2xt-t^2}=\sum_{n=0}^\infty H_n(x)\frac{t^n}{n!},\qquad
 H_n'=2nH_{n-1},\qquad 2xH_n=H_{n+1}+2nH_{n-1}.
\]
利用 Rodrigues 公式分部积分 $n$ 次，
\[
 \int_{-\infty}^{\infty}H_nH_m\e^{-x^2}\dd x
 =2^nn!\sqrt\pi\,\delta_{nm}.
\]
前六个为
$H_0=1$、$H_1=2x$、$H_2=4x^2-2$、$H_3=8x^3-12x$、
$H_4=16x^4-48x^2+12$、$H_5=32x^5-160x^3+120x$。

\begin{center}\begin{minipage}{\linewidth}\centering
\centering
\begin{tikzpicture}
\begin{axis}[width=.82\textwidth,height=6.2cm,domain=-2.1:2.1,xmin=-2.1,xmax=2.1,ymin=-45,ymax=50,
 axis lines=middle,grid=major,xlabel={$x$},ylabel={$H_n(x)$},restrict y to domain=-55:60,
 legend style={font=\scriptsize,at={(0.5,-0.18)},anchor=north,legend columns=3},samples=221]
 \addplot[DeepBlue] {1}; \addlegendentry{$H_0$}
 \addplot[orange!80!black] {2*x}; \addlegendentry{$H_1$}
 \addplot[green!55!black] {4*x^2-2}; \addlegendentry{$H_2$}
 \addplot[red!75!black] {8*x^3-12*x}; \addlegendentry{$H_3$}
 \addplot[violet] {16*x^4-48*x^2+12}; \addlegendentry{$H_4$}
 \addplot[brown!80!black] {32*x^5-160*x^3+120*x}; \addlegendentry{$H_5$}
\end{axis}
\end{tikzpicture}
\bookcaption{前六个物理学家约定的 Hermite 多项式；纵轴截断仅为同时显示低阶曲线。}\label{advA3:fig:hermite}
\end{minipage}\end{center}

$\mu$ 阶 Laguerre 方程给出
\[
 xL_{n,\mu}''+(\mu+1-x)L_{n,\mu}'+nL_{n,\mu}=0,
\]
\[
 L_{n,\mu}(x)=\frac{x^{-\mu}\e^x}{n!}
 \frac{\dd^n}{\dd x^n}(\e^{-x}x^{n+\mu}),\qquad
 \frac{\e^{-xt/(1-t)}}{(1-t)^{\mu+1}}
 =\sum_{n=0}^\infty L_{n,\mu}(x)t^n.
\]
其范数和递推关系为
\[
 \int_0^\infty L_{n,\mu}L_{m,\mu}x^\mu\e^{-x}\dd x
 =\frac{\Gamma(n+\mu+1)}{n!}\delta_{nm},
\]
\[
 xL_{n,\mu}=(2n+\mu+1)L_{n,\mu}
 -(n+\mu)L_{n-1,\mu}-(n+1)L_{n+1,\mu},
\]
\[
 \frac{\dd^k}{\dd x^k}L_{n,\mu}=(-1)^kL_{n-k,\mu+k}\quad(k\le n).
\]
前四个为
\[
 L_{0,\mu}=1,\qquad L_{1,\mu}=-x+\mu+1,
\]
\[
 L_{2,\mu}=\frac{x^2}{2}-(\mu+2)x+\frac{(\mu+1)(\mu+2)}2,
\]
\[
 L_{3,\mu}=-\frac{x^3}{6}+\frac{(\mu+3)x^2}{2}
 -\frac{(\mu+2)(\mu+3)x}{2}
 +\frac{(\mu+1)(\mu+2)(\mu+3)}6.
\]

\begin{center}\begin{minipage}{\linewidth}\centering
\centering
\begin{tikzpicture}
\begin{axis}[width=.82\textwidth,height=6.2cm,domain=0:8,xmin=0,xmax=8,ymin=-5,ymax=8,
 axis lines=middle,grid=major,xlabel={$x$},ylabel={$L_{n,\mu}(x)$},restrict y to domain=-6:9,
 legend style={font=\scriptsize,at={(0.5,-0.18)},anchor=north,legend columns=4},samples=221]
 \addplot[DeepBlue] {1}; \addlegendentry{$L_{0,0}$}
 \addplot[orange!80!black] {1-x}; \addlegendentry{$L_{1,0}$}
 \addplot[green!55!black] {x^2/2-2*x+1}; \addlegendentry{$L_{2,0}$}
 \addplot[red!75!black] {-x^3/6+1.5*x^2-3*x+1}; \addlegendentry{$L_{3,0}$}
 \addplot[violet,dashed] {2-x}; \addlegendentry{$L_{1,1}$}
 \addplot[brown!80!black,dashed] {x^2/2-3*x+3}; \addlegendentry{$L_{2,1}$}
 \addplot[cyan!60!black,dashed] {-x^3/6+2*x^2-6*x+4}; \addlegendentry{$L_{3,1}$}
\end{axis}
\end{tikzpicture}
\bookcaption{若干广义 Laguerre 多项式。实线为 $\mu=0$，虚线为 $\mu=1$；每条曲线的次数和零点数可直接核对。}\label{advA3:fig:laguerre}
\end{minipage}\end{center}

\begin{advexamplebox}{Legendre 多项式的导数递推}\label{advA3:ex:legendre-derivative-recurrence}
证明
\[
 P'_{n+1}(x)-xP'_n(x)=(n+1)P_n(x).
\]
\end{advexamplebox}
\begin{advsolution}
从 Bonnet 递推式
\[
 (n+1)P_{n+1}=(2n+1)xP_n-nP_{n-1}
\]
求导，再使用上一阶关系消去 $P'_{n-1}$ 也可证明。更直接地，对母函数
$G=(1-2xt+t^2)^{-1/2}$ 求 $x$ 导数：
\[
 \partial_xG=t(1-2xt+t^2)^{-3/2}.
\]
另一方面
\[
 (1-xt)\partial_xG-t^2\partial_tG=tG.
\]
比较 $t^{n+1}$ 系数得到
$P'_{n+1}-xP'_n=(n+1)P_n$。在 $n=0$ 时左边是 $P'_1=1$，右边也是1，起始阶检查通过。
\end{advsolution}

\begin{advexamplebox}{Legendre 多项式导数的范数}\label{advA3:ex:legendre-derivative-norm}
证明
\[
 \int_{-1}^1[P'_n(x)]^2\dd x=n(n+1).
\]
\end{advexamplebox}
\begin{advsolution}
由导数递推反复展开，
\[
 P'_n=(2n-1)P_{n-1}+(2n-5)P_{n-3}+\cdots.
\]
各项彼此正交，故
\[
 \int_{-1}^1(P'_n)^2\dd x
 =\sum_{j\ge0}(2n-4j-1)^2\frac{2}{2(n-2j-1)+1}.
\]
每项约去一个因子，成为
$2\sum_{j\ge0}(2n-4j-1)$，末项取到正数1或3。这个等差和等于 $n(n+1)$。

也可用归纳：由 $P'_{n+1}=xP'_n+(n+1)P_n$，平方积分。分部积分给出
$\int xP'_nP_n\dd x=1-\tfrac12\int P_n^2\dd x$，再代入
$\int P_n^2\dd x=2/(2n+1)$ 和由微分方程导出的
$\int x^2(P'_n)^2=\int(P'_n)^2-n(n+1)\int P_n^2$，即可由 $n$ 推到 $n+1$。结果为纯数，与 $x$ 无量纲一致。
\end{advsolution}

\begin{advexamplebox}{Legendre 多项式的偶极矩}\label{advA3:ex:legendre-dipole}
计算
\[
 \mathcal P_{ij}=\int_{-1}^1xP_i(x)P_j(x)\dd x.
\]
\end{advexamplebox}
\begin{advsolution}
由 Bonnet 递推
\[
 xP_i=\frac{i+1}{2i+1}P_{i+1}+\frac{i}{2i+1}P_{i-1}.
\]
与 $P_j$ 内积并用正交性，只有 $j=i\pm1$ 非零：
\[
 \boxed{\mathcal P_{ij}=
 \frac{2(i+1)}{(2i+1)(2i+3)}\delta_{j,i+1}
 +\frac{2i}{(2i+1)(2i-1)}\delta_{j,i-1}}.
\]
交换 $i,j$ 后数值不变，符合积分本身的对称性；$i=j$ 时被积函数奇偶性也立刻给出零。
\end{advsolution}

三族多项式的核心数据可压缩为下表，后续量子力学和势函数展开会反复使用：
\begingroup\setlength{\arraycolsep}{3pt}
\[
\begin{array}{c|c|c|c}
&\text{Legendre}&\text{Hermite}&\text{Laguerre }(\mu=0)\\\hline
\text{方程}&(1-x^2)y''-2xy'+n(n+1)y=0&y''-2xy'+2ny=0&xy''+(1-x)y'+ny=0\\
\text{权函数}&1&\e^{-x^2}&\e^{-x}\\
\text{范数平方}&2/(2n+1)&2^nn!\sqrt\pi&1
\end{array}
\]
\endgroup

\subsection{常见二阶偏微分方程及其分类}

物理中最常见的线性二阶方程包括
\[
 \text{波动方程：}\quad \partial_t^2u-c^2\nabla^2u=0,
\]
\[
 \text{热方程：}\quad \partial_tu-\kappa\nabla^2u=0,
\]
\[
 \text{Schr\"odinger 方程：}\quad
 \ii\hbar\partial_t\psi=H(\vect r,-\ii\hbar\nabla)\psi,
\]
\[
 \text{Laplace/Poisson/Helmholtz 方程：}\quad
 \nabla^2\phi=0,\quad \nabla^2\phi=-\rho/\varepsilon_0,\quad
 \nabla^2\phi+k^2\phi=0.
\]
柱坐标与球坐标下的 Laplace 算符分别为
\[
 \nabla^2u=\frac1r\partial_r(r\partial_ru)
 +\frac1{r^2}\partial_\varphi^2u+\partial_z^2u,
\]
\[
 \nabla^2u=\frac1{r^2}\partial_r(r^2\partial_ru)
 +\frac1{r^2\sin\theta}\partial_\theta(\sin\theta\,\partial_\theta u)
 +\frac1{r^2\sin^2\theta}\partial_\varphi^2u.
\]
这些方程都是线性的，最高偏导数为二阶；除 Poisson 方程外，上列方程在右端为零时齐次。坐标变换会让常系数方程出现变系数，但不会改变其几何类型。

\begin{advdefinition}{二阶线性偏微分方程的类型}\label{advA3:def:pde-classification}
在一点把二阶主部写成
\[
 \sum_{i,j=1}^n a_{ij}\partial_i\partial_j u,
\]
其中 $A=(a_{ij})$ 取对称部分。若 $A$ 非退化且全部特征值同号，称为椭圆型；若 $A$ 有且只有一个零特征值、其余同号，称为抛物型；若 $A$ 非退化，且一正其余负或一负其余正，称为双曲型；若正、负特征值各多于一个，则称为超双曲型。更高维中若零特征值多于一个，还需视作退化型单独说明。
\end{advdefinition}

\begin{advexamplebox}{典型物理方程的分类}\label{advA3:ex:pde-types}
判断波动方程、热方程、Schr\"odinger 方程、Laplace 方程、Poisson 方程和 Helmholtz 方程的类型。
\end{advexamplebox}
\begin{advsolution}
把时间也视为自变量。$n$ 维波动方程主部矩阵为
$\operatorname{diag}(1,-c^2,\ldots,-c^2)$，有一正 $n$ 负，故为双曲型。热方程关于 $(t,\vect x)$ 的二阶主部是
$\operatorname{diag}(0,-\kappa,\ldots,-\kappa)$，恰有一个零特征值，故为抛物型。

自由 Schr\"odinger 方程具有虚数时间系数，不能直接套用上述实系数标量分类；应单列为色散型的 Schr\"odinger 方程。虽然对时间一阶、对空间二阶，但它不是热方程式的抛物演化：空间傅里叶模随时间乘以模为1的相位，而不是热方程的衰减因子。Laplace 方程的主部矩阵为单位阵，故为椭圆型；Poisson 的非齐次项不改变主部，仍为椭圆型；Helmholtz 的 $k^2\phi$ 是零阶项，也不改变类型，仍为椭圆型。分类只看最高阶项，因此势能、源项和常数质量项不改变结论。
\end{advsolution}

\subsection{分离变量、傅里叶分析与初边值问题}

\paragraph{先检查边界条件是否允许解存在。}
对 $-\Delta u=f$、外法向导数 $\partial_nu=g$，散度定理要求
\[
 \int_\Omega f\,\dd V+\oint_{\partial\Omega}g\,\dd S=0.
\]
纯 Neumann 问题若不满足此式便无解；满足时仍可加任意常数，需指定平均值来唯一化。它是守恒条件，不是计算结束后的可选检查。
例如 $-u''=f(x)$、$u'(0)=u'(L)=0$，必须有 $\int_0^L f\dd x=0$。若 $f=\cos(\pi x/L)$，零平均解为
\[
 u=\frac{L^2}{\pi^2}\cos\frac{\pi x}{L};
\]
若 $f=1$ 则不可能满足两端零通量。谱展开中同一条件表现为零本征值对应的常数模必须与源正交，不能把其系数除以零。


对一维波动方程
\[
 \partial_t^2y-c^2\partial_x^2y=0
\]
令 $y=f(x)g(t)$，则
\[
 \frac{g''}{g}=c^2\frac{f''}{f}=-\omega^2,\qquad k=\omega/c.
\]
于是
\[
 g=A\cos\omega t+B\sin\omega t,\qquad
 f=C\cos kx+D\sin kx.
\]
复指数形式把两列行波清楚分开：
\[
 y_+=C_+\e^{\ii(kx-\omega t)},\qquad
 y_-=C_-\e^{\ii(-kx-\omega t)},\qquad \omega=c|k|.
\]
若分离常数取相反符号，会出现 $\cosh$ 或 $\e^{\pm\beta x}$。在无限空间中这类解通常因不平方可积而舍去；在有限区域或倏逝波问题中则可能正是物理解。

\begin{advtheorem}{Fourier 反演与 Dirac $\delta$ 分布}\label{advA3:thm:fourier-inversion}
采用约定
\[
 \widetilde f(k)=\frac1{\sqrt{2\pi}}\int_{-\infty}^{\infty}f(x)\e^{-\ii kx}\dd x,\qquad
 f(x)=\frac1{\sqrt{2\pi}}\int_{-\infty}^{\infty}\widetilde f(k)\e^{\ii kx}\dd k,
\]
则分布意义下
\[
 \int_{-\infty}^{\infty}\e^{\ii kx}\dd x=2\pi\delta(k),
\]
其中 $\int\delta(k)\varphi(k)\dd k=\varphi(0)$。
\end{advtheorem}
\noindent\textbf{证明}\quad 用 Gaussian 核正则化：
\[
 I_\varepsilon(k)=\int_{-\infty}^{\infty}
 \e^{-\varepsilon x^2+\ii kx}\dd x
 =\sqrt{\frac\pi\varepsilon}\e^{-k^2/(4\varepsilon)}.
\]
对任意光滑快速衰减的 $\varphi$，令 $k=2\sqrt\varepsilon,u$，
\[
 \frac1{2\pi}\int I_\varepsilon(k)\varphi(k)\dd k
 =\frac1{\sqrt\pi}\int\e^{-u^2}\varphi(2\sqrt\varepsilon u)\dd u
 \longrightarrow\varphi(0).
\]
故 $I_\varepsilon\to2\pi\delta$，再代入 Fourier 变换的两重积分即得反演。\hfill$\square$

\begin{advexamplebox}{Dirichlet 积分}\label{advA3:ex:dirichlet-integral}
证明
\[
 \int_{-\infty}^{\infty}\frac{\sin kx}{x}\dd x=\pi\operatorname{sgn}k.
\]
\end{advexamplebox}
\begin{advsolution}
设 $k>0$，先加收敛因子
\[
 F_\varepsilon(k)=\int_0^\infty\e^{-\varepsilon x}
 \frac{\sin kx}{x}\dd x.
\]
对 $k$ 求导可以消去 $1/x$：
\[
 F_\varepsilon'(k)=\int_0^\infty\e^{-\varepsilon x}\cos kx\dd x
 =\frac{\varepsilon}{\varepsilon^2+k^2}.
\]
由 $F_\varepsilon(0)=0$，
$F_\varepsilon(k)=\arctan(k/\varepsilon)\to\pi/2$。被积函数为偶函数，因此全轴积分为 $\pi$。$k<0$ 时对 $k$ 为奇函数，结果反号；$k=0$ 时积分为零。由此得到统一表达式 $\pi\operatorname{sgn}k$。
\end{advsolution}

\begin{advtheorem}{d'Alembert 公式}\label{advA3:thm:dalembert}
无界弦初值问题
\[
 y_{tt}-c^2y_{xx}=0,\qquad y(x,0)=f(x),\quad y_t(x,0)=g(x)
\]
的解为
\[
 \boxed{y(x,t)=\frac{f(x-ct)+f(x+ct)}2
 +\frac1{2c}\int_{x-ct}^{x+ct}g(\xi)\dd\xi}.
\]
\end{advtheorem}
\noindent\textbf{证明}\quad 改用特征变量 $\xi=x-ct$、$\eta=x+ct$，则
\[
 \partial_t^2-c^2\partial_x^2=-4c^2\partial_\xi\partial_\eta.
\]
方程变成 $y_{\xi\eta}=0$，积分得 $y=F(\xi)+G(\eta)$。初始条件给出
\[
 F(x)+G(x)=f(x),\qquad -cF'(x)+cG'(x)=g(x).
\]
把第二式积分并与第一式联立，常数在 $F,G$ 中相消，得到所列公式。直接代入 $t=0$ 得 $y=f$；对 $t$ 求导再令 $t=0$ 得 $y_t=g$，方向和端点符号均得到核验。\hfill$\square$

有界弦的边界条件决定允许的空间本征函数。例如
\[
\begin{array}{c|c}
y(0)=y(l)=0&\sin(n\pi x/l),\ n\ge1\\
y(0)=0,\ y'(l)=0&\sin[(n+\tfrac12)\pi x/l],\ n\ge0\\
y'(0)=0,\ y(l)=0&\cos[(n+\tfrac12)\pi x/l],\ n\ge0\\
y'(0)=y'(l)=0&\cos(n\pi x/l),\ n\ge0
\end{array}
\]
若 $y(0,t)=0$、$y_x(l,t)=0$，则 $k_nl=(n+\tfrac12)\pi$，展开
\[
 y=\sum_{n=0}^\infty[A_n\cos(k_nct)+B_n\sin(k_nct)]\sin k_nx,
\]
其中
\[
 A_n=\frac2l\int_0^lf(x)\sin k_nx\dd x,\qquad
 B_n=\frac2{k_ncl}\int_0^lg(x)\sin k_nx\dd x.
\]
周期边界则同时需要 $\sin(2n\pi x/l)$ 与 $\cos(2n\pi x/l)$ 两支。

\begin{advdefinition}{$n$ 维 Fourier 变换}\label{advA3:def:nd-fourier}
对 $\mathbb R^n$ 上平方可积函数，采用对称归一化
\[
 \widetilde f(\vect k)=\frac1{(2\pi)^{n/2}}
 \int_{\mathbb R^n}f(\vect x)\e^{-\ii\vect k\cdot\vect x}\dd^n x,
\]
\[
 f(\vect x)=\frac1{(2\pi)^{n/2}}
 \int_{\mathbb R^n}\widetilde f(\vect k)\e^{\ii\vect k\cdot\vect x}\dd^n k.
\]
在这一约定下 $\mathcal F^{-1}=\mathcal F^*$，并有
$\mathcal F^2f(\vect x)=f(-\vect x)$。
\end{advdefinition}

\begin{advexamplebox}{Parseval 与 Plancherel 定理}\label{advA3:ex:parseval-plancherel}
证明
\[
 \int_{-\infty}^{\infty}f^*(x)g(x)\dd x
 =\int_{-\infty}^{\infty}\widetilde f^*(k)\widetilde g(k)\dd k,
\]
并由此得到 $\|f\|_2=\|\widetilde f\|_2$。
\end{advexamplebox}
\begin{advsolution}
把 $g$ 的反变换代入左边并交换积分次序：
\[
 \int f^*(x)g(x)\dd x
 =\frac1{\sqrt{2\pi}}\int\!\int
 f^*(x)\widetilde g(k)\e^{\ii kx}\dd k\dd x.
\]
括号中的 $x$ 积分是
\[
 \frac1{\sqrt{2\pi}}\int f^*(x)\e^{\ii kx}\dd x
 =\widetilde f^*(k),
\]
所以得到 Parseval 等式。令 $g=f$ 即得
$\int|f|^2=\int|\widetilde f|^2$，这就是 Plancherel 定理。它说明 Fourier 变换是 $L^2$ 空间的酉变换，量纲上两边还通过 $\widetilde f$ 所含的长度因子严格匹配。
\end{advsolution}

\begin{advexamplebox}{标准正态分布的 Fourier 变换}\label{advA3:ex:gaussian-fourier}
求
\[
 f(x)=\frac1{\sqrt{2\pi}}\e^{-x^2/2}
\]
在本节对称归一化约定下的 Fourier 变换。
\end{advexamplebox}
\begin{advsolution}
配方
\[
 -\frac{x^2}{2}-\ii kx=-\frac{(x+\ii k)^2}{2}-\frac{k^2}{2}.
\]
整函数 $\e^{-z^2/2}$ 的积分路径可平移，或直接使用 Gaussian 积分，因此
\[
 \widetilde f(k)=\frac1{2\pi}\int_{-\infty}^{\infty}
 \e^{-x^2/2-\ii kx}\dd x
 =\boxed{\frac1{\sqrt{2\pi}}\e^{-k^2/2}}.
\]
标准 Gaussian 在这一归一化下是 Fourier 变换的不动函数。取 $k=0$ 得
$\widetilde f(0)=1/\sqrt{2\pi}$，也等于 $(2\pi)^{-1/2}\int f\dd x$，常数检查一致。
\end{advsolution}

\begin{advexamplebox}{非齐次边界条件的齐次化}\label{advA3:ex:inhomogeneous-boundary}
讨论一维波动方程
\[
 y_{tt}-c^2y_{xx}=0,\qquad y(0,t)=a(t),\quad y(l,t)=b(t)
\]
的求解方法。
\end{advexamplebox}
\begin{advsolution}
在固定边界 $x=l$ 上，边值只能随时间变化。构造满足两端边界条件的线性函数
\[
 h(x,t)=\left(1-\frac xl\right)a(t)+\frac xl b(t),
\qquad u=y-h.
\]
于是 $u(0,t)=u(l,t)=0$，而
\[
 u_{tt}-c^2u_{xx}=-h_{tt}equiv F(x,t).
\]
展开
\[
 u(x,t)=\sum_{n=1}^\infty q_n(t)\sin\frac{n\pi x}{l},\qquad
 F_n(t)=\frac2l\int_0^lF(x,t)\sin\frac{n\pi x}{l}\dd x,
\]
得到每个模态
\[
 q_n''+\omega_n^2q_n=F_n(t),\qquad \omega_n=cn\pi/l.
\]
给定初值后用受迫振子公式求 $q_n$。最后 $y=u+h$。这样边界条件被精确满足，且 $a=b=0$ 时退回齐次边界结果。
\end{advsolution}

\begin{advexamplebox}{非齐次波动方程的模态解}\label{advA3:ex:forced-wave}
求解思路：
\[
 y_{tt}-c^2y_{xx}=A(x,t),\qquad y(0,t)=y(l,t)=0,
\]
并给出零初值时的显式模态公式。
\end{advexamplebox}
\begin{advsolution}
令
\[
 y=\sum_{n=1}^\infty q_n(t)\sin k_nx,\qquad
 A=\sum_{n=1}^\infty A_n(t)\sin k_nx,\qquad k_n=n\pi/l,
\]
其中
$A_n(t)=\tfrac2l\int_0^lA(x,t)\sin k_nx\dd x$。正交投影后
\[
 q_n''+\omega_n^2q_n=A_n(t),\qquad \omega_n=ck_n.
\]
零初值的 Green 函数解为
\[
 \boxed{q_n(t)=\frac1{\omega_n}\int_0^t
 A_n(\tau)\sin[\omega_n(t-\tau)]\dd\tau}.
\]
所以
\[
 y(x,t)=\sum_{n=1}^\infty\frac{\sin k_nx}{\omega_n}
 \int_0^tA_n(\tau)\sin[\omega_n(t-\tau)]\dd\tau.
\]
逐项代回时，每个模态都产生 $A_n(t)$；在 $t=0$ 时积分及其一阶导数均为零，初边值条件全部满足。
\end{advsolution}

\begin{advexamplebox}{无界域热方程的热核}\label{advA3:ex:heat-kernel}
证明初值问题
\[
 u_t-\kappa u_{xx}=0,\qquad u(x,0)=f(x),\quad x\in\mathbb R
\]
的解为
\[
 u(x,t)=\frac1{\sqrt{4\pi\kappa t}}
 \int_{-\infty}^{\infty}
 \exp\!\left[-\frac{(x-\xi)^2}{4\kappa t}\right]f(\xi)\dd\xi.
\]
\end{advexamplebox}
\begin{advsolution}
对 $x$ 作 Fourier 变换，$\widetilde{u_{xx}}=-k^2\widetilde u$，故
\[
 \partial_t\widetilde u+\kappa k^2\widetilde u=0,\qquad
 \widetilde u(k,t)=\widetilde f(k)\e^{-\kappa k^2t}.
\]
反变换中乘积对应卷积，而
\[
 \frac1{2\pi}\int_{-\infty}^{\infty}
 \e^{-\kappa k^2t+\ii kx}\dd k
 =\frac1{\sqrt{4\pi\kappa t}}\e^{-x^2/(4\kappa t)}
 \equiv K_t(x).
\]
于是 $u=K_t*f$，即题给公式。核满足
$\int K_t\dd x=1$，所以总量 $\int u\dd x$ 守恒；$\kappa t$ 的量纲是长度平方，指数无量纲。$t\downarrow0$ 时 $K_t\to\delta$，故恢复初值。直接计算也有 $\partial_tK_t=\kappa\partial_x^2K_t$，方程校验完成。
\end{advsolution}

\subsection{球坐标分离与球谐函数}

对球坐标中的 Laplace 方程令 $\phi=R(r)Y(\theta,\varphi)$，得到
\[
 \frac1R\frac{\dd}{\dd r}\left(r^2\frac{\dd R}{\dd r}\right)
 =-\frac1Y\left[
 \frac1{\sin\theta}\partial_\theta(\sin\theta\,\partial_\theta Y)
 +\frac1{\sin^2\theta}\partial_\varphi^2Y\right]
 =l(l+1).
\]
径向解为
\[
 R(r)=A_lr^l+B_lr^{-l-1}.
\]
若轴对称，角方程在 $x=\cos\theta$ 后成为 Legendre 方程，要求两极有限迫使 $l$ 为非负整数，故
\[
 \phi(r,\theta)=\sum_{l=0}^\infty
 (A_lr^l+B_lr^{-l-1})P_l(\cos\theta).
\]
若还有 $\varphi$ 依赖，再令 $Y=y(x)z(\varphi)$。周期性给出
$z=\e^{\ii m\varphi}$、$m\in\mathbb Z$；$y$ 满足连带 Legendre 方程
\[
 \frac{\dd}{\dd x}\left[(1-x^2)y'\right]
 +\left[l(l+1)-\frac{m^2}{1-x^2}\right]y=0.
\]
其有限解可写为
\[
 P_l^m(x)=(-1)^m(1-x^2)^{m/2}\frac{\dd^mP_l}{\dd x^m},
 \qquad |m|\le l.
\]
这里采用含 Condon--Shortley 相位 $(-1)^m$ 的现代约定；若去掉该相位，下面负阶公式中的相应符号也要同步改变。

\begin{advexamplebox}{连带 Legendre 函数的负阶与宇称关系}\label{advA3:ex:associated-legendre}
证明
\[
 P_l^{-m}(x)=(-1)^m\frac{(l-m)!}{(l+m)!}P_l^m(x),
 \qquad P_l^m(-x)=(-1)^{l+m}P_l^m(x).
\]
\end{advexamplebox}
\begin{advsolution}
第二式最直接。由 Rodrigues 公式 $P_l(-x)=(-1)^lP_l(x)$，对 $x$ 求 $m$ 阶导数时每次链式法则再带一个负号，故
\[
 \frac{\dd^mP_l(-x)}{\dd(-x)^m}=(-1)^{l+m}
 \frac{\dd^mP_l(x)}{\dd x^m}.
\]
乘上偶函数 $(1-x^2)^{m/2}$，得到
$P_l^m(-x)=(-1)^{l+m}P_l^m(x)$。

对负阶，连带 Legendre 方程只含 $m^2$，所以 $P_l^{-m}$ 与 $P_l^m$ 是同一端点正则解的常数倍。用 Rodrigues 公式比较 $x\to1^-$ 的首项，或比较球谐函数归一化
\[
 Y_l^m=N_{lm}P_l^m(\cos\theta)\e^{\ii m\varphi},\qquad
 Y_l^{-m}=(-1)^m(Y_l^m)^*,
\]
可得比例常数
\[
 \boxed{P_l^{-m}=(-1)^m\frac{(l-m)!}{(l+m)!}P_l^m}.
\]
例如 $l=1,m=1$ 时 $P_1^1=-\sqrt{1-x^2}$、
$P_1^{-1}=\tfrac12\sqrt{1-x^2}$，正好满足公式，符号约定检查通过。
\end{advsolution}


\clearpage\endgroup

\begingroup
\newif\ifAdvHideHeading
\RenewDocumentEnvironment{advexamplebox}{m}{%
  \refstepcounter{advexample}%
  \begin{tcolorbox}[
    enhanced,unbreakable,colback=white,colframe=RuleGray,boxrule=0.5pt,arc=0mm,
    title={例~\theadvexample\quad #1},fonttitle=\bfseries,
    coltitle=Ink,colbacktitle=PaleBlue,borderline west={1.4pt}{0pt}{DeepBlue}]
}{\end{tcolorbox}}

\RenewDocumentEnvironment{advsolution}{+b}{%
  \par\smallskip\begingroup\small\noindent
  \ifAdvHideHeading\global\AdvHideHeadingfalse\else\LectureSolutionHeading\fi #1\par\endgroup\smallskip
}{}

\RenewDocumentEnvironment{advdefinition}{m +b}{%
  \refstepcounter{advdefinition}%
  \begin{tcolorbox}[
    enhanced,unbreakable,colback=white,colframe=RuleGray,boxrule=0.5pt,arc=0mm,
    title={定义~\theadvdefinition\quad #1},fonttitle=\bfseries,
    coltitle=Ink,colbacktitle=PaleGold,borderline west={1.2pt}{0pt}{Gold}]
  #2
  \end{tcolorbox}
}{}

\RenewDocumentEnvironment{advtheorem}{m +b}{%
  \refstepcounter{advtheorem}%
  \begin{tcolorbox}[
    enhanced,unbreakable,colback=white,colframe=RuleGray,boxrule=0.5pt,arc=0mm,
    title={定理~\theadvtheorem\quad #1},fonttitle=\bfseries,
    coltitle=Ink,colbacktitle=PaleBlue,borderline west={1.2pt}{0pt}{DeepBlue}]
  #2
  \end{tcolorbox}
}{}


\AdvancedUnit{4}{电磁学}
\section{局域规范对称性与电磁场}

电磁场可以从“带电物质场的局域相位对称性”出发统一引入。考虑复 Klein--Gordon 标量场 $\phi(\vect r,t)$，把 $\phi$ 与 $\bar\phi$ 在变分时视为独立变量，取
\[
 \mathcal L_m=\frac{\hbar^2}{m}
 \left(\frac1{c^2}\dot{\bar\phi}\dot\phi-
 \nabla\bar\phi\cdot\nabla\phi\right)-mc^2\bar\phi\phi.
\]
Euler--Lagrange 方程为
\[
 \frac1{c^2}\partial_t^2\phi-\nabla^2\phi
 =-\frac{m^2c^2}{\hbar^2}\phi,
\]
以及其复共轭。这里取 $|\phi|^2$ 具有数密度量纲；将方程除以动能前因子 $\hbar^2/m$ 可直接核对质量项为 $m^2c^2/\hbar^2$。代入平面波
$\phi=A\e^{-\ii\omega t+\ii\vect k\cdot\vect r}$，得到
\[
 (\hbar\omega)^2-(\hbar k)^2c^2=m^2c^4,
 \]
即 $E=\hbar\omega$、$\vect p=\hbar\vect k$ 与相对论能量--动量关系。

\begin{advexamplebox}{Klein--Gordon 平面波的能量与动量密度}\label{advA4:ex:kg-energy-momentum}
认为 $\bar\phi=\phi^*$。求平面波解的能量密度和动量密度，并检查它们与单粒子 $E,\vect p$ 的比例关系。
\end{advexamplebox}
\begin{advsolution}
共轭正则动量为
\[
 \pi_\phi=\frac{\partial\mathcal L}{\partial\dot\phi}
 =\frac{\hbar^2}{mc^2}\dot{\bar\phi},
 \qquad
 \pi_{\bar\phi}=\frac{\hbar^2}{mc^2}\dot\phi.
\]
因此 Hamilton 密度
\[
 \mathcal H=\pi_\phi\dot\phi+\pi_{\bar\phi}\dot{\bar\phi}-\mathcal L
 =\frac{\hbar^2}{mc^2}|\dot\phi|^2
 +\frac{\hbar^2}{m}|\nabla\phi|^2+mc^2|\phi|^2.
\]
对平面波，利用色散关系
$\hbar^2\omega^2=\hbar^2c^2k^2+m^2c^4$，得到
\[
 \boxed{u=\mathcal H=\frac{2\hbar^2\omega^2}{mc^2}|A|^2}.
\]
平移对称性的 Noether 动量密度（取物理空间协变分量的负号）为
\[
 \vect g=-\pi_\phi\nabla\phi-\pi_{\bar\phi}\nabla\bar\phi
 =\boxed{\frac{2\hbar^2\omega}{mc^2}\vect k\,|A|^2}.
\]
所以
\[
 \frac{\vect g}{u}=\frac{\vect k}{\omega}
 =\frac{\vect p}{E}.
\]
这说明 $|A|^2$ 决定物质量的密度因子，而每份物质的能量、动量之比正确。静止极限 $k=0$ 给出零动量密度；超相对论极限 $\hbar kc\gg mc^2$ 给出 $|\vect g|/u\to1/c$。本节场归一化含 $1/m$，不把固定振幅时 $m=0$ 代入密度公式。
\end{advsolution}

\subsection{从全局相位到局域相位}

物质拉格朗日量在全局变换
\[
 \phi\mapsto\phi\e^{\ii\lambda},\qquad
 \bar\phi\mapsto\bar\phi\e^{-\ii\lambda}
\]
下不变。Noether 流可写为
\[
 \varrho=\frac{\ii\hbar}{mc^2}
 (\bar\phi\dot\phi-\phi\dot{\bar\phi}),
 \qquad
 \vect j_N=-\frac{\ii\hbar}{m}
 (\bar\phi\nabla\phi-\phi\nabla\bar\phi),
\]
并满足 $\partial_t\varrho+\nabla\cdot\vect j_N=0$。

若相位参数变成时空函数，
\[
 \phi'=\phi\e^{\ii e\theta(x)},
\]
这里 $e$ 为带符号的电荷，$e\theta$ 无量纲。取 $x^\mu=(ct,\vect r)$、度规 $(-,+,+,+)$、$\partial_\mu=(c^{-1}\partial_t,\nabla)$，并取 $A_\mu=(-\varphi/c,\vect A)$。普通导数会多出相位导数；规定
\[
 A'_\mu=A_\mu+\hbar\partial_\mu\theta,
\]
并以协变导数
\[
 D_\mu\phi=\left(\partial_\mu-\frac{\ii e}{\hbar}A_\mu\right)\phi,
 \qquad
 D_\mu\bar\phi=\left(\partial_\mu+\frac{\ii e}{\hbar}A_\mu\right)\bar\phi
\]
替代普通导数，则 $D_\mu\phi$ 与 $\phi$ 具有相同的相位变换律。

规范场强
\[
 F_{\mu\nu}=\partial_\mu A_\nu-\partial_\nu A_\mu
\]
在 $A_\mu\mapsto A_\mu+\hbar\partial_\mu\theta$ 下保持不变。最简单的 Lorentz 标量自能为
\[
 \mathcal L_e=-\frac1{4\mu_0}F_{\mu\nu}F^{\mu\nu}.
\]
上述时间分量含 $1/c$，因此 $F_{0i}=-E_i/c$，而空间分量给出磁场。定义
\[
 \boxed{\vect E=-\partial_t\vect A-\nabla\varphi},
 \qquad
 \boxed{\vect B=\nabla\times\vect A}.
\]
对于给定的守恒外源 $\rho,\vect j$，势与外源的作用量可写成
\[
 \mathcal L=\frac{\varepsilon_0}{2}E^2-\frac1{2\mu_0}B^2
 -\rho\varphi+\vect j\cdot\vect A.
\]
由势的定义恒有
\[
 \nabla\cdot\vect B=0,
 \qquad \nabla\times\vect E+\partial_t\vect B=0.
\]
对 $\varphi,\vect A$ 作场变分得到另两式
\[
 \nabla\cdot\vect E=\frac\rho{\varepsilon_0},
 \qquad
 \nabla\times\vect B-\frac1{c^2}\partial_t\vect E=\mu_0\vect j.
\]
这四式构成有源真空 Maxwell 方程。若源本身就是上述带电标量场，应改用完整的 $\mathcal L_m(D_\mu\phi)+\mathcal L_e$ 并由变分求出电流；协变导数平方中的 $A^2|\phi|^2$ 项也须保留，不能把删去它的弱场表达冒充精确规范不变理论。

在三维记号下，$\vect A'=\vect A+\hbar\nabla\theta$、$\varphi'=\varphi-\hbar\partial_t\theta$。对点电荷的粒子拉格朗日量
\[
 L_p=-q\varphi(\vect r)+q\dot{\vect r}\cdot\vect A(\vect r)+T(\dot{\vect r})
\]
应用 Euler--Lagrange 方程，得到
\[
 \boxed{\dot{\vect p}=q(\vect E+\vect v\times\vect B)}.
\]
规范变换只让 $L_p$ 增加全时间导数 $q\hbar\dot\theta$，不改变运动方程；这同时说明局域相位对称性与电荷耦合的一致性。

\section{静电场}

\subsection{分布、电势与能量}

静电场满足
\[
 \nabla\times\vect E=0,
 \qquad \nabla\cdot\vect E=\frac\rho{\varepsilon_0}.
\]
电荷分布对区域 $\Omega$ 给出 $Q(\Omega)=\int_\Omega\rho\dd V$；更一般地，它对测试函数 $f$ 给出线性泛函
\[
 \varrho(f)=\int_{\mathbb R^3}\rho(\vect r)f(\vect r)\dd V.
\]
点值泛函 $f\mapsto f(0)$ 和导数泛函 $f\mapsto f'(0)$ 都是分布。这样的语言可以无歧义地处理点电荷和多极子。

\begin{advdefinition}{点电荷与 Dirac $\delta$ 分布}\label{advA4:def:point-charge}
位于 $\vect r_0$、电荷量为 $q$ 的点电荷定义为
\[
 \rho(\vect r)=q\delta^{(3)}(\vect r-\vect r_0),
 \qquad
 \int q\delta^{(3)}(\vect r-\vect r_0)f(\vect r)\dd V=qf(\vect r_0).
\]
一维中 $\delta$ 不是普通函数，而由
$\int\delta(x-x_0)f(x)\dd x=f(x_0)$ 定义；其导数满足
$\int\delta'(x)f(x)\dd x=-f'(0)$。
\end{advdefinition}

Lorentz 核
\[
 L_\epsilon(x)=\frac1\pi\frac\epsilon{x^2+\epsilon^2}
\]
在 $\epsilon\downarrow0$ 时收敛到 $\delta(x)$：它非负、积分为1，且任意固定远离零点的区域上质量趋于零。分部积分还给出 $L_\epsilon'\to\delta'$。

\begin{advexamplebox}{Dirac $\delta$ 的四种逼近表示}\label{advA4:ex:delta-representations}
证明分布意义下
\[
 \delta(x)=\lim_{\epsilon\downarrow0}\frac1{\sqrt{2\pi}\epsilon}
 \e^{-x^2/(2\epsilon^2)}
 =\lim_{K\to\infty}\frac{\sin Kx}{\pi x},
\]
并说明周期 Fourier 核和连续 Fourier 积分如何给出
\[
 \delta(x)=\frac1{2\pi}\int_{-\infty}^{\infty}\e^{\ii kx}\dd k.
\]
\end{advexamplebox}
\begin{advsolution}
对 Gaussian 核令 $x=\epsilon u$，则对任意连续快速衰减的 $f$，
\[
 \int\frac{\e^{-x^2/(2\epsilon^2)}}{\sqrt{2\pi}\epsilon}f(x)\dd x
 =\int\frac{\e^{-u^2/2}}{\sqrt{2\pi}}f(\epsilon u)\dd u
 \longrightarrow f(0).
\]
对 sinc 核，用 Fourier 反演：
\[
 \int\frac{\sin Kx}{\pi x}f(x)\dd x
 =\frac1{2\pi}\int_{-K}^{K}\!\int f(x)\e^{\ii kx}\dd x\dd k
 \longrightarrow f(0).
\]
长度为 $2\pi$ 的周期区间上，Dirichlet 核
$D_N(x)=\sum_{n=-N}^{N}\e^{\ii nx}$ 满足
$(2\pi)^{-1}\int_{-\pi}^{\pi}D_N(x)f(x)\dd x\to f(0)$（在满足 Dirichlet 条件时取左右极限平均）。令周期趋于无穷，离散波数和变成积分，即得连续 Fourier 表示。四种表达都只在分布意义成立，不能把 $\delta(0)=\infty$ 当普通函数值参与代数运算。
\end{advsolution}

\subsection{Cartesian 多极展开}

当场点满足 $r$ 大于电荷分布尺度时，平移算符给出
\[
 \frac1{|\vect r-\vect r'|}
 =\sum_{n=0}^\infty\frac{(-1)^n}{n!}
 x'_{i_1}\cdots x'_{i_n}
 \partial_{i_1}\cdots\partial_{i_n}\frac1r.
\]
定义位置张量与导数张量
\[
 R_n(\vect r)=\underbrace{\vect r\cdots\vect r}_{n},
 \qquad
 T_n(\vect r)=(-1)^nr^{2n+1}
 \underbrace{\nabla\cdots\nabla}_{n}\frac1r.
\]
递推式
\[
 T_{n+1}=(2n+1)\vect r,T_n-r^2\nabla T_n
\]
给出
\[
 T_0=1,\qquad T_1=\vect r,\qquad
 T_2=3\vect r\vect r-r^2\mathsf1,
\]
\[
 (T_3)_{abc}=15x_ax_bx_c
 -3r^2(x_a\delta_{bc}+x_b\delta_{ca}+x_c\delta_{ab}).
\]
由于 $\nabla^2(1/r)=0$（$r\ne0$），$T_n$ 对任意两指标缩并均为零，是完全对称无迹张量。定义普通与无迹多极矩
\[
 Q_n=\int\rho(\vect r')R_n(\vect r')\dd V',
 \qquad
 D_n=\int\rho(\vect r')T_n(\vect r')\dd V'.
\]
电势可写成
\[
 \boxed{\varphi(\vect r)=\frac1{4\pi\varepsilon_0}
 \sum_{n=0}^\infty\frac{(D_n,R_n(\vect r))}{n!r^{2n+1}}}.
\]
前几阶为总电荷 $D_0=Q$、偶极矩
$\vect p=D_1=\int\rho\vect r\dd V$ 和四极矩
\[
 \mathsf D=D_2=\int\rho(3\vect r\vect r-r^2\mathsf1)\dd V.
\]
相应的偶极与四极电势为
\[
 \varphi_1=\frac{\vect p\cdot\vect r}{4\pi\varepsilon_0r^3},
 \qquad
 \varphi_2=\frac{\vect r\cdot\mathsf D\cdot\vect r}
 {8\pi\varepsilon_0r^5}.
\]
偶极可由相反点电荷作无穷小分离得到，分布写成
$\rho=-\vect p\cdot\nabla\delta^{(3)}(\vect r)$；四极则由两个相反偶极再分离得到，写成
$\rho=\tfrac12Q_{ij}\partial_i\partial_j\delta^{(3)}$。在光滑外势中的相互作用能由 Taylor 展开直接给出
\[
 U=Q\varphi_{\rm ext}(0)+\vect p\cdot\nabla\varphi_{\rm ext}(0)
 +\frac1{2!}Q_{ij}\partial_i\partial_j\varphi_{\rm ext}(0)+\cdots.
\]

\begin{advexamplebox}{对称无迹张量的独立分量数}\label{advA4:ex:stf-dimension}
证明三维 $n$ 阶完全对称无迹张量有 $2n+1$ 个独立分量。
\end{advexamplebox}
\begin{advsolution}
$n$ 阶完全对称张量只由指标1、2、3各出现多少次决定，因此
\[
 \dim S^n(\mathbb R^3)=\binom{n+2}{2}.
\]
取迹把 $S^n$ 映到 $S^{n-2}$。这个映射是满射：给定任意
$B\in S^{n-2}$，先取对称化的 $\delta\odot B$，再减去有限个更高重迹项即可构造迹恰为 $B$ 的张量。因此无迹子空间是迹映射的核，秩--零化度定理给出
\[
 \dim\operatorname{STF}^n
 =\binom{n+2}{2}-\binom n2
 =\boxed{2n+1}.
\]
$n=1$ 得3个偶极分量，$n=2$ 得5个四极分量，也与“对称矩阵6个分量减去一个迹条件”一致。
\end{advsolution}

\begin{advtheorem}{点电偶极子的完整分布场}\label{advA4:thm:dipole-distribution-field}
点偶极子在全空间的电场是
\[
 \boxed{\vect E(\vect r)=\frac1{4\pi\varepsilon_0}
 \frac{3(\vect p\cdot\vect r)\vect r-r^2\vect p}{r^5}
 -\frac{\vect p}{3\varepsilon_0}\delta^{(3)}(\vect r)}.
\]
\end{advtheorem}
\noindent\textbf{证明}\quad 分布恒等式
\[
 \partial_i\partial_j\frac1r
 =\frac{3x_ix_j-r^2\delta_{ij}}{r^5}
 -\frac{4\pi}{3}\delta_{ij}\delta^{(3)}(\vect r)
\]
可通过挖去半径 $\epsilon$ 小球后分部积分两次证明；小球面项在
$\epsilon\downarrow0$ 时留下各向同性张量
$-(4\pi/3)\delta_{ij}f(0)$。偶极势
$\varphi=-(4\pi\varepsilon_0)^{-1}p_j\partial_j(1/r)$，故
$E_i=(4\pi\varepsilon_0)^{-1}p_j\partial_i\partial_j(1/r)$，代入即得。$\delta$ 项保证分布意义下
$\nabla\cdot\vect E=-(\vect p\cdot\nabla\delta)/\varepsilon_0$。\hfill$\square$

\begin{advexamplebox}{电偶极子的合力与力矩}\label{advA4:ex:dipole-force-torque}
证明在尺度上缓变的外电场中
\[
 \vect F=(\vect p\cdot\nabla)\vect E,
 \qquad \vect M=\vect p\times\vect E.
\]
\end{advexamplebox}
\begin{advsolution}
把偶极子看成位于 $\vect r\pm\vect a/2$ 的 $\pm q$，其中 $\vect p=q\vect a$。合力
\[
 \vect F=q\vect E(\vect r+\vect a/2)-q\vect E(\vect r-\vect a/2)
 =(\vect p\cdot\nabla)\vect E+O(a^2).
\]
两力对中心的力矩为
\[
 \vect M=\frac{\vect a}{2}\times q\vect E
 +\left(-\frac{\vect a}{2}\right)\times[-q\vect E]
 =\vect p\times\vect E.
\]
也可从势能 $U=-\vect p\cdot\vect E$ 得到
$\vect F=-\nabla U$；静电场无旋时
$\nabla(\vect p\cdot\vect E)=(\vect p\cdot\nabla)\vect E$。力矩使 $\vect p$ 转向 $\vect E$，稳定方向为同向。
\end{advsolution}

\begin{advexamplebox}{任意取向偶极子与孤立中性导体球}\label{advA4:ex:dipole-grounded-sphere}
半径为 $R$、总电荷为零的孤立导体球外，距球心 $a>R$ 处放置偶极矩 $\vect p$，它与连心线夹角为 $\theta$。令 $\vect p=p_\parallel\vect e_r+\vect p_\perp$。求像系、势能、力和力矩，并与接地球比较。
\end{advexamplebox}
\begin{advsolution}
先构造接地球的像系：令 $d=a^2-R^2$，像点位于 $b=R^2/a$，其上有
\[
 q'=\frac{Rp_\parallel}{a^2},\qquad
 \vect p'=\left(\frac Ra\right)^3
 (p_\parallel\vect e_r-\vect p_\perp).
\]
偶极像的总电荷为零，但像点电荷 $q'$ 一般非零。
原题要求球不带电，故还须在球心增加
\[
 q_0=-q'=-\frac{Rp_\parallel}{a^2}.
\]
中心像只在球面上增加常势，不破坏等势边界，同时使感应总电荷为零。
接地与孤立中性是两种边界条件，不可互换。

像场在真实偶极子位置为
\[
 \vect E_{\rm im}=\frac1{4\pi\varepsilon_0}
 (A p_\parallel\vect e_r+B\vect p_\perp),
\]
其中
\[
 A=\frac R{d^2}+\frac{2R^3}{d^3}-\frac R{a^4},
 \qquad B=\frac{R^3}{d^3}.
\]
因感应电荷随源线性建立，势能为
\[
 \boxed{U=-\frac12\vect p\cdot\vect E_{\rm im}
 =-\frac{A p_\parallel^2+B p_\perp^2}{8\pi\varepsilon_0}}.
\]
相对接地球，中心补偿像使 $U$ 增加 $Rp_\parallel^2/(8\pi\varepsilon_0a^4)$。
固定偶极矩大小和空间方向求 $-\nabla U$，径向及横向分量分别为
\[
 F_r=-\frac1{8\pi\varepsilon_0}
 \left[\frac{4Ra p_\parallel^2}{d^3}
 +\frac{6R^3a(2p_\parallel^2+p_\perp^2)}{d^4}
 -\frac{4R p_\parallel^2}{a^5}\right],
\]
\[
 \vect F_\perp=\frac{(A-B)p_\parallel\vect p_\perp}
 {4\pi\varepsilon_0a},\qquad
 \boxed{\vect F=F_r\vect e_r+\vect F_\perp}.
\]
任意倾斜偶极子还受到非径向力，不能只交出 $F_r$。
力矩可直接由 $\vect M=\vect p\times\vect E_{\rm im}$ 求出。
若以增加 $\theta$ 为正，其分量为
\[
 \boxed{M_\theta=-\frac{p^2\sin2\theta}{8\pi\varepsilon_0}
 \left[\frac R{d^2}+\frac{R^3}{d^3}-\frac R{a^4}\right]}.
\]
因为方括号为正，径向的 $\theta=0,\pi$ 是稳定取向。
若题目改为接地球，则删去 $q_0$，即所有式中删去来自 $R/a^4$ 的项。
令 $a=R+h$、$R\to\infty$，两种边界的局部极限都为
\[
 U\longrightarrow-\frac{2p_n^2+p_t^2}{64\pi\varepsilon_0h^3},
\]
其中 $p_n,p_t$ 是相对于平面的法向、切向分量。
\end{advsolution}

\subsection{电介质}

极化强度 $\vect P$ 是单位体积的电偶极矩。若每单位体积有 $n$ 个相同偶极子，$\vect P=n\vect p$。

\begin{advtheorem}{极化电荷}\label{advA4:thm:polarization-charge}
分段光滑极化分布产生的电场等效于体束缚电荷与界面束缚电荷
\[
 \boxed{\rho_b=-\nabla\cdot\vect P},
 \qquad
 \boxed{\sigma_b=(\vect P_1-\vect P_2)\cdot\vect n},
\]
其中 $\vect n$ 从介质1指向介质2。介质外 $\vect P_2=0$ 时
$\sigma_b=\vect P\cdot\vect n$。
\end{advtheorem}
\noindent\textbf{证明}\quad 偶极势叠加为
\[
 \varphi=\frac1{4\pi\varepsilon_0}
 \int_V\vect P(\vect r')\cdot
 \frac{\vect r-\vect r'}{|\vect r-\vect r'|^3}\dd V'
 =\frac1{4\pi\varepsilon_0}\int_V
 \vect P\cdot\nabla'\frac1{|\vect r-\vect r'|}\dd V'.
\]
对 $\vect r'$ 分部积分，体项为
$-\int(\nabla'\cdot\vect P)/R$，边界项为
$\oint(\vect P\cdot\vect n)/R$。界面两侧法向相反，合并即得跳跃式。\hfill$\square$

定义
\[
 \vect D=\varepsilon_0\vect E+\vect P,
 \qquad \nabla\cdot\vect D=\rho_f.
\]
线性各向同性介质中
$\vect P=\varepsilon_0\chi_e\vect E$，故
$\vect D=\varepsilon\vect E$、
$\varepsilon=\varepsilon_0(1+\chi_e)$，场能密度为
$w=\tfrac12\vect D\cdot\vect E$。

\begin{advexamplebox}{位移极化的极化率}\label{advA4:ex:displacement-polarization}
把分子视为电荷 $\pm Q$ 的一维谐振子，相对位移势能为 $kx^2/2$，数密度为 $n$。求自由能系数 $\alpha$ 与电极化率 $\chi_e$。
\end{advexamplebox}
\begin{advsolution}
平衡时 $kx=QE$，单分子偶极矩
$p=Qx=Q^2E/k$。所以
\[
 P=np=\frac{nQ^2}{k}E.
\]
若自由能写成 $f=f_0+\alpha P^2/2-\vect P\cdot\vect E$，极值条件为 $\alpha P=E$，因此
\[
 \boxed{\alpha=\frac{k}{nQ^2}},
 \qquad
 \boxed{\chi_e=\frac{P}{\varepsilon_0E}
 =\frac{nQ^2}{\varepsilon_0k}}.
\]
$k$ 越大越难极化，$n,Q$ 越大越易极化，趋势正确。三维各向同性束缚若三个方向劲度相同，结果逐方向相同。
\end{advsolution}

\begin{advexamplebox}{取向极化与 Langevin 函数}\label{advA4:ex:orientation-polarization}
数密度为 $n$ 的永久偶极矩 $p$ 在温度 $T$ 下达到热平衡。求极化强度的精确式和弱场极化率。
\end{advexamplebox}
\begin{advsolution}
取 $\vect E$ 为极轴，单偶极能量 $U=-pE\cos\theta$。令
$x=pE/(k_BT)$，配分函数的角部分为
\[
 Z_1=2\pi\int_0^\pi\e^{x\cos\theta}\sin\theta\dd\theta
 =4\pi\frac{\sinh x}{x}.
\]
平均投影
\[
 \langle p\cos\theta\rangle
 =p\partial_x\ln Z_1
 =p\left(\coth x-\frac1x\right)
 =p\mathcal L(x).
\]
因此
\[
 \boxed{P=np\mathcal L\!\left(\frac{pE}{k_BT}\right)}.
\]
弱场 $\mathcal L(x)=x/3+O(x^3)$，故
\[
 P=\frac{np^2}{3k_BT}E,\qquad
 \boxed{\alpha=\frac{3k_BT}{np^2}},
 \quad
 \boxed{\chi_e=\frac{np^2}{3\varepsilon_0k_BT}}.
\]
强场 $x\to\infty$ 时 $P\to np$，正确体现饱和而不是无限线性增长。
\end{advsolution}

\begin{advtheorem}{Helmholtz 分解}\label{advA4:thm:helmholtz}
设 $\vect F$ 在区域 $V$ 内足够光滑。定义
\[
 \Phi(\vect r)=\frac1{4\pi}\int_V
 \frac{\nabla'\cdot\vect F(\vect r')}{|\vect r-\vect r'|}\dd V'
 -\frac1{4\pi}\oint_{\partial V}
 \frac{\vect F(\vect r')\cdot\dd\vect S'}{|\vect r-\vect r'|},
\]
\[
 \vect A(\vect r)=\frac1{4\pi}\int_V
 \frac{\nabla'\times\vect F(\vect r')}{|\vect r-\vect r'|}\dd V'
 -\frac1{4\pi}\oint_{\partial V}
 \frac{\dd\vect S'\times\vect F(\vect r')}{|\vect r-\vect r'|}.
\]
则
\[
 \boxed{\vect F=-\nabla\Phi+\nabla\times\vect A}.
\]
在全空间且场充分衰减时边界项消失；无旋场可写成梯度，无散场可写成旋度。
\end{advtheorem}
\noindent\textbf{证明}\quad 对 $G=1/(4\pi|\vect r-\vect r'|)$ 使用
$\nabla^2G=-\delta^{(3)}(\vect r-\vect r')$ 和矢量恒等式
\[
 \nabla\times\nabla\times(G\vect F')
 =\nabla[\nabla\cdot(G\vect F')]-\nabla^2(G\vect F').
\]
对 $\vect r'$ 分部积分，把导数从 $\vect F'$ 移到 $G$，体积分的边界项恰好产生定理中的两个曲面积分；其余体积分合并为
$-\int_V(\nabla^2G)\vect F'\dd V'=\vect F(\vect r)$。这也说明若不规定边界行为，势并不唯一。\hfill$\square$

全空间静电势与电场因而为
\[
 \varphi(\vect r)=\frac1{4\pi\varepsilon_0}
 \int\frac{\rho(\vect r')}{|\vect r-\vect r'|}\dd V',
\]
\[
 \vect E(\vect r)=\frac1{4\pi\varepsilon_0}
 \int\frac{\rho(\vect r')(\vect r-\vect r')}{|\vect r-\vect r'|^3}\dd V'.
\]
环路积分 $\oint\vect E\cdot\dd\vect l=0$，通量满足
$\oint\vect E\cdot\dd\vect S=Q_{\rm in}/\varepsilon_0$。点电荷穿过任意开曲面的通量为
$q\Omega/(4\pi\varepsilon_0)$，其中
$\Omega=\int_S\dd\vect S\cdot\vect e_r/r^2$ 是有向立体角。

\begin{advexamplebox}{圆锥立体角与球坐标立体角元}\label{advA4:ex:solid-angle}
半顶角为 $\alpha$ 的圆锥从顶点截取单位球面。求其立体角、$\alpha$ 改变时的立体角元，并推广到球坐标一般面元。
\end{advexamplebox}
\begin{advsolution}
单位球面上圆锥对应 $0\le\theta\le\alpha$，故
\[
 \boxed{\Omega=\int_0^{2\pi}\!\int_0^\alpha
 \sin\theta\dd\theta\dd\phi=2\pi(1-\cos\alpha)}.
\]
微分得
$\dd\Omega=2\pi\sin\alpha\dd\alpha$。一般球坐标面元在半径 $r$ 的球面上为
\[
 \dd S=r^2\sin\theta\dd\theta\dd\phi,
 \qquad
 \boxed{\dd\Omega=\frac{\dd S}{r^2}
 =\sin\theta\dd\theta\dd\phi}.
\]
若开曲面法向与视线夹角为 $\beta$，则
$\dd\Omega=(\cos\beta/r^2)\dd S$。$\alpha=\pi$ 时得到 $4\pi$，全空间极限正确。
\end{advsolution}

\begin{advexamplebox}{匀强场与共轴电偶极子的电场线}\label{advA4:ex:dipole-uniform-field-lines}
匀强电场 $E_0\vect e_z$ 中放置位于原点、偶极矩平行或反平行于 $\vect e_z$ 的点偶极子。求轴对称电场线方程。
\end{advexamplebox}
\begin{advsolution}
令 $s=+1$ 表示同向，$s=-1$ 表示反向，$p>0$。总电势为
\[
 \varphi=-E_0r\cos\theta+s\frac{p\cos\theta}{4\pi\varepsilon_0r^2}.
\]
轴对称无旋场可用通量函数 $\Psi$ 表示
\[
 E_r=\frac1{r^2\sin\theta}\partial_\theta\Psi,
 \qquad
 E_\theta=-\frac1{r\sin\theta}\partial_r\Psi.
\]
对匀强场和偶极场分别积分，得到
\[
 \Psi_{\rm uni}=\frac12E_0r^2\sin^2\theta,
 \qquad
 \Psi_{\rm dip}=s\frac{p}{4\pi\varepsilon_0r}\sin^2\theta.
\]
电场线沿线 $\Psi$ 不变，因此
\[
 \boxed{\sin^2\theta\left(
 \frac12E_0r^2+s\frac{p}{4\pi\varepsilon_0r}
 \right)=C}.
\]
同向时两条轴上偶极场都与匀强场同向，因而轴上无零点；赤道面有半径 $r=(p/4\pi\varepsilon_0E_0)^{1/3}$ 的零场环。反向时零点在正、负两条轴上，均满足 $r=(2p/4\pi\varepsilon_0E_0)^{1/3}$。把 $p$ 置零退回平行直线，把 $E_0$ 置零得到偶极场线 $r\propto\sin^2\theta$。
\end{advsolution}

静电能可写成
\[
 U=\frac12\int\rho\varphi\dd V.
\]
若 $\rho=\rho_1+\rho_2$、$\varphi=\varphi_1+\varphi_2$，交叉项是两个体系的相互作用能。

\begin{advtheorem}{Green 互易关系}\label{advA4:thm:green-reciprocity}
对在同一边界条件下的两组静电解，若边界项为零，则
\[
 \boxed{\int_V\rho_1\varphi_2\dd V=
 \int_V\rho_2\varphi_1\dd V}.
\]
\end{advtheorem}
\noindent\textbf{证明}\quad Green 第二恒等式给出
\[
 \int_V(\varphi_1\nabla^2\varphi_2-\varphi_2\nabla^2\varphi_1)\dd V
 =\oint_{\partial V}(\varphi_1\partial_n\varphi_2-\varphi_2\partial_n\varphi_1)\dd S.
\]
共同的齐次 Dirichlet 或 Neumann 边界使右边为零；代入
$\nabla^2\varphi_i=-\rho_i/\varepsilon_0$ 即得。\hfill$\square$

\begin{advtheorem}{静电场能}\label{advA4:thm:electrostatic-energy}
在场充分衰减或边界项消失时，
\[
 \boxed{\frac12\int\rho\varphi\dd V
 =\frac{\varepsilon_0}{2}\int E^2\dd V}.
\]
\end{advtheorem}
\noindent\textbf{证明}\quad 用 $\rho=-\varepsilon_0\nabla^2\varphi$ 分部积分：
\[
 \frac12\int\rho\varphi\dd V
 =-\frac{\varepsilon_0}{2}\oint\varphi\nabla\varphi\cdot\dd\vect S
 +\frac{\varepsilon_0}{2}\int|\nabla\varphi|^2\dd V.
\]
边界项消失且 $\vect E=-\nabla\varphi$，即得结论。场能非负，但点电荷自能发散；比较不同构型时应分离并扣除与位置无关的自能。

\begin{advexamplebox}{带电线段、旋转椭球与等势面}\label{advA4:ex:spheroid}
均匀带电线段位于 $z$ 轴 $-c<z'<c$。完成以下讨论：
\begin{enumerate}
 \item 证明等势面是长旋转椭球，电场线是正交的双曲线；
 \item 求长半轴 $A$、短半轴 $B$ 的孤立导体长旋转椭球电容；
 \item 写出其表面电荷密度；
 \item 给出扁旋转椭球的对应结果和等效像电荷分布；
 \item 在长轴上距中心 $l>A$ 处放点电荷 $Q$ 并把椭球接地，求椭球总感应电荷。
\end{enumerate}
\end{advexamplebox}
\begin{advsolution}
记到线段两端的距离为 $r_\pm$。线电荷密度为 $\lambda$ 时
\[
 \varphi=\frac\lambda{4\pi\varepsilon_0}
 \int_{-c}^{c}\frac{\dd z'}{\sqrt{\rho^2+(z-z')^2}}
 =\frac\lambda{4\pi\varepsilon_0}
 \ln\frac{r_++r_-+2c}{r_++r_--2c}.
\]
令长旋转椭球坐标
\[
 \xi=\frac{r_++r_-}{2c},\qquad
 \eta=\frac{r_+-r_-}{2c},
\]
则 $\varphi$ 只依赖 $\xi$；$\xi=$ 常数是焦点为 $\pm c$ 的长旋转椭球，$\eta=$ 常数是与它正交的双曲面，因此后者给出电场线。

导体表面 $\xi=\xi_0=A/c$，其中 $c=\sqrt{A^2-B^2}$。令
$Q_0(\xi)=\tfrac12\ln[(\xi+1)/(\xi-1)]$，外部势为
\[
 \varphi(\xi)=V\frac{Q_0(\xi)}{Q_0(\xi_0)}.
\]
远处 $Q_0(\xi)\sim1/\xi=c/r$，与 $Q/(4\pi\varepsilon_0r)$ 比较得到
\[
 \boxed{C_{\rm prolate}=\frac QV
 =\frac{4\pi\varepsilon_0c}{Q_0(A/c)}
 =\frac{8\pi\varepsilon_0c}
 {\ln[(A+c)/(A-c)]}}.
\]
用尺度因子
$h_\xi=c\sqrt{(\xi^2-\eta^2)/(\xi^2-1)}$，表面法向场为
\[
 E_n=-\frac1{h_\xi}\partial_\xi\varphi
 =\frac{V}{cQ_0(\xi_0)
 \sqrt{(\xi_0^2-1)(\xi_0^2-\eta^2)}}.
\]
故
\[
 \boxed{\sigma(\eta)=\varepsilon_0E_n}.
\]
它在尖端 $|\eta|=1$ 较大，符合曲率越大电荷越集中的物理图像。

扁旋转椭球取极半轴 $A$、赤道半轴 $B>A$，令
$c=\sqrt{B^2-A^2}$。把 $c$ 作解析延拓可得
\[
 \boxed{C_{\rm oblate}=
 \frac{4\pi\varepsilon_0c}{\arctan(c/A)}}.
\]
它可由焦盘 $s<c$ 上的等效电荷产生，盘面密度可取
\[
 \sigma_{\rm disk}(s)=\frac{Q}{2\pi c\sqrt{c^2-s^2}},
 \qquad 0<s<c,
\]
其积分为 $Q$，所产生的等势面正是共焦扁旋转椭球。$A\to B$ 时两种电容都趋于 $4\pi\varepsilon_0A$，球极限通过。

最后用互易定理。令状态1为外点电荷 $Q$ 与接地椭球，状态2为椭球保持单位电势且无外点电荷。状态2在轴上 $l$ 处的势为
\[
 \varphi_2(l)=\frac{Q_0(l/c)}{Q_0(A/c)}.
\]
互易关系给出 $Q\varphi_2(l)+Q_{\rm ind}\cdot1=0$，所以
\[
 \boxed{Q_{\rm ind}=-Q\frac{Q_0(l/c)}{Q_0(A/c)}
 =-Q\frac{\ln[(l+c)/(l-c)]}{\ln[(A+c)/(A-c)]}}.
\]
$l\to\infty$ 时感应总电荷趋零，$l\downarrow A$ 时趋于 $-Q$，极限合理。
\end{advsolution}

\subsection{静电边值问题与格林函数}

Poisson 方程与无源区的 Laplace 方程为
\[
 \nabla^2\varphi=-\frac\rho{\varepsilon_0},
 \qquad \nabla^2\varphi=0.
\]
Dirichlet 问题给定 $\varphi|_{\partial V}$，Neumann 问题给定
$\partial_n\varphi|_{\partial V}$；后者还须满足由 Gauss 定理得到的相容条件，并且解只确定到一个加法常数。

Green 函数满足
\[
 \nabla'^2G(\vect r,\vect r')=-4\pi\delta^{(3)}(\vect r-\vect r').
\]
Dirichlet Green 函数取 $G_D=0$ 于边界，解为
\[
 \varphi(\vect r)=\frac1{4\pi\varepsilon_0}\int_V
 \rho(\vect r')G_D(\vect r,\vect r')\dd V'
 -\frac1{4\pi}\oint_{\partial V}
 \varphi(\vect r')\partial_{n'}G_D\dd S'.
\]
Neumann Green 函数取常数法向导数以满足积分相容条件；选定势的加权平均后，解由体电荷项和
$G_N\partial_n\varphi$ 的边界项给出。

\begin{advtheorem}{静电 Green 函数的对称性}\label{advA4:thm:green-symmetry}
在相同的齐次边界条件和同一规范常数约定下，
\[
 \boxed{G_D(\vect r,\vect r')=G_D(\vect r',\vect r)},
 \qquad
 \boxed{G_N(\vect r,\vect r')=G_N(\vect r',\vect r)}.
\]
\end{advtheorem}
\noindent\textbf{证明}\quad 在 Green 第二恒等式中取
$u(\vect x)=G(\vect r,\vect x)$、
$v(\vect x)=G(\vect r',\vect x)$。体积分利用各自的 $\delta$ 源，等于
$-4\pi G(\vect r',\vect r)+4\pi G(\vect r,\vect r')$。Dirichlet 情形边界上两函数都为零；Neumann 情形两者法向导数相同为规定常数，边界项在固定平均值规范下也相消，所以体积分为零。\hfill$\square$

\begin{advexamplebox}{上半空间的 Dirichlet 与 Neumann Green 函数}\label{advA4:ex:half-space-green}
区域为 $z>0$，源点为 $\vect r_0=(x_0,y_0,z_0)$，镜像点为
$\vect r_0^*=(x_0,y_0,-z_0)$。求两类 Green 函数并解释像电荷法。
\end{advexamplebox}
\begin{advsolution}
令 $R_-=|\vect r-\vect r_0|$、$R_+=|\vect r-\vect r_0^*|$。在平面 $z=0$ 上 $R_-=R_+$，故
\[
 \boxed{G_D=\frac1{R_-}-\frac1{R_+}}
\]
满足 $G_D=0$。它等价于在镜像点放置反号像电荷，适用于接地平面。

对 Neumann 条件，平面上两个核的法向导数互为相反数，所以
\[
 \boxed{G_N=\frac1{R_-}+\frac1{R_+}}
\]
满足 $\partial_nG_N=0$，等价于同号像电荷。前者使势为奇延拓，后者使势为偶延拓。代回
$\nabla^2G=-4\pi\delta$ 时，镜像点在区域外，不引入额外真实源。Green 函数均关于 $\vect r,\vect r_0$ 对称。
\end{advsolution}

\begin{advexamplebox}{球外区域的两类 Green 函数}\label{advA4:ex:sphere-green}
对半径 $R$ 的球外区域 $r>R$，源点 $r_0>R$。求 Dirichlet Green 函数；给出零法向导数 Neumann Green 函数的球谐展开，并说明为什么后者不能由单个像点电荷完成。
\end{advexamplebox}
\begin{advsolution}
令
\[
 \vect r_0^*=\frac{R^2}{r_0^2}\vect r_0.
\]
在球面上有
$|\vect r-\vect r_0^*|=(R/r_0)|\vect r-\vect r_0|$，故
\[
 \boxed{G_D(\vect r,\vect r_0)=
 \frac1{|\vect r-\vect r_0|}
 -\frac{R/r_0}{|\vect r-\vect r_0^*|}}.
\]
第二项是位于球内的像电荷 $q^*=-(R/r_0)q$，因此不在求解区域内。

对 Neumann 条件，把轴取在 $\vect r_0$ 方向。自由核在球面附近展开为
\[
 \frac1{|\vect r-\vect r_0|}
 =\sum_{l=0}^\infty\frac{r^l}{r_0^{l+1}}P_l(\cos\gamma)
 \quad (r<r_0).
\]
加入在区域内无奇点、在无穷远衰减的项 $a_lr^{-l-1}P_l$，令
$\partial_rG_N|_{r=R}=0$，逐阶得到
\[
 a_l=\frac{l}{l+1}\frac{R^{2l+1}}{r_0^{l+1}},\qquad l\ge1.
\]
因此可取
\[
 \boxed{G_N=\frac1{|\vect r-\vect r_0|}
 +\sum_{l=1}^\infty\frac{l}{l+1}
 \frac{R^{2l+1}}{r^{l+1}r_0^{l+1}}P_l(\cos\gamma)}.
\]
单个像点的展开系数只能按固定几何比变化，不可能产生逐阶的 $l/(l+1)$，所以 Neumann 球问题通常需要连续像分布或无穷级数。
\end{advsolution}

\begin{advexamplebox}{两半球维持相反电势}\label{advA4:ex:split-sphere-potential}
半径 $R$ 的球面上，北半球维持 $+V$、南半球维持 $-V$，赤道绝缘缝宽忽略。求球内电势。
\end{advexamplebox}
\begin{advsolution}
边值只依赖 $\theta$，球内正则轴对称解为
\[
 \varphi(r,\theta)=\sum_{l=0}^\infty A_l
 \left(\frac rR\right)^lP_l(\cos\theta).
\]
令 $u=\cos\theta$，边值 $f(u)=V\operatorname{sgn}u$ 是奇函数，所以偶数 $l$ 系数为零。由 Legendre 正交性，对奇数 $l$，
\[
 A_l=\frac{2l+1}{2}\int_{-1}^1f(u)P_l(u)\dd u
 =V(2l+1)\int_0^1P_l(u)\dd u.
\]
利用
$\int P_l\dd u=[P_{l+1}-P_{l-1}]/(2l+1)$，得到
\[
 \boxed{A_l=V[P_{l-1}(0)-P_{l+1}(0)]\quad(l\text{ 奇})}.
\]
故
\[
 \boxed{\varphi=V\sum_{n=0}^\infty
 [P_{2n}(0)-P_{2n+2}(0)]
 \left(\frac rR\right)^{2n+1}P_{2n+1}(\cos\theta)}.
\]
前两项是
$\varphi=(3V/2)(r/R)P_1-(7V/8)(r/R)^3P_3+\cdots$。中心势为0，符合南北反对称；在任意 $r<R$ 级数绝对收敛。
\end{advsolution}

\begin{advexamplebox}{六面给定电势的立方体中心}\label{advA4:ex:cube-center-potential}
立方体两相对面维持 $+V$，其余四面维持 $-V$，棱边窄缝忽略。求中心电势。
\end{advexamplebox}
\begin{advsolution}
Dirichlet 解对边值线性。中心关于立方体的六个面完全对称，因此从任一面传到中心的调和测度相同；六个权重之和为1，所以每个面权重是 $1/6$。棱和顶点的集合面积为零，不影响调和测度。于是
\[
 \boxed{\varphi(0)=\frac16[2(+V)+4(-V)]=-\frac V3}.
\]
也可把六个“单独一面为1、其余面为0”的解相加：常数边界1的解处处为1，中心对称性迫使六个单面解在中心相等，故各为 $1/6$。这避免了不必要的三重 Fourier 级数。
\end{advsolution}

\begin{advtheorem}{球谐函数加法定理}\label{advA4:thm:spherical-addition}
单位球面上两方向 $(\theta,\phi)$ 与 $(\theta',\phi')$ 夹角为 $\gamma$，则
\[
 \boxed{P_l(\cos\gamma)=\frac{4\pi}{2l+1}
 \sum_{m=-l}^{l}Y_{lm}(\theta,\phi)Y_{lm}^*(\theta',\phi')}.
\]
\end{advtheorem}
\noindent\textbf{证明}\quad 固定第二方向，右边是关于第一方向的 $l$ 阶球谐函数，并在同时旋转两方向时不变，因此只能是 $C_lP_l(\cos\gamma)$。令两方向重合并在球面上积分：正交归一性给出右边积分为
$\sum_m1=2l+1$；而旋转对称性说明
$\sum_m|Y_{lm}|^2$ 是常数，其球面积分为 $2l+1$，故常数为 $(2l+1)/(4\pi)$。于是 $C_l=(2l+1)/(4\pi)$ 的倒数，得到公式。\hfill$\square$

\begin{advexamplebox}{接地球内的均匀带电直径}\label{advA4:ex:charged-diameter-in-sphere}
半径 $R$ 的接地导体球内，沿 $z$ 轴直径放置总电荷 $Q$ 的均匀细线。求球内电势的一种完整表示。
\end{advexamplebox}
\begin{advsolution}
线密度 $\lambda=Q/(2R)$。球内 Dirichlet Green 函数对位于
$\vect r'=z'\vect e_z$ 的源点为
\[
 G_D(\vect r,\vect r')=
 \frac1{|\vect r-\vect r'|}
 -\frac{R/|z'|}{|\vect r-(R^2/z')\vect e_z|},
\]
其中 $z'=0$ 时第二项取连续极限 $1/R$。因此
\[
 \boxed{\varphi(\vect r)=\frac{\lambda}{4\pi\varepsilon_0}
 \int_{-R}^{R}\left[
 \frac1{|\vect r-z'\vect e_z|}
 -\frac{R/|z'|}{|\vect r-(R^2/z')\vect e_z|}
 \right]\dd z'}.
\]
第二项表示球外轴线上的连续像电荷分布。对任意球面点
$|\vect r|=R$，反演恒等式
\[
 \left|\vect r-\frac{R^2}{z'}\vect e_z\right|
 =\frac R{|z'|}|\vect r-z'\vect e_z|
\]
使方括号逐点为零，所以边界条件被精确满足。积分关于 $z'$ 对称，故势关于赤道对称；线端点的理想化会造成局部奇异，但球内离开线段处势有限且满足 Laplace 方程。
\end{advsolution}

\subsection{静电平衡与电容矩阵}

设有若干导体。电势与各导体电荷量线性相关，可写成
$\varphi_i=P_{ij}q_j$。其逆关系为
$q_i=C_{ij}\varphi_j$，$\mathsf C=\mathsf P^{-1}$ 称电容矩阵；$C_{ii}$ 是自容，非对角项描述互容。

\begin{advtheorem}{静电能是电荷的正定二次型}\label{advA4:thm:capacitance-energy}
导体体系的静电能为
\[
 \boxed{U=\frac12P_{ij}q_iq_j
 =\frac12C_{ij}\varphi_i\varphi_j\ge0},
\]
且
\[
 \varphi_i=\frac{\partial U}{\partial q_i}.
\]
\end{advtheorem}
\noindent\textbf{证明}\quad 准静态地把电荷从零增加到最终值。在线性体系中，中间态 $sq_i$ 的电势为 $s\varphi_i$，所做功
\[
 U=\int_0^1(s\varphi_i)q_i\dd s=\frac12q_i\varphi_i.
\]
代入两种线性关系得到两个二次型。另一方面
$U=(\varepsilon_0/2)\int E^2\dd V\ge0$，除所有电荷为零外严格为正，所以 $\mathsf P$（在去掉整体规范自由度后）正定。对 $q_i$ 求偏导得到导体电势。\hfill$\square$

\begin{advtheorem}{电容矩阵的互易性}\label{advA4:thm:capacitance-reciprocity}
\[
 \boxed{P_{ij}=P_{ji},\qquad C_{ij}=C_{ji}}.
\]
\end{advtheorem}
\noindent\textbf{证明}\quad 取状态1仅第 $i$ 个导体带单位电荷，状态2仅第 $j$ 个导体带单位电荷。Green 互易关系给出状态1的第 $j$ 个导体电势等于状态2的第 $i$ 个导体电势，即
$P_{ji}=P_{ij}$。实对称可逆矩阵的逆仍对称，所以
$C_{ij}=C_{ji}$。等价地，二次型 $U$ 的混合二阶导数可交换。\hfill$\square$


\paragraph{参考电势、互容的符号与两导体模型。}
以上正定性采用无穷远电势为零并保留对无穷远的场能。若只保留两导体间的理想电容 $C$，则
\[
 \binom{q_1}{q_2}=C\begin{pmatrix}1&-1\\-1&1\end{pmatrix}
 \binom{\varphi_1}{\varphi_2},\qquad
 U=\frac C2(\varphi_1-\varphi_2)^2.
\]
此矩阵只有半正定性，共同平移电势不改变场，故不能直接求逆；选第二导体接地后，剩余关系才是 $q_1=C\varphi_1$。矩阵非对角元通常为负，而工程上把正数 $-C_{12}$ 称互电容，必须先说明约定。
以半径 $a<b$、长度 $L\gg b$ 的同轴圆柱为例，忽略端部，Gauss 定律给 $E(r)=q/(2\pi\varepsilon Lr)$，从而
\[
 \varphi_1-\varphi_2=\frac q{2\pi\varepsilon L}\ln\frac ba,
 \qquad C=\frac{2\pi\varepsilon L}{\ln(b/a)}.
\]
将场能积分 $\int_a^b\varepsilon E^2(2\pi rL)\dd r/2$ 代入，恰得 $q^2/(2C)$；窄隙 $b-a\ll a$ 时又趋于平行板结果 $\varepsilon(2\pi aL)/(b-a)$。

% AdvancedEMQuality20261023
\section{静磁场}

\subsection{电流分布、矢势与磁能}

稳恒电流满足 $\nabla\cdot\vect j=0$，静磁场方程为
\[
 \nabla\times\vect B=\mu_0\vect j,
 \qquad \nabla\cdot\vect B=0.
\]
因此可引入矢势
\[
 \boxed{\vect A(\vect r)=\frac{\mu_0}{4\pi}
 \int\frac{\vect j(\vect r')}{|\vect r-\vect r'|}\dd V'},
 \qquad \vect B=\nabla\times\vect A,
\]
即 Biot--Savart 公式
\[
 \vect B=\frac{\mu_0}{4\pi}\int
 \frac{\vect j(\vect r')\times(\vect r-\vect r')}
 {|\vect r-\vect r'|^3}\dd V'.
\]
对线电流，$\vect j\dd V$ 替换为 $I\dd\vect l$。积分形式为
\[
 \oint\vect B\cdot\dd\vect S=0,
 \qquad \oint\vect B\cdot\dd\vect l=\mu_0I_{\rm enc},
 \qquad \int_S\vect B\cdot\dd\vect S=\oint_{\partial S}\vect A\cdot\dd\vect l.
\]

\begin{advexamplebox}{Biot--Savart 矢势的 Coulomb 规范}\label{advA4:ex:coulomb-gauge}
证明在稳恒局域电流条件下，上述 $\vect A$ 自动满足
$\nabla\cdot\vect A=0$。

\begin{center}\begin{minipage}{.96\linewidth}\centering
\begin{tikzpicture}[>=Latex,font=\small,scale=.85,every node/.style={fill=none}]

\draw[thick] (0,0) ellipse (1.2 and .6);\draw[DeepBlue,->] (1.2,0) arc[start angle=0,end angle=65,x radius=1.2,y radius=.6];
\fill (2.5,1.2) circle (1.7pt) node[right]{$\boldsymbol r$};\fill (0,.6) circle (1.5pt) node[above]{$\boldsymbol r'$};\draw[dashed] (0,.6)--(2.5,1.2) node[midway,above]{$|\boldsymbol r-\boldsymbol r'|$};
\node[below] at (0,-.9){局域稳恒电流，$\nabla'\cdot\boldsymbol j=0$};

\end{tikzpicture}
\figcap{题图：源点与场点的导数作用在同一个距离核上，源电流局域且无散。}\label{advA4:quality:A4-magnetostatics:q1}
\end{minipage}\end{center}
\end{advexamplebox}
\global\AdvHideHeadingtrue
\begin{advsolution}
\begin{center}\begin{minipage}{.96\linewidth}\centering
% EMAdvFigureHeading
\raggedright{\bfseries\color{DeepBlue}解答}\par\smallskip\centering

\begin{tikzpicture}[>=Latex,font=\small,scale=.85,every node/.style={fill=none}]

\draw[thick] (0,0) ellipse (1.2 and .6);\draw[DeepBlue,->] (1.2,0) arc[start angle=0,end angle=65,x radius=1.2,y radius=.6];
\fill (2.5,1.2) circle (1.7pt) node[right]{$\boldsymbol r$};\fill (0,.6) circle (1.5pt) node[above]{$\boldsymbol r'$};\draw[dashed] (0,.6)--(2.5,1.2) node[midway,above]{$|\boldsymbol r-\boldsymbol r'|$};
\node[below] at (0,-.9){局域稳恒电流，$\nabla'\cdot\boldsymbol j=0$};

\end{tikzpicture}
\figcap{解图：源点与场点的导数作用在同一个距离核上，源电流局域且无散。}\label{advA4:quality:A4-magnetostatics:a1}
\end{minipage}\end{center}

利用 $\nabla(1/R)=-\nabla'(1/R)$，
\[
 \nabla\cdot\vect A
 =\frac{\mu_0}{4\pi}\int\vect j'\cdot\nabla\frac1R\dd V'
 =-\frac{\mu_0}{4\pi}\int\vect j'\cdot\nabla'\frac1R\dd V'.
\]
把被积式写成
\[
 \vect j'\cdot\nabla'\frac1R
 =\nabla'\cdot\frac{\vect j'}R-\frac{\nabla'\cdot\vect j'}R.
\]
第二项因稳恒连续性 $\nabla'\cdot\vect j'=0$ 消失；第一项化为无穷远边界通量，局域电流使其为零。因此
\[
 \boxed{\nabla\cdot\vect A=0}.
\]
若电流穿过所选有限区域边界，边界项不能任意丢弃；应把回流路径一并计入。

图~\ref{advA4:quality:A4-magnetostatics:a1} 中，同一核满足 $\nabla(1/R)=-\nabla'(1/R)$；对源坐标分部积分时，体项由稳恒连续性为零，远处面项由局域性为零。
\end{advsolution}

静磁能为
\[
 U_B=\frac12\int\vect j\cdot\vect A\dd V.
\]

\begin{advtheorem}{静磁 Green 互易关系}\label{advA4:thm:magnetic-reciprocity}
两组局域稳恒电流及其 Coulomb 规范矢势满足
\[
 \boxed{\int\vect j_1\cdot\vect A_2\dd V
 =\int\vect j_2\cdot\vect A_1\dd V}.
\]
\end{advtheorem}
\noindent\textbf{证明}\quad 代入积分核，
\[
 \int\vect j_1(\vect r)\cdot\vect A_2(\vect r)\dd V
 =\frac{\mu_0}{4\pi}\iint
 \frac{\vect j_1(\vect r)\cdot\vect j_2(\vect r')}
 {|\vect r-\vect r'|}\dd V\dd V'.
\]
交换 $\vect r,\vect r'$ 及下标1、2，积分不变，即得结论。\hfill$\square$

\begin{advtheorem}{磁场能}\label{advA4:thm:magnetic-energy}
局域稳恒电流的能量也可写成
\[
 \boxed{\frac12\int\vect j\cdot\vect A\dd V
 =\frac1{2\mu_0}\int B^2\dd V}.
\]
\end{advtheorem}
\noindent\textbf{证明}\quad 由 $\vect j=\mu_0^{-1}\nabla\times\vect B$ 和
\[
 \nabla\cdot(\vect A\times\vect B)
 =\vect B\cdot\nabla\times\vect A-
 \vect A\cdot\nabla\times\vect B
\]
得到
\[
 \int\vect j\cdot\vect A\dd V
 =\frac1{\mu_0}\int B^2\dd V
 -\frac1{\mu_0}\oint(\vect A\times\vect B)\cdot\dd\vect S.
\]
局域场在无穷远的边界项消失，乘以 $1/2$ 即得。\hfill$\square$

\begin{advexamplebox}{线圈系统的电感矩阵}\label{advA4:ex:inductance-matrix}
证明真空或线性互易介质中 $N$ 个线圈的磁能为
\[
 U_B=\frac12\sum_{i,j}L_{ij}I_iI_j,
\]
其中穿过线圈 $i$ 的磁链 $\Phi_i=\sum_jL_{ij}I_j$，并证明
$L_{ij}=L_{ji}$。说明固定电流求机械力时为什么使用磁共能。

\begin{center}\begin{minipage}{.96\linewidth}\centering
\begin{tikzpicture}[>=Latex,font=\small,scale=.85,every node/.style={fill=none}]

\draw[thick] (0,0) ellipse (1.1 and .5);\draw[thick] (2.6,.7) ellipse (1.1 and .5);\node at (0,-.8){线圈 $i,I_i$};\node at (2.6,1.6){线圈 $j,I_j$};\draw[dashed] (.5,.35)--(2.1,.45) node[midway,above]{$M_{ij}$};

\end{tikzpicture}
\figcap{题图：相互交链的线圈系统，各电流和磁链都带同一个端口方向约定。}\label{advA4:quality:A4-magnetostatics:q2}
\end{minipage}\end{center}
\end{advexamplebox}
\global\AdvHideHeadingtrue
\begin{advsolution}
\begin{center}\begin{minipage}{.96\linewidth}\centering
% EMAdvFigureHeading
\raggedright{\bfseries\color{DeepBlue}解答}\par\smallskip\centering

\begin{tikzpicture}[>=Latex,font=\small,scale=.85,every node/.style={fill=none}]

\draw[thick] (0,0) ellipse (1.1 and .5);\draw[thick] (2.6,.7) ellipse (1.1 and .5);\node at (0,-.8){线圈 $i,I_i$};\node at (2.6,1.6){线圈 $j,I_j$};\draw[dashed] (.5,.35)--(2.1,.45) node[midway,above]{$M_{ij}$};

\end{tikzpicture}
\figcap{解图：相互交链的线圈系统，各电流和磁链都带同一个端口方向约定。}\label{advA4:quality:A4-magnetostatics:a2}
\end{minipage}\end{center}

% EMReciprocalMedium
“互易”是这里不可省略的条件。真空可直接使用前述积分核；对于位置可变但满足 $\mu_{ij}=\mu_{ji}$ 的线性磁介质，忽略无穷远面项后有
\[
 \int\boldsymbol j_1\cdot\boldsymbol A_2\,\mathrm dV
 =\int\boldsymbol H_1\cdot\boldsymbol B_2\,\mathrm dV
 =\int\boldsymbol H_2\cdot\boldsymbol B_1\,\mathrm dV
 =\int\boldsymbol j_2\cdot\boldsymbol A_1\,\mathrm dV.
\]
中间等号才使用介质张量的对称性。一般非互易介质不能仅凭“线性”就声称电感矩阵对称。

把总电流写成 $\vect j=\sum_iI_i\vect j_i^{(1)}$，矢势也线性叠加。于是
\[
 U_B=\frac12\sum_{i,j}I_iI_j
 \int\vect j_i^{(1)}\cdot\vect A_j^{(1)}\dd V
 \equiv\frac12\sum_{i,j}L_{ij}I_iI_j.
\]
对细线圈，$\int\vect j_i^{(1)}\cdot\vect A_j^{(1)}\dd V
=\oint_i\vect A_j^{(1)}\cdot\dd\vect l=\Phi_{i\leftarrow j}/I_j$，故
$\Phi_i=\partial U_B/\partial I_i=\sum_jL_{ij}I_j$。定理\ref{advA4:thm:magnetic-reciprocity}立即给出
$L_{ij}=L_{ji}$。

若外电源固定 $I_i$，构型改变时电源还会做功
$\sum_iI_i\dd\Phi_i$。Legendre 变换后的机械势为
\[
 V_{\rm mech}=U_B-\sum_iI_i\Phi_i=-U_B
\]
（线性体系），所以广义力
$Q_\alpha=\partial U_B/\partial q_\alpha|_I$。若超导回路固定磁链，则直接在固定 $\Phi$ 下用场能，两个条件不能混用。

图~\ref{advA4:quality:A4-magnetostatics:a2} 中，互感项连接两个不同线圈，能量双重求和中的 $1/2$ 防止把同一对耦合重复计数；电流固定的机械变化必须把电源做功一起核算。
\end{advsolution}

\subsection{磁多极矩}

远区矢势可按
$1/|\vect r-\vect r'|$ 展开。零阶项含
$\int\vect j\dd V=0$；其证明是对任意常矢量 $\vect c$ 写
\[
 \vect c\cdot\int\vect j\dd V
 =\int\nabla\cdot[(\vect c\cdot\vect r)\vect j]\dd V=0.
\]
一阶矩张量为
\[
 M_{ij}=\int x_i j_j\dd V.
\]

\begin{advtheorem}{稳恒电流的一阶矩反对称}\label{advA4:thm:moment-antisymmetric}
\[
 \boxed{M_{ij}+M_{ji}=0}.
\]
因而可用磁偶极矩
\[
 \boxed{\vect m=\frac12\int\vect r\times\vect j\dd V}
\]
完全表示它。
\end{advtheorem}
\noindent\textbf{证明}\quad 对局域电流积分
\[
 \partial_k(x_ix_jj_k)=x_jj_i+x_ij_j+x_ix_j\partial_kj_k.
\]
最后一项由 $\nabla\cdot\vect j=0$ 消失；左边体积分化为无穷远边界通量并为零，故
$\int(x_jj_i+x_ij_j)\dd V=0$。反对称二阶张量在三维与矢量对偶，正是所列 $\vect m$。\hfill$\square$

远区矢势和磁场为
\[
 \vect A=\frac{\mu_0}{4\pi}\frac{\vect m\times\vect r}{r^3},
 \qquad (r\ne0).
\]

\begin{advtheorem}{点磁偶极子的完整分布场}\label{advA4:thm:magnetic-dipole-field}
\[
 \boxed{\vect B(\vect r)=\frac{\mu_0}{4\pi}
 \frac{3(\vect m\cdot\vect r)\vect r-r^2\vect m}{r^5}
 +\frac{2\mu_0}{3}\vect m\,\delta^{(3)}(\vect r)}.
\]
\end{advtheorem}
\noindent\textbf{证明}\quad 由
$\vect A=-(\mu_0/4\pi)\vect m\times\nabla(1/r)$ 计算旋度，并使用分布 Hessian
\[
 \partial_i\partial_j\frac1r
 =\frac{3x_ix_j-r^2\delta_{ij}}{r^5}
 -\frac{4\pi}{3}\delta_{ij}\delta^{(3)}(\vect r).
\]
两个缩并项组合后，$\delta$ 系数为 $2\mu_0/3$。该项保证
$\nabla\cdot\vect B=0$ 在原点也成立。\hfill$\square$

外场缓变时
\[
 \vect M=\vect m\times\vect B,
 \qquad \vect F=\nabla(\vect m\cdot\vect B),
\]
可由小电流环或势能 $U=-\vect m\cdot\vect B$ 推得。

\begin{advexamplebox}{可伸缩超导螺线管}\label{advA4:ex:superconducting-solenoid}
轻质软螺线管截面积恒为 $S$、匝数 $N$、初长 $l$、初始电流 $I$。两端各接质量 $m$ 的质点，初态静止，释放后求长度 $x(t)$ 和电流 $i(t)$，忽略边缘效应。

\begin{center}
\begin{tikzpicture}[x=1cm,y=1cm,>=Latex]
 \draw[thick] (-3,0) circle(.32) node {$m$};
 \draw[thick] (3,0) circle(.32) node {$m$};
 \draw[thick] (-2.68,0) -- (-2.35,0);
 \draw[thick,decorate,decoration={coil,segment length=5.5pt,amplitude=5pt}]
   (-2.35,0)--(2.35,0);
 \draw[thick] (2.35,0)--(2.68,0);
 \draw[<->] (-2.35,-.72)--(2.35,-.72) node[midway,fill=none,below] {$x(t)$};
 \draw[->,DeepBlue] (-3.45,0)--(-3.9,0) node[left] {$+q_L$};
 \draw[->,DeepBlue] (3.45,0)--(3.9,0) node[right] {$+q_R$};
 \node[above] at (0,.48) {$N$ 匝，截面积 $S$};
\node at(0,-1.5){箭头规定两端向外位移的正向，非实际速度};
\end{tikzpicture}
\figcap{题图：端点质点、可变线圈长度和固定截面积，原图保留。}\label{advA4:quality:A4-magnetostatics:q3}

\end{center}
\end{advexamplebox}
\global\AdvHideHeadingtrue
\begin{advsolution}
\begin{center}\begin{minipage}{.96\linewidth}\centering
% EMAdvFigureHeading
\raggedright{\bfseries\color{DeepBlue}解答}\par\smallskip\centering

\begin{tikzpicture}[x=1cm,y=1cm,>=Latex]
 \draw[thick] (-3,0) circle(.32) node {$m$};
 \draw[thick] (3,0) circle(.32) node {$m$};
 \draw[thick] (-2.68,0) -- (-2.35,0);
 \draw[thick,decorate,decoration={coil,segment length=5.5pt,amplitude=5pt}]
   (-2.35,0)--(2.35,0);
 \draw[thick] (2.35,0)--(2.68,0);
 \draw[<->] (-2.35,-.72)--(2.35,-.72) node[midway,fill=none,below] {$x(t)$};
 \draw[->,DeepBlue] (-3.45,0)--(-3.9,0) node[left] {$+q_L$};
 \draw[->,DeepBlue] (3.45,0)--(3.9,0) node[right] {$+q_R$};
 \node[above] at (0,.48) {$N$ 匝，截面积 $S$};
\node at(0,-1.5){箭头规定两端向外位移的正向，非实际速度};
\end{tikzpicture}
\figcap{解图：端点质点、可变线圈长度和固定截面积，原图保留。}\label{advA4:quality:A4-magnetostatics:a3}
\end{minipage}\end{center}

长螺线管自感
\[
 L(x)=\frac{\mu_0N^2S}{x}.
\]
超导回路磁链守恒：
\[
 L(x)i(t)=L(l)I
 \quad\Longrightarrow\quad
 \boxed{i(t)=I\frac{x(t)}l}.
\]
固定磁链下的磁能
\[
 U_B=\frac12L(x)i^2
 =\frac{\mu_0N^2SI^2}{2l^2}x\equiv ax.
\]
两质点关于质心对称运动，各自速度大小 $|\dot x|/2$，总动能
$K=m\dot x^2/4$。能量守恒
\[
 \frac m4\dot x^2+ax=al
\]
或约化质量方程 $(m/2)\ddot x=-a$ 给出
\[
 \boxed{x(t)=l-\frac{\mu_0N^2SI^2}{2ml^2}t^2},
\]
\[
 \boxed{i(t)=I\left(1-\frac{\mu_0N^2SI^2}{2ml^3}t^2\right)}.
\]
此式适用于 $0\le t\le t_c$，其中
$t_c=\sqrt{2ml^3/(\mu_0N^2SI^2)}$。磁能减少恰等于两质点动能；$x\downarrow0$ 时理想模型的长螺线管近似失效，不能继续外推。

图~\ref{advA4:quality:A4-magnetostatics:a3} 中，长度变化改变自感 $L(x)$，理想超导闭回路保持磁链而不是电流；两端以相反速度运动，长度坐标的等效质量须由两质点动能相加得到。
\end{advsolution}

\subsection{磁介质与磁标势}

% CompletenessEM:a4-magnetic-medium
\noindent\textbf{从磁偶极矩密度到束缚电流。}
磁化强度 $\vect M$ 是单位体积的磁偶极矩。令
$R=|\vect r-\vect r'|$，体元的偶极矢势叠加为
\[
 \vect A_M(\vect r)=\frac{\mu_0}{4\pi}\int_V
 \vect M(\vect r')\times\nabla'\frac1R\dd V'.
\]
由 $\nabla'\times(\vect M/R)
=\nabla'(1/R)\times\vect M+(\nabla'\times\vect M)/R$，
并把旋度体积分转为边界积分，得到
\[
 \vect A_M=\frac{\mu_0}{4\pi}\left[
 \int_V\frac{\nabla'\times\vect M}{R}\dd V'
 +\oint_{\partial V}\frac{\vect M\times\vect n}{R}\dd S'\right].
\]
与电流的矢势比较，便识别出
\[
 \boxed{\vect j_M=\nabla\times\vect M,\qquad
 \vect K_M=\vect M\times\vect n}.
\]
若界面两侧都有磁化，$\vect n$ 从介质1指向介质2，则净面流为
$(\vect M_1-\vect M_2)\times\vect n$。
均匀磁化体内部 $\vect j_M=0$，并不意味着没有磁场，边界面流仍然存在。

\begin{qualityfigure}
\begin{tikzpicture}[>=Latex,font=\small,scale=.9]
 \fill[PaleBlue](-1.6,-1.1)rectangle(1.3,1.35);
 \draw[thick](1.3,-1.1)--(1.3,1.35);
 \foreach\x in{-.8,.2}{\draw[->,DeepBlue,thick](\x,-.65)--(\x,.7);}
 \node[left]at(-.8,.1){$\vect M$};
 \draw[->](1.3,.9)--(2.45,.9)node[right]{$\vect n=\vect e_x$};
 \node[draw,circle,inner sep=1pt]at(1.3,-.35){$\times$};
 \node[right]at(1.7,-.35){$\vect K_M=M\vect e_y$（入纸面）};
 \draw[->](-1.8,-1.65)--(-.65,-1.65)node[right]{$x$};
 \draw[->](-1.8,-1.65)--(-1.8,-.6)node[above]{$z$};
\end{tikzpicture}
\figcap{右手坐标的 $xz$ 截面：$+y$ 指向纸内，因此 $\vect e_z\times\vect e_x=\vect e_y$ 的面流画作叉号。\label{complete:em:bound-current}}
\end{qualityfigure}

图~\ref{complete:em:bound-current} 把 $\vect M$、外法线和面流三个方向同时固定，
可用于检查下一例磁化球的面流符号。定义
$\vect H=\vect B/\mu_0-\vect M$ 后，静磁 Maxwell 方程成为
\[
 \nabla\times\vect H=\vect j_f,\qquad
 \nabla\cdot\vect H=-\nabla\cdot\vect M.
\]
在无自由电流的单连通区域才能全局写 $\vect H=-\nabla\psi_M$；于是
\[
 \rho_M=-\nabla\cdot\vect M,\qquad
 \sigma_M=(\vect M_1-\vect M_2)\cdot\vect n,
\]
\[
 \nabla^2\psi_M=-\rho_M,\qquad
 \psi_M(\vect r)=\frac1{4\pi}\left[
 \int\frac{\rho_M}{R}\dd V'+\oint\frac{\sigma_M}{R}\dd S'\right].
\]
这些磁荷只是计算 $\vect H$ 的等效源；实际 $\vect B=\mu_0(\vect H+\vect M)$
仍满足 $\nabla\cdot\vect B=0$，不能把等效磁荷理解为真实磁单极。
若有自由电流，可先求其真空形式的场
\[
 \vect H_f=\frac1{4\pi}\int
 \frac{\vect j_f(\vect r')\times(\vect r-\vect r')}{R^3}\dd V',
 \qquad \vect H=\vect H_f-\nabla\psi_M,
\]
再联立磁化关系与边界条件确定 $\vect M$。

对准静态、无磁滞且可逆的介质，固定材料参数时增量能量为
$\dd w=\vect H\cdot\dd\vect B$。
若本构关系是通过原点的线性各向同性关系
$\vect M=\chi_m\vect H$、$\vect B=\mu\vect H$，从零场积分得到
\[
 w=\int_0^{\vect B}\vect H(\vect B')\cdot\dd\vect B'
 =\frac12\vect B\cdot\vect H.
\]
有永久磁化、磁滞或色散时，不能不加条件地套用这一简单能量式。



\begin{advexamplebox}{均匀磁化球的内外磁场}\label{advA4:ex:uniform-magnetized-sphere}
半径 $R$ 的球体具有均匀磁化强度 $\vect M=M\vect e_z$。求内外 $\vect H,\vect B$。

\begin{center}\begin{minipage}{.96\linewidth}\centering
\begin{tikzpicture}[>=Latex,font=\small,scale=.85,every node/.style={fill=none}]

\draw[thick] (0,0) circle (1.2);\foreach \x in {-.5,0,.5}{\draw[DeepBlue,->] (\x,-.6)--(\x,.65);}\node[above] at (0,1.5){$\boldsymbol M=M\boldsymbol e_z$};\draw[->] (0,0)--(1.2,0) node[midway,below]{$R$};\node[right] at (1.7,.4){球外无自由电流};

\end{tikzpicture}
\figcap{题图：均匀磁化球、磁化轴与半径，内外场均待求。}\label{advA4:quality:A4-magnetostatics:q4}
\end{minipage}\end{center}
\end{advexamplebox}
\global\AdvHideHeadingtrue
\begin{advsolution}
\begin{center}\begin{minipage}{.96\linewidth}\centering
% EMAdvFigureHeading
\raggedright{\bfseries\color{DeepBlue}解答}\par\smallskip\centering

\begin{tikzpicture}[>=Latex,font=\small,scale=.85,every node/.style={fill=none}]

\draw[thick] (0,0) circle (1.2);\foreach \x in {-.5,0,.5}{\draw[DeepBlue,->] (\x,-.6)--(\x,.65);}\node[above] at (0,1.5){$\boldsymbol M=M\boldsymbol e_z$};\draw[->] (0,0)--(1.2,0) node[midway,below]{$R$};\node[right] at (1.7,.4){球外无自由电流};

\end{tikzpicture}
\figcap{解图：均匀磁化球、磁化轴与半径，内外场均待求。}\label{advA4:quality:A4-magnetostatics:a4}
\end{minipage}\end{center}

体磁荷 $-\nabla\cdot\vect M=0$，表面磁荷
$\sigma_M=M\cos\theta$。轴对称标势取
\[
 \psi_{\rm in}=Ar\cos\theta,\qquad
 \psi_{\rm out}=\frac{B\cos\theta}{r^2}.
\]
切向 $H$ 连续，法向满足
$H_{n,\rm out}-H_{n,\rm in}=\sigma_M$。代入得到
\[
 \boxed{\vect H_{\rm in}=-\frac13\vect M},
 \qquad
 \boxed{\vect B_{\rm in}=\frac{2\mu_0}{3}\vect M}.
\]
球外等效于位于球心的磁偶极子
\[
 \vect m=\frac{4\pi R^3}{3}\vect M,
\]
所以
\[
 \boxed{\vect B_{\rm out}=\frac{\mu_0}{4\pi r^3}
 [3(\vect m\cdot\vect e_r)\vect e_r-\vect m]}.
\]
球外是真空，另有 $\boldsymbol H_{\rm out}=\boldsymbol B_{\rm out}/\mu_0$。在 $r=R$ 检查 $B_n$ 连续、$H_t$ 连续，边界条件满足。

图~\ref{advA4:quality:A4-magnetostatics:a4} 中，均匀磁化使体束缚电流为零、面束缚电流为 $\boldsymbol M\times\boldsymbol n$；内外标势的连续和法向跃变决定退磁场，不能把 $H$ 与 $B/\mu_0$ 在球内等同。
\end{advsolution}

\begin{advexamplebox}{自由电流圈的磁偶极层与磁标势}\label{advA4:ex:current-loop-scalar-potential}
证明载流 $I$ 的闭合线圈在无电流的单连通子区域内可由磁标势描述，并等效于跨任意张面 $S$ 的磁偶极层；说明跨张面的支路跳跃。

\begin{center}\begin{minipage}{.96\linewidth}\centering
\begin{tikzpicture}[>=Latex,font=\small,scale=.85,every node/.style={fill=none}]

\fill[DeepBlue!8] (0,0) ellipse (1.4 and .65);\draw[thick] (0,0) ellipse (1.4 and .65);\draw[DeepBlue,->] (1.4,0) arc[start angle=0,end angle=65,x radius=1.4,y radius=.65];\draw[->] (0,0)--(0,1.5) node[above]{$\boldsymbol n$};\node at (.3,-.2){张面 $S$};\node[right] at (1.7,0){边界电流 $I$};

\end{tikzpicture}
\figcap{题图：闭合电流圈与任取的跨面，张面法向和边界绕行方向按右手规则关联。}\label{advA4:quality:A4-magnetostatics:q5}
\end{minipage}\end{center}
\end{advexamplebox}
\global\AdvHideHeadingtrue
\begin{advsolution}
\begin{center}\begin{minipage}{.96\linewidth}\centering
% EMAdvFigureHeading
\raggedright{\bfseries\color{DeepBlue}解答}\par\smallskip\centering

\begin{tikzpicture}[>=Latex,font=\small,scale=.85,every node/.style={fill=none}]

\fill[DeepBlue!8] (0,0) ellipse (1.4 and .65);\draw[thick] (0,0) ellipse (1.4 and .65);\draw[DeepBlue,->] (1.4,0) arc[start angle=0,end angle=65,x radius=1.4,y radius=.65];\draw[->] (0,0)--(0,1.5) node[above]{$\boldsymbol n$};\node at (.3,-.2){张面 $S$};\node[right] at (1.7,0){边界电流 $I$};

\end{tikzpicture}
\figcap{解图：闭合电流圈与任取的跨面，张面法向和边界绕行方向按右手规则关联。}\label{advA4:quality:A4-magnetostatics:a5}
\end{minipage}\end{center}

张面法线与正电流方向满足右手关系，其有向立体角定义为
\[
 \Omega(\vect r)=\int_S
 \frac{(\vect r-\vect r')\cdot\dd\vect S'}{|\vect r-\vect r'|^3}.
\]
% CompletenessEM:a4-loop-stokes
下面把这个梯度具体化。记 $f=1/|\vect r-\vect r'|$，则
$\nabla'f=(\vect r-\vect r')/R^3$，且对核有 $\partial_i f=-\partial'_i f$。
在场点不在张面上时，$\nabla'^2f=0$。用分量形式的 Stokes 公式，
\[
 \begin{aligned}
 \left[\oint_{\partial S}\dd\vect l'\times\nabla'f\right]_i
 &=\int_S\left(n_j\partial'_i\partial'_jf
       -n_i\nabla'^2f\right)\dd S'\\
 &=\int_S n_j\partial'_i\partial'_jf\dd S'
 =-\partial_i\int_S\vect n\cdot\nabla'f\dd S'
 =-\partial_i\Omega.
 \end{aligned}
\]
第一行可由 $\varepsilon_{ijk}\varepsilon_{j\ell m}
=\delta_{im}\delta_{k\ell}-\delta_{i\ell}\delta_{km}$ 展开验证，
这也明确了源点导数与场点导数之间的负号。

\begin{qualityfigure}
\begin{tikzpicture}[>=Latex,font=\small,scale=.9]
 \fill[PaleBlue](0,0)ellipse(1.65 and .6);
 \draw[thick](0,0)ellipse(1.65 and .6);
 \draw[DeepBlue,->,thick](1.65,0)arc[start angle=0,end angle=65,x radius=1.65,y radius=.6];
 \node[right]at(1.65,.25){$I,\ \partial S$};
 \draw[->](0,-.95)--(0,1.6)node[above]{$\vect n$};
 \fill(0,.85)circle(1.7pt)node[right]{$h>0$};
 \fill(0,-.75)circle(1.7pt)node[right]{$h<0$};
 \draw[dashed](-.65,0)ellipse(.65 and .23);
 \draw[<->](-.65,-1.25)--(0,-1.25)node[midway,below]{$a_0$};
\end{tikzpicture}
\figcap{张面局部用半径 $a_0$ 的平面片核对支路跳跃；正法向与边界电流按右手规则关联。\label{complete:em:loop-cut}}
\end{qualityfigure}

图~\ref{complete:em:loop-cut} 中，令 $h$ 为场点相对局部平面片的有符号高度，
其立体角贡献为
\[
 \Omega_{\rm loc}(h)=2\pi\int_0^{a_0}
 \frac{hs\dd s}{(s^2+h^2)^{3/2}}
 =2\pi\left(\operatorname{sgn}h-\frac{h}{\sqrt{h^2+a_0^2}}\right).
\]
$h\to0^\pm$ 时分别为 $\pm2\pi$；远离这个小面片的部分连续，
所以跨割面的跳跃恰为 $4\pi$。这里不是说磁场在任意张面上产生真实面流奇点，
而是对两侧支路分别求梯度。综上得到


\[
 -\nabla\left(\frac I{4\pi}\Omega\right)
 =\frac I{4\pi}\oint_{\partial S}
 \frac{\dd\vect l'\times(\vect r-\vect r')}{R^3}
 =\vect H.
\]
因此
\[
 \boxed{\psi_H=\frac I{4\pi}\Omega},
 \qquad \vect B=-\mu_0\nabla\psi_H.
\]
按正法线方向穿过割面，$\Omega$ 从 $-2\pi$ 跳到 $+2\pi$，故选定支上的 $\psi_H$ 跳变 $+I$；沿正向闭合回路跟踪连续支，势累计下降 $I$，与环路积分
$\oint\vect H\cdot\dd\vect l=I$。面上单位面积磁偶极矩取
$I\vect n$ 就产生同一势，所以线圈等效于磁偶极层。标势是多值的，只有在割去张面后的单连通区域内才能单值使用。

图~\ref{advA4:quality:A4-magnetostatics:a5} 中，改变张面只改变标势所用的支路面，真实磁场由边界电流决定；跨面势函数可有跳跃，而绕过电流圈的环路不允许全局单值标势。
\end{advsolution}

\subsection{静磁边值问题}

无自由电流区域内 $\vect H=-\nabla\psi$ 且 $\nabla^2\psi=0$。理想超导体完全抗磁，内部 $\vect B=0$，表面给出
$\partial_n\psi=0$ 的 Neumann 条件。

\begin{advexamplebox}{反向平行电流、二维磁偶极与超导圆柱}\label{advA4:ex:parallel-currents-cylinder}
两条无限长 $z$ 向导线分别位于 $x=+d/2$ 和 $x=-d/2$，相应载流为 $+I$ 和 $-I$。求磁标势；取 $p=Id$、$d\to0$ 的二维偶极极限。再把该偶极置于半径 $R$ 的无限长超导圆柱空腔轴线上，求腔内平衡场。

\begin{center}\begin{minipage}{.96\linewidth}\centering
\begin{tikzpicture}[>=Latex,font=\small,scale=.85,every node/.style={fill=none}]

\draw[thick] (0,0) circle (1.4);\node at (-.4,0){$\otimes$};\node at (.4,0){$\odot$};\node[above] at (-.4,.3){$-I$};\node[above] at (.4,.3){$+I$};\draw[<->] (-.4,-.8)--(.4,-.8) node[midway,below]{$d$};\node[above] at (0,1.7){超导圆柱空腔横截面 $R$};

\end{tikzpicture}
\figcap{题图：两反向长电流及其二维偶极极限，柱面边界在横截面内取圆。}\label{advA4:quality:A4-magnetostatics:q6}
\end{minipage}\end{center}
\end{advexamplebox}
\global\AdvHideHeadingtrue
\begin{advsolution}
\begin{center}\begin{minipage}{.96\linewidth}\centering
% EMAdvFigureHeading
\raggedright{\bfseries\color{DeepBlue}解答}\par\smallskip\centering

\begin{tikzpicture}[>=Latex,font=\small,scale=.85,every node/.style={fill=none}]

\draw[thick] (0,0) circle (1.4);\node at (-.4,0){$\otimes$};\node at (.4,0){$\odot$};\node[above] at (-.4,.3){$-I$};\node[above] at (.4,.3){$+I$};\draw[<->] (-.4,-.8)--(.4,-.8) node[midway,below]{$d$};\node[above] at (0,1.7){超导圆柱空腔横截面 $R$};

\end{tikzpicture}
\figcap{解图：两反向长电流及其二维偶极极限，柱面边界在横截面内取圆。}\label{advA4:quality:A4-magnetostatics:a6}
\end{minipage}\end{center}

在不跨过割线的区域，单根 $+I$ 导线的标势是
$-I\arg(z-z_0)/(2\pi)$。因此在共同选定的割线之外，成对电流的标势为
\[
 \boxed{\psi=-\frac I{2\pi}
 [\arg(x-d/2+\ii y)-\arg(x+d/2+\ii y)]}.
\]
小 $d$ 展开得到
\[
 \boxed{\psi_0=-\frac p{2\pi}\frac y{x^2+y^2}
 =-\frac p{2\pi r}\sin\phi}.
\]
超导边界要求 $\partial_r\psi|_{R}=0$。加入在圆柱内正则的 $l=1$ 调和项，
\[
 \boxed{\psi=-\frac p{2\pi}\left(\frac{\sin\phi}{r}
 +\frac{r\sin\phi}{R^2}\right)}.
\]
于是 $\vect H=-\nabla\psi$，直角坐标分量为
\[
 H_x=-\frac{pxy}{\pi r^4},
 \qquad
 H_y=\frac p{2\pi}\left(\frac{x^2-y^2}{r^4}+\frac1{R^2}\right),
 \qquad \vect B=\mu_0\vect H.
\]
在 $r=R$ 直接计算 $H_r=0$；例如自由空间中 $x=0,y\ne0$ 处有 $H_y=-p/(2\pi y^2)$，也与右侧 $+I$、左侧 $-I$ 的 Biot--Savart 场相符。$R\to\infty$ 时附加均匀场消失，退回自由空间偶极。

图~\ref{advA4:quality:A4-magnetostatics:a6} 中，两导线的标势必须使用相同支路约定，取 $d\to0$ 时保持 $Id$ 固定；超导边界要求法向磁场为零，不能只规定标势本身为零。
\end{advsolution}

\begin{advexamplebox}{圆柱面电流与三圆柱相交体}\label{advA4:ex:steinmetz-currents}
无限薄圆柱面半径为 $R$，分别载
$\vect K=i_0\vect e_\phi$ 与
$\vect K=i_0\cos\phi\,\vect e_z$，求内外磁场。

三个互相垂直、半径均为 $R$ 的圆柱公共区域组成三圆柱相交体。其三类边界面上分别给定
\[
 \vect K_1=-i_0\cos\phi_y\vect e_x,\quad
 \vect K_2=i_0\cos\phi_x\vect e_y,\quad
 \vect K_3=i_0\vect e_{\phi_z}.
\]
证明电流连续并求远场。

\begin{center}
\begin{tikzpicture}[scale=.9]
 \draw[axis](-2,-1)--(2.4,1.2)node[right]{$x$};
 \draw[axis](2,-1)--(-2.4,1.2)node[left]{$y$};
 \draw[axis](0,-1.5)--(0,2.1)node[above]{$z$};
 \draw[DeepBlue] (0,0)ellipse(1.8 and .9);
 \draw[vector,Gold!80!black](.7,.8)--(-.3,.3);\node[anchor=west]at(.8,1.55){$\vect K_1\parallel-\vect e_x$};\draw[guide](.8,1.4)--(.2,.55);
 \draw[vector,DeepBlue](-.7,.8)--(-1.7,1.3)node[above]{$\vect K_2\parallel\vect e_y$};
 \draw[vector,TextGray](1.8,0)--(1.8,.65)node[right]{$\vect K_3$ 切向};
 \node at(0,-2){局部方向示意，非三圆柱相交体的投影边界};
\end{tikzpicture}
\figcap{题图：原圆柱与三圆柱相交边界的切向电流几何保留。}\label{advA4:quality:A4-magnetostatics:q7}

\end{center}
\end{advexamplebox}
\global\AdvHideHeadingtrue
\begin{advsolution}
\begin{center}\begin{minipage}{.96\linewidth}\centering
% EMAdvFigureHeading
\raggedright{\bfseries\color{DeepBlue}解答}\par\smallskip\centering

\begin{tikzpicture}[scale=.9]
 \draw[axis](-2,-1)--(2.4,1.2)node[right]{$x$};
 \draw[axis](2,-1)--(-2.4,1.2)node[left]{$y$};
 \draw[axis](0,-1.5)--(0,2.1)node[above]{$z$};
 \draw[DeepBlue] (0,0)ellipse(1.8 and .9);
 \draw[vector,Gold!80!black](.7,.8)--(-.3,.3);\node[anchor=west]at(.8,1.55){$\vect K_1\parallel-\vect e_x$};\draw[guide](.8,1.4)--(.2,.55);
 \draw[vector,DeepBlue](-.7,.8)--(-1.7,1.3)node[above]{$\vect K_2\parallel\vect e_y$};
 \draw[vector,TextGray](1.8,0)--(1.8,.65)node[right]{$\vect K_3$ 切向};
 \node at(0,-2){局部方向示意，非三圆柱相交体的投影边界};
\end{tikzpicture}
\figcap{解图：原圆柱与三圆柱相交边界的切向电流几何保留。}\label{advA4:quality:A4-magnetostatics:a7}
\end{minipage}\end{center}

环形面电流正是均匀磁化 $\vect M=i_0\vect e_z$ 的束缚面电流
$\vect K=\vect M\times\vect n$，所以
\[
 \boxed{\vect B_{\rm in}=\mu_0i_0\vect e_z,\qquad
 \vect B_{\rm out}=0}.
\]
对 $K_z=i_0\cos\phi$，取 $A_z$：
\[
 A_{z,\rm in}=ar\cos\phi,\qquad
 A_{z,\rm out}=\frac b r\cos\phi.
\]
$B_r$ 连续和
$B_{\phi,\rm out}-B_{\phi,\rm in}=\mu_0K_z$ 给出
$a=\mu_0i_0/2$、$b=aR^2$。于是
\[
 \boxed{\vect B_{\rm in}=-\frac{\mu_0i_0}{2}\vect e_y},
\]
\[
 \boxed{B_{r,\rm out}=-\frac{\mu_0i_0R^2}{2r^2}\sin\phi,\quad
 B_{\phi,\rm out}=\frac{\mu_0i_0R^2}{2r^2}\cos\phi}.
\]

三圆柱相交体的关键是三组电流并非彼此拼凑。取整个体积均匀磁化
$\vect M=i_0\vect e_z$，则任意边界面上
$\vect K=\vect M\times\vect n$。在沿 $x$ 的圆柱面，
$\vect n=\cos\phi_y\vect e_y+\sin\phi_y\vect e_z$，得到
$\vect K_1=-i_0\cos\phi_y\vect e_x$；沿 $y$、$z$ 的圆柱面同理恰得 $\vect K_2,\vect K_3$。因此接缝两侧是同一个连续面流的限制，电荷不在接缝堆积，连续性自动满足。

远场只保留总磁偶极矩
\[
 \vect m=\int\vect M\dd V=i_0V_S\vect e_z.
\]
% CompletenessEM:a4-triple-volume
为计算 $V_S$，把三圆柱公共区域写成
\[
 x^2+y^2\leq R^2,\quad x^2+z^2\leq R^2,\quad y^2+z^2\leq R^2.
\]
先限于第一卦限，再取其中的 $x\geq y$ 部分；其余由三坐标变号和 $x,y$ 互换生成，
总体积是这一部分的 $16$ 倍。对固定的 $x,y$，后两个不等式给出
$0\leq z\leq\sqrt{R^2-x^2}$；投影区域如图~\ref{complete:em:triple-projection}，
$y$ 的上界为 $\min(x,\sqrt{R^2-x^2})$，在 $x=R/\sqrt2$ 处分段。

\begin{qualityfigure}
\begin{tikzpicture}[>=Latex,font=\small,scale=1.05]
 \fill[PaleBlue](0,0)--(1.6971,1.6971)
 arc[start angle=45,end angle=0,radius=2.4]--cycle;
 \draw[->](0,0)--(2.85,0)node[right]{$x$};
 \draw[->](0,0)--(0,2.8)node[above]{$y$};
 \draw[DeepBlue,thick](0,0)--(1.6971,1.6971);
 \draw[DeepBlue,thick](1.6971,1.6971)arc[start angle=45,end angle=0,radius=2.4];
 \draw[TextGray](0,2.4)arc[start angle=90,end angle=45,radius=2.4];
 \draw[dashed](1.6971,0)--(1.6971,1.6971);
 \node[below left]at(1.6971,-.05){$R/\sqrt2$};\node[below right]at(2.4,-.05){$R$};
 \node[above left]at(1,1){$y=x$};
 \node[above right]at(2.05,1.35){$x^2+y^2=R^2$};
 \node at(1.25,.5){$D$};
\end{tikzpicture}
\figcap{积分域 $D$ 是第一卦限且 $x\geq y$ 的 $xy$ 投影；每点竖向高度为 $\sqrt{R^2-x^2}$。\label{complete:em:triple-projection}}
\end{qualityfigure}

因此不用预先记忆相交体积，也可直接积分：
\[
 \begin{aligned}
 V_S&=16\left[\int_0^{R/\sqrt2}x\sqrt{R^2-x^2}\dd x
 +\int_{R/\sqrt2}^{R}(R^2-x^2)\dd x\right]\\
 &=16R^3\left[\frac13-\frac1{6\sqrt2}
 +\frac23-\frac5{6\sqrt2}\right]
 =8(2-\sqrt2)R^3.
 \end{aligned}
\]
第一段来自 $0\leq y\leq x$，第二段来自 $0\leq y\leq\sqrt{R^2-x^2}$；
这就是数值系数中出现 $\sqrt2$ 的几何来源。


三个正交圆柱公共区域的体积为
$V_S=8(2-\sqrt2)R^3$，故
\[
 \boxed{\vect m=8(2-\sqrt2)i_0R^3\vect e_z},
\]
\[
 \boxed{\vect B(\vect r)\sim\frac{\mu_0}{4\pi r^3}
 [3(\vect m\cdot\vect e_r)\vect e_r-\vect m]}
 \qquad(r\gg R).
\]
图中只示三类面流的局部方向，不把椭圆当作三圆柱相交体的边界。各面电流均由同一公式
$\vect M\times\vect n$ 决定，所以方向关系不是装饰性示意，而直接参与连续性证明。

图~\ref{advA4:quality:A4-magnetostatics:a7} 中，面电流只能沿实际边界切向流动；接缝处各面流入、流出量须连续，远区偶极矩要对三类面电流共同积分，不能任意把某一完整圆柱场截断。
\end{advsolution}

\section{电磁场、电磁波与辐射}

\subsection{介质中的守恒定律与平面波}

介质中的宏观 Maxwell 方程为
\[
 \nabla\cdot\vect D=\rho_f,
 \qquad \nabla\times\vect H-\partial_t\vect D=\vect j_f,
\]
\[
 \nabla\cdot\vect B=0,
 \qquad \nabla\times\vect E+\partial_t\vect B=0.
\]
在线性、各向同性、无色散介质中
$\vect D=\varepsilon\vect E$、
$\vect B=\mu\vect H$。无源区内电、磁场都满足波动方程，波速
\[
 v=\frac1{\sqrt{\varepsilon\mu}}
 =\frac c{\sqrt{\varepsilon_r\mu_r}}.
\]

把 $\vect E$ 点乘 Ampere--Maxwell 方程、$\vect H$ 点乘 Faraday 方程并相减，得到 Poynting 定理
\[
 \boxed{\partial_tw+\nabla\cdot\vect S=-\vect j_f\cdot\vect E},
\]
\[
 w=\frac12(\vect D\cdot\vect E+\vect B\cdot\vect H),
 \qquad \vect S=\vect E\times\vect H.
\]
均匀介质中的动量守恒可写成
\[
 \partial_t\vect g+\nabla\cdot\mathsf T
 =-\rho_f\vect E-\vect j_f\times\vect B,
\]
\[
 \vect g=\vect D\times\vect B,
 \qquad
 \mathsf T=w\mathsf1-\vect D\vect E-\vect B\vect H.
\]
这里采用均匀、线性、各向同性、无色散且参数不随时空改变的介质模型，$\vect g$ 是 Minkowski 动量密度，$\mathsf T$ 为动量通量张量，等于通常应力张量的负号。在真空中它退化为 $\varepsilon_0\vect E\times\vect B$。若改用 Abraham 密度 $\vect E\times\vect H/c^2$，就必须把相应介质动量交换计入右侧，不能仅替换左侧密度。

平面波
$\vect E=\vect E_0\e^{\ii(\vect k\cdot\vect r-\omega t)}$ 满足
\[
 \vect k\cdot\vect E_0=\vect k\cdot\vect B_0=0,
 \qquad
 \vect B_0=\frac1\omega\vect k\times\vect E_0,
\]
\[
 \langle\vect S\rangle
 =\frac12\sqrt{\frac\varepsilon\mu}|E_0|^2\vect e_k.
\]
折射率 $n=\sqrt{\varepsilon_r\mu_r}$，Snell 定律为
$n_1\sin\theta_1=n_2\sin\theta_2$。令波阻抗
$Z_i=\sqrt{\mu_i/\varepsilon_i}$，电场振幅的 Fresnel 反射系数为
\[
 r_{\rm TE}=\frac{Z_2\cos\theta_1-Z_1\cos\theta_2}
 {Z_2\cos\theta_1+Z_1\cos\theta_2},
\]
\[
 r_{\rm TM}=\frac{Z_1\cos\theta_1-Z_2\cos\theta_2}
 {Z_1\cos\theta_1+Z_2\cos\theta_2}.
\]
透射振幅由切向 $\vect E,\vect H$ 连续确定；功率系数还必须乘上两侧法向能流因子，不能简单用 $|t|^2$。

导电介质在 $\e^{-\ii\omega t}$ 约定下具有复介电常数
\[
 \varepsilon_c=\varepsilon+\frac{\ii\sigma}{\omega},
 \qquad k^2=\omega^2\mu\varepsilon_c.
\]
良导体 $\sigma\gg\omega\varepsilon$ 时
\[
 k\simeq(1+\ii)\sqrt{\frac{\omega\mu\sigma}{2}},
 \qquad
 \boxed{\delta=\sqrt{\frac2{\omega\mu\sigma}}}
\]
是趋肤深度。

\begin{advexamplebox}{黄昆方程、离子晶体极化与 Lyddane--Sachs--Teller 关系}\label{advA4:ex:huang-kun}
黄昆先生是中国固体物理和半导体物理的重要奠基人之一。他关于晶格动力学、长光学波及声子极化激元的工作影响深远，与 Max Born 合著的《晶格动力学理论》成为这一领域的经典著作。

离子晶体的一个原胞体积为 $\Omega$。正、负离子质量分别为 $M,m$，电荷为 $\pm q$，相对位移为 $x$，谐回复劲度为 $K$。电子云极化满足
\[
 P_e=\alpha\left(E+\frac{P}{3\varepsilon_0}\right).
\]
定义质量归一化坐标 $W$ 使动能密度为 $\dot W^2/2$。
\begin{enumerate}
 \item 求 $W$ 与 $x$ 的关系；
 \item 推导黄昆第二方程 $P=a_{21}W+a_{22}E$；
 \item 推导黄昆第一方程 $\ddot W=-a_{11}W+a_{12}E$，并判断 $a_{12}$ 与 $a_{21}$ 的关系；
 \item 求静态与高频介电常数 $\varepsilon_{\rm st},\varepsilon_\infty$；
 \item 求长波极限的纵、横光学模频率 $\omega_{\rm LO},\omega_{\rm TO}$；
 \item 验证 Lyddane--Sachs--Teller 关系。
\end{enumerate}
\end{advexamplebox}
\begin{advsolution}
约化质量
\[
 \mu_r=\frac{Mm}{M+m}.
\]
质心固定时每原胞动能为 $\mu_r\dot x^2/2$，除以体积后与
$\dot W^2/2$ 比较，得到
\[
 \boxed{W=\sqrt{\frac{\mu_r}{\Omega}}x}.
\]

总极化包含离子位移项与电子项：
\[
 P=\frac{qx}{\Omega}+\alpha\left(E+\frac{P}{3\varepsilon_0}\right).
\]
令 $d=1-\alpha/(3\varepsilon_0)$，则
\[
 \boxed{a_{21}=\frac{q}{\sqrt{\mu_r\Omega}\,d}},
 \qquad
 \boxed{a_{22}=\frac\alpha d}.
\]

相对运动方程
\[
 \mu_r\ddot x=-Kx+q\left(E+\frac{P}{3\varepsilon_0}\right)
\]
化为
\[
 \ddot W=-\left(\frac K{\mu_r}
 -\frac{q a_{21}}{3\varepsilon_0\sqrt{\mu_r\Omega}}\right)W
 +\frac q{\sqrt{\mu_r\Omega}}
 \left(1+\frac{a_{22}}{3\varepsilon_0}\right)E.
\]
因此
\[
 \boxed{a_{11}=\frac K{\mu_r}
 -\frac{q a_{21}}{3\varepsilon_0\sqrt{\mu_r\Omega}}},
 \qquad
 \boxed{a_{12}=\frac q{\sqrt{\mu_r\Omega}\,d}=a_{21}}.
\]
末一等式体现微观可逆性。

取谐振时间因子，
\[
 W=\frac{a_{12}}{a_{11}-\omega^2}E,
\]
所以
\[
 \varepsilon(\omega)=\varepsilon_0+a_{22}
 +\frac{a_{12}^2}{a_{11}-\omega^2}.
\]
从而
\[
 \boxed{\varepsilon_\infty=\varepsilon_0+a_{22}},
 \qquad
 \boxed{\varepsilon_{\rm st}=\varepsilon_\infty+\frac{a_{12}^2}{a_{11}}}.
\]

长波横模没有宏观极化电荷，可取 $E=0$，故
\[
 \boxed{\omega_{\rm TO}^2=a_{11}}.
\]
这里的横光学模采用非迟延的长波晶格近似；若保留电磁场传播，应求声子极化激元色散，不能把其严格 $k\to0$ 极限与上述近似直接等同。

纵模无自由电荷，$D=\varepsilon_0E+P=0$。结合第二方程得
$E=-a_{21}W/\varepsilon_\infty$，代入第一方程：
\[
 \boxed{\omega_{\rm LO}^2=a_{11}+\frac{a_{12}^2}{\varepsilon_\infty}}.
\]
因此
\[
 \boxed{\frac{\omega_{\rm LO}^2}{\omega_{\rm TO}^2}
 =1+\frac{a_{12}^2}{a_{11}\varepsilon_\infty}
 =\frac{\varepsilon_{\rm st}}{\varepsilon_\infty}}.
\]
各 $a_{11},\omega^2$ 量纲相同，且稳定晶体要求
$a_{11}>0$、$\varepsilon_\infty>0$，从而
$\omega_{\rm LO}>\omega_{\rm TO}$。
\end{advsolution}

\subsection{矩形波导模式}

理想导体矩形波导沿 $z$ 方向，截面 $0<x<a$、$0<y<b$。设场依赖为
$\e^{\ii(\beta z-\omega t)}$，定义
\[
 k=\omega\sqrt{\mu\varepsilon},
 \qquad k_c^2=k_x^2+k_y^2,
 \qquad \beta^2=k^2-k_c^2.
\]
导体壁满足切向 $\vect E=0$、法向 $\vect B=0$；表面电荷和电流为
$\sigma_s=-\vect n\cdot\vect D$、
$\vect K_s=-\vect n\times\vect H$（$\vect n$ 从波导内部指向金属）。

\begin{advexamplebox}{为什么不存在 $\mathrm{TE}_{00}$ 模}\label{advA4:ex:no-te00}
证明单连通空心矩形波导不存在非零 $\mathrm{TE}_{00}$ 模，也不存在 TEM 模。
\end{advexamplebox}
\begin{advsolution}
TE 模以 $H_z$ 为母函数，边界条件为 $\partial_nH_z=0$。
$m=n=0$ 时它在截面内是常数。取 $\omega\ne0$，直接在截面使用 Faraday 定律：
\[
 \oint_{\partial S}\vect E_t\cdot\dd\vect l
 =\ii\omega\mu\int_S H_z\,\dd S.
\]
理想导体边界上切向电场为零，故左边为零；右边为 $\ii\omega\mu abH_z$，
所以 $H_z=0$。不能把 $k_c=0$ 直接代入含 $1/k_c^2$ 的横场公式来证明结论。
剩余横场若存在，只能属于 TEM 情形，下面排除它。

TEM 模要求 $E_z=H_z=0$，横向电场可写成二维静电势的梯度，势在整个单一导体边界上为同一常数。Laplace 方程唯一性迫使截面内势也为常数，故 $\vect E_t=0$。TEM 只有在同轴线等多连通截面中才可能存在。
\end{advsolution}

\begin{advexamplebox}{$\mathrm{TE}_{11}$ 与 $\mathrm{TM}_{11}$ 的完整场}\label{advA4:ex:waveguide-11}
写出矩形波导 TE$_{11}$、TM$_{11}$ 的内部电磁场及壁面电荷、电流分布，并给出截止条件。
\end{advexamplebox}
\begin{advsolution}
令 $k_x=\pi/a$、$k_y=\pi/b$、
$k_c^2=k_x^2+k_y^2$。TE$_{11}$ 取
\[
 H_z=H_0\cos k_xx\cos k_yy.
\]
其横向场为
\[
 \boxed{\vect E_t=-\frac{\ii\omega\mu}{k_c^2}
 \vect e_z\times\nabla_tH_z},
 \qquad
 \boxed{\vect H_t=\frac{\ii\beta}{k_c^2}\nabla_tH_z}.
\]
$E_z=0$。余弦保证壁面法向导数为零，因而切向电场为零。

TM$_{11}$ 取
\[
 E_z=E_0\sin k_xx\sin k_yy,
\]
其横向场为
\[
 \boxed{\vect E_t=\frac{\ii\beta}{k_c^2}\nabla_tE_z},
 \qquad
 \boxed{\vect H_t=\frac{\ii\omega\varepsilon}{k_c^2}
 \vect e_z\times\nabla_tE_z}.
\]
$H_z=0$。正弦保证导体壁上 $E_z=0$，而 $\vect E_t$ 的切向分量也为零。

两模的截止角频率相同：
\[
 \boxed{\omega_c=\frac1{\sqrt{\mu\varepsilon}}
 \sqrt{\left(\frac\pi a\right)^2+
 \left(\frac\pi b\right)^2}}.
\]
当 $\omega<\omega_c$ 时 $\beta$ 为纯虚数，模式沿 $z$ 倏逝。各壁面把相应场代入
$\sigma_s=-\vect n\cdot\varepsilon\vect E$、
$\vect K_s=-\vect n\times\vect H$ 即得：TE 模虽 $E_z=0$，仍因法向横场产生交变面电荷；TM 模的纵向电场在壁上为零但法向横场和切向磁场仍分别产生面电荷与面电流。本题单一模式是有限的正弦、余弦场，在矩形四角没有由这些公式产生的发散；不能把其他尖角静场的结论直接套来。
\end{advsolution}

\subsection{推迟势与电磁辐射}

Helmholtz 方程和波动方程的外向 Green 函数分别为
\[
 G_+(\vect r,\vect r')=\frac{\e^{\ii kR}}{4\pi R},
 \qquad
 G_{\rm ret}(t,\vect r;t',\vect r')
 =\frac{\delta(t-t'-R/c)}{4\pi R},
\]
其中 $R=|\vect r-\vect r'|$。因果性选择推迟解，四维势为
\[
 A^\nu(t,\vect r)=\frac{\mu_0}{4\pi}\int
 \frac{j^\nu(t-R/c,\vect r')}{R}\dd V'.
\]

\begin{advexamplebox}{电偶极辐射与 Larmor 公式}\label{advA4:ex:electric-dipole-radiation}
先证明局域电流体系满足
\[
 \dot{\vect p}=\int\vect j\dd V,
 \qquad \vect p=\int\rho\vect r\dd V.
\]
再求谐振偶极 $\vect p(t)=\vect p_0\cos\omega t$ 的远区平均能流、总功率，并推出非相对论点电荷的 Larmor 公式。
\end{advexamplebox}
\begin{advsolution}
由连续性方程和分部积分，
\[
 \dot p_i=\int x_i\partial_t\rho\dd V
 =-\int x_i\partial_jj_j\dd V
 =\int j_i\dd V,
\]
边界项因电流局域而消失。

远区 $r$ 远大于源尺度时
\[
 \vect A\simeq\frac{\mu_0}{4\pi r}
 \dot{\vect p}(t-r/c).
\]
只保留 $1/r$ 辐射项，
\[
 \vect E_{\rm rad}=\frac1{4\pi\varepsilon_0c^2r}
 \vect e_r\times(\vect e_r\times\ddot{\vect p}),
 \qquad
 \vect B_{\rm rad}=\frac1c\vect e_r\times\vect E_{\rm rad}.
\]
设 $\theta$ 是 $\vect p_0$ 与观察方向夹角，时间平均能流为
\[
 \boxed{\langle\vect S\rangle=
 \frac{\omega^4p_0^2\sin^2\theta}
 {32\pi^2\varepsilon_0c^3r^2}\vect e_r}.
\]
对球面积分，$\int\sin^2\theta\dd\Omega=8\pi/3$，得到
\[
 \boxed{\langle P\rangle=\frac{\omega^4p_0^2}
 {12\pi\varepsilon_0c^3}}.
\]
一般随时间变化的偶极瞬时辐射功率为
$P=|\ddot{\vect p}|^2/(6\pi\varepsilon_0c^3)$。点电荷
$\vect p=q\vect r$，所以
\[
 \boxed{P_{\rm Larmor}=\frac{q^2a^2}{6\pi\varepsilon_0c^3}}.
\]
量纲为能量/时间，且 $a=0$ 时无辐射。
\end{advsolution}

\begin{advexamplebox}{双星引力辐射与啁啾质量}\label{advA4:ex:binary-chirp}
引力波概念由 Poincar\'e 在1905年前后提出，Einstein 随后建立广义相对论中的波动理论。Hulse--Taylor 脉冲双星的轨道衰减为引力辐射提供了间接证据；LIGO 对双黑洞并合的直接探测使 Rainer Weiss、Barry C. Barish、Kip S. Thorne 获得2017年 Nobel 物理学奖。

质量 $m_1,m_2$ 的双星作圆轨道运动，间距为 $R$。求四极辐射功率；再定义引力波频率 $f=\Omega/\pi$，推导 $\dot f$ 与啁啾质量的关系，并确定
\[
 \mathcal M=\frac{m_1^\alpha m_2^\alpha}{(m_1+m_2)^\beta}
\]
中的指数。
\end{advexamplebox}
\begin{advsolution}
记总质量 $M=m_1+m_2$、约化质量
$\mu=m_1m_2/M$。圆轨道满足
\[
 \Omega^2=\frac{GM}{R^3}.
\]
质量四极矩的三阶时间导数代入引力四极辐射公式，得到
\[
 \boxed{P_{\rm GW}=\frac{32}{5}\frac{G}{c^5}
 \mu^2R^4\Omega^6
 =\frac{32}{5}\frac{G^4}{c^5}
 \frac{\mu^2M^3}{R^5}}.
\]
轨道机械能
\[
 E=-\frac{G\mu M}{2R}
 =-\frac12\mu(GM)^{2/3}\Omega^{2/3}.
\]
能量平衡 $-\dd E/\dd t=P_{\rm GW}$ 给出
\[
 \dot\Omega=\frac{96}{5}\frac{G^{5/3}}{c^5}
 \mu M^{2/3}\Omega^{11/3}.
\]
定义
\[
 \boxed{\mathcal M=\mu^{3/5}M^{2/5}
 =\frac{(m_1m_2)^{3/5}}{(m_1+m_2)^{1/5}}},
\]
所以
\[
 \boxed{\alpha=\frac35,\qquad\beta=\frac15}.
\]
由于 $f=\Omega/\pi$，
\[
 \boxed{\dot f=\frac{96}{5}\pi^{8/3}
 \left(\frac{G\mathcal M}{c^3}\right)^{5/3}f^{11/3}}.
\]
因此质量指数为 $\gamma=5/3$、频率指数为
$\delta=11/3$、无量纲常数
$\lambda=(96/5)\pi^{8/3}$。右端量纲为
$T^{5/3}T^{-11/3}=T^{-2}$，与 $\dot f$ 一致；$P>0$ 使
$R$ 变小、$f$ 增大，方向判断一致。
\end{advsolution}


\clearpage\endgroup

\begingroup
\newif\ifAdvHideHeading
\RenewDocumentEnvironment{advexamplebox}{m}{%
  \refstepcounter{advexample}%
  \begin{tcolorbox}[
    enhanced,unbreakable,colback=white,colframe=RuleGray,boxrule=0.5pt,arc=0mm,
    title={例~\theadvexample\quad #1},fonttitle=\bfseries,
    coltitle=Ink,colbacktitle=PaleBlue,borderline west={1.4pt}{0pt}{DeepBlue}]
}{\end{tcolorbox}}

\RenewDocumentEnvironment{advsolution}{+b}{%
  \par\smallskip\begingroup\small\noindent
  \ifAdvHideHeading\global\AdvHideHeadingfalse\else\LectureSolutionHeading\fi #1\par\endgroup\smallskip
}{}

\RenewDocumentEnvironment{advdefinition}{m +b}{%
  \refstepcounter{advdefinition}%
  \begin{tcolorbox}[
    enhanced,unbreakable,colback=white,colframe=RuleGray,boxrule=0.5pt,arc=0mm,
    title={定义~\theadvdefinition\quad #1},fonttitle=\bfseries,
    coltitle=Ink,colbacktitle=PaleGold,borderline west={1.2pt}{0pt}{Gold}]
  #2
  \end{tcolorbox}
}{}

\RenewDocumentEnvironment{advtheorem}{m +b}{%
  \refstepcounter{advtheorem}%
  \begin{tcolorbox}[
    enhanced,unbreakable,colback=white,colframe=RuleGray,boxrule=0.5pt,arc=0mm,
    title={定理~\theadvtheorem\quad #1},fonttitle=\bfseries,
    coltitle=Ink,colbacktitle=PaleBlue,borderline west={1.2pt}{0pt}{DeepBlue}]
  #2
  \end{tcolorbox}
}{}


\AdvancedUnit{5}{光学}
\section{几何光学}

\subsection{光线方程}

在折射率随位置连续变化的介质中，把光的局域传播方向记为单位切向量
$\vect\tau$，弧长记为 $s$。光线垂直于等相位面，满足
\begin{equation}\label{advA5:eq:ray}
 \boxed{\frac{\dd}{\dd s}(n\vect\tau)=\nabla n}.
\end{equation}
取 Frenet 标架 $\vect\tau,\vect\nu$，并记曲率
$\kappa=\dd\theta/\dd s$，式 \eqref{advA5:eq:ray} 的切向和法向分量分别为
\[
 \frac{\dd n}{\dd s}=(\nabla n)\cdot\vect\tau,\qquad
 n\kappa=(\nabla n)\cdot\vect\nu .
\]
前式是链式法则，后式则说明光线总向折射率增大的方向弯曲。图
\ref{advA5:fig:ray-equation} 把局部折射率分层、切向、法向和梯度放在同一幅图中；
其中所有标签都避开了曲线。

\begin{center}\begin{minipage}{\linewidth}\centering
\centering
\begin{tikzpicture}[scale=.9]
 \foreach \y in {-.5,.4,1.3}{\draw[RuleGray,dashed] (-3,\y)--(3.3,\y);}
 \draw[curve,domain=-2.7:2.7,samples=90] plot (\x,{.4+.13*\x*\x});
 \coordinate(P)at(0,.4);\fill(P)circle(2pt)node[below left]{$P$};
 \draw[vector] (P)--(1.4,.4) node[below]{$\vect\tau$};
 \draw[vector] (P)--(0,1.55) node[left]{$\vect\nu$};
 \draw[vector,Gold!80!black] (1.9,.4)--(1.9,1.55) node[above]{$\nabla n$};
 \node[right] at (3.35,1.3){等 $n$ 线};\node at (-1.6,1.9){光线向高折射率侧弯曲};
\end{tikzpicture}
\bookcaption{光线方程的局域几何：$\nabla n$ 的法向分量决定曲率}
\label{advA5:fig:ray-equation}
\end{minipage}\end{center}

若介质具有固定法线方向的平行分层（例如 $n=n(y)$），$\vect\tau$ 与 $\nabla n$ 夹角为
$\varphi$，则
\[
 \dd n=|\nabla n|\cos\varphi\,\dd s,\qquad
 \dd(n\sin\varphi)=0,
\]
因而
\[
 \left|\frac{\dd\varphi}{\dd s}\right|
 =\frac{|\nabla n|}{n}\sin\varphi,\qquad
 n\frac{\dd\theta}{\dd s}=(\nabla n)\cdot\vect\nu,
\]
这里梯度方向固定，上式才可由角度 $\varphi$ 直接积分。对一般 $n(\vect r)$，$\nabla n$ 自身方向也沿途改变，不能声称 $n\sin\varphi$ 守恒；球对称介质相应的不变量为 $nr\sin\varphi$。

二维情形若 $n=n(y)$，取 $\theta$ 为光线与 $y$ 轴的夹角，则平移对称性给出
\[
 \boxed{n(y)\sin\theta=\text{常量}}.
\]
它是分层介质 Snell 定律的连续极限。已知初始点与初始方向后，
$\dd x/\dd y=\tan\theta(y)$ 便给出轨迹。

\begin{advexamplebox}{费马原理与光线方程的等价性}\label{advA5:ex:fermat-ray}
从光程泛函
\[
 \mathcal S[\vect r]=\int_A^B n(\vect r)\,\dd s
\]
出发，证明其驻值条件与式 \eqref{advA5:eq:ray} 等价。
\end{advexamplebox}
\begin{advsolution}
以任意参数 $\lambda$ 描述路径，令
$\dot{\vect r}=\dd\vect r/\dd\lambda$，则
\[
 \mathcal S=\int n(\vect r)|\dot{\vect r}|\,\dd\lambda .
\]
Euler--Lagrange 方程为
\[
 \frac{\dd}{\dd\lambda}
 \left(n\frac{\dot{\vect r}}{|\dot{\vect r}|}\right)
 -|\dot{\vect r}|\,\nabla n=0.
\]
利用 $\vect\tau=\dot{\vect r}/|\dot{\vect r}|$ 以及
$\dd s=|\dot{\vect r}|\dd\lambda$，约去 $|\dot{\vect r}|$ 得
\[
 \frac{\dd}{\dd s}(n\vect\tau)=\nabla n.
\]
反过来，若每段光线满足此方程，乘以任意端点固定的变分
$\delta\vect r$ 并分部积分，边界项为零，得到
$\delta\mathcal S=0$。两者等价。量纲上两边均为
$\text{长度}^{-1}$；均匀介质 $\nabla n=0$ 时
$\vect\tau$ 不变，退化为直线传播。
\end{advsolution}

\subsection{光力类比与变折射率介质}

把质量取为 $1$ 的辅助粒子速度规定为 $\vect v=n\vect\tau$，则
\[
 \frac{\dd\vect v}{\dd t}=n\nabla n
 =-\nabla\!\left(-\frac{n^2}{2}\right).
\]
因此光线轨迹等价于辅助粒子在
\[
 V(\vect r)=-\frac{n^2(\vect r)}2,\qquad
 E=\frac{v^2}2+V=0
\]
中的零能量轨迹。这里的 $v=n$ 只是一种重新参数化，绝不是介质中的真实光速
$c/n$。平移对称性对应 $n\sin\theta$ 守恒，球对称性对应
$nr\sin\theta$ 守恒。

\begin{advexamplebox}{浅水波靠岸的轨迹}\label{advA5:ex:beach-ray}
海岸线为 $x$ 轴，水深 $H(y)=ay$，$y>0$。在浅水近似下
$c(y)=\sqrt{gH(y)}$。求一条波线的轨迹，并说明它如何靠岸。
\end{advexamplebox}
\begin{advsolution}
定频波的等效折射率 $n\propto1/c\propto y^{-1/2}$。
沿海岸方向平移对称，故
\[
 n\sin\theta=C,\qquad
 \sin\theta=\sqrt{\frac{y}{y_m}},
\]
其中 $y_m$ 是轨迹在 $\theta=\pi/2$ 处的最深位置。于是
\[
 \frac{\dd x}{\dd y}=\tan\theta
 =\sqrt{\frac{y}{y_m-y}}.
\]
积分得到
\[
 \boxed{x-x_c
 =y_m\arcsin\sqrt{\frac{y}{y_m}}
 -\sqrt{y(y_m-y)}}.
\]
令 $y=y_m\sin^2u$，可写成摆线参数式
\[
 y=\frac{y_m}{2}(1-\cos2u),\qquad
 x-x_c=y_m\!\left(u-\frac12\sin2u\right).
\]
图中靠岸端 $y\to0$ 时 $\theta\to0$，波线垂直海岸入射，符合海浪近岸转向的物理图像。

\begin{center}
\begin{tikzpicture}[>=Latex,scale=.75]
 \fill[DeepBlue!7] (-.4,0) rectangle (7,3.25);
 \draw[thick] (-.4,0)--(7,0);
 \node[below] at (5.8,0) {海岸};
 \draw[->] (0,0)--(0,3.45) node[above]{$y$};
 \draw[->] (0,0)--(7.2,0) node[right]{$x$};
 \draw[DeepBlue,very thick,->,domain=.06:3.05,samples=80]
   plot ({2.1*(\x-.5*sin(2*\x r))},{2.1*sin(\x r)^2});
 \node[TextGray] at (5.35,2.65) {$H(y)=ay$};
\end{tikzpicture}
\end{center}
\end{advsolution}

\begin{advexamplebox}{构造全方向无像差成像的折射率场}\label{advA5:ex:maxwell-fisheye}
构造一个球对称折射率场，使任一点发出的全部光线在另一点重新会聚。
\end{advexamplebox}
\begin{advsolution}
取 Maxwell 鱼眼分布
\[
 \boxed{n(r)=\frac{2n_0}{1+(r/a)^2}}.
\]
作球极投影 $r=a\tan(\chi/2)$，平面线元满足
\[
 n^2(r)(\dd r^2+r^2\dd\phi^2)
 =n_0^2a^2(\dd\chi^2+\sin^2\chi\,\dd\phi^2).
\]
右端正是半径 $n_0a$ 的球面度量，所以平面中的光线对应球面大圆。
任一点在球面上的对跖点由所有经过该点的大圆共同通过，投影回平面便得到唯一像点。
若物点为复坐标 $z_0$，像点为
\[
 \boxed{z_i=-\frac{a^2}{\bar z_0}}.
\]
等光程来自球面上两对跖点间大圆弧长恒为 $\pi n_0a$。
在 $r\ll a$ 的局部，$n\approx2n_0$，光线近似直线；在 $r\to\infty$ 时
$n\to0$，这说明理想分布需要有限边界或镜面截断才能工程实现。
\end{advsolution}

\begin{advexamplebox}{螺旋光线、出射双曲面与透镜束腰}\label{advA5:ex:helical-fiber}
圆柱光纤中，一条光线在半径 $r$ 的薄层上以螺距 $h$ 沿 $z$ 方向传播，
该层折射率为 $n$。

\begin{enumerate}
\item 求该处 $\dd n/\dd r$；
\item 光线从端面 $z=0$ 射入真空，求所有出射线构成的包络曲面；
\item 端面加焦距 $f$ 的薄凸透镜，在出射光满足近轴条件时求束腰半径和束腰到端面的距离。
\end{enumerate}

\begin{center}
\begin{tikzpicture}[>=Latex,scale=.72]
 \draw[RuleGray] (-2,-1.25) ellipse (1.15 and .38);
 \draw[RuleGray] (-2,2.1) ellipse (1.15 and .38);
 \draw[RuleGray] (-3.15,-1.25)--(-3.15,2.1) ( -0.85,-1.25)--(-.85,2.1);
 \draw[DeepBlue,thick,domain=0:720,samples=150]
   plot ({-2+1.15*cos(\x)},{-1.1+.0042*\x+.18*sin(\x)});
 \draw[->] (-2,-1.4)--(-2,2.75) node[above]{$z$};
 \draw[<->] (-2,-.82)--(-.85,-.82) node[midway,above]{$r$};
 \draw[<->] (-3.55,-1.08)--(-3.55,.44) node[midway,left]{$h$};
 \draw[->,DeepBlue] ({-2+1.15*cos(675)},{-1.1+.0042*675+.18*sin(675)})--++(.7,.316698) node[right]{$\vect\tau$};
\end{tikzpicture}
\end{center}
\end{advexamplebox}
\begin{advsolution}
令 $b=h/(2\pi)$。螺线可写为
$\vect r(\phi)=(r\cos\phi,r\sin\phi,b\phi)$，其曲率
\[
 \kappa=\frac{r}{r^2+b^2},
\]
主法向沿径向内侧。光线方程的法向分量给出
\[
 \boxed{\frac{\dd n}{\dd r}=-n\kappa
 =-\frac{nr}{r^2+b^2}}.
\]
负号表示折射率向轴心增大，才能提供向心弯曲。

螺线相对端面法线的入射角满足
$\sin\theta_i=r/\sqrt{r^2+b^2}$。由 Snell 定律
\[
 q\equiv\tan\theta_t,\qquad
 \sin\theta_t=\frac{nr}{\sqrt{r^2+b^2}},
\]
这要求右端不超过 $1$；否则发生全反射，题设的出射不成立。
出射线的横向分量沿圆周切向，故传播距离 $z$ 后
\[
 \boxed{\rho^2=r^2+q^2z^2},
\]
是一叶双曲面的外侧轮廓。

薄透镜只增加径向会聚斜率 $-r/f$，不改变切向斜率 $q$。因此
\[
 \rho^2(z)=\left(r-\frac{rz}{f}\right)^2+q^2z^2.
\]
令其对 $z$ 的导数为零，得到
\[
 \boxed{z_w=\frac{r^2f}{r^2+q^2f^2}},\qquad
 \boxed{\rho_w=\frac{rqf}{\sqrt{r^2+q^2f^2}}}.
\]
$q\to0$ 时回到普通聚焦 $z_w=f,\rho_w=0$；$f\to\infty$ 时
$z_w\to0,\rho_w\to r$，极限正确。
\end{advsolution}

\subsection{成像法则与傍轴系统}

点光源 $S$ 发出的全部光线组成物方光场；若一族真实光线在另一点 $S'$ 再次相交，
则 $S'$ 为像点。由费马原理，同一连续光线族的光程相等，这就是物像间的等光程原理。

\begin{center}\begin{minipage}{\linewidth}\centering
\centering
\begin{tikzpicture}[>=Latex,scale=.78]
 \coordinate (S) at (-3.4,0); \coordinate (I) at (3.4,0);
 \foreach \h in {-2,-1.25,-.55,.55,1.25,2}{
   \draw[DeepBlue!85,->] (S) .. controls (-1.2,\h) and (1.2,\h) .. (I);
 }
 \fill (S) circle(1.5pt) node[left]{$S$};
 \fill (I) circle(1.5pt) node[right]{$S'$};
 \coordinate (P) at (0,1.45);
 \fill (P) circle(1.3pt) node[above]{$P$};
\end{tikzpicture}
\bookcaption{物点与像点：连续光线族在两端共同相交且光程相等}
\label{advA5:fig:object-image}
\end{minipage}\end{center}

实际光具组通常在主光轴附近工作，满足傍轴条件
$|\theta|\ll1$ 或横向尺度与轴向尺度之比远小于 $1$。入射光发散时称实物，
入射光会聚时称虚物；出射光会聚形成实像，出射光发散而反向延长相交形成虚像。

\begin{center}\begin{minipage}{\linewidth}\centering
\centering
\begin{tikzpicture}[scale=.82]
 \begin{scope}
  \draw[axis](-2.7,0)--(2.7,0);\draw[very thick](0,-1.3)--(0,1.3)node[above]{系统};
  \foreach \h in {-.8,.8}{\draw[DeepBlue,->](-2.4,{2.2*\h})--(0,\h);\draw[DeepBlue,->](0,\h)--(2.4,\h);\draw[guide](0,\h)--(2,0);}
  \fill(2,0)circle(1.6pt)node[below]{$S$};\node at (0,-2){虚物：入射线向前延长相交};
 \end{scope}
 \begin{scope}[xshift=6.8cm]
  \draw[axis](-2.7,0)--(2.7,0);\draw[very thick](0,-1.3)--(0,1.3)node[above]{系统};
  \foreach \h in {-.65,.65}{\draw[DeepBlue,->](-2.4,\h)--(0,\h);\draw[DeepBlue,->](0,\h)--(2.4,{2.2*\h});\draw[guide](0,\h)--(-2,0);}
  \fill(-2,0)circle(1.6pt)node[below]{$S'$};\node at (0,-2){虚像：出射线向后延长相交};
 \end{scope}
\end{tikzpicture}
\bookcaption{实物、虚物、实像和虚像由光束的发散或会聚状态判定}
\label{advA5:fig:real-virtual}
\end{minipage}\end{center}

考虑曲率半径为 $R$ 的球面折射，物方、像方折射率分别为 $n,n'$。
按“实物距、实像距为正；从出射方看曲率中心所在侧 $R>0$”的约定，
傍轴展开等光程至二阶得到
\begin{equation}\label{advA5:eq:spherical-refraction}
 \boxed{\frac ns+\frac{n'}{s'}=\frac{n'-n}{R}},
 \qquad
 \boxed{\beta=\frac{y'}y=-\frac{ns'}{n's}}.
\end{equation}
反射可视为令像方传播方向反转；在统一符号约定下球面镜满足
$1/s+1/s'=2/R$。

\begin{center}\begin{minipage}{\linewidth}\centering
\centering
\begin{tikzpicture}[scale=.9]
\begin{scope}
\draw[axis](-3.3,0)--(3.5,0);\draw[thick,domain=-65:65,variable=\t] plot ({1-cos(\t)},{sin(\t)});
\fill(-3,0)circle(1.4pt)node[below]{$S$};\fill(2.771218,0)circle(1.4pt)node[below]{$S'$};
\draw[DeepBlue,->](-3,0)--(0.063250,0.35)--(2.771218,0);
\draw[guide](0.063250,0.35)--(1,0);\fill(1,0)circle(1.3pt)node[below]{$C$};
\node at (-1.8,1){$n=1$};\node at (2,1){$n'=2$};\node at (0,-1.3){球面折射};
\end{scope}
\begin{scope}[xshift=7.1cm]
\draw[axis](-3.4,0)--(.8,0);\draw[very thick,domain=-38:38,variable=\t] plot ({-2+2*cos(\t)},{2*sin(\t)});
\fill(-3,0)circle(1.4pt)node[below]{$S$};\fill(-1.486050,0)circle(1.4pt)node[above left]{$S'$};
\draw[DeepBlue,->](-3,0)--(-0.108572,0.65)--(-1.486050,0);
\draw[guide](-2,0)--(-0.108572,0.65);\fill(-2,0)circle(1.3pt)node[below]{$C$};
\node at (-1.3,-1.3){球面反射：法线通过 $C$};
\end{scope}\end{tikzpicture}
\bookcaption{球面折射（左）与球面反射（右）的物像几何和符号}
\label{advA5:fig:spherical-imaging}
\end{minipage}\end{center}

定义光焦度
\[
 \Phi=\frac{n'-n}{R},\qquad
 f=\frac n\Phi,\qquad f'=\frac{n'}\Phi ,
\]
于是
\[
 \frac fs+\frac{f'}{s'}=1,\qquad f:f'=n:n'.
\]
若改用从焦点量起的 Newton 坐标 $x=s-f$、$x'=s'-f'$，则
\[
 \boxed{xx'=ff'},\qquad
 \boxed{\beta=-\frac fx=-\frac{x'}{f'}}.
\]

\begin{center}\begin{minipage}{\linewidth}\centering
\centering
\begin{tikzpicture}[scale=.9]
 \draw[axis](-4.2,0)--(3.6,0);
 \filldraw[fill=DeepBlue!8,thick] (0,-1.3) .. controls (-0.22,-0.65) and (-0.22,0.65) .. (0,1.3) .. controls (0.22,0.65) and (0.22,-0.65) .. cycle;
 \foreach \x/\lab in {-2/F,0/O,2/F'}{\fill(\x,0)circle(1.5pt)node[below=4pt]{$\lab$};}
 \foreach \h in {-.8,.8}{\draw[DeepBlue,->](-4,\h)--(0,\h)--(3,{-.5*\h});}
 \draw[<->](-2,1.65)--(0,1.65)node[midway,above]{$f$};
 \draw[<->](0,1.65)--(2,1.65)node[midway,above]{$f'$};
 \node at (-.4,-1.7){同一介质中的薄会聚透镜：$f=f'$};
\end{tikzpicture}
\bookcaption{焦点、光心与物像方焦距；焦距符号服从实正虚负}
\label{advA5:fig:focal-points}
\end{minipage}\end{center}

薄透镜可看成两个密接球面，光焦度直接相加：
\[
 \Phi=\Phi_1+\Phi_2
 =\frac{n_L-n}{R_1}+\frac{n'-n_L}{R_2}.
\]
在空气中、采用几何曲率半径的常用符号，薄透镜制造者公式为
$1/f=(n_L-1)(1/R_1-1/R_2)$。

\begin{center}\begin{minipage}{\linewidth}\centering
\centering
\begin{tikzpicture}[scale=.95]
 \draw[axis](-4.4,0)--(3.1,0);\filldraw[fill=DeepBlue!8,thick] (0,-1.3) .. controls (-0.22,-0.65) and (-0.22,0.65) .. (0,1.3) .. controls (0.22,0.65) and (0.22,-0.65) .. cycle;
 \draw[very thick,->](-4,0)--(-4,1)node[above]{$y$};
 \draw[very thick,->](2.4,0)--(2.4,-.6)node[below]{$y'$};
 \draw[DeepBlue,->](-4,1)--(0,1)--(2.4,-.6);
 \draw[Gold!80!black,->](-4,1)--(0,-.6)--(2.4,-.6);
 \draw[TextGray,->](-4,1)--(2.4,-.6);
 \foreach \x/\lab in {-1.5/F,0/O,1.5/F'}{\fill(\x,0)circle(1.4pt)node[below=4pt]{$\lab$};}
\end{tikzpicture}
\bookcaption{薄透镜的三条傍轴特殊光线}
\label{advA5:fig:thin-lens}
\end{minipage}\end{center}

若两个等效薄透镜相距 $d$，定义其相邻焦点间光学距离
$\Delta=d-f_1'-f_2$，则复合系统的等效焦距为
\[
 f=-\frac{f_1f_2}{\Delta},\qquad
 f'=-\frac{f_1'f_2'}{\Delta}.
\]
其主点一般不在任何一个透镜处；从主点量起，Gaussian 与 Newton 公式仍保持不变。

\begin{center}\begin{minipage}{\linewidth}\centering
\centering
\begin{tikzpicture}[x=1.65cm,y=1cm]
 \draw[axis](-2.5,0)--(3.5,0);
 \draw[very thick](0,-1.25)--(0,1.3)node[above]{$O_1$};
 \draw[very thick](1,-1.25)--(1,1.3)node[above]{$O_2$};
 \foreach \x/\lab in {-2/F_1,2/F'_1,-1/F_2,3/F'_2}{\fill(\x,0)circle(1.4pt)node[below]{$\lab$};}
 \draw[guide](.666667,-1)--(.666667,1)node[above]{$H$};
 \draw[guide](.333333,-1)--(.333333,1)node[below]{$H'$};
 \draw[DeepBlue,->](-2.4,.8)--(0,.8)--(1,.4)--(2,-.2);
 \draw[<->](0,-1.65)--(1,-1.65)node[midway,below]{$d$};
 \node at (.5,1.95){示例：$f_1=f_2=2d$，主面由矩阵确定};
\end{tikzpicture}
\bookcaption{分离透镜组的光学距离、焦点与等效主面}
\label{advA5:fig:lens-group}
\end{minipage}\end{center}

\subsection{理想成像系统与基点}

\paragraph{用同一套矩阵核对光线和基点。}
取光线列向量 $(y,nu)^{\mathsf T}$，其中 $u$ 为相对正光轴的小倾角；传播距离 $\ell$ 与光焦度 $\Phi$ 的薄界面分别为
\[
 T(\ell,n)=\begin{pmatrix}1&\ell/n\\0&1\end{pmatrix},
 \qquad R(\Phi)=\begin{pmatrix}1&0\\-\Phi&1\end{pmatrix}.
\]
依照光实际通过的次序从右往左相乘。系统矩阵 $M=(\begin{smallmatrix}A&B\\C&D\end{smallmatrix})$ 的行列式为1，因而保存两条光线的拉格朗日不变量；共轭面之间 $B=0$，横向放大率为 $A$。
空气中两薄透镜相距 $d$ 时，$M=R(1/f_2)T(d,1)R(1/f_1)$，故
\[
 C=-\left(\frac1{f_1}+\frac1{f_2}-\frac d{f_1f_2}\right),
 \qquad f_{\rm eff}=-\frac1C.
\]
若 $d=f_1+f_2$，则 $C=0$，平行束入射后仍平行，是无焦系统，通常有限主面/有限焦距的构造不再适用。画图前应先据矩阵求焦点和主面，再选两条独立光线；焦点本身不是折射界面，光线不能在那里凭空拐折。


理想成像系统是扩充空间中的可逆点映射，保主光轴、绕主光轴对称，并把直线和平面仍映成直线和平面。
因此垂直主轴的平面彼此共轭，其间横向映射是统一放缩。除望远系统外，可用两对基点确定：
横向放大率为 $1$ 的主点 $H,H'$，以及物方焦点 $F$ 和像方焦点 $F'$。

\begin{center}\begin{minipage}{\linewidth}\centering
\centering
\begin{tikzpicture}[scale=.86]
 \draw[axis](-4.6,0)--(4.6,0);
 \foreach \x/\lab in {-2.5/F,-1/H,1/H',2.5/F'}{\draw[RuleGray](\x,-1.4)--(\x,1.4);\node[above] at(\x,1.4){$\lab$};}
 \draw[DeepBlue,->](-4.2,.7)--(-1,.7);\draw[guide,DeepBlue](-1,.7)--(1,.7);\draw[DeepBlue,->](1,.7)--(4,-.7);
 \draw[Gold!80!black,->](-4,-.6)--(-1,.6);\draw[guide,Gold!80!black](-1,.6)--(1,.6);\draw[Gold!80!black,->](1,.6)--(4.2,.6);
 \node at(0,-1.85){主面间虚线为等效对应，非真实内部光路};
\end{tikzpicture}
\bookcaption{主面和焦面决定的特殊光线}
\label{advA5:fig:special-rays}
\end{minipage}\end{center}

从图 \ref{advA5:fig:special-rays} 可见：入射到物方主面高度 $y$ 的光线从像方主面同高处出射；
平行光在像方焦面会聚，物方焦面上一点发出的光出射后平行。由相似三角形立刻得到
$xx'=ff'$ 与 $\beta=-f/x=-x'/f'$。

\begin{center}\begin{minipage}{\linewidth}\centering
\centering
\begin{tikzpicture}[scale=.86]
 \draw[axis](-5.1,0)--(3.8,0);
 \foreach \x/\lab in {-2.2/F,-.7/H,.7/H',2.2/F'}{\draw[RuleGray](\x,-1)--(\x,1.5);\node[above] at(\x,1.5){$\lab$};}
 \draw[very thick,->](-4.7,0)--(-4.7,1)node[above]{$S$};
 \draw[very thick,->](3.1,0)--(3.1,-.6)node[below right]{$S'$};
 \draw[DeepBlue,->](-4.7,1)--(-.7,1);\draw[guide,DeepBlue](-.7,1)--(.7,1);\draw[DeepBlue,->](.7,1)--(3.1,-.6);
 \draw[Gold!80!black,->](-4.7,1)--(-.7,-.6);\draw[guide,Gold!80!black](-.7,-.6)--(.7,-.6);\draw[Gold!80!black,->](.7,-.6)--(3.1,-.6);
 \draw[<->](-4.7,-1.35)--(-2.2,-1.35)node[midway,below]{$x$};
 \draw[<->](2.2,-1.35)--(3.1,-1.35)node[midway,below]{$x'$};
\end{tikzpicture}
\bookcaption{共轭物像点与 Newton 坐标的几何推导}
\label{advA5:fig:conjugate-points}
\end{minipage}\end{center}

角放大率定义为 $B=\tan u'/\tan u$，傍轴下为 $u'/u$。拉格朗日--亥姆霍兹不变量给出
\[
 nyu=n'y'u',\qquad
 \boxed{B=\frac{n}{n'\beta}}.
\]

\begin{center}\begin{minipage}{\linewidth}\centering
\centering
\begin{tikzpicture}[scale=.95]
 \draw[axis](-4.6,0)--(3.1,0);\draw[very thick](0,-1.4)--(0,1.4)node[above]{薄透镜};
 \foreach \h in {-.9,.9}{\draw[DeepBlue,->](-4,0)--(0,\h)--(2.4,0);}
 \fill(-4,0)circle(1.5pt)node[below]{$S$};\fill(2.4,0)circle(1.5pt)node[below]{$S'$};
 \draw[Gold!80!black](-2.65,0)arc[start angle=0,end angle=12.6804,radius=1.35];
 \node at(-2.30,.36){$u$};
 \draw[Gold!80!black](1.3,0)arc[start angle=180,end angle=159.444,radius=1.1];
 \node at(.85,.4){$|u'|$};
 \node at(-.8,-1.7){$u,u'$ 按有向光线对主轴的小倾角定义};
\end{tikzpicture}
\bookcaption{角放大率由共轭光锥的顶角比定义}
\label{advA5:fig:angular-magnification}
\end{minipage}\end{center}

除焦点外，三对常用基点为
\[
\begin{array}{c|c|c}
\text{基点}&(x,x')&\text{特性}\\ \hline
\text{主点}&(-f,-f')&\beta=1\\
\text{节点}&(-f',-f)&B=1\\
\text{二倍焦距点}&(f,f')&\beta=-1
\end{array}
\]
当物像方折射率相同，$f=f'$，主点与节点重合。

\begin{center}\begin{minipage}{\linewidth}\centering
\centering
\begin{tikzpicture}[scale=.87]
 \draw[axis](-5,0)--(5.2,0);
 \foreach \x/\lab/\pos in {-3.5/F/below,0/N/above,-1/H/below,1/H'/below,2/N'/above,4.5/F'/below}{\fill(\x,0)circle(1.5pt)node[\pos=4pt]{$\lab$};}
 \draw[RuleGray](-1,-1.2)--(-1,1.4)(1,-1.2)--(1,1.4);
 \draw[DeepBlue,->](-4,-.8)--(-1,-.2);\draw[guide,DeepBlue](-1,-.2)--(1,-.2);\draw[DeepBlue,->](1,-.2)--(4,.4);
 \draw[guide](1,-.2)--(2,0);
 \draw[<->](-3.5,1.8)--(-1,1.8)node[midway,above]{$f$};
 \draw[<->](1,1.8)--(4.5,1.8)node[midway,above]{$f'$};
 \node at(0,-1.65){指向 $N$ 的入射线与由 $N'$ 出射的光线平行};
\end{tikzpicture}
\bookcaption{焦点、主点和节点三对关键基点}
\label{advA5:fig:cardinal-points}
\end{minipage}\end{center}

空气中的厚透镜，折射率为 $n$、中心厚度为 $d$，常用制造者公式为
\[
 \Phi=\frac1f=(n-1)\left[
 \frac1{R_1}-\frac1{R_2}+\frac{(n-1)d}{nR_1R_2}\right],
\]
其中 $R_1,R_2$ 按几何曲率正负取值。正负并不由外形名称单独决定，而由 $\Phi$ 决定。

\begin{center}\begin{minipage}{\linewidth}\centering
\centering
\begin{tikzpicture}[scale=.85]
\begin{scope}[xshift=0.0cm]\filldraw[thick,fill=DeepBlue!10] (-.12,-1)..controls(-.7,-.4)and(-.7,.4)..(-.12,1)--(.12,1)..controls(.7,.4)and(.7,-.4)..(.12,-1)--cycle;\node at(0,-1.45){双凸};\end{scope}
\begin{scope}[xshift=1.7cm]\filldraw[thick,fill=DeepBlue!10] (-.3,-1)--(-.3,1)--(-.05,1)..controls(.7,.4)and(.7,-.4)..(-.05,-1)--cycle;\node at(0,-1.45){平凸};\end{scope}
\begin{scope}[xshift=3.4cm]\filldraw[thick,fill=DeepBlue!10] (-.45,-1)..controls(.15,-.4)and(.15,.4)..(-.45,1)--(-.25,1)..controls(.95,.4)and(.95,-.4)..(-.25,-1)--cycle;\node at(0,-1.45){正弯月};\end{scope}
\begin{scope}[xshift=5.1cm]\filldraw[thick,fill=DeepBlue!10] (-.65,-1)..controls(.3,-.4)and(.3,.4)..(-.65,1)--(.15,1)..controls(.55,.4)and(.55,-.4)..(.15,-1)--cycle;\node at(0,-1.45){负弯月};\end{scope}
\begin{scope}[xshift=6.8cm]\filldraw[thick,fill=DeepBlue!10] (-.5,-1)--(-.5,1)--(.5,1)..controls(-.3,.4)and(-.3,-.4)..(.5,-1)--cycle;\node at(0,-1.45){平凹};\end{scope}
\begin{scope}[xshift=8.5cm]\filldraw[thick,fill=DeepBlue!10] (-.6,-1)..controls(-.1,-.4)and(-.1,.4)..(-.6,1)--(.6,1)..controls(.1,.4)and(.1,-.4)..(.6,-1)--cycle;\node at(0,-1.45){双凹};\end{scope}
\end{tikzpicture}
\bookcaption{六种常见厚透镜外形；实际正负性由光焦度判定}
\label{advA5:fig:lens-types}
\end{minipage}\end{center}

厚透镜有四个值得单独核验的极限：两曲面共心、等效光焦度为零、一个曲面为平面、
两曲面均为平面。图 \ref{advA5:fig:four-thick-lens-cases} 给出界面、曲率中心和光线的位置关系；无焦系统没有有限的通常主面，
避免把曲率中心误画成焦点。

\begin{center}\begin{minipage}{\linewidth}\centering
\centering
\begin{tikzpicture}[scale=.85]
 \begin{scope}
  \draw[axis](-3.4,0)--(1,0);
  \fill[DeepBlue!8] (145:2.4)arc(145:215:2.4)--(215:3)arc(215:145:3)--cycle;
  \draw[thick](145:2.4)arc(145:215:2.4)(145:3)arc(145:215:3);
  \fill(0,0)circle(1.4pt)node[below]{$C$};
  \draw[DeepBlue,->](160:3.5)--(160:1.3);\draw[guide](160:1.3)--(0,0);
  \node at(-1.3,2){共心：法线均过 $C$};
 \end{scope}
 \begin{scope}[xshift=6cm]
  \draw[axis](-1,0)--(4.3,0);\filldraw[thick,fill=DeepBlue!8](0,-1)..controls(-0.2,-0.5)and(-0.2,0.5)..(0,1)..controls(0.2,0.5)and(0.2,-0.5)..cycle;\filldraw[thick,fill=DeepBlue!8](3,-1)..controls(2.8,-0.5)and(2.8,0.5)..(3,1)..controls(3.2,0.5)and(3.2,-0.5)..cycle;
  \foreach \h in {-.65,.65}{\draw[DeepBlue,->](-1,\h)--(0,\h)--(3,{-.5*\h})--(4.2,{-.5*\h});}
  \fill(2,0)circle(1.4pt)node[below=4pt]{$F'_1=F_2$};
  \node at(1.5,2){望远系统：无有限主面};
  \draw[<->](0,-1.35)--(3,-1.35)node[midway,below]{$d=f'_1+f_2$};
 \end{scope}
 \begin{scope}[yshift=-4.7cm]
  \draw[axis](-3.4,0)--(1.6,0);
  \fill[DeepBlue!8](-1.2,-.866025)--(-1.2,.866025)--(-1,.866025)arc(60:-60:1)--cycle;
  \draw[thick](-1.2,-.866025)--(-1.2,.866025)--(-1,.866025)arc(60:-60:1)--cycle;
  \foreach \sgn in {-1,1}{\draw[DeepBlue,->](-3.2,{\sgn*.18})--(-.516334,{\sgn*.18})--(1.426032,{0});}
  \fill(1.426032,0)circle(1.4pt)node[below=4pt]{$F'$};
  \node at(-.6,1.8){平凸：近轴附近的对称光线};\node[above left] at(-1.35,.5){$n=1.5$};\draw[guide](-1.4,.5)--(-.9,.4);
 \end{scope}
 \begin{scope}[shift={(6,-4.7)}]
  \draw[axis](-1.4,0)--(4,0);\filldraw[fill=DeepBlue!8] (0,-1.1)rectangle(1.2,1.1);
  \draw[DeepBlue,->](-1,1.07735)--(0,.5)--(1.2,.075736)--(3,-.963495);
  \draw[guide](0,.5)--(2.4,-.88564);\node[below] at(.6,-1.25){$n=1.5$};
  \node at(1.3,1.8){平行平板：出射平行、侧移};
 \end{scope}
\end{tikzpicture}
\bookcaption{厚透镜的四种特殊极限及相应光线关系}
\label{advA5:fig:four-thick-lens-cases}
\end{minipage}\end{center}

\subsection{射影变换的成像本质}

实射影几何基本定理指出，维数不小于 $2$ 的射影空间中，保持共线关系的非退化映射由半线性变换诱导。
在实数域且排除非连续域自同构后，可用齐次矩阵表示。对绕主轴对称的子午面成像，
\[
 (X,Y)\mapsto(X',Y'),\qquad
 X'=\frac{aX+c}{gX+i},\quad
 Y'=\frac{eY}{gX+i}.
\]
$g=0$ 时是望远系统，物像距线性且横向放大率恒定；$g\ne0$ 时平移坐标后必可化为
\[
 xx'=ff',\qquad \frac{y'}y=-\frac fx.
\]
因此 Gaussian 和 Newton 公式并非偶然的薄透镜特例，而是轴对称理想点成像的射影结构。

\begin{center}\begin{minipage}{\linewidth}\centering
\centering
\begin{tikzpicture}[>=Latex,scale=.72]
 \begin{scope}
  \draw[->] (-2.6,0)--(2.7,0) node[right]{$X$};
  \draw[->] (0,-1.45)--(0,1.55) node[above]{$Y$};
  \draw[DeepBlue,thick] (-2,-.65)--(1.85,.9);
  \draw[Gold!80!black,thick] (-1.8,-.95)--(2.0,.58);
  \node at (0,-1.8){物平面中的共线点};
 \end{scope}
 \begin{scope}[xshift=6.4cm]
  \draw[->] (-2.6,0)--(2.7,0) node[right]{$X'$};
  \draw[->] (0,-1.45)--(0,1.55) node[above]{$Y'$};
  \coordinate (O) at (-1.7,-1.05);
  \draw[RuleGray] (O)--(1.85,1.05) (O)--(1.35,-.15);
  \draw[DeepBlue,thick] (.65,-1.25)--(.65,1.35);
  \fill (O) circle(1.3pt) node[left]{$\mathcal O$};
  \node at (0,-1.8){齐次空间投影};
 \end{scope}
 \draw[->,very thick] (2.75,.65)--(3.65,.65) node[midway,above]{$\mathcal C$};
\end{tikzpicture}
\bookcaption{共线性映射可由高一维线性变换与射影实现}
\label{advA5:fig:projective-map}
\end{minipage}\end{center}

\begin{advexamplebox}{厚目镜望远镜：景物像、刻线与出瞳}\label{advA5:ex:thick-eyepiece}
如图，薄凸物镜 $L_O$ 在右，折射率 $n=3$ 的厚凸目镜 $L_E$ 在左，光从右向左传播；镜筒长度 $L=8\,\mathrm{cm}$。目镜右表面刻有基准线，用来测量远物张角。设计要求景物与刻线均可由放松眼看清，并把右侧物镜的孔径成像到目镜左出射面，使收集的光束进入眼瞳。角放大率大小 $M=5$。求目镜厚度 $x$、两镜空气间距 $y$、物镜焦距 $f_O$、目镜左、右表面曲率半径大小 $R,R'$，以及匹配人眼瞳孔 $d=4\,\mathrm{mm}$ 的物镜孔径 $D$。采用近轴球面折射模型。
\end{advexamplebox}
\begin{advsolution}
原图的“将右侧凸透镜成像至左侧单球面透镜”是\textbf{物镜孔径的瞳成像}，与景物成像同时发生。不能把它读成目镜右表面向左表面成像。

\begin{center}\begin{minipage}{\linewidth}\centering
% Design: 12 cm print width, 9 pt labels; source p175 right-to-left layout.
% [eye d] <- [left R | glass x,n | right R'刻线] <- air y <- [objective D].
% Exact first-order ray vertices; surface outlines only identify the curved faces.
\begin{tikzpicture}[>=Latex,font=\small,x=1.25cm,y=1.25cm]
\draw[gray,dashed](-.7,0)--(8.8,0);
\draw[thick,fill=DeepBlue!6](.15,-.9)..controls(-.05,-.4)and(-.05,.4)..(.15,.9)--(2.8,.9)..controls(3.08,.4)and(3.08,-.4)..(2.8,-.9)--cycle;
\draw[thick,<->](8,-1.15)--(8,1.15);
\node[above]at(1.5,1.1){目镜 $L_E$，$n=3$};\node[above]at(8,1.25){物镜 $L_O$};
\node[below left]at(0,-.8){$R$};\node[below right]at(3,-.8){$R'$};
\foreach \u in {-.7,.7}{
 \draw[DeepBlue,thick,->](8.65,\u)--(8,\u)--(3,0)--(0,{-.2*\u})--(-.55,{-.2*\u});
}
\fill(3,0)circle(1.3pt);\draw[gray,thin](3,0)--(4,1.15);\node[above]at(4,1.15){景物第一像与刻线};
\draw[<->](0,-1.4)--(3,-1.4)node[midway,below]{$x$};
\draw[<->](3,-1.4)--(8,-1.4)node[midway,below]{$y$};
\draw[<->](0,-2.1)--(8,-2.1)node[midway,below]{$L=x+y$};
\draw[<->](8.45,-.7)--(8.45,.7)node[midway,right]{$D$};
\draw[<->](-.35,-.14)--(-.35,.14);\node[left]at(-.45,.4){$d$};
\end{tikzpicture}
\figcap{轴上景物的两条边缘光线。目镜右面是景物像面，左面是物镜的孔径像面；两种共轭不能混用。\label{advA5:fig:eyepiece-source}}
\end{minipage}\end{center}
远物经薄物镜在后焦面成像；为使该像与刻线重合，应有
\[
 y=f_O,\qquad x+y=L.
\]
刻线在玻璃内距左出射面 $x$。经该球面折射后要成为平行光；在玻璃到空气、传播向左的约定下，左凸面的有向曲率为 $-R$，故
\[
 \frac n x+\frac1\infty=\frac{1-n}{-R}
 \quad\Longrightarrow\quad R=\frac{n-1}{n}x.
\]
另一方面，物镜孔径作为目镜右球面的物，物距 $y$，其像在玻璃内左表面，像距 $x$。空气到玻璃的球面式给出
\[
 \frac1y+\frac n x=\frac{n-1}{R'}.
\]
远物张角 $\theta$ 在刻线面形成高度大小 $f_O\theta$，左表面出射角大小为 $n f_O\theta/x$，所以
\[
 M=\frac{nf_O}{x}=5.
\]
与 $x+y=8\,\mathrm{cm}$、$y=f_O$ 联立，得
\[
 \boxed{x=3.00\,\mathrm{cm},\qquad y=f_O=5.00\,\mathrm{cm}},
 \qquad
 \boxed{R=2.00\,\mathrm{cm},\quad R'=\frac53\,\mathrm{cm}}.
\]
瞳成像的横向倍率是
\[
 \beta_p=-\frac{x}{ny}=-\frac15.
\]
使整个孔径像恰好落入眼瞳，有 $d=|\beta_p|D$，即
\[
 \boxed{D=5d=20\,\mathrm{mm}}.
\]
独立检验：令物镜面光线高度为 $u$、入射倾角为 $\theta$，连续应用两球面折射和玻璃内传播，可得目镜左面的高度 $-u/5$、出射倾角 $-5\theta$。前者与 $\theta$ 无关，说明孔径确已成像；后者与 $u$ 无关，说明远物确实给出平行出射光。负号表示景物倒像。这些是近轴设计参数，不能据此声称有限孔径下的球差、色差与边缘通光已经全部校正。
\end{advsolution}

\begin{advexamplebox}{潜望镜的五片透镜、镜筒长度与视场}\label{advA5:ex:periscope}
潜望镜有两块反射棱镜及五片薄凸透镜 $L_0,\ldots,L_4$，焦距和有效孔径分别为 $f_i,D_i$。两棱镜反射中心的竖直间距为 $H$。原图在 $L_1,L_3$ 处分别标出了远物的第一、第二实像 $y',y''$，$L_2$ 在二者之间。说明各镜作用，求适配镜筒的焦距及共轭条件、望远倍率，并讨论各孔径如何限制视场。
\end{advexamplebox}
\begin{advsolution}
$L_0$ 是物镜；$L_1$ 是第一像面上的场镜；$L_2$ 是将第一像再次成像的中继物镜；$L_3$ 是第二像面上的场镜；$L_4$ 是目镜。\textbf{不能把 $L_2,L_3$ 当作两片中继物镜}：光线在 $L_3$ 处已经会聚为第二像，场镜改变光束的去向，却不改变位于自身平面上的物像点。

\begin{center}\begin{minipage}{\linewidth}\centering
% Design: 12.5cm, 9pt; left folded source geometry, right unfolded distance map.
\begin{tikzpicture}[>=Latex,font=\small,scale=.95]
\begin{scope}
 \draw[gray,dashed](-2.2,5)--(0,5)--(0,-1)--(2.1,-1);
 \draw[thick](-.5,5.5)--(.5,4.5);\draw[thick](-.95,-.05)--(.95,-1.95);
 \node[right]at(.6,5){$P$};\node[below]at(-.75,-1.85){$P'$};
 \foreach \yy/\name in {4.4/0,3.3/1,1.9/2,.5/3}{
  \draw[thick,<->](-.75,\yy)--(.75,\yy);\node[right]at(.85,\yy){$L_{\name}$};}
 \draw[thick,<->](1.5,-2.25)--(1.5,.25);\node[right]at(1.7,-1){$L_4$};
 \draw[DeepBlue,thick,->](-2,5.4)--(-.4,5.4)--(-.4,4.4)--(0,3.3)--(.5090909,1.9)--(0,.5)--(-.4,-.6)--(1.5,.0909091)--(2.2,.0909091);
 \draw[DeepBlue,thick](-2,4.6)--(.4,4.6)--(.4,4.4)--(0,3.3)--(-.5090909,1.9)--(0,.5)--(.8571429,-1.8571429)--(1.5,-2.0909091)--(2.2,-2.0909091);
 \fill(0,3.3)circle(1.3pt);\fill(0,.5)circle(1.3pt);
 \node[left]at(-.9,3.3){$y'$};\node[left]at(-.9,.5){$y''$};
 \draw[<->](-2.6,-1)--(-2.6,5)node[midway,left]{$H$};
 \node[below,align=center]at(-.2,-2.6){原图折叠关系\\棱镜玻璃路程按薄棱镜忽略};
\end{scope}
\begin{scope}[xshift=7cm]
 \draw[gray,dashed](0,-1.2)--(0,5.5);
 \foreach \yy/\name in {5/0,3.8/1,1.8/2,-.2/3,-1.2/4}{
  \draw[thick,<->](-.8,\yy)--(.8,\yy);\node[left]at(-.9,\yy){$L_{\name}$};}
 \foreach \top/\bottom/\lab in {5/3.8/f_0,3.8/1.8/a,1.8/-.2/b,-.2/-1.2/f_4}{
  \draw[<->](1.3,\top)--(1.3,\bottom)node[midway,right]{$\lab$};}
 \node[below,align=center]at(0,-2.6){沿传播方向展开\\$1/a+1/b=1/f_2$};
\end{scope}
\end{tikzpicture}
\figcap{潜望镜原图的两次成像与展开距离。左图只表示轴上束与折叠关系，尺寸未按比例；右图明确两片场镜所在的像面。\label{advA5:fig:periscope-source}}
\end{minipage}\end{center}
令 $a=L_1L_2$、$b=L_2L_3$ 为展开后的正距离。景物从无穷远到无穷远的共轭条件为
\[
 L_0L_1=f_0,\qquad \frac1a+\frac1b=\frac1{f_2},\qquad L_3L_4=f_4.
\]
第一像高度为 $f_0\theta$，中继倍率为 $-b/a$，目镜将第二像转换为倾角 $-y''/f_4$，故
\[
 \boxed{M=+\frac{f_0b}{af_4}}.
\]
系统两次倒像给出正立像；只有等倍中继 $a=b=2f_2$ 时才有 $M=f_0/f_4$。

$H$ 是棱镜中心距，不应直接改称光程。设从上棱镜中心 $P$ 到 $L_0$ 的展开距离为 $p$，下棱镜中心 $P'$ 到 $L_4$ 的距离为 $q$，则
\[
 \boxed{H=f_0+a+b+f_4+p-q}.
\]
常用的对称端部近似 $p=q$ 下，令 $S=H-f_0-f_4$，有
\[
 a+b=S,\quad ab=f_2S,\quad
 a,b=\frac{S\pm\sqrt{S(S-4f_2)}}2.
\]
因此必须 $H\ge f_0+f_4+4f_2$；等号才是等倍中继。原图未给 $p,q$ 和等倍条件，不能无条件写出 $H=4f_2$。

进一步求\textbf{全照明视场}，即物镜全孔径都能通过。令入口瞳坐标 $|u|\le D_0/2$，定义
\[
 C_2=f_0+a-\frac{af_0}{f_1},\qquad
 C_4=-\frac{(b+f_4)f_0}{a}-\frac{f_4 C_2}{b}
       +\frac{f_4bf_0}{af_3}.
\]
依次用 $h_{\rm next}=h+\ell\theta$、$\theta_{\rm out}=\theta_{\rm in}-h/f$，得到
\[
 h_1=f_0\theta,\quad h_2=C_2\theta-\frac a{f_0}u,\quad
 h_3=-\frac ba f_0\theta,\quad
 h_4=C_4\theta+\frac{af_4}{bf_0}u.
\]
每个面都要满足 $|h_i|\le D_i/2$。对整个入口孔径取最大值，得
\[
\begin{split}
 2f_0|\theta|&\le D_1,\\
 2|C_2||\theta|+\frac a{f_0}D_0&\le D_2,\\
 2\frac ba f_0|\theta|&\le D_3,\\
 2|C_4||\theta|+\frac{af_4}{bf_0}D_0&\le D_4.
\end{split}
\]
满足所有不等式的最大 $|\theta|$ 为全照明半视场角。$D_1,D_3$ 是像面视场限制；一般情况下 $D_2,D_4$ 也会随视场截光。若允许渐晕，应要求存在\emph{同一个} $u$ 同时通过所有孔径，而不是分别挑选不同光线。

场镜常按瞳共轭条件设计：
\[
 \frac1{f_1}=\frac1{f_0}+\frac1a,\qquad
 \frac1{f_3}=\frac1b+\frac1{f_4}.
\]
这时 $C_2=C_4=0$，全照明条件简化为
\[
 |\theta|\le\min\left(\frac{D_1}{2f_0},\frac{aD_3}{2bf_0}\right),\quad
 D_2\ge\frac a{f_0}D_0,\quad D_4\ge\frac{af_4}{bf_0}D_0.
\]
在这一附加设计及 $p=q$ 条件下，各焦距满足
\[
 \frac1{f_2}=\frac1{f_1}+\frac1{f_3}-\frac1{f_0}-\frac1{f_4},\qquad
 H=\frac{f_0^2}{f_0-f_1}+\frac{f_4^2}{f_4-f_3}.
\]
棱镜的有效面宽若有限，也须加入相同的逐面截光检验；仅由五个透镜孔径无法保证真实棱镜不挡光。
\end{advsolution}

\subsection{光度学}

辐射度量学描述电磁辐射本身，光度学则把人眼的明视觉响应纳入测量。
国际单位制以发光强度为基本量，其单位为坎德拉（candela，$\mathrm{cd}$）。
按照现行定义，频率 $540\times10^{12}\,\mathrm{Hz}$ 的单色辐射，
光视效能固定为 $683\,\mathrm{lm/W}$。这一定义与旧的“特定光源”实物基准不同，
本质上把坎德拉归结为瓦特和人眼标准响应。

对波长为 $\lambda$ 的单色辐射，引入 CIE 明视觉光视效率函数 $V(\lambda)$。
若辐射强度为 $\mathcal I$，发光强度为 $I$，则
\[
 I=K_mV(\lambda)\mathcal I,\qquad
 K_m\simeq683\,\mathrm{lm/W}.
\]
对连续光谱，若辐射强度谱为 $\mathcal i(\lambda)$，
\[
 \mathcal I=\int\mathcal i(\lambda)\,\dd\lambda,\qquad
 I=K_m\int V(\lambda)\mathcal i(\lambda)\,\dd\lambda.
\]
图 \ref{advA5:fig:luminosity-function} 显示 $V(\lambda)$ 在约
$555\,\mathrm{nm}$ 附近达到峰值，紫端和红端迅速衰减。

\begin{center}\begin{minipage}{\linewidth}\centering
\centering
\begin{tikzpicture}
\begin{axis}[
 width=.82\textwidth,height=5.8cm,
 xmin=380,xmax=730,ymin=0,ymax=1.08,
 xlabel={波长 $\lambda/\mathrm{nm}$},
 ylabel={$V(\lambda)$},
 grid=major,
 axis line style={RuleGray},
 tick label style={font=\small},
 label style={font=\small},
 every axis plot/.append style={very thick,DeepBlue},
]
\addplot[smooth] coordinates {
(380,.000039)(400,.000396)(420,.004)(440,.023)(460,.060)
(480,.139)(500,.323)(520,.710)(530,.862)(540,.954)(550,.995)
(555,1)(560,.995)(570,.952)(580,.870)(590,.757)(600,.631)
(610,.503)(620,.381)(630,.265)(640,.175)(650,.107)(660,.061)
(670,.032)(680,.017)(690,.00821)(700,.00410)(710,.00209)
(720,.00105)(730,.00052)};
\addplot[Gold!85!black,dashed] coordinates {(555,0)(555,1)};
\node[anchor=south west,font=\small] at (axis cs:600,.94) {峰值约 $555\,\mathrm{nm}$};
\end{axis}
\end{tikzpicture}
\bookcaption{CIE 明视觉光视效率函数}
\label{advA5:fig:luminosity-function}
\end{minipage}\end{center}

常用光度量与辐射度量的对应关系为
\begin{table}[htbp]
\centering
\caption{光度量与辐射度量的对应}
\begin{tabular}{llll}
\toprule
光度量&单位&辐射度量&单位\\
\midrule
光能量&$\mathrm{lm\,s}$&辐射能量&$\mathrm J$\\
光通量 $\Phi$&$\mathrm{lm}$&辐射通量 $P$&$\mathrm W$\\
发光强度 $I$&$\mathrm{cd}$&辐射强度 $\mathcal I$&$\mathrm{W/sr}$\\
亮度 $B$&$\mathrm{cd/m^2}$&辐亮度 $\mathcal B$&$\mathrm{W/(sr\,m^2)}$\\
照度 $E$&$\mathrm{lx}$&辐照度&$\mathrm{W/m^2}$\\
\bottomrule
\end{tabular}
\end{table}

若辐射通量谱为 $p(\lambda)$，则
\[
 P=\int p(\lambda)\,\dd\lambda,\qquad
 \Phi=\int K_mV(\lambda)p(\lambda)\,\dd\lambda.
\]
光源的总光视效能定义为 $K=\Phi/P$。这里 $P$ 是辐射功率，故 $K$ 取决于光谱；若以输入电功率为分母评价灯具，则还须计入电能转化为辐射的效率，
不能只由某个“代表波长”估计。

点光源对某方向的发光强度、被照面的照度分别为
\[
 I=\frac{\dd\Phi}{\dd\Omega},\qquad
 E=\frac{\dd\Phi}{\dd S}.
\]
面光源沿与法线夹角为 $\theta$ 的方向的亮度为
\[
 B=\frac{\dd^2\Phi}{\dd\Omega\,\dd S\cos\theta}.
\]
若 $B$ 与方向无关，则称为 Lambert 发射体。此时
$\dd S\cos\theta$ 正是沿观察方向看到的投影面积。

无吸收的理想光学系统还满足三种彼此关联的不变量：
\[
 \dd G=n^2\cos\theta\,\dd S\,\dd\Omega,
 \qquad n(y_1u_2-y_2u_1)=n'(y'_1u'_2-y'_2u'_1),
 \qquad \frac{B}{n^2}=\frac{B'}{n'^2}.
\]
中间一式是同一傍轴系统内两条近轴光线的拉格朗日不变量，$u_i$ 为小倾角；不能把任意单条光线的 $nyu$ 当作处处不变。最后一式对同一频率、无吸收、无反射损失且无截光的理想传输通道成立，在物像方折射率相等时简化为亮度不变。
它说明被动光具可以放大像或增大收集立体角，却不能同时提高单位投影面积、单位立体角的辐射亮度；
这正是光学扩展量守恒的内容。

% AdvancedOpticsQuality20261023
\section{波动光学}

\subsection{标量波、复振幅与傍轴近似}

标量波满足
\[
 \left(\frac1{c^2}\frac{\partial^2}{\partial t^2}-\nabla^2\right)
 A(\vect r,t)=0.
\]
在直角、柱和球坐标中，典型外行波分别为
\[
\begin{aligned}
 A_{\rm p}&=A_0\cos(\omega t-\vect k\cdot\vect r+\varphi),\\
 A_{\rm c}&=\frac{A_0}{\sqrt{\rho}}\cos(\omega t-k\rho+\varphi),\\
 A_{\rm s}&=\frac{A_0}{r}\cos(\omega t-kr+\varphi),
\end{aligned}
\qquad \frac{\omega}{|k|}=c.
\]
上式柱面波是 $k\rho\gg1$ 时的远场渐近式，严格的单频径向解为 Hankel 函数 $H_0^{(1)}(k\rho)$；球面式在离开源点处是严格外行解。柱面波和球面波的幅度衰减分别保证单位长度圆柱面、球面上的平均能流守恒。

\begin{center}\begin{minipage}{\linewidth}\centering
\centering
\begin{tikzpicture}[>=Latex,scale=.72]
 \begin{scope}[xshift=-4.7cm]
  \foreach \y in {-.9,-.3,.3,.9}{\draw[DeepBlue,thick] (-1.45,\y)--(1.45,\y);}
  \foreach \x in {-1.0,0,1.0}{\draw[->,Gold!85!black] (\x,-1.15)--(\x,1.45);}
  \node at (0,-1.55){平面波};
 \end{scope}
 \begin{scope}
  \foreach \r in {.35,.72,1.1,1.48}{\draw[DeepBlue,thick] (25:\r) arc[start angle=25,end angle=155,radius=\r];}
  \foreach \a in {35,65,95,125,150}{\draw[->,Gold!85!black] (0,0)--(\a:1.75);}
  \draw[RuleGray] (-.22,-.2)--(.22,.2);
  \node at (0,-1.55){柱面波截面};
 \end{scope}
 \begin{scope}[xshift=4.7cm]
  \foreach \r in {.38,.76,1.14,1.52}{\draw[DeepBlue,thick] (0,0) circle (\r);}
  \foreach \a in {20,80,145,210,275,335}{\draw[->,Gold!85!black] (0,0)--(\a:1.82);}
  \fill (0,0) circle(1.2pt);
  \node at (0,-1.95){球面波截面};
 \end{scope}
\end{tikzpicture}
\bookcaption{平面波、柱面波和球面波的波前与传播方向}
\label{advA5:fig:three-waves}
\end{minipage}\end{center}

统一采用光学符号约定
\[
 A(\vect r,t)=A(\vect r)\e^{-\ii\omega t},\qquad
 A(\vect r)=|A|\e^{\ii\phi(\vect r)},\qquad
 \vect k=\nabla\phi.
\]
单色场的空间部分满足 Helmholtz 方程
\[
 \boxed{(\nabla^2+k^2)A(\vect r)=0},\qquad k=\omega/c.
\]
在均匀、各向同性、线性无源介质中，电磁场各分量满足相应标量 Helmholtz 方程，势函数还须选用相容规范；非均匀介质不能无条件逐分量套用同一标量方程，
但偏振问题仍需保留矢量方向。

取主传播轴为 $z$，傍轴平面波写成
\[
 A(x,y,z)=A_0\e^{\ii kz}
 \e^{\ii k(\theta_xx+\theta_yy)-\ii kz(\theta_x^2+\theta_y^2)/2},\qquad
 |\theta_x|,|\theta_y|\ll1.
\]
物平面 $z=-L$ 上 $(X,Y)$ 点发出的球面波传播到像平面 $(x,y,0)$ 时，
\[
\begin{aligned}
 R&=\sqrt{L^2+(x-X)^2+(y-Y)^2}\\
 &\simeq L+\frac{(x-X)^2+(y-Y)^2}{2L},
\end{aligned}
\]
所以
\[
 A(x,y)\propto
 \e^{\ii k(x^2+y^2)/(2L)}
 \e^{-\ii k(Xx+Yy)/L}.
\]
图 \ref{advA5:fig:spherical-plane-approx} 中虚线是局部切平面；
图中 $X,Y,x,y$ 的位置与公式保持同一符号约定。

\begin{center}\begin{minipage}{\linewidth}\centering
\centering
\begin{tikzpicture}[>=Latex,scale=.76]
 \draw[->] (-5,0)--(5.4,0) node[right]{$z$};
 \draw[thick] (-4,-1.6)--(-4,1.8) node[above]{物平面};
 \draw[thick] (3.8,-1.6)--(3.8,1.8) node[above]{像平面};
 \coordinate (S) at (-4,.8);
 \coordinate (P) at (3.8,.42);
 \fill (S) circle(1.3pt) node[left]{$(X,Y)$};
 \fill (P) circle(1.3pt) node[right]{$(x,y)$};
 \draw[DeepBlue,thick,->] (S)--(P) node[midway,above]{$R$};
 \draw[DeepBlue!65] (S) ++(-65:1.5) arc[start angle=-65,end angle=65,radius=1.5];
 \draw[DeepBlue!65] (S) ++(-42:3.0) arc[start angle=-42,end angle=42,radius=3.0];
 \draw[DeepBlue!55,dashed] (2.82,-1.15)--(3.05,1.55)
   node[above left,text=TextGray,xshift=-6pt]{局部切波前};
 \draw[<->] (-4,-1.95)--(3.8,-1.95) node[midway,below]{$L$};
 \draw[Gold!80!black] (-3.15,.8) arc[start angle=0,end angle=-3,radius=.85];
 \node[Gold!70!black] at (-3.0,.57) {$\theta$};
\end{tikzpicture}
\bookcaption{球面波在像平面附近的傍轴展开与局部平面波近似}
\label{advA5:fig:spherical-plane-approx}
\end{minipage}\end{center}

当观测孔径半径为 $\rho$ 时，若
\[
 \rho^2\ll\lambda L,
\]
二次相位在孔径内可以忽略，得到远场平面波近似。若不满足此式，
保留二次相位即为 Fresnel 传播；这不是“误差项”，而是近场衍射和散焦的主要来源。

\subsection{相干与多光束干涉}

两复振幅叠加为 $A=A_1+A_2$，探测器测得时间平均光强
\[
 I=I_1+I_2+2\sqrt{I_1I_2}\,
 \operatorname{Re}\gamma_{12},
\qquad
 \gamma_{12}=
 \frac{\langle A_1^*A_2\rangle}
 {\sqrt{\langle|A_1|^2\rangle\langle|A_2|^2\rangle}}.
\]
$|\gamma|=1$ 为完全相干，$|\gamma|=0$ 为完全不相干。扩展光源限制空间相干性，
典型相干角量级为 $\Delta\theta\sim\lambda/b$；有限谱宽限制时间相干性，
相干长度量级为 $L_c\sim\lambda^2/\Delta\lambda$。

\begin{advexamplebox}{两界面的 Stokes 关系与等效反射、透射}\label{advA5:ex:multibeam}
三层无吸收介质有两个平行界面。正向通过第 $i$ 个界面的振幅反射、透射系数为
$r_i,t_i$，反向为 $r_i',t_i'$；层内单程相位为 $\delta$。

\begin{enumerate}
\item 从能流守恒和时间反演说明 Stokes 关系的正确复数形式；
\item 求整个双界面的等效振幅反射系数与透射系数。
\end{enumerate}

各介质中的通道均为传播波，选定同一偏振并按单位能流归一化，排除全反射的倏逝通道情形。

\QAOFrame{\QAOInterfaces{0}}{两平行界面的传播通道及单程相位；幅度系数使用同一能流规范。}{quality:A5:q7}
\end{advexamplebox}
\begin{advsolution}
若先把行波振幅按单位能流归一化，无吸收界面的散射矩阵
\[
 S_i=\begin{pmatrix}r_i&t_i'\\ t_i&r_i'\end{pmatrix}
\]
满足 $S_i^\dagger S_i=I$。互易、时间反演对称时还可取 $S_i=S_i^{\mathsf T}$。
因此一般的复数 Stokes 关系是
\[
 |r_i|^2+|t_i|^2=1,\qquad
 r_i^*t_i'+t_i^*r_i'=0.
\]
只有在系数可同时取实数并采用特定相位规范时，才化成
$r_i+r_i'=0$、$t_it_i'+r_i^2=1$。这修正了直接把“实系数关系”用于任意复系数的错误。

从第一界面反射后的各束构成几何级数：
\[
 r=r_1+t_1t_1'r_2\e^{2\ii\delta}
 \sum_{m=0}^{\infty}(r_1'r_2\e^{2\ii\delta})^m.
\]
故
\[
 r=r_1+\frac{t_1t_1'r_2\e^{2\ii\delta}}
 {1-r_1'r_2\e^{2\ii\delta}},
\qquad
 t=\frac{t_1t_2\e^{\ii\delta}}
 {1-r_1'r_2\e^{2\ii\delta}}.
\]
若采用实系数相位规范 $r_1'=-r_1$、$t_1t_1'=1-r_1^2$，
便化为
\[
 \boxed{r=\frac{r_1+r_2\e^{2\ii\delta}}
 {1+r_1r_2\e^{2\ii\delta}}},
 \qquad
 \boxed{t=\frac{t_1t_2\e^{\ii\delta}}
 {1+r_1r_2\e^{2\ii\delta}}}.
\]
$r_2=0$ 时退化为单界面 $r=r_1$；相位平均后交叉项消失，回到非相干多次反射。

\QAOFrame{\QAOInterfaces{1}}{两条透射路径与第一次内部反射；每一新增往返都贡献相同的复数乘子。}{quality:A5:a7}

图~\ref{quality:A5:a7}中第一束在界面1直接反射，第二束先进入中间层，再经界面2反射并从界面1退出。因此第二束的系数同时含 $t_1,t'_1,r_2$ 和往返相位；后续往返再乘 $r'_1r_2\e^{2\ii\delta}$。能流归一化后透射率才可直接写为 $|t|^2$。
\end{advsolution}

\begin{advexamplebox}{周期双层介质的允许带与禁带}\label{advA5:ex:photonic-crystal}
折射率 $n_1,n_2$、厚度 $d_1,d_2$ 的两层介质周期交错。
分析哪些频率的正入射光可以传播，哪些频率落入禁带。

介质均匀、无吸收、非磁且折射率为正，先讨论无限周期结构的 Bloch 模式。

\QAOFrame{\QAOStackQuestion}{无限周期双层结构；每一周期的两个折射率与厚度必须分别保留。}{quality:A5:q8}
\end{advexamplebox}
\begin{advsolution}
一周期 $\Lambda=d_1+d_2$ 的特征矩阵为 $M=M_2M_1$。对非磁介质正入射，
令相位厚度 $\delta_j=k_0n_jd_j$，光学导纳 $\eta_j=n_j$，单层矩阵
\[
 M_j=
 \begin{pmatrix}
 \cos\delta_j&\ii\sin\delta_j/\eta_j\\
 \ii\eta_j\sin\delta_j&\cos\delta_j
 \end{pmatrix}.
\]
Bloch 条件要求本征值为 $\e^{\pm\ii K\Lambda}$，且 $\det M=1$，故
\[
 \boxed{\cos K\Lambda=
 \cos\delta_1\cos\delta_2
 -\frac12\left(\frac{\eta_1}{\eta_2}
 +\frac{\eta_2}{\eta_1}\right)
 \sin\delta_1\sin\delta_2}.
\]
当右端绝对值不超过 $1$ 时 $K$ 为实数，是允许带；超过 $1$ 时
$K$ 有非零虚部，场随周期数指数衰减，是禁带。四分之一波堆
$n_1d_1=n_2d_2=\lambda_0/4$ 在中心波长处右端
$=-(n_1/n_2+n_2/n_1)/2<-1$（只要 $n_1\ne n_2$），所以形成反射禁带。
$n_1=n_2$ 时右端化为 $\cos(\delta_1+\delta_2)$，禁带消失，极限正确。

\QAOFrame{\QAOStackAnswer}{定量示例：$n_1/n_2=2$、$n_1d_1=n_2d_2=\lambda_0/4$；阴影内矩阵迹超出传播区间。}{quality:A5:a8}

图~\ref{quality:A5:a8}的阴影由 $|\operatorname{tr}M/2|>1$ 确定，对应 Bloch 特征值一增一减。有限厚度的无吸收叠层仍可有隧穿透射，禁带不是吸收带；带边的重根也应单独看待，不能从 $K$ 实数就断言任意有限器件都完全透射。
\end{advsolution}

\subsection{衍射、路径积分与光栅}

\begin{advtheorem}{Kirchhoff 衍射积分}\label{advA5:thm:kirchhoff}
设 $A$ 在体积 $V$ 内满足 Helmholtz 方程，Green 函数
$G(\vect r,\vect r')=\e^{\ii kR}/R$，$R=|\vect r-\vect r'|$。
若边界由孔径面和远处闭合面组成，则
\[
 A(\vect r)=\frac1{4\pi}\iint_{\partial V}
 \left(G\frac{\partial A}{\partial n'}
 -A\frac{\partial G}{\partial n'}\right)\dd S'.
\]
在 Kirchhoff 边界近似下，孔径 $S$ 上可写成
\[
 \boxed{A(\vect r)\simeq\frac1{\ii\lambda}
 \iint_S A(\vect r')\frac{\e^{\ii kR}}{R}
 \frac{\cos\theta+\cos\theta'}2\,\dd S'}.
\]
\end{advtheorem}

\noindent\textit{证明。}
对 $A$ 与 $G$ 应用 Green 第二恒等式：
\[
 \iiint_V(A\nabla'^2G-G\nabla'^2A)\dd V'
 =\iint_{\partial V}\left(
 A\frac{\partial G}{\partial n'}
 -G\frac{\partial A}{\partial n'}\right)\dd S'.
\]
$G$ 满足 $(\nabla'^2+k^2)G=-4\pi\delta(\vect r-\vect r')$，
$A$ 满足齐次 Helmholtz 方程，代入后左端给出 $-4\pi A(\vect r)$，
移项即得第一式。远处半球在 Sommerfeld 辐射条件下无贡献。
在孔径上把 $A$ 和法向导数取为无屏时的值，在遮挡部分取零；
对 $G$ 的法向导数保留主导的 $\ii k\cos\theta/R$ 项，
整理得到倾斜因子 $(\cos\theta+\cos\theta')/2$。这也说明 Kirchhoff 边界条件是近似，
不能同时作为严格的 Dirichlet 与 Neumann 条件。

在透镜后焦面观察且忽略缓变倾斜因子时，得到 Fraunhofer 公式
\[
 \boxed{A(\theta_x,\theta_y)=
 \frac{\e^{\ii kf}}{\ii\lambda f}
 \iint A(x',y')\e^{-\ii k(\theta_xx'+\theta_yy')}
 \dd x'\dd y'}.
\]
因此远场衍射花样是孔径复振幅的二维 Fourier 变换。

\begin{advtheorem}{光传播的路径积分表述}\label{advA5:thm:path-integral}
在高频标量近似下，传播核可形式地写为
\[
 K(B,A)=\frac1{\mathcal N}\int_{\gamma:A\to B}
 \exp\!\left[\ii k_0\int_\gamma n(\vect r)\,\dd s\right]
 \mathcal D\gamma .
\]
短波极限 $k_0\to\infty$ 中，主要贡献来自光程驻值路径，即费马光线。
\end{advtheorem}

\noindent\textit{证明思路。}
把传播距离分成许多薄层，在每一层使用 Fresnel 核
\[
 K(\vect r_{j+1},\vect r_j)\propto
 \exp\!\left[\ii k_0n_j
 \left(\Delta z+\frac{|\Delta\vect r_\perp|^2}{2\Delta z}\right)\right],
\]
并对所有中间横向位置积分。令层数趋于无穷，就得到对路径的泛函积分。
当 $k_0$ 很大时，非驻相邻路径的相位快速转动而相消；驻相条件
$\delta\int n\,\dd s=0$ 正是费马原理。对驻值路径作二次展开，
高斯积分给出几何扩散振幅和 Maslov 相位，说明“只剩一条光线”只是领头阶近似。

\begin{advexamplebox}{2020级联双缝的中心光强与谱选择}\label{advA5:ex:cascade-slits}
单缝 $P_0$ 后依次放置 $P_1,\ldots,P_{2020}$ 共 $2020$ 个等间距双缝面，
相邻面间距 $L$，每个双缝间距 $D$，末端屏中心为 $O$。
所有缝均等宽，双缝关于主轴对称，满足傍轴条件。$P_0$ 到 $P_1$ 距离为 $L_0$，最后双缝到屏距离为 $L_1$。若把这些双缝全换成同宽、位于轴上的单缝，$O$ 点光强为 $I_0$。比较不同波长时取相同的参考强度 $I_0$，即比较相对增强因子 $I_O/I_0$。本题只比较理想传播模型的相对光强；实际多次遮挡会使绝对透射功率很小。

\begin{enumerate}
\item 求双缝级联时 $O$ 点光强；
\item 取 $L_0=L=L_1=5\,\mathrm{cm}$、$D=0.4\,\mathrm{mm}$，
求可见光中使中心光强最大的波长；
\item 求中心光强不小于最大值一半的波长范围；
\item 说明白光下条纹的颜色、间距和级次。
\end{enumerate}

\QAOFrame{\begin{tikzpicture}[>=Latex,font=\small]
 \draw[thick](-5,-1.1)--(-5,-.06) (-5,.06)--(-5,1.1)node[above]{$P_0$};
 \foreach \x/\t in {-3/P_1,-1/P_2,2/P_{2020}}{
  \draw[thick](\x,-1.1)--(\x,-.45) (\x,-.31)--(\x,.31) (\x,.45)--(\x,1.1)node[above]{$\t$};
 }
 \draw[thick](4,-1.1)--(4,1.1)node[above]{$P$};
 \fill(4,0)circle(1.4pt)node[right]{$O$};
 \draw[DeepBlue!60](-5,0)--(-3,.38)--(-1,-.38);
 \draw[DeepBlue!60](-5,0)--(-3,-.38)--(-1,.38);
 \draw[DeepBlue!60](2,-.38)--(4,0) (2,.38)--(4,0);
 \node at(.5,0){$\cdots$};
 \foreach \a/\b/\lab in {-5/-3/L_0,-3/-1/L,2/4/L_1}{
  \draw[<->](\a,-1.45)--(\b,-1.45)node[midway,below]{$\lab$};}
 \draw[<->](-3.5,-.38)--(-3.5,.38)node[midway,left]{$D$};
\end{tikzpicture}}{单缝入射、双缝级联和末端中心；省略号仅省略中间面，不改变面数。\label{advA5:fig:cascade-source}}{quality:A5:q9}
\end{advexamplebox}
\begin{advsolution}
采用窄缝离散传播模型：每缝内的振幅和相位近似均匀，忽略斜程的振幅差及单缝包络变化。若缝宽为 $a$，除 $a\ll D$ 外，还需 $2\pi aD/(\lambda L)\ll1$。
在相邻双缝面之间，同侧传播距离为 $L$，交叉传播距离在傍轴下为
$L+D^2/(2L)$。略去公共相位，单步传播矩阵为
\[
 T=\begin{pmatrix}1&\e^{\ii\delta}\\
 \e^{\ii\delta}&1\end{pmatrix},\qquad
 \delta=\frac{2\pi}{\lambda}\frac{D^2}{2L}.
\]
入射与中心读出都对应对称本征向量 $(1,1)^{\mathsf T}$，
其本征值为 $1+\e^{\ii\delta}
=2\e^{\ii\delta/2}\cos(\delta/2)$。
与全改单缝的振幅比较，第一处分束和末端合束再给出一个因子 $2$，
故
\[
 \frac{A_O}{A_{O,0}}
 =2^{2020}\e^{\ii\Phi}
 \cos^{2019}\!\frac{\delta}{2},
\]
\[
 \boxed{\frac{I_O}{I_0}
 =4^{2020}
 \cos^{4038}\!\left(\frac{\pi D^2}{2\lambda L}\right)}.
\]

数值上
\[
 \Delta\ell=\frac{D^2}{2L}=1.6000\,\mu\mathrm m.
\]
最大值要求 $\Delta\ell/\lambda=m$ 为整数。按可见范围约
$380$--$780\,\mathrm{nm}$，有
\[
 \boxed{\lambda=533.33\,\mathrm{nm}\ (m=3)},\qquad
 \boxed{\lambda=400.00\,\mathrm{nm}\ (m=4)}.
\]
若采用不同的可见光边界约定，边界附近的 $400\,\mathrm{nm}$ 是否计入应随约定说明。

半高条件为
\[
 \left|\cos\frac{\pi\Delta\ell}{\lambda}\right|
 \ge 2^{-1/4038}.
\]
令
$q=\arccos(2^{-1/4038})=0.018528\ldots$，每个整数 $m$ 周围满足
\[
 \left|\frac{\pi\Delta\ell}{\lambda}-m\pi\right|\le q.
\]
于是两段可见范围近似为
\[
 \boxed{532.29\text{--}534.38\,\mathrm{nm}},\qquad
 \boxed{399.41\text{--}400.59\,\mathrm{nm}}.
\]
\textbf{谱选择与空间条纹应分开。}最后一组双缝的两出射振幅始终相等。屏上纵坐标为 $y$ 时，末段光程差为 $Dy/L_1$，所以忽略单缝衍射包络的变化，有
\[
 \frac{I(y,\lambda)}{I_0(\lambda)}=
 4^{2020}\cos^{4038}\!\left(\frac{\pi\Delta\ell}{\lambda}\right)
 \cos^2\!\left(\frac{\pi Dy}{\lambda L_1}\right).
\]
第一个因子选择窄波段，第二个因子给出普通双缝空间条纹：
\[
 y_k=k\frac{\lambda L_1}{D},\qquad
 \Delta y_{533}=66.667\,\mu\mathrm m,\quad
 \Delta y_{400}=50.000\,\mu\mathrm m.
\]
绿光的第 $3j$ 级与紫光的第 $4j$ 级在 $y=200j\,\mu\mathrm m$ 处重合。白光经级联后主要留下这两个窄波段；各自形成相应间距的亮暗纹，重合处颜色叠加，其余处随两分量比例改变颜色。中心的实际颜色、相对亮度还取决于入射谱和 $I_0(\lambda)$，不能只由几何量确定。有限谱宽会使远处条纹逐渐模糊。级联幂次使\emph{谱峰}变窄，并不把最后双缝的单色空间因子变成同样的高次幂。

\QAOFrame{\QAOCascadeAnswer}{左为级联的谱选择，右为最终双缝的单色空间条纹；两个不同变量上的峰宽不能混同。}{quality:A5:a9}

图~\ref{quality:A5:a9}左图以某一谱峰附近的相位失谐为横坐标，展示高次幂选谱；右图固定波长，只剩末端双缝的 $\cos^2$ 空间因子。两者的自变量分别包含 $\lambda$ 和 $y$，因此窄谱峰不能直接当作窄空间条纹。题图~\ref{quality:A5:q9}中的2020个双缝面只有2019次双缝到双缝传播。
\end{advsolution}

对周期为 $d$、缝宽 $a$、缝数 $N$ 的矩形光栅，
\[
 I(\theta)=I_1
 \left(\frac{\sin\alpha}{\alpha}\right)^2
 \left(\frac{\sin N\beta}{\sin\beta}\right)^2,
\quad
 \alpha=\frac{\pi a}{\lambda}\sin\theta,\quad
 \beta=\frac{\pi d}{\lambda}\sin\theta.
\]
主极大满足 $d\sin\theta_j=j\lambda$，强度随 $N^2$ 增长，
相邻零点给出的角宽约为
$2\lambda/(Nd\cos\theta_j)$。

\begin{advexamplebox}{用 Fresnel 双棱镜阵列模拟光栅}\label{advA5:ex:biprism-grating}
实验室只有狭缝、对称双凸薄透镜（两面曲率半径 $R=30\,\mathrm{cm}$、折射率 $n_0=1.5$）和许多薄 Fresnel 双棱镜（单侧劈角 $\alpha=30'$、折射率 $n=1.005$）。拟在 $\lambda=500\,\mathrm{nm}$ 时获得一级角位置 $\theta_1=30^\circ$，且后焦面主极大两侧第一零点间宽度不大于 $0.1\,\mathrm{mm}$。按原题要求讨论至少需要多少个双棱镜以及各仪器的放置方法，并检查所需近似和器材条件。
\QAOFrame{\begin{tikzpicture}[scale=.8]
 \draw[thick](-4,-1.2)--(-4,1.2)node[above]{狭缝 $S$};
 \foreach \x in {-1.8,-1.4,-1}{\filldraw[fill=DeepBlue!10](\x,-1)--(\x+.2,0)--(\x,1)--(\x-.1,0)--cycle;}
 \node at(-1.4,1.5){双棱镜组};
 \filldraw[thick,fill=DeepBlue!8](1,-1.3)..controls(.7,-.65)and(.7,.65)..(1,1.3)..controls(1.3,.65)and(1.3,-.65)..cycle;
 \node at(1,1.65){透镜 $L$};
 \draw[thick](4,-1.3)--(4,1.3)node[above]{屏};
 \draw[guide](-4.3,0)--(4.3,0);
 \node at(0,-1.8){器件顺序示意；棱脊横移量与间距尚须确定};
\end{tikzpicture}}{狭缝、横移双棱镜组、透镜与屏的器件顺序；单元数仍待求。}{quality:A5:q10}
\end{advexamplebox}
\begin{advsolution}
\textbf{先区分题目所给条件与理想阵列模型。}
透镜焦距 $f=R/[2(n_0-1)]=0.300\,\mathrm m$；单侧小劈角的偏折量为
$\delta=(n-1)\alpha=4.363\times10^{-5}\,\mathrm{rad}$。
题给 $30^\circ$ 与全程傍轴的高精度使用并不相容。以下先给出原题形式近轴模型中的构造，再说明它不能证明原器材条件下的全局最小数目。

若等效透射场由 $N$ 个等宽、等幅、等相位的平移单元组成，周期为 $d$，则近轴阵列因子为
\[
 \left|\sum_{j=0}^{N-1}e^{-ikjd\theta}\right|^2
 =\frac{\sin^2(Nkd\theta/2)}{\sin^2(kd\theta/2)}.
\]
一级主峰位置与两侧第一零点间的宽度给出
\[
 d=\frac{\lambda}{\theta_1}=0.9549\,\mu\mathrm m,
 \qquad w=\frac{2f\lambda}{Nd}=\frac{2f\theta_1}{N},
 \qquad \boxed{N\ge3142}.
\]
这里 $\theta_1$ 以弧度代入。$N$ 是场单元数，尚不是棱镜数。

\textbf{一个有明确附加条件的 $N-1$ 片方案。}
将 $N-1$ 片薄双棱镜等效贴合于同一横截面，棱脊横移至
\[
 b_j=(j-N/2)d\quad(j=1,\ldots,N-1),
 \qquad a=\frac{d}{2\delta}=10.94\,\mathrm{mm},
\]
其中 $a$ 为狭缝到该截面的距离。另须有公共孔径选出总宽 $Nd$ 的 $N$ 个等宽区段，并在这些区段上可近似等幅照明。薄相位模型给出
\[
 \Phi(y)=\frac{ky^2}{2a}-k\delta\sum_{j=1}^{N-1}|y-b_j|.
\]
第 $q=0,\ldots,N-1$ 个区段的中心为 $c_q=[q-(N-1)/2]d$。因该区间内绝对值之和的斜率为 $2q-(N-1)$，配方得
\[
 \Phi(y)=\frac{k(y-c_q)^2}{2a}+C,
\]
各段常数 $C$ 相同：两相邻段在公共边界上连续，且边界距两段中心均为 $d/2$。故各区段确为同一场分布的平移，透镜后焦面得到“单元包络乘阵列因子”。透镜放在阵列后，屏置于后焦面；其有效口径还须接收所需全部单元。
在这一限定构型内，每片双棱镜只增加一个相位折点，要产生 $N$ 个区段至少需 $N-1$ 片，因此形式近轴结果为
\[
 \boxed{N-1=3141\ \text{片（限定构型及公共孔径假设下）}}.
\]
同轴重叠而不横移的任意多片双棱镜仍仅给两段相位，不能按片数声称形成 $2^n$ 束。

\textbf{为什么这不是原题无条件的唯一数值答案。}
原器材清单没有规定公共孔径、可否贴合及横移、照明均匀性和大角度成像规律，因而上述方案有明确的附加条件。若保留精确的 $30^\circ$，并指定平面屏坐标 $Y=f\tan\theta$，则
\[
 d=\frac{\lambda}{\sin\theta_1}=1.000\,\mu\mathrm m,
 \quad w\simeq\frac{2f\lambda}{Nd\cos^3\theta_1},
 \quad N\ge4619.
\]
此数值仅是指定大角度屏映射下的理想阵列要求，不能继续直接套用前述近轴棱镜构造。原题若要有可实现且唯一的“最少片数”，必须补足这些条件；不能混用精确光栅角与 $Y=f\theta$ 给出混合近似的束数结论。

\QAOFrame{\QAOBiprismAnswer}{四个平移相位单元的简化示意；各段抛物线相同，公共边界上相位连续。}{quality:A5:a10}

图~\ref{quality:A5:a10}中的棱脊 $b_j$ 是相位斜率的折点，而每个区段的中心是抛物线的极小点。公共孔径必须切出这些完整区段，才能将透射场写成相同单元的平移求和；仅把棱镜同轴堆放不会产生该图的多区段结构。
\end{advsolution}

Fresnel 波带法把波前按到观察点的光程差每增加 $\lambda/2$ 分带。
在傍轴球面上，带状面元面积的增加恰好补偿传播距离变化，使相邻半波带振幅近似相等、相位相反。

\begin{center}\begin{minipage}{\linewidth}\centering
\centering
\begin{tikzpicture}[scale=.9]
 \begin{scope}
  \draw[thick,domain=-40:40,variable=\t]plot({-3+2*cos(\t)},{2*sin(\t)});
  \coordinate(P)at(2,0);\fill(P)circle(1.6pt)node[right]{$P$};
  \foreach \t in {-35,-20,-7,7,20,35}{\draw[DeepBlue!70]({-3+2*cos(\t)},{2*sin(\t)})--(P);}
  \node at(0,-1.8){同一球面上的面元至 $P$};
 \end{scope}
 \begin{scope}[xshift=6cm]
  \foreach \m in {1,...,6}{\draw[DeepBlue] (0,0)circle({.62*sqrt(\m)});}
  \draw[Gold!80!black,->](0,0)--({.62*sqrt(4)*cos(35)},{.62*sqrt(4)*sin(35)})node[above right]{$r_4$};
  \node at(0,-1.8){近轴半波带：$r_m\propto\sqrt m$};
 \end{scope}
\end{tikzpicture}
\bookcaption{Fresnel 波带与圆孔半波带的几何关系}
\label{advA5:fig:fresnel-zones}
\end{minipage}\end{center}

若把每个半波带的复振幅首尾相接，离散求和在复平面上来回折返；
对于直边衍射的一维二次相位积分，连续累积振幅形成 Cornu 螺旋；圆孔环带的面积权重不同，不把两种几何的积分曲线等同。图
\ref{advA5:fig:phasor-sum} 同时显示这两种合成方法。

\begin{center}\begin{minipage}{\linewidth}\centering
\centering
\begin{tikzpicture}[scale=1]
\begin{scope}\draw[axis](-.3,0)--(2.7,0)node[right]{$\operatorname{Re}A$};\draw[axis](0,-.3)--(0,2)node[above]{$\operatorname{Im}A$};
\draw[DeepBlue,->](0.000000,0.000000)--(1.163562,1.570739);
\draw[DeepBlue,->](1.163562,1.570739)--(1.074157,0.755514);
\draw[DeepBlue,->](1.074157,0.755514)--(1.114611,1.383667);
\draw[DeepBlue,->](1.114611,1.383667)--(1.090363,0.853532);
\draw[DeepBlue,->](1.090363,0.853532)--(1.106948,1.320797);
\draw[DeepBlue,->](1.106948,1.320797)--(1.094692,0.898263);
\draw[Gold!80!black,very thick,->](0,0)--(1.094692,0.898263)node[right]{$\sum A_m$};\node at(1,-.8){相邻半周期分组};\end{scope}
\begin{scope}[xshift=6cm]\draw[axis](-2,0)--(2,0);\draw[axis](0,-2)--(0,2);\draw[DeepBlue,thick]plot coordinates {(-1.332586,-1.091889) (-1.326154,-1.037393) (-1.307278,-0.985865) (-1.277074,-0.940047) (-1.237228,-0.902309) (-1.189888,-0.874542) (-1.137532,-0.858078) (-1.082825,-0.853647) (-1.028482,-0.861363) (-0.977129,-0.880748) (-0.931180,-0.910780) (-0.892734,-0.949963) (-0.863487,-0.996418) (-0.844671,-1.047991) (-0.837025,-1.102356) (-0.840777,-1.157129) (-0.855662,-1.209977) (-0.880957,-1.258709) (-0.915525,-1.301368) (-0.957890,-1.336299) (-1.006309,-1.362200) (-1.058851,-1.378157) (-1.113489,-1.383665) (-1.168173,-1.378622) (-1.220915,-1.363318) (-1.269857,-1.338401) (-1.313330,-1.304839) (-1.349907,-1.263867) (-1.378436,-1.216934) (-1.398069,-1.165636) (-1.408273,-1.111665) (-1.408830,-1.056738) (-1.399829,-1.002550) (-1.381645,-0.950715) (-1.354913,-0.902724) (-1.320495,-0.859907) (-1.279441,-0.823401) (-1.232949,-0.794131) (-1.182323,-0.772791) (-1.128929,-0.759843) (-1.074157,-0.755514) (-1.019381,-0.759807) (-0.965925,-0.772512) (-0.915030,-0.793223) (-0.867835,-0.821364) (-0.825347,-0.856210) (-0.788431,-0.896914) (-0.757797,-0.942535) (-0.733992,-0.992065) (-0.717404,-1.044457) (-0.708257,-1.098646) (-0.706621,-1.153579) (-0.712419,-1.208231) (-0.725442,-1.261626) (-0.745356,-1.312852) (-0.771724,-1.361076) (-0.804016,-1.405553) (-0.841628,-1.445633) (-0.883900,-1.480766) (-0.930125,-1.510506) (-0.979575,-1.534511) (-1.031503,-1.552537) (-1.085169,-1.564441) (-1.139840,-1.570169) (-1.194811,-1.569755) (-1.249405,-1.563312) (-1.302987,-1.551021) (-1.354968,-1.533129) (-1.404811,-1.509933) (-1.452031,-1.481778) (-1.496200,-1.449042) (-1.536948,-1.412134) (-1.573963,-1.371482) (-1.606989,-1.327526) (-1.635823,-1.280713) (-1.660318,-1.231489) (-1.680375,-1.180295) (-1.695939,-1.127561) (-1.707001,-1.073701) (-1.713589,-1.019113) (-1.715765,-0.964170) (-1.713625,-0.909226) (-1.707287,-0.854605) (-1.696895,-0.800609) (-1.682611,-0.747508) (-1.664612,-0.695548) (-1.643087,-0.644947) (-1.618234,-0.595893) (-1.590257,-0.548551) (-1.559362,-0.503059) (-1.525757,-0.459530) (-1.489647,-0.418054) (-1.451235,-0.378700) (-1.410718,-0.341516) (-1.368287,-0.306531) (-1.324126,-0.273757) (-1.278410,-0.243188) (-1.231304,-0.214807) (-1.182964,-0.188582) (-1.133537,-0.164468) (-1.083157,-0.142411) (-1.031950,-0.122350) (-0.980030,-0.104212) (-0.927501,-0.087920) (-0.874458,-0.073391) (-0.820984,-0.060534) (-0.767154,-0.049258) (-0.713034,-0.039465) (-0.658682,-0.031057) (-0.604147,-0.023932) (-0.549470,-0.017986) (-0.494687,-0.013115) (-0.439826,-0.009213) (-0.384911,-0.006173) (-0.329959,-0.003887) (-0.274983,-0.002250) (-0.219995,-0.001152) (-0.164999,-0.000486) (-0.110000,-0.000144) (-0.055000,-0.000018) (0.000000,0.000000) (0.055000,0.000018) (0.110000,0.000144) (0.164999,0.000486) (0.219995,0.001152) (0.274983,0.002250) (0.329959,0.003887) (0.384911,0.006173) (0.439826,0.009213) (0.494687,0.013115) (0.549470,0.017986) (0.604147,0.023932) (0.658682,0.031057) (0.713034,0.039465) (0.767154,0.049258) (0.820984,0.060534) (0.874458,0.073391) (0.927501,0.087920) (0.980030,0.104212) (1.031950,0.122350) (1.083157,0.142411) (1.133537,0.164468) (1.182964,0.188582) (1.231304,0.214807) (1.278410,0.243188) (1.324126,0.273757) (1.368287,0.306531) (1.410718,0.341516) (1.451235,0.378700) (1.489647,0.418054) (1.525757,0.459530) (1.559362,0.503059) (1.590257,0.548551) (1.618234,0.595893) (1.643087,0.644947) (1.664612,0.695548) (1.682611,0.747508) (1.696895,0.800609) (1.707287,0.854605) (1.713625,0.909226) (1.715765,0.964170) (1.713589,1.019113) (1.707001,1.073701) (1.695939,1.127561) (1.680375,1.180295) (1.660318,1.231489) (1.635823,1.280713) (1.606989,1.327526) (1.573963,1.371482) (1.536948,1.412134) (1.496200,1.449042) (1.452031,1.481778) (1.404811,1.509933) (1.354968,1.533129) (1.302987,1.551021) (1.249405,1.563312) (1.194811,1.569755) (1.139840,1.570169) (1.085169,1.564441) (1.031503,1.552537) (0.979575,1.534511) (0.930125,1.510506) (0.883900,1.480766) (0.841628,1.445633) (0.804016,1.405553) (0.771724,1.361076) (0.745356,1.312852) (0.725442,1.261626) (0.712419,1.208231) (0.706621,1.153579) (0.708257,1.098646) (0.717404,1.044457) (0.733992,0.992065) (0.757797,0.942535) (0.788431,0.896914) (0.825347,0.856210) (0.867835,0.821364) (0.915030,0.793223) (0.965925,0.772512) (1.019381,0.759807) (1.074157,0.755514) (1.128929,0.759843) (1.182323,0.772791) (1.232949,0.794131) (1.279441,0.823401) (1.320495,0.859907) (1.354913,0.902724) (1.381645,0.950715) (1.399829,1.002550) (1.408830,1.056738) (1.408273,1.111665) (1.398069,1.165636) (1.378436,1.216934) (1.349907,1.263867) (1.313330,1.304839) (1.269857,1.338401) (1.220915,1.363318) (1.168173,1.378622) (1.113489,1.383665) (1.058851,1.378157) (1.006309,1.362200) (0.957890,1.336299) (0.915525,1.301368) (0.880957,1.258709) (0.855662,1.209977) (0.840777,1.157129) (0.837025,1.102356) (0.844671,1.047991) (0.863487,0.996418) (0.892734,0.949963) (0.931180,0.910780) (0.977129,0.880748) (1.028482,0.861363) (1.082825,0.853647) (1.137532,0.858078) (1.189888,0.874542) (1.237228,0.902309) (1.277074,0.940047) (1.307278,0.985865) (1.326154,1.037393) (1.332586,1.091889)};
\draw[Gold!80!black,very thick,->](-1.573963,-1.371482)--(0.712419,1.208231)node[pos=.55,above left=4pt]{$A$};\fill(-1.573963,-1.371482)circle(1.5pt)(0.712419,1.208231)circle(1.5pt);\node at(0,-2.35){直边衍射的 Fresnel 积分};\end{scope}\end{tikzpicture}
\bookcaption{离散半波带合成与连续复振幅积分}
\label{advA5:fig:phasor-sum}
\end{minipage}\end{center}

\begin{advexamplebox}{阶梯折射率相位光栅的后焦面分布}\label{advA5:ex:gradient-grating}
光栅由 $2N$ 个宽度为 $a$、厚度为 $d$ 的透明条带组成。每一条带内部
$\dd n/\dd y=0$，相邻条带折射率增加 $\alpha a$。
波数为 $k$ 的平面波正入射，随后通过焦距 $F$ 的透镜。
求后焦面坐标 $Y$ 上的光强因子，并判断何时花样关于焦点对称。

\QAOFrame{\begin{tikzpicture}[scale=.9]
 \foreach \i in {0,...,7}{\pgfmathtruncatemacro{\shade}{8+5*\i}\filldraw[fill=DeepBlue!\shade,draw=RuleGray](0,{-1.2+.3*\i})rectangle(.5,{-1.2+.3*(\i+1)});}
 \draw[axis](0,-1.4)--(0,1.7)node[above]{$y$};
 \foreach \h in {-.9,0,.9}{\draw[DeepBlue,->](-1.6,\h)--(-.15,\h);}
 \filldraw[thick,fill=DeepBlue!8](2.5,-1.4)..controls(2.3,-0.7)and(2.3,0.7)..(2.5,1.4)..controls(2.7,0.7)and(2.7,-0.7)..cycle;
 \draw[thick](5,-1.5)--(5,1.5)node[above]{后焦面};
 \draw[guide,DeepBlue](.55,.9)--(2.5,.9)--(5,0);
 \draw[guide,DeepBlue](.55,-.9)--(2.5,-.9)--(5,0);
 \draw[<->](.75,-1.2)--(.75,-.9)node[right]{$a$};
 \draw[<->](0,-1.7)--(.5,-1.7)node[midway,below]{$d$};
 \node at(2.5,2.2){条带沿 $y$ 并列；虚线仅示透镜轴上聚焦};
\end{tikzpicture}}{折射率台阶沿横向分布；每条带内部均匀，透镜用于后焦面观察。}{quality:A5:q11}
\end{advexamplebox}
\begin{advsolution}
令后焦面的横向空间频率
$q=kY/F$。第 $m$ 条带（$m=-N,\ldots,N-1$）的附加相位为
$kd(n_0+m\alpha a)$。略去公共相位后，
\[
\begin{aligned}
 A(Y)&\propto
 \sum_{m=-N}^{N-1}
 \int_{ma}^{(m+1)a}
 \e^{\ii kd\alpha am}\e^{-\ii qy}\,\dd y\\
 &\propto
 \sinc\!\left(\frac{qa}{2}\right)
 \sum_{m=-N}^{N-1}\e^{\ii m\xi},
\end{aligned}
\]
其中
\[
 \xi=kd\alpha a-qa.
\]
有限几何级数给出
\[
 \boxed{I(Y)=I_c
 \sinc^2\!\left(\frac{kYa}{2F}\right)
 \frac{\sin^2(N\xi)}{\sin^2(\xi/2)}}.
\]
第一因子是单条带包络，第二因子是 $2N$ 条带结构因子。
当每条折射率台阶带来的相位增量为半周的整数倍
\[
 \boxed{kd\alpha a=\pi\ell,\qquad \ell\in\mathbb Z}
\]
时，利用结构因子 $G(\xi)=\sin^2(N\xi)/\sin^2(\xi/2)$ 的偶性和 $2\pi$ 周期性，可得 $G(\pi\ell-qa)=G(\pi\ell+qa)$，故整个花样关于焦点对称。奇数半周与整周的峰位不同，但都具有该对称性。一般相位台阶使主峰偏向
$Y_0=F d\alpha$（模倒易周期），即形成闪耀光栅。
令 $\alpha=0$ 可恢复 $2N$ 个相邻同相条带，即总宽 $2Na$ 的矩形孔径，量纲和极限均正确。

\QAOFrame{\QAOGratingAnswer}{取 $2N=8$：同相时向量相加，第一零点时向量闭合；每条带的公共包络另乘。}{quality:A5:a11}

图~\ref{quality:A5:a11}用八条带说明结构因子：$\xi=0$ 时所有相量同向，$\xi=\pi/4$ 时围成闭合八边形。实际强度还要乘单条带的 $\sinc^2$ 包络，不能把相量合成的主峰位置直接当作全部花样的唯一最大值。$\alpha=0$ 时各条带首尾相接且同相，实际退化为总宽 $2Na$ 的单个矩形孔径，而非缝间含不透光间隔的光栅。
对 $2N\ge2$，上述对称条件也是必要的。令 $M=2N$、$\varphi=kd\alpha a$，结构因子展开为 $G(\xi)=M+2\sum_{j=1}^{M-1}(M-j)\cos(j\xi)$。要求 $G(\varphi-u)=G(\varphi+u)$ 对所有 $u$ 成立，正弦级数的每个系数必须为零；其中 $j=1$ 的系数非零，便有 $\sin\varphi=0$，即 $\varphi=\pi\ell$。公共包络仅在离散点为零，不改变这一整体对称条件。
\end{advsolution}

\subsection{偏振密度矩阵}

% Completeness20261024: coherency matrix, its decomposition and an analyzer calculation.
在固定横截面上，用复振幅描述准单色光的两个横向分量，写成
$\vect A=(A_x,A_y)^{\mathsf T}$。把强度的比例常数吸收进振幅定义，便有
$I=\langle|A_x|^2+|A_y|^2\rangle$；尖括号表示探测时间内的平均。
稳定偏振光的两个分量有确定的振幅比和相位差，Jones向量即可描述；
若它们随时间随机变化，必须先平均振幅的乘积，而不能先平均振幅再平方。
相干矩阵定义为
\[
 P=\langle\vect A\vect A^\dagger\rangle
 =\begin{pmatrix}
 \langle|A_x|^2\rangle&\langle A_xA_y^*\rangle\\
 \langle A_yA_x^*\rangle&\langle|A_y|^2\rangle
 \end{pmatrix},\qquad I=\operatorname{tr}P.
\]
非对角元保留两个方向的相关性及相对相位。由定义直接得到 $P^\dagger=P$。
对任意复向量 $\vect z$，
\[
 \vect z^\dagger P\vect z
 =\left\langle|\vect z^\dagger\vect A|^2\right\rangle\ge0,
\]
所以 $P$ 半正定，两个本征值均非负。这一性质保证任意检偏方向测出的强度不会为负。

\paragraph{为什么能够分成完全偏振光和非偏振光}
设 $I>0$，将 $P$ 对角化，记本征值 $\lambda_1\ge\lambda_2\ge0$，
对应单位正交Jones向量 $\vect e_1,\vect e_2$。利用
$\vect e_1\vect e_1^\dagger+\vect e_2\vect e_2^\dagger=I_2$，有
\[
 \begin{aligned}
 P&=\lambda_1\vect e_1\vect e_1^\dagger
       +\lambda_2\vect e_2\vect e_2^\dagger\\
  &=(\lambda_1-\lambda_2)\vect e_1\vect e_1^\dagger+\lambda_2 I_2.
 \end{aligned}
\]
第一项沿固定偏振态 $\vect e_1$，强度为 $I_{\rm p}=\lambda_1-\lambda_2$；
第二项在任何正交偏振基中都相同，是强度为 $I_{\rm u}=2\lambda_2$ 的完全非偏振分量。
两项的强度相加等于 $I$。因此偏振度为
\[
 \mathcal P=\frac{I_{\rm p}}I
 =\frac{\lambda_1-\lambda_2}{\lambda_1+\lambda_2}
 =\sqrt{1-\frac{4\det P}{(\operatorname{tr}P)^2}}.
\]
纯偏振光的 $\lambda_2=0$，即 $\det P=0$、$\mathcal P=1$；
完全非偏振光则有 $\lambda_1=\lambda_2=I/2$，即 $P=(I/2)I_2$。
零光强时偏振度不定义，不能把零矩阵仅凭行列式为零称为纯偏振光。

\paragraph{从矩阵算出检偏器的读数}
图~\ref{comp:A5:analyzer}在横截面内标出检偏器透光轴
$\vect e_\theta=(\cos\theta,\sin\theta)^{\mathsf T}$。
理想线偏振片透过的振幅是 $A_\theta=\vect e_\theta^\dagger\vect A$，故
\[
 I_\theta=\langle|A_\theta|^2\rangle
 =\vect e_\theta^\dagger P\vect e_\theta
 =P_{xx}\cos^2\theta+P_{yy}\sin^2\theta
   +2\operatorname{Re}P_{xy}\sin\theta\cos\theta.
\]
\begin{center}\begin{minipage}{.94\linewidth}\centering
\begin{tikzpicture}[>=Latex,font=\small]
 \draw[->] (-1.5,0)--(1.8,0) node[right]{$x$};
 \draw[->] (0,-1.05)--(0,1.65) node[above]{$y$};
 \draw[very thick,DeepBlue,<->] (-.9,-.63)--(1.25,.875);
 \draw[->] (.65,0) arc[start angle=0,end angle=35,radius=.65];
 \node at (.9,.25) {$\theta$};
 \node[above right] at (1.05,.86) {$\vect e_\theta$};
 \node[below,align=center] at (0,-1.25) {传播方向垂直纸面\\透光轴与 $x$ 轴夹角为 $\theta$};
 \draw[->,thick] (2.35,.35)--(3.65,.35);
 \node[draw,rounded corners=1pt,align=center,text width=3.6cm] at (5,.35)
 {$I_\theta=\vect e_\theta^\dagger P\vect e_\theta$\\探测器记录平均强度};
\end{tikzpicture}
\figcap{偏振矩阵的检偏测量：图中的方向向量同时进入振幅投影和强度二次型。}
\label{comp:A5:analyzer}
\end{minipage}\end{center}

例如，把总强度 $I_0$ 的光分解为强度 $0.6I_0$ 的 $x$ 线偏振光和强度
$0.4I_0$ 的非偏振光，两部分互不相干。各自的相干矩阵相加，得到
\[
 P=0.6I_0\begin{pmatrix}1&0\\0&0\end{pmatrix}
   +0.2I_0 I_2
 =I_0\begin{pmatrix}0.8&0\\0&0.2\end{pmatrix}.
\]
由图~\ref{comp:A5:analyzer}的投影式，
\[
 \frac{I_\theta}{I_0}=0.8\cos^2\theta+0.2\sin^2\theta
 =0.5+0.3\cos2\theta.
\]
图~\ref{comp:A5:analyzer-result}中 $\theta=0$ 透过 $0.8I_0$，
$\theta=\pi/2$ 仍透过 $0.2I_0$：后者来自非偏振分量的一半。
本例本征偏振是线偏振，故旋转线检偏器的反差
$(I_{\max}-I_{\min})/(I_{\max}+I_{\min})=0.6$ 恰等于偏振度。
这个结论不能推广为任意光的测量公式：纯圆偏振光通过所有线检偏方向都透过一半，
却有 $\mathcal P=1$；一般椭圆偏振态还需波片才能测出虚部相关信息。
\begin{center}\begin{minipage}{.94\linewidth}\centering
\begin{tikzpicture}[>=Latex,font=\small,x=1.6cm,y=2.5cm]
 \draw[->] (0,0)--(3.5,0) node[right]{$\theta$};
 \draw[->] (0,0)--(0,1.02) node[above]{$I_\theta/I_0$};
 \draw[DeepBlue,thick,domain=0:3.141593,samples=100]
   plot(\x,{.5+.3*cos(2*\x*180/pi)});
 \draw[dashed,TextGray] (0,.2)--(3.141593,.2);
 \draw[dashed,TextGray] (0,.8)--(3.141593,.8);
 \node[left] at (0,.2) {$0.2$};\node[left] at (0,.8) {$0.8$};
 \node[below] at (0,0) {$0$};
 \node[below] at (1.570796,0) {$\pi/2$};\node[below] at (3.141593,0) {$\pi$};
 \fill[DeepBlue] (1.570796,.2) circle[radius=1.2pt];
\end{tikzpicture}
\figcap{上述混合光的检偏曲线：非偏振分量使最小透射强度保持为正。}
\label{comp:A5:analyzer-result}
\end{minipage}\end{center}

\paragraph{偏振元件怎样作用于相干矩阵}
若线性光学元件的Jones矩阵为 $J$，则 $\vect A'=J\vect A$，从定义即得
$P'=JPJ^\dagger$。理想无损波片的 $J$ 为酉矩阵，保留 $P$ 的本征值、总强度和偏振度，
但可以改变偏振椭圆。以快慢轴为坐标轴、忽略公共相位，可按相位约定写成
\[
 J_{\rm ret}=\begin{pmatrix}\e^{\ii\delta}&0\\0&1\end{pmatrix}.
\]
$\delta=\pi$ 为半波片，$\delta=\pi/2$ 为四分之一波片。
例如等幅同相输入 $(1,1)^{\mathsf T}/\sqrt2$ 经过后一种矩阵变为
$(\ii,1)^{\mathsf T}/\sqrt2$，两个分量相差四分之一周期，成为圆偏振光。
其旋向还须连同传播方向和时间因子约定说明。


\subsection{晶体光学}

各向异性线性介质满足
\[
 \vect D=\boldsymbol\varepsilon\vect E.
\]

\begin{advtheorem}{无耗散介质的介电张量对称性}\label{advA5:thm:epsilon-symmetric}
在无色散或准静态、无耗散的可逆线性介质中，若电场能密度是状态函数，则
\[
 \boxed{\boldsymbol\varepsilon^{\mathsf T}
 =\boldsymbol\varepsilon}.
\]
\end{advtheorem}

\noindent\textit{证明。}
从零场准静态建立电场，电场能密度的增量为
$\dd u=\vect E\cdot\dd\vect D$。若 $u$ 是状态函数，定义电共能
$u^*(\vect E)=\vect E\cdot\vect D-u$，则
\[
 \dd u^*=\vect D\cdot\dd\vect E,\qquad
 D_i=\frac{\partial u^*}{\partial E_i}.
\]
因此二阶偏导可交换给出
\[
 \varepsilon_{ij}=\frac{\partial D_i}{\partial E_j}
 =\frac{\partial^2u^*}{\partial E_j\partial E_i}
 =\frac{\partial^2u^*}{\partial E_i\partial E_j}
 =\varepsilon_{ji}.
\]
对这种线性介质才进一步有 $u=u^*=\tfrac12\vect E\cdot\vect D$。
若存在吸收，频域介电张量可为复数；互易性与无耗散是不同条件。
反 Hermitian 部分决定单色场的平均吸收，色散介质的储能还涉及介电响应的频率导数，不能简单照搬上述实对称、无色散结论。

实对称介电张量可正交对角化为
$\operatorname{diag}(\varepsilon_x,\varepsilon_y,\varepsilon_z)$。
三者相等为各向同性晶体；两者相等为单轴晶体；三者互异为双轴晶体。
相位法线方向 $\hat{\vect p}=\vect k/k$ 与能流方向
$\hat{\vect r}\parallel\vect S$ 一般不重合，所以相速度和射线速度需分别讨论。

\begin{advtheorem}{Fresnel 相速度面与射线速度面}\label{advA5:thm:fresnel-crystal}
在主轴坐标中令
$v_i=1/\sqrt{\mu\varepsilon_i}$。平面波的相速度 $v_p$ 与法线方向余弦
$p_i$ 满足
\[
 \boxed{
 \frac{p_x^2}{v_p^2-v_x^2}
 +\frac{p_y^2}{v_p^2-v_y^2}
 +\frac{p_z^2}{v_p^2-v_z^2}=0}.
\]
射线速度 $v_r$ 与能流方向余弦 $r_i$ 满足对偶式
\[
 \boxed{
 \frac{r_x^2}{v_r^{-2}-v_x^{-2}}
 +\frac{r_y^2}{v_r^{-2}-v_y^{-2}}
 +\frac{r_z^2}{v_r^{-2}-v_z^{-2}}=0}.
\]
\end{advtheorem}

上述分式须先乘去分母再按连续极限解释。相速度的完整多项式为 $\sum_i p_i^2\prod_{j\ne i}(v_p^2-v_j^2)=0$；沿主轴或出现相等本征值时，不能在分式形式中直接约去零因子而丢失偏振支。单轴情形可因式分解为 $[v_p^2-v_o^2][v_p^2-v_o^2\cos^2\theta-v_e^2\sin^2\theta]=0$，两支都须保留。射线速度的对偶分式亦须作相同处理。

\noindent\textit{证明。}
对 $\e^{\ii(\vect k\cdot\vect r-\omega t)}$ 代入 Maxwell 方程，
\[
 \vect k\times(\vect k\times\vect E)
 +\omega^2\mu\boldsymbol\varepsilon\vect E=0.
\]
写成分量并令 $v_p=\omega/k$：
\[
 \left(\varepsilon_i-\frac1{\mu v_p^2}\right)E_i
 +\frac{p_i}{\mu v_p^2}(\hat{\vect p}\cdot\vect E)=0.
\]
消去 $E_i$，要求非零解的行列式为零，整理即得第一式。
能流方向与波矢法线面互为 Legendre 对偶；用群速度
$\vect v_g=\nabla_{\vect k}\omega$ 求相速度面的法向，
再消去 $p_i$ 得第二式。

对单轴晶体，取 $v_x=v_y=v_o$、$v_z=v_e$。一支为寻常光
$v_{po}=v_o$，电场垂直主截面；另一支为非寻常光
\[
 \boxed{v_{pe}^2=v_o^2\cos^2\theta+v_e^2\sin^2\theta},
\]
其中 $\theta$ 为波法线 $\vect k$ 与光轴的夹角；其 $\vect k$ 与能流一般不平行。界面折射时，频率和切向波矢连续；
因此应先在相速度面上求允许 $\vect k$，再由曲面的法向求射线方向，
不能把各向同性 Snell 定律直接套给能流方向。

% AdvancedOpticsQuality20261023
\section{干涉应用：镀膜、测距与曲面条纹}
以下四题讨论增透膜、干涉测距与曲面条纹。涉及反射相位、视场边界和测量误差的假设在解答中分别说明。

\begin{advexamplebox}{半球形砷化镓发光器件的增透膜}\label{advA5:ex:source-ar}
砷化镓发光器件制成半球形，发光区位于球心，真空波长为 $930\,\mathrm{nm}$。砷化镓折射率 $n_1=3.4$，外部空气折射率 $n_2=1$。求未镀膜时的表面反射率；若在半球面镀一层均匀增透膜，求消除该波长反射所需的膜折射率与最小厚度。忽略吸收，发光区视为点源。

球半径远大于膜厚和波长，采用局部平面膜近似。

\QAOFrame{\begin{tikzpicture}[>=Latex,font=\small]
 \fill[DeepBlue!5](-2,0)arc(180:0:2)--cycle;
 \draw[thick](-2,0)arc(180:0:2);\draw[thick](-2.2,0)arc(180:0:2.2);
 \draw(-2.6,0)--(2.8,0);\fill(0,0)circle(2pt)node[below]{$O$（发光区）};
 \foreach \a in {35,90,145}{\draw[DeepBlue,->,thick](0,0)--(\a:2.7);}
 \node at(-.5,.85){$n_1$};\node at(2.5,2){$n_2$};
 \draw(50:2.1)--(2.4,2.6)--(3.9,2.6)node[right]{膜：$n_f,t$};
 \draw[thin](2,0)--(2,-.8) (2.2,0)--(2.2,-.8);\draw[->](1.4,-.5)--(2,-.5);\draw[->](2.8,-.5)--(2.2,-.5);\node[below]at(2.1,-.8){$t$};
\end{tikzpicture}}{球心点源与同心半球薄膜；膜厚为放大的示意量。}{quality:A5:q12}
\end{advexamplebox}
\begin{advsolution}
球心发出的每条光线均沿局部法线入射，故整个半球面都可用正入射公式。
\begin{center}\begin{minipage}{\linewidth}\centering
% Design: 10cm optical section; 9pt labels; radial rays preserve normal incidence.
\begin{tikzpicture}[>=Latex,font=\small]
 \fill[DeepBlue!5](-2,0)arc(180:0:2)--cycle;
 \draw[thick](-2,0)arc(180:0:2);\draw[thick](-2.2,0)arc(180:0:2.2);
 \draw(-2.6,0)--(2.8,0);\fill(0,0)circle(2pt)node[below]{$O$（发光区）};
 \foreach \a in {35,90,145}{\draw[DeepBlue,->,thick](0,0)--(\a:2.7);}
 \node at(-.5,.85){$n_1$};\node at(2.5,2){$n_2$};
 \draw(50:2.1)--(2.4,2.6)--(3.9,2.6)node[right]{膜：$n_f,t$};
 \draw[thin](2,0)--(2,-.8) (2.2,0)--(2.2,-.8);\draw[->](1.4,-.5)--(2,-.5);\draw[->](2.8,-.5)--(2.2,-.5);\node[below]at(2.1,-.8){$t$};
\end{tikzpicture}
\figcap{原题半球镀膜截面。膜厚为示意放大；光源位于共同球心，各光线正入射。\label{advA5:fig:source-ar}}
\end{minipage}\end{center}
未镀膜时电场反射系数及能量反射率为
\[
 r=\frac{n_1-n_2}{n_1+n_2},\qquad
 \boxed{\mathcal R=|r|^2=\left(\frac{2.4}{4.4}\right)^2=0.29752\approx29.75\%}.
\]
镀膜后令 $r_{1f}=(n_1-n_f)/(n_1+n_f)$、$r_{f2}=(n_f-n_2)/(n_f+n_2)$，薄膜的精确总反射系数为
\[
 r_{\rm film}=\frac{r_{1f}+r_{f2}\e^{2\ii\beta}}{1+r_{1f}r_{f2}\e^{2\ii\beta}},\qquad
 \beta=\frac{2\pi n_ft}{\lambda_0}.
\]
要使分子为零，先匹配两反射振幅 $r_{1f}=r_{f2}$，得 $n_f^2=n_1n_2$；再使往返相位相差 $\pi$，得
\[
 \boxed{n_f=\sqrt{3.4}=1.8439},\qquad
 \boxed{t_{\min}=\frac{930}{4\sqrt{3.4}}\,\mathrm{nm}=126.09\,\mathrm{nm}}.
\]
一般厚度为 $(2k+1)\lambda_0/(4n_f)$。这里两次反射均从高折射率一侧射向低折射率一侧，反射相位相同；相消来自膜内往返相位。

图~\ref{advA5:fig:source-ar}中球心到每个界面交点的连线都沿法线，才可把平面膜正入射公式逐点用于半球面。膜厚须远小于球半径，使局部平面膜近似成立；偏离设计波长或把点源移离球心后，零反射条件一般不再同时成立。
\end{advsolution}

\begin{advexamplebox}{迈克尔逊干涉仪的环数、分辨率与量程}\label{advA5:ex:source-michelson}
用 $\lambda=632.8\,\mathrm{nm}$ 的 He--Ne 激光照明迈克尔逊干涉仪。望远镜固定视场内原有20个亮环，移动一臂的反射镜后，亮环向中央收拢，共吞入20环，视场内还剩10个亮环。求原来中心亮斑的级次及后来第五亮环的角半径。

再用此仪器测长。激光谱线宽度 $\Delta\lambda=10^{-3}\,\mathrm{nm}$，条纹读数可分辨 $1/8$ 条纹，求测长分辨率并讨论量程。

分别讨论视场边界恰落亮环与仅给定视场内环数两种计数含义。

\QAOFrame{\QAOMichelsonQuestion}{等倾干涉的两个等效镜面与同一视场内的吞环计数；虚线表示等效延长段。}{quality:A5:q13}
\end{advexamplebox}
\begin{advsolution}
设两臂等效几何长度差为 $d>0$，等倾干涉光程差为 $2d\cos\theta$。中心级次 $m=2d/\lambda$，第 $j$ 个非中心亮环满足 $m\cos\theta_j=m-j$。收拢吞环表明 $d$ 减小，故 $m_1=m_0-20$。
\begin{center}\begin{minipage}{\linewidth}\centering
% Design: 11cm two-panel illustration, no instrument beam paths invented.
\begin{tikzpicture}[>=Latex,font=\small,scale=.9]
 \draw[->](-2.6,0)--(2.8,0)node[right]{光轴};
 \draw[DeepBlue,->,thick](-2.2,0)--(1.6,1.7);
 \draw[DeepBlue,->,thick](-2.2,0)--(1.6,-1.7);
 \draw(-.9,0)arc(0:24:1.3);\node at(-.7,.33){$\theta$};
 \draw[dashed](1.6,-1.9)--(1.6,1.9);\node[above]at(1.6,1.9){固定视场};
 \begin{scope}[xshift=6cm]
  \foreach \r in {.45,.8,1.08,1.32,1.53}{\draw[DeepBlue](0,0)circle(\r);}
  \draw[gray,dashed](0,0)circle(1.75);\fill(0,0)circle(1.6pt);
  \draw[->,thick](-1.6,0)--(-.9,0);\node[above]at(0,1.9){亮环向中心收拢};
  \node[below]at(0,-1.9){仅示意收缩方向，非20环计数图};
 \end{scope}
\end{tikzpicture}
\figcap{等倾环的观察角与吞环方向。圆环从中心向外计数，中心亮斑另计。\label{advA5:fig:source-michelson}}
\end{minipage}\end{center}
\textbf{通常教材约定：初、末视场边界都恰落在亮环上。}若视场半角为 $\alpha$，则
\[
 20=m_0(1-\cos\alpha),\qquad10=(m_0-20)(1-\cos\alpha).
\]
相减并代回，得到 $m_0=40$、$m_1=20$、$\alpha=60^\circ$。于是
\[
 \boxed{\theta_5=\arccos\frac{20-5}{20}=41.410^\circ}.
\]
此角并不小，不能把 $\cos\theta\simeq1-\theta^2/2$ 当作精确关系。

\textbf{原题的严格边界问题。}只说“视场里有多少环”而未规定边界，计数应为
\[
 20=\lfloor m_0(1-\cos\alpha)\rfloor,\qquad
 10=\lfloor(m_0-20)(1-\cos\alpha)\rfloor.
\]
即便初末中心均为亮斑，也允许 $m_0=39,40,\ldots,44$，因此40是上述教材边界约定下的答案，并非原文条件下严格唯一的答案。

镜移 $\delta d$ 使光程差改变 $2\delta d$，一整条纹对应镜移 $\lambda/2$。所以读数分辨率为
\[
 \boxed{\delta d=\frac{\lambda}{16}=39.55\,\mathrm{nm}}.
\]
量程有两种不同含义。由线宽估算相干长度 $L_c\sim\lambda^2/\Delta\lambda=0.40044\,\mathrm m$，可维持干涉的臂差量程量级为
\[
 \boxed{|d|_{\rm coh}\sim\frac{\lambda^2}{2\Delta\lambda}=0.20022\,\mathrm m}.
\]
具体数值因谱线形状和最低可见度而异，不是一个严格截止长度。
若\emph{另将} $\Delta\lambda$ 作为波长标定不确定度，并要求 $d\Delta\lambda/\lambda\le\delta d$，则保证上述绝对误差的量程为
\[
 \boxed{|d|_{\rm err}\lesssim\frac{\lambda^2}{16\Delta\lambda}=2.5027\,\mathrm{cm}}.
\]
线宽本身不等于中心波长标定误差，二者不可无条件替换。

图~\ref{advA5:fig:source-michelson}中的固定视场边界决定环数的取整。镜移使中心级次连续改变，吞入20环只给出级次变化，严格唯一的整数答案还需要初末中心亮、视场边界亮等附加约定。大于40度的结果必须使用余弦精确关系。
\end{advsolution}

\begin{advexamplebox}{半球凸起附近的条纹间距与对比度}\label{advA5:ex:source-bump}
波长为 $\lambda$ 的平行光沿 $-z$ 方向入射，电场沿 $y$ 方向。在 $xy$ 平面上有半径为 $R$ 的半球凸起；光在球面理想反射后与原入射光在平面上干涉。原图截面为 $xz$ 平面，球面反射点的半径与 $+x$ 轴夹角为 $\theta$。求对应平面位置的局部条纹间距；将观察点绕 $z$ 轴转过方位角 $\phi$，讨论条纹对比度
\[
 \gamma=\frac{I_{\max}-I_{\min}}{I_{\max}+I_{\min}}.
\]
按局部平面波及电场平方强度模型计算，说明所用反射偏振模型。

\QAOFrame{\begin{tikzpicture}[>=Latex,font=\small,scale=.9]
 \fill[gray!8](-1.8,0)arc(180:0:1.8)--cycle;
 \draw[thick](-1.8,0)arc(180:0:1.8);
 \draw[->](-2.2,0)--(3.3,0)node[right]{$x$};\draw[->](0,0)--(0,2.7)node[above]{$z$};
 \coordinate(P)at(1.63135,.76071);\coordinate(Q)at(2.53787,0);
 \draw[gray,dashed](0,0)--(2.15,1.0026);
 \draw[DeepBlue,->,thick](1.63135,2.6)--(P);
 \draw[DeepBlue,->,thick](P)--(Q);
 \draw[DeepBlue,->,thick](2.53787,2.6)--(Q);
 \fill(P)circle(1.5pt);\draw[thin](P)--(1.95,1.12);\node[above]at(1.95,1.12){$P$};\fill(Q)circle(1.5pt)node[below]{$Q$};
 \node[below left]at(0,0){$O$};\node at(.7,.55){$R$};
 \draw(.75,0)arc(0:25:.75);\node at(1,.18){$\theta$};
 \node[left]at(1.6,2.3){$\boldsymbol E\parallel y$};
 \draw[<->](0,-.65)--(2.53787,-.65)node[midway,below]{$\rho$};
 \begin{scope}[xshift=7cm,yshift=.5cm]
  \draw[->](-1.4,0)--(1.8,0)node[right]{$x$};\draw[->](0,-1.3)--(0,1.8)node[above]{$y,\boldsymbol E_i$};
  \draw[gray](0,0)circle(.6);\draw[dashed](0,0)circle(1.3);
  \draw[DeepBlue,->](0,0)--(40:1.3)node[above right]{$Q_\phi$};
  \draw(.8,0)arc(0:40:.8);\node at(1,.3){$\phi$};
  \node[below]at(0,-1.6){俯视：同一半径，不同方位};
 \end{scope}
\end{tikzpicture}}{球面反射点 $P$、平面观察点 $Q$ 与俯视方位角；只示几何，不给条纹公式。}{quality:A5:q14}
\end{advexamplebox}
\AdvHideHeadingtrue
\begin{advsolution}
\noindent\begin{minipage}{\linewidth}{\bfseries\color{DeepBlue}解答}\par
\begin{center}\begin{minipage}{\linewidth}\centering
% Design: 12.5cm, 9pt. Geometry theta=25deg; C=cos50deg; Qx=Rcos25/C.
\begin{tikzpicture}[>=Latex,font=\small,scale=.9]
 \fill[gray!8](-1.8,0)arc(180:0:1.8)--cycle;
 \draw[thick](-1.8,0)arc(180:0:1.8);
 \draw[->](-2.2,0)--(3.3,0)node[right]{$x$};\draw[->](0,0)--(0,2.7)node[above]{$z$};
 \coordinate(P)at(1.63135,.76071);\coordinate(Q)at(2.53787,0);
 \draw[gray,dashed](0,0)--(2.15,1.0026);
 \draw[DeepBlue,->,thick](1.63135,2.6)--(P);
 \draw[DeepBlue,->,thick](P)--(Q);
 \draw[DeepBlue,->,thick](2.53787,2.6)--(Q);
 \fill(P)circle(1.5pt);\draw[thin](P)--(1.95,1.12);\node[above]at(1.95,1.12){$P$};\fill(Q)circle(1.5pt)node[below]{$Q$};
 \node[below left]at(0,0){$O$};\node at(.7,.55){$R$};
 \draw(.75,0)arc(0:25:.75);\node at(1,.18){$\theta$};
 \node[left]at(1.6,2.3){$\boldsymbol E\parallel y$};
 \draw[<->](0,-.65)--(2.53787,-.65)node[midway,below]{$\rho$};
 \begin{scope}[xshift=7cm,yshift=.5cm]
  \draw[->](-1.4,0)--(1.8,0)node[right]{$x$};\draw[->](0,-1.3)--(0,1.8)node[above]{$y,\boldsymbol E_i$};
  \draw[gray](0,0)circle(.6);\draw[dashed](0,0)circle(1.3);
  \draw[DeepBlue,->](0,0)--(40:1.3)node[above right]{$Q_\phi$};
  \draw(.8,0)arc(0:40:.8);\node at(1,.3){$\phi$};
  \node[below]at(0,-1.6){俯视：同一半径，不同方位};
 \end{scope}
\end{tikzpicture}
\figcap{恢复原图中从水平轴量起的 $\theta$。仅 $0<\theta<\pi/4$ 的反射束能向下到达外侧平面。\label{advA5:fig:source-bump}}
\end{minipage}\end{center}\end{minipage}

记 $C=\cos2\theta>0$、$S=\sin2\theta$。镜面反射定律给出
\[
 P=(R\cos\theta,0,R\sin\theta),\qquad
 \hat{\boldsymbol k}_r=S\boldsymbol e_\rho-C\boldsymbol e_z.
\]
因此从 $P$ 到平面的光程 $\ell=R\sin\theta/C$，落点半径及相对原入射波的几何光程差为
\[
 \rho=R\cos\theta+\ell S=\frac{R\cos\theta}{C},\qquad
 \Delta=\ell-R\sin\theta=\frac{2R\sin^3\theta}{C}.
\]
由参数微分得 $\dd\Delta/\dd\rho=S$。一个条纹周期要求相位变化 $2\pi$，故
\[
 \boxed{\Delta\rho\simeq\frac{\lambda}{\sin2\theta}}.
\]
这是局部条纹间距，要求波长远小于曲率尺度；反射相位若随位置变化，还需将其梯度计入。

\textbf{先处理扩束，再处理偏振。}入射环带半径为 $b=R\cos\theta$。无吸收时，入射环带功率等于落在平面环带上的功率：
\[
 \frac{I_{r,z}}{I_i}=\frac{b|\dd b|}{\rho\,\dd\rho}=\frac{C^3}{C+2}.
\]
$I_{r,z}$ 是穿过水平面的法向能流；反射光束自身截面上的强度为 $I_r=I_{r,z}/C$。电场平方干涉模型中所需的强度比因而是
\[
 q=\frac{I_r}{I_i}=\frac{C^2}{C+2},\qquad
 \gamma=\frac{2\sqrt q}{1+q}\,|\hat{\boldsymbol e}_i^*\!\cdot\hat{\boldsymbol e}_r|.
\]
令局部 $s$ 基矢为 $\boldsymbol e_\phi$，$p$ 基矢分别取 $\boldsymbol s\times\hat{\boldsymbol k}_{i,r}$。若单位模的反射系数为 $r_s,r_p$，直接投影原入射电场 $\boldsymbol E_i\parallel y$，得到
\[
 \boxed{\gamma(\phi)=\frac{2\sqrt q}{1+q}
 \left|r_s\cos^2\phi+r_p C\sin^2\phi\right|}.
\]
理想导体表面在这一基底约定下 $r_s=-1,r_p=+1$，因此
\[
 \boxed{\gamma_{\rm PEC}(\phi)=\frac{2\sqrt q}{1+q}
       |\cos^2\phi-C\sin^2\phi|}.
\]
若另采用 $r_s=r_p$ 的简化反射模型，绝对值内则为 $\cos^2\phi+C\sin^2\phi$。原文仅规定几何光学反射，不能据此唯一确定两偏振分量的相对相位。一般材料还须使用实际振幅反射系数。

原图方位 $\phi=0$ 时，$\gamma_0=2\sqrt q/(1+q)$；由于扩束，即使无吸收也通常小于1。靠近凸起边缘的几何光学极限为 $\sqrt3/2$，远离凸起时趋于0。若另人为设两束等强，应将上式的前因子改成1，并明确这是忽略扩束的近似。

图~\ref{advA5:fig:source-bump}的 $P\to Q$ 方向决定径向波矢 $k\sin2\theta$，它与入射波在平面内的零切向波矢之差，给出局部条纹间距。俯视图中的 $\phi$ 只改变偏振投影；扩束引起的强度比仍由相同 $\rho$ 的环带面积比计算。
\end{advsolution}

\begin{advexamplebox}{相互垂直的圆柱面所产生的干涉条纹}\label{advA5:ex:source-cylinders}
两相向凸圆柱面的半径分别为 $R_1,R_2$，轴线互相垂直，在原点 $O$ 处相切。竖直平行光向下入射，观察两相向表面反射光形成的等厚干涉。求条纹形状与级次的关系。约定下圆柱轴沿 $x$，半径为 $R_1$；上圆柱轴沿 $y$，半径为 $R_2$。两表面间为空气，按普通透明材料两反射光中只有一次半波损失计算。

\QAOFrame{\QAOCylindersQuestion}{两个正交截面中不同的曲率；相切点为原点，箭头表示竖直入射光。}{quality:A5:q15}
\end{advexamplebox}
\AdvHideHeadingtrue
\begin{advsolution}
\noindent\begin{minipage}{\linewidth}{\bfseries\color{DeepBlue}解答}\par
\begin{center}\begin{minipage}{\linewidth}\centering
% Design: 12cm three coordinated panels, 9pt. Curvature sum shown in two sections.
\begin{tikzpicture}[>=Latex,font=\small,scale=.95]
 % Source perspective reconstruction: lower axis x, upper axis y; tangent at O.
 \begin{scope}[shift={(4.2,3.8)},x={(-.55cm,-.33cm)},y={(.95cm,0cm)},z={(0cm,.85cm)}]
  \draw[thick,DeepBlue,domain=-.85:.85,samples=40]plot(-2.3,\x,{-.5*\x*\x});
  \draw[thick,DeepBlue,domain=-.85:.85,samples=40]plot(2.3,\x,{-.5*\x*\x});
  \foreach \v in {-.85,.85}{\draw[thick,DeepBlue](-2.3,\v,-.36125)--(2.3,\v,-.36125);}
  \draw[dashed,gray](-2.3,0,0)--(2.3,0,0);
  \draw[thick,domain=-.85:.85,samples=40]plot(\x,-2.3,{.5*\x*\x});
  \draw[thick,domain=-.85:.85,samples=40]plot(\x,2.3,{.5*\x*\x});
  \foreach \v in {-.85,.85}{\draw[thick](\v,-2.3,.36125)--(\v,2.3,.36125);}
  \draw[dashed,gray](0,-2.3,0)--(0,2.3,0);
  \draw[->](0,0,0)--(3,0,0)node[below]{$x$};
  \draw[->](0,0,0)--(0,3,0)node[right]{$y$};
  \draw[->](0,0,0)--(0,0,2)node[above]{$z$};
  \draw[DeepBlue,->,thick](0,-.3,1.8)--(0,-.3,.15);
  \fill(0,0,0)circle(1.5pt)node[below right]{$O$};
  \node[left]at(1.2,-1.7,-.3){下：$R_1$};\node[right]at(-.85,2.3,.5){上：$R_2$};
 \end{scope}
 \begin{scope}
  \draw[->](-1.4,0)--(1.5,0)node[right]{$x$};\draw[->](0,-.45)--(0,1.5)node[above]{$z$};
  \draw[DeepBlue,thick,domain=-1.2:1.2,samples=50]plot(\x,{.6*\x*\x});
  \draw[thick](-1.2,0)--(1.2,0);\node[above]at(0,1.9){$xz$ 截面};
  \node[below,align=center]at(0,-.55){上表面弯曲，半径 $R_2$\\下表面沿轴线平直};
 \end{scope}
 \begin{scope}[xshift=4.2cm]
  \draw[->](-1.4,0)--(1.5,0)node[right]{$y$};\draw[->](0,-1.15)--(0,1.5)node[above]{$z$};
  \draw[DeepBlue,thick,domain=-1.2:1.2,samples=50]plot(\x,{-.6*\x*\x});
  \draw[thick](-1.2,0)--(1.2,0);\node[above]at(0,1.9){$yz$ 截面};
  \node[below,align=center]at(0,-1.25){下表面弯曲，半径 $R_1$};
 \end{scope}
 \begin{scope}[xshift=8.5cm]
  \foreach \r in {.5,.7071,.8660,1}{\draw[DeepBlue](0,0)ellipse({1.3*\r} and {\r});}
  \draw[->](-1.5,0)--(1.7,0)node[right]{$x$};\draw[->](0,-1.2)--(0,1.5)node[above]{$y$};
  \node[above]at(0,1.9){近轴等厚暗纹};\node[below]at(0,-1.4){$R_2>R_1$ 的示例};
 \end{scope}
\end{tikzpicture}
\figcap{上：原题两相向圆柱面的透视关系，箭头为入射光；下：正交截面及条纹俯视图。两项曲率共同增加气隙厚度。\label{advA5:fig:source-cylinders}}
\end{minipage}\end{center}\end{minipage}

按题设比较两相向表面的反射光，采用空气膜等厚干涉模型；两束反射光必须在同一观察区域叠加。

取相切点为 $z=0$，两相向表面的高度分别为
\[
 z_-=-R_1+\sqrt{R_1^2-y^2},\qquad z_+=R_2-\sqrt{R_2^2-x^2}.
\]
两者之差才是膜厚；注意本题曲率\emph{相加}：
\[
 t=z_+-z_-=R_1-\sqrt{R_1^2-y^2}+R_2-\sqrt{R_2^2-x^2}
 \simeq\frac{y^2}{2R_1}+\frac{x^2}{2R_2}.
\]
近轴、局部近似正入射时，空气膜的几何光程差约为 $2t$。计入一次半波损失，暗纹满足 $2t=m\lambda$，亮纹满足 $2t=(m+\tfrac12)\lambda$，故近轴下
\[
 \boxed{\frac{x^2}{R_2}+\frac{y^2}{R_1}=
 \begin{cases}m\lambda,&\text{暗纹},\\(m+\tfrac12)\lambda,&\text{亮纹},\end{cases}
 \qquad m=0,1,2,\ldots}
\]
第 $m\ge1$ 级暗纹的半轴为 $\sqrt{m\lambda R_2}$ 与 $\sqrt{m\lambda R_1}$，轴比固定为 $\sqrt{R_2/R_1}$；相邻暗纹半轴平方之差为 $\lambda R_2$ 与 $\lambda R_1$。两半径相等时退化为圆环。

椭圆是近轴近似。根式给出精确几何膜厚；较大孔径时，还须计算局部折射角和实际反射光路，不能仅把近轴膜厚换成根式而继续使用 $2t$。若夹层折射率为 $n$，右端改为相应级次乘 $\lambda/n$；若两次反射相位相同，则明暗条件互换。原题未给出材料，以上明暗级次依赖题干中明示的普通空气膜假设。

图~\ref{advA5:fig:source-cylinders}的两个正交截面中，曲率分别来自不同表面；膜厚是上表面减下表面，故两项都为正。椭圆长轴沿半径较大的一侧，半轴平方随级次等差，提供了图形方向和系数的直接核对。
\end{advsolution}


\clearpage\endgroup

\begingroup
\newif\ifAdvHideHeading
\RenewDocumentEnvironment{advexamplebox}{m}{%
  \refstepcounter{advexample}%
  \begin{tcolorbox}[
    enhanced,unbreakable,colback=white,colframe=RuleGray,boxrule=0.5pt,arc=0mm,
    title={例~\theadvexample\quad #1},fonttitle=\bfseries,
    coltitle=Ink,colbacktitle=PaleBlue,borderline west={1.4pt}{0pt}{DeepBlue}]
}{\end{tcolorbox}}

\RenewDocumentEnvironment{advsolution}{+b}{%
  \par\smallskip\begingroup\small\noindent
  \ifAdvHideHeading\global\AdvHideHeadingfalse\else\LectureSolutionHeading\fi #1\par\endgroup\smallskip
}{}

\RenewDocumentEnvironment{advdefinition}{m +b}{%
  \refstepcounter{advdefinition}%
  \begin{tcolorbox}[
    enhanced,unbreakable,colback=white,colframe=RuleGray,boxrule=0.5pt,arc=0mm,
    title={定义~\theadvdefinition\quad #1},fonttitle=\bfseries,
    coltitle=Ink,colbacktitle=PaleGold,borderline west={1.2pt}{0pt}{Gold}]
  #2
  \end{tcolorbox}
}{}

\RenewDocumentEnvironment{advtheorem}{m +b}{%
  \refstepcounter{advtheorem}%
  \begin{tcolorbox}[
    enhanced,unbreakable,colback=white,colframe=RuleGray,boxrule=0.5pt,arc=0mm,
    title={定理~\theadvtheorem\quad #1},fonttitle=\bfseries,
    coltitle=Ink,colbacktitle=PaleBlue,borderline west={1.2pt}{0pt}{DeepBlue}]
  #2
  \end{tcolorbox}
}{}


\AdvancedUnit{6}{近代物理}
% AdvancedEMQuality20261023
\section{狭义相对论}

\subsection{时空观与相对性原理}

\begin{advtheorem}{Newton 运动定律的惯性系表述}\label{advA6:thm:newton-laws}
\begin{enumerate}
\item 在惯性系中，不受合外力的质点保持静止或匀速直线运动；
\item 质点动量的变化率等于所受合外力；质量恒定、低速时写成
$\vect F=m\vect a$；
\item 两质点间的相互作用力等大、反向；对满足中心力模型的相互作用还共线。
\end{enumerate}
\end{advtheorem}

\begin{advtheorem}{Einstein 狭义相对论公设}\label{advA6:thm:einstein-postulates}
\begin{enumerate}
\item 所有惯性系中的物理定律形式相同；
\item 真空光速 $c$ 与光源、观察者的惯性运动状态无关。
\end{enumerate}
\end{advtheorem}

Newton--Cartan 时空把绝对时间与三维 Euclid 空间分开；
Minkowski 时空则把事件写成
$x^\mu=(ct,x,y,z)$，以
\[
 \dd s^2=c^2\dd t^2-\dd x^2-\dd y^2-\dd z^2
\]
作为不变量。广义相对论进一步允许度规随位置变化，但本节只讨论平直时空。

沿 $x$ 方向、速度为 $v$ 的 Lorentz 推促写成
\[
\begin{pmatrix}ct'\\x'\\y'\\z'\end{pmatrix}
=
\begin{pmatrix}
\gamma&-\gamma\beta&0&0\\
-\gamma\beta&\gamma&0&0\\
0&0&1&0\\0&0&0&1
\end{pmatrix}
\begin{pmatrix}ct\\x\\y\\z\end{pmatrix},
\qquad
\beta=\frac vc,\quad \gamma=\frac1{\sqrt{1-\beta^2}}.
\]
任意四矢量 $A^\mu=(A^0,\vect A)$ 服从同一线性变换，伪内积
$A^\mu B_\mu=A^0B^0-\vect A\cdot\vect B$ 不变。类时四矢量可在某一惯性系中把空间分量变为零，
该系称为其本征系。

\subsection{常见四矢量与相对论动力学}

四位移、四速度与四动量为
\[
 \dd x^\mu=(c\dd t,\dd\vect r),\qquad
 U^\mu=\frac{\dd x^\mu}{\dd\tau}=\gamma(c,\vect v),
\]
\[
 p^\mu=mU^\mu=\left(\frac Ec,\vect p\right),\qquad
 E^2-p^2c^2=m^2c^4,
\]
\[
 \vect p=\gamma m\vect v,\qquad E=\gamma mc^2,\qquad
 \vect v=\frac{\vect p c^2}{E}.
\]
四加速度与四力
\[
 A^\mu=\frac{\dd U^\mu}{\dd\tau},\qquad
 K^\mu=\frac{\dd p^\mu}{\dd\tau}
=\gamma\left(\frac{\vect F\cdot\vect v}{c},\vect F\right)
\]
均与 $U^\mu$ 正交。三维形式仍为
\[
 \vect F=\frac{\dd\vect p}{\dd t},\qquad
 \vect F\cdot\dd\vect r=\dd E.
\]
把加速度分解为平行、垂直于瞬时速度的分量，
\[
 a_{\parallel,0}=\gamma^3a_\parallel,\qquad
 a_{\perp,0}=\gamma^2a_\perp,
\]
其中下标 $0$ 表示瞬时本征惯性系。快度 $\vartheta$ 满足
$\beta=\tanh\vartheta$、$\gamma=\cosh\vartheta$，共线速度合成变为快度相加。

\begin{advexamplebox}{加速杆的收缩率、本征加速度与信号频移}\label{advA6:ex:accelerating-rod}
本征长度为 $L$ 的杆以速度 $v$ 沿杆方向运动，随后发生缓慢纵向加速。
地面系瞬时加速度为 $a$，且 $aL\ll c^2$。涉及有限杆时采用局部短杆近似，
并分别说明保持刚性及频移展开所需的额外条件。

\begin{enumerate}
\item 求局部近似下地面系同时测得的杆长 $\ell\simeq L/\gamma$ 的变化率，
并写成 $-v^{-1}\dd\ell/\dd t=\lambda\gamma^\mu(aL/c^2)^\nu$，求三个无量纲常数；
\item 把杆视为质点，求纵向本征加速度 $a_0$、本征力 $F_0$ 与地面系量的关系；
\item 在保持 Born 刚性的缓慢加速近似下，尾端发出本征频率 $\nu_1$ 的信号，
求头端收到的频率 $\nu_2$。
\end{enumerate}

\begin{center}\begin{minipage}{.96\linewidth}\centering
\begin{tikzpicture}[>=Latex,font=\small,scale=.85,every node/.style={fill=none}]

\draw[very thick] (0,0)--(3,0);\fill (0,0) circle (1.5pt) node[below]{尾};\fill (3,0) circle (1.5pt) node[below]{头};\draw[->] (.3,1)--(1.6,1) node[midway,above]{$v$};\draw[->] (2,1)--(3.2,1) node[midway,above]{$a$};\node[below] at (1.5,-.6){本征长 $L$；地面同时长度 $\ell$ 待求};

\end{tikzpicture}
\figcap{题图：沿杆方向的速度与缓慢加速度；本征长和地面同时长度属于不同测量口径。}\label{advA6:quality:A6-relativity:q1}
\end{minipage}\end{center}
\end{advexamplebox}
\global\AdvHideHeadingtrue
\begin{advsolution}
\begin{center}\begin{minipage}{.96\linewidth}\centering
% EMAdvFigureHeading
\raggedright{\bfseries\color{DeepBlue}解答}\par\smallskip\centering

\begin{tikzpicture}[>=Latex,font=\small,scale=.85,every node/.style={fill=none}]

\draw[->] (-.3,0)--(4,0) node[right]{$x$};\draw[->] (0,-.2)--(0,3.3) node[above]{$ct$};
\draw[thick] (.3,0)..controls(.5,1.7)and(.9,2.5)..(1.3,3);\draw[thick] (2.1,0)..controls(2.3,1.7)and(2.7,2.5)..(3.1,3);
\draw[dashed] (0,1.4)--(3.5,1.4) node[right]{地面同时截面};\draw[DeepBlue,->] (.35,.2)--(2.5,2.35) node[midway,above left]{光信号};
\node[below] at (.3,0){尾};\node[below] at (2.1,0){头};

\end{tikzpicture}
\figcap{解图：局部短杆世界线的示意，水平截面用于测长，斜光线用于频移；曲线不作精确Born刚性比例图。}\label{advA6:quality:A6-relativity:a1}
\end{minipage}\end{center}

第一问保留杆长 $L$ 的一阶项，$v,a$ 取尾端瞬时值。
对保持Born刚性的加速杆，使用此近似需要更强条件
$\varepsilon=\gamma^3aL/c^2\ll1$，并要求加速度在跨杆的本征时间尺度 $L/c$ 内缓变。
由 $\ell\simeq L/\gamma$ 和
$\dot\gamma=\gamma^3va/c^2$，
\[
 \boxed{\frac{\dd\ell}{\dd t}
 \simeq-\frac{L\gamma va}{c^2}},\qquad
 -\frac1v\frac{\dd\ell}{\dd t}
 \simeq\gamma\frac{aL}{c^2}\quad(v\ne0).
\]
所以在这个局部近似中，$\lambda=1,\mu=1,\nu=1$。这里必须区分本征长度 $L$ 与实验室同时长度 $\ell$。
这不是有限长加速杆的精确洛伦兹收缩律。例如尾端恒定本征加速度为 $a_0$，
在尾端瞬时静止的时刻取惯性系的 $t=0$，令 $R=c^2/a_0$，则该惯性系同时杆长精确为
\[
 \ell(t)=\sqrt{(R+L)^2+c^2t^2}-\sqrt{R^2+c^2t^2}.
\]
对 $L/R\ll1$ 展开才得到 $L/\gamma$；二阶项为
$a_0L^2\beta^2/(2c^2\gamma)$。单凭 $aL\ll c^2$，不能在任意大的 $\gamma$ 下丢掉它。

本征加速度不能只凭量纲猜出。取与杆在考察瞬间共速的惯性系，其速度参数为 $v$，
并在这个固定的惯性系中比较相邻两个事件。由速度变换求微分，在粒子瞬时速度 $u=v$ 处有
\[
 \left.\frac{\partial u'}{\partial u}\right|_{u=v}
 =\frac{1-v^2/c^2}{(1-uv/c^2)^2}\bigg|_{u=v}=\gamma^2,
 \qquad \dd t'=\gamma(1-vu/c^2)\dd t=\dd t/\gamma.
\]
于是 $a_0=\dd u'/\dd t'=\gamma^3\dd u/\dd t$，即
\[
 \boxed{a_0=\gamma^3a}.
\]
又
\[
 F=\frac{\dd(\gamma mv)}{\dd t}
 =m\gamma^3a=ma_0,
\]
故纵向三力在瞬时本征系与地面系数值相同：
$\boxed{F=F_0}$。这并不意味着任意方向三力都不变。

对第三问，在信号通过期间把尾端本征加速度 $a_0$ 视为常量。
Born 刚性表示各点在共动测量中的距离不变，不表示各点的本征加速度相同。
尾端 $x=0$、头端 $x=L$ 的加速度分别为 $a_0$ 和
$a_0/(1+a_0L/c^2)$。以尾端钟读数 $T$ 作时间坐标，在这个模型中
\[
 \dd s^2=
 \left(1+\frac{a_0x}{c^2}\right)^2c^2\dd T^2-\dd x^2.
\]
这个式子可由惯性坐标
$ct=(c^2/a_0+x)\sinh(a_0T/c)$、
$X=(c^2/a_0+x)\cosh(a_0T/c)-c^2/a_0$ 微分后相减得到。
固定 $x$ 的钟因此记录 $\dd\tau_x=(1+a_0x/c^2)\dd T$。
光满足 $\dd s^2=0$，从尾到头的坐标传播时间为
\[
 T_{\rm rec}-T_{\rm em}=\int_0^L\frac{\dd x}{c(1+a_0x/c^2)}
 =\frac c{a_0}\ln(1+a_0L/c^2),
\]
它不依赖发射时刻。相邻两个波峰在尾端相隔 $\Delta\tau_1=\Delta T$，
在头端相隔 $\Delta\tau_2=(1+a_0L/c^2)\Delta T$。
频率是波峰间隔的倒数，所以
\[
 \boxed{\frac{\nu_2}{\nu_1}
 =\frac1{1+a_0L/c^2}
 \simeq1-\frac{a_0L}{c^2}}.
\]
头端收到红移。上式的最后一步近似还要求
$a_0L/c^2=\gamma^3aL/c^2\ll1$；高速时，题设 $aL/c^2\ll1$ 本身并不足以保证它。
频移比的未展开式适用于上述恒定本征加速度模型；加速度随时间变化时，
还需其在一次传播时间内的变化足够小，才能局部使用。
第一问使用的 $\ell=L/\gamma$ 同样是各点速度差可忽略的局部短杆表达。
若杆各点被强行赋予相同实验室加速度，杆将不保持 Born 刚性，
题目便需要另行给出内部变形模型。

图~\ref{advA6:quality:A6-relativity:a1} 中，地面长度由同一地面时刻截取两端世界线得到，信号传播却连接不同时刻的两个事件；Born刚性还要求两端本征加速度按位置改变，不能给整根杆同一个本征加速度。
\end{advsolution}

\begin{advexamplebox}{相对论电子束的电场加速与磁场回转}\label{advA6:ex:relativistic-accelerator}
电子从下方入口进入两平行极板间，在匀强电场 $E$ 中沿长度 $L$ 加速，
从 $X$ 点进入匀强磁场并沿半圆到 $Z$ 点，再在上方极板间减速离开。
上下通道间距为 $D$，稳恒束流电流的大小为 $I>0$，电子电荷为 $-e$、质量为 $m$。
若入口速度不可忽略，以入口动量大小 $p_0$ 为参数作答；初动能可忽略的结果取 $p_0\to0$ 极限。
取向右为 $+x$、向上为 $+y$、出纸面为 $+z$。

\begin{enumerate}
\item 求磁感应强度以及下电场段、磁场段、上电场段中的总束流电荷量；
\item 换到使 $X$ 点电子静止的惯性系，求 $Y,Z$ 点速度，再求三个区域中的总电荷量；
\item 不直接套用场变换公式，利用测试电荷受力确定新系中的匀强场。
\end{enumerate}

\begin{center}
% Contract: E=-E ex; B=+B ez; q=-e. R=D/2, X=(0,-R),Y=(R,0),Z=(0,R).
% Path is tangent to horizontal channels. Plate apertures remain visibly open.
\begin{tikzpicture}[>=Latex,font=\small]
 \foreach \x in {-3,0}{
  \draw[very thick] (\x,-2.05)--(\x,-1.6) (\x,-1.4)--(\x,1.4) (\x,1.6)--(\x,2.05);
 }
 \draw[DeepBlue,thick,->] (-4.1,-1.5)--(-.45,-1.5);
 \draw[DeepBlue,thick] (-.45,-1.5)--(0,-1.5) arc(-90:90:1.5)--(-.4,1.5);
 \draw[DeepBlue,thick,->] (-.4,1.5)--(-4.1,1.5);
 \draw[DeepBlue,thick,->] ({1.5*cos(-20)},{1.5*sin(-20)}) arc(-20:18:1.5);
 \node[below right=5pt] at (0,-1.5){X};\node[right=7pt] at (1.5,0){Y};
 \node[above right=5pt] at (0,1.5){Z};
 \foreach \y in {-.55,.55}{\draw[->] (-.5,\y)--(-2.5,\y);}
 \node at (-1.5,0){$\vect E$};
 \foreach \x/\y in {.45/-.65,.85/-.65,.45/0,.85/0,.45/.65,.85/.65}{
  \draw (\x,\y) circle(.075);\fill (\x,\y) circle(.018);
 }
 \node at (2.1,-.9){$\vect B$ 出纸面};
 \draw[<->] (-3.55,-1.5)--(-3.55,1.5) node[midway,left=4pt]{$D$};
 \draw[<->] (-3,-2.4)--(0,-2.4) node[midway,below=4pt]{$L$};
 \node[below=5pt] at (-4.1,-1.5){$p_0$};
\end{tikzpicture}

\figcap{题图：保留下加速段、磁场半圆与上减速段的原始TikZ几何。}\label{advA6:quality:A6-relativity:q2}

\par\smallskip
电子经下通道向右、沿半圆转回、经上通道向左。极板开孔处允许束流通过。
\end{center}
\end{advexamplebox}
\global\AdvHideHeadingtrue
\begin{advsolution}
\begin{center}\begin{minipage}{.96\linewidth}\centering
% EMAdvFigureHeading
\raggedright{\bfseries\color{DeepBlue}解答}\par\smallskip\centering

\begin{tikzpicture}[>=Latex,font=\small]
 \foreach \x in {-3,0}{
  \draw[very thick] (\x,-2.05)--(\x,-1.6) (\x,-1.4)--(\x,1.4) (\x,1.6)--(\x,2.05);
 }
 \draw[DeepBlue,thick,->] (-4.1,-1.5)--(-.45,-1.5);
 \draw[DeepBlue,thick] (-.45,-1.5)--(0,-1.5) arc(-90:90:1.5)--(-.4,1.5);
 \draw[DeepBlue,thick,->] (-.4,1.5)--(-4.1,1.5);
 \draw[DeepBlue,thick,->] ({1.5*cos(-20)},{1.5*sin(-20)}) arc(-20:18:1.5);
 \node[below right=5pt] at (0,-1.5){X};\node[right=7pt] at (1.5,0){Y};
 \node[above right=5pt] at (0,1.5){Z};
 \foreach \y in {-.55,.55}{\draw[->] (-.5,\y)--(-2.5,\y);}
 \node at (-1.5,0){$\vect E$};
 \foreach \x/\y in {.45/-.65,.85/-.65,.45/0,.85/0,.45/.65,.85/.65}{
  \draw (\x,\y) circle(.075);\fill (\x,\y) circle(.018);
 }
 \node at (2.1,-.9){$\vect B$ 出纸面};
 \draw[<->] (-3.55,-1.5)--(-3.55,1.5) node[midway,left=4pt]{$D$};
 \draw[<->] (-3,-2.4)--(0,-2.4) node[midway,below=4pt]{$L$};
 \node[below=5pt] at (-4.1,-1.5){$p_0$};
\end{tikzpicture}
\figcap{解图：保留下加速段、磁场半圆与上减速段的原始TikZ几何。}\label{advA6:quality:A6-relativity:a2}
\end{minipage}\end{center}

图中负电子在下通道向右加速，电场必须向左：$\vect E=-E\vect e_x$。
在X点速度向右而磁力向上，$-e\vect v\times\vect B$ 要指向半圆圆心，故 $\vect B=B\vect e_z$ 出纸面。
设入口 $\gamma_0=\sqrt{1+(p_0/mc)^2}$。电场做功给出
\[
 \gamma_X=\gamma_0+\frac{eEL}{mc^2},\quad
 p_X=mc\sqrt{\gamma_X^2-1},\quad
 v_X=c\sqrt{1-\gamma_X^{-2}}.
\]
半圆半径 $R=D/2=p_X/(eB)$，故
\[
 \boxed{B=\frac{2p_X}{eD}}.
\]
稳恒束流在某段的总电荷量等于电流乘以驻留时间。电场段中
$\dd p/\dd t=eE$，所以
\[
 Q_{\rm lower}=Q_{\rm upper}=-I t_E,\qquad
 t_E=\frac{p_X-p_0}{eE}.
\]
磁场段速率不变，半圆驻留时间为 $\pi D/(2v_X)$，故
\[
 \boxed{Q_B=-I t_B},\qquad t_B=\frac{\pi D}{2v_X}.
\]
负号不能省去；若只问电荷量大小，才对三式取绝对值。
上通道电场力仍向右，与向左速度相反，恰好减速至入口的动能。

取新系相对实验室沿 $+x$ 方向速度 $v_X$，则 $X$ 点静止。
速度变换
\[
 u_x'=\frac{u_x-v_X}{1-u_xv_X/c^2},\qquad
 u_y'=\frac{u_y}{\gamma_X(1-u_xv_X/c^2)}
\]
给出
\[
 \vect v_X'=0,\quad
 \vect v_Y'=(-v_X,+v_X/\gamma_X),\quad
 \vect v_Z'=\frac{-2v_X}{1+\beta_X^2}\,\vect e_x .
\]
这里 $Y$ 点实验室速度沿 $+y$，故第一式中的 $u_x=0$，
而 $y$ 分量的正号与图中的绕行方向一致。

现在真正计算换系后的三个电荷。这里“区域”指实验室装置指定的各段世界管，
不是预先选定的一组电子。新系的同时截面与实验室截面不同，选入的电子可以不同。
沿电子前进方向取路径微元 $\dd s>0$，局部速率为 $u$，线电荷密度为 $\lambda=-I/u$。
实验室稳恒流的四流变换给出 $\rho'=\gamma_X(\rho-v_Xj_x/c^2)$；
同一实验室固定体积在 $t'=\text{常量}$ 截面上有 $\dd V'=\dd V/\gamma_X$。
于是
\[
 Q'=\int\left(\rho-\frac{v_Xj_x}{c^2}\right)\dd V
 =-I\int\frac{\dd s}{u}+\frac{I v_X}{c^2}\int\dd x.
\]
上式用了 $j_x\dd V=\lambda u_x\dd s=-I\dd x$。
图中下段、半圆段、上段的端点横坐标差分别为 $L,0,-L$，故
\[
 \boxed{Q'_{\rm lower}=-I(t_E-v_XL/c^2)},\qquad
 \boxed{Q'_B=-It_B},
\]
\[
 \boxed{Q'_{\rm upper}=-I(t_E+v_XL/c^2)}.
\]
$t_E>L/c$ 保证下段仍带负电；前后两段的修正相消。这也表明不能把三个电荷统一乘一个 $\gamma_X$。

最后不用场变换公式，而用测试电荷的运动方程确定新场。
对任意测试粒子，四力变换与 $\dd t'/\dd t=\gamma_X(1-v_Xu_x/c^2)$ 给出
\[
 F'_x=\frac{F_x-v_X(\vect F\cdot\vect u)/c^2}{1-v_Xu_x/c^2},
 \qquad F'_y=\frac{F_y}{\gamma_X(1-v_Xu_x/c^2)}.
\]
电场区域内 $\vect F=q(-E,0,0)$：令测试粒子在新系静止，即 $\vect u=v_X\vect e_x$，
得到 $\vect E'=-E\vect e_x$。再取任意非零 $u_y$，上述力变换始终给 $F'_y=0$，
且 $F'_x=-qE$。若测试粒子还具有任意 $u_z$，同理 $F'_z=0$。
三方向的力对测试粒子速度均不依赖，因此不能存在任何非零磁场分量，$\vect B'=0$。
在纯磁场区域，$\vect F=qB(u_y,-u_x,0)$ 且 $\vect F\cdot\vect u=0$。
先令 $\vect u=v_X\vect e_x$，新系粒子静止，因而
\[
 E'_y=F'_y/q=-\gamma_Xv_XB.
\]
再令测试粒子在实验室静止，原力为零；它在新系速度为 $-v_X\vect e_x$，
于是 $0=E'_y+v_XB'_z$，给出 $B'_z=\gamma_XB$。
从而三个区域的场为
\[
 \boxed{\text{两电场段： }\vect E'=-E\vect e_x,\quad\vect B'=0},
\]
\[
 \boxed{\text{磁场段： }\vect E'=-\gamma_Xv_XB\vect e_y,
 \quad\vect B'=\gamma_XB\vect e_z}.
\]
原来无场的外部区域仍无场。这里仅讨论理想外加场，忽略束流自场及开孔边缘场。
在X处 $-e\vect E'$ 向上，与电子的弯转方向一致；$v_X\to0$ 时也退回原场。

图~\ref{advA6:quality:A6-relativity:a2} 中，$X,Y,Z$ 是不同空间位置的束流截面，换系后应在新系同一时刻重新截取每个区域的电荷；不能把实验室同时装入区域的全部粒子不加说明地当作新系同时的同一集合。
\end{advsolution}

\begin{advexamplebox}{运动椭圆柱的表面四流与内部电磁场}\label{advA6:ex:elliptic-cylinder}
在 $S_0$ 系中有无限长椭圆柱，截面半轴为 $a>b$。表面叠放
$+\sigma_0$ 与 $-\sigma_0$ 两层电荷，负层静止，正层以速率 $u$ 沿周界流动。
选取相对 $S_0$ 沿 $-x$ 方向运动的 $S$ 系，使椭圆在 $S$ 中成为向右运动的圆柱。

\begin{enumerate}
\item 求内部磁场与相对速度 $v$；
\item 求 $S$ 系表面电荷密度和面电流密度随极角 $\theta$ 的关系；
\item 不使用场变换公式求圆柱内外的电磁场；已知 $S$ 中柱内仍为匀强场、柱外无场。
\end{enumerate}

\begin{center}
% Contract: ellipse (a,b)=(2.0,1.2); boost V=-v ex, gamma=a/b.
% Circle is x-vt=b cos(theta),y=b sin(theta); arrows tangent, CCW.
\begin{tikzpicture}[>=Latex,font=\small]
 \begin{scope}
  \draw[thick] (0,0) ellipse(2 and 1.2);
  \draw[->] (-2.4,0)--(2.4,0) node[right]{$x_0$};
  \draw[->] (0,-1.5)--(0,1.75) node[above]{$y_0$};
  \node[below=4pt] at (1,0){$a$};\node[left=4pt] at (0,.7){$b$};
  \draw[DeepBlue,thick,->] (2,-.35)--(2,.35) node[right=4pt]{$u$};
  \draw[DeepBlue,thick,->] (.45,1.2)--(-.45,1.2) node[above=5pt]{$u$};
  \node at (1.1,.5){$\pm\sigma_0$};
  \node at (0,-2){$S_0$：截面静止};
 \end{scope}
 \begin{scope}[xshift=7cm]
  \draw[thick] (0,0) circle(1.2);
  \draw[->] (-1.6,0)--(2,0) node[right]{$x-vt$};
  \draw[->] (0,-1.5)--(0,1.75) node[above]{$y$};
  \coordinate (P) at (45:1.2);
  \draw (0,0)--(P);\node[above left=3pt] at (45:.8){$b$};
  \draw (.48,0) arc(0:45:.48);\node at (23:.75){$\theta$};
  \fill (P) circle(1.8pt);
  \draw[thin] (P)--(1.55,1.2);\node[right=3pt] at (1.55,1.2){$\sigma,\vect K$};
  \draw[DeepBlue,thick,->] (1.4,-.8)--(2.2,-.8) node[midway,below=3pt]{$v$};
  \node at (0,-2){$S$：圆截面右移};
 \end{scope}
 \draw[->] (4.3,.55)--(3.1,.55) node[midway,above=4pt]{$\vect V=-v\vect e_x$};
\end{tikzpicture}

\figcap{题图：保留椭圆柱的正负面电荷层、切向正电流与运动方向。}\label{advA6:quality:A6-relativity:q3}

\par\smallskip
左图正电荷逆时针流动、负电荷静止。右图的角度从圆心随动横轴量起，$\vect K$ 为实验室面电流。
\end{center}
\end{advexamplebox}
\global\AdvHideHeadingtrue
\begin{advsolution}
\begin{center}\begin{minipage}{.96\linewidth}\centering
% EMAdvFigureHeading
\raggedright{\bfseries\color{DeepBlue}解答}\par\smallskip\centering

\begin{tikzpicture}[>=Latex,font=\small]
 \begin{scope}
  \draw[thick] (0,0) ellipse(2 and 1.2);
  \draw[->] (-2.4,0)--(2.4,0) node[right]{$x_0$};
  \draw[->] (0,-1.5)--(0,1.75) node[above]{$y_0$};
  \node[below=4pt] at (1,0){$a$};\node[left=4pt] at (0,.7){$b$};
  \draw[DeepBlue,thick,->] (2,-.35)--(2,.35) node[right=4pt]{$u$};
  \draw[DeepBlue,thick,->] (.45,1.2)--(-.45,1.2) node[above=5pt]{$u$};
  \node at (1.1,.5){$\pm\sigma_0$};
  \node at (0,-2){$S_0$：截面静止};
 \end{scope}
 \begin{scope}[xshift=7cm]
  \draw[thick] (0,0) circle(1.2);
  \draw[->] (-1.6,0)--(2,0) node[right]{$x-vt$};
  \draw[->] (0,-1.5)--(0,1.75) node[above]{$y$};
  \coordinate (P) at (45:1.2);
  \draw (0,0)--(P);\node[above left=3pt] at (45:.8){$b$};
  \draw (.48,0) arc(0:45:.48);\node at (23:.75){$\theta$};
  \fill (P) circle(1.8pt);
  \draw[thin] (P)--(1.55,1.2);\node[right=3pt] at (1.55,1.2){$\sigma,\vect K$};
  \draw[DeepBlue,thick,->] (1.4,-.8)--(2.2,-.8) node[midway,below=3pt]{$v$};
  \node at (0,-2){$S$：圆截面右移};
 \end{scope}
 \draw[->] (4.3,.55)--(3.1,.55) node[midway,above=4pt]{$\vect V=-v\vect e_x$};
\end{tikzpicture}
\figcap{解图：保留椭圆柱的正负面电荷层、切向正电流与运动方向。}\label{advA6:quality:A6-relativity:a3}
\end{minipage}\end{center}

均匀磁化 $\vect M=\sigma_0u\,\vect e_z$ 的束缚面电流
$\vect K=\vect M\times\vect n$ 恰为大小 $\sigma_0u$ 的周向电流，
与截面形状无关。因此
\[
 \boxed{\vect B_0=\mu_0\sigma_0u\,\vect e_z},\qquad
 \vect B_{\rm out}=0.
\]
沿 $x$ 的 Lorentz 收缩把 $a$ 变为 $a/\gamma$。要求它等于 $b$：
\[
 \boxed{\gamma=\frac ab},\qquad
 \boxed{v=c\sqrt{1-\frac{b^2}{a^2}}}.
\]

令圆截面点为 $(b\cos\theta,b\sin\theta)$，它在 $S_0$ 同时对应
$(a\cos\theta,b\sin\theta)$。记
\[
 s_0(\theta)=\sqrt{a^2\sin^2\theta+b^2\cos^2\theta}.
\]
$S_0$ 中面电流切向分量为
\[
 \vect K_0=\sigma_0u
 \frac{(-a\sin\theta,b\cos\theta)}{s_0}.
\]
下面写出面元因子，避免把面密度误当体密度直接变换。
定义同一个柱面方程
\[
 f=\frac{x_0^2}{a^2}+\frac{y_0^2}{b^2}-1
 =\frac{(x-vt)^2+y^2}{b^2}-1.
\]
表面电流的体分布是 $\vect j_0=\vect K_0|\nabla_0 f|\delta(f)$。
因为 $\int|\nabla f|\delta(f)\dd n=1$，这个写法沿法向积分恰好返回面电流。
在图中对应角度 $\theta$ 处，
\[
 |\nabla_0 f|=\frac{2s_0}{ab},\qquad |\nabla f|=\frac2b,
 \qquad \frac{|\nabla_0f|}{|\nabla f|}=\frac{s_0}a.
\]
换系速度为 $-v\vect e_x$，而初始净体电荷密度为零，所以四流变换为
$\rho=\gamma vj_{0x}/c^2$、$j_x=\gamma j_{0x}$、$j_y=j_{0y}$。
把各体分布除以新系的 $|\nabla f|\delta(f)$，逐项得到
\[
 \sigma=\frac{\gamma v}{c^2}K_{0x}\frac{s_0}a,
 \quad K_x=\gamma K_{0x}\frac{s_0}a,
 \quad K_y=K_{0y}\frac{s_0}a.
\]
代入 $K_{0x}=-\sigma_0ua\sin\theta/s_0$、
$K_{0y}=\sigma_0ub\cos\theta/s_0$ 后才有
\[
 \boxed{\sigma(\theta)=
 -\frac{\sigma_0uva}{bc^2}\sin\theta},
\]
\[
 \boxed{\vect K(\theta)=
 \sigma_0u\left(
 -\frac ab\sin\theta\,\vect e_x
 +\frac ba\cos\theta\,\vect e_y\right)}.
\]
这里的 $\vect K$ 包含随整个柱面平移的对流电流。
圆面的单位法向为 $\vect n=(\cos\theta,\sin\theta,0)$，法向速度为 $v_n=v\cos\theta$。
直接代入可核验 $\vect K\cdot\vect n=\sigma v_n$；
相对于移动表面的电流 $\vect K-\sigma v\vect e_x$ 才纯粹沿圆周切向。
又 $\int_0^{2\pi}\sigma b\dd\theta=0$，每单位轴向长度的净电荷仍为零。

第三问不用场变换。记 $[\vect E]=\vect E_{\rm out}-\vect E_{\rm in}$，磁场同理。
高斯定律在很薄的跨界柱体上积分，仍给出
$\vect n\cdot[\vect E]=\sigma/\varepsilon_0$。
但安培--麦克斯韦定律还含移动边界扫过测量点产生的位移电流。
把电场写成内外两部分的阶跃函数，时间导数的表面项为
$-v_n[\vect E]\delta(n)$，旋度的表面项为 $\vect n\times[\vect B]\delta(n)$，
故正确跳跃条件是
\[
 \vect n\times[\vect B]=\mu_0\vect K-\frac{v_n}{c^2}[\vect E].
\]
已知外场为零、内场均匀，由对任意 $\theta$ 都成立的高斯式
\[
 -E_x\cos\theta-E_y\sin\theta
 =-\gamma v\mu_0\sigma_0u\sin\theta
\]
先读出 $E_x=0$、$E_y=\gamma vB_0$，且轴向对称性给 $E_z=0$。
再看安培式的两个分量：
\[
 -B_z\sin\theta=\mu_0K_x=-\gamma B_0\sin\theta,
\]
\[
 B_z\cos\theta=\mu_0K_y+\frac{v\cos\theta}{c^2}E_y
 =B_0\cos\theta\left(\frac1\gamma+\gamma\frac{v^2}{c^2}\right)
 =\gamma B_0\cos\theta.
\]
两式一致给出
\[
 \boxed{\vect E_{\rm in}=\gamma vB_0\,\vect e_y},\qquad
 \boxed{\vect B_{\rm in}=\gamma B_0\,\vect e_z},
\qquad
 \vect E_{\rm out}=\vect B_{\rm out}=0.
\]
若漏掉位移电流项，第二个分量就会错误地给出 $B_z=B_0/\gamma$，与第一个分量矛盾。
另外两条跳跃条件也应检查：磁场沿轴向，所以 $\vect n\cdot[\vect B]=0$；
法拉第定律在移动边界上给 $\vect n\times[\vect E]=v_n[\vect B]$，
其 $z$ 分量为 $-E_y\cos\theta=-vB_z\cos\theta$，由 $E_y=vB_z$ 满足。
最后检查 $a\to b$ 时 $v\to0$，电场消失、磁场回到 $B_0$；
电场尺度 $vB_0$ 的单位也正是 $\mathrm{V/m}$。

图~\ref{advA6:quality:A6-relativity:a3} 中，柱面法向随极角改变，所以面密度变换除四流外还须包含面积元素的变化；圆截面是在新系同一时刻取得，不能把椭圆照片与同时截面混用。
\end{advsolution}

\begingroup\tcbset{breakable/.default=false}
\begin{advexamplebox}{GRAND 中微子反应阈能与定角电子最高能量}\label{advA6:ex:grand-neutrino}
超高能 $\tau$ 中微子可在岩石中发生
\[
 \nu_\tau+n\longrightarrow\tau^-+p.
\]
取 $m_\tau=1777\,\mathrm{MeV}/c^2$、
$m_p=938.2\,\mathrm{MeV}/c^2$、
$m_n=939.6\,\mathrm{MeV}/c^2$。

\begin{enumerate}
\item 中子初始静止，求入射中微子阈能；
\item 产生的 $\tau^-$ 经
$\tau^-\to\nu_\tau+e^-+\bar\nu_e$ 衰变，
$m_e=0.5110\,\mathrm{MeV}/c^2$。在与入射束成 $30^\circ$ 的方向探测电子，
求其最大能量。先把入射能量 $E_\nu\ge E_{\nu,\rm th}$ 保留为参数，
再讨论未限制入射能量时的上界，以及阈值处的特例。
\end{enumerate}

\begin{center}
\includegraphics[width=.84\textwidth]{assets/generated/A6-grand-radio-array.png}
\par\smallskip
{\small\color{TextGray}山地射电天线阵列通过短暂大气簇射的低频脉冲寻找地球掠射中微子事件。}
\end{center}

\begin{center}\begin{minipage}{.96\linewidth}\centering
\begin{tikzpicture}[>=Latex,font=\small,scale=.85,every node/.style={fill=none}]

\draw[DeepBlue,->] (-2,0)--(0,0) node[midway,above]{$\nu_\tau,E_\nu$};\fill (0,0) circle (2pt) node[below]{$n$ 静止};\draw[->] (.4,0)--(1.5,0);\node[above] at (1.7,0){$\tau^-+p$};
\fill (2.4,0) circle (1.5pt);\draw[dashed] (2.4,0)--(4.6,0);\draw[DeepBlue,->] (2.4,0)--(4.3,1.097) node[right]{$e^-$};\node at (3.1,.18){$30^\circ$};\node[below] at (3.8,-.5){衰变另有两中微子};

\end{tikzpicture}
\figcap{题图：固定中子靶的产生反应与随后定角电子探测，两个过程分别使用能动量守恒。}\label{advA6:quality:A6-relativity:q4}
\end{minipage}\end{center}
\end{advexamplebox}
\endgroup\par\nopagebreak[3]
\global\AdvHideHeadingtrue
\begin{advsolution}
\begin{center}\begin{minipage}{.96\linewidth}\centering
% EMAdvFigureHeading
\raggedright{\bfseries\color{DeepBlue}解答}\par\smallskip\centering

\begin{tikzpicture}[>=Latex,font=\small,scale=.85,every node/.style={fill=none}]

\draw[DeepBlue,->] (-2,0)--(0,0) node[midway,above]{$\nu_\tau,E_\nu$};\fill (0,0) circle (2pt) node[below]{$n$ 静止};\draw[->] (.4,0)--(1.5,0);\node[above] at (1.7,0){$\tau^-+p$};
\fill (2.4,0) circle (1.5pt);\draw[dashed] (2.4,0)--(4.6,0);\draw[DeepBlue,->] (2.4,0)--(4.3,1.097) node[right]{$e^-$};\node at (3.1,.18){$30^\circ$};\node[below] at (3.8,-.5){衰变另有两中微子};

\end{tikzpicture}
\figcap{解图：固定中子靶的产生反应与随后定角电子探测，两个过程分别使用能动量守恒。}\label{advA6:quality:A6-relativity:a4}
\end{minipage}\end{center}

用 Mandelstam 不变量
\[
 s=(p_\nu+p_n)^2=m_n^2c^4+2m_nc^2E_\nu.
\]
阈值时末态在质心系无相对运动，故
$s_{\rm th}=(m_\tau+m_p)^2c^4$，于是
\[
 \boxed{E_{\nu,\rm th}=
 \frac{(m_\tau+m_p)^2-m_n^2}{2m_n}c^2
 =3.453\,\mathrm{GeV}}.
\]

第二问不能把第一问的阈值自动当作入射能量。下面统一以能量单位表示静质量，
令 $N=m_nc^2$、$P=m_pc^2$、$M=m_\tau c^2$、$\mu=m_ec^2$，并记
\[
 s=N^2+2NE_\nu,\qquad
 \Lambda=[s-(M+P)^2][s-(M-P)^2].
\]
$\sqrt s$ 是整个反应的质心能量。两体产生过程的能量、动量守恒给出
\[
 E_\tau^*=\frac{s+M^2-P^2}{2\sqrt s},\qquad
 cp_\tau^*=\frac{\sqrt\Lambda}{2\sqrt s}.
\]
星号指反应质心系，而下文小标 $0$ 指 $\tau$ 的静止系，二者不能混同。
把两个末态中微子合成四动量 $q$，其不变质量平方非负。由
$(p_\tau-p_e)^2=q^2\ge0$ 可知，在 $\tau$ 静止系中
\[
 E_{e0}\le\frac{M^2+\mu^2}{2M},\qquad
 cp_{e0,\max}=\frac{M^2-\mu^2}{2M}.
\]
两中微子同向时达到能谱端点。把端点电子沿 $\tau$ 在质心系中的运动方向发射，
便得到电子在反应质心系中的最高能量
\[
 H=E_{e,\max}^*
 =\frac{E_\tau^*(M^2+\mu^2)+cp_\tau^*(M^2-\mu^2)}{2M^2}.
\]
在仅讨论运动学允许性的意义下，产生的 $\tau$ 方向可选，故这一端点可沿任意质心方向取到；
实际反应概率分布不改变运动学边界。

实验室电子的方向固定为 $\theta=30^\circ$。总四动量与电子四动量的内积给出
\[
 E_e^*=\frac{(E_\nu+N)E_e-E_\nu\cos\theta\sqrt{E_e^2-\mu^2}}{\sqrt s}.
\]
令 $A=E_\nu+N$、$B=E_\nu\cos\theta$、$Y=\sqrt s H$。
边界上 $AE_e-B\sqrt{E_e^2-\mu^2}=Y$。
移项、平方解二次方程，并取满足原式的高能根：
\[
 \boxed{E_{e,\max}(E_\nu,\theta)=
 \frac{AY+B\sqrt{Y^2-(A^2-B^2)\mu^2}}{A^2-B^2}}.
\]
对本题 $\theta=30^\circ$ 和允许的入射能量，该根可达。忽略极小的电子质量修正，式子简化为
\[
 \boxed{E_{e,\max}\simeq
 \frac{s+M^2-P^2+\sqrt\Lambda}
 {4[N+E_\nu(1-\cos\theta)]}}.
\]
这才是保留一般入射能量的答案。若不设 $E_\nu$ 上限，其高能极限为
\[
 \sup E_e\simeq\frac{N}{1-\cos30^\circ}=7.013\,\mathrm{GeV}.
\]
它是可趋近的上确界，不是任一有限入射能量都能达到的电子能量。
在 $\mu=0$ 近似下，由 $\sqrt\Lambda<s-M^2-P^2$，前式严格小于
$(N^2-P^2+2NE_\nu)/[2N+2E_\nu(1-\cos\theta)]$，
后者又小于 $N/(1-\cos30^\circ)$；而 $E_\nu\to\infty$ 时两界合拢，故上界确实可趋近。

作为校验，单独把入射能量取到阈值。此时 $p_\tau^*=0$，质心整体因子和速度为
\[
 \gamma_c=\frac{E_{\nu,\rm th}+m_nc^2}
 {(m_\tau+m_p)c^2}=1.618,\qquad
 \beta_c=\frac{E_{\nu,\rm th}}
 {E_{\nu,\rm th}+m_nc^2}=0.7861.
\]
三体衰变中，当两个中微子平行且合成不变质量趋于零时，
电子在 $\tau$ 静止系的能量最大：
\[
 E_e^*=\frac{m_\tau^2+m_e^2}{2m_\tau}c^2
 =888.5\,\mathrm{MeV}.
\]
对给定实验室夹角 $\theta$，用不变量
\[
 p_\tau\cdot p_e
 =E_\tau E_e/c^2-P_\tau p_e\cos\theta
 =m_\tau E_e^*
\]
并取 $p_e=\sqrt{E_e^2/c^2-m_e^2c^2}$ 解正根。
由于 $m_e/E_e\ll1$，数值上等价于
\[
 E_{e,\max}\simeq
 \frac{E_e^*}{\gamma_c(1-\beta_c\cos30^\circ)}
 =\boxed{1.720\,\mathrm{GeV}}.
\]
使用电子质量求精确根只改变末位。结果小于入射能量且满足四动量守恒。

图~\ref{advA6:quality:A6-relativity:a4} 中，电子探测角从原入射束方向量起，未必等于相对 $\tau$ 运动方向的角；先保留入射能量参数并优化允许的 $\tau$ 运动，再讨论无界入射能量极限。
\end{advsolution}

\section{量子物理现象}

\subsection{真空光场的模式与黑体辐射}

在边长 $l_1,l_2,l_3$ 的周期盒中，真空四势满足波动方程和 Lorenz 规范
\[
 \Box A^\mu=0,\qquad \partial_\mu A^\mu=0,
\]
允许波矢
\[
 k_i=\frac{2\pi N_i}{l_i},\qquad
 \omega=ck.
\]
对每个非零 $\vect k$ 只有两种相互正交的横向偏振具有物理电磁场。
沿 $k^\mu$ 的规范模式虽满足形式方程，却给出 $\vect E=\vect B=0$，不应误称为可观测“纵光子”。

把磁矢势 Fourier 展开为
\[
 \vect A(\vect r,t)=
 \sum_{\vect k,r}q_{\vect k,r}(t)\,
 \vect e_{\vect k,r}\e^{\ii\vect k\cdot\vect r},
\qquad r=1,2.
\]
利用 Parseval 公式，电磁场 Hamilton 量成为独立谐振子之和
\[
 H=\sum_{\vect k,r}
 \left(\frac{|p_{\vect k,r}|^2}{2m}
 +\frac{m\omega_{\vect k}^2|q_{\vect k,r}|^2}{2}\right).
\]
半径为 $k$ 的波矢球计数并乘两种偏振，单位体积角频率态密度为
\[
 \boxed{g(\omega)=\frac{\omega^2}{\pi^2c^3}}.
\]

经典能量均分给每个模式平均能量 $k_{\rm B}T$，得到 Rayleigh--Jeans 谱
\[
 u_{\rm RJ}(\omega)=\frac{k_{\rm B}T\omega^2}{\pi^2c^3},
\]
其高频积分发散。Planck 量子化令谐振子能级为
$E_n=n\hbar\omega$（真空零点项不影响热谱），配分函数与平均热能为
\[
 Z_\omega=\frac1{1-\e^{-\beta\hbar\omega}},\qquad
 \bar E_\omega=\frac{\hbar\omega}{\e^{\beta\hbar\omega}-1}.
\]
因此
\[
 \boxed{u(\omega)=
 \frac{\hbar\omega^3}{\pi^2c^3}
 \frac1{\e^{\hbar\omega/(k_{\rm B}T)}-1}}.
\]
低频极限回到 Rayleigh--Jeans 公式，高频呈指数衰减。

用
\[
 \int_0^\infty\frac{x^{s-1}}{\e^x-1}\,\dd x
 =\Gamma(s)\zeta(s)
\]
且 $\zeta(4)=\pi^4/90$，得到
\[
 u=a_{\rm R}T^4,\qquad
 a_{\rm R}=\frac{\pi^2k_{\rm B}^4}{15\hbar^3c^3},
\]
\[
 j=\frac{cu}{4}=\sigma_{\rm SB}T^4,\qquad
 \sigma_{\rm SB}=
 \frac{\pi^2k_{\rm B}^4}{60\hbar^3c^2}.
\]
Bernoulli 数的母函数
$t/(\e^t-1)=\sum B_nt^n/n!$
还给出低频渐近展开。

\begin{advtheorem}{$\pi\cot\pi z$ 的 Mittag--Leffler 展开}\label{advA6:thm:mittag-leffler}
对 $z\notin\mathbb Z$，
\[
 \boxed{\pi\cot\pi z=
 \frac1z+\sum_{n=1}^\infty
 \left(\frac1{z-n}+\frac1{z+n}\right)}.
\]
从而
\[
 \zeta(2m)=
 (-1)^{m+1}\frac{(2\pi)^{2m}}{2(2m)!}B_{2m}.
\]
\end{advtheorem}

\noindent\textit{证明。}
右端按 $\pm n$ 成对求和后局部一致收敛，并与 $\pi\cot\pi z$ 具有相同的一阶极点和留数。
令两者之差为 $h(z)$，则 $h$ 是整函数且周期为 $1$。
在矩形
$[-N-\tfrac12,N+\tfrac12]\times[-N,N]$ 上，
$\cot\pi z$ 在远离整数极点处有统一界；成对级数的尾项为 $O(N^{-1})$。
由最大模原理，令 $N\to\infty$ 得 $h$ 为有界整函数，故为常数。
取 $z\to\ii\infty$ 或比较 $z=1/2$ 可知常数为零。

再写
\[
 \pi z\cot\pi z=1-2\sum_{s=1}^\infty\zeta(2s)z^{2s},
\]
与 Bernoulli 展开
\[
 \pi z\cot\pi z=
 \sum_{s=0}^\infty
 (-1)^s\frac{(2\pi)^{2s}B_{2s}}{(2s)!}z^{2s}
\]
比较系数即得第二式。$m=1$ 给出 Basel 结果
$\zeta(2)=\pi^2/6$。

\begin{advexamplebox}{充满热辐射的弹性气球}\label{advA6:ex:radiation-balloon}
质量为 $m$ 的球形薄膜包住温度为 $T$ 的热辐射，外界为真空且无泄漏。
薄膜弹性等效于表面张力系数为 $\alpha$ 的肥皂泡。

\begin{enumerate}
\item 求平衡半径；
\item 若薄膜作各向同性快速径向小振动，辐射与薄膜来不及交换热量，求周期；
\item 判断这种热平衡在与热库接触时是否稳定。
\end{enumerate}
\end{advexamplebox}
\begin{advsolution}
光子气体能量密度与压强为
\[
 u=a_{\rm R}T^4,\qquad p=\frac u3.
\]
肥皂泡有内外两表面，Laplace 压强为 $4\alpha/R$，故
\[
 \frac{a_{\rm R}T^4}{3}=\frac{4\alpha}{R},
\qquad
 \boxed{R_0=\frac{12\alpha}{a_{\rm R}T^4}}.
\]

快速绝热运动中光子气体满足 $TR=\text{常量}$，所以
$U=a_{\rm R}T^4(4\pi R^3/3)=C/R$。
两表面的表面能为 $8\pi\alpha R^2$，
\[
 V_{\rm eff}(R)=\frac CR+8\pi\alpha R^2.
\]
平衡条件给 $C=16\pi\alpha R_0^3$，二阶导数
\[
 V_{\rm eff}''(R_0)=\frac{2C}{R_0^3}+16\pi\alpha
 =48\pi\alpha.
\]
薄膜所有质点径向速度均为 $\dot R$，动能为 $m\dot R^2/2$，故
\[
 \boxed{\omega_R=\sqrt{\frac{48\pi\alpha}{m}}},
\qquad
 \boxed{\mathcal T=2\pi\sqrt{\frac{m}{48\pi\alpha}}}.
\]

在平衡族上 $R_0\propto T^{-4}$，总能量
$U+8\pi\alpha R_0^2\propto T^{-8}$，因此
$C_{\rm eff}=\dd E/\dd T<0$。孤立体系可在力学上稳定振动，
但与无限热库作正反馈换热时热平衡不稳定。这区分了力学稳定性和热稳定性。
\end{advsolution}

\subsection{介质响应、因果性与散射}

Lorentz 振子模型
\[
 m\ddot x+m\Gamma\dot x+m\omega_0^2x=-eE
\]
给出
\[
 \chi(\omega)=
 \frac{\omega_p^2}{\omega_0^2-\omega^2-\ii\Gamma\omega},
\qquad
 \omega_p^2=\frac{n_eZe^2}{m\varepsilon_0}.
\]
写 $\chi=\chi_1+\ii\chi_2$，实部控制色散，虚部控制吸收。

\begin{advtheorem}{Kramers--Kronig 关系}\label{advA6:thm:kramers-kronig}
若线性响应满足因果性，$\chi(t<0)=0$，且高频足够快衰减，则
\[
 \boxed{\chi_1(\omega)=\frac1\pi
 \operatorname P\!\int_{-\infty}^{\infty}
 \frac{\chi_2(\omega')}{\omega'-\omega}\,\dd\omega'},
\]
\[
 \boxed{\chi_2(\omega)=-\frac1\pi
 \operatorname P\!\int_{-\infty}^{\infty}
 \frac{\chi_1(\omega')}{\omega'-\omega}\,\dd\omega'}.
\]
\end{advtheorem}

\noindent\textit{证明。}
因果卷积
\[
 P(t)=\varepsilon_0\int_0^\infty\chi(\tau)E(t-\tau)\,\dd\tau
\]
给出
\[
 \chi(z)=\int_0^\infty\chi(\tau)\e^{\ii z\tau}\,\dd\tau .
\]
当 $\operatorname{Im}z>0$ 时指数衰减，故 $\chi(z)$ 在上半平面解析。
对
$\chi(z')/(z'-z)$ 沿实轴和上半平面大半圆应用 Cauchy 公式，
高频衰减使半圆积分消失：
\[
 \chi(z)=\frac1{2\pi\ii}
 \int_{-\infty}^{\infty}\frac{\chi(\omega')}{\omega'-z}\,\dd\omega'.
\]
令 $z=\omega+\ii0^+$，使用 Sokhotski--Plemelj 公式
\[
 \frac1{x-\ii0^+}=\operatorname P\frac1x+\ii\pi\delta(x),
\]
再分别取实部、虚部，得到两式。

\begin{advexamplebox}{Lorentz 振子的因果脉冲响应}\label{advA6:ex:lorentz-laplace}
对
\[
 \chi(\omega)=
 \frac{\omega_p^2}{\omega_0^2-\omega^2-\ii\Gamma\omega}
\]
用 Laplace 变换求时域响应 $\chi(t)$，并检验因果性。
\end{advexamplebox}
\begin{advsolution}
单位脉冲驱动的响应满足
\[
 \ddot\chi+\Gamma\dot\chi+\omega_0^2\chi
 =\omega_p^2\delta(t),\qquad \chi(t<0)=0.
\]
单边 Laplace 变换给
\[
 \tilde\chi(s)=\frac{\omega_p^2}
 {s^2+\Gamma s+\omega_0^2}.
\]
欠阻尼 $\Gamma<2\omega_0$ 时，令
$\Omega=\sqrt{\omega_0^2-\Gamma^2/4}$，逆变换得
\[
 \boxed{\chi(t)=
 \Theta(t)\frac{\omega_p^2}{\Omega}
 \e^{-\Gamma t/2}\sin\Omega t}.
\]
它在 $t<0$ 严格为零，并在 $t=0^+$ 满足
$\dot\chi(0^+)=\omega_p^2$。令 $s=-\ii\omega+0^+$ 便恢复给定的
$\chi(\omega)$，其极点均在频率下半平面，因果性检查通过。
\end{advsolution}

加速电荷的 Larmor 辐射功率为
\[
 P_{\rm rad}=\frac{q^2a^2}{6\pi\varepsilon_0c^3}.
\]
Lorentz--Abraham 辐射反作用的非相对论形式
\[
 \vect F_{\rm rad}=m\tau\dot{\vect a},\qquad
 \tau=\frac{q^2}{6\pi\varepsilon_0mc^3}
\]
只是在变化时间尺度远大于 $\tau$ 时的有效近似。
把它当作精确三阶运动方程会出现 runaway 解或预加速；正确做法是作降阶近似
\[
 m\vect a\simeq\vect F_{\rm ext}+\tau\dot{\vect F}_{\rm ext},
\]
并只保留 $O(\tau)$。对电子 $\tau$ 约为 $6.3\times10^{-24}\,\mathrm s$。

谐振电子的辐射散射截面可写为
\[
 \sigma(\omega)=\frac{8\pi r_e^2}{3}
 \frac1{(\omega_0^2/\omega^2-1)^2+\omega^2\tau^2},
\qquad
 r_e=\frac{e^2}{4\pi\varepsilon_0mc^2}.
\]
高于束缚频率且仍远低于 $1/\tau$ 时趋于 Thomson 截面
$8\pi r_e^2/3$；低频极限 $\sigma\propto\omega^4$，产生 Rayleigh 散射；
在 $\omega\simeq\omega_0$ 附近出现共振增强。经典模型能说明形状，
但离散原子谱和跃迁强度需要量子理论。

\subsection{Stern--Gerlach 实验与自旋态}

非均匀磁场对磁矩的 $z$ 分量施力
\[
 F_z=\mu_z\frac{\partial B_z}{\partial z}.
\]
碱金属原子束分成两束，说明电子自旋投影只有
$S_z=\pm\hbar/2$。电子磁矩
\[
 \vect\mu_s=-g_e\frac{e}{2m}\vect S,\qquad
 g_e\simeq2.002319\ldots
\]
的负号来自电子负电荷。

\begin{center}\begin{minipage}{\linewidth}\centering
\centering
\begin{tikzpicture}[>=Latex,scale=.70]
 \begin{scope}[yshift=2.1cm]
  \node[draw,minimum width=1.3cm,minimum height=.65cm] (O1) at (-5,0){炉};
  \node[draw,minimum width=1.3cm,minimum height=.65cm] (Z1) at (-2.8,0){SG$_z$};
  \node[draw,minimum width=1.3cm,minimum height=.65cm] (Z2) at (1.1,.55){SG$_z$};
  \draw[->] (O1)--(Z1);
  \draw[->] (Z1.east)--(Z2.west) node[midway,above]{$z+$};
  \draw[dashed] (Z1.east)--(1.1,-.55) node[midway,below]{$z-$ 被挡};
  \draw[->] (Z2.east)--(4.7,.55) node[right]{$z+$};
  \draw[dashed] (Z2.east)--(4.7,-.15) node[right]{无 $z-$};
  \node[left] at (-6.1,0){(a)};
 \end{scope}
 \begin{scope}
  \node[draw,minimum width=1.3cm,minimum height=.65cm] (O2) at (-5,0){炉};
  \node[draw,minimum width=1.3cm,minimum height=.65cm] (Z3) at (-2.8,0){SG$_z$};
  \node[draw,minimum width=1.3cm,minimum height=.65cm] (X1) at (1.1,.55){SG$_x$};
  \draw[->] (O2)--(Z3);
  \draw[->] (Z3.east)--(X1.west) node[midway,above]{$z+$};
  \draw[dashed] (Z3.east)--(1.1,-.55) node[midway,below]{$z-$ 被挡};
  \draw[->] (X1.east)--(4.7,.95) node[right]{$x+$};
  \draw[->] (X1.east)--(4.7,.15) node[right]{$x-$};
  \node[left] at (-6.1,0){(b)};
 \end{scope}
 \begin{scope}[yshift=-2.15cm]
  \node[draw,minimum width=1.3cm,minimum height=.65cm] (O3) at (-5,0){炉};
  \node[draw,minimum width=1.3cm,minimum height=.65cm] (Z4) at (-2.8,0){SG$_z$};
  \node[draw,minimum width=1.3cm,minimum height=.65cm] (X2) at (.2,.45){SG$_x$};
  \node[draw,minimum width=1.3cm,minimum height=.65cm] (Z5) at (3.2,.75){SG$_z$};
  \draw[->] (O3)--(Z4);
  \draw[->] (Z4.east)--(X2.west) node[midway,above]{$z+$};
  \draw[->] (X2.east)--(Z5.west) node[midway,above]{$x+$};
  \draw[->] (Z5.east)--(5.3,1.05) node[right]{$z+$};
  \draw[->] (Z5.east)--(5.3,.45) node[right]{$z-$};
  \node[left] at (-6.1,0){(c)};
 \end{scope}
\end{tikzpicture}
\bookcaption{序列 Stern--Gerlach 实验：选择一个方向的本征态会破坏另一方向的确定性}
\label{advA6:fig:sequential-sg}
\end{minipage}\end{center}

图 \ref{advA6:fig:sequential-sg}(a) 表明重复测量同一分量结果确定；
(b) 表明 $|z+\rangle$ 对 $S_x$ 是等概率叠加；
(c) 表明选择 $x+$ 后再次测量 $S_z$ 又分成两束。
因此态振幅必须是复数：
\[
 |\vect n,+\rangle=
 \cos\frac\theta2|z+\rangle
 +\e^{\ii\phi}\sin\frac\theta2|z-\rangle,
\]
概率由模方给出。相位 $\phi$ 不影响单次 $S_z$ 概率，却影响其他方向的干涉结果，
这正是实概率模型无法表达的结构。

% AdvancedEMQuality20261023
\section{量子力学初步}

\subsection{Schrödinger 方程与定态问题}

把经典 Hamilton 量中的
$\vect p$ 换成算符 $-\ii\hbar\nabla$，得到
\[
 \ii\hbar\frac{\partial\psi}{\partial t}
 =\widehat H(\vect r,-\ii\hbar\nabla)\psi.
\]
若势能不显含时间，可分离为
$\psi(\vect r,t)=\Psi(\vect r)\e^{-\ii Et/\hbar}$，空间波函数满足
\[
 \boxed{\widehat H\Psi=E\Psi}.
\]
自由粒子
$\widehat H=-\hbar^2\nabla^2/(2m)$
的平面波满足 $E=\hbar\omega$、$\vect p=\hbar\vect k$、
$E=p^2/(2m)$。

\begin{advexamplebox}{球形禁闭、动能下界与氢原子稳定性}\label{advA6:ex:hydrogen-confinement}
把电子波函数完全限制在半径为 $R$ 的球内并令边界上 $\psi=0$。

\begin{enumerate}
\item 证明动能存在下界，取下界的波函数满足
$\widehat T\psi=T_{\min}\psi$；
\item 求 $T_{\min}$ 与基态波函数；
\item 说明 Coulomb 势下电子不能无限塌缩到原点；
\item 用该球内基态作变分试探函数，估算氢原子基态能量。
\end{enumerate}

\begin{center}\begin{minipage}{.96\linewidth}\centering
\begin{tikzpicture}[>=Latex,font=\small,scale=.85,every node/.style={fill=none}]

\draw[thick] (0,0) circle (1.3);\fill (0,0) circle (1.5pt) node[below]{$O$};\draw[->] (0,0)--(1.3,0) node[midway,below]{$R$};\node[above] at (0,1.6){$\psi(R)=0$};\node[right] at (1.7,.5){球内允许，球外 $\psi=0$};

\end{tikzpicture}
\figcap{题图：球形禁闭区域与Dirichlet边界；Coulomb问题另以此态作变分试探。}\label{advA6:quality:A6-quantum-mechanics:q1}
\end{minipage}\end{center}
\end{advexamplebox}
\global\AdvHideHeadingtrue
\begin{advsolution}
\begin{center}\begin{minipage}{.96\linewidth}\centering
% EMAdvFigureHeading
\raggedright{\bfseries\color{DeepBlue}解答}\par\smallskip\centering

\begin{tikzpicture}[>=Latex,font=\small,scale=.85,every node/.style={fill=none}]

\draw[thick] (0,0) circle (1.3);\fill (0,0) circle (1.5pt) node[below]{$O$};\draw[->] (0,0)--(1.3,0) node[midway,below]{$R$};\node[above] at (0,1.6){$\psi(R)=0$};\node[right] at (1.7,.5){球内允许，球外 $\psi=0$};

\end{tikzpicture}
\figcap{解图：球形禁闭区域与Dirichlet边界；Coulomb问题另以此态作变分试探。}\label{advA6:quality:A6-quantum-mechanics:a1}
\end{minipage}\end{center}

归一化约束下令
\[
 \mathcal F[\psi]=
 \int_{r<R}\psi^*
 \left(-\frac{\hbar^2}{2m}\nabla^2\right)\psi\,\dd V
 -\lambda\left(\int|\psi|^2\dd V-1\right).
\]
对 $\psi^*$ 变分得
$\widehat T\psi=\lambda\psi$。因此最小值由 Dirichlet Laplacian 的最低本征值给出。

球对称基态令 $u=r\psi$，满足
$u''+k^2u=0$、$u(0)=u(R)=0$。故
\[
 k=\frac\pi R,\qquad
 \boxed{T_{\min}=\frac{\pi^2\hbar^2}{2mR^2}},
\]
\[
 \boxed{\psi_R(r)=
 \frac1{\sqrt{2\pi R}}\frac{\sin(\pi r/R)}r
 \quad(0<r<R)}.
\]
尺度压缩 $R\to0$ 时动能按 $R^{-2}$ 发散，而 Coulomb 势能只按
$-R^{-1}$ 发散，所以总能量必在有限尺度取最小值。

上述是尺度解释，不能只凭一族试探函数就断言所有态都有下界。对任意归一化、有限动能且使分部积分面项消失的波函数，令 $a_0=\hbar^2/(mk_ee^2)$，利用 $\nabla\cdot\hat{\boldsymbol r}=2/r$ 可直接配方：
\[
 \langle H\rangle=
 \frac{\hbar^2}{2m}\int\left|\nabla\psi+
 \frac{\hat{\boldsymbol r}}{a_0}\psi\right|^2\mathrm dV
 -\frac{\hbar^2}{2ma_0^2}
 \ge-\frac{mk_e^2e^4}{2\hbar^2}.
\]
这里 $k_e=1/(4\pi\varepsilon_0)$；平方项展开后的交叉项恰是 $-k_ee^2\langle1/r\rangle$。这个对所有允许态成立的界排除了无限塌缩。等号态的径向衰减为 $e^{-r/a_0}$，也说明下面的硬球试探态只能给出较粗的变分上界。


记 $k_e=1/(4\pi\varepsilon_0)$。对上述试探态
\[
 \left\langle\frac1r\right\rangle
 =\frac{C}{R},\qquad
 C=2\int_0^\pi\frac{\sin^2u}{u}\,\dd u
 =2.43765\ldots
\]
故
\[
 E(R)=\frac{\pi^2\hbar^2}{2mR^2}
 -\frac{Ck_ee^2}{R}.
\]
极小点与能量为
\[
 R_*=\frac{\pi^2}{C}a_0=4.049\,a_0,
\]
\[
 \boxed{E_{\rm var}=
 -\frac{C^2}{\pi^2}\frac{mk_e^2e^4}{2\hbar^2}
 =-8.19\,\mathrm{eV}}.
\]
它高于精确基态 $-13.6\,\mathrm{eV}$，符合变分上界原理；估算较粗是因硬球边界并非真实指数尾部。

图~\ref{advA6:quality:A6-quantum-mechanics:a1} 中，半径 $R$ 同时决定边界和径向驻波波数；球内基态用 $u=r\psi$ 消去球坐标的一阶导数，变分能量还须连同势能一起计算。
\end{advsolution}

\subsection{三个一维模型}

\begin{advexamplebox}{无限深方势阱}\label{advA6:ex:infinite-well}
粒子被限制在 $0<x<a$，阱外势能无穷大。证明能级和归一化波函数为
\[
 E_n=\frac{n^2\pi^2\hbar^2}{2ma^2},\qquad
 \Psi_n(x)=\sqrt{\frac2a}\sin\frac{n\pi x}{a}.
\]

\begin{center}\begin{minipage}{.96\linewidth}\centering
\begin{tikzpicture}[>=Latex,font=\small,scale=.85,every node/.style={fill=none}]

\draw[->] (-.7,0)--(4.2,0) node[right]{$x$};\draw[very thick] (0,2)--(0,0)--(3,0)--(3,2);\node[below] at (0,0){0};\node[below] at (3,0){$a$};\node at (1.5,1){$V=0$};\node[left] at (0,2){$\infty$};\node[right] at (3,2){$\infty$};

\end{tikzpicture}
\figcap{题图：无限深阱的两个硬壁位于 $0,a$，波函数在两壁为零。}\label{advA6:quality:A6-quantum-mechanics:q2}
\end{minipage}\end{center}
\end{advexamplebox}
\global\AdvHideHeadingtrue
\begin{advsolution}
\begin{center}\begin{minipage}{.96\linewidth}\centering
% EMAdvFigureHeading
\raggedright{\bfseries\color{DeepBlue}解答}\par\smallskip\centering

\begin{tikzpicture}[>=Latex,font=\small,scale=.85,every node/.style={fill=none}]

\foreach \shift/\n/\freq in {0/1/60,5/2/120}{\begin{scope}[xshift=\shift cm]
\draw[->] (-.3,0)--(3.5,0) node[right]{$x$};\draw[->] (0,-.9)--(0,1.1) node[above]{$\Psi_{\n}$};\draw[dashed] (3,-.9)--(3,1.1);\node[below] at (3,-.9){$a$};
\draw[DeepBlue,thick,domain=0:3,samples=80] plot(\x,{.7*sin(\freq*\x)});\end{scope}}

\end{tikzpicture}
\figcap{解图：$n=1,2$ 的驻波在两端均为零，第二态还在阱中点有节点；两图纵轴共同按波函数幅值取示意比例。}\label{advA6:quality:A6-quantum-mechanics:a2}
\end{minipage}\end{center}

阱内定态方程为
\[
 -\frac{\hbar^2}{2m}\Psi''=E\Psi,
\quad
 \Psi=A\sin kx+B\cos kx.
\]
边界 $\Psi(0)=0$ 给 $B=0$，$\Psi(a)=0$ 给
$ka=n\pi$，$n=1,2,\ldots$。由 $E=\hbar^2k^2/(2m)$ 得所求能级。
归一化
\[
 1=|A|^2\int_0^a\sin^2\frac{n\pi x}{a}\,\dd x
 =|A|^2\frac a2
\]
给 $A=\sqrt{2/a}$。$n=0$ 只给零波函数，不是物理态；
$E_n\propto a^{-2}$ 与禁闭动能尺度一致。

图~\ref{advA6:quality:A6-quantum-mechanics:a2} 中，两个端点同时为节点，所以只有整数个半波可放入宽度 $a$；$n=0$ 的曲线恒为零，不能归一化，因而不构成态。
\end{advsolution}

\begin{advexamplebox}{有限方势垒的隧穿系数}\label{advA6:ex:square-barrier}
势能
\[
 V(x)=\begin{cases}
 V_0,&0<x<a,\\0,&\text{其他},
 \end{cases}
\qquad 0<E<V_0.
\]
证明透射系数为
\[
 T=\left[
 1+\frac{V_0^2\sinh^2(\kappa a)}
 {4E(V_0-E)}\right]^{-1},
\qquad
 \kappa=\frac{\sqrt{2m(V_0-E)}}{\hbar}.
\]

\begin{center}\begin{minipage}{.96\linewidth}\centering
\begin{tikzpicture}[>=Latex,font=\small,scale=.85,every node/.style={fill=none}]

\draw[->] (-1.5,0)--(4.5,0) node[right]{$x$};\draw[thick] (-1.4,0)--(0,0)--(0,1.5)--(2.5,1.5)--(2.5,0)--(4.2,0);
\draw[dashed] (-1.3,.8)--(4,.8) node[right]{$E$};\node[above] at (1.2,1.5){$V_0$};\node[below] at (0,0){0};\node[below] at (2.5,0){$a$};\node[below] at (-.8,-.5){I};\node[below] at (1.2,-.5){II};\node[below] at (3.5,-.5){III};

\end{tikzpicture}
\figcap{题图：三区方势垒，能量位于阱外势能与垒顶之间。}\label{advA6:quality:A6-quantum-mechanics:q3}
\end{minipage}\end{center}
\end{advexamplebox}
\global\AdvHideHeadingtrue
\begin{advsolution}
\begin{center}\begin{minipage}{.96\linewidth}\centering
% EMAdvFigureHeading
\raggedright{\bfseries\color{DeepBlue}解答}\par\smallskip\centering

\begin{tikzpicture}[>=Latex,font=\small,scale=.85,every node/.style={fill=none}]

\draw[->] (-1.5,0)--(4.5,0) node[right]{$x$};\draw[thick] (-1.4,0)--(0,0)--(0,1.5)--(2.5,1.5)--(2.5,0)--(4.2,0);
\draw[dashed] (-1.3,.8)--(4,.8) node[right]{$E$};\node[above] at (1.2,1.5){$V_0$};\node[below] at (0,0){0};\node[below] at (2.5,0){$a$};\node[below] at (-.8,-.5){I};\node[below] at (1.2,-.5){II};\node[below] at (3.5,-.5){III};

\end{tikzpicture}
\figcap{解图：三区方势垒，能量位于阱外势能与垒顶之间。}\label{advA6:quality:A6-quantum-mechanics:a3}
\end{minipage}\end{center}

三区波函数取
\[
 \Psi_I=\e^{\ii kx}+r\e^{-\ii kx},\quad
 \Psi_{II}=A\e^{\kappa x}+B\e^{-\kappa x},\quad
 \Psi_{III}=t\e^{\ii kx},
\]
$k=\sqrt{2mE}/\hbar$。在 $x=0,a$ 连续匹配 $\Psi,\Psi'$。
为看清消元所依据的四个匹配条件，势垒内部的传播可写成
\[
 \begin{pmatrix}\Psi(a)\\\Psi'(a)\end{pmatrix}
 =\begin{pmatrix}
 \cosh\kappa a&\sinh(\kappa a)/\kappa\\
 \kappa\sinh\kappa a&\cosh\kappa a
 \end{pmatrix}
 \begin{pmatrix}\Psi(0)\\\Psi'(0)\end{pmatrix}.
\]
代入 $\Psi(0)=1+r$、$\Psi'(0)=\mathrm ik(1-r)$，以及 $\Psi(a)=te^{\mathrm ika}$、$\Psi'(a)=\mathrm ikte^{\mathrm ika}$，两式相减即可消去 $r$。整理后，
\[
 t=\frac{\e^{-\ii ka}}
 {\cosh\kappa a+
 \frac{\ii}{2}\left(\frac\kappa k-\frac k\kappa\right)
 \sinh\kappa a}.
\]
左右介质相同，流密度比就是 $T=|t|^2$。利用
\[
 \frac14\left(\frac\kappa k+\frac k\kappa\right)^2
 =\frac{V_0^2}{4E(V_0-E)}
\]
得到题中结果。高厚势垒 $\kappa a\gg1$ 时
\[
 T\simeq
 \frac{16E(V_0-E)}{V_0^2}\e^{-2\kappa a},
\]
显示指数隧穿；$a\to0$ 时 $T\to1$。分母中的 $4$ 不可遗漏。

图~\ref{advA6:quality:A6-quantum-mechanics:a3} 中，I、III区具有同一实波数 $k$，II区具有实衰减系数 $\kappa$；因此透射流比可直接取 $|t|^2$，而势垒内须同时保留两个指数项以满足两个界面的匹配。
\end{advsolution}

\begin{advexamplebox}{一维量子谐振子}\label{advA6:ex:harmonic-oscillator}
证明
\[
 E_n=\left(n+\frac12\right)\hbar\omega,
\]
\[
 \Psi_n(x)=
 \frac{1}{\sqrt{2^nn!}}
 \left(\frac{m\omega}{\pi\hbar}\right)^{1/4}
 H_n(\xi)\e^{-\xi^2/2},
\qquad
 \xi=\sqrt{\frac{m\omega}{\hbar}}x.
\]

\begin{center}\begin{minipage}{.96\linewidth}\centering
\begin{tikzpicture}[>=Latex,font=\small,scale=.85,every node/.style={fill=none}]

\draw[->] (-2,0)--(2.3,0) node[right]{$x$};\draw[->] (0,-.2)--(0,2.2) node[above]{$V$};\draw[DeepBlue,domain=-1.7:1.7,samples=70] plot(\x,{.65*\x*\x});\node[right] at (1.2,.8){$V=m\omega^2x^2/2$};

\end{tikzpicture}
\figcap{题图：一维谐势及坐标原点，束缚态需在左右无穷远可归一化。}\label{advA6:quality:A6-quantum-mechanics:q4}
\end{minipage}\end{center}
\end{advexamplebox}
\global\AdvHideHeadingtrue
\begin{advsolution}
\begin{center}\begin{minipage}{.96\linewidth}\centering
% EMAdvFigureHeading
\raggedright{\bfseries\color{DeepBlue}解答}\par\smallskip\centering

\begin{tikzpicture}[>=Latex,font=\small,scale=.85,every node/.style={fill=none}]

\draw[->] (-1.8,0)--(2,0) node[right]{$x$};\draw[DeepBlue,domain=-1.5:1.5] plot(\x,{.7*\x*\x});\foreach \y/\n in {.35/0,.85/1,1.35/2}{\draw[thick] (-1.5,\y)--(1.5,\y);\node[right] at (1.5,\y){$E_{\n}$};}\draw[->] (-1.7,.4)--(-1.7,.8) node[left]{$a^\dagger$};

\end{tikzpicture}
\figcap{解图：谐势内的等间隔能级，最低能级高于势能最低点。}\label{advA6:quality:A6-quantum-mechanics:a4}
\end{minipage}\end{center}

定义
\[
 \widehat a=
 \sqrt{\frac{m\omega}{2\hbar}}\widehat x
 +\frac{\ii}{\sqrt{2m\hbar\omega}}\widehat p,\qquad
 [\widehat a,\widehat a^\dagger]=1.
\]
Hamilton 量分解为
\[
 \widehat H=\hbar\omega
 \left(\widehat a^\dagger\widehat a+\frac12\right).
\]
数算符 $\widehat N=\widehat a^\dagger\widehat a$ 半正定。
若其最低本征值不是 $0$，连续施用 $\widehat a$ 会得到更低本征值，矛盾；
故 $\widehat a|0\rangle=0$，$E_0=\hbar\omega/2$。
再由
\[
 |n\rangle=\frac{(\widehat a^\dagger)^n}{\sqrt{n!}}|0\rangle
\]
得到 $E_n=(n+1/2)\hbar\omega$。

位置表象中 $\widehat a\Psi_0=0$ 给
$\Psi_0\propto\e^{-\xi^2/2}$；反复作用
$\widehat a^\dagger$ 并用 Rodrigues 公式
\[
 H_n(\xi)=(-1)^n\e^{\xi^2}
 \frac{\dd^n}{\dd\xi^n}\e^{-\xi^2}
\]
得到题中归一化波函数。$n$ 偶、奇时波函数宇称分别为偶、奇，
与对称势能相符。

图~\ref{advA6:quality:A6-quantum-mechanics:a4} 中，抛物势的两端都上升，允许态要同时满足两端衰减；引入升降算符后，基态被降算符湮灭，逐级升高能量而不改变相邻间隔 $\hbar\omega$。
\end{advsolution}

\subsection{Hilbert 空间与 Dirac 符号}

\begin{advdefinition}{Hilbert 空间}\label{advA6:def:hilbert}
Hilbert 空间 $\mathcal H$ 是带正定 Hermitian 内积的完备复线性空间。
内积对第二个变量线性、对第一个变量反线性：
\[
 \langle a|\lambda b+\mu c\rangle
 =\lambda\langle a|b\rangle+\mu\langle a|c\rangle,
\quad
 \langle b|a\rangle=\langle a|b\rangle^*.
\]
范数与距离为
\[
 \|a\|=\sqrt{\langle a|a\rangle},\qquad
 d(a,b)=\|a-b\|.
\]
完备性指每个 Cauchy 序列都收敛到空间内的元素。
\end{advdefinition}

\begin{advexamplebox}{Hilbert 空间中的 Cauchy--Schwarz 不等式}\label{advA6:ex:cauchy-schwarz}
证明
\[
 |\langle a|b\rangle|^2
 \le\langle a|a\rangle\langle b|b\rangle.
\]

\begin{center}\begin{minipage}{.96\linewidth}\centering
\begin{tikzpicture}[>=Latex,font=\small,scale=.85,every node/.style={fill=none}]

\draw[DeepBlue,thick,->] (0,0)--(2.7,0) node[right]{$|a\rangle$};\draw[DeepBlue,thick,->] (0,0)--(1.5,1.6) node[above]{$|b\rangle$};\fill (0,0) circle (1pt) node[below]{$0$};

\end{tikzpicture}
\figcap{题图：两个Hilbert空间向量；图示为它们张成子空间的示意，复相位仍由内积规定。}\label{advA6:quality:A6-quantum-mechanics:q5}
\end{minipage}\end{center}
\end{advexamplebox}
\global\AdvHideHeadingtrue
\begin{advsolution}
\begin{center}\begin{minipage}{.96\linewidth}\centering
% EMAdvFigureHeading
\raggedright{\bfseries\color{DeepBlue}解答}\par\smallskip\centering

\begin{tikzpicture}[>=Latex,font=\small,scale=.85,every node/.style={fill=none}]

\draw[DeepBlue,thick,->] (0,0)--(2.7,0) node[right]{$|a\rangle$};\draw[DeepBlue,thick,->] (0,0)--(1.5,1.6) node[above]{$|b\rangle$};\fill (0,0) circle (1pt) node[below]{$0$};
\draw[dashed] (1.5,0)--(1.5,1.6);\fill (1.5,0) circle (1.5pt) node[below]{$\lambda|a\rangle$};\node[right] at (1.5,.8){$|b\rangle-\lambda|a\rangle$};
\end{tikzpicture}
\figcap{解图：经适当相位选择后，投影与正交余量的实平面截面。}\label{advA6:quality:A6-quantum-mechanics:a5}
\end{minipage}\end{center}

按解图把 $b$ 分解为沿 $a$ 的投影和正交余量。若 $a=0$ 显然成立；
若 $a\ne0$，对任意复数 $\lambda$ 都有
\[
 0\le\|b-\lambda a\|^2.
\]
取 $\lambda=\langle a|b\rangle/\langle a|a\rangle$，余量与 $a$ 正交，且
\[
 0\le\langle b|b\rangle
 -\frac{|\langle a|b\rangle|^2}{\langle a|a\rangle}.
\]
乘以正数 $\langle a|a\rangle$ 即得不等式。等号要求余量为零，
即 $b=\lambda a$；连同零向量情形，等号当且仅当两个向量线性相关。

图~\ref{advA6:quality:A6-quantum-mechanics:a5} 中，把 $|b\rangle$ 分成沿 $|a\rangle$ 的投影与正交余量，余量范数非负给出不等式；等号要求余量为零，即两个向量线性相关。
\end{advsolution}

\begin{advexamplebox}{Hilbert 空间中的三角不等式}\label{advA6:ex:triangle}
证明
\[
 d(a,b)\le d(a,c)+d(c,b).
\]

\begin{center}\begin{minipage}{.96\linewidth}\centering
\begin{tikzpicture}[>=Latex,font=\small,scale=.85,every node/.style={fill=none}]

\coordinate (a) at (0,0);\coordinate (b) at (3,.3);\coordinate (c) at (1.3,1.7);\draw[thick] (a)--(b)--(c)--cycle;\node[left] at (a){$a$};\node[right] at (b){$b$};\node[above] at (c){$c$};\node[below] at (1.5,.1){$d(a,b)$};\node[left] at (.6,.9){$d(a,c)$};\node[right] at (2.2,1){$d(c,b)$};

\end{tikzpicture}
\figcap{题图：三个态向量在其仿射实平面中的距离三角形，距离由Hilbert范数定义。}\label{advA6:quality:A6-quantum-mechanics:q6}
\end{minipage}\end{center}
\end{advexamplebox}
\global\AdvHideHeadingtrue
\begin{advsolution}
\begin{center}\begin{minipage}{.96\linewidth}\centering
% EMAdvFigureHeading
\raggedright{\bfseries\color{DeepBlue}解答}\par\smallskip\centering

\begin{tikzpicture}[>=Latex,font=\small,scale=.85,every node/.style={fill=none}]

\coordinate (a) at (0,0);\coordinate (b) at (3,.3);\coordinate (c) at (1.3,1.7);\draw[thick] (a)--(b)--(c)--cycle;\node[left] at (a){$a$};\node[right] at (b){$b$};\node[above] at (c){$c$};\node[below] at (1.5,.1){$d(a,b)$};\node[left] at (.6,.9){$d(a,c)$};\node[right] at (2.2,1){$d(c,b)$};

\end{tikzpicture}
\figcap{解图：三个态向量在其仿射实平面中的距离三角形，距离由Hilbert范数定义。}\label{advA6:quality:A6-quantum-mechanics:a6}
\end{minipage}\end{center}

令 $x=a-c$、$y=c-b$，则 $a-b=x+y$。利用上题，
\[
\begin{aligned}
\|x+y\|^2
&=\|x\|^2+\|y\|^2+2\operatorname{Re}\langle x|y\rangle\\
&\le\|x\|^2+\|y\|^2+2\|x\|\|y\|
=(\|x\|+\|y\|)^2.
\end{aligned}
\]
两边取非负平方根即得
$\|a-b\|\le\|a-c\|+\|c-b\|$。

图~\ref{advA6:quality:A6-quantum-mechanics:a6} 中，沿 $a\to c\to b$ 的向量位移相加等于 $a\to b$，展开范数平方后用Cauchy--Schwarz控制交叉项，才得到长度的三角不等式。
\end{advsolution}

\begin{advdefinition}{共射线的归一化态矢}\label{advA6:def:same-ray}
若两个归一化态矢仅差一个全局相位，
\[
 |\psi'\rangle=\e^{\ii\theta}|\psi\rangle,
\]
则称它们在同一射线上。任何转移概率和可观测量期望都不依赖该全局相位。
\end{advdefinition}

\begin{advdefinition}{态等价与纯量子态}\label{advA6:def:state-equivalence}
任意非零共线态矢
$|\psi'\rangle=\lambda\e^{\ii\theta}|\psi\rangle$（$\lambda>0$）
代表同一纯态。纯量子态是非零 Hilbert 向量在此等价关系下的等价类；
计算时通常选取范数为 $1$ 的代表元。一般混态需用密度算符描述，不能都表示为单条态矢射线。
\end{advdefinition}

\begin{advdefinition}{左矢、右矢与对偶}\label{advA6:def:bra-ket}
空间中的向量记为右矢 $|a\rangle$。由 Riesz 表示定理，它唯一对应连续线性泛函
\[
 \langle a|:\ |\psi\rangle\mapsto\langle a|\psi\rangle,
\]
称为左矢。映射 $|a\rangle\mapsto\langle a|$ 是反线性的：
\[
 (\lambda|a\rangle+\mu|b\rangle)^\dagger
 =\lambda^*\langle a|+\mu^*\langle b|.
\]
\end{advdefinition}

\begin{advdefinition}{线性算符与反线性算符}\label{advA6:def:operator}
线性算符满足
\[
 A(\lambda|a\rangle+\mu|b\rangle)
 =\lambda A|a\rangle+\mu A|b\rangle.
\]
反线性算符则把系数取复共轭。时间反演是典型反线性对称操作；
一般可观测量由线性算符描述。
\end{advdefinition}

\begin{advexamplebox}{反线性算符诱导的对偶作用}\label{advA6:ex:antioperator}
设 $\Theta$ 是反线性算符。证明
\[
 \langle\Theta a|\Theta b\rangle
\]
对 $|a\rangle$ 线性、对 $|b\rangle$ 反线性，并说明反幺正条件如何保持概率。

\begin{center}\begin{minipage}{.96\linewidth}\centering
\begin{tikzpicture}[>=Latex,font=\small,scale=.85,every node/.style={fill=none}]
\draw[->] (-.4,0)--(3,0) node[right]{$x$};\draw[->] (0,-.4)--(0,2) node[above]{$y$};
\draw[DeepBlue,->] (0,0)--(1.3,1.1) node[above]{$\lambda$};\draw[dashed] (1.3,-1.1)--(1.3,1.1);\node[below] at (1.3,-1.1){$\lambda^*$};\node[right] at (2,1.5){复系数平面};

\end{tikzpicture}
\figcap{题图：反线性作用会共轭复系数，图只表示复系数而非Hilbert空间的物理位置。}\label{advA6:quality:A6-quantum-mechanics:q7}
\end{minipage}\end{center}
\end{advexamplebox}
\global\AdvHideHeadingtrue
\begin{advsolution}
\begin{center}\begin{minipage}{.96\linewidth}\centering
% EMAdvFigureHeading
\raggedright{\bfseries\color{DeepBlue}解答}\par\smallskip\centering

\begin{tikzpicture}[>=Latex,font=\small,scale=.85,every node/.style={fill=none}]
\draw[->] (-.4,0)--(3,0) node[right]{$x$};\draw[->] (0,-.4)--(0,2) node[above]{$y$};
\draw[DeepBlue,->] (0,0)--(1.3,1.1) node[above]{$\lambda$};\draw[dashed] (1.3,-1.1)--(1.3,1.1);\node[below] at (1.3,-1.1){$\lambda^*$};\node[right] at (2,1.5){复系数平面};

\end{tikzpicture}
\figcap{解图：反线性作用会共轭复系数，图只表示复系数而非Hilbert空间的物理位置。}\label{advA6:quality:A6-quantum-mechanics:a7}
\end{minipage}\end{center}

由反线性，
$\Theta(\lambda a+\mu b)=\lambda^*\Theta a+\mu^*\Theta b$。
因此
\[
 \langle\Theta(\lambda a+\mu b)|\Theta c\rangle
 =\lambda\langle\Theta a|\Theta c\rangle
 +\mu\langle\Theta b|\Theta c\rangle,
\]
而
\[
 \langle\Theta c|\Theta(\lambda a+\mu b)\rangle
 =\lambda^*\langle\Theta c|\Theta a\rangle
 +\mu^*\langle\Theta c|\Theta b\rangle.
\]
若
\[
 \langle\Theta a|\Theta b\rangle
 =\langle a|b\rangle^*=\langle b|a\rangle,
\]
若 $\Theta$ 还把整个空间映满，则称为反幺正；有限维等维情形下满射性由上述保内积条件推出。此时
\[
 |\langle\Theta a|\Theta b\rangle|^2
 =|\langle a|b\rangle|^2,
\]
所有转移概率保持不变。这里不能把反线性算符的“伴随”当作普通矩阵转置，
必须同时追踪复共轭的作用位置。

图~\ref{advA6:quality:A6-quantum-mechanics:a7} 中，$\lambda$ 与 $\lambda^*$ 在复系数平面关于实轴对称，故 $\Theta(\lambda|a\rangle)=\lambda^*\Theta|a\rangle$；再结合内积第一槽反线性，才能逐槽判断题中双线性类型。
\end{advsolution}

\begin{advdefinition}{伴随算符}\label{advA6:def:adjoint}
稠密定义的线性算符 $A$ 的伴随 $A^\dagger$ 由
\[
 \langle\psi|A\phi\rangle
 =\langle A^\dagger\psi|\phi\rangle
\]
定义，其中还必须指定两者的定义域。有限维中，
$[A^\dagger]=[A]^{*\,\mathsf T}$。
\end{advdefinition}

\begin{advdefinition}{可观测量}\label{advA6:def:observable}
可观测量由自伴算符表示。自伴不仅要求形式上
$A^\dagger=A$，还要求定义域相同。
其谱为实数，并由谱投影给出完备分解；有限维时即 Hermitian 矩阵。
\end{advdefinition}

\subsection{位置表象、动量表象和常用算符}

位置与动量广义本征态满足
\[
 \langle x'|x\rangle=\delta(x-x'),\qquad
 \langle p'|p\rangle=\delta(p-p'),
\]
\[
 \langle x|p\rangle=
 \frac1{\sqrt{2\pi\hbar}}\e^{\ii px/\hbar}.
\]
完备性为
\[
 \int|x\rangle\langle x|\,\dd x=I,\qquad
 \int|p\rangle\langle p|\,\dd p=I.
\]
在位置表象中
\[
 \widehat x=x,\qquad
 \widehat p=-\ii\hbar\frac{\dd}{\dd x},\qquad
 [\widehat x,\widehat p]=\ii\hbar.
\]
平移算符
\[
 T(\varepsilon)=\exp\!\left(-\frac{\ii}{\hbar}
 \varepsilon\widehat p\right)
\]
作用为
$(T(\varepsilon)\psi)(x)=\psi(x-\varepsilon)$；
正负号取决于是主动平移态还是被动平移坐标，必须保持约定一致。

\begin{advexamplebox}{投影、位置、动量和平移算符的伴随性质}\label{advA6:ex:operator-adjoints}
证明位置/动量谱投影、位置算符和动量算符为自伴对象，
而非零平移算符一般不是 Hermitian，却是幺正算符，且
$T(\varepsilon)^\dagger=T(-\varepsilon)$。

\begin{center}\begin{minipage}{.96\linewidth}\centering
\begin{tikzpicture}[>=Latex,font=\small,scale=.85,every node/.style={fill=none}]

\draw[->] (-1,0)--(4.5,0) node[right]{$x$};\draw[DeepBlue,domain=-.8:2,samples=60] plot(\x,{exp(-3*(\x-.5)^2)});\draw[DeepBlue,dashed,domain=.5:3.5,samples=60] plot(\x,{exp(-3*(\x-2)^2)});\draw[->] (.5,1.3)--(2,1.3) node[midway,above]{$+\epsilon$};\node[below] at (2,-.5){虚线为 $\psi(x-\epsilon)$};

\end{tikzpicture}
\figcap{题图：实轴上的波包与主动平移后的波包，图中取 $\epsilon>0$。}\label{advA6:quality:A6-quantum-mechanics:q8}
\end{minipage}\end{center}
\end{advexamplebox}
\global\AdvHideHeadingtrue
\begin{advsolution}
\begin{center}\begin{minipage}{.96\linewidth}\centering
% EMAdvFigureHeading
\raggedright{\bfseries\color{DeepBlue}解答}\par\smallskip\centering

\begin{tikzpicture}[>=Latex,font=\small,scale=.85,every node/.style={fill=none}]

\draw[->] (-1,0)--(4.5,0) node[right]{$x$};\draw[DeepBlue,domain=-.8:2,samples=60] plot(\x,{exp(-3*(\x-.5)^2)});\draw[DeepBlue,dashed,domain=.5:3.5,samples=60] plot(\x,{exp(-3*(\x-2)^2)});\draw[->] (.5,1.3)--(2,1.3) node[midway,above]{$+\epsilon$};\node[below] at (2,-.5){虚线为 $\psi(x-\epsilon)$};

\end{tikzpicture}
\figcap{解图：实轴上的波包与主动平移后的波包，图中取 $\epsilon>0$。}\label{advA6:quality:A6-quantum-mechanics:a8}
\end{minipage}\end{center}

先用有限测量区间 $\Delta$ 写位置谱投影：
\[
 (P_x(\Delta)\psi)(x)=\mathbf1_\Delta(x)\psi(x),\qquad
 P_x(\Delta)^\dagger=P_x(\Delta),\quad P_x(\Delta)^2=P_x(\Delta).
\]
动量谱投影在动量表象中有同样的乘法形式，或由 Fourier 酉变换从该形式得到。
$|x\rangle\langle x|$ 是形式上的谱密度，不是单个正规的秩一投影。
$\widehat x$ 在 $D(\widehat x)=\{\psi\in L^2:x\psi\in L^2\}$ 上是自伴乘法算符。对动量算符，
\[
 \langle\phi|\widehat p\psi\rangle
 =-\ii\hbar\int\phi^*\psi'\,\dd x
 =-\ii\hbar[\phi^*\psi]_{\rm bdry}
 +\ii\hbar\int(\phi^*)'\psi\,\dd x.
\]
在波函数于无穷远消失或采用相容边界条件时，边界项为零，得到
$\langle\phi|\widehat p\psi\rangle
=\langle\widehat p\phi|\psi\rangle$；
在整条实线上取 $D(\widehat p)=H^1(\mathbb R)$，可核验伴随定义域相同，因而自伴。有限区间时仅消掉边界项只能证明对称，还须另行指定自伴边界条件；例如两端都取Dirichlet条件的动量算符并不自动自伴。

因为 $\widehat p$ 自伴，
\[
 T(\varepsilon)^\dagger
 =\exp\!\left(+\frac{\ii\varepsilon\widehat p}{\hbar}\right)
 =T(-\varepsilon)=T(\varepsilon)^{-1}.
\]
所以 $T$ 幺正。除 $\varepsilon=0$ 或特殊离散谱情形外，
$T^\dagger\ne T$，本征值 $\e^{-\ii\varepsilon p/\hbar}$ 也不是实数，
故它本身不是可观测量。

图~\ref{advA6:quality:A6-quantum-mechanics:a8} 中，主动向右平移作用在函数上是自变量替换 $x\mapsto x-\epsilon$，其逆为 $-\epsilon$；内积保持说明幺正，不能由此直接断言平移算符自伴。
\end{advsolution}

\begin{advexamplebox}{位置、动量与平移的广义本征态}\label{advA6:ex:generalized-eigenstates}
说明位置算符、动量算符和平移算符的本征态与完备性；
解释连续谱密度、区间投影与位置算符三者的区别。

\begin{center}\begin{minipage}{.96\linewidth}\centering
\begin{tikzpicture}[>=Latex,font=\small,scale=.85,every node/.style={fill=none}]

\draw[->] (-1,0)--(4.7,0) node[right]{$x$};\draw[very thick,DeepBlue] (.5,0)--(3,0);\draw (.5,-.15)--(.5,.15) (3,-.15)--(3,.15);\node[below] at (.5,0){$a$};\node[below] at (3,0){$b$};\draw[<->] (1.5,.7)--(1.9,.7) node[midway,above]{$\mathrm dx$};\node[above] at (3.2,1.3){位置谱区间投影};

\end{tikzpicture}
\figcap{题图：连续位置谱中的一个区间及微元，区间投影与点本征态不同。}\label{advA6:quality:A6-quantum-mechanics:q9}
\end{minipage}\end{center}
\end{advexamplebox}
\global\AdvHideHeadingtrue
\begin{advsolution}
\begin{center}\begin{minipage}{.96\linewidth}\centering
% EMAdvFigureHeading
\raggedright{\bfseries\color{DeepBlue}解答}\par\smallskip\centering

\begin{tikzpicture}[>=Latex,font=\small,scale=.85,every node/.style={fill=none}]

\draw[->] (-1,0)--(4.7,0) node[right]{$x$};\draw[very thick,DeepBlue] (.5,0)--(3,0);\draw (.5,-.15)--(.5,.15) (3,-.15)--(3,.15);\node[below] at (.5,0){$a$};\node[below] at (3,0){$b$};\draw[<->] (1.5,.7)--(1.9,.7) node[midway,above]{$\mathrm dx$};\node[above] at (3.2,1.3){位置谱区间投影};

\end{tikzpicture}
\figcap{解图：连续位置谱中的一个区间及微元，区间投影与点本征态不同。}\label{advA6:quality:A6-quantum-mechanics:a9}
\end{minipage}\end{center}

\[
 \widehat x|x_0\rangle=x_0|x_0\rangle,\qquad
 \widehat p|p_0\rangle=p_0|p_0\rangle.
\]
$|x_0\rangle$ 和 $|p_0\rangle$ 不是 Hilbert 空间中的可归一向量，
而是 rigged Hilbert space 中的广义向量；其 $\delta$ 归一化与积分完备性如上。
任意态均可展开
\[
 |\psi\rangle=\int\psi(x)|x\rangle\,\dd x
 =\int\tilde\psi(p)|p\rangle\,\dd p.
\]
平移算符与动量算符具有同一组广义本征态：
\[
 T(\varepsilon)|p\rangle
 =\e^{-\ii\varepsilon p/\hbar}|p\rangle.
\]
严格位置测量由投影值测度
\[
 P_x(\Delta)=\int_\Delta|x\rangle\langle x|\,\dd x
\]
描述。非平凡区间的投影有 $0,1$ 两个谱值，回答“是否在这个区间”；
位置算符则是 $\widehat x=\int x\,\dd P_x(x)$，给出连续的位置值。
在 $L^2(\mathbb R)$ 的通常位置表象中，单点集合测度为零，所以
$P_x(\{x_0\})=0$。形式符号 $|x_0\rangle\langle x_0|$ 不能直接平方成自身，
否则会遇到无意义的 $\delta(0)$；它不是具有 $0,1$ 两个谱值的普通秩一投影。

图~\ref{advA6:quality:A6-quantum-mechanics:a9} 中，蓝色区间对应 $P_{[a,b]}=\int_a^b|x\rangle\langle x|\,\mathrm dx$；点核是算符值密度，位置算符还需再乘坐标 $x$ 积分，三者的量纲和作用都不同。
\end{advsolution}

\subsection{有限维态、矩阵与自旋}

在二维自旋空间中取基
$|1\rangle=|\uparrow\rangle$、
$|2\rangle=|\downarrow\rangle$，
态写为
\[
 |\psi\rangle=
 \psi_1|\uparrow\rangle+\psi_2|\downarrow\rangle,
\qquad
 |\psi_1|^2+|\psi_2|^2=1.
\]

\begin{advexamplebox}{算符矩阵与伴随矩阵}\label{advA6:ex:matrix-adjoint}
在正交归一基中证明
\[
 [A^\dagger]=[A]^{*\,\mathsf T}.
\]
并说明 Hermitian 矩阵的元素条件。

\begin{center}\begin{minipage}{.96\linewidth}\centering
\begin{tikzpicture}[>=Latex,font=\small,scale=.85,every node/.style={fill=none}]

\fill[DeepBlue!12] (.8,1.2) rectangle (1.6,1.8);\draw (0,0) grid[xstep=.8,ystep=.6] (2.4,1.8);\node at (1.2,1.5){$A_{ij}$};\node[above] at (1.2,2){正交归一基};\node[below] at (1.2,-.4){行 $i$，列 $j$（以非对角元示意）};

\end{tikzpicture}
\figcap{题图：算符矩阵的一个行列元素，索引由所选正交归一基决定。}\label{advA6:quality:A6-quantum-mechanics:q10}
\end{minipage}\end{center}
\end{advexamplebox}
\global\AdvHideHeadingtrue
\begin{advsolution}
\begin{center}\begin{minipage}{.96\linewidth}\centering
% EMAdvFigureHeading
\raggedright{\bfseries\color{DeepBlue}解答}\par\smallskip\centering

\begin{tikzpicture}[>=Latex,font=\small,scale=.85,every node/.style={fill=none}]

\fill[DeepBlue!12] (.8,1.2) rectangle (1.6,1.8);\draw (0,0) grid[xstep=.8,ystep=.6] (2.4,1.8);\node at (1.2,1.5){$A_{ij}$};\node[above] at (1.2,2){$[A]$};
\begin{scope}[xshift=4.2cm]\fill[DeepBlue!12] (0,.6) rectangle (.8,1.2);\draw (0,0) grid[xstep=.8,ystep=.6] (2.4,1.8);\node at (.4,.9){$A_{ij}^*$};\node[above] at (1.2,2){$[A^\dagger]$};\end{scope}
\node[below] at (3.3,-.5){原来的 $(i,j)$ 元素进入伴随矩阵的 $(j,i)$ 位置并取复共轭};

\end{tikzpicture}
\figcap{解图：非对角矩阵元的位置对换与复共轭同时进行。}\label{advA6:quality:A6-quantum-mechanics:a10}
\end{minipage}\end{center}

矩阵元 $A_{ij}=\langle i|A|j\rangle$。由伴随定义，
\[
 (A^\dagger)_{ij}
 =\langle i|A^\dagger|j\rangle
 =\langle j|A|i\rangle^*
 =A_{ji}^*.
\]
所以 $[A^\dagger]=[A]^{*\,\mathsf T}$。
若 $A=A^\dagger$，则
$A_{ij}=A_{ji}^*$，对角元 $A_{ii}$ 必为实数；
反之满足这些元素关系的矩阵就是 Hermitian 矩阵。
投影和 Pauli 矩阵均满足该条件；升、降算符互为伴随。

图~\ref{advA6:quality:A6-quantum-mechanics:a10} 中，行列交换来自交换内积的两个基向量，复共轭来自伴随定义；因此必须把转置和共轭同时执行，不能仅对每个元素取共轭而保持索引不变。
\end{advsolution}

Hermitian 算符的本征值实、不同本征值的本征态正交。
若 $A|a\rangle=a|a\rangle$，
\[
 a\langle a|a\rangle
 =\langle a|A|a\rangle
 =\langle a|A^\dagger|a\rangle
 =a^*\langle a|a\rangle,
\]
故 $a=a^*$。对 $a\ne b$，
$(a-b)\langle b|a\rangle=0$，故正交。
离散谱的完备性写成
$\sum_{i,j}|\psi_{ij}\rangle\langle\psi_{ij}|=I$；
连续谱则以积分和 Dirac $\delta$ 取代求和与 Kronecker $\delta$。

\begin{advtheorem}{投影测量规则}\label{advA6:thm:copenhagen}
设可观测量具有谱分解
\[
 A=\sum_iA_iP_i,
\]
其中 $P_i$ 是本征值 $A_i$ 的简并本征空间投影。归一化态
$|\psi\rangle$ 测得 $A_i$ 的概率为
\[
 \boxed{p_i=\langle\psi|P_i|\psi\rangle}.
\]
在理想选择性测量中，对 $p_i>0$ 的实际结果，得到 $A_i$ 后的态为
\[
 \boxed{|\psi_i'\rangle=
 \frac{P_i|\psi\rangle}{\sqrt{p_i}}}.
\]
连续谱用投影值测度替代离散求和。
\end{advtheorem}

\begin{advdefinition}{期望值、矩阵元与跃迁振幅}\label{advA6:def:expectation-transition}
归一化态中算符 $A$ 的期望值为 $\langle A\rangle_\psi=\langle\psi|A|\psi\rangle$。
矩阵元 $A_{\phi\psi}=\langle\phi|A|\psi\rangle$ 与 $A$ 具有相同量纲，
不能无条件把它的模方当作跃迁概率。真正的演化振幅是
\[
 \mathcal A_{\psi\to\phi}(t)=\langle\phi|U(t,0)|\psi\rangle,
\]
其中 $U$ 是时间演化算符。对时间无关的 $H_0$ 的不同本征态 $i,f$，
若附加弱扰动 $V(t)$，相互作用表象的一阶振幅为
\[
 c_f^{(1)}(t)=-\frac{\ii}{\hbar}\int_0^t
 \langle f|V(t')|i\rangle\,
 \e^{\ii(E_f-E_i)t'/\hbar}\,\dd t'.
\]
积分和 $1/\hbar$ 使振幅无量纲；在一阶微扰适用且该一阶振幅非零时，$|c_f^{(1)}|^2$ 给出相应跃迁概率的最低非零阶。若一阶振幅被选择定则禁止，则需继续计算高阶项。
\end{advdefinition}

\begin{advdefinition}{对易子}\label{advA6:def:commutator}
两个算符的对易子为
\[
 [A,B]=AB-BA.
\]
$[A,B]=0$ 时称两者对易或相容。对易关系需在共同稠密定义域上理解；
仅有形式微分表达式相等还不够。
\end{advdefinition}

\begin{advtheorem}{对易自伴算符的共同本征分解}\label{advA6:thm:commuting-observables}
有限维中，两个对易 Hermitian 算符可被同一幺正矩阵同时对角化，
因而存在共同正交归一本征基。无限维中，强对易的自伴算符具有共同谱测度。
\end{advtheorem}

\noindent\textit{证明（有限维）。}
取 $A$ 的某一本征空间 $V_a$。若 $|x\rangle\in V_a$，
\[
 A(B|x\rangle)=B(A|x\rangle)=aB|x\rangle,
\]
故 $B$ 把 $V_a$ 映入自身。在每个 $V_a$ 中限制的 $B$ 仍为 Hermitian，
可选取其正交归一本征基。对所有 $A$ 的本征空间合并这些基，
便得到 $A,B$ 的共同正交归一本征基。
一组对易可观测量若共同量子数能唯一标记每个共同本征态，称为对易可观测量完全集。

\begin{advdefinition}{Pauli 矩阵}\label{advA6:def:pauli}
\[
 \sigma_x=\begin{pmatrix}0&1\\1&0\end{pmatrix},\quad
 \sigma_y=\begin{pmatrix}0&-\ii\\\ii&0\end{pmatrix},\quad
 \sigma_z=\begin{pmatrix}1&0\\0&-1\end{pmatrix}.
\]
它们满足
\[
 \sigma_i\sigma_j=\delta_{ij}I+\ii\varepsilon_{ijk}\sigma_k,
\qquad
 [\sigma_i,\sigma_j]=2\ii\varepsilon_{ijk}\sigma_k.
\]
自旋算符为 $\vect S=(\hbar/2)\vect\sigma$。
\end{advdefinition}

\begin{advexamplebox}{Pauli 本征态、投影与升降算符}\label{advA6:ex:pauli-eigenstates}
\begin{enumerate}
\item 求 $P_\uparrow=|\uparrow\rangle\langle\uparrow|$、
$P_\downarrow$ 的本征值和本征空间；
\item 求 $\sigma_x,\sigma_y,\sigma_z$ 的本征态；
\item 定义
$\sigma_\pm=(\sigma_x\pm\ii\sigma_y)/2$，
说明其作用并证明没有非零本征值。
\end{enumerate}

\begin{center}\begin{minipage}{.96\linewidth}\centering
\begin{tikzpicture}[>=Latex,font=\small,scale=.85,every node/.style={fill=none}]

\draw[->] (0,-1.4)--(0,1.7) node[above]{$z$};\draw[->] (-1.5,0)--(1.7,0) node[right]{$x$};\draw[->] (-.9,-.7)--(.9,.7) node[above right]{$y$};\draw[dashed] (0,0) circle (1.1);\fill[DeepBlue] (0,1.1) circle (2pt) node[left]{$|\uparrow\rangle$};\fill[DeepBlue] (0,-1.1) circle (2pt) node[left]{$|\downarrow\rangle$};

\end{tikzpicture}
\figcap{题图：自旋1/2态的Bloch球示意，南北极对应 $\sigma_z$ 的两个本征方向。}\label{advA6:quality:A6-quantum-mechanics:q11}
\end{minipage}\end{center}
\end{advexamplebox}
\global\AdvHideHeadingtrue
\begin{advsolution}
\begin{center}\begin{minipage}{.96\linewidth}\centering
% EMAdvFigureHeading
\raggedright{\bfseries\color{DeepBlue}解答}\par\smallskip\centering

\begin{tikzpicture}[>=Latex,font=\small,scale=.85,every node/.style={fill=none}]

\draw[->] (0,-1.4)--(0,1.7) node[above]{$z$};\draw[->] (-1.5,0)--(1.7,0) node[right]{$x$};\draw[->] (-.9,-.7)--(.9,.7) node[above right]{$y$};\draw[dashed] (0,0) circle (1.1);\fill[DeepBlue] (0,1.1) circle (2pt) node[left]{$|\uparrow\rangle$};\fill[DeepBlue] (0,-1.1) circle (2pt) node[left]{$|\downarrow\rangle$};

\end{tikzpicture}
\figcap{解图：自旋1/2态的Bloch球示意，南北极对应 $\sigma_z$ 的两个本征方向。}\label{advA6:quality:A6-quantum-mechanics:a11}
\end{minipage}\end{center}

\[
 P_\uparrow=\begin{pmatrix}1&0\\0&0\end{pmatrix},\qquad
 P_\downarrow=\begin{pmatrix}0&0\\0&1\end{pmatrix}.
\]
两者本征值均为 $0,1$，各自的完整本征基是
$|\uparrow\rangle,|\downarrow\rangle$；
投影本身可作为二值可观测量，但不是“位置或自旋分量”本身。

三个 Pauli 矩阵的 $\pm1$ 本征态可取
\[
 |x\pm\rangle=\frac1{\sqrt2}
 (|\uparrow\rangle\pm|\downarrow\rangle),
\]
\[
 |y\pm\rangle=\frac1{\sqrt2}
 (|\uparrow\rangle\pm\ii|\downarrow\rangle),
\qquad
 |z+\rangle=|\uparrow\rangle,\quad
 |z-\rangle=|\downarrow\rangle.
\]
每组正交完备。

\[
 \sigma_+=
 \begin{pmatrix}0&1\\0&0\end{pmatrix},\qquad
 \sigma_-=
 \begin{pmatrix}0&0\\1&0\end{pmatrix}
\]
（若把基矢次序反置，矩阵位置相应互换）。它们把一个 $z$ 本征态映到另一个，
并把最高/最低态映为零。由于
$\sigma_\pm^2=0$，任一本征值 $\lambda$ 必满足 $\lambda^2=0$，
故只有 $\lambda=0$；核空间虽有非零向量，但不存在非零本征值。

图~\ref{advA6:quality:A6-quantum-mechanics:a11} 中，投影算符保留各自极点的态并湮灭正交态；$\sigma_+$ 从下态升到上态而在上态作用为零，$\sigma_-$ 相反。Bloch方向不等于一个粒子的空间轨道。
\end{advsolution}

\begin{advexamplebox}{Pauli 矩阵指数与自旋旋转}\label{advA6:ex:pauli-exponential}
对单位矢量 $\vect n$，求
\[
 U=\exp(\ii\theta\,\vect n\cdot\vect\sigma).
\]

\begin{center}\begin{minipage}{.96\linewidth}\centering
\begin{tikzpicture}[>=Latex,font=\small,scale=.85,every node/.style={fill=none}]

\draw[->] (0,-1.4)--(0,1.7) node[above]{$z$};\draw[->] (-1.5,0)--(1.7,0) node[right]{$x$};\draw[->] (-.9,-.7)--(.9,.7) node[above right]{$y$};\draw[dashed] (0,0) circle (1.1);\fill[DeepBlue] (0,1.1) circle (2pt) node[left]{$|\uparrow\rangle$};\fill[DeepBlue] (0,-1.1) circle (2pt) node[left]{$|\downarrow\rangle$};
\draw[DeepBlue,->,thick] (0,0)--(1,1.3) node[right]{$\boldsymbol n$};
\end{tikzpicture}
\figcap{题图：单位方向 $\boldsymbol n$ 决定自旋算符指数的轴，指数参数为 $\theta$。}\label{advA6:quality:A6-quantum-mechanics:q12}
\end{minipage}\end{center}
\end{advexamplebox}
\global\AdvHideHeadingtrue
\begin{advsolution}
\begin{center}\begin{minipage}{.96\linewidth}\centering
% EMAdvFigureHeading
\raggedright{\bfseries\color{DeepBlue}解答}\par\smallskip\centering

\begin{tikzpicture}[>=Latex,font=\small,scale=.85,every node/.style={fill=none}]

\draw[->] (0,-1.4)--(0,1.7) node[above]{$z$};\draw[->] (-1.5,0)--(1.7,0) node[right]{$x$};\draw[->] (-.9,-.7)--(.9,.7) node[above right]{$y$};\draw[dashed] (0,0) circle (1.1);\fill[DeepBlue] (0,1.1) circle (2pt) node[left]{$|\uparrow\rangle$};\fill[DeepBlue] (0,-1.1) circle (2pt) node[left]{$|\downarrow\rangle$};
\draw[DeepBlue,->,thick] (0,0)--(1,1.3) node[right]{$\boldsymbol n$};
\end{tikzpicture}
\figcap{解图：单位方向 $\boldsymbol n$ 决定自旋算符指数的轴，指数参数为 $\theta$。}\label{advA6:quality:A6-quantum-mechanics:a12}
\end{minipage}\end{center}

利用 Pauli 乘法，
\[
 (\vect n\cdot\vect\sigma)^2
 =n_in_j(\delta_{ij}I+\ii\varepsilon_{ijk}\sigma_k)
 =|\vect n|^2I=I.
\]
指数级数的偶次幂与奇次幂分别求和：
\[
\begin{aligned}
 U&=\sum_{m=0}^\infty
 \frac{(\ii\theta)^{2m}}{(2m)!}I
 +\sum_{m=0}^\infty
 \frac{(\ii\theta)^{2m+1}}{(2m+1)!}
 \vect n\cdot\vect\sigma\\
 &=\boxed{\cos\theta\,I+
 \ii\sin\theta\,\vect n\cdot\vect\sigma}.
\end{aligned}
\]
它满足 $U^\dagger U=I$、$\det U=1$。
物理空间把自旋旋转角 $\varphi$ 时，旋转算符为
\[
 U_{\rm rot}(\varphi)=
 \exp\!\left(-\frac{\ii\varphi}{2}
 \vect n\cdot\vect\sigma\right).
\]
因半角出现，自旋子旋转 $2\pi$ 后变号，旋转 $4\pi$ 才回到原态矢；
但全局负号不改变单次态的物理射线。

图~\ref{advA6:quality:A6-quantum-mechanics:a12} 中，沿 $\boldsymbol n$ 的两个本征态分别获得 $e^{\pm\mathrm i\theta}$ 相位；相对相位变为 $2\theta$，对 $U=e^{+\mathrm i\theta\boldsymbol n\cdot\boldsymbol\sigma}$，Bloch向量的主动转角为 $-2\theta$。
\end{advsolution}

\subsubsection*{Bloch 矢量上的旋转及可观测检验}
任意纯自旋态的密度矩阵都可写成
\[
 \rho=|\psi\rangle\langle\psi|
 =\frac12\bigl(I+\vect r\cdot\vect\sigma\bigr),
 \qquad |\vect r|=1.
\]
混态仍具有同一形式，只是 $|\vect r|<1$。对空间转角 $\varphi$，令
$U=\exp[-\ii\varphi\vect n\cdot\vect\sigma/2]$，则
\[
 \rho'=U\rho U^\dagger
 =\frac12\bigl(I+\vect r'\cdot\vect\sigma\bigr),
\]
其中直接展开 Pauli 乘积可得 Rodrigues 公式
\[
 \boxed{\vect r'
 =\vect r\cos\varphi
 +(\vect n\times\vect r)\sin\varphi
 +\vect n(\vect n\cdot\vect r)(1-\cos\varphi)}.
\]
这说明自旋态矢使用半角旋转，而所有可观测的 Bloch 矢量仍按真实空间角
$\varphi$ 旋转。沿单位方向 $\vect a$ 测量自旋，得到正结果的概率为
\[
 P_{\vect a,+}
 =\operatorname{tr}\!\left[
 \rho\,\frac{I+\vect a\cdot\vect\sigma}{2}\right]
 =\frac{1+\vect a\cdot\vect r}{2}.
\]
因此 $2\pi$ 旋转产生的整体负号在单一路径概率中消失；若把旋转过的态与
未旋转的参考路径相干叠加，相对负号则可转化为干涉条纹的半周期移动。

% AdvancedEMQuality20261023
\section{相对论时空与量子现象综合题}
\begingroup\tcbset{breakable/.default=false}

\begingroup\tcbset{breakable/.default=false}
\begin{advexamplebox}{闭合轨道上的固有时与局部同步}\label{advA6:ex:closed-clock}
静质量为 $m$ 的相对论性质点在外力下沿平面内一条光滑简单闭合曲线作周期运动。
沿轨道走过的弧长为 $s$，动量大小为连续函数 $p(s)$；从O出发，$s=0$，
走到 $s=S$ 时返回O，$S$ 是周长，$A$ 是曲线围成的面积。
\begin{enumerate}
\item 用积分表示质点携带的钟在一周期内的走时 $\tau_1$。
\item 大量全同质点先后从O出发，作完全相同的运动。沿从 $s=0$ 到 $s=S$ 的开线段校准各钟：
每个质点的局部惯性观察者认为足够近的前后邻钟与自己同步。求O两侧的钟读数差的大小 $\tau_2$，用积分表示。
\item 证明 $\tau_1\tau_2\ge4\pi A/c^2$，并讨论等号条件。
\end{enumerate}
\begin{center}
% Source: spring-II physical page193, lecture10 problem1.
% Contract: simple closed curve, O=start/end at phi=270deg; CCW tangent p at phi=345deg.
% Curve and tangent use x=1.7cos(t)+.45cos(2t),y=1.7sin(t); no hand-guessed arrow.
\begin{tikzpicture}[>=Latex,font=\small]
 \draw[thick,samples=161,domain=0:360,smooth,variable=\t]
  plot ({1.7*cos(\t)+.45*cos(2*\t)},{1.7*sin(\t)});
 \coordinate (O) at (-.45,-1.7);
 \coordinate (M) at ({1.7*cos(345)+.45*cos(690)},{1.7*sin(345)});
 \fill (O) circle(2pt) (M) circle(2pt);
 \node[below left=5pt] at (O){$O:\ s=0,S$};
 \node[below right=4pt] at (M){$m$};
 \draw[DeepBlue,very thick,->] (M)--++({.65*(-1.7*sin(345)-.9*sin(690))},{.65*1.7*cos(345)})
  node[above right=4pt]{$\vect p(s)$};
 \draw[DeepBlue,thick,->,samples=31,domain=283:317,smooth,variable=\t]
  plot ({1.7*cos(\t)+.45*cos(2*\t)},{1.7*sin(\t)});
 \node at (1,-1.95){$s$ 增大};
 \node at (0,.2){面积 $A$};
\end{tikzpicture}

\figcap{题图：闭合轨道、切口O与随轨道增大的弧长 $s$，原TikZ图保留。}\label{advA6:quality:A6-source-exercises:q1}

\par\smallskip
O既是弧长零点也是同步开线段的接缝；动量始终沿轨道切向。
\end{center}
\end{advexamplebox}
\endgroup\par\nopagebreak[3]
\global\AdvHideHeadingtrue
\begin{advsolution}
\begin{center}\begin{minipage}{.96\linewidth}\centering
% EMAdvFigureHeading
\raggedright{\bfseries\color{DeepBlue}解答}\par\smallskip\centering

\begin{tikzpicture}[>=Latex,font=\small]
 \draw[thick,samples=161,domain=0:360,smooth,variable=\t]
  plot ({1.7*cos(\t)+.45*cos(2*\t)},{1.7*sin(\t)});
 \coordinate (O) at (-.45,-1.7);
 \coordinate (M) at ({1.7*cos(345)+.45*cos(690)},{1.7*sin(345)});
 \fill (O) circle(2pt) (M) circle(2pt);
 \node[below left=5pt] at (O){$O:\ s=0,S$};
 \node[below right=4pt] at (M){$m$};
 \draw[DeepBlue,very thick,->] (M)--++({.65*(-1.7*sin(345)-.9*sin(690))},{.65*1.7*cos(345)})
  node[above right=4pt]{$\vect p(s)$};
 \draw[DeepBlue,thick,->,samples=31,domain=283:317,smooth,variable=\t]
  plot ({1.7*cos(\t)+.45*cos(2*\t)},{1.7*sin(\t)});
 \node at (1,-1.95){$s$ 增大};
 \node at (0,.2){面积 $A$};
\end{tikzpicture}
\figcap{解图：闭合轨道、切口O与随轨道增大的弧长 $s$，原TikZ图保留。}\label{advA6:quality:A6-source-exercises:a1}
\end{minipage}\end{center}

（1）图中 $s$ 按粒子前进方向增大，令速率 $u(s)>0$、
$\gamma(s)=1/\sqrt{1-u(s)^2/c^2}$。在微小弧段内，
\[
 \dd t=\frac{\dd s}{u},\qquad
 \dd\tau=\frac{\dd t}{\gamma}=\frac{\dd s}{\gamma u}
 =\frac{m\,\dd s}{p(s)}.
\]
所以
\[
 \boxed{\tau_1=m\int_0^S\frac{\dd s}{p(s)}}.
\]
若某个孤立点速率为零，应把上式理解为相应的广义积分；题称有限周期，要求该积分收敛。

（2）必须区分“某一时刻按局部同时性校准”与“以后所有时刻一直同步”。
先按前一种含义求校准后的接缝差。令 $C(t,s)$ 表示实验室事件 $(t,s)$ 处的钟读数。
钟随粒子移动，故其固有走时率为
\[
 C_t+u C_s=1/\gamma.
\]
局部惯性系的相邻同时事件满足 $\dd t-u\dd s/c^2=0$。
让这两个事件的钟读数相同，即 $\dd C=C_t\dd t+C_s\dd s=0$，得到
\[
 C_s+\frac u{c^2}C_t=0.
\]
联立两式，而不是把实验室同时误作局部同时：
\[
 C_t=\gamma,\qquad C_s=-\frac{\gamma u}{c^2}=-\frac{p(s)}{mc^2}.
\]
因此沿同一实验室校准时刻绕行，从 $0^+$ 到 $S^-$ 的钟读数逐渐减小，
\[
 C(t,S^-)-C(t,0^+)=-\frac1{mc^2}\int_0^S p(s)\,\dd s,
 \qquad
 \boxed{\tau_2=\frac1{mc^2}\int_0^S p(s)\,\dd s}.
\]
图中O是同一个位置，但把闭合线切开实施局部同步后，接缝两侧保留这个读数跳跃。
$\tau_2$ 是其绝对值；若交换两侧相减的顺序，符号随之改变。

如果第二问被理解为自由走时的钟必须在所有时刻永久保持上述局部同步，条件更强。
这时两式须在一段时间内同时成立，连续混合偏导要求
\[
 \partial_s C_t=\frac{\dd\gamma}{\dd s}
 =\partial_t C_s=0.
\]
也就是速率和动量大小必须沿轨道恒定。任意变化的 $p(s)$ 只能在选定时刻作此校准，
不能靠一次调整钟读数就永远同步。即使匀速，闭合一周仍有接缝跳跃，不能建立全局单值的同步读数；
自由运行的钟群所携带的切口随钟群移动，若坚持把切口固定在O，过缝的钟便需要重新校准。
下面的不等式对瞬时校准的任意正 $p(s)$ 成立，匀速的局部持续同步分支也包含在内。

（3）代入两个积分，质量约去。柯西不等式给出
\[
 \tau_1\tau_2
 =\frac1{c^2}\left(\int_0^S\frac{\dd s}{p}\right)
 \left(\int_0^S p\,\dd s\right)
 \ge\frac1{c^2}\left(\int_0^S1\,\dd s\right)^2
 =\frac{S^2}{c^2}.
\]
等号要求 $p(s)$ 恒定。再用平面简单闭合曲线的等周不等式 $S^2\ge4\pi A$，
即得题求结论。这里说明第二个不等式的依据：选曲线弧长平均位置为原点，
记 $\vect r(s)=(x(s),y(s))$，则 $\int\vect r\dd s=0$、$|\vect r'|=1$。
格林公式及柯西不等式给出
\[
 2A=\left|\int_0^S(xy'-yx')\dd s\right|
 \le\left(\int_0^S|\vect r|^2\dd s\right)^{1/2}S^{1/2}.
\]
把均值为零的周期函数 $x,y$ 展成正弦、余弦级数，每个非零模的波数至少为 $2\pi/S$，所以
\[
 \int_0^S|\vect r|^2\dd s
 \le\left(\frac S{2\pi}\right)^2\int_0^S|\vect r'|^2\dd s
 =\frac{S^3}{4\pi^2}.
\]
合并便有 $2A\le S^2/(2\pi)$。两步均取等号要求曲线为圆。
因此最终等号恰在沿圆周匀速运动时成立。
检验匀速圆：$\tau_1=2\pi R/(\gamma u)$、$\tau_2=2\pi R\gamma u/c^2$，
乘积为 $4\pi^2R^2/c^2=4\pi A/c^2$；两边的量纲都是时间平方。

图~\ref{advA6:quality:A6-source-exercises:a1} 中，切口O两侧是同一个空间位置的不同同步支路；每段固有走时用 $m\,\mathrm ds/p$，同一实验室校准时刻的读数梯度用 $-p/(mc^2)$，两种积分不能混为一个时间差。
\end{advsolution}

\begin{advexamplebox}{核辐射的反冲与运动共振}\label{advA6:ex:nuclear-resonance}
晶体(a)冷却到内部原子的热运动可以忽略，其中一个原子核 $N$ 向探测器方向发射光子，
自身向相反方向反冲，反冲速率远小于光速。光子在另一晶体(b)中的同种原子核 $N'$ 处被吸收。
晶体(b)相对(a)以速率 $v$ 接近。两个核发生同一种能级跃迁，但环境不同：
发射核的跃迁能为 $hc\sigma$，吸收核的跃迁能为 $hc\sigma'$；$\sigma,\sigma'$ 是相应辐射的波数，
这里“波数”采用波长倒数 $1/\lambda$，不是角波数 $2\pi/\lambda$。
已知原子核质量为 $M$，用 $v,\sigma,M$ 表示发生共振吸收所需的 $\sigma'$，并说明近似条件。
\begin{center}
% Source: spring-II physical p200, lecture11 problem1, two source panels.
% Contract: +x right; photon right; emitter recoils left; absorber approaches left.
% Crystal outlines are physical objects, not backgrounds masking other lines.
\begin{tikzpicture}[>=Latex,font=\small]
 \foreach \x in {0,4.2}{\draw[TextGray] (\x,0) rectangle ++(3.4,2);}
 \fill (.9,1) circle(4pt) (6.7,1) circle(4pt);
 \node[below=7pt] at (.9,1){$N$};
 \node[below=7pt] at (6.7,1){$N'$};
 \draw[->,thick] (.72,1)--(.22,1) node[above=5pt,anchor=south west]{反冲};
 \draw[->,thick,DeepBlue] (1.55,1)--(2.95,1) node[midway,above=6pt]{$\gamma$};
 \draw[->,thick,DeepBlue] (4.65,1)--(5.95,1) node[midway,above=6pt]{$\gamma$};
 \draw[->,thick] (6.52,1)--(6.07,1);
 \draw[->,thick] (7.45,2.45)--(6.2,2.45) node[midway,above=3pt]{$v$：接近};
 \node at (1.7,-.45){(a) 发射；跃迁能 $hc\sigma$};
 \node at (5.9,-.45){(b) 吸收；跃迁能 $hc\sigma'$};
 \draw[->] (3.05,2.5)--(3.85,2.5) node[right]{$x$};
\end{tikzpicture}

\figcap{题图：发射晶体、反冲核和迎面运动的吸收晶体，原TikZ图保留。}\label{advA6:quality:A6-source-exercises:q2}

\end{center}
\end{advexamplebox}
\global\AdvHideHeadingtrue
\begin{advsolution}
\begin{center}\begin{minipage}{.96\linewidth}\centering
% EMAdvFigureHeading
\raggedright{\bfseries\color{DeepBlue}解答}\par\smallskip\centering

\begin{tikzpicture}[>=Latex,font=\small]
 \foreach \x in {0,4.2}{\draw[TextGray] (\x,0) rectangle ++(3.4,2);}
 \fill (.9,1) circle(4pt) (6.7,1) circle(4pt);
 \node[below=7pt] at (.9,1){$N$};
 \node[below=7pt] at (6.7,1){$N'$};
 \draw[->,thick] (.72,1)--(.22,1) node[above=5pt,anchor=south west]{反冲};
 \draw[->,thick,DeepBlue] (1.55,1)--(2.95,1) node[midway,above=6pt]{$\gamma$};
 \draw[->,thick,DeepBlue] (4.65,1)--(5.95,1) node[midway,above=6pt]{$\gamma$};
 \draw[->,thick] (6.52,1)--(6.07,1);
 \draw[->,thick] (7.45,2.45)--(6.2,2.45) node[midway,above=3pt]{$v$：接近};
 \node at (1.7,-.45){(a) 发射；跃迁能 $hc\sigma$};
 \node at (5.9,-.45){(b) 吸收；跃迁能 $hc\sigma'$};
 \draw[->] (3.05,2.5)--(3.85,2.5) node[right]{$x$};
\end{tikzpicture}
\figcap{解图：发射晶体、反冲核和迎面运动的吸收晶体，原TikZ图保留。}\label{advA6:quality:A6-source-exercises:a2}
\end{minipage}\end{center}

图(a)中光子沿 $+x$ 传播，反冲核沿 $-x$ 运动；图(b)中的吸收核在吸收前沿 $-x$ 接近光源。
不能把跃迁能直接当成实际发射的光子能量，因为一部分能量变成了反冲动能。
先按照题设的单核反冲模型计算，记 $\Delta=hc\sigma$、$\Delta'=hc\sigma'$。

设光源静止系中的光子能量为 $\varepsilon$。动量守恒给核的反冲动量大小 $p_N=\varepsilon/c$。
由于反冲速度很小，可用 $p_N^2/(2M)$ 表示反冲动能。因此
\[
 \Delta=\varepsilon+\frac{\varepsilon^2}{2Mc^2}.
\]
令 $r=\Delta/(Mc^2)\ll1$。反冲项本身已是一阶小修正，故在该项中用 $\varepsilon\simeq\Delta$，得
\[
 \varepsilon=\Delta-\frac{\Delta^2}{2Mc^2}+O(r^2\Delta).
\]
在吸收核的吸收前静止系中，迎面而来的光子能量被多普勒增大为
\[
 \varepsilon_a=\gamma_v(1+v/c)\varepsilon,
 \qquad \gamma_v=(1-v^2/c^2)^{-1/2}.
\]
吸收后核获得光子动量，也要付出反冲动能。这次能量守恒为
\[
 \varepsilon_a=\Delta'+\frac{\varepsilon_a^2}{2Mc^2},
 \quad\text{即}\quad
 \Delta'=\varepsilon_a-\frac{\varepsilon_a^2}{2Mc^2}.
\]
若晶体的接近速度不取低速近似，记 $D=\gamma_v(1+v/c)$，则在
$r\ll1$ 且 $Dr\ll1$、发射和吸收反冲均非相对论的条件下，可先得到
\[
 \boxed{\sigma'\simeq D\sigma-\frac{h\sigma^2}{2Mc}(D+D^2)}.
\]
这里 $D$ 乘发射反冲修正，而吸收反冲修正含 $D^2$，两者在一般速度下不相等。
常见的小速度调谐再取 $\beta=v/c\ll1$。只保留 $\beta$ 和 $r$ 的一次项，舍去 $\beta^2,r^2,\beta r$：
\begin{align*}
 \Delta'
 &\simeq(1+\beta)\left(\Delta-\frac{\Delta^2}{2Mc^2}\right)
     -\frac{\Delta^2}{2Mc^2}\\
 &\simeq\Delta(1+\beta)-\frac{\Delta^2}{Mc^2}.
\end{align*}
除以 $hc$，得到
\[
 \boxed{\sigma'\simeq\sigma\left(1+\frac vc\right)-\frac{h\sigma^2}{Mc}}.
\]
最后一项包含发射和吸收两次反冲，不能只扣一次。$h\sigma^2/(Mc)$ 的量纲是长度的倒数，
与 $\sigma$ 一致。若 $M\to\infty$，反冲项消失，只剩迎面运动的多普勒蓝移；
若两环境完全相同而 $\sigma'=\sigma$，则一阶近似要求 $v=h\sigma/M$，即用蓝移补偿两次反冲损失。

还要区分本题模型与穆斯堡尔效应中的无反冲谱线。晶体中的核有一定概率不激发晶格振动，
此时动量由整个晶体承担，反冲有效质量远大于单个原子核的 $M$，上面的单核反冲修正便显著减小。
不能在同一步计算中既采用单核质量，又宣称反冲完全由晶体承担。

图~\ref{advA6:quality:A6-source-exercises:a2} 中，光子沿探测方向传播、发射核向反方向反冲，吸收晶体的迎面速度产生蓝移；发射与吸收各自的反冲能都需进入共振条件。
\end{advsolution}

\begin{advexamplebox}{两光子散射与出射波矢的取值范围}\label{advA6:ex:photon-scattering}
两光子可经四个电磁相互作用顶点发生散射，图(a)给出一幅电子闭合内线图：
两条入射光子外线对应 $\nu_1,\nu_2$，两条出射外线对应 $\nu_3,\nu_4$。
中间电子用虚粒子传播子描述，不能预先认为它们都是满足
$E^2=p^2c^2+m_e^2c^4$、能量为正的自由实粒子。
\begin{enumerate}
\item 证明不能把图中过程解释为一个自由光子先变成一对真实电子、正电子，再依次吸收或发射光子。
说明能量、动量守恒与中间电子为虚粒子之间的关系。
\item 图(b)中两入射光子的频率均为 $\nu_0$，入射方向夹角为 $120^\circ$。
求两出射光子的波矢在三维波矢空间中的取值范围、频率与出射方向的关系，以及最大和最小频率。
\end{enumerate}
\begin{center}
% Source: spring-II physical p201, lecture11 problem2, both source panels.
% Contract: closed fermion loop with four external photon legs; arrows show
% photon propagation. Right panel: equal incident vectors enclose 120deg,
% sum along +y; representative output momenta sum to the same vector.
\begin{tikzpicture}[>=Latex,font=\small]
 \begin{scope}[xshift=-3.0cm]
  \draw[thick,->] (-.55,-.55)--(-.55,.55);
  \draw[thick,->] (-.55,.55)--(.55,.55);
  \draw[thick,->] (.55,.55)--(.55,-.55);
  \draw[thick,->] (.55,-.55)--(-.55,-.55);
  \foreach \sx/\sy/\lab in {-1/1/1,-1/-1/2,1/1/3,1/-1/4}{
   \draw[DeepBlue,thick,domain=0:1,samples=70,variable=\t]
    plot ({\sx*(.55+1.0*\t)+.045*sin(\t*2160)},
          {\sy*(.55+.75*\t)-\sx*\sy*.06*sin(\t*2160)});
   \node[anchor=\ifnum\sx<0 east\else west\fi] at ({\sx*1.65},{\sy*1.32}){$\nu_{\lab}$};
  }
  \node at (0,0){$e$};
 \node[align=center] at (0,-1.95){(a) 四条光子外线\\电子闭合内线};
 \end{scope}
 \begin{scope}[xshift=3.1cm]
  \coordinate (O) at (0,0);
  \draw[->,thick,DeepBlue] ({-1.8*cos(30)},-.9)--(O);
  \draw[->,thick,DeepBlue] ({1.8*cos(30)},-.9)--(O);
  \node at (-1.72,-1.18){$\nu_0$};
  \node at (1.72,-1.18){$\nu_0$};
  \draw[->,thick] (O)--(0,1.8) node[above]{$\vect K=\vect k_1+\vect k_2$};
  % theta=40deg, k3/k0=3/[2(2-cos(theta))]; k4=K-k3.
  \pgfmathsetmacro{\kk}{3/(2*(2-cos(40)))}
  \draw[->,thick,DeepBlue] (O)--({1.8*\kk*sin(40)},{1.8*\kk*cos(40)}) node[right]{$\vect k_3$};
  \draw[->,thick,DeepBlue] (O)--({-1.8*\kk*sin(40)},{1.8*(1-\kk*cos(40))}) node[above left]{$\vect k_4$};
  \draw[->] (0,.6) arc[start angle=90,end angle=50,radius=.6];
  \node at (.29,.87){$\theta$};
  \draw (-.43,-.25) arc[start angle=210,end angle=330,radius=.5];
  \node at (0,-.78){$120^\circ$};
  \node at (0,-1.95){(b) 入射夹角与出射方向};
 \end{scope}
\end{tikzpicture}

\figcap{题图：两入射与两出射光子外线及题给角度，内部电子线表示虚传播子。}\label{advA6:quality:A6-source-exercises:q3}

\end{center}
\end{advexamplebox}
\begin{advsolution}
（1）一个自由光子的四动量平方为零。若它能在真空中单独变成一对正能量的真实电子和正电子，
则末态总四动量平方至少为 $(2m_ec)^2>0$：在二者的质心系，
总能量至少为 $2m_ec^2$，总动量为零。这与初态四动量平方为零矛盾。
也可以直接用三维不等式：
\[
 E_++E_->c(|\vect p_+|+|\vect p_-|)
 \ge c|\vect p_++\vect p_-|,
\]
而初态光子恰有 $E_\gamma=c|\vect p_\gamma|$，不能同时满足总能量和总动量守恒。

同样，一个自由实电子不能在保持同一静质量的同时单独吸收一个非零能量的实光子：
若电子初四动量为 $P$、光子为 $k$，假定吸收后的电子也在质量壳上，则
\[
 (P+k)^2=P^2\quad\Rightarrow\quad2P\mathbin{\cdot}k=0.
\]
但在初电子静止系，$P\mathbin{\cdot}k=m_e E_\gamma>0$，仍然矛盾。
因此图(a)不能解释为一系列经典的、可逐个观测的自由电子中间态。
闭合内线必须作为虚粒子传播子处理；每个顶点仍严格满足四动量守恒，
并不是“虚粒子可以暂时不守恒能量”。“虚”指它不是可作为初末态观测的自由在壳粒子。
这项论证排除把全部中间过程解释成实粒子的图景，并不要求在环动量积分的每个积分点，
每一条内线都绝不经过质量壳上的数值；不能把图中线段当成电子实际走过的空间轨迹。

（2）以下波矢采用角波数，$k_0=2\pi\nu_0/c$，光子能量为 $\hbar ck_0$。
图(b)的合波矢 $\vect K=\vect k_1+\vect k_2$ 沿两入射方向的角平分线，且
\[
 K^2=k_0^2+k_0^2+2k_0^2\cos120^\circ=k_0^2.
\]
由能量、动量守恒，两个出射波矢满足
\[
 \vect k_3+\vect k_4=\vect K,\qquad k_3+k_4=2k_0.
\]
把 $\vect k_3$ 的端点记为P，则 $OP=k_3$、从P到 $\vect K$ 端点的距离为 $k_4$。
所以P到两个定点O、$\vect K$ 的距离之和恒为 $2k_0$。
这不是椭球内部的任意点，而是旋转椭球的\emph{表面}；两个焦点相距 $k_0$，
中心为 $\vect K/2$，长半轴为 $k_0$，短半轴为 $\sqrt3k_0/2$。
下面画其包含 $\vect K$ 的一个截面；绕水平轴旋转即为完整允许面。
\begin{center}
% Deterministic momentum-space ellipse, meridional section of prolate ellipsoid.
% Foci z=0,k0; center k0/2; semimajor=k0, semiminor=sqrt(3)k0/2.
\begin{tikzpicture}[>=Latex,font=\small,x=2cm,y=2cm]
 \draw[->,TextGray] (-.85,0)--(1.85,0) node[right]{$k_{\parallel}/k_0$};
 \draw[->,TextGray] (0,-1.08)--(0,1.08) node[above]{$k_{\perp}/k_0$};
 \draw[DeepBlue,thick] (.5,0) ellipse[x radius=1,y radius=.8660254];
 \coordinate (P) at ({.5+cos(65)},{.8660254*sin(65)});
 \fill (0,0) circle(.023) (1,0) circle(.023) (P) circle(.023);
 \draw[->,thick] (0,0)--(P);
 \draw[->,thick] (P)--(1,0);
 \node[below left=4pt] at (0,0){$O$};
 \node[below=5pt] at (1,0){$\vect K/k_0$};
 \node[above right=4pt] at (P){$P$};
 \node at (.37,.61){$\vect k_3/k_0$};
 \node[anchor=west] at (1.6,.55){$\vect k_4/k_0$};
 \draw[TextGray] (1.55,.55)--(.98,.4);
 \draw[TextGray](-.5,-.04)--(-.5,.04);\node[below left=3pt] at (-.5,-.04){$-\tfrac12$};
 \draw[TextGray](1.5,-.04)--(1.5,.04);\node[below right=3pt] at (1.5,-.04){$\tfrac32$};
\end{tikzpicture}

\figcap{解图：出射波矢端点允许面是旋转椭球面；此图为含总波矢 $\boldsymbol K$ 轴的截面，O与第二焦点的间距为 $k_0$。}\label{advA6:quality:A6-source-exercises:a3}

\end{center}
图中的O到P代表 $\vect k_3$，P到第二焦点代表 $\vect k_4$，两矢量首尾相接恰等于 $\vect K$。
用 $k_\parallel$ 和 $k_\perp$ 表示任一出射波矢平行、垂直于 $\vect K$ 的分量，允许面满足
\[
 \frac{(k_\parallel-k_0/2)^2}{k_0^2}
 +\frac{k_\perp^2}{3k_0^2/4}=1.
\]
再求方向与频率。令 $\theta$ 为 $\vect k_3$ 与 $\vect K$ 的夹角，如题图(b)，并令 $k=k_3$。
另一光子的波数必须同时满足
\[
 k_4=2k_0-k,\qquad
 k_4^2=|\vect K-\vect k_3|^2=k_0^2+k^2-2k_0k\cos\theta.
\]
代入、平方展开并消去 $k^2$：
\[
 4k_0^2-4k_0k=k_0^2-2k_0k\cos\theta,
 \qquad
 \boxed{\nu(\theta)=\frac{3\nu_0}{2(2-\cos\theta)}}.
\]
两光子全同，交换下标后这一关系对另一光子也成立。由于 $-1\le\cos\theta\le1$，
所有方向都允许，分母恒正；频率范围是
\[
 \boxed{\frac{\nu_0}{2}\le\nu\le\frac{3\nu_0}{2}}.
\]
最大值沿合动量方向出现，最小值沿其反方向出现。极值时两光子反向传播，频率分别为
$3\nu_0/2$ 和 $\nu_0/2$：频率和为 $2\nu_0$，动量差为 $\hbar k_0$，正好复核两条守恒式。
允许面给出运动学范围，不意味着各方向的散射概率相同；概率还取决于量子散射振幅。

图~\ref{advA6:quality:A6-source-exercises:a3} 中，O到P及P到第二焦点的两矢量首尾相接为 $\boldsymbol K$，两段长度之和为 $2k_0$；这同时表达动量守恒与能量守恒，所以允许集合是椭球表面而非内部区域。
\end{advsolution}

\begin{advexamplebox}{有限质量氢靶的卢瑟福散射}\label{advA6:ex:rutherford-recoil}
动能为 $K=2\,\mathrm{MeV}$ 的 $\alpha$ 粒子入射到静止氢核上。
采用非相对论弹性库仑散射模型，取 $m_\alpha=4m_p=4m$。
图中 $\theta_L,\varphi_L$ 分别表示实验室系中 $\alpha$ 粒子的散射角、氢核的反冲角，
$\theta_c,\varphi_c$ 表示对应的质心系角。求实验室系中反冲氢核沿
$\varphi_L=30^\circ$ 方向的微分散射截面及其反冲动能。
\begin{center}
% Source: spring-II p219 lecture14 problem4. Velocity construction, not positions.
% m_alpha=4m_p; V_cm=.8v0; theta_cm=120deg. Scale 4cm/v0.
\begin{tikzpicture}[>=Latex,font=\small]
 \coordinate (O) at (0,0);\coordinate (C) at (3.2,0);
 \coordinate (A) at (2.8,{.8*sin(120)});
 \coordinate (P) at (4.8,{-3.2*sin(120)});
 \draw[->,TextGray] (-1.3,0)--(5.15,0);
 \draw[->,thick] (-1.25,0)--(-.2,0);
 \node[above=6pt] at (-.85,0){$\alpha$ 入射};
 \draw[->,thick,DeepBlue] (O)--(A);
 \draw[->,thick,DeepBlue] (O)--(P);
 \draw[dashed,->,thick] (C)--(A);
 \draw[dashed,->,thick] (C)--(P);
 \fill (O) circle(2pt) (C) circle(2pt) (A) circle(2pt) (P) circle(2pt);
 \node[below left=6pt] at (O){$O$};
 \node[below left=5pt] at (C){$C$};
 \node[above left=4pt] at (A){$\alpha$};
 \node[below right=4pt] at (P){$p$};
 \draw (.85,0) arc[start angle=0,end angle=13.8979,radius=.85];
 \node at (.8,.4){$\theta_L$};
 \draw (1.18,0) arc[start angle=0,end angle=-30,radius=1.18];
 \node at (1.45,-.33){$\varphi_L$};
 \draw (3.67,0) arc[start angle=0,end angle=120,radius=.47];
 \node at (3.56,.78){$\theta_c$};
 \draw (3.9,0) arc[start angle=0,end angle=-60,radius=.7];
 \node at (4.18,-.47){$\varphi_c$};
 \node[align=left,anchor=west] at (-1.25,-2.25){实线：实验室速度\\虚线：质心系速度};
\end{tikzpicture}

\figcap{题图：有限质量靶的散射几何与两个参考系的角度，原TikZ图保留。}\label{advA6:quality:A6-source-exercises:q4}

\end{center}
图为同一比例尺的速度矢量构图；C表示质心平移速度的端点，不是碰撞时第三个粒子的位置。
\end{advexamplebox}
\global\AdvHideHeadingtrue
\begin{advsolution}
\begin{center}\begin{minipage}{.96\linewidth}\centering
% EMAdvFigureHeading
\raggedright{\bfseries\color{DeepBlue}解答}\par\smallskip\centering

\begin{tikzpicture}[>=Latex,font=\small]
 \coordinate (O) at (0,0);\coordinate (C) at (3.2,0);
 \coordinate (A) at (2.8,{.8*sin(120)});
 \coordinate (P) at (4.8,{-3.2*sin(120)});
 \draw[->,TextGray] (-1.3,0)--(5.15,0);
 \draw[->,thick] (-1.25,0)--(-.2,0);
 \node[above=6pt] at (-.85,0){$\alpha$ 入射};
 \draw[->,thick,DeepBlue] (O)--(A);
 \draw[->,thick,DeepBlue] (O)--(P);
 \draw[dashed,->,thick] (C)--(A);
 \draw[dashed,->,thick] (C)--(P);
 \fill (O) circle(2pt) (C) circle(2pt) (A) circle(2pt) (P) circle(2pt);
 \node[below left=6pt] at (O){$O$};
 \node[below left=5pt] at (C){$C$};
 \node[above left=4pt] at (A){$\alpha$};
 \node[below right=4pt] at (P){$p$};
 \draw (.85,0) arc[start angle=0,end angle=13.8979,radius=.85];
 \node at (.8,.4){$\theta_L$};
 \draw (1.18,0) arc[start angle=0,end angle=-30,radius=1.18];
 \node at (1.45,-.33){$\varphi_L$};
 \draw (3.67,0) arc[start angle=0,end angle=120,radius=.47];
 \node at (3.56,.78){$\theta_c$};
 \draw (3.9,0) arc[start angle=0,end angle=-60,radius=.7];
 \node at (4.18,-.47){$\varphi_c$};
 \node[align=left,anchor=west] at (-1.25,-2.25){实线：实验室速度\\虚线：质心系速度};
\end{tikzpicture}
\figcap{解图：有限质量靶的散射几何与两个参考系的角度，原TikZ图保留。}\label{advA6:quality:A6-source-exercises:a4}
\end{minipage}\end{center}

设入射速率为 $v_0$，入射方向为 $+x$。质心速度和相对运动约化质量分别为
\[
 V=\frac{4m}{5m}v_0=\frac45v_0,\qquad
 \mu=\frac{4m\cdot m}{5m}=\frac45m.
\]
相对运动初速度仍为 $v_0$，所以用于卢瑟福公式的能量不是 $K$，而是
\[
 E_{\rm rel}=\frac12\mu v_0^2=\frac K5=0.4\,\mathrm{MeV}.
\]
质心系中 $\alpha$ 粒子的初末速率均为 $v_0/5$，氢核的初末速率均为 $4v_0/5$。
取 $\alpha$ 出射到 $x$ 轴上方，则它的质心系速度为
$(v_0/5)(\cos\theta_c,\sin\theta_c)$，氢核速度与之相反、大小为其四倍。
加上平移速度 $V\vect e_x$ 后，氢核在实验室系的速度为
\[
 \vect u_p=\frac{4v_0}5(1-\cos\theta_c,-\sin\theta_c).
\]
题图中的 $OC$ 是共同加上的平移速度，$C\alpha$ 与 $Cp$ 是反向的质心系速度，
而 $O\alpha$、$Op$ 为实验室速度。因此反冲角满足
\[
 \tan\varphi_L=\frac{\sin\theta_c}{1-\cos\theta_c}
 =\cot\frac{\theta_c}2,
 \quad \varphi_L=\frac\pi2-\frac{\theta_c}2.
\]
由 $\varphi_L=30^\circ$ 得 $\theta_c=120^\circ$，同时 $\varphi_c=60^\circ$。
图中的 $\alpha$ 实验室散射角则满足
$\tan\theta_L=\sin\theta_c/(4+\cos\theta_c)$，数值约为 $13.9^\circ$，
不能把 $\theta_L$、$\theta_c$ 或 $\varphi_L$ 混作同一个角。

记库仑常数参数 $\kappa=Z_\alpha Z_p e^2/(4\pi\varepsilon_0)$。
相对运动的排斥库仑双曲线给出冲量参数
\[
 b=\frac{\kappa}{2E_{\rm rel}}\cot\frac{\theta_c}2.
\]
下面把它转成截面，明确角度变换而不直接套用固定重靶公式。
轴对称入射中，$b$ 至 $b+\dd b$ 的环形面积为 $2\pi b|\dd b|$，
质心系对应立体角为 $\dd\Omega_c=2\pi\sin\theta_c|\dd\theta_c|$，故
\[
 \frac{\dd\sigma}{\dd\Omega_c}
 =\frac b{\sin\theta_c}\left|\frac{\dd b}{\dd\theta_c}\right|
 =\left(\frac{\kappa}{4E_{\rm rel}}\right)^2
   \frac1{\sin^4(\theta_c/2)}.
\]
探测器接收的是反冲氢核，所需立体角是
$\dd\Omega_L=2\pi\sin\varphi_L|\dd\varphi_L|$。
由 $\theta_c=\pi-2\varphi_L$，有
\[
 \left|\frac{\dd\Omega_c}{\dd\Omega_L}\right|
 =\frac{2\sin(2\varphi_L)}{\sin\varphi_L}=4\cos\varphi_L.
\]
同一批散射事件数不变，故
\[
 \boxed{\frac{\dd\sigma}{\dd\Omega_L}
 =\frac{\kappa^2}{4E_{\rm rel}^2\cos^3\varphi_L}}.
\]
对 $Z_\alpha=2,Z_p=1$，$\kappa\simeq2.880\,\mathrm{MeV\,fm}$，代入得到
\[
 \left.\frac{\dd\sigma}{\dd\Omega_L}\right|_{30^\circ}
 =\frac{(2.880)^2}{4(0.4)^2(\sqrt3/2)^3}\,\mathrm{fm^2/sr}
 \simeq19.95\,\mathrm{fm^2/sr}=0.1995\,\mathrm{b/sr}.
\]
这是\emph{反冲氢核}的微分截面；若观测散射 $\alpha$ 粒子，需另用 $\theta_L$ 的雅可比，
不能沿用这一结果。

反冲速率由速度矢量长度直接得到
\[
 u_p=\frac{8v_0}5\sin\frac{\theta_c}2
 =\frac{8v_0}5\cos\varphi_L,
\]
因此
\[
 \boxed{K_p=\frac12mu_p^2
 =\frac{16}{25}K\cos^2\varphi_L=0.96\,\mathrm{MeV}}.
\]
余下 $\alpha$ 粒子动能为 $1.04\,\mathrm{MeV}$；用其图示速度分量
$v_0(0.7,\sqrt3/10)$ 计算也得到 $K_\alpha=0.52K$，两者之和为 $2\,\mathrm{MeV}$。
当 $\varphi_L\to90^\circ$，反冲动能趋于零，而纯库仑模型的小角散射截面发散；
实际屏蔽会限制该理想模型在极小动量转移处的适用性。

图~\ref{advA6:quality:A6-source-exercises:a4} 中，实验室中的重入射粒子偏转角与动量中心系散射角不同，最大实验室偏转受靶反冲限制；角度分支应与速度矢量三角形一致。
\end{advsolution}

\begin{advexamplebox}{中子重力干涉与两个输出端口}\label{advA6:ex:neutron-gravity}
低能中子经单晶片发生布拉格衍射，可沿特定方向反射或透射。
四段光路构成边长均为 $l$ 的菱形 $OACB$，底边 $OB$ 水平，入射方向沿 $OB$。
晶体的布拉格角为 $\theta$，分束后的两方向夹角为 $2\theta$。
装置可绕平行于底边的轴转过倾角 $\varphi$。已知中子质量为 $m$、入射动能为 $E_k$，
入射束强度为1，求端口1、2的输出强度。
先给出理想等分束、两臂完全反射的结果，再说明分束比或中间晶体反射率不是理想值时怎样计算。
\begin{center}
% Source: spring-II p220 lecture14 problem5. Plan view plus tilt cross-section.
% OA=OB=AC=BC=l. OA makes 2theta to OB; OC makes theta. Rotation about OB.
% All joints come from the same rhombus coordinates. Crystal lines bisect rays.
\begin{tikzpicture}[>=Latex,font=\small]
 \coordinate (O) at (0,0);\coordinate (B) at (3.25,0);
 \coordinate (A) at ({3.25*cos(50)},{3.25*sin(50)});
 \coordinate (C) at ({3.25*(1+cos(50))},{3.25*sin(50)});
 \draw[->,thick,DeepBlue] (-1.2,0)--(O);
 \draw[->,thick,DeepBlue] (O)--(B);
 \draw[->,thick,DeepBlue] (O)--(A);
 \draw[->,thick,DeepBlue] (A)--(C);
 \draw[->,thick,DeepBlue] (B)--(C);
 \draw[->,thick,DeepBlue] (C)--++(1.2,0) node[right]{1};
 \draw[->,thick,DeepBlue] (C)--++({.9*cos(50)},{.9*sin(50)}) node[above]{2};
 \draw[dashed,TextGray] (O)--(C);
 \foreach \name in {O,A,B,C}{
  \draw[very thick,TextGray] ($ (\name)+({-.42*cos(25)},{-.42*sin(25)}) $)
     --($ (\name)+({.42*cos(25)},{.42*sin(25)}) $);
 }
 \node[below=7pt] at (O){$O$};\node[below=7pt] at (B){$B$};
 \node[above=8pt] at (A){$A$};\node[below right=6pt] at (C){$C$};
 \node[below=6pt] at (1.63,0){$l$};
 \node[above=6pt] at ({3.25*cos(50)+1.63},{3.25*sin(50)}){$l$};
 \node[left=6pt] at ({1.63*cos(50)},{1.63*sin(50)}){$l$};
 \node[right=6pt] at ({3.25+1.63*cos(50)},{1.63*sin(50)}){$l$};
 \draw (.9,0) arc[start angle=0,end angle=25,radius=.9];
 \node at (1.15,.23){$\theta$};
 \draw (1.45,0) arc[start angle=0,end angle=50,radius=1.45];
 \node at (1.5,1.35){$2\theta$};
 \node[anchor=west] at (-1.1,3.2){(a) 装置平面内的光路};
 \begin{scope}[shift={(0,-2.1)}]
  \draw[->,TextGray] (0,0)--(4,0) node[right]{水平};
  \draw[thick] (0,0)--(3.2,1.16);
  \draw[dashed] (3.2,0)--(3.2,1.16);
  \draw[<->] (3.65,0)--(3.65,1.16) node[midway,right=5pt]{$h=l\sin2\theta\sin\varphi$};
  \fill (0,0) circle(2pt);
  \node[below=5pt] at (0,0){转轴 $OB$};
  \draw (.9,0) arc[start angle=0,end angle=19.9256,radius=.9];
  \node at (1.17,.17){$\varphi$};
  \node[above,rotate=19.9256] at (1.95,.71){$l\sin2\theta$};
  \node[anchor=west] at (-1.1,-.75){(b) 垂直转轴的截面};
 \end{scope}
\end{tikzpicture}

\figcap{题图：中子菱形干涉仪的四段光路与两个输出端口，原TikZ图保留。}\label{advA6:quality:A6-source-exercises:q5}

\end{center}
\end{advexamplebox}
\global\AdvHideHeadingtrue
\begin{advsolution}
\begin{center}\begin{minipage}{.96\linewidth}\centering
% EMAdvFigureHeading
\raggedright{\bfseries\color{DeepBlue}解答}\par\smallskip\centering

\begin{tikzpicture}[>=Latex,font=\small]
 \coordinate (O) at (0,0);\coordinate (B) at (3.25,0);
 \coordinate (A) at ({3.25*cos(50)},{3.25*sin(50)});
 \coordinate (C) at ({3.25*(1+cos(50))},{3.25*sin(50)});
 \draw[->,thick,DeepBlue] (-1.2,0)--(O);
 \draw[->,thick,DeepBlue] (O)--(B);
 \draw[->,thick,DeepBlue] (O)--(A);
 \draw[->,thick,DeepBlue] (A)--(C);
 \draw[->,thick,DeepBlue] (B)--(C);
 \draw[->,thick,DeepBlue] (C)--++(1.2,0) node[right]{1};
 \draw[->,thick,DeepBlue] (C)--++({.9*cos(50)},{.9*sin(50)}) node[above]{2};
 \draw[dashed,TextGray] (O)--(C);
 \foreach \name in {O,A,B,C}{
  \draw[very thick,TextGray] ($ (\name)+({-.42*cos(25)},{-.42*sin(25)}) $)
     --($ (\name)+({.42*cos(25)},{.42*sin(25)}) $);
 }
 \node[below=7pt] at (O){$O$};\node[below=7pt] at (B){$B$};
 \node[above=8pt] at (A){$A$};\node[below right=6pt] at (C){$C$};
 \node[below=6pt] at (1.63,0){$l$};
 \node[above=6pt] at ({3.25*cos(50)+1.63},{3.25*sin(50)}){$l$};
 \node[left=6pt] at ({1.63*cos(50)},{1.63*sin(50)}){$l$};
 \node[right=6pt] at ({3.25+1.63*cos(50)},{1.63*sin(50)}){$l$};
 \draw (.9,0) arc[start angle=0,end angle=25,radius=.9];
 \node at (1.15,.23){$\theta$};
 \draw (1.45,0) arc[start angle=0,end angle=50,radius=1.45];
 \node at (1.5,1.35){$2\theta$};
 \node[anchor=west] at (-1.1,3.2){(a) 装置平面内的光路};
 \begin{scope}[shift={(0,-2.1)}]
  \draw[->,TextGray] (0,0)--(4,0) node[right]{水平};
  \draw[thick] (0,0)--(3.2,1.16);
  \draw[dashed] (3.2,0)--(3.2,1.16);
  \draw[<->] (3.65,0)--(3.65,1.16) node[midway,right=5pt]{$h=l\sin2\theta\sin\varphi$};
  \fill (0,0) circle(2pt);
  \node[below=5pt] at (0,0){转轴 $OB$};
  \draw (.9,0) arc[start angle=0,end angle=19.9256,radius=.9];
  \node at (1.17,.17){$\varphi$};
  \node[above,rotate=19.9256] at (1.95,.71){$l\sin2\theta$};
  \node[anchor=west] at (-1.1,-.75){(b) 垂直转轴的截面};
 \end{scope}
\end{tikzpicture}
\figcap{解图：中子菱形干涉仪的四段光路与两个输出端口，原TikZ图保留。}\label{advA6:quality:A6-source-exercises:a5}
\end{minipage}\end{center}

几何关系必须先从图(a)确定。$OB$、$AC$ 长为 $l$ 且水平，$OA$、$BC$ 与它们夹角为 $2\theta$，
对角线 $OC$ 才与底边夹角为 $\theta$。不能把整个菱形的顶角错当作布拉格角。
菱形面积为 $A=l^2\sin2\theta$，垂直于转轴的宽度为 $l\sin2\theta$。
图(b)表示绕 $OB$ 倾转后上下两条水平臂的高度差：
\[
 h=l\sin2\theta\sin\varphi.
\]

先求重力引起的相位差。中子非相对论速率为 $v=\sqrt{2E_k/m}$，
假定 $mgh\ll E_k$，忽略重力造成的轨迹偏折，只保留重力相位的一阶修正。
在固定总能量下，离底边高度 $z$ 处的局部动量为
\[
 p(z)=\sqrt{2m(E_k-mgz)}
 \simeq p_0-\frac{mgz}{v},\qquad p_0=mv.
\]
空间传播相位为 $\hbar^{-1}\int p\,\dd s$。
两条斜臂 $OA$、$BC$ 的高度分布完全相同，积分相消；
上方水平臂 $AC$ 比下方 $OB$ 高 $h$，且长度都是 $l$。于是
\[
 \delta\equiv\Phi_{OAC}-\Phi_{OBC}
 =-\frac{mghl}{\hbar v}
 =-\frac{mg l^2\sin2\theta\sin\varphi}{\hbar\sqrt{2E_k/m}}.
\]
高处动量较小，所以上臂相位少积累，负号由路径次序决定。
以下强度只含 $\cos\delta$，交换两臂或反向倾转不改变强度。

输出强度不能只看相位差，还要把晶片的透射、衍射振幅相加。
为把端口对应讲清，取无吸收、对称的第一和末片晶体的振幅系数
$t=\sqrt T$、$r=\mathrm i\sqrt R$，其中 $T+R=1$；这个 $\pi/2$ 相对相位使分束矩阵幺正。
两条中间臂各衍射一次，设其共同反射振幅为 $r_m$，$|r_m|^2=R_m$。
端口1沿 $AC$ 的延长线：上臂经历 $r,r_m,t$，下臂经历 $t,r_m,r$；
端口2沿 $BC$ 的延长线：上臂经历 $r,r_m,r$，下臂经历 $t,r_m,t$。
因此
\begin{align*}
 a_1&=r_mtr\left(e^{\mathrm i\Phi_{OAC}}+e^{\mathrm i\Phi_{OBC}}\right),\\
 a_2&=r_m\left(r^2e^{\mathrm i\Phi_{OAC}}+t^2e^{\mathrm i\Phi_{OBC}}\right),
\end{align*}
从而对单位入射强度，
\[
 \boxed{I_1=2R_mTR(1+\cos\delta)},\qquad
 \boxed{I_2=R_m(T^2+R^2-2TR\cos\delta)}.
\]
理想马赫--曾德尔模型取 $T=R=1/2$、$R_m=1$，得到
\[
 \boxed{I_1=\cos^2\frac\delta2,\qquad I_2=\sin^2\frac\delta2}.
\]
在无其他相移、$\varphi=0$ 时，端口1亮、端口2暗，二者之和为1。
若中间晶片本身也只是反射率 $1/2$ 的普通分束晶片，则 $R_m=1/2$，上述两个强度都须再乘 $1/2$；
其余中子从未进入这两个端口的旁路离开，不能仍宣称 $I_1+I_2=1$。
如果把进入最终重合两臂的中子束强度重新归一为1，才可约去共同因子 $R_m$。

题设没有给出晶片分束比、损耗和附加相位，故不附加理想假设时，
能够确定的是重力相位及上述参数化关系，而不是唯一的绝对强度曲线。
若两臂还有固定相位差 $\delta_0$，把公式中的 $\delta$ 换为 $\delta+\delta_0$。
检验：$\delta$ 的量纲为能量乘时间除以作用量，是无量纲；
$g\to0$ 或 $\varphi\to0$ 时重力条纹位移消失；一般无吸收模型恒有 $I_1+I_2=R_m$。

图~\ref{advA6:quality:A6-source-exercises:a5} 中，两条到达同一输出端口的路径须保持相干，重力相位取两路势能积分之差；另一个端口的相位与分束振幅相联，不能把某端口变暗解释成中子消失。
\end{advsolution}

\begin{advexamplebox}{发光杆、激光脉冲与恒星光行差}\label{advA6:ex:rod-pulse-aberration}
惯性系 $S'$ 相对 $S$ 以速率 $v$ 沿 $+x$ 运动，两系坐标轴平行。
\begin{enumerate}
\item $S$ 中有一根静止发光直杆，与 $x$ 轴夹角为 $\theta$，长度为 $l$。
求 $S'$ 中同时测量所得的杆长 $l_1$ 和放置方向 $\theta_1$。
\item $S$ 中沿与 $x$ 轴夹角 $\theta$ 的方向传播一个长为 $l$ 的激光脉冲，传播速度为 $c$。
求 $S'$ 中的传播方向 $\theta_2$ 和脉冲长度 $l_2$，说明长度的测量约定。
\item 在远方恒星静止系中，恒星位于 $\theta$ 方向，天文台向右以 $v$ 运动。
不换参考系，从运动学推导望远镜在该系中的朝向 $\theta_3$，使星像位于视野正中。
\item 再换到天文台参考系，分别计算望远镜轴与恒星的方向，验证二者重合。
\end{enumerate}
\begin{center}
% Source spring-II p194 lecture10 problem2, both original technical panels.
% Telescope: theta=60deg,beta=.4; theta3=atan(sin(theta)/(cos(theta)+beta)).
% Stellar direction at aperture is NOT telescope axis; both are drawn separately.
\begin{tikzpicture}[>=Latex,font=\small]
 \begin{scope}[shift={(-3.5,0)}]
  \draw[->] (0,0)--(3.1,0) node[right]{$x$};
  \draw[->] (0,0)--(0,3.25) node[above]{$y,S$};
  \draw[->,TextGray] (.28,-.28)--(3.1,-.28) node[right]{$x'$};
  \draw[->,TextGray] (.28,-.28)--(.28,2.9) node[above right]{$S'$};
  \draw[DeepBlue,line width=3pt] (0,0)--({2.8*cos(60)},{2.8*sin(60)});
  \node[right=7pt] at (.7,1.21){$l$};
  \draw (.72,0) arc[start angle=0,end angle=60,radius=.72];
  \node at (.88,.43){$\theta$};
  \draw[->,thick] (.65,2.75)--(1.45,2.75) node[midway,above]{$v$};
  \node at (1.35,-.85){(a) 静止发光杆与两个参考系};
 \end{scope}
 \begin{scope}[shift={(1.9,.25)}]
  \pgfmathsetmacro{\tt}{atan(sin(60)/(cos(60)+.4))}
  \coordinate (L) at ({2*cos(\tt)},{2*sin(\tt)});
  \draw[TextGray] (-.65,-.6)--(-.65,.35) arc[start angle=180,end angle=0,radius=1.25]--(1.85,-.6)--cycle;
  \begin{scope}[rotate=\tt]
   \draw[thick] (0,-.08) rectangle (2,.08);
   \draw[thick] (2,-.14)--(2,.14);
  \end{scope}
  \draw[dashed] (0,0)--(1.1,0);
  \draw (.67,0) arc[start angle=0,end angle=\tt,radius=.67];
  \node at (.95,.32){$\theta_3$};
  \draw[dashed,DeepBlue] (L)--++(60:1.55);
  \draw[dashed] (L)--++(1.15,0);
  \draw ($(L)+(.58,0)$) arc[start angle=0,end angle=60,radius=.58];
  \node at ($(L)+(.74,.4)$){$\theta$};
  \node at ($(L)+(60:1.7)$){$\ast$};
  \node[right=3pt] at ($(L)+(60:1.7)$){恒星};
  \draw[->,thick] (2,-.42)--(2.8,-.42) node[midway,below]{$v$};
  \node at (1.1,-1.1){(b) 运动天文台中的望远镜};
 \end{scope}
\end{tikzpicture}

\figcap{题图：发光杆、脉冲和观察方向的原TikZ配置保留。}\label{advA6:quality:A6-source-exercises:q6}

\end{center}
\end{advexamplebox}
\global\AdvHideHeadingtrue
\begin{advsolution}
\begin{center}\begin{minipage}{.96\linewidth}\centering
% EMAdvFigureHeading
\raggedright{\bfseries\color{DeepBlue}解答}\par\smallskip\centering

\begin{tikzpicture}[>=Latex,font=\small]
 \begin{scope}[shift={(-3.5,0)}]
  \draw[->] (0,0)--(3.1,0) node[right]{$x$};
  \draw[->] (0,0)--(0,3.25) node[above]{$y,S$};
  \draw[->,TextGray] (.28,-.28)--(3.1,-.28) node[right]{$x'$};
  \draw[->,TextGray] (.28,-.28)--(.28,2.9) node[above right]{$S'$};
  \draw[DeepBlue,line width=3pt] (0,0)--({2.8*cos(60)},{2.8*sin(60)});
  \node[right=7pt] at (.7,1.21){$l$};
  \draw (.72,0) arc[start angle=0,end angle=60,radius=.72];
  \node at (.88,.43){$\theta$};
  \draw[->,thick] (.65,2.75)--(1.45,2.75) node[midway,above]{$v$};
  \node at (1.35,-.85){(a) 静止发光杆与两个参考系};
 \end{scope}
 \begin{scope}[shift={(1.9,.25)}]
  \pgfmathsetmacro{\tt}{atan(sin(60)/(cos(60)+.4))}
  \coordinate (L) at ({2*cos(\tt)},{2*sin(\tt)});
  \draw[TextGray] (-.65,-.6)--(-.65,.35) arc[start angle=180,end angle=0,radius=1.25]--(1.85,-.6)--cycle;
  \begin{scope}[rotate=\tt]
   \draw[thick] (0,-.08) rectangle (2,.08);
   \draw[thick] (2,-.14)--(2,.14);
  \end{scope}
  \draw[dashed] (0,0)--(1.1,0);
  \draw (.67,0) arc[start angle=0,end angle=\tt,radius=.67];
  \node at (.95,.32){$\theta_3$};
  \draw[dashed,DeepBlue] (L)--++(60:1.55);
  \draw[dashed] (L)--++(1.15,0);
  \draw ($(L)+(.58,0)$) arc[start angle=0,end angle=60,radius=.58];
  \node at ($(L)+(.74,.4)$){$\theta$};
  \node at ($(L)+(60:1.7)$){$\ast$};
  \node[right=3pt] at ($(L)+(60:1.7)$){恒星};
  \draw[->,thick] (2,-.42)--(2.8,-.42) node[midway,below]{$v$};
  \node at (1.1,-1.1){(b) 运动天文台中的望远镜};
 \end{scope}
\end{tikzpicture}
\figcap{解图：发光杆、脉冲和观察方向的原TikZ配置保留。}\label{advA6:quality:A6-source-exercises:a6}
\end{minipage}\end{center}

令 $\beta=v/c$、$\gamma=(1-\beta^2)^{-1/2}$，所有角均从本系的正 $x$ 轴量起。
这里“方向”指坐标测量方向，不是未作光传播时延校正的杆的照片形状。

（1）杆在 $S$ 中静止，两端坐标差恒为 $(l\cos\theta,l\sin\theta)$。
$S'$ 的测长必须使用 $\Delta t'=0$ 的两端事件。
由洛伦兹变换 $\Delta t'=\gamma(\Delta t-v\Delta x/c^2)=0$，得到
$\Delta t=v\Delta x/c^2$，从而
\[
 \Delta x'=\gamma(\Delta x-v\Delta t)=\frac{l\cos\theta}{\gamma},
 \qquad \Delta y'=l\sin\theta.
\]
因此
\[
 \boxed{l_1=l\sqrt{\cos^2\theta/\gamma^2+\sin^2\theta}},\qquad
 \boxed{\tan\theta_1=\gamma\tan\theta}.
\]
正切式不能单独确定象限，须同时使用 $\Delta x',\Delta y'$ 的符号。
沿运动方向的杆缩短为 $l/\gamma$，横杆长度不变，作为两种极限检验。

（2）光速的分量为 $u_x=c\cos\theta,u_y=c\sin\theta$。速度变换给出
\[
 \boxed{\cos\theta_2=\frac{\cos\theta-\beta}{1-\beta\cos\theta}},\qquad
 \boxed{\sin\theta_2=\frac{\sin\theta}{\gamma(1-\beta\cos\theta)}}.
\]
两式平方和为1，表明合速度仍为 $c$。这与静止杆的方向变换不同：
杆的两端是两个物质点，光的传播方向是一个光子相邻两事件的速度方向。

通常把激光脉冲长度定义为近似平面的前、后包络波面之间\emph{沿传播方向的距离}；
有限横向光束的完整形状还需要横向包络信息。
在 $S$ 中令 $\vect n=(\cos\theta,\sin\theta)$，脉冲占据的相位区间宽度是
$c t-\vect n\cdot\vect r$ 的宽度 $l$。波四矢量变换给频率比
$D=\gamma(1-\beta\cos\theta)$；因同一组前后包络面的相位差不变，
而在 $S'$ 同时截面上波数扩大为 $D$ 倍，包络面的法向间距便为
\[
 \boxed{l_2=\frac{l}{\gamma(1-\beta\cos\theta)}}.
\]
也可取脉冲中相同的振荡周期数 $N$：$l=Nc/\nu$、$l_2=Nc/\nu'=l/D$。

如果“光束”指一串横向宽度趋于零、首尾光子在 $S$ 中相距 $l\vect n$ 的有限光子列，
那么 $S'$ 同时测量的\emph{两个端点的空间距离}是另一个量。
取后端在 $t=0$ 经过原点，前端轨迹为 $\vect r_f=l\vect n+c t\vect n$。
选后端事件 $t'=0$，前端同时事件满足
$t_f=v x_f/c^2$，解得 $t_f=\beta l\cos\theta/[c(1-\beta\cos\theta)]$。于是
\[
 \Delta\vect r'=
 \frac{l}{1-\beta\cos\theta}\left(\frac{\cos\theta}{\gamma},\sin\theta\right),
 \qquad
 |\Delta\vect r'|=l\sqrt{\frac{1+\beta\cos\theta}{1-\beta\cos\theta}}.
\]
这根首尾连线一般\emph{不沿}光子的传播方向，不能把其长度冒充包络的纵向长度；
把它投影到上述 $\vect n'$ 后，才得到 $\Delta\vect r'\cdot\vect n'=l/D$。
两种定义都明确后，便不会把不同测量过程的两个公式误认为矛盾。

（3）恒星的视线方向为 $\vect n=(\cos\theta,\sin\theta)$，
光实际沿 $-\vect n$ 传播，不能漏掉这个负号。
图(b)中镜筒由后端指向前端的向量记为 $\vect L$。
光在时刻0到达前端，在时刻 $\Delta t$ 到达后端；这段时间里后端向右移动 $v\Delta t$。
在同一个恒星静止系里，光的位移与两端的位移关系为
\[
 -c\vect n\Delta t=-\vect L+v\vect e_x\Delta t,
 \quad\Rightarrow\quad
 \vect L=c\Delta t(\vect n+\beta\vect e_x).
\]
因此镜筒朝向应满足
\[
 \boxed{\tan\theta_3=\frac{\sin\theta}{\cos\theta+\beta}}.
\]
仍用两个分量共同确定象限。图中恒星虚线与镜筒轴不平行，正是因为光穿过镜筒期间，
镜筒本身也向右移动；不能在恒星系强行把镜筒轴画成入射光线。

（4）镜筒在天文台系静止；第一问的长度收缩关系逆用时，
恒星系看到的纵向长度比天文台系短 $1/\gamma$，故
\[
 \tan\theta'_{\rm tube}=\frac{\tan\theta_3}{\gamma}
 =\frac{\sin\theta}{\gamma(\cos\theta+\beta)}.
\]
另一方面把传播方向 $-\vect n$ 代入第二问的速度变换，再取反便是恒星的观测方向：
\[
 \cos\theta'_{\rm star}=\frac{\cos\theta+\beta}{1+\beta\cos\theta},\qquad
 \sin\theta'_{\rm star}=\frac{\sin\theta}{\gamma(1+\beta\cos\theta)}.
\]
二者分量比和象限均与镜筒相同，所以星像位于镜筒正中。
当 $v\to0$，四个方向均回到 $\theta$；对 $\theta=90^\circ$，
恒星观测方向具有正 $x'$ 分量 $\beta$，即向天文台运动前方偏移。

图~\ref{advA6:quality:A6-source-exercises:a6} 中，杆长由新系同时截面测得，照片方向由同一接收时刻的迟延光决定；激光脉冲长度还涉及两个端点事件的选择，三个测量问题必须各自指定事件。
\end{advsolution}

\begin{advexamplebox}{等速追及曲线与迟延接收的多普勒频移}\label{advA6:ex:pursuit-doppler}
雷达鼠R以恒定速率 $v$ 沿直线运动，电波狗D也以相同速率运动，
且它每一时刻的速度都指向该时刻的R。初始时狗的速度与鼠的速度垂直，两者距离为 $p$。
\begin{enumerate}
\item 当瞬时连线与鼠运动方向的夹角为 $\theta$ 时，求两者距离。
\item 狗的静质量为 $m$，求维持该运动所需的力。
\item 狗在此时发射自身瞬时静止系频率为 $\nu_0$ 的短光脉冲，求鼠接收到的频率。
必须考虑传播期间鼠的位置变化。
\end{enumerate}
\begin{center}
% Source spring-II p202 lecture11 problem3; analytic pursuit with equal speeds.
% Parameter q=tan(theta/2), xD=p(q^2-1)/4-p ln(q)/2, yD=-p q.
% Tangent points at R_e. Radio ray points at later R_r, not instantaneous R_e.
\begin{tikzpicture}[>=Latex,font=\small]
 \pgfmathsetmacro{\pp}{3}\pgfmathsetmacro{\uu}{.45}
 \pgfmathsetmacro{\ct}{(1-\uu*\uu)/(1+\uu*\uu)}
 \pgfmathsetmacro{\st}{2*\uu/(1+\uu*\uu)}
 \pgfmathsetmacro{\ell}{\pp*(1+\uu*\uu)/2}
 \pgfmathsetmacro{\xd}{\pp*(\uu*\uu-1)/4-\pp*ln(\uu)/2}
 \pgfmathsetmacro{\xr}{\xd+\ell*\ct}
 \pgfmathsetmacro{\aa}{sqrt(1-.4*.4*\st*\st)-.4*\ct}
 \coordinate (D) at (\xd,{-\pp*\uu});
 \coordinate (R) at (\xr,0);
 \coordinate (H) at ({\xr+.4*\ell/\aa},0);
 \draw[->,thick] (-.2,0)--(4.5,0);
 \draw[DeepBlue,thick,domain=1:.05,samples=130,smooth,variable=\q]
  plot ({\pp*(\q*\q-1)/4-\pp*ln(\q)/2},{-\pp*\q});
 \draw[dashed] (0,0)--(0,-3);
 \draw[<->] (-.48,0)--(-.48,-3) node[midway,left]{$p$};
 \draw[dashed] (D)--(R);
 \draw[->,thick,DeepBlue] (D)--++({.8*\ct},{.8*\st}) node[left=6pt]{$\vect v_D$};
 \draw[->,TextGray] (D)--(H);
 \node[anchor=west] at (4,-1.35){光信号};
 \draw[TextGray] (3.95,-1.35)--($(D)!.5!(H)$);
 \draw[->,thick] (.25,.23)--(.9,.23) node[midway,above]{$v$};
 \draw[->,thick] (0,-3)--(0,-2.3);
 \foreach \name in {D,R,H}{\fill (\name) circle(2pt);}
 \fill (0,0) circle(2pt) (0,-3) circle(2pt);
 \node[above left] at (0,0){$R_0$};\node[below] at (0,-3){$D_0$};
 \node[left=6pt] at (D){$D_e$};
 \node[above=5pt] at (R){$R_e$};\node[above=5pt] at (H){$R_r$};
 \draw ($(R)+(-.55,0)$) arc[start angle=180,end angle={180+2*atan(\uu)},radius=.55];
 \node at ($(R)+(-.78,-.22)$){$\theta$};
 \node[align=left,anchor=west] at (2.1,-2.35){$e$：发射时刻\\$r$：接收时刻};
\end{tikzpicture}

\figcap{题图：雷达鼠和电波狗的初始垂直速度及连线，原TikZ图保留。}\label{advA6:quality:A6-source-exercises:q7}

\end{center}
\end{advexamplebox}
\begin{advsolution}
令鼠沿 $+x$ 运动，初态 $R_0=(0,0)$、$D_0=(0,-p)$。
图中 $D_e,R_e$ 是发射时刻的位置，$R_r$ 是接收时刻的位置。
用 $\ell$ 表示瞬时距离，$\vect e=(\cos\theta,\sin\theta)$ 指向由狗到鼠，
狗的速度为 $v\vect e$，鼠的速度为 $v\vect e_x$。

（1）相对位置满足 $\vect r_R-\vect r_D=\ell\vect e$。
对时间求导，并沿连线及其法向 $\vect e_\theta=(-\sin\theta,\cos\theta)$ 分解：
\[
 \dot\ell=v(\cos\theta-1),\qquad
 \ell\dot\theta=-v\sin\theta.
\]
两式相除以消去时间，得到
\[
 \frac{\dd\ell}{\ell}=\frac{1-\cos\theta}{\sin\theta}\dd\theta
 =\tan\frac\theta2\dd\theta=-2\dd\ln\cos\frac\theta2.
\]
用初值 $\theta=\pi/2,\ell=p$ 确定常数，得
\[
 \boxed{\ell=\frac{p}{1+\cos\theta}}.
\]
角度逐渐减小、距离也减小，但极限为 $p/2$，狗始终追不上同速的鼠。
题图轨迹由这一方程生成；若设 $q=\tan(\theta/2)$，可积分得
\[
 x_D=\frac p4(q^2-1)-\frac p2\ln q,\quad y_D=-pq,
 \quad x_R=\frac p4(1-q^2)-\frac p2\ln q.
\]
对这些表达式求导，轨迹切线恰指向 $R_e$，图中的速度箭头因而不是凭外观放置的。

（2）速率不变，所以 $\gamma=(1-v^2/c^2)^{-1/2}$ 为常数，
动量大小为 $\gamma mv$，只有方向转动。力为
\[
 \vect F=\frac{\dd}{\dd t}(\gamma mv\vect e)
 =\gamma mv\dot\theta\vect e_\theta
 =\boxed{\frac{\gamma mv^2\sin\theta}{\ell}(\sin\theta,-\cos\theta)}.
\]
其大小为 $\gamma mv^2\sin\theta(1+\cos\theta)/p$，
方向如下面的解图所示，从 $D_e$ 指向曲率中心一侧，与瞬时速度垂直，因此外力不做功、动能恒定。
\begin{center}\begin{minipage}{\linewidth}\centering
% Solution-only: equal-speed pursuit has force normal to velocity.
\begin{tikzpicture}[>=Latex,scale=.95]
 \fill(0,0)circle(1.6pt)node[left]{$D_e$};
 \draw[->,thick,DeepBlue](0,0)--(1.6,1.2)node[above right]{$\vect v_D$};
 \draw[->,thick](0,0)--(1.2,-1.6)node[below right]{$\vect F$};
 \draw[thin](.24,.18)--(.42,-.06)--(.18,-.24);
 \node[anchor=west]at(2.2,0){$\vect F\cdot\vect v_D=0$};
\end{tikzpicture}
\figcap{同速追逐的受力方向：力垂直于瞬时速度，改变动量方向而不改变其大小。}\label{advA6:quality:A6-source-exercises:a7}\label{advA6:fig:pursuit-force-solution}
\end{minipage}\end{center}

初始力向右，大小为 $\gamma mv^2/p$；当 $\theta\to0$ 时曲线趋于直线，所需力趋于零。

（3）不能把发射时刻的 $D_eR_e$ 当作光脉冲的传播方向：
鼠在光到达前继续向右移动，所以光应沿 $D_eR_r$。
令实验室光方向单位矢量为 $\vect n$，传播时间为 $\Delta t$，并记 $\beta=v/c$。
追光的几何条件是
\[
 c\Delta t\vect n=\ell\vect e+v\Delta t\vect e_x.
\]
设 $a=\ell/(c\Delta t)>0$，则 $\vect n=a\vect e+\beta\vect e_x$。
由 $|\vect n|=1$ 得
\[
 a^2+2\beta a\cos\theta+\beta^2=1,
 \quad a=Q-\beta\cos\theta,\quad
 Q=\sqrt{1-\beta^2\sin^2\theta}.
\]
取正根保证光向未来传播，且 $Q>\beta\cos\theta$。

设该光在实验室频率为 $\nu$。发射者、接收者所测频率分别由四波矢与各自四速度的内积给出：
\[
 \nu_0=\gamma\nu(1-\beta\vect n\cdot\vect e),\qquad
 \nu_r=\gamma\nu(1-\beta n_x).
\]
两者速率相同，$\gamma$ 因子相消，但方向因子不相消。用
$\vect n\cdot\vect e=a+\beta\cos\theta=Q$、$n_x=\beta+a\cos\theta$，得到
\[
 \boxed{\frac{\nu_r}{\nu_0}
 =\frac{1-\beta^2-\beta a\cos\theta}{1-\beta Q}
 =\frac{aQ}{1-\beta Q}},\qquad a=Q-\beta\cos\theta.
\]
这里的频率是短脉冲的局部载频，忽略脉冲持续期间狗的速度方向变化。
检验低速极限：$Q\simeq1$、$a\simeq1-\beta\cos\theta$，故
$\nu_r/\nu_0\simeq1+\beta(1-\cos\theta)$，与距离在减小产生的蓝移一致。
当 $\theta\to0$，两者同向同速，比例回到1。初始 $\theta=90^\circ$ 时也须使用上述迟延几何，
不能直接套“横向接收”公式，因为实际光线并不沿瞬时竖直连线。
低速蓝移也不能推广到所有高速情况。例如 $\beta=0.9,\theta=60^\circ$ 时，
本式给出 $\nu_r/\nu_0\simeq0.2535$：尽管瞬时间距减小，迟延接收仍出现红移。

图~\ref{advA6:quality:A6-source-exercises:a7} 中，追及曲线中的连线连接同一时刻的两个位置，接收频移中的光线却连接发射与接收两个迟延事件；频移必须用后者的传播方向。
\end{advsolution}
\endgroup

% AdvancedEMQuality20261023
\section{相对论辐射、散射与综合模型}
\noindent 本节把方向变换、能动量守恒和波相位联系起来。先明确观测量的时间口径与参考系，再列不变量，最后检查低速、无限重靶及零重力等极限；同一符号不能跨参考系沿用而省略变换。

\begin{advexamplebox}{运动光源的辐射分布与星空密度}\label{advA6:ex:radiation-angle}
\begin{enumerate}
\item 单色点光源在自身静止系向各方向均匀辐射，发光强度为 $I_0$。
光源相对实验室以速率 $v$ 向前匀速运动，求所测强度与方向的关系。
\item 静止观察者看到星云中恒星的单位立体角数密度均匀。
若观察者以速率 $v$ 匀速运动，他看到的星体数密度角分布如何变化？
这与同一时刻测定星体实际空间位置后得到的角分布是否相同？
\end{enumerate}

\begin{center}\begin{minipage}{.96\linewidth}\centering
\begin{tikzpicture}[>=Latex,font=\small,scale=.85,every node/.style={fill=none}]

\fill (0,0) circle (2pt) node[below]{光源};\draw[->] (0,-.5)--(2,-.5) node[midway,below]{$v$};\draw[dashed,->] (0,0)--(3.5,0);\draw[DeepBlue,->] (0,0)--(2.4,1.4) node[right]{观察光线};\node at (.9,.2){$\theta$};\draw[gray,dashed] (0,0) circle (1.1);

\end{tikzpicture}
\figcap{题图：运动方向、观察方向与源静止系各向同性条件；球形虚线不预设实验室角分布。}\label{advA6:quality:A6-source-more:q1}
\end{minipage}\end{center}
\end{advexamplebox}
\global\AdvHideHeadingtrue
\begin{advsolution}
\begin{center}\begin{minipage}{.96\linewidth}\centering
% EMAdvFigureHeading
\raggedright{\bfseries\color{DeepBlue}解答}\par\smallskip\centering

\begin{tikzpicture}[>=Latex,font=\small,scale=.85,every node/.style={fill=none}]

\fill (0,0) circle (2pt) node[below]{光源};\draw[->] (0,-.5)--(2,-.5) node[midway,below]{$v$};\draw[dashed,->] (0,0)--(3.5,0);\draw[DeepBlue,->] (0,0)--(2.4,1.4) node[right]{观察光线};\node at (.9,.2){$\theta$};\draw[gray,dashed] (0,0) circle (1.1);

\end{tikzpicture}
\figcap{解图：运动方向、观察方向与源静止系各向同性条件；球形虚线不预设实验室角分布。}\label{advA6:quality:A6-source-more:a1}
\end{minipage}\end{center}

（1）先明确强度的时间约定：取 $I_0=\dd E_0/(\dd\tau\dd\Omega_0)$，
即光源静止系每单位时间、每单位立体角辐射的能量；若给的是总功率，以下 $I_0$ 应换成总功率除以 $4\pi$。
实验室出射方向与光源速度的夹角为 $\theta$，令
\[
 D=\frac1{\gamma(1-\beta\cos\theta)},\quad \beta=v/c.
\]
光行差关系 $\cos\theta=(\cos\theta_0+\beta)/(1+\beta\cos\theta_0)$ 给出
$\dd\Omega=\dd\Omega_0/D^2$；每个光子的能量变为 $D$ 倍，
同一组光源事件的实验室发射时间间隔为 $\dd t=\gamma\dd\tau$。
所以按\emph{发射时刻}计的角功率为
\[
 \boxed{\frac{\dd P_{\rm em}}{\dd\Omega}=I_0\frac{D^3}{\gamma}}.
\]
远处固定探测器收到相邻光信号的时间还要扣除光源向它前进所节省的传播时间：
\[
 \dd t_{\rm arr}=\dd t(1-\beta\cos\theta)=\frac{\dd\tau}{D}.
\]
于是按\emph{到达时刻}计、已扣除几何距离平方衰减的观测角强度为
\[
 \boxed{I_{\rm obs}(\theta)=I_0D^4}.
\]
这两个结果的幂次不同，是时间分母不同，不是多普勒公式矛盾。
若需要有限频宽的谱强度，还须变换频率间隔；此处单色能量已对谱线积分。
$v\to0$ 时两式均回到 $I_0$；正前方 $D$ 大于正后方，辐射向前集中。

（2）记静止系中恒星的\emph{视线}角余弦为 $\mu$，运动观察者所见为 $\mu'$。
视线与光传播方向相反，故
\[
 \mu'=\frac{\mu+\beta}{1+\beta\mu},\qquad
 \frac{\dd\Omega}{\dd\Omega'}=
 \frac1{\gamma^2(1-\beta\mu')^2}.
\]
若统计同一组恒星、不另加亮度阈值，星数不因换系改变。
静止时单位立体角恒星数为 $n_0$，则
\[
 \boxed{n'_{\rm apparent}(\mu')=
 \frac{n_0}{\gamma^2(1-\beta\mu')^2}}.
\]
积分 $2\pi\int_{-1}^1n'_{\rm apparent}\dd\mu'=4\pi n_0$，总星数守恒；
光行差使更多视线集中到运动前方。

“实际位置角分布”还需给定恒星运动及取样方式。
若同一组恒星全在原系静止，并在两观察者于原点相遇的事件取 $t=t'=0$，
按新系 $t'=0$ 同时确定其坐标，则 $x'=x/\gamma,y'=y,z'=z$。
设实际位置方向余弦为 $\eta'$，原方向余弦为 $\eta$，有
\[
 \eta=\frac{\gamma\eta'}{\sqrt{1+(\gamma^2-1)\eta'^2}},\qquad
 \boxed{n'_{\rm position}(\eta')=
 \frac{n_0\gamma}{[1+(\gamma^2-1)\eta'^2]^{3/2}}}.
\]
这里使用的是空间位置而非迟延到达的光线，分布关于前后对称，与所见星空的前向聚集不同。
若恒星各有自行、按新系固定半径截取样本，或按探测亮度筛选，单凭原先的均匀角分布不能确定新的实际分布；
必须再给速度、距离或光度信息，不能把上述特定模型推广到任意星云。
即使恒星静止，离开这个相遇时刻后还有 $-vt'$ 的坐标平移；有限距离星云的实际角分布也会随取样时刻改变。

图~\ref{advA6:quality:A6-source-more:a1} 中，光行差改变立体角元，多普勒因子改变单光子能量；若强度按观察者接收时间定义，还需计入光行时压缩，不能与按发射实验室时间定义的功率混用。
\end{advsolution}

\begin{advexamplebox}{经典原子寿命与同步辐射}\label{advA6:ex:synchrotron}
\begin{enumerate}
\item 经典氢原子中电子绕固定原子核作圆周运动。用非相对论拉莫尔辐射公式估计其坍缩寿命。
\item 同步辐射装置中，近光速电子在偏转磁铁内改变方向并辐射。某存储环参数为：
电子能量 $3.5\,\mathrm{GeV}$、束流 $300\,\mathrm{mA}$、周长 $432\,\mathrm m$、
平均分配辐射的弯道单元数 $20$、偏转磁场 $1.27\,\mathrm T$。
估算平均每个光口的辐射总功率。
\item 辐射集中在沿瞬时切线、半顶角约 $1/\gamma$ 的光锥内。
估算远处一个固定方向收到单个电子脉冲的持续时间。
\end{enumerate}

\begin{center}\begin{minipage}{.96\linewidth}\centering
\begin{tikzpicture}[>=Latex,font=\small,scale=.85,every node/.style={fill=none}]

\draw[thick] (0,0) circle (.9);\fill (0,0) circle (2pt) node[below]{核};\fill (.9,0) circle (1.5pt) node[right]{$e^-$};\draw[DeepBlue,->] (.9,0)--(.9,.7);\node[below] at (0,-1.3){经典Coulomb圆轨道};
\begin{scope}[xshift=4.1cm]\draw[thick] (0,0) circle (1.1);\draw[DeepBlue,->] (1.1,0)--(1.1,.8);\node[right] at (1.1,.2){近光速电子};\node[below] at (0,-1.3){存储环弯转半径 $\rho$};\end{scope}

\end{tikzpicture}
\figcap{题图：原子圆轨道与存储环弯转轨道分别使用各自的动力学和辐射近似。}\label{advA6:quality:A6-source-more:q2}
\end{minipage}\end{center}
\end{advexamplebox}
\global\AdvHideHeadingtrue
\begin{advsolution}
\begin{center}\begin{minipage}{.96\linewidth}\centering
% EMAdvFigureHeading
\raggedright{\bfseries\color{DeepBlue}解答}\par\smallskip\centering

\begin{tikzpicture}[>=Latex,font=\small,scale=.85,every node/.style={fill=none}]

\draw[thick] (0,0) circle (.9);\fill (0,0) circle (2pt) node[below]{核};\fill (.9,0) circle (1.5pt) node[right]{$e^-$};\draw[DeepBlue,->] (.9,0)--(.9,.7);\node[below] at (0,-1.3){经典Coulomb圆轨道};
\begin{scope}[xshift=4.1cm]\draw[thick] (0,0) circle (1.1);\draw[DeepBlue,->] (1.1,0)--(1.1,.8);\node[right] at (1.1,.2){近光速电子};\node[below] at (0,-1.3){存储环弯转半径 $\rho$};\end{scope}

\end{tikzpicture}
\figcap{解图：原子圆轨道与存储环弯转轨道分别使用各自的动力学和辐射近似。}\label{advA6:quality:A6-source-more:a2}
\end{minipage}\end{center}

（1）令 $k_e=1/(4\pi\varepsilon_0)$，轨道半径为 $r$。
圆周运动给 $a=k_e e^2/(m_er^2)$，总机械能为 $E=-k_e e^2/(2r)$。
拉莫尔功率和能量损失方程为
\[
 P=\frac{2k_e e^2a^2}{3c^3}
 =\frac{2k_e^3e^6}{3m_e^2c^3r^4},\qquad
 \frac{k_e e^2}{2r^2}\frac{\dd r}{\dd t}=-P.
\]
故 $\dd r/\dd t=-4k_e^2e^4/(3m_e^2c^3r^2)$。
从初始半径 $r_0$ 积分到0，得到经典模型的时间
\[
 \boxed{\tau=\frac{m_e^2c^3r_0^3}{4k_e^2e^4}}.
\]
题设没有给 $r_0$，所以首先保留这一参数化结果；以典型原子尺度
$r_0=a_0=5.29\times10^{-11}\,\mathrm m$ 为例，$\tau\simeq1.56\times10^{-11}\,\mathrm s$。
接近核心时非相对论、缓慢缩轨近似会失效，但时间积分主要来自外圈，数量级仍说明经典模型无法解释稳定原子。

（2）取题给电子能量为总能量，$\gamma=E/(m_ec^2)\simeq6849$、$v\simeq c$。
弯道曲率半径由动量和磁力确定，而不能用整个环周长除以 $2\pi$：
\[
 \rho=\frac{p}{eB}=\frac{\gamma m_ev}{eB}\simeq9.19\,\mathrm m.
\]
加速度垂直速度时，李纳辐射公式给单电子\emph{在弯道内}的瞬时功率
\[
 P_e=\frac{e^2\gamma^4a^2}{6\pi\varepsilon_0c^3}
 =\frac{e^2\gamma^4v^4}{6\pi\varepsilon_0c^3\rho^2}
 \simeq1.20\times10^{-6}\,\mathrm W.
\]
假定闭合平面环各弯道曲率相同、总偏转角 $2\pi$，直线段辐射忽略，
一个电子每圈在弯道内的总时间为 $2\pi\rho/v$。
束中电子数由电流求得 $N_e=IL/(ev)$；处于弯道的比例是 $2\pi\rho/L$。
于是
\[
 P_{\rm total}=P_e\frac{IL}{ev}\frac{2\pi\rho}{L}
 =P_e\frac{2\pi I\rho}{ev}\simeq4.33\times10^5\,\mathrm W,
 \qquad
 \boxed{P_0=\frac{P_{\rm total}}{20}\simeq2.17\times10^4\,\mathrm W}.
\]
周长在最后约去，是因为环越长时同样电流对应更多电子，但电子在弯道中所占时间比例更小。
这是各单元总辐射平均分配的估计；实际光口受孔径、接受角和传输损耗限制，收集功率可更低。

（3）取光锥的半角 $\alpha=1/\gamma$，固定观察方向位于圆弧中点的切线方向。
电子从 $-\alpha$ 转到 $+\alpha$ 的发射时间为 $2\rho\alpha/v$。
末次发光点沿观察方向比首次前进了 $2\rho\sin\alpha$，光程缩短，故到达时间跨度是
\[
 \Delta T=\frac{2\rho\alpha}{v}-\frac{2\rho\sin\alpha}{c}
 =\frac{2\rho}{c}\left(\frac\alpha\beta-\sin\alpha\right).
\]
分别展开 $1/\beta\simeq1+1/(2\gamma^2)$、$\sin\alpha\simeq\alpha-\alpha^3/6$，得到
\[
 \boxed{\Delta T\simeq\frac{4\rho}{3c\gamma^3}
 \simeq1.27\times10^{-19}\,\mathrm s}.
\]
这个系数对应“半顶角严格取 $1/\gamma$”的几何宽度；真实脉冲无硬边界，不同宽度定义的数值系数不同，
但 $\rho/(c\gamma^3)$ 的时间尺度相同。漏掉光程变化就会误得大得多的发射时间尺度 $\rho/(c\gamma)$。

图~\ref{advA6:quality:A6-source-more:a2} 中，原子模型需用能量损失求半径随时间收缩，不能用初始能量除初始功率替代整个积分；存储环的近光速辐射含相对论增强，周长也不能直接替代偏转磁铁的弯曲半径。
\end{advsolution}

\begin{advexamplebox}{正交连续换系后的坐标轴方向}\label{advA6:ex:successive-boosts}
$S'$ 相对 $S$ 沿 $+x$ 以 $v_1$ 运动，$S''$ 相对 $S'$ 沿 $+y'$ 以 $v_2$ 运动。
在 $S$ 的同一时刻测量，求 $S''$ 的 $x''$ 轴与 $x$ 轴夹角。
分别用同步钟与测长的物理方法、直接洛伦兹变换两种方法求解。

\begin{center}\begin{minipage}{.96\linewidth}\centering
\begin{tikzpicture}[>=Latex,font=\small,scale=.85,every node/.style={fill=none}]

\draw[->] (0,0)--(3,0) node[right]{$x$};\draw[->] (0,0)--(0,2) node[above]{$y$};\draw[DeepBlue,->] (.4,.4)--(2,.4) node[midway,above]{$v_1\boldsymbol e_x$};\draw[DeepBlue,->] (2,.5)--(2,1.7) node[right]{$v_2\boldsymbol e_{y'}$};

\end{tikzpicture}
\figcap{题图：两次正交换系的速度分别相对各自前一参考系定义。}\label{advA6:quality:A6-source-more:q3}
\end{minipage}\end{center}
\end{advexamplebox}
\global\AdvHideHeadingtrue
\begin{advsolution}
\begin{center}\begin{minipage}{.96\linewidth}\centering
% EMAdvFigureHeading
\raggedright{\bfseries\color{DeepBlue}解答}\par\smallskip\centering

\begin{tikzpicture}[>=Latex,font=\small,scale=.85,every node/.style={fill=none}]

\draw[->] (-.3,0)--(3.2,0) node[right]{$x$};\draw[->] (0,-1.2)--(0,1.5) node[above]{$y$};\draw[DeepBlue,thick,->] (0,0)--(2.6,-.9) node[right]{$x''$（$S$ 同时截面）};\node at (1,-.12){$\psi<0$};\node[above] at (1.4,1){$v_1,v_2>0$ 的示意};

\end{tikzpicture}
\figcap{解图：同时截面上的轴倾斜向右下，斜率由两个推动共同决定。}\label{advA6:quality:A6-source-more:a3}
\end{minipage}\end{center}

这里的夹角指 $S$ 同时截面上坐标轴的空间朝向，不是直接把两个速度方向夹角称为坐标轴转角。
记 $\gamma_i=(1-v_i^2/c^2)^{-1/2}$。

第一种方法只用同步钟、时间膨胀和长度收缩，不先写坐标变换。
一排随 $S'$ 沿 $x$ 移动、彼此按本系同步的钟，在 $S$ 同一时刻并不读数相同。
取本系间距为 $L'$ 的两钟，在本系中点同时发光校准。
$S$ 中两端离中点的初距离为 $L'/(2\gamma_1)$，前后传播时间差为
\[
 \frac{L'}{2\gamma_1}\left(\frac1{c-v_1}-\frac1{c+v_1}\right)
 =\frac{\gamma_1v_1L'}{c^2}.
\]
运动钟走时慢 $\gamma_1$ 倍，故在 $S$ 同一时刻，前方钟比后方钟少读 $v_1L'/c^2$。
也就是两处的本系时间差 $\Delta t'=-v_1\Delta x'/c^2$。
$x''$ 轴在 $S'$ 中保持平行 $x'$、以 $v_2$ 向上平移，因此这两个事件的高度差为
$\Delta y=v_2\Delta t'=-v_1v_2\Delta x'/c^2$。
横向坐标不收缩，而 $S$ 同时测得 $\Delta x=\Delta x'/\gamma_1$，所以
\[
 \boxed{\tan\psi=\frac{\Delta y}{\Delta x}
 =-\frac{\gamma_1v_1v_2}{c^2}}.
\]
当 $v_1,v_2>0$ 时，$x''$ 轴的正方向在 $S$ 中向右下方倾斜。

第二种方法直接写两步变换：
\[
 t'=\gamma_1(t-v_1x/c^2),\quad y'=y,
 \qquad y''=\gamma_2(y'-v_2t').
\]
在 $S$ 中取 $t=0$，$x''$ 轴上 $y''=0$ 的条件为
$y+\gamma_1v_1v_2x/c^2=0$，立刻得到同一斜率。
任一速度为零时夹角为零；低速时夹角为二阶小量 $-v_1v_2/c^2$。
连续推动还可分解成一次标准推动加一个维格纳旋转，但该旋转角不应不加区分地等同于这里的同时截面轴倾角。

图~\ref{advA6:quality:A6-source-more:a3} 中，$S$ 同时测量要求 $t$ 固定，$x''$ 轴还须满足 $y''=0$；合并两条件才得到右下方倾斜的轴线，该倾角不应直接称为标准推动分解中的Wigner旋转角。
\end{advsolution}

\begin{advexamplebox}{载流导线、横向力与场的变换}\label{advA6:ex:wire-field-transform}
惯性系 $S$ 中，一根沿 $x$ 的无限长直导线载有正向电流 $I$，导线在该系电中性。
在 $y=a>0$ 处有电荷 $q$，速度为 $v_x\vect e_x$。
\begin{enumerate}
\item 求电荷所受的电磁力。
\item 换到以 $u\vect e_x$ 运动的 $S'$，求同一事件处的力。
\item 用导线的电荷、电流变换推导垂直于换系速度的电场和磁场变换。
\end{enumerate}

\begin{center}\begin{minipage}{.96\linewidth}\centering
\begin{tikzpicture}[>=Latex,font=\small,scale=.85,every node/.style={fill=none}]

\draw[very thick] (-1,0)--(4.5,0);\draw[->] (0,0)--(1.5,0) node[midway,below]{$I>0$};\fill (1.8,1.3) circle (2pt) node[above]{$q$};\draw[DeepBlue,->] (1.8,1.3)--(3.4,1.3) node[midway,above]{$v_x$};\draw[<->] (.4,0)--(.4,1.3) node[midway,left]{$a$};\node[right] at (4.5,0){$x$};

\end{tikzpicture}
\figcap{题图：中性载流导线、电荷高度与电荷速度；换系沿导线方向。}\label{advA6:quality:A6-source-more:q4}
\end{minipage}\end{center}
\end{advexamplebox}
\global\AdvHideHeadingtrue
\begin{advsolution}
\begin{center}\begin{minipage}{.96\linewidth}\centering
% EMAdvFigureHeading
\raggedright{\bfseries\color{DeepBlue}解答}\par\smallskip\centering

\begin{tikzpicture}[>=Latex,font=\small,scale=.85,every node/.style={fill=none}]

\draw[very thick] (-1,0)--(4.5,0);\draw[->] (0,0)--(1.5,0) node[midway,below]{$I>0$};\fill (1.8,1.3) circle (2pt) node[above]{$q$};\draw[DeepBlue,->] (1.8,1.3)--(3.4,1.3) node[midway,above]{$v_x$};\draw[<->] (.4,0)--(.4,1.3) node[midway,left]{$a$};\node[right] at (4.5,0){$x$};
\node at (1.1,1.1){$\boldsymbol B\;\odot$};\draw[->] (1.8,1.1)--(1.8,.4) node[right]{$\boldsymbol F\ (qv_x>0)$};
\end{tikzpicture}
\figcap{解图：实验室磁力方向及其符号条件，其他参考系的电磁力须重新分解。}\label{advA6:quality:A6-source-more:a4}
\end{minipage}\end{center}

（1）由右手定则，在导线上方 $\vect B=B\vect e_z$，$B=\mu_0I/(2\pi a)$，而 $\vect E=0$。
由于 $\vect e_x\times\vect e_z=-\vect e_y$，
\[
 \boxed{\vect F=-qv_xB\vect e_y}.
\]
若 $qv_x>0$，力指向导线；若为负，则背离导线。

（2）把正、负载流组分各自的线密度与电流相加后，线四流变换为
\[
 \lambda'=\gamma_u(\lambda-uI/c^2)=-\gamma_u uI/c^2,
 \qquad I'=\gamma_u(I-u\lambda)=\gamma_u I.
\]
距离 $a$ 垂直于换系方向不变。用新系高斯定律和安培定律，
\[
 E'_y=\frac{\lambda'}{2\pi\varepsilon_0a}=-\gamma_u uB,
 \qquad B'_z=\frac{\mu_0I'}{2\pi a}=\gamma_u B.
\]
电荷速度为 $v'_x=(v_x-u)/(1-uv_x/c^2)$，故
\begin{align*}
 F'_y&=q(E'_y-v'_xB'_z)=-q\gamma_u B(u+v'_x)\\
 &=-q\gamma_u B\frac{v_x(1-u^2/c^2)}{1-uv_x/c^2}
 =\boxed{\frac{F_y}{\gamma_u(1-uv_x/c^2)}}.
\end{align*}
特别取 $u=v_x$，电荷在新系静止，力全部来自电场，而不再有磁力；二者描述的是同一物理作用。

（3）为了得到一般变换，不能只用电中性导线的一种特例猜公式。
再允许导线在 $S$ 中有任意线密度 $\lambda$，则
$E_y=\lambda/(2\pi\varepsilon_0a)$、$B_z=\mu_0I/(2\pi a)$。
把一般四流变换代入场公式，利用 $\mu_0\varepsilon_0=1/c^2$，得到
\[
 E'_y=\gamma_u(E_y-uB_z),\qquad
 B'_z=\gamma_u(B_z-uE_y/c^2).
\]
绕 $x$ 轴转动这个装置，便得到另一组横向分量。
线性叠加后写成矢量形式为
\[
 \boxed{\vect E'_\perp=\gamma_u(\vect E_\perp+\vect u\times\vect B)},\qquad
 \boxed{\vect B'_\perp=\gamma_u(\vect B_\perp-\vect u\times\vect E/c^2)}.
\]
这里 $uB$ 与电场同量纲，$uE/c^2$ 与磁场同量纲；把 $u$ 换为 $-u$ 并交换两系即可得到逆变换。

图~\ref{advA6:quality:A6-source-more:a4} 中，电荷处的实验室磁场出纸面，取 $qv_x>0$ 时力指向导线；换系后同时出现电场和磁场，必须在同一事件相加电力与磁力后再比较三维力。
\end{advsolution}

\begin{advexamplebox}{逐次放出核子的最大末速度}\label{advA6:ex:nuclear-chain-speed}
设想一种只由质量为 $m$ 的核子构成的核物质，其结合不是放能，而是吸收能量；
对于含两个及以上核子的核，平均每核子的结合吸能恒为 $Q>0$。
初始有一个静止的20核子核，连续19次各放出一个核子，最后剩下一个核子。
求这个剩余核子的最大速度与光速之比。不考虑其他粒子或辐射带走能量。

\begin{center}\begin{minipage}{.96\linewidth}\centering
\begin{tikzpicture}[>=Latex,font=\small,scale=.85,every node/.style={fill=none}]

\node[draw,circle,minimum size=1.1cm] (a) at (0,0){$A_N$};\draw[->] (.7,0)--(1.6,0);\node[draw,circle,minimum size=.8cm] (b) at (2.4,0){$A_{N-1}$};\node[draw,circle,minimum size=.5cm] (c) at (4,0){$m$};\draw[DeepBlue,->] (b)--(1.6,-.6);\draw[DeepBlue,->] (c)--(4.8,.6);\node[below] at (2,-1.2){从20核子初态连续放出19个核子};

\end{tikzpicture}
\figcap{题图：每一步都是二体反冲，所画箭头在该步母核静止系中相反。}\label{advA6:quality:A6-source-more:q5}
\end{minipage}\end{center}
\end{advexamplebox}
\global\AdvHideHeadingtrue
\begin{advsolution}
\begin{center}\begin{minipage}{.96\linewidth}\centering
% EMAdvFigureHeading
\raggedright{\bfseries\color{DeepBlue}解答}\par\smallskip\centering

\begin{tikzpicture}[>=Latex,font=\small,scale=.85,every node/.style={fill=none}]

\node[draw,circle,minimum size=1.1cm] (a) at (0,0){$A_N$};\draw[->] (.7,0)--(1.6,0);\node[draw,circle,minimum size=.8cm] (b) at (2.4,0){$A_{N-1}$};\node[draw,circle,minimum size=.5cm] (c) at (4,0){$m$};\draw[DeepBlue,->] (b)--(1.6,-.6);\draw[DeepBlue,->] (c)--(4.8,.6);\node[below] at (2,-1.2){从20核子初态连续放出19个核子};

\end{tikzpicture}
\figcap{解图：每一步都是二体反冲，所画箭头在该步母核静止系中相反。}\label{advA6:quality:A6-source-more:a5}
\end{minipage}\end{center}

先把反常的结合能符号写进质量，而不能套用普通核的质量亏损：
\[
 M_j=j(m+Q/c^2)\quad(j\ge2),\qquad M_1=m.
\]
单个自由核子没有“每核子结合能”，故最后一步 $M_2\to m+m$ 必须单独使用 $M_1=m$。
每次在母核自身静止系作二体衰变，能量、动量守恒给剩余子核的能量
\[
 E_{j-1}^{(j)}=\frac{M_j^2+M_{j-1}^2-m^2}{2M_j}c^2.
\]
用快度 $\xi_j\ge0$ 表示它相对母核的运动，则
\[
 \cosh\xi_j=\frac{E_{j-1}^{(j)}}{M_{j-1}c^2}
 =\frac{M_j^2+M_{j-1}^2-m^2}{2M_jM_{j-1}}.
\]
母核若已在实验室中有快度 $\eta$，以它的实验室运动方向定义前向轴。
在\emph{本次母核的瞬时静止系}测量剩余子核发射方向与这条轴的夹角，记为 $\chi$，
实验室洛伦兹因子为
\[
 \gamma_{\rm new}=\cosh\eta\cosh\xi_j
 +\sinh\eta\sinh\xi_j\cos\chi.
\]
给定母核速度时最大值在 $\chi=0$ 取得，即在每次母核静止系中，剩余核向前、放出的核子向后。
换到最初的实验室系后，后者未必仍然向后，不能混淆这两个方向判断。
此时 $\gamma_{\rm new}=\cosh(\eta+\xi_j)$，且关于旧快度单调增加；
逐步归纳，最大末快度就是各次正快度之和。

为把答案只写成题给参数，令 $s=m/(m+Q/c^2)$，则 $0<s<1$。
对 $j\ge3$，代入质量得到
\[
 \cosh\xi_j=1+\frac{1-s^2}{2j(j-1)},\qquad
 \cosh\xi_2=\frac1s.
\]
最终
\[
 \boxed{\frac{v_{\max}}c=
 \tanh\left[\operatorname{arcosh}\frac1s+
 \sum_{j=3}^{20}\operatorname{arcosh}\left(1+\frac{1-s^2}{2j(j-1)}\right)\right]}.
\]
有限求和就是精确的参数化结果，不需要把19次释放近似成连续火箭。
各次二体公式均保持能量、动量守恒；快度有限时速度严格小于 $c$。
当 $Q\to0$，$s\to1$，所有快度都趋于零，末速度也为零，符合没有可释放能量的极限。

图~\ref{advA6:quality:A6-source-more:a5} 中，每次反冲应在当次母核静止系计算，再用快度相加组合到初始系；最大末速度要求后继核的反冲方向逐步同向，最后二核子到单核子的质量公式须单独处理。
\end{advsolution}

\begin{advexamplebox}{顶夸克对的玻尔模型估计}\label{advA6:ex:toponium-estimate}
1995年费米实验室CDF和D0实验组发现顶夸克。按题给估计参数，
顶夸克质量为 $m_t=175\,\mathrm{GeV}/c^2$，寿命为 $\tau=0.4\times10^{-24}\,\mathrm s$。
顶夸克与反顶夸克之间的势能取
$U=-4\alpha_s k/(3r)$，其中 $\alpha_s=0.12$、$k=199\,\mathrm{MeV\,fm}$。
类比玻尔模型求基态二者间距、基态周期以及更高能级周期与基态周期的关系。
根据寿命判断能否形成长寿命、具有清晰能级的束缚态。

\begin{center}\begin{minipage}{.96\linewidth}\centering
\begin{tikzpicture}[>=Latex,font=\small,scale=.85,every node/.style={fill=none}]

\fill (-1,0) circle (3pt) node[below]{$t$};\fill (1,0) circle (3pt) node[below]{$\bar t$};\draw[DeepBlue,->] (-1,0)--(-.3,0);\draw[DeepBlue,->] (1,0)--(.3,0);\draw[<->] (-1,.7)--(1,.7) node[midway,above]{$r$};\fill (0,0) circle (1pt) node[below]{质心};\node[above] at (0,1.4){$U(r)=-4\alpha_s k/(3r)$};

\end{tikzpicture}
\figcap{题图：等质量夸克对的相对距离与共同质心，库仑型势只作题设估算模型。}\label{advA6:quality:A6-source-more:q6}
\end{minipage}\end{center}
\end{advexamplebox}
\global\AdvHideHeadingtrue
\begin{advsolution}
\begin{center}\begin{minipage}{.96\linewidth}\centering
% EMAdvFigureHeading
\raggedright{\bfseries\color{DeepBlue}解答}\par\smallskip\centering

\begin{tikzpicture}[>=Latex,font=\small,scale=.85,every node/.style={fill=none}]

\fill (-1,0) circle (3pt) node[below]{$t$};\fill (1,0) circle (3pt) node[below]{$\bar t$};\draw[DeepBlue,->] (-1,0)--(-.3,0);\draw[DeepBlue,->] (1,0)--(.3,0);\draw[<->] (-1,.7)--(1,.7) node[midway,above]{$r$};\fill (0,0) circle (1pt) node[below]{质心};\node[above] at (0,1.4){$U(r)=-4\alpha_s k/(3r)$};

\end{tikzpicture}
\figcap{解图：等质量夸克对的相对距离与共同质心，库仑型势只作题设估算模型。}\label{advA6:quality:A6-source-more:a6}
\end{minipage}\end{center}

两个质量相同的粒子都在运动，应使用约化质量 $\mu=m_t/2$，而不是把反夸克固定后仍用 $m_t$。
令 $\kappa=4\alpha_s k/3=31.84\,\mathrm{MeV\,fm}$，相对圆周运动满足
\[
 \frac{\mu v_n^2}{r_n}=\frac\kappa{r_n^2},\qquad
 \mu v_nr_n=n\hbar\quad(n=1,2,\ldots).
\]
消去 $v_n$，得到
\[
 r_n=\frac{n^2\hbar^2}{\mu\kappa},\qquad
 v_n=\frac{\kappa}{n\hbar},\qquad
 T_n=\frac{2\pi n^3\hbar^3}{\mu\kappa^2}=n^3T_1.
\]
这里 $r_n$ 是两粒子之间的距离，每个粒子绕质心的半径为 $r_n/2$、速率为 $v_n/2$，
其转动周期仍为同一个 $T_n$。若把“第 $n$ 激发态”按基态之后的第 $n$ 个状态编号，
主量子数应为 $n+1$，相应周期比为 $(n+1)^3$。
用 $\hbar c=197.33\,\mathrm{MeV\,fm}$，数值为
\[
 \boxed{r_1=\frac{(197.33)^2}{87500\times31.84}\,\mathrm{fm}
 \simeq1.40\times10^{-2}\,\mathrm{fm}},\qquad
 \boxed{T_1\simeq1.82\times10^{-24}\,\mathrm s}.
\]
相对速率约为 $0.161c$，每个粒子的速率约 $0.0807c$，非相对论估计在这个粗略模型内合理。
但单粒子寿命只有基态周期约 $0.22$ 倍；若两者独立衰变，完整夸克对的平均存活时间更短，为 $\tau/2$。
因此不能形成这种玻尔图景下长寿命、能级清晰的原子式束缚系统。
这一寿命判据不等于否认阈值附近的短程库仑关联或宽共振效应。

图~\ref{advA6:quality:A6-source-more:a6} 中，$r$ 是两体相对距离，单体绕质心半径为 $r/2$，径向方程使用约化质量 $m_t/2$；能否形成轨道结构还要比较公转时间与题给寿命。
\end{advsolution}

\begin{advexamplebox}{一维有限粒子数费米气体}\label{advA6:ex:finite-fermi-gas}
\begin{enumerate}
\item 写出一维经典理想气体在绝热过程中边界力 $F$ 与容器长度 $L$ 的关系及绝热指数。
\item 对长度 $L$ 的一维刚性盒内、质量为 $m$ 的非相对论粒子，写出驻波量子化条件。
\item 零温下 $N$ 个自旋 $1/2$ 费米子从低能级依次填充，每个空间态最多容纳两个反向自旋粒子。
求边界力与 $L$ 的精确关系，不得假定 $N\gg1$。
\item 缓慢压缩到 $L/2$，使每个粒子保持原占据态的量子数，求过程的指数及外界做功，仍不得对 $N$ 近似。
\end{enumerate}

\begin{center}\begin{minipage}{.96\linewidth}\centering
\begin{tikzpicture}[>=Latex,font=\small,scale=.85,every node/.style={fill=none}]

\draw[very thick] (0,1.4)--(0,0)--(4,0)--(4,1.4);\draw[<->] (0,-.5)--(4,-.5) node[midway,below]{$L$};\node[above] at (2,1.6){$N$ 个自旋1/2费米子};\draw[->] (4,.6)--(4.9,.6) node[right]{$F$};

\end{tikzpicture}
\figcap{题图：一维刚性盒的长度、粒子数与右边界力，粒子数在求长度导数时保持不变。}\label{advA6:quality:A6-source-more:q7}
\end{minipage}\end{center}
\end{advexamplebox}
\global\AdvHideHeadingtrue
\begin{advsolution}
\begin{center}\begin{minipage}{.96\linewidth}\centering
% EMAdvFigureHeading
\raggedright{\bfseries\color{DeepBlue}解答}\par\smallskip\centering

\begin{tikzpicture}[>=Latex,font=\small,scale=.85,every node/.style={fill=none}]

\foreach \y/\n in {0/1,.6/2,1.6/3}{\draw[thick] (0,\y)--(3,\y);\node[right] at (3,\y){$n=\n$};}\foreach \y in {0,.6}{\draw[DeepBlue,->] (1,\y-.18)--(1,\y+.18);\draw[DeepBlue,->] (2,\y+.18)--(2,\y-.18);}\draw[DeepBlue,->] (1,1.42)--(1,1.78);\node[above] at (1.5,2.2){$N=5$ 的奇数填充示意};

\end{tikzpicture}
\figcap{解图：两个反向自旋占同一空间能级，奇数粒子的最后一级单占。}\label{advA6:quality:A6-source-more:a7}
\end{minipage}\end{center}

（1）一维气体只有一个平动自由度，$U=Nk_BT/2$、$FL=Nk_BT=2U$。
绝热准静态过程满足 $\dd U=-F\dd L=-2U\dd L/L$，积分得 $UL^2$ 恒定，故
\[
 \boxed{FL^3=\text{常量},\qquad\gamma_{\rm ad}=3}.
\]
这是用 $F,L$ 代替三维气体 $p,V$ 的一维版本。

（2）硬壁要求波函数在两端为零，盒内恰有整数个半波长：
\[
 k_jL=j\pi,\quad p_j=\frac{j\pi\hbar}{L},\quad
 \varepsilon_j=\frac{j^2\pi^2\hbar^2}{2mL^2},\quad j=1,2,\ldots.
\]
这里 $p_j$ 是动量\emph{大小}，正弦驻波并非确定带符号动量的本征态；
不能把 $+p_j,-p_j$ 再当成两个独立空间能级而重复计数。

（3）定义 $C=\pi^2\hbar^2/(2m)$。偶数 $N=2M$ 时，每个 $j=1,\ldots,M$ 占两个粒子，
\[
 U_N=\frac C{L^2}\,2\sum_{j=1}^M j^2
 =\frac C{L^2}\frac{M(M+1)(2M+1)}3.
\]
奇数 $N=2M+1$ 时，除上述填充外，第 $M+1$ 级占一个粒子，故
\[
 U_N=\frac C{L^2}\left[\frac{M(M+1)(2M+1)}3+(M+1)^2\right].
\]
这两式对有限 $N$ 精确，包含 $N=1$、$N=2$，没有用积分替代求和。
将方括号或偶数式系数统一记为 $A_N$，在固定占据数下
\[
 \boxed{F=-\frac{\partial U_N}{\partial L}=\frac{2CA_N}{L^3}=\frac{2U_N}{L}}.
\]
即使温度为零，量子驻波与泡利填充也使能量和边界力非零；不能用经典 $T=0$ 代入 $FL=Nk_BT$ 去否定它。

（4）量子数不变意味着 $A_N$ 不变，所以仍有 $FL^3$ 恒定、指数为3。
长度减半使每个能级增为4倍，从而
\[
 \boxed{W_{\rm on}=U_N(L/2)-U_N(L)=3U_N(L)=\frac{3CA_N}{L^2}}.
\]
也可积分外力功 $-\int_L^{L/2}F\dd L$ 得到同一结果。
压缩功为正，且量纲 $\hbar^2/(mL^2)$ 为能量。经典与量子指数相同源于能量的 $L^{-2}$ 标度，
并不意味着它们具有相同的微观分布或零温能量。

图~\ref{advA6:quality:A6-source-more:a7} 中，每个空间能级最多容纳两个相反自旋，零温填充从低能级向上进行；奇数粒子的最高空间态只有一个粒子，有限粒子数结果应区分奇偶。
\end{advsolution}

\begin{advexamplebox}{均匀磁场中的半经典量子化}\label{advA6:ex:landau-action}
带电粒子的正则动量与机械动量不同。用半经典闭轨道量子化条件
$\oint\vect P\cdot\dd\vect r=(n+1/2)h$，求电子在均匀磁场中的回旋运动量子化条件。
说明哪些量离散，角动量能否取 $n\hbar$ 的任意整数倍。
忽略自旋与相对论修正。

\begin{center}\begin{minipage}{.96\linewidth}\centering
\begin{tikzpicture}[>=Latex,font=\small,scale=.85,every node/.style={fill=none}]

\draw[thick] (0,0) circle (1.15);\draw[->] (0,0)--(0,1.15) node[midway,right]{$R$};\fill (0,1.15) circle (1.6pt);\draw[DeepBlue,->] (0,1.15)--(-1.1,1.15) node[left]{$\boldsymbol\pi=m\boldsymbol v$};\node[right] at (1.8,.4){电子 $q=-e$};\node[right] at (1.8,-.1){$\boldsymbol B=B\boldsymbol e_z\;\odot$};

\end{tikzpicture}
\figcap{题图：电子在出纸面磁场中的逆时针局部回旋，机械动量沿轨道切线。}\label{advA6:quality:A6-source-more:q8}
\end{minipage}\end{center}
\end{advexamplebox}
\global\AdvHideHeadingtrue
\begin{advsolution}
\begin{center}\begin{minipage}{.96\linewidth}\centering
% EMAdvFigureHeading
\raggedright{\bfseries\color{DeepBlue}解答}\par\smallskip\centering

\begin{tikzpicture}[>=Latex,font=\small,scale=.85,every node/.style={fill=none}]

\fill[DeepBlue!7] (0,0) circle (1.1);\draw[thick] (0,0) circle (1.1);\node at (0,0){$\Phi>0\;\odot$};
\draw[DeepBlue,thick,->] (0,1.1)--(-1.3,1.1) node[left]{$\boldsymbol\pi$};\draw[->] (0,1.5)--(-.65,1.5) node[left]{$\boldsymbol P$};\draw[->] (0,.75)--(.65,.75) node[right]{$q\boldsymbol A$};
\node[right,align=left] at (1.8,0){对称规范，原点取导引中心\\$\boldsymbol P=\boldsymbol\pi+q\boldsymbol A$\\上方三箭头仅平移分层以便比较};

\end{tikzpicture}
\figcap{解图：在对称规范下，电子的磁通项使闭轨正则作用量减小；各动量箭头按同一比例画出。}\label{advA6:quality:A6-source-more:a8}
\end{minipage}\end{center}

明确符号：机械动量 $\vect\pi=m\vect v$，正则动量为
$\vect P=\vect\pi+q\vect A$，即 $\vect\pi=\vect P-q\vect A$。
量子化的作用量积分必须用正则动量，不能把两者交换。
取 $\vect B=B\vect e_z$，电子 $q=-e$，其回旋方向从 $+z$ 看为逆时针。
半径 $R$ 的圆满足 $\pi=eBR$。沿运动方向积分，机械动量一项为
$2\pi R\pi=2\pi eBR^2$，磁通项为
\[
 q\oint\vect A\cdot\dd\vect r=q\Phi=-eB\pi R^2.
\]
两项合并，正则作用量只有机械项的一半：
\[
 J=\oint\vect P\cdot\dd\vect r=\pi eBR^2
 =\frac{2\pi E_\perp}{\omega_c},\qquad \omega_c=\frac{eB}m.
\]
因此
\[
 \boxed{E_{\perp,n}=\left(n+\frac12\right)\hbar\omega_c},\qquad
 \boxed{R_n^2=\frac{(2n+1)\hbar}{eB}},\qquad n=0,1,\ldots.
\]
这是不计自旋的朗道回旋能量。沿磁场的自由动能 $p_z^2/(2m)$ 仍连续，
总能量需再加这一项；横向轨道的导引中心位置也不是由 $n$ 唯一确定的。

对题中的半经典圆轨道，绕\emph{导引中心}的机械角动量大小为
\[
 |L_{\rm mech}|=m\omega_cR_n^2=(2n+1)\hbar,
\]
因此在这个圆轨道图景中是奇数倍，不是任意整数倍，也不能漏掉 $1/2$ 零点项。
这不等于说精确量子态绕任意固定原点的正则角动量只能取奇数倍。
在轴对称规范中，固定原点的正则 $L_z$ 可另有整数角量子数；
它与机械角动量不同，朗道能量又有导引中心简并，不能只给能级 $n$ 就唯一确定后一角量子数。
半经典的半径和机械角动量应理解为回旋量的对应关系，不是电子在精确量子态中具有可同时确定的位置轨迹。

图~\ref{advA6:quality:A6-source-more:a8} 中，轨道按电子运动方向逆时针积分时磁通为正，但电荷 $q=-e$ 使磁通项为负；正则作用量是机械项加 $q\Phi$，导引中心简并与沿场自由动能仍需保留。
\end{advsolution}

\begin{advexamplebox}{作用量相位、费马原理与圆轨道量子化}\label{advA6:ex:action-optics}
几何光学的费马原理为 $\delta\int_A^B n\,\dd l=0$，单色光的相位为
$\phi=(\omega/c)\int n\,\dd l-\omega t$。
非相对论力学的哈密顿原理为 $\delta S=0$，
$S=\int(\vect p\cdot\vect v-E)\dd t$，$E=mv^2/2+V(\vect r)$。
\begin{enumerate}
\item 若物质波相位写成 $\phi=kS$，利用德布罗意关系确定 $k$。
\item 把光线在折射率分布中的轨迹形式地类比为质量 $m$ 的经典粒子轨迹，求等效势能与折射率的关系。
\item 对氢原子库仑势 $V=-e^2/(4\pi\varepsilon_0r)$，用光学类比推导电子圆轨道的向心力条件。
\item 用环绕一周的相位相容性推导角动量量子化条件。
\end{enumerate}

\begin{center}\begin{minipage}{.96\linewidth}\centering
\begin{tikzpicture}[>=Latex,font=\small,scale=.85,every node/.style={fill=none}]

\draw[thick] (0,0) circle (1.15);\fill (0,0) circle (1.5pt) node[below]{$O$};\draw[DeepBlue,->] (1.15,0)--(1.15,.8) node[right]{$p$};\draw[->] (0,0)--(1.15,0) node[midway,below]{$r$};\node[right] at (2,0){$n=n(r)$ 或 $V=V(r)$};\draw[dashed] (0,0) circle (1.4);

\end{tikzpicture}
\figcap{题图：圆对称介质/势场中的候选圆轨道，虚圆表示半径的邻近变分。}\label{advA6:quality:A6-source-more:q9}
\end{minipage}\end{center}
\end{advexamplebox}
\global\AdvHideHeadingtrue
\begin{advsolution}
\begin{center}\begin{minipage}{.96\linewidth}\centering
% EMAdvFigureHeading
\raggedright{\bfseries\color{DeepBlue}解答}\par\smallskip\centering

\begin{tikzpicture}[>=Latex,font=\small,scale=.85,every node/.style={fill=none}]

\draw[thick] (0,0) circle (1.15);\fill (0,0) circle (1.5pt) node[below]{$O$};\draw[DeepBlue,->] (1.15,0)--(1.15,.8) node[right]{$p$};\draw[->] (0,0)--(1.15,0) node[midway,below]{$r$};\node[right] at (2,0){$n=n(r)$ 或 $V=V(r)$};\draw[dashed] (0,0) circle (1.4);

\end{tikzpicture}
\figcap{解图：圆对称介质/势场中的候选圆轨道，虚圆表示半径的邻近变分。}\label{advA6:quality:A6-source-more:a9}
\end{minipage}\end{center}

（1）自由粒子的作用量为 $S=\vect p\cdot\vect r-Et$。
德布罗意关系要求相位为 $(\vect p\cdot\vect r-Et)/\hbar$，故
\[
 \boxed{k=1/\hbar},\qquad \phi=S/\hbar.
\]
这里的 $k$ 是题设比例常量，量纲为作用量的倒数，不是通常有长度倒数量纲的波数。

（2）要比较空间轨迹，应在\emph{固定能量}下使用简化作用量
$W=\int\vect p\cdot\dd\vect r=\int p\,\dd l$ 的驻值原理。
这与固定端点时间的哈密顿原理有联系，但不能在任意变分过程中强行把时间也保持不变、能量也保持不变。
若取 $p(\vect r)=p_*n(\vect r)$，两种空间变分只差常数 $p_*$，产生相同轨迹。
由 $p^2/(2m)+V=E$，得到
\[
 \boxed{V(\vect r)=E-\frac{p_*^2}{2m}n^2(\vect r)}.
\]
$p_*$ 是任取的正动量尺度，$E$ 可随势能零点一同平移；
因此不是从折射率唯一确定一个绝对势能数值。
这只是轨迹等效模型，\emph{不是}给真实光子赋予静质量 $m$，也不要求两种模型的运行时间相同。

（3）圆对称折射率中，半径 $r$ 的圆光路光程为 $2\pi r n(r)$。
对半径作小变化，圆轨道的驻值条件是
\[
 \frac{\dd}{\dd r}[r n(r)]=0
 \quad\Rightarrow\quad \frac{n'}n=-\frac1r.
\]
由 $n\propto p=\sqrt{2m(E-V)}$，在固定 $E$ 下有
$n'/n=-V'/[2(E-V)]$。
联立得 $V'=2(E-V)/r=mv^2/r$，即径向力大小
\[
 \boxed{\frac{mv^2}{r}=\frac{\dd V}{\dd r}
 =\frac{e^2}{4\pi\varepsilon_0r^2}}.
\]
因为 $V'>0$，力 $-V'\vect e_r$ 指向圆心，方向正确。

（4）同一时刻绕圆一周，时间相位项相消，波函数应回到相同相位，要求
\[
 \Delta\phi=\frac1\hbar\oint\vect p\cdot\dd\vect r
 =\frac{2\pi r p}{\hbar}=2\pi n,
 \quad\Rightarrow\quad\boxed{L=rp=n\hbar}.
\]
对非零圆轨道取 $n=1,2,\ldots$。与上一问联立给出玻尔半径的 $n^2$ 标度，
但这是早期圆轨道半经典模型，不应把 $L=n\hbar$ 直接当作氢原子精确三维量子力学中
$|\vect L|=\sqrt{\ell(\ell+1)}\hbar$ 的替代式。
费马驻值光线也不表示单光子一定沿一条经典路径传播；它是波动相位在短波极限中相干叠加所给出的几何描述。

图~\ref{advA6:quality:A6-source-more:a9} 中，半径变分比较的是固定能量的空间简化作用量，故圆轨道条件来自 $\mathrm d[rn(r)]/\mathrm dr=0$；绕同一时刻的一周相位闭合再给出早期圆轨道量子化，两种条件不能互相替代。
\end{advsolution}


\clearpage\endgroup


\bookmarksetup{startatroot}
\backmatter
\chapter*{记号说明}
\addcontentsline{toc}{chapter}{记号说明}
除题目另有说明外，本书采用国际单位制。数值计算所用物理常数和近似值以题目或所在节的说明为准。

矢量通常以粗体表示，部分图示与推导使用箭头记号；两者均表示具有大小和方向的量。微分记号写作 $\dd$，虚数单位写作 $\ii$。
同一字母在不同专题中可能表示不同物理量，应结合本节定义及单位辨认；例如 $I$ 可表示电流、转动惯量或光强，$L$ 可表示长度、角动量大小或电感。

除专门说明角度单位的题目外，微积分与小角展开中的角度以弧度计。对数的底、正方向及参考系均以局部定义为准。

各章例题、图和公式均按“章号.顺序号”编号。
\end{document}
